% Options for packages loaded elsewhere
\PassOptionsToPackage{unicode}{hyperref}
\PassOptionsToPackage{hyphens}{url}
\PassOptionsToPackage{space}{xeCJK}
%
\documentclass[
  fontset=none]{ctexart}
\usepackage{amsmath,amssymb}
\usepackage{lmodern}
\usepackage{iftex}
\ifPDFTeX
  \usepackage[T1]{fontenc}
  \usepackage[utf8]{inputenc}
  \usepackage{textcomp} % provide euro and other symbols
\else % if luatex or xetex
  \usepackage{unicode-math}
  \defaultfontfeatures{Scale=MatchLowercase}
  \defaultfontfeatures[\rmfamily]{Ligatures=TeX,Scale=1}
  \setmainfont[]{Hiragino Sans GB}
  \setsansfont[]{Hiragino Sans GB}
  \setmonofont[]{Menlo}
  \ifXeTeX
    \usepackage{xeCJK}
    \setCJKmainfont[]{Hiragino Sans GB}
  \fi
  \ifLuaTeX
    \usepackage[]{luatexja-fontspec}
    \setmainjfont[]{Hiragino Sans GB}
  \fi
\fi
% Use upquote if available, for straight quotes in verbatim environments
\IfFileExists{upquote.sty}{\usepackage{upquote}}{}
\IfFileExists{microtype.sty}{% use microtype if available
  \usepackage[]{microtype}
  \UseMicrotypeSet[protrusion]{basicmath} % disable protrusion for tt fonts
}{}
\makeatletter
\@ifundefined{KOMAClassName}{% if non-KOMA class
  \IfFileExists{parskip.sty}{%
    \usepackage{parskip}
  }{% else
    \setlength{\parindent}{0pt}
    \setlength{\parskip}{6pt plus 2pt minus 1pt}}
}{% if KOMA class
  \KOMAoptions{parskip=half}}
\makeatother
\usepackage{xcolor}
\IfFileExists{xurl.sty}{\usepackage{xurl}}{} % add URL line breaks if available
\IfFileExists{bookmark.sty}{\usepackage{bookmark}}{\usepackage{hyperref}}
\hypersetup{
  hidelinks,
  pdfcreator={LaTeX via pandoc}}
\urlstyle{same} % disable monospaced font for URLs
\usepackage{listings}
\newcommand{\passthrough}[1]{#1}
\lstset{defaultdialect=[5.3]Lua}
\lstset{defaultdialect=[x86masm]Assembler}
\usepackage{longtable,booktabs,array}
\usepackage{calc} % for calculating minipage widths
% Correct order of tables after \paragraph or \subparagraph
\usepackage{etoolbox}
\makeatletter
\patchcmd\longtable{\par}{\if@noskipsec\mbox{}\fi\par}{}{}
\makeatother
% Allow footnotes in longtable head/foot
\IfFileExists{footnotehyper.sty}{\usepackage{footnotehyper}}{\usepackage{footnote}}
\makesavenoteenv{longtable}
\setlength{\emergencystretch}{3em} % prevent overfull lines
\providecommand{\tightlist}{%
  \setlength{\itemsep}{0pt}\setlength{\parskip}{0pt}}
\setcounter{secnumdepth}{-\maxdimen} % remove section numbering
\usepackage{bookmark}
\usepackage{etoolbox}
\usepackage{fancyhdr}
\usepackage{xcolor}
\usepackage{listings}

\definecolor{UnitreeBlue}{HTML}{145DA0}
\definecolor{UnitreeInk}{HTML}{17212B}
\definecolor{UnitreeCode}{HTML}{F3F5F7}
\definecolor{UnitreeRule}{HTML}{D7DEE5}

\hypersetup{
  colorlinks=true,
  linkcolor=UnitreeBlue,
  urlcolor=UnitreeBlue,
  citecolor=UnitreeBlue,
  pdfborder={0 0 0},
  bookmarksopen=true,
  bookmarksnumbered=false
}
\bookmarksetup{open,depth=3}

\pagestyle{fancy}
\fancyhf{}
\fancyhead[L]{\small\color{UnitreeInk}Unitree SDK2 Python Bindings}
\fancyhead[R]{\small\color{UnitreeInk}API Documentation}
\fancyfoot[C]{\small\color{UnitreeInk}\thepage}
\renewcommand{\headrulewidth}{0.4pt}
\renewcommand{\footrulewidth}{0pt}
\setlength{\headheight}{14pt}

\lstset{
  basicstyle=\ttfamily\small,
  backgroundcolor=\color{UnitreeCode},
  frame=single,
  rulecolor=\color{UnitreeRule},
  framesep=6pt,
  breaklines=true,
  breakatwhitespace=false,
  columns=fullflexible,
  keepspaces=true,
  showstringspaces=false,
  upquote=true
}

\AtBeginEnvironment{longtable}{%
  \small
  \setlength{\tabcolsep}{3pt}%
  \renewcommand{\arraystretch}{1.08}%
}

\setlength{\parindent}{0pt}
\setlength{\parskip}{0.45em}
\setlength{\emergencystretch}{3em}
\setcounter{secnumdepth}{0}
\sloppy

\AtBeginDocument{%
  \color{UnitreeInk}%
}
\ifLuaTeX
  \usepackage{selnolig}  % disable illegal ligatures
\fi

\author{}
\date{}

\begin{document}

\hypertarget{unitree-sdk2-python-ux5b8cux6574-api-ux53c2ux8003}{%
\section{Unitree SDK2 Python 完整 API
参考}\label{unitree-sdk2-python-ux5b8cux6574-api-ux53c2ux8003}}

本参考采用常见科学计算库的 API Reference
组织方式：先按模块分类索引，再为每个类、函数和属性提供签名、参数、返回值、可用性、C++
对应项和用法。学习路径和完整教程请先阅读
\href{BEGINNER_GUIDE_ZH.pdf}{从零开始指南}。

\begin{quote}
{[}!WARNING{]} 本文同时包含 \passthrough{\lstinline!AVAILABLE!} 和
\passthrough{\lstinline!SIGNATURE\_ONLY!}。后者只是设计期类型签名。所有运动方法目前均为
\passthrough{\lstinline!SIGNATURE\_ONLY!}，不得按当前可执行接口使用。
\end{quote}

\begin{quote}
{[}!NOTE{]} 本文件由
\passthrough{\lstinline!generator/generate\_api\_docs.py!} 从
\passthrough{\lstinline!.pyi!}、\passthrough{\lstinline!api\_manifest.json!}
和 Clang AST 清单生成。不要手工维护 API
条目；修改签名后应重新生成并运行测试。
\end{quote}

\hypertarget{ux5982ux4f55ux9605ux8bfbux672cux53c2ux8003}{%
\subsection{如何阅读本参考}\label{ux5982ux4f55ux9605ux8bfbux672cux53c2ux8003}}

每个 API 条目包含以下部分：

\begin{itemize}
\tightlist
\item
  \textbf{签名}：Python 类型检查器看到的准确调用形态；
\item
  \textbf{可用性与安全性}：区分当前实现和仅签名预览；
\item
  \textbf{参数}：Python 类型、默认值和能够从 SDK 定义可靠确认的含义；
\item
  \textbf{返回值}：包括 C++ 输出引用转换后的 Python 元组顺序；
\item
  \textbf{对应 C++}：原类、原签名、绑定策略和头文件位置；
\item
  \textbf{用法}：\passthrough{\lstinline!AVAILABLE!}
  给出调用形态，\passthrough{\lstinline!SIGNATURE\_ONLY!} 只给出
  \passthrough{\lstinline!TYPE\_CHECKING!} 计划代码。
\end{itemize}

各条目中的代码是局部调用片段，不是完整脚本。类章节会先给出导入语句；\passthrough{\lstinline!obj!}、\passthrough{\lstinline!helper!}、\passthrough{\lstinline!client!}、\passthrough{\lstinline!publisher!}、\passthrough{\lstinline!subscriber!}
和参数变量表示调用者已经按业务规则创建、初始化并校验好的对象。

参数名称无法表达单位、坐标系、范围或枚举值时，本文会明确要求查目标型号协议。这比根据名字猜测更可靠。

\hypertarget{ux8986ux76d6ux7edfux8ba1}{%
\subsection{覆盖统计}\label{ux8986ux76d6ux7edfux8ba1}}

\begin{longtable}[]{@{}
  >{\raggedright\arraybackslash}p{(\columnwidth - 2\tabcolsep) * \real{0.35}}
  >{\raggedleft\arraybackslash}p{(\columnwidth - 2\tabcolsep) * \real{0.65}}@{}}
\toprule
项目 & 数量 \\
\midrule
\endhead
有内容的 Python 模块 & 15 \\
类和枚举 & 189 \\
普通/重载函数签名 & 844 \\
Python 属性 & 346 \\
公开数据属性 & 173 \\
Manifest 条目 & 1190 \\
\passthrough{\lstinline!AVAILABLE!} & 650 \\
\passthrough{\lstinline!SIGNATURE\_ONLY!} & 540 \\
\bottomrule
\end{longtable}

一个 Python 属性在 manifest 中占一个条目，但下文会同时写出 getter 和
setter 签名。重载方法按不同 C++ 签名分别展开。

\hypertarget{ux5206ux7c7bux7d22ux5f15}{%
\subsection{分类索引}\label{ux5206ux7c7bux7d22ux5f15}}

\begin{longtable}[]{@{}
  >{\raggedright\arraybackslash}p{(\columnwidth - 6\tabcolsep) * \real{0.18}}
  >{\raggedright\arraybackslash}p{(\columnwidth - 6\tabcolsep) * \real{0.24}}
  >{\raggedleft\arraybackslash}p{(\columnwidth - 6\tabcolsep) * \real{0.18}}
  >{\raggedleft\arraybackslash}p{(\columnwidth - 6\tabcolsep) * \real{0.40}}@{}}
\toprule
分类 & Python 模块 & 函数 & 类 \\
\midrule
\endhead
系统辅助 API &
\protect\hyperlink{unitree-sdk2-cpp}{\passthrough{\lstinline!unitree\_sdk2\_cpp!}}
& 0 & 1 \\
Typed DDS Channel API &
\protect\hyperlink{unitree-sdk2-cpp-channel}{\passthrough{\lstinline!unitree\_sdk2\_cpp.channel!}}
& 4 & 2 \\
Go2 IDL 消息 &
\protect\hyperlink{unitree-sdk2-cpp-idl-go2}{\passthrough{\lstinline!unitree\_sdk2\_cpp.idl.go2!}}
& 0 & 26 \\
HG IDL 消息 &
\protect\hyperlink{unitree-sdk2-cpp-idl-hg}{\passthrough{\lstinline!unitree\_sdk2\_cpp.idl.hg!}}
& 0 & 13 \\
HG Double-IMU IDL 消息 &
\protect\hyperlink{unitree-sdk2-cpp-idl-hg-doubleimu}{\passthrough{\lstinline!unitree\_sdk2\_cpp.idl.hg\_doubleimu!}}
& 0 & 1 \\
ROS2 兼容 IDL 消息 &
\protect\hyperlink{unitree-sdk2-cpp-idl-ros2}{\passthrough{\lstinline!unitree\_sdk2\_cpp.idl.ros2!}}
& 0 & 24 \\
Robot 公共基础 API &
\protect\hyperlink{unitree-sdk2-cpp-robot}{\passthrough{\lstinline!unitree\_sdk2\_cpp.robot!}}
& 0 & 18 \\
A2 Robot API &
\protect\hyperlink{unitree-sdk2-cpp-robot-a2}{\passthrough{\lstinline!unitree\_sdk2\_cpp.robot.a2!}}
& 0 & 8 \\
AS2 Robot API &
\protect\hyperlink{unitree-sdk2-cpp-robot-as2}{\passthrough{\lstinline!unitree\_sdk2\_cpp.robot.as2!}}
& 0 & 2 \\
B2 Robot API &
\protect\hyperlink{unitree-sdk2-cpp-robot-b2}{\passthrough{\lstinline!unitree\_sdk2\_cpp.robot.b2!}}
& 0 & 25 \\
G1 Robot API &
\protect\hyperlink{unitree-sdk2-cpp-robot-g1}{\passthrough{\lstinline!unitree\_sdk2\_cpp.robot.g1!}}
& 0 & 12 \\
Go2 Robot API &
\protect\hyperlink{unitree-sdk2-cpp-robot-go2}{\passthrough{\lstinline!unitree\_sdk2\_cpp.robot.go2!}}
& 0 & 37 \\
H1 Robot API &
\protect\hyperlink{unitree-sdk2-cpp-robot-h1}{\passthrough{\lstinline!unitree\_sdk2\_cpp.robot.h1!}}
& 0 & 4 \\
H2 Robot API &
\protect\hyperlink{unitree-sdk2-cpp-robot-h2}{\passthrough{\lstinline!unitree\_sdk2\_cpp.robot.h2!}}
& 0 & 8 \\
R1 Robot API &
\protect\hyperlink{unitree-sdk2-cpp-robot-r1}{\passthrough{\lstinline!unitree\_sdk2\_cpp.robot.r1!}}
& 0 & 8 \\
\bottomrule
\end{longtable}

\begin{center}\rule{0.5\linewidth}{0.5pt}\end{center}

\hypertarget{unitree-sdk2-cpp}{%
\subsection{系统辅助 API}\label{unitree-sdk2-cpp}}

模块：\passthrough{\lstinline!unitree\_sdk2\_cpp!}

\hypertarget{ux7c7bux7d22ux5f15}{%
\subsubsection{类索引}\label{ux7c7bux7d22ux5f15}}

\begin{longtable}[]{@{}
  >{\raggedright\arraybackslash}p{(\columnwidth - 4\tabcolsep) * \real{0.18}}
  >{\raggedleft\arraybackslash}p{(\columnwidth - 4\tabcolsep) * \real{0.20}}
  >{\raggedleft\arraybackslash}p{(\columnwidth - 4\tabcolsep) * \real{0.62}}@{}}
\toprule
类 & 公开函数签名 & 属性 \\
\midrule
\endhead
\protect\hyperlink{unitree-sdk2-cpp-oshelper}{\passthrough{\lstinline!OsHelper!}}
& 7 & 0 \\
\bottomrule
\end{longtable}

\hypertarget{unitree-sdk2-cpp-oshelper}{%
\subsubsection{\texorpdfstring{\texttt{unitree\_sdk2\_cpp.OsHelper}}{unitree\_sdk2\_cpp.OsHelper}}\label{unitree-sdk2-cpp-oshelper}}

读取当前主机操作系统信息的单例辅助类。

\textbf{导入}

\begin{lstlisting}[language=Python]
from unitree_sdk2_cpp import OsHelper
\end{lstlisting}

\textbf{方法索引}

\begin{longtable}[]{@{}
  >{\raggedright\arraybackslash}p{(\columnwidth - 6\tabcolsep) * \real{0.18}}
  >{\raggedright\arraybackslash}p{(\columnwidth - 6\tabcolsep) * \real{0.24}}
  >{\raggedright\arraybackslash}p{(\columnwidth - 6\tabcolsep) * \real{0.18}}
  >{\raggedright\arraybackslash}p{(\columnwidth - 6\tabcolsep) * \real{0.40}}@{}}
\toprule
方法 & Python 签名 & 状态 & 安全分类 \\
\midrule
\endhead
\protect\hyperlink{unitree-sdk2-cpp-oshelper-instance-1}{\passthrough{\lstinline!instance!}}
& \passthrough{\lstinline!def instance() -> OsHelper!} &
\passthrough{\lstinline!AVAILABLE!} &
\passthrough{\lstinline!READ\_ONLY!} \\
\protect\hyperlink{unitree-sdk2-cpp-oshelper-get-uid-1}{\passthrough{\lstinline!get\_uid!}}
& \passthrough{\lstinline!def get\_uid(self) -> int!} &
\passthrough{\lstinline!AVAILABLE!} &
\passthrough{\lstinline!READ\_ONLY!} \\
\protect\hyperlink{unitree-sdk2-cpp-oshelper-get-gid-1}{\passthrough{\lstinline!get\_gid!}}
& \passthrough{\lstinline!def get\_gid(self) -> int!} &
\passthrough{\lstinline!AVAILABLE!} &
\passthrough{\lstinline!READ\_ONLY!} \\
\protect\hyperlink{unitree-sdk2-cpp-oshelper-get-user-1}{\passthrough{\lstinline!get\_user!}}
& \passthrough{\lstinline!def get\_user(self) -> str!} &
\passthrough{\lstinline!AVAILABLE!} &
\passthrough{\lstinline!READ\_ONLY!} \\
\protect\hyperlink{unitree-sdk2-cpp-oshelper-get-processor-number-1}{\passthrough{\lstinline!get\_processor\_number!}}
& \passthrough{\lstinline!def get\_processor\_number(self) -> int!} &
\passthrough{\lstinline!AVAILABLE!} &
\passthrough{\lstinline!READ\_ONLY!} \\
\protect\hyperlink{unitree-sdk2-cpp-oshelper-get-page-size-1}{\passthrough{\lstinline!get\_page\_size!}}
& \passthrough{\lstinline!def get\_page\_size(self) -> int!} &
\passthrough{\lstinline!AVAILABLE!} &
\passthrough{\lstinline!READ\_ONLY!} \\
\protect\hyperlink{unitree-sdk2-cpp-oshelper-get-hostname-1}{\passthrough{\lstinline!get\_hostname!}}
& \passthrough{\lstinline!def get\_hostname(self) -> str!} &
\passthrough{\lstinline!AVAILABLE!} &
\passthrough{\lstinline!READ\_ONLY!} \\
\bottomrule
\end{longtable}

\hypertarget{unitree-sdk2-cpp-oshelper-instance-1}{%
\paragraph{\texorpdfstring{\texttt{unitree\_sdk2\_cpp.OsHelper.instance}}{unitree\_sdk2\_cpp.OsHelper.instance}}\label{unitree-sdk2-cpp-oshelper-instance-1}}

返回进程内的 \passthrough{\lstinline!OsHelper!} 单例。

\textbf{签名}

\begin{lstlisting}[language=Python]
@staticmethod
def instance() -> OsHelper
\end{lstlisting}

\textbf{可用性与安全性}

\textbf{\passthrough{\lstinline!AVAILABLE!} /
\passthrough{\lstinline!READ\_ONLY!}}：当前绑定源码已实现本机信息读取。该方法不初始化
DDS，也不访问机器人。

\textbf{参数}

无显式参数。实例方法中的 \passthrough{\lstinline!self!} 由 Python
自动传入。

\textbf{返回值}

\begin{longtable}[]{@{}
  >{\raggedright\arraybackslash}p{(\columnwidth - 4\tabcolsep) * \real{0.18}}
  >{\raggedright\arraybackslash}p{(\columnwidth - 4\tabcolsep) * \real{0.20}}
  >{\raggedright\arraybackslash}p{(\columnwidth - 4\tabcolsep) * \real{0.62}}@{}}
\toprule
位置 & 类型 & 含义 \\
\midrule
\endhead
返回值 & \passthrough{\lstinline!OsHelper!} & 进程内的
\passthrough{\lstinline!OsHelper!} 单例。 \\
\bottomrule
\end{longtable}

\textbf{对应 C++}

\begin{itemize}
\tightlist
\item
  类：\passthrough{\lstinline!unitree\_sdk2\_cpp.OsHelper!}
\item
  签名：\passthrough{\lstinline!OsHelper::Instance()!}
\item
  绑定策略：\passthrough{\lstinline!未单独标注!}
\item
  声明位置：由 pybind11 手工绑定或生成注册表提供；manifest
  未记录独立头文件行号。
\end{itemize}

\textbf{说明}

\begin{itemize}
\tightlist
\item
  示例是局部调用片段；\passthrough{\lstinline!obj!}
  和参数变量表示已经按上层业务校验并准备好的对象。
\end{itemize}

\textbf{用法}

\begin{lstlisting}[language=Python]
helper = OsHelper.instance()
\end{lstlisting}

\hypertarget{unitree-sdk2-cpp-oshelper-get-uid-1}{%
\paragraph{\texorpdfstring{\texttt{unitree\_sdk2\_cpp.OsHelper.get\_uid}}{unitree\_sdk2\_cpp.OsHelper.get\_uid}}\label{unitree-sdk2-cpp-oshelper-get-uid-1}}

返回当前进程用户的 Unix UID。

\textbf{签名}

\begin{lstlisting}[language=Python]
def get_uid(self) -> int
\end{lstlisting}

\textbf{可用性与安全性}

\textbf{\passthrough{\lstinline!AVAILABLE!} /
\passthrough{\lstinline!READ\_ONLY!}}：当前绑定源码已实现本机信息读取。该方法不初始化
DDS，也不访问机器人。

\textbf{参数}

无显式参数。实例方法中的 \passthrough{\lstinline!self!} 由 Python
自动传入。

\textbf{返回值}

\begin{longtable}[]{@{}
  >{\raggedright\arraybackslash}p{(\columnwidth - 4\tabcolsep) * \real{0.18}}
  >{\raggedright\arraybackslash}p{(\columnwidth - 4\tabcolsep) * \real{0.20}}
  >{\raggedright\arraybackslash}p{(\columnwidth - 4\tabcolsep) * \real{0.62}}@{}}
\toprule
位置 & 类型 & 含义 \\
\midrule
\endhead
返回值 & \passthrough{\lstinline!int!} & 当前进程用户的 Unix UID。 \\
\bottomrule
\end{longtable}

\textbf{对应 C++}

\begin{itemize}
\tightlist
\item
  类：\passthrough{\lstinline!unitree\_sdk2\_cpp.OsHelper!}
\item
  签名：\passthrough{\lstinline!OsHelper::GetUID()!}
\item
  绑定策略：\passthrough{\lstinline!未单独标注!}
\item
  声明位置：由 pybind11 手工绑定或生成注册表提供；manifest
  未记录独立头文件行号。
\end{itemize}

\textbf{说明}

\begin{itemize}
\tightlist
\item
  示例是局部调用片段；\passthrough{\lstinline!obj!}
  和参数变量表示已经按上层业务校验并准备好的对象。
\end{itemize}

\textbf{用法}

\begin{lstlisting}[language=Python]
uid = helper.get_uid()
\end{lstlisting}

\hypertarget{unitree-sdk2-cpp-oshelper-get-gid-1}{%
\paragraph{\texorpdfstring{\texttt{unitree\_sdk2\_cpp.OsHelper.get\_gid}}{unitree\_sdk2\_cpp.OsHelper.get\_gid}}\label{unitree-sdk2-cpp-oshelper-get-gid-1}}

返回当前进程用户的 Unix GID。

\textbf{签名}

\begin{lstlisting}[language=Python]
def get_gid(self) -> int
\end{lstlisting}

\textbf{可用性与安全性}

\textbf{\passthrough{\lstinline!AVAILABLE!} /
\passthrough{\lstinline!READ\_ONLY!}}：当前绑定源码已实现本机信息读取。该方法不初始化
DDS，也不访问机器人。

\textbf{参数}

无显式参数。实例方法中的 \passthrough{\lstinline!self!} 由 Python
自动传入。

\textbf{返回值}

\begin{longtable}[]{@{}
  >{\raggedright\arraybackslash}p{(\columnwidth - 4\tabcolsep) * \real{0.18}}
  >{\raggedright\arraybackslash}p{(\columnwidth - 4\tabcolsep) * \real{0.20}}
  >{\raggedright\arraybackslash}p{(\columnwidth - 4\tabcolsep) * \real{0.62}}@{}}
\toprule
位置 & 类型 & 含义 \\
\midrule
\endhead
返回值 & \passthrough{\lstinline!int!} & 当前进程用户的 Unix GID。 \\
\bottomrule
\end{longtable}

\textbf{对应 C++}

\begin{itemize}
\tightlist
\item
  类：\passthrough{\lstinline!unitree\_sdk2\_cpp.OsHelper!}
\item
  签名：\passthrough{\lstinline!OsHelper::GetGID()!}
\item
  绑定策略：\passthrough{\lstinline!未单独标注!}
\item
  声明位置：由 pybind11 手工绑定或生成注册表提供；manifest
  未记录独立头文件行号。
\end{itemize}

\textbf{说明}

\begin{itemize}
\tightlist
\item
  示例是局部调用片段；\passthrough{\lstinline!obj!}
  和参数变量表示已经按上层业务校验并准备好的对象。
\end{itemize}

\textbf{用法}

\begin{lstlisting}[language=Python]
gid = helper.get_gid()
\end{lstlisting}

\hypertarget{unitree-sdk2-cpp-oshelper-get-user-1}{%
\paragraph{\texorpdfstring{\texttt{unitree\_sdk2\_cpp.OsHelper.get\_user}}{unitree\_sdk2\_cpp.OsHelper.get\_user}}\label{unitree-sdk2-cpp-oshelper-get-user-1}}

返回当前进程对应的用户名。

\textbf{签名}

\begin{lstlisting}[language=Python]
def get_user(self) -> str
\end{lstlisting}

\textbf{可用性与安全性}

\textbf{\passthrough{\lstinline!AVAILABLE!} /
\passthrough{\lstinline!READ\_ONLY!}}：当前绑定源码已实现本机信息读取。该方法不初始化
DDS，也不访问机器人。

\textbf{参数}

无显式参数。实例方法中的 \passthrough{\lstinline!self!} 由 Python
自动传入。

\textbf{返回值}

\begin{longtable}[]{@{}
  >{\raggedright\arraybackslash}p{(\columnwidth - 4\tabcolsep) * \real{0.18}}
  >{\raggedright\arraybackslash}p{(\columnwidth - 4\tabcolsep) * \real{0.20}}
  >{\raggedright\arraybackslash}p{(\columnwidth - 4\tabcolsep) * \real{0.62}}@{}}
\toprule
位置 & 类型 & 含义 \\
\midrule
\endhead
返回值 & \passthrough{\lstinline!str!} & 当前进程对应的用户名。 \\
\bottomrule
\end{longtable}

\textbf{对应 C++}

\begin{itemize}
\tightlist
\item
  类：\passthrough{\lstinline!unitree\_sdk2\_cpp.OsHelper!}
\item
  签名：\passthrough{\lstinline!OsHelper::GetUser()!}
\item
  绑定策略：\passthrough{\lstinline!未单独标注!}
\item
  声明位置：由 pybind11 手工绑定或生成注册表提供；manifest
  未记录独立头文件行号。
\end{itemize}

\textbf{说明}

\begin{itemize}
\tightlist
\item
  示例是局部调用片段；\passthrough{\lstinline!obj!}
  和参数变量表示已经按上层业务校验并准备好的对象。
\end{itemize}

\textbf{用法}

\begin{lstlisting}[language=Python]
user = helper.get_user()
\end{lstlisting}

\hypertarget{unitree-sdk2-cpp-oshelper-get-processor-number-1}{%
\paragraph{\texorpdfstring{\texttt{unitree\_sdk2\_cpp.OsHelper.get\_processor\_number}}{unitree\_sdk2\_cpp.OsHelper.get\_processor\_number}}\label{unitree-sdk2-cpp-oshelper-get-processor-number-1}}

返回当前主机可见的处理器数量。

\textbf{签名}

\begin{lstlisting}[language=Python]
def get_processor_number(self) -> int
\end{lstlisting}

\textbf{可用性与安全性}

\textbf{\passthrough{\lstinline!AVAILABLE!} /
\passthrough{\lstinline!READ\_ONLY!}}：当前绑定源码已实现本机信息读取。该方法不初始化
DDS，也不访问机器人。

\textbf{参数}

无显式参数。实例方法中的 \passthrough{\lstinline!self!} 由 Python
自动传入。

\textbf{返回值}

\begin{longtable}[]{@{}
  >{\raggedright\arraybackslash}p{(\columnwidth - 4\tabcolsep) * \real{0.18}}
  >{\raggedright\arraybackslash}p{(\columnwidth - 4\tabcolsep) * \real{0.20}}
  >{\raggedright\arraybackslash}p{(\columnwidth - 4\tabcolsep) * \real{0.62}}@{}}
\toprule
位置 & 类型 & 含义 \\
\midrule
\endhead
返回值 & \passthrough{\lstinline!int!} & 当前主机可见的处理器数量。 \\
\bottomrule
\end{longtable}

\textbf{对应 C++}

\begin{itemize}
\tightlist
\item
  类：\passthrough{\lstinline!unitree\_sdk2\_cpp.OsHelper!}
\item
  签名：\passthrough{\lstinline!OsHelper::GetProcessorNumber()!}
\item
  绑定策略：\passthrough{\lstinline!未单独标注!}
\item
  声明位置：由 pybind11 手工绑定或生成注册表提供；manifest
  未记录独立头文件行号。
\end{itemize}

\textbf{说明}

\begin{itemize}
\tightlist
\item
  示例是局部调用片段；\passthrough{\lstinline!obj!}
  和参数变量表示已经按上层业务校验并准备好的对象。
\end{itemize}

\textbf{用法}

\begin{lstlisting}[language=Python]
processor_count = helper.get_processor_number()
\end{lstlisting}

\hypertarget{unitree-sdk2-cpp-oshelper-get-page-size-1}{%
\paragraph{\texorpdfstring{\texttt{unitree\_sdk2\_cpp.OsHelper.get\_page\_size}}{unitree\_sdk2\_cpp.OsHelper.get\_page\_size}}\label{unitree-sdk2-cpp-oshelper-get-page-size-1}}

返回当前主机的操作系统内存页大小。

\textbf{签名}

\begin{lstlisting}[language=Python]
def get_page_size(self) -> int
\end{lstlisting}

\textbf{可用性与安全性}

\textbf{\passthrough{\lstinline!AVAILABLE!} /
\passthrough{\lstinline!READ\_ONLY!}}：当前绑定源码已实现本机信息读取。该方法不初始化
DDS，也不访问机器人。

\textbf{参数}

无显式参数。实例方法中的 \passthrough{\lstinline!self!} 由 Python
自动传入。

\textbf{返回值}

\begin{longtable}[]{@{}
  >{\raggedright\arraybackslash}p{(\columnwidth - 4\tabcolsep) * \real{0.18}}
  >{\raggedright\arraybackslash}p{(\columnwidth - 4\tabcolsep) * \real{0.20}}
  >{\raggedright\arraybackslash}p{(\columnwidth - 4\tabcolsep) * \real{0.62}}@{}}
\toprule
位置 & 类型 & 含义 \\
\midrule
\endhead
返回值 & \passthrough{\lstinline!int!} &
操作系统内存页大小，单位为字节。 \\
\bottomrule
\end{longtable}

\textbf{对应 C++}

\begin{itemize}
\tightlist
\item
  类：\passthrough{\lstinline!unitree\_sdk2\_cpp.OsHelper!}
\item
  签名：\passthrough{\lstinline!OsHelper::GetPageSize()!}
\item
  绑定策略：\passthrough{\lstinline!未单独标注!}
\item
  声明位置：由 pybind11 手工绑定或生成注册表提供；manifest
  未记录独立头文件行号。
\end{itemize}

\textbf{说明}

\begin{itemize}
\tightlist
\item
  示例是局部调用片段；\passthrough{\lstinline!obj!}
  和参数变量表示已经按上层业务校验并准备好的对象。
\end{itemize}

\textbf{用法}

\begin{lstlisting}[language=Python]
page_size = helper.get_page_size()
\end{lstlisting}

\hypertarget{unitree-sdk2-cpp-oshelper-get-hostname-1}{%
\paragraph{\texorpdfstring{\texttt{unitree\_sdk2\_cpp.OsHelper.get\_hostname}}{unitree\_sdk2\_cpp.OsHelper.get\_hostname}}\label{unitree-sdk2-cpp-oshelper-get-hostname-1}}

返回当前主机名。

\textbf{签名}

\begin{lstlisting}[language=Python]
def get_hostname(self) -> str
\end{lstlisting}

\textbf{可用性与安全性}

\textbf{\passthrough{\lstinline!AVAILABLE!} /
\passthrough{\lstinline!READ\_ONLY!}}：当前绑定源码已实现本机信息读取。该方法不初始化
DDS，也不访问机器人。

\textbf{参数}

无显式参数。实例方法中的 \passthrough{\lstinline!self!} 由 Python
自动传入。

\textbf{返回值}

\begin{longtable}[]{@{}
  >{\raggedright\arraybackslash}p{(\columnwidth - 4\tabcolsep) * \real{0.18}}
  >{\raggedright\arraybackslash}p{(\columnwidth - 4\tabcolsep) * \real{0.20}}
  >{\raggedright\arraybackslash}p{(\columnwidth - 4\tabcolsep) * \real{0.62}}@{}}
\toprule
位置 & 类型 & 含义 \\
\midrule
\endhead
返回值 & \passthrough{\lstinline!str!} & 当前主机名。 \\
\bottomrule
\end{longtable}

\textbf{对应 C++}

\begin{itemize}
\tightlist
\item
  类：\passthrough{\lstinline!unitree\_sdk2\_cpp.OsHelper!}
\item
  签名：\passthrough{\lstinline!OsHelper::GetHostname()!}
\item
  绑定策略：\passthrough{\lstinline!未单独标注!}
\item
  声明位置：由 pybind11 手工绑定或生成注册表提供；manifest
  未记录独立头文件行号。
\end{itemize}

\textbf{说明}

\begin{itemize}
\tightlist
\item
  示例是局部调用片段；\passthrough{\lstinline!obj!}
  和参数变量表示已经按上层业务校验并准备好的对象。
\end{itemize}

\textbf{用法}

\begin{lstlisting}[language=Python]
hostname = helper.get_hostname()
\end{lstlisting}

\begin{center}\rule{0.5\linewidth}{0.5pt}\end{center}

\hypertarget{unitree-sdk2-cpp-channel}{%
\subsection{Typed DDS Channel API}\label{unitree-sdk2-cpp-channel}}

模块：\passthrough{\lstinline!unitree\_sdk2\_cpp.channel!}

\hypertarget{ux6a21ux5757ux51fdux6570ux5bfcux5165}{%
\subsubsection{模块函数导入}\label{ux6a21ux5757ux51fdux6570ux5bfcux5165}}

\begin{lstlisting}[language=Python]
from unitree_sdk2_cpp.channel import registered_message_types, initialize, initialize_from_config, release
\end{lstlisting}

下面的函数用法片段假设已经完成上述导入。

\hypertarget{ux6a21ux5757ux51fdux6570ux7d22ux5f15}{%
\subsubsection{模块函数索引}\label{ux6a21ux5757ux51fdux6570ux7d22ux5f15}}

\begin{longtable}[]{@{}
  >{\raggedright\arraybackslash}p{(\columnwidth - 6\tabcolsep) * \real{0.18}}
  >{\raggedright\arraybackslash}p{(\columnwidth - 6\tabcolsep) * \real{0.24}}
  >{\raggedright\arraybackslash}p{(\columnwidth - 6\tabcolsep) * \real{0.18}}
  >{\raggedright\arraybackslash}p{(\columnwidth - 6\tabcolsep) * \real{0.40}}@{}}
\toprule
函数 & Python 签名 & 状态 & 安全分类 \\
\midrule
\endhead
\protect\hyperlink{unitree-sdk2-cpp-channel-registered-message-types-1}{\passthrough{\lstinline!registered\_message\_types()!}}
&
\passthrough{\lstinline!def registered\_message\_types() -> list[str]!}
& \passthrough{\lstinline!AVAILABLE!} &
\passthrough{\lstinline!DDS\_LIFECYCLE!} \\
\protect\hyperlink{unitree-sdk2-cpp-channel-initialize-1}{\passthrough{\lstinline!initialize()!}}
&
\passthrough{\lstinline!def initialize(domain\_id: int = 0, network\_interface: str = '') -> None!}
& \passthrough{\lstinline!AVAILABLE!} &
\passthrough{\lstinline!DDS\_LIFECYCLE!} \\
\protect\hyperlink{unitree-sdk2-cpp-channel-initialize-from-config-1}{\passthrough{\lstinline!initialize\_from\_config()!}}
&
\passthrough{\lstinline!def initialize\_from\_config(config\_file: str = '') -> None!}
& \passthrough{\lstinline!AVAILABLE!} &
\passthrough{\lstinline!DDS\_LIFECYCLE!} \\
\protect\hyperlink{unitree-sdk2-cpp-channel-release-1}{\passthrough{\lstinline!release()!}}
& \passthrough{\lstinline!def release() -> None!} &
\passthrough{\lstinline!AVAILABLE!} &
\passthrough{\lstinline!DDS\_LIFECYCLE!} \\
\bottomrule
\end{longtable}

\hypertarget{ux7c7bux7d22ux5f15-1}{%
\subsubsection{类索引}\label{ux7c7bux7d22ux5f15-1}}

\begin{longtable}[]{@{}
  >{\raggedright\arraybackslash}p{(\columnwidth - 4\tabcolsep) * \real{0.18}}
  >{\raggedleft\arraybackslash}p{(\columnwidth - 4\tabcolsep) * \real{0.20}}
  >{\raggedleft\arraybackslash}p{(\columnwidth - 4\tabcolsep) * \real{0.62}}@{}}
\toprule
类 & 公开函数签名 & 属性 \\
\midrule
\endhead
\protect\hyperlink{unitree-sdk2-cpp-channel-channelpublisher}{\passthrough{\lstinline!ChannelPublisher!}}
& 4 & 2 \\
\protect\hyperlink{unitree-sdk2-cpp-channel-channelsubscriber}{\passthrough{\lstinline!ChannelSubscriber!}}
& 3 & 3 \\
\bottomrule
\end{longtable}

\hypertarget{unitree-sdk2-cpp-channel-registered-message-types-1}{%
\paragraph{\texorpdfstring{\texttt{unitree\_sdk2\_cpp.channel.registered\_message\_types}}{unitree\_sdk2\_cpp.channel.registered\_message\_types}}\label{unitree-sdk2-cpp-channel-registered-message-types-1}}

返回当前二进制扩展已经注册的 typed DDS 消息类型名。

\textbf{签名}

\begin{lstlisting}[language=Python]
def registered_message_types() -> list[str]
\end{lstlisting}

\textbf{可用性与安全性}

\textbf{\passthrough{\lstinline!AVAILABLE!} /
\passthrough{\lstinline!DDS\_LIFECYCLE!}}：当前绑定源码已实现。该操作可能创建、使用或释放
DDS 资源。

\textbf{参数}

无参数。

\textbf{返回值}

\begin{longtable}[]{@{}
  >{\raggedright\arraybackslash}p{(\columnwidth - 4\tabcolsep) * \real{0.18}}
  >{\raggedright\arraybackslash}p{(\columnwidth - 4\tabcolsep) * \real{0.20}}
  >{\raggedright\arraybackslash}p{(\columnwidth - 4\tabcolsep) * \real{0.62}}@{}}
\toprule
位置 & 类型 & 含义 \\
\midrule
\endhead
返回值 & \passthrough{\lstinline!list[str]!} & 当前扩展已注册的 C++ DDS
消息类型名列表。 \\
\bottomrule
\end{longtable}

\textbf{对应 C++}

\begin{itemize}
\tightlist
\item
  类：\passthrough{\lstinline!unitree\_sdk2\_cpp.channel!}
\item
  签名：\passthrough{\lstinline!registered\_message\_types()!}
\item
  绑定策略：\passthrough{\lstinline!未单独标注!}
\item
  声明位置：由 pybind11 手工绑定或生成注册表提供；manifest
  未记录独立头文件行号。
\end{itemize}

\textbf{说明}

\begin{itemize}
\tightlist
\item
  示例是局部调用片段；\passthrough{\lstinline!obj!}
  和参数变量表示已经按上层业务校验并准备好的对象。
\end{itemize}

\textbf{用法}

\begin{lstlisting}[language=Python]
message_type_names = registered_message_types()
\end{lstlisting}

\hypertarget{unitree-sdk2-cpp-channel-initialize-1}{%
\paragraph{\texorpdfstring{\texttt{unitree\_sdk2\_cpp.channel.initialize}}{unitree\_sdk2\_cpp.channel.initialize}}\label{unitree-sdk2-cpp-channel-initialize-1}}

使用 Domain ID 和网络接口初始化全局 DDS ChannelFactory。

\textbf{签名}

\begin{lstlisting}[language=Python]
def initialize(domain_id: int = 0, network_interface: str = '') -> None
\end{lstlisting}

\textbf{可用性与安全性}

\textbf{\passthrough{\lstinline!AVAILABLE!} /
\passthrough{\lstinline!DDS\_LIFECYCLE!}}：当前绑定源码已实现。该操作可能创建、使用或释放
DDS 资源。

\textbf{参数}

\begin{longtable}[]{@{}
  >{\raggedright\arraybackslash}p{(\columnwidth - 6\tabcolsep) * \real{0.18}}
  >{\raggedright\arraybackslash}p{(\columnwidth - 6\tabcolsep) * \real{0.24}}
  >{\raggedright\arraybackslash}p{(\columnwidth - 6\tabcolsep) * \real{0.18}}
  >{\raggedright\arraybackslash}p{(\columnwidth - 6\tabcolsep) * \real{0.40}}@{}}
\toprule
名称 & Python 类型 & 默认值 & 含义 \\
\midrule
\endhead
\passthrough{\lstinline!domain\_id!} & \passthrough{\lstinline!int!} &
\passthrough{\lstinline!0!} & DDS Domain ID。只有使用兼容 domain
的参与者才能按预期互相发现。 \\
\passthrough{\lstinline!network\_interface!} &
\passthrough{\lstinline!str!} & \passthrough{\lstinline!''!} & DDS
使用的网络接口名称，例如 Linux 上的 \passthrough{\lstinline!eth0!} 或
\passthrough{\lstinline!enp3s0!}；必须按目标机实际网卡填写。 \\
\bottomrule
\end{longtable}

\textbf{返回值}

\begin{longtable}[]{@{}
  >{\raggedright\arraybackslash}p{(\columnwidth - 4\tabcolsep) * \real{0.18}}
  >{\raggedright\arraybackslash}p{(\columnwidth - 4\tabcolsep) * \real{0.20}}
  >{\raggedright\arraybackslash}p{(\columnwidth - 4\tabcolsep) * \real{0.62}}@{}}
\toprule
位置 & 类型 & 含义 \\
\midrule
\endhead
无 & \passthrough{\lstinline!None!} & 该方法不返回 Python
值；失败可能表现为异常或底层状态变化。 \\
\bottomrule
\end{longtable}

\textbf{对应 C++}

\begin{itemize}
\tightlist
\item
  类：\passthrough{\lstinline!unitree\_sdk2\_cpp.channel!}
\item
  签名：\passthrough{\lstinline!ChannelFactory::Init(int32\_t, const std::string \&)!}
\item
  绑定策略：\passthrough{\lstinline!未单独标注!}
\item
  声明位置：由 pybind11 手工绑定或生成注册表提供；manifest
  未记录独立头文件行号。
\end{itemize}

\textbf{说明}

\begin{itemize}
\tightlist
\item
  示例是局部调用片段；\passthrough{\lstinline!obj!}
  和参数变量表示已经按上层业务校验并准备好的对象。
\end{itemize}

\textbf{用法}

\begin{lstlisting}[language=Python]
initialize(domain_id=0, network_interface="eth0")
\end{lstlisting}

\hypertarget{unitree-sdk2-cpp-channel-initialize-from-config-1}{%
\paragraph{\texorpdfstring{\texttt{unitree\_sdk2\_cpp.channel.initialize\_from\_config}}{unitree\_sdk2\_cpp.channel.initialize\_from\_config}}\label{unitree-sdk2-cpp-channel-initialize-from-config-1}}

使用 CycloneDDS 配置文件初始化全局 DDS ChannelFactory。

\textbf{签名}

\begin{lstlisting}[language=Python]
def initialize_from_config(config_file: str = '') -> None
\end{lstlisting}

\textbf{可用性与安全性}

\textbf{\passthrough{\lstinline!AVAILABLE!} /
\passthrough{\lstinline!DDS\_LIFECYCLE!}}：当前绑定源码已实现。该操作可能创建、使用或释放
DDS 资源。

\textbf{参数}

\begin{longtable}[]{@{}
  >{\raggedright\arraybackslash}p{(\columnwidth - 6\tabcolsep) * \real{0.18}}
  >{\raggedright\arraybackslash}p{(\columnwidth - 6\tabcolsep) * \real{0.24}}
  >{\raggedright\arraybackslash}p{(\columnwidth - 6\tabcolsep) * \real{0.18}}
  >{\raggedright\arraybackslash}p{(\columnwidth - 6\tabcolsep) * \real{0.40}}@{}}
\toprule
名称 & Python 类型 & 默认值 & 含义 \\
\midrule
\endhead
\passthrough{\lstinline!config\_file!} & \passthrough{\lstinline!str!} &
\passthrough{\lstinline!''!} & CycloneDDS
配置文件路径；空字符串表示使用底层默认配置。 \\
\bottomrule
\end{longtable}

\textbf{返回值}

\begin{longtable}[]{@{}
  >{\raggedright\arraybackslash}p{(\columnwidth - 4\tabcolsep) * \real{0.18}}
  >{\raggedright\arraybackslash}p{(\columnwidth - 4\tabcolsep) * \real{0.20}}
  >{\raggedright\arraybackslash}p{(\columnwidth - 4\tabcolsep) * \real{0.62}}@{}}
\toprule
位置 & 类型 & 含义 \\
\midrule
\endhead
无 & \passthrough{\lstinline!None!} & 该方法不返回 Python
值；失败可能表现为异常或底层状态变化。 \\
\bottomrule
\end{longtable}

\textbf{对应 C++}

\begin{itemize}
\tightlist
\item
  类：\passthrough{\lstinline!unitree\_sdk2\_cpp.channel!}
\item
  签名：\passthrough{\lstinline!ChannelFactory::Init(const std::string \&)!}
\item
  绑定策略：\passthrough{\lstinline!未单独标注!}
\item
  声明位置：由 pybind11 手工绑定或生成注册表提供；manifest
  未记录独立头文件行号。
\end{itemize}

\textbf{说明}

\begin{itemize}
\tightlist
\item
  示例是局部调用片段；\passthrough{\lstinline!obj!}
  和参数变量表示已经按上层业务校验并准备好的对象。
\end{itemize}

\textbf{用法}

\begin{lstlisting}[language=Python]
initialize_from_config(config_file="cyclonedds.xml")
\end{lstlisting}

\hypertarget{unitree-sdk2-cpp-channel-release-1}{%
\paragraph{\texorpdfstring{\texttt{unitree\_sdk2\_cpp.channel.release}}{unitree\_sdk2\_cpp.channel.release}}\label{unitree-sdk2-cpp-channel-release-1}}

释放全局 DDS ChannelFactory 及其持有的底层资源。

\textbf{签名}

\begin{lstlisting}[language=Python]
def release() -> None
\end{lstlisting}

\textbf{可用性与安全性}

\textbf{\passthrough{\lstinline!AVAILABLE!} /
\passthrough{\lstinline!DDS\_LIFECYCLE!}}：当前绑定源码已实现。该操作可能创建、使用或释放
DDS 资源。

\textbf{参数}

无参数。

\textbf{返回值}

\begin{longtable}[]{@{}
  >{\raggedright\arraybackslash}p{(\columnwidth - 4\tabcolsep) * \real{0.18}}
  >{\raggedright\arraybackslash}p{(\columnwidth - 4\tabcolsep) * \real{0.20}}
  >{\raggedright\arraybackslash}p{(\columnwidth - 4\tabcolsep) * \real{0.62}}@{}}
\toprule
位置 & 类型 & 含义 \\
\midrule
\endhead
无 & \passthrough{\lstinline!None!} & 该方法不返回 Python
值；失败可能表现为异常或底层状态变化。 \\
\bottomrule
\end{longtable}

\textbf{对应 C++}

\begin{itemize}
\tightlist
\item
  类：\passthrough{\lstinline!unitree\_sdk2\_cpp.channel!}
\item
  签名：\passthrough{\lstinline!ChannelFactory::Release()!}
\item
  绑定策略：\passthrough{\lstinline!未单独标注!}
\item
  声明位置：由 pybind11 手工绑定或生成注册表提供；manifest
  未记录独立头文件行号。
\end{itemize}

\textbf{说明}

\begin{itemize}
\tightlist
\item
  示例是局部调用片段；\passthrough{\lstinline!obj!}
  和参数变量表示已经按上层业务校验并准备好的对象。
\end{itemize}

\textbf{用法}

\begin{lstlisting}[language=Python]
release()
\end{lstlisting}

\hypertarget{unitree-sdk2-cpp-channel-channelpublisher}{%
\subsubsection{\texorpdfstring{\texttt{unitree\_sdk2\_cpp.channel.ChannelPublisher}}{unitree\_sdk2\_cpp.channel.ChannelPublisher}}\label{unitree-sdk2-cpp-channel-channelpublisher}}

typed DDS 通道包装类。必须遵守初始化、启动、关闭和全局释放顺序。

\textbf{导入}

\begin{lstlisting}[language=Python]
from unitree_sdk2_cpp.channel import ChannelPublisher
\end{lstlisting}

\textbf{基类}：\passthrough{\lstinline!Generic[MessageT]!}

\textbf{方法索引}

\begin{longtable}[]{@{}
  >{\raggedright\arraybackslash}p{(\columnwidth - 6\tabcolsep) * \real{0.18}}
  >{\raggedright\arraybackslash}p{(\columnwidth - 6\tabcolsep) * \real{0.24}}
  >{\raggedright\arraybackslash}p{(\columnwidth - 6\tabcolsep) * \real{0.18}}
  >{\raggedright\arraybackslash}p{(\columnwidth - 6\tabcolsep) * \real{0.40}}@{}}
\toprule
方法 & Python 签名 & 状态 & 安全分类 \\
\midrule
\endhead
\protect\hyperlink{unitree-sdk2-cpp-channel-channelpublisher-dunder-init-1}{\passthrough{\lstinline!\_\_init\_\_!}}
&
\passthrough{\lstinline!def \_\_init\_\_(self, topic: str, message\_type: type[MessageT]) -> None!}
& \passthrough{\lstinline!AVAILABLE!} &
\passthrough{\lstinline!DDS\_LIFECYCLE!} \\
\protect\hyperlink{unitree-sdk2-cpp-channel-channelpublisher-init-channel-1}{\passthrough{\lstinline!init\_channel!}}
& \passthrough{\lstinline!def init\_channel(self) -> None!} &
\passthrough{\lstinline!AVAILABLE!} &
\passthrough{\lstinline!DDS\_LIFECYCLE!} \\
\protect\hyperlink{unitree-sdk2-cpp-channel-channelpublisher-close-channel-1}{\passthrough{\lstinline!close\_channel!}}
& \passthrough{\lstinline!def close\_channel(self) -> None!} &
\passthrough{\lstinline!AVAILABLE!} &
\passthrough{\lstinline!DDS\_LIFECYCLE!} \\
\protect\hyperlink{unitree-sdk2-cpp-channel-channelpublisher-write-1}{\passthrough{\lstinline!write!}}
&
\passthrough{\lstinline!def write(self, message: MessageT, wait\_microsec: int = 0) -> bool!}
& \passthrough{\lstinline!AVAILABLE!} &
\passthrough{\lstinline!DDS\_LIFECYCLE!} \\
\bottomrule
\end{longtable}

\textbf{属性索引}

\begin{longtable}[]{@{}
  >{\raggedright\arraybackslash}p{(\columnwidth - 6\tabcolsep) * \real{0.18}}
  >{\raggedright\arraybackslash}p{(\columnwidth - 6\tabcolsep) * \real{0.24}}
  >{\raggedright\arraybackslash}p{(\columnwidth - 6\tabcolsep) * \real{0.18}}
  >{\raggedright\arraybackslash}p{(\columnwidth - 6\tabcolsep) * \real{0.40}}@{}}
\toprule
属性 & 读取类型 & 写入类型 & 状态 \\
\midrule
\endhead
\protect\hyperlink{unitree-sdk2-cpp-channel-channelpublisher-topic}{\passthrough{\lstinline!topic!}}
& \passthrough{\lstinline!str!} & \passthrough{\lstinline!只读!} &
\passthrough{\lstinline!AVAILABLE!} \\
\protect\hyperlink{unitree-sdk2-cpp-channel-channelpublisher-message-type-name}{\passthrough{\lstinline!message\_type\_name!}}
& \passthrough{\lstinline!str!} & \passthrough{\lstinline!只读!} &
\passthrough{\lstinline!AVAILABLE!} \\
\bottomrule
\end{longtable}

\hypertarget{unitree-sdk2-cpp-channel-channelpublisher-dunder-init-1}{%
\paragraph{\texorpdfstring{\texttt{unitree\_sdk2\_cpp.channel.ChannelPublisher.\_\_init\_\_}}{unitree\_sdk2\_cpp.channel.ChannelPublisher.\_\_init\_\_}}\label{unitree-sdk2-cpp-channel-channelpublisher-dunder-init-1}}

初始化 \passthrough{\lstinline!ChannelPublisher!}
实例。是否能实际构造取决于下方可用性状态。

\textbf{签名}

\begin{lstlisting}[language=Python]
def __init__(self, topic: str, message_type: type[MessageT]) -> None
\end{lstlisting}

\textbf{可用性与安全性}

\textbf{\passthrough{\lstinline!AVAILABLE!} /
\passthrough{\lstinline!DDS\_LIFECYCLE!}}：当前绑定源码已实现。该操作可能创建、使用或释放
DDS 资源。

\textbf{参数}

\begin{longtable}[]{@{}
  >{\raggedright\arraybackslash}p{(\columnwidth - 6\tabcolsep) * \real{0.18}}
  >{\raggedright\arraybackslash}p{(\columnwidth - 6\tabcolsep) * \real{0.24}}
  >{\raggedright\arraybackslash}p{(\columnwidth - 6\tabcolsep) * \real{0.18}}
  >{\raggedright\arraybackslash}p{(\columnwidth - 6\tabcolsep) * \real{0.40}}@{}}
\toprule
名称 & Python 类型 & 默认值 & 含义 \\
\midrule
\endhead
\passthrough{\lstinline!topic!} & \passthrough{\lstinline!str!} & 必填 &
DDS topic 名称。发布端和订阅端必须同时使用兼容的 topic、消息类型和
QoS。 \\
\passthrough{\lstinline!message\_type!} &
\passthrough{\lstinline!type[MessageT]!} & 必填 & IDL 消息类本身，例如
\passthrough{\lstinline!LowState!}，不是
\passthrough{\lstinline!LowState()!}
实例。该类型必须存在于运行时注册表中。 \\
\bottomrule
\end{longtable}

\textbf{返回值}

\begin{longtable}[]{@{}
  >{\raggedright\arraybackslash}p{(\columnwidth - 4\tabcolsep) * \real{0.18}}
  >{\raggedright\arraybackslash}p{(\columnwidth - 4\tabcolsep) * \real{0.20}}
  >{\raggedright\arraybackslash}p{(\columnwidth - 4\tabcolsep) * \real{0.62}}@{}}
\toprule
位置 & 类型 & 含义 \\
\midrule
\endhead
无 & \passthrough{\lstinline!None!} &
\passthrough{\lstinline!\_\_init\_\_!}
本身不返回值；调用类对象时，在构造函数可用的前提下得到该类实例。 \\
\bottomrule
\end{longtable}

\textbf{对应 C++}

\begin{itemize}
\tightlist
\item
  类：\passthrough{\lstinline!unitree\_sdk2\_cpp.channel.ChannelPublisher!}
\item
  签名：\passthrough{\lstinline!ChannelPublisher(const std::string \&, type)!}
\item
  绑定策略：\passthrough{\lstinline!未单独标注!}
\item
  声明位置：由 pybind11 手工绑定或生成注册表提供；manifest
  未记录独立头文件行号。
\end{itemize}

\textbf{说明}

\begin{itemize}
\tightlist
\item
  示例是局部调用片段；\passthrough{\lstinline!obj!}
  和参数变量表示已经按上层业务校验并准备好的对象。
\end{itemize}

\textbf{用法}

\begin{lstlisting}[language=Python]
value = ChannelPublisher(topic=topic, message_type=message_type)
\end{lstlisting}

\hypertarget{unitree-sdk2-cpp-channel-channelpublisher-init-channel-1}{%
\paragraph{\texorpdfstring{\texttt{unitree\_sdk2\_cpp.channel.ChannelPublisher.init\_channel}}{unitree\_sdk2\_cpp.channel.ChannelPublisher.init\_channel}}\label{unitree-sdk2-cpp-channel-channelpublisher-init-channel-1}}

创建并启动 Publisher 的底层 DDS 写通道。

\textbf{签名}

\begin{lstlisting}[language=Python]
def init_channel(self) -> None
\end{lstlisting}

\textbf{可用性与安全性}

\textbf{\passthrough{\lstinline!AVAILABLE!} /
\passthrough{\lstinline!DDS\_LIFECYCLE!}}：当前绑定源码已实现。该操作可能创建、使用或释放
DDS 资源。

\textbf{参数}

无显式参数。实例方法中的 \passthrough{\lstinline!self!} 由 Python
自动传入。

\textbf{返回值}

\begin{longtable}[]{@{}
  >{\raggedright\arraybackslash}p{(\columnwidth - 4\tabcolsep) * \real{0.18}}
  >{\raggedright\arraybackslash}p{(\columnwidth - 4\tabcolsep) * \real{0.20}}
  >{\raggedright\arraybackslash}p{(\columnwidth - 4\tabcolsep) * \real{0.62}}@{}}
\toprule
位置 & 类型 & 含义 \\
\midrule
\endhead
无 & \passthrough{\lstinline!None!} & 该方法不返回 Python
值；失败可能表现为异常或底层状态变化。 \\
\bottomrule
\end{longtable}

\textbf{对应 C++}

\begin{itemize}
\tightlist
\item
  类：\passthrough{\lstinline!unitree\_sdk2\_cpp.channel.ChannelPublisher!}
\item
  签名：\passthrough{\lstinline!ChannelPublisher::init\_channel()!}
\item
  绑定策略：\passthrough{\lstinline!未单独标注!}
\item
  声明位置：由 pybind11 手工绑定或生成注册表提供；manifest
  未记录独立头文件行号。
\end{itemize}

\textbf{说明}

\begin{itemize}
\tightlist
\item
  示例是局部调用片段；\passthrough{\lstinline!obj!}
  和参数变量表示已经按上层业务校验并准备好的对象。
\end{itemize}

\textbf{用法}

\begin{lstlisting}[language=Python]
obj.init_channel()
\end{lstlisting}

\hypertarget{unitree-sdk2-cpp-channel-channelpublisher-close-channel-1}{%
\paragraph{\texorpdfstring{\texttt{unitree\_sdk2\_cpp.channel.ChannelPublisher.close\_channel}}{unitree\_sdk2\_cpp.channel.ChannelPublisher.close\_channel}}\label{unitree-sdk2-cpp-channel-channelpublisher-close-channel-1}}

关闭 Publisher 的底层 DDS 写通道。

\textbf{签名}

\begin{lstlisting}[language=Python]
def close_channel(self) -> None
\end{lstlisting}

\textbf{可用性与安全性}

\textbf{\passthrough{\lstinline!AVAILABLE!} /
\passthrough{\lstinline!DDS\_LIFECYCLE!}}：当前绑定源码已实现。该操作可能创建、使用或释放
DDS 资源。

\textbf{参数}

无显式参数。实例方法中的 \passthrough{\lstinline!self!} 由 Python
自动传入。

\textbf{返回值}

\begin{longtable}[]{@{}
  >{\raggedright\arraybackslash}p{(\columnwidth - 4\tabcolsep) * \real{0.18}}
  >{\raggedright\arraybackslash}p{(\columnwidth - 4\tabcolsep) * \real{0.20}}
  >{\raggedright\arraybackslash}p{(\columnwidth - 4\tabcolsep) * \real{0.62}}@{}}
\toprule
位置 & 类型 & 含义 \\
\midrule
\endhead
无 & \passthrough{\lstinline!None!} & 该方法不返回 Python
值；失败可能表现为异常或底层状态变化。 \\
\bottomrule
\end{longtable}

\textbf{对应 C++}

\begin{itemize}
\tightlist
\item
  类：\passthrough{\lstinline!unitree\_sdk2\_cpp.channel.ChannelPublisher!}
\item
  签名：\passthrough{\lstinline!ChannelPublisher::close\_channel()!}
\item
  绑定策略：\passthrough{\lstinline!未单独标注!}
\item
  声明位置：由 pybind11 手工绑定或生成注册表提供；manifest
  未记录独立头文件行号。
\end{itemize}

\textbf{说明}

\begin{itemize}
\tightlist
\item
  示例是局部调用片段；\passthrough{\lstinline!obj!}
  和参数变量表示已经按上层业务校验并准备好的对象。
\end{itemize}

\textbf{用法}

\begin{lstlisting}[language=Python]
obj.close_channel()
\end{lstlisting}

\hypertarget{unitree-sdk2-cpp-channel-channelpublisher-write-1}{%
\paragraph{\texorpdfstring{\texttt{unitree\_sdk2\_cpp.channel.ChannelPublisher.write}}{unitree\_sdk2\_cpp.channel.ChannelPublisher.write}}\label{unitree-sdk2-cpp-channel-channelpublisher-write-1}}

把一个类型匹配的消息样本写入 Publisher 对应的 DDS topic。

\textbf{签名}

\begin{lstlisting}[language=Python]
def write(self, message: MessageT, wait_microsec: int = 0) -> bool
\end{lstlisting}

\textbf{可用性与安全性}

\textbf{\passthrough{\lstinline!AVAILABLE!} /
\passthrough{\lstinline!DDS\_LIFECYCLE!}}：当前绑定源码已实现。该操作可能创建、使用或释放
DDS 资源。

\textbf{参数}

\begin{longtable}[]{@{}
  >{\raggedright\arraybackslash}p{(\columnwidth - 6\tabcolsep) * \real{0.18}}
  >{\raggedright\arraybackslash}p{(\columnwidth - 6\tabcolsep) * \real{0.24}}
  >{\raggedright\arraybackslash}p{(\columnwidth - 6\tabcolsep) * \real{0.18}}
  >{\raggedright\arraybackslash}p{(\columnwidth - 6\tabcolsep) * \real{0.40}}@{}}
\toprule
名称 & Python 类型 & 默认值 & 含义 \\
\midrule
\endhead
\passthrough{\lstinline!message!} & \passthrough{\lstinline!MessageT!} &
必填 & 要写入 DDS 的消息实例；类型必须与 Publisher
构造时声明的消息类完全一致。 \\
\passthrough{\lstinline!wait\_microsec!} & \passthrough{\lstinline!int!}
& \passthrough{\lstinline!0!} &
写操作允许等待的时间，单位为微秒；\passthrough{\lstinline!0!}
使用当前绑定的默认非额外等待行为。 \\
\bottomrule
\end{longtable}

\textbf{返回值}

\begin{longtable}[]{@{}
  >{\raggedright\arraybackslash}p{(\columnwidth - 4\tabcolsep) * \real{0.18}}
  >{\raggedright\arraybackslash}p{(\columnwidth - 4\tabcolsep) * \real{0.20}}
  >{\raggedright\arraybackslash}p{(\columnwidth - 4\tabcolsep) * \real{0.62}}@{}}
\toprule
位置 & 类型 & 含义 \\
\midrule
\endhead
返回值 & \passthrough{\lstinline!bool!} &
底层写操作是否被接受；不表示所有订阅者已处理，更不表示设备已经执行动作。 \\
\bottomrule
\end{longtable}

\textbf{对应 C++}

\begin{itemize}
\tightlist
\item
  类：\passthrough{\lstinline!unitree\_sdk2\_cpp.channel.ChannelPublisher!}
\item
  签名：\passthrough{\lstinline!ChannelPublisher::write(message, int64\_t)!}
\item
  绑定策略：\passthrough{\lstinline!未单独标注!}
\item
  声明位置：由 pybind11 手工绑定或生成注册表提供；manifest
  未记录独立头文件行号。
\end{itemize}

\textbf{说明}

\begin{itemize}
\tightlist
\item
  示例是局部调用片段；\passthrough{\lstinline!obj!}
  和参数变量表示已经按上层业务校验并准备好的对象。
\end{itemize}

\textbf{用法}

\begin{lstlisting}[language=Python]
result = obj.write(message=message, wait_microsec=wait_microsec)
\end{lstlisting}

\hypertarget{unitree-sdk2-cpp-channel-channelpublisher-topic}{%
\paragraph{\texorpdfstring{\texttt{unitree\_sdk2\_cpp.channel.ChannelPublisher.topic}}{unitree\_sdk2\_cpp.channel.ChannelPublisher.topic}}\label{unitree-sdk2-cpp-channel-channelpublisher-topic}}

Publisher 构造时保存的 DDS topic 名称，只读。

\textbf{签名}

\begin{lstlisting}[language=Python]
@property
def topic(self) -> str
\end{lstlisting}

\textbf{可用性与安全性}

\textbf{\passthrough{\lstinline!AVAILABLE!} /
\passthrough{\lstinline!DDS\_LIFECYCLE!}}：当前绑定源码已实现。该操作可能创建、使用或释放
DDS 资源。

\textbf{参数}

该属性只读，没有 setter 参数。

\textbf{返回值}

\begin{longtable}[]{@{}
  >{\raggedright\arraybackslash}p{(\columnwidth - 4\tabcolsep) * \real{0.18}}
  >{\raggedright\arraybackslash}p{(\columnwidth - 4\tabcolsep) * \real{0.20}}
  >{\raggedright\arraybackslash}p{(\columnwidth - 4\tabcolsep) * \real{0.62}}@{}}
\toprule
位置 & 类型 & 含义 \\
\midrule
\endhead
getter 返回值 & \passthrough{\lstinline!str!} & 返回属性当前值。 \\
\bottomrule
\end{longtable}

\textbf{对应 C++}

\begin{itemize}
\tightlist
\item
  类：\passthrough{\lstinline!unitree\_sdk2\_cpp.channel.ChannelPublisher!}
\item
  字段类型：\passthrough{\lstinline!ChannelPublisher::topic() const!}
\item
  声明位置：由 pybind11 手工绑定提供；manifest 未记录独立头文件行号。
\end{itemize}

\textbf{用法}

\begin{lstlisting}[language=Python]
current_value = obj.topic
\end{lstlisting}

\hypertarget{unitree-sdk2-cpp-channel-channelpublisher-message-type-name}{%
\paragraph{\texorpdfstring{\texttt{unitree\_sdk2\_cpp.channel.ChannelPublisher.message\_type\_name}}{unitree\_sdk2\_cpp.channel.ChannelPublisher.message\_type\_name}}\label{unitree-sdk2-cpp-channel-channelpublisher-message-type-name}}

Publisher 使用的已注册 C++ 消息类型名，只读。

\textbf{签名}

\begin{lstlisting}[language=Python]
@property
def message_type_name(self) -> str
\end{lstlisting}

\textbf{可用性与安全性}

\textbf{\passthrough{\lstinline!AVAILABLE!} /
\passthrough{\lstinline!DDS\_LIFECYCLE!}}：当前绑定源码已实现。该操作可能创建、使用或释放
DDS 资源。

\textbf{参数}

该属性只读，没有 setter 参数。

\textbf{返回值}

\begin{longtable}[]{@{}
  >{\raggedright\arraybackslash}p{(\columnwidth - 4\tabcolsep) * \real{0.18}}
  >{\raggedright\arraybackslash}p{(\columnwidth - 4\tabcolsep) * \real{0.20}}
  >{\raggedright\arraybackslash}p{(\columnwidth - 4\tabcolsep) * \real{0.62}}@{}}
\toprule
位置 & 类型 & 含义 \\
\midrule
\endhead
getter 返回值 & \passthrough{\lstinline!str!} & 返回属性当前值。 \\
\bottomrule
\end{longtable}

\textbf{对应 C++}

\begin{itemize}
\tightlist
\item
  类：\passthrough{\lstinline!unitree\_sdk2\_cpp.channel.ChannelPublisher!}
\item
  字段类型：\passthrough{\lstinline!ChannelPublisher::message\_type\_name() const!}
\item
  声明位置：由 pybind11 手工绑定提供；manifest 未记录独立头文件行号。
\end{itemize}

\textbf{用法}

\begin{lstlisting}[language=Python]
current_value = obj.message_type_name
\end{lstlisting}

\hypertarget{unitree-sdk2-cpp-channel-channelsubscriber}{%
\subsubsection{\texorpdfstring{\texttt{unitree\_sdk2\_cpp.channel.ChannelSubscriber}}{unitree\_sdk2\_cpp.channel.ChannelSubscriber}}\label{unitree-sdk2-cpp-channel-channelsubscriber}}

typed DDS 通道包装类。必须遵守初始化、启动、关闭和全局释放顺序。

\textbf{导入}

\begin{lstlisting}[language=Python]
from unitree_sdk2_cpp.channel import ChannelSubscriber
\end{lstlisting}

\textbf{基类}：\passthrough{\lstinline!Generic[MessageT]!}

\textbf{方法索引}

\begin{longtable}[]{@{}
  >{\raggedright\arraybackslash}p{(\columnwidth - 6\tabcolsep) * \real{0.18}}
  >{\raggedright\arraybackslash}p{(\columnwidth - 6\tabcolsep) * \real{0.24}}
  >{\raggedright\arraybackslash}p{(\columnwidth - 6\tabcolsep) * \real{0.18}}
  >{\raggedright\arraybackslash}p{(\columnwidth - 6\tabcolsep) * \real{0.40}}@{}}
\toprule
方法 & Python 签名 & 状态 & 安全分类 \\
\midrule
\endhead
\protect\hyperlink{unitree-sdk2-cpp-channel-channelsubscriber-dunder-init-1}{\passthrough{\lstinline!\_\_init\_\_!}}
&
\passthrough{\lstinline!def \_\_init\_\_(self, topic: str, message\_type: type[MessageT], callback: Callable[[MessageT], None], queue\_length: int = 0) -> None!}
& \passthrough{\lstinline!AVAILABLE!} &
\passthrough{\lstinline!DDS\_LIFECYCLE!} \\
\protect\hyperlink{unitree-sdk2-cpp-channel-channelsubscriber-init-channel-1}{\passthrough{\lstinline!init\_channel!}}
& \passthrough{\lstinline!def init\_channel(self) -> None!} &
\passthrough{\lstinline!AVAILABLE!} &
\passthrough{\lstinline!DDS\_LIFECYCLE!} \\
\protect\hyperlink{unitree-sdk2-cpp-channel-channelsubscriber-close-channel-1}{\passthrough{\lstinline!close\_channel!}}
& \passthrough{\lstinline!def close\_channel(self) -> None!} &
\passthrough{\lstinline!AVAILABLE!} &
\passthrough{\lstinline!DDS\_LIFECYCLE!} \\
\bottomrule
\end{longtable}

\textbf{属性索引}

\begin{longtable}[]{@{}
  >{\raggedright\arraybackslash}p{(\columnwidth - 6\tabcolsep) * \real{0.18}}
  >{\raggedright\arraybackslash}p{(\columnwidth - 6\tabcolsep) * \real{0.24}}
  >{\raggedright\arraybackslash}p{(\columnwidth - 6\tabcolsep) * \real{0.18}}
  >{\raggedright\arraybackslash}p{(\columnwidth - 6\tabcolsep) * \real{0.40}}@{}}
\toprule
属性 & 读取类型 & 写入类型 & 状态 \\
\midrule
\endhead
\protect\hyperlink{unitree-sdk2-cpp-channel-channelsubscriber-topic}{\passthrough{\lstinline!topic!}}
& \passthrough{\lstinline!str!} & \passthrough{\lstinline!只读!} &
\passthrough{\lstinline!AVAILABLE!} \\
\protect\hyperlink{unitree-sdk2-cpp-channel-channelsubscriber-message-type-name}{\passthrough{\lstinline!message\_type\_name!}}
& \passthrough{\lstinline!str!} & \passthrough{\lstinline!只读!} &
\passthrough{\lstinline!AVAILABLE!} \\
\protect\hyperlink{unitree-sdk2-cpp-channel-channelsubscriber-last-data-available-time}{\passthrough{\lstinline!last\_data\_available\_time!}}
& \passthrough{\lstinline!int!} & \passthrough{\lstinline!只读!} &
\passthrough{\lstinline!AVAILABLE!} \\
\bottomrule
\end{longtable}

\hypertarget{unitree-sdk2-cpp-channel-channelsubscriber-dunder-init-1}{%
\paragraph{\texorpdfstring{\texttt{unitree\_sdk2\_cpp.channel.ChannelSubscriber.\_\_init\_\_}}{unitree\_sdk2\_cpp.channel.ChannelSubscriber.\_\_init\_\_}}\label{unitree-sdk2-cpp-channel-channelsubscriber-dunder-init-1}}

初始化 \passthrough{\lstinline!ChannelSubscriber!}
实例。是否能实际构造取决于下方可用性状态。

\textbf{签名}

\begin{lstlisting}[language=Python]
def __init__(self, topic: str, message_type: type[MessageT], callback: Callable[[MessageT], None], queue_length: int = 0) -> None
\end{lstlisting}

\textbf{可用性与安全性}

\textbf{\passthrough{\lstinline!AVAILABLE!} /
\passthrough{\lstinline!DDS\_LIFECYCLE!}}：当前绑定源码已实现。该操作可能创建、使用或释放
DDS 资源。

\textbf{参数}

\begin{longtable}[]{@{}
  >{\raggedright\arraybackslash}p{(\columnwidth - 6\tabcolsep) * \real{0.18}}
  >{\raggedright\arraybackslash}p{(\columnwidth - 6\tabcolsep) * \real{0.24}}
  >{\raggedright\arraybackslash}p{(\columnwidth - 6\tabcolsep) * \real{0.18}}
  >{\raggedright\arraybackslash}p{(\columnwidth - 6\tabcolsep) * \real{0.40}}@{}}
\toprule
名称 & Python 类型 & 默认值 & 含义 \\
\midrule
\endhead
\passthrough{\lstinline!topic!} & \passthrough{\lstinline!str!} & 必填 &
DDS topic 名称。发布端和订阅端必须同时使用兼容的 topic、消息类型和
QoS。 \\
\passthrough{\lstinline!message\_type!} &
\passthrough{\lstinline!type[MessageT]!} & 必填 & IDL 消息类本身，例如
\passthrough{\lstinline!LowState!}，不是
\passthrough{\lstinline!LowState()!}
实例。该类型必须存在于运行时注册表中。 \\
\passthrough{\lstinline!callback!} &
\passthrough{\lstinline!Callable[[MessageT], None]!} & 必填 &
收到数据或状态变化时调用的 Python
回调。回调签名必须与类型提示一致，并应快速返回。 \\
\passthrough{\lstinline!queue\_length!} & \passthrough{\lstinline!int!}
& \passthrough{\lstinline!0!} &
订阅队列长度。\passthrough{\lstinline!0!}
使用底层默认行为；更大值会允许更多积压，但也可能增加延迟和内存占用。 \\
\bottomrule
\end{longtable}

\textbf{返回值}

\begin{longtable}[]{@{}
  >{\raggedright\arraybackslash}p{(\columnwidth - 4\tabcolsep) * \real{0.18}}
  >{\raggedright\arraybackslash}p{(\columnwidth - 4\tabcolsep) * \real{0.20}}
  >{\raggedright\arraybackslash}p{(\columnwidth - 4\tabcolsep) * \real{0.62}}@{}}
\toprule
位置 & 类型 & 含义 \\
\midrule
\endhead
无 & \passthrough{\lstinline!None!} &
\passthrough{\lstinline!\_\_init\_\_!}
本身不返回值；调用类对象时，在构造函数可用的前提下得到该类实例。 \\
\bottomrule
\end{longtable}

\textbf{对应 C++}

\begin{itemize}
\tightlist
\item
  类：\passthrough{\lstinline!unitree\_sdk2\_cpp.channel.ChannelSubscriber!}
\item
  签名：\passthrough{\lstinline!ChannelSubscriber(const std::string \&, type, callback, int64\_t)!}
\item
  绑定策略：\passthrough{\lstinline!未单独标注!}
\item
  声明位置：由 pybind11 手工绑定或生成注册表提供；manifest
  未记录独立头文件行号。
\end{itemize}

\textbf{说明}

\begin{itemize}
\tightlist
\item
  示例是局部调用片段；\passthrough{\lstinline!obj!}
  和参数变量表示已经按上层业务校验并准备好的对象。
\item
  回调可能由 SDK
  工作线程触发；回调应快速返回、捕获异常，并通过线程安全队列移交耗时工作。
\end{itemize}

\textbf{用法}

\begin{lstlisting}[language=Python]
value = ChannelSubscriber(topic=topic, message_type=message_type, callback=callback, queue_length=queue_length)
\end{lstlisting}

\hypertarget{unitree-sdk2-cpp-channel-channelsubscriber-init-channel-1}{%
\paragraph{\texorpdfstring{\texttt{unitree\_sdk2\_cpp.channel.ChannelSubscriber.init\_channel}}{unitree\_sdk2\_cpp.channel.ChannelSubscriber.init\_channel}}\label{unitree-sdk2-cpp-channel-channelsubscriber-init-channel-1}}

创建并启动 Subscriber 的底层 DDS 读通道。

\textbf{签名}

\begin{lstlisting}[language=Python]
def init_channel(self) -> None
\end{lstlisting}

\textbf{可用性与安全性}

\textbf{\passthrough{\lstinline!AVAILABLE!} /
\passthrough{\lstinline!DDS\_LIFECYCLE!}}：当前绑定源码已实现。该操作可能创建、使用或释放
DDS 资源。

\textbf{参数}

无显式参数。实例方法中的 \passthrough{\lstinline!self!} 由 Python
自动传入。

\textbf{返回值}

\begin{longtable}[]{@{}
  >{\raggedright\arraybackslash}p{(\columnwidth - 4\tabcolsep) * \real{0.18}}
  >{\raggedright\arraybackslash}p{(\columnwidth - 4\tabcolsep) * \real{0.20}}
  >{\raggedright\arraybackslash}p{(\columnwidth - 4\tabcolsep) * \real{0.62}}@{}}
\toprule
位置 & 类型 & 含义 \\
\midrule
\endhead
无 & \passthrough{\lstinline!None!} & 该方法不返回 Python
值；失败可能表现为异常或底层状态变化。 \\
\bottomrule
\end{longtable}

\textbf{对应 C++}

\begin{itemize}
\tightlist
\item
  类：\passthrough{\lstinline!unitree\_sdk2\_cpp.channel.ChannelSubscriber!}
\item
  签名：\passthrough{\lstinline!ChannelSubscriber::init\_channel()!}
\item
  绑定策略：\passthrough{\lstinline!未单独标注!}
\item
  声明位置：由 pybind11 手工绑定或生成注册表提供；manifest
  未记录独立头文件行号。
\end{itemize}

\textbf{说明}

\begin{itemize}
\tightlist
\item
  示例是局部调用片段；\passthrough{\lstinline!obj!}
  和参数变量表示已经按上层业务校验并准备好的对象。
\end{itemize}

\textbf{用法}

\begin{lstlisting}[language=Python]
obj.init_channel()
\end{lstlisting}

\hypertarget{unitree-sdk2-cpp-channel-channelsubscriber-close-channel-1}{%
\paragraph{\texorpdfstring{\texttt{unitree\_sdk2\_cpp.channel.ChannelSubscriber.close\_channel}}{unitree\_sdk2\_cpp.channel.ChannelSubscriber.close\_channel}}\label{unitree-sdk2-cpp-channel-channelsubscriber-close-channel-1}}

关闭 Subscriber 的底层 DDS 读通道，并停止后续回调。

\textbf{签名}

\begin{lstlisting}[language=Python]
def close_channel(self) -> None
\end{lstlisting}

\textbf{可用性与安全性}

\textbf{\passthrough{\lstinline!AVAILABLE!} /
\passthrough{\lstinline!DDS\_LIFECYCLE!}}：当前绑定源码已实现。该操作可能创建、使用或释放
DDS 资源。

\textbf{参数}

无显式参数。实例方法中的 \passthrough{\lstinline!self!} 由 Python
自动传入。

\textbf{返回值}

\begin{longtable}[]{@{}
  >{\raggedright\arraybackslash}p{(\columnwidth - 4\tabcolsep) * \real{0.18}}
  >{\raggedright\arraybackslash}p{(\columnwidth - 4\tabcolsep) * \real{0.20}}
  >{\raggedright\arraybackslash}p{(\columnwidth - 4\tabcolsep) * \real{0.62}}@{}}
\toprule
位置 & 类型 & 含义 \\
\midrule
\endhead
无 & \passthrough{\lstinline!None!} & 该方法不返回 Python
值；失败可能表现为异常或底层状态变化。 \\
\bottomrule
\end{longtable}

\textbf{对应 C++}

\begin{itemize}
\tightlist
\item
  类：\passthrough{\lstinline!unitree\_sdk2\_cpp.channel.ChannelSubscriber!}
\item
  签名：\passthrough{\lstinline!ChannelSubscriber::close\_channel()!}
\item
  绑定策略：\passthrough{\lstinline!未单独标注!}
\item
  声明位置：由 pybind11 手工绑定或生成注册表提供；manifest
  未记录独立头文件行号。
\end{itemize}

\textbf{说明}

\begin{itemize}
\tightlist
\item
  示例是局部调用片段；\passthrough{\lstinline!obj!}
  和参数变量表示已经按上层业务校验并准备好的对象。
\end{itemize}

\textbf{用法}

\begin{lstlisting}[language=Python]
obj.close_channel()
\end{lstlisting}

\hypertarget{unitree-sdk2-cpp-channel-channelsubscriber-topic}{%
\paragraph{\texorpdfstring{\texttt{unitree\_sdk2\_cpp.channel.ChannelSubscriber.topic}}{unitree\_sdk2\_cpp.channel.ChannelSubscriber.topic}}\label{unitree-sdk2-cpp-channel-channelsubscriber-topic}}

Subscriber 构造时保存的 DDS topic 名称，只读。

\textbf{签名}

\begin{lstlisting}[language=Python]
@property
def topic(self) -> str
\end{lstlisting}

\textbf{可用性与安全性}

\textbf{\passthrough{\lstinline!AVAILABLE!} /
\passthrough{\lstinline!DDS\_LIFECYCLE!}}：当前绑定源码已实现。该操作可能创建、使用或释放
DDS 资源。

\textbf{参数}

该属性只读，没有 setter 参数。

\textbf{返回值}

\begin{longtable}[]{@{}
  >{\raggedright\arraybackslash}p{(\columnwidth - 4\tabcolsep) * \real{0.18}}
  >{\raggedright\arraybackslash}p{(\columnwidth - 4\tabcolsep) * \real{0.20}}
  >{\raggedright\arraybackslash}p{(\columnwidth - 4\tabcolsep) * \real{0.62}}@{}}
\toprule
位置 & 类型 & 含义 \\
\midrule
\endhead
getter 返回值 & \passthrough{\lstinline!str!} & 返回属性当前值。 \\
\bottomrule
\end{longtable}

\textbf{对应 C++}

\begin{itemize}
\tightlist
\item
  类：\passthrough{\lstinline!unitree\_sdk2\_cpp.channel.ChannelSubscriber!}
\item
  字段类型：\passthrough{\lstinline!ChannelSubscriber::topic() const!}
\item
  声明位置：由 pybind11 手工绑定提供；manifest 未记录独立头文件行号。
\end{itemize}

\textbf{用法}

\begin{lstlisting}[language=Python]
current_value = obj.topic
\end{lstlisting}

\hypertarget{unitree-sdk2-cpp-channel-channelsubscriber-message-type-name}{%
\paragraph{\texorpdfstring{\texttt{unitree\_sdk2\_cpp.channel.ChannelSubscriber.message\_type\_name}}{unitree\_sdk2\_cpp.channel.ChannelSubscriber.message\_type\_name}}\label{unitree-sdk2-cpp-channel-channelsubscriber-message-type-name}}

Subscriber 使用的已注册 C++ 消息类型名，只读。

\textbf{签名}

\begin{lstlisting}[language=Python]
@property
def message_type_name(self) -> str
\end{lstlisting}

\textbf{可用性与安全性}

\textbf{\passthrough{\lstinline!AVAILABLE!} /
\passthrough{\lstinline!DDS\_LIFECYCLE!}}：当前绑定源码已实现。该操作可能创建、使用或释放
DDS 资源。

\textbf{参数}

该属性只读，没有 setter 参数。

\textbf{返回值}

\begin{longtable}[]{@{}
  >{\raggedright\arraybackslash}p{(\columnwidth - 4\tabcolsep) * \real{0.18}}
  >{\raggedright\arraybackslash}p{(\columnwidth - 4\tabcolsep) * \real{0.20}}
  >{\raggedright\arraybackslash}p{(\columnwidth - 4\tabcolsep) * \real{0.62}}@{}}
\toprule
位置 & 类型 & 含义 \\
\midrule
\endhead
getter 返回值 & \passthrough{\lstinline!str!} & 返回属性当前值。 \\
\bottomrule
\end{longtable}

\textbf{对应 C++}

\begin{itemize}
\tightlist
\item
  类：\passthrough{\lstinline!unitree\_sdk2\_cpp.channel.ChannelSubscriber!}
\item
  字段类型：\passthrough{\lstinline!ChannelSubscriber::message\_type\_name() const!}
\item
  声明位置：由 pybind11 手工绑定提供；manifest 未记录独立头文件行号。
\end{itemize}

\textbf{用法}

\begin{lstlisting}[language=Python]
current_value = obj.message_type_name
\end{lstlisting}

\hypertarget{unitree-sdk2-cpp-channel-channelsubscriber-last-data-available-time}{%
\paragraph{\texorpdfstring{\texttt{unitree\_sdk2\_cpp.channel.ChannelSubscriber.last\_data\_available\_time}}{unitree\_sdk2\_cpp.channel.ChannelSubscriber.last\_data\_available\_time}}\label{unitree-sdk2-cpp-channel-channelsubscriber-last-data-available-time}}

底层记录的最近一次数据可用时间。其单位和时间基准以 SDK 实现为准。

\textbf{签名}

\begin{lstlisting}[language=Python]
@property
def last_data_available_time(self) -> int
\end{lstlisting}

\textbf{可用性与安全性}

\textbf{\passthrough{\lstinline!AVAILABLE!} /
\passthrough{\lstinline!DDS\_LIFECYCLE!}}：当前绑定源码已实现。该操作可能创建、使用或释放
DDS 资源。

\textbf{参数}

该属性只读，没有 setter 参数。

\textbf{返回值}

\begin{longtable}[]{@{}
  >{\raggedright\arraybackslash}p{(\columnwidth - 4\tabcolsep) * \real{0.18}}
  >{\raggedright\arraybackslash}p{(\columnwidth - 4\tabcolsep) * \real{0.20}}
  >{\raggedright\arraybackslash}p{(\columnwidth - 4\tabcolsep) * \real{0.62}}@{}}
\toprule
位置 & 类型 & 含义 \\
\midrule
\endhead
getter 返回值 & \passthrough{\lstinline!int!} & 返回属性当前值。 \\
\bottomrule
\end{longtable}

\textbf{对应 C++}

\begin{itemize}
\tightlist
\item
  类：\passthrough{\lstinline!unitree\_sdk2\_cpp.channel.ChannelSubscriber!}
\item
  字段类型：\passthrough{\lstinline!ChannelSubscriber::last\_data\_available\_time() const!}
\item
  声明位置：由 pybind11 手工绑定提供；manifest 未记录独立头文件行号。
\end{itemize}

\textbf{用法}

\begin{lstlisting}[language=Python]
current_value = obj.last_data_available_time
\end{lstlisting}

\begin{center}\rule{0.5\linewidth}{0.5pt}\end{center}

\hypertarget{unitree-sdk2-cpp-idl-go2}{%
\subsection{Go2 IDL 消息}\label{unitree-sdk2-cpp-idl-go2}}

模块：\passthrough{\lstinline!unitree\_sdk2\_cpp.idl.go2!}

\hypertarget{ux7c7bux7d22ux5f15-2}{%
\subsubsection{类索引}\label{ux7c7bux7d22ux5f15-2}}

\begin{longtable}[]{@{}
  >{\raggedright\arraybackslash}p{(\columnwidth - 4\tabcolsep) * \real{0.18}}
  >{\raggedleft\arraybackslash}p{(\columnwidth - 4\tabcolsep) * \real{0.20}}
  >{\raggedleft\arraybackslash}p{(\columnwidth - 4\tabcolsep) * \real{0.62}}@{}}
\toprule
类 & 公开函数签名 & 属性 \\
\midrule
\endhead
\protect\hyperlink{unitree-sdk2-cpp-idl-go2-audiodata}{\passthrough{\lstinline!AudioData!}}
& 3 & 2 \\
\protect\hyperlink{unitree-sdk2-cpp-idl-go2-bmscmd}{\passthrough{\lstinline!BmsCmd!}}
& 3 & 2 \\
\protect\hyperlink{unitree-sdk2-cpp-idl-go2-bmsstate}{\passthrough{\lstinline!BmsState!}}
& 3 & 9 \\
\protect\hyperlink{unitree-sdk2-cpp-idl-go2-configchangestatus}{\passthrough{\lstinline!ConfigChangeStatus!}}
& 3 & 2 \\
\protect\hyperlink{unitree-sdk2-cpp-idl-go2-error}{\passthrough{\lstinline!Error!}}
& 3 & 2 \\
\protect\hyperlink{unitree-sdk2-cpp-idl-go2-go2frontvideodata}{\passthrough{\lstinline!Go2FrontVideoData!}}
& 3 & 4 \\
\protect\hyperlink{unitree-sdk2-cpp-idl-go2-heightmap}{\passthrough{\lstinline!HeightMap!}}
& 3 & 7 \\
\protect\hyperlink{unitree-sdk2-cpp-idl-go2-imustate}{\passthrough{\lstinline!IMUState!}}
& 3 & 5 \\
\protect\hyperlink{unitree-sdk2-cpp-idl-go2-interfaceconfig}{\passthrough{\lstinline!InterfaceConfig!}}
& 3 & 3 \\
\protect\hyperlink{unitree-sdk2-cpp-idl-go2-lidarstate}{\passthrough{\lstinline!LidarState!}}
& 3 & 18 \\
\protect\hyperlink{unitree-sdk2-cpp-idl-go2-motorcmd}{\passthrough{\lstinline!MotorCmd!}}
& 3 & 7 \\
\protect\hyperlink{unitree-sdk2-cpp-idl-go2-lowcmd}{\passthrough{\lstinline!LowCmd!}}
& 3 & 14 \\
\protect\hyperlink{unitree-sdk2-cpp-idl-go2-motorstate}{\passthrough{\lstinline!MotorState!}}
& 3 & 11 \\
\protect\hyperlink{unitree-sdk2-cpp-idl-go2-lowstate}{\passthrough{\lstinline!LowState!}}
& 3 & 22 \\
\protect\hyperlink{unitree-sdk2-cpp-idl-go2-motorcmds}{\passthrough{\lstinline!MotorCmds!}}
& 3 & 1 \\
\protect\hyperlink{unitree-sdk2-cpp-idl-go2-motorstates}{\passthrough{\lstinline!MotorStates!}}
& 3 & 1 \\
\protect\hyperlink{unitree-sdk2-cpp-idl-go2-pathpoint}{\passthrough{\lstinline!PathPoint!}}
& 3 & 7 \\
\protect\hyperlink{unitree-sdk2-cpp-idl-go2-req}{\passthrough{\lstinline!Req!}}
& 3 & 2 \\
\protect\hyperlink{unitree-sdk2-cpp-idl-go2-res}{\passthrough{\lstinline!Res!}}
& 3 & 3 \\
\protect\hyperlink{unitree-sdk2-cpp-idl-go2-sportmodecmd}{\passthrough{\lstinline!SportModeCmd!}}
& 3 & 11 \\
\protect\hyperlink{unitree-sdk2-cpp-idl-go2-timespec}{\passthrough{\lstinline!TimeSpec!}}
& 3 & 2 \\
\protect\hyperlink{unitree-sdk2-cpp-idl-go2-sportmodestate}{\passthrough{\lstinline!SportModeState!}}
& 3 & 16 \\
\protect\hyperlink{unitree-sdk2-cpp-idl-go2-uwbstate}{\passthrough{\lstinline!UwbState!}}
& 3 & 17 \\
\protect\hyperlink{unitree-sdk2-cpp-idl-go2-uwbswitch}{\passthrough{\lstinline!UwbSwitch!}}
& 3 & 1 \\
\protect\hyperlink{unitree-sdk2-cpp-idl-go2-voxelmapcompressed}{\passthrough{\lstinline!VoxelMapCompressed!}}
& 3 & 7 \\
\protect\hyperlink{unitree-sdk2-cpp-idl-go2-wirelesscontroller}{\passthrough{\lstinline!WirelessController!}}
& 3 & 5 \\
\bottomrule
\end{longtable}

\hypertarget{unitree-sdk2-cpp-idl-go2-audiodata}{%
\subsubsection{\texorpdfstring{\texttt{unitree\_sdk2\_cpp.idl.go2.AudioData}}{unitree\_sdk2\_cpp.idl.go2.AudioData}}\label{unitree-sdk2-cpp-idl-go2-audiodata}}

可由 Python 读写的 DDS/IDL 消息值类型。字段采用复制语义。

\textbf{导入}

\begin{lstlisting}[language=Python]
from unitree_sdk2_cpp.idl.go2 import AudioData
\end{lstlisting}

\textbf{方法索引}

\begin{longtable}[]{@{}
  >{\raggedright\arraybackslash}p{(\columnwidth - 6\tabcolsep) * \real{0.18}}
  >{\raggedright\arraybackslash}p{(\columnwidth - 6\tabcolsep) * \real{0.24}}
  >{\raggedright\arraybackslash}p{(\columnwidth - 6\tabcolsep) * \real{0.18}}
  >{\raggedright\arraybackslash}p{(\columnwidth - 6\tabcolsep) * \real{0.40}}@{}}
\toprule
方法 & Python 签名 & 状态 & 安全分类 \\
\midrule
\endhead
\protect\hyperlink{unitree-sdk2-cpp-idl-go2-audiodata-dunder-init-1}{\passthrough{\lstinline!\_\_init\_\_!}}
& \passthrough{\lstinline!def \_\_init\_\_(self) -> None!} &
\passthrough{\lstinline!AVAILABLE!} &
\passthrough{\lstinline!VALUE\_TYPE!} \\
\protect\hyperlink{unitree-sdk2-cpp-idl-go2-audiodata-dunder-eq-1}{\passthrough{\lstinline!\_\_eq\_\_!}}
& \passthrough{\lstinline!def \_\_eq\_\_(self, other: object) -> bool!}
& \passthrough{\lstinline!AVAILABLE!} &
\passthrough{\lstinline!VALUE\_TYPE!} \\
\protect\hyperlink{unitree-sdk2-cpp-idl-go2-audiodata-dunder-ne-1}{\passthrough{\lstinline!\_\_ne\_\_!}}
& \passthrough{\lstinline!def \_\_ne\_\_(self, other: object) -> bool!}
& \passthrough{\lstinline!AVAILABLE!} &
\passthrough{\lstinline!VALUE\_TYPE!} \\
\bottomrule
\end{longtable}

\textbf{属性索引}

\begin{longtable}[]{@{}
  >{\raggedright\arraybackslash}p{(\columnwidth - 6\tabcolsep) * \real{0.18}}
  >{\raggedright\arraybackslash}p{(\columnwidth - 6\tabcolsep) * \real{0.24}}
  >{\raggedright\arraybackslash}p{(\columnwidth - 6\tabcolsep) * \real{0.18}}
  >{\raggedright\arraybackslash}p{(\columnwidth - 6\tabcolsep) * \real{0.40}}@{}}
\toprule
属性 & 读取类型 & 写入类型 & 状态 \\
\midrule
\endhead
\protect\hyperlink{unitree-sdk2-cpp-idl-go2-audiodata-time-frame}{\passthrough{\lstinline!time\_frame!}}
& \passthrough{\lstinline!int!} & \passthrough{\lstinline!int!} &
\passthrough{\lstinline!AVAILABLE!} \\
\protect\hyperlink{unitree-sdk2-cpp-idl-go2-audiodata-data}{\passthrough{\lstinline!data!}}
& \passthrough{\lstinline!list[int]!} &
\passthrough{\lstinline!Sequence[int]!} &
\passthrough{\lstinline!AVAILABLE!} \\
\bottomrule
\end{longtable}

\hypertarget{unitree-sdk2-cpp-idl-go2-audiodata-dunder-init-1}{%
\paragraph{\texorpdfstring{\texttt{unitree\_sdk2\_cpp.idl.go2.AudioData.\_\_init\_\_}}{unitree\_sdk2\_cpp.idl.go2.AudioData.\_\_init\_\_}}\label{unitree-sdk2-cpp-idl-go2-audiodata-dunder-init-1}}

初始化 \passthrough{\lstinline!AudioData!}
实例。是否能实际构造取决于下方可用性状态。

\textbf{签名}

\begin{lstlisting}[language=Python]
def __init__(self) -> None
\end{lstlisting}

\textbf{可用性与安全性}

\textbf{\passthrough{\lstinline!AVAILABLE!} /
\passthrough{\lstinline!VALUE\_TYPE!}}：当前绑定源码已实现该消息值操作。

\textbf{参数}

无显式参数。实例方法中的 \passthrough{\lstinline!self!} 由 Python
自动传入。

\textbf{返回值}

\begin{longtable}[]{@{}
  >{\raggedright\arraybackslash}p{(\columnwidth - 4\tabcolsep) * \real{0.18}}
  >{\raggedright\arraybackslash}p{(\columnwidth - 4\tabcolsep) * \real{0.20}}
  >{\raggedright\arraybackslash}p{(\columnwidth - 4\tabcolsep) * \real{0.62}}@{}}
\toprule
位置 & 类型 & 含义 \\
\midrule
\endhead
无 & \passthrough{\lstinline!None!} &
\passthrough{\lstinline!\_\_init\_\_!}
本身不返回值；调用类对象时，在构造函数可用的前提下得到该类实例。 \\
\bottomrule
\end{longtable}

\textbf{对应 C++}

\begin{itemize}
\tightlist
\item
  类：\passthrough{\lstinline!unitree\_go::msg::dds\_::AudioData\_!}
\item
  签名：\passthrough{\lstinline!AudioData\_()!}
\item
  绑定策略：\passthrough{\lstinline!未单独标注!}
\item
  声明位置：\passthrough{\lstinline!include/unitree/idl/go2/AudioData\_.hpp:28!}
\end{itemize}

\textbf{说明}

\begin{itemize}
\tightlist
\item
  示例是局部调用片段；\passthrough{\lstinline!obj!}
  和参数变量表示已经按上层业务校验并准备好的对象。
\end{itemize}

\textbf{用法}

\begin{lstlisting}[language=Python]
value = AudioData()
\end{lstlisting}

\hypertarget{unitree-sdk2-cpp-idl-go2-audiodata-dunder-eq-1}{%
\paragraph{\texorpdfstring{\texttt{unitree\_sdk2\_cpp.idl.go2.AudioData.\_\_eq\_\_}}{unitree\_sdk2\_cpp.idl.go2.AudioData.\_\_eq\_\_}}\label{unitree-sdk2-cpp-idl-go2-audiodata-dunder-eq-1}}

按底层 IDL 消息内容比较两个对象是否相等。

\textbf{签名}

\begin{lstlisting}[language=Python]
def __eq__(self, other: object) -> bool
\end{lstlisting}

\textbf{可用性与安全性}

\textbf{\passthrough{\lstinline!AVAILABLE!} /
\passthrough{\lstinline!VALUE\_TYPE!}}：当前绑定源码已实现该消息值操作。

\textbf{参数}

\begin{longtable}[]{@{}
  >{\raggedright\arraybackslash}p{(\columnwidth - 6\tabcolsep) * \real{0.18}}
  >{\raggedright\arraybackslash}p{(\columnwidth - 6\tabcolsep) * \real{0.24}}
  >{\raggedright\arraybackslash}p{(\columnwidth - 6\tabcolsep) * \real{0.18}}
  >{\raggedright\arraybackslash}p{(\columnwidth - 6\tabcolsep) * \real{0.40}}@{}}
\toprule
名称 & Python 类型 & 默认值 & 含义 \\
\midrule
\endhead
\passthrough{\lstinline!other!} & \passthrough{\lstinline!object!} &
必填 & 用于比较的另一个 Python 对象。类型不兼容时比较结果通常为
\passthrough{\lstinline!False!}。 \\
\bottomrule
\end{longtable}

\textbf{返回值}

\begin{longtable}[]{@{}
  >{\raggedright\arraybackslash}p{(\columnwidth - 4\tabcolsep) * \real{0.18}}
  >{\raggedright\arraybackslash}p{(\columnwidth - 4\tabcolsep) * \real{0.20}}
  >{\raggedright\arraybackslash}p{(\columnwidth - 4\tabcolsep) * \real{0.62}}@{}}
\toprule
位置 & 类型 & 含义 \\
\midrule
\endhead
返回值 & \passthrough{\lstinline!bool!} & 返回
\passthrough{\lstinline!bool!}。更精确的业务含义见该方法用途、C++
签名和目标型号协议。 \\
\bottomrule
\end{longtable}

\textbf{对应 C++}

\begin{itemize}
\tightlist
\item
  类：\passthrough{\lstinline!unitree\_go::msg::dds\_::AudioData\_!}
\item
  签名：\passthrough{\lstinline!operator==(const value \&) const!}
\item
  绑定策略：\passthrough{\lstinline!未单独标注!}
\item
  声明位置：由 pybind11 手工绑定或生成注册表提供；manifest
  未记录独立头文件行号。
\end{itemize}

\textbf{说明}

\begin{itemize}
\tightlist
\item
  示例是局部调用片段；\passthrough{\lstinline!obj!}
  和参数变量表示已经按上层业务校验并准备好的对象。
\end{itemize}

\textbf{用法}

\begin{lstlisting}[language=Python]
same = left == right
\end{lstlisting}

\hypertarget{unitree-sdk2-cpp-idl-go2-audiodata-dunder-ne-1}{%
\paragraph{\texorpdfstring{\texttt{unitree\_sdk2\_cpp.idl.go2.AudioData.\_\_ne\_\_}}{unitree\_sdk2\_cpp.idl.go2.AudioData.\_\_ne\_\_}}\label{unitree-sdk2-cpp-idl-go2-audiodata-dunder-ne-1}}

按底层 IDL 消息内容比较两个对象是否不相等。

\textbf{签名}

\begin{lstlisting}[language=Python]
def __ne__(self, other: object) -> bool
\end{lstlisting}

\textbf{可用性与安全性}

\textbf{\passthrough{\lstinline!AVAILABLE!} /
\passthrough{\lstinline!VALUE\_TYPE!}}：当前绑定源码已实现该消息值操作。

\textbf{参数}

\begin{longtable}[]{@{}
  >{\raggedright\arraybackslash}p{(\columnwidth - 6\tabcolsep) * \real{0.18}}
  >{\raggedright\arraybackslash}p{(\columnwidth - 6\tabcolsep) * \real{0.24}}
  >{\raggedright\arraybackslash}p{(\columnwidth - 6\tabcolsep) * \real{0.18}}
  >{\raggedright\arraybackslash}p{(\columnwidth - 6\tabcolsep) * \real{0.40}}@{}}
\toprule
名称 & Python 类型 & 默认值 & 含义 \\
\midrule
\endhead
\passthrough{\lstinline!other!} & \passthrough{\lstinline!object!} &
必填 & 用于比较的另一个 Python 对象。类型不兼容时比较结果通常为
\passthrough{\lstinline!False!}。 \\
\bottomrule
\end{longtable}

\textbf{返回值}

\begin{longtable}[]{@{}
  >{\raggedright\arraybackslash}p{(\columnwidth - 4\tabcolsep) * \real{0.18}}
  >{\raggedright\arraybackslash}p{(\columnwidth - 4\tabcolsep) * \real{0.20}}
  >{\raggedright\arraybackslash}p{(\columnwidth - 4\tabcolsep) * \real{0.62}}@{}}
\toprule
位置 & 类型 & 含义 \\
\midrule
\endhead
返回值 & \passthrough{\lstinline!bool!} & 返回
\passthrough{\lstinline!bool!}。更精确的业务含义见该方法用途、C++
签名和目标型号协议。 \\
\bottomrule
\end{longtable}

\textbf{对应 C++}

\begin{itemize}
\tightlist
\item
  类：\passthrough{\lstinline!unitree\_go::msg::dds\_::AudioData\_!}
\item
  签名：\passthrough{\lstinline"operator!=(const value \&) const"}
\item
  绑定策略：\passthrough{\lstinline!未单独标注!}
\item
  声明位置：由 pybind11 手工绑定或生成注册表提供；manifest
  未记录独立头文件行号。
\end{itemize}

\textbf{说明}

\begin{itemize}
\tightlist
\item
  示例是局部调用片段；\passthrough{\lstinline!obj!}
  和参数变量表示已经按上层业务校验并准备好的对象。
\end{itemize}

\textbf{用法}

\begin{lstlisting}[language=Python]
different = left != right
\end{lstlisting}

\hypertarget{unitree-sdk2-cpp-idl-go2-audiodata-time-frame}{%
\paragraph{\texorpdfstring{\texttt{unitree\_sdk2\_cpp.idl.go2.AudioData.time\_frame}}{unitree\_sdk2\_cpp.idl.go2.AudioData.time\_frame}}\label{unitree-sdk2-cpp-idl-go2-audiodata-time-frame}}

底层 \passthrough{\lstinline!unitree\_go::msg::dds\_::AudioData\_!} 的
\passthrough{\lstinline!time\_frame!} 字段。类型签名只说明 Python
表示；单位、数值范围、数组固定长度和枚举含义需查对应 IDL 头文件。
取值范围为 0 到 18446744073709551615。

\textbf{签名}

\begin{lstlisting}[language=Python]
@property
def time_frame(self) -> int

@time_frame.setter
def time_frame(self, value: int) -> None
\end{lstlisting}

\textbf{可用性与安全性}

\textbf{\passthrough{\lstinline!AVAILABLE!} /
\passthrough{\lstinline!VALUE\_TYPE!}}：当前绑定源码已实现该消息值操作。

\textbf{参数}

\begin{longtable}[]{@{}
  >{\raggedright\arraybackslash}p{(\columnwidth - 4\tabcolsep) * \real{0.18}}
  >{\raggedright\arraybackslash}p{(\columnwidth - 4\tabcolsep) * \real{0.20}}
  >{\raggedright\arraybackslash}p{(\columnwidth - 4\tabcolsep) * \real{0.62}}@{}}
\toprule
名称 & Python 类型 & 含义 \\
\midrule
\endhead
\passthrough{\lstinline!value!} & \passthrough{\lstinline!int!} &
要写入或传递的新值，必须符合该参数的 Python 类型以及底层 C++
范围约束。 \\
\bottomrule
\end{longtable}

\textbf{返回值}

\begin{longtable}[]{@{}
  >{\raggedright\arraybackslash}p{(\columnwidth - 4\tabcolsep) * \real{0.18}}
  >{\raggedright\arraybackslash}p{(\columnwidth - 4\tabcolsep) * \real{0.20}}
  >{\raggedright\arraybackslash}p{(\columnwidth - 4\tabcolsep) * \real{0.62}}@{}}
\toprule
位置 & 类型 & 含义 \\
\midrule
\endhead
getter 返回值 & \passthrough{\lstinline!int!} & 返回属性当前值。 \\
\bottomrule
\end{longtable}

\textbf{对应 C++}

\begin{itemize}
\tightlist
\item
  类：\passthrough{\lstinline!unitree\_go::msg::dds\_::AudioData\_!}
\item
  字段类型：\passthrough{\lstinline!uint64\_t!}
\item
  声明位置：\passthrough{\lstinline!include/unitree/idl/go2/AudioData\_.hpp:24!}
\end{itemize}

\textbf{用法}

\begin{lstlisting}[language=Python]
current_value = obj.time_frame
obj.time_frame = new_value
\end{lstlisting}

\hypertarget{unitree-sdk2-cpp-idl-go2-audiodata-data}{%
\paragraph{\texorpdfstring{\texttt{unitree\_sdk2\_cpp.idl.go2.AudioData.data}}{unitree\_sdk2\_cpp.idl.go2.AudioData.data}}\label{unitree-sdk2-cpp-idl-go2-audiodata-data}}

底层 \passthrough{\lstinline!unitree\_go::msg::dds\_::AudioData\_!} 的
\passthrough{\lstinline!data!} 字段。类型签名只说明 Python
表示；单位、数值范围、数组固定长度和枚举含义需查对应 IDL 头文件。
底层是可变长度 vector；元素约束：取值范围为 0 到 255。

\textbf{签名}

\begin{lstlisting}[language=Python]
@property
def data(self) -> list[int]

@data.setter
def data(self, value: Sequence[int]) -> None
\end{lstlisting}

\textbf{可用性与安全性}

\textbf{\passthrough{\lstinline!AVAILABLE!} /
\passthrough{\lstinline!VALUE\_TYPE!}}：当前绑定源码已实现该消息值操作。

\textbf{参数}

\begin{longtable}[]{@{}
  >{\raggedright\arraybackslash}p{(\columnwidth - 4\tabcolsep) * \real{0.18}}
  >{\raggedright\arraybackslash}p{(\columnwidth - 4\tabcolsep) * \real{0.20}}
  >{\raggedright\arraybackslash}p{(\columnwidth - 4\tabcolsep) * \real{0.62}}@{}}
\toprule
名称 & Python 类型 & 含义 \\
\midrule
\endhead
\passthrough{\lstinline!value!} &
\passthrough{\lstinline!Sequence[int]!} &
要写入或传递的新值，必须符合该参数的 Python 类型以及底层 C++
范围约束。 \\
\bottomrule
\end{longtable}

\textbf{返回值}

\begin{longtable}[]{@{}
  >{\raggedright\arraybackslash}p{(\columnwidth - 4\tabcolsep) * \real{0.18}}
  >{\raggedright\arraybackslash}p{(\columnwidth - 4\tabcolsep) * \real{0.20}}
  >{\raggedright\arraybackslash}p{(\columnwidth - 4\tabcolsep) * \real{0.62}}@{}}
\toprule
位置 & 类型 & 含义 \\
\midrule
\endhead
getter 返回值 & \passthrough{\lstinline!list[int]!} &
返回属性当前值。数组、vector 和嵌套消息采用复制语义。 \\
\bottomrule
\end{longtable}

\textbf{对应 C++}

\begin{itemize}
\tightlist
\item
  类：\passthrough{\lstinline!unitree\_go::msg::dds\_::AudioData\_!}
\item
  字段类型：\passthrough{\lstinline!std::vector<uint8\_t>!}
\item
  声明位置：\passthrough{\lstinline!include/unitree/idl/go2/AudioData\_.hpp:25!}
\end{itemize}

\textbf{用法}

\begin{lstlisting}[language=Python]
current_value = obj.data
obj.data = new_value
\end{lstlisting}

\begin{quote}
{[}!TIP{]} 读取容器或嵌套值后应遵循``读取、修改、重新赋值''。只修改
getter 返回的副本不会可靠地更新原 C++ 消息。
\end{quote}

\hypertarget{unitree-sdk2-cpp-idl-go2-bmscmd}{%
\subsubsection{\texorpdfstring{\texttt{unitree\_sdk2\_cpp.idl.go2.BmsCmd}}{unitree\_sdk2\_cpp.idl.go2.BmsCmd}}\label{unitree-sdk2-cpp-idl-go2-bmscmd}}

可由 Python 读写的 DDS/IDL 消息值类型。字段采用复制语义。

\textbf{导入}

\begin{lstlisting}[language=Python]
from unitree_sdk2_cpp.idl.go2 import BmsCmd
\end{lstlisting}

\textbf{方法索引}

\begin{longtable}[]{@{}
  >{\raggedright\arraybackslash}p{(\columnwidth - 6\tabcolsep) * \real{0.18}}
  >{\raggedright\arraybackslash}p{(\columnwidth - 6\tabcolsep) * \real{0.24}}
  >{\raggedright\arraybackslash}p{(\columnwidth - 6\tabcolsep) * \real{0.18}}
  >{\raggedright\arraybackslash}p{(\columnwidth - 6\tabcolsep) * \real{0.40}}@{}}
\toprule
方法 & Python 签名 & 状态 & 安全分类 \\
\midrule
\endhead
\protect\hyperlink{unitree-sdk2-cpp-idl-go2-bmscmd-dunder-init-1}{\passthrough{\lstinline!\_\_init\_\_!}}
& \passthrough{\lstinline!def \_\_init\_\_(self) -> None!} &
\passthrough{\lstinline!AVAILABLE!} &
\passthrough{\lstinline!VALUE\_TYPE!} \\
\protect\hyperlink{unitree-sdk2-cpp-idl-go2-bmscmd-dunder-eq-1}{\passthrough{\lstinline!\_\_eq\_\_!}}
& \passthrough{\lstinline!def \_\_eq\_\_(self, other: object) -> bool!}
& \passthrough{\lstinline!AVAILABLE!} &
\passthrough{\lstinline!VALUE\_TYPE!} \\
\protect\hyperlink{unitree-sdk2-cpp-idl-go2-bmscmd-dunder-ne-1}{\passthrough{\lstinline!\_\_ne\_\_!}}
& \passthrough{\lstinline!def \_\_ne\_\_(self, other: object) -> bool!}
& \passthrough{\lstinline!AVAILABLE!} &
\passthrough{\lstinline!VALUE\_TYPE!} \\
\bottomrule
\end{longtable}

\textbf{属性索引}

\begin{longtable}[]{@{}
  >{\raggedright\arraybackslash}p{(\columnwidth - 6\tabcolsep) * \real{0.18}}
  >{\raggedright\arraybackslash}p{(\columnwidth - 6\tabcolsep) * \real{0.24}}
  >{\raggedright\arraybackslash}p{(\columnwidth - 6\tabcolsep) * \real{0.18}}
  >{\raggedright\arraybackslash}p{(\columnwidth - 6\tabcolsep) * \real{0.40}}@{}}
\toprule
属性 & 读取类型 & 写入类型 & 状态 \\
\midrule
\endhead
\protect\hyperlink{unitree-sdk2-cpp-idl-go2-bmscmd-off}{\passthrough{\lstinline!off!}}
& \passthrough{\lstinline!int!} & \passthrough{\lstinline!int!} &
\passthrough{\lstinline!AVAILABLE!} \\
\protect\hyperlink{unitree-sdk2-cpp-idl-go2-bmscmd-reserve}{\passthrough{\lstinline!reserve!}}
& \passthrough{\lstinline!list[int]!} &
\passthrough{\lstinline!Sequence[int]!} &
\passthrough{\lstinline!AVAILABLE!} \\
\bottomrule
\end{longtable}

\hypertarget{unitree-sdk2-cpp-idl-go2-bmscmd-dunder-init-1}{%
\paragraph{\texorpdfstring{\texttt{unitree\_sdk2\_cpp.idl.go2.BmsCmd.\_\_init\_\_}}{unitree\_sdk2\_cpp.idl.go2.BmsCmd.\_\_init\_\_}}\label{unitree-sdk2-cpp-idl-go2-bmscmd-dunder-init-1}}

初始化 \passthrough{\lstinline!BmsCmd!}
实例。是否能实际构造取决于下方可用性状态。

\textbf{签名}

\begin{lstlisting}[language=Python]
def __init__(self) -> None
\end{lstlisting}

\textbf{可用性与安全性}

\textbf{\passthrough{\lstinline!AVAILABLE!} /
\passthrough{\lstinline!VALUE\_TYPE!}}：当前绑定源码已实现该消息值操作。

\textbf{参数}

无显式参数。实例方法中的 \passthrough{\lstinline!self!} 由 Python
自动传入。

\textbf{返回值}

\begin{longtable}[]{@{}
  >{\raggedright\arraybackslash}p{(\columnwidth - 4\tabcolsep) * \real{0.18}}
  >{\raggedright\arraybackslash}p{(\columnwidth - 4\tabcolsep) * \real{0.20}}
  >{\raggedright\arraybackslash}p{(\columnwidth - 4\tabcolsep) * \real{0.62}}@{}}
\toprule
位置 & 类型 & 含义 \\
\midrule
\endhead
无 & \passthrough{\lstinline!None!} &
\passthrough{\lstinline!\_\_init\_\_!}
本身不返回值；调用类对象时，在构造函数可用的前提下得到该类实例。 \\
\bottomrule
\end{longtable}

\textbf{对应 C++}

\begin{itemize}
\tightlist
\item
  类：\passthrough{\lstinline!unitree\_go::msg::dds\_::BmsCmd\_!}
\item
  签名：\passthrough{\lstinline!BmsCmd\_()!}
\item
  绑定策略：\passthrough{\lstinline!未单独标注!}
\item
  声明位置：\passthrough{\lstinline!include/unitree/idl/go2/BmsCmd\_.hpp:28!}
\end{itemize}

\textbf{说明}

\begin{itemize}
\tightlist
\item
  示例是局部调用片段；\passthrough{\lstinline!obj!}
  和参数变量表示已经按上层业务校验并准备好的对象。
\end{itemize}

\textbf{用法}

\begin{lstlisting}[language=Python]
value = BmsCmd()
\end{lstlisting}

\hypertarget{unitree-sdk2-cpp-idl-go2-bmscmd-dunder-eq-1}{%
\paragraph{\texorpdfstring{\texttt{unitree\_sdk2\_cpp.idl.go2.BmsCmd.\_\_eq\_\_}}{unitree\_sdk2\_cpp.idl.go2.BmsCmd.\_\_eq\_\_}}\label{unitree-sdk2-cpp-idl-go2-bmscmd-dunder-eq-1}}

按底层 IDL 消息内容比较两个对象是否相等。

\textbf{签名}

\begin{lstlisting}[language=Python]
def __eq__(self, other: object) -> bool
\end{lstlisting}

\textbf{可用性与安全性}

\textbf{\passthrough{\lstinline!AVAILABLE!} /
\passthrough{\lstinline!VALUE\_TYPE!}}：当前绑定源码已实现该消息值操作。

\textbf{参数}

\begin{longtable}[]{@{}
  >{\raggedright\arraybackslash}p{(\columnwidth - 6\tabcolsep) * \real{0.18}}
  >{\raggedright\arraybackslash}p{(\columnwidth - 6\tabcolsep) * \real{0.24}}
  >{\raggedright\arraybackslash}p{(\columnwidth - 6\tabcolsep) * \real{0.18}}
  >{\raggedright\arraybackslash}p{(\columnwidth - 6\tabcolsep) * \real{0.40}}@{}}
\toprule
名称 & Python 类型 & 默认值 & 含义 \\
\midrule
\endhead
\passthrough{\lstinline!other!} & \passthrough{\lstinline!object!} &
必填 & 用于比较的另一个 Python 对象。类型不兼容时比较结果通常为
\passthrough{\lstinline!False!}。 \\
\bottomrule
\end{longtable}

\textbf{返回值}

\begin{longtable}[]{@{}
  >{\raggedright\arraybackslash}p{(\columnwidth - 4\tabcolsep) * \real{0.18}}
  >{\raggedright\arraybackslash}p{(\columnwidth - 4\tabcolsep) * \real{0.20}}
  >{\raggedright\arraybackslash}p{(\columnwidth - 4\tabcolsep) * \real{0.62}}@{}}
\toprule
位置 & 类型 & 含义 \\
\midrule
\endhead
返回值 & \passthrough{\lstinline!bool!} & 返回
\passthrough{\lstinline!bool!}。更精确的业务含义见该方法用途、C++
签名和目标型号协议。 \\
\bottomrule
\end{longtable}

\textbf{对应 C++}

\begin{itemize}
\tightlist
\item
  类：\passthrough{\lstinline!unitree\_go::msg::dds\_::BmsCmd\_!}
\item
  签名：\passthrough{\lstinline!operator==(const value \&) const!}
\item
  绑定策略：\passthrough{\lstinline!未单独标注!}
\item
  声明位置：由 pybind11 手工绑定或生成注册表提供；manifest
  未记录独立头文件行号。
\end{itemize}

\textbf{说明}

\begin{itemize}
\tightlist
\item
  示例是局部调用片段；\passthrough{\lstinline!obj!}
  和参数变量表示已经按上层业务校验并准备好的对象。
\end{itemize}

\textbf{用法}

\begin{lstlisting}[language=Python]
same = left == right
\end{lstlisting}

\hypertarget{unitree-sdk2-cpp-idl-go2-bmscmd-dunder-ne-1}{%
\paragraph{\texorpdfstring{\texttt{unitree\_sdk2\_cpp.idl.go2.BmsCmd.\_\_ne\_\_}}{unitree\_sdk2\_cpp.idl.go2.BmsCmd.\_\_ne\_\_}}\label{unitree-sdk2-cpp-idl-go2-bmscmd-dunder-ne-1}}

按底层 IDL 消息内容比较两个对象是否不相等。

\textbf{签名}

\begin{lstlisting}[language=Python]
def __ne__(self, other: object) -> bool
\end{lstlisting}

\textbf{可用性与安全性}

\textbf{\passthrough{\lstinline!AVAILABLE!} /
\passthrough{\lstinline!VALUE\_TYPE!}}：当前绑定源码已实现该消息值操作。

\textbf{参数}

\begin{longtable}[]{@{}
  >{\raggedright\arraybackslash}p{(\columnwidth - 6\tabcolsep) * \real{0.18}}
  >{\raggedright\arraybackslash}p{(\columnwidth - 6\tabcolsep) * \real{0.24}}
  >{\raggedright\arraybackslash}p{(\columnwidth - 6\tabcolsep) * \real{0.18}}
  >{\raggedright\arraybackslash}p{(\columnwidth - 6\tabcolsep) * \real{0.40}}@{}}
\toprule
名称 & Python 类型 & 默认值 & 含义 \\
\midrule
\endhead
\passthrough{\lstinline!other!} & \passthrough{\lstinline!object!} &
必填 & 用于比较的另一个 Python 对象。类型不兼容时比较结果通常为
\passthrough{\lstinline!False!}。 \\
\bottomrule
\end{longtable}

\textbf{返回值}

\begin{longtable}[]{@{}
  >{\raggedright\arraybackslash}p{(\columnwidth - 4\tabcolsep) * \real{0.18}}
  >{\raggedright\arraybackslash}p{(\columnwidth - 4\tabcolsep) * \real{0.20}}
  >{\raggedright\arraybackslash}p{(\columnwidth - 4\tabcolsep) * \real{0.62}}@{}}
\toprule
位置 & 类型 & 含义 \\
\midrule
\endhead
返回值 & \passthrough{\lstinline!bool!} & 返回
\passthrough{\lstinline!bool!}。更精确的业务含义见该方法用途、C++
签名和目标型号协议。 \\
\bottomrule
\end{longtable}

\textbf{对应 C++}

\begin{itemize}
\tightlist
\item
  类：\passthrough{\lstinline!unitree\_go::msg::dds\_::BmsCmd\_!}
\item
  签名：\passthrough{\lstinline"operator!=(const value \&) const"}
\item
  绑定策略：\passthrough{\lstinline!未单独标注!}
\item
  声明位置：由 pybind11 手工绑定或生成注册表提供；manifest
  未记录独立头文件行号。
\end{itemize}

\textbf{说明}

\begin{itemize}
\tightlist
\item
  示例是局部调用片段；\passthrough{\lstinline!obj!}
  和参数变量表示已经按上层业务校验并准备好的对象。
\end{itemize}

\textbf{用法}

\begin{lstlisting}[language=Python]
different = left != right
\end{lstlisting}

\hypertarget{unitree-sdk2-cpp-idl-go2-bmscmd-off}{%
\paragraph{\texorpdfstring{\texttt{unitree\_sdk2\_cpp.idl.go2.BmsCmd.off}}{unitree\_sdk2\_cpp.idl.go2.BmsCmd.off}}\label{unitree-sdk2-cpp-idl-go2-bmscmd-off}}

底层 \passthrough{\lstinline!unitree\_go::msg::dds\_::BmsCmd\_!} 的
\passthrough{\lstinline!off!} 字段。类型签名只说明 Python
表示；单位、数值范围、数组固定长度和枚举含义需查对应 IDL 头文件。
取值范围为 0 到 255。

\textbf{签名}

\begin{lstlisting}[language=Python]
@property
def off(self) -> int

@off.setter
def off(self, value: int) -> None
\end{lstlisting}

\textbf{可用性与安全性}

\textbf{\passthrough{\lstinline!AVAILABLE!} /
\passthrough{\lstinline!VALUE\_TYPE!}}：当前绑定源码已实现该消息值操作。

\textbf{参数}

\begin{longtable}[]{@{}
  >{\raggedright\arraybackslash}p{(\columnwidth - 4\tabcolsep) * \real{0.18}}
  >{\raggedright\arraybackslash}p{(\columnwidth - 4\tabcolsep) * \real{0.20}}
  >{\raggedright\arraybackslash}p{(\columnwidth - 4\tabcolsep) * \real{0.62}}@{}}
\toprule
名称 & Python 类型 & 含义 \\
\midrule
\endhead
\passthrough{\lstinline!value!} & \passthrough{\lstinline!int!} &
要写入或传递的新值，必须符合该参数的 Python 类型以及底层 C++
范围约束。 \\
\bottomrule
\end{longtable}

\textbf{返回值}

\begin{longtable}[]{@{}
  >{\raggedright\arraybackslash}p{(\columnwidth - 4\tabcolsep) * \real{0.18}}
  >{\raggedright\arraybackslash}p{(\columnwidth - 4\tabcolsep) * \real{0.20}}
  >{\raggedright\arraybackslash}p{(\columnwidth - 4\tabcolsep) * \real{0.62}}@{}}
\toprule
位置 & 类型 & 含义 \\
\midrule
\endhead
getter 返回值 & \passthrough{\lstinline!int!} & 返回属性当前值。 \\
\bottomrule
\end{longtable}

\textbf{对应 C++}

\begin{itemize}
\tightlist
\item
  类：\passthrough{\lstinline!unitree\_go::msg::dds\_::BmsCmd\_!}
\item
  字段类型：\passthrough{\lstinline!uint8\_t!}
\item
  声明位置：\passthrough{\lstinline!include/unitree/idl/go2/BmsCmd\_.hpp:24!}
\end{itemize}

\textbf{用法}

\begin{lstlisting}[language=Python]
current_value = obj.off
obj.off = new_value
\end{lstlisting}

\hypertarget{unitree-sdk2-cpp-idl-go2-bmscmd-reserve}{%
\paragraph{\texorpdfstring{\texttt{unitree\_sdk2\_cpp.idl.go2.BmsCmd.reserve}}{unitree\_sdk2\_cpp.idl.go2.BmsCmd.reserve}}\label{unitree-sdk2-cpp-idl-go2-bmscmd-reserve}}

底层 \passthrough{\lstinline!unitree\_go::msg::dds\_::BmsCmd\_!} 的
\passthrough{\lstinline!reserve!} 字段。类型签名只说明 Python
表示；单位、数值范围、数组固定长度和枚举含义需查对应 IDL 头文件。
底层是固定长度数组，写入序列必须正好包含 3 个元素；元素约束：取值范围为
0 到 255。

\textbf{签名}

\begin{lstlisting}[language=Python]
@property
def reserve(self) -> list[int]

@reserve.setter
def reserve(self, value: Sequence[int]) -> None
\end{lstlisting}

\textbf{可用性与安全性}

\textbf{\passthrough{\lstinline!AVAILABLE!} /
\passthrough{\lstinline!VALUE\_TYPE!}}：当前绑定源码已实现该消息值操作。

\textbf{参数}

\begin{longtable}[]{@{}
  >{\raggedright\arraybackslash}p{(\columnwidth - 4\tabcolsep) * \real{0.18}}
  >{\raggedright\arraybackslash}p{(\columnwidth - 4\tabcolsep) * \real{0.20}}
  >{\raggedright\arraybackslash}p{(\columnwidth - 4\tabcolsep) * \real{0.62}}@{}}
\toprule
名称 & Python 类型 & 含义 \\
\midrule
\endhead
\passthrough{\lstinline!value!} &
\passthrough{\lstinline!Sequence[int]!} &
要写入或传递的新值，必须符合该参数的 Python 类型以及底层 C++
范围约束。 \\
\bottomrule
\end{longtable}

\textbf{返回值}

\begin{longtable}[]{@{}
  >{\raggedright\arraybackslash}p{(\columnwidth - 4\tabcolsep) * \real{0.18}}
  >{\raggedright\arraybackslash}p{(\columnwidth - 4\tabcolsep) * \real{0.20}}
  >{\raggedright\arraybackslash}p{(\columnwidth - 4\tabcolsep) * \real{0.62}}@{}}
\toprule
位置 & 类型 & 含义 \\
\midrule
\endhead
getter 返回值 & \passthrough{\lstinline!list[int]!} &
返回属性当前值。数组、vector 和嵌套消息采用复制语义。 \\
\bottomrule
\end{longtable}

\textbf{对应 C++}

\begin{itemize}
\tightlist
\item
  类：\passthrough{\lstinline!unitree\_go::msg::dds\_::BmsCmd\_!}
\item
  字段类型：\passthrough{\lstinline!std::array<uint8\_t, 3>!}
\item
  声明位置：\passthrough{\lstinline!include/unitree/idl/go2/BmsCmd\_.hpp:25!}
\end{itemize}

\textbf{用法}

\begin{lstlisting}[language=Python]
current_value = obj.reserve
obj.reserve = new_value
\end{lstlisting}

\begin{quote}
{[}!TIP{]} 读取容器或嵌套值后应遵循``读取、修改、重新赋值''。只修改
getter 返回的副本不会可靠地更新原 C++ 消息。
\end{quote}

\hypertarget{unitree-sdk2-cpp-idl-go2-bmsstate}{%
\subsubsection{\texorpdfstring{\texttt{unitree\_sdk2\_cpp.idl.go2.BmsState}}{unitree\_sdk2\_cpp.idl.go2.BmsState}}\label{unitree-sdk2-cpp-idl-go2-bmsstate}}

可由 Python 读写的 DDS/IDL 消息值类型。字段采用复制语义。

\textbf{导入}

\begin{lstlisting}[language=Python]
from unitree_sdk2_cpp.idl.go2 import BmsState
\end{lstlisting}

\textbf{方法索引}

\begin{longtable}[]{@{}
  >{\raggedright\arraybackslash}p{(\columnwidth - 6\tabcolsep) * \real{0.18}}
  >{\raggedright\arraybackslash}p{(\columnwidth - 6\tabcolsep) * \real{0.24}}
  >{\raggedright\arraybackslash}p{(\columnwidth - 6\tabcolsep) * \real{0.18}}
  >{\raggedright\arraybackslash}p{(\columnwidth - 6\tabcolsep) * \real{0.40}}@{}}
\toprule
方法 & Python 签名 & 状态 & 安全分类 \\
\midrule
\endhead
\protect\hyperlink{unitree-sdk2-cpp-idl-go2-bmsstate-dunder-init-1}{\passthrough{\lstinline!\_\_init\_\_!}}
& \passthrough{\lstinline!def \_\_init\_\_(self) -> None!} &
\passthrough{\lstinline!AVAILABLE!} &
\passthrough{\lstinline!VALUE\_TYPE!} \\
\protect\hyperlink{unitree-sdk2-cpp-idl-go2-bmsstate-dunder-eq-1}{\passthrough{\lstinline!\_\_eq\_\_!}}
& \passthrough{\lstinline!def \_\_eq\_\_(self, other: object) -> bool!}
& \passthrough{\lstinline!AVAILABLE!} &
\passthrough{\lstinline!VALUE\_TYPE!} \\
\protect\hyperlink{unitree-sdk2-cpp-idl-go2-bmsstate-dunder-ne-1}{\passthrough{\lstinline!\_\_ne\_\_!}}
& \passthrough{\lstinline!def \_\_ne\_\_(self, other: object) -> bool!}
& \passthrough{\lstinline!AVAILABLE!} &
\passthrough{\lstinline!VALUE\_TYPE!} \\
\bottomrule
\end{longtable}

\textbf{属性索引}

\begin{longtable}[]{@{}
  >{\raggedright\arraybackslash}p{(\columnwidth - 6\tabcolsep) * \real{0.18}}
  >{\raggedright\arraybackslash}p{(\columnwidth - 6\tabcolsep) * \real{0.24}}
  >{\raggedright\arraybackslash}p{(\columnwidth - 6\tabcolsep) * \real{0.18}}
  >{\raggedright\arraybackslash}p{(\columnwidth - 6\tabcolsep) * \real{0.40}}@{}}
\toprule
属性 & 读取类型 & 写入类型 & 状态 \\
\midrule
\endhead
\protect\hyperlink{unitree-sdk2-cpp-idl-go2-bmsstate-version-high}{\passthrough{\lstinline!version\_high!}}
& \passthrough{\lstinline!int!} & \passthrough{\lstinline!int!} &
\passthrough{\lstinline!AVAILABLE!} \\
\protect\hyperlink{unitree-sdk2-cpp-idl-go2-bmsstate-version-low}{\passthrough{\lstinline!version\_low!}}
& \passthrough{\lstinline!int!} & \passthrough{\lstinline!int!} &
\passthrough{\lstinline!AVAILABLE!} \\
\protect\hyperlink{unitree-sdk2-cpp-idl-go2-bmsstate-status}{\passthrough{\lstinline!status!}}
& \passthrough{\lstinline!int!} & \passthrough{\lstinline!int!} &
\passthrough{\lstinline!AVAILABLE!} \\
\protect\hyperlink{unitree-sdk2-cpp-idl-go2-bmsstate-soc}{\passthrough{\lstinline!soc!}}
& \passthrough{\lstinline!int!} & \passthrough{\lstinline!int!} &
\passthrough{\lstinline!AVAILABLE!} \\
\protect\hyperlink{unitree-sdk2-cpp-idl-go2-bmsstate-current}{\passthrough{\lstinline!current!}}
& \passthrough{\lstinline!int!} & \passthrough{\lstinline!int!} &
\passthrough{\lstinline!AVAILABLE!} \\
\protect\hyperlink{unitree-sdk2-cpp-idl-go2-bmsstate-cycle}{\passthrough{\lstinline!cycle!}}
& \passthrough{\lstinline!int!} & \passthrough{\lstinline!int!} &
\passthrough{\lstinline!AVAILABLE!} \\
\protect\hyperlink{unitree-sdk2-cpp-idl-go2-bmsstate-bq-ntc}{\passthrough{\lstinline!bq\_ntc!}}
& \passthrough{\lstinline!list[int]!} &
\passthrough{\lstinline!Sequence[int]!} &
\passthrough{\lstinline!AVAILABLE!} \\
\protect\hyperlink{unitree-sdk2-cpp-idl-go2-bmsstate-mcu-ntc}{\passthrough{\lstinline!mcu\_ntc!}}
& \passthrough{\lstinline!list[int]!} &
\passthrough{\lstinline!Sequence[int]!} &
\passthrough{\lstinline!AVAILABLE!} \\
\protect\hyperlink{unitree-sdk2-cpp-idl-go2-bmsstate-cell-vol}{\passthrough{\lstinline!cell\_vol!}}
& \passthrough{\lstinline!list[int]!} &
\passthrough{\lstinline!Sequence[int]!} &
\passthrough{\lstinline!AVAILABLE!} \\
\bottomrule
\end{longtable}

\hypertarget{unitree-sdk2-cpp-idl-go2-bmsstate-dunder-init-1}{%
\paragraph{\texorpdfstring{\texttt{unitree\_sdk2\_cpp.idl.go2.BmsState.\_\_init\_\_}}{unitree\_sdk2\_cpp.idl.go2.BmsState.\_\_init\_\_}}\label{unitree-sdk2-cpp-idl-go2-bmsstate-dunder-init-1}}

初始化 \passthrough{\lstinline!BmsState!}
实例。是否能实际构造取决于下方可用性状态。

\textbf{签名}

\begin{lstlisting}[language=Python]
def __init__(self) -> None
\end{lstlisting}

\textbf{可用性与安全性}

\textbf{\passthrough{\lstinline!AVAILABLE!} /
\passthrough{\lstinline!VALUE\_TYPE!}}：当前绑定源码已实现该消息值操作。

\textbf{参数}

无显式参数。实例方法中的 \passthrough{\lstinline!self!} 由 Python
自动传入。

\textbf{返回值}

\begin{longtable}[]{@{}
  >{\raggedright\arraybackslash}p{(\columnwidth - 4\tabcolsep) * \real{0.18}}
  >{\raggedright\arraybackslash}p{(\columnwidth - 4\tabcolsep) * \real{0.20}}
  >{\raggedright\arraybackslash}p{(\columnwidth - 4\tabcolsep) * \real{0.62}}@{}}
\toprule
位置 & 类型 & 含义 \\
\midrule
\endhead
无 & \passthrough{\lstinline!None!} &
\passthrough{\lstinline!\_\_init\_\_!}
本身不返回值；调用类对象时，在构造函数可用的前提下得到该类实例。 \\
\bottomrule
\end{longtable}

\textbf{对应 C++}

\begin{itemize}
\tightlist
\item
  类：\passthrough{\lstinline!unitree\_go::msg::dds\_::BmsState\_!}
\item
  签名：\passthrough{\lstinline!BmsState\_()!}
\item
  绑定策略：\passthrough{\lstinline!未单独标注!}
\item
  声明位置：\passthrough{\lstinline!include/unitree/idl/go2/BmsState\_.hpp:35!}
\end{itemize}

\textbf{说明}

\begin{itemize}
\tightlist
\item
  示例是局部调用片段；\passthrough{\lstinline!obj!}
  和参数变量表示已经按上层业务校验并准备好的对象。
\end{itemize}

\textbf{用法}

\begin{lstlisting}[language=Python]
value = BmsState()
\end{lstlisting}

\hypertarget{unitree-sdk2-cpp-idl-go2-bmsstate-dunder-eq-1}{%
\paragraph{\texorpdfstring{\texttt{unitree\_sdk2\_cpp.idl.go2.BmsState.\_\_eq\_\_}}{unitree\_sdk2\_cpp.idl.go2.BmsState.\_\_eq\_\_}}\label{unitree-sdk2-cpp-idl-go2-bmsstate-dunder-eq-1}}

按底层 IDL 消息内容比较两个对象是否相等。

\textbf{签名}

\begin{lstlisting}[language=Python]
def __eq__(self, other: object) -> bool
\end{lstlisting}

\textbf{可用性与安全性}

\textbf{\passthrough{\lstinline!AVAILABLE!} /
\passthrough{\lstinline!VALUE\_TYPE!}}：当前绑定源码已实现该消息值操作。

\textbf{参数}

\begin{longtable}[]{@{}
  >{\raggedright\arraybackslash}p{(\columnwidth - 6\tabcolsep) * \real{0.18}}
  >{\raggedright\arraybackslash}p{(\columnwidth - 6\tabcolsep) * \real{0.24}}
  >{\raggedright\arraybackslash}p{(\columnwidth - 6\tabcolsep) * \real{0.18}}
  >{\raggedright\arraybackslash}p{(\columnwidth - 6\tabcolsep) * \real{0.40}}@{}}
\toprule
名称 & Python 类型 & 默认值 & 含义 \\
\midrule
\endhead
\passthrough{\lstinline!other!} & \passthrough{\lstinline!object!} &
必填 & 用于比较的另一个 Python 对象。类型不兼容时比较结果通常为
\passthrough{\lstinline!False!}。 \\
\bottomrule
\end{longtable}

\textbf{返回值}

\begin{longtable}[]{@{}
  >{\raggedright\arraybackslash}p{(\columnwidth - 4\tabcolsep) * \real{0.18}}
  >{\raggedright\arraybackslash}p{(\columnwidth - 4\tabcolsep) * \real{0.20}}
  >{\raggedright\arraybackslash}p{(\columnwidth - 4\tabcolsep) * \real{0.62}}@{}}
\toprule
位置 & 类型 & 含义 \\
\midrule
\endhead
返回值 & \passthrough{\lstinline!bool!} & 返回
\passthrough{\lstinline!bool!}。更精确的业务含义见该方法用途、C++
签名和目标型号协议。 \\
\bottomrule
\end{longtable}

\textbf{对应 C++}

\begin{itemize}
\tightlist
\item
  类：\passthrough{\lstinline!unitree\_go::msg::dds\_::BmsState\_!}
\item
  签名：\passthrough{\lstinline!operator==(const value \&) const!}
\item
  绑定策略：\passthrough{\lstinline!未单独标注!}
\item
  声明位置：由 pybind11 手工绑定或生成注册表提供；manifest
  未记录独立头文件行号。
\end{itemize}

\textbf{说明}

\begin{itemize}
\tightlist
\item
  示例是局部调用片段；\passthrough{\lstinline!obj!}
  和参数变量表示已经按上层业务校验并准备好的对象。
\end{itemize}

\textbf{用法}

\begin{lstlisting}[language=Python]
same = left == right
\end{lstlisting}

\hypertarget{unitree-sdk2-cpp-idl-go2-bmsstate-dunder-ne-1}{%
\paragraph{\texorpdfstring{\texttt{unitree\_sdk2\_cpp.idl.go2.BmsState.\_\_ne\_\_}}{unitree\_sdk2\_cpp.idl.go2.BmsState.\_\_ne\_\_}}\label{unitree-sdk2-cpp-idl-go2-bmsstate-dunder-ne-1}}

按底层 IDL 消息内容比较两个对象是否不相等。

\textbf{签名}

\begin{lstlisting}[language=Python]
def __ne__(self, other: object) -> bool
\end{lstlisting}

\textbf{可用性与安全性}

\textbf{\passthrough{\lstinline!AVAILABLE!} /
\passthrough{\lstinline!VALUE\_TYPE!}}：当前绑定源码已实现该消息值操作。

\textbf{参数}

\begin{longtable}[]{@{}
  >{\raggedright\arraybackslash}p{(\columnwidth - 6\tabcolsep) * \real{0.18}}
  >{\raggedright\arraybackslash}p{(\columnwidth - 6\tabcolsep) * \real{0.24}}
  >{\raggedright\arraybackslash}p{(\columnwidth - 6\tabcolsep) * \real{0.18}}
  >{\raggedright\arraybackslash}p{(\columnwidth - 6\tabcolsep) * \real{0.40}}@{}}
\toprule
名称 & Python 类型 & 默认值 & 含义 \\
\midrule
\endhead
\passthrough{\lstinline!other!} & \passthrough{\lstinline!object!} &
必填 & 用于比较的另一个 Python 对象。类型不兼容时比较结果通常为
\passthrough{\lstinline!False!}。 \\
\bottomrule
\end{longtable}

\textbf{返回值}

\begin{longtable}[]{@{}
  >{\raggedright\arraybackslash}p{(\columnwidth - 4\tabcolsep) * \real{0.18}}
  >{\raggedright\arraybackslash}p{(\columnwidth - 4\tabcolsep) * \real{0.20}}
  >{\raggedright\arraybackslash}p{(\columnwidth - 4\tabcolsep) * \real{0.62}}@{}}
\toprule
位置 & 类型 & 含义 \\
\midrule
\endhead
返回值 & \passthrough{\lstinline!bool!} & 返回
\passthrough{\lstinline!bool!}。更精确的业务含义见该方法用途、C++
签名和目标型号协议。 \\
\bottomrule
\end{longtable}

\textbf{对应 C++}

\begin{itemize}
\tightlist
\item
  类：\passthrough{\lstinline!unitree\_go::msg::dds\_::BmsState\_!}
\item
  签名：\passthrough{\lstinline"operator!=(const value \&) const"}
\item
  绑定策略：\passthrough{\lstinline!未单独标注!}
\item
  声明位置：由 pybind11 手工绑定或生成注册表提供；manifest
  未记录独立头文件行号。
\end{itemize}

\textbf{说明}

\begin{itemize}
\tightlist
\item
  示例是局部调用片段；\passthrough{\lstinline!obj!}
  和参数变量表示已经按上层业务校验并准备好的对象。
\end{itemize}

\textbf{用法}

\begin{lstlisting}[language=Python]
different = left != right
\end{lstlisting}

\hypertarget{unitree-sdk2-cpp-idl-go2-bmsstate-version-high}{%
\paragraph{\texorpdfstring{\texttt{unitree\_sdk2\_cpp.idl.go2.BmsState.version\_high}}{unitree\_sdk2\_cpp.idl.go2.BmsState.version\_high}}\label{unitree-sdk2-cpp-idl-go2-bmsstate-version-high}}

底层 \passthrough{\lstinline!unitree\_go::msg::dds\_::BmsState\_!} 的
\passthrough{\lstinline!version\_high!} 字段。类型签名只说明 Python
表示；单位、数值范围、数组固定长度和枚举含义需查对应 IDL 头文件。
取值范围为 0 到 255。

\textbf{签名}

\begin{lstlisting}[language=Python]
@property
def version_high(self) -> int

@version_high.setter
def version_high(self, value: int) -> None
\end{lstlisting}

\textbf{可用性与安全性}

\textbf{\passthrough{\lstinline!AVAILABLE!} /
\passthrough{\lstinline!VALUE\_TYPE!}}：当前绑定源码已实现该消息值操作。

\textbf{参数}

\begin{longtable}[]{@{}
  >{\raggedright\arraybackslash}p{(\columnwidth - 4\tabcolsep) * \real{0.18}}
  >{\raggedright\arraybackslash}p{(\columnwidth - 4\tabcolsep) * \real{0.20}}
  >{\raggedright\arraybackslash}p{(\columnwidth - 4\tabcolsep) * \real{0.62}}@{}}
\toprule
名称 & Python 类型 & 含义 \\
\midrule
\endhead
\passthrough{\lstinline!value!} & \passthrough{\lstinline!int!} &
要写入或传递的新值，必须符合该参数的 Python 类型以及底层 C++
范围约束。 \\
\bottomrule
\end{longtable}

\textbf{返回值}

\begin{longtable}[]{@{}
  >{\raggedright\arraybackslash}p{(\columnwidth - 4\tabcolsep) * \real{0.18}}
  >{\raggedright\arraybackslash}p{(\columnwidth - 4\tabcolsep) * \real{0.20}}
  >{\raggedright\arraybackslash}p{(\columnwidth - 4\tabcolsep) * \real{0.62}}@{}}
\toprule
位置 & 类型 & 含义 \\
\midrule
\endhead
getter 返回值 & \passthrough{\lstinline!int!} & 返回属性当前值。 \\
\bottomrule
\end{longtable}

\textbf{对应 C++}

\begin{itemize}
\tightlist
\item
  类：\passthrough{\lstinline!unitree\_go::msg::dds\_::BmsState\_!}
\item
  字段类型：\passthrough{\lstinline!uint8\_t!}
\item
  声明位置：\passthrough{\lstinline!include/unitree/idl/go2/BmsState\_.hpp:24!}
\end{itemize}

\textbf{用法}

\begin{lstlisting}[language=Python]
current_value = obj.version_high
obj.version_high = new_value
\end{lstlisting}

\hypertarget{unitree-sdk2-cpp-idl-go2-bmsstate-version-low}{%
\paragraph{\texorpdfstring{\texttt{unitree\_sdk2\_cpp.idl.go2.BmsState.version\_low}}{unitree\_sdk2\_cpp.idl.go2.BmsState.version\_low}}\label{unitree-sdk2-cpp-idl-go2-bmsstate-version-low}}

底层 \passthrough{\lstinline!unitree\_go::msg::dds\_::BmsState\_!} 的
\passthrough{\lstinline!version\_low!} 字段。类型签名只说明 Python
表示；单位、数值范围、数组固定长度和枚举含义需查对应 IDL 头文件。
取值范围为 0 到 255。

\textbf{签名}

\begin{lstlisting}[language=Python]
@property
def version_low(self) -> int

@version_low.setter
def version_low(self, value: int) -> None
\end{lstlisting}

\textbf{可用性与安全性}

\textbf{\passthrough{\lstinline!AVAILABLE!} /
\passthrough{\lstinline!VALUE\_TYPE!}}：当前绑定源码已实现该消息值操作。

\textbf{参数}

\begin{longtable}[]{@{}
  >{\raggedright\arraybackslash}p{(\columnwidth - 4\tabcolsep) * \real{0.18}}
  >{\raggedright\arraybackslash}p{(\columnwidth - 4\tabcolsep) * \real{0.20}}
  >{\raggedright\arraybackslash}p{(\columnwidth - 4\tabcolsep) * \real{0.62}}@{}}
\toprule
名称 & Python 类型 & 含义 \\
\midrule
\endhead
\passthrough{\lstinline!value!} & \passthrough{\lstinline!int!} &
要写入或传递的新值，必须符合该参数的 Python 类型以及底层 C++
范围约束。 \\
\bottomrule
\end{longtable}

\textbf{返回值}

\begin{longtable}[]{@{}
  >{\raggedright\arraybackslash}p{(\columnwidth - 4\tabcolsep) * \real{0.18}}
  >{\raggedright\arraybackslash}p{(\columnwidth - 4\tabcolsep) * \real{0.20}}
  >{\raggedright\arraybackslash}p{(\columnwidth - 4\tabcolsep) * \real{0.62}}@{}}
\toprule
位置 & 类型 & 含义 \\
\midrule
\endhead
getter 返回值 & \passthrough{\lstinline!int!} & 返回属性当前值。 \\
\bottomrule
\end{longtable}

\textbf{对应 C++}

\begin{itemize}
\tightlist
\item
  类：\passthrough{\lstinline!unitree\_go::msg::dds\_::BmsState\_!}
\item
  字段类型：\passthrough{\lstinline!uint8\_t!}
\item
  声明位置：\passthrough{\lstinline!include/unitree/idl/go2/BmsState\_.hpp:25!}
\end{itemize}

\textbf{用法}

\begin{lstlisting}[language=Python]
current_value = obj.version_low
obj.version_low = new_value
\end{lstlisting}

\hypertarget{unitree-sdk2-cpp-idl-go2-bmsstate-status}{%
\paragraph{\texorpdfstring{\texttt{unitree\_sdk2\_cpp.idl.go2.BmsState.status}}{unitree\_sdk2\_cpp.idl.go2.BmsState.status}}\label{unitree-sdk2-cpp-idl-go2-bmsstate-status}}

底层 \passthrough{\lstinline!unitree\_go::msg::dds\_::BmsState\_!} 的
\passthrough{\lstinline!status!} 字段。类型签名只说明 Python
表示；单位、数值范围、数组固定长度和枚举含义需查对应 IDL 头文件。
取值范围为 0 到 255。

\textbf{签名}

\begin{lstlisting}[language=Python]
@property
def status(self) -> int

@status.setter
def status(self, value: int) -> None
\end{lstlisting}

\textbf{可用性与安全性}

\textbf{\passthrough{\lstinline!AVAILABLE!} /
\passthrough{\lstinline!VALUE\_TYPE!}}：当前绑定源码已实现该消息值操作。

\textbf{参数}

\begin{longtable}[]{@{}
  >{\raggedright\arraybackslash}p{(\columnwidth - 4\tabcolsep) * \real{0.18}}
  >{\raggedright\arraybackslash}p{(\columnwidth - 4\tabcolsep) * \real{0.20}}
  >{\raggedright\arraybackslash}p{(\columnwidth - 4\tabcolsep) * \real{0.62}}@{}}
\toprule
名称 & Python 类型 & 含义 \\
\midrule
\endhead
\passthrough{\lstinline!value!} & \passthrough{\lstinline!int!} &
要写入或传递的新值，必须符合该参数的 Python 类型以及底层 C++
范围约束。 \\
\bottomrule
\end{longtable}

\textbf{返回值}

\begin{longtable}[]{@{}
  >{\raggedright\arraybackslash}p{(\columnwidth - 4\tabcolsep) * \real{0.18}}
  >{\raggedright\arraybackslash}p{(\columnwidth - 4\tabcolsep) * \real{0.20}}
  >{\raggedright\arraybackslash}p{(\columnwidth - 4\tabcolsep) * \real{0.62}}@{}}
\toprule
位置 & 类型 & 含义 \\
\midrule
\endhead
getter 返回值 & \passthrough{\lstinline!int!} & 返回属性当前值。 \\
\bottomrule
\end{longtable}

\textbf{对应 C++}

\begin{itemize}
\tightlist
\item
  类：\passthrough{\lstinline!unitree\_go::msg::dds\_::BmsState\_!}
\item
  字段类型：\passthrough{\lstinline!uint8\_t!}
\item
  声明位置：\passthrough{\lstinline!include/unitree/idl/go2/BmsState\_.hpp:26!}
\end{itemize}

\textbf{用法}

\begin{lstlisting}[language=Python]
current_value = obj.status
obj.status = new_value
\end{lstlisting}

\hypertarget{unitree-sdk2-cpp-idl-go2-bmsstate-soc}{%
\paragraph{\texorpdfstring{\texttt{unitree\_sdk2\_cpp.idl.go2.BmsState.soc}}{unitree\_sdk2\_cpp.idl.go2.BmsState.soc}}\label{unitree-sdk2-cpp-idl-go2-bmsstate-soc}}

底层 \passthrough{\lstinline!unitree\_go::msg::dds\_::BmsState\_!} 的
\passthrough{\lstinline!soc!} 字段。类型签名只说明 Python
表示；单位、数值范围、数组固定长度和枚举含义需查对应 IDL 头文件。
取值范围为 0 到 255。

\textbf{签名}

\begin{lstlisting}[language=Python]
@property
def soc(self) -> int

@soc.setter
def soc(self, value: int) -> None
\end{lstlisting}

\textbf{可用性与安全性}

\textbf{\passthrough{\lstinline!AVAILABLE!} /
\passthrough{\lstinline!VALUE\_TYPE!}}：当前绑定源码已实现该消息值操作。

\textbf{参数}

\begin{longtable}[]{@{}
  >{\raggedright\arraybackslash}p{(\columnwidth - 4\tabcolsep) * \real{0.18}}
  >{\raggedright\arraybackslash}p{(\columnwidth - 4\tabcolsep) * \real{0.20}}
  >{\raggedright\arraybackslash}p{(\columnwidth - 4\tabcolsep) * \real{0.62}}@{}}
\toprule
名称 & Python 类型 & 含义 \\
\midrule
\endhead
\passthrough{\lstinline!value!} & \passthrough{\lstinline!int!} &
要写入或传递的新值，必须符合该参数的 Python 类型以及底层 C++
范围约束。 \\
\bottomrule
\end{longtable}

\textbf{返回值}

\begin{longtable}[]{@{}
  >{\raggedright\arraybackslash}p{(\columnwidth - 4\tabcolsep) * \real{0.18}}
  >{\raggedright\arraybackslash}p{(\columnwidth - 4\tabcolsep) * \real{0.20}}
  >{\raggedright\arraybackslash}p{(\columnwidth - 4\tabcolsep) * \real{0.62}}@{}}
\toprule
位置 & 类型 & 含义 \\
\midrule
\endhead
getter 返回值 & \passthrough{\lstinline!int!} & 返回属性当前值。 \\
\bottomrule
\end{longtable}

\textbf{对应 C++}

\begin{itemize}
\tightlist
\item
  类：\passthrough{\lstinline!unitree\_go::msg::dds\_::BmsState\_!}
\item
  字段类型：\passthrough{\lstinline!uint8\_t!}
\item
  声明位置：\passthrough{\lstinline!include/unitree/idl/go2/BmsState\_.hpp:27!}
\end{itemize}

\textbf{用法}

\begin{lstlisting}[language=Python]
current_value = obj.soc
obj.soc = new_value
\end{lstlisting}

\hypertarget{unitree-sdk2-cpp-idl-go2-bmsstate-current}{%
\paragraph{\texorpdfstring{\texttt{unitree\_sdk2\_cpp.idl.go2.BmsState.current}}{unitree\_sdk2\_cpp.idl.go2.BmsState.current}}\label{unitree-sdk2-cpp-idl-go2-bmsstate-current}}

底层 \passthrough{\lstinline!unitree\_go::msg::dds\_::BmsState\_!} 的
\passthrough{\lstinline!current!} 字段。类型签名只说明 Python
表示；单位、数值范围、数组固定长度和枚举含义需查对应 IDL 头文件。
取值范围为 -2147483648 到 2147483647。

\textbf{签名}

\begin{lstlisting}[language=Python]
@property
def current(self) -> int

@current.setter
def current(self, value: int) -> None
\end{lstlisting}

\textbf{可用性与安全性}

\textbf{\passthrough{\lstinline!AVAILABLE!} /
\passthrough{\lstinline!VALUE\_TYPE!}}：当前绑定源码已实现该消息值操作。

\textbf{参数}

\begin{longtable}[]{@{}
  >{\raggedright\arraybackslash}p{(\columnwidth - 4\tabcolsep) * \real{0.18}}
  >{\raggedright\arraybackslash}p{(\columnwidth - 4\tabcolsep) * \real{0.20}}
  >{\raggedright\arraybackslash}p{(\columnwidth - 4\tabcolsep) * \real{0.62}}@{}}
\toprule
名称 & Python 类型 & 含义 \\
\midrule
\endhead
\passthrough{\lstinline!value!} & \passthrough{\lstinline!int!} &
要写入或传递的新值，必须符合该参数的 Python 类型以及底层 C++
范围约束。 \\
\bottomrule
\end{longtable}

\textbf{返回值}

\begin{longtable}[]{@{}
  >{\raggedright\arraybackslash}p{(\columnwidth - 4\tabcolsep) * \real{0.18}}
  >{\raggedright\arraybackslash}p{(\columnwidth - 4\tabcolsep) * \real{0.20}}
  >{\raggedright\arraybackslash}p{(\columnwidth - 4\tabcolsep) * \real{0.62}}@{}}
\toprule
位置 & 类型 & 含义 \\
\midrule
\endhead
getter 返回值 & \passthrough{\lstinline!int!} & 返回属性当前值。 \\
\bottomrule
\end{longtable}

\textbf{对应 C++}

\begin{itemize}
\tightlist
\item
  类：\passthrough{\lstinline!unitree\_go::msg::dds\_::BmsState\_!}
\item
  字段类型：\passthrough{\lstinline!int32\_t!}
\item
  声明位置：\passthrough{\lstinline!include/unitree/idl/go2/BmsState\_.hpp:28!}
\end{itemize}

\textbf{用法}

\begin{lstlisting}[language=Python]
current_value = obj.current
obj.current = new_value
\end{lstlisting}

\hypertarget{unitree-sdk2-cpp-idl-go2-bmsstate-cycle}{%
\paragraph{\texorpdfstring{\texttt{unitree\_sdk2\_cpp.idl.go2.BmsState.cycle}}{unitree\_sdk2\_cpp.idl.go2.BmsState.cycle}}\label{unitree-sdk2-cpp-idl-go2-bmsstate-cycle}}

底层 \passthrough{\lstinline!unitree\_go::msg::dds\_::BmsState\_!} 的
\passthrough{\lstinline!cycle!} 字段。类型签名只说明 Python
表示；单位、数值范围、数组固定长度和枚举含义需查对应 IDL 头文件。
取值范围为 0 到 65535。

\textbf{签名}

\begin{lstlisting}[language=Python]
@property
def cycle(self) -> int

@cycle.setter
def cycle(self, value: int) -> None
\end{lstlisting}

\textbf{可用性与安全性}

\textbf{\passthrough{\lstinline!AVAILABLE!} /
\passthrough{\lstinline!VALUE\_TYPE!}}：当前绑定源码已实现该消息值操作。

\textbf{参数}

\begin{longtable}[]{@{}
  >{\raggedright\arraybackslash}p{(\columnwidth - 4\tabcolsep) * \real{0.18}}
  >{\raggedright\arraybackslash}p{(\columnwidth - 4\tabcolsep) * \real{0.20}}
  >{\raggedright\arraybackslash}p{(\columnwidth - 4\tabcolsep) * \real{0.62}}@{}}
\toprule
名称 & Python 类型 & 含义 \\
\midrule
\endhead
\passthrough{\lstinline!value!} & \passthrough{\lstinline!int!} &
要写入或传递的新值，必须符合该参数的 Python 类型以及底层 C++
范围约束。 \\
\bottomrule
\end{longtable}

\textbf{返回值}

\begin{longtable}[]{@{}
  >{\raggedright\arraybackslash}p{(\columnwidth - 4\tabcolsep) * \real{0.18}}
  >{\raggedright\arraybackslash}p{(\columnwidth - 4\tabcolsep) * \real{0.20}}
  >{\raggedright\arraybackslash}p{(\columnwidth - 4\tabcolsep) * \real{0.62}}@{}}
\toprule
位置 & 类型 & 含义 \\
\midrule
\endhead
getter 返回值 & \passthrough{\lstinline!int!} & 返回属性当前值。 \\
\bottomrule
\end{longtable}

\textbf{对应 C++}

\begin{itemize}
\tightlist
\item
  类：\passthrough{\lstinline!unitree\_go::msg::dds\_::BmsState\_!}
\item
  字段类型：\passthrough{\lstinline!uint16\_t!}
\item
  声明位置：\passthrough{\lstinline!include/unitree/idl/go2/BmsState\_.hpp:29!}
\end{itemize}

\textbf{用法}

\begin{lstlisting}[language=Python]
current_value = obj.cycle
obj.cycle = new_value
\end{lstlisting}

\hypertarget{unitree-sdk2-cpp-idl-go2-bmsstate-bq-ntc}{%
\paragraph{\texorpdfstring{\texttt{unitree\_sdk2\_cpp.idl.go2.BmsState.bq\_ntc}}{unitree\_sdk2\_cpp.idl.go2.BmsState.bq\_ntc}}\label{unitree-sdk2-cpp-idl-go2-bmsstate-bq-ntc}}

底层 \passthrough{\lstinline!unitree\_go::msg::dds\_::BmsState\_!} 的
\passthrough{\lstinline!bq\_ntc!} 字段。类型签名只说明 Python
表示；单位、数值范围、数组固定长度和枚举含义需查对应 IDL 头文件。
底层是固定长度数组，写入序列必须正好包含 2 个元素；元素约束：取值范围为
0 到 255。

\textbf{签名}

\begin{lstlisting}[language=Python]
@property
def bq_ntc(self) -> list[int]

@bq_ntc.setter
def bq_ntc(self, value: Sequence[int]) -> None
\end{lstlisting}

\textbf{可用性与安全性}

\textbf{\passthrough{\lstinline!AVAILABLE!} /
\passthrough{\lstinline!VALUE\_TYPE!}}：当前绑定源码已实现该消息值操作。

\textbf{参数}

\begin{longtable}[]{@{}
  >{\raggedright\arraybackslash}p{(\columnwidth - 4\tabcolsep) * \real{0.18}}
  >{\raggedright\arraybackslash}p{(\columnwidth - 4\tabcolsep) * \real{0.20}}
  >{\raggedright\arraybackslash}p{(\columnwidth - 4\tabcolsep) * \real{0.62}}@{}}
\toprule
名称 & Python 类型 & 含义 \\
\midrule
\endhead
\passthrough{\lstinline!value!} &
\passthrough{\lstinline!Sequence[int]!} &
要写入或传递的新值，必须符合该参数的 Python 类型以及底层 C++
范围约束。 \\
\bottomrule
\end{longtable}

\textbf{返回值}

\begin{longtable}[]{@{}
  >{\raggedright\arraybackslash}p{(\columnwidth - 4\tabcolsep) * \real{0.18}}
  >{\raggedright\arraybackslash}p{(\columnwidth - 4\tabcolsep) * \real{0.20}}
  >{\raggedright\arraybackslash}p{(\columnwidth - 4\tabcolsep) * \real{0.62}}@{}}
\toprule
位置 & 类型 & 含义 \\
\midrule
\endhead
getter 返回值 & \passthrough{\lstinline!list[int]!} &
返回属性当前值。数组、vector 和嵌套消息采用复制语义。 \\
\bottomrule
\end{longtable}

\textbf{对应 C++}

\begin{itemize}
\tightlist
\item
  类：\passthrough{\lstinline!unitree\_go::msg::dds\_::BmsState\_!}
\item
  字段类型：\passthrough{\lstinline!std::array<uint8\_t, 2>!}
\item
  声明位置：\passthrough{\lstinline!include/unitree/idl/go2/BmsState\_.hpp:30!}
\end{itemize}

\textbf{用法}

\begin{lstlisting}[language=Python]
current_value = obj.bq_ntc
obj.bq_ntc = new_value
\end{lstlisting}

\begin{quote}
{[}!TIP{]} 读取容器或嵌套值后应遵循``读取、修改、重新赋值''。只修改
getter 返回的副本不会可靠地更新原 C++ 消息。
\end{quote}

\hypertarget{unitree-sdk2-cpp-idl-go2-bmsstate-mcu-ntc}{%
\paragraph{\texorpdfstring{\texttt{unitree\_sdk2\_cpp.idl.go2.BmsState.mcu\_ntc}}{unitree\_sdk2\_cpp.idl.go2.BmsState.mcu\_ntc}}\label{unitree-sdk2-cpp-idl-go2-bmsstate-mcu-ntc}}

底层 \passthrough{\lstinline!unitree\_go::msg::dds\_::BmsState\_!} 的
\passthrough{\lstinline!mcu\_ntc!} 字段。类型签名只说明 Python
表示；单位、数值范围、数组固定长度和枚举含义需查对应 IDL 头文件。
底层是固定长度数组，写入序列必须正好包含 2 个元素；元素约束：取值范围为
0 到 255。

\textbf{签名}

\begin{lstlisting}[language=Python]
@property
def mcu_ntc(self) -> list[int]

@mcu_ntc.setter
def mcu_ntc(self, value: Sequence[int]) -> None
\end{lstlisting}

\textbf{可用性与安全性}

\textbf{\passthrough{\lstinline!AVAILABLE!} /
\passthrough{\lstinline!VALUE\_TYPE!}}：当前绑定源码已实现该消息值操作。

\textbf{参数}

\begin{longtable}[]{@{}
  >{\raggedright\arraybackslash}p{(\columnwidth - 4\tabcolsep) * \real{0.18}}
  >{\raggedright\arraybackslash}p{(\columnwidth - 4\tabcolsep) * \real{0.20}}
  >{\raggedright\arraybackslash}p{(\columnwidth - 4\tabcolsep) * \real{0.62}}@{}}
\toprule
名称 & Python 类型 & 含义 \\
\midrule
\endhead
\passthrough{\lstinline!value!} &
\passthrough{\lstinline!Sequence[int]!} &
要写入或传递的新值，必须符合该参数的 Python 类型以及底层 C++
范围约束。 \\
\bottomrule
\end{longtable}

\textbf{返回值}

\begin{longtable}[]{@{}
  >{\raggedright\arraybackslash}p{(\columnwidth - 4\tabcolsep) * \real{0.18}}
  >{\raggedright\arraybackslash}p{(\columnwidth - 4\tabcolsep) * \real{0.20}}
  >{\raggedright\arraybackslash}p{(\columnwidth - 4\tabcolsep) * \real{0.62}}@{}}
\toprule
位置 & 类型 & 含义 \\
\midrule
\endhead
getter 返回值 & \passthrough{\lstinline!list[int]!} &
返回属性当前值。数组、vector 和嵌套消息采用复制语义。 \\
\bottomrule
\end{longtable}

\textbf{对应 C++}

\begin{itemize}
\tightlist
\item
  类：\passthrough{\lstinline!unitree\_go::msg::dds\_::BmsState\_!}
\item
  字段类型：\passthrough{\lstinline!std::array<uint8\_t, 2>!}
\item
  声明位置：\passthrough{\lstinline!include/unitree/idl/go2/BmsState\_.hpp:31!}
\end{itemize}

\textbf{用法}

\begin{lstlisting}[language=Python]
current_value = obj.mcu_ntc
obj.mcu_ntc = new_value
\end{lstlisting}

\begin{quote}
{[}!TIP{]} 读取容器或嵌套值后应遵循``读取、修改、重新赋值''。只修改
getter 返回的副本不会可靠地更新原 C++ 消息。
\end{quote}

\hypertarget{unitree-sdk2-cpp-idl-go2-bmsstate-cell-vol}{%
\paragraph{\texorpdfstring{\texttt{unitree\_sdk2\_cpp.idl.go2.BmsState.cell\_vol}}{unitree\_sdk2\_cpp.idl.go2.BmsState.cell\_vol}}\label{unitree-sdk2-cpp-idl-go2-bmsstate-cell-vol}}

底层 \passthrough{\lstinline!unitree\_go::msg::dds\_::BmsState\_!} 的
\passthrough{\lstinline!cell\_vol!} 字段。类型签名只说明 Python
表示；单位、数值范围、数组固定长度和枚举含义需查对应 IDL 头文件。
底层是固定长度数组，写入序列必须正好包含 15 个元素；元素约束：取值范围为
0 到 65535。

\textbf{签名}

\begin{lstlisting}[language=Python]
@property
def cell_vol(self) -> list[int]

@cell_vol.setter
def cell_vol(self, value: Sequence[int]) -> None
\end{lstlisting}

\textbf{可用性与安全性}

\textbf{\passthrough{\lstinline!AVAILABLE!} /
\passthrough{\lstinline!VALUE\_TYPE!}}：当前绑定源码已实现该消息值操作。

\textbf{参数}

\begin{longtable}[]{@{}
  >{\raggedright\arraybackslash}p{(\columnwidth - 4\tabcolsep) * \real{0.18}}
  >{\raggedright\arraybackslash}p{(\columnwidth - 4\tabcolsep) * \real{0.20}}
  >{\raggedright\arraybackslash}p{(\columnwidth - 4\tabcolsep) * \real{0.62}}@{}}
\toprule
名称 & Python 类型 & 含义 \\
\midrule
\endhead
\passthrough{\lstinline!value!} &
\passthrough{\lstinline!Sequence[int]!} &
要写入或传递的新值，必须符合该参数的 Python 类型以及底层 C++
范围约束。 \\
\bottomrule
\end{longtable}

\textbf{返回值}

\begin{longtable}[]{@{}
  >{\raggedright\arraybackslash}p{(\columnwidth - 4\tabcolsep) * \real{0.18}}
  >{\raggedright\arraybackslash}p{(\columnwidth - 4\tabcolsep) * \real{0.20}}
  >{\raggedright\arraybackslash}p{(\columnwidth - 4\tabcolsep) * \real{0.62}}@{}}
\toprule
位置 & 类型 & 含义 \\
\midrule
\endhead
getter 返回值 & \passthrough{\lstinline!list[int]!} &
返回属性当前值。数组、vector 和嵌套消息采用复制语义。 \\
\bottomrule
\end{longtable}

\textbf{对应 C++}

\begin{itemize}
\tightlist
\item
  类：\passthrough{\lstinline!unitree\_go::msg::dds\_::BmsState\_!}
\item
  字段类型：\passthrough{\lstinline!std::array<uint16\_t, 15>!}
\item
  声明位置：\passthrough{\lstinline!include/unitree/idl/go2/BmsState\_.hpp:32!}
\end{itemize}

\textbf{用法}

\begin{lstlisting}[language=Python]
current_value = obj.cell_vol
obj.cell_vol = new_value
\end{lstlisting}

\begin{quote}
{[}!TIP{]} 读取容器或嵌套值后应遵循``读取、修改、重新赋值''。只修改
getter 返回的副本不会可靠地更新原 C++ 消息。
\end{quote}

\hypertarget{unitree-sdk2-cpp-idl-go2-configchangestatus}{%
\subsubsection{\texorpdfstring{\texttt{unitree\_sdk2\_cpp.idl.go2.ConfigChangeStatus}}{unitree\_sdk2\_cpp.idl.go2.ConfigChangeStatus}}\label{unitree-sdk2-cpp-idl-go2-configchangestatus}}

可由 Python 读写的 DDS/IDL 消息值类型。字段采用复制语义。

\textbf{导入}

\begin{lstlisting}[language=Python]
from unitree_sdk2_cpp.idl.go2 import ConfigChangeStatus
\end{lstlisting}

\textbf{方法索引}

\begin{longtable}[]{@{}
  >{\raggedright\arraybackslash}p{(\columnwidth - 6\tabcolsep) * \real{0.18}}
  >{\raggedright\arraybackslash}p{(\columnwidth - 6\tabcolsep) * \real{0.24}}
  >{\raggedright\arraybackslash}p{(\columnwidth - 6\tabcolsep) * \real{0.18}}
  >{\raggedright\arraybackslash}p{(\columnwidth - 6\tabcolsep) * \real{0.40}}@{}}
\toprule
方法 & Python 签名 & 状态 & 安全分类 \\
\midrule
\endhead
\protect\hyperlink{unitree-sdk2-cpp-idl-go2-configchangestatus-dunder-init-1}{\passthrough{\lstinline!\_\_init\_\_!}}
& \passthrough{\lstinline!def \_\_init\_\_(self) -> None!} &
\passthrough{\lstinline!AVAILABLE!} &
\passthrough{\lstinline!VALUE\_TYPE!} \\
\protect\hyperlink{unitree-sdk2-cpp-idl-go2-configchangestatus-dunder-eq-1}{\passthrough{\lstinline!\_\_eq\_\_!}}
& \passthrough{\lstinline!def \_\_eq\_\_(self, other: object) -> bool!}
& \passthrough{\lstinline!AVAILABLE!} &
\passthrough{\lstinline!VALUE\_TYPE!} \\
\protect\hyperlink{unitree-sdk2-cpp-idl-go2-configchangestatus-dunder-ne-1}{\passthrough{\lstinline!\_\_ne\_\_!}}
& \passthrough{\lstinline!def \_\_ne\_\_(self, other: object) -> bool!}
& \passthrough{\lstinline!AVAILABLE!} &
\passthrough{\lstinline!VALUE\_TYPE!} \\
\bottomrule
\end{longtable}

\textbf{属性索引}

\begin{longtable}[]{@{}
  >{\raggedright\arraybackslash}p{(\columnwidth - 6\tabcolsep) * \real{0.18}}
  >{\raggedright\arraybackslash}p{(\columnwidth - 6\tabcolsep) * \real{0.24}}
  >{\raggedright\arraybackslash}p{(\columnwidth - 6\tabcolsep) * \real{0.18}}
  >{\raggedright\arraybackslash}p{(\columnwidth - 6\tabcolsep) * \real{0.40}}@{}}
\toprule
属性 & 读取类型 & 写入类型 & 状态 \\
\midrule
\endhead
\protect\hyperlink{unitree-sdk2-cpp-idl-go2-configchangestatus-name}{\passthrough{\lstinline!name!}}
& \passthrough{\lstinline!str!} & \passthrough{\lstinline!str!} &
\passthrough{\lstinline!AVAILABLE!} \\
\protect\hyperlink{unitree-sdk2-cpp-idl-go2-configchangestatus-content}{\passthrough{\lstinline!content!}}
& \passthrough{\lstinline!str!} & \passthrough{\lstinline!str!} &
\passthrough{\lstinline!AVAILABLE!} \\
\bottomrule
\end{longtable}

\hypertarget{unitree-sdk2-cpp-idl-go2-configchangestatus-dunder-init-1}{%
\paragraph{\texorpdfstring{\texttt{unitree\_sdk2\_cpp.idl.go2.ConfigChangeStatus.\_\_init\_\_}}{unitree\_sdk2\_cpp.idl.go2.ConfigChangeStatus.\_\_init\_\_}}\label{unitree-sdk2-cpp-idl-go2-configchangestatus-dunder-init-1}}

初始化 \passthrough{\lstinline!ConfigChangeStatus!}
实例。是否能实际构造取决于下方可用性状态。

\textbf{签名}

\begin{lstlisting}[language=Python]
def __init__(self) -> None
\end{lstlisting}

\textbf{可用性与安全性}

\textbf{\passthrough{\lstinline!AVAILABLE!} /
\passthrough{\lstinline!VALUE\_TYPE!}}：当前绑定源码已实现该消息值操作。

\textbf{参数}

无显式参数。实例方法中的 \passthrough{\lstinline!self!} 由 Python
自动传入。

\textbf{返回值}

\begin{longtable}[]{@{}
  >{\raggedright\arraybackslash}p{(\columnwidth - 4\tabcolsep) * \real{0.18}}
  >{\raggedright\arraybackslash}p{(\columnwidth - 4\tabcolsep) * \real{0.20}}
  >{\raggedright\arraybackslash}p{(\columnwidth - 4\tabcolsep) * \real{0.62}}@{}}
\toprule
位置 & 类型 & 含义 \\
\midrule
\endhead
无 & \passthrough{\lstinline!None!} &
\passthrough{\lstinline!\_\_init\_\_!}
本身不返回值；调用类对象时，在构造函数可用的前提下得到该类实例。 \\
\bottomrule
\end{longtable}

\textbf{对应 C++}

\begin{itemize}
\tightlist
\item
  类：\passthrough{\lstinline!unitree\_go::msg::dds\_::ConfigChangeStatus\_!}
\item
  签名：\passthrough{\lstinline!ConfigChangeStatus\_()!}
\item
  绑定策略：\passthrough{\lstinline!未单独标注!}
\item
  声明位置：\passthrough{\lstinline!include/unitree/idl/go2/ConfigChangeStatus\_.hpp:27!}
\end{itemize}

\textbf{说明}

\begin{itemize}
\tightlist
\item
  示例是局部调用片段；\passthrough{\lstinline!obj!}
  和参数变量表示已经按上层业务校验并准备好的对象。
\end{itemize}

\textbf{用法}

\begin{lstlisting}[language=Python]
value = ConfigChangeStatus()
\end{lstlisting}

\hypertarget{unitree-sdk2-cpp-idl-go2-configchangestatus-dunder-eq-1}{%
\paragraph{\texorpdfstring{\texttt{unitree\_sdk2\_cpp.idl.go2.ConfigChangeStatus.\_\_eq\_\_}}{unitree\_sdk2\_cpp.idl.go2.ConfigChangeStatus.\_\_eq\_\_}}\label{unitree-sdk2-cpp-idl-go2-configchangestatus-dunder-eq-1}}

按底层 IDL 消息内容比较两个对象是否相等。

\textbf{签名}

\begin{lstlisting}[language=Python]
def __eq__(self, other: object) -> bool
\end{lstlisting}

\textbf{可用性与安全性}

\textbf{\passthrough{\lstinline!AVAILABLE!} /
\passthrough{\lstinline!VALUE\_TYPE!}}：当前绑定源码已实现该消息值操作。

\textbf{参数}

\begin{longtable}[]{@{}
  >{\raggedright\arraybackslash}p{(\columnwidth - 6\tabcolsep) * \real{0.18}}
  >{\raggedright\arraybackslash}p{(\columnwidth - 6\tabcolsep) * \real{0.24}}
  >{\raggedright\arraybackslash}p{(\columnwidth - 6\tabcolsep) * \real{0.18}}
  >{\raggedright\arraybackslash}p{(\columnwidth - 6\tabcolsep) * \real{0.40}}@{}}
\toprule
名称 & Python 类型 & 默认值 & 含义 \\
\midrule
\endhead
\passthrough{\lstinline!other!} & \passthrough{\lstinline!object!} &
必填 & 用于比较的另一个 Python 对象。类型不兼容时比较结果通常为
\passthrough{\lstinline!False!}。 \\
\bottomrule
\end{longtable}

\textbf{返回值}

\begin{longtable}[]{@{}
  >{\raggedright\arraybackslash}p{(\columnwidth - 4\tabcolsep) * \real{0.18}}
  >{\raggedright\arraybackslash}p{(\columnwidth - 4\tabcolsep) * \real{0.20}}
  >{\raggedright\arraybackslash}p{(\columnwidth - 4\tabcolsep) * \real{0.62}}@{}}
\toprule
位置 & 类型 & 含义 \\
\midrule
\endhead
返回值 & \passthrough{\lstinline!bool!} & 返回
\passthrough{\lstinline!bool!}。更精确的业务含义见该方法用途、C++
签名和目标型号协议。 \\
\bottomrule
\end{longtable}

\textbf{对应 C++}

\begin{itemize}
\tightlist
\item
  类：\passthrough{\lstinline!unitree\_go::msg::dds\_::ConfigChangeStatus\_!}
\item
  签名：\passthrough{\lstinline!operator==(const value \&) const!}
\item
  绑定策略：\passthrough{\lstinline!未单独标注!}
\item
  声明位置：由 pybind11 手工绑定或生成注册表提供；manifest
  未记录独立头文件行号。
\end{itemize}

\textbf{说明}

\begin{itemize}
\tightlist
\item
  示例是局部调用片段；\passthrough{\lstinline!obj!}
  和参数变量表示已经按上层业务校验并准备好的对象。
\end{itemize}

\textbf{用法}

\begin{lstlisting}[language=Python]
same = left == right
\end{lstlisting}

\hypertarget{unitree-sdk2-cpp-idl-go2-configchangestatus-dunder-ne-1}{%
\paragraph{\texorpdfstring{\texttt{unitree\_sdk2\_cpp.idl.go2.ConfigChangeStatus.\_\_ne\_\_}}{unitree\_sdk2\_cpp.idl.go2.ConfigChangeStatus.\_\_ne\_\_}}\label{unitree-sdk2-cpp-idl-go2-configchangestatus-dunder-ne-1}}

按底层 IDL 消息内容比较两个对象是否不相等。

\textbf{签名}

\begin{lstlisting}[language=Python]
def __ne__(self, other: object) -> bool
\end{lstlisting}

\textbf{可用性与安全性}

\textbf{\passthrough{\lstinline!AVAILABLE!} /
\passthrough{\lstinline!VALUE\_TYPE!}}：当前绑定源码已实现该消息值操作。

\textbf{参数}

\begin{longtable}[]{@{}
  >{\raggedright\arraybackslash}p{(\columnwidth - 6\tabcolsep) * \real{0.18}}
  >{\raggedright\arraybackslash}p{(\columnwidth - 6\tabcolsep) * \real{0.24}}
  >{\raggedright\arraybackslash}p{(\columnwidth - 6\tabcolsep) * \real{0.18}}
  >{\raggedright\arraybackslash}p{(\columnwidth - 6\tabcolsep) * \real{0.40}}@{}}
\toprule
名称 & Python 类型 & 默认值 & 含义 \\
\midrule
\endhead
\passthrough{\lstinline!other!} & \passthrough{\lstinline!object!} &
必填 & 用于比较的另一个 Python 对象。类型不兼容时比较结果通常为
\passthrough{\lstinline!False!}。 \\
\bottomrule
\end{longtable}

\textbf{返回值}

\begin{longtable}[]{@{}
  >{\raggedright\arraybackslash}p{(\columnwidth - 4\tabcolsep) * \real{0.18}}
  >{\raggedright\arraybackslash}p{(\columnwidth - 4\tabcolsep) * \real{0.20}}
  >{\raggedright\arraybackslash}p{(\columnwidth - 4\tabcolsep) * \real{0.62}}@{}}
\toprule
位置 & 类型 & 含义 \\
\midrule
\endhead
返回值 & \passthrough{\lstinline!bool!} & 返回
\passthrough{\lstinline!bool!}。更精确的业务含义见该方法用途、C++
签名和目标型号协议。 \\
\bottomrule
\end{longtable}

\textbf{对应 C++}

\begin{itemize}
\tightlist
\item
  类：\passthrough{\lstinline!unitree\_go::msg::dds\_::ConfigChangeStatus\_!}
\item
  签名：\passthrough{\lstinline"operator!=(const value \&) const"}
\item
  绑定策略：\passthrough{\lstinline!未单独标注!}
\item
  声明位置：由 pybind11 手工绑定或生成注册表提供；manifest
  未记录独立头文件行号。
\end{itemize}

\textbf{说明}

\begin{itemize}
\tightlist
\item
  示例是局部调用片段；\passthrough{\lstinline!obj!}
  和参数变量表示已经按上层业务校验并准备好的对象。
\end{itemize}

\textbf{用法}

\begin{lstlisting}[language=Python]
different = left != right
\end{lstlisting}

\hypertarget{unitree-sdk2-cpp-idl-go2-configchangestatus-name}{%
\paragraph{\texorpdfstring{\texttt{unitree\_sdk2\_cpp.idl.go2.ConfigChangeStatus.name}}{unitree\_sdk2\_cpp.idl.go2.ConfigChangeStatus.name}}\label{unitree-sdk2-cpp-idl-go2-configchangestatus-name}}

底层
\passthrough{\lstinline!unitree\_go::msg::dds\_::ConfigChangeStatus\_!}
的 \passthrough{\lstinline!name!} 字段。类型签名只说明 Python
表示；单位、数值范围、数组固定长度和枚举含义需查对应 IDL 头文件。
底层为字符串；长度、编码和允许值由具体协议决定。

\textbf{签名}

\begin{lstlisting}[language=Python]
@property
def name(self) -> str

@name.setter
def name(self, value: str) -> None
\end{lstlisting}

\textbf{可用性与安全性}

\textbf{\passthrough{\lstinline!AVAILABLE!} /
\passthrough{\lstinline!VALUE\_TYPE!}}：当前绑定源码已实现该消息值操作。

\textbf{参数}

\begin{longtable}[]{@{}
  >{\raggedright\arraybackslash}p{(\columnwidth - 4\tabcolsep) * \real{0.18}}
  >{\raggedright\arraybackslash}p{(\columnwidth - 4\tabcolsep) * \real{0.20}}
  >{\raggedright\arraybackslash}p{(\columnwidth - 4\tabcolsep) * \real{0.62}}@{}}
\toprule
名称 & Python 类型 & 含义 \\
\midrule
\endhead
\passthrough{\lstinline!value!} & \passthrough{\lstinline!str!} &
要写入或传递的新值，必须符合该参数的 Python 类型以及底层 C++
范围约束。 \\
\bottomrule
\end{longtable}

\textbf{返回值}

\begin{longtable}[]{@{}
  >{\raggedright\arraybackslash}p{(\columnwidth - 4\tabcolsep) * \real{0.18}}
  >{\raggedright\arraybackslash}p{(\columnwidth - 4\tabcolsep) * \real{0.20}}
  >{\raggedright\arraybackslash}p{(\columnwidth - 4\tabcolsep) * \real{0.62}}@{}}
\toprule
位置 & 类型 & 含义 \\
\midrule
\endhead
getter 返回值 & \passthrough{\lstinline!str!} & 返回属性当前值。 \\
\bottomrule
\end{longtable}

\textbf{对应 C++}

\begin{itemize}
\tightlist
\item
  类：\passthrough{\lstinline!unitree\_go::msg::dds\_::ConfigChangeStatus\_!}
\item
  字段类型：\passthrough{\lstinline!std::string!}
\item
  声明位置：\passthrough{\lstinline!include/unitree/idl/go2/ConfigChangeStatus\_.hpp:23!}
\end{itemize}

\textbf{用法}

\begin{lstlisting}[language=Python]
current_value = obj.name
obj.name = new_value
\end{lstlisting}

\hypertarget{unitree-sdk2-cpp-idl-go2-configchangestatus-content}{%
\paragraph{\texorpdfstring{\texttt{unitree\_sdk2\_cpp.idl.go2.ConfigChangeStatus.content}}{unitree\_sdk2\_cpp.idl.go2.ConfigChangeStatus.content}}\label{unitree-sdk2-cpp-idl-go2-configchangestatus-content}}

底层
\passthrough{\lstinline!unitree\_go::msg::dds\_::ConfigChangeStatus\_!}
的 \passthrough{\lstinline!content!} 字段。类型签名只说明 Python
表示；单位、数值范围、数组固定长度和枚举含义需查对应 IDL 头文件。
底层为字符串；长度、编码和允许值由具体协议决定。

\textbf{签名}

\begin{lstlisting}[language=Python]
@property
def content(self) -> str

@content.setter
def content(self, value: str) -> None
\end{lstlisting}

\textbf{可用性与安全性}

\textbf{\passthrough{\lstinline!AVAILABLE!} /
\passthrough{\lstinline!VALUE\_TYPE!}}：当前绑定源码已实现该消息值操作。

\textbf{参数}

\begin{longtable}[]{@{}
  >{\raggedright\arraybackslash}p{(\columnwidth - 4\tabcolsep) * \real{0.18}}
  >{\raggedright\arraybackslash}p{(\columnwidth - 4\tabcolsep) * \real{0.20}}
  >{\raggedright\arraybackslash}p{(\columnwidth - 4\tabcolsep) * \real{0.62}}@{}}
\toprule
名称 & Python 类型 & 含义 \\
\midrule
\endhead
\passthrough{\lstinline!value!} & \passthrough{\lstinline!str!} &
要写入或传递的新值，必须符合该参数的 Python 类型以及底层 C++
范围约束。 \\
\bottomrule
\end{longtable}

\textbf{返回值}

\begin{longtable}[]{@{}
  >{\raggedright\arraybackslash}p{(\columnwidth - 4\tabcolsep) * \real{0.18}}
  >{\raggedright\arraybackslash}p{(\columnwidth - 4\tabcolsep) * \real{0.20}}
  >{\raggedright\arraybackslash}p{(\columnwidth - 4\tabcolsep) * \real{0.62}}@{}}
\toprule
位置 & 类型 & 含义 \\
\midrule
\endhead
getter 返回值 & \passthrough{\lstinline!str!} & 返回属性当前值。 \\
\bottomrule
\end{longtable}

\textbf{对应 C++}

\begin{itemize}
\tightlist
\item
  类：\passthrough{\lstinline!unitree\_go::msg::dds\_::ConfigChangeStatus\_!}
\item
  字段类型：\passthrough{\lstinline!std::string!}
\item
  声明位置：\passthrough{\lstinline!include/unitree/idl/go2/ConfigChangeStatus\_.hpp:24!}
\end{itemize}

\textbf{用法}

\begin{lstlisting}[language=Python]
current_value = obj.content
obj.content = new_value
\end{lstlisting}

\hypertarget{unitree-sdk2-cpp-idl-go2-error}{%
\subsubsection{\texorpdfstring{\texttt{unitree\_sdk2\_cpp.idl.go2.Error}}{unitree\_sdk2\_cpp.idl.go2.Error}}\label{unitree-sdk2-cpp-idl-go2-error}}

可由 Python 读写的 DDS/IDL 消息值类型。字段采用复制语义。

\textbf{导入}

\begin{lstlisting}[language=Python]
from unitree_sdk2_cpp.idl.go2 import Error
\end{lstlisting}

\textbf{方法索引}

\begin{longtable}[]{@{}
  >{\raggedright\arraybackslash}p{(\columnwidth - 6\tabcolsep) * \real{0.18}}
  >{\raggedright\arraybackslash}p{(\columnwidth - 6\tabcolsep) * \real{0.24}}
  >{\raggedright\arraybackslash}p{(\columnwidth - 6\tabcolsep) * \real{0.18}}
  >{\raggedright\arraybackslash}p{(\columnwidth - 6\tabcolsep) * \real{0.40}}@{}}
\toprule
方法 & Python 签名 & 状态 & 安全分类 \\
\midrule
\endhead
\protect\hyperlink{unitree-sdk2-cpp-idl-go2-error-dunder-init-1}{\passthrough{\lstinline!\_\_init\_\_!}}
& \passthrough{\lstinline!def \_\_init\_\_(self) -> None!} &
\passthrough{\lstinline!AVAILABLE!} &
\passthrough{\lstinline!VALUE\_TYPE!} \\
\protect\hyperlink{unitree-sdk2-cpp-idl-go2-error-dunder-eq-1}{\passthrough{\lstinline!\_\_eq\_\_!}}
& \passthrough{\lstinline!def \_\_eq\_\_(self, other: object) -> bool!}
& \passthrough{\lstinline!AVAILABLE!} &
\passthrough{\lstinline!VALUE\_TYPE!} \\
\protect\hyperlink{unitree-sdk2-cpp-idl-go2-error-dunder-ne-1}{\passthrough{\lstinline!\_\_ne\_\_!}}
& \passthrough{\lstinline!def \_\_ne\_\_(self, other: object) -> bool!}
& \passthrough{\lstinline!AVAILABLE!} &
\passthrough{\lstinline!VALUE\_TYPE!} \\
\bottomrule
\end{longtable}

\textbf{属性索引}

\begin{longtable}[]{@{}
  >{\raggedright\arraybackslash}p{(\columnwidth - 6\tabcolsep) * \real{0.18}}
  >{\raggedright\arraybackslash}p{(\columnwidth - 6\tabcolsep) * \real{0.24}}
  >{\raggedright\arraybackslash}p{(\columnwidth - 6\tabcolsep) * \real{0.18}}
  >{\raggedright\arraybackslash}p{(\columnwidth - 6\tabcolsep) * \real{0.40}}@{}}
\toprule
属性 & 读取类型 & 写入类型 & 状态 \\
\midrule
\endhead
\protect\hyperlink{unitree-sdk2-cpp-idl-go2-error-source}{\passthrough{\lstinline!source!}}
& \passthrough{\lstinline!int!} & \passthrough{\lstinline!int!} &
\passthrough{\lstinline!AVAILABLE!} \\
\protect\hyperlink{unitree-sdk2-cpp-idl-go2-error-state}{\passthrough{\lstinline!state!}}
& \passthrough{\lstinline!int!} & \passthrough{\lstinline!int!} &
\passthrough{\lstinline!AVAILABLE!} \\
\bottomrule
\end{longtable}

\hypertarget{unitree-sdk2-cpp-idl-go2-error-dunder-init-1}{%
\paragraph{\texorpdfstring{\texttt{unitree\_sdk2\_cpp.idl.go2.Error.\_\_init\_\_}}{unitree\_sdk2\_cpp.idl.go2.Error.\_\_init\_\_}}\label{unitree-sdk2-cpp-idl-go2-error-dunder-init-1}}

初始化 \passthrough{\lstinline!Error!}
实例。是否能实际构造取决于下方可用性状态。

\textbf{签名}

\begin{lstlisting}[language=Python]
def __init__(self) -> None
\end{lstlisting}

\textbf{可用性与安全性}

\textbf{\passthrough{\lstinline!AVAILABLE!} /
\passthrough{\lstinline!VALUE\_TYPE!}}：当前绑定源码已实现该消息值操作。

\textbf{参数}

无显式参数。实例方法中的 \passthrough{\lstinline!self!} 由 Python
自动传入。

\textbf{返回值}

\begin{longtable}[]{@{}
  >{\raggedright\arraybackslash}p{(\columnwidth - 4\tabcolsep) * \real{0.18}}
  >{\raggedright\arraybackslash}p{(\columnwidth - 4\tabcolsep) * \real{0.20}}
  >{\raggedright\arraybackslash}p{(\columnwidth - 4\tabcolsep) * \real{0.62}}@{}}
\toprule
位置 & 类型 & 含义 \\
\midrule
\endhead
无 & \passthrough{\lstinline!None!} &
\passthrough{\lstinline!\_\_init\_\_!}
本身不返回值；调用类对象时，在构造函数可用的前提下得到该类实例。 \\
\bottomrule
\end{longtable}

\textbf{对应 C++}

\begin{itemize}
\tightlist
\item
  类：\passthrough{\lstinline!unitree\_go::msg::dds\_::Error\_!}
\item
  签名：\passthrough{\lstinline!Error\_()!}
\item
  绑定策略：\passthrough{\lstinline!未单独标注!}
\item
  声明位置：\passthrough{\lstinline!include/unitree/idl/go2/Error\_.hpp:27!}
\end{itemize}

\textbf{说明}

\begin{itemize}
\tightlist
\item
  示例是局部调用片段；\passthrough{\lstinline!obj!}
  和参数变量表示已经按上层业务校验并准备好的对象。
\end{itemize}

\textbf{用法}

\begin{lstlisting}[language=Python]
value = Error()
\end{lstlisting}

\hypertarget{unitree-sdk2-cpp-idl-go2-error-dunder-eq-1}{%
\paragraph{\texorpdfstring{\texttt{unitree\_sdk2\_cpp.idl.go2.Error.\_\_eq\_\_}}{unitree\_sdk2\_cpp.idl.go2.Error.\_\_eq\_\_}}\label{unitree-sdk2-cpp-idl-go2-error-dunder-eq-1}}

按底层 IDL 消息内容比较两个对象是否相等。

\textbf{签名}

\begin{lstlisting}[language=Python]
def __eq__(self, other: object) -> bool
\end{lstlisting}

\textbf{可用性与安全性}

\textbf{\passthrough{\lstinline!AVAILABLE!} /
\passthrough{\lstinline!VALUE\_TYPE!}}：当前绑定源码已实现该消息值操作。

\textbf{参数}

\begin{longtable}[]{@{}
  >{\raggedright\arraybackslash}p{(\columnwidth - 6\tabcolsep) * \real{0.18}}
  >{\raggedright\arraybackslash}p{(\columnwidth - 6\tabcolsep) * \real{0.24}}
  >{\raggedright\arraybackslash}p{(\columnwidth - 6\tabcolsep) * \real{0.18}}
  >{\raggedright\arraybackslash}p{(\columnwidth - 6\tabcolsep) * \real{0.40}}@{}}
\toprule
名称 & Python 类型 & 默认值 & 含义 \\
\midrule
\endhead
\passthrough{\lstinline!other!} & \passthrough{\lstinline!object!} &
必填 & 用于比较的另一个 Python 对象。类型不兼容时比较结果通常为
\passthrough{\lstinline!False!}。 \\
\bottomrule
\end{longtable}

\textbf{返回值}

\begin{longtable}[]{@{}
  >{\raggedright\arraybackslash}p{(\columnwidth - 4\tabcolsep) * \real{0.18}}
  >{\raggedright\arraybackslash}p{(\columnwidth - 4\tabcolsep) * \real{0.20}}
  >{\raggedright\arraybackslash}p{(\columnwidth - 4\tabcolsep) * \real{0.62}}@{}}
\toprule
位置 & 类型 & 含义 \\
\midrule
\endhead
返回值 & \passthrough{\lstinline!bool!} & 返回
\passthrough{\lstinline!bool!}。更精确的业务含义见该方法用途、C++
签名和目标型号协议。 \\
\bottomrule
\end{longtable}

\textbf{对应 C++}

\begin{itemize}
\tightlist
\item
  类：\passthrough{\lstinline!unitree\_go::msg::dds\_::Error\_!}
\item
  签名：\passthrough{\lstinline!operator==(const value \&) const!}
\item
  绑定策略：\passthrough{\lstinline!未单独标注!}
\item
  声明位置：由 pybind11 手工绑定或生成注册表提供；manifest
  未记录独立头文件行号。
\end{itemize}

\textbf{说明}

\begin{itemize}
\tightlist
\item
  示例是局部调用片段；\passthrough{\lstinline!obj!}
  和参数变量表示已经按上层业务校验并准备好的对象。
\end{itemize}

\textbf{用法}

\begin{lstlisting}[language=Python]
same = left == right
\end{lstlisting}

\hypertarget{unitree-sdk2-cpp-idl-go2-error-dunder-ne-1}{%
\paragraph{\texorpdfstring{\texttt{unitree\_sdk2\_cpp.idl.go2.Error.\_\_ne\_\_}}{unitree\_sdk2\_cpp.idl.go2.Error.\_\_ne\_\_}}\label{unitree-sdk2-cpp-idl-go2-error-dunder-ne-1}}

按底层 IDL 消息内容比较两个对象是否不相等。

\textbf{签名}

\begin{lstlisting}[language=Python]
def __ne__(self, other: object) -> bool
\end{lstlisting}

\textbf{可用性与安全性}

\textbf{\passthrough{\lstinline!AVAILABLE!} /
\passthrough{\lstinline!VALUE\_TYPE!}}：当前绑定源码已实现该消息值操作。

\textbf{参数}

\begin{longtable}[]{@{}
  >{\raggedright\arraybackslash}p{(\columnwidth - 6\tabcolsep) * \real{0.18}}
  >{\raggedright\arraybackslash}p{(\columnwidth - 6\tabcolsep) * \real{0.24}}
  >{\raggedright\arraybackslash}p{(\columnwidth - 6\tabcolsep) * \real{0.18}}
  >{\raggedright\arraybackslash}p{(\columnwidth - 6\tabcolsep) * \real{0.40}}@{}}
\toprule
名称 & Python 类型 & 默认值 & 含义 \\
\midrule
\endhead
\passthrough{\lstinline!other!} & \passthrough{\lstinline!object!} &
必填 & 用于比较的另一个 Python 对象。类型不兼容时比较结果通常为
\passthrough{\lstinline!False!}。 \\
\bottomrule
\end{longtable}

\textbf{返回值}

\begin{longtable}[]{@{}
  >{\raggedright\arraybackslash}p{(\columnwidth - 4\tabcolsep) * \real{0.18}}
  >{\raggedright\arraybackslash}p{(\columnwidth - 4\tabcolsep) * \real{0.20}}
  >{\raggedright\arraybackslash}p{(\columnwidth - 4\tabcolsep) * \real{0.62}}@{}}
\toprule
位置 & 类型 & 含义 \\
\midrule
\endhead
返回值 & \passthrough{\lstinline!bool!} & 返回
\passthrough{\lstinline!bool!}。更精确的业务含义见该方法用途、C++
签名和目标型号协议。 \\
\bottomrule
\end{longtable}

\textbf{对应 C++}

\begin{itemize}
\tightlist
\item
  类：\passthrough{\lstinline!unitree\_go::msg::dds\_::Error\_!}
\item
  签名：\passthrough{\lstinline"operator!=(const value \&) const"}
\item
  绑定策略：\passthrough{\lstinline!未单独标注!}
\item
  声明位置：由 pybind11 手工绑定或生成注册表提供；manifest
  未记录独立头文件行号。
\end{itemize}

\textbf{说明}

\begin{itemize}
\tightlist
\item
  示例是局部调用片段；\passthrough{\lstinline!obj!}
  和参数变量表示已经按上层业务校验并准备好的对象。
\end{itemize}

\textbf{用法}

\begin{lstlisting}[language=Python]
different = left != right
\end{lstlisting}

\hypertarget{unitree-sdk2-cpp-idl-go2-error-source}{%
\paragraph{\texorpdfstring{\texttt{unitree\_sdk2\_cpp.idl.go2.Error.source}}{unitree\_sdk2\_cpp.idl.go2.Error.source}}\label{unitree-sdk2-cpp-idl-go2-error-source}}

底层 \passthrough{\lstinline!unitree\_go::msg::dds\_::Error\_!} 的
\passthrough{\lstinline!source!} 字段。类型签名只说明 Python
表示；单位、数值范围、数组固定长度和枚举含义需查对应 IDL 头文件。
取值范围为 0 到 4294967295。

\textbf{签名}

\begin{lstlisting}[language=Python]
@property
def source(self) -> int

@source.setter
def source(self, value: int) -> None
\end{lstlisting}

\textbf{可用性与安全性}

\textbf{\passthrough{\lstinline!AVAILABLE!} /
\passthrough{\lstinline!VALUE\_TYPE!}}：当前绑定源码已实现该消息值操作。

\textbf{参数}

\begin{longtable}[]{@{}
  >{\raggedright\arraybackslash}p{(\columnwidth - 4\tabcolsep) * \real{0.18}}
  >{\raggedright\arraybackslash}p{(\columnwidth - 4\tabcolsep) * \real{0.20}}
  >{\raggedright\arraybackslash}p{(\columnwidth - 4\tabcolsep) * \real{0.62}}@{}}
\toprule
名称 & Python 类型 & 含义 \\
\midrule
\endhead
\passthrough{\lstinline!value!} & \passthrough{\lstinline!int!} &
要写入或传递的新值，必须符合该参数的 Python 类型以及底层 C++
范围约束。 \\
\bottomrule
\end{longtable}

\textbf{返回值}

\begin{longtable}[]{@{}
  >{\raggedright\arraybackslash}p{(\columnwidth - 4\tabcolsep) * \real{0.18}}
  >{\raggedright\arraybackslash}p{(\columnwidth - 4\tabcolsep) * \real{0.20}}
  >{\raggedright\arraybackslash}p{(\columnwidth - 4\tabcolsep) * \real{0.62}}@{}}
\toprule
位置 & 类型 & 含义 \\
\midrule
\endhead
getter 返回值 & \passthrough{\lstinline!int!} & 返回属性当前值。 \\
\bottomrule
\end{longtable}

\textbf{对应 C++}

\begin{itemize}
\tightlist
\item
  类：\passthrough{\lstinline!unitree\_go::msg::dds\_::Error\_!}
\item
  字段类型：\passthrough{\lstinline!uint32\_t!}
\item
  声明位置：\passthrough{\lstinline!include/unitree/idl/go2/Error\_.hpp:23!}
\end{itemize}

\textbf{用法}

\begin{lstlisting}[language=Python]
current_value = obj.source
obj.source = new_value
\end{lstlisting}

\hypertarget{unitree-sdk2-cpp-idl-go2-error-state}{%
\paragraph{\texorpdfstring{\texttt{unitree\_sdk2\_cpp.idl.go2.Error.state}}{unitree\_sdk2\_cpp.idl.go2.Error.state}}\label{unitree-sdk2-cpp-idl-go2-error-state}}

底层 \passthrough{\lstinline!unitree\_go::msg::dds\_::Error\_!} 的
\passthrough{\lstinline!state!} 字段。类型签名只说明 Python
表示；单位、数值范围、数组固定长度和枚举含义需查对应 IDL 头文件。
取值范围为 0 到 4294967295。

\textbf{签名}

\begin{lstlisting}[language=Python]
@property
def state(self) -> int

@state.setter
def state(self, value: int) -> None
\end{lstlisting}

\textbf{可用性与安全性}

\textbf{\passthrough{\lstinline!AVAILABLE!} /
\passthrough{\lstinline!VALUE\_TYPE!}}：当前绑定源码已实现该消息值操作。

\textbf{参数}

\begin{longtable}[]{@{}
  >{\raggedright\arraybackslash}p{(\columnwidth - 4\tabcolsep) * \real{0.18}}
  >{\raggedright\arraybackslash}p{(\columnwidth - 4\tabcolsep) * \real{0.20}}
  >{\raggedright\arraybackslash}p{(\columnwidth - 4\tabcolsep) * \real{0.62}}@{}}
\toprule
名称 & Python 类型 & 含义 \\
\midrule
\endhead
\passthrough{\lstinline!value!} & \passthrough{\lstinline!int!} &
要写入或传递的新值，必须符合该参数的 Python 类型以及底层 C++
范围约束。 \\
\bottomrule
\end{longtable}

\textbf{返回值}

\begin{longtable}[]{@{}
  >{\raggedright\arraybackslash}p{(\columnwidth - 4\tabcolsep) * \real{0.18}}
  >{\raggedright\arraybackslash}p{(\columnwidth - 4\tabcolsep) * \real{0.20}}
  >{\raggedright\arraybackslash}p{(\columnwidth - 4\tabcolsep) * \real{0.62}}@{}}
\toprule
位置 & 类型 & 含义 \\
\midrule
\endhead
getter 返回值 & \passthrough{\lstinline!int!} & 返回属性当前值。 \\
\bottomrule
\end{longtable}

\textbf{对应 C++}

\begin{itemize}
\tightlist
\item
  类：\passthrough{\lstinline!unitree\_go::msg::dds\_::Error\_!}
\item
  字段类型：\passthrough{\lstinline!uint32\_t!}
\item
  声明位置：\passthrough{\lstinline!include/unitree/idl/go2/Error\_.hpp:24!}
\end{itemize}

\textbf{用法}

\begin{lstlisting}[language=Python]
current_value = obj.state
obj.state = new_value
\end{lstlisting}

\hypertarget{unitree-sdk2-cpp-idl-go2-go2frontvideodata}{%
\subsubsection{\texorpdfstring{\texttt{unitree\_sdk2\_cpp.idl.go2.Go2FrontVideoData}}{unitree\_sdk2\_cpp.idl.go2.Go2FrontVideoData}}\label{unitree-sdk2-cpp-idl-go2-go2frontvideodata}}

可由 Python 读写的 DDS/IDL 消息值类型。字段采用复制语义。

\textbf{导入}

\begin{lstlisting}[language=Python]
from unitree_sdk2_cpp.idl.go2 import Go2FrontVideoData
\end{lstlisting}

\textbf{方法索引}

\begin{longtable}[]{@{}
  >{\raggedright\arraybackslash}p{(\columnwidth - 6\tabcolsep) * \real{0.18}}
  >{\raggedright\arraybackslash}p{(\columnwidth - 6\tabcolsep) * \real{0.24}}
  >{\raggedright\arraybackslash}p{(\columnwidth - 6\tabcolsep) * \real{0.18}}
  >{\raggedright\arraybackslash}p{(\columnwidth - 6\tabcolsep) * \real{0.40}}@{}}
\toprule
方法 & Python 签名 & 状态 & 安全分类 \\
\midrule
\endhead
\protect\hyperlink{unitree-sdk2-cpp-idl-go2-go2frontvideodata-dunder-init-1}{\passthrough{\lstinline!\_\_init\_\_!}}
& \passthrough{\lstinline!def \_\_init\_\_(self) -> None!} &
\passthrough{\lstinline!AVAILABLE!} &
\passthrough{\lstinline!VALUE\_TYPE!} \\
\protect\hyperlink{unitree-sdk2-cpp-idl-go2-go2frontvideodata-dunder-eq-1}{\passthrough{\lstinline!\_\_eq\_\_!}}
& \passthrough{\lstinline!def \_\_eq\_\_(self, other: object) -> bool!}
& \passthrough{\lstinline!AVAILABLE!} &
\passthrough{\lstinline!VALUE\_TYPE!} \\
\protect\hyperlink{unitree-sdk2-cpp-idl-go2-go2frontvideodata-dunder-ne-1}{\passthrough{\lstinline!\_\_ne\_\_!}}
& \passthrough{\lstinline!def \_\_ne\_\_(self, other: object) -> bool!}
& \passthrough{\lstinline!AVAILABLE!} &
\passthrough{\lstinline!VALUE\_TYPE!} \\
\bottomrule
\end{longtable}

\textbf{属性索引}

\begin{longtable}[]{@{}
  >{\raggedright\arraybackslash}p{(\columnwidth - 6\tabcolsep) * \real{0.18}}
  >{\raggedright\arraybackslash}p{(\columnwidth - 6\tabcolsep) * \real{0.24}}
  >{\raggedright\arraybackslash}p{(\columnwidth - 6\tabcolsep) * \real{0.18}}
  >{\raggedright\arraybackslash}p{(\columnwidth - 6\tabcolsep) * \real{0.40}}@{}}
\toprule
属性 & 读取类型 & 写入类型 & 状态 \\
\midrule
\endhead
\protect\hyperlink{unitree-sdk2-cpp-idl-go2-go2frontvideodata-time-frame}{\passthrough{\lstinline!time\_frame!}}
& \passthrough{\lstinline!int!} & \passthrough{\lstinline!int!} &
\passthrough{\lstinline!AVAILABLE!} \\
\protect\hyperlink{unitree-sdk2-cpp-idl-go2-go2frontvideodata-video720p}{\passthrough{\lstinline!video720p!}}
& \passthrough{\lstinline!list[int]!} &
\passthrough{\lstinline!Sequence[int]!} &
\passthrough{\lstinline!AVAILABLE!} \\
\protect\hyperlink{unitree-sdk2-cpp-idl-go2-go2frontvideodata-video360p}{\passthrough{\lstinline!video360p!}}
& \passthrough{\lstinline!list[int]!} &
\passthrough{\lstinline!Sequence[int]!} &
\passthrough{\lstinline!AVAILABLE!} \\
\protect\hyperlink{unitree-sdk2-cpp-idl-go2-go2frontvideodata-video180p}{\passthrough{\lstinline!video180p!}}
& \passthrough{\lstinline!list[int]!} &
\passthrough{\lstinline!Sequence[int]!} &
\passthrough{\lstinline!AVAILABLE!} \\
\bottomrule
\end{longtable}

\hypertarget{unitree-sdk2-cpp-idl-go2-go2frontvideodata-dunder-init-1}{%
\paragraph{\texorpdfstring{\texttt{unitree\_sdk2\_cpp.idl.go2.Go2FrontVideoData.\_\_init\_\_}}{unitree\_sdk2\_cpp.idl.go2.Go2FrontVideoData.\_\_init\_\_}}\label{unitree-sdk2-cpp-idl-go2-go2frontvideodata-dunder-init-1}}

初始化 \passthrough{\lstinline!Go2FrontVideoData!}
实例。是否能实际构造取决于下方可用性状态。

\textbf{签名}

\begin{lstlisting}[language=Python]
def __init__(self) -> None
\end{lstlisting}

\textbf{可用性与安全性}

\textbf{\passthrough{\lstinline!AVAILABLE!} /
\passthrough{\lstinline!VALUE\_TYPE!}}：当前绑定源码已实现该消息值操作。

\textbf{参数}

无显式参数。实例方法中的 \passthrough{\lstinline!self!} 由 Python
自动传入。

\textbf{返回值}

\begin{longtable}[]{@{}
  >{\raggedright\arraybackslash}p{(\columnwidth - 4\tabcolsep) * \real{0.18}}
  >{\raggedright\arraybackslash}p{(\columnwidth - 4\tabcolsep) * \real{0.20}}
  >{\raggedright\arraybackslash}p{(\columnwidth - 4\tabcolsep) * \real{0.62}}@{}}
\toprule
位置 & 类型 & 含义 \\
\midrule
\endhead
无 & \passthrough{\lstinline!None!} &
\passthrough{\lstinline!\_\_init\_\_!}
本身不返回值；调用类对象时，在构造函数可用的前提下得到该类实例。 \\
\bottomrule
\end{longtable}

\textbf{对应 C++}

\begin{itemize}
\tightlist
\item
  类：\passthrough{\lstinline!unitree\_go::msg::dds\_::Go2FrontVideoData\_!}
\item
  签名：\passthrough{\lstinline!Go2FrontVideoData\_()!}
\item
  绑定策略：\passthrough{\lstinline!未单独标注!}
\item
  声明位置：\passthrough{\lstinline!include/unitree/idl/go2/Go2FrontVideoData\_.hpp:30!}
\end{itemize}

\textbf{说明}

\begin{itemize}
\tightlist
\item
  示例是局部调用片段；\passthrough{\lstinline!obj!}
  和参数变量表示已经按上层业务校验并准备好的对象。
\end{itemize}

\textbf{用法}

\begin{lstlisting}[language=Python]
value = Go2FrontVideoData()
\end{lstlisting}

\hypertarget{unitree-sdk2-cpp-idl-go2-go2frontvideodata-dunder-eq-1}{%
\paragraph{\texorpdfstring{\texttt{unitree\_sdk2\_cpp.idl.go2.Go2FrontVideoData.\_\_eq\_\_}}{unitree\_sdk2\_cpp.idl.go2.Go2FrontVideoData.\_\_eq\_\_}}\label{unitree-sdk2-cpp-idl-go2-go2frontvideodata-dunder-eq-1}}

按底层 IDL 消息内容比较两个对象是否相等。

\textbf{签名}

\begin{lstlisting}[language=Python]
def __eq__(self, other: object) -> bool
\end{lstlisting}

\textbf{可用性与安全性}

\textbf{\passthrough{\lstinline!AVAILABLE!} /
\passthrough{\lstinline!VALUE\_TYPE!}}：当前绑定源码已实现该消息值操作。

\textbf{参数}

\begin{longtable}[]{@{}
  >{\raggedright\arraybackslash}p{(\columnwidth - 6\tabcolsep) * \real{0.18}}
  >{\raggedright\arraybackslash}p{(\columnwidth - 6\tabcolsep) * \real{0.24}}
  >{\raggedright\arraybackslash}p{(\columnwidth - 6\tabcolsep) * \real{0.18}}
  >{\raggedright\arraybackslash}p{(\columnwidth - 6\tabcolsep) * \real{0.40}}@{}}
\toprule
名称 & Python 类型 & 默认值 & 含义 \\
\midrule
\endhead
\passthrough{\lstinline!other!} & \passthrough{\lstinline!object!} &
必填 & 用于比较的另一个 Python 对象。类型不兼容时比较结果通常为
\passthrough{\lstinline!False!}。 \\
\bottomrule
\end{longtable}

\textbf{返回值}

\begin{longtable}[]{@{}
  >{\raggedright\arraybackslash}p{(\columnwidth - 4\tabcolsep) * \real{0.18}}
  >{\raggedright\arraybackslash}p{(\columnwidth - 4\tabcolsep) * \real{0.20}}
  >{\raggedright\arraybackslash}p{(\columnwidth - 4\tabcolsep) * \real{0.62}}@{}}
\toprule
位置 & 类型 & 含义 \\
\midrule
\endhead
返回值 & \passthrough{\lstinline!bool!} & 返回
\passthrough{\lstinline!bool!}。更精确的业务含义见该方法用途、C++
签名和目标型号协议。 \\
\bottomrule
\end{longtable}

\textbf{对应 C++}

\begin{itemize}
\tightlist
\item
  类：\passthrough{\lstinline!unitree\_go::msg::dds\_::Go2FrontVideoData\_!}
\item
  签名：\passthrough{\lstinline!operator==(const value \&) const!}
\item
  绑定策略：\passthrough{\lstinline!未单独标注!}
\item
  声明位置：由 pybind11 手工绑定或生成注册表提供；manifest
  未记录独立头文件行号。
\end{itemize}

\textbf{说明}

\begin{itemize}
\tightlist
\item
  示例是局部调用片段；\passthrough{\lstinline!obj!}
  和参数变量表示已经按上层业务校验并准备好的对象。
\end{itemize}

\textbf{用法}

\begin{lstlisting}[language=Python]
same = left == right
\end{lstlisting}

\hypertarget{unitree-sdk2-cpp-idl-go2-go2frontvideodata-dunder-ne-1}{%
\paragraph{\texorpdfstring{\texttt{unitree\_sdk2\_cpp.idl.go2.Go2FrontVideoData.\_\_ne\_\_}}{unitree\_sdk2\_cpp.idl.go2.Go2FrontVideoData.\_\_ne\_\_}}\label{unitree-sdk2-cpp-idl-go2-go2frontvideodata-dunder-ne-1}}

按底层 IDL 消息内容比较两个对象是否不相等。

\textbf{签名}

\begin{lstlisting}[language=Python]
def __ne__(self, other: object) -> bool
\end{lstlisting}

\textbf{可用性与安全性}

\textbf{\passthrough{\lstinline!AVAILABLE!} /
\passthrough{\lstinline!VALUE\_TYPE!}}：当前绑定源码已实现该消息值操作。

\textbf{参数}

\begin{longtable}[]{@{}
  >{\raggedright\arraybackslash}p{(\columnwidth - 6\tabcolsep) * \real{0.18}}
  >{\raggedright\arraybackslash}p{(\columnwidth - 6\tabcolsep) * \real{0.24}}
  >{\raggedright\arraybackslash}p{(\columnwidth - 6\tabcolsep) * \real{0.18}}
  >{\raggedright\arraybackslash}p{(\columnwidth - 6\tabcolsep) * \real{0.40}}@{}}
\toprule
名称 & Python 类型 & 默认值 & 含义 \\
\midrule
\endhead
\passthrough{\lstinline!other!} & \passthrough{\lstinline!object!} &
必填 & 用于比较的另一个 Python 对象。类型不兼容时比较结果通常为
\passthrough{\lstinline!False!}。 \\
\bottomrule
\end{longtable}

\textbf{返回值}

\begin{longtable}[]{@{}
  >{\raggedright\arraybackslash}p{(\columnwidth - 4\tabcolsep) * \real{0.18}}
  >{\raggedright\arraybackslash}p{(\columnwidth - 4\tabcolsep) * \real{0.20}}
  >{\raggedright\arraybackslash}p{(\columnwidth - 4\tabcolsep) * \real{0.62}}@{}}
\toprule
位置 & 类型 & 含义 \\
\midrule
\endhead
返回值 & \passthrough{\lstinline!bool!} & 返回
\passthrough{\lstinline!bool!}。更精确的业务含义见该方法用途、C++
签名和目标型号协议。 \\
\bottomrule
\end{longtable}

\textbf{对应 C++}

\begin{itemize}
\tightlist
\item
  类：\passthrough{\lstinline!unitree\_go::msg::dds\_::Go2FrontVideoData\_!}
\item
  签名：\passthrough{\lstinline"operator!=(const value \&) const"}
\item
  绑定策略：\passthrough{\lstinline!未单独标注!}
\item
  声明位置：由 pybind11 手工绑定或生成注册表提供；manifest
  未记录独立头文件行号。
\end{itemize}

\textbf{说明}

\begin{itemize}
\tightlist
\item
  示例是局部调用片段；\passthrough{\lstinline!obj!}
  和参数变量表示已经按上层业务校验并准备好的对象。
\end{itemize}

\textbf{用法}

\begin{lstlisting}[language=Python]
different = left != right
\end{lstlisting}

\hypertarget{unitree-sdk2-cpp-idl-go2-go2frontvideodata-time-frame}{%
\paragraph{\texorpdfstring{\texttt{unitree\_sdk2\_cpp.idl.go2.Go2FrontVideoData.time\_frame}}{unitree\_sdk2\_cpp.idl.go2.Go2FrontVideoData.time\_frame}}\label{unitree-sdk2-cpp-idl-go2-go2frontvideodata-time-frame}}

底层
\passthrough{\lstinline!unitree\_go::msg::dds\_::Go2FrontVideoData\_!}
的 \passthrough{\lstinline!time\_frame!} 字段。类型签名只说明 Python
表示；单位、数值范围、数组固定长度和枚举含义需查对应 IDL 头文件。
取值范围为 0 到 18446744073709551615。

\textbf{签名}

\begin{lstlisting}[language=Python]
@property
def time_frame(self) -> int

@time_frame.setter
def time_frame(self, value: int) -> None
\end{lstlisting}

\textbf{可用性与安全性}

\textbf{\passthrough{\lstinline!AVAILABLE!} /
\passthrough{\lstinline!VALUE\_TYPE!}}：当前绑定源码已实现该消息值操作。

\textbf{参数}

\begin{longtable}[]{@{}
  >{\raggedright\arraybackslash}p{(\columnwidth - 4\tabcolsep) * \real{0.18}}
  >{\raggedright\arraybackslash}p{(\columnwidth - 4\tabcolsep) * \real{0.20}}
  >{\raggedright\arraybackslash}p{(\columnwidth - 4\tabcolsep) * \real{0.62}}@{}}
\toprule
名称 & Python 类型 & 含义 \\
\midrule
\endhead
\passthrough{\lstinline!value!} & \passthrough{\lstinline!int!} &
要写入或传递的新值，必须符合该参数的 Python 类型以及底层 C++
范围约束。 \\
\bottomrule
\end{longtable}

\textbf{返回值}

\begin{longtable}[]{@{}
  >{\raggedright\arraybackslash}p{(\columnwidth - 4\tabcolsep) * \real{0.18}}
  >{\raggedright\arraybackslash}p{(\columnwidth - 4\tabcolsep) * \real{0.20}}
  >{\raggedright\arraybackslash}p{(\columnwidth - 4\tabcolsep) * \real{0.62}}@{}}
\toprule
位置 & 类型 & 含义 \\
\midrule
\endhead
getter 返回值 & \passthrough{\lstinline!int!} & 返回属性当前值。 \\
\bottomrule
\end{longtable}

\textbf{对应 C++}

\begin{itemize}
\tightlist
\item
  类：\passthrough{\lstinline!unitree\_go::msg::dds\_::Go2FrontVideoData\_!}
\item
  字段类型：\passthrough{\lstinline!uint64\_t!}
\item
  声明位置：\passthrough{\lstinline!include/unitree/idl/go2/Go2FrontVideoData\_.hpp:24!}
\end{itemize}

\textbf{用法}

\begin{lstlisting}[language=Python]
current_value = obj.time_frame
obj.time_frame = new_value
\end{lstlisting}

\hypertarget{unitree-sdk2-cpp-idl-go2-go2frontvideodata-video720p}{%
\paragraph{\texorpdfstring{\texttt{unitree\_sdk2\_cpp.idl.go2.Go2FrontVideoData.video720p}}{unitree\_sdk2\_cpp.idl.go2.Go2FrontVideoData.video720p}}\label{unitree-sdk2-cpp-idl-go2-go2frontvideodata-video720p}}

底层
\passthrough{\lstinline!unitree\_go::msg::dds\_::Go2FrontVideoData\_!}
的 \passthrough{\lstinline!video720p!} 字段。类型签名只说明 Python
表示；单位、数值范围、数组固定长度和枚举含义需查对应 IDL 头文件。
底层是可变长度 vector；元素约束：取值范围为 0 到 255。

\textbf{签名}

\begin{lstlisting}[language=Python]
@property
def video720p(self) -> list[int]

@video720p.setter
def video720p(self, value: Sequence[int]) -> None
\end{lstlisting}

\textbf{可用性与安全性}

\textbf{\passthrough{\lstinline!AVAILABLE!} /
\passthrough{\lstinline!VALUE\_TYPE!}}：当前绑定源码已实现该消息值操作。

\textbf{参数}

\begin{longtable}[]{@{}
  >{\raggedright\arraybackslash}p{(\columnwidth - 4\tabcolsep) * \real{0.18}}
  >{\raggedright\arraybackslash}p{(\columnwidth - 4\tabcolsep) * \real{0.20}}
  >{\raggedright\arraybackslash}p{(\columnwidth - 4\tabcolsep) * \real{0.62}}@{}}
\toprule
名称 & Python 类型 & 含义 \\
\midrule
\endhead
\passthrough{\lstinline!value!} &
\passthrough{\lstinline!Sequence[int]!} &
要写入或传递的新值，必须符合该参数的 Python 类型以及底层 C++
范围约束。 \\
\bottomrule
\end{longtable}

\textbf{返回值}

\begin{longtable}[]{@{}
  >{\raggedright\arraybackslash}p{(\columnwidth - 4\tabcolsep) * \real{0.18}}
  >{\raggedright\arraybackslash}p{(\columnwidth - 4\tabcolsep) * \real{0.20}}
  >{\raggedright\arraybackslash}p{(\columnwidth - 4\tabcolsep) * \real{0.62}}@{}}
\toprule
位置 & 类型 & 含义 \\
\midrule
\endhead
getter 返回值 & \passthrough{\lstinline!list[int]!} &
返回属性当前值。数组、vector 和嵌套消息采用复制语义。 \\
\bottomrule
\end{longtable}

\textbf{对应 C++}

\begin{itemize}
\tightlist
\item
  类：\passthrough{\lstinline!unitree\_go::msg::dds\_::Go2FrontVideoData\_!}
\item
  字段类型：\passthrough{\lstinline!std::vector<uint8\_t>!}
\item
  声明位置：\passthrough{\lstinline!include/unitree/idl/go2/Go2FrontVideoData\_.hpp:25!}
\end{itemize}

\textbf{用法}

\begin{lstlisting}[language=Python]
current_value = obj.video720p
obj.video720p = new_value
\end{lstlisting}

\begin{quote}
{[}!TIP{]} 读取容器或嵌套值后应遵循``读取、修改、重新赋值''。只修改
getter 返回的副本不会可靠地更新原 C++ 消息。
\end{quote}

\hypertarget{unitree-sdk2-cpp-idl-go2-go2frontvideodata-video360p}{%
\paragraph{\texorpdfstring{\texttt{unitree\_sdk2\_cpp.idl.go2.Go2FrontVideoData.video360p}}{unitree\_sdk2\_cpp.idl.go2.Go2FrontVideoData.video360p}}\label{unitree-sdk2-cpp-idl-go2-go2frontvideodata-video360p}}

底层
\passthrough{\lstinline!unitree\_go::msg::dds\_::Go2FrontVideoData\_!}
的 \passthrough{\lstinline!video360p!} 字段。类型签名只说明 Python
表示；单位、数值范围、数组固定长度和枚举含义需查对应 IDL 头文件。
底层是可变长度 vector；元素约束：取值范围为 0 到 255。

\textbf{签名}

\begin{lstlisting}[language=Python]
@property
def video360p(self) -> list[int]

@video360p.setter
def video360p(self, value: Sequence[int]) -> None
\end{lstlisting}

\textbf{可用性与安全性}

\textbf{\passthrough{\lstinline!AVAILABLE!} /
\passthrough{\lstinline!VALUE\_TYPE!}}：当前绑定源码已实现该消息值操作。

\textbf{参数}

\begin{longtable}[]{@{}
  >{\raggedright\arraybackslash}p{(\columnwidth - 4\tabcolsep) * \real{0.18}}
  >{\raggedright\arraybackslash}p{(\columnwidth - 4\tabcolsep) * \real{0.20}}
  >{\raggedright\arraybackslash}p{(\columnwidth - 4\tabcolsep) * \real{0.62}}@{}}
\toprule
名称 & Python 类型 & 含义 \\
\midrule
\endhead
\passthrough{\lstinline!value!} &
\passthrough{\lstinline!Sequence[int]!} &
要写入或传递的新值，必须符合该参数的 Python 类型以及底层 C++
范围约束。 \\
\bottomrule
\end{longtable}

\textbf{返回值}

\begin{longtable}[]{@{}
  >{\raggedright\arraybackslash}p{(\columnwidth - 4\tabcolsep) * \real{0.18}}
  >{\raggedright\arraybackslash}p{(\columnwidth - 4\tabcolsep) * \real{0.20}}
  >{\raggedright\arraybackslash}p{(\columnwidth - 4\tabcolsep) * \real{0.62}}@{}}
\toprule
位置 & 类型 & 含义 \\
\midrule
\endhead
getter 返回值 & \passthrough{\lstinline!list[int]!} &
返回属性当前值。数组、vector 和嵌套消息采用复制语义。 \\
\bottomrule
\end{longtable}

\textbf{对应 C++}

\begin{itemize}
\tightlist
\item
  类：\passthrough{\lstinline!unitree\_go::msg::dds\_::Go2FrontVideoData\_!}
\item
  字段类型：\passthrough{\lstinline!std::vector<uint8\_t>!}
\item
  声明位置：\passthrough{\lstinline!include/unitree/idl/go2/Go2FrontVideoData\_.hpp:26!}
\end{itemize}

\textbf{用法}

\begin{lstlisting}[language=Python]
current_value = obj.video360p
obj.video360p = new_value
\end{lstlisting}

\begin{quote}
{[}!TIP{]} 读取容器或嵌套值后应遵循``读取、修改、重新赋值''。只修改
getter 返回的副本不会可靠地更新原 C++ 消息。
\end{quote}

\hypertarget{unitree-sdk2-cpp-idl-go2-go2frontvideodata-video180p}{%
\paragraph{\texorpdfstring{\texttt{unitree\_sdk2\_cpp.idl.go2.Go2FrontVideoData.video180p}}{unitree\_sdk2\_cpp.idl.go2.Go2FrontVideoData.video180p}}\label{unitree-sdk2-cpp-idl-go2-go2frontvideodata-video180p}}

底层
\passthrough{\lstinline!unitree\_go::msg::dds\_::Go2FrontVideoData\_!}
的 \passthrough{\lstinline!video180p!} 字段。类型签名只说明 Python
表示；单位、数值范围、数组固定长度和枚举含义需查对应 IDL 头文件。
底层是可变长度 vector；元素约束：取值范围为 0 到 255。

\textbf{签名}

\begin{lstlisting}[language=Python]
@property
def video180p(self) -> list[int]

@video180p.setter
def video180p(self, value: Sequence[int]) -> None
\end{lstlisting}

\textbf{可用性与安全性}

\textbf{\passthrough{\lstinline!AVAILABLE!} /
\passthrough{\lstinline!VALUE\_TYPE!}}：当前绑定源码已实现该消息值操作。

\textbf{参数}

\begin{longtable}[]{@{}
  >{\raggedright\arraybackslash}p{(\columnwidth - 4\tabcolsep) * \real{0.18}}
  >{\raggedright\arraybackslash}p{(\columnwidth - 4\tabcolsep) * \real{0.20}}
  >{\raggedright\arraybackslash}p{(\columnwidth - 4\tabcolsep) * \real{0.62}}@{}}
\toprule
名称 & Python 类型 & 含义 \\
\midrule
\endhead
\passthrough{\lstinline!value!} &
\passthrough{\lstinline!Sequence[int]!} &
要写入或传递的新值，必须符合该参数的 Python 类型以及底层 C++
范围约束。 \\
\bottomrule
\end{longtable}

\textbf{返回值}

\begin{longtable}[]{@{}
  >{\raggedright\arraybackslash}p{(\columnwidth - 4\tabcolsep) * \real{0.18}}
  >{\raggedright\arraybackslash}p{(\columnwidth - 4\tabcolsep) * \real{0.20}}
  >{\raggedright\arraybackslash}p{(\columnwidth - 4\tabcolsep) * \real{0.62}}@{}}
\toprule
位置 & 类型 & 含义 \\
\midrule
\endhead
getter 返回值 & \passthrough{\lstinline!list[int]!} &
返回属性当前值。数组、vector 和嵌套消息采用复制语义。 \\
\bottomrule
\end{longtable}

\textbf{对应 C++}

\begin{itemize}
\tightlist
\item
  类：\passthrough{\lstinline!unitree\_go::msg::dds\_::Go2FrontVideoData\_!}
\item
  字段类型：\passthrough{\lstinline!std::vector<uint8\_t>!}
\item
  声明位置：\passthrough{\lstinline!include/unitree/idl/go2/Go2FrontVideoData\_.hpp:27!}
\end{itemize}

\textbf{用法}

\begin{lstlisting}[language=Python]
current_value = obj.video180p
obj.video180p = new_value
\end{lstlisting}

\begin{quote}
{[}!TIP{]} 读取容器或嵌套值后应遵循``读取、修改、重新赋值''。只修改
getter 返回的副本不会可靠地更新原 C++ 消息。
\end{quote}

\hypertarget{unitree-sdk2-cpp-idl-go2-heightmap}{%
\subsubsection{\texorpdfstring{\texttt{unitree\_sdk2\_cpp.idl.go2.HeightMap}}{unitree\_sdk2\_cpp.idl.go2.HeightMap}}\label{unitree-sdk2-cpp-idl-go2-heightmap}}

可由 Python 读写的 DDS/IDL 消息值类型。字段采用复制语义。

\textbf{导入}

\begin{lstlisting}[language=Python]
from unitree_sdk2_cpp.idl.go2 import HeightMap
\end{lstlisting}

\textbf{方法索引}

\begin{longtable}[]{@{}
  >{\raggedright\arraybackslash}p{(\columnwidth - 6\tabcolsep) * \real{0.18}}
  >{\raggedright\arraybackslash}p{(\columnwidth - 6\tabcolsep) * \real{0.24}}
  >{\raggedright\arraybackslash}p{(\columnwidth - 6\tabcolsep) * \real{0.18}}
  >{\raggedright\arraybackslash}p{(\columnwidth - 6\tabcolsep) * \real{0.40}}@{}}
\toprule
方法 & Python 签名 & 状态 & 安全分类 \\
\midrule
\endhead
\protect\hyperlink{unitree-sdk2-cpp-idl-go2-heightmap-dunder-init-1}{\passthrough{\lstinline!\_\_init\_\_!}}
& \passthrough{\lstinline!def \_\_init\_\_(self) -> None!} &
\passthrough{\lstinline!AVAILABLE!} &
\passthrough{\lstinline!VALUE\_TYPE!} \\
\protect\hyperlink{unitree-sdk2-cpp-idl-go2-heightmap-dunder-eq-1}{\passthrough{\lstinline!\_\_eq\_\_!}}
& \passthrough{\lstinline!def \_\_eq\_\_(self, other: object) -> bool!}
& \passthrough{\lstinline!AVAILABLE!} &
\passthrough{\lstinline!VALUE\_TYPE!} \\
\protect\hyperlink{unitree-sdk2-cpp-idl-go2-heightmap-dunder-ne-1}{\passthrough{\lstinline!\_\_ne\_\_!}}
& \passthrough{\lstinline!def \_\_ne\_\_(self, other: object) -> bool!}
& \passthrough{\lstinline!AVAILABLE!} &
\passthrough{\lstinline!VALUE\_TYPE!} \\
\bottomrule
\end{longtable}

\textbf{属性索引}

\begin{longtable}[]{@{}
  >{\raggedright\arraybackslash}p{(\columnwidth - 6\tabcolsep) * \real{0.18}}
  >{\raggedright\arraybackslash}p{(\columnwidth - 6\tabcolsep) * \real{0.24}}
  >{\raggedright\arraybackslash}p{(\columnwidth - 6\tabcolsep) * \real{0.18}}
  >{\raggedright\arraybackslash}p{(\columnwidth - 6\tabcolsep) * \real{0.40}}@{}}
\toprule
属性 & 读取类型 & 写入类型 & 状态 \\
\midrule
\endhead
\protect\hyperlink{unitree-sdk2-cpp-idl-go2-heightmap-stamp}{\passthrough{\lstinline!stamp!}}
& \passthrough{\lstinline!float!} & \passthrough{\lstinline!float!} &
\passthrough{\lstinline!AVAILABLE!} \\
\protect\hyperlink{unitree-sdk2-cpp-idl-go2-heightmap-frame-id}{\passthrough{\lstinline!frame\_id!}}
& \passthrough{\lstinline!str!} & \passthrough{\lstinline!str!} &
\passthrough{\lstinline!AVAILABLE!} \\
\protect\hyperlink{unitree-sdk2-cpp-idl-go2-heightmap-resolution}{\passthrough{\lstinline!resolution!}}
& \passthrough{\lstinline!float!} & \passthrough{\lstinline!float!} &
\passthrough{\lstinline!AVAILABLE!} \\
\protect\hyperlink{unitree-sdk2-cpp-idl-go2-heightmap-width}{\passthrough{\lstinline!width!}}
& \passthrough{\lstinline!int!} & \passthrough{\lstinline!int!} &
\passthrough{\lstinline!AVAILABLE!} \\
\protect\hyperlink{unitree-sdk2-cpp-idl-go2-heightmap-height}{\passthrough{\lstinline!height!}}
& \passthrough{\lstinline!int!} & \passthrough{\lstinline!int!} &
\passthrough{\lstinline!AVAILABLE!} \\
\protect\hyperlink{unitree-sdk2-cpp-idl-go2-heightmap-origin}{\passthrough{\lstinline!origin!}}
& \passthrough{\lstinline!list[float]!} &
\passthrough{\lstinline!Sequence[float]!} &
\passthrough{\lstinline!AVAILABLE!} \\
\protect\hyperlink{unitree-sdk2-cpp-idl-go2-heightmap-data}{\passthrough{\lstinline!data!}}
& \passthrough{\lstinline!list[float]!} &
\passthrough{\lstinline!Sequence[float]!} &
\passthrough{\lstinline!AVAILABLE!} \\
\bottomrule
\end{longtable}

\hypertarget{unitree-sdk2-cpp-idl-go2-heightmap-dunder-init-1}{%
\paragraph{\texorpdfstring{\texttt{unitree\_sdk2\_cpp.idl.go2.HeightMap.\_\_init\_\_}}{unitree\_sdk2\_cpp.idl.go2.HeightMap.\_\_init\_\_}}\label{unitree-sdk2-cpp-idl-go2-heightmap-dunder-init-1}}

初始化 \passthrough{\lstinline!HeightMap!}
实例。是否能实际构造取决于下方可用性状态。

\textbf{签名}

\begin{lstlisting}[language=Python]
def __init__(self) -> None
\end{lstlisting}

\textbf{可用性与安全性}

\textbf{\passthrough{\lstinline!AVAILABLE!} /
\passthrough{\lstinline!VALUE\_TYPE!}}：当前绑定源码已实现该消息值操作。

\textbf{参数}

无显式参数。实例方法中的 \passthrough{\lstinline!self!} 由 Python
自动传入。

\textbf{返回值}

\begin{longtable}[]{@{}
  >{\raggedright\arraybackslash}p{(\columnwidth - 4\tabcolsep) * \real{0.18}}
  >{\raggedright\arraybackslash}p{(\columnwidth - 4\tabcolsep) * \real{0.20}}
  >{\raggedright\arraybackslash}p{(\columnwidth - 4\tabcolsep) * \real{0.62}}@{}}
\toprule
位置 & 类型 & 含义 \\
\midrule
\endhead
无 & \passthrough{\lstinline!None!} &
\passthrough{\lstinline!\_\_init\_\_!}
本身不返回值；调用类对象时，在构造函数可用的前提下得到该类实例。 \\
\bottomrule
\end{longtable}

\textbf{对应 C++}

\begin{itemize}
\tightlist
\item
  类：\passthrough{\lstinline!unitree\_go::msg::dds\_::HeightMap\_!}
\item
  签名：\passthrough{\lstinline!HeightMap\_()!}
\item
  绑定策略：\passthrough{\lstinline!未单独标注!}
\item
  声明位置：\passthrough{\lstinline!include/unitree/idl/go2/HeightMap\_.hpp:35!}
\end{itemize}

\textbf{说明}

\begin{itemize}
\tightlist
\item
  示例是局部调用片段；\passthrough{\lstinline!obj!}
  和参数变量表示已经按上层业务校验并准备好的对象。
\end{itemize}

\textbf{用法}

\begin{lstlisting}[language=Python]
value = HeightMap()
\end{lstlisting}

\hypertarget{unitree-sdk2-cpp-idl-go2-heightmap-dunder-eq-1}{%
\paragraph{\texorpdfstring{\texttt{unitree\_sdk2\_cpp.idl.go2.HeightMap.\_\_eq\_\_}}{unitree\_sdk2\_cpp.idl.go2.HeightMap.\_\_eq\_\_}}\label{unitree-sdk2-cpp-idl-go2-heightmap-dunder-eq-1}}

按底层 IDL 消息内容比较两个对象是否相等。

\textbf{签名}

\begin{lstlisting}[language=Python]
def __eq__(self, other: object) -> bool
\end{lstlisting}

\textbf{可用性与安全性}

\textbf{\passthrough{\lstinline!AVAILABLE!} /
\passthrough{\lstinline!VALUE\_TYPE!}}：当前绑定源码已实现该消息值操作。

\textbf{参数}

\begin{longtable}[]{@{}
  >{\raggedright\arraybackslash}p{(\columnwidth - 6\tabcolsep) * \real{0.18}}
  >{\raggedright\arraybackslash}p{(\columnwidth - 6\tabcolsep) * \real{0.24}}
  >{\raggedright\arraybackslash}p{(\columnwidth - 6\tabcolsep) * \real{0.18}}
  >{\raggedright\arraybackslash}p{(\columnwidth - 6\tabcolsep) * \real{0.40}}@{}}
\toprule
名称 & Python 类型 & 默认值 & 含义 \\
\midrule
\endhead
\passthrough{\lstinline!other!} & \passthrough{\lstinline!object!} &
必填 & 用于比较的另一个 Python 对象。类型不兼容时比较结果通常为
\passthrough{\lstinline!False!}。 \\
\bottomrule
\end{longtable}

\textbf{返回值}

\begin{longtable}[]{@{}
  >{\raggedright\arraybackslash}p{(\columnwidth - 4\tabcolsep) * \real{0.18}}
  >{\raggedright\arraybackslash}p{(\columnwidth - 4\tabcolsep) * \real{0.20}}
  >{\raggedright\arraybackslash}p{(\columnwidth - 4\tabcolsep) * \real{0.62}}@{}}
\toprule
位置 & 类型 & 含义 \\
\midrule
\endhead
返回值 & \passthrough{\lstinline!bool!} & 返回
\passthrough{\lstinline!bool!}。更精确的业务含义见该方法用途、C++
签名和目标型号协议。 \\
\bottomrule
\end{longtable}

\textbf{对应 C++}

\begin{itemize}
\tightlist
\item
  类：\passthrough{\lstinline!unitree\_go::msg::dds\_::HeightMap\_!}
\item
  签名：\passthrough{\lstinline!operator==(const value \&) const!}
\item
  绑定策略：\passthrough{\lstinline!未单独标注!}
\item
  声明位置：由 pybind11 手工绑定或生成注册表提供；manifest
  未记录独立头文件行号。
\end{itemize}

\textbf{说明}

\begin{itemize}
\tightlist
\item
  示例是局部调用片段；\passthrough{\lstinline!obj!}
  和参数变量表示已经按上层业务校验并准备好的对象。
\end{itemize}

\textbf{用法}

\begin{lstlisting}[language=Python]
same = left == right
\end{lstlisting}

\hypertarget{unitree-sdk2-cpp-idl-go2-heightmap-dunder-ne-1}{%
\paragraph{\texorpdfstring{\texttt{unitree\_sdk2\_cpp.idl.go2.HeightMap.\_\_ne\_\_}}{unitree\_sdk2\_cpp.idl.go2.HeightMap.\_\_ne\_\_}}\label{unitree-sdk2-cpp-idl-go2-heightmap-dunder-ne-1}}

按底层 IDL 消息内容比较两个对象是否不相等。

\textbf{签名}

\begin{lstlisting}[language=Python]
def __ne__(self, other: object) -> bool
\end{lstlisting}

\textbf{可用性与安全性}

\textbf{\passthrough{\lstinline!AVAILABLE!} /
\passthrough{\lstinline!VALUE\_TYPE!}}：当前绑定源码已实现该消息值操作。

\textbf{参数}

\begin{longtable}[]{@{}
  >{\raggedright\arraybackslash}p{(\columnwidth - 6\tabcolsep) * \real{0.18}}
  >{\raggedright\arraybackslash}p{(\columnwidth - 6\tabcolsep) * \real{0.24}}
  >{\raggedright\arraybackslash}p{(\columnwidth - 6\tabcolsep) * \real{0.18}}
  >{\raggedright\arraybackslash}p{(\columnwidth - 6\tabcolsep) * \real{0.40}}@{}}
\toprule
名称 & Python 类型 & 默认值 & 含义 \\
\midrule
\endhead
\passthrough{\lstinline!other!} & \passthrough{\lstinline!object!} &
必填 & 用于比较的另一个 Python 对象。类型不兼容时比较结果通常为
\passthrough{\lstinline!False!}。 \\
\bottomrule
\end{longtable}

\textbf{返回值}

\begin{longtable}[]{@{}
  >{\raggedright\arraybackslash}p{(\columnwidth - 4\tabcolsep) * \real{0.18}}
  >{\raggedright\arraybackslash}p{(\columnwidth - 4\tabcolsep) * \real{0.20}}
  >{\raggedright\arraybackslash}p{(\columnwidth - 4\tabcolsep) * \real{0.62}}@{}}
\toprule
位置 & 类型 & 含义 \\
\midrule
\endhead
返回值 & \passthrough{\lstinline!bool!} & 返回
\passthrough{\lstinline!bool!}。更精确的业务含义见该方法用途、C++
签名和目标型号协议。 \\
\bottomrule
\end{longtable}

\textbf{对应 C++}

\begin{itemize}
\tightlist
\item
  类：\passthrough{\lstinline!unitree\_go::msg::dds\_::HeightMap\_!}
\item
  签名：\passthrough{\lstinline"operator!=(const value \&) const"}
\item
  绑定策略：\passthrough{\lstinline!未单独标注!}
\item
  声明位置：由 pybind11 手工绑定或生成注册表提供；manifest
  未记录独立头文件行号。
\end{itemize}

\textbf{说明}

\begin{itemize}
\tightlist
\item
  示例是局部调用片段；\passthrough{\lstinline!obj!}
  和参数变量表示已经按上层业务校验并准备好的对象。
\end{itemize}

\textbf{用法}

\begin{lstlisting}[language=Python]
different = left != right
\end{lstlisting}

\hypertarget{unitree-sdk2-cpp-idl-go2-heightmap-stamp}{%
\paragraph{\texorpdfstring{\texttt{unitree\_sdk2\_cpp.idl.go2.HeightMap.stamp}}{unitree\_sdk2\_cpp.idl.go2.HeightMap.stamp}}\label{unitree-sdk2-cpp-idl-go2-heightmap-stamp}}

底层 \passthrough{\lstinline!unitree\_go::msg::dds\_::HeightMap\_!} 的
\passthrough{\lstinline!stamp!} 字段。类型签名只说明 Python
表示；单位、数值范围、数组固定长度和枚举含义需查对应 IDL 头文件。
底层为浮点数；类型本身不说明单位、坐标系或业务安全范围。

\textbf{签名}

\begin{lstlisting}[language=Python]
@property
def stamp(self) -> float

@stamp.setter
def stamp(self, value: float) -> None
\end{lstlisting}

\textbf{可用性与安全性}

\textbf{\passthrough{\lstinline!AVAILABLE!} /
\passthrough{\lstinline!VALUE\_TYPE!}}：当前绑定源码已实现该消息值操作。

\textbf{参数}

\begin{longtable}[]{@{}
  >{\raggedright\arraybackslash}p{(\columnwidth - 4\tabcolsep) * \real{0.18}}
  >{\raggedright\arraybackslash}p{(\columnwidth - 4\tabcolsep) * \real{0.20}}
  >{\raggedright\arraybackslash}p{(\columnwidth - 4\tabcolsep) * \real{0.62}}@{}}
\toprule
名称 & Python 类型 & 含义 \\
\midrule
\endhead
\passthrough{\lstinline!value!} & \passthrough{\lstinline!float!} &
要写入或传递的新值，必须符合该参数的 Python 类型以及底层 C++
范围约束。 \\
\bottomrule
\end{longtable}

\textbf{返回值}

\begin{longtable}[]{@{}
  >{\raggedright\arraybackslash}p{(\columnwidth - 4\tabcolsep) * \real{0.18}}
  >{\raggedright\arraybackslash}p{(\columnwidth - 4\tabcolsep) * \real{0.20}}
  >{\raggedright\arraybackslash}p{(\columnwidth - 4\tabcolsep) * \real{0.62}}@{}}
\toprule
位置 & 类型 & 含义 \\
\midrule
\endhead
getter 返回值 & \passthrough{\lstinline!float!} & 返回属性当前值。 \\
\bottomrule
\end{longtable}

\textbf{对应 C++}

\begin{itemize}
\tightlist
\item
  类：\passthrough{\lstinline!unitree\_go::msg::dds\_::HeightMap\_!}
\item
  字段类型：\passthrough{\lstinline!double!}
\item
  声明位置：\passthrough{\lstinline!include/unitree/idl/go2/HeightMap\_.hpp:26!}
\end{itemize}

\textbf{用法}

\begin{lstlisting}[language=Python]
current_value = obj.stamp
obj.stamp = new_value
\end{lstlisting}

\hypertarget{unitree-sdk2-cpp-idl-go2-heightmap-frame-id}{%
\paragraph{\texorpdfstring{\texttt{unitree\_sdk2\_cpp.idl.go2.HeightMap.frame\_id}}{unitree\_sdk2\_cpp.idl.go2.HeightMap.frame\_id}}\label{unitree-sdk2-cpp-idl-go2-heightmap-frame-id}}

底层 \passthrough{\lstinline!unitree\_go::msg::dds\_::HeightMap\_!} 的
\passthrough{\lstinline!frame\_id!} 字段。类型签名只说明 Python
表示；单位、数值范围、数组固定长度和枚举含义需查对应 IDL 头文件。
底层为字符串；长度、编码和允许值由具体协议决定。

\textbf{签名}

\begin{lstlisting}[language=Python]
@property
def frame_id(self) -> str

@frame_id.setter
def frame_id(self, value: str) -> None
\end{lstlisting}

\textbf{可用性与安全性}

\textbf{\passthrough{\lstinline!AVAILABLE!} /
\passthrough{\lstinline!VALUE\_TYPE!}}：当前绑定源码已实现该消息值操作。

\textbf{参数}

\begin{longtable}[]{@{}
  >{\raggedright\arraybackslash}p{(\columnwidth - 4\tabcolsep) * \real{0.18}}
  >{\raggedright\arraybackslash}p{(\columnwidth - 4\tabcolsep) * \real{0.20}}
  >{\raggedright\arraybackslash}p{(\columnwidth - 4\tabcolsep) * \real{0.62}}@{}}
\toprule
名称 & Python 类型 & 含义 \\
\midrule
\endhead
\passthrough{\lstinline!value!} & \passthrough{\lstinline!str!} &
要写入或传递的新值，必须符合该参数的 Python 类型以及底层 C++
范围约束。 \\
\bottomrule
\end{longtable}

\textbf{返回值}

\begin{longtable}[]{@{}
  >{\raggedright\arraybackslash}p{(\columnwidth - 4\tabcolsep) * \real{0.18}}
  >{\raggedright\arraybackslash}p{(\columnwidth - 4\tabcolsep) * \real{0.20}}
  >{\raggedright\arraybackslash}p{(\columnwidth - 4\tabcolsep) * \real{0.62}}@{}}
\toprule
位置 & 类型 & 含义 \\
\midrule
\endhead
getter 返回值 & \passthrough{\lstinline!str!} & 返回属性当前值。 \\
\bottomrule
\end{longtable}

\textbf{对应 C++}

\begin{itemize}
\tightlist
\item
  类：\passthrough{\lstinline!unitree\_go::msg::dds\_::HeightMap\_!}
\item
  字段类型：\passthrough{\lstinline!std::string!}
\item
  声明位置：\passthrough{\lstinline!include/unitree/idl/go2/HeightMap\_.hpp:27!}
\end{itemize}

\textbf{用法}

\begin{lstlisting}[language=Python]
current_value = obj.frame_id
obj.frame_id = new_value
\end{lstlisting}

\hypertarget{unitree-sdk2-cpp-idl-go2-heightmap-resolution}{%
\paragraph{\texorpdfstring{\texttt{unitree\_sdk2\_cpp.idl.go2.HeightMap.resolution}}{unitree\_sdk2\_cpp.idl.go2.HeightMap.resolution}}\label{unitree-sdk2-cpp-idl-go2-heightmap-resolution}}

底层 \passthrough{\lstinline!unitree\_go::msg::dds\_::HeightMap\_!} 的
\passthrough{\lstinline!resolution!} 字段。类型签名只说明 Python
表示；单位、数值范围、数组固定长度和枚举含义需查对应 IDL 头文件。
底层为浮点数；类型本身不说明单位、坐标系或业务安全范围。

\textbf{签名}

\begin{lstlisting}[language=Python]
@property
def resolution(self) -> float

@resolution.setter
def resolution(self, value: float) -> None
\end{lstlisting}

\textbf{可用性与安全性}

\textbf{\passthrough{\lstinline!AVAILABLE!} /
\passthrough{\lstinline!VALUE\_TYPE!}}：当前绑定源码已实现该消息值操作。

\textbf{参数}

\begin{longtable}[]{@{}
  >{\raggedright\arraybackslash}p{(\columnwidth - 4\tabcolsep) * \real{0.18}}
  >{\raggedright\arraybackslash}p{(\columnwidth - 4\tabcolsep) * \real{0.20}}
  >{\raggedright\arraybackslash}p{(\columnwidth - 4\tabcolsep) * \real{0.62}}@{}}
\toprule
名称 & Python 类型 & 含义 \\
\midrule
\endhead
\passthrough{\lstinline!value!} & \passthrough{\lstinline!float!} &
要写入或传递的新值，必须符合该参数的 Python 类型以及底层 C++
范围约束。 \\
\bottomrule
\end{longtable}

\textbf{返回值}

\begin{longtable}[]{@{}
  >{\raggedright\arraybackslash}p{(\columnwidth - 4\tabcolsep) * \real{0.18}}
  >{\raggedright\arraybackslash}p{(\columnwidth - 4\tabcolsep) * \real{0.20}}
  >{\raggedright\arraybackslash}p{(\columnwidth - 4\tabcolsep) * \real{0.62}}@{}}
\toprule
位置 & 类型 & 含义 \\
\midrule
\endhead
getter 返回值 & \passthrough{\lstinline!float!} & 返回属性当前值。 \\
\bottomrule
\end{longtable}

\textbf{对应 C++}

\begin{itemize}
\tightlist
\item
  类：\passthrough{\lstinline!unitree\_go::msg::dds\_::HeightMap\_!}
\item
  字段类型：\passthrough{\lstinline!float!}
\item
  声明位置：\passthrough{\lstinline!include/unitree/idl/go2/HeightMap\_.hpp:28!}
\end{itemize}

\textbf{用法}

\begin{lstlisting}[language=Python]
current_value = obj.resolution
obj.resolution = new_value
\end{lstlisting}

\hypertarget{unitree-sdk2-cpp-idl-go2-heightmap-width}{%
\paragraph{\texorpdfstring{\texttt{unitree\_sdk2\_cpp.idl.go2.HeightMap.width}}{unitree\_sdk2\_cpp.idl.go2.HeightMap.width}}\label{unitree-sdk2-cpp-idl-go2-heightmap-width}}

底层 \passthrough{\lstinline!unitree\_go::msg::dds\_::HeightMap\_!} 的
\passthrough{\lstinline!width!} 字段。类型签名只说明 Python
表示；单位、数值范围、数组固定长度和枚举含义需查对应 IDL 头文件。
取值范围为 0 到 4294967295。

\textbf{签名}

\begin{lstlisting}[language=Python]
@property
def width(self) -> int

@width.setter
def width(self, value: int) -> None
\end{lstlisting}

\textbf{可用性与安全性}

\textbf{\passthrough{\lstinline!AVAILABLE!} /
\passthrough{\lstinline!VALUE\_TYPE!}}：当前绑定源码已实现该消息值操作。

\textbf{参数}

\begin{longtable}[]{@{}
  >{\raggedright\arraybackslash}p{(\columnwidth - 4\tabcolsep) * \real{0.18}}
  >{\raggedright\arraybackslash}p{(\columnwidth - 4\tabcolsep) * \real{0.20}}
  >{\raggedright\arraybackslash}p{(\columnwidth - 4\tabcolsep) * \real{0.62}}@{}}
\toprule
名称 & Python 类型 & 含义 \\
\midrule
\endhead
\passthrough{\lstinline!value!} & \passthrough{\lstinline!int!} &
要写入或传递的新值，必须符合该参数的 Python 类型以及底层 C++
范围约束。 \\
\bottomrule
\end{longtable}

\textbf{返回值}

\begin{longtable}[]{@{}
  >{\raggedright\arraybackslash}p{(\columnwidth - 4\tabcolsep) * \real{0.18}}
  >{\raggedright\arraybackslash}p{(\columnwidth - 4\tabcolsep) * \real{0.20}}
  >{\raggedright\arraybackslash}p{(\columnwidth - 4\tabcolsep) * \real{0.62}}@{}}
\toprule
位置 & 类型 & 含义 \\
\midrule
\endhead
getter 返回值 & \passthrough{\lstinline!int!} & 返回属性当前值。 \\
\bottomrule
\end{longtable}

\textbf{对应 C++}

\begin{itemize}
\tightlist
\item
  类：\passthrough{\lstinline!unitree\_go::msg::dds\_::HeightMap\_!}
\item
  字段类型：\passthrough{\lstinline!uint32\_t!}
\item
  声明位置：\passthrough{\lstinline!include/unitree/idl/go2/HeightMap\_.hpp:29!}
\end{itemize}

\textbf{用法}

\begin{lstlisting}[language=Python]
current_value = obj.width
obj.width = new_value
\end{lstlisting}

\hypertarget{unitree-sdk2-cpp-idl-go2-heightmap-height}{%
\paragraph{\texorpdfstring{\texttt{unitree\_sdk2\_cpp.idl.go2.HeightMap.height}}{unitree\_sdk2\_cpp.idl.go2.HeightMap.height}}\label{unitree-sdk2-cpp-idl-go2-heightmap-height}}

底层 \passthrough{\lstinline!unitree\_go::msg::dds\_::HeightMap\_!} 的
\passthrough{\lstinline!height!} 字段。类型签名只说明 Python
表示；单位、数值范围、数组固定长度和枚举含义需查对应 IDL 头文件。
取值范围为 0 到 4294967295。

\textbf{签名}

\begin{lstlisting}[language=Python]
@property
def height(self) -> int

@height.setter
def height(self, value: int) -> None
\end{lstlisting}

\textbf{可用性与安全性}

\textbf{\passthrough{\lstinline!AVAILABLE!} /
\passthrough{\lstinline!VALUE\_TYPE!}}：当前绑定源码已实现该消息值操作。

\textbf{参数}

\begin{longtable}[]{@{}
  >{\raggedright\arraybackslash}p{(\columnwidth - 4\tabcolsep) * \real{0.18}}
  >{\raggedright\arraybackslash}p{(\columnwidth - 4\tabcolsep) * \real{0.20}}
  >{\raggedright\arraybackslash}p{(\columnwidth - 4\tabcolsep) * \real{0.62}}@{}}
\toprule
名称 & Python 类型 & 含义 \\
\midrule
\endhead
\passthrough{\lstinline!value!} & \passthrough{\lstinline!int!} &
要写入或传递的新值，必须符合该参数的 Python 类型以及底层 C++
范围约束。 \\
\bottomrule
\end{longtable}

\textbf{返回值}

\begin{longtable}[]{@{}
  >{\raggedright\arraybackslash}p{(\columnwidth - 4\tabcolsep) * \real{0.18}}
  >{\raggedright\arraybackslash}p{(\columnwidth - 4\tabcolsep) * \real{0.20}}
  >{\raggedright\arraybackslash}p{(\columnwidth - 4\tabcolsep) * \real{0.62}}@{}}
\toprule
位置 & 类型 & 含义 \\
\midrule
\endhead
getter 返回值 & \passthrough{\lstinline!int!} & 返回属性当前值。 \\
\bottomrule
\end{longtable}

\textbf{对应 C++}

\begin{itemize}
\tightlist
\item
  类：\passthrough{\lstinline!unitree\_go::msg::dds\_::HeightMap\_!}
\item
  字段类型：\passthrough{\lstinline!uint32\_t!}
\item
  声明位置：\passthrough{\lstinline!include/unitree/idl/go2/HeightMap\_.hpp:30!}
\end{itemize}

\textbf{用法}

\begin{lstlisting}[language=Python]
current_value = obj.height
obj.height = new_value
\end{lstlisting}

\hypertarget{unitree-sdk2-cpp-idl-go2-heightmap-origin}{%
\paragraph{\texorpdfstring{\texttt{unitree\_sdk2\_cpp.idl.go2.HeightMap.origin}}{unitree\_sdk2\_cpp.idl.go2.HeightMap.origin}}\label{unitree-sdk2-cpp-idl-go2-heightmap-origin}}

底层 \passthrough{\lstinline!unitree\_go::msg::dds\_::HeightMap\_!} 的
\passthrough{\lstinline!origin!} 字段。类型签名只说明 Python
表示；单位、数值范围、数组固定长度和枚举含义需查对应 IDL 头文件。
底层是固定长度数组，写入序列必须正好包含 2
个元素；元素约束：底层为浮点数；类型本身不说明单位、坐标系或业务安全范围。

\textbf{签名}

\begin{lstlisting}[language=Python]
@property
def origin(self) -> list[float]

@origin.setter
def origin(self, value: Sequence[float]) -> None
\end{lstlisting}

\textbf{可用性与安全性}

\textbf{\passthrough{\lstinline!AVAILABLE!} /
\passthrough{\lstinline!VALUE\_TYPE!}}：当前绑定源码已实现该消息值操作。

\textbf{参数}

\begin{longtable}[]{@{}
  >{\raggedright\arraybackslash}p{(\columnwidth - 4\tabcolsep) * \real{0.18}}
  >{\raggedright\arraybackslash}p{(\columnwidth - 4\tabcolsep) * \real{0.20}}
  >{\raggedright\arraybackslash}p{(\columnwidth - 4\tabcolsep) * \real{0.62}}@{}}
\toprule
名称 & Python 类型 & 含义 \\
\midrule
\endhead
\passthrough{\lstinline!value!} &
\passthrough{\lstinline!Sequence[float]!} &
要写入或传递的新值，必须符合该参数的 Python 类型以及底层 C++
范围约束。 \\
\bottomrule
\end{longtable}

\textbf{返回值}

\begin{longtable}[]{@{}
  >{\raggedright\arraybackslash}p{(\columnwidth - 4\tabcolsep) * \real{0.18}}
  >{\raggedright\arraybackslash}p{(\columnwidth - 4\tabcolsep) * \real{0.20}}
  >{\raggedright\arraybackslash}p{(\columnwidth - 4\tabcolsep) * \real{0.62}}@{}}
\toprule
位置 & 类型 & 含义 \\
\midrule
\endhead
getter 返回值 & \passthrough{\lstinline!list[float]!} &
返回属性当前值。数组、vector 和嵌套消息采用复制语义。 \\
\bottomrule
\end{longtable}

\textbf{对应 C++}

\begin{itemize}
\tightlist
\item
  类：\passthrough{\lstinline!unitree\_go::msg::dds\_::HeightMap\_!}
\item
  字段类型：\passthrough{\lstinline!std::array<float, 2>!}
\item
  声明位置：\passthrough{\lstinline!include/unitree/idl/go2/HeightMap\_.hpp:31!}
\end{itemize}

\textbf{用法}

\begin{lstlisting}[language=Python]
current_value = obj.origin
obj.origin = new_value
\end{lstlisting}

\begin{quote}
{[}!TIP{]} 读取容器或嵌套值后应遵循``读取、修改、重新赋值''。只修改
getter 返回的副本不会可靠地更新原 C++ 消息。
\end{quote}

\hypertarget{unitree-sdk2-cpp-idl-go2-heightmap-data}{%
\paragraph{\texorpdfstring{\texttt{unitree\_sdk2\_cpp.idl.go2.HeightMap.data}}{unitree\_sdk2\_cpp.idl.go2.HeightMap.data}}\label{unitree-sdk2-cpp-idl-go2-heightmap-data}}

底层 \passthrough{\lstinline!unitree\_go::msg::dds\_::HeightMap\_!} 的
\passthrough{\lstinline!data!} 字段。类型签名只说明 Python
表示；单位、数值范围、数组固定长度和枚举含义需查对应 IDL 头文件。
底层是可变长度
vector；元素约束：底层为浮点数；类型本身不说明单位、坐标系或业务安全范围。

\textbf{签名}

\begin{lstlisting}[language=Python]
@property
def data(self) -> list[float]

@data.setter
def data(self, value: Sequence[float]) -> None
\end{lstlisting}

\textbf{可用性与安全性}

\textbf{\passthrough{\lstinline!AVAILABLE!} /
\passthrough{\lstinline!VALUE\_TYPE!}}：当前绑定源码已实现该消息值操作。

\textbf{参数}

\begin{longtable}[]{@{}
  >{\raggedright\arraybackslash}p{(\columnwidth - 4\tabcolsep) * \real{0.18}}
  >{\raggedright\arraybackslash}p{(\columnwidth - 4\tabcolsep) * \real{0.20}}
  >{\raggedright\arraybackslash}p{(\columnwidth - 4\tabcolsep) * \real{0.62}}@{}}
\toprule
名称 & Python 类型 & 含义 \\
\midrule
\endhead
\passthrough{\lstinline!value!} &
\passthrough{\lstinline!Sequence[float]!} &
要写入或传递的新值，必须符合该参数的 Python 类型以及底层 C++
范围约束。 \\
\bottomrule
\end{longtable}

\textbf{返回值}

\begin{longtable}[]{@{}
  >{\raggedright\arraybackslash}p{(\columnwidth - 4\tabcolsep) * \real{0.18}}
  >{\raggedright\arraybackslash}p{(\columnwidth - 4\tabcolsep) * \real{0.20}}
  >{\raggedright\arraybackslash}p{(\columnwidth - 4\tabcolsep) * \real{0.62}}@{}}
\toprule
位置 & 类型 & 含义 \\
\midrule
\endhead
getter 返回值 & \passthrough{\lstinline!list[float]!} &
返回属性当前值。数组、vector 和嵌套消息采用复制语义。 \\
\bottomrule
\end{longtable}

\textbf{对应 C++}

\begin{itemize}
\tightlist
\item
  类：\passthrough{\lstinline!unitree\_go::msg::dds\_::HeightMap\_!}
\item
  字段类型：\passthrough{\lstinline!std::vector<float>!}
\item
  声明位置：\passthrough{\lstinline!include/unitree/idl/go2/HeightMap\_.hpp:32!}
\end{itemize}

\textbf{用法}

\begin{lstlisting}[language=Python]
current_value = obj.data
obj.data = new_value
\end{lstlisting}

\begin{quote}
{[}!TIP{]} 读取容器或嵌套值后应遵循``读取、修改、重新赋值''。只修改
getter 返回的副本不会可靠地更新原 C++ 消息。
\end{quote}

\hypertarget{unitree-sdk2-cpp-idl-go2-imustate}{%
\subsubsection{\texorpdfstring{\texttt{unitree\_sdk2\_cpp.idl.go2.IMUState}}{unitree\_sdk2\_cpp.idl.go2.IMUState}}\label{unitree-sdk2-cpp-idl-go2-imustate}}

可由 Python 读写的 DDS/IDL 消息值类型。字段采用复制语义。

\textbf{导入}

\begin{lstlisting}[language=Python]
from unitree_sdk2_cpp.idl.go2 import IMUState
\end{lstlisting}

\textbf{方法索引}

\begin{longtable}[]{@{}
  >{\raggedright\arraybackslash}p{(\columnwidth - 6\tabcolsep) * \real{0.18}}
  >{\raggedright\arraybackslash}p{(\columnwidth - 6\tabcolsep) * \real{0.24}}
  >{\raggedright\arraybackslash}p{(\columnwidth - 6\tabcolsep) * \real{0.18}}
  >{\raggedright\arraybackslash}p{(\columnwidth - 6\tabcolsep) * \real{0.40}}@{}}
\toprule
方法 & Python 签名 & 状态 & 安全分类 \\
\midrule
\endhead
\protect\hyperlink{unitree-sdk2-cpp-idl-go2-imustate-dunder-init-1}{\passthrough{\lstinline!\_\_init\_\_!}}
& \passthrough{\lstinline!def \_\_init\_\_(self) -> None!} &
\passthrough{\lstinline!AVAILABLE!} &
\passthrough{\lstinline!VALUE\_TYPE!} \\
\protect\hyperlink{unitree-sdk2-cpp-idl-go2-imustate-dunder-eq-1}{\passthrough{\lstinline!\_\_eq\_\_!}}
& \passthrough{\lstinline!def \_\_eq\_\_(self, other: object) -> bool!}
& \passthrough{\lstinline!AVAILABLE!} &
\passthrough{\lstinline!VALUE\_TYPE!} \\
\protect\hyperlink{unitree-sdk2-cpp-idl-go2-imustate-dunder-ne-1}{\passthrough{\lstinline!\_\_ne\_\_!}}
& \passthrough{\lstinline!def \_\_ne\_\_(self, other: object) -> bool!}
& \passthrough{\lstinline!AVAILABLE!} &
\passthrough{\lstinline!VALUE\_TYPE!} \\
\bottomrule
\end{longtable}

\textbf{属性索引}

\begin{longtable}[]{@{}
  >{\raggedright\arraybackslash}p{(\columnwidth - 6\tabcolsep) * \real{0.18}}
  >{\raggedright\arraybackslash}p{(\columnwidth - 6\tabcolsep) * \real{0.24}}
  >{\raggedright\arraybackslash}p{(\columnwidth - 6\tabcolsep) * \real{0.18}}
  >{\raggedright\arraybackslash}p{(\columnwidth - 6\tabcolsep) * \real{0.40}}@{}}
\toprule
属性 & 读取类型 & 写入类型 & 状态 \\
\midrule
\endhead
\protect\hyperlink{unitree-sdk2-cpp-idl-go2-imustate-quaternion}{\passthrough{\lstinline!quaternion!}}
& \passthrough{\lstinline!list[float]!} &
\passthrough{\lstinline!Sequence[float]!} &
\passthrough{\lstinline!AVAILABLE!} \\
\protect\hyperlink{unitree-sdk2-cpp-idl-go2-imustate-gyroscope}{\passthrough{\lstinline!gyroscope!}}
& \passthrough{\lstinline!list[float]!} &
\passthrough{\lstinline!Sequence[float]!} &
\passthrough{\lstinline!AVAILABLE!} \\
\protect\hyperlink{unitree-sdk2-cpp-idl-go2-imustate-accelerometer}{\passthrough{\lstinline!accelerometer!}}
& \passthrough{\lstinline!list[float]!} &
\passthrough{\lstinline!Sequence[float]!} &
\passthrough{\lstinline!AVAILABLE!} \\
\protect\hyperlink{unitree-sdk2-cpp-idl-go2-imustate-rpy}{\passthrough{\lstinline!rpy!}}
& \passthrough{\lstinline!list[float]!} &
\passthrough{\lstinline!Sequence[float]!} &
\passthrough{\lstinline!AVAILABLE!} \\
\protect\hyperlink{unitree-sdk2-cpp-idl-go2-imustate-temperature}{\passthrough{\lstinline!temperature!}}
& \passthrough{\lstinline!int!} & \passthrough{\lstinline!int!} &
\passthrough{\lstinline!AVAILABLE!} \\
\bottomrule
\end{longtable}

\hypertarget{unitree-sdk2-cpp-idl-go2-imustate-dunder-init-1}{%
\paragraph{\texorpdfstring{\texttt{unitree\_sdk2\_cpp.idl.go2.IMUState.\_\_init\_\_}}{unitree\_sdk2\_cpp.idl.go2.IMUState.\_\_init\_\_}}\label{unitree-sdk2-cpp-idl-go2-imustate-dunder-init-1}}

初始化 \passthrough{\lstinline!IMUState!}
实例。是否能实际构造取决于下方可用性状态。

\textbf{签名}

\begin{lstlisting}[language=Python]
def __init__(self) -> None
\end{lstlisting}

\textbf{可用性与安全性}

\textbf{\passthrough{\lstinline!AVAILABLE!} /
\passthrough{\lstinline!VALUE\_TYPE!}}：当前绑定源码已实现该消息值操作。

\textbf{参数}

无显式参数。实例方法中的 \passthrough{\lstinline!self!} 由 Python
自动传入。

\textbf{返回值}

\begin{longtable}[]{@{}
  >{\raggedright\arraybackslash}p{(\columnwidth - 4\tabcolsep) * \real{0.18}}
  >{\raggedright\arraybackslash}p{(\columnwidth - 4\tabcolsep) * \real{0.20}}
  >{\raggedright\arraybackslash}p{(\columnwidth - 4\tabcolsep) * \real{0.62}}@{}}
\toprule
位置 & 类型 & 含义 \\
\midrule
\endhead
无 & \passthrough{\lstinline!None!} &
\passthrough{\lstinline!\_\_init\_\_!}
本身不返回值；调用类对象时，在构造函数可用的前提下得到该类实例。 \\
\bottomrule
\end{longtable}

\textbf{对应 C++}

\begin{itemize}
\tightlist
\item
  类：\passthrough{\lstinline!unitree\_go::msg::dds\_::IMUState\_!}
\item
  签名：\passthrough{\lstinline!IMUState\_()!}
\item
  绑定策略：\passthrough{\lstinline!未单独标注!}
\item
  声明位置：\passthrough{\lstinline!include/unitree/idl/go2/IMUState\_.hpp:31!}
\end{itemize}

\textbf{说明}

\begin{itemize}
\tightlist
\item
  示例是局部调用片段；\passthrough{\lstinline!obj!}
  和参数变量表示已经按上层业务校验并准备好的对象。
\end{itemize}

\textbf{用法}

\begin{lstlisting}[language=Python]
value = IMUState()
\end{lstlisting}

\hypertarget{unitree-sdk2-cpp-idl-go2-imustate-dunder-eq-1}{%
\paragraph{\texorpdfstring{\texttt{unitree\_sdk2\_cpp.idl.go2.IMUState.\_\_eq\_\_}}{unitree\_sdk2\_cpp.idl.go2.IMUState.\_\_eq\_\_}}\label{unitree-sdk2-cpp-idl-go2-imustate-dunder-eq-1}}

按底层 IDL 消息内容比较两个对象是否相等。

\textbf{签名}

\begin{lstlisting}[language=Python]
def __eq__(self, other: object) -> bool
\end{lstlisting}

\textbf{可用性与安全性}

\textbf{\passthrough{\lstinline!AVAILABLE!} /
\passthrough{\lstinline!VALUE\_TYPE!}}：当前绑定源码已实现该消息值操作。

\textbf{参数}

\begin{longtable}[]{@{}
  >{\raggedright\arraybackslash}p{(\columnwidth - 6\tabcolsep) * \real{0.18}}
  >{\raggedright\arraybackslash}p{(\columnwidth - 6\tabcolsep) * \real{0.24}}
  >{\raggedright\arraybackslash}p{(\columnwidth - 6\tabcolsep) * \real{0.18}}
  >{\raggedright\arraybackslash}p{(\columnwidth - 6\tabcolsep) * \real{0.40}}@{}}
\toprule
名称 & Python 类型 & 默认值 & 含义 \\
\midrule
\endhead
\passthrough{\lstinline!other!} & \passthrough{\lstinline!object!} &
必填 & 用于比较的另一个 Python 对象。类型不兼容时比较结果通常为
\passthrough{\lstinline!False!}。 \\
\bottomrule
\end{longtable}

\textbf{返回值}

\begin{longtable}[]{@{}
  >{\raggedright\arraybackslash}p{(\columnwidth - 4\tabcolsep) * \real{0.18}}
  >{\raggedright\arraybackslash}p{(\columnwidth - 4\tabcolsep) * \real{0.20}}
  >{\raggedright\arraybackslash}p{(\columnwidth - 4\tabcolsep) * \real{0.62}}@{}}
\toprule
位置 & 类型 & 含义 \\
\midrule
\endhead
返回值 & \passthrough{\lstinline!bool!} & 返回
\passthrough{\lstinline!bool!}。更精确的业务含义见该方法用途、C++
签名和目标型号协议。 \\
\bottomrule
\end{longtable}

\textbf{对应 C++}

\begin{itemize}
\tightlist
\item
  类：\passthrough{\lstinline!unitree\_go::msg::dds\_::IMUState\_!}
\item
  签名：\passthrough{\lstinline!operator==(const value \&) const!}
\item
  绑定策略：\passthrough{\lstinline!未单独标注!}
\item
  声明位置：由 pybind11 手工绑定或生成注册表提供；manifest
  未记录独立头文件行号。
\end{itemize}

\textbf{说明}

\begin{itemize}
\tightlist
\item
  示例是局部调用片段；\passthrough{\lstinline!obj!}
  和参数变量表示已经按上层业务校验并准备好的对象。
\end{itemize}

\textbf{用法}

\begin{lstlisting}[language=Python]
same = left == right
\end{lstlisting}

\hypertarget{unitree-sdk2-cpp-idl-go2-imustate-dunder-ne-1}{%
\paragraph{\texorpdfstring{\texttt{unitree\_sdk2\_cpp.idl.go2.IMUState.\_\_ne\_\_}}{unitree\_sdk2\_cpp.idl.go2.IMUState.\_\_ne\_\_}}\label{unitree-sdk2-cpp-idl-go2-imustate-dunder-ne-1}}

按底层 IDL 消息内容比较两个对象是否不相等。

\textbf{签名}

\begin{lstlisting}[language=Python]
def __ne__(self, other: object) -> bool
\end{lstlisting}

\textbf{可用性与安全性}

\textbf{\passthrough{\lstinline!AVAILABLE!} /
\passthrough{\lstinline!VALUE\_TYPE!}}：当前绑定源码已实现该消息值操作。

\textbf{参数}

\begin{longtable}[]{@{}
  >{\raggedright\arraybackslash}p{(\columnwidth - 6\tabcolsep) * \real{0.18}}
  >{\raggedright\arraybackslash}p{(\columnwidth - 6\tabcolsep) * \real{0.24}}
  >{\raggedright\arraybackslash}p{(\columnwidth - 6\tabcolsep) * \real{0.18}}
  >{\raggedright\arraybackslash}p{(\columnwidth - 6\tabcolsep) * \real{0.40}}@{}}
\toprule
名称 & Python 类型 & 默认值 & 含义 \\
\midrule
\endhead
\passthrough{\lstinline!other!} & \passthrough{\lstinline!object!} &
必填 & 用于比较的另一个 Python 对象。类型不兼容时比较结果通常为
\passthrough{\lstinline!False!}。 \\
\bottomrule
\end{longtable}

\textbf{返回值}

\begin{longtable}[]{@{}
  >{\raggedright\arraybackslash}p{(\columnwidth - 4\tabcolsep) * \real{0.18}}
  >{\raggedright\arraybackslash}p{(\columnwidth - 4\tabcolsep) * \real{0.20}}
  >{\raggedright\arraybackslash}p{(\columnwidth - 4\tabcolsep) * \real{0.62}}@{}}
\toprule
位置 & 类型 & 含义 \\
\midrule
\endhead
返回值 & \passthrough{\lstinline!bool!} & 返回
\passthrough{\lstinline!bool!}。更精确的业务含义见该方法用途、C++
签名和目标型号协议。 \\
\bottomrule
\end{longtable}

\textbf{对应 C++}

\begin{itemize}
\tightlist
\item
  类：\passthrough{\lstinline!unitree\_go::msg::dds\_::IMUState\_!}
\item
  签名：\passthrough{\lstinline"operator!=(const value \&) const"}
\item
  绑定策略：\passthrough{\lstinline!未单独标注!}
\item
  声明位置：由 pybind11 手工绑定或生成注册表提供；manifest
  未记录独立头文件行号。
\end{itemize}

\textbf{说明}

\begin{itemize}
\tightlist
\item
  示例是局部调用片段；\passthrough{\lstinline!obj!}
  和参数变量表示已经按上层业务校验并准备好的对象。
\end{itemize}

\textbf{用法}

\begin{lstlisting}[language=Python]
different = left != right
\end{lstlisting}

\hypertarget{unitree-sdk2-cpp-idl-go2-imustate-quaternion}{%
\paragraph{\texorpdfstring{\texttt{unitree\_sdk2\_cpp.idl.go2.IMUState.quaternion}}{unitree\_sdk2\_cpp.idl.go2.IMUState.quaternion}}\label{unitree-sdk2-cpp-idl-go2-imustate-quaternion}}

底层 \passthrough{\lstinline!unitree\_go::msg::dds\_::IMUState\_!} 的
\passthrough{\lstinline!quaternion!} 字段。类型签名只说明 Python
表示；单位、数值范围、数组固定长度和枚举含义需查对应 IDL 头文件。
底层是固定长度数组，写入序列必须正好包含 4
个元素；元素约束：底层为浮点数；类型本身不说明单位、坐标系或业务安全范围。

\textbf{签名}

\begin{lstlisting}[language=Python]
@property
def quaternion(self) -> list[float]

@quaternion.setter
def quaternion(self, value: Sequence[float]) -> None
\end{lstlisting}

\textbf{可用性与安全性}

\textbf{\passthrough{\lstinline!AVAILABLE!} /
\passthrough{\lstinline!VALUE\_TYPE!}}：当前绑定源码已实现该消息值操作。

\textbf{参数}

\begin{longtable}[]{@{}
  >{\raggedright\arraybackslash}p{(\columnwidth - 4\tabcolsep) * \real{0.18}}
  >{\raggedright\arraybackslash}p{(\columnwidth - 4\tabcolsep) * \real{0.20}}
  >{\raggedright\arraybackslash}p{(\columnwidth - 4\tabcolsep) * \real{0.62}}@{}}
\toprule
名称 & Python 类型 & 含义 \\
\midrule
\endhead
\passthrough{\lstinline!value!} &
\passthrough{\lstinline!Sequence[float]!} &
要写入或传递的新值，必须符合该参数的 Python 类型以及底层 C++
范围约束。 \\
\bottomrule
\end{longtable}

\textbf{返回值}

\begin{longtable}[]{@{}
  >{\raggedright\arraybackslash}p{(\columnwidth - 4\tabcolsep) * \real{0.18}}
  >{\raggedright\arraybackslash}p{(\columnwidth - 4\tabcolsep) * \real{0.20}}
  >{\raggedright\arraybackslash}p{(\columnwidth - 4\tabcolsep) * \real{0.62}}@{}}
\toprule
位置 & 类型 & 含义 \\
\midrule
\endhead
getter 返回值 & \passthrough{\lstinline!list[float]!} &
返回属性当前值。数组、vector 和嵌套消息采用复制语义。 \\
\bottomrule
\end{longtable}

\textbf{对应 C++}

\begin{itemize}
\tightlist
\item
  类：\passthrough{\lstinline!unitree\_go::msg::dds\_::IMUState\_!}
\item
  字段类型：\passthrough{\lstinline!std::array<float, 4>!}
\item
  声明位置：\passthrough{\lstinline!include/unitree/idl/go2/IMUState\_.hpp:24!}
\end{itemize}

\textbf{用法}

\begin{lstlisting}[language=Python]
current_value = obj.quaternion
obj.quaternion = new_value
\end{lstlisting}

\begin{quote}
{[}!TIP{]} 读取容器或嵌套值后应遵循``读取、修改、重新赋值''。只修改
getter 返回的副本不会可靠地更新原 C++ 消息。
\end{quote}

\hypertarget{unitree-sdk2-cpp-idl-go2-imustate-gyroscope}{%
\paragraph{\texorpdfstring{\texttt{unitree\_sdk2\_cpp.idl.go2.IMUState.gyroscope}}{unitree\_sdk2\_cpp.idl.go2.IMUState.gyroscope}}\label{unitree-sdk2-cpp-idl-go2-imustate-gyroscope}}

底层 \passthrough{\lstinline!unitree\_go::msg::dds\_::IMUState\_!} 的
\passthrough{\lstinline!gyroscope!} 字段。类型签名只说明 Python
表示；单位、数值范围、数组固定长度和枚举含义需查对应 IDL 头文件。
底层是固定长度数组，写入序列必须正好包含 3
个元素；元素约束：底层为浮点数；类型本身不说明单位、坐标系或业务安全范围。

\textbf{签名}

\begin{lstlisting}[language=Python]
@property
def gyroscope(self) -> list[float]

@gyroscope.setter
def gyroscope(self, value: Sequence[float]) -> None
\end{lstlisting}

\textbf{可用性与安全性}

\textbf{\passthrough{\lstinline!AVAILABLE!} /
\passthrough{\lstinline!VALUE\_TYPE!}}：当前绑定源码已实现该消息值操作。

\textbf{参数}

\begin{longtable}[]{@{}
  >{\raggedright\arraybackslash}p{(\columnwidth - 4\tabcolsep) * \real{0.18}}
  >{\raggedright\arraybackslash}p{(\columnwidth - 4\tabcolsep) * \real{0.20}}
  >{\raggedright\arraybackslash}p{(\columnwidth - 4\tabcolsep) * \real{0.62}}@{}}
\toprule
名称 & Python 类型 & 含义 \\
\midrule
\endhead
\passthrough{\lstinline!value!} &
\passthrough{\lstinline!Sequence[float]!} &
要写入或传递的新值，必须符合该参数的 Python 类型以及底层 C++
范围约束。 \\
\bottomrule
\end{longtable}

\textbf{返回值}

\begin{longtable}[]{@{}
  >{\raggedright\arraybackslash}p{(\columnwidth - 4\tabcolsep) * \real{0.18}}
  >{\raggedright\arraybackslash}p{(\columnwidth - 4\tabcolsep) * \real{0.20}}
  >{\raggedright\arraybackslash}p{(\columnwidth - 4\tabcolsep) * \real{0.62}}@{}}
\toprule
位置 & 类型 & 含义 \\
\midrule
\endhead
getter 返回值 & \passthrough{\lstinline!list[float]!} &
返回属性当前值。数组、vector 和嵌套消息采用复制语义。 \\
\bottomrule
\end{longtable}

\textbf{对应 C++}

\begin{itemize}
\tightlist
\item
  类：\passthrough{\lstinline!unitree\_go::msg::dds\_::IMUState\_!}
\item
  字段类型：\passthrough{\lstinline!std::array<float, 3>!}
\item
  声明位置：\passthrough{\lstinline!include/unitree/idl/go2/IMUState\_.hpp:25!}
\end{itemize}

\textbf{用法}

\begin{lstlisting}[language=Python]
current_value = obj.gyroscope
obj.gyroscope = new_value
\end{lstlisting}

\begin{quote}
{[}!TIP{]} 读取容器或嵌套值后应遵循``读取、修改、重新赋值''。只修改
getter 返回的副本不会可靠地更新原 C++ 消息。
\end{quote}

\hypertarget{unitree-sdk2-cpp-idl-go2-imustate-accelerometer}{%
\paragraph{\texorpdfstring{\texttt{unitree\_sdk2\_cpp.idl.go2.IMUState.accelerometer}}{unitree\_sdk2\_cpp.idl.go2.IMUState.accelerometer}}\label{unitree-sdk2-cpp-idl-go2-imustate-accelerometer}}

底层 \passthrough{\lstinline!unitree\_go::msg::dds\_::IMUState\_!} 的
\passthrough{\lstinline!accelerometer!} 字段。类型签名只说明 Python
表示；单位、数值范围、数组固定长度和枚举含义需查对应 IDL 头文件。
底层是固定长度数组，写入序列必须正好包含 3
个元素；元素约束：底层为浮点数；类型本身不说明单位、坐标系或业务安全范围。

\textbf{签名}

\begin{lstlisting}[language=Python]
@property
def accelerometer(self) -> list[float]

@accelerometer.setter
def accelerometer(self, value: Sequence[float]) -> None
\end{lstlisting}

\textbf{可用性与安全性}

\textbf{\passthrough{\lstinline!AVAILABLE!} /
\passthrough{\lstinline!VALUE\_TYPE!}}：当前绑定源码已实现该消息值操作。

\textbf{参数}

\begin{longtable}[]{@{}
  >{\raggedright\arraybackslash}p{(\columnwidth - 4\tabcolsep) * \real{0.18}}
  >{\raggedright\arraybackslash}p{(\columnwidth - 4\tabcolsep) * \real{0.20}}
  >{\raggedright\arraybackslash}p{(\columnwidth - 4\tabcolsep) * \real{0.62}}@{}}
\toprule
名称 & Python 类型 & 含义 \\
\midrule
\endhead
\passthrough{\lstinline!value!} &
\passthrough{\lstinline!Sequence[float]!} &
要写入或传递的新值，必须符合该参数的 Python 类型以及底层 C++
范围约束。 \\
\bottomrule
\end{longtable}

\textbf{返回值}

\begin{longtable}[]{@{}
  >{\raggedright\arraybackslash}p{(\columnwidth - 4\tabcolsep) * \real{0.18}}
  >{\raggedright\arraybackslash}p{(\columnwidth - 4\tabcolsep) * \real{0.20}}
  >{\raggedright\arraybackslash}p{(\columnwidth - 4\tabcolsep) * \real{0.62}}@{}}
\toprule
位置 & 类型 & 含义 \\
\midrule
\endhead
getter 返回值 & \passthrough{\lstinline!list[float]!} &
返回属性当前值。数组、vector 和嵌套消息采用复制语义。 \\
\bottomrule
\end{longtable}

\textbf{对应 C++}

\begin{itemize}
\tightlist
\item
  类：\passthrough{\lstinline!unitree\_go::msg::dds\_::IMUState\_!}
\item
  字段类型：\passthrough{\lstinline!std::array<float, 3>!}
\item
  声明位置：\passthrough{\lstinline!include/unitree/idl/go2/IMUState\_.hpp:26!}
\end{itemize}

\textbf{用法}

\begin{lstlisting}[language=Python]
current_value = obj.accelerometer
obj.accelerometer = new_value
\end{lstlisting}

\begin{quote}
{[}!TIP{]} 读取容器或嵌套值后应遵循``读取、修改、重新赋值''。只修改
getter 返回的副本不会可靠地更新原 C++ 消息。
\end{quote}

\hypertarget{unitree-sdk2-cpp-idl-go2-imustate-rpy}{%
\paragraph{\texorpdfstring{\texttt{unitree\_sdk2\_cpp.idl.go2.IMUState.rpy}}{unitree\_sdk2\_cpp.idl.go2.IMUState.rpy}}\label{unitree-sdk2-cpp-idl-go2-imustate-rpy}}

底层 \passthrough{\lstinline!unitree\_go::msg::dds\_::IMUState\_!} 的
\passthrough{\lstinline!rpy!} 字段。类型签名只说明 Python
表示；单位、数值范围、数组固定长度和枚举含义需查对应 IDL 头文件。
底层是固定长度数组，写入序列必须正好包含 3
个元素；元素约束：底层为浮点数；类型本身不说明单位、坐标系或业务安全范围。

\textbf{签名}

\begin{lstlisting}[language=Python]
@property
def rpy(self) -> list[float]

@rpy.setter
def rpy(self, value: Sequence[float]) -> None
\end{lstlisting}

\textbf{可用性与安全性}

\textbf{\passthrough{\lstinline!AVAILABLE!} /
\passthrough{\lstinline!VALUE\_TYPE!}}：当前绑定源码已实现该消息值操作。

\textbf{参数}

\begin{longtable}[]{@{}
  >{\raggedright\arraybackslash}p{(\columnwidth - 4\tabcolsep) * \real{0.18}}
  >{\raggedright\arraybackslash}p{(\columnwidth - 4\tabcolsep) * \real{0.20}}
  >{\raggedright\arraybackslash}p{(\columnwidth - 4\tabcolsep) * \real{0.62}}@{}}
\toprule
名称 & Python 类型 & 含义 \\
\midrule
\endhead
\passthrough{\lstinline!value!} &
\passthrough{\lstinline!Sequence[float]!} &
要写入或传递的新值，必须符合该参数的 Python 类型以及底层 C++
范围约束。 \\
\bottomrule
\end{longtable}

\textbf{返回值}

\begin{longtable}[]{@{}
  >{\raggedright\arraybackslash}p{(\columnwidth - 4\tabcolsep) * \real{0.18}}
  >{\raggedright\arraybackslash}p{(\columnwidth - 4\tabcolsep) * \real{0.20}}
  >{\raggedright\arraybackslash}p{(\columnwidth - 4\tabcolsep) * \real{0.62}}@{}}
\toprule
位置 & 类型 & 含义 \\
\midrule
\endhead
getter 返回值 & \passthrough{\lstinline!list[float]!} &
返回属性当前值。数组、vector 和嵌套消息采用复制语义。 \\
\bottomrule
\end{longtable}

\textbf{对应 C++}

\begin{itemize}
\tightlist
\item
  类：\passthrough{\lstinline!unitree\_go::msg::dds\_::IMUState\_!}
\item
  字段类型：\passthrough{\lstinline!std::array<float, 3>!}
\item
  声明位置：\passthrough{\lstinline!include/unitree/idl/go2/IMUState\_.hpp:27!}
\end{itemize}

\textbf{用法}

\begin{lstlisting}[language=Python]
current_value = obj.rpy
obj.rpy = new_value
\end{lstlisting}

\begin{quote}
{[}!TIP{]} 读取容器或嵌套值后应遵循``读取、修改、重新赋值''。只修改
getter 返回的副本不会可靠地更新原 C++ 消息。
\end{quote}

\hypertarget{unitree-sdk2-cpp-idl-go2-imustate-temperature}{%
\paragraph{\texorpdfstring{\texttt{unitree\_sdk2\_cpp.idl.go2.IMUState.temperature}}{unitree\_sdk2\_cpp.idl.go2.IMUState.temperature}}\label{unitree-sdk2-cpp-idl-go2-imustate-temperature}}

底层 \passthrough{\lstinline!unitree\_go::msg::dds\_::IMUState\_!} 的
\passthrough{\lstinline!temperature!} 字段。类型签名只说明 Python
表示；单位、数值范围、数组固定长度和枚举含义需查对应 IDL 头文件。
取值范围为 0 到 255。

\textbf{签名}

\begin{lstlisting}[language=Python]
@property
def temperature(self) -> int

@temperature.setter
def temperature(self, value: int) -> None
\end{lstlisting}

\textbf{可用性与安全性}

\textbf{\passthrough{\lstinline!AVAILABLE!} /
\passthrough{\lstinline!VALUE\_TYPE!}}：当前绑定源码已实现该消息值操作。

\textbf{参数}

\begin{longtable}[]{@{}
  >{\raggedright\arraybackslash}p{(\columnwidth - 4\tabcolsep) * \real{0.18}}
  >{\raggedright\arraybackslash}p{(\columnwidth - 4\tabcolsep) * \real{0.20}}
  >{\raggedright\arraybackslash}p{(\columnwidth - 4\tabcolsep) * \real{0.62}}@{}}
\toprule
名称 & Python 类型 & 含义 \\
\midrule
\endhead
\passthrough{\lstinline!value!} & \passthrough{\lstinline!int!} &
要写入或传递的新值，必须符合该参数的 Python 类型以及底层 C++
范围约束。 \\
\bottomrule
\end{longtable}

\textbf{返回值}

\begin{longtable}[]{@{}
  >{\raggedright\arraybackslash}p{(\columnwidth - 4\tabcolsep) * \real{0.18}}
  >{\raggedright\arraybackslash}p{(\columnwidth - 4\tabcolsep) * \real{0.20}}
  >{\raggedright\arraybackslash}p{(\columnwidth - 4\tabcolsep) * \real{0.62}}@{}}
\toprule
位置 & 类型 & 含义 \\
\midrule
\endhead
getter 返回值 & \passthrough{\lstinline!int!} & 返回属性当前值。 \\
\bottomrule
\end{longtable}

\textbf{对应 C++}

\begin{itemize}
\tightlist
\item
  类：\passthrough{\lstinline!unitree\_go::msg::dds\_::IMUState\_!}
\item
  字段类型：\passthrough{\lstinline!uint8\_t!}
\item
  声明位置：\passthrough{\lstinline!include/unitree/idl/go2/IMUState\_.hpp:28!}
\end{itemize}

\textbf{用法}

\begin{lstlisting}[language=Python]
current_value = obj.temperature
obj.temperature = new_value
\end{lstlisting}

\hypertarget{unitree-sdk2-cpp-idl-go2-interfaceconfig}{%
\subsubsection{\texorpdfstring{\texttt{unitree\_sdk2\_cpp.idl.go2.InterfaceConfig}}{unitree\_sdk2\_cpp.idl.go2.InterfaceConfig}}\label{unitree-sdk2-cpp-idl-go2-interfaceconfig}}

可由 Python 读写的 DDS/IDL 消息值类型。字段采用复制语义。

\textbf{导入}

\begin{lstlisting}[language=Python]
from unitree_sdk2_cpp.idl.go2 import InterfaceConfig
\end{lstlisting}

\textbf{方法索引}

\begin{longtable}[]{@{}
  >{\raggedright\arraybackslash}p{(\columnwidth - 6\tabcolsep) * \real{0.18}}
  >{\raggedright\arraybackslash}p{(\columnwidth - 6\tabcolsep) * \real{0.24}}
  >{\raggedright\arraybackslash}p{(\columnwidth - 6\tabcolsep) * \real{0.18}}
  >{\raggedright\arraybackslash}p{(\columnwidth - 6\tabcolsep) * \real{0.40}}@{}}
\toprule
方法 & Python 签名 & 状态 & 安全分类 \\
\midrule
\endhead
\protect\hyperlink{unitree-sdk2-cpp-idl-go2-interfaceconfig-dunder-init-1}{\passthrough{\lstinline!\_\_init\_\_!}}
& \passthrough{\lstinline!def \_\_init\_\_(self) -> None!} &
\passthrough{\lstinline!AVAILABLE!} &
\passthrough{\lstinline!VALUE\_TYPE!} \\
\protect\hyperlink{unitree-sdk2-cpp-idl-go2-interfaceconfig-dunder-eq-1}{\passthrough{\lstinline!\_\_eq\_\_!}}
& \passthrough{\lstinline!def \_\_eq\_\_(self, other: object) -> bool!}
& \passthrough{\lstinline!AVAILABLE!} &
\passthrough{\lstinline!VALUE\_TYPE!} \\
\protect\hyperlink{unitree-sdk2-cpp-idl-go2-interfaceconfig-dunder-ne-1}{\passthrough{\lstinline!\_\_ne\_\_!}}
& \passthrough{\lstinline!def \_\_ne\_\_(self, other: object) -> bool!}
& \passthrough{\lstinline!AVAILABLE!} &
\passthrough{\lstinline!VALUE\_TYPE!} \\
\bottomrule
\end{longtable}

\textbf{属性索引}

\begin{longtable}[]{@{}
  >{\raggedright\arraybackslash}p{(\columnwidth - 6\tabcolsep) * \real{0.18}}
  >{\raggedright\arraybackslash}p{(\columnwidth - 6\tabcolsep) * \real{0.24}}
  >{\raggedright\arraybackslash}p{(\columnwidth - 6\tabcolsep) * \real{0.18}}
  >{\raggedright\arraybackslash}p{(\columnwidth - 6\tabcolsep) * \real{0.40}}@{}}
\toprule
属性 & 读取类型 & 写入类型 & 状态 \\
\midrule
\endhead
\protect\hyperlink{unitree-sdk2-cpp-idl-go2-interfaceconfig-mode}{\passthrough{\lstinline!mode!}}
& \passthrough{\lstinline!int!} & \passthrough{\lstinline!int!} &
\passthrough{\lstinline!AVAILABLE!} \\
\protect\hyperlink{unitree-sdk2-cpp-idl-go2-interfaceconfig-value}{\passthrough{\lstinline!value!}}
& \passthrough{\lstinline!int!} & \passthrough{\lstinline!int!} &
\passthrough{\lstinline!AVAILABLE!} \\
\protect\hyperlink{unitree-sdk2-cpp-idl-go2-interfaceconfig-reserve}{\passthrough{\lstinline!reserve!}}
& \passthrough{\lstinline!list[int]!} &
\passthrough{\lstinline!Sequence[int]!} &
\passthrough{\lstinline!AVAILABLE!} \\
\bottomrule
\end{longtable}

\hypertarget{unitree-sdk2-cpp-idl-go2-interfaceconfig-dunder-init-1}{%
\paragraph{\texorpdfstring{\texttt{unitree\_sdk2\_cpp.idl.go2.InterfaceConfig.\_\_init\_\_}}{unitree\_sdk2\_cpp.idl.go2.InterfaceConfig.\_\_init\_\_}}\label{unitree-sdk2-cpp-idl-go2-interfaceconfig-dunder-init-1}}

初始化 \passthrough{\lstinline!InterfaceConfig!}
实例。是否能实际构造取决于下方可用性状态。

\textbf{签名}

\begin{lstlisting}[language=Python]
def __init__(self) -> None
\end{lstlisting}

\textbf{可用性与安全性}

\textbf{\passthrough{\lstinline!AVAILABLE!} /
\passthrough{\lstinline!VALUE\_TYPE!}}：当前绑定源码已实现该消息值操作。

\textbf{参数}

无显式参数。实例方法中的 \passthrough{\lstinline!self!} 由 Python
自动传入。

\textbf{返回值}

\begin{longtable}[]{@{}
  >{\raggedright\arraybackslash}p{(\columnwidth - 4\tabcolsep) * \real{0.18}}
  >{\raggedright\arraybackslash}p{(\columnwidth - 4\tabcolsep) * \real{0.20}}
  >{\raggedright\arraybackslash}p{(\columnwidth - 4\tabcolsep) * \real{0.62}}@{}}
\toprule
位置 & 类型 & 含义 \\
\midrule
\endhead
无 & \passthrough{\lstinline!None!} &
\passthrough{\lstinline!\_\_init\_\_!}
本身不返回值；调用类对象时，在构造函数可用的前提下得到该类实例。 \\
\bottomrule
\end{longtable}

\textbf{对应 C++}

\begin{itemize}
\tightlist
\item
  类：\passthrough{\lstinline!unitree\_go::msg::dds\_::InterfaceConfig\_!}
\item
  签名：\passthrough{\lstinline!InterfaceConfig\_()!}
\item
  绑定策略：\passthrough{\lstinline!未单独标注!}
\item
  声明位置：\passthrough{\lstinline!include/unitree/idl/go2/InterfaceConfig\_.hpp:29!}
\end{itemize}

\textbf{说明}

\begin{itemize}
\tightlist
\item
  示例是局部调用片段；\passthrough{\lstinline!obj!}
  和参数变量表示已经按上层业务校验并准备好的对象。
\end{itemize}

\textbf{用法}

\begin{lstlisting}[language=Python]
value = InterfaceConfig()
\end{lstlisting}

\hypertarget{unitree-sdk2-cpp-idl-go2-interfaceconfig-dunder-eq-1}{%
\paragraph{\texorpdfstring{\texttt{unitree\_sdk2\_cpp.idl.go2.InterfaceConfig.\_\_eq\_\_}}{unitree\_sdk2\_cpp.idl.go2.InterfaceConfig.\_\_eq\_\_}}\label{unitree-sdk2-cpp-idl-go2-interfaceconfig-dunder-eq-1}}

按底层 IDL 消息内容比较两个对象是否相等。

\textbf{签名}

\begin{lstlisting}[language=Python]
def __eq__(self, other: object) -> bool
\end{lstlisting}

\textbf{可用性与安全性}

\textbf{\passthrough{\lstinline!AVAILABLE!} /
\passthrough{\lstinline!VALUE\_TYPE!}}：当前绑定源码已实现该消息值操作。

\textbf{参数}

\begin{longtable}[]{@{}
  >{\raggedright\arraybackslash}p{(\columnwidth - 6\tabcolsep) * \real{0.18}}
  >{\raggedright\arraybackslash}p{(\columnwidth - 6\tabcolsep) * \real{0.24}}
  >{\raggedright\arraybackslash}p{(\columnwidth - 6\tabcolsep) * \real{0.18}}
  >{\raggedright\arraybackslash}p{(\columnwidth - 6\tabcolsep) * \real{0.40}}@{}}
\toprule
名称 & Python 类型 & 默认值 & 含义 \\
\midrule
\endhead
\passthrough{\lstinline!other!} & \passthrough{\lstinline!object!} &
必填 & 用于比较的另一个 Python 对象。类型不兼容时比较结果通常为
\passthrough{\lstinline!False!}。 \\
\bottomrule
\end{longtable}

\textbf{返回值}

\begin{longtable}[]{@{}
  >{\raggedright\arraybackslash}p{(\columnwidth - 4\tabcolsep) * \real{0.18}}
  >{\raggedright\arraybackslash}p{(\columnwidth - 4\tabcolsep) * \real{0.20}}
  >{\raggedright\arraybackslash}p{(\columnwidth - 4\tabcolsep) * \real{0.62}}@{}}
\toprule
位置 & 类型 & 含义 \\
\midrule
\endhead
返回值 & \passthrough{\lstinline!bool!} & 返回
\passthrough{\lstinline!bool!}。更精确的业务含义见该方法用途、C++
签名和目标型号协议。 \\
\bottomrule
\end{longtable}

\textbf{对应 C++}

\begin{itemize}
\tightlist
\item
  类：\passthrough{\lstinline!unitree\_go::msg::dds\_::InterfaceConfig\_!}
\item
  签名：\passthrough{\lstinline!operator==(const value \&) const!}
\item
  绑定策略：\passthrough{\lstinline!未单独标注!}
\item
  声明位置：由 pybind11 手工绑定或生成注册表提供；manifest
  未记录独立头文件行号。
\end{itemize}

\textbf{说明}

\begin{itemize}
\tightlist
\item
  示例是局部调用片段；\passthrough{\lstinline!obj!}
  和参数变量表示已经按上层业务校验并准备好的对象。
\end{itemize}

\textbf{用法}

\begin{lstlisting}[language=Python]
same = left == right
\end{lstlisting}

\hypertarget{unitree-sdk2-cpp-idl-go2-interfaceconfig-dunder-ne-1}{%
\paragraph{\texorpdfstring{\texttt{unitree\_sdk2\_cpp.idl.go2.InterfaceConfig.\_\_ne\_\_}}{unitree\_sdk2\_cpp.idl.go2.InterfaceConfig.\_\_ne\_\_}}\label{unitree-sdk2-cpp-idl-go2-interfaceconfig-dunder-ne-1}}

按底层 IDL 消息内容比较两个对象是否不相等。

\textbf{签名}

\begin{lstlisting}[language=Python]
def __ne__(self, other: object) -> bool
\end{lstlisting}

\textbf{可用性与安全性}

\textbf{\passthrough{\lstinline!AVAILABLE!} /
\passthrough{\lstinline!VALUE\_TYPE!}}：当前绑定源码已实现该消息值操作。

\textbf{参数}

\begin{longtable}[]{@{}
  >{\raggedright\arraybackslash}p{(\columnwidth - 6\tabcolsep) * \real{0.18}}
  >{\raggedright\arraybackslash}p{(\columnwidth - 6\tabcolsep) * \real{0.24}}
  >{\raggedright\arraybackslash}p{(\columnwidth - 6\tabcolsep) * \real{0.18}}
  >{\raggedright\arraybackslash}p{(\columnwidth - 6\tabcolsep) * \real{0.40}}@{}}
\toprule
名称 & Python 类型 & 默认值 & 含义 \\
\midrule
\endhead
\passthrough{\lstinline!other!} & \passthrough{\lstinline!object!} &
必填 & 用于比较的另一个 Python 对象。类型不兼容时比较结果通常为
\passthrough{\lstinline!False!}。 \\
\bottomrule
\end{longtable}

\textbf{返回值}

\begin{longtable}[]{@{}
  >{\raggedright\arraybackslash}p{(\columnwidth - 4\tabcolsep) * \real{0.18}}
  >{\raggedright\arraybackslash}p{(\columnwidth - 4\tabcolsep) * \real{0.20}}
  >{\raggedright\arraybackslash}p{(\columnwidth - 4\tabcolsep) * \real{0.62}}@{}}
\toprule
位置 & 类型 & 含义 \\
\midrule
\endhead
返回值 & \passthrough{\lstinline!bool!} & 返回
\passthrough{\lstinline!bool!}。更精确的业务含义见该方法用途、C++
签名和目标型号协议。 \\
\bottomrule
\end{longtable}

\textbf{对应 C++}

\begin{itemize}
\tightlist
\item
  类：\passthrough{\lstinline!unitree\_go::msg::dds\_::InterfaceConfig\_!}
\item
  签名：\passthrough{\lstinline"operator!=(const value \&) const"}
\item
  绑定策略：\passthrough{\lstinline!未单独标注!}
\item
  声明位置：由 pybind11 手工绑定或生成注册表提供；manifest
  未记录独立头文件行号。
\end{itemize}

\textbf{说明}

\begin{itemize}
\tightlist
\item
  示例是局部调用片段；\passthrough{\lstinline!obj!}
  和参数变量表示已经按上层业务校验并准备好的对象。
\end{itemize}

\textbf{用法}

\begin{lstlisting}[language=Python]
different = left != right
\end{lstlisting}

\hypertarget{unitree-sdk2-cpp-idl-go2-interfaceconfig-mode}{%
\paragraph{\texorpdfstring{\texttt{unitree\_sdk2\_cpp.idl.go2.InterfaceConfig.mode}}{unitree\_sdk2\_cpp.idl.go2.InterfaceConfig.mode}}\label{unitree-sdk2-cpp-idl-go2-interfaceconfig-mode}}

底层
\passthrough{\lstinline!unitree\_go::msg::dds\_::InterfaceConfig\_!} 的
\passthrough{\lstinline!mode!} 字段。类型签名只说明 Python
表示；单位、数值范围、数组固定长度和枚举含义需查对应 IDL 头文件。
取值范围为 0 到 255。

\textbf{签名}

\begin{lstlisting}[language=Python]
@property
def mode(self) -> int

@mode.setter
def mode(self, value: int) -> None
\end{lstlisting}

\textbf{可用性与安全性}

\textbf{\passthrough{\lstinline!AVAILABLE!} /
\passthrough{\lstinline!VALUE\_TYPE!}}：当前绑定源码已实现该消息值操作。

\textbf{参数}

\begin{longtable}[]{@{}
  >{\raggedright\arraybackslash}p{(\columnwidth - 4\tabcolsep) * \real{0.18}}
  >{\raggedright\arraybackslash}p{(\columnwidth - 4\tabcolsep) * \real{0.20}}
  >{\raggedright\arraybackslash}p{(\columnwidth - 4\tabcolsep) * \real{0.62}}@{}}
\toprule
名称 & Python 类型 & 含义 \\
\midrule
\endhead
\passthrough{\lstinline!value!} & \passthrough{\lstinline!int!} &
要写入或传递的新值，必须符合该参数的 Python 类型以及底层 C++
范围约束。 \\
\bottomrule
\end{longtable}

\textbf{返回值}

\begin{longtable}[]{@{}
  >{\raggedright\arraybackslash}p{(\columnwidth - 4\tabcolsep) * \real{0.18}}
  >{\raggedright\arraybackslash}p{(\columnwidth - 4\tabcolsep) * \real{0.20}}
  >{\raggedright\arraybackslash}p{(\columnwidth - 4\tabcolsep) * \real{0.62}}@{}}
\toprule
位置 & 类型 & 含义 \\
\midrule
\endhead
getter 返回值 & \passthrough{\lstinline!int!} & 返回属性当前值。 \\
\bottomrule
\end{longtable}

\textbf{对应 C++}

\begin{itemize}
\tightlist
\item
  类：\passthrough{\lstinline!unitree\_go::msg::dds\_::InterfaceConfig\_!}
\item
  字段类型：\passthrough{\lstinline!uint8\_t!}
\item
  声明位置：\passthrough{\lstinline!include/unitree/idl/go2/InterfaceConfig\_.hpp:24!}
\end{itemize}

\textbf{用法}

\begin{lstlisting}[language=Python]
current_value = obj.mode
obj.mode = new_value
\end{lstlisting}

\hypertarget{unitree-sdk2-cpp-idl-go2-interfaceconfig-value}{%
\paragraph{\texorpdfstring{\texttt{unitree\_sdk2\_cpp.idl.go2.InterfaceConfig.value}}{unitree\_sdk2\_cpp.idl.go2.InterfaceConfig.value}}\label{unitree-sdk2-cpp-idl-go2-interfaceconfig-value}}

底层
\passthrough{\lstinline!unitree\_go::msg::dds\_::InterfaceConfig\_!} 的
\passthrough{\lstinline!value!} 字段。类型签名只说明 Python
表示；单位、数值范围、数组固定长度和枚举含义需查对应 IDL 头文件。
取值范围为 0 到 255。

\textbf{签名}

\begin{lstlisting}[language=Python]
@property
def value(self) -> int

@value.setter
def value(self, value: int) -> None
\end{lstlisting}

\textbf{可用性与安全性}

\textbf{\passthrough{\lstinline!AVAILABLE!} /
\passthrough{\lstinline!VALUE\_TYPE!}}：当前绑定源码已实现该消息值操作。

\textbf{参数}

\begin{longtable}[]{@{}
  >{\raggedright\arraybackslash}p{(\columnwidth - 4\tabcolsep) * \real{0.18}}
  >{\raggedright\arraybackslash}p{(\columnwidth - 4\tabcolsep) * \real{0.20}}
  >{\raggedright\arraybackslash}p{(\columnwidth - 4\tabcolsep) * \real{0.62}}@{}}
\toprule
名称 & Python 类型 & 含义 \\
\midrule
\endhead
\passthrough{\lstinline!value!} & \passthrough{\lstinline!int!} &
要写入或传递的新值，必须符合该参数的 Python 类型以及底层 C++
范围约束。 \\
\bottomrule
\end{longtable}

\textbf{返回值}

\begin{longtable}[]{@{}
  >{\raggedright\arraybackslash}p{(\columnwidth - 4\tabcolsep) * \real{0.18}}
  >{\raggedright\arraybackslash}p{(\columnwidth - 4\tabcolsep) * \real{0.20}}
  >{\raggedright\arraybackslash}p{(\columnwidth - 4\tabcolsep) * \real{0.62}}@{}}
\toprule
位置 & 类型 & 含义 \\
\midrule
\endhead
getter 返回值 & \passthrough{\lstinline!int!} & 返回属性当前值。 \\
\bottomrule
\end{longtable}

\textbf{对应 C++}

\begin{itemize}
\tightlist
\item
  类：\passthrough{\lstinline!unitree\_go::msg::dds\_::InterfaceConfig\_!}
\item
  字段类型：\passthrough{\lstinline!uint8\_t!}
\item
  声明位置：\passthrough{\lstinline!include/unitree/idl/go2/InterfaceConfig\_.hpp:25!}
\end{itemize}

\textbf{用法}

\begin{lstlisting}[language=Python]
current_value = obj.value
obj.value = new_value
\end{lstlisting}

\hypertarget{unitree-sdk2-cpp-idl-go2-interfaceconfig-reserve}{%
\paragraph{\texorpdfstring{\texttt{unitree\_sdk2\_cpp.idl.go2.InterfaceConfig.reserve}}{unitree\_sdk2\_cpp.idl.go2.InterfaceConfig.reserve}}\label{unitree-sdk2-cpp-idl-go2-interfaceconfig-reserve}}

底层
\passthrough{\lstinline!unitree\_go::msg::dds\_::InterfaceConfig\_!} 的
\passthrough{\lstinline!reserve!} 字段。类型签名只说明 Python
表示；单位、数值范围、数组固定长度和枚举含义需查对应 IDL 头文件。
底层是固定长度数组，写入序列必须正好包含 2 个元素；元素约束：取值范围为
0 到 255。

\textbf{签名}

\begin{lstlisting}[language=Python]
@property
def reserve(self) -> list[int]

@reserve.setter
def reserve(self, value: Sequence[int]) -> None
\end{lstlisting}

\textbf{可用性与安全性}

\textbf{\passthrough{\lstinline!AVAILABLE!} /
\passthrough{\lstinline!VALUE\_TYPE!}}：当前绑定源码已实现该消息值操作。

\textbf{参数}

\begin{longtable}[]{@{}
  >{\raggedright\arraybackslash}p{(\columnwidth - 4\tabcolsep) * \real{0.18}}
  >{\raggedright\arraybackslash}p{(\columnwidth - 4\tabcolsep) * \real{0.20}}
  >{\raggedright\arraybackslash}p{(\columnwidth - 4\tabcolsep) * \real{0.62}}@{}}
\toprule
名称 & Python 类型 & 含义 \\
\midrule
\endhead
\passthrough{\lstinline!value!} &
\passthrough{\lstinline!Sequence[int]!} &
要写入或传递的新值，必须符合该参数的 Python 类型以及底层 C++
范围约束。 \\
\bottomrule
\end{longtable}

\textbf{返回值}

\begin{longtable}[]{@{}
  >{\raggedright\arraybackslash}p{(\columnwidth - 4\tabcolsep) * \real{0.18}}
  >{\raggedright\arraybackslash}p{(\columnwidth - 4\tabcolsep) * \real{0.20}}
  >{\raggedright\arraybackslash}p{(\columnwidth - 4\tabcolsep) * \real{0.62}}@{}}
\toprule
位置 & 类型 & 含义 \\
\midrule
\endhead
getter 返回值 & \passthrough{\lstinline!list[int]!} &
返回属性当前值。数组、vector 和嵌套消息采用复制语义。 \\
\bottomrule
\end{longtable}

\textbf{对应 C++}

\begin{itemize}
\tightlist
\item
  类：\passthrough{\lstinline!unitree\_go::msg::dds\_::InterfaceConfig\_!}
\item
  字段类型：\passthrough{\lstinline!std::array<uint8\_t, 2>!}
\item
  声明位置：\passthrough{\lstinline!include/unitree/idl/go2/InterfaceConfig\_.hpp:26!}
\end{itemize}

\textbf{用法}

\begin{lstlisting}[language=Python]
current_value = obj.reserve
obj.reserve = new_value
\end{lstlisting}

\begin{quote}
{[}!TIP{]} 读取容器或嵌套值后应遵循``读取、修改、重新赋值''。只修改
getter 返回的副本不会可靠地更新原 C++ 消息。
\end{quote}

\hypertarget{unitree-sdk2-cpp-idl-go2-lidarstate}{%
\subsubsection{\texorpdfstring{\texttt{unitree\_sdk2\_cpp.idl.go2.LidarState}}{unitree\_sdk2\_cpp.idl.go2.LidarState}}\label{unitree-sdk2-cpp-idl-go2-lidarstate}}

可由 Python 读写的 DDS/IDL 消息值类型。字段采用复制语义。

\textbf{导入}

\begin{lstlisting}[language=Python]
from unitree_sdk2_cpp.idl.go2 import LidarState
\end{lstlisting}

\textbf{方法索引}

\begin{longtable}[]{@{}
  >{\raggedright\arraybackslash}p{(\columnwidth - 6\tabcolsep) * \real{0.18}}
  >{\raggedright\arraybackslash}p{(\columnwidth - 6\tabcolsep) * \real{0.24}}
  >{\raggedright\arraybackslash}p{(\columnwidth - 6\tabcolsep) * \real{0.18}}
  >{\raggedright\arraybackslash}p{(\columnwidth - 6\tabcolsep) * \real{0.40}}@{}}
\toprule
方法 & Python 签名 & 状态 & 安全分类 \\
\midrule
\endhead
\protect\hyperlink{unitree-sdk2-cpp-idl-go2-lidarstate-dunder-init-1}{\passthrough{\lstinline!\_\_init\_\_!}}
& \passthrough{\lstinline!def \_\_init\_\_(self) -> None!} &
\passthrough{\lstinline!AVAILABLE!} &
\passthrough{\lstinline!VALUE\_TYPE!} \\
\protect\hyperlink{unitree-sdk2-cpp-idl-go2-lidarstate-dunder-eq-1}{\passthrough{\lstinline!\_\_eq\_\_!}}
& \passthrough{\lstinline!def \_\_eq\_\_(self, other: object) -> bool!}
& \passthrough{\lstinline!AVAILABLE!} &
\passthrough{\lstinline!VALUE\_TYPE!} \\
\protect\hyperlink{unitree-sdk2-cpp-idl-go2-lidarstate-dunder-ne-1}{\passthrough{\lstinline!\_\_ne\_\_!}}
& \passthrough{\lstinline!def \_\_ne\_\_(self, other: object) -> bool!}
& \passthrough{\lstinline!AVAILABLE!} &
\passthrough{\lstinline!VALUE\_TYPE!} \\
\bottomrule
\end{longtable}

\textbf{属性索引}

\begin{longtable}[]{@{}
  >{\raggedright\arraybackslash}p{(\columnwidth - 6\tabcolsep) * \real{0.18}}
  >{\raggedright\arraybackslash}p{(\columnwidth - 6\tabcolsep) * \real{0.24}}
  >{\raggedright\arraybackslash}p{(\columnwidth - 6\tabcolsep) * \real{0.18}}
  >{\raggedright\arraybackslash}p{(\columnwidth - 6\tabcolsep) * \real{0.40}}@{}}
\toprule
属性 & 读取类型 & 写入类型 & 状态 \\
\midrule
\endhead
\protect\hyperlink{unitree-sdk2-cpp-idl-go2-lidarstate-stamp}{\passthrough{\lstinline!stamp!}}
& \passthrough{\lstinline!float!} & \passthrough{\lstinline!float!} &
\passthrough{\lstinline!AVAILABLE!} \\
\protect\hyperlink{unitree-sdk2-cpp-idl-go2-lidarstate-firmware-version}{\passthrough{\lstinline!firmware\_version!}}
& \passthrough{\lstinline!str!} & \passthrough{\lstinline!str!} &
\passthrough{\lstinline!AVAILABLE!} \\
\protect\hyperlink{unitree-sdk2-cpp-idl-go2-lidarstate-software-version}{\passthrough{\lstinline!software\_version!}}
& \passthrough{\lstinline!str!} & \passthrough{\lstinline!str!} &
\passthrough{\lstinline!AVAILABLE!} \\
\protect\hyperlink{unitree-sdk2-cpp-idl-go2-lidarstate-sdk-version}{\passthrough{\lstinline!sdk\_version!}}
& \passthrough{\lstinline!str!} & \passthrough{\lstinline!str!} &
\passthrough{\lstinline!AVAILABLE!} \\
\protect\hyperlink{unitree-sdk2-cpp-idl-go2-lidarstate-sys-rotation-speed}{\passthrough{\lstinline!sys\_rotation\_speed!}}
& \passthrough{\lstinline!float!} & \passthrough{\lstinline!float!} &
\passthrough{\lstinline!AVAILABLE!} \\
\protect\hyperlink{unitree-sdk2-cpp-idl-go2-lidarstate-com-rotation-speed}{\passthrough{\lstinline!com\_rotation\_speed!}}
& \passthrough{\lstinline!float!} & \passthrough{\lstinline!float!} &
\passthrough{\lstinline!AVAILABLE!} \\
\protect\hyperlink{unitree-sdk2-cpp-idl-go2-lidarstate-error-state}{\passthrough{\lstinline!error\_state!}}
& \passthrough{\lstinline!int!} & \passthrough{\lstinline!int!} &
\passthrough{\lstinline!AVAILABLE!} \\
\protect\hyperlink{unitree-sdk2-cpp-idl-go2-lidarstate-dirty-percentage}{\passthrough{\lstinline!dirty\_percentage!}}
& \passthrough{\lstinline!int!} & \passthrough{\lstinline!int!} &
\passthrough{\lstinline!AVAILABLE!} \\
\protect\hyperlink{unitree-sdk2-cpp-idl-go2-lidarstate-cloud-frequency}{\passthrough{\lstinline!cloud\_frequency!}}
& \passthrough{\lstinline!float!} & \passthrough{\lstinline!float!} &
\passthrough{\lstinline!AVAILABLE!} \\
\protect\hyperlink{unitree-sdk2-cpp-idl-go2-lidarstate-cloud-packet-loss-rate}{\passthrough{\lstinline!cloud\_packet\_loss\_rate!}}
& \passthrough{\lstinline!float!} & \passthrough{\lstinline!float!} &
\passthrough{\lstinline!AVAILABLE!} \\
\protect\hyperlink{unitree-sdk2-cpp-idl-go2-lidarstate-cloud-size}{\passthrough{\lstinline!cloud\_size!}}
& \passthrough{\lstinline!int!} & \passthrough{\lstinline!int!} &
\passthrough{\lstinline!AVAILABLE!} \\
\protect\hyperlink{unitree-sdk2-cpp-idl-go2-lidarstate-cloud-scan-num}{\passthrough{\lstinline!cloud\_scan\_num!}}
& \passthrough{\lstinline!int!} & \passthrough{\lstinline!int!} &
\passthrough{\lstinline!AVAILABLE!} \\
\protect\hyperlink{unitree-sdk2-cpp-idl-go2-lidarstate-imu-frequency}{\passthrough{\lstinline!imu\_frequency!}}
& \passthrough{\lstinline!float!} & \passthrough{\lstinline!float!} &
\passthrough{\lstinline!AVAILABLE!} \\
\protect\hyperlink{unitree-sdk2-cpp-idl-go2-lidarstate-imu-packet-loss-rate}{\passthrough{\lstinline!imu\_packet\_loss\_rate!}}
& \passthrough{\lstinline!float!} & \passthrough{\lstinline!float!} &
\passthrough{\lstinline!AVAILABLE!} \\
\protect\hyperlink{unitree-sdk2-cpp-idl-go2-lidarstate-imu-rpy}{\passthrough{\lstinline!imu\_rpy!}}
& \passthrough{\lstinline!list[float]!} &
\passthrough{\lstinline!Sequence[float]!} &
\passthrough{\lstinline!AVAILABLE!} \\
\protect\hyperlink{unitree-sdk2-cpp-idl-go2-lidarstate-serial-recv-stamp}{\passthrough{\lstinline!serial\_recv\_stamp!}}
& \passthrough{\lstinline!float!} & \passthrough{\lstinline!float!} &
\passthrough{\lstinline!AVAILABLE!} \\
\protect\hyperlink{unitree-sdk2-cpp-idl-go2-lidarstate-serial-buffer-size}{\passthrough{\lstinline!serial\_buffer\_size!}}
& \passthrough{\lstinline!int!} & \passthrough{\lstinline!int!} &
\passthrough{\lstinline!AVAILABLE!} \\
\protect\hyperlink{unitree-sdk2-cpp-idl-go2-lidarstate-serial-buffer-read}{\passthrough{\lstinline!serial\_buffer\_read!}}
& \passthrough{\lstinline!int!} & \passthrough{\lstinline!int!} &
\passthrough{\lstinline!AVAILABLE!} \\
\bottomrule
\end{longtable}

\hypertarget{unitree-sdk2-cpp-idl-go2-lidarstate-dunder-init-1}{%
\paragraph{\texorpdfstring{\texttt{unitree\_sdk2\_cpp.idl.go2.LidarState.\_\_init\_\_}}{unitree\_sdk2\_cpp.idl.go2.LidarState.\_\_init\_\_}}\label{unitree-sdk2-cpp-idl-go2-lidarstate-dunder-init-1}}

初始化 \passthrough{\lstinline!LidarState!}
实例。是否能实际构造取决于下方可用性状态。

\textbf{签名}

\begin{lstlisting}[language=Python]
def __init__(self) -> None
\end{lstlisting}

\textbf{可用性与安全性}

\textbf{\passthrough{\lstinline!AVAILABLE!} /
\passthrough{\lstinline!VALUE\_TYPE!}}：当前绑定源码已实现该消息值操作。

\textbf{参数}

无显式参数。实例方法中的 \passthrough{\lstinline!self!} 由 Python
自动传入。

\textbf{返回值}

\begin{longtable}[]{@{}
  >{\raggedright\arraybackslash}p{(\columnwidth - 4\tabcolsep) * \real{0.18}}
  >{\raggedright\arraybackslash}p{(\columnwidth - 4\tabcolsep) * \real{0.20}}
  >{\raggedright\arraybackslash}p{(\columnwidth - 4\tabcolsep) * \real{0.62}}@{}}
\toprule
位置 & 类型 & 含义 \\
\midrule
\endhead
无 & \passthrough{\lstinline!None!} &
\passthrough{\lstinline!\_\_init\_\_!}
本身不返回值；调用类对象时，在构造函数可用的前提下得到该类实例。 \\
\bottomrule
\end{longtable}

\textbf{对应 C++}

\begin{itemize}
\tightlist
\item
  类：\passthrough{\lstinline!unitree\_go::msg::dds\_::LidarState\_!}
\item
  签名：\passthrough{\lstinline!LidarState\_()!}
\item
  绑定策略：\passthrough{\lstinline!未单独标注!}
\item
  声明位置：\passthrough{\lstinline!include/unitree/idl/go2/LidarState\_.hpp:45!}
\end{itemize}

\textbf{说明}

\begin{itemize}
\tightlist
\item
  示例是局部调用片段；\passthrough{\lstinline!obj!}
  和参数变量表示已经按上层业务校验并准备好的对象。
\end{itemize}

\textbf{用法}

\begin{lstlisting}[language=Python]
value = LidarState()
\end{lstlisting}

\hypertarget{unitree-sdk2-cpp-idl-go2-lidarstate-dunder-eq-1}{%
\paragraph{\texorpdfstring{\texttt{unitree\_sdk2\_cpp.idl.go2.LidarState.\_\_eq\_\_}}{unitree\_sdk2\_cpp.idl.go2.LidarState.\_\_eq\_\_}}\label{unitree-sdk2-cpp-idl-go2-lidarstate-dunder-eq-1}}

按底层 IDL 消息内容比较两个对象是否相等。

\textbf{签名}

\begin{lstlisting}[language=Python]
def __eq__(self, other: object) -> bool
\end{lstlisting}

\textbf{可用性与安全性}

\textbf{\passthrough{\lstinline!AVAILABLE!} /
\passthrough{\lstinline!VALUE\_TYPE!}}：当前绑定源码已实现该消息值操作。

\textbf{参数}

\begin{longtable}[]{@{}
  >{\raggedright\arraybackslash}p{(\columnwidth - 6\tabcolsep) * \real{0.18}}
  >{\raggedright\arraybackslash}p{(\columnwidth - 6\tabcolsep) * \real{0.24}}
  >{\raggedright\arraybackslash}p{(\columnwidth - 6\tabcolsep) * \real{0.18}}
  >{\raggedright\arraybackslash}p{(\columnwidth - 6\tabcolsep) * \real{0.40}}@{}}
\toprule
名称 & Python 类型 & 默认值 & 含义 \\
\midrule
\endhead
\passthrough{\lstinline!other!} & \passthrough{\lstinline!object!} &
必填 & 用于比较的另一个 Python 对象。类型不兼容时比较结果通常为
\passthrough{\lstinline!False!}。 \\
\bottomrule
\end{longtable}

\textbf{返回值}

\begin{longtable}[]{@{}
  >{\raggedright\arraybackslash}p{(\columnwidth - 4\tabcolsep) * \real{0.18}}
  >{\raggedright\arraybackslash}p{(\columnwidth - 4\tabcolsep) * \real{0.20}}
  >{\raggedright\arraybackslash}p{(\columnwidth - 4\tabcolsep) * \real{0.62}}@{}}
\toprule
位置 & 类型 & 含义 \\
\midrule
\endhead
返回值 & \passthrough{\lstinline!bool!} & 返回
\passthrough{\lstinline!bool!}。更精确的业务含义见该方法用途、C++
签名和目标型号协议。 \\
\bottomrule
\end{longtable}

\textbf{对应 C++}

\begin{itemize}
\tightlist
\item
  类：\passthrough{\lstinline!unitree\_go::msg::dds\_::LidarState\_!}
\item
  签名：\passthrough{\lstinline!operator==(const value \&) const!}
\item
  绑定策略：\passthrough{\lstinline!未单独标注!}
\item
  声明位置：由 pybind11 手工绑定或生成注册表提供；manifest
  未记录独立头文件行号。
\end{itemize}

\textbf{说明}

\begin{itemize}
\tightlist
\item
  示例是局部调用片段；\passthrough{\lstinline!obj!}
  和参数变量表示已经按上层业务校验并准备好的对象。
\end{itemize}

\textbf{用法}

\begin{lstlisting}[language=Python]
same = left == right
\end{lstlisting}

\hypertarget{unitree-sdk2-cpp-idl-go2-lidarstate-dunder-ne-1}{%
\paragraph{\texorpdfstring{\texttt{unitree\_sdk2\_cpp.idl.go2.LidarState.\_\_ne\_\_}}{unitree\_sdk2\_cpp.idl.go2.LidarState.\_\_ne\_\_}}\label{unitree-sdk2-cpp-idl-go2-lidarstate-dunder-ne-1}}

按底层 IDL 消息内容比较两个对象是否不相等。

\textbf{签名}

\begin{lstlisting}[language=Python]
def __ne__(self, other: object) -> bool
\end{lstlisting}

\textbf{可用性与安全性}

\textbf{\passthrough{\lstinline!AVAILABLE!} /
\passthrough{\lstinline!VALUE\_TYPE!}}：当前绑定源码已实现该消息值操作。

\textbf{参数}

\begin{longtable}[]{@{}
  >{\raggedright\arraybackslash}p{(\columnwidth - 6\tabcolsep) * \real{0.18}}
  >{\raggedright\arraybackslash}p{(\columnwidth - 6\tabcolsep) * \real{0.24}}
  >{\raggedright\arraybackslash}p{(\columnwidth - 6\tabcolsep) * \real{0.18}}
  >{\raggedright\arraybackslash}p{(\columnwidth - 6\tabcolsep) * \real{0.40}}@{}}
\toprule
名称 & Python 类型 & 默认值 & 含义 \\
\midrule
\endhead
\passthrough{\lstinline!other!} & \passthrough{\lstinline!object!} &
必填 & 用于比较的另一个 Python 对象。类型不兼容时比较结果通常为
\passthrough{\lstinline!False!}。 \\
\bottomrule
\end{longtable}

\textbf{返回值}

\begin{longtable}[]{@{}
  >{\raggedright\arraybackslash}p{(\columnwidth - 4\tabcolsep) * \real{0.18}}
  >{\raggedright\arraybackslash}p{(\columnwidth - 4\tabcolsep) * \real{0.20}}
  >{\raggedright\arraybackslash}p{(\columnwidth - 4\tabcolsep) * \real{0.62}}@{}}
\toprule
位置 & 类型 & 含义 \\
\midrule
\endhead
返回值 & \passthrough{\lstinline!bool!} & 返回
\passthrough{\lstinline!bool!}。更精确的业务含义见该方法用途、C++
签名和目标型号协议。 \\
\bottomrule
\end{longtable}

\textbf{对应 C++}

\begin{itemize}
\tightlist
\item
  类：\passthrough{\lstinline!unitree\_go::msg::dds\_::LidarState\_!}
\item
  签名：\passthrough{\lstinline"operator!=(const value \&) const"}
\item
  绑定策略：\passthrough{\lstinline!未单独标注!}
\item
  声明位置：由 pybind11 手工绑定或生成注册表提供；manifest
  未记录独立头文件行号。
\end{itemize}

\textbf{说明}

\begin{itemize}
\tightlist
\item
  示例是局部调用片段；\passthrough{\lstinline!obj!}
  和参数变量表示已经按上层业务校验并准备好的对象。
\end{itemize}

\textbf{用法}

\begin{lstlisting}[language=Python]
different = left != right
\end{lstlisting}

\hypertarget{unitree-sdk2-cpp-idl-go2-lidarstate-stamp}{%
\paragraph{\texorpdfstring{\texttt{unitree\_sdk2\_cpp.idl.go2.LidarState.stamp}}{unitree\_sdk2\_cpp.idl.go2.LidarState.stamp}}\label{unitree-sdk2-cpp-idl-go2-lidarstate-stamp}}

底层 \passthrough{\lstinline!unitree\_go::msg::dds\_::LidarState\_!} 的
\passthrough{\lstinline!stamp!} 字段。类型签名只说明 Python
表示；单位、数值范围、数组固定长度和枚举含义需查对应 IDL 头文件。
底层为浮点数；类型本身不说明单位、坐标系或业务安全范围。

\textbf{签名}

\begin{lstlisting}[language=Python]
@property
def stamp(self) -> float

@stamp.setter
def stamp(self, value: float) -> None
\end{lstlisting}

\textbf{可用性与安全性}

\textbf{\passthrough{\lstinline!AVAILABLE!} /
\passthrough{\lstinline!VALUE\_TYPE!}}：当前绑定源码已实现该消息值操作。

\textbf{参数}

\begin{longtable}[]{@{}
  >{\raggedright\arraybackslash}p{(\columnwidth - 4\tabcolsep) * \real{0.18}}
  >{\raggedright\arraybackslash}p{(\columnwidth - 4\tabcolsep) * \real{0.20}}
  >{\raggedright\arraybackslash}p{(\columnwidth - 4\tabcolsep) * \real{0.62}}@{}}
\toprule
名称 & Python 类型 & 含义 \\
\midrule
\endhead
\passthrough{\lstinline!value!} & \passthrough{\lstinline!float!} &
要写入或传递的新值，必须符合该参数的 Python 类型以及底层 C++
范围约束。 \\
\bottomrule
\end{longtable}

\textbf{返回值}

\begin{longtable}[]{@{}
  >{\raggedright\arraybackslash}p{(\columnwidth - 4\tabcolsep) * \real{0.18}}
  >{\raggedright\arraybackslash}p{(\columnwidth - 4\tabcolsep) * \real{0.20}}
  >{\raggedright\arraybackslash}p{(\columnwidth - 4\tabcolsep) * \real{0.62}}@{}}
\toprule
位置 & 类型 & 含义 \\
\midrule
\endhead
getter 返回值 & \passthrough{\lstinline!float!} & 返回属性当前值。 \\
\bottomrule
\end{longtable}

\textbf{对应 C++}

\begin{itemize}
\tightlist
\item
  类：\passthrough{\lstinline!unitree\_go::msg::dds\_::LidarState\_!}
\item
  字段类型：\passthrough{\lstinline!double!}
\item
  声明位置：\passthrough{\lstinline!include/unitree/idl/go2/LidarState\_.hpp:25!}
\end{itemize}

\textbf{用法}

\begin{lstlisting}[language=Python]
current_value = obj.stamp
obj.stamp = new_value
\end{lstlisting}

\hypertarget{unitree-sdk2-cpp-idl-go2-lidarstate-firmware-version}{%
\paragraph{\texorpdfstring{\texttt{unitree\_sdk2\_cpp.idl.go2.LidarState.firmware\_version}}{unitree\_sdk2\_cpp.idl.go2.LidarState.firmware\_version}}\label{unitree-sdk2-cpp-idl-go2-lidarstate-firmware-version}}

底层 \passthrough{\lstinline!unitree\_go::msg::dds\_::LidarState\_!} 的
\passthrough{\lstinline!firmware\_version!} 字段。类型签名只说明 Python
表示；单位、数值范围、数组固定长度和枚举含义需查对应 IDL 头文件。
底层为字符串；长度、编码和允许值由具体协议决定。

\textbf{签名}

\begin{lstlisting}[language=Python]
@property
def firmware_version(self) -> str

@firmware_version.setter
def firmware_version(self, value: str) -> None
\end{lstlisting}

\textbf{可用性与安全性}

\textbf{\passthrough{\lstinline!AVAILABLE!} /
\passthrough{\lstinline!VALUE\_TYPE!}}：当前绑定源码已实现该消息值操作。

\textbf{参数}

\begin{longtable}[]{@{}
  >{\raggedright\arraybackslash}p{(\columnwidth - 4\tabcolsep) * \real{0.18}}
  >{\raggedright\arraybackslash}p{(\columnwidth - 4\tabcolsep) * \real{0.20}}
  >{\raggedright\arraybackslash}p{(\columnwidth - 4\tabcolsep) * \real{0.62}}@{}}
\toprule
名称 & Python 类型 & 含义 \\
\midrule
\endhead
\passthrough{\lstinline!value!} & \passthrough{\lstinline!str!} &
要写入或传递的新值，必须符合该参数的 Python 类型以及底层 C++
范围约束。 \\
\bottomrule
\end{longtable}

\textbf{返回值}

\begin{longtable}[]{@{}
  >{\raggedright\arraybackslash}p{(\columnwidth - 4\tabcolsep) * \real{0.18}}
  >{\raggedright\arraybackslash}p{(\columnwidth - 4\tabcolsep) * \real{0.20}}
  >{\raggedright\arraybackslash}p{(\columnwidth - 4\tabcolsep) * \real{0.62}}@{}}
\toprule
位置 & 类型 & 含义 \\
\midrule
\endhead
getter 返回值 & \passthrough{\lstinline!str!} & 返回属性当前值。 \\
\bottomrule
\end{longtable}

\textbf{对应 C++}

\begin{itemize}
\tightlist
\item
  类：\passthrough{\lstinline!unitree\_go::msg::dds\_::LidarState\_!}
\item
  字段类型：\passthrough{\lstinline!std::string!}
\item
  声明位置：\passthrough{\lstinline!include/unitree/idl/go2/LidarState\_.hpp:26!}
\end{itemize}

\textbf{用法}

\begin{lstlisting}[language=Python]
current_value = obj.firmware_version
obj.firmware_version = new_value
\end{lstlisting}

\hypertarget{unitree-sdk2-cpp-idl-go2-lidarstate-software-version}{%
\paragraph{\texorpdfstring{\texttt{unitree\_sdk2\_cpp.idl.go2.LidarState.software\_version}}{unitree\_sdk2\_cpp.idl.go2.LidarState.software\_version}}\label{unitree-sdk2-cpp-idl-go2-lidarstate-software-version}}

底层 \passthrough{\lstinline!unitree\_go::msg::dds\_::LidarState\_!} 的
\passthrough{\lstinline!software\_version!} 字段。类型签名只说明 Python
表示；单位、数值范围、数组固定长度和枚举含义需查对应 IDL 头文件。
底层为字符串；长度、编码和允许值由具体协议决定。

\textbf{签名}

\begin{lstlisting}[language=Python]
@property
def software_version(self) -> str

@software_version.setter
def software_version(self, value: str) -> None
\end{lstlisting}

\textbf{可用性与安全性}

\textbf{\passthrough{\lstinline!AVAILABLE!} /
\passthrough{\lstinline!VALUE\_TYPE!}}：当前绑定源码已实现该消息值操作。

\textbf{参数}

\begin{longtable}[]{@{}
  >{\raggedright\arraybackslash}p{(\columnwidth - 4\tabcolsep) * \real{0.18}}
  >{\raggedright\arraybackslash}p{(\columnwidth - 4\tabcolsep) * \real{0.20}}
  >{\raggedright\arraybackslash}p{(\columnwidth - 4\tabcolsep) * \real{0.62}}@{}}
\toprule
名称 & Python 类型 & 含义 \\
\midrule
\endhead
\passthrough{\lstinline!value!} & \passthrough{\lstinline!str!} &
要写入或传递的新值，必须符合该参数的 Python 类型以及底层 C++
范围约束。 \\
\bottomrule
\end{longtable}

\textbf{返回值}

\begin{longtable}[]{@{}
  >{\raggedright\arraybackslash}p{(\columnwidth - 4\tabcolsep) * \real{0.18}}
  >{\raggedright\arraybackslash}p{(\columnwidth - 4\tabcolsep) * \real{0.20}}
  >{\raggedright\arraybackslash}p{(\columnwidth - 4\tabcolsep) * \real{0.62}}@{}}
\toprule
位置 & 类型 & 含义 \\
\midrule
\endhead
getter 返回值 & \passthrough{\lstinline!str!} & 返回属性当前值。 \\
\bottomrule
\end{longtable}

\textbf{对应 C++}

\begin{itemize}
\tightlist
\item
  类：\passthrough{\lstinline!unitree\_go::msg::dds\_::LidarState\_!}
\item
  字段类型：\passthrough{\lstinline!std::string!}
\item
  声明位置：\passthrough{\lstinline!include/unitree/idl/go2/LidarState\_.hpp:27!}
\end{itemize}

\textbf{用法}

\begin{lstlisting}[language=Python]
current_value = obj.software_version
obj.software_version = new_value
\end{lstlisting}

\hypertarget{unitree-sdk2-cpp-idl-go2-lidarstate-sdk-version}{%
\paragraph{\texorpdfstring{\texttt{unitree\_sdk2\_cpp.idl.go2.LidarState.sdk\_version}}{unitree\_sdk2\_cpp.idl.go2.LidarState.sdk\_version}}\label{unitree-sdk2-cpp-idl-go2-lidarstate-sdk-version}}

底层 \passthrough{\lstinline!unitree\_go::msg::dds\_::LidarState\_!} 的
\passthrough{\lstinline!sdk\_version!} 字段。类型签名只说明 Python
表示；单位、数值范围、数组固定长度和枚举含义需查对应 IDL 头文件。
底层为字符串；长度、编码和允许值由具体协议决定。

\textbf{签名}

\begin{lstlisting}[language=Python]
@property
def sdk_version(self) -> str

@sdk_version.setter
def sdk_version(self, value: str) -> None
\end{lstlisting}

\textbf{可用性与安全性}

\textbf{\passthrough{\lstinline!AVAILABLE!} /
\passthrough{\lstinline!VALUE\_TYPE!}}：当前绑定源码已实现该消息值操作。

\textbf{参数}

\begin{longtable}[]{@{}
  >{\raggedright\arraybackslash}p{(\columnwidth - 4\tabcolsep) * \real{0.18}}
  >{\raggedright\arraybackslash}p{(\columnwidth - 4\tabcolsep) * \real{0.20}}
  >{\raggedright\arraybackslash}p{(\columnwidth - 4\tabcolsep) * \real{0.62}}@{}}
\toprule
名称 & Python 类型 & 含义 \\
\midrule
\endhead
\passthrough{\lstinline!value!} & \passthrough{\lstinline!str!} &
要写入或传递的新值，必须符合该参数的 Python 类型以及底层 C++
范围约束。 \\
\bottomrule
\end{longtable}

\textbf{返回值}

\begin{longtable}[]{@{}
  >{\raggedright\arraybackslash}p{(\columnwidth - 4\tabcolsep) * \real{0.18}}
  >{\raggedright\arraybackslash}p{(\columnwidth - 4\tabcolsep) * \real{0.20}}
  >{\raggedright\arraybackslash}p{(\columnwidth - 4\tabcolsep) * \real{0.62}}@{}}
\toprule
位置 & 类型 & 含义 \\
\midrule
\endhead
getter 返回值 & \passthrough{\lstinline!str!} & 返回属性当前值。 \\
\bottomrule
\end{longtable}

\textbf{对应 C++}

\begin{itemize}
\tightlist
\item
  类：\passthrough{\lstinline!unitree\_go::msg::dds\_::LidarState\_!}
\item
  字段类型：\passthrough{\lstinline!std::string!}
\item
  声明位置：\passthrough{\lstinline!include/unitree/idl/go2/LidarState\_.hpp:28!}
\end{itemize}

\textbf{用法}

\begin{lstlisting}[language=Python]
current_value = obj.sdk_version
obj.sdk_version = new_value
\end{lstlisting}

\hypertarget{unitree-sdk2-cpp-idl-go2-lidarstate-sys-rotation-speed}{%
\paragraph{\texorpdfstring{\texttt{unitree\_sdk2\_cpp.idl.go2.LidarState.sys\_rotation\_speed}}{unitree\_sdk2\_cpp.idl.go2.LidarState.sys\_rotation\_speed}}\label{unitree-sdk2-cpp-idl-go2-lidarstate-sys-rotation-speed}}

底层 \passthrough{\lstinline!unitree\_go::msg::dds\_::LidarState\_!} 的
\passthrough{\lstinline!sys\_rotation\_speed!} 字段。类型签名只说明
Python 表示；单位、数值范围、数组固定长度和枚举含义需查对应 IDL 头文件。
底层为浮点数；类型本身不说明单位、坐标系或业务安全范围。

\textbf{签名}

\begin{lstlisting}[language=Python]
@property
def sys_rotation_speed(self) -> float

@sys_rotation_speed.setter
def sys_rotation_speed(self, value: float) -> None
\end{lstlisting}

\textbf{可用性与安全性}

\textbf{\passthrough{\lstinline!AVAILABLE!} /
\passthrough{\lstinline!VALUE\_TYPE!}}：当前绑定源码已实现该消息值操作。

\textbf{参数}

\begin{longtable}[]{@{}
  >{\raggedright\arraybackslash}p{(\columnwidth - 4\tabcolsep) * \real{0.18}}
  >{\raggedright\arraybackslash}p{(\columnwidth - 4\tabcolsep) * \real{0.20}}
  >{\raggedright\arraybackslash}p{(\columnwidth - 4\tabcolsep) * \real{0.62}}@{}}
\toprule
名称 & Python 类型 & 含义 \\
\midrule
\endhead
\passthrough{\lstinline!value!} & \passthrough{\lstinline!float!} &
要写入或传递的新值，必须符合该参数的 Python 类型以及底层 C++
范围约束。 \\
\bottomrule
\end{longtable}

\textbf{返回值}

\begin{longtable}[]{@{}
  >{\raggedright\arraybackslash}p{(\columnwidth - 4\tabcolsep) * \real{0.18}}
  >{\raggedright\arraybackslash}p{(\columnwidth - 4\tabcolsep) * \real{0.20}}
  >{\raggedright\arraybackslash}p{(\columnwidth - 4\tabcolsep) * \real{0.62}}@{}}
\toprule
位置 & 类型 & 含义 \\
\midrule
\endhead
getter 返回值 & \passthrough{\lstinline!float!} & 返回属性当前值。 \\
\bottomrule
\end{longtable}

\textbf{对应 C++}

\begin{itemize}
\tightlist
\item
  类：\passthrough{\lstinline!unitree\_go::msg::dds\_::LidarState\_!}
\item
  字段类型：\passthrough{\lstinline!float!}
\item
  声明位置：\passthrough{\lstinline!include/unitree/idl/go2/LidarState\_.hpp:29!}
\end{itemize}

\textbf{用法}

\begin{lstlisting}[language=Python]
current_value = obj.sys_rotation_speed
obj.sys_rotation_speed = new_value
\end{lstlisting}

\hypertarget{unitree-sdk2-cpp-idl-go2-lidarstate-com-rotation-speed}{%
\paragraph{\texorpdfstring{\texttt{unitree\_sdk2\_cpp.idl.go2.LidarState.com\_rotation\_speed}}{unitree\_sdk2\_cpp.idl.go2.LidarState.com\_rotation\_speed}}\label{unitree-sdk2-cpp-idl-go2-lidarstate-com-rotation-speed}}

底层 \passthrough{\lstinline!unitree\_go::msg::dds\_::LidarState\_!} 的
\passthrough{\lstinline!com\_rotation\_speed!} 字段。类型签名只说明
Python 表示；单位、数值范围、数组固定长度和枚举含义需查对应 IDL 头文件。
底层为浮点数；类型本身不说明单位、坐标系或业务安全范围。

\textbf{签名}

\begin{lstlisting}[language=Python]
@property
def com_rotation_speed(self) -> float

@com_rotation_speed.setter
def com_rotation_speed(self, value: float) -> None
\end{lstlisting}

\textbf{可用性与安全性}

\textbf{\passthrough{\lstinline!AVAILABLE!} /
\passthrough{\lstinline!VALUE\_TYPE!}}：当前绑定源码已实现该消息值操作。

\textbf{参数}

\begin{longtable}[]{@{}
  >{\raggedright\arraybackslash}p{(\columnwidth - 4\tabcolsep) * \real{0.18}}
  >{\raggedright\arraybackslash}p{(\columnwidth - 4\tabcolsep) * \real{0.20}}
  >{\raggedright\arraybackslash}p{(\columnwidth - 4\tabcolsep) * \real{0.62}}@{}}
\toprule
名称 & Python 类型 & 含义 \\
\midrule
\endhead
\passthrough{\lstinline!value!} & \passthrough{\lstinline!float!} &
要写入或传递的新值，必须符合该参数的 Python 类型以及底层 C++
范围约束。 \\
\bottomrule
\end{longtable}

\textbf{返回值}

\begin{longtable}[]{@{}
  >{\raggedright\arraybackslash}p{(\columnwidth - 4\tabcolsep) * \real{0.18}}
  >{\raggedright\arraybackslash}p{(\columnwidth - 4\tabcolsep) * \real{0.20}}
  >{\raggedright\arraybackslash}p{(\columnwidth - 4\tabcolsep) * \real{0.62}}@{}}
\toprule
位置 & 类型 & 含义 \\
\midrule
\endhead
getter 返回值 & \passthrough{\lstinline!float!} & 返回属性当前值。 \\
\bottomrule
\end{longtable}

\textbf{对应 C++}

\begin{itemize}
\tightlist
\item
  类：\passthrough{\lstinline!unitree\_go::msg::dds\_::LidarState\_!}
\item
  字段类型：\passthrough{\lstinline!float!}
\item
  声明位置：\passthrough{\lstinline!include/unitree/idl/go2/LidarState\_.hpp:30!}
\end{itemize}

\textbf{用法}

\begin{lstlisting}[language=Python]
current_value = obj.com_rotation_speed
obj.com_rotation_speed = new_value
\end{lstlisting}

\hypertarget{unitree-sdk2-cpp-idl-go2-lidarstate-error-state}{%
\paragraph{\texorpdfstring{\texttt{unitree\_sdk2\_cpp.idl.go2.LidarState.error\_state}}{unitree\_sdk2\_cpp.idl.go2.LidarState.error\_state}}\label{unitree-sdk2-cpp-idl-go2-lidarstate-error-state}}

底层 \passthrough{\lstinline!unitree\_go::msg::dds\_::LidarState\_!} 的
\passthrough{\lstinline!error\_state!} 字段。类型签名只说明 Python
表示；单位、数值范围、数组固定长度和枚举含义需查对应 IDL 头文件。
取值范围为 0 到 255。

\textbf{签名}

\begin{lstlisting}[language=Python]
@property
def error_state(self) -> int

@error_state.setter
def error_state(self, value: int) -> None
\end{lstlisting}

\textbf{可用性与安全性}

\textbf{\passthrough{\lstinline!AVAILABLE!} /
\passthrough{\lstinline!VALUE\_TYPE!}}：当前绑定源码已实现该消息值操作。

\textbf{参数}

\begin{longtable}[]{@{}
  >{\raggedright\arraybackslash}p{(\columnwidth - 4\tabcolsep) * \real{0.18}}
  >{\raggedright\arraybackslash}p{(\columnwidth - 4\tabcolsep) * \real{0.20}}
  >{\raggedright\arraybackslash}p{(\columnwidth - 4\tabcolsep) * \real{0.62}}@{}}
\toprule
名称 & Python 类型 & 含义 \\
\midrule
\endhead
\passthrough{\lstinline!value!} & \passthrough{\lstinline!int!} &
要写入或传递的新值，必须符合该参数的 Python 类型以及底层 C++
范围约束。 \\
\bottomrule
\end{longtable}

\textbf{返回值}

\begin{longtable}[]{@{}
  >{\raggedright\arraybackslash}p{(\columnwidth - 4\tabcolsep) * \real{0.18}}
  >{\raggedright\arraybackslash}p{(\columnwidth - 4\tabcolsep) * \real{0.20}}
  >{\raggedright\arraybackslash}p{(\columnwidth - 4\tabcolsep) * \real{0.62}}@{}}
\toprule
位置 & 类型 & 含义 \\
\midrule
\endhead
getter 返回值 & \passthrough{\lstinline!int!} & 返回属性当前值。 \\
\bottomrule
\end{longtable}

\textbf{对应 C++}

\begin{itemize}
\tightlist
\item
  类：\passthrough{\lstinline!unitree\_go::msg::dds\_::LidarState\_!}
\item
  字段类型：\passthrough{\lstinline!uint8\_t!}
\item
  声明位置：\passthrough{\lstinline!include/unitree/idl/go2/LidarState\_.hpp:31!}
\end{itemize}

\textbf{用法}

\begin{lstlisting}[language=Python]
current_value = obj.error_state
obj.error_state = new_value
\end{lstlisting}

\hypertarget{unitree-sdk2-cpp-idl-go2-lidarstate-dirty-percentage}{%
\paragraph{\texorpdfstring{\texttt{unitree\_sdk2\_cpp.idl.go2.LidarState.dirty\_percentage}}{unitree\_sdk2\_cpp.idl.go2.LidarState.dirty\_percentage}}\label{unitree-sdk2-cpp-idl-go2-lidarstate-dirty-percentage}}

底层 \passthrough{\lstinline!unitree\_go::msg::dds\_::LidarState\_!} 的
\passthrough{\lstinline!dirty\_percentage!} 字段。类型签名只说明 Python
表示；单位、数值范围、数组固定长度和枚举含义需查对应 IDL 头文件。
取值范围为 0 到 255。

\textbf{签名}

\begin{lstlisting}[language=Python]
@property
def dirty_percentage(self) -> int

@dirty_percentage.setter
def dirty_percentage(self, value: int) -> None
\end{lstlisting}

\textbf{可用性与安全性}

\textbf{\passthrough{\lstinline!AVAILABLE!} /
\passthrough{\lstinline!VALUE\_TYPE!}}：当前绑定源码已实现该消息值操作。

\textbf{参数}

\begin{longtable}[]{@{}
  >{\raggedright\arraybackslash}p{(\columnwidth - 4\tabcolsep) * \real{0.18}}
  >{\raggedright\arraybackslash}p{(\columnwidth - 4\tabcolsep) * \real{0.20}}
  >{\raggedright\arraybackslash}p{(\columnwidth - 4\tabcolsep) * \real{0.62}}@{}}
\toprule
名称 & Python 类型 & 含义 \\
\midrule
\endhead
\passthrough{\lstinline!value!} & \passthrough{\lstinline!int!} &
要写入或传递的新值，必须符合该参数的 Python 类型以及底层 C++
范围约束。 \\
\bottomrule
\end{longtable}

\textbf{返回值}

\begin{longtable}[]{@{}
  >{\raggedright\arraybackslash}p{(\columnwidth - 4\tabcolsep) * \real{0.18}}
  >{\raggedright\arraybackslash}p{(\columnwidth - 4\tabcolsep) * \real{0.20}}
  >{\raggedright\arraybackslash}p{(\columnwidth - 4\tabcolsep) * \real{0.62}}@{}}
\toprule
位置 & 类型 & 含义 \\
\midrule
\endhead
getter 返回值 & \passthrough{\lstinline!int!} & 返回属性当前值。 \\
\bottomrule
\end{longtable}

\textbf{对应 C++}

\begin{itemize}
\tightlist
\item
  类：\passthrough{\lstinline!unitree\_go::msg::dds\_::LidarState\_!}
\item
  字段类型：\passthrough{\lstinline!uint8\_t!}
\item
  声明位置：\passthrough{\lstinline!include/unitree/idl/go2/LidarState\_.hpp:32!}
\end{itemize}

\textbf{用法}

\begin{lstlisting}[language=Python]
current_value = obj.dirty_percentage
obj.dirty_percentage = new_value
\end{lstlisting}

\hypertarget{unitree-sdk2-cpp-idl-go2-lidarstate-cloud-frequency}{%
\paragraph{\texorpdfstring{\texttt{unitree\_sdk2\_cpp.idl.go2.LidarState.cloud\_frequency}}{unitree\_sdk2\_cpp.idl.go2.LidarState.cloud\_frequency}}\label{unitree-sdk2-cpp-idl-go2-lidarstate-cloud-frequency}}

底层 \passthrough{\lstinline!unitree\_go::msg::dds\_::LidarState\_!} 的
\passthrough{\lstinline!cloud\_frequency!} 字段。类型签名只说明 Python
表示；单位、数值范围、数组固定长度和枚举含义需查对应 IDL 头文件。
底层为浮点数；类型本身不说明单位、坐标系或业务安全范围。

\textbf{签名}

\begin{lstlisting}[language=Python]
@property
def cloud_frequency(self) -> float

@cloud_frequency.setter
def cloud_frequency(self, value: float) -> None
\end{lstlisting}

\textbf{可用性与安全性}

\textbf{\passthrough{\lstinline!AVAILABLE!} /
\passthrough{\lstinline!VALUE\_TYPE!}}：当前绑定源码已实现该消息值操作。

\textbf{参数}

\begin{longtable}[]{@{}
  >{\raggedright\arraybackslash}p{(\columnwidth - 4\tabcolsep) * \real{0.18}}
  >{\raggedright\arraybackslash}p{(\columnwidth - 4\tabcolsep) * \real{0.20}}
  >{\raggedright\arraybackslash}p{(\columnwidth - 4\tabcolsep) * \real{0.62}}@{}}
\toprule
名称 & Python 类型 & 含义 \\
\midrule
\endhead
\passthrough{\lstinline!value!} & \passthrough{\lstinline!float!} &
要写入或传递的新值，必须符合该参数的 Python 类型以及底层 C++
范围约束。 \\
\bottomrule
\end{longtable}

\textbf{返回值}

\begin{longtable}[]{@{}
  >{\raggedright\arraybackslash}p{(\columnwidth - 4\tabcolsep) * \real{0.18}}
  >{\raggedright\arraybackslash}p{(\columnwidth - 4\tabcolsep) * \real{0.20}}
  >{\raggedright\arraybackslash}p{(\columnwidth - 4\tabcolsep) * \real{0.62}}@{}}
\toprule
位置 & 类型 & 含义 \\
\midrule
\endhead
getter 返回值 & \passthrough{\lstinline!float!} & 返回属性当前值。 \\
\bottomrule
\end{longtable}

\textbf{对应 C++}

\begin{itemize}
\tightlist
\item
  类：\passthrough{\lstinline!unitree\_go::msg::dds\_::LidarState\_!}
\item
  字段类型：\passthrough{\lstinline!float!}
\item
  声明位置：\passthrough{\lstinline!include/unitree/idl/go2/LidarState\_.hpp:33!}
\end{itemize}

\textbf{用法}

\begin{lstlisting}[language=Python]
current_value = obj.cloud_frequency
obj.cloud_frequency = new_value
\end{lstlisting}

\hypertarget{unitree-sdk2-cpp-idl-go2-lidarstate-cloud-packet-loss-rate}{%
\paragraph{\texorpdfstring{\texttt{unitree\_sdk2\_cpp.idl.go2.LidarState.cloud\_packet\_loss\_rate}}{unitree\_sdk2\_cpp.idl.go2.LidarState.cloud\_packet\_loss\_rate}}\label{unitree-sdk2-cpp-idl-go2-lidarstate-cloud-packet-loss-rate}}

底层 \passthrough{\lstinline!unitree\_go::msg::dds\_::LidarState\_!} 的
\passthrough{\lstinline!cloud\_packet\_loss\_rate!} 字段。类型签名只说明
Python 表示；单位、数值范围、数组固定长度和枚举含义需查对应 IDL 头文件。
底层为浮点数；类型本身不说明单位、坐标系或业务安全范围。

\textbf{签名}

\begin{lstlisting}[language=Python]
@property
def cloud_packet_loss_rate(self) -> float

@cloud_packet_loss_rate.setter
def cloud_packet_loss_rate(self, value: float) -> None
\end{lstlisting}

\textbf{可用性与安全性}

\textbf{\passthrough{\lstinline!AVAILABLE!} /
\passthrough{\lstinline!VALUE\_TYPE!}}：当前绑定源码已实现该消息值操作。

\textbf{参数}

\begin{longtable}[]{@{}
  >{\raggedright\arraybackslash}p{(\columnwidth - 4\tabcolsep) * \real{0.18}}
  >{\raggedright\arraybackslash}p{(\columnwidth - 4\tabcolsep) * \real{0.20}}
  >{\raggedright\arraybackslash}p{(\columnwidth - 4\tabcolsep) * \real{0.62}}@{}}
\toprule
名称 & Python 类型 & 含义 \\
\midrule
\endhead
\passthrough{\lstinline!value!} & \passthrough{\lstinline!float!} &
要写入或传递的新值，必须符合该参数的 Python 类型以及底层 C++
范围约束。 \\
\bottomrule
\end{longtable}

\textbf{返回值}

\begin{longtable}[]{@{}
  >{\raggedright\arraybackslash}p{(\columnwidth - 4\tabcolsep) * \real{0.18}}
  >{\raggedright\arraybackslash}p{(\columnwidth - 4\tabcolsep) * \real{0.20}}
  >{\raggedright\arraybackslash}p{(\columnwidth - 4\tabcolsep) * \real{0.62}}@{}}
\toprule
位置 & 类型 & 含义 \\
\midrule
\endhead
getter 返回值 & \passthrough{\lstinline!float!} & 返回属性当前值。 \\
\bottomrule
\end{longtable}

\textbf{对应 C++}

\begin{itemize}
\tightlist
\item
  类：\passthrough{\lstinline!unitree\_go::msg::dds\_::LidarState\_!}
\item
  字段类型：\passthrough{\lstinline!float!}
\item
  声明位置：\passthrough{\lstinline!include/unitree/idl/go2/LidarState\_.hpp:34!}
\end{itemize}

\textbf{用法}

\begin{lstlisting}[language=Python]
current_value = obj.cloud_packet_loss_rate
obj.cloud_packet_loss_rate = new_value
\end{lstlisting}

\hypertarget{unitree-sdk2-cpp-idl-go2-lidarstate-cloud-size}{%
\paragraph{\texorpdfstring{\texttt{unitree\_sdk2\_cpp.idl.go2.LidarState.cloud\_size}}{unitree\_sdk2\_cpp.idl.go2.LidarState.cloud\_size}}\label{unitree-sdk2-cpp-idl-go2-lidarstate-cloud-size}}

底层 \passthrough{\lstinline!unitree\_go::msg::dds\_::LidarState\_!} 的
\passthrough{\lstinline!cloud\_size!} 字段。类型签名只说明 Python
表示；单位、数值范围、数组固定长度和枚举含义需查对应 IDL 头文件。
取值范围为 0 到 4294967295。

\textbf{签名}

\begin{lstlisting}[language=Python]
@property
def cloud_size(self) -> int

@cloud_size.setter
def cloud_size(self, value: int) -> None
\end{lstlisting}

\textbf{可用性与安全性}

\textbf{\passthrough{\lstinline!AVAILABLE!} /
\passthrough{\lstinline!VALUE\_TYPE!}}：当前绑定源码已实现该消息值操作。

\textbf{参数}

\begin{longtable}[]{@{}
  >{\raggedright\arraybackslash}p{(\columnwidth - 4\tabcolsep) * \real{0.18}}
  >{\raggedright\arraybackslash}p{(\columnwidth - 4\tabcolsep) * \real{0.20}}
  >{\raggedright\arraybackslash}p{(\columnwidth - 4\tabcolsep) * \real{0.62}}@{}}
\toprule
名称 & Python 类型 & 含义 \\
\midrule
\endhead
\passthrough{\lstinline!value!} & \passthrough{\lstinline!int!} &
要写入或传递的新值，必须符合该参数的 Python 类型以及底层 C++
范围约束。 \\
\bottomrule
\end{longtable}

\textbf{返回值}

\begin{longtable}[]{@{}
  >{\raggedright\arraybackslash}p{(\columnwidth - 4\tabcolsep) * \real{0.18}}
  >{\raggedright\arraybackslash}p{(\columnwidth - 4\tabcolsep) * \real{0.20}}
  >{\raggedright\arraybackslash}p{(\columnwidth - 4\tabcolsep) * \real{0.62}}@{}}
\toprule
位置 & 类型 & 含义 \\
\midrule
\endhead
getter 返回值 & \passthrough{\lstinline!int!} & 返回属性当前值。 \\
\bottomrule
\end{longtable}

\textbf{对应 C++}

\begin{itemize}
\tightlist
\item
  类：\passthrough{\lstinline!unitree\_go::msg::dds\_::LidarState\_!}
\item
  字段类型：\passthrough{\lstinline!uint32\_t!}
\item
  声明位置：\passthrough{\lstinline!include/unitree/idl/go2/LidarState\_.hpp:35!}
\end{itemize}

\textbf{用法}

\begin{lstlisting}[language=Python]
current_value = obj.cloud_size
obj.cloud_size = new_value
\end{lstlisting}

\hypertarget{unitree-sdk2-cpp-idl-go2-lidarstate-cloud-scan-num}{%
\paragraph{\texorpdfstring{\texttt{unitree\_sdk2\_cpp.idl.go2.LidarState.cloud\_scan\_num}}{unitree\_sdk2\_cpp.idl.go2.LidarState.cloud\_scan\_num}}\label{unitree-sdk2-cpp-idl-go2-lidarstate-cloud-scan-num}}

底层 \passthrough{\lstinline!unitree\_go::msg::dds\_::LidarState\_!} 的
\passthrough{\lstinline!cloud\_scan\_num!} 字段。类型签名只说明 Python
表示；单位、数值范围、数组固定长度和枚举含义需查对应 IDL 头文件。
取值范围为 0 到 4294967295。

\textbf{签名}

\begin{lstlisting}[language=Python]
@property
def cloud_scan_num(self) -> int

@cloud_scan_num.setter
def cloud_scan_num(self, value: int) -> None
\end{lstlisting}

\textbf{可用性与安全性}

\textbf{\passthrough{\lstinline!AVAILABLE!} /
\passthrough{\lstinline!VALUE\_TYPE!}}：当前绑定源码已实现该消息值操作。

\textbf{参数}

\begin{longtable}[]{@{}
  >{\raggedright\arraybackslash}p{(\columnwidth - 4\tabcolsep) * \real{0.18}}
  >{\raggedright\arraybackslash}p{(\columnwidth - 4\tabcolsep) * \real{0.20}}
  >{\raggedright\arraybackslash}p{(\columnwidth - 4\tabcolsep) * \real{0.62}}@{}}
\toprule
名称 & Python 类型 & 含义 \\
\midrule
\endhead
\passthrough{\lstinline!value!} & \passthrough{\lstinline!int!} &
要写入或传递的新值，必须符合该参数的 Python 类型以及底层 C++
范围约束。 \\
\bottomrule
\end{longtable}

\textbf{返回值}

\begin{longtable}[]{@{}
  >{\raggedright\arraybackslash}p{(\columnwidth - 4\tabcolsep) * \real{0.18}}
  >{\raggedright\arraybackslash}p{(\columnwidth - 4\tabcolsep) * \real{0.20}}
  >{\raggedright\arraybackslash}p{(\columnwidth - 4\tabcolsep) * \real{0.62}}@{}}
\toprule
位置 & 类型 & 含义 \\
\midrule
\endhead
getter 返回值 & \passthrough{\lstinline!int!} & 返回属性当前值。 \\
\bottomrule
\end{longtable}

\textbf{对应 C++}

\begin{itemize}
\tightlist
\item
  类：\passthrough{\lstinline!unitree\_go::msg::dds\_::LidarState\_!}
\item
  字段类型：\passthrough{\lstinline!uint32\_t!}
\item
  声明位置：\passthrough{\lstinline!include/unitree/idl/go2/LidarState\_.hpp:36!}
\end{itemize}

\textbf{用法}

\begin{lstlisting}[language=Python]
current_value = obj.cloud_scan_num
obj.cloud_scan_num = new_value
\end{lstlisting}

\hypertarget{unitree-sdk2-cpp-idl-go2-lidarstate-imu-frequency}{%
\paragraph{\texorpdfstring{\texttt{unitree\_sdk2\_cpp.idl.go2.LidarState.imu\_frequency}}{unitree\_sdk2\_cpp.idl.go2.LidarState.imu\_frequency}}\label{unitree-sdk2-cpp-idl-go2-lidarstate-imu-frequency}}

底层 \passthrough{\lstinline!unitree\_go::msg::dds\_::LidarState\_!} 的
\passthrough{\lstinline!imu\_frequency!} 字段。类型签名只说明 Python
表示；单位、数值范围、数组固定长度和枚举含义需查对应 IDL 头文件。
底层为浮点数；类型本身不说明单位、坐标系或业务安全范围。

\textbf{签名}

\begin{lstlisting}[language=Python]
@property
def imu_frequency(self) -> float

@imu_frequency.setter
def imu_frequency(self, value: float) -> None
\end{lstlisting}

\textbf{可用性与安全性}

\textbf{\passthrough{\lstinline!AVAILABLE!} /
\passthrough{\lstinline!VALUE\_TYPE!}}：当前绑定源码已实现该消息值操作。

\textbf{参数}

\begin{longtable}[]{@{}
  >{\raggedright\arraybackslash}p{(\columnwidth - 4\tabcolsep) * \real{0.18}}
  >{\raggedright\arraybackslash}p{(\columnwidth - 4\tabcolsep) * \real{0.20}}
  >{\raggedright\arraybackslash}p{(\columnwidth - 4\tabcolsep) * \real{0.62}}@{}}
\toprule
名称 & Python 类型 & 含义 \\
\midrule
\endhead
\passthrough{\lstinline!value!} & \passthrough{\lstinline!float!} &
要写入或传递的新值，必须符合该参数的 Python 类型以及底层 C++
范围约束。 \\
\bottomrule
\end{longtable}

\textbf{返回值}

\begin{longtable}[]{@{}
  >{\raggedright\arraybackslash}p{(\columnwidth - 4\tabcolsep) * \real{0.18}}
  >{\raggedright\arraybackslash}p{(\columnwidth - 4\tabcolsep) * \real{0.20}}
  >{\raggedright\arraybackslash}p{(\columnwidth - 4\tabcolsep) * \real{0.62}}@{}}
\toprule
位置 & 类型 & 含义 \\
\midrule
\endhead
getter 返回值 & \passthrough{\lstinline!float!} & 返回属性当前值。 \\
\bottomrule
\end{longtable}

\textbf{对应 C++}

\begin{itemize}
\tightlist
\item
  类：\passthrough{\lstinline!unitree\_go::msg::dds\_::LidarState\_!}
\item
  字段类型：\passthrough{\lstinline!float!}
\item
  声明位置：\passthrough{\lstinline!include/unitree/idl/go2/LidarState\_.hpp:37!}
\end{itemize}

\textbf{用法}

\begin{lstlisting}[language=Python]
current_value = obj.imu_frequency
obj.imu_frequency = new_value
\end{lstlisting}

\hypertarget{unitree-sdk2-cpp-idl-go2-lidarstate-imu-packet-loss-rate}{%
\paragraph{\texorpdfstring{\texttt{unitree\_sdk2\_cpp.idl.go2.LidarState.imu\_packet\_loss\_rate}}{unitree\_sdk2\_cpp.idl.go2.LidarState.imu\_packet\_loss\_rate}}\label{unitree-sdk2-cpp-idl-go2-lidarstate-imu-packet-loss-rate}}

底层 \passthrough{\lstinline!unitree\_go::msg::dds\_::LidarState\_!} 的
\passthrough{\lstinline!imu\_packet\_loss\_rate!} 字段。类型签名只说明
Python 表示；单位、数值范围、数组固定长度和枚举含义需查对应 IDL 头文件。
底层为浮点数；类型本身不说明单位、坐标系或业务安全范围。

\textbf{签名}

\begin{lstlisting}[language=Python]
@property
def imu_packet_loss_rate(self) -> float

@imu_packet_loss_rate.setter
def imu_packet_loss_rate(self, value: float) -> None
\end{lstlisting}

\textbf{可用性与安全性}

\textbf{\passthrough{\lstinline!AVAILABLE!} /
\passthrough{\lstinline!VALUE\_TYPE!}}：当前绑定源码已实现该消息值操作。

\textbf{参数}

\begin{longtable}[]{@{}
  >{\raggedright\arraybackslash}p{(\columnwidth - 4\tabcolsep) * \real{0.18}}
  >{\raggedright\arraybackslash}p{(\columnwidth - 4\tabcolsep) * \real{0.20}}
  >{\raggedright\arraybackslash}p{(\columnwidth - 4\tabcolsep) * \real{0.62}}@{}}
\toprule
名称 & Python 类型 & 含义 \\
\midrule
\endhead
\passthrough{\lstinline!value!} & \passthrough{\lstinline!float!} &
要写入或传递的新值，必须符合该参数的 Python 类型以及底层 C++
范围约束。 \\
\bottomrule
\end{longtable}

\textbf{返回值}

\begin{longtable}[]{@{}
  >{\raggedright\arraybackslash}p{(\columnwidth - 4\tabcolsep) * \real{0.18}}
  >{\raggedright\arraybackslash}p{(\columnwidth - 4\tabcolsep) * \real{0.20}}
  >{\raggedright\arraybackslash}p{(\columnwidth - 4\tabcolsep) * \real{0.62}}@{}}
\toprule
位置 & 类型 & 含义 \\
\midrule
\endhead
getter 返回值 & \passthrough{\lstinline!float!} & 返回属性当前值。 \\
\bottomrule
\end{longtable}

\textbf{对应 C++}

\begin{itemize}
\tightlist
\item
  类：\passthrough{\lstinline!unitree\_go::msg::dds\_::LidarState\_!}
\item
  字段类型：\passthrough{\lstinline!float!}
\item
  声明位置：\passthrough{\lstinline!include/unitree/idl/go2/LidarState\_.hpp:38!}
\end{itemize}

\textbf{用法}

\begin{lstlisting}[language=Python]
current_value = obj.imu_packet_loss_rate
obj.imu_packet_loss_rate = new_value
\end{lstlisting}

\hypertarget{unitree-sdk2-cpp-idl-go2-lidarstate-imu-rpy}{%
\paragraph{\texorpdfstring{\texttt{unitree\_sdk2\_cpp.idl.go2.LidarState.imu\_rpy}}{unitree\_sdk2\_cpp.idl.go2.LidarState.imu\_rpy}}\label{unitree-sdk2-cpp-idl-go2-lidarstate-imu-rpy}}

底层 \passthrough{\lstinline!unitree\_go::msg::dds\_::LidarState\_!} 的
\passthrough{\lstinline!imu\_rpy!} 字段。类型签名只说明 Python
表示；单位、数值范围、数组固定长度和枚举含义需查对应 IDL 头文件。
底层是固定长度数组，写入序列必须正好包含 3
个元素；元素约束：底层为浮点数；类型本身不说明单位、坐标系或业务安全范围。

\textbf{签名}

\begin{lstlisting}[language=Python]
@property
def imu_rpy(self) -> list[float]

@imu_rpy.setter
def imu_rpy(self, value: Sequence[float]) -> None
\end{lstlisting}

\textbf{可用性与安全性}

\textbf{\passthrough{\lstinline!AVAILABLE!} /
\passthrough{\lstinline!VALUE\_TYPE!}}：当前绑定源码已实现该消息值操作。

\textbf{参数}

\begin{longtable}[]{@{}
  >{\raggedright\arraybackslash}p{(\columnwidth - 4\tabcolsep) * \real{0.18}}
  >{\raggedright\arraybackslash}p{(\columnwidth - 4\tabcolsep) * \real{0.20}}
  >{\raggedright\arraybackslash}p{(\columnwidth - 4\tabcolsep) * \real{0.62}}@{}}
\toprule
名称 & Python 类型 & 含义 \\
\midrule
\endhead
\passthrough{\lstinline!value!} &
\passthrough{\lstinline!Sequence[float]!} &
要写入或传递的新值，必须符合该参数的 Python 类型以及底层 C++
范围约束。 \\
\bottomrule
\end{longtable}

\textbf{返回值}

\begin{longtable}[]{@{}
  >{\raggedright\arraybackslash}p{(\columnwidth - 4\tabcolsep) * \real{0.18}}
  >{\raggedright\arraybackslash}p{(\columnwidth - 4\tabcolsep) * \real{0.20}}
  >{\raggedright\arraybackslash}p{(\columnwidth - 4\tabcolsep) * \real{0.62}}@{}}
\toprule
位置 & 类型 & 含义 \\
\midrule
\endhead
getter 返回值 & \passthrough{\lstinline!list[float]!} &
返回属性当前值。数组、vector 和嵌套消息采用复制语义。 \\
\bottomrule
\end{longtable}

\textbf{对应 C++}

\begin{itemize}
\tightlist
\item
  类：\passthrough{\lstinline!unitree\_go::msg::dds\_::LidarState\_!}
\item
  字段类型：\passthrough{\lstinline!std::array<float, 3>!}
\item
  声明位置：\passthrough{\lstinline!include/unitree/idl/go2/LidarState\_.hpp:39!}
\end{itemize}

\textbf{用法}

\begin{lstlisting}[language=Python]
current_value = obj.imu_rpy
obj.imu_rpy = new_value
\end{lstlisting}

\begin{quote}
{[}!TIP{]} 读取容器或嵌套值后应遵循``读取、修改、重新赋值''。只修改
getter 返回的副本不会可靠地更新原 C++ 消息。
\end{quote}

\hypertarget{unitree-sdk2-cpp-idl-go2-lidarstate-serial-recv-stamp}{%
\paragraph{\texorpdfstring{\texttt{unitree\_sdk2\_cpp.idl.go2.LidarState.serial\_recv\_stamp}}{unitree\_sdk2\_cpp.idl.go2.LidarState.serial\_recv\_stamp}}\label{unitree-sdk2-cpp-idl-go2-lidarstate-serial-recv-stamp}}

底层 \passthrough{\lstinline!unitree\_go::msg::dds\_::LidarState\_!} 的
\passthrough{\lstinline!serial\_recv\_stamp!} 字段。类型签名只说明
Python 表示；单位、数值范围、数组固定长度和枚举含义需查对应 IDL 头文件。
底层为浮点数；类型本身不说明单位、坐标系或业务安全范围。

\textbf{签名}

\begin{lstlisting}[language=Python]
@property
def serial_recv_stamp(self) -> float

@serial_recv_stamp.setter
def serial_recv_stamp(self, value: float) -> None
\end{lstlisting}

\textbf{可用性与安全性}

\textbf{\passthrough{\lstinline!AVAILABLE!} /
\passthrough{\lstinline!VALUE\_TYPE!}}：当前绑定源码已实现该消息值操作。

\textbf{参数}

\begin{longtable}[]{@{}
  >{\raggedright\arraybackslash}p{(\columnwidth - 4\tabcolsep) * \real{0.18}}
  >{\raggedright\arraybackslash}p{(\columnwidth - 4\tabcolsep) * \real{0.20}}
  >{\raggedright\arraybackslash}p{(\columnwidth - 4\tabcolsep) * \real{0.62}}@{}}
\toprule
名称 & Python 类型 & 含义 \\
\midrule
\endhead
\passthrough{\lstinline!value!} & \passthrough{\lstinline!float!} &
要写入或传递的新值，必须符合该参数的 Python 类型以及底层 C++
范围约束。 \\
\bottomrule
\end{longtable}

\textbf{返回值}

\begin{longtable}[]{@{}
  >{\raggedright\arraybackslash}p{(\columnwidth - 4\tabcolsep) * \real{0.18}}
  >{\raggedright\arraybackslash}p{(\columnwidth - 4\tabcolsep) * \real{0.20}}
  >{\raggedright\arraybackslash}p{(\columnwidth - 4\tabcolsep) * \real{0.62}}@{}}
\toprule
位置 & 类型 & 含义 \\
\midrule
\endhead
getter 返回值 & \passthrough{\lstinline!float!} & 返回属性当前值。 \\
\bottomrule
\end{longtable}

\textbf{对应 C++}

\begin{itemize}
\tightlist
\item
  类：\passthrough{\lstinline!unitree\_go::msg::dds\_::LidarState\_!}
\item
  字段类型：\passthrough{\lstinline!double!}
\item
  声明位置：\passthrough{\lstinline!include/unitree/idl/go2/LidarState\_.hpp:40!}
\end{itemize}

\textbf{用法}

\begin{lstlisting}[language=Python]
current_value = obj.serial_recv_stamp
obj.serial_recv_stamp = new_value
\end{lstlisting}

\hypertarget{unitree-sdk2-cpp-idl-go2-lidarstate-serial-buffer-size}{%
\paragraph{\texorpdfstring{\texttt{unitree\_sdk2\_cpp.idl.go2.LidarState.serial\_buffer\_size}}{unitree\_sdk2\_cpp.idl.go2.LidarState.serial\_buffer\_size}}\label{unitree-sdk2-cpp-idl-go2-lidarstate-serial-buffer-size}}

底层 \passthrough{\lstinline!unitree\_go::msg::dds\_::LidarState\_!} 的
\passthrough{\lstinline!serial\_buffer\_size!} 字段。类型签名只说明
Python 表示；单位、数值范围、数组固定长度和枚举含义需查对应 IDL 头文件。
取值范围为 0 到 4294967295。

\textbf{签名}

\begin{lstlisting}[language=Python]
@property
def serial_buffer_size(self) -> int

@serial_buffer_size.setter
def serial_buffer_size(self, value: int) -> None
\end{lstlisting}

\textbf{可用性与安全性}

\textbf{\passthrough{\lstinline!AVAILABLE!} /
\passthrough{\lstinline!VALUE\_TYPE!}}：当前绑定源码已实现该消息值操作。

\textbf{参数}

\begin{longtable}[]{@{}
  >{\raggedright\arraybackslash}p{(\columnwidth - 4\tabcolsep) * \real{0.18}}
  >{\raggedright\arraybackslash}p{(\columnwidth - 4\tabcolsep) * \real{0.20}}
  >{\raggedright\arraybackslash}p{(\columnwidth - 4\tabcolsep) * \real{0.62}}@{}}
\toprule
名称 & Python 类型 & 含义 \\
\midrule
\endhead
\passthrough{\lstinline!value!} & \passthrough{\lstinline!int!} &
要写入或传递的新值，必须符合该参数的 Python 类型以及底层 C++
范围约束。 \\
\bottomrule
\end{longtable}

\textbf{返回值}

\begin{longtable}[]{@{}
  >{\raggedright\arraybackslash}p{(\columnwidth - 4\tabcolsep) * \real{0.18}}
  >{\raggedright\arraybackslash}p{(\columnwidth - 4\tabcolsep) * \real{0.20}}
  >{\raggedright\arraybackslash}p{(\columnwidth - 4\tabcolsep) * \real{0.62}}@{}}
\toprule
位置 & 类型 & 含义 \\
\midrule
\endhead
getter 返回值 & \passthrough{\lstinline!int!} & 返回属性当前值。 \\
\bottomrule
\end{longtable}

\textbf{对应 C++}

\begin{itemize}
\tightlist
\item
  类：\passthrough{\lstinline!unitree\_go::msg::dds\_::LidarState\_!}
\item
  字段类型：\passthrough{\lstinline!uint32\_t!}
\item
  声明位置：\passthrough{\lstinline!include/unitree/idl/go2/LidarState\_.hpp:41!}
\end{itemize}

\textbf{用法}

\begin{lstlisting}[language=Python]
current_value = obj.serial_buffer_size
obj.serial_buffer_size = new_value
\end{lstlisting}

\hypertarget{unitree-sdk2-cpp-idl-go2-lidarstate-serial-buffer-read}{%
\paragraph{\texorpdfstring{\texttt{unitree\_sdk2\_cpp.idl.go2.LidarState.serial\_buffer\_read}}{unitree\_sdk2\_cpp.idl.go2.LidarState.serial\_buffer\_read}}\label{unitree-sdk2-cpp-idl-go2-lidarstate-serial-buffer-read}}

底层 \passthrough{\lstinline!unitree\_go::msg::dds\_::LidarState\_!} 的
\passthrough{\lstinline!serial\_buffer\_read!} 字段。类型签名只说明
Python 表示；单位、数值范围、数组固定长度和枚举含义需查对应 IDL 头文件。
取值范围为 0 到 4294967295。

\textbf{签名}

\begin{lstlisting}[language=Python]
@property
def serial_buffer_read(self) -> int

@serial_buffer_read.setter
def serial_buffer_read(self, value: int) -> None
\end{lstlisting}

\textbf{可用性与安全性}

\textbf{\passthrough{\lstinline!AVAILABLE!} /
\passthrough{\lstinline!VALUE\_TYPE!}}：当前绑定源码已实现该消息值操作。

\textbf{参数}

\begin{longtable}[]{@{}
  >{\raggedright\arraybackslash}p{(\columnwidth - 4\tabcolsep) * \real{0.18}}
  >{\raggedright\arraybackslash}p{(\columnwidth - 4\tabcolsep) * \real{0.20}}
  >{\raggedright\arraybackslash}p{(\columnwidth - 4\tabcolsep) * \real{0.62}}@{}}
\toprule
名称 & Python 类型 & 含义 \\
\midrule
\endhead
\passthrough{\lstinline!value!} & \passthrough{\lstinline!int!} &
要写入或传递的新值，必须符合该参数的 Python 类型以及底层 C++
范围约束。 \\
\bottomrule
\end{longtable}

\textbf{返回值}

\begin{longtable}[]{@{}
  >{\raggedright\arraybackslash}p{(\columnwidth - 4\tabcolsep) * \real{0.18}}
  >{\raggedright\arraybackslash}p{(\columnwidth - 4\tabcolsep) * \real{0.20}}
  >{\raggedright\arraybackslash}p{(\columnwidth - 4\tabcolsep) * \real{0.62}}@{}}
\toprule
位置 & 类型 & 含义 \\
\midrule
\endhead
getter 返回值 & \passthrough{\lstinline!int!} & 返回属性当前值。 \\
\bottomrule
\end{longtable}

\textbf{对应 C++}

\begin{itemize}
\tightlist
\item
  类：\passthrough{\lstinline!unitree\_go::msg::dds\_::LidarState\_!}
\item
  字段类型：\passthrough{\lstinline!uint32\_t!}
\item
  声明位置：\passthrough{\lstinline!include/unitree/idl/go2/LidarState\_.hpp:42!}
\end{itemize}

\textbf{用法}

\begin{lstlisting}[language=Python]
current_value = obj.serial_buffer_read
obj.serial_buffer_read = new_value
\end{lstlisting}

\hypertarget{unitree-sdk2-cpp-idl-go2-motorcmd}{%
\subsubsection{\texorpdfstring{\texttt{unitree\_sdk2\_cpp.idl.go2.MotorCmd}}{unitree\_sdk2\_cpp.idl.go2.MotorCmd}}\label{unitree-sdk2-cpp-idl-go2-motorcmd}}

可由 Python 读写的 DDS/IDL 消息值类型。字段采用复制语义。

\textbf{导入}

\begin{lstlisting}[language=Python]
from unitree_sdk2_cpp.idl.go2 import MotorCmd
\end{lstlisting}

\textbf{方法索引}

\begin{longtable}[]{@{}
  >{\raggedright\arraybackslash}p{(\columnwidth - 6\tabcolsep) * \real{0.18}}
  >{\raggedright\arraybackslash}p{(\columnwidth - 6\tabcolsep) * \real{0.24}}
  >{\raggedright\arraybackslash}p{(\columnwidth - 6\tabcolsep) * \real{0.18}}
  >{\raggedright\arraybackslash}p{(\columnwidth - 6\tabcolsep) * \real{0.40}}@{}}
\toprule
方法 & Python 签名 & 状态 & 安全分类 \\
\midrule
\endhead
\protect\hyperlink{unitree-sdk2-cpp-idl-go2-motorcmd-dunder-init-1}{\passthrough{\lstinline!\_\_init\_\_!}}
& \passthrough{\lstinline!def \_\_init\_\_(self) -> None!} &
\passthrough{\lstinline!AVAILABLE!} &
\passthrough{\lstinline!VALUE\_TYPE!} \\
\protect\hyperlink{unitree-sdk2-cpp-idl-go2-motorcmd-dunder-eq-1}{\passthrough{\lstinline!\_\_eq\_\_!}}
& \passthrough{\lstinline!def \_\_eq\_\_(self, other: object) -> bool!}
& \passthrough{\lstinline!AVAILABLE!} &
\passthrough{\lstinline!VALUE\_TYPE!} \\
\protect\hyperlink{unitree-sdk2-cpp-idl-go2-motorcmd-dunder-ne-1}{\passthrough{\lstinline!\_\_ne\_\_!}}
& \passthrough{\lstinline!def \_\_ne\_\_(self, other: object) -> bool!}
& \passthrough{\lstinline!AVAILABLE!} &
\passthrough{\lstinline!VALUE\_TYPE!} \\
\bottomrule
\end{longtable}

\textbf{属性索引}

\begin{longtable}[]{@{}
  >{\raggedright\arraybackslash}p{(\columnwidth - 6\tabcolsep) * \real{0.18}}
  >{\raggedright\arraybackslash}p{(\columnwidth - 6\tabcolsep) * \real{0.24}}
  >{\raggedright\arraybackslash}p{(\columnwidth - 6\tabcolsep) * \real{0.18}}
  >{\raggedright\arraybackslash}p{(\columnwidth - 6\tabcolsep) * \real{0.40}}@{}}
\toprule
属性 & 读取类型 & 写入类型 & 状态 \\
\midrule
\endhead
\protect\hyperlink{unitree-sdk2-cpp-idl-go2-motorcmd-mode}{\passthrough{\lstinline!mode!}}
& \passthrough{\lstinline!int!} & \passthrough{\lstinline!int!} &
\passthrough{\lstinline!AVAILABLE!} \\
\protect\hyperlink{unitree-sdk2-cpp-idl-go2-motorcmd-q}{\passthrough{\lstinline!q!}}
& \passthrough{\lstinline!float!} & \passthrough{\lstinline!float!} &
\passthrough{\lstinline!AVAILABLE!} \\
\protect\hyperlink{unitree-sdk2-cpp-idl-go2-motorcmd-dq}{\passthrough{\lstinline!dq!}}
& \passthrough{\lstinline!float!} & \passthrough{\lstinline!float!} &
\passthrough{\lstinline!AVAILABLE!} \\
\protect\hyperlink{unitree-sdk2-cpp-idl-go2-motorcmd-tau}{\passthrough{\lstinline!tau!}}
& \passthrough{\lstinline!float!} & \passthrough{\lstinline!float!} &
\passthrough{\lstinline!AVAILABLE!} \\
\protect\hyperlink{unitree-sdk2-cpp-idl-go2-motorcmd-kp}{\passthrough{\lstinline!kp!}}
& \passthrough{\lstinline!float!} & \passthrough{\lstinline!float!} &
\passthrough{\lstinline!AVAILABLE!} \\
\protect\hyperlink{unitree-sdk2-cpp-idl-go2-motorcmd-kd}{\passthrough{\lstinline!kd!}}
& \passthrough{\lstinline!float!} & \passthrough{\lstinline!float!} &
\passthrough{\lstinline!AVAILABLE!} \\
\protect\hyperlink{unitree-sdk2-cpp-idl-go2-motorcmd-reserve}{\passthrough{\lstinline!reserve!}}
& \passthrough{\lstinline!list[int]!} &
\passthrough{\lstinline!Sequence[int]!} &
\passthrough{\lstinline!AVAILABLE!} \\
\bottomrule
\end{longtable}

\hypertarget{unitree-sdk2-cpp-idl-go2-motorcmd-dunder-init-1}{%
\paragraph{\texorpdfstring{\texttt{unitree\_sdk2\_cpp.idl.go2.MotorCmd.\_\_init\_\_}}{unitree\_sdk2\_cpp.idl.go2.MotorCmd.\_\_init\_\_}}\label{unitree-sdk2-cpp-idl-go2-motorcmd-dunder-init-1}}

初始化 \passthrough{\lstinline!MotorCmd!}
实例。是否能实际构造取决于下方可用性状态。

\textbf{签名}

\begin{lstlisting}[language=Python]
def __init__(self) -> None
\end{lstlisting}

\textbf{可用性与安全性}

\textbf{\passthrough{\lstinline!AVAILABLE!} /
\passthrough{\lstinline!VALUE\_TYPE!}}：当前绑定源码已实现该消息值操作。

\textbf{参数}

无显式参数。实例方法中的 \passthrough{\lstinline!self!} 由 Python
自动传入。

\textbf{返回值}

\begin{longtable}[]{@{}
  >{\raggedright\arraybackslash}p{(\columnwidth - 4\tabcolsep) * \real{0.18}}
  >{\raggedright\arraybackslash}p{(\columnwidth - 4\tabcolsep) * \real{0.20}}
  >{\raggedright\arraybackslash}p{(\columnwidth - 4\tabcolsep) * \real{0.62}}@{}}
\toprule
位置 & 类型 & 含义 \\
\midrule
\endhead
无 & \passthrough{\lstinline!None!} &
\passthrough{\lstinline!\_\_init\_\_!}
本身不返回值；调用类对象时，在构造函数可用的前提下得到该类实例。 \\
\bottomrule
\end{longtable}

\textbf{对应 C++}

\begin{itemize}
\tightlist
\item
  类：\passthrough{\lstinline!unitree\_go::msg::dds\_::MotorCmd\_!}
\item
  签名：\passthrough{\lstinline!MotorCmd\_()!}
\item
  绑定策略：\passthrough{\lstinline!未单独标注!}
\item
  声明位置：\passthrough{\lstinline!include/unitree/idl/go2/MotorCmd\_.hpp:33!}
\end{itemize}

\textbf{说明}

\begin{itemize}
\tightlist
\item
  示例是局部调用片段；\passthrough{\lstinline!obj!}
  和参数变量表示已经按上层业务校验并准备好的对象。
\end{itemize}

\textbf{用法}

\begin{lstlisting}[language=Python]
value = MotorCmd()
\end{lstlisting}

\hypertarget{unitree-sdk2-cpp-idl-go2-motorcmd-dunder-eq-1}{%
\paragraph{\texorpdfstring{\texttt{unitree\_sdk2\_cpp.idl.go2.MotorCmd.\_\_eq\_\_}}{unitree\_sdk2\_cpp.idl.go2.MotorCmd.\_\_eq\_\_}}\label{unitree-sdk2-cpp-idl-go2-motorcmd-dunder-eq-1}}

按底层 IDL 消息内容比较两个对象是否相等。

\textbf{签名}

\begin{lstlisting}[language=Python]
def __eq__(self, other: object) -> bool
\end{lstlisting}

\textbf{可用性与安全性}

\textbf{\passthrough{\lstinline!AVAILABLE!} /
\passthrough{\lstinline!VALUE\_TYPE!}}：当前绑定源码已实现该消息值操作。

\textbf{参数}

\begin{longtable}[]{@{}
  >{\raggedright\arraybackslash}p{(\columnwidth - 6\tabcolsep) * \real{0.18}}
  >{\raggedright\arraybackslash}p{(\columnwidth - 6\tabcolsep) * \real{0.24}}
  >{\raggedright\arraybackslash}p{(\columnwidth - 6\tabcolsep) * \real{0.18}}
  >{\raggedright\arraybackslash}p{(\columnwidth - 6\tabcolsep) * \real{0.40}}@{}}
\toprule
名称 & Python 类型 & 默认值 & 含义 \\
\midrule
\endhead
\passthrough{\lstinline!other!} & \passthrough{\lstinline!object!} &
必填 & 用于比较的另一个 Python 对象。类型不兼容时比较结果通常为
\passthrough{\lstinline!False!}。 \\
\bottomrule
\end{longtable}

\textbf{返回值}

\begin{longtable}[]{@{}
  >{\raggedright\arraybackslash}p{(\columnwidth - 4\tabcolsep) * \real{0.18}}
  >{\raggedright\arraybackslash}p{(\columnwidth - 4\tabcolsep) * \real{0.20}}
  >{\raggedright\arraybackslash}p{(\columnwidth - 4\tabcolsep) * \real{0.62}}@{}}
\toprule
位置 & 类型 & 含义 \\
\midrule
\endhead
返回值 & \passthrough{\lstinline!bool!} & 返回
\passthrough{\lstinline!bool!}。更精确的业务含义见该方法用途、C++
签名和目标型号协议。 \\
\bottomrule
\end{longtable}

\textbf{对应 C++}

\begin{itemize}
\tightlist
\item
  类：\passthrough{\lstinline!unitree\_go::msg::dds\_::MotorCmd\_!}
\item
  签名：\passthrough{\lstinline!operator==(const value \&) const!}
\item
  绑定策略：\passthrough{\lstinline!未单独标注!}
\item
  声明位置：由 pybind11 手工绑定或生成注册表提供；manifest
  未记录独立头文件行号。
\end{itemize}

\textbf{说明}

\begin{itemize}
\tightlist
\item
  示例是局部调用片段；\passthrough{\lstinline!obj!}
  和参数变量表示已经按上层业务校验并准备好的对象。
\end{itemize}

\textbf{用法}

\begin{lstlisting}[language=Python]
same = left == right
\end{lstlisting}

\hypertarget{unitree-sdk2-cpp-idl-go2-motorcmd-dunder-ne-1}{%
\paragraph{\texorpdfstring{\texttt{unitree\_sdk2\_cpp.idl.go2.MotorCmd.\_\_ne\_\_}}{unitree\_sdk2\_cpp.idl.go2.MotorCmd.\_\_ne\_\_}}\label{unitree-sdk2-cpp-idl-go2-motorcmd-dunder-ne-1}}

按底层 IDL 消息内容比较两个对象是否不相等。

\textbf{签名}

\begin{lstlisting}[language=Python]
def __ne__(self, other: object) -> bool
\end{lstlisting}

\textbf{可用性与安全性}

\textbf{\passthrough{\lstinline!AVAILABLE!} /
\passthrough{\lstinline!VALUE\_TYPE!}}：当前绑定源码已实现该消息值操作。

\textbf{参数}

\begin{longtable}[]{@{}
  >{\raggedright\arraybackslash}p{(\columnwidth - 6\tabcolsep) * \real{0.18}}
  >{\raggedright\arraybackslash}p{(\columnwidth - 6\tabcolsep) * \real{0.24}}
  >{\raggedright\arraybackslash}p{(\columnwidth - 6\tabcolsep) * \real{0.18}}
  >{\raggedright\arraybackslash}p{(\columnwidth - 6\tabcolsep) * \real{0.40}}@{}}
\toprule
名称 & Python 类型 & 默认值 & 含义 \\
\midrule
\endhead
\passthrough{\lstinline!other!} & \passthrough{\lstinline!object!} &
必填 & 用于比较的另一个 Python 对象。类型不兼容时比较结果通常为
\passthrough{\lstinline!False!}。 \\
\bottomrule
\end{longtable}

\textbf{返回值}

\begin{longtable}[]{@{}
  >{\raggedright\arraybackslash}p{(\columnwidth - 4\tabcolsep) * \real{0.18}}
  >{\raggedright\arraybackslash}p{(\columnwidth - 4\tabcolsep) * \real{0.20}}
  >{\raggedright\arraybackslash}p{(\columnwidth - 4\tabcolsep) * \real{0.62}}@{}}
\toprule
位置 & 类型 & 含义 \\
\midrule
\endhead
返回值 & \passthrough{\lstinline!bool!} & 返回
\passthrough{\lstinline!bool!}。更精确的业务含义见该方法用途、C++
签名和目标型号协议。 \\
\bottomrule
\end{longtable}

\textbf{对应 C++}

\begin{itemize}
\tightlist
\item
  类：\passthrough{\lstinline!unitree\_go::msg::dds\_::MotorCmd\_!}
\item
  签名：\passthrough{\lstinline"operator!=(const value \&) const"}
\item
  绑定策略：\passthrough{\lstinline!未单独标注!}
\item
  声明位置：由 pybind11 手工绑定或生成注册表提供；manifest
  未记录独立头文件行号。
\end{itemize}

\textbf{说明}

\begin{itemize}
\tightlist
\item
  示例是局部调用片段；\passthrough{\lstinline!obj!}
  和参数变量表示已经按上层业务校验并准备好的对象。
\end{itemize}

\textbf{用法}

\begin{lstlisting}[language=Python]
different = left != right
\end{lstlisting}

\hypertarget{unitree-sdk2-cpp-idl-go2-motorcmd-mode}{%
\paragraph{\texorpdfstring{\texttt{unitree\_sdk2\_cpp.idl.go2.MotorCmd.mode}}{unitree\_sdk2\_cpp.idl.go2.MotorCmd.mode}}\label{unitree-sdk2-cpp-idl-go2-motorcmd-mode}}

底层 \passthrough{\lstinline!unitree\_go::msg::dds\_::MotorCmd\_!} 的
\passthrough{\lstinline!mode!} 字段。类型签名只说明 Python
表示；单位、数值范围、数组固定长度和枚举含义需查对应 IDL 头文件。
取值范围为 0 到 255。

\textbf{签名}

\begin{lstlisting}[language=Python]
@property
def mode(self) -> int

@mode.setter
def mode(self, value: int) -> None
\end{lstlisting}

\textbf{可用性与安全性}

\textbf{\passthrough{\lstinline!AVAILABLE!} /
\passthrough{\lstinline!VALUE\_TYPE!}}：当前绑定源码已实现该消息值操作。

\textbf{参数}

\begin{longtable}[]{@{}
  >{\raggedright\arraybackslash}p{(\columnwidth - 4\tabcolsep) * \real{0.18}}
  >{\raggedright\arraybackslash}p{(\columnwidth - 4\tabcolsep) * \real{0.20}}
  >{\raggedright\arraybackslash}p{(\columnwidth - 4\tabcolsep) * \real{0.62}}@{}}
\toprule
名称 & Python 类型 & 含义 \\
\midrule
\endhead
\passthrough{\lstinline!value!} & \passthrough{\lstinline!int!} &
要写入或传递的新值，必须符合该参数的 Python 类型以及底层 C++
范围约束。 \\
\bottomrule
\end{longtable}

\textbf{返回值}

\begin{longtable}[]{@{}
  >{\raggedright\arraybackslash}p{(\columnwidth - 4\tabcolsep) * \real{0.18}}
  >{\raggedright\arraybackslash}p{(\columnwidth - 4\tabcolsep) * \real{0.20}}
  >{\raggedright\arraybackslash}p{(\columnwidth - 4\tabcolsep) * \real{0.62}}@{}}
\toprule
位置 & 类型 & 含义 \\
\midrule
\endhead
getter 返回值 & \passthrough{\lstinline!int!} & 返回属性当前值。 \\
\bottomrule
\end{longtable}

\textbf{对应 C++}

\begin{itemize}
\tightlist
\item
  类：\passthrough{\lstinline!unitree\_go::msg::dds\_::MotorCmd\_!}
\item
  字段类型：\passthrough{\lstinline!uint8\_t!}
\item
  声明位置：\passthrough{\lstinline!include/unitree/idl/go2/MotorCmd\_.hpp:24!}
\end{itemize}

\textbf{用法}

\begin{lstlisting}[language=Python]
current_value = obj.mode
obj.mode = new_value
\end{lstlisting}

\hypertarget{unitree-sdk2-cpp-idl-go2-motorcmd-q}{%
\paragraph{\texorpdfstring{\texttt{unitree\_sdk2\_cpp.idl.go2.MotorCmd.q}}{unitree\_sdk2\_cpp.idl.go2.MotorCmd.q}}\label{unitree-sdk2-cpp-idl-go2-motorcmd-q}}

底层 \passthrough{\lstinline!unitree\_go::msg::dds\_::MotorCmd\_!} 的
\passthrough{\lstinline!q!} 字段。类型签名只说明 Python
表示；单位、数值范围、数组固定长度和枚举含义需查对应 IDL 头文件。
底层为浮点数；类型本身不说明单位、坐标系或业务安全范围。

\textbf{签名}

\begin{lstlisting}[language=Python]
@property
def q(self) -> float

@q.setter
def q(self, value: float) -> None
\end{lstlisting}

\textbf{可用性与安全性}

\textbf{\passthrough{\lstinline!AVAILABLE!} /
\passthrough{\lstinline!VALUE\_TYPE!}}：当前绑定源码已实现该消息值操作。

\textbf{参数}

\begin{longtable}[]{@{}
  >{\raggedright\arraybackslash}p{(\columnwidth - 4\tabcolsep) * \real{0.18}}
  >{\raggedright\arraybackslash}p{(\columnwidth - 4\tabcolsep) * \real{0.20}}
  >{\raggedright\arraybackslash}p{(\columnwidth - 4\tabcolsep) * \real{0.62}}@{}}
\toprule
名称 & Python 类型 & 含义 \\
\midrule
\endhead
\passthrough{\lstinline!value!} & \passthrough{\lstinline!float!} &
要写入或传递的新值，必须符合该参数的 Python 类型以及底层 C++
范围约束。 \\
\bottomrule
\end{longtable}

\textbf{返回值}

\begin{longtable}[]{@{}
  >{\raggedright\arraybackslash}p{(\columnwidth - 4\tabcolsep) * \real{0.18}}
  >{\raggedright\arraybackslash}p{(\columnwidth - 4\tabcolsep) * \real{0.20}}
  >{\raggedright\arraybackslash}p{(\columnwidth - 4\tabcolsep) * \real{0.62}}@{}}
\toprule
位置 & 类型 & 含义 \\
\midrule
\endhead
getter 返回值 & \passthrough{\lstinline!float!} & 返回属性当前值。 \\
\bottomrule
\end{longtable}

\textbf{对应 C++}

\begin{itemize}
\tightlist
\item
  类：\passthrough{\lstinline!unitree\_go::msg::dds\_::MotorCmd\_!}
\item
  字段类型：\passthrough{\lstinline!float!}
\item
  声明位置：\passthrough{\lstinline!include/unitree/idl/go2/MotorCmd\_.hpp:25!}
\end{itemize}

\textbf{用法}

\begin{lstlisting}[language=Python]
current_value = obj.q
obj.q = new_value
\end{lstlisting}

\hypertarget{unitree-sdk2-cpp-idl-go2-motorcmd-dq}{%
\paragraph{\texorpdfstring{\texttt{unitree\_sdk2\_cpp.idl.go2.MotorCmd.dq}}{unitree\_sdk2\_cpp.idl.go2.MotorCmd.dq}}\label{unitree-sdk2-cpp-idl-go2-motorcmd-dq}}

底层 \passthrough{\lstinline!unitree\_go::msg::dds\_::MotorCmd\_!} 的
\passthrough{\lstinline!dq!} 字段。类型签名只说明 Python
表示；单位、数值范围、数组固定长度和枚举含义需查对应 IDL 头文件。
底层为浮点数；类型本身不说明单位、坐标系或业务安全范围。

\textbf{签名}

\begin{lstlisting}[language=Python]
@property
def dq(self) -> float

@dq.setter
def dq(self, value: float) -> None
\end{lstlisting}

\textbf{可用性与安全性}

\textbf{\passthrough{\lstinline!AVAILABLE!} /
\passthrough{\lstinline!VALUE\_TYPE!}}：当前绑定源码已实现该消息值操作。

\textbf{参数}

\begin{longtable}[]{@{}
  >{\raggedright\arraybackslash}p{(\columnwidth - 4\tabcolsep) * \real{0.18}}
  >{\raggedright\arraybackslash}p{(\columnwidth - 4\tabcolsep) * \real{0.20}}
  >{\raggedright\arraybackslash}p{(\columnwidth - 4\tabcolsep) * \real{0.62}}@{}}
\toprule
名称 & Python 类型 & 含义 \\
\midrule
\endhead
\passthrough{\lstinline!value!} & \passthrough{\lstinline!float!} &
要写入或传递的新值，必须符合该参数的 Python 类型以及底层 C++
范围约束。 \\
\bottomrule
\end{longtable}

\textbf{返回值}

\begin{longtable}[]{@{}
  >{\raggedright\arraybackslash}p{(\columnwidth - 4\tabcolsep) * \real{0.18}}
  >{\raggedright\arraybackslash}p{(\columnwidth - 4\tabcolsep) * \real{0.20}}
  >{\raggedright\arraybackslash}p{(\columnwidth - 4\tabcolsep) * \real{0.62}}@{}}
\toprule
位置 & 类型 & 含义 \\
\midrule
\endhead
getter 返回值 & \passthrough{\lstinline!float!} & 返回属性当前值。 \\
\bottomrule
\end{longtable}

\textbf{对应 C++}

\begin{itemize}
\tightlist
\item
  类：\passthrough{\lstinline!unitree\_go::msg::dds\_::MotorCmd\_!}
\item
  字段类型：\passthrough{\lstinline!float!}
\item
  声明位置：\passthrough{\lstinline!include/unitree/idl/go2/MotorCmd\_.hpp:26!}
\end{itemize}

\textbf{用法}

\begin{lstlisting}[language=Python]
current_value = obj.dq
obj.dq = new_value
\end{lstlisting}

\hypertarget{unitree-sdk2-cpp-idl-go2-motorcmd-tau}{%
\paragraph{\texorpdfstring{\texttt{unitree\_sdk2\_cpp.idl.go2.MotorCmd.tau}}{unitree\_sdk2\_cpp.idl.go2.MotorCmd.tau}}\label{unitree-sdk2-cpp-idl-go2-motorcmd-tau}}

底层 \passthrough{\lstinline!unitree\_go::msg::dds\_::MotorCmd\_!} 的
\passthrough{\lstinline!tau!} 字段。类型签名只说明 Python
表示；单位、数值范围、数组固定长度和枚举含义需查对应 IDL 头文件。
底层为浮点数；类型本身不说明单位、坐标系或业务安全范围。

\textbf{签名}

\begin{lstlisting}[language=Python]
@property
def tau(self) -> float

@tau.setter
def tau(self, value: float) -> None
\end{lstlisting}

\textbf{可用性与安全性}

\textbf{\passthrough{\lstinline!AVAILABLE!} /
\passthrough{\lstinline!VALUE\_TYPE!}}：当前绑定源码已实现该消息值操作。

\textbf{参数}

\begin{longtable}[]{@{}
  >{\raggedright\arraybackslash}p{(\columnwidth - 4\tabcolsep) * \real{0.18}}
  >{\raggedright\arraybackslash}p{(\columnwidth - 4\tabcolsep) * \real{0.20}}
  >{\raggedright\arraybackslash}p{(\columnwidth - 4\tabcolsep) * \real{0.62}}@{}}
\toprule
名称 & Python 类型 & 含义 \\
\midrule
\endhead
\passthrough{\lstinline!value!} & \passthrough{\lstinline!float!} &
要写入或传递的新值，必须符合该参数的 Python 类型以及底层 C++
范围约束。 \\
\bottomrule
\end{longtable}

\textbf{返回值}

\begin{longtable}[]{@{}
  >{\raggedright\arraybackslash}p{(\columnwidth - 4\tabcolsep) * \real{0.18}}
  >{\raggedright\arraybackslash}p{(\columnwidth - 4\tabcolsep) * \real{0.20}}
  >{\raggedright\arraybackslash}p{(\columnwidth - 4\tabcolsep) * \real{0.62}}@{}}
\toprule
位置 & 类型 & 含义 \\
\midrule
\endhead
getter 返回值 & \passthrough{\lstinline!float!} & 返回属性当前值。 \\
\bottomrule
\end{longtable}

\textbf{对应 C++}

\begin{itemize}
\tightlist
\item
  类：\passthrough{\lstinline!unitree\_go::msg::dds\_::MotorCmd\_!}
\item
  字段类型：\passthrough{\lstinline!float!}
\item
  声明位置：\passthrough{\lstinline!include/unitree/idl/go2/MotorCmd\_.hpp:27!}
\end{itemize}

\textbf{用法}

\begin{lstlisting}[language=Python]
current_value = obj.tau
obj.tau = new_value
\end{lstlisting}

\hypertarget{unitree-sdk2-cpp-idl-go2-motorcmd-kp}{%
\paragraph{\texorpdfstring{\texttt{unitree\_sdk2\_cpp.idl.go2.MotorCmd.kp}}{unitree\_sdk2\_cpp.idl.go2.MotorCmd.kp}}\label{unitree-sdk2-cpp-idl-go2-motorcmd-kp}}

底层 \passthrough{\lstinline!unitree\_go::msg::dds\_::MotorCmd\_!} 的
\passthrough{\lstinline!kp!} 字段。类型签名只说明 Python
表示；单位、数值范围、数组固定长度和枚举含义需查对应 IDL 头文件。
底层为浮点数；类型本身不说明单位、坐标系或业务安全范围。

\textbf{签名}

\begin{lstlisting}[language=Python]
@property
def kp(self) -> float

@kp.setter
def kp(self, value: float) -> None
\end{lstlisting}

\textbf{可用性与安全性}

\textbf{\passthrough{\lstinline!AVAILABLE!} /
\passthrough{\lstinline!VALUE\_TYPE!}}：当前绑定源码已实现该消息值操作。

\textbf{参数}

\begin{longtable}[]{@{}
  >{\raggedright\arraybackslash}p{(\columnwidth - 4\tabcolsep) * \real{0.18}}
  >{\raggedright\arraybackslash}p{(\columnwidth - 4\tabcolsep) * \real{0.20}}
  >{\raggedright\arraybackslash}p{(\columnwidth - 4\tabcolsep) * \real{0.62}}@{}}
\toprule
名称 & Python 类型 & 含义 \\
\midrule
\endhead
\passthrough{\lstinline!value!} & \passthrough{\lstinline!float!} &
要写入或传递的新值，必须符合该参数的 Python 类型以及底层 C++
范围约束。 \\
\bottomrule
\end{longtable}

\textbf{返回值}

\begin{longtable}[]{@{}
  >{\raggedright\arraybackslash}p{(\columnwidth - 4\tabcolsep) * \real{0.18}}
  >{\raggedright\arraybackslash}p{(\columnwidth - 4\tabcolsep) * \real{0.20}}
  >{\raggedright\arraybackslash}p{(\columnwidth - 4\tabcolsep) * \real{0.62}}@{}}
\toprule
位置 & 类型 & 含义 \\
\midrule
\endhead
getter 返回值 & \passthrough{\lstinline!float!} & 返回属性当前值。 \\
\bottomrule
\end{longtable}

\textbf{对应 C++}

\begin{itemize}
\tightlist
\item
  类：\passthrough{\lstinline!unitree\_go::msg::dds\_::MotorCmd\_!}
\item
  字段类型：\passthrough{\lstinline!float!}
\item
  声明位置：\passthrough{\lstinline!include/unitree/idl/go2/MotorCmd\_.hpp:28!}
\end{itemize}

\textbf{用法}

\begin{lstlisting}[language=Python]
current_value = obj.kp
obj.kp = new_value
\end{lstlisting}

\hypertarget{unitree-sdk2-cpp-idl-go2-motorcmd-kd}{%
\paragraph{\texorpdfstring{\texttt{unitree\_sdk2\_cpp.idl.go2.MotorCmd.kd}}{unitree\_sdk2\_cpp.idl.go2.MotorCmd.kd}}\label{unitree-sdk2-cpp-idl-go2-motorcmd-kd}}

底层 \passthrough{\lstinline!unitree\_go::msg::dds\_::MotorCmd\_!} 的
\passthrough{\lstinline!kd!} 字段。类型签名只说明 Python
表示；单位、数值范围、数组固定长度和枚举含义需查对应 IDL 头文件。
底层为浮点数；类型本身不说明单位、坐标系或业务安全范围。

\textbf{签名}

\begin{lstlisting}[language=Python]
@property
def kd(self) -> float

@kd.setter
def kd(self, value: float) -> None
\end{lstlisting}

\textbf{可用性与安全性}

\textbf{\passthrough{\lstinline!AVAILABLE!} /
\passthrough{\lstinline!VALUE\_TYPE!}}：当前绑定源码已实现该消息值操作。

\textbf{参数}

\begin{longtable}[]{@{}
  >{\raggedright\arraybackslash}p{(\columnwidth - 4\tabcolsep) * \real{0.18}}
  >{\raggedright\arraybackslash}p{(\columnwidth - 4\tabcolsep) * \real{0.20}}
  >{\raggedright\arraybackslash}p{(\columnwidth - 4\tabcolsep) * \real{0.62}}@{}}
\toprule
名称 & Python 类型 & 含义 \\
\midrule
\endhead
\passthrough{\lstinline!value!} & \passthrough{\lstinline!float!} &
要写入或传递的新值，必须符合该参数的 Python 类型以及底层 C++
范围约束。 \\
\bottomrule
\end{longtable}

\textbf{返回值}

\begin{longtable}[]{@{}
  >{\raggedright\arraybackslash}p{(\columnwidth - 4\tabcolsep) * \real{0.18}}
  >{\raggedright\arraybackslash}p{(\columnwidth - 4\tabcolsep) * \real{0.20}}
  >{\raggedright\arraybackslash}p{(\columnwidth - 4\tabcolsep) * \real{0.62}}@{}}
\toprule
位置 & 类型 & 含义 \\
\midrule
\endhead
getter 返回值 & \passthrough{\lstinline!float!} & 返回属性当前值。 \\
\bottomrule
\end{longtable}

\textbf{对应 C++}

\begin{itemize}
\tightlist
\item
  类：\passthrough{\lstinline!unitree\_go::msg::dds\_::MotorCmd\_!}
\item
  字段类型：\passthrough{\lstinline!float!}
\item
  声明位置：\passthrough{\lstinline!include/unitree/idl/go2/MotorCmd\_.hpp:29!}
\end{itemize}

\textbf{用法}

\begin{lstlisting}[language=Python]
current_value = obj.kd
obj.kd = new_value
\end{lstlisting}

\hypertarget{unitree-sdk2-cpp-idl-go2-motorcmd-reserve}{%
\paragraph{\texorpdfstring{\texttt{unitree\_sdk2\_cpp.idl.go2.MotorCmd.reserve}}{unitree\_sdk2\_cpp.idl.go2.MotorCmd.reserve}}\label{unitree-sdk2-cpp-idl-go2-motorcmd-reserve}}

底层 \passthrough{\lstinline!unitree\_go::msg::dds\_::MotorCmd\_!} 的
\passthrough{\lstinline!reserve!} 字段。类型签名只说明 Python
表示；单位、数值范围、数组固定长度和枚举含义需查对应 IDL 头文件。
底层是固定长度数组，写入序列必须正好包含 3 个元素；元素约束：取值范围为
0 到 4294967295。

\textbf{签名}

\begin{lstlisting}[language=Python]
@property
def reserve(self) -> list[int]

@reserve.setter
def reserve(self, value: Sequence[int]) -> None
\end{lstlisting}

\textbf{可用性与安全性}

\textbf{\passthrough{\lstinline!AVAILABLE!} /
\passthrough{\lstinline!VALUE\_TYPE!}}：当前绑定源码已实现该消息值操作。

\textbf{参数}

\begin{longtable}[]{@{}
  >{\raggedright\arraybackslash}p{(\columnwidth - 4\tabcolsep) * \real{0.18}}
  >{\raggedright\arraybackslash}p{(\columnwidth - 4\tabcolsep) * \real{0.20}}
  >{\raggedright\arraybackslash}p{(\columnwidth - 4\tabcolsep) * \real{0.62}}@{}}
\toprule
名称 & Python 类型 & 含义 \\
\midrule
\endhead
\passthrough{\lstinline!value!} &
\passthrough{\lstinline!Sequence[int]!} &
要写入或传递的新值，必须符合该参数的 Python 类型以及底层 C++
范围约束。 \\
\bottomrule
\end{longtable}

\textbf{返回值}

\begin{longtable}[]{@{}
  >{\raggedright\arraybackslash}p{(\columnwidth - 4\tabcolsep) * \real{0.18}}
  >{\raggedright\arraybackslash}p{(\columnwidth - 4\tabcolsep) * \real{0.20}}
  >{\raggedright\arraybackslash}p{(\columnwidth - 4\tabcolsep) * \real{0.62}}@{}}
\toprule
位置 & 类型 & 含义 \\
\midrule
\endhead
getter 返回值 & \passthrough{\lstinline!list[int]!} &
返回属性当前值。数组、vector 和嵌套消息采用复制语义。 \\
\bottomrule
\end{longtable}

\textbf{对应 C++}

\begin{itemize}
\tightlist
\item
  类：\passthrough{\lstinline!unitree\_go::msg::dds\_::MotorCmd\_!}
\item
  字段类型：\passthrough{\lstinline!std::array<uint32\_t, 3>!}
\item
  声明位置：\passthrough{\lstinline!include/unitree/idl/go2/MotorCmd\_.hpp:30!}
\end{itemize}

\textbf{用法}

\begin{lstlisting}[language=Python]
current_value = obj.reserve
obj.reserve = new_value
\end{lstlisting}

\begin{quote}
{[}!TIP{]} 读取容器或嵌套值后应遵循``读取、修改、重新赋值''。只修改
getter 返回的副本不会可靠地更新原 C++ 消息。
\end{quote}

\hypertarget{unitree-sdk2-cpp-idl-go2-lowcmd}{%
\subsubsection{\texorpdfstring{\texttt{unitree\_sdk2\_cpp.idl.go2.LowCmd}}{unitree\_sdk2\_cpp.idl.go2.LowCmd}}\label{unitree-sdk2-cpp-idl-go2-lowcmd}}

可由 Python 读写的 DDS/IDL 消息值类型。字段采用复制语义。

\textbf{导入}

\begin{lstlisting}[language=Python]
from unitree_sdk2_cpp.idl.go2 import LowCmd
\end{lstlisting}

\textbf{方法索引}

\begin{longtable}[]{@{}
  >{\raggedright\arraybackslash}p{(\columnwidth - 6\tabcolsep) * \real{0.18}}
  >{\raggedright\arraybackslash}p{(\columnwidth - 6\tabcolsep) * \real{0.24}}
  >{\raggedright\arraybackslash}p{(\columnwidth - 6\tabcolsep) * \real{0.18}}
  >{\raggedright\arraybackslash}p{(\columnwidth - 6\tabcolsep) * \real{0.40}}@{}}
\toprule
方法 & Python 签名 & 状态 & 安全分类 \\
\midrule
\endhead
\protect\hyperlink{unitree-sdk2-cpp-idl-go2-lowcmd-dunder-init-1}{\passthrough{\lstinline!\_\_init\_\_!}}
& \passthrough{\lstinline!def \_\_init\_\_(self) -> None!} &
\passthrough{\lstinline!AVAILABLE!} &
\passthrough{\lstinline!VALUE\_TYPE!} \\
\protect\hyperlink{unitree-sdk2-cpp-idl-go2-lowcmd-dunder-eq-1}{\passthrough{\lstinline!\_\_eq\_\_!}}
& \passthrough{\lstinline!def \_\_eq\_\_(self, other: object) -> bool!}
& \passthrough{\lstinline!AVAILABLE!} &
\passthrough{\lstinline!VALUE\_TYPE!} \\
\protect\hyperlink{unitree-sdk2-cpp-idl-go2-lowcmd-dunder-ne-1}{\passthrough{\lstinline!\_\_ne\_\_!}}
& \passthrough{\lstinline!def \_\_ne\_\_(self, other: object) -> bool!}
& \passthrough{\lstinline!AVAILABLE!} &
\passthrough{\lstinline!VALUE\_TYPE!} \\
\bottomrule
\end{longtable}

\textbf{属性索引}

\begin{longtable}[]{@{}
  >{\raggedright\arraybackslash}p{(\columnwidth - 6\tabcolsep) * \real{0.18}}
  >{\raggedright\arraybackslash}p{(\columnwidth - 6\tabcolsep) * \real{0.24}}
  >{\raggedright\arraybackslash}p{(\columnwidth - 6\tabcolsep) * \real{0.18}}
  >{\raggedright\arraybackslash}p{(\columnwidth - 6\tabcolsep) * \real{0.40}}@{}}
\toprule
属性 & 读取类型 & 写入类型 & 状态 \\
\midrule
\endhead
\protect\hyperlink{unitree-sdk2-cpp-idl-go2-lowcmd-head}{\passthrough{\lstinline!head!}}
& \passthrough{\lstinline!list[int]!} &
\passthrough{\lstinline!Sequence[int]!} &
\passthrough{\lstinline!AVAILABLE!} \\
\protect\hyperlink{unitree-sdk2-cpp-idl-go2-lowcmd-level-flag}{\passthrough{\lstinline!level\_flag!}}
& \passthrough{\lstinline!int!} & \passthrough{\lstinline!int!} &
\passthrough{\lstinline!AVAILABLE!} \\
\protect\hyperlink{unitree-sdk2-cpp-idl-go2-lowcmd-frame-reserve}{\passthrough{\lstinline!frame\_reserve!}}
& \passthrough{\lstinline!int!} & \passthrough{\lstinline!int!} &
\passthrough{\lstinline!AVAILABLE!} \\
\protect\hyperlink{unitree-sdk2-cpp-idl-go2-lowcmd-sn}{\passthrough{\lstinline!sn!}}
& \passthrough{\lstinline!list[int]!} &
\passthrough{\lstinline!Sequence[int]!} &
\passthrough{\lstinline!AVAILABLE!} \\
\protect\hyperlink{unitree-sdk2-cpp-idl-go2-lowcmd-version}{\passthrough{\lstinline!version!}}
& \passthrough{\lstinline!list[int]!} &
\passthrough{\lstinline!Sequence[int]!} &
\passthrough{\lstinline!AVAILABLE!} \\
\protect\hyperlink{unitree-sdk2-cpp-idl-go2-lowcmd-bandwidth}{\passthrough{\lstinline!bandwidth!}}
& \passthrough{\lstinline!int!} & \passthrough{\lstinline!int!} &
\passthrough{\lstinline!AVAILABLE!} \\
\protect\hyperlink{unitree-sdk2-cpp-idl-go2-lowcmd-motor-cmd}{\passthrough{\lstinline!motor\_cmd!}}
& \passthrough{\lstinline!list[MotorCmd]!} &
\passthrough{\lstinline!Sequence[MotorCmd]!} &
\passthrough{\lstinline!AVAILABLE!} \\
\protect\hyperlink{unitree-sdk2-cpp-idl-go2-lowcmd-bms-cmd}{\passthrough{\lstinline!bms\_cmd!}}
& \passthrough{\lstinline!BmsCmd!} & \passthrough{\lstinline!BmsCmd!} &
\passthrough{\lstinline!AVAILABLE!} \\
\protect\hyperlink{unitree-sdk2-cpp-idl-go2-lowcmd-wireless-remote}{\passthrough{\lstinline!wireless\_remote!}}
& \passthrough{\lstinline!list[int]!} &
\passthrough{\lstinline!Sequence[int]!} &
\passthrough{\lstinline!AVAILABLE!} \\
\protect\hyperlink{unitree-sdk2-cpp-idl-go2-lowcmd-led}{\passthrough{\lstinline!led!}}
& \passthrough{\lstinline!list[int]!} &
\passthrough{\lstinline!Sequence[int]!} &
\passthrough{\lstinline!AVAILABLE!} \\
\protect\hyperlink{unitree-sdk2-cpp-idl-go2-lowcmd-fan}{\passthrough{\lstinline!fan!}}
& \passthrough{\lstinline!list[int]!} &
\passthrough{\lstinline!Sequence[int]!} &
\passthrough{\lstinline!AVAILABLE!} \\
\protect\hyperlink{unitree-sdk2-cpp-idl-go2-lowcmd-gpio}{\passthrough{\lstinline!gpio!}}
& \passthrough{\lstinline!int!} & \passthrough{\lstinline!int!} &
\passthrough{\lstinline!AVAILABLE!} \\
\protect\hyperlink{unitree-sdk2-cpp-idl-go2-lowcmd-reserve}{\passthrough{\lstinline!reserve!}}
& \passthrough{\lstinline!int!} & \passthrough{\lstinline!int!} &
\passthrough{\lstinline!AVAILABLE!} \\
\protect\hyperlink{unitree-sdk2-cpp-idl-go2-lowcmd-crc}{\passthrough{\lstinline!crc!}}
& \passthrough{\lstinline!int!} & \passthrough{\lstinline!int!} &
\passthrough{\lstinline!AVAILABLE!} \\
\bottomrule
\end{longtable}

\hypertarget{unitree-sdk2-cpp-idl-go2-lowcmd-dunder-init-1}{%
\paragraph{\texorpdfstring{\texttt{unitree\_sdk2\_cpp.idl.go2.LowCmd.\_\_init\_\_}}{unitree\_sdk2\_cpp.idl.go2.LowCmd.\_\_init\_\_}}\label{unitree-sdk2-cpp-idl-go2-lowcmd-dunder-init-1}}

初始化 \passthrough{\lstinline!LowCmd!}
实例。是否能实际构造取决于下方可用性状态。

\textbf{签名}

\begin{lstlisting}[language=Python]
def __init__(self) -> None
\end{lstlisting}

\textbf{可用性与安全性}

\textbf{\passthrough{\lstinline!AVAILABLE!} /
\passthrough{\lstinline!VALUE\_TYPE!}}：当前绑定源码已实现该消息值操作。

\textbf{参数}

无显式参数。实例方法中的 \passthrough{\lstinline!self!} 由 Python
自动传入。

\textbf{返回值}

\begin{longtable}[]{@{}
  >{\raggedright\arraybackslash}p{(\columnwidth - 4\tabcolsep) * \real{0.18}}
  >{\raggedright\arraybackslash}p{(\columnwidth - 4\tabcolsep) * \real{0.20}}
  >{\raggedright\arraybackslash}p{(\columnwidth - 4\tabcolsep) * \real{0.62}}@{}}
\toprule
位置 & 类型 & 含义 \\
\midrule
\endhead
无 & \passthrough{\lstinline!None!} &
\passthrough{\lstinline!\_\_init\_\_!}
本身不返回值；调用类对象时，在构造函数可用的前提下得到该类实例。 \\
\bottomrule
\end{longtable}

\textbf{对应 C++}

\begin{itemize}
\tightlist
\item
  类：\passthrough{\lstinline!unitree\_go::msg::dds\_::LowCmd\_!}
\item
  签名：\passthrough{\lstinline!LowCmd\_()!}
\item
  绑定策略：\passthrough{\lstinline!未单独标注!}
\item
  声明位置：\passthrough{\lstinline!include/unitree/idl/go2/LowCmd\_.hpp:44!}
\end{itemize}

\textbf{说明}

\begin{itemize}
\tightlist
\item
  示例是局部调用片段；\passthrough{\lstinline!obj!}
  和参数变量表示已经按上层业务校验并准备好的对象。
\end{itemize}

\textbf{用法}

\begin{lstlisting}[language=Python]
value = LowCmd()
\end{lstlisting}

\hypertarget{unitree-sdk2-cpp-idl-go2-lowcmd-dunder-eq-1}{%
\paragraph{\texorpdfstring{\texttt{unitree\_sdk2\_cpp.idl.go2.LowCmd.\_\_eq\_\_}}{unitree\_sdk2\_cpp.idl.go2.LowCmd.\_\_eq\_\_}}\label{unitree-sdk2-cpp-idl-go2-lowcmd-dunder-eq-1}}

按底层 IDL 消息内容比较两个对象是否相等。

\textbf{签名}

\begin{lstlisting}[language=Python]
def __eq__(self, other: object) -> bool
\end{lstlisting}

\textbf{可用性与安全性}

\textbf{\passthrough{\lstinline!AVAILABLE!} /
\passthrough{\lstinline!VALUE\_TYPE!}}：当前绑定源码已实现该消息值操作。

\textbf{参数}

\begin{longtable}[]{@{}
  >{\raggedright\arraybackslash}p{(\columnwidth - 6\tabcolsep) * \real{0.18}}
  >{\raggedright\arraybackslash}p{(\columnwidth - 6\tabcolsep) * \real{0.24}}
  >{\raggedright\arraybackslash}p{(\columnwidth - 6\tabcolsep) * \real{0.18}}
  >{\raggedright\arraybackslash}p{(\columnwidth - 6\tabcolsep) * \real{0.40}}@{}}
\toprule
名称 & Python 类型 & 默认值 & 含义 \\
\midrule
\endhead
\passthrough{\lstinline!other!} & \passthrough{\lstinline!object!} &
必填 & 用于比较的另一个 Python 对象。类型不兼容时比较结果通常为
\passthrough{\lstinline!False!}。 \\
\bottomrule
\end{longtable}

\textbf{返回值}

\begin{longtable}[]{@{}
  >{\raggedright\arraybackslash}p{(\columnwidth - 4\tabcolsep) * \real{0.18}}
  >{\raggedright\arraybackslash}p{(\columnwidth - 4\tabcolsep) * \real{0.20}}
  >{\raggedright\arraybackslash}p{(\columnwidth - 4\tabcolsep) * \real{0.62}}@{}}
\toprule
位置 & 类型 & 含义 \\
\midrule
\endhead
返回值 & \passthrough{\lstinline!bool!} & 返回
\passthrough{\lstinline!bool!}。更精确的业务含义见该方法用途、C++
签名和目标型号协议。 \\
\bottomrule
\end{longtable}

\textbf{对应 C++}

\begin{itemize}
\tightlist
\item
  类：\passthrough{\lstinline!unitree\_go::msg::dds\_::LowCmd\_!}
\item
  签名：\passthrough{\lstinline!operator==(const value \&) const!}
\item
  绑定策略：\passthrough{\lstinline!未单独标注!}
\item
  声明位置：由 pybind11 手工绑定或生成注册表提供；manifest
  未记录独立头文件行号。
\end{itemize}

\textbf{说明}

\begin{itemize}
\tightlist
\item
  示例是局部调用片段；\passthrough{\lstinline!obj!}
  和参数变量表示已经按上层业务校验并准备好的对象。
\end{itemize}

\textbf{用法}

\begin{lstlisting}[language=Python]
same = left == right
\end{lstlisting}

\hypertarget{unitree-sdk2-cpp-idl-go2-lowcmd-dunder-ne-1}{%
\paragraph{\texorpdfstring{\texttt{unitree\_sdk2\_cpp.idl.go2.LowCmd.\_\_ne\_\_}}{unitree\_sdk2\_cpp.idl.go2.LowCmd.\_\_ne\_\_}}\label{unitree-sdk2-cpp-idl-go2-lowcmd-dunder-ne-1}}

按底层 IDL 消息内容比较两个对象是否不相等。

\textbf{签名}

\begin{lstlisting}[language=Python]
def __ne__(self, other: object) -> bool
\end{lstlisting}

\textbf{可用性与安全性}

\textbf{\passthrough{\lstinline!AVAILABLE!} /
\passthrough{\lstinline!VALUE\_TYPE!}}：当前绑定源码已实现该消息值操作。

\textbf{参数}

\begin{longtable}[]{@{}
  >{\raggedright\arraybackslash}p{(\columnwidth - 6\tabcolsep) * \real{0.18}}
  >{\raggedright\arraybackslash}p{(\columnwidth - 6\tabcolsep) * \real{0.24}}
  >{\raggedright\arraybackslash}p{(\columnwidth - 6\tabcolsep) * \real{0.18}}
  >{\raggedright\arraybackslash}p{(\columnwidth - 6\tabcolsep) * \real{0.40}}@{}}
\toprule
名称 & Python 类型 & 默认值 & 含义 \\
\midrule
\endhead
\passthrough{\lstinline!other!} & \passthrough{\lstinline!object!} &
必填 & 用于比较的另一个 Python 对象。类型不兼容时比较结果通常为
\passthrough{\lstinline!False!}。 \\
\bottomrule
\end{longtable}

\textbf{返回值}

\begin{longtable}[]{@{}
  >{\raggedright\arraybackslash}p{(\columnwidth - 4\tabcolsep) * \real{0.18}}
  >{\raggedright\arraybackslash}p{(\columnwidth - 4\tabcolsep) * \real{0.20}}
  >{\raggedright\arraybackslash}p{(\columnwidth - 4\tabcolsep) * \real{0.62}}@{}}
\toprule
位置 & 类型 & 含义 \\
\midrule
\endhead
返回值 & \passthrough{\lstinline!bool!} & 返回
\passthrough{\lstinline!bool!}。更精确的业务含义见该方法用途、C++
签名和目标型号协议。 \\
\bottomrule
\end{longtable}

\textbf{对应 C++}

\begin{itemize}
\tightlist
\item
  类：\passthrough{\lstinline!unitree\_go::msg::dds\_::LowCmd\_!}
\item
  签名：\passthrough{\lstinline"operator!=(const value \&) const"}
\item
  绑定策略：\passthrough{\lstinline!未单独标注!}
\item
  声明位置：由 pybind11 手工绑定或生成注册表提供；manifest
  未记录独立头文件行号。
\end{itemize}

\textbf{说明}

\begin{itemize}
\tightlist
\item
  示例是局部调用片段；\passthrough{\lstinline!obj!}
  和参数变量表示已经按上层业务校验并准备好的对象。
\end{itemize}

\textbf{用法}

\begin{lstlisting}[language=Python]
different = left != right
\end{lstlisting}

\hypertarget{unitree-sdk2-cpp-idl-go2-lowcmd-head}{%
\paragraph{\texorpdfstring{\texttt{unitree\_sdk2\_cpp.idl.go2.LowCmd.head}}{unitree\_sdk2\_cpp.idl.go2.LowCmd.head}}\label{unitree-sdk2-cpp-idl-go2-lowcmd-head}}

底层 \passthrough{\lstinline!unitree\_go::msg::dds\_::LowCmd\_!} 的
\passthrough{\lstinline!head!} 字段。类型签名只说明 Python
表示；单位、数值范围、数组固定长度和枚举含义需查对应 IDL 头文件。
底层是固定长度数组，写入序列必须正好包含 2 个元素；元素约束：取值范围为
0 到 255。

\textbf{签名}

\begin{lstlisting}[language=Python]
@property
def head(self) -> list[int]

@head.setter
def head(self, value: Sequence[int]) -> None
\end{lstlisting}

\textbf{可用性与安全性}

\textbf{\passthrough{\lstinline!AVAILABLE!} /
\passthrough{\lstinline!VALUE\_TYPE!}}：当前绑定源码已实现该消息值操作。

\textbf{参数}

\begin{longtable}[]{@{}
  >{\raggedright\arraybackslash}p{(\columnwidth - 4\tabcolsep) * \real{0.18}}
  >{\raggedright\arraybackslash}p{(\columnwidth - 4\tabcolsep) * \real{0.20}}
  >{\raggedright\arraybackslash}p{(\columnwidth - 4\tabcolsep) * \real{0.62}}@{}}
\toprule
名称 & Python 类型 & 含义 \\
\midrule
\endhead
\passthrough{\lstinline!value!} &
\passthrough{\lstinline!Sequence[int]!} &
要写入或传递的新值，必须符合该参数的 Python 类型以及底层 C++
范围约束。 \\
\bottomrule
\end{longtable}

\textbf{返回值}

\begin{longtable}[]{@{}
  >{\raggedright\arraybackslash}p{(\columnwidth - 4\tabcolsep) * \real{0.18}}
  >{\raggedright\arraybackslash}p{(\columnwidth - 4\tabcolsep) * \real{0.20}}
  >{\raggedright\arraybackslash}p{(\columnwidth - 4\tabcolsep) * \real{0.62}}@{}}
\toprule
位置 & 类型 & 含义 \\
\midrule
\endhead
getter 返回值 & \passthrough{\lstinline!list[int]!} &
返回属性当前值。数组、vector 和嵌套消息采用复制语义。 \\
\bottomrule
\end{longtable}

\textbf{对应 C++}

\begin{itemize}
\tightlist
\item
  类：\passthrough{\lstinline!unitree\_go::msg::dds\_::LowCmd\_!}
\item
  字段类型：\passthrough{\lstinline!std::array<uint8\_t, 2>!}
\item
  声明位置：\passthrough{\lstinline!include/unitree/idl/go2/LowCmd\_.hpp:28!}
\end{itemize}

\textbf{用法}

\begin{lstlisting}[language=Python]
current_value = obj.head
obj.head = new_value
\end{lstlisting}

\begin{quote}
{[}!TIP{]} 读取容器或嵌套值后应遵循``读取、修改、重新赋值''。只修改
getter 返回的副本不会可靠地更新原 C++ 消息。
\end{quote}

\hypertarget{unitree-sdk2-cpp-idl-go2-lowcmd-level-flag}{%
\paragraph{\texorpdfstring{\texttt{unitree\_sdk2\_cpp.idl.go2.LowCmd.level\_flag}}{unitree\_sdk2\_cpp.idl.go2.LowCmd.level\_flag}}\label{unitree-sdk2-cpp-idl-go2-lowcmd-level-flag}}

底层 \passthrough{\lstinline!unitree\_go::msg::dds\_::LowCmd\_!} 的
\passthrough{\lstinline!level\_flag!} 字段。类型签名只说明 Python
表示；单位、数值范围、数组固定长度和枚举含义需查对应 IDL 头文件。
取值范围为 0 到 255。

\textbf{签名}

\begin{lstlisting}[language=Python]
@property
def level_flag(self) -> int

@level_flag.setter
def level_flag(self, value: int) -> None
\end{lstlisting}

\textbf{可用性与安全性}

\textbf{\passthrough{\lstinline!AVAILABLE!} /
\passthrough{\lstinline!VALUE\_TYPE!}}：当前绑定源码已实现该消息值操作。

\textbf{参数}

\begin{longtable}[]{@{}
  >{\raggedright\arraybackslash}p{(\columnwidth - 4\tabcolsep) * \real{0.18}}
  >{\raggedright\arraybackslash}p{(\columnwidth - 4\tabcolsep) * \real{0.20}}
  >{\raggedright\arraybackslash}p{(\columnwidth - 4\tabcolsep) * \real{0.62}}@{}}
\toprule
名称 & Python 类型 & 含义 \\
\midrule
\endhead
\passthrough{\lstinline!value!} & \passthrough{\lstinline!int!} &
要写入或传递的新值，必须符合该参数的 Python 类型以及底层 C++
范围约束。 \\
\bottomrule
\end{longtable}

\textbf{返回值}

\begin{longtable}[]{@{}
  >{\raggedright\arraybackslash}p{(\columnwidth - 4\tabcolsep) * \real{0.18}}
  >{\raggedright\arraybackslash}p{(\columnwidth - 4\tabcolsep) * \real{0.20}}
  >{\raggedright\arraybackslash}p{(\columnwidth - 4\tabcolsep) * \real{0.62}}@{}}
\toprule
位置 & 类型 & 含义 \\
\midrule
\endhead
getter 返回值 & \passthrough{\lstinline!int!} & 返回属性当前值。 \\
\bottomrule
\end{longtable}

\textbf{对应 C++}

\begin{itemize}
\tightlist
\item
  类：\passthrough{\lstinline!unitree\_go::msg::dds\_::LowCmd\_!}
\item
  字段类型：\passthrough{\lstinline!uint8\_t!}
\item
  声明位置：\passthrough{\lstinline!include/unitree/idl/go2/LowCmd\_.hpp:29!}
\end{itemize}

\textbf{用法}

\begin{lstlisting}[language=Python]
current_value = obj.level_flag
obj.level_flag = new_value
\end{lstlisting}

\hypertarget{unitree-sdk2-cpp-idl-go2-lowcmd-frame-reserve}{%
\paragraph{\texorpdfstring{\texttt{unitree\_sdk2\_cpp.idl.go2.LowCmd.frame\_reserve}}{unitree\_sdk2\_cpp.idl.go2.LowCmd.frame\_reserve}}\label{unitree-sdk2-cpp-idl-go2-lowcmd-frame-reserve}}

底层 \passthrough{\lstinline!unitree\_go::msg::dds\_::LowCmd\_!} 的
\passthrough{\lstinline!frame\_reserve!} 字段。类型签名只说明 Python
表示；单位、数值范围、数组固定长度和枚举含义需查对应 IDL 头文件。
取值范围为 0 到 255。

\textbf{签名}

\begin{lstlisting}[language=Python]
@property
def frame_reserve(self) -> int

@frame_reserve.setter
def frame_reserve(self, value: int) -> None
\end{lstlisting}

\textbf{可用性与安全性}

\textbf{\passthrough{\lstinline!AVAILABLE!} /
\passthrough{\lstinline!VALUE\_TYPE!}}：当前绑定源码已实现该消息值操作。

\textbf{参数}

\begin{longtable}[]{@{}
  >{\raggedright\arraybackslash}p{(\columnwidth - 4\tabcolsep) * \real{0.18}}
  >{\raggedright\arraybackslash}p{(\columnwidth - 4\tabcolsep) * \real{0.20}}
  >{\raggedright\arraybackslash}p{(\columnwidth - 4\tabcolsep) * \real{0.62}}@{}}
\toprule
名称 & Python 类型 & 含义 \\
\midrule
\endhead
\passthrough{\lstinline!value!} & \passthrough{\lstinline!int!} &
要写入或传递的新值，必须符合该参数的 Python 类型以及底层 C++
范围约束。 \\
\bottomrule
\end{longtable}

\textbf{返回值}

\begin{longtable}[]{@{}
  >{\raggedright\arraybackslash}p{(\columnwidth - 4\tabcolsep) * \real{0.18}}
  >{\raggedright\arraybackslash}p{(\columnwidth - 4\tabcolsep) * \real{0.20}}
  >{\raggedright\arraybackslash}p{(\columnwidth - 4\tabcolsep) * \real{0.62}}@{}}
\toprule
位置 & 类型 & 含义 \\
\midrule
\endhead
getter 返回值 & \passthrough{\lstinline!int!} & 返回属性当前值。 \\
\bottomrule
\end{longtable}

\textbf{对应 C++}

\begin{itemize}
\tightlist
\item
  类：\passthrough{\lstinline!unitree\_go::msg::dds\_::LowCmd\_!}
\item
  字段类型：\passthrough{\lstinline!uint8\_t!}
\item
  声明位置：\passthrough{\lstinline!include/unitree/idl/go2/LowCmd\_.hpp:30!}
\end{itemize}

\textbf{用法}

\begin{lstlisting}[language=Python]
current_value = obj.frame_reserve
obj.frame_reserve = new_value
\end{lstlisting}

\hypertarget{unitree-sdk2-cpp-idl-go2-lowcmd-sn}{%
\paragraph{\texorpdfstring{\texttt{unitree\_sdk2\_cpp.idl.go2.LowCmd.sn}}{unitree\_sdk2\_cpp.idl.go2.LowCmd.sn}}\label{unitree-sdk2-cpp-idl-go2-lowcmd-sn}}

底层 \passthrough{\lstinline!unitree\_go::msg::dds\_::LowCmd\_!} 的
\passthrough{\lstinline!sn!} 字段。类型签名只说明 Python
表示；单位、数值范围、数组固定长度和枚举含义需查对应 IDL 头文件。
底层是固定长度数组，写入序列必须正好包含 2 个元素；元素约束：取值范围为
0 到 4294967295。

\textbf{签名}

\begin{lstlisting}[language=Python]
@property
def sn(self) -> list[int]

@sn.setter
def sn(self, value: Sequence[int]) -> None
\end{lstlisting}

\textbf{可用性与安全性}

\textbf{\passthrough{\lstinline!AVAILABLE!} /
\passthrough{\lstinline!VALUE\_TYPE!}}：当前绑定源码已实现该消息值操作。

\textbf{参数}

\begin{longtable}[]{@{}
  >{\raggedright\arraybackslash}p{(\columnwidth - 4\tabcolsep) * \real{0.18}}
  >{\raggedright\arraybackslash}p{(\columnwidth - 4\tabcolsep) * \real{0.20}}
  >{\raggedright\arraybackslash}p{(\columnwidth - 4\tabcolsep) * \real{0.62}}@{}}
\toprule
名称 & Python 类型 & 含义 \\
\midrule
\endhead
\passthrough{\lstinline!value!} &
\passthrough{\lstinline!Sequence[int]!} &
要写入或传递的新值，必须符合该参数的 Python 类型以及底层 C++
范围约束。 \\
\bottomrule
\end{longtable}

\textbf{返回值}

\begin{longtable}[]{@{}
  >{\raggedright\arraybackslash}p{(\columnwidth - 4\tabcolsep) * \real{0.18}}
  >{\raggedright\arraybackslash}p{(\columnwidth - 4\tabcolsep) * \real{0.20}}
  >{\raggedright\arraybackslash}p{(\columnwidth - 4\tabcolsep) * \real{0.62}}@{}}
\toprule
位置 & 类型 & 含义 \\
\midrule
\endhead
getter 返回值 & \passthrough{\lstinline!list[int]!} &
返回属性当前值。数组、vector 和嵌套消息采用复制语义。 \\
\bottomrule
\end{longtable}

\textbf{对应 C++}

\begin{itemize}
\tightlist
\item
  类：\passthrough{\lstinline!unitree\_go::msg::dds\_::LowCmd\_!}
\item
  字段类型：\passthrough{\lstinline!std::array<uint32\_t, 2>!}
\item
  声明位置：\passthrough{\lstinline!include/unitree/idl/go2/LowCmd\_.hpp:31!}
\end{itemize}

\textbf{用法}

\begin{lstlisting}[language=Python]
current_value = obj.sn
obj.sn = new_value
\end{lstlisting}

\begin{quote}
{[}!TIP{]} 读取容器或嵌套值后应遵循``读取、修改、重新赋值''。只修改
getter 返回的副本不会可靠地更新原 C++ 消息。
\end{quote}

\hypertarget{unitree-sdk2-cpp-idl-go2-lowcmd-version}{%
\paragraph{\texorpdfstring{\texttt{unitree\_sdk2\_cpp.idl.go2.LowCmd.version}}{unitree\_sdk2\_cpp.idl.go2.LowCmd.version}}\label{unitree-sdk2-cpp-idl-go2-lowcmd-version}}

底层 \passthrough{\lstinline!unitree\_go::msg::dds\_::LowCmd\_!} 的
\passthrough{\lstinline!version!} 字段。类型签名只说明 Python
表示；单位、数值范围、数组固定长度和枚举含义需查对应 IDL 头文件。
底层是固定长度数组，写入序列必须正好包含 2 个元素；元素约束：取值范围为
0 到 4294967295。

\textbf{签名}

\begin{lstlisting}[language=Python]
@property
def version(self) -> list[int]

@version.setter
def version(self, value: Sequence[int]) -> None
\end{lstlisting}

\textbf{可用性与安全性}

\textbf{\passthrough{\lstinline!AVAILABLE!} /
\passthrough{\lstinline!VALUE\_TYPE!}}：当前绑定源码已实现该消息值操作。

\textbf{参数}

\begin{longtable}[]{@{}
  >{\raggedright\arraybackslash}p{(\columnwidth - 4\tabcolsep) * \real{0.18}}
  >{\raggedright\arraybackslash}p{(\columnwidth - 4\tabcolsep) * \real{0.20}}
  >{\raggedright\arraybackslash}p{(\columnwidth - 4\tabcolsep) * \real{0.62}}@{}}
\toprule
名称 & Python 类型 & 含义 \\
\midrule
\endhead
\passthrough{\lstinline!value!} &
\passthrough{\lstinline!Sequence[int]!} &
要写入或传递的新值，必须符合该参数的 Python 类型以及底层 C++
范围约束。 \\
\bottomrule
\end{longtable}

\textbf{返回值}

\begin{longtable}[]{@{}
  >{\raggedright\arraybackslash}p{(\columnwidth - 4\tabcolsep) * \real{0.18}}
  >{\raggedright\arraybackslash}p{(\columnwidth - 4\tabcolsep) * \real{0.20}}
  >{\raggedright\arraybackslash}p{(\columnwidth - 4\tabcolsep) * \real{0.62}}@{}}
\toprule
位置 & 类型 & 含义 \\
\midrule
\endhead
getter 返回值 & \passthrough{\lstinline!list[int]!} &
返回属性当前值。数组、vector 和嵌套消息采用复制语义。 \\
\bottomrule
\end{longtable}

\textbf{对应 C++}

\begin{itemize}
\tightlist
\item
  类：\passthrough{\lstinline!unitree\_go::msg::dds\_::LowCmd\_!}
\item
  字段类型：\passthrough{\lstinline!std::array<uint32\_t, 2>!}
\item
  声明位置：\passthrough{\lstinline!include/unitree/idl/go2/LowCmd\_.hpp:32!}
\end{itemize}

\textbf{用法}

\begin{lstlisting}[language=Python]
current_value = obj.version
obj.version = new_value
\end{lstlisting}

\begin{quote}
{[}!TIP{]} 读取容器或嵌套值后应遵循``读取、修改、重新赋值''。只修改
getter 返回的副本不会可靠地更新原 C++ 消息。
\end{quote}

\hypertarget{unitree-sdk2-cpp-idl-go2-lowcmd-bandwidth}{%
\paragraph{\texorpdfstring{\texttt{unitree\_sdk2\_cpp.idl.go2.LowCmd.bandwidth}}{unitree\_sdk2\_cpp.idl.go2.LowCmd.bandwidth}}\label{unitree-sdk2-cpp-idl-go2-lowcmd-bandwidth}}

底层 \passthrough{\lstinline!unitree\_go::msg::dds\_::LowCmd\_!} 的
\passthrough{\lstinline!bandwidth!} 字段。类型签名只说明 Python
表示；单位、数值范围、数组固定长度和枚举含义需查对应 IDL 头文件。
取值范围为 0 到 65535。

\textbf{签名}

\begin{lstlisting}[language=Python]
@property
def bandwidth(self) -> int

@bandwidth.setter
def bandwidth(self, value: int) -> None
\end{lstlisting}

\textbf{可用性与安全性}

\textbf{\passthrough{\lstinline!AVAILABLE!} /
\passthrough{\lstinline!VALUE\_TYPE!}}：当前绑定源码已实现该消息值操作。

\textbf{参数}

\begin{longtable}[]{@{}
  >{\raggedright\arraybackslash}p{(\columnwidth - 4\tabcolsep) * \real{0.18}}
  >{\raggedright\arraybackslash}p{(\columnwidth - 4\tabcolsep) * \real{0.20}}
  >{\raggedright\arraybackslash}p{(\columnwidth - 4\tabcolsep) * \real{0.62}}@{}}
\toprule
名称 & Python 类型 & 含义 \\
\midrule
\endhead
\passthrough{\lstinline!value!} & \passthrough{\lstinline!int!} &
要写入或传递的新值，必须符合该参数的 Python 类型以及底层 C++
范围约束。 \\
\bottomrule
\end{longtable}

\textbf{返回值}

\begin{longtable}[]{@{}
  >{\raggedright\arraybackslash}p{(\columnwidth - 4\tabcolsep) * \real{0.18}}
  >{\raggedright\arraybackslash}p{(\columnwidth - 4\tabcolsep) * \real{0.20}}
  >{\raggedright\arraybackslash}p{(\columnwidth - 4\tabcolsep) * \real{0.62}}@{}}
\toprule
位置 & 类型 & 含义 \\
\midrule
\endhead
getter 返回值 & \passthrough{\lstinline!int!} & 返回属性当前值。 \\
\bottomrule
\end{longtable}

\textbf{对应 C++}

\begin{itemize}
\tightlist
\item
  类：\passthrough{\lstinline!unitree\_go::msg::dds\_::LowCmd\_!}
\item
  字段类型：\passthrough{\lstinline!uint16\_t!}
\item
  声明位置：\passthrough{\lstinline!include/unitree/idl/go2/LowCmd\_.hpp:33!}
\end{itemize}

\textbf{用法}

\begin{lstlisting}[language=Python]
current_value = obj.bandwidth
obj.bandwidth = new_value
\end{lstlisting}

\hypertarget{unitree-sdk2-cpp-idl-go2-lowcmd-motor-cmd}{%
\paragraph{\texorpdfstring{\texttt{unitree\_sdk2\_cpp.idl.go2.LowCmd.motor\_cmd}}{unitree\_sdk2\_cpp.idl.go2.LowCmd.motor\_cmd}}\label{unitree-sdk2-cpp-idl-go2-lowcmd-motor-cmd}}

底层 \passthrough{\lstinline!unitree\_go::msg::dds\_::LowCmd\_!} 的
\passthrough{\lstinline!motor\_cmd!} 字段。类型签名只说明 Python
表示；单位、数值范围、数组固定长度和枚举含义需查对应 IDL 头文件。
底层是固定长度数组，写入序列必须正好包含 20 个元素

\textbf{签名}

\begin{lstlisting}[language=Python]
@property
def motor_cmd(self) -> list[MotorCmd]

@motor_cmd.setter
def motor_cmd(self, value: Sequence[MotorCmd]) -> None
\end{lstlisting}

\textbf{可用性与安全性}

\textbf{\passthrough{\lstinline!AVAILABLE!} /
\passthrough{\lstinline!VALUE\_TYPE!}}：当前绑定源码已实现该消息值操作。

\textbf{参数}

\begin{longtable}[]{@{}
  >{\raggedright\arraybackslash}p{(\columnwidth - 4\tabcolsep) * \real{0.18}}
  >{\raggedright\arraybackslash}p{(\columnwidth - 4\tabcolsep) * \real{0.20}}
  >{\raggedright\arraybackslash}p{(\columnwidth - 4\tabcolsep) * \real{0.62}}@{}}
\toprule
名称 & Python 类型 & 含义 \\
\midrule
\endhead
\passthrough{\lstinline!value!} &
\passthrough{\lstinline!Sequence[MotorCmd]!} &
要写入或传递的新值，必须符合该参数的 Python 类型以及底层 C++
范围约束。 \\
\bottomrule
\end{longtable}

\textbf{返回值}

\begin{longtable}[]{@{}
  >{\raggedright\arraybackslash}p{(\columnwidth - 4\tabcolsep) * \real{0.18}}
  >{\raggedright\arraybackslash}p{(\columnwidth - 4\tabcolsep) * \real{0.20}}
  >{\raggedright\arraybackslash}p{(\columnwidth - 4\tabcolsep) * \real{0.62}}@{}}
\toprule
位置 & 类型 & 含义 \\
\midrule
\endhead
getter 返回值 & \passthrough{\lstinline!list[MotorCmd]!} &
返回属性当前值。数组、vector 和嵌套消息采用复制语义。 \\
\bottomrule
\end{longtable}

\textbf{对应 C++}

\begin{itemize}
\tightlist
\item
  类：\passthrough{\lstinline!unitree\_go::msg::dds\_::LowCmd\_!}
\item
  字段类型：\passthrough{\lstinline!std::array< ::unitree\_go::msg::dds\_::MotorCmd\_, 20>!}
\item
  声明位置：\passthrough{\lstinline!include/unitree/idl/go2/LowCmd\_.hpp:34!}
\end{itemize}

\textbf{用法}

\begin{lstlisting}[language=Python]
current_value = obj.motor_cmd
obj.motor_cmd = new_value
\end{lstlisting}

\begin{quote}
{[}!TIP{]} 读取容器或嵌套值后应遵循``读取、修改、重新赋值''。只修改
getter 返回的副本不会可靠地更新原 C++ 消息。
\end{quote}

\hypertarget{unitree-sdk2-cpp-idl-go2-lowcmd-bms-cmd}{%
\paragraph{\texorpdfstring{\texttt{unitree\_sdk2\_cpp.idl.go2.LowCmd.bms\_cmd}}{unitree\_sdk2\_cpp.idl.go2.LowCmd.bms\_cmd}}\label{unitree-sdk2-cpp-idl-go2-lowcmd-bms-cmd}}

底层 \passthrough{\lstinline!unitree\_go::msg::dds\_::LowCmd\_!} 的
\passthrough{\lstinline!bms\_cmd!} 字段。类型签名只说明 Python
表示；单位、数值范围、数组固定长度和枚举含义需查对应 IDL 头文件。

\textbf{签名}

\begin{lstlisting}[language=Python]
@property
def bms_cmd(self) -> BmsCmd

@bms_cmd.setter
def bms_cmd(self, value: BmsCmd) -> None
\end{lstlisting}

\textbf{可用性与安全性}

\textbf{\passthrough{\lstinline!AVAILABLE!} /
\passthrough{\lstinline!VALUE\_TYPE!}}：当前绑定源码已实现该消息值操作。

\textbf{参数}

\begin{longtable}[]{@{}
  >{\raggedright\arraybackslash}p{(\columnwidth - 4\tabcolsep) * \real{0.18}}
  >{\raggedright\arraybackslash}p{(\columnwidth - 4\tabcolsep) * \real{0.20}}
  >{\raggedright\arraybackslash}p{(\columnwidth - 4\tabcolsep) * \real{0.62}}@{}}
\toprule
名称 & Python 类型 & 含义 \\
\midrule
\endhead
\passthrough{\lstinline!value!} & \passthrough{\lstinline!BmsCmd!} &
要写入或传递的新值，必须符合该参数的 Python 类型以及底层 C++
范围约束。 \\
\bottomrule
\end{longtable}

\textbf{返回值}

\begin{longtable}[]{@{}
  >{\raggedright\arraybackslash}p{(\columnwidth - 4\tabcolsep) * \real{0.18}}
  >{\raggedright\arraybackslash}p{(\columnwidth - 4\tabcolsep) * \real{0.20}}
  >{\raggedright\arraybackslash}p{(\columnwidth - 4\tabcolsep) * \real{0.62}}@{}}
\toprule
位置 & 类型 & 含义 \\
\midrule
\endhead
getter 返回值 & \passthrough{\lstinline!BmsCmd!} &
返回属性当前值。数组、vector 和嵌套消息采用复制语义。 \\
\bottomrule
\end{longtable}

\textbf{对应 C++}

\begin{itemize}
\tightlist
\item
  类：\passthrough{\lstinline!unitree\_go::msg::dds\_::LowCmd\_!}
\item
  字段类型：\passthrough{\lstinline!::unitree\_go::msg::dds\_::BmsCmd\_!}
\item
  声明位置：\passthrough{\lstinline!include/unitree/idl/go2/LowCmd\_.hpp:35!}
\end{itemize}

\textbf{用法}

\begin{lstlisting}[language=Python]
current_value = obj.bms_cmd
obj.bms_cmd = new_value
\end{lstlisting}

\hypertarget{unitree-sdk2-cpp-idl-go2-lowcmd-wireless-remote}{%
\paragraph{\texorpdfstring{\texttt{unitree\_sdk2\_cpp.idl.go2.LowCmd.wireless\_remote}}{unitree\_sdk2\_cpp.idl.go2.LowCmd.wireless\_remote}}\label{unitree-sdk2-cpp-idl-go2-lowcmd-wireless-remote}}

底层 \passthrough{\lstinline!unitree\_go::msg::dds\_::LowCmd\_!} 的
\passthrough{\lstinline!wireless\_remote!} 字段。类型签名只说明 Python
表示；单位、数值范围、数组固定长度和枚举含义需查对应 IDL 头文件。
底层是固定长度数组，写入序列必须正好包含 40 个元素；元素约束：取值范围为
0 到 255。

\textbf{签名}

\begin{lstlisting}[language=Python]
@property
def wireless_remote(self) -> list[int]

@wireless_remote.setter
def wireless_remote(self, value: Sequence[int]) -> None
\end{lstlisting}

\textbf{可用性与安全性}

\textbf{\passthrough{\lstinline!AVAILABLE!} /
\passthrough{\lstinline!VALUE\_TYPE!}}：当前绑定源码已实现该消息值操作。

\textbf{参数}

\begin{longtable}[]{@{}
  >{\raggedright\arraybackslash}p{(\columnwidth - 4\tabcolsep) * \real{0.18}}
  >{\raggedright\arraybackslash}p{(\columnwidth - 4\tabcolsep) * \real{0.20}}
  >{\raggedright\arraybackslash}p{(\columnwidth - 4\tabcolsep) * \real{0.62}}@{}}
\toprule
名称 & Python 类型 & 含义 \\
\midrule
\endhead
\passthrough{\lstinline!value!} &
\passthrough{\lstinline!Sequence[int]!} &
要写入或传递的新值，必须符合该参数的 Python 类型以及底层 C++
范围约束。 \\
\bottomrule
\end{longtable}

\textbf{返回值}

\begin{longtable}[]{@{}
  >{\raggedright\arraybackslash}p{(\columnwidth - 4\tabcolsep) * \real{0.18}}
  >{\raggedright\arraybackslash}p{(\columnwidth - 4\tabcolsep) * \real{0.20}}
  >{\raggedright\arraybackslash}p{(\columnwidth - 4\tabcolsep) * \real{0.62}}@{}}
\toprule
位置 & 类型 & 含义 \\
\midrule
\endhead
getter 返回值 & \passthrough{\lstinline!list[int]!} &
返回属性当前值。数组、vector 和嵌套消息采用复制语义。 \\
\bottomrule
\end{longtable}

\textbf{对应 C++}

\begin{itemize}
\tightlist
\item
  类：\passthrough{\lstinline!unitree\_go::msg::dds\_::LowCmd\_!}
\item
  字段类型：\passthrough{\lstinline!std::array<uint8\_t, 40>!}
\item
  声明位置：\passthrough{\lstinline!include/unitree/idl/go2/LowCmd\_.hpp:36!}
\end{itemize}

\textbf{用法}

\begin{lstlisting}[language=Python]
current_value = obj.wireless_remote
obj.wireless_remote = new_value
\end{lstlisting}

\begin{quote}
{[}!TIP{]} 读取容器或嵌套值后应遵循``读取、修改、重新赋值''。只修改
getter 返回的副本不会可靠地更新原 C++ 消息。
\end{quote}

\hypertarget{unitree-sdk2-cpp-idl-go2-lowcmd-led}{%
\paragraph{\texorpdfstring{\texttt{unitree\_sdk2\_cpp.idl.go2.LowCmd.led}}{unitree\_sdk2\_cpp.idl.go2.LowCmd.led}}\label{unitree-sdk2-cpp-idl-go2-lowcmd-led}}

底层 \passthrough{\lstinline!unitree\_go::msg::dds\_::LowCmd\_!} 的
\passthrough{\lstinline!led!} 字段。类型签名只说明 Python
表示；单位、数值范围、数组固定长度和枚举含义需查对应 IDL 头文件。
底层是固定长度数组，写入序列必须正好包含 12 个元素；元素约束：取值范围为
0 到 255。

\textbf{签名}

\begin{lstlisting}[language=Python]
@property
def led(self) -> list[int]

@led.setter
def led(self, value: Sequence[int]) -> None
\end{lstlisting}

\textbf{可用性与安全性}

\textbf{\passthrough{\lstinline!AVAILABLE!} /
\passthrough{\lstinline!VALUE\_TYPE!}}：当前绑定源码已实现该消息值操作。

\textbf{参数}

\begin{longtable}[]{@{}
  >{\raggedright\arraybackslash}p{(\columnwidth - 4\tabcolsep) * \real{0.18}}
  >{\raggedright\arraybackslash}p{(\columnwidth - 4\tabcolsep) * \real{0.20}}
  >{\raggedright\arraybackslash}p{(\columnwidth - 4\tabcolsep) * \real{0.62}}@{}}
\toprule
名称 & Python 类型 & 含义 \\
\midrule
\endhead
\passthrough{\lstinline!value!} &
\passthrough{\lstinline!Sequence[int]!} &
要写入或传递的新值，必须符合该参数的 Python 类型以及底层 C++
范围约束。 \\
\bottomrule
\end{longtable}

\textbf{返回值}

\begin{longtable}[]{@{}
  >{\raggedright\arraybackslash}p{(\columnwidth - 4\tabcolsep) * \real{0.18}}
  >{\raggedright\arraybackslash}p{(\columnwidth - 4\tabcolsep) * \real{0.20}}
  >{\raggedright\arraybackslash}p{(\columnwidth - 4\tabcolsep) * \real{0.62}}@{}}
\toprule
位置 & 类型 & 含义 \\
\midrule
\endhead
getter 返回值 & \passthrough{\lstinline!list[int]!} &
返回属性当前值。数组、vector 和嵌套消息采用复制语义。 \\
\bottomrule
\end{longtable}

\textbf{对应 C++}

\begin{itemize}
\tightlist
\item
  类：\passthrough{\lstinline!unitree\_go::msg::dds\_::LowCmd\_!}
\item
  字段类型：\passthrough{\lstinline!std::array<uint8\_t, 12>!}
\item
  声明位置：\passthrough{\lstinline!include/unitree/idl/go2/LowCmd\_.hpp:37!}
\end{itemize}

\textbf{用法}

\begin{lstlisting}[language=Python]
current_value = obj.led
obj.led = new_value
\end{lstlisting}

\begin{quote}
{[}!TIP{]} 读取容器或嵌套值后应遵循``读取、修改、重新赋值''。只修改
getter 返回的副本不会可靠地更新原 C++ 消息。
\end{quote}

\hypertarget{unitree-sdk2-cpp-idl-go2-lowcmd-fan}{%
\paragraph{\texorpdfstring{\texttt{unitree\_sdk2\_cpp.idl.go2.LowCmd.fan}}{unitree\_sdk2\_cpp.idl.go2.LowCmd.fan}}\label{unitree-sdk2-cpp-idl-go2-lowcmd-fan}}

底层 \passthrough{\lstinline!unitree\_go::msg::dds\_::LowCmd\_!} 的
\passthrough{\lstinline!fan!} 字段。类型签名只说明 Python
表示；单位、数值范围、数组固定长度和枚举含义需查对应 IDL 头文件。
底层是固定长度数组，写入序列必须正好包含 2 个元素；元素约束：取值范围为
0 到 255。

\textbf{签名}

\begin{lstlisting}[language=Python]
@property
def fan(self) -> list[int]

@fan.setter
def fan(self, value: Sequence[int]) -> None
\end{lstlisting}

\textbf{可用性与安全性}

\textbf{\passthrough{\lstinline!AVAILABLE!} /
\passthrough{\lstinline!VALUE\_TYPE!}}：当前绑定源码已实现该消息值操作。

\textbf{参数}

\begin{longtable}[]{@{}
  >{\raggedright\arraybackslash}p{(\columnwidth - 4\tabcolsep) * \real{0.18}}
  >{\raggedright\arraybackslash}p{(\columnwidth - 4\tabcolsep) * \real{0.20}}
  >{\raggedright\arraybackslash}p{(\columnwidth - 4\tabcolsep) * \real{0.62}}@{}}
\toprule
名称 & Python 类型 & 含义 \\
\midrule
\endhead
\passthrough{\lstinline!value!} &
\passthrough{\lstinline!Sequence[int]!} &
要写入或传递的新值，必须符合该参数的 Python 类型以及底层 C++
范围约束。 \\
\bottomrule
\end{longtable}

\textbf{返回值}

\begin{longtable}[]{@{}
  >{\raggedright\arraybackslash}p{(\columnwidth - 4\tabcolsep) * \real{0.18}}
  >{\raggedright\arraybackslash}p{(\columnwidth - 4\tabcolsep) * \real{0.20}}
  >{\raggedright\arraybackslash}p{(\columnwidth - 4\tabcolsep) * \real{0.62}}@{}}
\toprule
位置 & 类型 & 含义 \\
\midrule
\endhead
getter 返回值 & \passthrough{\lstinline!list[int]!} &
返回属性当前值。数组、vector 和嵌套消息采用复制语义。 \\
\bottomrule
\end{longtable}

\textbf{对应 C++}

\begin{itemize}
\tightlist
\item
  类：\passthrough{\lstinline!unitree\_go::msg::dds\_::LowCmd\_!}
\item
  字段类型：\passthrough{\lstinline!std::array<uint8\_t, 2>!}
\item
  声明位置：\passthrough{\lstinline!include/unitree/idl/go2/LowCmd\_.hpp:38!}
\end{itemize}

\textbf{用法}

\begin{lstlisting}[language=Python]
current_value = obj.fan
obj.fan = new_value
\end{lstlisting}

\begin{quote}
{[}!TIP{]} 读取容器或嵌套值后应遵循``读取、修改、重新赋值''。只修改
getter 返回的副本不会可靠地更新原 C++ 消息。
\end{quote}

\hypertarget{unitree-sdk2-cpp-idl-go2-lowcmd-gpio}{%
\paragraph{\texorpdfstring{\texttt{unitree\_sdk2\_cpp.idl.go2.LowCmd.gpio}}{unitree\_sdk2\_cpp.idl.go2.LowCmd.gpio}}\label{unitree-sdk2-cpp-idl-go2-lowcmd-gpio}}

底层 \passthrough{\lstinline!unitree\_go::msg::dds\_::LowCmd\_!} 的
\passthrough{\lstinline!gpio!} 字段。类型签名只说明 Python
表示；单位、数值范围、数组固定长度和枚举含义需查对应 IDL 头文件。
取值范围为 0 到 255。

\textbf{签名}

\begin{lstlisting}[language=Python]
@property
def gpio(self) -> int

@gpio.setter
def gpio(self, value: int) -> None
\end{lstlisting}

\textbf{可用性与安全性}

\textbf{\passthrough{\lstinline!AVAILABLE!} /
\passthrough{\lstinline!VALUE\_TYPE!}}：当前绑定源码已实现该消息值操作。

\textbf{参数}

\begin{longtable}[]{@{}
  >{\raggedright\arraybackslash}p{(\columnwidth - 4\tabcolsep) * \real{0.18}}
  >{\raggedright\arraybackslash}p{(\columnwidth - 4\tabcolsep) * \real{0.20}}
  >{\raggedright\arraybackslash}p{(\columnwidth - 4\tabcolsep) * \real{0.62}}@{}}
\toprule
名称 & Python 类型 & 含义 \\
\midrule
\endhead
\passthrough{\lstinline!value!} & \passthrough{\lstinline!int!} &
要写入或传递的新值，必须符合该参数的 Python 类型以及底层 C++
范围约束。 \\
\bottomrule
\end{longtable}

\textbf{返回值}

\begin{longtable}[]{@{}
  >{\raggedright\arraybackslash}p{(\columnwidth - 4\tabcolsep) * \real{0.18}}
  >{\raggedright\arraybackslash}p{(\columnwidth - 4\tabcolsep) * \real{0.20}}
  >{\raggedright\arraybackslash}p{(\columnwidth - 4\tabcolsep) * \real{0.62}}@{}}
\toprule
位置 & 类型 & 含义 \\
\midrule
\endhead
getter 返回值 & \passthrough{\lstinline!int!} & 返回属性当前值。 \\
\bottomrule
\end{longtable}

\textbf{对应 C++}

\begin{itemize}
\tightlist
\item
  类：\passthrough{\lstinline!unitree\_go::msg::dds\_::LowCmd\_!}
\item
  字段类型：\passthrough{\lstinline!uint8\_t!}
\item
  声明位置：\passthrough{\lstinline!include/unitree/idl/go2/LowCmd\_.hpp:39!}
\end{itemize}

\textbf{用法}

\begin{lstlisting}[language=Python]
current_value = obj.gpio
obj.gpio = new_value
\end{lstlisting}

\hypertarget{unitree-sdk2-cpp-idl-go2-lowcmd-reserve}{%
\paragraph{\texorpdfstring{\texttt{unitree\_sdk2\_cpp.idl.go2.LowCmd.reserve}}{unitree\_sdk2\_cpp.idl.go2.LowCmd.reserve}}\label{unitree-sdk2-cpp-idl-go2-lowcmd-reserve}}

底层 \passthrough{\lstinline!unitree\_go::msg::dds\_::LowCmd\_!} 的
\passthrough{\lstinline!reserve!} 字段。类型签名只说明 Python
表示；单位、数值范围、数组固定长度和枚举含义需查对应 IDL 头文件。
取值范围为 0 到 4294967295。

\textbf{签名}

\begin{lstlisting}[language=Python]
@property
def reserve(self) -> int

@reserve.setter
def reserve(self, value: int) -> None
\end{lstlisting}

\textbf{可用性与安全性}

\textbf{\passthrough{\lstinline!AVAILABLE!} /
\passthrough{\lstinline!VALUE\_TYPE!}}：当前绑定源码已实现该消息值操作。

\textbf{参数}

\begin{longtable}[]{@{}
  >{\raggedright\arraybackslash}p{(\columnwidth - 4\tabcolsep) * \real{0.18}}
  >{\raggedright\arraybackslash}p{(\columnwidth - 4\tabcolsep) * \real{0.20}}
  >{\raggedright\arraybackslash}p{(\columnwidth - 4\tabcolsep) * \real{0.62}}@{}}
\toprule
名称 & Python 类型 & 含义 \\
\midrule
\endhead
\passthrough{\lstinline!value!} & \passthrough{\lstinline!int!} &
要写入或传递的新值，必须符合该参数的 Python 类型以及底层 C++
范围约束。 \\
\bottomrule
\end{longtable}

\textbf{返回值}

\begin{longtable}[]{@{}
  >{\raggedright\arraybackslash}p{(\columnwidth - 4\tabcolsep) * \real{0.18}}
  >{\raggedright\arraybackslash}p{(\columnwidth - 4\tabcolsep) * \real{0.20}}
  >{\raggedright\arraybackslash}p{(\columnwidth - 4\tabcolsep) * \real{0.62}}@{}}
\toprule
位置 & 类型 & 含义 \\
\midrule
\endhead
getter 返回值 & \passthrough{\lstinline!int!} & 返回属性当前值。 \\
\bottomrule
\end{longtable}

\textbf{对应 C++}

\begin{itemize}
\tightlist
\item
  类：\passthrough{\lstinline!unitree\_go::msg::dds\_::LowCmd\_!}
\item
  字段类型：\passthrough{\lstinline!uint32\_t!}
\item
  声明位置：\passthrough{\lstinline!include/unitree/idl/go2/LowCmd\_.hpp:40!}
\end{itemize}

\textbf{用法}

\begin{lstlisting}[language=Python]
current_value = obj.reserve
obj.reserve = new_value
\end{lstlisting}

\hypertarget{unitree-sdk2-cpp-idl-go2-lowcmd-crc}{%
\paragraph{\texorpdfstring{\texttt{unitree\_sdk2\_cpp.idl.go2.LowCmd.crc}}{unitree\_sdk2\_cpp.idl.go2.LowCmd.crc}}\label{unitree-sdk2-cpp-idl-go2-lowcmd-crc}}

底层 \passthrough{\lstinline!unitree\_go::msg::dds\_::LowCmd\_!} 的
\passthrough{\lstinline!crc!} 字段。类型签名只说明 Python
表示；单位、数值范围、数组固定长度和枚举含义需查对应 IDL 头文件。
取值范围为 0 到 4294967295。

\textbf{签名}

\begin{lstlisting}[language=Python]
@property
def crc(self) -> int

@crc.setter
def crc(self, value: int) -> None
\end{lstlisting}

\textbf{可用性与安全性}

\textbf{\passthrough{\lstinline!AVAILABLE!} /
\passthrough{\lstinline!VALUE\_TYPE!}}：当前绑定源码已实现该消息值操作。

\textbf{参数}

\begin{longtable}[]{@{}
  >{\raggedright\arraybackslash}p{(\columnwidth - 4\tabcolsep) * \real{0.18}}
  >{\raggedright\arraybackslash}p{(\columnwidth - 4\tabcolsep) * \real{0.20}}
  >{\raggedright\arraybackslash}p{(\columnwidth - 4\tabcolsep) * \real{0.62}}@{}}
\toprule
名称 & Python 类型 & 含义 \\
\midrule
\endhead
\passthrough{\lstinline!value!} & \passthrough{\lstinline!int!} &
要写入或传递的新值，必须符合该参数的 Python 类型以及底层 C++
范围约束。 \\
\bottomrule
\end{longtable}

\textbf{返回值}

\begin{longtable}[]{@{}
  >{\raggedright\arraybackslash}p{(\columnwidth - 4\tabcolsep) * \real{0.18}}
  >{\raggedright\arraybackslash}p{(\columnwidth - 4\tabcolsep) * \real{0.20}}
  >{\raggedright\arraybackslash}p{(\columnwidth - 4\tabcolsep) * \real{0.62}}@{}}
\toprule
位置 & 类型 & 含义 \\
\midrule
\endhead
getter 返回值 & \passthrough{\lstinline!int!} & 返回属性当前值。 \\
\bottomrule
\end{longtable}

\textbf{对应 C++}

\begin{itemize}
\tightlist
\item
  类：\passthrough{\lstinline!unitree\_go::msg::dds\_::LowCmd\_!}
\item
  字段类型：\passthrough{\lstinline!uint32\_t!}
\item
  声明位置：\passthrough{\lstinline!include/unitree/idl/go2/LowCmd\_.hpp:41!}
\end{itemize}

\textbf{用法}

\begin{lstlisting}[language=Python]
current_value = obj.crc
obj.crc = new_value
\end{lstlisting}

\hypertarget{unitree-sdk2-cpp-idl-go2-motorstate}{%
\subsubsection{\texorpdfstring{\texttt{unitree\_sdk2\_cpp.idl.go2.MotorState}}{unitree\_sdk2\_cpp.idl.go2.MotorState}}\label{unitree-sdk2-cpp-idl-go2-motorstate}}

可由 Python 读写的 DDS/IDL 消息值类型。字段采用复制语义。

\textbf{导入}

\begin{lstlisting}[language=Python]
from unitree_sdk2_cpp.idl.go2 import MotorState
\end{lstlisting}

\textbf{方法索引}

\begin{longtable}[]{@{}
  >{\raggedright\arraybackslash}p{(\columnwidth - 6\tabcolsep) * \real{0.18}}
  >{\raggedright\arraybackslash}p{(\columnwidth - 6\tabcolsep) * \real{0.24}}
  >{\raggedright\arraybackslash}p{(\columnwidth - 6\tabcolsep) * \real{0.18}}
  >{\raggedright\arraybackslash}p{(\columnwidth - 6\tabcolsep) * \real{0.40}}@{}}
\toprule
方法 & Python 签名 & 状态 & 安全分类 \\
\midrule
\endhead
\protect\hyperlink{unitree-sdk2-cpp-idl-go2-motorstate-dunder-init-1}{\passthrough{\lstinline!\_\_init\_\_!}}
& \passthrough{\lstinline!def \_\_init\_\_(self) -> None!} &
\passthrough{\lstinline!AVAILABLE!} &
\passthrough{\lstinline!VALUE\_TYPE!} \\
\protect\hyperlink{unitree-sdk2-cpp-idl-go2-motorstate-dunder-eq-1}{\passthrough{\lstinline!\_\_eq\_\_!}}
& \passthrough{\lstinline!def \_\_eq\_\_(self, other: object) -> bool!}
& \passthrough{\lstinline!AVAILABLE!} &
\passthrough{\lstinline!VALUE\_TYPE!} \\
\protect\hyperlink{unitree-sdk2-cpp-idl-go2-motorstate-dunder-ne-1}{\passthrough{\lstinline!\_\_ne\_\_!}}
& \passthrough{\lstinline!def \_\_ne\_\_(self, other: object) -> bool!}
& \passthrough{\lstinline!AVAILABLE!} &
\passthrough{\lstinline!VALUE\_TYPE!} \\
\bottomrule
\end{longtable}

\textbf{属性索引}

\begin{longtable}[]{@{}
  >{\raggedright\arraybackslash}p{(\columnwidth - 6\tabcolsep) * \real{0.18}}
  >{\raggedright\arraybackslash}p{(\columnwidth - 6\tabcolsep) * \real{0.24}}
  >{\raggedright\arraybackslash}p{(\columnwidth - 6\tabcolsep) * \real{0.18}}
  >{\raggedright\arraybackslash}p{(\columnwidth - 6\tabcolsep) * \real{0.40}}@{}}
\toprule
属性 & 读取类型 & 写入类型 & 状态 \\
\midrule
\endhead
\protect\hyperlink{unitree-sdk2-cpp-idl-go2-motorstate-mode}{\passthrough{\lstinline!mode!}}
& \passthrough{\lstinline!int!} & \passthrough{\lstinline!int!} &
\passthrough{\lstinline!AVAILABLE!} \\
\protect\hyperlink{unitree-sdk2-cpp-idl-go2-motorstate-q}{\passthrough{\lstinline!q!}}
& \passthrough{\lstinline!float!} & \passthrough{\lstinline!float!} &
\passthrough{\lstinline!AVAILABLE!} \\
\protect\hyperlink{unitree-sdk2-cpp-idl-go2-motorstate-dq}{\passthrough{\lstinline!dq!}}
& \passthrough{\lstinline!float!} & \passthrough{\lstinline!float!} &
\passthrough{\lstinline!AVAILABLE!} \\
\protect\hyperlink{unitree-sdk2-cpp-idl-go2-motorstate-ddq}{\passthrough{\lstinline!ddq!}}
& \passthrough{\lstinline!float!} & \passthrough{\lstinline!float!} &
\passthrough{\lstinline!AVAILABLE!} \\
\protect\hyperlink{unitree-sdk2-cpp-idl-go2-motorstate-tau-est}{\passthrough{\lstinline!tau\_est!}}
& \passthrough{\lstinline!float!} & \passthrough{\lstinline!float!} &
\passthrough{\lstinline!AVAILABLE!} \\
\protect\hyperlink{unitree-sdk2-cpp-idl-go2-motorstate-q-raw}{\passthrough{\lstinline!q\_raw!}}
& \passthrough{\lstinline!float!} & \passthrough{\lstinline!float!} &
\passthrough{\lstinline!AVAILABLE!} \\
\protect\hyperlink{unitree-sdk2-cpp-idl-go2-motorstate-dq-raw}{\passthrough{\lstinline!dq\_raw!}}
& \passthrough{\lstinline!float!} & \passthrough{\lstinline!float!} &
\passthrough{\lstinline!AVAILABLE!} \\
\protect\hyperlink{unitree-sdk2-cpp-idl-go2-motorstate-ddq-raw}{\passthrough{\lstinline!ddq\_raw!}}
& \passthrough{\lstinline!float!} & \passthrough{\lstinline!float!} &
\passthrough{\lstinline!AVAILABLE!} \\
\protect\hyperlink{unitree-sdk2-cpp-idl-go2-motorstate-temperature}{\passthrough{\lstinline!temperature!}}
& \passthrough{\lstinline!int!} & \passthrough{\lstinline!int!} &
\passthrough{\lstinline!AVAILABLE!} \\
\protect\hyperlink{unitree-sdk2-cpp-idl-go2-motorstate-lost}{\passthrough{\lstinline!lost!}}
& \passthrough{\lstinline!int!} & \passthrough{\lstinline!int!} &
\passthrough{\lstinline!AVAILABLE!} \\
\protect\hyperlink{unitree-sdk2-cpp-idl-go2-motorstate-reserve}{\passthrough{\lstinline!reserve!}}
& \passthrough{\lstinline!list[int]!} &
\passthrough{\lstinline!Sequence[int]!} &
\passthrough{\lstinline!AVAILABLE!} \\
\bottomrule
\end{longtable}

\hypertarget{unitree-sdk2-cpp-idl-go2-motorstate-dunder-init-1}{%
\paragraph{\texorpdfstring{\texttt{unitree\_sdk2\_cpp.idl.go2.MotorState.\_\_init\_\_}}{unitree\_sdk2\_cpp.idl.go2.MotorState.\_\_init\_\_}}\label{unitree-sdk2-cpp-idl-go2-motorstate-dunder-init-1}}

初始化 \passthrough{\lstinline!MotorState!}
实例。是否能实际构造取决于下方可用性状态。

\textbf{签名}

\begin{lstlisting}[language=Python]
def __init__(self) -> None
\end{lstlisting}

\textbf{可用性与安全性}

\textbf{\passthrough{\lstinline!AVAILABLE!} /
\passthrough{\lstinline!VALUE\_TYPE!}}：当前绑定源码已实现该消息值操作。

\textbf{参数}

无显式参数。实例方法中的 \passthrough{\lstinline!self!} 由 Python
自动传入。

\textbf{返回值}

\begin{longtable}[]{@{}
  >{\raggedright\arraybackslash}p{(\columnwidth - 4\tabcolsep) * \real{0.18}}
  >{\raggedright\arraybackslash}p{(\columnwidth - 4\tabcolsep) * \real{0.20}}
  >{\raggedright\arraybackslash}p{(\columnwidth - 4\tabcolsep) * \real{0.62}}@{}}
\toprule
位置 & 类型 & 含义 \\
\midrule
\endhead
无 & \passthrough{\lstinline!None!} &
\passthrough{\lstinline!\_\_init\_\_!}
本身不返回值；调用类对象时，在构造函数可用的前提下得到该类实例。 \\
\bottomrule
\end{longtable}

\textbf{对应 C++}

\begin{itemize}
\tightlist
\item
  类：\passthrough{\lstinline!unitree\_go::msg::dds\_::MotorState\_!}
\item
  签名：\passthrough{\lstinline!MotorState\_()!}
\item
  绑定策略：\passthrough{\lstinline!未单独标注!}
\item
  声明位置：\passthrough{\lstinline!include/unitree/idl/go2/MotorState\_.hpp:37!}
\end{itemize}

\textbf{说明}

\begin{itemize}
\tightlist
\item
  示例是局部调用片段；\passthrough{\lstinline!obj!}
  和参数变量表示已经按上层业务校验并准备好的对象。
\end{itemize}

\textbf{用法}

\begin{lstlisting}[language=Python]
value = MotorState()
\end{lstlisting}

\hypertarget{unitree-sdk2-cpp-idl-go2-motorstate-dunder-eq-1}{%
\paragraph{\texorpdfstring{\texttt{unitree\_sdk2\_cpp.idl.go2.MotorState.\_\_eq\_\_}}{unitree\_sdk2\_cpp.idl.go2.MotorState.\_\_eq\_\_}}\label{unitree-sdk2-cpp-idl-go2-motorstate-dunder-eq-1}}

按底层 IDL 消息内容比较两个对象是否相等。

\textbf{签名}

\begin{lstlisting}[language=Python]
def __eq__(self, other: object) -> bool
\end{lstlisting}

\textbf{可用性与安全性}

\textbf{\passthrough{\lstinline!AVAILABLE!} /
\passthrough{\lstinline!VALUE\_TYPE!}}：当前绑定源码已实现该消息值操作。

\textbf{参数}

\begin{longtable}[]{@{}
  >{\raggedright\arraybackslash}p{(\columnwidth - 6\tabcolsep) * \real{0.18}}
  >{\raggedright\arraybackslash}p{(\columnwidth - 6\tabcolsep) * \real{0.24}}
  >{\raggedright\arraybackslash}p{(\columnwidth - 6\tabcolsep) * \real{0.18}}
  >{\raggedright\arraybackslash}p{(\columnwidth - 6\tabcolsep) * \real{0.40}}@{}}
\toprule
名称 & Python 类型 & 默认值 & 含义 \\
\midrule
\endhead
\passthrough{\lstinline!other!} & \passthrough{\lstinline!object!} &
必填 & 用于比较的另一个 Python 对象。类型不兼容时比较结果通常为
\passthrough{\lstinline!False!}。 \\
\bottomrule
\end{longtable}

\textbf{返回值}

\begin{longtable}[]{@{}
  >{\raggedright\arraybackslash}p{(\columnwidth - 4\tabcolsep) * \real{0.18}}
  >{\raggedright\arraybackslash}p{(\columnwidth - 4\tabcolsep) * \real{0.20}}
  >{\raggedright\arraybackslash}p{(\columnwidth - 4\tabcolsep) * \real{0.62}}@{}}
\toprule
位置 & 类型 & 含义 \\
\midrule
\endhead
返回值 & \passthrough{\lstinline!bool!} & 返回
\passthrough{\lstinline!bool!}。更精确的业务含义见该方法用途、C++
签名和目标型号协议。 \\
\bottomrule
\end{longtable}

\textbf{对应 C++}

\begin{itemize}
\tightlist
\item
  类：\passthrough{\lstinline!unitree\_go::msg::dds\_::MotorState\_!}
\item
  签名：\passthrough{\lstinline!operator==(const value \&) const!}
\item
  绑定策略：\passthrough{\lstinline!未单独标注!}
\item
  声明位置：由 pybind11 手工绑定或生成注册表提供；manifest
  未记录独立头文件行号。
\end{itemize}

\textbf{说明}

\begin{itemize}
\tightlist
\item
  示例是局部调用片段；\passthrough{\lstinline!obj!}
  和参数变量表示已经按上层业务校验并准备好的对象。
\end{itemize}

\textbf{用法}

\begin{lstlisting}[language=Python]
same = left == right
\end{lstlisting}

\hypertarget{unitree-sdk2-cpp-idl-go2-motorstate-dunder-ne-1}{%
\paragraph{\texorpdfstring{\texttt{unitree\_sdk2\_cpp.idl.go2.MotorState.\_\_ne\_\_}}{unitree\_sdk2\_cpp.idl.go2.MotorState.\_\_ne\_\_}}\label{unitree-sdk2-cpp-idl-go2-motorstate-dunder-ne-1}}

按底层 IDL 消息内容比较两个对象是否不相等。

\textbf{签名}

\begin{lstlisting}[language=Python]
def __ne__(self, other: object) -> bool
\end{lstlisting}

\textbf{可用性与安全性}

\textbf{\passthrough{\lstinline!AVAILABLE!} /
\passthrough{\lstinline!VALUE\_TYPE!}}：当前绑定源码已实现该消息值操作。

\textbf{参数}

\begin{longtable}[]{@{}
  >{\raggedright\arraybackslash}p{(\columnwidth - 6\tabcolsep) * \real{0.18}}
  >{\raggedright\arraybackslash}p{(\columnwidth - 6\tabcolsep) * \real{0.24}}
  >{\raggedright\arraybackslash}p{(\columnwidth - 6\tabcolsep) * \real{0.18}}
  >{\raggedright\arraybackslash}p{(\columnwidth - 6\tabcolsep) * \real{0.40}}@{}}
\toprule
名称 & Python 类型 & 默认值 & 含义 \\
\midrule
\endhead
\passthrough{\lstinline!other!} & \passthrough{\lstinline!object!} &
必填 & 用于比较的另一个 Python 对象。类型不兼容时比较结果通常为
\passthrough{\lstinline!False!}。 \\
\bottomrule
\end{longtable}

\textbf{返回值}

\begin{longtable}[]{@{}
  >{\raggedright\arraybackslash}p{(\columnwidth - 4\tabcolsep) * \real{0.18}}
  >{\raggedright\arraybackslash}p{(\columnwidth - 4\tabcolsep) * \real{0.20}}
  >{\raggedright\arraybackslash}p{(\columnwidth - 4\tabcolsep) * \real{0.62}}@{}}
\toprule
位置 & 类型 & 含义 \\
\midrule
\endhead
返回值 & \passthrough{\lstinline!bool!} & 返回
\passthrough{\lstinline!bool!}。更精确的业务含义见该方法用途、C++
签名和目标型号协议。 \\
\bottomrule
\end{longtable}

\textbf{对应 C++}

\begin{itemize}
\tightlist
\item
  类：\passthrough{\lstinline!unitree\_go::msg::dds\_::MotorState\_!}
\item
  签名：\passthrough{\lstinline"operator!=(const value \&) const"}
\item
  绑定策略：\passthrough{\lstinline!未单独标注!}
\item
  声明位置：由 pybind11 手工绑定或生成注册表提供；manifest
  未记录独立头文件行号。
\end{itemize}

\textbf{说明}

\begin{itemize}
\tightlist
\item
  示例是局部调用片段；\passthrough{\lstinline!obj!}
  和参数变量表示已经按上层业务校验并准备好的对象。
\end{itemize}

\textbf{用法}

\begin{lstlisting}[language=Python]
different = left != right
\end{lstlisting}

\hypertarget{unitree-sdk2-cpp-idl-go2-motorstate-mode}{%
\paragraph{\texorpdfstring{\texttt{unitree\_sdk2\_cpp.idl.go2.MotorState.mode}}{unitree\_sdk2\_cpp.idl.go2.MotorState.mode}}\label{unitree-sdk2-cpp-idl-go2-motorstate-mode}}

底层 \passthrough{\lstinline!unitree\_go::msg::dds\_::MotorState\_!} 的
\passthrough{\lstinline!mode!} 字段。类型签名只说明 Python
表示；单位、数值范围、数组固定长度和枚举含义需查对应 IDL 头文件。
取值范围为 0 到 255。

\textbf{签名}

\begin{lstlisting}[language=Python]
@property
def mode(self) -> int

@mode.setter
def mode(self, value: int) -> None
\end{lstlisting}

\textbf{可用性与安全性}

\textbf{\passthrough{\lstinline!AVAILABLE!} /
\passthrough{\lstinline!VALUE\_TYPE!}}：当前绑定源码已实现该消息值操作。

\textbf{参数}

\begin{longtable}[]{@{}
  >{\raggedright\arraybackslash}p{(\columnwidth - 4\tabcolsep) * \real{0.18}}
  >{\raggedright\arraybackslash}p{(\columnwidth - 4\tabcolsep) * \real{0.20}}
  >{\raggedright\arraybackslash}p{(\columnwidth - 4\tabcolsep) * \real{0.62}}@{}}
\toprule
名称 & Python 类型 & 含义 \\
\midrule
\endhead
\passthrough{\lstinline!value!} & \passthrough{\lstinline!int!} &
要写入或传递的新值，必须符合该参数的 Python 类型以及底层 C++
范围约束。 \\
\bottomrule
\end{longtable}

\textbf{返回值}

\begin{longtable}[]{@{}
  >{\raggedright\arraybackslash}p{(\columnwidth - 4\tabcolsep) * \real{0.18}}
  >{\raggedright\arraybackslash}p{(\columnwidth - 4\tabcolsep) * \real{0.20}}
  >{\raggedright\arraybackslash}p{(\columnwidth - 4\tabcolsep) * \real{0.62}}@{}}
\toprule
位置 & 类型 & 含义 \\
\midrule
\endhead
getter 返回值 & \passthrough{\lstinline!int!} & 返回属性当前值。 \\
\bottomrule
\end{longtable}

\textbf{对应 C++}

\begin{itemize}
\tightlist
\item
  类：\passthrough{\lstinline!unitree\_go::msg::dds\_::MotorState\_!}
\item
  字段类型：\passthrough{\lstinline!uint8\_t!}
\item
  声明位置：\passthrough{\lstinline!include/unitree/idl/go2/MotorState\_.hpp:24!}
\end{itemize}

\textbf{用法}

\begin{lstlisting}[language=Python]
current_value = obj.mode
obj.mode = new_value
\end{lstlisting}

\hypertarget{unitree-sdk2-cpp-idl-go2-motorstate-q}{%
\paragraph{\texorpdfstring{\texttt{unitree\_sdk2\_cpp.idl.go2.MotorState.q}}{unitree\_sdk2\_cpp.idl.go2.MotorState.q}}\label{unitree-sdk2-cpp-idl-go2-motorstate-q}}

底层 \passthrough{\lstinline!unitree\_go::msg::dds\_::MotorState\_!} 的
\passthrough{\lstinline!q!} 字段。类型签名只说明 Python
表示；单位、数值范围、数组固定长度和枚举含义需查对应 IDL 头文件。
底层为浮点数；类型本身不说明单位、坐标系或业务安全范围。

\textbf{签名}

\begin{lstlisting}[language=Python]
@property
def q(self) -> float

@q.setter
def q(self, value: float) -> None
\end{lstlisting}

\textbf{可用性与安全性}

\textbf{\passthrough{\lstinline!AVAILABLE!} /
\passthrough{\lstinline!VALUE\_TYPE!}}：当前绑定源码已实现该消息值操作。

\textbf{参数}

\begin{longtable}[]{@{}
  >{\raggedright\arraybackslash}p{(\columnwidth - 4\tabcolsep) * \real{0.18}}
  >{\raggedright\arraybackslash}p{(\columnwidth - 4\tabcolsep) * \real{0.20}}
  >{\raggedright\arraybackslash}p{(\columnwidth - 4\tabcolsep) * \real{0.62}}@{}}
\toprule
名称 & Python 类型 & 含义 \\
\midrule
\endhead
\passthrough{\lstinline!value!} & \passthrough{\lstinline!float!} &
要写入或传递的新值，必须符合该参数的 Python 类型以及底层 C++
范围约束。 \\
\bottomrule
\end{longtable}

\textbf{返回值}

\begin{longtable}[]{@{}
  >{\raggedright\arraybackslash}p{(\columnwidth - 4\tabcolsep) * \real{0.18}}
  >{\raggedright\arraybackslash}p{(\columnwidth - 4\tabcolsep) * \real{0.20}}
  >{\raggedright\arraybackslash}p{(\columnwidth - 4\tabcolsep) * \real{0.62}}@{}}
\toprule
位置 & 类型 & 含义 \\
\midrule
\endhead
getter 返回值 & \passthrough{\lstinline!float!} & 返回属性当前值。 \\
\bottomrule
\end{longtable}

\textbf{对应 C++}

\begin{itemize}
\tightlist
\item
  类：\passthrough{\lstinline!unitree\_go::msg::dds\_::MotorState\_!}
\item
  字段类型：\passthrough{\lstinline!float!}
\item
  声明位置：\passthrough{\lstinline!include/unitree/idl/go2/MotorState\_.hpp:25!}
\end{itemize}

\textbf{用法}

\begin{lstlisting}[language=Python]
current_value = obj.q
obj.q = new_value
\end{lstlisting}

\hypertarget{unitree-sdk2-cpp-idl-go2-motorstate-dq}{%
\paragraph{\texorpdfstring{\texttt{unitree\_sdk2\_cpp.idl.go2.MotorState.dq}}{unitree\_sdk2\_cpp.idl.go2.MotorState.dq}}\label{unitree-sdk2-cpp-idl-go2-motorstate-dq}}

底层 \passthrough{\lstinline!unitree\_go::msg::dds\_::MotorState\_!} 的
\passthrough{\lstinline!dq!} 字段。类型签名只说明 Python
表示；单位、数值范围、数组固定长度和枚举含义需查对应 IDL 头文件。
底层为浮点数；类型本身不说明单位、坐标系或业务安全范围。

\textbf{签名}

\begin{lstlisting}[language=Python]
@property
def dq(self) -> float

@dq.setter
def dq(self, value: float) -> None
\end{lstlisting}

\textbf{可用性与安全性}

\textbf{\passthrough{\lstinline!AVAILABLE!} /
\passthrough{\lstinline!VALUE\_TYPE!}}：当前绑定源码已实现该消息值操作。

\textbf{参数}

\begin{longtable}[]{@{}
  >{\raggedright\arraybackslash}p{(\columnwidth - 4\tabcolsep) * \real{0.18}}
  >{\raggedright\arraybackslash}p{(\columnwidth - 4\tabcolsep) * \real{0.20}}
  >{\raggedright\arraybackslash}p{(\columnwidth - 4\tabcolsep) * \real{0.62}}@{}}
\toprule
名称 & Python 类型 & 含义 \\
\midrule
\endhead
\passthrough{\lstinline!value!} & \passthrough{\lstinline!float!} &
要写入或传递的新值，必须符合该参数的 Python 类型以及底层 C++
范围约束。 \\
\bottomrule
\end{longtable}

\textbf{返回值}

\begin{longtable}[]{@{}
  >{\raggedright\arraybackslash}p{(\columnwidth - 4\tabcolsep) * \real{0.18}}
  >{\raggedright\arraybackslash}p{(\columnwidth - 4\tabcolsep) * \real{0.20}}
  >{\raggedright\arraybackslash}p{(\columnwidth - 4\tabcolsep) * \real{0.62}}@{}}
\toprule
位置 & 类型 & 含义 \\
\midrule
\endhead
getter 返回值 & \passthrough{\lstinline!float!} & 返回属性当前值。 \\
\bottomrule
\end{longtable}

\textbf{对应 C++}

\begin{itemize}
\tightlist
\item
  类：\passthrough{\lstinline!unitree\_go::msg::dds\_::MotorState\_!}
\item
  字段类型：\passthrough{\lstinline!float!}
\item
  声明位置：\passthrough{\lstinline!include/unitree/idl/go2/MotorState\_.hpp:26!}
\end{itemize}

\textbf{用法}

\begin{lstlisting}[language=Python]
current_value = obj.dq
obj.dq = new_value
\end{lstlisting}

\hypertarget{unitree-sdk2-cpp-idl-go2-motorstate-ddq}{%
\paragraph{\texorpdfstring{\texttt{unitree\_sdk2\_cpp.idl.go2.MotorState.ddq}}{unitree\_sdk2\_cpp.idl.go2.MotorState.ddq}}\label{unitree-sdk2-cpp-idl-go2-motorstate-ddq}}

底层 \passthrough{\lstinline!unitree\_go::msg::dds\_::MotorState\_!} 的
\passthrough{\lstinline!ddq!} 字段。类型签名只说明 Python
表示；单位、数值范围、数组固定长度和枚举含义需查对应 IDL 头文件。
底层为浮点数；类型本身不说明单位、坐标系或业务安全范围。

\textbf{签名}

\begin{lstlisting}[language=Python]
@property
def ddq(self) -> float

@ddq.setter
def ddq(self, value: float) -> None
\end{lstlisting}

\textbf{可用性与安全性}

\textbf{\passthrough{\lstinline!AVAILABLE!} /
\passthrough{\lstinline!VALUE\_TYPE!}}：当前绑定源码已实现该消息值操作。

\textbf{参数}

\begin{longtable}[]{@{}
  >{\raggedright\arraybackslash}p{(\columnwidth - 4\tabcolsep) * \real{0.18}}
  >{\raggedright\arraybackslash}p{(\columnwidth - 4\tabcolsep) * \real{0.20}}
  >{\raggedright\arraybackslash}p{(\columnwidth - 4\tabcolsep) * \real{0.62}}@{}}
\toprule
名称 & Python 类型 & 含义 \\
\midrule
\endhead
\passthrough{\lstinline!value!} & \passthrough{\lstinline!float!} &
要写入或传递的新值，必须符合该参数的 Python 类型以及底层 C++
范围约束。 \\
\bottomrule
\end{longtable}

\textbf{返回值}

\begin{longtable}[]{@{}
  >{\raggedright\arraybackslash}p{(\columnwidth - 4\tabcolsep) * \real{0.18}}
  >{\raggedright\arraybackslash}p{(\columnwidth - 4\tabcolsep) * \real{0.20}}
  >{\raggedright\arraybackslash}p{(\columnwidth - 4\tabcolsep) * \real{0.62}}@{}}
\toprule
位置 & 类型 & 含义 \\
\midrule
\endhead
getter 返回值 & \passthrough{\lstinline!float!} & 返回属性当前值。 \\
\bottomrule
\end{longtable}

\textbf{对应 C++}

\begin{itemize}
\tightlist
\item
  类：\passthrough{\lstinline!unitree\_go::msg::dds\_::MotorState\_!}
\item
  字段类型：\passthrough{\lstinline!float!}
\item
  声明位置：\passthrough{\lstinline!include/unitree/idl/go2/MotorState\_.hpp:27!}
\end{itemize}

\textbf{用法}

\begin{lstlisting}[language=Python]
current_value = obj.ddq
obj.ddq = new_value
\end{lstlisting}

\hypertarget{unitree-sdk2-cpp-idl-go2-motorstate-tau-est}{%
\paragraph{\texorpdfstring{\texttt{unitree\_sdk2\_cpp.idl.go2.MotorState.tau\_est}}{unitree\_sdk2\_cpp.idl.go2.MotorState.tau\_est}}\label{unitree-sdk2-cpp-idl-go2-motorstate-tau-est}}

底层 \passthrough{\lstinline!unitree\_go::msg::dds\_::MotorState\_!} 的
\passthrough{\lstinline!tau\_est!} 字段。类型签名只说明 Python
表示；单位、数值范围、数组固定长度和枚举含义需查对应 IDL 头文件。
底层为浮点数；类型本身不说明单位、坐标系或业务安全范围。

\textbf{签名}

\begin{lstlisting}[language=Python]
@property
def tau_est(self) -> float

@tau_est.setter
def tau_est(self, value: float) -> None
\end{lstlisting}

\textbf{可用性与安全性}

\textbf{\passthrough{\lstinline!AVAILABLE!} /
\passthrough{\lstinline!VALUE\_TYPE!}}：当前绑定源码已实现该消息值操作。

\textbf{参数}

\begin{longtable}[]{@{}
  >{\raggedright\arraybackslash}p{(\columnwidth - 4\tabcolsep) * \real{0.18}}
  >{\raggedright\arraybackslash}p{(\columnwidth - 4\tabcolsep) * \real{0.20}}
  >{\raggedright\arraybackslash}p{(\columnwidth - 4\tabcolsep) * \real{0.62}}@{}}
\toprule
名称 & Python 类型 & 含义 \\
\midrule
\endhead
\passthrough{\lstinline!value!} & \passthrough{\lstinline!float!} &
要写入或传递的新值，必须符合该参数的 Python 类型以及底层 C++
范围约束。 \\
\bottomrule
\end{longtable}

\textbf{返回值}

\begin{longtable}[]{@{}
  >{\raggedright\arraybackslash}p{(\columnwidth - 4\tabcolsep) * \real{0.18}}
  >{\raggedright\arraybackslash}p{(\columnwidth - 4\tabcolsep) * \real{0.20}}
  >{\raggedright\arraybackslash}p{(\columnwidth - 4\tabcolsep) * \real{0.62}}@{}}
\toprule
位置 & 类型 & 含义 \\
\midrule
\endhead
getter 返回值 & \passthrough{\lstinline!float!} & 返回属性当前值。 \\
\bottomrule
\end{longtable}

\textbf{对应 C++}

\begin{itemize}
\tightlist
\item
  类：\passthrough{\lstinline!unitree\_go::msg::dds\_::MotorState\_!}
\item
  字段类型：\passthrough{\lstinline!float!}
\item
  声明位置：\passthrough{\lstinline!include/unitree/idl/go2/MotorState\_.hpp:28!}
\end{itemize}

\textbf{用法}

\begin{lstlisting}[language=Python]
current_value = obj.tau_est
obj.tau_est = new_value
\end{lstlisting}

\hypertarget{unitree-sdk2-cpp-idl-go2-motorstate-q-raw}{%
\paragraph{\texorpdfstring{\texttt{unitree\_sdk2\_cpp.idl.go2.MotorState.q\_raw}}{unitree\_sdk2\_cpp.idl.go2.MotorState.q\_raw}}\label{unitree-sdk2-cpp-idl-go2-motorstate-q-raw}}

底层 \passthrough{\lstinline!unitree\_go::msg::dds\_::MotorState\_!} 的
\passthrough{\lstinline!q\_raw!} 字段。类型签名只说明 Python
表示；单位、数值范围、数组固定长度和枚举含义需查对应 IDL 头文件。
底层为浮点数；类型本身不说明单位、坐标系或业务安全范围。

\textbf{签名}

\begin{lstlisting}[language=Python]
@property
def q_raw(self) -> float

@q_raw.setter
def q_raw(self, value: float) -> None
\end{lstlisting}

\textbf{可用性与安全性}

\textbf{\passthrough{\lstinline!AVAILABLE!} /
\passthrough{\lstinline!VALUE\_TYPE!}}：当前绑定源码已实现该消息值操作。

\textbf{参数}

\begin{longtable}[]{@{}
  >{\raggedright\arraybackslash}p{(\columnwidth - 4\tabcolsep) * \real{0.18}}
  >{\raggedright\arraybackslash}p{(\columnwidth - 4\tabcolsep) * \real{0.20}}
  >{\raggedright\arraybackslash}p{(\columnwidth - 4\tabcolsep) * \real{0.62}}@{}}
\toprule
名称 & Python 类型 & 含义 \\
\midrule
\endhead
\passthrough{\lstinline!value!} & \passthrough{\lstinline!float!} &
要写入或传递的新值，必须符合该参数的 Python 类型以及底层 C++
范围约束。 \\
\bottomrule
\end{longtable}

\textbf{返回值}

\begin{longtable}[]{@{}
  >{\raggedright\arraybackslash}p{(\columnwidth - 4\tabcolsep) * \real{0.18}}
  >{\raggedright\arraybackslash}p{(\columnwidth - 4\tabcolsep) * \real{0.20}}
  >{\raggedright\arraybackslash}p{(\columnwidth - 4\tabcolsep) * \real{0.62}}@{}}
\toprule
位置 & 类型 & 含义 \\
\midrule
\endhead
getter 返回值 & \passthrough{\lstinline!float!} & 返回属性当前值。 \\
\bottomrule
\end{longtable}

\textbf{对应 C++}

\begin{itemize}
\tightlist
\item
  类：\passthrough{\lstinline!unitree\_go::msg::dds\_::MotorState\_!}
\item
  字段类型：\passthrough{\lstinline!float!}
\item
  声明位置：\passthrough{\lstinline!include/unitree/idl/go2/MotorState\_.hpp:29!}
\end{itemize}

\textbf{用法}

\begin{lstlisting}[language=Python]
current_value = obj.q_raw
obj.q_raw = new_value
\end{lstlisting}

\hypertarget{unitree-sdk2-cpp-idl-go2-motorstate-dq-raw}{%
\paragraph{\texorpdfstring{\texttt{unitree\_sdk2\_cpp.idl.go2.MotorState.dq\_raw}}{unitree\_sdk2\_cpp.idl.go2.MotorState.dq\_raw}}\label{unitree-sdk2-cpp-idl-go2-motorstate-dq-raw}}

底层 \passthrough{\lstinline!unitree\_go::msg::dds\_::MotorState\_!} 的
\passthrough{\lstinline!dq\_raw!} 字段。类型签名只说明 Python
表示；单位、数值范围、数组固定长度和枚举含义需查对应 IDL 头文件。
底层为浮点数；类型本身不说明单位、坐标系或业务安全范围。

\textbf{签名}

\begin{lstlisting}[language=Python]
@property
def dq_raw(self) -> float

@dq_raw.setter
def dq_raw(self, value: float) -> None
\end{lstlisting}

\textbf{可用性与安全性}

\textbf{\passthrough{\lstinline!AVAILABLE!} /
\passthrough{\lstinline!VALUE\_TYPE!}}：当前绑定源码已实现该消息值操作。

\textbf{参数}

\begin{longtable}[]{@{}
  >{\raggedright\arraybackslash}p{(\columnwidth - 4\tabcolsep) * \real{0.18}}
  >{\raggedright\arraybackslash}p{(\columnwidth - 4\tabcolsep) * \real{0.20}}
  >{\raggedright\arraybackslash}p{(\columnwidth - 4\tabcolsep) * \real{0.62}}@{}}
\toprule
名称 & Python 类型 & 含义 \\
\midrule
\endhead
\passthrough{\lstinline!value!} & \passthrough{\lstinline!float!} &
要写入或传递的新值，必须符合该参数的 Python 类型以及底层 C++
范围约束。 \\
\bottomrule
\end{longtable}

\textbf{返回值}

\begin{longtable}[]{@{}
  >{\raggedright\arraybackslash}p{(\columnwidth - 4\tabcolsep) * \real{0.18}}
  >{\raggedright\arraybackslash}p{(\columnwidth - 4\tabcolsep) * \real{0.20}}
  >{\raggedright\arraybackslash}p{(\columnwidth - 4\tabcolsep) * \real{0.62}}@{}}
\toprule
位置 & 类型 & 含义 \\
\midrule
\endhead
getter 返回值 & \passthrough{\lstinline!float!} & 返回属性当前值。 \\
\bottomrule
\end{longtable}

\textbf{对应 C++}

\begin{itemize}
\tightlist
\item
  类：\passthrough{\lstinline!unitree\_go::msg::dds\_::MotorState\_!}
\item
  字段类型：\passthrough{\lstinline!float!}
\item
  声明位置：\passthrough{\lstinline!include/unitree/idl/go2/MotorState\_.hpp:30!}
\end{itemize}

\textbf{用法}

\begin{lstlisting}[language=Python]
current_value = obj.dq_raw
obj.dq_raw = new_value
\end{lstlisting}

\hypertarget{unitree-sdk2-cpp-idl-go2-motorstate-ddq-raw}{%
\paragraph{\texorpdfstring{\texttt{unitree\_sdk2\_cpp.idl.go2.MotorState.ddq\_raw}}{unitree\_sdk2\_cpp.idl.go2.MotorState.ddq\_raw}}\label{unitree-sdk2-cpp-idl-go2-motorstate-ddq-raw}}

底层 \passthrough{\lstinline!unitree\_go::msg::dds\_::MotorState\_!} 的
\passthrough{\lstinline!ddq\_raw!} 字段。类型签名只说明 Python
表示；单位、数值范围、数组固定长度和枚举含义需查对应 IDL 头文件。
底层为浮点数；类型本身不说明单位、坐标系或业务安全范围。

\textbf{签名}

\begin{lstlisting}[language=Python]
@property
def ddq_raw(self) -> float

@ddq_raw.setter
def ddq_raw(self, value: float) -> None
\end{lstlisting}

\textbf{可用性与安全性}

\textbf{\passthrough{\lstinline!AVAILABLE!} /
\passthrough{\lstinline!VALUE\_TYPE!}}：当前绑定源码已实现该消息值操作。

\textbf{参数}

\begin{longtable}[]{@{}
  >{\raggedright\arraybackslash}p{(\columnwidth - 4\tabcolsep) * \real{0.18}}
  >{\raggedright\arraybackslash}p{(\columnwidth - 4\tabcolsep) * \real{0.20}}
  >{\raggedright\arraybackslash}p{(\columnwidth - 4\tabcolsep) * \real{0.62}}@{}}
\toprule
名称 & Python 类型 & 含义 \\
\midrule
\endhead
\passthrough{\lstinline!value!} & \passthrough{\lstinline!float!} &
要写入或传递的新值，必须符合该参数的 Python 类型以及底层 C++
范围约束。 \\
\bottomrule
\end{longtable}

\textbf{返回值}

\begin{longtable}[]{@{}
  >{\raggedright\arraybackslash}p{(\columnwidth - 4\tabcolsep) * \real{0.18}}
  >{\raggedright\arraybackslash}p{(\columnwidth - 4\tabcolsep) * \real{0.20}}
  >{\raggedright\arraybackslash}p{(\columnwidth - 4\tabcolsep) * \real{0.62}}@{}}
\toprule
位置 & 类型 & 含义 \\
\midrule
\endhead
getter 返回值 & \passthrough{\lstinline!float!} & 返回属性当前值。 \\
\bottomrule
\end{longtable}

\textbf{对应 C++}

\begin{itemize}
\tightlist
\item
  类：\passthrough{\lstinline!unitree\_go::msg::dds\_::MotorState\_!}
\item
  字段类型：\passthrough{\lstinline!float!}
\item
  声明位置：\passthrough{\lstinline!include/unitree/idl/go2/MotorState\_.hpp:31!}
\end{itemize}

\textbf{用法}

\begin{lstlisting}[language=Python]
current_value = obj.ddq_raw
obj.ddq_raw = new_value
\end{lstlisting}

\hypertarget{unitree-sdk2-cpp-idl-go2-motorstate-temperature}{%
\paragraph{\texorpdfstring{\texttt{unitree\_sdk2\_cpp.idl.go2.MotorState.temperature}}{unitree\_sdk2\_cpp.idl.go2.MotorState.temperature}}\label{unitree-sdk2-cpp-idl-go2-motorstate-temperature}}

底层 \passthrough{\lstinline!unitree\_go::msg::dds\_::MotorState\_!} 的
\passthrough{\lstinline!temperature!} 字段。类型签名只说明 Python
表示；单位、数值范围、数组固定长度和枚举含义需查对应 IDL 头文件。
取值范围为 0 到 255。

\textbf{签名}

\begin{lstlisting}[language=Python]
@property
def temperature(self) -> int

@temperature.setter
def temperature(self, value: int) -> None
\end{lstlisting}

\textbf{可用性与安全性}

\textbf{\passthrough{\lstinline!AVAILABLE!} /
\passthrough{\lstinline!VALUE\_TYPE!}}：当前绑定源码已实现该消息值操作。

\textbf{参数}

\begin{longtable}[]{@{}
  >{\raggedright\arraybackslash}p{(\columnwidth - 4\tabcolsep) * \real{0.18}}
  >{\raggedright\arraybackslash}p{(\columnwidth - 4\tabcolsep) * \real{0.20}}
  >{\raggedright\arraybackslash}p{(\columnwidth - 4\tabcolsep) * \real{0.62}}@{}}
\toprule
名称 & Python 类型 & 含义 \\
\midrule
\endhead
\passthrough{\lstinline!value!} & \passthrough{\lstinline!int!} &
要写入或传递的新值，必须符合该参数的 Python 类型以及底层 C++
范围约束。 \\
\bottomrule
\end{longtable}

\textbf{返回值}

\begin{longtable}[]{@{}
  >{\raggedright\arraybackslash}p{(\columnwidth - 4\tabcolsep) * \real{0.18}}
  >{\raggedright\arraybackslash}p{(\columnwidth - 4\tabcolsep) * \real{0.20}}
  >{\raggedright\arraybackslash}p{(\columnwidth - 4\tabcolsep) * \real{0.62}}@{}}
\toprule
位置 & 类型 & 含义 \\
\midrule
\endhead
getter 返回值 & \passthrough{\lstinline!int!} & 返回属性当前值。 \\
\bottomrule
\end{longtable}

\textbf{对应 C++}

\begin{itemize}
\tightlist
\item
  类：\passthrough{\lstinline!unitree\_go::msg::dds\_::MotorState\_!}
\item
  字段类型：\passthrough{\lstinline!uint8\_t!}
\item
  声明位置：\passthrough{\lstinline!include/unitree/idl/go2/MotorState\_.hpp:32!}
\end{itemize}

\textbf{用法}

\begin{lstlisting}[language=Python]
current_value = obj.temperature
obj.temperature = new_value
\end{lstlisting}

\hypertarget{unitree-sdk2-cpp-idl-go2-motorstate-lost}{%
\paragraph{\texorpdfstring{\texttt{unitree\_sdk2\_cpp.idl.go2.MotorState.lost}}{unitree\_sdk2\_cpp.idl.go2.MotorState.lost}}\label{unitree-sdk2-cpp-idl-go2-motorstate-lost}}

底层 \passthrough{\lstinline!unitree\_go::msg::dds\_::MotorState\_!} 的
\passthrough{\lstinline!lost!} 字段。类型签名只说明 Python
表示；单位、数值范围、数组固定长度和枚举含义需查对应 IDL 头文件。
取值范围为 0 到 4294967295。

\textbf{签名}

\begin{lstlisting}[language=Python]
@property
def lost(self) -> int

@lost.setter
def lost(self, value: int) -> None
\end{lstlisting}

\textbf{可用性与安全性}

\textbf{\passthrough{\lstinline!AVAILABLE!} /
\passthrough{\lstinline!VALUE\_TYPE!}}：当前绑定源码已实现该消息值操作。

\textbf{参数}

\begin{longtable}[]{@{}
  >{\raggedright\arraybackslash}p{(\columnwidth - 4\tabcolsep) * \real{0.18}}
  >{\raggedright\arraybackslash}p{(\columnwidth - 4\tabcolsep) * \real{0.20}}
  >{\raggedright\arraybackslash}p{(\columnwidth - 4\tabcolsep) * \real{0.62}}@{}}
\toprule
名称 & Python 类型 & 含义 \\
\midrule
\endhead
\passthrough{\lstinline!value!} & \passthrough{\lstinline!int!} &
要写入或传递的新值，必须符合该参数的 Python 类型以及底层 C++
范围约束。 \\
\bottomrule
\end{longtable}

\textbf{返回值}

\begin{longtable}[]{@{}
  >{\raggedright\arraybackslash}p{(\columnwidth - 4\tabcolsep) * \real{0.18}}
  >{\raggedright\arraybackslash}p{(\columnwidth - 4\tabcolsep) * \real{0.20}}
  >{\raggedright\arraybackslash}p{(\columnwidth - 4\tabcolsep) * \real{0.62}}@{}}
\toprule
位置 & 类型 & 含义 \\
\midrule
\endhead
getter 返回值 & \passthrough{\lstinline!int!} & 返回属性当前值。 \\
\bottomrule
\end{longtable}

\textbf{对应 C++}

\begin{itemize}
\tightlist
\item
  类：\passthrough{\lstinline!unitree\_go::msg::dds\_::MotorState\_!}
\item
  字段类型：\passthrough{\lstinline!uint32\_t!}
\item
  声明位置：\passthrough{\lstinline!include/unitree/idl/go2/MotorState\_.hpp:33!}
\end{itemize}

\textbf{用法}

\begin{lstlisting}[language=Python]
current_value = obj.lost
obj.lost = new_value
\end{lstlisting}

\hypertarget{unitree-sdk2-cpp-idl-go2-motorstate-reserve}{%
\paragraph{\texorpdfstring{\texttt{unitree\_sdk2\_cpp.idl.go2.MotorState.reserve}}{unitree\_sdk2\_cpp.idl.go2.MotorState.reserve}}\label{unitree-sdk2-cpp-idl-go2-motorstate-reserve}}

底层 \passthrough{\lstinline!unitree\_go::msg::dds\_::MotorState\_!} 的
\passthrough{\lstinline!reserve!} 字段。类型签名只说明 Python
表示；单位、数值范围、数组固定长度和枚举含义需查对应 IDL 头文件。
底层是固定长度数组，写入序列必须正好包含 2 个元素；元素约束：取值范围为
0 到 4294967295。

\textbf{签名}

\begin{lstlisting}[language=Python]
@property
def reserve(self) -> list[int]

@reserve.setter
def reserve(self, value: Sequence[int]) -> None
\end{lstlisting}

\textbf{可用性与安全性}

\textbf{\passthrough{\lstinline!AVAILABLE!} /
\passthrough{\lstinline!VALUE\_TYPE!}}：当前绑定源码已实现该消息值操作。

\textbf{参数}

\begin{longtable}[]{@{}
  >{\raggedright\arraybackslash}p{(\columnwidth - 4\tabcolsep) * \real{0.18}}
  >{\raggedright\arraybackslash}p{(\columnwidth - 4\tabcolsep) * \real{0.20}}
  >{\raggedright\arraybackslash}p{(\columnwidth - 4\tabcolsep) * \real{0.62}}@{}}
\toprule
名称 & Python 类型 & 含义 \\
\midrule
\endhead
\passthrough{\lstinline!value!} &
\passthrough{\lstinline!Sequence[int]!} &
要写入或传递的新值，必须符合该参数的 Python 类型以及底层 C++
范围约束。 \\
\bottomrule
\end{longtable}

\textbf{返回值}

\begin{longtable}[]{@{}
  >{\raggedright\arraybackslash}p{(\columnwidth - 4\tabcolsep) * \real{0.18}}
  >{\raggedright\arraybackslash}p{(\columnwidth - 4\tabcolsep) * \real{0.20}}
  >{\raggedright\arraybackslash}p{(\columnwidth - 4\tabcolsep) * \real{0.62}}@{}}
\toprule
位置 & 类型 & 含义 \\
\midrule
\endhead
getter 返回值 & \passthrough{\lstinline!list[int]!} &
返回属性当前值。数组、vector 和嵌套消息采用复制语义。 \\
\bottomrule
\end{longtable}

\textbf{对应 C++}

\begin{itemize}
\tightlist
\item
  类：\passthrough{\lstinline!unitree\_go::msg::dds\_::MotorState\_!}
\item
  字段类型：\passthrough{\lstinline!std::array<uint32\_t, 2>!}
\item
  声明位置：\passthrough{\lstinline!include/unitree/idl/go2/MotorState\_.hpp:34!}
\end{itemize}

\textbf{用法}

\begin{lstlisting}[language=Python]
current_value = obj.reserve
obj.reserve = new_value
\end{lstlisting}

\begin{quote}
{[}!TIP{]} 读取容器或嵌套值后应遵循``读取、修改、重新赋值''。只修改
getter 返回的副本不会可靠地更新原 C++ 消息。
\end{quote}

\hypertarget{unitree-sdk2-cpp-idl-go2-lowstate}{%
\subsubsection{\texorpdfstring{\texttt{unitree\_sdk2\_cpp.idl.go2.LowState}}{unitree\_sdk2\_cpp.idl.go2.LowState}}\label{unitree-sdk2-cpp-idl-go2-lowstate}}

可由 Python 读写的 DDS/IDL 消息值类型。字段采用复制语义。

\textbf{导入}

\begin{lstlisting}[language=Python]
from unitree_sdk2_cpp.idl.go2 import LowState
\end{lstlisting}

\textbf{方法索引}

\begin{longtable}[]{@{}
  >{\raggedright\arraybackslash}p{(\columnwidth - 6\tabcolsep) * \real{0.18}}
  >{\raggedright\arraybackslash}p{(\columnwidth - 6\tabcolsep) * \real{0.24}}
  >{\raggedright\arraybackslash}p{(\columnwidth - 6\tabcolsep) * \real{0.18}}
  >{\raggedright\arraybackslash}p{(\columnwidth - 6\tabcolsep) * \real{0.40}}@{}}
\toprule
方法 & Python 签名 & 状态 & 安全分类 \\
\midrule
\endhead
\protect\hyperlink{unitree-sdk2-cpp-idl-go2-lowstate-dunder-init-1}{\passthrough{\lstinline!\_\_init\_\_!}}
& \passthrough{\lstinline!def \_\_init\_\_(self) -> None!} &
\passthrough{\lstinline!AVAILABLE!} &
\passthrough{\lstinline!VALUE\_TYPE!} \\
\protect\hyperlink{unitree-sdk2-cpp-idl-go2-lowstate-dunder-eq-1}{\passthrough{\lstinline!\_\_eq\_\_!}}
& \passthrough{\lstinline!def \_\_eq\_\_(self, other: object) -> bool!}
& \passthrough{\lstinline!AVAILABLE!} &
\passthrough{\lstinline!VALUE\_TYPE!} \\
\protect\hyperlink{unitree-sdk2-cpp-idl-go2-lowstate-dunder-ne-1}{\passthrough{\lstinline!\_\_ne\_\_!}}
& \passthrough{\lstinline!def \_\_ne\_\_(self, other: object) -> bool!}
& \passthrough{\lstinline!AVAILABLE!} &
\passthrough{\lstinline!VALUE\_TYPE!} \\
\bottomrule
\end{longtable}

\textbf{属性索引}

\begin{longtable}[]{@{}
  >{\raggedright\arraybackslash}p{(\columnwidth - 6\tabcolsep) * \real{0.18}}
  >{\raggedright\arraybackslash}p{(\columnwidth - 6\tabcolsep) * \real{0.24}}
  >{\raggedright\arraybackslash}p{(\columnwidth - 6\tabcolsep) * \real{0.18}}
  >{\raggedright\arraybackslash}p{(\columnwidth - 6\tabcolsep) * \real{0.40}}@{}}
\toprule
属性 & 读取类型 & 写入类型 & 状态 \\
\midrule
\endhead
\protect\hyperlink{unitree-sdk2-cpp-idl-go2-lowstate-head}{\passthrough{\lstinline!head!}}
& \passthrough{\lstinline!list[int]!} &
\passthrough{\lstinline!Sequence[int]!} &
\passthrough{\lstinline!AVAILABLE!} \\
\protect\hyperlink{unitree-sdk2-cpp-idl-go2-lowstate-level-flag}{\passthrough{\lstinline!level\_flag!}}
& \passthrough{\lstinline!int!} & \passthrough{\lstinline!int!} &
\passthrough{\lstinline!AVAILABLE!} \\
\protect\hyperlink{unitree-sdk2-cpp-idl-go2-lowstate-frame-reserve}{\passthrough{\lstinline!frame\_reserve!}}
& \passthrough{\lstinline!int!} & \passthrough{\lstinline!int!} &
\passthrough{\lstinline!AVAILABLE!} \\
\protect\hyperlink{unitree-sdk2-cpp-idl-go2-lowstate-sn}{\passthrough{\lstinline!sn!}}
& \passthrough{\lstinline!list[int]!} &
\passthrough{\lstinline!Sequence[int]!} &
\passthrough{\lstinline!AVAILABLE!} \\
\protect\hyperlink{unitree-sdk2-cpp-idl-go2-lowstate-version}{\passthrough{\lstinline!version!}}
& \passthrough{\lstinline!list[int]!} &
\passthrough{\lstinline!Sequence[int]!} &
\passthrough{\lstinline!AVAILABLE!} \\
\protect\hyperlink{unitree-sdk2-cpp-idl-go2-lowstate-bandwidth}{\passthrough{\lstinline!bandwidth!}}
& \passthrough{\lstinline!int!} & \passthrough{\lstinline!int!} &
\passthrough{\lstinline!AVAILABLE!} \\
\protect\hyperlink{unitree-sdk2-cpp-idl-go2-lowstate-imu-state}{\passthrough{\lstinline!imu\_state!}}
& \passthrough{\lstinline!IMUState!} &
\passthrough{\lstinline!IMUState!} &
\passthrough{\lstinline!AVAILABLE!} \\
\protect\hyperlink{unitree-sdk2-cpp-idl-go2-lowstate-motor-state}{\passthrough{\lstinline!motor\_state!}}
& \passthrough{\lstinline!list[MotorState]!} &
\passthrough{\lstinline!Sequence[MotorState]!} &
\passthrough{\lstinline!AVAILABLE!} \\
\protect\hyperlink{unitree-sdk2-cpp-idl-go2-lowstate-bms-state}{\passthrough{\lstinline!bms\_state!}}
& \passthrough{\lstinline!BmsState!} &
\passthrough{\lstinline!BmsState!} &
\passthrough{\lstinline!AVAILABLE!} \\
\protect\hyperlink{unitree-sdk2-cpp-idl-go2-lowstate-foot-force}{\passthrough{\lstinline!foot\_force!}}
& \passthrough{\lstinline!list[int]!} &
\passthrough{\lstinline!Sequence[int]!} &
\passthrough{\lstinline!AVAILABLE!} \\
\protect\hyperlink{unitree-sdk2-cpp-idl-go2-lowstate-foot-force-est}{\passthrough{\lstinline!foot\_force\_est!}}
& \passthrough{\lstinline!list[int]!} &
\passthrough{\lstinline!Sequence[int]!} &
\passthrough{\lstinline!AVAILABLE!} \\
\protect\hyperlink{unitree-sdk2-cpp-idl-go2-lowstate-tick}{\passthrough{\lstinline!tick!}}
& \passthrough{\lstinline!int!} & \passthrough{\lstinline!int!} &
\passthrough{\lstinline!AVAILABLE!} \\
\protect\hyperlink{unitree-sdk2-cpp-idl-go2-lowstate-wireless-remote}{\passthrough{\lstinline!wireless\_remote!}}
& \passthrough{\lstinline!list[int]!} &
\passthrough{\lstinline!Sequence[int]!} &
\passthrough{\lstinline!AVAILABLE!} \\
\protect\hyperlink{unitree-sdk2-cpp-idl-go2-lowstate-bit-flag}{\passthrough{\lstinline!bit\_flag!}}
& \passthrough{\lstinline!int!} & \passthrough{\lstinline!int!} &
\passthrough{\lstinline!AVAILABLE!} \\
\protect\hyperlink{unitree-sdk2-cpp-idl-go2-lowstate-adc-reel}{\passthrough{\lstinline!adc\_reel!}}
& \passthrough{\lstinline!float!} & \passthrough{\lstinline!float!} &
\passthrough{\lstinline!AVAILABLE!} \\
\protect\hyperlink{unitree-sdk2-cpp-idl-go2-lowstate-temperature-ntc1}{\passthrough{\lstinline!temperature\_ntc1!}}
& \passthrough{\lstinline!int!} & \passthrough{\lstinline!int!} &
\passthrough{\lstinline!AVAILABLE!} \\
\protect\hyperlink{unitree-sdk2-cpp-idl-go2-lowstate-temperature-ntc2}{\passthrough{\lstinline!temperature\_ntc2!}}
& \passthrough{\lstinline!int!} & \passthrough{\lstinline!int!} &
\passthrough{\lstinline!AVAILABLE!} \\
\protect\hyperlink{unitree-sdk2-cpp-idl-go2-lowstate-power-v}{\passthrough{\lstinline!power\_v!}}
& \passthrough{\lstinline!float!} & \passthrough{\lstinline!float!} &
\passthrough{\lstinline!AVAILABLE!} \\
\protect\hyperlink{unitree-sdk2-cpp-idl-go2-lowstate-power-a}{\passthrough{\lstinline!power\_a!}}
& \passthrough{\lstinline!float!} & \passthrough{\lstinline!float!} &
\passthrough{\lstinline!AVAILABLE!} \\
\protect\hyperlink{unitree-sdk2-cpp-idl-go2-lowstate-fan-frequency}{\passthrough{\lstinline!fan\_frequency!}}
& \passthrough{\lstinline!list[int]!} &
\passthrough{\lstinline!Sequence[int]!} &
\passthrough{\lstinline!AVAILABLE!} \\
\protect\hyperlink{unitree-sdk2-cpp-idl-go2-lowstate-reserve}{\passthrough{\lstinline!reserve!}}
& \passthrough{\lstinline!int!} & \passthrough{\lstinline!int!} &
\passthrough{\lstinline!AVAILABLE!} \\
\protect\hyperlink{unitree-sdk2-cpp-idl-go2-lowstate-crc}{\passthrough{\lstinline!crc!}}
& \passthrough{\lstinline!int!} & \passthrough{\lstinline!int!} &
\passthrough{\lstinline!AVAILABLE!} \\
\bottomrule
\end{longtable}

\hypertarget{unitree-sdk2-cpp-idl-go2-lowstate-dunder-init-1}{%
\paragraph{\texorpdfstring{\texttt{unitree\_sdk2\_cpp.idl.go2.LowState.\_\_init\_\_}}{unitree\_sdk2\_cpp.idl.go2.LowState.\_\_init\_\_}}\label{unitree-sdk2-cpp-idl-go2-lowstate-dunder-init-1}}

初始化 \passthrough{\lstinline!LowState!}
实例。是否能实际构造取决于下方可用性状态。

\textbf{签名}

\begin{lstlisting}[language=Python]
def __init__(self) -> None
\end{lstlisting}

\textbf{可用性与安全性}

\textbf{\passthrough{\lstinline!AVAILABLE!} /
\passthrough{\lstinline!VALUE\_TYPE!}}：当前绑定源码已实现该消息值操作。

\textbf{参数}

无显式参数。实例方法中的 \passthrough{\lstinline!self!} 由 Python
自动传入。

\textbf{返回值}

\begin{longtable}[]{@{}
  >{\raggedright\arraybackslash}p{(\columnwidth - 4\tabcolsep) * \real{0.18}}
  >{\raggedright\arraybackslash}p{(\columnwidth - 4\tabcolsep) * \real{0.20}}
  >{\raggedright\arraybackslash}p{(\columnwidth - 4\tabcolsep) * \real{0.62}}@{}}
\toprule
位置 & 类型 & 含义 \\
\midrule
\endhead
无 & \passthrough{\lstinline!None!} &
\passthrough{\lstinline!\_\_init\_\_!}
本身不返回值；调用类对象时，在构造函数可用的前提下得到该类实例。 \\
\bottomrule
\end{longtable}

\textbf{对应 C++}

\begin{itemize}
\tightlist
\item
  类：\passthrough{\lstinline!unitree\_go::msg::dds\_::LowState\_!}
\item
  签名：\passthrough{\lstinline!LowState\_()!}
\item
  绑定策略：\passthrough{\lstinline!未单独标注!}
\item
  声明位置：\passthrough{\lstinline!include/unitree/idl/go2/LowState\_.hpp:54!}
\end{itemize}

\textbf{说明}

\begin{itemize}
\tightlist
\item
  示例是局部调用片段；\passthrough{\lstinline!obj!}
  和参数变量表示已经按上层业务校验并准备好的对象。
\end{itemize}

\textbf{用法}

\begin{lstlisting}[language=Python]
value = LowState()
\end{lstlisting}

\hypertarget{unitree-sdk2-cpp-idl-go2-lowstate-dunder-eq-1}{%
\paragraph{\texorpdfstring{\texttt{unitree\_sdk2\_cpp.idl.go2.LowState.\_\_eq\_\_}}{unitree\_sdk2\_cpp.idl.go2.LowState.\_\_eq\_\_}}\label{unitree-sdk2-cpp-idl-go2-lowstate-dunder-eq-1}}

按底层 IDL 消息内容比较两个对象是否相等。

\textbf{签名}

\begin{lstlisting}[language=Python]
def __eq__(self, other: object) -> bool
\end{lstlisting}

\textbf{可用性与安全性}

\textbf{\passthrough{\lstinline!AVAILABLE!} /
\passthrough{\lstinline!VALUE\_TYPE!}}：当前绑定源码已实现该消息值操作。

\textbf{参数}

\begin{longtable}[]{@{}
  >{\raggedright\arraybackslash}p{(\columnwidth - 6\tabcolsep) * \real{0.18}}
  >{\raggedright\arraybackslash}p{(\columnwidth - 6\tabcolsep) * \real{0.24}}
  >{\raggedright\arraybackslash}p{(\columnwidth - 6\tabcolsep) * \real{0.18}}
  >{\raggedright\arraybackslash}p{(\columnwidth - 6\tabcolsep) * \real{0.40}}@{}}
\toprule
名称 & Python 类型 & 默认值 & 含义 \\
\midrule
\endhead
\passthrough{\lstinline!other!} & \passthrough{\lstinline!object!} &
必填 & 用于比较的另一个 Python 对象。类型不兼容时比较结果通常为
\passthrough{\lstinline!False!}。 \\
\bottomrule
\end{longtable}

\textbf{返回值}

\begin{longtable}[]{@{}
  >{\raggedright\arraybackslash}p{(\columnwidth - 4\tabcolsep) * \real{0.18}}
  >{\raggedright\arraybackslash}p{(\columnwidth - 4\tabcolsep) * \real{0.20}}
  >{\raggedright\arraybackslash}p{(\columnwidth - 4\tabcolsep) * \real{0.62}}@{}}
\toprule
位置 & 类型 & 含义 \\
\midrule
\endhead
返回值 & \passthrough{\lstinline!bool!} & 返回
\passthrough{\lstinline!bool!}。更精确的业务含义见该方法用途、C++
签名和目标型号协议。 \\
\bottomrule
\end{longtable}

\textbf{对应 C++}

\begin{itemize}
\tightlist
\item
  类：\passthrough{\lstinline!unitree\_go::msg::dds\_::LowState\_!}
\item
  签名：\passthrough{\lstinline!operator==(const value \&) const!}
\item
  绑定策略：\passthrough{\lstinline!未单独标注!}
\item
  声明位置：由 pybind11 手工绑定或生成注册表提供；manifest
  未记录独立头文件行号。
\end{itemize}

\textbf{说明}

\begin{itemize}
\tightlist
\item
  示例是局部调用片段；\passthrough{\lstinline!obj!}
  和参数变量表示已经按上层业务校验并准备好的对象。
\end{itemize}

\textbf{用法}

\begin{lstlisting}[language=Python]
same = left == right
\end{lstlisting}

\hypertarget{unitree-sdk2-cpp-idl-go2-lowstate-dunder-ne-1}{%
\paragraph{\texorpdfstring{\texttt{unitree\_sdk2\_cpp.idl.go2.LowState.\_\_ne\_\_}}{unitree\_sdk2\_cpp.idl.go2.LowState.\_\_ne\_\_}}\label{unitree-sdk2-cpp-idl-go2-lowstate-dunder-ne-1}}

按底层 IDL 消息内容比较两个对象是否不相等。

\textbf{签名}

\begin{lstlisting}[language=Python]
def __ne__(self, other: object) -> bool
\end{lstlisting}

\textbf{可用性与安全性}

\textbf{\passthrough{\lstinline!AVAILABLE!} /
\passthrough{\lstinline!VALUE\_TYPE!}}：当前绑定源码已实现该消息值操作。

\textbf{参数}

\begin{longtable}[]{@{}
  >{\raggedright\arraybackslash}p{(\columnwidth - 6\tabcolsep) * \real{0.18}}
  >{\raggedright\arraybackslash}p{(\columnwidth - 6\tabcolsep) * \real{0.24}}
  >{\raggedright\arraybackslash}p{(\columnwidth - 6\tabcolsep) * \real{0.18}}
  >{\raggedright\arraybackslash}p{(\columnwidth - 6\tabcolsep) * \real{0.40}}@{}}
\toprule
名称 & Python 类型 & 默认值 & 含义 \\
\midrule
\endhead
\passthrough{\lstinline!other!} & \passthrough{\lstinline!object!} &
必填 & 用于比较的另一个 Python 对象。类型不兼容时比较结果通常为
\passthrough{\lstinline!False!}。 \\
\bottomrule
\end{longtable}

\textbf{返回值}

\begin{longtable}[]{@{}
  >{\raggedright\arraybackslash}p{(\columnwidth - 4\tabcolsep) * \real{0.18}}
  >{\raggedright\arraybackslash}p{(\columnwidth - 4\tabcolsep) * \real{0.20}}
  >{\raggedright\arraybackslash}p{(\columnwidth - 4\tabcolsep) * \real{0.62}}@{}}
\toprule
位置 & 类型 & 含义 \\
\midrule
\endhead
返回值 & \passthrough{\lstinline!bool!} & 返回
\passthrough{\lstinline!bool!}。更精确的业务含义见该方法用途、C++
签名和目标型号协议。 \\
\bottomrule
\end{longtable}

\textbf{对应 C++}

\begin{itemize}
\tightlist
\item
  类：\passthrough{\lstinline!unitree\_go::msg::dds\_::LowState\_!}
\item
  签名：\passthrough{\lstinline"operator!=(const value \&) const"}
\item
  绑定策略：\passthrough{\lstinline!未单独标注!}
\item
  声明位置：由 pybind11 手工绑定或生成注册表提供；manifest
  未记录独立头文件行号。
\end{itemize}

\textbf{说明}

\begin{itemize}
\tightlist
\item
  示例是局部调用片段；\passthrough{\lstinline!obj!}
  和参数变量表示已经按上层业务校验并准备好的对象。
\end{itemize}

\textbf{用法}

\begin{lstlisting}[language=Python]
different = left != right
\end{lstlisting}

\hypertarget{unitree-sdk2-cpp-idl-go2-lowstate-head}{%
\paragraph{\texorpdfstring{\texttt{unitree\_sdk2\_cpp.idl.go2.LowState.head}}{unitree\_sdk2\_cpp.idl.go2.LowState.head}}\label{unitree-sdk2-cpp-idl-go2-lowstate-head}}

底层 \passthrough{\lstinline!unitree\_go::msg::dds\_::LowState\_!} 的
\passthrough{\lstinline!head!} 字段。类型签名只说明 Python
表示；单位、数值范围、数组固定长度和枚举含义需查对应 IDL 头文件。
底层是固定长度数组，写入序列必须正好包含 2 个元素；元素约束：取值范围为
0 到 255。

\textbf{签名}

\begin{lstlisting}[language=Python]
@property
def head(self) -> list[int]

@head.setter
def head(self, value: Sequence[int]) -> None
\end{lstlisting}

\textbf{可用性与安全性}

\textbf{\passthrough{\lstinline!AVAILABLE!} /
\passthrough{\lstinline!VALUE\_TYPE!}}：当前绑定源码已实现该消息值操作。

\textbf{参数}

\begin{longtable}[]{@{}
  >{\raggedright\arraybackslash}p{(\columnwidth - 4\tabcolsep) * \real{0.18}}
  >{\raggedright\arraybackslash}p{(\columnwidth - 4\tabcolsep) * \real{0.20}}
  >{\raggedright\arraybackslash}p{(\columnwidth - 4\tabcolsep) * \real{0.62}}@{}}
\toprule
名称 & Python 类型 & 含义 \\
\midrule
\endhead
\passthrough{\lstinline!value!} &
\passthrough{\lstinline!Sequence[int]!} &
要写入或传递的新值，必须符合该参数的 Python 类型以及底层 C++
范围约束。 \\
\bottomrule
\end{longtable}

\textbf{返回值}

\begin{longtable}[]{@{}
  >{\raggedright\arraybackslash}p{(\columnwidth - 4\tabcolsep) * \real{0.18}}
  >{\raggedright\arraybackslash}p{(\columnwidth - 4\tabcolsep) * \real{0.20}}
  >{\raggedright\arraybackslash}p{(\columnwidth - 4\tabcolsep) * \real{0.62}}@{}}
\toprule
位置 & 类型 & 含义 \\
\midrule
\endhead
getter 返回值 & \passthrough{\lstinline!list[int]!} &
返回属性当前值。数组、vector 和嵌套消息采用复制语义。 \\
\bottomrule
\end{longtable}

\textbf{对应 C++}

\begin{itemize}
\tightlist
\item
  类：\passthrough{\lstinline!unitree\_go::msg::dds\_::LowState\_!}
\item
  字段类型：\passthrough{\lstinline!std::array<uint8\_t, 2>!}
\item
  声明位置：\passthrough{\lstinline!include/unitree/idl/go2/LowState\_.hpp:30!}
\end{itemize}

\textbf{用法}

\begin{lstlisting}[language=Python]
current_value = obj.head
obj.head = new_value
\end{lstlisting}

\begin{quote}
{[}!TIP{]} 读取容器或嵌套值后应遵循``读取、修改、重新赋值''。只修改
getter 返回的副本不会可靠地更新原 C++ 消息。
\end{quote}

\hypertarget{unitree-sdk2-cpp-idl-go2-lowstate-level-flag}{%
\paragraph{\texorpdfstring{\texttt{unitree\_sdk2\_cpp.idl.go2.LowState.level\_flag}}{unitree\_sdk2\_cpp.idl.go2.LowState.level\_flag}}\label{unitree-sdk2-cpp-idl-go2-lowstate-level-flag}}

底层 \passthrough{\lstinline!unitree\_go::msg::dds\_::LowState\_!} 的
\passthrough{\lstinline!level\_flag!} 字段。类型签名只说明 Python
表示；单位、数值范围、数组固定长度和枚举含义需查对应 IDL 头文件。
取值范围为 0 到 255。

\textbf{签名}

\begin{lstlisting}[language=Python]
@property
def level_flag(self) -> int

@level_flag.setter
def level_flag(self, value: int) -> None
\end{lstlisting}

\textbf{可用性与安全性}

\textbf{\passthrough{\lstinline!AVAILABLE!} /
\passthrough{\lstinline!VALUE\_TYPE!}}：当前绑定源码已实现该消息值操作。

\textbf{参数}

\begin{longtable}[]{@{}
  >{\raggedright\arraybackslash}p{(\columnwidth - 4\tabcolsep) * \real{0.18}}
  >{\raggedright\arraybackslash}p{(\columnwidth - 4\tabcolsep) * \real{0.20}}
  >{\raggedright\arraybackslash}p{(\columnwidth - 4\tabcolsep) * \real{0.62}}@{}}
\toprule
名称 & Python 类型 & 含义 \\
\midrule
\endhead
\passthrough{\lstinline!value!} & \passthrough{\lstinline!int!} &
要写入或传递的新值，必须符合该参数的 Python 类型以及底层 C++
范围约束。 \\
\bottomrule
\end{longtable}

\textbf{返回值}

\begin{longtable}[]{@{}
  >{\raggedright\arraybackslash}p{(\columnwidth - 4\tabcolsep) * \real{0.18}}
  >{\raggedright\arraybackslash}p{(\columnwidth - 4\tabcolsep) * \real{0.20}}
  >{\raggedright\arraybackslash}p{(\columnwidth - 4\tabcolsep) * \real{0.62}}@{}}
\toprule
位置 & 类型 & 含义 \\
\midrule
\endhead
getter 返回值 & \passthrough{\lstinline!int!} & 返回属性当前值。 \\
\bottomrule
\end{longtable}

\textbf{对应 C++}

\begin{itemize}
\tightlist
\item
  类：\passthrough{\lstinline!unitree\_go::msg::dds\_::LowState\_!}
\item
  字段类型：\passthrough{\lstinline!uint8\_t!}
\item
  声明位置：\passthrough{\lstinline!include/unitree/idl/go2/LowState\_.hpp:31!}
\end{itemize}

\textbf{用法}

\begin{lstlisting}[language=Python]
current_value = obj.level_flag
obj.level_flag = new_value
\end{lstlisting}

\hypertarget{unitree-sdk2-cpp-idl-go2-lowstate-frame-reserve}{%
\paragraph{\texorpdfstring{\texttt{unitree\_sdk2\_cpp.idl.go2.LowState.frame\_reserve}}{unitree\_sdk2\_cpp.idl.go2.LowState.frame\_reserve}}\label{unitree-sdk2-cpp-idl-go2-lowstate-frame-reserve}}

底层 \passthrough{\lstinline!unitree\_go::msg::dds\_::LowState\_!} 的
\passthrough{\lstinline!frame\_reserve!} 字段。类型签名只说明 Python
表示；单位、数值范围、数组固定长度和枚举含义需查对应 IDL 头文件。
取值范围为 0 到 255。

\textbf{签名}

\begin{lstlisting}[language=Python]
@property
def frame_reserve(self) -> int

@frame_reserve.setter
def frame_reserve(self, value: int) -> None
\end{lstlisting}

\textbf{可用性与安全性}

\textbf{\passthrough{\lstinline!AVAILABLE!} /
\passthrough{\lstinline!VALUE\_TYPE!}}：当前绑定源码已实现该消息值操作。

\textbf{参数}

\begin{longtable}[]{@{}
  >{\raggedright\arraybackslash}p{(\columnwidth - 4\tabcolsep) * \real{0.18}}
  >{\raggedright\arraybackslash}p{(\columnwidth - 4\tabcolsep) * \real{0.20}}
  >{\raggedright\arraybackslash}p{(\columnwidth - 4\tabcolsep) * \real{0.62}}@{}}
\toprule
名称 & Python 类型 & 含义 \\
\midrule
\endhead
\passthrough{\lstinline!value!} & \passthrough{\lstinline!int!} &
要写入或传递的新值，必须符合该参数的 Python 类型以及底层 C++
范围约束。 \\
\bottomrule
\end{longtable}

\textbf{返回值}

\begin{longtable}[]{@{}
  >{\raggedright\arraybackslash}p{(\columnwidth - 4\tabcolsep) * \real{0.18}}
  >{\raggedright\arraybackslash}p{(\columnwidth - 4\tabcolsep) * \real{0.20}}
  >{\raggedright\arraybackslash}p{(\columnwidth - 4\tabcolsep) * \real{0.62}}@{}}
\toprule
位置 & 类型 & 含义 \\
\midrule
\endhead
getter 返回值 & \passthrough{\lstinline!int!} & 返回属性当前值。 \\
\bottomrule
\end{longtable}

\textbf{对应 C++}

\begin{itemize}
\tightlist
\item
  类：\passthrough{\lstinline!unitree\_go::msg::dds\_::LowState\_!}
\item
  字段类型：\passthrough{\lstinline!uint8\_t!}
\item
  声明位置：\passthrough{\lstinline!include/unitree/idl/go2/LowState\_.hpp:32!}
\end{itemize}

\textbf{用法}

\begin{lstlisting}[language=Python]
current_value = obj.frame_reserve
obj.frame_reserve = new_value
\end{lstlisting}

\hypertarget{unitree-sdk2-cpp-idl-go2-lowstate-sn}{%
\paragraph{\texorpdfstring{\texttt{unitree\_sdk2\_cpp.idl.go2.LowState.sn}}{unitree\_sdk2\_cpp.idl.go2.LowState.sn}}\label{unitree-sdk2-cpp-idl-go2-lowstate-sn}}

底层 \passthrough{\lstinline!unitree\_go::msg::dds\_::LowState\_!} 的
\passthrough{\lstinline!sn!} 字段。类型签名只说明 Python
表示；单位、数值范围、数组固定长度和枚举含义需查对应 IDL 头文件。
底层是固定长度数组，写入序列必须正好包含 2 个元素；元素约束：取值范围为
0 到 4294967295。

\textbf{签名}

\begin{lstlisting}[language=Python]
@property
def sn(self) -> list[int]

@sn.setter
def sn(self, value: Sequence[int]) -> None
\end{lstlisting}

\textbf{可用性与安全性}

\textbf{\passthrough{\lstinline!AVAILABLE!} /
\passthrough{\lstinline!VALUE\_TYPE!}}：当前绑定源码已实现该消息值操作。

\textbf{参数}

\begin{longtable}[]{@{}
  >{\raggedright\arraybackslash}p{(\columnwidth - 4\tabcolsep) * \real{0.18}}
  >{\raggedright\arraybackslash}p{(\columnwidth - 4\tabcolsep) * \real{0.20}}
  >{\raggedright\arraybackslash}p{(\columnwidth - 4\tabcolsep) * \real{0.62}}@{}}
\toprule
名称 & Python 类型 & 含义 \\
\midrule
\endhead
\passthrough{\lstinline!value!} &
\passthrough{\lstinline!Sequence[int]!} &
要写入或传递的新值，必须符合该参数的 Python 类型以及底层 C++
范围约束。 \\
\bottomrule
\end{longtable}

\textbf{返回值}

\begin{longtable}[]{@{}
  >{\raggedright\arraybackslash}p{(\columnwidth - 4\tabcolsep) * \real{0.18}}
  >{\raggedright\arraybackslash}p{(\columnwidth - 4\tabcolsep) * \real{0.20}}
  >{\raggedright\arraybackslash}p{(\columnwidth - 4\tabcolsep) * \real{0.62}}@{}}
\toprule
位置 & 类型 & 含义 \\
\midrule
\endhead
getter 返回值 & \passthrough{\lstinline!list[int]!} &
返回属性当前值。数组、vector 和嵌套消息采用复制语义。 \\
\bottomrule
\end{longtable}

\textbf{对应 C++}

\begin{itemize}
\tightlist
\item
  类：\passthrough{\lstinline!unitree\_go::msg::dds\_::LowState\_!}
\item
  字段类型：\passthrough{\lstinline!std::array<uint32\_t, 2>!}
\item
  声明位置：\passthrough{\lstinline!include/unitree/idl/go2/LowState\_.hpp:33!}
\end{itemize}

\textbf{用法}

\begin{lstlisting}[language=Python]
current_value = obj.sn
obj.sn = new_value
\end{lstlisting}

\begin{quote}
{[}!TIP{]} 读取容器或嵌套值后应遵循``读取、修改、重新赋值''。只修改
getter 返回的副本不会可靠地更新原 C++ 消息。
\end{quote}

\hypertarget{unitree-sdk2-cpp-idl-go2-lowstate-version}{%
\paragraph{\texorpdfstring{\texttt{unitree\_sdk2\_cpp.idl.go2.LowState.version}}{unitree\_sdk2\_cpp.idl.go2.LowState.version}}\label{unitree-sdk2-cpp-idl-go2-lowstate-version}}

底层 \passthrough{\lstinline!unitree\_go::msg::dds\_::LowState\_!} 的
\passthrough{\lstinline!version!} 字段。类型签名只说明 Python
表示；单位、数值范围、数组固定长度和枚举含义需查对应 IDL 头文件。
底层是固定长度数组，写入序列必须正好包含 2 个元素；元素约束：取值范围为
0 到 4294967295。

\textbf{签名}

\begin{lstlisting}[language=Python]
@property
def version(self) -> list[int]

@version.setter
def version(self, value: Sequence[int]) -> None
\end{lstlisting}

\textbf{可用性与安全性}

\textbf{\passthrough{\lstinline!AVAILABLE!} /
\passthrough{\lstinline!VALUE\_TYPE!}}：当前绑定源码已实现该消息值操作。

\textbf{参数}

\begin{longtable}[]{@{}
  >{\raggedright\arraybackslash}p{(\columnwidth - 4\tabcolsep) * \real{0.18}}
  >{\raggedright\arraybackslash}p{(\columnwidth - 4\tabcolsep) * \real{0.20}}
  >{\raggedright\arraybackslash}p{(\columnwidth - 4\tabcolsep) * \real{0.62}}@{}}
\toprule
名称 & Python 类型 & 含义 \\
\midrule
\endhead
\passthrough{\lstinline!value!} &
\passthrough{\lstinline!Sequence[int]!} &
要写入或传递的新值，必须符合该参数的 Python 类型以及底层 C++
范围约束。 \\
\bottomrule
\end{longtable}

\textbf{返回值}

\begin{longtable}[]{@{}
  >{\raggedright\arraybackslash}p{(\columnwidth - 4\tabcolsep) * \real{0.18}}
  >{\raggedright\arraybackslash}p{(\columnwidth - 4\tabcolsep) * \real{0.20}}
  >{\raggedright\arraybackslash}p{(\columnwidth - 4\tabcolsep) * \real{0.62}}@{}}
\toprule
位置 & 类型 & 含义 \\
\midrule
\endhead
getter 返回值 & \passthrough{\lstinline!list[int]!} &
返回属性当前值。数组、vector 和嵌套消息采用复制语义。 \\
\bottomrule
\end{longtable}

\textbf{对应 C++}

\begin{itemize}
\tightlist
\item
  类：\passthrough{\lstinline!unitree\_go::msg::dds\_::LowState\_!}
\item
  字段类型：\passthrough{\lstinline!std::array<uint32\_t, 2>!}
\item
  声明位置：\passthrough{\lstinline!include/unitree/idl/go2/LowState\_.hpp:34!}
\end{itemize}

\textbf{用法}

\begin{lstlisting}[language=Python]
current_value = obj.version
obj.version = new_value
\end{lstlisting}

\begin{quote}
{[}!TIP{]} 读取容器或嵌套值后应遵循``读取、修改、重新赋值''。只修改
getter 返回的副本不会可靠地更新原 C++ 消息。
\end{quote}

\hypertarget{unitree-sdk2-cpp-idl-go2-lowstate-bandwidth}{%
\paragraph{\texorpdfstring{\texttt{unitree\_sdk2\_cpp.idl.go2.LowState.bandwidth}}{unitree\_sdk2\_cpp.idl.go2.LowState.bandwidth}}\label{unitree-sdk2-cpp-idl-go2-lowstate-bandwidth}}

底层 \passthrough{\lstinline!unitree\_go::msg::dds\_::LowState\_!} 的
\passthrough{\lstinline!bandwidth!} 字段。类型签名只说明 Python
表示；单位、数值范围、数组固定长度和枚举含义需查对应 IDL 头文件。
取值范围为 0 到 65535。

\textbf{签名}

\begin{lstlisting}[language=Python]
@property
def bandwidth(self) -> int

@bandwidth.setter
def bandwidth(self, value: int) -> None
\end{lstlisting}

\textbf{可用性与安全性}

\textbf{\passthrough{\lstinline!AVAILABLE!} /
\passthrough{\lstinline!VALUE\_TYPE!}}：当前绑定源码已实现该消息值操作。

\textbf{参数}

\begin{longtable}[]{@{}
  >{\raggedright\arraybackslash}p{(\columnwidth - 4\tabcolsep) * \real{0.18}}
  >{\raggedright\arraybackslash}p{(\columnwidth - 4\tabcolsep) * \real{0.20}}
  >{\raggedright\arraybackslash}p{(\columnwidth - 4\tabcolsep) * \real{0.62}}@{}}
\toprule
名称 & Python 类型 & 含义 \\
\midrule
\endhead
\passthrough{\lstinline!value!} & \passthrough{\lstinline!int!} &
要写入或传递的新值，必须符合该参数的 Python 类型以及底层 C++
范围约束。 \\
\bottomrule
\end{longtable}

\textbf{返回值}

\begin{longtable}[]{@{}
  >{\raggedright\arraybackslash}p{(\columnwidth - 4\tabcolsep) * \real{0.18}}
  >{\raggedright\arraybackslash}p{(\columnwidth - 4\tabcolsep) * \real{0.20}}
  >{\raggedright\arraybackslash}p{(\columnwidth - 4\tabcolsep) * \real{0.62}}@{}}
\toprule
位置 & 类型 & 含义 \\
\midrule
\endhead
getter 返回值 & \passthrough{\lstinline!int!} & 返回属性当前值。 \\
\bottomrule
\end{longtable}

\textbf{对应 C++}

\begin{itemize}
\tightlist
\item
  类：\passthrough{\lstinline!unitree\_go::msg::dds\_::LowState\_!}
\item
  字段类型：\passthrough{\lstinline!uint16\_t!}
\item
  声明位置：\passthrough{\lstinline!include/unitree/idl/go2/LowState\_.hpp:35!}
\end{itemize}

\textbf{用法}

\begin{lstlisting}[language=Python]
current_value = obj.bandwidth
obj.bandwidth = new_value
\end{lstlisting}

\hypertarget{unitree-sdk2-cpp-idl-go2-lowstate-imu-state}{%
\paragraph{\texorpdfstring{\texttt{unitree\_sdk2\_cpp.idl.go2.LowState.imu\_state}}{unitree\_sdk2\_cpp.idl.go2.LowState.imu\_state}}\label{unitree-sdk2-cpp-idl-go2-lowstate-imu-state}}

底层 \passthrough{\lstinline!unitree\_go::msg::dds\_::LowState\_!} 的
\passthrough{\lstinline!imu\_state!} 字段。类型签名只说明 Python
表示；单位、数值范围、数组固定长度和枚举含义需查对应 IDL 头文件。

\textbf{签名}

\begin{lstlisting}[language=Python]
@property
def imu_state(self) -> IMUState

@imu_state.setter
def imu_state(self, value: IMUState) -> None
\end{lstlisting}

\textbf{可用性与安全性}

\textbf{\passthrough{\lstinline!AVAILABLE!} /
\passthrough{\lstinline!VALUE\_TYPE!}}：当前绑定源码已实现该消息值操作。

\textbf{参数}

\begin{longtable}[]{@{}
  >{\raggedright\arraybackslash}p{(\columnwidth - 4\tabcolsep) * \real{0.18}}
  >{\raggedright\arraybackslash}p{(\columnwidth - 4\tabcolsep) * \real{0.20}}
  >{\raggedright\arraybackslash}p{(\columnwidth - 4\tabcolsep) * \real{0.62}}@{}}
\toprule
名称 & Python 类型 & 含义 \\
\midrule
\endhead
\passthrough{\lstinline!value!} & \passthrough{\lstinline!IMUState!} &
要写入或传递的新值，必须符合该参数的 Python 类型以及底层 C++
范围约束。 \\
\bottomrule
\end{longtable}

\textbf{返回值}

\begin{longtable}[]{@{}
  >{\raggedright\arraybackslash}p{(\columnwidth - 4\tabcolsep) * \real{0.18}}
  >{\raggedright\arraybackslash}p{(\columnwidth - 4\tabcolsep) * \real{0.20}}
  >{\raggedright\arraybackslash}p{(\columnwidth - 4\tabcolsep) * \real{0.62}}@{}}
\toprule
位置 & 类型 & 含义 \\
\midrule
\endhead
getter 返回值 & \passthrough{\lstinline!IMUState!} &
返回属性当前值。数组、vector 和嵌套消息采用复制语义。 \\
\bottomrule
\end{longtable}

\textbf{对应 C++}

\begin{itemize}
\tightlist
\item
  类：\passthrough{\lstinline!unitree\_go::msg::dds\_::LowState\_!}
\item
  字段类型：\passthrough{\lstinline!::unitree\_go::msg::dds\_::IMUState\_!}
\item
  声明位置：\passthrough{\lstinline!include/unitree/idl/go2/LowState\_.hpp:36!}
\end{itemize}

\textbf{用法}

\begin{lstlisting}[language=Python]
current_value = obj.imu_state
obj.imu_state = new_value
\end{lstlisting}

\hypertarget{unitree-sdk2-cpp-idl-go2-lowstate-motor-state}{%
\paragraph{\texorpdfstring{\texttt{unitree\_sdk2\_cpp.idl.go2.LowState.motor\_state}}{unitree\_sdk2\_cpp.idl.go2.LowState.motor\_state}}\label{unitree-sdk2-cpp-idl-go2-lowstate-motor-state}}

底层 \passthrough{\lstinline!unitree\_go::msg::dds\_::LowState\_!} 的
\passthrough{\lstinline!motor\_state!} 字段。类型签名只说明 Python
表示；单位、数值范围、数组固定长度和枚举含义需查对应 IDL 头文件。
底层是固定长度数组，写入序列必须正好包含 20 个元素

\textbf{签名}

\begin{lstlisting}[language=Python]
@property
def motor_state(self) -> list[MotorState]

@motor_state.setter
def motor_state(self, value: Sequence[MotorState]) -> None
\end{lstlisting}

\textbf{可用性与安全性}

\textbf{\passthrough{\lstinline!AVAILABLE!} /
\passthrough{\lstinline!VALUE\_TYPE!}}：当前绑定源码已实现该消息值操作。

\textbf{参数}

\begin{longtable}[]{@{}
  >{\raggedright\arraybackslash}p{(\columnwidth - 4\tabcolsep) * \real{0.18}}
  >{\raggedright\arraybackslash}p{(\columnwidth - 4\tabcolsep) * \real{0.20}}
  >{\raggedright\arraybackslash}p{(\columnwidth - 4\tabcolsep) * \real{0.62}}@{}}
\toprule
名称 & Python 类型 & 含义 \\
\midrule
\endhead
\passthrough{\lstinline!value!} &
\passthrough{\lstinline!Sequence[MotorState]!} &
要写入或传递的新值，必须符合该参数的 Python 类型以及底层 C++
范围约束。 \\
\bottomrule
\end{longtable}

\textbf{返回值}

\begin{longtable}[]{@{}
  >{\raggedright\arraybackslash}p{(\columnwidth - 4\tabcolsep) * \real{0.18}}
  >{\raggedright\arraybackslash}p{(\columnwidth - 4\tabcolsep) * \real{0.20}}
  >{\raggedright\arraybackslash}p{(\columnwidth - 4\tabcolsep) * \real{0.62}}@{}}
\toprule
位置 & 类型 & 含义 \\
\midrule
\endhead
getter 返回值 & \passthrough{\lstinline!list[MotorState]!} &
返回属性当前值。数组、vector 和嵌套消息采用复制语义。 \\
\bottomrule
\end{longtable}

\textbf{对应 C++}

\begin{itemize}
\tightlist
\item
  类：\passthrough{\lstinline!unitree\_go::msg::dds\_::LowState\_!}
\item
  字段类型：\passthrough{\lstinline!std::array< ::unitree\_go::msg::dds\_::MotorState\_, 20>!}
\item
  声明位置：\passthrough{\lstinline!include/unitree/idl/go2/LowState\_.hpp:37!}
\end{itemize}

\textbf{用法}

\begin{lstlisting}[language=Python]
current_value = obj.motor_state
obj.motor_state = new_value
\end{lstlisting}

\begin{quote}
{[}!TIP{]} 读取容器或嵌套值后应遵循``读取、修改、重新赋值''。只修改
getter 返回的副本不会可靠地更新原 C++ 消息。
\end{quote}

\hypertarget{unitree-sdk2-cpp-idl-go2-lowstate-bms-state}{%
\paragraph{\texorpdfstring{\texttt{unitree\_sdk2\_cpp.idl.go2.LowState.bms\_state}}{unitree\_sdk2\_cpp.idl.go2.LowState.bms\_state}}\label{unitree-sdk2-cpp-idl-go2-lowstate-bms-state}}

底层 \passthrough{\lstinline!unitree\_go::msg::dds\_::LowState\_!} 的
\passthrough{\lstinline!bms\_state!} 字段。类型签名只说明 Python
表示；单位、数值范围、数组固定长度和枚举含义需查对应 IDL 头文件。

\textbf{签名}

\begin{lstlisting}[language=Python]
@property
def bms_state(self) -> BmsState

@bms_state.setter
def bms_state(self, value: BmsState) -> None
\end{lstlisting}

\textbf{可用性与安全性}

\textbf{\passthrough{\lstinline!AVAILABLE!} /
\passthrough{\lstinline!VALUE\_TYPE!}}：当前绑定源码已实现该消息值操作。

\textbf{参数}

\begin{longtable}[]{@{}
  >{\raggedright\arraybackslash}p{(\columnwidth - 4\tabcolsep) * \real{0.18}}
  >{\raggedright\arraybackslash}p{(\columnwidth - 4\tabcolsep) * \real{0.20}}
  >{\raggedright\arraybackslash}p{(\columnwidth - 4\tabcolsep) * \real{0.62}}@{}}
\toprule
名称 & Python 类型 & 含义 \\
\midrule
\endhead
\passthrough{\lstinline!value!} & \passthrough{\lstinline!BmsState!} &
要写入或传递的新值，必须符合该参数的 Python 类型以及底层 C++
范围约束。 \\
\bottomrule
\end{longtable}

\textbf{返回值}

\begin{longtable}[]{@{}
  >{\raggedright\arraybackslash}p{(\columnwidth - 4\tabcolsep) * \real{0.18}}
  >{\raggedright\arraybackslash}p{(\columnwidth - 4\tabcolsep) * \real{0.20}}
  >{\raggedright\arraybackslash}p{(\columnwidth - 4\tabcolsep) * \real{0.62}}@{}}
\toprule
位置 & 类型 & 含义 \\
\midrule
\endhead
getter 返回值 & \passthrough{\lstinline!BmsState!} &
返回属性当前值。数组、vector 和嵌套消息采用复制语义。 \\
\bottomrule
\end{longtable}

\textbf{对应 C++}

\begin{itemize}
\tightlist
\item
  类：\passthrough{\lstinline!unitree\_go::msg::dds\_::LowState\_!}
\item
  字段类型：\passthrough{\lstinline!::unitree\_go::msg::dds\_::BmsState\_!}
\item
  声明位置：\passthrough{\lstinline!include/unitree/idl/go2/LowState\_.hpp:38!}
\end{itemize}

\textbf{用法}

\begin{lstlisting}[language=Python]
current_value = obj.bms_state
obj.bms_state = new_value
\end{lstlisting}

\hypertarget{unitree-sdk2-cpp-idl-go2-lowstate-foot-force}{%
\paragraph{\texorpdfstring{\texttt{unitree\_sdk2\_cpp.idl.go2.LowState.foot\_force}}{unitree\_sdk2\_cpp.idl.go2.LowState.foot\_force}}\label{unitree-sdk2-cpp-idl-go2-lowstate-foot-force}}

底层 \passthrough{\lstinline!unitree\_go::msg::dds\_::LowState\_!} 的
\passthrough{\lstinline!foot\_force!} 字段。类型签名只说明 Python
表示；单位、数值范围、数组固定长度和枚举含义需查对应 IDL 头文件。
底层是固定长度数组，写入序列必须正好包含 4 个元素；元素约束：取值范围为
-32768 到 32767。

\textbf{签名}

\begin{lstlisting}[language=Python]
@property
def foot_force(self) -> list[int]

@foot_force.setter
def foot_force(self, value: Sequence[int]) -> None
\end{lstlisting}

\textbf{可用性与安全性}

\textbf{\passthrough{\lstinline!AVAILABLE!} /
\passthrough{\lstinline!VALUE\_TYPE!}}：当前绑定源码已实现该消息值操作。

\textbf{参数}

\begin{longtable}[]{@{}
  >{\raggedright\arraybackslash}p{(\columnwidth - 4\tabcolsep) * \real{0.18}}
  >{\raggedright\arraybackslash}p{(\columnwidth - 4\tabcolsep) * \real{0.20}}
  >{\raggedright\arraybackslash}p{(\columnwidth - 4\tabcolsep) * \real{0.62}}@{}}
\toprule
名称 & Python 类型 & 含义 \\
\midrule
\endhead
\passthrough{\lstinline!value!} &
\passthrough{\lstinline!Sequence[int]!} &
要写入或传递的新值，必须符合该参数的 Python 类型以及底层 C++
范围约束。 \\
\bottomrule
\end{longtable}

\textbf{返回值}

\begin{longtable}[]{@{}
  >{\raggedright\arraybackslash}p{(\columnwidth - 4\tabcolsep) * \real{0.18}}
  >{\raggedright\arraybackslash}p{(\columnwidth - 4\tabcolsep) * \real{0.20}}
  >{\raggedright\arraybackslash}p{(\columnwidth - 4\tabcolsep) * \real{0.62}}@{}}
\toprule
位置 & 类型 & 含义 \\
\midrule
\endhead
getter 返回值 & \passthrough{\lstinline!list[int]!} &
返回属性当前值。数组、vector 和嵌套消息采用复制语义。 \\
\bottomrule
\end{longtable}

\textbf{对应 C++}

\begin{itemize}
\tightlist
\item
  类：\passthrough{\lstinline!unitree\_go::msg::dds\_::LowState\_!}
\item
  字段类型：\passthrough{\lstinline!std::array<int16\_t, 4>!}
\item
  声明位置：\passthrough{\lstinline!include/unitree/idl/go2/LowState\_.hpp:39!}
\end{itemize}

\textbf{用法}

\begin{lstlisting}[language=Python]
current_value = obj.foot_force
obj.foot_force = new_value
\end{lstlisting}

\begin{quote}
{[}!TIP{]} 读取容器或嵌套值后应遵循``读取、修改、重新赋值''。只修改
getter 返回的副本不会可靠地更新原 C++ 消息。
\end{quote}

\hypertarget{unitree-sdk2-cpp-idl-go2-lowstate-foot-force-est}{%
\paragraph{\texorpdfstring{\texttt{unitree\_sdk2\_cpp.idl.go2.LowState.foot\_force\_est}}{unitree\_sdk2\_cpp.idl.go2.LowState.foot\_force\_est}}\label{unitree-sdk2-cpp-idl-go2-lowstate-foot-force-est}}

底层 \passthrough{\lstinline!unitree\_go::msg::dds\_::LowState\_!} 的
\passthrough{\lstinline!foot\_force\_est!} 字段。类型签名只说明 Python
表示；单位、数值范围、数组固定长度和枚举含义需查对应 IDL 头文件。
底层是固定长度数组，写入序列必须正好包含 4 个元素；元素约束：取值范围为
-32768 到 32767。

\textbf{签名}

\begin{lstlisting}[language=Python]
@property
def foot_force_est(self) -> list[int]

@foot_force_est.setter
def foot_force_est(self, value: Sequence[int]) -> None
\end{lstlisting}

\textbf{可用性与安全性}

\textbf{\passthrough{\lstinline!AVAILABLE!} /
\passthrough{\lstinline!VALUE\_TYPE!}}：当前绑定源码已实现该消息值操作。

\textbf{参数}

\begin{longtable}[]{@{}
  >{\raggedright\arraybackslash}p{(\columnwidth - 4\tabcolsep) * \real{0.18}}
  >{\raggedright\arraybackslash}p{(\columnwidth - 4\tabcolsep) * \real{0.20}}
  >{\raggedright\arraybackslash}p{(\columnwidth - 4\tabcolsep) * \real{0.62}}@{}}
\toprule
名称 & Python 类型 & 含义 \\
\midrule
\endhead
\passthrough{\lstinline!value!} &
\passthrough{\lstinline!Sequence[int]!} &
要写入或传递的新值，必须符合该参数的 Python 类型以及底层 C++
范围约束。 \\
\bottomrule
\end{longtable}

\textbf{返回值}

\begin{longtable}[]{@{}
  >{\raggedright\arraybackslash}p{(\columnwidth - 4\tabcolsep) * \real{0.18}}
  >{\raggedright\arraybackslash}p{(\columnwidth - 4\tabcolsep) * \real{0.20}}
  >{\raggedright\arraybackslash}p{(\columnwidth - 4\tabcolsep) * \real{0.62}}@{}}
\toprule
位置 & 类型 & 含义 \\
\midrule
\endhead
getter 返回值 & \passthrough{\lstinline!list[int]!} &
返回属性当前值。数组、vector 和嵌套消息采用复制语义。 \\
\bottomrule
\end{longtable}

\textbf{对应 C++}

\begin{itemize}
\tightlist
\item
  类：\passthrough{\lstinline!unitree\_go::msg::dds\_::LowState\_!}
\item
  字段类型：\passthrough{\lstinline!std::array<int16\_t, 4>!}
\item
  声明位置：\passthrough{\lstinline!include/unitree/idl/go2/LowState\_.hpp:40!}
\end{itemize}

\textbf{用法}

\begin{lstlisting}[language=Python]
current_value = obj.foot_force_est
obj.foot_force_est = new_value
\end{lstlisting}

\begin{quote}
{[}!TIP{]} 读取容器或嵌套值后应遵循``读取、修改、重新赋值''。只修改
getter 返回的副本不会可靠地更新原 C++ 消息。
\end{quote}

\hypertarget{unitree-sdk2-cpp-idl-go2-lowstate-tick}{%
\paragraph{\texorpdfstring{\texttt{unitree\_sdk2\_cpp.idl.go2.LowState.tick}}{unitree\_sdk2\_cpp.idl.go2.LowState.tick}}\label{unitree-sdk2-cpp-idl-go2-lowstate-tick}}

底层 \passthrough{\lstinline!unitree\_go::msg::dds\_::LowState\_!} 的
\passthrough{\lstinline!tick!} 字段。类型签名只说明 Python
表示；单位、数值范围、数组固定长度和枚举含义需查对应 IDL 头文件。
取值范围为 0 到 4294967295。

\textbf{签名}

\begin{lstlisting}[language=Python]
@property
def tick(self) -> int

@tick.setter
def tick(self, value: int) -> None
\end{lstlisting}

\textbf{可用性与安全性}

\textbf{\passthrough{\lstinline!AVAILABLE!} /
\passthrough{\lstinline!VALUE\_TYPE!}}：当前绑定源码已实现该消息值操作。

\textbf{参数}

\begin{longtable}[]{@{}
  >{\raggedright\arraybackslash}p{(\columnwidth - 4\tabcolsep) * \real{0.18}}
  >{\raggedright\arraybackslash}p{(\columnwidth - 4\tabcolsep) * \real{0.20}}
  >{\raggedright\arraybackslash}p{(\columnwidth - 4\tabcolsep) * \real{0.62}}@{}}
\toprule
名称 & Python 类型 & 含义 \\
\midrule
\endhead
\passthrough{\lstinline!value!} & \passthrough{\lstinline!int!} &
要写入或传递的新值，必须符合该参数的 Python 类型以及底层 C++
范围约束。 \\
\bottomrule
\end{longtable}

\textbf{返回值}

\begin{longtable}[]{@{}
  >{\raggedright\arraybackslash}p{(\columnwidth - 4\tabcolsep) * \real{0.18}}
  >{\raggedright\arraybackslash}p{(\columnwidth - 4\tabcolsep) * \real{0.20}}
  >{\raggedright\arraybackslash}p{(\columnwidth - 4\tabcolsep) * \real{0.62}}@{}}
\toprule
位置 & 类型 & 含义 \\
\midrule
\endhead
getter 返回值 & \passthrough{\lstinline!int!} & 返回属性当前值。 \\
\bottomrule
\end{longtable}

\textbf{对应 C++}

\begin{itemize}
\tightlist
\item
  类：\passthrough{\lstinline!unitree\_go::msg::dds\_::LowState\_!}
\item
  字段类型：\passthrough{\lstinline!uint32\_t!}
\item
  声明位置：\passthrough{\lstinline!include/unitree/idl/go2/LowState\_.hpp:41!}
\end{itemize}

\textbf{用法}

\begin{lstlisting}[language=Python]
current_value = obj.tick
obj.tick = new_value
\end{lstlisting}

\hypertarget{unitree-sdk2-cpp-idl-go2-lowstate-wireless-remote}{%
\paragraph{\texorpdfstring{\texttt{unitree\_sdk2\_cpp.idl.go2.LowState.wireless\_remote}}{unitree\_sdk2\_cpp.idl.go2.LowState.wireless\_remote}}\label{unitree-sdk2-cpp-idl-go2-lowstate-wireless-remote}}

底层 \passthrough{\lstinline!unitree\_go::msg::dds\_::LowState\_!} 的
\passthrough{\lstinline!wireless\_remote!} 字段。类型签名只说明 Python
表示；单位、数值范围、数组固定长度和枚举含义需查对应 IDL 头文件。
底层是固定长度数组，写入序列必须正好包含 40 个元素；元素约束：取值范围为
0 到 255。

\textbf{签名}

\begin{lstlisting}[language=Python]
@property
def wireless_remote(self) -> list[int]

@wireless_remote.setter
def wireless_remote(self, value: Sequence[int]) -> None
\end{lstlisting}

\textbf{可用性与安全性}

\textbf{\passthrough{\lstinline!AVAILABLE!} /
\passthrough{\lstinline!VALUE\_TYPE!}}：当前绑定源码已实现该消息值操作。

\textbf{参数}

\begin{longtable}[]{@{}
  >{\raggedright\arraybackslash}p{(\columnwidth - 4\tabcolsep) * \real{0.18}}
  >{\raggedright\arraybackslash}p{(\columnwidth - 4\tabcolsep) * \real{0.20}}
  >{\raggedright\arraybackslash}p{(\columnwidth - 4\tabcolsep) * \real{0.62}}@{}}
\toprule
名称 & Python 类型 & 含义 \\
\midrule
\endhead
\passthrough{\lstinline!value!} &
\passthrough{\lstinline!Sequence[int]!} &
要写入或传递的新值，必须符合该参数的 Python 类型以及底层 C++
范围约束。 \\
\bottomrule
\end{longtable}

\textbf{返回值}

\begin{longtable}[]{@{}
  >{\raggedright\arraybackslash}p{(\columnwidth - 4\tabcolsep) * \real{0.18}}
  >{\raggedright\arraybackslash}p{(\columnwidth - 4\tabcolsep) * \real{0.20}}
  >{\raggedright\arraybackslash}p{(\columnwidth - 4\tabcolsep) * \real{0.62}}@{}}
\toprule
位置 & 类型 & 含义 \\
\midrule
\endhead
getter 返回值 & \passthrough{\lstinline!list[int]!} &
返回属性当前值。数组、vector 和嵌套消息采用复制语义。 \\
\bottomrule
\end{longtable}

\textbf{对应 C++}

\begin{itemize}
\tightlist
\item
  类：\passthrough{\lstinline!unitree\_go::msg::dds\_::LowState\_!}
\item
  字段类型：\passthrough{\lstinline!std::array<uint8\_t, 40>!}
\item
  声明位置：\passthrough{\lstinline!include/unitree/idl/go2/LowState\_.hpp:42!}
\end{itemize}

\textbf{用法}

\begin{lstlisting}[language=Python]
current_value = obj.wireless_remote
obj.wireless_remote = new_value
\end{lstlisting}

\begin{quote}
{[}!TIP{]} 读取容器或嵌套值后应遵循``读取、修改、重新赋值''。只修改
getter 返回的副本不会可靠地更新原 C++ 消息。
\end{quote}

\hypertarget{unitree-sdk2-cpp-idl-go2-lowstate-bit-flag}{%
\paragraph{\texorpdfstring{\texttt{unitree\_sdk2\_cpp.idl.go2.LowState.bit\_flag}}{unitree\_sdk2\_cpp.idl.go2.LowState.bit\_flag}}\label{unitree-sdk2-cpp-idl-go2-lowstate-bit-flag}}

底层 \passthrough{\lstinline!unitree\_go::msg::dds\_::LowState\_!} 的
\passthrough{\lstinline!bit\_flag!} 字段。类型签名只说明 Python
表示；单位、数值范围、数组固定长度和枚举含义需查对应 IDL 头文件。
取值范围为 0 到 255。

\textbf{签名}

\begin{lstlisting}[language=Python]
@property
def bit_flag(self) -> int

@bit_flag.setter
def bit_flag(self, value: int) -> None
\end{lstlisting}

\textbf{可用性与安全性}

\textbf{\passthrough{\lstinline!AVAILABLE!} /
\passthrough{\lstinline!VALUE\_TYPE!}}：当前绑定源码已实现该消息值操作。

\textbf{参数}

\begin{longtable}[]{@{}
  >{\raggedright\arraybackslash}p{(\columnwidth - 4\tabcolsep) * \real{0.18}}
  >{\raggedright\arraybackslash}p{(\columnwidth - 4\tabcolsep) * \real{0.20}}
  >{\raggedright\arraybackslash}p{(\columnwidth - 4\tabcolsep) * \real{0.62}}@{}}
\toprule
名称 & Python 类型 & 含义 \\
\midrule
\endhead
\passthrough{\lstinline!value!} & \passthrough{\lstinline!int!} &
要写入或传递的新值，必须符合该参数的 Python 类型以及底层 C++
范围约束。 \\
\bottomrule
\end{longtable}

\textbf{返回值}

\begin{longtable}[]{@{}
  >{\raggedright\arraybackslash}p{(\columnwidth - 4\tabcolsep) * \real{0.18}}
  >{\raggedright\arraybackslash}p{(\columnwidth - 4\tabcolsep) * \real{0.20}}
  >{\raggedright\arraybackslash}p{(\columnwidth - 4\tabcolsep) * \real{0.62}}@{}}
\toprule
位置 & 类型 & 含义 \\
\midrule
\endhead
getter 返回值 & \passthrough{\lstinline!int!} & 返回属性当前值。 \\
\bottomrule
\end{longtable}

\textbf{对应 C++}

\begin{itemize}
\tightlist
\item
  类：\passthrough{\lstinline!unitree\_go::msg::dds\_::LowState\_!}
\item
  字段类型：\passthrough{\lstinline!uint8\_t!}
\item
  声明位置：\passthrough{\lstinline!include/unitree/idl/go2/LowState\_.hpp:43!}
\end{itemize}

\textbf{用法}

\begin{lstlisting}[language=Python]
current_value = obj.bit_flag
obj.bit_flag = new_value
\end{lstlisting}

\hypertarget{unitree-sdk2-cpp-idl-go2-lowstate-adc-reel}{%
\paragraph{\texorpdfstring{\texttt{unitree\_sdk2\_cpp.idl.go2.LowState.adc\_reel}}{unitree\_sdk2\_cpp.idl.go2.LowState.adc\_reel}}\label{unitree-sdk2-cpp-idl-go2-lowstate-adc-reel}}

底层 \passthrough{\lstinline!unitree\_go::msg::dds\_::LowState\_!} 的
\passthrough{\lstinline!adc\_reel!} 字段。类型签名只说明 Python
表示；单位、数值范围、数组固定长度和枚举含义需查对应 IDL 头文件。
底层为浮点数；类型本身不说明单位、坐标系或业务安全范围。

\textbf{签名}

\begin{lstlisting}[language=Python]
@property
def adc_reel(self) -> float

@adc_reel.setter
def adc_reel(self, value: float) -> None
\end{lstlisting}

\textbf{可用性与安全性}

\textbf{\passthrough{\lstinline!AVAILABLE!} /
\passthrough{\lstinline!VALUE\_TYPE!}}：当前绑定源码已实现该消息值操作。

\textbf{参数}

\begin{longtable}[]{@{}
  >{\raggedright\arraybackslash}p{(\columnwidth - 4\tabcolsep) * \real{0.18}}
  >{\raggedright\arraybackslash}p{(\columnwidth - 4\tabcolsep) * \real{0.20}}
  >{\raggedright\arraybackslash}p{(\columnwidth - 4\tabcolsep) * \real{0.62}}@{}}
\toprule
名称 & Python 类型 & 含义 \\
\midrule
\endhead
\passthrough{\lstinline!value!} & \passthrough{\lstinline!float!} &
要写入或传递的新值，必须符合该参数的 Python 类型以及底层 C++
范围约束。 \\
\bottomrule
\end{longtable}

\textbf{返回值}

\begin{longtable}[]{@{}
  >{\raggedright\arraybackslash}p{(\columnwidth - 4\tabcolsep) * \real{0.18}}
  >{\raggedright\arraybackslash}p{(\columnwidth - 4\tabcolsep) * \real{0.20}}
  >{\raggedright\arraybackslash}p{(\columnwidth - 4\tabcolsep) * \real{0.62}}@{}}
\toprule
位置 & 类型 & 含义 \\
\midrule
\endhead
getter 返回值 & \passthrough{\lstinline!float!} & 返回属性当前值。 \\
\bottomrule
\end{longtable}

\textbf{对应 C++}

\begin{itemize}
\tightlist
\item
  类：\passthrough{\lstinline!unitree\_go::msg::dds\_::LowState\_!}
\item
  字段类型：\passthrough{\lstinline!float!}
\item
  声明位置：\passthrough{\lstinline!include/unitree/idl/go2/LowState\_.hpp:44!}
\end{itemize}

\textbf{用法}

\begin{lstlisting}[language=Python]
current_value = obj.adc_reel
obj.adc_reel = new_value
\end{lstlisting}

\hypertarget{unitree-sdk2-cpp-idl-go2-lowstate-temperature-ntc1}{%
\paragraph{\texorpdfstring{\texttt{unitree\_sdk2\_cpp.idl.go2.LowState.temperature\_ntc1}}{unitree\_sdk2\_cpp.idl.go2.LowState.temperature\_ntc1}}\label{unitree-sdk2-cpp-idl-go2-lowstate-temperature-ntc1}}

底层 \passthrough{\lstinline!unitree\_go::msg::dds\_::LowState\_!} 的
\passthrough{\lstinline!temperature\_ntc1!} 字段。类型签名只说明 Python
表示；单位、数值范围、数组固定长度和枚举含义需查对应 IDL 头文件。
取值范围为 0 到 255。

\textbf{签名}

\begin{lstlisting}[language=Python]
@property
def temperature_ntc1(self) -> int

@temperature_ntc1.setter
def temperature_ntc1(self, value: int) -> None
\end{lstlisting}

\textbf{可用性与安全性}

\textbf{\passthrough{\lstinline!AVAILABLE!} /
\passthrough{\lstinline!VALUE\_TYPE!}}：当前绑定源码已实现该消息值操作。

\textbf{参数}

\begin{longtable}[]{@{}
  >{\raggedright\arraybackslash}p{(\columnwidth - 4\tabcolsep) * \real{0.18}}
  >{\raggedright\arraybackslash}p{(\columnwidth - 4\tabcolsep) * \real{0.20}}
  >{\raggedright\arraybackslash}p{(\columnwidth - 4\tabcolsep) * \real{0.62}}@{}}
\toprule
名称 & Python 类型 & 含义 \\
\midrule
\endhead
\passthrough{\lstinline!value!} & \passthrough{\lstinline!int!} &
要写入或传递的新值，必须符合该参数的 Python 类型以及底层 C++
范围约束。 \\
\bottomrule
\end{longtable}

\textbf{返回值}

\begin{longtable}[]{@{}
  >{\raggedright\arraybackslash}p{(\columnwidth - 4\tabcolsep) * \real{0.18}}
  >{\raggedright\arraybackslash}p{(\columnwidth - 4\tabcolsep) * \real{0.20}}
  >{\raggedright\arraybackslash}p{(\columnwidth - 4\tabcolsep) * \real{0.62}}@{}}
\toprule
位置 & 类型 & 含义 \\
\midrule
\endhead
getter 返回值 & \passthrough{\lstinline!int!} & 返回属性当前值。 \\
\bottomrule
\end{longtable}

\textbf{对应 C++}

\begin{itemize}
\tightlist
\item
  类：\passthrough{\lstinline!unitree\_go::msg::dds\_::LowState\_!}
\item
  字段类型：\passthrough{\lstinline!uint8\_t!}
\item
  声明位置：\passthrough{\lstinline!include/unitree/idl/go2/LowState\_.hpp:45!}
\end{itemize}

\textbf{用法}

\begin{lstlisting}[language=Python]
current_value = obj.temperature_ntc1
obj.temperature_ntc1 = new_value
\end{lstlisting}

\hypertarget{unitree-sdk2-cpp-idl-go2-lowstate-temperature-ntc2}{%
\paragraph{\texorpdfstring{\texttt{unitree\_sdk2\_cpp.idl.go2.LowState.temperature\_ntc2}}{unitree\_sdk2\_cpp.idl.go2.LowState.temperature\_ntc2}}\label{unitree-sdk2-cpp-idl-go2-lowstate-temperature-ntc2}}

底层 \passthrough{\lstinline!unitree\_go::msg::dds\_::LowState\_!} 的
\passthrough{\lstinline!temperature\_ntc2!} 字段。类型签名只说明 Python
表示；单位、数值范围、数组固定长度和枚举含义需查对应 IDL 头文件。
取值范围为 0 到 255。

\textbf{签名}

\begin{lstlisting}[language=Python]
@property
def temperature_ntc2(self) -> int

@temperature_ntc2.setter
def temperature_ntc2(self, value: int) -> None
\end{lstlisting}

\textbf{可用性与安全性}

\textbf{\passthrough{\lstinline!AVAILABLE!} /
\passthrough{\lstinline!VALUE\_TYPE!}}：当前绑定源码已实现该消息值操作。

\textbf{参数}

\begin{longtable}[]{@{}
  >{\raggedright\arraybackslash}p{(\columnwidth - 4\tabcolsep) * \real{0.18}}
  >{\raggedright\arraybackslash}p{(\columnwidth - 4\tabcolsep) * \real{0.20}}
  >{\raggedright\arraybackslash}p{(\columnwidth - 4\tabcolsep) * \real{0.62}}@{}}
\toprule
名称 & Python 类型 & 含义 \\
\midrule
\endhead
\passthrough{\lstinline!value!} & \passthrough{\lstinline!int!} &
要写入或传递的新值，必须符合该参数的 Python 类型以及底层 C++
范围约束。 \\
\bottomrule
\end{longtable}

\textbf{返回值}

\begin{longtable}[]{@{}
  >{\raggedright\arraybackslash}p{(\columnwidth - 4\tabcolsep) * \real{0.18}}
  >{\raggedright\arraybackslash}p{(\columnwidth - 4\tabcolsep) * \real{0.20}}
  >{\raggedright\arraybackslash}p{(\columnwidth - 4\tabcolsep) * \real{0.62}}@{}}
\toprule
位置 & 类型 & 含义 \\
\midrule
\endhead
getter 返回值 & \passthrough{\lstinline!int!} & 返回属性当前值。 \\
\bottomrule
\end{longtable}

\textbf{对应 C++}

\begin{itemize}
\tightlist
\item
  类：\passthrough{\lstinline!unitree\_go::msg::dds\_::LowState\_!}
\item
  字段类型：\passthrough{\lstinline!uint8\_t!}
\item
  声明位置：\passthrough{\lstinline!include/unitree/idl/go2/LowState\_.hpp:46!}
\end{itemize}

\textbf{用法}

\begin{lstlisting}[language=Python]
current_value = obj.temperature_ntc2
obj.temperature_ntc2 = new_value
\end{lstlisting}

\hypertarget{unitree-sdk2-cpp-idl-go2-lowstate-power-v}{%
\paragraph{\texorpdfstring{\texttt{unitree\_sdk2\_cpp.idl.go2.LowState.power\_v}}{unitree\_sdk2\_cpp.idl.go2.LowState.power\_v}}\label{unitree-sdk2-cpp-idl-go2-lowstate-power-v}}

底层 \passthrough{\lstinline!unitree\_go::msg::dds\_::LowState\_!} 的
\passthrough{\lstinline!power\_v!} 字段。类型签名只说明 Python
表示；单位、数值范围、数组固定长度和枚举含义需查对应 IDL 头文件。
底层为浮点数；类型本身不说明单位、坐标系或业务安全范围。

\textbf{签名}

\begin{lstlisting}[language=Python]
@property
def power_v(self) -> float

@power_v.setter
def power_v(self, value: float) -> None
\end{lstlisting}

\textbf{可用性与安全性}

\textbf{\passthrough{\lstinline!AVAILABLE!} /
\passthrough{\lstinline!VALUE\_TYPE!}}：当前绑定源码已实现该消息值操作。

\textbf{参数}

\begin{longtable}[]{@{}
  >{\raggedright\arraybackslash}p{(\columnwidth - 4\tabcolsep) * \real{0.18}}
  >{\raggedright\arraybackslash}p{(\columnwidth - 4\tabcolsep) * \real{0.20}}
  >{\raggedright\arraybackslash}p{(\columnwidth - 4\tabcolsep) * \real{0.62}}@{}}
\toprule
名称 & Python 类型 & 含义 \\
\midrule
\endhead
\passthrough{\lstinline!value!} & \passthrough{\lstinline!float!} &
要写入或传递的新值，必须符合该参数的 Python 类型以及底层 C++
范围约束。 \\
\bottomrule
\end{longtable}

\textbf{返回值}

\begin{longtable}[]{@{}
  >{\raggedright\arraybackslash}p{(\columnwidth - 4\tabcolsep) * \real{0.18}}
  >{\raggedright\arraybackslash}p{(\columnwidth - 4\tabcolsep) * \real{0.20}}
  >{\raggedright\arraybackslash}p{(\columnwidth - 4\tabcolsep) * \real{0.62}}@{}}
\toprule
位置 & 类型 & 含义 \\
\midrule
\endhead
getter 返回值 & \passthrough{\lstinline!float!} & 返回属性当前值。 \\
\bottomrule
\end{longtable}

\textbf{对应 C++}

\begin{itemize}
\tightlist
\item
  类：\passthrough{\lstinline!unitree\_go::msg::dds\_::LowState\_!}
\item
  字段类型：\passthrough{\lstinline!float!}
\item
  声明位置：\passthrough{\lstinline!include/unitree/idl/go2/LowState\_.hpp:47!}
\end{itemize}

\textbf{用法}

\begin{lstlisting}[language=Python]
current_value = obj.power_v
obj.power_v = new_value
\end{lstlisting}

\hypertarget{unitree-sdk2-cpp-idl-go2-lowstate-power-a}{%
\paragraph{\texorpdfstring{\texttt{unitree\_sdk2\_cpp.idl.go2.LowState.power\_a}}{unitree\_sdk2\_cpp.idl.go2.LowState.power\_a}}\label{unitree-sdk2-cpp-idl-go2-lowstate-power-a}}

底层 \passthrough{\lstinline!unitree\_go::msg::dds\_::LowState\_!} 的
\passthrough{\lstinline!power\_a!} 字段。类型签名只说明 Python
表示；单位、数值范围、数组固定长度和枚举含义需查对应 IDL 头文件。
底层为浮点数；类型本身不说明单位、坐标系或业务安全范围。

\textbf{签名}

\begin{lstlisting}[language=Python]
@property
def power_a(self) -> float

@power_a.setter
def power_a(self, value: float) -> None
\end{lstlisting}

\textbf{可用性与安全性}

\textbf{\passthrough{\lstinline!AVAILABLE!} /
\passthrough{\lstinline!VALUE\_TYPE!}}：当前绑定源码已实现该消息值操作。

\textbf{参数}

\begin{longtable}[]{@{}
  >{\raggedright\arraybackslash}p{(\columnwidth - 4\tabcolsep) * \real{0.18}}
  >{\raggedright\arraybackslash}p{(\columnwidth - 4\tabcolsep) * \real{0.20}}
  >{\raggedright\arraybackslash}p{(\columnwidth - 4\tabcolsep) * \real{0.62}}@{}}
\toprule
名称 & Python 类型 & 含义 \\
\midrule
\endhead
\passthrough{\lstinline!value!} & \passthrough{\lstinline!float!} &
要写入或传递的新值，必须符合该参数的 Python 类型以及底层 C++
范围约束。 \\
\bottomrule
\end{longtable}

\textbf{返回值}

\begin{longtable}[]{@{}
  >{\raggedright\arraybackslash}p{(\columnwidth - 4\tabcolsep) * \real{0.18}}
  >{\raggedright\arraybackslash}p{(\columnwidth - 4\tabcolsep) * \real{0.20}}
  >{\raggedright\arraybackslash}p{(\columnwidth - 4\tabcolsep) * \real{0.62}}@{}}
\toprule
位置 & 类型 & 含义 \\
\midrule
\endhead
getter 返回值 & \passthrough{\lstinline!float!} & 返回属性当前值。 \\
\bottomrule
\end{longtable}

\textbf{对应 C++}

\begin{itemize}
\tightlist
\item
  类：\passthrough{\lstinline!unitree\_go::msg::dds\_::LowState\_!}
\item
  字段类型：\passthrough{\lstinline!float!}
\item
  声明位置：\passthrough{\lstinline!include/unitree/idl/go2/LowState\_.hpp:48!}
\end{itemize}

\textbf{用法}

\begin{lstlisting}[language=Python]
current_value = obj.power_a
obj.power_a = new_value
\end{lstlisting}

\hypertarget{unitree-sdk2-cpp-idl-go2-lowstate-fan-frequency}{%
\paragraph{\texorpdfstring{\texttt{unitree\_sdk2\_cpp.idl.go2.LowState.fan\_frequency}}{unitree\_sdk2\_cpp.idl.go2.LowState.fan\_frequency}}\label{unitree-sdk2-cpp-idl-go2-lowstate-fan-frequency}}

底层 \passthrough{\lstinline!unitree\_go::msg::dds\_::LowState\_!} 的
\passthrough{\lstinline!fan\_frequency!} 字段。类型签名只说明 Python
表示；单位、数值范围、数组固定长度和枚举含义需查对应 IDL 头文件。
底层是固定长度数组，写入序列必须正好包含 4 个元素；元素约束：取值范围为
0 到 65535。

\textbf{签名}

\begin{lstlisting}[language=Python]
@property
def fan_frequency(self) -> list[int]

@fan_frequency.setter
def fan_frequency(self, value: Sequence[int]) -> None
\end{lstlisting}

\textbf{可用性与安全性}

\textbf{\passthrough{\lstinline!AVAILABLE!} /
\passthrough{\lstinline!VALUE\_TYPE!}}：当前绑定源码已实现该消息值操作。

\textbf{参数}

\begin{longtable}[]{@{}
  >{\raggedright\arraybackslash}p{(\columnwidth - 4\tabcolsep) * \real{0.18}}
  >{\raggedright\arraybackslash}p{(\columnwidth - 4\tabcolsep) * \real{0.20}}
  >{\raggedright\arraybackslash}p{(\columnwidth - 4\tabcolsep) * \real{0.62}}@{}}
\toprule
名称 & Python 类型 & 含义 \\
\midrule
\endhead
\passthrough{\lstinline!value!} &
\passthrough{\lstinline!Sequence[int]!} &
要写入或传递的新值，必须符合该参数的 Python 类型以及底层 C++
范围约束。 \\
\bottomrule
\end{longtable}

\textbf{返回值}

\begin{longtable}[]{@{}
  >{\raggedright\arraybackslash}p{(\columnwidth - 4\tabcolsep) * \real{0.18}}
  >{\raggedright\arraybackslash}p{(\columnwidth - 4\tabcolsep) * \real{0.20}}
  >{\raggedright\arraybackslash}p{(\columnwidth - 4\tabcolsep) * \real{0.62}}@{}}
\toprule
位置 & 类型 & 含义 \\
\midrule
\endhead
getter 返回值 & \passthrough{\lstinline!list[int]!} &
返回属性当前值。数组、vector 和嵌套消息采用复制语义。 \\
\bottomrule
\end{longtable}

\textbf{对应 C++}

\begin{itemize}
\tightlist
\item
  类：\passthrough{\lstinline!unitree\_go::msg::dds\_::LowState\_!}
\item
  字段类型：\passthrough{\lstinline!std::array<uint16\_t, 4>!}
\item
  声明位置：\passthrough{\lstinline!include/unitree/idl/go2/LowState\_.hpp:49!}
\end{itemize}

\textbf{用法}

\begin{lstlisting}[language=Python]
current_value = obj.fan_frequency
obj.fan_frequency = new_value
\end{lstlisting}

\begin{quote}
{[}!TIP{]} 读取容器或嵌套值后应遵循``读取、修改、重新赋值''。只修改
getter 返回的副本不会可靠地更新原 C++ 消息。
\end{quote}

\hypertarget{unitree-sdk2-cpp-idl-go2-lowstate-reserve}{%
\paragraph{\texorpdfstring{\texttt{unitree\_sdk2\_cpp.idl.go2.LowState.reserve}}{unitree\_sdk2\_cpp.idl.go2.LowState.reserve}}\label{unitree-sdk2-cpp-idl-go2-lowstate-reserve}}

底层 \passthrough{\lstinline!unitree\_go::msg::dds\_::LowState\_!} 的
\passthrough{\lstinline!reserve!} 字段。类型签名只说明 Python
表示；单位、数值范围、数组固定长度和枚举含义需查对应 IDL 头文件。
取值范围为 0 到 4294967295。

\textbf{签名}

\begin{lstlisting}[language=Python]
@property
def reserve(self) -> int

@reserve.setter
def reserve(self, value: int) -> None
\end{lstlisting}

\textbf{可用性与安全性}

\textbf{\passthrough{\lstinline!AVAILABLE!} /
\passthrough{\lstinline!VALUE\_TYPE!}}：当前绑定源码已实现该消息值操作。

\textbf{参数}

\begin{longtable}[]{@{}
  >{\raggedright\arraybackslash}p{(\columnwidth - 4\tabcolsep) * \real{0.18}}
  >{\raggedright\arraybackslash}p{(\columnwidth - 4\tabcolsep) * \real{0.20}}
  >{\raggedright\arraybackslash}p{(\columnwidth - 4\tabcolsep) * \real{0.62}}@{}}
\toprule
名称 & Python 类型 & 含义 \\
\midrule
\endhead
\passthrough{\lstinline!value!} & \passthrough{\lstinline!int!} &
要写入或传递的新值，必须符合该参数的 Python 类型以及底层 C++
范围约束。 \\
\bottomrule
\end{longtable}

\textbf{返回值}

\begin{longtable}[]{@{}
  >{\raggedright\arraybackslash}p{(\columnwidth - 4\tabcolsep) * \real{0.18}}
  >{\raggedright\arraybackslash}p{(\columnwidth - 4\tabcolsep) * \real{0.20}}
  >{\raggedright\arraybackslash}p{(\columnwidth - 4\tabcolsep) * \real{0.62}}@{}}
\toprule
位置 & 类型 & 含义 \\
\midrule
\endhead
getter 返回值 & \passthrough{\lstinline!int!} & 返回属性当前值。 \\
\bottomrule
\end{longtable}

\textbf{对应 C++}

\begin{itemize}
\tightlist
\item
  类：\passthrough{\lstinline!unitree\_go::msg::dds\_::LowState\_!}
\item
  字段类型：\passthrough{\lstinline!uint32\_t!}
\item
  声明位置：\passthrough{\lstinline!include/unitree/idl/go2/LowState\_.hpp:50!}
\end{itemize}

\textbf{用法}

\begin{lstlisting}[language=Python]
current_value = obj.reserve
obj.reserve = new_value
\end{lstlisting}

\hypertarget{unitree-sdk2-cpp-idl-go2-lowstate-crc}{%
\paragraph{\texorpdfstring{\texttt{unitree\_sdk2\_cpp.idl.go2.LowState.crc}}{unitree\_sdk2\_cpp.idl.go2.LowState.crc}}\label{unitree-sdk2-cpp-idl-go2-lowstate-crc}}

底层 \passthrough{\lstinline!unitree\_go::msg::dds\_::LowState\_!} 的
\passthrough{\lstinline!crc!} 字段。类型签名只说明 Python
表示；单位、数值范围、数组固定长度和枚举含义需查对应 IDL 头文件。
取值范围为 0 到 4294967295。

\textbf{签名}

\begin{lstlisting}[language=Python]
@property
def crc(self) -> int

@crc.setter
def crc(self, value: int) -> None
\end{lstlisting}

\textbf{可用性与安全性}

\textbf{\passthrough{\lstinline!AVAILABLE!} /
\passthrough{\lstinline!VALUE\_TYPE!}}：当前绑定源码已实现该消息值操作。

\textbf{参数}

\begin{longtable}[]{@{}
  >{\raggedright\arraybackslash}p{(\columnwidth - 4\tabcolsep) * \real{0.18}}
  >{\raggedright\arraybackslash}p{(\columnwidth - 4\tabcolsep) * \real{0.20}}
  >{\raggedright\arraybackslash}p{(\columnwidth - 4\tabcolsep) * \real{0.62}}@{}}
\toprule
名称 & Python 类型 & 含义 \\
\midrule
\endhead
\passthrough{\lstinline!value!} & \passthrough{\lstinline!int!} &
要写入或传递的新值，必须符合该参数的 Python 类型以及底层 C++
范围约束。 \\
\bottomrule
\end{longtable}

\textbf{返回值}

\begin{longtable}[]{@{}
  >{\raggedright\arraybackslash}p{(\columnwidth - 4\tabcolsep) * \real{0.18}}
  >{\raggedright\arraybackslash}p{(\columnwidth - 4\tabcolsep) * \real{0.20}}
  >{\raggedright\arraybackslash}p{(\columnwidth - 4\tabcolsep) * \real{0.62}}@{}}
\toprule
位置 & 类型 & 含义 \\
\midrule
\endhead
getter 返回值 & \passthrough{\lstinline!int!} & 返回属性当前值。 \\
\bottomrule
\end{longtable}

\textbf{对应 C++}

\begin{itemize}
\tightlist
\item
  类：\passthrough{\lstinline!unitree\_go::msg::dds\_::LowState\_!}
\item
  字段类型：\passthrough{\lstinline!uint32\_t!}
\item
  声明位置：\passthrough{\lstinline!include/unitree/idl/go2/LowState\_.hpp:51!}
\end{itemize}

\textbf{用法}

\begin{lstlisting}[language=Python]
current_value = obj.crc
obj.crc = new_value
\end{lstlisting}

\hypertarget{unitree-sdk2-cpp-idl-go2-motorcmds}{%
\subsubsection{\texorpdfstring{\texttt{unitree\_sdk2\_cpp.idl.go2.MotorCmds}}{unitree\_sdk2\_cpp.idl.go2.MotorCmds}}\label{unitree-sdk2-cpp-idl-go2-motorcmds}}

可由 Python 读写的 DDS/IDL 消息值类型。字段采用复制语义。

\textbf{导入}

\begin{lstlisting}[language=Python]
from unitree_sdk2_cpp.idl.go2 import MotorCmds
\end{lstlisting}

\textbf{方法索引}

\begin{longtable}[]{@{}
  >{\raggedright\arraybackslash}p{(\columnwidth - 6\tabcolsep) * \real{0.18}}
  >{\raggedright\arraybackslash}p{(\columnwidth - 6\tabcolsep) * \real{0.24}}
  >{\raggedright\arraybackslash}p{(\columnwidth - 6\tabcolsep) * \real{0.18}}
  >{\raggedright\arraybackslash}p{(\columnwidth - 6\tabcolsep) * \real{0.40}}@{}}
\toprule
方法 & Python 签名 & 状态 & 安全分类 \\
\midrule
\endhead
\protect\hyperlink{unitree-sdk2-cpp-idl-go2-motorcmds-dunder-init-1}{\passthrough{\lstinline!\_\_init\_\_!}}
& \passthrough{\lstinline!def \_\_init\_\_(self) -> None!} &
\passthrough{\lstinline!AVAILABLE!} &
\passthrough{\lstinline!VALUE\_TYPE!} \\
\protect\hyperlink{unitree-sdk2-cpp-idl-go2-motorcmds-dunder-eq-1}{\passthrough{\lstinline!\_\_eq\_\_!}}
& \passthrough{\lstinline!def \_\_eq\_\_(self, other: object) -> bool!}
& \passthrough{\lstinline!AVAILABLE!} &
\passthrough{\lstinline!VALUE\_TYPE!} \\
\protect\hyperlink{unitree-sdk2-cpp-idl-go2-motorcmds-dunder-ne-1}{\passthrough{\lstinline!\_\_ne\_\_!}}
& \passthrough{\lstinline!def \_\_ne\_\_(self, other: object) -> bool!}
& \passthrough{\lstinline!AVAILABLE!} &
\passthrough{\lstinline!VALUE\_TYPE!} \\
\bottomrule
\end{longtable}

\textbf{属性索引}

\begin{longtable}[]{@{}
  >{\raggedright\arraybackslash}p{(\columnwidth - 6\tabcolsep) * \real{0.18}}
  >{\raggedright\arraybackslash}p{(\columnwidth - 6\tabcolsep) * \real{0.24}}
  >{\raggedright\arraybackslash}p{(\columnwidth - 6\tabcolsep) * \real{0.18}}
  >{\raggedright\arraybackslash}p{(\columnwidth - 6\tabcolsep) * \real{0.40}}@{}}
\toprule
属性 & 读取类型 & 写入类型 & 状态 \\
\midrule
\endhead
\protect\hyperlink{unitree-sdk2-cpp-idl-go2-motorcmds-cmds}{\passthrough{\lstinline!cmds!}}
& \passthrough{\lstinline!list[MotorCmd]!} &
\passthrough{\lstinline!Sequence[MotorCmd]!} &
\passthrough{\lstinline!AVAILABLE!} \\
\bottomrule
\end{longtable}

\hypertarget{unitree-sdk2-cpp-idl-go2-motorcmds-dunder-init-1}{%
\paragraph{\texorpdfstring{\texttt{unitree\_sdk2\_cpp.idl.go2.MotorCmds.\_\_init\_\_}}{unitree\_sdk2\_cpp.idl.go2.MotorCmds.\_\_init\_\_}}\label{unitree-sdk2-cpp-idl-go2-motorcmds-dunder-init-1}}

初始化 \passthrough{\lstinline!MotorCmds!}
实例。是否能实际构造取决于下方可用性状态。

\textbf{签名}

\begin{lstlisting}[language=Python]
def __init__(self) -> None
\end{lstlisting}

\textbf{可用性与安全性}

\textbf{\passthrough{\lstinline!AVAILABLE!} /
\passthrough{\lstinline!VALUE\_TYPE!}}：当前绑定源码已实现该消息值操作。

\textbf{参数}

无显式参数。实例方法中的 \passthrough{\lstinline!self!} 由 Python
自动传入。

\textbf{返回值}

\begin{longtable}[]{@{}
  >{\raggedright\arraybackslash}p{(\columnwidth - 4\tabcolsep) * \real{0.18}}
  >{\raggedright\arraybackslash}p{(\columnwidth - 4\tabcolsep) * \real{0.20}}
  >{\raggedright\arraybackslash}p{(\columnwidth - 4\tabcolsep) * \real{0.62}}@{}}
\toprule
位置 & 类型 & 含义 \\
\midrule
\endhead
无 & \passthrough{\lstinline!None!} &
\passthrough{\lstinline!\_\_init\_\_!}
本身不返回值；调用类对象时，在构造函数可用的前提下得到该类实例。 \\
\bottomrule
\end{longtable}

\textbf{对应 C++}

\begin{itemize}
\tightlist
\item
  类：\passthrough{\lstinline!unitree\_go::msg::dds\_::MotorCmds\_!}
\item
  签名：\passthrough{\lstinline!MotorCmds\_()!}
\item
  绑定策略：\passthrough{\lstinline!未单独标注!}
\item
  声明位置：\passthrough{\lstinline!include/unitree/idl/go2/MotorCmds\_.hpp:28!}
\end{itemize}

\textbf{说明}

\begin{itemize}
\tightlist
\item
  示例是局部调用片段；\passthrough{\lstinline!obj!}
  和参数变量表示已经按上层业务校验并准备好的对象。
\end{itemize}

\textbf{用法}

\begin{lstlisting}[language=Python]
value = MotorCmds()
\end{lstlisting}

\hypertarget{unitree-sdk2-cpp-idl-go2-motorcmds-dunder-eq-1}{%
\paragraph{\texorpdfstring{\texttt{unitree\_sdk2\_cpp.idl.go2.MotorCmds.\_\_eq\_\_}}{unitree\_sdk2\_cpp.idl.go2.MotorCmds.\_\_eq\_\_}}\label{unitree-sdk2-cpp-idl-go2-motorcmds-dunder-eq-1}}

按底层 IDL 消息内容比较两个对象是否相等。

\textbf{签名}

\begin{lstlisting}[language=Python]
def __eq__(self, other: object) -> bool
\end{lstlisting}

\textbf{可用性与安全性}

\textbf{\passthrough{\lstinline!AVAILABLE!} /
\passthrough{\lstinline!VALUE\_TYPE!}}：当前绑定源码已实现该消息值操作。

\textbf{参数}

\begin{longtable}[]{@{}
  >{\raggedright\arraybackslash}p{(\columnwidth - 6\tabcolsep) * \real{0.18}}
  >{\raggedright\arraybackslash}p{(\columnwidth - 6\tabcolsep) * \real{0.24}}
  >{\raggedright\arraybackslash}p{(\columnwidth - 6\tabcolsep) * \real{0.18}}
  >{\raggedright\arraybackslash}p{(\columnwidth - 6\tabcolsep) * \real{0.40}}@{}}
\toprule
名称 & Python 类型 & 默认值 & 含义 \\
\midrule
\endhead
\passthrough{\lstinline!other!} & \passthrough{\lstinline!object!} &
必填 & 用于比较的另一个 Python 对象。类型不兼容时比较结果通常为
\passthrough{\lstinline!False!}。 \\
\bottomrule
\end{longtable}

\textbf{返回值}

\begin{longtable}[]{@{}
  >{\raggedright\arraybackslash}p{(\columnwidth - 4\tabcolsep) * \real{0.18}}
  >{\raggedright\arraybackslash}p{(\columnwidth - 4\tabcolsep) * \real{0.20}}
  >{\raggedright\arraybackslash}p{(\columnwidth - 4\tabcolsep) * \real{0.62}}@{}}
\toprule
位置 & 类型 & 含义 \\
\midrule
\endhead
返回值 & \passthrough{\lstinline!bool!} & 返回
\passthrough{\lstinline!bool!}。更精确的业务含义见该方法用途、C++
签名和目标型号协议。 \\
\bottomrule
\end{longtable}

\textbf{对应 C++}

\begin{itemize}
\tightlist
\item
  类：\passthrough{\lstinline!unitree\_go::msg::dds\_::MotorCmds\_!}
\item
  签名：\passthrough{\lstinline!operator==(const value \&) const!}
\item
  绑定策略：\passthrough{\lstinline!未单独标注!}
\item
  声明位置：由 pybind11 手工绑定或生成注册表提供；manifest
  未记录独立头文件行号。
\end{itemize}

\textbf{说明}

\begin{itemize}
\tightlist
\item
  示例是局部调用片段；\passthrough{\lstinline!obj!}
  和参数变量表示已经按上层业务校验并准备好的对象。
\end{itemize}

\textbf{用法}

\begin{lstlisting}[language=Python]
same = left == right
\end{lstlisting}

\hypertarget{unitree-sdk2-cpp-idl-go2-motorcmds-dunder-ne-1}{%
\paragraph{\texorpdfstring{\texttt{unitree\_sdk2\_cpp.idl.go2.MotorCmds.\_\_ne\_\_}}{unitree\_sdk2\_cpp.idl.go2.MotorCmds.\_\_ne\_\_}}\label{unitree-sdk2-cpp-idl-go2-motorcmds-dunder-ne-1}}

按底层 IDL 消息内容比较两个对象是否不相等。

\textbf{签名}

\begin{lstlisting}[language=Python]
def __ne__(self, other: object) -> bool
\end{lstlisting}

\textbf{可用性与安全性}

\textbf{\passthrough{\lstinline!AVAILABLE!} /
\passthrough{\lstinline!VALUE\_TYPE!}}：当前绑定源码已实现该消息值操作。

\textbf{参数}

\begin{longtable}[]{@{}
  >{\raggedright\arraybackslash}p{(\columnwidth - 6\tabcolsep) * \real{0.18}}
  >{\raggedright\arraybackslash}p{(\columnwidth - 6\tabcolsep) * \real{0.24}}
  >{\raggedright\arraybackslash}p{(\columnwidth - 6\tabcolsep) * \real{0.18}}
  >{\raggedright\arraybackslash}p{(\columnwidth - 6\tabcolsep) * \real{0.40}}@{}}
\toprule
名称 & Python 类型 & 默认值 & 含义 \\
\midrule
\endhead
\passthrough{\lstinline!other!} & \passthrough{\lstinline!object!} &
必填 & 用于比较的另一个 Python 对象。类型不兼容时比较结果通常为
\passthrough{\lstinline!False!}。 \\
\bottomrule
\end{longtable}

\textbf{返回值}

\begin{longtable}[]{@{}
  >{\raggedright\arraybackslash}p{(\columnwidth - 4\tabcolsep) * \real{0.18}}
  >{\raggedright\arraybackslash}p{(\columnwidth - 4\tabcolsep) * \real{0.20}}
  >{\raggedright\arraybackslash}p{(\columnwidth - 4\tabcolsep) * \real{0.62}}@{}}
\toprule
位置 & 类型 & 含义 \\
\midrule
\endhead
返回值 & \passthrough{\lstinline!bool!} & 返回
\passthrough{\lstinline!bool!}。更精确的业务含义见该方法用途、C++
签名和目标型号协议。 \\
\bottomrule
\end{longtable}

\textbf{对应 C++}

\begin{itemize}
\tightlist
\item
  类：\passthrough{\lstinline!unitree\_go::msg::dds\_::MotorCmds\_!}
\item
  签名：\passthrough{\lstinline"operator!=(const value \&) const"}
\item
  绑定策略：\passthrough{\lstinline!未单独标注!}
\item
  声明位置：由 pybind11 手工绑定或生成注册表提供；manifest
  未记录独立头文件行号。
\end{itemize}

\textbf{说明}

\begin{itemize}
\tightlist
\item
  示例是局部调用片段；\passthrough{\lstinline!obj!}
  和参数变量表示已经按上层业务校验并准备好的对象。
\end{itemize}

\textbf{用法}

\begin{lstlisting}[language=Python]
different = left != right
\end{lstlisting}

\hypertarget{unitree-sdk2-cpp-idl-go2-motorcmds-cmds}{%
\paragraph{\texorpdfstring{\texttt{unitree\_sdk2\_cpp.idl.go2.MotorCmds.cmds}}{unitree\_sdk2\_cpp.idl.go2.MotorCmds.cmds}}\label{unitree-sdk2-cpp-idl-go2-motorcmds-cmds}}

底层 \passthrough{\lstinline!unitree\_go::msg::dds\_::MotorCmds\_!} 的
\passthrough{\lstinline!cmds!} 字段。类型签名只说明 Python
表示；单位、数值范围、数组固定长度和枚举含义需查对应 IDL 头文件。
底层是可变长度 vector

\textbf{签名}

\begin{lstlisting}[language=Python]
@property
def cmds(self) -> list[MotorCmd]

@cmds.setter
def cmds(self, value: Sequence[MotorCmd]) -> None
\end{lstlisting}

\textbf{可用性与安全性}

\textbf{\passthrough{\lstinline!AVAILABLE!} /
\passthrough{\lstinline!VALUE\_TYPE!}}：当前绑定源码已实现该消息值操作。

\textbf{参数}

\begin{longtable}[]{@{}
  >{\raggedright\arraybackslash}p{(\columnwidth - 4\tabcolsep) * \real{0.18}}
  >{\raggedright\arraybackslash}p{(\columnwidth - 4\tabcolsep) * \real{0.20}}
  >{\raggedright\arraybackslash}p{(\columnwidth - 4\tabcolsep) * \real{0.62}}@{}}
\toprule
名称 & Python 类型 & 含义 \\
\midrule
\endhead
\passthrough{\lstinline!value!} &
\passthrough{\lstinline!Sequence[MotorCmd]!} &
要写入或传递的新值，必须符合该参数的 Python 类型以及底层 C++
范围约束。 \\
\bottomrule
\end{longtable}

\textbf{返回值}

\begin{longtable}[]{@{}
  >{\raggedright\arraybackslash}p{(\columnwidth - 4\tabcolsep) * \real{0.18}}
  >{\raggedright\arraybackslash}p{(\columnwidth - 4\tabcolsep) * \real{0.20}}
  >{\raggedright\arraybackslash}p{(\columnwidth - 4\tabcolsep) * \real{0.62}}@{}}
\toprule
位置 & 类型 & 含义 \\
\midrule
\endhead
getter 返回值 & \passthrough{\lstinline!list[MotorCmd]!} &
返回属性当前值。数组、vector 和嵌套消息采用复制语义。 \\
\bottomrule
\end{longtable}

\textbf{对应 C++}

\begin{itemize}
\tightlist
\item
  类：\passthrough{\lstinline!unitree\_go::msg::dds\_::MotorCmds\_!}
\item
  字段类型：\passthrough{\lstinline!std::vector< ::unitree\_go::msg::dds\_::MotorCmd\_>!}
\item
  声明位置：\passthrough{\lstinline!include/unitree/idl/go2/MotorCmds\_.hpp:25!}
\end{itemize}

\textbf{用法}

\begin{lstlisting}[language=Python]
current_value = obj.cmds
obj.cmds = new_value
\end{lstlisting}

\begin{quote}
{[}!TIP{]} 读取容器或嵌套值后应遵循``读取、修改、重新赋值''。只修改
getter 返回的副本不会可靠地更新原 C++ 消息。
\end{quote}

\hypertarget{unitree-sdk2-cpp-idl-go2-motorstates}{%
\subsubsection{\texorpdfstring{\texttt{unitree\_sdk2\_cpp.idl.go2.MotorStates}}{unitree\_sdk2\_cpp.idl.go2.MotorStates}}\label{unitree-sdk2-cpp-idl-go2-motorstates}}

可由 Python 读写的 DDS/IDL 消息值类型。字段采用复制语义。

\textbf{导入}

\begin{lstlisting}[language=Python]
from unitree_sdk2_cpp.idl.go2 import MotorStates
\end{lstlisting}

\textbf{方法索引}

\begin{longtable}[]{@{}
  >{\raggedright\arraybackslash}p{(\columnwidth - 6\tabcolsep) * \real{0.18}}
  >{\raggedright\arraybackslash}p{(\columnwidth - 6\tabcolsep) * \real{0.24}}
  >{\raggedright\arraybackslash}p{(\columnwidth - 6\tabcolsep) * \real{0.18}}
  >{\raggedright\arraybackslash}p{(\columnwidth - 6\tabcolsep) * \real{0.40}}@{}}
\toprule
方法 & Python 签名 & 状态 & 安全分类 \\
\midrule
\endhead
\protect\hyperlink{unitree-sdk2-cpp-idl-go2-motorstates-dunder-init-1}{\passthrough{\lstinline!\_\_init\_\_!}}
& \passthrough{\lstinline!def \_\_init\_\_(self) -> None!} &
\passthrough{\lstinline!AVAILABLE!} &
\passthrough{\lstinline!VALUE\_TYPE!} \\
\protect\hyperlink{unitree-sdk2-cpp-idl-go2-motorstates-dunder-eq-1}{\passthrough{\lstinline!\_\_eq\_\_!}}
& \passthrough{\lstinline!def \_\_eq\_\_(self, other: object) -> bool!}
& \passthrough{\lstinline!AVAILABLE!} &
\passthrough{\lstinline!VALUE\_TYPE!} \\
\protect\hyperlink{unitree-sdk2-cpp-idl-go2-motorstates-dunder-ne-1}{\passthrough{\lstinline!\_\_ne\_\_!}}
& \passthrough{\lstinline!def \_\_ne\_\_(self, other: object) -> bool!}
& \passthrough{\lstinline!AVAILABLE!} &
\passthrough{\lstinline!VALUE\_TYPE!} \\
\bottomrule
\end{longtable}

\textbf{属性索引}

\begin{longtable}[]{@{}
  >{\raggedright\arraybackslash}p{(\columnwidth - 6\tabcolsep) * \real{0.18}}
  >{\raggedright\arraybackslash}p{(\columnwidth - 6\tabcolsep) * \real{0.24}}
  >{\raggedright\arraybackslash}p{(\columnwidth - 6\tabcolsep) * \real{0.18}}
  >{\raggedright\arraybackslash}p{(\columnwidth - 6\tabcolsep) * \real{0.40}}@{}}
\toprule
属性 & 读取类型 & 写入类型 & 状态 \\
\midrule
\endhead
\protect\hyperlink{unitree-sdk2-cpp-idl-go2-motorstates-states}{\passthrough{\lstinline!states!}}
& \passthrough{\lstinline!list[MotorState]!} &
\passthrough{\lstinline!Sequence[MotorState]!} &
\passthrough{\lstinline!AVAILABLE!} \\
\bottomrule
\end{longtable}

\hypertarget{unitree-sdk2-cpp-idl-go2-motorstates-dunder-init-1}{%
\paragraph{\texorpdfstring{\texttt{unitree\_sdk2\_cpp.idl.go2.MotorStates.\_\_init\_\_}}{unitree\_sdk2\_cpp.idl.go2.MotorStates.\_\_init\_\_}}\label{unitree-sdk2-cpp-idl-go2-motorstates-dunder-init-1}}

初始化 \passthrough{\lstinline!MotorStates!}
实例。是否能实际构造取决于下方可用性状态。

\textbf{签名}

\begin{lstlisting}[language=Python]
def __init__(self) -> None
\end{lstlisting}

\textbf{可用性与安全性}

\textbf{\passthrough{\lstinline!AVAILABLE!} /
\passthrough{\lstinline!VALUE\_TYPE!}}：当前绑定源码已实现该消息值操作。

\textbf{参数}

无显式参数。实例方法中的 \passthrough{\lstinline!self!} 由 Python
自动传入。

\textbf{返回值}

\begin{longtable}[]{@{}
  >{\raggedright\arraybackslash}p{(\columnwidth - 4\tabcolsep) * \real{0.18}}
  >{\raggedright\arraybackslash}p{(\columnwidth - 4\tabcolsep) * \real{0.20}}
  >{\raggedright\arraybackslash}p{(\columnwidth - 4\tabcolsep) * \real{0.62}}@{}}
\toprule
位置 & 类型 & 含义 \\
\midrule
\endhead
无 & \passthrough{\lstinline!None!} &
\passthrough{\lstinline!\_\_init\_\_!}
本身不返回值；调用类对象时，在构造函数可用的前提下得到该类实例。 \\
\bottomrule
\end{longtable}

\textbf{对应 C++}

\begin{itemize}
\tightlist
\item
  类：\passthrough{\lstinline!unitree\_go::msg::dds\_::MotorStates\_!}
\item
  签名：\passthrough{\lstinline!MotorStates\_()!}
\item
  绑定策略：\passthrough{\lstinline!未单独标注!}
\item
  声明位置：\passthrough{\lstinline!include/unitree/idl/go2/MotorStates\_.hpp:28!}
\end{itemize}

\textbf{说明}

\begin{itemize}
\tightlist
\item
  示例是局部调用片段；\passthrough{\lstinline!obj!}
  和参数变量表示已经按上层业务校验并准备好的对象。
\end{itemize}

\textbf{用法}

\begin{lstlisting}[language=Python]
value = MotorStates()
\end{lstlisting}

\hypertarget{unitree-sdk2-cpp-idl-go2-motorstates-dunder-eq-1}{%
\paragraph{\texorpdfstring{\texttt{unitree\_sdk2\_cpp.idl.go2.MotorStates.\_\_eq\_\_}}{unitree\_sdk2\_cpp.idl.go2.MotorStates.\_\_eq\_\_}}\label{unitree-sdk2-cpp-idl-go2-motorstates-dunder-eq-1}}

按底层 IDL 消息内容比较两个对象是否相等。

\textbf{签名}

\begin{lstlisting}[language=Python]
def __eq__(self, other: object) -> bool
\end{lstlisting}

\textbf{可用性与安全性}

\textbf{\passthrough{\lstinline!AVAILABLE!} /
\passthrough{\lstinline!VALUE\_TYPE!}}：当前绑定源码已实现该消息值操作。

\textbf{参数}

\begin{longtable}[]{@{}
  >{\raggedright\arraybackslash}p{(\columnwidth - 6\tabcolsep) * \real{0.18}}
  >{\raggedright\arraybackslash}p{(\columnwidth - 6\tabcolsep) * \real{0.24}}
  >{\raggedright\arraybackslash}p{(\columnwidth - 6\tabcolsep) * \real{0.18}}
  >{\raggedright\arraybackslash}p{(\columnwidth - 6\tabcolsep) * \real{0.40}}@{}}
\toprule
名称 & Python 类型 & 默认值 & 含义 \\
\midrule
\endhead
\passthrough{\lstinline!other!} & \passthrough{\lstinline!object!} &
必填 & 用于比较的另一个 Python 对象。类型不兼容时比较结果通常为
\passthrough{\lstinline!False!}。 \\
\bottomrule
\end{longtable}

\textbf{返回值}

\begin{longtable}[]{@{}
  >{\raggedright\arraybackslash}p{(\columnwidth - 4\tabcolsep) * \real{0.18}}
  >{\raggedright\arraybackslash}p{(\columnwidth - 4\tabcolsep) * \real{0.20}}
  >{\raggedright\arraybackslash}p{(\columnwidth - 4\tabcolsep) * \real{0.62}}@{}}
\toprule
位置 & 类型 & 含义 \\
\midrule
\endhead
返回值 & \passthrough{\lstinline!bool!} & 返回
\passthrough{\lstinline!bool!}。更精确的业务含义见该方法用途、C++
签名和目标型号协议。 \\
\bottomrule
\end{longtable}

\textbf{对应 C++}

\begin{itemize}
\tightlist
\item
  类：\passthrough{\lstinline!unitree\_go::msg::dds\_::MotorStates\_!}
\item
  签名：\passthrough{\lstinline!operator==(const value \&) const!}
\item
  绑定策略：\passthrough{\lstinline!未单独标注!}
\item
  声明位置：由 pybind11 手工绑定或生成注册表提供；manifest
  未记录独立头文件行号。
\end{itemize}

\textbf{说明}

\begin{itemize}
\tightlist
\item
  示例是局部调用片段；\passthrough{\lstinline!obj!}
  和参数变量表示已经按上层业务校验并准备好的对象。
\end{itemize}

\textbf{用法}

\begin{lstlisting}[language=Python]
same = left == right
\end{lstlisting}

\hypertarget{unitree-sdk2-cpp-idl-go2-motorstates-dunder-ne-1}{%
\paragraph{\texorpdfstring{\texttt{unitree\_sdk2\_cpp.idl.go2.MotorStates.\_\_ne\_\_}}{unitree\_sdk2\_cpp.idl.go2.MotorStates.\_\_ne\_\_}}\label{unitree-sdk2-cpp-idl-go2-motorstates-dunder-ne-1}}

按底层 IDL 消息内容比较两个对象是否不相等。

\textbf{签名}

\begin{lstlisting}[language=Python]
def __ne__(self, other: object) -> bool
\end{lstlisting}

\textbf{可用性与安全性}

\textbf{\passthrough{\lstinline!AVAILABLE!} /
\passthrough{\lstinline!VALUE\_TYPE!}}：当前绑定源码已实现该消息值操作。

\textbf{参数}

\begin{longtable}[]{@{}
  >{\raggedright\arraybackslash}p{(\columnwidth - 6\tabcolsep) * \real{0.18}}
  >{\raggedright\arraybackslash}p{(\columnwidth - 6\tabcolsep) * \real{0.24}}
  >{\raggedright\arraybackslash}p{(\columnwidth - 6\tabcolsep) * \real{0.18}}
  >{\raggedright\arraybackslash}p{(\columnwidth - 6\tabcolsep) * \real{0.40}}@{}}
\toprule
名称 & Python 类型 & 默认值 & 含义 \\
\midrule
\endhead
\passthrough{\lstinline!other!} & \passthrough{\lstinline!object!} &
必填 & 用于比较的另一个 Python 对象。类型不兼容时比较结果通常为
\passthrough{\lstinline!False!}。 \\
\bottomrule
\end{longtable}

\textbf{返回值}

\begin{longtable}[]{@{}
  >{\raggedright\arraybackslash}p{(\columnwidth - 4\tabcolsep) * \real{0.18}}
  >{\raggedright\arraybackslash}p{(\columnwidth - 4\tabcolsep) * \real{0.20}}
  >{\raggedright\arraybackslash}p{(\columnwidth - 4\tabcolsep) * \real{0.62}}@{}}
\toprule
位置 & 类型 & 含义 \\
\midrule
\endhead
返回值 & \passthrough{\lstinline!bool!} & 返回
\passthrough{\lstinline!bool!}。更精确的业务含义见该方法用途、C++
签名和目标型号协议。 \\
\bottomrule
\end{longtable}

\textbf{对应 C++}

\begin{itemize}
\tightlist
\item
  类：\passthrough{\lstinline!unitree\_go::msg::dds\_::MotorStates\_!}
\item
  签名：\passthrough{\lstinline"operator!=(const value \&) const"}
\item
  绑定策略：\passthrough{\lstinline!未单独标注!}
\item
  声明位置：由 pybind11 手工绑定或生成注册表提供；manifest
  未记录独立头文件行号。
\end{itemize}

\textbf{说明}

\begin{itemize}
\tightlist
\item
  示例是局部调用片段；\passthrough{\lstinline!obj!}
  和参数变量表示已经按上层业务校验并准备好的对象。
\end{itemize}

\textbf{用法}

\begin{lstlisting}[language=Python]
different = left != right
\end{lstlisting}

\hypertarget{unitree-sdk2-cpp-idl-go2-motorstates-states}{%
\paragraph{\texorpdfstring{\texttt{unitree\_sdk2\_cpp.idl.go2.MotorStates.states}}{unitree\_sdk2\_cpp.idl.go2.MotorStates.states}}\label{unitree-sdk2-cpp-idl-go2-motorstates-states}}

底层 \passthrough{\lstinline!unitree\_go::msg::dds\_::MotorStates\_!} 的
\passthrough{\lstinline!states!} 字段。类型签名只说明 Python
表示；单位、数值范围、数组固定长度和枚举含义需查对应 IDL 头文件。
底层是可变长度 vector

\textbf{签名}

\begin{lstlisting}[language=Python]
@property
def states(self) -> list[MotorState]

@states.setter
def states(self, value: Sequence[MotorState]) -> None
\end{lstlisting}

\textbf{可用性与安全性}

\textbf{\passthrough{\lstinline!AVAILABLE!} /
\passthrough{\lstinline!VALUE\_TYPE!}}：当前绑定源码已实现该消息值操作。

\textbf{参数}

\begin{longtable}[]{@{}
  >{\raggedright\arraybackslash}p{(\columnwidth - 4\tabcolsep) * \real{0.18}}
  >{\raggedright\arraybackslash}p{(\columnwidth - 4\tabcolsep) * \real{0.20}}
  >{\raggedright\arraybackslash}p{(\columnwidth - 4\tabcolsep) * \real{0.62}}@{}}
\toprule
名称 & Python 类型 & 含义 \\
\midrule
\endhead
\passthrough{\lstinline!value!} &
\passthrough{\lstinline!Sequence[MotorState]!} &
要写入或传递的新值，必须符合该参数的 Python 类型以及底层 C++
范围约束。 \\
\bottomrule
\end{longtable}

\textbf{返回值}

\begin{longtable}[]{@{}
  >{\raggedright\arraybackslash}p{(\columnwidth - 4\tabcolsep) * \real{0.18}}
  >{\raggedright\arraybackslash}p{(\columnwidth - 4\tabcolsep) * \real{0.20}}
  >{\raggedright\arraybackslash}p{(\columnwidth - 4\tabcolsep) * \real{0.62}}@{}}
\toprule
位置 & 类型 & 含义 \\
\midrule
\endhead
getter 返回值 & \passthrough{\lstinline!list[MotorState]!} &
返回属性当前值。数组、vector 和嵌套消息采用复制语义。 \\
\bottomrule
\end{longtable}

\textbf{对应 C++}

\begin{itemize}
\tightlist
\item
  类：\passthrough{\lstinline!unitree\_go::msg::dds\_::MotorStates\_!}
\item
  字段类型：\passthrough{\lstinline!std::vector< ::unitree\_go::msg::dds\_::MotorState\_>!}
\item
  声明位置：\passthrough{\lstinline!include/unitree/idl/go2/MotorStates\_.hpp:25!}
\end{itemize}

\textbf{用法}

\begin{lstlisting}[language=Python]
current_value = obj.states
obj.states = new_value
\end{lstlisting}

\begin{quote}
{[}!TIP{]} 读取容器或嵌套值后应遵循``读取、修改、重新赋值''。只修改
getter 返回的副本不会可靠地更新原 C++ 消息。
\end{quote}

\hypertarget{unitree-sdk2-cpp-idl-go2-pathpoint}{%
\subsubsection{\texorpdfstring{\texttt{unitree\_sdk2\_cpp.idl.go2.PathPoint}}{unitree\_sdk2\_cpp.idl.go2.PathPoint}}\label{unitree-sdk2-cpp-idl-go2-pathpoint}}

可由 Python 读写的 DDS/IDL 消息值类型。字段采用复制语义。

\textbf{导入}

\begin{lstlisting}[language=Python]
from unitree_sdk2_cpp.idl.go2 import PathPoint
\end{lstlisting}

\textbf{方法索引}

\begin{longtable}[]{@{}
  >{\raggedright\arraybackslash}p{(\columnwidth - 6\tabcolsep) * \real{0.18}}
  >{\raggedright\arraybackslash}p{(\columnwidth - 6\tabcolsep) * \real{0.24}}
  >{\raggedright\arraybackslash}p{(\columnwidth - 6\tabcolsep) * \real{0.18}}
  >{\raggedright\arraybackslash}p{(\columnwidth - 6\tabcolsep) * \real{0.40}}@{}}
\toprule
方法 & Python 签名 & 状态 & 安全分类 \\
\midrule
\endhead
\protect\hyperlink{unitree-sdk2-cpp-idl-go2-pathpoint-dunder-init-1}{\passthrough{\lstinline!\_\_init\_\_!}}
& \passthrough{\lstinline!def \_\_init\_\_(self) -> None!} &
\passthrough{\lstinline!AVAILABLE!} &
\passthrough{\lstinline!VALUE\_TYPE!} \\
\protect\hyperlink{unitree-sdk2-cpp-idl-go2-pathpoint-dunder-eq-1}{\passthrough{\lstinline!\_\_eq\_\_!}}
& \passthrough{\lstinline!def \_\_eq\_\_(self, other: object) -> bool!}
& \passthrough{\lstinline!AVAILABLE!} &
\passthrough{\lstinline!VALUE\_TYPE!} \\
\protect\hyperlink{unitree-sdk2-cpp-idl-go2-pathpoint-dunder-ne-1}{\passthrough{\lstinline!\_\_ne\_\_!}}
& \passthrough{\lstinline!def \_\_ne\_\_(self, other: object) -> bool!}
& \passthrough{\lstinline!AVAILABLE!} &
\passthrough{\lstinline!VALUE\_TYPE!} \\
\bottomrule
\end{longtable}

\textbf{属性索引}

\begin{longtable}[]{@{}
  >{\raggedright\arraybackslash}p{(\columnwidth - 6\tabcolsep) * \real{0.18}}
  >{\raggedright\arraybackslash}p{(\columnwidth - 6\tabcolsep) * \real{0.24}}
  >{\raggedright\arraybackslash}p{(\columnwidth - 6\tabcolsep) * \real{0.18}}
  >{\raggedright\arraybackslash}p{(\columnwidth - 6\tabcolsep) * \real{0.40}}@{}}
\toprule
属性 & 读取类型 & 写入类型 & 状态 \\
\midrule
\endhead
\protect\hyperlink{unitree-sdk2-cpp-idl-go2-pathpoint-t-from-start}{\passthrough{\lstinline!t\_from\_start!}}
& \passthrough{\lstinline!float!} & \passthrough{\lstinline!float!} &
\passthrough{\lstinline!AVAILABLE!} \\
\protect\hyperlink{unitree-sdk2-cpp-idl-go2-pathpoint-x}{\passthrough{\lstinline!x!}}
& \passthrough{\lstinline!float!} & \passthrough{\lstinline!float!} &
\passthrough{\lstinline!AVAILABLE!} \\
\protect\hyperlink{unitree-sdk2-cpp-idl-go2-pathpoint-y}{\passthrough{\lstinline!y!}}
& \passthrough{\lstinline!float!} & \passthrough{\lstinline!float!} &
\passthrough{\lstinline!AVAILABLE!} \\
\protect\hyperlink{unitree-sdk2-cpp-idl-go2-pathpoint-yaw}{\passthrough{\lstinline!yaw!}}
& \passthrough{\lstinline!float!} & \passthrough{\lstinline!float!} &
\passthrough{\lstinline!AVAILABLE!} \\
\protect\hyperlink{unitree-sdk2-cpp-idl-go2-pathpoint-vx}{\passthrough{\lstinline!vx!}}
& \passthrough{\lstinline!float!} & \passthrough{\lstinline!float!} &
\passthrough{\lstinline!AVAILABLE!} \\
\protect\hyperlink{unitree-sdk2-cpp-idl-go2-pathpoint-vy}{\passthrough{\lstinline!vy!}}
& \passthrough{\lstinline!float!} & \passthrough{\lstinline!float!} &
\passthrough{\lstinline!AVAILABLE!} \\
\protect\hyperlink{unitree-sdk2-cpp-idl-go2-pathpoint-vyaw}{\passthrough{\lstinline!vyaw!}}
& \passthrough{\lstinline!float!} & \passthrough{\lstinline!float!} &
\passthrough{\lstinline!AVAILABLE!} \\
\bottomrule
\end{longtable}

\hypertarget{unitree-sdk2-cpp-idl-go2-pathpoint-dunder-init-1}{%
\paragraph{\texorpdfstring{\texttt{unitree\_sdk2\_cpp.idl.go2.PathPoint.\_\_init\_\_}}{unitree\_sdk2\_cpp.idl.go2.PathPoint.\_\_init\_\_}}\label{unitree-sdk2-cpp-idl-go2-pathpoint-dunder-init-1}}

初始化 \passthrough{\lstinline!PathPoint!}
实例。是否能实际构造取决于下方可用性状态。

\textbf{签名}

\begin{lstlisting}[language=Python]
def __init__(self) -> None
\end{lstlisting}

\textbf{可用性与安全性}

\textbf{\passthrough{\lstinline!AVAILABLE!} /
\passthrough{\lstinline!VALUE\_TYPE!}}：当前绑定源码已实现该消息值操作。

\textbf{参数}

无显式参数。实例方法中的 \passthrough{\lstinline!self!} 由 Python
自动传入。

\textbf{返回值}

\begin{longtable}[]{@{}
  >{\raggedright\arraybackslash}p{(\columnwidth - 4\tabcolsep) * \real{0.18}}
  >{\raggedright\arraybackslash}p{(\columnwidth - 4\tabcolsep) * \real{0.20}}
  >{\raggedright\arraybackslash}p{(\columnwidth - 4\tabcolsep) * \real{0.62}}@{}}
\toprule
位置 & 类型 & 含义 \\
\midrule
\endhead
无 & \passthrough{\lstinline!None!} &
\passthrough{\lstinline!\_\_init\_\_!}
本身不返回值；调用类对象时，在构造函数可用的前提下得到该类实例。 \\
\bottomrule
\end{longtable}

\textbf{对应 C++}

\begin{itemize}
\tightlist
\item
  类：\passthrough{\lstinline!unitree\_go::msg::dds\_::PathPoint\_!}
\item
  签名：\passthrough{\lstinline!PathPoint\_()!}
\item
  绑定策略：\passthrough{\lstinline!未单独标注!}
\item
  声明位置：\passthrough{\lstinline!include/unitree/idl/go2/PathPoint\_.hpp:31!}
\end{itemize}

\textbf{说明}

\begin{itemize}
\tightlist
\item
  示例是局部调用片段；\passthrough{\lstinline!obj!}
  和参数变量表示已经按上层业务校验并准备好的对象。
\end{itemize}

\textbf{用法}

\begin{lstlisting}[language=Python]
value = PathPoint()
\end{lstlisting}

\hypertarget{unitree-sdk2-cpp-idl-go2-pathpoint-dunder-eq-1}{%
\paragraph{\texorpdfstring{\texttt{unitree\_sdk2\_cpp.idl.go2.PathPoint.\_\_eq\_\_}}{unitree\_sdk2\_cpp.idl.go2.PathPoint.\_\_eq\_\_}}\label{unitree-sdk2-cpp-idl-go2-pathpoint-dunder-eq-1}}

按底层 IDL 消息内容比较两个对象是否相等。

\textbf{签名}

\begin{lstlisting}[language=Python]
def __eq__(self, other: object) -> bool
\end{lstlisting}

\textbf{可用性与安全性}

\textbf{\passthrough{\lstinline!AVAILABLE!} /
\passthrough{\lstinline!VALUE\_TYPE!}}：当前绑定源码已实现该消息值操作。

\textbf{参数}

\begin{longtable}[]{@{}
  >{\raggedright\arraybackslash}p{(\columnwidth - 6\tabcolsep) * \real{0.18}}
  >{\raggedright\arraybackslash}p{(\columnwidth - 6\tabcolsep) * \real{0.24}}
  >{\raggedright\arraybackslash}p{(\columnwidth - 6\tabcolsep) * \real{0.18}}
  >{\raggedright\arraybackslash}p{(\columnwidth - 6\tabcolsep) * \real{0.40}}@{}}
\toprule
名称 & Python 类型 & 默认值 & 含义 \\
\midrule
\endhead
\passthrough{\lstinline!other!} & \passthrough{\lstinline!object!} &
必填 & 用于比较的另一个 Python 对象。类型不兼容时比较结果通常为
\passthrough{\lstinline!False!}。 \\
\bottomrule
\end{longtable}

\textbf{返回值}

\begin{longtable}[]{@{}
  >{\raggedright\arraybackslash}p{(\columnwidth - 4\tabcolsep) * \real{0.18}}
  >{\raggedright\arraybackslash}p{(\columnwidth - 4\tabcolsep) * \real{0.20}}
  >{\raggedright\arraybackslash}p{(\columnwidth - 4\tabcolsep) * \real{0.62}}@{}}
\toprule
位置 & 类型 & 含义 \\
\midrule
\endhead
返回值 & \passthrough{\lstinline!bool!} & 返回
\passthrough{\lstinline!bool!}。更精确的业务含义见该方法用途、C++
签名和目标型号协议。 \\
\bottomrule
\end{longtable}

\textbf{对应 C++}

\begin{itemize}
\tightlist
\item
  类：\passthrough{\lstinline!unitree\_go::msg::dds\_::PathPoint\_!}
\item
  签名：\passthrough{\lstinline!operator==(const value \&) const!}
\item
  绑定策略：\passthrough{\lstinline!未单独标注!}
\item
  声明位置：由 pybind11 手工绑定或生成注册表提供；manifest
  未记录独立头文件行号。
\end{itemize}

\textbf{说明}

\begin{itemize}
\tightlist
\item
  示例是局部调用片段；\passthrough{\lstinline!obj!}
  和参数变量表示已经按上层业务校验并准备好的对象。
\end{itemize}

\textbf{用法}

\begin{lstlisting}[language=Python]
same = left == right
\end{lstlisting}

\hypertarget{unitree-sdk2-cpp-idl-go2-pathpoint-dunder-ne-1}{%
\paragraph{\texorpdfstring{\texttt{unitree\_sdk2\_cpp.idl.go2.PathPoint.\_\_ne\_\_}}{unitree\_sdk2\_cpp.idl.go2.PathPoint.\_\_ne\_\_}}\label{unitree-sdk2-cpp-idl-go2-pathpoint-dunder-ne-1}}

按底层 IDL 消息内容比较两个对象是否不相等。

\textbf{签名}

\begin{lstlisting}[language=Python]
def __ne__(self, other: object) -> bool
\end{lstlisting}

\textbf{可用性与安全性}

\textbf{\passthrough{\lstinline!AVAILABLE!} /
\passthrough{\lstinline!VALUE\_TYPE!}}：当前绑定源码已实现该消息值操作。

\textbf{参数}

\begin{longtable}[]{@{}
  >{\raggedright\arraybackslash}p{(\columnwidth - 6\tabcolsep) * \real{0.18}}
  >{\raggedright\arraybackslash}p{(\columnwidth - 6\tabcolsep) * \real{0.24}}
  >{\raggedright\arraybackslash}p{(\columnwidth - 6\tabcolsep) * \real{0.18}}
  >{\raggedright\arraybackslash}p{(\columnwidth - 6\tabcolsep) * \real{0.40}}@{}}
\toprule
名称 & Python 类型 & 默认值 & 含义 \\
\midrule
\endhead
\passthrough{\lstinline!other!} & \passthrough{\lstinline!object!} &
必填 & 用于比较的另一个 Python 对象。类型不兼容时比较结果通常为
\passthrough{\lstinline!False!}。 \\
\bottomrule
\end{longtable}

\textbf{返回值}

\begin{longtable}[]{@{}
  >{\raggedright\arraybackslash}p{(\columnwidth - 4\tabcolsep) * \real{0.18}}
  >{\raggedright\arraybackslash}p{(\columnwidth - 4\tabcolsep) * \real{0.20}}
  >{\raggedright\arraybackslash}p{(\columnwidth - 4\tabcolsep) * \real{0.62}}@{}}
\toprule
位置 & 类型 & 含义 \\
\midrule
\endhead
返回值 & \passthrough{\lstinline!bool!} & 返回
\passthrough{\lstinline!bool!}。更精确的业务含义见该方法用途、C++
签名和目标型号协议。 \\
\bottomrule
\end{longtable}

\textbf{对应 C++}

\begin{itemize}
\tightlist
\item
  类：\passthrough{\lstinline!unitree\_go::msg::dds\_::PathPoint\_!}
\item
  签名：\passthrough{\lstinline"operator!=(const value \&) const"}
\item
  绑定策略：\passthrough{\lstinline!未单独标注!}
\item
  声明位置：由 pybind11 手工绑定或生成注册表提供；manifest
  未记录独立头文件行号。
\end{itemize}

\textbf{说明}

\begin{itemize}
\tightlist
\item
  示例是局部调用片段；\passthrough{\lstinline!obj!}
  和参数变量表示已经按上层业务校验并准备好的对象。
\end{itemize}

\textbf{用法}

\begin{lstlisting}[language=Python]
different = left != right
\end{lstlisting}

\hypertarget{unitree-sdk2-cpp-idl-go2-pathpoint-t-from-start}{%
\paragraph{\texorpdfstring{\texttt{unitree\_sdk2\_cpp.idl.go2.PathPoint.t\_from\_start}}{unitree\_sdk2\_cpp.idl.go2.PathPoint.t\_from\_start}}\label{unitree-sdk2-cpp-idl-go2-pathpoint-t-from-start}}

底层 \passthrough{\lstinline!unitree\_go::msg::dds\_::PathPoint\_!} 的
\passthrough{\lstinline!t\_from\_start!} 字段。类型签名只说明 Python
表示；单位、数值范围、数组固定长度和枚举含义需查对应 IDL 头文件。
底层为浮点数；类型本身不说明单位、坐标系或业务安全范围。

\textbf{签名}

\begin{lstlisting}[language=Python]
@property
def t_from_start(self) -> float

@t_from_start.setter
def t_from_start(self, value: float) -> None
\end{lstlisting}

\textbf{可用性与安全性}

\textbf{\passthrough{\lstinline!AVAILABLE!} /
\passthrough{\lstinline!VALUE\_TYPE!}}：当前绑定源码已实现该消息值操作。

\textbf{参数}

\begin{longtable}[]{@{}
  >{\raggedright\arraybackslash}p{(\columnwidth - 4\tabcolsep) * \real{0.18}}
  >{\raggedright\arraybackslash}p{(\columnwidth - 4\tabcolsep) * \real{0.20}}
  >{\raggedright\arraybackslash}p{(\columnwidth - 4\tabcolsep) * \real{0.62}}@{}}
\toprule
名称 & Python 类型 & 含义 \\
\midrule
\endhead
\passthrough{\lstinline!value!} & \passthrough{\lstinline!float!} &
要写入或传递的新值，必须符合该参数的 Python 类型以及底层 C++
范围约束。 \\
\bottomrule
\end{longtable}

\textbf{返回值}

\begin{longtable}[]{@{}
  >{\raggedright\arraybackslash}p{(\columnwidth - 4\tabcolsep) * \real{0.18}}
  >{\raggedright\arraybackslash}p{(\columnwidth - 4\tabcolsep) * \real{0.20}}
  >{\raggedright\arraybackslash}p{(\columnwidth - 4\tabcolsep) * \real{0.62}}@{}}
\toprule
位置 & 类型 & 含义 \\
\midrule
\endhead
getter 返回值 & \passthrough{\lstinline!float!} & 返回属性当前值。 \\
\bottomrule
\end{longtable}

\textbf{对应 C++}

\begin{itemize}
\tightlist
\item
  类：\passthrough{\lstinline!unitree\_go::msg::dds\_::PathPoint\_!}
\item
  字段类型：\passthrough{\lstinline!float!}
\item
  声明位置：\passthrough{\lstinline!include/unitree/idl/go2/PathPoint\_.hpp:22!}
\end{itemize}

\textbf{用法}

\begin{lstlisting}[language=Python]
current_value = obj.t_from_start
obj.t_from_start = new_value
\end{lstlisting}

\hypertarget{unitree-sdk2-cpp-idl-go2-pathpoint-x}{%
\paragraph{\texorpdfstring{\texttt{unitree\_sdk2\_cpp.idl.go2.PathPoint.x}}{unitree\_sdk2\_cpp.idl.go2.PathPoint.x}}\label{unitree-sdk2-cpp-idl-go2-pathpoint-x}}

底层 \passthrough{\lstinline!unitree\_go::msg::dds\_::PathPoint\_!} 的
\passthrough{\lstinline!x!} 字段。类型签名只说明 Python
表示；单位、数值范围、数组固定长度和枚举含义需查对应 IDL 头文件。
底层为浮点数；类型本身不说明单位、坐标系或业务安全范围。

\textbf{签名}

\begin{lstlisting}[language=Python]
@property
def x(self) -> float

@x.setter
def x(self, value: float) -> None
\end{lstlisting}

\textbf{可用性与安全性}

\textbf{\passthrough{\lstinline!AVAILABLE!} /
\passthrough{\lstinline!VALUE\_TYPE!}}：当前绑定源码已实现该消息值操作。

\textbf{参数}

\begin{longtable}[]{@{}
  >{\raggedright\arraybackslash}p{(\columnwidth - 4\tabcolsep) * \real{0.18}}
  >{\raggedright\arraybackslash}p{(\columnwidth - 4\tabcolsep) * \real{0.20}}
  >{\raggedright\arraybackslash}p{(\columnwidth - 4\tabcolsep) * \real{0.62}}@{}}
\toprule
名称 & Python 类型 & 含义 \\
\midrule
\endhead
\passthrough{\lstinline!value!} & \passthrough{\lstinline!float!} &
要写入或传递的新值，必须符合该参数的 Python 类型以及底层 C++
范围约束。 \\
\bottomrule
\end{longtable}

\textbf{返回值}

\begin{longtable}[]{@{}
  >{\raggedright\arraybackslash}p{(\columnwidth - 4\tabcolsep) * \real{0.18}}
  >{\raggedright\arraybackslash}p{(\columnwidth - 4\tabcolsep) * \real{0.20}}
  >{\raggedright\arraybackslash}p{(\columnwidth - 4\tabcolsep) * \real{0.62}}@{}}
\toprule
位置 & 类型 & 含义 \\
\midrule
\endhead
getter 返回值 & \passthrough{\lstinline!float!} & 返回属性当前值。 \\
\bottomrule
\end{longtable}

\textbf{对应 C++}

\begin{itemize}
\tightlist
\item
  类：\passthrough{\lstinline!unitree\_go::msg::dds\_::PathPoint\_!}
\item
  字段类型：\passthrough{\lstinline!float!}
\item
  声明位置：\passthrough{\lstinline!include/unitree/idl/go2/PathPoint\_.hpp:23!}
\end{itemize}

\textbf{用法}

\begin{lstlisting}[language=Python]
current_value = obj.x
obj.x = new_value
\end{lstlisting}

\hypertarget{unitree-sdk2-cpp-idl-go2-pathpoint-y}{%
\paragraph{\texorpdfstring{\texttt{unitree\_sdk2\_cpp.idl.go2.PathPoint.y}}{unitree\_sdk2\_cpp.idl.go2.PathPoint.y}}\label{unitree-sdk2-cpp-idl-go2-pathpoint-y}}

底层 \passthrough{\lstinline!unitree\_go::msg::dds\_::PathPoint\_!} 的
\passthrough{\lstinline!y!} 字段。类型签名只说明 Python
表示；单位、数值范围、数组固定长度和枚举含义需查对应 IDL 头文件。
底层为浮点数；类型本身不说明单位、坐标系或业务安全范围。

\textbf{签名}

\begin{lstlisting}[language=Python]
@property
def y(self) -> float

@y.setter
def y(self, value: float) -> None
\end{lstlisting}

\textbf{可用性与安全性}

\textbf{\passthrough{\lstinline!AVAILABLE!} /
\passthrough{\lstinline!VALUE\_TYPE!}}：当前绑定源码已实现该消息值操作。

\textbf{参数}

\begin{longtable}[]{@{}
  >{\raggedright\arraybackslash}p{(\columnwidth - 4\tabcolsep) * \real{0.18}}
  >{\raggedright\arraybackslash}p{(\columnwidth - 4\tabcolsep) * \real{0.20}}
  >{\raggedright\arraybackslash}p{(\columnwidth - 4\tabcolsep) * \real{0.62}}@{}}
\toprule
名称 & Python 类型 & 含义 \\
\midrule
\endhead
\passthrough{\lstinline!value!} & \passthrough{\lstinline!float!} &
要写入或传递的新值，必须符合该参数的 Python 类型以及底层 C++
范围约束。 \\
\bottomrule
\end{longtable}

\textbf{返回值}

\begin{longtable}[]{@{}
  >{\raggedright\arraybackslash}p{(\columnwidth - 4\tabcolsep) * \real{0.18}}
  >{\raggedright\arraybackslash}p{(\columnwidth - 4\tabcolsep) * \real{0.20}}
  >{\raggedright\arraybackslash}p{(\columnwidth - 4\tabcolsep) * \real{0.62}}@{}}
\toprule
位置 & 类型 & 含义 \\
\midrule
\endhead
getter 返回值 & \passthrough{\lstinline!float!} & 返回属性当前值。 \\
\bottomrule
\end{longtable}

\textbf{对应 C++}

\begin{itemize}
\tightlist
\item
  类：\passthrough{\lstinline!unitree\_go::msg::dds\_::PathPoint\_!}
\item
  字段类型：\passthrough{\lstinline!float!}
\item
  声明位置：\passthrough{\lstinline!include/unitree/idl/go2/PathPoint\_.hpp:24!}
\end{itemize}

\textbf{用法}

\begin{lstlisting}[language=Python]
current_value = obj.y
obj.y = new_value
\end{lstlisting}

\hypertarget{unitree-sdk2-cpp-idl-go2-pathpoint-yaw}{%
\paragraph{\texorpdfstring{\texttt{unitree\_sdk2\_cpp.idl.go2.PathPoint.yaw}}{unitree\_sdk2\_cpp.idl.go2.PathPoint.yaw}}\label{unitree-sdk2-cpp-idl-go2-pathpoint-yaw}}

底层 \passthrough{\lstinline!unitree\_go::msg::dds\_::PathPoint\_!} 的
\passthrough{\lstinline!yaw!} 字段。类型签名只说明 Python
表示；单位、数值范围、数组固定长度和枚举含义需查对应 IDL 头文件。
底层为浮点数；类型本身不说明单位、坐标系或业务安全范围。

\textbf{签名}

\begin{lstlisting}[language=Python]
@property
def yaw(self) -> float

@yaw.setter
def yaw(self, value: float) -> None
\end{lstlisting}

\textbf{可用性与安全性}

\textbf{\passthrough{\lstinline!AVAILABLE!} /
\passthrough{\lstinline!VALUE\_TYPE!}}：当前绑定源码已实现该消息值操作。

\textbf{参数}

\begin{longtable}[]{@{}
  >{\raggedright\arraybackslash}p{(\columnwidth - 4\tabcolsep) * \real{0.18}}
  >{\raggedright\arraybackslash}p{(\columnwidth - 4\tabcolsep) * \real{0.20}}
  >{\raggedright\arraybackslash}p{(\columnwidth - 4\tabcolsep) * \real{0.62}}@{}}
\toprule
名称 & Python 类型 & 含义 \\
\midrule
\endhead
\passthrough{\lstinline!value!} & \passthrough{\lstinline!float!} &
要写入或传递的新值，必须符合该参数的 Python 类型以及底层 C++
范围约束。 \\
\bottomrule
\end{longtable}

\textbf{返回值}

\begin{longtable}[]{@{}
  >{\raggedright\arraybackslash}p{(\columnwidth - 4\tabcolsep) * \real{0.18}}
  >{\raggedright\arraybackslash}p{(\columnwidth - 4\tabcolsep) * \real{0.20}}
  >{\raggedright\arraybackslash}p{(\columnwidth - 4\tabcolsep) * \real{0.62}}@{}}
\toprule
位置 & 类型 & 含义 \\
\midrule
\endhead
getter 返回值 & \passthrough{\lstinline!float!} & 返回属性当前值。 \\
\bottomrule
\end{longtable}

\textbf{对应 C++}

\begin{itemize}
\tightlist
\item
  类：\passthrough{\lstinline!unitree\_go::msg::dds\_::PathPoint\_!}
\item
  字段类型：\passthrough{\lstinline!float!}
\item
  声明位置：\passthrough{\lstinline!include/unitree/idl/go2/PathPoint\_.hpp:25!}
\end{itemize}

\textbf{用法}

\begin{lstlisting}[language=Python]
current_value = obj.yaw
obj.yaw = new_value
\end{lstlisting}

\hypertarget{unitree-sdk2-cpp-idl-go2-pathpoint-vx}{%
\paragraph{\texorpdfstring{\texttt{unitree\_sdk2\_cpp.idl.go2.PathPoint.vx}}{unitree\_sdk2\_cpp.idl.go2.PathPoint.vx}}\label{unitree-sdk2-cpp-idl-go2-pathpoint-vx}}

底层 \passthrough{\lstinline!unitree\_go::msg::dds\_::PathPoint\_!} 的
\passthrough{\lstinline!vx!} 字段。类型签名只说明 Python
表示；单位、数值范围、数组固定长度和枚举含义需查对应 IDL 头文件。
底层为浮点数；类型本身不说明单位、坐标系或业务安全范围。

\textbf{签名}

\begin{lstlisting}[language=Python]
@property
def vx(self) -> float

@vx.setter
def vx(self, value: float) -> None
\end{lstlisting}

\textbf{可用性与安全性}

\textbf{\passthrough{\lstinline!AVAILABLE!} /
\passthrough{\lstinline!VALUE\_TYPE!}}：当前绑定源码已实现该消息值操作。

\textbf{参数}

\begin{longtable}[]{@{}
  >{\raggedright\arraybackslash}p{(\columnwidth - 4\tabcolsep) * \real{0.18}}
  >{\raggedright\arraybackslash}p{(\columnwidth - 4\tabcolsep) * \real{0.20}}
  >{\raggedright\arraybackslash}p{(\columnwidth - 4\tabcolsep) * \real{0.62}}@{}}
\toprule
名称 & Python 类型 & 含义 \\
\midrule
\endhead
\passthrough{\lstinline!value!} & \passthrough{\lstinline!float!} &
要写入或传递的新值，必须符合该参数的 Python 类型以及底层 C++
范围约束。 \\
\bottomrule
\end{longtable}

\textbf{返回值}

\begin{longtable}[]{@{}
  >{\raggedright\arraybackslash}p{(\columnwidth - 4\tabcolsep) * \real{0.18}}
  >{\raggedright\arraybackslash}p{(\columnwidth - 4\tabcolsep) * \real{0.20}}
  >{\raggedright\arraybackslash}p{(\columnwidth - 4\tabcolsep) * \real{0.62}}@{}}
\toprule
位置 & 类型 & 含义 \\
\midrule
\endhead
getter 返回值 & \passthrough{\lstinline!float!} & 返回属性当前值。 \\
\bottomrule
\end{longtable}

\textbf{对应 C++}

\begin{itemize}
\tightlist
\item
  类：\passthrough{\lstinline!unitree\_go::msg::dds\_::PathPoint\_!}
\item
  字段类型：\passthrough{\lstinline!float!}
\item
  声明位置：\passthrough{\lstinline!include/unitree/idl/go2/PathPoint\_.hpp:26!}
\end{itemize}

\textbf{用法}

\begin{lstlisting}[language=Python]
current_value = obj.vx
obj.vx = new_value
\end{lstlisting}

\hypertarget{unitree-sdk2-cpp-idl-go2-pathpoint-vy}{%
\paragraph{\texorpdfstring{\texttt{unitree\_sdk2\_cpp.idl.go2.PathPoint.vy}}{unitree\_sdk2\_cpp.idl.go2.PathPoint.vy}}\label{unitree-sdk2-cpp-idl-go2-pathpoint-vy}}

底层 \passthrough{\lstinline!unitree\_go::msg::dds\_::PathPoint\_!} 的
\passthrough{\lstinline!vy!} 字段。类型签名只说明 Python
表示；单位、数值范围、数组固定长度和枚举含义需查对应 IDL 头文件。
底层为浮点数；类型本身不说明单位、坐标系或业务安全范围。

\textbf{签名}

\begin{lstlisting}[language=Python]
@property
def vy(self) -> float

@vy.setter
def vy(self, value: float) -> None
\end{lstlisting}

\textbf{可用性与安全性}

\textbf{\passthrough{\lstinline!AVAILABLE!} /
\passthrough{\lstinline!VALUE\_TYPE!}}：当前绑定源码已实现该消息值操作。

\textbf{参数}

\begin{longtable}[]{@{}
  >{\raggedright\arraybackslash}p{(\columnwidth - 4\tabcolsep) * \real{0.18}}
  >{\raggedright\arraybackslash}p{(\columnwidth - 4\tabcolsep) * \real{0.20}}
  >{\raggedright\arraybackslash}p{(\columnwidth - 4\tabcolsep) * \real{0.62}}@{}}
\toprule
名称 & Python 类型 & 含义 \\
\midrule
\endhead
\passthrough{\lstinline!value!} & \passthrough{\lstinline!float!} &
要写入或传递的新值，必须符合该参数的 Python 类型以及底层 C++
范围约束。 \\
\bottomrule
\end{longtable}

\textbf{返回值}

\begin{longtable}[]{@{}
  >{\raggedright\arraybackslash}p{(\columnwidth - 4\tabcolsep) * \real{0.18}}
  >{\raggedright\arraybackslash}p{(\columnwidth - 4\tabcolsep) * \real{0.20}}
  >{\raggedright\arraybackslash}p{(\columnwidth - 4\tabcolsep) * \real{0.62}}@{}}
\toprule
位置 & 类型 & 含义 \\
\midrule
\endhead
getter 返回值 & \passthrough{\lstinline!float!} & 返回属性当前值。 \\
\bottomrule
\end{longtable}

\textbf{对应 C++}

\begin{itemize}
\tightlist
\item
  类：\passthrough{\lstinline!unitree\_go::msg::dds\_::PathPoint\_!}
\item
  字段类型：\passthrough{\lstinline!float!}
\item
  声明位置：\passthrough{\lstinline!include/unitree/idl/go2/PathPoint\_.hpp:27!}
\end{itemize}

\textbf{用法}

\begin{lstlisting}[language=Python]
current_value = obj.vy
obj.vy = new_value
\end{lstlisting}

\hypertarget{unitree-sdk2-cpp-idl-go2-pathpoint-vyaw}{%
\paragraph{\texorpdfstring{\texttt{unitree\_sdk2\_cpp.idl.go2.PathPoint.vyaw}}{unitree\_sdk2\_cpp.idl.go2.PathPoint.vyaw}}\label{unitree-sdk2-cpp-idl-go2-pathpoint-vyaw}}

底层 \passthrough{\lstinline!unitree\_go::msg::dds\_::PathPoint\_!} 的
\passthrough{\lstinline!vyaw!} 字段。类型签名只说明 Python
表示；单位、数值范围、数组固定长度和枚举含义需查对应 IDL 头文件。
底层为浮点数；类型本身不说明单位、坐标系或业务安全范围。

\textbf{签名}

\begin{lstlisting}[language=Python]
@property
def vyaw(self) -> float

@vyaw.setter
def vyaw(self, value: float) -> None
\end{lstlisting}

\textbf{可用性与安全性}

\textbf{\passthrough{\lstinline!AVAILABLE!} /
\passthrough{\lstinline!VALUE\_TYPE!}}：当前绑定源码已实现该消息值操作。

\textbf{参数}

\begin{longtable}[]{@{}
  >{\raggedright\arraybackslash}p{(\columnwidth - 4\tabcolsep) * \real{0.18}}
  >{\raggedright\arraybackslash}p{(\columnwidth - 4\tabcolsep) * \real{0.20}}
  >{\raggedright\arraybackslash}p{(\columnwidth - 4\tabcolsep) * \real{0.62}}@{}}
\toprule
名称 & Python 类型 & 含义 \\
\midrule
\endhead
\passthrough{\lstinline!value!} & \passthrough{\lstinline!float!} &
要写入或传递的新值，必须符合该参数的 Python 类型以及底层 C++
范围约束。 \\
\bottomrule
\end{longtable}

\textbf{返回值}

\begin{longtable}[]{@{}
  >{\raggedright\arraybackslash}p{(\columnwidth - 4\tabcolsep) * \real{0.18}}
  >{\raggedright\arraybackslash}p{(\columnwidth - 4\tabcolsep) * \real{0.20}}
  >{\raggedright\arraybackslash}p{(\columnwidth - 4\tabcolsep) * \real{0.62}}@{}}
\toprule
位置 & 类型 & 含义 \\
\midrule
\endhead
getter 返回值 & \passthrough{\lstinline!float!} & 返回属性当前值。 \\
\bottomrule
\end{longtable}

\textbf{对应 C++}

\begin{itemize}
\tightlist
\item
  类：\passthrough{\lstinline!unitree\_go::msg::dds\_::PathPoint\_!}
\item
  字段类型：\passthrough{\lstinline!float!}
\item
  声明位置：\passthrough{\lstinline!include/unitree/idl/go2/PathPoint\_.hpp:28!}
\end{itemize}

\textbf{用法}

\begin{lstlisting}[language=Python]
current_value = obj.vyaw
obj.vyaw = new_value
\end{lstlisting}

\hypertarget{unitree-sdk2-cpp-idl-go2-req}{%
\subsubsection{\texorpdfstring{\texttt{unitree\_sdk2\_cpp.idl.go2.Req}}{unitree\_sdk2\_cpp.idl.go2.Req}}\label{unitree-sdk2-cpp-idl-go2-req}}

可由 Python 读写的 DDS/IDL 消息值类型。字段采用复制语义。

\textbf{导入}

\begin{lstlisting}[language=Python]
from unitree_sdk2_cpp.idl.go2 import Req
\end{lstlisting}

\textbf{方法索引}

\begin{longtable}[]{@{}
  >{\raggedright\arraybackslash}p{(\columnwidth - 6\tabcolsep) * \real{0.18}}
  >{\raggedright\arraybackslash}p{(\columnwidth - 6\tabcolsep) * \real{0.24}}
  >{\raggedright\arraybackslash}p{(\columnwidth - 6\tabcolsep) * \real{0.18}}
  >{\raggedright\arraybackslash}p{(\columnwidth - 6\tabcolsep) * \real{0.40}}@{}}
\toprule
方法 & Python 签名 & 状态 & 安全分类 \\
\midrule
\endhead
\protect\hyperlink{unitree-sdk2-cpp-idl-go2-req-dunder-init-1}{\passthrough{\lstinline!\_\_init\_\_!}}
& \passthrough{\lstinline!def \_\_init\_\_(self) -> None!} &
\passthrough{\lstinline!AVAILABLE!} &
\passthrough{\lstinline!VALUE\_TYPE!} \\
\protect\hyperlink{unitree-sdk2-cpp-idl-go2-req-dunder-eq-1}{\passthrough{\lstinline!\_\_eq\_\_!}}
& \passthrough{\lstinline!def \_\_eq\_\_(self, other: object) -> bool!}
& \passthrough{\lstinline!AVAILABLE!} &
\passthrough{\lstinline!VALUE\_TYPE!} \\
\protect\hyperlink{unitree-sdk2-cpp-idl-go2-req-dunder-ne-1}{\passthrough{\lstinline!\_\_ne\_\_!}}
& \passthrough{\lstinline!def \_\_ne\_\_(self, other: object) -> bool!}
& \passthrough{\lstinline!AVAILABLE!} &
\passthrough{\lstinline!VALUE\_TYPE!} \\
\bottomrule
\end{longtable}

\textbf{属性索引}

\begin{longtable}[]{@{}
  >{\raggedright\arraybackslash}p{(\columnwidth - 6\tabcolsep) * \real{0.18}}
  >{\raggedright\arraybackslash}p{(\columnwidth - 6\tabcolsep) * \real{0.24}}
  >{\raggedright\arraybackslash}p{(\columnwidth - 6\tabcolsep) * \real{0.18}}
  >{\raggedright\arraybackslash}p{(\columnwidth - 6\tabcolsep) * \real{0.40}}@{}}
\toprule
属性 & 读取类型 & 写入类型 & 状态 \\
\midrule
\endhead
\protect\hyperlink{unitree-sdk2-cpp-idl-go2-req-uuid}{\passthrough{\lstinline!uuid!}}
& \passthrough{\lstinline!str!} & \passthrough{\lstinline!str!} &
\passthrough{\lstinline!AVAILABLE!} \\
\protect\hyperlink{unitree-sdk2-cpp-idl-go2-req-body}{\passthrough{\lstinline!body!}}
& \passthrough{\lstinline!str!} & \passthrough{\lstinline!str!} &
\passthrough{\lstinline!AVAILABLE!} \\
\bottomrule
\end{longtable}

\hypertarget{unitree-sdk2-cpp-idl-go2-req-dunder-init-1}{%
\paragraph{\texorpdfstring{\texttt{unitree\_sdk2\_cpp.idl.go2.Req.\_\_init\_\_}}{unitree\_sdk2\_cpp.idl.go2.Req.\_\_init\_\_}}\label{unitree-sdk2-cpp-idl-go2-req-dunder-init-1}}

初始化 \passthrough{\lstinline!Req!}
实例。是否能实际构造取决于下方可用性状态。

\textbf{签名}

\begin{lstlisting}[language=Python]
def __init__(self) -> None
\end{lstlisting}

\textbf{可用性与安全性}

\textbf{\passthrough{\lstinline!AVAILABLE!} /
\passthrough{\lstinline!VALUE\_TYPE!}}：当前绑定源码已实现该消息值操作。

\textbf{参数}

无显式参数。实例方法中的 \passthrough{\lstinline!self!} 由 Python
自动传入。

\textbf{返回值}

\begin{longtable}[]{@{}
  >{\raggedright\arraybackslash}p{(\columnwidth - 4\tabcolsep) * \real{0.18}}
  >{\raggedright\arraybackslash}p{(\columnwidth - 4\tabcolsep) * \real{0.20}}
  >{\raggedright\arraybackslash}p{(\columnwidth - 4\tabcolsep) * \real{0.62}}@{}}
\toprule
位置 & 类型 & 含义 \\
\midrule
\endhead
无 & \passthrough{\lstinline!None!} &
\passthrough{\lstinline!\_\_init\_\_!}
本身不返回值；调用类对象时，在构造函数可用的前提下得到该类实例。 \\
\bottomrule
\end{longtable}

\textbf{对应 C++}

\begin{itemize}
\tightlist
\item
  类：\passthrough{\lstinline!unitree\_go::msg::dds\_::Req\_!}
\item
  签名：\passthrough{\lstinline!Req\_()!}
\item
  绑定策略：\passthrough{\lstinline!未单独标注!}
\item
  声明位置：\passthrough{\lstinline!include/unitree/idl/go2/Req\_.hpp:27!}
\end{itemize}

\textbf{说明}

\begin{itemize}
\tightlist
\item
  示例是局部调用片段；\passthrough{\lstinline!obj!}
  和参数变量表示已经按上层业务校验并准备好的对象。
\end{itemize}

\textbf{用法}

\begin{lstlisting}[language=Python]
value = Req()
\end{lstlisting}

\hypertarget{unitree-sdk2-cpp-idl-go2-req-dunder-eq-1}{%
\paragraph{\texorpdfstring{\texttt{unitree\_sdk2\_cpp.idl.go2.Req.\_\_eq\_\_}}{unitree\_sdk2\_cpp.idl.go2.Req.\_\_eq\_\_}}\label{unitree-sdk2-cpp-idl-go2-req-dunder-eq-1}}

按底层 IDL 消息内容比较两个对象是否相等。

\textbf{签名}

\begin{lstlisting}[language=Python]
def __eq__(self, other: object) -> bool
\end{lstlisting}

\textbf{可用性与安全性}

\textbf{\passthrough{\lstinline!AVAILABLE!} /
\passthrough{\lstinline!VALUE\_TYPE!}}：当前绑定源码已实现该消息值操作。

\textbf{参数}

\begin{longtable}[]{@{}
  >{\raggedright\arraybackslash}p{(\columnwidth - 6\tabcolsep) * \real{0.18}}
  >{\raggedright\arraybackslash}p{(\columnwidth - 6\tabcolsep) * \real{0.24}}
  >{\raggedright\arraybackslash}p{(\columnwidth - 6\tabcolsep) * \real{0.18}}
  >{\raggedright\arraybackslash}p{(\columnwidth - 6\tabcolsep) * \real{0.40}}@{}}
\toprule
名称 & Python 类型 & 默认值 & 含义 \\
\midrule
\endhead
\passthrough{\lstinline!other!} & \passthrough{\lstinline!object!} &
必填 & 用于比较的另一个 Python 对象。类型不兼容时比较结果通常为
\passthrough{\lstinline!False!}。 \\
\bottomrule
\end{longtable}

\textbf{返回值}

\begin{longtable}[]{@{}
  >{\raggedright\arraybackslash}p{(\columnwidth - 4\tabcolsep) * \real{0.18}}
  >{\raggedright\arraybackslash}p{(\columnwidth - 4\tabcolsep) * \real{0.20}}
  >{\raggedright\arraybackslash}p{(\columnwidth - 4\tabcolsep) * \real{0.62}}@{}}
\toprule
位置 & 类型 & 含义 \\
\midrule
\endhead
返回值 & \passthrough{\lstinline!bool!} & 返回
\passthrough{\lstinline!bool!}。更精确的业务含义见该方法用途、C++
签名和目标型号协议。 \\
\bottomrule
\end{longtable}

\textbf{对应 C++}

\begin{itemize}
\tightlist
\item
  类：\passthrough{\lstinline!unitree\_go::msg::dds\_::Req\_!}
\item
  签名：\passthrough{\lstinline!operator==(const value \&) const!}
\item
  绑定策略：\passthrough{\lstinline!未单独标注!}
\item
  声明位置：由 pybind11 手工绑定或生成注册表提供；manifest
  未记录独立头文件行号。
\end{itemize}

\textbf{说明}

\begin{itemize}
\tightlist
\item
  示例是局部调用片段；\passthrough{\lstinline!obj!}
  和参数变量表示已经按上层业务校验并准备好的对象。
\end{itemize}

\textbf{用法}

\begin{lstlisting}[language=Python]
same = left == right
\end{lstlisting}

\hypertarget{unitree-sdk2-cpp-idl-go2-req-dunder-ne-1}{%
\paragraph{\texorpdfstring{\texttt{unitree\_sdk2\_cpp.idl.go2.Req.\_\_ne\_\_}}{unitree\_sdk2\_cpp.idl.go2.Req.\_\_ne\_\_}}\label{unitree-sdk2-cpp-idl-go2-req-dunder-ne-1}}

按底层 IDL 消息内容比较两个对象是否不相等。

\textbf{签名}

\begin{lstlisting}[language=Python]
def __ne__(self, other: object) -> bool
\end{lstlisting}

\textbf{可用性与安全性}

\textbf{\passthrough{\lstinline!AVAILABLE!} /
\passthrough{\lstinline!VALUE\_TYPE!}}：当前绑定源码已实现该消息值操作。

\textbf{参数}

\begin{longtable}[]{@{}
  >{\raggedright\arraybackslash}p{(\columnwidth - 6\tabcolsep) * \real{0.18}}
  >{\raggedright\arraybackslash}p{(\columnwidth - 6\tabcolsep) * \real{0.24}}
  >{\raggedright\arraybackslash}p{(\columnwidth - 6\tabcolsep) * \real{0.18}}
  >{\raggedright\arraybackslash}p{(\columnwidth - 6\tabcolsep) * \real{0.40}}@{}}
\toprule
名称 & Python 类型 & 默认值 & 含义 \\
\midrule
\endhead
\passthrough{\lstinline!other!} & \passthrough{\lstinline!object!} &
必填 & 用于比较的另一个 Python 对象。类型不兼容时比较结果通常为
\passthrough{\lstinline!False!}。 \\
\bottomrule
\end{longtable}

\textbf{返回值}

\begin{longtable}[]{@{}
  >{\raggedright\arraybackslash}p{(\columnwidth - 4\tabcolsep) * \real{0.18}}
  >{\raggedright\arraybackslash}p{(\columnwidth - 4\tabcolsep) * \real{0.20}}
  >{\raggedright\arraybackslash}p{(\columnwidth - 4\tabcolsep) * \real{0.62}}@{}}
\toprule
位置 & 类型 & 含义 \\
\midrule
\endhead
返回值 & \passthrough{\lstinline!bool!} & 返回
\passthrough{\lstinline!bool!}。更精确的业务含义见该方法用途、C++
签名和目标型号协议。 \\
\bottomrule
\end{longtable}

\textbf{对应 C++}

\begin{itemize}
\tightlist
\item
  类：\passthrough{\lstinline!unitree\_go::msg::dds\_::Req\_!}
\item
  签名：\passthrough{\lstinline"operator!=(const value \&) const"}
\item
  绑定策略：\passthrough{\lstinline!未单独标注!}
\item
  声明位置：由 pybind11 手工绑定或生成注册表提供；manifest
  未记录独立头文件行号。
\end{itemize}

\textbf{说明}

\begin{itemize}
\tightlist
\item
  示例是局部调用片段；\passthrough{\lstinline!obj!}
  和参数变量表示已经按上层业务校验并准备好的对象。
\end{itemize}

\textbf{用法}

\begin{lstlisting}[language=Python]
different = left != right
\end{lstlisting}

\hypertarget{unitree-sdk2-cpp-idl-go2-req-uuid}{%
\paragraph{\texorpdfstring{\texttt{unitree\_sdk2\_cpp.idl.go2.Req.uuid}}{unitree\_sdk2\_cpp.idl.go2.Req.uuid}}\label{unitree-sdk2-cpp-idl-go2-req-uuid}}

底层 \passthrough{\lstinline!unitree\_go::msg::dds\_::Req\_!} 的
\passthrough{\lstinline!uuid!} 字段。类型签名只说明 Python
表示；单位、数值范围、数组固定长度和枚举含义需查对应 IDL 头文件。
底层为字符串；长度、编码和允许值由具体协议决定。

\textbf{签名}

\begin{lstlisting}[language=Python]
@property
def uuid(self) -> str

@uuid.setter
def uuid(self, value: str) -> None
\end{lstlisting}

\textbf{可用性与安全性}

\textbf{\passthrough{\lstinline!AVAILABLE!} /
\passthrough{\lstinline!VALUE\_TYPE!}}：当前绑定源码已实现该消息值操作。

\textbf{参数}

\begin{longtable}[]{@{}
  >{\raggedright\arraybackslash}p{(\columnwidth - 4\tabcolsep) * \real{0.18}}
  >{\raggedright\arraybackslash}p{(\columnwidth - 4\tabcolsep) * \real{0.20}}
  >{\raggedright\arraybackslash}p{(\columnwidth - 4\tabcolsep) * \real{0.62}}@{}}
\toprule
名称 & Python 类型 & 含义 \\
\midrule
\endhead
\passthrough{\lstinline!value!} & \passthrough{\lstinline!str!} &
要写入或传递的新值，必须符合该参数的 Python 类型以及底层 C++
范围约束。 \\
\bottomrule
\end{longtable}

\textbf{返回值}

\begin{longtable}[]{@{}
  >{\raggedright\arraybackslash}p{(\columnwidth - 4\tabcolsep) * \real{0.18}}
  >{\raggedright\arraybackslash}p{(\columnwidth - 4\tabcolsep) * \real{0.20}}
  >{\raggedright\arraybackslash}p{(\columnwidth - 4\tabcolsep) * \real{0.62}}@{}}
\toprule
位置 & 类型 & 含义 \\
\midrule
\endhead
getter 返回值 & \passthrough{\lstinline!str!} & 返回属性当前值。 \\
\bottomrule
\end{longtable}

\textbf{对应 C++}

\begin{itemize}
\tightlist
\item
  类：\passthrough{\lstinline!unitree\_go::msg::dds\_::Req\_!}
\item
  字段类型：\passthrough{\lstinline!std::string!}
\item
  声明位置：\passthrough{\lstinline!include/unitree/idl/go2/Req\_.hpp:23!}
\end{itemize}

\textbf{用法}

\begin{lstlisting}[language=Python]
current_value = obj.uuid
obj.uuid = new_value
\end{lstlisting}

\hypertarget{unitree-sdk2-cpp-idl-go2-req-body}{%
\paragraph{\texorpdfstring{\texttt{unitree\_sdk2\_cpp.idl.go2.Req.body}}{unitree\_sdk2\_cpp.idl.go2.Req.body}}\label{unitree-sdk2-cpp-idl-go2-req-body}}

底层 \passthrough{\lstinline!unitree\_go::msg::dds\_::Req\_!} 的
\passthrough{\lstinline!body!} 字段。类型签名只说明 Python
表示；单位、数值范围、数组固定长度和枚举含义需查对应 IDL 头文件。
底层为字符串；长度、编码和允许值由具体协议决定。

\textbf{签名}

\begin{lstlisting}[language=Python]
@property
def body(self) -> str

@body.setter
def body(self, value: str) -> None
\end{lstlisting}

\textbf{可用性与安全性}

\textbf{\passthrough{\lstinline!AVAILABLE!} /
\passthrough{\lstinline!VALUE\_TYPE!}}：当前绑定源码已实现该消息值操作。

\textbf{参数}

\begin{longtable}[]{@{}
  >{\raggedright\arraybackslash}p{(\columnwidth - 4\tabcolsep) * \real{0.18}}
  >{\raggedright\arraybackslash}p{(\columnwidth - 4\tabcolsep) * \real{0.20}}
  >{\raggedright\arraybackslash}p{(\columnwidth - 4\tabcolsep) * \real{0.62}}@{}}
\toprule
名称 & Python 类型 & 含义 \\
\midrule
\endhead
\passthrough{\lstinline!value!} & \passthrough{\lstinline!str!} &
要写入或传递的新值，必须符合该参数的 Python 类型以及底层 C++
范围约束。 \\
\bottomrule
\end{longtable}

\textbf{返回值}

\begin{longtable}[]{@{}
  >{\raggedright\arraybackslash}p{(\columnwidth - 4\tabcolsep) * \real{0.18}}
  >{\raggedright\arraybackslash}p{(\columnwidth - 4\tabcolsep) * \real{0.20}}
  >{\raggedright\arraybackslash}p{(\columnwidth - 4\tabcolsep) * \real{0.62}}@{}}
\toprule
位置 & 类型 & 含义 \\
\midrule
\endhead
getter 返回值 & \passthrough{\lstinline!str!} & 返回属性当前值。 \\
\bottomrule
\end{longtable}

\textbf{对应 C++}

\begin{itemize}
\tightlist
\item
  类：\passthrough{\lstinline!unitree\_go::msg::dds\_::Req\_!}
\item
  字段类型：\passthrough{\lstinline!std::string!}
\item
  声明位置：\passthrough{\lstinline!include/unitree/idl/go2/Req\_.hpp:24!}
\end{itemize}

\textbf{用法}

\begin{lstlisting}[language=Python]
current_value = obj.body
obj.body = new_value
\end{lstlisting}

\hypertarget{unitree-sdk2-cpp-idl-go2-res}{%
\subsubsection{\texorpdfstring{\texttt{unitree\_sdk2\_cpp.idl.go2.Res}}{unitree\_sdk2\_cpp.idl.go2.Res}}\label{unitree-sdk2-cpp-idl-go2-res}}

可由 Python 读写的 DDS/IDL 消息值类型。字段采用复制语义。

\textbf{导入}

\begin{lstlisting}[language=Python]
from unitree_sdk2_cpp.idl.go2 import Res
\end{lstlisting}

\textbf{方法索引}

\begin{longtable}[]{@{}
  >{\raggedright\arraybackslash}p{(\columnwidth - 6\tabcolsep) * \real{0.18}}
  >{\raggedright\arraybackslash}p{(\columnwidth - 6\tabcolsep) * \real{0.24}}
  >{\raggedright\arraybackslash}p{(\columnwidth - 6\tabcolsep) * \real{0.18}}
  >{\raggedright\arraybackslash}p{(\columnwidth - 6\tabcolsep) * \real{0.40}}@{}}
\toprule
方法 & Python 签名 & 状态 & 安全分类 \\
\midrule
\endhead
\protect\hyperlink{unitree-sdk2-cpp-idl-go2-res-dunder-init-1}{\passthrough{\lstinline!\_\_init\_\_!}}
& \passthrough{\lstinline!def \_\_init\_\_(self) -> None!} &
\passthrough{\lstinline!AVAILABLE!} &
\passthrough{\lstinline!VALUE\_TYPE!} \\
\protect\hyperlink{unitree-sdk2-cpp-idl-go2-res-dunder-eq-1}{\passthrough{\lstinline!\_\_eq\_\_!}}
& \passthrough{\lstinline!def \_\_eq\_\_(self, other: object) -> bool!}
& \passthrough{\lstinline!AVAILABLE!} &
\passthrough{\lstinline!VALUE\_TYPE!} \\
\protect\hyperlink{unitree-sdk2-cpp-idl-go2-res-dunder-ne-1}{\passthrough{\lstinline!\_\_ne\_\_!}}
& \passthrough{\lstinline!def \_\_ne\_\_(self, other: object) -> bool!}
& \passthrough{\lstinline!AVAILABLE!} &
\passthrough{\lstinline!VALUE\_TYPE!} \\
\bottomrule
\end{longtable}

\textbf{属性索引}

\begin{longtable}[]{@{}
  >{\raggedright\arraybackslash}p{(\columnwidth - 6\tabcolsep) * \real{0.18}}
  >{\raggedright\arraybackslash}p{(\columnwidth - 6\tabcolsep) * \real{0.24}}
  >{\raggedright\arraybackslash}p{(\columnwidth - 6\tabcolsep) * \real{0.18}}
  >{\raggedright\arraybackslash}p{(\columnwidth - 6\tabcolsep) * \real{0.40}}@{}}
\toprule
属性 & 读取类型 & 写入类型 & 状态 \\
\midrule
\endhead
\protect\hyperlink{unitree-sdk2-cpp-idl-go2-res-uuid}{\passthrough{\lstinline!uuid!}}
& \passthrough{\lstinline!str!} & \passthrough{\lstinline!str!} &
\passthrough{\lstinline!AVAILABLE!} \\
\protect\hyperlink{unitree-sdk2-cpp-idl-go2-res-data}{\passthrough{\lstinline!data!}}
& \passthrough{\lstinline!list[int]!} &
\passthrough{\lstinline!Sequence[int]!} &
\passthrough{\lstinline!AVAILABLE!} \\
\protect\hyperlink{unitree-sdk2-cpp-idl-go2-res-body}{\passthrough{\lstinline!body!}}
& \passthrough{\lstinline!str!} & \passthrough{\lstinline!str!} &
\passthrough{\lstinline!AVAILABLE!} \\
\bottomrule
\end{longtable}

\hypertarget{unitree-sdk2-cpp-idl-go2-res-dunder-init-1}{%
\paragraph{\texorpdfstring{\texttt{unitree\_sdk2\_cpp.idl.go2.Res.\_\_init\_\_}}{unitree\_sdk2\_cpp.idl.go2.Res.\_\_init\_\_}}\label{unitree-sdk2-cpp-idl-go2-res-dunder-init-1}}

初始化 \passthrough{\lstinline!Res!}
实例。是否能实际构造取决于下方可用性状态。

\textbf{签名}

\begin{lstlisting}[language=Python]
def __init__(self) -> None
\end{lstlisting}

\textbf{可用性与安全性}

\textbf{\passthrough{\lstinline!AVAILABLE!} /
\passthrough{\lstinline!VALUE\_TYPE!}}：当前绑定源码已实现该消息值操作。

\textbf{参数}

无显式参数。实例方法中的 \passthrough{\lstinline!self!} 由 Python
自动传入。

\textbf{返回值}

\begin{longtable}[]{@{}
  >{\raggedright\arraybackslash}p{(\columnwidth - 4\tabcolsep) * \real{0.18}}
  >{\raggedright\arraybackslash}p{(\columnwidth - 4\tabcolsep) * \real{0.20}}
  >{\raggedright\arraybackslash}p{(\columnwidth - 4\tabcolsep) * \real{0.62}}@{}}
\toprule
位置 & 类型 & 含义 \\
\midrule
\endhead
无 & \passthrough{\lstinline!None!} &
\passthrough{\lstinline!\_\_init\_\_!}
本身不返回值；调用类对象时，在构造函数可用的前提下得到该类实例。 \\
\bottomrule
\end{longtable}

\textbf{对应 C++}

\begin{itemize}
\tightlist
\item
  类：\passthrough{\lstinline!unitree\_go::msg::dds\_::Res\_!}
\item
  签名：\passthrough{\lstinline!Res\_()!}
\item
  绑定策略：\passthrough{\lstinline!未单独标注!}
\item
  声明位置：\passthrough{\lstinline!include/unitree/idl/go2/Res\_.hpp:30!}
\end{itemize}

\textbf{说明}

\begin{itemize}
\tightlist
\item
  示例是局部调用片段；\passthrough{\lstinline!obj!}
  和参数变量表示已经按上层业务校验并准备好的对象。
\end{itemize}

\textbf{用法}

\begin{lstlisting}[language=Python]
value = Res()
\end{lstlisting}

\hypertarget{unitree-sdk2-cpp-idl-go2-res-dunder-eq-1}{%
\paragraph{\texorpdfstring{\texttt{unitree\_sdk2\_cpp.idl.go2.Res.\_\_eq\_\_}}{unitree\_sdk2\_cpp.idl.go2.Res.\_\_eq\_\_}}\label{unitree-sdk2-cpp-idl-go2-res-dunder-eq-1}}

按底层 IDL 消息内容比较两个对象是否相等。

\textbf{签名}

\begin{lstlisting}[language=Python]
def __eq__(self, other: object) -> bool
\end{lstlisting}

\textbf{可用性与安全性}

\textbf{\passthrough{\lstinline!AVAILABLE!} /
\passthrough{\lstinline!VALUE\_TYPE!}}：当前绑定源码已实现该消息值操作。

\textbf{参数}

\begin{longtable}[]{@{}
  >{\raggedright\arraybackslash}p{(\columnwidth - 6\tabcolsep) * \real{0.18}}
  >{\raggedright\arraybackslash}p{(\columnwidth - 6\tabcolsep) * \real{0.24}}
  >{\raggedright\arraybackslash}p{(\columnwidth - 6\tabcolsep) * \real{0.18}}
  >{\raggedright\arraybackslash}p{(\columnwidth - 6\tabcolsep) * \real{0.40}}@{}}
\toprule
名称 & Python 类型 & 默认值 & 含义 \\
\midrule
\endhead
\passthrough{\lstinline!other!} & \passthrough{\lstinline!object!} &
必填 & 用于比较的另一个 Python 对象。类型不兼容时比较结果通常为
\passthrough{\lstinline!False!}。 \\
\bottomrule
\end{longtable}

\textbf{返回值}

\begin{longtable}[]{@{}
  >{\raggedright\arraybackslash}p{(\columnwidth - 4\tabcolsep) * \real{0.18}}
  >{\raggedright\arraybackslash}p{(\columnwidth - 4\tabcolsep) * \real{0.20}}
  >{\raggedright\arraybackslash}p{(\columnwidth - 4\tabcolsep) * \real{0.62}}@{}}
\toprule
位置 & 类型 & 含义 \\
\midrule
\endhead
返回值 & \passthrough{\lstinline!bool!} & 返回
\passthrough{\lstinline!bool!}。更精确的业务含义见该方法用途、C++
签名和目标型号协议。 \\
\bottomrule
\end{longtable}

\textbf{对应 C++}

\begin{itemize}
\tightlist
\item
  类：\passthrough{\lstinline!unitree\_go::msg::dds\_::Res\_!}
\item
  签名：\passthrough{\lstinline!operator==(const value \&) const!}
\item
  绑定策略：\passthrough{\lstinline!未单独标注!}
\item
  声明位置：由 pybind11 手工绑定或生成注册表提供；manifest
  未记录独立头文件行号。
\end{itemize}

\textbf{说明}

\begin{itemize}
\tightlist
\item
  示例是局部调用片段；\passthrough{\lstinline!obj!}
  和参数变量表示已经按上层业务校验并准备好的对象。
\end{itemize}

\textbf{用法}

\begin{lstlisting}[language=Python]
same = left == right
\end{lstlisting}

\hypertarget{unitree-sdk2-cpp-idl-go2-res-dunder-ne-1}{%
\paragraph{\texorpdfstring{\texttt{unitree\_sdk2\_cpp.idl.go2.Res.\_\_ne\_\_}}{unitree\_sdk2\_cpp.idl.go2.Res.\_\_ne\_\_}}\label{unitree-sdk2-cpp-idl-go2-res-dunder-ne-1}}

按底层 IDL 消息内容比较两个对象是否不相等。

\textbf{签名}

\begin{lstlisting}[language=Python]
def __ne__(self, other: object) -> bool
\end{lstlisting}

\textbf{可用性与安全性}

\textbf{\passthrough{\lstinline!AVAILABLE!} /
\passthrough{\lstinline!VALUE\_TYPE!}}：当前绑定源码已实现该消息值操作。

\textbf{参数}

\begin{longtable}[]{@{}
  >{\raggedright\arraybackslash}p{(\columnwidth - 6\tabcolsep) * \real{0.18}}
  >{\raggedright\arraybackslash}p{(\columnwidth - 6\tabcolsep) * \real{0.24}}
  >{\raggedright\arraybackslash}p{(\columnwidth - 6\tabcolsep) * \real{0.18}}
  >{\raggedright\arraybackslash}p{(\columnwidth - 6\tabcolsep) * \real{0.40}}@{}}
\toprule
名称 & Python 类型 & 默认值 & 含义 \\
\midrule
\endhead
\passthrough{\lstinline!other!} & \passthrough{\lstinline!object!} &
必填 & 用于比较的另一个 Python 对象。类型不兼容时比较结果通常为
\passthrough{\lstinline!False!}。 \\
\bottomrule
\end{longtable}

\textbf{返回值}

\begin{longtable}[]{@{}
  >{\raggedright\arraybackslash}p{(\columnwidth - 4\tabcolsep) * \real{0.18}}
  >{\raggedright\arraybackslash}p{(\columnwidth - 4\tabcolsep) * \real{0.20}}
  >{\raggedright\arraybackslash}p{(\columnwidth - 4\tabcolsep) * \real{0.62}}@{}}
\toprule
位置 & 类型 & 含义 \\
\midrule
\endhead
返回值 & \passthrough{\lstinline!bool!} & 返回
\passthrough{\lstinline!bool!}。更精确的业务含义见该方法用途、C++
签名和目标型号协议。 \\
\bottomrule
\end{longtable}

\textbf{对应 C++}

\begin{itemize}
\tightlist
\item
  类：\passthrough{\lstinline!unitree\_go::msg::dds\_::Res\_!}
\item
  签名：\passthrough{\lstinline"operator!=(const value \&) const"}
\item
  绑定策略：\passthrough{\lstinline!未单独标注!}
\item
  声明位置：由 pybind11 手工绑定或生成注册表提供；manifest
  未记录独立头文件行号。
\end{itemize}

\textbf{说明}

\begin{itemize}
\tightlist
\item
  示例是局部调用片段；\passthrough{\lstinline!obj!}
  和参数变量表示已经按上层业务校验并准备好的对象。
\end{itemize}

\textbf{用法}

\begin{lstlisting}[language=Python]
different = left != right
\end{lstlisting}

\hypertarget{unitree-sdk2-cpp-idl-go2-res-uuid}{%
\paragraph{\texorpdfstring{\texttt{unitree\_sdk2\_cpp.idl.go2.Res.uuid}}{unitree\_sdk2\_cpp.idl.go2.Res.uuid}}\label{unitree-sdk2-cpp-idl-go2-res-uuid}}

底层 \passthrough{\lstinline!unitree\_go::msg::dds\_::Res\_!} 的
\passthrough{\lstinline!uuid!} 字段。类型签名只说明 Python
表示；单位、数值范围、数组固定长度和枚举含义需查对应 IDL 头文件。
底层为字符串；长度、编码和允许值由具体协议决定。

\textbf{签名}

\begin{lstlisting}[language=Python]
@property
def uuid(self) -> str

@uuid.setter
def uuid(self, value: str) -> None
\end{lstlisting}

\textbf{可用性与安全性}

\textbf{\passthrough{\lstinline!AVAILABLE!} /
\passthrough{\lstinline!VALUE\_TYPE!}}：当前绑定源码已实现该消息值操作。

\textbf{参数}

\begin{longtable}[]{@{}
  >{\raggedright\arraybackslash}p{(\columnwidth - 4\tabcolsep) * \real{0.18}}
  >{\raggedright\arraybackslash}p{(\columnwidth - 4\tabcolsep) * \real{0.20}}
  >{\raggedright\arraybackslash}p{(\columnwidth - 4\tabcolsep) * \real{0.62}}@{}}
\toprule
名称 & Python 类型 & 含义 \\
\midrule
\endhead
\passthrough{\lstinline!value!} & \passthrough{\lstinline!str!} &
要写入或传递的新值，必须符合该参数的 Python 类型以及底层 C++
范围约束。 \\
\bottomrule
\end{longtable}

\textbf{返回值}

\begin{longtable}[]{@{}
  >{\raggedright\arraybackslash}p{(\columnwidth - 4\tabcolsep) * \real{0.18}}
  >{\raggedright\arraybackslash}p{(\columnwidth - 4\tabcolsep) * \real{0.20}}
  >{\raggedright\arraybackslash}p{(\columnwidth - 4\tabcolsep) * \real{0.62}}@{}}
\toprule
位置 & 类型 & 含义 \\
\midrule
\endhead
getter 返回值 & \passthrough{\lstinline!str!} & 返回属性当前值。 \\
\bottomrule
\end{longtable}

\textbf{对应 C++}

\begin{itemize}
\tightlist
\item
  类：\passthrough{\lstinline!unitree\_go::msg::dds\_::Res\_!}
\item
  字段类型：\passthrough{\lstinline!std::string!}
\item
  声明位置：\passthrough{\lstinline!include/unitree/idl/go2/Res\_.hpp:25!}
\end{itemize}

\textbf{用法}

\begin{lstlisting}[language=Python]
current_value = obj.uuid
obj.uuid = new_value
\end{lstlisting}

\hypertarget{unitree-sdk2-cpp-idl-go2-res-data}{%
\paragraph{\texorpdfstring{\texttt{unitree\_sdk2\_cpp.idl.go2.Res.data}}{unitree\_sdk2\_cpp.idl.go2.Res.data}}\label{unitree-sdk2-cpp-idl-go2-res-data}}

底层 \passthrough{\lstinline!unitree\_go::msg::dds\_::Res\_!} 的
\passthrough{\lstinline!data!} 字段。类型签名只说明 Python
表示；单位、数值范围、数组固定长度和枚举含义需查对应 IDL 头文件。
底层是可变长度 vector；元素约束：取值范围为 0 到 255。

\textbf{签名}

\begin{lstlisting}[language=Python]
@property
def data(self) -> list[int]

@data.setter
def data(self, value: Sequence[int]) -> None
\end{lstlisting}

\textbf{可用性与安全性}

\textbf{\passthrough{\lstinline!AVAILABLE!} /
\passthrough{\lstinline!VALUE\_TYPE!}}：当前绑定源码已实现该消息值操作。

\textbf{参数}

\begin{longtable}[]{@{}
  >{\raggedright\arraybackslash}p{(\columnwidth - 4\tabcolsep) * \real{0.18}}
  >{\raggedright\arraybackslash}p{(\columnwidth - 4\tabcolsep) * \real{0.20}}
  >{\raggedright\arraybackslash}p{(\columnwidth - 4\tabcolsep) * \real{0.62}}@{}}
\toprule
名称 & Python 类型 & 含义 \\
\midrule
\endhead
\passthrough{\lstinline!value!} &
\passthrough{\lstinline!Sequence[int]!} &
要写入或传递的新值，必须符合该参数的 Python 类型以及底层 C++
范围约束。 \\
\bottomrule
\end{longtable}

\textbf{返回值}

\begin{longtable}[]{@{}
  >{\raggedright\arraybackslash}p{(\columnwidth - 4\tabcolsep) * \real{0.18}}
  >{\raggedright\arraybackslash}p{(\columnwidth - 4\tabcolsep) * \real{0.20}}
  >{\raggedright\arraybackslash}p{(\columnwidth - 4\tabcolsep) * \real{0.62}}@{}}
\toprule
位置 & 类型 & 含义 \\
\midrule
\endhead
getter 返回值 & \passthrough{\lstinline!list[int]!} &
返回属性当前值。数组、vector 和嵌套消息采用复制语义。 \\
\bottomrule
\end{longtable}

\textbf{对应 C++}

\begin{itemize}
\tightlist
\item
  类：\passthrough{\lstinline!unitree\_go::msg::dds\_::Res\_!}
\item
  字段类型：\passthrough{\lstinline!std::vector<uint8\_t>!}
\item
  声明位置：\passthrough{\lstinline!include/unitree/idl/go2/Res\_.hpp:26!}
\end{itemize}

\textbf{用法}

\begin{lstlisting}[language=Python]
current_value = obj.data
obj.data = new_value
\end{lstlisting}

\begin{quote}
{[}!TIP{]} 读取容器或嵌套值后应遵循``读取、修改、重新赋值''。只修改
getter 返回的副本不会可靠地更新原 C++ 消息。
\end{quote}

\hypertarget{unitree-sdk2-cpp-idl-go2-res-body}{%
\paragraph{\texorpdfstring{\texttt{unitree\_sdk2\_cpp.idl.go2.Res.body}}{unitree\_sdk2\_cpp.idl.go2.Res.body}}\label{unitree-sdk2-cpp-idl-go2-res-body}}

底层 \passthrough{\lstinline!unitree\_go::msg::dds\_::Res\_!} 的
\passthrough{\lstinline!body!} 字段。类型签名只说明 Python
表示；单位、数值范围、数组固定长度和枚举含义需查对应 IDL 头文件。
底层为字符串；长度、编码和允许值由具体协议决定。

\textbf{签名}

\begin{lstlisting}[language=Python]
@property
def body(self) -> str

@body.setter
def body(self, value: str) -> None
\end{lstlisting}

\textbf{可用性与安全性}

\textbf{\passthrough{\lstinline!AVAILABLE!} /
\passthrough{\lstinline!VALUE\_TYPE!}}：当前绑定源码已实现该消息值操作。

\textbf{参数}

\begin{longtable}[]{@{}
  >{\raggedright\arraybackslash}p{(\columnwidth - 4\tabcolsep) * \real{0.18}}
  >{\raggedright\arraybackslash}p{(\columnwidth - 4\tabcolsep) * \real{0.20}}
  >{\raggedright\arraybackslash}p{(\columnwidth - 4\tabcolsep) * \real{0.62}}@{}}
\toprule
名称 & Python 类型 & 含义 \\
\midrule
\endhead
\passthrough{\lstinline!value!} & \passthrough{\lstinline!str!} &
要写入或传递的新值，必须符合该参数的 Python 类型以及底层 C++
范围约束。 \\
\bottomrule
\end{longtable}

\textbf{返回值}

\begin{longtable}[]{@{}
  >{\raggedright\arraybackslash}p{(\columnwidth - 4\tabcolsep) * \real{0.18}}
  >{\raggedright\arraybackslash}p{(\columnwidth - 4\tabcolsep) * \real{0.20}}
  >{\raggedright\arraybackslash}p{(\columnwidth - 4\tabcolsep) * \real{0.62}}@{}}
\toprule
位置 & 类型 & 含义 \\
\midrule
\endhead
getter 返回值 & \passthrough{\lstinline!str!} & 返回属性当前值。 \\
\bottomrule
\end{longtable}

\textbf{对应 C++}

\begin{itemize}
\tightlist
\item
  类：\passthrough{\lstinline!unitree\_go::msg::dds\_::Res\_!}
\item
  字段类型：\passthrough{\lstinline!std::string!}
\item
  声明位置：\passthrough{\lstinline!include/unitree/idl/go2/Res\_.hpp:27!}
\end{itemize}

\textbf{用法}

\begin{lstlisting}[language=Python]
current_value = obj.body
obj.body = new_value
\end{lstlisting}

\hypertarget{unitree-sdk2-cpp-idl-go2-sportmodecmd}{%
\subsubsection{\texorpdfstring{\texttt{unitree\_sdk2\_cpp.idl.go2.SportModeCmd}}{unitree\_sdk2\_cpp.idl.go2.SportModeCmd}}\label{unitree-sdk2-cpp-idl-go2-sportmodecmd}}

可由 Python 读写的 DDS/IDL 消息值类型。字段采用复制语义。

\textbf{导入}

\begin{lstlisting}[language=Python]
from unitree_sdk2_cpp.idl.go2 import SportModeCmd
\end{lstlisting}

\textbf{方法索引}

\begin{longtable}[]{@{}
  >{\raggedright\arraybackslash}p{(\columnwidth - 6\tabcolsep) * \real{0.18}}
  >{\raggedright\arraybackslash}p{(\columnwidth - 6\tabcolsep) * \real{0.24}}
  >{\raggedright\arraybackslash}p{(\columnwidth - 6\tabcolsep) * \real{0.18}}
  >{\raggedright\arraybackslash}p{(\columnwidth - 6\tabcolsep) * \real{0.40}}@{}}
\toprule
方法 & Python 签名 & 状态 & 安全分类 \\
\midrule
\endhead
\protect\hyperlink{unitree-sdk2-cpp-idl-go2-sportmodecmd-dunder-init-1}{\passthrough{\lstinline!\_\_init\_\_!}}
& \passthrough{\lstinline!def \_\_init\_\_(self) -> None!} &
\passthrough{\lstinline!AVAILABLE!} &
\passthrough{\lstinline!VALUE\_TYPE!} \\
\protect\hyperlink{unitree-sdk2-cpp-idl-go2-sportmodecmd-dunder-eq-1}{\passthrough{\lstinline!\_\_eq\_\_!}}
& \passthrough{\lstinline!def \_\_eq\_\_(self, other: object) -> bool!}
& \passthrough{\lstinline!AVAILABLE!} &
\passthrough{\lstinline!VALUE\_TYPE!} \\
\protect\hyperlink{unitree-sdk2-cpp-idl-go2-sportmodecmd-dunder-ne-1}{\passthrough{\lstinline!\_\_ne\_\_!}}
& \passthrough{\lstinline!def \_\_ne\_\_(self, other: object) -> bool!}
& \passthrough{\lstinline!AVAILABLE!} &
\passthrough{\lstinline!VALUE\_TYPE!} \\
\bottomrule
\end{longtable}

\textbf{属性索引}

\begin{longtable}[]{@{}
  >{\raggedright\arraybackslash}p{(\columnwidth - 6\tabcolsep) * \real{0.18}}
  >{\raggedright\arraybackslash}p{(\columnwidth - 6\tabcolsep) * \real{0.24}}
  >{\raggedright\arraybackslash}p{(\columnwidth - 6\tabcolsep) * \real{0.18}}
  >{\raggedright\arraybackslash}p{(\columnwidth - 6\tabcolsep) * \real{0.40}}@{}}
\toprule
属性 & 读取类型 & 写入类型 & 状态 \\
\midrule
\endhead
\protect\hyperlink{unitree-sdk2-cpp-idl-go2-sportmodecmd-mode}{\passthrough{\lstinline!mode!}}
& \passthrough{\lstinline!int!} & \passthrough{\lstinline!int!} &
\passthrough{\lstinline!AVAILABLE!} \\
\protect\hyperlink{unitree-sdk2-cpp-idl-go2-sportmodecmd-gait-type}{\passthrough{\lstinline!gait\_type!}}
& \passthrough{\lstinline!int!} & \passthrough{\lstinline!int!} &
\passthrough{\lstinline!AVAILABLE!} \\
\protect\hyperlink{unitree-sdk2-cpp-idl-go2-sportmodecmd-speed-level}{\passthrough{\lstinline!speed\_level!}}
& \passthrough{\lstinline!int!} & \passthrough{\lstinline!int!} &
\passthrough{\lstinline!AVAILABLE!} \\
\protect\hyperlink{unitree-sdk2-cpp-idl-go2-sportmodecmd-foot-raise-height}{\passthrough{\lstinline!foot\_raise\_height!}}
& \passthrough{\lstinline!float!} & \passthrough{\lstinline!float!} &
\passthrough{\lstinline!AVAILABLE!} \\
\protect\hyperlink{unitree-sdk2-cpp-idl-go2-sportmodecmd-body-height}{\passthrough{\lstinline!body\_height!}}
& \passthrough{\lstinline!float!} & \passthrough{\lstinline!float!} &
\passthrough{\lstinline!AVAILABLE!} \\
\protect\hyperlink{unitree-sdk2-cpp-idl-go2-sportmodecmd-position}{\passthrough{\lstinline!position!}}
& \passthrough{\lstinline!list[float]!} &
\passthrough{\lstinline!Sequence[float]!} &
\passthrough{\lstinline!AVAILABLE!} \\
\protect\hyperlink{unitree-sdk2-cpp-idl-go2-sportmodecmd-euler}{\passthrough{\lstinline!euler!}}
& \passthrough{\lstinline!list[float]!} &
\passthrough{\lstinline!Sequence[float]!} &
\passthrough{\lstinline!AVAILABLE!} \\
\protect\hyperlink{unitree-sdk2-cpp-idl-go2-sportmodecmd-velocity}{\passthrough{\lstinline!velocity!}}
& \passthrough{\lstinline!list[float]!} &
\passthrough{\lstinline!Sequence[float]!} &
\passthrough{\lstinline!AVAILABLE!} \\
\protect\hyperlink{unitree-sdk2-cpp-idl-go2-sportmodecmd-yaw-speed}{\passthrough{\lstinline!yaw\_speed!}}
& \passthrough{\lstinline!float!} & \passthrough{\lstinline!float!} &
\passthrough{\lstinline!AVAILABLE!} \\
\protect\hyperlink{unitree-sdk2-cpp-idl-go2-sportmodecmd-bms-cmd}{\passthrough{\lstinline!bms\_cmd!}}
& \passthrough{\lstinline!BmsCmd!} & \passthrough{\lstinline!BmsCmd!} &
\passthrough{\lstinline!AVAILABLE!} \\
\protect\hyperlink{unitree-sdk2-cpp-idl-go2-sportmodecmd-path-point}{\passthrough{\lstinline!path\_point!}}
& \passthrough{\lstinline!list[PathPoint]!} &
\passthrough{\lstinline!Sequence[PathPoint]!} &
\passthrough{\lstinline!AVAILABLE!} \\
\bottomrule
\end{longtable}

\hypertarget{unitree-sdk2-cpp-idl-go2-sportmodecmd-dunder-init-1}{%
\paragraph{\texorpdfstring{\texttt{unitree\_sdk2\_cpp.idl.go2.SportModeCmd.\_\_init\_\_}}{unitree\_sdk2\_cpp.idl.go2.SportModeCmd.\_\_init\_\_}}\label{unitree-sdk2-cpp-idl-go2-sportmodecmd-dunder-init-1}}

初始化 \passthrough{\lstinline!SportModeCmd!}
实例。是否能实际构造取决于下方可用性状态。

\textbf{签名}

\begin{lstlisting}[language=Python]
def __init__(self) -> None
\end{lstlisting}

\textbf{可用性与安全性}

\textbf{\passthrough{\lstinline!AVAILABLE!} /
\passthrough{\lstinline!VALUE\_TYPE!}}：当前绑定源码已实现该消息值操作。

\textbf{参数}

无显式参数。实例方法中的 \passthrough{\lstinline!self!} 由 Python
自动传入。

\textbf{返回值}

\begin{longtable}[]{@{}
  >{\raggedright\arraybackslash}p{(\columnwidth - 4\tabcolsep) * \real{0.18}}
  >{\raggedright\arraybackslash}p{(\columnwidth - 4\tabcolsep) * \real{0.20}}
  >{\raggedright\arraybackslash}p{(\columnwidth - 4\tabcolsep) * \real{0.62}}@{}}
\toprule
位置 & 类型 & 含义 \\
\midrule
\endhead
无 & \passthrough{\lstinline!None!} &
\passthrough{\lstinline!\_\_init\_\_!}
本身不返回值；调用类对象时，在构造函数可用的前提下得到该类实例。 \\
\bottomrule
\end{longtable}

\textbf{对应 C++}

\begin{itemize}
\tightlist
\item
  类：\passthrough{\lstinline!unitree\_go::msg::dds\_::SportModeCmd\_!}
\item
  签名：\passthrough{\lstinline!SportModeCmd\_()!}
\item
  绑定策略：\passthrough{\lstinline!未单独标注!}
\item
  声明位置：\passthrough{\lstinline!include/unitree/idl/go2/SportModeCmd\_.hpp:41!}
\end{itemize}

\textbf{说明}

\begin{itemize}
\tightlist
\item
  示例是局部调用片段；\passthrough{\lstinline!obj!}
  和参数变量表示已经按上层业务校验并准备好的对象。
\end{itemize}

\textbf{用法}

\begin{lstlisting}[language=Python]
value = SportModeCmd()
\end{lstlisting}

\hypertarget{unitree-sdk2-cpp-idl-go2-sportmodecmd-dunder-eq-1}{%
\paragraph{\texorpdfstring{\texttt{unitree\_sdk2\_cpp.idl.go2.SportModeCmd.\_\_eq\_\_}}{unitree\_sdk2\_cpp.idl.go2.SportModeCmd.\_\_eq\_\_}}\label{unitree-sdk2-cpp-idl-go2-sportmodecmd-dunder-eq-1}}

按底层 IDL 消息内容比较两个对象是否相等。

\textbf{签名}

\begin{lstlisting}[language=Python]
def __eq__(self, other: object) -> bool
\end{lstlisting}

\textbf{可用性与安全性}

\textbf{\passthrough{\lstinline!AVAILABLE!} /
\passthrough{\lstinline!VALUE\_TYPE!}}：当前绑定源码已实现该消息值操作。

\textbf{参数}

\begin{longtable}[]{@{}
  >{\raggedright\arraybackslash}p{(\columnwidth - 6\tabcolsep) * \real{0.18}}
  >{\raggedright\arraybackslash}p{(\columnwidth - 6\tabcolsep) * \real{0.24}}
  >{\raggedright\arraybackslash}p{(\columnwidth - 6\tabcolsep) * \real{0.18}}
  >{\raggedright\arraybackslash}p{(\columnwidth - 6\tabcolsep) * \real{0.40}}@{}}
\toprule
名称 & Python 类型 & 默认值 & 含义 \\
\midrule
\endhead
\passthrough{\lstinline!other!} & \passthrough{\lstinline!object!} &
必填 & 用于比较的另一个 Python 对象。类型不兼容时比较结果通常为
\passthrough{\lstinline!False!}。 \\
\bottomrule
\end{longtable}

\textbf{返回值}

\begin{longtable}[]{@{}
  >{\raggedright\arraybackslash}p{(\columnwidth - 4\tabcolsep) * \real{0.18}}
  >{\raggedright\arraybackslash}p{(\columnwidth - 4\tabcolsep) * \real{0.20}}
  >{\raggedright\arraybackslash}p{(\columnwidth - 4\tabcolsep) * \real{0.62}}@{}}
\toprule
位置 & 类型 & 含义 \\
\midrule
\endhead
返回值 & \passthrough{\lstinline!bool!} & 返回
\passthrough{\lstinline!bool!}。更精确的业务含义见该方法用途、C++
签名和目标型号协议。 \\
\bottomrule
\end{longtable}

\textbf{对应 C++}

\begin{itemize}
\tightlist
\item
  类：\passthrough{\lstinline!unitree\_go::msg::dds\_::SportModeCmd\_!}
\item
  签名：\passthrough{\lstinline!operator==(const value \&) const!}
\item
  绑定策略：\passthrough{\lstinline!未单独标注!}
\item
  声明位置：由 pybind11 手工绑定或生成注册表提供；manifest
  未记录独立头文件行号。
\end{itemize}

\textbf{说明}

\begin{itemize}
\tightlist
\item
  示例是局部调用片段；\passthrough{\lstinline!obj!}
  和参数变量表示已经按上层业务校验并准备好的对象。
\end{itemize}

\textbf{用法}

\begin{lstlisting}[language=Python]
same = left == right
\end{lstlisting}

\hypertarget{unitree-sdk2-cpp-idl-go2-sportmodecmd-dunder-ne-1}{%
\paragraph{\texorpdfstring{\texttt{unitree\_sdk2\_cpp.idl.go2.SportModeCmd.\_\_ne\_\_}}{unitree\_sdk2\_cpp.idl.go2.SportModeCmd.\_\_ne\_\_}}\label{unitree-sdk2-cpp-idl-go2-sportmodecmd-dunder-ne-1}}

按底层 IDL 消息内容比较两个对象是否不相等。

\textbf{签名}

\begin{lstlisting}[language=Python]
def __ne__(self, other: object) -> bool
\end{lstlisting}

\textbf{可用性与安全性}

\textbf{\passthrough{\lstinline!AVAILABLE!} /
\passthrough{\lstinline!VALUE\_TYPE!}}：当前绑定源码已实现该消息值操作。

\textbf{参数}

\begin{longtable}[]{@{}
  >{\raggedright\arraybackslash}p{(\columnwidth - 6\tabcolsep) * \real{0.18}}
  >{\raggedright\arraybackslash}p{(\columnwidth - 6\tabcolsep) * \real{0.24}}
  >{\raggedright\arraybackslash}p{(\columnwidth - 6\tabcolsep) * \real{0.18}}
  >{\raggedright\arraybackslash}p{(\columnwidth - 6\tabcolsep) * \real{0.40}}@{}}
\toprule
名称 & Python 类型 & 默认值 & 含义 \\
\midrule
\endhead
\passthrough{\lstinline!other!} & \passthrough{\lstinline!object!} &
必填 & 用于比较的另一个 Python 对象。类型不兼容时比较结果通常为
\passthrough{\lstinline!False!}。 \\
\bottomrule
\end{longtable}

\textbf{返回值}

\begin{longtable}[]{@{}
  >{\raggedright\arraybackslash}p{(\columnwidth - 4\tabcolsep) * \real{0.18}}
  >{\raggedright\arraybackslash}p{(\columnwidth - 4\tabcolsep) * \real{0.20}}
  >{\raggedright\arraybackslash}p{(\columnwidth - 4\tabcolsep) * \real{0.62}}@{}}
\toprule
位置 & 类型 & 含义 \\
\midrule
\endhead
返回值 & \passthrough{\lstinline!bool!} & 返回
\passthrough{\lstinline!bool!}。更精确的业务含义见该方法用途、C++
签名和目标型号协议。 \\
\bottomrule
\end{longtable}

\textbf{对应 C++}

\begin{itemize}
\tightlist
\item
  类：\passthrough{\lstinline!unitree\_go::msg::dds\_::SportModeCmd\_!}
\item
  签名：\passthrough{\lstinline"operator!=(const value \&) const"}
\item
  绑定策略：\passthrough{\lstinline!未单独标注!}
\item
  声明位置：由 pybind11 手工绑定或生成注册表提供；manifest
  未记录独立头文件行号。
\end{itemize}

\textbf{说明}

\begin{itemize}
\tightlist
\item
  示例是局部调用片段；\passthrough{\lstinline!obj!}
  和参数变量表示已经按上层业务校验并准备好的对象。
\end{itemize}

\textbf{用法}

\begin{lstlisting}[language=Python]
different = left != right
\end{lstlisting}

\hypertarget{unitree-sdk2-cpp-idl-go2-sportmodecmd-mode}{%
\paragraph{\texorpdfstring{\texttt{unitree\_sdk2\_cpp.idl.go2.SportModeCmd.mode}}{unitree\_sdk2\_cpp.idl.go2.SportModeCmd.mode}}\label{unitree-sdk2-cpp-idl-go2-sportmodecmd-mode}}

底层 \passthrough{\lstinline!unitree\_go::msg::dds\_::SportModeCmd\_!}
的 \passthrough{\lstinline!mode!} 字段。类型签名只说明 Python
表示；单位、数值范围、数组固定长度和枚举含义需查对应 IDL 头文件。
取值范围为 0 到 255。

\textbf{签名}

\begin{lstlisting}[language=Python]
@property
def mode(self) -> int

@mode.setter
def mode(self, value: int) -> None
\end{lstlisting}

\textbf{可用性与安全性}

\textbf{\passthrough{\lstinline!AVAILABLE!} /
\passthrough{\lstinline!VALUE\_TYPE!}}：当前绑定源码已实现该消息值操作。

\textbf{参数}

\begin{longtable}[]{@{}
  >{\raggedright\arraybackslash}p{(\columnwidth - 4\tabcolsep) * \real{0.18}}
  >{\raggedright\arraybackslash}p{(\columnwidth - 4\tabcolsep) * \real{0.20}}
  >{\raggedright\arraybackslash}p{(\columnwidth - 4\tabcolsep) * \real{0.62}}@{}}
\toprule
名称 & Python 类型 & 含义 \\
\midrule
\endhead
\passthrough{\lstinline!value!} & \passthrough{\lstinline!int!} &
要写入或传递的新值，必须符合该参数的 Python 类型以及底层 C++
范围约束。 \\
\bottomrule
\end{longtable}

\textbf{返回值}

\begin{longtable}[]{@{}
  >{\raggedright\arraybackslash}p{(\columnwidth - 4\tabcolsep) * \real{0.18}}
  >{\raggedright\arraybackslash}p{(\columnwidth - 4\tabcolsep) * \real{0.20}}
  >{\raggedright\arraybackslash}p{(\columnwidth - 4\tabcolsep) * \real{0.62}}@{}}
\toprule
位置 & 类型 & 含义 \\
\midrule
\endhead
getter 返回值 & \passthrough{\lstinline!int!} & 返回属性当前值。 \\
\bottomrule
\end{longtable}

\textbf{对应 C++}

\begin{itemize}
\tightlist
\item
  类：\passthrough{\lstinline!unitree\_go::msg::dds\_::SportModeCmd\_!}
\item
  字段类型：\passthrough{\lstinline!uint8\_t!}
\item
  声明位置：\passthrough{\lstinline!include/unitree/idl/go2/SportModeCmd\_.hpp:28!}
\end{itemize}

\textbf{用法}

\begin{lstlisting}[language=Python]
current_value = obj.mode
obj.mode = new_value
\end{lstlisting}

\hypertarget{unitree-sdk2-cpp-idl-go2-sportmodecmd-gait-type}{%
\paragraph{\texorpdfstring{\texttt{unitree\_sdk2\_cpp.idl.go2.SportModeCmd.gait\_type}}{unitree\_sdk2\_cpp.idl.go2.SportModeCmd.gait\_type}}\label{unitree-sdk2-cpp-idl-go2-sportmodecmd-gait-type}}

底层 \passthrough{\lstinline!unitree\_go::msg::dds\_::SportModeCmd\_!}
的 \passthrough{\lstinline!gait\_type!} 字段。类型签名只说明 Python
表示；单位、数值范围、数组固定长度和枚举含义需查对应 IDL 头文件。
取值范围为 0 到 255。

\textbf{签名}

\begin{lstlisting}[language=Python]
@property
def gait_type(self) -> int

@gait_type.setter
def gait_type(self, value: int) -> None
\end{lstlisting}

\textbf{可用性与安全性}

\textbf{\passthrough{\lstinline!AVAILABLE!} /
\passthrough{\lstinline!VALUE\_TYPE!}}：当前绑定源码已实现该消息值操作。

\textbf{参数}

\begin{longtable}[]{@{}
  >{\raggedright\arraybackslash}p{(\columnwidth - 4\tabcolsep) * \real{0.18}}
  >{\raggedright\arraybackslash}p{(\columnwidth - 4\tabcolsep) * \real{0.20}}
  >{\raggedright\arraybackslash}p{(\columnwidth - 4\tabcolsep) * \real{0.62}}@{}}
\toprule
名称 & Python 类型 & 含义 \\
\midrule
\endhead
\passthrough{\lstinline!value!} & \passthrough{\lstinline!int!} &
要写入或传递的新值，必须符合该参数的 Python 类型以及底层 C++
范围约束。 \\
\bottomrule
\end{longtable}

\textbf{返回值}

\begin{longtable}[]{@{}
  >{\raggedright\arraybackslash}p{(\columnwidth - 4\tabcolsep) * \real{0.18}}
  >{\raggedright\arraybackslash}p{(\columnwidth - 4\tabcolsep) * \real{0.20}}
  >{\raggedright\arraybackslash}p{(\columnwidth - 4\tabcolsep) * \real{0.62}}@{}}
\toprule
位置 & 类型 & 含义 \\
\midrule
\endhead
getter 返回值 & \passthrough{\lstinline!int!} & 返回属性当前值。 \\
\bottomrule
\end{longtable}

\textbf{对应 C++}

\begin{itemize}
\tightlist
\item
  类：\passthrough{\lstinline!unitree\_go::msg::dds\_::SportModeCmd\_!}
\item
  字段类型：\passthrough{\lstinline!uint8\_t!}
\item
  声明位置：\passthrough{\lstinline!include/unitree/idl/go2/SportModeCmd\_.hpp:29!}
\end{itemize}

\textbf{用法}

\begin{lstlisting}[language=Python]
current_value = obj.gait_type
obj.gait_type = new_value
\end{lstlisting}

\hypertarget{unitree-sdk2-cpp-idl-go2-sportmodecmd-speed-level}{%
\paragraph{\texorpdfstring{\texttt{unitree\_sdk2\_cpp.idl.go2.SportModeCmd.speed\_level}}{unitree\_sdk2\_cpp.idl.go2.SportModeCmd.speed\_level}}\label{unitree-sdk2-cpp-idl-go2-sportmodecmd-speed-level}}

底层 \passthrough{\lstinline!unitree\_go::msg::dds\_::SportModeCmd\_!}
的 \passthrough{\lstinline!speed\_level!} 字段。类型签名只说明 Python
表示；单位、数值范围、数组固定长度和枚举含义需查对应 IDL 头文件。
取值范围为 0 到 255。

\textbf{签名}

\begin{lstlisting}[language=Python]
@property
def speed_level(self) -> int

@speed_level.setter
def speed_level(self, value: int) -> None
\end{lstlisting}

\textbf{可用性与安全性}

\textbf{\passthrough{\lstinline!AVAILABLE!} /
\passthrough{\lstinline!VALUE\_TYPE!}}：当前绑定源码已实现该消息值操作。

\textbf{参数}

\begin{longtable}[]{@{}
  >{\raggedright\arraybackslash}p{(\columnwidth - 4\tabcolsep) * \real{0.18}}
  >{\raggedright\arraybackslash}p{(\columnwidth - 4\tabcolsep) * \real{0.20}}
  >{\raggedright\arraybackslash}p{(\columnwidth - 4\tabcolsep) * \real{0.62}}@{}}
\toprule
名称 & Python 类型 & 含义 \\
\midrule
\endhead
\passthrough{\lstinline!value!} & \passthrough{\lstinline!int!} &
要写入或传递的新值，必须符合该参数的 Python 类型以及底层 C++
范围约束。 \\
\bottomrule
\end{longtable}

\textbf{返回值}

\begin{longtable}[]{@{}
  >{\raggedright\arraybackslash}p{(\columnwidth - 4\tabcolsep) * \real{0.18}}
  >{\raggedright\arraybackslash}p{(\columnwidth - 4\tabcolsep) * \real{0.20}}
  >{\raggedright\arraybackslash}p{(\columnwidth - 4\tabcolsep) * \real{0.62}}@{}}
\toprule
位置 & 类型 & 含义 \\
\midrule
\endhead
getter 返回值 & \passthrough{\lstinline!int!} & 返回属性当前值。 \\
\bottomrule
\end{longtable}

\textbf{对应 C++}

\begin{itemize}
\tightlist
\item
  类：\passthrough{\lstinline!unitree\_go::msg::dds\_::SportModeCmd\_!}
\item
  字段类型：\passthrough{\lstinline!uint8\_t!}
\item
  声明位置：\passthrough{\lstinline!include/unitree/idl/go2/SportModeCmd\_.hpp:30!}
\end{itemize}

\textbf{用法}

\begin{lstlisting}[language=Python]
current_value = obj.speed_level
obj.speed_level = new_value
\end{lstlisting}

\hypertarget{unitree-sdk2-cpp-idl-go2-sportmodecmd-foot-raise-height}{%
\paragraph{\texorpdfstring{\texttt{unitree\_sdk2\_cpp.idl.go2.SportModeCmd.foot\_raise\_height}}{unitree\_sdk2\_cpp.idl.go2.SportModeCmd.foot\_raise\_height}}\label{unitree-sdk2-cpp-idl-go2-sportmodecmd-foot-raise-height}}

底层 \passthrough{\lstinline!unitree\_go::msg::dds\_::SportModeCmd\_!}
的 \passthrough{\lstinline!foot\_raise\_height!} 字段。类型签名只说明
Python 表示；单位、数值范围、数组固定长度和枚举含义需查对应 IDL 头文件。
底层为浮点数；类型本身不说明单位、坐标系或业务安全范围。

\textbf{签名}

\begin{lstlisting}[language=Python]
@property
def foot_raise_height(self) -> float

@foot_raise_height.setter
def foot_raise_height(self, value: float) -> None
\end{lstlisting}

\textbf{可用性与安全性}

\textbf{\passthrough{\lstinline!AVAILABLE!} /
\passthrough{\lstinline!VALUE\_TYPE!}}：当前绑定源码已实现该消息值操作。

\textbf{参数}

\begin{longtable}[]{@{}
  >{\raggedright\arraybackslash}p{(\columnwidth - 4\tabcolsep) * \real{0.18}}
  >{\raggedright\arraybackslash}p{(\columnwidth - 4\tabcolsep) * \real{0.20}}
  >{\raggedright\arraybackslash}p{(\columnwidth - 4\tabcolsep) * \real{0.62}}@{}}
\toprule
名称 & Python 类型 & 含义 \\
\midrule
\endhead
\passthrough{\lstinline!value!} & \passthrough{\lstinline!float!} &
要写入或传递的新值，必须符合该参数的 Python 类型以及底层 C++
范围约束。 \\
\bottomrule
\end{longtable}

\textbf{返回值}

\begin{longtable}[]{@{}
  >{\raggedright\arraybackslash}p{(\columnwidth - 4\tabcolsep) * \real{0.18}}
  >{\raggedright\arraybackslash}p{(\columnwidth - 4\tabcolsep) * \real{0.20}}
  >{\raggedright\arraybackslash}p{(\columnwidth - 4\tabcolsep) * \real{0.62}}@{}}
\toprule
位置 & 类型 & 含义 \\
\midrule
\endhead
getter 返回值 & \passthrough{\lstinline!float!} & 返回属性当前值。 \\
\bottomrule
\end{longtable}

\textbf{对应 C++}

\begin{itemize}
\tightlist
\item
  类：\passthrough{\lstinline!unitree\_go::msg::dds\_::SportModeCmd\_!}
\item
  字段类型：\passthrough{\lstinline!float!}
\item
  声明位置：\passthrough{\lstinline!include/unitree/idl/go2/SportModeCmd\_.hpp:31!}
\end{itemize}

\textbf{用法}

\begin{lstlisting}[language=Python]
current_value = obj.foot_raise_height
obj.foot_raise_height = new_value
\end{lstlisting}

\hypertarget{unitree-sdk2-cpp-idl-go2-sportmodecmd-body-height}{%
\paragraph{\texorpdfstring{\texttt{unitree\_sdk2\_cpp.idl.go2.SportModeCmd.body\_height}}{unitree\_sdk2\_cpp.idl.go2.SportModeCmd.body\_height}}\label{unitree-sdk2-cpp-idl-go2-sportmodecmd-body-height}}

底层 \passthrough{\lstinline!unitree\_go::msg::dds\_::SportModeCmd\_!}
的 \passthrough{\lstinline!body\_height!} 字段。类型签名只说明 Python
表示；单位、数值范围、数组固定长度和枚举含义需查对应 IDL 头文件。
底层为浮点数；类型本身不说明单位、坐标系或业务安全范围。

\textbf{签名}

\begin{lstlisting}[language=Python]
@property
def body_height(self) -> float

@body_height.setter
def body_height(self, value: float) -> None
\end{lstlisting}

\textbf{可用性与安全性}

\textbf{\passthrough{\lstinline!AVAILABLE!} /
\passthrough{\lstinline!VALUE\_TYPE!}}：当前绑定源码已实现该消息值操作。

\textbf{参数}

\begin{longtable}[]{@{}
  >{\raggedright\arraybackslash}p{(\columnwidth - 4\tabcolsep) * \real{0.18}}
  >{\raggedright\arraybackslash}p{(\columnwidth - 4\tabcolsep) * \real{0.20}}
  >{\raggedright\arraybackslash}p{(\columnwidth - 4\tabcolsep) * \real{0.62}}@{}}
\toprule
名称 & Python 类型 & 含义 \\
\midrule
\endhead
\passthrough{\lstinline!value!} & \passthrough{\lstinline!float!} &
要写入或传递的新值，必须符合该参数的 Python 类型以及底层 C++
范围约束。 \\
\bottomrule
\end{longtable}

\textbf{返回值}

\begin{longtable}[]{@{}
  >{\raggedright\arraybackslash}p{(\columnwidth - 4\tabcolsep) * \real{0.18}}
  >{\raggedright\arraybackslash}p{(\columnwidth - 4\tabcolsep) * \real{0.20}}
  >{\raggedright\arraybackslash}p{(\columnwidth - 4\tabcolsep) * \real{0.62}}@{}}
\toprule
位置 & 类型 & 含义 \\
\midrule
\endhead
getter 返回值 & \passthrough{\lstinline!float!} & 返回属性当前值。 \\
\bottomrule
\end{longtable}

\textbf{对应 C++}

\begin{itemize}
\tightlist
\item
  类：\passthrough{\lstinline!unitree\_go::msg::dds\_::SportModeCmd\_!}
\item
  字段类型：\passthrough{\lstinline!float!}
\item
  声明位置：\passthrough{\lstinline!include/unitree/idl/go2/SportModeCmd\_.hpp:32!}
\end{itemize}

\textbf{用法}

\begin{lstlisting}[language=Python]
current_value = obj.body_height
obj.body_height = new_value
\end{lstlisting}

\hypertarget{unitree-sdk2-cpp-idl-go2-sportmodecmd-position}{%
\paragraph{\texorpdfstring{\texttt{unitree\_sdk2\_cpp.idl.go2.SportModeCmd.position}}{unitree\_sdk2\_cpp.idl.go2.SportModeCmd.position}}\label{unitree-sdk2-cpp-idl-go2-sportmodecmd-position}}

底层 \passthrough{\lstinline!unitree\_go::msg::dds\_::SportModeCmd\_!}
的 \passthrough{\lstinline!position!} 字段。类型签名只说明 Python
表示；单位、数值范围、数组固定长度和枚举含义需查对应 IDL 头文件。
底层是固定长度数组，写入序列必须正好包含 2
个元素；元素约束：底层为浮点数；类型本身不说明单位、坐标系或业务安全范围。

\textbf{签名}

\begin{lstlisting}[language=Python]
@property
def position(self) -> list[float]

@position.setter
def position(self, value: Sequence[float]) -> None
\end{lstlisting}

\textbf{可用性与安全性}

\textbf{\passthrough{\lstinline!AVAILABLE!} /
\passthrough{\lstinline!VALUE\_TYPE!}}：当前绑定源码已实现该消息值操作。

\textbf{参数}

\begin{longtable}[]{@{}
  >{\raggedright\arraybackslash}p{(\columnwidth - 4\tabcolsep) * \real{0.18}}
  >{\raggedright\arraybackslash}p{(\columnwidth - 4\tabcolsep) * \real{0.20}}
  >{\raggedright\arraybackslash}p{(\columnwidth - 4\tabcolsep) * \real{0.62}}@{}}
\toprule
名称 & Python 类型 & 含义 \\
\midrule
\endhead
\passthrough{\lstinline!value!} &
\passthrough{\lstinline!Sequence[float]!} &
要写入或传递的新值，必须符合该参数的 Python 类型以及底层 C++
范围约束。 \\
\bottomrule
\end{longtable}

\textbf{返回值}

\begin{longtable}[]{@{}
  >{\raggedright\arraybackslash}p{(\columnwidth - 4\tabcolsep) * \real{0.18}}
  >{\raggedright\arraybackslash}p{(\columnwidth - 4\tabcolsep) * \real{0.20}}
  >{\raggedright\arraybackslash}p{(\columnwidth - 4\tabcolsep) * \real{0.62}}@{}}
\toprule
位置 & 类型 & 含义 \\
\midrule
\endhead
getter 返回值 & \passthrough{\lstinline!list[float]!} &
返回属性当前值。数组、vector 和嵌套消息采用复制语义。 \\
\bottomrule
\end{longtable}

\textbf{对应 C++}

\begin{itemize}
\tightlist
\item
  类：\passthrough{\lstinline!unitree\_go::msg::dds\_::SportModeCmd\_!}
\item
  字段类型：\passthrough{\lstinline!std::array<float, 2>!}
\item
  声明位置：\passthrough{\lstinline!include/unitree/idl/go2/SportModeCmd\_.hpp:33!}
\end{itemize}

\textbf{用法}

\begin{lstlisting}[language=Python]
current_value = obj.position
obj.position = new_value
\end{lstlisting}

\begin{quote}
{[}!TIP{]} 读取容器或嵌套值后应遵循``读取、修改、重新赋值''。只修改
getter 返回的副本不会可靠地更新原 C++ 消息。
\end{quote}

\hypertarget{unitree-sdk2-cpp-idl-go2-sportmodecmd-euler}{%
\paragraph{\texorpdfstring{\texttt{unitree\_sdk2\_cpp.idl.go2.SportModeCmd.euler}}{unitree\_sdk2\_cpp.idl.go2.SportModeCmd.euler}}\label{unitree-sdk2-cpp-idl-go2-sportmodecmd-euler}}

底层 \passthrough{\lstinline!unitree\_go::msg::dds\_::SportModeCmd\_!}
的 \passthrough{\lstinline!euler!} 字段。类型签名只说明 Python
表示；单位、数值范围、数组固定长度和枚举含义需查对应 IDL 头文件。
底层是固定长度数组，写入序列必须正好包含 3
个元素；元素约束：底层为浮点数；类型本身不说明单位、坐标系或业务安全范围。

\textbf{签名}

\begin{lstlisting}[language=Python]
@property
def euler(self) -> list[float]

@euler.setter
def euler(self, value: Sequence[float]) -> None
\end{lstlisting}

\textbf{可用性与安全性}

\textbf{\passthrough{\lstinline!AVAILABLE!} /
\passthrough{\lstinline!VALUE\_TYPE!}}：当前绑定源码已实现该消息值操作。

\textbf{参数}

\begin{longtable}[]{@{}
  >{\raggedright\arraybackslash}p{(\columnwidth - 4\tabcolsep) * \real{0.18}}
  >{\raggedright\arraybackslash}p{(\columnwidth - 4\tabcolsep) * \real{0.20}}
  >{\raggedright\arraybackslash}p{(\columnwidth - 4\tabcolsep) * \real{0.62}}@{}}
\toprule
名称 & Python 类型 & 含义 \\
\midrule
\endhead
\passthrough{\lstinline!value!} &
\passthrough{\lstinline!Sequence[float]!} &
要写入或传递的新值，必须符合该参数的 Python 类型以及底层 C++
范围约束。 \\
\bottomrule
\end{longtable}

\textbf{返回值}

\begin{longtable}[]{@{}
  >{\raggedright\arraybackslash}p{(\columnwidth - 4\tabcolsep) * \real{0.18}}
  >{\raggedright\arraybackslash}p{(\columnwidth - 4\tabcolsep) * \real{0.20}}
  >{\raggedright\arraybackslash}p{(\columnwidth - 4\tabcolsep) * \real{0.62}}@{}}
\toprule
位置 & 类型 & 含义 \\
\midrule
\endhead
getter 返回值 & \passthrough{\lstinline!list[float]!} &
返回属性当前值。数组、vector 和嵌套消息采用复制语义。 \\
\bottomrule
\end{longtable}

\textbf{对应 C++}

\begin{itemize}
\tightlist
\item
  类：\passthrough{\lstinline!unitree\_go::msg::dds\_::SportModeCmd\_!}
\item
  字段类型：\passthrough{\lstinline!std::array<float, 3>!}
\item
  声明位置：\passthrough{\lstinline!include/unitree/idl/go2/SportModeCmd\_.hpp:34!}
\end{itemize}

\textbf{用法}

\begin{lstlisting}[language=Python]
current_value = obj.euler
obj.euler = new_value
\end{lstlisting}

\begin{quote}
{[}!TIP{]} 读取容器或嵌套值后应遵循``读取、修改、重新赋值''。只修改
getter 返回的副本不会可靠地更新原 C++ 消息。
\end{quote}

\hypertarget{unitree-sdk2-cpp-idl-go2-sportmodecmd-velocity}{%
\paragraph{\texorpdfstring{\texttt{unitree\_sdk2\_cpp.idl.go2.SportModeCmd.velocity}}{unitree\_sdk2\_cpp.idl.go2.SportModeCmd.velocity}}\label{unitree-sdk2-cpp-idl-go2-sportmodecmd-velocity}}

底层 \passthrough{\lstinline!unitree\_go::msg::dds\_::SportModeCmd\_!}
的 \passthrough{\lstinline!velocity!} 字段。类型签名只说明 Python
表示；单位、数值范围、数组固定长度和枚举含义需查对应 IDL 头文件。
底层是固定长度数组，写入序列必须正好包含 2
个元素；元素约束：底层为浮点数；类型本身不说明单位、坐标系或业务安全范围。

\textbf{签名}

\begin{lstlisting}[language=Python]
@property
def velocity(self) -> list[float]

@velocity.setter
def velocity(self, value: Sequence[float]) -> None
\end{lstlisting}

\textbf{可用性与安全性}

\textbf{\passthrough{\lstinline!AVAILABLE!} /
\passthrough{\lstinline!VALUE\_TYPE!}}：当前绑定源码已实现该消息值操作。

\textbf{参数}

\begin{longtable}[]{@{}
  >{\raggedright\arraybackslash}p{(\columnwidth - 4\tabcolsep) * \real{0.18}}
  >{\raggedright\arraybackslash}p{(\columnwidth - 4\tabcolsep) * \real{0.20}}
  >{\raggedright\arraybackslash}p{(\columnwidth - 4\tabcolsep) * \real{0.62}}@{}}
\toprule
名称 & Python 类型 & 含义 \\
\midrule
\endhead
\passthrough{\lstinline!value!} &
\passthrough{\lstinline!Sequence[float]!} &
要写入或传递的新值，必须符合该参数的 Python 类型以及底层 C++
范围约束。 \\
\bottomrule
\end{longtable}

\textbf{返回值}

\begin{longtable}[]{@{}
  >{\raggedright\arraybackslash}p{(\columnwidth - 4\tabcolsep) * \real{0.18}}
  >{\raggedright\arraybackslash}p{(\columnwidth - 4\tabcolsep) * \real{0.20}}
  >{\raggedright\arraybackslash}p{(\columnwidth - 4\tabcolsep) * \real{0.62}}@{}}
\toprule
位置 & 类型 & 含义 \\
\midrule
\endhead
getter 返回值 & \passthrough{\lstinline!list[float]!} &
返回属性当前值。数组、vector 和嵌套消息采用复制语义。 \\
\bottomrule
\end{longtable}

\textbf{对应 C++}

\begin{itemize}
\tightlist
\item
  类：\passthrough{\lstinline!unitree\_go::msg::dds\_::SportModeCmd\_!}
\item
  字段类型：\passthrough{\lstinline!std::array<float, 2>!}
\item
  声明位置：\passthrough{\lstinline!include/unitree/idl/go2/SportModeCmd\_.hpp:35!}
\end{itemize}

\textbf{用法}

\begin{lstlisting}[language=Python]
current_value = obj.velocity
obj.velocity = new_value
\end{lstlisting}

\begin{quote}
{[}!TIP{]} 读取容器或嵌套值后应遵循``读取、修改、重新赋值''。只修改
getter 返回的副本不会可靠地更新原 C++ 消息。
\end{quote}

\hypertarget{unitree-sdk2-cpp-idl-go2-sportmodecmd-yaw-speed}{%
\paragraph{\texorpdfstring{\texttt{unitree\_sdk2\_cpp.idl.go2.SportModeCmd.yaw\_speed}}{unitree\_sdk2\_cpp.idl.go2.SportModeCmd.yaw\_speed}}\label{unitree-sdk2-cpp-idl-go2-sportmodecmd-yaw-speed}}

底层 \passthrough{\lstinline!unitree\_go::msg::dds\_::SportModeCmd\_!}
的 \passthrough{\lstinline!yaw\_speed!} 字段。类型签名只说明 Python
表示；单位、数值范围、数组固定长度和枚举含义需查对应 IDL 头文件。
底层为浮点数；类型本身不说明单位、坐标系或业务安全范围。

\textbf{签名}

\begin{lstlisting}[language=Python]
@property
def yaw_speed(self) -> float

@yaw_speed.setter
def yaw_speed(self, value: float) -> None
\end{lstlisting}

\textbf{可用性与安全性}

\textbf{\passthrough{\lstinline!AVAILABLE!} /
\passthrough{\lstinline!VALUE\_TYPE!}}：当前绑定源码已实现该消息值操作。

\textbf{参数}

\begin{longtable}[]{@{}
  >{\raggedright\arraybackslash}p{(\columnwidth - 4\tabcolsep) * \real{0.18}}
  >{\raggedright\arraybackslash}p{(\columnwidth - 4\tabcolsep) * \real{0.20}}
  >{\raggedright\arraybackslash}p{(\columnwidth - 4\tabcolsep) * \real{0.62}}@{}}
\toprule
名称 & Python 类型 & 含义 \\
\midrule
\endhead
\passthrough{\lstinline!value!} & \passthrough{\lstinline!float!} &
要写入或传递的新值，必须符合该参数的 Python 类型以及底层 C++
范围约束。 \\
\bottomrule
\end{longtable}

\textbf{返回值}

\begin{longtable}[]{@{}
  >{\raggedright\arraybackslash}p{(\columnwidth - 4\tabcolsep) * \real{0.18}}
  >{\raggedright\arraybackslash}p{(\columnwidth - 4\tabcolsep) * \real{0.20}}
  >{\raggedright\arraybackslash}p{(\columnwidth - 4\tabcolsep) * \real{0.62}}@{}}
\toprule
位置 & 类型 & 含义 \\
\midrule
\endhead
getter 返回值 & \passthrough{\lstinline!float!} & 返回属性当前值。 \\
\bottomrule
\end{longtable}

\textbf{对应 C++}

\begin{itemize}
\tightlist
\item
  类：\passthrough{\lstinline!unitree\_go::msg::dds\_::SportModeCmd\_!}
\item
  字段类型：\passthrough{\lstinline!float!}
\item
  声明位置：\passthrough{\lstinline!include/unitree/idl/go2/SportModeCmd\_.hpp:36!}
\end{itemize}

\textbf{用法}

\begin{lstlisting}[language=Python]
current_value = obj.yaw_speed
obj.yaw_speed = new_value
\end{lstlisting}

\hypertarget{unitree-sdk2-cpp-idl-go2-sportmodecmd-bms-cmd}{%
\paragraph{\texorpdfstring{\texttt{unitree\_sdk2\_cpp.idl.go2.SportModeCmd.bms\_cmd}}{unitree\_sdk2\_cpp.idl.go2.SportModeCmd.bms\_cmd}}\label{unitree-sdk2-cpp-idl-go2-sportmodecmd-bms-cmd}}

底层 \passthrough{\lstinline!unitree\_go::msg::dds\_::SportModeCmd\_!}
的 \passthrough{\lstinline!bms\_cmd!} 字段。类型签名只说明 Python
表示；单位、数值范围、数组固定长度和枚举含义需查对应 IDL 头文件。

\textbf{签名}

\begin{lstlisting}[language=Python]
@property
def bms_cmd(self) -> BmsCmd

@bms_cmd.setter
def bms_cmd(self, value: BmsCmd) -> None
\end{lstlisting}

\textbf{可用性与安全性}

\textbf{\passthrough{\lstinline!AVAILABLE!} /
\passthrough{\lstinline!VALUE\_TYPE!}}：当前绑定源码已实现该消息值操作。

\textbf{参数}

\begin{longtable}[]{@{}
  >{\raggedright\arraybackslash}p{(\columnwidth - 4\tabcolsep) * \real{0.18}}
  >{\raggedright\arraybackslash}p{(\columnwidth - 4\tabcolsep) * \real{0.20}}
  >{\raggedright\arraybackslash}p{(\columnwidth - 4\tabcolsep) * \real{0.62}}@{}}
\toprule
名称 & Python 类型 & 含义 \\
\midrule
\endhead
\passthrough{\lstinline!value!} & \passthrough{\lstinline!BmsCmd!} &
要写入或传递的新值，必须符合该参数的 Python 类型以及底层 C++
范围约束。 \\
\bottomrule
\end{longtable}

\textbf{返回值}

\begin{longtable}[]{@{}
  >{\raggedright\arraybackslash}p{(\columnwidth - 4\tabcolsep) * \real{0.18}}
  >{\raggedright\arraybackslash}p{(\columnwidth - 4\tabcolsep) * \real{0.20}}
  >{\raggedright\arraybackslash}p{(\columnwidth - 4\tabcolsep) * \real{0.62}}@{}}
\toprule
位置 & 类型 & 含义 \\
\midrule
\endhead
getter 返回值 & \passthrough{\lstinline!BmsCmd!} &
返回属性当前值。数组、vector 和嵌套消息采用复制语义。 \\
\bottomrule
\end{longtable}

\textbf{对应 C++}

\begin{itemize}
\tightlist
\item
  类：\passthrough{\lstinline!unitree\_go::msg::dds\_::SportModeCmd\_!}
\item
  字段类型：\passthrough{\lstinline!::unitree\_go::msg::dds\_::BmsCmd\_!}
\item
  声明位置：\passthrough{\lstinline!include/unitree/idl/go2/SportModeCmd\_.hpp:37!}
\end{itemize}

\textbf{用法}

\begin{lstlisting}[language=Python]
current_value = obj.bms_cmd
obj.bms_cmd = new_value
\end{lstlisting}

\hypertarget{unitree-sdk2-cpp-idl-go2-sportmodecmd-path-point}{%
\paragraph{\texorpdfstring{\texttt{unitree\_sdk2\_cpp.idl.go2.SportModeCmd.path\_point}}{unitree\_sdk2\_cpp.idl.go2.SportModeCmd.path\_point}}\label{unitree-sdk2-cpp-idl-go2-sportmodecmd-path-point}}

底层 \passthrough{\lstinline!unitree\_go::msg::dds\_::SportModeCmd\_!}
的 \passthrough{\lstinline!path\_point!} 字段。类型签名只说明 Python
表示；单位、数值范围、数组固定长度和枚举含义需查对应 IDL 头文件。
底层是固定长度数组，写入序列必须正好包含 30 个元素

\textbf{签名}

\begin{lstlisting}[language=Python]
@property
def path_point(self) -> list[PathPoint]

@path_point.setter
def path_point(self, value: Sequence[PathPoint]) -> None
\end{lstlisting}

\textbf{可用性与安全性}

\textbf{\passthrough{\lstinline!AVAILABLE!} /
\passthrough{\lstinline!VALUE\_TYPE!}}：当前绑定源码已实现该消息值操作。

\textbf{参数}

\begin{longtable}[]{@{}
  >{\raggedright\arraybackslash}p{(\columnwidth - 4\tabcolsep) * \real{0.18}}
  >{\raggedright\arraybackslash}p{(\columnwidth - 4\tabcolsep) * \real{0.20}}
  >{\raggedright\arraybackslash}p{(\columnwidth - 4\tabcolsep) * \real{0.62}}@{}}
\toprule
名称 & Python 类型 & 含义 \\
\midrule
\endhead
\passthrough{\lstinline!value!} &
\passthrough{\lstinline!Sequence[PathPoint]!} &
要写入或传递的新值，必须符合该参数的 Python 类型以及底层 C++
范围约束。 \\
\bottomrule
\end{longtable}

\textbf{返回值}

\begin{longtable}[]{@{}
  >{\raggedright\arraybackslash}p{(\columnwidth - 4\tabcolsep) * \real{0.18}}
  >{\raggedright\arraybackslash}p{(\columnwidth - 4\tabcolsep) * \real{0.20}}
  >{\raggedright\arraybackslash}p{(\columnwidth - 4\tabcolsep) * \real{0.62}}@{}}
\toprule
位置 & 类型 & 含义 \\
\midrule
\endhead
getter 返回值 & \passthrough{\lstinline!list[PathPoint]!} &
返回属性当前值。数组、vector 和嵌套消息采用复制语义。 \\
\bottomrule
\end{longtable}

\textbf{对应 C++}

\begin{itemize}
\tightlist
\item
  类：\passthrough{\lstinline!unitree\_go::msg::dds\_::SportModeCmd\_!}
\item
  字段类型：\passthrough{\lstinline!std::array< ::unitree\_go::msg::dds\_::PathPoint\_, 30>!}
\item
  声明位置：\passthrough{\lstinline!include/unitree/idl/go2/SportModeCmd\_.hpp:38!}
\end{itemize}

\textbf{用法}

\begin{lstlisting}[language=Python]
current_value = obj.path_point
obj.path_point = new_value
\end{lstlisting}

\begin{quote}
{[}!TIP{]} 读取容器或嵌套值后应遵循``读取、修改、重新赋值''。只修改
getter 返回的副本不会可靠地更新原 C++ 消息。
\end{quote}

\hypertarget{unitree-sdk2-cpp-idl-go2-timespec}{%
\subsubsection{\texorpdfstring{\texttt{unitree\_sdk2\_cpp.idl.go2.TimeSpec}}{unitree\_sdk2\_cpp.idl.go2.TimeSpec}}\label{unitree-sdk2-cpp-idl-go2-timespec}}

可由 Python 读写的 DDS/IDL 消息值类型。字段采用复制语义。

\textbf{导入}

\begin{lstlisting}[language=Python]
from unitree_sdk2_cpp.idl.go2 import TimeSpec
\end{lstlisting}

\textbf{方法索引}

\begin{longtable}[]{@{}
  >{\raggedright\arraybackslash}p{(\columnwidth - 6\tabcolsep) * \real{0.18}}
  >{\raggedright\arraybackslash}p{(\columnwidth - 6\tabcolsep) * \real{0.24}}
  >{\raggedright\arraybackslash}p{(\columnwidth - 6\tabcolsep) * \real{0.18}}
  >{\raggedright\arraybackslash}p{(\columnwidth - 6\tabcolsep) * \real{0.40}}@{}}
\toprule
方法 & Python 签名 & 状态 & 安全分类 \\
\midrule
\endhead
\protect\hyperlink{unitree-sdk2-cpp-idl-go2-timespec-dunder-init-1}{\passthrough{\lstinline!\_\_init\_\_!}}
& \passthrough{\lstinline!def \_\_init\_\_(self) -> None!} &
\passthrough{\lstinline!AVAILABLE!} &
\passthrough{\lstinline!VALUE\_TYPE!} \\
\protect\hyperlink{unitree-sdk2-cpp-idl-go2-timespec-dunder-eq-1}{\passthrough{\lstinline!\_\_eq\_\_!}}
& \passthrough{\lstinline!def \_\_eq\_\_(self, other: object) -> bool!}
& \passthrough{\lstinline!AVAILABLE!} &
\passthrough{\lstinline!VALUE\_TYPE!} \\
\protect\hyperlink{unitree-sdk2-cpp-idl-go2-timespec-dunder-ne-1}{\passthrough{\lstinline!\_\_ne\_\_!}}
& \passthrough{\lstinline!def \_\_ne\_\_(self, other: object) -> bool!}
& \passthrough{\lstinline!AVAILABLE!} &
\passthrough{\lstinline!VALUE\_TYPE!} \\
\bottomrule
\end{longtable}

\textbf{属性索引}

\begin{longtable}[]{@{}
  >{\raggedright\arraybackslash}p{(\columnwidth - 6\tabcolsep) * \real{0.18}}
  >{\raggedright\arraybackslash}p{(\columnwidth - 6\tabcolsep) * \real{0.24}}
  >{\raggedright\arraybackslash}p{(\columnwidth - 6\tabcolsep) * \real{0.18}}
  >{\raggedright\arraybackslash}p{(\columnwidth - 6\tabcolsep) * \real{0.40}}@{}}
\toprule
属性 & 读取类型 & 写入类型 & 状态 \\
\midrule
\endhead
\protect\hyperlink{unitree-sdk2-cpp-idl-go2-timespec-sec}{\passthrough{\lstinline!sec!}}
& \passthrough{\lstinline!int!} & \passthrough{\lstinline!int!} &
\passthrough{\lstinline!AVAILABLE!} \\
\protect\hyperlink{unitree-sdk2-cpp-idl-go2-timespec-nanosec}{\passthrough{\lstinline!nanosec!}}
& \passthrough{\lstinline!int!} & \passthrough{\lstinline!int!} &
\passthrough{\lstinline!AVAILABLE!} \\
\bottomrule
\end{longtable}

\hypertarget{unitree-sdk2-cpp-idl-go2-timespec-dunder-init-1}{%
\paragraph{\texorpdfstring{\texttt{unitree\_sdk2\_cpp.idl.go2.TimeSpec.\_\_init\_\_}}{unitree\_sdk2\_cpp.idl.go2.TimeSpec.\_\_init\_\_}}\label{unitree-sdk2-cpp-idl-go2-timespec-dunder-init-1}}

初始化 \passthrough{\lstinline!TimeSpec!}
实例。是否能实际构造取决于下方可用性状态。

\textbf{签名}

\begin{lstlisting}[language=Python]
def __init__(self) -> None
\end{lstlisting}

\textbf{可用性与安全性}

\textbf{\passthrough{\lstinline!AVAILABLE!} /
\passthrough{\lstinline!VALUE\_TYPE!}}：当前绑定源码已实现该消息值操作。

\textbf{参数}

无显式参数。实例方法中的 \passthrough{\lstinline!self!} 由 Python
自动传入。

\textbf{返回值}

\begin{longtable}[]{@{}
  >{\raggedright\arraybackslash}p{(\columnwidth - 4\tabcolsep) * \real{0.18}}
  >{\raggedright\arraybackslash}p{(\columnwidth - 4\tabcolsep) * \real{0.20}}
  >{\raggedright\arraybackslash}p{(\columnwidth - 4\tabcolsep) * \real{0.62}}@{}}
\toprule
位置 & 类型 & 含义 \\
\midrule
\endhead
无 & \passthrough{\lstinline!None!} &
\passthrough{\lstinline!\_\_init\_\_!}
本身不返回值；调用类对象时，在构造函数可用的前提下得到该类实例。 \\
\bottomrule
\end{longtable}

\textbf{对应 C++}

\begin{itemize}
\tightlist
\item
  类：\passthrough{\lstinline!unitree\_go::msg::dds\_::TimeSpec\_!}
\item
  签名：\passthrough{\lstinline!TimeSpec\_()!}
\item
  绑定策略：\passthrough{\lstinline!未单独标注!}
\item
  声明位置：\passthrough{\lstinline!include/unitree/idl/go2/TimeSpec\_.hpp:27!}
\end{itemize}

\textbf{说明}

\begin{itemize}
\tightlist
\item
  示例是局部调用片段；\passthrough{\lstinline!obj!}
  和参数变量表示已经按上层业务校验并准备好的对象。
\end{itemize}

\textbf{用法}

\begin{lstlisting}[language=Python]
value = TimeSpec()
\end{lstlisting}

\hypertarget{unitree-sdk2-cpp-idl-go2-timespec-dunder-eq-1}{%
\paragraph{\texorpdfstring{\texttt{unitree\_sdk2\_cpp.idl.go2.TimeSpec.\_\_eq\_\_}}{unitree\_sdk2\_cpp.idl.go2.TimeSpec.\_\_eq\_\_}}\label{unitree-sdk2-cpp-idl-go2-timespec-dunder-eq-1}}

按底层 IDL 消息内容比较两个对象是否相等。

\textbf{签名}

\begin{lstlisting}[language=Python]
def __eq__(self, other: object) -> bool
\end{lstlisting}

\textbf{可用性与安全性}

\textbf{\passthrough{\lstinline!AVAILABLE!} /
\passthrough{\lstinline!VALUE\_TYPE!}}：当前绑定源码已实现该消息值操作。

\textbf{参数}

\begin{longtable}[]{@{}
  >{\raggedright\arraybackslash}p{(\columnwidth - 6\tabcolsep) * \real{0.18}}
  >{\raggedright\arraybackslash}p{(\columnwidth - 6\tabcolsep) * \real{0.24}}
  >{\raggedright\arraybackslash}p{(\columnwidth - 6\tabcolsep) * \real{0.18}}
  >{\raggedright\arraybackslash}p{(\columnwidth - 6\tabcolsep) * \real{0.40}}@{}}
\toprule
名称 & Python 类型 & 默认值 & 含义 \\
\midrule
\endhead
\passthrough{\lstinline!other!} & \passthrough{\lstinline!object!} &
必填 & 用于比较的另一个 Python 对象。类型不兼容时比较结果通常为
\passthrough{\lstinline!False!}。 \\
\bottomrule
\end{longtable}

\textbf{返回值}

\begin{longtable}[]{@{}
  >{\raggedright\arraybackslash}p{(\columnwidth - 4\tabcolsep) * \real{0.18}}
  >{\raggedright\arraybackslash}p{(\columnwidth - 4\tabcolsep) * \real{0.20}}
  >{\raggedright\arraybackslash}p{(\columnwidth - 4\tabcolsep) * \real{0.62}}@{}}
\toprule
位置 & 类型 & 含义 \\
\midrule
\endhead
返回值 & \passthrough{\lstinline!bool!} & 返回
\passthrough{\lstinline!bool!}。更精确的业务含义见该方法用途、C++
签名和目标型号协议。 \\
\bottomrule
\end{longtable}

\textbf{对应 C++}

\begin{itemize}
\tightlist
\item
  类：\passthrough{\lstinline!unitree\_go::msg::dds\_::TimeSpec\_!}
\item
  签名：\passthrough{\lstinline!operator==(const value \&) const!}
\item
  绑定策略：\passthrough{\lstinline!未单独标注!}
\item
  声明位置：由 pybind11 手工绑定或生成注册表提供；manifest
  未记录独立头文件行号。
\end{itemize}

\textbf{说明}

\begin{itemize}
\tightlist
\item
  示例是局部调用片段；\passthrough{\lstinline!obj!}
  和参数变量表示已经按上层业务校验并准备好的对象。
\end{itemize}

\textbf{用法}

\begin{lstlisting}[language=Python]
same = left == right
\end{lstlisting}

\hypertarget{unitree-sdk2-cpp-idl-go2-timespec-dunder-ne-1}{%
\paragraph{\texorpdfstring{\texttt{unitree\_sdk2\_cpp.idl.go2.TimeSpec.\_\_ne\_\_}}{unitree\_sdk2\_cpp.idl.go2.TimeSpec.\_\_ne\_\_}}\label{unitree-sdk2-cpp-idl-go2-timespec-dunder-ne-1}}

按底层 IDL 消息内容比较两个对象是否不相等。

\textbf{签名}

\begin{lstlisting}[language=Python]
def __ne__(self, other: object) -> bool
\end{lstlisting}

\textbf{可用性与安全性}

\textbf{\passthrough{\lstinline!AVAILABLE!} /
\passthrough{\lstinline!VALUE\_TYPE!}}：当前绑定源码已实现该消息值操作。

\textbf{参数}

\begin{longtable}[]{@{}
  >{\raggedright\arraybackslash}p{(\columnwidth - 6\tabcolsep) * \real{0.18}}
  >{\raggedright\arraybackslash}p{(\columnwidth - 6\tabcolsep) * \real{0.24}}
  >{\raggedright\arraybackslash}p{(\columnwidth - 6\tabcolsep) * \real{0.18}}
  >{\raggedright\arraybackslash}p{(\columnwidth - 6\tabcolsep) * \real{0.40}}@{}}
\toprule
名称 & Python 类型 & 默认值 & 含义 \\
\midrule
\endhead
\passthrough{\lstinline!other!} & \passthrough{\lstinline!object!} &
必填 & 用于比较的另一个 Python 对象。类型不兼容时比较结果通常为
\passthrough{\lstinline!False!}。 \\
\bottomrule
\end{longtable}

\textbf{返回值}

\begin{longtable}[]{@{}
  >{\raggedright\arraybackslash}p{(\columnwidth - 4\tabcolsep) * \real{0.18}}
  >{\raggedright\arraybackslash}p{(\columnwidth - 4\tabcolsep) * \real{0.20}}
  >{\raggedright\arraybackslash}p{(\columnwidth - 4\tabcolsep) * \real{0.62}}@{}}
\toprule
位置 & 类型 & 含义 \\
\midrule
\endhead
返回值 & \passthrough{\lstinline!bool!} & 返回
\passthrough{\lstinline!bool!}。更精确的业务含义见该方法用途、C++
签名和目标型号协议。 \\
\bottomrule
\end{longtable}

\textbf{对应 C++}

\begin{itemize}
\tightlist
\item
  类：\passthrough{\lstinline!unitree\_go::msg::dds\_::TimeSpec\_!}
\item
  签名：\passthrough{\lstinline"operator!=(const value \&) const"}
\item
  绑定策略：\passthrough{\lstinline!未单独标注!}
\item
  声明位置：由 pybind11 手工绑定或生成注册表提供；manifest
  未记录独立头文件行号。
\end{itemize}

\textbf{说明}

\begin{itemize}
\tightlist
\item
  示例是局部调用片段；\passthrough{\lstinline!obj!}
  和参数变量表示已经按上层业务校验并准备好的对象。
\end{itemize}

\textbf{用法}

\begin{lstlisting}[language=Python]
different = left != right
\end{lstlisting}

\hypertarget{unitree-sdk2-cpp-idl-go2-timespec-sec}{%
\paragraph{\texorpdfstring{\texttt{unitree\_sdk2\_cpp.idl.go2.TimeSpec.sec}}{unitree\_sdk2\_cpp.idl.go2.TimeSpec.sec}}\label{unitree-sdk2-cpp-idl-go2-timespec-sec}}

底层 \passthrough{\lstinline!unitree\_go::msg::dds\_::TimeSpec\_!} 的
\passthrough{\lstinline!sec!} 字段。类型签名只说明 Python
表示；单位、数值范围、数组固定长度和枚举含义需查对应 IDL 头文件。
取值范围为 -2147483648 到 2147483647。

\textbf{签名}

\begin{lstlisting}[language=Python]
@property
def sec(self) -> int

@sec.setter
def sec(self, value: int) -> None
\end{lstlisting}

\textbf{可用性与安全性}

\textbf{\passthrough{\lstinline!AVAILABLE!} /
\passthrough{\lstinline!VALUE\_TYPE!}}：当前绑定源码已实现该消息值操作。

\textbf{参数}

\begin{longtable}[]{@{}
  >{\raggedright\arraybackslash}p{(\columnwidth - 4\tabcolsep) * \real{0.18}}
  >{\raggedright\arraybackslash}p{(\columnwidth - 4\tabcolsep) * \real{0.20}}
  >{\raggedright\arraybackslash}p{(\columnwidth - 4\tabcolsep) * \real{0.62}}@{}}
\toprule
名称 & Python 类型 & 含义 \\
\midrule
\endhead
\passthrough{\lstinline!value!} & \passthrough{\lstinline!int!} &
要写入或传递的新值，必须符合该参数的 Python 类型以及底层 C++
范围约束。 \\
\bottomrule
\end{longtable}

\textbf{返回值}

\begin{longtable}[]{@{}
  >{\raggedright\arraybackslash}p{(\columnwidth - 4\tabcolsep) * \real{0.18}}
  >{\raggedright\arraybackslash}p{(\columnwidth - 4\tabcolsep) * \real{0.20}}
  >{\raggedright\arraybackslash}p{(\columnwidth - 4\tabcolsep) * \real{0.62}}@{}}
\toprule
位置 & 类型 & 含义 \\
\midrule
\endhead
getter 返回值 & \passthrough{\lstinline!int!} & 返回属性当前值。 \\
\bottomrule
\end{longtable}

\textbf{对应 C++}

\begin{itemize}
\tightlist
\item
  类：\passthrough{\lstinline!unitree\_go::msg::dds\_::TimeSpec\_!}
\item
  字段类型：\passthrough{\lstinline!int32\_t!}
\item
  声明位置：\passthrough{\lstinline!include/unitree/idl/go2/TimeSpec\_.hpp:23!}
\end{itemize}

\textbf{用法}

\begin{lstlisting}[language=Python]
current_value = obj.sec
obj.sec = new_value
\end{lstlisting}

\hypertarget{unitree-sdk2-cpp-idl-go2-timespec-nanosec}{%
\paragraph{\texorpdfstring{\texttt{unitree\_sdk2\_cpp.idl.go2.TimeSpec.nanosec}}{unitree\_sdk2\_cpp.idl.go2.TimeSpec.nanosec}}\label{unitree-sdk2-cpp-idl-go2-timespec-nanosec}}

底层 \passthrough{\lstinline!unitree\_go::msg::dds\_::TimeSpec\_!} 的
\passthrough{\lstinline!nanosec!} 字段。类型签名只说明 Python
表示；单位、数值范围、数组固定长度和枚举含义需查对应 IDL 头文件。
取值范围为 0 到 4294967295。

\textbf{签名}

\begin{lstlisting}[language=Python]
@property
def nanosec(self) -> int

@nanosec.setter
def nanosec(self, value: int) -> None
\end{lstlisting}

\textbf{可用性与安全性}

\textbf{\passthrough{\lstinline!AVAILABLE!} /
\passthrough{\lstinline!VALUE\_TYPE!}}：当前绑定源码已实现该消息值操作。

\textbf{参数}

\begin{longtable}[]{@{}
  >{\raggedright\arraybackslash}p{(\columnwidth - 4\tabcolsep) * \real{0.18}}
  >{\raggedright\arraybackslash}p{(\columnwidth - 4\tabcolsep) * \real{0.20}}
  >{\raggedright\arraybackslash}p{(\columnwidth - 4\tabcolsep) * \real{0.62}}@{}}
\toprule
名称 & Python 类型 & 含义 \\
\midrule
\endhead
\passthrough{\lstinline!value!} & \passthrough{\lstinline!int!} &
要写入或传递的新值，必须符合该参数的 Python 类型以及底层 C++
范围约束。 \\
\bottomrule
\end{longtable}

\textbf{返回值}

\begin{longtable}[]{@{}
  >{\raggedright\arraybackslash}p{(\columnwidth - 4\tabcolsep) * \real{0.18}}
  >{\raggedright\arraybackslash}p{(\columnwidth - 4\tabcolsep) * \real{0.20}}
  >{\raggedright\arraybackslash}p{(\columnwidth - 4\tabcolsep) * \real{0.62}}@{}}
\toprule
位置 & 类型 & 含义 \\
\midrule
\endhead
getter 返回值 & \passthrough{\lstinline!int!} & 返回属性当前值。 \\
\bottomrule
\end{longtable}

\textbf{对应 C++}

\begin{itemize}
\tightlist
\item
  类：\passthrough{\lstinline!unitree\_go::msg::dds\_::TimeSpec\_!}
\item
  字段类型：\passthrough{\lstinline!uint32\_t!}
\item
  声明位置：\passthrough{\lstinline!include/unitree/idl/go2/TimeSpec\_.hpp:24!}
\end{itemize}

\textbf{用法}

\begin{lstlisting}[language=Python]
current_value = obj.nanosec
obj.nanosec = new_value
\end{lstlisting}

\hypertarget{unitree-sdk2-cpp-idl-go2-sportmodestate}{%
\subsubsection{\texorpdfstring{\texttt{unitree\_sdk2\_cpp.idl.go2.SportModeState}}{unitree\_sdk2\_cpp.idl.go2.SportModeState}}\label{unitree-sdk2-cpp-idl-go2-sportmodestate}}

可由 Python 读写的 DDS/IDL 消息值类型。字段采用复制语义。

\textbf{导入}

\begin{lstlisting}[language=Python]
from unitree_sdk2_cpp.idl.go2 import SportModeState
\end{lstlisting}

\textbf{方法索引}

\begin{longtable}[]{@{}
  >{\raggedright\arraybackslash}p{(\columnwidth - 6\tabcolsep) * \real{0.18}}
  >{\raggedright\arraybackslash}p{(\columnwidth - 6\tabcolsep) * \real{0.24}}
  >{\raggedright\arraybackslash}p{(\columnwidth - 6\tabcolsep) * \real{0.18}}
  >{\raggedright\arraybackslash}p{(\columnwidth - 6\tabcolsep) * \real{0.40}}@{}}
\toprule
方法 & Python 签名 & 状态 & 安全分类 \\
\midrule
\endhead
\protect\hyperlink{unitree-sdk2-cpp-idl-go2-sportmodestate-dunder-init-1}{\passthrough{\lstinline!\_\_init\_\_!}}
& \passthrough{\lstinline!def \_\_init\_\_(self) -> None!} &
\passthrough{\lstinline!AVAILABLE!} &
\passthrough{\lstinline!VALUE\_TYPE!} \\
\protect\hyperlink{unitree-sdk2-cpp-idl-go2-sportmodestate-dunder-eq-1}{\passthrough{\lstinline!\_\_eq\_\_!}}
& \passthrough{\lstinline!def \_\_eq\_\_(self, other: object) -> bool!}
& \passthrough{\lstinline!AVAILABLE!} &
\passthrough{\lstinline!VALUE\_TYPE!} \\
\protect\hyperlink{unitree-sdk2-cpp-idl-go2-sportmodestate-dunder-ne-1}{\passthrough{\lstinline!\_\_ne\_\_!}}
& \passthrough{\lstinline!def \_\_ne\_\_(self, other: object) -> bool!}
& \passthrough{\lstinline!AVAILABLE!} &
\passthrough{\lstinline!VALUE\_TYPE!} \\
\bottomrule
\end{longtable}

\textbf{属性索引}

\begin{longtable}[]{@{}
  >{\raggedright\arraybackslash}p{(\columnwidth - 6\tabcolsep) * \real{0.18}}
  >{\raggedright\arraybackslash}p{(\columnwidth - 6\tabcolsep) * \real{0.24}}
  >{\raggedright\arraybackslash}p{(\columnwidth - 6\tabcolsep) * \real{0.18}}
  >{\raggedright\arraybackslash}p{(\columnwidth - 6\tabcolsep) * \real{0.40}}@{}}
\toprule
属性 & 读取类型 & 写入类型 & 状态 \\
\midrule
\endhead
\protect\hyperlink{unitree-sdk2-cpp-idl-go2-sportmodestate-stamp}{\passthrough{\lstinline!stamp!}}
& \passthrough{\lstinline!TimeSpec!} &
\passthrough{\lstinline!TimeSpec!} &
\passthrough{\lstinline!AVAILABLE!} \\
\protect\hyperlink{unitree-sdk2-cpp-idl-go2-sportmodestate-error-code}{\passthrough{\lstinline!error\_code!}}
& \passthrough{\lstinline!int!} & \passthrough{\lstinline!int!} &
\passthrough{\lstinline!AVAILABLE!} \\
\protect\hyperlink{unitree-sdk2-cpp-idl-go2-sportmodestate-imu-state}{\passthrough{\lstinline!imu\_state!}}
& \passthrough{\lstinline!IMUState!} &
\passthrough{\lstinline!IMUState!} &
\passthrough{\lstinline!AVAILABLE!} \\
\protect\hyperlink{unitree-sdk2-cpp-idl-go2-sportmodestate-mode}{\passthrough{\lstinline!mode!}}
& \passthrough{\lstinline!int!} & \passthrough{\lstinline!int!} &
\passthrough{\lstinline!AVAILABLE!} \\
\protect\hyperlink{unitree-sdk2-cpp-idl-go2-sportmodestate-progress}{\passthrough{\lstinline!progress!}}
& \passthrough{\lstinline!float!} & \passthrough{\lstinline!float!} &
\passthrough{\lstinline!AVAILABLE!} \\
\protect\hyperlink{unitree-sdk2-cpp-idl-go2-sportmodestate-gait-type}{\passthrough{\lstinline!gait\_type!}}
& \passthrough{\lstinline!int!} & \passthrough{\lstinline!int!} &
\passthrough{\lstinline!AVAILABLE!} \\
\protect\hyperlink{unitree-sdk2-cpp-idl-go2-sportmodestate-foot-raise-height}{\passthrough{\lstinline!foot\_raise\_height!}}
& \passthrough{\lstinline!float!} & \passthrough{\lstinline!float!} &
\passthrough{\lstinline!AVAILABLE!} \\
\protect\hyperlink{unitree-sdk2-cpp-idl-go2-sportmodestate-position}{\passthrough{\lstinline!position!}}
& \passthrough{\lstinline!list[float]!} &
\passthrough{\lstinline!Sequence[float]!} &
\passthrough{\lstinline!AVAILABLE!} \\
\protect\hyperlink{unitree-sdk2-cpp-idl-go2-sportmodestate-body-height}{\passthrough{\lstinline!body\_height!}}
& \passthrough{\lstinline!float!} & \passthrough{\lstinline!float!} &
\passthrough{\lstinline!AVAILABLE!} \\
\protect\hyperlink{unitree-sdk2-cpp-idl-go2-sportmodestate-velocity}{\passthrough{\lstinline!velocity!}}
& \passthrough{\lstinline!list[float]!} &
\passthrough{\lstinline!Sequence[float]!} &
\passthrough{\lstinline!AVAILABLE!} \\
\protect\hyperlink{unitree-sdk2-cpp-idl-go2-sportmodestate-yaw-speed}{\passthrough{\lstinline!yaw\_speed!}}
& \passthrough{\lstinline!float!} & \passthrough{\lstinline!float!} &
\passthrough{\lstinline!AVAILABLE!} \\
\protect\hyperlink{unitree-sdk2-cpp-idl-go2-sportmodestate-range-obstacle}{\passthrough{\lstinline!range\_obstacle!}}
& \passthrough{\lstinline!list[float]!} &
\passthrough{\lstinline!Sequence[float]!} &
\passthrough{\lstinline!AVAILABLE!} \\
\protect\hyperlink{unitree-sdk2-cpp-idl-go2-sportmodestate-foot-force}{\passthrough{\lstinline!foot\_force!}}
& \passthrough{\lstinline!list[int]!} &
\passthrough{\lstinline!Sequence[int]!} &
\passthrough{\lstinline!AVAILABLE!} \\
\protect\hyperlink{unitree-sdk2-cpp-idl-go2-sportmodestate-foot-position-body}{\passthrough{\lstinline!foot\_position\_body!}}
& \passthrough{\lstinline!list[float]!} &
\passthrough{\lstinline!Sequence[float]!} &
\passthrough{\lstinline!AVAILABLE!} \\
\protect\hyperlink{unitree-sdk2-cpp-idl-go2-sportmodestate-foot-speed-body}{\passthrough{\lstinline!foot\_speed\_body!}}
& \passthrough{\lstinline!list[float]!} &
\passthrough{\lstinline!Sequence[float]!} &
\passthrough{\lstinline!AVAILABLE!} \\
\protect\hyperlink{unitree-sdk2-cpp-idl-go2-sportmodestate-path-point}{\passthrough{\lstinline!path\_point!}}
& \passthrough{\lstinline!list[PathPoint]!} &
\passthrough{\lstinline!Sequence[PathPoint]!} &
\passthrough{\lstinline!AVAILABLE!} \\
\bottomrule
\end{longtable}

\hypertarget{unitree-sdk2-cpp-idl-go2-sportmodestate-dunder-init-1}{%
\paragraph{\texorpdfstring{\texttt{unitree\_sdk2\_cpp.idl.go2.SportModeState.\_\_init\_\_}}{unitree\_sdk2\_cpp.idl.go2.SportModeState.\_\_init\_\_}}\label{unitree-sdk2-cpp-idl-go2-sportmodestate-dunder-init-1}}

初始化 \passthrough{\lstinline!SportModeState!}
实例。是否能实际构造取决于下方可用性状态。

\textbf{签名}

\begin{lstlisting}[language=Python]
def __init__(self) -> None
\end{lstlisting}

\textbf{可用性与安全性}

\textbf{\passthrough{\lstinline!AVAILABLE!} /
\passthrough{\lstinline!VALUE\_TYPE!}}：当前绑定源码已实现该消息值操作。

\textbf{参数}

无显式参数。实例方法中的 \passthrough{\lstinline!self!} 由 Python
自动传入。

\textbf{返回值}

\begin{longtable}[]{@{}
  >{\raggedright\arraybackslash}p{(\columnwidth - 4\tabcolsep) * \real{0.18}}
  >{\raggedright\arraybackslash}p{(\columnwidth - 4\tabcolsep) * \real{0.20}}
  >{\raggedright\arraybackslash}p{(\columnwidth - 4\tabcolsep) * \real{0.62}}@{}}
\toprule
位置 & 类型 & 含义 \\
\midrule
\endhead
无 & \passthrough{\lstinline!None!} &
\passthrough{\lstinline!\_\_init\_\_!}
本身不返回值；调用类对象时，在构造函数可用的前提下得到该类实例。 \\
\bottomrule
\end{longtable}

\textbf{对应 C++}

\begin{itemize}
\tightlist
\item
  类：\passthrough{\lstinline!unitree\_go::msg::dds\_::SportModeState\_!}
\item
  签名：\passthrough{\lstinline!SportModeState\_()!}
\item
  绑定策略：\passthrough{\lstinline!未单独标注!}
\item
  声明位置：\passthrough{\lstinline!include/unitree/idl/go2/SportModeState\_.hpp:48!}
\end{itemize}

\textbf{说明}

\begin{itemize}
\tightlist
\item
  示例是局部调用片段；\passthrough{\lstinline!obj!}
  和参数变量表示已经按上层业务校验并准备好的对象。
\end{itemize}

\textbf{用法}

\begin{lstlisting}[language=Python]
value = SportModeState()
\end{lstlisting}

\hypertarget{unitree-sdk2-cpp-idl-go2-sportmodestate-dunder-eq-1}{%
\paragraph{\texorpdfstring{\texttt{unitree\_sdk2\_cpp.idl.go2.SportModeState.\_\_eq\_\_}}{unitree\_sdk2\_cpp.idl.go2.SportModeState.\_\_eq\_\_}}\label{unitree-sdk2-cpp-idl-go2-sportmodestate-dunder-eq-1}}

按底层 IDL 消息内容比较两个对象是否相等。

\textbf{签名}

\begin{lstlisting}[language=Python]
def __eq__(self, other: object) -> bool
\end{lstlisting}

\textbf{可用性与安全性}

\textbf{\passthrough{\lstinline!AVAILABLE!} /
\passthrough{\lstinline!VALUE\_TYPE!}}：当前绑定源码已实现该消息值操作。

\textbf{参数}

\begin{longtable}[]{@{}
  >{\raggedright\arraybackslash}p{(\columnwidth - 6\tabcolsep) * \real{0.18}}
  >{\raggedright\arraybackslash}p{(\columnwidth - 6\tabcolsep) * \real{0.24}}
  >{\raggedright\arraybackslash}p{(\columnwidth - 6\tabcolsep) * \real{0.18}}
  >{\raggedright\arraybackslash}p{(\columnwidth - 6\tabcolsep) * \real{0.40}}@{}}
\toprule
名称 & Python 类型 & 默认值 & 含义 \\
\midrule
\endhead
\passthrough{\lstinline!other!} & \passthrough{\lstinline!object!} &
必填 & 用于比较的另一个 Python 对象。类型不兼容时比较结果通常为
\passthrough{\lstinline!False!}。 \\
\bottomrule
\end{longtable}

\textbf{返回值}

\begin{longtable}[]{@{}
  >{\raggedright\arraybackslash}p{(\columnwidth - 4\tabcolsep) * \real{0.18}}
  >{\raggedright\arraybackslash}p{(\columnwidth - 4\tabcolsep) * \real{0.20}}
  >{\raggedright\arraybackslash}p{(\columnwidth - 4\tabcolsep) * \real{0.62}}@{}}
\toprule
位置 & 类型 & 含义 \\
\midrule
\endhead
返回值 & \passthrough{\lstinline!bool!} & 返回
\passthrough{\lstinline!bool!}。更精确的业务含义见该方法用途、C++
签名和目标型号协议。 \\
\bottomrule
\end{longtable}

\textbf{对应 C++}

\begin{itemize}
\tightlist
\item
  类：\passthrough{\lstinline!unitree\_go::msg::dds\_::SportModeState\_!}
\item
  签名：\passthrough{\lstinline!operator==(const value \&) const!}
\item
  绑定策略：\passthrough{\lstinline!未单独标注!}
\item
  声明位置：由 pybind11 手工绑定或生成注册表提供；manifest
  未记录独立头文件行号。
\end{itemize}

\textbf{说明}

\begin{itemize}
\tightlist
\item
  示例是局部调用片段；\passthrough{\lstinline!obj!}
  和参数变量表示已经按上层业务校验并准备好的对象。
\end{itemize}

\textbf{用法}

\begin{lstlisting}[language=Python]
same = left == right
\end{lstlisting}

\hypertarget{unitree-sdk2-cpp-idl-go2-sportmodestate-dunder-ne-1}{%
\paragraph{\texorpdfstring{\texttt{unitree\_sdk2\_cpp.idl.go2.SportModeState.\_\_ne\_\_}}{unitree\_sdk2\_cpp.idl.go2.SportModeState.\_\_ne\_\_}}\label{unitree-sdk2-cpp-idl-go2-sportmodestate-dunder-ne-1}}

按底层 IDL 消息内容比较两个对象是否不相等。

\textbf{签名}

\begin{lstlisting}[language=Python]
def __ne__(self, other: object) -> bool
\end{lstlisting}

\textbf{可用性与安全性}

\textbf{\passthrough{\lstinline!AVAILABLE!} /
\passthrough{\lstinline!VALUE\_TYPE!}}：当前绑定源码已实现该消息值操作。

\textbf{参数}

\begin{longtable}[]{@{}
  >{\raggedright\arraybackslash}p{(\columnwidth - 6\tabcolsep) * \real{0.18}}
  >{\raggedright\arraybackslash}p{(\columnwidth - 6\tabcolsep) * \real{0.24}}
  >{\raggedright\arraybackslash}p{(\columnwidth - 6\tabcolsep) * \real{0.18}}
  >{\raggedright\arraybackslash}p{(\columnwidth - 6\tabcolsep) * \real{0.40}}@{}}
\toprule
名称 & Python 类型 & 默认值 & 含义 \\
\midrule
\endhead
\passthrough{\lstinline!other!} & \passthrough{\lstinline!object!} &
必填 & 用于比较的另一个 Python 对象。类型不兼容时比较结果通常为
\passthrough{\lstinline!False!}。 \\
\bottomrule
\end{longtable}

\textbf{返回值}

\begin{longtable}[]{@{}
  >{\raggedright\arraybackslash}p{(\columnwidth - 4\tabcolsep) * \real{0.18}}
  >{\raggedright\arraybackslash}p{(\columnwidth - 4\tabcolsep) * \real{0.20}}
  >{\raggedright\arraybackslash}p{(\columnwidth - 4\tabcolsep) * \real{0.62}}@{}}
\toprule
位置 & 类型 & 含义 \\
\midrule
\endhead
返回值 & \passthrough{\lstinline!bool!} & 返回
\passthrough{\lstinline!bool!}。更精确的业务含义见该方法用途、C++
签名和目标型号协议。 \\
\bottomrule
\end{longtable}

\textbf{对应 C++}

\begin{itemize}
\tightlist
\item
  类：\passthrough{\lstinline!unitree\_go::msg::dds\_::SportModeState\_!}
\item
  签名：\passthrough{\lstinline"operator!=(const value \&) const"}
\item
  绑定策略：\passthrough{\lstinline!未单独标注!}
\item
  声明位置：由 pybind11 手工绑定或生成注册表提供；manifest
  未记录独立头文件行号。
\end{itemize}

\textbf{说明}

\begin{itemize}
\tightlist
\item
  示例是局部调用片段；\passthrough{\lstinline!obj!}
  和参数变量表示已经按上层业务校验并准备好的对象。
\end{itemize}

\textbf{用法}

\begin{lstlisting}[language=Python]
different = left != right
\end{lstlisting}

\hypertarget{unitree-sdk2-cpp-idl-go2-sportmodestate-stamp}{%
\paragraph{\texorpdfstring{\texttt{unitree\_sdk2\_cpp.idl.go2.SportModeState.stamp}}{unitree\_sdk2\_cpp.idl.go2.SportModeState.stamp}}\label{unitree-sdk2-cpp-idl-go2-sportmodestate-stamp}}

底层 \passthrough{\lstinline!unitree\_go::msg::dds\_::SportModeState\_!}
的 \passthrough{\lstinline!stamp!} 字段。类型签名只说明 Python
表示；单位、数值范围、数组固定长度和枚举含义需查对应 IDL 头文件。

\textbf{签名}

\begin{lstlisting}[language=Python]
@property
def stamp(self) -> TimeSpec

@stamp.setter
def stamp(self, value: TimeSpec) -> None
\end{lstlisting}

\textbf{可用性与安全性}

\textbf{\passthrough{\lstinline!AVAILABLE!} /
\passthrough{\lstinline!VALUE\_TYPE!}}：当前绑定源码已实现该消息值操作。

\textbf{参数}

\begin{longtable}[]{@{}
  >{\raggedright\arraybackslash}p{(\columnwidth - 4\tabcolsep) * \real{0.18}}
  >{\raggedright\arraybackslash}p{(\columnwidth - 4\tabcolsep) * \real{0.20}}
  >{\raggedright\arraybackslash}p{(\columnwidth - 4\tabcolsep) * \real{0.62}}@{}}
\toprule
名称 & Python 类型 & 含义 \\
\midrule
\endhead
\passthrough{\lstinline!value!} & \passthrough{\lstinline!TimeSpec!} &
要写入或传递的新值，必须符合该参数的 Python 类型以及底层 C++
范围约束。 \\
\bottomrule
\end{longtable}

\textbf{返回值}

\begin{longtable}[]{@{}
  >{\raggedright\arraybackslash}p{(\columnwidth - 4\tabcolsep) * \real{0.18}}
  >{\raggedright\arraybackslash}p{(\columnwidth - 4\tabcolsep) * \real{0.20}}
  >{\raggedright\arraybackslash}p{(\columnwidth - 4\tabcolsep) * \real{0.62}}@{}}
\toprule
位置 & 类型 & 含义 \\
\midrule
\endhead
getter 返回值 & \passthrough{\lstinline!TimeSpec!} &
返回属性当前值。数组、vector 和嵌套消息采用复制语义。 \\
\bottomrule
\end{longtable}

\textbf{对应 C++}

\begin{itemize}
\tightlist
\item
  类：\passthrough{\lstinline!unitree\_go::msg::dds\_::SportModeState\_!}
\item
  字段类型：\passthrough{\lstinline!::unitree\_go::msg::dds\_::TimeSpec\_!}
\item
  声明位置：\passthrough{\lstinline!include/unitree/idl/go2/SportModeState\_.hpp:30!}
\end{itemize}

\textbf{用法}

\begin{lstlisting}[language=Python]
current_value = obj.stamp
obj.stamp = new_value
\end{lstlisting}

\hypertarget{unitree-sdk2-cpp-idl-go2-sportmodestate-error-code}{%
\paragraph{\texorpdfstring{\texttt{unitree\_sdk2\_cpp.idl.go2.SportModeState.error\_code}}{unitree\_sdk2\_cpp.idl.go2.SportModeState.error\_code}}\label{unitree-sdk2-cpp-idl-go2-sportmodestate-error-code}}

底层 \passthrough{\lstinline!unitree\_go::msg::dds\_::SportModeState\_!}
的 \passthrough{\lstinline!error\_code!} 字段。类型签名只说明 Python
表示；单位、数值范围、数组固定长度和枚举含义需查对应 IDL 头文件。
取值范围为 0 到 4294967295。

\textbf{签名}

\begin{lstlisting}[language=Python]
@property
def error_code(self) -> int

@error_code.setter
def error_code(self, value: int) -> None
\end{lstlisting}

\textbf{可用性与安全性}

\textbf{\passthrough{\lstinline!AVAILABLE!} /
\passthrough{\lstinline!VALUE\_TYPE!}}：当前绑定源码已实现该消息值操作。

\textbf{参数}

\begin{longtable}[]{@{}
  >{\raggedright\arraybackslash}p{(\columnwidth - 4\tabcolsep) * \real{0.18}}
  >{\raggedright\arraybackslash}p{(\columnwidth - 4\tabcolsep) * \real{0.20}}
  >{\raggedright\arraybackslash}p{(\columnwidth - 4\tabcolsep) * \real{0.62}}@{}}
\toprule
名称 & Python 类型 & 含义 \\
\midrule
\endhead
\passthrough{\lstinline!value!} & \passthrough{\lstinline!int!} &
要写入或传递的新值，必须符合该参数的 Python 类型以及底层 C++
范围约束。 \\
\bottomrule
\end{longtable}

\textbf{返回值}

\begin{longtable}[]{@{}
  >{\raggedright\arraybackslash}p{(\columnwidth - 4\tabcolsep) * \real{0.18}}
  >{\raggedright\arraybackslash}p{(\columnwidth - 4\tabcolsep) * \real{0.20}}
  >{\raggedright\arraybackslash}p{(\columnwidth - 4\tabcolsep) * \real{0.62}}@{}}
\toprule
位置 & 类型 & 含义 \\
\midrule
\endhead
getter 返回值 & \passthrough{\lstinline!int!} & 返回属性当前值。 \\
\bottomrule
\end{longtable}

\textbf{对应 C++}

\begin{itemize}
\tightlist
\item
  类：\passthrough{\lstinline!unitree\_go::msg::dds\_::SportModeState\_!}
\item
  字段类型：\passthrough{\lstinline!uint32\_t!}
\item
  声明位置：\passthrough{\lstinline!include/unitree/idl/go2/SportModeState\_.hpp:31!}
\end{itemize}

\textbf{用法}

\begin{lstlisting}[language=Python]
current_value = obj.error_code
obj.error_code = new_value
\end{lstlisting}

\hypertarget{unitree-sdk2-cpp-idl-go2-sportmodestate-imu-state}{%
\paragraph{\texorpdfstring{\texttt{unitree\_sdk2\_cpp.idl.go2.SportModeState.imu\_state}}{unitree\_sdk2\_cpp.idl.go2.SportModeState.imu\_state}}\label{unitree-sdk2-cpp-idl-go2-sportmodestate-imu-state}}

底层 \passthrough{\lstinline!unitree\_go::msg::dds\_::SportModeState\_!}
的 \passthrough{\lstinline!imu\_state!} 字段。类型签名只说明 Python
表示；单位、数值范围、数组固定长度和枚举含义需查对应 IDL 头文件。

\textbf{签名}

\begin{lstlisting}[language=Python]
@property
def imu_state(self) -> IMUState

@imu_state.setter
def imu_state(self, value: IMUState) -> None
\end{lstlisting}

\textbf{可用性与安全性}

\textbf{\passthrough{\lstinline!AVAILABLE!} /
\passthrough{\lstinline!VALUE\_TYPE!}}：当前绑定源码已实现该消息值操作。

\textbf{参数}

\begin{longtable}[]{@{}
  >{\raggedright\arraybackslash}p{(\columnwidth - 4\tabcolsep) * \real{0.18}}
  >{\raggedright\arraybackslash}p{(\columnwidth - 4\tabcolsep) * \real{0.20}}
  >{\raggedright\arraybackslash}p{(\columnwidth - 4\tabcolsep) * \real{0.62}}@{}}
\toprule
名称 & Python 类型 & 含义 \\
\midrule
\endhead
\passthrough{\lstinline!value!} & \passthrough{\lstinline!IMUState!} &
要写入或传递的新值，必须符合该参数的 Python 类型以及底层 C++
范围约束。 \\
\bottomrule
\end{longtable}

\textbf{返回值}

\begin{longtable}[]{@{}
  >{\raggedright\arraybackslash}p{(\columnwidth - 4\tabcolsep) * \real{0.18}}
  >{\raggedright\arraybackslash}p{(\columnwidth - 4\tabcolsep) * \real{0.20}}
  >{\raggedright\arraybackslash}p{(\columnwidth - 4\tabcolsep) * \real{0.62}}@{}}
\toprule
位置 & 类型 & 含义 \\
\midrule
\endhead
getter 返回值 & \passthrough{\lstinline!IMUState!} &
返回属性当前值。数组、vector 和嵌套消息采用复制语义。 \\
\bottomrule
\end{longtable}

\textbf{对应 C++}

\begin{itemize}
\tightlist
\item
  类：\passthrough{\lstinline!unitree\_go::msg::dds\_::SportModeState\_!}
\item
  字段类型：\passthrough{\lstinline!::unitree\_go::msg::dds\_::IMUState\_!}
\item
  声明位置：\passthrough{\lstinline!include/unitree/idl/go2/SportModeState\_.hpp:32!}
\end{itemize}

\textbf{用法}

\begin{lstlisting}[language=Python]
current_value = obj.imu_state
obj.imu_state = new_value
\end{lstlisting}

\hypertarget{unitree-sdk2-cpp-idl-go2-sportmodestate-mode}{%
\paragraph{\texorpdfstring{\texttt{unitree\_sdk2\_cpp.idl.go2.SportModeState.mode}}{unitree\_sdk2\_cpp.idl.go2.SportModeState.mode}}\label{unitree-sdk2-cpp-idl-go2-sportmodestate-mode}}

底层 \passthrough{\lstinline!unitree\_go::msg::dds\_::SportModeState\_!}
的 \passthrough{\lstinline!mode!} 字段。类型签名只说明 Python
表示；单位、数值范围、数组固定长度和枚举含义需查对应 IDL 头文件。
取值范围为 0 到 255。

\textbf{签名}

\begin{lstlisting}[language=Python]
@property
def mode(self) -> int

@mode.setter
def mode(self, value: int) -> None
\end{lstlisting}

\textbf{可用性与安全性}

\textbf{\passthrough{\lstinline!AVAILABLE!} /
\passthrough{\lstinline!VALUE\_TYPE!}}：当前绑定源码已实现该消息值操作。

\textbf{参数}

\begin{longtable}[]{@{}
  >{\raggedright\arraybackslash}p{(\columnwidth - 4\tabcolsep) * \real{0.18}}
  >{\raggedright\arraybackslash}p{(\columnwidth - 4\tabcolsep) * \real{0.20}}
  >{\raggedright\arraybackslash}p{(\columnwidth - 4\tabcolsep) * \real{0.62}}@{}}
\toprule
名称 & Python 类型 & 含义 \\
\midrule
\endhead
\passthrough{\lstinline!value!} & \passthrough{\lstinline!int!} &
要写入或传递的新值，必须符合该参数的 Python 类型以及底层 C++
范围约束。 \\
\bottomrule
\end{longtable}

\textbf{返回值}

\begin{longtable}[]{@{}
  >{\raggedright\arraybackslash}p{(\columnwidth - 4\tabcolsep) * \real{0.18}}
  >{\raggedright\arraybackslash}p{(\columnwidth - 4\tabcolsep) * \real{0.20}}
  >{\raggedright\arraybackslash}p{(\columnwidth - 4\tabcolsep) * \real{0.62}}@{}}
\toprule
位置 & 类型 & 含义 \\
\midrule
\endhead
getter 返回值 & \passthrough{\lstinline!int!} & 返回属性当前值。 \\
\bottomrule
\end{longtable}

\textbf{对应 C++}

\begin{itemize}
\tightlist
\item
  类：\passthrough{\lstinline!unitree\_go::msg::dds\_::SportModeState\_!}
\item
  字段类型：\passthrough{\lstinline!uint8\_t!}
\item
  声明位置：\passthrough{\lstinline!include/unitree/idl/go2/SportModeState\_.hpp:33!}
\end{itemize}

\textbf{用法}

\begin{lstlisting}[language=Python]
current_value = obj.mode
obj.mode = new_value
\end{lstlisting}

\hypertarget{unitree-sdk2-cpp-idl-go2-sportmodestate-progress}{%
\paragraph{\texorpdfstring{\texttt{unitree\_sdk2\_cpp.idl.go2.SportModeState.progress}}{unitree\_sdk2\_cpp.idl.go2.SportModeState.progress}}\label{unitree-sdk2-cpp-idl-go2-sportmodestate-progress}}

底层 \passthrough{\lstinline!unitree\_go::msg::dds\_::SportModeState\_!}
的 \passthrough{\lstinline!progress!} 字段。类型签名只说明 Python
表示；单位、数值范围、数组固定长度和枚举含义需查对应 IDL 头文件。
底层为浮点数；类型本身不说明单位、坐标系或业务安全范围。

\textbf{签名}

\begin{lstlisting}[language=Python]
@property
def progress(self) -> float

@progress.setter
def progress(self, value: float) -> None
\end{lstlisting}

\textbf{可用性与安全性}

\textbf{\passthrough{\lstinline!AVAILABLE!} /
\passthrough{\lstinline!VALUE\_TYPE!}}：当前绑定源码已实现该消息值操作。

\textbf{参数}

\begin{longtable}[]{@{}
  >{\raggedright\arraybackslash}p{(\columnwidth - 4\tabcolsep) * \real{0.18}}
  >{\raggedright\arraybackslash}p{(\columnwidth - 4\tabcolsep) * \real{0.20}}
  >{\raggedright\arraybackslash}p{(\columnwidth - 4\tabcolsep) * \real{0.62}}@{}}
\toprule
名称 & Python 类型 & 含义 \\
\midrule
\endhead
\passthrough{\lstinline!value!} & \passthrough{\lstinline!float!} &
要写入或传递的新值，必须符合该参数的 Python 类型以及底层 C++
范围约束。 \\
\bottomrule
\end{longtable}

\textbf{返回值}

\begin{longtable}[]{@{}
  >{\raggedright\arraybackslash}p{(\columnwidth - 4\tabcolsep) * \real{0.18}}
  >{\raggedright\arraybackslash}p{(\columnwidth - 4\tabcolsep) * \real{0.20}}
  >{\raggedright\arraybackslash}p{(\columnwidth - 4\tabcolsep) * \real{0.62}}@{}}
\toprule
位置 & 类型 & 含义 \\
\midrule
\endhead
getter 返回值 & \passthrough{\lstinline!float!} & 返回属性当前值。 \\
\bottomrule
\end{longtable}

\textbf{对应 C++}

\begin{itemize}
\tightlist
\item
  类：\passthrough{\lstinline!unitree\_go::msg::dds\_::SportModeState\_!}
\item
  字段类型：\passthrough{\lstinline!float!}
\item
  声明位置：\passthrough{\lstinline!include/unitree/idl/go2/SportModeState\_.hpp:34!}
\end{itemize}

\textbf{用法}

\begin{lstlisting}[language=Python]
current_value = obj.progress
obj.progress = new_value
\end{lstlisting}

\hypertarget{unitree-sdk2-cpp-idl-go2-sportmodestate-gait-type}{%
\paragraph{\texorpdfstring{\texttt{unitree\_sdk2\_cpp.idl.go2.SportModeState.gait\_type}}{unitree\_sdk2\_cpp.idl.go2.SportModeState.gait\_type}}\label{unitree-sdk2-cpp-idl-go2-sportmodestate-gait-type}}

底层 \passthrough{\lstinline!unitree\_go::msg::dds\_::SportModeState\_!}
的 \passthrough{\lstinline!gait\_type!} 字段。类型签名只说明 Python
表示；单位、数值范围、数组固定长度和枚举含义需查对应 IDL 头文件。
取值范围为 0 到 255。

\textbf{签名}

\begin{lstlisting}[language=Python]
@property
def gait_type(self) -> int

@gait_type.setter
def gait_type(self, value: int) -> None
\end{lstlisting}

\textbf{可用性与安全性}

\textbf{\passthrough{\lstinline!AVAILABLE!} /
\passthrough{\lstinline!VALUE\_TYPE!}}：当前绑定源码已实现该消息值操作。

\textbf{参数}

\begin{longtable}[]{@{}
  >{\raggedright\arraybackslash}p{(\columnwidth - 4\tabcolsep) * \real{0.18}}
  >{\raggedright\arraybackslash}p{(\columnwidth - 4\tabcolsep) * \real{0.20}}
  >{\raggedright\arraybackslash}p{(\columnwidth - 4\tabcolsep) * \real{0.62}}@{}}
\toprule
名称 & Python 类型 & 含义 \\
\midrule
\endhead
\passthrough{\lstinline!value!} & \passthrough{\lstinline!int!} &
要写入或传递的新值，必须符合该参数的 Python 类型以及底层 C++
范围约束。 \\
\bottomrule
\end{longtable}

\textbf{返回值}

\begin{longtable}[]{@{}
  >{\raggedright\arraybackslash}p{(\columnwidth - 4\tabcolsep) * \real{0.18}}
  >{\raggedright\arraybackslash}p{(\columnwidth - 4\tabcolsep) * \real{0.20}}
  >{\raggedright\arraybackslash}p{(\columnwidth - 4\tabcolsep) * \real{0.62}}@{}}
\toprule
位置 & 类型 & 含义 \\
\midrule
\endhead
getter 返回值 & \passthrough{\lstinline!int!} & 返回属性当前值。 \\
\bottomrule
\end{longtable}

\textbf{对应 C++}

\begin{itemize}
\tightlist
\item
  类：\passthrough{\lstinline!unitree\_go::msg::dds\_::SportModeState\_!}
\item
  字段类型：\passthrough{\lstinline!uint8\_t!}
\item
  声明位置：\passthrough{\lstinline!include/unitree/idl/go2/SportModeState\_.hpp:35!}
\end{itemize}

\textbf{用法}

\begin{lstlisting}[language=Python]
current_value = obj.gait_type
obj.gait_type = new_value
\end{lstlisting}

\hypertarget{unitree-sdk2-cpp-idl-go2-sportmodestate-foot-raise-height}{%
\paragraph{\texorpdfstring{\texttt{unitree\_sdk2\_cpp.idl.go2.SportModeState.foot\_raise\_height}}{unitree\_sdk2\_cpp.idl.go2.SportModeState.foot\_raise\_height}}\label{unitree-sdk2-cpp-idl-go2-sportmodestate-foot-raise-height}}

底层 \passthrough{\lstinline!unitree\_go::msg::dds\_::SportModeState\_!}
的 \passthrough{\lstinline!foot\_raise\_height!} 字段。类型签名只说明
Python 表示；单位、数值范围、数组固定长度和枚举含义需查对应 IDL 头文件。
底层为浮点数；类型本身不说明单位、坐标系或业务安全范围。

\textbf{签名}

\begin{lstlisting}[language=Python]
@property
def foot_raise_height(self) -> float

@foot_raise_height.setter
def foot_raise_height(self, value: float) -> None
\end{lstlisting}

\textbf{可用性与安全性}

\textbf{\passthrough{\lstinline!AVAILABLE!} /
\passthrough{\lstinline!VALUE\_TYPE!}}：当前绑定源码已实现该消息值操作。

\textbf{参数}

\begin{longtable}[]{@{}
  >{\raggedright\arraybackslash}p{(\columnwidth - 4\tabcolsep) * \real{0.18}}
  >{\raggedright\arraybackslash}p{(\columnwidth - 4\tabcolsep) * \real{0.20}}
  >{\raggedright\arraybackslash}p{(\columnwidth - 4\tabcolsep) * \real{0.62}}@{}}
\toprule
名称 & Python 类型 & 含义 \\
\midrule
\endhead
\passthrough{\lstinline!value!} & \passthrough{\lstinline!float!} &
要写入或传递的新值，必须符合该参数的 Python 类型以及底层 C++
范围约束。 \\
\bottomrule
\end{longtable}

\textbf{返回值}

\begin{longtable}[]{@{}
  >{\raggedright\arraybackslash}p{(\columnwidth - 4\tabcolsep) * \real{0.18}}
  >{\raggedright\arraybackslash}p{(\columnwidth - 4\tabcolsep) * \real{0.20}}
  >{\raggedright\arraybackslash}p{(\columnwidth - 4\tabcolsep) * \real{0.62}}@{}}
\toprule
位置 & 类型 & 含义 \\
\midrule
\endhead
getter 返回值 & \passthrough{\lstinline!float!} & 返回属性当前值。 \\
\bottomrule
\end{longtable}

\textbf{对应 C++}

\begin{itemize}
\tightlist
\item
  类：\passthrough{\lstinline!unitree\_go::msg::dds\_::SportModeState\_!}
\item
  字段类型：\passthrough{\lstinline!float!}
\item
  声明位置：\passthrough{\lstinline!include/unitree/idl/go2/SportModeState\_.hpp:36!}
\end{itemize}

\textbf{用法}

\begin{lstlisting}[language=Python]
current_value = obj.foot_raise_height
obj.foot_raise_height = new_value
\end{lstlisting}

\hypertarget{unitree-sdk2-cpp-idl-go2-sportmodestate-position}{%
\paragraph{\texorpdfstring{\texttt{unitree\_sdk2\_cpp.idl.go2.SportModeState.position}}{unitree\_sdk2\_cpp.idl.go2.SportModeState.position}}\label{unitree-sdk2-cpp-idl-go2-sportmodestate-position}}

底层 \passthrough{\lstinline!unitree\_go::msg::dds\_::SportModeState\_!}
的 \passthrough{\lstinline!position!} 字段。类型签名只说明 Python
表示；单位、数值范围、数组固定长度和枚举含义需查对应 IDL 头文件。
底层是固定长度数组，写入序列必须正好包含 3
个元素；元素约束：底层为浮点数；类型本身不说明单位、坐标系或业务安全范围。

\textbf{签名}

\begin{lstlisting}[language=Python]
@property
def position(self) -> list[float]

@position.setter
def position(self, value: Sequence[float]) -> None
\end{lstlisting}

\textbf{可用性与安全性}

\textbf{\passthrough{\lstinline!AVAILABLE!} /
\passthrough{\lstinline!VALUE\_TYPE!}}：当前绑定源码已实现该消息值操作。

\textbf{参数}

\begin{longtable}[]{@{}
  >{\raggedright\arraybackslash}p{(\columnwidth - 4\tabcolsep) * \real{0.18}}
  >{\raggedright\arraybackslash}p{(\columnwidth - 4\tabcolsep) * \real{0.20}}
  >{\raggedright\arraybackslash}p{(\columnwidth - 4\tabcolsep) * \real{0.62}}@{}}
\toprule
名称 & Python 类型 & 含义 \\
\midrule
\endhead
\passthrough{\lstinline!value!} &
\passthrough{\lstinline!Sequence[float]!} &
要写入或传递的新值，必须符合该参数的 Python 类型以及底层 C++
范围约束。 \\
\bottomrule
\end{longtable}

\textbf{返回值}

\begin{longtable}[]{@{}
  >{\raggedright\arraybackslash}p{(\columnwidth - 4\tabcolsep) * \real{0.18}}
  >{\raggedright\arraybackslash}p{(\columnwidth - 4\tabcolsep) * \real{0.20}}
  >{\raggedright\arraybackslash}p{(\columnwidth - 4\tabcolsep) * \real{0.62}}@{}}
\toprule
位置 & 类型 & 含义 \\
\midrule
\endhead
getter 返回值 & \passthrough{\lstinline!list[float]!} &
返回属性当前值。数组、vector 和嵌套消息采用复制语义。 \\
\bottomrule
\end{longtable}

\textbf{对应 C++}

\begin{itemize}
\tightlist
\item
  类：\passthrough{\lstinline!unitree\_go::msg::dds\_::SportModeState\_!}
\item
  字段类型：\passthrough{\lstinline!std::array<float, 3>!}
\item
  声明位置：\passthrough{\lstinline!include/unitree/idl/go2/SportModeState\_.hpp:37!}
\end{itemize}

\textbf{用法}

\begin{lstlisting}[language=Python]
current_value = obj.position
obj.position = new_value
\end{lstlisting}

\begin{quote}
{[}!TIP{]} 读取容器或嵌套值后应遵循``读取、修改、重新赋值''。只修改
getter 返回的副本不会可靠地更新原 C++ 消息。
\end{quote}

\hypertarget{unitree-sdk2-cpp-idl-go2-sportmodestate-body-height}{%
\paragraph{\texorpdfstring{\texttt{unitree\_sdk2\_cpp.idl.go2.SportModeState.body\_height}}{unitree\_sdk2\_cpp.idl.go2.SportModeState.body\_height}}\label{unitree-sdk2-cpp-idl-go2-sportmodestate-body-height}}

底层 \passthrough{\lstinline!unitree\_go::msg::dds\_::SportModeState\_!}
的 \passthrough{\lstinline!body\_height!} 字段。类型签名只说明 Python
表示；单位、数值范围、数组固定长度和枚举含义需查对应 IDL 头文件。
底层为浮点数；类型本身不说明单位、坐标系或业务安全范围。

\textbf{签名}

\begin{lstlisting}[language=Python]
@property
def body_height(self) -> float

@body_height.setter
def body_height(self, value: float) -> None
\end{lstlisting}

\textbf{可用性与安全性}

\textbf{\passthrough{\lstinline!AVAILABLE!} /
\passthrough{\lstinline!VALUE\_TYPE!}}：当前绑定源码已实现该消息值操作。

\textbf{参数}

\begin{longtable}[]{@{}
  >{\raggedright\arraybackslash}p{(\columnwidth - 4\tabcolsep) * \real{0.18}}
  >{\raggedright\arraybackslash}p{(\columnwidth - 4\tabcolsep) * \real{0.20}}
  >{\raggedright\arraybackslash}p{(\columnwidth - 4\tabcolsep) * \real{0.62}}@{}}
\toprule
名称 & Python 类型 & 含义 \\
\midrule
\endhead
\passthrough{\lstinline!value!} & \passthrough{\lstinline!float!} &
要写入或传递的新值，必须符合该参数的 Python 类型以及底层 C++
范围约束。 \\
\bottomrule
\end{longtable}

\textbf{返回值}

\begin{longtable}[]{@{}
  >{\raggedright\arraybackslash}p{(\columnwidth - 4\tabcolsep) * \real{0.18}}
  >{\raggedright\arraybackslash}p{(\columnwidth - 4\tabcolsep) * \real{0.20}}
  >{\raggedright\arraybackslash}p{(\columnwidth - 4\tabcolsep) * \real{0.62}}@{}}
\toprule
位置 & 类型 & 含义 \\
\midrule
\endhead
getter 返回值 & \passthrough{\lstinline!float!} & 返回属性当前值。 \\
\bottomrule
\end{longtable}

\textbf{对应 C++}

\begin{itemize}
\tightlist
\item
  类：\passthrough{\lstinline!unitree\_go::msg::dds\_::SportModeState\_!}
\item
  字段类型：\passthrough{\lstinline!float!}
\item
  声明位置：\passthrough{\lstinline!include/unitree/idl/go2/SportModeState\_.hpp:38!}
\end{itemize}

\textbf{用法}

\begin{lstlisting}[language=Python]
current_value = obj.body_height
obj.body_height = new_value
\end{lstlisting}

\hypertarget{unitree-sdk2-cpp-idl-go2-sportmodestate-velocity}{%
\paragraph{\texorpdfstring{\texttt{unitree\_sdk2\_cpp.idl.go2.SportModeState.velocity}}{unitree\_sdk2\_cpp.idl.go2.SportModeState.velocity}}\label{unitree-sdk2-cpp-idl-go2-sportmodestate-velocity}}

底层 \passthrough{\lstinline!unitree\_go::msg::dds\_::SportModeState\_!}
的 \passthrough{\lstinline!velocity!} 字段。类型签名只说明 Python
表示；单位、数值范围、数组固定长度和枚举含义需查对应 IDL 头文件。
底层是固定长度数组，写入序列必须正好包含 3
个元素；元素约束：底层为浮点数；类型本身不说明单位、坐标系或业务安全范围。

\textbf{签名}

\begin{lstlisting}[language=Python]
@property
def velocity(self) -> list[float]

@velocity.setter
def velocity(self, value: Sequence[float]) -> None
\end{lstlisting}

\textbf{可用性与安全性}

\textbf{\passthrough{\lstinline!AVAILABLE!} /
\passthrough{\lstinline!VALUE\_TYPE!}}：当前绑定源码已实现该消息值操作。

\textbf{参数}

\begin{longtable}[]{@{}
  >{\raggedright\arraybackslash}p{(\columnwidth - 4\tabcolsep) * \real{0.18}}
  >{\raggedright\arraybackslash}p{(\columnwidth - 4\tabcolsep) * \real{0.20}}
  >{\raggedright\arraybackslash}p{(\columnwidth - 4\tabcolsep) * \real{0.62}}@{}}
\toprule
名称 & Python 类型 & 含义 \\
\midrule
\endhead
\passthrough{\lstinline!value!} &
\passthrough{\lstinline!Sequence[float]!} &
要写入或传递的新值，必须符合该参数的 Python 类型以及底层 C++
范围约束。 \\
\bottomrule
\end{longtable}

\textbf{返回值}

\begin{longtable}[]{@{}
  >{\raggedright\arraybackslash}p{(\columnwidth - 4\tabcolsep) * \real{0.18}}
  >{\raggedright\arraybackslash}p{(\columnwidth - 4\tabcolsep) * \real{0.20}}
  >{\raggedright\arraybackslash}p{(\columnwidth - 4\tabcolsep) * \real{0.62}}@{}}
\toprule
位置 & 类型 & 含义 \\
\midrule
\endhead
getter 返回值 & \passthrough{\lstinline!list[float]!} &
返回属性当前值。数组、vector 和嵌套消息采用复制语义。 \\
\bottomrule
\end{longtable}

\textbf{对应 C++}

\begin{itemize}
\tightlist
\item
  类：\passthrough{\lstinline!unitree\_go::msg::dds\_::SportModeState\_!}
\item
  字段类型：\passthrough{\lstinline!std::array<float, 3>!}
\item
  声明位置：\passthrough{\lstinline!include/unitree/idl/go2/SportModeState\_.hpp:39!}
\end{itemize}

\textbf{用法}

\begin{lstlisting}[language=Python]
current_value = obj.velocity
obj.velocity = new_value
\end{lstlisting}

\begin{quote}
{[}!TIP{]} 读取容器或嵌套值后应遵循``读取、修改、重新赋值''。只修改
getter 返回的副本不会可靠地更新原 C++ 消息。
\end{quote}

\hypertarget{unitree-sdk2-cpp-idl-go2-sportmodestate-yaw-speed}{%
\paragraph{\texorpdfstring{\texttt{unitree\_sdk2\_cpp.idl.go2.SportModeState.yaw\_speed}}{unitree\_sdk2\_cpp.idl.go2.SportModeState.yaw\_speed}}\label{unitree-sdk2-cpp-idl-go2-sportmodestate-yaw-speed}}

底层 \passthrough{\lstinline!unitree\_go::msg::dds\_::SportModeState\_!}
的 \passthrough{\lstinline!yaw\_speed!} 字段。类型签名只说明 Python
表示；单位、数值范围、数组固定长度和枚举含义需查对应 IDL 头文件。
底层为浮点数；类型本身不说明单位、坐标系或业务安全范围。

\textbf{签名}

\begin{lstlisting}[language=Python]
@property
def yaw_speed(self) -> float

@yaw_speed.setter
def yaw_speed(self, value: float) -> None
\end{lstlisting}

\textbf{可用性与安全性}

\textbf{\passthrough{\lstinline!AVAILABLE!} /
\passthrough{\lstinline!VALUE\_TYPE!}}：当前绑定源码已实现该消息值操作。

\textbf{参数}

\begin{longtable}[]{@{}
  >{\raggedright\arraybackslash}p{(\columnwidth - 4\tabcolsep) * \real{0.18}}
  >{\raggedright\arraybackslash}p{(\columnwidth - 4\tabcolsep) * \real{0.20}}
  >{\raggedright\arraybackslash}p{(\columnwidth - 4\tabcolsep) * \real{0.62}}@{}}
\toprule
名称 & Python 类型 & 含义 \\
\midrule
\endhead
\passthrough{\lstinline!value!} & \passthrough{\lstinline!float!} &
要写入或传递的新值，必须符合该参数的 Python 类型以及底层 C++
范围约束。 \\
\bottomrule
\end{longtable}

\textbf{返回值}

\begin{longtable}[]{@{}
  >{\raggedright\arraybackslash}p{(\columnwidth - 4\tabcolsep) * \real{0.18}}
  >{\raggedright\arraybackslash}p{(\columnwidth - 4\tabcolsep) * \real{0.20}}
  >{\raggedright\arraybackslash}p{(\columnwidth - 4\tabcolsep) * \real{0.62}}@{}}
\toprule
位置 & 类型 & 含义 \\
\midrule
\endhead
getter 返回值 & \passthrough{\lstinline!float!} & 返回属性当前值。 \\
\bottomrule
\end{longtable}

\textbf{对应 C++}

\begin{itemize}
\tightlist
\item
  类：\passthrough{\lstinline!unitree\_go::msg::dds\_::SportModeState\_!}
\item
  字段类型：\passthrough{\lstinline!float!}
\item
  声明位置：\passthrough{\lstinline!include/unitree/idl/go2/SportModeState\_.hpp:40!}
\end{itemize}

\textbf{用法}

\begin{lstlisting}[language=Python]
current_value = obj.yaw_speed
obj.yaw_speed = new_value
\end{lstlisting}

\hypertarget{unitree-sdk2-cpp-idl-go2-sportmodestate-range-obstacle}{%
\paragraph{\texorpdfstring{\texttt{unitree\_sdk2\_cpp.idl.go2.SportModeState.range\_obstacle}}{unitree\_sdk2\_cpp.idl.go2.SportModeState.range\_obstacle}}\label{unitree-sdk2-cpp-idl-go2-sportmodestate-range-obstacle}}

底层 \passthrough{\lstinline!unitree\_go::msg::dds\_::SportModeState\_!}
的 \passthrough{\lstinline!range\_obstacle!} 字段。类型签名只说明 Python
表示；单位、数值范围、数组固定长度和枚举含义需查对应 IDL 头文件。
底层是固定长度数组，写入序列必须正好包含 4
个元素；元素约束：底层为浮点数；类型本身不说明单位、坐标系或业务安全范围。

\textbf{签名}

\begin{lstlisting}[language=Python]
@property
def range_obstacle(self) -> list[float]

@range_obstacle.setter
def range_obstacle(self, value: Sequence[float]) -> None
\end{lstlisting}

\textbf{可用性与安全性}

\textbf{\passthrough{\lstinline!AVAILABLE!} /
\passthrough{\lstinline!VALUE\_TYPE!}}：当前绑定源码已实现该消息值操作。

\textbf{参数}

\begin{longtable}[]{@{}
  >{\raggedright\arraybackslash}p{(\columnwidth - 4\tabcolsep) * \real{0.18}}
  >{\raggedright\arraybackslash}p{(\columnwidth - 4\tabcolsep) * \real{0.20}}
  >{\raggedright\arraybackslash}p{(\columnwidth - 4\tabcolsep) * \real{0.62}}@{}}
\toprule
名称 & Python 类型 & 含义 \\
\midrule
\endhead
\passthrough{\lstinline!value!} &
\passthrough{\lstinline!Sequence[float]!} &
要写入或传递的新值，必须符合该参数的 Python 类型以及底层 C++
范围约束。 \\
\bottomrule
\end{longtable}

\textbf{返回值}

\begin{longtable}[]{@{}
  >{\raggedright\arraybackslash}p{(\columnwidth - 4\tabcolsep) * \real{0.18}}
  >{\raggedright\arraybackslash}p{(\columnwidth - 4\tabcolsep) * \real{0.20}}
  >{\raggedright\arraybackslash}p{(\columnwidth - 4\tabcolsep) * \real{0.62}}@{}}
\toprule
位置 & 类型 & 含义 \\
\midrule
\endhead
getter 返回值 & \passthrough{\lstinline!list[float]!} &
返回属性当前值。数组、vector 和嵌套消息采用复制语义。 \\
\bottomrule
\end{longtable}

\textbf{对应 C++}

\begin{itemize}
\tightlist
\item
  类：\passthrough{\lstinline!unitree\_go::msg::dds\_::SportModeState\_!}
\item
  字段类型：\passthrough{\lstinline!std::array<float, 4>!}
\item
  声明位置：\passthrough{\lstinline!include/unitree/idl/go2/SportModeState\_.hpp:41!}
\end{itemize}

\textbf{用法}

\begin{lstlisting}[language=Python]
current_value = obj.range_obstacle
obj.range_obstacle = new_value
\end{lstlisting}

\begin{quote}
{[}!TIP{]} 读取容器或嵌套值后应遵循``读取、修改、重新赋值''。只修改
getter 返回的副本不会可靠地更新原 C++ 消息。
\end{quote}

\hypertarget{unitree-sdk2-cpp-idl-go2-sportmodestate-foot-force}{%
\paragraph{\texorpdfstring{\texttt{unitree\_sdk2\_cpp.idl.go2.SportModeState.foot\_force}}{unitree\_sdk2\_cpp.idl.go2.SportModeState.foot\_force}}\label{unitree-sdk2-cpp-idl-go2-sportmodestate-foot-force}}

底层 \passthrough{\lstinline!unitree\_go::msg::dds\_::SportModeState\_!}
的 \passthrough{\lstinline!foot\_force!} 字段。类型签名只说明 Python
表示；单位、数值范围、数组固定长度和枚举含义需查对应 IDL 头文件。
底层是固定长度数组，写入序列必须正好包含 4 个元素；元素约束：取值范围为
-32768 到 32767。

\textbf{签名}

\begin{lstlisting}[language=Python]
@property
def foot_force(self) -> list[int]

@foot_force.setter
def foot_force(self, value: Sequence[int]) -> None
\end{lstlisting}

\textbf{可用性与安全性}

\textbf{\passthrough{\lstinline!AVAILABLE!} /
\passthrough{\lstinline!VALUE\_TYPE!}}：当前绑定源码已实现该消息值操作。

\textbf{参数}

\begin{longtable}[]{@{}
  >{\raggedright\arraybackslash}p{(\columnwidth - 4\tabcolsep) * \real{0.18}}
  >{\raggedright\arraybackslash}p{(\columnwidth - 4\tabcolsep) * \real{0.20}}
  >{\raggedright\arraybackslash}p{(\columnwidth - 4\tabcolsep) * \real{0.62}}@{}}
\toprule
名称 & Python 类型 & 含义 \\
\midrule
\endhead
\passthrough{\lstinline!value!} &
\passthrough{\lstinline!Sequence[int]!} &
要写入或传递的新值，必须符合该参数的 Python 类型以及底层 C++
范围约束。 \\
\bottomrule
\end{longtable}

\textbf{返回值}

\begin{longtable}[]{@{}
  >{\raggedright\arraybackslash}p{(\columnwidth - 4\tabcolsep) * \real{0.18}}
  >{\raggedright\arraybackslash}p{(\columnwidth - 4\tabcolsep) * \real{0.20}}
  >{\raggedright\arraybackslash}p{(\columnwidth - 4\tabcolsep) * \real{0.62}}@{}}
\toprule
位置 & 类型 & 含义 \\
\midrule
\endhead
getter 返回值 & \passthrough{\lstinline!list[int]!} &
返回属性当前值。数组、vector 和嵌套消息采用复制语义。 \\
\bottomrule
\end{longtable}

\textbf{对应 C++}

\begin{itemize}
\tightlist
\item
  类：\passthrough{\lstinline!unitree\_go::msg::dds\_::SportModeState\_!}
\item
  字段类型：\passthrough{\lstinline!std::array<int16\_t, 4>!}
\item
  声明位置：\passthrough{\lstinline!include/unitree/idl/go2/SportModeState\_.hpp:42!}
\end{itemize}

\textbf{用法}

\begin{lstlisting}[language=Python]
current_value = obj.foot_force
obj.foot_force = new_value
\end{lstlisting}

\begin{quote}
{[}!TIP{]} 读取容器或嵌套值后应遵循``读取、修改、重新赋值''。只修改
getter 返回的副本不会可靠地更新原 C++ 消息。
\end{quote}

\hypertarget{unitree-sdk2-cpp-idl-go2-sportmodestate-foot-position-body}{%
\paragraph{\texorpdfstring{\texttt{unitree\_sdk2\_cpp.idl.go2.SportModeState.foot\_position\_body}}{unitree\_sdk2\_cpp.idl.go2.SportModeState.foot\_position\_body}}\label{unitree-sdk2-cpp-idl-go2-sportmodestate-foot-position-body}}

底层 \passthrough{\lstinline!unitree\_go::msg::dds\_::SportModeState\_!}
的 \passthrough{\lstinline!foot\_position\_body!} 字段。类型签名只说明
Python 表示；单位、数值范围、数组固定长度和枚举含义需查对应 IDL 头文件。
底层是固定长度数组，写入序列必须正好包含 12
个元素；元素约束：底层为浮点数；类型本身不说明单位、坐标系或业务安全范围。

\textbf{签名}

\begin{lstlisting}[language=Python]
@property
def foot_position_body(self) -> list[float]

@foot_position_body.setter
def foot_position_body(self, value: Sequence[float]) -> None
\end{lstlisting}

\textbf{可用性与安全性}

\textbf{\passthrough{\lstinline!AVAILABLE!} /
\passthrough{\lstinline!VALUE\_TYPE!}}：当前绑定源码已实现该消息值操作。

\textbf{参数}

\begin{longtable}[]{@{}
  >{\raggedright\arraybackslash}p{(\columnwidth - 4\tabcolsep) * \real{0.18}}
  >{\raggedright\arraybackslash}p{(\columnwidth - 4\tabcolsep) * \real{0.20}}
  >{\raggedright\arraybackslash}p{(\columnwidth - 4\tabcolsep) * \real{0.62}}@{}}
\toprule
名称 & Python 类型 & 含义 \\
\midrule
\endhead
\passthrough{\lstinline!value!} &
\passthrough{\lstinline!Sequence[float]!} &
要写入或传递的新值，必须符合该参数的 Python 类型以及底层 C++
范围约束。 \\
\bottomrule
\end{longtable}

\textbf{返回值}

\begin{longtable}[]{@{}
  >{\raggedright\arraybackslash}p{(\columnwidth - 4\tabcolsep) * \real{0.18}}
  >{\raggedright\arraybackslash}p{(\columnwidth - 4\tabcolsep) * \real{0.20}}
  >{\raggedright\arraybackslash}p{(\columnwidth - 4\tabcolsep) * \real{0.62}}@{}}
\toprule
位置 & 类型 & 含义 \\
\midrule
\endhead
getter 返回值 & \passthrough{\lstinline!list[float]!} &
返回属性当前值。数组、vector 和嵌套消息采用复制语义。 \\
\bottomrule
\end{longtable}

\textbf{对应 C++}

\begin{itemize}
\tightlist
\item
  类：\passthrough{\lstinline!unitree\_go::msg::dds\_::SportModeState\_!}
\item
  字段类型：\passthrough{\lstinline!std::array<float, 12>!}
\item
  声明位置：\passthrough{\lstinline!include/unitree/idl/go2/SportModeState\_.hpp:43!}
\end{itemize}

\textbf{用法}

\begin{lstlisting}[language=Python]
current_value = obj.foot_position_body
obj.foot_position_body = new_value
\end{lstlisting}

\begin{quote}
{[}!TIP{]} 读取容器或嵌套值后应遵循``读取、修改、重新赋值''。只修改
getter 返回的副本不会可靠地更新原 C++ 消息。
\end{quote}

\hypertarget{unitree-sdk2-cpp-idl-go2-sportmodestate-foot-speed-body}{%
\paragraph{\texorpdfstring{\texttt{unitree\_sdk2\_cpp.idl.go2.SportModeState.foot\_speed\_body}}{unitree\_sdk2\_cpp.idl.go2.SportModeState.foot\_speed\_body}}\label{unitree-sdk2-cpp-idl-go2-sportmodestate-foot-speed-body}}

底层 \passthrough{\lstinline!unitree\_go::msg::dds\_::SportModeState\_!}
的 \passthrough{\lstinline!foot\_speed\_body!} 字段。类型签名只说明
Python 表示；单位、数值范围、数组固定长度和枚举含义需查对应 IDL 头文件。
底层是固定长度数组，写入序列必须正好包含 12
个元素；元素约束：底层为浮点数；类型本身不说明单位、坐标系或业务安全范围。

\textbf{签名}

\begin{lstlisting}[language=Python]
@property
def foot_speed_body(self) -> list[float]

@foot_speed_body.setter
def foot_speed_body(self, value: Sequence[float]) -> None
\end{lstlisting}

\textbf{可用性与安全性}

\textbf{\passthrough{\lstinline!AVAILABLE!} /
\passthrough{\lstinline!VALUE\_TYPE!}}：当前绑定源码已实现该消息值操作。

\textbf{参数}

\begin{longtable}[]{@{}
  >{\raggedright\arraybackslash}p{(\columnwidth - 4\tabcolsep) * \real{0.18}}
  >{\raggedright\arraybackslash}p{(\columnwidth - 4\tabcolsep) * \real{0.20}}
  >{\raggedright\arraybackslash}p{(\columnwidth - 4\tabcolsep) * \real{0.62}}@{}}
\toprule
名称 & Python 类型 & 含义 \\
\midrule
\endhead
\passthrough{\lstinline!value!} &
\passthrough{\lstinline!Sequence[float]!} &
要写入或传递的新值，必须符合该参数的 Python 类型以及底层 C++
范围约束。 \\
\bottomrule
\end{longtable}

\textbf{返回值}

\begin{longtable}[]{@{}
  >{\raggedright\arraybackslash}p{(\columnwidth - 4\tabcolsep) * \real{0.18}}
  >{\raggedright\arraybackslash}p{(\columnwidth - 4\tabcolsep) * \real{0.20}}
  >{\raggedright\arraybackslash}p{(\columnwidth - 4\tabcolsep) * \real{0.62}}@{}}
\toprule
位置 & 类型 & 含义 \\
\midrule
\endhead
getter 返回值 & \passthrough{\lstinline!list[float]!} &
返回属性当前值。数组、vector 和嵌套消息采用复制语义。 \\
\bottomrule
\end{longtable}

\textbf{对应 C++}

\begin{itemize}
\tightlist
\item
  类：\passthrough{\lstinline!unitree\_go::msg::dds\_::SportModeState\_!}
\item
  字段类型：\passthrough{\lstinline!std::array<float, 12>!}
\item
  声明位置：\passthrough{\lstinline!include/unitree/idl/go2/SportModeState\_.hpp:44!}
\end{itemize}

\textbf{用法}

\begin{lstlisting}[language=Python]
current_value = obj.foot_speed_body
obj.foot_speed_body = new_value
\end{lstlisting}

\begin{quote}
{[}!TIP{]} 读取容器或嵌套值后应遵循``读取、修改、重新赋值''。只修改
getter 返回的副本不会可靠地更新原 C++ 消息。
\end{quote}

\hypertarget{unitree-sdk2-cpp-idl-go2-sportmodestate-path-point}{%
\paragraph{\texorpdfstring{\texttt{unitree\_sdk2\_cpp.idl.go2.SportModeState.path\_point}}{unitree\_sdk2\_cpp.idl.go2.SportModeState.path\_point}}\label{unitree-sdk2-cpp-idl-go2-sportmodestate-path-point}}

底层 \passthrough{\lstinline!unitree\_go::msg::dds\_::SportModeState\_!}
的 \passthrough{\lstinline!path\_point!} 字段。类型签名只说明 Python
表示；单位、数值范围、数组固定长度和枚举含义需查对应 IDL 头文件。
底层是固定长度数组，写入序列必须正好包含 10 个元素

\textbf{签名}

\begin{lstlisting}[language=Python]
@property
def path_point(self) -> list[PathPoint]

@path_point.setter
def path_point(self, value: Sequence[PathPoint]) -> None
\end{lstlisting}

\textbf{可用性与安全性}

\textbf{\passthrough{\lstinline!AVAILABLE!} /
\passthrough{\lstinline!VALUE\_TYPE!}}：当前绑定源码已实现该消息值操作。

\textbf{参数}

\begin{longtable}[]{@{}
  >{\raggedright\arraybackslash}p{(\columnwidth - 4\tabcolsep) * \real{0.18}}
  >{\raggedright\arraybackslash}p{(\columnwidth - 4\tabcolsep) * \real{0.20}}
  >{\raggedright\arraybackslash}p{(\columnwidth - 4\tabcolsep) * \real{0.62}}@{}}
\toprule
名称 & Python 类型 & 含义 \\
\midrule
\endhead
\passthrough{\lstinline!value!} &
\passthrough{\lstinline!Sequence[PathPoint]!} &
要写入或传递的新值，必须符合该参数的 Python 类型以及底层 C++
范围约束。 \\
\bottomrule
\end{longtable}

\textbf{返回值}

\begin{longtable}[]{@{}
  >{\raggedright\arraybackslash}p{(\columnwidth - 4\tabcolsep) * \real{0.18}}
  >{\raggedright\arraybackslash}p{(\columnwidth - 4\tabcolsep) * \real{0.20}}
  >{\raggedright\arraybackslash}p{(\columnwidth - 4\tabcolsep) * \real{0.62}}@{}}
\toprule
位置 & 类型 & 含义 \\
\midrule
\endhead
getter 返回值 & \passthrough{\lstinline!list[PathPoint]!} &
返回属性当前值。数组、vector 和嵌套消息采用复制语义。 \\
\bottomrule
\end{longtable}

\textbf{对应 C++}

\begin{itemize}
\tightlist
\item
  类：\passthrough{\lstinline!unitree\_go::msg::dds\_::SportModeState\_!}
\item
  字段类型：\passthrough{\lstinline!std::array< ::unitree\_go::msg::dds\_::PathPoint\_, 10>!}
\item
  声明位置：\passthrough{\lstinline!include/unitree/idl/go2/SportModeState\_.hpp:45!}
\end{itemize}

\textbf{用法}

\begin{lstlisting}[language=Python]
current_value = obj.path_point
obj.path_point = new_value
\end{lstlisting}

\begin{quote}
{[}!TIP{]} 读取容器或嵌套值后应遵循``读取、修改、重新赋值''。只修改
getter 返回的副本不会可靠地更新原 C++ 消息。
\end{quote}

\hypertarget{unitree-sdk2-cpp-idl-go2-uwbstate}{%
\subsubsection{\texorpdfstring{\texttt{unitree\_sdk2\_cpp.idl.go2.UwbState}}{unitree\_sdk2\_cpp.idl.go2.UwbState}}\label{unitree-sdk2-cpp-idl-go2-uwbstate}}

可由 Python 读写的 DDS/IDL 消息值类型。字段采用复制语义。

\textbf{导入}

\begin{lstlisting}[language=Python]
from unitree_sdk2_cpp.idl.go2 import UwbState
\end{lstlisting}

\textbf{方法索引}

\begin{longtable}[]{@{}
  >{\raggedright\arraybackslash}p{(\columnwidth - 6\tabcolsep) * \real{0.18}}
  >{\raggedright\arraybackslash}p{(\columnwidth - 6\tabcolsep) * \real{0.24}}
  >{\raggedright\arraybackslash}p{(\columnwidth - 6\tabcolsep) * \real{0.18}}
  >{\raggedright\arraybackslash}p{(\columnwidth - 6\tabcolsep) * \real{0.40}}@{}}
\toprule
方法 & Python 签名 & 状态 & 安全分类 \\
\midrule
\endhead
\protect\hyperlink{unitree-sdk2-cpp-idl-go2-uwbstate-dunder-init-1}{\passthrough{\lstinline!\_\_init\_\_!}}
& \passthrough{\lstinline!def \_\_init\_\_(self) -> None!} &
\passthrough{\lstinline!AVAILABLE!} &
\passthrough{\lstinline!VALUE\_TYPE!} \\
\protect\hyperlink{unitree-sdk2-cpp-idl-go2-uwbstate-dunder-eq-1}{\passthrough{\lstinline!\_\_eq\_\_!}}
& \passthrough{\lstinline!def \_\_eq\_\_(self, other: object) -> bool!}
& \passthrough{\lstinline!AVAILABLE!} &
\passthrough{\lstinline!VALUE\_TYPE!} \\
\protect\hyperlink{unitree-sdk2-cpp-idl-go2-uwbstate-dunder-ne-1}{\passthrough{\lstinline!\_\_ne\_\_!}}
& \passthrough{\lstinline!def \_\_ne\_\_(self, other: object) -> bool!}
& \passthrough{\lstinline!AVAILABLE!} &
\passthrough{\lstinline!VALUE\_TYPE!} \\
\bottomrule
\end{longtable}

\textbf{属性索引}

\begin{longtable}[]{@{}
  >{\raggedright\arraybackslash}p{(\columnwidth - 6\tabcolsep) * \real{0.18}}
  >{\raggedright\arraybackslash}p{(\columnwidth - 6\tabcolsep) * \real{0.24}}
  >{\raggedright\arraybackslash}p{(\columnwidth - 6\tabcolsep) * \real{0.18}}
  >{\raggedright\arraybackslash}p{(\columnwidth - 6\tabcolsep) * \real{0.40}}@{}}
\toprule
属性 & 读取类型 & 写入类型 & 状态 \\
\midrule
\endhead
\protect\hyperlink{unitree-sdk2-cpp-idl-go2-uwbstate-version}{\passthrough{\lstinline!version!}}
& \passthrough{\lstinline!list[int]!} &
\passthrough{\lstinline!Sequence[int]!} &
\passthrough{\lstinline!AVAILABLE!} \\
\protect\hyperlink{unitree-sdk2-cpp-idl-go2-uwbstate-channel}{\passthrough{\lstinline!channel!}}
& \passthrough{\lstinline!int!} & \passthrough{\lstinline!int!} &
\passthrough{\lstinline!AVAILABLE!} \\
\protect\hyperlink{unitree-sdk2-cpp-idl-go2-uwbstate-joy-mode}{\passthrough{\lstinline!joy\_mode!}}
& \passthrough{\lstinline!int!} & \passthrough{\lstinline!int!} &
\passthrough{\lstinline!AVAILABLE!} \\
\protect\hyperlink{unitree-sdk2-cpp-idl-go2-uwbstate-orientation-est}{\passthrough{\lstinline!orientation\_est!}}
& \passthrough{\lstinline!float!} & \passthrough{\lstinline!float!} &
\passthrough{\lstinline!AVAILABLE!} \\
\protect\hyperlink{unitree-sdk2-cpp-idl-go2-uwbstate-pitch-est}{\passthrough{\lstinline!pitch\_est!}}
& \passthrough{\lstinline!float!} & \passthrough{\lstinline!float!} &
\passthrough{\lstinline!AVAILABLE!} \\
\protect\hyperlink{unitree-sdk2-cpp-idl-go2-uwbstate-distance-est}{\passthrough{\lstinline!distance\_est!}}
& \passthrough{\lstinline!float!} & \passthrough{\lstinline!float!} &
\passthrough{\lstinline!AVAILABLE!} \\
\protect\hyperlink{unitree-sdk2-cpp-idl-go2-uwbstate-yaw-est}{\passthrough{\lstinline!yaw\_est!}}
& \passthrough{\lstinline!float!} & \passthrough{\lstinline!float!} &
\passthrough{\lstinline!AVAILABLE!} \\
\protect\hyperlink{unitree-sdk2-cpp-idl-go2-uwbstate-tag-roll}{\passthrough{\lstinline!tag\_roll!}}
& \passthrough{\lstinline!float!} & \passthrough{\lstinline!float!} &
\passthrough{\lstinline!AVAILABLE!} \\
\protect\hyperlink{unitree-sdk2-cpp-idl-go2-uwbstate-tag-pitch}{\passthrough{\lstinline!tag\_pitch!}}
& \passthrough{\lstinline!float!} & \passthrough{\lstinline!float!} &
\passthrough{\lstinline!AVAILABLE!} \\
\protect\hyperlink{unitree-sdk2-cpp-idl-go2-uwbstate-tag-yaw}{\passthrough{\lstinline!tag\_yaw!}}
& \passthrough{\lstinline!float!} & \passthrough{\lstinline!float!} &
\passthrough{\lstinline!AVAILABLE!} \\
\protect\hyperlink{unitree-sdk2-cpp-idl-go2-uwbstate-base-roll}{\passthrough{\lstinline!base\_roll!}}
& \passthrough{\lstinline!float!} & \passthrough{\lstinline!float!} &
\passthrough{\lstinline!AVAILABLE!} \\
\protect\hyperlink{unitree-sdk2-cpp-idl-go2-uwbstate-base-pitch}{\passthrough{\lstinline!base\_pitch!}}
& \passthrough{\lstinline!float!} & \passthrough{\lstinline!float!} &
\passthrough{\lstinline!AVAILABLE!} \\
\protect\hyperlink{unitree-sdk2-cpp-idl-go2-uwbstate-base-yaw}{\passthrough{\lstinline!base\_yaw!}}
& \passthrough{\lstinline!float!} & \passthrough{\lstinline!float!} &
\passthrough{\lstinline!AVAILABLE!} \\
\protect\hyperlink{unitree-sdk2-cpp-idl-go2-uwbstate-joystick}{\passthrough{\lstinline!joystick!}}
& \passthrough{\lstinline!list[float]!} &
\passthrough{\lstinline!Sequence[float]!} &
\passthrough{\lstinline!AVAILABLE!} \\
\protect\hyperlink{unitree-sdk2-cpp-idl-go2-uwbstate-error-state}{\passthrough{\lstinline!error\_state!}}
& \passthrough{\lstinline!int!} & \passthrough{\lstinline!int!} &
\passthrough{\lstinline!AVAILABLE!} \\
\protect\hyperlink{unitree-sdk2-cpp-idl-go2-uwbstate-buttons}{\passthrough{\lstinline!buttons!}}
& \passthrough{\lstinline!int!} & \passthrough{\lstinline!int!} &
\passthrough{\lstinline!AVAILABLE!} \\
\protect\hyperlink{unitree-sdk2-cpp-idl-go2-uwbstate-enabled-from-app}{\passthrough{\lstinline!enabled\_from\_app!}}
& \passthrough{\lstinline!int!} & \passthrough{\lstinline!int!} &
\passthrough{\lstinline!AVAILABLE!} \\
\bottomrule
\end{longtable}

\hypertarget{unitree-sdk2-cpp-idl-go2-uwbstate-dunder-init-1}{%
\paragraph{\texorpdfstring{\texttt{unitree\_sdk2\_cpp.idl.go2.UwbState.\_\_init\_\_}}{unitree\_sdk2\_cpp.idl.go2.UwbState.\_\_init\_\_}}\label{unitree-sdk2-cpp-idl-go2-uwbstate-dunder-init-1}}

初始化 \passthrough{\lstinline!UwbState!}
实例。是否能实际构造取决于下方可用性状态。

\textbf{签名}

\begin{lstlisting}[language=Python]
def __init__(self) -> None
\end{lstlisting}

\textbf{可用性与安全性}

\textbf{\passthrough{\lstinline!AVAILABLE!} /
\passthrough{\lstinline!VALUE\_TYPE!}}：当前绑定源码已实现该消息值操作。

\textbf{参数}

无显式参数。实例方法中的 \passthrough{\lstinline!self!} 由 Python
自动传入。

\textbf{返回值}

\begin{longtable}[]{@{}
  >{\raggedright\arraybackslash}p{(\columnwidth - 4\tabcolsep) * \real{0.18}}
  >{\raggedright\arraybackslash}p{(\columnwidth - 4\tabcolsep) * \real{0.20}}
  >{\raggedright\arraybackslash}p{(\columnwidth - 4\tabcolsep) * \real{0.62}}@{}}
\toprule
位置 & 类型 & 含义 \\
\midrule
\endhead
无 & \passthrough{\lstinline!None!} &
\passthrough{\lstinline!\_\_init\_\_!}
本身不返回值；调用类对象时，在构造函数可用的前提下得到该类实例。 \\
\bottomrule
\end{longtable}

\textbf{对应 C++}

\begin{itemize}
\tightlist
\item
  类：\passthrough{\lstinline!unitree\_go::msg::dds\_::UwbState\_!}
\item
  签名：\passthrough{\lstinline!UwbState\_()!}
\item
  绑定策略：\passthrough{\lstinline!未单独标注!}
\item
  声明位置：\passthrough{\lstinline!include/unitree/idl/go2/UwbState\_.hpp:43!}
\end{itemize}

\textbf{说明}

\begin{itemize}
\tightlist
\item
  示例是局部调用片段；\passthrough{\lstinline!obj!}
  和参数变量表示已经按上层业务校验并准备好的对象。
\end{itemize}

\textbf{用法}

\begin{lstlisting}[language=Python]
value = UwbState()
\end{lstlisting}

\hypertarget{unitree-sdk2-cpp-idl-go2-uwbstate-dunder-eq-1}{%
\paragraph{\texorpdfstring{\texttt{unitree\_sdk2\_cpp.idl.go2.UwbState.\_\_eq\_\_}}{unitree\_sdk2\_cpp.idl.go2.UwbState.\_\_eq\_\_}}\label{unitree-sdk2-cpp-idl-go2-uwbstate-dunder-eq-1}}

按底层 IDL 消息内容比较两个对象是否相等。

\textbf{签名}

\begin{lstlisting}[language=Python]
def __eq__(self, other: object) -> bool
\end{lstlisting}

\textbf{可用性与安全性}

\textbf{\passthrough{\lstinline!AVAILABLE!} /
\passthrough{\lstinline!VALUE\_TYPE!}}：当前绑定源码已实现该消息值操作。

\textbf{参数}

\begin{longtable}[]{@{}
  >{\raggedright\arraybackslash}p{(\columnwidth - 6\tabcolsep) * \real{0.18}}
  >{\raggedright\arraybackslash}p{(\columnwidth - 6\tabcolsep) * \real{0.24}}
  >{\raggedright\arraybackslash}p{(\columnwidth - 6\tabcolsep) * \real{0.18}}
  >{\raggedright\arraybackslash}p{(\columnwidth - 6\tabcolsep) * \real{0.40}}@{}}
\toprule
名称 & Python 类型 & 默认值 & 含义 \\
\midrule
\endhead
\passthrough{\lstinline!other!} & \passthrough{\lstinline!object!} &
必填 & 用于比较的另一个 Python 对象。类型不兼容时比较结果通常为
\passthrough{\lstinline!False!}。 \\
\bottomrule
\end{longtable}

\textbf{返回值}

\begin{longtable}[]{@{}
  >{\raggedright\arraybackslash}p{(\columnwidth - 4\tabcolsep) * \real{0.18}}
  >{\raggedright\arraybackslash}p{(\columnwidth - 4\tabcolsep) * \real{0.20}}
  >{\raggedright\arraybackslash}p{(\columnwidth - 4\tabcolsep) * \real{0.62}}@{}}
\toprule
位置 & 类型 & 含义 \\
\midrule
\endhead
返回值 & \passthrough{\lstinline!bool!} & 返回
\passthrough{\lstinline!bool!}。更精确的业务含义见该方法用途、C++
签名和目标型号协议。 \\
\bottomrule
\end{longtable}

\textbf{对应 C++}

\begin{itemize}
\tightlist
\item
  类：\passthrough{\lstinline!unitree\_go::msg::dds\_::UwbState\_!}
\item
  签名：\passthrough{\lstinline!operator==(const value \&) const!}
\item
  绑定策略：\passthrough{\lstinline!未单独标注!}
\item
  声明位置：由 pybind11 手工绑定或生成注册表提供；manifest
  未记录独立头文件行号。
\end{itemize}

\textbf{说明}

\begin{itemize}
\tightlist
\item
  示例是局部调用片段；\passthrough{\lstinline!obj!}
  和参数变量表示已经按上层业务校验并准备好的对象。
\end{itemize}

\textbf{用法}

\begin{lstlisting}[language=Python]
same = left == right
\end{lstlisting}

\hypertarget{unitree-sdk2-cpp-idl-go2-uwbstate-dunder-ne-1}{%
\paragraph{\texorpdfstring{\texttt{unitree\_sdk2\_cpp.idl.go2.UwbState.\_\_ne\_\_}}{unitree\_sdk2\_cpp.idl.go2.UwbState.\_\_ne\_\_}}\label{unitree-sdk2-cpp-idl-go2-uwbstate-dunder-ne-1}}

按底层 IDL 消息内容比较两个对象是否不相等。

\textbf{签名}

\begin{lstlisting}[language=Python]
def __ne__(self, other: object) -> bool
\end{lstlisting}

\textbf{可用性与安全性}

\textbf{\passthrough{\lstinline!AVAILABLE!} /
\passthrough{\lstinline!VALUE\_TYPE!}}：当前绑定源码已实现该消息值操作。

\textbf{参数}

\begin{longtable}[]{@{}
  >{\raggedright\arraybackslash}p{(\columnwidth - 6\tabcolsep) * \real{0.18}}
  >{\raggedright\arraybackslash}p{(\columnwidth - 6\tabcolsep) * \real{0.24}}
  >{\raggedright\arraybackslash}p{(\columnwidth - 6\tabcolsep) * \real{0.18}}
  >{\raggedright\arraybackslash}p{(\columnwidth - 6\tabcolsep) * \real{0.40}}@{}}
\toprule
名称 & Python 类型 & 默认值 & 含义 \\
\midrule
\endhead
\passthrough{\lstinline!other!} & \passthrough{\lstinline!object!} &
必填 & 用于比较的另一个 Python 对象。类型不兼容时比较结果通常为
\passthrough{\lstinline!False!}。 \\
\bottomrule
\end{longtable}

\textbf{返回值}

\begin{longtable}[]{@{}
  >{\raggedright\arraybackslash}p{(\columnwidth - 4\tabcolsep) * \real{0.18}}
  >{\raggedright\arraybackslash}p{(\columnwidth - 4\tabcolsep) * \real{0.20}}
  >{\raggedright\arraybackslash}p{(\columnwidth - 4\tabcolsep) * \real{0.62}}@{}}
\toprule
位置 & 类型 & 含义 \\
\midrule
\endhead
返回值 & \passthrough{\lstinline!bool!} & 返回
\passthrough{\lstinline!bool!}。更精确的业务含义见该方法用途、C++
签名和目标型号协议。 \\
\bottomrule
\end{longtable}

\textbf{对应 C++}

\begin{itemize}
\tightlist
\item
  类：\passthrough{\lstinline!unitree\_go::msg::dds\_::UwbState\_!}
\item
  签名：\passthrough{\lstinline"operator!=(const value \&) const"}
\item
  绑定策略：\passthrough{\lstinline!未单独标注!}
\item
  声明位置：由 pybind11 手工绑定或生成注册表提供；manifest
  未记录独立头文件行号。
\end{itemize}

\textbf{说明}

\begin{itemize}
\tightlist
\item
  示例是局部调用片段；\passthrough{\lstinline!obj!}
  和参数变量表示已经按上层业务校验并准备好的对象。
\end{itemize}

\textbf{用法}

\begin{lstlisting}[language=Python]
different = left != right
\end{lstlisting}

\hypertarget{unitree-sdk2-cpp-idl-go2-uwbstate-version}{%
\paragraph{\texorpdfstring{\texttt{unitree\_sdk2\_cpp.idl.go2.UwbState.version}}{unitree\_sdk2\_cpp.idl.go2.UwbState.version}}\label{unitree-sdk2-cpp-idl-go2-uwbstate-version}}

底层 \passthrough{\lstinline!unitree\_go::msg::dds\_::UwbState\_!} 的
\passthrough{\lstinline!version!} 字段。类型签名只说明 Python
表示；单位、数值范围、数组固定长度和枚举含义需查对应 IDL 头文件。
底层是固定长度数组，写入序列必须正好包含 2 个元素；元素约束：取值范围为
0 到 255。

\textbf{签名}

\begin{lstlisting}[language=Python]
@property
def version(self) -> list[int]

@version.setter
def version(self, value: Sequence[int]) -> None
\end{lstlisting}

\textbf{可用性与安全性}

\textbf{\passthrough{\lstinline!AVAILABLE!} /
\passthrough{\lstinline!VALUE\_TYPE!}}：当前绑定源码已实现该消息值操作。

\textbf{参数}

\begin{longtable}[]{@{}
  >{\raggedright\arraybackslash}p{(\columnwidth - 4\tabcolsep) * \real{0.18}}
  >{\raggedright\arraybackslash}p{(\columnwidth - 4\tabcolsep) * \real{0.20}}
  >{\raggedright\arraybackslash}p{(\columnwidth - 4\tabcolsep) * \real{0.62}}@{}}
\toprule
名称 & Python 类型 & 含义 \\
\midrule
\endhead
\passthrough{\lstinline!value!} &
\passthrough{\lstinline!Sequence[int]!} &
要写入或传递的新值，必须符合该参数的 Python 类型以及底层 C++
范围约束。 \\
\bottomrule
\end{longtable}

\textbf{返回值}

\begin{longtable}[]{@{}
  >{\raggedright\arraybackslash}p{(\columnwidth - 4\tabcolsep) * \real{0.18}}
  >{\raggedright\arraybackslash}p{(\columnwidth - 4\tabcolsep) * \real{0.20}}
  >{\raggedright\arraybackslash}p{(\columnwidth - 4\tabcolsep) * \real{0.62}}@{}}
\toprule
位置 & 类型 & 含义 \\
\midrule
\endhead
getter 返回值 & \passthrough{\lstinline!list[int]!} &
返回属性当前值。数组、vector 和嵌套消息采用复制语义。 \\
\bottomrule
\end{longtable}

\textbf{对应 C++}

\begin{itemize}
\tightlist
\item
  类：\passthrough{\lstinline!unitree\_go::msg::dds\_::UwbState\_!}
\item
  字段类型：\passthrough{\lstinline!std::array<uint8\_t, 2>!}
\item
  声明位置：\passthrough{\lstinline!include/unitree/idl/go2/UwbState\_.hpp:24!}
\end{itemize}

\textbf{用法}

\begin{lstlisting}[language=Python]
current_value = obj.version
obj.version = new_value
\end{lstlisting}

\begin{quote}
{[}!TIP{]} 读取容器或嵌套值后应遵循``读取、修改、重新赋值''。只修改
getter 返回的副本不会可靠地更新原 C++ 消息。
\end{quote}

\hypertarget{unitree-sdk2-cpp-idl-go2-uwbstate-channel}{%
\paragraph{\texorpdfstring{\texttt{unitree\_sdk2\_cpp.idl.go2.UwbState.channel}}{unitree\_sdk2\_cpp.idl.go2.UwbState.channel}}\label{unitree-sdk2-cpp-idl-go2-uwbstate-channel}}

底层 \passthrough{\lstinline!unitree\_go::msg::dds\_::UwbState\_!} 的
\passthrough{\lstinline!channel!} 字段。类型签名只说明 Python
表示；单位、数值范围、数组固定长度和枚举含义需查对应 IDL 头文件。
取值范围为 0 到 255。

\textbf{签名}

\begin{lstlisting}[language=Python]
@property
def channel(self) -> int

@channel.setter
def channel(self, value: int) -> None
\end{lstlisting}

\textbf{可用性与安全性}

\textbf{\passthrough{\lstinline!AVAILABLE!} /
\passthrough{\lstinline!VALUE\_TYPE!}}：当前绑定源码已实现该消息值操作。

\textbf{参数}

\begin{longtable}[]{@{}
  >{\raggedright\arraybackslash}p{(\columnwidth - 4\tabcolsep) * \real{0.18}}
  >{\raggedright\arraybackslash}p{(\columnwidth - 4\tabcolsep) * \real{0.20}}
  >{\raggedright\arraybackslash}p{(\columnwidth - 4\tabcolsep) * \real{0.62}}@{}}
\toprule
名称 & Python 类型 & 含义 \\
\midrule
\endhead
\passthrough{\lstinline!value!} & \passthrough{\lstinline!int!} &
要写入或传递的新值，必须符合该参数的 Python 类型以及底层 C++
范围约束。 \\
\bottomrule
\end{longtable}

\textbf{返回值}

\begin{longtable}[]{@{}
  >{\raggedright\arraybackslash}p{(\columnwidth - 4\tabcolsep) * \real{0.18}}
  >{\raggedright\arraybackslash}p{(\columnwidth - 4\tabcolsep) * \real{0.20}}
  >{\raggedright\arraybackslash}p{(\columnwidth - 4\tabcolsep) * \real{0.62}}@{}}
\toprule
位置 & 类型 & 含义 \\
\midrule
\endhead
getter 返回值 & \passthrough{\lstinline!int!} & 返回属性当前值。 \\
\bottomrule
\end{longtable}

\textbf{对应 C++}

\begin{itemize}
\tightlist
\item
  类：\passthrough{\lstinline!unitree\_go::msg::dds\_::UwbState\_!}
\item
  字段类型：\passthrough{\lstinline!uint8\_t!}
\item
  声明位置：\passthrough{\lstinline!include/unitree/idl/go2/UwbState\_.hpp:25!}
\end{itemize}

\textbf{用法}

\begin{lstlisting}[language=Python]
current_value = obj.channel
obj.channel = new_value
\end{lstlisting}

\hypertarget{unitree-sdk2-cpp-idl-go2-uwbstate-joy-mode}{%
\paragraph{\texorpdfstring{\texttt{unitree\_sdk2\_cpp.idl.go2.UwbState.joy\_mode}}{unitree\_sdk2\_cpp.idl.go2.UwbState.joy\_mode}}\label{unitree-sdk2-cpp-idl-go2-uwbstate-joy-mode}}

底层 \passthrough{\lstinline!unitree\_go::msg::dds\_::UwbState\_!} 的
\passthrough{\lstinline!joy\_mode!} 字段。类型签名只说明 Python
表示；单位、数值范围、数组固定长度和枚举含义需查对应 IDL 头文件。
取值范围为 0 到 255。

\textbf{签名}

\begin{lstlisting}[language=Python]
@property
def joy_mode(self) -> int

@joy_mode.setter
def joy_mode(self, value: int) -> None
\end{lstlisting}

\textbf{可用性与安全性}

\textbf{\passthrough{\lstinline!AVAILABLE!} /
\passthrough{\lstinline!VALUE\_TYPE!}}：当前绑定源码已实现该消息值操作。

\textbf{参数}

\begin{longtable}[]{@{}
  >{\raggedright\arraybackslash}p{(\columnwidth - 4\tabcolsep) * \real{0.18}}
  >{\raggedright\arraybackslash}p{(\columnwidth - 4\tabcolsep) * \real{0.20}}
  >{\raggedright\arraybackslash}p{(\columnwidth - 4\tabcolsep) * \real{0.62}}@{}}
\toprule
名称 & Python 类型 & 含义 \\
\midrule
\endhead
\passthrough{\lstinline!value!} & \passthrough{\lstinline!int!} &
要写入或传递的新值，必须符合该参数的 Python 类型以及底层 C++
范围约束。 \\
\bottomrule
\end{longtable}

\textbf{返回值}

\begin{longtable}[]{@{}
  >{\raggedright\arraybackslash}p{(\columnwidth - 4\tabcolsep) * \real{0.18}}
  >{\raggedright\arraybackslash}p{(\columnwidth - 4\tabcolsep) * \real{0.20}}
  >{\raggedright\arraybackslash}p{(\columnwidth - 4\tabcolsep) * \real{0.62}}@{}}
\toprule
位置 & 类型 & 含义 \\
\midrule
\endhead
getter 返回值 & \passthrough{\lstinline!int!} & 返回属性当前值。 \\
\bottomrule
\end{longtable}

\textbf{对应 C++}

\begin{itemize}
\tightlist
\item
  类：\passthrough{\lstinline!unitree\_go::msg::dds\_::UwbState\_!}
\item
  字段类型：\passthrough{\lstinline!uint8\_t!}
\item
  声明位置：\passthrough{\lstinline!include/unitree/idl/go2/UwbState\_.hpp:26!}
\end{itemize}

\textbf{用法}

\begin{lstlisting}[language=Python]
current_value = obj.joy_mode
obj.joy_mode = new_value
\end{lstlisting}

\hypertarget{unitree-sdk2-cpp-idl-go2-uwbstate-orientation-est}{%
\paragraph{\texorpdfstring{\texttt{unitree\_sdk2\_cpp.idl.go2.UwbState.orientation\_est}}{unitree\_sdk2\_cpp.idl.go2.UwbState.orientation\_est}}\label{unitree-sdk2-cpp-idl-go2-uwbstate-orientation-est}}

底层 \passthrough{\lstinline!unitree\_go::msg::dds\_::UwbState\_!} 的
\passthrough{\lstinline!orientation\_est!} 字段。类型签名只说明 Python
表示；单位、数值范围、数组固定长度和枚举含义需查对应 IDL 头文件。
底层为浮点数；类型本身不说明单位、坐标系或业务安全范围。

\textbf{签名}

\begin{lstlisting}[language=Python]
@property
def orientation_est(self) -> float

@orientation_est.setter
def orientation_est(self, value: float) -> None
\end{lstlisting}

\textbf{可用性与安全性}

\textbf{\passthrough{\lstinline!AVAILABLE!} /
\passthrough{\lstinline!VALUE\_TYPE!}}：当前绑定源码已实现该消息值操作。

\textbf{参数}

\begin{longtable}[]{@{}
  >{\raggedright\arraybackslash}p{(\columnwidth - 4\tabcolsep) * \real{0.18}}
  >{\raggedright\arraybackslash}p{(\columnwidth - 4\tabcolsep) * \real{0.20}}
  >{\raggedright\arraybackslash}p{(\columnwidth - 4\tabcolsep) * \real{0.62}}@{}}
\toprule
名称 & Python 类型 & 含义 \\
\midrule
\endhead
\passthrough{\lstinline!value!} & \passthrough{\lstinline!float!} &
要写入或传递的新值，必须符合该参数的 Python 类型以及底层 C++
范围约束。 \\
\bottomrule
\end{longtable}

\textbf{返回值}

\begin{longtable}[]{@{}
  >{\raggedright\arraybackslash}p{(\columnwidth - 4\tabcolsep) * \real{0.18}}
  >{\raggedright\arraybackslash}p{(\columnwidth - 4\tabcolsep) * \real{0.20}}
  >{\raggedright\arraybackslash}p{(\columnwidth - 4\tabcolsep) * \real{0.62}}@{}}
\toprule
位置 & 类型 & 含义 \\
\midrule
\endhead
getter 返回值 & \passthrough{\lstinline!float!} & 返回属性当前值。 \\
\bottomrule
\end{longtable}

\textbf{对应 C++}

\begin{itemize}
\tightlist
\item
  类：\passthrough{\lstinline!unitree\_go::msg::dds\_::UwbState\_!}
\item
  字段类型：\passthrough{\lstinline!float!}
\item
  声明位置：\passthrough{\lstinline!include/unitree/idl/go2/UwbState\_.hpp:27!}
\end{itemize}

\textbf{用法}

\begin{lstlisting}[language=Python]
current_value = obj.orientation_est
obj.orientation_est = new_value
\end{lstlisting}

\hypertarget{unitree-sdk2-cpp-idl-go2-uwbstate-pitch-est}{%
\paragraph{\texorpdfstring{\texttt{unitree\_sdk2\_cpp.idl.go2.UwbState.pitch\_est}}{unitree\_sdk2\_cpp.idl.go2.UwbState.pitch\_est}}\label{unitree-sdk2-cpp-idl-go2-uwbstate-pitch-est}}

底层 \passthrough{\lstinline!unitree\_go::msg::dds\_::UwbState\_!} 的
\passthrough{\lstinline!pitch\_est!} 字段。类型签名只说明 Python
表示；单位、数值范围、数组固定长度和枚举含义需查对应 IDL 头文件。
底层为浮点数；类型本身不说明单位、坐标系或业务安全范围。

\textbf{签名}

\begin{lstlisting}[language=Python]
@property
def pitch_est(self) -> float

@pitch_est.setter
def pitch_est(self, value: float) -> None
\end{lstlisting}

\textbf{可用性与安全性}

\textbf{\passthrough{\lstinline!AVAILABLE!} /
\passthrough{\lstinline!VALUE\_TYPE!}}：当前绑定源码已实现该消息值操作。

\textbf{参数}

\begin{longtable}[]{@{}
  >{\raggedright\arraybackslash}p{(\columnwidth - 4\tabcolsep) * \real{0.18}}
  >{\raggedright\arraybackslash}p{(\columnwidth - 4\tabcolsep) * \real{0.20}}
  >{\raggedright\arraybackslash}p{(\columnwidth - 4\tabcolsep) * \real{0.62}}@{}}
\toprule
名称 & Python 类型 & 含义 \\
\midrule
\endhead
\passthrough{\lstinline!value!} & \passthrough{\lstinline!float!} &
要写入或传递的新值，必须符合该参数的 Python 类型以及底层 C++
范围约束。 \\
\bottomrule
\end{longtable}

\textbf{返回值}

\begin{longtable}[]{@{}
  >{\raggedright\arraybackslash}p{(\columnwidth - 4\tabcolsep) * \real{0.18}}
  >{\raggedright\arraybackslash}p{(\columnwidth - 4\tabcolsep) * \real{0.20}}
  >{\raggedright\arraybackslash}p{(\columnwidth - 4\tabcolsep) * \real{0.62}}@{}}
\toprule
位置 & 类型 & 含义 \\
\midrule
\endhead
getter 返回值 & \passthrough{\lstinline!float!} & 返回属性当前值。 \\
\bottomrule
\end{longtable}

\textbf{对应 C++}

\begin{itemize}
\tightlist
\item
  类：\passthrough{\lstinline!unitree\_go::msg::dds\_::UwbState\_!}
\item
  字段类型：\passthrough{\lstinline!float!}
\item
  声明位置：\passthrough{\lstinline!include/unitree/idl/go2/UwbState\_.hpp:28!}
\end{itemize}

\textbf{用法}

\begin{lstlisting}[language=Python]
current_value = obj.pitch_est
obj.pitch_est = new_value
\end{lstlisting}

\hypertarget{unitree-sdk2-cpp-idl-go2-uwbstate-distance-est}{%
\paragraph{\texorpdfstring{\texttt{unitree\_sdk2\_cpp.idl.go2.UwbState.distance\_est}}{unitree\_sdk2\_cpp.idl.go2.UwbState.distance\_est}}\label{unitree-sdk2-cpp-idl-go2-uwbstate-distance-est}}

底层 \passthrough{\lstinline!unitree\_go::msg::dds\_::UwbState\_!} 的
\passthrough{\lstinline!distance\_est!} 字段。类型签名只说明 Python
表示；单位、数值范围、数组固定长度和枚举含义需查对应 IDL 头文件。
底层为浮点数；类型本身不说明单位、坐标系或业务安全范围。

\textbf{签名}

\begin{lstlisting}[language=Python]
@property
def distance_est(self) -> float

@distance_est.setter
def distance_est(self, value: float) -> None
\end{lstlisting}

\textbf{可用性与安全性}

\textbf{\passthrough{\lstinline!AVAILABLE!} /
\passthrough{\lstinline!VALUE\_TYPE!}}：当前绑定源码已实现该消息值操作。

\textbf{参数}

\begin{longtable}[]{@{}
  >{\raggedright\arraybackslash}p{(\columnwidth - 4\tabcolsep) * \real{0.18}}
  >{\raggedright\arraybackslash}p{(\columnwidth - 4\tabcolsep) * \real{0.20}}
  >{\raggedright\arraybackslash}p{(\columnwidth - 4\tabcolsep) * \real{0.62}}@{}}
\toprule
名称 & Python 类型 & 含义 \\
\midrule
\endhead
\passthrough{\lstinline!value!} & \passthrough{\lstinline!float!} &
要写入或传递的新值，必须符合该参数的 Python 类型以及底层 C++
范围约束。 \\
\bottomrule
\end{longtable}

\textbf{返回值}

\begin{longtable}[]{@{}
  >{\raggedright\arraybackslash}p{(\columnwidth - 4\tabcolsep) * \real{0.18}}
  >{\raggedright\arraybackslash}p{(\columnwidth - 4\tabcolsep) * \real{0.20}}
  >{\raggedright\arraybackslash}p{(\columnwidth - 4\tabcolsep) * \real{0.62}}@{}}
\toprule
位置 & 类型 & 含义 \\
\midrule
\endhead
getter 返回值 & \passthrough{\lstinline!float!} & 返回属性当前值。 \\
\bottomrule
\end{longtable}

\textbf{对应 C++}

\begin{itemize}
\tightlist
\item
  类：\passthrough{\lstinline!unitree\_go::msg::dds\_::UwbState\_!}
\item
  字段类型：\passthrough{\lstinline!float!}
\item
  声明位置：\passthrough{\lstinline!include/unitree/idl/go2/UwbState\_.hpp:29!}
\end{itemize}

\textbf{用法}

\begin{lstlisting}[language=Python]
current_value = obj.distance_est
obj.distance_est = new_value
\end{lstlisting}

\hypertarget{unitree-sdk2-cpp-idl-go2-uwbstate-yaw-est}{%
\paragraph{\texorpdfstring{\texttt{unitree\_sdk2\_cpp.idl.go2.UwbState.yaw\_est}}{unitree\_sdk2\_cpp.idl.go2.UwbState.yaw\_est}}\label{unitree-sdk2-cpp-idl-go2-uwbstate-yaw-est}}

底层 \passthrough{\lstinline!unitree\_go::msg::dds\_::UwbState\_!} 的
\passthrough{\lstinline!yaw\_est!} 字段。类型签名只说明 Python
表示；单位、数值范围、数组固定长度和枚举含义需查对应 IDL 头文件。
底层为浮点数；类型本身不说明单位、坐标系或业务安全范围。

\textbf{签名}

\begin{lstlisting}[language=Python]
@property
def yaw_est(self) -> float

@yaw_est.setter
def yaw_est(self, value: float) -> None
\end{lstlisting}

\textbf{可用性与安全性}

\textbf{\passthrough{\lstinline!AVAILABLE!} /
\passthrough{\lstinline!VALUE\_TYPE!}}：当前绑定源码已实现该消息值操作。

\textbf{参数}

\begin{longtable}[]{@{}
  >{\raggedright\arraybackslash}p{(\columnwidth - 4\tabcolsep) * \real{0.18}}
  >{\raggedright\arraybackslash}p{(\columnwidth - 4\tabcolsep) * \real{0.20}}
  >{\raggedright\arraybackslash}p{(\columnwidth - 4\tabcolsep) * \real{0.62}}@{}}
\toprule
名称 & Python 类型 & 含义 \\
\midrule
\endhead
\passthrough{\lstinline!value!} & \passthrough{\lstinline!float!} &
要写入或传递的新值，必须符合该参数的 Python 类型以及底层 C++
范围约束。 \\
\bottomrule
\end{longtable}

\textbf{返回值}

\begin{longtable}[]{@{}
  >{\raggedright\arraybackslash}p{(\columnwidth - 4\tabcolsep) * \real{0.18}}
  >{\raggedright\arraybackslash}p{(\columnwidth - 4\tabcolsep) * \real{0.20}}
  >{\raggedright\arraybackslash}p{(\columnwidth - 4\tabcolsep) * \real{0.62}}@{}}
\toprule
位置 & 类型 & 含义 \\
\midrule
\endhead
getter 返回值 & \passthrough{\lstinline!float!} & 返回属性当前值。 \\
\bottomrule
\end{longtable}

\textbf{对应 C++}

\begin{itemize}
\tightlist
\item
  类：\passthrough{\lstinline!unitree\_go::msg::dds\_::UwbState\_!}
\item
  字段类型：\passthrough{\lstinline!float!}
\item
  声明位置：\passthrough{\lstinline!include/unitree/idl/go2/UwbState\_.hpp:30!}
\end{itemize}

\textbf{用法}

\begin{lstlisting}[language=Python]
current_value = obj.yaw_est
obj.yaw_est = new_value
\end{lstlisting}

\hypertarget{unitree-sdk2-cpp-idl-go2-uwbstate-tag-roll}{%
\paragraph{\texorpdfstring{\texttt{unitree\_sdk2\_cpp.idl.go2.UwbState.tag\_roll}}{unitree\_sdk2\_cpp.idl.go2.UwbState.tag\_roll}}\label{unitree-sdk2-cpp-idl-go2-uwbstate-tag-roll}}

底层 \passthrough{\lstinline!unitree\_go::msg::dds\_::UwbState\_!} 的
\passthrough{\lstinline!tag\_roll!} 字段。类型签名只说明 Python
表示；单位、数值范围、数组固定长度和枚举含义需查对应 IDL 头文件。
底层为浮点数；类型本身不说明单位、坐标系或业务安全范围。

\textbf{签名}

\begin{lstlisting}[language=Python]
@property
def tag_roll(self) -> float

@tag_roll.setter
def tag_roll(self, value: float) -> None
\end{lstlisting}

\textbf{可用性与安全性}

\textbf{\passthrough{\lstinline!AVAILABLE!} /
\passthrough{\lstinline!VALUE\_TYPE!}}：当前绑定源码已实现该消息值操作。

\textbf{参数}

\begin{longtable}[]{@{}
  >{\raggedright\arraybackslash}p{(\columnwidth - 4\tabcolsep) * \real{0.18}}
  >{\raggedright\arraybackslash}p{(\columnwidth - 4\tabcolsep) * \real{0.20}}
  >{\raggedright\arraybackslash}p{(\columnwidth - 4\tabcolsep) * \real{0.62}}@{}}
\toprule
名称 & Python 类型 & 含义 \\
\midrule
\endhead
\passthrough{\lstinline!value!} & \passthrough{\lstinline!float!} &
要写入或传递的新值，必须符合该参数的 Python 类型以及底层 C++
范围约束。 \\
\bottomrule
\end{longtable}

\textbf{返回值}

\begin{longtable}[]{@{}
  >{\raggedright\arraybackslash}p{(\columnwidth - 4\tabcolsep) * \real{0.18}}
  >{\raggedright\arraybackslash}p{(\columnwidth - 4\tabcolsep) * \real{0.20}}
  >{\raggedright\arraybackslash}p{(\columnwidth - 4\tabcolsep) * \real{0.62}}@{}}
\toprule
位置 & 类型 & 含义 \\
\midrule
\endhead
getter 返回值 & \passthrough{\lstinline!float!} & 返回属性当前值。 \\
\bottomrule
\end{longtable}

\textbf{对应 C++}

\begin{itemize}
\tightlist
\item
  类：\passthrough{\lstinline!unitree\_go::msg::dds\_::UwbState\_!}
\item
  字段类型：\passthrough{\lstinline!float!}
\item
  声明位置：\passthrough{\lstinline!include/unitree/idl/go2/UwbState\_.hpp:31!}
\end{itemize}

\textbf{用法}

\begin{lstlisting}[language=Python]
current_value = obj.tag_roll
obj.tag_roll = new_value
\end{lstlisting}

\hypertarget{unitree-sdk2-cpp-idl-go2-uwbstate-tag-pitch}{%
\paragraph{\texorpdfstring{\texttt{unitree\_sdk2\_cpp.idl.go2.UwbState.tag\_pitch}}{unitree\_sdk2\_cpp.idl.go2.UwbState.tag\_pitch}}\label{unitree-sdk2-cpp-idl-go2-uwbstate-tag-pitch}}

底层 \passthrough{\lstinline!unitree\_go::msg::dds\_::UwbState\_!} 的
\passthrough{\lstinline!tag\_pitch!} 字段。类型签名只说明 Python
表示；单位、数值范围、数组固定长度和枚举含义需查对应 IDL 头文件。
底层为浮点数；类型本身不说明单位、坐标系或业务安全范围。

\textbf{签名}

\begin{lstlisting}[language=Python]
@property
def tag_pitch(self) -> float

@tag_pitch.setter
def tag_pitch(self, value: float) -> None
\end{lstlisting}

\textbf{可用性与安全性}

\textbf{\passthrough{\lstinline!AVAILABLE!} /
\passthrough{\lstinline!VALUE\_TYPE!}}：当前绑定源码已实现该消息值操作。

\textbf{参数}

\begin{longtable}[]{@{}
  >{\raggedright\arraybackslash}p{(\columnwidth - 4\tabcolsep) * \real{0.18}}
  >{\raggedright\arraybackslash}p{(\columnwidth - 4\tabcolsep) * \real{0.20}}
  >{\raggedright\arraybackslash}p{(\columnwidth - 4\tabcolsep) * \real{0.62}}@{}}
\toprule
名称 & Python 类型 & 含义 \\
\midrule
\endhead
\passthrough{\lstinline!value!} & \passthrough{\lstinline!float!} &
要写入或传递的新值，必须符合该参数的 Python 类型以及底层 C++
范围约束。 \\
\bottomrule
\end{longtable}

\textbf{返回值}

\begin{longtable}[]{@{}
  >{\raggedright\arraybackslash}p{(\columnwidth - 4\tabcolsep) * \real{0.18}}
  >{\raggedright\arraybackslash}p{(\columnwidth - 4\tabcolsep) * \real{0.20}}
  >{\raggedright\arraybackslash}p{(\columnwidth - 4\tabcolsep) * \real{0.62}}@{}}
\toprule
位置 & 类型 & 含义 \\
\midrule
\endhead
getter 返回值 & \passthrough{\lstinline!float!} & 返回属性当前值。 \\
\bottomrule
\end{longtable}

\textbf{对应 C++}

\begin{itemize}
\tightlist
\item
  类：\passthrough{\lstinline!unitree\_go::msg::dds\_::UwbState\_!}
\item
  字段类型：\passthrough{\lstinline!float!}
\item
  声明位置：\passthrough{\lstinline!include/unitree/idl/go2/UwbState\_.hpp:32!}
\end{itemize}

\textbf{用法}

\begin{lstlisting}[language=Python]
current_value = obj.tag_pitch
obj.tag_pitch = new_value
\end{lstlisting}

\hypertarget{unitree-sdk2-cpp-idl-go2-uwbstate-tag-yaw}{%
\paragraph{\texorpdfstring{\texttt{unitree\_sdk2\_cpp.idl.go2.UwbState.tag\_yaw}}{unitree\_sdk2\_cpp.idl.go2.UwbState.tag\_yaw}}\label{unitree-sdk2-cpp-idl-go2-uwbstate-tag-yaw}}

底层 \passthrough{\lstinline!unitree\_go::msg::dds\_::UwbState\_!} 的
\passthrough{\lstinline!tag\_yaw!} 字段。类型签名只说明 Python
表示；单位、数值范围、数组固定长度和枚举含义需查对应 IDL 头文件。
底层为浮点数；类型本身不说明单位、坐标系或业务安全范围。

\textbf{签名}

\begin{lstlisting}[language=Python]
@property
def tag_yaw(self) -> float

@tag_yaw.setter
def tag_yaw(self, value: float) -> None
\end{lstlisting}

\textbf{可用性与安全性}

\textbf{\passthrough{\lstinline!AVAILABLE!} /
\passthrough{\lstinline!VALUE\_TYPE!}}：当前绑定源码已实现该消息值操作。

\textbf{参数}

\begin{longtable}[]{@{}
  >{\raggedright\arraybackslash}p{(\columnwidth - 4\tabcolsep) * \real{0.18}}
  >{\raggedright\arraybackslash}p{(\columnwidth - 4\tabcolsep) * \real{0.20}}
  >{\raggedright\arraybackslash}p{(\columnwidth - 4\tabcolsep) * \real{0.62}}@{}}
\toprule
名称 & Python 类型 & 含义 \\
\midrule
\endhead
\passthrough{\lstinline!value!} & \passthrough{\lstinline!float!} &
要写入或传递的新值，必须符合该参数的 Python 类型以及底层 C++
范围约束。 \\
\bottomrule
\end{longtable}

\textbf{返回值}

\begin{longtable}[]{@{}
  >{\raggedright\arraybackslash}p{(\columnwidth - 4\tabcolsep) * \real{0.18}}
  >{\raggedright\arraybackslash}p{(\columnwidth - 4\tabcolsep) * \real{0.20}}
  >{\raggedright\arraybackslash}p{(\columnwidth - 4\tabcolsep) * \real{0.62}}@{}}
\toprule
位置 & 类型 & 含义 \\
\midrule
\endhead
getter 返回值 & \passthrough{\lstinline!float!} & 返回属性当前值。 \\
\bottomrule
\end{longtable}

\textbf{对应 C++}

\begin{itemize}
\tightlist
\item
  类：\passthrough{\lstinline!unitree\_go::msg::dds\_::UwbState\_!}
\item
  字段类型：\passthrough{\lstinline!float!}
\item
  声明位置：\passthrough{\lstinline!include/unitree/idl/go2/UwbState\_.hpp:33!}
\end{itemize}

\textbf{用法}

\begin{lstlisting}[language=Python]
current_value = obj.tag_yaw
obj.tag_yaw = new_value
\end{lstlisting}

\hypertarget{unitree-sdk2-cpp-idl-go2-uwbstate-base-roll}{%
\paragraph{\texorpdfstring{\texttt{unitree\_sdk2\_cpp.idl.go2.UwbState.base\_roll}}{unitree\_sdk2\_cpp.idl.go2.UwbState.base\_roll}}\label{unitree-sdk2-cpp-idl-go2-uwbstate-base-roll}}

底层 \passthrough{\lstinline!unitree\_go::msg::dds\_::UwbState\_!} 的
\passthrough{\lstinline!base\_roll!} 字段。类型签名只说明 Python
表示；单位、数值范围、数组固定长度和枚举含义需查对应 IDL 头文件。
底层为浮点数；类型本身不说明单位、坐标系或业务安全范围。

\textbf{签名}

\begin{lstlisting}[language=Python]
@property
def base_roll(self) -> float

@base_roll.setter
def base_roll(self, value: float) -> None
\end{lstlisting}

\textbf{可用性与安全性}

\textbf{\passthrough{\lstinline!AVAILABLE!} /
\passthrough{\lstinline!VALUE\_TYPE!}}：当前绑定源码已实现该消息值操作。

\textbf{参数}

\begin{longtable}[]{@{}
  >{\raggedright\arraybackslash}p{(\columnwidth - 4\tabcolsep) * \real{0.18}}
  >{\raggedright\arraybackslash}p{(\columnwidth - 4\tabcolsep) * \real{0.20}}
  >{\raggedright\arraybackslash}p{(\columnwidth - 4\tabcolsep) * \real{0.62}}@{}}
\toprule
名称 & Python 类型 & 含义 \\
\midrule
\endhead
\passthrough{\lstinline!value!} & \passthrough{\lstinline!float!} &
要写入或传递的新值，必须符合该参数的 Python 类型以及底层 C++
范围约束。 \\
\bottomrule
\end{longtable}

\textbf{返回值}

\begin{longtable}[]{@{}
  >{\raggedright\arraybackslash}p{(\columnwidth - 4\tabcolsep) * \real{0.18}}
  >{\raggedright\arraybackslash}p{(\columnwidth - 4\tabcolsep) * \real{0.20}}
  >{\raggedright\arraybackslash}p{(\columnwidth - 4\tabcolsep) * \real{0.62}}@{}}
\toprule
位置 & 类型 & 含义 \\
\midrule
\endhead
getter 返回值 & \passthrough{\lstinline!float!} & 返回属性当前值。 \\
\bottomrule
\end{longtable}

\textbf{对应 C++}

\begin{itemize}
\tightlist
\item
  类：\passthrough{\lstinline!unitree\_go::msg::dds\_::UwbState\_!}
\item
  字段类型：\passthrough{\lstinline!float!}
\item
  声明位置：\passthrough{\lstinline!include/unitree/idl/go2/UwbState\_.hpp:34!}
\end{itemize}

\textbf{用法}

\begin{lstlisting}[language=Python]
current_value = obj.base_roll
obj.base_roll = new_value
\end{lstlisting}

\hypertarget{unitree-sdk2-cpp-idl-go2-uwbstate-base-pitch}{%
\paragraph{\texorpdfstring{\texttt{unitree\_sdk2\_cpp.idl.go2.UwbState.base\_pitch}}{unitree\_sdk2\_cpp.idl.go2.UwbState.base\_pitch}}\label{unitree-sdk2-cpp-idl-go2-uwbstate-base-pitch}}

底层 \passthrough{\lstinline!unitree\_go::msg::dds\_::UwbState\_!} 的
\passthrough{\lstinline!base\_pitch!} 字段。类型签名只说明 Python
表示；单位、数值范围、数组固定长度和枚举含义需查对应 IDL 头文件。
底层为浮点数；类型本身不说明单位、坐标系或业务安全范围。

\textbf{签名}

\begin{lstlisting}[language=Python]
@property
def base_pitch(self) -> float

@base_pitch.setter
def base_pitch(self, value: float) -> None
\end{lstlisting}

\textbf{可用性与安全性}

\textbf{\passthrough{\lstinline!AVAILABLE!} /
\passthrough{\lstinline!VALUE\_TYPE!}}：当前绑定源码已实现该消息值操作。

\textbf{参数}

\begin{longtable}[]{@{}
  >{\raggedright\arraybackslash}p{(\columnwidth - 4\tabcolsep) * \real{0.18}}
  >{\raggedright\arraybackslash}p{(\columnwidth - 4\tabcolsep) * \real{0.20}}
  >{\raggedright\arraybackslash}p{(\columnwidth - 4\tabcolsep) * \real{0.62}}@{}}
\toprule
名称 & Python 类型 & 含义 \\
\midrule
\endhead
\passthrough{\lstinline!value!} & \passthrough{\lstinline!float!} &
要写入或传递的新值，必须符合该参数的 Python 类型以及底层 C++
范围约束。 \\
\bottomrule
\end{longtable}

\textbf{返回值}

\begin{longtable}[]{@{}
  >{\raggedright\arraybackslash}p{(\columnwidth - 4\tabcolsep) * \real{0.18}}
  >{\raggedright\arraybackslash}p{(\columnwidth - 4\tabcolsep) * \real{0.20}}
  >{\raggedright\arraybackslash}p{(\columnwidth - 4\tabcolsep) * \real{0.62}}@{}}
\toprule
位置 & 类型 & 含义 \\
\midrule
\endhead
getter 返回值 & \passthrough{\lstinline!float!} & 返回属性当前值。 \\
\bottomrule
\end{longtable}

\textbf{对应 C++}

\begin{itemize}
\tightlist
\item
  类：\passthrough{\lstinline!unitree\_go::msg::dds\_::UwbState\_!}
\item
  字段类型：\passthrough{\lstinline!float!}
\item
  声明位置：\passthrough{\lstinline!include/unitree/idl/go2/UwbState\_.hpp:35!}
\end{itemize}

\textbf{用法}

\begin{lstlisting}[language=Python]
current_value = obj.base_pitch
obj.base_pitch = new_value
\end{lstlisting}

\hypertarget{unitree-sdk2-cpp-idl-go2-uwbstate-base-yaw}{%
\paragraph{\texorpdfstring{\texttt{unitree\_sdk2\_cpp.idl.go2.UwbState.base\_yaw}}{unitree\_sdk2\_cpp.idl.go2.UwbState.base\_yaw}}\label{unitree-sdk2-cpp-idl-go2-uwbstate-base-yaw}}

底层 \passthrough{\lstinline!unitree\_go::msg::dds\_::UwbState\_!} 的
\passthrough{\lstinline!base\_yaw!} 字段。类型签名只说明 Python
表示；单位、数值范围、数组固定长度和枚举含义需查对应 IDL 头文件。
底层为浮点数；类型本身不说明单位、坐标系或业务安全范围。

\textbf{签名}

\begin{lstlisting}[language=Python]
@property
def base_yaw(self) -> float

@base_yaw.setter
def base_yaw(self, value: float) -> None
\end{lstlisting}

\textbf{可用性与安全性}

\textbf{\passthrough{\lstinline!AVAILABLE!} /
\passthrough{\lstinline!VALUE\_TYPE!}}：当前绑定源码已实现该消息值操作。

\textbf{参数}

\begin{longtable}[]{@{}
  >{\raggedright\arraybackslash}p{(\columnwidth - 4\tabcolsep) * \real{0.18}}
  >{\raggedright\arraybackslash}p{(\columnwidth - 4\tabcolsep) * \real{0.20}}
  >{\raggedright\arraybackslash}p{(\columnwidth - 4\tabcolsep) * \real{0.62}}@{}}
\toprule
名称 & Python 类型 & 含义 \\
\midrule
\endhead
\passthrough{\lstinline!value!} & \passthrough{\lstinline!float!} &
要写入或传递的新值，必须符合该参数的 Python 类型以及底层 C++
范围约束。 \\
\bottomrule
\end{longtable}

\textbf{返回值}

\begin{longtable}[]{@{}
  >{\raggedright\arraybackslash}p{(\columnwidth - 4\tabcolsep) * \real{0.18}}
  >{\raggedright\arraybackslash}p{(\columnwidth - 4\tabcolsep) * \real{0.20}}
  >{\raggedright\arraybackslash}p{(\columnwidth - 4\tabcolsep) * \real{0.62}}@{}}
\toprule
位置 & 类型 & 含义 \\
\midrule
\endhead
getter 返回值 & \passthrough{\lstinline!float!} & 返回属性当前值。 \\
\bottomrule
\end{longtable}

\textbf{对应 C++}

\begin{itemize}
\tightlist
\item
  类：\passthrough{\lstinline!unitree\_go::msg::dds\_::UwbState\_!}
\item
  字段类型：\passthrough{\lstinline!float!}
\item
  声明位置：\passthrough{\lstinline!include/unitree/idl/go2/UwbState\_.hpp:36!}
\end{itemize}

\textbf{用法}

\begin{lstlisting}[language=Python]
current_value = obj.base_yaw
obj.base_yaw = new_value
\end{lstlisting}

\hypertarget{unitree-sdk2-cpp-idl-go2-uwbstate-joystick}{%
\paragraph{\texorpdfstring{\texttt{unitree\_sdk2\_cpp.idl.go2.UwbState.joystick}}{unitree\_sdk2\_cpp.idl.go2.UwbState.joystick}}\label{unitree-sdk2-cpp-idl-go2-uwbstate-joystick}}

底层 \passthrough{\lstinline!unitree\_go::msg::dds\_::UwbState\_!} 的
\passthrough{\lstinline!joystick!} 字段。类型签名只说明 Python
表示；单位、数值范围、数组固定长度和枚举含义需查对应 IDL 头文件。
底层是固定长度数组，写入序列必须正好包含 2
个元素；元素约束：底层为浮点数；类型本身不说明单位、坐标系或业务安全范围。

\textbf{签名}

\begin{lstlisting}[language=Python]
@property
def joystick(self) -> list[float]

@joystick.setter
def joystick(self, value: Sequence[float]) -> None
\end{lstlisting}

\textbf{可用性与安全性}

\textbf{\passthrough{\lstinline!AVAILABLE!} /
\passthrough{\lstinline!VALUE\_TYPE!}}：当前绑定源码已实现该消息值操作。

\textbf{参数}

\begin{longtable}[]{@{}
  >{\raggedright\arraybackslash}p{(\columnwidth - 4\tabcolsep) * \real{0.18}}
  >{\raggedright\arraybackslash}p{(\columnwidth - 4\tabcolsep) * \real{0.20}}
  >{\raggedright\arraybackslash}p{(\columnwidth - 4\tabcolsep) * \real{0.62}}@{}}
\toprule
名称 & Python 类型 & 含义 \\
\midrule
\endhead
\passthrough{\lstinline!value!} &
\passthrough{\lstinline!Sequence[float]!} &
要写入或传递的新值，必须符合该参数的 Python 类型以及底层 C++
范围约束。 \\
\bottomrule
\end{longtable}

\textbf{返回值}

\begin{longtable}[]{@{}
  >{\raggedright\arraybackslash}p{(\columnwidth - 4\tabcolsep) * \real{0.18}}
  >{\raggedright\arraybackslash}p{(\columnwidth - 4\tabcolsep) * \real{0.20}}
  >{\raggedright\arraybackslash}p{(\columnwidth - 4\tabcolsep) * \real{0.62}}@{}}
\toprule
位置 & 类型 & 含义 \\
\midrule
\endhead
getter 返回值 & \passthrough{\lstinline!list[float]!} &
返回属性当前值。数组、vector 和嵌套消息采用复制语义。 \\
\bottomrule
\end{longtable}

\textbf{对应 C++}

\begin{itemize}
\tightlist
\item
  类：\passthrough{\lstinline!unitree\_go::msg::dds\_::UwbState\_!}
\item
  字段类型：\passthrough{\lstinline!std::array<float, 2>!}
\item
  声明位置：\passthrough{\lstinline!include/unitree/idl/go2/UwbState\_.hpp:37!}
\end{itemize}

\textbf{用法}

\begin{lstlisting}[language=Python]
current_value = obj.joystick
obj.joystick = new_value
\end{lstlisting}

\begin{quote}
{[}!TIP{]} 读取容器或嵌套值后应遵循``读取、修改、重新赋值''。只修改
getter 返回的副本不会可靠地更新原 C++ 消息。
\end{quote}

\hypertarget{unitree-sdk2-cpp-idl-go2-uwbstate-error-state}{%
\paragraph{\texorpdfstring{\texttt{unitree\_sdk2\_cpp.idl.go2.UwbState.error\_state}}{unitree\_sdk2\_cpp.idl.go2.UwbState.error\_state}}\label{unitree-sdk2-cpp-idl-go2-uwbstate-error-state}}

底层 \passthrough{\lstinline!unitree\_go::msg::dds\_::UwbState\_!} 的
\passthrough{\lstinline!error\_state!} 字段。类型签名只说明 Python
表示；单位、数值范围、数组固定长度和枚举含义需查对应 IDL 头文件。
取值范围为 0 到 255。

\textbf{签名}

\begin{lstlisting}[language=Python]
@property
def error_state(self) -> int

@error_state.setter
def error_state(self, value: int) -> None
\end{lstlisting}

\textbf{可用性与安全性}

\textbf{\passthrough{\lstinline!AVAILABLE!} /
\passthrough{\lstinline!VALUE\_TYPE!}}：当前绑定源码已实现该消息值操作。

\textbf{参数}

\begin{longtable}[]{@{}
  >{\raggedright\arraybackslash}p{(\columnwidth - 4\tabcolsep) * \real{0.18}}
  >{\raggedright\arraybackslash}p{(\columnwidth - 4\tabcolsep) * \real{0.20}}
  >{\raggedright\arraybackslash}p{(\columnwidth - 4\tabcolsep) * \real{0.62}}@{}}
\toprule
名称 & Python 类型 & 含义 \\
\midrule
\endhead
\passthrough{\lstinline!value!} & \passthrough{\lstinline!int!} &
要写入或传递的新值，必须符合该参数的 Python 类型以及底层 C++
范围约束。 \\
\bottomrule
\end{longtable}

\textbf{返回值}

\begin{longtable}[]{@{}
  >{\raggedright\arraybackslash}p{(\columnwidth - 4\tabcolsep) * \real{0.18}}
  >{\raggedright\arraybackslash}p{(\columnwidth - 4\tabcolsep) * \real{0.20}}
  >{\raggedright\arraybackslash}p{(\columnwidth - 4\tabcolsep) * \real{0.62}}@{}}
\toprule
位置 & 类型 & 含义 \\
\midrule
\endhead
getter 返回值 & \passthrough{\lstinline!int!} & 返回属性当前值。 \\
\bottomrule
\end{longtable}

\textbf{对应 C++}

\begin{itemize}
\tightlist
\item
  类：\passthrough{\lstinline!unitree\_go::msg::dds\_::UwbState\_!}
\item
  字段类型：\passthrough{\lstinline!uint8\_t!}
\item
  声明位置：\passthrough{\lstinline!include/unitree/idl/go2/UwbState\_.hpp:38!}
\end{itemize}

\textbf{用法}

\begin{lstlisting}[language=Python]
current_value = obj.error_state
obj.error_state = new_value
\end{lstlisting}

\hypertarget{unitree-sdk2-cpp-idl-go2-uwbstate-buttons}{%
\paragraph{\texorpdfstring{\texttt{unitree\_sdk2\_cpp.idl.go2.UwbState.buttons}}{unitree\_sdk2\_cpp.idl.go2.UwbState.buttons}}\label{unitree-sdk2-cpp-idl-go2-uwbstate-buttons}}

底层 \passthrough{\lstinline!unitree\_go::msg::dds\_::UwbState\_!} 的
\passthrough{\lstinline!buttons!} 字段。类型签名只说明 Python
表示；单位、数值范围、数组固定长度和枚举含义需查对应 IDL 头文件。
取值范围为 0 到 255。

\textbf{签名}

\begin{lstlisting}[language=Python]
@property
def buttons(self) -> int

@buttons.setter
def buttons(self, value: int) -> None
\end{lstlisting}

\textbf{可用性与安全性}

\textbf{\passthrough{\lstinline!AVAILABLE!} /
\passthrough{\lstinline!VALUE\_TYPE!}}：当前绑定源码已实现该消息值操作。

\textbf{参数}

\begin{longtable}[]{@{}
  >{\raggedright\arraybackslash}p{(\columnwidth - 4\tabcolsep) * \real{0.18}}
  >{\raggedright\arraybackslash}p{(\columnwidth - 4\tabcolsep) * \real{0.20}}
  >{\raggedright\arraybackslash}p{(\columnwidth - 4\tabcolsep) * \real{0.62}}@{}}
\toprule
名称 & Python 类型 & 含义 \\
\midrule
\endhead
\passthrough{\lstinline!value!} & \passthrough{\lstinline!int!} &
要写入或传递的新值，必须符合该参数的 Python 类型以及底层 C++
范围约束。 \\
\bottomrule
\end{longtable}

\textbf{返回值}

\begin{longtable}[]{@{}
  >{\raggedright\arraybackslash}p{(\columnwidth - 4\tabcolsep) * \real{0.18}}
  >{\raggedright\arraybackslash}p{(\columnwidth - 4\tabcolsep) * \real{0.20}}
  >{\raggedright\arraybackslash}p{(\columnwidth - 4\tabcolsep) * \real{0.62}}@{}}
\toprule
位置 & 类型 & 含义 \\
\midrule
\endhead
getter 返回值 & \passthrough{\lstinline!int!} & 返回属性当前值。 \\
\bottomrule
\end{longtable}

\textbf{对应 C++}

\begin{itemize}
\tightlist
\item
  类：\passthrough{\lstinline!unitree\_go::msg::dds\_::UwbState\_!}
\item
  字段类型：\passthrough{\lstinline!uint8\_t!}
\item
  声明位置：\passthrough{\lstinline!include/unitree/idl/go2/UwbState\_.hpp:39!}
\end{itemize}

\textbf{用法}

\begin{lstlisting}[language=Python]
current_value = obj.buttons
obj.buttons = new_value
\end{lstlisting}

\hypertarget{unitree-sdk2-cpp-idl-go2-uwbstate-enabled-from-app}{%
\paragraph{\texorpdfstring{\texttt{unitree\_sdk2\_cpp.idl.go2.UwbState.enabled\_from\_app}}{unitree\_sdk2\_cpp.idl.go2.UwbState.enabled\_from\_app}}\label{unitree-sdk2-cpp-idl-go2-uwbstate-enabled-from-app}}

底层 \passthrough{\lstinline!unitree\_go::msg::dds\_::UwbState\_!} 的
\passthrough{\lstinline!enabled\_from\_app!} 字段。类型签名只说明 Python
表示；单位、数值范围、数组固定长度和枚举含义需查对应 IDL 头文件。
取值范围为 0 到 255。

\textbf{签名}

\begin{lstlisting}[language=Python]
@property
def enabled_from_app(self) -> int

@enabled_from_app.setter
def enabled_from_app(self, value: int) -> None
\end{lstlisting}

\textbf{可用性与安全性}

\textbf{\passthrough{\lstinline!AVAILABLE!} /
\passthrough{\lstinline!VALUE\_TYPE!}}：当前绑定源码已实现该消息值操作。

\textbf{参数}

\begin{longtable}[]{@{}
  >{\raggedright\arraybackslash}p{(\columnwidth - 4\tabcolsep) * \real{0.18}}
  >{\raggedright\arraybackslash}p{(\columnwidth - 4\tabcolsep) * \real{0.20}}
  >{\raggedright\arraybackslash}p{(\columnwidth - 4\tabcolsep) * \real{0.62}}@{}}
\toprule
名称 & Python 类型 & 含义 \\
\midrule
\endhead
\passthrough{\lstinline!value!} & \passthrough{\lstinline!int!} &
要写入或传递的新值，必须符合该参数的 Python 类型以及底层 C++
范围约束。 \\
\bottomrule
\end{longtable}

\textbf{返回值}

\begin{longtable}[]{@{}
  >{\raggedright\arraybackslash}p{(\columnwidth - 4\tabcolsep) * \real{0.18}}
  >{\raggedright\arraybackslash}p{(\columnwidth - 4\tabcolsep) * \real{0.20}}
  >{\raggedright\arraybackslash}p{(\columnwidth - 4\tabcolsep) * \real{0.62}}@{}}
\toprule
位置 & 类型 & 含义 \\
\midrule
\endhead
getter 返回值 & \passthrough{\lstinline!int!} & 返回属性当前值。 \\
\bottomrule
\end{longtable}

\textbf{对应 C++}

\begin{itemize}
\tightlist
\item
  类：\passthrough{\lstinline!unitree\_go::msg::dds\_::UwbState\_!}
\item
  字段类型：\passthrough{\lstinline!uint8\_t!}
\item
  声明位置：\passthrough{\lstinline!include/unitree/idl/go2/UwbState\_.hpp:40!}
\end{itemize}

\textbf{用法}

\begin{lstlisting}[language=Python]
current_value = obj.enabled_from_app
obj.enabled_from_app = new_value
\end{lstlisting}

\hypertarget{unitree-sdk2-cpp-idl-go2-uwbswitch}{%
\subsubsection{\texorpdfstring{\texttt{unitree\_sdk2\_cpp.idl.go2.UwbSwitch}}{unitree\_sdk2\_cpp.idl.go2.UwbSwitch}}\label{unitree-sdk2-cpp-idl-go2-uwbswitch}}

可由 Python 读写的 DDS/IDL 消息值类型。字段采用复制语义。

\textbf{导入}

\begin{lstlisting}[language=Python]
from unitree_sdk2_cpp.idl.go2 import UwbSwitch
\end{lstlisting}

\textbf{方法索引}

\begin{longtable}[]{@{}
  >{\raggedright\arraybackslash}p{(\columnwidth - 6\tabcolsep) * \real{0.18}}
  >{\raggedright\arraybackslash}p{(\columnwidth - 6\tabcolsep) * \real{0.24}}
  >{\raggedright\arraybackslash}p{(\columnwidth - 6\tabcolsep) * \real{0.18}}
  >{\raggedright\arraybackslash}p{(\columnwidth - 6\tabcolsep) * \real{0.40}}@{}}
\toprule
方法 & Python 签名 & 状态 & 安全分类 \\
\midrule
\endhead
\protect\hyperlink{unitree-sdk2-cpp-idl-go2-uwbswitch-dunder-init-1}{\passthrough{\lstinline!\_\_init\_\_!}}
& \passthrough{\lstinline!def \_\_init\_\_(self) -> None!} &
\passthrough{\lstinline!AVAILABLE!} &
\passthrough{\lstinline!VALUE\_TYPE!} \\
\protect\hyperlink{unitree-sdk2-cpp-idl-go2-uwbswitch-dunder-eq-1}{\passthrough{\lstinline!\_\_eq\_\_!}}
& \passthrough{\lstinline!def \_\_eq\_\_(self, other: object) -> bool!}
& \passthrough{\lstinline!AVAILABLE!} &
\passthrough{\lstinline!VALUE\_TYPE!} \\
\protect\hyperlink{unitree-sdk2-cpp-idl-go2-uwbswitch-dunder-ne-1}{\passthrough{\lstinline!\_\_ne\_\_!}}
& \passthrough{\lstinline!def \_\_ne\_\_(self, other: object) -> bool!}
& \passthrough{\lstinline!AVAILABLE!} &
\passthrough{\lstinline!VALUE\_TYPE!} \\
\bottomrule
\end{longtable}

\textbf{属性索引}

\begin{longtable}[]{@{}
  >{\raggedright\arraybackslash}p{(\columnwidth - 6\tabcolsep) * \real{0.18}}
  >{\raggedright\arraybackslash}p{(\columnwidth - 6\tabcolsep) * \real{0.24}}
  >{\raggedright\arraybackslash}p{(\columnwidth - 6\tabcolsep) * \real{0.18}}
  >{\raggedright\arraybackslash}p{(\columnwidth - 6\tabcolsep) * \real{0.40}}@{}}
\toprule
属性 & 读取类型 & 写入类型 & 状态 \\
\midrule
\endhead
\protect\hyperlink{unitree-sdk2-cpp-idl-go2-uwbswitch-enabled}{\passthrough{\lstinline!enabled!}}
& \passthrough{\lstinline!int!} & \passthrough{\lstinline!int!} &
\passthrough{\lstinline!AVAILABLE!} \\
\bottomrule
\end{longtable}

\hypertarget{unitree-sdk2-cpp-idl-go2-uwbswitch-dunder-init-1}{%
\paragraph{\texorpdfstring{\texttt{unitree\_sdk2\_cpp.idl.go2.UwbSwitch.\_\_init\_\_}}{unitree\_sdk2\_cpp.idl.go2.UwbSwitch.\_\_init\_\_}}\label{unitree-sdk2-cpp-idl-go2-uwbswitch-dunder-init-1}}

初始化 \passthrough{\lstinline!UwbSwitch!}
实例。是否能实际构造取决于下方可用性状态。

\textbf{签名}

\begin{lstlisting}[language=Python]
def __init__(self) -> None
\end{lstlisting}

\textbf{可用性与安全性}

\textbf{\passthrough{\lstinline!AVAILABLE!} /
\passthrough{\lstinline!VALUE\_TYPE!}}：当前绑定源码已实现该消息值操作。

\textbf{参数}

无显式参数。实例方法中的 \passthrough{\lstinline!self!} 由 Python
自动传入。

\textbf{返回值}

\begin{longtable}[]{@{}
  >{\raggedright\arraybackslash}p{(\columnwidth - 4\tabcolsep) * \real{0.18}}
  >{\raggedright\arraybackslash}p{(\columnwidth - 4\tabcolsep) * \real{0.20}}
  >{\raggedright\arraybackslash}p{(\columnwidth - 4\tabcolsep) * \real{0.62}}@{}}
\toprule
位置 & 类型 & 含义 \\
\midrule
\endhead
无 & \passthrough{\lstinline!None!} &
\passthrough{\lstinline!\_\_init\_\_!}
本身不返回值；调用类对象时，在构造函数可用的前提下得到该类实例。 \\
\bottomrule
\end{longtable}

\textbf{对应 C++}

\begin{itemize}
\tightlist
\item
  类：\passthrough{\lstinline!unitree\_go::msg::dds\_::UwbSwitch\_!}
\item
  签名：\passthrough{\lstinline!UwbSwitch\_()!}
\item
  绑定策略：\passthrough{\lstinline!未单独标注!}
\item
  声明位置：\passthrough{\lstinline!include/unitree/idl/go2/UwbSwitch\_.hpp:26!}
\end{itemize}

\textbf{说明}

\begin{itemize}
\tightlist
\item
  示例是局部调用片段；\passthrough{\lstinline!obj!}
  和参数变量表示已经按上层业务校验并准备好的对象。
\end{itemize}

\textbf{用法}

\begin{lstlisting}[language=Python]
value = UwbSwitch()
\end{lstlisting}

\hypertarget{unitree-sdk2-cpp-idl-go2-uwbswitch-dunder-eq-1}{%
\paragraph{\texorpdfstring{\texttt{unitree\_sdk2\_cpp.idl.go2.UwbSwitch.\_\_eq\_\_}}{unitree\_sdk2\_cpp.idl.go2.UwbSwitch.\_\_eq\_\_}}\label{unitree-sdk2-cpp-idl-go2-uwbswitch-dunder-eq-1}}

按底层 IDL 消息内容比较两个对象是否相等。

\textbf{签名}

\begin{lstlisting}[language=Python]
def __eq__(self, other: object) -> bool
\end{lstlisting}

\textbf{可用性与安全性}

\textbf{\passthrough{\lstinline!AVAILABLE!} /
\passthrough{\lstinline!VALUE\_TYPE!}}：当前绑定源码已实现该消息值操作。

\textbf{参数}

\begin{longtable}[]{@{}
  >{\raggedright\arraybackslash}p{(\columnwidth - 6\tabcolsep) * \real{0.18}}
  >{\raggedright\arraybackslash}p{(\columnwidth - 6\tabcolsep) * \real{0.24}}
  >{\raggedright\arraybackslash}p{(\columnwidth - 6\tabcolsep) * \real{0.18}}
  >{\raggedright\arraybackslash}p{(\columnwidth - 6\tabcolsep) * \real{0.40}}@{}}
\toprule
名称 & Python 类型 & 默认值 & 含义 \\
\midrule
\endhead
\passthrough{\lstinline!other!} & \passthrough{\lstinline!object!} &
必填 & 用于比较的另一个 Python 对象。类型不兼容时比较结果通常为
\passthrough{\lstinline!False!}。 \\
\bottomrule
\end{longtable}

\textbf{返回值}

\begin{longtable}[]{@{}
  >{\raggedright\arraybackslash}p{(\columnwidth - 4\tabcolsep) * \real{0.18}}
  >{\raggedright\arraybackslash}p{(\columnwidth - 4\tabcolsep) * \real{0.20}}
  >{\raggedright\arraybackslash}p{(\columnwidth - 4\tabcolsep) * \real{0.62}}@{}}
\toprule
位置 & 类型 & 含义 \\
\midrule
\endhead
返回值 & \passthrough{\lstinline!bool!} & 返回
\passthrough{\lstinline!bool!}。更精确的业务含义见该方法用途、C++
签名和目标型号协议。 \\
\bottomrule
\end{longtable}

\textbf{对应 C++}

\begin{itemize}
\tightlist
\item
  类：\passthrough{\lstinline!unitree\_go::msg::dds\_::UwbSwitch\_!}
\item
  签名：\passthrough{\lstinline!operator==(const value \&) const!}
\item
  绑定策略：\passthrough{\lstinline!未单独标注!}
\item
  声明位置：由 pybind11 手工绑定或生成注册表提供；manifest
  未记录独立头文件行号。
\end{itemize}

\textbf{说明}

\begin{itemize}
\tightlist
\item
  示例是局部调用片段；\passthrough{\lstinline!obj!}
  和参数变量表示已经按上层业务校验并准备好的对象。
\end{itemize}

\textbf{用法}

\begin{lstlisting}[language=Python]
same = left == right
\end{lstlisting}

\hypertarget{unitree-sdk2-cpp-idl-go2-uwbswitch-dunder-ne-1}{%
\paragraph{\texorpdfstring{\texttt{unitree\_sdk2\_cpp.idl.go2.UwbSwitch.\_\_ne\_\_}}{unitree\_sdk2\_cpp.idl.go2.UwbSwitch.\_\_ne\_\_}}\label{unitree-sdk2-cpp-idl-go2-uwbswitch-dunder-ne-1}}

按底层 IDL 消息内容比较两个对象是否不相等。

\textbf{签名}

\begin{lstlisting}[language=Python]
def __ne__(self, other: object) -> bool
\end{lstlisting}

\textbf{可用性与安全性}

\textbf{\passthrough{\lstinline!AVAILABLE!} /
\passthrough{\lstinline!VALUE\_TYPE!}}：当前绑定源码已实现该消息值操作。

\textbf{参数}

\begin{longtable}[]{@{}
  >{\raggedright\arraybackslash}p{(\columnwidth - 6\tabcolsep) * \real{0.18}}
  >{\raggedright\arraybackslash}p{(\columnwidth - 6\tabcolsep) * \real{0.24}}
  >{\raggedright\arraybackslash}p{(\columnwidth - 6\tabcolsep) * \real{0.18}}
  >{\raggedright\arraybackslash}p{(\columnwidth - 6\tabcolsep) * \real{0.40}}@{}}
\toprule
名称 & Python 类型 & 默认值 & 含义 \\
\midrule
\endhead
\passthrough{\lstinline!other!} & \passthrough{\lstinline!object!} &
必填 & 用于比较的另一个 Python 对象。类型不兼容时比较结果通常为
\passthrough{\lstinline!False!}。 \\
\bottomrule
\end{longtable}

\textbf{返回值}

\begin{longtable}[]{@{}
  >{\raggedright\arraybackslash}p{(\columnwidth - 4\tabcolsep) * \real{0.18}}
  >{\raggedright\arraybackslash}p{(\columnwidth - 4\tabcolsep) * \real{0.20}}
  >{\raggedright\arraybackslash}p{(\columnwidth - 4\tabcolsep) * \real{0.62}}@{}}
\toprule
位置 & 类型 & 含义 \\
\midrule
\endhead
返回值 & \passthrough{\lstinline!bool!} & 返回
\passthrough{\lstinline!bool!}。更精确的业务含义见该方法用途、C++
签名和目标型号协议。 \\
\bottomrule
\end{longtable}

\textbf{对应 C++}

\begin{itemize}
\tightlist
\item
  类：\passthrough{\lstinline!unitree\_go::msg::dds\_::UwbSwitch\_!}
\item
  签名：\passthrough{\lstinline"operator!=(const value \&) const"}
\item
  绑定策略：\passthrough{\lstinline!未单独标注!}
\item
  声明位置：由 pybind11 手工绑定或生成注册表提供；manifest
  未记录独立头文件行号。
\end{itemize}

\textbf{说明}

\begin{itemize}
\tightlist
\item
  示例是局部调用片段；\passthrough{\lstinline!obj!}
  和参数变量表示已经按上层业务校验并准备好的对象。
\end{itemize}

\textbf{用法}

\begin{lstlisting}[language=Python]
different = left != right
\end{lstlisting}

\hypertarget{unitree-sdk2-cpp-idl-go2-uwbswitch-enabled}{%
\paragraph{\texorpdfstring{\texttt{unitree\_sdk2\_cpp.idl.go2.UwbSwitch.enabled}}{unitree\_sdk2\_cpp.idl.go2.UwbSwitch.enabled}}\label{unitree-sdk2-cpp-idl-go2-uwbswitch-enabled}}

底层 \passthrough{\lstinline!unitree\_go::msg::dds\_::UwbSwitch\_!} 的
\passthrough{\lstinline!enabled!} 字段。类型签名只说明 Python
表示；单位、数值范围、数组固定长度和枚举含义需查对应 IDL 头文件。
取值范围为 0 到 255。

\textbf{签名}

\begin{lstlisting}[language=Python]
@property
def enabled(self) -> int

@enabled.setter
def enabled(self, value: int) -> None
\end{lstlisting}

\textbf{可用性与安全性}

\textbf{\passthrough{\lstinline!AVAILABLE!} /
\passthrough{\lstinline!VALUE\_TYPE!}}：当前绑定源码已实现该消息值操作。

\textbf{参数}

\begin{longtable}[]{@{}
  >{\raggedright\arraybackslash}p{(\columnwidth - 4\tabcolsep) * \real{0.18}}
  >{\raggedright\arraybackslash}p{(\columnwidth - 4\tabcolsep) * \real{0.20}}
  >{\raggedright\arraybackslash}p{(\columnwidth - 4\tabcolsep) * \real{0.62}}@{}}
\toprule
名称 & Python 类型 & 含义 \\
\midrule
\endhead
\passthrough{\lstinline!value!} & \passthrough{\lstinline!int!} &
要写入或传递的新值，必须符合该参数的 Python 类型以及底层 C++
范围约束。 \\
\bottomrule
\end{longtable}

\textbf{返回值}

\begin{longtable}[]{@{}
  >{\raggedright\arraybackslash}p{(\columnwidth - 4\tabcolsep) * \real{0.18}}
  >{\raggedright\arraybackslash}p{(\columnwidth - 4\tabcolsep) * \real{0.20}}
  >{\raggedright\arraybackslash}p{(\columnwidth - 4\tabcolsep) * \real{0.62}}@{}}
\toprule
位置 & 类型 & 含义 \\
\midrule
\endhead
getter 返回值 & \passthrough{\lstinline!int!} & 返回属性当前值。 \\
\bottomrule
\end{longtable}

\textbf{对应 C++}

\begin{itemize}
\tightlist
\item
  类：\passthrough{\lstinline!unitree\_go::msg::dds\_::UwbSwitch\_!}
\item
  字段类型：\passthrough{\lstinline!uint8\_t!}
\item
  声明位置：\passthrough{\lstinline!include/unitree/idl/go2/UwbSwitch\_.hpp:23!}
\end{itemize}

\textbf{用法}

\begin{lstlisting}[language=Python]
current_value = obj.enabled
obj.enabled = new_value
\end{lstlisting}

\hypertarget{unitree-sdk2-cpp-idl-go2-voxelmapcompressed}{%
\subsubsection{\texorpdfstring{\texttt{unitree\_sdk2\_cpp.idl.go2.VoxelMapCompressed}}{unitree\_sdk2\_cpp.idl.go2.VoxelMapCompressed}}\label{unitree-sdk2-cpp-idl-go2-voxelmapcompressed}}

可由 Python 读写的 DDS/IDL 消息值类型。字段采用复制语义。

\textbf{导入}

\begin{lstlisting}[language=Python]
from unitree_sdk2_cpp.idl.go2 import VoxelMapCompressed
\end{lstlisting}

\textbf{方法索引}

\begin{longtable}[]{@{}
  >{\raggedright\arraybackslash}p{(\columnwidth - 6\tabcolsep) * \real{0.18}}
  >{\raggedright\arraybackslash}p{(\columnwidth - 6\tabcolsep) * \real{0.24}}
  >{\raggedright\arraybackslash}p{(\columnwidth - 6\tabcolsep) * \real{0.18}}
  >{\raggedright\arraybackslash}p{(\columnwidth - 6\tabcolsep) * \real{0.40}}@{}}
\toprule
方法 & Python 签名 & 状态 & 安全分类 \\
\midrule
\endhead
\protect\hyperlink{unitree-sdk2-cpp-idl-go2-voxelmapcompressed-dunder-init-1}{\passthrough{\lstinline!\_\_init\_\_!}}
& \passthrough{\lstinline!def \_\_init\_\_(self) -> None!} &
\passthrough{\lstinline!AVAILABLE!} &
\passthrough{\lstinline!VALUE\_TYPE!} \\
\protect\hyperlink{unitree-sdk2-cpp-idl-go2-voxelmapcompressed-dunder-eq-1}{\passthrough{\lstinline!\_\_eq\_\_!}}
& \passthrough{\lstinline!def \_\_eq\_\_(self, other: object) -> bool!}
& \passthrough{\lstinline!AVAILABLE!} &
\passthrough{\lstinline!VALUE\_TYPE!} \\
\protect\hyperlink{unitree-sdk2-cpp-idl-go2-voxelmapcompressed-dunder-ne-1}{\passthrough{\lstinline!\_\_ne\_\_!}}
& \passthrough{\lstinline!def \_\_ne\_\_(self, other: object) -> bool!}
& \passthrough{\lstinline!AVAILABLE!} &
\passthrough{\lstinline!VALUE\_TYPE!} \\
\bottomrule
\end{longtable}

\textbf{属性索引}

\begin{longtable}[]{@{}
  >{\raggedright\arraybackslash}p{(\columnwidth - 6\tabcolsep) * \real{0.18}}
  >{\raggedright\arraybackslash}p{(\columnwidth - 6\tabcolsep) * \real{0.24}}
  >{\raggedright\arraybackslash}p{(\columnwidth - 6\tabcolsep) * \real{0.18}}
  >{\raggedright\arraybackslash}p{(\columnwidth - 6\tabcolsep) * \real{0.40}}@{}}
\toprule
属性 & 读取类型 & 写入类型 & 状态 \\
\midrule
\endhead
\protect\hyperlink{unitree-sdk2-cpp-idl-go2-voxelmapcompressed-stamp}{\passthrough{\lstinline!stamp!}}
& \passthrough{\lstinline!float!} & \passthrough{\lstinline!float!} &
\passthrough{\lstinline!AVAILABLE!} \\
\protect\hyperlink{unitree-sdk2-cpp-idl-go2-voxelmapcompressed-frame-id}{\passthrough{\lstinline!frame\_id!}}
& \passthrough{\lstinline!str!} & \passthrough{\lstinline!str!} &
\passthrough{\lstinline!AVAILABLE!} \\
\protect\hyperlink{unitree-sdk2-cpp-idl-go2-voxelmapcompressed-resolution}{\passthrough{\lstinline!resolution!}}
& \passthrough{\lstinline!float!} & \passthrough{\lstinline!float!} &
\passthrough{\lstinline!AVAILABLE!} \\
\protect\hyperlink{unitree-sdk2-cpp-idl-go2-voxelmapcompressed-origin}{\passthrough{\lstinline!origin!}}
& \passthrough{\lstinline!list[float]!} &
\passthrough{\lstinline!Sequence[float]!} &
\passthrough{\lstinline!AVAILABLE!} \\
\protect\hyperlink{unitree-sdk2-cpp-idl-go2-voxelmapcompressed-width}{\passthrough{\lstinline!width!}}
& \passthrough{\lstinline!list[int]!} &
\passthrough{\lstinline!Sequence[int]!} &
\passthrough{\lstinline!AVAILABLE!} \\
\protect\hyperlink{unitree-sdk2-cpp-idl-go2-voxelmapcompressed-src-size}{\passthrough{\lstinline!src\_size!}}
& \passthrough{\lstinline!int!} & \passthrough{\lstinline!int!} &
\passthrough{\lstinline!AVAILABLE!} \\
\protect\hyperlink{unitree-sdk2-cpp-idl-go2-voxelmapcompressed-data}{\passthrough{\lstinline!data!}}
& \passthrough{\lstinline!list[int]!} &
\passthrough{\lstinline!Sequence[int]!} &
\passthrough{\lstinline!AVAILABLE!} \\
\bottomrule
\end{longtable}

\hypertarget{unitree-sdk2-cpp-idl-go2-voxelmapcompressed-dunder-init-1}{%
\paragraph{\texorpdfstring{\texttt{unitree\_sdk2\_cpp.idl.go2.VoxelMapCompressed.\_\_init\_\_}}{unitree\_sdk2\_cpp.idl.go2.VoxelMapCompressed.\_\_init\_\_}}\label{unitree-sdk2-cpp-idl-go2-voxelmapcompressed-dunder-init-1}}

初始化 \passthrough{\lstinline!VoxelMapCompressed!}
实例。是否能实际构造取决于下方可用性状态。

\textbf{签名}

\begin{lstlisting}[language=Python]
def __init__(self) -> None
\end{lstlisting}

\textbf{可用性与安全性}

\textbf{\passthrough{\lstinline!AVAILABLE!} /
\passthrough{\lstinline!VALUE\_TYPE!}}：当前绑定源码已实现该消息值操作。

\textbf{参数}

无显式参数。实例方法中的 \passthrough{\lstinline!self!} 由 Python
自动传入。

\textbf{返回值}

\begin{longtable}[]{@{}
  >{\raggedright\arraybackslash}p{(\columnwidth - 4\tabcolsep) * \real{0.18}}
  >{\raggedright\arraybackslash}p{(\columnwidth - 4\tabcolsep) * \real{0.20}}
  >{\raggedright\arraybackslash}p{(\columnwidth - 4\tabcolsep) * \real{0.62}}@{}}
\toprule
位置 & 类型 & 含义 \\
\midrule
\endhead
无 & \passthrough{\lstinline!None!} &
\passthrough{\lstinline!\_\_init\_\_!}
本身不返回值；调用类对象时，在构造函数可用的前提下得到该类实例。 \\
\bottomrule
\end{longtable}

\textbf{对应 C++}

\begin{itemize}
\tightlist
\item
  类：\passthrough{\lstinline!unitree\_go::msg::dds\_::VoxelMapCompressed\_!}
\item
  签名：\passthrough{\lstinline!VoxelMapCompressed\_()!}
\item
  绑定策略：\passthrough{\lstinline!未单独标注!}
\item
  声明位置：\passthrough{\lstinline!include/unitree/idl/go2/VoxelMapCompressed\_.hpp:35!}
\end{itemize}

\textbf{说明}

\begin{itemize}
\tightlist
\item
  示例是局部调用片段；\passthrough{\lstinline!obj!}
  和参数变量表示已经按上层业务校验并准备好的对象。
\end{itemize}

\textbf{用法}

\begin{lstlisting}[language=Python]
value = VoxelMapCompressed()
\end{lstlisting}

\hypertarget{unitree-sdk2-cpp-idl-go2-voxelmapcompressed-dunder-eq-1}{%
\paragraph{\texorpdfstring{\texttt{unitree\_sdk2\_cpp.idl.go2.VoxelMapCompressed.\_\_eq\_\_}}{unitree\_sdk2\_cpp.idl.go2.VoxelMapCompressed.\_\_eq\_\_}}\label{unitree-sdk2-cpp-idl-go2-voxelmapcompressed-dunder-eq-1}}

按底层 IDL 消息内容比较两个对象是否相等。

\textbf{签名}

\begin{lstlisting}[language=Python]
def __eq__(self, other: object) -> bool
\end{lstlisting}

\textbf{可用性与安全性}

\textbf{\passthrough{\lstinline!AVAILABLE!} /
\passthrough{\lstinline!VALUE\_TYPE!}}：当前绑定源码已实现该消息值操作。

\textbf{参数}

\begin{longtable}[]{@{}
  >{\raggedright\arraybackslash}p{(\columnwidth - 6\tabcolsep) * \real{0.18}}
  >{\raggedright\arraybackslash}p{(\columnwidth - 6\tabcolsep) * \real{0.24}}
  >{\raggedright\arraybackslash}p{(\columnwidth - 6\tabcolsep) * \real{0.18}}
  >{\raggedright\arraybackslash}p{(\columnwidth - 6\tabcolsep) * \real{0.40}}@{}}
\toprule
名称 & Python 类型 & 默认值 & 含义 \\
\midrule
\endhead
\passthrough{\lstinline!other!} & \passthrough{\lstinline!object!} &
必填 & 用于比较的另一个 Python 对象。类型不兼容时比较结果通常为
\passthrough{\lstinline!False!}。 \\
\bottomrule
\end{longtable}

\textbf{返回值}

\begin{longtable}[]{@{}
  >{\raggedright\arraybackslash}p{(\columnwidth - 4\tabcolsep) * \real{0.18}}
  >{\raggedright\arraybackslash}p{(\columnwidth - 4\tabcolsep) * \real{0.20}}
  >{\raggedright\arraybackslash}p{(\columnwidth - 4\tabcolsep) * \real{0.62}}@{}}
\toprule
位置 & 类型 & 含义 \\
\midrule
\endhead
返回值 & \passthrough{\lstinline!bool!} & 返回
\passthrough{\lstinline!bool!}。更精确的业务含义见该方法用途、C++
签名和目标型号协议。 \\
\bottomrule
\end{longtable}

\textbf{对应 C++}

\begin{itemize}
\tightlist
\item
  类：\passthrough{\lstinline!unitree\_go::msg::dds\_::VoxelMapCompressed\_!}
\item
  签名：\passthrough{\lstinline!operator==(const value \&) const!}
\item
  绑定策略：\passthrough{\lstinline!未单独标注!}
\item
  声明位置：由 pybind11 手工绑定或生成注册表提供；manifest
  未记录独立头文件行号。
\end{itemize}

\textbf{说明}

\begin{itemize}
\tightlist
\item
  示例是局部调用片段；\passthrough{\lstinline!obj!}
  和参数变量表示已经按上层业务校验并准备好的对象。
\end{itemize}

\textbf{用法}

\begin{lstlisting}[language=Python]
same = left == right
\end{lstlisting}

\hypertarget{unitree-sdk2-cpp-idl-go2-voxelmapcompressed-dunder-ne-1}{%
\paragraph{\texorpdfstring{\texttt{unitree\_sdk2\_cpp.idl.go2.VoxelMapCompressed.\_\_ne\_\_}}{unitree\_sdk2\_cpp.idl.go2.VoxelMapCompressed.\_\_ne\_\_}}\label{unitree-sdk2-cpp-idl-go2-voxelmapcompressed-dunder-ne-1}}

按底层 IDL 消息内容比较两个对象是否不相等。

\textbf{签名}

\begin{lstlisting}[language=Python]
def __ne__(self, other: object) -> bool
\end{lstlisting}

\textbf{可用性与安全性}

\textbf{\passthrough{\lstinline!AVAILABLE!} /
\passthrough{\lstinline!VALUE\_TYPE!}}：当前绑定源码已实现该消息值操作。

\textbf{参数}

\begin{longtable}[]{@{}
  >{\raggedright\arraybackslash}p{(\columnwidth - 6\tabcolsep) * \real{0.18}}
  >{\raggedright\arraybackslash}p{(\columnwidth - 6\tabcolsep) * \real{0.24}}
  >{\raggedright\arraybackslash}p{(\columnwidth - 6\tabcolsep) * \real{0.18}}
  >{\raggedright\arraybackslash}p{(\columnwidth - 6\tabcolsep) * \real{0.40}}@{}}
\toprule
名称 & Python 类型 & 默认值 & 含义 \\
\midrule
\endhead
\passthrough{\lstinline!other!} & \passthrough{\lstinline!object!} &
必填 & 用于比较的另一个 Python 对象。类型不兼容时比较结果通常为
\passthrough{\lstinline!False!}。 \\
\bottomrule
\end{longtable}

\textbf{返回值}

\begin{longtable}[]{@{}
  >{\raggedright\arraybackslash}p{(\columnwidth - 4\tabcolsep) * \real{0.18}}
  >{\raggedright\arraybackslash}p{(\columnwidth - 4\tabcolsep) * \real{0.20}}
  >{\raggedright\arraybackslash}p{(\columnwidth - 4\tabcolsep) * \real{0.62}}@{}}
\toprule
位置 & 类型 & 含义 \\
\midrule
\endhead
返回值 & \passthrough{\lstinline!bool!} & 返回
\passthrough{\lstinline!bool!}。更精确的业务含义见该方法用途、C++
签名和目标型号协议。 \\
\bottomrule
\end{longtable}

\textbf{对应 C++}

\begin{itemize}
\tightlist
\item
  类：\passthrough{\lstinline!unitree\_go::msg::dds\_::VoxelMapCompressed\_!}
\item
  签名：\passthrough{\lstinline"operator!=(const value \&) const"}
\item
  绑定策略：\passthrough{\lstinline!未单独标注!}
\item
  声明位置：由 pybind11 手工绑定或生成注册表提供；manifest
  未记录独立头文件行号。
\end{itemize}

\textbf{说明}

\begin{itemize}
\tightlist
\item
  示例是局部调用片段；\passthrough{\lstinline!obj!}
  和参数变量表示已经按上层业务校验并准备好的对象。
\end{itemize}

\textbf{用法}

\begin{lstlisting}[language=Python]
different = left != right
\end{lstlisting}

\hypertarget{unitree-sdk2-cpp-idl-go2-voxelmapcompressed-stamp}{%
\paragraph{\texorpdfstring{\texttt{unitree\_sdk2\_cpp.idl.go2.VoxelMapCompressed.stamp}}{unitree\_sdk2\_cpp.idl.go2.VoxelMapCompressed.stamp}}\label{unitree-sdk2-cpp-idl-go2-voxelmapcompressed-stamp}}

底层
\passthrough{\lstinline!unitree\_go::msg::dds\_::VoxelMapCompressed\_!}
的 \passthrough{\lstinline!stamp!} 字段。类型签名只说明 Python
表示；单位、数值范围、数组固定长度和枚举含义需查对应 IDL 头文件。
底层为浮点数；类型本身不说明单位、坐标系或业务安全范围。

\textbf{签名}

\begin{lstlisting}[language=Python]
@property
def stamp(self) -> float

@stamp.setter
def stamp(self, value: float) -> None
\end{lstlisting}

\textbf{可用性与安全性}

\textbf{\passthrough{\lstinline!AVAILABLE!} /
\passthrough{\lstinline!VALUE\_TYPE!}}：当前绑定源码已实现该消息值操作。

\textbf{参数}

\begin{longtable}[]{@{}
  >{\raggedright\arraybackslash}p{(\columnwidth - 4\tabcolsep) * \real{0.18}}
  >{\raggedright\arraybackslash}p{(\columnwidth - 4\tabcolsep) * \real{0.20}}
  >{\raggedright\arraybackslash}p{(\columnwidth - 4\tabcolsep) * \real{0.62}}@{}}
\toprule
名称 & Python 类型 & 含义 \\
\midrule
\endhead
\passthrough{\lstinline!value!} & \passthrough{\lstinline!float!} &
要写入或传递的新值，必须符合该参数的 Python 类型以及底层 C++
范围约束。 \\
\bottomrule
\end{longtable}

\textbf{返回值}

\begin{longtable}[]{@{}
  >{\raggedright\arraybackslash}p{(\columnwidth - 4\tabcolsep) * \real{0.18}}
  >{\raggedright\arraybackslash}p{(\columnwidth - 4\tabcolsep) * \real{0.20}}
  >{\raggedright\arraybackslash}p{(\columnwidth - 4\tabcolsep) * \real{0.62}}@{}}
\toprule
位置 & 类型 & 含义 \\
\midrule
\endhead
getter 返回值 & \passthrough{\lstinline!float!} & 返回属性当前值。 \\
\bottomrule
\end{longtable}

\textbf{对应 C++}

\begin{itemize}
\tightlist
\item
  类：\passthrough{\lstinline!unitree\_go::msg::dds\_::VoxelMapCompressed\_!}
\item
  字段类型：\passthrough{\lstinline!double!}
\item
  声明位置：\passthrough{\lstinline!include/unitree/idl/go2/VoxelMapCompressed\_.hpp:26!}
\end{itemize}

\textbf{用法}

\begin{lstlisting}[language=Python]
current_value = obj.stamp
obj.stamp = new_value
\end{lstlisting}

\hypertarget{unitree-sdk2-cpp-idl-go2-voxelmapcompressed-frame-id}{%
\paragraph{\texorpdfstring{\texttt{unitree\_sdk2\_cpp.idl.go2.VoxelMapCompressed.frame\_id}}{unitree\_sdk2\_cpp.idl.go2.VoxelMapCompressed.frame\_id}}\label{unitree-sdk2-cpp-idl-go2-voxelmapcompressed-frame-id}}

底层
\passthrough{\lstinline!unitree\_go::msg::dds\_::VoxelMapCompressed\_!}
的 \passthrough{\lstinline!frame\_id!} 字段。类型签名只说明 Python
表示；单位、数值范围、数组固定长度和枚举含义需查对应 IDL 头文件。
底层为字符串；长度、编码和允许值由具体协议决定。

\textbf{签名}

\begin{lstlisting}[language=Python]
@property
def frame_id(self) -> str

@frame_id.setter
def frame_id(self, value: str) -> None
\end{lstlisting}

\textbf{可用性与安全性}

\textbf{\passthrough{\lstinline!AVAILABLE!} /
\passthrough{\lstinline!VALUE\_TYPE!}}：当前绑定源码已实现该消息值操作。

\textbf{参数}

\begin{longtable}[]{@{}
  >{\raggedright\arraybackslash}p{(\columnwidth - 4\tabcolsep) * \real{0.18}}
  >{\raggedright\arraybackslash}p{(\columnwidth - 4\tabcolsep) * \real{0.20}}
  >{\raggedright\arraybackslash}p{(\columnwidth - 4\tabcolsep) * \real{0.62}}@{}}
\toprule
名称 & Python 类型 & 含义 \\
\midrule
\endhead
\passthrough{\lstinline!value!} & \passthrough{\lstinline!str!} &
要写入或传递的新值，必须符合该参数的 Python 类型以及底层 C++
范围约束。 \\
\bottomrule
\end{longtable}

\textbf{返回值}

\begin{longtable}[]{@{}
  >{\raggedright\arraybackslash}p{(\columnwidth - 4\tabcolsep) * \real{0.18}}
  >{\raggedright\arraybackslash}p{(\columnwidth - 4\tabcolsep) * \real{0.20}}
  >{\raggedright\arraybackslash}p{(\columnwidth - 4\tabcolsep) * \real{0.62}}@{}}
\toprule
位置 & 类型 & 含义 \\
\midrule
\endhead
getter 返回值 & \passthrough{\lstinline!str!} & 返回属性当前值。 \\
\bottomrule
\end{longtable}

\textbf{对应 C++}

\begin{itemize}
\tightlist
\item
  类：\passthrough{\lstinline!unitree\_go::msg::dds\_::VoxelMapCompressed\_!}
\item
  字段类型：\passthrough{\lstinline!std::string!}
\item
  声明位置：\passthrough{\lstinline!include/unitree/idl/go2/VoxelMapCompressed\_.hpp:27!}
\end{itemize}

\textbf{用法}

\begin{lstlisting}[language=Python]
current_value = obj.frame_id
obj.frame_id = new_value
\end{lstlisting}

\hypertarget{unitree-sdk2-cpp-idl-go2-voxelmapcompressed-resolution}{%
\paragraph{\texorpdfstring{\texttt{unitree\_sdk2\_cpp.idl.go2.VoxelMapCompressed.resolution}}{unitree\_sdk2\_cpp.idl.go2.VoxelMapCompressed.resolution}}\label{unitree-sdk2-cpp-idl-go2-voxelmapcompressed-resolution}}

底层
\passthrough{\lstinline!unitree\_go::msg::dds\_::VoxelMapCompressed\_!}
的 \passthrough{\lstinline!resolution!} 字段。类型签名只说明 Python
表示；单位、数值范围、数组固定长度和枚举含义需查对应 IDL 头文件。
底层为浮点数；类型本身不说明单位、坐标系或业务安全范围。

\textbf{签名}

\begin{lstlisting}[language=Python]
@property
def resolution(self) -> float

@resolution.setter
def resolution(self, value: float) -> None
\end{lstlisting}

\textbf{可用性与安全性}

\textbf{\passthrough{\lstinline!AVAILABLE!} /
\passthrough{\lstinline!VALUE\_TYPE!}}：当前绑定源码已实现该消息值操作。

\textbf{参数}

\begin{longtable}[]{@{}
  >{\raggedright\arraybackslash}p{(\columnwidth - 4\tabcolsep) * \real{0.18}}
  >{\raggedright\arraybackslash}p{(\columnwidth - 4\tabcolsep) * \real{0.20}}
  >{\raggedright\arraybackslash}p{(\columnwidth - 4\tabcolsep) * \real{0.62}}@{}}
\toprule
名称 & Python 类型 & 含义 \\
\midrule
\endhead
\passthrough{\lstinline!value!} & \passthrough{\lstinline!float!} &
要写入或传递的新值，必须符合该参数的 Python 类型以及底层 C++
范围约束。 \\
\bottomrule
\end{longtable}

\textbf{返回值}

\begin{longtable}[]{@{}
  >{\raggedright\arraybackslash}p{(\columnwidth - 4\tabcolsep) * \real{0.18}}
  >{\raggedright\arraybackslash}p{(\columnwidth - 4\tabcolsep) * \real{0.20}}
  >{\raggedright\arraybackslash}p{(\columnwidth - 4\tabcolsep) * \real{0.62}}@{}}
\toprule
位置 & 类型 & 含义 \\
\midrule
\endhead
getter 返回值 & \passthrough{\lstinline!float!} & 返回属性当前值。 \\
\bottomrule
\end{longtable}

\textbf{对应 C++}

\begin{itemize}
\tightlist
\item
  类：\passthrough{\lstinline!unitree\_go::msg::dds\_::VoxelMapCompressed\_!}
\item
  字段类型：\passthrough{\lstinline!double!}
\item
  声明位置：\passthrough{\lstinline!include/unitree/idl/go2/VoxelMapCompressed\_.hpp:28!}
\end{itemize}

\textbf{用法}

\begin{lstlisting}[language=Python]
current_value = obj.resolution
obj.resolution = new_value
\end{lstlisting}

\hypertarget{unitree-sdk2-cpp-idl-go2-voxelmapcompressed-origin}{%
\paragraph{\texorpdfstring{\texttt{unitree\_sdk2\_cpp.idl.go2.VoxelMapCompressed.origin}}{unitree\_sdk2\_cpp.idl.go2.VoxelMapCompressed.origin}}\label{unitree-sdk2-cpp-idl-go2-voxelmapcompressed-origin}}

底层
\passthrough{\lstinline!unitree\_go::msg::dds\_::VoxelMapCompressed\_!}
的 \passthrough{\lstinline!origin!} 字段。类型签名只说明 Python
表示；单位、数值范围、数组固定长度和枚举含义需查对应 IDL 头文件。
底层是固定长度数组，写入序列必须正好包含 3
个元素；元素约束：底层为浮点数；类型本身不说明单位、坐标系或业务安全范围。

\textbf{签名}

\begin{lstlisting}[language=Python]
@property
def origin(self) -> list[float]

@origin.setter
def origin(self, value: Sequence[float]) -> None
\end{lstlisting}

\textbf{可用性与安全性}

\textbf{\passthrough{\lstinline!AVAILABLE!} /
\passthrough{\lstinline!VALUE\_TYPE!}}：当前绑定源码已实现该消息值操作。

\textbf{参数}

\begin{longtable}[]{@{}
  >{\raggedright\arraybackslash}p{(\columnwidth - 4\tabcolsep) * \real{0.18}}
  >{\raggedright\arraybackslash}p{(\columnwidth - 4\tabcolsep) * \real{0.20}}
  >{\raggedright\arraybackslash}p{(\columnwidth - 4\tabcolsep) * \real{0.62}}@{}}
\toprule
名称 & Python 类型 & 含义 \\
\midrule
\endhead
\passthrough{\lstinline!value!} &
\passthrough{\lstinline!Sequence[float]!} &
要写入或传递的新值，必须符合该参数的 Python 类型以及底层 C++
范围约束。 \\
\bottomrule
\end{longtable}

\textbf{返回值}

\begin{longtable}[]{@{}
  >{\raggedright\arraybackslash}p{(\columnwidth - 4\tabcolsep) * \real{0.18}}
  >{\raggedright\arraybackslash}p{(\columnwidth - 4\tabcolsep) * \real{0.20}}
  >{\raggedright\arraybackslash}p{(\columnwidth - 4\tabcolsep) * \real{0.62}}@{}}
\toprule
位置 & 类型 & 含义 \\
\midrule
\endhead
getter 返回值 & \passthrough{\lstinline!list[float]!} &
返回属性当前值。数组、vector 和嵌套消息采用复制语义。 \\
\bottomrule
\end{longtable}

\textbf{对应 C++}

\begin{itemize}
\tightlist
\item
  类：\passthrough{\lstinline!unitree\_go::msg::dds\_::VoxelMapCompressed\_!}
\item
  字段类型：\passthrough{\lstinline!std::array<double, 3>!}
\item
  声明位置：\passthrough{\lstinline!include/unitree/idl/go2/VoxelMapCompressed\_.hpp:29!}
\end{itemize}

\textbf{用法}

\begin{lstlisting}[language=Python]
current_value = obj.origin
obj.origin = new_value
\end{lstlisting}

\begin{quote}
{[}!TIP{]} 读取容器或嵌套值后应遵循``读取、修改、重新赋值''。只修改
getter 返回的副本不会可靠地更新原 C++ 消息。
\end{quote}

\hypertarget{unitree-sdk2-cpp-idl-go2-voxelmapcompressed-width}{%
\paragraph{\texorpdfstring{\texttt{unitree\_sdk2\_cpp.idl.go2.VoxelMapCompressed.width}}{unitree\_sdk2\_cpp.idl.go2.VoxelMapCompressed.width}}\label{unitree-sdk2-cpp-idl-go2-voxelmapcompressed-width}}

底层
\passthrough{\lstinline!unitree\_go::msg::dds\_::VoxelMapCompressed\_!}
的 \passthrough{\lstinline!width!} 字段。类型签名只说明 Python
表示；单位、数值范围、数组固定长度和枚举含义需查对应 IDL 头文件。
底层是固定长度数组，写入序列必须正好包含 3 个元素；元素约束：取值范围为
0 到 65535。

\textbf{签名}

\begin{lstlisting}[language=Python]
@property
def width(self) -> list[int]

@width.setter
def width(self, value: Sequence[int]) -> None
\end{lstlisting}

\textbf{可用性与安全性}

\textbf{\passthrough{\lstinline!AVAILABLE!} /
\passthrough{\lstinline!VALUE\_TYPE!}}：当前绑定源码已实现该消息值操作。

\textbf{参数}

\begin{longtable}[]{@{}
  >{\raggedright\arraybackslash}p{(\columnwidth - 4\tabcolsep) * \real{0.18}}
  >{\raggedright\arraybackslash}p{(\columnwidth - 4\tabcolsep) * \real{0.20}}
  >{\raggedright\arraybackslash}p{(\columnwidth - 4\tabcolsep) * \real{0.62}}@{}}
\toprule
名称 & Python 类型 & 含义 \\
\midrule
\endhead
\passthrough{\lstinline!value!} &
\passthrough{\lstinline!Sequence[int]!} &
要写入或传递的新值，必须符合该参数的 Python 类型以及底层 C++
范围约束。 \\
\bottomrule
\end{longtable}

\textbf{返回值}

\begin{longtable}[]{@{}
  >{\raggedright\arraybackslash}p{(\columnwidth - 4\tabcolsep) * \real{0.18}}
  >{\raggedright\arraybackslash}p{(\columnwidth - 4\tabcolsep) * \real{0.20}}
  >{\raggedright\arraybackslash}p{(\columnwidth - 4\tabcolsep) * \real{0.62}}@{}}
\toprule
位置 & 类型 & 含义 \\
\midrule
\endhead
getter 返回值 & \passthrough{\lstinline!list[int]!} &
返回属性当前值。数组、vector 和嵌套消息采用复制语义。 \\
\bottomrule
\end{longtable}

\textbf{对应 C++}

\begin{itemize}
\tightlist
\item
  类：\passthrough{\lstinline!unitree\_go::msg::dds\_::VoxelMapCompressed\_!}
\item
  字段类型：\passthrough{\lstinline!std::array<uint16\_t, 3>!}
\item
  声明位置：\passthrough{\lstinline!include/unitree/idl/go2/VoxelMapCompressed\_.hpp:30!}
\end{itemize}

\textbf{用法}

\begin{lstlisting}[language=Python]
current_value = obj.width
obj.width = new_value
\end{lstlisting}

\begin{quote}
{[}!TIP{]} 读取容器或嵌套值后应遵循``读取、修改、重新赋值''。只修改
getter 返回的副本不会可靠地更新原 C++ 消息。
\end{quote}

\hypertarget{unitree-sdk2-cpp-idl-go2-voxelmapcompressed-src-size}{%
\paragraph{\texorpdfstring{\texttt{unitree\_sdk2\_cpp.idl.go2.VoxelMapCompressed.src\_size}}{unitree\_sdk2\_cpp.idl.go2.VoxelMapCompressed.src\_size}}\label{unitree-sdk2-cpp-idl-go2-voxelmapcompressed-src-size}}

底层
\passthrough{\lstinline!unitree\_go::msg::dds\_::VoxelMapCompressed\_!}
的 \passthrough{\lstinline!src\_size!} 字段。类型签名只说明 Python
表示；单位、数值范围、数组固定长度和枚举含义需查对应 IDL 头文件。
取值范围为 0 到 18446744073709551615。

\textbf{签名}

\begin{lstlisting}[language=Python]
@property
def src_size(self) -> int

@src_size.setter
def src_size(self, value: int) -> None
\end{lstlisting}

\textbf{可用性与安全性}

\textbf{\passthrough{\lstinline!AVAILABLE!} /
\passthrough{\lstinline!VALUE\_TYPE!}}：当前绑定源码已实现该消息值操作。

\textbf{参数}

\begin{longtable}[]{@{}
  >{\raggedright\arraybackslash}p{(\columnwidth - 4\tabcolsep) * \real{0.18}}
  >{\raggedright\arraybackslash}p{(\columnwidth - 4\tabcolsep) * \real{0.20}}
  >{\raggedright\arraybackslash}p{(\columnwidth - 4\tabcolsep) * \real{0.62}}@{}}
\toprule
名称 & Python 类型 & 含义 \\
\midrule
\endhead
\passthrough{\lstinline!value!} & \passthrough{\lstinline!int!} &
要写入或传递的新值，必须符合该参数的 Python 类型以及底层 C++
范围约束。 \\
\bottomrule
\end{longtable}

\textbf{返回值}

\begin{longtable}[]{@{}
  >{\raggedright\arraybackslash}p{(\columnwidth - 4\tabcolsep) * \real{0.18}}
  >{\raggedright\arraybackslash}p{(\columnwidth - 4\tabcolsep) * \real{0.20}}
  >{\raggedright\arraybackslash}p{(\columnwidth - 4\tabcolsep) * \real{0.62}}@{}}
\toprule
位置 & 类型 & 含义 \\
\midrule
\endhead
getter 返回值 & \passthrough{\lstinline!int!} & 返回属性当前值。 \\
\bottomrule
\end{longtable}

\textbf{对应 C++}

\begin{itemize}
\tightlist
\item
  类：\passthrough{\lstinline!unitree\_go::msg::dds\_::VoxelMapCompressed\_!}
\item
  字段类型：\passthrough{\lstinline!uint64\_t!}
\item
  声明位置：\passthrough{\lstinline!include/unitree/idl/go2/VoxelMapCompressed\_.hpp:31!}
\end{itemize}

\textbf{用法}

\begin{lstlisting}[language=Python]
current_value = obj.src_size
obj.src_size = new_value
\end{lstlisting}

\hypertarget{unitree-sdk2-cpp-idl-go2-voxelmapcompressed-data}{%
\paragraph{\texorpdfstring{\texttt{unitree\_sdk2\_cpp.idl.go2.VoxelMapCompressed.data}}{unitree\_sdk2\_cpp.idl.go2.VoxelMapCompressed.data}}\label{unitree-sdk2-cpp-idl-go2-voxelmapcompressed-data}}

底层
\passthrough{\lstinline!unitree\_go::msg::dds\_::VoxelMapCompressed\_!}
的 \passthrough{\lstinline!data!} 字段。类型签名只说明 Python
表示；单位、数值范围、数组固定长度和枚举含义需查对应 IDL 头文件。
底层是可变长度 vector；元素约束：取值范围为 0 到 255。

\textbf{签名}

\begin{lstlisting}[language=Python]
@property
def data(self) -> list[int]

@data.setter
def data(self, value: Sequence[int]) -> None
\end{lstlisting}

\textbf{可用性与安全性}

\textbf{\passthrough{\lstinline!AVAILABLE!} /
\passthrough{\lstinline!VALUE\_TYPE!}}：当前绑定源码已实现该消息值操作。

\textbf{参数}

\begin{longtable}[]{@{}
  >{\raggedright\arraybackslash}p{(\columnwidth - 4\tabcolsep) * \real{0.18}}
  >{\raggedright\arraybackslash}p{(\columnwidth - 4\tabcolsep) * \real{0.20}}
  >{\raggedright\arraybackslash}p{(\columnwidth - 4\tabcolsep) * \real{0.62}}@{}}
\toprule
名称 & Python 类型 & 含义 \\
\midrule
\endhead
\passthrough{\lstinline!value!} &
\passthrough{\lstinline!Sequence[int]!} &
要写入或传递的新值，必须符合该参数的 Python 类型以及底层 C++
范围约束。 \\
\bottomrule
\end{longtable}

\textbf{返回值}

\begin{longtable}[]{@{}
  >{\raggedright\arraybackslash}p{(\columnwidth - 4\tabcolsep) * \real{0.18}}
  >{\raggedright\arraybackslash}p{(\columnwidth - 4\tabcolsep) * \real{0.20}}
  >{\raggedright\arraybackslash}p{(\columnwidth - 4\tabcolsep) * \real{0.62}}@{}}
\toprule
位置 & 类型 & 含义 \\
\midrule
\endhead
getter 返回值 & \passthrough{\lstinline!list[int]!} &
返回属性当前值。数组、vector 和嵌套消息采用复制语义。 \\
\bottomrule
\end{longtable}

\textbf{对应 C++}

\begin{itemize}
\tightlist
\item
  类：\passthrough{\lstinline!unitree\_go::msg::dds\_::VoxelMapCompressed\_!}
\item
  字段类型：\passthrough{\lstinline!std::vector<uint8\_t>!}
\item
  声明位置：\passthrough{\lstinline!include/unitree/idl/go2/VoxelMapCompressed\_.hpp:32!}
\end{itemize}

\textbf{用法}

\begin{lstlisting}[language=Python]
current_value = obj.data
obj.data = new_value
\end{lstlisting}

\begin{quote}
{[}!TIP{]} 读取容器或嵌套值后应遵循``读取、修改、重新赋值''。只修改
getter 返回的副本不会可靠地更新原 C++ 消息。
\end{quote}

\hypertarget{unitree-sdk2-cpp-idl-go2-wirelesscontroller}{%
\subsubsection{\texorpdfstring{\texttt{unitree\_sdk2\_cpp.idl.go2.WirelessController}}{unitree\_sdk2\_cpp.idl.go2.WirelessController}}\label{unitree-sdk2-cpp-idl-go2-wirelesscontroller}}

可由 Python 读写的 DDS/IDL 消息值类型。字段采用复制语义。

\textbf{导入}

\begin{lstlisting}[language=Python]
from unitree_sdk2_cpp.idl.go2 import WirelessController
\end{lstlisting}

\textbf{方法索引}

\begin{longtable}[]{@{}
  >{\raggedright\arraybackslash}p{(\columnwidth - 6\tabcolsep) * \real{0.18}}
  >{\raggedright\arraybackslash}p{(\columnwidth - 6\tabcolsep) * \real{0.24}}
  >{\raggedright\arraybackslash}p{(\columnwidth - 6\tabcolsep) * \real{0.18}}
  >{\raggedright\arraybackslash}p{(\columnwidth - 6\tabcolsep) * \real{0.40}}@{}}
\toprule
方法 & Python 签名 & 状态 & 安全分类 \\
\midrule
\endhead
\protect\hyperlink{unitree-sdk2-cpp-idl-go2-wirelesscontroller-dunder-init-1}{\passthrough{\lstinline!\_\_init\_\_!}}
& \passthrough{\lstinline!def \_\_init\_\_(self) -> None!} &
\passthrough{\lstinline!AVAILABLE!} &
\passthrough{\lstinline!VALUE\_TYPE!} \\
\protect\hyperlink{unitree-sdk2-cpp-idl-go2-wirelesscontroller-dunder-eq-1}{\passthrough{\lstinline!\_\_eq\_\_!}}
& \passthrough{\lstinline!def \_\_eq\_\_(self, other: object) -> bool!}
& \passthrough{\lstinline!AVAILABLE!} &
\passthrough{\lstinline!VALUE\_TYPE!} \\
\protect\hyperlink{unitree-sdk2-cpp-idl-go2-wirelesscontroller-dunder-ne-1}{\passthrough{\lstinline!\_\_ne\_\_!}}
& \passthrough{\lstinline!def \_\_ne\_\_(self, other: object) -> bool!}
& \passthrough{\lstinline!AVAILABLE!} &
\passthrough{\lstinline!VALUE\_TYPE!} \\
\bottomrule
\end{longtable}

\textbf{属性索引}

\begin{longtable}[]{@{}
  >{\raggedright\arraybackslash}p{(\columnwidth - 6\tabcolsep) * \real{0.18}}
  >{\raggedright\arraybackslash}p{(\columnwidth - 6\tabcolsep) * \real{0.24}}
  >{\raggedright\arraybackslash}p{(\columnwidth - 6\tabcolsep) * \real{0.18}}
  >{\raggedright\arraybackslash}p{(\columnwidth - 6\tabcolsep) * \real{0.40}}@{}}
\toprule
属性 & 读取类型 & 写入类型 & 状态 \\
\midrule
\endhead
\protect\hyperlink{unitree-sdk2-cpp-idl-go2-wirelesscontroller-lx}{\passthrough{\lstinline!lx!}}
& \passthrough{\lstinline!float!} & \passthrough{\lstinline!float!} &
\passthrough{\lstinline!AVAILABLE!} \\
\protect\hyperlink{unitree-sdk2-cpp-idl-go2-wirelesscontroller-ly}{\passthrough{\lstinline!ly!}}
& \passthrough{\lstinline!float!} & \passthrough{\lstinline!float!} &
\passthrough{\lstinline!AVAILABLE!} \\
\protect\hyperlink{unitree-sdk2-cpp-idl-go2-wirelesscontroller-rx}{\passthrough{\lstinline!rx!}}
& \passthrough{\lstinline!float!} & \passthrough{\lstinline!float!} &
\passthrough{\lstinline!AVAILABLE!} \\
\protect\hyperlink{unitree-sdk2-cpp-idl-go2-wirelesscontroller-ry}{\passthrough{\lstinline!ry!}}
& \passthrough{\lstinline!float!} & \passthrough{\lstinline!float!} &
\passthrough{\lstinline!AVAILABLE!} \\
\protect\hyperlink{unitree-sdk2-cpp-idl-go2-wirelesscontroller-keys}{\passthrough{\lstinline!keys!}}
& \passthrough{\lstinline!int!} & \passthrough{\lstinline!int!} &
\passthrough{\lstinline!AVAILABLE!} \\
\bottomrule
\end{longtable}

\hypertarget{unitree-sdk2-cpp-idl-go2-wirelesscontroller-dunder-init-1}{%
\paragraph{\texorpdfstring{\texttt{unitree\_sdk2\_cpp.idl.go2.WirelessController.\_\_init\_\_}}{unitree\_sdk2\_cpp.idl.go2.WirelessController.\_\_init\_\_}}\label{unitree-sdk2-cpp-idl-go2-wirelesscontroller-dunder-init-1}}

初始化 \passthrough{\lstinline!WirelessController!}
实例。是否能实际构造取决于下方可用性状态。

\textbf{签名}

\begin{lstlisting}[language=Python]
def __init__(self) -> None
\end{lstlisting}

\textbf{可用性与安全性}

\textbf{\passthrough{\lstinline!AVAILABLE!} /
\passthrough{\lstinline!VALUE\_TYPE!}}：当前绑定源码已实现该消息值操作。

\textbf{参数}

无显式参数。实例方法中的 \passthrough{\lstinline!self!} 由 Python
自动传入。

\textbf{返回值}

\begin{longtable}[]{@{}
  >{\raggedright\arraybackslash}p{(\columnwidth - 4\tabcolsep) * \real{0.18}}
  >{\raggedright\arraybackslash}p{(\columnwidth - 4\tabcolsep) * \real{0.20}}
  >{\raggedright\arraybackslash}p{(\columnwidth - 4\tabcolsep) * \real{0.62}}@{}}
\toprule
位置 & 类型 & 含义 \\
\midrule
\endhead
无 & \passthrough{\lstinline!None!} &
\passthrough{\lstinline!\_\_init\_\_!}
本身不返回值；调用类对象时，在构造函数可用的前提下得到该类实例。 \\
\bottomrule
\end{longtable}

\textbf{对应 C++}

\begin{itemize}
\tightlist
\item
  类：\passthrough{\lstinline!unitree\_go::msg::dds\_::WirelessController\_!}
\item
  签名：\passthrough{\lstinline!WirelessController\_()!}
\item
  绑定策略：\passthrough{\lstinline!未单独标注!}
\item
  声明位置：\passthrough{\lstinline!include/unitree/idl/go2/WirelessController\_.hpp:30!}
\end{itemize}

\textbf{说明}

\begin{itemize}
\tightlist
\item
  示例是局部调用片段；\passthrough{\lstinline!obj!}
  和参数变量表示已经按上层业务校验并准备好的对象。
\end{itemize}

\textbf{用法}

\begin{lstlisting}[language=Python]
value = WirelessController()
\end{lstlisting}

\hypertarget{unitree-sdk2-cpp-idl-go2-wirelesscontroller-dunder-eq-1}{%
\paragraph{\texorpdfstring{\texttt{unitree\_sdk2\_cpp.idl.go2.WirelessController.\_\_eq\_\_}}{unitree\_sdk2\_cpp.idl.go2.WirelessController.\_\_eq\_\_}}\label{unitree-sdk2-cpp-idl-go2-wirelesscontroller-dunder-eq-1}}

按底层 IDL 消息内容比较两个对象是否相等。

\textbf{签名}

\begin{lstlisting}[language=Python]
def __eq__(self, other: object) -> bool
\end{lstlisting}

\textbf{可用性与安全性}

\textbf{\passthrough{\lstinline!AVAILABLE!} /
\passthrough{\lstinline!VALUE\_TYPE!}}：当前绑定源码已实现该消息值操作。

\textbf{参数}

\begin{longtable}[]{@{}
  >{\raggedright\arraybackslash}p{(\columnwidth - 6\tabcolsep) * \real{0.18}}
  >{\raggedright\arraybackslash}p{(\columnwidth - 6\tabcolsep) * \real{0.24}}
  >{\raggedright\arraybackslash}p{(\columnwidth - 6\tabcolsep) * \real{0.18}}
  >{\raggedright\arraybackslash}p{(\columnwidth - 6\tabcolsep) * \real{0.40}}@{}}
\toprule
名称 & Python 类型 & 默认值 & 含义 \\
\midrule
\endhead
\passthrough{\lstinline!other!} & \passthrough{\lstinline!object!} &
必填 & 用于比较的另一个 Python 对象。类型不兼容时比较结果通常为
\passthrough{\lstinline!False!}。 \\
\bottomrule
\end{longtable}

\textbf{返回值}

\begin{longtable}[]{@{}
  >{\raggedright\arraybackslash}p{(\columnwidth - 4\tabcolsep) * \real{0.18}}
  >{\raggedright\arraybackslash}p{(\columnwidth - 4\tabcolsep) * \real{0.20}}
  >{\raggedright\arraybackslash}p{(\columnwidth - 4\tabcolsep) * \real{0.62}}@{}}
\toprule
位置 & 类型 & 含义 \\
\midrule
\endhead
返回值 & \passthrough{\lstinline!bool!} & 返回
\passthrough{\lstinline!bool!}。更精确的业务含义见该方法用途、C++
签名和目标型号协议。 \\
\bottomrule
\end{longtable}

\textbf{对应 C++}

\begin{itemize}
\tightlist
\item
  类：\passthrough{\lstinline!unitree\_go::msg::dds\_::WirelessController\_!}
\item
  签名：\passthrough{\lstinline!operator==(const value \&) const!}
\item
  绑定策略：\passthrough{\lstinline!未单独标注!}
\item
  声明位置：由 pybind11 手工绑定或生成注册表提供；manifest
  未记录独立头文件行号。
\end{itemize}

\textbf{说明}

\begin{itemize}
\tightlist
\item
  示例是局部调用片段；\passthrough{\lstinline!obj!}
  和参数变量表示已经按上层业务校验并准备好的对象。
\end{itemize}

\textbf{用法}

\begin{lstlisting}[language=Python]
same = left == right
\end{lstlisting}

\hypertarget{unitree-sdk2-cpp-idl-go2-wirelesscontroller-dunder-ne-1}{%
\paragraph{\texorpdfstring{\texttt{unitree\_sdk2\_cpp.idl.go2.WirelessController.\_\_ne\_\_}}{unitree\_sdk2\_cpp.idl.go2.WirelessController.\_\_ne\_\_}}\label{unitree-sdk2-cpp-idl-go2-wirelesscontroller-dunder-ne-1}}

按底层 IDL 消息内容比较两个对象是否不相等。

\textbf{签名}

\begin{lstlisting}[language=Python]
def __ne__(self, other: object) -> bool
\end{lstlisting}

\textbf{可用性与安全性}

\textbf{\passthrough{\lstinline!AVAILABLE!} /
\passthrough{\lstinline!VALUE\_TYPE!}}：当前绑定源码已实现该消息值操作。

\textbf{参数}

\begin{longtable}[]{@{}
  >{\raggedright\arraybackslash}p{(\columnwidth - 6\tabcolsep) * \real{0.18}}
  >{\raggedright\arraybackslash}p{(\columnwidth - 6\tabcolsep) * \real{0.24}}
  >{\raggedright\arraybackslash}p{(\columnwidth - 6\tabcolsep) * \real{0.18}}
  >{\raggedright\arraybackslash}p{(\columnwidth - 6\tabcolsep) * \real{0.40}}@{}}
\toprule
名称 & Python 类型 & 默认值 & 含义 \\
\midrule
\endhead
\passthrough{\lstinline!other!} & \passthrough{\lstinline!object!} &
必填 & 用于比较的另一个 Python 对象。类型不兼容时比较结果通常为
\passthrough{\lstinline!False!}。 \\
\bottomrule
\end{longtable}

\textbf{返回值}

\begin{longtable}[]{@{}
  >{\raggedright\arraybackslash}p{(\columnwidth - 4\tabcolsep) * \real{0.18}}
  >{\raggedright\arraybackslash}p{(\columnwidth - 4\tabcolsep) * \real{0.20}}
  >{\raggedright\arraybackslash}p{(\columnwidth - 4\tabcolsep) * \real{0.62}}@{}}
\toprule
位置 & 类型 & 含义 \\
\midrule
\endhead
返回值 & \passthrough{\lstinline!bool!} & 返回
\passthrough{\lstinline!bool!}。更精确的业务含义见该方法用途、C++
签名和目标型号协议。 \\
\bottomrule
\end{longtable}

\textbf{对应 C++}

\begin{itemize}
\tightlist
\item
  类：\passthrough{\lstinline!unitree\_go::msg::dds\_::WirelessController\_!}
\item
  签名：\passthrough{\lstinline"operator!=(const value \&) const"}
\item
  绑定策略：\passthrough{\lstinline!未单独标注!}
\item
  声明位置：由 pybind11 手工绑定或生成注册表提供；manifest
  未记录独立头文件行号。
\end{itemize}

\textbf{说明}

\begin{itemize}
\tightlist
\item
  示例是局部调用片段；\passthrough{\lstinline!obj!}
  和参数变量表示已经按上层业务校验并准备好的对象。
\end{itemize}

\textbf{用法}

\begin{lstlisting}[language=Python]
different = left != right
\end{lstlisting}

\hypertarget{unitree-sdk2-cpp-idl-go2-wirelesscontroller-lx}{%
\paragraph{\texorpdfstring{\texttt{unitree\_sdk2\_cpp.idl.go2.WirelessController.lx}}{unitree\_sdk2\_cpp.idl.go2.WirelessController.lx}}\label{unitree-sdk2-cpp-idl-go2-wirelesscontroller-lx}}

底层
\passthrough{\lstinline!unitree\_go::msg::dds\_::WirelessController\_!}
的 \passthrough{\lstinline!lx!} 字段。类型签名只说明 Python
表示；单位、数值范围、数组固定长度和枚举含义需查对应 IDL 头文件。
底层为浮点数；类型本身不说明单位、坐标系或业务安全范围。

\textbf{签名}

\begin{lstlisting}[language=Python]
@property
def lx(self) -> float

@lx.setter
def lx(self, value: float) -> None
\end{lstlisting}

\textbf{可用性与安全性}

\textbf{\passthrough{\lstinline!AVAILABLE!} /
\passthrough{\lstinline!VALUE\_TYPE!}}：当前绑定源码已实现该消息值操作。

\textbf{参数}

\begin{longtable}[]{@{}
  >{\raggedright\arraybackslash}p{(\columnwidth - 4\tabcolsep) * \real{0.18}}
  >{\raggedright\arraybackslash}p{(\columnwidth - 4\tabcolsep) * \real{0.20}}
  >{\raggedright\arraybackslash}p{(\columnwidth - 4\tabcolsep) * \real{0.62}}@{}}
\toprule
名称 & Python 类型 & 含义 \\
\midrule
\endhead
\passthrough{\lstinline!value!} & \passthrough{\lstinline!float!} &
要写入或传递的新值，必须符合该参数的 Python 类型以及底层 C++
范围约束。 \\
\bottomrule
\end{longtable}

\textbf{返回值}

\begin{longtable}[]{@{}
  >{\raggedright\arraybackslash}p{(\columnwidth - 4\tabcolsep) * \real{0.18}}
  >{\raggedright\arraybackslash}p{(\columnwidth - 4\tabcolsep) * \real{0.20}}
  >{\raggedright\arraybackslash}p{(\columnwidth - 4\tabcolsep) * \real{0.62}}@{}}
\toprule
位置 & 类型 & 含义 \\
\midrule
\endhead
getter 返回值 & \passthrough{\lstinline!float!} & 返回属性当前值。 \\
\bottomrule
\end{longtable}

\textbf{对应 C++}

\begin{itemize}
\tightlist
\item
  类：\passthrough{\lstinline!unitree\_go::msg::dds\_::WirelessController\_!}
\item
  字段类型：\passthrough{\lstinline!float!}
\item
  声明位置：\passthrough{\lstinline!include/unitree/idl/go2/WirelessController\_.hpp:23!}
\end{itemize}

\textbf{用法}

\begin{lstlisting}[language=Python]
current_value = obj.lx
obj.lx = new_value
\end{lstlisting}

\hypertarget{unitree-sdk2-cpp-idl-go2-wirelesscontroller-ly}{%
\paragraph{\texorpdfstring{\texttt{unitree\_sdk2\_cpp.idl.go2.WirelessController.ly}}{unitree\_sdk2\_cpp.idl.go2.WirelessController.ly}}\label{unitree-sdk2-cpp-idl-go2-wirelesscontroller-ly}}

底层
\passthrough{\lstinline!unitree\_go::msg::dds\_::WirelessController\_!}
的 \passthrough{\lstinline!ly!} 字段。类型签名只说明 Python
表示；单位、数值范围、数组固定长度和枚举含义需查对应 IDL 头文件。
底层为浮点数；类型本身不说明单位、坐标系或业务安全范围。

\textbf{签名}

\begin{lstlisting}[language=Python]
@property
def ly(self) -> float

@ly.setter
def ly(self, value: float) -> None
\end{lstlisting}

\textbf{可用性与安全性}

\textbf{\passthrough{\lstinline!AVAILABLE!} /
\passthrough{\lstinline!VALUE\_TYPE!}}：当前绑定源码已实现该消息值操作。

\textbf{参数}

\begin{longtable}[]{@{}
  >{\raggedright\arraybackslash}p{(\columnwidth - 4\tabcolsep) * \real{0.18}}
  >{\raggedright\arraybackslash}p{(\columnwidth - 4\tabcolsep) * \real{0.20}}
  >{\raggedright\arraybackslash}p{(\columnwidth - 4\tabcolsep) * \real{0.62}}@{}}
\toprule
名称 & Python 类型 & 含义 \\
\midrule
\endhead
\passthrough{\lstinline!value!} & \passthrough{\lstinline!float!} &
要写入或传递的新值，必须符合该参数的 Python 类型以及底层 C++
范围约束。 \\
\bottomrule
\end{longtable}

\textbf{返回值}

\begin{longtable}[]{@{}
  >{\raggedright\arraybackslash}p{(\columnwidth - 4\tabcolsep) * \real{0.18}}
  >{\raggedright\arraybackslash}p{(\columnwidth - 4\tabcolsep) * \real{0.20}}
  >{\raggedright\arraybackslash}p{(\columnwidth - 4\tabcolsep) * \real{0.62}}@{}}
\toprule
位置 & 类型 & 含义 \\
\midrule
\endhead
getter 返回值 & \passthrough{\lstinline!float!} & 返回属性当前值。 \\
\bottomrule
\end{longtable}

\textbf{对应 C++}

\begin{itemize}
\tightlist
\item
  类：\passthrough{\lstinline!unitree\_go::msg::dds\_::WirelessController\_!}
\item
  字段类型：\passthrough{\lstinline!float!}
\item
  声明位置：\passthrough{\lstinline!include/unitree/idl/go2/WirelessController\_.hpp:24!}
\end{itemize}

\textbf{用法}

\begin{lstlisting}[language=Python]
current_value = obj.ly
obj.ly = new_value
\end{lstlisting}

\hypertarget{unitree-sdk2-cpp-idl-go2-wirelesscontroller-rx}{%
\paragraph{\texorpdfstring{\texttt{unitree\_sdk2\_cpp.idl.go2.WirelessController.rx}}{unitree\_sdk2\_cpp.idl.go2.WirelessController.rx}}\label{unitree-sdk2-cpp-idl-go2-wirelesscontroller-rx}}

底层
\passthrough{\lstinline!unitree\_go::msg::dds\_::WirelessController\_!}
的 \passthrough{\lstinline!rx!} 字段。类型签名只说明 Python
表示；单位、数值范围、数组固定长度和枚举含义需查对应 IDL 头文件。
底层为浮点数；类型本身不说明单位、坐标系或业务安全范围。

\textbf{签名}

\begin{lstlisting}[language=Python]
@property
def rx(self) -> float

@rx.setter
def rx(self, value: float) -> None
\end{lstlisting}

\textbf{可用性与安全性}

\textbf{\passthrough{\lstinline!AVAILABLE!} /
\passthrough{\lstinline!VALUE\_TYPE!}}：当前绑定源码已实现该消息值操作。

\textbf{参数}

\begin{longtable}[]{@{}
  >{\raggedright\arraybackslash}p{(\columnwidth - 4\tabcolsep) * \real{0.18}}
  >{\raggedright\arraybackslash}p{(\columnwidth - 4\tabcolsep) * \real{0.20}}
  >{\raggedright\arraybackslash}p{(\columnwidth - 4\tabcolsep) * \real{0.62}}@{}}
\toprule
名称 & Python 类型 & 含义 \\
\midrule
\endhead
\passthrough{\lstinline!value!} & \passthrough{\lstinline!float!} &
要写入或传递的新值，必须符合该参数的 Python 类型以及底层 C++
范围约束。 \\
\bottomrule
\end{longtable}

\textbf{返回值}

\begin{longtable}[]{@{}
  >{\raggedright\arraybackslash}p{(\columnwidth - 4\tabcolsep) * \real{0.18}}
  >{\raggedright\arraybackslash}p{(\columnwidth - 4\tabcolsep) * \real{0.20}}
  >{\raggedright\arraybackslash}p{(\columnwidth - 4\tabcolsep) * \real{0.62}}@{}}
\toprule
位置 & 类型 & 含义 \\
\midrule
\endhead
getter 返回值 & \passthrough{\lstinline!float!} & 返回属性当前值。 \\
\bottomrule
\end{longtable}

\textbf{对应 C++}

\begin{itemize}
\tightlist
\item
  类：\passthrough{\lstinline!unitree\_go::msg::dds\_::WirelessController\_!}
\item
  字段类型：\passthrough{\lstinline!float!}
\item
  声明位置：\passthrough{\lstinline!include/unitree/idl/go2/WirelessController\_.hpp:25!}
\end{itemize}

\textbf{用法}

\begin{lstlisting}[language=Python]
current_value = obj.rx
obj.rx = new_value
\end{lstlisting}

\hypertarget{unitree-sdk2-cpp-idl-go2-wirelesscontroller-ry}{%
\paragraph{\texorpdfstring{\texttt{unitree\_sdk2\_cpp.idl.go2.WirelessController.ry}}{unitree\_sdk2\_cpp.idl.go2.WirelessController.ry}}\label{unitree-sdk2-cpp-idl-go2-wirelesscontroller-ry}}

底层
\passthrough{\lstinline!unitree\_go::msg::dds\_::WirelessController\_!}
的 \passthrough{\lstinline!ry!} 字段。类型签名只说明 Python
表示；单位、数值范围、数组固定长度和枚举含义需查对应 IDL 头文件。
底层为浮点数；类型本身不说明单位、坐标系或业务安全范围。

\textbf{签名}

\begin{lstlisting}[language=Python]
@property
def ry(self) -> float

@ry.setter
def ry(self, value: float) -> None
\end{lstlisting}

\textbf{可用性与安全性}

\textbf{\passthrough{\lstinline!AVAILABLE!} /
\passthrough{\lstinline!VALUE\_TYPE!}}：当前绑定源码已实现该消息值操作。

\textbf{参数}

\begin{longtable}[]{@{}
  >{\raggedright\arraybackslash}p{(\columnwidth - 4\tabcolsep) * \real{0.18}}
  >{\raggedright\arraybackslash}p{(\columnwidth - 4\tabcolsep) * \real{0.20}}
  >{\raggedright\arraybackslash}p{(\columnwidth - 4\tabcolsep) * \real{0.62}}@{}}
\toprule
名称 & Python 类型 & 含义 \\
\midrule
\endhead
\passthrough{\lstinline!value!} & \passthrough{\lstinline!float!} &
要写入或传递的新值，必须符合该参数的 Python 类型以及底层 C++
范围约束。 \\
\bottomrule
\end{longtable}

\textbf{返回值}

\begin{longtable}[]{@{}
  >{\raggedright\arraybackslash}p{(\columnwidth - 4\tabcolsep) * \real{0.18}}
  >{\raggedright\arraybackslash}p{(\columnwidth - 4\tabcolsep) * \real{0.20}}
  >{\raggedright\arraybackslash}p{(\columnwidth - 4\tabcolsep) * \real{0.62}}@{}}
\toprule
位置 & 类型 & 含义 \\
\midrule
\endhead
getter 返回值 & \passthrough{\lstinline!float!} & 返回属性当前值。 \\
\bottomrule
\end{longtable}

\textbf{对应 C++}

\begin{itemize}
\tightlist
\item
  类：\passthrough{\lstinline!unitree\_go::msg::dds\_::WirelessController\_!}
\item
  字段类型：\passthrough{\lstinline!float!}
\item
  声明位置：\passthrough{\lstinline!include/unitree/idl/go2/WirelessController\_.hpp:26!}
\end{itemize}

\textbf{用法}

\begin{lstlisting}[language=Python]
current_value = obj.ry
obj.ry = new_value
\end{lstlisting}

\hypertarget{unitree-sdk2-cpp-idl-go2-wirelesscontroller-keys}{%
\paragraph{\texorpdfstring{\texttt{unitree\_sdk2\_cpp.idl.go2.WirelessController.keys}}{unitree\_sdk2\_cpp.idl.go2.WirelessController.keys}}\label{unitree-sdk2-cpp-idl-go2-wirelesscontroller-keys}}

底层
\passthrough{\lstinline!unitree\_go::msg::dds\_::WirelessController\_!}
的 \passthrough{\lstinline!keys!} 字段。类型签名只说明 Python
表示；单位、数值范围、数组固定长度和枚举含义需查对应 IDL 头文件。
取值范围为 0 到 65535。

\textbf{签名}

\begin{lstlisting}[language=Python]
@property
def keys(self) -> int

@keys.setter
def keys(self, value: int) -> None
\end{lstlisting}

\textbf{可用性与安全性}

\textbf{\passthrough{\lstinline!AVAILABLE!} /
\passthrough{\lstinline!VALUE\_TYPE!}}：当前绑定源码已实现该消息值操作。

\textbf{参数}

\begin{longtable}[]{@{}
  >{\raggedright\arraybackslash}p{(\columnwidth - 4\tabcolsep) * \real{0.18}}
  >{\raggedright\arraybackslash}p{(\columnwidth - 4\tabcolsep) * \real{0.20}}
  >{\raggedright\arraybackslash}p{(\columnwidth - 4\tabcolsep) * \real{0.62}}@{}}
\toprule
名称 & Python 类型 & 含义 \\
\midrule
\endhead
\passthrough{\lstinline!value!} & \passthrough{\lstinline!int!} &
要写入或传递的新值，必须符合该参数的 Python 类型以及底层 C++
范围约束。 \\
\bottomrule
\end{longtable}

\textbf{返回值}

\begin{longtable}[]{@{}
  >{\raggedright\arraybackslash}p{(\columnwidth - 4\tabcolsep) * \real{0.18}}
  >{\raggedright\arraybackslash}p{(\columnwidth - 4\tabcolsep) * \real{0.20}}
  >{\raggedright\arraybackslash}p{(\columnwidth - 4\tabcolsep) * \real{0.62}}@{}}
\toprule
位置 & 类型 & 含义 \\
\midrule
\endhead
getter 返回值 & \passthrough{\lstinline!int!} & 返回属性当前值。 \\
\bottomrule
\end{longtable}

\textbf{对应 C++}

\begin{itemize}
\tightlist
\item
  类：\passthrough{\lstinline!unitree\_go::msg::dds\_::WirelessController\_!}
\item
  字段类型：\passthrough{\lstinline!uint16\_t!}
\item
  声明位置：\passthrough{\lstinline!include/unitree/idl/go2/WirelessController\_.hpp:27!}
\end{itemize}

\textbf{用法}

\begin{lstlisting}[language=Python]
current_value = obj.keys
obj.keys = new_value
\end{lstlisting}

\begin{center}\rule{0.5\linewidth}{0.5pt}\end{center}

\hypertarget{unitree-sdk2-cpp-idl-hg}{%
\subsection{HG IDL 消息}\label{unitree-sdk2-cpp-idl-hg}}

模块：\passthrough{\lstinline!unitree\_sdk2\_cpp.idl.hg!}

\hypertarget{ux7c7bux7d22ux5f15-3}{%
\subsubsection{类索引}\label{ux7c7bux7d22ux5f15-3}}

\begin{longtable}[]{@{}
  >{\raggedright\arraybackslash}p{(\columnwidth - 4\tabcolsep) * \real{0.18}}
  >{\raggedleft\arraybackslash}p{(\columnwidth - 4\tabcolsep) * \real{0.20}}
  >{\raggedleft\arraybackslash}p{(\columnwidth - 4\tabcolsep) * \real{0.62}}@{}}
\toprule
类 & 公开函数签名 & 属性 \\
\midrule
\endhead
\protect\hyperlink{unitree-sdk2-cpp-idl-hg-agvbmsstate}{\passthrough{\lstinline!AgvBmsState!}}
& 3 & 7 \\
\protect\hyperlink{unitree-sdk2-cpp-idl-hg-bmscmd}{\passthrough{\lstinline!BmsCmd!}}
& 3 & 2 \\
\protect\hyperlink{unitree-sdk2-cpp-idl-hg-bmsstate}{\passthrough{\lstinline!BmsState!}}
& 3 & 13 \\
\protect\hyperlink{unitree-sdk2-cpp-idl-hg-motorcmd}{\passthrough{\lstinline!MotorCmd!}}
& 3 & 7 \\
\protect\hyperlink{unitree-sdk2-cpp-idl-hg-handcmd}{\passthrough{\lstinline!HandCmd!}}
& 3 & 2 \\
\protect\hyperlink{unitree-sdk2-cpp-idl-hg-imustate}{\passthrough{\lstinline!IMUState!}}
& 3 & 5 \\
\protect\hyperlink{unitree-sdk2-cpp-idl-hg-motorstate}{\passthrough{\lstinline!MotorState!}}
& 3 & 10 \\
\protect\hyperlink{unitree-sdk2-cpp-idl-hg-presssensorstate}{\passthrough{\lstinline!PressSensorState!}}
& 3 & 4 \\
\protect\hyperlink{unitree-sdk2-cpp-idl-hg-handstate}{\passthrough{\lstinline!HandState!}}
& 3 & 9 \\
\protect\hyperlink{unitree-sdk2-cpp-idl-hg-lowcmd}{\passthrough{\lstinline!LowCmd!}}
& 3 & 5 \\
\protect\hyperlink{unitree-sdk2-cpp-idl-hg-lowstate}{\passthrough{\lstinline!LowState!}}
& 3 & 9 \\
\protect\hyperlink{unitree-sdk2-cpp-idl-hg-mainboardstate}{\passthrough{\lstinline!MainBoardState!}}
& 3 & 4 \\
\protect\hyperlink{unitree-sdk2-cpp-idl-hg-sportmodestate}{\passthrough{\lstinline!SportModeState!}}
& 3 & 4 \\
\bottomrule
\end{longtable}

\hypertarget{unitree-sdk2-cpp-idl-hg-agvbmsstate}{%
\subsubsection{\texorpdfstring{\texttt{unitree\_sdk2\_cpp.idl.hg.AgvBmsState}}{unitree\_sdk2\_cpp.idl.hg.AgvBmsState}}\label{unitree-sdk2-cpp-idl-hg-agvbmsstate}}

可由 Python 读写的 DDS/IDL 消息值类型。字段采用复制语义。

\textbf{导入}

\begin{lstlisting}[language=Python]
from unitree_sdk2_cpp.idl.hg import AgvBmsState
\end{lstlisting}

\textbf{方法索引}

\begin{longtable}[]{@{}
  >{\raggedright\arraybackslash}p{(\columnwidth - 6\tabcolsep) * \real{0.18}}
  >{\raggedright\arraybackslash}p{(\columnwidth - 6\tabcolsep) * \real{0.24}}
  >{\raggedright\arraybackslash}p{(\columnwidth - 6\tabcolsep) * \real{0.18}}
  >{\raggedright\arraybackslash}p{(\columnwidth - 6\tabcolsep) * \real{0.40}}@{}}
\toprule
方法 & Python 签名 & 状态 & 安全分类 \\
\midrule
\endhead
\protect\hyperlink{unitree-sdk2-cpp-idl-hg-agvbmsstate-dunder-init-1}{\passthrough{\lstinline!\_\_init\_\_!}}
& \passthrough{\lstinline!def \_\_init\_\_(self) -> None!} &
\passthrough{\lstinline!AVAILABLE!} &
\passthrough{\lstinline!VALUE\_TYPE!} \\
\protect\hyperlink{unitree-sdk2-cpp-idl-hg-agvbmsstate-dunder-eq-1}{\passthrough{\lstinline!\_\_eq\_\_!}}
& \passthrough{\lstinline!def \_\_eq\_\_(self, other: object) -> bool!}
& \passthrough{\lstinline!AVAILABLE!} &
\passthrough{\lstinline!VALUE\_TYPE!} \\
\protect\hyperlink{unitree-sdk2-cpp-idl-hg-agvbmsstate-dunder-ne-1}{\passthrough{\lstinline!\_\_ne\_\_!}}
& \passthrough{\lstinline!def \_\_ne\_\_(self, other: object) -> bool!}
& \passthrough{\lstinline!AVAILABLE!} &
\passthrough{\lstinline!VALUE\_TYPE!} \\
\bottomrule
\end{longtable}

\textbf{属性索引}

\begin{longtable}[]{@{}
  >{\raggedright\arraybackslash}p{(\columnwidth - 6\tabcolsep) * \real{0.18}}
  >{\raggedright\arraybackslash}p{(\columnwidth - 6\tabcolsep) * \real{0.24}}
  >{\raggedright\arraybackslash}p{(\columnwidth - 6\tabcolsep) * \real{0.18}}
  >{\raggedright\arraybackslash}p{(\columnwidth - 6\tabcolsep) * \real{0.40}}@{}}
\toprule
属性 & 读取类型 & 写入类型 & 状态 \\
\midrule
\endhead
\protect\hyperlink{unitree-sdk2-cpp-idl-hg-agvbmsstate-software-version}{\passthrough{\lstinline!software\_version!}}
& \passthrough{\lstinline!str!} & \passthrough{\lstinline!str!} &
\passthrough{\lstinline!AVAILABLE!} \\
\protect\hyperlink{unitree-sdk2-cpp-idl-hg-agvbmsstate-battery-percentage}{\passthrough{\lstinline!battery\_percentage!}}
& \passthrough{\lstinline!int!} & \passthrough{\lstinline!int!} &
\passthrough{\lstinline!AVAILABLE!} \\
\protect\hyperlink{unitree-sdk2-cpp-idl-hg-agvbmsstate-current}{\passthrough{\lstinline!current!}}
& \passthrough{\lstinline!int!} & \passthrough{\lstinline!int!} &
\passthrough{\lstinline!AVAILABLE!} \\
\protect\hyperlink{unitree-sdk2-cpp-idl-hg-agvbmsstate-temperature}{\passthrough{\lstinline!temperature!}}
& \passthrough{\lstinline!list[int]!} &
\passthrough{\lstinline!Sequence[int]!} &
\passthrough{\lstinline!AVAILABLE!} \\
\protect\hyperlink{unitree-sdk2-cpp-idl-hg-agvbmsstate-docking-status}{\passthrough{\lstinline!docking\_status!}}
& \passthrough{\lstinline!str!} & \passthrough{\lstinline!str!} &
\passthrough{\lstinline!AVAILABLE!} \\
\protect\hyperlink{unitree-sdk2-cpp-idl-hg-agvbmsstate-is-charging}{\passthrough{\lstinline!is\_charging!}}
& \passthrough{\lstinline!bool!} & \passthrough{\lstinline!bool!} &
\passthrough{\lstinline!AVAILABLE!} \\
\protect\hyperlink{unitree-sdk2-cpp-idl-hg-agvbmsstate-is-dc-connected}{\passthrough{\lstinline!is\_dc\_connected!}}
& \passthrough{\lstinline!bool!} & \passthrough{\lstinline!bool!} &
\passthrough{\lstinline!AVAILABLE!} \\
\bottomrule
\end{longtable}

\hypertarget{unitree-sdk2-cpp-idl-hg-agvbmsstate-dunder-init-1}{%
\paragraph{\texorpdfstring{\texttt{unitree\_sdk2\_cpp.idl.hg.AgvBmsState.\_\_init\_\_}}{unitree\_sdk2\_cpp.idl.hg.AgvBmsState.\_\_init\_\_}}\label{unitree-sdk2-cpp-idl-hg-agvbmsstate-dunder-init-1}}

初始化 \passthrough{\lstinline!AgvBmsState!}
实例。是否能实际构造取决于下方可用性状态。

\textbf{签名}

\begin{lstlisting}[language=Python]
def __init__(self) -> None
\end{lstlisting}

\textbf{可用性与安全性}

\textbf{\passthrough{\lstinline!AVAILABLE!} /
\passthrough{\lstinline!VALUE\_TYPE!}}：当前绑定源码已实现该消息值操作。

\textbf{参数}

无显式参数。实例方法中的 \passthrough{\lstinline!self!} 由 Python
自动传入。

\textbf{返回值}

\begin{longtable}[]{@{}
  >{\raggedright\arraybackslash}p{(\columnwidth - 4\tabcolsep) * \real{0.18}}
  >{\raggedright\arraybackslash}p{(\columnwidth - 4\tabcolsep) * \real{0.20}}
  >{\raggedright\arraybackslash}p{(\columnwidth - 4\tabcolsep) * \real{0.62}}@{}}
\toprule
位置 & 类型 & 含义 \\
\midrule
\endhead
无 & \passthrough{\lstinline!None!} &
\passthrough{\lstinline!\_\_init\_\_!}
本身不返回值；调用类对象时，在构造函数可用的前提下得到该类实例。 \\
\bottomrule
\end{longtable}

\textbf{对应 C++}

\begin{itemize}
\tightlist
\item
  类：\passthrough{\lstinline!unitree\_hg::msg::dds\_::AgvBmsState\_!}
\item
  签名：\passthrough{\lstinline!AgvBmsState\_()!}
\item
  绑定策略：\passthrough{\lstinline!未单独标注!}
\item
  声明位置：\passthrough{\lstinline!include/unitree/idl/hg/AgvBmsState\_.hpp:34!}
\end{itemize}

\textbf{说明}

\begin{itemize}
\tightlist
\item
  示例是局部调用片段；\passthrough{\lstinline!obj!}
  和参数变量表示已经按上层业务校验并准备好的对象。
\end{itemize}

\textbf{用法}

\begin{lstlisting}[language=Python]
value = AgvBmsState()
\end{lstlisting}

\hypertarget{unitree-sdk2-cpp-idl-hg-agvbmsstate-dunder-eq-1}{%
\paragraph{\texorpdfstring{\texttt{unitree\_sdk2\_cpp.idl.hg.AgvBmsState.\_\_eq\_\_}}{unitree\_sdk2\_cpp.idl.hg.AgvBmsState.\_\_eq\_\_}}\label{unitree-sdk2-cpp-idl-hg-agvbmsstate-dunder-eq-1}}

按底层 IDL 消息内容比较两个对象是否相等。

\textbf{签名}

\begin{lstlisting}[language=Python]
def __eq__(self, other: object) -> bool
\end{lstlisting}

\textbf{可用性与安全性}

\textbf{\passthrough{\lstinline!AVAILABLE!} /
\passthrough{\lstinline!VALUE\_TYPE!}}：当前绑定源码已实现该消息值操作。

\textbf{参数}

\begin{longtable}[]{@{}
  >{\raggedright\arraybackslash}p{(\columnwidth - 6\tabcolsep) * \real{0.18}}
  >{\raggedright\arraybackslash}p{(\columnwidth - 6\tabcolsep) * \real{0.24}}
  >{\raggedright\arraybackslash}p{(\columnwidth - 6\tabcolsep) * \real{0.18}}
  >{\raggedright\arraybackslash}p{(\columnwidth - 6\tabcolsep) * \real{0.40}}@{}}
\toprule
名称 & Python 类型 & 默认值 & 含义 \\
\midrule
\endhead
\passthrough{\lstinline!other!} & \passthrough{\lstinline!object!} &
必填 & 用于比较的另一个 Python 对象。类型不兼容时比较结果通常为
\passthrough{\lstinline!False!}。 \\
\bottomrule
\end{longtable}

\textbf{返回值}

\begin{longtable}[]{@{}
  >{\raggedright\arraybackslash}p{(\columnwidth - 4\tabcolsep) * \real{0.18}}
  >{\raggedright\arraybackslash}p{(\columnwidth - 4\tabcolsep) * \real{0.20}}
  >{\raggedright\arraybackslash}p{(\columnwidth - 4\tabcolsep) * \real{0.62}}@{}}
\toprule
位置 & 类型 & 含义 \\
\midrule
\endhead
返回值 & \passthrough{\lstinline!bool!} & 返回
\passthrough{\lstinline!bool!}。更精确的业务含义见该方法用途、C++
签名和目标型号协议。 \\
\bottomrule
\end{longtable}

\textbf{对应 C++}

\begin{itemize}
\tightlist
\item
  类：\passthrough{\lstinline!unitree\_hg::msg::dds\_::AgvBmsState\_!}
\item
  签名：\passthrough{\lstinline!operator==(const value \&) const!}
\item
  绑定策略：\passthrough{\lstinline!未单独标注!}
\item
  声明位置：由 pybind11 手工绑定或生成注册表提供；manifest
  未记录独立头文件行号。
\end{itemize}

\textbf{说明}

\begin{itemize}
\tightlist
\item
  示例是局部调用片段；\passthrough{\lstinline!obj!}
  和参数变量表示已经按上层业务校验并准备好的对象。
\end{itemize}

\textbf{用法}

\begin{lstlisting}[language=Python]
same = left == right
\end{lstlisting}

\hypertarget{unitree-sdk2-cpp-idl-hg-agvbmsstate-dunder-ne-1}{%
\paragraph{\texorpdfstring{\texttt{unitree\_sdk2\_cpp.idl.hg.AgvBmsState.\_\_ne\_\_}}{unitree\_sdk2\_cpp.idl.hg.AgvBmsState.\_\_ne\_\_}}\label{unitree-sdk2-cpp-idl-hg-agvbmsstate-dunder-ne-1}}

按底层 IDL 消息内容比较两个对象是否不相等。

\textbf{签名}

\begin{lstlisting}[language=Python]
def __ne__(self, other: object) -> bool
\end{lstlisting}

\textbf{可用性与安全性}

\textbf{\passthrough{\lstinline!AVAILABLE!} /
\passthrough{\lstinline!VALUE\_TYPE!}}：当前绑定源码已实现该消息值操作。

\textbf{参数}

\begin{longtable}[]{@{}
  >{\raggedright\arraybackslash}p{(\columnwidth - 6\tabcolsep) * \real{0.18}}
  >{\raggedright\arraybackslash}p{(\columnwidth - 6\tabcolsep) * \real{0.24}}
  >{\raggedright\arraybackslash}p{(\columnwidth - 6\tabcolsep) * \real{0.18}}
  >{\raggedright\arraybackslash}p{(\columnwidth - 6\tabcolsep) * \real{0.40}}@{}}
\toprule
名称 & Python 类型 & 默认值 & 含义 \\
\midrule
\endhead
\passthrough{\lstinline!other!} & \passthrough{\lstinline!object!} &
必填 & 用于比较的另一个 Python 对象。类型不兼容时比较结果通常为
\passthrough{\lstinline!False!}。 \\
\bottomrule
\end{longtable}

\textbf{返回值}

\begin{longtable}[]{@{}
  >{\raggedright\arraybackslash}p{(\columnwidth - 4\tabcolsep) * \real{0.18}}
  >{\raggedright\arraybackslash}p{(\columnwidth - 4\tabcolsep) * \real{0.20}}
  >{\raggedright\arraybackslash}p{(\columnwidth - 4\tabcolsep) * \real{0.62}}@{}}
\toprule
位置 & 类型 & 含义 \\
\midrule
\endhead
返回值 & \passthrough{\lstinline!bool!} & 返回
\passthrough{\lstinline!bool!}。更精确的业务含义见该方法用途、C++
签名和目标型号协议。 \\
\bottomrule
\end{longtable}

\textbf{对应 C++}

\begin{itemize}
\tightlist
\item
  类：\passthrough{\lstinline!unitree\_hg::msg::dds\_::AgvBmsState\_!}
\item
  签名：\passthrough{\lstinline"operator!=(const value \&) const"}
\item
  绑定策略：\passthrough{\lstinline!未单独标注!}
\item
  声明位置：由 pybind11 手工绑定或生成注册表提供；manifest
  未记录独立头文件行号。
\end{itemize}

\textbf{说明}

\begin{itemize}
\tightlist
\item
  示例是局部调用片段；\passthrough{\lstinline!obj!}
  和参数变量表示已经按上层业务校验并准备好的对象。
\end{itemize}

\textbf{用法}

\begin{lstlisting}[language=Python]
different = left != right
\end{lstlisting}

\hypertarget{unitree-sdk2-cpp-idl-hg-agvbmsstate-software-version}{%
\paragraph{\texorpdfstring{\texttt{unitree\_sdk2\_cpp.idl.hg.AgvBmsState.software\_version}}{unitree\_sdk2\_cpp.idl.hg.AgvBmsState.software\_version}}\label{unitree-sdk2-cpp-idl-hg-agvbmsstate-software-version}}

底层 \passthrough{\lstinline!unitree\_hg::msg::dds\_::AgvBmsState\_!} 的
\passthrough{\lstinline!software\_version!} 字段。类型签名只说明 Python
表示；单位、数值范围、数组固定长度和枚举含义需查对应 IDL 头文件。
底层为字符串；长度、编码和允许值由具体协议决定。

\textbf{签名}

\begin{lstlisting}[language=Python]
@property
def software_version(self) -> str

@software_version.setter
def software_version(self, value: str) -> None
\end{lstlisting}

\textbf{可用性与安全性}

\textbf{\passthrough{\lstinline!AVAILABLE!} /
\passthrough{\lstinline!VALUE\_TYPE!}}：当前绑定源码已实现该消息值操作。

\textbf{参数}

\begin{longtable}[]{@{}
  >{\raggedright\arraybackslash}p{(\columnwidth - 4\tabcolsep) * \real{0.18}}
  >{\raggedright\arraybackslash}p{(\columnwidth - 4\tabcolsep) * \real{0.20}}
  >{\raggedright\arraybackslash}p{(\columnwidth - 4\tabcolsep) * \real{0.62}}@{}}
\toprule
名称 & Python 类型 & 含义 \\
\midrule
\endhead
\passthrough{\lstinline!value!} & \passthrough{\lstinline!str!} &
要写入或传递的新值，必须符合该参数的 Python 类型以及底层 C++
范围约束。 \\
\bottomrule
\end{longtable}

\textbf{返回值}

\begin{longtable}[]{@{}
  >{\raggedright\arraybackslash}p{(\columnwidth - 4\tabcolsep) * \real{0.18}}
  >{\raggedright\arraybackslash}p{(\columnwidth - 4\tabcolsep) * \real{0.20}}
  >{\raggedright\arraybackslash}p{(\columnwidth - 4\tabcolsep) * \real{0.62}}@{}}
\toprule
位置 & 类型 & 含义 \\
\midrule
\endhead
getter 返回值 & \passthrough{\lstinline!str!} & 返回属性当前值。 \\
\bottomrule
\end{longtable}

\textbf{对应 C++}

\begin{itemize}
\tightlist
\item
  类：\passthrough{\lstinline!unitree\_hg::msg::dds\_::AgvBmsState\_!}
\item
  字段类型：\passthrough{\lstinline!std::string!}
\item
  声明位置：\passthrough{\lstinline!include/unitree/idl/hg/AgvBmsState\_.hpp:25!}
\end{itemize}

\textbf{用法}

\begin{lstlisting}[language=Python]
current_value = obj.software_version
obj.software_version = new_value
\end{lstlisting}

\hypertarget{unitree-sdk2-cpp-idl-hg-agvbmsstate-battery-percentage}{%
\paragraph{\texorpdfstring{\texttt{unitree\_sdk2\_cpp.idl.hg.AgvBmsState.battery\_percentage}}{unitree\_sdk2\_cpp.idl.hg.AgvBmsState.battery\_percentage}}\label{unitree-sdk2-cpp-idl-hg-agvbmsstate-battery-percentage}}

底层 \passthrough{\lstinline!unitree\_hg::msg::dds\_::AgvBmsState\_!} 的
\passthrough{\lstinline!battery\_percentage!} 字段。类型签名只说明
Python 表示；单位、数值范围、数组固定长度和枚举含义需查对应 IDL 头文件。
取值范围为 0 到 255。

\textbf{签名}

\begin{lstlisting}[language=Python]
@property
def battery_percentage(self) -> int

@battery_percentage.setter
def battery_percentage(self, value: int) -> None
\end{lstlisting}

\textbf{可用性与安全性}

\textbf{\passthrough{\lstinline!AVAILABLE!} /
\passthrough{\lstinline!VALUE\_TYPE!}}：当前绑定源码已实现该消息值操作。

\textbf{参数}

\begin{longtable}[]{@{}
  >{\raggedright\arraybackslash}p{(\columnwidth - 4\tabcolsep) * \real{0.18}}
  >{\raggedright\arraybackslash}p{(\columnwidth - 4\tabcolsep) * \real{0.20}}
  >{\raggedright\arraybackslash}p{(\columnwidth - 4\tabcolsep) * \real{0.62}}@{}}
\toprule
名称 & Python 类型 & 含义 \\
\midrule
\endhead
\passthrough{\lstinline!value!} & \passthrough{\lstinline!int!} &
要写入或传递的新值，必须符合该参数的 Python 类型以及底层 C++
范围约束。 \\
\bottomrule
\end{longtable}

\textbf{返回值}

\begin{longtable}[]{@{}
  >{\raggedright\arraybackslash}p{(\columnwidth - 4\tabcolsep) * \real{0.18}}
  >{\raggedright\arraybackslash}p{(\columnwidth - 4\tabcolsep) * \real{0.20}}
  >{\raggedright\arraybackslash}p{(\columnwidth - 4\tabcolsep) * \real{0.62}}@{}}
\toprule
位置 & 类型 & 含义 \\
\midrule
\endhead
getter 返回值 & \passthrough{\lstinline!int!} & 返回属性当前值。 \\
\bottomrule
\end{longtable}

\textbf{对应 C++}

\begin{itemize}
\tightlist
\item
  类：\passthrough{\lstinline!unitree\_hg::msg::dds\_::AgvBmsState\_!}
\item
  字段类型：\passthrough{\lstinline!uint8\_t!}
\item
  声明位置：\passthrough{\lstinline!include/unitree/idl/hg/AgvBmsState\_.hpp:26!}
\end{itemize}

\textbf{用法}

\begin{lstlisting}[language=Python]
current_value = obj.battery_percentage
obj.battery_percentage = new_value
\end{lstlisting}

\hypertarget{unitree-sdk2-cpp-idl-hg-agvbmsstate-current}{%
\paragraph{\texorpdfstring{\texttt{unitree\_sdk2\_cpp.idl.hg.AgvBmsState.current}}{unitree\_sdk2\_cpp.idl.hg.AgvBmsState.current}}\label{unitree-sdk2-cpp-idl-hg-agvbmsstate-current}}

底层 \passthrough{\lstinline!unitree\_hg::msg::dds\_::AgvBmsState\_!} 的
\passthrough{\lstinline!current!} 字段。类型签名只说明 Python
表示；单位、数值范围、数组固定长度和枚举含义需查对应 IDL 头文件。
取值范围为 -2147483648 到 2147483647。

\textbf{签名}

\begin{lstlisting}[language=Python]
@property
def current(self) -> int

@current.setter
def current(self, value: int) -> None
\end{lstlisting}

\textbf{可用性与安全性}

\textbf{\passthrough{\lstinline!AVAILABLE!} /
\passthrough{\lstinline!VALUE\_TYPE!}}：当前绑定源码已实现该消息值操作。

\textbf{参数}

\begin{longtable}[]{@{}
  >{\raggedright\arraybackslash}p{(\columnwidth - 4\tabcolsep) * \real{0.18}}
  >{\raggedright\arraybackslash}p{(\columnwidth - 4\tabcolsep) * \real{0.20}}
  >{\raggedright\arraybackslash}p{(\columnwidth - 4\tabcolsep) * \real{0.62}}@{}}
\toprule
名称 & Python 类型 & 含义 \\
\midrule
\endhead
\passthrough{\lstinline!value!} & \passthrough{\lstinline!int!} &
要写入或传递的新值，必须符合该参数的 Python 类型以及底层 C++
范围约束。 \\
\bottomrule
\end{longtable}

\textbf{返回值}

\begin{longtable}[]{@{}
  >{\raggedright\arraybackslash}p{(\columnwidth - 4\tabcolsep) * \real{0.18}}
  >{\raggedright\arraybackslash}p{(\columnwidth - 4\tabcolsep) * \real{0.20}}
  >{\raggedright\arraybackslash}p{(\columnwidth - 4\tabcolsep) * \real{0.62}}@{}}
\toprule
位置 & 类型 & 含义 \\
\midrule
\endhead
getter 返回值 & \passthrough{\lstinline!int!} & 返回属性当前值。 \\
\bottomrule
\end{longtable}

\textbf{对应 C++}

\begin{itemize}
\tightlist
\item
  类：\passthrough{\lstinline!unitree\_hg::msg::dds\_::AgvBmsState\_!}
\item
  字段类型：\passthrough{\lstinline!int32\_t!}
\item
  声明位置：\passthrough{\lstinline!include/unitree/idl/hg/AgvBmsState\_.hpp:27!}
\end{itemize}

\textbf{用法}

\begin{lstlisting}[language=Python]
current_value = obj.current
obj.current = new_value
\end{lstlisting}

\hypertarget{unitree-sdk2-cpp-idl-hg-agvbmsstate-temperature}{%
\paragraph{\texorpdfstring{\texttt{unitree\_sdk2\_cpp.idl.hg.AgvBmsState.temperature}}{unitree\_sdk2\_cpp.idl.hg.AgvBmsState.temperature}}\label{unitree-sdk2-cpp-idl-hg-agvbmsstate-temperature}}

底层 \passthrough{\lstinline!unitree\_hg::msg::dds\_::AgvBmsState\_!} 的
\passthrough{\lstinline!temperature!} 字段。类型签名只说明 Python
表示；单位、数值范围、数组固定长度和枚举含义需查对应 IDL 头文件。
底层是固定长度数组，写入序列必须正好包含 3 个元素；元素约束：取值范围为
-32768 到 32767。

\textbf{签名}

\begin{lstlisting}[language=Python]
@property
def temperature(self) -> list[int]

@temperature.setter
def temperature(self, value: Sequence[int]) -> None
\end{lstlisting}

\textbf{可用性与安全性}

\textbf{\passthrough{\lstinline!AVAILABLE!} /
\passthrough{\lstinline!VALUE\_TYPE!}}：当前绑定源码已实现该消息值操作。

\textbf{参数}

\begin{longtable}[]{@{}
  >{\raggedright\arraybackslash}p{(\columnwidth - 4\tabcolsep) * \real{0.18}}
  >{\raggedright\arraybackslash}p{(\columnwidth - 4\tabcolsep) * \real{0.20}}
  >{\raggedright\arraybackslash}p{(\columnwidth - 4\tabcolsep) * \real{0.62}}@{}}
\toprule
名称 & Python 类型 & 含义 \\
\midrule
\endhead
\passthrough{\lstinline!value!} &
\passthrough{\lstinline!Sequence[int]!} &
要写入或传递的新值，必须符合该参数的 Python 类型以及底层 C++
范围约束。 \\
\bottomrule
\end{longtable}

\textbf{返回值}

\begin{longtable}[]{@{}
  >{\raggedright\arraybackslash}p{(\columnwidth - 4\tabcolsep) * \real{0.18}}
  >{\raggedright\arraybackslash}p{(\columnwidth - 4\tabcolsep) * \real{0.20}}
  >{\raggedright\arraybackslash}p{(\columnwidth - 4\tabcolsep) * \real{0.62}}@{}}
\toprule
位置 & 类型 & 含义 \\
\midrule
\endhead
getter 返回值 & \passthrough{\lstinline!list[int]!} &
返回属性当前值。数组、vector 和嵌套消息采用复制语义。 \\
\bottomrule
\end{longtable}

\textbf{对应 C++}

\begin{itemize}
\tightlist
\item
  类：\passthrough{\lstinline!unitree\_hg::msg::dds\_::AgvBmsState\_!}
\item
  字段类型：\passthrough{\lstinline!std::array<int16\_t, 3>!}
\item
  声明位置：\passthrough{\lstinline!include/unitree/idl/hg/AgvBmsState\_.hpp:28!}
\end{itemize}

\textbf{用法}

\begin{lstlisting}[language=Python]
current_value = obj.temperature
obj.temperature = new_value
\end{lstlisting}

\begin{quote}
{[}!TIP{]} 读取容器或嵌套值后应遵循``读取、修改、重新赋值''。只修改
getter 返回的副本不会可靠地更新原 C++ 消息。
\end{quote}

\hypertarget{unitree-sdk2-cpp-idl-hg-agvbmsstate-docking-status}{%
\paragraph{\texorpdfstring{\texttt{unitree\_sdk2\_cpp.idl.hg.AgvBmsState.docking\_status}}{unitree\_sdk2\_cpp.idl.hg.AgvBmsState.docking\_status}}\label{unitree-sdk2-cpp-idl-hg-agvbmsstate-docking-status}}

底层 \passthrough{\lstinline!unitree\_hg::msg::dds\_::AgvBmsState\_!} 的
\passthrough{\lstinline!docking\_status!} 字段。类型签名只说明 Python
表示；单位、数值范围、数组固定长度和枚举含义需查对应 IDL 头文件。
底层为字符串；长度、编码和允许值由具体协议决定。

\textbf{签名}

\begin{lstlisting}[language=Python]
@property
def docking_status(self) -> str

@docking_status.setter
def docking_status(self, value: str) -> None
\end{lstlisting}

\textbf{可用性与安全性}

\textbf{\passthrough{\lstinline!AVAILABLE!} /
\passthrough{\lstinline!VALUE\_TYPE!}}：当前绑定源码已实现该消息值操作。

\textbf{参数}

\begin{longtable}[]{@{}
  >{\raggedright\arraybackslash}p{(\columnwidth - 4\tabcolsep) * \real{0.18}}
  >{\raggedright\arraybackslash}p{(\columnwidth - 4\tabcolsep) * \real{0.20}}
  >{\raggedright\arraybackslash}p{(\columnwidth - 4\tabcolsep) * \real{0.62}}@{}}
\toprule
名称 & Python 类型 & 含义 \\
\midrule
\endhead
\passthrough{\lstinline!value!} & \passthrough{\lstinline!str!} &
要写入或传递的新值，必须符合该参数的 Python 类型以及底层 C++
范围约束。 \\
\bottomrule
\end{longtable}

\textbf{返回值}

\begin{longtable}[]{@{}
  >{\raggedright\arraybackslash}p{(\columnwidth - 4\tabcolsep) * \real{0.18}}
  >{\raggedright\arraybackslash}p{(\columnwidth - 4\tabcolsep) * \real{0.20}}
  >{\raggedright\arraybackslash}p{(\columnwidth - 4\tabcolsep) * \real{0.62}}@{}}
\toprule
位置 & 类型 & 含义 \\
\midrule
\endhead
getter 返回值 & \passthrough{\lstinline!str!} & 返回属性当前值。 \\
\bottomrule
\end{longtable}

\textbf{对应 C++}

\begin{itemize}
\tightlist
\item
  类：\passthrough{\lstinline!unitree\_hg::msg::dds\_::AgvBmsState\_!}
\item
  字段类型：\passthrough{\lstinline!std::string!}
\item
  声明位置：\passthrough{\lstinline!include/unitree/idl/hg/AgvBmsState\_.hpp:29!}
\end{itemize}

\textbf{用法}

\begin{lstlisting}[language=Python]
current_value = obj.docking_status
obj.docking_status = new_value
\end{lstlisting}

\hypertarget{unitree-sdk2-cpp-idl-hg-agvbmsstate-is-charging}{%
\paragraph{\texorpdfstring{\texttt{unitree\_sdk2\_cpp.idl.hg.AgvBmsState.is\_charging}}{unitree\_sdk2\_cpp.idl.hg.AgvBmsState.is\_charging}}\label{unitree-sdk2-cpp-idl-hg-agvbmsstate-is-charging}}

底层 \passthrough{\lstinline!unitree\_hg::msg::dds\_::AgvBmsState\_!} 的
\passthrough{\lstinline!is\_charging!} 字段。类型签名只说明 Python
表示；单位、数值范围、数组固定长度和枚举含义需查对应 IDL 头文件。
只接受布尔语义。

\textbf{签名}

\begin{lstlisting}[language=Python]
@property
def is_charging(self) -> bool

@is_charging.setter
def is_charging(self, value: bool) -> None
\end{lstlisting}

\textbf{可用性与安全性}

\textbf{\passthrough{\lstinline!AVAILABLE!} /
\passthrough{\lstinline!VALUE\_TYPE!}}：当前绑定源码已实现该消息值操作。

\textbf{参数}

\begin{longtable}[]{@{}
  >{\raggedright\arraybackslash}p{(\columnwidth - 4\tabcolsep) * \real{0.18}}
  >{\raggedright\arraybackslash}p{(\columnwidth - 4\tabcolsep) * \real{0.20}}
  >{\raggedright\arraybackslash}p{(\columnwidth - 4\tabcolsep) * \real{0.62}}@{}}
\toprule
名称 & Python 类型 & 含义 \\
\midrule
\endhead
\passthrough{\lstinline!value!} & \passthrough{\lstinline!bool!} &
要写入或传递的新值，必须符合该参数的 Python 类型以及底层 C++
范围约束。 \\
\bottomrule
\end{longtable}

\textbf{返回值}

\begin{longtable}[]{@{}
  >{\raggedright\arraybackslash}p{(\columnwidth - 4\tabcolsep) * \real{0.18}}
  >{\raggedright\arraybackslash}p{(\columnwidth - 4\tabcolsep) * \real{0.20}}
  >{\raggedright\arraybackslash}p{(\columnwidth - 4\tabcolsep) * \real{0.62}}@{}}
\toprule
位置 & 类型 & 含义 \\
\midrule
\endhead
getter 返回值 & \passthrough{\lstinline!bool!} & 返回属性当前值。 \\
\bottomrule
\end{longtable}

\textbf{对应 C++}

\begin{itemize}
\tightlist
\item
  类：\passthrough{\lstinline!unitree\_hg::msg::dds\_::AgvBmsState\_!}
\item
  字段类型：\passthrough{\lstinline!bool!}
\item
  声明位置：\passthrough{\lstinline!include/unitree/idl/hg/AgvBmsState\_.hpp:30!}
\end{itemize}

\textbf{用法}

\begin{lstlisting}[language=Python]
current_value = obj.is_charging
obj.is_charging = new_value
\end{lstlisting}

\hypertarget{unitree-sdk2-cpp-idl-hg-agvbmsstate-is-dc-connected}{%
\paragraph{\texorpdfstring{\texttt{unitree\_sdk2\_cpp.idl.hg.AgvBmsState.is\_dc\_connected}}{unitree\_sdk2\_cpp.idl.hg.AgvBmsState.is\_dc\_connected}}\label{unitree-sdk2-cpp-idl-hg-agvbmsstate-is-dc-connected}}

底层 \passthrough{\lstinline!unitree\_hg::msg::dds\_::AgvBmsState\_!} 的
\passthrough{\lstinline!is\_dc\_connected!} 字段。类型签名只说明 Python
表示；单位、数值范围、数组固定长度和枚举含义需查对应 IDL 头文件。
只接受布尔语义。

\textbf{签名}

\begin{lstlisting}[language=Python]
@property
def is_dc_connected(self) -> bool

@is_dc_connected.setter
def is_dc_connected(self, value: bool) -> None
\end{lstlisting}

\textbf{可用性与安全性}

\textbf{\passthrough{\lstinline!AVAILABLE!} /
\passthrough{\lstinline!VALUE\_TYPE!}}：当前绑定源码已实现该消息值操作。

\textbf{参数}

\begin{longtable}[]{@{}
  >{\raggedright\arraybackslash}p{(\columnwidth - 4\tabcolsep) * \real{0.18}}
  >{\raggedright\arraybackslash}p{(\columnwidth - 4\tabcolsep) * \real{0.20}}
  >{\raggedright\arraybackslash}p{(\columnwidth - 4\tabcolsep) * \real{0.62}}@{}}
\toprule
名称 & Python 类型 & 含义 \\
\midrule
\endhead
\passthrough{\lstinline!value!} & \passthrough{\lstinline!bool!} &
要写入或传递的新值，必须符合该参数的 Python 类型以及底层 C++
范围约束。 \\
\bottomrule
\end{longtable}

\textbf{返回值}

\begin{longtable}[]{@{}
  >{\raggedright\arraybackslash}p{(\columnwidth - 4\tabcolsep) * \real{0.18}}
  >{\raggedright\arraybackslash}p{(\columnwidth - 4\tabcolsep) * \real{0.20}}
  >{\raggedright\arraybackslash}p{(\columnwidth - 4\tabcolsep) * \real{0.62}}@{}}
\toprule
位置 & 类型 & 含义 \\
\midrule
\endhead
getter 返回值 & \passthrough{\lstinline!bool!} & 返回属性当前值。 \\
\bottomrule
\end{longtable}

\textbf{对应 C++}

\begin{itemize}
\tightlist
\item
  类：\passthrough{\lstinline!unitree\_hg::msg::dds\_::AgvBmsState\_!}
\item
  字段类型：\passthrough{\lstinline!bool!}
\item
  声明位置：\passthrough{\lstinline!include/unitree/idl/hg/AgvBmsState\_.hpp:31!}
\end{itemize}

\textbf{用法}

\begin{lstlisting}[language=Python]
current_value = obj.is_dc_connected
obj.is_dc_connected = new_value
\end{lstlisting}

\hypertarget{unitree-sdk2-cpp-idl-hg-bmscmd}{%
\subsubsection{\texorpdfstring{\texttt{unitree\_sdk2\_cpp.idl.hg.BmsCmd}}{unitree\_sdk2\_cpp.idl.hg.BmsCmd}}\label{unitree-sdk2-cpp-idl-hg-bmscmd}}

可由 Python 读写的 DDS/IDL 消息值类型。字段采用复制语义。

\textbf{导入}

\begin{lstlisting}[language=Python]
from unitree_sdk2_cpp.idl.hg import BmsCmd
\end{lstlisting}

\textbf{方法索引}

\begin{longtable}[]{@{}
  >{\raggedright\arraybackslash}p{(\columnwidth - 6\tabcolsep) * \real{0.18}}
  >{\raggedright\arraybackslash}p{(\columnwidth - 6\tabcolsep) * \real{0.24}}
  >{\raggedright\arraybackslash}p{(\columnwidth - 6\tabcolsep) * \real{0.18}}
  >{\raggedright\arraybackslash}p{(\columnwidth - 6\tabcolsep) * \real{0.40}}@{}}
\toprule
方法 & Python 签名 & 状态 & 安全分类 \\
\midrule
\endhead
\protect\hyperlink{unitree-sdk2-cpp-idl-hg-bmscmd-dunder-init-1}{\passthrough{\lstinline!\_\_init\_\_!}}
& \passthrough{\lstinline!def \_\_init\_\_(self) -> None!} &
\passthrough{\lstinline!AVAILABLE!} &
\passthrough{\lstinline!VALUE\_TYPE!} \\
\protect\hyperlink{unitree-sdk2-cpp-idl-hg-bmscmd-dunder-eq-1}{\passthrough{\lstinline!\_\_eq\_\_!}}
& \passthrough{\lstinline!def \_\_eq\_\_(self, other: object) -> bool!}
& \passthrough{\lstinline!AVAILABLE!} &
\passthrough{\lstinline!VALUE\_TYPE!} \\
\protect\hyperlink{unitree-sdk2-cpp-idl-hg-bmscmd-dunder-ne-1}{\passthrough{\lstinline!\_\_ne\_\_!}}
& \passthrough{\lstinline!def \_\_ne\_\_(self, other: object) -> bool!}
& \passthrough{\lstinline!AVAILABLE!} &
\passthrough{\lstinline!VALUE\_TYPE!} \\
\bottomrule
\end{longtable}

\textbf{属性索引}

\begin{longtable}[]{@{}
  >{\raggedright\arraybackslash}p{(\columnwidth - 6\tabcolsep) * \real{0.18}}
  >{\raggedright\arraybackslash}p{(\columnwidth - 6\tabcolsep) * \real{0.24}}
  >{\raggedright\arraybackslash}p{(\columnwidth - 6\tabcolsep) * \real{0.18}}
  >{\raggedright\arraybackslash}p{(\columnwidth - 6\tabcolsep) * \real{0.40}}@{}}
\toprule
属性 & 读取类型 & 写入类型 & 状态 \\
\midrule
\endhead
\protect\hyperlink{unitree-sdk2-cpp-idl-hg-bmscmd-cmd}{\passthrough{\lstinline!cmd!}}
& \passthrough{\lstinline!int!} & \passthrough{\lstinline!int!} &
\passthrough{\lstinline!AVAILABLE!} \\
\protect\hyperlink{unitree-sdk2-cpp-idl-hg-bmscmd-reserve}{\passthrough{\lstinline!reserve!}}
& \passthrough{\lstinline!list[int]!} &
\passthrough{\lstinline!Sequence[int]!} &
\passthrough{\lstinline!AVAILABLE!} \\
\bottomrule
\end{longtable}

\hypertarget{unitree-sdk2-cpp-idl-hg-bmscmd-dunder-init-1}{%
\paragraph{\texorpdfstring{\texttt{unitree\_sdk2\_cpp.idl.hg.BmsCmd.\_\_init\_\_}}{unitree\_sdk2\_cpp.idl.hg.BmsCmd.\_\_init\_\_}}\label{unitree-sdk2-cpp-idl-hg-bmscmd-dunder-init-1}}

初始化 \passthrough{\lstinline!BmsCmd!}
实例。是否能实际构造取决于下方可用性状态。

\textbf{签名}

\begin{lstlisting}[language=Python]
def __init__(self) -> None
\end{lstlisting}

\textbf{可用性与安全性}

\textbf{\passthrough{\lstinline!AVAILABLE!} /
\passthrough{\lstinline!VALUE\_TYPE!}}：当前绑定源码已实现该消息值操作。

\textbf{参数}

无显式参数。实例方法中的 \passthrough{\lstinline!self!} 由 Python
自动传入。

\textbf{返回值}

\begin{longtable}[]{@{}
  >{\raggedright\arraybackslash}p{(\columnwidth - 4\tabcolsep) * \real{0.18}}
  >{\raggedright\arraybackslash}p{(\columnwidth - 4\tabcolsep) * \real{0.20}}
  >{\raggedright\arraybackslash}p{(\columnwidth - 4\tabcolsep) * \real{0.62}}@{}}
\toprule
位置 & 类型 & 含义 \\
\midrule
\endhead
无 & \passthrough{\lstinline!None!} &
\passthrough{\lstinline!\_\_init\_\_!}
本身不返回值；调用类对象时，在构造函数可用的前提下得到该类实例。 \\
\bottomrule
\end{longtable}

\textbf{对应 C++}

\begin{itemize}
\tightlist
\item
  类：\passthrough{\lstinline!unitree\_hg::msg::dds\_::BmsCmd\_!}
\item
  签名：\passthrough{\lstinline!BmsCmd\_()!}
\item
  绑定策略：\passthrough{\lstinline!未单独标注!}
\item
  声明位置：\passthrough{\lstinline!include/unitree/idl/hg/BmsCmd\_.hpp:28!}
\end{itemize}

\textbf{说明}

\begin{itemize}
\tightlist
\item
  示例是局部调用片段；\passthrough{\lstinline!obj!}
  和参数变量表示已经按上层业务校验并准备好的对象。
\end{itemize}

\textbf{用法}

\begin{lstlisting}[language=Python]
value = BmsCmd()
\end{lstlisting}

\hypertarget{unitree-sdk2-cpp-idl-hg-bmscmd-dunder-eq-1}{%
\paragraph{\texorpdfstring{\texttt{unitree\_sdk2\_cpp.idl.hg.BmsCmd.\_\_eq\_\_}}{unitree\_sdk2\_cpp.idl.hg.BmsCmd.\_\_eq\_\_}}\label{unitree-sdk2-cpp-idl-hg-bmscmd-dunder-eq-1}}

按底层 IDL 消息内容比较两个对象是否相等。

\textbf{签名}

\begin{lstlisting}[language=Python]
def __eq__(self, other: object) -> bool
\end{lstlisting}

\textbf{可用性与安全性}

\textbf{\passthrough{\lstinline!AVAILABLE!} /
\passthrough{\lstinline!VALUE\_TYPE!}}：当前绑定源码已实现该消息值操作。

\textbf{参数}

\begin{longtable}[]{@{}
  >{\raggedright\arraybackslash}p{(\columnwidth - 6\tabcolsep) * \real{0.18}}
  >{\raggedright\arraybackslash}p{(\columnwidth - 6\tabcolsep) * \real{0.24}}
  >{\raggedright\arraybackslash}p{(\columnwidth - 6\tabcolsep) * \real{0.18}}
  >{\raggedright\arraybackslash}p{(\columnwidth - 6\tabcolsep) * \real{0.40}}@{}}
\toprule
名称 & Python 类型 & 默认值 & 含义 \\
\midrule
\endhead
\passthrough{\lstinline!other!} & \passthrough{\lstinline!object!} &
必填 & 用于比较的另一个 Python 对象。类型不兼容时比较结果通常为
\passthrough{\lstinline!False!}。 \\
\bottomrule
\end{longtable}

\textbf{返回值}

\begin{longtable}[]{@{}
  >{\raggedright\arraybackslash}p{(\columnwidth - 4\tabcolsep) * \real{0.18}}
  >{\raggedright\arraybackslash}p{(\columnwidth - 4\tabcolsep) * \real{0.20}}
  >{\raggedright\arraybackslash}p{(\columnwidth - 4\tabcolsep) * \real{0.62}}@{}}
\toprule
位置 & 类型 & 含义 \\
\midrule
\endhead
返回值 & \passthrough{\lstinline!bool!} & 返回
\passthrough{\lstinline!bool!}。更精确的业务含义见该方法用途、C++
签名和目标型号协议。 \\
\bottomrule
\end{longtable}

\textbf{对应 C++}

\begin{itemize}
\tightlist
\item
  类：\passthrough{\lstinline!unitree\_hg::msg::dds\_::BmsCmd\_!}
\item
  签名：\passthrough{\lstinline!operator==(const value \&) const!}
\item
  绑定策略：\passthrough{\lstinline!未单独标注!}
\item
  声明位置：由 pybind11 手工绑定或生成注册表提供；manifest
  未记录独立头文件行号。
\end{itemize}

\textbf{说明}

\begin{itemize}
\tightlist
\item
  示例是局部调用片段；\passthrough{\lstinline!obj!}
  和参数变量表示已经按上层业务校验并准备好的对象。
\end{itemize}

\textbf{用法}

\begin{lstlisting}[language=Python]
same = left == right
\end{lstlisting}

\hypertarget{unitree-sdk2-cpp-idl-hg-bmscmd-dunder-ne-1}{%
\paragraph{\texorpdfstring{\texttt{unitree\_sdk2\_cpp.idl.hg.BmsCmd.\_\_ne\_\_}}{unitree\_sdk2\_cpp.idl.hg.BmsCmd.\_\_ne\_\_}}\label{unitree-sdk2-cpp-idl-hg-bmscmd-dunder-ne-1}}

按底层 IDL 消息内容比较两个对象是否不相等。

\textbf{签名}

\begin{lstlisting}[language=Python]
def __ne__(self, other: object) -> bool
\end{lstlisting}

\textbf{可用性与安全性}

\textbf{\passthrough{\lstinline!AVAILABLE!} /
\passthrough{\lstinline!VALUE\_TYPE!}}：当前绑定源码已实现该消息值操作。

\textbf{参数}

\begin{longtable}[]{@{}
  >{\raggedright\arraybackslash}p{(\columnwidth - 6\tabcolsep) * \real{0.18}}
  >{\raggedright\arraybackslash}p{(\columnwidth - 6\tabcolsep) * \real{0.24}}
  >{\raggedright\arraybackslash}p{(\columnwidth - 6\tabcolsep) * \real{0.18}}
  >{\raggedright\arraybackslash}p{(\columnwidth - 6\tabcolsep) * \real{0.40}}@{}}
\toprule
名称 & Python 类型 & 默认值 & 含义 \\
\midrule
\endhead
\passthrough{\lstinline!other!} & \passthrough{\lstinline!object!} &
必填 & 用于比较的另一个 Python 对象。类型不兼容时比较结果通常为
\passthrough{\lstinline!False!}。 \\
\bottomrule
\end{longtable}

\textbf{返回值}

\begin{longtable}[]{@{}
  >{\raggedright\arraybackslash}p{(\columnwidth - 4\tabcolsep) * \real{0.18}}
  >{\raggedright\arraybackslash}p{(\columnwidth - 4\tabcolsep) * \real{0.20}}
  >{\raggedright\arraybackslash}p{(\columnwidth - 4\tabcolsep) * \real{0.62}}@{}}
\toprule
位置 & 类型 & 含义 \\
\midrule
\endhead
返回值 & \passthrough{\lstinline!bool!} & 返回
\passthrough{\lstinline!bool!}。更精确的业务含义见该方法用途、C++
签名和目标型号协议。 \\
\bottomrule
\end{longtable}

\textbf{对应 C++}

\begin{itemize}
\tightlist
\item
  类：\passthrough{\lstinline!unitree\_hg::msg::dds\_::BmsCmd\_!}
\item
  签名：\passthrough{\lstinline"operator!=(const value \&) const"}
\item
  绑定策略：\passthrough{\lstinline!未单独标注!}
\item
  声明位置：由 pybind11 手工绑定或生成注册表提供；manifest
  未记录独立头文件行号。
\end{itemize}

\textbf{说明}

\begin{itemize}
\tightlist
\item
  示例是局部调用片段；\passthrough{\lstinline!obj!}
  和参数变量表示已经按上层业务校验并准备好的对象。
\end{itemize}

\textbf{用法}

\begin{lstlisting}[language=Python]
different = left != right
\end{lstlisting}

\hypertarget{unitree-sdk2-cpp-idl-hg-bmscmd-cmd}{%
\paragraph{\texorpdfstring{\texttt{unitree\_sdk2\_cpp.idl.hg.BmsCmd.cmd}}{unitree\_sdk2\_cpp.idl.hg.BmsCmd.cmd}}\label{unitree-sdk2-cpp-idl-hg-bmscmd-cmd}}

底层 \passthrough{\lstinline!unitree\_hg::msg::dds\_::BmsCmd\_!} 的
\passthrough{\lstinline!cmd!} 字段。类型签名只说明 Python
表示；单位、数值范围、数组固定长度和枚举含义需查对应 IDL 头文件。
取值范围为 0 到 255。

\textbf{签名}

\begin{lstlisting}[language=Python]
@property
def cmd(self) -> int

@cmd.setter
def cmd(self, value: int) -> None
\end{lstlisting}

\textbf{可用性与安全性}

\textbf{\passthrough{\lstinline!AVAILABLE!} /
\passthrough{\lstinline!VALUE\_TYPE!}}：当前绑定源码已实现该消息值操作。

\textbf{参数}

\begin{longtable}[]{@{}
  >{\raggedright\arraybackslash}p{(\columnwidth - 4\tabcolsep) * \real{0.18}}
  >{\raggedright\arraybackslash}p{(\columnwidth - 4\tabcolsep) * \real{0.20}}
  >{\raggedright\arraybackslash}p{(\columnwidth - 4\tabcolsep) * \real{0.62}}@{}}
\toprule
名称 & Python 类型 & 含义 \\
\midrule
\endhead
\passthrough{\lstinline!value!} & \passthrough{\lstinline!int!} &
要写入或传递的新值，必须符合该参数的 Python 类型以及底层 C++
范围约束。 \\
\bottomrule
\end{longtable}

\textbf{返回值}

\begin{longtable}[]{@{}
  >{\raggedright\arraybackslash}p{(\columnwidth - 4\tabcolsep) * \real{0.18}}
  >{\raggedright\arraybackslash}p{(\columnwidth - 4\tabcolsep) * \real{0.20}}
  >{\raggedright\arraybackslash}p{(\columnwidth - 4\tabcolsep) * \real{0.62}}@{}}
\toprule
位置 & 类型 & 含义 \\
\midrule
\endhead
getter 返回值 & \passthrough{\lstinline!int!} & 返回属性当前值。 \\
\bottomrule
\end{longtable}

\textbf{对应 C++}

\begin{itemize}
\tightlist
\item
  类：\passthrough{\lstinline!unitree\_hg::msg::dds\_::BmsCmd\_!}
\item
  字段类型：\passthrough{\lstinline!uint8\_t!}
\item
  声明位置：\passthrough{\lstinline!include/unitree/idl/hg/BmsCmd\_.hpp:24!}
\end{itemize}

\textbf{用法}

\begin{lstlisting}[language=Python]
current_value = obj.cmd
obj.cmd = new_value
\end{lstlisting}

\hypertarget{unitree-sdk2-cpp-idl-hg-bmscmd-reserve}{%
\paragraph{\texorpdfstring{\texttt{unitree\_sdk2\_cpp.idl.hg.BmsCmd.reserve}}{unitree\_sdk2\_cpp.idl.hg.BmsCmd.reserve}}\label{unitree-sdk2-cpp-idl-hg-bmscmd-reserve}}

底层 \passthrough{\lstinline!unitree\_hg::msg::dds\_::BmsCmd\_!} 的
\passthrough{\lstinline!reserve!} 字段。类型签名只说明 Python
表示；单位、数值范围、数组固定长度和枚举含义需查对应 IDL 头文件。
底层是固定长度数组，写入序列必须正好包含 40 个元素；元素约束：取值范围为
0 到 255。

\textbf{签名}

\begin{lstlisting}[language=Python]
@property
def reserve(self) -> list[int]

@reserve.setter
def reserve(self, value: Sequence[int]) -> None
\end{lstlisting}

\textbf{可用性与安全性}

\textbf{\passthrough{\lstinline!AVAILABLE!} /
\passthrough{\lstinline!VALUE\_TYPE!}}：当前绑定源码已实现该消息值操作。

\textbf{参数}

\begin{longtable}[]{@{}
  >{\raggedright\arraybackslash}p{(\columnwidth - 4\tabcolsep) * \real{0.18}}
  >{\raggedright\arraybackslash}p{(\columnwidth - 4\tabcolsep) * \real{0.20}}
  >{\raggedright\arraybackslash}p{(\columnwidth - 4\tabcolsep) * \real{0.62}}@{}}
\toprule
名称 & Python 类型 & 含义 \\
\midrule
\endhead
\passthrough{\lstinline!value!} &
\passthrough{\lstinline!Sequence[int]!} &
要写入或传递的新值，必须符合该参数的 Python 类型以及底层 C++
范围约束。 \\
\bottomrule
\end{longtable}

\textbf{返回值}

\begin{longtable}[]{@{}
  >{\raggedright\arraybackslash}p{(\columnwidth - 4\tabcolsep) * \real{0.18}}
  >{\raggedright\arraybackslash}p{(\columnwidth - 4\tabcolsep) * \real{0.20}}
  >{\raggedright\arraybackslash}p{(\columnwidth - 4\tabcolsep) * \real{0.62}}@{}}
\toprule
位置 & 类型 & 含义 \\
\midrule
\endhead
getter 返回值 & \passthrough{\lstinline!list[int]!} &
返回属性当前值。数组、vector 和嵌套消息采用复制语义。 \\
\bottomrule
\end{longtable}

\textbf{对应 C++}

\begin{itemize}
\tightlist
\item
  类：\passthrough{\lstinline!unitree\_hg::msg::dds\_::BmsCmd\_!}
\item
  字段类型：\passthrough{\lstinline!std::array<uint8\_t, 40>!}
\item
  声明位置：\passthrough{\lstinline!include/unitree/idl/hg/BmsCmd\_.hpp:25!}
\end{itemize}

\textbf{用法}

\begin{lstlisting}[language=Python]
current_value = obj.reserve
obj.reserve = new_value
\end{lstlisting}

\begin{quote}
{[}!TIP{]} 读取容器或嵌套值后应遵循``读取、修改、重新赋值''。只修改
getter 返回的副本不会可靠地更新原 C++ 消息。
\end{quote}

\hypertarget{unitree-sdk2-cpp-idl-hg-bmsstate}{%
\subsubsection{\texorpdfstring{\texttt{unitree\_sdk2\_cpp.idl.hg.BmsState}}{unitree\_sdk2\_cpp.idl.hg.BmsState}}\label{unitree-sdk2-cpp-idl-hg-bmsstate}}

可由 Python 读写的 DDS/IDL 消息值类型。字段采用复制语义。

\textbf{导入}

\begin{lstlisting}[language=Python]
from unitree_sdk2_cpp.idl.hg import BmsState
\end{lstlisting}

\textbf{方法索引}

\begin{longtable}[]{@{}
  >{\raggedright\arraybackslash}p{(\columnwidth - 6\tabcolsep) * \real{0.18}}
  >{\raggedright\arraybackslash}p{(\columnwidth - 6\tabcolsep) * \real{0.24}}
  >{\raggedright\arraybackslash}p{(\columnwidth - 6\tabcolsep) * \real{0.18}}
  >{\raggedright\arraybackslash}p{(\columnwidth - 6\tabcolsep) * \real{0.40}}@{}}
\toprule
方法 & Python 签名 & 状态 & 安全分类 \\
\midrule
\endhead
\protect\hyperlink{unitree-sdk2-cpp-idl-hg-bmsstate-dunder-init-1}{\passthrough{\lstinline!\_\_init\_\_!}}
& \passthrough{\lstinline!def \_\_init\_\_(self) -> None!} &
\passthrough{\lstinline!AVAILABLE!} &
\passthrough{\lstinline!VALUE\_TYPE!} \\
\protect\hyperlink{unitree-sdk2-cpp-idl-hg-bmsstate-dunder-eq-1}{\passthrough{\lstinline!\_\_eq\_\_!}}
& \passthrough{\lstinline!def \_\_eq\_\_(self, other: object) -> bool!}
& \passthrough{\lstinline!AVAILABLE!} &
\passthrough{\lstinline!VALUE\_TYPE!} \\
\protect\hyperlink{unitree-sdk2-cpp-idl-hg-bmsstate-dunder-ne-1}{\passthrough{\lstinline!\_\_ne\_\_!}}
& \passthrough{\lstinline!def \_\_ne\_\_(self, other: object) -> bool!}
& \passthrough{\lstinline!AVAILABLE!} &
\passthrough{\lstinline!VALUE\_TYPE!} \\
\bottomrule
\end{longtable}

\textbf{属性索引}

\begin{longtable}[]{@{}
  >{\raggedright\arraybackslash}p{(\columnwidth - 6\tabcolsep) * \real{0.18}}
  >{\raggedright\arraybackslash}p{(\columnwidth - 6\tabcolsep) * \real{0.24}}
  >{\raggedright\arraybackslash}p{(\columnwidth - 6\tabcolsep) * \real{0.18}}
  >{\raggedright\arraybackslash}p{(\columnwidth - 6\tabcolsep) * \real{0.40}}@{}}
\toprule
属性 & 读取类型 & 写入类型 & 状态 \\
\midrule
\endhead
\protect\hyperlink{unitree-sdk2-cpp-idl-hg-bmsstate-version-high}{\passthrough{\lstinline!version\_high!}}
& \passthrough{\lstinline!int!} & \passthrough{\lstinline!int!} &
\passthrough{\lstinline!AVAILABLE!} \\
\protect\hyperlink{unitree-sdk2-cpp-idl-hg-bmsstate-version-low}{\passthrough{\lstinline!version\_low!}}
& \passthrough{\lstinline!int!} & \passthrough{\lstinline!int!} &
\passthrough{\lstinline!AVAILABLE!} \\
\protect\hyperlink{unitree-sdk2-cpp-idl-hg-bmsstate-fn}{\passthrough{\lstinline!fn!}}
& \passthrough{\lstinline!int!} & \passthrough{\lstinline!int!} &
\passthrough{\lstinline!AVAILABLE!} \\
\protect\hyperlink{unitree-sdk2-cpp-idl-hg-bmsstate-cell-vol}{\passthrough{\lstinline!cell\_vol!}}
& \passthrough{\lstinline!list[int]!} &
\passthrough{\lstinline!Sequence[int]!} &
\passthrough{\lstinline!AVAILABLE!} \\
\protect\hyperlink{unitree-sdk2-cpp-idl-hg-bmsstate-bmsvoltage}{\passthrough{\lstinline!bmsvoltage!}}
& \passthrough{\lstinline!list[int]!} &
\passthrough{\lstinline!Sequence[int]!} &
\passthrough{\lstinline!AVAILABLE!} \\
\protect\hyperlink{unitree-sdk2-cpp-idl-hg-bmsstate-current}{\passthrough{\lstinline!current!}}
& \passthrough{\lstinline!int!} & \passthrough{\lstinline!int!} &
\passthrough{\lstinline!AVAILABLE!} \\
\protect\hyperlink{unitree-sdk2-cpp-idl-hg-bmsstate-soc}{\passthrough{\lstinline!soc!}}
& \passthrough{\lstinline!int!} & \passthrough{\lstinline!int!} &
\passthrough{\lstinline!AVAILABLE!} \\
\protect\hyperlink{unitree-sdk2-cpp-idl-hg-bmsstate-soh}{\passthrough{\lstinline!soh!}}
& \passthrough{\lstinline!int!} & \passthrough{\lstinline!int!} &
\passthrough{\lstinline!AVAILABLE!} \\
\protect\hyperlink{unitree-sdk2-cpp-idl-hg-bmsstate-temperature}{\passthrough{\lstinline!temperature!}}
& \passthrough{\lstinline!list[int]!} &
\passthrough{\lstinline!Sequence[int]!} &
\passthrough{\lstinline!AVAILABLE!} \\
\protect\hyperlink{unitree-sdk2-cpp-idl-hg-bmsstate-cycle}{\passthrough{\lstinline!cycle!}}
& \passthrough{\lstinline!int!} & \passthrough{\lstinline!int!} &
\passthrough{\lstinline!AVAILABLE!} \\
\protect\hyperlink{unitree-sdk2-cpp-idl-hg-bmsstate-manufacturer-date}{\passthrough{\lstinline!manufacturer\_date!}}
& \passthrough{\lstinline!int!} & \passthrough{\lstinline!int!} &
\passthrough{\lstinline!AVAILABLE!} \\
\protect\hyperlink{unitree-sdk2-cpp-idl-hg-bmsstate-bmsstate}{\passthrough{\lstinline!bmsstate!}}
& \passthrough{\lstinline!list[int]!} &
\passthrough{\lstinline!Sequence[int]!} &
\passthrough{\lstinline!AVAILABLE!} \\
\protect\hyperlink{unitree-sdk2-cpp-idl-hg-bmsstate-reserve}{\passthrough{\lstinline!reserve!}}
& \passthrough{\lstinline!list[int]!} &
\passthrough{\lstinline!Sequence[int]!} &
\passthrough{\lstinline!AVAILABLE!} \\
\bottomrule
\end{longtable}

\hypertarget{unitree-sdk2-cpp-idl-hg-bmsstate-dunder-init-1}{%
\paragraph{\texorpdfstring{\texttt{unitree\_sdk2\_cpp.idl.hg.BmsState.\_\_init\_\_}}{unitree\_sdk2\_cpp.idl.hg.BmsState.\_\_init\_\_}}\label{unitree-sdk2-cpp-idl-hg-bmsstate-dunder-init-1}}

初始化 \passthrough{\lstinline!BmsState!}
实例。是否能实际构造取决于下方可用性状态。

\textbf{签名}

\begin{lstlisting}[language=Python]
def __init__(self) -> None
\end{lstlisting}

\textbf{可用性与安全性}

\textbf{\passthrough{\lstinline!AVAILABLE!} /
\passthrough{\lstinline!VALUE\_TYPE!}}：当前绑定源码已实现该消息值操作。

\textbf{参数}

无显式参数。实例方法中的 \passthrough{\lstinline!self!} 由 Python
自动传入。

\textbf{返回值}

\begin{longtable}[]{@{}
  >{\raggedright\arraybackslash}p{(\columnwidth - 4\tabcolsep) * \real{0.18}}
  >{\raggedright\arraybackslash}p{(\columnwidth - 4\tabcolsep) * \real{0.20}}
  >{\raggedright\arraybackslash}p{(\columnwidth - 4\tabcolsep) * \real{0.62}}@{}}
\toprule
位置 & 类型 & 含义 \\
\midrule
\endhead
无 & \passthrough{\lstinline!None!} &
\passthrough{\lstinline!\_\_init\_\_!}
本身不返回值；调用类对象时，在构造函数可用的前提下得到该类实例。 \\
\bottomrule
\end{longtable}

\textbf{对应 C++}

\begin{itemize}
\tightlist
\item
  类：\passthrough{\lstinline!unitree\_hg::msg::dds\_::BmsState\_!}
\item
  签名：\passthrough{\lstinline!BmsState\_()!}
\item
  绑定策略：\passthrough{\lstinline!未单独标注!}
\item
  声明位置：\passthrough{\lstinline!include/unitree/idl/hg/BmsState\_.hpp:39!}
\end{itemize}

\textbf{说明}

\begin{itemize}
\tightlist
\item
  示例是局部调用片段；\passthrough{\lstinline!obj!}
  和参数变量表示已经按上层业务校验并准备好的对象。
\end{itemize}

\textbf{用法}

\begin{lstlisting}[language=Python]
value = BmsState()
\end{lstlisting}

\hypertarget{unitree-sdk2-cpp-idl-hg-bmsstate-dunder-eq-1}{%
\paragraph{\texorpdfstring{\texttt{unitree\_sdk2\_cpp.idl.hg.BmsState.\_\_eq\_\_}}{unitree\_sdk2\_cpp.idl.hg.BmsState.\_\_eq\_\_}}\label{unitree-sdk2-cpp-idl-hg-bmsstate-dunder-eq-1}}

按底层 IDL 消息内容比较两个对象是否相等。

\textbf{签名}

\begin{lstlisting}[language=Python]
def __eq__(self, other: object) -> bool
\end{lstlisting}

\textbf{可用性与安全性}

\textbf{\passthrough{\lstinline!AVAILABLE!} /
\passthrough{\lstinline!VALUE\_TYPE!}}：当前绑定源码已实现该消息值操作。

\textbf{参数}

\begin{longtable}[]{@{}
  >{\raggedright\arraybackslash}p{(\columnwidth - 6\tabcolsep) * \real{0.18}}
  >{\raggedright\arraybackslash}p{(\columnwidth - 6\tabcolsep) * \real{0.24}}
  >{\raggedright\arraybackslash}p{(\columnwidth - 6\tabcolsep) * \real{0.18}}
  >{\raggedright\arraybackslash}p{(\columnwidth - 6\tabcolsep) * \real{0.40}}@{}}
\toprule
名称 & Python 类型 & 默认值 & 含义 \\
\midrule
\endhead
\passthrough{\lstinline!other!} & \passthrough{\lstinline!object!} &
必填 & 用于比较的另一个 Python 对象。类型不兼容时比较结果通常为
\passthrough{\lstinline!False!}。 \\
\bottomrule
\end{longtable}

\textbf{返回值}

\begin{longtable}[]{@{}
  >{\raggedright\arraybackslash}p{(\columnwidth - 4\tabcolsep) * \real{0.18}}
  >{\raggedright\arraybackslash}p{(\columnwidth - 4\tabcolsep) * \real{0.20}}
  >{\raggedright\arraybackslash}p{(\columnwidth - 4\tabcolsep) * \real{0.62}}@{}}
\toprule
位置 & 类型 & 含义 \\
\midrule
\endhead
返回值 & \passthrough{\lstinline!bool!} & 返回
\passthrough{\lstinline!bool!}。更精确的业务含义见该方法用途、C++
签名和目标型号协议。 \\
\bottomrule
\end{longtable}

\textbf{对应 C++}

\begin{itemize}
\tightlist
\item
  类：\passthrough{\lstinline!unitree\_hg::msg::dds\_::BmsState\_!}
\item
  签名：\passthrough{\lstinline!operator==(const value \&) const!}
\item
  绑定策略：\passthrough{\lstinline!未单独标注!}
\item
  声明位置：由 pybind11 手工绑定或生成注册表提供；manifest
  未记录独立头文件行号。
\end{itemize}

\textbf{说明}

\begin{itemize}
\tightlist
\item
  示例是局部调用片段；\passthrough{\lstinline!obj!}
  和参数变量表示已经按上层业务校验并准备好的对象。
\end{itemize}

\textbf{用法}

\begin{lstlisting}[language=Python]
same = left == right
\end{lstlisting}

\hypertarget{unitree-sdk2-cpp-idl-hg-bmsstate-dunder-ne-1}{%
\paragraph{\texorpdfstring{\texttt{unitree\_sdk2\_cpp.idl.hg.BmsState.\_\_ne\_\_}}{unitree\_sdk2\_cpp.idl.hg.BmsState.\_\_ne\_\_}}\label{unitree-sdk2-cpp-idl-hg-bmsstate-dunder-ne-1}}

按底层 IDL 消息内容比较两个对象是否不相等。

\textbf{签名}

\begin{lstlisting}[language=Python]
def __ne__(self, other: object) -> bool
\end{lstlisting}

\textbf{可用性与安全性}

\textbf{\passthrough{\lstinline!AVAILABLE!} /
\passthrough{\lstinline!VALUE\_TYPE!}}：当前绑定源码已实现该消息值操作。

\textbf{参数}

\begin{longtable}[]{@{}
  >{\raggedright\arraybackslash}p{(\columnwidth - 6\tabcolsep) * \real{0.18}}
  >{\raggedright\arraybackslash}p{(\columnwidth - 6\tabcolsep) * \real{0.24}}
  >{\raggedright\arraybackslash}p{(\columnwidth - 6\tabcolsep) * \real{0.18}}
  >{\raggedright\arraybackslash}p{(\columnwidth - 6\tabcolsep) * \real{0.40}}@{}}
\toprule
名称 & Python 类型 & 默认值 & 含义 \\
\midrule
\endhead
\passthrough{\lstinline!other!} & \passthrough{\lstinline!object!} &
必填 & 用于比较的另一个 Python 对象。类型不兼容时比较结果通常为
\passthrough{\lstinline!False!}。 \\
\bottomrule
\end{longtable}

\textbf{返回值}

\begin{longtable}[]{@{}
  >{\raggedright\arraybackslash}p{(\columnwidth - 4\tabcolsep) * \real{0.18}}
  >{\raggedright\arraybackslash}p{(\columnwidth - 4\tabcolsep) * \real{0.20}}
  >{\raggedright\arraybackslash}p{(\columnwidth - 4\tabcolsep) * \real{0.62}}@{}}
\toprule
位置 & 类型 & 含义 \\
\midrule
\endhead
返回值 & \passthrough{\lstinline!bool!} & 返回
\passthrough{\lstinline!bool!}。更精确的业务含义见该方法用途、C++
签名和目标型号协议。 \\
\bottomrule
\end{longtable}

\textbf{对应 C++}

\begin{itemize}
\tightlist
\item
  类：\passthrough{\lstinline!unitree\_hg::msg::dds\_::BmsState\_!}
\item
  签名：\passthrough{\lstinline"operator!=(const value \&) const"}
\item
  绑定策略：\passthrough{\lstinline!未单独标注!}
\item
  声明位置：由 pybind11 手工绑定或生成注册表提供；manifest
  未记录独立头文件行号。
\end{itemize}

\textbf{说明}

\begin{itemize}
\tightlist
\item
  示例是局部调用片段；\passthrough{\lstinline!obj!}
  和参数变量表示已经按上层业务校验并准备好的对象。
\end{itemize}

\textbf{用法}

\begin{lstlisting}[language=Python]
different = left != right
\end{lstlisting}

\hypertarget{unitree-sdk2-cpp-idl-hg-bmsstate-version-high}{%
\paragraph{\texorpdfstring{\texttt{unitree\_sdk2\_cpp.idl.hg.BmsState.version\_high}}{unitree\_sdk2\_cpp.idl.hg.BmsState.version\_high}}\label{unitree-sdk2-cpp-idl-hg-bmsstate-version-high}}

底层 \passthrough{\lstinline!unitree\_hg::msg::dds\_::BmsState\_!} 的
\passthrough{\lstinline!version\_high!} 字段。类型签名只说明 Python
表示；单位、数值范围、数组固定长度和枚举含义需查对应 IDL 头文件。
取值范围为 0 到 255。

\textbf{签名}

\begin{lstlisting}[language=Python]
@property
def version_high(self) -> int

@version_high.setter
def version_high(self, value: int) -> None
\end{lstlisting}

\textbf{可用性与安全性}

\textbf{\passthrough{\lstinline!AVAILABLE!} /
\passthrough{\lstinline!VALUE\_TYPE!}}：当前绑定源码已实现该消息值操作。

\textbf{参数}

\begin{longtable}[]{@{}
  >{\raggedright\arraybackslash}p{(\columnwidth - 4\tabcolsep) * \real{0.18}}
  >{\raggedright\arraybackslash}p{(\columnwidth - 4\tabcolsep) * \real{0.20}}
  >{\raggedright\arraybackslash}p{(\columnwidth - 4\tabcolsep) * \real{0.62}}@{}}
\toprule
名称 & Python 类型 & 含义 \\
\midrule
\endhead
\passthrough{\lstinline!value!} & \passthrough{\lstinline!int!} &
要写入或传递的新值，必须符合该参数的 Python 类型以及底层 C++
范围约束。 \\
\bottomrule
\end{longtable}

\textbf{返回值}

\begin{longtable}[]{@{}
  >{\raggedright\arraybackslash}p{(\columnwidth - 4\tabcolsep) * \real{0.18}}
  >{\raggedright\arraybackslash}p{(\columnwidth - 4\tabcolsep) * \real{0.20}}
  >{\raggedright\arraybackslash}p{(\columnwidth - 4\tabcolsep) * \real{0.62}}@{}}
\toprule
位置 & 类型 & 含义 \\
\midrule
\endhead
getter 返回值 & \passthrough{\lstinline!int!} & 返回属性当前值。 \\
\bottomrule
\end{longtable}

\textbf{对应 C++}

\begin{itemize}
\tightlist
\item
  类：\passthrough{\lstinline!unitree\_hg::msg::dds\_::BmsState\_!}
\item
  字段类型：\passthrough{\lstinline!uint8\_t!}
\item
  声明位置：\passthrough{\lstinline!include/unitree/idl/hg/BmsState\_.hpp:24!}
\end{itemize}

\textbf{用法}

\begin{lstlisting}[language=Python]
current_value = obj.version_high
obj.version_high = new_value
\end{lstlisting}

\hypertarget{unitree-sdk2-cpp-idl-hg-bmsstate-version-low}{%
\paragraph{\texorpdfstring{\texttt{unitree\_sdk2\_cpp.idl.hg.BmsState.version\_low}}{unitree\_sdk2\_cpp.idl.hg.BmsState.version\_low}}\label{unitree-sdk2-cpp-idl-hg-bmsstate-version-low}}

底层 \passthrough{\lstinline!unitree\_hg::msg::dds\_::BmsState\_!} 的
\passthrough{\lstinline!version\_low!} 字段。类型签名只说明 Python
表示；单位、数值范围、数组固定长度和枚举含义需查对应 IDL 头文件。
取值范围为 0 到 255。

\textbf{签名}

\begin{lstlisting}[language=Python]
@property
def version_low(self) -> int

@version_low.setter
def version_low(self, value: int) -> None
\end{lstlisting}

\textbf{可用性与安全性}

\textbf{\passthrough{\lstinline!AVAILABLE!} /
\passthrough{\lstinline!VALUE\_TYPE!}}：当前绑定源码已实现该消息值操作。

\textbf{参数}

\begin{longtable}[]{@{}
  >{\raggedright\arraybackslash}p{(\columnwidth - 4\tabcolsep) * \real{0.18}}
  >{\raggedright\arraybackslash}p{(\columnwidth - 4\tabcolsep) * \real{0.20}}
  >{\raggedright\arraybackslash}p{(\columnwidth - 4\tabcolsep) * \real{0.62}}@{}}
\toprule
名称 & Python 类型 & 含义 \\
\midrule
\endhead
\passthrough{\lstinline!value!} & \passthrough{\lstinline!int!} &
要写入或传递的新值，必须符合该参数的 Python 类型以及底层 C++
范围约束。 \\
\bottomrule
\end{longtable}

\textbf{返回值}

\begin{longtable}[]{@{}
  >{\raggedright\arraybackslash}p{(\columnwidth - 4\tabcolsep) * \real{0.18}}
  >{\raggedright\arraybackslash}p{(\columnwidth - 4\tabcolsep) * \real{0.20}}
  >{\raggedright\arraybackslash}p{(\columnwidth - 4\tabcolsep) * \real{0.62}}@{}}
\toprule
位置 & 类型 & 含义 \\
\midrule
\endhead
getter 返回值 & \passthrough{\lstinline!int!} & 返回属性当前值。 \\
\bottomrule
\end{longtable}

\textbf{对应 C++}

\begin{itemize}
\tightlist
\item
  类：\passthrough{\lstinline!unitree\_hg::msg::dds\_::BmsState\_!}
\item
  字段类型：\passthrough{\lstinline!uint8\_t!}
\item
  声明位置：\passthrough{\lstinline!include/unitree/idl/hg/BmsState\_.hpp:25!}
\end{itemize}

\textbf{用法}

\begin{lstlisting}[language=Python]
current_value = obj.version_low
obj.version_low = new_value
\end{lstlisting}

\hypertarget{unitree-sdk2-cpp-idl-hg-bmsstate-fn}{%
\paragraph{\texorpdfstring{\texttt{unitree\_sdk2\_cpp.idl.hg.BmsState.fn}}{unitree\_sdk2\_cpp.idl.hg.BmsState.fn}}\label{unitree-sdk2-cpp-idl-hg-bmsstate-fn}}

底层 \passthrough{\lstinline!unitree\_hg::msg::dds\_::BmsState\_!} 的
\passthrough{\lstinline!fn!} 字段。类型签名只说明 Python
表示；单位、数值范围、数组固定长度和枚举含义需查对应 IDL 头文件。
取值范围为 0 到 255。

\textbf{签名}

\begin{lstlisting}[language=Python]
@property
def fn(self) -> int

@fn.setter
def fn(self, value: int) -> None
\end{lstlisting}

\textbf{可用性与安全性}

\textbf{\passthrough{\lstinline!AVAILABLE!} /
\passthrough{\lstinline!VALUE\_TYPE!}}：当前绑定源码已实现该消息值操作。

\textbf{参数}

\begin{longtable}[]{@{}
  >{\raggedright\arraybackslash}p{(\columnwidth - 4\tabcolsep) * \real{0.18}}
  >{\raggedright\arraybackslash}p{(\columnwidth - 4\tabcolsep) * \real{0.20}}
  >{\raggedright\arraybackslash}p{(\columnwidth - 4\tabcolsep) * \real{0.62}}@{}}
\toprule
名称 & Python 类型 & 含义 \\
\midrule
\endhead
\passthrough{\lstinline!value!} & \passthrough{\lstinline!int!} &
要写入或传递的新值，必须符合该参数的 Python 类型以及底层 C++
范围约束。 \\
\bottomrule
\end{longtable}

\textbf{返回值}

\begin{longtable}[]{@{}
  >{\raggedright\arraybackslash}p{(\columnwidth - 4\tabcolsep) * \real{0.18}}
  >{\raggedright\arraybackslash}p{(\columnwidth - 4\tabcolsep) * \real{0.20}}
  >{\raggedright\arraybackslash}p{(\columnwidth - 4\tabcolsep) * \real{0.62}}@{}}
\toprule
位置 & 类型 & 含义 \\
\midrule
\endhead
getter 返回值 & \passthrough{\lstinline!int!} & 返回属性当前值。 \\
\bottomrule
\end{longtable}

\textbf{对应 C++}

\begin{itemize}
\tightlist
\item
  类：\passthrough{\lstinline!unitree\_hg::msg::dds\_::BmsState\_!}
\item
  字段类型：\passthrough{\lstinline!uint8\_t!}
\item
  声明位置：\passthrough{\lstinline!include/unitree/idl/hg/BmsState\_.hpp:26!}
\end{itemize}

\textbf{用法}

\begin{lstlisting}[language=Python]
current_value = obj.fn
obj.fn = new_value
\end{lstlisting}

\hypertarget{unitree-sdk2-cpp-idl-hg-bmsstate-cell-vol}{%
\paragraph{\texorpdfstring{\texttt{unitree\_sdk2\_cpp.idl.hg.BmsState.cell\_vol}}{unitree\_sdk2\_cpp.idl.hg.BmsState.cell\_vol}}\label{unitree-sdk2-cpp-idl-hg-bmsstate-cell-vol}}

底层 \passthrough{\lstinline!unitree\_hg::msg::dds\_::BmsState\_!} 的
\passthrough{\lstinline!cell\_vol!} 字段。类型签名只说明 Python
表示；单位、数值范围、数组固定长度和枚举含义需查对应 IDL 头文件。
底层是固定长度数组，写入序列必须正好包含 40 个元素；元素约束：取值范围为
0 到 65535。

\textbf{签名}

\begin{lstlisting}[language=Python]
@property
def cell_vol(self) -> list[int]

@cell_vol.setter
def cell_vol(self, value: Sequence[int]) -> None
\end{lstlisting}

\textbf{可用性与安全性}

\textbf{\passthrough{\lstinline!AVAILABLE!} /
\passthrough{\lstinline!VALUE\_TYPE!}}：当前绑定源码已实现该消息值操作。

\textbf{参数}

\begin{longtable}[]{@{}
  >{\raggedright\arraybackslash}p{(\columnwidth - 4\tabcolsep) * \real{0.18}}
  >{\raggedright\arraybackslash}p{(\columnwidth - 4\tabcolsep) * \real{0.20}}
  >{\raggedright\arraybackslash}p{(\columnwidth - 4\tabcolsep) * \real{0.62}}@{}}
\toprule
名称 & Python 类型 & 含义 \\
\midrule
\endhead
\passthrough{\lstinline!value!} &
\passthrough{\lstinline!Sequence[int]!} &
要写入或传递的新值，必须符合该参数的 Python 类型以及底层 C++
范围约束。 \\
\bottomrule
\end{longtable}

\textbf{返回值}

\begin{longtable}[]{@{}
  >{\raggedright\arraybackslash}p{(\columnwidth - 4\tabcolsep) * \real{0.18}}
  >{\raggedright\arraybackslash}p{(\columnwidth - 4\tabcolsep) * \real{0.20}}
  >{\raggedright\arraybackslash}p{(\columnwidth - 4\tabcolsep) * \real{0.62}}@{}}
\toprule
位置 & 类型 & 含义 \\
\midrule
\endhead
getter 返回值 & \passthrough{\lstinline!list[int]!} &
返回属性当前值。数组、vector 和嵌套消息采用复制语义。 \\
\bottomrule
\end{longtable}

\textbf{对应 C++}

\begin{itemize}
\tightlist
\item
  类：\passthrough{\lstinline!unitree\_hg::msg::dds\_::BmsState\_!}
\item
  字段类型：\passthrough{\lstinline!std::array<uint16\_t, 40>!}
\item
  声明位置：\passthrough{\lstinline!include/unitree/idl/hg/BmsState\_.hpp:27!}
\end{itemize}

\textbf{用法}

\begin{lstlisting}[language=Python]
current_value = obj.cell_vol
obj.cell_vol = new_value
\end{lstlisting}

\begin{quote}
{[}!TIP{]} 读取容器或嵌套值后应遵循``读取、修改、重新赋值''。只修改
getter 返回的副本不会可靠地更新原 C++ 消息。
\end{quote}

\hypertarget{unitree-sdk2-cpp-idl-hg-bmsstate-bmsvoltage}{%
\paragraph{\texorpdfstring{\texttt{unitree\_sdk2\_cpp.idl.hg.BmsState.bmsvoltage}}{unitree\_sdk2\_cpp.idl.hg.BmsState.bmsvoltage}}\label{unitree-sdk2-cpp-idl-hg-bmsstate-bmsvoltage}}

底层 \passthrough{\lstinline!unitree\_hg::msg::dds\_::BmsState\_!} 的
\passthrough{\lstinline!bmsvoltage!} 字段。类型签名只说明 Python
表示；单位、数值范围、数组固定长度和枚举含义需查对应 IDL 头文件。
底层是固定长度数组，写入序列必须正好包含 3 个元素；元素约束：取值范围为
0 到 4294967295。

\textbf{签名}

\begin{lstlisting}[language=Python]
@property
def bmsvoltage(self) -> list[int]

@bmsvoltage.setter
def bmsvoltage(self, value: Sequence[int]) -> None
\end{lstlisting}

\textbf{可用性与安全性}

\textbf{\passthrough{\lstinline!AVAILABLE!} /
\passthrough{\lstinline!VALUE\_TYPE!}}：当前绑定源码已实现该消息值操作。

\textbf{参数}

\begin{longtable}[]{@{}
  >{\raggedright\arraybackslash}p{(\columnwidth - 4\tabcolsep) * \real{0.18}}
  >{\raggedright\arraybackslash}p{(\columnwidth - 4\tabcolsep) * \real{0.20}}
  >{\raggedright\arraybackslash}p{(\columnwidth - 4\tabcolsep) * \real{0.62}}@{}}
\toprule
名称 & Python 类型 & 含义 \\
\midrule
\endhead
\passthrough{\lstinline!value!} &
\passthrough{\lstinline!Sequence[int]!} &
要写入或传递的新值，必须符合该参数的 Python 类型以及底层 C++
范围约束。 \\
\bottomrule
\end{longtable}

\textbf{返回值}

\begin{longtable}[]{@{}
  >{\raggedright\arraybackslash}p{(\columnwidth - 4\tabcolsep) * \real{0.18}}
  >{\raggedright\arraybackslash}p{(\columnwidth - 4\tabcolsep) * \real{0.20}}
  >{\raggedright\arraybackslash}p{(\columnwidth - 4\tabcolsep) * \real{0.62}}@{}}
\toprule
位置 & 类型 & 含义 \\
\midrule
\endhead
getter 返回值 & \passthrough{\lstinline!list[int]!} &
返回属性当前值。数组、vector 和嵌套消息采用复制语义。 \\
\bottomrule
\end{longtable}

\textbf{对应 C++}

\begin{itemize}
\tightlist
\item
  类：\passthrough{\lstinline!unitree\_hg::msg::dds\_::BmsState\_!}
\item
  字段类型：\passthrough{\lstinline!std::array<uint32\_t, 3>!}
\item
  声明位置：\passthrough{\lstinline!include/unitree/idl/hg/BmsState\_.hpp:28!}
\end{itemize}

\textbf{用法}

\begin{lstlisting}[language=Python]
current_value = obj.bmsvoltage
obj.bmsvoltage = new_value
\end{lstlisting}

\begin{quote}
{[}!TIP{]} 读取容器或嵌套值后应遵循``读取、修改、重新赋值''。只修改
getter 返回的副本不会可靠地更新原 C++ 消息。
\end{quote}

\hypertarget{unitree-sdk2-cpp-idl-hg-bmsstate-current}{%
\paragraph{\texorpdfstring{\texttt{unitree\_sdk2\_cpp.idl.hg.BmsState.current}}{unitree\_sdk2\_cpp.idl.hg.BmsState.current}}\label{unitree-sdk2-cpp-idl-hg-bmsstate-current}}

底层 \passthrough{\lstinline!unitree\_hg::msg::dds\_::BmsState\_!} 的
\passthrough{\lstinline!current!} 字段。类型签名只说明 Python
表示；单位、数值范围、数组固定长度和枚举含义需查对应 IDL 头文件。
取值范围为 -2147483648 到 2147483647。

\textbf{签名}

\begin{lstlisting}[language=Python]
@property
def current(self) -> int

@current.setter
def current(self, value: int) -> None
\end{lstlisting}

\textbf{可用性与安全性}

\textbf{\passthrough{\lstinline!AVAILABLE!} /
\passthrough{\lstinline!VALUE\_TYPE!}}：当前绑定源码已实现该消息值操作。

\textbf{参数}

\begin{longtable}[]{@{}
  >{\raggedright\arraybackslash}p{(\columnwidth - 4\tabcolsep) * \real{0.18}}
  >{\raggedright\arraybackslash}p{(\columnwidth - 4\tabcolsep) * \real{0.20}}
  >{\raggedright\arraybackslash}p{(\columnwidth - 4\tabcolsep) * \real{0.62}}@{}}
\toprule
名称 & Python 类型 & 含义 \\
\midrule
\endhead
\passthrough{\lstinline!value!} & \passthrough{\lstinline!int!} &
要写入或传递的新值，必须符合该参数的 Python 类型以及底层 C++
范围约束。 \\
\bottomrule
\end{longtable}

\textbf{返回值}

\begin{longtable}[]{@{}
  >{\raggedright\arraybackslash}p{(\columnwidth - 4\tabcolsep) * \real{0.18}}
  >{\raggedright\arraybackslash}p{(\columnwidth - 4\tabcolsep) * \real{0.20}}
  >{\raggedright\arraybackslash}p{(\columnwidth - 4\tabcolsep) * \real{0.62}}@{}}
\toprule
位置 & 类型 & 含义 \\
\midrule
\endhead
getter 返回值 & \passthrough{\lstinline!int!} & 返回属性当前值。 \\
\bottomrule
\end{longtable}

\textbf{对应 C++}

\begin{itemize}
\tightlist
\item
  类：\passthrough{\lstinline!unitree\_hg::msg::dds\_::BmsState\_!}
\item
  字段类型：\passthrough{\lstinline!int32\_t!}
\item
  声明位置：\passthrough{\lstinline!include/unitree/idl/hg/BmsState\_.hpp:29!}
\end{itemize}

\textbf{用法}

\begin{lstlisting}[language=Python]
current_value = obj.current
obj.current = new_value
\end{lstlisting}

\hypertarget{unitree-sdk2-cpp-idl-hg-bmsstate-soc}{%
\paragraph{\texorpdfstring{\texttt{unitree\_sdk2\_cpp.idl.hg.BmsState.soc}}{unitree\_sdk2\_cpp.idl.hg.BmsState.soc}}\label{unitree-sdk2-cpp-idl-hg-bmsstate-soc}}

底层 \passthrough{\lstinline!unitree\_hg::msg::dds\_::BmsState\_!} 的
\passthrough{\lstinline!soc!} 字段。类型签名只说明 Python
表示；单位、数值范围、数组固定长度和枚举含义需查对应 IDL 头文件。
取值范围为 0 到 255。

\textbf{签名}

\begin{lstlisting}[language=Python]
@property
def soc(self) -> int

@soc.setter
def soc(self, value: int) -> None
\end{lstlisting}

\textbf{可用性与安全性}

\textbf{\passthrough{\lstinline!AVAILABLE!} /
\passthrough{\lstinline!VALUE\_TYPE!}}：当前绑定源码已实现该消息值操作。

\textbf{参数}

\begin{longtable}[]{@{}
  >{\raggedright\arraybackslash}p{(\columnwidth - 4\tabcolsep) * \real{0.18}}
  >{\raggedright\arraybackslash}p{(\columnwidth - 4\tabcolsep) * \real{0.20}}
  >{\raggedright\arraybackslash}p{(\columnwidth - 4\tabcolsep) * \real{0.62}}@{}}
\toprule
名称 & Python 类型 & 含义 \\
\midrule
\endhead
\passthrough{\lstinline!value!} & \passthrough{\lstinline!int!} &
要写入或传递的新值，必须符合该参数的 Python 类型以及底层 C++
范围约束。 \\
\bottomrule
\end{longtable}

\textbf{返回值}

\begin{longtable}[]{@{}
  >{\raggedright\arraybackslash}p{(\columnwidth - 4\tabcolsep) * \real{0.18}}
  >{\raggedright\arraybackslash}p{(\columnwidth - 4\tabcolsep) * \real{0.20}}
  >{\raggedright\arraybackslash}p{(\columnwidth - 4\tabcolsep) * \real{0.62}}@{}}
\toprule
位置 & 类型 & 含义 \\
\midrule
\endhead
getter 返回值 & \passthrough{\lstinline!int!} & 返回属性当前值。 \\
\bottomrule
\end{longtable}

\textbf{对应 C++}

\begin{itemize}
\tightlist
\item
  类：\passthrough{\lstinline!unitree\_hg::msg::dds\_::BmsState\_!}
\item
  字段类型：\passthrough{\lstinline!uint8\_t!}
\item
  声明位置：\passthrough{\lstinline!include/unitree/idl/hg/BmsState\_.hpp:30!}
\end{itemize}

\textbf{用法}

\begin{lstlisting}[language=Python]
current_value = obj.soc
obj.soc = new_value
\end{lstlisting}

\hypertarget{unitree-sdk2-cpp-idl-hg-bmsstate-soh}{%
\paragraph{\texorpdfstring{\texttt{unitree\_sdk2\_cpp.idl.hg.BmsState.soh}}{unitree\_sdk2\_cpp.idl.hg.BmsState.soh}}\label{unitree-sdk2-cpp-idl-hg-bmsstate-soh}}

底层 \passthrough{\lstinline!unitree\_hg::msg::dds\_::BmsState\_!} 的
\passthrough{\lstinline!soh!} 字段。类型签名只说明 Python
表示；单位、数值范围、数组固定长度和枚举含义需查对应 IDL 头文件。
取值范围为 0 到 255。

\textbf{签名}

\begin{lstlisting}[language=Python]
@property
def soh(self) -> int

@soh.setter
def soh(self, value: int) -> None
\end{lstlisting}

\textbf{可用性与安全性}

\textbf{\passthrough{\lstinline!AVAILABLE!} /
\passthrough{\lstinline!VALUE\_TYPE!}}：当前绑定源码已实现该消息值操作。

\textbf{参数}

\begin{longtable}[]{@{}
  >{\raggedright\arraybackslash}p{(\columnwidth - 4\tabcolsep) * \real{0.18}}
  >{\raggedright\arraybackslash}p{(\columnwidth - 4\tabcolsep) * \real{0.20}}
  >{\raggedright\arraybackslash}p{(\columnwidth - 4\tabcolsep) * \real{0.62}}@{}}
\toprule
名称 & Python 类型 & 含义 \\
\midrule
\endhead
\passthrough{\lstinline!value!} & \passthrough{\lstinline!int!} &
要写入或传递的新值，必须符合该参数的 Python 类型以及底层 C++
范围约束。 \\
\bottomrule
\end{longtable}

\textbf{返回值}

\begin{longtable}[]{@{}
  >{\raggedright\arraybackslash}p{(\columnwidth - 4\tabcolsep) * \real{0.18}}
  >{\raggedright\arraybackslash}p{(\columnwidth - 4\tabcolsep) * \real{0.20}}
  >{\raggedright\arraybackslash}p{(\columnwidth - 4\tabcolsep) * \real{0.62}}@{}}
\toprule
位置 & 类型 & 含义 \\
\midrule
\endhead
getter 返回值 & \passthrough{\lstinline!int!} & 返回属性当前值。 \\
\bottomrule
\end{longtable}

\textbf{对应 C++}

\begin{itemize}
\tightlist
\item
  类：\passthrough{\lstinline!unitree\_hg::msg::dds\_::BmsState\_!}
\item
  字段类型：\passthrough{\lstinline!uint8\_t!}
\item
  声明位置：\passthrough{\lstinline!include/unitree/idl/hg/BmsState\_.hpp:31!}
\end{itemize}

\textbf{用法}

\begin{lstlisting}[language=Python]
current_value = obj.soh
obj.soh = new_value
\end{lstlisting}

\hypertarget{unitree-sdk2-cpp-idl-hg-bmsstate-temperature}{%
\paragraph{\texorpdfstring{\texttt{unitree\_sdk2\_cpp.idl.hg.BmsState.temperature}}{unitree\_sdk2\_cpp.idl.hg.BmsState.temperature}}\label{unitree-sdk2-cpp-idl-hg-bmsstate-temperature}}

底层 \passthrough{\lstinline!unitree\_hg::msg::dds\_::BmsState\_!} 的
\passthrough{\lstinline!temperature!} 字段。类型签名只说明 Python
表示；单位、数值范围、数组固定长度和枚举含义需查对应 IDL 头文件。
底层是固定长度数组，写入序列必须正好包含 12 个元素；元素约束：取值范围为
-32768 到 32767。

\textbf{签名}

\begin{lstlisting}[language=Python]
@property
def temperature(self) -> list[int]

@temperature.setter
def temperature(self, value: Sequence[int]) -> None
\end{lstlisting}

\textbf{可用性与安全性}

\textbf{\passthrough{\lstinline!AVAILABLE!} /
\passthrough{\lstinline!VALUE\_TYPE!}}：当前绑定源码已实现该消息值操作。

\textbf{参数}

\begin{longtable}[]{@{}
  >{\raggedright\arraybackslash}p{(\columnwidth - 4\tabcolsep) * \real{0.18}}
  >{\raggedright\arraybackslash}p{(\columnwidth - 4\tabcolsep) * \real{0.20}}
  >{\raggedright\arraybackslash}p{(\columnwidth - 4\tabcolsep) * \real{0.62}}@{}}
\toprule
名称 & Python 类型 & 含义 \\
\midrule
\endhead
\passthrough{\lstinline!value!} &
\passthrough{\lstinline!Sequence[int]!} &
要写入或传递的新值，必须符合该参数的 Python 类型以及底层 C++
范围约束。 \\
\bottomrule
\end{longtable}

\textbf{返回值}

\begin{longtable}[]{@{}
  >{\raggedright\arraybackslash}p{(\columnwidth - 4\tabcolsep) * \real{0.18}}
  >{\raggedright\arraybackslash}p{(\columnwidth - 4\tabcolsep) * \real{0.20}}
  >{\raggedright\arraybackslash}p{(\columnwidth - 4\tabcolsep) * \real{0.62}}@{}}
\toprule
位置 & 类型 & 含义 \\
\midrule
\endhead
getter 返回值 & \passthrough{\lstinline!list[int]!} &
返回属性当前值。数组、vector 和嵌套消息采用复制语义。 \\
\bottomrule
\end{longtable}

\textbf{对应 C++}

\begin{itemize}
\tightlist
\item
  类：\passthrough{\lstinline!unitree\_hg::msg::dds\_::BmsState\_!}
\item
  字段类型：\passthrough{\lstinline!std::array<int16\_t, 12>!}
\item
  声明位置：\passthrough{\lstinline!include/unitree/idl/hg/BmsState\_.hpp:32!}
\end{itemize}

\textbf{用法}

\begin{lstlisting}[language=Python]
current_value = obj.temperature
obj.temperature = new_value
\end{lstlisting}

\begin{quote}
{[}!TIP{]} 读取容器或嵌套值后应遵循``读取、修改、重新赋值''。只修改
getter 返回的副本不会可靠地更新原 C++ 消息。
\end{quote}

\hypertarget{unitree-sdk2-cpp-idl-hg-bmsstate-cycle}{%
\paragraph{\texorpdfstring{\texttt{unitree\_sdk2\_cpp.idl.hg.BmsState.cycle}}{unitree\_sdk2\_cpp.idl.hg.BmsState.cycle}}\label{unitree-sdk2-cpp-idl-hg-bmsstate-cycle}}

底层 \passthrough{\lstinline!unitree\_hg::msg::dds\_::BmsState\_!} 的
\passthrough{\lstinline!cycle!} 字段。类型签名只说明 Python
表示；单位、数值范围、数组固定长度和枚举含义需查对应 IDL 头文件。
取值范围为 0 到 65535。

\textbf{签名}

\begin{lstlisting}[language=Python]
@property
def cycle(self) -> int

@cycle.setter
def cycle(self, value: int) -> None
\end{lstlisting}

\textbf{可用性与安全性}

\textbf{\passthrough{\lstinline!AVAILABLE!} /
\passthrough{\lstinline!VALUE\_TYPE!}}：当前绑定源码已实现该消息值操作。

\textbf{参数}

\begin{longtable}[]{@{}
  >{\raggedright\arraybackslash}p{(\columnwidth - 4\tabcolsep) * \real{0.18}}
  >{\raggedright\arraybackslash}p{(\columnwidth - 4\tabcolsep) * \real{0.20}}
  >{\raggedright\arraybackslash}p{(\columnwidth - 4\tabcolsep) * \real{0.62}}@{}}
\toprule
名称 & Python 类型 & 含义 \\
\midrule
\endhead
\passthrough{\lstinline!value!} & \passthrough{\lstinline!int!} &
要写入或传递的新值，必须符合该参数的 Python 类型以及底层 C++
范围约束。 \\
\bottomrule
\end{longtable}

\textbf{返回值}

\begin{longtable}[]{@{}
  >{\raggedright\arraybackslash}p{(\columnwidth - 4\tabcolsep) * \real{0.18}}
  >{\raggedright\arraybackslash}p{(\columnwidth - 4\tabcolsep) * \real{0.20}}
  >{\raggedright\arraybackslash}p{(\columnwidth - 4\tabcolsep) * \real{0.62}}@{}}
\toprule
位置 & 类型 & 含义 \\
\midrule
\endhead
getter 返回值 & \passthrough{\lstinline!int!} & 返回属性当前值。 \\
\bottomrule
\end{longtable}

\textbf{对应 C++}

\begin{itemize}
\tightlist
\item
  类：\passthrough{\lstinline!unitree\_hg::msg::dds\_::BmsState\_!}
\item
  字段类型：\passthrough{\lstinline!uint16\_t!}
\item
  声明位置：\passthrough{\lstinline!include/unitree/idl/hg/BmsState\_.hpp:33!}
\end{itemize}

\textbf{用法}

\begin{lstlisting}[language=Python]
current_value = obj.cycle
obj.cycle = new_value
\end{lstlisting}

\hypertarget{unitree-sdk2-cpp-idl-hg-bmsstate-manufacturer-date}{%
\paragraph{\texorpdfstring{\texttt{unitree\_sdk2\_cpp.idl.hg.BmsState.manufacturer\_date}}{unitree\_sdk2\_cpp.idl.hg.BmsState.manufacturer\_date}}\label{unitree-sdk2-cpp-idl-hg-bmsstate-manufacturer-date}}

底层 \passthrough{\lstinline!unitree\_hg::msg::dds\_::BmsState\_!} 的
\passthrough{\lstinline!manufacturer\_date!} 字段。类型签名只说明 Python
表示；单位、数值范围、数组固定长度和枚举含义需查对应 IDL 头文件。
取值范围为 0 到 65535。

\textbf{签名}

\begin{lstlisting}[language=Python]
@property
def manufacturer_date(self) -> int

@manufacturer_date.setter
def manufacturer_date(self, value: int) -> None
\end{lstlisting}

\textbf{可用性与安全性}

\textbf{\passthrough{\lstinline!AVAILABLE!} /
\passthrough{\lstinline!VALUE\_TYPE!}}：当前绑定源码已实现该消息值操作。

\textbf{参数}

\begin{longtable}[]{@{}
  >{\raggedright\arraybackslash}p{(\columnwidth - 4\tabcolsep) * \real{0.18}}
  >{\raggedright\arraybackslash}p{(\columnwidth - 4\tabcolsep) * \real{0.20}}
  >{\raggedright\arraybackslash}p{(\columnwidth - 4\tabcolsep) * \real{0.62}}@{}}
\toprule
名称 & Python 类型 & 含义 \\
\midrule
\endhead
\passthrough{\lstinline!value!} & \passthrough{\lstinline!int!} &
要写入或传递的新值，必须符合该参数的 Python 类型以及底层 C++
范围约束。 \\
\bottomrule
\end{longtable}

\textbf{返回值}

\begin{longtable}[]{@{}
  >{\raggedright\arraybackslash}p{(\columnwidth - 4\tabcolsep) * \real{0.18}}
  >{\raggedright\arraybackslash}p{(\columnwidth - 4\tabcolsep) * \real{0.20}}
  >{\raggedright\arraybackslash}p{(\columnwidth - 4\tabcolsep) * \real{0.62}}@{}}
\toprule
位置 & 类型 & 含义 \\
\midrule
\endhead
getter 返回值 & \passthrough{\lstinline!int!} & 返回属性当前值。 \\
\bottomrule
\end{longtable}

\textbf{对应 C++}

\begin{itemize}
\tightlist
\item
  类：\passthrough{\lstinline!unitree\_hg::msg::dds\_::BmsState\_!}
\item
  字段类型：\passthrough{\lstinline!uint16\_t!}
\item
  声明位置：\passthrough{\lstinline!include/unitree/idl/hg/BmsState\_.hpp:34!}
\end{itemize}

\textbf{用法}

\begin{lstlisting}[language=Python]
current_value = obj.manufacturer_date
obj.manufacturer_date = new_value
\end{lstlisting}

\hypertarget{unitree-sdk2-cpp-idl-hg-bmsstate-bmsstate}{%
\paragraph{\texorpdfstring{\texttt{unitree\_sdk2\_cpp.idl.hg.BmsState.bmsstate}}{unitree\_sdk2\_cpp.idl.hg.BmsState.bmsstate}}\label{unitree-sdk2-cpp-idl-hg-bmsstate-bmsstate}}

底层 \passthrough{\lstinline!unitree\_hg::msg::dds\_::BmsState\_!} 的
\passthrough{\lstinline!bmsstate!} 字段。类型签名只说明 Python
表示；单位、数值范围、数组固定长度和枚举含义需查对应 IDL 头文件。
底层是固定长度数组，写入序列必须正好包含 5 个元素；元素约束：取值范围为
0 到 4294967295。

\textbf{签名}

\begin{lstlisting}[language=Python]
@property
def bmsstate(self) -> list[int]

@bmsstate.setter
def bmsstate(self, value: Sequence[int]) -> None
\end{lstlisting}

\textbf{可用性与安全性}

\textbf{\passthrough{\lstinline!AVAILABLE!} /
\passthrough{\lstinline!VALUE\_TYPE!}}：当前绑定源码已实现该消息值操作。

\textbf{参数}

\begin{longtable}[]{@{}
  >{\raggedright\arraybackslash}p{(\columnwidth - 4\tabcolsep) * \real{0.18}}
  >{\raggedright\arraybackslash}p{(\columnwidth - 4\tabcolsep) * \real{0.20}}
  >{\raggedright\arraybackslash}p{(\columnwidth - 4\tabcolsep) * \real{0.62}}@{}}
\toprule
名称 & Python 类型 & 含义 \\
\midrule
\endhead
\passthrough{\lstinline!value!} &
\passthrough{\lstinline!Sequence[int]!} &
要写入或传递的新值，必须符合该参数的 Python 类型以及底层 C++
范围约束。 \\
\bottomrule
\end{longtable}

\textbf{返回值}

\begin{longtable}[]{@{}
  >{\raggedright\arraybackslash}p{(\columnwidth - 4\tabcolsep) * \real{0.18}}
  >{\raggedright\arraybackslash}p{(\columnwidth - 4\tabcolsep) * \real{0.20}}
  >{\raggedright\arraybackslash}p{(\columnwidth - 4\tabcolsep) * \real{0.62}}@{}}
\toprule
位置 & 类型 & 含义 \\
\midrule
\endhead
getter 返回值 & \passthrough{\lstinline!list[int]!} &
返回属性当前值。数组、vector 和嵌套消息采用复制语义。 \\
\bottomrule
\end{longtable}

\textbf{对应 C++}

\begin{itemize}
\tightlist
\item
  类：\passthrough{\lstinline!unitree\_hg::msg::dds\_::BmsState\_!}
\item
  字段类型：\passthrough{\lstinline!std::array<uint32\_t, 5>!}
\item
  声明位置：\passthrough{\lstinline!include/unitree/idl/hg/BmsState\_.hpp:35!}
\end{itemize}

\textbf{用法}

\begin{lstlisting}[language=Python]
current_value = obj.bmsstate
obj.bmsstate = new_value
\end{lstlisting}

\begin{quote}
{[}!TIP{]} 读取容器或嵌套值后应遵循``读取、修改、重新赋值''。只修改
getter 返回的副本不会可靠地更新原 C++ 消息。
\end{quote}

\hypertarget{unitree-sdk2-cpp-idl-hg-bmsstate-reserve}{%
\paragraph{\texorpdfstring{\texttt{unitree\_sdk2\_cpp.idl.hg.BmsState.reserve}}{unitree\_sdk2\_cpp.idl.hg.BmsState.reserve}}\label{unitree-sdk2-cpp-idl-hg-bmsstate-reserve}}

底层 \passthrough{\lstinline!unitree\_hg::msg::dds\_::BmsState\_!} 的
\passthrough{\lstinline!reserve!} 字段。类型签名只说明 Python
表示；单位、数值范围、数组固定长度和枚举含义需查对应 IDL 头文件。
底层是固定长度数组，写入序列必须正好包含 3 个元素；元素约束：取值范围为
0 到 4294967295。

\textbf{签名}

\begin{lstlisting}[language=Python]
@property
def reserve(self) -> list[int]

@reserve.setter
def reserve(self, value: Sequence[int]) -> None
\end{lstlisting}

\textbf{可用性与安全性}

\textbf{\passthrough{\lstinline!AVAILABLE!} /
\passthrough{\lstinline!VALUE\_TYPE!}}：当前绑定源码已实现该消息值操作。

\textbf{参数}

\begin{longtable}[]{@{}
  >{\raggedright\arraybackslash}p{(\columnwidth - 4\tabcolsep) * \real{0.18}}
  >{\raggedright\arraybackslash}p{(\columnwidth - 4\tabcolsep) * \real{0.20}}
  >{\raggedright\arraybackslash}p{(\columnwidth - 4\tabcolsep) * \real{0.62}}@{}}
\toprule
名称 & Python 类型 & 含义 \\
\midrule
\endhead
\passthrough{\lstinline!value!} &
\passthrough{\lstinline!Sequence[int]!} &
要写入或传递的新值，必须符合该参数的 Python 类型以及底层 C++
范围约束。 \\
\bottomrule
\end{longtable}

\textbf{返回值}

\begin{longtable}[]{@{}
  >{\raggedright\arraybackslash}p{(\columnwidth - 4\tabcolsep) * \real{0.18}}
  >{\raggedright\arraybackslash}p{(\columnwidth - 4\tabcolsep) * \real{0.20}}
  >{\raggedright\arraybackslash}p{(\columnwidth - 4\tabcolsep) * \real{0.62}}@{}}
\toprule
位置 & 类型 & 含义 \\
\midrule
\endhead
getter 返回值 & \passthrough{\lstinline!list[int]!} &
返回属性当前值。数组、vector 和嵌套消息采用复制语义。 \\
\bottomrule
\end{longtable}

\textbf{对应 C++}

\begin{itemize}
\tightlist
\item
  类：\passthrough{\lstinline!unitree\_hg::msg::dds\_::BmsState\_!}
\item
  字段类型：\passthrough{\lstinline!std::array<uint32\_t, 3>!}
\item
  声明位置：\passthrough{\lstinline!include/unitree/idl/hg/BmsState\_.hpp:36!}
\end{itemize}

\textbf{用法}

\begin{lstlisting}[language=Python]
current_value = obj.reserve
obj.reserve = new_value
\end{lstlisting}

\begin{quote}
{[}!TIP{]} 读取容器或嵌套值后应遵循``读取、修改、重新赋值''。只修改
getter 返回的副本不会可靠地更新原 C++ 消息。
\end{quote}

\hypertarget{unitree-sdk2-cpp-idl-hg-motorcmd}{%
\subsubsection{\texorpdfstring{\texttt{unitree\_sdk2\_cpp.idl.hg.MotorCmd}}{unitree\_sdk2\_cpp.idl.hg.MotorCmd}}\label{unitree-sdk2-cpp-idl-hg-motorcmd}}

可由 Python 读写的 DDS/IDL 消息值类型。字段采用复制语义。

\textbf{导入}

\begin{lstlisting}[language=Python]
from unitree_sdk2_cpp.idl.hg import MotorCmd
\end{lstlisting}

\textbf{方法索引}

\begin{longtable}[]{@{}
  >{\raggedright\arraybackslash}p{(\columnwidth - 6\tabcolsep) * \real{0.18}}
  >{\raggedright\arraybackslash}p{(\columnwidth - 6\tabcolsep) * \real{0.24}}
  >{\raggedright\arraybackslash}p{(\columnwidth - 6\tabcolsep) * \real{0.18}}
  >{\raggedright\arraybackslash}p{(\columnwidth - 6\tabcolsep) * \real{0.40}}@{}}
\toprule
方法 & Python 签名 & 状态 & 安全分类 \\
\midrule
\endhead
\protect\hyperlink{unitree-sdk2-cpp-idl-hg-motorcmd-dunder-init-1}{\passthrough{\lstinline!\_\_init\_\_!}}
& \passthrough{\lstinline!def \_\_init\_\_(self) -> None!} &
\passthrough{\lstinline!AVAILABLE!} &
\passthrough{\lstinline!VALUE\_TYPE!} \\
\protect\hyperlink{unitree-sdk2-cpp-idl-hg-motorcmd-dunder-eq-1}{\passthrough{\lstinline!\_\_eq\_\_!}}
& \passthrough{\lstinline!def \_\_eq\_\_(self, other: object) -> bool!}
& \passthrough{\lstinline!AVAILABLE!} &
\passthrough{\lstinline!VALUE\_TYPE!} \\
\protect\hyperlink{unitree-sdk2-cpp-idl-hg-motorcmd-dunder-ne-1}{\passthrough{\lstinline!\_\_ne\_\_!}}
& \passthrough{\lstinline!def \_\_ne\_\_(self, other: object) -> bool!}
& \passthrough{\lstinline!AVAILABLE!} &
\passthrough{\lstinline!VALUE\_TYPE!} \\
\bottomrule
\end{longtable}

\textbf{属性索引}

\begin{longtable}[]{@{}
  >{\raggedright\arraybackslash}p{(\columnwidth - 6\tabcolsep) * \real{0.18}}
  >{\raggedright\arraybackslash}p{(\columnwidth - 6\tabcolsep) * \real{0.24}}
  >{\raggedright\arraybackslash}p{(\columnwidth - 6\tabcolsep) * \real{0.18}}
  >{\raggedright\arraybackslash}p{(\columnwidth - 6\tabcolsep) * \real{0.40}}@{}}
\toprule
属性 & 读取类型 & 写入类型 & 状态 \\
\midrule
\endhead
\protect\hyperlink{unitree-sdk2-cpp-idl-hg-motorcmd-mode}{\passthrough{\lstinline!mode!}}
& \passthrough{\lstinline!int!} & \passthrough{\lstinline!int!} &
\passthrough{\lstinline!AVAILABLE!} \\
\protect\hyperlink{unitree-sdk2-cpp-idl-hg-motorcmd-q}{\passthrough{\lstinline!q!}}
& \passthrough{\lstinline!float!} & \passthrough{\lstinline!float!} &
\passthrough{\lstinline!AVAILABLE!} \\
\protect\hyperlink{unitree-sdk2-cpp-idl-hg-motorcmd-dq}{\passthrough{\lstinline!dq!}}
& \passthrough{\lstinline!float!} & \passthrough{\lstinline!float!} &
\passthrough{\lstinline!AVAILABLE!} \\
\protect\hyperlink{unitree-sdk2-cpp-idl-hg-motorcmd-tau}{\passthrough{\lstinline!tau!}}
& \passthrough{\lstinline!float!} & \passthrough{\lstinline!float!} &
\passthrough{\lstinline!AVAILABLE!} \\
\protect\hyperlink{unitree-sdk2-cpp-idl-hg-motorcmd-kp}{\passthrough{\lstinline!kp!}}
& \passthrough{\lstinline!float!} & \passthrough{\lstinline!float!} &
\passthrough{\lstinline!AVAILABLE!} \\
\protect\hyperlink{unitree-sdk2-cpp-idl-hg-motorcmd-kd}{\passthrough{\lstinline!kd!}}
& \passthrough{\lstinline!float!} & \passthrough{\lstinline!float!} &
\passthrough{\lstinline!AVAILABLE!} \\
\protect\hyperlink{unitree-sdk2-cpp-idl-hg-motorcmd-reserve}{\passthrough{\lstinline!reserve!}}
& \passthrough{\lstinline!int!} & \passthrough{\lstinline!int!} &
\passthrough{\lstinline!AVAILABLE!} \\
\bottomrule
\end{longtable}

\hypertarget{unitree-sdk2-cpp-idl-hg-motorcmd-dunder-init-1}{%
\paragraph{\texorpdfstring{\texttt{unitree\_sdk2\_cpp.idl.hg.MotorCmd.\_\_init\_\_}}{unitree\_sdk2\_cpp.idl.hg.MotorCmd.\_\_init\_\_}}\label{unitree-sdk2-cpp-idl-hg-motorcmd-dunder-init-1}}

初始化 \passthrough{\lstinline!MotorCmd!}
实例。是否能实际构造取决于下方可用性状态。

\textbf{签名}

\begin{lstlisting}[language=Python]
def __init__(self) -> None
\end{lstlisting}

\textbf{可用性与安全性}

\textbf{\passthrough{\lstinline!AVAILABLE!} /
\passthrough{\lstinline!VALUE\_TYPE!}}：当前绑定源码已实现该消息值操作。

\textbf{参数}

无显式参数。实例方法中的 \passthrough{\lstinline!self!} 由 Python
自动传入。

\textbf{返回值}

\begin{longtable}[]{@{}
  >{\raggedright\arraybackslash}p{(\columnwidth - 4\tabcolsep) * \real{0.18}}
  >{\raggedright\arraybackslash}p{(\columnwidth - 4\tabcolsep) * \real{0.20}}
  >{\raggedright\arraybackslash}p{(\columnwidth - 4\tabcolsep) * \real{0.62}}@{}}
\toprule
位置 & 类型 & 含义 \\
\midrule
\endhead
无 & \passthrough{\lstinline!None!} &
\passthrough{\lstinline!\_\_init\_\_!}
本身不返回值；调用类对象时，在构造函数可用的前提下得到该类实例。 \\
\bottomrule
\end{longtable}

\textbf{对应 C++}

\begin{itemize}
\tightlist
\item
  类：\passthrough{\lstinline!unitree\_hg::msg::dds\_::MotorCmd\_!}
\item
  签名：\passthrough{\lstinline!MotorCmd\_()!}
\item
  绑定策略：\passthrough{\lstinline!未单独标注!}
\item
  声明位置：\passthrough{\lstinline!include/unitree/idl/hg/MotorCmd\_.hpp:32!}
\end{itemize}

\textbf{说明}

\begin{itemize}
\tightlist
\item
  示例是局部调用片段；\passthrough{\lstinline!obj!}
  和参数变量表示已经按上层业务校验并准备好的对象。
\end{itemize}

\textbf{用法}

\begin{lstlisting}[language=Python]
value = MotorCmd()
\end{lstlisting}

\hypertarget{unitree-sdk2-cpp-idl-hg-motorcmd-dunder-eq-1}{%
\paragraph{\texorpdfstring{\texttt{unitree\_sdk2\_cpp.idl.hg.MotorCmd.\_\_eq\_\_}}{unitree\_sdk2\_cpp.idl.hg.MotorCmd.\_\_eq\_\_}}\label{unitree-sdk2-cpp-idl-hg-motorcmd-dunder-eq-1}}

按底层 IDL 消息内容比较两个对象是否相等。

\textbf{签名}

\begin{lstlisting}[language=Python]
def __eq__(self, other: object) -> bool
\end{lstlisting}

\textbf{可用性与安全性}

\textbf{\passthrough{\lstinline!AVAILABLE!} /
\passthrough{\lstinline!VALUE\_TYPE!}}：当前绑定源码已实现该消息值操作。

\textbf{参数}

\begin{longtable}[]{@{}
  >{\raggedright\arraybackslash}p{(\columnwidth - 6\tabcolsep) * \real{0.18}}
  >{\raggedright\arraybackslash}p{(\columnwidth - 6\tabcolsep) * \real{0.24}}
  >{\raggedright\arraybackslash}p{(\columnwidth - 6\tabcolsep) * \real{0.18}}
  >{\raggedright\arraybackslash}p{(\columnwidth - 6\tabcolsep) * \real{0.40}}@{}}
\toprule
名称 & Python 类型 & 默认值 & 含义 \\
\midrule
\endhead
\passthrough{\lstinline!other!} & \passthrough{\lstinline!object!} &
必填 & 用于比较的另一个 Python 对象。类型不兼容时比较结果通常为
\passthrough{\lstinline!False!}。 \\
\bottomrule
\end{longtable}

\textbf{返回值}

\begin{longtable}[]{@{}
  >{\raggedright\arraybackslash}p{(\columnwidth - 4\tabcolsep) * \real{0.18}}
  >{\raggedright\arraybackslash}p{(\columnwidth - 4\tabcolsep) * \real{0.20}}
  >{\raggedright\arraybackslash}p{(\columnwidth - 4\tabcolsep) * \real{0.62}}@{}}
\toprule
位置 & 类型 & 含义 \\
\midrule
\endhead
返回值 & \passthrough{\lstinline!bool!} & 返回
\passthrough{\lstinline!bool!}。更精确的业务含义见该方法用途、C++
签名和目标型号协议。 \\
\bottomrule
\end{longtable}

\textbf{对应 C++}

\begin{itemize}
\tightlist
\item
  类：\passthrough{\lstinline!unitree\_hg::msg::dds\_::MotorCmd\_!}
\item
  签名：\passthrough{\lstinline!operator==(const value \&) const!}
\item
  绑定策略：\passthrough{\lstinline!未单独标注!}
\item
  声明位置：由 pybind11 手工绑定或生成注册表提供；manifest
  未记录独立头文件行号。
\end{itemize}

\textbf{说明}

\begin{itemize}
\tightlist
\item
  示例是局部调用片段；\passthrough{\lstinline!obj!}
  和参数变量表示已经按上层业务校验并准备好的对象。
\end{itemize}

\textbf{用法}

\begin{lstlisting}[language=Python]
same = left == right
\end{lstlisting}

\hypertarget{unitree-sdk2-cpp-idl-hg-motorcmd-dunder-ne-1}{%
\paragraph{\texorpdfstring{\texttt{unitree\_sdk2\_cpp.idl.hg.MotorCmd.\_\_ne\_\_}}{unitree\_sdk2\_cpp.idl.hg.MotorCmd.\_\_ne\_\_}}\label{unitree-sdk2-cpp-idl-hg-motorcmd-dunder-ne-1}}

按底层 IDL 消息内容比较两个对象是否不相等。

\textbf{签名}

\begin{lstlisting}[language=Python]
def __ne__(self, other: object) -> bool
\end{lstlisting}

\textbf{可用性与安全性}

\textbf{\passthrough{\lstinline!AVAILABLE!} /
\passthrough{\lstinline!VALUE\_TYPE!}}：当前绑定源码已实现该消息值操作。

\textbf{参数}

\begin{longtable}[]{@{}
  >{\raggedright\arraybackslash}p{(\columnwidth - 6\tabcolsep) * \real{0.18}}
  >{\raggedright\arraybackslash}p{(\columnwidth - 6\tabcolsep) * \real{0.24}}
  >{\raggedright\arraybackslash}p{(\columnwidth - 6\tabcolsep) * \real{0.18}}
  >{\raggedright\arraybackslash}p{(\columnwidth - 6\tabcolsep) * \real{0.40}}@{}}
\toprule
名称 & Python 类型 & 默认值 & 含义 \\
\midrule
\endhead
\passthrough{\lstinline!other!} & \passthrough{\lstinline!object!} &
必填 & 用于比较的另一个 Python 对象。类型不兼容时比较结果通常为
\passthrough{\lstinline!False!}。 \\
\bottomrule
\end{longtable}

\textbf{返回值}

\begin{longtable}[]{@{}
  >{\raggedright\arraybackslash}p{(\columnwidth - 4\tabcolsep) * \real{0.18}}
  >{\raggedright\arraybackslash}p{(\columnwidth - 4\tabcolsep) * \real{0.20}}
  >{\raggedright\arraybackslash}p{(\columnwidth - 4\tabcolsep) * \real{0.62}}@{}}
\toprule
位置 & 类型 & 含义 \\
\midrule
\endhead
返回值 & \passthrough{\lstinline!bool!} & 返回
\passthrough{\lstinline!bool!}。更精确的业务含义见该方法用途、C++
签名和目标型号协议。 \\
\bottomrule
\end{longtable}

\textbf{对应 C++}

\begin{itemize}
\tightlist
\item
  类：\passthrough{\lstinline!unitree\_hg::msg::dds\_::MotorCmd\_!}
\item
  签名：\passthrough{\lstinline"operator!=(const value \&) const"}
\item
  绑定策略：\passthrough{\lstinline!未单独标注!}
\item
  声明位置：由 pybind11 手工绑定或生成注册表提供；manifest
  未记录独立头文件行号。
\end{itemize}

\textbf{说明}

\begin{itemize}
\tightlist
\item
  示例是局部调用片段；\passthrough{\lstinline!obj!}
  和参数变量表示已经按上层业务校验并准备好的对象。
\end{itemize}

\textbf{用法}

\begin{lstlisting}[language=Python]
different = left != right
\end{lstlisting}

\hypertarget{unitree-sdk2-cpp-idl-hg-motorcmd-mode}{%
\paragraph{\texorpdfstring{\texttt{unitree\_sdk2\_cpp.idl.hg.MotorCmd.mode}}{unitree\_sdk2\_cpp.idl.hg.MotorCmd.mode}}\label{unitree-sdk2-cpp-idl-hg-motorcmd-mode}}

底层 \passthrough{\lstinline!unitree\_hg::msg::dds\_::MotorCmd\_!} 的
\passthrough{\lstinline!mode!} 字段。类型签名只说明 Python
表示；单位、数值范围、数组固定长度和枚举含义需查对应 IDL 头文件。
取值范围为 0 到 255。

\textbf{签名}

\begin{lstlisting}[language=Python]
@property
def mode(self) -> int

@mode.setter
def mode(self, value: int) -> None
\end{lstlisting}

\textbf{可用性与安全性}

\textbf{\passthrough{\lstinline!AVAILABLE!} /
\passthrough{\lstinline!VALUE\_TYPE!}}：当前绑定源码已实现该消息值操作。

\textbf{参数}

\begin{longtable}[]{@{}
  >{\raggedright\arraybackslash}p{(\columnwidth - 4\tabcolsep) * \real{0.18}}
  >{\raggedright\arraybackslash}p{(\columnwidth - 4\tabcolsep) * \real{0.20}}
  >{\raggedright\arraybackslash}p{(\columnwidth - 4\tabcolsep) * \real{0.62}}@{}}
\toprule
名称 & Python 类型 & 含义 \\
\midrule
\endhead
\passthrough{\lstinline!value!} & \passthrough{\lstinline!int!} &
要写入或传递的新值，必须符合该参数的 Python 类型以及底层 C++
范围约束。 \\
\bottomrule
\end{longtable}

\textbf{返回值}

\begin{longtable}[]{@{}
  >{\raggedright\arraybackslash}p{(\columnwidth - 4\tabcolsep) * \real{0.18}}
  >{\raggedright\arraybackslash}p{(\columnwidth - 4\tabcolsep) * \real{0.20}}
  >{\raggedright\arraybackslash}p{(\columnwidth - 4\tabcolsep) * \real{0.62}}@{}}
\toprule
位置 & 类型 & 含义 \\
\midrule
\endhead
getter 返回值 & \passthrough{\lstinline!int!} & 返回属性当前值。 \\
\bottomrule
\end{longtable}

\textbf{对应 C++}

\begin{itemize}
\tightlist
\item
  类：\passthrough{\lstinline!unitree\_hg::msg::dds\_::MotorCmd\_!}
\item
  字段类型：\passthrough{\lstinline!uint8\_t!}
\item
  声明位置：\passthrough{\lstinline!include/unitree/idl/hg/MotorCmd\_.hpp:23!}
\end{itemize}

\textbf{用法}

\begin{lstlisting}[language=Python]
current_value = obj.mode
obj.mode = new_value
\end{lstlisting}

\hypertarget{unitree-sdk2-cpp-idl-hg-motorcmd-q}{%
\paragraph{\texorpdfstring{\texttt{unitree\_sdk2\_cpp.idl.hg.MotorCmd.q}}{unitree\_sdk2\_cpp.idl.hg.MotorCmd.q}}\label{unitree-sdk2-cpp-idl-hg-motorcmd-q}}

底层 \passthrough{\lstinline!unitree\_hg::msg::dds\_::MotorCmd\_!} 的
\passthrough{\lstinline!q!} 字段。类型签名只说明 Python
表示；单位、数值范围、数组固定长度和枚举含义需查对应 IDL 头文件。
底层为浮点数；类型本身不说明单位、坐标系或业务安全范围。

\textbf{签名}

\begin{lstlisting}[language=Python]
@property
def q(self) -> float

@q.setter
def q(self, value: float) -> None
\end{lstlisting}

\textbf{可用性与安全性}

\textbf{\passthrough{\lstinline!AVAILABLE!} /
\passthrough{\lstinline!VALUE\_TYPE!}}：当前绑定源码已实现该消息值操作。

\textbf{参数}

\begin{longtable}[]{@{}
  >{\raggedright\arraybackslash}p{(\columnwidth - 4\tabcolsep) * \real{0.18}}
  >{\raggedright\arraybackslash}p{(\columnwidth - 4\tabcolsep) * \real{0.20}}
  >{\raggedright\arraybackslash}p{(\columnwidth - 4\tabcolsep) * \real{0.62}}@{}}
\toprule
名称 & Python 类型 & 含义 \\
\midrule
\endhead
\passthrough{\lstinline!value!} & \passthrough{\lstinline!float!} &
要写入或传递的新值，必须符合该参数的 Python 类型以及底层 C++
范围约束。 \\
\bottomrule
\end{longtable}

\textbf{返回值}

\begin{longtable}[]{@{}
  >{\raggedright\arraybackslash}p{(\columnwidth - 4\tabcolsep) * \real{0.18}}
  >{\raggedright\arraybackslash}p{(\columnwidth - 4\tabcolsep) * \real{0.20}}
  >{\raggedright\arraybackslash}p{(\columnwidth - 4\tabcolsep) * \real{0.62}}@{}}
\toprule
位置 & 类型 & 含义 \\
\midrule
\endhead
getter 返回值 & \passthrough{\lstinline!float!} & 返回属性当前值。 \\
\bottomrule
\end{longtable}

\textbf{对应 C++}

\begin{itemize}
\tightlist
\item
  类：\passthrough{\lstinline!unitree\_hg::msg::dds\_::MotorCmd\_!}
\item
  字段类型：\passthrough{\lstinline!float!}
\item
  声明位置：\passthrough{\lstinline!include/unitree/idl/hg/MotorCmd\_.hpp:24!}
\end{itemize}

\textbf{用法}

\begin{lstlisting}[language=Python]
current_value = obj.q
obj.q = new_value
\end{lstlisting}

\hypertarget{unitree-sdk2-cpp-idl-hg-motorcmd-dq}{%
\paragraph{\texorpdfstring{\texttt{unitree\_sdk2\_cpp.idl.hg.MotorCmd.dq}}{unitree\_sdk2\_cpp.idl.hg.MotorCmd.dq}}\label{unitree-sdk2-cpp-idl-hg-motorcmd-dq}}

底层 \passthrough{\lstinline!unitree\_hg::msg::dds\_::MotorCmd\_!} 的
\passthrough{\lstinline!dq!} 字段。类型签名只说明 Python
表示；单位、数值范围、数组固定长度和枚举含义需查对应 IDL 头文件。
底层为浮点数；类型本身不说明单位、坐标系或业务安全范围。

\textbf{签名}

\begin{lstlisting}[language=Python]
@property
def dq(self) -> float

@dq.setter
def dq(self, value: float) -> None
\end{lstlisting}

\textbf{可用性与安全性}

\textbf{\passthrough{\lstinline!AVAILABLE!} /
\passthrough{\lstinline!VALUE\_TYPE!}}：当前绑定源码已实现该消息值操作。

\textbf{参数}

\begin{longtable}[]{@{}
  >{\raggedright\arraybackslash}p{(\columnwidth - 4\tabcolsep) * \real{0.18}}
  >{\raggedright\arraybackslash}p{(\columnwidth - 4\tabcolsep) * \real{0.20}}
  >{\raggedright\arraybackslash}p{(\columnwidth - 4\tabcolsep) * \real{0.62}}@{}}
\toprule
名称 & Python 类型 & 含义 \\
\midrule
\endhead
\passthrough{\lstinline!value!} & \passthrough{\lstinline!float!} &
要写入或传递的新值，必须符合该参数的 Python 类型以及底层 C++
范围约束。 \\
\bottomrule
\end{longtable}

\textbf{返回值}

\begin{longtable}[]{@{}
  >{\raggedright\arraybackslash}p{(\columnwidth - 4\tabcolsep) * \real{0.18}}
  >{\raggedright\arraybackslash}p{(\columnwidth - 4\tabcolsep) * \real{0.20}}
  >{\raggedright\arraybackslash}p{(\columnwidth - 4\tabcolsep) * \real{0.62}}@{}}
\toprule
位置 & 类型 & 含义 \\
\midrule
\endhead
getter 返回值 & \passthrough{\lstinline!float!} & 返回属性当前值。 \\
\bottomrule
\end{longtable}

\textbf{对应 C++}

\begin{itemize}
\tightlist
\item
  类：\passthrough{\lstinline!unitree\_hg::msg::dds\_::MotorCmd\_!}
\item
  字段类型：\passthrough{\lstinline!float!}
\item
  声明位置：\passthrough{\lstinline!include/unitree/idl/hg/MotorCmd\_.hpp:25!}
\end{itemize}

\textbf{用法}

\begin{lstlisting}[language=Python]
current_value = obj.dq
obj.dq = new_value
\end{lstlisting}

\hypertarget{unitree-sdk2-cpp-idl-hg-motorcmd-tau}{%
\paragraph{\texorpdfstring{\texttt{unitree\_sdk2\_cpp.idl.hg.MotorCmd.tau}}{unitree\_sdk2\_cpp.idl.hg.MotorCmd.tau}}\label{unitree-sdk2-cpp-idl-hg-motorcmd-tau}}

底层 \passthrough{\lstinline!unitree\_hg::msg::dds\_::MotorCmd\_!} 的
\passthrough{\lstinline!tau!} 字段。类型签名只说明 Python
表示；单位、数值范围、数组固定长度和枚举含义需查对应 IDL 头文件。
底层为浮点数；类型本身不说明单位、坐标系或业务安全范围。

\textbf{签名}

\begin{lstlisting}[language=Python]
@property
def tau(self) -> float

@tau.setter
def tau(self, value: float) -> None
\end{lstlisting}

\textbf{可用性与安全性}

\textbf{\passthrough{\lstinline!AVAILABLE!} /
\passthrough{\lstinline!VALUE\_TYPE!}}：当前绑定源码已实现该消息值操作。

\textbf{参数}

\begin{longtable}[]{@{}
  >{\raggedright\arraybackslash}p{(\columnwidth - 4\tabcolsep) * \real{0.18}}
  >{\raggedright\arraybackslash}p{(\columnwidth - 4\tabcolsep) * \real{0.20}}
  >{\raggedright\arraybackslash}p{(\columnwidth - 4\tabcolsep) * \real{0.62}}@{}}
\toprule
名称 & Python 类型 & 含义 \\
\midrule
\endhead
\passthrough{\lstinline!value!} & \passthrough{\lstinline!float!} &
要写入或传递的新值，必须符合该参数的 Python 类型以及底层 C++
范围约束。 \\
\bottomrule
\end{longtable}

\textbf{返回值}

\begin{longtable}[]{@{}
  >{\raggedright\arraybackslash}p{(\columnwidth - 4\tabcolsep) * \real{0.18}}
  >{\raggedright\arraybackslash}p{(\columnwidth - 4\tabcolsep) * \real{0.20}}
  >{\raggedright\arraybackslash}p{(\columnwidth - 4\tabcolsep) * \real{0.62}}@{}}
\toprule
位置 & 类型 & 含义 \\
\midrule
\endhead
getter 返回值 & \passthrough{\lstinline!float!} & 返回属性当前值。 \\
\bottomrule
\end{longtable}

\textbf{对应 C++}

\begin{itemize}
\tightlist
\item
  类：\passthrough{\lstinline!unitree\_hg::msg::dds\_::MotorCmd\_!}
\item
  字段类型：\passthrough{\lstinline!float!}
\item
  声明位置：\passthrough{\lstinline!include/unitree/idl/hg/MotorCmd\_.hpp:26!}
\end{itemize}

\textbf{用法}

\begin{lstlisting}[language=Python]
current_value = obj.tau
obj.tau = new_value
\end{lstlisting}

\hypertarget{unitree-sdk2-cpp-idl-hg-motorcmd-kp}{%
\paragraph{\texorpdfstring{\texttt{unitree\_sdk2\_cpp.idl.hg.MotorCmd.kp}}{unitree\_sdk2\_cpp.idl.hg.MotorCmd.kp}}\label{unitree-sdk2-cpp-idl-hg-motorcmd-kp}}

底层 \passthrough{\lstinline!unitree\_hg::msg::dds\_::MotorCmd\_!} 的
\passthrough{\lstinline!kp!} 字段。类型签名只说明 Python
表示；单位、数值范围、数组固定长度和枚举含义需查对应 IDL 头文件。
底层为浮点数；类型本身不说明单位、坐标系或业务安全范围。

\textbf{签名}

\begin{lstlisting}[language=Python]
@property
def kp(self) -> float

@kp.setter
def kp(self, value: float) -> None
\end{lstlisting}

\textbf{可用性与安全性}

\textbf{\passthrough{\lstinline!AVAILABLE!} /
\passthrough{\lstinline!VALUE\_TYPE!}}：当前绑定源码已实现该消息值操作。

\textbf{参数}

\begin{longtable}[]{@{}
  >{\raggedright\arraybackslash}p{(\columnwidth - 4\tabcolsep) * \real{0.18}}
  >{\raggedright\arraybackslash}p{(\columnwidth - 4\tabcolsep) * \real{0.20}}
  >{\raggedright\arraybackslash}p{(\columnwidth - 4\tabcolsep) * \real{0.62}}@{}}
\toprule
名称 & Python 类型 & 含义 \\
\midrule
\endhead
\passthrough{\lstinline!value!} & \passthrough{\lstinline!float!} &
要写入或传递的新值，必须符合该参数的 Python 类型以及底层 C++
范围约束。 \\
\bottomrule
\end{longtable}

\textbf{返回值}

\begin{longtable}[]{@{}
  >{\raggedright\arraybackslash}p{(\columnwidth - 4\tabcolsep) * \real{0.18}}
  >{\raggedright\arraybackslash}p{(\columnwidth - 4\tabcolsep) * \real{0.20}}
  >{\raggedright\arraybackslash}p{(\columnwidth - 4\tabcolsep) * \real{0.62}}@{}}
\toprule
位置 & 类型 & 含义 \\
\midrule
\endhead
getter 返回值 & \passthrough{\lstinline!float!} & 返回属性当前值。 \\
\bottomrule
\end{longtable}

\textbf{对应 C++}

\begin{itemize}
\tightlist
\item
  类：\passthrough{\lstinline!unitree\_hg::msg::dds\_::MotorCmd\_!}
\item
  字段类型：\passthrough{\lstinline!float!}
\item
  声明位置：\passthrough{\lstinline!include/unitree/idl/hg/MotorCmd\_.hpp:27!}
\end{itemize}

\textbf{用法}

\begin{lstlisting}[language=Python]
current_value = obj.kp
obj.kp = new_value
\end{lstlisting}

\hypertarget{unitree-sdk2-cpp-idl-hg-motorcmd-kd}{%
\paragraph{\texorpdfstring{\texttt{unitree\_sdk2\_cpp.idl.hg.MotorCmd.kd}}{unitree\_sdk2\_cpp.idl.hg.MotorCmd.kd}}\label{unitree-sdk2-cpp-idl-hg-motorcmd-kd}}

底层 \passthrough{\lstinline!unitree\_hg::msg::dds\_::MotorCmd\_!} 的
\passthrough{\lstinline!kd!} 字段。类型签名只说明 Python
表示；单位、数值范围、数组固定长度和枚举含义需查对应 IDL 头文件。
底层为浮点数；类型本身不说明单位、坐标系或业务安全范围。

\textbf{签名}

\begin{lstlisting}[language=Python]
@property
def kd(self) -> float

@kd.setter
def kd(self, value: float) -> None
\end{lstlisting}

\textbf{可用性与安全性}

\textbf{\passthrough{\lstinline!AVAILABLE!} /
\passthrough{\lstinline!VALUE\_TYPE!}}：当前绑定源码已实现该消息值操作。

\textbf{参数}

\begin{longtable}[]{@{}
  >{\raggedright\arraybackslash}p{(\columnwidth - 4\tabcolsep) * \real{0.18}}
  >{\raggedright\arraybackslash}p{(\columnwidth - 4\tabcolsep) * \real{0.20}}
  >{\raggedright\arraybackslash}p{(\columnwidth - 4\tabcolsep) * \real{0.62}}@{}}
\toprule
名称 & Python 类型 & 含义 \\
\midrule
\endhead
\passthrough{\lstinline!value!} & \passthrough{\lstinline!float!} &
要写入或传递的新值，必须符合该参数的 Python 类型以及底层 C++
范围约束。 \\
\bottomrule
\end{longtable}

\textbf{返回值}

\begin{longtable}[]{@{}
  >{\raggedright\arraybackslash}p{(\columnwidth - 4\tabcolsep) * \real{0.18}}
  >{\raggedright\arraybackslash}p{(\columnwidth - 4\tabcolsep) * \real{0.20}}
  >{\raggedright\arraybackslash}p{(\columnwidth - 4\tabcolsep) * \real{0.62}}@{}}
\toprule
位置 & 类型 & 含义 \\
\midrule
\endhead
getter 返回值 & \passthrough{\lstinline!float!} & 返回属性当前值。 \\
\bottomrule
\end{longtable}

\textbf{对应 C++}

\begin{itemize}
\tightlist
\item
  类：\passthrough{\lstinline!unitree\_hg::msg::dds\_::MotorCmd\_!}
\item
  字段类型：\passthrough{\lstinline!float!}
\item
  声明位置：\passthrough{\lstinline!include/unitree/idl/hg/MotorCmd\_.hpp:28!}
\end{itemize}

\textbf{用法}

\begin{lstlisting}[language=Python]
current_value = obj.kd
obj.kd = new_value
\end{lstlisting}

\hypertarget{unitree-sdk2-cpp-idl-hg-motorcmd-reserve}{%
\paragraph{\texorpdfstring{\texttt{unitree\_sdk2\_cpp.idl.hg.MotorCmd.reserve}}{unitree\_sdk2\_cpp.idl.hg.MotorCmd.reserve}}\label{unitree-sdk2-cpp-idl-hg-motorcmd-reserve}}

底层 \passthrough{\lstinline!unitree\_hg::msg::dds\_::MotorCmd\_!} 的
\passthrough{\lstinline!reserve!} 字段。类型签名只说明 Python
表示；单位、数值范围、数组固定长度和枚举含义需查对应 IDL 头文件。
取值范围为 0 到 4294967295。

\textbf{签名}

\begin{lstlisting}[language=Python]
@property
def reserve(self) -> int

@reserve.setter
def reserve(self, value: int) -> None
\end{lstlisting}

\textbf{可用性与安全性}

\textbf{\passthrough{\lstinline!AVAILABLE!} /
\passthrough{\lstinline!VALUE\_TYPE!}}：当前绑定源码已实现该消息值操作。

\textbf{参数}

\begin{longtable}[]{@{}
  >{\raggedright\arraybackslash}p{(\columnwidth - 4\tabcolsep) * \real{0.18}}
  >{\raggedright\arraybackslash}p{(\columnwidth - 4\tabcolsep) * \real{0.20}}
  >{\raggedright\arraybackslash}p{(\columnwidth - 4\tabcolsep) * \real{0.62}}@{}}
\toprule
名称 & Python 类型 & 含义 \\
\midrule
\endhead
\passthrough{\lstinline!value!} & \passthrough{\lstinline!int!} &
要写入或传递的新值，必须符合该参数的 Python 类型以及底层 C++
范围约束。 \\
\bottomrule
\end{longtable}

\textbf{返回值}

\begin{longtable}[]{@{}
  >{\raggedright\arraybackslash}p{(\columnwidth - 4\tabcolsep) * \real{0.18}}
  >{\raggedright\arraybackslash}p{(\columnwidth - 4\tabcolsep) * \real{0.20}}
  >{\raggedright\arraybackslash}p{(\columnwidth - 4\tabcolsep) * \real{0.62}}@{}}
\toprule
位置 & 类型 & 含义 \\
\midrule
\endhead
getter 返回值 & \passthrough{\lstinline!int!} & 返回属性当前值。 \\
\bottomrule
\end{longtable}

\textbf{对应 C++}

\begin{itemize}
\tightlist
\item
  类：\passthrough{\lstinline!unitree\_hg::msg::dds\_::MotorCmd\_!}
\item
  字段类型：\passthrough{\lstinline!uint32\_t!}
\item
  声明位置：\passthrough{\lstinline!include/unitree/idl/hg/MotorCmd\_.hpp:29!}
\end{itemize}

\textbf{用法}

\begin{lstlisting}[language=Python]
current_value = obj.reserve
obj.reserve = new_value
\end{lstlisting}

\hypertarget{unitree-sdk2-cpp-idl-hg-handcmd}{%
\subsubsection{\texorpdfstring{\texttt{unitree\_sdk2\_cpp.idl.hg.HandCmd}}{unitree\_sdk2\_cpp.idl.hg.HandCmd}}\label{unitree-sdk2-cpp-idl-hg-handcmd}}

可由 Python 读写的 DDS/IDL 消息值类型。字段采用复制语义。

\textbf{导入}

\begin{lstlisting}[language=Python]
from unitree_sdk2_cpp.idl.hg import HandCmd
\end{lstlisting}

\textbf{方法索引}

\begin{longtable}[]{@{}
  >{\raggedright\arraybackslash}p{(\columnwidth - 6\tabcolsep) * \real{0.18}}
  >{\raggedright\arraybackslash}p{(\columnwidth - 6\tabcolsep) * \real{0.24}}
  >{\raggedright\arraybackslash}p{(\columnwidth - 6\tabcolsep) * \real{0.18}}
  >{\raggedright\arraybackslash}p{(\columnwidth - 6\tabcolsep) * \real{0.40}}@{}}
\toprule
方法 & Python 签名 & 状态 & 安全分类 \\
\midrule
\endhead
\protect\hyperlink{unitree-sdk2-cpp-idl-hg-handcmd-dunder-init-1}{\passthrough{\lstinline!\_\_init\_\_!}}
& \passthrough{\lstinline!def \_\_init\_\_(self) -> None!} &
\passthrough{\lstinline!AVAILABLE!} &
\passthrough{\lstinline!VALUE\_TYPE!} \\
\protect\hyperlink{unitree-sdk2-cpp-idl-hg-handcmd-dunder-eq-1}{\passthrough{\lstinline!\_\_eq\_\_!}}
& \passthrough{\lstinline!def \_\_eq\_\_(self, other: object) -> bool!}
& \passthrough{\lstinline!AVAILABLE!} &
\passthrough{\lstinline!VALUE\_TYPE!} \\
\protect\hyperlink{unitree-sdk2-cpp-idl-hg-handcmd-dunder-ne-1}{\passthrough{\lstinline!\_\_ne\_\_!}}
& \passthrough{\lstinline!def \_\_ne\_\_(self, other: object) -> bool!}
& \passthrough{\lstinline!AVAILABLE!} &
\passthrough{\lstinline!VALUE\_TYPE!} \\
\bottomrule
\end{longtable}

\textbf{属性索引}

\begin{longtable}[]{@{}
  >{\raggedright\arraybackslash}p{(\columnwidth - 6\tabcolsep) * \real{0.18}}
  >{\raggedright\arraybackslash}p{(\columnwidth - 6\tabcolsep) * \real{0.24}}
  >{\raggedright\arraybackslash}p{(\columnwidth - 6\tabcolsep) * \real{0.18}}
  >{\raggedright\arraybackslash}p{(\columnwidth - 6\tabcolsep) * \real{0.40}}@{}}
\toprule
属性 & 读取类型 & 写入类型 & 状态 \\
\midrule
\endhead
\protect\hyperlink{unitree-sdk2-cpp-idl-hg-handcmd-motor-cmd}{\passthrough{\lstinline!motor\_cmd!}}
& \passthrough{\lstinline!list[MotorCmd]!} &
\passthrough{\lstinline!Sequence[MotorCmd]!} &
\passthrough{\lstinline!AVAILABLE!} \\
\protect\hyperlink{unitree-sdk2-cpp-idl-hg-handcmd-reserve}{\passthrough{\lstinline!reserve!}}
& \passthrough{\lstinline!list[int]!} &
\passthrough{\lstinline!Sequence[int]!} &
\passthrough{\lstinline!AVAILABLE!} \\
\bottomrule
\end{longtable}

\hypertarget{unitree-sdk2-cpp-idl-hg-handcmd-dunder-init-1}{%
\paragraph{\texorpdfstring{\texttt{unitree\_sdk2\_cpp.idl.hg.HandCmd.\_\_init\_\_}}{unitree\_sdk2\_cpp.idl.hg.HandCmd.\_\_init\_\_}}\label{unitree-sdk2-cpp-idl-hg-handcmd-dunder-init-1}}

初始化 \passthrough{\lstinline!HandCmd!}
实例。是否能实际构造取决于下方可用性状态。

\textbf{签名}

\begin{lstlisting}[language=Python]
def __init__(self) -> None
\end{lstlisting}

\textbf{可用性与安全性}

\textbf{\passthrough{\lstinline!AVAILABLE!} /
\passthrough{\lstinline!VALUE\_TYPE!}}：当前绑定源码已实现该消息值操作。

\textbf{参数}

无显式参数。实例方法中的 \passthrough{\lstinline!self!} 由 Python
自动传入。

\textbf{返回值}

\begin{longtable}[]{@{}
  >{\raggedright\arraybackslash}p{(\columnwidth - 4\tabcolsep) * \real{0.18}}
  >{\raggedright\arraybackslash}p{(\columnwidth - 4\tabcolsep) * \real{0.20}}
  >{\raggedright\arraybackslash}p{(\columnwidth - 4\tabcolsep) * \real{0.62}}@{}}
\toprule
位置 & 类型 & 含义 \\
\midrule
\endhead
无 & \passthrough{\lstinline!None!} &
\passthrough{\lstinline!\_\_init\_\_!}
本身不返回值；调用类对象时，在构造函数可用的前提下得到该类实例。 \\
\bottomrule
\end{longtable}

\textbf{对应 C++}

\begin{itemize}
\tightlist
\item
  类：\passthrough{\lstinline!unitree\_hg::msg::dds\_::HandCmd\_!}
\item
  签名：\passthrough{\lstinline!HandCmd\_()!}
\item
  绑定策略：\passthrough{\lstinline!未单独标注!}
\item
  声明位置：\passthrough{\lstinline!include/unitree/idl/hg/HandCmd\_.hpp:31!}
\end{itemize}

\textbf{说明}

\begin{itemize}
\tightlist
\item
  示例是局部调用片段；\passthrough{\lstinline!obj!}
  和参数变量表示已经按上层业务校验并准备好的对象。
\end{itemize}

\textbf{用法}

\begin{lstlisting}[language=Python]
value = HandCmd()
\end{lstlisting}

\hypertarget{unitree-sdk2-cpp-idl-hg-handcmd-dunder-eq-1}{%
\paragraph{\texorpdfstring{\texttt{unitree\_sdk2\_cpp.idl.hg.HandCmd.\_\_eq\_\_}}{unitree\_sdk2\_cpp.idl.hg.HandCmd.\_\_eq\_\_}}\label{unitree-sdk2-cpp-idl-hg-handcmd-dunder-eq-1}}

按底层 IDL 消息内容比较两个对象是否相等。

\textbf{签名}

\begin{lstlisting}[language=Python]
def __eq__(self, other: object) -> bool
\end{lstlisting}

\textbf{可用性与安全性}

\textbf{\passthrough{\lstinline!AVAILABLE!} /
\passthrough{\lstinline!VALUE\_TYPE!}}：当前绑定源码已实现该消息值操作。

\textbf{参数}

\begin{longtable}[]{@{}
  >{\raggedright\arraybackslash}p{(\columnwidth - 6\tabcolsep) * \real{0.18}}
  >{\raggedright\arraybackslash}p{(\columnwidth - 6\tabcolsep) * \real{0.24}}
  >{\raggedright\arraybackslash}p{(\columnwidth - 6\tabcolsep) * \real{0.18}}
  >{\raggedright\arraybackslash}p{(\columnwidth - 6\tabcolsep) * \real{0.40}}@{}}
\toprule
名称 & Python 类型 & 默认值 & 含义 \\
\midrule
\endhead
\passthrough{\lstinline!other!} & \passthrough{\lstinline!object!} &
必填 & 用于比较的另一个 Python 对象。类型不兼容时比较结果通常为
\passthrough{\lstinline!False!}。 \\
\bottomrule
\end{longtable}

\textbf{返回值}

\begin{longtable}[]{@{}
  >{\raggedright\arraybackslash}p{(\columnwidth - 4\tabcolsep) * \real{0.18}}
  >{\raggedright\arraybackslash}p{(\columnwidth - 4\tabcolsep) * \real{0.20}}
  >{\raggedright\arraybackslash}p{(\columnwidth - 4\tabcolsep) * \real{0.62}}@{}}
\toprule
位置 & 类型 & 含义 \\
\midrule
\endhead
返回值 & \passthrough{\lstinline!bool!} & 返回
\passthrough{\lstinline!bool!}。更精确的业务含义见该方法用途、C++
签名和目标型号协议。 \\
\bottomrule
\end{longtable}

\textbf{对应 C++}

\begin{itemize}
\tightlist
\item
  类：\passthrough{\lstinline!unitree\_hg::msg::dds\_::HandCmd\_!}
\item
  签名：\passthrough{\lstinline!operator==(const value \&) const!}
\item
  绑定策略：\passthrough{\lstinline!未单独标注!}
\item
  声明位置：由 pybind11 手工绑定或生成注册表提供；manifest
  未记录独立头文件行号。
\end{itemize}

\textbf{说明}

\begin{itemize}
\tightlist
\item
  示例是局部调用片段；\passthrough{\lstinline!obj!}
  和参数变量表示已经按上层业务校验并准备好的对象。
\end{itemize}

\textbf{用法}

\begin{lstlisting}[language=Python]
same = left == right
\end{lstlisting}

\hypertarget{unitree-sdk2-cpp-idl-hg-handcmd-dunder-ne-1}{%
\paragraph{\texorpdfstring{\texttt{unitree\_sdk2\_cpp.idl.hg.HandCmd.\_\_ne\_\_}}{unitree\_sdk2\_cpp.idl.hg.HandCmd.\_\_ne\_\_}}\label{unitree-sdk2-cpp-idl-hg-handcmd-dunder-ne-1}}

按底层 IDL 消息内容比较两个对象是否不相等。

\textbf{签名}

\begin{lstlisting}[language=Python]
def __ne__(self, other: object) -> bool
\end{lstlisting}

\textbf{可用性与安全性}

\textbf{\passthrough{\lstinline!AVAILABLE!} /
\passthrough{\lstinline!VALUE\_TYPE!}}：当前绑定源码已实现该消息值操作。

\textbf{参数}

\begin{longtable}[]{@{}
  >{\raggedright\arraybackslash}p{(\columnwidth - 6\tabcolsep) * \real{0.18}}
  >{\raggedright\arraybackslash}p{(\columnwidth - 6\tabcolsep) * \real{0.24}}
  >{\raggedright\arraybackslash}p{(\columnwidth - 6\tabcolsep) * \real{0.18}}
  >{\raggedright\arraybackslash}p{(\columnwidth - 6\tabcolsep) * \real{0.40}}@{}}
\toprule
名称 & Python 类型 & 默认值 & 含义 \\
\midrule
\endhead
\passthrough{\lstinline!other!} & \passthrough{\lstinline!object!} &
必填 & 用于比较的另一个 Python 对象。类型不兼容时比较结果通常为
\passthrough{\lstinline!False!}。 \\
\bottomrule
\end{longtable}

\textbf{返回值}

\begin{longtable}[]{@{}
  >{\raggedright\arraybackslash}p{(\columnwidth - 4\tabcolsep) * \real{0.18}}
  >{\raggedright\arraybackslash}p{(\columnwidth - 4\tabcolsep) * \real{0.20}}
  >{\raggedright\arraybackslash}p{(\columnwidth - 4\tabcolsep) * \real{0.62}}@{}}
\toprule
位置 & 类型 & 含义 \\
\midrule
\endhead
返回值 & \passthrough{\lstinline!bool!} & 返回
\passthrough{\lstinline!bool!}。更精确的业务含义见该方法用途、C++
签名和目标型号协议。 \\
\bottomrule
\end{longtable}

\textbf{对应 C++}

\begin{itemize}
\tightlist
\item
  类：\passthrough{\lstinline!unitree\_hg::msg::dds\_::HandCmd\_!}
\item
  签名：\passthrough{\lstinline"operator!=(const value \&) const"}
\item
  绑定策略：\passthrough{\lstinline!未单独标注!}
\item
  声明位置：由 pybind11 手工绑定或生成注册表提供；manifest
  未记录独立头文件行号。
\end{itemize}

\textbf{说明}

\begin{itemize}
\tightlist
\item
  示例是局部调用片段；\passthrough{\lstinline!obj!}
  和参数变量表示已经按上层业务校验并准备好的对象。
\end{itemize}

\textbf{用法}

\begin{lstlisting}[language=Python]
different = left != right
\end{lstlisting}

\hypertarget{unitree-sdk2-cpp-idl-hg-handcmd-motor-cmd}{%
\paragraph{\texorpdfstring{\texttt{unitree\_sdk2\_cpp.idl.hg.HandCmd.motor\_cmd}}{unitree\_sdk2\_cpp.idl.hg.HandCmd.motor\_cmd}}\label{unitree-sdk2-cpp-idl-hg-handcmd-motor-cmd}}

底层 \passthrough{\lstinline!unitree\_hg::msg::dds\_::HandCmd\_!} 的
\passthrough{\lstinline!motor\_cmd!} 字段。类型签名只说明 Python
表示；单位、数值范围、数组固定长度和枚举含义需查对应 IDL 头文件。
底层是可变长度 vector

\textbf{签名}

\begin{lstlisting}[language=Python]
@property
def motor_cmd(self) -> list[MotorCmd]

@motor_cmd.setter
def motor_cmd(self, value: Sequence[MotorCmd]) -> None
\end{lstlisting}

\textbf{可用性与安全性}

\textbf{\passthrough{\lstinline!AVAILABLE!} /
\passthrough{\lstinline!VALUE\_TYPE!}}：当前绑定源码已实现该消息值操作。

\textbf{参数}

\begin{longtable}[]{@{}
  >{\raggedright\arraybackslash}p{(\columnwidth - 4\tabcolsep) * \real{0.18}}
  >{\raggedright\arraybackslash}p{(\columnwidth - 4\tabcolsep) * \real{0.20}}
  >{\raggedright\arraybackslash}p{(\columnwidth - 4\tabcolsep) * \real{0.62}}@{}}
\toprule
名称 & Python 类型 & 含义 \\
\midrule
\endhead
\passthrough{\lstinline!value!} &
\passthrough{\lstinline!Sequence[MotorCmd]!} &
要写入或传递的新值，必须符合该参数的 Python 类型以及底层 C++
范围约束。 \\
\bottomrule
\end{longtable}

\textbf{返回值}

\begin{longtable}[]{@{}
  >{\raggedright\arraybackslash}p{(\columnwidth - 4\tabcolsep) * \real{0.18}}
  >{\raggedright\arraybackslash}p{(\columnwidth - 4\tabcolsep) * \real{0.20}}
  >{\raggedright\arraybackslash}p{(\columnwidth - 4\tabcolsep) * \real{0.62}}@{}}
\toprule
位置 & 类型 & 含义 \\
\midrule
\endhead
getter 返回值 & \passthrough{\lstinline!list[MotorCmd]!} &
返回属性当前值。数组、vector 和嵌套消息采用复制语义。 \\
\bottomrule
\end{longtable}

\textbf{对应 C++}

\begin{itemize}
\tightlist
\item
  类：\passthrough{\lstinline!unitree\_hg::msg::dds\_::HandCmd\_!}
\item
  字段类型：\passthrough{\lstinline!std::vector< ::unitree\_hg::msg::dds\_::MotorCmd\_>!}
\item
  声明位置：\passthrough{\lstinline!include/unitree/idl/hg/HandCmd\_.hpp:27!}
\end{itemize}

\textbf{用法}

\begin{lstlisting}[language=Python]
current_value = obj.motor_cmd
obj.motor_cmd = new_value
\end{lstlisting}

\begin{quote}
{[}!TIP{]} 读取容器或嵌套值后应遵循``读取、修改、重新赋值''。只修改
getter 返回的副本不会可靠地更新原 C++ 消息。
\end{quote}

\hypertarget{unitree-sdk2-cpp-idl-hg-handcmd-reserve}{%
\paragraph{\texorpdfstring{\texttt{unitree\_sdk2\_cpp.idl.hg.HandCmd.reserve}}{unitree\_sdk2\_cpp.idl.hg.HandCmd.reserve}}\label{unitree-sdk2-cpp-idl-hg-handcmd-reserve}}

底层 \passthrough{\lstinline!unitree\_hg::msg::dds\_::HandCmd\_!} 的
\passthrough{\lstinline!reserve!} 字段。类型签名只说明 Python
表示；单位、数值范围、数组固定长度和枚举含义需查对应 IDL 头文件。
底层是固定长度数组，写入序列必须正好包含 4 个元素；元素约束：取值范围为
0 到 4294967295。

\textbf{签名}

\begin{lstlisting}[language=Python]
@property
def reserve(self) -> list[int]

@reserve.setter
def reserve(self, value: Sequence[int]) -> None
\end{lstlisting}

\textbf{可用性与安全性}

\textbf{\passthrough{\lstinline!AVAILABLE!} /
\passthrough{\lstinline!VALUE\_TYPE!}}：当前绑定源码已实现该消息值操作。

\textbf{参数}

\begin{longtable}[]{@{}
  >{\raggedright\arraybackslash}p{(\columnwidth - 4\tabcolsep) * \real{0.18}}
  >{\raggedright\arraybackslash}p{(\columnwidth - 4\tabcolsep) * \real{0.20}}
  >{\raggedright\arraybackslash}p{(\columnwidth - 4\tabcolsep) * \real{0.62}}@{}}
\toprule
名称 & Python 类型 & 含义 \\
\midrule
\endhead
\passthrough{\lstinline!value!} &
\passthrough{\lstinline!Sequence[int]!} &
要写入或传递的新值，必须符合该参数的 Python 类型以及底层 C++
范围约束。 \\
\bottomrule
\end{longtable}

\textbf{返回值}

\begin{longtable}[]{@{}
  >{\raggedright\arraybackslash}p{(\columnwidth - 4\tabcolsep) * \real{0.18}}
  >{\raggedright\arraybackslash}p{(\columnwidth - 4\tabcolsep) * \real{0.20}}
  >{\raggedright\arraybackslash}p{(\columnwidth - 4\tabcolsep) * \real{0.62}}@{}}
\toprule
位置 & 类型 & 含义 \\
\midrule
\endhead
getter 返回值 & \passthrough{\lstinline!list[int]!} &
返回属性当前值。数组、vector 和嵌套消息采用复制语义。 \\
\bottomrule
\end{longtable}

\textbf{对应 C++}

\begin{itemize}
\tightlist
\item
  类：\passthrough{\lstinline!unitree\_hg::msg::dds\_::HandCmd\_!}
\item
  字段类型：\passthrough{\lstinline!std::array<uint32\_t, 4>!}
\item
  声明位置：\passthrough{\lstinline!include/unitree/idl/hg/HandCmd\_.hpp:28!}
\end{itemize}

\textbf{用法}

\begin{lstlisting}[language=Python]
current_value = obj.reserve
obj.reserve = new_value
\end{lstlisting}

\begin{quote}
{[}!TIP{]} 读取容器或嵌套值后应遵循``读取、修改、重新赋值''。只修改
getter 返回的副本不会可靠地更新原 C++ 消息。
\end{quote}

\hypertarget{unitree-sdk2-cpp-idl-hg-imustate}{%
\subsubsection{\texorpdfstring{\texttt{unitree\_sdk2\_cpp.idl.hg.IMUState}}{unitree\_sdk2\_cpp.idl.hg.IMUState}}\label{unitree-sdk2-cpp-idl-hg-imustate}}

可由 Python 读写的 DDS/IDL 消息值类型。字段采用复制语义。

\textbf{导入}

\begin{lstlisting}[language=Python]
from unitree_sdk2_cpp.idl.hg import IMUState
\end{lstlisting}

\textbf{方法索引}

\begin{longtable}[]{@{}
  >{\raggedright\arraybackslash}p{(\columnwidth - 6\tabcolsep) * \real{0.18}}
  >{\raggedright\arraybackslash}p{(\columnwidth - 6\tabcolsep) * \real{0.24}}
  >{\raggedright\arraybackslash}p{(\columnwidth - 6\tabcolsep) * \real{0.18}}
  >{\raggedright\arraybackslash}p{(\columnwidth - 6\tabcolsep) * \real{0.40}}@{}}
\toprule
方法 & Python 签名 & 状态 & 安全分类 \\
\midrule
\endhead
\protect\hyperlink{unitree-sdk2-cpp-idl-hg-imustate-dunder-init-1}{\passthrough{\lstinline!\_\_init\_\_!}}
& \passthrough{\lstinline!def \_\_init\_\_(self) -> None!} &
\passthrough{\lstinline!AVAILABLE!} &
\passthrough{\lstinline!VALUE\_TYPE!} \\
\protect\hyperlink{unitree-sdk2-cpp-idl-hg-imustate-dunder-eq-1}{\passthrough{\lstinline!\_\_eq\_\_!}}
& \passthrough{\lstinline!def \_\_eq\_\_(self, other: object) -> bool!}
& \passthrough{\lstinline!AVAILABLE!} &
\passthrough{\lstinline!VALUE\_TYPE!} \\
\protect\hyperlink{unitree-sdk2-cpp-idl-hg-imustate-dunder-ne-1}{\passthrough{\lstinline!\_\_ne\_\_!}}
& \passthrough{\lstinline!def \_\_ne\_\_(self, other: object) -> bool!}
& \passthrough{\lstinline!AVAILABLE!} &
\passthrough{\lstinline!VALUE\_TYPE!} \\
\bottomrule
\end{longtable}

\textbf{属性索引}

\begin{longtable}[]{@{}
  >{\raggedright\arraybackslash}p{(\columnwidth - 6\tabcolsep) * \real{0.18}}
  >{\raggedright\arraybackslash}p{(\columnwidth - 6\tabcolsep) * \real{0.24}}
  >{\raggedright\arraybackslash}p{(\columnwidth - 6\tabcolsep) * \real{0.18}}
  >{\raggedright\arraybackslash}p{(\columnwidth - 6\tabcolsep) * \real{0.40}}@{}}
\toprule
属性 & 读取类型 & 写入类型 & 状态 \\
\midrule
\endhead
\protect\hyperlink{unitree-sdk2-cpp-idl-hg-imustate-quaternion}{\passthrough{\lstinline!quaternion!}}
& \passthrough{\lstinline!list[float]!} &
\passthrough{\lstinline!Sequence[float]!} &
\passthrough{\lstinline!AVAILABLE!} \\
\protect\hyperlink{unitree-sdk2-cpp-idl-hg-imustate-gyroscope}{\passthrough{\lstinline!gyroscope!}}
& \passthrough{\lstinline!list[float]!} &
\passthrough{\lstinline!Sequence[float]!} &
\passthrough{\lstinline!AVAILABLE!} \\
\protect\hyperlink{unitree-sdk2-cpp-idl-hg-imustate-accelerometer}{\passthrough{\lstinline!accelerometer!}}
& \passthrough{\lstinline!list[float]!} &
\passthrough{\lstinline!Sequence[float]!} &
\passthrough{\lstinline!AVAILABLE!} \\
\protect\hyperlink{unitree-sdk2-cpp-idl-hg-imustate-rpy}{\passthrough{\lstinline!rpy!}}
& \passthrough{\lstinline!list[float]!} &
\passthrough{\lstinline!Sequence[float]!} &
\passthrough{\lstinline!AVAILABLE!} \\
\protect\hyperlink{unitree-sdk2-cpp-idl-hg-imustate-temperature}{\passthrough{\lstinline!temperature!}}
& \passthrough{\lstinline!int!} & \passthrough{\lstinline!int!} &
\passthrough{\lstinline!AVAILABLE!} \\
\bottomrule
\end{longtable}

\hypertarget{unitree-sdk2-cpp-idl-hg-imustate-dunder-init-1}{%
\paragraph{\texorpdfstring{\texttt{unitree\_sdk2\_cpp.idl.hg.IMUState.\_\_init\_\_}}{unitree\_sdk2\_cpp.idl.hg.IMUState.\_\_init\_\_}}\label{unitree-sdk2-cpp-idl-hg-imustate-dunder-init-1}}

初始化 \passthrough{\lstinline!IMUState!}
实例。是否能实际构造取决于下方可用性状态。

\textbf{签名}

\begin{lstlisting}[language=Python]
def __init__(self) -> None
\end{lstlisting}

\textbf{可用性与安全性}

\textbf{\passthrough{\lstinline!AVAILABLE!} /
\passthrough{\lstinline!VALUE\_TYPE!}}：当前绑定源码已实现该消息值操作。

\textbf{参数}

无显式参数。实例方法中的 \passthrough{\lstinline!self!} 由 Python
自动传入。

\textbf{返回值}

\begin{longtable}[]{@{}
  >{\raggedright\arraybackslash}p{(\columnwidth - 4\tabcolsep) * \real{0.18}}
  >{\raggedright\arraybackslash}p{(\columnwidth - 4\tabcolsep) * \real{0.20}}
  >{\raggedright\arraybackslash}p{(\columnwidth - 4\tabcolsep) * \real{0.62}}@{}}
\toprule
位置 & 类型 & 含义 \\
\midrule
\endhead
无 & \passthrough{\lstinline!None!} &
\passthrough{\lstinline!\_\_init\_\_!}
本身不返回值；调用类对象时，在构造函数可用的前提下得到该类实例。 \\
\bottomrule
\end{longtable}

\textbf{对应 C++}

\begin{itemize}
\tightlist
\item
  类：\passthrough{\lstinline!unitree\_hg::msg::dds\_::IMUState\_!}
\item
  签名：\passthrough{\lstinline!IMUState\_()!}
\item
  绑定策略：\passthrough{\lstinline!未单独标注!}
\item
  声明位置：\passthrough{\lstinline!include/unitree/idl/hg/IMUState\_.hpp:31!}
\end{itemize}

\textbf{说明}

\begin{itemize}
\tightlist
\item
  示例是局部调用片段；\passthrough{\lstinline!obj!}
  和参数变量表示已经按上层业务校验并准备好的对象。
\end{itemize}

\textbf{用法}

\begin{lstlisting}[language=Python]
value = IMUState()
\end{lstlisting}

\hypertarget{unitree-sdk2-cpp-idl-hg-imustate-dunder-eq-1}{%
\paragraph{\texorpdfstring{\texttt{unitree\_sdk2\_cpp.idl.hg.IMUState.\_\_eq\_\_}}{unitree\_sdk2\_cpp.idl.hg.IMUState.\_\_eq\_\_}}\label{unitree-sdk2-cpp-idl-hg-imustate-dunder-eq-1}}

按底层 IDL 消息内容比较两个对象是否相等。

\textbf{签名}

\begin{lstlisting}[language=Python]
def __eq__(self, other: object) -> bool
\end{lstlisting}

\textbf{可用性与安全性}

\textbf{\passthrough{\lstinline!AVAILABLE!} /
\passthrough{\lstinline!VALUE\_TYPE!}}：当前绑定源码已实现该消息值操作。

\textbf{参数}

\begin{longtable}[]{@{}
  >{\raggedright\arraybackslash}p{(\columnwidth - 6\tabcolsep) * \real{0.18}}
  >{\raggedright\arraybackslash}p{(\columnwidth - 6\tabcolsep) * \real{0.24}}
  >{\raggedright\arraybackslash}p{(\columnwidth - 6\tabcolsep) * \real{0.18}}
  >{\raggedright\arraybackslash}p{(\columnwidth - 6\tabcolsep) * \real{0.40}}@{}}
\toprule
名称 & Python 类型 & 默认值 & 含义 \\
\midrule
\endhead
\passthrough{\lstinline!other!} & \passthrough{\lstinline!object!} &
必填 & 用于比较的另一个 Python 对象。类型不兼容时比较结果通常为
\passthrough{\lstinline!False!}。 \\
\bottomrule
\end{longtable}

\textbf{返回值}

\begin{longtable}[]{@{}
  >{\raggedright\arraybackslash}p{(\columnwidth - 4\tabcolsep) * \real{0.18}}
  >{\raggedright\arraybackslash}p{(\columnwidth - 4\tabcolsep) * \real{0.20}}
  >{\raggedright\arraybackslash}p{(\columnwidth - 4\tabcolsep) * \real{0.62}}@{}}
\toprule
位置 & 类型 & 含义 \\
\midrule
\endhead
返回值 & \passthrough{\lstinline!bool!} & 返回
\passthrough{\lstinline!bool!}。更精确的业务含义见该方法用途、C++
签名和目标型号协议。 \\
\bottomrule
\end{longtable}

\textbf{对应 C++}

\begin{itemize}
\tightlist
\item
  类：\passthrough{\lstinline!unitree\_hg::msg::dds\_::IMUState\_!}
\item
  签名：\passthrough{\lstinline!operator==(const value \&) const!}
\item
  绑定策略：\passthrough{\lstinline!未单独标注!}
\item
  声明位置：由 pybind11 手工绑定或生成注册表提供；manifest
  未记录独立头文件行号。
\end{itemize}

\textbf{说明}

\begin{itemize}
\tightlist
\item
  示例是局部调用片段；\passthrough{\lstinline!obj!}
  和参数变量表示已经按上层业务校验并准备好的对象。
\end{itemize}

\textbf{用法}

\begin{lstlisting}[language=Python]
same = left == right
\end{lstlisting}

\hypertarget{unitree-sdk2-cpp-idl-hg-imustate-dunder-ne-1}{%
\paragraph{\texorpdfstring{\texttt{unitree\_sdk2\_cpp.idl.hg.IMUState.\_\_ne\_\_}}{unitree\_sdk2\_cpp.idl.hg.IMUState.\_\_ne\_\_}}\label{unitree-sdk2-cpp-idl-hg-imustate-dunder-ne-1}}

按底层 IDL 消息内容比较两个对象是否不相等。

\textbf{签名}

\begin{lstlisting}[language=Python]
def __ne__(self, other: object) -> bool
\end{lstlisting}

\textbf{可用性与安全性}

\textbf{\passthrough{\lstinline!AVAILABLE!} /
\passthrough{\lstinline!VALUE\_TYPE!}}：当前绑定源码已实现该消息值操作。

\textbf{参数}

\begin{longtable}[]{@{}
  >{\raggedright\arraybackslash}p{(\columnwidth - 6\tabcolsep) * \real{0.18}}
  >{\raggedright\arraybackslash}p{(\columnwidth - 6\tabcolsep) * \real{0.24}}
  >{\raggedright\arraybackslash}p{(\columnwidth - 6\tabcolsep) * \real{0.18}}
  >{\raggedright\arraybackslash}p{(\columnwidth - 6\tabcolsep) * \real{0.40}}@{}}
\toprule
名称 & Python 类型 & 默认值 & 含义 \\
\midrule
\endhead
\passthrough{\lstinline!other!} & \passthrough{\lstinline!object!} &
必填 & 用于比较的另一个 Python 对象。类型不兼容时比较结果通常为
\passthrough{\lstinline!False!}。 \\
\bottomrule
\end{longtable}

\textbf{返回值}

\begin{longtable}[]{@{}
  >{\raggedright\arraybackslash}p{(\columnwidth - 4\tabcolsep) * \real{0.18}}
  >{\raggedright\arraybackslash}p{(\columnwidth - 4\tabcolsep) * \real{0.20}}
  >{\raggedright\arraybackslash}p{(\columnwidth - 4\tabcolsep) * \real{0.62}}@{}}
\toprule
位置 & 类型 & 含义 \\
\midrule
\endhead
返回值 & \passthrough{\lstinline!bool!} & 返回
\passthrough{\lstinline!bool!}。更精确的业务含义见该方法用途、C++
签名和目标型号协议。 \\
\bottomrule
\end{longtable}

\textbf{对应 C++}

\begin{itemize}
\tightlist
\item
  类：\passthrough{\lstinline!unitree\_hg::msg::dds\_::IMUState\_!}
\item
  签名：\passthrough{\lstinline"operator!=(const value \&) const"}
\item
  绑定策略：\passthrough{\lstinline!未单独标注!}
\item
  声明位置：由 pybind11 手工绑定或生成注册表提供；manifest
  未记录独立头文件行号。
\end{itemize}

\textbf{说明}

\begin{itemize}
\tightlist
\item
  示例是局部调用片段；\passthrough{\lstinline!obj!}
  和参数变量表示已经按上层业务校验并准备好的对象。
\end{itemize}

\textbf{用法}

\begin{lstlisting}[language=Python]
different = left != right
\end{lstlisting}

\hypertarget{unitree-sdk2-cpp-idl-hg-imustate-quaternion}{%
\paragraph{\texorpdfstring{\texttt{unitree\_sdk2\_cpp.idl.hg.IMUState.quaternion}}{unitree\_sdk2\_cpp.idl.hg.IMUState.quaternion}}\label{unitree-sdk2-cpp-idl-hg-imustate-quaternion}}

底层 \passthrough{\lstinline!unitree\_hg::msg::dds\_::IMUState\_!} 的
\passthrough{\lstinline!quaternion!} 字段。类型签名只说明 Python
表示；单位、数值范围、数组固定长度和枚举含义需查对应 IDL 头文件。
底层是固定长度数组，写入序列必须正好包含 4
个元素；元素约束：底层为浮点数；类型本身不说明单位、坐标系或业务安全范围。

\textbf{签名}

\begin{lstlisting}[language=Python]
@property
def quaternion(self) -> list[float]

@quaternion.setter
def quaternion(self, value: Sequence[float]) -> None
\end{lstlisting}

\textbf{可用性与安全性}

\textbf{\passthrough{\lstinline!AVAILABLE!} /
\passthrough{\lstinline!VALUE\_TYPE!}}：当前绑定源码已实现该消息值操作。

\textbf{参数}

\begin{longtable}[]{@{}
  >{\raggedright\arraybackslash}p{(\columnwidth - 4\tabcolsep) * \real{0.18}}
  >{\raggedright\arraybackslash}p{(\columnwidth - 4\tabcolsep) * \real{0.20}}
  >{\raggedright\arraybackslash}p{(\columnwidth - 4\tabcolsep) * \real{0.62}}@{}}
\toprule
名称 & Python 类型 & 含义 \\
\midrule
\endhead
\passthrough{\lstinline!value!} &
\passthrough{\lstinline!Sequence[float]!} &
要写入或传递的新值，必须符合该参数的 Python 类型以及底层 C++
范围约束。 \\
\bottomrule
\end{longtable}

\textbf{返回值}

\begin{longtable}[]{@{}
  >{\raggedright\arraybackslash}p{(\columnwidth - 4\tabcolsep) * \real{0.18}}
  >{\raggedright\arraybackslash}p{(\columnwidth - 4\tabcolsep) * \real{0.20}}
  >{\raggedright\arraybackslash}p{(\columnwidth - 4\tabcolsep) * \real{0.62}}@{}}
\toprule
位置 & 类型 & 含义 \\
\midrule
\endhead
getter 返回值 & \passthrough{\lstinline!list[float]!} &
返回属性当前值。数组、vector 和嵌套消息采用复制语义。 \\
\bottomrule
\end{longtable}

\textbf{对应 C++}

\begin{itemize}
\tightlist
\item
  类：\passthrough{\lstinline!unitree\_hg::msg::dds\_::IMUState\_!}
\item
  字段类型：\passthrough{\lstinline!std::array<float, 4>!}
\item
  声明位置：\passthrough{\lstinline!include/unitree/idl/hg/IMUState\_.hpp:24!}
\end{itemize}

\textbf{用法}

\begin{lstlisting}[language=Python]
current_value = obj.quaternion
obj.quaternion = new_value
\end{lstlisting}

\begin{quote}
{[}!TIP{]} 读取容器或嵌套值后应遵循``读取、修改、重新赋值''。只修改
getter 返回的副本不会可靠地更新原 C++ 消息。
\end{quote}

\hypertarget{unitree-sdk2-cpp-idl-hg-imustate-gyroscope}{%
\paragraph{\texorpdfstring{\texttt{unitree\_sdk2\_cpp.idl.hg.IMUState.gyroscope}}{unitree\_sdk2\_cpp.idl.hg.IMUState.gyroscope}}\label{unitree-sdk2-cpp-idl-hg-imustate-gyroscope}}

底层 \passthrough{\lstinline!unitree\_hg::msg::dds\_::IMUState\_!} 的
\passthrough{\lstinline!gyroscope!} 字段。类型签名只说明 Python
表示；单位、数值范围、数组固定长度和枚举含义需查对应 IDL 头文件。
底层是固定长度数组，写入序列必须正好包含 3
个元素；元素约束：底层为浮点数；类型本身不说明单位、坐标系或业务安全范围。

\textbf{签名}

\begin{lstlisting}[language=Python]
@property
def gyroscope(self) -> list[float]

@gyroscope.setter
def gyroscope(self, value: Sequence[float]) -> None
\end{lstlisting}

\textbf{可用性与安全性}

\textbf{\passthrough{\lstinline!AVAILABLE!} /
\passthrough{\lstinline!VALUE\_TYPE!}}：当前绑定源码已实现该消息值操作。

\textbf{参数}

\begin{longtable}[]{@{}
  >{\raggedright\arraybackslash}p{(\columnwidth - 4\tabcolsep) * \real{0.18}}
  >{\raggedright\arraybackslash}p{(\columnwidth - 4\tabcolsep) * \real{0.20}}
  >{\raggedright\arraybackslash}p{(\columnwidth - 4\tabcolsep) * \real{0.62}}@{}}
\toprule
名称 & Python 类型 & 含义 \\
\midrule
\endhead
\passthrough{\lstinline!value!} &
\passthrough{\lstinline!Sequence[float]!} &
要写入或传递的新值，必须符合该参数的 Python 类型以及底层 C++
范围约束。 \\
\bottomrule
\end{longtable}

\textbf{返回值}

\begin{longtable}[]{@{}
  >{\raggedright\arraybackslash}p{(\columnwidth - 4\tabcolsep) * \real{0.18}}
  >{\raggedright\arraybackslash}p{(\columnwidth - 4\tabcolsep) * \real{0.20}}
  >{\raggedright\arraybackslash}p{(\columnwidth - 4\tabcolsep) * \real{0.62}}@{}}
\toprule
位置 & 类型 & 含义 \\
\midrule
\endhead
getter 返回值 & \passthrough{\lstinline!list[float]!} &
返回属性当前值。数组、vector 和嵌套消息采用复制语义。 \\
\bottomrule
\end{longtable}

\textbf{对应 C++}

\begin{itemize}
\tightlist
\item
  类：\passthrough{\lstinline!unitree\_hg::msg::dds\_::IMUState\_!}
\item
  字段类型：\passthrough{\lstinline!std::array<float, 3>!}
\item
  声明位置：\passthrough{\lstinline!include/unitree/idl/hg/IMUState\_.hpp:25!}
\end{itemize}

\textbf{用法}

\begin{lstlisting}[language=Python]
current_value = obj.gyroscope
obj.gyroscope = new_value
\end{lstlisting}

\begin{quote}
{[}!TIP{]} 读取容器或嵌套值后应遵循``读取、修改、重新赋值''。只修改
getter 返回的副本不会可靠地更新原 C++ 消息。
\end{quote}

\hypertarget{unitree-sdk2-cpp-idl-hg-imustate-accelerometer}{%
\paragraph{\texorpdfstring{\texttt{unitree\_sdk2\_cpp.idl.hg.IMUState.accelerometer}}{unitree\_sdk2\_cpp.idl.hg.IMUState.accelerometer}}\label{unitree-sdk2-cpp-idl-hg-imustate-accelerometer}}

底层 \passthrough{\lstinline!unitree\_hg::msg::dds\_::IMUState\_!} 的
\passthrough{\lstinline!accelerometer!} 字段。类型签名只说明 Python
表示；单位、数值范围、数组固定长度和枚举含义需查对应 IDL 头文件。
底层是固定长度数组，写入序列必须正好包含 3
个元素；元素约束：底层为浮点数；类型本身不说明单位、坐标系或业务安全范围。

\textbf{签名}

\begin{lstlisting}[language=Python]
@property
def accelerometer(self) -> list[float]

@accelerometer.setter
def accelerometer(self, value: Sequence[float]) -> None
\end{lstlisting}

\textbf{可用性与安全性}

\textbf{\passthrough{\lstinline!AVAILABLE!} /
\passthrough{\lstinline!VALUE\_TYPE!}}：当前绑定源码已实现该消息值操作。

\textbf{参数}

\begin{longtable}[]{@{}
  >{\raggedright\arraybackslash}p{(\columnwidth - 4\tabcolsep) * \real{0.18}}
  >{\raggedright\arraybackslash}p{(\columnwidth - 4\tabcolsep) * \real{0.20}}
  >{\raggedright\arraybackslash}p{(\columnwidth - 4\tabcolsep) * \real{0.62}}@{}}
\toprule
名称 & Python 类型 & 含义 \\
\midrule
\endhead
\passthrough{\lstinline!value!} &
\passthrough{\lstinline!Sequence[float]!} &
要写入或传递的新值，必须符合该参数的 Python 类型以及底层 C++
范围约束。 \\
\bottomrule
\end{longtable}

\textbf{返回值}

\begin{longtable}[]{@{}
  >{\raggedright\arraybackslash}p{(\columnwidth - 4\tabcolsep) * \real{0.18}}
  >{\raggedright\arraybackslash}p{(\columnwidth - 4\tabcolsep) * \real{0.20}}
  >{\raggedright\arraybackslash}p{(\columnwidth - 4\tabcolsep) * \real{0.62}}@{}}
\toprule
位置 & 类型 & 含义 \\
\midrule
\endhead
getter 返回值 & \passthrough{\lstinline!list[float]!} &
返回属性当前值。数组、vector 和嵌套消息采用复制语义。 \\
\bottomrule
\end{longtable}

\textbf{对应 C++}

\begin{itemize}
\tightlist
\item
  类：\passthrough{\lstinline!unitree\_hg::msg::dds\_::IMUState\_!}
\item
  字段类型：\passthrough{\lstinline!std::array<float, 3>!}
\item
  声明位置：\passthrough{\lstinline!include/unitree/idl/hg/IMUState\_.hpp:26!}
\end{itemize}

\textbf{用法}

\begin{lstlisting}[language=Python]
current_value = obj.accelerometer
obj.accelerometer = new_value
\end{lstlisting}

\begin{quote}
{[}!TIP{]} 读取容器或嵌套值后应遵循``读取、修改、重新赋值''。只修改
getter 返回的副本不会可靠地更新原 C++ 消息。
\end{quote}

\hypertarget{unitree-sdk2-cpp-idl-hg-imustate-rpy}{%
\paragraph{\texorpdfstring{\texttt{unitree\_sdk2\_cpp.idl.hg.IMUState.rpy}}{unitree\_sdk2\_cpp.idl.hg.IMUState.rpy}}\label{unitree-sdk2-cpp-idl-hg-imustate-rpy}}

底层 \passthrough{\lstinline!unitree\_hg::msg::dds\_::IMUState\_!} 的
\passthrough{\lstinline!rpy!} 字段。类型签名只说明 Python
表示；单位、数值范围、数组固定长度和枚举含义需查对应 IDL 头文件。
底层是固定长度数组，写入序列必须正好包含 3
个元素；元素约束：底层为浮点数；类型本身不说明单位、坐标系或业务安全范围。

\textbf{签名}

\begin{lstlisting}[language=Python]
@property
def rpy(self) -> list[float]

@rpy.setter
def rpy(self, value: Sequence[float]) -> None
\end{lstlisting}

\textbf{可用性与安全性}

\textbf{\passthrough{\lstinline!AVAILABLE!} /
\passthrough{\lstinline!VALUE\_TYPE!}}：当前绑定源码已实现该消息值操作。

\textbf{参数}

\begin{longtable}[]{@{}
  >{\raggedright\arraybackslash}p{(\columnwidth - 4\tabcolsep) * \real{0.18}}
  >{\raggedright\arraybackslash}p{(\columnwidth - 4\tabcolsep) * \real{0.20}}
  >{\raggedright\arraybackslash}p{(\columnwidth - 4\tabcolsep) * \real{0.62}}@{}}
\toprule
名称 & Python 类型 & 含义 \\
\midrule
\endhead
\passthrough{\lstinline!value!} &
\passthrough{\lstinline!Sequence[float]!} &
要写入或传递的新值，必须符合该参数的 Python 类型以及底层 C++
范围约束。 \\
\bottomrule
\end{longtable}

\textbf{返回值}

\begin{longtable}[]{@{}
  >{\raggedright\arraybackslash}p{(\columnwidth - 4\tabcolsep) * \real{0.18}}
  >{\raggedright\arraybackslash}p{(\columnwidth - 4\tabcolsep) * \real{0.20}}
  >{\raggedright\arraybackslash}p{(\columnwidth - 4\tabcolsep) * \real{0.62}}@{}}
\toprule
位置 & 类型 & 含义 \\
\midrule
\endhead
getter 返回值 & \passthrough{\lstinline!list[float]!} &
返回属性当前值。数组、vector 和嵌套消息采用复制语义。 \\
\bottomrule
\end{longtable}

\textbf{对应 C++}

\begin{itemize}
\tightlist
\item
  类：\passthrough{\lstinline!unitree\_hg::msg::dds\_::IMUState\_!}
\item
  字段类型：\passthrough{\lstinline!std::array<float, 3>!}
\item
  声明位置：\passthrough{\lstinline!include/unitree/idl/hg/IMUState\_.hpp:27!}
\end{itemize}

\textbf{用法}

\begin{lstlisting}[language=Python]
current_value = obj.rpy
obj.rpy = new_value
\end{lstlisting}

\begin{quote}
{[}!TIP{]} 读取容器或嵌套值后应遵循``读取、修改、重新赋值''。只修改
getter 返回的副本不会可靠地更新原 C++ 消息。
\end{quote}

\hypertarget{unitree-sdk2-cpp-idl-hg-imustate-temperature}{%
\paragraph{\texorpdfstring{\texttt{unitree\_sdk2\_cpp.idl.hg.IMUState.temperature}}{unitree\_sdk2\_cpp.idl.hg.IMUState.temperature}}\label{unitree-sdk2-cpp-idl-hg-imustate-temperature}}

底层 \passthrough{\lstinline!unitree\_hg::msg::dds\_::IMUState\_!} 的
\passthrough{\lstinline!temperature!} 字段。类型签名只说明 Python
表示；单位、数值范围、数组固定长度和枚举含义需查对应 IDL 头文件。
取值范围为 -32768 到 32767。

\textbf{签名}

\begin{lstlisting}[language=Python]
@property
def temperature(self) -> int

@temperature.setter
def temperature(self, value: int) -> None
\end{lstlisting}

\textbf{可用性与安全性}

\textbf{\passthrough{\lstinline!AVAILABLE!} /
\passthrough{\lstinline!VALUE\_TYPE!}}：当前绑定源码已实现该消息值操作。

\textbf{参数}

\begin{longtable}[]{@{}
  >{\raggedright\arraybackslash}p{(\columnwidth - 4\tabcolsep) * \real{0.18}}
  >{\raggedright\arraybackslash}p{(\columnwidth - 4\tabcolsep) * \real{0.20}}
  >{\raggedright\arraybackslash}p{(\columnwidth - 4\tabcolsep) * \real{0.62}}@{}}
\toprule
名称 & Python 类型 & 含义 \\
\midrule
\endhead
\passthrough{\lstinline!value!} & \passthrough{\lstinline!int!} &
要写入或传递的新值，必须符合该参数的 Python 类型以及底层 C++
范围约束。 \\
\bottomrule
\end{longtable}

\textbf{返回值}

\begin{longtable}[]{@{}
  >{\raggedright\arraybackslash}p{(\columnwidth - 4\tabcolsep) * \real{0.18}}
  >{\raggedright\arraybackslash}p{(\columnwidth - 4\tabcolsep) * \real{0.20}}
  >{\raggedright\arraybackslash}p{(\columnwidth - 4\tabcolsep) * \real{0.62}}@{}}
\toprule
位置 & 类型 & 含义 \\
\midrule
\endhead
getter 返回值 & \passthrough{\lstinline!int!} & 返回属性当前值。 \\
\bottomrule
\end{longtable}

\textbf{对应 C++}

\begin{itemize}
\tightlist
\item
  类：\passthrough{\lstinline!unitree\_hg::msg::dds\_::IMUState\_!}
\item
  字段类型：\passthrough{\lstinline!int16\_t!}
\item
  声明位置：\passthrough{\lstinline!include/unitree/idl/hg/IMUState\_.hpp:28!}
\end{itemize}

\textbf{用法}

\begin{lstlisting}[language=Python]
current_value = obj.temperature
obj.temperature = new_value
\end{lstlisting}

\hypertarget{unitree-sdk2-cpp-idl-hg-motorstate}{%
\subsubsection{\texorpdfstring{\texttt{unitree\_sdk2\_cpp.idl.hg.MotorState}}{unitree\_sdk2\_cpp.idl.hg.MotorState}}\label{unitree-sdk2-cpp-idl-hg-motorstate}}

可由 Python 读写的 DDS/IDL 消息值类型。字段采用复制语义。

\textbf{导入}

\begin{lstlisting}[language=Python]
from unitree_sdk2_cpp.idl.hg import MotorState
\end{lstlisting}

\textbf{方法索引}

\begin{longtable}[]{@{}
  >{\raggedright\arraybackslash}p{(\columnwidth - 6\tabcolsep) * \real{0.18}}
  >{\raggedright\arraybackslash}p{(\columnwidth - 6\tabcolsep) * \real{0.24}}
  >{\raggedright\arraybackslash}p{(\columnwidth - 6\tabcolsep) * \real{0.18}}
  >{\raggedright\arraybackslash}p{(\columnwidth - 6\tabcolsep) * \real{0.40}}@{}}
\toprule
方法 & Python 签名 & 状态 & 安全分类 \\
\midrule
\endhead
\protect\hyperlink{unitree-sdk2-cpp-idl-hg-motorstate-dunder-init-1}{\passthrough{\lstinline!\_\_init\_\_!}}
& \passthrough{\lstinline!def \_\_init\_\_(self) -> None!} &
\passthrough{\lstinline!AVAILABLE!} &
\passthrough{\lstinline!VALUE\_TYPE!} \\
\protect\hyperlink{unitree-sdk2-cpp-idl-hg-motorstate-dunder-eq-1}{\passthrough{\lstinline!\_\_eq\_\_!}}
& \passthrough{\lstinline!def \_\_eq\_\_(self, other: object) -> bool!}
& \passthrough{\lstinline!AVAILABLE!} &
\passthrough{\lstinline!VALUE\_TYPE!} \\
\protect\hyperlink{unitree-sdk2-cpp-idl-hg-motorstate-dunder-ne-1}{\passthrough{\lstinline!\_\_ne\_\_!}}
& \passthrough{\lstinline!def \_\_ne\_\_(self, other: object) -> bool!}
& \passthrough{\lstinline!AVAILABLE!} &
\passthrough{\lstinline!VALUE\_TYPE!} \\
\bottomrule
\end{longtable}

\textbf{属性索引}

\begin{longtable}[]{@{}
  >{\raggedright\arraybackslash}p{(\columnwidth - 6\tabcolsep) * \real{0.18}}
  >{\raggedright\arraybackslash}p{(\columnwidth - 6\tabcolsep) * \real{0.24}}
  >{\raggedright\arraybackslash}p{(\columnwidth - 6\tabcolsep) * \real{0.18}}
  >{\raggedright\arraybackslash}p{(\columnwidth - 6\tabcolsep) * \real{0.40}}@{}}
\toprule
属性 & 读取类型 & 写入类型 & 状态 \\
\midrule
\endhead
\protect\hyperlink{unitree-sdk2-cpp-idl-hg-motorstate-mode}{\passthrough{\lstinline!mode!}}
& \passthrough{\lstinline!int!} & \passthrough{\lstinline!int!} &
\passthrough{\lstinline!AVAILABLE!} \\
\protect\hyperlink{unitree-sdk2-cpp-idl-hg-motorstate-q}{\passthrough{\lstinline!q!}}
& \passthrough{\lstinline!float!} & \passthrough{\lstinline!float!} &
\passthrough{\lstinline!AVAILABLE!} \\
\protect\hyperlink{unitree-sdk2-cpp-idl-hg-motorstate-dq}{\passthrough{\lstinline!dq!}}
& \passthrough{\lstinline!float!} & \passthrough{\lstinline!float!} &
\passthrough{\lstinline!AVAILABLE!} \\
\protect\hyperlink{unitree-sdk2-cpp-idl-hg-motorstate-ddq}{\passthrough{\lstinline!ddq!}}
& \passthrough{\lstinline!float!} & \passthrough{\lstinline!float!} &
\passthrough{\lstinline!AVAILABLE!} \\
\protect\hyperlink{unitree-sdk2-cpp-idl-hg-motorstate-tau-est}{\passthrough{\lstinline!tau\_est!}}
& \passthrough{\lstinline!float!} & \passthrough{\lstinline!float!} &
\passthrough{\lstinline!AVAILABLE!} \\
\protect\hyperlink{unitree-sdk2-cpp-idl-hg-motorstate-temperature}{\passthrough{\lstinline!temperature!}}
& \passthrough{\lstinline!list[int]!} &
\passthrough{\lstinline!Sequence[int]!} &
\passthrough{\lstinline!AVAILABLE!} \\
\protect\hyperlink{unitree-sdk2-cpp-idl-hg-motorstate-vol}{\passthrough{\lstinline!vol!}}
& \passthrough{\lstinline!float!} & \passthrough{\lstinline!float!} &
\passthrough{\lstinline!AVAILABLE!} \\
\protect\hyperlink{unitree-sdk2-cpp-idl-hg-motorstate-sensor}{\passthrough{\lstinline!sensor!}}
& \passthrough{\lstinline!list[int]!} &
\passthrough{\lstinline!Sequence[int]!} &
\passthrough{\lstinline!AVAILABLE!} \\
\protect\hyperlink{unitree-sdk2-cpp-idl-hg-motorstate-motorstate}{\passthrough{\lstinline!motorstate!}}
& \passthrough{\lstinline!int!} & \passthrough{\lstinline!int!} &
\passthrough{\lstinline!AVAILABLE!} \\
\protect\hyperlink{unitree-sdk2-cpp-idl-hg-motorstate-reserve}{\passthrough{\lstinline!reserve!}}
& \passthrough{\lstinline!list[int]!} &
\passthrough{\lstinline!Sequence[int]!} &
\passthrough{\lstinline!AVAILABLE!} \\
\bottomrule
\end{longtable}

\hypertarget{unitree-sdk2-cpp-idl-hg-motorstate-dunder-init-1}{%
\paragraph{\texorpdfstring{\texttt{unitree\_sdk2\_cpp.idl.hg.MotorState.\_\_init\_\_}}{unitree\_sdk2\_cpp.idl.hg.MotorState.\_\_init\_\_}}\label{unitree-sdk2-cpp-idl-hg-motorstate-dunder-init-1}}

初始化 \passthrough{\lstinline!MotorState!}
实例。是否能实际构造取决于下方可用性状态。

\textbf{签名}

\begin{lstlisting}[language=Python]
def __init__(self) -> None
\end{lstlisting}

\textbf{可用性与安全性}

\textbf{\passthrough{\lstinline!AVAILABLE!} /
\passthrough{\lstinline!VALUE\_TYPE!}}：当前绑定源码已实现该消息值操作。

\textbf{参数}

无显式参数。实例方法中的 \passthrough{\lstinline!self!} 由 Python
自动传入。

\textbf{返回值}

\begin{longtable}[]{@{}
  >{\raggedright\arraybackslash}p{(\columnwidth - 4\tabcolsep) * \real{0.18}}
  >{\raggedright\arraybackslash}p{(\columnwidth - 4\tabcolsep) * \real{0.20}}
  >{\raggedright\arraybackslash}p{(\columnwidth - 4\tabcolsep) * \real{0.62}}@{}}
\toprule
位置 & 类型 & 含义 \\
\midrule
\endhead
无 & \passthrough{\lstinline!None!} &
\passthrough{\lstinline!\_\_init\_\_!}
本身不返回值；调用类对象时，在构造函数可用的前提下得到该类实例。 \\
\bottomrule
\end{longtable}

\textbf{对应 C++}

\begin{itemize}
\tightlist
\item
  类：\passthrough{\lstinline!unitree\_hg::msg::dds\_::MotorState\_!}
\item
  签名：\passthrough{\lstinline!MotorState\_()!}
\item
  绑定策略：\passthrough{\lstinline!未单独标注!}
\item
  声明位置：\passthrough{\lstinline!include/unitree/idl/hg/MotorState\_.hpp:36!}
\end{itemize}

\textbf{说明}

\begin{itemize}
\tightlist
\item
  示例是局部调用片段；\passthrough{\lstinline!obj!}
  和参数变量表示已经按上层业务校验并准备好的对象。
\end{itemize}

\textbf{用法}

\begin{lstlisting}[language=Python]
value = MotorState()
\end{lstlisting}

\hypertarget{unitree-sdk2-cpp-idl-hg-motorstate-dunder-eq-1}{%
\paragraph{\texorpdfstring{\texttt{unitree\_sdk2\_cpp.idl.hg.MotorState.\_\_eq\_\_}}{unitree\_sdk2\_cpp.idl.hg.MotorState.\_\_eq\_\_}}\label{unitree-sdk2-cpp-idl-hg-motorstate-dunder-eq-1}}

按底层 IDL 消息内容比较两个对象是否相等。

\textbf{签名}

\begin{lstlisting}[language=Python]
def __eq__(self, other: object) -> bool
\end{lstlisting}

\textbf{可用性与安全性}

\textbf{\passthrough{\lstinline!AVAILABLE!} /
\passthrough{\lstinline!VALUE\_TYPE!}}：当前绑定源码已实现该消息值操作。

\textbf{参数}

\begin{longtable}[]{@{}
  >{\raggedright\arraybackslash}p{(\columnwidth - 6\tabcolsep) * \real{0.18}}
  >{\raggedright\arraybackslash}p{(\columnwidth - 6\tabcolsep) * \real{0.24}}
  >{\raggedright\arraybackslash}p{(\columnwidth - 6\tabcolsep) * \real{0.18}}
  >{\raggedright\arraybackslash}p{(\columnwidth - 6\tabcolsep) * \real{0.40}}@{}}
\toprule
名称 & Python 类型 & 默认值 & 含义 \\
\midrule
\endhead
\passthrough{\lstinline!other!} & \passthrough{\lstinline!object!} &
必填 & 用于比较的另一个 Python 对象。类型不兼容时比较结果通常为
\passthrough{\lstinline!False!}。 \\
\bottomrule
\end{longtable}

\textbf{返回值}

\begin{longtable}[]{@{}
  >{\raggedright\arraybackslash}p{(\columnwidth - 4\tabcolsep) * \real{0.18}}
  >{\raggedright\arraybackslash}p{(\columnwidth - 4\tabcolsep) * \real{0.20}}
  >{\raggedright\arraybackslash}p{(\columnwidth - 4\tabcolsep) * \real{0.62}}@{}}
\toprule
位置 & 类型 & 含义 \\
\midrule
\endhead
返回值 & \passthrough{\lstinline!bool!} & 返回
\passthrough{\lstinline!bool!}。更精确的业务含义见该方法用途、C++
签名和目标型号协议。 \\
\bottomrule
\end{longtable}

\textbf{对应 C++}

\begin{itemize}
\tightlist
\item
  类：\passthrough{\lstinline!unitree\_hg::msg::dds\_::MotorState\_!}
\item
  签名：\passthrough{\lstinline!operator==(const value \&) const!}
\item
  绑定策略：\passthrough{\lstinline!未单独标注!}
\item
  声明位置：由 pybind11 手工绑定或生成注册表提供；manifest
  未记录独立头文件行号。
\end{itemize}

\textbf{说明}

\begin{itemize}
\tightlist
\item
  示例是局部调用片段；\passthrough{\lstinline!obj!}
  和参数变量表示已经按上层业务校验并准备好的对象。
\end{itemize}

\textbf{用法}

\begin{lstlisting}[language=Python]
same = left == right
\end{lstlisting}

\hypertarget{unitree-sdk2-cpp-idl-hg-motorstate-dunder-ne-1}{%
\paragraph{\texorpdfstring{\texttt{unitree\_sdk2\_cpp.idl.hg.MotorState.\_\_ne\_\_}}{unitree\_sdk2\_cpp.idl.hg.MotorState.\_\_ne\_\_}}\label{unitree-sdk2-cpp-idl-hg-motorstate-dunder-ne-1}}

按底层 IDL 消息内容比较两个对象是否不相等。

\textbf{签名}

\begin{lstlisting}[language=Python]
def __ne__(self, other: object) -> bool
\end{lstlisting}

\textbf{可用性与安全性}

\textbf{\passthrough{\lstinline!AVAILABLE!} /
\passthrough{\lstinline!VALUE\_TYPE!}}：当前绑定源码已实现该消息值操作。

\textbf{参数}

\begin{longtable}[]{@{}
  >{\raggedright\arraybackslash}p{(\columnwidth - 6\tabcolsep) * \real{0.18}}
  >{\raggedright\arraybackslash}p{(\columnwidth - 6\tabcolsep) * \real{0.24}}
  >{\raggedright\arraybackslash}p{(\columnwidth - 6\tabcolsep) * \real{0.18}}
  >{\raggedright\arraybackslash}p{(\columnwidth - 6\tabcolsep) * \real{0.40}}@{}}
\toprule
名称 & Python 类型 & 默认值 & 含义 \\
\midrule
\endhead
\passthrough{\lstinline!other!} & \passthrough{\lstinline!object!} &
必填 & 用于比较的另一个 Python 对象。类型不兼容时比较结果通常为
\passthrough{\lstinline!False!}。 \\
\bottomrule
\end{longtable}

\textbf{返回值}

\begin{longtable}[]{@{}
  >{\raggedright\arraybackslash}p{(\columnwidth - 4\tabcolsep) * \real{0.18}}
  >{\raggedright\arraybackslash}p{(\columnwidth - 4\tabcolsep) * \real{0.20}}
  >{\raggedright\arraybackslash}p{(\columnwidth - 4\tabcolsep) * \real{0.62}}@{}}
\toprule
位置 & 类型 & 含义 \\
\midrule
\endhead
返回值 & \passthrough{\lstinline!bool!} & 返回
\passthrough{\lstinline!bool!}。更精确的业务含义见该方法用途、C++
签名和目标型号协议。 \\
\bottomrule
\end{longtable}

\textbf{对应 C++}

\begin{itemize}
\tightlist
\item
  类：\passthrough{\lstinline!unitree\_hg::msg::dds\_::MotorState\_!}
\item
  签名：\passthrough{\lstinline"operator!=(const value \&) const"}
\item
  绑定策略：\passthrough{\lstinline!未单独标注!}
\item
  声明位置：由 pybind11 手工绑定或生成注册表提供；manifest
  未记录独立头文件行号。
\end{itemize}

\textbf{说明}

\begin{itemize}
\tightlist
\item
  示例是局部调用片段；\passthrough{\lstinline!obj!}
  和参数变量表示已经按上层业务校验并准备好的对象。
\end{itemize}

\textbf{用法}

\begin{lstlisting}[language=Python]
different = left != right
\end{lstlisting}

\hypertarget{unitree-sdk2-cpp-idl-hg-motorstate-mode}{%
\paragraph{\texorpdfstring{\texttt{unitree\_sdk2\_cpp.idl.hg.MotorState.mode}}{unitree\_sdk2\_cpp.idl.hg.MotorState.mode}}\label{unitree-sdk2-cpp-idl-hg-motorstate-mode}}

底层 \passthrough{\lstinline!unitree\_hg::msg::dds\_::MotorState\_!} 的
\passthrough{\lstinline!mode!} 字段。类型签名只说明 Python
表示；单位、数值范围、数组固定长度和枚举含义需查对应 IDL 头文件。
取值范围为 0 到 255。

\textbf{签名}

\begin{lstlisting}[language=Python]
@property
def mode(self) -> int

@mode.setter
def mode(self, value: int) -> None
\end{lstlisting}

\textbf{可用性与安全性}

\textbf{\passthrough{\lstinline!AVAILABLE!} /
\passthrough{\lstinline!VALUE\_TYPE!}}：当前绑定源码已实现该消息值操作。

\textbf{参数}

\begin{longtable}[]{@{}
  >{\raggedright\arraybackslash}p{(\columnwidth - 4\tabcolsep) * \real{0.18}}
  >{\raggedright\arraybackslash}p{(\columnwidth - 4\tabcolsep) * \real{0.20}}
  >{\raggedright\arraybackslash}p{(\columnwidth - 4\tabcolsep) * \real{0.62}}@{}}
\toprule
名称 & Python 类型 & 含义 \\
\midrule
\endhead
\passthrough{\lstinline!value!} & \passthrough{\lstinline!int!} &
要写入或传递的新值，必须符合该参数的 Python 类型以及底层 C++
范围约束。 \\
\bottomrule
\end{longtable}

\textbf{返回值}

\begin{longtable}[]{@{}
  >{\raggedright\arraybackslash}p{(\columnwidth - 4\tabcolsep) * \real{0.18}}
  >{\raggedright\arraybackslash}p{(\columnwidth - 4\tabcolsep) * \real{0.20}}
  >{\raggedright\arraybackslash}p{(\columnwidth - 4\tabcolsep) * \real{0.62}}@{}}
\toprule
位置 & 类型 & 含义 \\
\midrule
\endhead
getter 返回值 & \passthrough{\lstinline!int!} & 返回属性当前值。 \\
\bottomrule
\end{longtable}

\textbf{对应 C++}

\begin{itemize}
\tightlist
\item
  类：\passthrough{\lstinline!unitree\_hg::msg::dds\_::MotorState\_!}
\item
  字段类型：\passthrough{\lstinline!uint8\_t!}
\item
  声明位置：\passthrough{\lstinline!include/unitree/idl/hg/MotorState\_.hpp:24!}
\end{itemize}

\textbf{用法}

\begin{lstlisting}[language=Python]
current_value = obj.mode
obj.mode = new_value
\end{lstlisting}

\hypertarget{unitree-sdk2-cpp-idl-hg-motorstate-q}{%
\paragraph{\texorpdfstring{\texttt{unitree\_sdk2\_cpp.idl.hg.MotorState.q}}{unitree\_sdk2\_cpp.idl.hg.MotorState.q}}\label{unitree-sdk2-cpp-idl-hg-motorstate-q}}

底层 \passthrough{\lstinline!unitree\_hg::msg::dds\_::MotorState\_!} 的
\passthrough{\lstinline!q!} 字段。类型签名只说明 Python
表示；单位、数值范围、数组固定长度和枚举含义需查对应 IDL 头文件。
底层为浮点数；类型本身不说明单位、坐标系或业务安全范围。

\textbf{签名}

\begin{lstlisting}[language=Python]
@property
def q(self) -> float

@q.setter
def q(self, value: float) -> None
\end{lstlisting}

\textbf{可用性与安全性}

\textbf{\passthrough{\lstinline!AVAILABLE!} /
\passthrough{\lstinline!VALUE\_TYPE!}}：当前绑定源码已实现该消息值操作。

\textbf{参数}

\begin{longtable}[]{@{}
  >{\raggedright\arraybackslash}p{(\columnwidth - 4\tabcolsep) * \real{0.18}}
  >{\raggedright\arraybackslash}p{(\columnwidth - 4\tabcolsep) * \real{0.20}}
  >{\raggedright\arraybackslash}p{(\columnwidth - 4\tabcolsep) * \real{0.62}}@{}}
\toprule
名称 & Python 类型 & 含义 \\
\midrule
\endhead
\passthrough{\lstinline!value!} & \passthrough{\lstinline!float!} &
要写入或传递的新值，必须符合该参数的 Python 类型以及底层 C++
范围约束。 \\
\bottomrule
\end{longtable}

\textbf{返回值}

\begin{longtable}[]{@{}
  >{\raggedright\arraybackslash}p{(\columnwidth - 4\tabcolsep) * \real{0.18}}
  >{\raggedright\arraybackslash}p{(\columnwidth - 4\tabcolsep) * \real{0.20}}
  >{\raggedright\arraybackslash}p{(\columnwidth - 4\tabcolsep) * \real{0.62}}@{}}
\toprule
位置 & 类型 & 含义 \\
\midrule
\endhead
getter 返回值 & \passthrough{\lstinline!float!} & 返回属性当前值。 \\
\bottomrule
\end{longtable}

\textbf{对应 C++}

\begin{itemize}
\tightlist
\item
  类：\passthrough{\lstinline!unitree\_hg::msg::dds\_::MotorState\_!}
\item
  字段类型：\passthrough{\lstinline!float!}
\item
  声明位置：\passthrough{\lstinline!include/unitree/idl/hg/MotorState\_.hpp:25!}
\end{itemize}

\textbf{用法}

\begin{lstlisting}[language=Python]
current_value = obj.q
obj.q = new_value
\end{lstlisting}

\hypertarget{unitree-sdk2-cpp-idl-hg-motorstate-dq}{%
\paragraph{\texorpdfstring{\texttt{unitree\_sdk2\_cpp.idl.hg.MotorState.dq}}{unitree\_sdk2\_cpp.idl.hg.MotorState.dq}}\label{unitree-sdk2-cpp-idl-hg-motorstate-dq}}

底层 \passthrough{\lstinline!unitree\_hg::msg::dds\_::MotorState\_!} 的
\passthrough{\lstinline!dq!} 字段。类型签名只说明 Python
表示；单位、数值范围、数组固定长度和枚举含义需查对应 IDL 头文件。
底层为浮点数；类型本身不说明单位、坐标系或业务安全范围。

\textbf{签名}

\begin{lstlisting}[language=Python]
@property
def dq(self) -> float

@dq.setter
def dq(self, value: float) -> None
\end{lstlisting}

\textbf{可用性与安全性}

\textbf{\passthrough{\lstinline!AVAILABLE!} /
\passthrough{\lstinline!VALUE\_TYPE!}}：当前绑定源码已实现该消息值操作。

\textbf{参数}

\begin{longtable}[]{@{}
  >{\raggedright\arraybackslash}p{(\columnwidth - 4\tabcolsep) * \real{0.18}}
  >{\raggedright\arraybackslash}p{(\columnwidth - 4\tabcolsep) * \real{0.20}}
  >{\raggedright\arraybackslash}p{(\columnwidth - 4\tabcolsep) * \real{0.62}}@{}}
\toprule
名称 & Python 类型 & 含义 \\
\midrule
\endhead
\passthrough{\lstinline!value!} & \passthrough{\lstinline!float!} &
要写入或传递的新值，必须符合该参数的 Python 类型以及底层 C++
范围约束。 \\
\bottomrule
\end{longtable}

\textbf{返回值}

\begin{longtable}[]{@{}
  >{\raggedright\arraybackslash}p{(\columnwidth - 4\tabcolsep) * \real{0.18}}
  >{\raggedright\arraybackslash}p{(\columnwidth - 4\tabcolsep) * \real{0.20}}
  >{\raggedright\arraybackslash}p{(\columnwidth - 4\tabcolsep) * \real{0.62}}@{}}
\toprule
位置 & 类型 & 含义 \\
\midrule
\endhead
getter 返回值 & \passthrough{\lstinline!float!} & 返回属性当前值。 \\
\bottomrule
\end{longtable}

\textbf{对应 C++}

\begin{itemize}
\tightlist
\item
  类：\passthrough{\lstinline!unitree\_hg::msg::dds\_::MotorState\_!}
\item
  字段类型：\passthrough{\lstinline!float!}
\item
  声明位置：\passthrough{\lstinline!include/unitree/idl/hg/MotorState\_.hpp:26!}
\end{itemize}

\textbf{用法}

\begin{lstlisting}[language=Python]
current_value = obj.dq
obj.dq = new_value
\end{lstlisting}

\hypertarget{unitree-sdk2-cpp-idl-hg-motorstate-ddq}{%
\paragraph{\texorpdfstring{\texttt{unitree\_sdk2\_cpp.idl.hg.MotorState.ddq}}{unitree\_sdk2\_cpp.idl.hg.MotorState.ddq}}\label{unitree-sdk2-cpp-idl-hg-motorstate-ddq}}

底层 \passthrough{\lstinline!unitree\_hg::msg::dds\_::MotorState\_!} 的
\passthrough{\lstinline!ddq!} 字段。类型签名只说明 Python
表示；单位、数值范围、数组固定长度和枚举含义需查对应 IDL 头文件。
底层为浮点数；类型本身不说明单位、坐标系或业务安全范围。

\textbf{签名}

\begin{lstlisting}[language=Python]
@property
def ddq(self) -> float

@ddq.setter
def ddq(self, value: float) -> None
\end{lstlisting}

\textbf{可用性与安全性}

\textbf{\passthrough{\lstinline!AVAILABLE!} /
\passthrough{\lstinline!VALUE\_TYPE!}}：当前绑定源码已实现该消息值操作。

\textbf{参数}

\begin{longtable}[]{@{}
  >{\raggedright\arraybackslash}p{(\columnwidth - 4\tabcolsep) * \real{0.18}}
  >{\raggedright\arraybackslash}p{(\columnwidth - 4\tabcolsep) * \real{0.20}}
  >{\raggedright\arraybackslash}p{(\columnwidth - 4\tabcolsep) * \real{0.62}}@{}}
\toprule
名称 & Python 类型 & 含义 \\
\midrule
\endhead
\passthrough{\lstinline!value!} & \passthrough{\lstinline!float!} &
要写入或传递的新值，必须符合该参数的 Python 类型以及底层 C++
范围约束。 \\
\bottomrule
\end{longtable}

\textbf{返回值}

\begin{longtable}[]{@{}
  >{\raggedright\arraybackslash}p{(\columnwidth - 4\tabcolsep) * \real{0.18}}
  >{\raggedright\arraybackslash}p{(\columnwidth - 4\tabcolsep) * \real{0.20}}
  >{\raggedright\arraybackslash}p{(\columnwidth - 4\tabcolsep) * \real{0.62}}@{}}
\toprule
位置 & 类型 & 含义 \\
\midrule
\endhead
getter 返回值 & \passthrough{\lstinline!float!} & 返回属性当前值。 \\
\bottomrule
\end{longtable}

\textbf{对应 C++}

\begin{itemize}
\tightlist
\item
  类：\passthrough{\lstinline!unitree\_hg::msg::dds\_::MotorState\_!}
\item
  字段类型：\passthrough{\lstinline!float!}
\item
  声明位置：\passthrough{\lstinline!include/unitree/idl/hg/MotorState\_.hpp:27!}
\end{itemize}

\textbf{用法}

\begin{lstlisting}[language=Python]
current_value = obj.ddq
obj.ddq = new_value
\end{lstlisting}

\hypertarget{unitree-sdk2-cpp-idl-hg-motorstate-tau-est}{%
\paragraph{\texorpdfstring{\texttt{unitree\_sdk2\_cpp.idl.hg.MotorState.tau\_est}}{unitree\_sdk2\_cpp.idl.hg.MotorState.tau\_est}}\label{unitree-sdk2-cpp-idl-hg-motorstate-tau-est}}

底层 \passthrough{\lstinline!unitree\_hg::msg::dds\_::MotorState\_!} 的
\passthrough{\lstinline!tau\_est!} 字段。类型签名只说明 Python
表示；单位、数值范围、数组固定长度和枚举含义需查对应 IDL 头文件。
底层为浮点数；类型本身不说明单位、坐标系或业务安全范围。

\textbf{签名}

\begin{lstlisting}[language=Python]
@property
def tau_est(self) -> float

@tau_est.setter
def tau_est(self, value: float) -> None
\end{lstlisting}

\textbf{可用性与安全性}

\textbf{\passthrough{\lstinline!AVAILABLE!} /
\passthrough{\lstinline!VALUE\_TYPE!}}：当前绑定源码已实现该消息值操作。

\textbf{参数}

\begin{longtable}[]{@{}
  >{\raggedright\arraybackslash}p{(\columnwidth - 4\tabcolsep) * \real{0.18}}
  >{\raggedright\arraybackslash}p{(\columnwidth - 4\tabcolsep) * \real{0.20}}
  >{\raggedright\arraybackslash}p{(\columnwidth - 4\tabcolsep) * \real{0.62}}@{}}
\toprule
名称 & Python 类型 & 含义 \\
\midrule
\endhead
\passthrough{\lstinline!value!} & \passthrough{\lstinline!float!} &
要写入或传递的新值，必须符合该参数的 Python 类型以及底层 C++
范围约束。 \\
\bottomrule
\end{longtable}

\textbf{返回值}

\begin{longtable}[]{@{}
  >{\raggedright\arraybackslash}p{(\columnwidth - 4\tabcolsep) * \real{0.18}}
  >{\raggedright\arraybackslash}p{(\columnwidth - 4\tabcolsep) * \real{0.20}}
  >{\raggedright\arraybackslash}p{(\columnwidth - 4\tabcolsep) * \real{0.62}}@{}}
\toprule
位置 & 类型 & 含义 \\
\midrule
\endhead
getter 返回值 & \passthrough{\lstinline!float!} & 返回属性当前值。 \\
\bottomrule
\end{longtable}

\textbf{对应 C++}

\begin{itemize}
\tightlist
\item
  类：\passthrough{\lstinline!unitree\_hg::msg::dds\_::MotorState\_!}
\item
  字段类型：\passthrough{\lstinline!float!}
\item
  声明位置：\passthrough{\lstinline!include/unitree/idl/hg/MotorState\_.hpp:28!}
\end{itemize}

\textbf{用法}

\begin{lstlisting}[language=Python]
current_value = obj.tau_est
obj.tau_est = new_value
\end{lstlisting}

\hypertarget{unitree-sdk2-cpp-idl-hg-motorstate-temperature}{%
\paragraph{\texorpdfstring{\texttt{unitree\_sdk2\_cpp.idl.hg.MotorState.temperature}}{unitree\_sdk2\_cpp.idl.hg.MotorState.temperature}}\label{unitree-sdk2-cpp-idl-hg-motorstate-temperature}}

底层 \passthrough{\lstinline!unitree\_hg::msg::dds\_::MotorState\_!} 的
\passthrough{\lstinline!temperature!} 字段。类型签名只说明 Python
表示；单位、数值范围、数组固定长度和枚举含义需查对应 IDL 头文件。
底层是固定长度数组，写入序列必须正好包含 2 个元素；元素约束：取值范围为
-32768 到 32767。

\textbf{签名}

\begin{lstlisting}[language=Python]
@property
def temperature(self) -> list[int]

@temperature.setter
def temperature(self, value: Sequence[int]) -> None
\end{lstlisting}

\textbf{可用性与安全性}

\textbf{\passthrough{\lstinline!AVAILABLE!} /
\passthrough{\lstinline!VALUE\_TYPE!}}：当前绑定源码已实现该消息值操作。

\textbf{参数}

\begin{longtable}[]{@{}
  >{\raggedright\arraybackslash}p{(\columnwidth - 4\tabcolsep) * \real{0.18}}
  >{\raggedright\arraybackslash}p{(\columnwidth - 4\tabcolsep) * \real{0.20}}
  >{\raggedright\arraybackslash}p{(\columnwidth - 4\tabcolsep) * \real{0.62}}@{}}
\toprule
名称 & Python 类型 & 含义 \\
\midrule
\endhead
\passthrough{\lstinline!value!} &
\passthrough{\lstinline!Sequence[int]!} &
要写入或传递的新值，必须符合该参数的 Python 类型以及底层 C++
范围约束。 \\
\bottomrule
\end{longtable}

\textbf{返回值}

\begin{longtable}[]{@{}
  >{\raggedright\arraybackslash}p{(\columnwidth - 4\tabcolsep) * \real{0.18}}
  >{\raggedright\arraybackslash}p{(\columnwidth - 4\tabcolsep) * \real{0.20}}
  >{\raggedright\arraybackslash}p{(\columnwidth - 4\tabcolsep) * \real{0.62}}@{}}
\toprule
位置 & 类型 & 含义 \\
\midrule
\endhead
getter 返回值 & \passthrough{\lstinline!list[int]!} &
返回属性当前值。数组、vector 和嵌套消息采用复制语义。 \\
\bottomrule
\end{longtable}

\textbf{对应 C++}

\begin{itemize}
\tightlist
\item
  类：\passthrough{\lstinline!unitree\_hg::msg::dds\_::MotorState\_!}
\item
  字段类型：\passthrough{\lstinline!std::array<int16\_t, 2>!}
\item
  声明位置：\passthrough{\lstinline!include/unitree/idl/hg/MotorState\_.hpp:29!}
\end{itemize}

\textbf{用法}

\begin{lstlisting}[language=Python]
current_value = obj.temperature
obj.temperature = new_value
\end{lstlisting}

\begin{quote}
{[}!TIP{]} 读取容器或嵌套值后应遵循``读取、修改、重新赋值''。只修改
getter 返回的副本不会可靠地更新原 C++ 消息。
\end{quote}

\hypertarget{unitree-sdk2-cpp-idl-hg-motorstate-vol}{%
\paragraph{\texorpdfstring{\texttt{unitree\_sdk2\_cpp.idl.hg.MotorState.vol}}{unitree\_sdk2\_cpp.idl.hg.MotorState.vol}}\label{unitree-sdk2-cpp-idl-hg-motorstate-vol}}

底层 \passthrough{\lstinline!unitree\_hg::msg::dds\_::MotorState\_!} 的
\passthrough{\lstinline!vol!} 字段。类型签名只说明 Python
表示；单位、数值范围、数组固定长度和枚举含义需查对应 IDL 头文件。
底层为浮点数；类型本身不说明单位、坐标系或业务安全范围。

\textbf{签名}

\begin{lstlisting}[language=Python]
@property
def vol(self) -> float

@vol.setter
def vol(self, value: float) -> None
\end{lstlisting}

\textbf{可用性与安全性}

\textbf{\passthrough{\lstinline!AVAILABLE!} /
\passthrough{\lstinline!VALUE\_TYPE!}}：当前绑定源码已实现该消息值操作。

\textbf{参数}

\begin{longtable}[]{@{}
  >{\raggedright\arraybackslash}p{(\columnwidth - 4\tabcolsep) * \real{0.18}}
  >{\raggedright\arraybackslash}p{(\columnwidth - 4\tabcolsep) * \real{0.20}}
  >{\raggedright\arraybackslash}p{(\columnwidth - 4\tabcolsep) * \real{0.62}}@{}}
\toprule
名称 & Python 类型 & 含义 \\
\midrule
\endhead
\passthrough{\lstinline!value!} & \passthrough{\lstinline!float!} &
要写入或传递的新值，必须符合该参数的 Python 类型以及底层 C++
范围约束。 \\
\bottomrule
\end{longtable}

\textbf{返回值}

\begin{longtable}[]{@{}
  >{\raggedright\arraybackslash}p{(\columnwidth - 4\tabcolsep) * \real{0.18}}
  >{\raggedright\arraybackslash}p{(\columnwidth - 4\tabcolsep) * \real{0.20}}
  >{\raggedright\arraybackslash}p{(\columnwidth - 4\tabcolsep) * \real{0.62}}@{}}
\toprule
位置 & 类型 & 含义 \\
\midrule
\endhead
getter 返回值 & \passthrough{\lstinline!float!} & 返回属性当前值。 \\
\bottomrule
\end{longtable}

\textbf{对应 C++}

\begin{itemize}
\tightlist
\item
  类：\passthrough{\lstinline!unitree\_hg::msg::dds\_::MotorState\_!}
\item
  字段类型：\passthrough{\lstinline!float!}
\item
  声明位置：\passthrough{\lstinline!include/unitree/idl/hg/MotorState\_.hpp:30!}
\end{itemize}

\textbf{用法}

\begin{lstlisting}[language=Python]
current_value = obj.vol
obj.vol = new_value
\end{lstlisting}

\hypertarget{unitree-sdk2-cpp-idl-hg-motorstate-sensor}{%
\paragraph{\texorpdfstring{\texttt{unitree\_sdk2\_cpp.idl.hg.MotorState.sensor}}{unitree\_sdk2\_cpp.idl.hg.MotorState.sensor}}\label{unitree-sdk2-cpp-idl-hg-motorstate-sensor}}

底层 \passthrough{\lstinline!unitree\_hg::msg::dds\_::MotorState\_!} 的
\passthrough{\lstinline!sensor!} 字段。类型签名只说明 Python
表示；单位、数值范围、数组固定长度和枚举含义需查对应 IDL 头文件。
底层是固定长度数组，写入序列必须正好包含 2 个元素；元素约束：取值范围为
0 到 4294967295。

\textbf{签名}

\begin{lstlisting}[language=Python]
@property
def sensor(self) -> list[int]

@sensor.setter
def sensor(self, value: Sequence[int]) -> None
\end{lstlisting}

\textbf{可用性与安全性}

\textbf{\passthrough{\lstinline!AVAILABLE!} /
\passthrough{\lstinline!VALUE\_TYPE!}}：当前绑定源码已实现该消息值操作。

\textbf{参数}

\begin{longtable}[]{@{}
  >{\raggedright\arraybackslash}p{(\columnwidth - 4\tabcolsep) * \real{0.18}}
  >{\raggedright\arraybackslash}p{(\columnwidth - 4\tabcolsep) * \real{0.20}}
  >{\raggedright\arraybackslash}p{(\columnwidth - 4\tabcolsep) * \real{0.62}}@{}}
\toprule
名称 & Python 类型 & 含义 \\
\midrule
\endhead
\passthrough{\lstinline!value!} &
\passthrough{\lstinline!Sequence[int]!} &
要写入或传递的新值，必须符合该参数的 Python 类型以及底层 C++
范围约束。 \\
\bottomrule
\end{longtable}

\textbf{返回值}

\begin{longtable}[]{@{}
  >{\raggedright\arraybackslash}p{(\columnwidth - 4\tabcolsep) * \real{0.18}}
  >{\raggedright\arraybackslash}p{(\columnwidth - 4\tabcolsep) * \real{0.20}}
  >{\raggedright\arraybackslash}p{(\columnwidth - 4\tabcolsep) * \real{0.62}}@{}}
\toprule
位置 & 类型 & 含义 \\
\midrule
\endhead
getter 返回值 & \passthrough{\lstinline!list[int]!} &
返回属性当前值。数组、vector 和嵌套消息采用复制语义。 \\
\bottomrule
\end{longtable}

\textbf{对应 C++}

\begin{itemize}
\tightlist
\item
  类：\passthrough{\lstinline!unitree\_hg::msg::dds\_::MotorState\_!}
\item
  字段类型：\passthrough{\lstinline!std::array<uint32\_t, 2>!}
\item
  声明位置：\passthrough{\lstinline!include/unitree/idl/hg/MotorState\_.hpp:31!}
\end{itemize}

\textbf{用法}

\begin{lstlisting}[language=Python]
current_value = obj.sensor
obj.sensor = new_value
\end{lstlisting}

\begin{quote}
{[}!TIP{]} 读取容器或嵌套值后应遵循``读取、修改、重新赋值''。只修改
getter 返回的副本不会可靠地更新原 C++ 消息。
\end{quote}

\hypertarget{unitree-sdk2-cpp-idl-hg-motorstate-motorstate}{%
\paragraph{\texorpdfstring{\texttt{unitree\_sdk2\_cpp.idl.hg.MotorState.motorstate}}{unitree\_sdk2\_cpp.idl.hg.MotorState.motorstate}}\label{unitree-sdk2-cpp-idl-hg-motorstate-motorstate}}

底层 \passthrough{\lstinline!unitree\_hg::msg::dds\_::MotorState\_!} 的
\passthrough{\lstinline!motorstate!} 字段。类型签名只说明 Python
表示；单位、数值范围、数组固定长度和枚举含义需查对应 IDL 头文件。
取值范围为 0 到 4294967295。

\textbf{签名}

\begin{lstlisting}[language=Python]
@property
def motorstate(self) -> int

@motorstate.setter
def motorstate(self, value: int) -> None
\end{lstlisting}

\textbf{可用性与安全性}

\textbf{\passthrough{\lstinline!AVAILABLE!} /
\passthrough{\lstinline!VALUE\_TYPE!}}：当前绑定源码已实现该消息值操作。

\textbf{参数}

\begin{longtable}[]{@{}
  >{\raggedright\arraybackslash}p{(\columnwidth - 4\tabcolsep) * \real{0.18}}
  >{\raggedright\arraybackslash}p{(\columnwidth - 4\tabcolsep) * \real{0.20}}
  >{\raggedright\arraybackslash}p{(\columnwidth - 4\tabcolsep) * \real{0.62}}@{}}
\toprule
名称 & Python 类型 & 含义 \\
\midrule
\endhead
\passthrough{\lstinline!value!} & \passthrough{\lstinline!int!} &
要写入或传递的新值，必须符合该参数的 Python 类型以及底层 C++
范围约束。 \\
\bottomrule
\end{longtable}

\textbf{返回值}

\begin{longtable}[]{@{}
  >{\raggedright\arraybackslash}p{(\columnwidth - 4\tabcolsep) * \real{0.18}}
  >{\raggedright\arraybackslash}p{(\columnwidth - 4\tabcolsep) * \real{0.20}}
  >{\raggedright\arraybackslash}p{(\columnwidth - 4\tabcolsep) * \real{0.62}}@{}}
\toprule
位置 & 类型 & 含义 \\
\midrule
\endhead
getter 返回值 & \passthrough{\lstinline!int!} & 返回属性当前值。 \\
\bottomrule
\end{longtable}

\textbf{对应 C++}

\begin{itemize}
\tightlist
\item
  类：\passthrough{\lstinline!unitree\_hg::msg::dds\_::MotorState\_!}
\item
  字段类型：\passthrough{\lstinline!uint32\_t!}
\item
  声明位置：\passthrough{\lstinline!include/unitree/idl/hg/MotorState\_.hpp:32!}
\end{itemize}

\textbf{用法}

\begin{lstlisting}[language=Python]
current_value = obj.motorstate
obj.motorstate = new_value
\end{lstlisting}

\hypertarget{unitree-sdk2-cpp-idl-hg-motorstate-reserve}{%
\paragraph{\texorpdfstring{\texttt{unitree\_sdk2\_cpp.idl.hg.MotorState.reserve}}{unitree\_sdk2\_cpp.idl.hg.MotorState.reserve}}\label{unitree-sdk2-cpp-idl-hg-motorstate-reserve}}

底层 \passthrough{\lstinline!unitree\_hg::msg::dds\_::MotorState\_!} 的
\passthrough{\lstinline!reserve!} 字段。类型签名只说明 Python
表示；单位、数值范围、数组固定长度和枚举含义需查对应 IDL 头文件。
底层是固定长度数组，写入序列必须正好包含 4 个元素；元素约束：取值范围为
0 到 4294967295。

\textbf{签名}

\begin{lstlisting}[language=Python]
@property
def reserve(self) -> list[int]

@reserve.setter
def reserve(self, value: Sequence[int]) -> None
\end{lstlisting}

\textbf{可用性与安全性}

\textbf{\passthrough{\lstinline!AVAILABLE!} /
\passthrough{\lstinline!VALUE\_TYPE!}}：当前绑定源码已实现该消息值操作。

\textbf{参数}

\begin{longtable}[]{@{}
  >{\raggedright\arraybackslash}p{(\columnwidth - 4\tabcolsep) * \real{0.18}}
  >{\raggedright\arraybackslash}p{(\columnwidth - 4\tabcolsep) * \real{0.20}}
  >{\raggedright\arraybackslash}p{(\columnwidth - 4\tabcolsep) * \real{0.62}}@{}}
\toprule
名称 & Python 类型 & 含义 \\
\midrule
\endhead
\passthrough{\lstinline!value!} &
\passthrough{\lstinline!Sequence[int]!} &
要写入或传递的新值，必须符合该参数的 Python 类型以及底层 C++
范围约束。 \\
\bottomrule
\end{longtable}

\textbf{返回值}

\begin{longtable}[]{@{}
  >{\raggedright\arraybackslash}p{(\columnwidth - 4\tabcolsep) * \real{0.18}}
  >{\raggedright\arraybackslash}p{(\columnwidth - 4\tabcolsep) * \real{0.20}}
  >{\raggedright\arraybackslash}p{(\columnwidth - 4\tabcolsep) * \real{0.62}}@{}}
\toprule
位置 & 类型 & 含义 \\
\midrule
\endhead
getter 返回值 & \passthrough{\lstinline!list[int]!} &
返回属性当前值。数组、vector 和嵌套消息采用复制语义。 \\
\bottomrule
\end{longtable}

\textbf{对应 C++}

\begin{itemize}
\tightlist
\item
  类：\passthrough{\lstinline!unitree\_hg::msg::dds\_::MotorState\_!}
\item
  字段类型：\passthrough{\lstinline!std::array<uint32\_t, 4>!}
\item
  声明位置：\passthrough{\lstinline!include/unitree/idl/hg/MotorState\_.hpp:33!}
\end{itemize}

\textbf{用法}

\begin{lstlisting}[language=Python]
current_value = obj.reserve
obj.reserve = new_value
\end{lstlisting}

\begin{quote}
{[}!TIP{]} 读取容器或嵌套值后应遵循``读取、修改、重新赋值''。只修改
getter 返回的副本不会可靠地更新原 C++ 消息。
\end{quote}

\hypertarget{unitree-sdk2-cpp-idl-hg-presssensorstate}{%
\subsubsection{\texorpdfstring{\texttt{unitree\_sdk2\_cpp.idl.hg.PressSensorState}}{unitree\_sdk2\_cpp.idl.hg.PressSensorState}}\label{unitree-sdk2-cpp-idl-hg-presssensorstate}}

可由 Python 读写的 DDS/IDL 消息值类型。字段采用复制语义。

\textbf{导入}

\begin{lstlisting}[language=Python]
from unitree_sdk2_cpp.idl.hg import PressSensorState
\end{lstlisting}

\textbf{方法索引}

\begin{longtable}[]{@{}
  >{\raggedright\arraybackslash}p{(\columnwidth - 6\tabcolsep) * \real{0.18}}
  >{\raggedright\arraybackslash}p{(\columnwidth - 6\tabcolsep) * \real{0.24}}
  >{\raggedright\arraybackslash}p{(\columnwidth - 6\tabcolsep) * \real{0.18}}
  >{\raggedright\arraybackslash}p{(\columnwidth - 6\tabcolsep) * \real{0.40}}@{}}
\toprule
方法 & Python 签名 & 状态 & 安全分类 \\
\midrule
\endhead
\protect\hyperlink{unitree-sdk2-cpp-idl-hg-presssensorstate-dunder-init-1}{\passthrough{\lstinline!\_\_init\_\_!}}
& \passthrough{\lstinline!def \_\_init\_\_(self) -> None!} &
\passthrough{\lstinline!AVAILABLE!} &
\passthrough{\lstinline!VALUE\_TYPE!} \\
\protect\hyperlink{unitree-sdk2-cpp-idl-hg-presssensorstate-dunder-eq-1}{\passthrough{\lstinline!\_\_eq\_\_!}}
& \passthrough{\lstinline!def \_\_eq\_\_(self, other: object) -> bool!}
& \passthrough{\lstinline!AVAILABLE!} &
\passthrough{\lstinline!VALUE\_TYPE!} \\
\protect\hyperlink{unitree-sdk2-cpp-idl-hg-presssensorstate-dunder-ne-1}{\passthrough{\lstinline!\_\_ne\_\_!}}
& \passthrough{\lstinline!def \_\_ne\_\_(self, other: object) -> bool!}
& \passthrough{\lstinline!AVAILABLE!} &
\passthrough{\lstinline!VALUE\_TYPE!} \\
\bottomrule
\end{longtable}

\textbf{属性索引}

\begin{longtable}[]{@{}
  >{\raggedright\arraybackslash}p{(\columnwidth - 6\tabcolsep) * \real{0.18}}
  >{\raggedright\arraybackslash}p{(\columnwidth - 6\tabcolsep) * \real{0.24}}
  >{\raggedright\arraybackslash}p{(\columnwidth - 6\tabcolsep) * \real{0.18}}
  >{\raggedright\arraybackslash}p{(\columnwidth - 6\tabcolsep) * \real{0.40}}@{}}
\toprule
属性 & 读取类型 & 写入类型 & 状态 \\
\midrule
\endhead
\protect\hyperlink{unitree-sdk2-cpp-idl-hg-presssensorstate-pressure}{\passthrough{\lstinline!pressure!}}
& \passthrough{\lstinline!list[float]!} &
\passthrough{\lstinline!Sequence[float]!} &
\passthrough{\lstinline!AVAILABLE!} \\
\protect\hyperlink{unitree-sdk2-cpp-idl-hg-presssensorstate-temperature}{\passthrough{\lstinline!temperature!}}
& \passthrough{\lstinline!list[float]!} &
\passthrough{\lstinline!Sequence[float]!} &
\passthrough{\lstinline!AVAILABLE!} \\
\protect\hyperlink{unitree-sdk2-cpp-idl-hg-presssensorstate-lost}{\passthrough{\lstinline!lost!}}
& \passthrough{\lstinline!int!} & \passthrough{\lstinline!int!} &
\passthrough{\lstinline!AVAILABLE!} \\
\protect\hyperlink{unitree-sdk2-cpp-idl-hg-presssensorstate-reserve}{\passthrough{\lstinline!reserve!}}
& \passthrough{\lstinline!int!} & \passthrough{\lstinline!int!} &
\passthrough{\lstinline!AVAILABLE!} \\
\bottomrule
\end{longtable}

\hypertarget{unitree-sdk2-cpp-idl-hg-presssensorstate-dunder-init-1}{%
\paragraph{\texorpdfstring{\texttt{unitree\_sdk2\_cpp.idl.hg.PressSensorState.\_\_init\_\_}}{unitree\_sdk2\_cpp.idl.hg.PressSensorState.\_\_init\_\_}}\label{unitree-sdk2-cpp-idl-hg-presssensorstate-dunder-init-1}}

初始化 \passthrough{\lstinline!PressSensorState!}
实例。是否能实际构造取决于下方可用性状态。

\textbf{签名}

\begin{lstlisting}[language=Python]
def __init__(self) -> None
\end{lstlisting}

\textbf{可用性与安全性}

\textbf{\passthrough{\lstinline!AVAILABLE!} /
\passthrough{\lstinline!VALUE\_TYPE!}}：当前绑定源码已实现该消息值操作。

\textbf{参数}

无显式参数。实例方法中的 \passthrough{\lstinline!self!} 由 Python
自动传入。

\textbf{返回值}

\begin{longtable}[]{@{}
  >{\raggedright\arraybackslash}p{(\columnwidth - 4\tabcolsep) * \real{0.18}}
  >{\raggedright\arraybackslash}p{(\columnwidth - 4\tabcolsep) * \real{0.20}}
  >{\raggedright\arraybackslash}p{(\columnwidth - 4\tabcolsep) * \real{0.62}}@{}}
\toprule
位置 & 类型 & 含义 \\
\midrule
\endhead
无 & \passthrough{\lstinline!None!} &
\passthrough{\lstinline!\_\_init\_\_!}
本身不返回值；调用类对象时，在构造函数可用的前提下得到该类实例。 \\
\bottomrule
\end{longtable}

\textbf{对应 C++}

\begin{itemize}
\tightlist
\item
  类：\passthrough{\lstinline!unitree\_hg::msg::dds\_::PressSensorState\_!}
\item
  签名：\passthrough{\lstinline!PressSensorState\_()!}
\item
  绑定策略：\passthrough{\lstinline!未单独标注!}
\item
  声明位置：\passthrough{\lstinline!include/unitree/idl/hg/PressSensorState\_.hpp:30!}
\end{itemize}

\textbf{说明}

\begin{itemize}
\tightlist
\item
  示例是局部调用片段；\passthrough{\lstinline!obj!}
  和参数变量表示已经按上层业务校验并准备好的对象。
\end{itemize}

\textbf{用法}

\begin{lstlisting}[language=Python]
value = PressSensorState()
\end{lstlisting}

\hypertarget{unitree-sdk2-cpp-idl-hg-presssensorstate-dunder-eq-1}{%
\paragraph{\texorpdfstring{\texttt{unitree\_sdk2\_cpp.idl.hg.PressSensorState.\_\_eq\_\_}}{unitree\_sdk2\_cpp.idl.hg.PressSensorState.\_\_eq\_\_}}\label{unitree-sdk2-cpp-idl-hg-presssensorstate-dunder-eq-1}}

按底层 IDL 消息内容比较两个对象是否相等。

\textbf{签名}

\begin{lstlisting}[language=Python]
def __eq__(self, other: object) -> bool
\end{lstlisting}

\textbf{可用性与安全性}

\textbf{\passthrough{\lstinline!AVAILABLE!} /
\passthrough{\lstinline!VALUE\_TYPE!}}：当前绑定源码已实现该消息值操作。

\textbf{参数}

\begin{longtable}[]{@{}
  >{\raggedright\arraybackslash}p{(\columnwidth - 6\tabcolsep) * \real{0.18}}
  >{\raggedright\arraybackslash}p{(\columnwidth - 6\tabcolsep) * \real{0.24}}
  >{\raggedright\arraybackslash}p{(\columnwidth - 6\tabcolsep) * \real{0.18}}
  >{\raggedright\arraybackslash}p{(\columnwidth - 6\tabcolsep) * \real{0.40}}@{}}
\toprule
名称 & Python 类型 & 默认值 & 含义 \\
\midrule
\endhead
\passthrough{\lstinline!other!} & \passthrough{\lstinline!object!} &
必填 & 用于比较的另一个 Python 对象。类型不兼容时比较结果通常为
\passthrough{\lstinline!False!}。 \\
\bottomrule
\end{longtable}

\textbf{返回值}

\begin{longtable}[]{@{}
  >{\raggedright\arraybackslash}p{(\columnwidth - 4\tabcolsep) * \real{0.18}}
  >{\raggedright\arraybackslash}p{(\columnwidth - 4\tabcolsep) * \real{0.20}}
  >{\raggedright\arraybackslash}p{(\columnwidth - 4\tabcolsep) * \real{0.62}}@{}}
\toprule
位置 & 类型 & 含义 \\
\midrule
\endhead
返回值 & \passthrough{\lstinline!bool!} & 返回
\passthrough{\lstinline!bool!}。更精确的业务含义见该方法用途、C++
签名和目标型号协议。 \\
\bottomrule
\end{longtable}

\textbf{对应 C++}

\begin{itemize}
\tightlist
\item
  类：\passthrough{\lstinline!unitree\_hg::msg::dds\_::PressSensorState\_!}
\item
  签名：\passthrough{\lstinline!operator==(const value \&) const!}
\item
  绑定策略：\passthrough{\lstinline!未单独标注!}
\item
  声明位置：由 pybind11 手工绑定或生成注册表提供；manifest
  未记录独立头文件行号。
\end{itemize}

\textbf{说明}

\begin{itemize}
\tightlist
\item
  示例是局部调用片段；\passthrough{\lstinline!obj!}
  和参数变量表示已经按上层业务校验并准备好的对象。
\end{itemize}

\textbf{用法}

\begin{lstlisting}[language=Python]
same = left == right
\end{lstlisting}

\hypertarget{unitree-sdk2-cpp-idl-hg-presssensorstate-dunder-ne-1}{%
\paragraph{\texorpdfstring{\texttt{unitree\_sdk2\_cpp.idl.hg.PressSensorState.\_\_ne\_\_}}{unitree\_sdk2\_cpp.idl.hg.PressSensorState.\_\_ne\_\_}}\label{unitree-sdk2-cpp-idl-hg-presssensorstate-dunder-ne-1}}

按底层 IDL 消息内容比较两个对象是否不相等。

\textbf{签名}

\begin{lstlisting}[language=Python]
def __ne__(self, other: object) -> bool
\end{lstlisting}

\textbf{可用性与安全性}

\textbf{\passthrough{\lstinline!AVAILABLE!} /
\passthrough{\lstinline!VALUE\_TYPE!}}：当前绑定源码已实现该消息值操作。

\textbf{参数}

\begin{longtable}[]{@{}
  >{\raggedright\arraybackslash}p{(\columnwidth - 6\tabcolsep) * \real{0.18}}
  >{\raggedright\arraybackslash}p{(\columnwidth - 6\tabcolsep) * \real{0.24}}
  >{\raggedright\arraybackslash}p{(\columnwidth - 6\tabcolsep) * \real{0.18}}
  >{\raggedright\arraybackslash}p{(\columnwidth - 6\tabcolsep) * \real{0.40}}@{}}
\toprule
名称 & Python 类型 & 默认值 & 含义 \\
\midrule
\endhead
\passthrough{\lstinline!other!} & \passthrough{\lstinline!object!} &
必填 & 用于比较的另一个 Python 对象。类型不兼容时比较结果通常为
\passthrough{\lstinline!False!}。 \\
\bottomrule
\end{longtable}

\textbf{返回值}

\begin{longtable}[]{@{}
  >{\raggedright\arraybackslash}p{(\columnwidth - 4\tabcolsep) * \real{0.18}}
  >{\raggedright\arraybackslash}p{(\columnwidth - 4\tabcolsep) * \real{0.20}}
  >{\raggedright\arraybackslash}p{(\columnwidth - 4\tabcolsep) * \real{0.62}}@{}}
\toprule
位置 & 类型 & 含义 \\
\midrule
\endhead
返回值 & \passthrough{\lstinline!bool!} & 返回
\passthrough{\lstinline!bool!}。更精确的业务含义见该方法用途、C++
签名和目标型号协议。 \\
\bottomrule
\end{longtable}

\textbf{对应 C++}

\begin{itemize}
\tightlist
\item
  类：\passthrough{\lstinline!unitree\_hg::msg::dds\_::PressSensorState\_!}
\item
  签名：\passthrough{\lstinline"operator!=(const value \&) const"}
\item
  绑定策略：\passthrough{\lstinline!未单独标注!}
\item
  声明位置：由 pybind11 手工绑定或生成注册表提供；manifest
  未记录独立头文件行号。
\end{itemize}

\textbf{说明}

\begin{itemize}
\tightlist
\item
  示例是局部调用片段；\passthrough{\lstinline!obj!}
  和参数变量表示已经按上层业务校验并准备好的对象。
\end{itemize}

\textbf{用法}

\begin{lstlisting}[language=Python]
different = left != right
\end{lstlisting}

\hypertarget{unitree-sdk2-cpp-idl-hg-presssensorstate-pressure}{%
\paragraph{\texorpdfstring{\texttt{unitree\_sdk2\_cpp.idl.hg.PressSensorState.pressure}}{unitree\_sdk2\_cpp.idl.hg.PressSensorState.pressure}}\label{unitree-sdk2-cpp-idl-hg-presssensorstate-pressure}}

底层
\passthrough{\lstinline!unitree\_hg::msg::dds\_::PressSensorState\_!} 的
\passthrough{\lstinline!pressure!} 字段。类型签名只说明 Python
表示；单位、数值范围、数组固定长度和枚举含义需查对应 IDL 头文件。
底层是固定长度数组，写入序列必须正好包含 12
个元素；元素约束：底层为浮点数；类型本身不说明单位、坐标系或业务安全范围。

\textbf{签名}

\begin{lstlisting}[language=Python]
@property
def pressure(self) -> list[float]

@pressure.setter
def pressure(self, value: Sequence[float]) -> None
\end{lstlisting}

\textbf{可用性与安全性}

\textbf{\passthrough{\lstinline!AVAILABLE!} /
\passthrough{\lstinline!VALUE\_TYPE!}}：当前绑定源码已实现该消息值操作。

\textbf{参数}

\begin{longtable}[]{@{}
  >{\raggedright\arraybackslash}p{(\columnwidth - 4\tabcolsep) * \real{0.18}}
  >{\raggedright\arraybackslash}p{(\columnwidth - 4\tabcolsep) * \real{0.20}}
  >{\raggedright\arraybackslash}p{(\columnwidth - 4\tabcolsep) * \real{0.62}}@{}}
\toprule
名称 & Python 类型 & 含义 \\
\midrule
\endhead
\passthrough{\lstinline!value!} &
\passthrough{\lstinline!Sequence[float]!} &
要写入或传递的新值，必须符合该参数的 Python 类型以及底层 C++
范围约束。 \\
\bottomrule
\end{longtable}

\textbf{返回值}

\begin{longtable}[]{@{}
  >{\raggedright\arraybackslash}p{(\columnwidth - 4\tabcolsep) * \real{0.18}}
  >{\raggedright\arraybackslash}p{(\columnwidth - 4\tabcolsep) * \real{0.20}}
  >{\raggedright\arraybackslash}p{(\columnwidth - 4\tabcolsep) * \real{0.62}}@{}}
\toprule
位置 & 类型 & 含义 \\
\midrule
\endhead
getter 返回值 & \passthrough{\lstinline!list[float]!} &
返回属性当前值。数组、vector 和嵌套消息采用复制语义。 \\
\bottomrule
\end{longtable}

\textbf{对应 C++}

\begin{itemize}
\tightlist
\item
  类：\passthrough{\lstinline!unitree\_hg::msg::dds\_::PressSensorState\_!}
\item
  字段类型：\passthrough{\lstinline!std::array<float, 12>!}
\item
  声明位置：\passthrough{\lstinline!include/unitree/idl/hg/PressSensorState\_.hpp:24!}
\end{itemize}

\textbf{用法}

\begin{lstlisting}[language=Python]
current_value = obj.pressure
obj.pressure = new_value
\end{lstlisting}

\begin{quote}
{[}!TIP{]} 读取容器或嵌套值后应遵循``读取、修改、重新赋值''。只修改
getter 返回的副本不会可靠地更新原 C++ 消息。
\end{quote}

\hypertarget{unitree-sdk2-cpp-idl-hg-presssensorstate-temperature}{%
\paragraph{\texorpdfstring{\texttt{unitree\_sdk2\_cpp.idl.hg.PressSensorState.temperature}}{unitree\_sdk2\_cpp.idl.hg.PressSensorState.temperature}}\label{unitree-sdk2-cpp-idl-hg-presssensorstate-temperature}}

底层
\passthrough{\lstinline!unitree\_hg::msg::dds\_::PressSensorState\_!} 的
\passthrough{\lstinline!temperature!} 字段。类型签名只说明 Python
表示；单位、数值范围、数组固定长度和枚举含义需查对应 IDL 头文件。
底层是固定长度数组，写入序列必须正好包含 12
个元素；元素约束：底层为浮点数；类型本身不说明单位、坐标系或业务安全范围。

\textbf{签名}

\begin{lstlisting}[language=Python]
@property
def temperature(self) -> list[float]

@temperature.setter
def temperature(self, value: Sequence[float]) -> None
\end{lstlisting}

\textbf{可用性与安全性}

\textbf{\passthrough{\lstinline!AVAILABLE!} /
\passthrough{\lstinline!VALUE\_TYPE!}}：当前绑定源码已实现该消息值操作。

\textbf{参数}

\begin{longtable}[]{@{}
  >{\raggedright\arraybackslash}p{(\columnwidth - 4\tabcolsep) * \real{0.18}}
  >{\raggedright\arraybackslash}p{(\columnwidth - 4\tabcolsep) * \real{0.20}}
  >{\raggedright\arraybackslash}p{(\columnwidth - 4\tabcolsep) * \real{0.62}}@{}}
\toprule
名称 & Python 类型 & 含义 \\
\midrule
\endhead
\passthrough{\lstinline!value!} &
\passthrough{\lstinline!Sequence[float]!} &
要写入或传递的新值，必须符合该参数的 Python 类型以及底层 C++
范围约束。 \\
\bottomrule
\end{longtable}

\textbf{返回值}

\begin{longtable}[]{@{}
  >{\raggedright\arraybackslash}p{(\columnwidth - 4\tabcolsep) * \real{0.18}}
  >{\raggedright\arraybackslash}p{(\columnwidth - 4\tabcolsep) * \real{0.20}}
  >{\raggedright\arraybackslash}p{(\columnwidth - 4\tabcolsep) * \real{0.62}}@{}}
\toprule
位置 & 类型 & 含义 \\
\midrule
\endhead
getter 返回值 & \passthrough{\lstinline!list[float]!} &
返回属性当前值。数组、vector 和嵌套消息采用复制语义。 \\
\bottomrule
\end{longtable}

\textbf{对应 C++}

\begin{itemize}
\tightlist
\item
  类：\passthrough{\lstinline!unitree\_hg::msg::dds\_::PressSensorState\_!}
\item
  字段类型：\passthrough{\lstinline!std::array<float, 12>!}
\item
  声明位置：\passthrough{\lstinline!include/unitree/idl/hg/PressSensorState\_.hpp:25!}
\end{itemize}

\textbf{用法}

\begin{lstlisting}[language=Python]
current_value = obj.temperature
obj.temperature = new_value
\end{lstlisting}

\begin{quote}
{[}!TIP{]} 读取容器或嵌套值后应遵循``读取、修改、重新赋值''。只修改
getter 返回的副本不会可靠地更新原 C++ 消息。
\end{quote}

\hypertarget{unitree-sdk2-cpp-idl-hg-presssensorstate-lost}{%
\paragraph{\texorpdfstring{\texttt{unitree\_sdk2\_cpp.idl.hg.PressSensorState.lost}}{unitree\_sdk2\_cpp.idl.hg.PressSensorState.lost}}\label{unitree-sdk2-cpp-idl-hg-presssensorstate-lost}}

底层
\passthrough{\lstinline!unitree\_hg::msg::dds\_::PressSensorState\_!} 的
\passthrough{\lstinline!lost!} 字段。类型签名只说明 Python
表示；单位、数值范围、数组固定长度和枚举含义需查对应 IDL 头文件。
取值范围为 0 到 4294967295。

\textbf{签名}

\begin{lstlisting}[language=Python]
@property
def lost(self) -> int

@lost.setter
def lost(self, value: int) -> None
\end{lstlisting}

\textbf{可用性与安全性}

\textbf{\passthrough{\lstinline!AVAILABLE!} /
\passthrough{\lstinline!VALUE\_TYPE!}}：当前绑定源码已实现该消息值操作。

\textbf{参数}

\begin{longtable}[]{@{}
  >{\raggedright\arraybackslash}p{(\columnwidth - 4\tabcolsep) * \real{0.18}}
  >{\raggedright\arraybackslash}p{(\columnwidth - 4\tabcolsep) * \real{0.20}}
  >{\raggedright\arraybackslash}p{(\columnwidth - 4\tabcolsep) * \real{0.62}}@{}}
\toprule
名称 & Python 类型 & 含义 \\
\midrule
\endhead
\passthrough{\lstinline!value!} & \passthrough{\lstinline!int!} &
要写入或传递的新值，必须符合该参数的 Python 类型以及底层 C++
范围约束。 \\
\bottomrule
\end{longtable}

\textbf{返回值}

\begin{longtable}[]{@{}
  >{\raggedright\arraybackslash}p{(\columnwidth - 4\tabcolsep) * \real{0.18}}
  >{\raggedright\arraybackslash}p{(\columnwidth - 4\tabcolsep) * \real{0.20}}
  >{\raggedright\arraybackslash}p{(\columnwidth - 4\tabcolsep) * \real{0.62}}@{}}
\toprule
位置 & 类型 & 含义 \\
\midrule
\endhead
getter 返回值 & \passthrough{\lstinline!int!} & 返回属性当前值。 \\
\bottomrule
\end{longtable}

\textbf{对应 C++}

\begin{itemize}
\tightlist
\item
  类：\passthrough{\lstinline!unitree\_hg::msg::dds\_::PressSensorState\_!}
\item
  字段类型：\passthrough{\lstinline!uint32\_t!}
\item
  声明位置：\passthrough{\lstinline!include/unitree/idl/hg/PressSensorState\_.hpp:26!}
\end{itemize}

\textbf{用法}

\begin{lstlisting}[language=Python]
current_value = obj.lost
obj.lost = new_value
\end{lstlisting}

\hypertarget{unitree-sdk2-cpp-idl-hg-presssensorstate-reserve}{%
\paragraph{\texorpdfstring{\texttt{unitree\_sdk2\_cpp.idl.hg.PressSensorState.reserve}}{unitree\_sdk2\_cpp.idl.hg.PressSensorState.reserve}}\label{unitree-sdk2-cpp-idl-hg-presssensorstate-reserve}}

底层
\passthrough{\lstinline!unitree\_hg::msg::dds\_::PressSensorState\_!} 的
\passthrough{\lstinline!reserve!} 字段。类型签名只说明 Python
表示；单位、数值范围、数组固定长度和枚举含义需查对应 IDL 头文件。
取值范围为 0 到 4294967295。

\textbf{签名}

\begin{lstlisting}[language=Python]
@property
def reserve(self) -> int

@reserve.setter
def reserve(self, value: int) -> None
\end{lstlisting}

\textbf{可用性与安全性}

\textbf{\passthrough{\lstinline!AVAILABLE!} /
\passthrough{\lstinline!VALUE\_TYPE!}}：当前绑定源码已实现该消息值操作。

\textbf{参数}

\begin{longtable}[]{@{}
  >{\raggedright\arraybackslash}p{(\columnwidth - 4\tabcolsep) * \real{0.18}}
  >{\raggedright\arraybackslash}p{(\columnwidth - 4\tabcolsep) * \real{0.20}}
  >{\raggedright\arraybackslash}p{(\columnwidth - 4\tabcolsep) * \real{0.62}}@{}}
\toprule
名称 & Python 类型 & 含义 \\
\midrule
\endhead
\passthrough{\lstinline!value!} & \passthrough{\lstinline!int!} &
要写入或传递的新值，必须符合该参数的 Python 类型以及底层 C++
范围约束。 \\
\bottomrule
\end{longtable}

\textbf{返回值}

\begin{longtable}[]{@{}
  >{\raggedright\arraybackslash}p{(\columnwidth - 4\tabcolsep) * \real{0.18}}
  >{\raggedright\arraybackslash}p{(\columnwidth - 4\tabcolsep) * \real{0.20}}
  >{\raggedright\arraybackslash}p{(\columnwidth - 4\tabcolsep) * \real{0.62}}@{}}
\toprule
位置 & 类型 & 含义 \\
\midrule
\endhead
getter 返回值 & \passthrough{\lstinline!int!} & 返回属性当前值。 \\
\bottomrule
\end{longtable}

\textbf{对应 C++}

\begin{itemize}
\tightlist
\item
  类：\passthrough{\lstinline!unitree\_hg::msg::dds\_::PressSensorState\_!}
\item
  字段类型：\passthrough{\lstinline!uint32\_t!}
\item
  声明位置：\passthrough{\lstinline!include/unitree/idl/hg/PressSensorState\_.hpp:27!}
\end{itemize}

\textbf{用法}

\begin{lstlisting}[language=Python]
current_value = obj.reserve
obj.reserve = new_value
\end{lstlisting}

\hypertarget{unitree-sdk2-cpp-idl-hg-handstate}{%
\subsubsection{\texorpdfstring{\texttt{unitree\_sdk2\_cpp.idl.hg.HandState}}{unitree\_sdk2\_cpp.idl.hg.HandState}}\label{unitree-sdk2-cpp-idl-hg-handstate}}

可由 Python 读写的 DDS/IDL 消息值类型。字段采用复制语义。

\textbf{导入}

\begin{lstlisting}[language=Python]
from unitree_sdk2_cpp.idl.hg import HandState
\end{lstlisting}

\textbf{方法索引}

\begin{longtable}[]{@{}
  >{\raggedright\arraybackslash}p{(\columnwidth - 6\tabcolsep) * \real{0.18}}
  >{\raggedright\arraybackslash}p{(\columnwidth - 6\tabcolsep) * \real{0.24}}
  >{\raggedright\arraybackslash}p{(\columnwidth - 6\tabcolsep) * \real{0.18}}
  >{\raggedright\arraybackslash}p{(\columnwidth - 6\tabcolsep) * \real{0.40}}@{}}
\toprule
方法 & Python 签名 & 状态 & 安全分类 \\
\midrule
\endhead
\protect\hyperlink{unitree-sdk2-cpp-idl-hg-handstate-dunder-init-1}{\passthrough{\lstinline!\_\_init\_\_!}}
& \passthrough{\lstinline!def \_\_init\_\_(self) -> None!} &
\passthrough{\lstinline!AVAILABLE!} &
\passthrough{\lstinline!VALUE\_TYPE!} \\
\protect\hyperlink{unitree-sdk2-cpp-idl-hg-handstate-dunder-eq-1}{\passthrough{\lstinline!\_\_eq\_\_!}}
& \passthrough{\lstinline!def \_\_eq\_\_(self, other: object) -> bool!}
& \passthrough{\lstinline!AVAILABLE!} &
\passthrough{\lstinline!VALUE\_TYPE!} \\
\protect\hyperlink{unitree-sdk2-cpp-idl-hg-handstate-dunder-ne-1}{\passthrough{\lstinline!\_\_ne\_\_!}}
& \passthrough{\lstinline!def \_\_ne\_\_(self, other: object) -> bool!}
& \passthrough{\lstinline!AVAILABLE!} &
\passthrough{\lstinline!VALUE\_TYPE!} \\
\bottomrule
\end{longtable}

\textbf{属性索引}

\begin{longtable}[]{@{}
  >{\raggedright\arraybackslash}p{(\columnwidth - 6\tabcolsep) * \real{0.18}}
  >{\raggedright\arraybackslash}p{(\columnwidth - 6\tabcolsep) * \real{0.24}}
  >{\raggedright\arraybackslash}p{(\columnwidth - 6\tabcolsep) * \real{0.18}}
  >{\raggedright\arraybackslash}p{(\columnwidth - 6\tabcolsep) * \real{0.40}}@{}}
\toprule
属性 & 读取类型 & 写入类型 & 状态 \\
\midrule
\endhead
\protect\hyperlink{unitree-sdk2-cpp-idl-hg-handstate-motor-state}{\passthrough{\lstinline!motor\_state!}}
& \passthrough{\lstinline!list[MotorState]!} &
\passthrough{\lstinline!Sequence[MotorState]!} &
\passthrough{\lstinline!AVAILABLE!} \\
\protect\hyperlink{unitree-sdk2-cpp-idl-hg-handstate-press-sensor-state}{\passthrough{\lstinline!press\_sensor\_state!}}
& \passthrough{\lstinline!list[PressSensorState]!} &
\passthrough{\lstinline!Sequence[PressSensorState]!} &
\passthrough{\lstinline!AVAILABLE!} \\
\protect\hyperlink{unitree-sdk2-cpp-idl-hg-handstate-imu-state}{\passthrough{\lstinline!imu\_state!}}
& \passthrough{\lstinline!IMUState!} &
\passthrough{\lstinline!IMUState!} &
\passthrough{\lstinline!AVAILABLE!} \\
\protect\hyperlink{unitree-sdk2-cpp-idl-hg-handstate-power-v}{\passthrough{\lstinline!power\_v!}}
& \passthrough{\lstinline!float!} & \passthrough{\lstinline!float!} &
\passthrough{\lstinline!AVAILABLE!} \\
\protect\hyperlink{unitree-sdk2-cpp-idl-hg-handstate-power-a}{\passthrough{\lstinline!power\_a!}}
& \passthrough{\lstinline!float!} & \passthrough{\lstinline!float!} &
\passthrough{\lstinline!AVAILABLE!} \\
\protect\hyperlink{unitree-sdk2-cpp-idl-hg-handstate-system-v}{\passthrough{\lstinline!system\_v!}}
& \passthrough{\lstinline!float!} & \passthrough{\lstinline!float!} &
\passthrough{\lstinline!AVAILABLE!} \\
\protect\hyperlink{unitree-sdk2-cpp-idl-hg-handstate-device-v}{\passthrough{\lstinline!device\_v!}}
& \passthrough{\lstinline!float!} & \passthrough{\lstinline!float!} &
\passthrough{\lstinline!AVAILABLE!} \\
\protect\hyperlink{unitree-sdk2-cpp-idl-hg-handstate-error}{\passthrough{\lstinline!error!}}
& \passthrough{\lstinline!list[int]!} &
\passthrough{\lstinline!Sequence[int]!} &
\passthrough{\lstinline!AVAILABLE!} \\
\protect\hyperlink{unitree-sdk2-cpp-idl-hg-handstate-reserve}{\passthrough{\lstinline!reserve!}}
& \passthrough{\lstinline!list[int]!} &
\passthrough{\lstinline!Sequence[int]!} &
\passthrough{\lstinline!AVAILABLE!} \\
\bottomrule
\end{longtable}

\hypertarget{unitree-sdk2-cpp-idl-hg-handstate-dunder-init-1}{%
\paragraph{\texorpdfstring{\texttt{unitree\_sdk2\_cpp.idl.hg.HandState.\_\_init\_\_}}{unitree\_sdk2\_cpp.idl.hg.HandState.\_\_init\_\_}}\label{unitree-sdk2-cpp-idl-hg-handstate-dunder-init-1}}

初始化 \passthrough{\lstinline!HandState!}
实例。是否能实际构造取决于下方可用性状态。

\textbf{签名}

\begin{lstlisting}[language=Python]
def __init__(self) -> None
\end{lstlisting}

\textbf{可用性与安全性}

\textbf{\passthrough{\lstinline!AVAILABLE!} /
\passthrough{\lstinline!VALUE\_TYPE!}}：当前绑定源码已实现该消息值操作。

\textbf{参数}

无显式参数。实例方法中的 \passthrough{\lstinline!self!} 由 Python
自动传入。

\textbf{返回值}

\begin{longtable}[]{@{}
  >{\raggedright\arraybackslash}p{(\columnwidth - 4\tabcolsep) * \real{0.18}}
  >{\raggedright\arraybackslash}p{(\columnwidth - 4\tabcolsep) * \real{0.20}}
  >{\raggedright\arraybackslash}p{(\columnwidth - 4\tabcolsep) * \real{0.62}}@{}}
\toprule
位置 & 类型 & 含义 \\
\midrule
\endhead
无 & \passthrough{\lstinline!None!} &
\passthrough{\lstinline!\_\_init\_\_!}
本身不返回值；调用类对象时，在构造函数可用的前提下得到该类实例。 \\
\bottomrule
\end{longtable}

\textbf{对应 C++}

\begin{itemize}
\tightlist
\item
  类：\passthrough{\lstinline!unitree\_hg::msg::dds\_::HandState\_!}
\item
  签名：\passthrough{\lstinline!HandState\_()!}
\item
  绑定策略：\passthrough{\lstinline!未单独标注!}
\item
  声明位置：\passthrough{\lstinline!include/unitree/idl/hg/HandState\_.hpp:42!}
\end{itemize}

\textbf{说明}

\begin{itemize}
\tightlist
\item
  示例是局部调用片段；\passthrough{\lstinline!obj!}
  和参数变量表示已经按上层业务校验并准备好的对象。
\end{itemize}

\textbf{用法}

\begin{lstlisting}[language=Python]
value = HandState()
\end{lstlisting}

\hypertarget{unitree-sdk2-cpp-idl-hg-handstate-dunder-eq-1}{%
\paragraph{\texorpdfstring{\texttt{unitree\_sdk2\_cpp.idl.hg.HandState.\_\_eq\_\_}}{unitree\_sdk2\_cpp.idl.hg.HandState.\_\_eq\_\_}}\label{unitree-sdk2-cpp-idl-hg-handstate-dunder-eq-1}}

按底层 IDL 消息内容比较两个对象是否相等。

\textbf{签名}

\begin{lstlisting}[language=Python]
def __eq__(self, other: object) -> bool
\end{lstlisting}

\textbf{可用性与安全性}

\textbf{\passthrough{\lstinline!AVAILABLE!} /
\passthrough{\lstinline!VALUE\_TYPE!}}：当前绑定源码已实现该消息值操作。

\textbf{参数}

\begin{longtable}[]{@{}
  >{\raggedright\arraybackslash}p{(\columnwidth - 6\tabcolsep) * \real{0.18}}
  >{\raggedright\arraybackslash}p{(\columnwidth - 6\tabcolsep) * \real{0.24}}
  >{\raggedright\arraybackslash}p{(\columnwidth - 6\tabcolsep) * \real{0.18}}
  >{\raggedright\arraybackslash}p{(\columnwidth - 6\tabcolsep) * \real{0.40}}@{}}
\toprule
名称 & Python 类型 & 默认值 & 含义 \\
\midrule
\endhead
\passthrough{\lstinline!other!} & \passthrough{\lstinline!object!} &
必填 & 用于比较的另一个 Python 对象。类型不兼容时比较结果通常为
\passthrough{\lstinline!False!}。 \\
\bottomrule
\end{longtable}

\textbf{返回值}

\begin{longtable}[]{@{}
  >{\raggedright\arraybackslash}p{(\columnwidth - 4\tabcolsep) * \real{0.18}}
  >{\raggedright\arraybackslash}p{(\columnwidth - 4\tabcolsep) * \real{0.20}}
  >{\raggedright\arraybackslash}p{(\columnwidth - 4\tabcolsep) * \real{0.62}}@{}}
\toprule
位置 & 类型 & 含义 \\
\midrule
\endhead
返回值 & \passthrough{\lstinline!bool!} & 返回
\passthrough{\lstinline!bool!}。更精确的业务含义见该方法用途、C++
签名和目标型号协议。 \\
\bottomrule
\end{longtable}

\textbf{对应 C++}

\begin{itemize}
\tightlist
\item
  类：\passthrough{\lstinline!unitree\_hg::msg::dds\_::HandState\_!}
\item
  签名：\passthrough{\lstinline!operator==(const value \&) const!}
\item
  绑定策略：\passthrough{\lstinline!未单独标注!}
\item
  声明位置：由 pybind11 手工绑定或生成注册表提供；manifest
  未记录独立头文件行号。
\end{itemize}

\textbf{说明}

\begin{itemize}
\tightlist
\item
  示例是局部调用片段；\passthrough{\lstinline!obj!}
  和参数变量表示已经按上层业务校验并准备好的对象。
\end{itemize}

\textbf{用法}

\begin{lstlisting}[language=Python]
same = left == right
\end{lstlisting}

\hypertarget{unitree-sdk2-cpp-idl-hg-handstate-dunder-ne-1}{%
\paragraph{\texorpdfstring{\texttt{unitree\_sdk2\_cpp.idl.hg.HandState.\_\_ne\_\_}}{unitree\_sdk2\_cpp.idl.hg.HandState.\_\_ne\_\_}}\label{unitree-sdk2-cpp-idl-hg-handstate-dunder-ne-1}}

按底层 IDL 消息内容比较两个对象是否不相等。

\textbf{签名}

\begin{lstlisting}[language=Python]
def __ne__(self, other: object) -> bool
\end{lstlisting}

\textbf{可用性与安全性}

\textbf{\passthrough{\lstinline!AVAILABLE!} /
\passthrough{\lstinline!VALUE\_TYPE!}}：当前绑定源码已实现该消息值操作。

\textbf{参数}

\begin{longtable}[]{@{}
  >{\raggedright\arraybackslash}p{(\columnwidth - 6\tabcolsep) * \real{0.18}}
  >{\raggedright\arraybackslash}p{(\columnwidth - 6\tabcolsep) * \real{0.24}}
  >{\raggedright\arraybackslash}p{(\columnwidth - 6\tabcolsep) * \real{0.18}}
  >{\raggedright\arraybackslash}p{(\columnwidth - 6\tabcolsep) * \real{0.40}}@{}}
\toprule
名称 & Python 类型 & 默认值 & 含义 \\
\midrule
\endhead
\passthrough{\lstinline!other!} & \passthrough{\lstinline!object!} &
必填 & 用于比较的另一个 Python 对象。类型不兼容时比较结果通常为
\passthrough{\lstinline!False!}。 \\
\bottomrule
\end{longtable}

\textbf{返回值}

\begin{longtable}[]{@{}
  >{\raggedright\arraybackslash}p{(\columnwidth - 4\tabcolsep) * \real{0.18}}
  >{\raggedright\arraybackslash}p{(\columnwidth - 4\tabcolsep) * \real{0.20}}
  >{\raggedright\arraybackslash}p{(\columnwidth - 4\tabcolsep) * \real{0.62}}@{}}
\toprule
位置 & 类型 & 含义 \\
\midrule
\endhead
返回值 & \passthrough{\lstinline!bool!} & 返回
\passthrough{\lstinline!bool!}。更精确的业务含义见该方法用途、C++
签名和目标型号协议。 \\
\bottomrule
\end{longtable}

\textbf{对应 C++}

\begin{itemize}
\tightlist
\item
  类：\passthrough{\lstinline!unitree\_hg::msg::dds\_::HandState\_!}
\item
  签名：\passthrough{\lstinline"operator!=(const value \&) const"}
\item
  绑定策略：\passthrough{\lstinline!未单独标注!}
\item
  声明位置：由 pybind11 手工绑定或生成注册表提供；manifest
  未记录独立头文件行号。
\end{itemize}

\textbf{说明}

\begin{itemize}
\tightlist
\item
  示例是局部调用片段；\passthrough{\lstinline!obj!}
  和参数变量表示已经按上层业务校验并准备好的对象。
\end{itemize}

\textbf{用法}

\begin{lstlisting}[language=Python]
different = left != right
\end{lstlisting}

\hypertarget{unitree-sdk2-cpp-idl-hg-handstate-motor-state}{%
\paragraph{\texorpdfstring{\texttt{unitree\_sdk2\_cpp.idl.hg.HandState.motor\_state}}{unitree\_sdk2\_cpp.idl.hg.HandState.motor\_state}}\label{unitree-sdk2-cpp-idl-hg-handstate-motor-state}}

底层 \passthrough{\lstinline!unitree\_hg::msg::dds\_::HandState\_!} 的
\passthrough{\lstinline!motor\_state!} 字段。类型签名只说明 Python
表示；单位、数值范围、数组固定长度和枚举含义需查对应 IDL 头文件。
底层是可变长度 vector

\textbf{签名}

\begin{lstlisting}[language=Python]
@property
def motor_state(self) -> list[MotorState]

@motor_state.setter
def motor_state(self, value: Sequence[MotorState]) -> None
\end{lstlisting}

\textbf{可用性与安全性}

\textbf{\passthrough{\lstinline!AVAILABLE!} /
\passthrough{\lstinline!VALUE\_TYPE!}}：当前绑定源码已实现该消息值操作。

\textbf{参数}

\begin{longtable}[]{@{}
  >{\raggedright\arraybackslash}p{(\columnwidth - 4\tabcolsep) * \real{0.18}}
  >{\raggedright\arraybackslash}p{(\columnwidth - 4\tabcolsep) * \real{0.20}}
  >{\raggedright\arraybackslash}p{(\columnwidth - 4\tabcolsep) * \real{0.62}}@{}}
\toprule
名称 & Python 类型 & 含义 \\
\midrule
\endhead
\passthrough{\lstinline!value!} &
\passthrough{\lstinline!Sequence[MotorState]!} &
要写入或传递的新值，必须符合该参数的 Python 类型以及底层 C++
范围约束。 \\
\bottomrule
\end{longtable}

\textbf{返回值}

\begin{longtable}[]{@{}
  >{\raggedright\arraybackslash}p{(\columnwidth - 4\tabcolsep) * \real{0.18}}
  >{\raggedright\arraybackslash}p{(\columnwidth - 4\tabcolsep) * \real{0.20}}
  >{\raggedright\arraybackslash}p{(\columnwidth - 4\tabcolsep) * \real{0.62}}@{}}
\toprule
位置 & 类型 & 含义 \\
\midrule
\endhead
getter 返回值 & \passthrough{\lstinline!list[MotorState]!} &
返回属性当前值。数组、vector 和嵌套消息采用复制语义。 \\
\bottomrule
\end{longtable}

\textbf{对应 C++}

\begin{itemize}
\tightlist
\item
  类：\passthrough{\lstinline!unitree\_hg::msg::dds\_::HandState\_!}
\item
  字段类型：\passthrough{\lstinline!std::vector< ::unitree\_hg::msg::dds\_::MotorState\_>!}
\item
  声明位置：\passthrough{\lstinline!include/unitree/idl/hg/HandState\_.hpp:31!}
\end{itemize}

\textbf{用法}

\begin{lstlisting}[language=Python]
current_value = obj.motor_state
obj.motor_state = new_value
\end{lstlisting}

\begin{quote}
{[}!TIP{]} 读取容器或嵌套值后应遵循``读取、修改、重新赋值''。只修改
getter 返回的副本不会可靠地更新原 C++ 消息。
\end{quote}

\hypertarget{unitree-sdk2-cpp-idl-hg-handstate-press-sensor-state}{%
\paragraph{\texorpdfstring{\texttt{unitree\_sdk2\_cpp.idl.hg.HandState.press\_sensor\_state}}{unitree\_sdk2\_cpp.idl.hg.HandState.press\_sensor\_state}}\label{unitree-sdk2-cpp-idl-hg-handstate-press-sensor-state}}

底层 \passthrough{\lstinline!unitree\_hg::msg::dds\_::HandState\_!} 的
\passthrough{\lstinline!press\_sensor\_state!} 字段。类型签名只说明
Python 表示；单位、数值范围、数组固定长度和枚举含义需查对应 IDL 头文件。
底层是可变长度 vector

\textbf{签名}

\begin{lstlisting}[language=Python]
@property
def press_sensor_state(self) -> list[PressSensorState]

@press_sensor_state.setter
def press_sensor_state(self, value: Sequence[PressSensorState]) -> None
\end{lstlisting}

\textbf{可用性与安全性}

\textbf{\passthrough{\lstinline!AVAILABLE!} /
\passthrough{\lstinline!VALUE\_TYPE!}}：当前绑定源码已实现该消息值操作。

\textbf{参数}

\begin{longtable}[]{@{}
  >{\raggedright\arraybackslash}p{(\columnwidth - 4\tabcolsep) * \real{0.18}}
  >{\raggedright\arraybackslash}p{(\columnwidth - 4\tabcolsep) * \real{0.20}}
  >{\raggedright\arraybackslash}p{(\columnwidth - 4\tabcolsep) * \real{0.62}}@{}}
\toprule
名称 & Python 类型 & 含义 \\
\midrule
\endhead
\passthrough{\lstinline!value!} &
\passthrough{\lstinline!Sequence[PressSensorState]!} &
要写入或传递的新值，必须符合该参数的 Python 类型以及底层 C++
范围约束。 \\
\bottomrule
\end{longtable}

\textbf{返回值}

\begin{longtable}[]{@{}
  >{\raggedright\arraybackslash}p{(\columnwidth - 4\tabcolsep) * \real{0.18}}
  >{\raggedright\arraybackslash}p{(\columnwidth - 4\tabcolsep) * \real{0.20}}
  >{\raggedright\arraybackslash}p{(\columnwidth - 4\tabcolsep) * \real{0.62}}@{}}
\toprule
位置 & 类型 & 含义 \\
\midrule
\endhead
getter 返回值 & \passthrough{\lstinline!list[PressSensorState]!} &
返回属性当前值。数组、vector 和嵌套消息采用复制语义。 \\
\bottomrule
\end{longtable}

\textbf{对应 C++}

\begin{itemize}
\tightlist
\item
  类：\passthrough{\lstinline!unitree\_hg::msg::dds\_::HandState\_!}
\item
  字段类型：\passthrough{\lstinline!std::vector< ::unitree\_hg::msg::dds\_::PressSensorState\_>!}
\item
  声明位置：\passthrough{\lstinline!include/unitree/idl/hg/HandState\_.hpp:32!}
\end{itemize}

\textbf{用法}

\begin{lstlisting}[language=Python]
current_value = obj.press_sensor_state
obj.press_sensor_state = new_value
\end{lstlisting}

\begin{quote}
{[}!TIP{]} 读取容器或嵌套值后应遵循``读取、修改、重新赋值''。只修改
getter 返回的副本不会可靠地更新原 C++ 消息。
\end{quote}

\hypertarget{unitree-sdk2-cpp-idl-hg-handstate-imu-state}{%
\paragraph{\texorpdfstring{\texttt{unitree\_sdk2\_cpp.idl.hg.HandState.imu\_state}}{unitree\_sdk2\_cpp.idl.hg.HandState.imu\_state}}\label{unitree-sdk2-cpp-idl-hg-handstate-imu-state}}

底层 \passthrough{\lstinline!unitree\_hg::msg::dds\_::HandState\_!} 的
\passthrough{\lstinline!imu\_state!} 字段。类型签名只说明 Python
表示；单位、数值范围、数组固定长度和枚举含义需查对应 IDL 头文件。

\textbf{签名}

\begin{lstlisting}[language=Python]
@property
def imu_state(self) -> IMUState

@imu_state.setter
def imu_state(self, value: IMUState) -> None
\end{lstlisting}

\textbf{可用性与安全性}

\textbf{\passthrough{\lstinline!AVAILABLE!} /
\passthrough{\lstinline!VALUE\_TYPE!}}：当前绑定源码已实现该消息值操作。

\textbf{参数}

\begin{longtable}[]{@{}
  >{\raggedright\arraybackslash}p{(\columnwidth - 4\tabcolsep) * \real{0.18}}
  >{\raggedright\arraybackslash}p{(\columnwidth - 4\tabcolsep) * \real{0.20}}
  >{\raggedright\arraybackslash}p{(\columnwidth - 4\tabcolsep) * \real{0.62}}@{}}
\toprule
名称 & Python 类型 & 含义 \\
\midrule
\endhead
\passthrough{\lstinline!value!} & \passthrough{\lstinline!IMUState!} &
要写入或传递的新值，必须符合该参数的 Python 类型以及底层 C++
范围约束。 \\
\bottomrule
\end{longtable}

\textbf{返回值}

\begin{longtable}[]{@{}
  >{\raggedright\arraybackslash}p{(\columnwidth - 4\tabcolsep) * \real{0.18}}
  >{\raggedright\arraybackslash}p{(\columnwidth - 4\tabcolsep) * \real{0.20}}
  >{\raggedright\arraybackslash}p{(\columnwidth - 4\tabcolsep) * \real{0.62}}@{}}
\toprule
位置 & 类型 & 含义 \\
\midrule
\endhead
getter 返回值 & \passthrough{\lstinline!IMUState!} &
返回属性当前值。数组、vector 和嵌套消息采用复制语义。 \\
\bottomrule
\end{longtable}

\textbf{对应 C++}

\begin{itemize}
\tightlist
\item
  类：\passthrough{\lstinline!unitree\_hg::msg::dds\_::HandState\_!}
\item
  字段类型：\passthrough{\lstinline!::unitree\_hg::msg::dds\_::IMUState\_!}
\item
  声明位置：\passthrough{\lstinline!include/unitree/idl/hg/HandState\_.hpp:33!}
\end{itemize}

\textbf{用法}

\begin{lstlisting}[language=Python]
current_value = obj.imu_state
obj.imu_state = new_value
\end{lstlisting}

\hypertarget{unitree-sdk2-cpp-idl-hg-handstate-power-v}{%
\paragraph{\texorpdfstring{\texttt{unitree\_sdk2\_cpp.idl.hg.HandState.power\_v}}{unitree\_sdk2\_cpp.idl.hg.HandState.power\_v}}\label{unitree-sdk2-cpp-idl-hg-handstate-power-v}}

底层 \passthrough{\lstinline!unitree\_hg::msg::dds\_::HandState\_!} 的
\passthrough{\lstinline!power\_v!} 字段。类型签名只说明 Python
表示；单位、数值范围、数组固定长度和枚举含义需查对应 IDL 头文件。
底层为浮点数；类型本身不说明单位、坐标系或业务安全范围。

\textbf{签名}

\begin{lstlisting}[language=Python]
@property
def power_v(self) -> float

@power_v.setter
def power_v(self, value: float) -> None
\end{lstlisting}

\textbf{可用性与安全性}

\textbf{\passthrough{\lstinline!AVAILABLE!} /
\passthrough{\lstinline!VALUE\_TYPE!}}：当前绑定源码已实现该消息值操作。

\textbf{参数}

\begin{longtable}[]{@{}
  >{\raggedright\arraybackslash}p{(\columnwidth - 4\tabcolsep) * \real{0.18}}
  >{\raggedright\arraybackslash}p{(\columnwidth - 4\tabcolsep) * \real{0.20}}
  >{\raggedright\arraybackslash}p{(\columnwidth - 4\tabcolsep) * \real{0.62}}@{}}
\toprule
名称 & Python 类型 & 含义 \\
\midrule
\endhead
\passthrough{\lstinline!value!} & \passthrough{\lstinline!float!} &
要写入或传递的新值，必须符合该参数的 Python 类型以及底层 C++
范围约束。 \\
\bottomrule
\end{longtable}

\textbf{返回值}

\begin{longtable}[]{@{}
  >{\raggedright\arraybackslash}p{(\columnwidth - 4\tabcolsep) * \real{0.18}}
  >{\raggedright\arraybackslash}p{(\columnwidth - 4\tabcolsep) * \real{0.20}}
  >{\raggedright\arraybackslash}p{(\columnwidth - 4\tabcolsep) * \real{0.62}}@{}}
\toprule
位置 & 类型 & 含义 \\
\midrule
\endhead
getter 返回值 & \passthrough{\lstinline!float!} & 返回属性当前值。 \\
\bottomrule
\end{longtable}

\textbf{对应 C++}

\begin{itemize}
\tightlist
\item
  类：\passthrough{\lstinline!unitree\_hg::msg::dds\_::HandState\_!}
\item
  字段类型：\passthrough{\lstinline!float!}
\item
  声明位置：\passthrough{\lstinline!include/unitree/idl/hg/HandState\_.hpp:34!}
\end{itemize}

\textbf{用法}

\begin{lstlisting}[language=Python]
current_value = obj.power_v
obj.power_v = new_value
\end{lstlisting}

\hypertarget{unitree-sdk2-cpp-idl-hg-handstate-power-a}{%
\paragraph{\texorpdfstring{\texttt{unitree\_sdk2\_cpp.idl.hg.HandState.power\_a}}{unitree\_sdk2\_cpp.idl.hg.HandState.power\_a}}\label{unitree-sdk2-cpp-idl-hg-handstate-power-a}}

底层 \passthrough{\lstinline!unitree\_hg::msg::dds\_::HandState\_!} 的
\passthrough{\lstinline!power\_a!} 字段。类型签名只说明 Python
表示；单位、数值范围、数组固定长度和枚举含义需查对应 IDL 头文件。
底层为浮点数；类型本身不说明单位、坐标系或业务安全范围。

\textbf{签名}

\begin{lstlisting}[language=Python]
@property
def power_a(self) -> float

@power_a.setter
def power_a(self, value: float) -> None
\end{lstlisting}

\textbf{可用性与安全性}

\textbf{\passthrough{\lstinline!AVAILABLE!} /
\passthrough{\lstinline!VALUE\_TYPE!}}：当前绑定源码已实现该消息值操作。

\textbf{参数}

\begin{longtable}[]{@{}
  >{\raggedright\arraybackslash}p{(\columnwidth - 4\tabcolsep) * \real{0.18}}
  >{\raggedright\arraybackslash}p{(\columnwidth - 4\tabcolsep) * \real{0.20}}
  >{\raggedright\arraybackslash}p{(\columnwidth - 4\tabcolsep) * \real{0.62}}@{}}
\toprule
名称 & Python 类型 & 含义 \\
\midrule
\endhead
\passthrough{\lstinline!value!} & \passthrough{\lstinline!float!} &
要写入或传递的新值，必须符合该参数的 Python 类型以及底层 C++
范围约束。 \\
\bottomrule
\end{longtable}

\textbf{返回值}

\begin{longtable}[]{@{}
  >{\raggedright\arraybackslash}p{(\columnwidth - 4\tabcolsep) * \real{0.18}}
  >{\raggedright\arraybackslash}p{(\columnwidth - 4\tabcolsep) * \real{0.20}}
  >{\raggedright\arraybackslash}p{(\columnwidth - 4\tabcolsep) * \real{0.62}}@{}}
\toprule
位置 & 类型 & 含义 \\
\midrule
\endhead
getter 返回值 & \passthrough{\lstinline!float!} & 返回属性当前值。 \\
\bottomrule
\end{longtable}

\textbf{对应 C++}

\begin{itemize}
\tightlist
\item
  类：\passthrough{\lstinline!unitree\_hg::msg::dds\_::HandState\_!}
\item
  字段类型：\passthrough{\lstinline!float!}
\item
  声明位置：\passthrough{\lstinline!include/unitree/idl/hg/HandState\_.hpp:35!}
\end{itemize}

\textbf{用法}

\begin{lstlisting}[language=Python]
current_value = obj.power_a
obj.power_a = new_value
\end{lstlisting}

\hypertarget{unitree-sdk2-cpp-idl-hg-handstate-system-v}{%
\paragraph{\texorpdfstring{\texttt{unitree\_sdk2\_cpp.idl.hg.HandState.system\_v}}{unitree\_sdk2\_cpp.idl.hg.HandState.system\_v}}\label{unitree-sdk2-cpp-idl-hg-handstate-system-v}}

底层 \passthrough{\lstinline!unitree\_hg::msg::dds\_::HandState\_!} 的
\passthrough{\lstinline!system\_v!} 字段。类型签名只说明 Python
表示；单位、数值范围、数组固定长度和枚举含义需查对应 IDL 头文件。
底层为浮点数；类型本身不说明单位、坐标系或业务安全范围。

\textbf{签名}

\begin{lstlisting}[language=Python]
@property
def system_v(self) -> float

@system_v.setter
def system_v(self, value: float) -> None
\end{lstlisting}

\textbf{可用性与安全性}

\textbf{\passthrough{\lstinline!AVAILABLE!} /
\passthrough{\lstinline!VALUE\_TYPE!}}：当前绑定源码已实现该消息值操作。

\textbf{参数}

\begin{longtable}[]{@{}
  >{\raggedright\arraybackslash}p{(\columnwidth - 4\tabcolsep) * \real{0.18}}
  >{\raggedright\arraybackslash}p{(\columnwidth - 4\tabcolsep) * \real{0.20}}
  >{\raggedright\arraybackslash}p{(\columnwidth - 4\tabcolsep) * \real{0.62}}@{}}
\toprule
名称 & Python 类型 & 含义 \\
\midrule
\endhead
\passthrough{\lstinline!value!} & \passthrough{\lstinline!float!} &
要写入或传递的新值，必须符合该参数的 Python 类型以及底层 C++
范围约束。 \\
\bottomrule
\end{longtable}

\textbf{返回值}

\begin{longtable}[]{@{}
  >{\raggedright\arraybackslash}p{(\columnwidth - 4\tabcolsep) * \real{0.18}}
  >{\raggedright\arraybackslash}p{(\columnwidth - 4\tabcolsep) * \real{0.20}}
  >{\raggedright\arraybackslash}p{(\columnwidth - 4\tabcolsep) * \real{0.62}}@{}}
\toprule
位置 & 类型 & 含义 \\
\midrule
\endhead
getter 返回值 & \passthrough{\lstinline!float!} & 返回属性当前值。 \\
\bottomrule
\end{longtable}

\textbf{对应 C++}

\begin{itemize}
\tightlist
\item
  类：\passthrough{\lstinline!unitree\_hg::msg::dds\_::HandState\_!}
\item
  字段类型：\passthrough{\lstinline!float!}
\item
  声明位置：\passthrough{\lstinline!include/unitree/idl/hg/HandState\_.hpp:36!}
\end{itemize}

\textbf{用法}

\begin{lstlisting}[language=Python]
current_value = obj.system_v
obj.system_v = new_value
\end{lstlisting}

\hypertarget{unitree-sdk2-cpp-idl-hg-handstate-device-v}{%
\paragraph{\texorpdfstring{\texttt{unitree\_sdk2\_cpp.idl.hg.HandState.device\_v}}{unitree\_sdk2\_cpp.idl.hg.HandState.device\_v}}\label{unitree-sdk2-cpp-idl-hg-handstate-device-v}}

底层 \passthrough{\lstinline!unitree\_hg::msg::dds\_::HandState\_!} 的
\passthrough{\lstinline!device\_v!} 字段。类型签名只说明 Python
表示；单位、数值范围、数组固定长度和枚举含义需查对应 IDL 头文件。
底层为浮点数；类型本身不说明单位、坐标系或业务安全范围。

\textbf{签名}

\begin{lstlisting}[language=Python]
@property
def device_v(self) -> float

@device_v.setter
def device_v(self, value: float) -> None
\end{lstlisting}

\textbf{可用性与安全性}

\textbf{\passthrough{\lstinline!AVAILABLE!} /
\passthrough{\lstinline!VALUE\_TYPE!}}：当前绑定源码已实现该消息值操作。

\textbf{参数}

\begin{longtable}[]{@{}
  >{\raggedright\arraybackslash}p{(\columnwidth - 4\tabcolsep) * \real{0.18}}
  >{\raggedright\arraybackslash}p{(\columnwidth - 4\tabcolsep) * \real{0.20}}
  >{\raggedright\arraybackslash}p{(\columnwidth - 4\tabcolsep) * \real{0.62}}@{}}
\toprule
名称 & Python 类型 & 含义 \\
\midrule
\endhead
\passthrough{\lstinline!value!} & \passthrough{\lstinline!float!} &
要写入或传递的新值，必须符合该参数的 Python 类型以及底层 C++
范围约束。 \\
\bottomrule
\end{longtable}

\textbf{返回值}

\begin{longtable}[]{@{}
  >{\raggedright\arraybackslash}p{(\columnwidth - 4\tabcolsep) * \real{0.18}}
  >{\raggedright\arraybackslash}p{(\columnwidth - 4\tabcolsep) * \real{0.20}}
  >{\raggedright\arraybackslash}p{(\columnwidth - 4\tabcolsep) * \real{0.62}}@{}}
\toprule
位置 & 类型 & 含义 \\
\midrule
\endhead
getter 返回值 & \passthrough{\lstinline!float!} & 返回属性当前值。 \\
\bottomrule
\end{longtable}

\textbf{对应 C++}

\begin{itemize}
\tightlist
\item
  类：\passthrough{\lstinline!unitree\_hg::msg::dds\_::HandState\_!}
\item
  字段类型：\passthrough{\lstinline!float!}
\item
  声明位置：\passthrough{\lstinline!include/unitree/idl/hg/HandState\_.hpp:37!}
\end{itemize}

\textbf{用法}

\begin{lstlisting}[language=Python]
current_value = obj.device_v
obj.device_v = new_value
\end{lstlisting}

\hypertarget{unitree-sdk2-cpp-idl-hg-handstate-error}{%
\paragraph{\texorpdfstring{\texttt{unitree\_sdk2\_cpp.idl.hg.HandState.error}}{unitree\_sdk2\_cpp.idl.hg.HandState.error}}\label{unitree-sdk2-cpp-idl-hg-handstate-error}}

底层 \passthrough{\lstinline!unitree\_hg::msg::dds\_::HandState\_!} 的
\passthrough{\lstinline!error!} 字段。类型签名只说明 Python
表示；单位、数值范围、数组固定长度和枚举含义需查对应 IDL 头文件。
底层是固定长度数组，写入序列必须正好包含 2 个元素；元素约束：取值范围为
0 到 4294967295。

\textbf{签名}

\begin{lstlisting}[language=Python]
@property
def error(self) -> list[int]

@error.setter
def error(self, value: Sequence[int]) -> None
\end{lstlisting}

\textbf{可用性与安全性}

\textbf{\passthrough{\lstinline!AVAILABLE!} /
\passthrough{\lstinline!VALUE\_TYPE!}}：当前绑定源码已实现该消息值操作。

\textbf{参数}

\begin{longtable}[]{@{}
  >{\raggedright\arraybackslash}p{(\columnwidth - 4\tabcolsep) * \real{0.18}}
  >{\raggedright\arraybackslash}p{(\columnwidth - 4\tabcolsep) * \real{0.20}}
  >{\raggedright\arraybackslash}p{(\columnwidth - 4\tabcolsep) * \real{0.62}}@{}}
\toprule
名称 & Python 类型 & 含义 \\
\midrule
\endhead
\passthrough{\lstinline!value!} &
\passthrough{\lstinline!Sequence[int]!} &
要写入或传递的新值，必须符合该参数的 Python 类型以及底层 C++
范围约束。 \\
\bottomrule
\end{longtable}

\textbf{返回值}

\begin{longtable}[]{@{}
  >{\raggedright\arraybackslash}p{(\columnwidth - 4\tabcolsep) * \real{0.18}}
  >{\raggedright\arraybackslash}p{(\columnwidth - 4\tabcolsep) * \real{0.20}}
  >{\raggedright\arraybackslash}p{(\columnwidth - 4\tabcolsep) * \real{0.62}}@{}}
\toprule
位置 & 类型 & 含义 \\
\midrule
\endhead
getter 返回值 & \passthrough{\lstinline!list[int]!} &
返回属性当前值。数组、vector 和嵌套消息采用复制语义。 \\
\bottomrule
\end{longtable}

\textbf{对应 C++}

\begin{itemize}
\tightlist
\item
  类：\passthrough{\lstinline!unitree\_hg::msg::dds\_::HandState\_!}
\item
  字段类型：\passthrough{\lstinline!std::array<uint32\_t, 2>!}
\item
  声明位置：\passthrough{\lstinline!include/unitree/idl/hg/HandState\_.hpp:38!}
\end{itemize}

\textbf{用法}

\begin{lstlisting}[language=Python]
current_value = obj.error
obj.error = new_value
\end{lstlisting}

\begin{quote}
{[}!TIP{]} 读取容器或嵌套值后应遵循``读取、修改、重新赋值''。只修改
getter 返回的副本不会可靠地更新原 C++ 消息。
\end{quote}

\hypertarget{unitree-sdk2-cpp-idl-hg-handstate-reserve}{%
\paragraph{\texorpdfstring{\texttt{unitree\_sdk2\_cpp.idl.hg.HandState.reserve}}{unitree\_sdk2\_cpp.idl.hg.HandState.reserve}}\label{unitree-sdk2-cpp-idl-hg-handstate-reserve}}

底层 \passthrough{\lstinline!unitree\_hg::msg::dds\_::HandState\_!} 的
\passthrough{\lstinline!reserve!} 字段。类型签名只说明 Python
表示；单位、数值范围、数组固定长度和枚举含义需查对应 IDL 头文件。
底层是固定长度数组，写入序列必须正好包含 2 个元素；元素约束：取值范围为
0 到 4294967295。

\textbf{签名}

\begin{lstlisting}[language=Python]
@property
def reserve(self) -> list[int]

@reserve.setter
def reserve(self, value: Sequence[int]) -> None
\end{lstlisting}

\textbf{可用性与安全性}

\textbf{\passthrough{\lstinline!AVAILABLE!} /
\passthrough{\lstinline!VALUE\_TYPE!}}：当前绑定源码已实现该消息值操作。

\textbf{参数}

\begin{longtable}[]{@{}
  >{\raggedright\arraybackslash}p{(\columnwidth - 4\tabcolsep) * \real{0.18}}
  >{\raggedright\arraybackslash}p{(\columnwidth - 4\tabcolsep) * \real{0.20}}
  >{\raggedright\arraybackslash}p{(\columnwidth - 4\tabcolsep) * \real{0.62}}@{}}
\toprule
名称 & Python 类型 & 含义 \\
\midrule
\endhead
\passthrough{\lstinline!value!} &
\passthrough{\lstinline!Sequence[int]!} &
要写入或传递的新值，必须符合该参数的 Python 类型以及底层 C++
范围约束。 \\
\bottomrule
\end{longtable}

\textbf{返回值}

\begin{longtable}[]{@{}
  >{\raggedright\arraybackslash}p{(\columnwidth - 4\tabcolsep) * \real{0.18}}
  >{\raggedright\arraybackslash}p{(\columnwidth - 4\tabcolsep) * \real{0.20}}
  >{\raggedright\arraybackslash}p{(\columnwidth - 4\tabcolsep) * \real{0.62}}@{}}
\toprule
位置 & 类型 & 含义 \\
\midrule
\endhead
getter 返回值 & \passthrough{\lstinline!list[int]!} &
返回属性当前值。数组、vector 和嵌套消息采用复制语义。 \\
\bottomrule
\end{longtable}

\textbf{对应 C++}

\begin{itemize}
\tightlist
\item
  类：\passthrough{\lstinline!unitree\_hg::msg::dds\_::HandState\_!}
\item
  字段类型：\passthrough{\lstinline!std::array<uint32\_t, 2>!}
\item
  声明位置：\passthrough{\lstinline!include/unitree/idl/hg/HandState\_.hpp:39!}
\end{itemize}

\textbf{用法}

\begin{lstlisting}[language=Python]
current_value = obj.reserve
obj.reserve = new_value
\end{lstlisting}

\begin{quote}
{[}!TIP{]} 读取容器或嵌套值后应遵循``读取、修改、重新赋值''。只修改
getter 返回的副本不会可靠地更新原 C++ 消息。
\end{quote}

\hypertarget{unitree-sdk2-cpp-idl-hg-lowcmd}{%
\subsubsection{\texorpdfstring{\texttt{unitree\_sdk2\_cpp.idl.hg.LowCmd}}{unitree\_sdk2\_cpp.idl.hg.LowCmd}}\label{unitree-sdk2-cpp-idl-hg-lowcmd}}

可由 Python 读写的 DDS/IDL 消息值类型。字段采用复制语义。

\textbf{导入}

\begin{lstlisting}[language=Python]
from unitree_sdk2_cpp.idl.hg import LowCmd
\end{lstlisting}

\textbf{方法索引}

\begin{longtable}[]{@{}
  >{\raggedright\arraybackslash}p{(\columnwidth - 6\tabcolsep) * \real{0.18}}
  >{\raggedright\arraybackslash}p{(\columnwidth - 6\tabcolsep) * \real{0.24}}
  >{\raggedright\arraybackslash}p{(\columnwidth - 6\tabcolsep) * \real{0.18}}
  >{\raggedright\arraybackslash}p{(\columnwidth - 6\tabcolsep) * \real{0.40}}@{}}
\toprule
方法 & Python 签名 & 状态 & 安全分类 \\
\midrule
\endhead
\protect\hyperlink{unitree-sdk2-cpp-idl-hg-lowcmd-dunder-init-1}{\passthrough{\lstinline!\_\_init\_\_!}}
& \passthrough{\lstinline!def \_\_init\_\_(self) -> None!} &
\passthrough{\lstinline!AVAILABLE!} &
\passthrough{\lstinline!VALUE\_TYPE!} \\
\protect\hyperlink{unitree-sdk2-cpp-idl-hg-lowcmd-dunder-eq-1}{\passthrough{\lstinline!\_\_eq\_\_!}}
& \passthrough{\lstinline!def \_\_eq\_\_(self, other: object) -> bool!}
& \passthrough{\lstinline!AVAILABLE!} &
\passthrough{\lstinline!VALUE\_TYPE!} \\
\protect\hyperlink{unitree-sdk2-cpp-idl-hg-lowcmd-dunder-ne-1}{\passthrough{\lstinline!\_\_ne\_\_!}}
& \passthrough{\lstinline!def \_\_ne\_\_(self, other: object) -> bool!}
& \passthrough{\lstinline!AVAILABLE!} &
\passthrough{\lstinline!VALUE\_TYPE!} \\
\bottomrule
\end{longtable}

\textbf{属性索引}

\begin{longtable}[]{@{}
  >{\raggedright\arraybackslash}p{(\columnwidth - 6\tabcolsep) * \real{0.18}}
  >{\raggedright\arraybackslash}p{(\columnwidth - 6\tabcolsep) * \real{0.24}}
  >{\raggedright\arraybackslash}p{(\columnwidth - 6\tabcolsep) * \real{0.18}}
  >{\raggedright\arraybackslash}p{(\columnwidth - 6\tabcolsep) * \real{0.40}}@{}}
\toprule
属性 & 读取类型 & 写入类型 & 状态 \\
\midrule
\endhead
\protect\hyperlink{unitree-sdk2-cpp-idl-hg-lowcmd-mode-pr}{\passthrough{\lstinline!mode\_pr!}}
& \passthrough{\lstinline!int!} & \passthrough{\lstinline!int!} &
\passthrough{\lstinline!AVAILABLE!} \\
\protect\hyperlink{unitree-sdk2-cpp-idl-hg-lowcmd-mode-machine}{\passthrough{\lstinline!mode\_machine!}}
& \passthrough{\lstinline!int!} & \passthrough{\lstinline!int!} &
\passthrough{\lstinline!AVAILABLE!} \\
\protect\hyperlink{unitree-sdk2-cpp-idl-hg-lowcmd-motor-cmd}{\passthrough{\lstinline!motor\_cmd!}}
& \passthrough{\lstinline!list[MotorCmd]!} &
\passthrough{\lstinline!Sequence[MotorCmd]!} &
\passthrough{\lstinline!AVAILABLE!} \\
\protect\hyperlink{unitree-sdk2-cpp-idl-hg-lowcmd-reserve}{\passthrough{\lstinline!reserve!}}
& \passthrough{\lstinline!list[int]!} &
\passthrough{\lstinline!Sequence[int]!} &
\passthrough{\lstinline!AVAILABLE!} \\
\protect\hyperlink{unitree-sdk2-cpp-idl-hg-lowcmd-crc}{\passthrough{\lstinline!crc!}}
& \passthrough{\lstinline!int!} & \passthrough{\lstinline!int!} &
\passthrough{\lstinline!AVAILABLE!} \\
\bottomrule
\end{longtable}

\hypertarget{unitree-sdk2-cpp-idl-hg-lowcmd-dunder-init-1}{%
\paragraph{\texorpdfstring{\texttt{unitree\_sdk2\_cpp.idl.hg.LowCmd.\_\_init\_\_}}{unitree\_sdk2\_cpp.idl.hg.LowCmd.\_\_init\_\_}}\label{unitree-sdk2-cpp-idl-hg-lowcmd-dunder-init-1}}

初始化 \passthrough{\lstinline!LowCmd!}
实例。是否能实际构造取决于下方可用性状态。

\textbf{签名}

\begin{lstlisting}[language=Python]
def __init__(self) -> None
\end{lstlisting}

\textbf{可用性与安全性}

\textbf{\passthrough{\lstinline!AVAILABLE!} /
\passthrough{\lstinline!VALUE\_TYPE!}}：当前绑定源码已实现该消息值操作。

\textbf{参数}

无显式参数。实例方法中的 \passthrough{\lstinline!self!} 由 Python
自动传入。

\textbf{返回值}

\begin{longtable}[]{@{}
  >{\raggedright\arraybackslash}p{(\columnwidth - 4\tabcolsep) * \real{0.18}}
  >{\raggedright\arraybackslash}p{(\columnwidth - 4\tabcolsep) * \real{0.20}}
  >{\raggedright\arraybackslash}p{(\columnwidth - 4\tabcolsep) * \real{0.62}}@{}}
\toprule
位置 & 类型 & 含义 \\
\midrule
\endhead
无 & \passthrough{\lstinline!None!} &
\passthrough{\lstinline!\_\_init\_\_!}
本身不返回值；调用类对象时，在构造函数可用的前提下得到该类实例。 \\
\bottomrule
\end{longtable}

\textbf{对应 C++}

\begin{itemize}
\tightlist
\item
  类：\passthrough{\lstinline!unitree\_hg::msg::dds\_::LowCmd\_!}
\item
  签名：\passthrough{\lstinline!LowCmd\_()!}
\item
  绑定策略：\passthrough{\lstinline!未单独标注!}
\item
  声明位置：\passthrough{\lstinline!include/unitree/idl/hg/LowCmd\_.hpp:33!}
\end{itemize}

\textbf{说明}

\begin{itemize}
\tightlist
\item
  示例是局部调用片段；\passthrough{\lstinline!obj!}
  和参数变量表示已经按上层业务校验并准备好的对象。
\end{itemize}

\textbf{用法}

\begin{lstlisting}[language=Python]
value = LowCmd()
\end{lstlisting}

\hypertarget{unitree-sdk2-cpp-idl-hg-lowcmd-dunder-eq-1}{%
\paragraph{\texorpdfstring{\texttt{unitree\_sdk2\_cpp.idl.hg.LowCmd.\_\_eq\_\_}}{unitree\_sdk2\_cpp.idl.hg.LowCmd.\_\_eq\_\_}}\label{unitree-sdk2-cpp-idl-hg-lowcmd-dunder-eq-1}}

按底层 IDL 消息内容比较两个对象是否相等。

\textbf{签名}

\begin{lstlisting}[language=Python]
def __eq__(self, other: object) -> bool
\end{lstlisting}

\textbf{可用性与安全性}

\textbf{\passthrough{\lstinline!AVAILABLE!} /
\passthrough{\lstinline!VALUE\_TYPE!}}：当前绑定源码已实现该消息值操作。

\textbf{参数}

\begin{longtable}[]{@{}
  >{\raggedright\arraybackslash}p{(\columnwidth - 6\tabcolsep) * \real{0.18}}
  >{\raggedright\arraybackslash}p{(\columnwidth - 6\tabcolsep) * \real{0.24}}
  >{\raggedright\arraybackslash}p{(\columnwidth - 6\tabcolsep) * \real{0.18}}
  >{\raggedright\arraybackslash}p{(\columnwidth - 6\tabcolsep) * \real{0.40}}@{}}
\toprule
名称 & Python 类型 & 默认值 & 含义 \\
\midrule
\endhead
\passthrough{\lstinline!other!} & \passthrough{\lstinline!object!} &
必填 & 用于比较的另一个 Python 对象。类型不兼容时比较结果通常为
\passthrough{\lstinline!False!}。 \\
\bottomrule
\end{longtable}

\textbf{返回值}

\begin{longtable}[]{@{}
  >{\raggedright\arraybackslash}p{(\columnwidth - 4\tabcolsep) * \real{0.18}}
  >{\raggedright\arraybackslash}p{(\columnwidth - 4\tabcolsep) * \real{0.20}}
  >{\raggedright\arraybackslash}p{(\columnwidth - 4\tabcolsep) * \real{0.62}}@{}}
\toprule
位置 & 类型 & 含义 \\
\midrule
\endhead
返回值 & \passthrough{\lstinline!bool!} & 返回
\passthrough{\lstinline!bool!}。更精确的业务含义见该方法用途、C++
签名和目标型号协议。 \\
\bottomrule
\end{longtable}

\textbf{对应 C++}

\begin{itemize}
\tightlist
\item
  类：\passthrough{\lstinline!unitree\_hg::msg::dds\_::LowCmd\_!}
\item
  签名：\passthrough{\lstinline!operator==(const value \&) const!}
\item
  绑定策略：\passthrough{\lstinline!未单独标注!}
\item
  声明位置：由 pybind11 手工绑定或生成注册表提供；manifest
  未记录独立头文件行号。
\end{itemize}

\textbf{说明}

\begin{itemize}
\tightlist
\item
  示例是局部调用片段；\passthrough{\lstinline!obj!}
  和参数变量表示已经按上层业务校验并准备好的对象。
\end{itemize}

\textbf{用法}

\begin{lstlisting}[language=Python]
same = left == right
\end{lstlisting}

\hypertarget{unitree-sdk2-cpp-idl-hg-lowcmd-dunder-ne-1}{%
\paragraph{\texorpdfstring{\texttt{unitree\_sdk2\_cpp.idl.hg.LowCmd.\_\_ne\_\_}}{unitree\_sdk2\_cpp.idl.hg.LowCmd.\_\_ne\_\_}}\label{unitree-sdk2-cpp-idl-hg-lowcmd-dunder-ne-1}}

按底层 IDL 消息内容比较两个对象是否不相等。

\textbf{签名}

\begin{lstlisting}[language=Python]
def __ne__(self, other: object) -> bool
\end{lstlisting}

\textbf{可用性与安全性}

\textbf{\passthrough{\lstinline!AVAILABLE!} /
\passthrough{\lstinline!VALUE\_TYPE!}}：当前绑定源码已实现该消息值操作。

\textbf{参数}

\begin{longtable}[]{@{}
  >{\raggedright\arraybackslash}p{(\columnwidth - 6\tabcolsep) * \real{0.18}}
  >{\raggedright\arraybackslash}p{(\columnwidth - 6\tabcolsep) * \real{0.24}}
  >{\raggedright\arraybackslash}p{(\columnwidth - 6\tabcolsep) * \real{0.18}}
  >{\raggedright\arraybackslash}p{(\columnwidth - 6\tabcolsep) * \real{0.40}}@{}}
\toprule
名称 & Python 类型 & 默认值 & 含义 \\
\midrule
\endhead
\passthrough{\lstinline!other!} & \passthrough{\lstinline!object!} &
必填 & 用于比较的另一个 Python 对象。类型不兼容时比较结果通常为
\passthrough{\lstinline!False!}。 \\
\bottomrule
\end{longtable}

\textbf{返回值}

\begin{longtable}[]{@{}
  >{\raggedright\arraybackslash}p{(\columnwidth - 4\tabcolsep) * \real{0.18}}
  >{\raggedright\arraybackslash}p{(\columnwidth - 4\tabcolsep) * \real{0.20}}
  >{\raggedright\arraybackslash}p{(\columnwidth - 4\tabcolsep) * \real{0.62}}@{}}
\toprule
位置 & 类型 & 含义 \\
\midrule
\endhead
返回值 & \passthrough{\lstinline!bool!} & 返回
\passthrough{\lstinline!bool!}。更精确的业务含义见该方法用途、C++
签名和目标型号协议。 \\
\bottomrule
\end{longtable}

\textbf{对应 C++}

\begin{itemize}
\tightlist
\item
  类：\passthrough{\lstinline!unitree\_hg::msg::dds\_::LowCmd\_!}
\item
  签名：\passthrough{\lstinline"operator!=(const value \&) const"}
\item
  绑定策略：\passthrough{\lstinline!未单独标注!}
\item
  声明位置：由 pybind11 手工绑定或生成注册表提供；manifest
  未记录独立头文件行号。
\end{itemize}

\textbf{说明}

\begin{itemize}
\tightlist
\item
  示例是局部调用片段；\passthrough{\lstinline!obj!}
  和参数变量表示已经按上层业务校验并准备好的对象。
\end{itemize}

\textbf{用法}

\begin{lstlisting}[language=Python]
different = left != right
\end{lstlisting}

\hypertarget{unitree-sdk2-cpp-idl-hg-lowcmd-mode-pr}{%
\paragraph{\texorpdfstring{\texttt{unitree\_sdk2\_cpp.idl.hg.LowCmd.mode\_pr}}{unitree\_sdk2\_cpp.idl.hg.LowCmd.mode\_pr}}\label{unitree-sdk2-cpp-idl-hg-lowcmd-mode-pr}}

底层 \passthrough{\lstinline!unitree\_hg::msg::dds\_::LowCmd\_!} 的
\passthrough{\lstinline!mode\_pr!} 字段。类型签名只说明 Python
表示；单位、数值范围、数组固定长度和枚举含义需查对应 IDL 头文件。
取值范围为 0 到 255。

\textbf{签名}

\begin{lstlisting}[language=Python]
@property
def mode_pr(self) -> int

@mode_pr.setter
def mode_pr(self, value: int) -> None
\end{lstlisting}

\textbf{可用性与安全性}

\textbf{\passthrough{\lstinline!AVAILABLE!} /
\passthrough{\lstinline!VALUE\_TYPE!}}：当前绑定源码已实现该消息值操作。

\textbf{参数}

\begin{longtable}[]{@{}
  >{\raggedright\arraybackslash}p{(\columnwidth - 4\tabcolsep) * \real{0.18}}
  >{\raggedright\arraybackslash}p{(\columnwidth - 4\tabcolsep) * \real{0.20}}
  >{\raggedright\arraybackslash}p{(\columnwidth - 4\tabcolsep) * \real{0.62}}@{}}
\toprule
名称 & Python 类型 & 含义 \\
\midrule
\endhead
\passthrough{\lstinline!value!} & \passthrough{\lstinline!int!} &
要写入或传递的新值，必须符合该参数的 Python 类型以及底层 C++
范围约束。 \\
\bottomrule
\end{longtable}

\textbf{返回值}

\begin{longtable}[]{@{}
  >{\raggedright\arraybackslash}p{(\columnwidth - 4\tabcolsep) * \real{0.18}}
  >{\raggedright\arraybackslash}p{(\columnwidth - 4\tabcolsep) * \real{0.20}}
  >{\raggedright\arraybackslash}p{(\columnwidth - 4\tabcolsep) * \real{0.62}}@{}}
\toprule
位置 & 类型 & 含义 \\
\midrule
\endhead
getter 返回值 & \passthrough{\lstinline!int!} & 返回属性当前值。 \\
\bottomrule
\end{longtable}

\textbf{对应 C++}

\begin{itemize}
\tightlist
\item
  类：\passthrough{\lstinline!unitree\_hg::msg::dds\_::LowCmd\_!}
\item
  字段类型：\passthrough{\lstinline!uint8\_t!}
\item
  声明位置：\passthrough{\lstinline!include/unitree/idl/hg/LowCmd\_.hpp:26!}
\end{itemize}

\textbf{用法}

\begin{lstlisting}[language=Python]
current_value = obj.mode_pr
obj.mode_pr = new_value
\end{lstlisting}

\hypertarget{unitree-sdk2-cpp-idl-hg-lowcmd-mode-machine}{%
\paragraph{\texorpdfstring{\texttt{unitree\_sdk2\_cpp.idl.hg.LowCmd.mode\_machine}}{unitree\_sdk2\_cpp.idl.hg.LowCmd.mode\_machine}}\label{unitree-sdk2-cpp-idl-hg-lowcmd-mode-machine}}

底层 \passthrough{\lstinline!unitree\_hg::msg::dds\_::LowCmd\_!} 的
\passthrough{\lstinline!mode\_machine!} 字段。类型签名只说明 Python
表示；单位、数值范围、数组固定长度和枚举含义需查对应 IDL 头文件。
取值范围为 0 到 255。

\textbf{签名}

\begin{lstlisting}[language=Python]
@property
def mode_machine(self) -> int

@mode_machine.setter
def mode_machine(self, value: int) -> None
\end{lstlisting}

\textbf{可用性与安全性}

\textbf{\passthrough{\lstinline!AVAILABLE!} /
\passthrough{\lstinline!VALUE\_TYPE!}}：当前绑定源码已实现该消息值操作。

\textbf{参数}

\begin{longtable}[]{@{}
  >{\raggedright\arraybackslash}p{(\columnwidth - 4\tabcolsep) * \real{0.18}}
  >{\raggedright\arraybackslash}p{(\columnwidth - 4\tabcolsep) * \real{0.20}}
  >{\raggedright\arraybackslash}p{(\columnwidth - 4\tabcolsep) * \real{0.62}}@{}}
\toprule
名称 & Python 类型 & 含义 \\
\midrule
\endhead
\passthrough{\lstinline!value!} & \passthrough{\lstinline!int!} &
要写入或传递的新值，必须符合该参数的 Python 类型以及底层 C++
范围约束。 \\
\bottomrule
\end{longtable}

\textbf{返回值}

\begin{longtable}[]{@{}
  >{\raggedright\arraybackslash}p{(\columnwidth - 4\tabcolsep) * \real{0.18}}
  >{\raggedright\arraybackslash}p{(\columnwidth - 4\tabcolsep) * \real{0.20}}
  >{\raggedright\arraybackslash}p{(\columnwidth - 4\tabcolsep) * \real{0.62}}@{}}
\toprule
位置 & 类型 & 含义 \\
\midrule
\endhead
getter 返回值 & \passthrough{\lstinline!int!} & 返回属性当前值。 \\
\bottomrule
\end{longtable}

\textbf{对应 C++}

\begin{itemize}
\tightlist
\item
  类：\passthrough{\lstinline!unitree\_hg::msg::dds\_::LowCmd\_!}
\item
  字段类型：\passthrough{\lstinline!uint8\_t!}
\item
  声明位置：\passthrough{\lstinline!include/unitree/idl/hg/LowCmd\_.hpp:27!}
\end{itemize}

\textbf{用法}

\begin{lstlisting}[language=Python]
current_value = obj.mode_machine
obj.mode_machine = new_value
\end{lstlisting}

\hypertarget{unitree-sdk2-cpp-idl-hg-lowcmd-motor-cmd}{%
\paragraph{\texorpdfstring{\texttt{unitree\_sdk2\_cpp.idl.hg.LowCmd.motor\_cmd}}{unitree\_sdk2\_cpp.idl.hg.LowCmd.motor\_cmd}}\label{unitree-sdk2-cpp-idl-hg-lowcmd-motor-cmd}}

底层 \passthrough{\lstinline!unitree\_hg::msg::dds\_::LowCmd\_!} 的
\passthrough{\lstinline!motor\_cmd!} 字段。类型签名只说明 Python
表示；单位、数值范围、数组固定长度和枚举含义需查对应 IDL 头文件。
底层是固定长度数组，写入序列必须正好包含 35 个元素

\textbf{签名}

\begin{lstlisting}[language=Python]
@property
def motor_cmd(self) -> list[MotorCmd]

@motor_cmd.setter
def motor_cmd(self, value: Sequence[MotorCmd]) -> None
\end{lstlisting}

\textbf{可用性与安全性}

\textbf{\passthrough{\lstinline!AVAILABLE!} /
\passthrough{\lstinline!VALUE\_TYPE!}}：当前绑定源码已实现该消息值操作。

\textbf{参数}

\begin{longtable}[]{@{}
  >{\raggedright\arraybackslash}p{(\columnwidth - 4\tabcolsep) * \real{0.18}}
  >{\raggedright\arraybackslash}p{(\columnwidth - 4\tabcolsep) * \real{0.20}}
  >{\raggedright\arraybackslash}p{(\columnwidth - 4\tabcolsep) * \real{0.62}}@{}}
\toprule
名称 & Python 类型 & 含义 \\
\midrule
\endhead
\passthrough{\lstinline!value!} &
\passthrough{\lstinline!Sequence[MotorCmd]!} &
要写入或传递的新值，必须符合该参数的 Python 类型以及底层 C++
范围约束。 \\
\bottomrule
\end{longtable}

\textbf{返回值}

\begin{longtable}[]{@{}
  >{\raggedright\arraybackslash}p{(\columnwidth - 4\tabcolsep) * \real{0.18}}
  >{\raggedright\arraybackslash}p{(\columnwidth - 4\tabcolsep) * \real{0.20}}
  >{\raggedright\arraybackslash}p{(\columnwidth - 4\tabcolsep) * \real{0.62}}@{}}
\toprule
位置 & 类型 & 含义 \\
\midrule
\endhead
getter 返回值 & \passthrough{\lstinline!list[MotorCmd]!} &
返回属性当前值。数组、vector 和嵌套消息采用复制语义。 \\
\bottomrule
\end{longtable}

\textbf{对应 C++}

\begin{itemize}
\tightlist
\item
  类：\passthrough{\lstinline!unitree\_hg::msg::dds\_::LowCmd\_!}
\item
  字段类型：\passthrough{\lstinline!std::array< ::unitree\_hg::msg::dds\_::MotorCmd\_, 35>!}
\item
  声明位置：\passthrough{\lstinline!include/unitree/idl/hg/LowCmd\_.hpp:28!}
\end{itemize}

\textbf{用法}

\begin{lstlisting}[language=Python]
current_value = obj.motor_cmd
obj.motor_cmd = new_value
\end{lstlisting}

\begin{quote}
{[}!TIP{]} 读取容器或嵌套值后应遵循``读取、修改、重新赋值''。只修改
getter 返回的副本不会可靠地更新原 C++ 消息。
\end{quote}

\hypertarget{unitree-sdk2-cpp-idl-hg-lowcmd-reserve}{%
\paragraph{\texorpdfstring{\texttt{unitree\_sdk2\_cpp.idl.hg.LowCmd.reserve}}{unitree\_sdk2\_cpp.idl.hg.LowCmd.reserve}}\label{unitree-sdk2-cpp-idl-hg-lowcmd-reserve}}

底层 \passthrough{\lstinline!unitree\_hg::msg::dds\_::LowCmd\_!} 的
\passthrough{\lstinline!reserve!} 字段。类型签名只说明 Python
表示；单位、数值范围、数组固定长度和枚举含义需查对应 IDL 头文件。
底层是固定长度数组，写入序列必须正好包含 4 个元素；元素约束：取值范围为
0 到 4294967295。

\textbf{签名}

\begin{lstlisting}[language=Python]
@property
def reserve(self) -> list[int]

@reserve.setter
def reserve(self, value: Sequence[int]) -> None
\end{lstlisting}

\textbf{可用性与安全性}

\textbf{\passthrough{\lstinline!AVAILABLE!} /
\passthrough{\lstinline!VALUE\_TYPE!}}：当前绑定源码已实现该消息值操作。

\textbf{参数}

\begin{longtable}[]{@{}
  >{\raggedright\arraybackslash}p{(\columnwidth - 4\tabcolsep) * \real{0.18}}
  >{\raggedright\arraybackslash}p{(\columnwidth - 4\tabcolsep) * \real{0.20}}
  >{\raggedright\arraybackslash}p{(\columnwidth - 4\tabcolsep) * \real{0.62}}@{}}
\toprule
名称 & Python 类型 & 含义 \\
\midrule
\endhead
\passthrough{\lstinline!value!} &
\passthrough{\lstinline!Sequence[int]!} &
要写入或传递的新值，必须符合该参数的 Python 类型以及底层 C++
范围约束。 \\
\bottomrule
\end{longtable}

\textbf{返回值}

\begin{longtable}[]{@{}
  >{\raggedright\arraybackslash}p{(\columnwidth - 4\tabcolsep) * \real{0.18}}
  >{\raggedright\arraybackslash}p{(\columnwidth - 4\tabcolsep) * \real{0.20}}
  >{\raggedright\arraybackslash}p{(\columnwidth - 4\tabcolsep) * \real{0.62}}@{}}
\toprule
位置 & 类型 & 含义 \\
\midrule
\endhead
getter 返回值 & \passthrough{\lstinline!list[int]!} &
返回属性当前值。数组、vector 和嵌套消息采用复制语义。 \\
\bottomrule
\end{longtable}

\textbf{对应 C++}

\begin{itemize}
\tightlist
\item
  类：\passthrough{\lstinline!unitree\_hg::msg::dds\_::LowCmd\_!}
\item
  字段类型：\passthrough{\lstinline!std::array<uint32\_t, 4>!}
\item
  声明位置：\passthrough{\lstinline!include/unitree/idl/hg/LowCmd\_.hpp:29!}
\end{itemize}

\textbf{用法}

\begin{lstlisting}[language=Python]
current_value = obj.reserve
obj.reserve = new_value
\end{lstlisting}

\begin{quote}
{[}!TIP{]} 读取容器或嵌套值后应遵循``读取、修改、重新赋值''。只修改
getter 返回的副本不会可靠地更新原 C++ 消息。
\end{quote}

\hypertarget{unitree-sdk2-cpp-idl-hg-lowcmd-crc}{%
\paragraph{\texorpdfstring{\texttt{unitree\_sdk2\_cpp.idl.hg.LowCmd.crc}}{unitree\_sdk2\_cpp.idl.hg.LowCmd.crc}}\label{unitree-sdk2-cpp-idl-hg-lowcmd-crc}}

底层 \passthrough{\lstinline!unitree\_hg::msg::dds\_::LowCmd\_!} 的
\passthrough{\lstinline!crc!} 字段。类型签名只说明 Python
表示；单位、数值范围、数组固定长度和枚举含义需查对应 IDL 头文件。
取值范围为 0 到 4294967295。

\textbf{签名}

\begin{lstlisting}[language=Python]
@property
def crc(self) -> int

@crc.setter
def crc(self, value: int) -> None
\end{lstlisting}

\textbf{可用性与安全性}

\textbf{\passthrough{\lstinline!AVAILABLE!} /
\passthrough{\lstinline!VALUE\_TYPE!}}：当前绑定源码已实现该消息值操作。

\textbf{参数}

\begin{longtable}[]{@{}
  >{\raggedright\arraybackslash}p{(\columnwidth - 4\tabcolsep) * \real{0.18}}
  >{\raggedright\arraybackslash}p{(\columnwidth - 4\tabcolsep) * \real{0.20}}
  >{\raggedright\arraybackslash}p{(\columnwidth - 4\tabcolsep) * \real{0.62}}@{}}
\toprule
名称 & Python 类型 & 含义 \\
\midrule
\endhead
\passthrough{\lstinline!value!} & \passthrough{\lstinline!int!} &
要写入或传递的新值，必须符合该参数的 Python 类型以及底层 C++
范围约束。 \\
\bottomrule
\end{longtable}

\textbf{返回值}

\begin{longtable}[]{@{}
  >{\raggedright\arraybackslash}p{(\columnwidth - 4\tabcolsep) * \real{0.18}}
  >{\raggedright\arraybackslash}p{(\columnwidth - 4\tabcolsep) * \real{0.20}}
  >{\raggedright\arraybackslash}p{(\columnwidth - 4\tabcolsep) * \real{0.62}}@{}}
\toprule
位置 & 类型 & 含义 \\
\midrule
\endhead
getter 返回值 & \passthrough{\lstinline!int!} & 返回属性当前值。 \\
\bottomrule
\end{longtable}

\textbf{对应 C++}

\begin{itemize}
\tightlist
\item
  类：\passthrough{\lstinline!unitree\_hg::msg::dds\_::LowCmd\_!}
\item
  字段类型：\passthrough{\lstinline!uint32\_t!}
\item
  声明位置：\passthrough{\lstinline!include/unitree/idl/hg/LowCmd\_.hpp:30!}
\end{itemize}

\textbf{用法}

\begin{lstlisting}[language=Python]
current_value = obj.crc
obj.crc = new_value
\end{lstlisting}

\hypertarget{unitree-sdk2-cpp-idl-hg-lowstate}{%
\subsubsection{\texorpdfstring{\texttt{unitree\_sdk2\_cpp.idl.hg.LowState}}{unitree\_sdk2\_cpp.idl.hg.LowState}}\label{unitree-sdk2-cpp-idl-hg-lowstate}}

可由 Python 读写的 DDS/IDL 消息值类型。字段采用复制语义。

\textbf{导入}

\begin{lstlisting}[language=Python]
from unitree_sdk2_cpp.idl.hg import LowState
\end{lstlisting}

\textbf{方法索引}

\begin{longtable}[]{@{}
  >{\raggedright\arraybackslash}p{(\columnwidth - 6\tabcolsep) * \real{0.18}}
  >{\raggedright\arraybackslash}p{(\columnwidth - 6\tabcolsep) * \real{0.24}}
  >{\raggedright\arraybackslash}p{(\columnwidth - 6\tabcolsep) * \real{0.18}}
  >{\raggedright\arraybackslash}p{(\columnwidth - 6\tabcolsep) * \real{0.40}}@{}}
\toprule
方法 & Python 签名 & 状态 & 安全分类 \\
\midrule
\endhead
\protect\hyperlink{unitree-sdk2-cpp-idl-hg-lowstate-dunder-init-1}{\passthrough{\lstinline!\_\_init\_\_!}}
& \passthrough{\lstinline!def \_\_init\_\_(self) -> None!} &
\passthrough{\lstinline!AVAILABLE!} &
\passthrough{\lstinline!VALUE\_TYPE!} \\
\protect\hyperlink{unitree-sdk2-cpp-idl-hg-lowstate-dunder-eq-1}{\passthrough{\lstinline!\_\_eq\_\_!}}
& \passthrough{\lstinline!def \_\_eq\_\_(self, other: object) -> bool!}
& \passthrough{\lstinline!AVAILABLE!} &
\passthrough{\lstinline!VALUE\_TYPE!} \\
\protect\hyperlink{unitree-sdk2-cpp-idl-hg-lowstate-dunder-ne-1}{\passthrough{\lstinline!\_\_ne\_\_!}}
& \passthrough{\lstinline!def \_\_ne\_\_(self, other: object) -> bool!}
& \passthrough{\lstinline!AVAILABLE!} &
\passthrough{\lstinline!VALUE\_TYPE!} \\
\bottomrule
\end{longtable}

\textbf{属性索引}

\begin{longtable}[]{@{}
  >{\raggedright\arraybackslash}p{(\columnwidth - 6\tabcolsep) * \real{0.18}}
  >{\raggedright\arraybackslash}p{(\columnwidth - 6\tabcolsep) * \real{0.24}}
  >{\raggedright\arraybackslash}p{(\columnwidth - 6\tabcolsep) * \real{0.18}}
  >{\raggedright\arraybackslash}p{(\columnwidth - 6\tabcolsep) * \real{0.40}}@{}}
\toprule
属性 & 读取类型 & 写入类型 & 状态 \\
\midrule
\endhead
\protect\hyperlink{unitree-sdk2-cpp-idl-hg-lowstate-version}{\passthrough{\lstinline!version!}}
& \passthrough{\lstinline!list[int]!} &
\passthrough{\lstinline!Sequence[int]!} &
\passthrough{\lstinline!AVAILABLE!} \\
\protect\hyperlink{unitree-sdk2-cpp-idl-hg-lowstate-mode-pr}{\passthrough{\lstinline!mode\_pr!}}
& \passthrough{\lstinline!int!} & \passthrough{\lstinline!int!} &
\passthrough{\lstinline!AVAILABLE!} \\
\protect\hyperlink{unitree-sdk2-cpp-idl-hg-lowstate-mode-machine}{\passthrough{\lstinline!mode\_machine!}}
& \passthrough{\lstinline!int!} & \passthrough{\lstinline!int!} &
\passthrough{\lstinline!AVAILABLE!} \\
\protect\hyperlink{unitree-sdk2-cpp-idl-hg-lowstate-tick}{\passthrough{\lstinline!tick!}}
& \passthrough{\lstinline!int!} & \passthrough{\lstinline!int!} &
\passthrough{\lstinline!AVAILABLE!} \\
\protect\hyperlink{unitree-sdk2-cpp-idl-hg-lowstate-imu-state}{\passthrough{\lstinline!imu\_state!}}
& \passthrough{\lstinline!IMUState!} &
\passthrough{\lstinline!IMUState!} &
\passthrough{\lstinline!AVAILABLE!} \\
\protect\hyperlink{unitree-sdk2-cpp-idl-hg-lowstate-motor-state}{\passthrough{\lstinline!motor\_state!}}
& \passthrough{\lstinline!list[MotorState]!} &
\passthrough{\lstinline!Sequence[MotorState]!} &
\passthrough{\lstinline!AVAILABLE!} \\
\protect\hyperlink{unitree-sdk2-cpp-idl-hg-lowstate-wireless-remote}{\passthrough{\lstinline!wireless\_remote!}}
& \passthrough{\lstinline!list[int]!} &
\passthrough{\lstinline!Sequence[int]!} &
\passthrough{\lstinline!AVAILABLE!} \\
\protect\hyperlink{unitree-sdk2-cpp-idl-hg-lowstate-reserve}{\passthrough{\lstinline!reserve!}}
& \passthrough{\lstinline!list[int]!} &
\passthrough{\lstinline!Sequence[int]!} &
\passthrough{\lstinline!AVAILABLE!} \\
\protect\hyperlink{unitree-sdk2-cpp-idl-hg-lowstate-crc}{\passthrough{\lstinline!crc!}}
& \passthrough{\lstinline!int!} & \passthrough{\lstinline!int!} &
\passthrough{\lstinline!AVAILABLE!} \\
\bottomrule
\end{longtable}

\hypertarget{unitree-sdk2-cpp-idl-hg-lowstate-dunder-init-1}{%
\paragraph{\texorpdfstring{\texttt{unitree\_sdk2\_cpp.idl.hg.LowState.\_\_init\_\_}}{unitree\_sdk2\_cpp.idl.hg.LowState.\_\_init\_\_}}\label{unitree-sdk2-cpp-idl-hg-lowstate-dunder-init-1}}

初始化 \passthrough{\lstinline!LowState!}
实例。是否能实际构造取决于下方可用性状态。

\textbf{签名}

\begin{lstlisting}[language=Python]
def __init__(self) -> None
\end{lstlisting}

\textbf{可用性与安全性}

\textbf{\passthrough{\lstinline!AVAILABLE!} /
\passthrough{\lstinline!VALUE\_TYPE!}}：当前绑定源码已实现该消息值操作。

\textbf{参数}

无显式参数。实例方法中的 \passthrough{\lstinline!self!} 由 Python
自动传入。

\textbf{返回值}

\begin{longtable}[]{@{}
  >{\raggedright\arraybackslash}p{(\columnwidth - 4\tabcolsep) * \real{0.18}}
  >{\raggedright\arraybackslash}p{(\columnwidth - 4\tabcolsep) * \real{0.20}}
  >{\raggedright\arraybackslash}p{(\columnwidth - 4\tabcolsep) * \real{0.62}}@{}}
\toprule
位置 & 类型 & 含义 \\
\midrule
\endhead
无 & \passthrough{\lstinline!None!} &
\passthrough{\lstinline!\_\_init\_\_!}
本身不返回值；调用类对象时，在构造函数可用的前提下得到该类实例。 \\
\bottomrule
\end{longtable}

\textbf{对应 C++}

\begin{itemize}
\tightlist
\item
  类：\passthrough{\lstinline!unitree\_hg::msg::dds\_::LowState\_!}
\item
  签名：\passthrough{\lstinline!LowState\_()!}
\item
  绑定策略：\passthrough{\lstinline!未单独标注!}
\item
  声明位置：\passthrough{\lstinline!include/unitree/idl/hg/LowState\_.hpp:39!}
\end{itemize}

\textbf{说明}

\begin{itemize}
\tightlist
\item
  示例是局部调用片段；\passthrough{\lstinline!obj!}
  和参数变量表示已经按上层业务校验并准备好的对象。
\end{itemize}

\textbf{用法}

\begin{lstlisting}[language=Python]
value = LowState()
\end{lstlisting}

\hypertarget{unitree-sdk2-cpp-idl-hg-lowstate-dunder-eq-1}{%
\paragraph{\texorpdfstring{\texttt{unitree\_sdk2\_cpp.idl.hg.LowState.\_\_eq\_\_}}{unitree\_sdk2\_cpp.idl.hg.LowState.\_\_eq\_\_}}\label{unitree-sdk2-cpp-idl-hg-lowstate-dunder-eq-1}}

按底层 IDL 消息内容比较两个对象是否相等。

\textbf{签名}

\begin{lstlisting}[language=Python]
def __eq__(self, other: object) -> bool
\end{lstlisting}

\textbf{可用性与安全性}

\textbf{\passthrough{\lstinline!AVAILABLE!} /
\passthrough{\lstinline!VALUE\_TYPE!}}：当前绑定源码已实现该消息值操作。

\textbf{参数}

\begin{longtable}[]{@{}
  >{\raggedright\arraybackslash}p{(\columnwidth - 6\tabcolsep) * \real{0.18}}
  >{\raggedright\arraybackslash}p{(\columnwidth - 6\tabcolsep) * \real{0.24}}
  >{\raggedright\arraybackslash}p{(\columnwidth - 6\tabcolsep) * \real{0.18}}
  >{\raggedright\arraybackslash}p{(\columnwidth - 6\tabcolsep) * \real{0.40}}@{}}
\toprule
名称 & Python 类型 & 默认值 & 含义 \\
\midrule
\endhead
\passthrough{\lstinline!other!} & \passthrough{\lstinline!object!} &
必填 & 用于比较的另一个 Python 对象。类型不兼容时比较结果通常为
\passthrough{\lstinline!False!}。 \\
\bottomrule
\end{longtable}

\textbf{返回值}

\begin{longtable}[]{@{}
  >{\raggedright\arraybackslash}p{(\columnwidth - 4\tabcolsep) * \real{0.18}}
  >{\raggedright\arraybackslash}p{(\columnwidth - 4\tabcolsep) * \real{0.20}}
  >{\raggedright\arraybackslash}p{(\columnwidth - 4\tabcolsep) * \real{0.62}}@{}}
\toprule
位置 & 类型 & 含义 \\
\midrule
\endhead
返回值 & \passthrough{\lstinline!bool!} & 返回
\passthrough{\lstinline!bool!}。更精确的业务含义见该方法用途、C++
签名和目标型号协议。 \\
\bottomrule
\end{longtable}

\textbf{对应 C++}

\begin{itemize}
\tightlist
\item
  类：\passthrough{\lstinline!unitree\_hg::msg::dds\_::LowState\_!}
\item
  签名：\passthrough{\lstinline!operator==(const value \&) const!}
\item
  绑定策略：\passthrough{\lstinline!未单独标注!}
\item
  声明位置：由 pybind11 手工绑定或生成注册表提供；manifest
  未记录独立头文件行号。
\end{itemize}

\textbf{说明}

\begin{itemize}
\tightlist
\item
  示例是局部调用片段；\passthrough{\lstinline!obj!}
  和参数变量表示已经按上层业务校验并准备好的对象。
\end{itemize}

\textbf{用法}

\begin{lstlisting}[language=Python]
same = left == right
\end{lstlisting}

\hypertarget{unitree-sdk2-cpp-idl-hg-lowstate-dunder-ne-1}{%
\paragraph{\texorpdfstring{\texttt{unitree\_sdk2\_cpp.idl.hg.LowState.\_\_ne\_\_}}{unitree\_sdk2\_cpp.idl.hg.LowState.\_\_ne\_\_}}\label{unitree-sdk2-cpp-idl-hg-lowstate-dunder-ne-1}}

按底层 IDL 消息内容比较两个对象是否不相等。

\textbf{签名}

\begin{lstlisting}[language=Python]
def __ne__(self, other: object) -> bool
\end{lstlisting}

\textbf{可用性与安全性}

\textbf{\passthrough{\lstinline!AVAILABLE!} /
\passthrough{\lstinline!VALUE\_TYPE!}}：当前绑定源码已实现该消息值操作。

\textbf{参数}

\begin{longtable}[]{@{}
  >{\raggedright\arraybackslash}p{(\columnwidth - 6\tabcolsep) * \real{0.18}}
  >{\raggedright\arraybackslash}p{(\columnwidth - 6\tabcolsep) * \real{0.24}}
  >{\raggedright\arraybackslash}p{(\columnwidth - 6\tabcolsep) * \real{0.18}}
  >{\raggedright\arraybackslash}p{(\columnwidth - 6\tabcolsep) * \real{0.40}}@{}}
\toprule
名称 & Python 类型 & 默认值 & 含义 \\
\midrule
\endhead
\passthrough{\lstinline!other!} & \passthrough{\lstinline!object!} &
必填 & 用于比较的另一个 Python 对象。类型不兼容时比较结果通常为
\passthrough{\lstinline!False!}。 \\
\bottomrule
\end{longtable}

\textbf{返回值}

\begin{longtable}[]{@{}
  >{\raggedright\arraybackslash}p{(\columnwidth - 4\tabcolsep) * \real{0.18}}
  >{\raggedright\arraybackslash}p{(\columnwidth - 4\tabcolsep) * \real{0.20}}
  >{\raggedright\arraybackslash}p{(\columnwidth - 4\tabcolsep) * \real{0.62}}@{}}
\toprule
位置 & 类型 & 含义 \\
\midrule
\endhead
返回值 & \passthrough{\lstinline!bool!} & 返回
\passthrough{\lstinline!bool!}。更精确的业务含义见该方法用途、C++
签名和目标型号协议。 \\
\bottomrule
\end{longtable}

\textbf{对应 C++}

\begin{itemize}
\tightlist
\item
  类：\passthrough{\lstinline!unitree\_hg::msg::dds\_::LowState\_!}
\item
  签名：\passthrough{\lstinline"operator!=(const value \&) const"}
\item
  绑定策略：\passthrough{\lstinline!未单独标注!}
\item
  声明位置：由 pybind11 手工绑定或生成注册表提供；manifest
  未记录独立头文件行号。
\end{itemize}

\textbf{说明}

\begin{itemize}
\tightlist
\item
  示例是局部调用片段；\passthrough{\lstinline!obj!}
  和参数变量表示已经按上层业务校验并准备好的对象。
\end{itemize}

\textbf{用法}

\begin{lstlisting}[language=Python]
different = left != right
\end{lstlisting}

\hypertarget{unitree-sdk2-cpp-idl-hg-lowstate-version}{%
\paragraph{\texorpdfstring{\texttt{unitree\_sdk2\_cpp.idl.hg.LowState.version}}{unitree\_sdk2\_cpp.idl.hg.LowState.version}}\label{unitree-sdk2-cpp-idl-hg-lowstate-version}}

底层 \passthrough{\lstinline!unitree\_hg::msg::dds\_::LowState\_!} 的
\passthrough{\lstinline!version!} 字段。类型签名只说明 Python
表示；单位、数值范围、数组固定长度和枚举含义需查对应 IDL 头文件。
底层是固定长度数组，写入序列必须正好包含 2 个元素；元素约束：取值范围为
0 到 4294967295。

\textbf{签名}

\begin{lstlisting}[language=Python]
@property
def version(self) -> list[int]

@version.setter
def version(self, value: Sequence[int]) -> None
\end{lstlisting}

\textbf{可用性与安全性}

\textbf{\passthrough{\lstinline!AVAILABLE!} /
\passthrough{\lstinline!VALUE\_TYPE!}}：当前绑定源码已实现该消息值操作。

\textbf{参数}

\begin{longtable}[]{@{}
  >{\raggedright\arraybackslash}p{(\columnwidth - 4\tabcolsep) * \real{0.18}}
  >{\raggedright\arraybackslash}p{(\columnwidth - 4\tabcolsep) * \real{0.20}}
  >{\raggedright\arraybackslash}p{(\columnwidth - 4\tabcolsep) * \real{0.62}}@{}}
\toprule
名称 & Python 类型 & 含义 \\
\midrule
\endhead
\passthrough{\lstinline!value!} &
\passthrough{\lstinline!Sequence[int]!} &
要写入或传递的新值，必须符合该参数的 Python 类型以及底层 C++
范围约束。 \\
\bottomrule
\end{longtable}

\textbf{返回值}

\begin{longtable}[]{@{}
  >{\raggedright\arraybackslash}p{(\columnwidth - 4\tabcolsep) * \real{0.18}}
  >{\raggedright\arraybackslash}p{(\columnwidth - 4\tabcolsep) * \real{0.20}}
  >{\raggedright\arraybackslash}p{(\columnwidth - 4\tabcolsep) * \real{0.62}}@{}}
\toprule
位置 & 类型 & 含义 \\
\midrule
\endhead
getter 返回值 & \passthrough{\lstinline!list[int]!} &
返回属性当前值。数组、vector 和嵌套消息采用复制语义。 \\
\bottomrule
\end{longtable}

\textbf{对应 C++}

\begin{itemize}
\tightlist
\item
  类：\passthrough{\lstinline!unitree\_hg::msg::dds\_::LowState\_!}
\item
  字段类型：\passthrough{\lstinline!std::array<uint32\_t, 2>!}
\item
  声明位置：\passthrough{\lstinline!include/unitree/idl/hg/LowState\_.hpp:28!}
\end{itemize}

\textbf{用法}

\begin{lstlisting}[language=Python]
current_value = obj.version
obj.version = new_value
\end{lstlisting}

\begin{quote}
{[}!TIP{]} 读取容器或嵌套值后应遵循``读取、修改、重新赋值''。只修改
getter 返回的副本不会可靠地更新原 C++ 消息。
\end{quote}

\hypertarget{unitree-sdk2-cpp-idl-hg-lowstate-mode-pr}{%
\paragraph{\texorpdfstring{\texttt{unitree\_sdk2\_cpp.idl.hg.LowState.mode\_pr}}{unitree\_sdk2\_cpp.idl.hg.LowState.mode\_pr}}\label{unitree-sdk2-cpp-idl-hg-lowstate-mode-pr}}

底层 \passthrough{\lstinline!unitree\_hg::msg::dds\_::LowState\_!} 的
\passthrough{\lstinline!mode\_pr!} 字段。类型签名只说明 Python
表示；单位、数值范围、数组固定长度和枚举含义需查对应 IDL 头文件。
取值范围为 0 到 255。

\textbf{签名}

\begin{lstlisting}[language=Python]
@property
def mode_pr(self) -> int

@mode_pr.setter
def mode_pr(self, value: int) -> None
\end{lstlisting}

\textbf{可用性与安全性}

\textbf{\passthrough{\lstinline!AVAILABLE!} /
\passthrough{\lstinline!VALUE\_TYPE!}}：当前绑定源码已实现该消息值操作。

\textbf{参数}

\begin{longtable}[]{@{}
  >{\raggedright\arraybackslash}p{(\columnwidth - 4\tabcolsep) * \real{0.18}}
  >{\raggedright\arraybackslash}p{(\columnwidth - 4\tabcolsep) * \real{0.20}}
  >{\raggedright\arraybackslash}p{(\columnwidth - 4\tabcolsep) * \real{0.62}}@{}}
\toprule
名称 & Python 类型 & 含义 \\
\midrule
\endhead
\passthrough{\lstinline!value!} & \passthrough{\lstinline!int!} &
要写入或传递的新值，必须符合该参数的 Python 类型以及底层 C++
范围约束。 \\
\bottomrule
\end{longtable}

\textbf{返回值}

\begin{longtable}[]{@{}
  >{\raggedright\arraybackslash}p{(\columnwidth - 4\tabcolsep) * \real{0.18}}
  >{\raggedright\arraybackslash}p{(\columnwidth - 4\tabcolsep) * \real{0.20}}
  >{\raggedright\arraybackslash}p{(\columnwidth - 4\tabcolsep) * \real{0.62}}@{}}
\toprule
位置 & 类型 & 含义 \\
\midrule
\endhead
getter 返回值 & \passthrough{\lstinline!int!} & 返回属性当前值。 \\
\bottomrule
\end{longtable}

\textbf{对应 C++}

\begin{itemize}
\tightlist
\item
  类：\passthrough{\lstinline!unitree\_hg::msg::dds\_::LowState\_!}
\item
  字段类型：\passthrough{\lstinline!uint8\_t!}
\item
  声明位置：\passthrough{\lstinline!include/unitree/idl/hg/LowState\_.hpp:29!}
\end{itemize}

\textbf{用法}

\begin{lstlisting}[language=Python]
current_value = obj.mode_pr
obj.mode_pr = new_value
\end{lstlisting}

\hypertarget{unitree-sdk2-cpp-idl-hg-lowstate-mode-machine}{%
\paragraph{\texorpdfstring{\texttt{unitree\_sdk2\_cpp.idl.hg.LowState.mode\_machine}}{unitree\_sdk2\_cpp.idl.hg.LowState.mode\_machine}}\label{unitree-sdk2-cpp-idl-hg-lowstate-mode-machine}}

底层 \passthrough{\lstinline!unitree\_hg::msg::dds\_::LowState\_!} 的
\passthrough{\lstinline!mode\_machine!} 字段。类型签名只说明 Python
表示；单位、数值范围、数组固定长度和枚举含义需查对应 IDL 头文件。
取值范围为 0 到 255。

\textbf{签名}

\begin{lstlisting}[language=Python]
@property
def mode_machine(self) -> int

@mode_machine.setter
def mode_machine(self, value: int) -> None
\end{lstlisting}

\textbf{可用性与安全性}

\textbf{\passthrough{\lstinline!AVAILABLE!} /
\passthrough{\lstinline!VALUE\_TYPE!}}：当前绑定源码已实现该消息值操作。

\textbf{参数}

\begin{longtable}[]{@{}
  >{\raggedright\arraybackslash}p{(\columnwidth - 4\tabcolsep) * \real{0.18}}
  >{\raggedright\arraybackslash}p{(\columnwidth - 4\tabcolsep) * \real{0.20}}
  >{\raggedright\arraybackslash}p{(\columnwidth - 4\tabcolsep) * \real{0.62}}@{}}
\toprule
名称 & Python 类型 & 含义 \\
\midrule
\endhead
\passthrough{\lstinline!value!} & \passthrough{\lstinline!int!} &
要写入或传递的新值，必须符合该参数的 Python 类型以及底层 C++
范围约束。 \\
\bottomrule
\end{longtable}

\textbf{返回值}

\begin{longtable}[]{@{}
  >{\raggedright\arraybackslash}p{(\columnwidth - 4\tabcolsep) * \real{0.18}}
  >{\raggedright\arraybackslash}p{(\columnwidth - 4\tabcolsep) * \real{0.20}}
  >{\raggedright\arraybackslash}p{(\columnwidth - 4\tabcolsep) * \real{0.62}}@{}}
\toprule
位置 & 类型 & 含义 \\
\midrule
\endhead
getter 返回值 & \passthrough{\lstinline!int!} & 返回属性当前值。 \\
\bottomrule
\end{longtable}

\textbf{对应 C++}

\begin{itemize}
\tightlist
\item
  类：\passthrough{\lstinline!unitree\_hg::msg::dds\_::LowState\_!}
\item
  字段类型：\passthrough{\lstinline!uint8\_t!}
\item
  声明位置：\passthrough{\lstinline!include/unitree/idl/hg/LowState\_.hpp:30!}
\end{itemize}

\textbf{用法}

\begin{lstlisting}[language=Python]
current_value = obj.mode_machine
obj.mode_machine = new_value
\end{lstlisting}

\hypertarget{unitree-sdk2-cpp-idl-hg-lowstate-tick}{%
\paragraph{\texorpdfstring{\texttt{unitree\_sdk2\_cpp.idl.hg.LowState.tick}}{unitree\_sdk2\_cpp.idl.hg.LowState.tick}}\label{unitree-sdk2-cpp-idl-hg-lowstate-tick}}

底层 \passthrough{\lstinline!unitree\_hg::msg::dds\_::LowState\_!} 的
\passthrough{\lstinline!tick!} 字段。类型签名只说明 Python
表示；单位、数值范围、数组固定长度和枚举含义需查对应 IDL 头文件。
取值范围为 0 到 4294967295。

\textbf{签名}

\begin{lstlisting}[language=Python]
@property
def tick(self) -> int

@tick.setter
def tick(self, value: int) -> None
\end{lstlisting}

\textbf{可用性与安全性}

\textbf{\passthrough{\lstinline!AVAILABLE!} /
\passthrough{\lstinline!VALUE\_TYPE!}}：当前绑定源码已实现该消息值操作。

\textbf{参数}

\begin{longtable}[]{@{}
  >{\raggedright\arraybackslash}p{(\columnwidth - 4\tabcolsep) * \real{0.18}}
  >{\raggedright\arraybackslash}p{(\columnwidth - 4\tabcolsep) * \real{0.20}}
  >{\raggedright\arraybackslash}p{(\columnwidth - 4\tabcolsep) * \real{0.62}}@{}}
\toprule
名称 & Python 类型 & 含义 \\
\midrule
\endhead
\passthrough{\lstinline!value!} & \passthrough{\lstinline!int!} &
要写入或传递的新值，必须符合该参数的 Python 类型以及底层 C++
范围约束。 \\
\bottomrule
\end{longtable}

\textbf{返回值}

\begin{longtable}[]{@{}
  >{\raggedright\arraybackslash}p{(\columnwidth - 4\tabcolsep) * \real{0.18}}
  >{\raggedright\arraybackslash}p{(\columnwidth - 4\tabcolsep) * \real{0.20}}
  >{\raggedright\arraybackslash}p{(\columnwidth - 4\tabcolsep) * \real{0.62}}@{}}
\toprule
位置 & 类型 & 含义 \\
\midrule
\endhead
getter 返回值 & \passthrough{\lstinline!int!} & 返回属性当前值。 \\
\bottomrule
\end{longtable}

\textbf{对应 C++}

\begin{itemize}
\tightlist
\item
  类：\passthrough{\lstinline!unitree\_hg::msg::dds\_::LowState\_!}
\item
  字段类型：\passthrough{\lstinline!uint32\_t!}
\item
  声明位置：\passthrough{\lstinline!include/unitree/idl/hg/LowState\_.hpp:31!}
\end{itemize}

\textbf{用法}

\begin{lstlisting}[language=Python]
current_value = obj.tick
obj.tick = new_value
\end{lstlisting}

\hypertarget{unitree-sdk2-cpp-idl-hg-lowstate-imu-state}{%
\paragraph{\texorpdfstring{\texttt{unitree\_sdk2\_cpp.idl.hg.LowState.imu\_state}}{unitree\_sdk2\_cpp.idl.hg.LowState.imu\_state}}\label{unitree-sdk2-cpp-idl-hg-lowstate-imu-state}}

底层 \passthrough{\lstinline!unitree\_hg::msg::dds\_::LowState\_!} 的
\passthrough{\lstinline!imu\_state!} 字段。类型签名只说明 Python
表示；单位、数值范围、数组固定长度和枚举含义需查对应 IDL 头文件。

\textbf{签名}

\begin{lstlisting}[language=Python]
@property
def imu_state(self) -> IMUState

@imu_state.setter
def imu_state(self, value: IMUState) -> None
\end{lstlisting}

\textbf{可用性与安全性}

\textbf{\passthrough{\lstinline!AVAILABLE!} /
\passthrough{\lstinline!VALUE\_TYPE!}}：当前绑定源码已实现该消息值操作。

\textbf{参数}

\begin{longtable}[]{@{}
  >{\raggedright\arraybackslash}p{(\columnwidth - 4\tabcolsep) * \real{0.18}}
  >{\raggedright\arraybackslash}p{(\columnwidth - 4\tabcolsep) * \real{0.20}}
  >{\raggedright\arraybackslash}p{(\columnwidth - 4\tabcolsep) * \real{0.62}}@{}}
\toprule
名称 & Python 类型 & 含义 \\
\midrule
\endhead
\passthrough{\lstinline!value!} & \passthrough{\lstinline!IMUState!} &
要写入或传递的新值，必须符合该参数的 Python 类型以及底层 C++
范围约束。 \\
\bottomrule
\end{longtable}

\textbf{返回值}

\begin{longtable}[]{@{}
  >{\raggedright\arraybackslash}p{(\columnwidth - 4\tabcolsep) * \real{0.18}}
  >{\raggedright\arraybackslash}p{(\columnwidth - 4\tabcolsep) * \real{0.20}}
  >{\raggedright\arraybackslash}p{(\columnwidth - 4\tabcolsep) * \real{0.62}}@{}}
\toprule
位置 & 类型 & 含义 \\
\midrule
\endhead
getter 返回值 & \passthrough{\lstinline!IMUState!} &
返回属性当前值。数组、vector 和嵌套消息采用复制语义。 \\
\bottomrule
\end{longtable}

\textbf{对应 C++}

\begin{itemize}
\tightlist
\item
  类：\passthrough{\lstinline!unitree\_hg::msg::dds\_::LowState\_!}
\item
  字段类型：\passthrough{\lstinline!::unitree\_hg::msg::dds\_::IMUState\_!}
\item
  声明位置：\passthrough{\lstinline!include/unitree/idl/hg/LowState\_.hpp:32!}
\end{itemize}

\textbf{用法}

\begin{lstlisting}[language=Python]
current_value = obj.imu_state
obj.imu_state = new_value
\end{lstlisting}

\hypertarget{unitree-sdk2-cpp-idl-hg-lowstate-motor-state}{%
\paragraph{\texorpdfstring{\texttt{unitree\_sdk2\_cpp.idl.hg.LowState.motor\_state}}{unitree\_sdk2\_cpp.idl.hg.LowState.motor\_state}}\label{unitree-sdk2-cpp-idl-hg-lowstate-motor-state}}

底层 \passthrough{\lstinline!unitree\_hg::msg::dds\_::LowState\_!} 的
\passthrough{\lstinline!motor\_state!} 字段。类型签名只说明 Python
表示；单位、数值范围、数组固定长度和枚举含义需查对应 IDL 头文件。
底层是固定长度数组，写入序列必须正好包含 35 个元素

\textbf{签名}

\begin{lstlisting}[language=Python]
@property
def motor_state(self) -> list[MotorState]

@motor_state.setter
def motor_state(self, value: Sequence[MotorState]) -> None
\end{lstlisting}

\textbf{可用性与安全性}

\textbf{\passthrough{\lstinline!AVAILABLE!} /
\passthrough{\lstinline!VALUE\_TYPE!}}：当前绑定源码已实现该消息值操作。

\textbf{参数}

\begin{longtable}[]{@{}
  >{\raggedright\arraybackslash}p{(\columnwidth - 4\tabcolsep) * \real{0.18}}
  >{\raggedright\arraybackslash}p{(\columnwidth - 4\tabcolsep) * \real{0.20}}
  >{\raggedright\arraybackslash}p{(\columnwidth - 4\tabcolsep) * \real{0.62}}@{}}
\toprule
名称 & Python 类型 & 含义 \\
\midrule
\endhead
\passthrough{\lstinline!value!} &
\passthrough{\lstinline!Sequence[MotorState]!} &
要写入或传递的新值，必须符合该参数的 Python 类型以及底层 C++
范围约束。 \\
\bottomrule
\end{longtable}

\textbf{返回值}

\begin{longtable}[]{@{}
  >{\raggedright\arraybackslash}p{(\columnwidth - 4\tabcolsep) * \real{0.18}}
  >{\raggedright\arraybackslash}p{(\columnwidth - 4\tabcolsep) * \real{0.20}}
  >{\raggedright\arraybackslash}p{(\columnwidth - 4\tabcolsep) * \real{0.62}}@{}}
\toprule
位置 & 类型 & 含义 \\
\midrule
\endhead
getter 返回值 & \passthrough{\lstinline!list[MotorState]!} &
返回属性当前值。数组、vector 和嵌套消息采用复制语义。 \\
\bottomrule
\end{longtable}

\textbf{对应 C++}

\begin{itemize}
\tightlist
\item
  类：\passthrough{\lstinline!unitree\_hg::msg::dds\_::LowState\_!}
\item
  字段类型：\passthrough{\lstinline!std::array< ::unitree\_hg::msg::dds\_::MotorState\_, 35>!}
\item
  声明位置：\passthrough{\lstinline!include/unitree/idl/hg/LowState\_.hpp:33!}
\end{itemize}

\textbf{用法}

\begin{lstlisting}[language=Python]
current_value = obj.motor_state
obj.motor_state = new_value
\end{lstlisting}

\begin{quote}
{[}!TIP{]} 读取容器或嵌套值后应遵循``读取、修改、重新赋值''。只修改
getter 返回的副本不会可靠地更新原 C++ 消息。
\end{quote}

\hypertarget{unitree-sdk2-cpp-idl-hg-lowstate-wireless-remote}{%
\paragraph{\texorpdfstring{\texttt{unitree\_sdk2\_cpp.idl.hg.LowState.wireless\_remote}}{unitree\_sdk2\_cpp.idl.hg.LowState.wireless\_remote}}\label{unitree-sdk2-cpp-idl-hg-lowstate-wireless-remote}}

底层 \passthrough{\lstinline!unitree\_hg::msg::dds\_::LowState\_!} 的
\passthrough{\lstinline!wireless\_remote!} 字段。类型签名只说明 Python
表示；单位、数值范围、数组固定长度和枚举含义需查对应 IDL 头文件。
底层是固定长度数组，写入序列必须正好包含 40 个元素；元素约束：取值范围为
0 到 255。

\textbf{签名}

\begin{lstlisting}[language=Python]
@property
def wireless_remote(self) -> list[int]

@wireless_remote.setter
def wireless_remote(self, value: Sequence[int]) -> None
\end{lstlisting}

\textbf{可用性与安全性}

\textbf{\passthrough{\lstinline!AVAILABLE!} /
\passthrough{\lstinline!VALUE\_TYPE!}}：当前绑定源码已实现该消息值操作。

\textbf{参数}

\begin{longtable}[]{@{}
  >{\raggedright\arraybackslash}p{(\columnwidth - 4\tabcolsep) * \real{0.18}}
  >{\raggedright\arraybackslash}p{(\columnwidth - 4\tabcolsep) * \real{0.20}}
  >{\raggedright\arraybackslash}p{(\columnwidth - 4\tabcolsep) * \real{0.62}}@{}}
\toprule
名称 & Python 类型 & 含义 \\
\midrule
\endhead
\passthrough{\lstinline!value!} &
\passthrough{\lstinline!Sequence[int]!} &
要写入或传递的新值，必须符合该参数的 Python 类型以及底层 C++
范围约束。 \\
\bottomrule
\end{longtable}

\textbf{返回值}

\begin{longtable}[]{@{}
  >{\raggedright\arraybackslash}p{(\columnwidth - 4\tabcolsep) * \real{0.18}}
  >{\raggedright\arraybackslash}p{(\columnwidth - 4\tabcolsep) * \real{0.20}}
  >{\raggedright\arraybackslash}p{(\columnwidth - 4\tabcolsep) * \real{0.62}}@{}}
\toprule
位置 & 类型 & 含义 \\
\midrule
\endhead
getter 返回值 & \passthrough{\lstinline!list[int]!} &
返回属性当前值。数组、vector 和嵌套消息采用复制语义。 \\
\bottomrule
\end{longtable}

\textbf{对应 C++}

\begin{itemize}
\tightlist
\item
  类：\passthrough{\lstinline!unitree\_hg::msg::dds\_::LowState\_!}
\item
  字段类型：\passthrough{\lstinline!std::array<uint8\_t, 40>!}
\item
  声明位置：\passthrough{\lstinline!include/unitree/idl/hg/LowState\_.hpp:34!}
\end{itemize}

\textbf{用法}

\begin{lstlisting}[language=Python]
current_value = obj.wireless_remote
obj.wireless_remote = new_value
\end{lstlisting}

\begin{quote}
{[}!TIP{]} 读取容器或嵌套值后应遵循``读取、修改、重新赋值''。只修改
getter 返回的副本不会可靠地更新原 C++ 消息。
\end{quote}

\hypertarget{unitree-sdk2-cpp-idl-hg-lowstate-reserve}{%
\paragraph{\texorpdfstring{\texttt{unitree\_sdk2\_cpp.idl.hg.LowState.reserve}}{unitree\_sdk2\_cpp.idl.hg.LowState.reserve}}\label{unitree-sdk2-cpp-idl-hg-lowstate-reserve}}

底层 \passthrough{\lstinline!unitree\_hg::msg::dds\_::LowState\_!} 的
\passthrough{\lstinline!reserve!} 字段。类型签名只说明 Python
表示；单位、数值范围、数组固定长度和枚举含义需查对应 IDL 头文件。
底层是固定长度数组，写入序列必须正好包含 4 个元素；元素约束：取值范围为
0 到 4294967295。

\textbf{签名}

\begin{lstlisting}[language=Python]
@property
def reserve(self) -> list[int]

@reserve.setter
def reserve(self, value: Sequence[int]) -> None
\end{lstlisting}

\textbf{可用性与安全性}

\textbf{\passthrough{\lstinline!AVAILABLE!} /
\passthrough{\lstinline!VALUE\_TYPE!}}：当前绑定源码已实现该消息值操作。

\textbf{参数}

\begin{longtable}[]{@{}
  >{\raggedright\arraybackslash}p{(\columnwidth - 4\tabcolsep) * \real{0.18}}
  >{\raggedright\arraybackslash}p{(\columnwidth - 4\tabcolsep) * \real{0.20}}
  >{\raggedright\arraybackslash}p{(\columnwidth - 4\tabcolsep) * \real{0.62}}@{}}
\toprule
名称 & Python 类型 & 含义 \\
\midrule
\endhead
\passthrough{\lstinline!value!} &
\passthrough{\lstinline!Sequence[int]!} &
要写入或传递的新值，必须符合该参数的 Python 类型以及底层 C++
范围约束。 \\
\bottomrule
\end{longtable}

\textbf{返回值}

\begin{longtable}[]{@{}
  >{\raggedright\arraybackslash}p{(\columnwidth - 4\tabcolsep) * \real{0.18}}
  >{\raggedright\arraybackslash}p{(\columnwidth - 4\tabcolsep) * \real{0.20}}
  >{\raggedright\arraybackslash}p{(\columnwidth - 4\tabcolsep) * \real{0.62}}@{}}
\toprule
位置 & 类型 & 含义 \\
\midrule
\endhead
getter 返回值 & \passthrough{\lstinline!list[int]!} &
返回属性当前值。数组、vector 和嵌套消息采用复制语义。 \\
\bottomrule
\end{longtable}

\textbf{对应 C++}

\begin{itemize}
\tightlist
\item
  类：\passthrough{\lstinline!unitree\_hg::msg::dds\_::LowState\_!}
\item
  字段类型：\passthrough{\lstinline!std::array<uint32\_t, 4>!}
\item
  声明位置：\passthrough{\lstinline!include/unitree/idl/hg/LowState\_.hpp:35!}
\end{itemize}

\textbf{用法}

\begin{lstlisting}[language=Python]
current_value = obj.reserve
obj.reserve = new_value
\end{lstlisting}

\begin{quote}
{[}!TIP{]} 读取容器或嵌套值后应遵循``读取、修改、重新赋值''。只修改
getter 返回的副本不会可靠地更新原 C++ 消息。
\end{quote}

\hypertarget{unitree-sdk2-cpp-idl-hg-lowstate-crc}{%
\paragraph{\texorpdfstring{\texttt{unitree\_sdk2\_cpp.idl.hg.LowState.crc}}{unitree\_sdk2\_cpp.idl.hg.LowState.crc}}\label{unitree-sdk2-cpp-idl-hg-lowstate-crc}}

底层 \passthrough{\lstinline!unitree\_hg::msg::dds\_::LowState\_!} 的
\passthrough{\lstinline!crc!} 字段。类型签名只说明 Python
表示；单位、数值范围、数组固定长度和枚举含义需查对应 IDL 头文件。
取值范围为 0 到 4294967295。

\textbf{签名}

\begin{lstlisting}[language=Python]
@property
def crc(self) -> int

@crc.setter
def crc(self, value: int) -> None
\end{lstlisting}

\textbf{可用性与安全性}

\textbf{\passthrough{\lstinline!AVAILABLE!} /
\passthrough{\lstinline!VALUE\_TYPE!}}：当前绑定源码已实现该消息值操作。

\textbf{参数}

\begin{longtable}[]{@{}
  >{\raggedright\arraybackslash}p{(\columnwidth - 4\tabcolsep) * \real{0.18}}
  >{\raggedright\arraybackslash}p{(\columnwidth - 4\tabcolsep) * \real{0.20}}
  >{\raggedright\arraybackslash}p{(\columnwidth - 4\tabcolsep) * \real{0.62}}@{}}
\toprule
名称 & Python 类型 & 含义 \\
\midrule
\endhead
\passthrough{\lstinline!value!} & \passthrough{\lstinline!int!} &
要写入或传递的新值，必须符合该参数的 Python 类型以及底层 C++
范围约束。 \\
\bottomrule
\end{longtable}

\textbf{返回值}

\begin{longtable}[]{@{}
  >{\raggedright\arraybackslash}p{(\columnwidth - 4\tabcolsep) * \real{0.18}}
  >{\raggedright\arraybackslash}p{(\columnwidth - 4\tabcolsep) * \real{0.20}}
  >{\raggedright\arraybackslash}p{(\columnwidth - 4\tabcolsep) * \real{0.62}}@{}}
\toprule
位置 & 类型 & 含义 \\
\midrule
\endhead
getter 返回值 & \passthrough{\lstinline!int!} & 返回属性当前值。 \\
\bottomrule
\end{longtable}

\textbf{对应 C++}

\begin{itemize}
\tightlist
\item
  类：\passthrough{\lstinline!unitree\_hg::msg::dds\_::LowState\_!}
\item
  字段类型：\passthrough{\lstinline!uint32\_t!}
\item
  声明位置：\passthrough{\lstinline!include/unitree/idl/hg/LowState\_.hpp:36!}
\end{itemize}

\textbf{用法}

\begin{lstlisting}[language=Python]
current_value = obj.crc
obj.crc = new_value
\end{lstlisting}

\hypertarget{unitree-sdk2-cpp-idl-hg-mainboardstate}{%
\subsubsection{\texorpdfstring{\texttt{unitree\_sdk2\_cpp.idl.hg.MainBoardState}}{unitree\_sdk2\_cpp.idl.hg.MainBoardState}}\label{unitree-sdk2-cpp-idl-hg-mainboardstate}}

可由 Python 读写的 DDS/IDL 消息值类型。字段采用复制语义。

\textbf{导入}

\begin{lstlisting}[language=Python]
from unitree_sdk2_cpp.idl.hg import MainBoardState
\end{lstlisting}

\textbf{方法索引}

\begin{longtable}[]{@{}
  >{\raggedright\arraybackslash}p{(\columnwidth - 6\tabcolsep) * \real{0.18}}
  >{\raggedright\arraybackslash}p{(\columnwidth - 6\tabcolsep) * \real{0.24}}
  >{\raggedright\arraybackslash}p{(\columnwidth - 6\tabcolsep) * \real{0.18}}
  >{\raggedright\arraybackslash}p{(\columnwidth - 6\tabcolsep) * \real{0.40}}@{}}
\toprule
方法 & Python 签名 & 状态 & 安全分类 \\
\midrule
\endhead
\protect\hyperlink{unitree-sdk2-cpp-idl-hg-mainboardstate-dunder-init-1}{\passthrough{\lstinline!\_\_init\_\_!}}
& \passthrough{\lstinline!def \_\_init\_\_(self) -> None!} &
\passthrough{\lstinline!AVAILABLE!} &
\passthrough{\lstinline!VALUE\_TYPE!} \\
\protect\hyperlink{unitree-sdk2-cpp-idl-hg-mainboardstate-dunder-eq-1}{\passthrough{\lstinline!\_\_eq\_\_!}}
& \passthrough{\lstinline!def \_\_eq\_\_(self, other: object) -> bool!}
& \passthrough{\lstinline!AVAILABLE!} &
\passthrough{\lstinline!VALUE\_TYPE!} \\
\protect\hyperlink{unitree-sdk2-cpp-idl-hg-mainboardstate-dunder-ne-1}{\passthrough{\lstinline!\_\_ne\_\_!}}
& \passthrough{\lstinline!def \_\_ne\_\_(self, other: object) -> bool!}
& \passthrough{\lstinline!AVAILABLE!} &
\passthrough{\lstinline!VALUE\_TYPE!} \\
\bottomrule
\end{longtable}

\textbf{属性索引}

\begin{longtable}[]{@{}
  >{\raggedright\arraybackslash}p{(\columnwidth - 6\tabcolsep) * \real{0.18}}
  >{\raggedright\arraybackslash}p{(\columnwidth - 6\tabcolsep) * \real{0.24}}
  >{\raggedright\arraybackslash}p{(\columnwidth - 6\tabcolsep) * \real{0.18}}
  >{\raggedright\arraybackslash}p{(\columnwidth - 6\tabcolsep) * \real{0.40}}@{}}
\toprule
属性 & 读取类型 & 写入类型 & 状态 \\
\midrule
\endhead
\protect\hyperlink{unitree-sdk2-cpp-idl-hg-mainboardstate-fan-state}{\passthrough{\lstinline!fan\_state!}}
& \passthrough{\lstinline!list[int]!} &
\passthrough{\lstinline!Sequence[int]!} &
\passthrough{\lstinline!AVAILABLE!} \\
\protect\hyperlink{unitree-sdk2-cpp-idl-hg-mainboardstate-temperature}{\passthrough{\lstinline!temperature!}}
& \passthrough{\lstinline!list[int]!} &
\passthrough{\lstinline!Sequence[int]!} &
\passthrough{\lstinline!AVAILABLE!} \\
\protect\hyperlink{unitree-sdk2-cpp-idl-hg-mainboardstate-value}{\passthrough{\lstinline!value!}}
& \passthrough{\lstinline!list[float]!} &
\passthrough{\lstinline!Sequence[float]!} &
\passthrough{\lstinline!AVAILABLE!} \\
\protect\hyperlink{unitree-sdk2-cpp-idl-hg-mainboardstate-state}{\passthrough{\lstinline!state!}}
& \passthrough{\lstinline!list[int]!} &
\passthrough{\lstinline!Sequence[int]!} &
\passthrough{\lstinline!AVAILABLE!} \\
\bottomrule
\end{longtable}

\hypertarget{unitree-sdk2-cpp-idl-hg-mainboardstate-dunder-init-1}{%
\paragraph{\texorpdfstring{\texttt{unitree\_sdk2\_cpp.idl.hg.MainBoardState.\_\_init\_\_}}{unitree\_sdk2\_cpp.idl.hg.MainBoardState.\_\_init\_\_}}\label{unitree-sdk2-cpp-idl-hg-mainboardstate-dunder-init-1}}

初始化 \passthrough{\lstinline!MainBoardState!}
实例。是否能实际构造取决于下方可用性状态。

\textbf{签名}

\begin{lstlisting}[language=Python]
def __init__(self) -> None
\end{lstlisting}

\textbf{可用性与安全性}

\textbf{\passthrough{\lstinline!AVAILABLE!} /
\passthrough{\lstinline!VALUE\_TYPE!}}：当前绑定源码已实现该消息值操作。

\textbf{参数}

无显式参数。实例方法中的 \passthrough{\lstinline!self!} 由 Python
自动传入。

\textbf{返回值}

\begin{longtable}[]{@{}
  >{\raggedright\arraybackslash}p{(\columnwidth - 4\tabcolsep) * \real{0.18}}
  >{\raggedright\arraybackslash}p{(\columnwidth - 4\tabcolsep) * \real{0.20}}
  >{\raggedright\arraybackslash}p{(\columnwidth - 4\tabcolsep) * \real{0.62}}@{}}
\toprule
位置 & 类型 & 含义 \\
\midrule
\endhead
无 & \passthrough{\lstinline!None!} &
\passthrough{\lstinline!\_\_init\_\_!}
本身不返回值；调用类对象时，在构造函数可用的前提下得到该类实例。 \\
\bottomrule
\end{longtable}

\textbf{对应 C++}

\begin{itemize}
\tightlist
\item
  类：\passthrough{\lstinline!unitree\_hg::msg::dds\_::MainBoardState\_!}
\item
  签名：\passthrough{\lstinline!MainBoardState\_()!}
\item
  绑定策略：\passthrough{\lstinline!未单独标注!}
\item
  声明位置：\passthrough{\lstinline!include/unitree/idl/hg/MainBoardState\_.hpp:30!}
\end{itemize}

\textbf{说明}

\begin{itemize}
\tightlist
\item
  示例是局部调用片段；\passthrough{\lstinline!obj!}
  和参数变量表示已经按上层业务校验并准备好的对象。
\end{itemize}

\textbf{用法}

\begin{lstlisting}[language=Python]
value = MainBoardState()
\end{lstlisting}

\hypertarget{unitree-sdk2-cpp-idl-hg-mainboardstate-dunder-eq-1}{%
\paragraph{\texorpdfstring{\texttt{unitree\_sdk2\_cpp.idl.hg.MainBoardState.\_\_eq\_\_}}{unitree\_sdk2\_cpp.idl.hg.MainBoardState.\_\_eq\_\_}}\label{unitree-sdk2-cpp-idl-hg-mainboardstate-dunder-eq-1}}

按底层 IDL 消息内容比较两个对象是否相等。

\textbf{签名}

\begin{lstlisting}[language=Python]
def __eq__(self, other: object) -> bool
\end{lstlisting}

\textbf{可用性与安全性}

\textbf{\passthrough{\lstinline!AVAILABLE!} /
\passthrough{\lstinline!VALUE\_TYPE!}}：当前绑定源码已实现该消息值操作。

\textbf{参数}

\begin{longtable}[]{@{}
  >{\raggedright\arraybackslash}p{(\columnwidth - 6\tabcolsep) * \real{0.18}}
  >{\raggedright\arraybackslash}p{(\columnwidth - 6\tabcolsep) * \real{0.24}}
  >{\raggedright\arraybackslash}p{(\columnwidth - 6\tabcolsep) * \real{0.18}}
  >{\raggedright\arraybackslash}p{(\columnwidth - 6\tabcolsep) * \real{0.40}}@{}}
\toprule
名称 & Python 类型 & 默认值 & 含义 \\
\midrule
\endhead
\passthrough{\lstinline!other!} & \passthrough{\lstinline!object!} &
必填 & 用于比较的另一个 Python 对象。类型不兼容时比较结果通常为
\passthrough{\lstinline!False!}。 \\
\bottomrule
\end{longtable}

\textbf{返回值}

\begin{longtable}[]{@{}
  >{\raggedright\arraybackslash}p{(\columnwidth - 4\tabcolsep) * \real{0.18}}
  >{\raggedright\arraybackslash}p{(\columnwidth - 4\tabcolsep) * \real{0.20}}
  >{\raggedright\arraybackslash}p{(\columnwidth - 4\tabcolsep) * \real{0.62}}@{}}
\toprule
位置 & 类型 & 含义 \\
\midrule
\endhead
返回值 & \passthrough{\lstinline!bool!} & 返回
\passthrough{\lstinline!bool!}。更精确的业务含义见该方法用途、C++
签名和目标型号协议。 \\
\bottomrule
\end{longtable}

\textbf{对应 C++}

\begin{itemize}
\tightlist
\item
  类：\passthrough{\lstinline!unitree\_hg::msg::dds\_::MainBoardState\_!}
\item
  签名：\passthrough{\lstinline!operator==(const value \&) const!}
\item
  绑定策略：\passthrough{\lstinline!未单独标注!}
\item
  声明位置：由 pybind11 手工绑定或生成注册表提供；manifest
  未记录独立头文件行号。
\end{itemize}

\textbf{说明}

\begin{itemize}
\tightlist
\item
  示例是局部调用片段；\passthrough{\lstinline!obj!}
  和参数变量表示已经按上层业务校验并准备好的对象。
\end{itemize}

\textbf{用法}

\begin{lstlisting}[language=Python]
same = left == right
\end{lstlisting}

\hypertarget{unitree-sdk2-cpp-idl-hg-mainboardstate-dunder-ne-1}{%
\paragraph{\texorpdfstring{\texttt{unitree\_sdk2\_cpp.idl.hg.MainBoardState.\_\_ne\_\_}}{unitree\_sdk2\_cpp.idl.hg.MainBoardState.\_\_ne\_\_}}\label{unitree-sdk2-cpp-idl-hg-mainboardstate-dunder-ne-1}}

按底层 IDL 消息内容比较两个对象是否不相等。

\textbf{签名}

\begin{lstlisting}[language=Python]
def __ne__(self, other: object) -> bool
\end{lstlisting}

\textbf{可用性与安全性}

\textbf{\passthrough{\lstinline!AVAILABLE!} /
\passthrough{\lstinline!VALUE\_TYPE!}}：当前绑定源码已实现该消息值操作。

\textbf{参数}

\begin{longtable}[]{@{}
  >{\raggedright\arraybackslash}p{(\columnwidth - 6\tabcolsep) * \real{0.18}}
  >{\raggedright\arraybackslash}p{(\columnwidth - 6\tabcolsep) * \real{0.24}}
  >{\raggedright\arraybackslash}p{(\columnwidth - 6\tabcolsep) * \real{0.18}}
  >{\raggedright\arraybackslash}p{(\columnwidth - 6\tabcolsep) * \real{0.40}}@{}}
\toprule
名称 & Python 类型 & 默认值 & 含义 \\
\midrule
\endhead
\passthrough{\lstinline!other!} & \passthrough{\lstinline!object!} &
必填 & 用于比较的另一个 Python 对象。类型不兼容时比较结果通常为
\passthrough{\lstinline!False!}。 \\
\bottomrule
\end{longtable}

\textbf{返回值}

\begin{longtable}[]{@{}
  >{\raggedright\arraybackslash}p{(\columnwidth - 4\tabcolsep) * \real{0.18}}
  >{\raggedright\arraybackslash}p{(\columnwidth - 4\tabcolsep) * \real{0.20}}
  >{\raggedright\arraybackslash}p{(\columnwidth - 4\tabcolsep) * \real{0.62}}@{}}
\toprule
位置 & 类型 & 含义 \\
\midrule
\endhead
返回值 & \passthrough{\lstinline!bool!} & 返回
\passthrough{\lstinline!bool!}。更精确的业务含义见该方法用途、C++
签名和目标型号协议。 \\
\bottomrule
\end{longtable}

\textbf{对应 C++}

\begin{itemize}
\tightlist
\item
  类：\passthrough{\lstinline!unitree\_hg::msg::dds\_::MainBoardState\_!}
\item
  签名：\passthrough{\lstinline"operator!=(const value \&) const"}
\item
  绑定策略：\passthrough{\lstinline!未单独标注!}
\item
  声明位置：由 pybind11 手工绑定或生成注册表提供；manifest
  未记录独立头文件行号。
\end{itemize}

\textbf{说明}

\begin{itemize}
\tightlist
\item
  示例是局部调用片段；\passthrough{\lstinline!obj!}
  和参数变量表示已经按上层业务校验并准备好的对象。
\end{itemize}

\textbf{用法}

\begin{lstlisting}[language=Python]
different = left != right
\end{lstlisting}

\hypertarget{unitree-sdk2-cpp-idl-hg-mainboardstate-fan-state}{%
\paragraph{\texorpdfstring{\texttt{unitree\_sdk2\_cpp.idl.hg.MainBoardState.fan\_state}}{unitree\_sdk2\_cpp.idl.hg.MainBoardState.fan\_state}}\label{unitree-sdk2-cpp-idl-hg-mainboardstate-fan-state}}

底层 \passthrough{\lstinline!unitree\_hg::msg::dds\_::MainBoardState\_!}
的 \passthrough{\lstinline!fan\_state!} 字段。类型签名只说明 Python
表示；单位、数值范围、数组固定长度和枚举含义需查对应 IDL 头文件。
底层是固定长度数组，写入序列必须正好包含 6 个元素；元素约束：取值范围为
0 到 65535。

\textbf{签名}

\begin{lstlisting}[language=Python]
@property
def fan_state(self) -> list[int]

@fan_state.setter
def fan_state(self, value: Sequence[int]) -> None
\end{lstlisting}

\textbf{可用性与安全性}

\textbf{\passthrough{\lstinline!AVAILABLE!} /
\passthrough{\lstinline!VALUE\_TYPE!}}：当前绑定源码已实现该消息值操作。

\textbf{参数}

\begin{longtable}[]{@{}
  >{\raggedright\arraybackslash}p{(\columnwidth - 4\tabcolsep) * \real{0.18}}
  >{\raggedright\arraybackslash}p{(\columnwidth - 4\tabcolsep) * \real{0.20}}
  >{\raggedright\arraybackslash}p{(\columnwidth - 4\tabcolsep) * \real{0.62}}@{}}
\toprule
名称 & Python 类型 & 含义 \\
\midrule
\endhead
\passthrough{\lstinline!value!} &
\passthrough{\lstinline!Sequence[int]!} &
要写入或传递的新值，必须符合该参数的 Python 类型以及底层 C++
范围约束。 \\
\bottomrule
\end{longtable}

\textbf{返回值}

\begin{longtable}[]{@{}
  >{\raggedright\arraybackslash}p{(\columnwidth - 4\tabcolsep) * \real{0.18}}
  >{\raggedright\arraybackslash}p{(\columnwidth - 4\tabcolsep) * \real{0.20}}
  >{\raggedright\arraybackslash}p{(\columnwidth - 4\tabcolsep) * \real{0.62}}@{}}
\toprule
位置 & 类型 & 含义 \\
\midrule
\endhead
getter 返回值 & \passthrough{\lstinline!list[int]!} &
返回属性当前值。数组、vector 和嵌套消息采用复制语义。 \\
\bottomrule
\end{longtable}

\textbf{对应 C++}

\begin{itemize}
\tightlist
\item
  类：\passthrough{\lstinline!unitree\_hg::msg::dds\_::MainBoardState\_!}
\item
  字段类型：\passthrough{\lstinline!std::array<uint16\_t, 6>!}
\item
  声明位置：\passthrough{\lstinline!include/unitree/idl/hg/MainBoardState\_.hpp:24!}
\end{itemize}

\textbf{用法}

\begin{lstlisting}[language=Python]
current_value = obj.fan_state
obj.fan_state = new_value
\end{lstlisting}

\begin{quote}
{[}!TIP{]} 读取容器或嵌套值后应遵循``读取、修改、重新赋值''。只修改
getter 返回的副本不会可靠地更新原 C++ 消息。
\end{quote}

\hypertarget{unitree-sdk2-cpp-idl-hg-mainboardstate-temperature}{%
\paragraph{\texorpdfstring{\texttt{unitree\_sdk2\_cpp.idl.hg.MainBoardState.temperature}}{unitree\_sdk2\_cpp.idl.hg.MainBoardState.temperature}}\label{unitree-sdk2-cpp-idl-hg-mainboardstate-temperature}}

底层 \passthrough{\lstinline!unitree\_hg::msg::dds\_::MainBoardState\_!}
的 \passthrough{\lstinline!temperature!} 字段。类型签名只说明 Python
表示；单位、数值范围、数组固定长度和枚举含义需查对应 IDL 头文件。
底层是固定长度数组，写入序列必须正好包含 6 个元素；元素约束：取值范围为
-32768 到 32767。

\textbf{签名}

\begin{lstlisting}[language=Python]
@property
def temperature(self) -> list[int]

@temperature.setter
def temperature(self, value: Sequence[int]) -> None
\end{lstlisting}

\textbf{可用性与安全性}

\textbf{\passthrough{\lstinline!AVAILABLE!} /
\passthrough{\lstinline!VALUE\_TYPE!}}：当前绑定源码已实现该消息值操作。

\textbf{参数}

\begin{longtable}[]{@{}
  >{\raggedright\arraybackslash}p{(\columnwidth - 4\tabcolsep) * \real{0.18}}
  >{\raggedright\arraybackslash}p{(\columnwidth - 4\tabcolsep) * \real{0.20}}
  >{\raggedright\arraybackslash}p{(\columnwidth - 4\tabcolsep) * \real{0.62}}@{}}
\toprule
名称 & Python 类型 & 含义 \\
\midrule
\endhead
\passthrough{\lstinline!value!} &
\passthrough{\lstinline!Sequence[int]!} &
要写入或传递的新值，必须符合该参数的 Python 类型以及底层 C++
范围约束。 \\
\bottomrule
\end{longtable}

\textbf{返回值}

\begin{longtable}[]{@{}
  >{\raggedright\arraybackslash}p{(\columnwidth - 4\tabcolsep) * \real{0.18}}
  >{\raggedright\arraybackslash}p{(\columnwidth - 4\tabcolsep) * \real{0.20}}
  >{\raggedright\arraybackslash}p{(\columnwidth - 4\tabcolsep) * \real{0.62}}@{}}
\toprule
位置 & 类型 & 含义 \\
\midrule
\endhead
getter 返回值 & \passthrough{\lstinline!list[int]!} &
返回属性当前值。数组、vector 和嵌套消息采用复制语义。 \\
\bottomrule
\end{longtable}

\textbf{对应 C++}

\begin{itemize}
\tightlist
\item
  类：\passthrough{\lstinline!unitree\_hg::msg::dds\_::MainBoardState\_!}
\item
  字段类型：\passthrough{\lstinline!std::array<int16\_t, 6>!}
\item
  声明位置：\passthrough{\lstinline!include/unitree/idl/hg/MainBoardState\_.hpp:25!}
\end{itemize}

\textbf{用法}

\begin{lstlisting}[language=Python]
current_value = obj.temperature
obj.temperature = new_value
\end{lstlisting}

\begin{quote}
{[}!TIP{]} 读取容器或嵌套值后应遵循``读取、修改、重新赋值''。只修改
getter 返回的副本不会可靠地更新原 C++ 消息。
\end{quote}

\hypertarget{unitree-sdk2-cpp-idl-hg-mainboardstate-value}{%
\paragraph{\texorpdfstring{\texttt{unitree\_sdk2\_cpp.idl.hg.MainBoardState.value}}{unitree\_sdk2\_cpp.idl.hg.MainBoardState.value}}\label{unitree-sdk2-cpp-idl-hg-mainboardstate-value}}

底层 \passthrough{\lstinline!unitree\_hg::msg::dds\_::MainBoardState\_!}
的 \passthrough{\lstinline!value!} 字段。类型签名只说明 Python
表示；单位、数值范围、数组固定长度和枚举含义需查对应 IDL 头文件。
底层是固定长度数组，写入序列必须正好包含 6
个元素；元素约束：底层为浮点数；类型本身不说明单位、坐标系或业务安全范围。

\textbf{签名}

\begin{lstlisting}[language=Python]
@property
def value(self) -> list[float]

@value.setter
def value(self, value: Sequence[float]) -> None
\end{lstlisting}

\textbf{可用性与安全性}

\textbf{\passthrough{\lstinline!AVAILABLE!} /
\passthrough{\lstinline!VALUE\_TYPE!}}：当前绑定源码已实现该消息值操作。

\textbf{参数}

\begin{longtable}[]{@{}
  >{\raggedright\arraybackslash}p{(\columnwidth - 4\tabcolsep) * \real{0.18}}
  >{\raggedright\arraybackslash}p{(\columnwidth - 4\tabcolsep) * \real{0.20}}
  >{\raggedright\arraybackslash}p{(\columnwidth - 4\tabcolsep) * \real{0.62}}@{}}
\toprule
名称 & Python 类型 & 含义 \\
\midrule
\endhead
\passthrough{\lstinline!value!} &
\passthrough{\lstinline!Sequence[float]!} &
要写入或传递的新值，必须符合该参数的 Python 类型以及底层 C++
范围约束。 \\
\bottomrule
\end{longtable}

\textbf{返回值}

\begin{longtable}[]{@{}
  >{\raggedright\arraybackslash}p{(\columnwidth - 4\tabcolsep) * \real{0.18}}
  >{\raggedright\arraybackslash}p{(\columnwidth - 4\tabcolsep) * \real{0.20}}
  >{\raggedright\arraybackslash}p{(\columnwidth - 4\tabcolsep) * \real{0.62}}@{}}
\toprule
位置 & 类型 & 含义 \\
\midrule
\endhead
getter 返回值 & \passthrough{\lstinline!list[float]!} &
返回属性当前值。数组、vector 和嵌套消息采用复制语义。 \\
\bottomrule
\end{longtable}

\textbf{对应 C++}

\begin{itemize}
\tightlist
\item
  类：\passthrough{\lstinline!unitree\_hg::msg::dds\_::MainBoardState\_!}
\item
  字段类型：\passthrough{\lstinline!std::array<float, 6>!}
\item
  声明位置：\passthrough{\lstinline!include/unitree/idl/hg/MainBoardState\_.hpp:26!}
\end{itemize}

\textbf{用法}

\begin{lstlisting}[language=Python]
current_value = obj.value
obj.value = new_value
\end{lstlisting}

\begin{quote}
{[}!TIP{]} 读取容器或嵌套值后应遵循``读取、修改、重新赋值''。只修改
getter 返回的副本不会可靠地更新原 C++ 消息。
\end{quote}

\hypertarget{unitree-sdk2-cpp-idl-hg-mainboardstate-state}{%
\paragraph{\texorpdfstring{\texttt{unitree\_sdk2\_cpp.idl.hg.MainBoardState.state}}{unitree\_sdk2\_cpp.idl.hg.MainBoardState.state}}\label{unitree-sdk2-cpp-idl-hg-mainboardstate-state}}

底层 \passthrough{\lstinline!unitree\_hg::msg::dds\_::MainBoardState\_!}
的 \passthrough{\lstinline!state!} 字段。类型签名只说明 Python
表示；单位、数值范围、数组固定长度和枚举含义需查对应 IDL 头文件。
底层是固定长度数组，写入序列必须正好包含 6 个元素；元素约束：取值范围为
0 到 4294967295。

\textbf{签名}

\begin{lstlisting}[language=Python]
@property
def state(self) -> list[int]

@state.setter
def state(self, value: Sequence[int]) -> None
\end{lstlisting}

\textbf{可用性与安全性}

\textbf{\passthrough{\lstinline!AVAILABLE!} /
\passthrough{\lstinline!VALUE\_TYPE!}}：当前绑定源码已实现该消息值操作。

\textbf{参数}

\begin{longtable}[]{@{}
  >{\raggedright\arraybackslash}p{(\columnwidth - 4\tabcolsep) * \real{0.18}}
  >{\raggedright\arraybackslash}p{(\columnwidth - 4\tabcolsep) * \real{0.20}}
  >{\raggedright\arraybackslash}p{(\columnwidth - 4\tabcolsep) * \real{0.62}}@{}}
\toprule
名称 & Python 类型 & 含义 \\
\midrule
\endhead
\passthrough{\lstinline!value!} &
\passthrough{\lstinline!Sequence[int]!} &
要写入或传递的新值，必须符合该参数的 Python 类型以及底层 C++
范围约束。 \\
\bottomrule
\end{longtable}

\textbf{返回值}

\begin{longtable}[]{@{}
  >{\raggedright\arraybackslash}p{(\columnwidth - 4\tabcolsep) * \real{0.18}}
  >{\raggedright\arraybackslash}p{(\columnwidth - 4\tabcolsep) * \real{0.20}}
  >{\raggedright\arraybackslash}p{(\columnwidth - 4\tabcolsep) * \real{0.62}}@{}}
\toprule
位置 & 类型 & 含义 \\
\midrule
\endhead
getter 返回值 & \passthrough{\lstinline!list[int]!} &
返回属性当前值。数组、vector 和嵌套消息采用复制语义。 \\
\bottomrule
\end{longtable}

\textbf{对应 C++}

\begin{itemize}
\tightlist
\item
  类：\passthrough{\lstinline!unitree\_hg::msg::dds\_::MainBoardState\_!}
\item
  字段类型：\passthrough{\lstinline!std::array<uint32\_t, 6>!}
\item
  声明位置：\passthrough{\lstinline!include/unitree/idl/hg/MainBoardState\_.hpp:27!}
\end{itemize}

\textbf{用法}

\begin{lstlisting}[language=Python]
current_value = obj.state
obj.state = new_value
\end{lstlisting}

\begin{quote}
{[}!TIP{]} 读取容器或嵌套值后应遵循``读取、修改、重新赋值''。只修改
getter 返回的副本不会可靠地更新原 C++ 消息。
\end{quote}

\hypertarget{unitree-sdk2-cpp-idl-hg-sportmodestate}{%
\subsubsection{\texorpdfstring{\texttt{unitree\_sdk2\_cpp.idl.hg.SportModeState}}{unitree\_sdk2\_cpp.idl.hg.SportModeState}}\label{unitree-sdk2-cpp-idl-hg-sportmodestate}}

可由 Python 读写的 DDS/IDL 消息值类型。字段采用复制语义。

\textbf{导入}

\begin{lstlisting}[language=Python]
from unitree_sdk2_cpp.idl.hg import SportModeState
\end{lstlisting}

\textbf{方法索引}

\begin{longtable}[]{@{}
  >{\raggedright\arraybackslash}p{(\columnwidth - 6\tabcolsep) * \real{0.18}}
  >{\raggedright\arraybackslash}p{(\columnwidth - 6\tabcolsep) * \real{0.24}}
  >{\raggedright\arraybackslash}p{(\columnwidth - 6\tabcolsep) * \real{0.18}}
  >{\raggedright\arraybackslash}p{(\columnwidth - 6\tabcolsep) * \real{0.40}}@{}}
\toprule
方法 & Python 签名 & 状态 & 安全分类 \\
\midrule
\endhead
\protect\hyperlink{unitree-sdk2-cpp-idl-hg-sportmodestate-dunder-init-1}{\passthrough{\lstinline!\_\_init\_\_!}}
& \passthrough{\lstinline!def \_\_init\_\_(self) -> None!} &
\passthrough{\lstinline!AVAILABLE!} &
\passthrough{\lstinline!VALUE\_TYPE!} \\
\protect\hyperlink{unitree-sdk2-cpp-idl-hg-sportmodestate-dunder-eq-1}{\passthrough{\lstinline!\_\_eq\_\_!}}
& \passthrough{\lstinline!def \_\_eq\_\_(self, other: object) -> bool!}
& \passthrough{\lstinline!AVAILABLE!} &
\passthrough{\lstinline!VALUE\_TYPE!} \\
\protect\hyperlink{unitree-sdk2-cpp-idl-hg-sportmodestate-dunder-ne-1}{\passthrough{\lstinline!\_\_ne\_\_!}}
& \passthrough{\lstinline!def \_\_ne\_\_(self, other: object) -> bool!}
& \passthrough{\lstinline!AVAILABLE!} &
\passthrough{\lstinline!VALUE\_TYPE!} \\
\bottomrule
\end{longtable}

\textbf{属性索引}

\begin{longtable}[]{@{}
  >{\raggedright\arraybackslash}p{(\columnwidth - 6\tabcolsep) * \real{0.18}}
  >{\raggedright\arraybackslash}p{(\columnwidth - 6\tabcolsep) * \real{0.24}}
  >{\raggedright\arraybackslash}p{(\columnwidth - 6\tabcolsep) * \real{0.18}}
  >{\raggedright\arraybackslash}p{(\columnwidth - 6\tabcolsep) * \real{0.40}}@{}}
\toprule
属性 & 读取类型 & 写入类型 & 状态 \\
\midrule
\endhead
\protect\hyperlink{unitree-sdk2-cpp-idl-hg-sportmodestate-fsm-id}{\passthrough{\lstinline!fsm\_id!}}
& \passthrough{\lstinline!int!} & \passthrough{\lstinline!int!} &
\passthrough{\lstinline!AVAILABLE!} \\
\protect\hyperlink{unitree-sdk2-cpp-idl-hg-sportmodestate-fsm-mode}{\passthrough{\lstinline!fsm\_mode!}}
& \passthrough{\lstinline!int!} & \passthrough{\lstinline!int!} &
\passthrough{\lstinline!AVAILABLE!} \\
\protect\hyperlink{unitree-sdk2-cpp-idl-hg-sportmodestate-task-id}{\passthrough{\lstinline!task\_id!}}
& \passthrough{\lstinline!int!} & \passthrough{\lstinline!int!} &
\passthrough{\lstinline!AVAILABLE!} \\
\protect\hyperlink{unitree-sdk2-cpp-idl-hg-sportmodestate-task-time}{\passthrough{\lstinline!task\_time!}}
& \passthrough{\lstinline!float!} & \passthrough{\lstinline!float!} &
\passthrough{\lstinline!AVAILABLE!} \\
\bottomrule
\end{longtable}

\hypertarget{unitree-sdk2-cpp-idl-hg-sportmodestate-dunder-init-1}{%
\paragraph{\texorpdfstring{\texttt{unitree\_sdk2\_cpp.idl.hg.SportModeState.\_\_init\_\_}}{unitree\_sdk2\_cpp.idl.hg.SportModeState.\_\_init\_\_}}\label{unitree-sdk2-cpp-idl-hg-sportmodestate-dunder-init-1}}

初始化 \passthrough{\lstinline!SportModeState!}
实例。是否能实际构造取决于下方可用性状态。

\textbf{签名}

\begin{lstlisting}[language=Python]
def __init__(self) -> None
\end{lstlisting}

\textbf{可用性与安全性}

\textbf{\passthrough{\lstinline!AVAILABLE!} /
\passthrough{\lstinline!VALUE\_TYPE!}}：当前绑定源码已实现该消息值操作。

\textbf{参数}

无显式参数。实例方法中的 \passthrough{\lstinline!self!} 由 Python
自动传入。

\textbf{返回值}

\begin{longtable}[]{@{}
  >{\raggedright\arraybackslash}p{(\columnwidth - 4\tabcolsep) * \real{0.18}}
  >{\raggedright\arraybackslash}p{(\columnwidth - 4\tabcolsep) * \real{0.20}}
  >{\raggedright\arraybackslash}p{(\columnwidth - 4\tabcolsep) * \real{0.62}}@{}}
\toprule
位置 & 类型 & 含义 \\
\midrule
\endhead
无 & \passthrough{\lstinline!None!} &
\passthrough{\lstinline!\_\_init\_\_!}
本身不返回值；调用类对象时，在构造函数可用的前提下得到该类实例。 \\
\bottomrule
\end{longtable}

\textbf{对应 C++}

\begin{itemize}
\tightlist
\item
  类：\passthrough{\lstinline!unitree\_hg::msg::dds\_::SportModeState\_!}
\item
  签名：\passthrough{\lstinline!SportModeState\_()!}
\item
  绑定策略：\passthrough{\lstinline!未单独标注!}
\item
  声明位置：\passthrough{\lstinline!include/unitree/idl/hg/SportModeState\_.hpp:29!}
\end{itemize}

\textbf{说明}

\begin{itemize}
\tightlist
\item
  示例是局部调用片段；\passthrough{\lstinline!obj!}
  和参数变量表示已经按上层业务校验并准备好的对象。
\end{itemize}

\textbf{用法}

\begin{lstlisting}[language=Python]
value = SportModeState()
\end{lstlisting}

\hypertarget{unitree-sdk2-cpp-idl-hg-sportmodestate-dunder-eq-1}{%
\paragraph{\texorpdfstring{\texttt{unitree\_sdk2\_cpp.idl.hg.SportModeState.\_\_eq\_\_}}{unitree\_sdk2\_cpp.idl.hg.SportModeState.\_\_eq\_\_}}\label{unitree-sdk2-cpp-idl-hg-sportmodestate-dunder-eq-1}}

按底层 IDL 消息内容比较两个对象是否相等。

\textbf{签名}

\begin{lstlisting}[language=Python]
def __eq__(self, other: object) -> bool
\end{lstlisting}

\textbf{可用性与安全性}

\textbf{\passthrough{\lstinline!AVAILABLE!} /
\passthrough{\lstinline!VALUE\_TYPE!}}：当前绑定源码已实现该消息值操作。

\textbf{参数}

\begin{longtable}[]{@{}
  >{\raggedright\arraybackslash}p{(\columnwidth - 6\tabcolsep) * \real{0.18}}
  >{\raggedright\arraybackslash}p{(\columnwidth - 6\tabcolsep) * \real{0.24}}
  >{\raggedright\arraybackslash}p{(\columnwidth - 6\tabcolsep) * \real{0.18}}
  >{\raggedright\arraybackslash}p{(\columnwidth - 6\tabcolsep) * \real{0.40}}@{}}
\toprule
名称 & Python 类型 & 默认值 & 含义 \\
\midrule
\endhead
\passthrough{\lstinline!other!} & \passthrough{\lstinline!object!} &
必填 & 用于比较的另一个 Python 对象。类型不兼容时比较结果通常为
\passthrough{\lstinline!False!}。 \\
\bottomrule
\end{longtable}

\textbf{返回值}

\begin{longtable}[]{@{}
  >{\raggedright\arraybackslash}p{(\columnwidth - 4\tabcolsep) * \real{0.18}}
  >{\raggedright\arraybackslash}p{(\columnwidth - 4\tabcolsep) * \real{0.20}}
  >{\raggedright\arraybackslash}p{(\columnwidth - 4\tabcolsep) * \real{0.62}}@{}}
\toprule
位置 & 类型 & 含义 \\
\midrule
\endhead
返回值 & \passthrough{\lstinline!bool!} & 返回
\passthrough{\lstinline!bool!}。更精确的业务含义见该方法用途、C++
签名和目标型号协议。 \\
\bottomrule
\end{longtable}

\textbf{对应 C++}

\begin{itemize}
\tightlist
\item
  类：\passthrough{\lstinline!unitree\_hg::msg::dds\_::SportModeState\_!}
\item
  签名：\passthrough{\lstinline!operator==(const value \&) const!}
\item
  绑定策略：\passthrough{\lstinline!未单独标注!}
\item
  声明位置：由 pybind11 手工绑定或生成注册表提供；manifest
  未记录独立头文件行号。
\end{itemize}

\textbf{说明}

\begin{itemize}
\tightlist
\item
  示例是局部调用片段；\passthrough{\lstinline!obj!}
  和参数变量表示已经按上层业务校验并准备好的对象。
\end{itemize}

\textbf{用法}

\begin{lstlisting}[language=Python]
same = left == right
\end{lstlisting}

\hypertarget{unitree-sdk2-cpp-idl-hg-sportmodestate-dunder-ne-1}{%
\paragraph{\texorpdfstring{\texttt{unitree\_sdk2\_cpp.idl.hg.SportModeState.\_\_ne\_\_}}{unitree\_sdk2\_cpp.idl.hg.SportModeState.\_\_ne\_\_}}\label{unitree-sdk2-cpp-idl-hg-sportmodestate-dunder-ne-1}}

按底层 IDL 消息内容比较两个对象是否不相等。

\textbf{签名}

\begin{lstlisting}[language=Python]
def __ne__(self, other: object) -> bool
\end{lstlisting}

\textbf{可用性与安全性}

\textbf{\passthrough{\lstinline!AVAILABLE!} /
\passthrough{\lstinline!VALUE\_TYPE!}}：当前绑定源码已实现该消息值操作。

\textbf{参数}

\begin{longtable}[]{@{}
  >{\raggedright\arraybackslash}p{(\columnwidth - 6\tabcolsep) * \real{0.18}}
  >{\raggedright\arraybackslash}p{(\columnwidth - 6\tabcolsep) * \real{0.24}}
  >{\raggedright\arraybackslash}p{(\columnwidth - 6\tabcolsep) * \real{0.18}}
  >{\raggedright\arraybackslash}p{(\columnwidth - 6\tabcolsep) * \real{0.40}}@{}}
\toprule
名称 & Python 类型 & 默认值 & 含义 \\
\midrule
\endhead
\passthrough{\lstinline!other!} & \passthrough{\lstinline!object!} &
必填 & 用于比较的另一个 Python 对象。类型不兼容时比较结果通常为
\passthrough{\lstinline!False!}。 \\
\bottomrule
\end{longtable}

\textbf{返回值}

\begin{longtable}[]{@{}
  >{\raggedright\arraybackslash}p{(\columnwidth - 4\tabcolsep) * \real{0.18}}
  >{\raggedright\arraybackslash}p{(\columnwidth - 4\tabcolsep) * \real{0.20}}
  >{\raggedright\arraybackslash}p{(\columnwidth - 4\tabcolsep) * \real{0.62}}@{}}
\toprule
位置 & 类型 & 含义 \\
\midrule
\endhead
返回值 & \passthrough{\lstinline!bool!} & 返回
\passthrough{\lstinline!bool!}。更精确的业务含义见该方法用途、C++
签名和目标型号协议。 \\
\bottomrule
\end{longtable}

\textbf{对应 C++}

\begin{itemize}
\tightlist
\item
  类：\passthrough{\lstinline!unitree\_hg::msg::dds\_::SportModeState\_!}
\item
  签名：\passthrough{\lstinline"operator!=(const value \&) const"}
\item
  绑定策略：\passthrough{\lstinline!未单独标注!}
\item
  声明位置：由 pybind11 手工绑定或生成注册表提供；manifest
  未记录独立头文件行号。
\end{itemize}

\textbf{说明}

\begin{itemize}
\tightlist
\item
  示例是局部调用片段；\passthrough{\lstinline!obj!}
  和参数变量表示已经按上层业务校验并准备好的对象。
\end{itemize}

\textbf{用法}

\begin{lstlisting}[language=Python]
different = left != right
\end{lstlisting}

\hypertarget{unitree-sdk2-cpp-idl-hg-sportmodestate-fsm-id}{%
\paragraph{\texorpdfstring{\texttt{unitree\_sdk2\_cpp.idl.hg.SportModeState.fsm\_id}}{unitree\_sdk2\_cpp.idl.hg.SportModeState.fsm\_id}}\label{unitree-sdk2-cpp-idl-hg-sportmodestate-fsm-id}}

底层 \passthrough{\lstinline!unitree\_hg::msg::dds\_::SportModeState\_!}
的 \passthrough{\lstinline!fsm\_id!} 字段。类型签名只说明 Python
表示；单位、数值范围、数组固定长度和枚举含义需查对应 IDL 头文件。
取值范围为 0 到 4294967295。

\textbf{签名}

\begin{lstlisting}[language=Python]
@property
def fsm_id(self) -> int

@fsm_id.setter
def fsm_id(self, value: int) -> None
\end{lstlisting}

\textbf{可用性与安全性}

\textbf{\passthrough{\lstinline!AVAILABLE!} /
\passthrough{\lstinline!VALUE\_TYPE!}}：当前绑定源码已实现该消息值操作。

\textbf{参数}

\begin{longtable}[]{@{}
  >{\raggedright\arraybackslash}p{(\columnwidth - 4\tabcolsep) * \real{0.18}}
  >{\raggedright\arraybackslash}p{(\columnwidth - 4\tabcolsep) * \real{0.20}}
  >{\raggedright\arraybackslash}p{(\columnwidth - 4\tabcolsep) * \real{0.62}}@{}}
\toprule
名称 & Python 类型 & 含义 \\
\midrule
\endhead
\passthrough{\lstinline!value!} & \passthrough{\lstinline!int!} &
要写入或传递的新值，必须符合该参数的 Python 类型以及底层 C++
范围约束。 \\
\bottomrule
\end{longtable}

\textbf{返回值}

\begin{longtable}[]{@{}
  >{\raggedright\arraybackslash}p{(\columnwidth - 4\tabcolsep) * \real{0.18}}
  >{\raggedright\arraybackslash}p{(\columnwidth - 4\tabcolsep) * \real{0.20}}
  >{\raggedright\arraybackslash}p{(\columnwidth - 4\tabcolsep) * \real{0.62}}@{}}
\toprule
位置 & 类型 & 含义 \\
\midrule
\endhead
getter 返回值 & \passthrough{\lstinline!int!} & 返回属性当前值。 \\
\bottomrule
\end{longtable}

\textbf{对应 C++}

\begin{itemize}
\tightlist
\item
  类：\passthrough{\lstinline!unitree\_hg::msg::dds\_::SportModeState\_!}
\item
  字段类型：\passthrough{\lstinline!uint32\_t!}
\item
  声明位置：\passthrough{\lstinline!include/unitree/idl/hg/SportModeState\_.hpp:23!}
\end{itemize}

\textbf{用法}

\begin{lstlisting}[language=Python]
current_value = obj.fsm_id
obj.fsm_id = new_value
\end{lstlisting}

\hypertarget{unitree-sdk2-cpp-idl-hg-sportmodestate-fsm-mode}{%
\paragraph{\texorpdfstring{\texttt{unitree\_sdk2\_cpp.idl.hg.SportModeState.fsm\_mode}}{unitree\_sdk2\_cpp.idl.hg.SportModeState.fsm\_mode}}\label{unitree-sdk2-cpp-idl-hg-sportmodestate-fsm-mode}}

底层 \passthrough{\lstinline!unitree\_hg::msg::dds\_::SportModeState\_!}
的 \passthrough{\lstinline!fsm\_mode!} 字段。类型签名只说明 Python
表示；单位、数值范围、数组固定长度和枚举含义需查对应 IDL 头文件。
取值范围为 0 到 4294967295。

\textbf{签名}

\begin{lstlisting}[language=Python]
@property
def fsm_mode(self) -> int

@fsm_mode.setter
def fsm_mode(self, value: int) -> None
\end{lstlisting}

\textbf{可用性与安全性}

\textbf{\passthrough{\lstinline!AVAILABLE!} /
\passthrough{\lstinline!VALUE\_TYPE!}}：当前绑定源码已实现该消息值操作。

\textbf{参数}

\begin{longtable}[]{@{}
  >{\raggedright\arraybackslash}p{(\columnwidth - 4\tabcolsep) * \real{0.18}}
  >{\raggedright\arraybackslash}p{(\columnwidth - 4\tabcolsep) * \real{0.20}}
  >{\raggedright\arraybackslash}p{(\columnwidth - 4\tabcolsep) * \real{0.62}}@{}}
\toprule
名称 & Python 类型 & 含义 \\
\midrule
\endhead
\passthrough{\lstinline!value!} & \passthrough{\lstinline!int!} &
要写入或传递的新值，必须符合该参数的 Python 类型以及底层 C++
范围约束。 \\
\bottomrule
\end{longtable}

\textbf{返回值}

\begin{longtable}[]{@{}
  >{\raggedright\arraybackslash}p{(\columnwidth - 4\tabcolsep) * \real{0.18}}
  >{\raggedright\arraybackslash}p{(\columnwidth - 4\tabcolsep) * \real{0.20}}
  >{\raggedright\arraybackslash}p{(\columnwidth - 4\tabcolsep) * \real{0.62}}@{}}
\toprule
位置 & 类型 & 含义 \\
\midrule
\endhead
getter 返回值 & \passthrough{\lstinline!int!} & 返回属性当前值。 \\
\bottomrule
\end{longtable}

\textbf{对应 C++}

\begin{itemize}
\tightlist
\item
  类：\passthrough{\lstinline!unitree\_hg::msg::dds\_::SportModeState\_!}
\item
  字段类型：\passthrough{\lstinline!uint32\_t!}
\item
  声明位置：\passthrough{\lstinline!include/unitree/idl/hg/SportModeState\_.hpp:24!}
\end{itemize}

\textbf{用法}

\begin{lstlisting}[language=Python]
current_value = obj.fsm_mode
obj.fsm_mode = new_value
\end{lstlisting}

\hypertarget{unitree-sdk2-cpp-idl-hg-sportmodestate-task-id}{%
\paragraph{\texorpdfstring{\texttt{unitree\_sdk2\_cpp.idl.hg.SportModeState.task\_id}}{unitree\_sdk2\_cpp.idl.hg.SportModeState.task\_id}}\label{unitree-sdk2-cpp-idl-hg-sportmodestate-task-id}}

底层 \passthrough{\lstinline!unitree\_hg::msg::dds\_::SportModeState\_!}
的 \passthrough{\lstinline!task\_id!} 字段。类型签名只说明 Python
表示；单位、数值范围、数组固定长度和枚举含义需查对应 IDL 头文件。
取值范围为 0 到 4294967295。

\textbf{签名}

\begin{lstlisting}[language=Python]
@property
def task_id(self) -> int

@task_id.setter
def task_id(self, value: int) -> None
\end{lstlisting}

\textbf{可用性与安全性}

\textbf{\passthrough{\lstinline!AVAILABLE!} /
\passthrough{\lstinline!VALUE\_TYPE!}}：当前绑定源码已实现该消息值操作。

\textbf{参数}

\begin{longtable}[]{@{}
  >{\raggedright\arraybackslash}p{(\columnwidth - 4\tabcolsep) * \real{0.18}}
  >{\raggedright\arraybackslash}p{(\columnwidth - 4\tabcolsep) * \real{0.20}}
  >{\raggedright\arraybackslash}p{(\columnwidth - 4\tabcolsep) * \real{0.62}}@{}}
\toprule
名称 & Python 类型 & 含义 \\
\midrule
\endhead
\passthrough{\lstinline!value!} & \passthrough{\lstinline!int!} &
要写入或传递的新值，必须符合该参数的 Python 类型以及底层 C++
范围约束。 \\
\bottomrule
\end{longtable}

\textbf{返回值}

\begin{longtable}[]{@{}
  >{\raggedright\arraybackslash}p{(\columnwidth - 4\tabcolsep) * \real{0.18}}
  >{\raggedright\arraybackslash}p{(\columnwidth - 4\tabcolsep) * \real{0.20}}
  >{\raggedright\arraybackslash}p{(\columnwidth - 4\tabcolsep) * \real{0.62}}@{}}
\toprule
位置 & 类型 & 含义 \\
\midrule
\endhead
getter 返回值 & \passthrough{\lstinline!int!} & 返回属性当前值。 \\
\bottomrule
\end{longtable}

\textbf{对应 C++}

\begin{itemize}
\tightlist
\item
  类：\passthrough{\lstinline!unitree\_hg::msg::dds\_::SportModeState\_!}
\item
  字段类型：\passthrough{\lstinline!uint32\_t!}
\item
  声明位置：\passthrough{\lstinline!include/unitree/idl/hg/SportModeState\_.hpp:25!}
\end{itemize}

\textbf{用法}

\begin{lstlisting}[language=Python]
current_value = obj.task_id
obj.task_id = new_value
\end{lstlisting}

\hypertarget{unitree-sdk2-cpp-idl-hg-sportmodestate-task-time}{%
\paragraph{\texorpdfstring{\texttt{unitree\_sdk2\_cpp.idl.hg.SportModeState.task\_time}}{unitree\_sdk2\_cpp.idl.hg.SportModeState.task\_time}}\label{unitree-sdk2-cpp-idl-hg-sportmodestate-task-time}}

底层 \passthrough{\lstinline!unitree\_hg::msg::dds\_::SportModeState\_!}
的 \passthrough{\lstinline!task\_time!} 字段。类型签名只说明 Python
表示；单位、数值范围、数组固定长度和枚举含义需查对应 IDL 头文件。
底层为浮点数；类型本身不说明单位、坐标系或业务安全范围。

\textbf{签名}

\begin{lstlisting}[language=Python]
@property
def task_time(self) -> float

@task_time.setter
def task_time(self, value: float) -> None
\end{lstlisting}

\textbf{可用性与安全性}

\textbf{\passthrough{\lstinline!AVAILABLE!} /
\passthrough{\lstinline!VALUE\_TYPE!}}：当前绑定源码已实现该消息值操作。

\textbf{参数}

\begin{longtable}[]{@{}
  >{\raggedright\arraybackslash}p{(\columnwidth - 4\tabcolsep) * \real{0.18}}
  >{\raggedright\arraybackslash}p{(\columnwidth - 4\tabcolsep) * \real{0.20}}
  >{\raggedright\arraybackslash}p{(\columnwidth - 4\tabcolsep) * \real{0.62}}@{}}
\toprule
名称 & Python 类型 & 含义 \\
\midrule
\endhead
\passthrough{\lstinline!value!} & \passthrough{\lstinline!float!} &
要写入或传递的新值，必须符合该参数的 Python 类型以及底层 C++
范围约束。 \\
\bottomrule
\end{longtable}

\textbf{返回值}

\begin{longtable}[]{@{}
  >{\raggedright\arraybackslash}p{(\columnwidth - 4\tabcolsep) * \real{0.18}}
  >{\raggedright\arraybackslash}p{(\columnwidth - 4\tabcolsep) * \real{0.20}}
  >{\raggedright\arraybackslash}p{(\columnwidth - 4\tabcolsep) * \real{0.62}}@{}}
\toprule
位置 & 类型 & 含义 \\
\midrule
\endhead
getter 返回值 & \passthrough{\lstinline!float!} & 返回属性当前值。 \\
\bottomrule
\end{longtable}

\textbf{对应 C++}

\begin{itemize}
\tightlist
\item
  类：\passthrough{\lstinline!unitree\_hg::msg::dds\_::SportModeState\_!}
\item
  字段类型：\passthrough{\lstinline!float!}
\item
  声明位置：\passthrough{\lstinline!include/unitree/idl/hg/SportModeState\_.hpp:26!}
\end{itemize}

\textbf{用法}

\begin{lstlisting}[language=Python]
current_value = obj.task_time
obj.task_time = new_value
\end{lstlisting}

\begin{center}\rule{0.5\linewidth}{0.5pt}\end{center}

\hypertarget{unitree-sdk2-cpp-idl-hg-doubleimu}{%
\subsection{HG Double-IMU IDL
消息}\label{unitree-sdk2-cpp-idl-hg-doubleimu}}

模块：\passthrough{\lstinline!unitree\_sdk2\_cpp.idl.hg\_doubleimu!}

\hypertarget{ux7c7bux7d22ux5f15-4}{%
\subsubsection{类索引}\label{ux7c7bux7d22ux5f15-4}}

\begin{longtable}[]{@{}
  >{\raggedright\arraybackslash}p{(\columnwidth - 4\tabcolsep) * \real{0.18}}
  >{\raggedleft\arraybackslash}p{(\columnwidth - 4\tabcolsep) * \real{0.20}}
  >{\raggedleft\arraybackslash}p{(\columnwidth - 4\tabcolsep) * \real{0.62}}@{}}
\toprule
类 & 公开函数签名 & 属性 \\
\midrule
\endhead
\protect\hyperlink{unitree-sdk2-cpp-idl-hg-doubleimu-doubleimustate}{\passthrough{\lstinline!doubleIMUState!}}
& 3 & 6 \\
\bottomrule
\end{longtable}

\hypertarget{unitree-sdk2-cpp-idl-hg-doubleimu-doubleimustate}{%
\subsubsection{\texorpdfstring{\texttt{unitree\_sdk2\_cpp.idl.hg\_doubleimu.doubleIMUState}}{unitree\_sdk2\_cpp.idl.hg\_doubleimu.doubleIMUState}}\label{unitree-sdk2-cpp-idl-hg-doubleimu-doubleimustate}}

可由 Python 读写的 DDS/IDL 消息值类型。字段采用复制语义。

\textbf{导入}

\begin{lstlisting}[language=Python]
from unitree_sdk2_cpp.idl.hg_doubleimu import doubleIMUState
\end{lstlisting}

\textbf{方法索引}

\begin{longtable}[]{@{}
  >{\raggedright\arraybackslash}p{(\columnwidth - 6\tabcolsep) * \real{0.18}}
  >{\raggedright\arraybackslash}p{(\columnwidth - 6\tabcolsep) * \real{0.24}}
  >{\raggedright\arraybackslash}p{(\columnwidth - 6\tabcolsep) * \real{0.18}}
  >{\raggedright\arraybackslash}p{(\columnwidth - 6\tabcolsep) * \real{0.40}}@{}}
\toprule
方法 & Python 签名 & 状态 & 安全分类 \\
\midrule
\endhead
\protect\hyperlink{unitree-sdk2-cpp-idl-hg-doubleimu-doubleimustate-dunder-init-1}{\passthrough{\lstinline!\_\_init\_\_!}}
& \passthrough{\lstinline!def \_\_init\_\_(self) -> None!} &
\passthrough{\lstinline!AVAILABLE!} &
\passthrough{\lstinline!VALUE\_TYPE!} \\
\protect\hyperlink{unitree-sdk2-cpp-idl-hg-doubleimu-doubleimustate-dunder-eq-1}{\passthrough{\lstinline!\_\_eq\_\_!}}
& \passthrough{\lstinline!def \_\_eq\_\_(self, other: object) -> bool!}
& \passthrough{\lstinline!AVAILABLE!} &
\passthrough{\lstinline!VALUE\_TYPE!} \\
\protect\hyperlink{unitree-sdk2-cpp-idl-hg-doubleimu-doubleimustate-dunder-ne-1}{\passthrough{\lstinline!\_\_ne\_\_!}}
& \passthrough{\lstinline!def \_\_ne\_\_(self, other: object) -> bool!}
& \passthrough{\lstinline!AVAILABLE!} &
\passthrough{\lstinline!VALUE\_TYPE!} \\
\bottomrule
\end{longtable}

\textbf{属性索引}

\begin{longtable}[]{@{}
  >{\raggedright\arraybackslash}p{(\columnwidth - 6\tabcolsep) * \real{0.18}}
  >{\raggedright\arraybackslash}p{(\columnwidth - 6\tabcolsep) * \real{0.24}}
  >{\raggedright\arraybackslash}p{(\columnwidth - 6\tabcolsep) * \real{0.18}}
  >{\raggedright\arraybackslash}p{(\columnwidth - 6\tabcolsep) * \real{0.40}}@{}}
\toprule
属性 & 读取类型 & 写入类型 & 状态 \\
\midrule
\endhead
\protect\hyperlink{unitree-sdk2-cpp-idl-hg-doubleimu-doubleimustate-quaternion}{\passthrough{\lstinline!quaternion!}}
& \passthrough{\lstinline!list[float]!} &
\passthrough{\lstinline!Sequence[float]!} &
\passthrough{\lstinline!AVAILABLE!} \\
\protect\hyperlink{unitree-sdk2-cpp-idl-hg-doubleimu-doubleimustate-gyroscope}{\passthrough{\lstinline!gyroscope!}}
& \passthrough{\lstinline!list[float]!} &
\passthrough{\lstinline!Sequence[float]!} &
\passthrough{\lstinline!AVAILABLE!} \\
\protect\hyperlink{unitree-sdk2-cpp-idl-hg-doubleimu-doubleimustate-accelerometer}{\passthrough{\lstinline!accelerometer!}}
& \passthrough{\lstinline!list[float]!} &
\passthrough{\lstinline!Sequence[float]!} &
\passthrough{\lstinline!AVAILABLE!} \\
\protect\hyperlink{unitree-sdk2-cpp-idl-hg-doubleimu-doubleimustate-rpy}{\passthrough{\lstinline!rpy!}}
& \passthrough{\lstinline!list[float]!} &
\passthrough{\lstinline!Sequence[float]!} &
\passthrough{\lstinline!AVAILABLE!} \\
\protect\hyperlink{unitree-sdk2-cpp-idl-hg-doubleimu-doubleimustate-temperature}{\passthrough{\lstinline!temperature!}}
& \passthrough{\lstinline!int!} & \passthrough{\lstinline!int!} &
\passthrough{\lstinline!AVAILABLE!} \\
\protect\hyperlink{unitree-sdk2-cpp-idl-hg-doubleimu-doubleimustate-tick}{\passthrough{\lstinline!tick!}}
& \passthrough{\lstinline!int!} & \passthrough{\lstinline!int!} &
\passthrough{\lstinline!AVAILABLE!} \\
\bottomrule
\end{longtable}

\hypertarget{unitree-sdk2-cpp-idl-hg-doubleimu-doubleimustate-dunder-init-1}{%
\paragraph{\texorpdfstring{\texttt{unitree\_sdk2\_cpp.idl.hg\_doubleimu.doubleIMUState.\_\_init\_\_}}{unitree\_sdk2\_cpp.idl.hg\_doubleimu.doubleIMUState.\_\_init\_\_}}\label{unitree-sdk2-cpp-idl-hg-doubleimu-doubleimustate-dunder-init-1}}

初始化 \passthrough{\lstinline!doubleIMUState!}
实例。是否能实际构造取决于下方可用性状态。

\textbf{签名}

\begin{lstlisting}[language=Python]
def __init__(self) -> None
\end{lstlisting}

\textbf{可用性与安全性}

\textbf{\passthrough{\lstinline!AVAILABLE!} /
\passthrough{\lstinline!VALUE\_TYPE!}}：当前绑定源码已实现该消息值操作。

\textbf{参数}

无显式参数。实例方法中的 \passthrough{\lstinline!self!} 由 Python
自动传入。

\textbf{返回值}

\begin{longtable}[]{@{}
  >{\raggedright\arraybackslash}p{(\columnwidth - 4\tabcolsep) * \real{0.18}}
  >{\raggedright\arraybackslash}p{(\columnwidth - 4\tabcolsep) * \real{0.20}}
  >{\raggedright\arraybackslash}p{(\columnwidth - 4\tabcolsep) * \real{0.62}}@{}}
\toprule
位置 & 类型 & 含义 \\
\midrule
\endhead
无 & \passthrough{\lstinline!None!} &
\passthrough{\lstinline!\_\_init\_\_!}
本身不返回值；调用类对象时，在构造函数可用的前提下得到该类实例。 \\
\bottomrule
\end{longtable}

\textbf{对应 C++}

\begin{itemize}
\tightlist
\item
  类：\passthrough{\lstinline!unitree\_hg\_doubleimu::msg::dds\_::doubleIMUState\_!}
\item
  签名：\passthrough{\lstinline!doubleIMUState\_()!}
\item
  绑定策略：\passthrough{\lstinline!未单独标注!}
\item
  声明位置：\passthrough{\lstinline!include/unitree/idl/hg\_doubleimu/doubleIMUState\_.hpp:32!}
\end{itemize}

\textbf{说明}

\begin{itemize}
\tightlist
\item
  示例是局部调用片段；\passthrough{\lstinline!obj!}
  和参数变量表示已经按上层业务校验并准备好的对象。
\end{itemize}

\textbf{用法}

\begin{lstlisting}[language=Python]
value = doubleIMUState()
\end{lstlisting}

\hypertarget{unitree-sdk2-cpp-idl-hg-doubleimu-doubleimustate-dunder-eq-1}{%
\paragraph{\texorpdfstring{\texttt{unitree\_sdk2\_cpp.idl.hg\_doubleimu.doubleIMUState.\_\_eq\_\_}}{unitree\_sdk2\_cpp.idl.hg\_doubleimu.doubleIMUState.\_\_eq\_\_}}\label{unitree-sdk2-cpp-idl-hg-doubleimu-doubleimustate-dunder-eq-1}}

按底层 IDL 消息内容比较两个对象是否相等。

\textbf{签名}

\begin{lstlisting}[language=Python]
def __eq__(self, other: object) -> bool
\end{lstlisting}

\textbf{可用性与安全性}

\textbf{\passthrough{\lstinline!AVAILABLE!} /
\passthrough{\lstinline!VALUE\_TYPE!}}：当前绑定源码已实现该消息值操作。

\textbf{参数}

\begin{longtable}[]{@{}
  >{\raggedright\arraybackslash}p{(\columnwidth - 6\tabcolsep) * \real{0.18}}
  >{\raggedright\arraybackslash}p{(\columnwidth - 6\tabcolsep) * \real{0.24}}
  >{\raggedright\arraybackslash}p{(\columnwidth - 6\tabcolsep) * \real{0.18}}
  >{\raggedright\arraybackslash}p{(\columnwidth - 6\tabcolsep) * \real{0.40}}@{}}
\toprule
名称 & Python 类型 & 默认值 & 含义 \\
\midrule
\endhead
\passthrough{\lstinline!other!} & \passthrough{\lstinline!object!} &
必填 & 用于比较的另一个 Python 对象。类型不兼容时比较结果通常为
\passthrough{\lstinline!False!}。 \\
\bottomrule
\end{longtable}

\textbf{返回值}

\begin{longtable}[]{@{}
  >{\raggedright\arraybackslash}p{(\columnwidth - 4\tabcolsep) * \real{0.18}}
  >{\raggedright\arraybackslash}p{(\columnwidth - 4\tabcolsep) * \real{0.20}}
  >{\raggedright\arraybackslash}p{(\columnwidth - 4\tabcolsep) * \real{0.62}}@{}}
\toprule
位置 & 类型 & 含义 \\
\midrule
\endhead
返回值 & \passthrough{\lstinline!bool!} & 返回
\passthrough{\lstinline!bool!}。更精确的业务含义见该方法用途、C++
签名和目标型号协议。 \\
\bottomrule
\end{longtable}

\textbf{对应 C++}

\begin{itemize}
\tightlist
\item
  类：\passthrough{\lstinline!unitree\_hg\_doubleimu::msg::dds\_::doubleIMUState\_!}
\item
  签名：\passthrough{\lstinline!operator==(const value \&) const!}
\item
  绑定策略：\passthrough{\lstinline!未单独标注!}
\item
  声明位置：由 pybind11 手工绑定或生成注册表提供；manifest
  未记录独立头文件行号。
\end{itemize}

\textbf{说明}

\begin{itemize}
\tightlist
\item
  示例是局部调用片段；\passthrough{\lstinline!obj!}
  和参数变量表示已经按上层业务校验并准备好的对象。
\end{itemize}

\textbf{用法}

\begin{lstlisting}[language=Python]
same = left == right
\end{lstlisting}

\hypertarget{unitree-sdk2-cpp-idl-hg-doubleimu-doubleimustate-dunder-ne-1}{%
\paragraph{\texorpdfstring{\texttt{unitree\_sdk2\_cpp.idl.hg\_doubleimu.doubleIMUState.\_\_ne\_\_}}{unitree\_sdk2\_cpp.idl.hg\_doubleimu.doubleIMUState.\_\_ne\_\_}}\label{unitree-sdk2-cpp-idl-hg-doubleimu-doubleimustate-dunder-ne-1}}

按底层 IDL 消息内容比较两个对象是否不相等。

\textbf{签名}

\begin{lstlisting}[language=Python]
def __ne__(self, other: object) -> bool
\end{lstlisting}

\textbf{可用性与安全性}

\textbf{\passthrough{\lstinline!AVAILABLE!} /
\passthrough{\lstinline!VALUE\_TYPE!}}：当前绑定源码已实现该消息值操作。

\textbf{参数}

\begin{longtable}[]{@{}
  >{\raggedright\arraybackslash}p{(\columnwidth - 6\tabcolsep) * \real{0.18}}
  >{\raggedright\arraybackslash}p{(\columnwidth - 6\tabcolsep) * \real{0.24}}
  >{\raggedright\arraybackslash}p{(\columnwidth - 6\tabcolsep) * \real{0.18}}
  >{\raggedright\arraybackslash}p{(\columnwidth - 6\tabcolsep) * \real{0.40}}@{}}
\toprule
名称 & Python 类型 & 默认值 & 含义 \\
\midrule
\endhead
\passthrough{\lstinline!other!} & \passthrough{\lstinline!object!} &
必填 & 用于比较的另一个 Python 对象。类型不兼容时比较结果通常为
\passthrough{\lstinline!False!}。 \\
\bottomrule
\end{longtable}

\textbf{返回值}

\begin{longtable}[]{@{}
  >{\raggedright\arraybackslash}p{(\columnwidth - 4\tabcolsep) * \real{0.18}}
  >{\raggedright\arraybackslash}p{(\columnwidth - 4\tabcolsep) * \real{0.20}}
  >{\raggedright\arraybackslash}p{(\columnwidth - 4\tabcolsep) * \real{0.62}}@{}}
\toprule
位置 & 类型 & 含义 \\
\midrule
\endhead
返回值 & \passthrough{\lstinline!bool!} & 返回
\passthrough{\lstinline!bool!}。更精确的业务含义见该方法用途、C++
签名和目标型号协议。 \\
\bottomrule
\end{longtable}

\textbf{对应 C++}

\begin{itemize}
\tightlist
\item
  类：\passthrough{\lstinline!unitree\_hg\_doubleimu::msg::dds\_::doubleIMUState\_!}
\item
  签名：\passthrough{\lstinline"operator!=(const value \&) const"}
\item
  绑定策略：\passthrough{\lstinline!未单独标注!}
\item
  声明位置：由 pybind11 手工绑定或生成注册表提供；manifest
  未记录独立头文件行号。
\end{itemize}

\textbf{说明}

\begin{itemize}
\tightlist
\item
  示例是局部调用片段；\passthrough{\lstinline!obj!}
  和参数变量表示已经按上层业务校验并准备好的对象。
\end{itemize}

\textbf{用法}

\begin{lstlisting}[language=Python]
different = left != right
\end{lstlisting}

\hypertarget{unitree-sdk2-cpp-idl-hg-doubleimu-doubleimustate-quaternion}{%
\paragraph{\texorpdfstring{\texttt{unitree\_sdk2\_cpp.idl.hg\_doubleimu.doubleIMUState.quaternion}}{unitree\_sdk2\_cpp.idl.hg\_doubleimu.doubleIMUState.quaternion}}\label{unitree-sdk2-cpp-idl-hg-doubleimu-doubleimustate-quaternion}}

底层
\passthrough{\lstinline!unitree\_hg\_doubleimu::msg::dds\_::doubleIMUState\_!}
的 \passthrough{\lstinline!quaternion!} 字段。类型签名只说明 Python
表示；单位、数值范围、数组固定长度和枚举含义需查对应 IDL 头文件。
底层是固定长度数组，写入序列必须正好包含 4
个元素；元素约束：底层为浮点数；类型本身不说明单位、坐标系或业务安全范围。

\textbf{签名}

\begin{lstlisting}[language=Python]
@property
def quaternion(self) -> list[float]

@quaternion.setter
def quaternion(self, value: Sequence[float]) -> None
\end{lstlisting}

\textbf{可用性与安全性}

\textbf{\passthrough{\lstinline!AVAILABLE!} /
\passthrough{\lstinline!VALUE\_TYPE!}}：当前绑定源码已实现该消息值操作。

\textbf{参数}

\begin{longtable}[]{@{}
  >{\raggedright\arraybackslash}p{(\columnwidth - 4\tabcolsep) * \real{0.18}}
  >{\raggedright\arraybackslash}p{(\columnwidth - 4\tabcolsep) * \real{0.20}}
  >{\raggedright\arraybackslash}p{(\columnwidth - 4\tabcolsep) * \real{0.62}}@{}}
\toprule
名称 & Python 类型 & 含义 \\
\midrule
\endhead
\passthrough{\lstinline!value!} &
\passthrough{\lstinline!Sequence[float]!} &
要写入或传递的新值，必须符合该参数的 Python 类型以及底层 C++
范围约束。 \\
\bottomrule
\end{longtable}

\textbf{返回值}

\begin{longtable}[]{@{}
  >{\raggedright\arraybackslash}p{(\columnwidth - 4\tabcolsep) * \real{0.18}}
  >{\raggedright\arraybackslash}p{(\columnwidth - 4\tabcolsep) * \real{0.20}}
  >{\raggedright\arraybackslash}p{(\columnwidth - 4\tabcolsep) * \real{0.62}}@{}}
\toprule
位置 & 类型 & 含义 \\
\midrule
\endhead
getter 返回值 & \passthrough{\lstinline!list[float]!} &
返回属性当前值。数组、vector 和嵌套消息采用复制语义。 \\
\bottomrule
\end{longtable}

\textbf{对应 C++}

\begin{itemize}
\tightlist
\item
  类：\passthrough{\lstinline!unitree\_hg\_doubleimu::msg::dds\_::doubleIMUState\_!}
\item
  字段类型：\passthrough{\lstinline!std::array<float, 4>!}
\item
  声明位置：\passthrough{\lstinline!include/unitree/idl/hg\_doubleimu/doubleIMUState\_.hpp:24!}
\end{itemize}

\textbf{用法}

\begin{lstlisting}[language=Python]
current_value = obj.quaternion
obj.quaternion = new_value
\end{lstlisting}

\begin{quote}
{[}!TIP{]} 读取容器或嵌套值后应遵循``读取、修改、重新赋值''。只修改
getter 返回的副本不会可靠地更新原 C++ 消息。
\end{quote}

\hypertarget{unitree-sdk2-cpp-idl-hg-doubleimu-doubleimustate-gyroscope}{%
\paragraph{\texorpdfstring{\texttt{unitree\_sdk2\_cpp.idl.hg\_doubleimu.doubleIMUState.gyroscope}}{unitree\_sdk2\_cpp.idl.hg\_doubleimu.doubleIMUState.gyroscope}}\label{unitree-sdk2-cpp-idl-hg-doubleimu-doubleimustate-gyroscope}}

底层
\passthrough{\lstinline!unitree\_hg\_doubleimu::msg::dds\_::doubleIMUState\_!}
的 \passthrough{\lstinline!gyroscope!} 字段。类型签名只说明 Python
表示；单位、数值范围、数组固定长度和枚举含义需查对应 IDL 头文件。
底层是固定长度数组，写入序列必须正好包含 3
个元素；元素约束：底层为浮点数；类型本身不说明单位、坐标系或业务安全范围。

\textbf{签名}

\begin{lstlisting}[language=Python]
@property
def gyroscope(self) -> list[float]

@gyroscope.setter
def gyroscope(self, value: Sequence[float]) -> None
\end{lstlisting}

\textbf{可用性与安全性}

\textbf{\passthrough{\lstinline!AVAILABLE!} /
\passthrough{\lstinline!VALUE\_TYPE!}}：当前绑定源码已实现该消息值操作。

\textbf{参数}

\begin{longtable}[]{@{}
  >{\raggedright\arraybackslash}p{(\columnwidth - 4\tabcolsep) * \real{0.18}}
  >{\raggedright\arraybackslash}p{(\columnwidth - 4\tabcolsep) * \real{0.20}}
  >{\raggedright\arraybackslash}p{(\columnwidth - 4\tabcolsep) * \real{0.62}}@{}}
\toprule
名称 & Python 类型 & 含义 \\
\midrule
\endhead
\passthrough{\lstinline!value!} &
\passthrough{\lstinline!Sequence[float]!} &
要写入或传递的新值，必须符合该参数的 Python 类型以及底层 C++
范围约束。 \\
\bottomrule
\end{longtable}

\textbf{返回值}

\begin{longtable}[]{@{}
  >{\raggedright\arraybackslash}p{(\columnwidth - 4\tabcolsep) * \real{0.18}}
  >{\raggedright\arraybackslash}p{(\columnwidth - 4\tabcolsep) * \real{0.20}}
  >{\raggedright\arraybackslash}p{(\columnwidth - 4\tabcolsep) * \real{0.62}}@{}}
\toprule
位置 & 类型 & 含义 \\
\midrule
\endhead
getter 返回值 & \passthrough{\lstinline!list[float]!} &
返回属性当前值。数组、vector 和嵌套消息采用复制语义。 \\
\bottomrule
\end{longtable}

\textbf{对应 C++}

\begin{itemize}
\tightlist
\item
  类：\passthrough{\lstinline!unitree\_hg\_doubleimu::msg::dds\_::doubleIMUState\_!}
\item
  字段类型：\passthrough{\lstinline!std::array<float, 3>!}
\item
  声明位置：\passthrough{\lstinline!include/unitree/idl/hg\_doubleimu/doubleIMUState\_.hpp:25!}
\end{itemize}

\textbf{用法}

\begin{lstlisting}[language=Python]
current_value = obj.gyroscope
obj.gyroscope = new_value
\end{lstlisting}

\begin{quote}
{[}!TIP{]} 读取容器或嵌套值后应遵循``读取、修改、重新赋值''。只修改
getter 返回的副本不会可靠地更新原 C++ 消息。
\end{quote}

\hypertarget{unitree-sdk2-cpp-idl-hg-doubleimu-doubleimustate-accelerometer}{%
\paragraph{\texorpdfstring{\texttt{unitree\_sdk2\_cpp.idl.hg\_doubleimu.doubleIMUState.accelerometer}}{unitree\_sdk2\_cpp.idl.hg\_doubleimu.doubleIMUState.accelerometer}}\label{unitree-sdk2-cpp-idl-hg-doubleimu-doubleimustate-accelerometer}}

底层
\passthrough{\lstinline!unitree\_hg\_doubleimu::msg::dds\_::doubleIMUState\_!}
的 \passthrough{\lstinline!accelerometer!} 字段。类型签名只说明 Python
表示；单位、数值范围、数组固定长度和枚举含义需查对应 IDL 头文件。
底层是固定长度数组，写入序列必须正好包含 3
个元素；元素约束：底层为浮点数；类型本身不说明单位、坐标系或业务安全范围。

\textbf{签名}

\begin{lstlisting}[language=Python]
@property
def accelerometer(self) -> list[float]

@accelerometer.setter
def accelerometer(self, value: Sequence[float]) -> None
\end{lstlisting}

\textbf{可用性与安全性}

\textbf{\passthrough{\lstinline!AVAILABLE!} /
\passthrough{\lstinline!VALUE\_TYPE!}}：当前绑定源码已实现该消息值操作。

\textbf{参数}

\begin{longtable}[]{@{}
  >{\raggedright\arraybackslash}p{(\columnwidth - 4\tabcolsep) * \real{0.18}}
  >{\raggedright\arraybackslash}p{(\columnwidth - 4\tabcolsep) * \real{0.20}}
  >{\raggedright\arraybackslash}p{(\columnwidth - 4\tabcolsep) * \real{0.62}}@{}}
\toprule
名称 & Python 类型 & 含义 \\
\midrule
\endhead
\passthrough{\lstinline!value!} &
\passthrough{\lstinline!Sequence[float]!} &
要写入或传递的新值，必须符合该参数的 Python 类型以及底层 C++
范围约束。 \\
\bottomrule
\end{longtable}

\textbf{返回值}

\begin{longtable}[]{@{}
  >{\raggedright\arraybackslash}p{(\columnwidth - 4\tabcolsep) * \real{0.18}}
  >{\raggedright\arraybackslash}p{(\columnwidth - 4\tabcolsep) * \real{0.20}}
  >{\raggedright\arraybackslash}p{(\columnwidth - 4\tabcolsep) * \real{0.62}}@{}}
\toprule
位置 & 类型 & 含义 \\
\midrule
\endhead
getter 返回值 & \passthrough{\lstinline!list[float]!} &
返回属性当前值。数组、vector 和嵌套消息采用复制语义。 \\
\bottomrule
\end{longtable}

\textbf{对应 C++}

\begin{itemize}
\tightlist
\item
  类：\passthrough{\lstinline!unitree\_hg\_doubleimu::msg::dds\_::doubleIMUState\_!}
\item
  字段类型：\passthrough{\lstinline!std::array<float, 3>!}
\item
  声明位置：\passthrough{\lstinline!include/unitree/idl/hg\_doubleimu/doubleIMUState\_.hpp:26!}
\end{itemize}

\textbf{用法}

\begin{lstlisting}[language=Python]
current_value = obj.accelerometer
obj.accelerometer = new_value
\end{lstlisting}

\begin{quote}
{[}!TIP{]} 读取容器或嵌套值后应遵循``读取、修改、重新赋值''。只修改
getter 返回的副本不会可靠地更新原 C++ 消息。
\end{quote}

\hypertarget{unitree-sdk2-cpp-idl-hg-doubleimu-doubleimustate-rpy}{%
\paragraph{\texorpdfstring{\texttt{unitree\_sdk2\_cpp.idl.hg\_doubleimu.doubleIMUState.rpy}}{unitree\_sdk2\_cpp.idl.hg\_doubleimu.doubleIMUState.rpy}}\label{unitree-sdk2-cpp-idl-hg-doubleimu-doubleimustate-rpy}}

底层
\passthrough{\lstinline!unitree\_hg\_doubleimu::msg::dds\_::doubleIMUState\_!}
的 \passthrough{\lstinline!rpy!} 字段。类型签名只说明 Python
表示；单位、数值范围、数组固定长度和枚举含义需查对应 IDL 头文件。
底层是固定长度数组，写入序列必须正好包含 3
个元素；元素约束：底层为浮点数；类型本身不说明单位、坐标系或业务安全范围。

\textbf{签名}

\begin{lstlisting}[language=Python]
@property
def rpy(self) -> list[float]

@rpy.setter
def rpy(self, value: Sequence[float]) -> None
\end{lstlisting}

\textbf{可用性与安全性}

\textbf{\passthrough{\lstinline!AVAILABLE!} /
\passthrough{\lstinline!VALUE\_TYPE!}}：当前绑定源码已实现该消息值操作。

\textbf{参数}

\begin{longtable}[]{@{}
  >{\raggedright\arraybackslash}p{(\columnwidth - 4\tabcolsep) * \real{0.18}}
  >{\raggedright\arraybackslash}p{(\columnwidth - 4\tabcolsep) * \real{0.20}}
  >{\raggedright\arraybackslash}p{(\columnwidth - 4\tabcolsep) * \real{0.62}}@{}}
\toprule
名称 & Python 类型 & 含义 \\
\midrule
\endhead
\passthrough{\lstinline!value!} &
\passthrough{\lstinline!Sequence[float]!} &
要写入或传递的新值，必须符合该参数的 Python 类型以及底层 C++
范围约束。 \\
\bottomrule
\end{longtable}

\textbf{返回值}

\begin{longtable}[]{@{}
  >{\raggedright\arraybackslash}p{(\columnwidth - 4\tabcolsep) * \real{0.18}}
  >{\raggedright\arraybackslash}p{(\columnwidth - 4\tabcolsep) * \real{0.20}}
  >{\raggedright\arraybackslash}p{(\columnwidth - 4\tabcolsep) * \real{0.62}}@{}}
\toprule
位置 & 类型 & 含义 \\
\midrule
\endhead
getter 返回值 & \passthrough{\lstinline!list[float]!} &
返回属性当前值。数组、vector 和嵌套消息采用复制语义。 \\
\bottomrule
\end{longtable}

\textbf{对应 C++}

\begin{itemize}
\tightlist
\item
  类：\passthrough{\lstinline!unitree\_hg\_doubleimu::msg::dds\_::doubleIMUState\_!}
\item
  字段类型：\passthrough{\lstinline!std::array<float, 3>!}
\item
  声明位置：\passthrough{\lstinline!include/unitree/idl/hg\_doubleimu/doubleIMUState\_.hpp:27!}
\end{itemize}

\textbf{用法}

\begin{lstlisting}[language=Python]
current_value = obj.rpy
obj.rpy = new_value
\end{lstlisting}

\begin{quote}
{[}!TIP{]} 读取容器或嵌套值后应遵循``读取、修改、重新赋值''。只修改
getter 返回的副本不会可靠地更新原 C++ 消息。
\end{quote}

\hypertarget{unitree-sdk2-cpp-idl-hg-doubleimu-doubleimustate-temperature}{%
\paragraph{\texorpdfstring{\texttt{unitree\_sdk2\_cpp.idl.hg\_doubleimu.doubleIMUState.temperature}}{unitree\_sdk2\_cpp.idl.hg\_doubleimu.doubleIMUState.temperature}}\label{unitree-sdk2-cpp-idl-hg-doubleimu-doubleimustate-temperature}}

底层
\passthrough{\lstinline!unitree\_hg\_doubleimu::msg::dds\_::doubleIMUState\_!}
的 \passthrough{\lstinline!temperature!} 字段。类型签名只说明 Python
表示；单位、数值范围、数组固定长度和枚举含义需查对应 IDL 头文件。
取值范围为 -32768 到 32767。

\textbf{签名}

\begin{lstlisting}[language=Python]
@property
def temperature(self) -> int

@temperature.setter
def temperature(self, value: int) -> None
\end{lstlisting}

\textbf{可用性与安全性}

\textbf{\passthrough{\lstinline!AVAILABLE!} /
\passthrough{\lstinline!VALUE\_TYPE!}}：当前绑定源码已实现该消息值操作。

\textbf{参数}

\begin{longtable}[]{@{}
  >{\raggedright\arraybackslash}p{(\columnwidth - 4\tabcolsep) * \real{0.18}}
  >{\raggedright\arraybackslash}p{(\columnwidth - 4\tabcolsep) * \real{0.20}}
  >{\raggedright\arraybackslash}p{(\columnwidth - 4\tabcolsep) * \real{0.62}}@{}}
\toprule
名称 & Python 类型 & 含义 \\
\midrule
\endhead
\passthrough{\lstinline!value!} & \passthrough{\lstinline!int!} &
要写入或传递的新值，必须符合该参数的 Python 类型以及底层 C++
范围约束。 \\
\bottomrule
\end{longtable}

\textbf{返回值}

\begin{longtable}[]{@{}
  >{\raggedright\arraybackslash}p{(\columnwidth - 4\tabcolsep) * \real{0.18}}
  >{\raggedright\arraybackslash}p{(\columnwidth - 4\tabcolsep) * \real{0.20}}
  >{\raggedright\arraybackslash}p{(\columnwidth - 4\tabcolsep) * \real{0.62}}@{}}
\toprule
位置 & 类型 & 含义 \\
\midrule
\endhead
getter 返回值 & \passthrough{\lstinline!int!} & 返回属性当前值。 \\
\bottomrule
\end{longtable}

\textbf{对应 C++}

\begin{itemize}
\tightlist
\item
  类：\passthrough{\lstinline!unitree\_hg\_doubleimu::msg::dds\_::doubleIMUState\_!}
\item
  字段类型：\passthrough{\lstinline!int16\_t!}
\item
  声明位置：\passthrough{\lstinline!include/unitree/idl/hg\_doubleimu/doubleIMUState\_.hpp:28!}
\end{itemize}

\textbf{用法}

\begin{lstlisting}[language=Python]
current_value = obj.temperature
obj.temperature = new_value
\end{lstlisting}

\hypertarget{unitree-sdk2-cpp-idl-hg-doubleimu-doubleimustate-tick}{%
\paragraph{\texorpdfstring{\texttt{unitree\_sdk2\_cpp.idl.hg\_doubleimu.doubleIMUState.tick}}{unitree\_sdk2\_cpp.idl.hg\_doubleimu.doubleIMUState.tick}}\label{unitree-sdk2-cpp-idl-hg-doubleimu-doubleimustate-tick}}

底层
\passthrough{\lstinline!unitree\_hg\_doubleimu::msg::dds\_::doubleIMUState\_!}
的 \passthrough{\lstinline!tick!} 字段。类型签名只说明 Python
表示；单位、数值范围、数组固定长度和枚举含义需查对应 IDL 头文件。
取值范围为 0 到 4294967295。

\textbf{签名}

\begin{lstlisting}[language=Python]
@property
def tick(self) -> int

@tick.setter
def tick(self, value: int) -> None
\end{lstlisting}

\textbf{可用性与安全性}

\textbf{\passthrough{\lstinline!AVAILABLE!} /
\passthrough{\lstinline!VALUE\_TYPE!}}：当前绑定源码已实现该消息值操作。

\textbf{参数}

\begin{longtable}[]{@{}
  >{\raggedright\arraybackslash}p{(\columnwidth - 4\tabcolsep) * \real{0.18}}
  >{\raggedright\arraybackslash}p{(\columnwidth - 4\tabcolsep) * \real{0.20}}
  >{\raggedright\arraybackslash}p{(\columnwidth - 4\tabcolsep) * \real{0.62}}@{}}
\toprule
名称 & Python 类型 & 含义 \\
\midrule
\endhead
\passthrough{\lstinline!value!} & \passthrough{\lstinline!int!} &
要写入或传递的新值，必须符合该参数的 Python 类型以及底层 C++
范围约束。 \\
\bottomrule
\end{longtable}

\textbf{返回值}

\begin{longtable}[]{@{}
  >{\raggedright\arraybackslash}p{(\columnwidth - 4\tabcolsep) * \real{0.18}}
  >{\raggedright\arraybackslash}p{(\columnwidth - 4\tabcolsep) * \real{0.20}}
  >{\raggedright\arraybackslash}p{(\columnwidth - 4\tabcolsep) * \real{0.62}}@{}}
\toprule
位置 & 类型 & 含义 \\
\midrule
\endhead
getter 返回值 & \passthrough{\lstinline!int!} & 返回属性当前值。 \\
\bottomrule
\end{longtable}

\textbf{对应 C++}

\begin{itemize}
\tightlist
\item
  类：\passthrough{\lstinline!unitree\_hg\_doubleimu::msg::dds\_::doubleIMUState\_!}
\item
  字段类型：\passthrough{\lstinline!uint32\_t!}
\item
  声明位置：\passthrough{\lstinline!include/unitree/idl/hg\_doubleimu/doubleIMUState\_.hpp:29!}
\end{itemize}

\textbf{用法}

\begin{lstlisting}[language=Python]
current_value = obj.tick
obj.tick = new_value
\end{lstlisting}

\begin{center}\rule{0.5\linewidth}{0.5pt}\end{center}

\hypertarget{unitree-sdk2-cpp-idl-ros2}{%
\subsection{ROS2 兼容 IDL 消息}\label{unitree-sdk2-cpp-idl-ros2}}

模块：\passthrough{\lstinline!unitree\_sdk2\_cpp.idl.ros2!}

\hypertarget{ux7c7bux7d22ux5f15-5}{%
\subsubsection{类索引}\label{ux7c7bux7d22ux5f15-5}}

\begin{longtable}[]{@{}
  >{\raggedright\arraybackslash}p{(\columnwidth - 4\tabcolsep) * \real{0.18}}
  >{\raggedleft\arraybackslash}p{(\columnwidth - 4\tabcolsep) * \real{0.20}}
  >{\raggedleft\arraybackslash}p{(\columnwidth - 4\tabcolsep) * \real{0.62}}@{}}
\toprule
类 & 公开函数签名 & 属性 \\
\midrule
\endhead
\protect\hyperlink{unitree-sdk2-cpp-idl-ros2-time}{\passthrough{\lstinline!Time!}}
& 3 & 2 \\
\protect\hyperlink{unitree-sdk2-cpp-idl-ros2-header}{\passthrough{\lstinline!Header!}}
& 3 & 2 \\
\protect\hyperlink{unitree-sdk2-cpp-idl-ros2-quaternion}{\passthrough{\lstinline!Quaternion!}}
& 3 & 4 \\
\protect\hyperlink{unitree-sdk2-cpp-idl-ros2-vector3}{\passthrough{\lstinline!Vector3!}}
& 3 & 3 \\
\protect\hyperlink{unitree-sdk2-cpp-idl-ros2-imu}{\passthrough{\lstinline!Imu!}}
& 3 & 7 \\
\protect\hyperlink{unitree-sdk2-cpp-idl-ros2-point}{\passthrough{\lstinline!Point!}}
& 3 & 3 \\
\protect\hyperlink{unitree-sdk2-cpp-idl-ros2-pose}{\passthrough{\lstinline!Pose!}}
& 3 & 2 \\
\protect\hyperlink{unitree-sdk2-cpp-idl-ros2-mapmetadata}{\passthrough{\lstinline!MapMetaData!}}
& 3 & 5 \\
\protect\hyperlink{unitree-sdk2-cpp-idl-ros2-occupancygrid}{\passthrough{\lstinline!OccupancyGrid!}}
& 3 & 3 \\
\protect\hyperlink{unitree-sdk2-cpp-idl-ros2-posewithcovariance}{\passthrough{\lstinline!PoseWithCovariance!}}
& 3 & 2 \\
\protect\hyperlink{unitree-sdk2-cpp-idl-ros2-twist}{\passthrough{\lstinline!Twist!}}
& 3 & 2 \\
\protect\hyperlink{unitree-sdk2-cpp-idl-ros2-twistwithcovariance}{\passthrough{\lstinline!TwistWithCovariance!}}
& 3 & 2 \\
\protect\hyperlink{unitree-sdk2-cpp-idl-ros2-odometry}{\passthrough{\lstinline!Odometry!}}
& 3 & 4 \\
\protect\hyperlink{unitree-sdk2-cpp-idl-ros2-point32}{\passthrough{\lstinline!Point32!}}
& 3 & 3 \\
\protect\hyperlink{unitree-sdk2-cpp-idl-ros2-pointfield}{\passthrough{\lstinline!PointField!}}
& 3 & 4 \\
\protect\hyperlink{unitree-sdk2-cpp-idl-ros2-pointcloud2}{\passthrough{\lstinline!PointCloud2!}}
& 3 & 9 \\
\protect\hyperlink{unitree-sdk2-cpp-idl-ros2-pointstamped}{\passthrough{\lstinline!PointStamped!}}
& 3 & 2 \\
\protect\hyperlink{unitree-sdk2-cpp-idl-ros2-pose2d}{\passthrough{\lstinline!Pose2D!}}
& 3 & 3 \\
\protect\hyperlink{unitree-sdk2-cpp-idl-ros2-posestamped}{\passthrough{\lstinline!PoseStamped!}}
& 3 & 2 \\
\protect\hyperlink{unitree-sdk2-cpp-idl-ros2-posewithcovariancestamped}{\passthrough{\lstinline!PoseWithCovarianceStamped!}}
& 3 & 2 \\
\protect\hyperlink{unitree-sdk2-cpp-idl-ros2-quaternionstamped}{\passthrough{\lstinline!QuaternionStamped!}}
& 3 & 2 \\
\protect\hyperlink{unitree-sdk2-cpp-idl-ros2-string}{\passthrough{\lstinline!String!}}
& 3 & 1 \\
\protect\hyperlink{unitree-sdk2-cpp-idl-ros2-twiststamped}{\passthrough{\lstinline!TwistStamped!}}
& 3 & 2 \\
\protect\hyperlink{unitree-sdk2-cpp-idl-ros2-twistwithcovariancestamped}{\passthrough{\lstinline!TwistWithCovarianceStamped!}}
& 3 & 2 \\
\bottomrule
\end{longtable}

\hypertarget{unitree-sdk2-cpp-idl-ros2-time}{%
\subsubsection{\texorpdfstring{\texttt{unitree\_sdk2\_cpp.idl.ros2.Time}}{unitree\_sdk2\_cpp.idl.ros2.Time}}\label{unitree-sdk2-cpp-idl-ros2-time}}

可由 Python 读写的 DDS/IDL 消息值类型。字段采用复制语义。

\textbf{导入}

\begin{lstlisting}[language=Python]
from unitree_sdk2_cpp.idl.ros2 import Time
\end{lstlisting}

\textbf{方法索引}

\begin{longtable}[]{@{}
  >{\raggedright\arraybackslash}p{(\columnwidth - 6\tabcolsep) * \real{0.18}}
  >{\raggedright\arraybackslash}p{(\columnwidth - 6\tabcolsep) * \real{0.24}}
  >{\raggedright\arraybackslash}p{(\columnwidth - 6\tabcolsep) * \real{0.18}}
  >{\raggedright\arraybackslash}p{(\columnwidth - 6\tabcolsep) * \real{0.40}}@{}}
\toprule
方法 & Python 签名 & 状态 & 安全分类 \\
\midrule
\endhead
\protect\hyperlink{unitree-sdk2-cpp-idl-ros2-time-dunder-init-1}{\passthrough{\lstinline!\_\_init\_\_!}}
& \passthrough{\lstinline!def \_\_init\_\_(self) -> None!} &
\passthrough{\lstinline!AVAILABLE!} &
\passthrough{\lstinline!VALUE\_TYPE!} \\
\protect\hyperlink{unitree-sdk2-cpp-idl-ros2-time-dunder-eq-1}{\passthrough{\lstinline!\_\_eq\_\_!}}
& \passthrough{\lstinline!def \_\_eq\_\_(self, other: object) -> bool!}
& \passthrough{\lstinline!AVAILABLE!} &
\passthrough{\lstinline!VALUE\_TYPE!} \\
\protect\hyperlink{unitree-sdk2-cpp-idl-ros2-time-dunder-ne-1}{\passthrough{\lstinline!\_\_ne\_\_!}}
& \passthrough{\lstinline!def \_\_ne\_\_(self, other: object) -> bool!}
& \passthrough{\lstinline!AVAILABLE!} &
\passthrough{\lstinline!VALUE\_TYPE!} \\
\bottomrule
\end{longtable}

\textbf{属性索引}

\begin{longtable}[]{@{}
  >{\raggedright\arraybackslash}p{(\columnwidth - 6\tabcolsep) * \real{0.18}}
  >{\raggedright\arraybackslash}p{(\columnwidth - 6\tabcolsep) * \real{0.24}}
  >{\raggedright\arraybackslash}p{(\columnwidth - 6\tabcolsep) * \real{0.18}}
  >{\raggedright\arraybackslash}p{(\columnwidth - 6\tabcolsep) * \real{0.40}}@{}}
\toprule
属性 & 读取类型 & 写入类型 & 状态 \\
\midrule
\endhead
\protect\hyperlink{unitree-sdk2-cpp-idl-ros2-time-sec}{\passthrough{\lstinline!sec!}}
& \passthrough{\lstinline!int!} & \passthrough{\lstinline!int!} &
\passthrough{\lstinline!AVAILABLE!} \\
\protect\hyperlink{unitree-sdk2-cpp-idl-ros2-time-nanosec}{\passthrough{\lstinline!nanosec!}}
& \passthrough{\lstinline!int!} & \passthrough{\lstinline!int!} &
\passthrough{\lstinline!AVAILABLE!} \\
\bottomrule
\end{longtable}

\hypertarget{unitree-sdk2-cpp-idl-ros2-time-dunder-init-1}{%
\paragraph{\texorpdfstring{\texttt{unitree\_sdk2\_cpp.idl.ros2.Time.\_\_init\_\_}}{unitree\_sdk2\_cpp.idl.ros2.Time.\_\_init\_\_}}\label{unitree-sdk2-cpp-idl-ros2-time-dunder-init-1}}

初始化 \passthrough{\lstinline!Time!}
实例。是否能实际构造取决于下方可用性状态。

\textbf{签名}

\begin{lstlisting}[language=Python]
def __init__(self) -> None
\end{lstlisting}

\textbf{可用性与安全性}

\textbf{\passthrough{\lstinline!AVAILABLE!} /
\passthrough{\lstinline!VALUE\_TYPE!}}：当前绑定源码已实现该消息值操作。

\textbf{参数}

无显式参数。实例方法中的 \passthrough{\lstinline!self!} 由 Python
自动传入。

\textbf{返回值}

\begin{longtable}[]{@{}
  >{\raggedright\arraybackslash}p{(\columnwidth - 4\tabcolsep) * \real{0.18}}
  >{\raggedright\arraybackslash}p{(\columnwidth - 4\tabcolsep) * \real{0.20}}
  >{\raggedright\arraybackslash}p{(\columnwidth - 4\tabcolsep) * \real{0.62}}@{}}
\toprule
位置 & 类型 & 含义 \\
\midrule
\endhead
无 & \passthrough{\lstinline!None!} &
\passthrough{\lstinline!\_\_init\_\_!}
本身不返回值；调用类对象时，在构造函数可用的前提下得到该类实例。 \\
\bottomrule
\end{longtable}

\textbf{对应 C++}

\begin{itemize}
\tightlist
\item
  类：\passthrough{\lstinline!builtin\_interfaces::msg::dds\_::Time\_!}
\item
  签名：\passthrough{\lstinline!Time\_()!}
\item
  绑定策略：\passthrough{\lstinline!未单独标注!}
\item
  声明位置：\passthrough{\lstinline!include/unitree/idl/ros2/Time\_.hpp:27!}
\end{itemize}

\textbf{说明}

\begin{itemize}
\tightlist
\item
  示例是局部调用片段；\passthrough{\lstinline!obj!}
  和参数变量表示已经按上层业务校验并准备好的对象。
\end{itemize}

\textbf{用法}

\begin{lstlisting}[language=Python]
value = Time()
\end{lstlisting}

\hypertarget{unitree-sdk2-cpp-idl-ros2-time-dunder-eq-1}{%
\paragraph{\texorpdfstring{\texttt{unitree\_sdk2\_cpp.idl.ros2.Time.\_\_eq\_\_}}{unitree\_sdk2\_cpp.idl.ros2.Time.\_\_eq\_\_}}\label{unitree-sdk2-cpp-idl-ros2-time-dunder-eq-1}}

按底层 IDL 消息内容比较两个对象是否相等。

\textbf{签名}

\begin{lstlisting}[language=Python]
def __eq__(self, other: object) -> bool
\end{lstlisting}

\textbf{可用性与安全性}

\textbf{\passthrough{\lstinline!AVAILABLE!} /
\passthrough{\lstinline!VALUE\_TYPE!}}：当前绑定源码已实现该消息值操作。

\textbf{参数}

\begin{longtable}[]{@{}
  >{\raggedright\arraybackslash}p{(\columnwidth - 6\tabcolsep) * \real{0.18}}
  >{\raggedright\arraybackslash}p{(\columnwidth - 6\tabcolsep) * \real{0.24}}
  >{\raggedright\arraybackslash}p{(\columnwidth - 6\tabcolsep) * \real{0.18}}
  >{\raggedright\arraybackslash}p{(\columnwidth - 6\tabcolsep) * \real{0.40}}@{}}
\toprule
名称 & Python 类型 & 默认值 & 含义 \\
\midrule
\endhead
\passthrough{\lstinline!other!} & \passthrough{\lstinline!object!} &
必填 & 用于比较的另一个 Python 对象。类型不兼容时比较结果通常为
\passthrough{\lstinline!False!}。 \\
\bottomrule
\end{longtable}

\textbf{返回值}

\begin{longtable}[]{@{}
  >{\raggedright\arraybackslash}p{(\columnwidth - 4\tabcolsep) * \real{0.18}}
  >{\raggedright\arraybackslash}p{(\columnwidth - 4\tabcolsep) * \real{0.20}}
  >{\raggedright\arraybackslash}p{(\columnwidth - 4\tabcolsep) * \real{0.62}}@{}}
\toprule
位置 & 类型 & 含义 \\
\midrule
\endhead
返回值 & \passthrough{\lstinline!bool!} & 返回
\passthrough{\lstinline!bool!}。更精确的业务含义见该方法用途、C++
签名和目标型号协议。 \\
\bottomrule
\end{longtable}

\textbf{对应 C++}

\begin{itemize}
\tightlist
\item
  类：\passthrough{\lstinline!builtin\_interfaces::msg::dds\_::Time\_!}
\item
  签名：\passthrough{\lstinline!operator==(const value \&) const!}
\item
  绑定策略：\passthrough{\lstinline!未单独标注!}
\item
  声明位置：由 pybind11 手工绑定或生成注册表提供；manifest
  未记录独立头文件行号。
\end{itemize}

\textbf{说明}

\begin{itemize}
\tightlist
\item
  示例是局部调用片段；\passthrough{\lstinline!obj!}
  和参数变量表示已经按上层业务校验并准备好的对象。
\end{itemize}

\textbf{用法}

\begin{lstlisting}[language=Python]
same = left == right
\end{lstlisting}

\hypertarget{unitree-sdk2-cpp-idl-ros2-time-dunder-ne-1}{%
\paragraph{\texorpdfstring{\texttt{unitree\_sdk2\_cpp.idl.ros2.Time.\_\_ne\_\_}}{unitree\_sdk2\_cpp.idl.ros2.Time.\_\_ne\_\_}}\label{unitree-sdk2-cpp-idl-ros2-time-dunder-ne-1}}

按底层 IDL 消息内容比较两个对象是否不相等。

\textbf{签名}

\begin{lstlisting}[language=Python]
def __ne__(self, other: object) -> bool
\end{lstlisting}

\textbf{可用性与安全性}

\textbf{\passthrough{\lstinline!AVAILABLE!} /
\passthrough{\lstinline!VALUE\_TYPE!}}：当前绑定源码已实现该消息值操作。

\textbf{参数}

\begin{longtable}[]{@{}
  >{\raggedright\arraybackslash}p{(\columnwidth - 6\tabcolsep) * \real{0.18}}
  >{\raggedright\arraybackslash}p{(\columnwidth - 6\tabcolsep) * \real{0.24}}
  >{\raggedright\arraybackslash}p{(\columnwidth - 6\tabcolsep) * \real{0.18}}
  >{\raggedright\arraybackslash}p{(\columnwidth - 6\tabcolsep) * \real{0.40}}@{}}
\toprule
名称 & Python 类型 & 默认值 & 含义 \\
\midrule
\endhead
\passthrough{\lstinline!other!} & \passthrough{\lstinline!object!} &
必填 & 用于比较的另一个 Python 对象。类型不兼容时比较结果通常为
\passthrough{\lstinline!False!}。 \\
\bottomrule
\end{longtable}

\textbf{返回值}

\begin{longtable}[]{@{}
  >{\raggedright\arraybackslash}p{(\columnwidth - 4\tabcolsep) * \real{0.18}}
  >{\raggedright\arraybackslash}p{(\columnwidth - 4\tabcolsep) * \real{0.20}}
  >{\raggedright\arraybackslash}p{(\columnwidth - 4\tabcolsep) * \real{0.62}}@{}}
\toprule
位置 & 类型 & 含义 \\
\midrule
\endhead
返回值 & \passthrough{\lstinline!bool!} & 返回
\passthrough{\lstinline!bool!}。更精确的业务含义见该方法用途、C++
签名和目标型号协议。 \\
\bottomrule
\end{longtable}

\textbf{对应 C++}

\begin{itemize}
\tightlist
\item
  类：\passthrough{\lstinline!builtin\_interfaces::msg::dds\_::Time\_!}
\item
  签名：\passthrough{\lstinline"operator!=(const value \&) const"}
\item
  绑定策略：\passthrough{\lstinline!未单独标注!}
\item
  声明位置：由 pybind11 手工绑定或生成注册表提供；manifest
  未记录独立头文件行号。
\end{itemize}

\textbf{说明}

\begin{itemize}
\tightlist
\item
  示例是局部调用片段；\passthrough{\lstinline!obj!}
  和参数变量表示已经按上层业务校验并准备好的对象。
\end{itemize}

\textbf{用法}

\begin{lstlisting}[language=Python]
different = left != right
\end{lstlisting}

\hypertarget{unitree-sdk2-cpp-idl-ros2-time-sec}{%
\paragraph{\texorpdfstring{\texttt{unitree\_sdk2\_cpp.idl.ros2.Time.sec}}{unitree\_sdk2\_cpp.idl.ros2.Time.sec}}\label{unitree-sdk2-cpp-idl-ros2-time-sec}}

底层 \passthrough{\lstinline!builtin\_interfaces::msg::dds\_::Time\_!}
的 \passthrough{\lstinline!sec!} 字段。类型签名只说明 Python
表示；单位、数值范围、数组固定长度和枚举含义需查对应 IDL 头文件。
取值范围为 -2147483648 到 2147483647。

\textbf{签名}

\begin{lstlisting}[language=Python]
@property
def sec(self) -> int

@sec.setter
def sec(self, value: int) -> None
\end{lstlisting}

\textbf{可用性与安全性}

\textbf{\passthrough{\lstinline!AVAILABLE!} /
\passthrough{\lstinline!VALUE\_TYPE!}}：当前绑定源码已实现该消息值操作。

\textbf{参数}

\begin{longtable}[]{@{}
  >{\raggedright\arraybackslash}p{(\columnwidth - 4\tabcolsep) * \real{0.18}}
  >{\raggedright\arraybackslash}p{(\columnwidth - 4\tabcolsep) * \real{0.20}}
  >{\raggedright\arraybackslash}p{(\columnwidth - 4\tabcolsep) * \real{0.62}}@{}}
\toprule
名称 & Python 类型 & 含义 \\
\midrule
\endhead
\passthrough{\lstinline!value!} & \passthrough{\lstinline!int!} &
要写入或传递的新值，必须符合该参数的 Python 类型以及底层 C++
范围约束。 \\
\bottomrule
\end{longtable}

\textbf{返回值}

\begin{longtable}[]{@{}
  >{\raggedright\arraybackslash}p{(\columnwidth - 4\tabcolsep) * \real{0.18}}
  >{\raggedright\arraybackslash}p{(\columnwidth - 4\tabcolsep) * \real{0.20}}
  >{\raggedright\arraybackslash}p{(\columnwidth - 4\tabcolsep) * \real{0.62}}@{}}
\toprule
位置 & 类型 & 含义 \\
\midrule
\endhead
getter 返回值 & \passthrough{\lstinline!int!} & 返回属性当前值。 \\
\bottomrule
\end{longtable}

\textbf{对应 C++}

\begin{itemize}
\tightlist
\item
  类：\passthrough{\lstinline!builtin\_interfaces::msg::dds\_::Time\_!}
\item
  字段类型：\passthrough{\lstinline!int32\_t!}
\item
  声明位置：\passthrough{\lstinline!include/unitree/idl/ros2/Time\_.hpp:23!}
\end{itemize}

\textbf{用法}

\begin{lstlisting}[language=Python]
current_value = obj.sec
obj.sec = new_value
\end{lstlisting}

\hypertarget{unitree-sdk2-cpp-idl-ros2-time-nanosec}{%
\paragraph{\texorpdfstring{\texttt{unitree\_sdk2\_cpp.idl.ros2.Time.nanosec}}{unitree\_sdk2\_cpp.idl.ros2.Time.nanosec}}\label{unitree-sdk2-cpp-idl-ros2-time-nanosec}}

底层 \passthrough{\lstinline!builtin\_interfaces::msg::dds\_::Time\_!}
的 \passthrough{\lstinline!nanosec!} 字段。类型签名只说明 Python
表示；单位、数值范围、数组固定长度和枚举含义需查对应 IDL 头文件。
取值范围为 0 到 4294967295。

\textbf{签名}

\begin{lstlisting}[language=Python]
@property
def nanosec(self) -> int

@nanosec.setter
def nanosec(self, value: int) -> None
\end{lstlisting}

\textbf{可用性与安全性}

\textbf{\passthrough{\lstinline!AVAILABLE!} /
\passthrough{\lstinline!VALUE\_TYPE!}}：当前绑定源码已实现该消息值操作。

\textbf{参数}

\begin{longtable}[]{@{}
  >{\raggedright\arraybackslash}p{(\columnwidth - 4\tabcolsep) * \real{0.18}}
  >{\raggedright\arraybackslash}p{(\columnwidth - 4\tabcolsep) * \real{0.20}}
  >{\raggedright\arraybackslash}p{(\columnwidth - 4\tabcolsep) * \real{0.62}}@{}}
\toprule
名称 & Python 类型 & 含义 \\
\midrule
\endhead
\passthrough{\lstinline!value!} & \passthrough{\lstinline!int!} &
要写入或传递的新值，必须符合该参数的 Python 类型以及底层 C++
范围约束。 \\
\bottomrule
\end{longtable}

\textbf{返回值}

\begin{longtable}[]{@{}
  >{\raggedright\arraybackslash}p{(\columnwidth - 4\tabcolsep) * \real{0.18}}
  >{\raggedright\arraybackslash}p{(\columnwidth - 4\tabcolsep) * \real{0.20}}
  >{\raggedright\arraybackslash}p{(\columnwidth - 4\tabcolsep) * \real{0.62}}@{}}
\toprule
位置 & 类型 & 含义 \\
\midrule
\endhead
getter 返回值 & \passthrough{\lstinline!int!} & 返回属性当前值。 \\
\bottomrule
\end{longtable}

\textbf{对应 C++}

\begin{itemize}
\tightlist
\item
  类：\passthrough{\lstinline!builtin\_interfaces::msg::dds\_::Time\_!}
\item
  字段类型：\passthrough{\lstinline!uint32\_t!}
\item
  声明位置：\passthrough{\lstinline!include/unitree/idl/ros2/Time\_.hpp:24!}
\end{itemize}

\textbf{用法}

\begin{lstlisting}[language=Python]
current_value = obj.nanosec
obj.nanosec = new_value
\end{lstlisting}

\hypertarget{unitree-sdk2-cpp-idl-ros2-header}{%
\subsubsection{\texorpdfstring{\texttt{unitree\_sdk2\_cpp.idl.ros2.Header}}{unitree\_sdk2\_cpp.idl.ros2.Header}}\label{unitree-sdk2-cpp-idl-ros2-header}}

可由 Python 读写的 DDS/IDL 消息值类型。字段采用复制语义。

\textbf{导入}

\begin{lstlisting}[language=Python]
from unitree_sdk2_cpp.idl.ros2 import Header
\end{lstlisting}

\textbf{方法索引}

\begin{longtable}[]{@{}
  >{\raggedright\arraybackslash}p{(\columnwidth - 6\tabcolsep) * \real{0.18}}
  >{\raggedright\arraybackslash}p{(\columnwidth - 6\tabcolsep) * \real{0.24}}
  >{\raggedright\arraybackslash}p{(\columnwidth - 6\tabcolsep) * \real{0.18}}
  >{\raggedright\arraybackslash}p{(\columnwidth - 6\tabcolsep) * \real{0.40}}@{}}
\toprule
方法 & Python 签名 & 状态 & 安全分类 \\
\midrule
\endhead
\protect\hyperlink{unitree-sdk2-cpp-idl-ros2-header-dunder-init-1}{\passthrough{\lstinline!\_\_init\_\_!}}
& \passthrough{\lstinline!def \_\_init\_\_(self) -> None!} &
\passthrough{\lstinline!AVAILABLE!} &
\passthrough{\lstinline!VALUE\_TYPE!} \\
\protect\hyperlink{unitree-sdk2-cpp-idl-ros2-header-dunder-eq-1}{\passthrough{\lstinline!\_\_eq\_\_!}}
& \passthrough{\lstinline!def \_\_eq\_\_(self, other: object) -> bool!}
& \passthrough{\lstinline!AVAILABLE!} &
\passthrough{\lstinline!VALUE\_TYPE!} \\
\protect\hyperlink{unitree-sdk2-cpp-idl-ros2-header-dunder-ne-1}{\passthrough{\lstinline!\_\_ne\_\_!}}
& \passthrough{\lstinline!def \_\_ne\_\_(self, other: object) -> bool!}
& \passthrough{\lstinline!AVAILABLE!} &
\passthrough{\lstinline!VALUE\_TYPE!} \\
\bottomrule
\end{longtable}

\textbf{属性索引}

\begin{longtable}[]{@{}
  >{\raggedright\arraybackslash}p{(\columnwidth - 6\tabcolsep) * \real{0.18}}
  >{\raggedright\arraybackslash}p{(\columnwidth - 6\tabcolsep) * \real{0.24}}
  >{\raggedright\arraybackslash}p{(\columnwidth - 6\tabcolsep) * \real{0.18}}
  >{\raggedright\arraybackslash}p{(\columnwidth - 6\tabcolsep) * \real{0.40}}@{}}
\toprule
属性 & 读取类型 & 写入类型 & 状态 \\
\midrule
\endhead
\protect\hyperlink{unitree-sdk2-cpp-idl-ros2-header-stamp}{\passthrough{\lstinline!stamp!}}
& \passthrough{\lstinline!Any!} & \passthrough{\lstinline!Any!} &
\passthrough{\lstinline!AVAILABLE!} \\
\protect\hyperlink{unitree-sdk2-cpp-idl-ros2-header-frame-id}{\passthrough{\lstinline!frame\_id!}}
& \passthrough{\lstinline!str!} & \passthrough{\lstinline!str!} &
\passthrough{\lstinline!AVAILABLE!} \\
\bottomrule
\end{longtable}

\hypertarget{unitree-sdk2-cpp-idl-ros2-header-dunder-init-1}{%
\paragraph{\texorpdfstring{\texttt{unitree\_sdk2\_cpp.idl.ros2.Header.\_\_init\_\_}}{unitree\_sdk2\_cpp.idl.ros2.Header.\_\_init\_\_}}\label{unitree-sdk2-cpp-idl-ros2-header-dunder-init-1}}

初始化 \passthrough{\lstinline!Header!}
实例。是否能实际构造取决于下方可用性状态。

\textbf{签名}

\begin{lstlisting}[language=Python]
def __init__(self) -> None
\end{lstlisting}

\textbf{可用性与安全性}

\textbf{\passthrough{\lstinline!AVAILABLE!} /
\passthrough{\lstinline!VALUE\_TYPE!}}：当前绑定源码已实现该消息值操作。

\textbf{参数}

无显式参数。实例方法中的 \passthrough{\lstinline!self!} 由 Python
自动传入。

\textbf{返回值}

\begin{longtable}[]{@{}
  >{\raggedright\arraybackslash}p{(\columnwidth - 4\tabcolsep) * \real{0.18}}
  >{\raggedright\arraybackslash}p{(\columnwidth - 4\tabcolsep) * \real{0.20}}
  >{\raggedright\arraybackslash}p{(\columnwidth - 4\tabcolsep) * \real{0.62}}@{}}
\toprule
位置 & 类型 & 含义 \\
\midrule
\endhead
无 & \passthrough{\lstinline!None!} &
\passthrough{\lstinline!\_\_init\_\_!}
本身不返回值；调用类对象时，在构造函数可用的前提下得到该类实例。 \\
\bottomrule
\end{longtable}

\textbf{对应 C++}

\begin{itemize}
\tightlist
\item
  类：\passthrough{\lstinline!std\_msgs::msg::dds\_::Header\_!}
\item
  签名：\passthrough{\lstinline!Header\_()!}
\item
  绑定策略：\passthrough{\lstinline!未单独标注!}
\item
  声明位置：\passthrough{\lstinline!include/unitree/idl/ros2/Header\_.hpp:29!}
\end{itemize}

\textbf{说明}

\begin{itemize}
\tightlist
\item
  示例是局部调用片段；\passthrough{\lstinline!obj!}
  和参数变量表示已经按上层业务校验并准备好的对象。
\end{itemize}

\textbf{用法}

\begin{lstlisting}[language=Python]
value = Header()
\end{lstlisting}

\hypertarget{unitree-sdk2-cpp-idl-ros2-header-dunder-eq-1}{%
\paragraph{\texorpdfstring{\texttt{unitree\_sdk2\_cpp.idl.ros2.Header.\_\_eq\_\_}}{unitree\_sdk2\_cpp.idl.ros2.Header.\_\_eq\_\_}}\label{unitree-sdk2-cpp-idl-ros2-header-dunder-eq-1}}

按底层 IDL 消息内容比较两个对象是否相等。

\textbf{签名}

\begin{lstlisting}[language=Python]
def __eq__(self, other: object) -> bool
\end{lstlisting}

\textbf{可用性与安全性}

\textbf{\passthrough{\lstinline!AVAILABLE!} /
\passthrough{\lstinline!VALUE\_TYPE!}}：当前绑定源码已实现该消息值操作。

\textbf{参数}

\begin{longtable}[]{@{}
  >{\raggedright\arraybackslash}p{(\columnwidth - 6\tabcolsep) * \real{0.18}}
  >{\raggedright\arraybackslash}p{(\columnwidth - 6\tabcolsep) * \real{0.24}}
  >{\raggedright\arraybackslash}p{(\columnwidth - 6\tabcolsep) * \real{0.18}}
  >{\raggedright\arraybackslash}p{(\columnwidth - 6\tabcolsep) * \real{0.40}}@{}}
\toprule
名称 & Python 类型 & 默认值 & 含义 \\
\midrule
\endhead
\passthrough{\lstinline!other!} & \passthrough{\lstinline!object!} &
必填 & 用于比较的另一个 Python 对象。类型不兼容时比较结果通常为
\passthrough{\lstinline!False!}。 \\
\bottomrule
\end{longtable}

\textbf{返回值}

\begin{longtable}[]{@{}
  >{\raggedright\arraybackslash}p{(\columnwidth - 4\tabcolsep) * \real{0.18}}
  >{\raggedright\arraybackslash}p{(\columnwidth - 4\tabcolsep) * \real{0.20}}
  >{\raggedright\arraybackslash}p{(\columnwidth - 4\tabcolsep) * \real{0.62}}@{}}
\toprule
位置 & 类型 & 含义 \\
\midrule
\endhead
返回值 & \passthrough{\lstinline!bool!} & 返回
\passthrough{\lstinline!bool!}。更精确的业务含义见该方法用途、C++
签名和目标型号协议。 \\
\bottomrule
\end{longtable}

\textbf{对应 C++}

\begin{itemize}
\tightlist
\item
  类：\passthrough{\lstinline!std\_msgs::msg::dds\_::Header\_!}
\item
  签名：\passthrough{\lstinline!operator==(const value \&) const!}
\item
  绑定策略：\passthrough{\lstinline!未单独标注!}
\item
  声明位置：由 pybind11 手工绑定或生成注册表提供；manifest
  未记录独立头文件行号。
\end{itemize}

\textbf{说明}

\begin{itemize}
\tightlist
\item
  示例是局部调用片段；\passthrough{\lstinline!obj!}
  和参数变量表示已经按上层业务校验并准备好的对象。
\end{itemize}

\textbf{用法}

\begin{lstlisting}[language=Python]
same = left == right
\end{lstlisting}

\hypertarget{unitree-sdk2-cpp-idl-ros2-header-dunder-ne-1}{%
\paragraph{\texorpdfstring{\texttt{unitree\_sdk2\_cpp.idl.ros2.Header.\_\_ne\_\_}}{unitree\_sdk2\_cpp.idl.ros2.Header.\_\_ne\_\_}}\label{unitree-sdk2-cpp-idl-ros2-header-dunder-ne-1}}

按底层 IDL 消息内容比较两个对象是否不相等。

\textbf{签名}

\begin{lstlisting}[language=Python]
def __ne__(self, other: object) -> bool
\end{lstlisting}

\textbf{可用性与安全性}

\textbf{\passthrough{\lstinline!AVAILABLE!} /
\passthrough{\lstinline!VALUE\_TYPE!}}：当前绑定源码已实现该消息值操作。

\textbf{参数}

\begin{longtable}[]{@{}
  >{\raggedright\arraybackslash}p{(\columnwidth - 6\tabcolsep) * \real{0.18}}
  >{\raggedright\arraybackslash}p{(\columnwidth - 6\tabcolsep) * \real{0.24}}
  >{\raggedright\arraybackslash}p{(\columnwidth - 6\tabcolsep) * \real{0.18}}
  >{\raggedright\arraybackslash}p{(\columnwidth - 6\tabcolsep) * \real{0.40}}@{}}
\toprule
名称 & Python 类型 & 默认值 & 含义 \\
\midrule
\endhead
\passthrough{\lstinline!other!} & \passthrough{\lstinline!object!} &
必填 & 用于比较的另一个 Python 对象。类型不兼容时比较结果通常为
\passthrough{\lstinline!False!}。 \\
\bottomrule
\end{longtable}

\textbf{返回值}

\begin{longtable}[]{@{}
  >{\raggedright\arraybackslash}p{(\columnwidth - 4\tabcolsep) * \real{0.18}}
  >{\raggedright\arraybackslash}p{(\columnwidth - 4\tabcolsep) * \real{0.20}}
  >{\raggedright\arraybackslash}p{(\columnwidth - 4\tabcolsep) * \real{0.62}}@{}}
\toprule
位置 & 类型 & 含义 \\
\midrule
\endhead
返回值 & \passthrough{\lstinline!bool!} & 返回
\passthrough{\lstinline!bool!}。更精确的业务含义见该方法用途、C++
签名和目标型号协议。 \\
\bottomrule
\end{longtable}

\textbf{对应 C++}

\begin{itemize}
\tightlist
\item
  类：\passthrough{\lstinline!std\_msgs::msg::dds\_::Header\_!}
\item
  签名：\passthrough{\lstinline"operator!=(const value \&) const"}
\item
  绑定策略：\passthrough{\lstinline!未单独标注!}
\item
  声明位置：由 pybind11 手工绑定或生成注册表提供；manifest
  未记录独立头文件行号。
\end{itemize}

\textbf{说明}

\begin{itemize}
\tightlist
\item
  示例是局部调用片段；\passthrough{\lstinline!obj!}
  和参数变量表示已经按上层业务校验并准备好的对象。
\end{itemize}

\textbf{用法}

\begin{lstlisting}[language=Python]
different = left != right
\end{lstlisting}

\hypertarget{unitree-sdk2-cpp-idl-ros2-header-stamp}{%
\paragraph{\texorpdfstring{\texttt{unitree\_sdk2\_cpp.idl.ros2.Header.stamp}}{unitree\_sdk2\_cpp.idl.ros2.Header.stamp}}\label{unitree-sdk2-cpp-idl-ros2-header-stamp}}

底层 \passthrough{\lstinline!std\_msgs::msg::dds\_::Header\_!} 的
\passthrough{\lstinline!stamp!} 字段。类型签名只说明 Python
表示；单位、数值范围、数组固定长度和枚举含义需查对应 IDL 头文件。

\textbf{签名}

\begin{lstlisting}[language=Python]
@property
def stamp(self) -> Any

@stamp.setter
def stamp(self, value: Any) -> None
\end{lstlisting}

\textbf{可用性与安全性}

\textbf{\passthrough{\lstinline!AVAILABLE!} /
\passthrough{\lstinline!VALUE\_TYPE!}}：当前绑定源码已实现该消息值操作。

\textbf{参数}

\begin{longtable}[]{@{}
  >{\raggedright\arraybackslash}p{(\columnwidth - 4\tabcolsep) * \real{0.18}}
  >{\raggedright\arraybackslash}p{(\columnwidth - 4\tabcolsep) * \real{0.20}}
  >{\raggedright\arraybackslash}p{(\columnwidth - 4\tabcolsep) * \real{0.62}}@{}}
\toprule
名称 & Python 类型 & 含义 \\
\midrule
\endhead
\passthrough{\lstinline!value!} & \passthrough{\lstinline!Any!} &
要写入或传递的新值，必须符合该参数的 Python 类型以及底层 C++
范围约束。 \\
\bottomrule
\end{longtable}

\textbf{返回值}

\begin{longtable}[]{@{}
  >{\raggedright\arraybackslash}p{(\columnwidth - 4\tabcolsep) * \real{0.18}}
  >{\raggedright\arraybackslash}p{(\columnwidth - 4\tabcolsep) * \real{0.20}}
  >{\raggedright\arraybackslash}p{(\columnwidth - 4\tabcolsep) * \real{0.62}}@{}}
\toprule
位置 & 类型 & 含义 \\
\midrule
\endhead
getter 返回值 & \passthrough{\lstinline!Any!} & 返回属性当前值。 \\
\bottomrule
\end{longtable}

\textbf{对应 C++}

\begin{itemize}
\tightlist
\item
  类：\passthrough{\lstinline!std\_msgs::msg::dds\_::Header\_!}
\item
  字段类型：\passthrough{\lstinline!::builtin\_interfaces::msg::dds\_::Time\_!}
\item
  声明位置：\passthrough{\lstinline!include/unitree/idl/ros2/Header\_.hpp:25!}
\end{itemize}

\textbf{用法}

\begin{lstlisting}[language=Python]
current_value = obj.stamp
obj.stamp = new_value
\end{lstlisting}

\hypertarget{unitree-sdk2-cpp-idl-ros2-header-frame-id}{%
\paragraph{\texorpdfstring{\texttt{unitree\_sdk2\_cpp.idl.ros2.Header.frame\_id}}{unitree\_sdk2\_cpp.idl.ros2.Header.frame\_id}}\label{unitree-sdk2-cpp-idl-ros2-header-frame-id}}

底层 \passthrough{\lstinline!std\_msgs::msg::dds\_::Header\_!} 的
\passthrough{\lstinline!frame\_id!} 字段。类型签名只说明 Python
表示；单位、数值范围、数组固定长度和枚举含义需查对应 IDL 头文件。
底层为字符串；长度、编码和允许值由具体协议决定。

\textbf{签名}

\begin{lstlisting}[language=Python]
@property
def frame_id(self) -> str

@frame_id.setter
def frame_id(self, value: str) -> None
\end{lstlisting}

\textbf{可用性与安全性}

\textbf{\passthrough{\lstinline!AVAILABLE!} /
\passthrough{\lstinline!VALUE\_TYPE!}}：当前绑定源码已实现该消息值操作。

\textbf{参数}

\begin{longtable}[]{@{}
  >{\raggedright\arraybackslash}p{(\columnwidth - 4\tabcolsep) * \real{0.18}}
  >{\raggedright\arraybackslash}p{(\columnwidth - 4\tabcolsep) * \real{0.20}}
  >{\raggedright\arraybackslash}p{(\columnwidth - 4\tabcolsep) * \real{0.62}}@{}}
\toprule
名称 & Python 类型 & 含义 \\
\midrule
\endhead
\passthrough{\lstinline!value!} & \passthrough{\lstinline!str!} &
要写入或传递的新值，必须符合该参数的 Python 类型以及底层 C++
范围约束。 \\
\bottomrule
\end{longtable}

\textbf{返回值}

\begin{longtable}[]{@{}
  >{\raggedright\arraybackslash}p{(\columnwidth - 4\tabcolsep) * \real{0.18}}
  >{\raggedright\arraybackslash}p{(\columnwidth - 4\tabcolsep) * \real{0.20}}
  >{\raggedright\arraybackslash}p{(\columnwidth - 4\tabcolsep) * \real{0.62}}@{}}
\toprule
位置 & 类型 & 含义 \\
\midrule
\endhead
getter 返回值 & \passthrough{\lstinline!str!} & 返回属性当前值。 \\
\bottomrule
\end{longtable}

\textbf{对应 C++}

\begin{itemize}
\tightlist
\item
  类：\passthrough{\lstinline!std\_msgs::msg::dds\_::Header\_!}
\item
  字段类型：\passthrough{\lstinline!std::string!}
\item
  声明位置：\passthrough{\lstinline!include/unitree/idl/ros2/Header\_.hpp:26!}
\end{itemize}

\textbf{用法}

\begin{lstlisting}[language=Python]
current_value = obj.frame_id
obj.frame_id = new_value
\end{lstlisting}

\hypertarget{unitree-sdk2-cpp-idl-ros2-quaternion}{%
\subsubsection{\texorpdfstring{\texttt{unitree\_sdk2\_cpp.idl.ros2.Quaternion}}{unitree\_sdk2\_cpp.idl.ros2.Quaternion}}\label{unitree-sdk2-cpp-idl-ros2-quaternion}}

可由 Python 读写的 DDS/IDL 消息值类型。字段采用复制语义。

\textbf{导入}

\begin{lstlisting}[language=Python]
from unitree_sdk2_cpp.idl.ros2 import Quaternion
\end{lstlisting}

\textbf{方法索引}

\begin{longtable}[]{@{}
  >{\raggedright\arraybackslash}p{(\columnwidth - 6\tabcolsep) * \real{0.18}}
  >{\raggedright\arraybackslash}p{(\columnwidth - 6\tabcolsep) * \real{0.24}}
  >{\raggedright\arraybackslash}p{(\columnwidth - 6\tabcolsep) * \real{0.18}}
  >{\raggedright\arraybackslash}p{(\columnwidth - 6\tabcolsep) * \real{0.40}}@{}}
\toprule
方法 & Python 签名 & 状态 & 安全分类 \\
\midrule
\endhead
\protect\hyperlink{unitree-sdk2-cpp-idl-ros2-quaternion-dunder-init-1}{\passthrough{\lstinline!\_\_init\_\_!}}
& \passthrough{\lstinline!def \_\_init\_\_(self) -> None!} &
\passthrough{\lstinline!AVAILABLE!} &
\passthrough{\lstinline!VALUE\_TYPE!} \\
\protect\hyperlink{unitree-sdk2-cpp-idl-ros2-quaternion-dunder-eq-1}{\passthrough{\lstinline!\_\_eq\_\_!}}
& \passthrough{\lstinline!def \_\_eq\_\_(self, other: object) -> bool!}
& \passthrough{\lstinline!AVAILABLE!} &
\passthrough{\lstinline!VALUE\_TYPE!} \\
\protect\hyperlink{unitree-sdk2-cpp-idl-ros2-quaternion-dunder-ne-1}{\passthrough{\lstinline!\_\_ne\_\_!}}
& \passthrough{\lstinline!def \_\_ne\_\_(self, other: object) -> bool!}
& \passthrough{\lstinline!AVAILABLE!} &
\passthrough{\lstinline!VALUE\_TYPE!} \\
\bottomrule
\end{longtable}

\textbf{属性索引}

\begin{longtable}[]{@{}
  >{\raggedright\arraybackslash}p{(\columnwidth - 6\tabcolsep) * \real{0.18}}
  >{\raggedright\arraybackslash}p{(\columnwidth - 6\tabcolsep) * \real{0.24}}
  >{\raggedright\arraybackslash}p{(\columnwidth - 6\tabcolsep) * \real{0.18}}
  >{\raggedright\arraybackslash}p{(\columnwidth - 6\tabcolsep) * \real{0.40}}@{}}
\toprule
属性 & 读取类型 & 写入类型 & 状态 \\
\midrule
\endhead
\protect\hyperlink{unitree-sdk2-cpp-idl-ros2-quaternion-x}{\passthrough{\lstinline!x!}}
& \passthrough{\lstinline!float!} & \passthrough{\lstinline!float!} &
\passthrough{\lstinline!AVAILABLE!} \\
\protect\hyperlink{unitree-sdk2-cpp-idl-ros2-quaternion-y}{\passthrough{\lstinline!y!}}
& \passthrough{\lstinline!float!} & \passthrough{\lstinline!float!} &
\passthrough{\lstinline!AVAILABLE!} \\
\protect\hyperlink{unitree-sdk2-cpp-idl-ros2-quaternion-z}{\passthrough{\lstinline!z!}}
& \passthrough{\lstinline!float!} & \passthrough{\lstinline!float!} &
\passthrough{\lstinline!AVAILABLE!} \\
\protect\hyperlink{unitree-sdk2-cpp-idl-ros2-quaternion-w}{\passthrough{\lstinline!w!}}
& \passthrough{\lstinline!float!} & \passthrough{\lstinline!float!} &
\passthrough{\lstinline!AVAILABLE!} \\
\bottomrule
\end{longtable}

\hypertarget{unitree-sdk2-cpp-idl-ros2-quaternion-dunder-init-1}{%
\paragraph{\texorpdfstring{\texttt{unitree\_sdk2\_cpp.idl.ros2.Quaternion.\_\_init\_\_}}{unitree\_sdk2\_cpp.idl.ros2.Quaternion.\_\_init\_\_}}\label{unitree-sdk2-cpp-idl-ros2-quaternion-dunder-init-1}}

初始化 \passthrough{\lstinline!Quaternion!}
实例。是否能实际构造取决于下方可用性状态。

\textbf{签名}

\begin{lstlisting}[language=Python]
def __init__(self) -> None
\end{lstlisting}

\textbf{可用性与安全性}

\textbf{\passthrough{\lstinline!AVAILABLE!} /
\passthrough{\lstinline!VALUE\_TYPE!}}：当前绑定源码已实现该消息值操作。

\textbf{参数}

无显式参数。实例方法中的 \passthrough{\lstinline!self!} 由 Python
自动传入。

\textbf{返回值}

\begin{longtable}[]{@{}
  >{\raggedright\arraybackslash}p{(\columnwidth - 4\tabcolsep) * \real{0.18}}
  >{\raggedright\arraybackslash}p{(\columnwidth - 4\tabcolsep) * \real{0.20}}
  >{\raggedright\arraybackslash}p{(\columnwidth - 4\tabcolsep) * \real{0.62}}@{}}
\toprule
位置 & 类型 & 含义 \\
\midrule
\endhead
无 & \passthrough{\lstinline!None!} &
\passthrough{\lstinline!\_\_init\_\_!}
本身不返回值；调用类对象时，在构造函数可用的前提下得到该类实例。 \\
\bottomrule
\end{longtable}

\textbf{对应 C++}

\begin{itemize}
\tightlist
\item
  类：\passthrough{\lstinline!geometry\_msgs::msg::dds\_::Quaternion\_!}
\item
  签名：\passthrough{\lstinline!Quaternion\_()!}
\item
  绑定策略：\passthrough{\lstinline!未单独标注!}
\item
  声明位置：\passthrough{\lstinline!include/unitree/idl/ros2/Quaternion\_.hpp:28!}
\end{itemize}

\textbf{说明}

\begin{itemize}
\tightlist
\item
  示例是局部调用片段；\passthrough{\lstinline!obj!}
  和参数变量表示已经按上层业务校验并准备好的对象。
\end{itemize}

\textbf{用法}

\begin{lstlisting}[language=Python]
value = Quaternion()
\end{lstlisting}

\hypertarget{unitree-sdk2-cpp-idl-ros2-quaternion-dunder-eq-1}{%
\paragraph{\texorpdfstring{\texttt{unitree\_sdk2\_cpp.idl.ros2.Quaternion.\_\_eq\_\_}}{unitree\_sdk2\_cpp.idl.ros2.Quaternion.\_\_eq\_\_}}\label{unitree-sdk2-cpp-idl-ros2-quaternion-dunder-eq-1}}

按底层 IDL 消息内容比较两个对象是否相等。

\textbf{签名}

\begin{lstlisting}[language=Python]
def __eq__(self, other: object) -> bool
\end{lstlisting}

\textbf{可用性与安全性}

\textbf{\passthrough{\lstinline!AVAILABLE!} /
\passthrough{\lstinline!VALUE\_TYPE!}}：当前绑定源码已实现该消息值操作。

\textbf{参数}

\begin{longtable}[]{@{}
  >{\raggedright\arraybackslash}p{(\columnwidth - 6\tabcolsep) * \real{0.18}}
  >{\raggedright\arraybackslash}p{(\columnwidth - 6\tabcolsep) * \real{0.24}}
  >{\raggedright\arraybackslash}p{(\columnwidth - 6\tabcolsep) * \real{0.18}}
  >{\raggedright\arraybackslash}p{(\columnwidth - 6\tabcolsep) * \real{0.40}}@{}}
\toprule
名称 & Python 类型 & 默认值 & 含义 \\
\midrule
\endhead
\passthrough{\lstinline!other!} & \passthrough{\lstinline!object!} &
必填 & 用于比较的另一个 Python 对象。类型不兼容时比较结果通常为
\passthrough{\lstinline!False!}。 \\
\bottomrule
\end{longtable}

\textbf{返回值}

\begin{longtable}[]{@{}
  >{\raggedright\arraybackslash}p{(\columnwidth - 4\tabcolsep) * \real{0.18}}
  >{\raggedright\arraybackslash}p{(\columnwidth - 4\tabcolsep) * \real{0.20}}
  >{\raggedright\arraybackslash}p{(\columnwidth - 4\tabcolsep) * \real{0.62}}@{}}
\toprule
位置 & 类型 & 含义 \\
\midrule
\endhead
返回值 & \passthrough{\lstinline!bool!} & 返回
\passthrough{\lstinline!bool!}。更精确的业务含义见该方法用途、C++
签名和目标型号协议。 \\
\bottomrule
\end{longtable}

\textbf{对应 C++}

\begin{itemize}
\tightlist
\item
  类：\passthrough{\lstinline!geometry\_msgs::msg::dds\_::Quaternion\_!}
\item
  签名：\passthrough{\lstinline!operator==(const value \&) const!}
\item
  绑定策略：\passthrough{\lstinline!未单独标注!}
\item
  声明位置：由 pybind11 手工绑定或生成注册表提供；manifest
  未记录独立头文件行号。
\end{itemize}

\textbf{说明}

\begin{itemize}
\tightlist
\item
  示例是局部调用片段；\passthrough{\lstinline!obj!}
  和参数变量表示已经按上层业务校验并准备好的对象。
\end{itemize}

\textbf{用法}

\begin{lstlisting}[language=Python]
same = left == right
\end{lstlisting}

\hypertarget{unitree-sdk2-cpp-idl-ros2-quaternion-dunder-ne-1}{%
\paragraph{\texorpdfstring{\texttt{unitree\_sdk2\_cpp.idl.ros2.Quaternion.\_\_ne\_\_}}{unitree\_sdk2\_cpp.idl.ros2.Quaternion.\_\_ne\_\_}}\label{unitree-sdk2-cpp-idl-ros2-quaternion-dunder-ne-1}}

按底层 IDL 消息内容比较两个对象是否不相等。

\textbf{签名}

\begin{lstlisting}[language=Python]
def __ne__(self, other: object) -> bool
\end{lstlisting}

\textbf{可用性与安全性}

\textbf{\passthrough{\lstinline!AVAILABLE!} /
\passthrough{\lstinline!VALUE\_TYPE!}}：当前绑定源码已实现该消息值操作。

\textbf{参数}

\begin{longtable}[]{@{}
  >{\raggedright\arraybackslash}p{(\columnwidth - 6\tabcolsep) * \real{0.18}}
  >{\raggedright\arraybackslash}p{(\columnwidth - 6\tabcolsep) * \real{0.24}}
  >{\raggedright\arraybackslash}p{(\columnwidth - 6\tabcolsep) * \real{0.18}}
  >{\raggedright\arraybackslash}p{(\columnwidth - 6\tabcolsep) * \real{0.40}}@{}}
\toprule
名称 & Python 类型 & 默认值 & 含义 \\
\midrule
\endhead
\passthrough{\lstinline!other!} & \passthrough{\lstinline!object!} &
必填 & 用于比较的另一个 Python 对象。类型不兼容时比较结果通常为
\passthrough{\lstinline!False!}。 \\
\bottomrule
\end{longtable}

\textbf{返回值}

\begin{longtable}[]{@{}
  >{\raggedright\arraybackslash}p{(\columnwidth - 4\tabcolsep) * \real{0.18}}
  >{\raggedright\arraybackslash}p{(\columnwidth - 4\tabcolsep) * \real{0.20}}
  >{\raggedright\arraybackslash}p{(\columnwidth - 4\tabcolsep) * \real{0.62}}@{}}
\toprule
位置 & 类型 & 含义 \\
\midrule
\endhead
返回值 & \passthrough{\lstinline!bool!} & 返回
\passthrough{\lstinline!bool!}。更精确的业务含义见该方法用途、C++
签名和目标型号协议。 \\
\bottomrule
\end{longtable}

\textbf{对应 C++}

\begin{itemize}
\tightlist
\item
  类：\passthrough{\lstinline!geometry\_msgs::msg::dds\_::Quaternion\_!}
\item
  签名：\passthrough{\lstinline"operator!=(const value \&) const"}
\item
  绑定策略：\passthrough{\lstinline!未单独标注!}
\item
  声明位置：由 pybind11 手工绑定或生成注册表提供；manifest
  未记录独立头文件行号。
\end{itemize}

\textbf{说明}

\begin{itemize}
\tightlist
\item
  示例是局部调用片段；\passthrough{\lstinline!obj!}
  和参数变量表示已经按上层业务校验并准备好的对象。
\end{itemize}

\textbf{用法}

\begin{lstlisting}[language=Python]
different = left != right
\end{lstlisting}

\hypertarget{unitree-sdk2-cpp-idl-ros2-quaternion-x}{%
\paragraph{\texorpdfstring{\texttt{unitree\_sdk2\_cpp.idl.ros2.Quaternion.x}}{unitree\_sdk2\_cpp.idl.ros2.Quaternion.x}}\label{unitree-sdk2-cpp-idl-ros2-quaternion-x}}

底层 \passthrough{\lstinline!geometry\_msgs::msg::dds\_::Quaternion\_!}
的 \passthrough{\lstinline!x!} 字段。类型签名只说明 Python
表示；单位、数值范围、数组固定长度和枚举含义需查对应 IDL 头文件。
底层为浮点数；类型本身不说明单位、坐标系或业务安全范围。

\textbf{签名}

\begin{lstlisting}[language=Python]
@property
def x(self) -> float

@x.setter
def x(self, value: float) -> None
\end{lstlisting}

\textbf{可用性与安全性}

\textbf{\passthrough{\lstinline!AVAILABLE!} /
\passthrough{\lstinline!VALUE\_TYPE!}}：当前绑定源码已实现该消息值操作。

\textbf{参数}

\begin{longtable}[]{@{}
  >{\raggedright\arraybackslash}p{(\columnwidth - 4\tabcolsep) * \real{0.18}}
  >{\raggedright\arraybackslash}p{(\columnwidth - 4\tabcolsep) * \real{0.20}}
  >{\raggedright\arraybackslash}p{(\columnwidth - 4\tabcolsep) * \real{0.62}}@{}}
\toprule
名称 & Python 类型 & 含义 \\
\midrule
\endhead
\passthrough{\lstinline!value!} & \passthrough{\lstinline!float!} &
要写入或传递的新值，必须符合该参数的 Python 类型以及底层 C++
范围约束。 \\
\bottomrule
\end{longtable}

\textbf{返回值}

\begin{longtable}[]{@{}
  >{\raggedright\arraybackslash}p{(\columnwidth - 4\tabcolsep) * \real{0.18}}
  >{\raggedright\arraybackslash}p{(\columnwidth - 4\tabcolsep) * \real{0.20}}
  >{\raggedright\arraybackslash}p{(\columnwidth - 4\tabcolsep) * \real{0.62}}@{}}
\toprule
位置 & 类型 & 含义 \\
\midrule
\endhead
getter 返回值 & \passthrough{\lstinline!float!} & 返回属性当前值。 \\
\bottomrule
\end{longtable}

\textbf{对应 C++}

\begin{itemize}
\tightlist
\item
  类：\passthrough{\lstinline!geometry\_msgs::msg::dds\_::Quaternion\_!}
\item
  字段类型：\passthrough{\lstinline!double!}
\item
  声明位置：\passthrough{\lstinline!include/unitree/idl/ros2/Quaternion\_.hpp:22!}
\end{itemize}

\textbf{用法}

\begin{lstlisting}[language=Python]
current_value = obj.x
obj.x = new_value
\end{lstlisting}

\hypertarget{unitree-sdk2-cpp-idl-ros2-quaternion-y}{%
\paragraph{\texorpdfstring{\texttt{unitree\_sdk2\_cpp.idl.ros2.Quaternion.y}}{unitree\_sdk2\_cpp.idl.ros2.Quaternion.y}}\label{unitree-sdk2-cpp-idl-ros2-quaternion-y}}

底层 \passthrough{\lstinline!geometry\_msgs::msg::dds\_::Quaternion\_!}
的 \passthrough{\lstinline!y!} 字段。类型签名只说明 Python
表示；单位、数值范围、数组固定长度和枚举含义需查对应 IDL 头文件。
底层为浮点数；类型本身不说明单位、坐标系或业务安全范围。

\textbf{签名}

\begin{lstlisting}[language=Python]
@property
def y(self) -> float

@y.setter
def y(self, value: float) -> None
\end{lstlisting}

\textbf{可用性与安全性}

\textbf{\passthrough{\lstinline!AVAILABLE!} /
\passthrough{\lstinline!VALUE\_TYPE!}}：当前绑定源码已实现该消息值操作。

\textbf{参数}

\begin{longtable}[]{@{}
  >{\raggedright\arraybackslash}p{(\columnwidth - 4\tabcolsep) * \real{0.18}}
  >{\raggedright\arraybackslash}p{(\columnwidth - 4\tabcolsep) * \real{0.20}}
  >{\raggedright\arraybackslash}p{(\columnwidth - 4\tabcolsep) * \real{0.62}}@{}}
\toprule
名称 & Python 类型 & 含义 \\
\midrule
\endhead
\passthrough{\lstinline!value!} & \passthrough{\lstinline!float!} &
要写入或传递的新值，必须符合该参数的 Python 类型以及底层 C++
范围约束。 \\
\bottomrule
\end{longtable}

\textbf{返回值}

\begin{longtable}[]{@{}
  >{\raggedright\arraybackslash}p{(\columnwidth - 4\tabcolsep) * \real{0.18}}
  >{\raggedright\arraybackslash}p{(\columnwidth - 4\tabcolsep) * \real{0.20}}
  >{\raggedright\arraybackslash}p{(\columnwidth - 4\tabcolsep) * \real{0.62}}@{}}
\toprule
位置 & 类型 & 含义 \\
\midrule
\endhead
getter 返回值 & \passthrough{\lstinline!float!} & 返回属性当前值。 \\
\bottomrule
\end{longtable}

\textbf{对应 C++}

\begin{itemize}
\tightlist
\item
  类：\passthrough{\lstinline!geometry\_msgs::msg::dds\_::Quaternion\_!}
\item
  字段类型：\passthrough{\lstinline!double!}
\item
  声明位置：\passthrough{\lstinline!include/unitree/idl/ros2/Quaternion\_.hpp:23!}
\end{itemize}

\textbf{用法}

\begin{lstlisting}[language=Python]
current_value = obj.y
obj.y = new_value
\end{lstlisting}

\hypertarget{unitree-sdk2-cpp-idl-ros2-quaternion-z}{%
\paragraph{\texorpdfstring{\texttt{unitree\_sdk2\_cpp.idl.ros2.Quaternion.z}}{unitree\_sdk2\_cpp.idl.ros2.Quaternion.z}}\label{unitree-sdk2-cpp-idl-ros2-quaternion-z}}

底层 \passthrough{\lstinline!geometry\_msgs::msg::dds\_::Quaternion\_!}
的 \passthrough{\lstinline!z!} 字段。类型签名只说明 Python
表示；单位、数值范围、数组固定长度和枚举含义需查对应 IDL 头文件。
底层为浮点数；类型本身不说明单位、坐标系或业务安全范围。

\textbf{签名}

\begin{lstlisting}[language=Python]
@property
def z(self) -> float

@z.setter
def z(self, value: float) -> None
\end{lstlisting}

\textbf{可用性与安全性}

\textbf{\passthrough{\lstinline!AVAILABLE!} /
\passthrough{\lstinline!VALUE\_TYPE!}}：当前绑定源码已实现该消息值操作。

\textbf{参数}

\begin{longtable}[]{@{}
  >{\raggedright\arraybackslash}p{(\columnwidth - 4\tabcolsep) * \real{0.18}}
  >{\raggedright\arraybackslash}p{(\columnwidth - 4\tabcolsep) * \real{0.20}}
  >{\raggedright\arraybackslash}p{(\columnwidth - 4\tabcolsep) * \real{0.62}}@{}}
\toprule
名称 & Python 类型 & 含义 \\
\midrule
\endhead
\passthrough{\lstinline!value!} & \passthrough{\lstinline!float!} &
要写入或传递的新值，必须符合该参数的 Python 类型以及底层 C++
范围约束。 \\
\bottomrule
\end{longtable}

\textbf{返回值}

\begin{longtable}[]{@{}
  >{\raggedright\arraybackslash}p{(\columnwidth - 4\tabcolsep) * \real{0.18}}
  >{\raggedright\arraybackslash}p{(\columnwidth - 4\tabcolsep) * \real{0.20}}
  >{\raggedright\arraybackslash}p{(\columnwidth - 4\tabcolsep) * \real{0.62}}@{}}
\toprule
位置 & 类型 & 含义 \\
\midrule
\endhead
getter 返回值 & \passthrough{\lstinline!float!} & 返回属性当前值。 \\
\bottomrule
\end{longtable}

\textbf{对应 C++}

\begin{itemize}
\tightlist
\item
  类：\passthrough{\lstinline!geometry\_msgs::msg::dds\_::Quaternion\_!}
\item
  字段类型：\passthrough{\lstinline!double!}
\item
  声明位置：\passthrough{\lstinline!include/unitree/idl/ros2/Quaternion\_.hpp:24!}
\end{itemize}

\textbf{用法}

\begin{lstlisting}[language=Python]
current_value = obj.z
obj.z = new_value
\end{lstlisting}

\hypertarget{unitree-sdk2-cpp-idl-ros2-quaternion-w}{%
\paragraph{\texorpdfstring{\texttt{unitree\_sdk2\_cpp.idl.ros2.Quaternion.w}}{unitree\_sdk2\_cpp.idl.ros2.Quaternion.w}}\label{unitree-sdk2-cpp-idl-ros2-quaternion-w}}

底层 \passthrough{\lstinline!geometry\_msgs::msg::dds\_::Quaternion\_!}
的 \passthrough{\lstinline!w!} 字段。类型签名只说明 Python
表示；单位、数值范围、数组固定长度和枚举含义需查对应 IDL 头文件。
底层为浮点数；类型本身不说明单位、坐标系或业务安全范围。

\textbf{签名}

\begin{lstlisting}[language=Python]
@property
def w(self) -> float

@w.setter
def w(self, value: float) -> None
\end{lstlisting}

\textbf{可用性与安全性}

\textbf{\passthrough{\lstinline!AVAILABLE!} /
\passthrough{\lstinline!VALUE\_TYPE!}}：当前绑定源码已实现该消息值操作。

\textbf{参数}

\begin{longtable}[]{@{}
  >{\raggedright\arraybackslash}p{(\columnwidth - 4\tabcolsep) * \real{0.18}}
  >{\raggedright\arraybackslash}p{(\columnwidth - 4\tabcolsep) * \real{0.20}}
  >{\raggedright\arraybackslash}p{(\columnwidth - 4\tabcolsep) * \real{0.62}}@{}}
\toprule
名称 & Python 类型 & 含义 \\
\midrule
\endhead
\passthrough{\lstinline!value!} & \passthrough{\lstinline!float!} &
要写入或传递的新值，必须符合该参数的 Python 类型以及底层 C++
范围约束。 \\
\bottomrule
\end{longtable}

\textbf{返回值}

\begin{longtable}[]{@{}
  >{\raggedright\arraybackslash}p{(\columnwidth - 4\tabcolsep) * \real{0.18}}
  >{\raggedright\arraybackslash}p{(\columnwidth - 4\tabcolsep) * \real{0.20}}
  >{\raggedright\arraybackslash}p{(\columnwidth - 4\tabcolsep) * \real{0.62}}@{}}
\toprule
位置 & 类型 & 含义 \\
\midrule
\endhead
getter 返回值 & \passthrough{\lstinline!float!} & 返回属性当前值。 \\
\bottomrule
\end{longtable}

\textbf{对应 C++}

\begin{itemize}
\tightlist
\item
  类：\passthrough{\lstinline!geometry\_msgs::msg::dds\_::Quaternion\_!}
\item
  字段类型：\passthrough{\lstinline!double!}
\item
  声明位置：\passthrough{\lstinline!include/unitree/idl/ros2/Quaternion\_.hpp:25!}
\end{itemize}

\textbf{用法}

\begin{lstlisting}[language=Python]
current_value = obj.w
obj.w = new_value
\end{lstlisting}

\hypertarget{unitree-sdk2-cpp-idl-ros2-vector3}{%
\subsubsection{\texorpdfstring{\texttt{unitree\_sdk2\_cpp.idl.ros2.Vector3}}{unitree\_sdk2\_cpp.idl.ros2.Vector3}}\label{unitree-sdk2-cpp-idl-ros2-vector3}}

可由 Python 读写的 DDS/IDL 消息值类型。字段采用复制语义。

\textbf{导入}

\begin{lstlisting}[language=Python]
from unitree_sdk2_cpp.idl.ros2 import Vector3
\end{lstlisting}

\textbf{方法索引}

\begin{longtable}[]{@{}
  >{\raggedright\arraybackslash}p{(\columnwidth - 6\tabcolsep) * \real{0.18}}
  >{\raggedright\arraybackslash}p{(\columnwidth - 6\tabcolsep) * \real{0.24}}
  >{\raggedright\arraybackslash}p{(\columnwidth - 6\tabcolsep) * \real{0.18}}
  >{\raggedright\arraybackslash}p{(\columnwidth - 6\tabcolsep) * \real{0.40}}@{}}
\toprule
方法 & Python 签名 & 状态 & 安全分类 \\
\midrule
\endhead
\protect\hyperlink{unitree-sdk2-cpp-idl-ros2-vector3-dunder-init-1}{\passthrough{\lstinline!\_\_init\_\_!}}
& \passthrough{\lstinline!def \_\_init\_\_(self) -> None!} &
\passthrough{\lstinline!AVAILABLE!} &
\passthrough{\lstinline!VALUE\_TYPE!} \\
\protect\hyperlink{unitree-sdk2-cpp-idl-ros2-vector3-dunder-eq-1}{\passthrough{\lstinline!\_\_eq\_\_!}}
& \passthrough{\lstinline!def \_\_eq\_\_(self, other: object) -> bool!}
& \passthrough{\lstinline!AVAILABLE!} &
\passthrough{\lstinline!VALUE\_TYPE!} \\
\protect\hyperlink{unitree-sdk2-cpp-idl-ros2-vector3-dunder-ne-1}{\passthrough{\lstinline!\_\_ne\_\_!}}
& \passthrough{\lstinline!def \_\_ne\_\_(self, other: object) -> bool!}
& \passthrough{\lstinline!AVAILABLE!} &
\passthrough{\lstinline!VALUE\_TYPE!} \\
\bottomrule
\end{longtable}

\textbf{属性索引}

\begin{longtable}[]{@{}
  >{\raggedright\arraybackslash}p{(\columnwidth - 6\tabcolsep) * \real{0.18}}
  >{\raggedright\arraybackslash}p{(\columnwidth - 6\tabcolsep) * \real{0.24}}
  >{\raggedright\arraybackslash}p{(\columnwidth - 6\tabcolsep) * \real{0.18}}
  >{\raggedright\arraybackslash}p{(\columnwidth - 6\tabcolsep) * \real{0.40}}@{}}
\toprule
属性 & 读取类型 & 写入类型 & 状态 \\
\midrule
\endhead
\protect\hyperlink{unitree-sdk2-cpp-idl-ros2-vector3-x}{\passthrough{\lstinline!x!}}
& \passthrough{\lstinline!float!} & \passthrough{\lstinline!float!} &
\passthrough{\lstinline!AVAILABLE!} \\
\protect\hyperlink{unitree-sdk2-cpp-idl-ros2-vector3-y}{\passthrough{\lstinline!y!}}
& \passthrough{\lstinline!float!} & \passthrough{\lstinline!float!} &
\passthrough{\lstinline!AVAILABLE!} \\
\protect\hyperlink{unitree-sdk2-cpp-idl-ros2-vector3-z}{\passthrough{\lstinline!z!}}
& \passthrough{\lstinline!float!} & \passthrough{\lstinline!float!} &
\passthrough{\lstinline!AVAILABLE!} \\
\bottomrule
\end{longtable}

\hypertarget{unitree-sdk2-cpp-idl-ros2-vector3-dunder-init-1}{%
\paragraph{\texorpdfstring{\texttt{unitree\_sdk2\_cpp.idl.ros2.Vector3.\_\_init\_\_}}{unitree\_sdk2\_cpp.idl.ros2.Vector3.\_\_init\_\_}}\label{unitree-sdk2-cpp-idl-ros2-vector3-dunder-init-1}}

初始化 \passthrough{\lstinline!Vector3!}
实例。是否能实际构造取决于下方可用性状态。

\textbf{签名}

\begin{lstlisting}[language=Python]
def __init__(self) -> None
\end{lstlisting}

\textbf{可用性与安全性}

\textbf{\passthrough{\lstinline!AVAILABLE!} /
\passthrough{\lstinline!VALUE\_TYPE!}}：当前绑定源码已实现该消息值操作。

\textbf{参数}

无显式参数。实例方法中的 \passthrough{\lstinline!self!} 由 Python
自动传入。

\textbf{返回值}

\begin{longtable}[]{@{}
  >{\raggedright\arraybackslash}p{(\columnwidth - 4\tabcolsep) * \real{0.18}}
  >{\raggedright\arraybackslash}p{(\columnwidth - 4\tabcolsep) * \real{0.20}}
  >{\raggedright\arraybackslash}p{(\columnwidth - 4\tabcolsep) * \real{0.62}}@{}}
\toprule
位置 & 类型 & 含义 \\
\midrule
\endhead
无 & \passthrough{\lstinline!None!} &
\passthrough{\lstinline!\_\_init\_\_!}
本身不返回值；调用类对象时，在构造函数可用的前提下得到该类实例。 \\
\bottomrule
\end{longtable}

\textbf{对应 C++}

\begin{itemize}
\tightlist
\item
  类：\passthrough{\lstinline!geometry\_msgs::msg::dds\_::Vector3\_!}
\item
  签名：\passthrough{\lstinline!Vector3\_()!}
\item
  绑定策略：\passthrough{\lstinline!未单独标注!}
\item
  声明位置：\passthrough{\lstinline!include/unitree/idl/ros2/Vector3\_.hpp:27!}
\end{itemize}

\textbf{说明}

\begin{itemize}
\tightlist
\item
  示例是局部调用片段；\passthrough{\lstinline!obj!}
  和参数变量表示已经按上层业务校验并准备好的对象。
\end{itemize}

\textbf{用法}

\begin{lstlisting}[language=Python]
value = Vector3()
\end{lstlisting}

\hypertarget{unitree-sdk2-cpp-idl-ros2-vector3-dunder-eq-1}{%
\paragraph{\texorpdfstring{\texttt{unitree\_sdk2\_cpp.idl.ros2.Vector3.\_\_eq\_\_}}{unitree\_sdk2\_cpp.idl.ros2.Vector3.\_\_eq\_\_}}\label{unitree-sdk2-cpp-idl-ros2-vector3-dunder-eq-1}}

按底层 IDL 消息内容比较两个对象是否相等。

\textbf{签名}

\begin{lstlisting}[language=Python]
def __eq__(self, other: object) -> bool
\end{lstlisting}

\textbf{可用性与安全性}

\textbf{\passthrough{\lstinline!AVAILABLE!} /
\passthrough{\lstinline!VALUE\_TYPE!}}：当前绑定源码已实现该消息值操作。

\textbf{参数}

\begin{longtable}[]{@{}
  >{\raggedright\arraybackslash}p{(\columnwidth - 6\tabcolsep) * \real{0.18}}
  >{\raggedright\arraybackslash}p{(\columnwidth - 6\tabcolsep) * \real{0.24}}
  >{\raggedright\arraybackslash}p{(\columnwidth - 6\tabcolsep) * \real{0.18}}
  >{\raggedright\arraybackslash}p{(\columnwidth - 6\tabcolsep) * \real{0.40}}@{}}
\toprule
名称 & Python 类型 & 默认值 & 含义 \\
\midrule
\endhead
\passthrough{\lstinline!other!} & \passthrough{\lstinline!object!} &
必填 & 用于比较的另一个 Python 对象。类型不兼容时比较结果通常为
\passthrough{\lstinline!False!}。 \\
\bottomrule
\end{longtable}

\textbf{返回值}

\begin{longtable}[]{@{}
  >{\raggedright\arraybackslash}p{(\columnwidth - 4\tabcolsep) * \real{0.18}}
  >{\raggedright\arraybackslash}p{(\columnwidth - 4\tabcolsep) * \real{0.20}}
  >{\raggedright\arraybackslash}p{(\columnwidth - 4\tabcolsep) * \real{0.62}}@{}}
\toprule
位置 & 类型 & 含义 \\
\midrule
\endhead
返回值 & \passthrough{\lstinline!bool!} & 返回
\passthrough{\lstinline!bool!}。更精确的业务含义见该方法用途、C++
签名和目标型号协议。 \\
\bottomrule
\end{longtable}

\textbf{对应 C++}

\begin{itemize}
\tightlist
\item
  类：\passthrough{\lstinline!geometry\_msgs::msg::dds\_::Vector3\_!}
\item
  签名：\passthrough{\lstinline!operator==(const value \&) const!}
\item
  绑定策略：\passthrough{\lstinline!未单独标注!}
\item
  声明位置：由 pybind11 手工绑定或生成注册表提供；manifest
  未记录独立头文件行号。
\end{itemize}

\textbf{说明}

\begin{itemize}
\tightlist
\item
  示例是局部调用片段；\passthrough{\lstinline!obj!}
  和参数变量表示已经按上层业务校验并准备好的对象。
\end{itemize}

\textbf{用法}

\begin{lstlisting}[language=Python]
same = left == right
\end{lstlisting}

\hypertarget{unitree-sdk2-cpp-idl-ros2-vector3-dunder-ne-1}{%
\paragraph{\texorpdfstring{\texttt{unitree\_sdk2\_cpp.idl.ros2.Vector3.\_\_ne\_\_}}{unitree\_sdk2\_cpp.idl.ros2.Vector3.\_\_ne\_\_}}\label{unitree-sdk2-cpp-idl-ros2-vector3-dunder-ne-1}}

按底层 IDL 消息内容比较两个对象是否不相等。

\textbf{签名}

\begin{lstlisting}[language=Python]
def __ne__(self, other: object) -> bool
\end{lstlisting}

\textbf{可用性与安全性}

\textbf{\passthrough{\lstinline!AVAILABLE!} /
\passthrough{\lstinline!VALUE\_TYPE!}}：当前绑定源码已实现该消息值操作。

\textbf{参数}

\begin{longtable}[]{@{}
  >{\raggedright\arraybackslash}p{(\columnwidth - 6\tabcolsep) * \real{0.18}}
  >{\raggedright\arraybackslash}p{(\columnwidth - 6\tabcolsep) * \real{0.24}}
  >{\raggedright\arraybackslash}p{(\columnwidth - 6\tabcolsep) * \real{0.18}}
  >{\raggedright\arraybackslash}p{(\columnwidth - 6\tabcolsep) * \real{0.40}}@{}}
\toprule
名称 & Python 类型 & 默认值 & 含义 \\
\midrule
\endhead
\passthrough{\lstinline!other!} & \passthrough{\lstinline!object!} &
必填 & 用于比较的另一个 Python 对象。类型不兼容时比较结果通常为
\passthrough{\lstinline!False!}。 \\
\bottomrule
\end{longtable}

\textbf{返回值}

\begin{longtable}[]{@{}
  >{\raggedright\arraybackslash}p{(\columnwidth - 4\tabcolsep) * \real{0.18}}
  >{\raggedright\arraybackslash}p{(\columnwidth - 4\tabcolsep) * \real{0.20}}
  >{\raggedright\arraybackslash}p{(\columnwidth - 4\tabcolsep) * \real{0.62}}@{}}
\toprule
位置 & 类型 & 含义 \\
\midrule
\endhead
返回值 & \passthrough{\lstinline!bool!} & 返回
\passthrough{\lstinline!bool!}。更精确的业务含义见该方法用途、C++
签名和目标型号协议。 \\
\bottomrule
\end{longtable}

\textbf{对应 C++}

\begin{itemize}
\tightlist
\item
  类：\passthrough{\lstinline!geometry\_msgs::msg::dds\_::Vector3\_!}
\item
  签名：\passthrough{\lstinline"operator!=(const value \&) const"}
\item
  绑定策略：\passthrough{\lstinline!未单独标注!}
\item
  声明位置：由 pybind11 手工绑定或生成注册表提供；manifest
  未记录独立头文件行号。
\end{itemize}

\textbf{说明}

\begin{itemize}
\tightlist
\item
  示例是局部调用片段；\passthrough{\lstinline!obj!}
  和参数变量表示已经按上层业务校验并准备好的对象。
\end{itemize}

\textbf{用法}

\begin{lstlisting}[language=Python]
different = left != right
\end{lstlisting}

\hypertarget{unitree-sdk2-cpp-idl-ros2-vector3-x}{%
\paragraph{\texorpdfstring{\texttt{unitree\_sdk2\_cpp.idl.ros2.Vector3.x}}{unitree\_sdk2\_cpp.idl.ros2.Vector3.x}}\label{unitree-sdk2-cpp-idl-ros2-vector3-x}}

底层 \passthrough{\lstinline!geometry\_msgs::msg::dds\_::Vector3\_!} 的
\passthrough{\lstinline!x!} 字段。类型签名只说明 Python
表示；单位、数值范围、数组固定长度和枚举含义需查对应 IDL 头文件。
底层为浮点数；类型本身不说明单位、坐标系或业务安全范围。

\textbf{签名}

\begin{lstlisting}[language=Python]
@property
def x(self) -> float

@x.setter
def x(self, value: float) -> None
\end{lstlisting}

\textbf{可用性与安全性}

\textbf{\passthrough{\lstinline!AVAILABLE!} /
\passthrough{\lstinline!VALUE\_TYPE!}}：当前绑定源码已实现该消息值操作。

\textbf{参数}

\begin{longtable}[]{@{}
  >{\raggedright\arraybackslash}p{(\columnwidth - 4\tabcolsep) * \real{0.18}}
  >{\raggedright\arraybackslash}p{(\columnwidth - 4\tabcolsep) * \real{0.20}}
  >{\raggedright\arraybackslash}p{(\columnwidth - 4\tabcolsep) * \real{0.62}}@{}}
\toprule
名称 & Python 类型 & 含义 \\
\midrule
\endhead
\passthrough{\lstinline!value!} & \passthrough{\lstinline!float!} &
要写入或传递的新值，必须符合该参数的 Python 类型以及底层 C++
范围约束。 \\
\bottomrule
\end{longtable}

\textbf{返回值}

\begin{longtable}[]{@{}
  >{\raggedright\arraybackslash}p{(\columnwidth - 4\tabcolsep) * \real{0.18}}
  >{\raggedright\arraybackslash}p{(\columnwidth - 4\tabcolsep) * \real{0.20}}
  >{\raggedright\arraybackslash}p{(\columnwidth - 4\tabcolsep) * \real{0.62}}@{}}
\toprule
位置 & 类型 & 含义 \\
\midrule
\endhead
getter 返回值 & \passthrough{\lstinline!float!} & 返回属性当前值。 \\
\bottomrule
\end{longtable}

\textbf{对应 C++}

\begin{itemize}
\tightlist
\item
  类：\passthrough{\lstinline!geometry\_msgs::msg::dds\_::Vector3\_!}
\item
  字段类型：\passthrough{\lstinline!double!}
\item
  声明位置：\passthrough{\lstinline!include/unitree/idl/ros2/Vector3\_.hpp:22!}
\end{itemize}

\textbf{用法}

\begin{lstlisting}[language=Python]
current_value = obj.x
obj.x = new_value
\end{lstlisting}

\hypertarget{unitree-sdk2-cpp-idl-ros2-vector3-y}{%
\paragraph{\texorpdfstring{\texttt{unitree\_sdk2\_cpp.idl.ros2.Vector3.y}}{unitree\_sdk2\_cpp.idl.ros2.Vector3.y}}\label{unitree-sdk2-cpp-idl-ros2-vector3-y}}

底层 \passthrough{\lstinline!geometry\_msgs::msg::dds\_::Vector3\_!} 的
\passthrough{\lstinline!y!} 字段。类型签名只说明 Python
表示；单位、数值范围、数组固定长度和枚举含义需查对应 IDL 头文件。
底层为浮点数；类型本身不说明单位、坐标系或业务安全范围。

\textbf{签名}

\begin{lstlisting}[language=Python]
@property
def y(self) -> float

@y.setter
def y(self, value: float) -> None
\end{lstlisting}

\textbf{可用性与安全性}

\textbf{\passthrough{\lstinline!AVAILABLE!} /
\passthrough{\lstinline!VALUE\_TYPE!}}：当前绑定源码已实现该消息值操作。

\textbf{参数}

\begin{longtable}[]{@{}
  >{\raggedright\arraybackslash}p{(\columnwidth - 4\tabcolsep) * \real{0.18}}
  >{\raggedright\arraybackslash}p{(\columnwidth - 4\tabcolsep) * \real{0.20}}
  >{\raggedright\arraybackslash}p{(\columnwidth - 4\tabcolsep) * \real{0.62}}@{}}
\toprule
名称 & Python 类型 & 含义 \\
\midrule
\endhead
\passthrough{\lstinline!value!} & \passthrough{\lstinline!float!} &
要写入或传递的新值，必须符合该参数的 Python 类型以及底层 C++
范围约束。 \\
\bottomrule
\end{longtable}

\textbf{返回值}

\begin{longtable}[]{@{}
  >{\raggedright\arraybackslash}p{(\columnwidth - 4\tabcolsep) * \real{0.18}}
  >{\raggedright\arraybackslash}p{(\columnwidth - 4\tabcolsep) * \real{0.20}}
  >{\raggedright\arraybackslash}p{(\columnwidth - 4\tabcolsep) * \real{0.62}}@{}}
\toprule
位置 & 类型 & 含义 \\
\midrule
\endhead
getter 返回值 & \passthrough{\lstinline!float!} & 返回属性当前值。 \\
\bottomrule
\end{longtable}

\textbf{对应 C++}

\begin{itemize}
\tightlist
\item
  类：\passthrough{\lstinline!geometry\_msgs::msg::dds\_::Vector3\_!}
\item
  字段类型：\passthrough{\lstinline!double!}
\item
  声明位置：\passthrough{\lstinline!include/unitree/idl/ros2/Vector3\_.hpp:23!}
\end{itemize}

\textbf{用法}

\begin{lstlisting}[language=Python]
current_value = obj.y
obj.y = new_value
\end{lstlisting}

\hypertarget{unitree-sdk2-cpp-idl-ros2-vector3-z}{%
\paragraph{\texorpdfstring{\texttt{unitree\_sdk2\_cpp.idl.ros2.Vector3.z}}{unitree\_sdk2\_cpp.idl.ros2.Vector3.z}}\label{unitree-sdk2-cpp-idl-ros2-vector3-z}}

底层 \passthrough{\lstinline!geometry\_msgs::msg::dds\_::Vector3\_!} 的
\passthrough{\lstinline!z!} 字段。类型签名只说明 Python
表示；单位、数值范围、数组固定长度和枚举含义需查对应 IDL 头文件。
底层为浮点数；类型本身不说明单位、坐标系或业务安全范围。

\textbf{签名}

\begin{lstlisting}[language=Python]
@property
def z(self) -> float

@z.setter
def z(self, value: float) -> None
\end{lstlisting}

\textbf{可用性与安全性}

\textbf{\passthrough{\lstinline!AVAILABLE!} /
\passthrough{\lstinline!VALUE\_TYPE!}}：当前绑定源码已实现该消息值操作。

\textbf{参数}

\begin{longtable}[]{@{}
  >{\raggedright\arraybackslash}p{(\columnwidth - 4\tabcolsep) * \real{0.18}}
  >{\raggedright\arraybackslash}p{(\columnwidth - 4\tabcolsep) * \real{0.20}}
  >{\raggedright\arraybackslash}p{(\columnwidth - 4\tabcolsep) * \real{0.62}}@{}}
\toprule
名称 & Python 类型 & 含义 \\
\midrule
\endhead
\passthrough{\lstinline!value!} & \passthrough{\lstinline!float!} &
要写入或传递的新值，必须符合该参数的 Python 类型以及底层 C++
范围约束。 \\
\bottomrule
\end{longtable}

\textbf{返回值}

\begin{longtable}[]{@{}
  >{\raggedright\arraybackslash}p{(\columnwidth - 4\tabcolsep) * \real{0.18}}
  >{\raggedright\arraybackslash}p{(\columnwidth - 4\tabcolsep) * \real{0.20}}
  >{\raggedright\arraybackslash}p{(\columnwidth - 4\tabcolsep) * \real{0.62}}@{}}
\toprule
位置 & 类型 & 含义 \\
\midrule
\endhead
getter 返回值 & \passthrough{\lstinline!float!} & 返回属性当前值。 \\
\bottomrule
\end{longtable}

\textbf{对应 C++}

\begin{itemize}
\tightlist
\item
  类：\passthrough{\lstinline!geometry\_msgs::msg::dds\_::Vector3\_!}
\item
  字段类型：\passthrough{\lstinline!double!}
\item
  声明位置：\passthrough{\lstinline!include/unitree/idl/ros2/Vector3\_.hpp:24!}
\end{itemize}

\textbf{用法}

\begin{lstlisting}[language=Python]
current_value = obj.z
obj.z = new_value
\end{lstlisting}

\hypertarget{unitree-sdk2-cpp-idl-ros2-imu}{%
\subsubsection{\texorpdfstring{\texttt{unitree\_sdk2\_cpp.idl.ros2.Imu}}{unitree\_sdk2\_cpp.idl.ros2.Imu}}\label{unitree-sdk2-cpp-idl-ros2-imu}}

可由 Python 读写的 DDS/IDL 消息值类型。字段采用复制语义。

\textbf{导入}

\begin{lstlisting}[language=Python]
from unitree_sdk2_cpp.idl.ros2 import Imu
\end{lstlisting}

\textbf{方法索引}

\begin{longtable}[]{@{}
  >{\raggedright\arraybackslash}p{(\columnwidth - 6\tabcolsep) * \real{0.18}}
  >{\raggedright\arraybackslash}p{(\columnwidth - 6\tabcolsep) * \real{0.24}}
  >{\raggedright\arraybackslash}p{(\columnwidth - 6\tabcolsep) * \real{0.18}}
  >{\raggedright\arraybackslash}p{(\columnwidth - 6\tabcolsep) * \real{0.40}}@{}}
\toprule
方法 & Python 签名 & 状态 & 安全分类 \\
\midrule
\endhead
\protect\hyperlink{unitree-sdk2-cpp-idl-ros2-imu-dunder-init-1}{\passthrough{\lstinline!\_\_init\_\_!}}
& \passthrough{\lstinline!def \_\_init\_\_(self) -> None!} &
\passthrough{\lstinline!AVAILABLE!} &
\passthrough{\lstinline!VALUE\_TYPE!} \\
\protect\hyperlink{unitree-sdk2-cpp-idl-ros2-imu-dunder-eq-1}{\passthrough{\lstinline!\_\_eq\_\_!}}
& \passthrough{\lstinline!def \_\_eq\_\_(self, other: object) -> bool!}
& \passthrough{\lstinline!AVAILABLE!} &
\passthrough{\lstinline!VALUE\_TYPE!} \\
\protect\hyperlink{unitree-sdk2-cpp-idl-ros2-imu-dunder-ne-1}{\passthrough{\lstinline!\_\_ne\_\_!}}
& \passthrough{\lstinline!def \_\_ne\_\_(self, other: object) -> bool!}
& \passthrough{\lstinline!AVAILABLE!} &
\passthrough{\lstinline!VALUE\_TYPE!} \\
\bottomrule
\end{longtable}

\textbf{属性索引}

\begin{longtable}[]{@{}
  >{\raggedright\arraybackslash}p{(\columnwidth - 6\tabcolsep) * \real{0.18}}
  >{\raggedright\arraybackslash}p{(\columnwidth - 6\tabcolsep) * \real{0.24}}
  >{\raggedright\arraybackslash}p{(\columnwidth - 6\tabcolsep) * \real{0.18}}
  >{\raggedright\arraybackslash}p{(\columnwidth - 6\tabcolsep) * \real{0.40}}@{}}
\toprule
属性 & 读取类型 & 写入类型 & 状态 \\
\midrule
\endhead
\protect\hyperlink{unitree-sdk2-cpp-idl-ros2-imu-header}{\passthrough{\lstinline!header!}}
& \passthrough{\lstinline!Any!} & \passthrough{\lstinline!Any!} &
\passthrough{\lstinline!AVAILABLE!} \\
\protect\hyperlink{unitree-sdk2-cpp-idl-ros2-imu-orientation}{\passthrough{\lstinline!orientation!}}
& \passthrough{\lstinline!Any!} & \passthrough{\lstinline!Any!} &
\passthrough{\lstinline!AVAILABLE!} \\
\protect\hyperlink{unitree-sdk2-cpp-idl-ros2-imu-orientation-covariance}{\passthrough{\lstinline!orientation\_covariance!}}
& \passthrough{\lstinline!list[float]!} &
\passthrough{\lstinline!Sequence[float]!} &
\passthrough{\lstinline!AVAILABLE!} \\
\protect\hyperlink{unitree-sdk2-cpp-idl-ros2-imu-angular-velocity}{\passthrough{\lstinline!angular\_velocity!}}
& \passthrough{\lstinline!Any!} & \passthrough{\lstinline!Any!} &
\passthrough{\lstinline!AVAILABLE!} \\
\protect\hyperlink{unitree-sdk2-cpp-idl-ros2-imu-angular-velocity-covariance}{\passthrough{\lstinline!angular\_velocity\_covariance!}}
& \passthrough{\lstinline!list[float]!} &
\passthrough{\lstinline!Sequence[float]!} &
\passthrough{\lstinline!AVAILABLE!} \\
\protect\hyperlink{unitree-sdk2-cpp-idl-ros2-imu-linear-acceleration}{\passthrough{\lstinline!linear\_acceleration!}}
& \passthrough{\lstinline!Any!} & \passthrough{\lstinline!Any!} &
\passthrough{\lstinline!AVAILABLE!} \\
\protect\hyperlink{unitree-sdk2-cpp-idl-ros2-imu-linear-acceleration-covariance}{\passthrough{\lstinline!linear\_acceleration\_covariance!}}
& \passthrough{\lstinline!list[float]!} &
\passthrough{\lstinline!Sequence[float]!} &
\passthrough{\lstinline!AVAILABLE!} \\
\bottomrule
\end{longtable}

\hypertarget{unitree-sdk2-cpp-idl-ros2-imu-dunder-init-1}{%
\paragraph{\texorpdfstring{\texttt{unitree\_sdk2\_cpp.idl.ros2.Imu.\_\_init\_\_}}{unitree\_sdk2\_cpp.idl.ros2.Imu.\_\_init\_\_}}\label{unitree-sdk2-cpp-idl-ros2-imu-dunder-init-1}}

初始化 \passthrough{\lstinline!Imu!}
实例。是否能实际构造取决于下方可用性状态。

\textbf{签名}

\begin{lstlisting}[language=Python]
def __init__(self) -> None
\end{lstlisting}

\textbf{可用性与安全性}

\textbf{\passthrough{\lstinline!AVAILABLE!} /
\passthrough{\lstinline!VALUE\_TYPE!}}：当前绑定源码已实现该消息值操作。

\textbf{参数}

无显式参数。实例方法中的 \passthrough{\lstinline!self!} 由 Python
自动传入。

\textbf{返回值}

\begin{longtable}[]{@{}
  >{\raggedright\arraybackslash}p{(\columnwidth - 4\tabcolsep) * \real{0.18}}
  >{\raggedright\arraybackslash}p{(\columnwidth - 4\tabcolsep) * \real{0.20}}
  >{\raggedright\arraybackslash}p{(\columnwidth - 4\tabcolsep) * \real{0.62}}@{}}
\toprule
位置 & 类型 & 含义 \\
\midrule
\endhead
无 & \passthrough{\lstinline!None!} &
\passthrough{\lstinline!\_\_init\_\_!}
本身不返回值；调用类对象时，在构造函数可用的前提下得到该类实例。 \\
\bottomrule
\end{longtable}

\textbf{对应 C++}

\begin{itemize}
\tightlist
\item
  类：\passthrough{\lstinline!sensor\_msgs::msg::dds\_::Imu\_!}
\item
  签名：\passthrough{\lstinline!Imu\_()!}
\item
  绑定策略：\passthrough{\lstinline!未单独标注!}
\item
  声明位置：\passthrough{\lstinline!include/unitree/idl/ros2/Imu\_.hpp:38!}
\end{itemize}

\textbf{说明}

\begin{itemize}
\tightlist
\item
  示例是局部调用片段；\passthrough{\lstinline!obj!}
  和参数变量表示已经按上层业务校验并准备好的对象。
\end{itemize}

\textbf{用法}

\begin{lstlisting}[language=Python]
value = Imu()
\end{lstlisting}

\hypertarget{unitree-sdk2-cpp-idl-ros2-imu-dunder-eq-1}{%
\paragraph{\texorpdfstring{\texttt{unitree\_sdk2\_cpp.idl.ros2.Imu.\_\_eq\_\_}}{unitree\_sdk2\_cpp.idl.ros2.Imu.\_\_eq\_\_}}\label{unitree-sdk2-cpp-idl-ros2-imu-dunder-eq-1}}

按底层 IDL 消息内容比较两个对象是否相等。

\textbf{签名}

\begin{lstlisting}[language=Python]
def __eq__(self, other: object) -> bool
\end{lstlisting}

\textbf{可用性与安全性}

\textbf{\passthrough{\lstinline!AVAILABLE!} /
\passthrough{\lstinline!VALUE\_TYPE!}}：当前绑定源码已实现该消息值操作。

\textbf{参数}

\begin{longtable}[]{@{}
  >{\raggedright\arraybackslash}p{(\columnwidth - 6\tabcolsep) * \real{0.18}}
  >{\raggedright\arraybackslash}p{(\columnwidth - 6\tabcolsep) * \real{0.24}}
  >{\raggedright\arraybackslash}p{(\columnwidth - 6\tabcolsep) * \real{0.18}}
  >{\raggedright\arraybackslash}p{(\columnwidth - 6\tabcolsep) * \real{0.40}}@{}}
\toprule
名称 & Python 类型 & 默认值 & 含义 \\
\midrule
\endhead
\passthrough{\lstinline!other!} & \passthrough{\lstinline!object!} &
必填 & 用于比较的另一个 Python 对象。类型不兼容时比较结果通常为
\passthrough{\lstinline!False!}。 \\
\bottomrule
\end{longtable}

\textbf{返回值}

\begin{longtable}[]{@{}
  >{\raggedright\arraybackslash}p{(\columnwidth - 4\tabcolsep) * \real{0.18}}
  >{\raggedright\arraybackslash}p{(\columnwidth - 4\tabcolsep) * \real{0.20}}
  >{\raggedright\arraybackslash}p{(\columnwidth - 4\tabcolsep) * \real{0.62}}@{}}
\toprule
位置 & 类型 & 含义 \\
\midrule
\endhead
返回值 & \passthrough{\lstinline!bool!} & 返回
\passthrough{\lstinline!bool!}。更精确的业务含义见该方法用途、C++
签名和目标型号协议。 \\
\bottomrule
\end{longtable}

\textbf{对应 C++}

\begin{itemize}
\tightlist
\item
  类：\passthrough{\lstinline!sensor\_msgs::msg::dds\_::Imu\_!}
\item
  签名：\passthrough{\lstinline!operator==(const value \&) const!}
\item
  绑定策略：\passthrough{\lstinline!未单独标注!}
\item
  声明位置：由 pybind11 手工绑定或生成注册表提供；manifest
  未记录独立头文件行号。
\end{itemize}

\textbf{说明}

\begin{itemize}
\tightlist
\item
  示例是局部调用片段；\passthrough{\lstinline!obj!}
  和参数变量表示已经按上层业务校验并准备好的对象。
\end{itemize}

\textbf{用法}

\begin{lstlisting}[language=Python]
same = left == right
\end{lstlisting}

\hypertarget{unitree-sdk2-cpp-idl-ros2-imu-dunder-ne-1}{%
\paragraph{\texorpdfstring{\texttt{unitree\_sdk2\_cpp.idl.ros2.Imu.\_\_ne\_\_}}{unitree\_sdk2\_cpp.idl.ros2.Imu.\_\_ne\_\_}}\label{unitree-sdk2-cpp-idl-ros2-imu-dunder-ne-1}}

按底层 IDL 消息内容比较两个对象是否不相等。

\textbf{签名}

\begin{lstlisting}[language=Python]
def __ne__(self, other: object) -> bool
\end{lstlisting}

\textbf{可用性与安全性}

\textbf{\passthrough{\lstinline!AVAILABLE!} /
\passthrough{\lstinline!VALUE\_TYPE!}}：当前绑定源码已实现该消息值操作。

\textbf{参数}

\begin{longtable}[]{@{}
  >{\raggedright\arraybackslash}p{(\columnwidth - 6\tabcolsep) * \real{0.18}}
  >{\raggedright\arraybackslash}p{(\columnwidth - 6\tabcolsep) * \real{0.24}}
  >{\raggedright\arraybackslash}p{(\columnwidth - 6\tabcolsep) * \real{0.18}}
  >{\raggedright\arraybackslash}p{(\columnwidth - 6\tabcolsep) * \real{0.40}}@{}}
\toprule
名称 & Python 类型 & 默认值 & 含义 \\
\midrule
\endhead
\passthrough{\lstinline!other!} & \passthrough{\lstinline!object!} &
必填 & 用于比较的另一个 Python 对象。类型不兼容时比较结果通常为
\passthrough{\lstinline!False!}。 \\
\bottomrule
\end{longtable}

\textbf{返回值}

\begin{longtable}[]{@{}
  >{\raggedright\arraybackslash}p{(\columnwidth - 4\tabcolsep) * \real{0.18}}
  >{\raggedright\arraybackslash}p{(\columnwidth - 4\tabcolsep) * \real{0.20}}
  >{\raggedright\arraybackslash}p{(\columnwidth - 4\tabcolsep) * \real{0.62}}@{}}
\toprule
位置 & 类型 & 含义 \\
\midrule
\endhead
返回值 & \passthrough{\lstinline!bool!} & 返回
\passthrough{\lstinline!bool!}。更精确的业务含义见该方法用途、C++
签名和目标型号协议。 \\
\bottomrule
\end{longtable}

\textbf{对应 C++}

\begin{itemize}
\tightlist
\item
  类：\passthrough{\lstinline!sensor\_msgs::msg::dds\_::Imu\_!}
\item
  签名：\passthrough{\lstinline"operator!=(const value \&) const"}
\item
  绑定策略：\passthrough{\lstinline!未单独标注!}
\item
  声明位置：由 pybind11 手工绑定或生成注册表提供；manifest
  未记录独立头文件行号。
\end{itemize}

\textbf{说明}

\begin{itemize}
\tightlist
\item
  示例是局部调用片段；\passthrough{\lstinline!obj!}
  和参数变量表示已经按上层业务校验并准备好的对象。
\end{itemize}

\textbf{用法}

\begin{lstlisting}[language=Python]
different = left != right
\end{lstlisting}

\hypertarget{unitree-sdk2-cpp-idl-ros2-imu-header}{%
\paragraph{\texorpdfstring{\texttt{unitree\_sdk2\_cpp.idl.ros2.Imu.header}}{unitree\_sdk2\_cpp.idl.ros2.Imu.header}}\label{unitree-sdk2-cpp-idl-ros2-imu-header}}

底层 \passthrough{\lstinline!sensor\_msgs::msg::dds\_::Imu\_!} 的
\passthrough{\lstinline!header!} 字段。类型签名只说明 Python
表示；单位、数值范围、数组固定长度和枚举含义需查对应 IDL 头文件。

\textbf{签名}

\begin{lstlisting}[language=Python]
@property
def header(self) -> Any

@header.setter
def header(self, value: Any) -> None
\end{lstlisting}

\textbf{可用性与安全性}

\textbf{\passthrough{\lstinline!AVAILABLE!} /
\passthrough{\lstinline!VALUE\_TYPE!}}：当前绑定源码已实现该消息值操作。

\textbf{参数}

\begin{longtable}[]{@{}
  >{\raggedright\arraybackslash}p{(\columnwidth - 4\tabcolsep) * \real{0.18}}
  >{\raggedright\arraybackslash}p{(\columnwidth - 4\tabcolsep) * \real{0.20}}
  >{\raggedright\arraybackslash}p{(\columnwidth - 4\tabcolsep) * \real{0.62}}@{}}
\toprule
名称 & Python 类型 & 含义 \\
\midrule
\endhead
\passthrough{\lstinline!value!} & \passthrough{\lstinline!Any!} &
要写入或传递的新值，必须符合该参数的 Python 类型以及底层 C++
范围约束。 \\
\bottomrule
\end{longtable}

\textbf{返回值}

\begin{longtable}[]{@{}
  >{\raggedright\arraybackslash}p{(\columnwidth - 4\tabcolsep) * \real{0.18}}
  >{\raggedright\arraybackslash}p{(\columnwidth - 4\tabcolsep) * \real{0.20}}
  >{\raggedright\arraybackslash}p{(\columnwidth - 4\tabcolsep) * \real{0.62}}@{}}
\toprule
位置 & 类型 & 含义 \\
\midrule
\endhead
getter 返回值 & \passthrough{\lstinline!Any!} & 返回属性当前值。 \\
\bottomrule
\end{longtable}

\textbf{对应 C++}

\begin{itemize}
\tightlist
\item
  类：\passthrough{\lstinline!sensor\_msgs::msg::dds\_::Imu\_!}
\item
  字段类型：\passthrough{\lstinline!::std\_msgs::msg::dds\_::Header\_!}
\item
  声明位置：\passthrough{\lstinline!include/unitree/idl/ros2/Imu\_.hpp:29!}
\end{itemize}

\textbf{用法}

\begin{lstlisting}[language=Python]
current_value = obj.header
obj.header = new_value
\end{lstlisting}

\hypertarget{unitree-sdk2-cpp-idl-ros2-imu-orientation}{%
\paragraph{\texorpdfstring{\texttt{unitree\_sdk2\_cpp.idl.ros2.Imu.orientation}}{unitree\_sdk2\_cpp.idl.ros2.Imu.orientation}}\label{unitree-sdk2-cpp-idl-ros2-imu-orientation}}

底层 \passthrough{\lstinline!sensor\_msgs::msg::dds\_::Imu\_!} 的
\passthrough{\lstinline!orientation!} 字段。类型签名只说明 Python
表示；单位、数值范围、数组固定长度和枚举含义需查对应 IDL 头文件。

\textbf{签名}

\begin{lstlisting}[language=Python]
@property
def orientation(self) -> Any

@orientation.setter
def orientation(self, value: Any) -> None
\end{lstlisting}

\textbf{可用性与安全性}

\textbf{\passthrough{\lstinline!AVAILABLE!} /
\passthrough{\lstinline!VALUE\_TYPE!}}：当前绑定源码已实现该消息值操作。

\textbf{参数}

\begin{longtable}[]{@{}
  >{\raggedright\arraybackslash}p{(\columnwidth - 4\tabcolsep) * \real{0.18}}
  >{\raggedright\arraybackslash}p{(\columnwidth - 4\tabcolsep) * \real{0.20}}
  >{\raggedright\arraybackslash}p{(\columnwidth - 4\tabcolsep) * \real{0.62}}@{}}
\toprule
名称 & Python 类型 & 含义 \\
\midrule
\endhead
\passthrough{\lstinline!value!} & \passthrough{\lstinline!Any!} &
要写入或传递的新值，必须符合该参数的 Python 类型以及底层 C++
范围约束。 \\
\bottomrule
\end{longtable}

\textbf{返回值}

\begin{longtable}[]{@{}
  >{\raggedright\arraybackslash}p{(\columnwidth - 4\tabcolsep) * \real{0.18}}
  >{\raggedright\arraybackslash}p{(\columnwidth - 4\tabcolsep) * \real{0.20}}
  >{\raggedright\arraybackslash}p{(\columnwidth - 4\tabcolsep) * \real{0.62}}@{}}
\toprule
位置 & 类型 & 含义 \\
\midrule
\endhead
getter 返回值 & \passthrough{\lstinline!Any!} & 返回属性当前值。 \\
\bottomrule
\end{longtable}

\textbf{对应 C++}

\begin{itemize}
\tightlist
\item
  类：\passthrough{\lstinline!sensor\_msgs::msg::dds\_::Imu\_!}
\item
  字段类型：\passthrough{\lstinline!::geometry\_msgs::msg::dds\_::Quaternion\_!}
\item
  声明位置：\passthrough{\lstinline!include/unitree/idl/ros2/Imu\_.hpp:30!}
\end{itemize}

\textbf{用法}

\begin{lstlisting}[language=Python]
current_value = obj.orientation
obj.orientation = new_value
\end{lstlisting}

\hypertarget{unitree-sdk2-cpp-idl-ros2-imu-orientation-covariance}{%
\paragraph{\texorpdfstring{\texttt{unitree\_sdk2\_cpp.idl.ros2.Imu.orientation\_covariance}}{unitree\_sdk2\_cpp.idl.ros2.Imu.orientation\_covariance}}\label{unitree-sdk2-cpp-idl-ros2-imu-orientation-covariance}}

底层 \passthrough{\lstinline!sensor\_msgs::msg::dds\_::Imu\_!} 的
\passthrough{\lstinline!orientation\_covariance!} 字段。类型签名只说明
Python 表示；单位、数值范围、数组固定长度和枚举含义需查对应 IDL 头文件。
底层是固定长度数组，写入序列必须正好包含 9
个元素；元素约束：底层为浮点数；类型本身不说明单位、坐标系或业务安全范围。

\textbf{签名}

\begin{lstlisting}[language=Python]
@property
def orientation_covariance(self) -> list[float]

@orientation_covariance.setter
def orientation_covariance(self, value: Sequence[float]) -> None
\end{lstlisting}

\textbf{可用性与安全性}

\textbf{\passthrough{\lstinline!AVAILABLE!} /
\passthrough{\lstinline!VALUE\_TYPE!}}：当前绑定源码已实现该消息值操作。

\textbf{参数}

\begin{longtable}[]{@{}
  >{\raggedright\arraybackslash}p{(\columnwidth - 4\tabcolsep) * \real{0.18}}
  >{\raggedright\arraybackslash}p{(\columnwidth - 4\tabcolsep) * \real{0.20}}
  >{\raggedright\arraybackslash}p{(\columnwidth - 4\tabcolsep) * \real{0.62}}@{}}
\toprule
名称 & Python 类型 & 含义 \\
\midrule
\endhead
\passthrough{\lstinline!value!} &
\passthrough{\lstinline!Sequence[float]!} &
要写入或传递的新值，必须符合该参数的 Python 类型以及底层 C++
范围约束。 \\
\bottomrule
\end{longtable}

\textbf{返回值}

\begin{longtable}[]{@{}
  >{\raggedright\arraybackslash}p{(\columnwidth - 4\tabcolsep) * \real{0.18}}
  >{\raggedright\arraybackslash}p{(\columnwidth - 4\tabcolsep) * \real{0.20}}
  >{\raggedright\arraybackslash}p{(\columnwidth - 4\tabcolsep) * \real{0.62}}@{}}
\toprule
位置 & 类型 & 含义 \\
\midrule
\endhead
getter 返回值 & \passthrough{\lstinline!list[float]!} &
返回属性当前值。数组、vector 和嵌套消息采用复制语义。 \\
\bottomrule
\end{longtable}

\textbf{对应 C++}

\begin{itemize}
\tightlist
\item
  类：\passthrough{\lstinline!sensor\_msgs::msg::dds\_::Imu\_!}
\item
  字段类型：\passthrough{\lstinline!std::array<double, 9>!}
\item
  声明位置：\passthrough{\lstinline!include/unitree/idl/ros2/Imu\_.hpp:31!}
\end{itemize}

\textbf{用法}

\begin{lstlisting}[language=Python]
current_value = obj.orientation_covariance
obj.orientation_covariance = new_value
\end{lstlisting}

\begin{quote}
{[}!TIP{]} 读取容器或嵌套值后应遵循``读取、修改、重新赋值''。只修改
getter 返回的副本不会可靠地更新原 C++ 消息。
\end{quote}

\hypertarget{unitree-sdk2-cpp-idl-ros2-imu-angular-velocity}{%
\paragraph{\texorpdfstring{\texttt{unitree\_sdk2\_cpp.idl.ros2.Imu.angular\_velocity}}{unitree\_sdk2\_cpp.idl.ros2.Imu.angular\_velocity}}\label{unitree-sdk2-cpp-idl-ros2-imu-angular-velocity}}

底层 \passthrough{\lstinline!sensor\_msgs::msg::dds\_::Imu\_!} 的
\passthrough{\lstinline!angular\_velocity!} 字段。类型签名只说明 Python
表示；单位、数值范围、数组固定长度和枚举含义需查对应 IDL 头文件。

\textbf{签名}

\begin{lstlisting}[language=Python]
@property
def angular_velocity(self) -> Any

@angular_velocity.setter
def angular_velocity(self, value: Any) -> None
\end{lstlisting}

\textbf{可用性与安全性}

\textbf{\passthrough{\lstinline!AVAILABLE!} /
\passthrough{\lstinline!VALUE\_TYPE!}}：当前绑定源码已实现该消息值操作。

\textbf{参数}

\begin{longtable}[]{@{}
  >{\raggedright\arraybackslash}p{(\columnwidth - 4\tabcolsep) * \real{0.18}}
  >{\raggedright\arraybackslash}p{(\columnwidth - 4\tabcolsep) * \real{0.20}}
  >{\raggedright\arraybackslash}p{(\columnwidth - 4\tabcolsep) * \real{0.62}}@{}}
\toprule
名称 & Python 类型 & 含义 \\
\midrule
\endhead
\passthrough{\lstinline!value!} & \passthrough{\lstinline!Any!} &
要写入或传递的新值，必须符合该参数的 Python 类型以及底层 C++
范围约束。 \\
\bottomrule
\end{longtable}

\textbf{返回值}

\begin{longtable}[]{@{}
  >{\raggedright\arraybackslash}p{(\columnwidth - 4\tabcolsep) * \real{0.18}}
  >{\raggedright\arraybackslash}p{(\columnwidth - 4\tabcolsep) * \real{0.20}}
  >{\raggedright\arraybackslash}p{(\columnwidth - 4\tabcolsep) * \real{0.62}}@{}}
\toprule
位置 & 类型 & 含义 \\
\midrule
\endhead
getter 返回值 & \passthrough{\lstinline!Any!} & 返回属性当前值。 \\
\bottomrule
\end{longtable}

\textbf{对应 C++}

\begin{itemize}
\tightlist
\item
  类：\passthrough{\lstinline!sensor\_msgs::msg::dds\_::Imu\_!}
\item
  字段类型：\passthrough{\lstinline!::geometry\_msgs::msg::dds\_::Vector3\_!}
\item
  声明位置：\passthrough{\lstinline!include/unitree/idl/ros2/Imu\_.hpp:32!}
\end{itemize}

\textbf{用法}

\begin{lstlisting}[language=Python]
current_value = obj.angular_velocity
obj.angular_velocity = new_value
\end{lstlisting}

\hypertarget{unitree-sdk2-cpp-idl-ros2-imu-angular-velocity-covariance}{%
\paragraph{\texorpdfstring{\texttt{unitree\_sdk2\_cpp.idl.ros2.Imu.angular\_velocity\_covariance}}{unitree\_sdk2\_cpp.idl.ros2.Imu.angular\_velocity\_covariance}}\label{unitree-sdk2-cpp-idl-ros2-imu-angular-velocity-covariance}}

底层 \passthrough{\lstinline!sensor\_msgs::msg::dds\_::Imu\_!} 的
\passthrough{\lstinline!angular\_velocity\_covariance!}
字段。类型签名只说明 Python
表示；单位、数值范围、数组固定长度和枚举含义需查对应 IDL 头文件。
底层是固定长度数组，写入序列必须正好包含 9
个元素；元素约束：底层为浮点数；类型本身不说明单位、坐标系或业务安全范围。

\textbf{签名}

\begin{lstlisting}[language=Python]
@property
def angular_velocity_covariance(self) -> list[float]

@angular_velocity_covariance.setter
def angular_velocity_covariance(self, value: Sequence[float]) -> None
\end{lstlisting}

\textbf{可用性与安全性}

\textbf{\passthrough{\lstinline!AVAILABLE!} /
\passthrough{\lstinline!VALUE\_TYPE!}}：当前绑定源码已实现该消息值操作。

\textbf{参数}

\begin{longtable}[]{@{}
  >{\raggedright\arraybackslash}p{(\columnwidth - 4\tabcolsep) * \real{0.18}}
  >{\raggedright\arraybackslash}p{(\columnwidth - 4\tabcolsep) * \real{0.20}}
  >{\raggedright\arraybackslash}p{(\columnwidth - 4\tabcolsep) * \real{0.62}}@{}}
\toprule
名称 & Python 类型 & 含义 \\
\midrule
\endhead
\passthrough{\lstinline!value!} &
\passthrough{\lstinline!Sequence[float]!} &
要写入或传递的新值，必须符合该参数的 Python 类型以及底层 C++
范围约束。 \\
\bottomrule
\end{longtable}

\textbf{返回值}

\begin{longtable}[]{@{}
  >{\raggedright\arraybackslash}p{(\columnwidth - 4\tabcolsep) * \real{0.18}}
  >{\raggedright\arraybackslash}p{(\columnwidth - 4\tabcolsep) * \real{0.20}}
  >{\raggedright\arraybackslash}p{(\columnwidth - 4\tabcolsep) * \real{0.62}}@{}}
\toprule
位置 & 类型 & 含义 \\
\midrule
\endhead
getter 返回值 & \passthrough{\lstinline!list[float]!} &
返回属性当前值。数组、vector 和嵌套消息采用复制语义。 \\
\bottomrule
\end{longtable}

\textbf{对应 C++}

\begin{itemize}
\tightlist
\item
  类：\passthrough{\lstinline!sensor\_msgs::msg::dds\_::Imu\_!}
\item
  字段类型：\passthrough{\lstinline!std::array<double, 9>!}
\item
  声明位置：\passthrough{\lstinline!include/unitree/idl/ros2/Imu\_.hpp:33!}
\end{itemize}

\textbf{用法}

\begin{lstlisting}[language=Python]
current_value = obj.angular_velocity_covariance
obj.angular_velocity_covariance = new_value
\end{lstlisting}

\begin{quote}
{[}!TIP{]} 读取容器或嵌套值后应遵循``读取、修改、重新赋值''。只修改
getter 返回的副本不会可靠地更新原 C++ 消息。
\end{quote}

\hypertarget{unitree-sdk2-cpp-idl-ros2-imu-linear-acceleration}{%
\paragraph{\texorpdfstring{\texttt{unitree\_sdk2\_cpp.idl.ros2.Imu.linear\_acceleration}}{unitree\_sdk2\_cpp.idl.ros2.Imu.linear\_acceleration}}\label{unitree-sdk2-cpp-idl-ros2-imu-linear-acceleration}}

底层 \passthrough{\lstinline!sensor\_msgs::msg::dds\_::Imu\_!} 的
\passthrough{\lstinline!linear\_acceleration!} 字段。类型签名只说明
Python 表示；单位、数值范围、数组固定长度和枚举含义需查对应 IDL 头文件。

\textbf{签名}

\begin{lstlisting}[language=Python]
@property
def linear_acceleration(self) -> Any

@linear_acceleration.setter
def linear_acceleration(self, value: Any) -> None
\end{lstlisting}

\textbf{可用性与安全性}

\textbf{\passthrough{\lstinline!AVAILABLE!} /
\passthrough{\lstinline!VALUE\_TYPE!}}：当前绑定源码已实现该消息值操作。

\textbf{参数}

\begin{longtable}[]{@{}
  >{\raggedright\arraybackslash}p{(\columnwidth - 4\tabcolsep) * \real{0.18}}
  >{\raggedright\arraybackslash}p{(\columnwidth - 4\tabcolsep) * \real{0.20}}
  >{\raggedright\arraybackslash}p{(\columnwidth - 4\tabcolsep) * \real{0.62}}@{}}
\toprule
名称 & Python 类型 & 含义 \\
\midrule
\endhead
\passthrough{\lstinline!value!} & \passthrough{\lstinline!Any!} &
要写入或传递的新值，必须符合该参数的 Python 类型以及底层 C++
范围约束。 \\
\bottomrule
\end{longtable}

\textbf{返回值}

\begin{longtable}[]{@{}
  >{\raggedright\arraybackslash}p{(\columnwidth - 4\tabcolsep) * \real{0.18}}
  >{\raggedright\arraybackslash}p{(\columnwidth - 4\tabcolsep) * \real{0.20}}
  >{\raggedright\arraybackslash}p{(\columnwidth - 4\tabcolsep) * \real{0.62}}@{}}
\toprule
位置 & 类型 & 含义 \\
\midrule
\endhead
getter 返回值 & \passthrough{\lstinline!Any!} & 返回属性当前值。 \\
\bottomrule
\end{longtable}

\textbf{对应 C++}

\begin{itemize}
\tightlist
\item
  类：\passthrough{\lstinline!sensor\_msgs::msg::dds\_::Imu\_!}
\item
  字段类型：\passthrough{\lstinline!::geometry\_msgs::msg::dds\_::Vector3\_!}
\item
  声明位置：\passthrough{\lstinline!include/unitree/idl/ros2/Imu\_.hpp:34!}
\end{itemize}

\textbf{用法}

\begin{lstlisting}[language=Python]
current_value = obj.linear_acceleration
obj.linear_acceleration = new_value
\end{lstlisting}

\hypertarget{unitree-sdk2-cpp-idl-ros2-imu-linear-acceleration-covariance}{%
\paragraph{\texorpdfstring{\texttt{unitree\_sdk2\_cpp.idl.ros2.Imu.linear\_acceleration\_covariance}}{unitree\_sdk2\_cpp.idl.ros2.Imu.linear\_acceleration\_covariance}}\label{unitree-sdk2-cpp-idl-ros2-imu-linear-acceleration-covariance}}

底层 \passthrough{\lstinline!sensor\_msgs::msg::dds\_::Imu\_!} 的
\passthrough{\lstinline!linear\_acceleration\_covariance!}
字段。类型签名只说明 Python
表示；单位、数值范围、数组固定长度和枚举含义需查对应 IDL 头文件。
底层是固定长度数组，写入序列必须正好包含 9
个元素；元素约束：底层为浮点数；类型本身不说明单位、坐标系或业务安全范围。

\textbf{签名}

\begin{lstlisting}[language=Python]
@property
def linear_acceleration_covariance(self) -> list[float]

@linear_acceleration_covariance.setter
def linear_acceleration_covariance(self, value: Sequence[float]) -> None
\end{lstlisting}

\textbf{可用性与安全性}

\textbf{\passthrough{\lstinline!AVAILABLE!} /
\passthrough{\lstinline!VALUE\_TYPE!}}：当前绑定源码已实现该消息值操作。

\textbf{参数}

\begin{longtable}[]{@{}
  >{\raggedright\arraybackslash}p{(\columnwidth - 4\tabcolsep) * \real{0.18}}
  >{\raggedright\arraybackslash}p{(\columnwidth - 4\tabcolsep) * \real{0.20}}
  >{\raggedright\arraybackslash}p{(\columnwidth - 4\tabcolsep) * \real{0.62}}@{}}
\toprule
名称 & Python 类型 & 含义 \\
\midrule
\endhead
\passthrough{\lstinline!value!} &
\passthrough{\lstinline!Sequence[float]!} &
要写入或传递的新值，必须符合该参数的 Python 类型以及底层 C++
范围约束。 \\
\bottomrule
\end{longtable}

\textbf{返回值}

\begin{longtable}[]{@{}
  >{\raggedright\arraybackslash}p{(\columnwidth - 4\tabcolsep) * \real{0.18}}
  >{\raggedright\arraybackslash}p{(\columnwidth - 4\tabcolsep) * \real{0.20}}
  >{\raggedright\arraybackslash}p{(\columnwidth - 4\tabcolsep) * \real{0.62}}@{}}
\toprule
位置 & 类型 & 含义 \\
\midrule
\endhead
getter 返回值 & \passthrough{\lstinline!list[float]!} &
返回属性当前值。数组、vector 和嵌套消息采用复制语义。 \\
\bottomrule
\end{longtable}

\textbf{对应 C++}

\begin{itemize}
\tightlist
\item
  类：\passthrough{\lstinline!sensor\_msgs::msg::dds\_::Imu\_!}
\item
  字段类型：\passthrough{\lstinline!std::array<double, 9>!}
\item
  声明位置：\passthrough{\lstinline!include/unitree/idl/ros2/Imu\_.hpp:35!}
\end{itemize}

\textbf{用法}

\begin{lstlisting}[language=Python]
current_value = obj.linear_acceleration_covariance
obj.linear_acceleration_covariance = new_value
\end{lstlisting}

\begin{quote}
{[}!TIP{]} 读取容器或嵌套值后应遵循``读取、修改、重新赋值''。只修改
getter 返回的副本不会可靠地更新原 C++ 消息。
\end{quote}

\hypertarget{unitree-sdk2-cpp-idl-ros2-point}{%
\subsubsection{\texorpdfstring{\texttt{unitree\_sdk2\_cpp.idl.ros2.Point}}{unitree\_sdk2\_cpp.idl.ros2.Point}}\label{unitree-sdk2-cpp-idl-ros2-point}}

可由 Python 读写的 DDS/IDL 消息值类型。字段采用复制语义。

\textbf{导入}

\begin{lstlisting}[language=Python]
from unitree_sdk2_cpp.idl.ros2 import Point
\end{lstlisting}

\textbf{方法索引}

\begin{longtable}[]{@{}
  >{\raggedright\arraybackslash}p{(\columnwidth - 6\tabcolsep) * \real{0.18}}
  >{\raggedright\arraybackslash}p{(\columnwidth - 6\tabcolsep) * \real{0.24}}
  >{\raggedright\arraybackslash}p{(\columnwidth - 6\tabcolsep) * \real{0.18}}
  >{\raggedright\arraybackslash}p{(\columnwidth - 6\tabcolsep) * \real{0.40}}@{}}
\toprule
方法 & Python 签名 & 状态 & 安全分类 \\
\midrule
\endhead
\protect\hyperlink{unitree-sdk2-cpp-idl-ros2-point-dunder-init-1}{\passthrough{\lstinline!\_\_init\_\_!}}
& \passthrough{\lstinline!def \_\_init\_\_(self) -> None!} &
\passthrough{\lstinline!AVAILABLE!} &
\passthrough{\lstinline!VALUE\_TYPE!} \\
\protect\hyperlink{unitree-sdk2-cpp-idl-ros2-point-dunder-eq-1}{\passthrough{\lstinline!\_\_eq\_\_!}}
& \passthrough{\lstinline!def \_\_eq\_\_(self, other: object) -> bool!}
& \passthrough{\lstinline!AVAILABLE!} &
\passthrough{\lstinline!VALUE\_TYPE!} \\
\protect\hyperlink{unitree-sdk2-cpp-idl-ros2-point-dunder-ne-1}{\passthrough{\lstinline!\_\_ne\_\_!}}
& \passthrough{\lstinline!def \_\_ne\_\_(self, other: object) -> bool!}
& \passthrough{\lstinline!AVAILABLE!} &
\passthrough{\lstinline!VALUE\_TYPE!} \\
\bottomrule
\end{longtable}

\textbf{属性索引}

\begin{longtable}[]{@{}
  >{\raggedright\arraybackslash}p{(\columnwidth - 6\tabcolsep) * \real{0.18}}
  >{\raggedright\arraybackslash}p{(\columnwidth - 6\tabcolsep) * \real{0.24}}
  >{\raggedright\arraybackslash}p{(\columnwidth - 6\tabcolsep) * \real{0.18}}
  >{\raggedright\arraybackslash}p{(\columnwidth - 6\tabcolsep) * \real{0.40}}@{}}
\toprule
属性 & 读取类型 & 写入类型 & 状态 \\
\midrule
\endhead
\protect\hyperlink{unitree-sdk2-cpp-idl-ros2-point-x}{\passthrough{\lstinline!x!}}
& \passthrough{\lstinline!float!} & \passthrough{\lstinline!float!} &
\passthrough{\lstinline!AVAILABLE!} \\
\protect\hyperlink{unitree-sdk2-cpp-idl-ros2-point-y}{\passthrough{\lstinline!y!}}
& \passthrough{\lstinline!float!} & \passthrough{\lstinline!float!} &
\passthrough{\lstinline!AVAILABLE!} \\
\protect\hyperlink{unitree-sdk2-cpp-idl-ros2-point-z}{\passthrough{\lstinline!z!}}
& \passthrough{\lstinline!float!} & \passthrough{\lstinline!float!} &
\passthrough{\lstinline!AVAILABLE!} \\
\bottomrule
\end{longtable}

\hypertarget{unitree-sdk2-cpp-idl-ros2-point-dunder-init-1}{%
\paragraph{\texorpdfstring{\texttt{unitree\_sdk2\_cpp.idl.ros2.Point.\_\_init\_\_}}{unitree\_sdk2\_cpp.idl.ros2.Point.\_\_init\_\_}}\label{unitree-sdk2-cpp-idl-ros2-point-dunder-init-1}}

初始化 \passthrough{\lstinline!Point!}
实例。是否能实际构造取决于下方可用性状态。

\textbf{签名}

\begin{lstlisting}[language=Python]
def __init__(self) -> None
\end{lstlisting}

\textbf{可用性与安全性}

\textbf{\passthrough{\lstinline!AVAILABLE!} /
\passthrough{\lstinline!VALUE\_TYPE!}}：当前绑定源码已实现该消息值操作。

\textbf{参数}

无显式参数。实例方法中的 \passthrough{\lstinline!self!} 由 Python
自动传入。

\textbf{返回值}

\begin{longtable}[]{@{}
  >{\raggedright\arraybackslash}p{(\columnwidth - 4\tabcolsep) * \real{0.18}}
  >{\raggedright\arraybackslash}p{(\columnwidth - 4\tabcolsep) * \real{0.20}}
  >{\raggedright\arraybackslash}p{(\columnwidth - 4\tabcolsep) * \real{0.62}}@{}}
\toprule
位置 & 类型 & 含义 \\
\midrule
\endhead
无 & \passthrough{\lstinline!None!} &
\passthrough{\lstinline!\_\_init\_\_!}
本身不返回值；调用类对象时，在构造函数可用的前提下得到该类实例。 \\
\bottomrule
\end{longtable}

\textbf{对应 C++}

\begin{itemize}
\tightlist
\item
  类：\passthrough{\lstinline!geometry\_msgs::msg::dds\_::Point\_!}
\item
  签名：\passthrough{\lstinline!Point\_()!}
\item
  绑定策略：\passthrough{\lstinline!未单独标注!}
\item
  声明位置：\passthrough{\lstinline!include/unitree/idl/ros2/Point\_.hpp:27!}
\end{itemize}

\textbf{说明}

\begin{itemize}
\tightlist
\item
  示例是局部调用片段；\passthrough{\lstinline!obj!}
  和参数变量表示已经按上层业务校验并准备好的对象。
\end{itemize}

\textbf{用法}

\begin{lstlisting}[language=Python]
value = Point()
\end{lstlisting}

\hypertarget{unitree-sdk2-cpp-idl-ros2-point-dunder-eq-1}{%
\paragraph{\texorpdfstring{\texttt{unitree\_sdk2\_cpp.idl.ros2.Point.\_\_eq\_\_}}{unitree\_sdk2\_cpp.idl.ros2.Point.\_\_eq\_\_}}\label{unitree-sdk2-cpp-idl-ros2-point-dunder-eq-1}}

按底层 IDL 消息内容比较两个对象是否相等。

\textbf{签名}

\begin{lstlisting}[language=Python]
def __eq__(self, other: object) -> bool
\end{lstlisting}

\textbf{可用性与安全性}

\textbf{\passthrough{\lstinline!AVAILABLE!} /
\passthrough{\lstinline!VALUE\_TYPE!}}：当前绑定源码已实现该消息值操作。

\textbf{参数}

\begin{longtable}[]{@{}
  >{\raggedright\arraybackslash}p{(\columnwidth - 6\tabcolsep) * \real{0.18}}
  >{\raggedright\arraybackslash}p{(\columnwidth - 6\tabcolsep) * \real{0.24}}
  >{\raggedright\arraybackslash}p{(\columnwidth - 6\tabcolsep) * \real{0.18}}
  >{\raggedright\arraybackslash}p{(\columnwidth - 6\tabcolsep) * \real{0.40}}@{}}
\toprule
名称 & Python 类型 & 默认值 & 含义 \\
\midrule
\endhead
\passthrough{\lstinline!other!} & \passthrough{\lstinline!object!} &
必填 & 用于比较的另一个 Python 对象。类型不兼容时比较结果通常为
\passthrough{\lstinline!False!}。 \\
\bottomrule
\end{longtable}

\textbf{返回值}

\begin{longtable}[]{@{}
  >{\raggedright\arraybackslash}p{(\columnwidth - 4\tabcolsep) * \real{0.18}}
  >{\raggedright\arraybackslash}p{(\columnwidth - 4\tabcolsep) * \real{0.20}}
  >{\raggedright\arraybackslash}p{(\columnwidth - 4\tabcolsep) * \real{0.62}}@{}}
\toprule
位置 & 类型 & 含义 \\
\midrule
\endhead
返回值 & \passthrough{\lstinline!bool!} & 返回
\passthrough{\lstinline!bool!}。更精确的业务含义见该方法用途、C++
签名和目标型号协议。 \\
\bottomrule
\end{longtable}

\textbf{对应 C++}

\begin{itemize}
\tightlist
\item
  类：\passthrough{\lstinline!geometry\_msgs::msg::dds\_::Point\_!}
\item
  签名：\passthrough{\lstinline!operator==(const value \&) const!}
\item
  绑定策略：\passthrough{\lstinline!未单独标注!}
\item
  声明位置：由 pybind11 手工绑定或生成注册表提供；manifest
  未记录独立头文件行号。
\end{itemize}

\textbf{说明}

\begin{itemize}
\tightlist
\item
  示例是局部调用片段；\passthrough{\lstinline!obj!}
  和参数变量表示已经按上层业务校验并准备好的对象。
\end{itemize}

\textbf{用法}

\begin{lstlisting}[language=Python]
same = left == right
\end{lstlisting}

\hypertarget{unitree-sdk2-cpp-idl-ros2-point-dunder-ne-1}{%
\paragraph{\texorpdfstring{\texttt{unitree\_sdk2\_cpp.idl.ros2.Point.\_\_ne\_\_}}{unitree\_sdk2\_cpp.idl.ros2.Point.\_\_ne\_\_}}\label{unitree-sdk2-cpp-idl-ros2-point-dunder-ne-1}}

按底层 IDL 消息内容比较两个对象是否不相等。

\textbf{签名}

\begin{lstlisting}[language=Python]
def __ne__(self, other: object) -> bool
\end{lstlisting}

\textbf{可用性与安全性}

\textbf{\passthrough{\lstinline!AVAILABLE!} /
\passthrough{\lstinline!VALUE\_TYPE!}}：当前绑定源码已实现该消息值操作。

\textbf{参数}

\begin{longtable}[]{@{}
  >{\raggedright\arraybackslash}p{(\columnwidth - 6\tabcolsep) * \real{0.18}}
  >{\raggedright\arraybackslash}p{(\columnwidth - 6\tabcolsep) * \real{0.24}}
  >{\raggedright\arraybackslash}p{(\columnwidth - 6\tabcolsep) * \real{0.18}}
  >{\raggedright\arraybackslash}p{(\columnwidth - 6\tabcolsep) * \real{0.40}}@{}}
\toprule
名称 & Python 类型 & 默认值 & 含义 \\
\midrule
\endhead
\passthrough{\lstinline!other!} & \passthrough{\lstinline!object!} &
必填 & 用于比较的另一个 Python 对象。类型不兼容时比较结果通常为
\passthrough{\lstinline!False!}。 \\
\bottomrule
\end{longtable}

\textbf{返回值}

\begin{longtable}[]{@{}
  >{\raggedright\arraybackslash}p{(\columnwidth - 4\tabcolsep) * \real{0.18}}
  >{\raggedright\arraybackslash}p{(\columnwidth - 4\tabcolsep) * \real{0.20}}
  >{\raggedright\arraybackslash}p{(\columnwidth - 4\tabcolsep) * \real{0.62}}@{}}
\toprule
位置 & 类型 & 含义 \\
\midrule
\endhead
返回值 & \passthrough{\lstinline!bool!} & 返回
\passthrough{\lstinline!bool!}。更精确的业务含义见该方法用途、C++
签名和目标型号协议。 \\
\bottomrule
\end{longtable}

\textbf{对应 C++}

\begin{itemize}
\tightlist
\item
  类：\passthrough{\lstinline!geometry\_msgs::msg::dds\_::Point\_!}
\item
  签名：\passthrough{\lstinline"operator!=(const value \&) const"}
\item
  绑定策略：\passthrough{\lstinline!未单独标注!}
\item
  声明位置：由 pybind11 手工绑定或生成注册表提供；manifest
  未记录独立头文件行号。
\end{itemize}

\textbf{说明}

\begin{itemize}
\tightlist
\item
  示例是局部调用片段；\passthrough{\lstinline!obj!}
  和参数变量表示已经按上层业务校验并准备好的对象。
\end{itemize}

\textbf{用法}

\begin{lstlisting}[language=Python]
different = left != right
\end{lstlisting}

\hypertarget{unitree-sdk2-cpp-idl-ros2-point-x}{%
\paragraph{\texorpdfstring{\texttt{unitree\_sdk2\_cpp.idl.ros2.Point.x}}{unitree\_sdk2\_cpp.idl.ros2.Point.x}}\label{unitree-sdk2-cpp-idl-ros2-point-x}}

底层 \passthrough{\lstinline!geometry\_msgs::msg::dds\_::Point\_!} 的
\passthrough{\lstinline!x!} 字段。类型签名只说明 Python
表示；单位、数值范围、数组固定长度和枚举含义需查对应 IDL 头文件。
底层为浮点数；类型本身不说明单位、坐标系或业务安全范围。

\textbf{签名}

\begin{lstlisting}[language=Python]
@property
def x(self) -> float

@x.setter
def x(self, value: float) -> None
\end{lstlisting}

\textbf{可用性与安全性}

\textbf{\passthrough{\lstinline!AVAILABLE!} /
\passthrough{\lstinline!VALUE\_TYPE!}}：当前绑定源码已实现该消息值操作。

\textbf{参数}

\begin{longtable}[]{@{}
  >{\raggedright\arraybackslash}p{(\columnwidth - 4\tabcolsep) * \real{0.18}}
  >{\raggedright\arraybackslash}p{(\columnwidth - 4\tabcolsep) * \real{0.20}}
  >{\raggedright\arraybackslash}p{(\columnwidth - 4\tabcolsep) * \real{0.62}}@{}}
\toprule
名称 & Python 类型 & 含义 \\
\midrule
\endhead
\passthrough{\lstinline!value!} & \passthrough{\lstinline!float!} &
要写入或传递的新值，必须符合该参数的 Python 类型以及底层 C++
范围约束。 \\
\bottomrule
\end{longtable}

\textbf{返回值}

\begin{longtable}[]{@{}
  >{\raggedright\arraybackslash}p{(\columnwidth - 4\tabcolsep) * \real{0.18}}
  >{\raggedright\arraybackslash}p{(\columnwidth - 4\tabcolsep) * \real{0.20}}
  >{\raggedright\arraybackslash}p{(\columnwidth - 4\tabcolsep) * \real{0.62}}@{}}
\toprule
位置 & 类型 & 含义 \\
\midrule
\endhead
getter 返回值 & \passthrough{\lstinline!float!} & 返回属性当前值。 \\
\bottomrule
\end{longtable}

\textbf{对应 C++}

\begin{itemize}
\tightlist
\item
  类：\passthrough{\lstinline!geometry\_msgs::msg::dds\_::Point\_!}
\item
  字段类型：\passthrough{\lstinline!double!}
\item
  声明位置：\passthrough{\lstinline!include/unitree/idl/ros2/Point\_.hpp:22!}
\end{itemize}

\textbf{用法}

\begin{lstlisting}[language=Python]
current_value = obj.x
obj.x = new_value
\end{lstlisting}

\hypertarget{unitree-sdk2-cpp-idl-ros2-point-y}{%
\paragraph{\texorpdfstring{\texttt{unitree\_sdk2\_cpp.idl.ros2.Point.y}}{unitree\_sdk2\_cpp.idl.ros2.Point.y}}\label{unitree-sdk2-cpp-idl-ros2-point-y}}

底层 \passthrough{\lstinline!geometry\_msgs::msg::dds\_::Point\_!} 的
\passthrough{\lstinline!y!} 字段。类型签名只说明 Python
表示；单位、数值范围、数组固定长度和枚举含义需查对应 IDL 头文件。
底层为浮点数；类型本身不说明单位、坐标系或业务安全范围。

\textbf{签名}

\begin{lstlisting}[language=Python]
@property
def y(self) -> float

@y.setter
def y(self, value: float) -> None
\end{lstlisting}

\textbf{可用性与安全性}

\textbf{\passthrough{\lstinline!AVAILABLE!} /
\passthrough{\lstinline!VALUE\_TYPE!}}：当前绑定源码已实现该消息值操作。

\textbf{参数}

\begin{longtable}[]{@{}
  >{\raggedright\arraybackslash}p{(\columnwidth - 4\tabcolsep) * \real{0.18}}
  >{\raggedright\arraybackslash}p{(\columnwidth - 4\tabcolsep) * \real{0.20}}
  >{\raggedright\arraybackslash}p{(\columnwidth - 4\tabcolsep) * \real{0.62}}@{}}
\toprule
名称 & Python 类型 & 含义 \\
\midrule
\endhead
\passthrough{\lstinline!value!} & \passthrough{\lstinline!float!} &
要写入或传递的新值，必须符合该参数的 Python 类型以及底层 C++
范围约束。 \\
\bottomrule
\end{longtable}

\textbf{返回值}

\begin{longtable}[]{@{}
  >{\raggedright\arraybackslash}p{(\columnwidth - 4\tabcolsep) * \real{0.18}}
  >{\raggedright\arraybackslash}p{(\columnwidth - 4\tabcolsep) * \real{0.20}}
  >{\raggedright\arraybackslash}p{(\columnwidth - 4\tabcolsep) * \real{0.62}}@{}}
\toprule
位置 & 类型 & 含义 \\
\midrule
\endhead
getter 返回值 & \passthrough{\lstinline!float!} & 返回属性当前值。 \\
\bottomrule
\end{longtable}

\textbf{对应 C++}

\begin{itemize}
\tightlist
\item
  类：\passthrough{\lstinline!geometry\_msgs::msg::dds\_::Point\_!}
\item
  字段类型：\passthrough{\lstinline!double!}
\item
  声明位置：\passthrough{\lstinline!include/unitree/idl/ros2/Point\_.hpp:23!}
\end{itemize}

\textbf{用法}

\begin{lstlisting}[language=Python]
current_value = obj.y
obj.y = new_value
\end{lstlisting}

\hypertarget{unitree-sdk2-cpp-idl-ros2-point-z}{%
\paragraph{\texorpdfstring{\texttt{unitree\_sdk2\_cpp.idl.ros2.Point.z}}{unitree\_sdk2\_cpp.idl.ros2.Point.z}}\label{unitree-sdk2-cpp-idl-ros2-point-z}}

底层 \passthrough{\lstinline!geometry\_msgs::msg::dds\_::Point\_!} 的
\passthrough{\lstinline!z!} 字段。类型签名只说明 Python
表示；单位、数值范围、数组固定长度和枚举含义需查对应 IDL 头文件。
底层为浮点数；类型本身不说明单位、坐标系或业务安全范围。

\textbf{签名}

\begin{lstlisting}[language=Python]
@property
def z(self) -> float

@z.setter
def z(self, value: float) -> None
\end{lstlisting}

\textbf{可用性与安全性}

\textbf{\passthrough{\lstinline!AVAILABLE!} /
\passthrough{\lstinline!VALUE\_TYPE!}}：当前绑定源码已实现该消息值操作。

\textbf{参数}

\begin{longtable}[]{@{}
  >{\raggedright\arraybackslash}p{(\columnwidth - 4\tabcolsep) * \real{0.18}}
  >{\raggedright\arraybackslash}p{(\columnwidth - 4\tabcolsep) * \real{0.20}}
  >{\raggedright\arraybackslash}p{(\columnwidth - 4\tabcolsep) * \real{0.62}}@{}}
\toprule
名称 & Python 类型 & 含义 \\
\midrule
\endhead
\passthrough{\lstinline!value!} & \passthrough{\lstinline!float!} &
要写入或传递的新值，必须符合该参数的 Python 类型以及底层 C++
范围约束。 \\
\bottomrule
\end{longtable}

\textbf{返回值}

\begin{longtable}[]{@{}
  >{\raggedright\arraybackslash}p{(\columnwidth - 4\tabcolsep) * \real{0.18}}
  >{\raggedright\arraybackslash}p{(\columnwidth - 4\tabcolsep) * \real{0.20}}
  >{\raggedright\arraybackslash}p{(\columnwidth - 4\tabcolsep) * \real{0.62}}@{}}
\toprule
位置 & 类型 & 含义 \\
\midrule
\endhead
getter 返回值 & \passthrough{\lstinline!float!} & 返回属性当前值。 \\
\bottomrule
\end{longtable}

\textbf{对应 C++}

\begin{itemize}
\tightlist
\item
  类：\passthrough{\lstinline!geometry\_msgs::msg::dds\_::Point\_!}
\item
  字段类型：\passthrough{\lstinline!double!}
\item
  声明位置：\passthrough{\lstinline!include/unitree/idl/ros2/Point\_.hpp:24!}
\end{itemize}

\textbf{用法}

\begin{lstlisting}[language=Python]
current_value = obj.z
obj.z = new_value
\end{lstlisting}

\hypertarget{unitree-sdk2-cpp-idl-ros2-pose}{%
\subsubsection{\texorpdfstring{\texttt{unitree\_sdk2\_cpp.idl.ros2.Pose}}{unitree\_sdk2\_cpp.idl.ros2.Pose}}\label{unitree-sdk2-cpp-idl-ros2-pose}}

可由 Python 读写的 DDS/IDL 消息值类型。字段采用复制语义。

\textbf{导入}

\begin{lstlisting}[language=Python]
from unitree_sdk2_cpp.idl.ros2 import Pose
\end{lstlisting}

\textbf{方法索引}

\begin{longtable}[]{@{}
  >{\raggedright\arraybackslash}p{(\columnwidth - 6\tabcolsep) * \real{0.18}}
  >{\raggedright\arraybackslash}p{(\columnwidth - 6\tabcolsep) * \real{0.24}}
  >{\raggedright\arraybackslash}p{(\columnwidth - 6\tabcolsep) * \real{0.18}}
  >{\raggedright\arraybackslash}p{(\columnwidth - 6\tabcolsep) * \real{0.40}}@{}}
\toprule
方法 & Python 签名 & 状态 & 安全分类 \\
\midrule
\endhead
\protect\hyperlink{unitree-sdk2-cpp-idl-ros2-pose-dunder-init-1}{\passthrough{\lstinline!\_\_init\_\_!}}
& \passthrough{\lstinline!def \_\_init\_\_(self) -> None!} &
\passthrough{\lstinline!AVAILABLE!} &
\passthrough{\lstinline!VALUE\_TYPE!} \\
\protect\hyperlink{unitree-sdk2-cpp-idl-ros2-pose-dunder-eq-1}{\passthrough{\lstinline!\_\_eq\_\_!}}
& \passthrough{\lstinline!def \_\_eq\_\_(self, other: object) -> bool!}
& \passthrough{\lstinline!AVAILABLE!} &
\passthrough{\lstinline!VALUE\_TYPE!} \\
\protect\hyperlink{unitree-sdk2-cpp-idl-ros2-pose-dunder-ne-1}{\passthrough{\lstinline!\_\_ne\_\_!}}
& \passthrough{\lstinline!def \_\_ne\_\_(self, other: object) -> bool!}
& \passthrough{\lstinline!AVAILABLE!} &
\passthrough{\lstinline!VALUE\_TYPE!} \\
\bottomrule
\end{longtable}

\textbf{属性索引}

\begin{longtable}[]{@{}
  >{\raggedright\arraybackslash}p{(\columnwidth - 6\tabcolsep) * \real{0.18}}
  >{\raggedright\arraybackslash}p{(\columnwidth - 6\tabcolsep) * \real{0.24}}
  >{\raggedright\arraybackslash}p{(\columnwidth - 6\tabcolsep) * \real{0.18}}
  >{\raggedright\arraybackslash}p{(\columnwidth - 6\tabcolsep) * \real{0.40}}@{}}
\toprule
属性 & 读取类型 & 写入类型 & 状态 \\
\midrule
\endhead
\protect\hyperlink{unitree-sdk2-cpp-idl-ros2-pose-position}{\passthrough{\lstinline!position!}}
& \passthrough{\lstinline!Point!} & \passthrough{\lstinline!Point!} &
\passthrough{\lstinline!AVAILABLE!} \\
\protect\hyperlink{unitree-sdk2-cpp-idl-ros2-pose-orientation}{\passthrough{\lstinline!orientation!}}
& \passthrough{\lstinline!Quaternion!} &
\passthrough{\lstinline!Quaternion!} &
\passthrough{\lstinline!AVAILABLE!} \\
\bottomrule
\end{longtable}

\hypertarget{unitree-sdk2-cpp-idl-ros2-pose-dunder-init-1}{%
\paragraph{\texorpdfstring{\texttt{unitree\_sdk2\_cpp.idl.ros2.Pose.\_\_init\_\_}}{unitree\_sdk2\_cpp.idl.ros2.Pose.\_\_init\_\_}}\label{unitree-sdk2-cpp-idl-ros2-pose-dunder-init-1}}

初始化 \passthrough{\lstinline!Pose!}
实例。是否能实际构造取决于下方可用性状态。

\textbf{签名}

\begin{lstlisting}[language=Python]
def __init__(self) -> None
\end{lstlisting}

\textbf{可用性与安全性}

\textbf{\passthrough{\lstinline!AVAILABLE!} /
\passthrough{\lstinline!VALUE\_TYPE!}}：当前绑定源码已实现该消息值操作。

\textbf{参数}

无显式参数。实例方法中的 \passthrough{\lstinline!self!} 由 Python
自动传入。

\textbf{返回值}

\begin{longtable}[]{@{}
  >{\raggedright\arraybackslash}p{(\columnwidth - 4\tabcolsep) * \real{0.18}}
  >{\raggedright\arraybackslash}p{(\columnwidth - 4\tabcolsep) * \real{0.20}}
  >{\raggedright\arraybackslash}p{(\columnwidth - 4\tabcolsep) * \real{0.62}}@{}}
\toprule
位置 & 类型 & 含义 \\
\midrule
\endhead
无 & \passthrough{\lstinline!None!} &
\passthrough{\lstinline!\_\_init\_\_!}
本身不返回值；调用类对象时，在构造函数可用的前提下得到该类实例。 \\
\bottomrule
\end{longtable}

\textbf{对应 C++}

\begin{itemize}
\tightlist
\item
  类：\passthrough{\lstinline!geometry\_msgs::msg::dds\_::Pose\_!}
\item
  签名：\passthrough{\lstinline!Pose\_()!}
\item
  绑定策略：\passthrough{\lstinline!未单独标注!}
\item
  声明位置：\passthrough{\lstinline!include/unitree/idl/ros2/Pose\_.hpp:30!}
\end{itemize}

\textbf{说明}

\begin{itemize}
\tightlist
\item
  示例是局部调用片段；\passthrough{\lstinline!obj!}
  和参数变量表示已经按上层业务校验并准备好的对象。
\end{itemize}

\textbf{用法}

\begin{lstlisting}[language=Python]
value = Pose()
\end{lstlisting}

\hypertarget{unitree-sdk2-cpp-idl-ros2-pose-dunder-eq-1}{%
\paragraph{\texorpdfstring{\texttt{unitree\_sdk2\_cpp.idl.ros2.Pose.\_\_eq\_\_}}{unitree\_sdk2\_cpp.idl.ros2.Pose.\_\_eq\_\_}}\label{unitree-sdk2-cpp-idl-ros2-pose-dunder-eq-1}}

按底层 IDL 消息内容比较两个对象是否相等。

\textbf{签名}

\begin{lstlisting}[language=Python]
def __eq__(self, other: object) -> bool
\end{lstlisting}

\textbf{可用性与安全性}

\textbf{\passthrough{\lstinline!AVAILABLE!} /
\passthrough{\lstinline!VALUE\_TYPE!}}：当前绑定源码已实现该消息值操作。

\textbf{参数}

\begin{longtable}[]{@{}
  >{\raggedright\arraybackslash}p{(\columnwidth - 6\tabcolsep) * \real{0.18}}
  >{\raggedright\arraybackslash}p{(\columnwidth - 6\tabcolsep) * \real{0.24}}
  >{\raggedright\arraybackslash}p{(\columnwidth - 6\tabcolsep) * \real{0.18}}
  >{\raggedright\arraybackslash}p{(\columnwidth - 6\tabcolsep) * \real{0.40}}@{}}
\toprule
名称 & Python 类型 & 默认值 & 含义 \\
\midrule
\endhead
\passthrough{\lstinline!other!} & \passthrough{\lstinline!object!} &
必填 & 用于比较的另一个 Python 对象。类型不兼容时比较结果通常为
\passthrough{\lstinline!False!}。 \\
\bottomrule
\end{longtable}

\textbf{返回值}

\begin{longtable}[]{@{}
  >{\raggedright\arraybackslash}p{(\columnwidth - 4\tabcolsep) * \real{0.18}}
  >{\raggedright\arraybackslash}p{(\columnwidth - 4\tabcolsep) * \real{0.20}}
  >{\raggedright\arraybackslash}p{(\columnwidth - 4\tabcolsep) * \real{0.62}}@{}}
\toprule
位置 & 类型 & 含义 \\
\midrule
\endhead
返回值 & \passthrough{\lstinline!bool!} & 返回
\passthrough{\lstinline!bool!}。更精确的业务含义见该方法用途、C++
签名和目标型号协议。 \\
\bottomrule
\end{longtable}

\textbf{对应 C++}

\begin{itemize}
\tightlist
\item
  类：\passthrough{\lstinline!geometry\_msgs::msg::dds\_::Pose\_!}
\item
  签名：\passthrough{\lstinline!operator==(const value \&) const!}
\item
  绑定策略：\passthrough{\lstinline!未单独标注!}
\item
  声明位置：由 pybind11 手工绑定或生成注册表提供；manifest
  未记录独立头文件行号。
\end{itemize}

\textbf{说明}

\begin{itemize}
\tightlist
\item
  示例是局部调用片段；\passthrough{\lstinline!obj!}
  和参数变量表示已经按上层业务校验并准备好的对象。
\end{itemize}

\textbf{用法}

\begin{lstlisting}[language=Python]
same = left == right
\end{lstlisting}

\hypertarget{unitree-sdk2-cpp-idl-ros2-pose-dunder-ne-1}{%
\paragraph{\texorpdfstring{\texttt{unitree\_sdk2\_cpp.idl.ros2.Pose.\_\_ne\_\_}}{unitree\_sdk2\_cpp.idl.ros2.Pose.\_\_ne\_\_}}\label{unitree-sdk2-cpp-idl-ros2-pose-dunder-ne-1}}

按底层 IDL 消息内容比较两个对象是否不相等。

\textbf{签名}

\begin{lstlisting}[language=Python]
def __ne__(self, other: object) -> bool
\end{lstlisting}

\textbf{可用性与安全性}

\textbf{\passthrough{\lstinline!AVAILABLE!} /
\passthrough{\lstinline!VALUE\_TYPE!}}：当前绑定源码已实现该消息值操作。

\textbf{参数}

\begin{longtable}[]{@{}
  >{\raggedright\arraybackslash}p{(\columnwidth - 6\tabcolsep) * \real{0.18}}
  >{\raggedright\arraybackslash}p{(\columnwidth - 6\tabcolsep) * \real{0.24}}
  >{\raggedright\arraybackslash}p{(\columnwidth - 6\tabcolsep) * \real{0.18}}
  >{\raggedright\arraybackslash}p{(\columnwidth - 6\tabcolsep) * \real{0.40}}@{}}
\toprule
名称 & Python 类型 & 默认值 & 含义 \\
\midrule
\endhead
\passthrough{\lstinline!other!} & \passthrough{\lstinline!object!} &
必填 & 用于比较的另一个 Python 对象。类型不兼容时比较结果通常为
\passthrough{\lstinline!False!}。 \\
\bottomrule
\end{longtable}

\textbf{返回值}

\begin{longtable}[]{@{}
  >{\raggedright\arraybackslash}p{(\columnwidth - 4\tabcolsep) * \real{0.18}}
  >{\raggedright\arraybackslash}p{(\columnwidth - 4\tabcolsep) * \real{0.20}}
  >{\raggedright\arraybackslash}p{(\columnwidth - 4\tabcolsep) * \real{0.62}}@{}}
\toprule
位置 & 类型 & 含义 \\
\midrule
\endhead
返回值 & \passthrough{\lstinline!bool!} & 返回
\passthrough{\lstinline!bool!}。更精确的业务含义见该方法用途、C++
签名和目标型号协议。 \\
\bottomrule
\end{longtable}

\textbf{对应 C++}

\begin{itemize}
\tightlist
\item
  类：\passthrough{\lstinline!geometry\_msgs::msg::dds\_::Pose\_!}
\item
  签名：\passthrough{\lstinline"operator!=(const value \&) const"}
\item
  绑定策略：\passthrough{\lstinline!未单独标注!}
\item
  声明位置：由 pybind11 手工绑定或生成注册表提供；manifest
  未记录独立头文件行号。
\end{itemize}

\textbf{说明}

\begin{itemize}
\tightlist
\item
  示例是局部调用片段；\passthrough{\lstinline!obj!}
  和参数变量表示已经按上层业务校验并准备好的对象。
\end{itemize}

\textbf{用法}

\begin{lstlisting}[language=Python]
different = left != right
\end{lstlisting}

\hypertarget{unitree-sdk2-cpp-idl-ros2-pose-position}{%
\paragraph{\texorpdfstring{\texttt{unitree\_sdk2\_cpp.idl.ros2.Pose.position}}{unitree\_sdk2\_cpp.idl.ros2.Pose.position}}\label{unitree-sdk2-cpp-idl-ros2-pose-position}}

底层 \passthrough{\lstinline!geometry\_msgs::msg::dds\_::Pose\_!} 的
\passthrough{\lstinline!position!} 字段。类型签名只说明 Python
表示；单位、数值范围、数组固定长度和枚举含义需查对应 IDL 头文件。

\textbf{签名}

\begin{lstlisting}[language=Python]
@property
def position(self) -> Point

@position.setter
def position(self, value: Point) -> None
\end{lstlisting}

\textbf{可用性与安全性}

\textbf{\passthrough{\lstinline!AVAILABLE!} /
\passthrough{\lstinline!VALUE\_TYPE!}}：当前绑定源码已实现该消息值操作。

\textbf{参数}

\begin{longtable}[]{@{}
  >{\raggedright\arraybackslash}p{(\columnwidth - 4\tabcolsep) * \real{0.18}}
  >{\raggedright\arraybackslash}p{(\columnwidth - 4\tabcolsep) * \real{0.20}}
  >{\raggedright\arraybackslash}p{(\columnwidth - 4\tabcolsep) * \real{0.62}}@{}}
\toprule
名称 & Python 类型 & 含义 \\
\midrule
\endhead
\passthrough{\lstinline!value!} & \passthrough{\lstinline!Point!} &
要写入或传递的新值，必须符合该参数的 Python 类型以及底层 C++
范围约束。 \\
\bottomrule
\end{longtable}

\textbf{返回值}

\begin{longtable}[]{@{}
  >{\raggedright\arraybackslash}p{(\columnwidth - 4\tabcolsep) * \real{0.18}}
  >{\raggedright\arraybackslash}p{(\columnwidth - 4\tabcolsep) * \real{0.20}}
  >{\raggedright\arraybackslash}p{(\columnwidth - 4\tabcolsep) * \real{0.62}}@{}}
\toprule
位置 & 类型 & 含义 \\
\midrule
\endhead
getter 返回值 & \passthrough{\lstinline!Point!} &
返回属性当前值。数组、vector 和嵌套消息采用复制语义。 \\
\bottomrule
\end{longtable}

\textbf{对应 C++}

\begin{itemize}
\tightlist
\item
  类：\passthrough{\lstinline!geometry\_msgs::msg::dds\_::Pose\_!}
\item
  字段类型：\passthrough{\lstinline!::geometry\_msgs::msg::dds\_::Point\_!}
\item
  声明位置：\passthrough{\lstinline!include/unitree/idl/ros2/Pose\_.hpp:26!}
\end{itemize}

\textbf{用法}

\begin{lstlisting}[language=Python]
current_value = obj.position
obj.position = new_value
\end{lstlisting}

\hypertarget{unitree-sdk2-cpp-idl-ros2-pose-orientation}{%
\paragraph{\texorpdfstring{\texttt{unitree\_sdk2\_cpp.idl.ros2.Pose.orientation}}{unitree\_sdk2\_cpp.idl.ros2.Pose.orientation}}\label{unitree-sdk2-cpp-idl-ros2-pose-orientation}}

底层 \passthrough{\lstinline!geometry\_msgs::msg::dds\_::Pose\_!} 的
\passthrough{\lstinline!orientation!} 字段。类型签名只说明 Python
表示；单位、数值范围、数组固定长度和枚举含义需查对应 IDL 头文件。

\textbf{签名}

\begin{lstlisting}[language=Python]
@property
def orientation(self) -> Quaternion

@orientation.setter
def orientation(self, value: Quaternion) -> None
\end{lstlisting}

\textbf{可用性与安全性}

\textbf{\passthrough{\lstinline!AVAILABLE!} /
\passthrough{\lstinline!VALUE\_TYPE!}}：当前绑定源码已实现该消息值操作。

\textbf{参数}

\begin{longtable}[]{@{}
  >{\raggedright\arraybackslash}p{(\columnwidth - 4\tabcolsep) * \real{0.18}}
  >{\raggedright\arraybackslash}p{(\columnwidth - 4\tabcolsep) * \real{0.20}}
  >{\raggedright\arraybackslash}p{(\columnwidth - 4\tabcolsep) * \real{0.62}}@{}}
\toprule
名称 & Python 类型 & 含义 \\
\midrule
\endhead
\passthrough{\lstinline!value!} & \passthrough{\lstinline!Quaternion!} &
要写入或传递的新值，必须符合该参数的 Python 类型以及底层 C++
范围约束。 \\
\bottomrule
\end{longtable}

\textbf{返回值}

\begin{longtable}[]{@{}
  >{\raggedright\arraybackslash}p{(\columnwidth - 4\tabcolsep) * \real{0.18}}
  >{\raggedright\arraybackslash}p{(\columnwidth - 4\tabcolsep) * \real{0.20}}
  >{\raggedright\arraybackslash}p{(\columnwidth - 4\tabcolsep) * \real{0.62}}@{}}
\toprule
位置 & 类型 & 含义 \\
\midrule
\endhead
getter 返回值 & \passthrough{\lstinline!Quaternion!} &
返回属性当前值。数组、vector 和嵌套消息采用复制语义。 \\
\bottomrule
\end{longtable}

\textbf{对应 C++}

\begin{itemize}
\tightlist
\item
  类：\passthrough{\lstinline!geometry\_msgs::msg::dds\_::Pose\_!}
\item
  字段类型：\passthrough{\lstinline!::geometry\_msgs::msg::dds\_::Quaternion\_!}
\item
  声明位置：\passthrough{\lstinline!include/unitree/idl/ros2/Pose\_.hpp:27!}
\end{itemize}

\textbf{用法}

\begin{lstlisting}[language=Python]
current_value = obj.orientation
obj.orientation = new_value
\end{lstlisting}

\hypertarget{unitree-sdk2-cpp-idl-ros2-mapmetadata}{%
\subsubsection{\texorpdfstring{\texttt{unitree\_sdk2\_cpp.idl.ros2.MapMetaData}}{unitree\_sdk2\_cpp.idl.ros2.MapMetaData}}\label{unitree-sdk2-cpp-idl-ros2-mapmetadata}}

可由 Python 读写的 DDS/IDL 消息值类型。字段采用复制语义。

\textbf{导入}

\begin{lstlisting}[language=Python]
from unitree_sdk2_cpp.idl.ros2 import MapMetaData
\end{lstlisting}

\textbf{方法索引}

\begin{longtable}[]{@{}
  >{\raggedright\arraybackslash}p{(\columnwidth - 6\tabcolsep) * \real{0.18}}
  >{\raggedright\arraybackslash}p{(\columnwidth - 6\tabcolsep) * \real{0.24}}
  >{\raggedright\arraybackslash}p{(\columnwidth - 6\tabcolsep) * \real{0.18}}
  >{\raggedright\arraybackslash}p{(\columnwidth - 6\tabcolsep) * \real{0.40}}@{}}
\toprule
方法 & Python 签名 & 状态 & 安全分类 \\
\midrule
\endhead
\protect\hyperlink{unitree-sdk2-cpp-idl-ros2-mapmetadata-dunder-init-1}{\passthrough{\lstinline!\_\_init\_\_!}}
& \passthrough{\lstinline!def \_\_init\_\_(self) -> None!} &
\passthrough{\lstinline!AVAILABLE!} &
\passthrough{\lstinline!VALUE\_TYPE!} \\
\protect\hyperlink{unitree-sdk2-cpp-idl-ros2-mapmetadata-dunder-eq-1}{\passthrough{\lstinline!\_\_eq\_\_!}}
& \passthrough{\lstinline!def \_\_eq\_\_(self, other: object) -> bool!}
& \passthrough{\lstinline!AVAILABLE!} &
\passthrough{\lstinline!VALUE\_TYPE!} \\
\protect\hyperlink{unitree-sdk2-cpp-idl-ros2-mapmetadata-dunder-ne-1}{\passthrough{\lstinline!\_\_ne\_\_!}}
& \passthrough{\lstinline!def \_\_ne\_\_(self, other: object) -> bool!}
& \passthrough{\lstinline!AVAILABLE!} &
\passthrough{\lstinline!VALUE\_TYPE!} \\
\bottomrule
\end{longtable}

\textbf{属性索引}

\begin{longtable}[]{@{}
  >{\raggedright\arraybackslash}p{(\columnwidth - 6\tabcolsep) * \real{0.18}}
  >{\raggedright\arraybackslash}p{(\columnwidth - 6\tabcolsep) * \real{0.24}}
  >{\raggedright\arraybackslash}p{(\columnwidth - 6\tabcolsep) * \real{0.18}}
  >{\raggedright\arraybackslash}p{(\columnwidth - 6\tabcolsep) * \real{0.40}}@{}}
\toprule
属性 & 读取类型 & 写入类型 & 状态 \\
\midrule
\endhead
\protect\hyperlink{unitree-sdk2-cpp-idl-ros2-mapmetadata-map-load-time}{\passthrough{\lstinline!map\_load\_time!}}
& \passthrough{\lstinline!Any!} & \passthrough{\lstinline!Any!} &
\passthrough{\lstinline!AVAILABLE!} \\
\protect\hyperlink{unitree-sdk2-cpp-idl-ros2-mapmetadata-resolution}{\passthrough{\lstinline!resolution!}}
& \passthrough{\lstinline!float!} & \passthrough{\lstinline!float!} &
\passthrough{\lstinline!AVAILABLE!} \\
\protect\hyperlink{unitree-sdk2-cpp-idl-ros2-mapmetadata-width}{\passthrough{\lstinline!width!}}
& \passthrough{\lstinline!int!} & \passthrough{\lstinline!int!} &
\passthrough{\lstinline!AVAILABLE!} \\
\protect\hyperlink{unitree-sdk2-cpp-idl-ros2-mapmetadata-height}{\passthrough{\lstinline!height!}}
& \passthrough{\lstinline!int!} & \passthrough{\lstinline!int!} &
\passthrough{\lstinline!AVAILABLE!} \\
\protect\hyperlink{unitree-sdk2-cpp-idl-ros2-mapmetadata-origin}{\passthrough{\lstinline!origin!}}
& \passthrough{\lstinline!Any!} & \passthrough{\lstinline!Any!} &
\passthrough{\lstinline!AVAILABLE!} \\
\bottomrule
\end{longtable}

\hypertarget{unitree-sdk2-cpp-idl-ros2-mapmetadata-dunder-init-1}{%
\paragraph{\texorpdfstring{\texttt{unitree\_sdk2\_cpp.idl.ros2.MapMetaData.\_\_init\_\_}}{unitree\_sdk2\_cpp.idl.ros2.MapMetaData.\_\_init\_\_}}\label{unitree-sdk2-cpp-idl-ros2-mapmetadata-dunder-init-1}}

初始化 \passthrough{\lstinline!MapMetaData!}
实例。是否能实际构造取决于下方可用性状态。

\textbf{签名}

\begin{lstlisting}[language=Python]
def __init__(self) -> None
\end{lstlisting}

\textbf{可用性与安全性}

\textbf{\passthrough{\lstinline!AVAILABLE!} /
\passthrough{\lstinline!VALUE\_TYPE!}}：当前绑定源码已实现该消息值操作。

\textbf{参数}

无显式参数。实例方法中的 \passthrough{\lstinline!self!} 由 Python
自动传入。

\textbf{返回值}

\begin{longtable}[]{@{}
  >{\raggedright\arraybackslash}p{(\columnwidth - 4\tabcolsep) * \real{0.18}}
  >{\raggedright\arraybackslash}p{(\columnwidth - 4\tabcolsep) * \real{0.20}}
  >{\raggedright\arraybackslash}p{(\columnwidth - 4\tabcolsep) * \real{0.62}}@{}}
\toprule
位置 & 类型 & 含义 \\
\midrule
\endhead
无 & \passthrough{\lstinline!None!} &
\passthrough{\lstinline!\_\_init\_\_!}
本身不返回值；调用类对象时，在构造函数可用的前提下得到该类实例。 \\
\bottomrule
\end{longtable}

\textbf{对应 C++}

\begin{itemize}
\tightlist
\item
  类：\passthrough{\lstinline!nav\_msgs::msg::dds\_::MapMetaData\_!}
\item
  签名：\passthrough{\lstinline!MapMetaData\_()!}
\item
  绑定策略：\passthrough{\lstinline!未单独标注!}
\item
  声明位置：\passthrough{\lstinline!include/unitree/idl/ros2/MapMetaData\_.hpp:34!}
\end{itemize}

\textbf{说明}

\begin{itemize}
\tightlist
\item
  示例是局部调用片段；\passthrough{\lstinline!obj!}
  和参数变量表示已经按上层业务校验并准备好的对象。
\end{itemize}

\textbf{用法}

\begin{lstlisting}[language=Python]
value = MapMetaData()
\end{lstlisting}

\hypertarget{unitree-sdk2-cpp-idl-ros2-mapmetadata-dunder-eq-1}{%
\paragraph{\texorpdfstring{\texttt{unitree\_sdk2\_cpp.idl.ros2.MapMetaData.\_\_eq\_\_}}{unitree\_sdk2\_cpp.idl.ros2.MapMetaData.\_\_eq\_\_}}\label{unitree-sdk2-cpp-idl-ros2-mapmetadata-dunder-eq-1}}

按底层 IDL 消息内容比较两个对象是否相等。

\textbf{签名}

\begin{lstlisting}[language=Python]
def __eq__(self, other: object) -> bool
\end{lstlisting}

\textbf{可用性与安全性}

\textbf{\passthrough{\lstinline!AVAILABLE!} /
\passthrough{\lstinline!VALUE\_TYPE!}}：当前绑定源码已实现该消息值操作。

\textbf{参数}

\begin{longtable}[]{@{}
  >{\raggedright\arraybackslash}p{(\columnwidth - 6\tabcolsep) * \real{0.18}}
  >{\raggedright\arraybackslash}p{(\columnwidth - 6\tabcolsep) * \real{0.24}}
  >{\raggedright\arraybackslash}p{(\columnwidth - 6\tabcolsep) * \real{0.18}}
  >{\raggedright\arraybackslash}p{(\columnwidth - 6\tabcolsep) * \real{0.40}}@{}}
\toprule
名称 & Python 类型 & 默认值 & 含义 \\
\midrule
\endhead
\passthrough{\lstinline!other!} & \passthrough{\lstinline!object!} &
必填 & 用于比较的另一个 Python 对象。类型不兼容时比较结果通常为
\passthrough{\lstinline!False!}。 \\
\bottomrule
\end{longtable}

\textbf{返回值}

\begin{longtable}[]{@{}
  >{\raggedright\arraybackslash}p{(\columnwidth - 4\tabcolsep) * \real{0.18}}
  >{\raggedright\arraybackslash}p{(\columnwidth - 4\tabcolsep) * \real{0.20}}
  >{\raggedright\arraybackslash}p{(\columnwidth - 4\tabcolsep) * \real{0.62}}@{}}
\toprule
位置 & 类型 & 含义 \\
\midrule
\endhead
返回值 & \passthrough{\lstinline!bool!} & 返回
\passthrough{\lstinline!bool!}。更精确的业务含义见该方法用途、C++
签名和目标型号协议。 \\
\bottomrule
\end{longtable}

\textbf{对应 C++}

\begin{itemize}
\tightlist
\item
  类：\passthrough{\lstinline!nav\_msgs::msg::dds\_::MapMetaData\_!}
\item
  签名：\passthrough{\lstinline!operator==(const value \&) const!}
\item
  绑定策略：\passthrough{\lstinline!未单独标注!}
\item
  声明位置：由 pybind11 手工绑定或生成注册表提供；manifest
  未记录独立头文件行号。
\end{itemize}

\textbf{说明}

\begin{itemize}
\tightlist
\item
  示例是局部调用片段；\passthrough{\lstinline!obj!}
  和参数变量表示已经按上层业务校验并准备好的对象。
\end{itemize}

\textbf{用法}

\begin{lstlisting}[language=Python]
same = left == right
\end{lstlisting}

\hypertarget{unitree-sdk2-cpp-idl-ros2-mapmetadata-dunder-ne-1}{%
\paragraph{\texorpdfstring{\texttt{unitree\_sdk2\_cpp.idl.ros2.MapMetaData.\_\_ne\_\_}}{unitree\_sdk2\_cpp.idl.ros2.MapMetaData.\_\_ne\_\_}}\label{unitree-sdk2-cpp-idl-ros2-mapmetadata-dunder-ne-1}}

按底层 IDL 消息内容比较两个对象是否不相等。

\textbf{签名}

\begin{lstlisting}[language=Python]
def __ne__(self, other: object) -> bool
\end{lstlisting}

\textbf{可用性与安全性}

\textbf{\passthrough{\lstinline!AVAILABLE!} /
\passthrough{\lstinline!VALUE\_TYPE!}}：当前绑定源码已实现该消息值操作。

\textbf{参数}

\begin{longtable}[]{@{}
  >{\raggedright\arraybackslash}p{(\columnwidth - 6\tabcolsep) * \real{0.18}}
  >{\raggedright\arraybackslash}p{(\columnwidth - 6\tabcolsep) * \real{0.24}}
  >{\raggedright\arraybackslash}p{(\columnwidth - 6\tabcolsep) * \real{0.18}}
  >{\raggedright\arraybackslash}p{(\columnwidth - 6\tabcolsep) * \real{0.40}}@{}}
\toprule
名称 & Python 类型 & 默认值 & 含义 \\
\midrule
\endhead
\passthrough{\lstinline!other!} & \passthrough{\lstinline!object!} &
必填 & 用于比较的另一个 Python 对象。类型不兼容时比较结果通常为
\passthrough{\lstinline!False!}。 \\
\bottomrule
\end{longtable}

\textbf{返回值}

\begin{longtable}[]{@{}
  >{\raggedright\arraybackslash}p{(\columnwidth - 4\tabcolsep) * \real{0.18}}
  >{\raggedright\arraybackslash}p{(\columnwidth - 4\tabcolsep) * \real{0.20}}
  >{\raggedright\arraybackslash}p{(\columnwidth - 4\tabcolsep) * \real{0.62}}@{}}
\toprule
位置 & 类型 & 含义 \\
\midrule
\endhead
返回值 & \passthrough{\lstinline!bool!} & 返回
\passthrough{\lstinline!bool!}。更精确的业务含义见该方法用途、C++
签名和目标型号协议。 \\
\bottomrule
\end{longtable}

\textbf{对应 C++}

\begin{itemize}
\tightlist
\item
  类：\passthrough{\lstinline!nav\_msgs::msg::dds\_::MapMetaData\_!}
\item
  签名：\passthrough{\lstinline"operator!=(const value \&) const"}
\item
  绑定策略：\passthrough{\lstinline!未单独标注!}
\item
  声明位置：由 pybind11 手工绑定或生成注册表提供；manifest
  未记录独立头文件行号。
\end{itemize}

\textbf{说明}

\begin{itemize}
\tightlist
\item
  示例是局部调用片段；\passthrough{\lstinline!obj!}
  和参数变量表示已经按上层业务校验并准备好的对象。
\end{itemize}

\textbf{用法}

\begin{lstlisting}[language=Python]
different = left != right
\end{lstlisting}

\hypertarget{unitree-sdk2-cpp-idl-ros2-mapmetadata-map-load-time}{%
\paragraph{\texorpdfstring{\texttt{unitree\_sdk2\_cpp.idl.ros2.MapMetaData.map\_load\_time}}{unitree\_sdk2\_cpp.idl.ros2.MapMetaData.map\_load\_time}}\label{unitree-sdk2-cpp-idl-ros2-mapmetadata-map-load-time}}

底层 \passthrough{\lstinline!nav\_msgs::msg::dds\_::MapMetaData\_!} 的
\passthrough{\lstinline!map\_load\_time!} 字段。类型签名只说明 Python
表示；单位、数值范围、数组固定长度和枚举含义需查对应 IDL 头文件。

\textbf{签名}

\begin{lstlisting}[language=Python]
@property
def map_load_time(self) -> Any

@map_load_time.setter
def map_load_time(self, value: Any) -> None
\end{lstlisting}

\textbf{可用性与安全性}

\textbf{\passthrough{\lstinline!AVAILABLE!} /
\passthrough{\lstinline!VALUE\_TYPE!}}：当前绑定源码已实现该消息值操作。

\textbf{参数}

\begin{longtable}[]{@{}
  >{\raggedright\arraybackslash}p{(\columnwidth - 4\tabcolsep) * \real{0.18}}
  >{\raggedright\arraybackslash}p{(\columnwidth - 4\tabcolsep) * \real{0.20}}
  >{\raggedright\arraybackslash}p{(\columnwidth - 4\tabcolsep) * \real{0.62}}@{}}
\toprule
名称 & Python 类型 & 含义 \\
\midrule
\endhead
\passthrough{\lstinline!value!} & \passthrough{\lstinline!Any!} &
要写入或传递的新值，必须符合该参数的 Python 类型以及底层 C++
范围约束。 \\
\bottomrule
\end{longtable}

\textbf{返回值}

\begin{longtable}[]{@{}
  >{\raggedright\arraybackslash}p{(\columnwidth - 4\tabcolsep) * \real{0.18}}
  >{\raggedright\arraybackslash}p{(\columnwidth - 4\tabcolsep) * \real{0.20}}
  >{\raggedright\arraybackslash}p{(\columnwidth - 4\tabcolsep) * \real{0.62}}@{}}
\toprule
位置 & 类型 & 含义 \\
\midrule
\endhead
getter 返回值 & \passthrough{\lstinline!Any!} & 返回属性当前值。 \\
\bottomrule
\end{longtable}

\textbf{对应 C++}

\begin{itemize}
\tightlist
\item
  类：\passthrough{\lstinline!nav\_msgs::msg::dds\_::MapMetaData\_!}
\item
  字段类型：\passthrough{\lstinline!::builtin\_interfaces::msg::dds\_::Time\_!}
\item
  声明位置：\passthrough{\lstinline!include/unitree/idl/ros2/MapMetaData\_.hpp:27!}
\end{itemize}

\textbf{用法}

\begin{lstlisting}[language=Python]
current_value = obj.map_load_time
obj.map_load_time = new_value
\end{lstlisting}

\hypertarget{unitree-sdk2-cpp-idl-ros2-mapmetadata-resolution}{%
\paragraph{\texorpdfstring{\texttt{unitree\_sdk2\_cpp.idl.ros2.MapMetaData.resolution}}{unitree\_sdk2\_cpp.idl.ros2.MapMetaData.resolution}}\label{unitree-sdk2-cpp-idl-ros2-mapmetadata-resolution}}

底层 \passthrough{\lstinline!nav\_msgs::msg::dds\_::MapMetaData\_!} 的
\passthrough{\lstinline!resolution!} 字段。类型签名只说明 Python
表示；单位、数值范围、数组固定长度和枚举含义需查对应 IDL 头文件。
底层为浮点数；类型本身不说明单位、坐标系或业务安全范围。

\textbf{签名}

\begin{lstlisting}[language=Python]
@property
def resolution(self) -> float

@resolution.setter
def resolution(self, value: float) -> None
\end{lstlisting}

\textbf{可用性与安全性}

\textbf{\passthrough{\lstinline!AVAILABLE!} /
\passthrough{\lstinline!VALUE\_TYPE!}}：当前绑定源码已实现该消息值操作。

\textbf{参数}

\begin{longtable}[]{@{}
  >{\raggedright\arraybackslash}p{(\columnwidth - 4\tabcolsep) * \real{0.18}}
  >{\raggedright\arraybackslash}p{(\columnwidth - 4\tabcolsep) * \real{0.20}}
  >{\raggedright\arraybackslash}p{(\columnwidth - 4\tabcolsep) * \real{0.62}}@{}}
\toprule
名称 & Python 类型 & 含义 \\
\midrule
\endhead
\passthrough{\lstinline!value!} & \passthrough{\lstinline!float!} &
要写入或传递的新值，必须符合该参数的 Python 类型以及底层 C++
范围约束。 \\
\bottomrule
\end{longtable}

\textbf{返回值}

\begin{longtable}[]{@{}
  >{\raggedright\arraybackslash}p{(\columnwidth - 4\tabcolsep) * \real{0.18}}
  >{\raggedright\arraybackslash}p{(\columnwidth - 4\tabcolsep) * \real{0.20}}
  >{\raggedright\arraybackslash}p{(\columnwidth - 4\tabcolsep) * \real{0.62}}@{}}
\toprule
位置 & 类型 & 含义 \\
\midrule
\endhead
getter 返回值 & \passthrough{\lstinline!float!} & 返回属性当前值。 \\
\bottomrule
\end{longtable}

\textbf{对应 C++}

\begin{itemize}
\tightlist
\item
  类：\passthrough{\lstinline!nav\_msgs::msg::dds\_::MapMetaData\_!}
\item
  字段类型：\passthrough{\lstinline!float!}
\item
  声明位置：\passthrough{\lstinline!include/unitree/idl/ros2/MapMetaData\_.hpp:28!}
\end{itemize}

\textbf{用法}

\begin{lstlisting}[language=Python]
current_value = obj.resolution
obj.resolution = new_value
\end{lstlisting}

\hypertarget{unitree-sdk2-cpp-idl-ros2-mapmetadata-width}{%
\paragraph{\texorpdfstring{\texttt{unitree\_sdk2\_cpp.idl.ros2.MapMetaData.width}}{unitree\_sdk2\_cpp.idl.ros2.MapMetaData.width}}\label{unitree-sdk2-cpp-idl-ros2-mapmetadata-width}}

底层 \passthrough{\lstinline!nav\_msgs::msg::dds\_::MapMetaData\_!} 的
\passthrough{\lstinline!width!} 字段。类型签名只说明 Python
表示；单位、数值范围、数组固定长度和枚举含义需查对应 IDL 头文件。
取值范围为 0 到 4294967295。

\textbf{签名}

\begin{lstlisting}[language=Python]
@property
def width(self) -> int

@width.setter
def width(self, value: int) -> None
\end{lstlisting}

\textbf{可用性与安全性}

\textbf{\passthrough{\lstinline!AVAILABLE!} /
\passthrough{\lstinline!VALUE\_TYPE!}}：当前绑定源码已实现该消息值操作。

\textbf{参数}

\begin{longtable}[]{@{}
  >{\raggedright\arraybackslash}p{(\columnwidth - 4\tabcolsep) * \real{0.18}}
  >{\raggedright\arraybackslash}p{(\columnwidth - 4\tabcolsep) * \real{0.20}}
  >{\raggedright\arraybackslash}p{(\columnwidth - 4\tabcolsep) * \real{0.62}}@{}}
\toprule
名称 & Python 类型 & 含义 \\
\midrule
\endhead
\passthrough{\lstinline!value!} & \passthrough{\lstinline!int!} &
要写入或传递的新值，必须符合该参数的 Python 类型以及底层 C++
范围约束。 \\
\bottomrule
\end{longtable}

\textbf{返回值}

\begin{longtable}[]{@{}
  >{\raggedright\arraybackslash}p{(\columnwidth - 4\tabcolsep) * \real{0.18}}
  >{\raggedright\arraybackslash}p{(\columnwidth - 4\tabcolsep) * \real{0.20}}
  >{\raggedright\arraybackslash}p{(\columnwidth - 4\tabcolsep) * \real{0.62}}@{}}
\toprule
位置 & 类型 & 含义 \\
\midrule
\endhead
getter 返回值 & \passthrough{\lstinline!int!} & 返回属性当前值。 \\
\bottomrule
\end{longtable}

\textbf{对应 C++}

\begin{itemize}
\tightlist
\item
  类：\passthrough{\lstinline!nav\_msgs::msg::dds\_::MapMetaData\_!}
\item
  字段类型：\passthrough{\lstinline!uint32\_t!}
\item
  声明位置：\passthrough{\lstinline!include/unitree/idl/ros2/MapMetaData\_.hpp:29!}
\end{itemize}

\textbf{用法}

\begin{lstlisting}[language=Python]
current_value = obj.width
obj.width = new_value
\end{lstlisting}

\hypertarget{unitree-sdk2-cpp-idl-ros2-mapmetadata-height}{%
\paragraph{\texorpdfstring{\texttt{unitree\_sdk2\_cpp.idl.ros2.MapMetaData.height}}{unitree\_sdk2\_cpp.idl.ros2.MapMetaData.height}}\label{unitree-sdk2-cpp-idl-ros2-mapmetadata-height}}

底层 \passthrough{\lstinline!nav\_msgs::msg::dds\_::MapMetaData\_!} 的
\passthrough{\lstinline!height!} 字段。类型签名只说明 Python
表示；单位、数值范围、数组固定长度和枚举含义需查对应 IDL 头文件。
取值范围为 0 到 4294967295。

\textbf{签名}

\begin{lstlisting}[language=Python]
@property
def height(self) -> int

@height.setter
def height(self, value: int) -> None
\end{lstlisting}

\textbf{可用性与安全性}

\textbf{\passthrough{\lstinline!AVAILABLE!} /
\passthrough{\lstinline!VALUE\_TYPE!}}：当前绑定源码已实现该消息值操作。

\textbf{参数}

\begin{longtable}[]{@{}
  >{\raggedright\arraybackslash}p{(\columnwidth - 4\tabcolsep) * \real{0.18}}
  >{\raggedright\arraybackslash}p{(\columnwidth - 4\tabcolsep) * \real{0.20}}
  >{\raggedright\arraybackslash}p{(\columnwidth - 4\tabcolsep) * \real{0.62}}@{}}
\toprule
名称 & Python 类型 & 含义 \\
\midrule
\endhead
\passthrough{\lstinline!value!} & \passthrough{\lstinline!int!} &
要写入或传递的新值，必须符合该参数的 Python 类型以及底层 C++
范围约束。 \\
\bottomrule
\end{longtable}

\textbf{返回值}

\begin{longtable}[]{@{}
  >{\raggedright\arraybackslash}p{(\columnwidth - 4\tabcolsep) * \real{0.18}}
  >{\raggedright\arraybackslash}p{(\columnwidth - 4\tabcolsep) * \real{0.20}}
  >{\raggedright\arraybackslash}p{(\columnwidth - 4\tabcolsep) * \real{0.62}}@{}}
\toprule
位置 & 类型 & 含义 \\
\midrule
\endhead
getter 返回值 & \passthrough{\lstinline!int!} & 返回属性当前值。 \\
\bottomrule
\end{longtable}

\textbf{对应 C++}

\begin{itemize}
\tightlist
\item
  类：\passthrough{\lstinline!nav\_msgs::msg::dds\_::MapMetaData\_!}
\item
  字段类型：\passthrough{\lstinline!uint32\_t!}
\item
  声明位置：\passthrough{\lstinline!include/unitree/idl/ros2/MapMetaData\_.hpp:30!}
\end{itemize}

\textbf{用法}

\begin{lstlisting}[language=Python]
current_value = obj.height
obj.height = new_value
\end{lstlisting}

\hypertarget{unitree-sdk2-cpp-idl-ros2-mapmetadata-origin}{%
\paragraph{\texorpdfstring{\texttt{unitree\_sdk2\_cpp.idl.ros2.MapMetaData.origin}}{unitree\_sdk2\_cpp.idl.ros2.MapMetaData.origin}}\label{unitree-sdk2-cpp-idl-ros2-mapmetadata-origin}}

底层 \passthrough{\lstinline!nav\_msgs::msg::dds\_::MapMetaData\_!} 的
\passthrough{\lstinline!origin!} 字段。类型签名只说明 Python
表示；单位、数值范围、数组固定长度和枚举含义需查对应 IDL 头文件。

\textbf{签名}

\begin{lstlisting}[language=Python]
@property
def origin(self) -> Any

@origin.setter
def origin(self, value: Any) -> None
\end{lstlisting}

\textbf{可用性与安全性}

\textbf{\passthrough{\lstinline!AVAILABLE!} /
\passthrough{\lstinline!VALUE\_TYPE!}}：当前绑定源码已实现该消息值操作。

\textbf{参数}

\begin{longtable}[]{@{}
  >{\raggedright\arraybackslash}p{(\columnwidth - 4\tabcolsep) * \real{0.18}}
  >{\raggedright\arraybackslash}p{(\columnwidth - 4\tabcolsep) * \real{0.20}}
  >{\raggedright\arraybackslash}p{(\columnwidth - 4\tabcolsep) * \real{0.62}}@{}}
\toprule
名称 & Python 类型 & 含义 \\
\midrule
\endhead
\passthrough{\lstinline!value!} & \passthrough{\lstinline!Any!} &
要写入或传递的新值，必须符合该参数的 Python 类型以及底层 C++
范围约束。 \\
\bottomrule
\end{longtable}

\textbf{返回值}

\begin{longtable}[]{@{}
  >{\raggedright\arraybackslash}p{(\columnwidth - 4\tabcolsep) * \real{0.18}}
  >{\raggedright\arraybackslash}p{(\columnwidth - 4\tabcolsep) * \real{0.20}}
  >{\raggedright\arraybackslash}p{(\columnwidth - 4\tabcolsep) * \real{0.62}}@{}}
\toprule
位置 & 类型 & 含义 \\
\midrule
\endhead
getter 返回值 & \passthrough{\lstinline!Any!} & 返回属性当前值。 \\
\bottomrule
\end{longtable}

\textbf{对应 C++}

\begin{itemize}
\tightlist
\item
  类：\passthrough{\lstinline!nav\_msgs::msg::dds\_::MapMetaData\_!}
\item
  字段类型：\passthrough{\lstinline!::geometry\_msgs::msg::dds\_::Pose\_!}
\item
  声明位置：\passthrough{\lstinline!include/unitree/idl/ros2/MapMetaData\_.hpp:31!}
\end{itemize}

\textbf{用法}

\begin{lstlisting}[language=Python]
current_value = obj.origin
obj.origin = new_value
\end{lstlisting}

\hypertarget{unitree-sdk2-cpp-idl-ros2-occupancygrid}{%
\subsubsection{\texorpdfstring{\texttt{unitree\_sdk2\_cpp.idl.ros2.OccupancyGrid}}{unitree\_sdk2\_cpp.idl.ros2.OccupancyGrid}}\label{unitree-sdk2-cpp-idl-ros2-occupancygrid}}

可由 Python 读写的 DDS/IDL 消息值类型。字段采用复制语义。

\textbf{导入}

\begin{lstlisting}[language=Python]
from unitree_sdk2_cpp.idl.ros2 import OccupancyGrid
\end{lstlisting}

\textbf{方法索引}

\begin{longtable}[]{@{}
  >{\raggedright\arraybackslash}p{(\columnwidth - 6\tabcolsep) * \real{0.18}}
  >{\raggedright\arraybackslash}p{(\columnwidth - 6\tabcolsep) * \real{0.24}}
  >{\raggedright\arraybackslash}p{(\columnwidth - 6\tabcolsep) * \real{0.18}}
  >{\raggedright\arraybackslash}p{(\columnwidth - 6\tabcolsep) * \real{0.40}}@{}}
\toprule
方法 & Python 签名 & 状态 & 安全分类 \\
\midrule
\endhead
\protect\hyperlink{unitree-sdk2-cpp-idl-ros2-occupancygrid-dunder-init-1}{\passthrough{\lstinline!\_\_init\_\_!}}
& \passthrough{\lstinline!def \_\_init\_\_(self) -> None!} &
\passthrough{\lstinline!AVAILABLE!} &
\passthrough{\lstinline!VALUE\_TYPE!} \\
\protect\hyperlink{unitree-sdk2-cpp-idl-ros2-occupancygrid-dunder-eq-1}{\passthrough{\lstinline!\_\_eq\_\_!}}
& \passthrough{\lstinline!def \_\_eq\_\_(self, other: object) -> bool!}
& \passthrough{\lstinline!AVAILABLE!} &
\passthrough{\lstinline!VALUE\_TYPE!} \\
\protect\hyperlink{unitree-sdk2-cpp-idl-ros2-occupancygrid-dunder-ne-1}{\passthrough{\lstinline!\_\_ne\_\_!}}
& \passthrough{\lstinline!def \_\_ne\_\_(self, other: object) -> bool!}
& \passthrough{\lstinline!AVAILABLE!} &
\passthrough{\lstinline!VALUE\_TYPE!} \\
\bottomrule
\end{longtable}

\textbf{属性索引}

\begin{longtable}[]{@{}
  >{\raggedright\arraybackslash}p{(\columnwidth - 6\tabcolsep) * \real{0.18}}
  >{\raggedright\arraybackslash}p{(\columnwidth - 6\tabcolsep) * \real{0.24}}
  >{\raggedright\arraybackslash}p{(\columnwidth - 6\tabcolsep) * \real{0.18}}
  >{\raggedright\arraybackslash}p{(\columnwidth - 6\tabcolsep) * \real{0.40}}@{}}
\toprule
属性 & 读取类型 & 写入类型 & 状态 \\
\midrule
\endhead
\protect\hyperlink{unitree-sdk2-cpp-idl-ros2-occupancygrid-header}{\passthrough{\lstinline!header!}}
& \passthrough{\lstinline!Any!} & \passthrough{\lstinline!Any!} &
\passthrough{\lstinline!AVAILABLE!} \\
\protect\hyperlink{unitree-sdk2-cpp-idl-ros2-occupancygrid-info}{\passthrough{\lstinline!info!}}
& \passthrough{\lstinline!MapMetaData!} &
\passthrough{\lstinline!MapMetaData!} &
\passthrough{\lstinline!AVAILABLE!} \\
\protect\hyperlink{unitree-sdk2-cpp-idl-ros2-occupancygrid-data}{\passthrough{\lstinline!data!}}
& \passthrough{\lstinline!list[int]!} &
\passthrough{\lstinline!Sequence[int]!} &
\passthrough{\lstinline!AVAILABLE!} \\
\bottomrule
\end{longtable}

\hypertarget{unitree-sdk2-cpp-idl-ros2-occupancygrid-dunder-init-1}{%
\paragraph{\texorpdfstring{\texttt{unitree\_sdk2\_cpp.idl.ros2.OccupancyGrid.\_\_init\_\_}}{unitree\_sdk2\_cpp.idl.ros2.OccupancyGrid.\_\_init\_\_}}\label{unitree-sdk2-cpp-idl-ros2-occupancygrid-dunder-init-1}}

初始化 \passthrough{\lstinline!OccupancyGrid!}
实例。是否能实际构造取决于下方可用性状态。

\textbf{签名}

\begin{lstlisting}[language=Python]
def __init__(self) -> None
\end{lstlisting}

\textbf{可用性与安全性}

\textbf{\passthrough{\lstinline!AVAILABLE!} /
\passthrough{\lstinline!VALUE\_TYPE!}}：当前绑定源码已实现该消息值操作。

\textbf{参数}

无显式参数。实例方法中的 \passthrough{\lstinline!self!} 由 Python
自动传入。

\textbf{返回值}

\begin{longtable}[]{@{}
  >{\raggedright\arraybackslash}p{(\columnwidth - 4\tabcolsep) * \real{0.18}}
  >{\raggedright\arraybackslash}p{(\columnwidth - 4\tabcolsep) * \real{0.20}}
  >{\raggedright\arraybackslash}p{(\columnwidth - 4\tabcolsep) * \real{0.62}}@{}}
\toprule
位置 & 类型 & 含义 \\
\midrule
\endhead
无 & \passthrough{\lstinline!None!} &
\passthrough{\lstinline!\_\_init\_\_!}
本身不返回值；调用类对象时，在构造函数可用的前提下得到该类实例。 \\
\bottomrule
\end{longtable}

\textbf{对应 C++}

\begin{itemize}
\tightlist
\item
  类：\passthrough{\lstinline!nav\_msgs::msg::dds\_::OccupancyGrid\_!}
\item
  签名：\passthrough{\lstinline!OccupancyGrid\_()!}
\item
  绑定策略：\passthrough{\lstinline!未单独标注!}
\item
  声明位置：\passthrough{\lstinline!include/unitree/idl/ros2/OccupancyGrid\_.hpp:33!}
\end{itemize}

\textbf{说明}

\begin{itemize}
\tightlist
\item
  示例是局部调用片段；\passthrough{\lstinline!obj!}
  和参数变量表示已经按上层业务校验并准备好的对象。
\end{itemize}

\textbf{用法}

\begin{lstlisting}[language=Python]
value = OccupancyGrid()
\end{lstlisting}

\hypertarget{unitree-sdk2-cpp-idl-ros2-occupancygrid-dunder-eq-1}{%
\paragraph{\texorpdfstring{\texttt{unitree\_sdk2\_cpp.idl.ros2.OccupancyGrid.\_\_eq\_\_}}{unitree\_sdk2\_cpp.idl.ros2.OccupancyGrid.\_\_eq\_\_}}\label{unitree-sdk2-cpp-idl-ros2-occupancygrid-dunder-eq-1}}

按底层 IDL 消息内容比较两个对象是否相等。

\textbf{签名}

\begin{lstlisting}[language=Python]
def __eq__(self, other: object) -> bool
\end{lstlisting}

\textbf{可用性与安全性}

\textbf{\passthrough{\lstinline!AVAILABLE!} /
\passthrough{\lstinline!VALUE\_TYPE!}}：当前绑定源码已实现该消息值操作。

\textbf{参数}

\begin{longtable}[]{@{}
  >{\raggedright\arraybackslash}p{(\columnwidth - 6\tabcolsep) * \real{0.18}}
  >{\raggedright\arraybackslash}p{(\columnwidth - 6\tabcolsep) * \real{0.24}}
  >{\raggedright\arraybackslash}p{(\columnwidth - 6\tabcolsep) * \real{0.18}}
  >{\raggedright\arraybackslash}p{(\columnwidth - 6\tabcolsep) * \real{0.40}}@{}}
\toprule
名称 & Python 类型 & 默认值 & 含义 \\
\midrule
\endhead
\passthrough{\lstinline!other!} & \passthrough{\lstinline!object!} &
必填 & 用于比较的另一个 Python 对象。类型不兼容时比较结果通常为
\passthrough{\lstinline!False!}。 \\
\bottomrule
\end{longtable}

\textbf{返回值}

\begin{longtable}[]{@{}
  >{\raggedright\arraybackslash}p{(\columnwidth - 4\tabcolsep) * \real{0.18}}
  >{\raggedright\arraybackslash}p{(\columnwidth - 4\tabcolsep) * \real{0.20}}
  >{\raggedright\arraybackslash}p{(\columnwidth - 4\tabcolsep) * \real{0.62}}@{}}
\toprule
位置 & 类型 & 含义 \\
\midrule
\endhead
返回值 & \passthrough{\lstinline!bool!} & 返回
\passthrough{\lstinline!bool!}。更精确的业务含义见该方法用途、C++
签名和目标型号协议。 \\
\bottomrule
\end{longtable}

\textbf{对应 C++}

\begin{itemize}
\tightlist
\item
  类：\passthrough{\lstinline!nav\_msgs::msg::dds\_::OccupancyGrid\_!}
\item
  签名：\passthrough{\lstinline!operator==(const value \&) const!}
\item
  绑定策略：\passthrough{\lstinline!未单独标注!}
\item
  声明位置：由 pybind11 手工绑定或生成注册表提供；manifest
  未记录独立头文件行号。
\end{itemize}

\textbf{说明}

\begin{itemize}
\tightlist
\item
  示例是局部调用片段；\passthrough{\lstinline!obj!}
  和参数变量表示已经按上层业务校验并准备好的对象。
\end{itemize}

\textbf{用法}

\begin{lstlisting}[language=Python]
same = left == right
\end{lstlisting}

\hypertarget{unitree-sdk2-cpp-idl-ros2-occupancygrid-dunder-ne-1}{%
\paragraph{\texorpdfstring{\texttt{unitree\_sdk2\_cpp.idl.ros2.OccupancyGrid.\_\_ne\_\_}}{unitree\_sdk2\_cpp.idl.ros2.OccupancyGrid.\_\_ne\_\_}}\label{unitree-sdk2-cpp-idl-ros2-occupancygrid-dunder-ne-1}}

按底层 IDL 消息内容比较两个对象是否不相等。

\textbf{签名}

\begin{lstlisting}[language=Python]
def __ne__(self, other: object) -> bool
\end{lstlisting}

\textbf{可用性与安全性}

\textbf{\passthrough{\lstinline!AVAILABLE!} /
\passthrough{\lstinline!VALUE\_TYPE!}}：当前绑定源码已实现该消息值操作。

\textbf{参数}

\begin{longtable}[]{@{}
  >{\raggedright\arraybackslash}p{(\columnwidth - 6\tabcolsep) * \real{0.18}}
  >{\raggedright\arraybackslash}p{(\columnwidth - 6\tabcolsep) * \real{0.24}}
  >{\raggedright\arraybackslash}p{(\columnwidth - 6\tabcolsep) * \real{0.18}}
  >{\raggedright\arraybackslash}p{(\columnwidth - 6\tabcolsep) * \real{0.40}}@{}}
\toprule
名称 & Python 类型 & 默认值 & 含义 \\
\midrule
\endhead
\passthrough{\lstinline!other!} & \passthrough{\lstinline!object!} &
必填 & 用于比较的另一个 Python 对象。类型不兼容时比较结果通常为
\passthrough{\lstinline!False!}。 \\
\bottomrule
\end{longtable}

\textbf{返回值}

\begin{longtable}[]{@{}
  >{\raggedright\arraybackslash}p{(\columnwidth - 4\tabcolsep) * \real{0.18}}
  >{\raggedright\arraybackslash}p{(\columnwidth - 4\tabcolsep) * \real{0.20}}
  >{\raggedright\arraybackslash}p{(\columnwidth - 4\tabcolsep) * \real{0.62}}@{}}
\toprule
位置 & 类型 & 含义 \\
\midrule
\endhead
返回值 & \passthrough{\lstinline!bool!} & 返回
\passthrough{\lstinline!bool!}。更精确的业务含义见该方法用途、C++
签名和目标型号协议。 \\
\bottomrule
\end{longtable}

\textbf{对应 C++}

\begin{itemize}
\tightlist
\item
  类：\passthrough{\lstinline!nav\_msgs::msg::dds\_::OccupancyGrid\_!}
\item
  签名：\passthrough{\lstinline"operator!=(const value \&) const"}
\item
  绑定策略：\passthrough{\lstinline!未单独标注!}
\item
  声明位置：由 pybind11 手工绑定或生成注册表提供；manifest
  未记录独立头文件行号。
\end{itemize}

\textbf{说明}

\begin{itemize}
\tightlist
\item
  示例是局部调用片段；\passthrough{\lstinline!obj!}
  和参数变量表示已经按上层业务校验并准备好的对象。
\end{itemize}

\textbf{用法}

\begin{lstlisting}[language=Python]
different = left != right
\end{lstlisting}

\hypertarget{unitree-sdk2-cpp-idl-ros2-occupancygrid-header}{%
\paragraph{\texorpdfstring{\texttt{unitree\_sdk2\_cpp.idl.ros2.OccupancyGrid.header}}{unitree\_sdk2\_cpp.idl.ros2.OccupancyGrid.header}}\label{unitree-sdk2-cpp-idl-ros2-occupancygrid-header}}

底层 \passthrough{\lstinline!nav\_msgs::msg::dds\_::OccupancyGrid\_!} 的
\passthrough{\lstinline!header!} 字段。类型签名只说明 Python
表示；单位、数值范围、数组固定长度和枚举含义需查对应 IDL 头文件。

\textbf{签名}

\begin{lstlisting}[language=Python]
@property
def header(self) -> Any

@header.setter
def header(self, value: Any) -> None
\end{lstlisting}

\textbf{可用性与安全性}

\textbf{\passthrough{\lstinline!AVAILABLE!} /
\passthrough{\lstinline!VALUE\_TYPE!}}：当前绑定源码已实现该消息值操作。

\textbf{参数}

\begin{longtable}[]{@{}
  >{\raggedright\arraybackslash}p{(\columnwidth - 4\tabcolsep) * \real{0.18}}
  >{\raggedright\arraybackslash}p{(\columnwidth - 4\tabcolsep) * \real{0.20}}
  >{\raggedright\arraybackslash}p{(\columnwidth - 4\tabcolsep) * \real{0.62}}@{}}
\toprule
名称 & Python 类型 & 含义 \\
\midrule
\endhead
\passthrough{\lstinline!value!} & \passthrough{\lstinline!Any!} &
要写入或传递的新值，必须符合该参数的 Python 类型以及底层 C++
范围约束。 \\
\bottomrule
\end{longtable}

\textbf{返回值}

\begin{longtable}[]{@{}
  >{\raggedright\arraybackslash}p{(\columnwidth - 4\tabcolsep) * \real{0.18}}
  >{\raggedright\arraybackslash}p{(\columnwidth - 4\tabcolsep) * \real{0.20}}
  >{\raggedright\arraybackslash}p{(\columnwidth - 4\tabcolsep) * \real{0.62}}@{}}
\toprule
位置 & 类型 & 含义 \\
\midrule
\endhead
getter 返回值 & \passthrough{\lstinline!Any!} & 返回属性当前值。 \\
\bottomrule
\end{longtable}

\textbf{对应 C++}

\begin{itemize}
\tightlist
\item
  类：\passthrough{\lstinline!nav\_msgs::msg::dds\_::OccupancyGrid\_!}
\item
  字段类型：\passthrough{\lstinline!::std\_msgs::msg::dds\_::Header\_!}
\item
  声明位置：\passthrough{\lstinline!include/unitree/idl/ros2/OccupancyGrid\_.hpp:28!}
\end{itemize}

\textbf{用法}

\begin{lstlisting}[language=Python]
current_value = obj.header
obj.header = new_value
\end{lstlisting}

\hypertarget{unitree-sdk2-cpp-idl-ros2-occupancygrid-info}{%
\paragraph{\texorpdfstring{\texttt{unitree\_sdk2\_cpp.idl.ros2.OccupancyGrid.info}}{unitree\_sdk2\_cpp.idl.ros2.OccupancyGrid.info}}\label{unitree-sdk2-cpp-idl-ros2-occupancygrid-info}}

底层 \passthrough{\lstinline!nav\_msgs::msg::dds\_::OccupancyGrid\_!} 的
\passthrough{\lstinline!info!} 字段。类型签名只说明 Python
表示；单位、数值范围、数组固定长度和枚举含义需查对应 IDL 头文件。

\textbf{签名}

\begin{lstlisting}[language=Python]
@property
def info(self) -> MapMetaData

@info.setter
def info(self, value: MapMetaData) -> None
\end{lstlisting}

\textbf{可用性与安全性}

\textbf{\passthrough{\lstinline!AVAILABLE!} /
\passthrough{\lstinline!VALUE\_TYPE!}}：当前绑定源码已实现该消息值操作。

\textbf{参数}

\begin{longtable}[]{@{}
  >{\raggedright\arraybackslash}p{(\columnwidth - 4\tabcolsep) * \real{0.18}}
  >{\raggedright\arraybackslash}p{(\columnwidth - 4\tabcolsep) * \real{0.20}}
  >{\raggedright\arraybackslash}p{(\columnwidth - 4\tabcolsep) * \real{0.62}}@{}}
\toprule
名称 & Python 类型 & 含义 \\
\midrule
\endhead
\passthrough{\lstinline!value!} & \passthrough{\lstinline!MapMetaData!}
& 要写入或传递的新值，必须符合该参数的 Python 类型以及底层 C++
范围约束。 \\
\bottomrule
\end{longtable}

\textbf{返回值}

\begin{longtable}[]{@{}
  >{\raggedright\arraybackslash}p{(\columnwidth - 4\tabcolsep) * \real{0.18}}
  >{\raggedright\arraybackslash}p{(\columnwidth - 4\tabcolsep) * \real{0.20}}
  >{\raggedright\arraybackslash}p{(\columnwidth - 4\tabcolsep) * \real{0.62}}@{}}
\toprule
位置 & 类型 & 含义 \\
\midrule
\endhead
getter 返回值 & \passthrough{\lstinline!MapMetaData!} &
返回属性当前值。数组、vector 和嵌套消息采用复制语义。 \\
\bottomrule
\end{longtable}

\textbf{对应 C++}

\begin{itemize}
\tightlist
\item
  类：\passthrough{\lstinline!nav\_msgs::msg::dds\_::OccupancyGrid\_!}
\item
  字段类型：\passthrough{\lstinline!::nav\_msgs::msg::dds\_::MapMetaData\_!}
\item
  声明位置：\passthrough{\lstinline!include/unitree/idl/ros2/OccupancyGrid\_.hpp:29!}
\end{itemize}

\textbf{用法}

\begin{lstlisting}[language=Python]
current_value = obj.info
obj.info = new_value
\end{lstlisting}

\hypertarget{unitree-sdk2-cpp-idl-ros2-occupancygrid-data}{%
\paragraph{\texorpdfstring{\texttt{unitree\_sdk2\_cpp.idl.ros2.OccupancyGrid.data}}{unitree\_sdk2\_cpp.idl.ros2.OccupancyGrid.data}}\label{unitree-sdk2-cpp-idl-ros2-occupancygrid-data}}

底层 \passthrough{\lstinline!nav\_msgs::msg::dds\_::OccupancyGrid\_!} 的
\passthrough{\lstinline!data!} 字段。类型签名只说明 Python
表示；单位、数值范围、数组固定长度和枚举含义需查对应 IDL 头文件。
底层是可变长度 vector；元素约束：取值范围为 0 到 255。

\textbf{签名}

\begin{lstlisting}[language=Python]
@property
def data(self) -> list[int]

@data.setter
def data(self, value: Sequence[int]) -> None
\end{lstlisting}

\textbf{可用性与安全性}

\textbf{\passthrough{\lstinline!AVAILABLE!} /
\passthrough{\lstinline!VALUE\_TYPE!}}：当前绑定源码已实现该消息值操作。

\textbf{参数}

\begin{longtable}[]{@{}
  >{\raggedright\arraybackslash}p{(\columnwidth - 4\tabcolsep) * \real{0.18}}
  >{\raggedright\arraybackslash}p{(\columnwidth - 4\tabcolsep) * \real{0.20}}
  >{\raggedright\arraybackslash}p{(\columnwidth - 4\tabcolsep) * \real{0.62}}@{}}
\toprule
名称 & Python 类型 & 含义 \\
\midrule
\endhead
\passthrough{\lstinline!value!} &
\passthrough{\lstinline!Sequence[int]!} &
要写入或传递的新值，必须符合该参数的 Python 类型以及底层 C++
范围约束。 \\
\bottomrule
\end{longtable}

\textbf{返回值}

\begin{longtable}[]{@{}
  >{\raggedright\arraybackslash}p{(\columnwidth - 4\tabcolsep) * \real{0.18}}
  >{\raggedright\arraybackslash}p{(\columnwidth - 4\tabcolsep) * \real{0.20}}
  >{\raggedright\arraybackslash}p{(\columnwidth - 4\tabcolsep) * \real{0.62}}@{}}
\toprule
位置 & 类型 & 含义 \\
\midrule
\endhead
getter 返回值 & \passthrough{\lstinline!list[int]!} &
返回属性当前值。数组、vector 和嵌套消息采用复制语义。 \\
\bottomrule
\end{longtable}

\textbf{对应 C++}

\begin{itemize}
\tightlist
\item
  类：\passthrough{\lstinline!nav\_msgs::msg::dds\_::OccupancyGrid\_!}
\item
  字段类型：\passthrough{\lstinline!std::vector<uint8\_t>!}
\item
  声明位置：\passthrough{\lstinline!include/unitree/idl/ros2/OccupancyGrid\_.hpp:30!}
\end{itemize}

\textbf{用法}

\begin{lstlisting}[language=Python]
current_value = obj.data
obj.data = new_value
\end{lstlisting}

\begin{quote}
{[}!TIP{]} 读取容器或嵌套值后应遵循``读取、修改、重新赋值''。只修改
getter 返回的副本不会可靠地更新原 C++ 消息。
\end{quote}

\hypertarget{unitree-sdk2-cpp-idl-ros2-posewithcovariance}{%
\subsubsection{\texorpdfstring{\texttt{unitree\_sdk2\_cpp.idl.ros2.PoseWithCovariance}}{unitree\_sdk2\_cpp.idl.ros2.PoseWithCovariance}}\label{unitree-sdk2-cpp-idl-ros2-posewithcovariance}}

可由 Python 读写的 DDS/IDL 消息值类型。字段采用复制语义。

\textbf{导入}

\begin{lstlisting}[language=Python]
from unitree_sdk2_cpp.idl.ros2 import PoseWithCovariance
\end{lstlisting}

\textbf{方法索引}

\begin{longtable}[]{@{}
  >{\raggedright\arraybackslash}p{(\columnwidth - 6\tabcolsep) * \real{0.18}}
  >{\raggedright\arraybackslash}p{(\columnwidth - 6\tabcolsep) * \real{0.24}}
  >{\raggedright\arraybackslash}p{(\columnwidth - 6\tabcolsep) * \real{0.18}}
  >{\raggedright\arraybackslash}p{(\columnwidth - 6\tabcolsep) * \real{0.40}}@{}}
\toprule
方法 & Python 签名 & 状态 & 安全分类 \\
\midrule
\endhead
\protect\hyperlink{unitree-sdk2-cpp-idl-ros2-posewithcovariance-dunder-init-1}{\passthrough{\lstinline!\_\_init\_\_!}}
& \passthrough{\lstinline!def \_\_init\_\_(self) -> None!} &
\passthrough{\lstinline!AVAILABLE!} &
\passthrough{\lstinline!VALUE\_TYPE!} \\
\protect\hyperlink{unitree-sdk2-cpp-idl-ros2-posewithcovariance-dunder-eq-1}{\passthrough{\lstinline!\_\_eq\_\_!}}
& \passthrough{\lstinline!def \_\_eq\_\_(self, other: object) -> bool!}
& \passthrough{\lstinline!AVAILABLE!} &
\passthrough{\lstinline!VALUE\_TYPE!} \\
\protect\hyperlink{unitree-sdk2-cpp-idl-ros2-posewithcovariance-dunder-ne-1}{\passthrough{\lstinline!\_\_ne\_\_!}}
& \passthrough{\lstinline!def \_\_ne\_\_(self, other: object) -> bool!}
& \passthrough{\lstinline!AVAILABLE!} &
\passthrough{\lstinline!VALUE\_TYPE!} \\
\bottomrule
\end{longtable}

\textbf{属性索引}

\begin{longtable}[]{@{}
  >{\raggedright\arraybackslash}p{(\columnwidth - 6\tabcolsep) * \real{0.18}}
  >{\raggedright\arraybackslash}p{(\columnwidth - 6\tabcolsep) * \real{0.24}}
  >{\raggedright\arraybackslash}p{(\columnwidth - 6\tabcolsep) * \real{0.18}}
  >{\raggedright\arraybackslash}p{(\columnwidth - 6\tabcolsep) * \real{0.40}}@{}}
\toprule
属性 & 读取类型 & 写入类型 & 状态 \\
\midrule
\endhead
\protect\hyperlink{unitree-sdk2-cpp-idl-ros2-posewithcovariance-pose}{\passthrough{\lstinline!pose!}}
& \passthrough{\lstinline!Pose!} & \passthrough{\lstinline!Pose!} &
\passthrough{\lstinline!AVAILABLE!} \\
\protect\hyperlink{unitree-sdk2-cpp-idl-ros2-posewithcovariance-covariance}{\passthrough{\lstinline!covariance!}}
& \passthrough{\lstinline!list[float]!} &
\passthrough{\lstinline!Sequence[float]!} &
\passthrough{\lstinline!AVAILABLE!} \\
\bottomrule
\end{longtable}

\hypertarget{unitree-sdk2-cpp-idl-ros2-posewithcovariance-dunder-init-1}{%
\paragraph{\texorpdfstring{\texttt{unitree\_sdk2\_cpp.idl.ros2.PoseWithCovariance.\_\_init\_\_}}{unitree\_sdk2\_cpp.idl.ros2.PoseWithCovariance.\_\_init\_\_}}\label{unitree-sdk2-cpp-idl-ros2-posewithcovariance-dunder-init-1}}

初始化 \passthrough{\lstinline!PoseWithCovariance!}
实例。是否能实际构造取决于下方可用性状态。

\textbf{签名}

\begin{lstlisting}[language=Python]
def __init__(self) -> None
\end{lstlisting}

\textbf{可用性与安全性}

\textbf{\passthrough{\lstinline!AVAILABLE!} /
\passthrough{\lstinline!VALUE\_TYPE!}}：当前绑定源码已实现该消息值操作。

\textbf{参数}

无显式参数。实例方法中的 \passthrough{\lstinline!self!} 由 Python
自动传入。

\textbf{返回值}

\begin{longtable}[]{@{}
  >{\raggedright\arraybackslash}p{(\columnwidth - 4\tabcolsep) * \real{0.18}}
  >{\raggedright\arraybackslash}p{(\columnwidth - 4\tabcolsep) * \real{0.20}}
  >{\raggedright\arraybackslash}p{(\columnwidth - 4\tabcolsep) * \real{0.62}}@{}}
\toprule
位置 & 类型 & 含义 \\
\midrule
\endhead
无 & \passthrough{\lstinline!None!} &
\passthrough{\lstinline!\_\_init\_\_!}
本身不返回值；调用类对象时，在构造函数可用的前提下得到该类实例。 \\
\bottomrule
\end{longtable}

\textbf{对应 C++}

\begin{itemize}
\tightlist
\item
  类：\passthrough{\lstinline!geometry\_msgs::msg::dds\_::PoseWithCovariance\_!}
\item
  签名：\passthrough{\lstinline!PoseWithCovariance\_()!}
\item
  绑定策略：\passthrough{\lstinline!未单独标注!}
\item
  声明位置：\passthrough{\lstinline!include/unitree/idl/ros2/PoseWithCovariance\_.hpp:29!}
\end{itemize}

\textbf{说明}

\begin{itemize}
\tightlist
\item
  示例是局部调用片段；\passthrough{\lstinline!obj!}
  和参数变量表示已经按上层业务校验并准备好的对象。
\end{itemize}

\textbf{用法}

\begin{lstlisting}[language=Python]
value = PoseWithCovariance()
\end{lstlisting}

\hypertarget{unitree-sdk2-cpp-idl-ros2-posewithcovariance-dunder-eq-1}{%
\paragraph{\texorpdfstring{\texttt{unitree\_sdk2\_cpp.idl.ros2.PoseWithCovariance.\_\_eq\_\_}}{unitree\_sdk2\_cpp.idl.ros2.PoseWithCovariance.\_\_eq\_\_}}\label{unitree-sdk2-cpp-idl-ros2-posewithcovariance-dunder-eq-1}}

按底层 IDL 消息内容比较两个对象是否相等。

\textbf{签名}

\begin{lstlisting}[language=Python]
def __eq__(self, other: object) -> bool
\end{lstlisting}

\textbf{可用性与安全性}

\textbf{\passthrough{\lstinline!AVAILABLE!} /
\passthrough{\lstinline!VALUE\_TYPE!}}：当前绑定源码已实现该消息值操作。

\textbf{参数}

\begin{longtable}[]{@{}
  >{\raggedright\arraybackslash}p{(\columnwidth - 6\tabcolsep) * \real{0.18}}
  >{\raggedright\arraybackslash}p{(\columnwidth - 6\tabcolsep) * \real{0.24}}
  >{\raggedright\arraybackslash}p{(\columnwidth - 6\tabcolsep) * \real{0.18}}
  >{\raggedright\arraybackslash}p{(\columnwidth - 6\tabcolsep) * \real{0.40}}@{}}
\toprule
名称 & Python 类型 & 默认值 & 含义 \\
\midrule
\endhead
\passthrough{\lstinline!other!} & \passthrough{\lstinline!object!} &
必填 & 用于比较的另一个 Python 对象。类型不兼容时比较结果通常为
\passthrough{\lstinline!False!}。 \\
\bottomrule
\end{longtable}

\textbf{返回值}

\begin{longtable}[]{@{}
  >{\raggedright\arraybackslash}p{(\columnwidth - 4\tabcolsep) * \real{0.18}}
  >{\raggedright\arraybackslash}p{(\columnwidth - 4\tabcolsep) * \real{0.20}}
  >{\raggedright\arraybackslash}p{(\columnwidth - 4\tabcolsep) * \real{0.62}}@{}}
\toprule
位置 & 类型 & 含义 \\
\midrule
\endhead
返回值 & \passthrough{\lstinline!bool!} & 返回
\passthrough{\lstinline!bool!}。更精确的业务含义见该方法用途、C++
签名和目标型号协议。 \\
\bottomrule
\end{longtable}

\textbf{对应 C++}

\begin{itemize}
\tightlist
\item
  类：\passthrough{\lstinline!geometry\_msgs::msg::dds\_::PoseWithCovariance\_!}
\item
  签名：\passthrough{\lstinline!operator==(const value \&) const!}
\item
  绑定策略：\passthrough{\lstinline!未单独标注!}
\item
  声明位置：由 pybind11 手工绑定或生成注册表提供；manifest
  未记录独立头文件行号。
\end{itemize}

\textbf{说明}

\begin{itemize}
\tightlist
\item
  示例是局部调用片段；\passthrough{\lstinline!obj!}
  和参数变量表示已经按上层业务校验并准备好的对象。
\end{itemize}

\textbf{用法}

\begin{lstlisting}[language=Python]
same = left == right
\end{lstlisting}

\hypertarget{unitree-sdk2-cpp-idl-ros2-posewithcovariance-dunder-ne-1}{%
\paragraph{\texorpdfstring{\texttt{unitree\_sdk2\_cpp.idl.ros2.PoseWithCovariance.\_\_ne\_\_}}{unitree\_sdk2\_cpp.idl.ros2.PoseWithCovariance.\_\_ne\_\_}}\label{unitree-sdk2-cpp-idl-ros2-posewithcovariance-dunder-ne-1}}

按底层 IDL 消息内容比较两个对象是否不相等。

\textbf{签名}

\begin{lstlisting}[language=Python]
def __ne__(self, other: object) -> bool
\end{lstlisting}

\textbf{可用性与安全性}

\textbf{\passthrough{\lstinline!AVAILABLE!} /
\passthrough{\lstinline!VALUE\_TYPE!}}：当前绑定源码已实现该消息值操作。

\textbf{参数}

\begin{longtable}[]{@{}
  >{\raggedright\arraybackslash}p{(\columnwidth - 6\tabcolsep) * \real{0.18}}
  >{\raggedright\arraybackslash}p{(\columnwidth - 6\tabcolsep) * \real{0.24}}
  >{\raggedright\arraybackslash}p{(\columnwidth - 6\tabcolsep) * \real{0.18}}
  >{\raggedright\arraybackslash}p{(\columnwidth - 6\tabcolsep) * \real{0.40}}@{}}
\toprule
名称 & Python 类型 & 默认值 & 含义 \\
\midrule
\endhead
\passthrough{\lstinline!other!} & \passthrough{\lstinline!object!} &
必填 & 用于比较的另一个 Python 对象。类型不兼容时比较结果通常为
\passthrough{\lstinline!False!}。 \\
\bottomrule
\end{longtable}

\textbf{返回值}

\begin{longtable}[]{@{}
  >{\raggedright\arraybackslash}p{(\columnwidth - 4\tabcolsep) * \real{0.18}}
  >{\raggedright\arraybackslash}p{(\columnwidth - 4\tabcolsep) * \real{0.20}}
  >{\raggedright\arraybackslash}p{(\columnwidth - 4\tabcolsep) * \real{0.62}}@{}}
\toprule
位置 & 类型 & 含义 \\
\midrule
\endhead
返回值 & \passthrough{\lstinline!bool!} & 返回
\passthrough{\lstinline!bool!}。更精确的业务含义见该方法用途、C++
签名和目标型号协议。 \\
\bottomrule
\end{longtable}

\textbf{对应 C++}

\begin{itemize}
\tightlist
\item
  类：\passthrough{\lstinline!geometry\_msgs::msg::dds\_::PoseWithCovariance\_!}
\item
  签名：\passthrough{\lstinline"operator!=(const value \&) const"}
\item
  绑定策略：\passthrough{\lstinline!未单独标注!}
\item
  声明位置：由 pybind11 手工绑定或生成注册表提供；manifest
  未记录独立头文件行号。
\end{itemize}

\textbf{说明}

\begin{itemize}
\tightlist
\item
  示例是局部调用片段；\passthrough{\lstinline!obj!}
  和参数变量表示已经按上层业务校验并准备好的对象。
\end{itemize}

\textbf{用法}

\begin{lstlisting}[language=Python]
different = left != right
\end{lstlisting}

\hypertarget{unitree-sdk2-cpp-idl-ros2-posewithcovariance-pose}{%
\paragraph{\texorpdfstring{\texttt{unitree\_sdk2\_cpp.idl.ros2.PoseWithCovariance.pose}}{unitree\_sdk2\_cpp.idl.ros2.PoseWithCovariance.pose}}\label{unitree-sdk2-cpp-idl-ros2-posewithcovariance-pose}}

底层
\passthrough{\lstinline!geometry\_msgs::msg::dds\_::PoseWithCovariance\_!}
的 \passthrough{\lstinline!pose!} 字段。类型签名只说明 Python
表示；单位、数值范围、数组固定长度和枚举含义需查对应 IDL 头文件。

\textbf{签名}

\begin{lstlisting}[language=Python]
@property
def pose(self) -> Pose

@pose.setter
def pose(self, value: Pose) -> None
\end{lstlisting}

\textbf{可用性与安全性}

\textbf{\passthrough{\lstinline!AVAILABLE!} /
\passthrough{\lstinline!VALUE\_TYPE!}}：当前绑定源码已实现该消息值操作。

\textbf{参数}

\begin{longtable}[]{@{}
  >{\raggedright\arraybackslash}p{(\columnwidth - 4\tabcolsep) * \real{0.18}}
  >{\raggedright\arraybackslash}p{(\columnwidth - 4\tabcolsep) * \real{0.20}}
  >{\raggedright\arraybackslash}p{(\columnwidth - 4\tabcolsep) * \real{0.62}}@{}}
\toprule
名称 & Python 类型 & 含义 \\
\midrule
\endhead
\passthrough{\lstinline!value!} & \passthrough{\lstinline!Pose!} &
要写入或传递的新值，必须符合该参数的 Python 类型以及底层 C++
范围约束。 \\
\bottomrule
\end{longtable}

\textbf{返回值}

\begin{longtable}[]{@{}
  >{\raggedright\arraybackslash}p{(\columnwidth - 4\tabcolsep) * \real{0.18}}
  >{\raggedright\arraybackslash}p{(\columnwidth - 4\tabcolsep) * \real{0.20}}
  >{\raggedright\arraybackslash}p{(\columnwidth - 4\tabcolsep) * \real{0.62}}@{}}
\toprule
位置 & 类型 & 含义 \\
\midrule
\endhead
getter 返回值 & \passthrough{\lstinline!Pose!} &
返回属性当前值。数组、vector 和嵌套消息采用复制语义。 \\
\bottomrule
\end{longtable}

\textbf{对应 C++}

\begin{itemize}
\tightlist
\item
  类：\passthrough{\lstinline!geometry\_msgs::msg::dds\_::PoseWithCovariance\_!}
\item
  字段类型：\passthrough{\lstinline!::geometry\_msgs::msg::dds\_::Pose\_!}
\item
  声明位置：\passthrough{\lstinline!include/unitree/idl/ros2/PoseWithCovariance\_.hpp:25!}
\end{itemize}

\textbf{用法}

\begin{lstlisting}[language=Python]
current_value = obj.pose
obj.pose = new_value
\end{lstlisting}

\hypertarget{unitree-sdk2-cpp-idl-ros2-posewithcovariance-covariance}{%
\paragraph{\texorpdfstring{\texttt{unitree\_sdk2\_cpp.idl.ros2.PoseWithCovariance.covariance}}{unitree\_sdk2\_cpp.idl.ros2.PoseWithCovariance.covariance}}\label{unitree-sdk2-cpp-idl-ros2-posewithcovariance-covariance}}

底层
\passthrough{\lstinline!geometry\_msgs::msg::dds\_::PoseWithCovariance\_!}
的 \passthrough{\lstinline!covariance!} 字段。类型签名只说明 Python
表示；单位、数值范围、数组固定长度和枚举含义需查对应 IDL 头文件。
底层是固定长度数组，写入序列必须正好包含 36
个元素；元素约束：底层为浮点数；类型本身不说明单位、坐标系或业务安全范围。

\textbf{签名}

\begin{lstlisting}[language=Python]
@property
def covariance(self) -> list[float]

@covariance.setter
def covariance(self, value: Sequence[float]) -> None
\end{lstlisting}

\textbf{可用性与安全性}

\textbf{\passthrough{\lstinline!AVAILABLE!} /
\passthrough{\lstinline!VALUE\_TYPE!}}：当前绑定源码已实现该消息值操作。

\textbf{参数}

\begin{longtable}[]{@{}
  >{\raggedright\arraybackslash}p{(\columnwidth - 4\tabcolsep) * \real{0.18}}
  >{\raggedright\arraybackslash}p{(\columnwidth - 4\tabcolsep) * \real{0.20}}
  >{\raggedright\arraybackslash}p{(\columnwidth - 4\tabcolsep) * \real{0.62}}@{}}
\toprule
名称 & Python 类型 & 含义 \\
\midrule
\endhead
\passthrough{\lstinline!value!} &
\passthrough{\lstinline!Sequence[float]!} &
要写入或传递的新值，必须符合该参数的 Python 类型以及底层 C++
范围约束。 \\
\bottomrule
\end{longtable}

\textbf{返回值}

\begin{longtable}[]{@{}
  >{\raggedright\arraybackslash}p{(\columnwidth - 4\tabcolsep) * \real{0.18}}
  >{\raggedright\arraybackslash}p{(\columnwidth - 4\tabcolsep) * \real{0.20}}
  >{\raggedright\arraybackslash}p{(\columnwidth - 4\tabcolsep) * \real{0.62}}@{}}
\toprule
位置 & 类型 & 含义 \\
\midrule
\endhead
getter 返回值 & \passthrough{\lstinline!list[float]!} &
返回属性当前值。数组、vector 和嵌套消息采用复制语义。 \\
\bottomrule
\end{longtable}

\textbf{对应 C++}

\begin{itemize}
\tightlist
\item
  类：\passthrough{\lstinline!geometry\_msgs::msg::dds\_::PoseWithCovariance\_!}
\item
  字段类型：\passthrough{\lstinline!std::array<double, 36>!}
\item
  声明位置：\passthrough{\lstinline!include/unitree/idl/ros2/PoseWithCovariance\_.hpp:26!}
\end{itemize}

\textbf{用法}

\begin{lstlisting}[language=Python]
current_value = obj.covariance
obj.covariance = new_value
\end{lstlisting}

\begin{quote}
{[}!TIP{]} 读取容器或嵌套值后应遵循``读取、修改、重新赋值''。只修改
getter 返回的副本不会可靠地更新原 C++ 消息。
\end{quote}

\hypertarget{unitree-sdk2-cpp-idl-ros2-twist}{%
\subsubsection{\texorpdfstring{\texttt{unitree\_sdk2\_cpp.idl.ros2.Twist}}{unitree\_sdk2\_cpp.idl.ros2.Twist}}\label{unitree-sdk2-cpp-idl-ros2-twist}}

可由 Python 读写的 DDS/IDL 消息值类型。字段采用复制语义。

\textbf{导入}

\begin{lstlisting}[language=Python]
from unitree_sdk2_cpp.idl.ros2 import Twist
\end{lstlisting}

\textbf{方法索引}

\begin{longtable}[]{@{}
  >{\raggedright\arraybackslash}p{(\columnwidth - 6\tabcolsep) * \real{0.18}}
  >{\raggedright\arraybackslash}p{(\columnwidth - 6\tabcolsep) * \real{0.24}}
  >{\raggedright\arraybackslash}p{(\columnwidth - 6\tabcolsep) * \real{0.18}}
  >{\raggedright\arraybackslash}p{(\columnwidth - 6\tabcolsep) * \real{0.40}}@{}}
\toprule
方法 & Python 签名 & 状态 & 安全分类 \\
\midrule
\endhead
\protect\hyperlink{unitree-sdk2-cpp-idl-ros2-twist-dunder-init-1}{\passthrough{\lstinline!\_\_init\_\_!}}
& \passthrough{\lstinline!def \_\_init\_\_(self) -> None!} &
\passthrough{\lstinline!AVAILABLE!} &
\passthrough{\lstinline!VALUE\_TYPE!} \\
\protect\hyperlink{unitree-sdk2-cpp-idl-ros2-twist-dunder-eq-1}{\passthrough{\lstinline!\_\_eq\_\_!}}
& \passthrough{\lstinline!def \_\_eq\_\_(self, other: object) -> bool!}
& \passthrough{\lstinline!AVAILABLE!} &
\passthrough{\lstinline!VALUE\_TYPE!} \\
\protect\hyperlink{unitree-sdk2-cpp-idl-ros2-twist-dunder-ne-1}{\passthrough{\lstinline!\_\_ne\_\_!}}
& \passthrough{\lstinline!def \_\_ne\_\_(self, other: object) -> bool!}
& \passthrough{\lstinline!AVAILABLE!} &
\passthrough{\lstinline!VALUE\_TYPE!} \\
\bottomrule
\end{longtable}

\textbf{属性索引}

\begin{longtable}[]{@{}
  >{\raggedright\arraybackslash}p{(\columnwidth - 6\tabcolsep) * \real{0.18}}
  >{\raggedright\arraybackslash}p{(\columnwidth - 6\tabcolsep) * \real{0.24}}
  >{\raggedright\arraybackslash}p{(\columnwidth - 6\tabcolsep) * \real{0.18}}
  >{\raggedright\arraybackslash}p{(\columnwidth - 6\tabcolsep) * \real{0.40}}@{}}
\toprule
属性 & 读取类型 & 写入类型 & 状态 \\
\midrule
\endhead
\protect\hyperlink{unitree-sdk2-cpp-idl-ros2-twist-linear}{\passthrough{\lstinline!linear!}}
& \passthrough{\lstinline!Vector3!} & \passthrough{\lstinline!Vector3!}
& \passthrough{\lstinline!AVAILABLE!} \\
\protect\hyperlink{unitree-sdk2-cpp-idl-ros2-twist-angular}{\passthrough{\lstinline!angular!}}
& \passthrough{\lstinline!Vector3!} & \passthrough{\lstinline!Vector3!}
& \passthrough{\lstinline!AVAILABLE!} \\
\bottomrule
\end{longtable}

\hypertarget{unitree-sdk2-cpp-idl-ros2-twist-dunder-init-1}{%
\paragraph{\texorpdfstring{\texttt{unitree\_sdk2\_cpp.idl.ros2.Twist.\_\_init\_\_}}{unitree\_sdk2\_cpp.idl.ros2.Twist.\_\_init\_\_}}\label{unitree-sdk2-cpp-idl-ros2-twist-dunder-init-1}}

初始化 \passthrough{\lstinline!Twist!}
实例。是否能实际构造取决于下方可用性状态。

\textbf{签名}

\begin{lstlisting}[language=Python]
def __init__(self) -> None
\end{lstlisting}

\textbf{可用性与安全性}

\textbf{\passthrough{\lstinline!AVAILABLE!} /
\passthrough{\lstinline!VALUE\_TYPE!}}：当前绑定源码已实现该消息值操作。

\textbf{参数}

无显式参数。实例方法中的 \passthrough{\lstinline!self!} 由 Python
自动传入。

\textbf{返回值}

\begin{longtable}[]{@{}
  >{\raggedright\arraybackslash}p{(\columnwidth - 4\tabcolsep) * \real{0.18}}
  >{\raggedright\arraybackslash}p{(\columnwidth - 4\tabcolsep) * \real{0.20}}
  >{\raggedright\arraybackslash}p{(\columnwidth - 4\tabcolsep) * \real{0.62}}@{}}
\toprule
位置 & 类型 & 含义 \\
\midrule
\endhead
无 & \passthrough{\lstinline!None!} &
\passthrough{\lstinline!\_\_init\_\_!}
本身不返回值；调用类对象时，在构造函数可用的前提下得到该类实例。 \\
\bottomrule
\end{longtable}

\textbf{对应 C++}

\begin{itemize}
\tightlist
\item
  类：\passthrough{\lstinline!geometry\_msgs::msg::dds\_::Twist\_!}
\item
  签名：\passthrough{\lstinline!Twist\_()!}
\item
  绑定策略：\passthrough{\lstinline!未单独标注!}
\item
  声明位置：\passthrough{\lstinline!include/unitree/idl/ros2/Twist\_.hpp:28!}
\end{itemize}

\textbf{说明}

\begin{itemize}
\tightlist
\item
  示例是局部调用片段；\passthrough{\lstinline!obj!}
  和参数变量表示已经按上层业务校验并准备好的对象。
\end{itemize}

\textbf{用法}

\begin{lstlisting}[language=Python]
value = Twist()
\end{lstlisting}

\hypertarget{unitree-sdk2-cpp-idl-ros2-twist-dunder-eq-1}{%
\paragraph{\texorpdfstring{\texttt{unitree\_sdk2\_cpp.idl.ros2.Twist.\_\_eq\_\_}}{unitree\_sdk2\_cpp.idl.ros2.Twist.\_\_eq\_\_}}\label{unitree-sdk2-cpp-idl-ros2-twist-dunder-eq-1}}

按底层 IDL 消息内容比较两个对象是否相等。

\textbf{签名}

\begin{lstlisting}[language=Python]
def __eq__(self, other: object) -> bool
\end{lstlisting}

\textbf{可用性与安全性}

\textbf{\passthrough{\lstinline!AVAILABLE!} /
\passthrough{\lstinline!VALUE\_TYPE!}}：当前绑定源码已实现该消息值操作。

\textbf{参数}

\begin{longtable}[]{@{}
  >{\raggedright\arraybackslash}p{(\columnwidth - 6\tabcolsep) * \real{0.18}}
  >{\raggedright\arraybackslash}p{(\columnwidth - 6\tabcolsep) * \real{0.24}}
  >{\raggedright\arraybackslash}p{(\columnwidth - 6\tabcolsep) * \real{0.18}}
  >{\raggedright\arraybackslash}p{(\columnwidth - 6\tabcolsep) * \real{0.40}}@{}}
\toprule
名称 & Python 类型 & 默认值 & 含义 \\
\midrule
\endhead
\passthrough{\lstinline!other!} & \passthrough{\lstinline!object!} &
必填 & 用于比较的另一个 Python 对象。类型不兼容时比较结果通常为
\passthrough{\lstinline!False!}。 \\
\bottomrule
\end{longtable}

\textbf{返回值}

\begin{longtable}[]{@{}
  >{\raggedright\arraybackslash}p{(\columnwidth - 4\tabcolsep) * \real{0.18}}
  >{\raggedright\arraybackslash}p{(\columnwidth - 4\tabcolsep) * \real{0.20}}
  >{\raggedright\arraybackslash}p{(\columnwidth - 4\tabcolsep) * \real{0.62}}@{}}
\toprule
位置 & 类型 & 含义 \\
\midrule
\endhead
返回值 & \passthrough{\lstinline!bool!} & 返回
\passthrough{\lstinline!bool!}。更精确的业务含义见该方法用途、C++
签名和目标型号协议。 \\
\bottomrule
\end{longtable}

\textbf{对应 C++}

\begin{itemize}
\tightlist
\item
  类：\passthrough{\lstinline!geometry\_msgs::msg::dds\_::Twist\_!}
\item
  签名：\passthrough{\lstinline!operator==(const value \&) const!}
\item
  绑定策略：\passthrough{\lstinline!未单独标注!}
\item
  声明位置：由 pybind11 手工绑定或生成注册表提供；manifest
  未记录独立头文件行号。
\end{itemize}

\textbf{说明}

\begin{itemize}
\tightlist
\item
  示例是局部调用片段；\passthrough{\lstinline!obj!}
  和参数变量表示已经按上层业务校验并准备好的对象。
\end{itemize}

\textbf{用法}

\begin{lstlisting}[language=Python]
same = left == right
\end{lstlisting}

\hypertarget{unitree-sdk2-cpp-idl-ros2-twist-dunder-ne-1}{%
\paragraph{\texorpdfstring{\texttt{unitree\_sdk2\_cpp.idl.ros2.Twist.\_\_ne\_\_}}{unitree\_sdk2\_cpp.idl.ros2.Twist.\_\_ne\_\_}}\label{unitree-sdk2-cpp-idl-ros2-twist-dunder-ne-1}}

按底层 IDL 消息内容比较两个对象是否不相等。

\textbf{签名}

\begin{lstlisting}[language=Python]
def __ne__(self, other: object) -> bool
\end{lstlisting}

\textbf{可用性与安全性}

\textbf{\passthrough{\lstinline!AVAILABLE!} /
\passthrough{\lstinline!VALUE\_TYPE!}}：当前绑定源码已实现该消息值操作。

\textbf{参数}

\begin{longtable}[]{@{}
  >{\raggedright\arraybackslash}p{(\columnwidth - 6\tabcolsep) * \real{0.18}}
  >{\raggedright\arraybackslash}p{(\columnwidth - 6\tabcolsep) * \real{0.24}}
  >{\raggedright\arraybackslash}p{(\columnwidth - 6\tabcolsep) * \real{0.18}}
  >{\raggedright\arraybackslash}p{(\columnwidth - 6\tabcolsep) * \real{0.40}}@{}}
\toprule
名称 & Python 类型 & 默认值 & 含义 \\
\midrule
\endhead
\passthrough{\lstinline!other!} & \passthrough{\lstinline!object!} &
必填 & 用于比较的另一个 Python 对象。类型不兼容时比较结果通常为
\passthrough{\lstinline!False!}。 \\
\bottomrule
\end{longtable}

\textbf{返回值}

\begin{longtable}[]{@{}
  >{\raggedright\arraybackslash}p{(\columnwidth - 4\tabcolsep) * \real{0.18}}
  >{\raggedright\arraybackslash}p{(\columnwidth - 4\tabcolsep) * \real{0.20}}
  >{\raggedright\arraybackslash}p{(\columnwidth - 4\tabcolsep) * \real{0.62}}@{}}
\toprule
位置 & 类型 & 含义 \\
\midrule
\endhead
返回值 & \passthrough{\lstinline!bool!} & 返回
\passthrough{\lstinline!bool!}。更精确的业务含义见该方法用途、C++
签名和目标型号协议。 \\
\bottomrule
\end{longtable}

\textbf{对应 C++}

\begin{itemize}
\tightlist
\item
  类：\passthrough{\lstinline!geometry\_msgs::msg::dds\_::Twist\_!}
\item
  签名：\passthrough{\lstinline"operator!=(const value \&) const"}
\item
  绑定策略：\passthrough{\lstinline!未单独标注!}
\item
  声明位置：由 pybind11 手工绑定或生成注册表提供；manifest
  未记录独立头文件行号。
\end{itemize}

\textbf{说明}

\begin{itemize}
\tightlist
\item
  示例是局部调用片段；\passthrough{\lstinline!obj!}
  和参数变量表示已经按上层业务校验并准备好的对象。
\end{itemize}

\textbf{用法}

\begin{lstlisting}[language=Python]
different = left != right
\end{lstlisting}

\hypertarget{unitree-sdk2-cpp-idl-ros2-twist-linear}{%
\paragraph{\texorpdfstring{\texttt{unitree\_sdk2\_cpp.idl.ros2.Twist.linear}}{unitree\_sdk2\_cpp.idl.ros2.Twist.linear}}\label{unitree-sdk2-cpp-idl-ros2-twist-linear}}

底层 \passthrough{\lstinline!geometry\_msgs::msg::dds\_::Twist\_!} 的
\passthrough{\lstinline!linear!} 字段。类型签名只说明 Python
表示；单位、数值范围、数组固定长度和枚举含义需查对应 IDL 头文件。

\textbf{签名}

\begin{lstlisting}[language=Python]
@property
def linear(self) -> Vector3

@linear.setter
def linear(self, value: Vector3) -> None
\end{lstlisting}

\textbf{可用性与安全性}

\textbf{\passthrough{\lstinline!AVAILABLE!} /
\passthrough{\lstinline!VALUE\_TYPE!}}：当前绑定源码已实现该消息值操作。

\textbf{参数}

\begin{longtable}[]{@{}
  >{\raggedright\arraybackslash}p{(\columnwidth - 4\tabcolsep) * \real{0.18}}
  >{\raggedright\arraybackslash}p{(\columnwidth - 4\tabcolsep) * \real{0.20}}
  >{\raggedright\arraybackslash}p{(\columnwidth - 4\tabcolsep) * \real{0.62}}@{}}
\toprule
名称 & Python 类型 & 含义 \\
\midrule
\endhead
\passthrough{\lstinline!value!} & \passthrough{\lstinline!Vector3!} &
要写入或传递的新值，必须符合该参数的 Python 类型以及底层 C++
范围约束。 \\
\bottomrule
\end{longtable}

\textbf{返回值}

\begin{longtable}[]{@{}
  >{\raggedright\arraybackslash}p{(\columnwidth - 4\tabcolsep) * \real{0.18}}
  >{\raggedright\arraybackslash}p{(\columnwidth - 4\tabcolsep) * \real{0.20}}
  >{\raggedright\arraybackslash}p{(\columnwidth - 4\tabcolsep) * \real{0.62}}@{}}
\toprule
位置 & 类型 & 含义 \\
\midrule
\endhead
getter 返回值 & \passthrough{\lstinline!Vector3!} &
返回属性当前值。数组、vector 和嵌套消息采用复制语义。 \\
\bottomrule
\end{longtable}

\textbf{对应 C++}

\begin{itemize}
\tightlist
\item
  类：\passthrough{\lstinline!geometry\_msgs::msg::dds\_::Twist\_!}
\item
  字段类型：\passthrough{\lstinline!::geometry\_msgs::msg::dds\_::Vector3\_!}
\item
  声明位置：\passthrough{\lstinline!include/unitree/idl/ros2/Twist\_.hpp:24!}
\end{itemize}

\textbf{用法}

\begin{lstlisting}[language=Python]
current_value = obj.linear
obj.linear = new_value
\end{lstlisting}

\hypertarget{unitree-sdk2-cpp-idl-ros2-twist-angular}{%
\paragraph{\texorpdfstring{\texttt{unitree\_sdk2\_cpp.idl.ros2.Twist.angular}}{unitree\_sdk2\_cpp.idl.ros2.Twist.angular}}\label{unitree-sdk2-cpp-idl-ros2-twist-angular}}

底层 \passthrough{\lstinline!geometry\_msgs::msg::dds\_::Twist\_!} 的
\passthrough{\lstinline!angular!} 字段。类型签名只说明 Python
表示；单位、数值范围、数组固定长度和枚举含义需查对应 IDL 头文件。

\textbf{签名}

\begin{lstlisting}[language=Python]
@property
def angular(self) -> Vector3

@angular.setter
def angular(self, value: Vector3) -> None
\end{lstlisting}

\textbf{可用性与安全性}

\textbf{\passthrough{\lstinline!AVAILABLE!} /
\passthrough{\lstinline!VALUE\_TYPE!}}：当前绑定源码已实现该消息值操作。

\textbf{参数}

\begin{longtable}[]{@{}
  >{\raggedright\arraybackslash}p{(\columnwidth - 4\tabcolsep) * \real{0.18}}
  >{\raggedright\arraybackslash}p{(\columnwidth - 4\tabcolsep) * \real{0.20}}
  >{\raggedright\arraybackslash}p{(\columnwidth - 4\tabcolsep) * \real{0.62}}@{}}
\toprule
名称 & Python 类型 & 含义 \\
\midrule
\endhead
\passthrough{\lstinline!value!} & \passthrough{\lstinline!Vector3!} &
要写入或传递的新值，必须符合该参数的 Python 类型以及底层 C++
范围约束。 \\
\bottomrule
\end{longtable}

\textbf{返回值}

\begin{longtable}[]{@{}
  >{\raggedright\arraybackslash}p{(\columnwidth - 4\tabcolsep) * \real{0.18}}
  >{\raggedright\arraybackslash}p{(\columnwidth - 4\tabcolsep) * \real{0.20}}
  >{\raggedright\arraybackslash}p{(\columnwidth - 4\tabcolsep) * \real{0.62}}@{}}
\toprule
位置 & 类型 & 含义 \\
\midrule
\endhead
getter 返回值 & \passthrough{\lstinline!Vector3!} &
返回属性当前值。数组、vector 和嵌套消息采用复制语义。 \\
\bottomrule
\end{longtable}

\textbf{对应 C++}

\begin{itemize}
\tightlist
\item
  类：\passthrough{\lstinline!geometry\_msgs::msg::dds\_::Twist\_!}
\item
  字段类型：\passthrough{\lstinline!::geometry\_msgs::msg::dds\_::Vector3\_!}
\item
  声明位置：\passthrough{\lstinline!include/unitree/idl/ros2/Twist\_.hpp:25!}
\end{itemize}

\textbf{用法}

\begin{lstlisting}[language=Python]
current_value = obj.angular
obj.angular = new_value
\end{lstlisting}

\hypertarget{unitree-sdk2-cpp-idl-ros2-twistwithcovariance}{%
\subsubsection{\texorpdfstring{\texttt{unitree\_sdk2\_cpp.idl.ros2.TwistWithCovariance}}{unitree\_sdk2\_cpp.idl.ros2.TwistWithCovariance}}\label{unitree-sdk2-cpp-idl-ros2-twistwithcovariance}}

可由 Python 读写的 DDS/IDL 消息值类型。字段采用复制语义。

\textbf{导入}

\begin{lstlisting}[language=Python]
from unitree_sdk2_cpp.idl.ros2 import TwistWithCovariance
\end{lstlisting}

\textbf{方法索引}

\begin{longtable}[]{@{}
  >{\raggedright\arraybackslash}p{(\columnwidth - 6\tabcolsep) * \real{0.18}}
  >{\raggedright\arraybackslash}p{(\columnwidth - 6\tabcolsep) * \real{0.24}}
  >{\raggedright\arraybackslash}p{(\columnwidth - 6\tabcolsep) * \real{0.18}}
  >{\raggedright\arraybackslash}p{(\columnwidth - 6\tabcolsep) * \real{0.40}}@{}}
\toprule
方法 & Python 签名 & 状态 & 安全分类 \\
\midrule
\endhead
\protect\hyperlink{unitree-sdk2-cpp-idl-ros2-twistwithcovariance-dunder-init-1}{\passthrough{\lstinline!\_\_init\_\_!}}
& \passthrough{\lstinline!def \_\_init\_\_(self) -> None!} &
\passthrough{\lstinline!AVAILABLE!} &
\passthrough{\lstinline!VALUE\_TYPE!} \\
\protect\hyperlink{unitree-sdk2-cpp-idl-ros2-twistwithcovariance-dunder-eq-1}{\passthrough{\lstinline!\_\_eq\_\_!}}
& \passthrough{\lstinline!def \_\_eq\_\_(self, other: object) -> bool!}
& \passthrough{\lstinline!AVAILABLE!} &
\passthrough{\lstinline!VALUE\_TYPE!} \\
\protect\hyperlink{unitree-sdk2-cpp-idl-ros2-twistwithcovariance-dunder-ne-1}{\passthrough{\lstinline!\_\_ne\_\_!}}
& \passthrough{\lstinline!def \_\_ne\_\_(self, other: object) -> bool!}
& \passthrough{\lstinline!AVAILABLE!} &
\passthrough{\lstinline!VALUE\_TYPE!} \\
\bottomrule
\end{longtable}

\textbf{属性索引}

\begin{longtable}[]{@{}
  >{\raggedright\arraybackslash}p{(\columnwidth - 6\tabcolsep) * \real{0.18}}
  >{\raggedright\arraybackslash}p{(\columnwidth - 6\tabcolsep) * \real{0.24}}
  >{\raggedright\arraybackslash}p{(\columnwidth - 6\tabcolsep) * \real{0.18}}
  >{\raggedright\arraybackslash}p{(\columnwidth - 6\tabcolsep) * \real{0.40}}@{}}
\toprule
属性 & 读取类型 & 写入类型 & 状态 \\
\midrule
\endhead
\protect\hyperlink{unitree-sdk2-cpp-idl-ros2-twistwithcovariance-twist}{\passthrough{\lstinline!twist!}}
& \passthrough{\lstinline!Twist!} & \passthrough{\lstinline!Twist!} &
\passthrough{\lstinline!AVAILABLE!} \\
\protect\hyperlink{unitree-sdk2-cpp-idl-ros2-twistwithcovariance-covariance}{\passthrough{\lstinline!covariance!}}
& \passthrough{\lstinline!list[float]!} &
\passthrough{\lstinline!Sequence[float]!} &
\passthrough{\lstinline!AVAILABLE!} \\
\bottomrule
\end{longtable}

\hypertarget{unitree-sdk2-cpp-idl-ros2-twistwithcovariance-dunder-init-1}{%
\paragraph{\texorpdfstring{\texttt{unitree\_sdk2\_cpp.idl.ros2.TwistWithCovariance.\_\_init\_\_}}{unitree\_sdk2\_cpp.idl.ros2.TwistWithCovariance.\_\_init\_\_}}\label{unitree-sdk2-cpp-idl-ros2-twistwithcovariance-dunder-init-1}}

初始化 \passthrough{\lstinline!TwistWithCovariance!}
实例。是否能实际构造取决于下方可用性状态。

\textbf{签名}

\begin{lstlisting}[language=Python]
def __init__(self) -> None
\end{lstlisting}

\textbf{可用性与安全性}

\textbf{\passthrough{\lstinline!AVAILABLE!} /
\passthrough{\lstinline!VALUE\_TYPE!}}：当前绑定源码已实现该消息值操作。

\textbf{参数}

无显式参数。实例方法中的 \passthrough{\lstinline!self!} 由 Python
自动传入。

\textbf{返回值}

\begin{longtable}[]{@{}
  >{\raggedright\arraybackslash}p{(\columnwidth - 4\tabcolsep) * \real{0.18}}
  >{\raggedright\arraybackslash}p{(\columnwidth - 4\tabcolsep) * \real{0.20}}
  >{\raggedright\arraybackslash}p{(\columnwidth - 4\tabcolsep) * \real{0.62}}@{}}
\toprule
位置 & 类型 & 含义 \\
\midrule
\endhead
无 & \passthrough{\lstinline!None!} &
\passthrough{\lstinline!\_\_init\_\_!}
本身不返回值；调用类对象时，在构造函数可用的前提下得到该类实例。 \\
\bottomrule
\end{longtable}

\textbf{对应 C++}

\begin{itemize}
\tightlist
\item
  类：\passthrough{\lstinline!geometry\_msgs::msg::dds\_::TwistWithCovariance\_!}
\item
  签名：\passthrough{\lstinline!TwistWithCovariance\_()!}
\item
  绑定策略：\passthrough{\lstinline!未单独标注!}
\item
  声明位置：\passthrough{\lstinline!include/unitree/idl/ros2/TwistWithCovariance\_.hpp:29!}
\end{itemize}

\textbf{说明}

\begin{itemize}
\tightlist
\item
  示例是局部调用片段；\passthrough{\lstinline!obj!}
  和参数变量表示已经按上层业务校验并准备好的对象。
\end{itemize}

\textbf{用法}

\begin{lstlisting}[language=Python]
value = TwistWithCovariance()
\end{lstlisting}

\hypertarget{unitree-sdk2-cpp-idl-ros2-twistwithcovariance-dunder-eq-1}{%
\paragraph{\texorpdfstring{\texttt{unitree\_sdk2\_cpp.idl.ros2.TwistWithCovariance.\_\_eq\_\_}}{unitree\_sdk2\_cpp.idl.ros2.TwistWithCovariance.\_\_eq\_\_}}\label{unitree-sdk2-cpp-idl-ros2-twistwithcovariance-dunder-eq-1}}

按底层 IDL 消息内容比较两个对象是否相等。

\textbf{签名}

\begin{lstlisting}[language=Python]
def __eq__(self, other: object) -> bool
\end{lstlisting}

\textbf{可用性与安全性}

\textbf{\passthrough{\lstinline!AVAILABLE!} /
\passthrough{\lstinline!VALUE\_TYPE!}}：当前绑定源码已实现该消息值操作。

\textbf{参数}

\begin{longtable}[]{@{}
  >{\raggedright\arraybackslash}p{(\columnwidth - 6\tabcolsep) * \real{0.18}}
  >{\raggedright\arraybackslash}p{(\columnwidth - 6\tabcolsep) * \real{0.24}}
  >{\raggedright\arraybackslash}p{(\columnwidth - 6\tabcolsep) * \real{0.18}}
  >{\raggedright\arraybackslash}p{(\columnwidth - 6\tabcolsep) * \real{0.40}}@{}}
\toprule
名称 & Python 类型 & 默认值 & 含义 \\
\midrule
\endhead
\passthrough{\lstinline!other!} & \passthrough{\lstinline!object!} &
必填 & 用于比较的另一个 Python 对象。类型不兼容时比较结果通常为
\passthrough{\lstinline!False!}。 \\
\bottomrule
\end{longtable}

\textbf{返回值}

\begin{longtable}[]{@{}
  >{\raggedright\arraybackslash}p{(\columnwidth - 4\tabcolsep) * \real{0.18}}
  >{\raggedright\arraybackslash}p{(\columnwidth - 4\tabcolsep) * \real{0.20}}
  >{\raggedright\arraybackslash}p{(\columnwidth - 4\tabcolsep) * \real{0.62}}@{}}
\toprule
位置 & 类型 & 含义 \\
\midrule
\endhead
返回值 & \passthrough{\lstinline!bool!} & 返回
\passthrough{\lstinline!bool!}。更精确的业务含义见该方法用途、C++
签名和目标型号协议。 \\
\bottomrule
\end{longtable}

\textbf{对应 C++}

\begin{itemize}
\tightlist
\item
  类：\passthrough{\lstinline!geometry\_msgs::msg::dds\_::TwistWithCovariance\_!}
\item
  签名：\passthrough{\lstinline!operator==(const value \&) const!}
\item
  绑定策略：\passthrough{\lstinline!未单独标注!}
\item
  声明位置：由 pybind11 手工绑定或生成注册表提供；manifest
  未记录独立头文件行号。
\end{itemize}

\textbf{说明}

\begin{itemize}
\tightlist
\item
  示例是局部调用片段；\passthrough{\lstinline!obj!}
  和参数变量表示已经按上层业务校验并准备好的对象。
\end{itemize}

\textbf{用法}

\begin{lstlisting}[language=Python]
same = left == right
\end{lstlisting}

\hypertarget{unitree-sdk2-cpp-idl-ros2-twistwithcovariance-dunder-ne-1}{%
\paragraph{\texorpdfstring{\texttt{unitree\_sdk2\_cpp.idl.ros2.TwistWithCovariance.\_\_ne\_\_}}{unitree\_sdk2\_cpp.idl.ros2.TwistWithCovariance.\_\_ne\_\_}}\label{unitree-sdk2-cpp-idl-ros2-twistwithcovariance-dunder-ne-1}}

按底层 IDL 消息内容比较两个对象是否不相等。

\textbf{签名}

\begin{lstlisting}[language=Python]
def __ne__(self, other: object) -> bool
\end{lstlisting}

\textbf{可用性与安全性}

\textbf{\passthrough{\lstinline!AVAILABLE!} /
\passthrough{\lstinline!VALUE\_TYPE!}}：当前绑定源码已实现该消息值操作。

\textbf{参数}

\begin{longtable}[]{@{}
  >{\raggedright\arraybackslash}p{(\columnwidth - 6\tabcolsep) * \real{0.18}}
  >{\raggedright\arraybackslash}p{(\columnwidth - 6\tabcolsep) * \real{0.24}}
  >{\raggedright\arraybackslash}p{(\columnwidth - 6\tabcolsep) * \real{0.18}}
  >{\raggedright\arraybackslash}p{(\columnwidth - 6\tabcolsep) * \real{0.40}}@{}}
\toprule
名称 & Python 类型 & 默认值 & 含义 \\
\midrule
\endhead
\passthrough{\lstinline!other!} & \passthrough{\lstinline!object!} &
必填 & 用于比较的另一个 Python 对象。类型不兼容时比较结果通常为
\passthrough{\lstinline!False!}。 \\
\bottomrule
\end{longtable}

\textbf{返回值}

\begin{longtable}[]{@{}
  >{\raggedright\arraybackslash}p{(\columnwidth - 4\tabcolsep) * \real{0.18}}
  >{\raggedright\arraybackslash}p{(\columnwidth - 4\tabcolsep) * \real{0.20}}
  >{\raggedright\arraybackslash}p{(\columnwidth - 4\tabcolsep) * \real{0.62}}@{}}
\toprule
位置 & 类型 & 含义 \\
\midrule
\endhead
返回值 & \passthrough{\lstinline!bool!} & 返回
\passthrough{\lstinline!bool!}。更精确的业务含义见该方法用途、C++
签名和目标型号协议。 \\
\bottomrule
\end{longtable}

\textbf{对应 C++}

\begin{itemize}
\tightlist
\item
  类：\passthrough{\lstinline!geometry\_msgs::msg::dds\_::TwistWithCovariance\_!}
\item
  签名：\passthrough{\lstinline"operator!=(const value \&) const"}
\item
  绑定策略：\passthrough{\lstinline!未单独标注!}
\item
  声明位置：由 pybind11 手工绑定或生成注册表提供；manifest
  未记录独立头文件行号。
\end{itemize}

\textbf{说明}

\begin{itemize}
\tightlist
\item
  示例是局部调用片段；\passthrough{\lstinline!obj!}
  和参数变量表示已经按上层业务校验并准备好的对象。
\end{itemize}

\textbf{用法}

\begin{lstlisting}[language=Python]
different = left != right
\end{lstlisting}

\hypertarget{unitree-sdk2-cpp-idl-ros2-twistwithcovariance-twist}{%
\paragraph{\texorpdfstring{\texttt{unitree\_sdk2\_cpp.idl.ros2.TwistWithCovariance.twist}}{unitree\_sdk2\_cpp.idl.ros2.TwistWithCovariance.twist}}\label{unitree-sdk2-cpp-idl-ros2-twistwithcovariance-twist}}

底层
\passthrough{\lstinline!geometry\_msgs::msg::dds\_::TwistWithCovariance\_!}
的 \passthrough{\lstinline!twist!} 字段。类型签名只说明 Python
表示；单位、数值范围、数组固定长度和枚举含义需查对应 IDL 头文件。

\textbf{签名}

\begin{lstlisting}[language=Python]
@property
def twist(self) -> Twist

@twist.setter
def twist(self, value: Twist) -> None
\end{lstlisting}

\textbf{可用性与安全性}

\textbf{\passthrough{\lstinline!AVAILABLE!} /
\passthrough{\lstinline!VALUE\_TYPE!}}：当前绑定源码已实现该消息值操作。

\textbf{参数}

\begin{longtable}[]{@{}
  >{\raggedright\arraybackslash}p{(\columnwidth - 4\tabcolsep) * \real{0.18}}
  >{\raggedright\arraybackslash}p{(\columnwidth - 4\tabcolsep) * \real{0.20}}
  >{\raggedright\arraybackslash}p{(\columnwidth - 4\tabcolsep) * \real{0.62}}@{}}
\toprule
名称 & Python 类型 & 含义 \\
\midrule
\endhead
\passthrough{\lstinline!value!} & \passthrough{\lstinline!Twist!} &
要写入或传递的新值，必须符合该参数的 Python 类型以及底层 C++
范围约束。 \\
\bottomrule
\end{longtable}

\textbf{返回值}

\begin{longtable}[]{@{}
  >{\raggedright\arraybackslash}p{(\columnwidth - 4\tabcolsep) * \real{0.18}}
  >{\raggedright\arraybackslash}p{(\columnwidth - 4\tabcolsep) * \real{0.20}}
  >{\raggedright\arraybackslash}p{(\columnwidth - 4\tabcolsep) * \real{0.62}}@{}}
\toprule
位置 & 类型 & 含义 \\
\midrule
\endhead
getter 返回值 & \passthrough{\lstinline!Twist!} &
返回属性当前值。数组、vector 和嵌套消息采用复制语义。 \\
\bottomrule
\end{longtable}

\textbf{对应 C++}

\begin{itemize}
\tightlist
\item
  类：\passthrough{\lstinline!geometry\_msgs::msg::dds\_::TwistWithCovariance\_!}
\item
  字段类型：\passthrough{\lstinline!::geometry\_msgs::msg::dds\_::Twist\_!}
\item
  声明位置：\passthrough{\lstinline!include/unitree/idl/ros2/TwistWithCovariance\_.hpp:25!}
\end{itemize}

\textbf{用法}

\begin{lstlisting}[language=Python]
current_value = obj.twist
obj.twist = new_value
\end{lstlisting}

\hypertarget{unitree-sdk2-cpp-idl-ros2-twistwithcovariance-covariance}{%
\paragraph{\texorpdfstring{\texttt{unitree\_sdk2\_cpp.idl.ros2.TwistWithCovariance.covariance}}{unitree\_sdk2\_cpp.idl.ros2.TwistWithCovariance.covariance}}\label{unitree-sdk2-cpp-idl-ros2-twistwithcovariance-covariance}}

底层
\passthrough{\lstinline!geometry\_msgs::msg::dds\_::TwistWithCovariance\_!}
的 \passthrough{\lstinline!covariance!} 字段。类型签名只说明 Python
表示；单位、数值范围、数组固定长度和枚举含义需查对应 IDL 头文件。
底层是固定长度数组，写入序列必须正好包含 36
个元素；元素约束：底层为浮点数；类型本身不说明单位、坐标系或业务安全范围。

\textbf{签名}

\begin{lstlisting}[language=Python]
@property
def covariance(self) -> list[float]

@covariance.setter
def covariance(self, value: Sequence[float]) -> None
\end{lstlisting}

\textbf{可用性与安全性}

\textbf{\passthrough{\lstinline!AVAILABLE!} /
\passthrough{\lstinline!VALUE\_TYPE!}}：当前绑定源码已实现该消息值操作。

\textbf{参数}

\begin{longtable}[]{@{}
  >{\raggedright\arraybackslash}p{(\columnwidth - 4\tabcolsep) * \real{0.18}}
  >{\raggedright\arraybackslash}p{(\columnwidth - 4\tabcolsep) * \real{0.20}}
  >{\raggedright\arraybackslash}p{(\columnwidth - 4\tabcolsep) * \real{0.62}}@{}}
\toprule
名称 & Python 类型 & 含义 \\
\midrule
\endhead
\passthrough{\lstinline!value!} &
\passthrough{\lstinline!Sequence[float]!} &
要写入或传递的新值，必须符合该参数的 Python 类型以及底层 C++
范围约束。 \\
\bottomrule
\end{longtable}

\textbf{返回值}

\begin{longtable}[]{@{}
  >{\raggedright\arraybackslash}p{(\columnwidth - 4\tabcolsep) * \real{0.18}}
  >{\raggedright\arraybackslash}p{(\columnwidth - 4\tabcolsep) * \real{0.20}}
  >{\raggedright\arraybackslash}p{(\columnwidth - 4\tabcolsep) * \real{0.62}}@{}}
\toprule
位置 & 类型 & 含义 \\
\midrule
\endhead
getter 返回值 & \passthrough{\lstinline!list[float]!} &
返回属性当前值。数组、vector 和嵌套消息采用复制语义。 \\
\bottomrule
\end{longtable}

\textbf{对应 C++}

\begin{itemize}
\tightlist
\item
  类：\passthrough{\lstinline!geometry\_msgs::msg::dds\_::TwistWithCovariance\_!}
\item
  字段类型：\passthrough{\lstinline!std::array<double, 36>!}
\item
  声明位置：\passthrough{\lstinline!include/unitree/idl/ros2/TwistWithCovariance\_.hpp:26!}
\end{itemize}

\textbf{用法}

\begin{lstlisting}[language=Python]
current_value = obj.covariance
obj.covariance = new_value
\end{lstlisting}

\begin{quote}
{[}!TIP{]} 读取容器或嵌套值后应遵循``读取、修改、重新赋值''。只修改
getter 返回的副本不会可靠地更新原 C++ 消息。
\end{quote}

\hypertarget{unitree-sdk2-cpp-idl-ros2-odometry}{%
\subsubsection{\texorpdfstring{\texttt{unitree\_sdk2\_cpp.idl.ros2.Odometry}}{unitree\_sdk2\_cpp.idl.ros2.Odometry}}\label{unitree-sdk2-cpp-idl-ros2-odometry}}

可由 Python 读写的 DDS/IDL 消息值类型。字段采用复制语义。

\textbf{导入}

\begin{lstlisting}[language=Python]
from unitree_sdk2_cpp.idl.ros2 import Odometry
\end{lstlisting}

\textbf{方法索引}

\begin{longtable}[]{@{}
  >{\raggedright\arraybackslash}p{(\columnwidth - 6\tabcolsep) * \real{0.18}}
  >{\raggedright\arraybackslash}p{(\columnwidth - 6\tabcolsep) * \real{0.24}}
  >{\raggedright\arraybackslash}p{(\columnwidth - 6\tabcolsep) * \real{0.18}}
  >{\raggedright\arraybackslash}p{(\columnwidth - 6\tabcolsep) * \real{0.40}}@{}}
\toprule
方法 & Python 签名 & 状态 & 安全分类 \\
\midrule
\endhead
\protect\hyperlink{unitree-sdk2-cpp-idl-ros2-odometry-dunder-init-1}{\passthrough{\lstinline!\_\_init\_\_!}}
& \passthrough{\lstinline!def \_\_init\_\_(self) -> None!} &
\passthrough{\lstinline!AVAILABLE!} &
\passthrough{\lstinline!VALUE\_TYPE!} \\
\protect\hyperlink{unitree-sdk2-cpp-idl-ros2-odometry-dunder-eq-1}{\passthrough{\lstinline!\_\_eq\_\_!}}
& \passthrough{\lstinline!def \_\_eq\_\_(self, other: object) -> bool!}
& \passthrough{\lstinline!AVAILABLE!} &
\passthrough{\lstinline!VALUE\_TYPE!} \\
\protect\hyperlink{unitree-sdk2-cpp-idl-ros2-odometry-dunder-ne-1}{\passthrough{\lstinline!\_\_ne\_\_!}}
& \passthrough{\lstinline!def \_\_ne\_\_(self, other: object) -> bool!}
& \passthrough{\lstinline!AVAILABLE!} &
\passthrough{\lstinline!VALUE\_TYPE!} \\
\bottomrule
\end{longtable}

\textbf{属性索引}

\begin{longtable}[]{@{}
  >{\raggedright\arraybackslash}p{(\columnwidth - 6\tabcolsep) * \real{0.18}}
  >{\raggedright\arraybackslash}p{(\columnwidth - 6\tabcolsep) * \real{0.24}}
  >{\raggedright\arraybackslash}p{(\columnwidth - 6\tabcolsep) * \real{0.18}}
  >{\raggedright\arraybackslash}p{(\columnwidth - 6\tabcolsep) * \real{0.40}}@{}}
\toprule
属性 & 读取类型 & 写入类型 & 状态 \\
\midrule
\endhead
\protect\hyperlink{unitree-sdk2-cpp-idl-ros2-odometry-header}{\passthrough{\lstinline!header!}}
& \passthrough{\lstinline!Any!} & \passthrough{\lstinline!Any!} &
\passthrough{\lstinline!AVAILABLE!} \\
\protect\hyperlink{unitree-sdk2-cpp-idl-ros2-odometry-child-frame-id}{\passthrough{\lstinline!child\_frame\_id!}}
& \passthrough{\lstinline!str!} & \passthrough{\lstinline!str!} &
\passthrough{\lstinline!AVAILABLE!} \\
\protect\hyperlink{unitree-sdk2-cpp-idl-ros2-odometry-pose}{\passthrough{\lstinline!pose!}}
& \passthrough{\lstinline!Any!} & \passthrough{\lstinline!Any!} &
\passthrough{\lstinline!AVAILABLE!} \\
\protect\hyperlink{unitree-sdk2-cpp-idl-ros2-odometry-twist}{\passthrough{\lstinline!twist!}}
& \passthrough{\lstinline!Any!} & \passthrough{\lstinline!Any!} &
\passthrough{\lstinline!AVAILABLE!} \\
\bottomrule
\end{longtable}

\hypertarget{unitree-sdk2-cpp-idl-ros2-odometry-dunder-init-1}{%
\paragraph{\texorpdfstring{\texttt{unitree\_sdk2\_cpp.idl.ros2.Odometry.\_\_init\_\_}}{unitree\_sdk2\_cpp.idl.ros2.Odometry.\_\_init\_\_}}\label{unitree-sdk2-cpp-idl-ros2-odometry-dunder-init-1}}

初始化 \passthrough{\lstinline!Odometry!}
实例。是否能实际构造取决于下方可用性状态。

\textbf{签名}

\begin{lstlisting}[language=Python]
def __init__(self) -> None
\end{lstlisting}

\textbf{可用性与安全性}

\textbf{\passthrough{\lstinline!AVAILABLE!} /
\passthrough{\lstinline!VALUE\_TYPE!}}：当前绑定源码已实现该消息值操作。

\textbf{参数}

无显式参数。实例方法中的 \passthrough{\lstinline!self!} 由 Python
自动传入。

\textbf{返回值}

\begin{longtable}[]{@{}
  >{\raggedright\arraybackslash}p{(\columnwidth - 4\tabcolsep) * \real{0.18}}
  >{\raggedright\arraybackslash}p{(\columnwidth - 4\tabcolsep) * \real{0.20}}
  >{\raggedright\arraybackslash}p{(\columnwidth - 4\tabcolsep) * \real{0.62}}@{}}
\toprule
位置 & 类型 & 含义 \\
\midrule
\endhead
无 & \passthrough{\lstinline!None!} &
\passthrough{\lstinline!\_\_init\_\_!}
本身不返回值；调用类对象时，在构造函数可用的前提下得到该类实例。 \\
\bottomrule
\end{longtable}

\textbf{对应 C++}

\begin{itemize}
\tightlist
\item
  类：\passthrough{\lstinline!nav\_msgs::msg::dds\_::Odometry\_!}
\item
  签名：\passthrough{\lstinline!Odometry\_()!}
\item
  绑定策略：\passthrough{\lstinline!未单独标注!}
\item
  声明位置：\passthrough{\lstinline!include/unitree/idl/ros2/Odometry\_.hpp:35!}
\end{itemize}

\textbf{说明}

\begin{itemize}
\tightlist
\item
  示例是局部调用片段；\passthrough{\lstinline!obj!}
  和参数变量表示已经按上层业务校验并准备好的对象。
\end{itemize}

\textbf{用法}

\begin{lstlisting}[language=Python]
value = Odometry()
\end{lstlisting}

\hypertarget{unitree-sdk2-cpp-idl-ros2-odometry-dunder-eq-1}{%
\paragraph{\texorpdfstring{\texttt{unitree\_sdk2\_cpp.idl.ros2.Odometry.\_\_eq\_\_}}{unitree\_sdk2\_cpp.idl.ros2.Odometry.\_\_eq\_\_}}\label{unitree-sdk2-cpp-idl-ros2-odometry-dunder-eq-1}}

按底层 IDL 消息内容比较两个对象是否相等。

\textbf{签名}

\begin{lstlisting}[language=Python]
def __eq__(self, other: object) -> bool
\end{lstlisting}

\textbf{可用性与安全性}

\textbf{\passthrough{\lstinline!AVAILABLE!} /
\passthrough{\lstinline!VALUE\_TYPE!}}：当前绑定源码已实现该消息值操作。

\textbf{参数}

\begin{longtable}[]{@{}
  >{\raggedright\arraybackslash}p{(\columnwidth - 6\tabcolsep) * \real{0.18}}
  >{\raggedright\arraybackslash}p{(\columnwidth - 6\tabcolsep) * \real{0.24}}
  >{\raggedright\arraybackslash}p{(\columnwidth - 6\tabcolsep) * \real{0.18}}
  >{\raggedright\arraybackslash}p{(\columnwidth - 6\tabcolsep) * \real{0.40}}@{}}
\toprule
名称 & Python 类型 & 默认值 & 含义 \\
\midrule
\endhead
\passthrough{\lstinline!other!} & \passthrough{\lstinline!object!} &
必填 & 用于比较的另一个 Python 对象。类型不兼容时比较结果通常为
\passthrough{\lstinline!False!}。 \\
\bottomrule
\end{longtable}

\textbf{返回值}

\begin{longtable}[]{@{}
  >{\raggedright\arraybackslash}p{(\columnwidth - 4\tabcolsep) * \real{0.18}}
  >{\raggedright\arraybackslash}p{(\columnwidth - 4\tabcolsep) * \real{0.20}}
  >{\raggedright\arraybackslash}p{(\columnwidth - 4\tabcolsep) * \real{0.62}}@{}}
\toprule
位置 & 类型 & 含义 \\
\midrule
\endhead
返回值 & \passthrough{\lstinline!bool!} & 返回
\passthrough{\lstinline!bool!}。更精确的业务含义见该方法用途、C++
签名和目标型号协议。 \\
\bottomrule
\end{longtable}

\textbf{对应 C++}

\begin{itemize}
\tightlist
\item
  类：\passthrough{\lstinline!nav\_msgs::msg::dds\_::Odometry\_!}
\item
  签名：\passthrough{\lstinline!operator==(const value \&) const!}
\item
  绑定策略：\passthrough{\lstinline!未单独标注!}
\item
  声明位置：由 pybind11 手工绑定或生成注册表提供；manifest
  未记录独立头文件行号。
\end{itemize}

\textbf{说明}

\begin{itemize}
\tightlist
\item
  示例是局部调用片段；\passthrough{\lstinline!obj!}
  和参数变量表示已经按上层业务校验并准备好的对象。
\end{itemize}

\textbf{用法}

\begin{lstlisting}[language=Python]
same = left == right
\end{lstlisting}

\hypertarget{unitree-sdk2-cpp-idl-ros2-odometry-dunder-ne-1}{%
\paragraph{\texorpdfstring{\texttt{unitree\_sdk2\_cpp.idl.ros2.Odometry.\_\_ne\_\_}}{unitree\_sdk2\_cpp.idl.ros2.Odometry.\_\_ne\_\_}}\label{unitree-sdk2-cpp-idl-ros2-odometry-dunder-ne-1}}

按底层 IDL 消息内容比较两个对象是否不相等。

\textbf{签名}

\begin{lstlisting}[language=Python]
def __ne__(self, other: object) -> bool
\end{lstlisting}

\textbf{可用性与安全性}

\textbf{\passthrough{\lstinline!AVAILABLE!} /
\passthrough{\lstinline!VALUE\_TYPE!}}：当前绑定源码已实现该消息值操作。

\textbf{参数}

\begin{longtable}[]{@{}
  >{\raggedright\arraybackslash}p{(\columnwidth - 6\tabcolsep) * \real{0.18}}
  >{\raggedright\arraybackslash}p{(\columnwidth - 6\tabcolsep) * \real{0.24}}
  >{\raggedright\arraybackslash}p{(\columnwidth - 6\tabcolsep) * \real{0.18}}
  >{\raggedright\arraybackslash}p{(\columnwidth - 6\tabcolsep) * \real{0.40}}@{}}
\toprule
名称 & Python 类型 & 默认值 & 含义 \\
\midrule
\endhead
\passthrough{\lstinline!other!} & \passthrough{\lstinline!object!} &
必填 & 用于比较的另一个 Python 对象。类型不兼容时比较结果通常为
\passthrough{\lstinline!False!}。 \\
\bottomrule
\end{longtable}

\textbf{返回值}

\begin{longtable}[]{@{}
  >{\raggedright\arraybackslash}p{(\columnwidth - 4\tabcolsep) * \real{0.18}}
  >{\raggedright\arraybackslash}p{(\columnwidth - 4\tabcolsep) * \real{0.20}}
  >{\raggedright\arraybackslash}p{(\columnwidth - 4\tabcolsep) * \real{0.62}}@{}}
\toprule
位置 & 类型 & 含义 \\
\midrule
\endhead
返回值 & \passthrough{\lstinline!bool!} & 返回
\passthrough{\lstinline!bool!}。更精确的业务含义见该方法用途、C++
签名和目标型号协议。 \\
\bottomrule
\end{longtable}

\textbf{对应 C++}

\begin{itemize}
\tightlist
\item
  类：\passthrough{\lstinline!nav\_msgs::msg::dds\_::Odometry\_!}
\item
  签名：\passthrough{\lstinline"operator!=(const value \&) const"}
\item
  绑定策略：\passthrough{\lstinline!未单独标注!}
\item
  声明位置：由 pybind11 手工绑定或生成注册表提供；manifest
  未记录独立头文件行号。
\end{itemize}

\textbf{说明}

\begin{itemize}
\tightlist
\item
  示例是局部调用片段；\passthrough{\lstinline!obj!}
  和参数变量表示已经按上层业务校验并准备好的对象。
\end{itemize}

\textbf{用法}

\begin{lstlisting}[language=Python]
different = left != right
\end{lstlisting}

\hypertarget{unitree-sdk2-cpp-idl-ros2-odometry-header}{%
\paragraph{\texorpdfstring{\texttt{unitree\_sdk2\_cpp.idl.ros2.Odometry.header}}{unitree\_sdk2\_cpp.idl.ros2.Odometry.header}}\label{unitree-sdk2-cpp-idl-ros2-odometry-header}}

底层 \passthrough{\lstinline!nav\_msgs::msg::dds\_::Odometry\_!} 的
\passthrough{\lstinline!header!} 字段。类型签名只说明 Python
表示；单位、数值范围、数组固定长度和枚举含义需查对应 IDL 头文件。

\textbf{签名}

\begin{lstlisting}[language=Python]
@property
def header(self) -> Any

@header.setter
def header(self, value: Any) -> None
\end{lstlisting}

\textbf{可用性与安全性}

\textbf{\passthrough{\lstinline!AVAILABLE!} /
\passthrough{\lstinline!VALUE\_TYPE!}}：当前绑定源码已实现该消息值操作。

\textbf{参数}

\begin{longtable}[]{@{}
  >{\raggedright\arraybackslash}p{(\columnwidth - 4\tabcolsep) * \real{0.18}}
  >{\raggedright\arraybackslash}p{(\columnwidth - 4\tabcolsep) * \real{0.20}}
  >{\raggedright\arraybackslash}p{(\columnwidth - 4\tabcolsep) * \real{0.62}}@{}}
\toprule
名称 & Python 类型 & 含义 \\
\midrule
\endhead
\passthrough{\lstinline!value!} & \passthrough{\lstinline!Any!} &
要写入或传递的新值，必须符合该参数的 Python 类型以及底层 C++
范围约束。 \\
\bottomrule
\end{longtable}

\textbf{返回值}

\begin{longtable}[]{@{}
  >{\raggedright\arraybackslash}p{(\columnwidth - 4\tabcolsep) * \real{0.18}}
  >{\raggedright\arraybackslash}p{(\columnwidth - 4\tabcolsep) * \real{0.20}}
  >{\raggedright\arraybackslash}p{(\columnwidth - 4\tabcolsep) * \real{0.62}}@{}}
\toprule
位置 & 类型 & 含义 \\
\midrule
\endhead
getter 返回值 & \passthrough{\lstinline!Any!} & 返回属性当前值。 \\
\bottomrule
\end{longtable}

\textbf{对应 C++}

\begin{itemize}
\tightlist
\item
  类：\passthrough{\lstinline!nav\_msgs::msg::dds\_::Odometry\_!}
\item
  字段类型：\passthrough{\lstinline!::std\_msgs::msg::dds\_::Header\_!}
\item
  声明位置：\passthrough{\lstinline!include/unitree/idl/ros2/Odometry\_.hpp:29!}
\end{itemize}

\textbf{用法}

\begin{lstlisting}[language=Python]
current_value = obj.header
obj.header = new_value
\end{lstlisting}

\hypertarget{unitree-sdk2-cpp-idl-ros2-odometry-child-frame-id}{%
\paragraph{\texorpdfstring{\texttt{unitree\_sdk2\_cpp.idl.ros2.Odometry.child\_frame\_id}}{unitree\_sdk2\_cpp.idl.ros2.Odometry.child\_frame\_id}}\label{unitree-sdk2-cpp-idl-ros2-odometry-child-frame-id}}

底层 \passthrough{\lstinline!nav\_msgs::msg::dds\_::Odometry\_!} 的
\passthrough{\lstinline!child\_frame\_id!} 字段。类型签名只说明 Python
表示；单位、数值范围、数组固定长度和枚举含义需查对应 IDL 头文件。
底层为字符串；长度、编码和允许值由具体协议决定。

\textbf{签名}

\begin{lstlisting}[language=Python]
@property
def child_frame_id(self) -> str

@child_frame_id.setter
def child_frame_id(self, value: str) -> None
\end{lstlisting}

\textbf{可用性与安全性}

\textbf{\passthrough{\lstinline!AVAILABLE!} /
\passthrough{\lstinline!VALUE\_TYPE!}}：当前绑定源码已实现该消息值操作。

\textbf{参数}

\begin{longtable}[]{@{}
  >{\raggedright\arraybackslash}p{(\columnwidth - 4\tabcolsep) * \real{0.18}}
  >{\raggedright\arraybackslash}p{(\columnwidth - 4\tabcolsep) * \real{0.20}}
  >{\raggedright\arraybackslash}p{(\columnwidth - 4\tabcolsep) * \real{0.62}}@{}}
\toprule
名称 & Python 类型 & 含义 \\
\midrule
\endhead
\passthrough{\lstinline!value!} & \passthrough{\lstinline!str!} &
要写入或传递的新值，必须符合该参数的 Python 类型以及底层 C++
范围约束。 \\
\bottomrule
\end{longtable}

\textbf{返回值}

\begin{longtable}[]{@{}
  >{\raggedright\arraybackslash}p{(\columnwidth - 4\tabcolsep) * \real{0.18}}
  >{\raggedright\arraybackslash}p{(\columnwidth - 4\tabcolsep) * \real{0.20}}
  >{\raggedright\arraybackslash}p{(\columnwidth - 4\tabcolsep) * \real{0.62}}@{}}
\toprule
位置 & 类型 & 含义 \\
\midrule
\endhead
getter 返回值 & \passthrough{\lstinline!str!} & 返回属性当前值。 \\
\bottomrule
\end{longtable}

\textbf{对应 C++}

\begin{itemize}
\tightlist
\item
  类：\passthrough{\lstinline!nav\_msgs::msg::dds\_::Odometry\_!}
\item
  字段类型：\passthrough{\lstinline!std::string!}
\item
  声明位置：\passthrough{\lstinline!include/unitree/idl/ros2/Odometry\_.hpp:30!}
\end{itemize}

\textbf{用法}

\begin{lstlisting}[language=Python]
current_value = obj.child_frame_id
obj.child_frame_id = new_value
\end{lstlisting}

\hypertarget{unitree-sdk2-cpp-idl-ros2-odometry-pose}{%
\paragraph{\texorpdfstring{\texttt{unitree\_sdk2\_cpp.idl.ros2.Odometry.pose}}{unitree\_sdk2\_cpp.idl.ros2.Odometry.pose}}\label{unitree-sdk2-cpp-idl-ros2-odometry-pose}}

底层 \passthrough{\lstinline!nav\_msgs::msg::dds\_::Odometry\_!} 的
\passthrough{\lstinline!pose!} 字段。类型签名只说明 Python
表示；单位、数值范围、数组固定长度和枚举含义需查对应 IDL 头文件。

\textbf{签名}

\begin{lstlisting}[language=Python]
@property
def pose(self) -> Any

@pose.setter
def pose(self, value: Any) -> None
\end{lstlisting}

\textbf{可用性与安全性}

\textbf{\passthrough{\lstinline!AVAILABLE!} /
\passthrough{\lstinline!VALUE\_TYPE!}}：当前绑定源码已实现该消息值操作。

\textbf{参数}

\begin{longtable}[]{@{}
  >{\raggedright\arraybackslash}p{(\columnwidth - 4\tabcolsep) * \real{0.18}}
  >{\raggedright\arraybackslash}p{(\columnwidth - 4\tabcolsep) * \real{0.20}}
  >{\raggedright\arraybackslash}p{(\columnwidth - 4\tabcolsep) * \real{0.62}}@{}}
\toprule
名称 & Python 类型 & 含义 \\
\midrule
\endhead
\passthrough{\lstinline!value!} & \passthrough{\lstinline!Any!} &
要写入或传递的新值，必须符合该参数的 Python 类型以及底层 C++
范围约束。 \\
\bottomrule
\end{longtable}

\textbf{返回值}

\begin{longtable}[]{@{}
  >{\raggedright\arraybackslash}p{(\columnwidth - 4\tabcolsep) * \real{0.18}}
  >{\raggedright\arraybackslash}p{(\columnwidth - 4\tabcolsep) * \real{0.20}}
  >{\raggedright\arraybackslash}p{(\columnwidth - 4\tabcolsep) * \real{0.62}}@{}}
\toprule
位置 & 类型 & 含义 \\
\midrule
\endhead
getter 返回值 & \passthrough{\lstinline!Any!} & 返回属性当前值。 \\
\bottomrule
\end{longtable}

\textbf{对应 C++}

\begin{itemize}
\tightlist
\item
  类：\passthrough{\lstinline!nav\_msgs::msg::dds\_::Odometry\_!}
\item
  字段类型：\passthrough{\lstinline!::geometry\_msgs::msg::dds\_::PoseWithCovariance\_!}
\item
  声明位置：\passthrough{\lstinline!include/unitree/idl/ros2/Odometry\_.hpp:31!}
\end{itemize}

\textbf{用法}

\begin{lstlisting}[language=Python]
current_value = obj.pose
obj.pose = new_value
\end{lstlisting}

\hypertarget{unitree-sdk2-cpp-idl-ros2-odometry-twist}{%
\paragraph{\texorpdfstring{\texttt{unitree\_sdk2\_cpp.idl.ros2.Odometry.twist}}{unitree\_sdk2\_cpp.idl.ros2.Odometry.twist}}\label{unitree-sdk2-cpp-idl-ros2-odometry-twist}}

底层 \passthrough{\lstinline!nav\_msgs::msg::dds\_::Odometry\_!} 的
\passthrough{\lstinline!twist!} 字段。类型签名只说明 Python
表示；单位、数值范围、数组固定长度和枚举含义需查对应 IDL 头文件。

\textbf{签名}

\begin{lstlisting}[language=Python]
@property
def twist(self) -> Any

@twist.setter
def twist(self, value: Any) -> None
\end{lstlisting}

\textbf{可用性与安全性}

\textbf{\passthrough{\lstinline!AVAILABLE!} /
\passthrough{\lstinline!VALUE\_TYPE!}}：当前绑定源码已实现该消息值操作。

\textbf{参数}

\begin{longtable}[]{@{}
  >{\raggedright\arraybackslash}p{(\columnwidth - 4\tabcolsep) * \real{0.18}}
  >{\raggedright\arraybackslash}p{(\columnwidth - 4\tabcolsep) * \real{0.20}}
  >{\raggedright\arraybackslash}p{(\columnwidth - 4\tabcolsep) * \real{0.62}}@{}}
\toprule
名称 & Python 类型 & 含义 \\
\midrule
\endhead
\passthrough{\lstinline!value!} & \passthrough{\lstinline!Any!} &
要写入或传递的新值，必须符合该参数的 Python 类型以及底层 C++
范围约束。 \\
\bottomrule
\end{longtable}

\textbf{返回值}

\begin{longtable}[]{@{}
  >{\raggedright\arraybackslash}p{(\columnwidth - 4\tabcolsep) * \real{0.18}}
  >{\raggedright\arraybackslash}p{(\columnwidth - 4\tabcolsep) * \real{0.20}}
  >{\raggedright\arraybackslash}p{(\columnwidth - 4\tabcolsep) * \real{0.62}}@{}}
\toprule
位置 & 类型 & 含义 \\
\midrule
\endhead
getter 返回值 & \passthrough{\lstinline!Any!} & 返回属性当前值。 \\
\bottomrule
\end{longtable}

\textbf{对应 C++}

\begin{itemize}
\tightlist
\item
  类：\passthrough{\lstinline!nav\_msgs::msg::dds\_::Odometry\_!}
\item
  字段类型：\passthrough{\lstinline!::geometry\_msgs::msg::dds\_::TwistWithCovariance\_!}
\item
  声明位置：\passthrough{\lstinline!include/unitree/idl/ros2/Odometry\_.hpp:32!}
\end{itemize}

\textbf{用法}

\begin{lstlisting}[language=Python]
current_value = obj.twist
obj.twist = new_value
\end{lstlisting}

\hypertarget{unitree-sdk2-cpp-idl-ros2-point32}{%
\subsubsection{\texorpdfstring{\texttt{unitree\_sdk2\_cpp.idl.ros2.Point32}}{unitree\_sdk2\_cpp.idl.ros2.Point32}}\label{unitree-sdk2-cpp-idl-ros2-point32}}

可由 Python 读写的 DDS/IDL 消息值类型。字段采用复制语义。

\textbf{导入}

\begin{lstlisting}[language=Python]
from unitree_sdk2_cpp.idl.ros2 import Point32
\end{lstlisting}

\textbf{方法索引}

\begin{longtable}[]{@{}
  >{\raggedright\arraybackslash}p{(\columnwidth - 6\tabcolsep) * \real{0.18}}
  >{\raggedright\arraybackslash}p{(\columnwidth - 6\tabcolsep) * \real{0.24}}
  >{\raggedright\arraybackslash}p{(\columnwidth - 6\tabcolsep) * \real{0.18}}
  >{\raggedright\arraybackslash}p{(\columnwidth - 6\tabcolsep) * \real{0.40}}@{}}
\toprule
方法 & Python 签名 & 状态 & 安全分类 \\
\midrule
\endhead
\protect\hyperlink{unitree-sdk2-cpp-idl-ros2-point32-dunder-init-1}{\passthrough{\lstinline!\_\_init\_\_!}}
& \passthrough{\lstinline!def \_\_init\_\_(self) -> None!} &
\passthrough{\lstinline!AVAILABLE!} &
\passthrough{\lstinline!VALUE\_TYPE!} \\
\protect\hyperlink{unitree-sdk2-cpp-idl-ros2-point32-dunder-eq-1}{\passthrough{\lstinline!\_\_eq\_\_!}}
& \passthrough{\lstinline!def \_\_eq\_\_(self, other: object) -> bool!}
& \passthrough{\lstinline!AVAILABLE!} &
\passthrough{\lstinline!VALUE\_TYPE!} \\
\protect\hyperlink{unitree-sdk2-cpp-idl-ros2-point32-dunder-ne-1}{\passthrough{\lstinline!\_\_ne\_\_!}}
& \passthrough{\lstinline!def \_\_ne\_\_(self, other: object) -> bool!}
& \passthrough{\lstinline!AVAILABLE!} &
\passthrough{\lstinline!VALUE\_TYPE!} \\
\bottomrule
\end{longtable}

\textbf{属性索引}

\begin{longtable}[]{@{}
  >{\raggedright\arraybackslash}p{(\columnwidth - 6\tabcolsep) * \real{0.18}}
  >{\raggedright\arraybackslash}p{(\columnwidth - 6\tabcolsep) * \real{0.24}}
  >{\raggedright\arraybackslash}p{(\columnwidth - 6\tabcolsep) * \real{0.18}}
  >{\raggedright\arraybackslash}p{(\columnwidth - 6\tabcolsep) * \real{0.40}}@{}}
\toprule
属性 & 读取类型 & 写入类型 & 状态 \\
\midrule
\endhead
\protect\hyperlink{unitree-sdk2-cpp-idl-ros2-point32-x}{\passthrough{\lstinline!x!}}
& \passthrough{\lstinline!float!} & \passthrough{\lstinline!float!} &
\passthrough{\lstinline!AVAILABLE!} \\
\protect\hyperlink{unitree-sdk2-cpp-idl-ros2-point32-y}{\passthrough{\lstinline!y!}}
& \passthrough{\lstinline!float!} & \passthrough{\lstinline!float!} &
\passthrough{\lstinline!AVAILABLE!} \\
\protect\hyperlink{unitree-sdk2-cpp-idl-ros2-point32-z}{\passthrough{\lstinline!z!}}
& \passthrough{\lstinline!float!} & \passthrough{\lstinline!float!} &
\passthrough{\lstinline!AVAILABLE!} \\
\bottomrule
\end{longtable}

\hypertarget{unitree-sdk2-cpp-idl-ros2-point32-dunder-init-1}{%
\paragraph{\texorpdfstring{\texttt{unitree\_sdk2\_cpp.idl.ros2.Point32.\_\_init\_\_}}{unitree\_sdk2\_cpp.idl.ros2.Point32.\_\_init\_\_}}\label{unitree-sdk2-cpp-idl-ros2-point32-dunder-init-1}}

初始化 \passthrough{\lstinline!Point32!}
实例。是否能实际构造取决于下方可用性状态。

\textbf{签名}

\begin{lstlisting}[language=Python]
def __init__(self) -> None
\end{lstlisting}

\textbf{可用性与安全性}

\textbf{\passthrough{\lstinline!AVAILABLE!} /
\passthrough{\lstinline!VALUE\_TYPE!}}：当前绑定源码已实现该消息值操作。

\textbf{参数}

无显式参数。实例方法中的 \passthrough{\lstinline!self!} 由 Python
自动传入。

\textbf{返回值}

\begin{longtable}[]{@{}
  >{\raggedright\arraybackslash}p{(\columnwidth - 4\tabcolsep) * \real{0.18}}
  >{\raggedright\arraybackslash}p{(\columnwidth - 4\tabcolsep) * \real{0.20}}
  >{\raggedright\arraybackslash}p{(\columnwidth - 4\tabcolsep) * \real{0.62}}@{}}
\toprule
位置 & 类型 & 含义 \\
\midrule
\endhead
无 & \passthrough{\lstinline!None!} &
\passthrough{\lstinline!\_\_init\_\_!}
本身不返回值；调用类对象时，在构造函数可用的前提下得到该类实例。 \\
\bottomrule
\end{longtable}

\textbf{对应 C++}

\begin{itemize}
\tightlist
\item
  类：\passthrough{\lstinline!geometry\_msgs::msg::dds\_::Point32\_!}
\item
  签名：\passthrough{\lstinline!Point32\_()!}
\item
  绑定策略：\passthrough{\lstinline!未单独标注!}
\item
  声明位置：\passthrough{\lstinline!include/unitree/idl/ros2/Point32\_.hpp:27!}
\end{itemize}

\textbf{说明}

\begin{itemize}
\tightlist
\item
  示例是局部调用片段；\passthrough{\lstinline!obj!}
  和参数变量表示已经按上层业务校验并准备好的对象。
\end{itemize}

\textbf{用法}

\begin{lstlisting}[language=Python]
value = Point32()
\end{lstlisting}

\hypertarget{unitree-sdk2-cpp-idl-ros2-point32-dunder-eq-1}{%
\paragraph{\texorpdfstring{\texttt{unitree\_sdk2\_cpp.idl.ros2.Point32.\_\_eq\_\_}}{unitree\_sdk2\_cpp.idl.ros2.Point32.\_\_eq\_\_}}\label{unitree-sdk2-cpp-idl-ros2-point32-dunder-eq-1}}

按底层 IDL 消息内容比较两个对象是否相等。

\textbf{签名}

\begin{lstlisting}[language=Python]
def __eq__(self, other: object) -> bool
\end{lstlisting}

\textbf{可用性与安全性}

\textbf{\passthrough{\lstinline!AVAILABLE!} /
\passthrough{\lstinline!VALUE\_TYPE!}}：当前绑定源码已实现该消息值操作。

\textbf{参数}

\begin{longtable}[]{@{}
  >{\raggedright\arraybackslash}p{(\columnwidth - 6\tabcolsep) * \real{0.18}}
  >{\raggedright\arraybackslash}p{(\columnwidth - 6\tabcolsep) * \real{0.24}}
  >{\raggedright\arraybackslash}p{(\columnwidth - 6\tabcolsep) * \real{0.18}}
  >{\raggedright\arraybackslash}p{(\columnwidth - 6\tabcolsep) * \real{0.40}}@{}}
\toprule
名称 & Python 类型 & 默认值 & 含义 \\
\midrule
\endhead
\passthrough{\lstinline!other!} & \passthrough{\lstinline!object!} &
必填 & 用于比较的另一个 Python 对象。类型不兼容时比较结果通常为
\passthrough{\lstinline!False!}。 \\
\bottomrule
\end{longtable}

\textbf{返回值}

\begin{longtable}[]{@{}
  >{\raggedright\arraybackslash}p{(\columnwidth - 4\tabcolsep) * \real{0.18}}
  >{\raggedright\arraybackslash}p{(\columnwidth - 4\tabcolsep) * \real{0.20}}
  >{\raggedright\arraybackslash}p{(\columnwidth - 4\tabcolsep) * \real{0.62}}@{}}
\toprule
位置 & 类型 & 含义 \\
\midrule
\endhead
返回值 & \passthrough{\lstinline!bool!} & 返回
\passthrough{\lstinline!bool!}。更精确的业务含义见该方法用途、C++
签名和目标型号协议。 \\
\bottomrule
\end{longtable}

\textbf{对应 C++}

\begin{itemize}
\tightlist
\item
  类：\passthrough{\lstinline!geometry\_msgs::msg::dds\_::Point32\_!}
\item
  签名：\passthrough{\lstinline!operator==(const value \&) const!}
\item
  绑定策略：\passthrough{\lstinline!未单独标注!}
\item
  声明位置：由 pybind11 手工绑定或生成注册表提供；manifest
  未记录独立头文件行号。
\end{itemize}

\textbf{说明}

\begin{itemize}
\tightlist
\item
  示例是局部调用片段；\passthrough{\lstinline!obj!}
  和参数变量表示已经按上层业务校验并准备好的对象。
\end{itemize}

\textbf{用法}

\begin{lstlisting}[language=Python]
same = left == right
\end{lstlisting}

\hypertarget{unitree-sdk2-cpp-idl-ros2-point32-dunder-ne-1}{%
\paragraph{\texorpdfstring{\texttt{unitree\_sdk2\_cpp.idl.ros2.Point32.\_\_ne\_\_}}{unitree\_sdk2\_cpp.idl.ros2.Point32.\_\_ne\_\_}}\label{unitree-sdk2-cpp-idl-ros2-point32-dunder-ne-1}}

按底层 IDL 消息内容比较两个对象是否不相等。

\textbf{签名}

\begin{lstlisting}[language=Python]
def __ne__(self, other: object) -> bool
\end{lstlisting}

\textbf{可用性与安全性}

\textbf{\passthrough{\lstinline!AVAILABLE!} /
\passthrough{\lstinline!VALUE\_TYPE!}}：当前绑定源码已实现该消息值操作。

\textbf{参数}

\begin{longtable}[]{@{}
  >{\raggedright\arraybackslash}p{(\columnwidth - 6\tabcolsep) * \real{0.18}}
  >{\raggedright\arraybackslash}p{(\columnwidth - 6\tabcolsep) * \real{0.24}}
  >{\raggedright\arraybackslash}p{(\columnwidth - 6\tabcolsep) * \real{0.18}}
  >{\raggedright\arraybackslash}p{(\columnwidth - 6\tabcolsep) * \real{0.40}}@{}}
\toprule
名称 & Python 类型 & 默认值 & 含义 \\
\midrule
\endhead
\passthrough{\lstinline!other!} & \passthrough{\lstinline!object!} &
必填 & 用于比较的另一个 Python 对象。类型不兼容时比较结果通常为
\passthrough{\lstinline!False!}。 \\
\bottomrule
\end{longtable}

\textbf{返回值}

\begin{longtable}[]{@{}
  >{\raggedright\arraybackslash}p{(\columnwidth - 4\tabcolsep) * \real{0.18}}
  >{\raggedright\arraybackslash}p{(\columnwidth - 4\tabcolsep) * \real{0.20}}
  >{\raggedright\arraybackslash}p{(\columnwidth - 4\tabcolsep) * \real{0.62}}@{}}
\toprule
位置 & 类型 & 含义 \\
\midrule
\endhead
返回值 & \passthrough{\lstinline!bool!} & 返回
\passthrough{\lstinline!bool!}。更精确的业务含义见该方法用途、C++
签名和目标型号协议。 \\
\bottomrule
\end{longtable}

\textbf{对应 C++}

\begin{itemize}
\tightlist
\item
  类：\passthrough{\lstinline!geometry\_msgs::msg::dds\_::Point32\_!}
\item
  签名：\passthrough{\lstinline"operator!=(const value \&) const"}
\item
  绑定策略：\passthrough{\lstinline!未单独标注!}
\item
  声明位置：由 pybind11 手工绑定或生成注册表提供；manifest
  未记录独立头文件行号。
\end{itemize}

\textbf{说明}

\begin{itemize}
\tightlist
\item
  示例是局部调用片段；\passthrough{\lstinline!obj!}
  和参数变量表示已经按上层业务校验并准备好的对象。
\end{itemize}

\textbf{用法}

\begin{lstlisting}[language=Python]
different = left != right
\end{lstlisting}

\hypertarget{unitree-sdk2-cpp-idl-ros2-point32-x}{%
\paragraph{\texorpdfstring{\texttt{unitree\_sdk2\_cpp.idl.ros2.Point32.x}}{unitree\_sdk2\_cpp.idl.ros2.Point32.x}}\label{unitree-sdk2-cpp-idl-ros2-point32-x}}

底层 \passthrough{\lstinline!geometry\_msgs::msg::dds\_::Point32\_!} 的
\passthrough{\lstinline!x!} 字段。类型签名只说明 Python
表示；单位、数值范围、数组固定长度和枚举含义需查对应 IDL 头文件。
底层为浮点数；类型本身不说明单位、坐标系或业务安全范围。

\textbf{签名}

\begin{lstlisting}[language=Python]
@property
def x(self) -> float

@x.setter
def x(self, value: float) -> None
\end{lstlisting}

\textbf{可用性与安全性}

\textbf{\passthrough{\lstinline!AVAILABLE!} /
\passthrough{\lstinline!VALUE\_TYPE!}}：当前绑定源码已实现该消息值操作。

\textbf{参数}

\begin{longtable}[]{@{}
  >{\raggedright\arraybackslash}p{(\columnwidth - 4\tabcolsep) * \real{0.18}}
  >{\raggedright\arraybackslash}p{(\columnwidth - 4\tabcolsep) * \real{0.20}}
  >{\raggedright\arraybackslash}p{(\columnwidth - 4\tabcolsep) * \real{0.62}}@{}}
\toprule
名称 & Python 类型 & 含义 \\
\midrule
\endhead
\passthrough{\lstinline!value!} & \passthrough{\lstinline!float!} &
要写入或传递的新值，必须符合该参数的 Python 类型以及底层 C++
范围约束。 \\
\bottomrule
\end{longtable}

\textbf{返回值}

\begin{longtable}[]{@{}
  >{\raggedright\arraybackslash}p{(\columnwidth - 4\tabcolsep) * \real{0.18}}
  >{\raggedright\arraybackslash}p{(\columnwidth - 4\tabcolsep) * \real{0.20}}
  >{\raggedright\arraybackslash}p{(\columnwidth - 4\tabcolsep) * \real{0.62}}@{}}
\toprule
位置 & 类型 & 含义 \\
\midrule
\endhead
getter 返回值 & \passthrough{\lstinline!float!} & 返回属性当前值。 \\
\bottomrule
\end{longtable}

\textbf{对应 C++}

\begin{itemize}
\tightlist
\item
  类：\passthrough{\lstinline!geometry\_msgs::msg::dds\_::Point32\_!}
\item
  字段类型：\passthrough{\lstinline!float!}
\item
  声明位置：\passthrough{\lstinline!include/unitree/idl/ros2/Point32\_.hpp:22!}
\end{itemize}

\textbf{用法}

\begin{lstlisting}[language=Python]
current_value = obj.x
obj.x = new_value
\end{lstlisting}

\hypertarget{unitree-sdk2-cpp-idl-ros2-point32-y}{%
\paragraph{\texorpdfstring{\texttt{unitree\_sdk2\_cpp.idl.ros2.Point32.y}}{unitree\_sdk2\_cpp.idl.ros2.Point32.y}}\label{unitree-sdk2-cpp-idl-ros2-point32-y}}

底层 \passthrough{\lstinline!geometry\_msgs::msg::dds\_::Point32\_!} 的
\passthrough{\lstinline!y!} 字段。类型签名只说明 Python
表示；单位、数值范围、数组固定长度和枚举含义需查对应 IDL 头文件。
底层为浮点数；类型本身不说明单位、坐标系或业务安全范围。

\textbf{签名}

\begin{lstlisting}[language=Python]
@property
def y(self) -> float

@y.setter
def y(self, value: float) -> None
\end{lstlisting}

\textbf{可用性与安全性}

\textbf{\passthrough{\lstinline!AVAILABLE!} /
\passthrough{\lstinline!VALUE\_TYPE!}}：当前绑定源码已实现该消息值操作。

\textbf{参数}

\begin{longtable}[]{@{}
  >{\raggedright\arraybackslash}p{(\columnwidth - 4\tabcolsep) * \real{0.18}}
  >{\raggedright\arraybackslash}p{(\columnwidth - 4\tabcolsep) * \real{0.20}}
  >{\raggedright\arraybackslash}p{(\columnwidth - 4\tabcolsep) * \real{0.62}}@{}}
\toprule
名称 & Python 类型 & 含义 \\
\midrule
\endhead
\passthrough{\lstinline!value!} & \passthrough{\lstinline!float!} &
要写入或传递的新值，必须符合该参数的 Python 类型以及底层 C++
范围约束。 \\
\bottomrule
\end{longtable}

\textbf{返回值}

\begin{longtable}[]{@{}
  >{\raggedright\arraybackslash}p{(\columnwidth - 4\tabcolsep) * \real{0.18}}
  >{\raggedright\arraybackslash}p{(\columnwidth - 4\tabcolsep) * \real{0.20}}
  >{\raggedright\arraybackslash}p{(\columnwidth - 4\tabcolsep) * \real{0.62}}@{}}
\toprule
位置 & 类型 & 含义 \\
\midrule
\endhead
getter 返回值 & \passthrough{\lstinline!float!} & 返回属性当前值。 \\
\bottomrule
\end{longtable}

\textbf{对应 C++}

\begin{itemize}
\tightlist
\item
  类：\passthrough{\lstinline!geometry\_msgs::msg::dds\_::Point32\_!}
\item
  字段类型：\passthrough{\lstinline!float!}
\item
  声明位置：\passthrough{\lstinline!include/unitree/idl/ros2/Point32\_.hpp:23!}
\end{itemize}

\textbf{用法}

\begin{lstlisting}[language=Python]
current_value = obj.y
obj.y = new_value
\end{lstlisting}

\hypertarget{unitree-sdk2-cpp-idl-ros2-point32-z}{%
\paragraph{\texorpdfstring{\texttt{unitree\_sdk2\_cpp.idl.ros2.Point32.z}}{unitree\_sdk2\_cpp.idl.ros2.Point32.z}}\label{unitree-sdk2-cpp-idl-ros2-point32-z}}

底层 \passthrough{\lstinline!geometry\_msgs::msg::dds\_::Point32\_!} 的
\passthrough{\lstinline!z!} 字段。类型签名只说明 Python
表示；单位、数值范围、数组固定长度和枚举含义需查对应 IDL 头文件。
底层为浮点数；类型本身不说明单位、坐标系或业务安全范围。

\textbf{签名}

\begin{lstlisting}[language=Python]
@property
def z(self) -> float

@z.setter
def z(self, value: float) -> None
\end{lstlisting}

\textbf{可用性与安全性}

\textbf{\passthrough{\lstinline!AVAILABLE!} /
\passthrough{\lstinline!VALUE\_TYPE!}}：当前绑定源码已实现该消息值操作。

\textbf{参数}

\begin{longtable}[]{@{}
  >{\raggedright\arraybackslash}p{(\columnwidth - 4\tabcolsep) * \real{0.18}}
  >{\raggedright\arraybackslash}p{(\columnwidth - 4\tabcolsep) * \real{0.20}}
  >{\raggedright\arraybackslash}p{(\columnwidth - 4\tabcolsep) * \real{0.62}}@{}}
\toprule
名称 & Python 类型 & 含义 \\
\midrule
\endhead
\passthrough{\lstinline!value!} & \passthrough{\lstinline!float!} &
要写入或传递的新值，必须符合该参数的 Python 类型以及底层 C++
范围约束。 \\
\bottomrule
\end{longtable}

\textbf{返回值}

\begin{longtable}[]{@{}
  >{\raggedright\arraybackslash}p{(\columnwidth - 4\tabcolsep) * \real{0.18}}
  >{\raggedright\arraybackslash}p{(\columnwidth - 4\tabcolsep) * \real{0.20}}
  >{\raggedright\arraybackslash}p{(\columnwidth - 4\tabcolsep) * \real{0.62}}@{}}
\toprule
位置 & 类型 & 含义 \\
\midrule
\endhead
getter 返回值 & \passthrough{\lstinline!float!} & 返回属性当前值。 \\
\bottomrule
\end{longtable}

\textbf{对应 C++}

\begin{itemize}
\tightlist
\item
  类：\passthrough{\lstinline!geometry\_msgs::msg::dds\_::Point32\_!}
\item
  字段类型：\passthrough{\lstinline!float!}
\item
  声明位置：\passthrough{\lstinline!include/unitree/idl/ros2/Point32\_.hpp:24!}
\end{itemize}

\textbf{用法}

\begin{lstlisting}[language=Python]
current_value = obj.z
obj.z = new_value
\end{lstlisting}

\hypertarget{unitree-sdk2-cpp-idl-ros2-pointfield}{%
\subsubsection{\texorpdfstring{\texttt{unitree\_sdk2\_cpp.idl.ros2.PointField}}{unitree\_sdk2\_cpp.idl.ros2.PointField}}\label{unitree-sdk2-cpp-idl-ros2-pointfield}}

可由 Python 读写的 DDS/IDL 消息值类型。字段采用复制语义。

\textbf{导入}

\begin{lstlisting}[language=Python]
from unitree_sdk2_cpp.idl.ros2 import PointField
\end{lstlisting}

\textbf{方法索引}

\begin{longtable}[]{@{}
  >{\raggedright\arraybackslash}p{(\columnwidth - 6\tabcolsep) * \real{0.18}}
  >{\raggedright\arraybackslash}p{(\columnwidth - 6\tabcolsep) * \real{0.24}}
  >{\raggedright\arraybackslash}p{(\columnwidth - 6\tabcolsep) * \real{0.18}}
  >{\raggedright\arraybackslash}p{(\columnwidth - 6\tabcolsep) * \real{0.40}}@{}}
\toprule
方法 & Python 签名 & 状态 & 安全分类 \\
\midrule
\endhead
\protect\hyperlink{unitree-sdk2-cpp-idl-ros2-pointfield-dunder-init-1}{\passthrough{\lstinline!\_\_init\_\_!}}
& \passthrough{\lstinline!def \_\_init\_\_(self) -> None!} &
\passthrough{\lstinline!AVAILABLE!} &
\passthrough{\lstinline!VALUE\_TYPE!} \\
\protect\hyperlink{unitree-sdk2-cpp-idl-ros2-pointfield-dunder-eq-1}{\passthrough{\lstinline!\_\_eq\_\_!}}
& \passthrough{\lstinline!def \_\_eq\_\_(self, other: object) -> bool!}
& \passthrough{\lstinline!AVAILABLE!} &
\passthrough{\lstinline!VALUE\_TYPE!} \\
\protect\hyperlink{unitree-sdk2-cpp-idl-ros2-pointfield-dunder-ne-1}{\passthrough{\lstinline!\_\_ne\_\_!}}
& \passthrough{\lstinline!def \_\_ne\_\_(self, other: object) -> bool!}
& \passthrough{\lstinline!AVAILABLE!} &
\passthrough{\lstinline!VALUE\_TYPE!} \\
\bottomrule
\end{longtable}

\textbf{属性索引}

\begin{longtable}[]{@{}
  >{\raggedright\arraybackslash}p{(\columnwidth - 6\tabcolsep) * \real{0.18}}
  >{\raggedright\arraybackslash}p{(\columnwidth - 6\tabcolsep) * \real{0.24}}
  >{\raggedright\arraybackslash}p{(\columnwidth - 6\tabcolsep) * \real{0.18}}
  >{\raggedright\arraybackslash}p{(\columnwidth - 6\tabcolsep) * \real{0.40}}@{}}
\toprule
属性 & 读取类型 & 写入类型 & 状态 \\
\midrule
\endhead
\protect\hyperlink{unitree-sdk2-cpp-idl-ros2-pointfield-name}{\passthrough{\lstinline!name!}}
& \passthrough{\lstinline!str!} & \passthrough{\lstinline!str!} &
\passthrough{\lstinline!AVAILABLE!} \\
\protect\hyperlink{unitree-sdk2-cpp-idl-ros2-pointfield-offset}{\passthrough{\lstinline!offset!}}
& \passthrough{\lstinline!int!} & \passthrough{\lstinline!int!} &
\passthrough{\lstinline!AVAILABLE!} \\
\protect\hyperlink{unitree-sdk2-cpp-idl-ros2-pointfield-datatype}{\passthrough{\lstinline!datatype!}}
& \passthrough{\lstinline!int!} & \passthrough{\lstinline!int!} &
\passthrough{\lstinline!AVAILABLE!} \\
\protect\hyperlink{unitree-sdk2-cpp-idl-ros2-pointfield-count}{\passthrough{\lstinline!count!}}
& \passthrough{\lstinline!int!} & \passthrough{\lstinline!int!} &
\passthrough{\lstinline!AVAILABLE!} \\
\bottomrule
\end{longtable}

\hypertarget{unitree-sdk2-cpp-idl-ros2-pointfield-dunder-init-1}{%
\paragraph{\texorpdfstring{\texttt{unitree\_sdk2\_cpp.idl.ros2.PointField.\_\_init\_\_}}{unitree\_sdk2\_cpp.idl.ros2.PointField.\_\_init\_\_}}\label{unitree-sdk2-cpp-idl-ros2-pointfield-dunder-init-1}}

初始化 \passthrough{\lstinline!PointField!}
实例。是否能实际构造取决于下方可用性状态。

\textbf{签名}

\begin{lstlisting}[language=Python]
def __init__(self) -> None
\end{lstlisting}

\textbf{可用性与安全性}

\textbf{\passthrough{\lstinline!AVAILABLE!} /
\passthrough{\lstinline!VALUE\_TYPE!}}：当前绑定源码已实现该消息值操作。

\textbf{参数}

无显式参数。实例方法中的 \passthrough{\lstinline!self!} 由 Python
自动传入。

\textbf{返回值}

\begin{longtable}[]{@{}
  >{\raggedright\arraybackslash}p{(\columnwidth - 4\tabcolsep) * \real{0.18}}
  >{\raggedright\arraybackslash}p{(\columnwidth - 4\tabcolsep) * \real{0.20}}
  >{\raggedright\arraybackslash}p{(\columnwidth - 4\tabcolsep) * \real{0.62}}@{}}
\toprule
位置 & 类型 & 含义 \\
\midrule
\endhead
无 & \passthrough{\lstinline!None!} &
\passthrough{\lstinline!\_\_init\_\_!}
本身不返回值；调用类对象时，在构造函数可用的前提下得到该类实例。 \\
\bottomrule
\end{longtable}

\textbf{对应 C++}

\begin{itemize}
\tightlist
\item
  类：\passthrough{\lstinline!sensor\_msgs::msg::dds\_::PointField\_!}
\item
  签名：\passthrough{\lstinline!PointField\_()!}
\item
  绑定策略：\passthrough{\lstinline!未单独标注!}
\item
  声明位置：\passthrough{\lstinline!include/unitree/idl/ros2/PointField\_.hpp:50!}
\end{itemize}

\textbf{说明}

\begin{itemize}
\tightlist
\item
  示例是局部调用片段；\passthrough{\lstinline!obj!}
  和参数变量表示已经按上层业务校验并准备好的对象。
\end{itemize}

\textbf{用法}

\begin{lstlisting}[language=Python]
value = PointField()
\end{lstlisting}

\hypertarget{unitree-sdk2-cpp-idl-ros2-pointfield-dunder-eq-1}{%
\paragraph{\texorpdfstring{\texttt{unitree\_sdk2\_cpp.idl.ros2.PointField.\_\_eq\_\_}}{unitree\_sdk2\_cpp.idl.ros2.PointField.\_\_eq\_\_}}\label{unitree-sdk2-cpp-idl-ros2-pointfield-dunder-eq-1}}

按底层 IDL 消息内容比较两个对象是否相等。

\textbf{签名}

\begin{lstlisting}[language=Python]
def __eq__(self, other: object) -> bool
\end{lstlisting}

\textbf{可用性与安全性}

\textbf{\passthrough{\lstinline!AVAILABLE!} /
\passthrough{\lstinline!VALUE\_TYPE!}}：当前绑定源码已实现该消息值操作。

\textbf{参数}

\begin{longtable}[]{@{}
  >{\raggedright\arraybackslash}p{(\columnwidth - 6\tabcolsep) * \real{0.18}}
  >{\raggedright\arraybackslash}p{(\columnwidth - 6\tabcolsep) * \real{0.24}}
  >{\raggedright\arraybackslash}p{(\columnwidth - 6\tabcolsep) * \real{0.18}}
  >{\raggedright\arraybackslash}p{(\columnwidth - 6\tabcolsep) * \real{0.40}}@{}}
\toprule
名称 & Python 类型 & 默认值 & 含义 \\
\midrule
\endhead
\passthrough{\lstinline!other!} & \passthrough{\lstinline!object!} &
必填 & 用于比较的另一个 Python 对象。类型不兼容时比较结果通常为
\passthrough{\lstinline!False!}。 \\
\bottomrule
\end{longtable}

\textbf{返回值}

\begin{longtable}[]{@{}
  >{\raggedright\arraybackslash}p{(\columnwidth - 4\tabcolsep) * \real{0.18}}
  >{\raggedright\arraybackslash}p{(\columnwidth - 4\tabcolsep) * \real{0.20}}
  >{\raggedright\arraybackslash}p{(\columnwidth - 4\tabcolsep) * \real{0.62}}@{}}
\toprule
位置 & 类型 & 含义 \\
\midrule
\endhead
返回值 & \passthrough{\lstinline!bool!} & 返回
\passthrough{\lstinline!bool!}。更精确的业务含义见该方法用途、C++
签名和目标型号协议。 \\
\bottomrule
\end{longtable}

\textbf{对应 C++}

\begin{itemize}
\tightlist
\item
  类：\passthrough{\lstinline!sensor\_msgs::msg::dds\_::PointField\_!}
\item
  签名：\passthrough{\lstinline!operator==(const value \&) const!}
\item
  绑定策略：\passthrough{\lstinline!未单独标注!}
\item
  声明位置：由 pybind11 手工绑定或生成注册表提供；manifest
  未记录独立头文件行号。
\end{itemize}

\textbf{说明}

\begin{itemize}
\tightlist
\item
  示例是局部调用片段；\passthrough{\lstinline!obj!}
  和参数变量表示已经按上层业务校验并准备好的对象。
\end{itemize}

\textbf{用法}

\begin{lstlisting}[language=Python]
same = left == right
\end{lstlisting}

\hypertarget{unitree-sdk2-cpp-idl-ros2-pointfield-dunder-ne-1}{%
\paragraph{\texorpdfstring{\texttt{unitree\_sdk2\_cpp.idl.ros2.PointField.\_\_ne\_\_}}{unitree\_sdk2\_cpp.idl.ros2.PointField.\_\_ne\_\_}}\label{unitree-sdk2-cpp-idl-ros2-pointfield-dunder-ne-1}}

按底层 IDL 消息内容比较两个对象是否不相等。

\textbf{签名}

\begin{lstlisting}[language=Python]
def __ne__(self, other: object) -> bool
\end{lstlisting}

\textbf{可用性与安全性}

\textbf{\passthrough{\lstinline!AVAILABLE!} /
\passthrough{\lstinline!VALUE\_TYPE!}}：当前绑定源码已实现该消息值操作。

\textbf{参数}

\begin{longtable}[]{@{}
  >{\raggedright\arraybackslash}p{(\columnwidth - 6\tabcolsep) * \real{0.18}}
  >{\raggedright\arraybackslash}p{(\columnwidth - 6\tabcolsep) * \real{0.24}}
  >{\raggedright\arraybackslash}p{(\columnwidth - 6\tabcolsep) * \real{0.18}}
  >{\raggedright\arraybackslash}p{(\columnwidth - 6\tabcolsep) * \real{0.40}}@{}}
\toprule
名称 & Python 类型 & 默认值 & 含义 \\
\midrule
\endhead
\passthrough{\lstinline!other!} & \passthrough{\lstinline!object!} &
必填 & 用于比较的另一个 Python 对象。类型不兼容时比较结果通常为
\passthrough{\lstinline!False!}。 \\
\bottomrule
\end{longtable}

\textbf{返回值}

\begin{longtable}[]{@{}
  >{\raggedright\arraybackslash}p{(\columnwidth - 4\tabcolsep) * \real{0.18}}
  >{\raggedright\arraybackslash}p{(\columnwidth - 4\tabcolsep) * \real{0.20}}
  >{\raggedright\arraybackslash}p{(\columnwidth - 4\tabcolsep) * \real{0.62}}@{}}
\toprule
位置 & 类型 & 含义 \\
\midrule
\endhead
返回值 & \passthrough{\lstinline!bool!} & 返回
\passthrough{\lstinline!bool!}。更精确的业务含义见该方法用途、C++
签名和目标型号协议。 \\
\bottomrule
\end{longtable}

\textbf{对应 C++}

\begin{itemize}
\tightlist
\item
  类：\passthrough{\lstinline!sensor\_msgs::msg::dds\_::PointField\_!}
\item
  签名：\passthrough{\lstinline"operator!=(const value \&) const"}
\item
  绑定策略：\passthrough{\lstinline!未单独标注!}
\item
  声明位置：由 pybind11 手工绑定或生成注册表提供；manifest
  未记录独立头文件行号。
\end{itemize}

\textbf{说明}

\begin{itemize}
\tightlist
\item
  示例是局部调用片段；\passthrough{\lstinline!obj!}
  和参数变量表示已经按上层业务校验并准备好的对象。
\end{itemize}

\textbf{用法}

\begin{lstlisting}[language=Python]
different = left != right
\end{lstlisting}

\hypertarget{unitree-sdk2-cpp-idl-ros2-pointfield-name}{%
\paragraph{\texorpdfstring{\texttt{unitree\_sdk2\_cpp.idl.ros2.PointField.name}}{unitree\_sdk2\_cpp.idl.ros2.PointField.name}}\label{unitree-sdk2-cpp-idl-ros2-pointfield-name}}

底层 \passthrough{\lstinline!sensor\_msgs::msg::dds\_::PointField\_!} 的
\passthrough{\lstinline!name!} 字段。类型签名只说明 Python
表示；单位、数值范围、数组固定长度和枚举含义需查对应 IDL 头文件。
底层为字符串；长度、编码和允许值由具体协议决定。

\textbf{签名}

\begin{lstlisting}[language=Python]
@property
def name(self) -> str

@name.setter
def name(self, value: str) -> None
\end{lstlisting}

\textbf{可用性与安全性}

\textbf{\passthrough{\lstinline!AVAILABLE!} /
\passthrough{\lstinline!VALUE\_TYPE!}}：当前绑定源码已实现该消息值操作。

\textbf{参数}

\begin{longtable}[]{@{}
  >{\raggedright\arraybackslash}p{(\columnwidth - 4\tabcolsep) * \real{0.18}}
  >{\raggedright\arraybackslash}p{(\columnwidth - 4\tabcolsep) * \real{0.20}}
  >{\raggedright\arraybackslash}p{(\columnwidth - 4\tabcolsep) * \real{0.62}}@{}}
\toprule
名称 & Python 类型 & 含义 \\
\midrule
\endhead
\passthrough{\lstinline!value!} & \passthrough{\lstinline!str!} &
要写入或传递的新值，必须符合该参数的 Python 类型以及底层 C++
范围约束。 \\
\bottomrule
\end{longtable}

\textbf{返回值}

\begin{longtable}[]{@{}
  >{\raggedright\arraybackslash}p{(\columnwidth - 4\tabcolsep) * \real{0.18}}
  >{\raggedright\arraybackslash}p{(\columnwidth - 4\tabcolsep) * \real{0.20}}
  >{\raggedright\arraybackslash}p{(\columnwidth - 4\tabcolsep) * \real{0.62}}@{}}
\toprule
位置 & 类型 & 含义 \\
\midrule
\endhead
getter 返回值 & \passthrough{\lstinline!str!} & 返回属性当前值。 \\
\bottomrule
\end{longtable}

\textbf{对应 C++}

\begin{itemize}
\tightlist
\item
  类：\passthrough{\lstinline!sensor\_msgs::msg::dds\_::PointField\_!}
\item
  字段类型：\passthrough{\lstinline!std::string!}
\item
  声明位置：\passthrough{\lstinline!include/unitree/idl/ros2/PointField\_.hpp:44!}
\end{itemize}

\textbf{用法}

\begin{lstlisting}[language=Python]
current_value = obj.name
obj.name = new_value
\end{lstlisting}

\hypertarget{unitree-sdk2-cpp-idl-ros2-pointfield-offset}{%
\paragraph{\texorpdfstring{\texttt{unitree\_sdk2\_cpp.idl.ros2.PointField.offset}}{unitree\_sdk2\_cpp.idl.ros2.PointField.offset}}\label{unitree-sdk2-cpp-idl-ros2-pointfield-offset}}

底层 \passthrough{\lstinline!sensor\_msgs::msg::dds\_::PointField\_!} 的
\passthrough{\lstinline!offset!} 字段。类型签名只说明 Python
表示；单位、数值范围、数组固定长度和枚举含义需查对应 IDL 头文件。
取值范围为 0 到 4294967295。

\textbf{签名}

\begin{lstlisting}[language=Python]
@property
def offset(self) -> int

@offset.setter
def offset(self, value: int) -> None
\end{lstlisting}

\textbf{可用性与安全性}

\textbf{\passthrough{\lstinline!AVAILABLE!} /
\passthrough{\lstinline!VALUE\_TYPE!}}：当前绑定源码已实现该消息值操作。

\textbf{参数}

\begin{longtable}[]{@{}
  >{\raggedright\arraybackslash}p{(\columnwidth - 4\tabcolsep) * \real{0.18}}
  >{\raggedright\arraybackslash}p{(\columnwidth - 4\tabcolsep) * \real{0.20}}
  >{\raggedright\arraybackslash}p{(\columnwidth - 4\tabcolsep) * \real{0.62}}@{}}
\toprule
名称 & Python 类型 & 含义 \\
\midrule
\endhead
\passthrough{\lstinline!value!} & \passthrough{\lstinline!int!} &
要写入或传递的新值，必须符合该参数的 Python 类型以及底层 C++
范围约束。 \\
\bottomrule
\end{longtable}

\textbf{返回值}

\begin{longtable}[]{@{}
  >{\raggedright\arraybackslash}p{(\columnwidth - 4\tabcolsep) * \real{0.18}}
  >{\raggedright\arraybackslash}p{(\columnwidth - 4\tabcolsep) * \real{0.20}}
  >{\raggedright\arraybackslash}p{(\columnwidth - 4\tabcolsep) * \real{0.62}}@{}}
\toprule
位置 & 类型 & 含义 \\
\midrule
\endhead
getter 返回值 & \passthrough{\lstinline!int!} & 返回属性当前值。 \\
\bottomrule
\end{longtable}

\textbf{对应 C++}

\begin{itemize}
\tightlist
\item
  类：\passthrough{\lstinline!sensor\_msgs::msg::dds\_::PointField\_!}
\item
  字段类型：\passthrough{\lstinline!uint32\_t!}
\item
  声明位置：\passthrough{\lstinline!include/unitree/idl/ros2/PointField\_.hpp:45!}
\end{itemize}

\textbf{用法}

\begin{lstlisting}[language=Python]
current_value = obj.offset
obj.offset = new_value
\end{lstlisting}

\hypertarget{unitree-sdk2-cpp-idl-ros2-pointfield-datatype}{%
\paragraph{\texorpdfstring{\texttt{unitree\_sdk2\_cpp.idl.ros2.PointField.datatype}}{unitree\_sdk2\_cpp.idl.ros2.PointField.datatype}}\label{unitree-sdk2-cpp-idl-ros2-pointfield-datatype}}

底层 \passthrough{\lstinline!sensor\_msgs::msg::dds\_::PointField\_!} 的
\passthrough{\lstinline!datatype!} 字段。类型签名只说明 Python
表示；单位、数值范围、数组固定长度和枚举含义需查对应 IDL 头文件。
取值范围为 0 到 255。

\textbf{签名}

\begin{lstlisting}[language=Python]
@property
def datatype(self) -> int

@datatype.setter
def datatype(self, value: int) -> None
\end{lstlisting}

\textbf{可用性与安全性}

\textbf{\passthrough{\lstinline!AVAILABLE!} /
\passthrough{\lstinline!VALUE\_TYPE!}}：当前绑定源码已实现该消息值操作。

\textbf{参数}

\begin{longtable}[]{@{}
  >{\raggedright\arraybackslash}p{(\columnwidth - 4\tabcolsep) * \real{0.18}}
  >{\raggedright\arraybackslash}p{(\columnwidth - 4\tabcolsep) * \real{0.20}}
  >{\raggedright\arraybackslash}p{(\columnwidth - 4\tabcolsep) * \real{0.62}}@{}}
\toprule
名称 & Python 类型 & 含义 \\
\midrule
\endhead
\passthrough{\lstinline!value!} & \passthrough{\lstinline!int!} &
要写入或传递的新值，必须符合该参数的 Python 类型以及底层 C++
范围约束。 \\
\bottomrule
\end{longtable}

\textbf{返回值}

\begin{longtable}[]{@{}
  >{\raggedright\arraybackslash}p{(\columnwidth - 4\tabcolsep) * \real{0.18}}
  >{\raggedright\arraybackslash}p{(\columnwidth - 4\tabcolsep) * \real{0.20}}
  >{\raggedright\arraybackslash}p{(\columnwidth - 4\tabcolsep) * \real{0.62}}@{}}
\toprule
位置 & 类型 & 含义 \\
\midrule
\endhead
getter 返回值 & \passthrough{\lstinline!int!} & 返回属性当前值。 \\
\bottomrule
\end{longtable}

\textbf{对应 C++}

\begin{itemize}
\tightlist
\item
  类：\passthrough{\lstinline!sensor\_msgs::msg::dds\_::PointField\_!}
\item
  字段类型：\passthrough{\lstinline!uint8\_t!}
\item
  声明位置：\passthrough{\lstinline!include/unitree/idl/ros2/PointField\_.hpp:46!}
\end{itemize}

\textbf{用法}

\begin{lstlisting}[language=Python]
current_value = obj.datatype
obj.datatype = new_value
\end{lstlisting}

\hypertarget{unitree-sdk2-cpp-idl-ros2-pointfield-count}{%
\paragraph{\texorpdfstring{\texttt{unitree\_sdk2\_cpp.idl.ros2.PointField.count}}{unitree\_sdk2\_cpp.idl.ros2.PointField.count}}\label{unitree-sdk2-cpp-idl-ros2-pointfield-count}}

底层 \passthrough{\lstinline!sensor\_msgs::msg::dds\_::PointField\_!} 的
\passthrough{\lstinline!count!} 字段。类型签名只说明 Python
表示；单位、数值范围、数组固定长度和枚举含义需查对应 IDL 头文件。
取值范围为 0 到 4294967295。

\textbf{签名}

\begin{lstlisting}[language=Python]
@property
def count(self) -> int

@count.setter
def count(self, value: int) -> None
\end{lstlisting}

\textbf{可用性与安全性}

\textbf{\passthrough{\lstinline!AVAILABLE!} /
\passthrough{\lstinline!VALUE\_TYPE!}}：当前绑定源码已实现该消息值操作。

\textbf{参数}

\begin{longtable}[]{@{}
  >{\raggedright\arraybackslash}p{(\columnwidth - 4\tabcolsep) * \real{0.18}}
  >{\raggedright\arraybackslash}p{(\columnwidth - 4\tabcolsep) * \real{0.20}}
  >{\raggedright\arraybackslash}p{(\columnwidth - 4\tabcolsep) * \real{0.62}}@{}}
\toprule
名称 & Python 类型 & 含义 \\
\midrule
\endhead
\passthrough{\lstinline!value!} & \passthrough{\lstinline!int!} &
要写入或传递的新值，必须符合该参数的 Python 类型以及底层 C++
范围约束。 \\
\bottomrule
\end{longtable}

\textbf{返回值}

\begin{longtable}[]{@{}
  >{\raggedright\arraybackslash}p{(\columnwidth - 4\tabcolsep) * \real{0.18}}
  >{\raggedright\arraybackslash}p{(\columnwidth - 4\tabcolsep) * \real{0.20}}
  >{\raggedright\arraybackslash}p{(\columnwidth - 4\tabcolsep) * \real{0.62}}@{}}
\toprule
位置 & 类型 & 含义 \\
\midrule
\endhead
getter 返回值 & \passthrough{\lstinline!int!} & 返回属性当前值。 \\
\bottomrule
\end{longtable}

\textbf{对应 C++}

\begin{itemize}
\tightlist
\item
  类：\passthrough{\lstinline!sensor\_msgs::msg::dds\_::PointField\_!}
\item
  字段类型：\passthrough{\lstinline!uint32\_t!}
\item
  声明位置：\passthrough{\lstinline!include/unitree/idl/ros2/PointField\_.hpp:47!}
\end{itemize}

\textbf{用法}

\begin{lstlisting}[language=Python]
current_value = obj.count
obj.count = new_value
\end{lstlisting}

\hypertarget{unitree-sdk2-cpp-idl-ros2-pointcloud2}{%
\subsubsection{\texorpdfstring{\texttt{unitree\_sdk2\_cpp.idl.ros2.PointCloud2}}{unitree\_sdk2\_cpp.idl.ros2.PointCloud2}}\label{unitree-sdk2-cpp-idl-ros2-pointcloud2}}

可由 Python 读写的 DDS/IDL 消息值类型。字段采用复制语义。

\textbf{导入}

\begin{lstlisting}[language=Python]
from unitree_sdk2_cpp.idl.ros2 import PointCloud2
\end{lstlisting}

\textbf{方法索引}

\begin{longtable}[]{@{}
  >{\raggedright\arraybackslash}p{(\columnwidth - 6\tabcolsep) * \real{0.18}}
  >{\raggedright\arraybackslash}p{(\columnwidth - 6\tabcolsep) * \real{0.24}}
  >{\raggedright\arraybackslash}p{(\columnwidth - 6\tabcolsep) * \real{0.18}}
  >{\raggedright\arraybackslash}p{(\columnwidth - 6\tabcolsep) * \real{0.40}}@{}}
\toprule
方法 & Python 签名 & 状态 & 安全分类 \\
\midrule
\endhead
\protect\hyperlink{unitree-sdk2-cpp-idl-ros2-pointcloud2-dunder-init-1}{\passthrough{\lstinline!\_\_init\_\_!}}
& \passthrough{\lstinline!def \_\_init\_\_(self) -> None!} &
\passthrough{\lstinline!AVAILABLE!} &
\passthrough{\lstinline!VALUE\_TYPE!} \\
\protect\hyperlink{unitree-sdk2-cpp-idl-ros2-pointcloud2-dunder-eq-1}{\passthrough{\lstinline!\_\_eq\_\_!}}
& \passthrough{\lstinline!def \_\_eq\_\_(self, other: object) -> bool!}
& \passthrough{\lstinline!AVAILABLE!} &
\passthrough{\lstinline!VALUE\_TYPE!} \\
\protect\hyperlink{unitree-sdk2-cpp-idl-ros2-pointcloud2-dunder-ne-1}{\passthrough{\lstinline!\_\_ne\_\_!}}
& \passthrough{\lstinline!def \_\_ne\_\_(self, other: object) -> bool!}
& \passthrough{\lstinline!AVAILABLE!} &
\passthrough{\lstinline!VALUE\_TYPE!} \\
\bottomrule
\end{longtable}

\textbf{属性索引}

\begin{longtable}[]{@{}
  >{\raggedright\arraybackslash}p{(\columnwidth - 6\tabcolsep) * \real{0.18}}
  >{\raggedright\arraybackslash}p{(\columnwidth - 6\tabcolsep) * \real{0.24}}
  >{\raggedright\arraybackslash}p{(\columnwidth - 6\tabcolsep) * \real{0.18}}
  >{\raggedright\arraybackslash}p{(\columnwidth - 6\tabcolsep) * \real{0.40}}@{}}
\toprule
属性 & 读取类型 & 写入类型 & 状态 \\
\midrule
\endhead
\protect\hyperlink{unitree-sdk2-cpp-idl-ros2-pointcloud2-header}{\passthrough{\lstinline!header!}}
& \passthrough{\lstinline!Any!} & \passthrough{\lstinline!Any!} &
\passthrough{\lstinline!AVAILABLE!} \\
\protect\hyperlink{unitree-sdk2-cpp-idl-ros2-pointcloud2-height}{\passthrough{\lstinline!height!}}
& \passthrough{\lstinline!int!} & \passthrough{\lstinline!int!} &
\passthrough{\lstinline!AVAILABLE!} \\
\protect\hyperlink{unitree-sdk2-cpp-idl-ros2-pointcloud2-width}{\passthrough{\lstinline!width!}}
& \passthrough{\lstinline!int!} & \passthrough{\lstinline!int!} &
\passthrough{\lstinline!AVAILABLE!} \\
\protect\hyperlink{unitree-sdk2-cpp-idl-ros2-pointcloud2-fields}{\passthrough{\lstinline!fields!}}
& \passthrough{\lstinline!list[PointField]!} &
\passthrough{\lstinline!Sequence[PointField]!} &
\passthrough{\lstinline!AVAILABLE!} \\
\protect\hyperlink{unitree-sdk2-cpp-idl-ros2-pointcloud2-is-bigendian}{\passthrough{\lstinline!is\_bigendian!}}
& \passthrough{\lstinline!bool!} & \passthrough{\lstinline!bool!} &
\passthrough{\lstinline!AVAILABLE!} \\
\protect\hyperlink{unitree-sdk2-cpp-idl-ros2-pointcloud2-point-step}{\passthrough{\lstinline!point\_step!}}
& \passthrough{\lstinline!int!} & \passthrough{\lstinline!int!} &
\passthrough{\lstinline!AVAILABLE!} \\
\protect\hyperlink{unitree-sdk2-cpp-idl-ros2-pointcloud2-row-step}{\passthrough{\lstinline!row\_step!}}
& \passthrough{\lstinline!int!} & \passthrough{\lstinline!int!} &
\passthrough{\lstinline!AVAILABLE!} \\
\protect\hyperlink{unitree-sdk2-cpp-idl-ros2-pointcloud2-data}{\passthrough{\lstinline!data!}}
& \passthrough{\lstinline!list[int]!} &
\passthrough{\lstinline!Sequence[int]!} &
\passthrough{\lstinline!AVAILABLE!} \\
\protect\hyperlink{unitree-sdk2-cpp-idl-ros2-pointcloud2-is-dense}{\passthrough{\lstinline!is\_dense!}}
& \passthrough{\lstinline!bool!} & \passthrough{\lstinline!bool!} &
\passthrough{\lstinline!AVAILABLE!} \\
\bottomrule
\end{longtable}

\hypertarget{unitree-sdk2-cpp-idl-ros2-pointcloud2-dunder-init-1}{%
\paragraph{\texorpdfstring{\texttt{unitree\_sdk2\_cpp.idl.ros2.PointCloud2.\_\_init\_\_}}{unitree\_sdk2\_cpp.idl.ros2.PointCloud2.\_\_init\_\_}}\label{unitree-sdk2-cpp-idl-ros2-pointcloud2-dunder-init-1}}

初始化 \passthrough{\lstinline!PointCloud2!}
实例。是否能实际构造取决于下方可用性状态。

\textbf{签名}

\begin{lstlisting}[language=Python]
def __init__(self) -> None
\end{lstlisting}

\textbf{可用性与安全性}

\textbf{\passthrough{\lstinline!AVAILABLE!} /
\passthrough{\lstinline!VALUE\_TYPE!}}：当前绑定源码已实现该消息值操作。

\textbf{参数}

无显式参数。实例方法中的 \passthrough{\lstinline!self!} 由 Python
自动传入。

\textbf{返回值}

\begin{longtable}[]{@{}
  >{\raggedright\arraybackslash}p{(\columnwidth - 4\tabcolsep) * \real{0.18}}
  >{\raggedright\arraybackslash}p{(\columnwidth - 4\tabcolsep) * \real{0.20}}
  >{\raggedright\arraybackslash}p{(\columnwidth - 4\tabcolsep) * \real{0.62}}@{}}
\toprule
位置 & 类型 & 含义 \\
\midrule
\endhead
无 & \passthrough{\lstinline!None!} &
\passthrough{\lstinline!\_\_init\_\_!}
本身不返回值；调用类对象时，在构造函数可用的前提下得到该类实例。 \\
\bottomrule
\end{longtable}

\textbf{对应 C++}

\begin{itemize}
\tightlist
\item
  类：\passthrough{\lstinline!sensor\_msgs::msg::dds\_::PointCloud2\_!}
\item
  签名：\passthrough{\lstinline!PointCloud2\_()!}
\item
  绑定策略：\passthrough{\lstinline!未单独标注!}
\item
  声明位置：\passthrough{\lstinline!include/unitree/idl/ros2/PointCloud2\_.hpp:39!}
\end{itemize}

\textbf{说明}

\begin{itemize}
\tightlist
\item
  示例是局部调用片段；\passthrough{\lstinline!obj!}
  和参数变量表示已经按上层业务校验并准备好的对象。
\end{itemize}

\textbf{用法}

\begin{lstlisting}[language=Python]
value = PointCloud2()
\end{lstlisting}

\hypertarget{unitree-sdk2-cpp-idl-ros2-pointcloud2-dunder-eq-1}{%
\paragraph{\texorpdfstring{\texttt{unitree\_sdk2\_cpp.idl.ros2.PointCloud2.\_\_eq\_\_}}{unitree\_sdk2\_cpp.idl.ros2.PointCloud2.\_\_eq\_\_}}\label{unitree-sdk2-cpp-idl-ros2-pointcloud2-dunder-eq-1}}

按底层 IDL 消息内容比较两个对象是否相等。

\textbf{签名}

\begin{lstlisting}[language=Python]
def __eq__(self, other: object) -> bool
\end{lstlisting}

\textbf{可用性与安全性}

\textbf{\passthrough{\lstinline!AVAILABLE!} /
\passthrough{\lstinline!VALUE\_TYPE!}}：当前绑定源码已实现该消息值操作。

\textbf{参数}

\begin{longtable}[]{@{}
  >{\raggedright\arraybackslash}p{(\columnwidth - 6\tabcolsep) * \real{0.18}}
  >{\raggedright\arraybackslash}p{(\columnwidth - 6\tabcolsep) * \real{0.24}}
  >{\raggedright\arraybackslash}p{(\columnwidth - 6\tabcolsep) * \real{0.18}}
  >{\raggedright\arraybackslash}p{(\columnwidth - 6\tabcolsep) * \real{0.40}}@{}}
\toprule
名称 & Python 类型 & 默认值 & 含义 \\
\midrule
\endhead
\passthrough{\lstinline!other!} & \passthrough{\lstinline!object!} &
必填 & 用于比较的另一个 Python 对象。类型不兼容时比较结果通常为
\passthrough{\lstinline!False!}。 \\
\bottomrule
\end{longtable}

\textbf{返回值}

\begin{longtable}[]{@{}
  >{\raggedright\arraybackslash}p{(\columnwidth - 4\tabcolsep) * \real{0.18}}
  >{\raggedright\arraybackslash}p{(\columnwidth - 4\tabcolsep) * \real{0.20}}
  >{\raggedright\arraybackslash}p{(\columnwidth - 4\tabcolsep) * \real{0.62}}@{}}
\toprule
位置 & 类型 & 含义 \\
\midrule
\endhead
返回值 & \passthrough{\lstinline!bool!} & 返回
\passthrough{\lstinline!bool!}。更精确的业务含义见该方法用途、C++
签名和目标型号协议。 \\
\bottomrule
\end{longtable}

\textbf{对应 C++}

\begin{itemize}
\tightlist
\item
  类：\passthrough{\lstinline!sensor\_msgs::msg::dds\_::PointCloud2\_!}
\item
  签名：\passthrough{\lstinline!operator==(const value \&) const!}
\item
  绑定策略：\passthrough{\lstinline!未单独标注!}
\item
  声明位置：由 pybind11 手工绑定或生成注册表提供；manifest
  未记录独立头文件行号。
\end{itemize}

\textbf{说明}

\begin{itemize}
\tightlist
\item
  示例是局部调用片段；\passthrough{\lstinline!obj!}
  和参数变量表示已经按上层业务校验并准备好的对象。
\end{itemize}

\textbf{用法}

\begin{lstlisting}[language=Python]
same = left == right
\end{lstlisting}

\hypertarget{unitree-sdk2-cpp-idl-ros2-pointcloud2-dunder-ne-1}{%
\paragraph{\texorpdfstring{\texttt{unitree\_sdk2\_cpp.idl.ros2.PointCloud2.\_\_ne\_\_}}{unitree\_sdk2\_cpp.idl.ros2.PointCloud2.\_\_ne\_\_}}\label{unitree-sdk2-cpp-idl-ros2-pointcloud2-dunder-ne-1}}

按底层 IDL 消息内容比较两个对象是否不相等。

\textbf{签名}

\begin{lstlisting}[language=Python]
def __ne__(self, other: object) -> bool
\end{lstlisting}

\textbf{可用性与安全性}

\textbf{\passthrough{\lstinline!AVAILABLE!} /
\passthrough{\lstinline!VALUE\_TYPE!}}：当前绑定源码已实现该消息值操作。

\textbf{参数}

\begin{longtable}[]{@{}
  >{\raggedright\arraybackslash}p{(\columnwidth - 6\tabcolsep) * \real{0.18}}
  >{\raggedright\arraybackslash}p{(\columnwidth - 6\tabcolsep) * \real{0.24}}
  >{\raggedright\arraybackslash}p{(\columnwidth - 6\tabcolsep) * \real{0.18}}
  >{\raggedright\arraybackslash}p{(\columnwidth - 6\tabcolsep) * \real{0.40}}@{}}
\toprule
名称 & Python 类型 & 默认值 & 含义 \\
\midrule
\endhead
\passthrough{\lstinline!other!} & \passthrough{\lstinline!object!} &
必填 & 用于比较的另一个 Python 对象。类型不兼容时比较结果通常为
\passthrough{\lstinline!False!}。 \\
\bottomrule
\end{longtable}

\textbf{返回值}

\begin{longtable}[]{@{}
  >{\raggedright\arraybackslash}p{(\columnwidth - 4\tabcolsep) * \real{0.18}}
  >{\raggedright\arraybackslash}p{(\columnwidth - 4\tabcolsep) * \real{0.20}}
  >{\raggedright\arraybackslash}p{(\columnwidth - 4\tabcolsep) * \real{0.62}}@{}}
\toprule
位置 & 类型 & 含义 \\
\midrule
\endhead
返回值 & \passthrough{\lstinline!bool!} & 返回
\passthrough{\lstinline!bool!}。更精确的业务含义见该方法用途、C++
签名和目标型号协议。 \\
\bottomrule
\end{longtable}

\textbf{对应 C++}

\begin{itemize}
\tightlist
\item
  类：\passthrough{\lstinline!sensor\_msgs::msg::dds\_::PointCloud2\_!}
\item
  签名：\passthrough{\lstinline"operator!=(const value \&) const"}
\item
  绑定策略：\passthrough{\lstinline!未单独标注!}
\item
  声明位置：由 pybind11 手工绑定或生成注册表提供；manifest
  未记录独立头文件行号。
\end{itemize}

\textbf{说明}

\begin{itemize}
\tightlist
\item
  示例是局部调用片段；\passthrough{\lstinline!obj!}
  和参数变量表示已经按上层业务校验并准备好的对象。
\end{itemize}

\textbf{用法}

\begin{lstlisting}[language=Python]
different = left != right
\end{lstlisting}

\hypertarget{unitree-sdk2-cpp-idl-ros2-pointcloud2-header}{%
\paragraph{\texorpdfstring{\texttt{unitree\_sdk2\_cpp.idl.ros2.PointCloud2.header}}{unitree\_sdk2\_cpp.idl.ros2.PointCloud2.header}}\label{unitree-sdk2-cpp-idl-ros2-pointcloud2-header}}

底层 \passthrough{\lstinline!sensor\_msgs::msg::dds\_::PointCloud2\_!}
的 \passthrough{\lstinline!header!} 字段。类型签名只说明 Python
表示；单位、数值范围、数组固定长度和枚举含义需查对应 IDL 头文件。

\textbf{签名}

\begin{lstlisting}[language=Python]
@property
def header(self) -> Any

@header.setter
def header(self, value: Any) -> None
\end{lstlisting}

\textbf{可用性与安全性}

\textbf{\passthrough{\lstinline!AVAILABLE!} /
\passthrough{\lstinline!VALUE\_TYPE!}}：当前绑定源码已实现该消息值操作。

\textbf{参数}

\begin{longtable}[]{@{}
  >{\raggedright\arraybackslash}p{(\columnwidth - 4\tabcolsep) * \real{0.18}}
  >{\raggedright\arraybackslash}p{(\columnwidth - 4\tabcolsep) * \real{0.20}}
  >{\raggedright\arraybackslash}p{(\columnwidth - 4\tabcolsep) * \real{0.62}}@{}}
\toprule
名称 & Python 类型 & 含义 \\
\midrule
\endhead
\passthrough{\lstinline!value!} & \passthrough{\lstinline!Any!} &
要写入或传递的新值，必须符合该参数的 Python 类型以及底层 C++
范围约束。 \\
\bottomrule
\end{longtable}

\textbf{返回值}

\begin{longtable}[]{@{}
  >{\raggedright\arraybackslash}p{(\columnwidth - 4\tabcolsep) * \real{0.18}}
  >{\raggedright\arraybackslash}p{(\columnwidth - 4\tabcolsep) * \real{0.20}}
  >{\raggedright\arraybackslash}p{(\columnwidth - 4\tabcolsep) * \real{0.62}}@{}}
\toprule
位置 & 类型 & 含义 \\
\midrule
\endhead
getter 返回值 & \passthrough{\lstinline!Any!} & 返回属性当前值。 \\
\bottomrule
\end{longtable}

\textbf{对应 C++}

\begin{itemize}
\tightlist
\item
  类：\passthrough{\lstinline!sensor\_msgs::msg::dds\_::PointCloud2\_!}
\item
  字段类型：\passthrough{\lstinline!::std\_msgs::msg::dds\_::Header\_!}
\item
  声明位置：\passthrough{\lstinline!include/unitree/idl/ros2/PointCloud2\_.hpp:28!}
\end{itemize}

\textbf{用法}

\begin{lstlisting}[language=Python]
current_value = obj.header
obj.header = new_value
\end{lstlisting}

\hypertarget{unitree-sdk2-cpp-idl-ros2-pointcloud2-height}{%
\paragraph{\texorpdfstring{\texttt{unitree\_sdk2\_cpp.idl.ros2.PointCloud2.height}}{unitree\_sdk2\_cpp.idl.ros2.PointCloud2.height}}\label{unitree-sdk2-cpp-idl-ros2-pointcloud2-height}}

底层 \passthrough{\lstinline!sensor\_msgs::msg::dds\_::PointCloud2\_!}
的 \passthrough{\lstinline!height!} 字段。类型签名只说明 Python
表示；单位、数值范围、数组固定长度和枚举含义需查对应 IDL 头文件。
取值范围为 0 到 4294967295。

\textbf{签名}

\begin{lstlisting}[language=Python]
@property
def height(self) -> int

@height.setter
def height(self, value: int) -> None
\end{lstlisting}

\textbf{可用性与安全性}

\textbf{\passthrough{\lstinline!AVAILABLE!} /
\passthrough{\lstinline!VALUE\_TYPE!}}：当前绑定源码已实现该消息值操作。

\textbf{参数}

\begin{longtable}[]{@{}
  >{\raggedright\arraybackslash}p{(\columnwidth - 4\tabcolsep) * \real{0.18}}
  >{\raggedright\arraybackslash}p{(\columnwidth - 4\tabcolsep) * \real{0.20}}
  >{\raggedright\arraybackslash}p{(\columnwidth - 4\tabcolsep) * \real{0.62}}@{}}
\toprule
名称 & Python 类型 & 含义 \\
\midrule
\endhead
\passthrough{\lstinline!value!} & \passthrough{\lstinline!int!} &
要写入或传递的新值，必须符合该参数的 Python 类型以及底层 C++
范围约束。 \\
\bottomrule
\end{longtable}

\textbf{返回值}

\begin{longtable}[]{@{}
  >{\raggedright\arraybackslash}p{(\columnwidth - 4\tabcolsep) * \real{0.18}}
  >{\raggedright\arraybackslash}p{(\columnwidth - 4\tabcolsep) * \real{0.20}}
  >{\raggedright\arraybackslash}p{(\columnwidth - 4\tabcolsep) * \real{0.62}}@{}}
\toprule
位置 & 类型 & 含义 \\
\midrule
\endhead
getter 返回值 & \passthrough{\lstinline!int!} & 返回属性当前值。 \\
\bottomrule
\end{longtable}

\textbf{对应 C++}

\begin{itemize}
\tightlist
\item
  类：\passthrough{\lstinline!sensor\_msgs::msg::dds\_::PointCloud2\_!}
\item
  字段类型：\passthrough{\lstinline!uint32\_t!}
\item
  声明位置：\passthrough{\lstinline!include/unitree/idl/ros2/PointCloud2\_.hpp:29!}
\end{itemize}

\textbf{用法}

\begin{lstlisting}[language=Python]
current_value = obj.height
obj.height = new_value
\end{lstlisting}

\hypertarget{unitree-sdk2-cpp-idl-ros2-pointcloud2-width}{%
\paragraph{\texorpdfstring{\texttt{unitree\_sdk2\_cpp.idl.ros2.PointCloud2.width}}{unitree\_sdk2\_cpp.idl.ros2.PointCloud2.width}}\label{unitree-sdk2-cpp-idl-ros2-pointcloud2-width}}

底层 \passthrough{\lstinline!sensor\_msgs::msg::dds\_::PointCloud2\_!}
的 \passthrough{\lstinline!width!} 字段。类型签名只说明 Python
表示；单位、数值范围、数组固定长度和枚举含义需查对应 IDL 头文件。
取值范围为 0 到 4294967295。

\textbf{签名}

\begin{lstlisting}[language=Python]
@property
def width(self) -> int

@width.setter
def width(self, value: int) -> None
\end{lstlisting}

\textbf{可用性与安全性}

\textbf{\passthrough{\lstinline!AVAILABLE!} /
\passthrough{\lstinline!VALUE\_TYPE!}}：当前绑定源码已实现该消息值操作。

\textbf{参数}

\begin{longtable}[]{@{}
  >{\raggedright\arraybackslash}p{(\columnwidth - 4\tabcolsep) * \real{0.18}}
  >{\raggedright\arraybackslash}p{(\columnwidth - 4\tabcolsep) * \real{0.20}}
  >{\raggedright\arraybackslash}p{(\columnwidth - 4\tabcolsep) * \real{0.62}}@{}}
\toprule
名称 & Python 类型 & 含义 \\
\midrule
\endhead
\passthrough{\lstinline!value!} & \passthrough{\lstinline!int!} &
要写入或传递的新值，必须符合该参数的 Python 类型以及底层 C++
范围约束。 \\
\bottomrule
\end{longtable}

\textbf{返回值}

\begin{longtable}[]{@{}
  >{\raggedright\arraybackslash}p{(\columnwidth - 4\tabcolsep) * \real{0.18}}
  >{\raggedright\arraybackslash}p{(\columnwidth - 4\tabcolsep) * \real{0.20}}
  >{\raggedright\arraybackslash}p{(\columnwidth - 4\tabcolsep) * \real{0.62}}@{}}
\toprule
位置 & 类型 & 含义 \\
\midrule
\endhead
getter 返回值 & \passthrough{\lstinline!int!} & 返回属性当前值。 \\
\bottomrule
\end{longtable}

\textbf{对应 C++}

\begin{itemize}
\tightlist
\item
  类：\passthrough{\lstinline!sensor\_msgs::msg::dds\_::PointCloud2\_!}
\item
  字段类型：\passthrough{\lstinline!uint32\_t!}
\item
  声明位置：\passthrough{\lstinline!include/unitree/idl/ros2/PointCloud2\_.hpp:30!}
\end{itemize}

\textbf{用法}

\begin{lstlisting}[language=Python]
current_value = obj.width
obj.width = new_value
\end{lstlisting}

\hypertarget{unitree-sdk2-cpp-idl-ros2-pointcloud2-fields}{%
\paragraph{\texorpdfstring{\texttt{unitree\_sdk2\_cpp.idl.ros2.PointCloud2.fields}}{unitree\_sdk2\_cpp.idl.ros2.PointCloud2.fields}}\label{unitree-sdk2-cpp-idl-ros2-pointcloud2-fields}}

底层 \passthrough{\lstinline!sensor\_msgs::msg::dds\_::PointCloud2\_!}
的 \passthrough{\lstinline!fields!} 字段。类型签名只说明 Python
表示；单位、数值范围、数组固定长度和枚举含义需查对应 IDL 头文件。
底层是可变长度 vector

\textbf{签名}

\begin{lstlisting}[language=Python]
@property
def fields(self) -> list[PointField]

@fields.setter
def fields(self, value: Sequence[PointField]) -> None
\end{lstlisting}

\textbf{可用性与安全性}

\textbf{\passthrough{\lstinline!AVAILABLE!} /
\passthrough{\lstinline!VALUE\_TYPE!}}：当前绑定源码已实现该消息值操作。

\textbf{参数}

\begin{longtable}[]{@{}
  >{\raggedright\arraybackslash}p{(\columnwidth - 4\tabcolsep) * \real{0.18}}
  >{\raggedright\arraybackslash}p{(\columnwidth - 4\tabcolsep) * \real{0.20}}
  >{\raggedright\arraybackslash}p{(\columnwidth - 4\tabcolsep) * \real{0.62}}@{}}
\toprule
名称 & Python 类型 & 含义 \\
\midrule
\endhead
\passthrough{\lstinline!value!} &
\passthrough{\lstinline!Sequence[PointField]!} &
要写入或传递的新值，必须符合该参数的 Python 类型以及底层 C++
范围约束。 \\
\bottomrule
\end{longtable}

\textbf{返回值}

\begin{longtable}[]{@{}
  >{\raggedright\arraybackslash}p{(\columnwidth - 4\tabcolsep) * \real{0.18}}
  >{\raggedright\arraybackslash}p{(\columnwidth - 4\tabcolsep) * \real{0.20}}
  >{\raggedright\arraybackslash}p{(\columnwidth - 4\tabcolsep) * \real{0.62}}@{}}
\toprule
位置 & 类型 & 含义 \\
\midrule
\endhead
getter 返回值 & \passthrough{\lstinline!list[PointField]!} &
返回属性当前值。数组、vector 和嵌套消息采用复制语义。 \\
\bottomrule
\end{longtable}

\textbf{对应 C++}

\begin{itemize}
\tightlist
\item
  类：\passthrough{\lstinline!sensor\_msgs::msg::dds\_::PointCloud2\_!}
\item
  字段类型：\passthrough{\lstinline!std::vector< ::sensor\_msgs::msg::dds\_::PointField\_>!}
\item
  声明位置：\passthrough{\lstinline!include/unitree/idl/ros2/PointCloud2\_.hpp:31!}
\end{itemize}

\textbf{用法}

\begin{lstlisting}[language=Python]
current_value = obj.fields
obj.fields = new_value
\end{lstlisting}

\begin{quote}
{[}!TIP{]} 读取容器或嵌套值后应遵循``读取、修改、重新赋值''。只修改
getter 返回的副本不会可靠地更新原 C++ 消息。
\end{quote}

\hypertarget{unitree-sdk2-cpp-idl-ros2-pointcloud2-is-bigendian}{%
\paragraph{\texorpdfstring{\texttt{unitree\_sdk2\_cpp.idl.ros2.PointCloud2.is\_bigendian}}{unitree\_sdk2\_cpp.idl.ros2.PointCloud2.is\_bigendian}}\label{unitree-sdk2-cpp-idl-ros2-pointcloud2-is-bigendian}}

底层 \passthrough{\lstinline!sensor\_msgs::msg::dds\_::PointCloud2\_!}
的 \passthrough{\lstinline!is\_bigendian!} 字段。类型签名只说明 Python
表示；单位、数值范围、数组固定长度和枚举含义需查对应 IDL 头文件。
只接受布尔语义。

\textbf{签名}

\begin{lstlisting}[language=Python]
@property
def is_bigendian(self) -> bool

@is_bigendian.setter
def is_bigendian(self, value: bool) -> None
\end{lstlisting}

\textbf{可用性与安全性}

\textbf{\passthrough{\lstinline!AVAILABLE!} /
\passthrough{\lstinline!VALUE\_TYPE!}}：当前绑定源码已实现该消息值操作。

\textbf{参数}

\begin{longtable}[]{@{}
  >{\raggedright\arraybackslash}p{(\columnwidth - 4\tabcolsep) * \real{0.18}}
  >{\raggedright\arraybackslash}p{(\columnwidth - 4\tabcolsep) * \real{0.20}}
  >{\raggedright\arraybackslash}p{(\columnwidth - 4\tabcolsep) * \real{0.62}}@{}}
\toprule
名称 & Python 类型 & 含义 \\
\midrule
\endhead
\passthrough{\lstinline!value!} & \passthrough{\lstinline!bool!} &
要写入或传递的新值，必须符合该参数的 Python 类型以及底层 C++
范围约束。 \\
\bottomrule
\end{longtable}

\textbf{返回值}

\begin{longtable}[]{@{}
  >{\raggedright\arraybackslash}p{(\columnwidth - 4\tabcolsep) * \real{0.18}}
  >{\raggedright\arraybackslash}p{(\columnwidth - 4\tabcolsep) * \real{0.20}}
  >{\raggedright\arraybackslash}p{(\columnwidth - 4\tabcolsep) * \real{0.62}}@{}}
\toprule
位置 & 类型 & 含义 \\
\midrule
\endhead
getter 返回值 & \passthrough{\lstinline!bool!} & 返回属性当前值。 \\
\bottomrule
\end{longtable}

\textbf{对应 C++}

\begin{itemize}
\tightlist
\item
  类：\passthrough{\lstinline!sensor\_msgs::msg::dds\_::PointCloud2\_!}
\item
  字段类型：\passthrough{\lstinline!bool!}
\item
  声明位置：\passthrough{\lstinline!include/unitree/idl/ros2/PointCloud2\_.hpp:32!}
\end{itemize}

\textbf{用法}

\begin{lstlisting}[language=Python]
current_value = obj.is_bigendian
obj.is_bigendian = new_value
\end{lstlisting}

\hypertarget{unitree-sdk2-cpp-idl-ros2-pointcloud2-point-step}{%
\paragraph{\texorpdfstring{\texttt{unitree\_sdk2\_cpp.idl.ros2.PointCloud2.point\_step}}{unitree\_sdk2\_cpp.idl.ros2.PointCloud2.point\_step}}\label{unitree-sdk2-cpp-idl-ros2-pointcloud2-point-step}}

底层 \passthrough{\lstinline!sensor\_msgs::msg::dds\_::PointCloud2\_!}
的 \passthrough{\lstinline!point\_step!} 字段。类型签名只说明 Python
表示；单位、数值范围、数组固定长度和枚举含义需查对应 IDL 头文件。
取值范围为 0 到 4294967295。

\textbf{签名}

\begin{lstlisting}[language=Python]
@property
def point_step(self) -> int

@point_step.setter
def point_step(self, value: int) -> None
\end{lstlisting}

\textbf{可用性与安全性}

\textbf{\passthrough{\lstinline!AVAILABLE!} /
\passthrough{\lstinline!VALUE\_TYPE!}}：当前绑定源码已实现该消息值操作。

\textbf{参数}

\begin{longtable}[]{@{}
  >{\raggedright\arraybackslash}p{(\columnwidth - 4\tabcolsep) * \real{0.18}}
  >{\raggedright\arraybackslash}p{(\columnwidth - 4\tabcolsep) * \real{0.20}}
  >{\raggedright\arraybackslash}p{(\columnwidth - 4\tabcolsep) * \real{0.62}}@{}}
\toprule
名称 & Python 类型 & 含义 \\
\midrule
\endhead
\passthrough{\lstinline!value!} & \passthrough{\lstinline!int!} &
要写入或传递的新值，必须符合该参数的 Python 类型以及底层 C++
范围约束。 \\
\bottomrule
\end{longtable}

\textbf{返回值}

\begin{longtable}[]{@{}
  >{\raggedright\arraybackslash}p{(\columnwidth - 4\tabcolsep) * \real{0.18}}
  >{\raggedright\arraybackslash}p{(\columnwidth - 4\tabcolsep) * \real{0.20}}
  >{\raggedright\arraybackslash}p{(\columnwidth - 4\tabcolsep) * \real{0.62}}@{}}
\toprule
位置 & 类型 & 含义 \\
\midrule
\endhead
getter 返回值 & \passthrough{\lstinline!int!} & 返回属性当前值。 \\
\bottomrule
\end{longtable}

\textbf{对应 C++}

\begin{itemize}
\tightlist
\item
  类：\passthrough{\lstinline!sensor\_msgs::msg::dds\_::PointCloud2\_!}
\item
  字段类型：\passthrough{\lstinline!uint32\_t!}
\item
  声明位置：\passthrough{\lstinline!include/unitree/idl/ros2/PointCloud2\_.hpp:33!}
\end{itemize}

\textbf{用法}

\begin{lstlisting}[language=Python]
current_value = obj.point_step
obj.point_step = new_value
\end{lstlisting}

\hypertarget{unitree-sdk2-cpp-idl-ros2-pointcloud2-row-step}{%
\paragraph{\texorpdfstring{\texttt{unitree\_sdk2\_cpp.idl.ros2.PointCloud2.row\_step}}{unitree\_sdk2\_cpp.idl.ros2.PointCloud2.row\_step}}\label{unitree-sdk2-cpp-idl-ros2-pointcloud2-row-step}}

底层 \passthrough{\lstinline!sensor\_msgs::msg::dds\_::PointCloud2\_!}
的 \passthrough{\lstinline!row\_step!} 字段。类型签名只说明 Python
表示；单位、数值范围、数组固定长度和枚举含义需查对应 IDL 头文件。
取值范围为 0 到 4294967295。

\textbf{签名}

\begin{lstlisting}[language=Python]
@property
def row_step(self) -> int

@row_step.setter
def row_step(self, value: int) -> None
\end{lstlisting}

\textbf{可用性与安全性}

\textbf{\passthrough{\lstinline!AVAILABLE!} /
\passthrough{\lstinline!VALUE\_TYPE!}}：当前绑定源码已实现该消息值操作。

\textbf{参数}

\begin{longtable}[]{@{}
  >{\raggedright\arraybackslash}p{(\columnwidth - 4\tabcolsep) * \real{0.18}}
  >{\raggedright\arraybackslash}p{(\columnwidth - 4\tabcolsep) * \real{0.20}}
  >{\raggedright\arraybackslash}p{(\columnwidth - 4\tabcolsep) * \real{0.62}}@{}}
\toprule
名称 & Python 类型 & 含义 \\
\midrule
\endhead
\passthrough{\lstinline!value!} & \passthrough{\lstinline!int!} &
要写入或传递的新值，必须符合该参数的 Python 类型以及底层 C++
范围约束。 \\
\bottomrule
\end{longtable}

\textbf{返回值}

\begin{longtable}[]{@{}
  >{\raggedright\arraybackslash}p{(\columnwidth - 4\tabcolsep) * \real{0.18}}
  >{\raggedright\arraybackslash}p{(\columnwidth - 4\tabcolsep) * \real{0.20}}
  >{\raggedright\arraybackslash}p{(\columnwidth - 4\tabcolsep) * \real{0.62}}@{}}
\toprule
位置 & 类型 & 含义 \\
\midrule
\endhead
getter 返回值 & \passthrough{\lstinline!int!} & 返回属性当前值。 \\
\bottomrule
\end{longtable}

\textbf{对应 C++}

\begin{itemize}
\tightlist
\item
  类：\passthrough{\lstinline!sensor\_msgs::msg::dds\_::PointCloud2\_!}
\item
  字段类型：\passthrough{\lstinline!uint32\_t!}
\item
  声明位置：\passthrough{\lstinline!include/unitree/idl/ros2/PointCloud2\_.hpp:34!}
\end{itemize}

\textbf{用法}

\begin{lstlisting}[language=Python]
current_value = obj.row_step
obj.row_step = new_value
\end{lstlisting}

\hypertarget{unitree-sdk2-cpp-idl-ros2-pointcloud2-data}{%
\paragraph{\texorpdfstring{\texttt{unitree\_sdk2\_cpp.idl.ros2.PointCloud2.data}}{unitree\_sdk2\_cpp.idl.ros2.PointCloud2.data}}\label{unitree-sdk2-cpp-idl-ros2-pointcloud2-data}}

底层 \passthrough{\lstinline!sensor\_msgs::msg::dds\_::PointCloud2\_!}
的 \passthrough{\lstinline!data!} 字段。类型签名只说明 Python
表示；单位、数值范围、数组固定长度和枚举含义需查对应 IDL 头文件。
底层是可变长度 vector；元素约束：取值范围为 0 到 255。

\textbf{签名}

\begin{lstlisting}[language=Python]
@property
def data(self) -> list[int]

@data.setter
def data(self, value: Sequence[int]) -> None
\end{lstlisting}

\textbf{可用性与安全性}

\textbf{\passthrough{\lstinline!AVAILABLE!} /
\passthrough{\lstinline!VALUE\_TYPE!}}：当前绑定源码已实现该消息值操作。

\textbf{参数}

\begin{longtable}[]{@{}
  >{\raggedright\arraybackslash}p{(\columnwidth - 4\tabcolsep) * \real{0.18}}
  >{\raggedright\arraybackslash}p{(\columnwidth - 4\tabcolsep) * \real{0.20}}
  >{\raggedright\arraybackslash}p{(\columnwidth - 4\tabcolsep) * \real{0.62}}@{}}
\toprule
名称 & Python 类型 & 含义 \\
\midrule
\endhead
\passthrough{\lstinline!value!} &
\passthrough{\lstinline!Sequence[int]!} &
要写入或传递的新值，必须符合该参数的 Python 类型以及底层 C++
范围约束。 \\
\bottomrule
\end{longtable}

\textbf{返回值}

\begin{longtable}[]{@{}
  >{\raggedright\arraybackslash}p{(\columnwidth - 4\tabcolsep) * \real{0.18}}
  >{\raggedright\arraybackslash}p{(\columnwidth - 4\tabcolsep) * \real{0.20}}
  >{\raggedright\arraybackslash}p{(\columnwidth - 4\tabcolsep) * \real{0.62}}@{}}
\toprule
位置 & 类型 & 含义 \\
\midrule
\endhead
getter 返回值 & \passthrough{\lstinline!list[int]!} &
返回属性当前值。数组、vector 和嵌套消息采用复制语义。 \\
\bottomrule
\end{longtable}

\textbf{对应 C++}

\begin{itemize}
\tightlist
\item
  类：\passthrough{\lstinline!sensor\_msgs::msg::dds\_::PointCloud2\_!}
\item
  字段类型：\passthrough{\lstinline!std::vector<uint8\_t>!}
\item
  声明位置：\passthrough{\lstinline!include/unitree/idl/ros2/PointCloud2\_.hpp:35!}
\end{itemize}

\textbf{用法}

\begin{lstlisting}[language=Python]
current_value = obj.data
obj.data = new_value
\end{lstlisting}

\begin{quote}
{[}!TIP{]} 读取容器或嵌套值后应遵循``读取、修改、重新赋值''。只修改
getter 返回的副本不会可靠地更新原 C++ 消息。
\end{quote}

\hypertarget{unitree-sdk2-cpp-idl-ros2-pointcloud2-is-dense}{%
\paragraph{\texorpdfstring{\texttt{unitree\_sdk2\_cpp.idl.ros2.PointCloud2.is\_dense}}{unitree\_sdk2\_cpp.idl.ros2.PointCloud2.is\_dense}}\label{unitree-sdk2-cpp-idl-ros2-pointcloud2-is-dense}}

底层 \passthrough{\lstinline!sensor\_msgs::msg::dds\_::PointCloud2\_!}
的 \passthrough{\lstinline!is\_dense!} 字段。类型签名只说明 Python
表示；单位、数值范围、数组固定长度和枚举含义需查对应 IDL 头文件。
只接受布尔语义。

\textbf{签名}

\begin{lstlisting}[language=Python]
@property
def is_dense(self) -> bool

@is_dense.setter
def is_dense(self, value: bool) -> None
\end{lstlisting}

\textbf{可用性与安全性}

\textbf{\passthrough{\lstinline!AVAILABLE!} /
\passthrough{\lstinline!VALUE\_TYPE!}}：当前绑定源码已实现该消息值操作。

\textbf{参数}

\begin{longtable}[]{@{}
  >{\raggedright\arraybackslash}p{(\columnwidth - 4\tabcolsep) * \real{0.18}}
  >{\raggedright\arraybackslash}p{(\columnwidth - 4\tabcolsep) * \real{0.20}}
  >{\raggedright\arraybackslash}p{(\columnwidth - 4\tabcolsep) * \real{0.62}}@{}}
\toprule
名称 & Python 类型 & 含义 \\
\midrule
\endhead
\passthrough{\lstinline!value!} & \passthrough{\lstinline!bool!} &
要写入或传递的新值，必须符合该参数的 Python 类型以及底层 C++
范围约束。 \\
\bottomrule
\end{longtable}

\textbf{返回值}

\begin{longtable}[]{@{}
  >{\raggedright\arraybackslash}p{(\columnwidth - 4\tabcolsep) * \real{0.18}}
  >{\raggedright\arraybackslash}p{(\columnwidth - 4\tabcolsep) * \real{0.20}}
  >{\raggedright\arraybackslash}p{(\columnwidth - 4\tabcolsep) * \real{0.62}}@{}}
\toprule
位置 & 类型 & 含义 \\
\midrule
\endhead
getter 返回值 & \passthrough{\lstinline!bool!} & 返回属性当前值。 \\
\bottomrule
\end{longtable}

\textbf{对应 C++}

\begin{itemize}
\tightlist
\item
  类：\passthrough{\lstinline!sensor\_msgs::msg::dds\_::PointCloud2\_!}
\item
  字段类型：\passthrough{\lstinline!bool!}
\item
  声明位置：\passthrough{\lstinline!include/unitree/idl/ros2/PointCloud2\_.hpp:36!}
\end{itemize}

\textbf{用法}

\begin{lstlisting}[language=Python]
current_value = obj.is_dense
obj.is_dense = new_value
\end{lstlisting}

\hypertarget{unitree-sdk2-cpp-idl-ros2-pointstamped}{%
\subsubsection{\texorpdfstring{\texttt{unitree\_sdk2\_cpp.idl.ros2.PointStamped}}{unitree\_sdk2\_cpp.idl.ros2.PointStamped}}\label{unitree-sdk2-cpp-idl-ros2-pointstamped}}

可由 Python 读写的 DDS/IDL 消息值类型。字段采用复制语义。

\textbf{导入}

\begin{lstlisting}[language=Python]
from unitree_sdk2_cpp.idl.ros2 import PointStamped
\end{lstlisting}

\textbf{方法索引}

\begin{longtable}[]{@{}
  >{\raggedright\arraybackslash}p{(\columnwidth - 6\tabcolsep) * \real{0.18}}
  >{\raggedright\arraybackslash}p{(\columnwidth - 6\tabcolsep) * \real{0.24}}
  >{\raggedright\arraybackslash}p{(\columnwidth - 6\tabcolsep) * \real{0.18}}
  >{\raggedright\arraybackslash}p{(\columnwidth - 6\tabcolsep) * \real{0.40}}@{}}
\toprule
方法 & Python 签名 & 状态 & 安全分类 \\
\midrule
\endhead
\protect\hyperlink{unitree-sdk2-cpp-idl-ros2-pointstamped-dunder-init-1}{\passthrough{\lstinline!\_\_init\_\_!}}
& \passthrough{\lstinline!def \_\_init\_\_(self) -> None!} &
\passthrough{\lstinline!AVAILABLE!} &
\passthrough{\lstinline!VALUE\_TYPE!} \\
\protect\hyperlink{unitree-sdk2-cpp-idl-ros2-pointstamped-dunder-eq-1}{\passthrough{\lstinline!\_\_eq\_\_!}}
& \passthrough{\lstinline!def \_\_eq\_\_(self, other: object) -> bool!}
& \passthrough{\lstinline!AVAILABLE!} &
\passthrough{\lstinline!VALUE\_TYPE!} \\
\protect\hyperlink{unitree-sdk2-cpp-idl-ros2-pointstamped-dunder-ne-1}{\passthrough{\lstinline!\_\_ne\_\_!}}
& \passthrough{\lstinline!def \_\_ne\_\_(self, other: object) -> bool!}
& \passthrough{\lstinline!AVAILABLE!} &
\passthrough{\lstinline!VALUE\_TYPE!} \\
\bottomrule
\end{longtable}

\textbf{属性索引}

\begin{longtable}[]{@{}
  >{\raggedright\arraybackslash}p{(\columnwidth - 6\tabcolsep) * \real{0.18}}
  >{\raggedright\arraybackslash}p{(\columnwidth - 6\tabcolsep) * \real{0.24}}
  >{\raggedright\arraybackslash}p{(\columnwidth - 6\tabcolsep) * \real{0.18}}
  >{\raggedright\arraybackslash}p{(\columnwidth - 6\tabcolsep) * \real{0.40}}@{}}
\toprule
属性 & 读取类型 & 写入类型 & 状态 \\
\midrule
\endhead
\protect\hyperlink{unitree-sdk2-cpp-idl-ros2-pointstamped-header}{\passthrough{\lstinline!header!}}
& \passthrough{\lstinline!Any!} & \passthrough{\lstinline!Any!} &
\passthrough{\lstinline!AVAILABLE!} \\
\protect\hyperlink{unitree-sdk2-cpp-idl-ros2-pointstamped-point}{\passthrough{\lstinline!point!}}
& \passthrough{\lstinline!Point!} & \passthrough{\lstinline!Point!} &
\passthrough{\lstinline!AVAILABLE!} \\
\bottomrule
\end{longtable}

\hypertarget{unitree-sdk2-cpp-idl-ros2-pointstamped-dunder-init-1}{%
\paragraph{\texorpdfstring{\texttt{unitree\_sdk2\_cpp.idl.ros2.PointStamped.\_\_init\_\_}}{unitree\_sdk2\_cpp.idl.ros2.PointStamped.\_\_init\_\_}}\label{unitree-sdk2-cpp-idl-ros2-pointstamped-dunder-init-1}}

初始化 \passthrough{\lstinline!PointStamped!}
实例。是否能实际构造取决于下方可用性状态。

\textbf{签名}

\begin{lstlisting}[language=Python]
def __init__(self) -> None
\end{lstlisting}

\textbf{可用性与安全性}

\textbf{\passthrough{\lstinline!AVAILABLE!} /
\passthrough{\lstinline!VALUE\_TYPE!}}：当前绑定源码已实现该消息值操作。

\textbf{参数}

无显式参数。实例方法中的 \passthrough{\lstinline!self!} 由 Python
自动传入。

\textbf{返回值}

\begin{longtable}[]{@{}
  >{\raggedright\arraybackslash}p{(\columnwidth - 4\tabcolsep) * \real{0.18}}
  >{\raggedright\arraybackslash}p{(\columnwidth - 4\tabcolsep) * \real{0.20}}
  >{\raggedright\arraybackslash}p{(\columnwidth - 4\tabcolsep) * \real{0.62}}@{}}
\toprule
位置 & 类型 & 含义 \\
\midrule
\endhead
无 & \passthrough{\lstinline!None!} &
\passthrough{\lstinline!\_\_init\_\_!}
本身不返回值；调用类对象时，在构造函数可用的前提下得到该类实例。 \\
\bottomrule
\end{longtable}

\textbf{对应 C++}

\begin{itemize}
\tightlist
\item
  类：\passthrough{\lstinline!geometry\_msgs::msg::dds\_::PointStamped\_!}
\item
  签名：\passthrough{\lstinline!PointStamped\_()!}
\item
  绑定策略：\passthrough{\lstinline!未单独标注!}
\item
  声明位置：\passthrough{\lstinline!include/unitree/idl/ros2/PointStamped\_.hpp:30!}
\end{itemize}

\textbf{说明}

\begin{itemize}
\tightlist
\item
  示例是局部调用片段；\passthrough{\lstinline!obj!}
  和参数变量表示已经按上层业务校验并准备好的对象。
\end{itemize}

\textbf{用法}

\begin{lstlisting}[language=Python]
value = PointStamped()
\end{lstlisting}

\hypertarget{unitree-sdk2-cpp-idl-ros2-pointstamped-dunder-eq-1}{%
\paragraph{\texorpdfstring{\texttt{unitree\_sdk2\_cpp.idl.ros2.PointStamped.\_\_eq\_\_}}{unitree\_sdk2\_cpp.idl.ros2.PointStamped.\_\_eq\_\_}}\label{unitree-sdk2-cpp-idl-ros2-pointstamped-dunder-eq-1}}

按底层 IDL 消息内容比较两个对象是否相等。

\textbf{签名}

\begin{lstlisting}[language=Python]
def __eq__(self, other: object) -> bool
\end{lstlisting}

\textbf{可用性与安全性}

\textbf{\passthrough{\lstinline!AVAILABLE!} /
\passthrough{\lstinline!VALUE\_TYPE!}}：当前绑定源码已实现该消息值操作。

\textbf{参数}

\begin{longtable}[]{@{}
  >{\raggedright\arraybackslash}p{(\columnwidth - 6\tabcolsep) * \real{0.18}}
  >{\raggedright\arraybackslash}p{(\columnwidth - 6\tabcolsep) * \real{0.24}}
  >{\raggedright\arraybackslash}p{(\columnwidth - 6\tabcolsep) * \real{0.18}}
  >{\raggedright\arraybackslash}p{(\columnwidth - 6\tabcolsep) * \real{0.40}}@{}}
\toprule
名称 & Python 类型 & 默认值 & 含义 \\
\midrule
\endhead
\passthrough{\lstinline!other!} & \passthrough{\lstinline!object!} &
必填 & 用于比较的另一个 Python 对象。类型不兼容时比较结果通常为
\passthrough{\lstinline!False!}。 \\
\bottomrule
\end{longtable}

\textbf{返回值}

\begin{longtable}[]{@{}
  >{\raggedright\arraybackslash}p{(\columnwidth - 4\tabcolsep) * \real{0.18}}
  >{\raggedright\arraybackslash}p{(\columnwidth - 4\tabcolsep) * \real{0.20}}
  >{\raggedright\arraybackslash}p{(\columnwidth - 4\tabcolsep) * \real{0.62}}@{}}
\toprule
位置 & 类型 & 含义 \\
\midrule
\endhead
返回值 & \passthrough{\lstinline!bool!} & 返回
\passthrough{\lstinline!bool!}。更精确的业务含义见该方法用途、C++
签名和目标型号协议。 \\
\bottomrule
\end{longtable}

\textbf{对应 C++}

\begin{itemize}
\tightlist
\item
  类：\passthrough{\lstinline!geometry\_msgs::msg::dds\_::PointStamped\_!}
\item
  签名：\passthrough{\lstinline!operator==(const value \&) const!}
\item
  绑定策略：\passthrough{\lstinline!未单独标注!}
\item
  声明位置：由 pybind11 手工绑定或生成注册表提供；manifest
  未记录独立头文件行号。
\end{itemize}

\textbf{说明}

\begin{itemize}
\tightlist
\item
  示例是局部调用片段；\passthrough{\lstinline!obj!}
  和参数变量表示已经按上层业务校验并准备好的对象。
\end{itemize}

\textbf{用法}

\begin{lstlisting}[language=Python]
same = left == right
\end{lstlisting}

\hypertarget{unitree-sdk2-cpp-idl-ros2-pointstamped-dunder-ne-1}{%
\paragraph{\texorpdfstring{\texttt{unitree\_sdk2\_cpp.idl.ros2.PointStamped.\_\_ne\_\_}}{unitree\_sdk2\_cpp.idl.ros2.PointStamped.\_\_ne\_\_}}\label{unitree-sdk2-cpp-idl-ros2-pointstamped-dunder-ne-1}}

按底层 IDL 消息内容比较两个对象是否不相等。

\textbf{签名}

\begin{lstlisting}[language=Python]
def __ne__(self, other: object) -> bool
\end{lstlisting}

\textbf{可用性与安全性}

\textbf{\passthrough{\lstinline!AVAILABLE!} /
\passthrough{\lstinline!VALUE\_TYPE!}}：当前绑定源码已实现该消息值操作。

\textbf{参数}

\begin{longtable}[]{@{}
  >{\raggedright\arraybackslash}p{(\columnwidth - 6\tabcolsep) * \real{0.18}}
  >{\raggedright\arraybackslash}p{(\columnwidth - 6\tabcolsep) * \real{0.24}}
  >{\raggedright\arraybackslash}p{(\columnwidth - 6\tabcolsep) * \real{0.18}}
  >{\raggedright\arraybackslash}p{(\columnwidth - 6\tabcolsep) * \real{0.40}}@{}}
\toprule
名称 & Python 类型 & 默认值 & 含义 \\
\midrule
\endhead
\passthrough{\lstinline!other!} & \passthrough{\lstinline!object!} &
必填 & 用于比较的另一个 Python 对象。类型不兼容时比较结果通常为
\passthrough{\lstinline!False!}。 \\
\bottomrule
\end{longtable}

\textbf{返回值}

\begin{longtable}[]{@{}
  >{\raggedright\arraybackslash}p{(\columnwidth - 4\tabcolsep) * \real{0.18}}
  >{\raggedright\arraybackslash}p{(\columnwidth - 4\tabcolsep) * \real{0.20}}
  >{\raggedright\arraybackslash}p{(\columnwidth - 4\tabcolsep) * \real{0.62}}@{}}
\toprule
位置 & 类型 & 含义 \\
\midrule
\endhead
返回值 & \passthrough{\lstinline!bool!} & 返回
\passthrough{\lstinline!bool!}。更精确的业务含义见该方法用途、C++
签名和目标型号协议。 \\
\bottomrule
\end{longtable}

\textbf{对应 C++}

\begin{itemize}
\tightlist
\item
  类：\passthrough{\lstinline!geometry\_msgs::msg::dds\_::PointStamped\_!}
\item
  签名：\passthrough{\lstinline"operator!=(const value \&) const"}
\item
  绑定策略：\passthrough{\lstinline!未单独标注!}
\item
  声明位置：由 pybind11 手工绑定或生成注册表提供；manifest
  未记录独立头文件行号。
\end{itemize}

\textbf{说明}

\begin{itemize}
\tightlist
\item
  示例是局部调用片段；\passthrough{\lstinline!obj!}
  和参数变量表示已经按上层业务校验并准备好的对象。
\end{itemize}

\textbf{用法}

\begin{lstlisting}[language=Python]
different = left != right
\end{lstlisting}

\hypertarget{unitree-sdk2-cpp-idl-ros2-pointstamped-header}{%
\paragraph{\texorpdfstring{\texttt{unitree\_sdk2\_cpp.idl.ros2.PointStamped.header}}{unitree\_sdk2\_cpp.idl.ros2.PointStamped.header}}\label{unitree-sdk2-cpp-idl-ros2-pointstamped-header}}

底层
\passthrough{\lstinline!geometry\_msgs::msg::dds\_::PointStamped\_!} 的
\passthrough{\lstinline!header!} 字段。类型签名只说明 Python
表示；单位、数值范围、数组固定长度和枚举含义需查对应 IDL 头文件。

\textbf{签名}

\begin{lstlisting}[language=Python]
@property
def header(self) -> Any

@header.setter
def header(self, value: Any) -> None
\end{lstlisting}

\textbf{可用性与安全性}

\textbf{\passthrough{\lstinline!AVAILABLE!} /
\passthrough{\lstinline!VALUE\_TYPE!}}：当前绑定源码已实现该消息值操作。

\textbf{参数}

\begin{longtable}[]{@{}
  >{\raggedright\arraybackslash}p{(\columnwidth - 4\tabcolsep) * \real{0.18}}
  >{\raggedright\arraybackslash}p{(\columnwidth - 4\tabcolsep) * \real{0.20}}
  >{\raggedright\arraybackslash}p{(\columnwidth - 4\tabcolsep) * \real{0.62}}@{}}
\toprule
名称 & Python 类型 & 含义 \\
\midrule
\endhead
\passthrough{\lstinline!value!} & \passthrough{\lstinline!Any!} &
要写入或传递的新值，必须符合该参数的 Python 类型以及底层 C++
范围约束。 \\
\bottomrule
\end{longtable}

\textbf{返回值}

\begin{longtable}[]{@{}
  >{\raggedright\arraybackslash}p{(\columnwidth - 4\tabcolsep) * \real{0.18}}
  >{\raggedright\arraybackslash}p{(\columnwidth - 4\tabcolsep) * \real{0.20}}
  >{\raggedright\arraybackslash}p{(\columnwidth - 4\tabcolsep) * \real{0.62}}@{}}
\toprule
位置 & 类型 & 含义 \\
\midrule
\endhead
getter 返回值 & \passthrough{\lstinline!Any!} & 返回属性当前值。 \\
\bottomrule
\end{longtable}

\textbf{对应 C++}

\begin{itemize}
\tightlist
\item
  类：\passthrough{\lstinline!geometry\_msgs::msg::dds\_::PointStamped\_!}
\item
  字段类型：\passthrough{\lstinline!::std\_msgs::msg::dds\_::Header\_!}
\item
  声明位置：\passthrough{\lstinline!include/unitree/idl/ros2/PointStamped\_.hpp:26!}
\end{itemize}

\textbf{用法}

\begin{lstlisting}[language=Python]
current_value = obj.header
obj.header = new_value
\end{lstlisting}

\hypertarget{unitree-sdk2-cpp-idl-ros2-pointstamped-point}{%
\paragraph{\texorpdfstring{\texttt{unitree\_sdk2\_cpp.idl.ros2.PointStamped.point}}{unitree\_sdk2\_cpp.idl.ros2.PointStamped.point}}\label{unitree-sdk2-cpp-idl-ros2-pointstamped-point}}

底层
\passthrough{\lstinline!geometry\_msgs::msg::dds\_::PointStamped\_!} 的
\passthrough{\lstinline!point!} 字段。类型签名只说明 Python
表示；单位、数值范围、数组固定长度和枚举含义需查对应 IDL 头文件。

\textbf{签名}

\begin{lstlisting}[language=Python]
@property
def point(self) -> Point

@point.setter
def point(self, value: Point) -> None
\end{lstlisting}

\textbf{可用性与安全性}

\textbf{\passthrough{\lstinline!AVAILABLE!} /
\passthrough{\lstinline!VALUE\_TYPE!}}：当前绑定源码已实现该消息值操作。

\textbf{参数}

\begin{longtable}[]{@{}
  >{\raggedright\arraybackslash}p{(\columnwidth - 4\tabcolsep) * \real{0.18}}
  >{\raggedright\arraybackslash}p{(\columnwidth - 4\tabcolsep) * \real{0.20}}
  >{\raggedright\arraybackslash}p{(\columnwidth - 4\tabcolsep) * \real{0.62}}@{}}
\toprule
名称 & Python 类型 & 含义 \\
\midrule
\endhead
\passthrough{\lstinline!value!} & \passthrough{\lstinline!Point!} &
要写入或传递的新值，必须符合该参数的 Python 类型以及底层 C++
范围约束。 \\
\bottomrule
\end{longtable}

\textbf{返回值}

\begin{longtable}[]{@{}
  >{\raggedright\arraybackslash}p{(\columnwidth - 4\tabcolsep) * \real{0.18}}
  >{\raggedright\arraybackslash}p{(\columnwidth - 4\tabcolsep) * \real{0.20}}
  >{\raggedright\arraybackslash}p{(\columnwidth - 4\tabcolsep) * \real{0.62}}@{}}
\toprule
位置 & 类型 & 含义 \\
\midrule
\endhead
getter 返回值 & \passthrough{\lstinline!Point!} &
返回属性当前值。数组、vector 和嵌套消息采用复制语义。 \\
\bottomrule
\end{longtable}

\textbf{对应 C++}

\begin{itemize}
\tightlist
\item
  类：\passthrough{\lstinline!geometry\_msgs::msg::dds\_::PointStamped\_!}
\item
  字段类型：\passthrough{\lstinline!::geometry\_msgs::msg::dds\_::Point\_!}
\item
  声明位置：\passthrough{\lstinline!include/unitree/idl/ros2/PointStamped\_.hpp:27!}
\end{itemize}

\textbf{用法}

\begin{lstlisting}[language=Python]
current_value = obj.point
obj.point = new_value
\end{lstlisting}

\hypertarget{unitree-sdk2-cpp-idl-ros2-pose2d}{%
\subsubsection{\texorpdfstring{\texttt{unitree\_sdk2\_cpp.idl.ros2.Pose2D}}{unitree\_sdk2\_cpp.idl.ros2.Pose2D}}\label{unitree-sdk2-cpp-idl-ros2-pose2d}}

可由 Python 读写的 DDS/IDL 消息值类型。字段采用复制语义。

\textbf{导入}

\begin{lstlisting}[language=Python]
from unitree_sdk2_cpp.idl.ros2 import Pose2D
\end{lstlisting}

\textbf{方法索引}

\begin{longtable}[]{@{}
  >{\raggedright\arraybackslash}p{(\columnwidth - 6\tabcolsep) * \real{0.18}}
  >{\raggedright\arraybackslash}p{(\columnwidth - 6\tabcolsep) * \real{0.24}}
  >{\raggedright\arraybackslash}p{(\columnwidth - 6\tabcolsep) * \real{0.18}}
  >{\raggedright\arraybackslash}p{(\columnwidth - 6\tabcolsep) * \real{0.40}}@{}}
\toprule
方法 & Python 签名 & 状态 & 安全分类 \\
\midrule
\endhead
\protect\hyperlink{unitree-sdk2-cpp-idl-ros2-pose2d-dunder-init-1}{\passthrough{\lstinline!\_\_init\_\_!}}
& \passthrough{\lstinline!def \_\_init\_\_(self) -> None!} &
\passthrough{\lstinline!AVAILABLE!} &
\passthrough{\lstinline!VALUE\_TYPE!} \\
\protect\hyperlink{unitree-sdk2-cpp-idl-ros2-pose2d-dunder-eq-1}{\passthrough{\lstinline!\_\_eq\_\_!}}
& \passthrough{\lstinline!def \_\_eq\_\_(self, other: object) -> bool!}
& \passthrough{\lstinline!AVAILABLE!} &
\passthrough{\lstinline!VALUE\_TYPE!} \\
\protect\hyperlink{unitree-sdk2-cpp-idl-ros2-pose2d-dunder-ne-1}{\passthrough{\lstinline!\_\_ne\_\_!}}
& \passthrough{\lstinline!def \_\_ne\_\_(self, other: object) -> bool!}
& \passthrough{\lstinline!AVAILABLE!} &
\passthrough{\lstinline!VALUE\_TYPE!} \\
\bottomrule
\end{longtable}

\textbf{属性索引}

\begin{longtable}[]{@{}
  >{\raggedright\arraybackslash}p{(\columnwidth - 6\tabcolsep) * \real{0.18}}
  >{\raggedright\arraybackslash}p{(\columnwidth - 6\tabcolsep) * \real{0.24}}
  >{\raggedright\arraybackslash}p{(\columnwidth - 6\tabcolsep) * \real{0.18}}
  >{\raggedright\arraybackslash}p{(\columnwidth - 6\tabcolsep) * \real{0.40}}@{}}
\toprule
属性 & 读取类型 & 写入类型 & 状态 \\
\midrule
\endhead
\protect\hyperlink{unitree-sdk2-cpp-idl-ros2-pose2d-x}{\passthrough{\lstinline!x!}}
& \passthrough{\lstinline!float!} & \passthrough{\lstinline!float!} &
\passthrough{\lstinline!AVAILABLE!} \\
\protect\hyperlink{unitree-sdk2-cpp-idl-ros2-pose2d-y}{\passthrough{\lstinline!y!}}
& \passthrough{\lstinline!float!} & \passthrough{\lstinline!float!} &
\passthrough{\lstinline!AVAILABLE!} \\
\protect\hyperlink{unitree-sdk2-cpp-idl-ros2-pose2d-theta}{\passthrough{\lstinline!theta!}}
& \passthrough{\lstinline!float!} & \passthrough{\lstinline!float!} &
\passthrough{\lstinline!AVAILABLE!} \\
\bottomrule
\end{longtable}

\hypertarget{unitree-sdk2-cpp-idl-ros2-pose2d-dunder-init-1}{%
\paragraph{\texorpdfstring{\texttt{unitree\_sdk2\_cpp.idl.ros2.Pose2D.\_\_init\_\_}}{unitree\_sdk2\_cpp.idl.ros2.Pose2D.\_\_init\_\_}}\label{unitree-sdk2-cpp-idl-ros2-pose2d-dunder-init-1}}

初始化 \passthrough{\lstinline!Pose2D!}
实例。是否能实际构造取决于下方可用性状态。

\textbf{签名}

\begin{lstlisting}[language=Python]
def __init__(self) -> None
\end{lstlisting}

\textbf{可用性与安全性}

\textbf{\passthrough{\lstinline!AVAILABLE!} /
\passthrough{\lstinline!VALUE\_TYPE!}}：当前绑定源码已实现该消息值操作。

\textbf{参数}

无显式参数。实例方法中的 \passthrough{\lstinline!self!} 由 Python
自动传入。

\textbf{返回值}

\begin{longtable}[]{@{}
  >{\raggedright\arraybackslash}p{(\columnwidth - 4\tabcolsep) * \real{0.18}}
  >{\raggedright\arraybackslash}p{(\columnwidth - 4\tabcolsep) * \real{0.20}}
  >{\raggedright\arraybackslash}p{(\columnwidth - 4\tabcolsep) * \real{0.62}}@{}}
\toprule
位置 & 类型 & 含义 \\
\midrule
\endhead
无 & \passthrough{\lstinline!None!} &
\passthrough{\lstinline!\_\_init\_\_!}
本身不返回值；调用类对象时，在构造函数可用的前提下得到该类实例。 \\
\bottomrule
\end{longtable}

\textbf{对应 C++}

\begin{itemize}
\tightlist
\item
  类：\passthrough{\lstinline!geometry\_msgs::msg::dds\_::Pose2D\_!}
\item
  签名：\passthrough{\lstinline!Pose2D\_()!}
\item
  绑定策略：\passthrough{\lstinline!未单独标注!}
\item
  声明位置：\passthrough{\lstinline!include/unitree/idl/ros2/Pose2D\_.hpp:27!}
\end{itemize}

\textbf{说明}

\begin{itemize}
\tightlist
\item
  示例是局部调用片段；\passthrough{\lstinline!obj!}
  和参数变量表示已经按上层业务校验并准备好的对象。
\end{itemize}

\textbf{用法}

\begin{lstlisting}[language=Python]
value = Pose2D()
\end{lstlisting}

\hypertarget{unitree-sdk2-cpp-idl-ros2-pose2d-dunder-eq-1}{%
\paragraph{\texorpdfstring{\texttt{unitree\_sdk2\_cpp.idl.ros2.Pose2D.\_\_eq\_\_}}{unitree\_sdk2\_cpp.idl.ros2.Pose2D.\_\_eq\_\_}}\label{unitree-sdk2-cpp-idl-ros2-pose2d-dunder-eq-1}}

按底层 IDL 消息内容比较两个对象是否相等。

\textbf{签名}

\begin{lstlisting}[language=Python]
def __eq__(self, other: object) -> bool
\end{lstlisting}

\textbf{可用性与安全性}

\textbf{\passthrough{\lstinline!AVAILABLE!} /
\passthrough{\lstinline!VALUE\_TYPE!}}：当前绑定源码已实现该消息值操作。

\textbf{参数}

\begin{longtable}[]{@{}
  >{\raggedright\arraybackslash}p{(\columnwidth - 6\tabcolsep) * \real{0.18}}
  >{\raggedright\arraybackslash}p{(\columnwidth - 6\tabcolsep) * \real{0.24}}
  >{\raggedright\arraybackslash}p{(\columnwidth - 6\tabcolsep) * \real{0.18}}
  >{\raggedright\arraybackslash}p{(\columnwidth - 6\tabcolsep) * \real{0.40}}@{}}
\toprule
名称 & Python 类型 & 默认值 & 含义 \\
\midrule
\endhead
\passthrough{\lstinline!other!} & \passthrough{\lstinline!object!} &
必填 & 用于比较的另一个 Python 对象。类型不兼容时比较结果通常为
\passthrough{\lstinline!False!}。 \\
\bottomrule
\end{longtable}

\textbf{返回值}

\begin{longtable}[]{@{}
  >{\raggedright\arraybackslash}p{(\columnwidth - 4\tabcolsep) * \real{0.18}}
  >{\raggedright\arraybackslash}p{(\columnwidth - 4\tabcolsep) * \real{0.20}}
  >{\raggedright\arraybackslash}p{(\columnwidth - 4\tabcolsep) * \real{0.62}}@{}}
\toprule
位置 & 类型 & 含义 \\
\midrule
\endhead
返回值 & \passthrough{\lstinline!bool!} & 返回
\passthrough{\lstinline!bool!}。更精确的业务含义见该方法用途、C++
签名和目标型号协议。 \\
\bottomrule
\end{longtable}

\textbf{对应 C++}

\begin{itemize}
\tightlist
\item
  类：\passthrough{\lstinline!geometry\_msgs::msg::dds\_::Pose2D\_!}
\item
  签名：\passthrough{\lstinline!operator==(const value \&) const!}
\item
  绑定策略：\passthrough{\lstinline!未单独标注!}
\item
  声明位置：由 pybind11 手工绑定或生成注册表提供；manifest
  未记录独立头文件行号。
\end{itemize}

\textbf{说明}

\begin{itemize}
\tightlist
\item
  示例是局部调用片段；\passthrough{\lstinline!obj!}
  和参数变量表示已经按上层业务校验并准备好的对象。
\end{itemize}

\textbf{用法}

\begin{lstlisting}[language=Python]
same = left == right
\end{lstlisting}

\hypertarget{unitree-sdk2-cpp-idl-ros2-pose2d-dunder-ne-1}{%
\paragraph{\texorpdfstring{\texttt{unitree\_sdk2\_cpp.idl.ros2.Pose2D.\_\_ne\_\_}}{unitree\_sdk2\_cpp.idl.ros2.Pose2D.\_\_ne\_\_}}\label{unitree-sdk2-cpp-idl-ros2-pose2d-dunder-ne-1}}

按底层 IDL 消息内容比较两个对象是否不相等。

\textbf{签名}

\begin{lstlisting}[language=Python]
def __ne__(self, other: object) -> bool
\end{lstlisting}

\textbf{可用性与安全性}

\textbf{\passthrough{\lstinline!AVAILABLE!} /
\passthrough{\lstinline!VALUE\_TYPE!}}：当前绑定源码已实现该消息值操作。

\textbf{参数}

\begin{longtable}[]{@{}
  >{\raggedright\arraybackslash}p{(\columnwidth - 6\tabcolsep) * \real{0.18}}
  >{\raggedright\arraybackslash}p{(\columnwidth - 6\tabcolsep) * \real{0.24}}
  >{\raggedright\arraybackslash}p{(\columnwidth - 6\tabcolsep) * \real{0.18}}
  >{\raggedright\arraybackslash}p{(\columnwidth - 6\tabcolsep) * \real{0.40}}@{}}
\toprule
名称 & Python 类型 & 默认值 & 含义 \\
\midrule
\endhead
\passthrough{\lstinline!other!} & \passthrough{\lstinline!object!} &
必填 & 用于比较的另一个 Python 对象。类型不兼容时比较结果通常为
\passthrough{\lstinline!False!}。 \\
\bottomrule
\end{longtable}

\textbf{返回值}

\begin{longtable}[]{@{}
  >{\raggedright\arraybackslash}p{(\columnwidth - 4\tabcolsep) * \real{0.18}}
  >{\raggedright\arraybackslash}p{(\columnwidth - 4\tabcolsep) * \real{0.20}}
  >{\raggedright\arraybackslash}p{(\columnwidth - 4\tabcolsep) * \real{0.62}}@{}}
\toprule
位置 & 类型 & 含义 \\
\midrule
\endhead
返回值 & \passthrough{\lstinline!bool!} & 返回
\passthrough{\lstinline!bool!}。更精确的业务含义见该方法用途、C++
签名和目标型号协议。 \\
\bottomrule
\end{longtable}

\textbf{对应 C++}

\begin{itemize}
\tightlist
\item
  类：\passthrough{\lstinline!geometry\_msgs::msg::dds\_::Pose2D\_!}
\item
  签名：\passthrough{\lstinline"operator!=(const value \&) const"}
\item
  绑定策略：\passthrough{\lstinline!未单独标注!}
\item
  声明位置：由 pybind11 手工绑定或生成注册表提供；manifest
  未记录独立头文件行号。
\end{itemize}

\textbf{说明}

\begin{itemize}
\tightlist
\item
  示例是局部调用片段；\passthrough{\lstinline!obj!}
  和参数变量表示已经按上层业务校验并准备好的对象。
\end{itemize}

\textbf{用法}

\begin{lstlisting}[language=Python]
different = left != right
\end{lstlisting}

\hypertarget{unitree-sdk2-cpp-idl-ros2-pose2d-x}{%
\paragraph{\texorpdfstring{\texttt{unitree\_sdk2\_cpp.idl.ros2.Pose2D.x}}{unitree\_sdk2\_cpp.idl.ros2.Pose2D.x}}\label{unitree-sdk2-cpp-idl-ros2-pose2d-x}}

底层 \passthrough{\lstinline!geometry\_msgs::msg::dds\_::Pose2D\_!} 的
\passthrough{\lstinline!x!} 字段。类型签名只说明 Python
表示；单位、数值范围、数组固定长度和枚举含义需查对应 IDL 头文件。
底层为浮点数；类型本身不说明单位、坐标系或业务安全范围。

\textbf{签名}

\begin{lstlisting}[language=Python]
@property
def x(self) -> float

@x.setter
def x(self, value: float) -> None
\end{lstlisting}

\textbf{可用性与安全性}

\textbf{\passthrough{\lstinline!AVAILABLE!} /
\passthrough{\lstinline!VALUE\_TYPE!}}：当前绑定源码已实现该消息值操作。

\textbf{参数}

\begin{longtable}[]{@{}
  >{\raggedright\arraybackslash}p{(\columnwidth - 4\tabcolsep) * \real{0.18}}
  >{\raggedright\arraybackslash}p{(\columnwidth - 4\tabcolsep) * \real{0.20}}
  >{\raggedright\arraybackslash}p{(\columnwidth - 4\tabcolsep) * \real{0.62}}@{}}
\toprule
名称 & Python 类型 & 含义 \\
\midrule
\endhead
\passthrough{\lstinline!value!} & \passthrough{\lstinline!float!} &
要写入或传递的新值，必须符合该参数的 Python 类型以及底层 C++
范围约束。 \\
\bottomrule
\end{longtable}

\textbf{返回值}

\begin{longtable}[]{@{}
  >{\raggedright\arraybackslash}p{(\columnwidth - 4\tabcolsep) * \real{0.18}}
  >{\raggedright\arraybackslash}p{(\columnwidth - 4\tabcolsep) * \real{0.20}}
  >{\raggedright\arraybackslash}p{(\columnwidth - 4\tabcolsep) * \real{0.62}}@{}}
\toprule
位置 & 类型 & 含义 \\
\midrule
\endhead
getter 返回值 & \passthrough{\lstinline!float!} & 返回属性当前值。 \\
\bottomrule
\end{longtable}

\textbf{对应 C++}

\begin{itemize}
\tightlist
\item
  类：\passthrough{\lstinline!geometry\_msgs::msg::dds\_::Pose2D\_!}
\item
  字段类型：\passthrough{\lstinline!double!}
\item
  声明位置：\passthrough{\lstinline!include/unitree/idl/ros2/Pose2D\_.hpp:22!}
\end{itemize}

\textbf{用法}

\begin{lstlisting}[language=Python]
current_value = obj.x
obj.x = new_value
\end{lstlisting}

\hypertarget{unitree-sdk2-cpp-idl-ros2-pose2d-y}{%
\paragraph{\texorpdfstring{\texttt{unitree\_sdk2\_cpp.idl.ros2.Pose2D.y}}{unitree\_sdk2\_cpp.idl.ros2.Pose2D.y}}\label{unitree-sdk2-cpp-idl-ros2-pose2d-y}}

底层 \passthrough{\lstinline!geometry\_msgs::msg::dds\_::Pose2D\_!} 的
\passthrough{\lstinline!y!} 字段。类型签名只说明 Python
表示；单位、数值范围、数组固定长度和枚举含义需查对应 IDL 头文件。
底层为浮点数；类型本身不说明单位、坐标系或业务安全范围。

\textbf{签名}

\begin{lstlisting}[language=Python]
@property
def y(self) -> float

@y.setter
def y(self, value: float) -> None
\end{lstlisting}

\textbf{可用性与安全性}

\textbf{\passthrough{\lstinline!AVAILABLE!} /
\passthrough{\lstinline!VALUE\_TYPE!}}：当前绑定源码已实现该消息值操作。

\textbf{参数}

\begin{longtable}[]{@{}
  >{\raggedright\arraybackslash}p{(\columnwidth - 4\tabcolsep) * \real{0.18}}
  >{\raggedright\arraybackslash}p{(\columnwidth - 4\tabcolsep) * \real{0.20}}
  >{\raggedright\arraybackslash}p{(\columnwidth - 4\tabcolsep) * \real{0.62}}@{}}
\toprule
名称 & Python 类型 & 含义 \\
\midrule
\endhead
\passthrough{\lstinline!value!} & \passthrough{\lstinline!float!} &
要写入或传递的新值，必须符合该参数的 Python 类型以及底层 C++
范围约束。 \\
\bottomrule
\end{longtable}

\textbf{返回值}

\begin{longtable}[]{@{}
  >{\raggedright\arraybackslash}p{(\columnwidth - 4\tabcolsep) * \real{0.18}}
  >{\raggedright\arraybackslash}p{(\columnwidth - 4\tabcolsep) * \real{0.20}}
  >{\raggedright\arraybackslash}p{(\columnwidth - 4\tabcolsep) * \real{0.62}}@{}}
\toprule
位置 & 类型 & 含义 \\
\midrule
\endhead
getter 返回值 & \passthrough{\lstinline!float!} & 返回属性当前值。 \\
\bottomrule
\end{longtable}

\textbf{对应 C++}

\begin{itemize}
\tightlist
\item
  类：\passthrough{\lstinline!geometry\_msgs::msg::dds\_::Pose2D\_!}
\item
  字段类型：\passthrough{\lstinline!double!}
\item
  声明位置：\passthrough{\lstinline!include/unitree/idl/ros2/Pose2D\_.hpp:23!}
\end{itemize}

\textbf{用法}

\begin{lstlisting}[language=Python]
current_value = obj.y
obj.y = new_value
\end{lstlisting}

\hypertarget{unitree-sdk2-cpp-idl-ros2-pose2d-theta}{%
\paragraph{\texorpdfstring{\texttt{unitree\_sdk2\_cpp.idl.ros2.Pose2D.theta}}{unitree\_sdk2\_cpp.idl.ros2.Pose2D.theta}}\label{unitree-sdk2-cpp-idl-ros2-pose2d-theta}}

底层 \passthrough{\lstinline!geometry\_msgs::msg::dds\_::Pose2D\_!} 的
\passthrough{\lstinline!theta!} 字段。类型签名只说明 Python
表示；单位、数值范围、数组固定长度和枚举含义需查对应 IDL 头文件。
底层为浮点数；类型本身不说明单位、坐标系或业务安全范围。

\textbf{签名}

\begin{lstlisting}[language=Python]
@property
def theta(self) -> float

@theta.setter
def theta(self, value: float) -> None
\end{lstlisting}

\textbf{可用性与安全性}

\textbf{\passthrough{\lstinline!AVAILABLE!} /
\passthrough{\lstinline!VALUE\_TYPE!}}：当前绑定源码已实现该消息值操作。

\textbf{参数}

\begin{longtable}[]{@{}
  >{\raggedright\arraybackslash}p{(\columnwidth - 4\tabcolsep) * \real{0.18}}
  >{\raggedright\arraybackslash}p{(\columnwidth - 4\tabcolsep) * \real{0.20}}
  >{\raggedright\arraybackslash}p{(\columnwidth - 4\tabcolsep) * \real{0.62}}@{}}
\toprule
名称 & Python 类型 & 含义 \\
\midrule
\endhead
\passthrough{\lstinline!value!} & \passthrough{\lstinline!float!} &
要写入或传递的新值，必须符合该参数的 Python 类型以及底层 C++
范围约束。 \\
\bottomrule
\end{longtable}

\textbf{返回值}

\begin{longtable}[]{@{}
  >{\raggedright\arraybackslash}p{(\columnwidth - 4\tabcolsep) * \real{0.18}}
  >{\raggedright\arraybackslash}p{(\columnwidth - 4\tabcolsep) * \real{0.20}}
  >{\raggedright\arraybackslash}p{(\columnwidth - 4\tabcolsep) * \real{0.62}}@{}}
\toprule
位置 & 类型 & 含义 \\
\midrule
\endhead
getter 返回值 & \passthrough{\lstinline!float!} & 返回属性当前值。 \\
\bottomrule
\end{longtable}

\textbf{对应 C++}

\begin{itemize}
\tightlist
\item
  类：\passthrough{\lstinline!geometry\_msgs::msg::dds\_::Pose2D\_!}
\item
  字段类型：\passthrough{\lstinline!double!}
\item
  声明位置：\passthrough{\lstinline!include/unitree/idl/ros2/Pose2D\_.hpp:24!}
\end{itemize}

\textbf{用法}

\begin{lstlisting}[language=Python]
current_value = obj.theta
obj.theta = new_value
\end{lstlisting}

\hypertarget{unitree-sdk2-cpp-idl-ros2-posestamped}{%
\subsubsection{\texorpdfstring{\texttt{unitree\_sdk2\_cpp.idl.ros2.PoseStamped}}{unitree\_sdk2\_cpp.idl.ros2.PoseStamped}}\label{unitree-sdk2-cpp-idl-ros2-posestamped}}

可由 Python 读写的 DDS/IDL 消息值类型。字段采用复制语义。

\textbf{导入}

\begin{lstlisting}[language=Python]
from unitree_sdk2_cpp.idl.ros2 import PoseStamped
\end{lstlisting}

\textbf{方法索引}

\begin{longtable}[]{@{}
  >{\raggedright\arraybackslash}p{(\columnwidth - 6\tabcolsep) * \real{0.18}}
  >{\raggedright\arraybackslash}p{(\columnwidth - 6\tabcolsep) * \real{0.24}}
  >{\raggedright\arraybackslash}p{(\columnwidth - 6\tabcolsep) * \real{0.18}}
  >{\raggedright\arraybackslash}p{(\columnwidth - 6\tabcolsep) * \real{0.40}}@{}}
\toprule
方法 & Python 签名 & 状态 & 安全分类 \\
\midrule
\endhead
\protect\hyperlink{unitree-sdk2-cpp-idl-ros2-posestamped-dunder-init-1}{\passthrough{\lstinline!\_\_init\_\_!}}
& \passthrough{\lstinline!def \_\_init\_\_(self) -> None!} &
\passthrough{\lstinline!AVAILABLE!} &
\passthrough{\lstinline!VALUE\_TYPE!} \\
\protect\hyperlink{unitree-sdk2-cpp-idl-ros2-posestamped-dunder-eq-1}{\passthrough{\lstinline!\_\_eq\_\_!}}
& \passthrough{\lstinline!def \_\_eq\_\_(self, other: object) -> bool!}
& \passthrough{\lstinline!AVAILABLE!} &
\passthrough{\lstinline!VALUE\_TYPE!} \\
\protect\hyperlink{unitree-sdk2-cpp-idl-ros2-posestamped-dunder-ne-1}{\passthrough{\lstinline!\_\_ne\_\_!}}
& \passthrough{\lstinline!def \_\_ne\_\_(self, other: object) -> bool!}
& \passthrough{\lstinline!AVAILABLE!} &
\passthrough{\lstinline!VALUE\_TYPE!} \\
\bottomrule
\end{longtable}

\textbf{属性索引}

\begin{longtable}[]{@{}
  >{\raggedright\arraybackslash}p{(\columnwidth - 6\tabcolsep) * \real{0.18}}
  >{\raggedright\arraybackslash}p{(\columnwidth - 6\tabcolsep) * \real{0.24}}
  >{\raggedright\arraybackslash}p{(\columnwidth - 6\tabcolsep) * \real{0.18}}
  >{\raggedright\arraybackslash}p{(\columnwidth - 6\tabcolsep) * \real{0.40}}@{}}
\toprule
属性 & 读取类型 & 写入类型 & 状态 \\
\midrule
\endhead
\protect\hyperlink{unitree-sdk2-cpp-idl-ros2-posestamped-header}{\passthrough{\lstinline!header!}}
& \passthrough{\lstinline!Any!} & \passthrough{\lstinline!Any!} &
\passthrough{\lstinline!AVAILABLE!} \\
\protect\hyperlink{unitree-sdk2-cpp-idl-ros2-posestamped-pose}{\passthrough{\lstinline!pose!}}
& \passthrough{\lstinline!Pose!} & \passthrough{\lstinline!Pose!} &
\passthrough{\lstinline!AVAILABLE!} \\
\bottomrule
\end{longtable}

\hypertarget{unitree-sdk2-cpp-idl-ros2-posestamped-dunder-init-1}{%
\paragraph{\texorpdfstring{\texttt{unitree\_sdk2\_cpp.idl.ros2.PoseStamped.\_\_init\_\_}}{unitree\_sdk2\_cpp.idl.ros2.PoseStamped.\_\_init\_\_}}\label{unitree-sdk2-cpp-idl-ros2-posestamped-dunder-init-1}}

初始化 \passthrough{\lstinline!PoseStamped!}
实例。是否能实际构造取决于下方可用性状态。

\textbf{签名}

\begin{lstlisting}[language=Python]
def __init__(self) -> None
\end{lstlisting}

\textbf{可用性与安全性}

\textbf{\passthrough{\lstinline!AVAILABLE!} /
\passthrough{\lstinline!VALUE\_TYPE!}}：当前绑定源码已实现该消息值操作。

\textbf{参数}

无显式参数。实例方法中的 \passthrough{\lstinline!self!} 由 Python
自动传入。

\textbf{返回值}

\begin{longtable}[]{@{}
  >{\raggedright\arraybackslash}p{(\columnwidth - 4\tabcolsep) * \real{0.18}}
  >{\raggedright\arraybackslash}p{(\columnwidth - 4\tabcolsep) * \real{0.20}}
  >{\raggedright\arraybackslash}p{(\columnwidth - 4\tabcolsep) * \real{0.62}}@{}}
\toprule
位置 & 类型 & 含义 \\
\midrule
\endhead
无 & \passthrough{\lstinline!None!} &
\passthrough{\lstinline!\_\_init\_\_!}
本身不返回值；调用类对象时，在构造函数可用的前提下得到该类实例。 \\
\bottomrule
\end{longtable}

\textbf{对应 C++}

\begin{itemize}
\tightlist
\item
  类：\passthrough{\lstinline!geometry\_msgs::msg::dds\_::PoseStamped\_!}
\item
  签名：\passthrough{\lstinline!PoseStamped\_()!}
\item
  绑定策略：\passthrough{\lstinline!未单独标注!}
\item
  声明位置：\passthrough{\lstinline!include/unitree/idl/ros2/PoseStamped\_.hpp:30!}
\end{itemize}

\textbf{说明}

\begin{itemize}
\tightlist
\item
  示例是局部调用片段；\passthrough{\lstinline!obj!}
  和参数变量表示已经按上层业务校验并准备好的对象。
\end{itemize}

\textbf{用法}

\begin{lstlisting}[language=Python]
value = PoseStamped()
\end{lstlisting}

\hypertarget{unitree-sdk2-cpp-idl-ros2-posestamped-dunder-eq-1}{%
\paragraph{\texorpdfstring{\texttt{unitree\_sdk2\_cpp.idl.ros2.PoseStamped.\_\_eq\_\_}}{unitree\_sdk2\_cpp.idl.ros2.PoseStamped.\_\_eq\_\_}}\label{unitree-sdk2-cpp-idl-ros2-posestamped-dunder-eq-1}}

按底层 IDL 消息内容比较两个对象是否相等。

\textbf{签名}

\begin{lstlisting}[language=Python]
def __eq__(self, other: object) -> bool
\end{lstlisting}

\textbf{可用性与安全性}

\textbf{\passthrough{\lstinline!AVAILABLE!} /
\passthrough{\lstinline!VALUE\_TYPE!}}：当前绑定源码已实现该消息值操作。

\textbf{参数}

\begin{longtable}[]{@{}
  >{\raggedright\arraybackslash}p{(\columnwidth - 6\tabcolsep) * \real{0.18}}
  >{\raggedright\arraybackslash}p{(\columnwidth - 6\tabcolsep) * \real{0.24}}
  >{\raggedright\arraybackslash}p{(\columnwidth - 6\tabcolsep) * \real{0.18}}
  >{\raggedright\arraybackslash}p{(\columnwidth - 6\tabcolsep) * \real{0.40}}@{}}
\toprule
名称 & Python 类型 & 默认值 & 含义 \\
\midrule
\endhead
\passthrough{\lstinline!other!} & \passthrough{\lstinline!object!} &
必填 & 用于比较的另一个 Python 对象。类型不兼容时比较结果通常为
\passthrough{\lstinline!False!}。 \\
\bottomrule
\end{longtable}

\textbf{返回值}

\begin{longtable}[]{@{}
  >{\raggedright\arraybackslash}p{(\columnwidth - 4\tabcolsep) * \real{0.18}}
  >{\raggedright\arraybackslash}p{(\columnwidth - 4\tabcolsep) * \real{0.20}}
  >{\raggedright\arraybackslash}p{(\columnwidth - 4\tabcolsep) * \real{0.62}}@{}}
\toprule
位置 & 类型 & 含义 \\
\midrule
\endhead
返回值 & \passthrough{\lstinline!bool!} & 返回
\passthrough{\lstinline!bool!}。更精确的业务含义见该方法用途、C++
签名和目标型号协议。 \\
\bottomrule
\end{longtable}

\textbf{对应 C++}

\begin{itemize}
\tightlist
\item
  类：\passthrough{\lstinline!geometry\_msgs::msg::dds\_::PoseStamped\_!}
\item
  签名：\passthrough{\lstinline!operator==(const value \&) const!}
\item
  绑定策略：\passthrough{\lstinline!未单独标注!}
\item
  声明位置：由 pybind11 手工绑定或生成注册表提供；manifest
  未记录独立头文件行号。
\end{itemize}

\textbf{说明}

\begin{itemize}
\tightlist
\item
  示例是局部调用片段；\passthrough{\lstinline!obj!}
  和参数变量表示已经按上层业务校验并准备好的对象。
\end{itemize}

\textbf{用法}

\begin{lstlisting}[language=Python]
same = left == right
\end{lstlisting}

\hypertarget{unitree-sdk2-cpp-idl-ros2-posestamped-dunder-ne-1}{%
\paragraph{\texorpdfstring{\texttt{unitree\_sdk2\_cpp.idl.ros2.PoseStamped.\_\_ne\_\_}}{unitree\_sdk2\_cpp.idl.ros2.PoseStamped.\_\_ne\_\_}}\label{unitree-sdk2-cpp-idl-ros2-posestamped-dunder-ne-1}}

按底层 IDL 消息内容比较两个对象是否不相等。

\textbf{签名}

\begin{lstlisting}[language=Python]
def __ne__(self, other: object) -> bool
\end{lstlisting}

\textbf{可用性与安全性}

\textbf{\passthrough{\lstinline!AVAILABLE!} /
\passthrough{\lstinline!VALUE\_TYPE!}}：当前绑定源码已实现该消息值操作。

\textbf{参数}

\begin{longtable}[]{@{}
  >{\raggedright\arraybackslash}p{(\columnwidth - 6\tabcolsep) * \real{0.18}}
  >{\raggedright\arraybackslash}p{(\columnwidth - 6\tabcolsep) * \real{0.24}}
  >{\raggedright\arraybackslash}p{(\columnwidth - 6\tabcolsep) * \real{0.18}}
  >{\raggedright\arraybackslash}p{(\columnwidth - 6\tabcolsep) * \real{0.40}}@{}}
\toprule
名称 & Python 类型 & 默认值 & 含义 \\
\midrule
\endhead
\passthrough{\lstinline!other!} & \passthrough{\lstinline!object!} &
必填 & 用于比较的另一个 Python 对象。类型不兼容时比较结果通常为
\passthrough{\lstinline!False!}。 \\
\bottomrule
\end{longtable}

\textbf{返回值}

\begin{longtable}[]{@{}
  >{\raggedright\arraybackslash}p{(\columnwidth - 4\tabcolsep) * \real{0.18}}
  >{\raggedright\arraybackslash}p{(\columnwidth - 4\tabcolsep) * \real{0.20}}
  >{\raggedright\arraybackslash}p{(\columnwidth - 4\tabcolsep) * \real{0.62}}@{}}
\toprule
位置 & 类型 & 含义 \\
\midrule
\endhead
返回值 & \passthrough{\lstinline!bool!} & 返回
\passthrough{\lstinline!bool!}。更精确的业务含义见该方法用途、C++
签名和目标型号协议。 \\
\bottomrule
\end{longtable}

\textbf{对应 C++}

\begin{itemize}
\tightlist
\item
  类：\passthrough{\lstinline!geometry\_msgs::msg::dds\_::PoseStamped\_!}
\item
  签名：\passthrough{\lstinline"operator!=(const value \&) const"}
\item
  绑定策略：\passthrough{\lstinline!未单独标注!}
\item
  声明位置：由 pybind11 手工绑定或生成注册表提供；manifest
  未记录独立头文件行号。
\end{itemize}

\textbf{说明}

\begin{itemize}
\tightlist
\item
  示例是局部调用片段；\passthrough{\lstinline!obj!}
  和参数变量表示已经按上层业务校验并准备好的对象。
\end{itemize}

\textbf{用法}

\begin{lstlisting}[language=Python]
different = left != right
\end{lstlisting}

\hypertarget{unitree-sdk2-cpp-idl-ros2-posestamped-header}{%
\paragraph{\texorpdfstring{\texttt{unitree\_sdk2\_cpp.idl.ros2.PoseStamped.header}}{unitree\_sdk2\_cpp.idl.ros2.PoseStamped.header}}\label{unitree-sdk2-cpp-idl-ros2-posestamped-header}}

底层 \passthrough{\lstinline!geometry\_msgs::msg::dds\_::PoseStamped\_!}
的 \passthrough{\lstinline!header!} 字段。类型签名只说明 Python
表示；单位、数值范围、数组固定长度和枚举含义需查对应 IDL 头文件。

\textbf{签名}

\begin{lstlisting}[language=Python]
@property
def header(self) -> Any

@header.setter
def header(self, value: Any) -> None
\end{lstlisting}

\textbf{可用性与安全性}

\textbf{\passthrough{\lstinline!AVAILABLE!} /
\passthrough{\lstinline!VALUE\_TYPE!}}：当前绑定源码已实现该消息值操作。

\textbf{参数}

\begin{longtable}[]{@{}
  >{\raggedright\arraybackslash}p{(\columnwidth - 4\tabcolsep) * \real{0.18}}
  >{\raggedright\arraybackslash}p{(\columnwidth - 4\tabcolsep) * \real{0.20}}
  >{\raggedright\arraybackslash}p{(\columnwidth - 4\tabcolsep) * \real{0.62}}@{}}
\toprule
名称 & Python 类型 & 含义 \\
\midrule
\endhead
\passthrough{\lstinline!value!} & \passthrough{\lstinline!Any!} &
要写入或传递的新值，必须符合该参数的 Python 类型以及底层 C++
范围约束。 \\
\bottomrule
\end{longtable}

\textbf{返回值}

\begin{longtable}[]{@{}
  >{\raggedright\arraybackslash}p{(\columnwidth - 4\tabcolsep) * \real{0.18}}
  >{\raggedright\arraybackslash}p{(\columnwidth - 4\tabcolsep) * \real{0.20}}
  >{\raggedright\arraybackslash}p{(\columnwidth - 4\tabcolsep) * \real{0.62}}@{}}
\toprule
位置 & 类型 & 含义 \\
\midrule
\endhead
getter 返回值 & \passthrough{\lstinline!Any!} & 返回属性当前值。 \\
\bottomrule
\end{longtable}

\textbf{对应 C++}

\begin{itemize}
\tightlist
\item
  类：\passthrough{\lstinline!geometry\_msgs::msg::dds\_::PoseStamped\_!}
\item
  字段类型：\passthrough{\lstinline!::std\_msgs::msg::dds\_::Header\_!}
\item
  声明位置：\passthrough{\lstinline!include/unitree/idl/ros2/PoseStamped\_.hpp:26!}
\end{itemize}

\textbf{用法}

\begin{lstlisting}[language=Python]
current_value = obj.header
obj.header = new_value
\end{lstlisting}

\hypertarget{unitree-sdk2-cpp-idl-ros2-posestamped-pose}{%
\paragraph{\texorpdfstring{\texttt{unitree\_sdk2\_cpp.idl.ros2.PoseStamped.pose}}{unitree\_sdk2\_cpp.idl.ros2.PoseStamped.pose}}\label{unitree-sdk2-cpp-idl-ros2-posestamped-pose}}

底层 \passthrough{\lstinline!geometry\_msgs::msg::dds\_::PoseStamped\_!}
的 \passthrough{\lstinline!pose!} 字段。类型签名只说明 Python
表示；单位、数值范围、数组固定长度和枚举含义需查对应 IDL 头文件。

\textbf{签名}

\begin{lstlisting}[language=Python]
@property
def pose(self) -> Pose

@pose.setter
def pose(self, value: Pose) -> None
\end{lstlisting}

\textbf{可用性与安全性}

\textbf{\passthrough{\lstinline!AVAILABLE!} /
\passthrough{\lstinline!VALUE\_TYPE!}}：当前绑定源码已实现该消息值操作。

\textbf{参数}

\begin{longtable}[]{@{}
  >{\raggedright\arraybackslash}p{(\columnwidth - 4\tabcolsep) * \real{0.18}}
  >{\raggedright\arraybackslash}p{(\columnwidth - 4\tabcolsep) * \real{0.20}}
  >{\raggedright\arraybackslash}p{(\columnwidth - 4\tabcolsep) * \real{0.62}}@{}}
\toprule
名称 & Python 类型 & 含义 \\
\midrule
\endhead
\passthrough{\lstinline!value!} & \passthrough{\lstinline!Pose!} &
要写入或传递的新值，必须符合该参数的 Python 类型以及底层 C++
范围约束。 \\
\bottomrule
\end{longtable}

\textbf{返回值}

\begin{longtable}[]{@{}
  >{\raggedright\arraybackslash}p{(\columnwidth - 4\tabcolsep) * \real{0.18}}
  >{\raggedright\arraybackslash}p{(\columnwidth - 4\tabcolsep) * \real{0.20}}
  >{\raggedright\arraybackslash}p{(\columnwidth - 4\tabcolsep) * \real{0.62}}@{}}
\toprule
位置 & 类型 & 含义 \\
\midrule
\endhead
getter 返回值 & \passthrough{\lstinline!Pose!} &
返回属性当前值。数组、vector 和嵌套消息采用复制语义。 \\
\bottomrule
\end{longtable}

\textbf{对应 C++}

\begin{itemize}
\tightlist
\item
  类：\passthrough{\lstinline!geometry\_msgs::msg::dds\_::PoseStamped\_!}
\item
  字段类型：\passthrough{\lstinline!::geometry\_msgs::msg::dds\_::Pose\_!}
\item
  声明位置：\passthrough{\lstinline!include/unitree/idl/ros2/PoseStamped\_.hpp:27!}
\end{itemize}

\textbf{用法}

\begin{lstlisting}[language=Python]
current_value = obj.pose
obj.pose = new_value
\end{lstlisting}

\hypertarget{unitree-sdk2-cpp-idl-ros2-posewithcovariancestamped}{%
\subsubsection{\texorpdfstring{\texttt{unitree\_sdk2\_cpp.idl.ros2.PoseWithCovarianceStamped}}{unitree\_sdk2\_cpp.idl.ros2.PoseWithCovarianceStamped}}\label{unitree-sdk2-cpp-idl-ros2-posewithcovariancestamped}}

可由 Python 读写的 DDS/IDL 消息值类型。字段采用复制语义。

\textbf{导入}

\begin{lstlisting}[language=Python]
from unitree_sdk2_cpp.idl.ros2 import PoseWithCovarianceStamped
\end{lstlisting}

\textbf{方法索引}

\begin{longtable}[]{@{}
  >{\raggedright\arraybackslash}p{(\columnwidth - 6\tabcolsep) * \real{0.18}}
  >{\raggedright\arraybackslash}p{(\columnwidth - 6\tabcolsep) * \real{0.24}}
  >{\raggedright\arraybackslash}p{(\columnwidth - 6\tabcolsep) * \real{0.18}}
  >{\raggedright\arraybackslash}p{(\columnwidth - 6\tabcolsep) * \real{0.40}}@{}}
\toprule
方法 & Python 签名 & 状态 & 安全分类 \\
\midrule
\endhead
\protect\hyperlink{unitree-sdk2-cpp-idl-ros2-posewithcovariancestamped-dunder-init-1}{\passthrough{\lstinline!\_\_init\_\_!}}
& \passthrough{\lstinline!def \_\_init\_\_(self) -> None!} &
\passthrough{\lstinline!AVAILABLE!} &
\passthrough{\lstinline!VALUE\_TYPE!} \\
\protect\hyperlink{unitree-sdk2-cpp-idl-ros2-posewithcovariancestamped-dunder-eq-1}{\passthrough{\lstinline!\_\_eq\_\_!}}
& \passthrough{\lstinline!def \_\_eq\_\_(self, other: object) -> bool!}
& \passthrough{\lstinline!AVAILABLE!} &
\passthrough{\lstinline!VALUE\_TYPE!} \\
\protect\hyperlink{unitree-sdk2-cpp-idl-ros2-posewithcovariancestamped-dunder-ne-1}{\passthrough{\lstinline!\_\_ne\_\_!}}
& \passthrough{\lstinline!def \_\_ne\_\_(self, other: object) -> bool!}
& \passthrough{\lstinline!AVAILABLE!} &
\passthrough{\lstinline!VALUE\_TYPE!} \\
\bottomrule
\end{longtable}

\textbf{属性索引}

\begin{longtable}[]{@{}
  >{\raggedright\arraybackslash}p{(\columnwidth - 6\tabcolsep) * \real{0.18}}
  >{\raggedright\arraybackslash}p{(\columnwidth - 6\tabcolsep) * \real{0.24}}
  >{\raggedright\arraybackslash}p{(\columnwidth - 6\tabcolsep) * \real{0.18}}
  >{\raggedright\arraybackslash}p{(\columnwidth - 6\tabcolsep) * \real{0.40}}@{}}
\toprule
属性 & 读取类型 & 写入类型 & 状态 \\
\midrule
\endhead
\protect\hyperlink{unitree-sdk2-cpp-idl-ros2-posewithcovariancestamped-header}{\passthrough{\lstinline!header!}}
& \passthrough{\lstinline!Any!} & \passthrough{\lstinline!Any!} &
\passthrough{\lstinline!AVAILABLE!} \\
\protect\hyperlink{unitree-sdk2-cpp-idl-ros2-posewithcovariancestamped-pose}{\passthrough{\lstinline!pose!}}
& \passthrough{\lstinline!PoseWithCovariance!} &
\passthrough{\lstinline!PoseWithCovariance!} &
\passthrough{\lstinline!AVAILABLE!} \\
\bottomrule
\end{longtable}

\hypertarget{unitree-sdk2-cpp-idl-ros2-posewithcovariancestamped-dunder-init-1}{%
\paragraph{\texorpdfstring{\texttt{unitree\_sdk2\_cpp.idl.ros2.PoseWithCovarianceStamped.\_\_init\_\_}}{unitree\_sdk2\_cpp.idl.ros2.PoseWithCovarianceStamped.\_\_init\_\_}}\label{unitree-sdk2-cpp-idl-ros2-posewithcovariancestamped-dunder-init-1}}

初始化 \passthrough{\lstinline!PoseWithCovarianceStamped!}
实例。是否能实际构造取决于下方可用性状态。

\textbf{签名}

\begin{lstlisting}[language=Python]
def __init__(self) -> None
\end{lstlisting}

\textbf{可用性与安全性}

\textbf{\passthrough{\lstinline!AVAILABLE!} /
\passthrough{\lstinline!VALUE\_TYPE!}}：当前绑定源码已实现该消息值操作。

\textbf{参数}

无显式参数。实例方法中的 \passthrough{\lstinline!self!} 由 Python
自动传入。

\textbf{返回值}

\begin{longtable}[]{@{}
  >{\raggedright\arraybackslash}p{(\columnwidth - 4\tabcolsep) * \real{0.18}}
  >{\raggedright\arraybackslash}p{(\columnwidth - 4\tabcolsep) * \real{0.20}}
  >{\raggedright\arraybackslash}p{(\columnwidth - 4\tabcolsep) * \real{0.62}}@{}}
\toprule
位置 & 类型 & 含义 \\
\midrule
\endhead
无 & \passthrough{\lstinline!None!} &
\passthrough{\lstinline!\_\_init\_\_!}
本身不返回值；调用类对象时，在构造函数可用的前提下得到该类实例。 \\
\bottomrule
\end{longtable}

\textbf{对应 C++}

\begin{itemize}
\tightlist
\item
  类：\passthrough{\lstinline!geometry\_msgs::msg::dds\_::PoseWithCovarianceStamped\_!}
\item
  签名：\passthrough{\lstinline!PoseWithCovarianceStamped\_()!}
\item
  绑定策略：\passthrough{\lstinline!未单独标注!}
\item
  声明位置：\passthrough{\lstinline!include/unitree/idl/ros2/PoseWithCovarianceStamped\_.hpp:30!}
\end{itemize}

\textbf{说明}

\begin{itemize}
\tightlist
\item
  示例是局部调用片段；\passthrough{\lstinline!obj!}
  和参数变量表示已经按上层业务校验并准备好的对象。
\end{itemize}

\textbf{用法}

\begin{lstlisting}[language=Python]
value = PoseWithCovarianceStamped()
\end{lstlisting}

\hypertarget{unitree-sdk2-cpp-idl-ros2-posewithcovariancestamped-dunder-eq-1}{%
\paragraph{\texorpdfstring{\texttt{unitree\_sdk2\_cpp.idl.ros2.PoseWithCovarianceStamped.\_\_eq\_\_}}{unitree\_sdk2\_cpp.idl.ros2.PoseWithCovarianceStamped.\_\_eq\_\_}}\label{unitree-sdk2-cpp-idl-ros2-posewithcovariancestamped-dunder-eq-1}}

按底层 IDL 消息内容比较两个对象是否相等。

\textbf{签名}

\begin{lstlisting}[language=Python]
def __eq__(self, other: object) -> bool
\end{lstlisting}

\textbf{可用性与安全性}

\textbf{\passthrough{\lstinline!AVAILABLE!} /
\passthrough{\lstinline!VALUE\_TYPE!}}：当前绑定源码已实现该消息值操作。

\textbf{参数}

\begin{longtable}[]{@{}
  >{\raggedright\arraybackslash}p{(\columnwidth - 6\tabcolsep) * \real{0.18}}
  >{\raggedright\arraybackslash}p{(\columnwidth - 6\tabcolsep) * \real{0.24}}
  >{\raggedright\arraybackslash}p{(\columnwidth - 6\tabcolsep) * \real{0.18}}
  >{\raggedright\arraybackslash}p{(\columnwidth - 6\tabcolsep) * \real{0.40}}@{}}
\toprule
名称 & Python 类型 & 默认值 & 含义 \\
\midrule
\endhead
\passthrough{\lstinline!other!} & \passthrough{\lstinline!object!} &
必填 & 用于比较的另一个 Python 对象。类型不兼容时比较结果通常为
\passthrough{\lstinline!False!}。 \\
\bottomrule
\end{longtable}

\textbf{返回值}

\begin{longtable}[]{@{}
  >{\raggedright\arraybackslash}p{(\columnwidth - 4\tabcolsep) * \real{0.18}}
  >{\raggedright\arraybackslash}p{(\columnwidth - 4\tabcolsep) * \real{0.20}}
  >{\raggedright\arraybackslash}p{(\columnwidth - 4\tabcolsep) * \real{0.62}}@{}}
\toprule
位置 & 类型 & 含义 \\
\midrule
\endhead
返回值 & \passthrough{\lstinline!bool!} & 返回
\passthrough{\lstinline!bool!}。更精确的业务含义见该方法用途、C++
签名和目标型号协议。 \\
\bottomrule
\end{longtable}

\textbf{对应 C++}

\begin{itemize}
\tightlist
\item
  类：\passthrough{\lstinline!geometry\_msgs::msg::dds\_::PoseWithCovarianceStamped\_!}
\item
  签名：\passthrough{\lstinline!operator==(const value \&) const!}
\item
  绑定策略：\passthrough{\lstinline!未单独标注!}
\item
  声明位置：由 pybind11 手工绑定或生成注册表提供；manifest
  未记录独立头文件行号。
\end{itemize}

\textbf{说明}

\begin{itemize}
\tightlist
\item
  示例是局部调用片段；\passthrough{\lstinline!obj!}
  和参数变量表示已经按上层业务校验并准备好的对象。
\end{itemize}

\textbf{用法}

\begin{lstlisting}[language=Python]
same = left == right
\end{lstlisting}

\hypertarget{unitree-sdk2-cpp-idl-ros2-posewithcovariancestamped-dunder-ne-1}{%
\paragraph{\texorpdfstring{\texttt{unitree\_sdk2\_cpp.idl.ros2.PoseWithCovarianceStamped.\_\_ne\_\_}}{unitree\_sdk2\_cpp.idl.ros2.PoseWithCovarianceStamped.\_\_ne\_\_}}\label{unitree-sdk2-cpp-idl-ros2-posewithcovariancestamped-dunder-ne-1}}

按底层 IDL 消息内容比较两个对象是否不相等。

\textbf{签名}

\begin{lstlisting}[language=Python]
def __ne__(self, other: object) -> bool
\end{lstlisting}

\textbf{可用性与安全性}

\textbf{\passthrough{\lstinline!AVAILABLE!} /
\passthrough{\lstinline!VALUE\_TYPE!}}：当前绑定源码已实现该消息值操作。

\textbf{参数}

\begin{longtable}[]{@{}
  >{\raggedright\arraybackslash}p{(\columnwidth - 6\tabcolsep) * \real{0.18}}
  >{\raggedright\arraybackslash}p{(\columnwidth - 6\tabcolsep) * \real{0.24}}
  >{\raggedright\arraybackslash}p{(\columnwidth - 6\tabcolsep) * \real{0.18}}
  >{\raggedright\arraybackslash}p{(\columnwidth - 6\tabcolsep) * \real{0.40}}@{}}
\toprule
名称 & Python 类型 & 默认值 & 含义 \\
\midrule
\endhead
\passthrough{\lstinline!other!} & \passthrough{\lstinline!object!} &
必填 & 用于比较的另一个 Python 对象。类型不兼容时比较结果通常为
\passthrough{\lstinline!False!}。 \\
\bottomrule
\end{longtable}

\textbf{返回值}

\begin{longtable}[]{@{}
  >{\raggedright\arraybackslash}p{(\columnwidth - 4\tabcolsep) * \real{0.18}}
  >{\raggedright\arraybackslash}p{(\columnwidth - 4\tabcolsep) * \real{0.20}}
  >{\raggedright\arraybackslash}p{(\columnwidth - 4\tabcolsep) * \real{0.62}}@{}}
\toprule
位置 & 类型 & 含义 \\
\midrule
\endhead
返回值 & \passthrough{\lstinline!bool!} & 返回
\passthrough{\lstinline!bool!}。更精确的业务含义见该方法用途、C++
签名和目标型号协议。 \\
\bottomrule
\end{longtable}

\textbf{对应 C++}

\begin{itemize}
\tightlist
\item
  类：\passthrough{\lstinline!geometry\_msgs::msg::dds\_::PoseWithCovarianceStamped\_!}
\item
  签名：\passthrough{\lstinline"operator!=(const value \&) const"}
\item
  绑定策略：\passthrough{\lstinline!未单独标注!}
\item
  声明位置：由 pybind11 手工绑定或生成注册表提供；manifest
  未记录独立头文件行号。
\end{itemize}

\textbf{说明}

\begin{itemize}
\tightlist
\item
  示例是局部调用片段；\passthrough{\lstinline!obj!}
  和参数变量表示已经按上层业务校验并准备好的对象。
\end{itemize}

\textbf{用法}

\begin{lstlisting}[language=Python]
different = left != right
\end{lstlisting}

\hypertarget{unitree-sdk2-cpp-idl-ros2-posewithcovariancestamped-header}{%
\paragraph{\texorpdfstring{\texttt{unitree\_sdk2\_cpp.idl.ros2.PoseWithCovarianceStamped.header}}{unitree\_sdk2\_cpp.idl.ros2.PoseWithCovarianceStamped.header}}\label{unitree-sdk2-cpp-idl-ros2-posewithcovariancestamped-header}}

底层
\passthrough{\lstinline!geometry\_msgs::msg::dds\_::PoseWithCovarianceStamped\_!}
的 \passthrough{\lstinline!header!} 字段。类型签名只说明 Python
表示；单位、数值范围、数组固定长度和枚举含义需查对应 IDL 头文件。

\textbf{签名}

\begin{lstlisting}[language=Python]
@property
def header(self) -> Any

@header.setter
def header(self, value: Any) -> None
\end{lstlisting}

\textbf{可用性与安全性}

\textbf{\passthrough{\lstinline!AVAILABLE!} /
\passthrough{\lstinline!VALUE\_TYPE!}}：当前绑定源码已实现该消息值操作。

\textbf{参数}

\begin{longtable}[]{@{}
  >{\raggedright\arraybackslash}p{(\columnwidth - 4\tabcolsep) * \real{0.18}}
  >{\raggedright\arraybackslash}p{(\columnwidth - 4\tabcolsep) * \real{0.20}}
  >{\raggedright\arraybackslash}p{(\columnwidth - 4\tabcolsep) * \real{0.62}}@{}}
\toprule
名称 & Python 类型 & 含义 \\
\midrule
\endhead
\passthrough{\lstinline!value!} & \passthrough{\lstinline!Any!} &
要写入或传递的新值，必须符合该参数的 Python 类型以及底层 C++
范围约束。 \\
\bottomrule
\end{longtable}

\textbf{返回值}

\begin{longtable}[]{@{}
  >{\raggedright\arraybackslash}p{(\columnwidth - 4\tabcolsep) * \real{0.18}}
  >{\raggedright\arraybackslash}p{(\columnwidth - 4\tabcolsep) * \real{0.20}}
  >{\raggedright\arraybackslash}p{(\columnwidth - 4\tabcolsep) * \real{0.62}}@{}}
\toprule
位置 & 类型 & 含义 \\
\midrule
\endhead
getter 返回值 & \passthrough{\lstinline!Any!} & 返回属性当前值。 \\
\bottomrule
\end{longtable}

\textbf{对应 C++}

\begin{itemize}
\tightlist
\item
  类：\passthrough{\lstinline!geometry\_msgs::msg::dds\_::PoseWithCovarianceStamped\_!}
\item
  字段类型：\passthrough{\lstinline!::std\_msgs::msg::dds\_::Header\_!}
\item
  声明位置：\passthrough{\lstinline!include/unitree/idl/ros2/PoseWithCovarianceStamped\_.hpp:26!}
\end{itemize}

\textbf{用法}

\begin{lstlisting}[language=Python]
current_value = obj.header
obj.header = new_value
\end{lstlisting}

\hypertarget{unitree-sdk2-cpp-idl-ros2-posewithcovariancestamped-pose}{%
\paragraph{\texorpdfstring{\texttt{unitree\_sdk2\_cpp.idl.ros2.PoseWithCovarianceStamped.pose}}{unitree\_sdk2\_cpp.idl.ros2.PoseWithCovarianceStamped.pose}}\label{unitree-sdk2-cpp-idl-ros2-posewithcovariancestamped-pose}}

底层
\passthrough{\lstinline!geometry\_msgs::msg::dds\_::PoseWithCovarianceStamped\_!}
的 \passthrough{\lstinline!pose!} 字段。类型签名只说明 Python
表示；单位、数值范围、数组固定长度和枚举含义需查对应 IDL 头文件。

\textbf{签名}

\begin{lstlisting}[language=Python]
@property
def pose(self) -> PoseWithCovariance

@pose.setter
def pose(self, value: PoseWithCovariance) -> None
\end{lstlisting}

\textbf{可用性与安全性}

\textbf{\passthrough{\lstinline!AVAILABLE!} /
\passthrough{\lstinline!VALUE\_TYPE!}}：当前绑定源码已实现该消息值操作。

\textbf{参数}

\begin{longtable}[]{@{}
  >{\raggedright\arraybackslash}p{(\columnwidth - 4\tabcolsep) * \real{0.18}}
  >{\raggedright\arraybackslash}p{(\columnwidth - 4\tabcolsep) * \real{0.20}}
  >{\raggedright\arraybackslash}p{(\columnwidth - 4\tabcolsep) * \real{0.62}}@{}}
\toprule
名称 & Python 类型 & 含义 \\
\midrule
\endhead
\passthrough{\lstinline!value!} &
\passthrough{\lstinline!PoseWithCovariance!} &
要写入或传递的新值，必须符合该参数的 Python 类型以及底层 C++
范围约束。 \\
\bottomrule
\end{longtable}

\textbf{返回值}

\begin{longtable}[]{@{}
  >{\raggedright\arraybackslash}p{(\columnwidth - 4\tabcolsep) * \real{0.18}}
  >{\raggedright\arraybackslash}p{(\columnwidth - 4\tabcolsep) * \real{0.20}}
  >{\raggedright\arraybackslash}p{(\columnwidth - 4\tabcolsep) * \real{0.62}}@{}}
\toprule
位置 & 类型 & 含义 \\
\midrule
\endhead
getter 返回值 & \passthrough{\lstinline!PoseWithCovariance!} &
返回属性当前值。数组、vector 和嵌套消息采用复制语义。 \\
\bottomrule
\end{longtable}

\textbf{对应 C++}

\begin{itemize}
\tightlist
\item
  类：\passthrough{\lstinline!geometry\_msgs::msg::dds\_::PoseWithCovarianceStamped\_!}
\item
  字段类型：\passthrough{\lstinline!::geometry\_msgs::msg::dds\_::PoseWithCovariance\_!}
\item
  声明位置：\passthrough{\lstinline!include/unitree/idl/ros2/PoseWithCovarianceStamped\_.hpp:27!}
\end{itemize}

\textbf{用法}

\begin{lstlisting}[language=Python]
current_value = obj.pose
obj.pose = new_value
\end{lstlisting}

\hypertarget{unitree-sdk2-cpp-idl-ros2-quaternionstamped}{%
\subsubsection{\texorpdfstring{\texttt{unitree\_sdk2\_cpp.idl.ros2.QuaternionStamped}}{unitree\_sdk2\_cpp.idl.ros2.QuaternionStamped}}\label{unitree-sdk2-cpp-idl-ros2-quaternionstamped}}

可由 Python 读写的 DDS/IDL 消息值类型。字段采用复制语义。

\textbf{导入}

\begin{lstlisting}[language=Python]
from unitree_sdk2_cpp.idl.ros2 import QuaternionStamped
\end{lstlisting}

\textbf{方法索引}

\begin{longtable}[]{@{}
  >{\raggedright\arraybackslash}p{(\columnwidth - 6\tabcolsep) * \real{0.18}}
  >{\raggedright\arraybackslash}p{(\columnwidth - 6\tabcolsep) * \real{0.24}}
  >{\raggedright\arraybackslash}p{(\columnwidth - 6\tabcolsep) * \real{0.18}}
  >{\raggedright\arraybackslash}p{(\columnwidth - 6\tabcolsep) * \real{0.40}}@{}}
\toprule
方法 & Python 签名 & 状态 & 安全分类 \\
\midrule
\endhead
\protect\hyperlink{unitree-sdk2-cpp-idl-ros2-quaternionstamped-dunder-init-1}{\passthrough{\lstinline!\_\_init\_\_!}}
& \passthrough{\lstinline!def \_\_init\_\_(self) -> None!} &
\passthrough{\lstinline!AVAILABLE!} &
\passthrough{\lstinline!VALUE\_TYPE!} \\
\protect\hyperlink{unitree-sdk2-cpp-idl-ros2-quaternionstamped-dunder-eq-1}{\passthrough{\lstinline!\_\_eq\_\_!}}
& \passthrough{\lstinline!def \_\_eq\_\_(self, other: object) -> bool!}
& \passthrough{\lstinline!AVAILABLE!} &
\passthrough{\lstinline!VALUE\_TYPE!} \\
\protect\hyperlink{unitree-sdk2-cpp-idl-ros2-quaternionstamped-dunder-ne-1}{\passthrough{\lstinline!\_\_ne\_\_!}}
& \passthrough{\lstinline!def \_\_ne\_\_(self, other: object) -> bool!}
& \passthrough{\lstinline!AVAILABLE!} &
\passthrough{\lstinline!VALUE\_TYPE!} \\
\bottomrule
\end{longtable}

\textbf{属性索引}

\begin{longtable}[]{@{}
  >{\raggedright\arraybackslash}p{(\columnwidth - 6\tabcolsep) * \real{0.18}}
  >{\raggedright\arraybackslash}p{(\columnwidth - 6\tabcolsep) * \real{0.24}}
  >{\raggedright\arraybackslash}p{(\columnwidth - 6\tabcolsep) * \real{0.18}}
  >{\raggedright\arraybackslash}p{(\columnwidth - 6\tabcolsep) * \real{0.40}}@{}}
\toprule
属性 & 读取类型 & 写入类型 & 状态 \\
\midrule
\endhead
\protect\hyperlink{unitree-sdk2-cpp-idl-ros2-quaternionstamped-header}{\passthrough{\lstinline!header!}}
& \passthrough{\lstinline!Any!} & \passthrough{\lstinline!Any!} &
\passthrough{\lstinline!AVAILABLE!} \\
\protect\hyperlink{unitree-sdk2-cpp-idl-ros2-quaternionstamped-quaternion}{\passthrough{\lstinline!quaternion!}}
& \passthrough{\lstinline!Quaternion!} &
\passthrough{\lstinline!Quaternion!} &
\passthrough{\lstinline!AVAILABLE!} \\
\bottomrule
\end{longtable}

\hypertarget{unitree-sdk2-cpp-idl-ros2-quaternionstamped-dunder-init-1}{%
\paragraph{\texorpdfstring{\texttt{unitree\_sdk2\_cpp.idl.ros2.QuaternionStamped.\_\_init\_\_}}{unitree\_sdk2\_cpp.idl.ros2.QuaternionStamped.\_\_init\_\_}}\label{unitree-sdk2-cpp-idl-ros2-quaternionstamped-dunder-init-1}}

初始化 \passthrough{\lstinline!QuaternionStamped!}
实例。是否能实际构造取决于下方可用性状态。

\textbf{签名}

\begin{lstlisting}[language=Python]
def __init__(self) -> None
\end{lstlisting}

\textbf{可用性与安全性}

\textbf{\passthrough{\lstinline!AVAILABLE!} /
\passthrough{\lstinline!VALUE\_TYPE!}}：当前绑定源码已实现该消息值操作。

\textbf{参数}

无显式参数。实例方法中的 \passthrough{\lstinline!self!} 由 Python
自动传入。

\textbf{返回值}

\begin{longtable}[]{@{}
  >{\raggedright\arraybackslash}p{(\columnwidth - 4\tabcolsep) * \real{0.18}}
  >{\raggedright\arraybackslash}p{(\columnwidth - 4\tabcolsep) * \real{0.20}}
  >{\raggedright\arraybackslash}p{(\columnwidth - 4\tabcolsep) * \real{0.62}}@{}}
\toprule
位置 & 类型 & 含义 \\
\midrule
\endhead
无 & \passthrough{\lstinline!None!} &
\passthrough{\lstinline!\_\_init\_\_!}
本身不返回值；调用类对象时，在构造函数可用的前提下得到该类实例。 \\
\bottomrule
\end{longtable}

\textbf{对应 C++}

\begin{itemize}
\tightlist
\item
  类：\passthrough{\lstinline!geometry\_msgs::msg::dds\_::QuaternionStamped\_!}
\item
  签名：\passthrough{\lstinline!QuaternionStamped\_()!}
\item
  绑定策略：\passthrough{\lstinline!未单独标注!}
\item
  声明位置：\passthrough{\lstinline!include/unitree/idl/ros2/QuaternionStamped\_.hpp:30!}
\end{itemize}

\textbf{说明}

\begin{itemize}
\tightlist
\item
  示例是局部调用片段；\passthrough{\lstinline!obj!}
  和参数变量表示已经按上层业务校验并准备好的对象。
\end{itemize}

\textbf{用法}

\begin{lstlisting}[language=Python]
value = QuaternionStamped()
\end{lstlisting}

\hypertarget{unitree-sdk2-cpp-idl-ros2-quaternionstamped-dunder-eq-1}{%
\paragraph{\texorpdfstring{\texttt{unitree\_sdk2\_cpp.idl.ros2.QuaternionStamped.\_\_eq\_\_}}{unitree\_sdk2\_cpp.idl.ros2.QuaternionStamped.\_\_eq\_\_}}\label{unitree-sdk2-cpp-idl-ros2-quaternionstamped-dunder-eq-1}}

按底层 IDL 消息内容比较两个对象是否相等。

\textbf{签名}

\begin{lstlisting}[language=Python]
def __eq__(self, other: object) -> bool
\end{lstlisting}

\textbf{可用性与安全性}

\textbf{\passthrough{\lstinline!AVAILABLE!} /
\passthrough{\lstinline!VALUE\_TYPE!}}：当前绑定源码已实现该消息值操作。

\textbf{参数}

\begin{longtable}[]{@{}
  >{\raggedright\arraybackslash}p{(\columnwidth - 6\tabcolsep) * \real{0.18}}
  >{\raggedright\arraybackslash}p{(\columnwidth - 6\tabcolsep) * \real{0.24}}
  >{\raggedright\arraybackslash}p{(\columnwidth - 6\tabcolsep) * \real{0.18}}
  >{\raggedright\arraybackslash}p{(\columnwidth - 6\tabcolsep) * \real{0.40}}@{}}
\toprule
名称 & Python 类型 & 默认值 & 含义 \\
\midrule
\endhead
\passthrough{\lstinline!other!} & \passthrough{\lstinline!object!} &
必填 & 用于比较的另一个 Python 对象。类型不兼容时比较结果通常为
\passthrough{\lstinline!False!}。 \\
\bottomrule
\end{longtable}

\textbf{返回值}

\begin{longtable}[]{@{}
  >{\raggedright\arraybackslash}p{(\columnwidth - 4\tabcolsep) * \real{0.18}}
  >{\raggedright\arraybackslash}p{(\columnwidth - 4\tabcolsep) * \real{0.20}}
  >{\raggedright\arraybackslash}p{(\columnwidth - 4\tabcolsep) * \real{0.62}}@{}}
\toprule
位置 & 类型 & 含义 \\
\midrule
\endhead
返回值 & \passthrough{\lstinline!bool!} & 返回
\passthrough{\lstinline!bool!}。更精确的业务含义见该方法用途、C++
签名和目标型号协议。 \\
\bottomrule
\end{longtable}

\textbf{对应 C++}

\begin{itemize}
\tightlist
\item
  类：\passthrough{\lstinline!geometry\_msgs::msg::dds\_::QuaternionStamped\_!}
\item
  签名：\passthrough{\lstinline!operator==(const value \&) const!}
\item
  绑定策略：\passthrough{\lstinline!未单独标注!}
\item
  声明位置：由 pybind11 手工绑定或生成注册表提供；manifest
  未记录独立头文件行号。
\end{itemize}

\textbf{说明}

\begin{itemize}
\tightlist
\item
  示例是局部调用片段；\passthrough{\lstinline!obj!}
  和参数变量表示已经按上层业务校验并准备好的对象。
\end{itemize}

\textbf{用法}

\begin{lstlisting}[language=Python]
same = left == right
\end{lstlisting}

\hypertarget{unitree-sdk2-cpp-idl-ros2-quaternionstamped-dunder-ne-1}{%
\paragraph{\texorpdfstring{\texttt{unitree\_sdk2\_cpp.idl.ros2.QuaternionStamped.\_\_ne\_\_}}{unitree\_sdk2\_cpp.idl.ros2.QuaternionStamped.\_\_ne\_\_}}\label{unitree-sdk2-cpp-idl-ros2-quaternionstamped-dunder-ne-1}}

按底层 IDL 消息内容比较两个对象是否不相等。

\textbf{签名}

\begin{lstlisting}[language=Python]
def __ne__(self, other: object) -> bool
\end{lstlisting}

\textbf{可用性与安全性}

\textbf{\passthrough{\lstinline!AVAILABLE!} /
\passthrough{\lstinline!VALUE\_TYPE!}}：当前绑定源码已实现该消息值操作。

\textbf{参数}

\begin{longtable}[]{@{}
  >{\raggedright\arraybackslash}p{(\columnwidth - 6\tabcolsep) * \real{0.18}}
  >{\raggedright\arraybackslash}p{(\columnwidth - 6\tabcolsep) * \real{0.24}}
  >{\raggedright\arraybackslash}p{(\columnwidth - 6\tabcolsep) * \real{0.18}}
  >{\raggedright\arraybackslash}p{(\columnwidth - 6\tabcolsep) * \real{0.40}}@{}}
\toprule
名称 & Python 类型 & 默认值 & 含义 \\
\midrule
\endhead
\passthrough{\lstinline!other!} & \passthrough{\lstinline!object!} &
必填 & 用于比较的另一个 Python 对象。类型不兼容时比较结果通常为
\passthrough{\lstinline!False!}。 \\
\bottomrule
\end{longtable}

\textbf{返回值}

\begin{longtable}[]{@{}
  >{\raggedright\arraybackslash}p{(\columnwidth - 4\tabcolsep) * \real{0.18}}
  >{\raggedright\arraybackslash}p{(\columnwidth - 4\tabcolsep) * \real{0.20}}
  >{\raggedright\arraybackslash}p{(\columnwidth - 4\tabcolsep) * \real{0.62}}@{}}
\toprule
位置 & 类型 & 含义 \\
\midrule
\endhead
返回值 & \passthrough{\lstinline!bool!} & 返回
\passthrough{\lstinline!bool!}。更精确的业务含义见该方法用途、C++
签名和目标型号协议。 \\
\bottomrule
\end{longtable}

\textbf{对应 C++}

\begin{itemize}
\tightlist
\item
  类：\passthrough{\lstinline!geometry\_msgs::msg::dds\_::QuaternionStamped\_!}
\item
  签名：\passthrough{\lstinline"operator!=(const value \&) const"}
\item
  绑定策略：\passthrough{\lstinline!未单独标注!}
\item
  声明位置：由 pybind11 手工绑定或生成注册表提供；manifest
  未记录独立头文件行号。
\end{itemize}

\textbf{说明}

\begin{itemize}
\tightlist
\item
  示例是局部调用片段；\passthrough{\lstinline!obj!}
  和参数变量表示已经按上层业务校验并准备好的对象。
\end{itemize}

\textbf{用法}

\begin{lstlisting}[language=Python]
different = left != right
\end{lstlisting}

\hypertarget{unitree-sdk2-cpp-idl-ros2-quaternionstamped-header}{%
\paragraph{\texorpdfstring{\texttt{unitree\_sdk2\_cpp.idl.ros2.QuaternionStamped.header}}{unitree\_sdk2\_cpp.idl.ros2.QuaternionStamped.header}}\label{unitree-sdk2-cpp-idl-ros2-quaternionstamped-header}}

底层
\passthrough{\lstinline!geometry\_msgs::msg::dds\_::QuaternionStamped\_!}
的 \passthrough{\lstinline!header!} 字段。类型签名只说明 Python
表示；单位、数值范围、数组固定长度和枚举含义需查对应 IDL 头文件。

\textbf{签名}

\begin{lstlisting}[language=Python]
@property
def header(self) -> Any

@header.setter
def header(self, value: Any) -> None
\end{lstlisting}

\textbf{可用性与安全性}

\textbf{\passthrough{\lstinline!AVAILABLE!} /
\passthrough{\lstinline!VALUE\_TYPE!}}：当前绑定源码已实现该消息值操作。

\textbf{参数}

\begin{longtable}[]{@{}
  >{\raggedright\arraybackslash}p{(\columnwidth - 4\tabcolsep) * \real{0.18}}
  >{\raggedright\arraybackslash}p{(\columnwidth - 4\tabcolsep) * \real{0.20}}
  >{\raggedright\arraybackslash}p{(\columnwidth - 4\tabcolsep) * \real{0.62}}@{}}
\toprule
名称 & Python 类型 & 含义 \\
\midrule
\endhead
\passthrough{\lstinline!value!} & \passthrough{\lstinline!Any!} &
要写入或传递的新值，必须符合该参数的 Python 类型以及底层 C++
范围约束。 \\
\bottomrule
\end{longtable}

\textbf{返回值}

\begin{longtable}[]{@{}
  >{\raggedright\arraybackslash}p{(\columnwidth - 4\tabcolsep) * \real{0.18}}
  >{\raggedright\arraybackslash}p{(\columnwidth - 4\tabcolsep) * \real{0.20}}
  >{\raggedright\arraybackslash}p{(\columnwidth - 4\tabcolsep) * \real{0.62}}@{}}
\toprule
位置 & 类型 & 含义 \\
\midrule
\endhead
getter 返回值 & \passthrough{\lstinline!Any!} & 返回属性当前值。 \\
\bottomrule
\end{longtable}

\textbf{对应 C++}

\begin{itemize}
\tightlist
\item
  类：\passthrough{\lstinline!geometry\_msgs::msg::dds\_::QuaternionStamped\_!}
\item
  字段类型：\passthrough{\lstinline!::std\_msgs::msg::dds\_::Header\_!}
\item
  声明位置：\passthrough{\lstinline!include/unitree/idl/ros2/QuaternionStamped\_.hpp:26!}
\end{itemize}

\textbf{用法}

\begin{lstlisting}[language=Python]
current_value = obj.header
obj.header = new_value
\end{lstlisting}

\hypertarget{unitree-sdk2-cpp-idl-ros2-quaternionstamped-quaternion}{%
\paragraph{\texorpdfstring{\texttt{unitree\_sdk2\_cpp.idl.ros2.QuaternionStamped.quaternion}}{unitree\_sdk2\_cpp.idl.ros2.QuaternionStamped.quaternion}}\label{unitree-sdk2-cpp-idl-ros2-quaternionstamped-quaternion}}

底层
\passthrough{\lstinline!geometry\_msgs::msg::dds\_::QuaternionStamped\_!}
的 \passthrough{\lstinline!quaternion!} 字段。类型签名只说明 Python
表示；单位、数值范围、数组固定长度和枚举含义需查对应 IDL 头文件。

\textbf{签名}

\begin{lstlisting}[language=Python]
@property
def quaternion(self) -> Quaternion

@quaternion.setter
def quaternion(self, value: Quaternion) -> None
\end{lstlisting}

\textbf{可用性与安全性}

\textbf{\passthrough{\lstinline!AVAILABLE!} /
\passthrough{\lstinline!VALUE\_TYPE!}}：当前绑定源码已实现该消息值操作。

\textbf{参数}

\begin{longtable}[]{@{}
  >{\raggedright\arraybackslash}p{(\columnwidth - 4\tabcolsep) * \real{0.18}}
  >{\raggedright\arraybackslash}p{(\columnwidth - 4\tabcolsep) * \real{0.20}}
  >{\raggedright\arraybackslash}p{(\columnwidth - 4\tabcolsep) * \real{0.62}}@{}}
\toprule
名称 & Python 类型 & 含义 \\
\midrule
\endhead
\passthrough{\lstinline!value!} & \passthrough{\lstinline!Quaternion!} &
要写入或传递的新值，必须符合该参数的 Python 类型以及底层 C++
范围约束。 \\
\bottomrule
\end{longtable}

\textbf{返回值}

\begin{longtable}[]{@{}
  >{\raggedright\arraybackslash}p{(\columnwidth - 4\tabcolsep) * \real{0.18}}
  >{\raggedright\arraybackslash}p{(\columnwidth - 4\tabcolsep) * \real{0.20}}
  >{\raggedright\arraybackslash}p{(\columnwidth - 4\tabcolsep) * \real{0.62}}@{}}
\toprule
位置 & 类型 & 含义 \\
\midrule
\endhead
getter 返回值 & \passthrough{\lstinline!Quaternion!} &
返回属性当前值。数组、vector 和嵌套消息采用复制语义。 \\
\bottomrule
\end{longtable}

\textbf{对应 C++}

\begin{itemize}
\tightlist
\item
  类：\passthrough{\lstinline!geometry\_msgs::msg::dds\_::QuaternionStamped\_!}
\item
  字段类型：\passthrough{\lstinline!::geometry\_msgs::msg::dds\_::Quaternion\_!}
\item
  声明位置：\passthrough{\lstinline!include/unitree/idl/ros2/QuaternionStamped\_.hpp:27!}
\end{itemize}

\textbf{用法}

\begin{lstlisting}[language=Python]
current_value = obj.quaternion
obj.quaternion = new_value
\end{lstlisting}

\hypertarget{unitree-sdk2-cpp-idl-ros2-string}{%
\subsubsection{\texorpdfstring{\texttt{unitree\_sdk2\_cpp.idl.ros2.String}}{unitree\_sdk2\_cpp.idl.ros2.String}}\label{unitree-sdk2-cpp-idl-ros2-string}}

可由 Python 读写的 DDS/IDL 消息值类型。字段采用复制语义。

\textbf{导入}

\begin{lstlisting}[language=Python]
from unitree_sdk2_cpp.idl.ros2 import String
\end{lstlisting}

\textbf{方法索引}

\begin{longtable}[]{@{}
  >{\raggedright\arraybackslash}p{(\columnwidth - 6\tabcolsep) * \real{0.18}}
  >{\raggedright\arraybackslash}p{(\columnwidth - 6\tabcolsep) * \real{0.24}}
  >{\raggedright\arraybackslash}p{(\columnwidth - 6\tabcolsep) * \real{0.18}}
  >{\raggedright\arraybackslash}p{(\columnwidth - 6\tabcolsep) * \real{0.40}}@{}}
\toprule
方法 & Python 签名 & 状态 & 安全分类 \\
\midrule
\endhead
\protect\hyperlink{unitree-sdk2-cpp-idl-ros2-string-dunder-init-1}{\passthrough{\lstinline!\_\_init\_\_!}}
& \passthrough{\lstinline!def \_\_init\_\_(self) -> None!} &
\passthrough{\lstinline!AVAILABLE!} &
\passthrough{\lstinline!VALUE\_TYPE!} \\
\protect\hyperlink{unitree-sdk2-cpp-idl-ros2-string-dunder-eq-1}{\passthrough{\lstinline!\_\_eq\_\_!}}
& \passthrough{\lstinline!def \_\_eq\_\_(self, other: object) -> bool!}
& \passthrough{\lstinline!AVAILABLE!} &
\passthrough{\lstinline!VALUE\_TYPE!} \\
\protect\hyperlink{unitree-sdk2-cpp-idl-ros2-string-dunder-ne-1}{\passthrough{\lstinline!\_\_ne\_\_!}}
& \passthrough{\lstinline!def \_\_ne\_\_(self, other: object) -> bool!}
& \passthrough{\lstinline!AVAILABLE!} &
\passthrough{\lstinline!VALUE\_TYPE!} \\
\bottomrule
\end{longtable}

\textbf{属性索引}

\begin{longtable}[]{@{}
  >{\raggedright\arraybackslash}p{(\columnwidth - 6\tabcolsep) * \real{0.18}}
  >{\raggedright\arraybackslash}p{(\columnwidth - 6\tabcolsep) * \real{0.24}}
  >{\raggedright\arraybackslash}p{(\columnwidth - 6\tabcolsep) * \real{0.18}}
  >{\raggedright\arraybackslash}p{(\columnwidth - 6\tabcolsep) * \real{0.40}}@{}}
\toprule
属性 & 读取类型 & 写入类型 & 状态 \\
\midrule
\endhead
\protect\hyperlink{unitree-sdk2-cpp-idl-ros2-string-data}{\passthrough{\lstinline!data!}}
& \passthrough{\lstinline!str!} & \passthrough{\lstinline!str!} &
\passthrough{\lstinline!AVAILABLE!} \\
\bottomrule
\end{longtable}

\hypertarget{unitree-sdk2-cpp-idl-ros2-string-dunder-init-1}{%
\paragraph{\texorpdfstring{\texttt{unitree\_sdk2\_cpp.idl.ros2.String.\_\_init\_\_}}{unitree\_sdk2\_cpp.idl.ros2.String.\_\_init\_\_}}\label{unitree-sdk2-cpp-idl-ros2-string-dunder-init-1}}

初始化 \passthrough{\lstinline!String!}
实例。是否能实际构造取决于下方可用性状态。

\textbf{签名}

\begin{lstlisting}[language=Python]
def __init__(self) -> None
\end{lstlisting}

\textbf{可用性与安全性}

\textbf{\passthrough{\lstinline!AVAILABLE!} /
\passthrough{\lstinline!VALUE\_TYPE!}}：当前绑定源码已实现该消息值操作。

\textbf{参数}

无显式参数。实例方法中的 \passthrough{\lstinline!self!} 由 Python
自动传入。

\textbf{返回值}

\begin{longtable}[]{@{}
  >{\raggedright\arraybackslash}p{(\columnwidth - 4\tabcolsep) * \real{0.18}}
  >{\raggedright\arraybackslash}p{(\columnwidth - 4\tabcolsep) * \real{0.20}}
  >{\raggedright\arraybackslash}p{(\columnwidth - 4\tabcolsep) * \real{0.62}}@{}}
\toprule
位置 & 类型 & 含义 \\
\midrule
\endhead
无 & \passthrough{\lstinline!None!} &
\passthrough{\lstinline!\_\_init\_\_!}
本身不返回值；调用类对象时，在构造函数可用的前提下得到该类实例。 \\
\bottomrule
\end{longtable}

\textbf{对应 C++}

\begin{itemize}
\tightlist
\item
  类：\passthrough{\lstinline!std\_msgs::msg::dds\_::String\_!}
\item
  签名：\passthrough{\lstinline!String\_()!}
\item
  绑定策略：\passthrough{\lstinline!未单独标注!}
\item
  声明位置：\passthrough{\lstinline!include/unitree/idl/ros2/String\_.hpp:26!}
\end{itemize}

\textbf{说明}

\begin{itemize}
\tightlist
\item
  示例是局部调用片段；\passthrough{\lstinline!obj!}
  和参数变量表示已经按上层业务校验并准备好的对象。
\end{itemize}

\textbf{用法}

\begin{lstlisting}[language=Python]
value = String()
\end{lstlisting}

\hypertarget{unitree-sdk2-cpp-idl-ros2-string-dunder-eq-1}{%
\paragraph{\texorpdfstring{\texttt{unitree\_sdk2\_cpp.idl.ros2.String.\_\_eq\_\_}}{unitree\_sdk2\_cpp.idl.ros2.String.\_\_eq\_\_}}\label{unitree-sdk2-cpp-idl-ros2-string-dunder-eq-1}}

按底层 IDL 消息内容比较两个对象是否相等。

\textbf{签名}

\begin{lstlisting}[language=Python]
def __eq__(self, other: object) -> bool
\end{lstlisting}

\textbf{可用性与安全性}

\textbf{\passthrough{\lstinline!AVAILABLE!} /
\passthrough{\lstinline!VALUE\_TYPE!}}：当前绑定源码已实现该消息值操作。

\textbf{参数}

\begin{longtable}[]{@{}
  >{\raggedright\arraybackslash}p{(\columnwidth - 6\tabcolsep) * \real{0.18}}
  >{\raggedright\arraybackslash}p{(\columnwidth - 6\tabcolsep) * \real{0.24}}
  >{\raggedright\arraybackslash}p{(\columnwidth - 6\tabcolsep) * \real{0.18}}
  >{\raggedright\arraybackslash}p{(\columnwidth - 6\tabcolsep) * \real{0.40}}@{}}
\toprule
名称 & Python 类型 & 默认值 & 含义 \\
\midrule
\endhead
\passthrough{\lstinline!other!} & \passthrough{\lstinline!object!} &
必填 & 用于比较的另一个 Python 对象。类型不兼容时比较结果通常为
\passthrough{\lstinline!False!}。 \\
\bottomrule
\end{longtable}

\textbf{返回值}

\begin{longtable}[]{@{}
  >{\raggedright\arraybackslash}p{(\columnwidth - 4\tabcolsep) * \real{0.18}}
  >{\raggedright\arraybackslash}p{(\columnwidth - 4\tabcolsep) * \real{0.20}}
  >{\raggedright\arraybackslash}p{(\columnwidth - 4\tabcolsep) * \real{0.62}}@{}}
\toprule
位置 & 类型 & 含义 \\
\midrule
\endhead
返回值 & \passthrough{\lstinline!bool!} & 返回
\passthrough{\lstinline!bool!}。更精确的业务含义见该方法用途、C++
签名和目标型号协议。 \\
\bottomrule
\end{longtable}

\textbf{对应 C++}

\begin{itemize}
\tightlist
\item
  类：\passthrough{\lstinline!std\_msgs::msg::dds\_::String\_!}
\item
  签名：\passthrough{\lstinline!operator==(const value \&) const!}
\item
  绑定策略：\passthrough{\lstinline!未单独标注!}
\item
  声明位置：由 pybind11 手工绑定或生成注册表提供；manifest
  未记录独立头文件行号。
\end{itemize}

\textbf{说明}

\begin{itemize}
\tightlist
\item
  示例是局部调用片段；\passthrough{\lstinline!obj!}
  和参数变量表示已经按上层业务校验并准备好的对象。
\end{itemize}

\textbf{用法}

\begin{lstlisting}[language=Python]
same = left == right
\end{lstlisting}

\hypertarget{unitree-sdk2-cpp-idl-ros2-string-dunder-ne-1}{%
\paragraph{\texorpdfstring{\texttt{unitree\_sdk2\_cpp.idl.ros2.String.\_\_ne\_\_}}{unitree\_sdk2\_cpp.idl.ros2.String.\_\_ne\_\_}}\label{unitree-sdk2-cpp-idl-ros2-string-dunder-ne-1}}

按底层 IDL 消息内容比较两个对象是否不相等。

\textbf{签名}

\begin{lstlisting}[language=Python]
def __ne__(self, other: object) -> bool
\end{lstlisting}

\textbf{可用性与安全性}

\textbf{\passthrough{\lstinline!AVAILABLE!} /
\passthrough{\lstinline!VALUE\_TYPE!}}：当前绑定源码已实现该消息值操作。

\textbf{参数}

\begin{longtable}[]{@{}
  >{\raggedright\arraybackslash}p{(\columnwidth - 6\tabcolsep) * \real{0.18}}
  >{\raggedright\arraybackslash}p{(\columnwidth - 6\tabcolsep) * \real{0.24}}
  >{\raggedright\arraybackslash}p{(\columnwidth - 6\tabcolsep) * \real{0.18}}
  >{\raggedright\arraybackslash}p{(\columnwidth - 6\tabcolsep) * \real{0.40}}@{}}
\toprule
名称 & Python 类型 & 默认值 & 含义 \\
\midrule
\endhead
\passthrough{\lstinline!other!} & \passthrough{\lstinline!object!} &
必填 & 用于比较的另一个 Python 对象。类型不兼容时比较结果通常为
\passthrough{\lstinline!False!}。 \\
\bottomrule
\end{longtable}

\textbf{返回值}

\begin{longtable}[]{@{}
  >{\raggedright\arraybackslash}p{(\columnwidth - 4\tabcolsep) * \real{0.18}}
  >{\raggedright\arraybackslash}p{(\columnwidth - 4\tabcolsep) * \real{0.20}}
  >{\raggedright\arraybackslash}p{(\columnwidth - 4\tabcolsep) * \real{0.62}}@{}}
\toprule
位置 & 类型 & 含义 \\
\midrule
\endhead
返回值 & \passthrough{\lstinline!bool!} & 返回
\passthrough{\lstinline!bool!}。更精确的业务含义见该方法用途、C++
签名和目标型号协议。 \\
\bottomrule
\end{longtable}

\textbf{对应 C++}

\begin{itemize}
\tightlist
\item
  类：\passthrough{\lstinline!std\_msgs::msg::dds\_::String\_!}
\item
  签名：\passthrough{\lstinline"operator!=(const value \&) const"}
\item
  绑定策略：\passthrough{\lstinline!未单独标注!}
\item
  声明位置：由 pybind11 手工绑定或生成注册表提供；manifest
  未记录独立头文件行号。
\end{itemize}

\textbf{说明}

\begin{itemize}
\tightlist
\item
  示例是局部调用片段；\passthrough{\lstinline!obj!}
  和参数变量表示已经按上层业务校验并准备好的对象。
\end{itemize}

\textbf{用法}

\begin{lstlisting}[language=Python]
different = left != right
\end{lstlisting}

\hypertarget{unitree-sdk2-cpp-idl-ros2-string-data}{%
\paragraph{\texorpdfstring{\texttt{unitree\_sdk2\_cpp.idl.ros2.String.data}}{unitree\_sdk2\_cpp.idl.ros2.String.data}}\label{unitree-sdk2-cpp-idl-ros2-string-data}}

底层 \passthrough{\lstinline!std\_msgs::msg::dds\_::String\_!} 的
\passthrough{\lstinline!data!} 字段。类型签名只说明 Python
表示；单位、数值范围、数组固定长度和枚举含义需查对应 IDL 头文件。
底层为字符串；长度、编码和允许值由具体协议决定。

\textbf{签名}

\begin{lstlisting}[language=Python]
@property
def data(self) -> str

@data.setter
def data(self, value: str) -> None
\end{lstlisting}

\textbf{可用性与安全性}

\textbf{\passthrough{\lstinline!AVAILABLE!} /
\passthrough{\lstinline!VALUE\_TYPE!}}：当前绑定源码已实现该消息值操作。

\textbf{参数}

\begin{longtable}[]{@{}
  >{\raggedright\arraybackslash}p{(\columnwidth - 4\tabcolsep) * \real{0.18}}
  >{\raggedright\arraybackslash}p{(\columnwidth - 4\tabcolsep) * \real{0.20}}
  >{\raggedright\arraybackslash}p{(\columnwidth - 4\tabcolsep) * \real{0.62}}@{}}
\toprule
名称 & Python 类型 & 含义 \\
\midrule
\endhead
\passthrough{\lstinline!value!} & \passthrough{\lstinline!str!} &
要写入或传递的新值，必须符合该参数的 Python 类型以及底层 C++
范围约束。 \\
\bottomrule
\end{longtable}

\textbf{返回值}

\begin{longtable}[]{@{}
  >{\raggedright\arraybackslash}p{(\columnwidth - 4\tabcolsep) * \real{0.18}}
  >{\raggedright\arraybackslash}p{(\columnwidth - 4\tabcolsep) * \real{0.20}}
  >{\raggedright\arraybackslash}p{(\columnwidth - 4\tabcolsep) * \real{0.62}}@{}}
\toprule
位置 & 类型 & 含义 \\
\midrule
\endhead
getter 返回值 & \passthrough{\lstinline!str!} & 返回属性当前值。 \\
\bottomrule
\end{longtable}

\textbf{对应 C++}

\begin{itemize}
\tightlist
\item
  类：\passthrough{\lstinline!std\_msgs::msg::dds\_::String\_!}
\item
  字段类型：\passthrough{\lstinline!std::string!}
\item
  声明位置：\passthrough{\lstinline!include/unitree/idl/ros2/String\_.hpp:23!}
\end{itemize}

\textbf{用法}

\begin{lstlisting}[language=Python]
current_value = obj.data
obj.data = new_value
\end{lstlisting}

\hypertarget{unitree-sdk2-cpp-idl-ros2-twiststamped}{%
\subsubsection{\texorpdfstring{\texttt{unitree\_sdk2\_cpp.idl.ros2.TwistStamped}}{unitree\_sdk2\_cpp.idl.ros2.TwistStamped}}\label{unitree-sdk2-cpp-idl-ros2-twiststamped}}

可由 Python 读写的 DDS/IDL 消息值类型。字段采用复制语义。

\textbf{导入}

\begin{lstlisting}[language=Python]
from unitree_sdk2_cpp.idl.ros2 import TwistStamped
\end{lstlisting}

\textbf{方法索引}

\begin{longtable}[]{@{}
  >{\raggedright\arraybackslash}p{(\columnwidth - 6\tabcolsep) * \real{0.18}}
  >{\raggedright\arraybackslash}p{(\columnwidth - 6\tabcolsep) * \real{0.24}}
  >{\raggedright\arraybackslash}p{(\columnwidth - 6\tabcolsep) * \real{0.18}}
  >{\raggedright\arraybackslash}p{(\columnwidth - 6\tabcolsep) * \real{0.40}}@{}}
\toprule
方法 & Python 签名 & 状态 & 安全分类 \\
\midrule
\endhead
\protect\hyperlink{unitree-sdk2-cpp-idl-ros2-twiststamped-dunder-init-1}{\passthrough{\lstinline!\_\_init\_\_!}}
& \passthrough{\lstinline!def \_\_init\_\_(self) -> None!} &
\passthrough{\lstinline!AVAILABLE!} &
\passthrough{\lstinline!VALUE\_TYPE!} \\
\protect\hyperlink{unitree-sdk2-cpp-idl-ros2-twiststamped-dunder-eq-1}{\passthrough{\lstinline!\_\_eq\_\_!}}
& \passthrough{\lstinline!def \_\_eq\_\_(self, other: object) -> bool!}
& \passthrough{\lstinline!AVAILABLE!} &
\passthrough{\lstinline!VALUE\_TYPE!} \\
\protect\hyperlink{unitree-sdk2-cpp-idl-ros2-twiststamped-dunder-ne-1}{\passthrough{\lstinline!\_\_ne\_\_!}}
& \passthrough{\lstinline!def \_\_ne\_\_(self, other: object) -> bool!}
& \passthrough{\lstinline!AVAILABLE!} &
\passthrough{\lstinline!VALUE\_TYPE!} \\
\bottomrule
\end{longtable}

\textbf{属性索引}

\begin{longtable}[]{@{}
  >{\raggedright\arraybackslash}p{(\columnwidth - 6\tabcolsep) * \real{0.18}}
  >{\raggedright\arraybackslash}p{(\columnwidth - 6\tabcolsep) * \real{0.24}}
  >{\raggedright\arraybackslash}p{(\columnwidth - 6\tabcolsep) * \real{0.18}}
  >{\raggedright\arraybackslash}p{(\columnwidth - 6\tabcolsep) * \real{0.40}}@{}}
\toprule
属性 & 读取类型 & 写入类型 & 状态 \\
\midrule
\endhead
\protect\hyperlink{unitree-sdk2-cpp-idl-ros2-twiststamped-header}{\passthrough{\lstinline!header!}}
& \passthrough{\lstinline!Any!} & \passthrough{\lstinline!Any!} &
\passthrough{\lstinline!AVAILABLE!} \\
\protect\hyperlink{unitree-sdk2-cpp-idl-ros2-twiststamped-twist}{\passthrough{\lstinline!twist!}}
& \passthrough{\lstinline!Twist!} & \passthrough{\lstinline!Twist!} &
\passthrough{\lstinline!AVAILABLE!} \\
\bottomrule
\end{longtable}

\hypertarget{unitree-sdk2-cpp-idl-ros2-twiststamped-dunder-init-1}{%
\paragraph{\texorpdfstring{\texttt{unitree\_sdk2\_cpp.idl.ros2.TwistStamped.\_\_init\_\_}}{unitree\_sdk2\_cpp.idl.ros2.TwistStamped.\_\_init\_\_}}\label{unitree-sdk2-cpp-idl-ros2-twiststamped-dunder-init-1}}

初始化 \passthrough{\lstinline!TwistStamped!}
实例。是否能实际构造取决于下方可用性状态。

\textbf{签名}

\begin{lstlisting}[language=Python]
def __init__(self) -> None
\end{lstlisting}

\textbf{可用性与安全性}

\textbf{\passthrough{\lstinline!AVAILABLE!} /
\passthrough{\lstinline!VALUE\_TYPE!}}：当前绑定源码已实现该消息值操作。

\textbf{参数}

无显式参数。实例方法中的 \passthrough{\lstinline!self!} 由 Python
自动传入。

\textbf{返回值}

\begin{longtable}[]{@{}
  >{\raggedright\arraybackslash}p{(\columnwidth - 4\tabcolsep) * \real{0.18}}
  >{\raggedright\arraybackslash}p{(\columnwidth - 4\tabcolsep) * \real{0.20}}
  >{\raggedright\arraybackslash}p{(\columnwidth - 4\tabcolsep) * \real{0.62}}@{}}
\toprule
位置 & 类型 & 含义 \\
\midrule
\endhead
无 & \passthrough{\lstinline!None!} &
\passthrough{\lstinline!\_\_init\_\_!}
本身不返回值；调用类对象时，在构造函数可用的前提下得到该类实例。 \\
\bottomrule
\end{longtable}

\textbf{对应 C++}

\begin{itemize}
\tightlist
\item
  类：\passthrough{\lstinline!geometry\_msgs::msg::dds\_::TwistStamped\_!}
\item
  签名：\passthrough{\lstinline!TwistStamped\_()!}
\item
  绑定策略：\passthrough{\lstinline!未单独标注!}
\item
  声明位置：\passthrough{\lstinline!include/unitree/idl/ros2/TwistStamped\_.hpp:30!}
\end{itemize}

\textbf{说明}

\begin{itemize}
\tightlist
\item
  示例是局部调用片段；\passthrough{\lstinline!obj!}
  和参数变量表示已经按上层业务校验并准备好的对象。
\end{itemize}

\textbf{用法}

\begin{lstlisting}[language=Python]
value = TwistStamped()
\end{lstlisting}

\hypertarget{unitree-sdk2-cpp-idl-ros2-twiststamped-dunder-eq-1}{%
\paragraph{\texorpdfstring{\texttt{unitree\_sdk2\_cpp.idl.ros2.TwistStamped.\_\_eq\_\_}}{unitree\_sdk2\_cpp.idl.ros2.TwistStamped.\_\_eq\_\_}}\label{unitree-sdk2-cpp-idl-ros2-twiststamped-dunder-eq-1}}

按底层 IDL 消息内容比较两个对象是否相等。

\textbf{签名}

\begin{lstlisting}[language=Python]
def __eq__(self, other: object) -> bool
\end{lstlisting}

\textbf{可用性与安全性}

\textbf{\passthrough{\lstinline!AVAILABLE!} /
\passthrough{\lstinline!VALUE\_TYPE!}}：当前绑定源码已实现该消息值操作。

\textbf{参数}

\begin{longtable}[]{@{}
  >{\raggedright\arraybackslash}p{(\columnwidth - 6\tabcolsep) * \real{0.18}}
  >{\raggedright\arraybackslash}p{(\columnwidth - 6\tabcolsep) * \real{0.24}}
  >{\raggedright\arraybackslash}p{(\columnwidth - 6\tabcolsep) * \real{0.18}}
  >{\raggedright\arraybackslash}p{(\columnwidth - 6\tabcolsep) * \real{0.40}}@{}}
\toprule
名称 & Python 类型 & 默认值 & 含义 \\
\midrule
\endhead
\passthrough{\lstinline!other!} & \passthrough{\lstinline!object!} &
必填 & 用于比较的另一个 Python 对象。类型不兼容时比较结果通常为
\passthrough{\lstinline!False!}。 \\
\bottomrule
\end{longtable}

\textbf{返回值}

\begin{longtable}[]{@{}
  >{\raggedright\arraybackslash}p{(\columnwidth - 4\tabcolsep) * \real{0.18}}
  >{\raggedright\arraybackslash}p{(\columnwidth - 4\tabcolsep) * \real{0.20}}
  >{\raggedright\arraybackslash}p{(\columnwidth - 4\tabcolsep) * \real{0.62}}@{}}
\toprule
位置 & 类型 & 含义 \\
\midrule
\endhead
返回值 & \passthrough{\lstinline!bool!} & 返回
\passthrough{\lstinline!bool!}。更精确的业务含义见该方法用途、C++
签名和目标型号协议。 \\
\bottomrule
\end{longtable}

\textbf{对应 C++}

\begin{itemize}
\tightlist
\item
  类：\passthrough{\lstinline!geometry\_msgs::msg::dds\_::TwistStamped\_!}
\item
  签名：\passthrough{\lstinline!operator==(const value \&) const!}
\item
  绑定策略：\passthrough{\lstinline!未单独标注!}
\item
  声明位置：由 pybind11 手工绑定或生成注册表提供；manifest
  未记录独立头文件行号。
\end{itemize}

\textbf{说明}

\begin{itemize}
\tightlist
\item
  示例是局部调用片段；\passthrough{\lstinline!obj!}
  和参数变量表示已经按上层业务校验并准备好的对象。
\end{itemize}

\textbf{用法}

\begin{lstlisting}[language=Python]
same = left == right
\end{lstlisting}

\hypertarget{unitree-sdk2-cpp-idl-ros2-twiststamped-dunder-ne-1}{%
\paragraph{\texorpdfstring{\texttt{unitree\_sdk2\_cpp.idl.ros2.TwistStamped.\_\_ne\_\_}}{unitree\_sdk2\_cpp.idl.ros2.TwistStamped.\_\_ne\_\_}}\label{unitree-sdk2-cpp-idl-ros2-twiststamped-dunder-ne-1}}

按底层 IDL 消息内容比较两个对象是否不相等。

\textbf{签名}

\begin{lstlisting}[language=Python]
def __ne__(self, other: object) -> bool
\end{lstlisting}

\textbf{可用性与安全性}

\textbf{\passthrough{\lstinline!AVAILABLE!} /
\passthrough{\lstinline!VALUE\_TYPE!}}：当前绑定源码已实现该消息值操作。

\textbf{参数}

\begin{longtable}[]{@{}
  >{\raggedright\arraybackslash}p{(\columnwidth - 6\tabcolsep) * \real{0.18}}
  >{\raggedright\arraybackslash}p{(\columnwidth - 6\tabcolsep) * \real{0.24}}
  >{\raggedright\arraybackslash}p{(\columnwidth - 6\tabcolsep) * \real{0.18}}
  >{\raggedright\arraybackslash}p{(\columnwidth - 6\tabcolsep) * \real{0.40}}@{}}
\toprule
名称 & Python 类型 & 默认值 & 含义 \\
\midrule
\endhead
\passthrough{\lstinline!other!} & \passthrough{\lstinline!object!} &
必填 & 用于比较的另一个 Python 对象。类型不兼容时比较结果通常为
\passthrough{\lstinline!False!}。 \\
\bottomrule
\end{longtable}

\textbf{返回值}

\begin{longtable}[]{@{}
  >{\raggedright\arraybackslash}p{(\columnwidth - 4\tabcolsep) * \real{0.18}}
  >{\raggedright\arraybackslash}p{(\columnwidth - 4\tabcolsep) * \real{0.20}}
  >{\raggedright\arraybackslash}p{(\columnwidth - 4\tabcolsep) * \real{0.62}}@{}}
\toprule
位置 & 类型 & 含义 \\
\midrule
\endhead
返回值 & \passthrough{\lstinline!bool!} & 返回
\passthrough{\lstinline!bool!}。更精确的业务含义见该方法用途、C++
签名和目标型号协议。 \\
\bottomrule
\end{longtable}

\textbf{对应 C++}

\begin{itemize}
\tightlist
\item
  类：\passthrough{\lstinline!geometry\_msgs::msg::dds\_::TwistStamped\_!}
\item
  签名：\passthrough{\lstinline"operator!=(const value \&) const"}
\item
  绑定策略：\passthrough{\lstinline!未单独标注!}
\item
  声明位置：由 pybind11 手工绑定或生成注册表提供；manifest
  未记录独立头文件行号。
\end{itemize}

\textbf{说明}

\begin{itemize}
\tightlist
\item
  示例是局部调用片段；\passthrough{\lstinline!obj!}
  和参数变量表示已经按上层业务校验并准备好的对象。
\end{itemize}

\textbf{用法}

\begin{lstlisting}[language=Python]
different = left != right
\end{lstlisting}

\hypertarget{unitree-sdk2-cpp-idl-ros2-twiststamped-header}{%
\paragraph{\texorpdfstring{\texttt{unitree\_sdk2\_cpp.idl.ros2.TwistStamped.header}}{unitree\_sdk2\_cpp.idl.ros2.TwistStamped.header}}\label{unitree-sdk2-cpp-idl-ros2-twiststamped-header}}

底层
\passthrough{\lstinline!geometry\_msgs::msg::dds\_::TwistStamped\_!} 的
\passthrough{\lstinline!header!} 字段。类型签名只说明 Python
表示；单位、数值范围、数组固定长度和枚举含义需查对应 IDL 头文件。

\textbf{签名}

\begin{lstlisting}[language=Python]
@property
def header(self) -> Any

@header.setter
def header(self, value: Any) -> None
\end{lstlisting}

\textbf{可用性与安全性}

\textbf{\passthrough{\lstinline!AVAILABLE!} /
\passthrough{\lstinline!VALUE\_TYPE!}}：当前绑定源码已实现该消息值操作。

\textbf{参数}

\begin{longtable}[]{@{}
  >{\raggedright\arraybackslash}p{(\columnwidth - 4\tabcolsep) * \real{0.18}}
  >{\raggedright\arraybackslash}p{(\columnwidth - 4\tabcolsep) * \real{0.20}}
  >{\raggedright\arraybackslash}p{(\columnwidth - 4\tabcolsep) * \real{0.62}}@{}}
\toprule
名称 & Python 类型 & 含义 \\
\midrule
\endhead
\passthrough{\lstinline!value!} & \passthrough{\lstinline!Any!} &
要写入或传递的新值，必须符合该参数的 Python 类型以及底层 C++
范围约束。 \\
\bottomrule
\end{longtable}

\textbf{返回值}

\begin{longtable}[]{@{}
  >{\raggedright\arraybackslash}p{(\columnwidth - 4\tabcolsep) * \real{0.18}}
  >{\raggedright\arraybackslash}p{(\columnwidth - 4\tabcolsep) * \real{0.20}}
  >{\raggedright\arraybackslash}p{(\columnwidth - 4\tabcolsep) * \real{0.62}}@{}}
\toprule
位置 & 类型 & 含义 \\
\midrule
\endhead
getter 返回值 & \passthrough{\lstinline!Any!} & 返回属性当前值。 \\
\bottomrule
\end{longtable}

\textbf{对应 C++}

\begin{itemize}
\tightlist
\item
  类：\passthrough{\lstinline!geometry\_msgs::msg::dds\_::TwistStamped\_!}
\item
  字段类型：\passthrough{\lstinline!::std\_msgs::msg::dds\_::Header\_!}
\item
  声明位置：\passthrough{\lstinline!include/unitree/idl/ros2/TwistStamped\_.hpp:26!}
\end{itemize}

\textbf{用法}

\begin{lstlisting}[language=Python]
current_value = obj.header
obj.header = new_value
\end{lstlisting}

\hypertarget{unitree-sdk2-cpp-idl-ros2-twiststamped-twist}{%
\paragraph{\texorpdfstring{\texttt{unitree\_sdk2\_cpp.idl.ros2.TwistStamped.twist}}{unitree\_sdk2\_cpp.idl.ros2.TwistStamped.twist}}\label{unitree-sdk2-cpp-idl-ros2-twiststamped-twist}}

底层
\passthrough{\lstinline!geometry\_msgs::msg::dds\_::TwistStamped\_!} 的
\passthrough{\lstinline!twist!} 字段。类型签名只说明 Python
表示；单位、数值范围、数组固定长度和枚举含义需查对应 IDL 头文件。

\textbf{签名}

\begin{lstlisting}[language=Python]
@property
def twist(self) -> Twist

@twist.setter
def twist(self, value: Twist) -> None
\end{lstlisting}

\textbf{可用性与安全性}

\textbf{\passthrough{\lstinline!AVAILABLE!} /
\passthrough{\lstinline!VALUE\_TYPE!}}：当前绑定源码已实现该消息值操作。

\textbf{参数}

\begin{longtable}[]{@{}
  >{\raggedright\arraybackslash}p{(\columnwidth - 4\tabcolsep) * \real{0.18}}
  >{\raggedright\arraybackslash}p{(\columnwidth - 4\tabcolsep) * \real{0.20}}
  >{\raggedright\arraybackslash}p{(\columnwidth - 4\tabcolsep) * \real{0.62}}@{}}
\toprule
名称 & Python 类型 & 含义 \\
\midrule
\endhead
\passthrough{\lstinline!value!} & \passthrough{\lstinline!Twist!} &
要写入或传递的新值，必须符合该参数的 Python 类型以及底层 C++
范围约束。 \\
\bottomrule
\end{longtable}

\textbf{返回值}

\begin{longtable}[]{@{}
  >{\raggedright\arraybackslash}p{(\columnwidth - 4\tabcolsep) * \real{0.18}}
  >{\raggedright\arraybackslash}p{(\columnwidth - 4\tabcolsep) * \real{0.20}}
  >{\raggedright\arraybackslash}p{(\columnwidth - 4\tabcolsep) * \real{0.62}}@{}}
\toprule
位置 & 类型 & 含义 \\
\midrule
\endhead
getter 返回值 & \passthrough{\lstinline!Twist!} &
返回属性当前值。数组、vector 和嵌套消息采用复制语义。 \\
\bottomrule
\end{longtable}

\textbf{对应 C++}

\begin{itemize}
\tightlist
\item
  类：\passthrough{\lstinline!geometry\_msgs::msg::dds\_::TwistStamped\_!}
\item
  字段类型：\passthrough{\lstinline!::geometry\_msgs::msg::dds\_::Twist\_!}
\item
  声明位置：\passthrough{\lstinline!include/unitree/idl/ros2/TwistStamped\_.hpp:27!}
\end{itemize}

\textbf{用法}

\begin{lstlisting}[language=Python]
current_value = obj.twist
obj.twist = new_value
\end{lstlisting}

\hypertarget{unitree-sdk2-cpp-idl-ros2-twistwithcovariancestamped}{%
\subsubsection{\texorpdfstring{\texttt{unitree\_sdk2\_cpp.idl.ros2.TwistWithCovarianceStamped}}{unitree\_sdk2\_cpp.idl.ros2.TwistWithCovarianceStamped}}\label{unitree-sdk2-cpp-idl-ros2-twistwithcovariancestamped}}

可由 Python 读写的 DDS/IDL 消息值类型。字段采用复制语义。

\textbf{导入}

\begin{lstlisting}[language=Python]
from unitree_sdk2_cpp.idl.ros2 import TwistWithCovarianceStamped
\end{lstlisting}

\textbf{方法索引}

\begin{longtable}[]{@{}
  >{\raggedright\arraybackslash}p{(\columnwidth - 6\tabcolsep) * \real{0.18}}
  >{\raggedright\arraybackslash}p{(\columnwidth - 6\tabcolsep) * \real{0.24}}
  >{\raggedright\arraybackslash}p{(\columnwidth - 6\tabcolsep) * \real{0.18}}
  >{\raggedright\arraybackslash}p{(\columnwidth - 6\tabcolsep) * \real{0.40}}@{}}
\toprule
方法 & Python 签名 & 状态 & 安全分类 \\
\midrule
\endhead
\protect\hyperlink{unitree-sdk2-cpp-idl-ros2-twistwithcovariancestamped-dunder-init-1}{\passthrough{\lstinline!\_\_init\_\_!}}
& \passthrough{\lstinline!def \_\_init\_\_(self) -> None!} &
\passthrough{\lstinline!AVAILABLE!} &
\passthrough{\lstinline!VALUE\_TYPE!} \\
\protect\hyperlink{unitree-sdk2-cpp-idl-ros2-twistwithcovariancestamped-dunder-eq-1}{\passthrough{\lstinline!\_\_eq\_\_!}}
& \passthrough{\lstinline!def \_\_eq\_\_(self, other: object) -> bool!}
& \passthrough{\lstinline!AVAILABLE!} &
\passthrough{\lstinline!VALUE\_TYPE!} \\
\protect\hyperlink{unitree-sdk2-cpp-idl-ros2-twistwithcovariancestamped-dunder-ne-1}{\passthrough{\lstinline!\_\_ne\_\_!}}
& \passthrough{\lstinline!def \_\_ne\_\_(self, other: object) -> bool!}
& \passthrough{\lstinline!AVAILABLE!} &
\passthrough{\lstinline!VALUE\_TYPE!} \\
\bottomrule
\end{longtable}

\textbf{属性索引}

\begin{longtable}[]{@{}
  >{\raggedright\arraybackslash}p{(\columnwidth - 6\tabcolsep) * \real{0.18}}
  >{\raggedright\arraybackslash}p{(\columnwidth - 6\tabcolsep) * \real{0.24}}
  >{\raggedright\arraybackslash}p{(\columnwidth - 6\tabcolsep) * \real{0.18}}
  >{\raggedright\arraybackslash}p{(\columnwidth - 6\tabcolsep) * \real{0.40}}@{}}
\toprule
属性 & 读取类型 & 写入类型 & 状态 \\
\midrule
\endhead
\protect\hyperlink{unitree-sdk2-cpp-idl-ros2-twistwithcovariancestamped-header}{\passthrough{\lstinline!header!}}
& \passthrough{\lstinline!Any!} & \passthrough{\lstinline!Any!} &
\passthrough{\lstinline!AVAILABLE!} \\
\protect\hyperlink{unitree-sdk2-cpp-idl-ros2-twistwithcovariancestamped-twist}{\passthrough{\lstinline!twist!}}
& \passthrough{\lstinline!TwistWithCovariance!} &
\passthrough{\lstinline!TwistWithCovariance!} &
\passthrough{\lstinline!AVAILABLE!} \\
\bottomrule
\end{longtable}

\hypertarget{unitree-sdk2-cpp-idl-ros2-twistwithcovariancestamped-dunder-init-1}{%
\paragraph{\texorpdfstring{\texttt{unitree\_sdk2\_cpp.idl.ros2.TwistWithCovarianceStamped.\_\_init\_\_}}{unitree\_sdk2\_cpp.idl.ros2.TwistWithCovarianceStamped.\_\_init\_\_}}\label{unitree-sdk2-cpp-idl-ros2-twistwithcovariancestamped-dunder-init-1}}

初始化 \passthrough{\lstinline!TwistWithCovarianceStamped!}
实例。是否能实际构造取决于下方可用性状态。

\textbf{签名}

\begin{lstlisting}[language=Python]
def __init__(self) -> None
\end{lstlisting}

\textbf{可用性与安全性}

\textbf{\passthrough{\lstinline!AVAILABLE!} /
\passthrough{\lstinline!VALUE\_TYPE!}}：当前绑定源码已实现该消息值操作。

\textbf{参数}

无显式参数。实例方法中的 \passthrough{\lstinline!self!} 由 Python
自动传入。

\textbf{返回值}

\begin{longtable}[]{@{}
  >{\raggedright\arraybackslash}p{(\columnwidth - 4\tabcolsep) * \real{0.18}}
  >{\raggedright\arraybackslash}p{(\columnwidth - 4\tabcolsep) * \real{0.20}}
  >{\raggedright\arraybackslash}p{(\columnwidth - 4\tabcolsep) * \real{0.62}}@{}}
\toprule
位置 & 类型 & 含义 \\
\midrule
\endhead
无 & \passthrough{\lstinline!None!} &
\passthrough{\lstinline!\_\_init\_\_!}
本身不返回值；调用类对象时，在构造函数可用的前提下得到该类实例。 \\
\bottomrule
\end{longtable}

\textbf{对应 C++}

\begin{itemize}
\tightlist
\item
  类：\passthrough{\lstinline!geometry\_msgs::msg::dds\_::TwistWithCovarianceStamped\_!}
\item
  签名：\passthrough{\lstinline!TwistWithCovarianceStamped\_()!}
\item
  绑定策略：\passthrough{\lstinline!未单独标注!}
\item
  声明位置：\passthrough{\lstinline!include/unitree/idl/ros2/TwistWithCovarianceStamped\_.hpp:30!}
\end{itemize}

\textbf{说明}

\begin{itemize}
\tightlist
\item
  示例是局部调用片段；\passthrough{\lstinline!obj!}
  和参数变量表示已经按上层业务校验并准备好的对象。
\end{itemize}

\textbf{用法}

\begin{lstlisting}[language=Python]
value = TwistWithCovarianceStamped()
\end{lstlisting}

\hypertarget{unitree-sdk2-cpp-idl-ros2-twistwithcovariancestamped-dunder-eq-1}{%
\paragraph{\texorpdfstring{\texttt{unitree\_sdk2\_cpp.idl.ros2.TwistWithCovarianceStamped.\_\_eq\_\_}}{unitree\_sdk2\_cpp.idl.ros2.TwistWithCovarianceStamped.\_\_eq\_\_}}\label{unitree-sdk2-cpp-idl-ros2-twistwithcovariancestamped-dunder-eq-1}}

按底层 IDL 消息内容比较两个对象是否相等。

\textbf{签名}

\begin{lstlisting}[language=Python]
def __eq__(self, other: object) -> bool
\end{lstlisting}

\textbf{可用性与安全性}

\textbf{\passthrough{\lstinline!AVAILABLE!} /
\passthrough{\lstinline!VALUE\_TYPE!}}：当前绑定源码已实现该消息值操作。

\textbf{参数}

\begin{longtable}[]{@{}
  >{\raggedright\arraybackslash}p{(\columnwidth - 6\tabcolsep) * \real{0.18}}
  >{\raggedright\arraybackslash}p{(\columnwidth - 6\tabcolsep) * \real{0.24}}
  >{\raggedright\arraybackslash}p{(\columnwidth - 6\tabcolsep) * \real{0.18}}
  >{\raggedright\arraybackslash}p{(\columnwidth - 6\tabcolsep) * \real{0.40}}@{}}
\toprule
名称 & Python 类型 & 默认值 & 含义 \\
\midrule
\endhead
\passthrough{\lstinline!other!} & \passthrough{\lstinline!object!} &
必填 & 用于比较的另一个 Python 对象。类型不兼容时比较结果通常为
\passthrough{\lstinline!False!}。 \\
\bottomrule
\end{longtable}

\textbf{返回值}

\begin{longtable}[]{@{}
  >{\raggedright\arraybackslash}p{(\columnwidth - 4\tabcolsep) * \real{0.18}}
  >{\raggedright\arraybackslash}p{(\columnwidth - 4\tabcolsep) * \real{0.20}}
  >{\raggedright\arraybackslash}p{(\columnwidth - 4\tabcolsep) * \real{0.62}}@{}}
\toprule
位置 & 类型 & 含义 \\
\midrule
\endhead
返回值 & \passthrough{\lstinline!bool!} & 返回
\passthrough{\lstinline!bool!}。更精确的业务含义见该方法用途、C++
签名和目标型号协议。 \\
\bottomrule
\end{longtable}

\textbf{对应 C++}

\begin{itemize}
\tightlist
\item
  类：\passthrough{\lstinline!geometry\_msgs::msg::dds\_::TwistWithCovarianceStamped\_!}
\item
  签名：\passthrough{\lstinline!operator==(const value \&) const!}
\item
  绑定策略：\passthrough{\lstinline!未单独标注!}
\item
  声明位置：由 pybind11 手工绑定或生成注册表提供；manifest
  未记录独立头文件行号。
\end{itemize}

\textbf{说明}

\begin{itemize}
\tightlist
\item
  示例是局部调用片段；\passthrough{\lstinline!obj!}
  和参数变量表示已经按上层业务校验并准备好的对象。
\end{itemize}

\textbf{用法}

\begin{lstlisting}[language=Python]
same = left == right
\end{lstlisting}

\hypertarget{unitree-sdk2-cpp-idl-ros2-twistwithcovariancestamped-dunder-ne-1}{%
\paragraph{\texorpdfstring{\texttt{unitree\_sdk2\_cpp.idl.ros2.TwistWithCovarianceStamped.\_\_ne\_\_}}{unitree\_sdk2\_cpp.idl.ros2.TwistWithCovarianceStamped.\_\_ne\_\_}}\label{unitree-sdk2-cpp-idl-ros2-twistwithcovariancestamped-dunder-ne-1}}

按底层 IDL 消息内容比较两个对象是否不相等。

\textbf{签名}

\begin{lstlisting}[language=Python]
def __ne__(self, other: object) -> bool
\end{lstlisting}

\textbf{可用性与安全性}

\textbf{\passthrough{\lstinline!AVAILABLE!} /
\passthrough{\lstinline!VALUE\_TYPE!}}：当前绑定源码已实现该消息值操作。

\textbf{参数}

\begin{longtable}[]{@{}
  >{\raggedright\arraybackslash}p{(\columnwidth - 6\tabcolsep) * \real{0.18}}
  >{\raggedright\arraybackslash}p{(\columnwidth - 6\tabcolsep) * \real{0.24}}
  >{\raggedright\arraybackslash}p{(\columnwidth - 6\tabcolsep) * \real{0.18}}
  >{\raggedright\arraybackslash}p{(\columnwidth - 6\tabcolsep) * \real{0.40}}@{}}
\toprule
名称 & Python 类型 & 默认值 & 含义 \\
\midrule
\endhead
\passthrough{\lstinline!other!} & \passthrough{\lstinline!object!} &
必填 & 用于比较的另一个 Python 对象。类型不兼容时比较结果通常为
\passthrough{\lstinline!False!}。 \\
\bottomrule
\end{longtable}

\textbf{返回值}

\begin{longtable}[]{@{}
  >{\raggedright\arraybackslash}p{(\columnwidth - 4\tabcolsep) * \real{0.18}}
  >{\raggedright\arraybackslash}p{(\columnwidth - 4\tabcolsep) * \real{0.20}}
  >{\raggedright\arraybackslash}p{(\columnwidth - 4\tabcolsep) * \real{0.62}}@{}}
\toprule
位置 & 类型 & 含义 \\
\midrule
\endhead
返回值 & \passthrough{\lstinline!bool!} & 返回
\passthrough{\lstinline!bool!}。更精确的业务含义见该方法用途、C++
签名和目标型号协议。 \\
\bottomrule
\end{longtable}

\textbf{对应 C++}

\begin{itemize}
\tightlist
\item
  类：\passthrough{\lstinline!geometry\_msgs::msg::dds\_::TwistWithCovarianceStamped\_!}
\item
  签名：\passthrough{\lstinline"operator!=(const value \&) const"}
\item
  绑定策略：\passthrough{\lstinline!未单独标注!}
\item
  声明位置：由 pybind11 手工绑定或生成注册表提供；manifest
  未记录独立头文件行号。
\end{itemize}

\textbf{说明}

\begin{itemize}
\tightlist
\item
  示例是局部调用片段；\passthrough{\lstinline!obj!}
  和参数变量表示已经按上层业务校验并准备好的对象。
\end{itemize}

\textbf{用法}

\begin{lstlisting}[language=Python]
different = left != right
\end{lstlisting}

\hypertarget{unitree-sdk2-cpp-idl-ros2-twistwithcovariancestamped-header}{%
\paragraph{\texorpdfstring{\texttt{unitree\_sdk2\_cpp.idl.ros2.TwistWithCovarianceStamped.header}}{unitree\_sdk2\_cpp.idl.ros2.TwistWithCovarianceStamped.header}}\label{unitree-sdk2-cpp-idl-ros2-twistwithcovariancestamped-header}}

底层
\passthrough{\lstinline!geometry\_msgs::msg::dds\_::TwistWithCovarianceStamped\_!}
的 \passthrough{\lstinline!header!} 字段。类型签名只说明 Python
表示；单位、数值范围、数组固定长度和枚举含义需查对应 IDL 头文件。

\textbf{签名}

\begin{lstlisting}[language=Python]
@property
def header(self) -> Any

@header.setter
def header(self, value: Any) -> None
\end{lstlisting}

\textbf{可用性与安全性}

\textbf{\passthrough{\lstinline!AVAILABLE!} /
\passthrough{\lstinline!VALUE\_TYPE!}}：当前绑定源码已实现该消息值操作。

\textbf{参数}

\begin{longtable}[]{@{}
  >{\raggedright\arraybackslash}p{(\columnwidth - 4\tabcolsep) * \real{0.18}}
  >{\raggedright\arraybackslash}p{(\columnwidth - 4\tabcolsep) * \real{0.20}}
  >{\raggedright\arraybackslash}p{(\columnwidth - 4\tabcolsep) * \real{0.62}}@{}}
\toprule
名称 & Python 类型 & 含义 \\
\midrule
\endhead
\passthrough{\lstinline!value!} & \passthrough{\lstinline!Any!} &
要写入或传递的新值，必须符合该参数的 Python 类型以及底层 C++
范围约束。 \\
\bottomrule
\end{longtable}

\textbf{返回值}

\begin{longtable}[]{@{}
  >{\raggedright\arraybackslash}p{(\columnwidth - 4\tabcolsep) * \real{0.18}}
  >{\raggedright\arraybackslash}p{(\columnwidth - 4\tabcolsep) * \real{0.20}}
  >{\raggedright\arraybackslash}p{(\columnwidth - 4\tabcolsep) * \real{0.62}}@{}}
\toprule
位置 & 类型 & 含义 \\
\midrule
\endhead
getter 返回值 & \passthrough{\lstinline!Any!} & 返回属性当前值。 \\
\bottomrule
\end{longtable}

\textbf{对应 C++}

\begin{itemize}
\tightlist
\item
  类：\passthrough{\lstinline!geometry\_msgs::msg::dds\_::TwistWithCovarianceStamped\_!}
\item
  字段类型：\passthrough{\lstinline!::std\_msgs::msg::dds\_::Header\_!}
\item
  声明位置：\passthrough{\lstinline!include/unitree/idl/ros2/TwistWithCovarianceStamped\_.hpp:26!}
\end{itemize}

\textbf{用法}

\begin{lstlisting}[language=Python]
current_value = obj.header
obj.header = new_value
\end{lstlisting}

\hypertarget{unitree-sdk2-cpp-idl-ros2-twistwithcovariancestamped-twist}{%
\paragraph{\texorpdfstring{\texttt{unitree\_sdk2\_cpp.idl.ros2.TwistWithCovarianceStamped.twist}}{unitree\_sdk2\_cpp.idl.ros2.TwistWithCovarianceStamped.twist}}\label{unitree-sdk2-cpp-idl-ros2-twistwithcovariancestamped-twist}}

底层
\passthrough{\lstinline!geometry\_msgs::msg::dds\_::TwistWithCovarianceStamped\_!}
的 \passthrough{\lstinline!twist!} 字段。类型签名只说明 Python
表示；单位、数值范围、数组固定长度和枚举含义需查对应 IDL 头文件。

\textbf{签名}

\begin{lstlisting}[language=Python]
@property
def twist(self) -> TwistWithCovariance

@twist.setter
def twist(self, value: TwistWithCovariance) -> None
\end{lstlisting}

\textbf{可用性与安全性}

\textbf{\passthrough{\lstinline!AVAILABLE!} /
\passthrough{\lstinline!VALUE\_TYPE!}}：当前绑定源码已实现该消息值操作。

\textbf{参数}

\begin{longtable}[]{@{}
  >{\raggedright\arraybackslash}p{(\columnwidth - 4\tabcolsep) * \real{0.18}}
  >{\raggedright\arraybackslash}p{(\columnwidth - 4\tabcolsep) * \real{0.20}}
  >{\raggedright\arraybackslash}p{(\columnwidth - 4\tabcolsep) * \real{0.62}}@{}}
\toprule
名称 & Python 类型 & 含义 \\
\midrule
\endhead
\passthrough{\lstinline!value!} &
\passthrough{\lstinline!TwistWithCovariance!} &
要写入或传递的新值，必须符合该参数的 Python 类型以及底层 C++
范围约束。 \\
\bottomrule
\end{longtable}

\textbf{返回值}

\begin{longtable}[]{@{}
  >{\raggedright\arraybackslash}p{(\columnwidth - 4\tabcolsep) * \real{0.18}}
  >{\raggedright\arraybackslash}p{(\columnwidth - 4\tabcolsep) * \real{0.20}}
  >{\raggedright\arraybackslash}p{(\columnwidth - 4\tabcolsep) * \real{0.62}}@{}}
\toprule
位置 & 类型 & 含义 \\
\midrule
\endhead
getter 返回值 & \passthrough{\lstinline!TwistWithCovariance!} &
返回属性当前值。数组、vector 和嵌套消息采用复制语义。 \\
\bottomrule
\end{longtable}

\textbf{对应 C++}

\begin{itemize}
\tightlist
\item
  类：\passthrough{\lstinline!geometry\_msgs::msg::dds\_::TwistWithCovarianceStamped\_!}
\item
  字段类型：\passthrough{\lstinline!::geometry\_msgs::msg::dds\_::TwistWithCovariance\_!}
\item
  声明位置：\passthrough{\lstinline!include/unitree/idl/ros2/TwistWithCovarianceStamped\_.hpp:27!}
\end{itemize}

\textbf{用法}

\begin{lstlisting}[language=Python]
current_value = obj.twist
obj.twist = new_value
\end{lstlisting}

\begin{center}\rule{0.5\linewidth}{0.5pt}\end{center}

\hypertarget{unitree-sdk2-cpp-robot}{%
\subsection{Robot 公共基础 API}\label{unitree-sdk2-cpp-robot}}

模块：\passthrough{\lstinline!unitree\_sdk2\_cpp.robot!}

\hypertarget{ux7c7bux7d22ux5f15-6}{%
\subsubsection{类索引}\label{ux7c7bux7d22ux5f15-6}}

\begin{longtable}[]{@{}
  >{\raggedright\arraybackslash}p{(\columnwidth - 4\tabcolsep) * \real{0.18}}
  >{\raggedleft\arraybackslash}p{(\columnwidth - 4\tabcolsep) * \real{0.20}}
  >{\raggedleft\arraybackslash}p{(\columnwidth - 4\tabcolsep) * \real{0.62}}@{}}
\toprule
类 & 公开函数签名 & 属性 \\
\midrule
\endhead
\protect\hyperlink{unitree-sdk2-cpp-robot-applyleasedata}{\passthrough{\lstinline!ApplyLeaseData!}}
& 3 & 0 \\
\protect\hyperlink{unitree-sdk2-cpp-robot-applyleaseparameter}{\passthrough{\lstinline!ApplyLeaseParameter!}}
& 3 & 0 \\
\protect\hyperlink{unitree-sdk2-cpp-robot-channelfactory}{\passthrough{\lstinline!ChannelFactory!}}
& 5 & 0 \\
\protect\hyperlink{unitree-sdk2-cpp-robot-channelnamer}{\passthrough{\lstinline!ChannelNamer!}}
& 2 & 0 \\
\protect\hyperlink{unitree-sdk2-cpp-robot-client}{\passthrough{\lstinline!Client!}}
& 5 & 0 \\
\protect\hyperlink{unitree-sdk2-cpp-robot-clientbase}{\passthrough{\lstinline!ClientBase!}}
& 4 & 0 \\
\protect\hyperlink{unitree-sdk2-cpp-robot-clientchannelnamer}{\passthrough{\lstinline!ClientChannelNamer!}}
& 1 & 0 \\
\protect\hyperlink{unitree-sdk2-cpp-robot-clientstub}{\passthrough{\lstinline!ClientStub!}}
& 4 & 0 \\
\protect\hyperlink{unitree-sdk2-cpp-robot-leasecache}{\passthrough{\lstinline!LeaseCache!}}
& 7 & 0 \\
\protect\hyperlink{unitree-sdk2-cpp-robot-leaseclient}{\passthrough{\lstinline!LeaseClient!}}
& 5 & 0 \\
\protect\hyperlink{unitree-sdk2-cpp-robot-leasecontext}{\passthrough{\lstinline!LeaseContext!}}
& 6 & 0 \\
\protect\hyperlink{unitree-sdk2-cpp-robot-leaseserver}{\passthrough{\lstinline!LeaseServer!}}
& 3 & 0 \\
\protect\hyperlink{unitree-sdk2-cpp-robot-requestfuture}{\passthrough{\lstinline!RequestFuture!}}
& 7 & 0 \\
\protect\hyperlink{unitree-sdk2-cpp-robot-requestfuturequeue}{\passthrough{\lstinline!RequestFutureQueue!}}
& 5 & 0 \\
\protect\hyperlink{unitree-sdk2-cpp-robot-server}{\passthrough{\lstinline!Server!}}
& 7 & 0 \\
\protect\hyperlink{unitree-sdk2-cpp-robot-serverbase}{\passthrough{\lstinline!ServerBase!}}
& 4 & 0 \\
\protect\hyperlink{unitree-sdk2-cpp-robot-serverchannelnamer}{\passthrough{\lstinline!ServerChannelNamer!}}
& 1 & 0 \\
\protect\hyperlink{unitree-sdk2-cpp-robot-serverstub}{\passthrough{\lstinline!ServerStub!}}
& 3 & 0 \\
\bottomrule
\end{longtable}

\hypertarget{unitree-sdk2-cpp-robot-applyleasedata}{%
\subsubsection{\texorpdfstring{\texttt{unitree\_sdk2\_cpp.robot.ApplyLeaseData}}{unitree\_sdk2\_cpp.robot.ApplyLeaseData}}\label{unitree-sdk2-cpp-robot-applyleasedata}}

SDK 数据值类型；公开字段和转换方法见下文。

\textbf{导入}

\begin{lstlisting}[language=Python]
from unitree_sdk2_cpp.robot import ApplyLeaseData
\end{lstlisting}

\textbf{基类}：\passthrough{\lstinline!object!}

\textbf{公开属性}

\begin{longtable}[]{@{}
  >{\raggedright\arraybackslash}p{(\columnwidth - 6\tabcolsep) * \real{0.18}}
  >{\raggedright\arraybackslash}p{(\columnwidth - 6\tabcolsep) * \real{0.24}}
  >{\raggedright\arraybackslash}p{(\columnwidth - 6\tabcolsep) * \real{0.18}}
  >{\raggedright\arraybackslash}p{(\columnwidth - 6\tabcolsep) * \real{0.40}}@{}}
\toprule
名称 & Python 类型 & 含义 & 用法 \\
\midrule
\endhead
\passthrough{\lstinline!id!} & \passthrough{\lstinline!int!} &
传给该接口的 ID 参数，Python 类型为
\passthrough{\lstinline!int!}。精确含义、单位、范围和枚举值需查目标型号对应头文件与协议。
& \passthrough{\lstinline!value = obj.id!} /
\passthrough{\lstinline!obj.id = value!} \\
\passthrough{\lstinline!term!} & \passthrough{\lstinline!int!} &
传给该接口的 \passthrough{\lstinline!term!} 参数，Python 类型为
\passthrough{\lstinline!int!}。精确含义、单位、范围和枚举值需查目标型号对应头文件与协议。
& \passthrough{\lstinline!value = obj.term!} /
\passthrough{\lstinline!obj.term = value!} \\
\bottomrule
\end{longtable}

\textbf{方法索引}

\begin{longtable}[]{@{}
  >{\raggedright\arraybackslash}p{(\columnwidth - 6\tabcolsep) * \real{0.18}}
  >{\raggedright\arraybackslash}p{(\columnwidth - 6\tabcolsep) * \real{0.24}}
  >{\raggedright\arraybackslash}p{(\columnwidth - 6\tabcolsep) * \real{0.18}}
  >{\raggedright\arraybackslash}p{(\columnwidth - 6\tabcolsep) * \real{0.40}}@{}}
\toprule
方法 & Python 签名 & 状态 & 安全分类 \\
\midrule
\endhead
\protect\hyperlink{unitree-sdk2-cpp-robot-applyleasedata-dunder-init-1}{\passthrough{\lstinline!\_\_init\_\_!}}
& \passthrough{\lstinline!def \_\_init\_\_(self) -> None!} &
\passthrough{\lstinline!SIGNATURE\_ONLY!} &
\passthrough{\lstinline!CONSTRUCTION!} \\
\protect\hyperlink{unitree-sdk2-cpp-robot-applyleasedata-from-json-1}{\passthrough{\lstinline!from\_json!}}
&
\passthrough{\lstinline!def from\_json(self, json: dict[str, Any]) -> None!}
& \passthrough{\lstinline!SIGNATURE\_ONLY!} &
\passthrough{\lstinline!UNCLASSIFIED!} \\
\protect\hyperlink{unitree-sdk2-cpp-robot-applyleasedata-to-json-1}{\passthrough{\lstinline!to\_json!}}
&
\passthrough{\lstinline!def to\_json(self, json: dict[str, Any]) -> None!}
& \passthrough{\lstinline!SIGNATURE\_ONLY!} &
\passthrough{\lstinline!UNCLASSIFIED!} \\
\bottomrule
\end{longtable}

\hypertarget{unitree-sdk2-cpp-robot-applyleasedata-dunder-init-1}{%
\paragraph{\texorpdfstring{\texttt{unitree\_sdk2\_cpp.robot.ApplyLeaseData.\_\_init\_\_}}{unitree\_sdk2\_cpp.robot.ApplyLeaseData.\_\_init\_\_}}\label{unitree-sdk2-cpp-robot-applyleasedata-dunder-init-1}}

初始化 \passthrough{\lstinline!ApplyLeaseData!}
实例。是否能实际构造取决于下方可用性状态。

\textbf{签名}

\begin{lstlisting}[language=Python]
def __init__(self) -> None
\end{lstlisting}

\textbf{可用性与安全性}

\textbf{\passthrough{\lstinline!SIGNATURE\_ONLY!} /
\passthrough{\lstinline!CONSTRUCTION!}}：当前只有设计期签名，不能假设已存在于运行时扩展。

\textbf{参数}

无显式参数。实例方法中的 \passthrough{\lstinline!self!} 由 Python
自动传入。

\textbf{返回值}

\begin{longtable}[]{@{}
  >{\raggedright\arraybackslash}p{(\columnwidth - 4\tabcolsep) * \real{0.18}}
  >{\raggedright\arraybackslash}p{(\columnwidth - 4\tabcolsep) * \real{0.20}}
  >{\raggedright\arraybackslash}p{(\columnwidth - 4\tabcolsep) * \real{0.62}}@{}}
\toprule
位置 & 类型 & 含义 \\
\midrule
\endhead
无 & \passthrough{\lstinline!None!} &
\passthrough{\lstinline!\_\_init\_\_!}
本身不返回值；调用类对象时，在构造函数可用的前提下得到该类实例。 \\
\bottomrule
\end{longtable}

\textbf{对应 C++}

\begin{itemize}
\tightlist
\item
  类：\passthrough{\lstinline!unitree::robot::ApplyLeaseData!}
\item
  签名：\passthrough{\lstinline!ApplyLeaseData()!}
\item
  绑定策略：\passthrough{\lstinline!未单独标注!}
\item
  声明位置：\passthrough{\lstinline!include/unitree/robot/internal/internal\_api.hpp:96!}
\end{itemize}

\textbf{说明}

\begin{itemize}
\tightlist
\item
  示例是局部调用片段；\passthrough{\lstinline!obj!}
  和参数变量表示已经按上层业务校验并准备好的对象。
\item
  该条目可以被编辑器和类型检查器识别，但当前运行时可能在属性查找阶段直接失败。
\end{itemize}

\textbf{用法}

\begin{lstlisting}[language=Python]
from typing import TYPE_CHECKING

if TYPE_CHECKING:
    def planned_construct() -> ApplyLeaseData:
        return ApplyLeaseData()
\end{lstlisting}

\begin{quote}
{[}!CAUTION{]} 上面的 \passthrough{\lstinline!TYPE\_CHECKING!}
示例只表达计划调用形态。该分支在普通 Python
运行时为假，不代表当前扩展能执行此接口。
\end{quote}

\hypertarget{unitree-sdk2-cpp-robot-applyleasedata-from-json-1}{%
\paragraph{\texorpdfstring{\texttt{unitree\_sdk2\_cpp.robot.ApplyLeaseData.from\_json}}{unitree\_sdk2\_cpp.robot.ApplyLeaseData.from\_json}}\label{unitree-sdk2-cpp-robot-applyleasedata-from-json-1}}

计划从 JSON 风格字典读取字段并更新当前 SDK 值对象。

\textbf{签名}

\begin{lstlisting}[language=Python]
def from_json(self, json: dict[str, Any]) -> None
\end{lstlisting}

\textbf{可用性与安全性}

\textbf{\passthrough{\lstinline!SIGNATURE\_ONLY!} /
\passthrough{\lstinline!UNCLASSIFIED!}}：当前只有设计期签名，不能假设已存在于运行时扩展。

\textbf{参数}

\begin{longtable}[]{@{}
  >{\raggedright\arraybackslash}p{(\columnwidth - 6\tabcolsep) * \real{0.18}}
  >{\raggedright\arraybackslash}p{(\columnwidth - 6\tabcolsep) * \real{0.24}}
  >{\raggedright\arraybackslash}p{(\columnwidth - 6\tabcolsep) * \real{0.18}}
  >{\raggedright\arraybackslash}p{(\columnwidth - 6\tabcolsep) * \real{0.40}}@{}}
\toprule
名称 & Python 类型 & 默认值 & 含义 \\
\midrule
\endhead
\passthrough{\lstinline!json!} &
\passthrough{\lstinline!dict[str, Any]!} & 必填 & JSON
风格字典，键和值必须满足对应 C++ \passthrough{\lstinline!JsonMap!}
数据结构的约束。 对应 C++ 参数
\passthrough{\lstinline!json: common::JsonMap \&!}。 \\
\bottomrule
\end{longtable}

\textbf{返回值}

\begin{longtable}[]{@{}
  >{\raggedright\arraybackslash}p{(\columnwidth - 4\tabcolsep) * \real{0.18}}
  >{\raggedright\arraybackslash}p{(\columnwidth - 4\tabcolsep) * \real{0.20}}
  >{\raggedright\arraybackslash}p{(\columnwidth - 4\tabcolsep) * \real{0.62}}@{}}
\toprule
位置 & 类型 & 含义 \\
\midrule
\endhead
无 & \passthrough{\lstinline!None!} & 该方法不返回 Python
值；失败可能表现为异常或底层状态变化。 \\
\bottomrule
\end{longtable}

\textbf{对应 C++}

\begin{itemize}
\tightlist
\item
  类：\passthrough{\lstinline!unitree::robot::ApplyLeaseData!}
\item
  签名：\passthrough{\lstinline!fromJson(common::JsonMap \&)!}
\item
  绑定策略：\passthrough{\lstinline!SIGNATURE\_PREVIEW!}
\item
  声明位置：\passthrough{\lstinline!include/unitree/robot/internal/internal\_api.hpp:102!}
\end{itemize}

\textbf{说明}

\begin{itemize}
\tightlist
\item
  示例是局部调用片段；\passthrough{\lstinline!obj!}
  和参数变量表示已经按上层业务校验并准备好的对象。
\item
  该条目可以被编辑器和类型检查器识别，但当前运行时可能在属性查找阶段直接失败。
\end{itemize}

\textbf{用法}

\begin{lstlisting}[language=Python]
from typing import TYPE_CHECKING

if TYPE_CHECKING:
    def planned_from_json(obj: ApplyLeaseData, json: dict[str, Any]) -> None:
        obj.from_json(json=json)
\end{lstlisting}

\begin{quote}
{[}!CAUTION{]} 上面的 \passthrough{\lstinline!TYPE\_CHECKING!}
示例只表达计划调用形态。该分支在普通 Python
运行时为假，不代表当前扩展能执行此接口。
\end{quote}

\hypertarget{unitree-sdk2-cpp-robot-applyleasedata-to-json-1}{%
\paragraph{\texorpdfstring{\texttt{unitree\_sdk2\_cpp.robot.ApplyLeaseData.to\_json}}{unitree\_sdk2\_cpp.robot.ApplyLeaseData.to\_json}}\label{unitree-sdk2-cpp-robot-applyleasedata-to-json-1}}

计划把当前 SDK 值对象写入 JSON 风格字典。

\textbf{签名}

\begin{lstlisting}[language=Python]
def to_json(self, json: dict[str, Any]) -> None
\end{lstlisting}

\textbf{可用性与安全性}

\textbf{\passthrough{\lstinline!SIGNATURE\_ONLY!} /
\passthrough{\lstinline!UNCLASSIFIED!}}：当前只有设计期签名，不能假设已存在于运行时扩展。

\textbf{参数}

\begin{longtable}[]{@{}
  >{\raggedright\arraybackslash}p{(\columnwidth - 6\tabcolsep) * \real{0.18}}
  >{\raggedright\arraybackslash}p{(\columnwidth - 6\tabcolsep) * \real{0.24}}
  >{\raggedright\arraybackslash}p{(\columnwidth - 6\tabcolsep) * \real{0.18}}
  >{\raggedright\arraybackslash}p{(\columnwidth - 6\tabcolsep) * \real{0.40}}@{}}
\toprule
名称 & Python 类型 & 默认值 & 含义 \\
\midrule
\endhead
\passthrough{\lstinline!json!} &
\passthrough{\lstinline!dict[str, Any]!} & 必填 & JSON
风格字典，键和值必须满足对应 C++ \passthrough{\lstinline!JsonMap!}
数据结构的约束。 对应 C++ 参数
\passthrough{\lstinline!json: common::JsonMap \&!}。 \\
\bottomrule
\end{longtable}

\textbf{返回值}

\begin{longtable}[]{@{}
  >{\raggedright\arraybackslash}p{(\columnwidth - 4\tabcolsep) * \real{0.18}}
  >{\raggedright\arraybackslash}p{(\columnwidth - 4\tabcolsep) * \real{0.20}}
  >{\raggedright\arraybackslash}p{(\columnwidth - 4\tabcolsep) * \real{0.62}}@{}}
\toprule
位置 & 类型 & 含义 \\
\midrule
\endhead
无 & \passthrough{\lstinline!None!} & 该方法不返回 Python
值；失败可能表现为异常或底层状态变化。 \\
\bottomrule
\end{longtable}

\textbf{对应 C++}

\begin{itemize}
\tightlist
\item
  类：\passthrough{\lstinline!unitree::robot::ApplyLeaseData!}
\item
  签名：\passthrough{\lstinline!toJson(common::JsonMap \&) const!}
\item
  绑定策略：\passthrough{\lstinline!SIGNATURE\_PREVIEW!}
\item
  声明位置：\passthrough{\lstinline!include/unitree/robot/internal/internal\_api.hpp:108!}
\end{itemize}

\textbf{说明}

\begin{itemize}
\tightlist
\item
  示例是局部调用片段；\passthrough{\lstinline!obj!}
  和参数变量表示已经按上层业务校验并准备好的对象。
\item
  该条目可以被编辑器和类型检查器识别，但当前运行时可能在属性查找阶段直接失败。
\end{itemize}

\textbf{用法}

\begin{lstlisting}[language=Python]
from typing import TYPE_CHECKING

if TYPE_CHECKING:
    def planned_to_json(obj: ApplyLeaseData, json: dict[str, Any]) -> None:
        obj.to_json(json=json)
\end{lstlisting}

\begin{quote}
{[}!CAUTION{]} 上面的 \passthrough{\lstinline!TYPE\_CHECKING!}
示例只表达计划调用形态。该分支在普通 Python
运行时为假，不代表当前扩展能执行此接口。
\end{quote}

\hypertarget{unitree-sdk2-cpp-robot-applyleaseparameter}{%
\subsubsection{\texorpdfstring{\texttt{unitree\_sdk2\_cpp.robot.ApplyLeaseParameter}}{unitree\_sdk2\_cpp.robot.ApplyLeaseParameter}}\label{unitree-sdk2-cpp-robot-applyleaseparameter}}

SDK 请求参数值类型；当前可能仅提供设计期签名。

\textbf{导入}

\begin{lstlisting}[language=Python]
from unitree_sdk2_cpp.robot import ApplyLeaseParameter
\end{lstlisting}

\textbf{基类}：\passthrough{\lstinline!object!}

\textbf{公开属性}

\begin{longtable}[]{@{}
  >{\raggedright\arraybackslash}p{(\columnwidth - 6\tabcolsep) * \real{0.18}}
  >{\raggedright\arraybackslash}p{(\columnwidth - 6\tabcolsep) * \real{0.24}}
  >{\raggedright\arraybackslash}p{(\columnwidth - 6\tabcolsep) * \real{0.18}}
  >{\raggedright\arraybackslash}p{(\columnwidth - 6\tabcolsep) * \real{0.40}}@{}}
\toprule
名称 & Python 类型 & 含义 & 用法 \\
\midrule
\endhead
\passthrough{\lstinline!name!} & \passthrough{\lstinline!str!} & SDK
对象、配置项、服务或动作的名称。允许值由调用它的具体接口定义。 &
\passthrough{\lstinline!value = obj.name!} /
\passthrough{\lstinline!obj.name = value!} \\
\bottomrule
\end{longtable}

\textbf{方法索引}

\begin{longtable}[]{@{}
  >{\raggedright\arraybackslash}p{(\columnwidth - 6\tabcolsep) * \real{0.18}}
  >{\raggedright\arraybackslash}p{(\columnwidth - 6\tabcolsep) * \real{0.24}}
  >{\raggedright\arraybackslash}p{(\columnwidth - 6\tabcolsep) * \real{0.18}}
  >{\raggedright\arraybackslash}p{(\columnwidth - 6\tabcolsep) * \real{0.40}}@{}}
\toprule
方法 & Python 签名 & 状态 & 安全分类 \\
\midrule
\endhead
\protect\hyperlink{unitree-sdk2-cpp-robot-applyleaseparameter-dunder-init-1}{\passthrough{\lstinline!\_\_init\_\_!}}
& \passthrough{\lstinline!def \_\_init\_\_(self) -> None!} &
\passthrough{\lstinline!SIGNATURE\_ONLY!} &
\passthrough{\lstinline!CONSTRUCTION!} \\
\protect\hyperlink{unitree-sdk2-cpp-robot-applyleaseparameter-from-json-1}{\passthrough{\lstinline!from\_json!}}
&
\passthrough{\lstinline!def from\_json(self, json: dict[str, Any]) -> None!}
& \passthrough{\lstinline!SIGNATURE\_ONLY!} &
\passthrough{\lstinline!UNCLASSIFIED!} \\
\protect\hyperlink{unitree-sdk2-cpp-robot-applyleaseparameter-to-json-1}{\passthrough{\lstinline!to\_json!}}
&
\passthrough{\lstinline!def to\_json(self, json: dict[str, Any]) -> None!}
& \passthrough{\lstinline!SIGNATURE\_ONLY!} &
\passthrough{\lstinline!UNCLASSIFIED!} \\
\bottomrule
\end{longtable}

\hypertarget{unitree-sdk2-cpp-robot-applyleaseparameter-dunder-init-1}{%
\paragraph{\texorpdfstring{\texttt{unitree\_sdk2\_cpp.robot.ApplyLeaseParameter.\_\_init\_\_}}{unitree\_sdk2\_cpp.robot.ApplyLeaseParameter.\_\_init\_\_}}\label{unitree-sdk2-cpp-robot-applyleaseparameter-dunder-init-1}}

初始化 \passthrough{\lstinline!ApplyLeaseParameter!}
实例。是否能实际构造取决于下方可用性状态。

\textbf{签名}

\begin{lstlisting}[language=Python]
def __init__(self) -> None
\end{lstlisting}

\textbf{可用性与安全性}

\textbf{\passthrough{\lstinline!SIGNATURE\_ONLY!} /
\passthrough{\lstinline!CONSTRUCTION!}}：当前只有设计期签名，不能假设已存在于运行时扩展。

\textbf{参数}

无显式参数。实例方法中的 \passthrough{\lstinline!self!} 由 Python
自动传入。

\textbf{返回值}

\begin{longtable}[]{@{}
  >{\raggedright\arraybackslash}p{(\columnwidth - 4\tabcolsep) * \real{0.18}}
  >{\raggedright\arraybackslash}p{(\columnwidth - 4\tabcolsep) * \real{0.20}}
  >{\raggedright\arraybackslash}p{(\columnwidth - 4\tabcolsep) * \real{0.62}}@{}}
\toprule
位置 & 类型 & 含义 \\
\midrule
\endhead
无 & \passthrough{\lstinline!None!} &
\passthrough{\lstinline!\_\_init\_\_!}
本身不返回值；调用类对象时，在构造函数可用的前提下得到该类实例。 \\
\bottomrule
\end{longtable}

\textbf{对应 C++}

\begin{itemize}
\tightlist
\item
  类：\passthrough{\lstinline!unitree::robot::ApplyLeaseParameter!}
\item
  签名：\passthrough{\lstinline!ApplyLeaseParameter()!}
\item
  绑定策略：\passthrough{\lstinline!未单独标注!}
\item
  声明位置：\passthrough{\lstinline!include/unitree/robot/internal/internal\_api.hpp:69!}
\end{itemize}

\textbf{说明}

\begin{itemize}
\tightlist
\item
  示例是局部调用片段；\passthrough{\lstinline!obj!}
  和参数变量表示已经按上层业务校验并准备好的对象。
\item
  该条目可以被编辑器和类型检查器识别，但当前运行时可能在属性查找阶段直接失败。
\end{itemize}

\textbf{用法}

\begin{lstlisting}[language=Python]
from typing import TYPE_CHECKING

if TYPE_CHECKING:
    def planned_construct() -> ApplyLeaseParameter:
        return ApplyLeaseParameter()
\end{lstlisting}

\begin{quote}
{[}!CAUTION{]} 上面的 \passthrough{\lstinline!TYPE\_CHECKING!}
示例只表达计划调用形态。该分支在普通 Python
运行时为假，不代表当前扩展能执行此接口。
\end{quote}

\hypertarget{unitree-sdk2-cpp-robot-applyleaseparameter-from-json-1}{%
\paragraph{\texorpdfstring{\texttt{unitree\_sdk2\_cpp.robot.ApplyLeaseParameter.from\_json}}{unitree\_sdk2\_cpp.robot.ApplyLeaseParameter.from\_json}}\label{unitree-sdk2-cpp-robot-applyleaseparameter-from-json-1}}

计划从 JSON 风格字典读取字段并更新当前 SDK 值对象。

\textbf{签名}

\begin{lstlisting}[language=Python]
def from_json(self, json: dict[str, Any]) -> None
\end{lstlisting}

\textbf{可用性与安全性}

\textbf{\passthrough{\lstinline!SIGNATURE\_ONLY!} /
\passthrough{\lstinline!UNCLASSIFIED!}}：当前只有设计期签名，不能假设已存在于运行时扩展。

\textbf{参数}

\begin{longtable}[]{@{}
  >{\raggedright\arraybackslash}p{(\columnwidth - 6\tabcolsep) * \real{0.18}}
  >{\raggedright\arraybackslash}p{(\columnwidth - 6\tabcolsep) * \real{0.24}}
  >{\raggedright\arraybackslash}p{(\columnwidth - 6\tabcolsep) * \real{0.18}}
  >{\raggedright\arraybackslash}p{(\columnwidth - 6\tabcolsep) * \real{0.40}}@{}}
\toprule
名称 & Python 类型 & 默认值 & 含义 \\
\midrule
\endhead
\passthrough{\lstinline!json!} &
\passthrough{\lstinline!dict[str, Any]!} & 必填 & JSON
风格字典，键和值必须满足对应 C++ \passthrough{\lstinline!JsonMap!}
数据结构的约束。 对应 C++ 参数
\passthrough{\lstinline!json: common::JsonMap \&!}。 \\
\bottomrule
\end{longtable}

\textbf{返回值}

\begin{longtable}[]{@{}
  >{\raggedright\arraybackslash}p{(\columnwidth - 4\tabcolsep) * \real{0.18}}
  >{\raggedright\arraybackslash}p{(\columnwidth - 4\tabcolsep) * \real{0.20}}
  >{\raggedright\arraybackslash}p{(\columnwidth - 4\tabcolsep) * \real{0.62}}@{}}
\toprule
位置 & 类型 & 含义 \\
\midrule
\endhead
无 & \passthrough{\lstinline!None!} & 该方法不返回 Python
值；失败可能表现为异常或底层状态变化。 \\
\bottomrule
\end{longtable}

\textbf{对应 C++}

\begin{itemize}
\tightlist
\item
  类：\passthrough{\lstinline!unitree::robot::ApplyLeaseParameter!}
\item
  签名：\passthrough{\lstinline!fromJson(common::JsonMap \&)!}
\item
  绑定策略：\passthrough{\lstinline!SIGNATURE\_PREVIEW!}
\item
  声明位置：\passthrough{\lstinline!include/unitree/robot/internal/internal\_api.hpp:75!}
\end{itemize}

\textbf{说明}

\begin{itemize}
\tightlist
\item
  示例是局部调用片段；\passthrough{\lstinline!obj!}
  和参数变量表示已经按上层业务校验并准备好的对象。
\item
  该条目可以被编辑器和类型检查器识别，但当前运行时可能在属性查找阶段直接失败。
\end{itemize}

\textbf{用法}

\begin{lstlisting}[language=Python]
from typing import TYPE_CHECKING

if TYPE_CHECKING:
    def planned_from_json(obj: ApplyLeaseParameter, json: dict[str, Any]) -> None:
        obj.from_json(json=json)
\end{lstlisting}

\begin{quote}
{[}!CAUTION{]} 上面的 \passthrough{\lstinline!TYPE\_CHECKING!}
示例只表达计划调用形态。该分支在普通 Python
运行时为假，不代表当前扩展能执行此接口。
\end{quote}

\hypertarget{unitree-sdk2-cpp-robot-applyleaseparameter-to-json-1}{%
\paragraph{\texorpdfstring{\texttt{unitree\_sdk2\_cpp.robot.ApplyLeaseParameter.to\_json}}{unitree\_sdk2\_cpp.robot.ApplyLeaseParameter.to\_json}}\label{unitree-sdk2-cpp-robot-applyleaseparameter-to-json-1}}

计划把当前 SDK 值对象写入 JSON 风格字典。

\textbf{签名}

\begin{lstlisting}[language=Python]
def to_json(self, json: dict[str, Any]) -> None
\end{lstlisting}

\textbf{可用性与安全性}

\textbf{\passthrough{\lstinline!SIGNATURE\_ONLY!} /
\passthrough{\lstinline!UNCLASSIFIED!}}：当前只有设计期签名，不能假设已存在于运行时扩展。

\textbf{参数}

\begin{longtable}[]{@{}
  >{\raggedright\arraybackslash}p{(\columnwidth - 6\tabcolsep) * \real{0.18}}
  >{\raggedright\arraybackslash}p{(\columnwidth - 6\tabcolsep) * \real{0.24}}
  >{\raggedright\arraybackslash}p{(\columnwidth - 6\tabcolsep) * \real{0.18}}
  >{\raggedright\arraybackslash}p{(\columnwidth - 6\tabcolsep) * \real{0.40}}@{}}
\toprule
名称 & Python 类型 & 默认值 & 含义 \\
\midrule
\endhead
\passthrough{\lstinline!json!} &
\passthrough{\lstinline!dict[str, Any]!} & 必填 & JSON
风格字典，键和值必须满足对应 C++ \passthrough{\lstinline!JsonMap!}
数据结构的约束。 对应 C++ 参数
\passthrough{\lstinline!json: common::JsonMap \&!}。 \\
\bottomrule
\end{longtable}

\textbf{返回值}

\begin{longtable}[]{@{}
  >{\raggedright\arraybackslash}p{(\columnwidth - 4\tabcolsep) * \real{0.18}}
  >{\raggedright\arraybackslash}p{(\columnwidth - 4\tabcolsep) * \real{0.20}}
  >{\raggedright\arraybackslash}p{(\columnwidth - 4\tabcolsep) * \real{0.62}}@{}}
\toprule
位置 & 类型 & 含义 \\
\midrule
\endhead
无 & \passthrough{\lstinline!None!} & 该方法不返回 Python
值；失败可能表现为异常或底层状态变化。 \\
\bottomrule
\end{longtable}

\textbf{对应 C++}

\begin{itemize}
\tightlist
\item
  类：\passthrough{\lstinline!unitree::robot::ApplyLeaseParameter!}
\item
  签名：\passthrough{\lstinline!toJson(common::JsonMap \&) const!}
\item
  绑定策略：\passthrough{\lstinline!SIGNATURE\_PREVIEW!}
\item
  声明位置：\passthrough{\lstinline!include/unitree/robot/internal/internal\_api.hpp:80!}
\end{itemize}

\textbf{说明}

\begin{itemize}
\tightlist
\item
  示例是局部调用片段；\passthrough{\lstinline!obj!}
  和参数变量表示已经按上层业务校验并准备好的对象。
\item
  该条目可以被编辑器和类型检查器识别，但当前运行时可能在属性查找阶段直接失败。
\end{itemize}

\textbf{用法}

\begin{lstlisting}[language=Python]
from typing import TYPE_CHECKING

if TYPE_CHECKING:
    def planned_to_json(obj: ApplyLeaseParameter, json: dict[str, Any]) -> None:
        obj.to_json(json=json)
\end{lstlisting}

\begin{quote}
{[}!CAUTION{]} 上面的 \passthrough{\lstinline!TYPE\_CHECKING!}
示例只表达计划调用形态。该分支在普通 Python
运行时为假，不代表当前扩展能执行此接口。
\end{quote}

\hypertarget{unitree-sdk2-cpp-robot-channelfactory}{%
\subsubsection{\texorpdfstring{\texttt{unitree\_sdk2\_cpp.robot.ChannelFactory}}{unitree\_sdk2\_cpp.robot.ChannelFactory}}\label{unitree-sdk2-cpp-robot-channelfactory}}

SDK 类型签名预览。请逐项查看构造函数和方法的可用性。

\textbf{导入}

\begin{lstlisting}[language=Python]
from unitree_sdk2_cpp.robot import ChannelFactory
\end{lstlisting}

\textbf{基类}：\passthrough{\lstinline!object!}

\textbf{方法索引}

\begin{longtable}[]{@{}
  >{\raggedright\arraybackslash}p{(\columnwidth - 6\tabcolsep) * \real{0.18}}
  >{\raggedright\arraybackslash}p{(\columnwidth - 6\tabcolsep) * \real{0.24}}
  >{\raggedright\arraybackslash}p{(\columnwidth - 6\tabcolsep) * \real{0.18}}
  >{\raggedright\arraybackslash}p{(\columnwidth - 6\tabcolsep) * \real{0.40}}@{}}
\toprule
方法 & Python 签名 & 状态 & 安全分类 \\
\midrule
\endhead
\protect\hyperlink{unitree-sdk2-cpp-robot-channelfactory-instance-1}{\passthrough{\lstinline!instance!}}
& \passthrough{\lstinline!def instance(self) -> ChannelFactory!} &
\passthrough{\lstinline!SIGNATURE\_ONLY!} &
\passthrough{\lstinline!UNCLASSIFIED!} \\
\protect\hyperlink{unitree-sdk2-cpp-robot-channelfactory-init-1}{\passthrough{\lstinline!init（重载 1/3）!}}
&
\passthrough{\lstinline!def init(self, domain\_id: int, network\_interface: str = ...) -> None!}
& \passthrough{\lstinline!SIGNATURE\_ONLY!} &
\passthrough{\lstinline!UNCLASSIFIED!} \\
\protect\hyperlink{unitree-sdk2-cpp-robot-channelfactory-init-2}{\passthrough{\lstinline!init（重载 2/3）!}}
&
\passthrough{\lstinline!def init(self, config\_file\_name: str = ...) -> None!}
& \passthrough{\lstinline!SIGNATURE\_ONLY!} &
\passthrough{\lstinline!UNCLASSIFIED!} \\
\protect\hyperlink{unitree-sdk2-cpp-robot-channelfactory-init-3}{\passthrough{\lstinline!init（重载 3/3）!}}
&
\passthrough{\lstinline!def init(self, json\_map: dict[str, Any]) -> None!}
& \passthrough{\lstinline!SIGNATURE\_ONLY!} &
\passthrough{\lstinline!UNCLASSIFIED!} \\
\protect\hyperlink{unitree-sdk2-cpp-robot-channelfactory-release-1}{\passthrough{\lstinline!release!}}
& \passthrough{\lstinline!def release(self) -> None!} &
\passthrough{\lstinline!SIGNATURE\_ONLY!} &
\passthrough{\lstinline!UNCLASSIFIED!} \\
\bottomrule
\end{longtable}

\hypertarget{unitree-sdk2-cpp-robot-channelfactory-instance-1}{%
\paragraph{\texorpdfstring{\texttt{unitree\_sdk2\_cpp.robot.ChannelFactory.instance}}{unitree\_sdk2\_cpp.robot.ChannelFactory.instance}}\label{unitree-sdk2-cpp-robot-channelfactory-instance-1}}

对应 C++ SDK 操作
\passthrough{\lstinline!Instance()!}。上游头文件没有可直接生成的业务说明时，本参考只保证签名映射，精确语义需查目标型号协议。

\textbf{签名}

\begin{lstlisting}[language=Python]
def instance(self) -> ChannelFactory
\end{lstlisting}

\textbf{可用性与安全性}

\textbf{\passthrough{\lstinline!SIGNATURE\_ONLY!} /
\passthrough{\lstinline!UNCLASSIFIED!}}：当前只有设计期签名，不能假设已存在于运行时扩展。

\textbf{参数}

无显式参数。实例方法中的 \passthrough{\lstinline!self!} 由 Python
自动传入。

\textbf{返回值}

\begin{longtable}[]{@{}
  >{\raggedright\arraybackslash}p{(\columnwidth - 4\tabcolsep) * \real{0.18}}
  >{\raggedright\arraybackslash}p{(\columnwidth - 4\tabcolsep) * \real{0.20}}
  >{\raggedright\arraybackslash}p{(\columnwidth - 4\tabcolsep) * \real{0.62}}@{}}
\toprule
位置 & 类型 & 含义 \\
\midrule
\endhead
返回值 & \passthrough{\lstinline!ChannelFactory!} & 返回
\passthrough{\lstinline!ChannelFactory!}。更精确的业务含义见该方法用途、C++
签名和目标型号协议。 \\
\bottomrule
\end{longtable}

\textbf{对应 C++}

\begin{itemize}
\tightlist
\item
  类：\passthrough{\lstinline!unitree::robot::ChannelFactory!}
\item
  签名：\passthrough{\lstinline!Instance()!}
\item
  绑定策略：\passthrough{\lstinline!SIGNATURE\_PREVIEW!}
\item
  声明位置：\passthrough{\lstinline!include/unitree/robot/channel/channel\_factory.hpp:19!}
\end{itemize}

\textbf{说明}

\begin{itemize}
\tightlist
\item
  示例是局部调用片段；\passthrough{\lstinline!obj!}
  和参数变量表示已经按上层业务校验并准备好的对象。
\item
  该条目可以被编辑器和类型检查器识别，但当前运行时可能在属性查找阶段直接失败。
\end{itemize}

\textbf{用法}

\begin{lstlisting}[language=Python]
from typing import TYPE_CHECKING

if TYPE_CHECKING:
    def planned_instance(obj: ChannelFactory) -> ChannelFactory:
        return obj.instance()
\end{lstlisting}

\begin{quote}
{[}!CAUTION{]} 上面的 \passthrough{\lstinline!TYPE\_CHECKING!}
示例只表达计划调用形态。该分支在普通 Python
运行时为假，不代表当前扩展能执行此接口。
\end{quote}

\hypertarget{unitree-sdk2-cpp-robot-channelfactory-init-1}{%
\paragraph{\texorpdfstring{\texttt{unitree\_sdk2\_cpp.robot.ChannelFactory.init}（重载
1/3）}{unitree\_sdk2\_cpp.robot.ChannelFactory.init（重载 1/3）}}\label{unitree-sdk2-cpp-robot-channelfactory-init-1}}

初始化当前 SDK 对象所需的底层通道或服务资源。

\textbf{签名}

\begin{lstlisting}[language=Python]
@overload
def init(self, domain_id: int, network_interface: str = ...) -> None
\end{lstlisting}

\textbf{可用性与安全性}

\textbf{\passthrough{\lstinline!SIGNATURE\_ONLY!} /
\passthrough{\lstinline!UNCLASSIFIED!}}：当前只有设计期签名，不能假设已存在于运行时扩展。

\textbf{参数}

\begin{longtable}[]{@{}
  >{\raggedright\arraybackslash}p{(\columnwidth - 6\tabcolsep) * \real{0.18}}
  >{\raggedright\arraybackslash}p{(\columnwidth - 6\tabcolsep) * \real{0.24}}
  >{\raggedright\arraybackslash}p{(\columnwidth - 6\tabcolsep) * \real{0.18}}
  >{\raggedright\arraybackslash}p{(\columnwidth - 6\tabcolsep) * \real{0.40}}@{}}
\toprule
名称 & Python 类型 & 默认值 & 含义 \\
\midrule
\endhead
\passthrough{\lstinline!domain\_id!} & \passthrough{\lstinline!int!} &
必填 & DDS Domain ID。只有使用兼容 domain 的参与者才能按预期互相发现。
对应 C++ 参数 \passthrough{\lstinline!domainId: int32\_t!}。 取值范围为
-2147483648 到 2147483647。 \\
\passthrough{\lstinline!network\_interface!} &
\passthrough{\lstinline!str!} & \passthrough{\lstinline!...!} & DDS
使用的网络接口名称，例如 Linux 上的 \passthrough{\lstinline!eth0!} 或
\passthrough{\lstinline!enp3s0!}；必须按目标机实际网卡填写。 对应 C++
参数 \passthrough{\lstinline!networkInterface: const std::string \&!}。
底层为字符串；长度、编码和允许值由具体协议决定。 \\
\bottomrule
\end{longtable}

\textbf{返回值}

\begin{longtable}[]{@{}
  >{\raggedright\arraybackslash}p{(\columnwidth - 4\tabcolsep) * \real{0.18}}
  >{\raggedright\arraybackslash}p{(\columnwidth - 4\tabcolsep) * \real{0.20}}
  >{\raggedright\arraybackslash}p{(\columnwidth - 4\tabcolsep) * \real{0.62}}@{}}
\toprule
位置 & 类型 & 含义 \\
\midrule
\endhead
无 & \passthrough{\lstinline!None!} & 该方法不返回 Python
值；失败可能表现为异常或底层状态变化。 \\
\bottomrule
\end{longtable}

\textbf{对应 C++}

\begin{itemize}
\tightlist
\item
  类：\passthrough{\lstinline!unitree::robot::ChannelFactory!}
\item
  签名：\passthrough{\lstinline!Init(int32\_t, const std::string \&)!}
\item
  绑定策略：\passthrough{\lstinline!SIGNATURE\_PREVIEW!}
\item
  声明位置：\passthrough{\lstinline!include/unitree/robot/channel/channel\_factory.hpp:25!}
\end{itemize}

\textbf{说明}

\begin{itemize}
\tightlist
\item
  示例是局部调用片段；\passthrough{\lstinline!obj!}
  和参数变量表示已经按上层业务校验并准备好的对象。
\item
  默认值显示为 \passthrough{\lstinline!...!}，表示 C++
  声明存在默认参数，但当前 AST
  清单没有保存其字面量；省略参数可使用上游默认行为。
\item
  该条目可以被编辑器和类型检查器识别，但当前运行时可能在属性查找阶段直接失败。
\end{itemize}

\textbf{用法}

\begin{lstlisting}[language=Python]
from typing import TYPE_CHECKING

if TYPE_CHECKING:
    def planned_init(obj: ChannelFactory, domain_id: int, network_interface: str) -> None:
        obj.init(domain_id=domain_id, network_interface=network_interface)
\end{lstlisting}

\begin{quote}
{[}!CAUTION{]} 上面的 \passthrough{\lstinline!TYPE\_CHECKING!}
示例只表达计划调用形态。该分支在普通 Python
运行时为假，不代表当前扩展能执行此接口。
\end{quote}

\hypertarget{unitree-sdk2-cpp-robot-channelfactory-init-2}{%
\paragraph{\texorpdfstring{\texttt{unitree\_sdk2\_cpp.robot.ChannelFactory.init}（重载
2/3）}{unitree\_sdk2\_cpp.robot.ChannelFactory.init（重载 2/3）}}\label{unitree-sdk2-cpp-robot-channelfactory-init-2}}

初始化当前 SDK 对象所需的底层通道或服务资源。

\textbf{签名}

\begin{lstlisting}[language=Python]
@overload
def init(self, config_file_name: str = ...) -> None
\end{lstlisting}

\textbf{可用性与安全性}

\textbf{\passthrough{\lstinline!SIGNATURE\_ONLY!} /
\passthrough{\lstinline!UNCLASSIFIED!}}：当前只有设计期签名，不能假设已存在于运行时扩展。

\textbf{参数}

\begin{longtable}[]{@{}
  >{\raggedright\arraybackslash}p{(\columnwidth - 6\tabcolsep) * \real{0.18}}
  >{\raggedright\arraybackslash}p{(\columnwidth - 6\tabcolsep) * \real{0.24}}
  >{\raggedright\arraybackslash}p{(\columnwidth - 6\tabcolsep) * \real{0.18}}
  >{\raggedright\arraybackslash}p{(\columnwidth - 6\tabcolsep) * \real{0.40}}@{}}
\toprule
名称 & Python 类型 & 默认值 & 含义 \\
\midrule
\endhead
\passthrough{\lstinline!config\_file\_name!} &
\passthrough{\lstinline!str!} & \passthrough{\lstinline!...!} &
传给该接口的 配置 \passthrough{\lstinline!file!} 名称 参数，Python
类型为
\passthrough{\lstinline!str!}。精确含义、单位、范围和枚举值需查目标型号对应头文件与协议。
对应 C++ 参数
\passthrough{\lstinline!configFileName: const std::string \&!}。
底层为字符串；长度、编码和允许值由具体协议决定。 \\
\bottomrule
\end{longtable}

\textbf{返回值}

\begin{longtable}[]{@{}
  >{\raggedright\arraybackslash}p{(\columnwidth - 4\tabcolsep) * \real{0.18}}
  >{\raggedright\arraybackslash}p{(\columnwidth - 4\tabcolsep) * \real{0.20}}
  >{\raggedright\arraybackslash}p{(\columnwidth - 4\tabcolsep) * \real{0.62}}@{}}
\toprule
位置 & 类型 & 含义 \\
\midrule
\endhead
无 & \passthrough{\lstinline!None!} & 该方法不返回 Python
值；失败可能表现为异常或底层状态变化。 \\
\bottomrule
\end{longtable}

\textbf{对应 C++}

\begin{itemize}
\tightlist
\item
  类：\passthrough{\lstinline!unitree::robot::ChannelFactory!}
\item
  签名：\passthrough{\lstinline!Init(const std::string \&)!}
\item
  绑定策略：\passthrough{\lstinline!SIGNATURE\_PREVIEW!}
\item
  声明位置：\passthrough{\lstinline!include/unitree/robot/channel/channel\_factory.hpp:26!}
\end{itemize}

\textbf{说明}

\begin{itemize}
\tightlist
\item
  示例是局部调用片段；\passthrough{\lstinline!obj!}
  和参数变量表示已经按上层业务校验并准备好的对象。
\item
  默认值显示为 \passthrough{\lstinline!...!}，表示 C++
  声明存在默认参数，但当前 AST
  清单没有保存其字面量；省略参数可使用上游默认行为。
\item
  该条目可以被编辑器和类型检查器识别，但当前运行时可能在属性查找阶段直接失败。
\end{itemize}

\textbf{用法}

\begin{lstlisting}[language=Python]
from typing import TYPE_CHECKING

if TYPE_CHECKING:
    def planned_init(obj: ChannelFactory, config_file_name: str) -> None:
        obj.init(config_file_name=config_file_name)
\end{lstlisting}

\begin{quote}
{[}!CAUTION{]} 上面的 \passthrough{\lstinline!TYPE\_CHECKING!}
示例只表达计划调用形态。该分支在普通 Python
运行时为假，不代表当前扩展能执行此接口。
\end{quote}

\hypertarget{unitree-sdk2-cpp-robot-channelfactory-init-3}{%
\paragraph{\texorpdfstring{\texttt{unitree\_sdk2\_cpp.robot.ChannelFactory.init}（重载
3/3）}{unitree\_sdk2\_cpp.robot.ChannelFactory.init（重载 3/3）}}\label{unitree-sdk2-cpp-robot-channelfactory-init-3}}

初始化当前 SDK 对象所需的底层通道或服务资源。

\textbf{签名}

\begin{lstlisting}[language=Python]
@overload
def init(self, json_map: dict[str, Any]) -> None
\end{lstlisting}

\textbf{可用性与安全性}

\textbf{\passthrough{\lstinline!SIGNATURE\_ONLY!} /
\passthrough{\lstinline!UNCLASSIFIED!}}：当前只有设计期签名，不能假设已存在于运行时扩展。

\textbf{参数}

\begin{longtable}[]{@{}
  >{\raggedright\arraybackslash}p{(\columnwidth - 6\tabcolsep) * \real{0.18}}
  >{\raggedright\arraybackslash}p{(\columnwidth - 6\tabcolsep) * \real{0.24}}
  >{\raggedright\arraybackslash}p{(\columnwidth - 6\tabcolsep) * \real{0.18}}
  >{\raggedright\arraybackslash}p{(\columnwidth - 6\tabcolsep) * \real{0.40}}@{}}
\toprule
名称 & Python 类型 & 默认值 & 含义 \\
\midrule
\endhead
\passthrough{\lstinline!json\_map!} &
\passthrough{\lstinline!dict[str, Any]!} & 必填 & 传给该接口的
\passthrough{\lstinline!json!} \passthrough{\lstinline!map!}
参数，Python 类型为
\passthrough{\lstinline!dict[str, Any]!}。精确含义、单位、范围和枚举值需查目标型号对应头文件与协议。
对应 C++ 参数
\passthrough{\lstinline!jsonMap: const common::JsonMap \&!}。 \\
\bottomrule
\end{longtable}

\textbf{返回值}

\begin{longtable}[]{@{}
  >{\raggedright\arraybackslash}p{(\columnwidth - 4\tabcolsep) * \real{0.18}}
  >{\raggedright\arraybackslash}p{(\columnwidth - 4\tabcolsep) * \real{0.20}}
  >{\raggedright\arraybackslash}p{(\columnwidth - 4\tabcolsep) * \real{0.62}}@{}}
\toprule
位置 & 类型 & 含义 \\
\midrule
\endhead
无 & \passthrough{\lstinline!None!} & 该方法不返回 Python
值；失败可能表现为异常或底层状态变化。 \\
\bottomrule
\end{longtable}

\textbf{对应 C++}

\begin{itemize}
\tightlist
\item
  类：\passthrough{\lstinline!unitree::robot::ChannelFactory!}
\item
  签名：\passthrough{\lstinline!Init(const common::JsonMap \&)!}
\item
  绑定策略：\passthrough{\lstinline!SIGNATURE\_PREVIEW!}
\item
  声明位置：\passthrough{\lstinline!include/unitree/robot/channel/channel\_factory.hpp:27!}
\end{itemize}

\textbf{说明}

\begin{itemize}
\tightlist
\item
  示例是局部调用片段；\passthrough{\lstinline!obj!}
  和参数变量表示已经按上层业务校验并准备好的对象。
\item
  该条目可以被编辑器和类型检查器识别，但当前运行时可能在属性查找阶段直接失败。
\end{itemize}

\textbf{用法}

\begin{lstlisting}[language=Python]
from typing import TYPE_CHECKING

if TYPE_CHECKING:
    def planned_init(obj: ChannelFactory, json_map: dict[str, Any]) -> None:
        obj.init(json_map=json_map)
\end{lstlisting}

\begin{quote}
{[}!CAUTION{]} 上面的 \passthrough{\lstinline!TYPE\_CHECKING!}
示例只表达计划调用形态。该分支在普通 Python
运行时为假，不代表当前扩展能执行此接口。
\end{quote}

\hypertarget{unitree-sdk2-cpp-robot-channelfactory-release-1}{%
\paragraph{\texorpdfstring{\texttt{unitree\_sdk2\_cpp.robot.ChannelFactory.release}}{unitree\_sdk2\_cpp.robot.ChannelFactory.release}}\label{unitree-sdk2-cpp-robot-channelfactory-release-1}}

对应 C++ SDK 操作
\passthrough{\lstinline!Release()!}。上游头文件没有可直接生成的业务说明时，本参考只保证签名映射，精确语义需查目标型号协议。

\textbf{签名}

\begin{lstlisting}[language=Python]
def release(self) -> None
\end{lstlisting}

\textbf{可用性与安全性}

\textbf{\passthrough{\lstinline!SIGNATURE\_ONLY!} /
\passthrough{\lstinline!UNCLASSIFIED!}}：当前只有设计期签名，不能假设已存在于运行时扩展。

\textbf{参数}

无显式参数。实例方法中的 \passthrough{\lstinline!self!} 由 Python
自动传入。

\textbf{返回值}

\begin{longtable}[]{@{}
  >{\raggedright\arraybackslash}p{(\columnwidth - 4\tabcolsep) * \real{0.18}}
  >{\raggedright\arraybackslash}p{(\columnwidth - 4\tabcolsep) * \real{0.20}}
  >{\raggedright\arraybackslash}p{(\columnwidth - 4\tabcolsep) * \real{0.62}}@{}}
\toprule
位置 & 类型 & 含义 \\
\midrule
\endhead
无 & \passthrough{\lstinline!None!} & 该方法不返回 Python
值；失败可能表现为异常或底层状态变化。 \\
\bottomrule
\end{longtable}

\textbf{对应 C++}

\begin{itemize}
\tightlist
\item
  类：\passthrough{\lstinline!unitree::robot::ChannelFactory!}
\item
  签名：\passthrough{\lstinline!Release()!}
\item
  绑定策略：\passthrough{\lstinline!SIGNATURE\_PREVIEW!}
\item
  声明位置：\passthrough{\lstinline!include/unitree/robot/channel/channel\_factory.hpp:29!}
\end{itemize}

\textbf{说明}

\begin{itemize}
\tightlist
\item
  示例是局部调用片段；\passthrough{\lstinline!obj!}
  和参数变量表示已经按上层业务校验并准备好的对象。
\item
  该条目可以被编辑器和类型检查器识别，但当前运行时可能在属性查找阶段直接失败。
\end{itemize}

\textbf{用法}

\begin{lstlisting}[language=Python]
from typing import TYPE_CHECKING

if TYPE_CHECKING:
    def planned_release(obj: ChannelFactory) -> None:
        obj.release()
\end{lstlisting}

\begin{quote}
{[}!CAUTION{]} 上面的 \passthrough{\lstinline!TYPE\_CHECKING!}
示例只表达计划调用形态。该分支在普通 Python
运行时为假，不代表当前扩展能执行此接口。
\end{quote}

\hypertarget{unitree-sdk2-cpp-robot-channelnamer}{%
\subsubsection{\texorpdfstring{\texttt{unitree\_sdk2\_cpp.robot.ChannelNamer}}{unitree\_sdk2\_cpp.robot.ChannelNamer}}\label{unitree-sdk2-cpp-robot-channelnamer}}

SDK 类型签名预览。请逐项查看构造函数和方法的可用性。

\textbf{导入}

\begin{lstlisting}[language=Python]
from unitree_sdk2_cpp.robot import ChannelNamer
\end{lstlisting}

\textbf{基类}：\passthrough{\lstinline!object!}

\textbf{方法索引}

\begin{longtable}[]{@{}
  >{\raggedright\arraybackslash}p{(\columnwidth - 6\tabcolsep) * \real{0.18}}
  >{\raggedright\arraybackslash}p{(\columnwidth - 6\tabcolsep) * \real{0.24}}
  >{\raggedright\arraybackslash}p{(\columnwidth - 6\tabcolsep) * \real{0.18}}
  >{\raggedright\arraybackslash}p{(\columnwidth - 6\tabcolsep) * \real{0.40}}@{}}
\toprule
方法 & Python 签名 & 状态 & 安全分类 \\
\midrule
\endhead
\protect\hyperlink{unitree-sdk2-cpp-robot-channelnamer-get-send-channel-name-1}{\passthrough{\lstinline!get\_send\_channel\_name!}}
&
\passthrough{\lstinline!def get\_send\_channel\_name(self, name: str) -> str!}
& \passthrough{\lstinline!SIGNATURE\_ONLY!} &
\passthrough{\lstinline!UNCLASSIFIED!} \\
\protect\hyperlink{unitree-sdk2-cpp-robot-channelnamer-get-recv-channel-name-1}{\passthrough{\lstinline!get\_recv\_channel\_name!}}
&
\passthrough{\lstinline!def get\_recv\_channel\_name(self, name: str) -> str!}
& \passthrough{\lstinline!SIGNATURE\_ONLY!} &
\passthrough{\lstinline!UNCLASSIFIED!} \\
\bottomrule
\end{longtable}

\hypertarget{unitree-sdk2-cpp-robot-channelnamer-get-send-channel-name-1}{%
\paragraph{\texorpdfstring{\texttt{unitree\_sdk2\_cpp.robot.ChannelNamer.get\_send\_channel\_name}}{unitree\_sdk2\_cpp.robot.ChannelNamer.get\_send\_channel\_name}}\label{unitree-sdk2-cpp-robot-channelnamer-get-send-channel-name-1}}

查询或检查 \passthrough{\lstinline!send!}
\passthrough{\lstinline!channel!} 名称。

\textbf{签名}

\begin{lstlisting}[language=Python]
def get_send_channel_name(self, name: str) -> str
\end{lstlisting}

\textbf{可用性与安全性}

\textbf{\passthrough{\lstinline!SIGNATURE\_ONLY!} /
\passthrough{\lstinline!UNCLASSIFIED!}}：当前只有设计期签名，不能假设已存在于运行时扩展。

\textbf{参数}

\begin{longtable}[]{@{}
  >{\raggedright\arraybackslash}p{(\columnwidth - 6\tabcolsep) * \real{0.18}}
  >{\raggedright\arraybackslash}p{(\columnwidth - 6\tabcolsep) * \real{0.24}}
  >{\raggedright\arraybackslash}p{(\columnwidth - 6\tabcolsep) * \real{0.18}}
  >{\raggedright\arraybackslash}p{(\columnwidth - 6\tabcolsep) * \real{0.40}}@{}}
\toprule
名称 & Python 类型 & 默认值 & 含义 \\
\midrule
\endhead
\passthrough{\lstinline!name!} & \passthrough{\lstinline!str!} & 必填 &
SDK 对象、配置项、服务或动作的名称。允许值由调用它的具体接口定义。 对应
C++ 参数 \passthrough{\lstinline!name: const std::string \&!}。
底层为字符串；长度、编码和允许值由具体协议决定。 \\
\bottomrule
\end{longtable}

\textbf{返回值}

\begin{longtable}[]{@{}
  >{\raggedright\arraybackslash}p{(\columnwidth - 4\tabcolsep) * \real{0.18}}
  >{\raggedright\arraybackslash}p{(\columnwidth - 4\tabcolsep) * \real{0.20}}
  >{\raggedright\arraybackslash}p{(\columnwidth - 4\tabcolsep) * \real{0.62}}@{}}
\toprule
位置 & 类型 & 含义 \\
\midrule
\endhead
返回值 & \passthrough{\lstinline!str!} & 返回
\passthrough{\lstinline!str!}。更精确的业务含义见该方法用途、C++
签名和目标型号协议。 \\
\bottomrule
\end{longtable}

\textbf{对应 C++}

\begin{itemize}
\tightlist
\item
  类：\passthrough{\lstinline!unitree::robot::ChannelNamer!}
\item
  签名：\passthrough{\lstinline!GetSendChannelName(const std::string \&)!}
\item
  绑定策略：\passthrough{\lstinline!SIGNATURE\_PREVIEW!}
\item
  声明位置：\passthrough{\lstinline!include/unitree/robot/channel/channel\_namer.hpp:21!}
\end{itemize}

\textbf{说明}

\begin{itemize}
\tightlist
\item
  示例是局部调用片段；\passthrough{\lstinline!obj!}
  和参数变量表示已经按上层业务校验并准备好的对象。
\item
  该条目可以被编辑器和类型检查器识别，但当前运行时可能在属性查找阶段直接失败。
\end{itemize}

\textbf{用法}

\begin{lstlisting}[language=Python]
from typing import TYPE_CHECKING

if TYPE_CHECKING:
    def planned_get_send_channel_name(obj: ChannelNamer, name: str) -> str:
        return obj.get_send_channel_name(name=name)
\end{lstlisting}

\begin{quote}
{[}!CAUTION{]} 上面的 \passthrough{\lstinline!TYPE\_CHECKING!}
示例只表达计划调用形态。该分支在普通 Python
运行时为假，不代表当前扩展能执行此接口。
\end{quote}

\hypertarget{unitree-sdk2-cpp-robot-channelnamer-get-recv-channel-name-1}{%
\paragraph{\texorpdfstring{\texttt{unitree\_sdk2\_cpp.robot.ChannelNamer.get\_recv\_channel\_name}}{unitree\_sdk2\_cpp.robot.ChannelNamer.get\_recv\_channel\_name}}\label{unitree-sdk2-cpp-robot-channelnamer-get-recv-channel-name-1}}

查询或检查 \passthrough{\lstinline!recv!}
\passthrough{\lstinline!channel!} 名称。

\textbf{签名}

\begin{lstlisting}[language=Python]
def get_recv_channel_name(self, name: str) -> str
\end{lstlisting}

\textbf{可用性与安全性}

\textbf{\passthrough{\lstinline!SIGNATURE\_ONLY!} /
\passthrough{\lstinline!UNCLASSIFIED!}}：当前只有设计期签名，不能假设已存在于运行时扩展。

\textbf{参数}

\begin{longtable}[]{@{}
  >{\raggedright\arraybackslash}p{(\columnwidth - 6\tabcolsep) * \real{0.18}}
  >{\raggedright\arraybackslash}p{(\columnwidth - 6\tabcolsep) * \real{0.24}}
  >{\raggedright\arraybackslash}p{(\columnwidth - 6\tabcolsep) * \real{0.18}}
  >{\raggedright\arraybackslash}p{(\columnwidth - 6\tabcolsep) * \real{0.40}}@{}}
\toprule
名称 & Python 类型 & 默认值 & 含义 \\
\midrule
\endhead
\passthrough{\lstinline!name!} & \passthrough{\lstinline!str!} & 必填 &
SDK 对象、配置项、服务或动作的名称。允许值由调用它的具体接口定义。 对应
C++ 参数 \passthrough{\lstinline!name: const std::string \&!}。
底层为字符串；长度、编码和允许值由具体协议决定。 \\
\bottomrule
\end{longtable}

\textbf{返回值}

\begin{longtable}[]{@{}
  >{\raggedright\arraybackslash}p{(\columnwidth - 4\tabcolsep) * \real{0.18}}
  >{\raggedright\arraybackslash}p{(\columnwidth - 4\tabcolsep) * \real{0.20}}
  >{\raggedright\arraybackslash}p{(\columnwidth - 4\tabcolsep) * \real{0.62}}@{}}
\toprule
位置 & 类型 & 含义 \\
\midrule
\endhead
返回值 & \passthrough{\lstinline!str!} & 返回
\passthrough{\lstinline!str!}。更精确的业务含义见该方法用途、C++
签名和目标型号协议。 \\
\bottomrule
\end{longtable}

\textbf{对应 C++}

\begin{itemize}
\tightlist
\item
  类：\passthrough{\lstinline!unitree::robot::ChannelNamer!}
\item
  签名：\passthrough{\lstinline!GetRecvChannelName(const std::string \&)!}
\item
  绑定策略：\passthrough{\lstinline!SIGNATURE\_PREVIEW!}
\item
  声明位置：\passthrough{\lstinline!include/unitree/robot/channel/channel\_namer.hpp:22!}
\end{itemize}

\textbf{说明}

\begin{itemize}
\tightlist
\item
  示例是局部调用片段；\passthrough{\lstinline!obj!}
  和参数变量表示已经按上层业务校验并准备好的对象。
\item
  该条目可以被编辑器和类型检查器识别，但当前运行时可能在属性查找阶段直接失败。
\end{itemize}

\textbf{用法}

\begin{lstlisting}[language=Python]
from typing import TYPE_CHECKING

if TYPE_CHECKING:
    def planned_get_recv_channel_name(obj: ChannelNamer, name: str) -> str:
        return obj.get_recv_channel_name(name=name)
\end{lstlisting}

\begin{quote}
{[}!CAUTION{]} 上面的 \passthrough{\lstinline!TYPE\_CHECKING!}
示例只表达计划调用形态。该分支在普通 Python
运行时为假，不代表当前扩展能执行此接口。
\end{quote}

\hypertarget{unitree-sdk2-cpp-robot-client}{%
\subsubsection{\texorpdfstring{\texttt{unitree\_sdk2\_cpp.robot.Client}}{unitree\_sdk2\_cpp.robot.Client}}\label{unitree-sdk2-cpp-robot-client}}

Robot 服务客户端。构造、初始化和具体方法的可用性必须分别检查。

\textbf{导入}

\begin{lstlisting}[language=Python]
from unitree_sdk2_cpp.robot import Client
\end{lstlisting}

\textbf{基类}：\passthrough{\lstinline!ClientBase!}

\textbf{方法索引}

\begin{longtable}[]{@{}
  >{\raggedright\arraybackslash}p{(\columnwidth - 6\tabcolsep) * \real{0.18}}
  >{\raggedright\arraybackslash}p{(\columnwidth - 6\tabcolsep) * \real{0.24}}
  >{\raggedright\arraybackslash}p{(\columnwidth - 6\tabcolsep) * \real{0.18}}
  >{\raggedright\arraybackslash}p{(\columnwidth - 6\tabcolsep) * \real{0.40}}@{}}
\toprule
方法 & Python 签名 & 状态 & 安全分类 \\
\midrule
\endhead
\protect\hyperlink{unitree-sdk2-cpp-robot-client-dunder-init-1}{\passthrough{\lstinline!\_\_init\_\_!}}
&
\passthrough{\lstinline!def \_\_init\_\_(self, name: str, enable\_lease: bool = ...) -> None!}
& \passthrough{\lstinline!SIGNATURE\_ONLY!} &
\passthrough{\lstinline!CONSTRUCTION!} \\
\protect\hyperlink{unitree-sdk2-cpp-robot-client-wait-lease-applied-1}{\passthrough{\lstinline!wait\_lease\_applied!}}
& \passthrough{\lstinline!def wait\_lease\_applied(self) -> None!} &
\passthrough{\lstinline!SIGNATURE\_ONLY!} &
\passthrough{\lstinline!INITIALIZATION!} \\
\protect\hyperlink{unitree-sdk2-cpp-robot-client-get-lease-id-1}{\passthrough{\lstinline!get\_lease\_id!}}
& \passthrough{\lstinline!def get\_lease\_id(self) -> int!} &
\passthrough{\lstinline!AVAILABLE!} &
\passthrough{\lstinline!READ\_ONLY!} \\
\protect\hyperlink{unitree-sdk2-cpp-robot-client-get-api-version-1}{\passthrough{\lstinline!get\_api\_version!}}
& \passthrough{\lstinline!def get\_api\_version(self) -> str!} &
\passthrough{\lstinline!AVAILABLE!} &
\passthrough{\lstinline!READ\_ONLY!} \\
\protect\hyperlink{unitree-sdk2-cpp-robot-client-get-server-api-version-1}{\passthrough{\lstinline!get\_server\_api\_version!}}
& \passthrough{\lstinline!def get\_server\_api\_version(self) -> str!} &
\passthrough{\lstinline!AVAILABLE!} &
\passthrough{\lstinline!READ\_ONLY!} \\
\bottomrule
\end{longtable}

\hypertarget{unitree-sdk2-cpp-robot-client-dunder-init-1}{%
\paragraph{\texorpdfstring{\texttt{unitree\_sdk2\_cpp.robot.Client.\_\_init\_\_}}{unitree\_sdk2\_cpp.robot.Client.\_\_init\_\_}}\label{unitree-sdk2-cpp-robot-client-dunder-init-1}}

初始化 \passthrough{\lstinline!Client!}
实例。是否能实际构造取决于下方可用性状态。

\textbf{签名}

\begin{lstlisting}[language=Python]
def __init__(self, name: str, enable_lease: bool = ...) -> None
\end{lstlisting}

\textbf{可用性与安全性}

\textbf{\passthrough{\lstinline!SIGNATURE\_ONLY!} /
\passthrough{\lstinline!CONSTRUCTION!}}：当前只有设计期签名，不能假设已存在于运行时扩展。

\textbf{参数}

\begin{longtable}[]{@{}
  >{\raggedright\arraybackslash}p{(\columnwidth - 6\tabcolsep) * \real{0.18}}
  >{\raggedright\arraybackslash}p{(\columnwidth - 6\tabcolsep) * \real{0.24}}
  >{\raggedright\arraybackslash}p{(\columnwidth - 6\tabcolsep) * \real{0.18}}
  >{\raggedright\arraybackslash}p{(\columnwidth - 6\tabcolsep) * \real{0.40}}@{}}
\toprule
名称 & Python 类型 & 默认值 & 含义 \\
\midrule
\endhead
\passthrough{\lstinline!name!} & \passthrough{\lstinline!str!} & 必填 &
SDK 对象、配置项、服务或动作的名称。允许值由调用它的具体接口定义。 对应
C++ 参数 \passthrough{\lstinline!name: const std::string \&!}。
底层为字符串；长度、编码和允许值由具体协议决定。 \\
\passthrough{\lstinline!enable\_lease!} & \passthrough{\lstinline!bool!}
& \passthrough{\lstinline!...!} & 是否启用 SDK lease 机制；lease
的获得、续期和释放规则以服务协议为准。 对应 C++ 参数
\passthrough{\lstinline!enableLease: bool!}。 只接受布尔语义。 \\
\bottomrule
\end{longtable}

\textbf{返回值}

\begin{longtable}[]{@{}
  >{\raggedright\arraybackslash}p{(\columnwidth - 4\tabcolsep) * \real{0.18}}
  >{\raggedright\arraybackslash}p{(\columnwidth - 4\tabcolsep) * \real{0.20}}
  >{\raggedright\arraybackslash}p{(\columnwidth - 4\tabcolsep) * \real{0.62}}@{}}
\toprule
位置 & 类型 & 含义 \\
\midrule
\endhead
无 & \passthrough{\lstinline!None!} &
\passthrough{\lstinline!\_\_init\_\_!}
本身不返回值；调用类对象时，在构造函数可用的前提下得到该类实例。 \\
\bottomrule
\end{longtable}

\textbf{对应 C++}

\begin{itemize}
\tightlist
\item
  类：\passthrough{\lstinline!unitree::robot::Client!}
\item
  签名：\passthrough{\lstinline!Client(const std::string \&, bool)!}
\item
  绑定策略：\passthrough{\lstinline!未单独标注!}
\item
  声明位置：\passthrough{\lstinline!include/unitree/robot/client/client.hpp:24!}
\end{itemize}

\textbf{说明}

\begin{itemize}
\tightlist
\item
  示例是局部调用片段；\passthrough{\lstinline!obj!}
  和参数变量表示已经按上层业务校验并准备好的对象。
\item
  默认值显示为 \passthrough{\lstinline!...!}，表示 C++
  声明存在默认参数，但当前 AST
  清单没有保存其字面量；省略参数可使用上游默认行为。
\item
  该条目可以被编辑器和类型检查器识别，但当前运行时可能在属性查找阶段直接失败。
\end{itemize}

\textbf{用法}

\begin{lstlisting}[language=Python]
from typing import TYPE_CHECKING

if TYPE_CHECKING:
    def planned_construct(name: str, enable_lease: bool) -> Client:
        return Client(name=name, enable_lease=enable_lease)
\end{lstlisting}

\begin{quote}
{[}!CAUTION{]} 上面的 \passthrough{\lstinline!TYPE\_CHECKING!}
示例只表达计划调用形态。该分支在普通 Python
运行时为假，不代表当前扩展能执行此接口。
\end{quote}

\hypertarget{unitree-sdk2-cpp-robot-client-wait-lease-applied-1}{%
\paragraph{\texorpdfstring{\texttt{unitree\_sdk2\_cpp.robot.Client.wait\_lease\_applied}}{unitree\_sdk2\_cpp.robot.Client.wait\_lease\_applied}}\label{unitree-sdk2-cpp-robot-client-wait-lease-applied-1}}

对应 C++ SDK 操作
\passthrough{\lstinline!WaitLeaseApplied()!}。上游头文件没有可直接生成的业务说明时，本参考只保证签名映射，精确语义需查目标型号协议。

\textbf{签名}

\begin{lstlisting}[language=Python]
def wait_lease_applied(self) -> None
\end{lstlisting}

\textbf{可用性与安全性}

\textbf{\passthrough{\lstinline!SIGNATURE\_ONLY!} /
\passthrough{\lstinline!INITIALIZATION!}}：当前只有设计期签名，不能假设已存在于运行时扩展。

\textbf{参数}

无显式参数。实例方法中的 \passthrough{\lstinline!self!} 由 Python
自动传入。

\textbf{返回值}

\begin{longtable}[]{@{}
  >{\raggedright\arraybackslash}p{(\columnwidth - 4\tabcolsep) * \real{0.18}}
  >{\raggedright\arraybackslash}p{(\columnwidth - 4\tabcolsep) * \real{0.20}}
  >{\raggedright\arraybackslash}p{(\columnwidth - 4\tabcolsep) * \real{0.62}}@{}}
\toprule
位置 & 类型 & 含义 \\
\midrule
\endhead
无 & \passthrough{\lstinline!None!} & 该方法不返回 Python
值；失败可能表现为异常或底层状态变化。 \\
\bottomrule
\end{longtable}

\textbf{对应 C++}

\begin{itemize}
\tightlist
\item
  类：\passthrough{\lstinline!unitree::robot::Client!}
\item
  签名：\passthrough{\lstinline!WaitLeaseApplied()!}
\item
  绑定策略：\passthrough{\lstinline!DIRECT!}
\item
  声明位置：\passthrough{\lstinline!include/unitree/robot/client/client.hpp:27!}
\end{itemize}

\textbf{说明}

\begin{itemize}
\tightlist
\item
  示例是局部调用片段；\passthrough{\lstinline!obj!}
  和参数变量表示已经按上层业务校验并准备好的对象。
\item
  该条目可以被编辑器和类型检查器识别，但当前运行时可能在属性查找阶段直接失败。
\end{itemize}

\textbf{用法}

\begin{lstlisting}[language=Python]
from typing import TYPE_CHECKING

if TYPE_CHECKING:
    def planned_wait_lease_applied(obj: Client) -> None:
        obj.wait_lease_applied()
\end{lstlisting}

\begin{quote}
{[}!CAUTION{]} 上面的 \passthrough{\lstinline!TYPE\_CHECKING!}
示例只表达计划调用形态。该分支在普通 Python
运行时为假，不代表当前扩展能执行此接口。
\end{quote}

\hypertarget{unitree-sdk2-cpp-robot-client-get-lease-id-1}{%
\paragraph{\texorpdfstring{\texttt{unitree\_sdk2\_cpp.robot.Client.get\_lease\_id}}{unitree\_sdk2\_cpp.robot.Client.get\_lease\_id}}\label{unitree-sdk2-cpp-robot-client-get-lease-id-1}}

查询或检查 lease ID。

\textbf{签名}

\begin{lstlisting}[language=Python]
def get_lease_id(self) -> int
\end{lstlisting}

\textbf{可用性与安全性}

\textbf{\passthrough{\lstinline!AVAILABLE!} /
\passthrough{\lstinline!READ\_ONLY!}}：当前绑定源码已实现只读查询。Robot
Client 的服务端查询通常仍需匹配的 Linux
扩展、DDS、网络、机器人和服务；纯本地 getter 则不一定需要全部条件。

\textbf{参数}

无显式参数。实例方法中的 \passthrough{\lstinline!self!} 由 Python
自动传入。

\textbf{返回值}

\begin{longtable}[]{@{}
  >{\raggedright\arraybackslash}p{(\columnwidth - 4\tabcolsep) * \real{0.18}}
  >{\raggedright\arraybackslash}p{(\columnwidth - 4\tabcolsep) * \real{0.20}}
  >{\raggedright\arraybackslash}p{(\columnwidth - 4\tabcolsep) * \real{0.62}}@{}}
\toprule
位置 & 类型 & 含义 \\
\midrule
\endhead
返回值 & \passthrough{\lstinline!int!} & 返回
\passthrough{\lstinline!int!}。更精确的业务含义见该方法用途、C++
签名和目标型号协议。 \\
\bottomrule
\end{longtable}

\textbf{对应 C++}

\begin{itemize}
\tightlist
\item
  类：\passthrough{\lstinline!unitree::robot::Client!}
\item
  签名：\passthrough{\lstinline!GetLeaseId()!}
\item
  绑定策略：\passthrough{\lstinline!DIRECT!}
\item
  声明位置：\passthrough{\lstinline!include/unitree/robot/client/client.hpp:28!}
\end{itemize}

\textbf{说明}

\begin{itemize}
\tightlist
\item
  示例是局部调用片段；\passthrough{\lstinline!obj!}
  和参数变量表示已经按上层业务校验并准备好的对象。
\end{itemize}

\textbf{用法}

\begin{lstlisting}[language=Python]
result = obj.get_lease_id()
\end{lstlisting}

\begin{quote}
{[}!NOTE{]} 只读查询仍会访问
DDS/机器人服务。必须检查返回状态码，并处理超时和网络中断。
\end{quote}

\hypertarget{unitree-sdk2-cpp-robot-client-get-api-version-1}{%
\paragraph{\texorpdfstring{\texttt{unitree\_sdk2\_cpp.robot.Client.get\_api\_version}}{unitree\_sdk2\_cpp.robot.Client.get\_api\_version}}\label{unitree-sdk2-cpp-robot-client-get-api-version-1}}

查询或检查 API \passthrough{\lstinline!version!}。

\textbf{签名}

\begin{lstlisting}[language=Python]
def get_api_version(self) -> str
\end{lstlisting}

\textbf{可用性与安全性}

\textbf{\passthrough{\lstinline!AVAILABLE!} /
\passthrough{\lstinline!READ\_ONLY!}}：当前绑定源码已实现只读查询。Robot
Client 的服务端查询通常仍需匹配的 Linux
扩展、DDS、网络、机器人和服务；纯本地 getter 则不一定需要全部条件。

\textbf{参数}

无显式参数。实例方法中的 \passthrough{\lstinline!self!} 由 Python
自动传入。

\textbf{返回值}

\begin{longtable}[]{@{}
  >{\raggedright\arraybackslash}p{(\columnwidth - 4\tabcolsep) * \real{0.18}}
  >{\raggedright\arraybackslash}p{(\columnwidth - 4\tabcolsep) * \real{0.20}}
  >{\raggedright\arraybackslash}p{(\columnwidth - 4\tabcolsep) * \real{0.62}}@{}}
\toprule
位置 & 类型 & 含义 \\
\midrule
\endhead
返回值 & \passthrough{\lstinline!str!} & 返回
\passthrough{\lstinline!str!}。更精确的业务含义见该方法用途、C++
签名和目标型号协议。 \\
\bottomrule
\end{longtable}

\textbf{对应 C++}

\begin{itemize}
\tightlist
\item
  类：\passthrough{\lstinline!unitree::robot::Client!}
\item
  签名：\passthrough{\lstinline!GetApiVersion() const!}
\item
  绑定策略：\passthrough{\lstinline!REFERENCE\_POLICY!}
\item
  声明位置：\passthrough{\lstinline!include/unitree/robot/client/client.hpp:30!}
\end{itemize}

\textbf{说明}

\begin{itemize}
\tightlist
\item
  示例是局部调用片段；\passthrough{\lstinline!obj!}
  和参数变量表示已经按上层业务校验并准备好的对象。
\end{itemize}

\textbf{用法}

\begin{lstlisting}[language=Python]
result = obj.get_api_version()
\end{lstlisting}

\begin{quote}
{[}!NOTE{]} 只读查询仍会访问
DDS/机器人服务。必须检查返回状态码，并处理超时和网络中断。
\end{quote}

\hypertarget{unitree-sdk2-cpp-robot-client-get-server-api-version-1}{%
\paragraph{\texorpdfstring{\texttt{unitree\_sdk2\_cpp.robot.Client.get\_server\_api\_version}}{unitree\_sdk2\_cpp.robot.Client.get\_server\_api\_version}}\label{unitree-sdk2-cpp-robot-client-get-server-api-version-1}}

查询或检查 \passthrough{\lstinline!server!} API
\passthrough{\lstinline!version!}。

\textbf{签名}

\begin{lstlisting}[language=Python]
def get_server_api_version(self) -> str
\end{lstlisting}

\textbf{可用性与安全性}

\textbf{\passthrough{\lstinline!AVAILABLE!} /
\passthrough{\lstinline!READ\_ONLY!}}：当前绑定源码已实现只读查询。Robot
Client 的服务端查询通常仍需匹配的 Linux
扩展、DDS、网络、机器人和服务；纯本地 getter 则不一定需要全部条件。

\textbf{参数}

无显式参数。实例方法中的 \passthrough{\lstinline!self!} 由 Python
自动传入。

\textbf{返回值}

\begin{longtable}[]{@{}
  >{\raggedright\arraybackslash}p{(\columnwidth - 4\tabcolsep) * \real{0.18}}
  >{\raggedright\arraybackslash}p{(\columnwidth - 4\tabcolsep) * \real{0.20}}
  >{\raggedright\arraybackslash}p{(\columnwidth - 4\tabcolsep) * \real{0.62}}@{}}
\toprule
位置 & 类型 & 含义 \\
\midrule
\endhead
返回值 & \passthrough{\lstinline!str!} & 返回
\passthrough{\lstinline!str!}。更精确的业务含义见该方法用途、C++
签名和目标型号协议。 \\
\bottomrule
\end{longtable}

\textbf{对应 C++}

\begin{itemize}
\tightlist
\item
  类：\passthrough{\lstinline!unitree::robot::Client!}
\item
  签名：\passthrough{\lstinline!GetServerApiVersion()!}
\item
  绑定策略：\passthrough{\lstinline!DIRECT!}
\item
  声明位置：\passthrough{\lstinline!include/unitree/robot/client/client.hpp:31!}
\end{itemize}

\textbf{说明}

\begin{itemize}
\tightlist
\item
  示例是局部调用片段；\passthrough{\lstinline!obj!}
  和参数变量表示已经按上层业务校验并准备好的对象。
\end{itemize}

\textbf{用法}

\begin{lstlisting}[language=Python]
result = obj.get_server_api_version()
\end{lstlisting}

\begin{quote}
{[}!NOTE{]} 只读查询仍会访问
DDS/机器人服务。必须检查返回状态码，并处理超时和网络中断。
\end{quote}

\hypertarget{unitree-sdk2-cpp-robot-clientbase}{%
\subsubsection{\texorpdfstring{\texttt{unitree\_sdk2\_cpp.robot.ClientBase}}{unitree\_sdk2\_cpp.robot.ClientBase}}\label{unitree-sdk2-cpp-robot-clientbase}}

SDK 类型签名预览。请逐项查看构造函数和方法的可用性。

\textbf{导入}

\begin{lstlisting}[language=Python]
from unitree_sdk2_cpp.robot import ClientBase
\end{lstlisting}

\textbf{基类}：\passthrough{\lstinline!object!}

\textbf{方法索引}

\begin{longtable}[]{@{}
  >{\raggedright\arraybackslash}p{(\columnwidth - 6\tabcolsep) * \real{0.18}}
  >{\raggedright\arraybackslash}p{(\columnwidth - 6\tabcolsep) * \real{0.24}}
  >{\raggedright\arraybackslash}p{(\columnwidth - 6\tabcolsep) * \real{0.18}}
  >{\raggedright\arraybackslash}p{(\columnwidth - 6\tabcolsep) * \real{0.40}}@{}}
\toprule
方法 & Python 签名 & 状态 & 安全分类 \\
\midrule
\endhead
\protect\hyperlink{unitree-sdk2-cpp-robot-clientbase-dunder-init-1}{\passthrough{\lstinline!\_\_init\_\_!}}
& \passthrough{\lstinline!def \_\_init\_\_(self, name: str) -> None!} &
\passthrough{\lstinline!SIGNATURE\_ONLY!} &
\passthrough{\lstinline!CONSTRUCTION!} \\
\protect\hyperlink{unitree-sdk2-cpp-robot-clientbase-init-1}{\passthrough{\lstinline!init!}}
& \passthrough{\lstinline!def init(self) -> None!} &
\passthrough{\lstinline!SIGNATURE\_ONLY!} &
\passthrough{\lstinline!INITIALIZATION!} \\
\protect\hyperlink{unitree-sdk2-cpp-robot-clientbase-set-timeout-microseconds-1}{\passthrough{\lstinline!set\_timeout\_microseconds!}}
&
\passthrough{\lstinline!def set\_timeout\_microseconds(self, microseconds: int) -> None!}
& \passthrough{\lstinline!AVAILABLE!} &
\passthrough{\lstinline!INITIALIZATION!} \\
\protect\hyperlink{unitree-sdk2-cpp-robot-clientbase-set-timeout-1}{\passthrough{\lstinline!set\_timeout!}}
&
\passthrough{\lstinline!def set\_timeout(self, seconds: float) -> None!}
& \passthrough{\lstinline!AVAILABLE!} &
\passthrough{\lstinline!INITIALIZATION!} \\
\bottomrule
\end{longtable}

\hypertarget{unitree-sdk2-cpp-robot-clientbase-dunder-init-1}{%
\paragraph{\texorpdfstring{\texttt{unitree\_sdk2\_cpp.robot.ClientBase.\_\_init\_\_}}{unitree\_sdk2\_cpp.robot.ClientBase.\_\_init\_\_}}\label{unitree-sdk2-cpp-robot-clientbase-dunder-init-1}}

初始化 \passthrough{\lstinline!ClientBase!}
实例。是否能实际构造取决于下方可用性状态。

\textbf{签名}

\begin{lstlisting}[language=Python]
def __init__(self, name: str) -> None
\end{lstlisting}

\textbf{可用性与安全性}

\textbf{\passthrough{\lstinline!SIGNATURE\_ONLY!} /
\passthrough{\lstinline!CONSTRUCTION!}}：当前只有设计期签名，不能假设已存在于运行时扩展。

\textbf{参数}

\begin{longtable}[]{@{}
  >{\raggedright\arraybackslash}p{(\columnwidth - 6\tabcolsep) * \real{0.18}}
  >{\raggedright\arraybackslash}p{(\columnwidth - 6\tabcolsep) * \real{0.24}}
  >{\raggedright\arraybackslash}p{(\columnwidth - 6\tabcolsep) * \real{0.18}}
  >{\raggedright\arraybackslash}p{(\columnwidth - 6\tabcolsep) * \real{0.40}}@{}}
\toprule
名称 & Python 类型 & 默认值 & 含义 \\
\midrule
\endhead
\passthrough{\lstinline!name!} & \passthrough{\lstinline!str!} & 必填 &
SDK 对象、配置项、服务或动作的名称。允许值由调用它的具体接口定义。 对应
C++ 参数 \passthrough{\lstinline!name: const std::string \&!}。
底层为字符串；长度、编码和允许值由具体协议决定。 \\
\bottomrule
\end{longtable}

\textbf{返回值}

\begin{longtable}[]{@{}
  >{\raggedright\arraybackslash}p{(\columnwidth - 4\tabcolsep) * \real{0.18}}
  >{\raggedright\arraybackslash}p{(\columnwidth - 4\tabcolsep) * \real{0.20}}
  >{\raggedright\arraybackslash}p{(\columnwidth - 4\tabcolsep) * \real{0.62}}@{}}
\toprule
位置 & 类型 & 含义 \\
\midrule
\endhead
无 & \passthrough{\lstinline!None!} &
\passthrough{\lstinline!\_\_init\_\_!}
本身不返回值；调用类对象时，在构造函数可用的前提下得到该类实例。 \\
\bottomrule
\end{longtable}

\textbf{对应 C++}

\begin{itemize}
\tightlist
\item
  类：\passthrough{\lstinline!unitree::robot::ClientBase!}
\item
  签名：\passthrough{\lstinline!ClientBase(const std::string \&)!}
\item
  绑定策略：\passthrough{\lstinline!未单独标注!}
\item
  声明位置：\passthrough{\lstinline!include/unitree/robot/client/client\_base.hpp:23!}
\end{itemize}

\textbf{说明}

\begin{itemize}
\tightlist
\item
  示例是局部调用片段；\passthrough{\lstinline!obj!}
  和参数变量表示已经按上层业务校验并准备好的对象。
\item
  该条目可以被编辑器和类型检查器识别，但当前运行时可能在属性查找阶段直接失败。
\end{itemize}

\textbf{用法}

\begin{lstlisting}[language=Python]
from typing import TYPE_CHECKING

if TYPE_CHECKING:
    def planned_construct(name: str) -> ClientBase:
        return ClientBase(name=name)
\end{lstlisting}

\begin{quote}
{[}!CAUTION{]} 上面的 \passthrough{\lstinline!TYPE\_CHECKING!}
示例只表达计划调用形态。该分支在普通 Python
运行时为假，不代表当前扩展能执行此接口。
\end{quote}

\hypertarget{unitree-sdk2-cpp-robot-clientbase-init-1}{%
\paragraph{\texorpdfstring{\texttt{unitree\_sdk2\_cpp.robot.ClientBase.init}}{unitree\_sdk2\_cpp.robot.ClientBase.init}}\label{unitree-sdk2-cpp-robot-clientbase-init-1}}

初始化当前 SDK 对象所需的底层通道或服务资源。

\textbf{签名}

\begin{lstlisting}[language=Python]
def init(self) -> None
\end{lstlisting}

\textbf{可用性与安全性}

\textbf{\passthrough{\lstinline!SIGNATURE\_ONLY!} /
\passthrough{\lstinline!INITIALIZATION!}}：当前只有设计期签名，不能假设已存在于运行时扩展。

\textbf{参数}

无显式参数。实例方法中的 \passthrough{\lstinline!self!} 由 Python
自动传入。

\textbf{返回值}

\begin{longtable}[]{@{}
  >{\raggedright\arraybackslash}p{(\columnwidth - 4\tabcolsep) * \real{0.18}}
  >{\raggedright\arraybackslash}p{(\columnwidth - 4\tabcolsep) * \real{0.20}}
  >{\raggedright\arraybackslash}p{(\columnwidth - 4\tabcolsep) * \real{0.62}}@{}}
\toprule
位置 & 类型 & 含义 \\
\midrule
\endhead
无 & \passthrough{\lstinline!None!} & 该方法不返回 Python
值；失败可能表现为异常或底层状态变化。 \\
\bottomrule
\end{longtable}

\textbf{对应 C++}

\begin{itemize}
\tightlist
\item
  类：\passthrough{\lstinline!unitree::robot::ClientBase!}
\item
  签名：\passthrough{\lstinline!Init()!}
\item
  绑定策略：\passthrough{\lstinline!DIRECT!}
\item
  声明位置：\passthrough{\lstinline!include/unitree/robot/client/client\_base.hpp:26!}
\end{itemize}

\textbf{说明}

\begin{itemize}
\tightlist
\item
  示例是局部调用片段；\passthrough{\lstinline!obj!}
  和参数变量表示已经按上层业务校验并准备好的对象。
\item
  该条目可以被编辑器和类型检查器识别，但当前运行时可能在属性查找阶段直接失败。
\end{itemize}

\textbf{用法}

\begin{lstlisting}[language=Python]
from typing import TYPE_CHECKING

if TYPE_CHECKING:
    def planned_init(obj: ClientBase) -> None:
        obj.init()
\end{lstlisting}

\begin{quote}
{[}!CAUTION{]} 上面的 \passthrough{\lstinline!TYPE\_CHECKING!}
示例只表达计划调用形态。该分支在普通 Python
运行时为假，不代表当前扩展能执行此接口。
\end{quote}

\hypertarget{unitree-sdk2-cpp-robot-clientbase-set-timeout-microseconds-1}{%
\paragraph{\texorpdfstring{\texttt{unitree\_sdk2\_cpp.robot.ClientBase.set\_timeout\_microseconds}}{unitree\_sdk2\_cpp.robot.ClientBase.set\_timeout\_microseconds}}\label{unitree-sdk2-cpp-robot-clientbase-set-timeout-microseconds-1}}

设置 超时
\passthrough{\lstinline!microseconds!}。具体副作用和安全边界见下方状态。

\textbf{签名}

\begin{lstlisting}[language=Python]
def set_timeout_microseconds(self, microseconds: int) -> None
\end{lstlisting}

\textbf{可用性与安全性}

\textbf{\passthrough{\lstinline!AVAILABLE!} /
\passthrough{\lstinline!INITIALIZATION!}}：当前绑定源码已实现；实际执行条件仍取决于平台和对应
SDK 环境。

\textbf{参数}

\begin{longtable}[]{@{}
  >{\raggedright\arraybackslash}p{(\columnwidth - 6\tabcolsep) * \real{0.18}}
  >{\raggedright\arraybackslash}p{(\columnwidth - 6\tabcolsep) * \real{0.24}}
  >{\raggedright\arraybackslash}p{(\columnwidth - 6\tabcolsep) * \real{0.18}}
  >{\raggedright\arraybackslash}p{(\columnwidth - 6\tabcolsep) * \real{0.40}}@{}}
\toprule
名称 & Python 类型 & 默认值 & 含义 \\
\midrule
\endhead
\passthrough{\lstinline!microseconds!} & \passthrough{\lstinline!int!} &
必填 & 客户端请求超时，单位为微秒。 对应 C++ 参数
\passthrough{\lstinline!timeout: int64\_t!}。 取值范围为
-9223372036854775808 到 9223372036854775807。 \\
\bottomrule
\end{longtable}

\textbf{返回值}

\begin{longtable}[]{@{}
  >{\raggedright\arraybackslash}p{(\columnwidth - 4\tabcolsep) * \real{0.18}}
  >{\raggedright\arraybackslash}p{(\columnwidth - 4\tabcolsep) * \real{0.20}}
  >{\raggedright\arraybackslash}p{(\columnwidth - 4\tabcolsep) * \real{0.62}}@{}}
\toprule
位置 & 类型 & 含义 \\
\midrule
\endhead
无 & \passthrough{\lstinline!None!} & 该方法不返回 Python
值；失败可能表现为异常或底层状态变化。 \\
\bottomrule
\end{longtable}

\textbf{对应 C++}

\begin{itemize}
\tightlist
\item
  类：\passthrough{\lstinline!unitree::robot::ClientBase!}
\item
  签名：\passthrough{\lstinline!SetTimeout(int64\_t)!}
\item
  绑定策略：\passthrough{\lstinline!DIRECT!}
\item
  声明位置：\passthrough{\lstinline!include/unitree/robot/client/client\_base.hpp:28!}
\end{itemize}

\textbf{说明}

\begin{itemize}
\tightlist
\item
  示例是局部调用片段；\passthrough{\lstinline!obj!}
  和参数变量表示已经按上层业务校验并准备好的对象。
\end{itemize}

\textbf{用法}

\begin{lstlisting}[language=Python]
obj.set_timeout_microseconds(microseconds=microseconds)
\end{lstlisting}

\hypertarget{unitree-sdk2-cpp-robot-clientbase-set-timeout-1}{%
\paragraph{\texorpdfstring{\texttt{unitree\_sdk2\_cpp.robot.ClientBase.set\_timeout}}{unitree\_sdk2\_cpp.robot.ClientBase.set\_timeout}}\label{unitree-sdk2-cpp-robot-clientbase-set-timeout-1}}

设置 超时。具体副作用和安全边界见下方状态。

\textbf{签名}

\begin{lstlisting}[language=Python]
def set_timeout(self, seconds: float) -> None
\end{lstlisting}

\textbf{可用性与安全性}

\textbf{\passthrough{\lstinline!AVAILABLE!} /
\passthrough{\lstinline!INITIALIZATION!}}：当前绑定源码已实现；实际执行条件仍取决于平台和对应
SDK 环境。

\textbf{参数}

\begin{longtable}[]{@{}
  >{\raggedright\arraybackslash}p{(\columnwidth - 6\tabcolsep) * \real{0.18}}
  >{\raggedright\arraybackslash}p{(\columnwidth - 6\tabcolsep) * \real{0.24}}
  >{\raggedright\arraybackslash}p{(\columnwidth - 6\tabcolsep) * \real{0.18}}
  >{\raggedright\arraybackslash}p{(\columnwidth - 6\tabcolsep) * \real{0.40}}@{}}
\toprule
名称 & Python 类型 & 默认值 & 含义 \\
\midrule
\endhead
\passthrough{\lstinline!seconds!} & \passthrough{\lstinline!float!} &
必填 & 客户端请求超时，单位为秒。 对应 C++ 参数
\passthrough{\lstinline!timeout: float!}。
底层为浮点数；类型本身不说明单位、坐标系或业务安全范围。 \\
\bottomrule
\end{longtable}

\textbf{返回值}

\begin{longtable}[]{@{}
  >{\raggedright\arraybackslash}p{(\columnwidth - 4\tabcolsep) * \real{0.18}}
  >{\raggedright\arraybackslash}p{(\columnwidth - 4\tabcolsep) * \real{0.20}}
  >{\raggedright\arraybackslash}p{(\columnwidth - 4\tabcolsep) * \real{0.62}}@{}}
\toprule
位置 & 类型 & 含义 \\
\midrule
\endhead
无 & \passthrough{\lstinline!None!} & 该方法不返回 Python
值；失败可能表现为异常或底层状态变化。 \\
\bottomrule
\end{longtable}

\textbf{对应 C++}

\begin{itemize}
\tightlist
\item
  类：\passthrough{\lstinline!unitree::robot::ClientBase!}
\item
  签名：\passthrough{\lstinline!SetTimeout(float)!}
\item
  绑定策略：\passthrough{\lstinline!DIRECT!}
\item
  声明位置：\passthrough{\lstinline!include/unitree/robot/client/client\_base.hpp:29!}
\end{itemize}

\textbf{说明}

\begin{itemize}
\tightlist
\item
  示例是局部调用片段；\passthrough{\lstinline!obj!}
  和参数变量表示已经按上层业务校验并准备好的对象。
\end{itemize}

\textbf{用法}

\begin{lstlisting}[language=Python]
obj.set_timeout(seconds=seconds)
\end{lstlisting}

\hypertarget{unitree-sdk2-cpp-robot-clientchannelnamer}{%
\subsubsection{\texorpdfstring{\texttt{unitree\_sdk2\_cpp.robot.ClientChannelNamer}}{unitree\_sdk2\_cpp.robot.ClientChannelNamer}}\label{unitree-sdk2-cpp-robot-clientchannelnamer}}

SDK 类型签名预览。请逐项查看构造函数和方法的可用性。

\textbf{导入}

\begin{lstlisting}[language=Python]
from unitree_sdk2_cpp.robot import ClientChannelNamer
\end{lstlisting}

\textbf{基类}：\passthrough{\lstinline!ChannelNamer!}

\textbf{方法索引}

\begin{longtable}[]{@{}
  >{\raggedright\arraybackslash}p{(\columnwidth - 6\tabcolsep) * \real{0.18}}
  >{\raggedright\arraybackslash}p{(\columnwidth - 6\tabcolsep) * \real{0.24}}
  >{\raggedright\arraybackslash}p{(\columnwidth - 6\tabcolsep) * \real{0.18}}
  >{\raggedright\arraybackslash}p{(\columnwidth - 6\tabcolsep) * \real{0.40}}@{}}
\toprule
方法 & Python 签名 & 状态 & 安全分类 \\
\midrule
\endhead
\protect\hyperlink{unitree-sdk2-cpp-robot-clientchannelnamer-dunder-init-1}{\passthrough{\lstinline!\_\_init\_\_!}}
& \passthrough{\lstinline!def \_\_init\_\_(self) -> None!} &
\passthrough{\lstinline!SIGNATURE\_ONLY!} &
\passthrough{\lstinline!CONSTRUCTION!} \\
\bottomrule
\end{longtable}

\hypertarget{unitree-sdk2-cpp-robot-clientchannelnamer-dunder-init-1}{%
\paragraph{\texorpdfstring{\texttt{unitree\_sdk2\_cpp.robot.ClientChannelNamer.\_\_init\_\_}}{unitree\_sdk2\_cpp.robot.ClientChannelNamer.\_\_init\_\_}}\label{unitree-sdk2-cpp-robot-clientchannelnamer-dunder-init-1}}

初始化 \passthrough{\lstinline!ClientChannelNamer!}
实例。是否能实际构造取决于下方可用性状态。

\textbf{签名}

\begin{lstlisting}[language=Python]
def __init__(self) -> None
\end{lstlisting}

\textbf{可用性与安全性}

\textbf{\passthrough{\lstinline!SIGNATURE\_ONLY!} /
\passthrough{\lstinline!CONSTRUCTION!}}：当前只有设计期签名，不能假设已存在于运行时扩展。

\textbf{参数}

无显式参数。实例方法中的 \passthrough{\lstinline!self!} 由 Python
自动传入。

\textbf{返回值}

\begin{longtable}[]{@{}
  >{\raggedright\arraybackslash}p{(\columnwidth - 4\tabcolsep) * \real{0.18}}
  >{\raggedright\arraybackslash}p{(\columnwidth - 4\tabcolsep) * \real{0.20}}
  >{\raggedright\arraybackslash}p{(\columnwidth - 4\tabcolsep) * \real{0.62}}@{}}
\toprule
位置 & 类型 & 含义 \\
\midrule
\endhead
无 & \passthrough{\lstinline!None!} &
\passthrough{\lstinline!\_\_init\_\_!}
本身不返回值；调用类对象时，在构造函数可用的前提下得到该类实例。 \\
\bottomrule
\end{longtable}

\textbf{对应 C++}

\begin{itemize}
\tightlist
\item
  类：\passthrough{\lstinline!unitree::robot::ClientChannelNamer!}
\item
  签名：\passthrough{\lstinline!ClientChannelNamer()!}
\item
  绑定策略：\passthrough{\lstinline!未单独标注!}
\item
  声明位置：\passthrough{\lstinline!include/unitree/robot/channel/channel\_namer.hpp:34!}
\end{itemize}

\textbf{说明}

\begin{itemize}
\tightlist
\item
  示例是局部调用片段；\passthrough{\lstinline!obj!}
  和参数变量表示已经按上层业务校验并准备好的对象。
\item
  该条目可以被编辑器和类型检查器识别，但当前运行时可能在属性查找阶段直接失败。
\end{itemize}

\textbf{用法}

\begin{lstlisting}[language=Python]
from typing import TYPE_CHECKING

if TYPE_CHECKING:
    def planned_construct() -> ClientChannelNamer:
        return ClientChannelNamer()
\end{lstlisting}

\begin{quote}
{[}!CAUTION{]} 上面的 \passthrough{\lstinline!TYPE\_CHECKING!}
示例只表达计划调用形态。该分支在普通 Python
运行时为假，不代表当前扩展能执行此接口。
\end{quote}

\hypertarget{unitree-sdk2-cpp-robot-clientstub}{%
\subsubsection{\texorpdfstring{\texttt{unitree\_sdk2\_cpp.robot.ClientStub}}{unitree\_sdk2\_cpp.robot.ClientStub}}\label{unitree-sdk2-cpp-robot-clientstub}}

SDK 类型签名预览。请逐项查看构造函数和方法的可用性。

\textbf{导入}

\begin{lstlisting}[language=Python]
from unitree_sdk2_cpp.robot import ClientStub
\end{lstlisting}

\textbf{基类}：\passthrough{\lstinline!object!}

\textbf{方法索引}

\begin{longtable}[]{@{}
  >{\raggedright\arraybackslash}p{(\columnwidth - 6\tabcolsep) * \real{0.18}}
  >{\raggedright\arraybackslash}p{(\columnwidth - 6\tabcolsep) * \real{0.24}}
  >{\raggedright\arraybackslash}p{(\columnwidth - 6\tabcolsep) * \real{0.18}}
  >{\raggedright\arraybackslash}p{(\columnwidth - 6\tabcolsep) * \real{0.40}}@{}}
\toprule
方法 & Python 签名 & 状态 & 安全分类 \\
\midrule
\endhead
\protect\hyperlink{unitree-sdk2-cpp-robot-clientstub-dunder-init-1}{\passthrough{\lstinline!\_\_init\_\_!}}
& \passthrough{\lstinline!def \_\_init\_\_(self) -> None!} &
\passthrough{\lstinline!SIGNATURE\_ONLY!} &
\passthrough{\lstinline!CONSTRUCTION!} \\
\protect\hyperlink{unitree-sdk2-cpp-robot-clientstub-init-1}{\passthrough{\lstinline!init!}}
& \passthrough{\lstinline!def init(self, name: str) -> None!} &
\passthrough{\lstinline!SIGNATURE\_ONLY!} &
\passthrough{\lstinline!INITIALIZATION!} \\
\protect\hyperlink{unitree-sdk2-cpp-robot-clientstub-send-1}{\passthrough{\lstinline!send!}}
&
\passthrough{\lstinline!def send(self, req: Any, wait\_timeout: int) -> bool!}
& \passthrough{\lstinline!SIGNATURE\_ONLY!} &
\passthrough{\lstinline!HARDWARE\_SIDE\_EFFECT!} \\
\protect\hyperlink{unitree-sdk2-cpp-robot-clientstub-send-request-1}{\passthrough{\lstinline!send\_request!}}
&
\passthrough{\lstinline!def send\_request(self, req: Any, wait\_timeout: int) -> RequestFuture!}
& \passthrough{\lstinline!SIGNATURE\_ONLY!} &
\passthrough{\lstinline!HARDWARE\_SIDE\_EFFECT!} \\
\bottomrule
\end{longtable}

\hypertarget{unitree-sdk2-cpp-robot-clientstub-dunder-init-1}{%
\paragraph{\texorpdfstring{\texttt{unitree\_sdk2\_cpp.robot.ClientStub.\_\_init\_\_}}{unitree\_sdk2\_cpp.robot.ClientStub.\_\_init\_\_}}\label{unitree-sdk2-cpp-robot-clientstub-dunder-init-1}}

初始化 \passthrough{\lstinline!ClientStub!}
实例。是否能实际构造取决于下方可用性状态。

\textbf{签名}

\begin{lstlisting}[language=Python]
def __init__(self) -> None
\end{lstlisting}

\textbf{可用性与安全性}

\textbf{\passthrough{\lstinline!SIGNATURE\_ONLY!} /
\passthrough{\lstinline!CONSTRUCTION!}}：当前只有设计期签名，不能假设已存在于运行时扩展。

\textbf{参数}

无显式参数。实例方法中的 \passthrough{\lstinline!self!} 由 Python
自动传入。

\textbf{返回值}

\begin{longtable}[]{@{}
  >{\raggedright\arraybackslash}p{(\columnwidth - 4\tabcolsep) * \real{0.18}}
  >{\raggedright\arraybackslash}p{(\columnwidth - 4\tabcolsep) * \real{0.20}}
  >{\raggedright\arraybackslash}p{(\columnwidth - 4\tabcolsep) * \real{0.62}}@{}}
\toprule
位置 & 类型 & 含义 \\
\midrule
\endhead
无 & \passthrough{\lstinline!None!} &
\passthrough{\lstinline!\_\_init\_\_!}
本身不返回值；调用类对象时，在构造函数可用的前提下得到该类实例。 \\
\bottomrule
\end{longtable}

\textbf{对应 C++}

\begin{itemize}
\tightlist
\item
  类：\passthrough{\lstinline!unitree::robot::ClientStub!}
\item
  签名：\passthrough{\lstinline!ClientStub()!}
\item
  绑定策略：\passthrough{\lstinline!未单独标注!}
\item
  声明位置：\passthrough{\lstinline!include/unitree/robot/client/client\_stub.hpp:15!}
\end{itemize}

\textbf{说明}

\begin{itemize}
\tightlist
\item
  示例是局部调用片段；\passthrough{\lstinline!obj!}
  和参数变量表示已经按上层业务校验并准备好的对象。
\item
  该条目可以被编辑器和类型检查器识别，但当前运行时可能在属性查找阶段直接失败。
\end{itemize}

\textbf{用法}

\begin{lstlisting}[language=Python]
from typing import TYPE_CHECKING

if TYPE_CHECKING:
    def planned_construct() -> ClientStub:
        return ClientStub()
\end{lstlisting}

\begin{quote}
{[}!CAUTION{]} 上面的 \passthrough{\lstinline!TYPE\_CHECKING!}
示例只表达计划调用形态。该分支在普通 Python
运行时为假，不代表当前扩展能执行此接口。
\end{quote}

\hypertarget{unitree-sdk2-cpp-robot-clientstub-init-1}{%
\paragraph{\texorpdfstring{\texttt{unitree\_sdk2\_cpp.robot.ClientStub.init}}{unitree\_sdk2\_cpp.robot.ClientStub.init}}\label{unitree-sdk2-cpp-robot-clientstub-init-1}}

初始化当前 SDK 对象所需的底层通道或服务资源。

\textbf{签名}

\begin{lstlisting}[language=Python]
def init(self, name: str) -> None
\end{lstlisting}

\textbf{可用性与安全性}

\textbf{\passthrough{\lstinline!SIGNATURE\_ONLY!} /
\passthrough{\lstinline!INITIALIZATION!}}：当前只有设计期签名，不能假设已存在于运行时扩展。

\textbf{参数}

\begin{longtable}[]{@{}
  >{\raggedright\arraybackslash}p{(\columnwidth - 6\tabcolsep) * \real{0.18}}
  >{\raggedright\arraybackslash}p{(\columnwidth - 6\tabcolsep) * \real{0.24}}
  >{\raggedright\arraybackslash}p{(\columnwidth - 6\tabcolsep) * \real{0.18}}
  >{\raggedright\arraybackslash}p{(\columnwidth - 6\tabcolsep) * \real{0.40}}@{}}
\toprule
名称 & Python 类型 & 默认值 & 含义 \\
\midrule
\endhead
\passthrough{\lstinline!name!} & \passthrough{\lstinline!str!} & 必填 &
SDK 对象、配置项、服务或动作的名称。允许值由调用它的具体接口定义。 对应
C++ 参数 \passthrough{\lstinline!name: const std::string \&!}。
底层为字符串；长度、编码和允许值由具体协议决定。 \\
\bottomrule
\end{longtable}

\textbf{返回值}

\begin{longtable}[]{@{}
  >{\raggedright\arraybackslash}p{(\columnwidth - 4\tabcolsep) * \real{0.18}}
  >{\raggedright\arraybackslash}p{(\columnwidth - 4\tabcolsep) * \real{0.20}}
  >{\raggedright\arraybackslash}p{(\columnwidth - 4\tabcolsep) * \real{0.62}}@{}}
\toprule
位置 & 类型 & 含义 \\
\midrule
\endhead
无 & \passthrough{\lstinline!None!} & 该方法不返回 Python
值；失败可能表现为异常或底层状态变化。 \\
\bottomrule
\end{longtable}

\textbf{对应 C++}

\begin{itemize}
\tightlist
\item
  类：\passthrough{\lstinline!unitree::robot::ClientStub!}
\item
  签名：\passthrough{\lstinline!Init(const std::string \&)!}
\item
  绑定策略：\passthrough{\lstinline!DIRECT!}
\item
  声明位置：\passthrough{\lstinline!include/unitree/robot/client/client\_stub.hpp:18!}
\end{itemize}

\textbf{说明}

\begin{itemize}
\tightlist
\item
  示例是局部调用片段；\passthrough{\lstinline!obj!}
  和参数变量表示已经按上层业务校验并准备好的对象。
\item
  该条目可以被编辑器和类型检查器识别，但当前运行时可能在属性查找阶段直接失败。
\end{itemize}

\textbf{用法}

\begin{lstlisting}[language=Python]
from typing import TYPE_CHECKING

if TYPE_CHECKING:
    def planned_init(obj: ClientStub, name: str) -> None:
        obj.init(name=name)
\end{lstlisting}

\begin{quote}
{[}!CAUTION{]} 上面的 \passthrough{\lstinline!TYPE\_CHECKING!}
示例只表达计划调用形态。该分支在普通 Python
运行时为假，不代表当前扩展能执行此接口。
\end{quote}

\hypertarget{unitree-sdk2-cpp-robot-clientstub-send-1}{%
\paragraph{\texorpdfstring{\texttt{unitree\_sdk2\_cpp.robot.ClientStub.send}}{unitree\_sdk2\_cpp.robot.ClientStub.send}}\label{unitree-sdk2-cpp-robot-clientstub-send-1}}

对应 C++ SDK 操作
\passthrough{\lstinline!Send(const unitree::robot::Request \&, int64\_t)!}。上游头文件没有可直接生成的业务说明时，本参考只保证签名映射，精确语义需查目标型号协议。

\textbf{签名}

\begin{lstlisting}[language=Python]
def send(self, req: Any, wait_timeout: int) -> bool
\end{lstlisting}

\textbf{可用性与安全性}

\textbf{\passthrough{\lstinline!SIGNATURE\_ONLY!} /
\passthrough{\lstinline!HARDWARE\_SIDE\_EFFECT!}}：仅用于补全和静态检查。当前二进制没有该
Python 方法，未来实现还需单独评审硬件副作用。

\textbf{参数}

\begin{longtable}[]{@{}
  >{\raggedright\arraybackslash}p{(\columnwidth - 6\tabcolsep) * \real{0.18}}
  >{\raggedright\arraybackslash}p{(\columnwidth - 6\tabcolsep) * \real{0.24}}
  >{\raggedright\arraybackslash}p{(\columnwidth - 6\tabcolsep) * \real{0.18}}
  >{\raggedright\arraybackslash}p{(\columnwidth - 6\tabcolsep) * \real{0.40}}@{}}
\toprule
名称 & Python 类型 & 默认值 & 含义 \\
\midrule
\endhead
\passthrough{\lstinline!req!} & \passthrough{\lstinline!Any!} & 必填 &
传给该接口的 \passthrough{\lstinline!req!} 参数，Python 类型为
\passthrough{\lstinline!Any!}。精确含义、单位、范围和枚举值需查目标型号对应头文件与协议。
对应 C++ 参数
\passthrough{\lstinline!req: const unitree::robot::Request \&!}。 \\
\passthrough{\lstinline!wait\_timeout!} & \passthrough{\lstinline!int!}
& 必填 & 请求等待超时参数；精确单位沿用对应 C++ 接口定义。 对应 C++ 参数
\passthrough{\lstinline!waitTimeout: int64\_t!}。 取值范围为
-9223372036854775808 到 9223372036854775807。 \\
\bottomrule
\end{longtable}

\textbf{返回值}

\begin{longtable}[]{@{}
  >{\raggedright\arraybackslash}p{(\columnwidth - 4\tabcolsep) * \real{0.18}}
  >{\raggedright\arraybackslash}p{(\columnwidth - 4\tabcolsep) * \real{0.20}}
  >{\raggedright\arraybackslash}p{(\columnwidth - 4\tabcolsep) * \real{0.62}}@{}}
\toprule
位置 & 类型 & 含义 \\
\midrule
\endhead
返回值 & \passthrough{\lstinline!bool!} & 返回
\passthrough{\lstinline!bool!}。更精确的业务含义见该方法用途、C++
签名和目标型号协议。 \\
\bottomrule
\end{longtable}

\textbf{对应 C++}

\begin{itemize}
\tightlist
\item
  类：\passthrough{\lstinline!unitree::robot::ClientStub!}
\item
  签名：\passthrough{\lstinline!Send(const unitree::robot::Request \&, int64\_t)!}
\item
  绑定策略：\passthrough{\lstinline!DIRECT!}
\item
  声明位置：\passthrough{\lstinline!include/unitree/robot/client/client\_stub.hpp:20!}
\end{itemize}

\textbf{说明}

\begin{itemize}
\tightlist
\item
  示例是局部调用片段；\passthrough{\lstinline!obj!}
  和参数变量表示已经按上层业务校验并准备好的对象。
\item
  该条目可以被编辑器和类型检查器识别，但当前运行时可能在属性查找阶段直接失败。
\end{itemize}

\textbf{用法}

\begin{lstlisting}[language=Python]
from typing import TYPE_CHECKING

if TYPE_CHECKING:
    def planned_send(obj: ClientStub, req: Any, wait_timeout: int) -> bool:
        return obj.send(req=req, wait_timeout=wait_timeout)
\end{lstlisting}

\begin{quote}
{[}!CAUTION{]} 上面的 \passthrough{\lstinline!TYPE\_CHECKING!}
示例只表达计划调用形态。该分支在普通 Python
运行时为假，不代表当前扩展能执行此接口。
\end{quote}

\hypertarget{unitree-sdk2-cpp-robot-clientstub-send-request-1}{%
\paragraph{\texorpdfstring{\texttt{unitree\_sdk2\_cpp.robot.ClientStub.send\_request}}{unitree\_sdk2\_cpp.robot.ClientStub.send\_request}}\label{unitree-sdk2-cpp-robot-clientstub-send-request-1}}

对应 C++ SDK 操作
\passthrough{\lstinline!SendRequest(const unitree::robot::Request \&, int64\_t)!}。上游头文件没有可直接生成的业务说明时，本参考只保证签名映射，精确语义需查目标型号协议。

\textbf{签名}

\begin{lstlisting}[language=Python]
def send_request(self, req: Any, wait_timeout: int) -> RequestFuture
\end{lstlisting}

\textbf{可用性与安全性}

\textbf{\passthrough{\lstinline!SIGNATURE\_ONLY!} /
\passthrough{\lstinline!HARDWARE\_SIDE\_EFFECT!}}：仅用于补全和静态检查。当前二进制没有该
Python 方法，未来实现还需单独评审硬件副作用。

\textbf{参数}

\begin{longtable}[]{@{}
  >{\raggedright\arraybackslash}p{(\columnwidth - 6\tabcolsep) * \real{0.18}}
  >{\raggedright\arraybackslash}p{(\columnwidth - 6\tabcolsep) * \real{0.24}}
  >{\raggedright\arraybackslash}p{(\columnwidth - 6\tabcolsep) * \real{0.18}}
  >{\raggedright\arraybackslash}p{(\columnwidth - 6\tabcolsep) * \real{0.40}}@{}}
\toprule
名称 & Python 类型 & 默认值 & 含义 \\
\midrule
\endhead
\passthrough{\lstinline!req!} & \passthrough{\lstinline!Any!} & 必填 &
传给该接口的 \passthrough{\lstinline!req!} 参数，Python 类型为
\passthrough{\lstinline!Any!}。精确含义、单位、范围和枚举值需查目标型号对应头文件与协议。
对应 C++ 参数
\passthrough{\lstinline!req: const unitree::robot::Request \&!}。 \\
\passthrough{\lstinline!wait\_timeout!} & \passthrough{\lstinline!int!}
& 必填 & 请求等待超时参数；精确单位沿用对应 C++ 接口定义。 对应 C++ 参数
\passthrough{\lstinline!waitTimeout: int64\_t!}。 取值范围为
-9223372036854775808 到 9223372036854775807。 \\
\bottomrule
\end{longtable}

\textbf{返回值}

\begin{longtable}[]{@{}
  >{\raggedright\arraybackslash}p{(\columnwidth - 4\tabcolsep) * \real{0.18}}
  >{\raggedright\arraybackslash}p{(\columnwidth - 4\tabcolsep) * \real{0.20}}
  >{\raggedright\arraybackslash}p{(\columnwidth - 4\tabcolsep) * \real{0.62}}@{}}
\toprule
位置 & 类型 & 含义 \\
\midrule
\endhead
返回值 & \passthrough{\lstinline!RequestFuture!} & 返回
\passthrough{\lstinline!RequestFuture!}。更精确的业务含义见该方法用途、C++
签名和目标型号协议。 \\
\bottomrule
\end{longtable}

\textbf{对应 C++}

\begin{itemize}
\tightlist
\item
  类：\passthrough{\lstinline!unitree::robot::ClientStub!}
\item
  签名：\passthrough{\lstinline!SendRequest(const unitree::robot::Request \&, int64\_t)!}
\item
  绑定策略：\passthrough{\lstinline!DIRECT!}
\item
  声明位置：\passthrough{\lstinline!include/unitree/robot/client/client\_stub.hpp:21!}
\end{itemize}

\textbf{说明}

\begin{itemize}
\tightlist
\item
  示例是局部调用片段；\passthrough{\lstinline!obj!}
  和参数变量表示已经按上层业务校验并准备好的对象。
\item
  该条目可以被编辑器和类型检查器识别，但当前运行时可能在属性查找阶段直接失败。
\end{itemize}

\textbf{用法}

\begin{lstlisting}[language=Python]
from typing import TYPE_CHECKING

if TYPE_CHECKING:
    def planned_send_request(obj: ClientStub, req: Any, wait_timeout: int) -> RequestFuture:
        return obj.send_request(req=req, wait_timeout=wait_timeout)
\end{lstlisting}

\begin{quote}
{[}!CAUTION{]} 上面的 \passthrough{\lstinline!TYPE\_CHECKING!}
示例只表达计划调用形态。该分支在普通 Python
运行时为假，不代表当前扩展能执行此接口。
\end{quote}

\hypertarget{unitree-sdk2-cpp-robot-leasecache}{%
\subsubsection{\texorpdfstring{\texttt{unitree\_sdk2\_cpp.robot.LeaseCache}}{unitree\_sdk2\_cpp.robot.LeaseCache}}\label{unitree-sdk2-cpp-robot-leasecache}}

SDK 类型签名预览。请逐项查看构造函数和方法的可用性。

\textbf{导入}

\begin{lstlisting}[language=Python]
from unitree_sdk2_cpp.robot import LeaseCache
\end{lstlisting}

\textbf{基类}：\passthrough{\lstinline!object!}

\textbf{方法索引}

\begin{longtable}[]{@{}
  >{\raggedright\arraybackslash}p{(\columnwidth - 6\tabcolsep) * \real{0.18}}
  >{\raggedright\arraybackslash}p{(\columnwidth - 6\tabcolsep) * \real{0.24}}
  >{\raggedright\arraybackslash}p{(\columnwidth - 6\tabcolsep) * \real{0.18}}
  >{\raggedright\arraybackslash}p{(\columnwidth - 6\tabcolsep) * \real{0.40}}@{}}
\toprule
方法 & Python 签名 & 状态 & 安全分类 \\
\midrule
\endhead
\protect\hyperlink{unitree-sdk2-cpp-robot-leasecache-dunder-init-1}{\passthrough{\lstinline!\_\_init\_\_!}}
& \passthrough{\lstinline!def \_\_init\_\_(self) -> None!} &
\passthrough{\lstinline!SIGNATURE\_ONLY!} &
\passthrough{\lstinline!CONSTRUCTION!} \\
\protect\hyperlink{unitree-sdk2-cpp-robot-leasecache-set-1}{\passthrough{\lstinline!set!}}
&
\passthrough{\lstinline!def set(self, id: int, m\_name: str, last\_modified: int = ...) -> None!}
& \passthrough{\lstinline!SIGNATURE\_ONLY!} &
\passthrough{\lstinline!UNCLASSIFIED!} \\
\protect\hyperlink{unitree-sdk2-cpp-robot-leasecache-renewal-1}{\passthrough{\lstinline!renewal!}}
&
\passthrough{\lstinline!def renewal(self, last\_modified: int = ...) -> None!}
& \passthrough{\lstinline!SIGNATURE\_ONLY!} &
\passthrough{\lstinline!UNCLASSIFIED!} \\
\protect\hyperlink{unitree-sdk2-cpp-robot-leasecache-clear-1}{\passthrough{\lstinline!clear!}}
& \passthrough{\lstinline!def clear(self) -> None!} &
\passthrough{\lstinline!SIGNATURE\_ONLY!} &
\passthrough{\lstinline!UNCLASSIFIED!} \\
\protect\hyperlink{unitree-sdk2-cpp-robot-leasecache-get-last-modified-1}{\passthrough{\lstinline!get\_last\_modified!}}
& \passthrough{\lstinline!def get\_last\_modified(self) -> int!} &
\passthrough{\lstinline!SIGNATURE\_ONLY!} &
\passthrough{\lstinline!UNCLASSIFIED!} \\
\protect\hyperlink{unitree-sdk2-cpp-robot-leasecache-get-id-1}{\passthrough{\lstinline!get\_id!}}
& \passthrough{\lstinline!def get\_id(self) -> int!} &
\passthrough{\lstinline!SIGNATURE\_ONLY!} &
\passthrough{\lstinline!UNCLASSIFIED!} \\
\protect\hyperlink{unitree-sdk2-cpp-robot-leasecache-get-name-1}{\passthrough{\lstinline!get\_name!}}
& \passthrough{\lstinline!def get\_name(self) -> str!} &
\passthrough{\lstinline!SIGNATURE\_ONLY!} &
\passthrough{\lstinline!UNCLASSIFIED!} \\
\bottomrule
\end{longtable}

\hypertarget{unitree-sdk2-cpp-robot-leasecache-dunder-init-1}{%
\paragraph{\texorpdfstring{\texttt{unitree\_sdk2\_cpp.robot.LeaseCache.\_\_init\_\_}}{unitree\_sdk2\_cpp.robot.LeaseCache.\_\_init\_\_}}\label{unitree-sdk2-cpp-robot-leasecache-dunder-init-1}}

初始化 \passthrough{\lstinline!LeaseCache!}
实例。是否能实际构造取决于下方可用性状态。

\textbf{签名}

\begin{lstlisting}[language=Python]
def __init__(self) -> None
\end{lstlisting}

\textbf{可用性与安全性}

\textbf{\passthrough{\lstinline!SIGNATURE\_ONLY!} /
\passthrough{\lstinline!CONSTRUCTION!}}：当前只有设计期签名，不能假设已存在于运行时扩展。

\textbf{参数}

无显式参数。实例方法中的 \passthrough{\lstinline!self!} 由 Python
自动传入。

\textbf{返回值}

\begin{longtable}[]{@{}
  >{\raggedright\arraybackslash}p{(\columnwidth - 4\tabcolsep) * \real{0.18}}
  >{\raggedright\arraybackslash}p{(\columnwidth - 4\tabcolsep) * \real{0.20}}
  >{\raggedright\arraybackslash}p{(\columnwidth - 4\tabcolsep) * \real{0.62}}@{}}
\toprule
位置 & 类型 & 含义 \\
\midrule
\endhead
无 & \passthrough{\lstinline!None!} &
\passthrough{\lstinline!\_\_init\_\_!}
本身不返回值；调用类对象时，在构造函数可用的前提下得到该类实例。 \\
\bottomrule
\end{longtable}

\textbf{对应 C++}

\begin{itemize}
\tightlist
\item
  类：\passthrough{\lstinline!unitree::robot::LeaseCache!}
\item
  签名：\passthrough{\lstinline!LeaseCache()!}
\item
  绑定策略：\passthrough{\lstinline!未单独标注!}
\item
  声明位置：\passthrough{\lstinline!include/unitree/robot/server/lease\_server.hpp:13!}
\end{itemize}

\textbf{说明}

\begin{itemize}
\tightlist
\item
  示例是局部调用片段；\passthrough{\lstinline!obj!}
  和参数变量表示已经按上层业务校验并准备好的对象。
\item
  该条目可以被编辑器和类型检查器识别，但当前运行时可能在属性查找阶段直接失败。
\end{itemize}

\textbf{用法}

\begin{lstlisting}[language=Python]
from typing import TYPE_CHECKING

if TYPE_CHECKING:
    def planned_construct() -> LeaseCache:
        return LeaseCache()
\end{lstlisting}

\begin{quote}
{[}!CAUTION{]} 上面的 \passthrough{\lstinline!TYPE\_CHECKING!}
示例只表达计划调用形态。该分支在普通 Python
运行时为假，不代表当前扩展能执行此接口。
\end{quote}

\hypertarget{unitree-sdk2-cpp-robot-leasecache-set-1}{%
\paragraph{\texorpdfstring{\texttt{unitree\_sdk2\_cpp.robot.LeaseCache.set}}{unitree\_sdk2\_cpp.robot.LeaseCache.set}}\label{unitree-sdk2-cpp-robot-leasecache-set-1}}

对应 C++ SDK 操作
\passthrough{\lstinline!Set(int64\_t, const std::string \&, int64\_t)!}。上游头文件没有可直接生成的业务说明时，本参考只保证签名映射，精确语义需查目标型号协议。

\textbf{签名}

\begin{lstlisting}[language=Python]
def set(self, id: int, m_name: str, last_modified: int = ...) -> None
\end{lstlisting}

\textbf{可用性与安全性}

\textbf{\passthrough{\lstinline!SIGNATURE\_ONLY!} /
\passthrough{\lstinline!UNCLASSIFIED!}}：当前只有设计期签名，不能假设已存在于运行时扩展。

\textbf{参数}

\begin{longtable}[]{@{}
  >{\raggedright\arraybackslash}p{(\columnwidth - 6\tabcolsep) * \real{0.18}}
  >{\raggedright\arraybackslash}p{(\columnwidth - 6\tabcolsep) * \real{0.24}}
  >{\raggedright\arraybackslash}p{(\columnwidth - 6\tabcolsep) * \real{0.18}}
  >{\raggedright\arraybackslash}p{(\columnwidth - 6\tabcolsep) * \real{0.40}}@{}}
\toprule
名称 & Python 类型 & 默认值 & 含义 \\
\midrule
\endhead
\passthrough{\lstinline!id!} & \passthrough{\lstinline!int!} & 必填 &
传给该接口的 ID 参数，Python 类型为
\passthrough{\lstinline!int!}。精确含义、单位、范围和枚举值需查目标型号对应头文件与协议。
对应 C++ 参数 \passthrough{\lstinline!id: int64\_t!}。 取值范围为
-9223372036854775808 到 9223372036854775807。 \\
\passthrough{\lstinline!m\_name!} & \passthrough{\lstinline!str!} & 必填
& 传给该接口的 \passthrough{\lstinline!m!} 名称 参数，Python 类型为
\passthrough{\lstinline!str!}。精确含义、单位、范围和枚举值需查目标型号对应头文件与协议。
对应 C++ 参数 \passthrough{\lstinline!mName: const std::string \&!}。
底层为字符串；长度、编码和允许值由具体协议决定。 \\
\passthrough{\lstinline!last\_modified!} & \passthrough{\lstinline!int!}
& \passthrough{\lstinline!...!} & 传给该接口的
\passthrough{\lstinline!last!} \passthrough{\lstinline!modified!}
参数，Python 类型为
\passthrough{\lstinline!int!}。精确含义、单位、范围和枚举值需查目标型号对应头文件与协议。
对应 C++ 参数 \passthrough{\lstinline!lastModified: int64\_t!}。
取值范围为 -9223372036854775808 到 9223372036854775807。 \\
\bottomrule
\end{longtable}

\textbf{返回值}

\begin{longtable}[]{@{}
  >{\raggedright\arraybackslash}p{(\columnwidth - 4\tabcolsep) * \real{0.18}}
  >{\raggedright\arraybackslash}p{(\columnwidth - 4\tabcolsep) * \real{0.20}}
  >{\raggedright\arraybackslash}p{(\columnwidth - 4\tabcolsep) * \real{0.62}}@{}}
\toprule
位置 & 类型 & 含义 \\
\midrule
\endhead
无 & \passthrough{\lstinline!None!} & 该方法不返回 Python
值；失败可能表现为异常或底层状态变化。 \\
\bottomrule
\end{longtable}

\textbf{对应 C++}

\begin{itemize}
\tightlist
\item
  类：\passthrough{\lstinline!unitree::robot::LeaseCache!}
\item
  签名：\passthrough{\lstinline!Set(int64\_t, const std::string \&, int64\_t)!}
\item
  绑定策略：\passthrough{\lstinline!SIGNATURE\_PREVIEW!}
\item
  声明位置：\passthrough{\lstinline!include/unitree/robot/server/lease\_server.hpp:16!}
\end{itemize}

\textbf{说明}

\begin{itemize}
\tightlist
\item
  示例是局部调用片段；\passthrough{\lstinline!obj!}
  和参数变量表示已经按上层业务校验并准备好的对象。
\item
  默认值显示为 \passthrough{\lstinline!...!}，表示 C++
  声明存在默认参数，但当前 AST
  清单没有保存其字面量；省略参数可使用上游默认行为。
\item
  该条目可以被编辑器和类型检查器识别，但当前运行时可能在属性查找阶段直接失败。
\end{itemize}

\textbf{用法}

\begin{lstlisting}[language=Python]
from typing import TYPE_CHECKING

if TYPE_CHECKING:
    def planned_set(obj: LeaseCache, id: int, m_name: str, last_modified: int) -> None:
        obj.set(id=id, m_name=m_name, last_modified=last_modified)
\end{lstlisting}

\begin{quote}
{[}!CAUTION{]} 上面的 \passthrough{\lstinline!TYPE\_CHECKING!}
示例只表达计划调用形态。该分支在普通 Python
运行时为假，不代表当前扩展能执行此接口。
\end{quote}

\hypertarget{unitree-sdk2-cpp-robot-leasecache-renewal-1}{%
\paragraph{\texorpdfstring{\texttt{unitree\_sdk2\_cpp.robot.LeaseCache.renewal}}{unitree\_sdk2\_cpp.robot.LeaseCache.renewal}}\label{unitree-sdk2-cpp-robot-leasecache-renewal-1}}

对应 C++ SDK 操作
\passthrough{\lstinline!Renewal(int64\_t)!}。上游头文件没有可直接生成的业务说明时，本参考只保证签名映射，精确语义需查目标型号协议。

\textbf{签名}

\begin{lstlisting}[language=Python]
def renewal(self, last_modified: int = ...) -> None
\end{lstlisting}

\textbf{可用性与安全性}

\textbf{\passthrough{\lstinline!SIGNATURE\_ONLY!} /
\passthrough{\lstinline!UNCLASSIFIED!}}：当前只有设计期签名，不能假设已存在于运行时扩展。

\textbf{参数}

\begin{longtable}[]{@{}
  >{\raggedright\arraybackslash}p{(\columnwidth - 6\tabcolsep) * \real{0.18}}
  >{\raggedright\arraybackslash}p{(\columnwidth - 6\tabcolsep) * \real{0.24}}
  >{\raggedright\arraybackslash}p{(\columnwidth - 6\tabcolsep) * \real{0.18}}
  >{\raggedright\arraybackslash}p{(\columnwidth - 6\tabcolsep) * \real{0.40}}@{}}
\toprule
名称 & Python 类型 & 默认值 & 含义 \\
\midrule
\endhead
\passthrough{\lstinline!last\_modified!} & \passthrough{\lstinline!int!}
& \passthrough{\lstinline!...!} & 传给该接口的
\passthrough{\lstinline!last!} \passthrough{\lstinline!modified!}
参数，Python 类型为
\passthrough{\lstinline!int!}。精确含义、单位、范围和枚举值需查目标型号对应头文件与协议。
对应 C++ 参数 \passthrough{\lstinline!lastModified: int64\_t!}。
取值范围为 -9223372036854775808 到 9223372036854775807。 \\
\bottomrule
\end{longtable}

\textbf{返回值}

\begin{longtable}[]{@{}
  >{\raggedright\arraybackslash}p{(\columnwidth - 4\tabcolsep) * \real{0.18}}
  >{\raggedright\arraybackslash}p{(\columnwidth - 4\tabcolsep) * \real{0.20}}
  >{\raggedright\arraybackslash}p{(\columnwidth - 4\tabcolsep) * \real{0.62}}@{}}
\toprule
位置 & 类型 & 含义 \\
\midrule
\endhead
无 & \passthrough{\lstinline!None!} & 该方法不返回 Python
值；失败可能表现为异常或底层状态变化。 \\
\bottomrule
\end{longtable}

\textbf{对应 C++}

\begin{itemize}
\tightlist
\item
  类：\passthrough{\lstinline!unitree::robot::LeaseCache!}
\item
  签名：\passthrough{\lstinline!Renewal(int64\_t)!}
\item
  绑定策略：\passthrough{\lstinline!SIGNATURE\_PREVIEW!}
\item
  声明位置：\passthrough{\lstinline!include/unitree/robot/server/lease\_server.hpp:17!}
\end{itemize}

\textbf{说明}

\begin{itemize}
\tightlist
\item
  示例是局部调用片段；\passthrough{\lstinline!obj!}
  和参数变量表示已经按上层业务校验并准备好的对象。
\item
  默认值显示为 \passthrough{\lstinline!...!}，表示 C++
  声明存在默认参数，但当前 AST
  清单没有保存其字面量；省略参数可使用上游默认行为。
\item
  该条目可以被编辑器和类型检查器识别，但当前运行时可能在属性查找阶段直接失败。
\end{itemize}

\textbf{用法}

\begin{lstlisting}[language=Python]
from typing import TYPE_CHECKING

if TYPE_CHECKING:
    def planned_renewal(obj: LeaseCache, last_modified: int) -> None:
        obj.renewal(last_modified=last_modified)
\end{lstlisting}

\begin{quote}
{[}!CAUTION{]} 上面的 \passthrough{\lstinline!TYPE\_CHECKING!}
示例只表达计划调用形态。该分支在普通 Python
运行时为假，不代表当前扩展能执行此接口。
\end{quote}

\hypertarget{unitree-sdk2-cpp-robot-leasecache-clear-1}{%
\paragraph{\texorpdfstring{\texttt{unitree\_sdk2\_cpp.robot.LeaseCache.clear}}{unitree\_sdk2\_cpp.robot.LeaseCache.clear}}\label{unitree-sdk2-cpp-robot-leasecache-clear-1}}

对应 C++ SDK 操作
\passthrough{\lstinline!Clear()!}。上游头文件没有可直接生成的业务说明时，本参考只保证签名映射，精确语义需查目标型号协议。

\textbf{签名}

\begin{lstlisting}[language=Python]
def clear(self) -> None
\end{lstlisting}

\textbf{可用性与安全性}

\textbf{\passthrough{\lstinline!SIGNATURE\_ONLY!} /
\passthrough{\lstinline!UNCLASSIFIED!}}：当前只有设计期签名，不能假设已存在于运行时扩展。

\textbf{参数}

无显式参数。实例方法中的 \passthrough{\lstinline!self!} 由 Python
自动传入。

\textbf{返回值}

\begin{longtable}[]{@{}
  >{\raggedright\arraybackslash}p{(\columnwidth - 4\tabcolsep) * \real{0.18}}
  >{\raggedright\arraybackslash}p{(\columnwidth - 4\tabcolsep) * \real{0.20}}
  >{\raggedright\arraybackslash}p{(\columnwidth - 4\tabcolsep) * \real{0.62}}@{}}
\toprule
位置 & 类型 & 含义 \\
\midrule
\endhead
无 & \passthrough{\lstinline!None!} & 该方法不返回 Python
值；失败可能表现为异常或底层状态变化。 \\
\bottomrule
\end{longtable}

\textbf{对应 C++}

\begin{itemize}
\tightlist
\item
  类：\passthrough{\lstinline!unitree::robot::LeaseCache!}
\item
  签名：\passthrough{\lstinline!Clear()!}
\item
  绑定策略：\passthrough{\lstinline!SIGNATURE\_PREVIEW!}
\item
  声明位置：\passthrough{\lstinline!include/unitree/robot/server/lease\_server.hpp:18!}
\end{itemize}

\textbf{说明}

\begin{itemize}
\tightlist
\item
  示例是局部调用片段；\passthrough{\lstinline!obj!}
  和参数变量表示已经按上层业务校验并准备好的对象。
\item
  该条目可以被编辑器和类型检查器识别，但当前运行时可能在属性查找阶段直接失败。
\end{itemize}

\textbf{用法}

\begin{lstlisting}[language=Python]
from typing import TYPE_CHECKING

if TYPE_CHECKING:
    def planned_clear(obj: LeaseCache) -> None:
        obj.clear()
\end{lstlisting}

\begin{quote}
{[}!CAUTION{]} 上面的 \passthrough{\lstinline!TYPE\_CHECKING!}
示例只表达计划调用形态。该分支在普通 Python
运行时为假，不代表当前扩展能执行此接口。
\end{quote}

\hypertarget{unitree-sdk2-cpp-robot-leasecache-get-last-modified-1}{%
\paragraph{\texorpdfstring{\texttt{unitree\_sdk2\_cpp.robot.LeaseCache.get\_last\_modified}}{unitree\_sdk2\_cpp.robot.LeaseCache.get\_last\_modified}}\label{unitree-sdk2-cpp-robot-leasecache-get-last-modified-1}}

查询或检查 \passthrough{\lstinline!last!}
\passthrough{\lstinline!modified!}。

\textbf{签名}

\begin{lstlisting}[language=Python]
def get_last_modified(self) -> int
\end{lstlisting}

\textbf{可用性与安全性}

\textbf{\passthrough{\lstinline!SIGNATURE\_ONLY!} /
\passthrough{\lstinline!UNCLASSIFIED!}}：当前只有设计期签名，不能假设已存在于运行时扩展。

\textbf{参数}

无显式参数。实例方法中的 \passthrough{\lstinline!self!} 由 Python
自动传入。

\textbf{返回值}

\begin{longtable}[]{@{}
  >{\raggedright\arraybackslash}p{(\columnwidth - 4\tabcolsep) * \real{0.18}}
  >{\raggedright\arraybackslash}p{(\columnwidth - 4\tabcolsep) * \real{0.20}}
  >{\raggedright\arraybackslash}p{(\columnwidth - 4\tabcolsep) * \real{0.62}}@{}}
\toprule
位置 & 类型 & 含义 \\
\midrule
\endhead
返回值 & \passthrough{\lstinline!int!} & SDK 状态码；Unitree 示例通常以
\passthrough{\lstinline!0!} 表示成功，非零值需按具体服务定义解释。 \\
\bottomrule
\end{longtable}

\textbf{对应 C++}

\begin{itemize}
\tightlist
\item
  类：\passthrough{\lstinline!unitree::robot::LeaseCache!}
\item
  签名：\passthrough{\lstinline!GetLastModified() const!}
\item
  绑定策略：\passthrough{\lstinline!SIGNATURE\_PREVIEW!}
\item
  声明位置：\passthrough{\lstinline!include/unitree/robot/server/lease\_server.hpp:20!}
\end{itemize}

\textbf{说明}

\begin{itemize}
\tightlist
\item
  示例是局部调用片段；\passthrough{\lstinline!obj!}
  和参数变量表示已经按上层业务校验并准备好的对象。
\item
  该条目可以被编辑器和类型检查器识别，但当前运行时可能在属性查找阶段直接失败。
\end{itemize}

\textbf{用法}

\begin{lstlisting}[language=Python]
from typing import TYPE_CHECKING

if TYPE_CHECKING:
    def planned_get_last_modified(obj: LeaseCache) -> int:
        return obj.get_last_modified()
\end{lstlisting}

\begin{quote}
{[}!CAUTION{]} 上面的 \passthrough{\lstinline!TYPE\_CHECKING!}
示例只表达计划调用形态。该分支在普通 Python
运行时为假，不代表当前扩展能执行此接口。
\end{quote}

\hypertarget{unitree-sdk2-cpp-robot-leasecache-get-id-1}{%
\paragraph{\texorpdfstring{\texttt{unitree\_sdk2\_cpp.robot.LeaseCache.get\_id}}{unitree\_sdk2\_cpp.robot.LeaseCache.get\_id}}\label{unitree-sdk2-cpp-robot-leasecache-get-id-1}}

查询或检查 ID。

\textbf{签名}

\begin{lstlisting}[language=Python]
def get_id(self) -> int
\end{lstlisting}

\textbf{可用性与安全性}

\textbf{\passthrough{\lstinline!SIGNATURE\_ONLY!} /
\passthrough{\lstinline!UNCLASSIFIED!}}：当前只有设计期签名，不能假设已存在于运行时扩展。

\textbf{参数}

无显式参数。实例方法中的 \passthrough{\lstinline!self!} 由 Python
自动传入。

\textbf{返回值}

\begin{longtable}[]{@{}
  >{\raggedright\arraybackslash}p{(\columnwidth - 4\tabcolsep) * \real{0.18}}
  >{\raggedright\arraybackslash}p{(\columnwidth - 4\tabcolsep) * \real{0.20}}
  >{\raggedright\arraybackslash}p{(\columnwidth - 4\tabcolsep) * \real{0.62}}@{}}
\toprule
位置 & 类型 & 含义 \\
\midrule
\endhead
返回值 & \passthrough{\lstinline!int!} & SDK 状态码；Unitree 示例通常以
\passthrough{\lstinline!0!} 表示成功，非零值需按具体服务定义解释。 \\
\bottomrule
\end{longtable}

\textbf{对应 C++}

\begin{itemize}
\tightlist
\item
  类：\passthrough{\lstinline!unitree::robot::LeaseCache!}
\item
  签名：\passthrough{\lstinline!GetId() const!}
\item
  绑定策略：\passthrough{\lstinline!SIGNATURE\_PREVIEW!}
\item
  声明位置：\passthrough{\lstinline!include/unitree/robot/server/lease\_server.hpp:21!}
\end{itemize}

\textbf{说明}

\begin{itemize}
\tightlist
\item
  示例是局部调用片段；\passthrough{\lstinline!obj!}
  和参数变量表示已经按上层业务校验并准备好的对象。
\item
  该条目可以被编辑器和类型检查器识别，但当前运行时可能在属性查找阶段直接失败。
\end{itemize}

\textbf{用法}

\begin{lstlisting}[language=Python]
from typing import TYPE_CHECKING

if TYPE_CHECKING:
    def planned_get_id(obj: LeaseCache) -> int:
        return obj.get_id()
\end{lstlisting}

\begin{quote}
{[}!CAUTION{]} 上面的 \passthrough{\lstinline!TYPE\_CHECKING!}
示例只表达计划调用形态。该分支在普通 Python
运行时为假，不代表当前扩展能执行此接口。
\end{quote}

\hypertarget{unitree-sdk2-cpp-robot-leasecache-get-name-1}{%
\paragraph{\texorpdfstring{\texttt{unitree\_sdk2\_cpp.robot.LeaseCache.get\_name}}{unitree\_sdk2\_cpp.robot.LeaseCache.get\_name}}\label{unitree-sdk2-cpp-robot-leasecache-get-name-1}}

查询或检查 名称。

\textbf{签名}

\begin{lstlisting}[language=Python]
def get_name(self) -> str
\end{lstlisting}

\textbf{可用性与安全性}

\textbf{\passthrough{\lstinline!SIGNATURE\_ONLY!} /
\passthrough{\lstinline!UNCLASSIFIED!}}：当前只有设计期签名，不能假设已存在于运行时扩展。

\textbf{参数}

无显式参数。实例方法中的 \passthrough{\lstinline!self!} 由 Python
自动传入。

\textbf{返回值}

\begin{longtable}[]{@{}
  >{\raggedright\arraybackslash}p{(\columnwidth - 4\tabcolsep) * \real{0.18}}
  >{\raggedright\arraybackslash}p{(\columnwidth - 4\tabcolsep) * \real{0.20}}
  >{\raggedright\arraybackslash}p{(\columnwidth - 4\tabcolsep) * \real{0.62}}@{}}
\toprule
位置 & 类型 & 含义 \\
\midrule
\endhead
返回值 & \passthrough{\lstinline!str!} & 返回
\passthrough{\lstinline!str!}。更精确的业务含义见该方法用途、C++
签名和目标型号协议。 \\
\bottomrule
\end{longtable}

\textbf{对应 C++}

\begin{itemize}
\tightlist
\item
  类：\passthrough{\lstinline!unitree::robot::LeaseCache!}
\item
  签名：\passthrough{\lstinline!GetName() const!}
\item
  绑定策略：\passthrough{\lstinline!SIGNATURE\_PREVIEW!}
\item
  声明位置：\passthrough{\lstinline!include/unitree/robot/server/lease\_server.hpp:22!}
\end{itemize}

\textbf{说明}

\begin{itemize}
\tightlist
\item
  示例是局部调用片段；\passthrough{\lstinline!obj!}
  和参数变量表示已经按上层业务校验并准备好的对象。
\item
  该条目可以被编辑器和类型检查器识别，但当前运行时可能在属性查找阶段直接失败。
\end{itemize}

\textbf{用法}

\begin{lstlisting}[language=Python]
from typing import TYPE_CHECKING

if TYPE_CHECKING:
    def planned_get_name(obj: LeaseCache) -> str:
        return obj.get_name()
\end{lstlisting}

\begin{quote}
{[}!CAUTION{]} 上面的 \passthrough{\lstinline!TYPE\_CHECKING!}
示例只表达计划调用形态。该分支在普通 Python
运行时为假，不代表当前扩展能执行此接口。
\end{quote}

\hypertarget{unitree-sdk2-cpp-robot-leaseclient}{%
\subsubsection{\texorpdfstring{\texttt{unitree\_sdk2\_cpp.robot.LeaseClient}}{unitree\_sdk2\_cpp.robot.LeaseClient}}\label{unitree-sdk2-cpp-robot-leaseclient}}

Robot 服务客户端。构造、初始化和具体方法的可用性必须分别检查。

\textbf{导入}

\begin{lstlisting}[language=Python]
from unitree_sdk2_cpp.robot import LeaseClient
\end{lstlisting}

\textbf{基类}：\passthrough{\lstinline!ClientBase!}

\textbf{方法索引}

\begin{longtable}[]{@{}
  >{\raggedright\arraybackslash}p{(\columnwidth - 6\tabcolsep) * \real{0.18}}
  >{\raggedright\arraybackslash}p{(\columnwidth - 6\tabcolsep) * \real{0.24}}
  >{\raggedright\arraybackslash}p{(\columnwidth - 6\tabcolsep) * \real{0.18}}
  >{\raggedright\arraybackslash}p{(\columnwidth - 6\tabcolsep) * \real{0.40}}@{}}
\toprule
方法 & Python 签名 & 状态 & 安全分类 \\
\midrule
\endhead
\protect\hyperlink{unitree-sdk2-cpp-robot-leaseclient-dunder-init-1}{\passthrough{\lstinline!\_\_init\_\_!}}
& \passthrough{\lstinline!def \_\_init\_\_(self, name: str) -> None!} &
\passthrough{\lstinline!SIGNATURE\_ONLY!} &
\passthrough{\lstinline!CONSTRUCTION!} \\
\protect\hyperlink{unitree-sdk2-cpp-robot-leaseclient-init-1}{\passthrough{\lstinline!init!}}
& \passthrough{\lstinline!def init(self) -> None!} &
\passthrough{\lstinline!SIGNATURE\_ONLY!} &
\passthrough{\lstinline!INITIALIZATION!} \\
\protect\hyperlink{unitree-sdk2-cpp-robot-leaseclient-wait-applied-1}{\passthrough{\lstinline!wait\_applied!}}
& \passthrough{\lstinline!def wait\_applied(self) -> None!} &
\passthrough{\lstinline!SIGNATURE\_ONLY!} &
\passthrough{\lstinline!INITIALIZATION!} \\
\protect\hyperlink{unitree-sdk2-cpp-robot-leaseclient-get-id-1}{\passthrough{\lstinline!get\_id!}}
& \passthrough{\lstinline!def get\_id(self) -> int!} &
\passthrough{\lstinline!SIGNATURE\_ONLY!} &
\passthrough{\lstinline!READ\_ONLY!} \\
\protect\hyperlink{unitree-sdk2-cpp-robot-leaseclient-applied-1}{\passthrough{\lstinline!applied!}}
& \passthrough{\lstinline!def applied(self) -> bool!} &
\passthrough{\lstinline!SIGNATURE\_ONLY!} &
\passthrough{\lstinline!READ\_ONLY!} \\
\bottomrule
\end{longtable}

\hypertarget{unitree-sdk2-cpp-robot-leaseclient-dunder-init-1}{%
\paragraph{\texorpdfstring{\texttt{unitree\_sdk2\_cpp.robot.LeaseClient.\_\_init\_\_}}{unitree\_sdk2\_cpp.robot.LeaseClient.\_\_init\_\_}}\label{unitree-sdk2-cpp-robot-leaseclient-dunder-init-1}}

初始化 \passthrough{\lstinline!LeaseClient!}
实例。是否能实际构造取决于下方可用性状态。

\textbf{签名}

\begin{lstlisting}[language=Python]
def __init__(self, name: str) -> None
\end{lstlisting}

\textbf{可用性与安全性}

\textbf{\passthrough{\lstinline!SIGNATURE\_ONLY!} /
\passthrough{\lstinline!CONSTRUCTION!}}：当前只有设计期签名，不能假设已存在于运行时扩展。

\textbf{参数}

\begin{longtable}[]{@{}
  >{\raggedright\arraybackslash}p{(\columnwidth - 6\tabcolsep) * \real{0.18}}
  >{\raggedright\arraybackslash}p{(\columnwidth - 6\tabcolsep) * \real{0.24}}
  >{\raggedright\arraybackslash}p{(\columnwidth - 6\tabcolsep) * \real{0.18}}
  >{\raggedright\arraybackslash}p{(\columnwidth - 6\tabcolsep) * \real{0.40}}@{}}
\toprule
名称 & Python 类型 & 默认值 & 含义 \\
\midrule
\endhead
\passthrough{\lstinline!name!} & \passthrough{\lstinline!str!} & 必填 &
SDK 对象、配置项、服务或动作的名称。允许值由调用它的具体接口定义。 对应
C++ 参数 \passthrough{\lstinline!name: const std::string \&!}。
底层为字符串；长度、编码和允许值由具体协议决定。 \\
\bottomrule
\end{longtable}

\textbf{返回值}

\begin{longtable}[]{@{}
  >{\raggedright\arraybackslash}p{(\columnwidth - 4\tabcolsep) * \real{0.18}}
  >{\raggedright\arraybackslash}p{(\columnwidth - 4\tabcolsep) * \real{0.20}}
  >{\raggedright\arraybackslash}p{(\columnwidth - 4\tabcolsep) * \real{0.62}}@{}}
\toprule
位置 & 类型 & 含义 \\
\midrule
\endhead
无 & \passthrough{\lstinline!None!} &
\passthrough{\lstinline!\_\_init\_\_!}
本身不返回值；调用类对象时，在构造函数可用的前提下得到该类实例。 \\
\bottomrule
\end{longtable}

\textbf{对应 C++}

\begin{itemize}
\tightlist
\item
  类：\passthrough{\lstinline!unitree::robot::LeaseClient!}
\item
  签名：\passthrough{\lstinline!LeaseClient(const std::string \&)!}
\item
  绑定策略：\passthrough{\lstinline!未单独标注!}
\item
  声明位置：\passthrough{\lstinline!include/unitree/robot/client/lease\_client.hpp:35!}
\end{itemize}

\textbf{说明}

\begin{itemize}
\tightlist
\item
  示例是局部调用片段；\passthrough{\lstinline!obj!}
  和参数变量表示已经按上层业务校验并准备好的对象。
\item
  该条目可以被编辑器和类型检查器识别，但当前运行时可能在属性查找阶段直接失败。
\end{itemize}

\textbf{用法}

\begin{lstlisting}[language=Python]
from typing import TYPE_CHECKING

if TYPE_CHECKING:
    def planned_construct(name: str) -> LeaseClient:
        return LeaseClient(name=name)
\end{lstlisting}

\begin{quote}
{[}!CAUTION{]} 上面的 \passthrough{\lstinline!TYPE\_CHECKING!}
示例只表达计划调用形态。该分支在普通 Python
运行时为假，不代表当前扩展能执行此接口。
\end{quote}

\hypertarget{unitree-sdk2-cpp-robot-leaseclient-init-1}{%
\paragraph{\texorpdfstring{\texttt{unitree\_sdk2\_cpp.robot.LeaseClient.init}}{unitree\_sdk2\_cpp.robot.LeaseClient.init}}\label{unitree-sdk2-cpp-robot-leaseclient-init-1}}

初始化当前 SDK 对象所需的底层通道或服务资源。

\textbf{签名}

\begin{lstlisting}[language=Python]
def init(self) -> None
\end{lstlisting}

\textbf{可用性与安全性}

\textbf{\passthrough{\lstinline!SIGNATURE\_ONLY!} /
\passthrough{\lstinline!INITIALIZATION!}}：当前只有设计期签名，不能假设已存在于运行时扩展。

\textbf{参数}

无显式参数。实例方法中的 \passthrough{\lstinline!self!} 由 Python
自动传入。

\textbf{返回值}

\begin{longtable}[]{@{}
  >{\raggedright\arraybackslash}p{(\columnwidth - 4\tabcolsep) * \real{0.18}}
  >{\raggedright\arraybackslash}p{(\columnwidth - 4\tabcolsep) * \real{0.20}}
  >{\raggedright\arraybackslash}p{(\columnwidth - 4\tabcolsep) * \real{0.62}}@{}}
\toprule
位置 & 类型 & 含义 \\
\midrule
\endhead
无 & \passthrough{\lstinline!None!} & 该方法不返回 Python
值；失败可能表现为异常或底层状态变化。 \\
\bottomrule
\end{longtable}

\textbf{对应 C++}

\begin{itemize}
\tightlist
\item
  类：\passthrough{\lstinline!unitree::robot::LeaseClient!}
\item
  签名：\passthrough{\lstinline!Init()!}
\item
  绑定策略：\passthrough{\lstinline!DIRECT!}
\item
  声明位置：\passthrough{\lstinline!include/unitree/robot/client/lease\_client.hpp:38!}
\end{itemize}

\textbf{说明}

\begin{itemize}
\tightlist
\item
  示例是局部调用片段；\passthrough{\lstinline!obj!}
  和参数变量表示已经按上层业务校验并准备好的对象。
\item
  该条目可以被编辑器和类型检查器识别，但当前运行时可能在属性查找阶段直接失败。
\end{itemize}

\textbf{用法}

\begin{lstlisting}[language=Python]
from typing import TYPE_CHECKING

if TYPE_CHECKING:
    def planned_init(obj: LeaseClient) -> None:
        obj.init()
\end{lstlisting}

\begin{quote}
{[}!CAUTION{]} 上面的 \passthrough{\lstinline!TYPE\_CHECKING!}
示例只表达计划调用形态。该分支在普通 Python
运行时为假，不代表当前扩展能执行此接口。
\end{quote}

\hypertarget{unitree-sdk2-cpp-robot-leaseclient-wait-applied-1}{%
\paragraph{\texorpdfstring{\texttt{unitree\_sdk2\_cpp.robot.LeaseClient.wait\_applied}}{unitree\_sdk2\_cpp.robot.LeaseClient.wait\_applied}}\label{unitree-sdk2-cpp-robot-leaseclient-wait-applied-1}}

对应 C++ SDK 操作
\passthrough{\lstinline!WaitApplied()!}。上游头文件没有可直接生成的业务说明时，本参考只保证签名映射，精确语义需查目标型号协议。

\textbf{签名}

\begin{lstlisting}[language=Python]
def wait_applied(self) -> None
\end{lstlisting}

\textbf{可用性与安全性}

\textbf{\passthrough{\lstinline!SIGNATURE\_ONLY!} /
\passthrough{\lstinline!INITIALIZATION!}}：当前只有设计期签名，不能假设已存在于运行时扩展。

\textbf{参数}

无显式参数。实例方法中的 \passthrough{\lstinline!self!} 由 Python
自动传入。

\textbf{返回值}

\begin{longtable}[]{@{}
  >{\raggedright\arraybackslash}p{(\columnwidth - 4\tabcolsep) * \real{0.18}}
  >{\raggedright\arraybackslash}p{(\columnwidth - 4\tabcolsep) * \real{0.20}}
  >{\raggedright\arraybackslash}p{(\columnwidth - 4\tabcolsep) * \real{0.62}}@{}}
\toprule
位置 & 类型 & 含义 \\
\midrule
\endhead
无 & \passthrough{\lstinline!None!} & 该方法不返回 Python
值；失败可能表现为异常或底层状态变化。 \\
\bottomrule
\end{longtable}

\textbf{对应 C++}

\begin{itemize}
\tightlist
\item
  类：\passthrough{\lstinline!unitree::robot::LeaseClient!}
\item
  签名：\passthrough{\lstinline!WaitApplied()!}
\item
  绑定策略：\passthrough{\lstinline!DIRECT!}
\item
  声明位置：\passthrough{\lstinline!include/unitree/robot/client/lease\_client.hpp:40!}
\end{itemize}

\textbf{说明}

\begin{itemize}
\tightlist
\item
  示例是局部调用片段；\passthrough{\lstinline!obj!}
  和参数变量表示已经按上层业务校验并准备好的对象。
\item
  该条目可以被编辑器和类型检查器识别，但当前运行时可能在属性查找阶段直接失败。
\end{itemize}

\textbf{用法}

\begin{lstlisting}[language=Python]
from typing import TYPE_CHECKING

if TYPE_CHECKING:
    def planned_wait_applied(obj: LeaseClient) -> None:
        obj.wait_applied()
\end{lstlisting}

\begin{quote}
{[}!CAUTION{]} 上面的 \passthrough{\lstinline!TYPE\_CHECKING!}
示例只表达计划调用形态。该分支在普通 Python
运行时为假，不代表当前扩展能执行此接口。
\end{quote}

\hypertarget{unitree-sdk2-cpp-robot-leaseclient-get-id-1}{%
\paragraph{\texorpdfstring{\texttt{unitree\_sdk2\_cpp.robot.LeaseClient.get\_id}}{unitree\_sdk2\_cpp.robot.LeaseClient.get\_id}}\label{unitree-sdk2-cpp-robot-leaseclient-get-id-1}}

查询或检查 ID。

\textbf{签名}

\begin{lstlisting}[language=Python]
def get_id(self) -> int
\end{lstlisting}

\textbf{可用性与安全性}

\textbf{\passthrough{\lstinline!SIGNATURE\_ONLY!} /
\passthrough{\lstinline!READ\_ONLY!}}：当前只有设计期签名，不能假设已存在于运行时扩展。

\textbf{参数}

无显式参数。实例方法中的 \passthrough{\lstinline!self!} 由 Python
自动传入。

\textbf{返回值}

\begin{longtable}[]{@{}
  >{\raggedright\arraybackslash}p{(\columnwidth - 4\tabcolsep) * \real{0.18}}
  >{\raggedright\arraybackslash}p{(\columnwidth - 4\tabcolsep) * \real{0.20}}
  >{\raggedright\arraybackslash}p{(\columnwidth - 4\tabcolsep) * \real{0.62}}@{}}
\toprule
位置 & 类型 & 含义 \\
\midrule
\endhead
返回值 & \passthrough{\lstinline!int!} & SDK 状态码；Unitree 示例通常以
\passthrough{\lstinline!0!} 表示成功，非零值需按具体服务定义解释。 \\
\bottomrule
\end{longtable}

\textbf{对应 C++}

\begin{itemize}
\tightlist
\item
  类：\passthrough{\lstinline!unitree::robot::LeaseClient!}
\item
  签名：\passthrough{\lstinline!GetId()!}
\item
  绑定策略：\passthrough{\lstinline!DIRECT!}
\item
  声明位置：\passthrough{\lstinline!include/unitree/robot/client/lease\_client.hpp:41!}
\end{itemize}

\textbf{说明}

\begin{itemize}
\tightlist
\item
  示例是局部调用片段；\passthrough{\lstinline!obj!}
  和参数变量表示已经按上层业务校验并准备好的对象。
\item
  该条目可以被编辑器和类型检查器识别，但当前运行时可能在属性查找阶段直接失败。
\end{itemize}

\textbf{用法}

\begin{lstlisting}[language=Python]
from typing import TYPE_CHECKING

if TYPE_CHECKING:
    def planned_get_id(obj: LeaseClient) -> int:
        return obj.get_id()
\end{lstlisting}

\begin{quote}
{[}!CAUTION{]} 上面的 \passthrough{\lstinline!TYPE\_CHECKING!}
示例只表达计划调用形态。该分支在普通 Python
运行时为假，不代表当前扩展能执行此接口。
\end{quote}

\hypertarget{unitree-sdk2-cpp-robot-leaseclient-applied-1}{%
\paragraph{\texorpdfstring{\texttt{unitree\_sdk2\_cpp.robot.LeaseClient.applied}}{unitree\_sdk2\_cpp.robot.LeaseClient.applied}}\label{unitree-sdk2-cpp-robot-leaseclient-applied-1}}

对应 C++ SDK 操作
\passthrough{\lstinline!Applied()!}。上游头文件没有可直接生成的业务说明时，本参考只保证签名映射，精确语义需查目标型号协议。

\textbf{签名}

\begin{lstlisting}[language=Python]
def applied(self) -> bool
\end{lstlisting}

\textbf{可用性与安全性}

\textbf{\passthrough{\lstinline!SIGNATURE\_ONLY!} /
\passthrough{\lstinline!READ\_ONLY!}}：当前只有设计期签名，不能假设已存在于运行时扩展。

\textbf{参数}

无显式参数。实例方法中的 \passthrough{\lstinline!self!} 由 Python
自动传入。

\textbf{返回值}

\begin{longtable}[]{@{}
  >{\raggedright\arraybackslash}p{(\columnwidth - 4\tabcolsep) * \real{0.18}}
  >{\raggedright\arraybackslash}p{(\columnwidth - 4\tabcolsep) * \real{0.20}}
  >{\raggedright\arraybackslash}p{(\columnwidth - 4\tabcolsep) * \real{0.62}}@{}}
\toprule
位置 & 类型 & 含义 \\
\midrule
\endhead
返回值 & \passthrough{\lstinline!bool!} & 返回
\passthrough{\lstinline!bool!}。更精确的业务含义见该方法用途、C++
签名和目标型号协议。 \\
\bottomrule
\end{longtable}

\textbf{对应 C++}

\begin{itemize}
\tightlist
\item
  类：\passthrough{\lstinline!unitree::robot::LeaseClient!}
\item
  签名：\passthrough{\lstinline!Applied()!}
\item
  绑定策略：\passthrough{\lstinline!DIRECT!}
\item
  声明位置：\passthrough{\lstinline!include/unitree/robot/client/lease\_client.hpp:42!}
\end{itemize}

\textbf{说明}

\begin{itemize}
\tightlist
\item
  示例是局部调用片段；\passthrough{\lstinline!obj!}
  和参数变量表示已经按上层业务校验并准备好的对象。
\item
  该条目可以被编辑器和类型检查器识别，但当前运行时可能在属性查找阶段直接失败。
\end{itemize}

\textbf{用法}

\begin{lstlisting}[language=Python]
from typing import TYPE_CHECKING

if TYPE_CHECKING:
    def planned_applied(obj: LeaseClient) -> bool:
        return obj.applied()
\end{lstlisting}

\begin{quote}
{[}!CAUTION{]} 上面的 \passthrough{\lstinline!TYPE\_CHECKING!}
示例只表达计划调用形态。该分支在普通 Python
运行时为假，不代表当前扩展能执行此接口。
\end{quote}

\hypertarget{unitree-sdk2-cpp-robot-leasecontext}{%
\subsubsection{\texorpdfstring{\texttt{unitree\_sdk2\_cpp.robot.LeaseContext}}{unitree\_sdk2\_cpp.robot.LeaseContext}}\label{unitree-sdk2-cpp-robot-leasecontext}}

SDK 类型签名预览。请逐项查看构造函数和方法的可用性。

\textbf{导入}

\begin{lstlisting}[language=Python]
from unitree_sdk2_cpp.robot import LeaseContext
\end{lstlisting}

\textbf{基类}：\passthrough{\lstinline!object!}

\textbf{方法索引}

\begin{longtable}[]{@{}
  >{\raggedright\arraybackslash}p{(\columnwidth - 6\tabcolsep) * \real{0.18}}
  >{\raggedright\arraybackslash}p{(\columnwidth - 6\tabcolsep) * \real{0.24}}
  >{\raggedright\arraybackslash}p{(\columnwidth - 6\tabcolsep) * \real{0.18}}
  >{\raggedright\arraybackslash}p{(\columnwidth - 6\tabcolsep) * \real{0.40}}@{}}
\toprule
方法 & Python 签名 & 状态 & 安全分类 \\
\midrule
\endhead
\protect\hyperlink{unitree-sdk2-cpp-robot-leasecontext-dunder-init-1}{\passthrough{\lstinline!\_\_init\_\_!}}
& \passthrough{\lstinline!def \_\_init\_\_(self) -> None!} &
\passthrough{\lstinline!SIGNATURE\_ONLY!} &
\passthrough{\lstinline!CONSTRUCTION!} \\
\protect\hyperlink{unitree-sdk2-cpp-robot-leasecontext-update-1}{\passthrough{\lstinline!update!}}
& \passthrough{\lstinline!def update(self, id: int, term: int) -> None!}
& \passthrough{\lstinline!SIGNATURE\_ONLY!} &
\passthrough{\lstinline!UNCLASSIFIED!} \\
\protect\hyperlink{unitree-sdk2-cpp-robot-leasecontext-reset-1}{\passthrough{\lstinline!reset!}}
& \passthrough{\lstinline!def reset(self) -> None!} &
\passthrough{\lstinline!SIGNATURE\_ONLY!} &
\passthrough{\lstinline!UNCLASSIFIED!} \\
\protect\hyperlink{unitree-sdk2-cpp-robot-leasecontext-valid-1}{\passthrough{\lstinline!valid!}}
& \passthrough{\lstinline!def valid(self) -> bool!} &
\passthrough{\lstinline!SIGNATURE\_ONLY!} &
\passthrough{\lstinline!UNCLASSIFIED!} \\
\protect\hyperlink{unitree-sdk2-cpp-robot-leasecontext-get-id-1}{\passthrough{\lstinline!get\_id!}}
& \passthrough{\lstinline!def get\_id(self) -> int!} &
\passthrough{\lstinline!SIGNATURE\_ONLY!} &
\passthrough{\lstinline!UNCLASSIFIED!} \\
\protect\hyperlink{unitree-sdk2-cpp-robot-leasecontext-get-term-1}{\passthrough{\lstinline!get\_term!}}
& \passthrough{\lstinline!def get\_term(self) -> int!} &
\passthrough{\lstinline!SIGNATURE\_ONLY!} &
\passthrough{\lstinline!UNCLASSIFIED!} \\
\bottomrule
\end{longtable}

\hypertarget{unitree-sdk2-cpp-robot-leasecontext-dunder-init-1}{%
\paragraph{\texorpdfstring{\texttt{unitree\_sdk2\_cpp.robot.LeaseContext.\_\_init\_\_}}{unitree\_sdk2\_cpp.robot.LeaseContext.\_\_init\_\_}}\label{unitree-sdk2-cpp-robot-leasecontext-dunder-init-1}}

初始化 \passthrough{\lstinline!LeaseContext!}
实例。是否能实际构造取决于下方可用性状态。

\textbf{签名}

\begin{lstlisting}[language=Python]
def __init__(self) -> None
\end{lstlisting}

\textbf{可用性与安全性}

\textbf{\passthrough{\lstinline!SIGNATURE\_ONLY!} /
\passthrough{\lstinline!CONSTRUCTION!}}：当前只有设计期签名，不能假设已存在于运行时扩展。

\textbf{参数}

无显式参数。实例方法中的 \passthrough{\lstinline!self!} 由 Python
自动传入。

\textbf{返回值}

\begin{longtable}[]{@{}
  >{\raggedright\arraybackslash}p{(\columnwidth - 4\tabcolsep) * \real{0.18}}
  >{\raggedright\arraybackslash}p{(\columnwidth - 4\tabcolsep) * \real{0.20}}
  >{\raggedright\arraybackslash}p{(\columnwidth - 4\tabcolsep) * \real{0.62}}@{}}
\toprule
位置 & 类型 & 含义 \\
\midrule
\endhead
无 & \passthrough{\lstinline!None!} &
\passthrough{\lstinline!\_\_init\_\_!}
本身不返回值；调用类对象时，在构造函数可用的前提下得到该类实例。 \\
\bottomrule
\end{longtable}

\textbf{对应 C++}

\begin{itemize}
\tightlist
\item
  类：\passthrough{\lstinline!unitree::robot::LeaseContext!}
\item
  签名：\passthrough{\lstinline!LeaseContext()!}
\item
  绑定策略：\passthrough{\lstinline!未单独标注!}
\item
  声明位置：\passthrough{\lstinline!include/unitree/robot/client/lease\_client.hpp:14!}
\end{itemize}

\textbf{说明}

\begin{itemize}
\tightlist
\item
  示例是局部调用片段；\passthrough{\lstinline!obj!}
  和参数变量表示已经按上层业务校验并准备好的对象。
\item
  该条目可以被编辑器和类型检查器识别，但当前运行时可能在属性查找阶段直接失败。
\end{itemize}

\textbf{用法}

\begin{lstlisting}[language=Python]
from typing import TYPE_CHECKING

if TYPE_CHECKING:
    def planned_construct() -> LeaseContext:
        return LeaseContext()
\end{lstlisting}

\begin{quote}
{[}!CAUTION{]} 上面的 \passthrough{\lstinline!TYPE\_CHECKING!}
示例只表达计划调用形态。该分支在普通 Python
运行时为假，不代表当前扩展能执行此接口。
\end{quote}

\hypertarget{unitree-sdk2-cpp-robot-leasecontext-update-1}{%
\paragraph{\texorpdfstring{\texttt{unitree\_sdk2\_cpp.robot.LeaseContext.update}}{unitree\_sdk2\_cpp.robot.LeaseContext.update}}\label{unitree-sdk2-cpp-robot-leasecontext-update-1}}

对应 C++ SDK 操作
\passthrough{\lstinline!Update(int64\_t, int64\_t)!}。上游头文件没有可直接生成的业务说明时，本参考只保证签名映射，精确语义需查目标型号协议。

\textbf{签名}

\begin{lstlisting}[language=Python]
def update(self, id: int, term: int) -> None
\end{lstlisting}

\textbf{可用性与安全性}

\textbf{\passthrough{\lstinline!SIGNATURE\_ONLY!} /
\passthrough{\lstinline!UNCLASSIFIED!}}：当前只有设计期签名，不能假设已存在于运行时扩展。

\textbf{参数}

\begin{longtable}[]{@{}
  >{\raggedright\arraybackslash}p{(\columnwidth - 6\tabcolsep) * \real{0.18}}
  >{\raggedright\arraybackslash}p{(\columnwidth - 6\tabcolsep) * \real{0.24}}
  >{\raggedright\arraybackslash}p{(\columnwidth - 6\tabcolsep) * \real{0.18}}
  >{\raggedright\arraybackslash}p{(\columnwidth - 6\tabcolsep) * \real{0.40}}@{}}
\toprule
名称 & Python 类型 & 默认值 & 含义 \\
\midrule
\endhead
\passthrough{\lstinline!id!} & \passthrough{\lstinline!int!} & 必填 &
传给该接口的 ID 参数，Python 类型为
\passthrough{\lstinline!int!}。精确含义、单位、范围和枚举值需查目标型号对应头文件与协议。
对应 C++ 参数 \passthrough{\lstinline!id: int64\_t!}。 取值范围为
-9223372036854775808 到 9223372036854775807。 \\
\passthrough{\lstinline!term!} & \passthrough{\lstinline!int!} & 必填 &
传给该接口的 \passthrough{\lstinline!term!} 参数，Python 类型为
\passthrough{\lstinline!int!}。精确含义、单位、范围和枚举值需查目标型号对应头文件与协议。
对应 C++ 参数 \passthrough{\lstinline!term: int64\_t!}。 取值范围为
-9223372036854775808 到 9223372036854775807。 \\
\bottomrule
\end{longtable}

\textbf{返回值}

\begin{longtable}[]{@{}
  >{\raggedright\arraybackslash}p{(\columnwidth - 4\tabcolsep) * \real{0.18}}
  >{\raggedright\arraybackslash}p{(\columnwidth - 4\tabcolsep) * \real{0.20}}
  >{\raggedright\arraybackslash}p{(\columnwidth - 4\tabcolsep) * \real{0.62}}@{}}
\toprule
位置 & 类型 & 含义 \\
\midrule
\endhead
无 & \passthrough{\lstinline!None!} & 该方法不返回 Python
值；失败可能表现为异常或底层状态变化。 \\
\bottomrule
\end{longtable}

\textbf{对应 C++}

\begin{itemize}
\tightlist
\item
  类：\passthrough{\lstinline!unitree::robot::LeaseContext!}
\item
  签名：\passthrough{\lstinline!Update(int64\_t, int64\_t)!}
\item
  绑定策略：\passthrough{\lstinline!SIGNATURE\_PREVIEW!}
\item
  声明位置：\passthrough{\lstinline!include/unitree/robot/client/lease\_client.hpp:17!}
\end{itemize}

\textbf{说明}

\begin{itemize}
\tightlist
\item
  示例是局部调用片段；\passthrough{\lstinline!obj!}
  和参数变量表示已经按上层业务校验并准备好的对象。
\item
  该条目可以被编辑器和类型检查器识别，但当前运行时可能在属性查找阶段直接失败。
\end{itemize}

\textbf{用法}

\begin{lstlisting}[language=Python]
from typing import TYPE_CHECKING

if TYPE_CHECKING:
    def planned_update(obj: LeaseContext, id: int, term: int) -> None:
        obj.update(id=id, term=term)
\end{lstlisting}

\begin{quote}
{[}!CAUTION{]} 上面的 \passthrough{\lstinline!TYPE\_CHECKING!}
示例只表达计划调用形态。该分支在普通 Python
运行时为假，不代表当前扩展能执行此接口。
\end{quote}

\hypertarget{unitree-sdk2-cpp-robot-leasecontext-reset-1}{%
\paragraph{\texorpdfstring{\texttt{unitree\_sdk2\_cpp.robot.LeaseContext.reset}}{unitree\_sdk2\_cpp.robot.LeaseContext.reset}}\label{unitree-sdk2-cpp-robot-leasecontext-reset-1}}

对应 C++ SDK 操作
\passthrough{\lstinline!Reset()!}。上游头文件没有可直接生成的业务说明时，本参考只保证签名映射，精确语义需查目标型号协议。

\textbf{签名}

\begin{lstlisting}[language=Python]
def reset(self) -> None
\end{lstlisting}

\textbf{可用性与安全性}

\textbf{\passthrough{\lstinline!SIGNATURE\_ONLY!} /
\passthrough{\lstinline!UNCLASSIFIED!}}：当前只有设计期签名，不能假设已存在于运行时扩展。

\textbf{参数}

无显式参数。实例方法中的 \passthrough{\lstinline!self!} 由 Python
自动传入。

\textbf{返回值}

\begin{longtable}[]{@{}
  >{\raggedright\arraybackslash}p{(\columnwidth - 4\tabcolsep) * \real{0.18}}
  >{\raggedright\arraybackslash}p{(\columnwidth - 4\tabcolsep) * \real{0.20}}
  >{\raggedright\arraybackslash}p{(\columnwidth - 4\tabcolsep) * \real{0.62}}@{}}
\toprule
位置 & 类型 & 含义 \\
\midrule
\endhead
无 & \passthrough{\lstinline!None!} & 该方法不返回 Python
值；失败可能表现为异常或底层状态变化。 \\
\bottomrule
\end{longtable}

\textbf{对应 C++}

\begin{itemize}
\tightlist
\item
  类：\passthrough{\lstinline!unitree::robot::LeaseContext!}
\item
  签名：\passthrough{\lstinline!Reset()!}
\item
  绑定策略：\passthrough{\lstinline!SIGNATURE\_PREVIEW!}
\item
  声明位置：\passthrough{\lstinline!include/unitree/robot/client/lease\_client.hpp:18!}
\end{itemize}

\textbf{说明}

\begin{itemize}
\tightlist
\item
  示例是局部调用片段；\passthrough{\lstinline!obj!}
  和参数变量表示已经按上层业务校验并准备好的对象。
\item
  该条目可以被编辑器和类型检查器识别，但当前运行时可能在属性查找阶段直接失败。
\end{itemize}

\textbf{用法}

\begin{lstlisting}[language=Python]
from typing import TYPE_CHECKING

if TYPE_CHECKING:
    def planned_reset(obj: LeaseContext) -> None:
        obj.reset()
\end{lstlisting}

\begin{quote}
{[}!CAUTION{]} 上面的 \passthrough{\lstinline!TYPE\_CHECKING!}
示例只表达计划调用形态。该分支在普通 Python
运行时为假，不代表当前扩展能执行此接口。
\end{quote}

\hypertarget{unitree-sdk2-cpp-robot-leasecontext-valid-1}{%
\paragraph{\texorpdfstring{\texttt{unitree\_sdk2\_cpp.robot.LeaseContext.valid}}{unitree\_sdk2\_cpp.robot.LeaseContext.valid}}\label{unitree-sdk2-cpp-robot-leasecontext-valid-1}}

对应 C++ SDK 操作
\passthrough{\lstinline!Valid() const!}。上游头文件没有可直接生成的业务说明时，本参考只保证签名映射，精确语义需查目标型号协议。

\textbf{签名}

\begin{lstlisting}[language=Python]
def valid(self) -> bool
\end{lstlisting}

\textbf{可用性与安全性}

\textbf{\passthrough{\lstinline!SIGNATURE\_ONLY!} /
\passthrough{\lstinline!UNCLASSIFIED!}}：当前只有设计期签名，不能假设已存在于运行时扩展。

\textbf{参数}

无显式参数。实例方法中的 \passthrough{\lstinline!self!} 由 Python
自动传入。

\textbf{返回值}

\begin{longtable}[]{@{}
  >{\raggedright\arraybackslash}p{(\columnwidth - 4\tabcolsep) * \real{0.18}}
  >{\raggedright\arraybackslash}p{(\columnwidth - 4\tabcolsep) * \real{0.20}}
  >{\raggedright\arraybackslash}p{(\columnwidth - 4\tabcolsep) * \real{0.62}}@{}}
\toprule
位置 & 类型 & 含义 \\
\midrule
\endhead
返回值 & \passthrough{\lstinline!bool!} & 返回
\passthrough{\lstinline!bool!}。更精确的业务含义见该方法用途、C++
签名和目标型号协议。 \\
\bottomrule
\end{longtable}

\textbf{对应 C++}

\begin{itemize}
\tightlist
\item
  类：\passthrough{\lstinline!unitree::robot::LeaseContext!}
\item
  签名：\passthrough{\lstinline!Valid() const!}
\item
  绑定策略：\passthrough{\lstinline!SIGNATURE\_PREVIEW!}
\item
  声明位置：\passthrough{\lstinline!include/unitree/robot/client/lease\_client.hpp:20!}
\end{itemize}

\textbf{说明}

\begin{itemize}
\tightlist
\item
  示例是局部调用片段；\passthrough{\lstinline!obj!}
  和参数变量表示已经按上层业务校验并准备好的对象。
\item
  该条目可以被编辑器和类型检查器识别，但当前运行时可能在属性查找阶段直接失败。
\end{itemize}

\textbf{用法}

\begin{lstlisting}[language=Python]
from typing import TYPE_CHECKING

if TYPE_CHECKING:
    def planned_valid(obj: LeaseContext) -> bool:
        return obj.valid()
\end{lstlisting}

\begin{quote}
{[}!CAUTION{]} 上面的 \passthrough{\lstinline!TYPE\_CHECKING!}
示例只表达计划调用形态。该分支在普通 Python
运行时为假，不代表当前扩展能执行此接口。
\end{quote}

\hypertarget{unitree-sdk2-cpp-robot-leasecontext-get-id-1}{%
\paragraph{\texorpdfstring{\texttt{unitree\_sdk2\_cpp.robot.LeaseContext.get\_id}}{unitree\_sdk2\_cpp.robot.LeaseContext.get\_id}}\label{unitree-sdk2-cpp-robot-leasecontext-get-id-1}}

查询或检查 ID。

\textbf{签名}

\begin{lstlisting}[language=Python]
def get_id(self) -> int
\end{lstlisting}

\textbf{可用性与安全性}

\textbf{\passthrough{\lstinline!SIGNATURE\_ONLY!} /
\passthrough{\lstinline!UNCLASSIFIED!}}：当前只有设计期签名，不能假设已存在于运行时扩展。

\textbf{参数}

无显式参数。实例方法中的 \passthrough{\lstinline!self!} 由 Python
自动传入。

\textbf{返回值}

\begin{longtable}[]{@{}
  >{\raggedright\arraybackslash}p{(\columnwidth - 4\tabcolsep) * \real{0.18}}
  >{\raggedright\arraybackslash}p{(\columnwidth - 4\tabcolsep) * \real{0.20}}
  >{\raggedright\arraybackslash}p{(\columnwidth - 4\tabcolsep) * \real{0.62}}@{}}
\toprule
位置 & 类型 & 含义 \\
\midrule
\endhead
返回值 & \passthrough{\lstinline!int!} & SDK 状态码；Unitree 示例通常以
\passthrough{\lstinline!0!} 表示成功，非零值需按具体服务定义解释。 \\
\bottomrule
\end{longtable}

\textbf{对应 C++}

\begin{itemize}
\tightlist
\item
  类：\passthrough{\lstinline!unitree::robot::LeaseContext!}
\item
  签名：\passthrough{\lstinline!GetId() const!}
\item
  绑定策略：\passthrough{\lstinline!SIGNATURE\_PREVIEW!}
\item
  声明位置：\passthrough{\lstinline!include/unitree/robot/client/lease\_client.hpp:22!}
\end{itemize}

\textbf{说明}

\begin{itemize}
\tightlist
\item
  示例是局部调用片段；\passthrough{\lstinline!obj!}
  和参数变量表示已经按上层业务校验并准备好的对象。
\item
  该条目可以被编辑器和类型检查器识别，但当前运行时可能在属性查找阶段直接失败。
\end{itemize}

\textbf{用法}

\begin{lstlisting}[language=Python]
from typing import TYPE_CHECKING

if TYPE_CHECKING:
    def planned_get_id(obj: LeaseContext) -> int:
        return obj.get_id()
\end{lstlisting}

\begin{quote}
{[}!CAUTION{]} 上面的 \passthrough{\lstinline!TYPE\_CHECKING!}
示例只表达计划调用形态。该分支在普通 Python
运行时为假，不代表当前扩展能执行此接口。
\end{quote}

\hypertarget{unitree-sdk2-cpp-robot-leasecontext-get-term-1}{%
\paragraph{\texorpdfstring{\texttt{unitree\_sdk2\_cpp.robot.LeaseContext.get\_term}}{unitree\_sdk2\_cpp.robot.LeaseContext.get\_term}}\label{unitree-sdk2-cpp-robot-leasecontext-get-term-1}}

查询或检查 \passthrough{\lstinline!term!}。

\textbf{签名}

\begin{lstlisting}[language=Python]
def get_term(self) -> int
\end{lstlisting}

\textbf{可用性与安全性}

\textbf{\passthrough{\lstinline!SIGNATURE\_ONLY!} /
\passthrough{\lstinline!UNCLASSIFIED!}}：当前只有设计期签名，不能假设已存在于运行时扩展。

\textbf{参数}

无显式参数。实例方法中的 \passthrough{\lstinline!self!} 由 Python
自动传入。

\textbf{返回值}

\begin{longtable}[]{@{}
  >{\raggedright\arraybackslash}p{(\columnwidth - 4\tabcolsep) * \real{0.18}}
  >{\raggedright\arraybackslash}p{(\columnwidth - 4\tabcolsep) * \real{0.20}}
  >{\raggedright\arraybackslash}p{(\columnwidth - 4\tabcolsep) * \real{0.62}}@{}}
\toprule
位置 & 类型 & 含义 \\
\midrule
\endhead
返回值 & \passthrough{\lstinline!int!} & SDK 状态码；Unitree 示例通常以
\passthrough{\lstinline!0!} 表示成功，非零值需按具体服务定义解释。 \\
\bottomrule
\end{longtable}

\textbf{对应 C++}

\begin{itemize}
\tightlist
\item
  类：\passthrough{\lstinline!unitree::robot::LeaseContext!}
\item
  签名：\passthrough{\lstinline!GetTerm() const!}
\item
  绑定策略：\passthrough{\lstinline!SIGNATURE\_PREVIEW!}
\item
  声明位置：\passthrough{\lstinline!include/unitree/robot/client/lease\_client.hpp:23!}
\end{itemize}

\textbf{说明}

\begin{itemize}
\tightlist
\item
  示例是局部调用片段；\passthrough{\lstinline!obj!}
  和参数变量表示已经按上层业务校验并准备好的对象。
\item
  该条目可以被编辑器和类型检查器识别，但当前运行时可能在属性查找阶段直接失败。
\end{itemize}

\textbf{用法}

\begin{lstlisting}[language=Python]
from typing import TYPE_CHECKING

if TYPE_CHECKING:
    def planned_get_term(obj: LeaseContext) -> int:
        return obj.get_term()
\end{lstlisting}

\begin{quote}
{[}!CAUTION{]} 上面的 \passthrough{\lstinline!TYPE\_CHECKING!}
示例只表达计划调用形态。该分支在普通 Python
运行时为假，不代表当前扩展能执行此接口。
\end{quote}

\hypertarget{unitree-sdk2-cpp-robot-leaseserver}{%
\subsubsection{\texorpdfstring{\texttt{unitree\_sdk2\_cpp.robot.LeaseServer}}{unitree\_sdk2\_cpp.robot.LeaseServer}}\label{unitree-sdk2-cpp-robot-leaseserver}}

SDK 类型签名预览。请逐项查看构造函数和方法的可用性。

\textbf{导入}

\begin{lstlisting}[language=Python]
from unitree_sdk2_cpp.robot import LeaseServer
\end{lstlisting}

\textbf{基类}：\passthrough{\lstinline!ServerBase!}

\textbf{方法索引}

\begin{longtable}[]{@{}
  >{\raggedright\arraybackslash}p{(\columnwidth - 6\tabcolsep) * \real{0.18}}
  >{\raggedright\arraybackslash}p{(\columnwidth - 6\tabcolsep) * \real{0.24}}
  >{\raggedright\arraybackslash}p{(\columnwidth - 6\tabcolsep) * \real{0.18}}
  >{\raggedright\arraybackslash}p{(\columnwidth - 6\tabcolsep) * \real{0.40}}@{}}
\toprule
方法 & Python 签名 & 状态 & 安全分类 \\
\midrule
\endhead
\protect\hyperlink{unitree-sdk2-cpp-robot-leaseserver-dunder-init-1}{\passthrough{\lstinline!\_\_init\_\_!}}
&
\passthrough{\lstinline!def \_\_init\_\_(self, name: str, term: int) -> None!}
& \passthrough{\lstinline!SIGNATURE\_ONLY!} &
\passthrough{\lstinline!CONSTRUCTION!} \\
\protect\hyperlink{unitree-sdk2-cpp-robot-leaseserver-init-1}{\passthrough{\lstinline!init!}}
& \passthrough{\lstinline!def init(self) -> None!} &
\passthrough{\lstinline!SIGNATURE\_ONLY!} &
\passthrough{\lstinline!UNCLASSIFIED!} \\
\protect\hyperlink{unitree-sdk2-cpp-robot-leaseserver-check-request-lease-denied-1}{\passthrough{\lstinline!check\_request\_lease\_denied!}}
&
\passthrough{\lstinline!def check\_request\_lease\_denied(self, lease\_id: int) -> bool!}
& \passthrough{\lstinline!SIGNATURE\_ONLY!} &
\passthrough{\lstinline!UNCLASSIFIED!} \\
\bottomrule
\end{longtable}

\hypertarget{unitree-sdk2-cpp-robot-leaseserver-dunder-init-1}{%
\paragraph{\texorpdfstring{\texttt{unitree\_sdk2\_cpp.robot.LeaseServer.\_\_init\_\_}}{unitree\_sdk2\_cpp.robot.LeaseServer.\_\_init\_\_}}\label{unitree-sdk2-cpp-robot-leaseserver-dunder-init-1}}

初始化 \passthrough{\lstinline!LeaseServer!}
实例。是否能实际构造取决于下方可用性状态。

\textbf{签名}

\begin{lstlisting}[language=Python]
def __init__(self, name: str, term: int) -> None
\end{lstlisting}

\textbf{可用性与安全性}

\textbf{\passthrough{\lstinline!SIGNATURE\_ONLY!} /
\passthrough{\lstinline!CONSTRUCTION!}}：当前只有设计期签名，不能假设已存在于运行时扩展。

\textbf{参数}

\begin{longtable}[]{@{}
  >{\raggedright\arraybackslash}p{(\columnwidth - 6\tabcolsep) * \real{0.18}}
  >{\raggedright\arraybackslash}p{(\columnwidth - 6\tabcolsep) * \real{0.24}}
  >{\raggedright\arraybackslash}p{(\columnwidth - 6\tabcolsep) * \real{0.18}}
  >{\raggedright\arraybackslash}p{(\columnwidth - 6\tabcolsep) * \real{0.40}}@{}}
\toprule
名称 & Python 类型 & 默认值 & 含义 \\
\midrule
\endhead
\passthrough{\lstinline!name!} & \passthrough{\lstinline!str!} & 必填 &
SDK 对象、配置项、服务或动作的名称。允许值由调用它的具体接口定义。 对应
C++ 参数 \passthrough{\lstinline!name: const std::string \&!}。
底层为字符串；长度、编码和允许值由具体协议决定。 \\
\passthrough{\lstinline!term!} & \passthrough{\lstinline!int!} & 必填 &
传给该接口的 \passthrough{\lstinline!term!} 参数，Python 类型为
\passthrough{\lstinline!int!}。精确含义、单位、范围和枚举值需查目标型号对应头文件与协议。
对应 C++ 参数 \passthrough{\lstinline!term: int64\_t!}。 取值范围为
-9223372036854775808 到 9223372036854775807。 \\
\bottomrule
\end{longtable}

\textbf{返回值}

\begin{longtable}[]{@{}
  >{\raggedright\arraybackslash}p{(\columnwidth - 4\tabcolsep) * \real{0.18}}
  >{\raggedright\arraybackslash}p{(\columnwidth - 4\tabcolsep) * \real{0.20}}
  >{\raggedright\arraybackslash}p{(\columnwidth - 4\tabcolsep) * \real{0.62}}@{}}
\toprule
位置 & 类型 & 含义 \\
\midrule
\endhead
无 & \passthrough{\lstinline!None!} &
\passthrough{\lstinline!\_\_init\_\_!}
本身不返回值；调用类对象时，在构造函数可用的前提下得到该类实例。 \\
\bottomrule
\end{longtable}

\textbf{对应 C++}

\begin{itemize}
\tightlist
\item
  类：\passthrough{\lstinline!unitree::robot::LeaseServer!}
\item
  签名：\passthrough{\lstinline!LeaseServer(const std::string \&, int64\_t)!}
\item
  绑定策略：\passthrough{\lstinline!未单独标注!}
\item
  声明位置：\passthrough{\lstinline!include/unitree/robot/server/lease\_server.hpp:33!}
\end{itemize}

\textbf{说明}

\begin{itemize}
\tightlist
\item
  示例是局部调用片段；\passthrough{\lstinline!obj!}
  和参数变量表示已经按上层业务校验并准备好的对象。
\item
  该条目可以被编辑器和类型检查器识别，但当前运行时可能在属性查找阶段直接失败。
\end{itemize}

\textbf{用法}

\begin{lstlisting}[language=Python]
from typing import TYPE_CHECKING

if TYPE_CHECKING:
    def planned_construct(name: str, term: int) -> LeaseServer:
        return LeaseServer(name=name, term=term)
\end{lstlisting}

\begin{quote}
{[}!CAUTION{]} 上面的 \passthrough{\lstinline!TYPE\_CHECKING!}
示例只表达计划调用形态。该分支在普通 Python
运行时为假，不代表当前扩展能执行此接口。
\end{quote}

\hypertarget{unitree-sdk2-cpp-robot-leaseserver-init-1}{%
\paragraph{\texorpdfstring{\texttt{unitree\_sdk2\_cpp.robot.LeaseServer.init}}{unitree\_sdk2\_cpp.robot.LeaseServer.init}}\label{unitree-sdk2-cpp-robot-leaseserver-init-1}}

初始化当前 SDK 对象所需的底层通道或服务资源。

\textbf{签名}

\begin{lstlisting}[language=Python]
def init(self) -> None
\end{lstlisting}

\textbf{可用性与安全性}

\textbf{\passthrough{\lstinline!SIGNATURE\_ONLY!} /
\passthrough{\lstinline!UNCLASSIFIED!}}：当前只有设计期签名，不能假设已存在于运行时扩展。

\textbf{参数}

无显式参数。实例方法中的 \passthrough{\lstinline!self!} 由 Python
自动传入。

\textbf{返回值}

\begin{longtable}[]{@{}
  >{\raggedright\arraybackslash}p{(\columnwidth - 4\tabcolsep) * \real{0.18}}
  >{\raggedright\arraybackslash}p{(\columnwidth - 4\tabcolsep) * \real{0.20}}
  >{\raggedright\arraybackslash}p{(\columnwidth - 4\tabcolsep) * \real{0.62}}@{}}
\toprule
位置 & 类型 & 含义 \\
\midrule
\endhead
无 & \passthrough{\lstinline!None!} & 该方法不返回 Python
值；失败可能表现为异常或底层状态变化。 \\
\bottomrule
\end{longtable}

\textbf{对应 C++}

\begin{itemize}
\tightlist
\item
  类：\passthrough{\lstinline!unitree::robot::LeaseServer!}
\item
  签名：\passthrough{\lstinline!Init()!}
\item
  绑定策略：\passthrough{\lstinline!SIGNATURE\_PREVIEW!}
\item
  声明位置：\passthrough{\lstinline!include/unitree/robot/server/lease\_server.hpp:36!}
\end{itemize}

\textbf{说明}

\begin{itemize}
\tightlist
\item
  示例是局部调用片段；\passthrough{\lstinline!obj!}
  和参数变量表示已经按上层业务校验并准备好的对象。
\item
  该条目可以被编辑器和类型检查器识别，但当前运行时可能在属性查找阶段直接失败。
\end{itemize}

\textbf{用法}

\begin{lstlisting}[language=Python]
from typing import TYPE_CHECKING

if TYPE_CHECKING:
    def planned_init(obj: LeaseServer) -> None:
        obj.init()
\end{lstlisting}

\begin{quote}
{[}!CAUTION{]} 上面的 \passthrough{\lstinline!TYPE\_CHECKING!}
示例只表达计划调用形态。该分支在普通 Python
运行时为假，不代表当前扩展能执行此接口。
\end{quote}

\hypertarget{unitree-sdk2-cpp-robot-leaseserver-check-request-lease-denied-1}{%
\paragraph{\texorpdfstring{\texttt{unitree\_sdk2\_cpp.robot.LeaseServer.check\_request\_lease\_denied}}{unitree\_sdk2\_cpp.robot.LeaseServer.check\_request\_lease\_denied}}\label{unitree-sdk2-cpp-robot-leaseserver-check-request-lease-denied-1}}

查询或检查 请求 lease \passthrough{\lstinline!denied!}。

\textbf{签名}

\begin{lstlisting}[language=Python]
def check_request_lease_denied(self, lease_id: int) -> bool
\end{lstlisting}

\textbf{可用性与安全性}

\textbf{\passthrough{\lstinline!SIGNATURE\_ONLY!} /
\passthrough{\lstinline!UNCLASSIFIED!}}：当前只有设计期签名，不能假设已存在于运行时扩展。

\textbf{参数}

\begin{longtable}[]{@{}
  >{\raggedright\arraybackslash}p{(\columnwidth - 6\tabcolsep) * \real{0.18}}
  >{\raggedright\arraybackslash}p{(\columnwidth - 6\tabcolsep) * \real{0.24}}
  >{\raggedright\arraybackslash}p{(\columnwidth - 6\tabcolsep) * \real{0.18}}
  >{\raggedright\arraybackslash}p{(\columnwidth - 6\tabcolsep) * \real{0.40}}@{}}
\toprule
名称 & Python 类型 & 默认值 & 含义 \\
\midrule
\endhead
\passthrough{\lstinline!lease\_id!} & \passthrough{\lstinline!int!} &
必填 & 传给该接口的 lease ID 参数，Python 类型为
\passthrough{\lstinline!int!}。精确含义、单位、范围和枚举值需查目标型号对应头文件与协议。
对应 C++ 参数 \passthrough{\lstinline!leaseId: int64\_t!}。 取值范围为
-9223372036854775808 到 9223372036854775807。 \\
\bottomrule
\end{longtable}

\textbf{返回值}

\begin{longtable}[]{@{}
  >{\raggedright\arraybackslash}p{(\columnwidth - 4\tabcolsep) * \real{0.18}}
  >{\raggedright\arraybackslash}p{(\columnwidth - 4\tabcolsep) * \real{0.20}}
  >{\raggedright\arraybackslash}p{(\columnwidth - 4\tabcolsep) * \real{0.62}}@{}}
\toprule
位置 & 类型 & 含义 \\
\midrule
\endhead
返回值 & \passthrough{\lstinline!bool!} & 返回
\passthrough{\lstinline!bool!}。更精确的业务含义见该方法用途、C++
签名和目标型号协议。 \\
\bottomrule
\end{longtable}

\textbf{对应 C++}

\begin{itemize}
\tightlist
\item
  类：\passthrough{\lstinline!unitree::robot::LeaseServer!}
\item
  签名：\passthrough{\lstinline!CheckRequestLeaseDenied(int64\_t)!}
\item
  绑定策略：\passthrough{\lstinline!SIGNATURE\_PREVIEW!}
\item
  声明位置：\passthrough{\lstinline!include/unitree/robot/server/lease\_server.hpp:38!}
\end{itemize}

\textbf{说明}

\begin{itemize}
\tightlist
\item
  示例是局部调用片段；\passthrough{\lstinline!obj!}
  和参数变量表示已经按上层业务校验并准备好的对象。
\item
  该条目可以被编辑器和类型检查器识别，但当前运行时可能在属性查找阶段直接失败。
\end{itemize}

\textbf{用法}

\begin{lstlisting}[language=Python]
from typing import TYPE_CHECKING

if TYPE_CHECKING:
    def planned_check_request_lease_denied(obj: LeaseServer, lease_id: int) -> bool:
        return obj.check_request_lease_denied(lease_id=lease_id)
\end{lstlisting}

\begin{quote}
{[}!CAUTION{]} 上面的 \passthrough{\lstinline!TYPE\_CHECKING!}
示例只表达计划调用形态。该分支在普通 Python
运行时为假，不代表当前扩展能执行此接口。
\end{quote}

\hypertarget{unitree-sdk2-cpp-robot-requestfuture}{%
\subsubsection{\texorpdfstring{\texttt{unitree\_sdk2\_cpp.robot.RequestFuture}}{unitree\_sdk2\_cpp.robot.RequestFuture}}\label{unitree-sdk2-cpp-robot-requestfuture}}

SDK 类型签名预览。请逐项查看构造函数和方法的可用性。

\textbf{导入}

\begin{lstlisting}[language=Python]
from unitree_sdk2_cpp.robot import RequestFuture
\end{lstlisting}

\textbf{基类}：\passthrough{\lstinline!object!}

\textbf{方法索引}

\begin{longtable}[]{@{}
  >{\raggedright\arraybackslash}p{(\columnwidth - 6\tabcolsep) * \real{0.18}}
  >{\raggedright\arraybackslash}p{(\columnwidth - 6\tabcolsep) * \real{0.24}}
  >{\raggedright\arraybackslash}p{(\columnwidth - 6\tabcolsep) * \real{0.18}}
  >{\raggedright\arraybackslash}p{(\columnwidth - 6\tabcolsep) * \real{0.40}}@{}}
\toprule
方法 & Python 签名 & 状态 & 安全分类 \\
\midrule
\endhead
\protect\hyperlink{unitree-sdk2-cpp-robot-requestfuture-dunder-init-1}{\passthrough{\lstinline!\_\_init\_\_（重载 1/2）!}}
& \passthrough{\lstinline!def \_\_init\_\_(self) -> None!} &
\passthrough{\lstinline!SIGNATURE\_ONLY!} &
\passthrough{\lstinline!CONSTRUCTION!} \\
\protect\hyperlink{unitree-sdk2-cpp-robot-requestfuture-dunder-init-2}{\passthrough{\lstinline!\_\_init\_\_（重载 2/2）!}}
&
\passthrough{\lstinline!def \_\_init\_\_(self, request\_id: int) -> None!}
& \passthrough{\lstinline!SIGNATURE\_ONLY!} &
\passthrough{\lstinline!CONSTRUCTION!} \\
\protect\hyperlink{unitree-sdk2-cpp-robot-requestfuture-set-request-id-1}{\passthrough{\lstinline!set\_request\_id!}}
&
\passthrough{\lstinline!def set\_request\_id(self, request\_id: int) -> None!}
& \passthrough{\lstinline!SIGNATURE\_ONLY!} &
\passthrough{\lstinline!UNCLASSIFIED!} \\
\protect\hyperlink{unitree-sdk2-cpp-robot-requestfuture-get-request-id-1}{\passthrough{\lstinline!get\_request\_id!}}
& \passthrough{\lstinline!def get\_request\_id(self) -> int!} &
\passthrough{\lstinline!SIGNATURE\_ONLY!} &
\passthrough{\lstinline!UNCLASSIFIED!} \\
\protect\hyperlink{unitree-sdk2-cpp-robot-requestfuture-set-queue-1}{\passthrough{\lstinline!set\_queue!}}
&
\passthrough{\lstinline!def set\_queue(self, future\_queue\_ptr: RequestFutureQueue) -> bool!}
& \passthrough{\lstinline!SIGNATURE\_ONLY!} &
\passthrough{\lstinline!UNCLASSIFIED!} \\
\protect\hyperlink{unitree-sdk2-cpp-robot-requestfuture-get-response-1}{\passthrough{\lstinline!get\_response!}}
&
\passthrough{\lstinline!def get\_response(self, microsec: int) -> Any!}
& \passthrough{\lstinline!SIGNATURE\_ONLY!} &
\passthrough{\lstinline!UNCLASSIFIED!} \\
\protect\hyperlink{unitree-sdk2-cpp-robot-requestfuture-ready-1}{\passthrough{\lstinline!ready!}}
& \passthrough{\lstinline!def ready(self, request\_ptr: Any) -> None!} &
\passthrough{\lstinline!SIGNATURE\_ONLY!} &
\passthrough{\lstinline!UNCLASSIFIED!} \\
\bottomrule
\end{longtable}

\hypertarget{unitree-sdk2-cpp-robot-requestfuture-dunder-init-1}{%
\paragraph{\texorpdfstring{\texttt{unitree\_sdk2\_cpp.robot.RequestFuture.\_\_init\_\_}（重载
1/2）}{unitree\_sdk2\_cpp.robot.RequestFuture.\_\_init\_\_（重载 1/2）}}\label{unitree-sdk2-cpp-robot-requestfuture-dunder-init-1}}

初始化 \passthrough{\lstinline!RequestFuture!}
实例。是否能实际构造取决于下方可用性状态。

\textbf{签名}

\begin{lstlisting}[language=Python]
@overload
def __init__(self) -> None
\end{lstlisting}

\textbf{可用性与安全性}

\textbf{\passthrough{\lstinline!SIGNATURE\_ONLY!} /
\passthrough{\lstinline!CONSTRUCTION!}}：当前只有设计期签名，不能假设已存在于运行时扩展。

\textbf{参数}

无显式参数。实例方法中的 \passthrough{\lstinline!self!} 由 Python
自动传入。

\textbf{返回值}

\begin{longtable}[]{@{}
  >{\raggedright\arraybackslash}p{(\columnwidth - 4\tabcolsep) * \real{0.18}}
  >{\raggedright\arraybackslash}p{(\columnwidth - 4\tabcolsep) * \real{0.20}}
  >{\raggedright\arraybackslash}p{(\columnwidth - 4\tabcolsep) * \real{0.62}}@{}}
\toprule
位置 & 类型 & 含义 \\
\midrule
\endhead
无 & \passthrough{\lstinline!None!} &
\passthrough{\lstinline!\_\_init\_\_!}
本身不返回值；调用类对象时，在构造函数可用的前提下得到该类实例。 \\
\bottomrule
\end{longtable}

\textbf{对应 C++}

\begin{itemize}
\tightlist
\item
  类：\passthrough{\lstinline!unitree::robot::RequestFuture!}
\item
  签名：\passthrough{\lstinline!RequestFuture()!}
\item
  绑定策略：\passthrough{\lstinline!未单独标注!}
\item
  声明位置：\passthrough{\lstinline!include/unitree/robot/future/request\_future.hpp:27!}
\end{itemize}

\textbf{说明}

\begin{itemize}
\tightlist
\item
  示例是局部调用片段；\passthrough{\lstinline!obj!}
  和参数变量表示已经按上层业务校验并准备好的对象。
\item
  该条目可以被编辑器和类型检查器识别，但当前运行时可能在属性查找阶段直接失败。
\end{itemize}

\textbf{用法}

\begin{lstlisting}[language=Python]
from typing import TYPE_CHECKING

if TYPE_CHECKING:
    def planned_construct() -> RequestFuture:
        return RequestFuture()
\end{lstlisting}

\begin{quote}
{[}!CAUTION{]} 上面的 \passthrough{\lstinline!TYPE\_CHECKING!}
示例只表达计划调用形态。该分支在普通 Python
运行时为假，不代表当前扩展能执行此接口。
\end{quote}

\hypertarget{unitree-sdk2-cpp-robot-requestfuture-dunder-init-2}{%
\paragraph{\texorpdfstring{\texttt{unitree\_sdk2\_cpp.robot.RequestFuture.\_\_init\_\_}（重载
2/2）}{unitree\_sdk2\_cpp.robot.RequestFuture.\_\_init\_\_（重载 2/2）}}\label{unitree-sdk2-cpp-robot-requestfuture-dunder-init-2}}

初始化 \passthrough{\lstinline!RequestFuture!}
实例。是否能实际构造取决于下方可用性状态。

\textbf{签名}

\begin{lstlisting}[language=Python]
@overload
def __init__(self, request_id: int) -> None
\end{lstlisting}

\textbf{可用性与安全性}

\textbf{\passthrough{\lstinline!SIGNATURE\_ONLY!} /
\passthrough{\lstinline!CONSTRUCTION!}}：当前只有设计期签名，不能假设已存在于运行时扩展。

\textbf{参数}

\begin{longtable}[]{@{}
  >{\raggedright\arraybackslash}p{(\columnwidth - 6\tabcolsep) * \real{0.18}}
  >{\raggedright\arraybackslash}p{(\columnwidth - 6\tabcolsep) * \real{0.24}}
  >{\raggedright\arraybackslash}p{(\columnwidth - 6\tabcolsep) * \real{0.18}}
  >{\raggedright\arraybackslash}p{(\columnwidth - 6\tabcolsep) * \real{0.40}}@{}}
\toprule
名称 & Python 类型 & 默认值 & 含义 \\
\midrule
\endhead
\passthrough{\lstinline!request\_id!} & \passthrough{\lstinline!int!} &
必填 & 传给该接口的 请求 ID 参数，Python 类型为
\passthrough{\lstinline!int!}。精确含义、单位、范围和枚举值需查目标型号对应头文件与协议。
对应 C++ 参数 \passthrough{\lstinline!requestId: int64\_t!}。 取值范围为
-9223372036854775808 到 9223372036854775807。 \\
\bottomrule
\end{longtable}

\textbf{返回值}

\begin{longtable}[]{@{}
  >{\raggedright\arraybackslash}p{(\columnwidth - 4\tabcolsep) * \real{0.18}}
  >{\raggedright\arraybackslash}p{(\columnwidth - 4\tabcolsep) * \real{0.20}}
  >{\raggedright\arraybackslash}p{(\columnwidth - 4\tabcolsep) * \real{0.62}}@{}}
\toprule
位置 & 类型 & 含义 \\
\midrule
\endhead
无 & \passthrough{\lstinline!None!} &
\passthrough{\lstinline!\_\_init\_\_!}
本身不返回值；调用类对象时，在构造函数可用的前提下得到该类实例。 \\
\bottomrule
\end{longtable}

\textbf{对应 C++}

\begin{itemize}
\tightlist
\item
  类：\passthrough{\lstinline!unitree::robot::RequestFuture!}
\item
  签名：\passthrough{\lstinline!RequestFuture(int64\_t)!}
\item
  绑定策略：\passthrough{\lstinline!未单独标注!}
\item
  声明位置：\passthrough{\lstinline!include/unitree/robot/future/request\_future.hpp:28!}
\end{itemize}

\textbf{说明}

\begin{itemize}
\tightlist
\item
  示例是局部调用片段；\passthrough{\lstinline!obj!}
  和参数变量表示已经按上层业务校验并准备好的对象。
\item
  该条目可以被编辑器和类型检查器识别，但当前运行时可能在属性查找阶段直接失败。
\end{itemize}

\textbf{用法}

\begin{lstlisting}[language=Python]
from typing import TYPE_CHECKING

if TYPE_CHECKING:
    def planned_construct(request_id: int) -> RequestFuture:
        return RequestFuture(request_id=request_id)
\end{lstlisting}

\begin{quote}
{[}!CAUTION{]} 上面的 \passthrough{\lstinline!TYPE\_CHECKING!}
示例只表达计划调用形态。该分支在普通 Python
运行时为假，不代表当前扩展能执行此接口。
\end{quote}

\hypertarget{unitree-sdk2-cpp-robot-requestfuture-set-request-id-1}{%
\paragraph{\texorpdfstring{\texttt{unitree\_sdk2\_cpp.robot.RequestFuture.set\_request\_id}}{unitree\_sdk2\_cpp.robot.RequestFuture.set\_request\_id}}\label{unitree-sdk2-cpp-robot-requestfuture-set-request-id-1}}

设置 请求 ID。具体副作用和安全边界见下方状态。

\textbf{签名}

\begin{lstlisting}[language=Python]
def set_request_id(self, request_id: int) -> None
\end{lstlisting}

\textbf{可用性与安全性}

\textbf{\passthrough{\lstinline!SIGNATURE\_ONLY!} /
\passthrough{\lstinline!UNCLASSIFIED!}}：当前只有设计期签名，不能假设已存在于运行时扩展。

\textbf{参数}

\begin{longtable}[]{@{}
  >{\raggedright\arraybackslash}p{(\columnwidth - 6\tabcolsep) * \real{0.18}}
  >{\raggedright\arraybackslash}p{(\columnwidth - 6\tabcolsep) * \real{0.24}}
  >{\raggedright\arraybackslash}p{(\columnwidth - 6\tabcolsep) * \real{0.18}}
  >{\raggedright\arraybackslash}p{(\columnwidth - 6\tabcolsep) * \real{0.40}}@{}}
\toprule
名称 & Python 类型 & 默认值 & 含义 \\
\midrule
\endhead
\passthrough{\lstinline!request\_id!} & \passthrough{\lstinline!int!} &
必填 & 传给该接口的 请求 ID 参数，Python 类型为
\passthrough{\lstinline!int!}。精确含义、单位、范围和枚举值需查目标型号对应头文件与协议。
对应 C++ 参数 \passthrough{\lstinline!requestId: int64\_t!}。 取值范围为
-9223372036854775808 到 9223372036854775807。 \\
\bottomrule
\end{longtable}

\textbf{返回值}

\begin{longtable}[]{@{}
  >{\raggedright\arraybackslash}p{(\columnwidth - 4\tabcolsep) * \real{0.18}}
  >{\raggedright\arraybackslash}p{(\columnwidth - 4\tabcolsep) * \real{0.20}}
  >{\raggedright\arraybackslash}p{(\columnwidth - 4\tabcolsep) * \real{0.62}}@{}}
\toprule
位置 & 类型 & 含义 \\
\midrule
\endhead
无 & \passthrough{\lstinline!None!} & 该方法不返回 Python
值；失败可能表现为异常或底层状态变化。 \\
\bottomrule
\end{longtable}

\textbf{对应 C++}

\begin{itemize}
\tightlist
\item
  类：\passthrough{\lstinline!unitree::robot::RequestFuture!}
\item
  签名：\passthrough{\lstinline!SetRequestId(int64\_t)!}
\item
  绑定策略：\passthrough{\lstinline!SIGNATURE\_PREVIEW!}
\item
  声明位置：\passthrough{\lstinline!include/unitree/robot/future/request\_future.hpp:31!}
\end{itemize}

\textbf{说明}

\begin{itemize}
\tightlist
\item
  示例是局部调用片段；\passthrough{\lstinline!obj!}
  和参数变量表示已经按上层业务校验并准备好的对象。
\item
  该条目可以被编辑器和类型检查器识别，但当前运行时可能在属性查找阶段直接失败。
\end{itemize}

\textbf{用法}

\begin{lstlisting}[language=Python]
from typing import TYPE_CHECKING

if TYPE_CHECKING:
    def planned_set_request_id(obj: RequestFuture, request_id: int) -> None:
        obj.set_request_id(request_id=request_id)
\end{lstlisting}

\begin{quote}
{[}!CAUTION{]} 上面的 \passthrough{\lstinline!TYPE\_CHECKING!}
示例只表达计划调用形态。该分支在普通 Python
运行时为假，不代表当前扩展能执行此接口。
\end{quote}

\hypertarget{unitree-sdk2-cpp-robot-requestfuture-get-request-id-1}{%
\paragraph{\texorpdfstring{\texttt{unitree\_sdk2\_cpp.robot.RequestFuture.get\_request\_id}}{unitree\_sdk2\_cpp.robot.RequestFuture.get\_request\_id}}\label{unitree-sdk2-cpp-robot-requestfuture-get-request-id-1}}

查询或检查 请求 ID。

\textbf{签名}

\begin{lstlisting}[language=Python]
def get_request_id(self) -> int
\end{lstlisting}

\textbf{可用性与安全性}

\textbf{\passthrough{\lstinline!SIGNATURE\_ONLY!} /
\passthrough{\lstinline!UNCLASSIFIED!}}：当前只有设计期签名，不能假设已存在于运行时扩展。

\textbf{参数}

无显式参数。实例方法中的 \passthrough{\lstinline!self!} 由 Python
自动传入。

\textbf{返回值}

\begin{longtable}[]{@{}
  >{\raggedright\arraybackslash}p{(\columnwidth - 4\tabcolsep) * \real{0.18}}
  >{\raggedright\arraybackslash}p{(\columnwidth - 4\tabcolsep) * \real{0.20}}
  >{\raggedright\arraybackslash}p{(\columnwidth - 4\tabcolsep) * \real{0.62}}@{}}
\toprule
位置 & 类型 & 含义 \\
\midrule
\endhead
返回值 & \passthrough{\lstinline!int!} & SDK 状态码；Unitree 示例通常以
\passthrough{\lstinline!0!} 表示成功，非零值需按具体服务定义解释。 \\
\bottomrule
\end{longtable}

\textbf{对应 C++}

\begin{itemize}
\tightlist
\item
  类：\passthrough{\lstinline!unitree::robot::RequestFuture!}
\item
  签名：\passthrough{\lstinline!GetRequestId() const!}
\item
  绑定策略：\passthrough{\lstinline!SIGNATURE\_PREVIEW!}
\item
  声明位置：\passthrough{\lstinline!include/unitree/robot/future/request\_future.hpp:32!}
\end{itemize}

\textbf{说明}

\begin{itemize}
\tightlist
\item
  示例是局部调用片段；\passthrough{\lstinline!obj!}
  和参数变量表示已经按上层业务校验并准备好的对象。
\item
  该条目可以被编辑器和类型检查器识别，但当前运行时可能在属性查找阶段直接失败。
\end{itemize}

\textbf{用法}

\begin{lstlisting}[language=Python]
from typing import TYPE_CHECKING

if TYPE_CHECKING:
    def planned_get_request_id(obj: RequestFuture) -> int:
        return obj.get_request_id()
\end{lstlisting}

\begin{quote}
{[}!CAUTION{]} 上面的 \passthrough{\lstinline!TYPE\_CHECKING!}
示例只表达计划调用形态。该分支在普通 Python
运行时为假，不代表当前扩展能执行此接口。
\end{quote}

\hypertarget{unitree-sdk2-cpp-robot-requestfuture-set-queue-1}{%
\paragraph{\texorpdfstring{\texttt{unitree\_sdk2\_cpp.robot.RequestFuture.set\_queue}}{unitree\_sdk2\_cpp.robot.RequestFuture.set\_queue}}\label{unitree-sdk2-cpp-robot-requestfuture-set-queue-1}}

设置 \passthrough{\lstinline!queue!}。具体副作用和安全边界见下方状态。

\textbf{签名}

\begin{lstlisting}[language=Python]
def set_queue(self, future_queue_ptr: RequestFutureQueue) -> bool
\end{lstlisting}

\textbf{可用性与安全性}

\textbf{\passthrough{\lstinline!SIGNATURE\_ONLY!} /
\passthrough{\lstinline!UNCLASSIFIED!}}：当前只有设计期签名，不能假设已存在于运行时扩展。

\textbf{参数}

\begin{longtable}[]{@{}
  >{\raggedright\arraybackslash}p{(\columnwidth - 6\tabcolsep) * \real{0.18}}
  >{\raggedright\arraybackslash}p{(\columnwidth - 6\tabcolsep) * \real{0.24}}
  >{\raggedright\arraybackslash}p{(\columnwidth - 6\tabcolsep) * \real{0.18}}
  >{\raggedright\arraybackslash}p{(\columnwidth - 6\tabcolsep) * \real{0.40}}@{}}
\toprule
名称 & Python 类型 & 默认值 & 含义 \\
\midrule
\endhead
\passthrough{\lstinline!future\_queue\_ptr!} &
\passthrough{\lstinline!RequestFutureQueue!} & 必填 & 传给该接口的
\passthrough{\lstinline!future!} \passthrough{\lstinline!queue!}
\passthrough{\lstinline!ptr!} 参数，Python 类型为
\passthrough{\lstinline!RequestFutureQueue!}。精确含义、单位、范围和枚举值需查目标型号对应头文件与协议。
对应 C++ 参数
\passthrough{\lstinline!futureQueuePtr: const std::shared\_ptr<RequestFutureQueue> \&!}。 \\
\bottomrule
\end{longtable}

\textbf{返回值}

\begin{longtable}[]{@{}
  >{\raggedright\arraybackslash}p{(\columnwidth - 4\tabcolsep) * \real{0.18}}
  >{\raggedright\arraybackslash}p{(\columnwidth - 4\tabcolsep) * \real{0.20}}
  >{\raggedright\arraybackslash}p{(\columnwidth - 4\tabcolsep) * \real{0.62}}@{}}
\toprule
位置 & 类型 & 含义 \\
\midrule
\endhead
返回值 & \passthrough{\lstinline!bool!} & 返回
\passthrough{\lstinline!bool!}。更精确的业务含义见该方法用途、C++
签名和目标型号协议。 \\
\bottomrule
\end{longtable}

\textbf{对应 C++}

\begin{itemize}
\tightlist
\item
  类：\passthrough{\lstinline!unitree::robot::RequestFuture!}
\item
  签名：\passthrough{\lstinline!SetQueue(const std::shared\_ptr<RequestFutureQueue> \&)!}
\item
  绑定策略：\passthrough{\lstinline!SIGNATURE\_PREVIEW!}
\item
  声明位置：\passthrough{\lstinline!include/unitree/robot/future/request\_future.hpp:34!}
\end{itemize}

\textbf{说明}

\begin{itemize}
\tightlist
\item
  示例是局部调用片段；\passthrough{\lstinline!obj!}
  和参数变量表示已经按上层业务校验并准备好的对象。
\item
  该条目可以被编辑器和类型检查器识别，但当前运行时可能在属性查找阶段直接失败。
\end{itemize}

\textbf{用法}

\begin{lstlisting}[language=Python]
from typing import TYPE_CHECKING

if TYPE_CHECKING:
    def planned_set_queue(obj: RequestFuture, future_queue_ptr: RequestFutureQueue) -> bool:
        return obj.set_queue(future_queue_ptr=future_queue_ptr)
\end{lstlisting}

\begin{quote}
{[}!CAUTION{]} 上面的 \passthrough{\lstinline!TYPE\_CHECKING!}
示例只表达计划调用形态。该分支在普通 Python
运行时为假，不代表当前扩展能执行此接口。
\end{quote}

\hypertarget{unitree-sdk2-cpp-robot-requestfuture-get-response-1}{%
\paragraph{\texorpdfstring{\texttt{unitree\_sdk2\_cpp.robot.RequestFuture.get\_response}}{unitree\_sdk2\_cpp.robot.RequestFuture.get\_response}}\label{unitree-sdk2-cpp-robot-requestfuture-get-response-1}}

查询或检查 响应。

\textbf{签名}

\begin{lstlisting}[language=Python]
def get_response(self, microsec: int) -> Any
\end{lstlisting}

\textbf{可用性与安全性}

\textbf{\passthrough{\lstinline!SIGNATURE\_ONLY!} /
\passthrough{\lstinline!UNCLASSIFIED!}}：当前只有设计期签名，不能假设已存在于运行时扩展。

\textbf{参数}

\begin{longtable}[]{@{}
  >{\raggedright\arraybackslash}p{(\columnwidth - 6\tabcolsep) * \real{0.18}}
  >{\raggedright\arraybackslash}p{(\columnwidth - 6\tabcolsep) * \real{0.24}}
  >{\raggedright\arraybackslash}p{(\columnwidth - 6\tabcolsep) * \real{0.18}}
  >{\raggedright\arraybackslash}p{(\columnwidth - 6\tabcolsep) * \real{0.40}}@{}}
\toprule
名称 & Python 类型 & 默认值 & 含义 \\
\midrule
\endhead
\passthrough{\lstinline!microsec!} & \passthrough{\lstinline!int!} &
必填 & 等待时间，单位为微秒。具体超时结果以 SDK 方法约定为准。 对应 C++
参数 \passthrough{\lstinline!microsec: int64\_t!}。 取值范围为
-9223372036854775808 到 9223372036854775807。 \\
\bottomrule
\end{longtable}

\textbf{返回值}

\begin{longtable}[]{@{}
  >{\raggedright\arraybackslash}p{(\columnwidth - 4\tabcolsep) * \real{0.18}}
  >{\raggedright\arraybackslash}p{(\columnwidth - 4\tabcolsep) * \real{0.20}}
  >{\raggedright\arraybackslash}p{(\columnwidth - 4\tabcolsep) * \real{0.62}}@{}}
\toprule
位置 & 类型 & 含义 \\
\midrule
\endhead
返回值 & \passthrough{\lstinline!Any!} & 返回
\passthrough{\lstinline!Any!}。更精确的业务含义见该方法用途、C++
签名和目标型号协议。 \\
\bottomrule
\end{longtable}

\textbf{对应 C++}

\begin{itemize}
\tightlist
\item
  类：\passthrough{\lstinline!unitree::robot::RequestFuture!}
\item
  签名：\passthrough{\lstinline!GetResponse(int64\_t)!}
\item
  绑定策略：\passthrough{\lstinline!SIGNATURE\_PREVIEW!}
\item
  声明位置：\passthrough{\lstinline!include/unitree/robot/future/request\_future.hpp:36!}
\end{itemize}

\textbf{说明}

\begin{itemize}
\tightlist
\item
  示例是局部调用片段；\passthrough{\lstinline!obj!}
  和参数变量表示已经按上层业务校验并准备好的对象。
\item
  该条目可以被编辑器和类型检查器识别，但当前运行时可能在属性查找阶段直接失败。
\end{itemize}

\textbf{用法}

\begin{lstlisting}[language=Python]
from typing import TYPE_CHECKING

if TYPE_CHECKING:
    def planned_get_response(obj: RequestFuture, microsec: int) -> Any:
        return obj.get_response(microsec=microsec)
\end{lstlisting}

\begin{quote}
{[}!CAUTION{]} 上面的 \passthrough{\lstinline!TYPE\_CHECKING!}
示例只表达计划调用形态。该分支在普通 Python
运行时为假，不代表当前扩展能执行此接口。
\end{quote}

\hypertarget{unitree-sdk2-cpp-robot-requestfuture-ready-1}{%
\paragraph{\texorpdfstring{\texttt{unitree\_sdk2\_cpp.robot.RequestFuture.ready}}{unitree\_sdk2\_cpp.robot.RequestFuture.ready}}\label{unitree-sdk2-cpp-robot-requestfuture-ready-1}}

对应 C++ SDK 操作
\passthrough{\lstinline!Ready(const unitree::robot::ResponsePtr \&)!}。上游头文件没有可直接生成的业务说明时，本参考只保证签名映射，精确语义需查目标型号协议。

\textbf{签名}

\begin{lstlisting}[language=Python]
def ready(self, request_ptr: Any) -> None
\end{lstlisting}

\textbf{可用性与安全性}

\textbf{\passthrough{\lstinline!SIGNATURE\_ONLY!} /
\passthrough{\lstinline!UNCLASSIFIED!}}：当前只有设计期签名，不能假设已存在于运行时扩展。

\textbf{参数}

\begin{longtable}[]{@{}
  >{\raggedright\arraybackslash}p{(\columnwidth - 6\tabcolsep) * \real{0.18}}
  >{\raggedright\arraybackslash}p{(\columnwidth - 6\tabcolsep) * \real{0.24}}
  >{\raggedright\arraybackslash}p{(\columnwidth - 6\tabcolsep) * \real{0.18}}
  >{\raggedright\arraybackslash}p{(\columnwidth - 6\tabcolsep) * \real{0.40}}@{}}
\toprule
名称 & Python 类型 & 默认值 & 含义 \\
\midrule
\endhead
\passthrough{\lstinline!request\_ptr!} & \passthrough{\lstinline!Any!} &
必填 & 传给该接口的 请求 \passthrough{\lstinline!ptr!} 参数，Python
类型为
\passthrough{\lstinline!Any!}。精确含义、单位、范围和枚举值需查目标型号对应头文件与协议。
对应 C++ 参数
\passthrough{\lstinline!requestPtr: const unitree::robot::ResponsePtr \&!}。 \\
\bottomrule
\end{longtable}

\textbf{返回值}

\begin{longtable}[]{@{}
  >{\raggedright\arraybackslash}p{(\columnwidth - 4\tabcolsep) * \real{0.18}}
  >{\raggedright\arraybackslash}p{(\columnwidth - 4\tabcolsep) * \real{0.20}}
  >{\raggedright\arraybackslash}p{(\columnwidth - 4\tabcolsep) * \real{0.62}}@{}}
\toprule
位置 & 类型 & 含义 \\
\midrule
\endhead
无 & \passthrough{\lstinline!None!} & 该方法不返回 Python
值；失败可能表现为异常或底层状态变化。 \\
\bottomrule
\end{longtable}

\textbf{对应 C++}

\begin{itemize}
\tightlist
\item
  类：\passthrough{\lstinline!unitree::robot::RequestFuture!}
\item
  签名：\passthrough{\lstinline!Ready(const unitree::robot::ResponsePtr \&)!}
\item
  绑定策略：\passthrough{\lstinline!SIGNATURE\_PREVIEW!}
\item
  声明位置：\passthrough{\lstinline!include/unitree/robot/future/request\_future.hpp:39!}
\end{itemize}

\textbf{说明}

\begin{itemize}
\tightlist
\item
  示例是局部调用片段；\passthrough{\lstinline!obj!}
  和参数变量表示已经按上层业务校验并准备好的对象。
\item
  该条目可以被编辑器和类型检查器识别，但当前运行时可能在属性查找阶段直接失败。
\end{itemize}

\textbf{用法}

\begin{lstlisting}[language=Python]
from typing import TYPE_CHECKING

if TYPE_CHECKING:
    def planned_ready(obj: RequestFuture, request_ptr: Any) -> None:
        obj.ready(request_ptr=request_ptr)
\end{lstlisting}

\begin{quote}
{[}!CAUTION{]} 上面的 \passthrough{\lstinline!TYPE\_CHECKING!}
示例只表达计划调用形态。该分支在普通 Python
运行时为假，不代表当前扩展能执行此接口。
\end{quote}

\hypertarget{unitree-sdk2-cpp-robot-requestfuturequeue}{%
\subsubsection{\texorpdfstring{\texttt{unitree\_sdk2\_cpp.robot.RequestFutureQueue}}{unitree\_sdk2\_cpp.robot.RequestFutureQueue}}\label{unitree-sdk2-cpp-robot-requestfuturequeue}}

SDK 类型签名预览。请逐项查看构造函数和方法的可用性。

\textbf{导入}

\begin{lstlisting}[language=Python]
from unitree_sdk2_cpp.robot import RequestFutureQueue
\end{lstlisting}

\textbf{基类}：\passthrough{\lstinline!object!}

\textbf{方法索引}

\begin{longtable}[]{@{}
  >{\raggedright\arraybackslash}p{(\columnwidth - 6\tabcolsep) * \real{0.18}}
  >{\raggedright\arraybackslash}p{(\columnwidth - 6\tabcolsep) * \real{0.24}}
  >{\raggedright\arraybackslash}p{(\columnwidth - 6\tabcolsep) * \real{0.18}}
  >{\raggedright\arraybackslash}p{(\columnwidth - 6\tabcolsep) * \real{0.40}}@{}}
\toprule
方法 & Python 签名 & 状态 & 安全分类 \\
\midrule
\endhead
\protect\hyperlink{unitree-sdk2-cpp-robot-requestfuturequeue-dunder-init-1}{\passthrough{\lstinline!\_\_init\_\_!}}
& \passthrough{\lstinline!def \_\_init\_\_(self) -> None!} &
\passthrough{\lstinline!SIGNATURE\_ONLY!} &
\passthrough{\lstinline!CONSTRUCTION!} \\
\protect\hyperlink{unitree-sdk2-cpp-robot-requestfuturequeue-get-1}{\passthrough{\lstinline!get!}}
&
\passthrough{\lstinline!def get(self, request\_id: int) -> RequestFuture!}
& \passthrough{\lstinline!SIGNATURE\_ONLY!} &
\passthrough{\lstinline!UNCLASSIFIED!} \\
\protect\hyperlink{unitree-sdk2-cpp-robot-requestfuturequeue-put-1}{\passthrough{\lstinline!put!}}
&
\passthrough{\lstinline!def put(self, request\_id: int, future\_ptr: RequestFuture) -> bool!}
& \passthrough{\lstinline!SIGNATURE\_ONLY!} &
\passthrough{\lstinline!UNCLASSIFIED!} \\
\protect\hyperlink{unitree-sdk2-cpp-robot-requestfuturequeue-remove-1}{\passthrough{\lstinline!remove!}}
& \passthrough{\lstinline!def remove(self, request\_id: int) -> None!} &
\passthrough{\lstinline!SIGNATURE\_ONLY!} &
\passthrough{\lstinline!UNCLASSIFIED!} \\
\protect\hyperlink{unitree-sdk2-cpp-robot-requestfuturequeue-size-1}{\passthrough{\lstinline!size!}}
& \passthrough{\lstinline!def size(self) -> int!} &
\passthrough{\lstinline!SIGNATURE\_ONLY!} &
\passthrough{\lstinline!UNCLASSIFIED!} \\
\bottomrule
\end{longtable}

\hypertarget{unitree-sdk2-cpp-robot-requestfuturequeue-dunder-init-1}{%
\paragraph{\texorpdfstring{\texttt{unitree\_sdk2\_cpp.robot.RequestFutureQueue.\_\_init\_\_}}{unitree\_sdk2\_cpp.robot.RequestFutureQueue.\_\_init\_\_}}\label{unitree-sdk2-cpp-robot-requestfuturequeue-dunder-init-1}}

初始化 \passthrough{\lstinline!RequestFutureQueue!}
实例。是否能实际构造取决于下方可用性状态。

\textbf{签名}

\begin{lstlisting}[language=Python]
def __init__(self) -> None
\end{lstlisting}

\textbf{可用性与安全性}

\textbf{\passthrough{\lstinline!SIGNATURE\_ONLY!} /
\passthrough{\lstinline!CONSTRUCTION!}}：当前只有设计期签名，不能假设已存在于运行时扩展。

\textbf{参数}

无显式参数。实例方法中的 \passthrough{\lstinline!self!} 由 Python
自动传入。

\textbf{返回值}

\begin{longtable}[]{@{}
  >{\raggedright\arraybackslash}p{(\columnwidth - 4\tabcolsep) * \real{0.18}}
  >{\raggedright\arraybackslash}p{(\columnwidth - 4\tabcolsep) * \real{0.20}}
  >{\raggedright\arraybackslash}p{(\columnwidth - 4\tabcolsep) * \real{0.62}}@{}}
\toprule
位置 & 类型 & 含义 \\
\midrule
\endhead
无 & \passthrough{\lstinline!None!} &
\passthrough{\lstinline!\_\_init\_\_!}
本身不返回值；调用类对象时，在构造函数可用的前提下得到该类实例。 \\
\bottomrule
\end{longtable}

\textbf{对应 C++}

\begin{itemize}
\tightlist
\item
  类：\passthrough{\lstinline!unitree::robot::RequestFutureQueue!}
\item
  签名：\passthrough{\lstinline!RequestFutureQueue()!}
\item
  绑定策略：\passthrough{\lstinline!未单独标注!}
\item
  声明位置：\passthrough{\lstinline!include/unitree/robot/future/request\_future.hpp:65!}
\end{itemize}

\textbf{说明}

\begin{itemize}
\tightlist
\item
  示例是局部调用片段；\passthrough{\lstinline!obj!}
  和参数变量表示已经按上层业务校验并准备好的对象。
\item
  该条目可以被编辑器和类型检查器识别，但当前运行时可能在属性查找阶段直接失败。
\end{itemize}

\textbf{用法}

\begin{lstlisting}[language=Python]
from typing import TYPE_CHECKING

if TYPE_CHECKING:
    def planned_construct() -> RequestFutureQueue:
        return RequestFutureQueue()
\end{lstlisting}

\begin{quote}
{[}!CAUTION{]} 上面的 \passthrough{\lstinline!TYPE\_CHECKING!}
示例只表达计划调用形态。该分支在普通 Python
运行时为假，不代表当前扩展能执行此接口。
\end{quote}

\hypertarget{unitree-sdk2-cpp-robot-requestfuturequeue-get-1}{%
\paragraph{\texorpdfstring{\texttt{unitree\_sdk2\_cpp.robot.RequestFutureQueue.get}}{unitree\_sdk2\_cpp.robot.RequestFutureQueue.get}}\label{unitree-sdk2-cpp-robot-requestfuturequeue-get-1}}

对应 C++ SDK 操作
\passthrough{\lstinline!Get(int64\_t)!}。上游头文件没有可直接生成的业务说明时，本参考只保证签名映射，精确语义需查目标型号协议。

\textbf{签名}

\begin{lstlisting}[language=Python]
def get(self, request_id: int) -> RequestFuture
\end{lstlisting}

\textbf{可用性与安全性}

\textbf{\passthrough{\lstinline!SIGNATURE\_ONLY!} /
\passthrough{\lstinline!UNCLASSIFIED!}}：当前只有设计期签名，不能假设已存在于运行时扩展。

\textbf{参数}

\begin{longtable}[]{@{}
  >{\raggedright\arraybackslash}p{(\columnwidth - 6\tabcolsep) * \real{0.18}}
  >{\raggedright\arraybackslash}p{(\columnwidth - 6\tabcolsep) * \real{0.24}}
  >{\raggedright\arraybackslash}p{(\columnwidth - 6\tabcolsep) * \real{0.18}}
  >{\raggedright\arraybackslash}p{(\columnwidth - 6\tabcolsep) * \real{0.40}}@{}}
\toprule
名称 & Python 类型 & 默认值 & 含义 \\
\midrule
\endhead
\passthrough{\lstinline!request\_id!} & \passthrough{\lstinline!int!} &
必填 & 传给该接口的 请求 ID 参数，Python 类型为
\passthrough{\lstinline!int!}。精确含义、单位、范围和枚举值需查目标型号对应头文件与协议。
对应 C++ 参数 \passthrough{\lstinline!requestId: int64\_t!}。 取值范围为
-9223372036854775808 到 9223372036854775807。 \\
\bottomrule
\end{longtable}

\textbf{返回值}

\begin{longtable}[]{@{}
  >{\raggedright\arraybackslash}p{(\columnwidth - 4\tabcolsep) * \real{0.18}}
  >{\raggedright\arraybackslash}p{(\columnwidth - 4\tabcolsep) * \real{0.20}}
  >{\raggedright\arraybackslash}p{(\columnwidth - 4\tabcolsep) * \real{0.62}}@{}}
\toprule
位置 & 类型 & 含义 \\
\midrule
\endhead
返回值 & \passthrough{\lstinline!RequestFuture!} & 返回
\passthrough{\lstinline!RequestFuture!}。更精确的业务含义见该方法用途、C++
签名和目标型号协议。 \\
\bottomrule
\end{longtable}

\textbf{对应 C++}

\begin{itemize}
\tightlist
\item
  类：\passthrough{\lstinline!unitree::robot::RequestFutureQueue!}
\item
  签名：\passthrough{\lstinline!Get(int64\_t)!}
\item
  绑定策略：\passthrough{\lstinline!SIGNATURE\_PREVIEW!}
\item
  声明位置：\passthrough{\lstinline!include/unitree/robot/future/request\_future.hpp:68!}
\end{itemize}

\textbf{说明}

\begin{itemize}
\tightlist
\item
  示例是局部调用片段；\passthrough{\lstinline!obj!}
  和参数变量表示已经按上层业务校验并准备好的对象。
\item
  该条目可以被编辑器和类型检查器识别，但当前运行时可能在属性查找阶段直接失败。
\end{itemize}

\textbf{用法}

\begin{lstlisting}[language=Python]
from typing import TYPE_CHECKING

if TYPE_CHECKING:
    def planned_get(obj: RequestFutureQueue, request_id: int) -> RequestFuture:
        return obj.get(request_id=request_id)
\end{lstlisting}

\begin{quote}
{[}!CAUTION{]} 上面的 \passthrough{\lstinline!TYPE\_CHECKING!}
示例只表达计划调用形态。该分支在普通 Python
运行时为假，不代表当前扩展能执行此接口。
\end{quote}

\hypertarget{unitree-sdk2-cpp-robot-requestfuturequeue-put-1}{%
\paragraph{\texorpdfstring{\texttt{unitree\_sdk2\_cpp.robot.RequestFutureQueue.put}}{unitree\_sdk2\_cpp.robot.RequestFutureQueue.put}}\label{unitree-sdk2-cpp-robot-requestfuturequeue-put-1}}

对应 C++ SDK 操作
\passthrough{\lstinline!Put(int64\_t, const unitree::robot::RequestFuturePtr \&)!}。上游头文件没有可直接生成的业务说明时，本参考只保证签名映射，精确语义需查目标型号协议。

\textbf{签名}

\begin{lstlisting}[language=Python]
def put(self, request_id: int, future_ptr: RequestFuture) -> bool
\end{lstlisting}

\textbf{可用性与安全性}

\textbf{\passthrough{\lstinline!SIGNATURE\_ONLY!} /
\passthrough{\lstinline!UNCLASSIFIED!}}：当前只有设计期签名，不能假设已存在于运行时扩展。

\textbf{参数}

\begin{longtable}[]{@{}
  >{\raggedright\arraybackslash}p{(\columnwidth - 6\tabcolsep) * \real{0.18}}
  >{\raggedright\arraybackslash}p{(\columnwidth - 6\tabcolsep) * \real{0.24}}
  >{\raggedright\arraybackslash}p{(\columnwidth - 6\tabcolsep) * \real{0.18}}
  >{\raggedright\arraybackslash}p{(\columnwidth - 6\tabcolsep) * \real{0.40}}@{}}
\toprule
名称 & Python 类型 & 默认值 & 含义 \\
\midrule
\endhead
\passthrough{\lstinline!request\_id!} & \passthrough{\lstinline!int!} &
必填 & 传给该接口的 请求 ID 参数，Python 类型为
\passthrough{\lstinline!int!}。精确含义、单位、范围和枚举值需查目标型号对应头文件与协议。
对应 C++ 参数 \passthrough{\lstinline!requestId: int64\_t!}。 取值范围为
-9223372036854775808 到 9223372036854775807。 \\
\passthrough{\lstinline!future\_ptr!} &
\passthrough{\lstinline!RequestFuture!} & 必填 & 传给该接口的
\passthrough{\lstinline!future!} \passthrough{\lstinline!ptr!}
参数，Python 类型为
\passthrough{\lstinline!RequestFuture!}。精确含义、单位、范围和枚举值需查目标型号对应头文件与协议。
对应 C++ 参数
\passthrough{\lstinline!futurePtr: const unitree::robot::RequestFuturePtr \&!}。 \\
\bottomrule
\end{longtable}

\textbf{返回值}

\begin{longtable}[]{@{}
  >{\raggedright\arraybackslash}p{(\columnwidth - 4\tabcolsep) * \real{0.18}}
  >{\raggedright\arraybackslash}p{(\columnwidth - 4\tabcolsep) * \real{0.20}}
  >{\raggedright\arraybackslash}p{(\columnwidth - 4\tabcolsep) * \real{0.62}}@{}}
\toprule
位置 & 类型 & 含义 \\
\midrule
\endhead
返回值 & \passthrough{\lstinline!bool!} & 返回
\passthrough{\lstinline!bool!}。更精确的业务含义见该方法用途、C++
签名和目标型号协议。 \\
\bottomrule
\end{longtable}

\textbf{对应 C++}

\begin{itemize}
\tightlist
\item
  类：\passthrough{\lstinline!unitree::robot::RequestFutureQueue!}
\item
  签名：\passthrough{\lstinline!Put(int64\_t, const unitree::robot::RequestFuturePtr \&)!}
\item
  绑定策略：\passthrough{\lstinline!SIGNATURE\_PREVIEW!}
\item
  声明位置：\passthrough{\lstinline!include/unitree/robot/future/request\_future.hpp:69!}
\end{itemize}

\textbf{说明}

\begin{itemize}
\tightlist
\item
  示例是局部调用片段；\passthrough{\lstinline!obj!}
  和参数变量表示已经按上层业务校验并准备好的对象。
\item
  该条目可以被编辑器和类型检查器识别，但当前运行时可能在属性查找阶段直接失败。
\end{itemize}

\textbf{用法}

\begin{lstlisting}[language=Python]
from typing import TYPE_CHECKING

if TYPE_CHECKING:
    def planned_put(obj: RequestFutureQueue, request_id: int, future_ptr: RequestFuture) -> bool:
        return obj.put(request_id=request_id, future_ptr=future_ptr)
\end{lstlisting}

\begin{quote}
{[}!CAUTION{]} 上面的 \passthrough{\lstinline!TYPE\_CHECKING!}
示例只表达计划调用形态。该分支在普通 Python
运行时为假，不代表当前扩展能执行此接口。
\end{quote}

\hypertarget{unitree-sdk2-cpp-robot-requestfuturequeue-remove-1}{%
\paragraph{\texorpdfstring{\texttt{unitree\_sdk2\_cpp.robot.RequestFutureQueue.remove}}{unitree\_sdk2\_cpp.robot.RequestFutureQueue.remove}}\label{unitree-sdk2-cpp-robot-requestfuturequeue-remove-1}}

对应 C++ SDK 操作
\passthrough{\lstinline!Remove(int64\_t)!}。上游头文件没有可直接生成的业务说明时，本参考只保证签名映射，精确语义需查目标型号协议。

\textbf{签名}

\begin{lstlisting}[language=Python]
def remove(self, request_id: int) -> None
\end{lstlisting}

\textbf{可用性与安全性}

\textbf{\passthrough{\lstinline!SIGNATURE\_ONLY!} /
\passthrough{\lstinline!UNCLASSIFIED!}}：当前只有设计期签名，不能假设已存在于运行时扩展。

\textbf{参数}

\begin{longtable}[]{@{}
  >{\raggedright\arraybackslash}p{(\columnwidth - 6\tabcolsep) * \real{0.18}}
  >{\raggedright\arraybackslash}p{(\columnwidth - 6\tabcolsep) * \real{0.24}}
  >{\raggedright\arraybackslash}p{(\columnwidth - 6\tabcolsep) * \real{0.18}}
  >{\raggedright\arraybackslash}p{(\columnwidth - 6\tabcolsep) * \real{0.40}}@{}}
\toprule
名称 & Python 类型 & 默认值 & 含义 \\
\midrule
\endhead
\passthrough{\lstinline!request\_id!} & \passthrough{\lstinline!int!} &
必填 & 传给该接口的 请求 ID 参数，Python 类型为
\passthrough{\lstinline!int!}。精确含义、单位、范围和枚举值需查目标型号对应头文件与协议。
对应 C++ 参数 \passthrough{\lstinline!requestId: int64\_t!}。 取值范围为
-9223372036854775808 到 9223372036854775807。 \\
\bottomrule
\end{longtable}

\textbf{返回值}

\begin{longtable}[]{@{}
  >{\raggedright\arraybackslash}p{(\columnwidth - 4\tabcolsep) * \real{0.18}}
  >{\raggedright\arraybackslash}p{(\columnwidth - 4\tabcolsep) * \real{0.20}}
  >{\raggedright\arraybackslash}p{(\columnwidth - 4\tabcolsep) * \real{0.62}}@{}}
\toprule
位置 & 类型 & 含义 \\
\midrule
\endhead
无 & \passthrough{\lstinline!None!} & 该方法不返回 Python
值；失败可能表现为异常或底层状态变化。 \\
\bottomrule
\end{longtable}

\textbf{对应 C++}

\begin{itemize}
\tightlist
\item
  类：\passthrough{\lstinline!unitree::robot::RequestFutureQueue!}
\item
  签名：\passthrough{\lstinline!Remove(int64\_t)!}
\item
  绑定策略：\passthrough{\lstinline!SIGNATURE\_PREVIEW!}
\item
  声明位置：\passthrough{\lstinline!include/unitree/robot/future/request\_future.hpp:70!}
\end{itemize}

\textbf{说明}

\begin{itemize}
\tightlist
\item
  示例是局部调用片段；\passthrough{\lstinline!obj!}
  和参数变量表示已经按上层业务校验并准备好的对象。
\item
  该条目可以被编辑器和类型检查器识别，但当前运行时可能在属性查找阶段直接失败。
\end{itemize}

\textbf{用法}

\begin{lstlisting}[language=Python]
from typing import TYPE_CHECKING

if TYPE_CHECKING:
    def planned_remove(obj: RequestFutureQueue, request_id: int) -> None:
        obj.remove(request_id=request_id)
\end{lstlisting}

\begin{quote}
{[}!CAUTION{]} 上面的 \passthrough{\lstinline!TYPE\_CHECKING!}
示例只表达计划调用形态。该分支在普通 Python
运行时为假，不代表当前扩展能执行此接口。
\end{quote}

\hypertarget{unitree-sdk2-cpp-robot-requestfuturequeue-size-1}{%
\paragraph{\texorpdfstring{\texttt{unitree\_sdk2\_cpp.robot.RequestFutureQueue.size}}{unitree\_sdk2\_cpp.robot.RequestFutureQueue.size}}\label{unitree-sdk2-cpp-robot-requestfuturequeue-size-1}}

对应 C++ SDK 操作
\passthrough{\lstinline!Size()!}。上游头文件没有可直接生成的业务说明时，本参考只保证签名映射，精确语义需查目标型号协议。

\textbf{签名}

\begin{lstlisting}[language=Python]
def size(self) -> int
\end{lstlisting}

\textbf{可用性与安全性}

\textbf{\passthrough{\lstinline!SIGNATURE\_ONLY!} /
\passthrough{\lstinline!UNCLASSIFIED!}}：当前只有设计期签名，不能假设已存在于运行时扩展。

\textbf{参数}

无显式参数。实例方法中的 \passthrough{\lstinline!self!} 由 Python
自动传入。

\textbf{返回值}

\begin{longtable}[]{@{}
  >{\raggedright\arraybackslash}p{(\columnwidth - 4\tabcolsep) * \real{0.18}}
  >{\raggedright\arraybackslash}p{(\columnwidth - 4\tabcolsep) * \real{0.20}}
  >{\raggedright\arraybackslash}p{(\columnwidth - 4\tabcolsep) * \real{0.62}}@{}}
\toprule
位置 & 类型 & 含义 \\
\midrule
\endhead
返回值 & \passthrough{\lstinline!int!} & SDK 状态码；Unitree 示例通常以
\passthrough{\lstinline!0!} 表示成功，非零值需按具体服务定义解释。 \\
\bottomrule
\end{longtable}

\textbf{对应 C++}

\begin{itemize}
\tightlist
\item
  类：\passthrough{\lstinline!unitree::robot::RequestFutureQueue!}
\item
  签名：\passthrough{\lstinline!Size()!}
\item
  绑定策略：\passthrough{\lstinline!SIGNATURE\_PREVIEW!}
\item
  声明位置：\passthrough{\lstinline!include/unitree/robot/future/request\_future.hpp:72!}
\end{itemize}

\textbf{说明}

\begin{itemize}
\tightlist
\item
  示例是局部调用片段；\passthrough{\lstinline!obj!}
  和参数变量表示已经按上层业务校验并准备好的对象。
\item
  该条目可以被编辑器和类型检查器识别，但当前运行时可能在属性查找阶段直接失败。
\end{itemize}

\textbf{用法}

\begin{lstlisting}[language=Python]
from typing import TYPE_CHECKING

if TYPE_CHECKING:
    def planned_size(obj: RequestFutureQueue) -> int:
        return obj.size()
\end{lstlisting}

\begin{quote}
{[}!CAUTION{]} 上面的 \passthrough{\lstinline!TYPE\_CHECKING!}
示例只表达计划调用形态。该分支在普通 Python
运行时为假，不代表当前扩展能执行此接口。
\end{quote}

\hypertarget{unitree-sdk2-cpp-robot-server}{%
\subsubsection{\texorpdfstring{\texttt{unitree\_sdk2\_cpp.robot.Server}}{unitree\_sdk2\_cpp.robot.Server}}\label{unitree-sdk2-cpp-robot-server}}

SDK 类型签名预览。请逐项查看构造函数和方法的可用性。

\textbf{导入}

\begin{lstlisting}[language=Python]
from unitree_sdk2_cpp.robot import Server
\end{lstlisting}

\textbf{基类}：\passthrough{\lstinline!ServerBase!}

\textbf{方法索引}

\begin{longtable}[]{@{}
  >{\raggedright\arraybackslash}p{(\columnwidth - 6\tabcolsep) * \real{0.18}}
  >{\raggedright\arraybackslash}p{(\columnwidth - 6\tabcolsep) * \real{0.24}}
  >{\raggedright\arraybackslash}p{(\columnwidth - 6\tabcolsep) * \real{0.18}}
  >{\raggedright\arraybackslash}p{(\columnwidth - 6\tabcolsep) * \real{0.40}}@{}}
\toprule
方法 & Python 签名 & 状态 & 安全分类 \\
\midrule
\endhead
\protect\hyperlink{unitree-sdk2-cpp-robot-server-dunder-init-1}{\passthrough{\lstinline!\_\_init\_\_!}}
& \passthrough{\lstinline!def \_\_init\_\_(self, name: str) -> None!} &
\passthrough{\lstinline!SIGNATURE\_ONLY!} &
\passthrough{\lstinline!CONSTRUCTION!} \\
\protect\hyperlink{unitree-sdk2-cpp-robot-server-init-1}{\passthrough{\lstinline!init!}}
& \passthrough{\lstinline!def init(self) -> None!} &
\passthrough{\lstinline!SIGNATURE\_ONLY!} &
\passthrough{\lstinline!UNCLASSIFIED!} \\
\protect\hyperlink{unitree-sdk2-cpp-robot-server-start-lease-1}{\passthrough{\lstinline!start\_lease（重载 1/2）!}}
&
\passthrough{\lstinline!def start\_lease(self, lease\_term: int) -> None!}
& \passthrough{\lstinline!SIGNATURE\_ONLY!} &
\passthrough{\lstinline!UNCLASSIFIED!} \\
\protect\hyperlink{unitree-sdk2-cpp-robot-server-start-lease-2}{\passthrough{\lstinline!start\_lease（重载 2/2）!}}
&
\passthrough{\lstinline!def start\_lease(self, lease\_term: float) -> None!}
& \passthrough{\lstinline!SIGNATURE\_ONLY!} &
\passthrough{\lstinline!UNCLASSIFIED!} \\
\protect\hyperlink{unitree-sdk2-cpp-robot-server-get-name-1}{\passthrough{\lstinline!get\_name!}}
& \passthrough{\lstinline!def get\_name(self) -> str!} &
\passthrough{\lstinline!SIGNATURE\_ONLY!} &
\passthrough{\lstinline!UNCLASSIFIED!} \\
\protect\hyperlink{unitree-sdk2-cpp-robot-server-get-api-version-1}{\passthrough{\lstinline!get\_api\_version!}}
& \passthrough{\lstinline!def get\_api\_version(self) -> str!} &
\passthrough{\lstinline!SIGNATURE\_ONLY!} &
\passthrough{\lstinline!UNCLASSIFIED!} \\
\protect\hyperlink{unitree-sdk2-cpp-robot-server-get-current-api-id-1}{\passthrough{\lstinline!get\_current\_api\_id!}}
& \passthrough{\lstinline!def get\_current\_api\_id(self) -> int!} &
\passthrough{\lstinline!SIGNATURE\_ONLY!} &
\passthrough{\lstinline!UNCLASSIFIED!} \\
\bottomrule
\end{longtable}

\hypertarget{unitree-sdk2-cpp-robot-server-dunder-init-1}{%
\paragraph{\texorpdfstring{\texttt{unitree\_sdk2\_cpp.robot.Server.\_\_init\_\_}}{unitree\_sdk2\_cpp.robot.Server.\_\_init\_\_}}\label{unitree-sdk2-cpp-robot-server-dunder-init-1}}

初始化 \passthrough{\lstinline!Server!}
实例。是否能实际构造取决于下方可用性状态。

\textbf{签名}

\begin{lstlisting}[language=Python]
def __init__(self, name: str) -> None
\end{lstlisting}

\textbf{可用性与安全性}

\textbf{\passthrough{\lstinline!SIGNATURE\_ONLY!} /
\passthrough{\lstinline!CONSTRUCTION!}}：当前只有设计期签名，不能假设已存在于运行时扩展。

\textbf{参数}

\begin{longtable}[]{@{}
  >{\raggedright\arraybackslash}p{(\columnwidth - 6\tabcolsep) * \real{0.18}}
  >{\raggedright\arraybackslash}p{(\columnwidth - 6\tabcolsep) * \real{0.24}}
  >{\raggedright\arraybackslash}p{(\columnwidth - 6\tabcolsep) * \real{0.18}}
  >{\raggedright\arraybackslash}p{(\columnwidth - 6\tabcolsep) * \real{0.40}}@{}}
\toprule
名称 & Python 类型 & 默认值 & 含义 \\
\midrule
\endhead
\passthrough{\lstinline!name!} & \passthrough{\lstinline!str!} & 必填 &
SDK 对象、配置项、服务或动作的名称。允许值由调用它的具体接口定义。 对应
C++ 参数 \passthrough{\lstinline!name: const std::string \&!}。
底层为字符串；长度、编码和允许值由具体协议决定。 \\
\bottomrule
\end{longtable}

\textbf{返回值}

\begin{longtable}[]{@{}
  >{\raggedright\arraybackslash}p{(\columnwidth - 4\tabcolsep) * \real{0.18}}
  >{\raggedright\arraybackslash}p{(\columnwidth - 4\tabcolsep) * \real{0.20}}
  >{\raggedright\arraybackslash}p{(\columnwidth - 4\tabcolsep) * \real{0.62}}@{}}
\toprule
位置 & 类型 & 含义 \\
\midrule
\endhead
无 & \passthrough{\lstinline!None!} &
\passthrough{\lstinline!\_\_init\_\_!}
本身不返回值；调用类对象时，在构造函数可用的前提下得到该类实例。 \\
\bottomrule
\end{longtable}

\textbf{对应 C++}

\begin{itemize}
\tightlist
\item
  类：\passthrough{\lstinline!unitree::robot::Server!}
\item
  签名：\passthrough{\lstinline!Server(const std::string \&)!}
\item
  绑定策略：\passthrough{\lstinline!未单独标注!}
\item
  声明位置：\passthrough{\lstinline!include/unitree/robot/server/server.hpp:29!}
\end{itemize}

\textbf{说明}

\begin{itemize}
\tightlist
\item
  示例是局部调用片段；\passthrough{\lstinline!obj!}
  和参数变量表示已经按上层业务校验并准备好的对象。
\item
  该条目可以被编辑器和类型检查器识别，但当前运行时可能在属性查找阶段直接失败。
\end{itemize}

\textbf{用法}

\begin{lstlisting}[language=Python]
from typing import TYPE_CHECKING

if TYPE_CHECKING:
    def planned_construct(name: str) -> Server:
        return Server(name=name)
\end{lstlisting}

\begin{quote}
{[}!CAUTION{]} 上面的 \passthrough{\lstinline!TYPE\_CHECKING!}
示例只表达计划调用形态。该分支在普通 Python
运行时为假，不代表当前扩展能执行此接口。
\end{quote}

\hypertarget{unitree-sdk2-cpp-robot-server-init-1}{%
\paragraph{\texorpdfstring{\texttt{unitree\_sdk2\_cpp.robot.Server.init}}{unitree\_sdk2\_cpp.robot.Server.init}}\label{unitree-sdk2-cpp-robot-server-init-1}}

初始化当前 SDK 对象所需的底层通道或服务资源。

\textbf{签名}

\begin{lstlisting}[language=Python]
def init(self) -> None
\end{lstlisting}

\textbf{可用性与安全性}

\textbf{\passthrough{\lstinline!SIGNATURE\_ONLY!} /
\passthrough{\lstinline!UNCLASSIFIED!}}：当前只有设计期签名，不能假设已存在于运行时扩展。

\textbf{参数}

无显式参数。实例方法中的 \passthrough{\lstinline!self!} 由 Python
自动传入。

\textbf{返回值}

\begin{longtable}[]{@{}
  >{\raggedright\arraybackslash}p{(\columnwidth - 4\tabcolsep) * \real{0.18}}
  >{\raggedright\arraybackslash}p{(\columnwidth - 4\tabcolsep) * \real{0.20}}
  >{\raggedright\arraybackslash}p{(\columnwidth - 4\tabcolsep) * \real{0.62}}@{}}
\toprule
位置 & 类型 & 含义 \\
\midrule
\endhead
无 & \passthrough{\lstinline!None!} & 该方法不返回 Python
值；失败可能表现为异常或底层状态变化。 \\
\bottomrule
\end{longtable}

\textbf{对应 C++}

\begin{itemize}
\tightlist
\item
  类：\passthrough{\lstinline!unitree::robot::Server!}
\item
  签名：\passthrough{\lstinline!Init()!}
\item
  绑定策略：\passthrough{\lstinline!SIGNATURE\_PREVIEW!}
\item
  声明位置：\passthrough{\lstinline!include/unitree/robot/server/server.hpp:32!}
\end{itemize}

\textbf{说明}

\begin{itemize}
\tightlist
\item
  示例是局部调用片段；\passthrough{\lstinline!obj!}
  和参数变量表示已经按上层业务校验并准备好的对象。
\item
  该条目可以被编辑器和类型检查器识别，但当前运行时可能在属性查找阶段直接失败。
\end{itemize}

\textbf{用法}

\begin{lstlisting}[language=Python]
from typing import TYPE_CHECKING

if TYPE_CHECKING:
    def planned_init(obj: Server) -> None:
        obj.init()
\end{lstlisting}

\begin{quote}
{[}!CAUTION{]} 上面的 \passthrough{\lstinline!TYPE\_CHECKING!}
示例只表达计划调用形态。该分支在普通 Python
运行时为假，不代表当前扩展能执行此接口。
\end{quote}

\hypertarget{unitree-sdk2-cpp-robot-server-start-lease-1}{%
\paragraph{\texorpdfstring{\texttt{unitree\_sdk2\_cpp.robot.Server.start\_lease}（重载
1/2）}{unitree\_sdk2\_cpp.robot.Server.start\_lease（重载 1/2）}}\label{unitree-sdk2-cpp-robot-server-start-lease-1}}

启动对应 SDK 操作。具体副作用和安全边界见下方状态。

\textbf{签名}

\begin{lstlisting}[language=Python]
@overload
def start_lease(self, lease_term: int) -> None
\end{lstlisting}

\textbf{可用性与安全性}

\textbf{\passthrough{\lstinline!SIGNATURE\_ONLY!} /
\passthrough{\lstinline!UNCLASSIFIED!}}：当前只有设计期签名，不能假设已存在于运行时扩展。

\textbf{参数}

\begin{longtable}[]{@{}
  >{\raggedright\arraybackslash}p{(\columnwidth - 6\tabcolsep) * \real{0.18}}
  >{\raggedright\arraybackslash}p{(\columnwidth - 6\tabcolsep) * \real{0.24}}
  >{\raggedright\arraybackslash}p{(\columnwidth - 6\tabcolsep) * \real{0.18}}
  >{\raggedright\arraybackslash}p{(\columnwidth - 6\tabcolsep) * \real{0.40}}@{}}
\toprule
名称 & Python 类型 & 默认值 & 含义 \\
\midrule
\endhead
\passthrough{\lstinline!lease\_term!} & \passthrough{\lstinline!int!} &
必填 & 传给该接口的 lease \passthrough{\lstinline!term!} 参数，Python
类型为
\passthrough{\lstinline!int!}。精确含义、单位、范围和枚举值需查目标型号对应头文件与协议。
对应 C++ 参数 \passthrough{\lstinline!leaseTerm: int64\_t!}。 取值范围为
-9223372036854775808 到 9223372036854775807。 \\
\bottomrule
\end{longtable}

\textbf{返回值}

\begin{longtable}[]{@{}
  >{\raggedright\arraybackslash}p{(\columnwidth - 4\tabcolsep) * \real{0.18}}
  >{\raggedright\arraybackslash}p{(\columnwidth - 4\tabcolsep) * \real{0.20}}
  >{\raggedright\arraybackslash}p{(\columnwidth - 4\tabcolsep) * \real{0.62}}@{}}
\toprule
位置 & 类型 & 含义 \\
\midrule
\endhead
无 & \passthrough{\lstinline!None!} & 该方法不返回 Python
值；失败可能表现为异常或底层状态变化。 \\
\bottomrule
\end{longtable}

\textbf{对应 C++}

\begin{itemize}
\tightlist
\item
  类：\passthrough{\lstinline!unitree::robot::Server!}
\item
  签名：\passthrough{\lstinline!StartLease(int64\_t)!}
\item
  绑定策略：\passthrough{\lstinline!SIGNATURE\_PREVIEW!}
\item
  声明位置：\passthrough{\lstinline!include/unitree/robot/server/server.hpp:34!}
\end{itemize}

\textbf{说明}

\begin{itemize}
\tightlist
\item
  示例是局部调用片段；\passthrough{\lstinline!obj!}
  和参数变量表示已经按上层业务校验并准备好的对象。
\item
  该条目可以被编辑器和类型检查器识别，但当前运行时可能在属性查找阶段直接失败。
\end{itemize}

\textbf{用法}

\begin{lstlisting}[language=Python]
from typing import TYPE_CHECKING

if TYPE_CHECKING:
    def planned_start_lease(obj: Server, lease_term: int) -> None:
        obj.start_lease(lease_term=lease_term)
\end{lstlisting}

\begin{quote}
{[}!CAUTION{]} 上面的 \passthrough{\lstinline!TYPE\_CHECKING!}
示例只表达计划调用形态。该分支在普通 Python
运行时为假，不代表当前扩展能执行此接口。
\end{quote}

\hypertarget{unitree-sdk2-cpp-robot-server-start-lease-2}{%
\paragraph{\texorpdfstring{\texttt{unitree\_sdk2\_cpp.robot.Server.start\_lease}（重载
2/2）}{unitree\_sdk2\_cpp.robot.Server.start\_lease（重载 2/2）}}\label{unitree-sdk2-cpp-robot-server-start-lease-2}}

启动对应 SDK 操作。具体副作用和安全边界见下方状态。

\textbf{签名}

\begin{lstlisting}[language=Python]
@overload
def start_lease(self, lease_term: float) -> None
\end{lstlisting}

\textbf{可用性与安全性}

\textbf{\passthrough{\lstinline!SIGNATURE\_ONLY!} /
\passthrough{\lstinline!UNCLASSIFIED!}}：当前只有设计期签名，不能假设已存在于运行时扩展。

\textbf{参数}

\begin{longtable}[]{@{}
  >{\raggedright\arraybackslash}p{(\columnwidth - 6\tabcolsep) * \real{0.18}}
  >{\raggedright\arraybackslash}p{(\columnwidth - 6\tabcolsep) * \real{0.24}}
  >{\raggedright\arraybackslash}p{(\columnwidth - 6\tabcolsep) * \real{0.18}}
  >{\raggedright\arraybackslash}p{(\columnwidth - 6\tabcolsep) * \real{0.40}}@{}}
\toprule
名称 & Python 类型 & 默认值 & 含义 \\
\midrule
\endhead
\passthrough{\lstinline!lease\_term!} & \passthrough{\lstinline!float!}
& 必填 & 传给该接口的 lease \passthrough{\lstinline!term!} 参数，Python
类型为
\passthrough{\lstinline!float!}。精确含义、单位、范围和枚举值需查目标型号对应头文件与协议。
对应 C++ 参数 \passthrough{\lstinline!leaseTerm: float!}。
底层为浮点数；类型本身不说明单位、坐标系或业务安全范围。 \\
\bottomrule
\end{longtable}

\textbf{返回值}

\begin{longtable}[]{@{}
  >{\raggedright\arraybackslash}p{(\columnwidth - 4\tabcolsep) * \real{0.18}}
  >{\raggedright\arraybackslash}p{(\columnwidth - 4\tabcolsep) * \real{0.20}}
  >{\raggedright\arraybackslash}p{(\columnwidth - 4\tabcolsep) * \real{0.62}}@{}}
\toprule
位置 & 类型 & 含义 \\
\midrule
\endhead
无 & \passthrough{\lstinline!None!} & 该方法不返回 Python
值；失败可能表现为异常或底层状态变化。 \\
\bottomrule
\end{longtable}

\textbf{对应 C++}

\begin{itemize}
\tightlist
\item
  类：\passthrough{\lstinline!unitree::robot::Server!}
\item
  签名：\passthrough{\lstinline!StartLease(float)!}
\item
  绑定策略：\passthrough{\lstinline!SIGNATURE\_PREVIEW!}
\item
  声明位置：\passthrough{\lstinline!include/unitree/robot/server/server.hpp:35!}
\end{itemize}

\textbf{说明}

\begin{itemize}
\tightlist
\item
  示例是局部调用片段；\passthrough{\lstinline!obj!}
  和参数变量表示已经按上层业务校验并准备好的对象。
\item
  该条目可以被编辑器和类型检查器识别，但当前运行时可能在属性查找阶段直接失败。
\end{itemize}

\textbf{用法}

\begin{lstlisting}[language=Python]
from typing import TYPE_CHECKING

if TYPE_CHECKING:
    def planned_start_lease(obj: Server, lease_term: float) -> None:
        obj.start_lease(lease_term=lease_term)
\end{lstlisting}

\begin{quote}
{[}!CAUTION{]} 上面的 \passthrough{\lstinline!TYPE\_CHECKING!}
示例只表达计划调用形态。该分支在普通 Python
运行时为假，不代表当前扩展能执行此接口。
\end{quote}

\hypertarget{unitree-sdk2-cpp-robot-server-get-name-1}{%
\paragraph{\texorpdfstring{\texttt{unitree\_sdk2\_cpp.robot.Server.get\_name}}{unitree\_sdk2\_cpp.robot.Server.get\_name}}\label{unitree-sdk2-cpp-robot-server-get-name-1}}

查询或检查 名称。

\textbf{签名}

\begin{lstlisting}[language=Python]
def get_name(self) -> str
\end{lstlisting}

\textbf{可用性与安全性}

\textbf{\passthrough{\lstinline!SIGNATURE\_ONLY!} /
\passthrough{\lstinline!UNCLASSIFIED!}}：当前只有设计期签名，不能假设已存在于运行时扩展。

\textbf{参数}

无显式参数。实例方法中的 \passthrough{\lstinline!self!} 由 Python
自动传入。

\textbf{返回值}

\begin{longtable}[]{@{}
  >{\raggedright\arraybackslash}p{(\columnwidth - 4\tabcolsep) * \real{0.18}}
  >{\raggedright\arraybackslash}p{(\columnwidth - 4\tabcolsep) * \real{0.20}}
  >{\raggedright\arraybackslash}p{(\columnwidth - 4\tabcolsep) * \real{0.62}}@{}}
\toprule
位置 & 类型 & 含义 \\
\midrule
\endhead
返回值 & \passthrough{\lstinline!str!} & 返回
\passthrough{\lstinline!str!}。更精确的业务含义见该方法用途、C++
签名和目标型号协议。 \\
\bottomrule
\end{longtable}

\textbf{对应 C++}

\begin{itemize}
\tightlist
\item
  类：\passthrough{\lstinline!unitree::robot::Server!}
\item
  签名：\passthrough{\lstinline!GetName()!}
\item
  绑定策略：\passthrough{\lstinline!SIGNATURE\_PREVIEW!}
\item
  声明位置：\passthrough{\lstinline!include/unitree/robot/server/server.hpp:37!}
\end{itemize}

\textbf{说明}

\begin{itemize}
\tightlist
\item
  示例是局部调用片段；\passthrough{\lstinline!obj!}
  和参数变量表示已经按上层业务校验并准备好的对象。
\item
  该条目可以被编辑器和类型检查器识别，但当前运行时可能在属性查找阶段直接失败。
\end{itemize}

\textbf{用法}

\begin{lstlisting}[language=Python]
from typing import TYPE_CHECKING

if TYPE_CHECKING:
    def planned_get_name(obj: Server) -> str:
        return obj.get_name()
\end{lstlisting}

\begin{quote}
{[}!CAUTION{]} 上面的 \passthrough{\lstinline!TYPE\_CHECKING!}
示例只表达计划调用形态。该分支在普通 Python
运行时为假，不代表当前扩展能执行此接口。
\end{quote}

\hypertarget{unitree-sdk2-cpp-robot-server-get-api-version-1}{%
\paragraph{\texorpdfstring{\texttt{unitree\_sdk2\_cpp.robot.Server.get\_api\_version}}{unitree\_sdk2\_cpp.robot.Server.get\_api\_version}}\label{unitree-sdk2-cpp-robot-server-get-api-version-1}}

查询或检查 API \passthrough{\lstinline!version!}。

\textbf{签名}

\begin{lstlisting}[language=Python]
def get_api_version(self) -> str
\end{lstlisting}

\textbf{可用性与安全性}

\textbf{\passthrough{\lstinline!SIGNATURE\_ONLY!} /
\passthrough{\lstinline!UNCLASSIFIED!}}：当前只有设计期签名，不能假设已存在于运行时扩展。

\textbf{参数}

无显式参数。实例方法中的 \passthrough{\lstinline!self!} 由 Python
自动传入。

\textbf{返回值}

\begin{longtable}[]{@{}
  >{\raggedright\arraybackslash}p{(\columnwidth - 4\tabcolsep) * \real{0.18}}
  >{\raggedright\arraybackslash}p{(\columnwidth - 4\tabcolsep) * \real{0.20}}
  >{\raggedright\arraybackslash}p{(\columnwidth - 4\tabcolsep) * \real{0.62}}@{}}
\toprule
位置 & 类型 & 含义 \\
\midrule
\endhead
返回值 & \passthrough{\lstinline!str!} & 返回
\passthrough{\lstinline!str!}。更精确的业务含义见该方法用途、C++
签名和目标型号协议。 \\
\bottomrule
\end{longtable}

\textbf{对应 C++}

\begin{itemize}
\tightlist
\item
  类：\passthrough{\lstinline!unitree::robot::Server!}
\item
  签名：\passthrough{\lstinline!GetApiVersion() const!}
\item
  绑定策略：\passthrough{\lstinline!SIGNATURE\_PREVIEW!}
\item
  声明位置：\passthrough{\lstinline!include/unitree/robot/server/server.hpp:38!}
\end{itemize}

\textbf{说明}

\begin{itemize}
\tightlist
\item
  示例是局部调用片段；\passthrough{\lstinline!obj!}
  和参数变量表示已经按上层业务校验并准备好的对象。
\item
  该条目可以被编辑器和类型检查器识别，但当前运行时可能在属性查找阶段直接失败。
\end{itemize}

\textbf{用法}

\begin{lstlisting}[language=Python]
from typing import TYPE_CHECKING

if TYPE_CHECKING:
    def planned_get_api_version(obj: Server) -> str:
        return obj.get_api_version()
\end{lstlisting}

\begin{quote}
{[}!CAUTION{]} 上面的 \passthrough{\lstinline!TYPE\_CHECKING!}
示例只表达计划调用形态。该分支在普通 Python
运行时为假，不代表当前扩展能执行此接口。
\end{quote}

\hypertarget{unitree-sdk2-cpp-robot-server-get-current-api-id-1}{%
\paragraph{\texorpdfstring{\texttt{unitree\_sdk2\_cpp.robot.Server.get\_current\_api\_id}}{unitree\_sdk2\_cpp.robot.Server.get\_current\_api\_id}}\label{unitree-sdk2-cpp-robot-server-get-current-api-id-1}}

查询或检查 \passthrough{\lstinline!current!} API ID。

\textbf{签名}

\begin{lstlisting}[language=Python]
def get_current_api_id(self) -> int
\end{lstlisting}

\textbf{可用性与安全性}

\textbf{\passthrough{\lstinline!SIGNATURE\_ONLY!} /
\passthrough{\lstinline!UNCLASSIFIED!}}：当前只有设计期签名，不能假设已存在于运行时扩展。

\textbf{参数}

无显式参数。实例方法中的 \passthrough{\lstinline!self!} 由 Python
自动传入。

\textbf{返回值}

\begin{longtable}[]{@{}
  >{\raggedright\arraybackslash}p{(\columnwidth - 4\tabcolsep) * \real{0.18}}
  >{\raggedright\arraybackslash}p{(\columnwidth - 4\tabcolsep) * \real{0.20}}
  >{\raggedright\arraybackslash}p{(\columnwidth - 4\tabcolsep) * \real{0.62}}@{}}
\toprule
位置 & 类型 & 含义 \\
\midrule
\endhead
返回值 & \passthrough{\lstinline!int!} & SDK 状态码；Unitree 示例通常以
\passthrough{\lstinline!0!} 表示成功，非零值需按具体服务定义解释。 \\
\bottomrule
\end{longtable}

\textbf{对应 C++}

\begin{itemize}
\tightlist
\item
  类：\passthrough{\lstinline!unitree::robot::Server!}
\item
  签名：\passthrough{\lstinline!GetCurrentApiId() const!}
\item
  绑定策略：\passthrough{\lstinline!SIGNATURE\_PREVIEW!}
\item
  声明位置：\passthrough{\lstinline!include/unitree/robot/server/server.hpp:40!}
\end{itemize}

\textbf{说明}

\begin{itemize}
\tightlist
\item
  示例是局部调用片段；\passthrough{\lstinline!obj!}
  和参数变量表示已经按上层业务校验并准备好的对象。
\item
  该条目可以被编辑器和类型检查器识别，但当前运行时可能在属性查找阶段直接失败。
\end{itemize}

\textbf{用法}

\begin{lstlisting}[language=Python]
from typing import TYPE_CHECKING

if TYPE_CHECKING:
    def planned_get_current_api_id(obj: Server) -> int:
        return obj.get_current_api_id()
\end{lstlisting}

\begin{quote}
{[}!CAUTION{]} 上面的 \passthrough{\lstinline!TYPE\_CHECKING!}
示例只表达计划调用形态。该分支在普通 Python
运行时为假，不代表当前扩展能执行此接口。
\end{quote}

\hypertarget{unitree-sdk2-cpp-robot-serverbase}{%
\subsubsection{\texorpdfstring{\texttt{unitree\_sdk2\_cpp.robot.ServerBase}}{unitree\_sdk2\_cpp.robot.ServerBase}}\label{unitree-sdk2-cpp-robot-serverbase}}

SDK 类型签名预览。请逐项查看构造函数和方法的可用性。

\textbf{导入}

\begin{lstlisting}[language=Python]
from unitree_sdk2_cpp.robot import ServerBase
\end{lstlisting}

\textbf{基类}：\passthrough{\lstinline!object!}

\textbf{方法索引}

\begin{longtable}[]{@{}
  >{\raggedright\arraybackslash}p{(\columnwidth - 6\tabcolsep) * \real{0.18}}
  >{\raggedright\arraybackslash}p{(\columnwidth - 6\tabcolsep) * \real{0.24}}
  >{\raggedright\arraybackslash}p{(\columnwidth - 6\tabcolsep) * \real{0.18}}
  >{\raggedright\arraybackslash}p{(\columnwidth - 6\tabcolsep) * \real{0.40}}@{}}
\toprule
方法 & Python 签名 & 状态 & 安全分类 \\
\midrule
\endhead
\protect\hyperlink{unitree-sdk2-cpp-robot-serverbase-dunder-init-1}{\passthrough{\lstinline!\_\_init\_\_!}}
& \passthrough{\lstinline!def \_\_init\_\_(self, name: str) -> None!} &
\passthrough{\lstinline!SIGNATURE\_ONLY!} &
\passthrough{\lstinline!CONSTRUCTION!} \\
\protect\hyperlink{unitree-sdk2-cpp-robot-serverbase-init-1}{\passthrough{\lstinline!init!}}
& \passthrough{\lstinline!def init(self) -> None!} &
\passthrough{\lstinline!SIGNATURE\_ONLY!} &
\passthrough{\lstinline!UNCLASSIFIED!} \\
\protect\hyperlink{unitree-sdk2-cpp-robot-serverbase-start-1}{\passthrough{\lstinline!start!}}
&
\passthrough{\lstinline!def start(self, enable\_proi\_queue: bool = ...) -> None!}
& \passthrough{\lstinline!SIGNATURE\_ONLY!} &
\passthrough{\lstinline!UNCLASSIFIED!} \\
\protect\hyperlink{unitree-sdk2-cpp-robot-serverbase-get-name-1}{\passthrough{\lstinline!get\_name!}}
& \passthrough{\lstinline!def get\_name(self) -> str!} &
\passthrough{\lstinline!SIGNATURE\_ONLY!} &
\passthrough{\lstinline!UNCLASSIFIED!} \\
\bottomrule
\end{longtable}

\hypertarget{unitree-sdk2-cpp-robot-serverbase-dunder-init-1}{%
\paragraph{\texorpdfstring{\texttt{unitree\_sdk2\_cpp.robot.ServerBase.\_\_init\_\_}}{unitree\_sdk2\_cpp.robot.ServerBase.\_\_init\_\_}}\label{unitree-sdk2-cpp-robot-serverbase-dunder-init-1}}

初始化 \passthrough{\lstinline!ServerBase!}
实例。是否能实际构造取决于下方可用性状态。

\textbf{签名}

\begin{lstlisting}[language=Python]
def __init__(self, name: str) -> None
\end{lstlisting}

\textbf{可用性与安全性}

\textbf{\passthrough{\lstinline!SIGNATURE\_ONLY!} /
\passthrough{\lstinline!CONSTRUCTION!}}：当前只有设计期签名，不能假设已存在于运行时扩展。

\textbf{参数}

\begin{longtable}[]{@{}
  >{\raggedright\arraybackslash}p{(\columnwidth - 6\tabcolsep) * \real{0.18}}
  >{\raggedright\arraybackslash}p{(\columnwidth - 6\tabcolsep) * \real{0.24}}
  >{\raggedright\arraybackslash}p{(\columnwidth - 6\tabcolsep) * \real{0.18}}
  >{\raggedright\arraybackslash}p{(\columnwidth - 6\tabcolsep) * \real{0.40}}@{}}
\toprule
名称 & Python 类型 & 默认值 & 含义 \\
\midrule
\endhead
\passthrough{\lstinline!name!} & \passthrough{\lstinline!str!} & 必填 &
SDK 对象、配置项、服务或动作的名称。允许值由调用它的具体接口定义。 对应
C++ 参数 \passthrough{\lstinline!name: const std::string \&!}。
底层为字符串；长度、编码和允许值由具体协议决定。 \\
\bottomrule
\end{longtable}

\textbf{返回值}

\begin{longtable}[]{@{}
  >{\raggedright\arraybackslash}p{(\columnwidth - 4\tabcolsep) * \real{0.18}}
  >{\raggedright\arraybackslash}p{(\columnwidth - 4\tabcolsep) * \real{0.20}}
  >{\raggedright\arraybackslash}p{(\columnwidth - 4\tabcolsep) * \real{0.62}}@{}}
\toprule
位置 & 类型 & 含义 \\
\midrule
\endhead
无 & \passthrough{\lstinline!None!} &
\passthrough{\lstinline!\_\_init\_\_!}
本身不返回值；调用类对象时，在构造函数可用的前提下得到该类实例。 \\
\bottomrule
\end{longtable}

\textbf{对应 C++}

\begin{itemize}
\tightlist
\item
  类：\passthrough{\lstinline!unitree::robot::ServerBase!}
\item
  签名：\passthrough{\lstinline!ServerBase(const std::string \&)!}
\item
  绑定策略：\passthrough{\lstinline!未单独标注!}
\item
  声明位置：\passthrough{\lstinline!include/unitree/robot/server/server\_base.hpp:13!}
\end{itemize}

\textbf{说明}

\begin{itemize}
\tightlist
\item
  示例是局部调用片段；\passthrough{\lstinline!obj!}
  和参数变量表示已经按上层业务校验并准备好的对象。
\item
  该条目可以被编辑器和类型检查器识别，但当前运行时可能在属性查找阶段直接失败。
\end{itemize}

\textbf{用法}

\begin{lstlisting}[language=Python]
from typing import TYPE_CHECKING

if TYPE_CHECKING:
    def planned_construct(name: str) -> ServerBase:
        return ServerBase(name=name)
\end{lstlisting}

\begin{quote}
{[}!CAUTION{]} 上面的 \passthrough{\lstinline!TYPE\_CHECKING!}
示例只表达计划调用形态。该分支在普通 Python
运行时为假，不代表当前扩展能执行此接口。
\end{quote}

\hypertarget{unitree-sdk2-cpp-robot-serverbase-init-1}{%
\paragraph{\texorpdfstring{\texttt{unitree\_sdk2\_cpp.robot.ServerBase.init}}{unitree\_sdk2\_cpp.robot.ServerBase.init}}\label{unitree-sdk2-cpp-robot-serverbase-init-1}}

初始化当前 SDK 对象所需的底层通道或服务资源。

\textbf{签名}

\begin{lstlisting}[language=Python]
def init(self) -> None
\end{lstlisting}

\textbf{可用性与安全性}

\textbf{\passthrough{\lstinline!SIGNATURE\_ONLY!} /
\passthrough{\lstinline!UNCLASSIFIED!}}：当前只有设计期签名，不能假设已存在于运行时扩展。

\textbf{参数}

无显式参数。实例方法中的 \passthrough{\lstinline!self!} 由 Python
自动传入。

\textbf{返回值}

\begin{longtable}[]{@{}
  >{\raggedright\arraybackslash}p{(\columnwidth - 4\tabcolsep) * \real{0.18}}
  >{\raggedright\arraybackslash}p{(\columnwidth - 4\tabcolsep) * \real{0.20}}
  >{\raggedright\arraybackslash}p{(\columnwidth - 4\tabcolsep) * \real{0.62}}@{}}
\toprule
位置 & 类型 & 含义 \\
\midrule
\endhead
无 & \passthrough{\lstinline!None!} & 该方法不返回 Python
值；失败可能表现为异常或底层状态变化。 \\
\bottomrule
\end{longtable}

\textbf{对应 C++}

\begin{itemize}
\tightlist
\item
  类：\passthrough{\lstinline!unitree::robot::ServerBase!}
\item
  签名：\passthrough{\lstinline!Init()!}
\item
  绑定策略：\passthrough{\lstinline!SIGNATURE\_PREVIEW!}
\item
  声明位置：\passthrough{\lstinline!include/unitree/robot/server/server\_base.hpp:16!}
\end{itemize}

\textbf{说明}

\begin{itemize}
\tightlist
\item
  示例是局部调用片段；\passthrough{\lstinline!obj!}
  和参数变量表示已经按上层业务校验并准备好的对象。
\item
  该条目可以被编辑器和类型检查器识别，但当前运行时可能在属性查找阶段直接失败。
\end{itemize}

\textbf{用法}

\begin{lstlisting}[language=Python]
from typing import TYPE_CHECKING

if TYPE_CHECKING:
    def planned_init(obj: ServerBase) -> None:
        obj.init()
\end{lstlisting}

\begin{quote}
{[}!CAUTION{]} 上面的 \passthrough{\lstinline!TYPE\_CHECKING!}
示例只表达计划调用形态。该分支在普通 Python
运行时为假，不代表当前扩展能执行此接口。
\end{quote}

\hypertarget{unitree-sdk2-cpp-robot-serverbase-start-1}{%
\paragraph{\texorpdfstring{\texttt{unitree\_sdk2\_cpp.robot.ServerBase.start}}{unitree\_sdk2\_cpp.robot.ServerBase.start}}\label{unitree-sdk2-cpp-robot-serverbase-start-1}}

启动对应 SDK 操作。具体副作用和安全边界见下方状态。

\textbf{签名}

\begin{lstlisting}[language=Python]
def start(self, enable_proi_queue: bool = ...) -> None
\end{lstlisting}

\textbf{可用性与安全性}

\textbf{\passthrough{\lstinline!SIGNATURE\_ONLY!} /
\passthrough{\lstinline!UNCLASSIFIED!}}：当前只有设计期签名，不能假设已存在于运行时扩展。

\textbf{参数}

\begin{longtable}[]{@{}
  >{\raggedright\arraybackslash}p{(\columnwidth - 6\tabcolsep) * \real{0.18}}
  >{\raggedright\arraybackslash}p{(\columnwidth - 6\tabcolsep) * \real{0.24}}
  >{\raggedright\arraybackslash}p{(\columnwidth - 6\tabcolsep) * \real{0.18}}
  >{\raggedright\arraybackslash}p{(\columnwidth - 6\tabcolsep) * \real{0.40}}@{}}
\toprule
名称 & Python 类型 & 默认值 & 含义 \\
\midrule
\endhead
\passthrough{\lstinline!enable\_proi\_queue!} &
\passthrough{\lstinline!bool!} & \passthrough{\lstinline!...!} &
是否启用 SDK 中名为 \passthrough{\lstinline!proi!}
的队列选项；该名称沿用上游接口，具体行为需查对应头文件。 对应 C++ 参数
\passthrough{\lstinline!enableProiQueue: bool!}。 只接受布尔语义。 \\
\bottomrule
\end{longtable}

\textbf{返回值}

\begin{longtable}[]{@{}
  >{\raggedright\arraybackslash}p{(\columnwidth - 4\tabcolsep) * \real{0.18}}
  >{\raggedright\arraybackslash}p{(\columnwidth - 4\tabcolsep) * \real{0.20}}
  >{\raggedright\arraybackslash}p{(\columnwidth - 4\tabcolsep) * \real{0.62}}@{}}
\toprule
位置 & 类型 & 含义 \\
\midrule
\endhead
无 & \passthrough{\lstinline!None!} & 该方法不返回 Python
值；失败可能表现为异常或底层状态变化。 \\
\bottomrule
\end{longtable}

\textbf{对应 C++}

\begin{itemize}
\tightlist
\item
  类：\passthrough{\lstinline!unitree::robot::ServerBase!}
\item
  签名：\passthrough{\lstinline!Start(bool)!}
\item
  绑定策略：\passthrough{\lstinline!SIGNATURE\_PREVIEW!}
\item
  声明位置：\passthrough{\lstinline!include/unitree/robot/server/server\_base.hpp:17!}
\end{itemize}

\textbf{说明}

\begin{itemize}
\tightlist
\item
  示例是局部调用片段；\passthrough{\lstinline!obj!}
  和参数变量表示已经按上层业务校验并准备好的对象。
\item
  默认值显示为 \passthrough{\lstinline!...!}，表示 C++
  声明存在默认参数，但当前 AST
  清单没有保存其字面量；省略参数可使用上游默认行为。
\item
  该条目可以被编辑器和类型检查器识别，但当前运行时可能在属性查找阶段直接失败。
\end{itemize}

\textbf{用法}

\begin{lstlisting}[language=Python]
from typing import TYPE_CHECKING

if TYPE_CHECKING:
    def planned_start(obj: ServerBase, enable_proi_queue: bool) -> None:
        obj.start(enable_proi_queue=enable_proi_queue)
\end{lstlisting}

\begin{quote}
{[}!CAUTION{]} 上面的 \passthrough{\lstinline!TYPE\_CHECKING!}
示例只表达计划调用形态。该分支在普通 Python
运行时为假，不代表当前扩展能执行此接口。
\end{quote}

\hypertarget{unitree-sdk2-cpp-robot-serverbase-get-name-1}{%
\paragraph{\texorpdfstring{\texttt{unitree\_sdk2\_cpp.robot.ServerBase.get\_name}}{unitree\_sdk2\_cpp.robot.ServerBase.get\_name}}\label{unitree-sdk2-cpp-robot-serverbase-get-name-1}}

查询或检查 名称。

\textbf{签名}

\begin{lstlisting}[language=Python]
def get_name(self) -> str
\end{lstlisting}

\textbf{可用性与安全性}

\textbf{\passthrough{\lstinline!SIGNATURE\_ONLY!} /
\passthrough{\lstinline!UNCLASSIFIED!}}：当前只有设计期签名，不能假设已存在于运行时扩展。

\textbf{参数}

无显式参数。实例方法中的 \passthrough{\lstinline!self!} 由 Python
自动传入。

\textbf{返回值}

\begin{longtable}[]{@{}
  >{\raggedright\arraybackslash}p{(\columnwidth - 4\tabcolsep) * \real{0.18}}
  >{\raggedright\arraybackslash}p{(\columnwidth - 4\tabcolsep) * \real{0.20}}
  >{\raggedright\arraybackslash}p{(\columnwidth - 4\tabcolsep) * \real{0.62}}@{}}
\toprule
位置 & 类型 & 含义 \\
\midrule
\endhead
返回值 & \passthrough{\lstinline!str!} & 返回
\passthrough{\lstinline!str!}。更精确的业务含义见该方法用途、C++
签名和目标型号协议。 \\
\bottomrule
\end{longtable}

\textbf{对应 C++}

\begin{itemize}
\tightlist
\item
  类：\passthrough{\lstinline!unitree::robot::ServerBase!}
\item
  签名：\passthrough{\lstinline!GetName() const!}
\item
  绑定策略：\passthrough{\lstinline!SIGNATURE\_PREVIEW!}
\item
  声明位置：\passthrough{\lstinline!include/unitree/robot/server/server\_base.hpp:19!}
\end{itemize}

\textbf{说明}

\begin{itemize}
\tightlist
\item
  示例是局部调用片段；\passthrough{\lstinline!obj!}
  和参数变量表示已经按上层业务校验并准备好的对象。
\item
  该条目可以被编辑器和类型检查器识别，但当前运行时可能在属性查找阶段直接失败。
\end{itemize}

\textbf{用法}

\begin{lstlisting}[language=Python]
from typing import TYPE_CHECKING

if TYPE_CHECKING:
    def planned_get_name(obj: ServerBase) -> str:
        return obj.get_name()
\end{lstlisting}

\begin{quote}
{[}!CAUTION{]} 上面的 \passthrough{\lstinline!TYPE\_CHECKING!}
示例只表达计划调用形态。该分支在普通 Python
运行时为假，不代表当前扩展能执行此接口。
\end{quote}

\hypertarget{unitree-sdk2-cpp-robot-serverchannelnamer}{%
\subsubsection{\texorpdfstring{\texttt{unitree\_sdk2\_cpp.robot.ServerChannelNamer}}{unitree\_sdk2\_cpp.robot.ServerChannelNamer}}\label{unitree-sdk2-cpp-robot-serverchannelnamer}}

SDK 类型签名预览。请逐项查看构造函数和方法的可用性。

\textbf{导入}

\begin{lstlisting}[language=Python]
from unitree_sdk2_cpp.robot import ServerChannelNamer
\end{lstlisting}

\textbf{基类}：\passthrough{\lstinline!ChannelNamer!}

\textbf{方法索引}

\begin{longtable}[]{@{}
  >{\raggedright\arraybackslash}p{(\columnwidth - 6\tabcolsep) * \real{0.18}}
  >{\raggedright\arraybackslash}p{(\columnwidth - 6\tabcolsep) * \real{0.24}}
  >{\raggedright\arraybackslash}p{(\columnwidth - 6\tabcolsep) * \real{0.18}}
  >{\raggedright\arraybackslash}p{(\columnwidth - 6\tabcolsep) * \real{0.40}}@{}}
\toprule
方法 & Python 签名 & 状态 & 安全分类 \\
\midrule
\endhead
\protect\hyperlink{unitree-sdk2-cpp-robot-serverchannelnamer-dunder-init-1}{\passthrough{\lstinline!\_\_init\_\_!}}
& \passthrough{\lstinline!def \_\_init\_\_(self) -> None!} &
\passthrough{\lstinline!SIGNATURE\_ONLY!} &
\passthrough{\lstinline!CONSTRUCTION!} \\
\bottomrule
\end{longtable}

\hypertarget{unitree-sdk2-cpp-robot-serverchannelnamer-dunder-init-1}{%
\paragraph{\texorpdfstring{\texttt{unitree\_sdk2\_cpp.robot.ServerChannelNamer.\_\_init\_\_}}{unitree\_sdk2\_cpp.robot.ServerChannelNamer.\_\_init\_\_}}\label{unitree-sdk2-cpp-robot-serverchannelnamer-dunder-init-1}}

初始化 \passthrough{\lstinline!ServerChannelNamer!}
实例。是否能实际构造取决于下方可用性状态。

\textbf{签名}

\begin{lstlisting}[language=Python]
def __init__(self) -> None
\end{lstlisting}

\textbf{可用性与安全性}

\textbf{\passthrough{\lstinline!SIGNATURE\_ONLY!} /
\passthrough{\lstinline!CONSTRUCTION!}}：当前只有设计期签名，不能假设已存在于运行时扩展。

\textbf{参数}

无显式参数。实例方法中的 \passthrough{\lstinline!self!} 由 Python
自动传入。

\textbf{返回值}

\begin{longtable}[]{@{}
  >{\raggedright\arraybackslash}p{(\columnwidth - 4\tabcolsep) * \real{0.18}}
  >{\raggedright\arraybackslash}p{(\columnwidth - 4\tabcolsep) * \real{0.20}}
  >{\raggedright\arraybackslash}p{(\columnwidth - 4\tabcolsep) * \real{0.62}}@{}}
\toprule
位置 & 类型 & 含义 \\
\midrule
\endhead
无 & \passthrough{\lstinline!None!} &
\passthrough{\lstinline!\_\_init\_\_!}
本身不返回值；调用类对象时，在构造函数可用的前提下得到该类实例。 \\
\bottomrule
\end{longtable}

\textbf{对应 C++}

\begin{itemize}
\tightlist
\item
  类：\passthrough{\lstinline!unitree::robot::ServerChannelNamer!}
\item
  签名：\passthrough{\lstinline!ServerChannelNamer()!}
\item
  绑定策略：\passthrough{\lstinline!未单独标注!}
\item
  声明位置：\passthrough{\lstinline!include/unitree/robot/channel/channel\_namer.hpp:51!}
\end{itemize}

\textbf{说明}

\begin{itemize}
\tightlist
\item
  示例是局部调用片段；\passthrough{\lstinline!obj!}
  和参数变量表示已经按上层业务校验并准备好的对象。
\item
  该条目可以被编辑器和类型检查器识别，但当前运行时可能在属性查找阶段直接失败。
\end{itemize}

\textbf{用法}

\begin{lstlisting}[language=Python]
from typing import TYPE_CHECKING

if TYPE_CHECKING:
    def planned_construct() -> ServerChannelNamer:
        return ServerChannelNamer()
\end{lstlisting}

\begin{quote}
{[}!CAUTION{]} 上面的 \passthrough{\lstinline!TYPE\_CHECKING!}
示例只表达计划调用形态。该分支在普通 Python
运行时为假，不代表当前扩展能执行此接口。
\end{quote}

\hypertarget{unitree-sdk2-cpp-robot-serverstub}{%
\subsubsection{\texorpdfstring{\texttt{unitree\_sdk2\_cpp.robot.ServerStub}}{unitree\_sdk2\_cpp.robot.ServerStub}}\label{unitree-sdk2-cpp-robot-serverstub}}

SDK 类型签名预览。请逐项查看构造函数和方法的可用性。

\textbf{导入}

\begin{lstlisting}[language=Python]
from unitree_sdk2_cpp.robot import ServerStub
\end{lstlisting}

\textbf{基类}：\passthrough{\lstinline!object!}

\textbf{方法索引}

\begin{longtable}[]{@{}
  >{\raggedright\arraybackslash}p{(\columnwidth - 6\tabcolsep) * \real{0.18}}
  >{\raggedright\arraybackslash}p{(\columnwidth - 6\tabcolsep) * \real{0.24}}
  >{\raggedright\arraybackslash}p{(\columnwidth - 6\tabcolsep) * \real{0.18}}
  >{\raggedright\arraybackslash}p{(\columnwidth - 6\tabcolsep) * \real{0.40}}@{}}
\toprule
方法 & Python 签名 & 状态 & 安全分类 \\
\midrule
\endhead
\protect\hyperlink{unitree-sdk2-cpp-robot-serverstub-dunder-init-1}{\passthrough{\lstinline!\_\_init\_\_!}}
& \passthrough{\lstinline!def \_\_init\_\_(self) -> None!} &
\passthrough{\lstinline!SIGNATURE\_ONLY!} &
\passthrough{\lstinline!CONSTRUCTION!} \\
\protect\hyperlink{unitree-sdk2-cpp-robot-serverstub-init-1}{\passthrough{\lstinline!init!}}
&
\passthrough{\lstinline!def init(self, name: str, handler: Callable[..., Any], enable\_proi\_queue: bool) -> None!}
& \passthrough{\lstinline!SIGNATURE\_ONLY!} &
\passthrough{\lstinline!UNCLASSIFIED!} \\
\protect\hyperlink{unitree-sdk2-cpp-robot-serverstub-send-1}{\passthrough{\lstinline!send!}}
&
\passthrough{\lstinline!def send(self, response: Any, timeout: int = ...) -> bool!}
& \passthrough{\lstinline!SIGNATURE\_ONLY!} &
\passthrough{\lstinline!UNCLASSIFIED!} \\
\bottomrule
\end{longtable}

\hypertarget{unitree-sdk2-cpp-robot-serverstub-dunder-init-1}{%
\paragraph{\texorpdfstring{\texttt{unitree\_sdk2\_cpp.robot.ServerStub.\_\_init\_\_}}{unitree\_sdk2\_cpp.robot.ServerStub.\_\_init\_\_}}\label{unitree-sdk2-cpp-robot-serverstub-dunder-init-1}}

初始化 \passthrough{\lstinline!ServerStub!}
实例。是否能实际构造取决于下方可用性状态。

\textbf{签名}

\begin{lstlisting}[language=Python]
def __init__(self) -> None
\end{lstlisting}

\textbf{可用性与安全性}

\textbf{\passthrough{\lstinline!SIGNATURE\_ONLY!} /
\passthrough{\lstinline!CONSTRUCTION!}}：当前只有设计期签名，不能假设已存在于运行时扩展。

\textbf{参数}

无显式参数。实例方法中的 \passthrough{\lstinline!self!} 由 Python
自动传入。

\textbf{返回值}

\begin{longtable}[]{@{}
  >{\raggedright\arraybackslash}p{(\columnwidth - 4\tabcolsep) * \real{0.18}}
  >{\raggedright\arraybackslash}p{(\columnwidth - 4\tabcolsep) * \real{0.20}}
  >{\raggedright\arraybackslash}p{(\columnwidth - 4\tabcolsep) * \real{0.62}}@{}}
\toprule
位置 & 类型 & 含义 \\
\midrule
\endhead
无 & \passthrough{\lstinline!None!} &
\passthrough{\lstinline!\_\_init\_\_!}
本身不返回值；调用类对象时，在构造函数可用的前提下得到该类实例。 \\
\bottomrule
\end{longtable}

\textbf{对应 C++}

\begin{itemize}
\tightlist
\item
  类：\passthrough{\lstinline!unitree::robot::ServerStub!}
\item
  签名：\passthrough{\lstinline!ServerStub()!}
\item
  绑定策略：\passthrough{\lstinline!未单独标注!}
\item
  声明位置：\passthrough{\lstinline!include/unitree/robot/server/server\_stub.hpp:18!}
\end{itemize}

\textbf{说明}

\begin{itemize}
\tightlist
\item
  示例是局部调用片段；\passthrough{\lstinline!obj!}
  和参数变量表示已经按上层业务校验并准备好的对象。
\item
  该条目可以被编辑器和类型检查器识别，但当前运行时可能在属性查找阶段直接失败。
\end{itemize}

\textbf{用法}

\begin{lstlisting}[language=Python]
from typing import TYPE_CHECKING

if TYPE_CHECKING:
    def planned_construct() -> ServerStub:
        return ServerStub()
\end{lstlisting}

\begin{quote}
{[}!CAUTION{]} 上面的 \passthrough{\lstinline!TYPE\_CHECKING!}
示例只表达计划调用形态。该分支在普通 Python
运行时为假，不代表当前扩展能执行此接口。
\end{quote}

\hypertarget{unitree-sdk2-cpp-robot-serverstub-init-1}{%
\paragraph{\texorpdfstring{\texttt{unitree\_sdk2\_cpp.robot.ServerStub.init}}{unitree\_sdk2\_cpp.robot.ServerStub.init}}\label{unitree-sdk2-cpp-robot-serverstub-init-1}}

初始化当前 SDK 对象所需的底层通道或服务资源。

\textbf{签名}

\begin{lstlisting}[language=Python]
def init(self, name: str, handler: Callable[..., Any], enable_proi_queue: bool) -> None
\end{lstlisting}

\textbf{可用性与安全性}

\textbf{\passthrough{\lstinline!SIGNATURE\_ONLY!} /
\passthrough{\lstinline!UNCLASSIFIED!}}：当前只有设计期签名，不能假设已存在于运行时扩展。

\textbf{参数}

\begin{longtable}[]{@{}
  >{\raggedright\arraybackslash}p{(\columnwidth - 6\tabcolsep) * \real{0.18}}
  >{\raggedright\arraybackslash}p{(\columnwidth - 6\tabcolsep) * \real{0.24}}
  >{\raggedright\arraybackslash}p{(\columnwidth - 6\tabcolsep) * \real{0.18}}
  >{\raggedright\arraybackslash}p{(\columnwidth - 6\tabcolsep) * \real{0.40}}@{}}
\toprule
名称 & Python 类型 & 默认值 & 含义 \\
\midrule
\endhead
\passthrough{\lstinline!name!} & \passthrough{\lstinline!str!} & 必填 &
SDK 对象、配置项、服务或动作的名称。允许值由调用它的具体接口定义。 对应
C++ 参数 \passthrough{\lstinline!name: const std::string \&!}。
底层为字符串；长度、编码和允许值由具体协议决定。 \\
\passthrough{\lstinline!handler!} &
\passthrough{\lstinline!Callable[..., Any]!} & 必填 & 传给该接口的
\passthrough{\lstinline!handler!} 参数，Python 类型为
\passthrough{\lstinline!Callable[..., Any]!}。精确含义、单位、范围和枚举值需查目标型号对应头文件与协议。
对应 C++ 参数
\passthrough{\lstinline!handler: const unitree::robot::ServerRequestHandler \&!}。 \\
\passthrough{\lstinline!enable\_proi\_queue!} &
\passthrough{\lstinline!bool!} & 必填 & 是否启用 SDK 中名为
\passthrough{\lstinline!proi!}
的队列选项；该名称沿用上游接口，具体行为需查对应头文件。 对应 C++ 参数
\passthrough{\lstinline!enableProiQueue: bool!}。 只接受布尔语义。 \\
\bottomrule
\end{longtable}

\textbf{返回值}

\begin{longtable}[]{@{}
  >{\raggedright\arraybackslash}p{(\columnwidth - 4\tabcolsep) * \real{0.18}}
  >{\raggedright\arraybackslash}p{(\columnwidth - 4\tabcolsep) * \real{0.20}}
  >{\raggedright\arraybackslash}p{(\columnwidth - 4\tabcolsep) * \real{0.62}}@{}}
\toprule
位置 & 类型 & 含义 \\
\midrule
\endhead
无 & \passthrough{\lstinline!None!} & 该方法不返回 Python
值；失败可能表现为异常或底层状态变化。 \\
\bottomrule
\end{longtable}

\textbf{对应 C++}

\begin{itemize}
\tightlist
\item
  类：\passthrough{\lstinline!unitree::robot::ServerStub!}
\item
  签名：\passthrough{\lstinline!Init(const std::string \&, const unitree::robot::ServerRequestHandler \&, bool)!}
\item
  绑定策略：\passthrough{\lstinline!SIGNATURE\_PREVIEW!}
\item
  声明位置：\passthrough{\lstinline!include/unitree/robot/server/server\_stub.hpp:21!}
\end{itemize}

\textbf{说明}

\begin{itemize}
\tightlist
\item
  示例是局部调用片段；\passthrough{\lstinline!obj!}
  和参数变量表示已经按上层业务校验并准备好的对象。
\item
  该条目可以被编辑器和类型检查器识别，但当前运行时可能在属性查找阶段直接失败。
\item
  回调可能由 SDK
  工作线程触发；回调应快速返回、捕获异常，并通过线程安全队列移交耗时工作。
\end{itemize}

\textbf{用法}

\begin{lstlisting}[language=Python]
from typing import TYPE_CHECKING

if TYPE_CHECKING:
    def planned_init(obj: ServerStub, name: str, handler: Callable[..., Any], enable_proi_queue: bool) -> None:
        obj.init(name=name, handler=handler, enable_proi_queue=enable_proi_queue)
\end{lstlisting}

\begin{quote}
{[}!CAUTION{]} 上面的 \passthrough{\lstinline!TYPE\_CHECKING!}
示例只表达计划调用形态。该分支在普通 Python
运行时为假，不代表当前扩展能执行此接口。
\end{quote}

\hypertarget{unitree-sdk2-cpp-robot-serverstub-send-1}{%
\paragraph{\texorpdfstring{\texttt{unitree\_sdk2\_cpp.robot.ServerStub.send}}{unitree\_sdk2\_cpp.robot.ServerStub.send}}\label{unitree-sdk2-cpp-robot-serverstub-send-1}}

对应 C++ SDK 操作
\passthrough{\lstinline!Send(const unitree::robot::Response \&, int64\_t)!}。上游头文件没有可直接生成的业务说明时，本参考只保证签名映射，精确语义需查目标型号协议。

\textbf{签名}

\begin{lstlisting}[language=Python]
def send(self, response: Any, timeout: int = ...) -> bool
\end{lstlisting}

\textbf{可用性与安全性}

\textbf{\passthrough{\lstinline!SIGNATURE\_ONLY!} /
\passthrough{\lstinline!UNCLASSIFIED!}}：当前只有设计期签名，不能假设已存在于运行时扩展。

\textbf{参数}

\begin{longtable}[]{@{}
  >{\raggedright\arraybackslash}p{(\columnwidth - 6\tabcolsep) * \real{0.18}}
  >{\raggedright\arraybackslash}p{(\columnwidth - 6\tabcolsep) * \real{0.24}}
  >{\raggedright\arraybackslash}p{(\columnwidth - 6\tabcolsep) * \real{0.18}}
  >{\raggedright\arraybackslash}p{(\columnwidth - 6\tabcolsep) * \real{0.40}}@{}}
\toprule
名称 & Python 类型 & 默认值 & 含义 \\
\midrule
\endhead
\passthrough{\lstinline!response!} & \passthrough{\lstinline!Any!} &
必填 & 传给该接口的 响应 参数，Python 类型为
\passthrough{\lstinline!Any!}。精确含义、单位、范围和枚举值需查目标型号对应头文件与协议。
对应 C++ 参数
\passthrough{\lstinline!response: const unitree::robot::Response \&!}。 \\
\passthrough{\lstinline!timeout!} & \passthrough{\lstinline!int!} &
\passthrough{\lstinline!...!} & 传给该接口的 超时 参数，Python 类型为
\passthrough{\lstinline!int!}。精确含义、单位、范围和枚举值需查目标型号对应头文件与协议。
对应 C++ 参数 \passthrough{\lstinline!timeout: int64\_t!}。 取值范围为
-9223372036854775808 到 9223372036854775807。 \\
\bottomrule
\end{longtable}

\textbf{返回值}

\begin{longtable}[]{@{}
  >{\raggedright\arraybackslash}p{(\columnwidth - 4\tabcolsep) * \real{0.18}}
  >{\raggedright\arraybackslash}p{(\columnwidth - 4\tabcolsep) * \real{0.20}}
  >{\raggedright\arraybackslash}p{(\columnwidth - 4\tabcolsep) * \real{0.62}}@{}}
\toprule
位置 & 类型 & 含义 \\
\midrule
\endhead
返回值 & \passthrough{\lstinline!bool!} & 返回
\passthrough{\lstinline!bool!}。更精确的业务含义见该方法用途、C++
签名和目标型号协议。 \\
\bottomrule
\end{longtable}

\textbf{对应 C++}

\begin{itemize}
\tightlist
\item
  类：\passthrough{\lstinline!unitree::robot::ServerStub!}
\item
  签名：\passthrough{\lstinline!Send(const unitree::robot::Response \&, int64\_t)!}
\item
  绑定策略：\passthrough{\lstinline!SIGNATURE\_PREVIEW!}
\item
  声明位置：\passthrough{\lstinline!include/unitree/robot/server/server\_stub.hpp:22!}
\end{itemize}

\textbf{说明}

\begin{itemize}
\tightlist
\item
  示例是局部调用片段；\passthrough{\lstinline!obj!}
  和参数变量表示已经按上层业务校验并准备好的对象。
\item
  默认值显示为 \passthrough{\lstinline!...!}，表示 C++
  声明存在默认参数，但当前 AST
  清单没有保存其字面量；省略参数可使用上游默认行为。
\item
  该条目可以被编辑器和类型检查器识别，但当前运行时可能在属性查找阶段直接失败。
\end{itemize}

\textbf{用法}

\begin{lstlisting}[language=Python]
from typing import TYPE_CHECKING

if TYPE_CHECKING:
    def planned_send(obj: ServerStub, response: Any, timeout: int) -> bool:
        return obj.send(response=response, timeout=timeout)
\end{lstlisting}

\begin{quote}
{[}!CAUTION{]} 上面的 \passthrough{\lstinline!TYPE\_CHECKING!}
示例只表达计划调用形态。该分支在普通 Python
运行时为假，不代表当前扩展能执行此接口。
\end{quote}

\begin{center}\rule{0.5\linewidth}{0.5pt}\end{center}

\hypertarget{unitree-sdk2-cpp-robot-a2}{%
\subsection{A2 Robot API}\label{unitree-sdk2-cpp-robot-a2}}

模块：\passthrough{\lstinline!unitree\_sdk2\_cpp.robot.a2!}

\hypertarget{ux7c7bux7d22ux5f15-7}{%
\subsubsection{类索引}\label{ux7c7bux7d22ux5f15-7}}

\begin{longtable}[]{@{}
  >{\raggedright\arraybackslash}p{(\columnwidth - 4\tabcolsep) * \real{0.18}}
  >{\raggedleft\arraybackslash}p{(\columnwidth - 4\tabcolsep) * \real{0.20}}
  >{\raggedleft\arraybackslash}p{(\columnwidth - 4\tabcolsep) * \real{0.62}}@{}}
\toprule
类 & 公开函数签名 & 属性 \\
\midrule
\endhead
\protect\hyperlink{unitree-sdk2-cpp-robot-a2-audioclient}{\passthrough{\lstinline!AudioClient!}}
& 8 & 0 \\
\protect\hyperlink{unitree-sdk2-cpp-robot-a2-ledcontrolparameter}{\passthrough{\lstinline!LedControlParameter!}}
& 3 & 0 \\
\protect\hyperlink{unitree-sdk2-cpp-robot-a2-pathpoint}{\passthrough{\lstinline!PathPoint!}}
& 0 & 0 \\
\protect\hyperlink{unitree-sdk2-cpp-robot-a2-playstopparameter}{\passthrough{\lstinline!PlayStopParameter!}}
& 3 & 0 \\
\protect\hyperlink{unitree-sdk2-cpp-robot-a2-playstreamparameter}{\passthrough{\lstinline!PlayStreamParameter!}}
& 3 & 0 \\
\protect\hyperlink{unitree-sdk2-cpp-robot-a2-posevec4}{\passthrough{\lstinline!PoseVec4!}}
& 3 & 0 \\
\protect\hyperlink{unitree-sdk2-cpp-robot-a2-sportclient}{\passthrough{\lstinline!SportClient!}}
& 24 & 0 \\
\protect\hyperlink{unitree-sdk2-cpp-robot-a2-ttsmakerparameter}{\passthrough{\lstinline!TtsMakerParameter!}}
& 3 & 0 \\
\bottomrule
\end{longtable}

\hypertarget{unitree-sdk2-cpp-robot-a2-audioclient}{%
\subsubsection{\texorpdfstring{\texttt{unitree\_sdk2\_cpp.robot.a2.AudioClient}}{unitree\_sdk2\_cpp.robot.a2.AudioClient}}\label{unitree-sdk2-cpp-robot-a2-audioclient}}

Robot 服务客户端。构造、初始化和具体方法的可用性必须分别检查。

\textbf{导入}

\begin{lstlisting}[language=Python]
from unitree_sdk2_cpp.robot.a2 import AudioClient
\end{lstlisting}

\textbf{基类}：\passthrough{\lstinline!Client!}

\textbf{方法索引}

\begin{longtable}[]{@{}
  >{\raggedright\arraybackslash}p{(\columnwidth - 6\tabcolsep) * \real{0.18}}
  >{\raggedright\arraybackslash}p{(\columnwidth - 6\tabcolsep) * \real{0.24}}
  >{\raggedright\arraybackslash}p{(\columnwidth - 6\tabcolsep) * \real{0.18}}
  >{\raggedright\arraybackslash}p{(\columnwidth - 6\tabcolsep) * \real{0.40}}@{}}
\toprule
方法 & Python 签名 & 状态 & 安全分类 \\
\midrule
\endhead
\protect\hyperlink{unitree-sdk2-cpp-robot-a2-audioclient-dunder-init-1}{\passthrough{\lstinline!\_\_init\_\_!}}
& \passthrough{\lstinline!def \_\_init\_\_(self) -> None!} &
\passthrough{\lstinline!AVAILABLE!} &
\passthrough{\lstinline!CONSTRUCTION!} \\
\protect\hyperlink{unitree-sdk2-cpp-robot-a2-audioclient-init-1}{\passthrough{\lstinline!init!}}
& \passthrough{\lstinline!def init(self) -> None!} &
\passthrough{\lstinline!AVAILABLE!} &
\passthrough{\lstinline!INITIALIZATION!} \\
\protect\hyperlink{unitree-sdk2-cpp-robot-a2-audioclient-tts-maker-1}{\passthrough{\lstinline!tts\_maker!}}
&
\passthrough{\lstinline!def tts\_maker(self, text: str, speaker\_id: int) -> int!}
& \passthrough{\lstinline!SIGNATURE\_ONLY!} &
\passthrough{\lstinline!HARDWARE\_SIDE\_EFFECT!} \\
\protect\hyperlink{unitree-sdk2-cpp-robot-a2-audioclient-get-volume-1}{\passthrough{\lstinline!get\_volume!}}
& \passthrough{\lstinline!def get\_volume(self) -> tuple[int, int]!} &
\passthrough{\lstinline!AVAILABLE!} &
\passthrough{\lstinline!READ\_ONLY!} \\
\protect\hyperlink{unitree-sdk2-cpp-robot-a2-audioclient-set-volume-1}{\passthrough{\lstinline!set\_volume!}}
& \passthrough{\lstinline!def set\_volume(self, volume: int) -> int!} &
\passthrough{\lstinline!SIGNATURE\_ONLY!} &
\passthrough{\lstinline!HARDWARE\_SIDE\_EFFECT!} \\
\protect\hyperlink{unitree-sdk2-cpp-robot-a2-audioclient-play-stream-1}{\passthrough{\lstinline!play\_stream!}}
&
\passthrough{\lstinline!def play\_stream(self, app\_name: str, stream\_id: str, pcm\_data: Sequence[int]) -> int!}
& \passthrough{\lstinline!SIGNATURE\_ONLY!} &
\passthrough{\lstinline!HARDWARE\_SIDE\_EFFECT!} \\
\protect\hyperlink{unitree-sdk2-cpp-robot-a2-audioclient-play-stop-1}{\passthrough{\lstinline!play\_stop!}}
& \passthrough{\lstinline!def play\_stop(self, app\_name: str) -> int!}
& \passthrough{\lstinline!SIGNATURE\_ONLY!} &
\passthrough{\lstinline!HARDWARE\_SIDE\_EFFECT!} \\
\protect\hyperlink{unitree-sdk2-cpp-robot-a2-audioclient-led-control-1}{\passthrough{\lstinline!led\_control!}}
&
\passthrough{\lstinline!def led\_control(self, r: int, g: int, b: int) -> int!}
& \passthrough{\lstinline!SIGNATURE\_ONLY!} &
\passthrough{\lstinline!HARDWARE\_SIDE\_EFFECT!} \\
\bottomrule
\end{longtable}

\hypertarget{unitree-sdk2-cpp-robot-a2-audioclient-dunder-init-1}{%
\paragraph{\texorpdfstring{\texttt{unitree\_sdk2\_cpp.robot.a2.AudioClient.\_\_init\_\_}}{unitree\_sdk2\_cpp.robot.a2.AudioClient.\_\_init\_\_}}\label{unitree-sdk2-cpp-robot-a2-audioclient-dunder-init-1}}

初始化 \passthrough{\lstinline!AudioClient!}
实例。是否能实际构造取决于下方可用性状态。

\textbf{签名}

\begin{lstlisting}[language=Python]
def __init__(self) -> None
\end{lstlisting}

\textbf{可用性与安全性}

\textbf{\passthrough{\lstinline!AVAILABLE!} /
\passthrough{\lstinline!CONSTRUCTION!}}：当前绑定源码已实现；实际执行条件仍取决于平台和对应
SDK 环境。

\textbf{参数}

无显式参数。实例方法中的 \passthrough{\lstinline!self!} 由 Python
自动传入。

\textbf{返回值}

\begin{longtable}[]{@{}
  >{\raggedright\arraybackslash}p{(\columnwidth - 4\tabcolsep) * \real{0.18}}
  >{\raggedright\arraybackslash}p{(\columnwidth - 4\tabcolsep) * \real{0.20}}
  >{\raggedright\arraybackslash}p{(\columnwidth - 4\tabcolsep) * \real{0.62}}@{}}
\toprule
位置 & 类型 & 含义 \\
\midrule
\endhead
无 & \passthrough{\lstinline!None!} &
\passthrough{\lstinline!\_\_init\_\_!}
本身不返回值；调用类对象时，在构造函数可用的前提下得到该类实例。 \\
\bottomrule
\end{longtable}

\textbf{对应 C++}

\begin{itemize}
\tightlist
\item
  类：\passthrough{\lstinline!unitree::robot::a2::AudioClient!}
\item
  签名：\passthrough{\lstinline!AudioClient()!}
\item
  绑定策略：\passthrough{\lstinline!未单独标注!}
\item
  声明位置：\passthrough{\lstinline!include/unitree/robot/a2/audio/audio\_client.hpp:15!}
\end{itemize}

\textbf{说明}

\begin{itemize}
\tightlist
\item
  示例是局部调用片段；\passthrough{\lstinline!obj!}
  和参数变量表示已经按上层业务校验并准备好的对象。
\end{itemize}

\textbf{用法}

\begin{lstlisting}[language=Python]
value = AudioClient()
\end{lstlisting}

\hypertarget{unitree-sdk2-cpp-robot-a2-audioclient-init-1}{%
\paragraph{\texorpdfstring{\texttt{unitree\_sdk2\_cpp.robot.a2.AudioClient.init}}{unitree\_sdk2\_cpp.robot.a2.AudioClient.init}}\label{unitree-sdk2-cpp-robot-a2-audioclient-init-1}}

初始化当前 SDK 对象所需的底层通道或服务资源。

\textbf{签名}

\begin{lstlisting}[language=Python]
def init(self) -> None
\end{lstlisting}

\textbf{可用性与安全性}

\textbf{\passthrough{\lstinline!AVAILABLE!} /
\passthrough{\lstinline!INITIALIZATION!}}：当前绑定源码已实现；实际执行条件仍取决于平台和对应
SDK 环境。

\textbf{参数}

无显式参数。实例方法中的 \passthrough{\lstinline!self!} 由 Python
自动传入。

\textbf{返回值}

\begin{longtable}[]{@{}
  >{\raggedright\arraybackslash}p{(\columnwidth - 4\tabcolsep) * \real{0.18}}
  >{\raggedright\arraybackslash}p{(\columnwidth - 4\tabcolsep) * \real{0.20}}
  >{\raggedright\arraybackslash}p{(\columnwidth - 4\tabcolsep) * \real{0.62}}@{}}
\toprule
位置 & 类型 & 含义 \\
\midrule
\endhead
无 & \passthrough{\lstinline!None!} & 该方法不返回 Python
值；失败可能表现为异常或底层状态变化。 \\
\bottomrule
\end{longtable}

\textbf{对应 C++}

\begin{itemize}
\tightlist
\item
  类：\passthrough{\lstinline!unitree::robot::a2::AudioClient!}
\item
  签名：\passthrough{\lstinline!Init()!}
\item
  绑定策略：\passthrough{\lstinline!DIRECT!}
\item
  声明位置：\passthrough{\lstinline!include/unitree/robot/a2/audio/audio\_client.hpp:19!}
\end{itemize}

\textbf{说明}

\begin{itemize}
\tightlist
\item
  示例是局部调用片段；\passthrough{\lstinline!obj!}
  和参数变量表示已经按上层业务校验并准备好的对象。
\end{itemize}

\textbf{用法}

\begin{lstlisting}[language=Python]
obj.init()
\end{lstlisting}

\hypertarget{unitree-sdk2-cpp-robot-a2-audioclient-tts-maker-1}{%
\paragraph{\texorpdfstring{\texttt{unitree\_sdk2\_cpp.robot.a2.AudioClient.tts\_maker}}{unitree\_sdk2\_cpp.robot.a2.AudioClient.tts\_maker}}\label{unitree-sdk2-cpp-robot-a2-audioclient-tts-maker-1}}

对应 C++ SDK 操作
\passthrough{\lstinline!TtsMaker(const std::string \&, int32\_t)!}。上游头文件没有可直接生成的业务说明时，本参考只保证签名映射，精确语义需查目标型号协议。

\textbf{签名}

\begin{lstlisting}[language=Python]
def tts_maker(self, text: str, speaker_id: int) -> int
\end{lstlisting}

\textbf{可用性与安全性}

\textbf{\passthrough{\lstinline!SIGNATURE\_ONLY!} /
\passthrough{\lstinline!HARDWARE\_SIDE\_EFFECT!}}：仅用于补全和静态检查。当前二进制没有该
Python 方法，未来实现还需单独评审硬件副作用。

\textbf{参数}

\begin{longtable}[]{@{}
  >{\raggedright\arraybackslash}p{(\columnwidth - 6\tabcolsep) * \real{0.18}}
  >{\raggedright\arraybackslash}p{(\columnwidth - 6\tabcolsep) * \real{0.24}}
  >{\raggedright\arraybackslash}p{(\columnwidth - 6\tabcolsep) * \real{0.18}}
  >{\raggedright\arraybackslash}p{(\columnwidth - 6\tabcolsep) * \real{0.40}}@{}}
\toprule
名称 & Python 类型 & 默认值 & 含义 \\
\midrule
\endhead
\passthrough{\lstinline!text!} & \passthrough{\lstinline!str!} & 必填 &
要处理的文本内容；编码、长度和语言支持由目标服务决定。 对应 C++ 参数
\passthrough{\lstinline!text: const std::string \&!}。
底层为字符串；长度、编码和允许值由具体协议决定。 \\
\passthrough{\lstinline!speaker\_id!} & \passthrough{\lstinline!int!} &
必填 & 语音合成说话人 ID。有效编号由机器人音频服务和固件决定。 对应 C++
参数 \passthrough{\lstinline!speaker\_id: int32\_t!}。 取值范围为
-2147483648 到 2147483647。 \\
\bottomrule
\end{longtable}

\textbf{返回值}

\begin{longtable}[]{@{}
  >{\raggedright\arraybackslash}p{(\columnwidth - 4\tabcolsep) * \real{0.18}}
  >{\raggedright\arraybackslash}p{(\columnwidth - 4\tabcolsep) * \real{0.20}}
  >{\raggedright\arraybackslash}p{(\columnwidth - 4\tabcolsep) * \real{0.62}}@{}}
\toprule
位置 & 类型 & 含义 \\
\midrule
\endhead
返回值 & \passthrough{\lstinline!int!} & SDK 状态码；Unitree 示例通常以
\passthrough{\lstinline!0!} 表示成功，非零值需按具体服务定义解释。 \\
\bottomrule
\end{longtable}

\textbf{对应 C++}

\begin{itemize}
\tightlist
\item
  类：\passthrough{\lstinline!unitree::robot::a2::AudioClient!}
\item
  签名：\passthrough{\lstinline!TtsMaker(const std::string \&, int32\_t)!}
\item
  绑定策略：\passthrough{\lstinline!DIRECT!}
\item
  声明位置：\passthrough{\lstinline!include/unitree/robot/a2/audio/audio\_client.hpp:31!}
\end{itemize}

\textbf{说明}

\begin{itemize}
\tightlist
\item
  示例是局部调用片段；\passthrough{\lstinline!obj!}
  和参数变量表示已经按上层业务校验并准备好的对象。
\item
  该条目可以被编辑器和类型检查器识别，但当前运行时可能在属性查找阶段直接失败。
\end{itemize}

\textbf{用法}

\begin{lstlisting}[language=Python]
from typing import TYPE_CHECKING

if TYPE_CHECKING:
    def planned_tts_maker(obj: AudioClient, text: str, speaker_id: int) -> int:
        return obj.tts_maker(text=text, speaker_id=speaker_id)
\end{lstlisting}

\begin{quote}
{[}!CAUTION{]} 上面的 \passthrough{\lstinline!TYPE\_CHECKING!}
示例只表达计划调用形态。该分支在普通 Python
运行时为假，不代表当前扩展能执行此接口。
\end{quote}

\hypertarget{unitree-sdk2-cpp-robot-a2-audioclient-get-volume-1}{%
\paragraph{\texorpdfstring{\texttt{unitree\_sdk2\_cpp.robot.a2.AudioClient.get\_volume}}{unitree\_sdk2\_cpp.robot.a2.AudioClient.get\_volume}}\label{unitree-sdk2-cpp-robot-a2-audioclient-get-volume-1}}

查询或检查 音量。

\textbf{签名}

\begin{lstlisting}[language=Python]
def get_volume(self) -> tuple[int, int]
\end{lstlisting}

\textbf{可用性与安全性}

\textbf{\passthrough{\lstinline!AVAILABLE!} /
\passthrough{\lstinline!READ\_ONLY!}}：当前绑定源码已实现只读查询。Robot
Client 的服务端查询通常仍需匹配的 Linux
扩展、DDS、网络、机器人和服务；纯本地 getter 则不一定需要全部条件。

\textbf{参数}

无显式参数。实例方法中的 \passthrough{\lstinline!self!} 由 Python
自动传入。

\textbf{返回值}

\begin{longtable}[]{@{}
  >{\raggedright\arraybackslash}p{(\columnwidth - 4\tabcolsep) * \real{0.18}}
  >{\raggedright\arraybackslash}p{(\columnwidth - 4\tabcolsep) * \real{0.20}}
  >{\raggedright\arraybackslash}p{(\columnwidth - 4\tabcolsep) * \real{0.62}}@{}}
\toprule
位置 & 类型 & 含义 \\
\midrule
\endhead
{[}0{]} \passthrough{\lstinline!status!} & \passthrough{\lstinline!int!}
& SDK 状态码；Unitree 示例通常以 \passthrough{\lstinline!0!}
表示成功，非零值需按具体服务错误码解释。 \\
{[}1{]} \passthrough{\lstinline!volume!} & \passthrough{\lstinline!int!}
& 传给该接口的 音量 参数，Python 类型为
\passthrough{\lstinline!int!}。精确含义、单位、范围和枚举值需查目标型号对应头文件与协议。
对应 C++ 参数 \passthrough{\lstinline!volume: uint8\_t \&!}。 取值范围为
0 到 255。 \\
\bottomrule
\end{longtable}

\textbf{对应 C++}

\begin{itemize}
\tightlist
\item
  类：\passthrough{\lstinline!unitree::robot::a2::AudioClient!}
\item
  签名：\passthrough{\lstinline!GetVolume(uint8\_t \&)!}
\item
  绑定策略：\passthrough{\lstinline!OUTPUT\_WRAPPER!}
\item
  声明位置：\passthrough{\lstinline!include/unitree/robot/a2/audio/audio\_client.hpp:43!}
\end{itemize}

\textbf{说明}

\begin{itemize}
\tightlist
\item
  示例是局部调用片段；\passthrough{\lstinline!obj!}
  和参数变量表示已经按上层业务校验并准备好的对象。
\item
  C++ 的可修改输出引用不会作为 Python 入参出现，而是按 C++
  参数顺序追加到返回元组中。
\end{itemize}

\textbf{用法}

\begin{lstlisting}[language=Python]
status, volume = obj.get_volume()
\end{lstlisting}

\begin{quote}
{[}!NOTE{]} 只读查询仍会访问
DDS/机器人服务。必须检查返回状态码，并处理超时和网络中断。
\end{quote}

\hypertarget{unitree-sdk2-cpp-robot-a2-audioclient-set-volume-1}{%
\paragraph{\texorpdfstring{\texttt{unitree\_sdk2\_cpp.robot.a2.AudioClient.set\_volume}}{unitree\_sdk2\_cpp.robot.a2.AudioClient.set\_volume}}\label{unitree-sdk2-cpp-robot-a2-audioclient-set-volume-1}}

设置 音量。具体副作用和安全边界见下方状态。

\textbf{签名}

\begin{lstlisting}[language=Python]
def set_volume(self, volume: int) -> int
\end{lstlisting}

\textbf{可用性与安全性}

\textbf{\passthrough{\lstinline!SIGNATURE\_ONLY!} /
\passthrough{\lstinline!HARDWARE\_SIDE\_EFFECT!}}：仅用于补全和静态检查。当前二进制没有该
Python 方法，未来实现还需单独评审硬件副作用。

\textbf{参数}

\begin{longtable}[]{@{}
  >{\raggedright\arraybackslash}p{(\columnwidth - 6\tabcolsep) * \real{0.18}}
  >{\raggedright\arraybackslash}p{(\columnwidth - 6\tabcolsep) * \real{0.24}}
  >{\raggedright\arraybackslash}p{(\columnwidth - 6\tabcolsep) * \real{0.18}}
  >{\raggedright\arraybackslash}p{(\columnwidth - 6\tabcolsep) * \real{0.40}}@{}}
\toprule
名称 & Python 类型 & 默认值 & 含义 \\
\midrule
\endhead
\passthrough{\lstinline!volume!} & \passthrough{\lstinline!int!} & 必填
& 传给该接口的 音量 参数，Python 类型为
\passthrough{\lstinline!int!}。精确含义、单位、范围和枚举值需查目标型号对应头文件与协议。
对应 C++ 参数 \passthrough{\lstinline!volume: uint8\_t!}。 取值范围为 0
到 255。 \\
\bottomrule
\end{longtable}

\textbf{返回值}

\begin{longtable}[]{@{}
  >{\raggedright\arraybackslash}p{(\columnwidth - 4\tabcolsep) * \real{0.18}}
  >{\raggedright\arraybackslash}p{(\columnwidth - 4\tabcolsep) * \real{0.20}}
  >{\raggedright\arraybackslash}p{(\columnwidth - 4\tabcolsep) * \real{0.62}}@{}}
\toprule
位置 & 类型 & 含义 \\
\midrule
\endhead
返回值 & \passthrough{\lstinline!int!} & SDK 状态码；Unitree 示例通常以
\passthrough{\lstinline!0!} 表示成功，非零值需按具体服务定义解释。 \\
\bottomrule
\end{longtable}

\textbf{对应 C++}

\begin{itemize}
\tightlist
\item
  类：\passthrough{\lstinline!unitree::robot::a2::AudioClient!}
\item
  签名：\passthrough{\lstinline!SetVolume(uint8\_t)!}
\item
  绑定策略：\passthrough{\lstinline!DIRECT!}
\item
  声明位置：\passthrough{\lstinline!include/unitree/robot/a2/audio/audio\_client.hpp:57!}
\end{itemize}

\textbf{说明}

\begin{itemize}
\tightlist
\item
  示例是局部调用片段；\passthrough{\lstinline!obj!}
  和参数变量表示已经按上层业务校验并准备好的对象。
\item
  该条目可以被编辑器和类型检查器识别，但当前运行时可能在属性查找阶段直接失败。
\end{itemize}

\textbf{用法}

\begin{lstlisting}[language=Python]
from typing import TYPE_CHECKING

if TYPE_CHECKING:
    def planned_set_volume(obj: AudioClient, volume: int) -> int:
        return obj.set_volume(volume=volume)
\end{lstlisting}

\begin{quote}
{[}!CAUTION{]} 上面的 \passthrough{\lstinline!TYPE\_CHECKING!}
示例只表达计划调用形态。该分支在普通 Python
运行时为假，不代表当前扩展能执行此接口。
\end{quote}

\hypertarget{unitree-sdk2-cpp-robot-a2-audioclient-play-stream-1}{%
\paragraph{\texorpdfstring{\texttt{unitree\_sdk2\_cpp.robot.a2.AudioClient.play\_stream}}{unitree\_sdk2\_cpp.robot.a2.AudioClient.play\_stream}}\label{unitree-sdk2-cpp-robot-a2-audioclient-play-stream-1}}

对应 C++ SDK 操作
\passthrough{\lstinline!PlayStream(std::string, std::string, std::vector<uint8\_t>)!}。上游头文件没有可直接生成的业务说明时，本参考只保证签名映射，精确语义需查目标型号协议。

\textbf{签名}

\begin{lstlisting}[language=Python]
def play_stream(self, app_name: str, stream_id: str, pcm_data: Sequence[int]) -> int
\end{lstlisting}

\textbf{可用性与安全性}

\textbf{\passthrough{\lstinline!SIGNATURE\_ONLY!} /
\passthrough{\lstinline!HARDWARE\_SIDE\_EFFECT!}}：仅用于补全和静态检查。当前二进制没有该
Python 方法，未来实现还需单独评审硬件副作用。

\textbf{参数}

\begin{longtable}[]{@{}
  >{\raggedright\arraybackslash}p{(\columnwidth - 6\tabcolsep) * \real{0.18}}
  >{\raggedright\arraybackslash}p{(\columnwidth - 6\tabcolsep) * \real{0.24}}
  >{\raggedright\arraybackslash}p{(\columnwidth - 6\tabcolsep) * \real{0.18}}
  >{\raggedright\arraybackslash}p{(\columnwidth - 6\tabcolsep) * \real{0.40}}@{}}
\toprule
名称 & Python 类型 & 默认值 & 含义 \\
\midrule
\endhead
\passthrough{\lstinline!app\_name!} & \passthrough{\lstinline!str!} &
必填 &
应用名称，用于标识音频流或播放会话。允许值和命名规则以目标服务协议为准。
对应 C++ 参数 \passthrough{\lstinline!app\_name: std::string!}。
底层为字符串；长度、编码和允许值由具体协议决定。 \\
\passthrough{\lstinline!stream\_id!} & \passthrough{\lstinline!str!} &
必填 & 音频流标识符，用于区分同一应用下的播放流。 对应 C++ 参数
\passthrough{\lstinline!stream\_id: std::string!}。
底层为字符串；长度、编码和允许值由具体协议决定。 \\
\passthrough{\lstinline!pcm\_data!} &
\passthrough{\lstinline!Sequence[int]!} & 必填 & 传给该接口的
\passthrough{\lstinline!pcm!} 数据 参数，Python 类型为
\passthrough{\lstinline!Sequence[int]!}。精确含义、单位、范围和枚举值需查目标型号对应头文件与协议。
对应 C++ 参数
\passthrough{\lstinline!pcm\_data: std::vector<uint8\_t>!}。
底层是可变长度 vector；元素约束：取值范围为 0 到 255。 \\
\bottomrule
\end{longtable}

\textbf{返回值}

\begin{longtable}[]{@{}
  >{\raggedright\arraybackslash}p{(\columnwidth - 4\tabcolsep) * \real{0.18}}
  >{\raggedright\arraybackslash}p{(\columnwidth - 4\tabcolsep) * \real{0.20}}
  >{\raggedright\arraybackslash}p{(\columnwidth - 4\tabcolsep) * \real{0.62}}@{}}
\toprule
位置 & 类型 & 含义 \\
\midrule
\endhead
返回值 & \passthrough{\lstinline!int!} & SDK 状态码；Unitree 示例通常以
\passthrough{\lstinline!0!} 表示成功，非零值需按具体服务定义解释。 \\
\bottomrule
\end{longtable}

\textbf{对应 C++}

\begin{itemize}
\tightlist
\item
  类：\passthrough{\lstinline!unitree::robot::a2::AudioClient!}
\item
  签名：\passthrough{\lstinline!PlayStream(std::string, std::string, std::vector<uint8\_t>)!}
\item
  绑定策略：\passthrough{\lstinline!DIRECT!}
\item
  声明位置：\passthrough{\lstinline!include/unitree/robot/a2/audio/audio\_client.hpp:68!}
\end{itemize}

\textbf{说明}

\begin{itemize}
\tightlist
\item
  示例是局部调用片段；\passthrough{\lstinline!obj!}
  和参数变量表示已经按上层业务校验并准备好的对象。
\item
  该条目可以被编辑器和类型检查器识别，但当前运行时可能在属性查找阶段直接失败。
\end{itemize}

\textbf{用法}

\begin{lstlisting}[language=Python]
from typing import TYPE_CHECKING

if TYPE_CHECKING:
    def planned_play_stream(obj: AudioClient, app_name: str, stream_id: str, pcm_data: Sequence[int]) -> int:
        return obj.play_stream(app_name=app_name, stream_id=stream_id, pcm_data=pcm_data)
\end{lstlisting}

\begin{quote}
{[}!CAUTION{]} 上面的 \passthrough{\lstinline!TYPE\_CHECKING!}
示例只表达计划调用形态。该分支在普通 Python
运行时为假，不代表当前扩展能执行此接口。
\end{quote}

\hypertarget{unitree-sdk2-cpp-robot-a2-audioclient-play-stop-1}{%
\paragraph{\texorpdfstring{\texttt{unitree\_sdk2\_cpp.robot.a2.AudioClient.play\_stop}}{unitree\_sdk2\_cpp.robot.a2.AudioClient.play\_stop}}\label{unitree-sdk2-cpp-robot-a2-audioclient-play-stop-1}}

对应 C++ SDK 操作
\passthrough{\lstinline!PlayStop(std::string)!}。上游头文件没有可直接生成的业务说明时，本参考只保证签名映射，精确语义需查目标型号协议。

\textbf{签名}

\begin{lstlisting}[language=Python]
def play_stop(self, app_name: str) -> int
\end{lstlisting}

\textbf{可用性与安全性}

\textbf{\passthrough{\lstinline!SIGNATURE\_ONLY!} /
\passthrough{\lstinline!HARDWARE\_SIDE\_EFFECT!}}：仅用于补全和静态检查。当前二进制没有该
Python 方法，未来实现还需单独评审硬件副作用。

\textbf{参数}

\begin{longtable}[]{@{}
  >{\raggedright\arraybackslash}p{(\columnwidth - 6\tabcolsep) * \real{0.18}}
  >{\raggedright\arraybackslash}p{(\columnwidth - 6\tabcolsep) * \real{0.24}}
  >{\raggedright\arraybackslash}p{(\columnwidth - 6\tabcolsep) * \real{0.18}}
  >{\raggedright\arraybackslash}p{(\columnwidth - 6\tabcolsep) * \real{0.40}}@{}}
\toprule
名称 & Python 类型 & 默认值 & 含义 \\
\midrule
\endhead
\passthrough{\lstinline!app\_name!} & \passthrough{\lstinline!str!} &
必填 &
应用名称，用于标识音频流或播放会话。允许值和命名规则以目标服务协议为准。
对应 C++ 参数 \passthrough{\lstinline!app\_name: std::string!}。
底层为字符串；长度、编码和允许值由具体协议决定。 \\
\bottomrule
\end{longtable}

\textbf{返回值}

\begin{longtable}[]{@{}
  >{\raggedright\arraybackslash}p{(\columnwidth - 4\tabcolsep) * \real{0.18}}
  >{\raggedright\arraybackslash}p{(\columnwidth - 4\tabcolsep) * \real{0.20}}
  >{\raggedright\arraybackslash}p{(\columnwidth - 4\tabcolsep) * \real{0.62}}@{}}
\toprule
位置 & 类型 & 含义 \\
\midrule
\endhead
返回值 & \passthrough{\lstinline!int!} & SDK 状态码；Unitree 示例通常以
\passthrough{\lstinline!0!} 表示成功，非零值需按具体服务定义解释。 \\
\bottomrule
\end{longtable}

\textbf{对应 C++}

\begin{itemize}
\tightlist
\item
  类：\passthrough{\lstinline!unitree::robot::a2::AudioClient!}
\item
  签名：\passthrough{\lstinline!PlayStop(std::string)!}
\item
  绑定策略：\passthrough{\lstinline!DIRECT!}
\item
  声明位置：\passthrough{\lstinline!include/unitree/robot/a2/audio/audio\_client.hpp:80!}
\end{itemize}

\textbf{说明}

\begin{itemize}
\tightlist
\item
  示例是局部调用片段；\passthrough{\lstinline!obj!}
  和参数变量表示已经按上层业务校验并准备好的对象。
\item
  该条目可以被编辑器和类型检查器识别，但当前运行时可能在属性查找阶段直接失败。
\end{itemize}

\textbf{用法}

\begin{lstlisting}[language=Python]
from typing import TYPE_CHECKING

if TYPE_CHECKING:
    def planned_play_stop(obj: AudioClient, app_name: str) -> int:
        return obj.play_stop(app_name=app_name)
\end{lstlisting}

\begin{quote}
{[}!CAUTION{]} 上面的 \passthrough{\lstinline!TYPE\_CHECKING!}
示例只表达计划调用形态。该分支在普通 Python
运行时为假，不代表当前扩展能执行此接口。
\end{quote}

\hypertarget{unitree-sdk2-cpp-robot-a2-audioclient-led-control-1}{%
\paragraph{\texorpdfstring{\texttt{unitree\_sdk2\_cpp.robot.a2.AudioClient.led\_control}}{unitree\_sdk2\_cpp.robot.a2.AudioClient.led\_control}}\label{unitree-sdk2-cpp-robot-a2-audioclient-led-control-1}}

对应 C++ SDK 操作
\passthrough{\lstinline!LedControl(uint8\_t, uint8\_t, uint8\_t)!}。上游头文件没有可直接生成的业务说明时，本参考只保证签名映射，精确语义需查目标型号协议。

\textbf{签名}

\begin{lstlisting}[language=Python]
def led_control(self, r: int, g: int, b: int) -> int
\end{lstlisting}

\textbf{可用性与安全性}

\textbf{\passthrough{\lstinline!SIGNATURE\_ONLY!} /
\passthrough{\lstinline!HARDWARE\_SIDE\_EFFECT!}}：仅用于补全和静态检查。当前二进制没有该
Python 方法，未来实现还需单独评审硬件副作用。

\textbf{参数}

\begin{longtable}[]{@{}
  >{\raggedright\arraybackslash}p{(\columnwidth - 6\tabcolsep) * \real{0.18}}
  >{\raggedright\arraybackslash}p{(\columnwidth - 6\tabcolsep) * \real{0.24}}
  >{\raggedright\arraybackslash}p{(\columnwidth - 6\tabcolsep) * \real{0.18}}
  >{\raggedright\arraybackslash}p{(\columnwidth - 6\tabcolsep) * \real{0.40}}@{}}
\toprule
名称 & Python 类型 & 默认值 & 含义 \\
\midrule
\endhead
\passthrough{\lstinline!r!} & \passthrough{\lstinline!int!} & 必填 &
传给该接口的 \passthrough{\lstinline!r!} 参数，Python 类型为
\passthrough{\lstinline!int!}。精确含义、单位、范围和枚举值需查目标型号对应头文件与协议。
对应 C++ 参数 \passthrough{\lstinline!R: uint8\_t!}。 取值范围为 0 到
255。 \\
\passthrough{\lstinline!g!} & \passthrough{\lstinline!int!} & 必填 &
传给该接口的 \passthrough{\lstinline!g!} 参数，Python 类型为
\passthrough{\lstinline!int!}。精确含义、单位、范围和枚举值需查目标型号对应头文件与协议。
对应 C++ 参数 \passthrough{\lstinline!G: uint8\_t!}。 取值范围为 0 到
255。 \\
\passthrough{\lstinline!b!} & \passthrough{\lstinline!int!} & 必填 &
传给该接口的 \passthrough{\lstinline!b!} 参数，Python 类型为
\passthrough{\lstinline!int!}。精确含义、单位、范围和枚举值需查目标型号对应头文件与协议。
对应 C++ 参数 \passthrough{\lstinline!B: uint8\_t!}。 取值范围为 0 到
255。 \\
\bottomrule
\end{longtable}

\textbf{返回值}

\begin{longtable}[]{@{}
  >{\raggedright\arraybackslash}p{(\columnwidth - 4\tabcolsep) * \real{0.18}}
  >{\raggedright\arraybackslash}p{(\columnwidth - 4\tabcolsep) * \real{0.20}}
  >{\raggedright\arraybackslash}p{(\columnwidth - 4\tabcolsep) * \real{0.62}}@{}}
\toprule
位置 & 类型 & 含义 \\
\midrule
\endhead
返回值 & \passthrough{\lstinline!int!} & SDK 状态码；Unitree 示例通常以
\passthrough{\lstinline!0!} 表示成功，非零值需按具体服务定义解释。 \\
\bottomrule
\end{longtable}

\textbf{对应 C++}

\begin{itemize}
\tightlist
\item
  类：\passthrough{\lstinline!unitree::robot::a2::AudioClient!}
\item
  签名：\passthrough{\lstinline!LedControl(uint8\_t, uint8\_t, uint8\_t)!}
\item
  绑定策略：\passthrough{\lstinline!DIRECT!}
\item
  声明位置：\passthrough{\lstinline!include/unitree/robot/a2/audio/audio\_client.hpp:91!}
\end{itemize}

\textbf{说明}

\begin{itemize}
\tightlist
\item
  示例是局部调用片段；\passthrough{\lstinline!obj!}
  和参数变量表示已经按上层业务校验并准备好的对象。
\item
  该条目可以被编辑器和类型检查器识别，但当前运行时可能在属性查找阶段直接失败。
\end{itemize}

\textbf{用法}

\begin{lstlisting}[language=Python]
from typing import TYPE_CHECKING

if TYPE_CHECKING:
    def planned_led_control(obj: AudioClient, r: int, g: int, b: int) -> int:
        return obj.led_control(r=r, g=g, b=b)
\end{lstlisting}

\begin{quote}
{[}!CAUTION{]} 上面的 \passthrough{\lstinline!TYPE\_CHECKING!}
示例只表达计划调用形态。该分支在普通 Python
运行时为假，不代表当前扩展能执行此接口。
\end{quote}

\hypertarget{unitree-sdk2-cpp-robot-a2-ledcontrolparameter}{%
\subsubsection{\texorpdfstring{\texttt{unitree\_sdk2\_cpp.robot.a2.LedControlParameter}}{unitree\_sdk2\_cpp.robot.a2.LedControlParameter}}\label{unitree-sdk2-cpp-robot-a2-ledcontrolparameter}}

SDK 请求参数值类型；当前可能仅提供设计期签名。

\textbf{导入}

\begin{lstlisting}[language=Python]
from unitree_sdk2_cpp.robot.a2 import LedControlParameter
\end{lstlisting}

\textbf{基类}：\passthrough{\lstinline!object!}

\textbf{公开属性}

\begin{longtable}[]{@{}
  >{\raggedright\arraybackslash}p{(\columnwidth - 6\tabcolsep) * \real{0.18}}
  >{\raggedright\arraybackslash}p{(\columnwidth - 6\tabcolsep) * \real{0.24}}
  >{\raggedright\arraybackslash}p{(\columnwidth - 6\tabcolsep) * \real{0.18}}
  >{\raggedright\arraybackslash}p{(\columnwidth - 6\tabcolsep) * \real{0.40}}@{}}
\toprule
名称 & Python 类型 & 含义 & 用法 \\
\midrule
\endhead
\passthrough{\lstinline!r!} & \passthrough{\lstinline!int!} &
传给该接口的 \passthrough{\lstinline!r!} 参数，Python 类型为
\passthrough{\lstinline!int!}。精确含义、单位、范围和枚举值需查目标型号对应头文件与协议。
& \passthrough{\lstinline!value = obj.r!} /
\passthrough{\lstinline!obj.r = value!} \\
\passthrough{\lstinline!g!} & \passthrough{\lstinline!int!} &
传给该接口的 \passthrough{\lstinline!g!} 参数，Python 类型为
\passthrough{\lstinline!int!}。精确含义、单位、范围和枚举值需查目标型号对应头文件与协议。
& \passthrough{\lstinline!value = obj.g!} /
\passthrough{\lstinline!obj.g = value!} \\
\passthrough{\lstinline!b!} & \passthrough{\lstinline!int!} &
传给该接口的 \passthrough{\lstinline!b!} 参数，Python 类型为
\passthrough{\lstinline!int!}。精确含义、单位、范围和枚举值需查目标型号对应头文件与协议。
& \passthrough{\lstinline!value = obj.b!} /
\passthrough{\lstinline!obj.b = value!} \\
\bottomrule
\end{longtable}

\textbf{方法索引}

\begin{longtable}[]{@{}
  >{\raggedright\arraybackslash}p{(\columnwidth - 6\tabcolsep) * \real{0.18}}
  >{\raggedright\arraybackslash}p{(\columnwidth - 6\tabcolsep) * \real{0.24}}
  >{\raggedright\arraybackslash}p{(\columnwidth - 6\tabcolsep) * \real{0.18}}
  >{\raggedright\arraybackslash}p{(\columnwidth - 6\tabcolsep) * \real{0.40}}@{}}
\toprule
方法 & Python 签名 & 状态 & 安全分类 \\
\midrule
\endhead
\protect\hyperlink{unitree-sdk2-cpp-robot-a2-ledcontrolparameter-dunder-init-1}{\passthrough{\lstinline!\_\_init\_\_!}}
& \passthrough{\lstinline!def \_\_init\_\_(self) -> None!} &
\passthrough{\lstinline!SIGNATURE\_ONLY!} &
\passthrough{\lstinline!CONSTRUCTION!} \\
\protect\hyperlink{unitree-sdk2-cpp-robot-a2-ledcontrolparameter-from-json-1}{\passthrough{\lstinline!from\_json!}}
&
\passthrough{\lstinline!def from\_json(self, json: dict[str, Any]) -> None!}
& \passthrough{\lstinline!SIGNATURE\_ONLY!} &
\passthrough{\lstinline!UNCLASSIFIED!} \\
\protect\hyperlink{unitree-sdk2-cpp-robot-a2-ledcontrolparameter-to-json-1}{\passthrough{\lstinline!to\_json!}}
&
\passthrough{\lstinline!def to\_json(self, json: dict[str, Any]) -> None!}
& \passthrough{\lstinline!SIGNATURE\_ONLY!} &
\passthrough{\lstinline!UNCLASSIFIED!} \\
\bottomrule
\end{longtable}

\hypertarget{unitree-sdk2-cpp-robot-a2-ledcontrolparameter-dunder-init-1}{%
\paragraph{\texorpdfstring{\texttt{unitree\_sdk2\_cpp.robot.a2.LedControlParameter.\_\_init\_\_}}{unitree\_sdk2\_cpp.robot.a2.LedControlParameter.\_\_init\_\_}}\label{unitree-sdk2-cpp-robot-a2-ledcontrolparameter-dunder-init-1}}

初始化 \passthrough{\lstinline!LedControlParameter!}
实例。是否能实际构造取决于下方可用性状态。

\textbf{签名}

\begin{lstlisting}[language=Python]
def __init__(self) -> None
\end{lstlisting}

\textbf{可用性与安全性}

\textbf{\passthrough{\lstinline!SIGNATURE\_ONLY!} /
\passthrough{\lstinline!CONSTRUCTION!}}：当前只有设计期签名，不能假设已存在于运行时扩展。

\textbf{参数}

无显式参数。实例方法中的 \passthrough{\lstinline!self!} 由 Python
自动传入。

\textbf{返回值}

\begin{longtable}[]{@{}
  >{\raggedright\arraybackslash}p{(\columnwidth - 4\tabcolsep) * \real{0.18}}
  >{\raggedright\arraybackslash}p{(\columnwidth - 4\tabcolsep) * \real{0.20}}
  >{\raggedright\arraybackslash}p{(\columnwidth - 4\tabcolsep) * \real{0.62}}@{}}
\toprule
位置 & 类型 & 含义 \\
\midrule
\endhead
无 & \passthrough{\lstinline!None!} &
\passthrough{\lstinline!\_\_init\_\_!}
本身不返回值；调用类对象时，在构造函数可用的前提下得到该类实例。 \\
\bottomrule
\end{longtable}

\textbf{对应 C++}

\begin{itemize}
\tightlist
\item
  类：\passthrough{\lstinline!unitree::robot::a2::LedControlParameter!}
\item
  签名：\passthrough{\lstinline!LedControlParameter()!}
\item
  绑定策略：\passthrough{\lstinline!未单独标注!}
\item
  声明位置：\passthrough{\lstinline!include/unitree/robot/a2/audio/audio\_api.hpp:75!}
\end{itemize}

\textbf{说明}

\begin{itemize}
\tightlist
\item
  示例是局部调用片段；\passthrough{\lstinline!obj!}
  和参数变量表示已经按上层业务校验并准备好的对象。
\item
  该条目可以被编辑器和类型检查器识别，但当前运行时可能在属性查找阶段直接失败。
\end{itemize}

\textbf{用法}

\begin{lstlisting}[language=Python]
from typing import TYPE_CHECKING

if TYPE_CHECKING:
    def planned_construct() -> LedControlParameter:
        return LedControlParameter()
\end{lstlisting}

\begin{quote}
{[}!CAUTION{]} 上面的 \passthrough{\lstinline!TYPE\_CHECKING!}
示例只表达计划调用形态。该分支在普通 Python
运行时为假，不代表当前扩展能执行此接口。
\end{quote}

\hypertarget{unitree-sdk2-cpp-robot-a2-ledcontrolparameter-from-json-1}{%
\paragraph{\texorpdfstring{\texttt{unitree\_sdk2\_cpp.robot.a2.LedControlParameter.from\_json}}{unitree\_sdk2\_cpp.robot.a2.LedControlParameter.from\_json}}\label{unitree-sdk2-cpp-robot-a2-ledcontrolparameter-from-json-1}}

计划从 JSON 风格字典读取字段并更新当前 SDK 值对象。

\textbf{签名}

\begin{lstlisting}[language=Python]
def from_json(self, json: dict[str, Any]) -> None
\end{lstlisting}

\textbf{可用性与安全性}

\textbf{\passthrough{\lstinline!SIGNATURE\_ONLY!} /
\passthrough{\lstinline!UNCLASSIFIED!}}：当前只有设计期签名，不能假设已存在于运行时扩展。

\textbf{参数}

\begin{longtable}[]{@{}
  >{\raggedright\arraybackslash}p{(\columnwidth - 6\tabcolsep) * \real{0.18}}
  >{\raggedright\arraybackslash}p{(\columnwidth - 6\tabcolsep) * \real{0.24}}
  >{\raggedright\arraybackslash}p{(\columnwidth - 6\tabcolsep) * \real{0.18}}
  >{\raggedright\arraybackslash}p{(\columnwidth - 6\tabcolsep) * \real{0.40}}@{}}
\toprule
名称 & Python 类型 & 默认值 & 含义 \\
\midrule
\endhead
\passthrough{\lstinline!json!} &
\passthrough{\lstinline!dict[str, Any]!} & 必填 & JSON
风格字典，键和值必须满足对应 C++ \passthrough{\lstinline!JsonMap!}
数据结构的约束。 对应 C++ 参数
\passthrough{\lstinline!json: common::JsonMap \&!}。 \\
\bottomrule
\end{longtable}

\textbf{返回值}

\begin{longtable}[]{@{}
  >{\raggedright\arraybackslash}p{(\columnwidth - 4\tabcolsep) * \real{0.18}}
  >{\raggedright\arraybackslash}p{(\columnwidth - 4\tabcolsep) * \real{0.20}}
  >{\raggedright\arraybackslash}p{(\columnwidth - 4\tabcolsep) * \real{0.62}}@{}}
\toprule
位置 & 类型 & 含义 \\
\midrule
\endhead
无 & \passthrough{\lstinline!None!} & 该方法不返回 Python
值；失败可能表现为异常或底层状态变化。 \\
\bottomrule
\end{longtable}

\textbf{对应 C++}

\begin{itemize}
\tightlist
\item
  类：\passthrough{\lstinline!unitree::robot::a2::LedControlParameter!}
\item
  签名：\passthrough{\lstinline!fromJson(common::JsonMap \&)!}
\item
  绑定策略：\passthrough{\lstinline!SIGNATURE\_PREVIEW!}
\item
  声明位置：\passthrough{\lstinline!include/unitree/robot/a2/audio/audio\_api.hpp:78!}
\end{itemize}

\textbf{说明}

\begin{itemize}
\tightlist
\item
  示例是局部调用片段；\passthrough{\lstinline!obj!}
  和参数变量表示已经按上层业务校验并准备好的对象。
\item
  该条目可以被编辑器和类型检查器识别，但当前运行时可能在属性查找阶段直接失败。
\end{itemize}

\textbf{用法}

\begin{lstlisting}[language=Python]
from typing import TYPE_CHECKING

if TYPE_CHECKING:
    def planned_from_json(obj: LedControlParameter, json: dict[str, Any]) -> None:
        obj.from_json(json=json)
\end{lstlisting}

\begin{quote}
{[}!CAUTION{]} 上面的 \passthrough{\lstinline!TYPE\_CHECKING!}
示例只表达计划调用形态。该分支在普通 Python
运行时为假，不代表当前扩展能执行此接口。
\end{quote}

\hypertarget{unitree-sdk2-cpp-robot-a2-ledcontrolparameter-to-json-1}{%
\paragraph{\texorpdfstring{\texttt{unitree\_sdk2\_cpp.robot.a2.LedControlParameter.to\_json}}{unitree\_sdk2\_cpp.robot.a2.LedControlParameter.to\_json}}\label{unitree-sdk2-cpp-robot-a2-ledcontrolparameter-to-json-1}}

计划把当前 SDK 值对象写入 JSON 风格字典。

\textbf{签名}

\begin{lstlisting}[language=Python]
def to_json(self, json: dict[str, Any]) -> None
\end{lstlisting}

\textbf{可用性与安全性}

\textbf{\passthrough{\lstinline!SIGNATURE\_ONLY!} /
\passthrough{\lstinline!UNCLASSIFIED!}}：当前只有设计期签名，不能假设已存在于运行时扩展。

\textbf{参数}

\begin{longtable}[]{@{}
  >{\raggedright\arraybackslash}p{(\columnwidth - 6\tabcolsep) * \real{0.18}}
  >{\raggedright\arraybackslash}p{(\columnwidth - 6\tabcolsep) * \real{0.24}}
  >{\raggedright\arraybackslash}p{(\columnwidth - 6\tabcolsep) * \real{0.18}}
  >{\raggedright\arraybackslash}p{(\columnwidth - 6\tabcolsep) * \real{0.40}}@{}}
\toprule
名称 & Python 类型 & 默认值 & 含义 \\
\midrule
\endhead
\passthrough{\lstinline!json!} &
\passthrough{\lstinline!dict[str, Any]!} & 必填 & JSON
风格字典，键和值必须满足对应 C++ \passthrough{\lstinline!JsonMap!}
数据结构的约束。 对应 C++ 参数
\passthrough{\lstinline!json: common::JsonMap \&!}。 \\
\bottomrule
\end{longtable}

\textbf{返回值}

\begin{longtable}[]{@{}
  >{\raggedright\arraybackslash}p{(\columnwidth - 4\tabcolsep) * \real{0.18}}
  >{\raggedright\arraybackslash}p{(\columnwidth - 4\tabcolsep) * \real{0.20}}
  >{\raggedright\arraybackslash}p{(\columnwidth - 4\tabcolsep) * \real{0.62}}@{}}
\toprule
位置 & 类型 & 含义 \\
\midrule
\endhead
无 & \passthrough{\lstinline!None!} & 该方法不返回 Python
值；失败可能表现为异常或底层状态变化。 \\
\bottomrule
\end{longtable}

\textbf{对应 C++}

\begin{itemize}
\tightlist
\item
  类：\passthrough{\lstinline!unitree::robot::a2::LedControlParameter!}
\item
  签名：\passthrough{\lstinline!toJson(common::JsonMap \&) const!}
\item
  绑定策略：\passthrough{\lstinline!SIGNATURE\_PREVIEW!}
\item
  声明位置：\passthrough{\lstinline!include/unitree/robot/a2/audio/audio\_api.hpp:80!}
\end{itemize}

\textbf{说明}

\begin{itemize}
\tightlist
\item
  示例是局部调用片段；\passthrough{\lstinline!obj!}
  和参数变量表示已经按上层业务校验并准备好的对象。
\item
  该条目可以被编辑器和类型检查器识别，但当前运行时可能在属性查找阶段直接失败。
\end{itemize}

\textbf{用法}

\begin{lstlisting}[language=Python]
from typing import TYPE_CHECKING

if TYPE_CHECKING:
    def planned_to_json(obj: LedControlParameter, json: dict[str, Any]) -> None:
        obj.to_json(json=json)
\end{lstlisting}

\begin{quote}
{[}!CAUTION{]} 上面的 \passthrough{\lstinline!TYPE\_CHECKING!}
示例只表达计划调用形态。该分支在普通 Python
运行时为假，不代表当前扩展能执行此接口。
\end{quote}

\hypertarget{unitree-sdk2-cpp-robot-a2-pathpoint}{%
\subsubsection{\texorpdfstring{\texttt{unitree\_sdk2\_cpp.robot.a2.PathPoint}}{unitree\_sdk2\_cpp.robot.a2.PathPoint}}\label{unitree-sdk2-cpp-robot-a2-pathpoint}}

SDK 数据值类型；公开字段和转换方法见下文。

\textbf{导入}

\begin{lstlisting}[language=Python]
from unitree_sdk2_cpp.robot.a2 import PathPoint
\end{lstlisting}

\textbf{基类}：\passthrough{\lstinline!object!}

\textbf{公开属性}

\begin{longtable}[]{@{}
  >{\raggedright\arraybackslash}p{(\columnwidth - 6\tabcolsep) * \real{0.18}}
  >{\raggedright\arraybackslash}p{(\columnwidth - 6\tabcolsep) * \real{0.24}}
  >{\raggedright\arraybackslash}p{(\columnwidth - 6\tabcolsep) * \real{0.18}}
  >{\raggedright\arraybackslash}p{(\columnwidth - 6\tabcolsep) * \real{0.40}}@{}}
\toprule
名称 & Python 类型 & 含义 & 用法 \\
\midrule
\endhead
\passthrough{\lstinline!t\_from\_start!} &
\passthrough{\lstinline!float!} & 传给该接口的
\passthrough{\lstinline!t!} \passthrough{\lstinline!from!}
\passthrough{\lstinline!start!} 参数，Python 类型为
\passthrough{\lstinline!float!}。精确含义、单位、范围和枚举值需查目标型号对应头文件与协议。
& \passthrough{\lstinline!value = obj.t\_from\_start!} /
\passthrough{\lstinline!obj.t\_from\_start = value!} \\
\passthrough{\lstinline!x!} & \passthrough{\lstinline!float!} & X
方向位置或分量参数。参考系、单位和范围必须查当前方法及目标型号协议。 &
\passthrough{\lstinline!value = obj.x!} /
\passthrough{\lstinline!obj.x = value!} \\
\passthrough{\lstinline!y!} & \passthrough{\lstinline!float!} & Y
方向位置或分量参数。参考系、单位和范围必须查当前方法及目标型号协议。 &
\passthrough{\lstinline!value = obj.y!} /
\passthrough{\lstinline!obj.y = value!} \\
\passthrough{\lstinline!yaw!} & \passthrough{\lstinline!float!} &
偏航角或偏航目标参数。参考系、单位和范围必须查目标型号协议。 &
\passthrough{\lstinline!value = obj.yaw!} /
\passthrough{\lstinline!obj.yaw = value!} \\
\bottomrule
\end{longtable}

\hypertarget{unitree-sdk2-cpp-robot-a2-playstopparameter}{%
\subsubsection{\texorpdfstring{\texttt{unitree\_sdk2\_cpp.robot.a2.PlayStopParameter}}{unitree\_sdk2\_cpp.robot.a2.PlayStopParameter}}\label{unitree-sdk2-cpp-robot-a2-playstopparameter}}

SDK 请求参数值类型；当前可能仅提供设计期签名。

\textbf{导入}

\begin{lstlisting}[language=Python]
from unitree_sdk2_cpp.robot.a2 import PlayStopParameter
\end{lstlisting}

\textbf{基类}：\passthrough{\lstinline!object!}

\textbf{公开属性}

\begin{longtable}[]{@{}
  >{\raggedright\arraybackslash}p{(\columnwidth - 6\tabcolsep) * \real{0.18}}
  >{\raggedright\arraybackslash}p{(\columnwidth - 6\tabcolsep) * \real{0.24}}
  >{\raggedright\arraybackslash}p{(\columnwidth - 6\tabcolsep) * \real{0.18}}
  >{\raggedright\arraybackslash}p{(\columnwidth - 6\tabcolsep) * \real{0.40}}@{}}
\toprule
名称 & Python 类型 & 含义 & 用法 \\
\midrule
\endhead
\passthrough{\lstinline!app\_name!} & \passthrough{\lstinline!str!} &
应用名称，用于标识音频流或播放会话。允许值和命名规则以目标服务协议为准。
& \passthrough{\lstinline!value = obj.app\_name!} /
\passthrough{\lstinline!obj.app\_name = value!} \\
\bottomrule
\end{longtable}

\textbf{方法索引}

\begin{longtable}[]{@{}
  >{\raggedright\arraybackslash}p{(\columnwidth - 6\tabcolsep) * \real{0.18}}
  >{\raggedright\arraybackslash}p{(\columnwidth - 6\tabcolsep) * \real{0.24}}
  >{\raggedright\arraybackslash}p{(\columnwidth - 6\tabcolsep) * \real{0.18}}
  >{\raggedright\arraybackslash}p{(\columnwidth - 6\tabcolsep) * \real{0.40}}@{}}
\toprule
方法 & Python 签名 & 状态 & 安全分类 \\
\midrule
\endhead
\protect\hyperlink{unitree-sdk2-cpp-robot-a2-playstopparameter-dunder-init-1}{\passthrough{\lstinline!\_\_init\_\_!}}
& \passthrough{\lstinline!def \_\_init\_\_(self) -> None!} &
\passthrough{\lstinline!SIGNATURE\_ONLY!} &
\passthrough{\lstinline!CONSTRUCTION!} \\
\protect\hyperlink{unitree-sdk2-cpp-robot-a2-playstopparameter-from-json-1}{\passthrough{\lstinline!from\_json!}}
&
\passthrough{\lstinline!def from\_json(self, json: dict[str, Any]) -> None!}
& \passthrough{\lstinline!SIGNATURE\_ONLY!} &
\passthrough{\lstinline!UNCLASSIFIED!} \\
\protect\hyperlink{unitree-sdk2-cpp-robot-a2-playstopparameter-to-json-1}{\passthrough{\lstinline!to\_json!}}
&
\passthrough{\lstinline!def to\_json(self, json: dict[str, Any]) -> None!}
& \passthrough{\lstinline!SIGNATURE\_ONLY!} &
\passthrough{\lstinline!UNCLASSIFIED!} \\
\bottomrule
\end{longtable}

\hypertarget{unitree-sdk2-cpp-robot-a2-playstopparameter-dunder-init-1}{%
\paragraph{\texorpdfstring{\texttt{unitree\_sdk2\_cpp.robot.a2.PlayStopParameter.\_\_init\_\_}}{unitree\_sdk2\_cpp.robot.a2.PlayStopParameter.\_\_init\_\_}}\label{unitree-sdk2-cpp-robot-a2-playstopparameter-dunder-init-1}}

初始化 \passthrough{\lstinline!PlayStopParameter!}
实例。是否能实际构造取决于下方可用性状态。

\textbf{签名}

\begin{lstlisting}[language=Python]
def __init__(self) -> None
\end{lstlisting}

\textbf{可用性与安全性}

\textbf{\passthrough{\lstinline!SIGNATURE\_ONLY!} /
\passthrough{\lstinline!CONSTRUCTION!}}：当前只有设计期签名，不能假设已存在于运行时扩展。

\textbf{参数}

无显式参数。实例方法中的 \passthrough{\lstinline!self!} 由 Python
自动传入。

\textbf{返回值}

\begin{longtable}[]{@{}
  >{\raggedright\arraybackslash}p{(\columnwidth - 4\tabcolsep) * \real{0.18}}
  >{\raggedright\arraybackslash}p{(\columnwidth - 4\tabcolsep) * \real{0.20}}
  >{\raggedright\arraybackslash}p{(\columnwidth - 4\tabcolsep) * \real{0.62}}@{}}
\toprule
位置 & 类型 & 含义 \\
\midrule
\endhead
无 & \passthrough{\lstinline!None!} &
\passthrough{\lstinline!\_\_init\_\_!}
本身不返回值；调用类对象时，在构造函数可用的前提下得到该类实例。 \\
\bottomrule
\end{longtable}

\textbf{对应 C++}

\begin{itemize}
\tightlist
\item
  类：\passthrough{\lstinline!unitree::robot::a2::PlayStopParameter!}
\item
  签名：\passthrough{\lstinline!PlayStopParameter()!}
\item
  绑定策略：\passthrough{\lstinline!未单独标注!}
\item
  声明位置：\passthrough{\lstinline!include/unitree/robot/a2/audio/audio\_api.hpp:61!}
\end{itemize}

\textbf{说明}

\begin{itemize}
\tightlist
\item
  示例是局部调用片段；\passthrough{\lstinline!obj!}
  和参数变量表示已经按上层业务校验并准备好的对象。
\item
  该条目可以被编辑器和类型检查器识别，但当前运行时可能在属性查找阶段直接失败。
\end{itemize}

\textbf{用法}

\begin{lstlisting}[language=Python]
from typing import TYPE_CHECKING

if TYPE_CHECKING:
    def planned_construct() -> PlayStopParameter:
        return PlayStopParameter()
\end{lstlisting}

\begin{quote}
{[}!CAUTION{]} 上面的 \passthrough{\lstinline!TYPE\_CHECKING!}
示例只表达计划调用形态。该分支在普通 Python
运行时为假，不代表当前扩展能执行此接口。
\end{quote}

\hypertarget{unitree-sdk2-cpp-robot-a2-playstopparameter-from-json-1}{%
\paragraph{\texorpdfstring{\texttt{unitree\_sdk2\_cpp.robot.a2.PlayStopParameter.from\_json}}{unitree\_sdk2\_cpp.robot.a2.PlayStopParameter.from\_json}}\label{unitree-sdk2-cpp-robot-a2-playstopparameter-from-json-1}}

计划从 JSON 风格字典读取字段并更新当前 SDK 值对象。

\textbf{签名}

\begin{lstlisting}[language=Python]
def from_json(self, json: dict[str, Any]) -> None
\end{lstlisting}

\textbf{可用性与安全性}

\textbf{\passthrough{\lstinline!SIGNATURE\_ONLY!} /
\passthrough{\lstinline!UNCLASSIFIED!}}：当前只有设计期签名，不能假设已存在于运行时扩展。

\textbf{参数}

\begin{longtable}[]{@{}
  >{\raggedright\arraybackslash}p{(\columnwidth - 6\tabcolsep) * \real{0.18}}
  >{\raggedright\arraybackslash}p{(\columnwidth - 6\tabcolsep) * \real{0.24}}
  >{\raggedright\arraybackslash}p{(\columnwidth - 6\tabcolsep) * \real{0.18}}
  >{\raggedright\arraybackslash}p{(\columnwidth - 6\tabcolsep) * \real{0.40}}@{}}
\toprule
名称 & Python 类型 & 默认值 & 含义 \\
\midrule
\endhead
\passthrough{\lstinline!json!} &
\passthrough{\lstinline!dict[str, Any]!} & 必填 & JSON
风格字典，键和值必须满足对应 C++ \passthrough{\lstinline!JsonMap!}
数据结构的约束。 对应 C++ 参数
\passthrough{\lstinline!json: common::JsonMap \&!}。 \\
\bottomrule
\end{longtable}

\textbf{返回值}

\begin{longtable}[]{@{}
  >{\raggedright\arraybackslash}p{(\columnwidth - 4\tabcolsep) * \real{0.18}}
  >{\raggedright\arraybackslash}p{(\columnwidth - 4\tabcolsep) * \real{0.20}}
  >{\raggedright\arraybackslash}p{(\columnwidth - 4\tabcolsep) * \real{0.62}}@{}}
\toprule
位置 & 类型 & 含义 \\
\midrule
\endhead
无 & \passthrough{\lstinline!None!} & 该方法不返回 Python
值；失败可能表现为异常或底层状态变化。 \\
\bottomrule
\end{longtable}

\textbf{对应 C++}

\begin{itemize}
\tightlist
\item
  类：\passthrough{\lstinline!unitree::robot::a2::PlayStopParameter!}
\item
  签名：\passthrough{\lstinline!fromJson(common::JsonMap \&)!}
\item
  绑定策略：\passthrough{\lstinline!SIGNATURE\_PREVIEW!}
\item
  声明位置：\passthrough{\lstinline!include/unitree/robot/a2/audio/audio\_api.hpp:64!}
\end{itemize}

\textbf{说明}

\begin{itemize}
\tightlist
\item
  示例是局部调用片段；\passthrough{\lstinline!obj!}
  和参数变量表示已经按上层业务校验并准备好的对象。
\item
  该条目可以被编辑器和类型检查器识别，但当前运行时可能在属性查找阶段直接失败。
\end{itemize}

\textbf{用法}

\begin{lstlisting}[language=Python]
from typing import TYPE_CHECKING

if TYPE_CHECKING:
    def planned_from_json(obj: PlayStopParameter, json: dict[str, Any]) -> None:
        obj.from_json(json=json)
\end{lstlisting}

\begin{quote}
{[}!CAUTION{]} 上面的 \passthrough{\lstinline!TYPE\_CHECKING!}
示例只表达计划调用形态。该分支在普通 Python
运行时为假，不代表当前扩展能执行此接口。
\end{quote}

\hypertarget{unitree-sdk2-cpp-robot-a2-playstopparameter-to-json-1}{%
\paragraph{\texorpdfstring{\texttt{unitree\_sdk2\_cpp.robot.a2.PlayStopParameter.to\_json}}{unitree\_sdk2\_cpp.robot.a2.PlayStopParameter.to\_json}}\label{unitree-sdk2-cpp-robot-a2-playstopparameter-to-json-1}}

计划把当前 SDK 值对象写入 JSON 风格字典。

\textbf{签名}

\begin{lstlisting}[language=Python]
def to_json(self, json: dict[str, Any]) -> None
\end{lstlisting}

\textbf{可用性与安全性}

\textbf{\passthrough{\lstinline!SIGNATURE\_ONLY!} /
\passthrough{\lstinline!UNCLASSIFIED!}}：当前只有设计期签名，不能假设已存在于运行时扩展。

\textbf{参数}

\begin{longtable}[]{@{}
  >{\raggedright\arraybackslash}p{(\columnwidth - 6\tabcolsep) * \real{0.18}}
  >{\raggedright\arraybackslash}p{(\columnwidth - 6\tabcolsep) * \real{0.24}}
  >{\raggedright\arraybackslash}p{(\columnwidth - 6\tabcolsep) * \real{0.18}}
  >{\raggedright\arraybackslash}p{(\columnwidth - 6\tabcolsep) * \real{0.40}}@{}}
\toprule
名称 & Python 类型 & 默认值 & 含义 \\
\midrule
\endhead
\passthrough{\lstinline!json!} &
\passthrough{\lstinline!dict[str, Any]!} & 必填 & JSON
风格字典，键和值必须满足对应 C++ \passthrough{\lstinline!JsonMap!}
数据结构的约束。 对应 C++ 参数
\passthrough{\lstinline!json: common::JsonMap \&!}。 \\
\bottomrule
\end{longtable}

\textbf{返回值}

\begin{longtable}[]{@{}
  >{\raggedright\arraybackslash}p{(\columnwidth - 4\tabcolsep) * \real{0.18}}
  >{\raggedright\arraybackslash}p{(\columnwidth - 4\tabcolsep) * \real{0.20}}
  >{\raggedright\arraybackslash}p{(\columnwidth - 4\tabcolsep) * \real{0.62}}@{}}
\toprule
位置 & 类型 & 含义 \\
\midrule
\endhead
无 & \passthrough{\lstinline!None!} & 该方法不返回 Python
值；失败可能表现为异常或底层状态变化。 \\
\bottomrule
\end{longtable}

\textbf{对应 C++}

\begin{itemize}
\tightlist
\item
  类：\passthrough{\lstinline!unitree::robot::a2::PlayStopParameter!}
\item
  签名：\passthrough{\lstinline!toJson(common::JsonMap \&) const!}
\item
  绑定策略：\passthrough{\lstinline!SIGNATURE\_PREVIEW!}
\item
  声明位置：\passthrough{\lstinline!include/unitree/robot/a2/audio/audio\_api.hpp:66!}
\end{itemize}

\textbf{说明}

\begin{itemize}
\tightlist
\item
  示例是局部调用片段；\passthrough{\lstinline!obj!}
  和参数变量表示已经按上层业务校验并准备好的对象。
\item
  该条目可以被编辑器和类型检查器识别，但当前运行时可能在属性查找阶段直接失败。
\end{itemize}

\textbf{用法}

\begin{lstlisting}[language=Python]
from typing import TYPE_CHECKING

if TYPE_CHECKING:
    def planned_to_json(obj: PlayStopParameter, json: dict[str, Any]) -> None:
        obj.to_json(json=json)
\end{lstlisting}

\begin{quote}
{[}!CAUTION{]} 上面的 \passthrough{\lstinline!TYPE\_CHECKING!}
示例只表达计划调用形态。该分支在普通 Python
运行时为假，不代表当前扩展能执行此接口。
\end{quote}

\hypertarget{unitree-sdk2-cpp-robot-a2-playstreamparameter}{%
\subsubsection{\texorpdfstring{\texttt{unitree\_sdk2\_cpp.robot.a2.PlayStreamParameter}}{unitree\_sdk2\_cpp.robot.a2.PlayStreamParameter}}\label{unitree-sdk2-cpp-robot-a2-playstreamparameter}}

SDK 请求参数值类型；当前可能仅提供设计期签名。

\textbf{导入}

\begin{lstlisting}[language=Python]
from unitree_sdk2_cpp.robot.a2 import PlayStreamParameter
\end{lstlisting}

\textbf{基类}：\passthrough{\lstinline!object!}

\textbf{公开属性}

\begin{longtable}[]{@{}
  >{\raggedright\arraybackslash}p{(\columnwidth - 6\tabcolsep) * \real{0.18}}
  >{\raggedright\arraybackslash}p{(\columnwidth - 6\tabcolsep) * \real{0.24}}
  >{\raggedright\arraybackslash}p{(\columnwidth - 6\tabcolsep) * \real{0.18}}
  >{\raggedright\arraybackslash}p{(\columnwidth - 6\tabcolsep) * \real{0.40}}@{}}
\toprule
名称 & Python 类型 & 含义 & 用法 \\
\midrule
\endhead
\passthrough{\lstinline!app\_name!} & \passthrough{\lstinline!str!} &
应用名称，用于标识音频流或播放会话。允许值和命名规则以目标服务协议为准。
& \passthrough{\lstinline!value = obj.app\_name!} /
\passthrough{\lstinline!obj.app\_name = value!} \\
\passthrough{\lstinline!stream\_id!} & \passthrough{\lstinline!str!} &
音频流标识符，用于区分同一应用下的播放流。 &
\passthrough{\lstinline!value = obj.stream\_id!} /
\passthrough{\lstinline!obj.stream\_id = value!} \\
\bottomrule
\end{longtable}

\textbf{方法索引}

\begin{longtable}[]{@{}
  >{\raggedright\arraybackslash}p{(\columnwidth - 6\tabcolsep) * \real{0.18}}
  >{\raggedright\arraybackslash}p{(\columnwidth - 6\tabcolsep) * \real{0.24}}
  >{\raggedright\arraybackslash}p{(\columnwidth - 6\tabcolsep) * \real{0.18}}
  >{\raggedright\arraybackslash}p{(\columnwidth - 6\tabcolsep) * \real{0.40}}@{}}
\toprule
方法 & Python 签名 & 状态 & 安全分类 \\
\midrule
\endhead
\protect\hyperlink{unitree-sdk2-cpp-robot-a2-playstreamparameter-dunder-init-1}{\passthrough{\lstinline!\_\_init\_\_!}}
& \passthrough{\lstinline!def \_\_init\_\_(self) -> None!} &
\passthrough{\lstinline!SIGNATURE\_ONLY!} &
\passthrough{\lstinline!CONSTRUCTION!} \\
\protect\hyperlink{unitree-sdk2-cpp-robot-a2-playstreamparameter-from-json-1}{\passthrough{\lstinline!from\_json!}}
&
\passthrough{\lstinline!def from\_json(self, json: dict[str, Any]) -> None!}
& \passthrough{\lstinline!SIGNATURE\_ONLY!} &
\passthrough{\lstinline!UNCLASSIFIED!} \\
\protect\hyperlink{unitree-sdk2-cpp-robot-a2-playstreamparameter-to-json-1}{\passthrough{\lstinline!to\_json!}}
&
\passthrough{\lstinline!def to\_json(self, json: dict[str, Any]) -> None!}
& \passthrough{\lstinline!SIGNATURE\_ONLY!} &
\passthrough{\lstinline!UNCLASSIFIED!} \\
\bottomrule
\end{longtable}

\hypertarget{unitree-sdk2-cpp-robot-a2-playstreamparameter-dunder-init-1}{%
\paragraph{\texorpdfstring{\texttt{unitree\_sdk2\_cpp.robot.a2.PlayStreamParameter.\_\_init\_\_}}{unitree\_sdk2\_cpp.robot.a2.PlayStreamParameter.\_\_init\_\_}}\label{unitree-sdk2-cpp-robot-a2-playstreamparameter-dunder-init-1}}

初始化 \passthrough{\lstinline!PlayStreamParameter!}
实例。是否能实际构造取决于下方可用性状态。

\textbf{签名}

\begin{lstlisting}[language=Python]
def __init__(self) -> None
\end{lstlisting}

\textbf{可用性与安全性}

\textbf{\passthrough{\lstinline!SIGNATURE\_ONLY!} /
\passthrough{\lstinline!CONSTRUCTION!}}：当前只有设计期签名，不能假设已存在于运行时扩展。

\textbf{参数}

无显式参数。实例方法中的 \passthrough{\lstinline!self!} 由 Python
自动传入。

\textbf{返回值}

\begin{longtable}[]{@{}
  >{\raggedright\arraybackslash}p{(\columnwidth - 4\tabcolsep) * \real{0.18}}
  >{\raggedright\arraybackslash}p{(\columnwidth - 4\tabcolsep) * \real{0.20}}
  >{\raggedright\arraybackslash}p{(\columnwidth - 4\tabcolsep) * \real{0.62}}@{}}
\toprule
位置 & 类型 & 含义 \\
\midrule
\endhead
无 & \passthrough{\lstinline!None!} &
\passthrough{\lstinline!\_\_init\_\_!}
本身不返回值；调用类对象时，在构造函数可用的前提下得到该类实例。 \\
\bottomrule
\end{longtable}

\textbf{对应 C++}

\begin{itemize}
\tightlist
\item
  类：\passthrough{\lstinline!unitree::robot::a2::PlayStreamParameter!}
\item
  签名：\passthrough{\lstinline!PlayStreamParameter()!}
\item
  绑定策略：\passthrough{\lstinline!未单独标注!}
\item
  声明位置：\passthrough{\lstinline!include/unitree/robot/a2/audio/audio\_api.hpp:45!}
\end{itemize}

\textbf{说明}

\begin{itemize}
\tightlist
\item
  示例是局部调用片段；\passthrough{\lstinline!obj!}
  和参数变量表示已经按上层业务校验并准备好的对象。
\item
  该条目可以被编辑器和类型检查器识别，但当前运行时可能在属性查找阶段直接失败。
\end{itemize}

\textbf{用法}

\begin{lstlisting}[language=Python]
from typing import TYPE_CHECKING

if TYPE_CHECKING:
    def planned_construct() -> PlayStreamParameter:
        return PlayStreamParameter()
\end{lstlisting}

\begin{quote}
{[}!CAUTION{]} 上面的 \passthrough{\lstinline!TYPE\_CHECKING!}
示例只表达计划调用形态。该分支在普通 Python
运行时为假，不代表当前扩展能执行此接口。
\end{quote}

\hypertarget{unitree-sdk2-cpp-robot-a2-playstreamparameter-from-json-1}{%
\paragraph{\texorpdfstring{\texttt{unitree\_sdk2\_cpp.robot.a2.PlayStreamParameter.from\_json}}{unitree\_sdk2\_cpp.robot.a2.PlayStreamParameter.from\_json}}\label{unitree-sdk2-cpp-robot-a2-playstreamparameter-from-json-1}}

计划从 JSON 风格字典读取字段并更新当前 SDK 值对象。

\textbf{签名}

\begin{lstlisting}[language=Python]
def from_json(self, json: dict[str, Any]) -> None
\end{lstlisting}

\textbf{可用性与安全性}

\textbf{\passthrough{\lstinline!SIGNATURE\_ONLY!} /
\passthrough{\lstinline!UNCLASSIFIED!}}：当前只有设计期签名，不能假设已存在于运行时扩展。

\textbf{参数}

\begin{longtable}[]{@{}
  >{\raggedright\arraybackslash}p{(\columnwidth - 6\tabcolsep) * \real{0.18}}
  >{\raggedright\arraybackslash}p{(\columnwidth - 6\tabcolsep) * \real{0.24}}
  >{\raggedright\arraybackslash}p{(\columnwidth - 6\tabcolsep) * \real{0.18}}
  >{\raggedright\arraybackslash}p{(\columnwidth - 6\tabcolsep) * \real{0.40}}@{}}
\toprule
名称 & Python 类型 & 默认值 & 含义 \\
\midrule
\endhead
\passthrough{\lstinline!json!} &
\passthrough{\lstinline!dict[str, Any]!} & 必填 & JSON
风格字典，键和值必须满足对应 C++ \passthrough{\lstinline!JsonMap!}
数据结构的约束。 对应 C++ 参数
\passthrough{\lstinline!json: common::JsonMap \&!}。 \\
\bottomrule
\end{longtable}

\textbf{返回值}

\begin{longtable}[]{@{}
  >{\raggedright\arraybackslash}p{(\columnwidth - 4\tabcolsep) * \real{0.18}}
  >{\raggedright\arraybackslash}p{(\columnwidth - 4\tabcolsep) * \real{0.20}}
  >{\raggedright\arraybackslash}p{(\columnwidth - 4\tabcolsep) * \real{0.62}}@{}}
\toprule
位置 & 类型 & 含义 \\
\midrule
\endhead
无 & \passthrough{\lstinline!None!} & 该方法不返回 Python
值；失败可能表现为异常或底层状态变化。 \\
\bottomrule
\end{longtable}

\textbf{对应 C++}

\begin{itemize}
\tightlist
\item
  类：\passthrough{\lstinline!unitree::robot::a2::PlayStreamParameter!}
\item
  签名：\passthrough{\lstinline!fromJson(common::JsonMap \&)!}
\item
  绑定策略：\passthrough{\lstinline!SIGNATURE\_PREVIEW!}
\item
  声明位置：\passthrough{\lstinline!include/unitree/robot/a2/audio/audio\_api.hpp:48!}
\end{itemize}

\textbf{说明}

\begin{itemize}
\tightlist
\item
  示例是局部调用片段；\passthrough{\lstinline!obj!}
  和参数变量表示已经按上层业务校验并准备好的对象。
\item
  该条目可以被编辑器和类型检查器识别，但当前运行时可能在属性查找阶段直接失败。
\end{itemize}

\textbf{用法}

\begin{lstlisting}[language=Python]
from typing import TYPE_CHECKING

if TYPE_CHECKING:
    def planned_from_json(obj: PlayStreamParameter, json: dict[str, Any]) -> None:
        obj.from_json(json=json)
\end{lstlisting}

\begin{quote}
{[}!CAUTION{]} 上面的 \passthrough{\lstinline!TYPE\_CHECKING!}
示例只表达计划调用形态。该分支在普通 Python
运行时为假，不代表当前扩展能执行此接口。
\end{quote}

\hypertarget{unitree-sdk2-cpp-robot-a2-playstreamparameter-to-json-1}{%
\paragraph{\texorpdfstring{\texttt{unitree\_sdk2\_cpp.robot.a2.PlayStreamParameter.to\_json}}{unitree\_sdk2\_cpp.robot.a2.PlayStreamParameter.to\_json}}\label{unitree-sdk2-cpp-robot-a2-playstreamparameter-to-json-1}}

计划把当前 SDK 值对象写入 JSON 风格字典。

\textbf{签名}

\begin{lstlisting}[language=Python]
def to_json(self, json: dict[str, Any]) -> None
\end{lstlisting}

\textbf{可用性与安全性}

\textbf{\passthrough{\lstinline!SIGNATURE\_ONLY!} /
\passthrough{\lstinline!UNCLASSIFIED!}}：当前只有设计期签名，不能假设已存在于运行时扩展。

\textbf{参数}

\begin{longtable}[]{@{}
  >{\raggedright\arraybackslash}p{(\columnwidth - 6\tabcolsep) * \real{0.18}}
  >{\raggedright\arraybackslash}p{(\columnwidth - 6\tabcolsep) * \real{0.24}}
  >{\raggedright\arraybackslash}p{(\columnwidth - 6\tabcolsep) * \real{0.18}}
  >{\raggedright\arraybackslash}p{(\columnwidth - 6\tabcolsep) * \real{0.40}}@{}}
\toprule
名称 & Python 类型 & 默认值 & 含义 \\
\midrule
\endhead
\passthrough{\lstinline!json!} &
\passthrough{\lstinline!dict[str, Any]!} & 必填 & JSON
风格字典，键和值必须满足对应 C++ \passthrough{\lstinline!JsonMap!}
数据结构的约束。 对应 C++ 参数
\passthrough{\lstinline!json: common::JsonMap \&!}。 \\
\bottomrule
\end{longtable}

\textbf{返回值}

\begin{longtable}[]{@{}
  >{\raggedright\arraybackslash}p{(\columnwidth - 4\tabcolsep) * \real{0.18}}
  >{\raggedright\arraybackslash}p{(\columnwidth - 4\tabcolsep) * \real{0.20}}
  >{\raggedright\arraybackslash}p{(\columnwidth - 4\tabcolsep) * \real{0.62}}@{}}
\toprule
位置 & 类型 & 含义 \\
\midrule
\endhead
无 & \passthrough{\lstinline!None!} & 该方法不返回 Python
值；失败可能表现为异常或底层状态变化。 \\
\bottomrule
\end{longtable}

\textbf{对应 C++}

\begin{itemize}
\tightlist
\item
  类：\passthrough{\lstinline!unitree::robot::a2::PlayStreamParameter!}
\item
  签名：\passthrough{\lstinline!toJson(common::JsonMap \&) const!}
\item
  绑定策略：\passthrough{\lstinline!SIGNATURE\_PREVIEW!}
\item
  声明位置：\passthrough{\lstinline!include/unitree/robot/a2/audio/audio\_api.hpp:50!}
\end{itemize}

\textbf{说明}

\begin{itemize}
\tightlist
\item
  示例是局部调用片段；\passthrough{\lstinline!obj!}
  和参数变量表示已经按上层业务校验并准备好的对象。
\item
  该条目可以被编辑器和类型检查器识别，但当前运行时可能在属性查找阶段直接失败。
\end{itemize}

\textbf{用法}

\begin{lstlisting}[language=Python]
from typing import TYPE_CHECKING

if TYPE_CHECKING:
    def planned_to_json(obj: PlayStreamParameter, json: dict[str, Any]) -> None:
        obj.to_json(json=json)
\end{lstlisting}

\begin{quote}
{[}!CAUTION{]} 上面的 \passthrough{\lstinline!TYPE\_CHECKING!}
示例只表达计划调用形态。该分支在普通 Python
运行时为假，不代表当前扩展能执行此接口。
\end{quote}

\hypertarget{unitree-sdk2-cpp-robot-a2-posevec4}{%
\subsubsection{\texorpdfstring{\texttt{unitree\_sdk2\_cpp.robot.a2.PoseVec4}}{unitree\_sdk2\_cpp.robot.a2.PoseVec4}}\label{unitree-sdk2-cpp-robot-a2-posevec4}}

SDK 类型签名预览。请逐项查看构造函数和方法的可用性。

\textbf{导入}

\begin{lstlisting}[language=Python]
from unitree_sdk2_cpp.robot.a2 import PoseVec4
\end{lstlisting}

\textbf{基类}：\passthrough{\lstinline!object!}

\textbf{公开属性}

\begin{longtable}[]{@{}
  >{\raggedright\arraybackslash}p{(\columnwidth - 6\tabcolsep) * \real{0.18}}
  >{\raggedright\arraybackslash}p{(\columnwidth - 6\tabcolsep) * \real{0.24}}
  >{\raggedright\arraybackslash}p{(\columnwidth - 6\tabcolsep) * \real{0.18}}
  >{\raggedright\arraybackslash}p{(\columnwidth - 6\tabcolsep) * \real{0.40}}@{}}
\toprule
名称 & Python 类型 & 含义 & 用法 \\
\midrule
\endhead
\passthrough{\lstinline!x!} & \passthrough{\lstinline!float!} & X
方向位置或分量参数。参考系、单位和范围必须查当前方法及目标型号协议。 &
\passthrough{\lstinline!value = obj.x!} /
\passthrough{\lstinline!obj.x = value!} \\
\passthrough{\lstinline!y!} & \passthrough{\lstinline!float!} & Y
方向位置或分量参数。参考系、单位和范围必须查当前方法及目标型号协议。 &
\passthrough{\lstinline!value = obj.y!} /
\passthrough{\lstinline!obj.y = value!} \\
\passthrough{\lstinline!z!} & \passthrough{\lstinline!float!} & Z
方向位置或分量参数。参考系、单位和范围必须查当前方法及目标型号协议。 &
\passthrough{\lstinline!value = obj.z!} /
\passthrough{\lstinline!obj.z = value!} \\
\passthrough{\lstinline!yaw!} & \passthrough{\lstinline!float!} &
偏航角或偏航目标参数。参考系、单位和范围必须查目标型号协议。 &
\passthrough{\lstinline!value = obj.yaw!} /
\passthrough{\lstinline!obj.yaw = value!} \\
\bottomrule
\end{longtable}

\textbf{方法索引}

\begin{longtable}[]{@{}
  >{\raggedright\arraybackslash}p{(\columnwidth - 6\tabcolsep) * \real{0.18}}
  >{\raggedright\arraybackslash}p{(\columnwidth - 6\tabcolsep) * \real{0.24}}
  >{\raggedright\arraybackslash}p{(\columnwidth - 6\tabcolsep) * \real{0.18}}
  >{\raggedright\arraybackslash}p{(\columnwidth - 6\tabcolsep) * \real{0.40}}@{}}
\toprule
方法 & Python 签名 & 状态 & 安全分类 \\
\midrule
\endhead
\protect\hyperlink{unitree-sdk2-cpp-robot-a2-posevec4-dunder-init-1}{\passthrough{\lstinline!\_\_init\_\_!}}
& \passthrough{\lstinline!def \_\_init\_\_(self) -> None!} &
\passthrough{\lstinline!SIGNATURE\_ONLY!} &
\passthrough{\lstinline!CONSTRUCTION!} \\
\protect\hyperlink{unitree-sdk2-cpp-robot-a2-posevec4-from-json-1}{\passthrough{\lstinline!from\_json!}}
&
\passthrough{\lstinline!def from\_json(self, json: dict[str, Any]) -> None!}
& \passthrough{\lstinline!SIGNATURE\_ONLY!} &
\passthrough{\lstinline!UNCLASSIFIED!} \\
\protect\hyperlink{unitree-sdk2-cpp-robot-a2-posevec4-to-json-1}{\passthrough{\lstinline!to\_json!}}
&
\passthrough{\lstinline!def to\_json(self, json: dict[str, Any]) -> None!}
& \passthrough{\lstinline!SIGNATURE\_ONLY!} &
\passthrough{\lstinline!UNCLASSIFIED!} \\
\bottomrule
\end{longtable}

\hypertarget{unitree-sdk2-cpp-robot-a2-posevec4-dunder-init-1}{%
\paragraph{\texorpdfstring{\texttt{unitree\_sdk2\_cpp.robot.a2.PoseVec4.\_\_init\_\_}}{unitree\_sdk2\_cpp.robot.a2.PoseVec4.\_\_init\_\_}}\label{unitree-sdk2-cpp-robot-a2-posevec4-dunder-init-1}}

初始化 \passthrough{\lstinline!PoseVec4!}
实例。是否能实际构造取决于下方可用性状态。

\textbf{签名}

\begin{lstlisting}[language=Python]
def __init__(self) -> None
\end{lstlisting}

\textbf{可用性与安全性}

\textbf{\passthrough{\lstinline!SIGNATURE\_ONLY!} /
\passthrough{\lstinline!CONSTRUCTION!}}：当前只有设计期签名，不能假设已存在于运行时扩展。

\textbf{参数}

无显式参数。实例方法中的 \passthrough{\lstinline!self!} 由 Python
自动传入。

\textbf{返回值}

\begin{longtable}[]{@{}
  >{\raggedright\arraybackslash}p{(\columnwidth - 4\tabcolsep) * \real{0.18}}
  >{\raggedright\arraybackslash}p{(\columnwidth - 4\tabcolsep) * \real{0.20}}
  >{\raggedright\arraybackslash}p{(\columnwidth - 4\tabcolsep) * \real{0.62}}@{}}
\toprule
位置 & 类型 & 含义 \\
\midrule
\endhead
无 & \passthrough{\lstinline!None!} &
\passthrough{\lstinline!\_\_init\_\_!}
本身不返回值；调用类对象时，在构造函数可用的前提下得到该类实例。 \\
\bottomrule
\end{longtable}

\textbf{对应 C++}

\begin{itemize}
\tightlist
\item
  类：\passthrough{\lstinline!unitree::robot::a2::PoseVec4!}
\item
  签名：\passthrough{\lstinline!PoseVec4()!}
\item
  绑定策略：\passthrough{\lstinline!未单独标注!}
\item
  声明位置：\passthrough{\lstinline!include/unitree/robot/a2/sport/sport\_client.hpp:35!}
\end{itemize}

\textbf{说明}

\begin{itemize}
\tightlist
\item
  示例是局部调用片段；\passthrough{\lstinline!obj!}
  和参数变量表示已经按上层业务校验并准备好的对象。
\item
  该条目可以被编辑器和类型检查器识别，但当前运行时可能在属性查找阶段直接失败。
\end{itemize}

\textbf{用法}

\begin{lstlisting}[language=Python]
from typing import TYPE_CHECKING

if TYPE_CHECKING:
    def planned_construct() -> PoseVec4:
        return PoseVec4()
\end{lstlisting}

\begin{quote}
{[}!CAUTION{]} 上面的 \passthrough{\lstinline!TYPE\_CHECKING!}
示例只表达计划调用形态。该分支在普通 Python
运行时为假，不代表当前扩展能执行此接口。
\end{quote}

\hypertarget{unitree-sdk2-cpp-robot-a2-posevec4-from-json-1}{%
\paragraph{\texorpdfstring{\texttt{unitree\_sdk2\_cpp.robot.a2.PoseVec4.from\_json}}{unitree\_sdk2\_cpp.robot.a2.PoseVec4.from\_json}}\label{unitree-sdk2-cpp-robot-a2-posevec4-from-json-1}}

计划从 JSON 风格字典读取字段并更新当前 SDK 值对象。

\textbf{签名}

\begin{lstlisting}[language=Python]
def from_json(self, json: dict[str, Any]) -> None
\end{lstlisting}

\textbf{可用性与安全性}

\textbf{\passthrough{\lstinline!SIGNATURE\_ONLY!} /
\passthrough{\lstinline!UNCLASSIFIED!}}：当前只有设计期签名，不能假设已存在于运行时扩展。

\textbf{参数}

\begin{longtable}[]{@{}
  >{\raggedright\arraybackslash}p{(\columnwidth - 6\tabcolsep) * \real{0.18}}
  >{\raggedright\arraybackslash}p{(\columnwidth - 6\tabcolsep) * \real{0.24}}
  >{\raggedright\arraybackslash}p{(\columnwidth - 6\tabcolsep) * \real{0.18}}
  >{\raggedright\arraybackslash}p{(\columnwidth - 6\tabcolsep) * \real{0.40}}@{}}
\toprule
名称 & Python 类型 & 默认值 & 含义 \\
\midrule
\endhead
\passthrough{\lstinline!json!} &
\passthrough{\lstinline!dict[str, Any]!} & 必填 & JSON
风格字典，键和值必须满足对应 C++ \passthrough{\lstinline!JsonMap!}
数据结构的约束。 对应 C++ 参数
\passthrough{\lstinline!json: common::JsonMap \&!}。 \\
\bottomrule
\end{longtable}

\textbf{返回值}

\begin{longtable}[]{@{}
  >{\raggedright\arraybackslash}p{(\columnwidth - 4\tabcolsep) * \real{0.18}}
  >{\raggedright\arraybackslash}p{(\columnwidth - 4\tabcolsep) * \real{0.20}}
  >{\raggedright\arraybackslash}p{(\columnwidth - 4\tabcolsep) * \real{0.62}}@{}}
\toprule
位置 & 类型 & 含义 \\
\midrule
\endhead
无 & \passthrough{\lstinline!None!} & 该方法不返回 Python
值；失败可能表现为异常或底层状态变化。 \\
\bottomrule
\end{longtable}

\textbf{对应 C++}

\begin{itemize}
\tightlist
\item
  类：\passthrough{\lstinline!unitree::robot::a2::PoseVec4!}
\item
  签名：\passthrough{\lstinline!fromJson(common::JsonMap \&)!}
\item
  绑定策略：\passthrough{\lstinline!SIGNATURE\_PREVIEW!}
\item
  声明位置：\passthrough{\lstinline!include/unitree/robot/a2/sport/sport\_client.hpp:43!}
\end{itemize}

\textbf{说明}

\begin{itemize}
\tightlist
\item
  示例是局部调用片段；\passthrough{\lstinline!obj!}
  和参数变量表示已经按上层业务校验并准备好的对象。
\item
  该条目可以被编辑器和类型检查器识别，但当前运行时可能在属性查找阶段直接失败。
\end{itemize}

\textbf{用法}

\begin{lstlisting}[language=Python]
from typing import TYPE_CHECKING

if TYPE_CHECKING:
    def planned_from_json(obj: PoseVec4, json: dict[str, Any]) -> None:
        obj.from_json(json=json)
\end{lstlisting}

\begin{quote}
{[}!CAUTION{]} 上面的 \passthrough{\lstinline!TYPE\_CHECKING!}
示例只表达计划调用形态。该分支在普通 Python
运行时为假，不代表当前扩展能执行此接口。
\end{quote}

\hypertarget{unitree-sdk2-cpp-robot-a2-posevec4-to-json-1}{%
\paragraph{\texorpdfstring{\texttt{unitree\_sdk2\_cpp.robot.a2.PoseVec4.to\_json}}{unitree\_sdk2\_cpp.robot.a2.PoseVec4.to\_json}}\label{unitree-sdk2-cpp-robot-a2-posevec4-to-json-1}}

计划把当前 SDK 值对象写入 JSON 风格字典。

\textbf{签名}

\begin{lstlisting}[language=Python]
def to_json(self, json: dict[str, Any]) -> None
\end{lstlisting}

\textbf{可用性与安全性}

\textbf{\passthrough{\lstinline!SIGNATURE\_ONLY!} /
\passthrough{\lstinline!UNCLASSIFIED!}}：当前只有设计期签名，不能假设已存在于运行时扩展。

\textbf{参数}

\begin{longtable}[]{@{}
  >{\raggedright\arraybackslash}p{(\columnwidth - 6\tabcolsep) * \real{0.18}}
  >{\raggedright\arraybackslash}p{(\columnwidth - 6\tabcolsep) * \real{0.24}}
  >{\raggedright\arraybackslash}p{(\columnwidth - 6\tabcolsep) * \real{0.18}}
  >{\raggedright\arraybackslash}p{(\columnwidth - 6\tabcolsep) * \real{0.40}}@{}}
\toprule
名称 & Python 类型 & 默认值 & 含义 \\
\midrule
\endhead
\passthrough{\lstinline!json!} &
\passthrough{\lstinline!dict[str, Any]!} & 必填 & JSON
风格字典，键和值必须满足对应 C++ \passthrough{\lstinline!JsonMap!}
数据结构的约束。 对应 C++ 参数
\passthrough{\lstinline!json: common::JsonMap \&!}。 \\
\bottomrule
\end{longtable}

\textbf{返回值}

\begin{longtable}[]{@{}
  >{\raggedright\arraybackslash}p{(\columnwidth - 4\tabcolsep) * \real{0.18}}
  >{\raggedright\arraybackslash}p{(\columnwidth - 4\tabcolsep) * \real{0.20}}
  >{\raggedright\arraybackslash}p{(\columnwidth - 4\tabcolsep) * \real{0.62}}@{}}
\toprule
位置 & 类型 & 含义 \\
\midrule
\endhead
无 & \passthrough{\lstinline!None!} & 该方法不返回 Python
值；失败可能表现为异常或底层状态变化。 \\
\bottomrule
\end{longtable}

\textbf{对应 C++}

\begin{itemize}
\tightlist
\item
  类：\passthrough{\lstinline!unitree::robot::a2::PoseVec4!}
\item
  签名：\passthrough{\lstinline!toJson(common::JsonMap \&) const!}
\item
  绑定策略：\passthrough{\lstinline!SIGNATURE\_PREVIEW!}
\item
  声明位置：\passthrough{\lstinline!include/unitree/robot/a2/sport/sport\_client.hpp:51!}
\end{itemize}

\textbf{说明}

\begin{itemize}
\tightlist
\item
  示例是局部调用片段；\passthrough{\lstinline!obj!}
  和参数变量表示已经按上层业务校验并准备好的对象。
\item
  该条目可以被编辑器和类型检查器识别，但当前运行时可能在属性查找阶段直接失败。
\end{itemize}

\textbf{用法}

\begin{lstlisting}[language=Python]
from typing import TYPE_CHECKING

if TYPE_CHECKING:
    def planned_to_json(obj: PoseVec4, json: dict[str, Any]) -> None:
        obj.to_json(json=json)
\end{lstlisting}

\begin{quote}
{[}!CAUTION{]} 上面的 \passthrough{\lstinline!TYPE\_CHECKING!}
示例只表达计划调用形态。该分支在普通 Python
运行时为假，不代表当前扩展能执行此接口。
\end{quote}

\hypertarget{unitree-sdk2-cpp-robot-a2-sportclient}{%
\subsubsection{\texorpdfstring{\texttt{unitree\_sdk2\_cpp.robot.a2.SportClient}}{unitree\_sdk2\_cpp.robot.a2.SportClient}}\label{unitree-sdk2-cpp-robot-a2-sportclient}}

Robot 服务客户端。构造、初始化和具体方法的可用性必须分别检查。

\textbf{导入}

\begin{lstlisting}[language=Python]
from unitree_sdk2_cpp.robot.a2 import SportClient
\end{lstlisting}

\textbf{基类}：\passthrough{\lstinline!Client!}

\textbf{方法索引}

\begin{longtable}[]{@{}
  >{\raggedright\arraybackslash}p{(\columnwidth - 6\tabcolsep) * \real{0.18}}
  >{\raggedright\arraybackslash}p{(\columnwidth - 6\tabcolsep) * \real{0.24}}
  >{\raggedright\arraybackslash}p{(\columnwidth - 6\tabcolsep) * \real{0.18}}
  >{\raggedright\arraybackslash}p{(\columnwidth - 6\tabcolsep) * \real{0.40}}@{}}
\toprule
方法 & Python 签名 & 状态 & 安全分类 \\
\midrule
\endhead
\protect\hyperlink{unitree-sdk2-cpp-robot-a2-sportclient-dunder-init-1}{\passthrough{\lstinline!\_\_init\_\_!}}
& \passthrough{\lstinline!def \_\_init\_\_(self) -> None!} &
\passthrough{\lstinline!AVAILABLE!} &
\passthrough{\lstinline!CONSTRUCTION!} \\
\protect\hyperlink{unitree-sdk2-cpp-robot-a2-sportclient-init-1}{\passthrough{\lstinline!init!}}
& \passthrough{\lstinline!def init(self) -> None!} &
\passthrough{\lstinline!AVAILABLE!} &
\passthrough{\lstinline!INITIALIZATION!} \\
\protect\hyperlink{unitree-sdk2-cpp-robot-a2-sportclient-damp-1}{\passthrough{\lstinline!damp!}}
& \passthrough{\lstinline!def damp(self) -> int!} &
\passthrough{\lstinline!SIGNATURE\_ONLY!} &
\passthrough{\lstinline!MOTION\_COMMAND!} \\
\protect\hyperlink{unitree-sdk2-cpp-robot-a2-sportclient-balance-stand-1}{\passthrough{\lstinline!balance\_stand!}}
& \passthrough{\lstinline!def balance\_stand(self) -> int!} &
\passthrough{\lstinline!SIGNATURE\_ONLY!} &
\passthrough{\lstinline!MOTION\_COMMAND!} \\
\protect\hyperlink{unitree-sdk2-cpp-robot-a2-sportclient-stop-move-1}{\passthrough{\lstinline!stop\_move!}}
& \passthrough{\lstinline!def stop\_move(self) -> int!} &
\passthrough{\lstinline!SIGNATURE\_ONLY!} &
\passthrough{\lstinline!MOTION\_COMMAND!} \\
\protect\hyperlink{unitree-sdk2-cpp-robot-a2-sportclient-stand-up-1}{\passthrough{\lstinline!stand\_up!}}
& \passthrough{\lstinline!def stand\_up(self) -> int!} &
\passthrough{\lstinline!SIGNATURE\_ONLY!} &
\passthrough{\lstinline!MOTION\_COMMAND!} \\
\protect\hyperlink{unitree-sdk2-cpp-robot-a2-sportclient-stand-down-1}{\passthrough{\lstinline!stand\_down!}}
& \passthrough{\lstinline!def stand\_down(self) -> int!} &
\passthrough{\lstinline!SIGNATURE\_ONLY!} &
\passthrough{\lstinline!MOTION\_COMMAND!} \\
\protect\hyperlink{unitree-sdk2-cpp-robot-a2-sportclient-recovery-stand-1}{\passthrough{\lstinline!recovery\_stand!}}
& \passthrough{\lstinline!def recovery\_stand(self) -> int!} &
\passthrough{\lstinline!SIGNATURE\_ONLY!} &
\passthrough{\lstinline!MOTION\_COMMAND!} \\
\protect\hyperlink{unitree-sdk2-cpp-robot-a2-sportclient-euler-1}{\passthrough{\lstinline!euler!}}
&
\passthrough{\lstinline!def euler(self, roll: float, pitch: float, yaw: float) -> int!}
& \passthrough{\lstinline!SIGNATURE\_ONLY!} &
\passthrough{\lstinline!MOTION\_COMMAND!} \\
\protect\hyperlink{unitree-sdk2-cpp-robot-a2-sportclient-move-1}{\passthrough{\lstinline!move!}}
&
\passthrough{\lstinline!def move(self, vx: float, vy: float, vyaw: float) -> int!}
& \passthrough{\lstinline!SIGNATURE\_ONLY!} &
\passthrough{\lstinline!MOTION\_COMMAND!} \\
\protect\hyperlink{unitree-sdk2-cpp-robot-a2-sportclient-switch-gait-1}{\passthrough{\lstinline!switch\_gait!}}
&
\passthrough{\lstinline!def switch\_gait(self, gait\_type: int) -> int!}
& \passthrough{\lstinline!SIGNATURE\_ONLY!} &
\passthrough{\lstinline!MOTION\_COMMAND!} \\
\protect\hyperlink{unitree-sdk2-cpp-robot-a2-sportclient-body-height-1}{\passthrough{\lstinline!body\_height!}}
& \passthrough{\lstinline!def body\_height(self, height: float) -> int!}
& \passthrough{\lstinline!SIGNATURE\_ONLY!} &
\passthrough{\lstinline!MOTION\_COMMAND!} \\
\protect\hyperlink{unitree-sdk2-cpp-robot-a2-sportclient-speed-level-1}{\passthrough{\lstinline!speed\_level!}}
& \passthrough{\lstinline!def speed\_level(self, level: int) -> int!} &
\passthrough{\lstinline!SIGNATURE\_ONLY!} &
\passthrough{\lstinline!MOTION\_COMMAND!} \\
\protect\hyperlink{unitree-sdk2-cpp-robot-a2-sportclient-body-position-1}{\passthrough{\lstinline!body\_position!}}
&
\passthrough{\lstinline!def body\_position(self, x: float, y: float, z: float, yaw: float) -> int!}
& \passthrough{\lstinline!SIGNATURE\_ONLY!} &
\passthrough{\lstinline!MOTION\_COMMAND!} \\
\protect\hyperlink{unitree-sdk2-cpp-robot-a2-sportclient-left-side-gait-1}{\passthrough{\lstinline!left\_side\_gait!}}
&
\passthrough{\lstinline!def left\_side\_gait(self, enter: int) -> int!}
& \passthrough{\lstinline!SIGNATURE\_ONLY!} &
\passthrough{\lstinline!MOTION\_COMMAND!} \\
\protect\hyperlink{unitree-sdk2-cpp-robot-a2-sportclient-right-side-gait-1}{\passthrough{\lstinline!right\_side\_gait!}}
&
\passthrough{\lstinline!def right\_side\_gait(self, enter: int) -> int!}
& \passthrough{\lstinline!SIGNATURE\_ONLY!} &
\passthrough{\lstinline!MOTION\_COMMAND!} \\
\protect\hyperlink{unitree-sdk2-cpp-robot-a2-sportclient-hand-stand-1}{\passthrough{\lstinline!hand\_stand!}}
& \passthrough{\lstinline!def hand\_stand(self, enter: int) -> int!} &
\passthrough{\lstinline!SIGNATURE\_ONLY!} &
\passthrough{\lstinline!MOTION\_COMMAND!} \\
\protect\hyperlink{unitree-sdk2-cpp-robot-a2-sportclient-biped-stand-1}{\passthrough{\lstinline!biped\_stand!}}
& \passthrough{\lstinline!def biped\_stand(self, enter: int) -> int!} &
\passthrough{\lstinline!SIGNATURE\_ONLY!} &
\passthrough{\lstinline!MOTION\_COMMAND!} \\
\protect\hyperlink{unitree-sdk2-cpp-robot-a2-sportclient-front-flip-1}{\passthrough{\lstinline!front\_flip!}}
& \passthrough{\lstinline!def front\_flip(self) -> int!} &
\passthrough{\lstinline!SIGNATURE\_ONLY!} &
\passthrough{\lstinline!MOTION\_COMMAND!} \\
\protect\hyperlink{unitree-sdk2-cpp-robot-a2-sportclient-back-flip-1}{\passthrough{\lstinline!back\_flip!}}
& \passthrough{\lstinline!def back\_flip(self) -> int!} &
\passthrough{\lstinline!SIGNATURE\_ONLY!} &
\passthrough{\lstinline!MOTION\_COMMAND!} \\
\protect\hyperlink{unitree-sdk2-cpp-robot-a2-sportclient-reset-estimator-1}{\passthrough{\lstinline!reset\_estimator!}}
& \passthrough{\lstinline!def reset\_estimator(self) -> int!} &
\passthrough{\lstinline!SIGNATURE\_ONLY!} &
\passthrough{\lstinline!MOTION\_COMMAND!} \\
\protect\hyperlink{unitree-sdk2-cpp-robot-a2-sportclient-trajectory-1}{\passthrough{\lstinline!trajectory!}}
&
\passthrough{\lstinline!def trajectory(self, path: Sequence[PathPoint], feedback\_mode: int = ..., external\_x: float = ..., external\_y: float = ..., external\_yaw: float = ...) -> int!}
& \passthrough{\lstinline!SIGNATURE\_ONLY!} &
\passthrough{\lstinline!MOTION\_COMMAND!} \\
\protect\hyperlink{unitree-sdk2-cpp-robot-a2-sportclient-set-auto-recovery-1}{\passthrough{\lstinline!set\_auto\_recovery!}}
&
\passthrough{\lstinline!def set\_auto\_recovery(self, switch\_on: int) -> int!}
& \passthrough{\lstinline!SIGNATURE\_ONLY!} &
\passthrough{\lstinline!MOTION\_COMMAND!} \\
\protect\hyperlink{unitree-sdk2-cpp-robot-a2-sportclient-get-state-1}{\passthrough{\lstinline!get\_state!}}
&
\passthrough{\lstinline!def get\_state(self) -> tuple[int, dict[str, str]]!}
& \passthrough{\lstinline!AVAILABLE!} &
\passthrough{\lstinline!READ\_ONLY!} \\
\bottomrule
\end{longtable}

\hypertarget{unitree-sdk2-cpp-robot-a2-sportclient-dunder-init-1}{%
\paragraph{\texorpdfstring{\texttt{unitree\_sdk2\_cpp.robot.a2.SportClient.\_\_init\_\_}}{unitree\_sdk2\_cpp.robot.a2.SportClient.\_\_init\_\_}}\label{unitree-sdk2-cpp-robot-a2-sportclient-dunder-init-1}}

初始化 \passthrough{\lstinline!SportClient!}
实例。是否能实际构造取决于下方可用性状态。

\textbf{签名}

\begin{lstlisting}[language=Python]
def __init__(self) -> None
\end{lstlisting}

\textbf{可用性与安全性}

\textbf{\passthrough{\lstinline!AVAILABLE!} /
\passthrough{\lstinline!CONSTRUCTION!}}：当前绑定源码已实现；实际执行条件仍取决于平台和对应
SDK 环境。

\textbf{参数}

无显式参数。实例方法中的 \passthrough{\lstinline!self!} 由 Python
自动传入。

\textbf{返回值}

\begin{longtable}[]{@{}
  >{\raggedright\arraybackslash}p{(\columnwidth - 4\tabcolsep) * \real{0.18}}
  >{\raggedright\arraybackslash}p{(\columnwidth - 4\tabcolsep) * \real{0.20}}
  >{\raggedright\arraybackslash}p{(\columnwidth - 4\tabcolsep) * \real{0.62}}@{}}
\toprule
位置 & 类型 & 含义 \\
\midrule
\endhead
无 & \passthrough{\lstinline!None!} &
\passthrough{\lstinline!\_\_init\_\_!}
本身不返回值；调用类对象时，在构造函数可用的前提下得到该类实例。 \\
\bottomrule
\end{longtable}

\textbf{对应 C++}

\begin{itemize}
\tightlist
\item
  类：\passthrough{\lstinline!unitree::robot::a2::SportClient!}
\item
  签名：\passthrough{\lstinline!SportClient()!}
\item
  绑定策略：\passthrough{\lstinline!未单独标注!}
\item
  声明位置：\passthrough{\lstinline!include/unitree/robot/a2/sport/sport\_client.hpp:77!}
\end{itemize}

\textbf{说明}

\begin{itemize}
\tightlist
\item
  示例是局部调用片段；\passthrough{\lstinline!obj!}
  和参数变量表示已经按上层业务校验并准备好的对象。
\end{itemize}

\textbf{用法}

\begin{lstlisting}[language=Python]
value = SportClient()
\end{lstlisting}

\hypertarget{unitree-sdk2-cpp-robot-a2-sportclient-init-1}{%
\paragraph{\texorpdfstring{\texttt{unitree\_sdk2\_cpp.robot.a2.SportClient.init}}{unitree\_sdk2\_cpp.robot.a2.SportClient.init}}\label{unitree-sdk2-cpp-robot-a2-sportclient-init-1}}

初始化当前 SDK 对象所需的底层通道或服务资源。

\textbf{签名}

\begin{lstlisting}[language=Python]
def init(self) -> None
\end{lstlisting}

\textbf{可用性与安全性}

\textbf{\passthrough{\lstinline!AVAILABLE!} /
\passthrough{\lstinline!INITIALIZATION!}}：当前绑定源码已实现；实际执行条件仍取决于平台和对应
SDK 环境。

\textbf{参数}

无显式参数。实例方法中的 \passthrough{\lstinline!self!} 由 Python
自动传入。

\textbf{返回值}

\begin{longtable}[]{@{}
  >{\raggedright\arraybackslash}p{(\columnwidth - 4\tabcolsep) * \real{0.18}}
  >{\raggedright\arraybackslash}p{(\columnwidth - 4\tabcolsep) * \real{0.20}}
  >{\raggedright\arraybackslash}p{(\columnwidth - 4\tabcolsep) * \real{0.62}}@{}}
\toprule
位置 & 类型 & 含义 \\
\midrule
\endhead
无 & \passthrough{\lstinline!None!} & 该方法不返回 Python
值；失败可能表现为异常或底层状态变化。 \\
\bottomrule
\end{longtable}

\textbf{对应 C++}

\begin{itemize}
\tightlist
\item
  类：\passthrough{\lstinline!unitree::robot::a2::SportClient!}
\item
  签名：\passthrough{\lstinline!Init()!}
\item
  绑定策略：\passthrough{\lstinline!DIRECT!}
\item
  声明位置：\passthrough{\lstinline!include/unitree/robot/a2/sport/sport\_client.hpp:81!}
\end{itemize}

\textbf{说明}

\begin{itemize}
\tightlist
\item
  示例是局部调用片段；\passthrough{\lstinline!obj!}
  和参数变量表示已经按上层业务校验并准备好的对象。
\end{itemize}

\textbf{用法}

\begin{lstlisting}[language=Python]
obj.init()
\end{lstlisting}

\hypertarget{unitree-sdk2-cpp-robot-a2-sportclient-damp-1}{%
\paragraph{\texorpdfstring{\texttt{unitree\_sdk2\_cpp.robot.a2.SportClient.damp}}{unitree\_sdk2\_cpp.robot.a2.SportClient.damp}}\label{unitree-sdk2-cpp-robot-a2-sportclient-damp-1}}

对应 C++ SDK 操作
\passthrough{\lstinline!Damp()!}。上游头文件没有可直接生成的业务说明时，本参考只保证签名映射，精确语义需查目标型号协议。

\textbf{签名}

\begin{lstlisting}[language=Python]
def damp(self) -> int
\end{lstlisting}

\textbf{可用性与安全性}

\textbf{\passthrough{\lstinline!SIGNATURE\_ONLY!} /
\passthrough{\lstinline!MOTION\_COMMAND!}}：仅用于补全和静态检查。当前二进制没有该
Python 方法；不得按可执行运动接口使用。

\textbf{参数}

无显式参数。实例方法中的 \passthrough{\lstinline!self!} 由 Python
自动传入。

\textbf{返回值}

\begin{longtable}[]{@{}
  >{\raggedright\arraybackslash}p{(\columnwidth - 4\tabcolsep) * \real{0.18}}
  >{\raggedright\arraybackslash}p{(\columnwidth - 4\tabcolsep) * \real{0.20}}
  >{\raggedright\arraybackslash}p{(\columnwidth - 4\tabcolsep) * \real{0.62}}@{}}
\toprule
位置 & 类型 & 含义 \\
\midrule
\endhead
返回值 & \passthrough{\lstinline!int!} & SDK 状态码；Unitree 示例通常以
\passthrough{\lstinline!0!} 表示成功，非零值需按具体服务定义解释。 \\
\bottomrule
\end{longtable}

\textbf{对应 C++}

\begin{itemize}
\tightlist
\item
  类：\passthrough{\lstinline!unitree::robot::a2::SportClient!}
\item
  签名：\passthrough{\lstinline!Damp()!}
\item
  绑定策略：\passthrough{\lstinline!DIRECT!}
\item
  声明位置：\passthrough{\lstinline!include/unitree/robot/a2/sport/sport\_client.hpp:112!}
\end{itemize}

\textbf{说明}

\begin{itemize}
\tightlist
\item
  示例是局部调用片段；\passthrough{\lstinline!obj!}
  和参数变量表示已经按上层业务校验并准备好的对象。
\item
  该条目可以被编辑器和类型检查器识别，但当前运行时可能在属性查找阶段直接失败。
\item
  该方法属于运动命令。方法名包含
  \passthrough{\lstinline!stop!}、\passthrough{\lstinline!damp!} 或
  \passthrough{\lstinline!zero!} 也不代表它是独立物理急停。
\end{itemize}

\textbf{用法}

\begin{lstlisting}[language=Python]
from typing import TYPE_CHECKING

if TYPE_CHECKING:
    def planned_damp(obj: SportClient) -> int:
        return obj.damp()
\end{lstlisting}

\begin{quote}
{[}!CAUTION{]} 上面的 \passthrough{\lstinline!TYPE\_CHECKING!}
示例只表达计划调用形态。该分支在普通 Python
运行时为假，不代表当前扩展能执行此接口。
\end{quote}

\hypertarget{unitree-sdk2-cpp-robot-a2-sportclient-balance-stand-1}{%
\paragraph{\texorpdfstring{\texttt{unitree\_sdk2\_cpp.robot.a2.SportClient.balance\_stand}}{unitree\_sdk2\_cpp.robot.a2.SportClient.balance\_stand}}\label{unitree-sdk2-cpp-robot-a2-sportclient-balance-stand-1}}

对应 C++ SDK 操作
\passthrough{\lstinline!BalanceStand()!}。上游头文件没有可直接生成的业务说明时，本参考只保证签名映射，精确语义需查目标型号协议。

\textbf{签名}

\begin{lstlisting}[language=Python]
def balance_stand(self) -> int
\end{lstlisting}

\textbf{可用性与安全性}

\textbf{\passthrough{\lstinline!SIGNATURE\_ONLY!} /
\passthrough{\lstinline!MOTION\_COMMAND!}}：仅用于补全和静态检查。当前二进制没有该
Python 方法；不得按可执行运动接口使用。

\textbf{参数}

无显式参数。实例方法中的 \passthrough{\lstinline!self!} 由 Python
自动传入。

\textbf{返回值}

\begin{longtable}[]{@{}
  >{\raggedright\arraybackslash}p{(\columnwidth - 4\tabcolsep) * \real{0.18}}
  >{\raggedright\arraybackslash}p{(\columnwidth - 4\tabcolsep) * \real{0.20}}
  >{\raggedright\arraybackslash}p{(\columnwidth - 4\tabcolsep) * \real{0.62}}@{}}
\toprule
位置 & 类型 & 含义 \\
\midrule
\endhead
返回值 & \passthrough{\lstinline!int!} & SDK 状态码；Unitree 示例通常以
\passthrough{\lstinline!0!} 表示成功，非零值需按具体服务定义解释。 \\
\bottomrule
\end{longtable}

\textbf{对应 C++}

\begin{itemize}
\tightlist
\item
  类：\passthrough{\lstinline!unitree::robot::a2::SportClient!}
\item
  签名：\passthrough{\lstinline!BalanceStand()!}
\item
  绑定策略：\passthrough{\lstinline!DIRECT!}
\item
  声明位置：\passthrough{\lstinline!include/unitree/robot/a2/sport/sport\_client.hpp:118!}
\end{itemize}

\textbf{说明}

\begin{itemize}
\tightlist
\item
  示例是局部调用片段；\passthrough{\lstinline!obj!}
  和参数变量表示已经按上层业务校验并准备好的对象。
\item
  该条目可以被编辑器和类型检查器识别，但当前运行时可能在属性查找阶段直接失败。
\item
  该方法属于运动命令。方法名包含
  \passthrough{\lstinline!stop!}、\passthrough{\lstinline!damp!} 或
  \passthrough{\lstinline!zero!} 也不代表它是独立物理急停。
\end{itemize}

\textbf{用法}

\begin{lstlisting}[language=Python]
from typing import TYPE_CHECKING

if TYPE_CHECKING:
    def planned_balance_stand(obj: SportClient) -> int:
        return obj.balance_stand()
\end{lstlisting}

\begin{quote}
{[}!CAUTION{]} 上面的 \passthrough{\lstinline!TYPE\_CHECKING!}
示例只表达计划调用形态。该分支在普通 Python
运行时为假，不代表当前扩展能执行此接口。
\end{quote}

\hypertarget{unitree-sdk2-cpp-robot-a2-sportclient-stop-move-1}{%
\paragraph{\texorpdfstring{\texttt{unitree\_sdk2\_cpp.robot.a2.SportClient.stop\_move}}{unitree\_sdk2\_cpp.robot.a2.SportClient.stop\_move}}\label{unitree-sdk2-cpp-robot-a2-sportclient-stop-move-1}}

请求停止对应 SDK 操作。该名称不等同于经过验证的物理急停。

\textbf{签名}

\begin{lstlisting}[language=Python]
def stop_move(self) -> int
\end{lstlisting}

\textbf{可用性与安全性}

\textbf{\passthrough{\lstinline!SIGNATURE\_ONLY!} /
\passthrough{\lstinline!MOTION\_COMMAND!}}：仅用于补全和静态检查。当前二进制没有该
Python 方法；不得按可执行运动接口使用。

\textbf{参数}

无显式参数。实例方法中的 \passthrough{\lstinline!self!} 由 Python
自动传入。

\textbf{返回值}

\begin{longtable}[]{@{}
  >{\raggedright\arraybackslash}p{(\columnwidth - 4\tabcolsep) * \real{0.18}}
  >{\raggedright\arraybackslash}p{(\columnwidth - 4\tabcolsep) * \real{0.20}}
  >{\raggedright\arraybackslash}p{(\columnwidth - 4\tabcolsep) * \real{0.62}}@{}}
\toprule
位置 & 类型 & 含义 \\
\midrule
\endhead
返回值 & \passthrough{\lstinline!int!} & SDK 状态码；Unitree 示例通常以
\passthrough{\lstinline!0!} 表示成功，非零值需按具体服务定义解释。 \\
\bottomrule
\end{longtable}

\textbf{对应 C++}

\begin{itemize}
\tightlist
\item
  类：\passthrough{\lstinline!unitree::robot::a2::SportClient!}
\item
  签名：\passthrough{\lstinline!StopMove()!}
\item
  绑定策略：\passthrough{\lstinline!DIRECT!}
\item
  声明位置：\passthrough{\lstinline!include/unitree/robot/a2/sport/sport\_client.hpp:123!}
\end{itemize}

\textbf{说明}

\begin{itemize}
\tightlist
\item
  示例是局部调用片段；\passthrough{\lstinline!obj!}
  和参数变量表示已经按上层业务校验并准备好的对象。
\item
  该条目可以被编辑器和类型检查器识别，但当前运行时可能在属性查找阶段直接失败。
\item
  该方法属于运动命令。方法名包含
  \passthrough{\lstinline!stop!}、\passthrough{\lstinline!damp!} 或
  \passthrough{\lstinline!zero!} 也不代表它是独立物理急停。
\end{itemize}

\textbf{用法}

\begin{lstlisting}[language=Python]
from typing import TYPE_CHECKING

if TYPE_CHECKING:
    def planned_stop_move(obj: SportClient) -> int:
        return obj.stop_move()
\end{lstlisting}

\begin{quote}
{[}!CAUTION{]} 上面的 \passthrough{\lstinline!TYPE\_CHECKING!}
示例只表达计划调用形态。该分支在普通 Python
运行时为假，不代表当前扩展能执行此接口。
\end{quote}

\hypertarget{unitree-sdk2-cpp-robot-a2-sportclient-stand-up-1}{%
\paragraph{\texorpdfstring{\texttt{unitree\_sdk2\_cpp.robot.a2.SportClient.stand\_up}}{unitree\_sdk2\_cpp.robot.a2.SportClient.stand\_up}}\label{unitree-sdk2-cpp-robot-a2-sportclient-stand-up-1}}

对应 C++ SDK 操作
\passthrough{\lstinline!StandUp()!}。上游头文件没有可直接生成的业务说明时，本参考只保证签名映射，精确语义需查目标型号协议。

\textbf{签名}

\begin{lstlisting}[language=Python]
def stand_up(self) -> int
\end{lstlisting}

\textbf{可用性与安全性}

\textbf{\passthrough{\lstinline!SIGNATURE\_ONLY!} /
\passthrough{\lstinline!MOTION\_COMMAND!}}：仅用于补全和静态检查。当前二进制没有该
Python 方法；不得按可执行运动接口使用。

\textbf{参数}

无显式参数。实例方法中的 \passthrough{\lstinline!self!} 由 Python
自动传入。

\textbf{返回值}

\begin{longtable}[]{@{}
  >{\raggedright\arraybackslash}p{(\columnwidth - 4\tabcolsep) * \real{0.18}}
  >{\raggedright\arraybackslash}p{(\columnwidth - 4\tabcolsep) * \real{0.20}}
  >{\raggedright\arraybackslash}p{(\columnwidth - 4\tabcolsep) * \real{0.62}}@{}}
\toprule
位置 & 类型 & 含义 \\
\midrule
\endhead
返回值 & \passthrough{\lstinline!int!} & SDK 状态码；Unitree 示例通常以
\passthrough{\lstinline!0!} 表示成功，非零值需按具体服务定义解释。 \\
\bottomrule
\end{longtable}

\textbf{对应 C++}

\begin{itemize}
\tightlist
\item
  类：\passthrough{\lstinline!unitree::robot::a2::SportClient!}
\item
  签名：\passthrough{\lstinline!StandUp()!}
\item
  绑定策略：\passthrough{\lstinline!DIRECT!}
\item
  声明位置：\passthrough{\lstinline!include/unitree/robot/a2/sport/sport\_client.hpp:129!}
\end{itemize}

\textbf{说明}

\begin{itemize}
\tightlist
\item
  示例是局部调用片段；\passthrough{\lstinline!obj!}
  和参数变量表示已经按上层业务校验并准备好的对象。
\item
  该条目可以被编辑器和类型检查器识别，但当前运行时可能在属性查找阶段直接失败。
\item
  该方法属于运动命令。方法名包含
  \passthrough{\lstinline!stop!}、\passthrough{\lstinline!damp!} 或
  \passthrough{\lstinline!zero!} 也不代表它是独立物理急停。
\end{itemize}

\textbf{用法}

\begin{lstlisting}[language=Python]
from typing import TYPE_CHECKING

if TYPE_CHECKING:
    def planned_stand_up(obj: SportClient) -> int:
        return obj.stand_up()
\end{lstlisting}

\begin{quote}
{[}!CAUTION{]} 上面的 \passthrough{\lstinline!TYPE\_CHECKING!}
示例只表达计划调用形态。该分支在普通 Python
运行时为假，不代表当前扩展能执行此接口。
\end{quote}

\hypertarget{unitree-sdk2-cpp-robot-a2-sportclient-stand-down-1}{%
\paragraph{\texorpdfstring{\texttt{unitree\_sdk2\_cpp.robot.a2.SportClient.stand\_down}}{unitree\_sdk2\_cpp.robot.a2.SportClient.stand\_down}}\label{unitree-sdk2-cpp-robot-a2-sportclient-stand-down-1}}

对应 C++ SDK 操作
\passthrough{\lstinline!StandDown()!}。上游头文件没有可直接生成的业务说明时，本参考只保证签名映射，精确语义需查目标型号协议。

\textbf{签名}

\begin{lstlisting}[language=Python]
def stand_down(self) -> int
\end{lstlisting}

\textbf{可用性与安全性}

\textbf{\passthrough{\lstinline!SIGNATURE\_ONLY!} /
\passthrough{\lstinline!MOTION\_COMMAND!}}：仅用于补全和静态检查。当前二进制没有该
Python 方法；不得按可执行运动接口使用。

\textbf{参数}

无显式参数。实例方法中的 \passthrough{\lstinline!self!} 由 Python
自动传入。

\textbf{返回值}

\begin{longtable}[]{@{}
  >{\raggedright\arraybackslash}p{(\columnwidth - 4\tabcolsep) * \real{0.18}}
  >{\raggedright\arraybackslash}p{(\columnwidth - 4\tabcolsep) * \real{0.20}}
  >{\raggedright\arraybackslash}p{(\columnwidth - 4\tabcolsep) * \real{0.62}}@{}}
\toprule
位置 & 类型 & 含义 \\
\midrule
\endhead
返回值 & \passthrough{\lstinline!int!} & SDK 状态码；Unitree 示例通常以
\passthrough{\lstinline!0!} 表示成功，非零值需按具体服务定义解释。 \\
\bottomrule
\end{longtable}

\textbf{对应 C++}

\begin{itemize}
\tightlist
\item
  类：\passthrough{\lstinline!unitree::robot::a2::SportClient!}
\item
  签名：\passthrough{\lstinline!StandDown()!}
\item
  绑定策略：\passthrough{\lstinline!DIRECT!}
\item
  声明位置：\passthrough{\lstinline!include/unitree/robot/a2/sport/sport\_client.hpp:135!}
\end{itemize}

\textbf{说明}

\begin{itemize}
\tightlist
\item
  示例是局部调用片段；\passthrough{\lstinline!obj!}
  和参数变量表示已经按上层业务校验并准备好的对象。
\item
  该条目可以被编辑器和类型检查器识别，但当前运行时可能在属性查找阶段直接失败。
\item
  该方法属于运动命令。方法名包含
  \passthrough{\lstinline!stop!}、\passthrough{\lstinline!damp!} 或
  \passthrough{\lstinline!zero!} 也不代表它是独立物理急停。
\end{itemize}

\textbf{用法}

\begin{lstlisting}[language=Python]
from typing import TYPE_CHECKING

if TYPE_CHECKING:
    def planned_stand_down(obj: SportClient) -> int:
        return obj.stand_down()
\end{lstlisting}

\begin{quote}
{[}!CAUTION{]} 上面的 \passthrough{\lstinline!TYPE\_CHECKING!}
示例只表达计划调用形态。该分支在普通 Python
运行时为假，不代表当前扩展能执行此接口。
\end{quote}

\hypertarget{unitree-sdk2-cpp-robot-a2-sportclient-recovery-stand-1}{%
\paragraph{\texorpdfstring{\texttt{unitree\_sdk2\_cpp.robot.a2.SportClient.recovery\_stand}}{unitree\_sdk2\_cpp.robot.a2.SportClient.recovery\_stand}}\label{unitree-sdk2-cpp-robot-a2-sportclient-recovery-stand-1}}

对应 C++ SDK 操作
\passthrough{\lstinline!RecoveryStand()!}。上游头文件没有可直接生成的业务说明时，本参考只保证签名映射，精确语义需查目标型号协议。

\textbf{签名}

\begin{lstlisting}[language=Python]
def recovery_stand(self) -> int
\end{lstlisting}

\textbf{可用性与安全性}

\textbf{\passthrough{\lstinline!SIGNATURE\_ONLY!} /
\passthrough{\lstinline!MOTION\_COMMAND!}}：仅用于补全和静态检查。当前二进制没有该
Python 方法；不得按可执行运动接口使用。

\textbf{参数}

无显式参数。实例方法中的 \passthrough{\lstinline!self!} 由 Python
自动传入。

\textbf{返回值}

\begin{longtable}[]{@{}
  >{\raggedright\arraybackslash}p{(\columnwidth - 4\tabcolsep) * \real{0.18}}
  >{\raggedright\arraybackslash}p{(\columnwidth - 4\tabcolsep) * \real{0.20}}
  >{\raggedright\arraybackslash}p{(\columnwidth - 4\tabcolsep) * \real{0.62}}@{}}
\toprule
位置 & 类型 & 含义 \\
\midrule
\endhead
返回值 & \passthrough{\lstinline!int!} & SDK 状态码；Unitree 示例通常以
\passthrough{\lstinline!0!} 表示成功，非零值需按具体服务定义解释。 \\
\bottomrule
\end{longtable}

\textbf{对应 C++}

\begin{itemize}
\tightlist
\item
  类：\passthrough{\lstinline!unitree::robot::a2::SportClient!}
\item
  签名：\passthrough{\lstinline!RecoveryStand()!}
\item
  绑定策略：\passthrough{\lstinline!DIRECT!}
\item
  声明位置：\passthrough{\lstinline!include/unitree/robot/a2/sport/sport\_client.hpp:141!}
\end{itemize}

\textbf{说明}

\begin{itemize}
\tightlist
\item
  示例是局部调用片段；\passthrough{\lstinline!obj!}
  和参数变量表示已经按上层业务校验并准备好的对象。
\item
  该条目可以被编辑器和类型检查器识别，但当前运行时可能在属性查找阶段直接失败。
\item
  该方法属于运动命令。方法名包含
  \passthrough{\lstinline!stop!}、\passthrough{\lstinline!damp!} 或
  \passthrough{\lstinline!zero!} 也不代表它是独立物理急停。
\end{itemize}

\textbf{用法}

\begin{lstlisting}[language=Python]
from typing import TYPE_CHECKING

if TYPE_CHECKING:
    def planned_recovery_stand(obj: SportClient) -> int:
        return obj.recovery_stand()
\end{lstlisting}

\begin{quote}
{[}!CAUTION{]} 上面的 \passthrough{\lstinline!TYPE\_CHECKING!}
示例只表达计划调用形态。该分支在普通 Python
运行时为假，不代表当前扩展能执行此接口。
\end{quote}

\hypertarget{unitree-sdk2-cpp-robot-a2-sportclient-euler-1}{%
\paragraph{\texorpdfstring{\texttt{unitree\_sdk2\_cpp.robot.a2.SportClient.euler}}{unitree\_sdk2\_cpp.robot.a2.SportClient.euler}}\label{unitree-sdk2-cpp-robot-a2-sportclient-euler-1}}

对应 C++ SDK 操作
\passthrough{\lstinline!Euler(float, float, float)!}。上游头文件没有可直接生成的业务说明时，本参考只保证签名映射，精确语义需查目标型号协议。

\textbf{签名}

\begin{lstlisting}[language=Python]
def euler(self, roll: float, pitch: float, yaw: float) -> int
\end{lstlisting}

\textbf{可用性与安全性}

\textbf{\passthrough{\lstinline!SIGNATURE\_ONLY!} /
\passthrough{\lstinline!MOTION\_COMMAND!}}：仅用于补全和静态检查。当前二进制没有该
Python 方法；不得按可执行运动接口使用。

\textbf{参数}

\begin{longtable}[]{@{}
  >{\raggedright\arraybackslash}p{(\columnwidth - 6\tabcolsep) * \real{0.18}}
  >{\raggedright\arraybackslash}p{(\columnwidth - 6\tabcolsep) * \real{0.24}}
  >{\raggedright\arraybackslash}p{(\columnwidth - 6\tabcolsep) * \real{0.18}}
  >{\raggedright\arraybackslash}p{(\columnwidth - 6\tabcolsep) * \real{0.40}}@{}}
\toprule
名称 & Python 类型 & 默认值 & 含义 \\
\midrule
\endhead
\passthrough{\lstinline!roll!} & \passthrough{\lstinline!float!} & 必填
& 传给该接口的 \passthrough{\lstinline!roll!} 参数，Python 类型为
\passthrough{\lstinline!float!}。精确含义、单位、范围和枚举值需查目标型号对应头文件与协议。
对应 C++ 参数 \passthrough{\lstinline!roll: float!}。
底层为浮点数；类型本身不说明单位、坐标系或业务安全范围。 \\
\passthrough{\lstinline!pitch!} & \passthrough{\lstinline!float!} & 必填
& 传给该接口的 \passthrough{\lstinline!pitch!} 参数，Python 类型为
\passthrough{\lstinline!float!}。精确含义、单位、范围和枚举值需查目标型号对应头文件与协议。
对应 C++ 参数 \passthrough{\lstinline!pitch: float!}。
底层为浮点数；类型本身不说明单位、坐标系或业务安全范围。 \\
\passthrough{\lstinline!yaw!} & \passthrough{\lstinline!float!} & 必填 &
偏航角或偏航目标参数。参考系、单位和范围必须查目标型号协议。 对应 C++
参数 \passthrough{\lstinline!yaw: float!}。
底层为浮点数；类型本身不说明单位、坐标系或业务安全范围。 \\
\bottomrule
\end{longtable}

\textbf{返回值}

\begin{longtable}[]{@{}
  >{\raggedright\arraybackslash}p{(\columnwidth - 4\tabcolsep) * \real{0.18}}
  >{\raggedright\arraybackslash}p{(\columnwidth - 4\tabcolsep) * \real{0.20}}
  >{\raggedright\arraybackslash}p{(\columnwidth - 4\tabcolsep) * \real{0.62}}@{}}
\toprule
位置 & 类型 & 含义 \\
\midrule
\endhead
返回值 & \passthrough{\lstinline!int!} & SDK 状态码；Unitree 示例通常以
\passthrough{\lstinline!0!} 表示成功，非零值需按具体服务定义解释。 \\
\bottomrule
\end{longtable}

\textbf{对应 C++}

\begin{itemize}
\tightlist
\item
  类：\passthrough{\lstinline!unitree::robot::a2::SportClient!}
\item
  签名：\passthrough{\lstinline!Euler(float, float, float)!}
\item
  绑定策略：\passthrough{\lstinline!DIRECT!}
\item
  声明位置：\passthrough{\lstinline!include/unitree/robot/a2/sport/sport\_client.hpp:147!}
\end{itemize}

\textbf{说明}

\begin{itemize}
\tightlist
\item
  示例是局部调用片段；\passthrough{\lstinline!obj!}
  和参数变量表示已经按上层业务校验并准备好的对象。
\item
  该条目可以被编辑器和类型检查器识别，但当前运行时可能在属性查找阶段直接失败。
\item
  该方法属于运动命令。方法名包含
  \passthrough{\lstinline!stop!}、\passthrough{\lstinline!damp!} 或
  \passthrough{\lstinline!zero!} 也不代表它是独立物理急停。
\end{itemize}

\textbf{用法}

\begin{lstlisting}[language=Python]
from typing import TYPE_CHECKING

if TYPE_CHECKING:
    def planned_euler(obj: SportClient, roll: float, pitch: float, yaw: float) -> int:
        return obj.euler(roll=roll, pitch=pitch, yaw=yaw)
\end{lstlisting}

\begin{quote}
{[}!CAUTION{]} 上面的 \passthrough{\lstinline!TYPE\_CHECKING!}
示例只表达计划调用形态。该分支在普通 Python
运行时为假，不代表当前扩展能执行此接口。
\end{quote}

\hypertarget{unitree-sdk2-cpp-robot-a2-sportclient-move-1}{%
\paragraph{\texorpdfstring{\texttt{unitree\_sdk2\_cpp.robot.a2.SportClient.move}}{unitree\_sdk2\_cpp.robot.a2.SportClient.move}}\label{unitree-sdk2-cpp-robot-a2-sportclient-move-1}}

对应 C++ SDK 操作
\passthrough{\lstinline!Move(float, float, float)!}。上游头文件没有可直接生成的业务说明时，本参考只保证签名映射，精确语义需查目标型号协议。

\textbf{签名}

\begin{lstlisting}[language=Python]
def move(self, vx: float, vy: float, vyaw: float) -> int
\end{lstlisting}

\textbf{可用性与安全性}

\textbf{\passthrough{\lstinline!SIGNATURE\_ONLY!} /
\passthrough{\lstinline!MOTION\_COMMAND!}}：仅用于补全和静态检查。当前二进制没有该
Python 方法；不得按可执行运动接口使用。

\textbf{参数}

\begin{longtable}[]{@{}
  >{\raggedright\arraybackslash}p{(\columnwidth - 6\tabcolsep) * \real{0.18}}
  >{\raggedright\arraybackslash}p{(\columnwidth - 6\tabcolsep) * \real{0.24}}
  >{\raggedright\arraybackslash}p{(\columnwidth - 6\tabcolsep) * \real{0.18}}
  >{\raggedright\arraybackslash}p{(\columnwidth - 6\tabcolsep) * \real{0.40}}@{}}
\toprule
名称 & Python 类型 & 默认值 & 含义 \\
\midrule
\endhead
\passthrough{\lstinline!vx!} & \passthrough{\lstinline!float!} & 必填 &
X 方向速度参数。坐标系、单位、符号和安全范围必须查目标型号运动协议。
对应 C++ 参数 \passthrough{\lstinline!vx: float!}。
底层为浮点数；类型本身不说明单位、坐标系或业务安全范围。 \\
\passthrough{\lstinline!vy!} & \passthrough{\lstinline!float!} & 必填 &
Y 方向速度参数。坐标系、单位、符号和安全范围必须查目标型号运动协议。
对应 C++ 参数 \passthrough{\lstinline!vy: float!}。
底层为浮点数；类型本身不说明单位、坐标系或业务安全范围。 \\
\passthrough{\lstinline!vyaw!} & \passthrough{\lstinline!float!} & 必填
& 偏航角速度参数。单位、符号和安全范围必须查目标型号运动协议。 对应 C++
参数 \passthrough{\lstinline!vyaw: float!}。
底层为浮点数；类型本身不说明单位、坐标系或业务安全范围。 \\
\bottomrule
\end{longtable}

\textbf{返回值}

\begin{longtable}[]{@{}
  >{\raggedright\arraybackslash}p{(\columnwidth - 4\tabcolsep) * \real{0.18}}
  >{\raggedright\arraybackslash}p{(\columnwidth - 4\tabcolsep) * \real{0.20}}
  >{\raggedright\arraybackslash}p{(\columnwidth - 4\tabcolsep) * \real{0.62}}@{}}
\toprule
位置 & 类型 & 含义 \\
\midrule
\endhead
返回值 & \passthrough{\lstinline!int!} & SDK 状态码；Unitree 示例通常以
\passthrough{\lstinline!0!} 表示成功，非零值需按具体服务定义解释。 \\
\bottomrule
\end{longtable}

\textbf{对应 C++}

\begin{itemize}
\tightlist
\item
  类：\passthrough{\lstinline!unitree::robot::a2::SportClient!}
\item
  签名：\passthrough{\lstinline!Move(float, float, float)!}
\item
  绑定策略：\passthrough{\lstinline!DIRECT!}
\item
  声明位置：\passthrough{\lstinline!include/unitree/robot/a2/sport/sport\_client.hpp:158!}
\end{itemize}

\textbf{说明}

\begin{itemize}
\tightlist
\item
  示例是局部调用片段；\passthrough{\lstinline!obj!}
  和参数变量表示已经按上层业务校验并准备好的对象。
\item
  该条目可以被编辑器和类型检查器识别，但当前运行时可能在属性查找阶段直接失败。
\item
  该方法属于运动命令。方法名包含
  \passthrough{\lstinline!stop!}、\passthrough{\lstinline!damp!} 或
  \passthrough{\lstinline!zero!} 也不代表它是独立物理急停。
\end{itemize}

\textbf{用法}

\begin{lstlisting}[language=Python]
from typing import TYPE_CHECKING

if TYPE_CHECKING:
    def planned_move(obj: SportClient, vx: float, vy: float, vyaw: float) -> int:
        return obj.move(vx=vx, vy=vy, vyaw=vyaw)
\end{lstlisting}

\begin{quote}
{[}!CAUTION{]} 上面的 \passthrough{\lstinline!TYPE\_CHECKING!}
示例只表达计划调用形态。该分支在普通 Python
运行时为假，不代表当前扩展能执行此接口。
\end{quote}

\hypertarget{unitree-sdk2-cpp-robot-a2-sportclient-switch-gait-1}{%
\paragraph{\texorpdfstring{\texttt{unitree\_sdk2\_cpp.robot.a2.SportClient.switch\_gait}}{unitree\_sdk2\_cpp.robot.a2.SportClient.switch\_gait}}\label{unitree-sdk2-cpp-robot-a2-sportclient-switch-gait-1}}

对应 C++ SDK 操作
\passthrough{\lstinline!SwitchGait(int)!}。上游头文件没有可直接生成的业务说明时，本参考只保证签名映射，精确语义需查目标型号协议。

\textbf{签名}

\begin{lstlisting}[language=Python]
def switch_gait(self, gait_type: int) -> int
\end{lstlisting}

\textbf{可用性与安全性}

\textbf{\passthrough{\lstinline!SIGNATURE\_ONLY!} /
\passthrough{\lstinline!MOTION\_COMMAND!}}：仅用于补全和静态检查。当前二进制没有该
Python 方法；不得按可执行运动接口使用。

\textbf{参数}

\begin{longtable}[]{@{}
  >{\raggedright\arraybackslash}p{(\columnwidth - 6\tabcolsep) * \real{0.18}}
  >{\raggedright\arraybackslash}p{(\columnwidth - 6\tabcolsep) * \real{0.24}}
  >{\raggedright\arraybackslash}p{(\columnwidth - 6\tabcolsep) * \real{0.18}}
  >{\raggedright\arraybackslash}p{(\columnwidth - 6\tabcolsep) * \real{0.40}}@{}}
\toprule
名称 & Python 类型 & 默认值 & 含义 \\
\midrule
\endhead
\passthrough{\lstinline!gait\_type!} & \passthrough{\lstinline!int!} &
必填 & 传给该接口的 \passthrough{\lstinline!gait!}
\passthrough{\lstinline!type!} 参数，Python 类型为
\passthrough{\lstinline!int!}。精确含义、单位、范围和枚举值需查目标型号对应头文件与协议。
对应 C++ 参数 \passthrough{\lstinline!gait\_type: int!}。 \\
\bottomrule
\end{longtable}

\textbf{返回值}

\begin{longtable}[]{@{}
  >{\raggedright\arraybackslash}p{(\columnwidth - 4\tabcolsep) * \real{0.18}}
  >{\raggedright\arraybackslash}p{(\columnwidth - 4\tabcolsep) * \real{0.20}}
  >{\raggedright\arraybackslash}p{(\columnwidth - 4\tabcolsep) * \real{0.62}}@{}}
\toprule
位置 & 类型 & 含义 \\
\midrule
\endhead
返回值 & \passthrough{\lstinline!int!} & SDK 状态码；Unitree 示例通常以
\passthrough{\lstinline!0!} 表示成功，非零值需按具体服务定义解释。 \\
\bottomrule
\end{longtable}

\textbf{对应 C++}

\begin{itemize}
\tightlist
\item
  类：\passthrough{\lstinline!unitree::robot::a2::SportClient!}
\item
  签名：\passthrough{\lstinline!SwitchGait(int)!}
\item
  绑定策略：\passthrough{\lstinline!DIRECT!}
\item
  声明位置：\passthrough{\lstinline!include/unitree/robot/a2/sport/sport\_client.hpp:169!}
\end{itemize}

\textbf{说明}

\begin{itemize}
\tightlist
\item
  示例是局部调用片段；\passthrough{\lstinline!obj!}
  和参数变量表示已经按上层业务校验并准备好的对象。
\item
  该条目可以被编辑器和类型检查器识别，但当前运行时可能在属性查找阶段直接失败。
\item
  该方法属于运动命令。方法名包含
  \passthrough{\lstinline!stop!}、\passthrough{\lstinline!damp!} 或
  \passthrough{\lstinline!zero!} 也不代表它是独立物理急停。
\end{itemize}

\textbf{用法}

\begin{lstlisting}[language=Python]
from typing import TYPE_CHECKING

if TYPE_CHECKING:
    def planned_switch_gait(obj: SportClient, gait_type: int) -> int:
        return obj.switch_gait(gait_type=gait_type)
\end{lstlisting}

\begin{quote}
{[}!CAUTION{]} 上面的 \passthrough{\lstinline!TYPE\_CHECKING!}
示例只表达计划调用形态。该分支在普通 Python
运行时为假，不代表当前扩展能执行此接口。
\end{quote}

\hypertarget{unitree-sdk2-cpp-robot-a2-sportclient-body-height-1}{%
\paragraph{\texorpdfstring{\texttt{unitree\_sdk2\_cpp.robot.a2.SportClient.body\_height}}{unitree\_sdk2\_cpp.robot.a2.SportClient.body\_height}}\label{unitree-sdk2-cpp-robot-a2-sportclient-body-height-1}}

对应 C++ SDK 操作
\passthrough{\lstinline!BodyHeight(float)!}。上游头文件没有可直接生成的业务说明时，本参考只保证签名映射，精确语义需查目标型号协议。

\textbf{签名}

\begin{lstlisting}[language=Python]
def body_height(self, height: float) -> int
\end{lstlisting}

\textbf{可用性与安全性}

\textbf{\passthrough{\lstinline!SIGNATURE\_ONLY!} /
\passthrough{\lstinline!MOTION\_COMMAND!}}：仅用于补全和静态检查。当前二进制没有该
Python 方法；不得按可执行运动接口使用。

\textbf{参数}

\begin{longtable}[]{@{}
  >{\raggedright\arraybackslash}p{(\columnwidth - 6\tabcolsep) * \real{0.18}}
  >{\raggedright\arraybackslash}p{(\columnwidth - 6\tabcolsep) * \real{0.24}}
  >{\raggedright\arraybackslash}p{(\columnwidth - 6\tabcolsep) * \real{0.18}}
  >{\raggedright\arraybackslash}p{(\columnwidth - 6\tabcolsep) * \real{0.40}}@{}}
\toprule
名称 & Python 类型 & 默认值 & 含义 \\
\midrule
\endhead
\passthrough{\lstinline!height!} & \passthrough{\lstinline!float!} &
必填 & 传给该接口的 高度 参数，Python 类型为
\passthrough{\lstinline!float!}。精确含义、单位、范围和枚举值需查目标型号对应头文件与协议。
对应 C++ 参数 \passthrough{\lstinline!height: float!}。
底层为浮点数；类型本身不说明单位、坐标系或业务安全范围。 \\
\bottomrule
\end{longtable}

\textbf{返回值}

\begin{longtable}[]{@{}
  >{\raggedright\arraybackslash}p{(\columnwidth - 4\tabcolsep) * \real{0.18}}
  >{\raggedright\arraybackslash}p{(\columnwidth - 4\tabcolsep) * \real{0.20}}
  >{\raggedright\arraybackslash}p{(\columnwidth - 4\tabcolsep) * \real{0.62}}@{}}
\toprule
位置 & 类型 & 含义 \\
\midrule
\endhead
返回值 & \passthrough{\lstinline!int!} & SDK 状态码；Unitree 示例通常以
\passthrough{\lstinline!0!} 表示成功，非零值需按具体服务定义解释。 \\
\bottomrule
\end{longtable}

\textbf{对应 C++}

\begin{itemize}
\tightlist
\item
  类：\passthrough{\lstinline!unitree::robot::a2::SportClient!}
\item
  签名：\passthrough{\lstinline!BodyHeight(float)!}
\item
  绑定策略：\passthrough{\lstinline!DIRECT!}
\item
  声明位置：\passthrough{\lstinline!include/unitree/robot/a2/sport/sport\_client.hpp:178!}
\end{itemize}

\textbf{说明}

\begin{itemize}
\tightlist
\item
  示例是局部调用片段；\passthrough{\lstinline!obj!}
  和参数变量表示已经按上层业务校验并准备好的对象。
\item
  该条目可以被编辑器和类型检查器识别，但当前运行时可能在属性查找阶段直接失败。
\item
  该方法属于运动命令。方法名包含
  \passthrough{\lstinline!stop!}、\passthrough{\lstinline!damp!} 或
  \passthrough{\lstinline!zero!} 也不代表它是独立物理急停。
\end{itemize}

\textbf{用法}

\begin{lstlisting}[language=Python]
from typing import TYPE_CHECKING

if TYPE_CHECKING:
    def planned_body_height(obj: SportClient, height: float) -> int:
        return obj.body_height(height=height)
\end{lstlisting}

\begin{quote}
{[}!CAUTION{]} 上面的 \passthrough{\lstinline!TYPE\_CHECKING!}
示例只表达计划调用形态。该分支在普通 Python
运行时为假，不代表当前扩展能执行此接口。
\end{quote}

\hypertarget{unitree-sdk2-cpp-robot-a2-sportclient-speed-level-1}{%
\paragraph{\texorpdfstring{\texttt{unitree\_sdk2\_cpp.robot.a2.SportClient.speed\_level}}{unitree\_sdk2\_cpp.robot.a2.SportClient.speed\_level}}\label{unitree-sdk2-cpp-robot-a2-sportclient-speed-level-1}}

对应 C++ SDK 操作
\passthrough{\lstinline!SpeedLevel(int)!}。上游头文件没有可直接生成的业务说明时，本参考只保证签名映射，精确语义需查目标型号协议。

\textbf{签名}

\begin{lstlisting}[language=Python]
def speed_level(self, level: int) -> int
\end{lstlisting}

\textbf{可用性与安全性}

\textbf{\passthrough{\lstinline!SIGNATURE\_ONLY!} /
\passthrough{\lstinline!MOTION\_COMMAND!}}：仅用于补全和静态检查。当前二进制没有该
Python 方法；不得按可执行运动接口使用。

\textbf{参数}

\begin{longtable}[]{@{}
  >{\raggedright\arraybackslash}p{(\columnwidth - 6\tabcolsep) * \real{0.18}}
  >{\raggedright\arraybackslash}p{(\columnwidth - 6\tabcolsep) * \real{0.24}}
  >{\raggedright\arraybackslash}p{(\columnwidth - 6\tabcolsep) * \real{0.18}}
  >{\raggedright\arraybackslash}p{(\columnwidth - 6\tabcolsep) * \real{0.40}}@{}}
\toprule
名称 & Python 类型 & 默认值 & 含义 \\
\midrule
\endhead
\passthrough{\lstinline!level!} & \passthrough{\lstinline!int!} & 必填 &
传给该接口的 \passthrough{\lstinline!level!} 参数，Python 类型为
\passthrough{\lstinline!int!}。精确含义、单位、范围和枚举值需查目标型号对应头文件与协议。
对应 C++ 参数 \passthrough{\lstinline!level: int!}。 \\
\bottomrule
\end{longtable}

\textbf{返回值}

\begin{longtable}[]{@{}
  >{\raggedright\arraybackslash}p{(\columnwidth - 4\tabcolsep) * \real{0.18}}
  >{\raggedright\arraybackslash}p{(\columnwidth - 4\tabcolsep) * \real{0.20}}
  >{\raggedright\arraybackslash}p{(\columnwidth - 4\tabcolsep) * \real{0.62}}@{}}
\toprule
位置 & 类型 & 含义 \\
\midrule
\endhead
返回值 & \passthrough{\lstinline!int!} & SDK 状态码；Unitree 示例通常以
\passthrough{\lstinline!0!} 表示成功，非零值需按具体服务定义解释。 \\
\bottomrule
\end{longtable}

\textbf{对应 C++}

\begin{itemize}
\tightlist
\item
  类：\passthrough{\lstinline!unitree::robot::a2::SportClient!}
\item
  签名：\passthrough{\lstinline!SpeedLevel(int)!}
\item
  绑定策略：\passthrough{\lstinline!DIRECT!}
\item
  声明位置：\passthrough{\lstinline!include/unitree/robot/a2/sport/sport\_client.hpp:187!}
\end{itemize}

\textbf{说明}

\begin{itemize}
\tightlist
\item
  示例是局部调用片段；\passthrough{\lstinline!obj!}
  和参数变量表示已经按上层业务校验并准备好的对象。
\item
  该条目可以被编辑器和类型检查器识别，但当前运行时可能在属性查找阶段直接失败。
\item
  该方法属于运动命令。方法名包含
  \passthrough{\lstinline!stop!}、\passthrough{\lstinline!damp!} 或
  \passthrough{\lstinline!zero!} 也不代表它是独立物理急停。
\end{itemize}

\textbf{用法}

\begin{lstlisting}[language=Python]
from typing import TYPE_CHECKING

if TYPE_CHECKING:
    def planned_speed_level(obj: SportClient, level: int) -> int:
        return obj.speed_level(level=level)
\end{lstlisting}

\begin{quote}
{[}!CAUTION{]} 上面的 \passthrough{\lstinline!TYPE\_CHECKING!}
示例只表达计划调用形态。该分支在普通 Python
运行时为假，不代表当前扩展能执行此接口。
\end{quote}

\hypertarget{unitree-sdk2-cpp-robot-a2-sportclient-body-position-1}{%
\paragraph{\texorpdfstring{\texttt{unitree\_sdk2\_cpp.robot.a2.SportClient.body\_position}}{unitree\_sdk2\_cpp.robot.a2.SportClient.body\_position}}\label{unitree-sdk2-cpp-robot-a2-sportclient-body-position-1}}

对应 C++ SDK 操作
\passthrough{\lstinline!BodyPosition(float, float, float, float)!}。上游头文件没有可直接生成的业务说明时，本参考只保证签名映射，精确语义需查目标型号协议。

\textbf{签名}

\begin{lstlisting}[language=Python]
def body_position(self, x: float, y: float, z: float, yaw: float) -> int
\end{lstlisting}

\textbf{可用性与安全性}

\textbf{\passthrough{\lstinline!SIGNATURE\_ONLY!} /
\passthrough{\lstinline!MOTION\_COMMAND!}}：仅用于补全和静态检查。当前二进制没有该
Python 方法；不得按可执行运动接口使用。

\textbf{参数}

\begin{longtable}[]{@{}
  >{\raggedright\arraybackslash}p{(\columnwidth - 6\tabcolsep) * \real{0.18}}
  >{\raggedright\arraybackslash}p{(\columnwidth - 6\tabcolsep) * \real{0.24}}
  >{\raggedright\arraybackslash}p{(\columnwidth - 6\tabcolsep) * \real{0.18}}
  >{\raggedright\arraybackslash}p{(\columnwidth - 6\tabcolsep) * \real{0.40}}@{}}
\toprule
名称 & Python 类型 & 默认值 & 含义 \\
\midrule
\endhead
\passthrough{\lstinline!x!} & \passthrough{\lstinline!float!} & 必填 & X
方向位置或分量参数。参考系、单位和范围必须查当前方法及目标型号协议。
对应 C++ 参数 \passthrough{\lstinline!x: float!}。
底层为浮点数；类型本身不说明单位、坐标系或业务安全范围。 \\
\passthrough{\lstinline!y!} & \passthrough{\lstinline!float!} & 必填 & Y
方向位置或分量参数。参考系、单位和范围必须查当前方法及目标型号协议。
对应 C++ 参数 \passthrough{\lstinline!y: float!}。
底层为浮点数；类型本身不说明单位、坐标系或业务安全范围。 \\
\passthrough{\lstinline!z!} & \passthrough{\lstinline!float!} & 必填 & Z
方向位置或分量参数。参考系、单位和范围必须查当前方法及目标型号协议。
对应 C++ 参数 \passthrough{\lstinline!z: float!}。
底层为浮点数；类型本身不说明单位、坐标系或业务安全范围。 \\
\passthrough{\lstinline!yaw!} & \passthrough{\lstinline!float!} & 必填 &
偏航角或偏航目标参数。参考系、单位和范围必须查目标型号协议。 对应 C++
参数 \passthrough{\lstinline!yaw: float!}。
底层为浮点数；类型本身不说明单位、坐标系或业务安全范围。 \\
\bottomrule
\end{longtable}

\textbf{返回值}

\begin{longtable}[]{@{}
  >{\raggedright\arraybackslash}p{(\columnwidth - 4\tabcolsep) * \real{0.18}}
  >{\raggedright\arraybackslash}p{(\columnwidth - 4\tabcolsep) * \real{0.20}}
  >{\raggedright\arraybackslash}p{(\columnwidth - 4\tabcolsep) * \real{0.62}}@{}}
\toprule
位置 & 类型 & 含义 \\
\midrule
\endhead
返回值 & \passthrough{\lstinline!int!} & SDK 状态码；Unitree 示例通常以
\passthrough{\lstinline!0!} 表示成功，非零值需按具体服务定义解释。 \\
\bottomrule
\end{longtable}

\textbf{对应 C++}

\begin{itemize}
\tightlist
\item
  类：\passthrough{\lstinline!unitree::robot::a2::SportClient!}
\item
  签名：\passthrough{\lstinline!BodyPosition(float, float, float, float)!}
\item
  绑定策略：\passthrough{\lstinline!DIRECT!}
\item
  声明位置：\passthrough{\lstinline!include/unitree/robot/a2/sport/sport\_client.hpp:196!}
\end{itemize}

\textbf{说明}

\begin{itemize}
\tightlist
\item
  示例是局部调用片段；\passthrough{\lstinline!obj!}
  和参数变量表示已经按上层业务校验并准备好的对象。
\item
  该条目可以被编辑器和类型检查器识别，但当前运行时可能在属性查找阶段直接失败。
\item
  该方法属于运动命令。方法名包含
  \passthrough{\lstinline!stop!}、\passthrough{\lstinline!damp!} 或
  \passthrough{\lstinline!zero!} 也不代表它是独立物理急停。
\end{itemize}

\textbf{用法}

\begin{lstlisting}[language=Python]
from typing import TYPE_CHECKING

if TYPE_CHECKING:
    def planned_body_position(obj: SportClient, x: float, y: float, z: float, yaw: float) -> int:
        return obj.body_position(x=x, y=y, z=z, yaw=yaw)
\end{lstlisting}

\begin{quote}
{[}!CAUTION{]} 上面的 \passthrough{\lstinline!TYPE\_CHECKING!}
示例只表达计划调用形态。该分支在普通 Python
运行时为假，不代表当前扩展能执行此接口。
\end{quote}

\hypertarget{unitree-sdk2-cpp-robot-a2-sportclient-left-side-gait-1}{%
\paragraph{\texorpdfstring{\texttt{unitree\_sdk2\_cpp.robot.a2.SportClient.left\_side\_gait}}{unitree\_sdk2\_cpp.robot.a2.SportClient.left\_side\_gait}}\label{unitree-sdk2-cpp-robot-a2-sportclient-left-side-gait-1}}

对应 C++ SDK 操作
\passthrough{\lstinline!LeftSideGait(int)!}。上游头文件没有可直接生成的业务说明时，本参考只保证签名映射，精确语义需查目标型号协议。

\textbf{签名}

\begin{lstlisting}[language=Python]
def left_side_gait(self, enter: int) -> int
\end{lstlisting}

\textbf{可用性与安全性}

\textbf{\passthrough{\lstinline!SIGNATURE\_ONLY!} /
\passthrough{\lstinline!MOTION\_COMMAND!}}：仅用于补全和静态检查。当前二进制没有该
Python 方法；不得按可执行运动接口使用。

\textbf{参数}

\begin{longtable}[]{@{}
  >{\raggedright\arraybackslash}p{(\columnwidth - 6\tabcolsep) * \real{0.18}}
  >{\raggedright\arraybackslash}p{(\columnwidth - 6\tabcolsep) * \real{0.24}}
  >{\raggedright\arraybackslash}p{(\columnwidth - 6\tabcolsep) * \real{0.18}}
  >{\raggedright\arraybackslash}p{(\columnwidth - 6\tabcolsep) * \real{0.40}}@{}}
\toprule
名称 & Python 类型 & 默认值 & 含义 \\
\midrule
\endhead
\passthrough{\lstinline!enter!} & \passthrough{\lstinline!int!} & 必填 &
传给该接口的 \passthrough{\lstinline!enter!} 参数，Python 类型为
\passthrough{\lstinline!int!}。精确含义、单位、范围和枚举值需查目标型号对应头文件与协议。
对应 C++ 参数 \passthrough{\lstinline!enter: int!}。 \\
\bottomrule
\end{longtable}

\textbf{返回值}

\begin{longtable}[]{@{}
  >{\raggedright\arraybackslash}p{(\columnwidth - 4\tabcolsep) * \real{0.18}}
  >{\raggedright\arraybackslash}p{(\columnwidth - 4\tabcolsep) * \real{0.20}}
  >{\raggedright\arraybackslash}p{(\columnwidth - 4\tabcolsep) * \real{0.62}}@{}}
\toprule
位置 & 类型 & 含义 \\
\midrule
\endhead
返回值 & \passthrough{\lstinline!int!} & SDK 状态码；Unitree 示例通常以
\passthrough{\lstinline!0!} 表示成功，非零值需按具体服务定义解释。 \\
\bottomrule
\end{longtable}

\textbf{对应 C++}

\begin{itemize}
\tightlist
\item
  类：\passthrough{\lstinline!unitree::robot::a2::SportClient!}
\item
  签名：\passthrough{\lstinline!LeftSideGait(int)!}
\item
  绑定策略：\passthrough{\lstinline!DIRECT!}
\item
  声明位置：\passthrough{\lstinline!include/unitree/robot/a2/sport/sport\_client.hpp:208!}
\end{itemize}

\textbf{说明}

\begin{itemize}
\tightlist
\item
  示例是局部调用片段；\passthrough{\lstinline!obj!}
  和参数变量表示已经按上层业务校验并准备好的对象。
\item
  该条目可以被编辑器和类型检查器识别，但当前运行时可能在属性查找阶段直接失败。
\item
  该方法属于运动命令。方法名包含
  \passthrough{\lstinline!stop!}、\passthrough{\lstinline!damp!} 或
  \passthrough{\lstinline!zero!} 也不代表它是独立物理急停。
\end{itemize}

\textbf{用法}

\begin{lstlisting}[language=Python]
from typing import TYPE_CHECKING

if TYPE_CHECKING:
    def planned_left_side_gait(obj: SportClient, enter: int) -> int:
        return obj.left_side_gait(enter=enter)
\end{lstlisting}

\begin{quote}
{[}!CAUTION{]} 上面的 \passthrough{\lstinline!TYPE\_CHECKING!}
示例只表达计划调用形态。该分支在普通 Python
运行时为假，不代表当前扩展能执行此接口。
\end{quote}

\hypertarget{unitree-sdk2-cpp-robot-a2-sportclient-right-side-gait-1}{%
\paragraph{\texorpdfstring{\texttt{unitree\_sdk2\_cpp.robot.a2.SportClient.right\_side\_gait}}{unitree\_sdk2\_cpp.robot.a2.SportClient.right\_side\_gait}}\label{unitree-sdk2-cpp-robot-a2-sportclient-right-side-gait-1}}

对应 C++ SDK 操作
\passthrough{\lstinline!RightSideGait(int)!}。上游头文件没有可直接生成的业务说明时，本参考只保证签名映射，精确语义需查目标型号协议。

\textbf{签名}

\begin{lstlisting}[language=Python]
def right_side_gait(self, enter: int) -> int
\end{lstlisting}

\textbf{可用性与安全性}

\textbf{\passthrough{\lstinline!SIGNATURE\_ONLY!} /
\passthrough{\lstinline!MOTION\_COMMAND!}}：仅用于补全和静态检查。当前二进制没有该
Python 方法；不得按可执行运动接口使用。

\textbf{参数}

\begin{longtable}[]{@{}
  >{\raggedright\arraybackslash}p{(\columnwidth - 6\tabcolsep) * \real{0.18}}
  >{\raggedright\arraybackslash}p{(\columnwidth - 6\tabcolsep) * \real{0.24}}
  >{\raggedright\arraybackslash}p{(\columnwidth - 6\tabcolsep) * \real{0.18}}
  >{\raggedright\arraybackslash}p{(\columnwidth - 6\tabcolsep) * \real{0.40}}@{}}
\toprule
名称 & Python 类型 & 默认值 & 含义 \\
\midrule
\endhead
\passthrough{\lstinline!enter!} & \passthrough{\lstinline!int!} & 必填 &
传给该接口的 \passthrough{\lstinline!enter!} 参数，Python 类型为
\passthrough{\lstinline!int!}。精确含义、单位、范围和枚举值需查目标型号对应头文件与协议。
对应 C++ 参数 \passthrough{\lstinline!enter: int!}。 \\
\bottomrule
\end{longtable}

\textbf{返回值}

\begin{longtable}[]{@{}
  >{\raggedright\arraybackslash}p{(\columnwidth - 4\tabcolsep) * \real{0.18}}
  >{\raggedright\arraybackslash}p{(\columnwidth - 4\tabcolsep) * \real{0.20}}
  >{\raggedright\arraybackslash}p{(\columnwidth - 4\tabcolsep) * \real{0.62}}@{}}
\toprule
位置 & 类型 & 含义 \\
\midrule
\endhead
返回值 & \passthrough{\lstinline!int!} & SDK 状态码；Unitree 示例通常以
\passthrough{\lstinline!0!} 表示成功，非零值需按具体服务定义解释。 \\
\bottomrule
\end{longtable}

\textbf{对应 C++}

\begin{itemize}
\tightlist
\item
  类：\passthrough{\lstinline!unitree::robot::a2::SportClient!}
\item
  签名：\passthrough{\lstinline!RightSideGait(int)!}
\item
  绑定策略：\passthrough{\lstinline!DIRECT!}
\item
  声明位置：\passthrough{\lstinline!include/unitree/robot/a2/sport/sport\_client.hpp:217!}
\end{itemize}

\textbf{说明}

\begin{itemize}
\tightlist
\item
  示例是局部调用片段；\passthrough{\lstinline!obj!}
  和参数变量表示已经按上层业务校验并准备好的对象。
\item
  该条目可以被编辑器和类型检查器识别，但当前运行时可能在属性查找阶段直接失败。
\item
  该方法属于运动命令。方法名包含
  \passthrough{\lstinline!stop!}、\passthrough{\lstinline!damp!} 或
  \passthrough{\lstinline!zero!} 也不代表它是独立物理急停。
\end{itemize}

\textbf{用法}

\begin{lstlisting}[language=Python]
from typing import TYPE_CHECKING

if TYPE_CHECKING:
    def planned_right_side_gait(obj: SportClient, enter: int) -> int:
        return obj.right_side_gait(enter=enter)
\end{lstlisting}

\begin{quote}
{[}!CAUTION{]} 上面的 \passthrough{\lstinline!TYPE\_CHECKING!}
示例只表达计划调用形态。该分支在普通 Python
运行时为假，不代表当前扩展能执行此接口。
\end{quote}

\hypertarget{unitree-sdk2-cpp-robot-a2-sportclient-hand-stand-1}{%
\paragraph{\texorpdfstring{\texttt{unitree\_sdk2\_cpp.robot.a2.SportClient.hand\_stand}}{unitree\_sdk2\_cpp.robot.a2.SportClient.hand\_stand}}\label{unitree-sdk2-cpp-robot-a2-sportclient-hand-stand-1}}

对应 C++ SDK 操作
\passthrough{\lstinline!HandStand(int)!}。上游头文件没有可直接生成的业务说明时，本参考只保证签名映射，精确语义需查目标型号协议。

\textbf{签名}

\begin{lstlisting}[language=Python]
def hand_stand(self, enter: int) -> int
\end{lstlisting}

\textbf{可用性与安全性}

\textbf{\passthrough{\lstinline!SIGNATURE\_ONLY!} /
\passthrough{\lstinline!MOTION\_COMMAND!}}：仅用于补全和静态检查。当前二进制没有该
Python 方法；不得按可执行运动接口使用。

\textbf{参数}

\begin{longtable}[]{@{}
  >{\raggedright\arraybackslash}p{(\columnwidth - 6\tabcolsep) * \real{0.18}}
  >{\raggedright\arraybackslash}p{(\columnwidth - 6\tabcolsep) * \real{0.24}}
  >{\raggedright\arraybackslash}p{(\columnwidth - 6\tabcolsep) * \real{0.18}}
  >{\raggedright\arraybackslash}p{(\columnwidth - 6\tabcolsep) * \real{0.40}}@{}}
\toprule
名称 & Python 类型 & 默认值 & 含义 \\
\midrule
\endhead
\passthrough{\lstinline!enter!} & \passthrough{\lstinline!int!} & 必填 &
传给该接口的 \passthrough{\lstinline!enter!} 参数，Python 类型为
\passthrough{\lstinline!int!}。精确含义、单位、范围和枚举值需查目标型号对应头文件与协议。
对应 C++ 参数 \passthrough{\lstinline!enter: int!}。 \\
\bottomrule
\end{longtable}

\textbf{返回值}

\begin{longtable}[]{@{}
  >{\raggedright\arraybackslash}p{(\columnwidth - 4\tabcolsep) * \real{0.18}}
  >{\raggedright\arraybackslash}p{(\columnwidth - 4\tabcolsep) * \real{0.20}}
  >{\raggedright\arraybackslash}p{(\columnwidth - 4\tabcolsep) * \real{0.62}}@{}}
\toprule
位置 & 类型 & 含义 \\
\midrule
\endhead
返回值 & \passthrough{\lstinline!int!} & SDK 状态码；Unitree 示例通常以
\passthrough{\lstinline!0!} 表示成功，非零值需按具体服务定义解释。 \\
\bottomrule
\end{longtable}

\textbf{对应 C++}

\begin{itemize}
\tightlist
\item
  类：\passthrough{\lstinline!unitree::robot::a2::SportClient!}
\item
  签名：\passthrough{\lstinline!HandStand(int)!}
\item
  绑定策略：\passthrough{\lstinline!DIRECT!}
\item
  声明位置：\passthrough{\lstinline!include/unitree/robot/a2/sport/sport\_client.hpp:226!}
\end{itemize}

\textbf{说明}

\begin{itemize}
\tightlist
\item
  示例是局部调用片段；\passthrough{\lstinline!obj!}
  和参数变量表示已经按上层业务校验并准备好的对象。
\item
  该条目可以被编辑器和类型检查器识别，但当前运行时可能在属性查找阶段直接失败。
\item
  该方法属于运动命令。方法名包含
  \passthrough{\lstinline!stop!}、\passthrough{\lstinline!damp!} 或
  \passthrough{\lstinline!zero!} 也不代表它是独立物理急停。
\end{itemize}

\textbf{用法}

\begin{lstlisting}[language=Python]
from typing import TYPE_CHECKING

if TYPE_CHECKING:
    def planned_hand_stand(obj: SportClient, enter: int) -> int:
        return obj.hand_stand(enter=enter)
\end{lstlisting}

\begin{quote}
{[}!CAUTION{]} 上面的 \passthrough{\lstinline!TYPE\_CHECKING!}
示例只表达计划调用形态。该分支在普通 Python
运行时为假，不代表当前扩展能执行此接口。
\end{quote}

\hypertarget{unitree-sdk2-cpp-robot-a2-sportclient-biped-stand-1}{%
\paragraph{\texorpdfstring{\texttt{unitree\_sdk2\_cpp.robot.a2.SportClient.biped\_stand}}{unitree\_sdk2\_cpp.robot.a2.SportClient.biped\_stand}}\label{unitree-sdk2-cpp-robot-a2-sportclient-biped-stand-1}}

对应 C++ SDK 操作
\passthrough{\lstinline!BipedStand(int)!}。上游头文件没有可直接生成的业务说明时，本参考只保证签名映射，精确语义需查目标型号协议。

\textbf{签名}

\begin{lstlisting}[language=Python]
def biped_stand(self, enter: int) -> int
\end{lstlisting}

\textbf{可用性与安全性}

\textbf{\passthrough{\lstinline!SIGNATURE\_ONLY!} /
\passthrough{\lstinline!MOTION\_COMMAND!}}：仅用于补全和静态检查。当前二进制没有该
Python 方法；不得按可执行运动接口使用。

\textbf{参数}

\begin{longtable}[]{@{}
  >{\raggedright\arraybackslash}p{(\columnwidth - 6\tabcolsep) * \real{0.18}}
  >{\raggedright\arraybackslash}p{(\columnwidth - 6\tabcolsep) * \real{0.24}}
  >{\raggedright\arraybackslash}p{(\columnwidth - 6\tabcolsep) * \real{0.18}}
  >{\raggedright\arraybackslash}p{(\columnwidth - 6\tabcolsep) * \real{0.40}}@{}}
\toprule
名称 & Python 类型 & 默认值 & 含义 \\
\midrule
\endhead
\passthrough{\lstinline!enter!} & \passthrough{\lstinline!int!} & 必填 &
传给该接口的 \passthrough{\lstinline!enter!} 参数，Python 类型为
\passthrough{\lstinline!int!}。精确含义、单位、范围和枚举值需查目标型号对应头文件与协议。
对应 C++ 参数 \passthrough{\lstinline!enter: int!}。 \\
\bottomrule
\end{longtable}

\textbf{返回值}

\begin{longtable}[]{@{}
  >{\raggedright\arraybackslash}p{(\columnwidth - 4\tabcolsep) * \real{0.18}}
  >{\raggedright\arraybackslash}p{(\columnwidth - 4\tabcolsep) * \real{0.20}}
  >{\raggedright\arraybackslash}p{(\columnwidth - 4\tabcolsep) * \real{0.62}}@{}}
\toprule
位置 & 类型 & 含义 \\
\midrule
\endhead
返回值 & \passthrough{\lstinline!int!} & SDK 状态码；Unitree 示例通常以
\passthrough{\lstinline!0!} 表示成功，非零值需按具体服务定义解释。 \\
\bottomrule
\end{longtable}

\textbf{对应 C++}

\begin{itemize}
\tightlist
\item
  类：\passthrough{\lstinline!unitree::robot::a2::SportClient!}
\item
  签名：\passthrough{\lstinline!BipedStand(int)!}
\item
  绑定策略：\passthrough{\lstinline!DIRECT!}
\item
  声明位置：\passthrough{\lstinline!include/unitree/robot/a2/sport/sport\_client.hpp:235!}
\end{itemize}

\textbf{说明}

\begin{itemize}
\tightlist
\item
  示例是局部调用片段；\passthrough{\lstinline!obj!}
  和参数变量表示已经按上层业务校验并准备好的对象。
\item
  该条目可以被编辑器和类型检查器识别，但当前运行时可能在属性查找阶段直接失败。
\item
  该方法属于运动命令。方法名包含
  \passthrough{\lstinline!stop!}、\passthrough{\lstinline!damp!} 或
  \passthrough{\lstinline!zero!} 也不代表它是独立物理急停。
\end{itemize}

\textbf{用法}

\begin{lstlisting}[language=Python]
from typing import TYPE_CHECKING

if TYPE_CHECKING:
    def planned_biped_stand(obj: SportClient, enter: int) -> int:
        return obj.biped_stand(enter=enter)
\end{lstlisting}

\begin{quote}
{[}!CAUTION{]} 上面的 \passthrough{\lstinline!TYPE\_CHECKING!}
示例只表达计划调用形态。该分支在普通 Python
运行时为假，不代表当前扩展能执行此接口。
\end{quote}

\hypertarget{unitree-sdk2-cpp-robot-a2-sportclient-front-flip-1}{%
\paragraph{\texorpdfstring{\texttt{unitree\_sdk2\_cpp.robot.a2.SportClient.front\_flip}}{unitree\_sdk2\_cpp.robot.a2.SportClient.front\_flip}}\label{unitree-sdk2-cpp-robot-a2-sportclient-front-flip-1}}

对应 C++ SDK 操作
\passthrough{\lstinline!FrontFlip()!}。上游头文件没有可直接生成的业务说明时，本参考只保证签名映射，精确语义需查目标型号协议。

\textbf{签名}

\begin{lstlisting}[language=Python]
def front_flip(self) -> int
\end{lstlisting}

\textbf{可用性与安全性}

\textbf{\passthrough{\lstinline!SIGNATURE\_ONLY!} /
\passthrough{\lstinline!MOTION\_COMMAND!}}：仅用于补全和静态检查。当前二进制没有该
Python 方法；不得按可执行运动接口使用。

\textbf{参数}

无显式参数。实例方法中的 \passthrough{\lstinline!self!} 由 Python
自动传入。

\textbf{返回值}

\begin{longtable}[]{@{}
  >{\raggedright\arraybackslash}p{(\columnwidth - 4\tabcolsep) * \real{0.18}}
  >{\raggedright\arraybackslash}p{(\columnwidth - 4\tabcolsep) * \real{0.20}}
  >{\raggedright\arraybackslash}p{(\columnwidth - 4\tabcolsep) * \real{0.62}}@{}}
\toprule
位置 & 类型 & 含义 \\
\midrule
\endhead
返回值 & \passthrough{\lstinline!int!} & SDK 状态码；Unitree 示例通常以
\passthrough{\lstinline!0!} 表示成功，非零值需按具体服务定义解释。 \\
\bottomrule
\end{longtable}

\textbf{对应 C++}

\begin{itemize}
\tightlist
\item
  类：\passthrough{\lstinline!unitree::robot::a2::SportClient!}
\item
  签名：\passthrough{\lstinline!FrontFlip()!}
\item
  绑定策略：\passthrough{\lstinline!DIRECT!}
\item
  声明位置：\passthrough{\lstinline!include/unitree/robot/a2/sport/sport\_client.hpp:244!}
\end{itemize}

\textbf{说明}

\begin{itemize}
\tightlist
\item
  示例是局部调用片段；\passthrough{\lstinline!obj!}
  和参数变量表示已经按上层业务校验并准备好的对象。
\item
  该条目可以被编辑器和类型检查器识别，但当前运行时可能在属性查找阶段直接失败。
\item
  该方法属于运动命令。方法名包含
  \passthrough{\lstinline!stop!}、\passthrough{\lstinline!damp!} 或
  \passthrough{\lstinline!zero!} 也不代表它是独立物理急停。
\end{itemize}

\textbf{用法}

\begin{lstlisting}[language=Python]
from typing import TYPE_CHECKING

if TYPE_CHECKING:
    def planned_front_flip(obj: SportClient) -> int:
        return obj.front_flip()
\end{lstlisting}

\begin{quote}
{[}!CAUTION{]} 上面的 \passthrough{\lstinline!TYPE\_CHECKING!}
示例只表达计划调用形态。该分支在普通 Python
运行时为假，不代表当前扩展能执行此接口。
\end{quote}

\hypertarget{unitree-sdk2-cpp-robot-a2-sportclient-back-flip-1}{%
\paragraph{\texorpdfstring{\texttt{unitree\_sdk2\_cpp.robot.a2.SportClient.back\_flip}}{unitree\_sdk2\_cpp.robot.a2.SportClient.back\_flip}}\label{unitree-sdk2-cpp-robot-a2-sportclient-back-flip-1}}

对应 C++ SDK 操作
\passthrough{\lstinline!BackFlip()!}。上游头文件没有可直接生成的业务说明时，本参考只保证签名映射，精确语义需查目标型号协议。

\textbf{签名}

\begin{lstlisting}[language=Python]
def back_flip(self) -> int
\end{lstlisting}

\textbf{可用性与安全性}

\textbf{\passthrough{\lstinline!SIGNATURE\_ONLY!} /
\passthrough{\lstinline!MOTION\_COMMAND!}}：仅用于补全和静态检查。当前二进制没有该
Python 方法；不得按可执行运动接口使用。

\textbf{参数}

无显式参数。实例方法中的 \passthrough{\lstinline!self!} 由 Python
自动传入。

\textbf{返回值}

\begin{longtable}[]{@{}
  >{\raggedright\arraybackslash}p{(\columnwidth - 4\tabcolsep) * \real{0.18}}
  >{\raggedright\arraybackslash}p{(\columnwidth - 4\tabcolsep) * \real{0.20}}
  >{\raggedright\arraybackslash}p{(\columnwidth - 4\tabcolsep) * \real{0.62}}@{}}
\toprule
位置 & 类型 & 含义 \\
\midrule
\endhead
返回值 & \passthrough{\lstinline!int!} & SDK 状态码；Unitree 示例通常以
\passthrough{\lstinline!0!} 表示成功，非零值需按具体服务定义解释。 \\
\bottomrule
\end{longtable}

\textbf{对应 C++}

\begin{itemize}
\tightlist
\item
  类：\passthrough{\lstinline!unitree::robot::a2::SportClient!}
\item
  签名：\passthrough{\lstinline!BackFlip()!}
\item
  绑定策略：\passthrough{\lstinline!DIRECT!}
\item
  声明位置：\passthrough{\lstinline!include/unitree/robot/a2/sport/sport\_client.hpp:250!}
\end{itemize}

\textbf{说明}

\begin{itemize}
\tightlist
\item
  示例是局部调用片段；\passthrough{\lstinline!obj!}
  和参数变量表示已经按上层业务校验并准备好的对象。
\item
  该条目可以被编辑器和类型检查器识别，但当前运行时可能在属性查找阶段直接失败。
\item
  该方法属于运动命令。方法名包含
  \passthrough{\lstinline!stop!}、\passthrough{\lstinline!damp!} 或
  \passthrough{\lstinline!zero!} 也不代表它是独立物理急停。
\end{itemize}

\textbf{用法}

\begin{lstlisting}[language=Python]
from typing import TYPE_CHECKING

if TYPE_CHECKING:
    def planned_back_flip(obj: SportClient) -> int:
        return obj.back_flip()
\end{lstlisting}

\begin{quote}
{[}!CAUTION{]} 上面的 \passthrough{\lstinline!TYPE\_CHECKING!}
示例只表达计划调用形态。该分支在普通 Python
运行时为假，不代表当前扩展能执行此接口。
\end{quote}

\hypertarget{unitree-sdk2-cpp-robot-a2-sportclient-reset-estimator-1}{%
\paragraph{\texorpdfstring{\texttt{unitree\_sdk2\_cpp.robot.a2.SportClient.reset\_estimator}}{unitree\_sdk2\_cpp.robot.a2.SportClient.reset\_estimator}}\label{unitree-sdk2-cpp-robot-a2-sportclient-reset-estimator-1}}

对应 C++ SDK 操作
\passthrough{\lstinline!ResetEstimator()!}。上游头文件没有可直接生成的业务说明时，本参考只保证签名映射，精确语义需查目标型号协议。

\textbf{签名}

\begin{lstlisting}[language=Python]
def reset_estimator(self) -> int
\end{lstlisting}

\textbf{可用性与安全性}

\textbf{\passthrough{\lstinline!SIGNATURE\_ONLY!} /
\passthrough{\lstinline!MOTION\_COMMAND!}}：仅用于补全和静态检查。当前二进制没有该
Python 方法；不得按可执行运动接口使用。

\textbf{参数}

无显式参数。实例方法中的 \passthrough{\lstinline!self!} 由 Python
自动传入。

\textbf{返回值}

\begin{longtable}[]{@{}
  >{\raggedright\arraybackslash}p{(\columnwidth - 4\tabcolsep) * \real{0.18}}
  >{\raggedright\arraybackslash}p{(\columnwidth - 4\tabcolsep) * \real{0.20}}
  >{\raggedright\arraybackslash}p{(\columnwidth - 4\tabcolsep) * \real{0.62}}@{}}
\toprule
位置 & 类型 & 含义 \\
\midrule
\endhead
返回值 & \passthrough{\lstinline!int!} & SDK 状态码；Unitree 示例通常以
\passthrough{\lstinline!0!} 表示成功，非零值需按具体服务定义解释。 \\
\bottomrule
\end{longtable}

\textbf{对应 C++}

\begin{itemize}
\tightlist
\item
  类：\passthrough{\lstinline!unitree::robot::a2::SportClient!}
\item
  签名：\passthrough{\lstinline!ResetEstimator()!}
\item
  绑定策略：\passthrough{\lstinline!DIRECT!}
\item
  声明位置：\passthrough{\lstinline!include/unitree/robot/a2/sport/sport\_client.hpp:256!}
\end{itemize}

\textbf{说明}

\begin{itemize}
\tightlist
\item
  示例是局部调用片段；\passthrough{\lstinline!obj!}
  和参数变量表示已经按上层业务校验并准备好的对象。
\item
  该条目可以被编辑器和类型检查器识别，但当前运行时可能在属性查找阶段直接失败。
\item
  该方法属于运动命令。方法名包含
  \passthrough{\lstinline!stop!}、\passthrough{\lstinline!damp!} 或
  \passthrough{\lstinline!zero!} 也不代表它是独立物理急停。
\end{itemize}

\textbf{用法}

\begin{lstlisting}[language=Python]
from typing import TYPE_CHECKING

if TYPE_CHECKING:
    def planned_reset_estimator(obj: SportClient) -> int:
        return obj.reset_estimator()
\end{lstlisting}

\begin{quote}
{[}!CAUTION{]} 上面的 \passthrough{\lstinline!TYPE\_CHECKING!}
示例只表达计划调用形态。该分支在普通 Python
运行时为假，不代表当前扩展能执行此接口。
\end{quote}

\hypertarget{unitree-sdk2-cpp-robot-a2-sportclient-trajectory-1}{%
\paragraph{\texorpdfstring{\texttt{unitree\_sdk2\_cpp.robot.a2.SportClient.trajectory}}{unitree\_sdk2\_cpp.robot.a2.SportClient.trajectory}}\label{unitree-sdk2-cpp-robot-a2-sportclient-trajectory-1}}

对应 C++ SDK 操作
\passthrough{\lstinline!Trajectory(const std::vector<PathPoint> \&, int, float, float, float)!}。上游头文件没有可直接生成的业务说明时，本参考只保证签名映射，精确语义需查目标型号协议。

\textbf{签名}

\begin{lstlisting}[language=Python]
def trajectory(self, path: Sequence[PathPoint], feedback_mode: int = ..., external_x: float = ..., external_y: float = ..., external_yaw: float = ...) -> int
\end{lstlisting}

\textbf{可用性与安全性}

\textbf{\passthrough{\lstinline!SIGNATURE\_ONLY!} /
\passthrough{\lstinline!MOTION\_COMMAND!}}：仅用于补全和静态检查。当前二进制没有该
Python 方法；不得按可执行运动接口使用。

\textbf{参数}

\begin{longtable}[]{@{}
  >{\raggedright\arraybackslash}p{(\columnwidth - 6\tabcolsep) * \real{0.18}}
  >{\raggedright\arraybackslash}p{(\columnwidth - 6\tabcolsep) * \real{0.24}}
  >{\raggedright\arraybackslash}p{(\columnwidth - 6\tabcolsep) * \real{0.18}}
  >{\raggedright\arraybackslash}p{(\columnwidth - 6\tabcolsep) * \real{0.40}}@{}}
\toprule
名称 & Python 类型 & 默认值 & 含义 \\
\midrule
\endhead
\passthrough{\lstinline!path!} &
\passthrough{\lstinline!Sequence[PathPoint]!} & 必填 & 传给该接口的
\passthrough{\lstinline!path!} 参数，Python 类型为
\passthrough{\lstinline!Sequence[PathPoint]!}。精确含义、单位、范围和枚举值需查目标型号对应头文件与协议。
对应 C++ 参数
\passthrough{\lstinline!path: const std::vector<PathPoint> \&!}。
底层是可变长度 vector \\
\passthrough{\lstinline!feedback\_mode!} & \passthrough{\lstinline!int!}
& \passthrough{\lstinline!...!} & 传给该接口的
\passthrough{\lstinline!feedback!} 模式 参数，Python 类型为
\passthrough{\lstinline!int!}。精确含义、单位、范围和枚举值需查目标型号对应头文件与协议。
对应 C++ 参数 \passthrough{\lstinline!feedback\_mode: int!}。 \\
\passthrough{\lstinline!external\_x!} & \passthrough{\lstinline!float!}
& \passthrough{\lstinline!...!} & 传给该接口的
\passthrough{\lstinline!external!} \passthrough{\lstinline!x!}
参数，Python 类型为
\passthrough{\lstinline!float!}。精确含义、单位、范围和枚举值需查目标型号对应头文件与协议。
对应 C++ 参数 \passthrough{\lstinline!external\_x: float!}。
底层为浮点数；类型本身不说明单位、坐标系或业务安全范围。 \\
\passthrough{\lstinline!external\_y!} & \passthrough{\lstinline!float!}
& \passthrough{\lstinline!...!} & 传给该接口的
\passthrough{\lstinline!external!} \passthrough{\lstinline!y!}
参数，Python 类型为
\passthrough{\lstinline!float!}。精确含义、单位、范围和枚举值需查目标型号对应头文件与协议。
对应 C++ 参数 \passthrough{\lstinline!external\_y: float!}。
底层为浮点数；类型本身不说明单位、坐标系或业务安全范围。 \\
\passthrough{\lstinline!external\_yaw!} &
\passthrough{\lstinline!float!} & \passthrough{\lstinline!...!} &
传给该接口的 \passthrough{\lstinline!external!}
\passthrough{\lstinline!yaw!} 参数，Python 类型为
\passthrough{\lstinline!float!}。精确含义、单位、范围和枚举值需查目标型号对应头文件与协议。
对应 C++ 参数 \passthrough{\lstinline!external\_yaw: float!}。
底层为浮点数；类型本身不说明单位、坐标系或业务安全范围。 \\
\bottomrule
\end{longtable}

\textbf{返回值}

\begin{longtable}[]{@{}
  >{\raggedright\arraybackslash}p{(\columnwidth - 4\tabcolsep) * \real{0.18}}
  >{\raggedright\arraybackslash}p{(\columnwidth - 4\tabcolsep) * \real{0.20}}
  >{\raggedright\arraybackslash}p{(\columnwidth - 4\tabcolsep) * \real{0.62}}@{}}
\toprule
位置 & 类型 & 含义 \\
\midrule
\endhead
返回值 & \passthrough{\lstinline!int!} & SDK 状态码；Unitree 示例通常以
\passthrough{\lstinline!0!} 表示成功，非零值需按具体服务定义解释。 \\
\bottomrule
\end{longtable}

\textbf{对应 C++}

\begin{itemize}
\tightlist
\item
  类：\passthrough{\lstinline!unitree::robot::a2::SportClient!}
\item
  签名：\passthrough{\lstinline!Trajectory(const std::vector<PathPoint> \&, int, float, float, float)!}
\item
  绑定策略：\passthrough{\lstinline!DIRECT!}
\item
  声明位置：\passthrough{\lstinline!include/unitree/robot/a2/sport/sport\_client.hpp:262!}
\end{itemize}

\textbf{说明}

\begin{itemize}
\tightlist
\item
  示例是局部调用片段；\passthrough{\lstinline!obj!}
  和参数变量表示已经按上层业务校验并准备好的对象。
\item
  默认值显示为 \passthrough{\lstinline!...!}，表示 C++
  声明存在默认参数，但当前 AST
  清单没有保存其字面量；省略参数可使用上游默认行为。
\item
  该条目可以被编辑器和类型检查器识别，但当前运行时可能在属性查找阶段直接失败。
\item
  该方法属于运动命令。方法名包含
  \passthrough{\lstinline!stop!}、\passthrough{\lstinline!damp!} 或
  \passthrough{\lstinline!zero!} 也不代表它是独立物理急停。
\end{itemize}

\textbf{用法}

\begin{lstlisting}[language=Python]
from typing import TYPE_CHECKING

if TYPE_CHECKING:
    def planned_trajectory(obj: SportClient, path: Sequence[PathPoint], feedback_mode: int, external_x: float, external_y: float, external_yaw: float) -> int:
        return obj.trajectory(path=path, feedback_mode=feedback_mode, external_x=external_x, external_y=external_y, external_yaw=external_yaw)
\end{lstlisting}

\begin{quote}
{[}!CAUTION{]} 上面的 \passthrough{\lstinline!TYPE\_CHECKING!}
示例只表达计划调用形态。该分支在普通 Python
运行时为假，不代表当前扩展能执行此接口。
\end{quote}

\hypertarget{unitree-sdk2-cpp-robot-a2-sportclient-set-auto-recovery-1}{%
\paragraph{\texorpdfstring{\texttt{unitree\_sdk2\_cpp.robot.a2.SportClient.set\_auto\_recovery}}{unitree\_sdk2\_cpp.robot.a2.SportClient.set\_auto\_recovery}}\label{unitree-sdk2-cpp-robot-a2-sportclient-set-auto-recovery-1}}

设置 \passthrough{\lstinline!auto!}
\passthrough{\lstinline!recovery!}。具体副作用和安全边界见下方状态。

\textbf{签名}

\begin{lstlisting}[language=Python]
def set_auto_recovery(self, switch_on: int) -> int
\end{lstlisting}

\textbf{可用性与安全性}

\textbf{\passthrough{\lstinline!SIGNATURE\_ONLY!} /
\passthrough{\lstinline!MOTION\_COMMAND!}}：仅用于补全和静态检查。当前二进制没有该
Python 方法；不得按可执行运动接口使用。

\textbf{参数}

\begin{longtable}[]{@{}
  >{\raggedright\arraybackslash}p{(\columnwidth - 6\tabcolsep) * \real{0.18}}
  >{\raggedright\arraybackslash}p{(\columnwidth - 6\tabcolsep) * \real{0.24}}
  >{\raggedright\arraybackslash}p{(\columnwidth - 6\tabcolsep) * \real{0.18}}
  >{\raggedright\arraybackslash}p{(\columnwidth - 6\tabcolsep) * \real{0.40}}@{}}
\toprule
名称 & Python 类型 & 默认值 & 含义 \\
\midrule
\endhead
\passthrough{\lstinline!switch\_on!} & \passthrough{\lstinline!int!} &
必填 & 传给该接口的 开关 \passthrough{\lstinline!on!} 参数，Python
类型为
\passthrough{\lstinline!int!}。精确含义、单位、范围和枚举值需查目标型号对应头文件与协议。
对应 C++ 参数 \passthrough{\lstinline!switch\_on: int!}。 \\
\bottomrule
\end{longtable}

\textbf{返回值}

\begin{longtable}[]{@{}
  >{\raggedright\arraybackslash}p{(\columnwidth - 4\tabcolsep) * \real{0.18}}
  >{\raggedright\arraybackslash}p{(\columnwidth - 4\tabcolsep) * \real{0.20}}
  >{\raggedright\arraybackslash}p{(\columnwidth - 4\tabcolsep) * \real{0.62}}@{}}
\toprule
位置 & 类型 & 含义 \\
\midrule
\endhead
返回值 & \passthrough{\lstinline!int!} & SDK 状态码；Unitree 示例通常以
\passthrough{\lstinline!0!} 表示成功，非零值需按具体服务定义解释。 \\
\bottomrule
\end{longtable}

\textbf{对应 C++}

\begin{itemize}
\tightlist
\item
  类：\passthrough{\lstinline!unitree::robot::a2::SportClient!}
\item
  签名：\passthrough{\lstinline!SetAutoRecovery(int)!}
\item
  绑定策略：\passthrough{\lstinline!DIRECT!}
\item
  声明位置：\passthrough{\lstinline!include/unitree/robot/a2/sport/sport\_client.hpp:286!}
\end{itemize}

\textbf{说明}

\begin{itemize}
\tightlist
\item
  示例是局部调用片段；\passthrough{\lstinline!obj!}
  和参数变量表示已经按上层业务校验并准备好的对象。
\item
  该条目可以被编辑器和类型检查器识别，但当前运行时可能在属性查找阶段直接失败。
\item
  该方法属于运动命令。方法名包含
  \passthrough{\lstinline!stop!}、\passthrough{\lstinline!damp!} 或
  \passthrough{\lstinline!zero!} 也不代表它是独立物理急停。
\end{itemize}

\textbf{用法}

\begin{lstlisting}[language=Python]
from typing import TYPE_CHECKING

if TYPE_CHECKING:
    def planned_set_auto_recovery(obj: SportClient, switch_on: int) -> int:
        return obj.set_auto_recovery(switch_on=switch_on)
\end{lstlisting}

\begin{quote}
{[}!CAUTION{]} 上面的 \passthrough{\lstinline!TYPE\_CHECKING!}
示例只表达计划调用形态。该分支在普通 Python
运行时为假，不代表当前扩展能执行此接口。
\end{quote}

\hypertarget{unitree-sdk2-cpp-robot-a2-sportclient-get-state-1}{%
\paragraph{\texorpdfstring{\texttt{unitree\_sdk2\_cpp.robot.a2.SportClient.get\_state}}{unitree\_sdk2\_cpp.robot.a2.SportClient.get\_state}}\label{unitree-sdk2-cpp-robot-a2-sportclient-get-state-1}}

查询或检查 状态。

\textbf{签名}

\begin{lstlisting}[language=Python]
def get_state(self) -> tuple[int, dict[str, str]]
\end{lstlisting}

\textbf{可用性与安全性}

\textbf{\passthrough{\lstinline!AVAILABLE!} /
\passthrough{\lstinline!READ\_ONLY!}}：当前绑定源码已实现只读查询。Robot
Client 的服务端查询通常仍需匹配的 Linux
扩展、DDS、网络、机器人和服务；纯本地 getter 则不一定需要全部条件。

\textbf{参数}

无显式参数。实例方法中的 \passthrough{\lstinline!self!} 由 Python
自动传入。

\textbf{返回值}

\begin{longtable}[]{@{}
  >{\raggedright\arraybackslash}p{(\columnwidth - 4\tabcolsep) * \real{0.18}}
  >{\raggedright\arraybackslash}p{(\columnwidth - 4\tabcolsep) * \real{0.20}}
  >{\raggedright\arraybackslash}p{(\columnwidth - 4\tabcolsep) * \real{0.62}}@{}}
\toprule
位置 & 类型 & 含义 \\
\midrule
\endhead
{[}0{]} \passthrough{\lstinline!status!} & \passthrough{\lstinline!int!}
& SDK 状态码；Unitree 示例通常以 \passthrough{\lstinline!0!}
表示成功，非零值需按具体服务错误码解释。 \\
{[}1{]} \passthrough{\lstinline!state\_map!} &
\passthrough{\lstinline!dict[str, str]!} & 传给该接口的 状态
\passthrough{\lstinline!map!} 参数，Python 类型为
\passthrough{\lstinline!dict[str, str]!}。精确含义、单位、范围和枚举值需查目标型号对应头文件与协议。
对应 C++ 参数
\passthrough{\lstinline!state\_map: std::map<std::string, std::string> \&!}。 \\
\bottomrule
\end{longtable}

\textbf{对应 C++}

\begin{itemize}
\tightlist
\item
  类：\passthrough{\lstinline!unitree::robot::a2::SportClient!}
\item
  签名：\passthrough{\lstinline!GetState(std::map<std::string, std::string> \&)!}
\item
  绑定策略：\passthrough{\lstinline!OUTPUT\_WRAPPER!}
\item
  声明位置：\passthrough{\lstinline!include/unitree/robot/a2/sport/sport\_client.hpp:295!}
\end{itemize}

\textbf{说明}

\begin{itemize}
\tightlist
\item
  示例是局部调用片段；\passthrough{\lstinline!obj!}
  和参数变量表示已经按上层业务校验并准备好的对象。
\item
  C++ 的可修改输出引用不会作为 Python 入参出现，而是按 C++
  参数顺序追加到返回元组中。
\end{itemize}

\textbf{用法}

\begin{lstlisting}[language=Python]
status, state_map = obj.get_state()
\end{lstlisting}

\begin{quote}
{[}!NOTE{]} 只读查询仍会访问
DDS/机器人服务。必须检查返回状态码，并处理超时和网络中断。
\end{quote}

\hypertarget{unitree-sdk2-cpp-robot-a2-ttsmakerparameter}{%
\subsubsection{\texorpdfstring{\texttt{unitree\_sdk2\_cpp.robot.a2.TtsMakerParameter}}{unitree\_sdk2\_cpp.robot.a2.TtsMakerParameter}}\label{unitree-sdk2-cpp-robot-a2-ttsmakerparameter}}

SDK 请求参数值类型；当前可能仅提供设计期签名。

\textbf{导入}

\begin{lstlisting}[language=Python]
from unitree_sdk2_cpp.robot.a2 import TtsMakerParameter
\end{lstlisting}

\textbf{基类}：\passthrough{\lstinline!object!}

\textbf{公开属性}

\begin{longtable}[]{@{}
  >{\raggedright\arraybackslash}p{(\columnwidth - 6\tabcolsep) * \real{0.18}}
  >{\raggedright\arraybackslash}p{(\columnwidth - 6\tabcolsep) * \real{0.24}}
  >{\raggedright\arraybackslash}p{(\columnwidth - 6\tabcolsep) * \real{0.18}}
  >{\raggedright\arraybackslash}p{(\columnwidth - 6\tabcolsep) * \real{0.40}}@{}}
\toprule
名称 & Python 类型 & 含义 & 用法 \\
\midrule
\endhead
\passthrough{\lstinline!index!} & \passthrough{\lstinline!int!} &
传给该接口的 \passthrough{\lstinline!index!} 参数，Python 类型为
\passthrough{\lstinline!int!}。精确含义、单位、范围和枚举值需查目标型号对应头文件与协议。
& \passthrough{\lstinline!value = obj.index!} /
\passthrough{\lstinline!obj.index = value!} \\
\passthrough{\lstinline!speaker\_id!} & \passthrough{\lstinline!int!} &
语音合成说话人 ID。有效编号由机器人音频服务和固件决定。 &
\passthrough{\lstinline!value = obj.speaker\_id!} /
\passthrough{\lstinline!obj.speaker\_id = value!} \\
\passthrough{\lstinline!text!} & \passthrough{\lstinline!str!} &
要处理的文本内容；编码、长度和语言支持由目标服务决定。 &
\passthrough{\lstinline!value = obj.text!} /
\passthrough{\lstinline!obj.text = value!} \\
\bottomrule
\end{longtable}

\textbf{方法索引}

\begin{longtable}[]{@{}
  >{\raggedright\arraybackslash}p{(\columnwidth - 6\tabcolsep) * \real{0.18}}
  >{\raggedright\arraybackslash}p{(\columnwidth - 6\tabcolsep) * \real{0.24}}
  >{\raggedright\arraybackslash}p{(\columnwidth - 6\tabcolsep) * \real{0.18}}
  >{\raggedright\arraybackslash}p{(\columnwidth - 6\tabcolsep) * \real{0.40}}@{}}
\toprule
方法 & Python 签名 & 状态 & 安全分类 \\
\midrule
\endhead
\protect\hyperlink{unitree-sdk2-cpp-robot-a2-ttsmakerparameter-dunder-init-1}{\passthrough{\lstinline!\_\_init\_\_!}}
& \passthrough{\lstinline!def \_\_init\_\_(self) -> None!} &
\passthrough{\lstinline!SIGNATURE\_ONLY!} &
\passthrough{\lstinline!CONSTRUCTION!} \\
\protect\hyperlink{unitree-sdk2-cpp-robot-a2-ttsmakerparameter-from-json-1}{\passthrough{\lstinline!from\_json!}}
&
\passthrough{\lstinline!def from\_json(self, json: dict[str, Any]) -> None!}
& \passthrough{\lstinline!SIGNATURE\_ONLY!} &
\passthrough{\lstinline!UNCLASSIFIED!} \\
\protect\hyperlink{unitree-sdk2-cpp-robot-a2-ttsmakerparameter-to-json-1}{\passthrough{\lstinline!to\_json!}}
&
\passthrough{\lstinline!def to\_json(self, json: dict[str, Any]) -> None!}
& \passthrough{\lstinline!SIGNATURE\_ONLY!} &
\passthrough{\lstinline!UNCLASSIFIED!} \\
\bottomrule
\end{longtable}

\hypertarget{unitree-sdk2-cpp-robot-a2-ttsmakerparameter-dunder-init-1}{%
\paragraph{\texorpdfstring{\texttt{unitree\_sdk2\_cpp.robot.a2.TtsMakerParameter.\_\_init\_\_}}{unitree\_sdk2\_cpp.robot.a2.TtsMakerParameter.\_\_init\_\_}}\label{unitree-sdk2-cpp-robot-a2-ttsmakerparameter-dunder-init-1}}

初始化 \passthrough{\lstinline!TtsMakerParameter!}
实例。是否能实际构造取决于下方可用性状态。

\textbf{签名}

\begin{lstlisting}[language=Python]
def __init__(self) -> None
\end{lstlisting}

\textbf{可用性与安全性}

\textbf{\passthrough{\lstinline!SIGNATURE\_ONLY!} /
\passthrough{\lstinline!CONSTRUCTION!}}：当前只有设计期签名，不能假设已存在于运行时扩展。

\textbf{参数}

无显式参数。实例方法中的 \passthrough{\lstinline!self!} 由 Python
自动传入。

\textbf{返回值}

\begin{longtable}[]{@{}
  >{\raggedright\arraybackslash}p{(\columnwidth - 4\tabcolsep) * \real{0.18}}
  >{\raggedright\arraybackslash}p{(\columnwidth - 4\tabcolsep) * \real{0.20}}
  >{\raggedright\arraybackslash}p{(\columnwidth - 4\tabcolsep) * \real{0.62}}@{}}
\toprule
位置 & 类型 & 含义 \\
\midrule
\endhead
无 & \passthrough{\lstinline!None!} &
\passthrough{\lstinline!\_\_init\_\_!}
本身不返回值；调用类对象时，在构造函数可用的前提下得到该类实例。 \\
\bottomrule
\end{longtable}

\textbf{对应 C++}

\begin{itemize}
\tightlist
\item
  类：\passthrough{\lstinline!unitree::robot::a2::TtsMakerParameter!}
\item
  签名：\passthrough{\lstinline!TtsMakerParameter()!}
\item
  绑定策略：\passthrough{\lstinline!未单独标注!}
\item
  声明位置：\passthrough{\lstinline!include/unitree/robot/a2/audio/audio\_api.hpp:27!}
\end{itemize}

\textbf{说明}

\begin{itemize}
\tightlist
\item
  示例是局部调用片段；\passthrough{\lstinline!obj!}
  和参数变量表示已经按上层业务校验并准备好的对象。
\item
  该条目可以被编辑器和类型检查器识别，但当前运行时可能在属性查找阶段直接失败。
\end{itemize}

\textbf{用法}

\begin{lstlisting}[language=Python]
from typing import TYPE_CHECKING

if TYPE_CHECKING:
    def planned_construct() -> TtsMakerParameter:
        return TtsMakerParameter()
\end{lstlisting}

\begin{quote}
{[}!CAUTION{]} 上面的 \passthrough{\lstinline!TYPE\_CHECKING!}
示例只表达计划调用形态。该分支在普通 Python
运行时为假，不代表当前扩展能执行此接口。
\end{quote}

\hypertarget{unitree-sdk2-cpp-robot-a2-ttsmakerparameter-from-json-1}{%
\paragraph{\texorpdfstring{\texttt{unitree\_sdk2\_cpp.robot.a2.TtsMakerParameter.from\_json}}{unitree\_sdk2\_cpp.robot.a2.TtsMakerParameter.from\_json}}\label{unitree-sdk2-cpp-robot-a2-ttsmakerparameter-from-json-1}}

计划从 JSON 风格字典读取字段并更新当前 SDK 值对象。

\textbf{签名}

\begin{lstlisting}[language=Python]
def from_json(self, json: dict[str, Any]) -> None
\end{lstlisting}

\textbf{可用性与安全性}

\textbf{\passthrough{\lstinline!SIGNATURE\_ONLY!} /
\passthrough{\lstinline!UNCLASSIFIED!}}：当前只有设计期签名，不能假设已存在于运行时扩展。

\textbf{参数}

\begin{longtable}[]{@{}
  >{\raggedright\arraybackslash}p{(\columnwidth - 6\tabcolsep) * \real{0.18}}
  >{\raggedright\arraybackslash}p{(\columnwidth - 6\tabcolsep) * \real{0.24}}
  >{\raggedright\arraybackslash}p{(\columnwidth - 6\tabcolsep) * \real{0.18}}
  >{\raggedright\arraybackslash}p{(\columnwidth - 6\tabcolsep) * \real{0.40}}@{}}
\toprule
名称 & Python 类型 & 默认值 & 含义 \\
\midrule
\endhead
\passthrough{\lstinline!json!} &
\passthrough{\lstinline!dict[str, Any]!} & 必填 & JSON
风格字典，键和值必须满足对应 C++ \passthrough{\lstinline!JsonMap!}
数据结构的约束。 对应 C++ 参数
\passthrough{\lstinline!json: common::JsonMap \&!}。 \\
\bottomrule
\end{longtable}

\textbf{返回值}

\begin{longtable}[]{@{}
  >{\raggedright\arraybackslash}p{(\columnwidth - 4\tabcolsep) * \real{0.18}}
  >{\raggedright\arraybackslash}p{(\columnwidth - 4\tabcolsep) * \real{0.20}}
  >{\raggedright\arraybackslash}p{(\columnwidth - 4\tabcolsep) * \real{0.62}}@{}}
\toprule
位置 & 类型 & 含义 \\
\midrule
\endhead
无 & \passthrough{\lstinline!None!} & 该方法不返回 Python
值；失败可能表现为异常或底层状态变化。 \\
\bottomrule
\end{longtable}

\textbf{对应 C++}

\begin{itemize}
\tightlist
\item
  类：\passthrough{\lstinline!unitree::robot::a2::TtsMakerParameter!}
\item
  签名：\passthrough{\lstinline!fromJson(common::JsonMap \&)!}
\item
  绑定策略：\passthrough{\lstinline!SIGNATURE\_PREVIEW!}
\item
  声明位置：\passthrough{\lstinline!include/unitree/robot/a2/audio/audio\_api.hpp:30!}
\end{itemize}

\textbf{说明}

\begin{itemize}
\tightlist
\item
  示例是局部调用片段；\passthrough{\lstinline!obj!}
  和参数变量表示已经按上层业务校验并准备好的对象。
\item
  该条目可以被编辑器和类型检查器识别，但当前运行时可能在属性查找阶段直接失败。
\end{itemize}

\textbf{用法}

\begin{lstlisting}[language=Python]
from typing import TYPE_CHECKING

if TYPE_CHECKING:
    def planned_from_json(obj: TtsMakerParameter, json: dict[str, Any]) -> None:
        obj.from_json(json=json)
\end{lstlisting}

\begin{quote}
{[}!CAUTION{]} 上面的 \passthrough{\lstinline!TYPE\_CHECKING!}
示例只表达计划调用形态。该分支在普通 Python
运行时为假，不代表当前扩展能执行此接口。
\end{quote}

\hypertarget{unitree-sdk2-cpp-robot-a2-ttsmakerparameter-to-json-1}{%
\paragraph{\texorpdfstring{\texttt{unitree\_sdk2\_cpp.robot.a2.TtsMakerParameter.to\_json}}{unitree\_sdk2\_cpp.robot.a2.TtsMakerParameter.to\_json}}\label{unitree-sdk2-cpp-robot-a2-ttsmakerparameter-to-json-1}}

计划把当前 SDK 值对象写入 JSON 风格字典。

\textbf{签名}

\begin{lstlisting}[language=Python]
def to_json(self, json: dict[str, Any]) -> None
\end{lstlisting}

\textbf{可用性与安全性}

\textbf{\passthrough{\lstinline!SIGNATURE\_ONLY!} /
\passthrough{\lstinline!UNCLASSIFIED!}}：当前只有设计期签名，不能假设已存在于运行时扩展。

\textbf{参数}

\begin{longtable}[]{@{}
  >{\raggedright\arraybackslash}p{(\columnwidth - 6\tabcolsep) * \real{0.18}}
  >{\raggedright\arraybackslash}p{(\columnwidth - 6\tabcolsep) * \real{0.24}}
  >{\raggedright\arraybackslash}p{(\columnwidth - 6\tabcolsep) * \real{0.18}}
  >{\raggedright\arraybackslash}p{(\columnwidth - 6\tabcolsep) * \real{0.40}}@{}}
\toprule
名称 & Python 类型 & 默认值 & 含义 \\
\midrule
\endhead
\passthrough{\lstinline!json!} &
\passthrough{\lstinline!dict[str, Any]!} & 必填 & JSON
风格字典，键和值必须满足对应 C++ \passthrough{\lstinline!JsonMap!}
数据结构的约束。 对应 C++ 参数
\passthrough{\lstinline!json: common::JsonMap \&!}。 \\
\bottomrule
\end{longtable}

\textbf{返回值}

\begin{longtable}[]{@{}
  >{\raggedright\arraybackslash}p{(\columnwidth - 4\tabcolsep) * \real{0.18}}
  >{\raggedright\arraybackslash}p{(\columnwidth - 4\tabcolsep) * \real{0.20}}
  >{\raggedright\arraybackslash}p{(\columnwidth - 4\tabcolsep) * \real{0.62}}@{}}
\toprule
位置 & 类型 & 含义 \\
\midrule
\endhead
无 & \passthrough{\lstinline!None!} & 该方法不返回 Python
值；失败可能表现为异常或底层状态变化。 \\
\bottomrule
\end{longtable}

\textbf{对应 C++}

\begin{itemize}
\tightlist
\item
  类：\passthrough{\lstinline!unitree::robot::a2::TtsMakerParameter!}
\item
  签名：\passthrough{\lstinline!toJson(common::JsonMap \&) const!}
\item
  绑定策略：\passthrough{\lstinline!SIGNATURE\_PREVIEW!}
\item
  声明位置：\passthrough{\lstinline!include/unitree/robot/a2/audio/audio\_api.hpp:32!}
\end{itemize}

\textbf{说明}

\begin{itemize}
\tightlist
\item
  示例是局部调用片段；\passthrough{\lstinline!obj!}
  和参数变量表示已经按上层业务校验并准备好的对象。
\item
  该条目可以被编辑器和类型检查器识别，但当前运行时可能在属性查找阶段直接失败。
\end{itemize}

\textbf{用法}

\begin{lstlisting}[language=Python]
from typing import TYPE_CHECKING

if TYPE_CHECKING:
    def planned_to_json(obj: TtsMakerParameter, json: dict[str, Any]) -> None:
        obj.to_json(json=json)
\end{lstlisting}

\begin{quote}
{[}!CAUTION{]} 上面的 \passthrough{\lstinline!TYPE\_CHECKING!}
示例只表达计划调用形态。该分支在普通 Python
运行时为假，不代表当前扩展能执行此接口。
\end{quote}

\begin{center}\rule{0.5\linewidth}{0.5pt}\end{center}

\hypertarget{unitree-sdk2-cpp-robot-as2}{%
\subsection{AS2 Robot API}\label{unitree-sdk2-cpp-robot-as2}}

模块：\passthrough{\lstinline!unitree\_sdk2\_cpp.robot.as2!}

\hypertarget{ux7c7bux7d22ux5f15-8}{%
\subsubsection{类索引}\label{ux7c7bux7d22ux5f15-8}}

\begin{longtable}[]{@{}
  >{\raggedright\arraybackslash}p{(\columnwidth - 4\tabcolsep) * \real{0.18}}
  >{\raggedleft\arraybackslash}p{(\columnwidth - 4\tabcolsep) * \real{0.20}}
  >{\raggedleft\arraybackslash}p{(\columnwidth - 4\tabcolsep) * \real{0.62}}@{}}
\toprule
类 & 公开函数签名 & 属性 \\
\midrule
\endhead
\protect\hyperlink{unitree-sdk2-cpp-robot-as2-posevec4}{\passthrough{\lstinline!PoseVec4!}}
& 3 & 0 \\
\protect\hyperlink{unitree-sdk2-cpp-robot-as2-sportclient}{\passthrough{\lstinline!SportClient!}}
& 34 & 0 \\
\bottomrule
\end{longtable}

\hypertarget{unitree-sdk2-cpp-robot-as2-posevec4}{%
\subsubsection{\texorpdfstring{\texttt{unitree\_sdk2\_cpp.robot.as2.PoseVec4}}{unitree\_sdk2\_cpp.robot.as2.PoseVec4}}\label{unitree-sdk2-cpp-robot-as2-posevec4}}

SDK 类型签名预览。请逐项查看构造函数和方法的可用性。

\textbf{导入}

\begin{lstlisting}[language=Python]
from unitree_sdk2_cpp.robot.as2 import PoseVec4
\end{lstlisting}

\textbf{基类}：\passthrough{\lstinline!object!}

\textbf{公开属性}

\begin{longtable}[]{@{}
  >{\raggedright\arraybackslash}p{(\columnwidth - 6\tabcolsep) * \real{0.18}}
  >{\raggedright\arraybackslash}p{(\columnwidth - 6\tabcolsep) * \real{0.24}}
  >{\raggedright\arraybackslash}p{(\columnwidth - 6\tabcolsep) * \real{0.18}}
  >{\raggedright\arraybackslash}p{(\columnwidth - 6\tabcolsep) * \real{0.40}}@{}}
\toprule
名称 & Python 类型 & 含义 & 用法 \\
\midrule
\endhead
\passthrough{\lstinline!x!} & \passthrough{\lstinline!float!} & X
方向位置或分量参数。参考系、单位和范围必须查当前方法及目标型号协议。 &
\passthrough{\lstinline!value = obj.x!} /
\passthrough{\lstinline!obj.x = value!} \\
\passthrough{\lstinline!y!} & \passthrough{\lstinline!float!} & Y
方向位置或分量参数。参考系、单位和范围必须查当前方法及目标型号协议。 &
\passthrough{\lstinline!value = obj.y!} /
\passthrough{\lstinline!obj.y = value!} \\
\passthrough{\lstinline!z!} & \passthrough{\lstinline!float!} & Z
方向位置或分量参数。参考系、单位和范围必须查当前方法及目标型号协议。 &
\passthrough{\lstinline!value = obj.z!} /
\passthrough{\lstinline!obj.z = value!} \\
\passthrough{\lstinline!yaw!} & \passthrough{\lstinline!float!} &
偏航角或偏航目标参数。参考系、单位和范围必须查目标型号协议。 &
\passthrough{\lstinline!value = obj.yaw!} /
\passthrough{\lstinline!obj.yaw = value!} \\
\bottomrule
\end{longtable}

\textbf{方法索引}

\begin{longtable}[]{@{}
  >{\raggedright\arraybackslash}p{(\columnwidth - 6\tabcolsep) * \real{0.18}}
  >{\raggedright\arraybackslash}p{(\columnwidth - 6\tabcolsep) * \real{0.24}}
  >{\raggedright\arraybackslash}p{(\columnwidth - 6\tabcolsep) * \real{0.18}}
  >{\raggedright\arraybackslash}p{(\columnwidth - 6\tabcolsep) * \real{0.40}}@{}}
\toprule
方法 & Python 签名 & 状态 & 安全分类 \\
\midrule
\endhead
\protect\hyperlink{unitree-sdk2-cpp-robot-as2-posevec4-dunder-init-1}{\passthrough{\lstinline!\_\_init\_\_!}}
& \passthrough{\lstinline!def \_\_init\_\_(self) -> None!} &
\passthrough{\lstinline!SIGNATURE\_ONLY!} &
\passthrough{\lstinline!CONSTRUCTION!} \\
\protect\hyperlink{unitree-sdk2-cpp-robot-as2-posevec4-from-json-1}{\passthrough{\lstinline!from\_json!}}
&
\passthrough{\lstinline!def from\_json(self, json: dict[str, Any]) -> None!}
& \passthrough{\lstinline!SIGNATURE\_ONLY!} &
\passthrough{\lstinline!UNCLASSIFIED!} \\
\protect\hyperlink{unitree-sdk2-cpp-robot-as2-posevec4-to-json-1}{\passthrough{\lstinline!to\_json!}}
&
\passthrough{\lstinline!def to\_json(self, json: dict[str, Any]) -> None!}
& \passthrough{\lstinline!SIGNATURE\_ONLY!} &
\passthrough{\lstinline!UNCLASSIFIED!} \\
\bottomrule
\end{longtable}

\hypertarget{unitree-sdk2-cpp-robot-as2-posevec4-dunder-init-1}{%
\paragraph{\texorpdfstring{\texttt{unitree\_sdk2\_cpp.robot.as2.PoseVec4.\_\_init\_\_}}{unitree\_sdk2\_cpp.robot.as2.PoseVec4.\_\_init\_\_}}\label{unitree-sdk2-cpp-robot-as2-posevec4-dunder-init-1}}

初始化 \passthrough{\lstinline!PoseVec4!}
实例。是否能实际构造取决于下方可用性状态。

\textbf{签名}

\begin{lstlisting}[language=Python]
def __init__(self) -> None
\end{lstlisting}

\textbf{可用性与安全性}

\textbf{\passthrough{\lstinline!SIGNATURE\_ONLY!} /
\passthrough{\lstinline!CONSTRUCTION!}}：当前只有设计期签名，不能假设已存在于运行时扩展。

\textbf{参数}

无显式参数。实例方法中的 \passthrough{\lstinline!self!} 由 Python
自动传入。

\textbf{返回值}

\begin{longtable}[]{@{}
  >{\raggedright\arraybackslash}p{(\columnwidth - 4\tabcolsep) * \real{0.18}}
  >{\raggedright\arraybackslash}p{(\columnwidth - 4\tabcolsep) * \real{0.20}}
  >{\raggedright\arraybackslash}p{(\columnwidth - 4\tabcolsep) * \real{0.62}}@{}}
\toprule
位置 & 类型 & 含义 \\
\midrule
\endhead
无 & \passthrough{\lstinline!None!} &
\passthrough{\lstinline!\_\_init\_\_!}
本身不返回值；调用类对象时，在构造函数可用的前提下得到该类实例。 \\
\bottomrule
\end{longtable}

\textbf{对应 C++}

\begin{itemize}
\tightlist
\item
  类：\passthrough{\lstinline!unitree::robot::as2::PoseVec4!}
\item
  签名：\passthrough{\lstinline!PoseVec4()!}
\item
  绑定策略：\passthrough{\lstinline!未单独标注!}
\item
  声明位置：\passthrough{\lstinline!include/unitree/robot/as2/sport/sport\_client.hpp:18!}
\end{itemize}

\textbf{说明}

\begin{itemize}
\tightlist
\item
  示例是局部调用片段；\passthrough{\lstinline!obj!}
  和参数变量表示已经按上层业务校验并准备好的对象。
\item
  该条目可以被编辑器和类型检查器识别，但当前运行时可能在属性查找阶段直接失败。
\end{itemize}

\textbf{用法}

\begin{lstlisting}[language=Python]
from typing import TYPE_CHECKING

if TYPE_CHECKING:
    def planned_construct() -> PoseVec4:
        return PoseVec4()
\end{lstlisting}

\begin{quote}
{[}!CAUTION{]} 上面的 \passthrough{\lstinline!TYPE\_CHECKING!}
示例只表达计划调用形态。该分支在普通 Python
运行时为假，不代表当前扩展能执行此接口。
\end{quote}

\hypertarget{unitree-sdk2-cpp-robot-as2-posevec4-from-json-1}{%
\paragraph{\texorpdfstring{\texttt{unitree\_sdk2\_cpp.robot.as2.PoseVec4.from\_json}}{unitree\_sdk2\_cpp.robot.as2.PoseVec4.from\_json}}\label{unitree-sdk2-cpp-robot-as2-posevec4-from-json-1}}

计划从 JSON 风格字典读取字段并更新当前 SDK 值对象。

\textbf{签名}

\begin{lstlisting}[language=Python]
def from_json(self, json: dict[str, Any]) -> None
\end{lstlisting}

\textbf{可用性与安全性}

\textbf{\passthrough{\lstinline!SIGNATURE\_ONLY!} /
\passthrough{\lstinline!UNCLASSIFIED!}}：当前只有设计期签名，不能假设已存在于运行时扩展。

\textbf{参数}

\begin{longtable}[]{@{}
  >{\raggedright\arraybackslash}p{(\columnwidth - 6\tabcolsep) * \real{0.18}}
  >{\raggedright\arraybackslash}p{(\columnwidth - 6\tabcolsep) * \real{0.24}}
  >{\raggedright\arraybackslash}p{(\columnwidth - 6\tabcolsep) * \real{0.18}}
  >{\raggedright\arraybackslash}p{(\columnwidth - 6\tabcolsep) * \real{0.40}}@{}}
\toprule
名称 & Python 类型 & 默认值 & 含义 \\
\midrule
\endhead
\passthrough{\lstinline!json!} &
\passthrough{\lstinline!dict[str, Any]!} & 必填 & JSON
风格字典，键和值必须满足对应 C++ \passthrough{\lstinline!JsonMap!}
数据结构的约束。 对应 C++ 参数
\passthrough{\lstinline!json: common::JsonMap \&!}。 \\
\bottomrule
\end{longtable}

\textbf{返回值}

\begin{longtable}[]{@{}
  >{\raggedright\arraybackslash}p{(\columnwidth - 4\tabcolsep) * \real{0.18}}
  >{\raggedright\arraybackslash}p{(\columnwidth - 4\tabcolsep) * \real{0.20}}
  >{\raggedright\arraybackslash}p{(\columnwidth - 4\tabcolsep) * \real{0.62}}@{}}
\toprule
位置 & 类型 & 含义 \\
\midrule
\endhead
无 & \passthrough{\lstinline!None!} & 该方法不返回 Python
值；失败可能表现为异常或底层状态变化。 \\
\bottomrule
\end{longtable}

\textbf{对应 C++}

\begin{itemize}
\tightlist
\item
  类：\passthrough{\lstinline!unitree::robot::as2::PoseVec4!}
\item
  签名：\passthrough{\lstinline!fromJson(common::JsonMap \&)!}
\item
  绑定策略：\passthrough{\lstinline!SIGNATURE\_PREVIEW!}
\item
  声明位置：\passthrough{\lstinline!include/unitree/robot/as2/sport/sport\_client.hpp:26!}
\end{itemize}

\textbf{说明}

\begin{itemize}
\tightlist
\item
  示例是局部调用片段；\passthrough{\lstinline!obj!}
  和参数变量表示已经按上层业务校验并准备好的对象。
\item
  该条目可以被编辑器和类型检查器识别，但当前运行时可能在属性查找阶段直接失败。
\end{itemize}

\textbf{用法}

\begin{lstlisting}[language=Python]
from typing import TYPE_CHECKING

if TYPE_CHECKING:
    def planned_from_json(obj: PoseVec4, json: dict[str, Any]) -> None:
        obj.from_json(json=json)
\end{lstlisting}

\begin{quote}
{[}!CAUTION{]} 上面的 \passthrough{\lstinline!TYPE\_CHECKING!}
示例只表达计划调用形态。该分支在普通 Python
运行时为假，不代表当前扩展能执行此接口。
\end{quote}

\hypertarget{unitree-sdk2-cpp-robot-as2-posevec4-to-json-1}{%
\paragraph{\texorpdfstring{\texttt{unitree\_sdk2\_cpp.robot.as2.PoseVec4.to\_json}}{unitree\_sdk2\_cpp.robot.as2.PoseVec4.to\_json}}\label{unitree-sdk2-cpp-robot-as2-posevec4-to-json-1}}

计划把当前 SDK 值对象写入 JSON 风格字典。

\textbf{签名}

\begin{lstlisting}[language=Python]
def to_json(self, json: dict[str, Any]) -> None
\end{lstlisting}

\textbf{可用性与安全性}

\textbf{\passthrough{\lstinline!SIGNATURE\_ONLY!} /
\passthrough{\lstinline!UNCLASSIFIED!}}：当前只有设计期签名，不能假设已存在于运行时扩展。

\textbf{参数}

\begin{longtable}[]{@{}
  >{\raggedright\arraybackslash}p{(\columnwidth - 6\tabcolsep) * \real{0.18}}
  >{\raggedright\arraybackslash}p{(\columnwidth - 6\tabcolsep) * \real{0.24}}
  >{\raggedright\arraybackslash}p{(\columnwidth - 6\tabcolsep) * \real{0.18}}
  >{\raggedright\arraybackslash}p{(\columnwidth - 6\tabcolsep) * \real{0.40}}@{}}
\toprule
名称 & Python 类型 & 默认值 & 含义 \\
\midrule
\endhead
\passthrough{\lstinline!json!} &
\passthrough{\lstinline!dict[str, Any]!} & 必填 & JSON
风格字典，键和值必须满足对应 C++ \passthrough{\lstinline!JsonMap!}
数据结构的约束。 对应 C++ 参数
\passthrough{\lstinline!json: common::JsonMap \&!}。 \\
\bottomrule
\end{longtable}

\textbf{返回值}

\begin{longtable}[]{@{}
  >{\raggedright\arraybackslash}p{(\columnwidth - 4\tabcolsep) * \real{0.18}}
  >{\raggedright\arraybackslash}p{(\columnwidth - 4\tabcolsep) * \real{0.20}}
  >{\raggedright\arraybackslash}p{(\columnwidth - 4\tabcolsep) * \real{0.62}}@{}}
\toprule
位置 & 类型 & 含义 \\
\midrule
\endhead
无 & \passthrough{\lstinline!None!} & 该方法不返回 Python
值；失败可能表现为异常或底层状态变化。 \\
\bottomrule
\end{longtable}

\textbf{对应 C++}

\begin{itemize}
\tightlist
\item
  类：\passthrough{\lstinline!unitree::robot::as2::PoseVec4!}
\item
  签名：\passthrough{\lstinline!toJson(common::JsonMap \&) const!}
\item
  绑定策略：\passthrough{\lstinline!SIGNATURE\_PREVIEW!}
\item
  声明位置：\passthrough{\lstinline!include/unitree/robot/as2/sport/sport\_client.hpp:34!}
\end{itemize}

\textbf{说明}

\begin{itemize}
\tightlist
\item
  示例是局部调用片段；\passthrough{\lstinline!obj!}
  和参数变量表示已经按上层业务校验并准备好的对象。
\item
  该条目可以被编辑器和类型检查器识别，但当前运行时可能在属性查找阶段直接失败。
\end{itemize}

\textbf{用法}

\begin{lstlisting}[language=Python]
from typing import TYPE_CHECKING

if TYPE_CHECKING:
    def planned_to_json(obj: PoseVec4, json: dict[str, Any]) -> None:
        obj.to_json(json=json)
\end{lstlisting}

\begin{quote}
{[}!CAUTION{]} 上面的 \passthrough{\lstinline!TYPE\_CHECKING!}
示例只表达计划调用形态。该分支在普通 Python
运行时为假，不代表当前扩展能执行此接口。
\end{quote}

\hypertarget{unitree-sdk2-cpp-robot-as2-sportclient}{%
\subsubsection{\texorpdfstring{\texttt{unitree\_sdk2\_cpp.robot.as2.SportClient}}{unitree\_sdk2\_cpp.robot.as2.SportClient}}\label{unitree-sdk2-cpp-robot-as2-sportclient}}

Robot 服务客户端。构造、初始化和具体方法的可用性必须分别检查。

\textbf{导入}

\begin{lstlisting}[language=Python]
from unitree_sdk2_cpp.robot.as2 import SportClient
\end{lstlisting}

\textbf{基类}：\passthrough{\lstinline!Client!}

\textbf{方法索引}

\begin{longtable}[]{@{}
  >{\raggedright\arraybackslash}p{(\columnwidth - 6\tabcolsep) * \real{0.18}}
  >{\raggedright\arraybackslash}p{(\columnwidth - 6\tabcolsep) * \real{0.24}}
  >{\raggedright\arraybackslash}p{(\columnwidth - 6\tabcolsep) * \real{0.18}}
  >{\raggedright\arraybackslash}p{(\columnwidth - 6\tabcolsep) * \real{0.40}}@{}}
\toprule
方法 & Python 签名 & 状态 & 安全分类 \\
\midrule
\endhead
\protect\hyperlink{unitree-sdk2-cpp-robot-as2-sportclient-dunder-init-1}{\passthrough{\lstinline!\_\_init\_\_!}}
& \passthrough{\lstinline!def \_\_init\_\_(self) -> None!} &
\passthrough{\lstinline!AVAILABLE!} &
\passthrough{\lstinline!CONSTRUCTION!} \\
\protect\hyperlink{unitree-sdk2-cpp-robot-as2-sportclient-init-1}{\passthrough{\lstinline!init!}}
& \passthrough{\lstinline!def init(self) -> None!} &
\passthrough{\lstinline!AVAILABLE!} &
\passthrough{\lstinline!INITIALIZATION!} \\
\protect\hyperlink{unitree-sdk2-cpp-robot-as2-sportclient-damp-1}{\passthrough{\lstinline!damp!}}
& \passthrough{\lstinline!def damp(self) -> int!} &
\passthrough{\lstinline!SIGNATURE\_ONLY!} &
\passthrough{\lstinline!MOTION\_COMMAND!} \\
\protect\hyperlink{unitree-sdk2-cpp-robot-as2-sportclient-balance-stand-1}{\passthrough{\lstinline!balance\_stand!}}
& \passthrough{\lstinline!def balance\_stand(self) -> int!} &
\passthrough{\lstinline!SIGNATURE\_ONLY!} &
\passthrough{\lstinline!MOTION\_COMMAND!} \\
\protect\hyperlink{unitree-sdk2-cpp-robot-as2-sportclient-stop-move-1}{\passthrough{\lstinline!stop\_move!}}
& \passthrough{\lstinline!def stop\_move(self) -> int!} &
\passthrough{\lstinline!SIGNATURE\_ONLY!} &
\passthrough{\lstinline!MOTION\_COMMAND!} \\
\protect\hyperlink{unitree-sdk2-cpp-robot-as2-sportclient-stand-up-1}{\passthrough{\lstinline!stand\_up!}}
& \passthrough{\lstinline!def stand\_up(self) -> int!} &
\passthrough{\lstinline!SIGNATURE\_ONLY!} &
\passthrough{\lstinline!MOTION\_COMMAND!} \\
\protect\hyperlink{unitree-sdk2-cpp-robot-as2-sportclient-stand-down-1}{\passthrough{\lstinline!stand\_down!}}
& \passthrough{\lstinline!def stand\_down(self) -> int!} &
\passthrough{\lstinline!SIGNATURE\_ONLY!} &
\passthrough{\lstinline!MOTION\_COMMAND!} \\
\protect\hyperlink{unitree-sdk2-cpp-robot-as2-sportclient-recovery-stand-1}{\passthrough{\lstinline!recovery\_stand!}}
& \passthrough{\lstinline!def recovery\_stand(self) -> int!} &
\passthrough{\lstinline!SIGNATURE\_ONLY!} &
\passthrough{\lstinline!MOTION\_COMMAND!} \\
\protect\hyperlink{unitree-sdk2-cpp-robot-as2-sportclient-euler-1}{\passthrough{\lstinline!euler!}}
&
\passthrough{\lstinline!def euler(self, roll: float, pitch: float, yaw: float) -> int!}
& \passthrough{\lstinline!SIGNATURE\_ONLY!} &
\passthrough{\lstinline!MOTION\_COMMAND!} \\
\protect\hyperlink{unitree-sdk2-cpp-robot-as2-sportclient-move-1}{\passthrough{\lstinline!move!}}
&
\passthrough{\lstinline!def move(self, vx: float, vy: float, vyaw: float) -> int!}
& \passthrough{\lstinline!SIGNATURE\_ONLY!} &
\passthrough{\lstinline!MOTION\_COMMAND!} \\
\protect\hyperlink{unitree-sdk2-cpp-robot-as2-sportclient-switch-gait-1}{\passthrough{\lstinline!switch\_gait!}}
&
\passthrough{\lstinline!def switch\_gait(self, gait\_type: int) -> int!}
& \passthrough{\lstinline!SIGNATURE\_ONLY!} &
\passthrough{\lstinline!MOTION\_COMMAND!} \\
\protect\hyperlink{unitree-sdk2-cpp-robot-as2-sportclient-body-height-1}{\passthrough{\lstinline!body\_height!}}
& \passthrough{\lstinline!def body\_height(self, height: float) -> int!}
& \passthrough{\lstinline!SIGNATURE\_ONLY!} &
\passthrough{\lstinline!MOTION\_COMMAND!} \\
\protect\hyperlink{unitree-sdk2-cpp-robot-as2-sportclient-speed-level-1}{\passthrough{\lstinline!speed\_level!}}
& \passthrough{\lstinline!def speed\_level(self, level: int) -> int!} &
\passthrough{\lstinline!SIGNATURE\_ONLY!} &
\passthrough{\lstinline!MOTION\_COMMAND!} \\
\protect\hyperlink{unitree-sdk2-cpp-robot-as2-sportclient-body-position-1}{\passthrough{\lstinline!body\_position!}}
&
\passthrough{\lstinline!def body\_position(self, x: float, y: float, z: float, yaw: float) -> int!}
& \passthrough{\lstinline!SIGNATURE\_ONLY!} &
\passthrough{\lstinline!MOTION\_COMMAND!} \\
\protect\hyperlink{unitree-sdk2-cpp-robot-as2-sportclient-left-side-gait-1}{\passthrough{\lstinline!left\_side\_gait!}}
&
\passthrough{\lstinline!def left\_side\_gait(self, enter: int) -> int!}
& \passthrough{\lstinline!SIGNATURE\_ONLY!} &
\passthrough{\lstinline!MOTION\_COMMAND!} \\
\protect\hyperlink{unitree-sdk2-cpp-robot-as2-sportclient-right-side-gait-1}{\passthrough{\lstinline!right\_side\_gait!}}
&
\passthrough{\lstinline!def right\_side\_gait(self, enter: int) -> int!}
& \passthrough{\lstinline!SIGNATURE\_ONLY!} &
\passthrough{\lstinline!MOTION\_COMMAND!} \\
\protect\hyperlink{unitree-sdk2-cpp-robot-as2-sportclient-hand-stand-1}{\passthrough{\lstinline!hand\_stand!}}
& \passthrough{\lstinline!def hand\_stand(self, enter: int) -> int!} &
\passthrough{\lstinline!SIGNATURE\_ONLY!} &
\passthrough{\lstinline!MOTION\_COMMAND!} \\
\protect\hyperlink{unitree-sdk2-cpp-robot-as2-sportclient-biped-stand-1}{\passthrough{\lstinline!biped\_stand!}}
& \passthrough{\lstinline!def biped\_stand(self, enter: int) -> int!} &
\passthrough{\lstinline!SIGNATURE\_ONLY!} &
\passthrough{\lstinline!MOTION\_COMMAND!} \\
\protect\hyperlink{unitree-sdk2-cpp-robot-as2-sportclient-front-flip-1}{\passthrough{\lstinline!front\_flip!}}
& \passthrough{\lstinline!def front\_flip(self) -> int!} &
\passthrough{\lstinline!SIGNATURE\_ONLY!} &
\passthrough{\lstinline!MOTION\_COMMAND!} \\
\protect\hyperlink{unitree-sdk2-cpp-robot-as2-sportclient-back-flip-1}{\passthrough{\lstinline!back\_flip!}}
& \passthrough{\lstinline!def back\_flip(self) -> int!} &
\passthrough{\lstinline!SIGNATURE\_ONLY!} &
\passthrough{\lstinline!MOTION\_COMMAND!} \\
\protect\hyperlink{unitree-sdk2-cpp-robot-as2-sportclient-greeting-1}{\passthrough{\lstinline!greeting!}}
& \passthrough{\lstinline!def greeting(self) -> int!} &
\passthrough{\lstinline!SIGNATURE\_ONLY!} &
\passthrough{\lstinline!MOTION\_COMMAND!} \\
\protect\hyperlink{unitree-sdk2-cpp-robot-as2-sportclient-heart-1}{\passthrough{\lstinline!heart!}}
& \passthrough{\lstinline!def heart(self) -> int!} &
\passthrough{\lstinline!SIGNATURE\_ONLY!} &
\passthrough{\lstinline!MOTION\_COMMAND!} \\
\protect\hyperlink{unitree-sdk2-cpp-robot-as2-sportclient-content-1}{\passthrough{\lstinline!content!}}
& \passthrough{\lstinline!def content(self) -> int!} &
\passthrough{\lstinline!SIGNATURE\_ONLY!} &
\passthrough{\lstinline!MOTION\_COMMAND!} \\
\protect\hyperlink{unitree-sdk2-cpp-robot-as2-sportclient-dance1-1}{\passthrough{\lstinline!dance1!}}
& \passthrough{\lstinline!def dance1(self) -> int!} &
\passthrough{\lstinline!SIGNATURE\_ONLY!} &
\passthrough{\lstinline!MOTION\_COMMAND!} \\
\protect\hyperlink{unitree-sdk2-cpp-robot-as2-sportclient-dance2-1}{\passthrough{\lstinline!dance2!}}
& \passthrough{\lstinline!def dance2(self) -> int!} &
\passthrough{\lstinline!SIGNATURE\_ONLY!} &
\passthrough{\lstinline!MOTION\_COMMAND!} \\
\protect\hyperlink{unitree-sdk2-cpp-robot-as2-sportclient-handshake-1}{\passthrough{\lstinline!handshake!}}
& \passthrough{\lstinline!def handshake(self) -> int!} &
\passthrough{\lstinline!SIGNATURE\_ONLY!} &
\passthrough{\lstinline!MOTION\_COMMAND!} \\
\protect\hyperlink{unitree-sdk2-cpp-robot-as2-sportclient-stretch-1}{\passthrough{\lstinline!stretch!}}
& \passthrough{\lstinline!def stretch(self) -> int!} &
\passthrough{\lstinline!SIGNATURE\_ONLY!} &
\passthrough{\lstinline!MOTION\_COMMAND!} \\
\protect\hyperlink{unitree-sdk2-cpp-robot-as2-sportclient-sit-1}{\passthrough{\lstinline!sit!}}
& \passthrough{\lstinline!def sit(self, enter: int) -> int!} &
\passthrough{\lstinline!SIGNATURE\_ONLY!} &
\passthrough{\lstinline!MOTION\_COMMAND!} \\
\protect\hyperlink{unitree-sdk2-cpp-robot-as2-sportclient-front-jump-1}{\passthrough{\lstinline!front\_jump!}}
& \passthrough{\lstinline!def front\_jump(self) -> int!} &
\passthrough{\lstinline!SIGNATURE\_ONLY!} &
\passthrough{\lstinline!MOTION\_COMMAND!} \\
\protect\hyperlink{unitree-sdk2-cpp-robot-as2-sportclient-push-up-1}{\passthrough{\lstinline!push\_up!}}
& \passthrough{\lstinline!def push\_up(self) -> int!} &
\passthrough{\lstinline!SIGNATURE\_ONLY!} &
\passthrough{\lstinline!MOTION\_COMMAND!} \\
\protect\hyperlink{unitree-sdk2-cpp-robot-as2-sportclient-up-jump-1}{\passthrough{\lstinline!up\_jump!}}
& \passthrough{\lstinline!def up\_jump(self) -> int!} &
\passthrough{\lstinline!SIGNATURE\_ONLY!} &
\passthrough{\lstinline!MOTION\_COMMAND!} \\
\protect\hyperlink{unitree-sdk2-cpp-robot-as2-sportclient-set-auto-recovery-1}{\passthrough{\lstinline!set\_auto\_recovery!}}
&
\passthrough{\lstinline!def set\_auto\_recovery(self, switch\_on: int) -> int!}
& \passthrough{\lstinline!SIGNATURE\_ONLY!} &
\passthrough{\lstinline!MOTION\_COMMAND!} \\
\protect\hyperlink{unitree-sdk2-cpp-robot-as2-sportclient-switch-joystick-1}{\passthrough{\lstinline!switch\_joystick!}}
&
\passthrough{\lstinline!def switch\_joystick(self, switch\_on: int) -> int!}
& \passthrough{\lstinline!SIGNATURE\_ONLY!} &
\passthrough{\lstinline!MOTION\_COMMAND!} \\
\protect\hyperlink{unitree-sdk2-cpp-robot-as2-sportclient-get-state-1}{\passthrough{\lstinline!get\_state!}}
&
\passthrough{\lstinline!def get\_state(self) -> tuple[int, dict[str, str]]!}
& \passthrough{\lstinline!AVAILABLE!} &
\passthrough{\lstinline!READ\_ONLY!} \\
\bottomrule
\end{longtable}

\hypertarget{unitree-sdk2-cpp-robot-as2-sportclient-dunder-init-1}{%
\paragraph{\texorpdfstring{\texttt{unitree\_sdk2\_cpp.robot.as2.SportClient.\_\_init\_\_}}{unitree\_sdk2\_cpp.robot.as2.SportClient.\_\_init\_\_}}\label{unitree-sdk2-cpp-robot-as2-sportclient-dunder-init-1}}

初始化 \passthrough{\lstinline!SportClient!}
实例。是否能实际构造取决于下方可用性状态。

\textbf{签名}

\begin{lstlisting}[language=Python]
def __init__(self) -> None
\end{lstlisting}

\textbf{可用性与安全性}

\textbf{\passthrough{\lstinline!AVAILABLE!} /
\passthrough{\lstinline!CONSTRUCTION!}}：当前绑定源码已实现；实际执行条件仍取决于平台和对应
SDK 环境。

\textbf{参数}

无显式参数。实例方法中的 \passthrough{\lstinline!self!} 由 Python
自动传入。

\textbf{返回值}

\begin{longtable}[]{@{}
  >{\raggedright\arraybackslash}p{(\columnwidth - 4\tabcolsep) * \real{0.18}}
  >{\raggedright\arraybackslash}p{(\columnwidth - 4\tabcolsep) * \real{0.20}}
  >{\raggedright\arraybackslash}p{(\columnwidth - 4\tabcolsep) * \real{0.62}}@{}}
\toprule
位置 & 类型 & 含义 \\
\midrule
\endhead
无 & \passthrough{\lstinline!None!} &
\passthrough{\lstinline!\_\_init\_\_!}
本身不返回值；调用类对象时，在构造函数可用的前提下得到该类实例。 \\
\bottomrule
\end{longtable}

\textbf{对应 C++}

\begin{itemize}
\tightlist
\item
  类：\passthrough{\lstinline!unitree::robot::as2::SportClient!}
\item
  签名：\passthrough{\lstinline!SportClient()!}
\item
  绑定策略：\passthrough{\lstinline!未单独标注!}
\item
  声明位置：\passthrough{\lstinline!include/unitree/robot/as2/sport/sport\_client.hpp:52!}
\end{itemize}

\textbf{说明}

\begin{itemize}
\tightlist
\item
  示例是局部调用片段；\passthrough{\lstinline!obj!}
  和参数变量表示已经按上层业务校验并准备好的对象。
\end{itemize}

\textbf{用法}

\begin{lstlisting}[language=Python]
value = SportClient()
\end{lstlisting}

\hypertarget{unitree-sdk2-cpp-robot-as2-sportclient-init-1}{%
\paragraph{\texorpdfstring{\texttt{unitree\_sdk2\_cpp.robot.as2.SportClient.init}}{unitree\_sdk2\_cpp.robot.as2.SportClient.init}}\label{unitree-sdk2-cpp-robot-as2-sportclient-init-1}}

初始化当前 SDK 对象所需的底层通道或服务资源。

\textbf{签名}

\begin{lstlisting}[language=Python]
def init(self) -> None
\end{lstlisting}

\textbf{可用性与安全性}

\textbf{\passthrough{\lstinline!AVAILABLE!} /
\passthrough{\lstinline!INITIALIZATION!}}：当前绑定源码已实现；实际执行条件仍取决于平台和对应
SDK 环境。

\textbf{参数}

无显式参数。实例方法中的 \passthrough{\lstinline!self!} 由 Python
自动传入。

\textbf{返回值}

\begin{longtable}[]{@{}
  >{\raggedright\arraybackslash}p{(\columnwidth - 4\tabcolsep) * \real{0.18}}
  >{\raggedright\arraybackslash}p{(\columnwidth - 4\tabcolsep) * \real{0.20}}
  >{\raggedright\arraybackslash}p{(\columnwidth - 4\tabcolsep) * \real{0.62}}@{}}
\toprule
位置 & 类型 & 含义 \\
\midrule
\endhead
无 & \passthrough{\lstinline!None!} & 该方法不返回 Python
值；失败可能表现为异常或底层状态变化。 \\
\bottomrule
\end{longtable}

\textbf{对应 C++}

\begin{itemize}
\tightlist
\item
  类：\passthrough{\lstinline!unitree::robot::as2::SportClient!}
\item
  签名：\passthrough{\lstinline!Init()!}
\item
  绑定策略：\passthrough{\lstinline!DIRECT!}
\item
  声明位置：\passthrough{\lstinline!include/unitree/robot/as2/sport/sport\_client.hpp:56!}
\end{itemize}

\textbf{说明}

\begin{itemize}
\tightlist
\item
  示例是局部调用片段；\passthrough{\lstinline!obj!}
  和参数变量表示已经按上层业务校验并准备好的对象。
\end{itemize}

\textbf{用法}

\begin{lstlisting}[language=Python]
obj.init()
\end{lstlisting}

\hypertarget{unitree-sdk2-cpp-robot-as2-sportclient-damp-1}{%
\paragraph{\texorpdfstring{\texttt{unitree\_sdk2\_cpp.robot.as2.SportClient.damp}}{unitree\_sdk2\_cpp.robot.as2.SportClient.damp}}\label{unitree-sdk2-cpp-robot-as2-sportclient-damp-1}}

对应 C++ SDK 操作
\passthrough{\lstinline!Damp()!}。上游头文件没有可直接生成的业务说明时，本参考只保证签名映射，精确语义需查目标型号协议。

\textbf{签名}

\begin{lstlisting}[language=Python]
def damp(self) -> int
\end{lstlisting}

\textbf{可用性与安全性}

\textbf{\passthrough{\lstinline!SIGNATURE\_ONLY!} /
\passthrough{\lstinline!MOTION\_COMMAND!}}：仅用于补全和静态检查。当前二进制没有该
Python 方法；不得按可执行运动接口使用。

\textbf{参数}

无显式参数。实例方法中的 \passthrough{\lstinline!self!} 由 Python
自动传入。

\textbf{返回值}

\begin{longtable}[]{@{}
  >{\raggedright\arraybackslash}p{(\columnwidth - 4\tabcolsep) * \real{0.18}}
  >{\raggedright\arraybackslash}p{(\columnwidth - 4\tabcolsep) * \real{0.20}}
  >{\raggedright\arraybackslash}p{(\columnwidth - 4\tabcolsep) * \real{0.62}}@{}}
\toprule
位置 & 类型 & 含义 \\
\midrule
\endhead
返回值 & \passthrough{\lstinline!int!} & SDK 状态码；Unitree 示例通常以
\passthrough{\lstinline!0!} 表示成功，非零值需按具体服务定义解释。 \\
\bottomrule
\end{longtable}

\textbf{对应 C++}

\begin{itemize}
\tightlist
\item
  类：\passthrough{\lstinline!unitree::robot::as2::SportClient!}
\item
  签名：\passthrough{\lstinline!Damp()!}
\item
  绑定策略：\passthrough{\lstinline!DIRECT!}
\item
  声明位置：\passthrough{\lstinline!include/unitree/robot/as2/sport/sport\_client.hpp:97!}
\end{itemize}

\textbf{说明}

\begin{itemize}
\tightlist
\item
  示例是局部调用片段；\passthrough{\lstinline!obj!}
  和参数变量表示已经按上层业务校验并准备好的对象。
\item
  该条目可以被编辑器和类型检查器识别，但当前运行时可能在属性查找阶段直接失败。
\item
  该方法属于运动命令。方法名包含
  \passthrough{\lstinline!stop!}、\passthrough{\lstinline!damp!} 或
  \passthrough{\lstinline!zero!} 也不代表它是独立物理急停。
\end{itemize}

\textbf{用法}

\begin{lstlisting}[language=Python]
from typing import TYPE_CHECKING

if TYPE_CHECKING:
    def planned_damp(obj: SportClient) -> int:
        return obj.damp()
\end{lstlisting}

\begin{quote}
{[}!CAUTION{]} 上面的 \passthrough{\lstinline!TYPE\_CHECKING!}
示例只表达计划调用形态。该分支在普通 Python
运行时为假，不代表当前扩展能执行此接口。
\end{quote}

\hypertarget{unitree-sdk2-cpp-robot-as2-sportclient-balance-stand-1}{%
\paragraph{\texorpdfstring{\texttt{unitree\_sdk2\_cpp.robot.as2.SportClient.balance\_stand}}{unitree\_sdk2\_cpp.robot.as2.SportClient.balance\_stand}}\label{unitree-sdk2-cpp-robot-as2-sportclient-balance-stand-1}}

对应 C++ SDK 操作
\passthrough{\lstinline!BalanceStand()!}。上游头文件没有可直接生成的业务说明时，本参考只保证签名映射，精确语义需查目标型号协议。

\textbf{签名}

\begin{lstlisting}[language=Python]
def balance_stand(self) -> int
\end{lstlisting}

\textbf{可用性与安全性}

\textbf{\passthrough{\lstinline!SIGNATURE\_ONLY!} /
\passthrough{\lstinline!MOTION\_COMMAND!}}：仅用于补全和静态检查。当前二进制没有该
Python 方法；不得按可执行运动接口使用。

\textbf{参数}

无显式参数。实例方法中的 \passthrough{\lstinline!self!} 由 Python
自动传入。

\textbf{返回值}

\begin{longtable}[]{@{}
  >{\raggedright\arraybackslash}p{(\columnwidth - 4\tabcolsep) * \real{0.18}}
  >{\raggedright\arraybackslash}p{(\columnwidth - 4\tabcolsep) * \real{0.20}}
  >{\raggedright\arraybackslash}p{(\columnwidth - 4\tabcolsep) * \real{0.62}}@{}}
\toprule
位置 & 类型 & 含义 \\
\midrule
\endhead
返回值 & \passthrough{\lstinline!int!} & SDK 状态码；Unitree 示例通常以
\passthrough{\lstinline!0!} 表示成功，非零值需按具体服务定义解释。 \\
\bottomrule
\end{longtable}

\textbf{对应 C++}

\begin{itemize}
\tightlist
\item
  类：\passthrough{\lstinline!unitree::robot::as2::SportClient!}
\item
  签名：\passthrough{\lstinline!BalanceStand()!}
\item
  绑定策略：\passthrough{\lstinline!DIRECT!}
\item
  声明位置：\passthrough{\lstinline!include/unitree/robot/as2/sport/sport\_client.hpp:103!}
\end{itemize}

\textbf{说明}

\begin{itemize}
\tightlist
\item
  示例是局部调用片段；\passthrough{\lstinline!obj!}
  和参数变量表示已经按上层业务校验并准备好的对象。
\item
  该条目可以被编辑器和类型检查器识别，但当前运行时可能在属性查找阶段直接失败。
\item
  该方法属于运动命令。方法名包含
  \passthrough{\lstinline!stop!}、\passthrough{\lstinline!damp!} 或
  \passthrough{\lstinline!zero!} 也不代表它是独立物理急停。
\end{itemize}

\textbf{用法}

\begin{lstlisting}[language=Python]
from typing import TYPE_CHECKING

if TYPE_CHECKING:
    def planned_balance_stand(obj: SportClient) -> int:
        return obj.balance_stand()
\end{lstlisting}

\begin{quote}
{[}!CAUTION{]} 上面的 \passthrough{\lstinline!TYPE\_CHECKING!}
示例只表达计划调用形态。该分支在普通 Python
运行时为假，不代表当前扩展能执行此接口。
\end{quote}

\hypertarget{unitree-sdk2-cpp-robot-as2-sportclient-stop-move-1}{%
\paragraph{\texorpdfstring{\texttt{unitree\_sdk2\_cpp.robot.as2.SportClient.stop\_move}}{unitree\_sdk2\_cpp.robot.as2.SportClient.stop\_move}}\label{unitree-sdk2-cpp-robot-as2-sportclient-stop-move-1}}

请求停止对应 SDK 操作。该名称不等同于经过验证的物理急停。

\textbf{签名}

\begin{lstlisting}[language=Python]
def stop_move(self) -> int
\end{lstlisting}

\textbf{可用性与安全性}

\textbf{\passthrough{\lstinline!SIGNATURE\_ONLY!} /
\passthrough{\lstinline!MOTION\_COMMAND!}}：仅用于补全和静态检查。当前二进制没有该
Python 方法；不得按可执行运动接口使用。

\textbf{参数}

无显式参数。实例方法中的 \passthrough{\lstinline!self!} 由 Python
自动传入。

\textbf{返回值}

\begin{longtable}[]{@{}
  >{\raggedright\arraybackslash}p{(\columnwidth - 4\tabcolsep) * \real{0.18}}
  >{\raggedright\arraybackslash}p{(\columnwidth - 4\tabcolsep) * \real{0.20}}
  >{\raggedright\arraybackslash}p{(\columnwidth - 4\tabcolsep) * \real{0.62}}@{}}
\toprule
位置 & 类型 & 含义 \\
\midrule
\endhead
返回值 & \passthrough{\lstinline!int!} & SDK 状态码；Unitree 示例通常以
\passthrough{\lstinline!0!} 表示成功，非零值需按具体服务定义解释。 \\
\bottomrule
\end{longtable}

\textbf{对应 C++}

\begin{itemize}
\tightlist
\item
  类：\passthrough{\lstinline!unitree::robot::as2::SportClient!}
\item
  签名：\passthrough{\lstinline!StopMove()!}
\item
  绑定策略：\passthrough{\lstinline!DIRECT!}
\item
  声明位置：\passthrough{\lstinline!include/unitree/robot/as2/sport/sport\_client.hpp:108!}
\end{itemize}

\textbf{说明}

\begin{itemize}
\tightlist
\item
  示例是局部调用片段；\passthrough{\lstinline!obj!}
  和参数变量表示已经按上层业务校验并准备好的对象。
\item
  该条目可以被编辑器和类型检查器识别，但当前运行时可能在属性查找阶段直接失败。
\item
  该方法属于运动命令。方法名包含
  \passthrough{\lstinline!stop!}、\passthrough{\lstinline!damp!} 或
  \passthrough{\lstinline!zero!} 也不代表它是独立物理急停。
\end{itemize}

\textbf{用法}

\begin{lstlisting}[language=Python]
from typing import TYPE_CHECKING

if TYPE_CHECKING:
    def planned_stop_move(obj: SportClient) -> int:
        return obj.stop_move()
\end{lstlisting}

\begin{quote}
{[}!CAUTION{]} 上面的 \passthrough{\lstinline!TYPE\_CHECKING!}
示例只表达计划调用形态。该分支在普通 Python
运行时为假，不代表当前扩展能执行此接口。
\end{quote}

\hypertarget{unitree-sdk2-cpp-robot-as2-sportclient-stand-up-1}{%
\paragraph{\texorpdfstring{\texttt{unitree\_sdk2\_cpp.robot.as2.SportClient.stand\_up}}{unitree\_sdk2\_cpp.robot.as2.SportClient.stand\_up}}\label{unitree-sdk2-cpp-robot-as2-sportclient-stand-up-1}}

对应 C++ SDK 操作
\passthrough{\lstinline!StandUp()!}。上游头文件没有可直接生成的业务说明时，本参考只保证签名映射，精确语义需查目标型号协议。

\textbf{签名}

\begin{lstlisting}[language=Python]
def stand_up(self) -> int
\end{lstlisting}

\textbf{可用性与安全性}

\textbf{\passthrough{\lstinline!SIGNATURE\_ONLY!} /
\passthrough{\lstinline!MOTION\_COMMAND!}}：仅用于补全和静态检查。当前二进制没有该
Python 方法；不得按可执行运动接口使用。

\textbf{参数}

无显式参数。实例方法中的 \passthrough{\lstinline!self!} 由 Python
自动传入。

\textbf{返回值}

\begin{longtable}[]{@{}
  >{\raggedright\arraybackslash}p{(\columnwidth - 4\tabcolsep) * \real{0.18}}
  >{\raggedright\arraybackslash}p{(\columnwidth - 4\tabcolsep) * \real{0.20}}
  >{\raggedright\arraybackslash}p{(\columnwidth - 4\tabcolsep) * \real{0.62}}@{}}
\toprule
位置 & 类型 & 含义 \\
\midrule
\endhead
返回值 & \passthrough{\lstinline!int!} & SDK 状态码；Unitree 示例通常以
\passthrough{\lstinline!0!} 表示成功，非零值需按具体服务定义解释。 \\
\bottomrule
\end{longtable}

\textbf{对应 C++}

\begin{itemize}
\tightlist
\item
  类：\passthrough{\lstinline!unitree::robot::as2::SportClient!}
\item
  签名：\passthrough{\lstinline!StandUp()!}
\item
  绑定策略：\passthrough{\lstinline!DIRECT!}
\item
  声明位置：\passthrough{\lstinline!include/unitree/robot/as2/sport/sport\_client.hpp:114!}
\end{itemize}

\textbf{说明}

\begin{itemize}
\tightlist
\item
  示例是局部调用片段；\passthrough{\lstinline!obj!}
  和参数变量表示已经按上层业务校验并准备好的对象。
\item
  该条目可以被编辑器和类型检查器识别，但当前运行时可能在属性查找阶段直接失败。
\item
  该方法属于运动命令。方法名包含
  \passthrough{\lstinline!stop!}、\passthrough{\lstinline!damp!} 或
  \passthrough{\lstinline!zero!} 也不代表它是独立物理急停。
\end{itemize}

\textbf{用法}

\begin{lstlisting}[language=Python]
from typing import TYPE_CHECKING

if TYPE_CHECKING:
    def planned_stand_up(obj: SportClient) -> int:
        return obj.stand_up()
\end{lstlisting}

\begin{quote}
{[}!CAUTION{]} 上面的 \passthrough{\lstinline!TYPE\_CHECKING!}
示例只表达计划调用形态。该分支在普通 Python
运行时为假，不代表当前扩展能执行此接口。
\end{quote}

\hypertarget{unitree-sdk2-cpp-robot-as2-sportclient-stand-down-1}{%
\paragraph{\texorpdfstring{\texttt{unitree\_sdk2\_cpp.robot.as2.SportClient.stand\_down}}{unitree\_sdk2\_cpp.robot.as2.SportClient.stand\_down}}\label{unitree-sdk2-cpp-robot-as2-sportclient-stand-down-1}}

对应 C++ SDK 操作
\passthrough{\lstinline!StandDown()!}。上游头文件没有可直接生成的业务说明时，本参考只保证签名映射，精确语义需查目标型号协议。

\textbf{签名}

\begin{lstlisting}[language=Python]
def stand_down(self) -> int
\end{lstlisting}

\textbf{可用性与安全性}

\textbf{\passthrough{\lstinline!SIGNATURE\_ONLY!} /
\passthrough{\lstinline!MOTION\_COMMAND!}}：仅用于补全和静态检查。当前二进制没有该
Python 方法；不得按可执行运动接口使用。

\textbf{参数}

无显式参数。实例方法中的 \passthrough{\lstinline!self!} 由 Python
自动传入。

\textbf{返回值}

\begin{longtable}[]{@{}
  >{\raggedright\arraybackslash}p{(\columnwidth - 4\tabcolsep) * \real{0.18}}
  >{\raggedright\arraybackslash}p{(\columnwidth - 4\tabcolsep) * \real{0.20}}
  >{\raggedright\arraybackslash}p{(\columnwidth - 4\tabcolsep) * \real{0.62}}@{}}
\toprule
位置 & 类型 & 含义 \\
\midrule
\endhead
返回值 & \passthrough{\lstinline!int!} & SDK 状态码；Unitree 示例通常以
\passthrough{\lstinline!0!} 表示成功，非零值需按具体服务定义解释。 \\
\bottomrule
\end{longtable}

\textbf{对应 C++}

\begin{itemize}
\tightlist
\item
  类：\passthrough{\lstinline!unitree::robot::as2::SportClient!}
\item
  签名：\passthrough{\lstinline!StandDown()!}
\item
  绑定策略：\passthrough{\lstinline!DIRECT!}
\item
  声明位置：\passthrough{\lstinline!include/unitree/robot/as2/sport/sport\_client.hpp:120!}
\end{itemize}

\textbf{说明}

\begin{itemize}
\tightlist
\item
  示例是局部调用片段；\passthrough{\lstinline!obj!}
  和参数变量表示已经按上层业务校验并准备好的对象。
\item
  该条目可以被编辑器和类型检查器识别，但当前运行时可能在属性查找阶段直接失败。
\item
  该方法属于运动命令。方法名包含
  \passthrough{\lstinline!stop!}、\passthrough{\lstinline!damp!} 或
  \passthrough{\lstinline!zero!} 也不代表它是独立物理急停。
\end{itemize}

\textbf{用法}

\begin{lstlisting}[language=Python]
from typing import TYPE_CHECKING

if TYPE_CHECKING:
    def planned_stand_down(obj: SportClient) -> int:
        return obj.stand_down()
\end{lstlisting}

\begin{quote}
{[}!CAUTION{]} 上面的 \passthrough{\lstinline!TYPE\_CHECKING!}
示例只表达计划调用形态。该分支在普通 Python
运行时为假，不代表当前扩展能执行此接口。
\end{quote}

\hypertarget{unitree-sdk2-cpp-robot-as2-sportclient-recovery-stand-1}{%
\paragraph{\texorpdfstring{\texttt{unitree\_sdk2\_cpp.robot.as2.SportClient.recovery\_stand}}{unitree\_sdk2\_cpp.robot.as2.SportClient.recovery\_stand}}\label{unitree-sdk2-cpp-robot-as2-sportclient-recovery-stand-1}}

对应 C++ SDK 操作
\passthrough{\lstinline!RecoveryStand()!}。上游头文件没有可直接生成的业务说明时，本参考只保证签名映射，精确语义需查目标型号协议。

\textbf{签名}

\begin{lstlisting}[language=Python]
def recovery_stand(self) -> int
\end{lstlisting}

\textbf{可用性与安全性}

\textbf{\passthrough{\lstinline!SIGNATURE\_ONLY!} /
\passthrough{\lstinline!MOTION\_COMMAND!}}：仅用于补全和静态检查。当前二进制没有该
Python 方法；不得按可执行运动接口使用。

\textbf{参数}

无显式参数。实例方法中的 \passthrough{\lstinline!self!} 由 Python
自动传入。

\textbf{返回值}

\begin{longtable}[]{@{}
  >{\raggedright\arraybackslash}p{(\columnwidth - 4\tabcolsep) * \real{0.18}}
  >{\raggedright\arraybackslash}p{(\columnwidth - 4\tabcolsep) * \real{0.20}}
  >{\raggedright\arraybackslash}p{(\columnwidth - 4\tabcolsep) * \real{0.62}}@{}}
\toprule
位置 & 类型 & 含义 \\
\midrule
\endhead
返回值 & \passthrough{\lstinline!int!} & SDK 状态码；Unitree 示例通常以
\passthrough{\lstinline!0!} 表示成功，非零值需按具体服务定义解释。 \\
\bottomrule
\end{longtable}

\textbf{对应 C++}

\begin{itemize}
\tightlist
\item
  类：\passthrough{\lstinline!unitree::robot::as2::SportClient!}
\item
  签名：\passthrough{\lstinline!RecoveryStand()!}
\item
  绑定策略：\passthrough{\lstinline!DIRECT!}
\item
  声明位置：\passthrough{\lstinline!include/unitree/robot/as2/sport/sport\_client.hpp:126!}
\end{itemize}

\textbf{说明}

\begin{itemize}
\tightlist
\item
  示例是局部调用片段；\passthrough{\lstinline!obj!}
  和参数变量表示已经按上层业务校验并准备好的对象。
\item
  该条目可以被编辑器和类型检查器识别，但当前运行时可能在属性查找阶段直接失败。
\item
  该方法属于运动命令。方法名包含
  \passthrough{\lstinline!stop!}、\passthrough{\lstinline!damp!} 或
  \passthrough{\lstinline!zero!} 也不代表它是独立物理急停。
\end{itemize}

\textbf{用法}

\begin{lstlisting}[language=Python]
from typing import TYPE_CHECKING

if TYPE_CHECKING:
    def planned_recovery_stand(obj: SportClient) -> int:
        return obj.recovery_stand()
\end{lstlisting}

\begin{quote}
{[}!CAUTION{]} 上面的 \passthrough{\lstinline!TYPE\_CHECKING!}
示例只表达计划调用形态。该分支在普通 Python
运行时为假，不代表当前扩展能执行此接口。
\end{quote}

\hypertarget{unitree-sdk2-cpp-robot-as2-sportclient-euler-1}{%
\paragraph{\texorpdfstring{\texttt{unitree\_sdk2\_cpp.robot.as2.SportClient.euler}}{unitree\_sdk2\_cpp.robot.as2.SportClient.euler}}\label{unitree-sdk2-cpp-robot-as2-sportclient-euler-1}}

对应 C++ SDK 操作
\passthrough{\lstinline!Euler(float, float, float)!}。上游头文件没有可直接生成的业务说明时，本参考只保证签名映射，精确语义需查目标型号协议。

\textbf{签名}

\begin{lstlisting}[language=Python]
def euler(self, roll: float, pitch: float, yaw: float) -> int
\end{lstlisting}

\textbf{可用性与安全性}

\textbf{\passthrough{\lstinline!SIGNATURE\_ONLY!} /
\passthrough{\lstinline!MOTION\_COMMAND!}}：仅用于补全和静态检查。当前二进制没有该
Python 方法；不得按可执行运动接口使用。

\textbf{参数}

\begin{longtable}[]{@{}
  >{\raggedright\arraybackslash}p{(\columnwidth - 6\tabcolsep) * \real{0.18}}
  >{\raggedright\arraybackslash}p{(\columnwidth - 6\tabcolsep) * \real{0.24}}
  >{\raggedright\arraybackslash}p{(\columnwidth - 6\tabcolsep) * \real{0.18}}
  >{\raggedright\arraybackslash}p{(\columnwidth - 6\tabcolsep) * \real{0.40}}@{}}
\toprule
名称 & Python 类型 & 默认值 & 含义 \\
\midrule
\endhead
\passthrough{\lstinline!roll!} & \passthrough{\lstinline!float!} & 必填
& 传给该接口的 \passthrough{\lstinline!roll!} 参数，Python 类型为
\passthrough{\lstinline!float!}。精确含义、单位、范围和枚举值需查目标型号对应头文件与协议。
对应 C++ 参数 \passthrough{\lstinline!roll: float!}。
底层为浮点数；类型本身不说明单位、坐标系或业务安全范围。 \\
\passthrough{\lstinline!pitch!} & \passthrough{\lstinline!float!} & 必填
& 传给该接口的 \passthrough{\lstinline!pitch!} 参数，Python 类型为
\passthrough{\lstinline!float!}。精确含义、单位、范围和枚举值需查目标型号对应头文件与协议。
对应 C++ 参数 \passthrough{\lstinline!pitch: float!}。
底层为浮点数；类型本身不说明单位、坐标系或业务安全范围。 \\
\passthrough{\lstinline!yaw!} & \passthrough{\lstinline!float!} & 必填 &
偏航角或偏航目标参数。参考系、单位和范围必须查目标型号协议。 对应 C++
参数 \passthrough{\lstinline!yaw: float!}。
底层为浮点数；类型本身不说明单位、坐标系或业务安全范围。 \\
\bottomrule
\end{longtable}

\textbf{返回值}

\begin{longtable}[]{@{}
  >{\raggedright\arraybackslash}p{(\columnwidth - 4\tabcolsep) * \real{0.18}}
  >{\raggedright\arraybackslash}p{(\columnwidth - 4\tabcolsep) * \real{0.20}}
  >{\raggedright\arraybackslash}p{(\columnwidth - 4\tabcolsep) * \real{0.62}}@{}}
\toprule
位置 & 类型 & 含义 \\
\midrule
\endhead
返回值 & \passthrough{\lstinline!int!} & SDK 状态码；Unitree 示例通常以
\passthrough{\lstinline!0!} 表示成功，非零值需按具体服务定义解释。 \\
\bottomrule
\end{longtable}

\textbf{对应 C++}

\begin{itemize}
\tightlist
\item
  类：\passthrough{\lstinline!unitree::robot::as2::SportClient!}
\item
  签名：\passthrough{\lstinline!Euler(float, float, float)!}
\item
  绑定策略：\passthrough{\lstinline!DIRECT!}
\item
  声明位置：\passthrough{\lstinline!include/unitree/robot/as2/sport/sport\_client.hpp:132!}
\end{itemize}

\textbf{说明}

\begin{itemize}
\tightlist
\item
  示例是局部调用片段；\passthrough{\lstinline!obj!}
  和参数变量表示已经按上层业务校验并准备好的对象。
\item
  该条目可以被编辑器和类型检查器识别，但当前运行时可能在属性查找阶段直接失败。
\item
  该方法属于运动命令。方法名包含
  \passthrough{\lstinline!stop!}、\passthrough{\lstinline!damp!} 或
  \passthrough{\lstinline!zero!} 也不代表它是独立物理急停。
\end{itemize}

\textbf{用法}

\begin{lstlisting}[language=Python]
from typing import TYPE_CHECKING

if TYPE_CHECKING:
    def planned_euler(obj: SportClient, roll: float, pitch: float, yaw: float) -> int:
        return obj.euler(roll=roll, pitch=pitch, yaw=yaw)
\end{lstlisting}

\begin{quote}
{[}!CAUTION{]} 上面的 \passthrough{\lstinline!TYPE\_CHECKING!}
示例只表达计划调用形态。该分支在普通 Python
运行时为假，不代表当前扩展能执行此接口。
\end{quote}

\hypertarget{unitree-sdk2-cpp-robot-as2-sportclient-move-1}{%
\paragraph{\texorpdfstring{\texttt{unitree\_sdk2\_cpp.robot.as2.SportClient.move}}{unitree\_sdk2\_cpp.robot.as2.SportClient.move}}\label{unitree-sdk2-cpp-robot-as2-sportclient-move-1}}

对应 C++ SDK 操作
\passthrough{\lstinline!Move(float, float, float)!}。上游头文件没有可直接生成的业务说明时，本参考只保证签名映射，精确语义需查目标型号协议。

\textbf{签名}

\begin{lstlisting}[language=Python]
def move(self, vx: float, vy: float, vyaw: float) -> int
\end{lstlisting}

\textbf{可用性与安全性}

\textbf{\passthrough{\lstinline!SIGNATURE\_ONLY!} /
\passthrough{\lstinline!MOTION\_COMMAND!}}：仅用于补全和静态检查。当前二进制没有该
Python 方法；不得按可执行运动接口使用。

\textbf{参数}

\begin{longtable}[]{@{}
  >{\raggedright\arraybackslash}p{(\columnwidth - 6\tabcolsep) * \real{0.18}}
  >{\raggedright\arraybackslash}p{(\columnwidth - 6\tabcolsep) * \real{0.24}}
  >{\raggedright\arraybackslash}p{(\columnwidth - 6\tabcolsep) * \real{0.18}}
  >{\raggedright\arraybackslash}p{(\columnwidth - 6\tabcolsep) * \real{0.40}}@{}}
\toprule
名称 & Python 类型 & 默认值 & 含义 \\
\midrule
\endhead
\passthrough{\lstinline!vx!} & \passthrough{\lstinline!float!} & 必填 &
X 方向速度参数。坐标系、单位、符号和安全范围必须查目标型号运动协议。
对应 C++ 参数 \passthrough{\lstinline!vx: float!}。
底层为浮点数；类型本身不说明单位、坐标系或业务安全范围。 \\
\passthrough{\lstinline!vy!} & \passthrough{\lstinline!float!} & 必填 &
Y 方向速度参数。坐标系、单位、符号和安全范围必须查目标型号运动协议。
对应 C++ 参数 \passthrough{\lstinline!vy: float!}。
底层为浮点数；类型本身不说明单位、坐标系或业务安全范围。 \\
\passthrough{\lstinline!vyaw!} & \passthrough{\lstinline!float!} & 必填
& 偏航角速度参数。单位、符号和安全范围必须查目标型号运动协议。 对应 C++
参数 \passthrough{\lstinline!vyaw: float!}。
底层为浮点数；类型本身不说明单位、坐标系或业务安全范围。 \\
\bottomrule
\end{longtable}

\textbf{返回值}

\begin{longtable}[]{@{}
  >{\raggedright\arraybackslash}p{(\columnwidth - 4\tabcolsep) * \real{0.18}}
  >{\raggedright\arraybackslash}p{(\columnwidth - 4\tabcolsep) * \real{0.20}}
  >{\raggedright\arraybackslash}p{(\columnwidth - 4\tabcolsep) * \real{0.62}}@{}}
\toprule
位置 & 类型 & 含义 \\
\midrule
\endhead
返回值 & \passthrough{\lstinline!int!} & SDK 状态码；Unitree 示例通常以
\passthrough{\lstinline!0!} 表示成功，非零值需按具体服务定义解释。 \\
\bottomrule
\end{longtable}

\textbf{对应 C++}

\begin{itemize}
\tightlist
\item
  类：\passthrough{\lstinline!unitree::robot::as2::SportClient!}
\item
  签名：\passthrough{\lstinline!Move(float, float, float)!}
\item
  绑定策略：\passthrough{\lstinline!DIRECT!}
\item
  声明位置：\passthrough{\lstinline!include/unitree/robot/as2/sport/sport\_client.hpp:143!}
\end{itemize}

\textbf{说明}

\begin{itemize}
\tightlist
\item
  示例是局部调用片段；\passthrough{\lstinline!obj!}
  和参数变量表示已经按上层业务校验并准备好的对象。
\item
  该条目可以被编辑器和类型检查器识别，但当前运行时可能在属性查找阶段直接失败。
\item
  该方法属于运动命令。方法名包含
  \passthrough{\lstinline!stop!}、\passthrough{\lstinline!damp!} 或
  \passthrough{\lstinline!zero!} 也不代表它是独立物理急停。
\end{itemize}

\textbf{用法}

\begin{lstlisting}[language=Python]
from typing import TYPE_CHECKING

if TYPE_CHECKING:
    def planned_move(obj: SportClient, vx: float, vy: float, vyaw: float) -> int:
        return obj.move(vx=vx, vy=vy, vyaw=vyaw)
\end{lstlisting}

\begin{quote}
{[}!CAUTION{]} 上面的 \passthrough{\lstinline!TYPE\_CHECKING!}
示例只表达计划调用形态。该分支在普通 Python
运行时为假，不代表当前扩展能执行此接口。
\end{quote}

\hypertarget{unitree-sdk2-cpp-robot-as2-sportclient-switch-gait-1}{%
\paragraph{\texorpdfstring{\texttt{unitree\_sdk2\_cpp.robot.as2.SportClient.switch\_gait}}{unitree\_sdk2\_cpp.robot.as2.SportClient.switch\_gait}}\label{unitree-sdk2-cpp-robot-as2-sportclient-switch-gait-1}}

对应 C++ SDK 操作
\passthrough{\lstinline!SwitchGait(int)!}。上游头文件没有可直接生成的业务说明时，本参考只保证签名映射，精确语义需查目标型号协议。

\textbf{签名}

\begin{lstlisting}[language=Python]
def switch_gait(self, gait_type: int) -> int
\end{lstlisting}

\textbf{可用性与安全性}

\textbf{\passthrough{\lstinline!SIGNATURE\_ONLY!} /
\passthrough{\lstinline!MOTION\_COMMAND!}}：仅用于补全和静态检查。当前二进制没有该
Python 方法；不得按可执行运动接口使用。

\textbf{参数}

\begin{longtable}[]{@{}
  >{\raggedright\arraybackslash}p{(\columnwidth - 6\tabcolsep) * \real{0.18}}
  >{\raggedright\arraybackslash}p{(\columnwidth - 6\tabcolsep) * \real{0.24}}
  >{\raggedright\arraybackslash}p{(\columnwidth - 6\tabcolsep) * \real{0.18}}
  >{\raggedright\arraybackslash}p{(\columnwidth - 6\tabcolsep) * \real{0.40}}@{}}
\toprule
名称 & Python 类型 & 默认值 & 含义 \\
\midrule
\endhead
\passthrough{\lstinline!gait\_type!} & \passthrough{\lstinline!int!} &
必填 & 传给该接口的 \passthrough{\lstinline!gait!}
\passthrough{\lstinline!type!} 参数，Python 类型为
\passthrough{\lstinline!int!}。精确含义、单位、范围和枚举值需查目标型号对应头文件与协议。
对应 C++ 参数 \passthrough{\lstinline!gait\_type: int!}。 \\
\bottomrule
\end{longtable}

\textbf{返回值}

\begin{longtable}[]{@{}
  >{\raggedright\arraybackslash}p{(\columnwidth - 4\tabcolsep) * \real{0.18}}
  >{\raggedright\arraybackslash}p{(\columnwidth - 4\tabcolsep) * \real{0.20}}
  >{\raggedright\arraybackslash}p{(\columnwidth - 4\tabcolsep) * \real{0.62}}@{}}
\toprule
位置 & 类型 & 含义 \\
\midrule
\endhead
返回值 & \passthrough{\lstinline!int!} & SDK 状态码；Unitree 示例通常以
\passthrough{\lstinline!0!} 表示成功，非零值需按具体服务定义解释。 \\
\bottomrule
\end{longtable}

\textbf{对应 C++}

\begin{itemize}
\tightlist
\item
  类：\passthrough{\lstinline!unitree::robot::as2::SportClient!}
\item
  签名：\passthrough{\lstinline!SwitchGait(int)!}
\item
  绑定策略：\passthrough{\lstinline!DIRECT!}
\item
  声明位置：\passthrough{\lstinline!include/unitree/robot/as2/sport/sport\_client.hpp:154!}
\end{itemize}

\textbf{说明}

\begin{itemize}
\tightlist
\item
  示例是局部调用片段；\passthrough{\lstinline!obj!}
  和参数变量表示已经按上层业务校验并准备好的对象。
\item
  该条目可以被编辑器和类型检查器识别，但当前运行时可能在属性查找阶段直接失败。
\item
  该方法属于运动命令。方法名包含
  \passthrough{\lstinline!stop!}、\passthrough{\lstinline!damp!} 或
  \passthrough{\lstinline!zero!} 也不代表它是独立物理急停。
\end{itemize}

\textbf{用法}

\begin{lstlisting}[language=Python]
from typing import TYPE_CHECKING

if TYPE_CHECKING:
    def planned_switch_gait(obj: SportClient, gait_type: int) -> int:
        return obj.switch_gait(gait_type=gait_type)
\end{lstlisting}

\begin{quote}
{[}!CAUTION{]} 上面的 \passthrough{\lstinline!TYPE\_CHECKING!}
示例只表达计划调用形态。该分支在普通 Python
运行时为假，不代表当前扩展能执行此接口。
\end{quote}

\hypertarget{unitree-sdk2-cpp-robot-as2-sportclient-body-height-1}{%
\paragraph{\texorpdfstring{\texttt{unitree\_sdk2\_cpp.robot.as2.SportClient.body\_height}}{unitree\_sdk2\_cpp.robot.as2.SportClient.body\_height}}\label{unitree-sdk2-cpp-robot-as2-sportclient-body-height-1}}

对应 C++ SDK 操作
\passthrough{\lstinline!BodyHeight(float)!}。上游头文件没有可直接生成的业务说明时，本参考只保证签名映射，精确语义需查目标型号协议。

\textbf{签名}

\begin{lstlisting}[language=Python]
def body_height(self, height: float) -> int
\end{lstlisting}

\textbf{可用性与安全性}

\textbf{\passthrough{\lstinline!SIGNATURE\_ONLY!} /
\passthrough{\lstinline!MOTION\_COMMAND!}}：仅用于补全和静态检查。当前二进制没有该
Python 方法；不得按可执行运动接口使用。

\textbf{参数}

\begin{longtable}[]{@{}
  >{\raggedright\arraybackslash}p{(\columnwidth - 6\tabcolsep) * \real{0.18}}
  >{\raggedright\arraybackslash}p{(\columnwidth - 6\tabcolsep) * \real{0.24}}
  >{\raggedright\arraybackslash}p{(\columnwidth - 6\tabcolsep) * \real{0.18}}
  >{\raggedright\arraybackslash}p{(\columnwidth - 6\tabcolsep) * \real{0.40}}@{}}
\toprule
名称 & Python 类型 & 默认值 & 含义 \\
\midrule
\endhead
\passthrough{\lstinline!height!} & \passthrough{\lstinline!float!} &
必填 & 传给该接口的 高度 参数，Python 类型为
\passthrough{\lstinline!float!}。精确含义、单位、范围和枚举值需查目标型号对应头文件与协议。
对应 C++ 参数 \passthrough{\lstinline!height: float!}。
底层为浮点数；类型本身不说明单位、坐标系或业务安全范围。 \\
\bottomrule
\end{longtable}

\textbf{返回值}

\begin{longtable}[]{@{}
  >{\raggedright\arraybackslash}p{(\columnwidth - 4\tabcolsep) * \real{0.18}}
  >{\raggedright\arraybackslash}p{(\columnwidth - 4\tabcolsep) * \real{0.20}}
  >{\raggedright\arraybackslash}p{(\columnwidth - 4\tabcolsep) * \real{0.62}}@{}}
\toprule
位置 & 类型 & 含义 \\
\midrule
\endhead
返回值 & \passthrough{\lstinline!int!} & SDK 状态码；Unitree 示例通常以
\passthrough{\lstinline!0!} 表示成功，非零值需按具体服务定义解释。 \\
\bottomrule
\end{longtable}

\textbf{对应 C++}

\begin{itemize}
\tightlist
\item
  类：\passthrough{\lstinline!unitree::robot::as2::SportClient!}
\item
  签名：\passthrough{\lstinline!BodyHeight(float)!}
\item
  绑定策略：\passthrough{\lstinline!DIRECT!}
\item
  声明位置：\passthrough{\lstinline!include/unitree/robot/as2/sport/sport\_client.hpp:163!}
\end{itemize}

\textbf{说明}

\begin{itemize}
\tightlist
\item
  示例是局部调用片段；\passthrough{\lstinline!obj!}
  和参数变量表示已经按上层业务校验并准备好的对象。
\item
  该条目可以被编辑器和类型检查器识别，但当前运行时可能在属性查找阶段直接失败。
\item
  该方法属于运动命令。方法名包含
  \passthrough{\lstinline!stop!}、\passthrough{\lstinline!damp!} 或
  \passthrough{\lstinline!zero!} 也不代表它是独立物理急停。
\end{itemize}

\textbf{用法}

\begin{lstlisting}[language=Python]
from typing import TYPE_CHECKING

if TYPE_CHECKING:
    def planned_body_height(obj: SportClient, height: float) -> int:
        return obj.body_height(height=height)
\end{lstlisting}

\begin{quote}
{[}!CAUTION{]} 上面的 \passthrough{\lstinline!TYPE\_CHECKING!}
示例只表达计划调用形态。该分支在普通 Python
运行时为假，不代表当前扩展能执行此接口。
\end{quote}

\hypertarget{unitree-sdk2-cpp-robot-as2-sportclient-speed-level-1}{%
\paragraph{\texorpdfstring{\texttt{unitree\_sdk2\_cpp.robot.as2.SportClient.speed\_level}}{unitree\_sdk2\_cpp.robot.as2.SportClient.speed\_level}}\label{unitree-sdk2-cpp-robot-as2-sportclient-speed-level-1}}

对应 C++ SDK 操作
\passthrough{\lstinline!SpeedLevel(int)!}。上游头文件没有可直接生成的业务说明时，本参考只保证签名映射，精确语义需查目标型号协议。

\textbf{签名}

\begin{lstlisting}[language=Python]
def speed_level(self, level: int) -> int
\end{lstlisting}

\textbf{可用性与安全性}

\textbf{\passthrough{\lstinline!SIGNATURE\_ONLY!} /
\passthrough{\lstinline!MOTION\_COMMAND!}}：仅用于补全和静态检查。当前二进制没有该
Python 方法；不得按可执行运动接口使用。

\textbf{参数}

\begin{longtable}[]{@{}
  >{\raggedright\arraybackslash}p{(\columnwidth - 6\tabcolsep) * \real{0.18}}
  >{\raggedright\arraybackslash}p{(\columnwidth - 6\tabcolsep) * \real{0.24}}
  >{\raggedright\arraybackslash}p{(\columnwidth - 6\tabcolsep) * \real{0.18}}
  >{\raggedright\arraybackslash}p{(\columnwidth - 6\tabcolsep) * \real{0.40}}@{}}
\toprule
名称 & Python 类型 & 默认值 & 含义 \\
\midrule
\endhead
\passthrough{\lstinline!level!} & \passthrough{\lstinline!int!} & 必填 &
传给该接口的 \passthrough{\lstinline!level!} 参数，Python 类型为
\passthrough{\lstinline!int!}。精确含义、单位、范围和枚举值需查目标型号对应头文件与协议。
对应 C++ 参数 \passthrough{\lstinline!level: int!}。 \\
\bottomrule
\end{longtable}

\textbf{返回值}

\begin{longtable}[]{@{}
  >{\raggedright\arraybackslash}p{(\columnwidth - 4\tabcolsep) * \real{0.18}}
  >{\raggedright\arraybackslash}p{(\columnwidth - 4\tabcolsep) * \real{0.20}}
  >{\raggedright\arraybackslash}p{(\columnwidth - 4\tabcolsep) * \real{0.62}}@{}}
\toprule
位置 & 类型 & 含义 \\
\midrule
\endhead
返回值 & \passthrough{\lstinline!int!} & SDK 状态码；Unitree 示例通常以
\passthrough{\lstinline!0!} 表示成功，非零值需按具体服务定义解释。 \\
\bottomrule
\end{longtable}

\textbf{对应 C++}

\begin{itemize}
\tightlist
\item
  类：\passthrough{\lstinline!unitree::robot::as2::SportClient!}
\item
  签名：\passthrough{\lstinline!SpeedLevel(int)!}
\item
  绑定策略：\passthrough{\lstinline!DIRECT!}
\item
  声明位置：\passthrough{\lstinline!include/unitree/robot/as2/sport/sport\_client.hpp:172!}
\end{itemize}

\textbf{说明}

\begin{itemize}
\tightlist
\item
  示例是局部调用片段；\passthrough{\lstinline!obj!}
  和参数变量表示已经按上层业务校验并准备好的对象。
\item
  该条目可以被编辑器和类型检查器识别，但当前运行时可能在属性查找阶段直接失败。
\item
  该方法属于运动命令。方法名包含
  \passthrough{\lstinline!stop!}、\passthrough{\lstinline!damp!} 或
  \passthrough{\lstinline!zero!} 也不代表它是独立物理急停。
\end{itemize}

\textbf{用法}

\begin{lstlisting}[language=Python]
from typing import TYPE_CHECKING

if TYPE_CHECKING:
    def planned_speed_level(obj: SportClient, level: int) -> int:
        return obj.speed_level(level=level)
\end{lstlisting}

\begin{quote}
{[}!CAUTION{]} 上面的 \passthrough{\lstinline!TYPE\_CHECKING!}
示例只表达计划调用形态。该分支在普通 Python
运行时为假，不代表当前扩展能执行此接口。
\end{quote}

\hypertarget{unitree-sdk2-cpp-robot-as2-sportclient-body-position-1}{%
\paragraph{\texorpdfstring{\texttt{unitree\_sdk2\_cpp.robot.as2.SportClient.body\_position}}{unitree\_sdk2\_cpp.robot.as2.SportClient.body\_position}}\label{unitree-sdk2-cpp-robot-as2-sportclient-body-position-1}}

对应 C++ SDK 操作
\passthrough{\lstinline!BodyPosition(float, float, float, float)!}。上游头文件没有可直接生成的业务说明时，本参考只保证签名映射，精确语义需查目标型号协议。

\textbf{签名}

\begin{lstlisting}[language=Python]
def body_position(self, x: float, y: float, z: float, yaw: float) -> int
\end{lstlisting}

\textbf{可用性与安全性}

\textbf{\passthrough{\lstinline!SIGNATURE\_ONLY!} /
\passthrough{\lstinline!MOTION\_COMMAND!}}：仅用于补全和静态检查。当前二进制没有该
Python 方法；不得按可执行运动接口使用。

\textbf{参数}

\begin{longtable}[]{@{}
  >{\raggedright\arraybackslash}p{(\columnwidth - 6\tabcolsep) * \real{0.18}}
  >{\raggedright\arraybackslash}p{(\columnwidth - 6\tabcolsep) * \real{0.24}}
  >{\raggedright\arraybackslash}p{(\columnwidth - 6\tabcolsep) * \real{0.18}}
  >{\raggedright\arraybackslash}p{(\columnwidth - 6\tabcolsep) * \real{0.40}}@{}}
\toprule
名称 & Python 类型 & 默认值 & 含义 \\
\midrule
\endhead
\passthrough{\lstinline!x!} & \passthrough{\lstinline!float!} & 必填 & X
方向位置或分量参数。参考系、单位和范围必须查当前方法及目标型号协议。
对应 C++ 参数 \passthrough{\lstinline!x: float!}。
底层为浮点数；类型本身不说明单位、坐标系或业务安全范围。 \\
\passthrough{\lstinline!y!} & \passthrough{\lstinline!float!} & 必填 & Y
方向位置或分量参数。参考系、单位和范围必须查当前方法及目标型号协议。
对应 C++ 参数 \passthrough{\lstinline!y: float!}。
底层为浮点数；类型本身不说明单位、坐标系或业务安全范围。 \\
\passthrough{\lstinline!z!} & \passthrough{\lstinline!float!} & 必填 & Z
方向位置或分量参数。参考系、单位和范围必须查当前方法及目标型号协议。
对应 C++ 参数 \passthrough{\lstinline!z: float!}。
底层为浮点数；类型本身不说明单位、坐标系或业务安全范围。 \\
\passthrough{\lstinline!yaw!} & \passthrough{\lstinline!float!} & 必填 &
偏航角或偏航目标参数。参考系、单位和范围必须查目标型号协议。 对应 C++
参数 \passthrough{\lstinline!yaw: float!}。
底层为浮点数；类型本身不说明单位、坐标系或业务安全范围。 \\
\bottomrule
\end{longtable}

\textbf{返回值}

\begin{longtable}[]{@{}
  >{\raggedright\arraybackslash}p{(\columnwidth - 4\tabcolsep) * \real{0.18}}
  >{\raggedright\arraybackslash}p{(\columnwidth - 4\tabcolsep) * \real{0.20}}
  >{\raggedright\arraybackslash}p{(\columnwidth - 4\tabcolsep) * \real{0.62}}@{}}
\toprule
位置 & 类型 & 含义 \\
\midrule
\endhead
返回值 & \passthrough{\lstinline!int!} & SDK 状态码；Unitree 示例通常以
\passthrough{\lstinline!0!} 表示成功，非零值需按具体服务定义解释。 \\
\bottomrule
\end{longtable}

\textbf{对应 C++}

\begin{itemize}
\tightlist
\item
  类：\passthrough{\lstinline!unitree::robot::as2::SportClient!}
\item
  签名：\passthrough{\lstinline!BodyPosition(float, float, float, float)!}
\item
  绑定策略：\passthrough{\lstinline!DIRECT!}
\item
  声明位置：\passthrough{\lstinline!include/unitree/robot/as2/sport/sport\_client.hpp:181!}
\end{itemize}

\textbf{说明}

\begin{itemize}
\tightlist
\item
  示例是局部调用片段；\passthrough{\lstinline!obj!}
  和参数变量表示已经按上层业务校验并准备好的对象。
\item
  该条目可以被编辑器和类型检查器识别，但当前运行时可能在属性查找阶段直接失败。
\item
  该方法属于运动命令。方法名包含
  \passthrough{\lstinline!stop!}、\passthrough{\lstinline!damp!} 或
  \passthrough{\lstinline!zero!} 也不代表它是独立物理急停。
\end{itemize}

\textbf{用法}

\begin{lstlisting}[language=Python]
from typing import TYPE_CHECKING

if TYPE_CHECKING:
    def planned_body_position(obj: SportClient, x: float, y: float, z: float, yaw: float) -> int:
        return obj.body_position(x=x, y=y, z=z, yaw=yaw)
\end{lstlisting}

\begin{quote}
{[}!CAUTION{]} 上面的 \passthrough{\lstinline!TYPE\_CHECKING!}
示例只表达计划调用形态。该分支在普通 Python
运行时为假，不代表当前扩展能执行此接口。
\end{quote}

\hypertarget{unitree-sdk2-cpp-robot-as2-sportclient-left-side-gait-1}{%
\paragraph{\texorpdfstring{\texttt{unitree\_sdk2\_cpp.robot.as2.SportClient.left\_side\_gait}}{unitree\_sdk2\_cpp.robot.as2.SportClient.left\_side\_gait}}\label{unitree-sdk2-cpp-robot-as2-sportclient-left-side-gait-1}}

对应 C++ SDK 操作
\passthrough{\lstinline!LeftSideGait(int)!}。上游头文件没有可直接生成的业务说明时，本参考只保证签名映射，精确语义需查目标型号协议。

\textbf{签名}

\begin{lstlisting}[language=Python]
def left_side_gait(self, enter: int) -> int
\end{lstlisting}

\textbf{可用性与安全性}

\textbf{\passthrough{\lstinline!SIGNATURE\_ONLY!} /
\passthrough{\lstinline!MOTION\_COMMAND!}}：仅用于补全和静态检查。当前二进制没有该
Python 方法；不得按可执行运动接口使用。

\textbf{参数}

\begin{longtable}[]{@{}
  >{\raggedright\arraybackslash}p{(\columnwidth - 6\tabcolsep) * \real{0.18}}
  >{\raggedright\arraybackslash}p{(\columnwidth - 6\tabcolsep) * \real{0.24}}
  >{\raggedright\arraybackslash}p{(\columnwidth - 6\tabcolsep) * \real{0.18}}
  >{\raggedright\arraybackslash}p{(\columnwidth - 6\tabcolsep) * \real{0.40}}@{}}
\toprule
名称 & Python 类型 & 默认值 & 含义 \\
\midrule
\endhead
\passthrough{\lstinline!enter!} & \passthrough{\lstinline!int!} & 必填 &
传给该接口的 \passthrough{\lstinline!enter!} 参数，Python 类型为
\passthrough{\lstinline!int!}。精确含义、单位、范围和枚举值需查目标型号对应头文件与协议。
对应 C++ 参数 \passthrough{\lstinline!enter: int!}。 \\
\bottomrule
\end{longtable}

\textbf{返回值}

\begin{longtable}[]{@{}
  >{\raggedright\arraybackslash}p{(\columnwidth - 4\tabcolsep) * \real{0.18}}
  >{\raggedright\arraybackslash}p{(\columnwidth - 4\tabcolsep) * \real{0.20}}
  >{\raggedright\arraybackslash}p{(\columnwidth - 4\tabcolsep) * \real{0.62}}@{}}
\toprule
位置 & 类型 & 含义 \\
\midrule
\endhead
返回值 & \passthrough{\lstinline!int!} & SDK 状态码；Unitree 示例通常以
\passthrough{\lstinline!0!} 表示成功，非零值需按具体服务定义解释。 \\
\bottomrule
\end{longtable}

\textbf{对应 C++}

\begin{itemize}
\tightlist
\item
  类：\passthrough{\lstinline!unitree::robot::as2::SportClient!}
\item
  签名：\passthrough{\lstinline!LeftSideGait(int)!}
\item
  绑定策略：\passthrough{\lstinline!DIRECT!}
\item
  声明位置：\passthrough{\lstinline!include/unitree/robot/as2/sport/sport\_client.hpp:193!}
\end{itemize}

\textbf{说明}

\begin{itemize}
\tightlist
\item
  示例是局部调用片段；\passthrough{\lstinline!obj!}
  和参数变量表示已经按上层业务校验并准备好的对象。
\item
  该条目可以被编辑器和类型检查器识别，但当前运行时可能在属性查找阶段直接失败。
\item
  该方法属于运动命令。方法名包含
  \passthrough{\lstinline!stop!}、\passthrough{\lstinline!damp!} 或
  \passthrough{\lstinline!zero!} 也不代表它是独立物理急停。
\end{itemize}

\textbf{用法}

\begin{lstlisting}[language=Python]
from typing import TYPE_CHECKING

if TYPE_CHECKING:
    def planned_left_side_gait(obj: SportClient, enter: int) -> int:
        return obj.left_side_gait(enter=enter)
\end{lstlisting}

\begin{quote}
{[}!CAUTION{]} 上面的 \passthrough{\lstinline!TYPE\_CHECKING!}
示例只表达计划调用形态。该分支在普通 Python
运行时为假，不代表当前扩展能执行此接口。
\end{quote}

\hypertarget{unitree-sdk2-cpp-robot-as2-sportclient-right-side-gait-1}{%
\paragraph{\texorpdfstring{\texttt{unitree\_sdk2\_cpp.robot.as2.SportClient.right\_side\_gait}}{unitree\_sdk2\_cpp.robot.as2.SportClient.right\_side\_gait}}\label{unitree-sdk2-cpp-robot-as2-sportclient-right-side-gait-1}}

对应 C++ SDK 操作
\passthrough{\lstinline!RightSideGait(int)!}。上游头文件没有可直接生成的业务说明时，本参考只保证签名映射，精确语义需查目标型号协议。

\textbf{签名}

\begin{lstlisting}[language=Python]
def right_side_gait(self, enter: int) -> int
\end{lstlisting}

\textbf{可用性与安全性}

\textbf{\passthrough{\lstinline!SIGNATURE\_ONLY!} /
\passthrough{\lstinline!MOTION\_COMMAND!}}：仅用于补全和静态检查。当前二进制没有该
Python 方法；不得按可执行运动接口使用。

\textbf{参数}

\begin{longtable}[]{@{}
  >{\raggedright\arraybackslash}p{(\columnwidth - 6\tabcolsep) * \real{0.18}}
  >{\raggedright\arraybackslash}p{(\columnwidth - 6\tabcolsep) * \real{0.24}}
  >{\raggedright\arraybackslash}p{(\columnwidth - 6\tabcolsep) * \real{0.18}}
  >{\raggedright\arraybackslash}p{(\columnwidth - 6\tabcolsep) * \real{0.40}}@{}}
\toprule
名称 & Python 类型 & 默认值 & 含义 \\
\midrule
\endhead
\passthrough{\lstinline!enter!} & \passthrough{\lstinline!int!} & 必填 &
传给该接口的 \passthrough{\lstinline!enter!} 参数，Python 类型为
\passthrough{\lstinline!int!}。精确含义、单位、范围和枚举值需查目标型号对应头文件与协议。
对应 C++ 参数 \passthrough{\lstinline!enter: int!}。 \\
\bottomrule
\end{longtable}

\textbf{返回值}

\begin{longtable}[]{@{}
  >{\raggedright\arraybackslash}p{(\columnwidth - 4\tabcolsep) * \real{0.18}}
  >{\raggedright\arraybackslash}p{(\columnwidth - 4\tabcolsep) * \real{0.20}}
  >{\raggedright\arraybackslash}p{(\columnwidth - 4\tabcolsep) * \real{0.62}}@{}}
\toprule
位置 & 类型 & 含义 \\
\midrule
\endhead
返回值 & \passthrough{\lstinline!int!} & SDK 状态码；Unitree 示例通常以
\passthrough{\lstinline!0!} 表示成功，非零值需按具体服务定义解释。 \\
\bottomrule
\end{longtable}

\textbf{对应 C++}

\begin{itemize}
\tightlist
\item
  类：\passthrough{\lstinline!unitree::robot::as2::SportClient!}
\item
  签名：\passthrough{\lstinline!RightSideGait(int)!}
\item
  绑定策略：\passthrough{\lstinline!DIRECT!}
\item
  声明位置：\passthrough{\lstinline!include/unitree/robot/as2/sport/sport\_client.hpp:202!}
\end{itemize}

\textbf{说明}

\begin{itemize}
\tightlist
\item
  示例是局部调用片段；\passthrough{\lstinline!obj!}
  和参数变量表示已经按上层业务校验并准备好的对象。
\item
  该条目可以被编辑器和类型检查器识别，但当前运行时可能在属性查找阶段直接失败。
\item
  该方法属于运动命令。方法名包含
  \passthrough{\lstinline!stop!}、\passthrough{\lstinline!damp!} 或
  \passthrough{\lstinline!zero!} 也不代表它是独立物理急停。
\end{itemize}

\textbf{用法}

\begin{lstlisting}[language=Python]
from typing import TYPE_CHECKING

if TYPE_CHECKING:
    def planned_right_side_gait(obj: SportClient, enter: int) -> int:
        return obj.right_side_gait(enter=enter)
\end{lstlisting}

\begin{quote}
{[}!CAUTION{]} 上面的 \passthrough{\lstinline!TYPE\_CHECKING!}
示例只表达计划调用形态。该分支在普通 Python
运行时为假，不代表当前扩展能执行此接口。
\end{quote}

\hypertarget{unitree-sdk2-cpp-robot-as2-sportclient-hand-stand-1}{%
\paragraph{\texorpdfstring{\texttt{unitree\_sdk2\_cpp.robot.as2.SportClient.hand\_stand}}{unitree\_sdk2\_cpp.robot.as2.SportClient.hand\_stand}}\label{unitree-sdk2-cpp-robot-as2-sportclient-hand-stand-1}}

对应 C++ SDK 操作
\passthrough{\lstinline!HandStand(int)!}。上游头文件没有可直接生成的业务说明时，本参考只保证签名映射，精确语义需查目标型号协议。

\textbf{签名}

\begin{lstlisting}[language=Python]
def hand_stand(self, enter: int) -> int
\end{lstlisting}

\textbf{可用性与安全性}

\textbf{\passthrough{\lstinline!SIGNATURE\_ONLY!} /
\passthrough{\lstinline!MOTION\_COMMAND!}}：仅用于补全和静态检查。当前二进制没有该
Python 方法；不得按可执行运动接口使用。

\textbf{参数}

\begin{longtable}[]{@{}
  >{\raggedright\arraybackslash}p{(\columnwidth - 6\tabcolsep) * \real{0.18}}
  >{\raggedright\arraybackslash}p{(\columnwidth - 6\tabcolsep) * \real{0.24}}
  >{\raggedright\arraybackslash}p{(\columnwidth - 6\tabcolsep) * \real{0.18}}
  >{\raggedright\arraybackslash}p{(\columnwidth - 6\tabcolsep) * \real{0.40}}@{}}
\toprule
名称 & Python 类型 & 默认值 & 含义 \\
\midrule
\endhead
\passthrough{\lstinline!enter!} & \passthrough{\lstinline!int!} & 必填 &
传给该接口的 \passthrough{\lstinline!enter!} 参数，Python 类型为
\passthrough{\lstinline!int!}。精确含义、单位、范围和枚举值需查目标型号对应头文件与协议。
对应 C++ 参数 \passthrough{\lstinline!enter: int!}。 \\
\bottomrule
\end{longtable}

\textbf{返回值}

\begin{longtable}[]{@{}
  >{\raggedright\arraybackslash}p{(\columnwidth - 4\tabcolsep) * \real{0.18}}
  >{\raggedright\arraybackslash}p{(\columnwidth - 4\tabcolsep) * \real{0.20}}
  >{\raggedright\arraybackslash}p{(\columnwidth - 4\tabcolsep) * \real{0.62}}@{}}
\toprule
位置 & 类型 & 含义 \\
\midrule
\endhead
返回值 & \passthrough{\lstinline!int!} & SDK 状态码；Unitree 示例通常以
\passthrough{\lstinline!0!} 表示成功，非零值需按具体服务定义解释。 \\
\bottomrule
\end{longtable}

\textbf{对应 C++}

\begin{itemize}
\tightlist
\item
  类：\passthrough{\lstinline!unitree::robot::as2::SportClient!}
\item
  签名：\passthrough{\lstinline!HandStand(int)!}
\item
  绑定策略：\passthrough{\lstinline!DIRECT!}
\item
  声明位置：\passthrough{\lstinline!include/unitree/robot/as2/sport/sport\_client.hpp:211!}
\end{itemize}

\textbf{说明}

\begin{itemize}
\tightlist
\item
  示例是局部调用片段；\passthrough{\lstinline!obj!}
  和参数变量表示已经按上层业务校验并准备好的对象。
\item
  该条目可以被编辑器和类型检查器识别，但当前运行时可能在属性查找阶段直接失败。
\item
  该方法属于运动命令。方法名包含
  \passthrough{\lstinline!stop!}、\passthrough{\lstinline!damp!} 或
  \passthrough{\lstinline!zero!} 也不代表它是独立物理急停。
\end{itemize}

\textbf{用法}

\begin{lstlisting}[language=Python]
from typing import TYPE_CHECKING

if TYPE_CHECKING:
    def planned_hand_stand(obj: SportClient, enter: int) -> int:
        return obj.hand_stand(enter=enter)
\end{lstlisting}

\begin{quote}
{[}!CAUTION{]} 上面的 \passthrough{\lstinline!TYPE\_CHECKING!}
示例只表达计划调用形态。该分支在普通 Python
运行时为假，不代表当前扩展能执行此接口。
\end{quote}

\hypertarget{unitree-sdk2-cpp-robot-as2-sportclient-biped-stand-1}{%
\paragraph{\texorpdfstring{\texttt{unitree\_sdk2\_cpp.robot.as2.SportClient.biped\_stand}}{unitree\_sdk2\_cpp.robot.as2.SportClient.biped\_stand}}\label{unitree-sdk2-cpp-robot-as2-sportclient-biped-stand-1}}

对应 C++ SDK 操作
\passthrough{\lstinline!BipedStand(int)!}。上游头文件没有可直接生成的业务说明时，本参考只保证签名映射，精确语义需查目标型号协议。

\textbf{签名}

\begin{lstlisting}[language=Python]
def biped_stand(self, enter: int) -> int
\end{lstlisting}

\textbf{可用性与安全性}

\textbf{\passthrough{\lstinline!SIGNATURE\_ONLY!} /
\passthrough{\lstinline!MOTION\_COMMAND!}}：仅用于补全和静态检查。当前二进制没有该
Python 方法；不得按可执行运动接口使用。

\textbf{参数}

\begin{longtable}[]{@{}
  >{\raggedright\arraybackslash}p{(\columnwidth - 6\tabcolsep) * \real{0.18}}
  >{\raggedright\arraybackslash}p{(\columnwidth - 6\tabcolsep) * \real{0.24}}
  >{\raggedright\arraybackslash}p{(\columnwidth - 6\tabcolsep) * \real{0.18}}
  >{\raggedright\arraybackslash}p{(\columnwidth - 6\tabcolsep) * \real{0.40}}@{}}
\toprule
名称 & Python 类型 & 默认值 & 含义 \\
\midrule
\endhead
\passthrough{\lstinline!enter!} & \passthrough{\lstinline!int!} & 必填 &
传给该接口的 \passthrough{\lstinline!enter!} 参数，Python 类型为
\passthrough{\lstinline!int!}。精确含义、单位、范围和枚举值需查目标型号对应头文件与协议。
对应 C++ 参数 \passthrough{\lstinline!enter: int!}。 \\
\bottomrule
\end{longtable}

\textbf{返回值}

\begin{longtable}[]{@{}
  >{\raggedright\arraybackslash}p{(\columnwidth - 4\tabcolsep) * \real{0.18}}
  >{\raggedright\arraybackslash}p{(\columnwidth - 4\tabcolsep) * \real{0.20}}
  >{\raggedright\arraybackslash}p{(\columnwidth - 4\tabcolsep) * \real{0.62}}@{}}
\toprule
位置 & 类型 & 含义 \\
\midrule
\endhead
返回值 & \passthrough{\lstinline!int!} & SDK 状态码；Unitree 示例通常以
\passthrough{\lstinline!0!} 表示成功，非零值需按具体服务定义解释。 \\
\bottomrule
\end{longtable}

\textbf{对应 C++}

\begin{itemize}
\tightlist
\item
  类：\passthrough{\lstinline!unitree::robot::as2::SportClient!}
\item
  签名：\passthrough{\lstinline!BipedStand(int)!}
\item
  绑定策略：\passthrough{\lstinline!DIRECT!}
\item
  声明位置：\passthrough{\lstinline!include/unitree/robot/as2/sport/sport\_client.hpp:220!}
\end{itemize}

\textbf{说明}

\begin{itemize}
\tightlist
\item
  示例是局部调用片段；\passthrough{\lstinline!obj!}
  和参数变量表示已经按上层业务校验并准备好的对象。
\item
  该条目可以被编辑器和类型检查器识别，但当前运行时可能在属性查找阶段直接失败。
\item
  该方法属于运动命令。方法名包含
  \passthrough{\lstinline!stop!}、\passthrough{\lstinline!damp!} 或
  \passthrough{\lstinline!zero!} 也不代表它是独立物理急停。
\end{itemize}

\textbf{用法}

\begin{lstlisting}[language=Python]
from typing import TYPE_CHECKING

if TYPE_CHECKING:
    def planned_biped_stand(obj: SportClient, enter: int) -> int:
        return obj.biped_stand(enter=enter)
\end{lstlisting}

\begin{quote}
{[}!CAUTION{]} 上面的 \passthrough{\lstinline!TYPE\_CHECKING!}
示例只表达计划调用形态。该分支在普通 Python
运行时为假，不代表当前扩展能执行此接口。
\end{quote}

\hypertarget{unitree-sdk2-cpp-robot-as2-sportclient-front-flip-1}{%
\paragraph{\texorpdfstring{\texttt{unitree\_sdk2\_cpp.robot.as2.SportClient.front\_flip}}{unitree\_sdk2\_cpp.robot.as2.SportClient.front\_flip}}\label{unitree-sdk2-cpp-robot-as2-sportclient-front-flip-1}}

对应 C++ SDK 操作
\passthrough{\lstinline!FrontFlip()!}。上游头文件没有可直接生成的业务说明时，本参考只保证签名映射，精确语义需查目标型号协议。

\textbf{签名}

\begin{lstlisting}[language=Python]
def front_flip(self) -> int
\end{lstlisting}

\textbf{可用性与安全性}

\textbf{\passthrough{\lstinline!SIGNATURE\_ONLY!} /
\passthrough{\lstinline!MOTION\_COMMAND!}}：仅用于补全和静态检查。当前二进制没有该
Python 方法；不得按可执行运动接口使用。

\textbf{参数}

无显式参数。实例方法中的 \passthrough{\lstinline!self!} 由 Python
自动传入。

\textbf{返回值}

\begin{longtable}[]{@{}
  >{\raggedright\arraybackslash}p{(\columnwidth - 4\tabcolsep) * \real{0.18}}
  >{\raggedright\arraybackslash}p{(\columnwidth - 4\tabcolsep) * \real{0.20}}
  >{\raggedright\arraybackslash}p{(\columnwidth - 4\tabcolsep) * \real{0.62}}@{}}
\toprule
位置 & 类型 & 含义 \\
\midrule
\endhead
返回值 & \passthrough{\lstinline!int!} & SDK 状态码；Unitree 示例通常以
\passthrough{\lstinline!0!} 表示成功，非零值需按具体服务定义解释。 \\
\bottomrule
\end{longtable}

\textbf{对应 C++}

\begin{itemize}
\tightlist
\item
  类：\passthrough{\lstinline!unitree::robot::as2::SportClient!}
\item
  签名：\passthrough{\lstinline!FrontFlip()!}
\item
  绑定策略：\passthrough{\lstinline!DIRECT!}
\item
  声明位置：\passthrough{\lstinline!include/unitree/robot/as2/sport/sport\_client.hpp:229!}
\end{itemize}

\textbf{说明}

\begin{itemize}
\tightlist
\item
  示例是局部调用片段；\passthrough{\lstinline!obj!}
  和参数变量表示已经按上层业务校验并准备好的对象。
\item
  该条目可以被编辑器和类型检查器识别，但当前运行时可能在属性查找阶段直接失败。
\item
  该方法属于运动命令。方法名包含
  \passthrough{\lstinline!stop!}、\passthrough{\lstinline!damp!} 或
  \passthrough{\lstinline!zero!} 也不代表它是独立物理急停。
\end{itemize}

\textbf{用法}

\begin{lstlisting}[language=Python]
from typing import TYPE_CHECKING

if TYPE_CHECKING:
    def planned_front_flip(obj: SportClient) -> int:
        return obj.front_flip()
\end{lstlisting}

\begin{quote}
{[}!CAUTION{]} 上面的 \passthrough{\lstinline!TYPE\_CHECKING!}
示例只表达计划调用形态。该分支在普通 Python
运行时为假，不代表当前扩展能执行此接口。
\end{quote}

\hypertarget{unitree-sdk2-cpp-robot-as2-sportclient-back-flip-1}{%
\paragraph{\texorpdfstring{\texttt{unitree\_sdk2\_cpp.robot.as2.SportClient.back\_flip}}{unitree\_sdk2\_cpp.robot.as2.SportClient.back\_flip}}\label{unitree-sdk2-cpp-robot-as2-sportclient-back-flip-1}}

对应 C++ SDK 操作
\passthrough{\lstinline!BackFlip()!}。上游头文件没有可直接生成的业务说明时，本参考只保证签名映射，精确语义需查目标型号协议。

\textbf{签名}

\begin{lstlisting}[language=Python]
def back_flip(self) -> int
\end{lstlisting}

\textbf{可用性与安全性}

\textbf{\passthrough{\lstinline!SIGNATURE\_ONLY!} /
\passthrough{\lstinline!MOTION\_COMMAND!}}：仅用于补全和静态检查。当前二进制没有该
Python 方法；不得按可执行运动接口使用。

\textbf{参数}

无显式参数。实例方法中的 \passthrough{\lstinline!self!} 由 Python
自动传入。

\textbf{返回值}

\begin{longtable}[]{@{}
  >{\raggedright\arraybackslash}p{(\columnwidth - 4\tabcolsep) * \real{0.18}}
  >{\raggedright\arraybackslash}p{(\columnwidth - 4\tabcolsep) * \real{0.20}}
  >{\raggedright\arraybackslash}p{(\columnwidth - 4\tabcolsep) * \real{0.62}}@{}}
\toprule
位置 & 类型 & 含义 \\
\midrule
\endhead
返回值 & \passthrough{\lstinline!int!} & SDK 状态码；Unitree 示例通常以
\passthrough{\lstinline!0!} 表示成功，非零值需按具体服务定义解释。 \\
\bottomrule
\end{longtable}

\textbf{对应 C++}

\begin{itemize}
\tightlist
\item
  类：\passthrough{\lstinline!unitree::robot::as2::SportClient!}
\item
  签名：\passthrough{\lstinline!BackFlip()!}
\item
  绑定策略：\passthrough{\lstinline!DIRECT!}
\item
  声明位置：\passthrough{\lstinline!include/unitree/robot/as2/sport/sport\_client.hpp:235!}
\end{itemize}

\textbf{说明}

\begin{itemize}
\tightlist
\item
  示例是局部调用片段；\passthrough{\lstinline!obj!}
  和参数变量表示已经按上层业务校验并准备好的对象。
\item
  该条目可以被编辑器和类型检查器识别，但当前运行时可能在属性查找阶段直接失败。
\item
  该方法属于运动命令。方法名包含
  \passthrough{\lstinline!stop!}、\passthrough{\lstinline!damp!} 或
  \passthrough{\lstinline!zero!} 也不代表它是独立物理急停。
\end{itemize}

\textbf{用法}

\begin{lstlisting}[language=Python]
from typing import TYPE_CHECKING

if TYPE_CHECKING:
    def planned_back_flip(obj: SportClient) -> int:
        return obj.back_flip()
\end{lstlisting}

\begin{quote}
{[}!CAUTION{]} 上面的 \passthrough{\lstinline!TYPE\_CHECKING!}
示例只表达计划调用形态。该分支在普通 Python
运行时为假，不代表当前扩展能执行此接口。
\end{quote}

\hypertarget{unitree-sdk2-cpp-robot-as2-sportclient-greeting-1}{%
\paragraph{\texorpdfstring{\texttt{unitree\_sdk2\_cpp.robot.as2.SportClient.greeting}}{unitree\_sdk2\_cpp.robot.as2.SportClient.greeting}}\label{unitree-sdk2-cpp-robot-as2-sportclient-greeting-1}}

对应 C++ SDK 操作
\passthrough{\lstinline!Greeting()!}。上游头文件没有可直接生成的业务说明时，本参考只保证签名映射，精确语义需查目标型号协议。

\textbf{签名}

\begin{lstlisting}[language=Python]
def greeting(self) -> int
\end{lstlisting}

\textbf{可用性与安全性}

\textbf{\passthrough{\lstinline!SIGNATURE\_ONLY!} /
\passthrough{\lstinline!MOTION\_COMMAND!}}：仅用于补全和静态检查。当前二进制没有该
Python 方法；不得按可执行运动接口使用。

\textbf{参数}

无显式参数。实例方法中的 \passthrough{\lstinline!self!} 由 Python
自动传入。

\textbf{返回值}

\begin{longtable}[]{@{}
  >{\raggedright\arraybackslash}p{(\columnwidth - 4\tabcolsep) * \real{0.18}}
  >{\raggedright\arraybackslash}p{(\columnwidth - 4\tabcolsep) * \real{0.20}}
  >{\raggedright\arraybackslash}p{(\columnwidth - 4\tabcolsep) * \real{0.62}}@{}}
\toprule
位置 & 类型 & 含义 \\
\midrule
\endhead
返回值 & \passthrough{\lstinline!int!} & SDK 状态码；Unitree 示例通常以
\passthrough{\lstinline!0!} 表示成功，非零值需按具体服务定义解释。 \\
\bottomrule
\end{longtable}

\textbf{对应 C++}

\begin{itemize}
\tightlist
\item
  类：\passthrough{\lstinline!unitree::robot::as2::SportClient!}
\item
  签名：\passthrough{\lstinline!Greeting()!}
\item
  绑定策略：\passthrough{\lstinline!DIRECT!}
\item
  声明位置：\passthrough{\lstinline!include/unitree/robot/as2/sport/sport\_client.hpp:241!}
\end{itemize}

\textbf{说明}

\begin{itemize}
\tightlist
\item
  示例是局部调用片段；\passthrough{\lstinline!obj!}
  和参数变量表示已经按上层业务校验并准备好的对象。
\item
  该条目可以被编辑器和类型检查器识别，但当前运行时可能在属性查找阶段直接失败。
\item
  该方法属于运动命令。方法名包含
  \passthrough{\lstinline!stop!}、\passthrough{\lstinline!damp!} 或
  \passthrough{\lstinline!zero!} 也不代表它是独立物理急停。
\end{itemize}

\textbf{用法}

\begin{lstlisting}[language=Python]
from typing import TYPE_CHECKING

if TYPE_CHECKING:
    def planned_greeting(obj: SportClient) -> int:
        return obj.greeting()
\end{lstlisting}

\begin{quote}
{[}!CAUTION{]} 上面的 \passthrough{\lstinline!TYPE\_CHECKING!}
示例只表达计划调用形态。该分支在普通 Python
运行时为假，不代表当前扩展能执行此接口。
\end{quote}

\hypertarget{unitree-sdk2-cpp-robot-as2-sportclient-heart-1}{%
\paragraph{\texorpdfstring{\texttt{unitree\_sdk2\_cpp.robot.as2.SportClient.heart}}{unitree\_sdk2\_cpp.robot.as2.SportClient.heart}}\label{unitree-sdk2-cpp-robot-as2-sportclient-heart-1}}

对应 C++ SDK 操作
\passthrough{\lstinline!Heart()!}。上游头文件没有可直接生成的业务说明时，本参考只保证签名映射，精确语义需查目标型号协议。

\textbf{签名}

\begin{lstlisting}[language=Python]
def heart(self) -> int
\end{lstlisting}

\textbf{可用性与安全性}

\textbf{\passthrough{\lstinline!SIGNATURE\_ONLY!} /
\passthrough{\lstinline!MOTION\_COMMAND!}}：仅用于补全和静态检查。当前二进制没有该
Python 方法；不得按可执行运动接口使用。

\textbf{参数}

无显式参数。实例方法中的 \passthrough{\lstinline!self!} 由 Python
自动传入。

\textbf{返回值}

\begin{longtable}[]{@{}
  >{\raggedright\arraybackslash}p{(\columnwidth - 4\tabcolsep) * \real{0.18}}
  >{\raggedright\arraybackslash}p{(\columnwidth - 4\tabcolsep) * \real{0.20}}
  >{\raggedright\arraybackslash}p{(\columnwidth - 4\tabcolsep) * \real{0.62}}@{}}
\toprule
位置 & 类型 & 含义 \\
\midrule
\endhead
返回值 & \passthrough{\lstinline!int!} & SDK 状态码；Unitree 示例通常以
\passthrough{\lstinline!0!} 表示成功，非零值需按具体服务定义解释。 \\
\bottomrule
\end{longtable}

\textbf{对应 C++}

\begin{itemize}
\tightlist
\item
  类：\passthrough{\lstinline!unitree::robot::as2::SportClient!}
\item
  签名：\passthrough{\lstinline!Heart()!}
\item
  绑定策略：\passthrough{\lstinline!DIRECT!}
\item
  声明位置：\passthrough{\lstinline!include/unitree/robot/as2/sport/sport\_client.hpp:247!}
\end{itemize}

\textbf{说明}

\begin{itemize}
\tightlist
\item
  示例是局部调用片段；\passthrough{\lstinline!obj!}
  和参数变量表示已经按上层业务校验并准备好的对象。
\item
  该条目可以被编辑器和类型检查器识别，但当前运行时可能在属性查找阶段直接失败。
\item
  该方法属于运动命令。方法名包含
  \passthrough{\lstinline!stop!}、\passthrough{\lstinline!damp!} 或
  \passthrough{\lstinline!zero!} 也不代表它是独立物理急停。
\end{itemize}

\textbf{用法}

\begin{lstlisting}[language=Python]
from typing import TYPE_CHECKING

if TYPE_CHECKING:
    def planned_heart(obj: SportClient) -> int:
        return obj.heart()
\end{lstlisting}

\begin{quote}
{[}!CAUTION{]} 上面的 \passthrough{\lstinline!TYPE\_CHECKING!}
示例只表达计划调用形态。该分支在普通 Python
运行时为假，不代表当前扩展能执行此接口。
\end{quote}

\hypertarget{unitree-sdk2-cpp-robot-as2-sportclient-content-1}{%
\paragraph{\texorpdfstring{\texttt{unitree\_sdk2\_cpp.robot.as2.SportClient.content}}{unitree\_sdk2\_cpp.robot.as2.SportClient.content}}\label{unitree-sdk2-cpp-robot-as2-sportclient-content-1}}

对应 C++ SDK 操作
\passthrough{\lstinline!Content()!}。上游头文件没有可直接生成的业务说明时，本参考只保证签名映射，精确语义需查目标型号协议。

\textbf{签名}

\begin{lstlisting}[language=Python]
def content(self) -> int
\end{lstlisting}

\textbf{可用性与安全性}

\textbf{\passthrough{\lstinline!SIGNATURE\_ONLY!} /
\passthrough{\lstinline!MOTION\_COMMAND!}}：仅用于补全和静态检查。当前二进制没有该
Python 方法；不得按可执行运动接口使用。

\textbf{参数}

无显式参数。实例方法中的 \passthrough{\lstinline!self!} 由 Python
自动传入。

\textbf{返回值}

\begin{longtable}[]{@{}
  >{\raggedright\arraybackslash}p{(\columnwidth - 4\tabcolsep) * \real{0.18}}
  >{\raggedright\arraybackslash}p{(\columnwidth - 4\tabcolsep) * \real{0.20}}
  >{\raggedright\arraybackslash}p{(\columnwidth - 4\tabcolsep) * \real{0.62}}@{}}
\toprule
位置 & 类型 & 含义 \\
\midrule
\endhead
返回值 & \passthrough{\lstinline!int!} & SDK 状态码；Unitree 示例通常以
\passthrough{\lstinline!0!} 表示成功，非零值需按具体服务定义解释。 \\
\bottomrule
\end{longtable}

\textbf{对应 C++}

\begin{itemize}
\tightlist
\item
  类：\passthrough{\lstinline!unitree::robot::as2::SportClient!}
\item
  签名：\passthrough{\lstinline!Content()!}
\item
  绑定策略：\passthrough{\lstinline!DIRECT!}
\item
  声明位置：\passthrough{\lstinline!include/unitree/robot/as2/sport/sport\_client.hpp:253!}
\end{itemize}

\textbf{说明}

\begin{itemize}
\tightlist
\item
  示例是局部调用片段；\passthrough{\lstinline!obj!}
  和参数变量表示已经按上层业务校验并准备好的对象。
\item
  该条目可以被编辑器和类型检查器识别，但当前运行时可能在属性查找阶段直接失败。
\item
  该方法属于运动命令。方法名包含
  \passthrough{\lstinline!stop!}、\passthrough{\lstinline!damp!} 或
  \passthrough{\lstinline!zero!} 也不代表它是独立物理急停。
\end{itemize}

\textbf{用法}

\begin{lstlisting}[language=Python]
from typing import TYPE_CHECKING

if TYPE_CHECKING:
    def planned_content(obj: SportClient) -> int:
        return obj.content()
\end{lstlisting}

\begin{quote}
{[}!CAUTION{]} 上面的 \passthrough{\lstinline!TYPE\_CHECKING!}
示例只表达计划调用形态。该分支在普通 Python
运行时为假，不代表当前扩展能执行此接口。
\end{quote}

\hypertarget{unitree-sdk2-cpp-robot-as2-sportclient-dance1-1}{%
\paragraph{\texorpdfstring{\texttt{unitree\_sdk2\_cpp.robot.as2.SportClient.dance1}}{unitree\_sdk2\_cpp.robot.as2.SportClient.dance1}}\label{unitree-sdk2-cpp-robot-as2-sportclient-dance1-1}}

对应 C++ SDK 操作
\passthrough{\lstinline!Dance1()!}。上游头文件没有可直接生成的业务说明时，本参考只保证签名映射，精确语义需查目标型号协议。

\textbf{签名}

\begin{lstlisting}[language=Python]
def dance1(self) -> int
\end{lstlisting}

\textbf{可用性与安全性}

\textbf{\passthrough{\lstinline!SIGNATURE\_ONLY!} /
\passthrough{\lstinline!MOTION\_COMMAND!}}：仅用于补全和静态检查。当前二进制没有该
Python 方法；不得按可执行运动接口使用。

\textbf{参数}

无显式参数。实例方法中的 \passthrough{\lstinline!self!} 由 Python
自动传入。

\textbf{返回值}

\begin{longtable}[]{@{}
  >{\raggedright\arraybackslash}p{(\columnwidth - 4\tabcolsep) * \real{0.18}}
  >{\raggedright\arraybackslash}p{(\columnwidth - 4\tabcolsep) * \real{0.20}}
  >{\raggedright\arraybackslash}p{(\columnwidth - 4\tabcolsep) * \real{0.62}}@{}}
\toprule
位置 & 类型 & 含义 \\
\midrule
\endhead
返回值 & \passthrough{\lstinline!int!} & SDK 状态码；Unitree 示例通常以
\passthrough{\lstinline!0!} 表示成功，非零值需按具体服务定义解释。 \\
\bottomrule
\end{longtable}

\textbf{对应 C++}

\begin{itemize}
\tightlist
\item
  类：\passthrough{\lstinline!unitree::robot::as2::SportClient!}
\item
  签名：\passthrough{\lstinline!Dance1()!}
\item
  绑定策略：\passthrough{\lstinline!DIRECT!}
\item
  声明位置：\passthrough{\lstinline!include/unitree/robot/as2/sport/sport\_client.hpp:259!}
\end{itemize}

\textbf{说明}

\begin{itemize}
\tightlist
\item
  示例是局部调用片段；\passthrough{\lstinline!obj!}
  和参数变量表示已经按上层业务校验并准备好的对象。
\item
  该条目可以被编辑器和类型检查器识别，但当前运行时可能在属性查找阶段直接失败。
\item
  该方法属于运动命令。方法名包含
  \passthrough{\lstinline!stop!}、\passthrough{\lstinline!damp!} 或
  \passthrough{\lstinline!zero!} 也不代表它是独立物理急停。
\end{itemize}

\textbf{用法}

\begin{lstlisting}[language=Python]
from typing import TYPE_CHECKING

if TYPE_CHECKING:
    def planned_dance1(obj: SportClient) -> int:
        return obj.dance1()
\end{lstlisting}

\begin{quote}
{[}!CAUTION{]} 上面的 \passthrough{\lstinline!TYPE\_CHECKING!}
示例只表达计划调用形态。该分支在普通 Python
运行时为假，不代表当前扩展能执行此接口。
\end{quote}

\hypertarget{unitree-sdk2-cpp-robot-as2-sportclient-dance2-1}{%
\paragraph{\texorpdfstring{\texttt{unitree\_sdk2\_cpp.robot.as2.SportClient.dance2}}{unitree\_sdk2\_cpp.robot.as2.SportClient.dance2}}\label{unitree-sdk2-cpp-robot-as2-sportclient-dance2-1}}

对应 C++ SDK 操作
\passthrough{\lstinline!Dance2()!}。上游头文件没有可直接生成的业务说明时，本参考只保证签名映射，精确语义需查目标型号协议。

\textbf{签名}

\begin{lstlisting}[language=Python]
def dance2(self) -> int
\end{lstlisting}

\textbf{可用性与安全性}

\textbf{\passthrough{\lstinline!SIGNATURE\_ONLY!} /
\passthrough{\lstinline!MOTION\_COMMAND!}}：仅用于补全和静态检查。当前二进制没有该
Python 方法；不得按可执行运动接口使用。

\textbf{参数}

无显式参数。实例方法中的 \passthrough{\lstinline!self!} 由 Python
自动传入。

\textbf{返回值}

\begin{longtable}[]{@{}
  >{\raggedright\arraybackslash}p{(\columnwidth - 4\tabcolsep) * \real{0.18}}
  >{\raggedright\arraybackslash}p{(\columnwidth - 4\tabcolsep) * \real{0.20}}
  >{\raggedright\arraybackslash}p{(\columnwidth - 4\tabcolsep) * \real{0.62}}@{}}
\toprule
位置 & 类型 & 含义 \\
\midrule
\endhead
返回值 & \passthrough{\lstinline!int!} & SDK 状态码；Unitree 示例通常以
\passthrough{\lstinline!0!} 表示成功，非零值需按具体服务定义解释。 \\
\bottomrule
\end{longtable}

\textbf{对应 C++}

\begin{itemize}
\tightlist
\item
  类：\passthrough{\lstinline!unitree::robot::as2::SportClient!}
\item
  签名：\passthrough{\lstinline!Dance2()!}
\item
  绑定策略：\passthrough{\lstinline!DIRECT!}
\item
  声明位置：\passthrough{\lstinline!include/unitree/robot/as2/sport/sport\_client.hpp:265!}
\end{itemize}

\textbf{说明}

\begin{itemize}
\tightlist
\item
  示例是局部调用片段；\passthrough{\lstinline!obj!}
  和参数变量表示已经按上层业务校验并准备好的对象。
\item
  该条目可以被编辑器和类型检查器识别，但当前运行时可能在属性查找阶段直接失败。
\item
  该方法属于运动命令。方法名包含
  \passthrough{\lstinline!stop!}、\passthrough{\lstinline!damp!} 或
  \passthrough{\lstinline!zero!} 也不代表它是独立物理急停。
\end{itemize}

\textbf{用法}

\begin{lstlisting}[language=Python]
from typing import TYPE_CHECKING

if TYPE_CHECKING:
    def planned_dance2(obj: SportClient) -> int:
        return obj.dance2()
\end{lstlisting}

\begin{quote}
{[}!CAUTION{]} 上面的 \passthrough{\lstinline!TYPE\_CHECKING!}
示例只表达计划调用形态。该分支在普通 Python
运行时为假，不代表当前扩展能执行此接口。
\end{quote}

\hypertarget{unitree-sdk2-cpp-robot-as2-sportclient-handshake-1}{%
\paragraph{\texorpdfstring{\texttt{unitree\_sdk2\_cpp.robot.as2.SportClient.handshake}}{unitree\_sdk2\_cpp.robot.as2.SportClient.handshake}}\label{unitree-sdk2-cpp-robot-as2-sportclient-handshake-1}}

对应 C++ SDK 操作
\passthrough{\lstinline!Handshake()!}。上游头文件没有可直接生成的业务说明时，本参考只保证签名映射，精确语义需查目标型号协议。

\textbf{签名}

\begin{lstlisting}[language=Python]
def handshake(self) -> int
\end{lstlisting}

\textbf{可用性与安全性}

\textbf{\passthrough{\lstinline!SIGNATURE\_ONLY!} /
\passthrough{\lstinline!MOTION\_COMMAND!}}：仅用于补全和静态检查。当前二进制没有该
Python 方法；不得按可执行运动接口使用。

\textbf{参数}

无显式参数。实例方法中的 \passthrough{\lstinline!self!} 由 Python
自动传入。

\textbf{返回值}

\begin{longtable}[]{@{}
  >{\raggedright\arraybackslash}p{(\columnwidth - 4\tabcolsep) * \real{0.18}}
  >{\raggedright\arraybackslash}p{(\columnwidth - 4\tabcolsep) * \real{0.20}}
  >{\raggedright\arraybackslash}p{(\columnwidth - 4\tabcolsep) * \real{0.62}}@{}}
\toprule
位置 & 类型 & 含义 \\
\midrule
\endhead
返回值 & \passthrough{\lstinline!int!} & SDK 状态码；Unitree 示例通常以
\passthrough{\lstinline!0!} 表示成功，非零值需按具体服务定义解释。 \\
\bottomrule
\end{longtable}

\textbf{对应 C++}

\begin{itemize}
\tightlist
\item
  类：\passthrough{\lstinline!unitree::robot::as2::SportClient!}
\item
  签名：\passthrough{\lstinline!Handshake()!}
\item
  绑定策略：\passthrough{\lstinline!DIRECT!}
\item
  声明位置：\passthrough{\lstinline!include/unitree/robot/as2/sport/sport\_client.hpp:271!}
\end{itemize}

\textbf{说明}

\begin{itemize}
\tightlist
\item
  示例是局部调用片段；\passthrough{\lstinline!obj!}
  和参数变量表示已经按上层业务校验并准备好的对象。
\item
  该条目可以被编辑器和类型检查器识别，但当前运行时可能在属性查找阶段直接失败。
\item
  该方法属于运动命令。方法名包含
  \passthrough{\lstinline!stop!}、\passthrough{\lstinline!damp!} 或
  \passthrough{\lstinline!zero!} 也不代表它是独立物理急停。
\end{itemize}

\textbf{用法}

\begin{lstlisting}[language=Python]
from typing import TYPE_CHECKING

if TYPE_CHECKING:
    def planned_handshake(obj: SportClient) -> int:
        return obj.handshake()
\end{lstlisting}

\begin{quote}
{[}!CAUTION{]} 上面的 \passthrough{\lstinline!TYPE\_CHECKING!}
示例只表达计划调用形态。该分支在普通 Python
运行时为假，不代表当前扩展能执行此接口。
\end{quote}

\hypertarget{unitree-sdk2-cpp-robot-as2-sportclient-stretch-1}{%
\paragraph{\texorpdfstring{\texttt{unitree\_sdk2\_cpp.robot.as2.SportClient.stretch}}{unitree\_sdk2\_cpp.robot.as2.SportClient.stretch}}\label{unitree-sdk2-cpp-robot-as2-sportclient-stretch-1}}

对应 C++ SDK 操作
\passthrough{\lstinline!Stretch()!}。上游头文件没有可直接生成的业务说明时，本参考只保证签名映射，精确语义需查目标型号协议。

\textbf{签名}

\begin{lstlisting}[language=Python]
def stretch(self) -> int
\end{lstlisting}

\textbf{可用性与安全性}

\textbf{\passthrough{\lstinline!SIGNATURE\_ONLY!} /
\passthrough{\lstinline!MOTION\_COMMAND!}}：仅用于补全和静态检查。当前二进制没有该
Python 方法；不得按可执行运动接口使用。

\textbf{参数}

无显式参数。实例方法中的 \passthrough{\lstinline!self!} 由 Python
自动传入。

\textbf{返回值}

\begin{longtable}[]{@{}
  >{\raggedright\arraybackslash}p{(\columnwidth - 4\tabcolsep) * \real{0.18}}
  >{\raggedright\arraybackslash}p{(\columnwidth - 4\tabcolsep) * \real{0.20}}
  >{\raggedright\arraybackslash}p{(\columnwidth - 4\tabcolsep) * \real{0.62}}@{}}
\toprule
位置 & 类型 & 含义 \\
\midrule
\endhead
返回值 & \passthrough{\lstinline!int!} & SDK 状态码；Unitree 示例通常以
\passthrough{\lstinline!0!} 表示成功，非零值需按具体服务定义解释。 \\
\bottomrule
\end{longtable}

\textbf{对应 C++}

\begin{itemize}
\tightlist
\item
  类：\passthrough{\lstinline!unitree::robot::as2::SportClient!}
\item
  签名：\passthrough{\lstinline!Stretch()!}
\item
  绑定策略：\passthrough{\lstinline!DIRECT!}
\item
  声明位置：\passthrough{\lstinline!include/unitree/robot/as2/sport/sport\_client.hpp:277!}
\end{itemize}

\textbf{说明}

\begin{itemize}
\tightlist
\item
  示例是局部调用片段；\passthrough{\lstinline!obj!}
  和参数变量表示已经按上层业务校验并准备好的对象。
\item
  该条目可以被编辑器和类型检查器识别，但当前运行时可能在属性查找阶段直接失败。
\item
  该方法属于运动命令。方法名包含
  \passthrough{\lstinline!stop!}、\passthrough{\lstinline!damp!} 或
  \passthrough{\lstinline!zero!} 也不代表它是独立物理急停。
\end{itemize}

\textbf{用法}

\begin{lstlisting}[language=Python]
from typing import TYPE_CHECKING

if TYPE_CHECKING:
    def planned_stretch(obj: SportClient) -> int:
        return obj.stretch()
\end{lstlisting}

\begin{quote}
{[}!CAUTION{]} 上面的 \passthrough{\lstinline!TYPE\_CHECKING!}
示例只表达计划调用形态。该分支在普通 Python
运行时为假，不代表当前扩展能执行此接口。
\end{quote}

\hypertarget{unitree-sdk2-cpp-robot-as2-sportclient-sit-1}{%
\paragraph{\texorpdfstring{\texttt{unitree\_sdk2\_cpp.robot.as2.SportClient.sit}}{unitree\_sdk2\_cpp.robot.as2.SportClient.sit}}\label{unitree-sdk2-cpp-robot-as2-sportclient-sit-1}}

对应 C++ SDK 操作
\passthrough{\lstinline!Sit(int)!}。上游头文件没有可直接生成的业务说明时，本参考只保证签名映射，精确语义需查目标型号协议。

\textbf{签名}

\begin{lstlisting}[language=Python]
def sit(self, enter: int) -> int
\end{lstlisting}

\textbf{可用性与安全性}

\textbf{\passthrough{\lstinline!SIGNATURE\_ONLY!} /
\passthrough{\lstinline!MOTION\_COMMAND!}}：仅用于补全和静态检查。当前二进制没有该
Python 方法；不得按可执行运动接口使用。

\textbf{参数}

\begin{longtable}[]{@{}
  >{\raggedright\arraybackslash}p{(\columnwidth - 6\tabcolsep) * \real{0.18}}
  >{\raggedright\arraybackslash}p{(\columnwidth - 6\tabcolsep) * \real{0.24}}
  >{\raggedright\arraybackslash}p{(\columnwidth - 6\tabcolsep) * \real{0.18}}
  >{\raggedright\arraybackslash}p{(\columnwidth - 6\tabcolsep) * \real{0.40}}@{}}
\toprule
名称 & Python 类型 & 默认值 & 含义 \\
\midrule
\endhead
\passthrough{\lstinline!enter!} & \passthrough{\lstinline!int!} & 必填 &
传给该接口的 \passthrough{\lstinline!enter!} 参数，Python 类型为
\passthrough{\lstinline!int!}。精确含义、单位、范围和枚举值需查目标型号对应头文件与协议。
对应 C++ 参数 \passthrough{\lstinline!enter: int!}。 \\
\bottomrule
\end{longtable}

\textbf{返回值}

\begin{longtable}[]{@{}
  >{\raggedright\arraybackslash}p{(\columnwidth - 4\tabcolsep) * \real{0.18}}
  >{\raggedright\arraybackslash}p{(\columnwidth - 4\tabcolsep) * \real{0.20}}
  >{\raggedright\arraybackslash}p{(\columnwidth - 4\tabcolsep) * \real{0.62}}@{}}
\toprule
位置 & 类型 & 含义 \\
\midrule
\endhead
返回值 & \passthrough{\lstinline!int!} & SDK 状态码；Unitree 示例通常以
\passthrough{\lstinline!0!} 表示成功，非零值需按具体服务定义解释。 \\
\bottomrule
\end{longtable}

\textbf{对应 C++}

\begin{itemize}
\tightlist
\item
  类：\passthrough{\lstinline!unitree::robot::as2::SportClient!}
\item
  签名：\passthrough{\lstinline!Sit(int)!}
\item
  绑定策略：\passthrough{\lstinline!DIRECT!}
\item
  声明位置：\passthrough{\lstinline!include/unitree/robot/as2/sport/sport\_client.hpp:283!}
\end{itemize}

\textbf{说明}

\begin{itemize}
\tightlist
\item
  示例是局部调用片段；\passthrough{\lstinline!obj!}
  和参数变量表示已经按上层业务校验并准备好的对象。
\item
  该条目可以被编辑器和类型检查器识别，但当前运行时可能在属性查找阶段直接失败。
\item
  该方法属于运动命令。方法名包含
  \passthrough{\lstinline!stop!}、\passthrough{\lstinline!damp!} 或
  \passthrough{\lstinline!zero!} 也不代表它是独立物理急停。
\end{itemize}

\textbf{用法}

\begin{lstlisting}[language=Python]
from typing import TYPE_CHECKING

if TYPE_CHECKING:
    def planned_sit(obj: SportClient, enter: int) -> int:
        return obj.sit(enter=enter)
\end{lstlisting}

\begin{quote}
{[}!CAUTION{]} 上面的 \passthrough{\lstinline!TYPE\_CHECKING!}
示例只表达计划调用形态。该分支在普通 Python
运行时为假，不代表当前扩展能执行此接口。
\end{quote}

\hypertarget{unitree-sdk2-cpp-robot-as2-sportclient-front-jump-1}{%
\paragraph{\texorpdfstring{\texttt{unitree\_sdk2\_cpp.robot.as2.SportClient.front\_jump}}{unitree\_sdk2\_cpp.robot.as2.SportClient.front\_jump}}\label{unitree-sdk2-cpp-robot-as2-sportclient-front-jump-1}}

对应 C++ SDK 操作
\passthrough{\lstinline!FrontJump()!}。上游头文件没有可直接生成的业务说明时，本参考只保证签名映射，精确语义需查目标型号协议。

\textbf{签名}

\begin{lstlisting}[language=Python]
def front_jump(self) -> int
\end{lstlisting}

\textbf{可用性与安全性}

\textbf{\passthrough{\lstinline!SIGNATURE\_ONLY!} /
\passthrough{\lstinline!MOTION\_COMMAND!}}：仅用于补全和静态检查。当前二进制没有该
Python 方法；不得按可执行运动接口使用。

\textbf{参数}

无显式参数。实例方法中的 \passthrough{\lstinline!self!} 由 Python
自动传入。

\textbf{返回值}

\begin{longtable}[]{@{}
  >{\raggedright\arraybackslash}p{(\columnwidth - 4\tabcolsep) * \real{0.18}}
  >{\raggedright\arraybackslash}p{(\columnwidth - 4\tabcolsep) * \real{0.20}}
  >{\raggedright\arraybackslash}p{(\columnwidth - 4\tabcolsep) * \real{0.62}}@{}}
\toprule
位置 & 类型 & 含义 \\
\midrule
\endhead
返回值 & \passthrough{\lstinline!int!} & SDK 状态码；Unitree 示例通常以
\passthrough{\lstinline!0!} 表示成功，非零值需按具体服务定义解释。 \\
\bottomrule
\end{longtable}

\textbf{对应 C++}

\begin{itemize}
\tightlist
\item
  类：\passthrough{\lstinline!unitree::robot::as2::SportClient!}
\item
  签名：\passthrough{\lstinline!FrontJump()!}
\item
  绑定策略：\passthrough{\lstinline!DIRECT!}
\item
  声明位置：\passthrough{\lstinline!include/unitree/robot/as2/sport/sport\_client.hpp:292!}
\end{itemize}

\textbf{说明}

\begin{itemize}
\tightlist
\item
  示例是局部调用片段；\passthrough{\lstinline!obj!}
  和参数变量表示已经按上层业务校验并准备好的对象。
\item
  该条目可以被编辑器和类型检查器识别，但当前运行时可能在属性查找阶段直接失败。
\item
  该方法属于运动命令。方法名包含
  \passthrough{\lstinline!stop!}、\passthrough{\lstinline!damp!} 或
  \passthrough{\lstinline!zero!} 也不代表它是独立物理急停。
\end{itemize}

\textbf{用法}

\begin{lstlisting}[language=Python]
from typing import TYPE_CHECKING

if TYPE_CHECKING:
    def planned_front_jump(obj: SportClient) -> int:
        return obj.front_jump()
\end{lstlisting}

\begin{quote}
{[}!CAUTION{]} 上面的 \passthrough{\lstinline!TYPE\_CHECKING!}
示例只表达计划调用形态。该分支在普通 Python
运行时为假，不代表当前扩展能执行此接口。
\end{quote}

\hypertarget{unitree-sdk2-cpp-robot-as2-sportclient-push-up-1}{%
\paragraph{\texorpdfstring{\texttt{unitree\_sdk2\_cpp.robot.as2.SportClient.push\_up}}{unitree\_sdk2\_cpp.robot.as2.SportClient.push\_up}}\label{unitree-sdk2-cpp-robot-as2-sportclient-push-up-1}}

对应 C++ SDK 操作
\passthrough{\lstinline!PushUp()!}。上游头文件没有可直接生成的业务说明时，本参考只保证签名映射，精确语义需查目标型号协议。

\textbf{签名}

\begin{lstlisting}[language=Python]
def push_up(self) -> int
\end{lstlisting}

\textbf{可用性与安全性}

\textbf{\passthrough{\lstinline!SIGNATURE\_ONLY!} /
\passthrough{\lstinline!MOTION\_COMMAND!}}：仅用于补全和静态检查。当前二进制没有该
Python 方法；不得按可执行运动接口使用。

\textbf{参数}

无显式参数。实例方法中的 \passthrough{\lstinline!self!} 由 Python
自动传入。

\textbf{返回值}

\begin{longtable}[]{@{}
  >{\raggedright\arraybackslash}p{(\columnwidth - 4\tabcolsep) * \real{0.18}}
  >{\raggedright\arraybackslash}p{(\columnwidth - 4\tabcolsep) * \real{0.20}}
  >{\raggedright\arraybackslash}p{(\columnwidth - 4\tabcolsep) * \real{0.62}}@{}}
\toprule
位置 & 类型 & 含义 \\
\midrule
\endhead
返回值 & \passthrough{\lstinline!int!} & SDK 状态码；Unitree 示例通常以
\passthrough{\lstinline!0!} 表示成功，非零值需按具体服务定义解释。 \\
\bottomrule
\end{longtable}

\textbf{对应 C++}

\begin{itemize}
\tightlist
\item
  类：\passthrough{\lstinline!unitree::robot::as2::SportClient!}
\item
  签名：\passthrough{\lstinline!PushUp()!}
\item
  绑定策略：\passthrough{\lstinline!DIRECT!}
\item
  声明位置：\passthrough{\lstinline!include/unitree/robot/as2/sport/sport\_client.hpp:298!}
\end{itemize}

\textbf{说明}

\begin{itemize}
\tightlist
\item
  示例是局部调用片段；\passthrough{\lstinline!obj!}
  和参数变量表示已经按上层业务校验并准备好的对象。
\item
  该条目可以被编辑器和类型检查器识别，但当前运行时可能在属性查找阶段直接失败。
\item
  该方法属于运动命令。方法名包含
  \passthrough{\lstinline!stop!}、\passthrough{\lstinline!damp!} 或
  \passthrough{\lstinline!zero!} 也不代表它是独立物理急停。
\end{itemize}

\textbf{用法}

\begin{lstlisting}[language=Python]
from typing import TYPE_CHECKING

if TYPE_CHECKING:
    def planned_push_up(obj: SportClient) -> int:
        return obj.push_up()
\end{lstlisting}

\begin{quote}
{[}!CAUTION{]} 上面的 \passthrough{\lstinline!TYPE\_CHECKING!}
示例只表达计划调用形态。该分支在普通 Python
运行时为假，不代表当前扩展能执行此接口。
\end{quote}

\hypertarget{unitree-sdk2-cpp-robot-as2-sportclient-up-jump-1}{%
\paragraph{\texorpdfstring{\texttt{unitree\_sdk2\_cpp.robot.as2.SportClient.up\_jump}}{unitree\_sdk2\_cpp.robot.as2.SportClient.up\_jump}}\label{unitree-sdk2-cpp-robot-as2-sportclient-up-jump-1}}

对应 C++ SDK 操作
\passthrough{\lstinline!UpJump()!}。上游头文件没有可直接生成的业务说明时，本参考只保证签名映射，精确语义需查目标型号协议。

\textbf{签名}

\begin{lstlisting}[language=Python]
def up_jump(self) -> int
\end{lstlisting}

\textbf{可用性与安全性}

\textbf{\passthrough{\lstinline!SIGNATURE\_ONLY!} /
\passthrough{\lstinline!MOTION\_COMMAND!}}：仅用于补全和静态检查。当前二进制没有该
Python 方法；不得按可执行运动接口使用。

\textbf{参数}

无显式参数。实例方法中的 \passthrough{\lstinline!self!} 由 Python
自动传入。

\textbf{返回值}

\begin{longtable}[]{@{}
  >{\raggedright\arraybackslash}p{(\columnwidth - 4\tabcolsep) * \real{0.18}}
  >{\raggedright\arraybackslash}p{(\columnwidth - 4\tabcolsep) * \real{0.20}}
  >{\raggedright\arraybackslash}p{(\columnwidth - 4\tabcolsep) * \real{0.62}}@{}}
\toprule
位置 & 类型 & 含义 \\
\midrule
\endhead
返回值 & \passthrough{\lstinline!int!} & SDK 状态码；Unitree 示例通常以
\passthrough{\lstinline!0!} 表示成功，非零值需按具体服务定义解释。 \\
\bottomrule
\end{longtable}

\textbf{对应 C++}

\begin{itemize}
\tightlist
\item
  类：\passthrough{\lstinline!unitree::robot::as2::SportClient!}
\item
  签名：\passthrough{\lstinline!UpJump()!}
\item
  绑定策略：\passthrough{\lstinline!DIRECT!}
\item
  声明位置：\passthrough{\lstinline!include/unitree/robot/as2/sport/sport\_client.hpp:304!}
\end{itemize}

\textbf{说明}

\begin{itemize}
\tightlist
\item
  示例是局部调用片段；\passthrough{\lstinline!obj!}
  和参数变量表示已经按上层业务校验并准备好的对象。
\item
  该条目可以被编辑器和类型检查器识别，但当前运行时可能在属性查找阶段直接失败。
\item
  该方法属于运动命令。方法名包含
  \passthrough{\lstinline!stop!}、\passthrough{\lstinline!damp!} 或
  \passthrough{\lstinline!zero!} 也不代表它是独立物理急停。
\end{itemize}

\textbf{用法}

\begin{lstlisting}[language=Python]
from typing import TYPE_CHECKING

if TYPE_CHECKING:
    def planned_up_jump(obj: SportClient) -> int:
        return obj.up_jump()
\end{lstlisting}

\begin{quote}
{[}!CAUTION{]} 上面的 \passthrough{\lstinline!TYPE\_CHECKING!}
示例只表达计划调用形态。该分支在普通 Python
运行时为假，不代表当前扩展能执行此接口。
\end{quote}

\hypertarget{unitree-sdk2-cpp-robot-as2-sportclient-set-auto-recovery-1}{%
\paragraph{\texorpdfstring{\texttt{unitree\_sdk2\_cpp.robot.as2.SportClient.set\_auto\_recovery}}{unitree\_sdk2\_cpp.robot.as2.SportClient.set\_auto\_recovery}}\label{unitree-sdk2-cpp-robot-as2-sportclient-set-auto-recovery-1}}

设置 \passthrough{\lstinline!auto!}
\passthrough{\lstinline!recovery!}。具体副作用和安全边界见下方状态。

\textbf{签名}

\begin{lstlisting}[language=Python]
def set_auto_recovery(self, switch_on: int) -> int
\end{lstlisting}

\textbf{可用性与安全性}

\textbf{\passthrough{\lstinline!SIGNATURE\_ONLY!} /
\passthrough{\lstinline!MOTION\_COMMAND!}}：仅用于补全和静态检查。当前二进制没有该
Python 方法；不得按可执行运动接口使用。

\textbf{参数}

\begin{longtable}[]{@{}
  >{\raggedright\arraybackslash}p{(\columnwidth - 6\tabcolsep) * \real{0.18}}
  >{\raggedright\arraybackslash}p{(\columnwidth - 6\tabcolsep) * \real{0.24}}
  >{\raggedright\arraybackslash}p{(\columnwidth - 6\tabcolsep) * \real{0.18}}
  >{\raggedright\arraybackslash}p{(\columnwidth - 6\tabcolsep) * \real{0.40}}@{}}
\toprule
名称 & Python 类型 & 默认值 & 含义 \\
\midrule
\endhead
\passthrough{\lstinline!switch\_on!} & \passthrough{\lstinline!int!} &
必填 & 传给该接口的 开关 \passthrough{\lstinline!on!} 参数，Python
类型为
\passthrough{\lstinline!int!}。精确含义、单位、范围和枚举值需查目标型号对应头文件与协议。
对应 C++ 参数 \passthrough{\lstinline!switch\_on: int!}。 \\
\bottomrule
\end{longtable}

\textbf{返回值}

\begin{longtable}[]{@{}
  >{\raggedright\arraybackslash}p{(\columnwidth - 4\tabcolsep) * \real{0.18}}
  >{\raggedright\arraybackslash}p{(\columnwidth - 4\tabcolsep) * \real{0.20}}
  >{\raggedright\arraybackslash}p{(\columnwidth - 4\tabcolsep) * \real{0.62}}@{}}
\toprule
位置 & 类型 & 含义 \\
\midrule
\endhead
返回值 & \passthrough{\lstinline!int!} & SDK 状态码；Unitree 示例通常以
\passthrough{\lstinline!0!} 表示成功，非零值需按具体服务定义解释。 \\
\bottomrule
\end{longtable}

\textbf{对应 C++}

\begin{itemize}
\tightlist
\item
  类：\passthrough{\lstinline!unitree::robot::as2::SportClient!}
\item
  签名：\passthrough{\lstinline!SetAutoRecovery(int)!}
\item
  绑定策略：\passthrough{\lstinline!DIRECT!}
\item
  声明位置：\passthrough{\lstinline!include/unitree/robot/as2/sport/sport\_client.hpp:310!}
\end{itemize}

\textbf{说明}

\begin{itemize}
\tightlist
\item
  示例是局部调用片段；\passthrough{\lstinline!obj!}
  和参数变量表示已经按上层业务校验并准备好的对象。
\item
  该条目可以被编辑器和类型检查器识别，但当前运行时可能在属性查找阶段直接失败。
\item
  该方法属于运动命令。方法名包含
  \passthrough{\lstinline!stop!}、\passthrough{\lstinline!damp!} 或
  \passthrough{\lstinline!zero!} 也不代表它是独立物理急停。
\end{itemize}

\textbf{用法}

\begin{lstlisting}[language=Python]
from typing import TYPE_CHECKING

if TYPE_CHECKING:
    def planned_set_auto_recovery(obj: SportClient, switch_on: int) -> int:
        return obj.set_auto_recovery(switch_on=switch_on)
\end{lstlisting}

\begin{quote}
{[}!CAUTION{]} 上面的 \passthrough{\lstinline!TYPE\_CHECKING!}
示例只表达计划调用形态。该分支在普通 Python
运行时为假，不代表当前扩展能执行此接口。
\end{quote}

\hypertarget{unitree-sdk2-cpp-robot-as2-sportclient-switch-joystick-1}{%
\paragraph{\texorpdfstring{\texttt{unitree\_sdk2\_cpp.robot.as2.SportClient.switch\_joystick}}{unitree\_sdk2\_cpp.robot.as2.SportClient.switch\_joystick}}\label{unitree-sdk2-cpp-robot-as2-sportclient-switch-joystick-1}}

对应 C++ SDK 操作
\passthrough{\lstinline!SwitchJoystick(int)!}。上游头文件没有可直接生成的业务说明时，本参考只保证签名映射，精确语义需查目标型号协议。

\textbf{签名}

\begin{lstlisting}[language=Python]
def switch_joystick(self, switch_on: int) -> int
\end{lstlisting}

\textbf{可用性与安全性}

\textbf{\passthrough{\lstinline!SIGNATURE\_ONLY!} /
\passthrough{\lstinline!MOTION\_COMMAND!}}：仅用于补全和静态检查。当前二进制没有该
Python 方法；不得按可执行运动接口使用。

\textbf{参数}

\begin{longtable}[]{@{}
  >{\raggedright\arraybackslash}p{(\columnwidth - 6\tabcolsep) * \real{0.18}}
  >{\raggedright\arraybackslash}p{(\columnwidth - 6\tabcolsep) * \real{0.24}}
  >{\raggedright\arraybackslash}p{(\columnwidth - 6\tabcolsep) * \real{0.18}}
  >{\raggedright\arraybackslash}p{(\columnwidth - 6\tabcolsep) * \real{0.40}}@{}}
\toprule
名称 & Python 类型 & 默认值 & 含义 \\
\midrule
\endhead
\passthrough{\lstinline!switch\_on!} & \passthrough{\lstinline!int!} &
必填 & 传给该接口的 开关 \passthrough{\lstinline!on!} 参数，Python
类型为
\passthrough{\lstinline!int!}。精确含义、单位、范围和枚举值需查目标型号对应头文件与协议。
对应 C++ 参数 \passthrough{\lstinline!switch\_on: int!}。 \\
\bottomrule
\end{longtable}

\textbf{返回值}

\begin{longtable}[]{@{}
  >{\raggedright\arraybackslash}p{(\columnwidth - 4\tabcolsep) * \real{0.18}}
  >{\raggedright\arraybackslash}p{(\columnwidth - 4\tabcolsep) * \real{0.20}}
  >{\raggedright\arraybackslash}p{(\columnwidth - 4\tabcolsep) * \real{0.62}}@{}}
\toprule
位置 & 类型 & 含义 \\
\midrule
\endhead
返回值 & \passthrough{\lstinline!int!} & SDK 状态码；Unitree 示例通常以
\passthrough{\lstinline!0!} 表示成功，非零值需按具体服务定义解释。 \\
\bottomrule
\end{longtable}

\textbf{对应 C++}

\begin{itemize}
\tightlist
\item
  类：\passthrough{\lstinline!unitree::robot::as2::SportClient!}
\item
  签名：\passthrough{\lstinline!SwitchJoystick(int)!}
\item
  绑定策略：\passthrough{\lstinline!DIRECT!}
\item
  声明位置：\passthrough{\lstinline!include/unitree/robot/as2/sport/sport\_client.hpp:319!}
\end{itemize}

\textbf{说明}

\begin{itemize}
\tightlist
\item
  示例是局部调用片段；\passthrough{\lstinline!obj!}
  和参数变量表示已经按上层业务校验并准备好的对象。
\item
  该条目可以被编辑器和类型检查器识别，但当前运行时可能在属性查找阶段直接失败。
\item
  该方法属于运动命令。方法名包含
  \passthrough{\lstinline!stop!}、\passthrough{\lstinline!damp!} 或
  \passthrough{\lstinline!zero!} 也不代表它是独立物理急停。
\end{itemize}

\textbf{用法}

\begin{lstlisting}[language=Python]
from typing import TYPE_CHECKING

if TYPE_CHECKING:
    def planned_switch_joystick(obj: SportClient, switch_on: int) -> int:
        return obj.switch_joystick(switch_on=switch_on)
\end{lstlisting}

\begin{quote}
{[}!CAUTION{]} 上面的 \passthrough{\lstinline!TYPE\_CHECKING!}
示例只表达计划调用形态。该分支在普通 Python
运行时为假，不代表当前扩展能执行此接口。
\end{quote}

\hypertarget{unitree-sdk2-cpp-robot-as2-sportclient-get-state-1}{%
\paragraph{\texorpdfstring{\texttt{unitree\_sdk2\_cpp.robot.as2.SportClient.get\_state}}{unitree\_sdk2\_cpp.robot.as2.SportClient.get\_state}}\label{unitree-sdk2-cpp-robot-as2-sportclient-get-state-1}}

查询或检查 状态。

\textbf{签名}

\begin{lstlisting}[language=Python]
def get_state(self) -> tuple[int, dict[str, str]]
\end{lstlisting}

\textbf{可用性与安全性}

\textbf{\passthrough{\lstinline!AVAILABLE!} /
\passthrough{\lstinline!READ\_ONLY!}}：当前绑定源码已实现只读查询。Robot
Client 的服务端查询通常仍需匹配的 Linux
扩展、DDS、网络、机器人和服务；纯本地 getter 则不一定需要全部条件。

\textbf{参数}

无显式参数。实例方法中的 \passthrough{\lstinline!self!} 由 Python
自动传入。

\textbf{返回值}

\begin{longtable}[]{@{}
  >{\raggedright\arraybackslash}p{(\columnwidth - 4\tabcolsep) * \real{0.18}}
  >{\raggedright\arraybackslash}p{(\columnwidth - 4\tabcolsep) * \real{0.20}}
  >{\raggedright\arraybackslash}p{(\columnwidth - 4\tabcolsep) * \real{0.62}}@{}}
\toprule
位置 & 类型 & 含义 \\
\midrule
\endhead
{[}0{]} \passthrough{\lstinline!status!} & \passthrough{\lstinline!int!}
& SDK 状态码；Unitree 示例通常以 \passthrough{\lstinline!0!}
表示成功，非零值需按具体服务错误码解释。 \\
{[}1{]} \passthrough{\lstinline!state\_map!} &
\passthrough{\lstinline!dict[str, str]!} & 传给该接口的 状态
\passthrough{\lstinline!map!} 参数，Python 类型为
\passthrough{\lstinline!dict[str, str]!}。精确含义、单位、范围和枚举值需查目标型号对应头文件与协议。
对应 C++ 参数
\passthrough{\lstinline!state\_map: std::map<std::string, std::string> \&!}。 \\
\bottomrule
\end{longtable}

\textbf{对应 C++}

\begin{itemize}
\tightlist
\item
  类：\passthrough{\lstinline!unitree::robot::as2::SportClient!}
\item
  签名：\passthrough{\lstinline!GetState(std::map<std::string, std::string> \&)!}
\item
  绑定策略：\passthrough{\lstinline!OUTPUT\_WRAPPER!}
\item
  声明位置：\passthrough{\lstinline!include/unitree/robot/as2/sport/sport\_client.hpp:328!}
\end{itemize}

\textbf{说明}

\begin{itemize}
\tightlist
\item
  示例是局部调用片段；\passthrough{\lstinline!obj!}
  和参数变量表示已经按上层业务校验并准备好的对象。
\item
  C++ 的可修改输出引用不会作为 Python 入参出现，而是按 C++
  参数顺序追加到返回元组中。
\end{itemize}

\textbf{用法}

\begin{lstlisting}[language=Python]
status, state_map = obj.get_state()
\end{lstlisting}

\begin{quote}
{[}!NOTE{]} 只读查询仍会访问
DDS/机器人服务。必须检查返回状态码，并处理超时和网络中断。
\end{quote}

\begin{center}\rule{0.5\linewidth}{0.5pt}\end{center}

\hypertarget{unitree-sdk2-cpp-robot-b2}{%
\subsection{B2 Robot API}\label{unitree-sdk2-cpp-robot-b2}}

模块：\passthrough{\lstinline!unitree\_sdk2\_cpp.robot.b2!}

\hypertarget{ux7c7bux7d22ux5f15-9}{%
\subsubsection{类索引}\label{ux7c7bux7d22ux5f15-9}}

\begin{longtable}[]{@{}
  >{\raggedright\arraybackslash}p{(\columnwidth - 4\tabcolsep) * \real{0.18}}
  >{\raggedleft\arraybackslash}p{(\columnwidth - 4\tabcolsep) * \real{0.20}}
  >{\raggedleft\arraybackslash}p{(\columnwidth - 4\tabcolsep) * \real{0.62}}@{}}
\toprule
类 & 公开函数签名 & 属性 \\
\midrule
\endhead
\protect\hyperlink{unitree-sdk2-cpp-robot-b2-backvideoclient}{\passthrough{\lstinline!BackVideoClient!}}
& 3 & 0 \\
\protect\hyperlink{unitree-sdk2-cpp-robot-b2-configclient}{\passthrough{\lstinline!ConfigClient!}}
& 8 & 0 \\
\protect\hyperlink{unitree-sdk2-cpp-robot-b2-configdelparameter}{\passthrough{\lstinline!ConfigDelParameter!}}
& 3 & 0 \\
\protect\hyperlink{unitree-sdk2-cpp-robot-b2-configgetdata}{\passthrough{\lstinline!ConfigGetData!}}
& 3 & 0 \\
\protect\hyperlink{unitree-sdk2-cpp-robot-b2-configgetparameter}{\passthrough{\lstinline!ConfigGetParameter!}}
& 3 & 0 \\
\protect\hyperlink{unitree-sdk2-cpp-robot-b2-configmeta}{\passthrough{\lstinline!ConfigMeta!}}
& 1 & 0 \\
\protect\hyperlink{unitree-sdk2-cpp-robot-b2-configmetadata}{\passthrough{\lstinline!ConfigMetaData!}}
& 3 & 0 \\
\protect\hyperlink{unitree-sdk2-cpp-robot-b2-configmetaparameter}{\passthrough{\lstinline!ConfigMetaParameter!}}
& 3 & 0 \\
\protect\hyperlink{unitree-sdk2-cpp-robot-b2-configsetparameter}{\passthrough{\lstinline!ConfigSetParameter!}}
& 3 & 0 \\
\protect\hyperlink{unitree-sdk2-cpp-robot-b2-frontvideoclient}{\passthrough{\lstinline!FrontVideoClient!}}
& 3 & 0 \\
\protect\hyperlink{unitree-sdk2-cpp-robot-b2-jsonizeconfigmeta}{\passthrough{\lstinline!JsonizeConfigMeta!}}
& 3 & 0 \\
\protect\hyperlink{unitree-sdk2-cpp-robot-b2-jsonizemodename}{\passthrough{\lstinline!JsonizeModeName!}}
& 3 & 0 \\
\protect\hyperlink{unitree-sdk2-cpp-robot-b2-jsonizesilent}{\passthrough{\lstinline!JsonizeSilent!}}
& 3 & 0 \\
\protect\hyperlink{unitree-sdk2-cpp-robot-b2-lowpowerstatusdata}{\passthrough{\lstinline!LowPowerStatusData!}}
& 3 & 0 \\
\protect\hyperlink{unitree-sdk2-cpp-robot-b2-lowpowerswitchparameter}{\passthrough{\lstinline!LowPowerSwitchParameter!}}
& 3 & 0 \\
\protect\hyperlink{unitree-sdk2-cpp-robot-b2-motionswitcherclient}{\passthrough{\lstinline!MotionSwitcherClient!}}
& 7 & 0 \\
\protect\hyperlink{unitree-sdk2-cpp-robot-b2-pkgversiondata}{\passthrough{\lstinline!PkgVersionData!}}
& 3 & 0 \\
\protect\hyperlink{unitree-sdk2-cpp-robot-b2-robotstateclient}{\passthrough{\lstinline!RobotStateClient!}}
& 8 & 0 \\
\protect\hyperlink{unitree-sdk2-cpp-robot-b2-servicestate}{\passthrough{\lstinline!ServiceState!}}
& 1 & 0 \\
\protect\hyperlink{unitree-sdk2-cpp-robot-b2-servicestatedata}{\passthrough{\lstinline!ServiceStateData!}}
& 3 & 0 \\
\protect\hyperlink{unitree-sdk2-cpp-robot-b2-serviceswitchdata}{\passthrough{\lstinline!ServiceSwitchData!}}
& 3 & 0 \\
\protect\hyperlink{unitree-sdk2-cpp-robot-b2-serviceswitchparameter}{\passthrough{\lstinline!ServiceSwitchParameter!}}
& 3 & 0 \\
\protect\hyperlink{unitree-sdk2-cpp-robot-b2-setreportfreqparameter}{\passthrough{\lstinline!SetReportFreqParameter!}}
& 3 & 0 \\
\protect\hyperlink{unitree-sdk2-cpp-robot-b2-sportclient}{\passthrough{\lstinline!SportClient!}}
& 25 & 0 \\
\protect\hyperlink{unitree-sdk2-cpp-robot-b2-stpathpoint}{\passthrough{\lstinline!stPathPoint!}}
& 0 & 0 \\
\bottomrule
\end{longtable}

\hypertarget{unitree-sdk2-cpp-robot-b2-backvideoclient}{%
\subsubsection{\texorpdfstring{\texttt{unitree\_sdk2\_cpp.robot.b2.BackVideoClient}}{unitree\_sdk2\_cpp.robot.b2.BackVideoClient}}\label{unitree-sdk2-cpp-robot-b2-backvideoclient}}

Robot 服务客户端。构造、初始化和具体方法的可用性必须分别检查。

\textbf{导入}

\begin{lstlisting}[language=Python]
from unitree_sdk2_cpp.robot.b2 import BackVideoClient
\end{lstlisting}

\textbf{基类}：\passthrough{\lstinline!Client!}

\textbf{方法索引}

\begin{longtable}[]{@{}
  >{\raggedright\arraybackslash}p{(\columnwidth - 6\tabcolsep) * \real{0.18}}
  >{\raggedright\arraybackslash}p{(\columnwidth - 6\tabcolsep) * \real{0.24}}
  >{\raggedright\arraybackslash}p{(\columnwidth - 6\tabcolsep) * \real{0.18}}
  >{\raggedright\arraybackslash}p{(\columnwidth - 6\tabcolsep) * \real{0.40}}@{}}
\toprule
方法 & Python 签名 & 状态 & 安全分类 \\
\midrule
\endhead
\protect\hyperlink{unitree-sdk2-cpp-robot-b2-backvideoclient-dunder-init-1}{\passthrough{\lstinline!\_\_init\_\_!}}
& \passthrough{\lstinline!def \_\_init\_\_(self) -> None!} &
\passthrough{\lstinline!AVAILABLE!} &
\passthrough{\lstinline!CONSTRUCTION!} \\
\protect\hyperlink{unitree-sdk2-cpp-robot-b2-backvideoclient-init-1}{\passthrough{\lstinline!init!}}
& \passthrough{\lstinline!def init(self) -> None!} &
\passthrough{\lstinline!AVAILABLE!} &
\passthrough{\lstinline!INITIALIZATION!} \\
\protect\hyperlink{unitree-sdk2-cpp-robot-b2-backvideoclient-get-image-sample-1}{\passthrough{\lstinline!get\_image\_sample!}}
&
\passthrough{\lstinline!def get\_image\_sample(self) -> tuple[int, list[int]]!}
& \passthrough{\lstinline!AVAILABLE!} &
\passthrough{\lstinline!READ\_ONLY!} \\
\bottomrule
\end{longtable}

\hypertarget{unitree-sdk2-cpp-robot-b2-backvideoclient-dunder-init-1}{%
\paragraph{\texorpdfstring{\texttt{unitree\_sdk2\_cpp.robot.b2.BackVideoClient.\_\_init\_\_}}{unitree\_sdk2\_cpp.robot.b2.BackVideoClient.\_\_init\_\_}}\label{unitree-sdk2-cpp-robot-b2-backvideoclient-dunder-init-1}}

初始化 \passthrough{\lstinline!BackVideoClient!}
实例。是否能实际构造取决于下方可用性状态。

\textbf{签名}

\begin{lstlisting}[language=Python]
def __init__(self) -> None
\end{lstlisting}

\textbf{可用性与安全性}

\textbf{\passthrough{\lstinline!AVAILABLE!} /
\passthrough{\lstinline!CONSTRUCTION!}}：当前绑定源码已实现；实际执行条件仍取决于平台和对应
SDK 环境。

\textbf{参数}

无显式参数。实例方法中的 \passthrough{\lstinline!self!} 由 Python
自动传入。

\textbf{返回值}

\begin{longtable}[]{@{}
  >{\raggedright\arraybackslash}p{(\columnwidth - 4\tabcolsep) * \real{0.18}}
  >{\raggedright\arraybackslash}p{(\columnwidth - 4\tabcolsep) * \real{0.20}}
  >{\raggedright\arraybackslash}p{(\columnwidth - 4\tabcolsep) * \real{0.62}}@{}}
\toprule
位置 & 类型 & 含义 \\
\midrule
\endhead
无 & \passthrough{\lstinline!None!} &
\passthrough{\lstinline!\_\_init\_\_!}
本身不返回值；调用类对象时，在构造函数可用的前提下得到该类实例。 \\
\bottomrule
\end{longtable}

\textbf{对应 C++}

\begin{itemize}
\tightlist
\item
  类：\passthrough{\lstinline!unitree::robot::b2::BackVideoClient!}
\item
  签名：\passthrough{\lstinline!BackVideoClient()!}
\item
  绑定策略：\passthrough{\lstinline!未单独标注!}
\item
  声明位置：\passthrough{\lstinline!include/unitree/robot/b2/back\_video/back\_video\_client.hpp:18!}
\end{itemize}

\textbf{说明}

\begin{itemize}
\tightlist
\item
  示例是局部调用片段；\passthrough{\lstinline!obj!}
  和参数变量表示已经按上层业务校验并准备好的对象。
\end{itemize}

\textbf{用法}

\begin{lstlisting}[language=Python]
value = BackVideoClient()
\end{lstlisting}

\hypertarget{unitree-sdk2-cpp-robot-b2-backvideoclient-init-1}{%
\paragraph{\texorpdfstring{\texttt{unitree\_sdk2\_cpp.robot.b2.BackVideoClient.init}}{unitree\_sdk2\_cpp.robot.b2.BackVideoClient.init}}\label{unitree-sdk2-cpp-robot-b2-backvideoclient-init-1}}

初始化当前 SDK 对象所需的底层通道或服务资源。

\textbf{签名}

\begin{lstlisting}[language=Python]
def init(self) -> None
\end{lstlisting}

\textbf{可用性与安全性}

\textbf{\passthrough{\lstinline!AVAILABLE!} /
\passthrough{\lstinline!INITIALIZATION!}}：当前绑定源码已实现；实际执行条件仍取决于平台和对应
SDK 环境。

\textbf{参数}

无显式参数。实例方法中的 \passthrough{\lstinline!self!} 由 Python
自动传入。

\textbf{返回值}

\begin{longtable}[]{@{}
  >{\raggedright\arraybackslash}p{(\columnwidth - 4\tabcolsep) * \real{0.18}}
  >{\raggedright\arraybackslash}p{(\columnwidth - 4\tabcolsep) * \real{0.20}}
  >{\raggedright\arraybackslash}p{(\columnwidth - 4\tabcolsep) * \real{0.62}}@{}}
\toprule
位置 & 类型 & 含义 \\
\midrule
\endhead
无 & \passthrough{\lstinline!None!} & 该方法不返回 Python
值；失败可能表现为异常或底层状态变化。 \\
\bottomrule
\end{longtable}

\textbf{对应 C++}

\begin{itemize}
\tightlist
\item
  类：\passthrough{\lstinline!unitree::robot::b2::BackVideoClient!}
\item
  签名：\passthrough{\lstinline!Init()!}
\item
  绑定策略：\passthrough{\lstinline!DIRECT!}
\item
  声明位置：\passthrough{\lstinline!include/unitree/robot/b2/back\_video/back\_video\_client.hpp:21!}
\end{itemize}

\textbf{说明}

\begin{itemize}
\tightlist
\item
  示例是局部调用片段；\passthrough{\lstinline!obj!}
  和参数变量表示已经按上层业务校验并准备好的对象。
\end{itemize}

\textbf{用法}

\begin{lstlisting}[language=Python]
obj.init()
\end{lstlisting}

\hypertarget{unitree-sdk2-cpp-robot-b2-backvideoclient-get-image-sample-1}{%
\paragraph{\texorpdfstring{\texttt{unitree\_sdk2\_cpp.robot.b2.BackVideoClient.get\_image\_sample}}{unitree\_sdk2\_cpp.robot.b2.BackVideoClient.get\_image\_sample}}\label{unitree-sdk2-cpp-robot-b2-backvideoclient-get-image-sample-1}}

查询或检查 图像 \passthrough{\lstinline!sample!}。

\textbf{签名}

\begin{lstlisting}[language=Python]
def get_image_sample(self) -> tuple[int, list[int]]
\end{lstlisting}

\textbf{可用性与安全性}

\textbf{\passthrough{\lstinline!AVAILABLE!} /
\passthrough{\lstinline!READ\_ONLY!}}：当前绑定源码已实现只读查询。Robot
Client 的服务端查询通常仍需匹配的 Linux
扩展、DDS、网络、机器人和服务；纯本地 getter 则不一定需要全部条件。

\textbf{参数}

无显式参数。实例方法中的 \passthrough{\lstinline!self!} 由 Python
自动传入。

\textbf{返回值}

\begin{longtable}[]{@{}
  >{\raggedright\arraybackslash}p{(\columnwidth - 4\tabcolsep) * \real{0.18}}
  >{\raggedright\arraybackslash}p{(\columnwidth - 4\tabcolsep) * \real{0.20}}
  >{\raggedright\arraybackslash}p{(\columnwidth - 4\tabcolsep) * \real{0.62}}@{}}
\toprule
位置 & 类型 & 含义 \\
\midrule
\endhead
{[}0{]} \passthrough{\lstinline!status!} & \passthrough{\lstinline!int!}
& SDK 状态码；Unitree 示例通常以 \passthrough{\lstinline!0!}
表示成功，非零值需按具体服务错误码解释。 \\
{[}1{]} \passthrough{\lstinline!output\_1!} &
\passthrough{\lstinline!list[int]!} & 传给该接口的
\passthrough{\lstinline!output!} \passthrough{\lstinline!1!}
参数，Python 类型为
\passthrough{\lstinline!list[int]!}。精确含义、单位、范围和枚举值需查目标型号对应头文件与协议。
对应 C++ 参数
\passthrough{\lstinline!output\_1: std::vector<uint8\_t> \&!}。
底层是可变长度 vector；元素约束：取值范围为 0 到 255。 \\
\bottomrule
\end{longtable}

\textbf{对应 C++}

\begin{itemize}
\tightlist
\item
  类：\passthrough{\lstinline!unitree::robot::b2::BackVideoClient!}
\item
  签名：\passthrough{\lstinline!GetImageSample(std::vector<uint8\_t> \&)!}
\item
  绑定策略：\passthrough{\lstinline!OUTPUT\_WRAPPER!}
\item
  声明位置：\passthrough{\lstinline!include/unitree/robot/b2/back\_video/back\_video\_client.hpp:27!}
\end{itemize}

\textbf{说明}

\begin{itemize}
\tightlist
\item
  示例是局部调用片段；\passthrough{\lstinline!obj!}
  和参数变量表示已经按上层业务校验并准备好的对象。
\item
  C++ 的可修改输出引用不会作为 Python 入参出现，而是按 C++
  参数顺序追加到返回元组中。
\end{itemize}

\textbf{用法}

\begin{lstlisting}[language=Python]
status, output_1 = obj.get_image_sample()
\end{lstlisting}

\begin{quote}
{[}!NOTE{]} 只读查询仍会访问
DDS/机器人服务。必须检查返回状态码，并处理超时和网络中断。
\end{quote}

\hypertarget{unitree-sdk2-cpp-robot-b2-configclient}{%
\subsubsection{\texorpdfstring{\texttt{unitree\_sdk2\_cpp.robot.b2.ConfigClient}}{unitree\_sdk2\_cpp.robot.b2.ConfigClient}}\label{unitree-sdk2-cpp-robot-b2-configclient}}

Robot 服务客户端。构造、初始化和具体方法的可用性必须分别检查。

\textbf{导入}

\begin{lstlisting}[language=Python]
from unitree_sdk2_cpp.robot.b2 import ConfigClient
\end{lstlisting}

\textbf{基类}：\passthrough{\lstinline!Client!}

\textbf{方法索引}

\begin{longtable}[]{@{}
  >{\raggedright\arraybackslash}p{(\columnwidth - 6\tabcolsep) * \real{0.18}}
  >{\raggedright\arraybackslash}p{(\columnwidth - 6\tabcolsep) * \real{0.24}}
  >{\raggedright\arraybackslash}p{(\columnwidth - 6\tabcolsep) * \real{0.18}}
  >{\raggedright\arraybackslash}p{(\columnwidth - 6\tabcolsep) * \real{0.40}}@{}}
\toprule
方法 & Python 签名 & 状态 & 安全分类 \\
\midrule
\endhead
\protect\hyperlink{unitree-sdk2-cpp-robot-b2-configclient-dunder-init-1}{\passthrough{\lstinline!\_\_init\_\_!}}
& \passthrough{\lstinline!def \_\_init\_\_(self) -> None!} &
\passthrough{\lstinline!AVAILABLE!} &
\passthrough{\lstinline!CONSTRUCTION!} \\
\protect\hyperlink{unitree-sdk2-cpp-robot-b2-configclient-init-1}{\passthrough{\lstinline!init!}}
& \passthrough{\lstinline!def init(self) -> None!} &
\passthrough{\lstinline!AVAILABLE!} &
\passthrough{\lstinline!INITIALIZATION!} \\
\protect\hyperlink{unitree-sdk2-cpp-robot-b2-configclient-set-1}{\passthrough{\lstinline!set!}}
&
\passthrough{\lstinline!def set(self, name: str, content: str) -> int!}
& \passthrough{\lstinline!SIGNATURE\_ONLY!} &
\passthrough{\lstinline!HARDWARE\_SIDE\_EFFECT!} \\
\protect\hyperlink{unitree-sdk2-cpp-robot-b2-configclient-get-1}{\passthrough{\lstinline!get!}}
& \passthrough{\lstinline!def get(self, name: str) -> tuple[int, str]!}
& \passthrough{\lstinline!AVAILABLE!} &
\passthrough{\lstinline!READ\_ONLY!} \\
\protect\hyperlink{unitree-sdk2-cpp-robot-b2-configclient-del-1}{\passthrough{\lstinline!del\_!}}
& \passthrough{\lstinline!def del\_(self, name: str) -> int!} &
\passthrough{\lstinline!SIGNATURE\_ONLY!} &
\passthrough{\lstinline!HARDWARE\_SIDE\_EFFECT!} \\
\protect\hyperlink{unitree-sdk2-cpp-robot-b2-configclient-meta-config-meta-1}{\passthrough{\lstinline!meta\_config\_meta!}}
&
\passthrough{\lstinline!def meta\_config\_meta(self, name: str) -> tuple[int, ConfigMeta]!}
& \passthrough{\lstinline!SIGNATURE\_ONLY!} &
\passthrough{\lstinline!HARDWARE\_SIDE\_EFFECT!} \\
\protect\hyperlink{unitree-sdk2-cpp-robot-b2-configclient-meta-string-1}{\passthrough{\lstinline!meta\_string!}}
&
\passthrough{\lstinline!def meta\_string(self, name: str) -> tuple[int, str]!}
& \passthrough{\lstinline!SIGNATURE\_ONLY!} &
\passthrough{\lstinline!HARDWARE\_SIDE\_EFFECT!} \\
\protect\hyperlink{unitree-sdk2-cpp-robot-b2-configclient-subscribe-change-status-1}{\passthrough{\lstinline!subscribe\_change\_status!}}
&
\passthrough{\lstinline!def subscribe\_change\_status(self, name: str, callback: Callable[[str, str], None]) -> None!}
& \passthrough{\lstinline!SIGNATURE\_ONLY!} &
\passthrough{\lstinline!HARDWARE\_SIDE\_EFFECT!} \\
\bottomrule
\end{longtable}

\hypertarget{unitree-sdk2-cpp-robot-b2-configclient-dunder-init-1}{%
\paragraph{\texorpdfstring{\texttt{unitree\_sdk2\_cpp.robot.b2.ConfigClient.\_\_init\_\_}}{unitree\_sdk2\_cpp.robot.b2.ConfigClient.\_\_init\_\_}}\label{unitree-sdk2-cpp-robot-b2-configclient-dunder-init-1}}

初始化 \passthrough{\lstinline!ConfigClient!}
实例。是否能实际构造取决于下方可用性状态。

\textbf{签名}

\begin{lstlisting}[language=Python]
def __init__(self) -> None
\end{lstlisting}

\textbf{可用性与安全性}

\textbf{\passthrough{\lstinline!AVAILABLE!} /
\passthrough{\lstinline!CONSTRUCTION!}}：当前绑定源码已实现；实际执行条件仍取决于平台和对应
SDK 环境。

\textbf{参数}

无显式参数。实例方法中的 \passthrough{\lstinline!self!} 由 Python
自动传入。

\textbf{返回值}

\begin{longtable}[]{@{}
  >{\raggedright\arraybackslash}p{(\columnwidth - 4\tabcolsep) * \real{0.18}}
  >{\raggedright\arraybackslash}p{(\columnwidth - 4\tabcolsep) * \real{0.20}}
  >{\raggedright\arraybackslash}p{(\columnwidth - 4\tabcolsep) * \real{0.62}}@{}}
\toprule
位置 & 类型 & 含义 \\
\midrule
\endhead
无 & \passthrough{\lstinline!None!} &
\passthrough{\lstinline!\_\_init\_\_!}
本身不返回值；调用类对象时，在构造函数可用的前提下得到该类实例。 \\
\bottomrule
\end{longtable}

\textbf{对应 C++}

\begin{itemize}
\tightlist
\item
  类：\passthrough{\lstinline!unitree::robot::b2::ConfigClient!}
\item
  签名：\passthrough{\lstinline!ConfigClient()!}
\item
  绑定策略：\passthrough{\lstinline!未单独标注!}
\item
  声明位置：\passthrough{\lstinline!include/unitree/robot/b2/config/config\_client.hpp:41!}
\end{itemize}

\textbf{说明}

\begin{itemize}
\tightlist
\item
  示例是局部调用片段；\passthrough{\lstinline!obj!}
  和参数变量表示已经按上层业务校验并准备好的对象。
\end{itemize}

\textbf{用法}

\begin{lstlisting}[language=Python]
value = ConfigClient()
\end{lstlisting}

\hypertarget{unitree-sdk2-cpp-robot-b2-configclient-init-1}{%
\paragraph{\texorpdfstring{\texttt{unitree\_sdk2\_cpp.robot.b2.ConfigClient.init}}{unitree\_sdk2\_cpp.robot.b2.ConfigClient.init}}\label{unitree-sdk2-cpp-robot-b2-configclient-init-1}}

初始化当前 SDK 对象所需的底层通道或服务资源。

\textbf{签名}

\begin{lstlisting}[language=Python]
def init(self) -> None
\end{lstlisting}

\textbf{可用性与安全性}

\textbf{\passthrough{\lstinline!AVAILABLE!} /
\passthrough{\lstinline!INITIALIZATION!}}：当前绑定源码已实现；实际执行条件仍取决于平台和对应
SDK 环境。

\textbf{参数}

无显式参数。实例方法中的 \passthrough{\lstinline!self!} 由 Python
自动传入。

\textbf{返回值}

\begin{longtable}[]{@{}
  >{\raggedright\arraybackslash}p{(\columnwidth - 4\tabcolsep) * \real{0.18}}
  >{\raggedright\arraybackslash}p{(\columnwidth - 4\tabcolsep) * \real{0.20}}
  >{\raggedright\arraybackslash}p{(\columnwidth - 4\tabcolsep) * \real{0.62}}@{}}
\toprule
位置 & 类型 & 含义 \\
\midrule
\endhead
无 & \passthrough{\lstinline!None!} & 该方法不返回 Python
值；失败可能表现为异常或底层状态变化。 \\
\bottomrule
\end{longtable}

\textbf{对应 C++}

\begin{itemize}
\tightlist
\item
  类：\passthrough{\lstinline!unitree::robot::b2::ConfigClient!}
\item
  签名：\passthrough{\lstinline!Init()!}
\item
  绑定策略：\passthrough{\lstinline!DIRECT!}
\item
  声明位置：\passthrough{\lstinline!include/unitree/robot/b2/config/config\_client.hpp:44!}
\end{itemize}

\textbf{说明}

\begin{itemize}
\tightlist
\item
  示例是局部调用片段；\passthrough{\lstinline!obj!}
  和参数变量表示已经按上层业务校验并准备好的对象。
\end{itemize}

\textbf{用法}

\begin{lstlisting}[language=Python]
obj.init()
\end{lstlisting}

\hypertarget{unitree-sdk2-cpp-robot-b2-configclient-set-1}{%
\paragraph{\texorpdfstring{\texttt{unitree\_sdk2\_cpp.robot.b2.ConfigClient.set}}{unitree\_sdk2\_cpp.robot.b2.ConfigClient.set}}\label{unitree-sdk2-cpp-robot-b2-configclient-set-1}}

对应 C++ SDK 操作
\passthrough{\lstinline!Set(const std::string \&, const std::string \&)!}。上游头文件没有可直接生成的业务说明时，本参考只保证签名映射，精确语义需查目标型号协议。

\textbf{签名}

\begin{lstlisting}[language=Python]
def set(self, name: str, content: str) -> int
\end{lstlisting}

\textbf{可用性与安全性}

\textbf{\passthrough{\lstinline!SIGNATURE\_ONLY!} /
\passthrough{\lstinline!HARDWARE\_SIDE\_EFFECT!}}：仅用于补全和静态检查。当前二进制没有该
Python 方法，未来实现还需单独评审硬件副作用。

\textbf{参数}

\begin{longtable}[]{@{}
  >{\raggedright\arraybackslash}p{(\columnwidth - 6\tabcolsep) * \real{0.18}}
  >{\raggedright\arraybackslash}p{(\columnwidth - 6\tabcolsep) * \real{0.24}}
  >{\raggedright\arraybackslash}p{(\columnwidth - 6\tabcolsep) * \real{0.18}}
  >{\raggedright\arraybackslash}p{(\columnwidth - 6\tabcolsep) * \real{0.40}}@{}}
\toprule
名称 & Python 类型 & 默认值 & 含义 \\
\midrule
\endhead
\passthrough{\lstinline!name!} & \passthrough{\lstinline!str!} & 必填 &
SDK 对象、配置项、服务或动作的名称。允许值由调用它的具体接口定义。 对应
C++ 参数 \passthrough{\lstinline!name: const std::string \&!}。
底层为字符串；长度、编码和允许值由具体协议决定。 \\
\passthrough{\lstinline!content!} & \passthrough{\lstinline!str!} & 必填
& 配置、请求或序列化内容。具体格式由对应服务协议定义。 对应 C++ 参数
\passthrough{\lstinline!content: const std::string \&!}。
底层为字符串；长度、编码和允许值由具体协议决定。 \\
\bottomrule
\end{longtable}

\textbf{返回值}

\begin{longtable}[]{@{}
  >{\raggedright\arraybackslash}p{(\columnwidth - 4\tabcolsep) * \real{0.18}}
  >{\raggedright\arraybackslash}p{(\columnwidth - 4\tabcolsep) * \real{0.20}}
  >{\raggedright\arraybackslash}p{(\columnwidth - 4\tabcolsep) * \real{0.62}}@{}}
\toprule
位置 & 类型 & 含义 \\
\midrule
\endhead
返回值 & \passthrough{\lstinline!int!} & SDK 状态码；Unitree 示例通常以
\passthrough{\lstinline!0!} 表示成功，非零值需按具体服务定义解释。 \\
\bottomrule
\end{longtable}

\textbf{对应 C++}

\begin{itemize}
\tightlist
\item
  类：\passthrough{\lstinline!unitree::robot::b2::ConfigClient!}
\item
  签名：\passthrough{\lstinline!Set(const std::string \&, const std::string \&)!}
\item
  绑定策略：\passthrough{\lstinline!DIRECT!}
\item
  声明位置：\passthrough{\lstinline!include/unitree/robot/b2/config/config\_client.hpp:46!}
\end{itemize}

\textbf{说明}

\begin{itemize}
\tightlist
\item
  示例是局部调用片段；\passthrough{\lstinline!obj!}
  和参数变量表示已经按上层业务校验并准备好的对象。
\item
  该条目可以被编辑器和类型检查器识别，但当前运行时可能在属性查找阶段直接失败。
\end{itemize}

\textbf{用法}

\begin{lstlisting}[language=Python]
from typing import TYPE_CHECKING

if TYPE_CHECKING:
    def planned_set(obj: ConfigClient, name: str, content: str) -> int:
        return obj.set(name=name, content=content)
\end{lstlisting}

\begin{quote}
{[}!CAUTION{]} 上面的 \passthrough{\lstinline!TYPE\_CHECKING!}
示例只表达计划调用形态。该分支在普通 Python
运行时为假，不代表当前扩展能执行此接口。
\end{quote}

\hypertarget{unitree-sdk2-cpp-robot-b2-configclient-get-1}{%
\paragraph{\texorpdfstring{\texttt{unitree\_sdk2\_cpp.robot.b2.ConfigClient.get}}{unitree\_sdk2\_cpp.robot.b2.ConfigClient.get}}\label{unitree-sdk2-cpp-robot-b2-configclient-get-1}}

对应 C++ SDK 操作
\passthrough{\lstinline!Get(const std::string \&, std::string \&)!}。上游头文件没有可直接生成的业务说明时，本参考只保证签名映射，精确语义需查目标型号协议。

\textbf{签名}

\begin{lstlisting}[language=Python]
def get(self, name: str) -> tuple[int, str]
\end{lstlisting}

\textbf{可用性与安全性}

\textbf{\passthrough{\lstinline!AVAILABLE!} /
\passthrough{\lstinline!READ\_ONLY!}}：当前绑定源码已实现只读查询。Robot
Client 的服务端查询通常仍需匹配的 Linux
扩展、DDS、网络、机器人和服务；纯本地 getter 则不一定需要全部条件。

\textbf{参数}

\begin{longtable}[]{@{}
  >{\raggedright\arraybackslash}p{(\columnwidth - 6\tabcolsep) * \real{0.18}}
  >{\raggedright\arraybackslash}p{(\columnwidth - 6\tabcolsep) * \real{0.24}}
  >{\raggedright\arraybackslash}p{(\columnwidth - 6\tabcolsep) * \real{0.18}}
  >{\raggedright\arraybackslash}p{(\columnwidth - 6\tabcolsep) * \real{0.40}}@{}}
\toprule
名称 & Python 类型 & 默认值 & 含义 \\
\midrule
\endhead
\passthrough{\lstinline!name!} & \passthrough{\lstinline!str!} & 必填 &
SDK 对象、配置项、服务或动作的名称。允许值由调用它的具体接口定义。 对应
C++ 参数 \passthrough{\lstinline!name: const std::string \&!}。
底层为字符串；长度、编码和允许值由具体协议决定。 \\
\bottomrule
\end{longtable}

\textbf{返回值}

\begin{longtable}[]{@{}
  >{\raggedright\arraybackslash}p{(\columnwidth - 4\tabcolsep) * \real{0.18}}
  >{\raggedright\arraybackslash}p{(\columnwidth - 4\tabcolsep) * \real{0.20}}
  >{\raggedright\arraybackslash}p{(\columnwidth - 4\tabcolsep) * \real{0.62}}@{}}
\toprule
位置 & 类型 & 含义 \\
\midrule
\endhead
{[}0{]} \passthrough{\lstinline!status!} & \passthrough{\lstinline!int!}
& SDK 状态码；Unitree 示例通常以 \passthrough{\lstinline!0!}
表示成功，非零值需按具体服务错误码解释。 \\
{[}1{]} \passthrough{\lstinline!content!} &
\passthrough{\lstinline!str!} &
配置、请求或序列化内容。具体格式由对应服务协议定义。 对应 C++ 参数
\passthrough{\lstinline!content: std::string \&!}。
底层为字符串；长度、编码和允许值由具体协议决定。 \\
\bottomrule
\end{longtable}

\textbf{对应 C++}

\begin{itemize}
\tightlist
\item
  类：\passthrough{\lstinline!unitree::robot::b2::ConfigClient!}
\item
  签名：\passthrough{\lstinline!Get(const std::string \&, std::string \&)!}
\item
  绑定策略：\passthrough{\lstinline!OUTPUT\_WRAPPER!}
\item
  声明位置：\passthrough{\lstinline!include/unitree/robot/b2/config/config\_client.hpp:47!}
\end{itemize}

\textbf{说明}

\begin{itemize}
\tightlist
\item
  示例是局部调用片段；\passthrough{\lstinline!obj!}
  和参数变量表示已经按上层业务校验并准备好的对象。
\item
  C++ 的可修改输出引用不会作为 Python 入参出现，而是按 C++
  参数顺序追加到返回元组中。
\end{itemize}

\textbf{用法}

\begin{lstlisting}[language=Python]
status, content = obj.get(name=name)
\end{lstlisting}

\begin{quote}
{[}!NOTE{]} 只读查询仍会访问
DDS/机器人服务。必须检查返回状态码，并处理超时和网络中断。
\end{quote}

\hypertarget{unitree-sdk2-cpp-robot-b2-configclient-del-1}{%
\paragraph{\texorpdfstring{\texttt{unitree\_sdk2\_cpp.robot.b2.ConfigClient.del\_}}{unitree\_sdk2\_cpp.robot.b2.ConfigClient.del\_}}\label{unitree-sdk2-cpp-robot-b2-configclient-del-1}}

对应 C++ SDK 操作
\passthrough{\lstinline!Del(const std::string \&)!}。上游头文件没有可直接生成的业务说明时，本参考只保证签名映射，精确语义需查目标型号协议。

\textbf{签名}

\begin{lstlisting}[language=Python]
def del_(self, name: str) -> int
\end{lstlisting}

\textbf{可用性与安全性}

\textbf{\passthrough{\lstinline!SIGNATURE\_ONLY!} /
\passthrough{\lstinline!HARDWARE\_SIDE\_EFFECT!}}：仅用于补全和静态检查。当前二进制没有该
Python 方法，未来实现还需单独评审硬件副作用。

\textbf{参数}

\begin{longtable}[]{@{}
  >{\raggedright\arraybackslash}p{(\columnwidth - 6\tabcolsep) * \real{0.18}}
  >{\raggedright\arraybackslash}p{(\columnwidth - 6\tabcolsep) * \real{0.24}}
  >{\raggedright\arraybackslash}p{(\columnwidth - 6\tabcolsep) * \real{0.18}}
  >{\raggedright\arraybackslash}p{(\columnwidth - 6\tabcolsep) * \real{0.40}}@{}}
\toprule
名称 & Python 类型 & 默认值 & 含义 \\
\midrule
\endhead
\passthrough{\lstinline!name!} & \passthrough{\lstinline!str!} & 必填 &
SDK 对象、配置项、服务或动作的名称。允许值由调用它的具体接口定义。 对应
C++ 参数 \passthrough{\lstinline!name: const std::string \&!}。
底层为字符串；长度、编码和允许值由具体协议决定。 \\
\bottomrule
\end{longtable}

\textbf{返回值}

\begin{longtable}[]{@{}
  >{\raggedright\arraybackslash}p{(\columnwidth - 4\tabcolsep) * \real{0.18}}
  >{\raggedright\arraybackslash}p{(\columnwidth - 4\tabcolsep) * \real{0.20}}
  >{\raggedright\arraybackslash}p{(\columnwidth - 4\tabcolsep) * \real{0.62}}@{}}
\toprule
位置 & 类型 & 含义 \\
\midrule
\endhead
返回值 & \passthrough{\lstinline!int!} & SDK 状态码；Unitree 示例通常以
\passthrough{\lstinline!0!} 表示成功，非零值需按具体服务定义解释。 \\
\bottomrule
\end{longtable}

\textbf{对应 C++}

\begin{itemize}
\tightlist
\item
  类：\passthrough{\lstinline!unitree::robot::b2::ConfigClient!}
\item
  签名：\passthrough{\lstinline!Del(const std::string \&)!}
\item
  绑定策略：\passthrough{\lstinline!DIRECT!}
\item
  声明位置：\passthrough{\lstinline!include/unitree/robot/b2/config/config\_client.hpp:48!}
\end{itemize}

\textbf{说明}

\begin{itemize}
\tightlist
\item
  示例是局部调用片段；\passthrough{\lstinline!obj!}
  和参数变量表示已经按上层业务校验并准备好的对象。
\item
  该条目可以被编辑器和类型检查器识别，但当前运行时可能在属性查找阶段直接失败。
\end{itemize}

\textbf{用法}

\begin{lstlisting}[language=Python]
from typing import TYPE_CHECKING

if TYPE_CHECKING:
    def planned_del_(obj: ConfigClient, name: str) -> int:
        return obj.del_(name=name)
\end{lstlisting}

\begin{quote}
{[}!CAUTION{]} 上面的 \passthrough{\lstinline!TYPE\_CHECKING!}
示例只表达计划调用形态。该分支在普通 Python
运行时为假，不代表当前扩展能执行此接口。
\end{quote}

\hypertarget{unitree-sdk2-cpp-robot-b2-configclient-meta-config-meta-1}{%
\paragraph{\texorpdfstring{\texttt{unitree\_sdk2\_cpp.robot.b2.ConfigClient.meta\_config\_meta}}{unitree\_sdk2\_cpp.robot.b2.ConfigClient.meta\_config\_meta}}\label{unitree-sdk2-cpp-robot-b2-configclient-meta-config-meta-1}}

对应 C++ SDK 操作
\passthrough{\lstinline!Meta(const std::string \&, unitree::robot::b2::ConfigMeta \&)!}。上游头文件没有可直接生成的业务说明时，本参考只保证签名映射，精确语义需查目标型号协议。

\textbf{签名}

\begin{lstlisting}[language=Python]
def meta_config_meta(self, name: str) -> tuple[int, ConfigMeta]
\end{lstlisting}

\textbf{可用性与安全性}

\textbf{\passthrough{\lstinline!SIGNATURE\_ONLY!} /
\passthrough{\lstinline!HARDWARE\_SIDE\_EFFECT!}}：仅用于补全和静态检查。当前二进制没有该
Python 方法，未来实现还需单独评审硬件副作用。

\textbf{参数}

\begin{longtable}[]{@{}
  >{\raggedright\arraybackslash}p{(\columnwidth - 6\tabcolsep) * \real{0.18}}
  >{\raggedright\arraybackslash}p{(\columnwidth - 6\tabcolsep) * \real{0.24}}
  >{\raggedright\arraybackslash}p{(\columnwidth - 6\tabcolsep) * \real{0.18}}
  >{\raggedright\arraybackslash}p{(\columnwidth - 6\tabcolsep) * \real{0.40}}@{}}
\toprule
名称 & Python 类型 & 默认值 & 含义 \\
\midrule
\endhead
\passthrough{\lstinline!name!} & \passthrough{\lstinline!str!} & 必填 &
SDK 对象、配置项、服务或动作的名称。允许值由调用它的具体接口定义。 对应
C++ 参数 \passthrough{\lstinline!name: const std::string \&!}。
底层为字符串；长度、编码和允许值由具体协议决定。 \\
\bottomrule
\end{longtable}

\textbf{返回值}

\begin{longtable}[]{@{}
  >{\raggedright\arraybackslash}p{(\columnwidth - 4\tabcolsep) * \real{0.18}}
  >{\raggedright\arraybackslash}p{(\columnwidth - 4\tabcolsep) * \real{0.20}}
  >{\raggedright\arraybackslash}p{(\columnwidth - 4\tabcolsep) * \real{0.62}}@{}}
\toprule
位置 & 类型 & 含义 \\
\midrule
\endhead
{[}0{]} \passthrough{\lstinline!status!} & \passthrough{\lstinline!int!}
& SDK 状态码；Unitree 示例通常以 \passthrough{\lstinline!0!}
表示成功，非零值需按具体服务错误码解释。 \\
{[}1{]} \passthrough{\lstinline!meta!} &
\passthrough{\lstinline!ConfigMeta!} & 传给该接口的
\passthrough{\lstinline!meta!} 参数，Python 类型为
\passthrough{\lstinline!ConfigMeta!}。精确含义、单位、范围和枚举值需查目标型号对应头文件与协议。
对应 C++ 参数
\passthrough{\lstinline!meta: unitree::robot::b2::ConfigMeta \&!}。 \\
\bottomrule
\end{longtable}

\textbf{对应 C++}

\begin{itemize}
\tightlist
\item
  类：\passthrough{\lstinline!unitree::robot::b2::ConfigClient!}
\item
  签名：\passthrough{\lstinline!Meta(const std::string \&, unitree::robot::b2::ConfigMeta \&)!}
\item
  绑定策略：\passthrough{\lstinline!OUTPUT\_WRAPPER!}
\item
  声明位置：\passthrough{\lstinline!include/unitree/robot/b2/config/config\_client.hpp:49!}
\end{itemize}

\textbf{说明}

\begin{itemize}
\tightlist
\item
  示例是局部调用片段；\passthrough{\lstinline!obj!}
  和参数变量表示已经按上层业务校验并准备好的对象。
\item
  C++ 的可修改输出引用不会作为 Python 入参出现，而是按 C++
  参数顺序追加到返回元组中。
\item
  该条目可以被编辑器和类型检查器识别，但当前运行时可能在属性查找阶段直接失败。
\end{itemize}

\textbf{用法}

\begin{lstlisting}[language=Python]
from typing import TYPE_CHECKING

if TYPE_CHECKING:
    def planned_meta_config_meta(obj: ConfigClient, name: str) -> tuple[int, ConfigMeta]:
        return obj.meta_config_meta(name=name)
\end{lstlisting}

\begin{quote}
{[}!CAUTION{]} 上面的 \passthrough{\lstinline!TYPE\_CHECKING!}
示例只表达计划调用形态。该分支在普通 Python
运行时为假，不代表当前扩展能执行此接口。
\end{quote}

\hypertarget{unitree-sdk2-cpp-robot-b2-configclient-meta-string-1}{%
\paragraph{\texorpdfstring{\texttt{unitree\_sdk2\_cpp.robot.b2.ConfigClient.meta\_string}}{unitree\_sdk2\_cpp.robot.b2.ConfigClient.meta\_string}}\label{unitree-sdk2-cpp-robot-b2-configclient-meta-string-1}}

对应 C++ SDK 操作
\passthrough{\lstinline!Meta(const std::string \&, std::string \&)!}。上游头文件没有可直接生成的业务说明时，本参考只保证签名映射，精确语义需查目标型号协议。

\textbf{签名}

\begin{lstlisting}[language=Python]
def meta_string(self, name: str) -> tuple[int, str]
\end{lstlisting}

\textbf{可用性与安全性}

\textbf{\passthrough{\lstinline!SIGNATURE\_ONLY!} /
\passthrough{\lstinline!HARDWARE\_SIDE\_EFFECT!}}：仅用于补全和静态检查。当前二进制没有该
Python 方法，未来实现还需单独评审硬件副作用。

\textbf{参数}

\begin{longtable}[]{@{}
  >{\raggedright\arraybackslash}p{(\columnwidth - 6\tabcolsep) * \real{0.18}}
  >{\raggedright\arraybackslash}p{(\columnwidth - 6\tabcolsep) * \real{0.24}}
  >{\raggedright\arraybackslash}p{(\columnwidth - 6\tabcolsep) * \real{0.18}}
  >{\raggedright\arraybackslash}p{(\columnwidth - 6\tabcolsep) * \real{0.40}}@{}}
\toprule
名称 & Python 类型 & 默认值 & 含义 \\
\midrule
\endhead
\passthrough{\lstinline!name!} & \passthrough{\lstinline!str!} & 必填 &
SDK 对象、配置项、服务或动作的名称。允许值由调用它的具体接口定义。 对应
C++ 参数 \passthrough{\lstinline!name: const std::string \&!}。
底层为字符串；长度、编码和允许值由具体协议决定。 \\
\bottomrule
\end{longtable}

\textbf{返回值}

\begin{longtable}[]{@{}
  >{\raggedright\arraybackslash}p{(\columnwidth - 4\tabcolsep) * \real{0.18}}
  >{\raggedright\arraybackslash}p{(\columnwidth - 4\tabcolsep) * \real{0.20}}
  >{\raggedright\arraybackslash}p{(\columnwidth - 4\tabcolsep) * \real{0.62}}@{}}
\toprule
位置 & 类型 & 含义 \\
\midrule
\endhead
{[}0{]} \passthrough{\lstinline!status!} & \passthrough{\lstinline!int!}
& SDK 状态码；Unitree 示例通常以 \passthrough{\lstinline!0!}
表示成功，非零值需按具体服务错误码解释。 \\
{[}1{]} \passthrough{\lstinline!meta!} & \passthrough{\lstinline!str!} &
传给该接口的 \passthrough{\lstinline!meta!} 参数，Python 类型为
\passthrough{\lstinline!str!}。精确含义、单位、范围和枚举值需查目标型号对应头文件与协议。
对应 C++ 参数 \passthrough{\lstinline!meta: std::string \&!}。
底层为字符串；长度、编码和允许值由具体协议决定。 \\
\bottomrule
\end{longtable}

\textbf{对应 C++}

\begin{itemize}
\tightlist
\item
  类：\passthrough{\lstinline!unitree::robot::b2::ConfigClient!}
\item
  签名：\passthrough{\lstinline!Meta(const std::string \&, std::string \&)!}
\item
  绑定策略：\passthrough{\lstinline!OUTPUT\_WRAPPER!}
\item
  声明位置：\passthrough{\lstinline!include/unitree/robot/b2/config/config\_client.hpp:50!}
\end{itemize}

\textbf{说明}

\begin{itemize}
\tightlist
\item
  示例是局部调用片段；\passthrough{\lstinline!obj!}
  和参数变量表示已经按上层业务校验并准备好的对象。
\item
  C++ 的可修改输出引用不会作为 Python 入参出现，而是按 C++
  参数顺序追加到返回元组中。
\item
  该条目可以被编辑器和类型检查器识别，但当前运行时可能在属性查找阶段直接失败。
\end{itemize}

\textbf{用法}

\begin{lstlisting}[language=Python]
from typing import TYPE_CHECKING

if TYPE_CHECKING:
    def planned_meta_string(obj: ConfigClient, name: str) -> tuple[int, str]:
        return obj.meta_string(name=name)
\end{lstlisting}

\begin{quote}
{[}!CAUTION{]} 上面的 \passthrough{\lstinline!TYPE\_CHECKING!}
示例只表达计划调用形态。该分支在普通 Python
运行时为假，不代表当前扩展能执行此接口。
\end{quote}

\hypertarget{unitree-sdk2-cpp-robot-b2-configclient-subscribe-change-status-1}{%
\paragraph{\texorpdfstring{\texttt{unitree\_sdk2\_cpp.robot.b2.ConfigClient.subscribe\_change\_status}}{unitree\_sdk2\_cpp.robot.b2.ConfigClient.subscribe\_change\_status}}\label{unitree-sdk2-cpp-robot-b2-configclient-subscribe-change-status-1}}

对应 C++ SDK 操作
\passthrough{\lstinline!SubscribeChangeStatus(const std::string \&, const unitree::robot::b2::ConfigChangeStatusCallback \&)!}。上游头文件没有可直接生成的业务说明时，本参考只保证签名映射，精确语义需查目标型号协议。

\textbf{签名}

\begin{lstlisting}[language=Python]
def subscribe_change_status(self, name: str, callback: Callable[[str, str], None]) -> None
\end{lstlisting}

\textbf{可用性与安全性}

\textbf{\passthrough{\lstinline!SIGNATURE\_ONLY!} /
\passthrough{\lstinline!HARDWARE\_SIDE\_EFFECT!}}：仅用于补全和静态检查。当前二进制没有该
Python 方法，未来实现还需单独评审硬件副作用。

\textbf{参数}

\begin{longtable}[]{@{}
  >{\raggedright\arraybackslash}p{(\columnwidth - 6\tabcolsep) * \real{0.18}}
  >{\raggedright\arraybackslash}p{(\columnwidth - 6\tabcolsep) * \real{0.24}}
  >{\raggedright\arraybackslash}p{(\columnwidth - 6\tabcolsep) * \real{0.18}}
  >{\raggedright\arraybackslash}p{(\columnwidth - 6\tabcolsep) * \real{0.40}}@{}}
\toprule
名称 & Python 类型 & 默认值 & 含义 \\
\midrule
\endhead
\passthrough{\lstinline!name!} & \passthrough{\lstinline!str!} & 必填 &
SDK 对象、配置项、服务或动作的名称。允许值由调用它的具体接口定义。 对应
C++ 参数 \passthrough{\lstinline!name: const std::string \&!}。
底层为字符串；长度、编码和允许值由具体协议决定。 \\
\passthrough{\lstinline!callback!} &
\passthrough{\lstinline!Callable[[str, str], None]!} & 必填 &
收到数据或状态变化时调用的 Python
回调。回调签名必须与类型提示一致，并应快速返回。 对应 C++ 参数
\passthrough{\lstinline!callback: const unitree::robot::b2::ConfigChangeStatusCallback \&!}。 \\
\bottomrule
\end{longtable}

\textbf{返回值}

\begin{longtable}[]{@{}
  >{\raggedright\arraybackslash}p{(\columnwidth - 4\tabcolsep) * \real{0.18}}
  >{\raggedright\arraybackslash}p{(\columnwidth - 4\tabcolsep) * \real{0.20}}
  >{\raggedright\arraybackslash}p{(\columnwidth - 4\tabcolsep) * \real{0.62}}@{}}
\toprule
位置 & 类型 & 含义 \\
\midrule
\endhead
无 & \passthrough{\lstinline!None!} & 该方法不返回 Python
值；失败可能表现为异常或底层状态变化。 \\
\bottomrule
\end{longtable}

\textbf{对应 C++}

\begin{itemize}
\tightlist
\item
  类：\passthrough{\lstinline!unitree::robot::b2::ConfigClient!}
\item
  签名：\passthrough{\lstinline!SubscribeChangeStatus(const std::string \&, const unitree::robot::b2::ConfigChangeStatusCallback \&)!}
\item
  绑定策略：\passthrough{\lstinline!CALLBACK\_MANUAL!}
\item
  声明位置：\passthrough{\lstinline!include/unitree/robot/b2/config/config\_client.hpp:52!}
\end{itemize}

\textbf{说明}

\begin{itemize}
\tightlist
\item
  示例是局部调用片段；\passthrough{\lstinline!obj!}
  和参数变量表示已经按上层业务校验并准备好的对象。
\item
  该条目可以被编辑器和类型检查器识别，但当前运行时可能在属性查找阶段直接失败。
\item
  回调可能由 SDK
  工作线程触发；回调应快速返回、捕获异常，并通过线程安全队列移交耗时工作。
\end{itemize}

\textbf{用法}

\begin{lstlisting}[language=Python]
from typing import TYPE_CHECKING

if TYPE_CHECKING:
    def planned_subscribe_change_status(obj: ConfigClient, name: str, callback: Callable[[str, str], None]) -> None:
        obj.subscribe_change_status(name=name, callback=callback)
\end{lstlisting}

\begin{quote}
{[}!CAUTION{]} 上面的 \passthrough{\lstinline!TYPE\_CHECKING!}
示例只表达计划调用形态。该分支在普通 Python
运行时为假，不代表当前扩展能执行此接口。
\end{quote}

\hypertarget{unitree-sdk2-cpp-robot-b2-configdelparameter}{%
\subsubsection{\texorpdfstring{\texttt{unitree\_sdk2\_cpp.robot.b2.ConfigDelParameter}}{unitree\_sdk2\_cpp.robot.b2.ConfigDelParameter}}\label{unitree-sdk2-cpp-robot-b2-configdelparameter}}

SDK 请求参数值类型；当前可能仅提供设计期签名。

\textbf{导入}

\begin{lstlisting}[language=Python]
from unitree_sdk2_cpp.robot.b2 import ConfigDelParameter
\end{lstlisting}

\textbf{基类}：\passthrough{\lstinline!object!}

\textbf{公开属性}

\begin{longtable}[]{@{}
  >{\raggedright\arraybackslash}p{(\columnwidth - 6\tabcolsep) * \real{0.18}}
  >{\raggedright\arraybackslash}p{(\columnwidth - 6\tabcolsep) * \real{0.24}}
  >{\raggedright\arraybackslash}p{(\columnwidth - 6\tabcolsep) * \real{0.18}}
  >{\raggedright\arraybackslash}p{(\columnwidth - 6\tabcolsep) * \real{0.40}}@{}}
\toprule
名称 & Python 类型 & 含义 & 用法 \\
\midrule
\endhead
\passthrough{\lstinline!name!} & \passthrough{\lstinline!str!} & SDK
对象、配置项、服务或动作的名称。允许值由调用它的具体接口定义。 &
\passthrough{\lstinline!value = obj.name!} /
\passthrough{\lstinline!obj.name = value!} \\
\bottomrule
\end{longtable}

\textbf{方法索引}

\begin{longtable}[]{@{}
  >{\raggedright\arraybackslash}p{(\columnwidth - 6\tabcolsep) * \real{0.18}}
  >{\raggedright\arraybackslash}p{(\columnwidth - 6\tabcolsep) * \real{0.24}}
  >{\raggedright\arraybackslash}p{(\columnwidth - 6\tabcolsep) * \real{0.18}}
  >{\raggedright\arraybackslash}p{(\columnwidth - 6\tabcolsep) * \real{0.40}}@{}}
\toprule
方法 & Python 签名 & 状态 & 安全分类 \\
\midrule
\endhead
\protect\hyperlink{unitree-sdk2-cpp-robot-b2-configdelparameter-dunder-init-1}{\passthrough{\lstinline!\_\_init\_\_!}}
& \passthrough{\lstinline!def \_\_init\_\_(self) -> None!} &
\passthrough{\lstinline!SIGNATURE\_ONLY!} &
\passthrough{\lstinline!CONSTRUCTION!} \\
\protect\hyperlink{unitree-sdk2-cpp-robot-b2-configdelparameter-from-json-1}{\passthrough{\lstinline!from\_json!}}
&
\passthrough{\lstinline!def from\_json(self, json: dict[str, Any]) -> None!}
& \passthrough{\lstinline!SIGNATURE\_ONLY!} &
\passthrough{\lstinline!UNCLASSIFIED!} \\
\protect\hyperlink{unitree-sdk2-cpp-robot-b2-configdelparameter-to-json-1}{\passthrough{\lstinline!to\_json!}}
&
\passthrough{\lstinline!def to\_json(self, json: dict[str, Any]) -> None!}
& \passthrough{\lstinline!SIGNATURE\_ONLY!} &
\passthrough{\lstinline!UNCLASSIFIED!} \\
\bottomrule
\end{longtable}

\hypertarget{unitree-sdk2-cpp-robot-b2-configdelparameter-dunder-init-1}{%
\paragraph{\texorpdfstring{\texttt{unitree\_sdk2\_cpp.robot.b2.ConfigDelParameter.\_\_init\_\_}}{unitree\_sdk2\_cpp.robot.b2.ConfigDelParameter.\_\_init\_\_}}\label{unitree-sdk2-cpp-robot-b2-configdelparameter-dunder-init-1}}

初始化 \passthrough{\lstinline!ConfigDelParameter!}
实例。是否能实际构造取决于下方可用性状态。

\textbf{签名}

\begin{lstlisting}[language=Python]
def __init__(self) -> None
\end{lstlisting}

\textbf{可用性与安全性}

\textbf{\passthrough{\lstinline!SIGNATURE\_ONLY!} /
\passthrough{\lstinline!CONSTRUCTION!}}：当前只有设计期签名，不能假设已存在于运行时扩展。

\textbf{参数}

无显式参数。实例方法中的 \passthrough{\lstinline!self!} 由 Python
自动传入。

\textbf{返回值}

\begin{longtable}[]{@{}
  >{\raggedright\arraybackslash}p{(\columnwidth - 4\tabcolsep) * \real{0.18}}
  >{\raggedright\arraybackslash}p{(\columnwidth - 4\tabcolsep) * \real{0.20}}
  >{\raggedright\arraybackslash}p{(\columnwidth - 4\tabcolsep) * \real{0.62}}@{}}
\toprule
位置 & 类型 & 含义 \\
\midrule
\endhead
无 & \passthrough{\lstinline!None!} &
\passthrough{\lstinline!\_\_init\_\_!}
本身不返回值；调用类对象时，在构造函数可用的前提下得到该类实例。 \\
\bottomrule
\end{longtable}

\textbf{对应 C++}

\begin{itemize}
\tightlist
\item
  类：\passthrough{\lstinline!unitree::robot::b2::ConfigDelParameter!}
\item
  签名：\passthrough{\lstinline!ConfigDelParameter()!}
\item
  绑定策略：\passthrough{\lstinline!未单独标注!}
\item
  声明位置：\passthrough{\lstinline!include/unitree/robot/b2/config/config\_api.hpp:140!}
\end{itemize}

\textbf{说明}

\begin{itemize}
\tightlist
\item
  示例是局部调用片段；\passthrough{\lstinline!obj!}
  和参数变量表示已经按上层业务校验并准备好的对象。
\item
  该条目可以被编辑器和类型检查器识别，但当前运行时可能在属性查找阶段直接失败。
\end{itemize}

\textbf{用法}

\begin{lstlisting}[language=Python]
from typing import TYPE_CHECKING

if TYPE_CHECKING:
    def planned_construct() -> ConfigDelParameter:
        return ConfigDelParameter()
\end{lstlisting}

\begin{quote}
{[}!CAUTION{]} 上面的 \passthrough{\lstinline!TYPE\_CHECKING!}
示例只表达计划调用形态。该分支在普通 Python
运行时为假，不代表当前扩展能执行此接口。
\end{quote}

\hypertarget{unitree-sdk2-cpp-robot-b2-configdelparameter-from-json-1}{%
\paragraph{\texorpdfstring{\texttt{unitree\_sdk2\_cpp.robot.b2.ConfigDelParameter.from\_json}}{unitree\_sdk2\_cpp.robot.b2.ConfigDelParameter.from\_json}}\label{unitree-sdk2-cpp-robot-b2-configdelparameter-from-json-1}}

计划从 JSON 风格字典读取字段并更新当前 SDK 值对象。

\textbf{签名}

\begin{lstlisting}[language=Python]
def from_json(self, json: dict[str, Any]) -> None
\end{lstlisting}

\textbf{可用性与安全性}

\textbf{\passthrough{\lstinline!SIGNATURE\_ONLY!} /
\passthrough{\lstinline!UNCLASSIFIED!}}：当前只有设计期签名，不能假设已存在于运行时扩展。

\textbf{参数}

\begin{longtable}[]{@{}
  >{\raggedright\arraybackslash}p{(\columnwidth - 6\tabcolsep) * \real{0.18}}
  >{\raggedright\arraybackslash}p{(\columnwidth - 6\tabcolsep) * \real{0.24}}
  >{\raggedright\arraybackslash}p{(\columnwidth - 6\tabcolsep) * \real{0.18}}
  >{\raggedright\arraybackslash}p{(\columnwidth - 6\tabcolsep) * \real{0.40}}@{}}
\toprule
名称 & Python 类型 & 默认值 & 含义 \\
\midrule
\endhead
\passthrough{\lstinline!json!} &
\passthrough{\lstinline!dict[str, Any]!} & 必填 & JSON
风格字典，键和值必须满足对应 C++ \passthrough{\lstinline!JsonMap!}
数据结构的约束。 对应 C++ 参数
\passthrough{\lstinline!json: common::JsonMap \&!}。 \\
\bottomrule
\end{longtable}

\textbf{返回值}

\begin{longtable}[]{@{}
  >{\raggedright\arraybackslash}p{(\columnwidth - 4\tabcolsep) * \real{0.18}}
  >{\raggedright\arraybackslash}p{(\columnwidth - 4\tabcolsep) * \real{0.20}}
  >{\raggedright\arraybackslash}p{(\columnwidth - 4\tabcolsep) * \real{0.62}}@{}}
\toprule
位置 & 类型 & 含义 \\
\midrule
\endhead
无 & \passthrough{\lstinline!None!} & 该方法不返回 Python
值；失败可能表现为异常或底层状态变化。 \\
\bottomrule
\end{longtable}

\textbf{对应 C++}

\begin{itemize}
\tightlist
\item
  类：\passthrough{\lstinline!unitree::robot::b2::ConfigDelParameter!}
\item
  签名：\passthrough{\lstinline!fromJson(common::JsonMap \&)!}
\item
  绑定策略：\passthrough{\lstinline!SIGNATURE\_PREVIEW!}
\item
  声明位置：\passthrough{\lstinline!include/unitree/robot/b2/config/config\_api.hpp:146!}
\end{itemize}

\textbf{说明}

\begin{itemize}
\tightlist
\item
  示例是局部调用片段；\passthrough{\lstinline!obj!}
  和参数变量表示已经按上层业务校验并准备好的对象。
\item
  该条目可以被编辑器和类型检查器识别，但当前运行时可能在属性查找阶段直接失败。
\end{itemize}

\textbf{用法}

\begin{lstlisting}[language=Python]
from typing import TYPE_CHECKING

if TYPE_CHECKING:
    def planned_from_json(obj: ConfigDelParameter, json: dict[str, Any]) -> None:
        obj.from_json(json=json)
\end{lstlisting}

\begin{quote}
{[}!CAUTION{]} 上面的 \passthrough{\lstinline!TYPE\_CHECKING!}
示例只表达计划调用形态。该分支在普通 Python
运行时为假，不代表当前扩展能执行此接口。
\end{quote}

\hypertarget{unitree-sdk2-cpp-robot-b2-configdelparameter-to-json-1}{%
\paragraph{\texorpdfstring{\texttt{unitree\_sdk2\_cpp.robot.b2.ConfigDelParameter.to\_json}}{unitree\_sdk2\_cpp.robot.b2.ConfigDelParameter.to\_json}}\label{unitree-sdk2-cpp-robot-b2-configdelparameter-to-json-1}}

计划把当前 SDK 值对象写入 JSON 风格字典。

\textbf{签名}

\begin{lstlisting}[language=Python]
def to_json(self, json: dict[str, Any]) -> None
\end{lstlisting}

\textbf{可用性与安全性}

\textbf{\passthrough{\lstinline!SIGNATURE\_ONLY!} /
\passthrough{\lstinline!UNCLASSIFIED!}}：当前只有设计期签名，不能假设已存在于运行时扩展。

\textbf{参数}

\begin{longtable}[]{@{}
  >{\raggedright\arraybackslash}p{(\columnwidth - 6\tabcolsep) * \real{0.18}}
  >{\raggedright\arraybackslash}p{(\columnwidth - 6\tabcolsep) * \real{0.24}}
  >{\raggedright\arraybackslash}p{(\columnwidth - 6\tabcolsep) * \real{0.18}}
  >{\raggedright\arraybackslash}p{(\columnwidth - 6\tabcolsep) * \real{0.40}}@{}}
\toprule
名称 & Python 类型 & 默认值 & 含义 \\
\midrule
\endhead
\passthrough{\lstinline!json!} &
\passthrough{\lstinline!dict[str, Any]!} & 必填 & JSON
风格字典，键和值必须满足对应 C++ \passthrough{\lstinline!JsonMap!}
数据结构的约束。 对应 C++ 参数
\passthrough{\lstinline!json: common::JsonMap \&!}。 \\
\bottomrule
\end{longtable}

\textbf{返回值}

\begin{longtable}[]{@{}
  >{\raggedright\arraybackslash}p{(\columnwidth - 4\tabcolsep) * \real{0.18}}
  >{\raggedright\arraybackslash}p{(\columnwidth - 4\tabcolsep) * \real{0.20}}
  >{\raggedright\arraybackslash}p{(\columnwidth - 4\tabcolsep) * \real{0.62}}@{}}
\toprule
位置 & 类型 & 含义 \\
\midrule
\endhead
无 & \passthrough{\lstinline!None!} & 该方法不返回 Python
值；失败可能表现为异常或底层状态变化。 \\
\bottomrule
\end{longtable}

\textbf{对应 C++}

\begin{itemize}
\tightlist
\item
  类：\passthrough{\lstinline!unitree::robot::b2::ConfigDelParameter!}
\item
  签名：\passthrough{\lstinline!toJson(common::JsonMap \&) const!}
\item
  绑定策略：\passthrough{\lstinline!SIGNATURE\_PREVIEW!}
\item
  声明位置：\passthrough{\lstinline!include/unitree/robot/b2/config/config\_api.hpp:151!}
\end{itemize}

\textbf{说明}

\begin{itemize}
\tightlist
\item
  示例是局部调用片段；\passthrough{\lstinline!obj!}
  和参数变量表示已经按上层业务校验并准备好的对象。
\item
  该条目可以被编辑器和类型检查器识别，但当前运行时可能在属性查找阶段直接失败。
\end{itemize}

\textbf{用法}

\begin{lstlisting}[language=Python]
from typing import TYPE_CHECKING

if TYPE_CHECKING:
    def planned_to_json(obj: ConfigDelParameter, json: dict[str, Any]) -> None:
        obj.to_json(json=json)
\end{lstlisting}

\begin{quote}
{[}!CAUTION{]} 上面的 \passthrough{\lstinline!TYPE\_CHECKING!}
示例只表达计划调用形态。该分支在普通 Python
运行时为假，不代表当前扩展能执行此接口。
\end{quote}

\hypertarget{unitree-sdk2-cpp-robot-b2-configgetdata}{%
\subsubsection{\texorpdfstring{\texttt{unitree\_sdk2\_cpp.robot.b2.ConfigGetData}}{unitree\_sdk2\_cpp.robot.b2.ConfigGetData}}\label{unitree-sdk2-cpp-robot-b2-configgetdata}}

SDK 数据值类型；公开字段和转换方法见下文。

\textbf{导入}

\begin{lstlisting}[language=Python]
from unitree_sdk2_cpp.robot.b2 import ConfigGetData
\end{lstlisting}

\textbf{基类}：\passthrough{\lstinline!object!}

\textbf{公开属性}

\begin{longtable}[]{@{}
  >{\raggedright\arraybackslash}p{(\columnwidth - 6\tabcolsep) * \real{0.18}}
  >{\raggedright\arraybackslash}p{(\columnwidth - 6\tabcolsep) * \real{0.24}}
  >{\raggedright\arraybackslash}p{(\columnwidth - 6\tabcolsep) * \real{0.18}}
  >{\raggedright\arraybackslash}p{(\columnwidth - 6\tabcolsep) * \real{0.40}}@{}}
\toprule
名称 & Python 类型 & 含义 & 用法 \\
\midrule
\endhead
\passthrough{\lstinline!content!} & \passthrough{\lstinline!str!} &
配置、请求或序列化内容。具体格式由对应服务协议定义。 &
\passthrough{\lstinline!value = obj.content!} /
\passthrough{\lstinline!obj.content = value!} \\
\bottomrule
\end{longtable}

\textbf{方法索引}

\begin{longtable}[]{@{}
  >{\raggedright\arraybackslash}p{(\columnwidth - 6\tabcolsep) * \real{0.18}}
  >{\raggedright\arraybackslash}p{(\columnwidth - 6\tabcolsep) * \real{0.24}}
  >{\raggedright\arraybackslash}p{(\columnwidth - 6\tabcolsep) * \real{0.18}}
  >{\raggedright\arraybackslash}p{(\columnwidth - 6\tabcolsep) * \real{0.40}}@{}}
\toprule
方法 & Python 签名 & 状态 & 安全分类 \\
\midrule
\endhead
\protect\hyperlink{unitree-sdk2-cpp-robot-b2-configgetdata-dunder-init-1}{\passthrough{\lstinline!\_\_init\_\_!}}
& \passthrough{\lstinline!def \_\_init\_\_(self) -> None!} &
\passthrough{\lstinline!SIGNATURE\_ONLY!} &
\passthrough{\lstinline!CONSTRUCTION!} \\
\protect\hyperlink{unitree-sdk2-cpp-robot-b2-configgetdata-from-json-1}{\passthrough{\lstinline!from\_json!}}
&
\passthrough{\lstinline!def from\_json(self, json: dict[str, Any]) -> None!}
& \passthrough{\lstinline!SIGNATURE\_ONLY!} &
\passthrough{\lstinline!UNCLASSIFIED!} \\
\protect\hyperlink{unitree-sdk2-cpp-robot-b2-configgetdata-to-json-1}{\passthrough{\lstinline!to\_json!}}
&
\passthrough{\lstinline!def to\_json(self, json: dict[str, Any]) -> None!}
& \passthrough{\lstinline!SIGNATURE\_ONLY!} &
\passthrough{\lstinline!UNCLASSIFIED!} \\
\bottomrule
\end{longtable}

\hypertarget{unitree-sdk2-cpp-robot-b2-configgetdata-dunder-init-1}{%
\paragraph{\texorpdfstring{\texttt{unitree\_sdk2\_cpp.robot.b2.ConfigGetData.\_\_init\_\_}}{unitree\_sdk2\_cpp.robot.b2.ConfigGetData.\_\_init\_\_}}\label{unitree-sdk2-cpp-robot-b2-configgetdata-dunder-init-1}}

初始化 \passthrough{\lstinline!ConfigGetData!}
实例。是否能实际构造取决于下方可用性状态。

\textbf{签名}

\begin{lstlisting}[language=Python]
def __init__(self) -> None
\end{lstlisting}

\textbf{可用性与安全性}

\textbf{\passthrough{\lstinline!SIGNATURE\_ONLY!} /
\passthrough{\lstinline!CONSTRUCTION!}}：当前只有设计期签名，不能假设已存在于运行时扩展。

\textbf{参数}

无显式参数。实例方法中的 \passthrough{\lstinline!self!} 由 Python
自动传入。

\textbf{返回值}

\begin{longtable}[]{@{}
  >{\raggedright\arraybackslash}p{(\columnwidth - 4\tabcolsep) * \real{0.18}}
  >{\raggedright\arraybackslash}p{(\columnwidth - 4\tabcolsep) * \real{0.20}}
  >{\raggedright\arraybackslash}p{(\columnwidth - 4\tabcolsep) * \real{0.62}}@{}}
\toprule
位置 & 类型 & 含义 \\
\midrule
\endhead
无 & \passthrough{\lstinline!None!} &
\passthrough{\lstinline!\_\_init\_\_!}
本身不返回值；调用类对象时，在构造函数可用的前提下得到该类实例。 \\
\bottomrule
\end{longtable}

\textbf{对应 C++}

\begin{itemize}
\tightlist
\item
  类：\passthrough{\lstinline!unitree::robot::b2::ConfigGetData!}
\item
  签名：\passthrough{\lstinline!ConfigGetData()!}
\item
  绑定策略：\passthrough{\lstinline!未单独标注!}
\item
  声明位置：\passthrough{\lstinline!include/unitree/robot/b2/config/config\_api.hpp:117!}
\end{itemize}

\textbf{说明}

\begin{itemize}
\tightlist
\item
  示例是局部调用片段；\passthrough{\lstinline!obj!}
  和参数变量表示已经按上层业务校验并准备好的对象。
\item
  该条目可以被编辑器和类型检查器识别，但当前运行时可能在属性查找阶段直接失败。
\end{itemize}

\textbf{用法}

\begin{lstlisting}[language=Python]
from typing import TYPE_CHECKING

if TYPE_CHECKING:
    def planned_construct() -> ConfigGetData:
        return ConfigGetData()
\end{lstlisting}

\begin{quote}
{[}!CAUTION{]} 上面的 \passthrough{\lstinline!TYPE\_CHECKING!}
示例只表达计划调用形态。该分支在普通 Python
运行时为假，不代表当前扩展能执行此接口。
\end{quote}

\hypertarget{unitree-sdk2-cpp-robot-b2-configgetdata-from-json-1}{%
\paragraph{\texorpdfstring{\texttt{unitree\_sdk2\_cpp.robot.b2.ConfigGetData.from\_json}}{unitree\_sdk2\_cpp.robot.b2.ConfigGetData.from\_json}}\label{unitree-sdk2-cpp-robot-b2-configgetdata-from-json-1}}

计划从 JSON 风格字典读取字段并更新当前 SDK 值对象。

\textbf{签名}

\begin{lstlisting}[language=Python]
def from_json(self, json: dict[str, Any]) -> None
\end{lstlisting}

\textbf{可用性与安全性}

\textbf{\passthrough{\lstinline!SIGNATURE\_ONLY!} /
\passthrough{\lstinline!UNCLASSIFIED!}}：当前只有设计期签名，不能假设已存在于运行时扩展。

\textbf{参数}

\begin{longtable}[]{@{}
  >{\raggedright\arraybackslash}p{(\columnwidth - 6\tabcolsep) * \real{0.18}}
  >{\raggedright\arraybackslash}p{(\columnwidth - 6\tabcolsep) * \real{0.24}}
  >{\raggedright\arraybackslash}p{(\columnwidth - 6\tabcolsep) * \real{0.18}}
  >{\raggedright\arraybackslash}p{(\columnwidth - 6\tabcolsep) * \real{0.40}}@{}}
\toprule
名称 & Python 类型 & 默认值 & 含义 \\
\midrule
\endhead
\passthrough{\lstinline!json!} &
\passthrough{\lstinline!dict[str, Any]!} & 必填 & JSON
风格字典，键和值必须满足对应 C++ \passthrough{\lstinline!JsonMap!}
数据结构的约束。 对应 C++ 参数
\passthrough{\lstinline!json: common::JsonMap \&!}。 \\
\bottomrule
\end{longtable}

\textbf{返回值}

\begin{longtable}[]{@{}
  >{\raggedright\arraybackslash}p{(\columnwidth - 4\tabcolsep) * \real{0.18}}
  >{\raggedright\arraybackslash}p{(\columnwidth - 4\tabcolsep) * \real{0.20}}
  >{\raggedright\arraybackslash}p{(\columnwidth - 4\tabcolsep) * \real{0.62}}@{}}
\toprule
位置 & 类型 & 含义 \\
\midrule
\endhead
无 & \passthrough{\lstinline!None!} & 该方法不返回 Python
值；失败可能表现为异常或底层状态变化。 \\
\bottomrule
\end{longtable}

\textbf{对应 C++}

\begin{itemize}
\tightlist
\item
  类：\passthrough{\lstinline!unitree::robot::b2::ConfigGetData!}
\item
  签名：\passthrough{\lstinline!fromJson(common::JsonMap \&)!}
\item
  绑定策略：\passthrough{\lstinline!SIGNATURE\_PREVIEW!}
\item
  声明位置：\passthrough{\lstinline!include/unitree/robot/b2/config/config\_api.hpp:123!}
\end{itemize}

\textbf{说明}

\begin{itemize}
\tightlist
\item
  示例是局部调用片段；\passthrough{\lstinline!obj!}
  和参数变量表示已经按上层业务校验并准备好的对象。
\item
  该条目可以被编辑器和类型检查器识别，但当前运行时可能在属性查找阶段直接失败。
\end{itemize}

\textbf{用法}

\begin{lstlisting}[language=Python]
from typing import TYPE_CHECKING

if TYPE_CHECKING:
    def planned_from_json(obj: ConfigGetData, json: dict[str, Any]) -> None:
        obj.from_json(json=json)
\end{lstlisting}

\begin{quote}
{[}!CAUTION{]} 上面的 \passthrough{\lstinline!TYPE\_CHECKING!}
示例只表达计划调用形态。该分支在普通 Python
运行时为假，不代表当前扩展能执行此接口。
\end{quote}

\hypertarget{unitree-sdk2-cpp-robot-b2-configgetdata-to-json-1}{%
\paragraph{\texorpdfstring{\texttt{unitree\_sdk2\_cpp.robot.b2.ConfigGetData.to\_json}}{unitree\_sdk2\_cpp.robot.b2.ConfigGetData.to\_json}}\label{unitree-sdk2-cpp-robot-b2-configgetdata-to-json-1}}

计划把当前 SDK 值对象写入 JSON 风格字典。

\textbf{签名}

\begin{lstlisting}[language=Python]
def to_json(self, json: dict[str, Any]) -> None
\end{lstlisting}

\textbf{可用性与安全性}

\textbf{\passthrough{\lstinline!SIGNATURE\_ONLY!} /
\passthrough{\lstinline!UNCLASSIFIED!}}：当前只有设计期签名，不能假设已存在于运行时扩展。

\textbf{参数}

\begin{longtable}[]{@{}
  >{\raggedright\arraybackslash}p{(\columnwidth - 6\tabcolsep) * \real{0.18}}
  >{\raggedright\arraybackslash}p{(\columnwidth - 6\tabcolsep) * \real{0.24}}
  >{\raggedright\arraybackslash}p{(\columnwidth - 6\tabcolsep) * \real{0.18}}
  >{\raggedright\arraybackslash}p{(\columnwidth - 6\tabcolsep) * \real{0.40}}@{}}
\toprule
名称 & Python 类型 & 默认值 & 含义 \\
\midrule
\endhead
\passthrough{\lstinline!json!} &
\passthrough{\lstinline!dict[str, Any]!} & 必填 & JSON
风格字典，键和值必须满足对应 C++ \passthrough{\lstinline!JsonMap!}
数据结构的约束。 对应 C++ 参数
\passthrough{\lstinline!json: common::JsonMap \&!}。 \\
\bottomrule
\end{longtable}

\textbf{返回值}

\begin{longtable}[]{@{}
  >{\raggedright\arraybackslash}p{(\columnwidth - 4\tabcolsep) * \real{0.18}}
  >{\raggedright\arraybackslash}p{(\columnwidth - 4\tabcolsep) * \real{0.20}}
  >{\raggedright\arraybackslash}p{(\columnwidth - 4\tabcolsep) * \real{0.62}}@{}}
\toprule
位置 & 类型 & 含义 \\
\midrule
\endhead
无 & \passthrough{\lstinline!None!} & 该方法不返回 Python
值；失败可能表现为异常或底层状态变化。 \\
\bottomrule
\end{longtable}

\textbf{对应 C++}

\begin{itemize}
\tightlist
\item
  类：\passthrough{\lstinline!unitree::robot::b2::ConfigGetData!}
\item
  签名：\passthrough{\lstinline!toJson(common::JsonMap \&) const!}
\item
  绑定策略：\passthrough{\lstinline!SIGNATURE\_PREVIEW!}
\item
  声明位置：\passthrough{\lstinline!include/unitree/robot/b2/config/config\_api.hpp:128!}
\end{itemize}

\textbf{说明}

\begin{itemize}
\tightlist
\item
  示例是局部调用片段；\passthrough{\lstinline!obj!}
  和参数变量表示已经按上层业务校验并准备好的对象。
\item
  该条目可以被编辑器和类型检查器识别，但当前运行时可能在属性查找阶段直接失败。
\end{itemize}

\textbf{用法}

\begin{lstlisting}[language=Python]
from typing import TYPE_CHECKING

if TYPE_CHECKING:
    def planned_to_json(obj: ConfigGetData, json: dict[str, Any]) -> None:
        obj.to_json(json=json)
\end{lstlisting}

\begin{quote}
{[}!CAUTION{]} 上面的 \passthrough{\lstinline!TYPE\_CHECKING!}
示例只表达计划调用形态。该分支在普通 Python
运行时为假，不代表当前扩展能执行此接口。
\end{quote}

\hypertarget{unitree-sdk2-cpp-robot-b2-configgetparameter}{%
\subsubsection{\texorpdfstring{\texttt{unitree\_sdk2\_cpp.robot.b2.ConfigGetParameter}}{unitree\_sdk2\_cpp.robot.b2.ConfigGetParameter}}\label{unitree-sdk2-cpp-robot-b2-configgetparameter}}

SDK 请求参数值类型；当前可能仅提供设计期签名。

\textbf{导入}

\begin{lstlisting}[language=Python]
from unitree_sdk2_cpp.robot.b2 import ConfigGetParameter
\end{lstlisting}

\textbf{基类}：\passthrough{\lstinline!object!}

\textbf{公开属性}

\begin{longtable}[]{@{}
  >{\raggedright\arraybackslash}p{(\columnwidth - 6\tabcolsep) * \real{0.18}}
  >{\raggedright\arraybackslash}p{(\columnwidth - 6\tabcolsep) * \real{0.24}}
  >{\raggedright\arraybackslash}p{(\columnwidth - 6\tabcolsep) * \real{0.18}}
  >{\raggedright\arraybackslash}p{(\columnwidth - 6\tabcolsep) * \real{0.40}}@{}}
\toprule
名称 & Python 类型 & 含义 & 用法 \\
\midrule
\endhead
\passthrough{\lstinline!name!} & \passthrough{\lstinline!str!} & SDK
对象、配置项、服务或动作的名称。允许值由调用它的具体接口定义。 &
\passthrough{\lstinline!value = obj.name!} /
\passthrough{\lstinline!obj.name = value!} \\
\bottomrule
\end{longtable}

\textbf{方法索引}

\begin{longtable}[]{@{}
  >{\raggedright\arraybackslash}p{(\columnwidth - 6\tabcolsep) * \real{0.18}}
  >{\raggedright\arraybackslash}p{(\columnwidth - 6\tabcolsep) * \real{0.24}}
  >{\raggedright\arraybackslash}p{(\columnwidth - 6\tabcolsep) * \real{0.18}}
  >{\raggedright\arraybackslash}p{(\columnwidth - 6\tabcolsep) * \real{0.40}}@{}}
\toprule
方法 & Python 签名 & 状态 & 安全分类 \\
\midrule
\endhead
\protect\hyperlink{unitree-sdk2-cpp-robot-b2-configgetparameter-dunder-init-1}{\passthrough{\lstinline!\_\_init\_\_!}}
& \passthrough{\lstinline!def \_\_init\_\_(self) -> None!} &
\passthrough{\lstinline!SIGNATURE\_ONLY!} &
\passthrough{\lstinline!CONSTRUCTION!} \\
\protect\hyperlink{unitree-sdk2-cpp-robot-b2-configgetparameter-from-json-1}{\passthrough{\lstinline!from\_json!}}
&
\passthrough{\lstinline!def from\_json(self, json: dict[str, Any]) -> None!}
& \passthrough{\lstinline!SIGNATURE\_ONLY!} &
\passthrough{\lstinline!UNCLASSIFIED!} \\
\protect\hyperlink{unitree-sdk2-cpp-robot-b2-configgetparameter-to-json-1}{\passthrough{\lstinline!to\_json!}}
&
\passthrough{\lstinline!def to\_json(self, json: dict[str, Any]) -> None!}
& \passthrough{\lstinline!SIGNATURE\_ONLY!} &
\passthrough{\lstinline!UNCLASSIFIED!} \\
\bottomrule
\end{longtable}

\hypertarget{unitree-sdk2-cpp-robot-b2-configgetparameter-dunder-init-1}{%
\paragraph{\texorpdfstring{\texttt{unitree\_sdk2\_cpp.robot.b2.ConfigGetParameter.\_\_init\_\_}}{unitree\_sdk2\_cpp.robot.b2.ConfigGetParameter.\_\_init\_\_}}\label{unitree-sdk2-cpp-robot-b2-configgetparameter-dunder-init-1}}

初始化 \passthrough{\lstinline!ConfigGetParameter!}
实例。是否能实际构造取决于下方可用性状态。

\textbf{签名}

\begin{lstlisting}[language=Python]
def __init__(self) -> None
\end{lstlisting}

\textbf{可用性与安全性}

\textbf{\passthrough{\lstinline!SIGNATURE\_ONLY!} /
\passthrough{\lstinline!CONSTRUCTION!}}：当前只有设计期签名，不能假设已存在于运行时扩展。

\textbf{参数}

无显式参数。实例方法中的 \passthrough{\lstinline!self!} 由 Python
自动传入。

\textbf{返回值}

\begin{longtable}[]{@{}
  >{\raggedright\arraybackslash}p{(\columnwidth - 4\tabcolsep) * \real{0.18}}
  >{\raggedright\arraybackslash}p{(\columnwidth - 4\tabcolsep) * \real{0.20}}
  >{\raggedright\arraybackslash}p{(\columnwidth - 4\tabcolsep) * \real{0.62}}@{}}
\toprule
位置 & 类型 & 含义 \\
\midrule
\endhead
无 & \passthrough{\lstinline!None!} &
\passthrough{\lstinline!\_\_init\_\_!}
本身不返回值；调用类对象时，在构造函数可用的前提下得到该类实例。 \\
\bottomrule
\end{longtable}

\textbf{对应 C++}

\begin{itemize}
\tightlist
\item
  类：\passthrough{\lstinline!unitree::robot::b2::ConfigGetParameter!}
\item
  签名：\passthrough{\lstinline!ConfigGetParameter()!}
\item
  绑定策略：\passthrough{\lstinline!未单独标注!}
\item
  声明位置：\passthrough{\lstinline!include/unitree/robot/b2/config/config\_api.hpp:94!}
\end{itemize}

\textbf{说明}

\begin{itemize}
\tightlist
\item
  示例是局部调用片段；\passthrough{\lstinline!obj!}
  和参数变量表示已经按上层业务校验并准备好的对象。
\item
  该条目可以被编辑器和类型检查器识别，但当前运行时可能在属性查找阶段直接失败。
\end{itemize}

\textbf{用法}

\begin{lstlisting}[language=Python]
from typing import TYPE_CHECKING

if TYPE_CHECKING:
    def planned_construct() -> ConfigGetParameter:
        return ConfigGetParameter()
\end{lstlisting}

\begin{quote}
{[}!CAUTION{]} 上面的 \passthrough{\lstinline!TYPE\_CHECKING!}
示例只表达计划调用形态。该分支在普通 Python
运行时为假，不代表当前扩展能执行此接口。
\end{quote}

\hypertarget{unitree-sdk2-cpp-robot-b2-configgetparameter-from-json-1}{%
\paragraph{\texorpdfstring{\texttt{unitree\_sdk2\_cpp.robot.b2.ConfigGetParameter.from\_json}}{unitree\_sdk2\_cpp.robot.b2.ConfigGetParameter.from\_json}}\label{unitree-sdk2-cpp-robot-b2-configgetparameter-from-json-1}}

计划从 JSON 风格字典读取字段并更新当前 SDK 值对象。

\textbf{签名}

\begin{lstlisting}[language=Python]
def from_json(self, json: dict[str, Any]) -> None
\end{lstlisting}

\textbf{可用性与安全性}

\textbf{\passthrough{\lstinline!SIGNATURE\_ONLY!} /
\passthrough{\lstinline!UNCLASSIFIED!}}：当前只有设计期签名，不能假设已存在于运行时扩展。

\textbf{参数}

\begin{longtable}[]{@{}
  >{\raggedright\arraybackslash}p{(\columnwidth - 6\tabcolsep) * \real{0.18}}
  >{\raggedright\arraybackslash}p{(\columnwidth - 6\tabcolsep) * \real{0.24}}
  >{\raggedright\arraybackslash}p{(\columnwidth - 6\tabcolsep) * \real{0.18}}
  >{\raggedright\arraybackslash}p{(\columnwidth - 6\tabcolsep) * \real{0.40}}@{}}
\toprule
名称 & Python 类型 & 默认值 & 含义 \\
\midrule
\endhead
\passthrough{\lstinline!json!} &
\passthrough{\lstinline!dict[str, Any]!} & 必填 & JSON
风格字典，键和值必须满足对应 C++ \passthrough{\lstinline!JsonMap!}
数据结构的约束。 对应 C++ 参数
\passthrough{\lstinline!json: common::JsonMap \&!}。 \\
\bottomrule
\end{longtable}

\textbf{返回值}

\begin{longtable}[]{@{}
  >{\raggedright\arraybackslash}p{(\columnwidth - 4\tabcolsep) * \real{0.18}}
  >{\raggedright\arraybackslash}p{(\columnwidth - 4\tabcolsep) * \real{0.20}}
  >{\raggedright\arraybackslash}p{(\columnwidth - 4\tabcolsep) * \real{0.62}}@{}}
\toprule
位置 & 类型 & 含义 \\
\midrule
\endhead
无 & \passthrough{\lstinline!None!} & 该方法不返回 Python
值；失败可能表现为异常或底层状态变化。 \\
\bottomrule
\end{longtable}

\textbf{对应 C++}

\begin{itemize}
\tightlist
\item
  类：\passthrough{\lstinline!unitree::robot::b2::ConfigGetParameter!}
\item
  签名：\passthrough{\lstinline!fromJson(common::JsonMap \&)!}
\item
  绑定策略：\passthrough{\lstinline!SIGNATURE\_PREVIEW!}
\item
  声明位置：\passthrough{\lstinline!include/unitree/robot/b2/config/config\_api.hpp:100!}
\end{itemize}

\textbf{说明}

\begin{itemize}
\tightlist
\item
  示例是局部调用片段；\passthrough{\lstinline!obj!}
  和参数变量表示已经按上层业务校验并准备好的对象。
\item
  该条目可以被编辑器和类型检查器识别，但当前运行时可能在属性查找阶段直接失败。
\end{itemize}

\textbf{用法}

\begin{lstlisting}[language=Python]
from typing import TYPE_CHECKING

if TYPE_CHECKING:
    def planned_from_json(obj: ConfigGetParameter, json: dict[str, Any]) -> None:
        obj.from_json(json=json)
\end{lstlisting}

\begin{quote}
{[}!CAUTION{]} 上面的 \passthrough{\lstinline!TYPE\_CHECKING!}
示例只表达计划调用形态。该分支在普通 Python
运行时为假，不代表当前扩展能执行此接口。
\end{quote}

\hypertarget{unitree-sdk2-cpp-robot-b2-configgetparameter-to-json-1}{%
\paragraph{\texorpdfstring{\texttt{unitree\_sdk2\_cpp.robot.b2.ConfigGetParameter.to\_json}}{unitree\_sdk2\_cpp.robot.b2.ConfigGetParameter.to\_json}}\label{unitree-sdk2-cpp-robot-b2-configgetparameter-to-json-1}}

计划把当前 SDK 值对象写入 JSON 风格字典。

\textbf{签名}

\begin{lstlisting}[language=Python]
def to_json(self, json: dict[str, Any]) -> None
\end{lstlisting}

\textbf{可用性与安全性}

\textbf{\passthrough{\lstinline!SIGNATURE\_ONLY!} /
\passthrough{\lstinline!UNCLASSIFIED!}}：当前只有设计期签名，不能假设已存在于运行时扩展。

\textbf{参数}

\begin{longtable}[]{@{}
  >{\raggedright\arraybackslash}p{(\columnwidth - 6\tabcolsep) * \real{0.18}}
  >{\raggedright\arraybackslash}p{(\columnwidth - 6\tabcolsep) * \real{0.24}}
  >{\raggedright\arraybackslash}p{(\columnwidth - 6\tabcolsep) * \real{0.18}}
  >{\raggedright\arraybackslash}p{(\columnwidth - 6\tabcolsep) * \real{0.40}}@{}}
\toprule
名称 & Python 类型 & 默认值 & 含义 \\
\midrule
\endhead
\passthrough{\lstinline!json!} &
\passthrough{\lstinline!dict[str, Any]!} & 必填 & JSON
风格字典，键和值必须满足对应 C++ \passthrough{\lstinline!JsonMap!}
数据结构的约束。 对应 C++ 参数
\passthrough{\lstinline!json: common::JsonMap \&!}。 \\
\bottomrule
\end{longtable}

\textbf{返回值}

\begin{longtable}[]{@{}
  >{\raggedright\arraybackslash}p{(\columnwidth - 4\tabcolsep) * \real{0.18}}
  >{\raggedright\arraybackslash}p{(\columnwidth - 4\tabcolsep) * \real{0.20}}
  >{\raggedright\arraybackslash}p{(\columnwidth - 4\tabcolsep) * \real{0.62}}@{}}
\toprule
位置 & 类型 & 含义 \\
\midrule
\endhead
无 & \passthrough{\lstinline!None!} & 该方法不返回 Python
值；失败可能表现为异常或底层状态变化。 \\
\bottomrule
\end{longtable}

\textbf{对应 C++}

\begin{itemize}
\tightlist
\item
  类：\passthrough{\lstinline!unitree::robot::b2::ConfigGetParameter!}
\item
  签名：\passthrough{\lstinline!toJson(common::JsonMap \&) const!}
\item
  绑定策略：\passthrough{\lstinline!SIGNATURE\_PREVIEW!}
\item
  声明位置：\passthrough{\lstinline!include/unitree/robot/b2/config/config\_api.hpp:105!}
\end{itemize}

\textbf{说明}

\begin{itemize}
\tightlist
\item
  示例是局部调用片段；\passthrough{\lstinline!obj!}
  和参数变量表示已经按上层业务校验并准备好的对象。
\item
  该条目可以被编辑器和类型检查器识别，但当前运行时可能在属性查找阶段直接失败。
\end{itemize}

\textbf{用法}

\begin{lstlisting}[language=Python]
from typing import TYPE_CHECKING

if TYPE_CHECKING:
    def planned_to_json(obj: ConfigGetParameter, json: dict[str, Any]) -> None:
        obj.to_json(json=json)
\end{lstlisting}

\begin{quote}
{[}!CAUTION{]} 上面的 \passthrough{\lstinline!TYPE\_CHECKING!}
示例只表达计划调用形态。该分支在普通 Python
运行时为假，不代表当前扩展能执行此接口。
\end{quote}

\hypertarget{unitree-sdk2-cpp-robot-b2-configmeta}{%
\subsubsection{\texorpdfstring{\texttt{unitree\_sdk2\_cpp.robot.b2.ConfigMeta}}{unitree\_sdk2\_cpp.robot.b2.ConfigMeta}}\label{unitree-sdk2-cpp-robot-b2-configmeta}}

SDK 数据值类型；公开字段和转换方法见下文。

\textbf{导入}

\begin{lstlisting}[language=Python]
from unitree_sdk2_cpp.robot.b2 import ConfigMeta
\end{lstlisting}

\textbf{基类}：\passthrough{\lstinline!object!}

\textbf{公开属性}

\begin{longtable}[]{@{}
  >{\raggedright\arraybackslash}p{(\columnwidth - 6\tabcolsep) * \real{0.18}}
  >{\raggedright\arraybackslash}p{(\columnwidth - 6\tabcolsep) * \real{0.24}}
  >{\raggedright\arraybackslash}p{(\columnwidth - 6\tabcolsep) * \real{0.18}}
  >{\raggedright\arraybackslash}p{(\columnwidth - 6\tabcolsep) * \real{0.40}}@{}}
\toprule
名称 & Python 类型 & 含义 & 用法 \\
\midrule
\endhead
\passthrough{\lstinline!name!} & \passthrough{\lstinline!str!} & SDK
对象、配置项、服务或动作的名称。允许值由调用它的具体接口定义。 &
\passthrough{\lstinline!value = obj.name!} /
\passthrough{\lstinline!obj.name = value!} \\
\passthrough{\lstinline!last\_modified!} & \passthrough{\lstinline!str!}
& 传给该接口的 \passthrough{\lstinline!last!}
\passthrough{\lstinline!modified!} 参数，Python 类型为
\passthrough{\lstinline!str!}。精确含义、单位、范围和枚举值需查目标型号对应头文件与协议。
& \passthrough{\lstinline!value = obj.last\_modified!} /
\passthrough{\lstinline!obj.last\_modified = value!} \\
\passthrough{\lstinline!size!} & \passthrough{\lstinline!int!} &
传给该接口的 \passthrough{\lstinline!size!} 参数，Python 类型为
\passthrough{\lstinline!int!}。精确含义、单位、范围和枚举值需查目标型号对应头文件与协议。
& \passthrough{\lstinline!value = obj.size!} /
\passthrough{\lstinline!obj.size = value!} \\
\passthrough{\lstinline!epoch!} & \passthrough{\lstinline!int!} &
传给该接口的 \passthrough{\lstinline!epoch!} 参数，Python 类型为
\passthrough{\lstinline!int!}。精确含义、单位、范围和枚举值需查目标型号对应头文件与协议。
& \passthrough{\lstinline!value = obj.epoch!} /
\passthrough{\lstinline!obj.epoch = value!} \\
\bottomrule
\end{longtable}

\textbf{方法索引}

\begin{longtable}[]{@{}
  >{\raggedright\arraybackslash}p{(\columnwidth - 6\tabcolsep) * \real{0.18}}
  >{\raggedright\arraybackslash}p{(\columnwidth - 6\tabcolsep) * \real{0.24}}
  >{\raggedright\arraybackslash}p{(\columnwidth - 6\tabcolsep) * \real{0.18}}
  >{\raggedright\arraybackslash}p{(\columnwidth - 6\tabcolsep) * \real{0.40}}@{}}
\toprule
方法 & Python 签名 & 状态 & 安全分类 \\
\midrule
\endhead
\protect\hyperlink{unitree-sdk2-cpp-robot-b2-configmeta-dunder-init-1}{\passthrough{\lstinline!\_\_init\_\_!}}
& \passthrough{\lstinline!def \_\_init\_\_(self) -> None!} &
\passthrough{\lstinline!SIGNATURE\_ONLY!} &
\passthrough{\lstinline!CONSTRUCTION!} \\
\bottomrule
\end{longtable}

\hypertarget{unitree-sdk2-cpp-robot-b2-configmeta-dunder-init-1}{%
\paragraph{\texorpdfstring{\texttt{unitree\_sdk2\_cpp.robot.b2.ConfigMeta.\_\_init\_\_}}{unitree\_sdk2\_cpp.robot.b2.ConfigMeta.\_\_init\_\_}}\label{unitree-sdk2-cpp-robot-b2-configmeta-dunder-init-1}}

初始化 \passthrough{\lstinline!ConfigMeta!}
实例。是否能实际构造取决于下方可用性状态。

\textbf{签名}

\begin{lstlisting}[language=Python]
def __init__(self) -> None
\end{lstlisting}

\textbf{可用性与安全性}

\textbf{\passthrough{\lstinline!SIGNATURE\_ONLY!} /
\passthrough{\lstinline!CONSTRUCTION!}}：当前只有设计期签名，不能假设已存在于运行时扩展。

\textbf{参数}

无显式参数。实例方法中的 \passthrough{\lstinline!self!} 由 Python
自动传入。

\textbf{返回值}

\begin{longtable}[]{@{}
  >{\raggedright\arraybackslash}p{(\columnwidth - 4\tabcolsep) * \real{0.18}}
  >{\raggedright\arraybackslash}p{(\columnwidth - 4\tabcolsep) * \real{0.20}}
  >{\raggedright\arraybackslash}p{(\columnwidth - 4\tabcolsep) * \real{0.62}}@{}}
\toprule
位置 & 类型 & 含义 \\
\midrule
\endhead
无 & \passthrough{\lstinline!None!} &
\passthrough{\lstinline!\_\_init\_\_!}
本身不返回值；调用类对象时，在构造函数可用的前提下得到该类实例。 \\
\bottomrule
\end{longtable}

\textbf{对应 C++}

\begin{itemize}
\tightlist
\item
  类：\passthrough{\lstinline!unitree::robot::b2::ConfigMeta!}
\item
  签名：\passthrough{\lstinline!ConfigMeta()!}
\item
  绑定策略：\passthrough{\lstinline!未单独标注!}
\item
  声明位置：\passthrough{\lstinline!include/unitree/robot/b2/config/config\_client.hpp:20!}
\end{itemize}

\textbf{说明}

\begin{itemize}
\tightlist
\item
  示例是局部调用片段；\passthrough{\lstinline!obj!}
  和参数变量表示已经按上层业务校验并准备好的对象。
\item
  该条目可以被编辑器和类型检查器识别，但当前运行时可能在属性查找阶段直接失败。
\end{itemize}

\textbf{用法}

\begin{lstlisting}[language=Python]
from typing import TYPE_CHECKING

if TYPE_CHECKING:
    def planned_construct() -> ConfigMeta:
        return ConfigMeta()
\end{lstlisting}

\begin{quote}
{[}!CAUTION{]} 上面的 \passthrough{\lstinline!TYPE\_CHECKING!}
示例只表达计划调用形态。该分支在普通 Python
运行时为假，不代表当前扩展能执行此接口。
\end{quote}

\hypertarget{unitree-sdk2-cpp-robot-b2-configmetadata}{%
\subsubsection{\texorpdfstring{\texttt{unitree\_sdk2\_cpp.robot.b2.ConfigMetaData}}{unitree\_sdk2\_cpp.robot.b2.ConfigMetaData}}\label{unitree-sdk2-cpp-robot-b2-configmetadata}}

SDK 数据值类型；公开字段和转换方法见下文。

\textbf{导入}

\begin{lstlisting}[language=Python]
from unitree_sdk2_cpp.robot.b2 import ConfigMetaData
\end{lstlisting}

\textbf{基类}：\passthrough{\lstinline!object!}

\textbf{公开属性}

\begin{longtable}[]{@{}
  >{\raggedright\arraybackslash}p{(\columnwidth - 6\tabcolsep) * \real{0.18}}
  >{\raggedright\arraybackslash}p{(\columnwidth - 6\tabcolsep) * \real{0.24}}
  >{\raggedright\arraybackslash}p{(\columnwidth - 6\tabcolsep) * \real{0.18}}
  >{\raggedright\arraybackslash}p{(\columnwidth - 6\tabcolsep) * \real{0.40}}@{}}
\toprule
名称 & Python 类型 & 含义 & 用法 \\
\midrule
\endhead
\passthrough{\lstinline!meta!} &
\passthrough{\lstinline!JsonizeConfigMeta!} & 传给该接口的
\passthrough{\lstinline!meta!} 参数，Python 类型为
\passthrough{\lstinline!JsonizeConfigMeta!}。精确含义、单位、范围和枚举值需查目标型号对应头文件与协议。
& \passthrough{\lstinline!value = obj.meta!} /
\passthrough{\lstinline!obj.meta = value!} \\
\bottomrule
\end{longtable}

\textbf{方法索引}

\begin{longtable}[]{@{}
  >{\raggedright\arraybackslash}p{(\columnwidth - 6\tabcolsep) * \real{0.18}}
  >{\raggedright\arraybackslash}p{(\columnwidth - 6\tabcolsep) * \real{0.24}}
  >{\raggedright\arraybackslash}p{(\columnwidth - 6\tabcolsep) * \real{0.18}}
  >{\raggedright\arraybackslash}p{(\columnwidth - 6\tabcolsep) * \real{0.40}}@{}}
\toprule
方法 & Python 签名 & 状态 & 安全分类 \\
\midrule
\endhead
\protect\hyperlink{unitree-sdk2-cpp-robot-b2-configmetadata-dunder-init-1}{\passthrough{\lstinline!\_\_init\_\_!}}
& \passthrough{\lstinline!def \_\_init\_\_(self) -> None!} &
\passthrough{\lstinline!SIGNATURE\_ONLY!} &
\passthrough{\lstinline!CONSTRUCTION!} \\
\protect\hyperlink{unitree-sdk2-cpp-robot-b2-configmetadata-from-json-1}{\passthrough{\lstinline!from\_json!}}
&
\passthrough{\lstinline!def from\_json(self, json: dict[str, Any]) -> None!}
& \passthrough{\lstinline!SIGNATURE\_ONLY!} &
\passthrough{\lstinline!UNCLASSIFIED!} \\
\protect\hyperlink{unitree-sdk2-cpp-robot-b2-configmetadata-to-json-1}{\passthrough{\lstinline!to\_json!}}
&
\passthrough{\lstinline!def to\_json(self, json: dict[str, Any]) -> None!}
& \passthrough{\lstinline!SIGNATURE\_ONLY!} &
\passthrough{\lstinline!UNCLASSIFIED!} \\
\bottomrule
\end{longtable}

\hypertarget{unitree-sdk2-cpp-robot-b2-configmetadata-dunder-init-1}{%
\paragraph{\texorpdfstring{\texttt{unitree\_sdk2\_cpp.robot.b2.ConfigMetaData.\_\_init\_\_}}{unitree\_sdk2\_cpp.robot.b2.ConfigMetaData.\_\_init\_\_}}\label{unitree-sdk2-cpp-robot-b2-configmetadata-dunder-init-1}}

初始化 \passthrough{\lstinline!ConfigMetaData!}
实例。是否能实际构造取决于下方可用性状态。

\textbf{签名}

\begin{lstlisting}[language=Python]
def __init__(self) -> None
\end{lstlisting}

\textbf{可用性与安全性}

\textbf{\passthrough{\lstinline!SIGNATURE\_ONLY!} /
\passthrough{\lstinline!CONSTRUCTION!}}：当前只有设计期签名，不能假设已存在于运行时扩展。

\textbf{参数}

无显式参数。实例方法中的 \passthrough{\lstinline!self!} 由 Python
自动传入。

\textbf{返回值}

\begin{longtable}[]{@{}
  >{\raggedright\arraybackslash}p{(\columnwidth - 4\tabcolsep) * \real{0.18}}
  >{\raggedright\arraybackslash}p{(\columnwidth - 4\tabcolsep) * \real{0.20}}
  >{\raggedright\arraybackslash}p{(\columnwidth - 4\tabcolsep) * \real{0.62}}@{}}
\toprule
位置 & 类型 & 含义 \\
\midrule
\endhead
无 & \passthrough{\lstinline!None!} &
\passthrough{\lstinline!\_\_init\_\_!}
本身不返回值；调用类对象时，在构造函数可用的前提下得到该类实例。 \\
\bottomrule
\end{longtable}

\textbf{对应 C++}

\begin{itemize}
\tightlist
\item
  类：\passthrough{\lstinline!unitree::robot::b2::ConfigMetaData!}
\item
  签名：\passthrough{\lstinline!ConfigMetaData()!}
\item
  绑定策略：\passthrough{\lstinline!未单独标注!}
\item
  声明位置：\passthrough{\lstinline!include/unitree/robot/b2/config/config\_api.hpp:186!}
\end{itemize}

\textbf{说明}

\begin{itemize}
\tightlist
\item
  示例是局部调用片段；\passthrough{\lstinline!obj!}
  和参数变量表示已经按上层业务校验并准备好的对象。
\item
  该条目可以被编辑器和类型检查器识别，但当前运行时可能在属性查找阶段直接失败。
\end{itemize}

\textbf{用法}

\begin{lstlisting}[language=Python]
from typing import TYPE_CHECKING

if TYPE_CHECKING:
    def planned_construct() -> ConfigMetaData:
        return ConfigMetaData()
\end{lstlisting}

\begin{quote}
{[}!CAUTION{]} 上面的 \passthrough{\lstinline!TYPE\_CHECKING!}
示例只表达计划调用形态。该分支在普通 Python
运行时为假，不代表当前扩展能执行此接口。
\end{quote}

\hypertarget{unitree-sdk2-cpp-robot-b2-configmetadata-from-json-1}{%
\paragraph{\texorpdfstring{\texttt{unitree\_sdk2\_cpp.robot.b2.ConfigMetaData.from\_json}}{unitree\_sdk2\_cpp.robot.b2.ConfigMetaData.from\_json}}\label{unitree-sdk2-cpp-robot-b2-configmetadata-from-json-1}}

计划从 JSON 风格字典读取字段并更新当前 SDK 值对象。

\textbf{签名}

\begin{lstlisting}[language=Python]
def from_json(self, json: dict[str, Any]) -> None
\end{lstlisting}

\textbf{可用性与安全性}

\textbf{\passthrough{\lstinline!SIGNATURE\_ONLY!} /
\passthrough{\lstinline!UNCLASSIFIED!}}：当前只有设计期签名，不能假设已存在于运行时扩展。

\textbf{参数}

\begin{longtable}[]{@{}
  >{\raggedright\arraybackslash}p{(\columnwidth - 6\tabcolsep) * \real{0.18}}
  >{\raggedright\arraybackslash}p{(\columnwidth - 6\tabcolsep) * \real{0.24}}
  >{\raggedright\arraybackslash}p{(\columnwidth - 6\tabcolsep) * \real{0.18}}
  >{\raggedright\arraybackslash}p{(\columnwidth - 6\tabcolsep) * \real{0.40}}@{}}
\toprule
名称 & Python 类型 & 默认值 & 含义 \\
\midrule
\endhead
\passthrough{\lstinline!json!} &
\passthrough{\lstinline!dict[str, Any]!} & 必填 & JSON
风格字典，键和值必须满足对应 C++ \passthrough{\lstinline!JsonMap!}
数据结构的约束。 对应 C++ 参数
\passthrough{\lstinline!json: common::JsonMap \&!}。 \\
\bottomrule
\end{longtable}

\textbf{返回值}

\begin{longtable}[]{@{}
  >{\raggedright\arraybackslash}p{(\columnwidth - 4\tabcolsep) * \real{0.18}}
  >{\raggedright\arraybackslash}p{(\columnwidth - 4\tabcolsep) * \real{0.20}}
  >{\raggedright\arraybackslash}p{(\columnwidth - 4\tabcolsep) * \real{0.62}}@{}}
\toprule
位置 & 类型 & 含义 \\
\midrule
\endhead
无 & \passthrough{\lstinline!None!} & 该方法不返回 Python
值；失败可能表现为异常或底层状态变化。 \\
\bottomrule
\end{longtable}

\textbf{对应 C++}

\begin{itemize}
\tightlist
\item
  类：\passthrough{\lstinline!unitree::robot::b2::ConfigMetaData!}
\item
  签名：\passthrough{\lstinline!fromJson(common::JsonMap \&)!}
\item
  绑定策略：\passthrough{\lstinline!SIGNATURE\_PREVIEW!}
\item
  声明位置：\passthrough{\lstinline!include/unitree/robot/b2/config/config\_api.hpp:192!}
\end{itemize}

\textbf{说明}

\begin{itemize}
\tightlist
\item
  示例是局部调用片段；\passthrough{\lstinline!obj!}
  和参数变量表示已经按上层业务校验并准备好的对象。
\item
  该条目可以被编辑器和类型检查器识别，但当前运行时可能在属性查找阶段直接失败。
\end{itemize}

\textbf{用法}

\begin{lstlisting}[language=Python]
from typing import TYPE_CHECKING

if TYPE_CHECKING:
    def planned_from_json(obj: ConfigMetaData, json: dict[str, Any]) -> None:
        obj.from_json(json=json)
\end{lstlisting}

\begin{quote}
{[}!CAUTION{]} 上面的 \passthrough{\lstinline!TYPE\_CHECKING!}
示例只表达计划调用形态。该分支在普通 Python
运行时为假，不代表当前扩展能执行此接口。
\end{quote}

\hypertarget{unitree-sdk2-cpp-robot-b2-configmetadata-to-json-1}{%
\paragraph{\texorpdfstring{\texttt{unitree\_sdk2\_cpp.robot.b2.ConfigMetaData.to\_json}}{unitree\_sdk2\_cpp.robot.b2.ConfigMetaData.to\_json}}\label{unitree-sdk2-cpp-robot-b2-configmetadata-to-json-1}}

计划把当前 SDK 值对象写入 JSON 风格字典。

\textbf{签名}

\begin{lstlisting}[language=Python]
def to_json(self, json: dict[str, Any]) -> None
\end{lstlisting}

\textbf{可用性与安全性}

\textbf{\passthrough{\lstinline!SIGNATURE\_ONLY!} /
\passthrough{\lstinline!UNCLASSIFIED!}}：当前只有设计期签名，不能假设已存在于运行时扩展。

\textbf{参数}

\begin{longtable}[]{@{}
  >{\raggedright\arraybackslash}p{(\columnwidth - 6\tabcolsep) * \real{0.18}}
  >{\raggedright\arraybackslash}p{(\columnwidth - 6\tabcolsep) * \real{0.24}}
  >{\raggedright\arraybackslash}p{(\columnwidth - 6\tabcolsep) * \real{0.18}}
  >{\raggedright\arraybackslash}p{(\columnwidth - 6\tabcolsep) * \real{0.40}}@{}}
\toprule
名称 & Python 类型 & 默认值 & 含义 \\
\midrule
\endhead
\passthrough{\lstinline!json!} &
\passthrough{\lstinline!dict[str, Any]!} & 必填 & JSON
风格字典，键和值必须满足对应 C++ \passthrough{\lstinline!JsonMap!}
数据结构的约束。 对应 C++ 参数
\passthrough{\lstinline!json: common::JsonMap \&!}。 \\
\bottomrule
\end{longtable}

\textbf{返回值}

\begin{longtable}[]{@{}
  >{\raggedright\arraybackslash}p{(\columnwidth - 4\tabcolsep) * \real{0.18}}
  >{\raggedright\arraybackslash}p{(\columnwidth - 4\tabcolsep) * \real{0.20}}
  >{\raggedright\arraybackslash}p{(\columnwidth - 4\tabcolsep) * \real{0.62}}@{}}
\toprule
位置 & 类型 & 含义 \\
\midrule
\endhead
无 & \passthrough{\lstinline!None!} & 该方法不返回 Python
值；失败可能表现为异常或底层状态变化。 \\
\bottomrule
\end{longtable}

\textbf{对应 C++}

\begin{itemize}
\tightlist
\item
  类：\passthrough{\lstinline!unitree::robot::b2::ConfigMetaData!}
\item
  签名：\passthrough{\lstinline!toJson(common::JsonMap \&) const!}
\item
  绑定策略：\passthrough{\lstinline!SIGNATURE\_PREVIEW!}
\item
  声明位置：\passthrough{\lstinline!include/unitree/robot/b2/config/config\_api.hpp:197!}
\end{itemize}

\textbf{说明}

\begin{itemize}
\tightlist
\item
  示例是局部调用片段；\passthrough{\lstinline!obj!}
  和参数变量表示已经按上层业务校验并准备好的对象。
\item
  该条目可以被编辑器和类型检查器识别，但当前运行时可能在属性查找阶段直接失败。
\end{itemize}

\textbf{用法}

\begin{lstlisting}[language=Python]
from typing import TYPE_CHECKING

if TYPE_CHECKING:
    def planned_to_json(obj: ConfigMetaData, json: dict[str, Any]) -> None:
        obj.to_json(json=json)
\end{lstlisting}

\begin{quote}
{[}!CAUTION{]} 上面的 \passthrough{\lstinline!TYPE\_CHECKING!}
示例只表达计划调用形态。该分支在普通 Python
运行时为假，不代表当前扩展能执行此接口。
\end{quote}

\hypertarget{unitree-sdk2-cpp-robot-b2-configmetaparameter}{%
\subsubsection{\texorpdfstring{\texttt{unitree\_sdk2\_cpp.robot.b2.ConfigMetaParameter}}{unitree\_sdk2\_cpp.robot.b2.ConfigMetaParameter}}\label{unitree-sdk2-cpp-robot-b2-configmetaparameter}}

SDK 请求参数值类型；当前可能仅提供设计期签名。

\textbf{导入}

\begin{lstlisting}[language=Python]
from unitree_sdk2_cpp.robot.b2 import ConfigMetaParameter
\end{lstlisting}

\textbf{基类}：\passthrough{\lstinline!object!}

\textbf{公开属性}

\begin{longtable}[]{@{}
  >{\raggedright\arraybackslash}p{(\columnwidth - 6\tabcolsep) * \real{0.18}}
  >{\raggedright\arraybackslash}p{(\columnwidth - 6\tabcolsep) * \real{0.24}}
  >{\raggedright\arraybackslash}p{(\columnwidth - 6\tabcolsep) * \real{0.18}}
  >{\raggedright\arraybackslash}p{(\columnwidth - 6\tabcolsep) * \real{0.40}}@{}}
\toprule
名称 & Python 类型 & 含义 & 用法 \\
\midrule
\endhead
\passthrough{\lstinline!name!} & \passthrough{\lstinline!str!} & SDK
对象、配置项、服务或动作的名称。允许值由调用它的具体接口定义。 &
\passthrough{\lstinline!value = obj.name!} /
\passthrough{\lstinline!obj.name = value!} \\
\bottomrule
\end{longtable}

\textbf{方法索引}

\begin{longtable}[]{@{}
  >{\raggedright\arraybackslash}p{(\columnwidth - 6\tabcolsep) * \real{0.18}}
  >{\raggedright\arraybackslash}p{(\columnwidth - 6\tabcolsep) * \real{0.24}}
  >{\raggedright\arraybackslash}p{(\columnwidth - 6\tabcolsep) * \real{0.18}}
  >{\raggedright\arraybackslash}p{(\columnwidth - 6\tabcolsep) * \real{0.40}}@{}}
\toprule
方法 & Python 签名 & 状态 & 安全分类 \\
\midrule
\endhead
\protect\hyperlink{unitree-sdk2-cpp-robot-b2-configmetaparameter-dunder-init-1}{\passthrough{\lstinline!\_\_init\_\_!}}
& \passthrough{\lstinline!def \_\_init\_\_(self) -> None!} &
\passthrough{\lstinline!SIGNATURE\_ONLY!} &
\passthrough{\lstinline!CONSTRUCTION!} \\
\protect\hyperlink{unitree-sdk2-cpp-robot-b2-configmetaparameter-from-json-1}{\passthrough{\lstinline!from\_json!}}
&
\passthrough{\lstinline!def from\_json(self, json: dict[str, Any]) -> None!}
& \passthrough{\lstinline!SIGNATURE\_ONLY!} &
\passthrough{\lstinline!UNCLASSIFIED!} \\
\protect\hyperlink{unitree-sdk2-cpp-robot-b2-configmetaparameter-to-json-1}{\passthrough{\lstinline!to\_json!}}
&
\passthrough{\lstinline!def to\_json(self, json: dict[str, Any]) -> None!}
& \passthrough{\lstinline!SIGNATURE\_ONLY!} &
\passthrough{\lstinline!UNCLASSIFIED!} \\
\bottomrule
\end{longtable}

\hypertarget{unitree-sdk2-cpp-robot-b2-configmetaparameter-dunder-init-1}{%
\paragraph{\texorpdfstring{\texttt{unitree\_sdk2\_cpp.robot.b2.ConfigMetaParameter.\_\_init\_\_}}{unitree\_sdk2\_cpp.robot.b2.ConfigMetaParameter.\_\_init\_\_}}\label{unitree-sdk2-cpp-robot-b2-configmetaparameter-dunder-init-1}}

初始化 \passthrough{\lstinline!ConfigMetaParameter!}
实例。是否能实际构造取决于下方可用性状态。

\textbf{签名}

\begin{lstlisting}[language=Python]
def __init__(self) -> None
\end{lstlisting}

\textbf{可用性与安全性}

\textbf{\passthrough{\lstinline!SIGNATURE\_ONLY!} /
\passthrough{\lstinline!CONSTRUCTION!}}：当前只有设计期签名，不能假设已存在于运行时扩展。

\textbf{参数}

无显式参数。实例方法中的 \passthrough{\lstinline!self!} 由 Python
自动传入。

\textbf{返回值}

\begin{longtable}[]{@{}
  >{\raggedright\arraybackslash}p{(\columnwidth - 4\tabcolsep) * \real{0.18}}
  >{\raggedright\arraybackslash}p{(\columnwidth - 4\tabcolsep) * \real{0.20}}
  >{\raggedright\arraybackslash}p{(\columnwidth - 4\tabcolsep) * \real{0.62}}@{}}
\toprule
位置 & 类型 & 含义 \\
\midrule
\endhead
无 & \passthrough{\lstinline!None!} &
\passthrough{\lstinline!\_\_init\_\_!}
本身不返回值；调用类对象时，在构造函数可用的前提下得到该类实例。 \\
\bottomrule
\end{longtable}

\textbf{对应 C++}

\begin{itemize}
\tightlist
\item
  类：\passthrough{\lstinline!unitree::robot::b2::ConfigMetaParameter!}
\item
  签名：\passthrough{\lstinline!ConfigMetaParameter()!}
\item
  绑定策略：\passthrough{\lstinline!未单独标注!}
\item
  声明位置：\passthrough{\lstinline!include/unitree/robot/b2/config/config\_api.hpp:163!}
\end{itemize}

\textbf{说明}

\begin{itemize}
\tightlist
\item
  示例是局部调用片段；\passthrough{\lstinline!obj!}
  和参数变量表示已经按上层业务校验并准备好的对象。
\item
  该条目可以被编辑器和类型检查器识别，但当前运行时可能在属性查找阶段直接失败。
\end{itemize}

\textbf{用法}

\begin{lstlisting}[language=Python]
from typing import TYPE_CHECKING

if TYPE_CHECKING:
    def planned_construct() -> ConfigMetaParameter:
        return ConfigMetaParameter()
\end{lstlisting}

\begin{quote}
{[}!CAUTION{]} 上面的 \passthrough{\lstinline!TYPE\_CHECKING!}
示例只表达计划调用形态。该分支在普通 Python
运行时为假，不代表当前扩展能执行此接口。
\end{quote}

\hypertarget{unitree-sdk2-cpp-robot-b2-configmetaparameter-from-json-1}{%
\paragraph{\texorpdfstring{\texttt{unitree\_sdk2\_cpp.robot.b2.ConfigMetaParameter.from\_json}}{unitree\_sdk2\_cpp.robot.b2.ConfigMetaParameter.from\_json}}\label{unitree-sdk2-cpp-robot-b2-configmetaparameter-from-json-1}}

计划从 JSON 风格字典读取字段并更新当前 SDK 值对象。

\textbf{签名}

\begin{lstlisting}[language=Python]
def from_json(self, json: dict[str, Any]) -> None
\end{lstlisting}

\textbf{可用性与安全性}

\textbf{\passthrough{\lstinline!SIGNATURE\_ONLY!} /
\passthrough{\lstinline!UNCLASSIFIED!}}：当前只有设计期签名，不能假设已存在于运行时扩展。

\textbf{参数}

\begin{longtable}[]{@{}
  >{\raggedright\arraybackslash}p{(\columnwidth - 6\tabcolsep) * \real{0.18}}
  >{\raggedright\arraybackslash}p{(\columnwidth - 6\tabcolsep) * \real{0.24}}
  >{\raggedright\arraybackslash}p{(\columnwidth - 6\tabcolsep) * \real{0.18}}
  >{\raggedright\arraybackslash}p{(\columnwidth - 6\tabcolsep) * \real{0.40}}@{}}
\toprule
名称 & Python 类型 & 默认值 & 含义 \\
\midrule
\endhead
\passthrough{\lstinline!json!} &
\passthrough{\lstinline!dict[str, Any]!} & 必填 & JSON
风格字典，键和值必须满足对应 C++ \passthrough{\lstinline!JsonMap!}
数据结构的约束。 对应 C++ 参数
\passthrough{\lstinline!json: common::JsonMap \&!}。 \\
\bottomrule
\end{longtable}

\textbf{返回值}

\begin{longtable}[]{@{}
  >{\raggedright\arraybackslash}p{(\columnwidth - 4\tabcolsep) * \real{0.18}}
  >{\raggedright\arraybackslash}p{(\columnwidth - 4\tabcolsep) * \real{0.20}}
  >{\raggedright\arraybackslash}p{(\columnwidth - 4\tabcolsep) * \real{0.62}}@{}}
\toprule
位置 & 类型 & 含义 \\
\midrule
\endhead
无 & \passthrough{\lstinline!None!} & 该方法不返回 Python
值；失败可能表现为异常或底层状态变化。 \\
\bottomrule
\end{longtable}

\textbf{对应 C++}

\begin{itemize}
\tightlist
\item
  类：\passthrough{\lstinline!unitree::robot::b2::ConfigMetaParameter!}
\item
  签名：\passthrough{\lstinline!fromJson(common::JsonMap \&)!}
\item
  绑定策略：\passthrough{\lstinline!SIGNATURE\_PREVIEW!}
\item
  声明位置：\passthrough{\lstinline!include/unitree/robot/b2/config/config\_api.hpp:169!}
\end{itemize}

\textbf{说明}

\begin{itemize}
\tightlist
\item
  示例是局部调用片段；\passthrough{\lstinline!obj!}
  和参数变量表示已经按上层业务校验并准备好的对象。
\item
  该条目可以被编辑器和类型检查器识别，但当前运行时可能在属性查找阶段直接失败。
\end{itemize}

\textbf{用法}

\begin{lstlisting}[language=Python]
from typing import TYPE_CHECKING

if TYPE_CHECKING:
    def planned_from_json(obj: ConfigMetaParameter, json: dict[str, Any]) -> None:
        obj.from_json(json=json)
\end{lstlisting}

\begin{quote}
{[}!CAUTION{]} 上面的 \passthrough{\lstinline!TYPE\_CHECKING!}
示例只表达计划调用形态。该分支在普通 Python
运行时为假，不代表当前扩展能执行此接口。
\end{quote}

\hypertarget{unitree-sdk2-cpp-robot-b2-configmetaparameter-to-json-1}{%
\paragraph{\texorpdfstring{\texttt{unitree\_sdk2\_cpp.robot.b2.ConfigMetaParameter.to\_json}}{unitree\_sdk2\_cpp.robot.b2.ConfigMetaParameter.to\_json}}\label{unitree-sdk2-cpp-robot-b2-configmetaparameter-to-json-1}}

计划把当前 SDK 值对象写入 JSON 风格字典。

\textbf{签名}

\begin{lstlisting}[language=Python]
def to_json(self, json: dict[str, Any]) -> None
\end{lstlisting}

\textbf{可用性与安全性}

\textbf{\passthrough{\lstinline!SIGNATURE\_ONLY!} /
\passthrough{\lstinline!UNCLASSIFIED!}}：当前只有设计期签名，不能假设已存在于运行时扩展。

\textbf{参数}

\begin{longtable}[]{@{}
  >{\raggedright\arraybackslash}p{(\columnwidth - 6\tabcolsep) * \real{0.18}}
  >{\raggedright\arraybackslash}p{(\columnwidth - 6\tabcolsep) * \real{0.24}}
  >{\raggedright\arraybackslash}p{(\columnwidth - 6\tabcolsep) * \real{0.18}}
  >{\raggedright\arraybackslash}p{(\columnwidth - 6\tabcolsep) * \real{0.40}}@{}}
\toprule
名称 & Python 类型 & 默认值 & 含义 \\
\midrule
\endhead
\passthrough{\lstinline!json!} &
\passthrough{\lstinline!dict[str, Any]!} & 必填 & JSON
风格字典，键和值必须满足对应 C++ \passthrough{\lstinline!JsonMap!}
数据结构的约束。 对应 C++ 参数
\passthrough{\lstinline!json: common::JsonMap \&!}。 \\
\bottomrule
\end{longtable}

\textbf{返回值}

\begin{longtable}[]{@{}
  >{\raggedright\arraybackslash}p{(\columnwidth - 4\tabcolsep) * \real{0.18}}
  >{\raggedright\arraybackslash}p{(\columnwidth - 4\tabcolsep) * \real{0.20}}
  >{\raggedright\arraybackslash}p{(\columnwidth - 4\tabcolsep) * \real{0.62}}@{}}
\toprule
位置 & 类型 & 含义 \\
\midrule
\endhead
无 & \passthrough{\lstinline!None!} & 该方法不返回 Python
值；失败可能表现为异常或底层状态变化。 \\
\bottomrule
\end{longtable}

\textbf{对应 C++}

\begin{itemize}
\tightlist
\item
  类：\passthrough{\lstinline!unitree::robot::b2::ConfigMetaParameter!}
\item
  签名：\passthrough{\lstinline!toJson(common::JsonMap \&) const!}
\item
  绑定策略：\passthrough{\lstinline!SIGNATURE\_PREVIEW!}
\item
  声明位置：\passthrough{\lstinline!include/unitree/robot/b2/config/config\_api.hpp:174!}
\end{itemize}

\textbf{说明}

\begin{itemize}
\tightlist
\item
  示例是局部调用片段；\passthrough{\lstinline!obj!}
  和参数变量表示已经按上层业务校验并准备好的对象。
\item
  该条目可以被编辑器和类型检查器识别，但当前运行时可能在属性查找阶段直接失败。
\end{itemize}

\textbf{用法}

\begin{lstlisting}[language=Python]
from typing import TYPE_CHECKING

if TYPE_CHECKING:
    def planned_to_json(obj: ConfigMetaParameter, json: dict[str, Any]) -> None:
        obj.to_json(json=json)
\end{lstlisting}

\begin{quote}
{[}!CAUTION{]} 上面的 \passthrough{\lstinline!TYPE\_CHECKING!}
示例只表达计划调用形态。该分支在普通 Python
运行时为假，不代表当前扩展能执行此接口。
\end{quote}

\hypertarget{unitree-sdk2-cpp-robot-b2-configsetparameter}{%
\subsubsection{\texorpdfstring{\texttt{unitree\_sdk2\_cpp.robot.b2.ConfigSetParameter}}{unitree\_sdk2\_cpp.robot.b2.ConfigSetParameter}}\label{unitree-sdk2-cpp-robot-b2-configsetparameter}}

SDK 请求参数值类型；当前可能仅提供设计期签名。

\textbf{导入}

\begin{lstlisting}[language=Python]
from unitree_sdk2_cpp.robot.b2 import ConfigSetParameter
\end{lstlisting}

\textbf{基类}：\passthrough{\lstinline!object!}

\textbf{公开属性}

\begin{longtable}[]{@{}
  >{\raggedright\arraybackslash}p{(\columnwidth - 6\tabcolsep) * \real{0.18}}
  >{\raggedright\arraybackslash}p{(\columnwidth - 6\tabcolsep) * \real{0.24}}
  >{\raggedright\arraybackslash}p{(\columnwidth - 6\tabcolsep) * \real{0.18}}
  >{\raggedright\arraybackslash}p{(\columnwidth - 6\tabcolsep) * \real{0.40}}@{}}
\toprule
名称 & Python 类型 & 含义 & 用法 \\
\midrule
\endhead
\passthrough{\lstinline!name!} & \passthrough{\lstinline!str!} & SDK
对象、配置项、服务或动作的名称。允许值由调用它的具体接口定义。 &
\passthrough{\lstinline!value = obj.name!} /
\passthrough{\lstinline!obj.name = value!} \\
\passthrough{\lstinline!content!} & \passthrough{\lstinline!str!} &
配置、请求或序列化内容。具体格式由对应服务协议定义。 &
\passthrough{\lstinline!value = obj.content!} /
\passthrough{\lstinline!obj.content = value!} \\
\bottomrule
\end{longtable}

\textbf{方法索引}

\begin{longtable}[]{@{}
  >{\raggedright\arraybackslash}p{(\columnwidth - 6\tabcolsep) * \real{0.18}}
  >{\raggedright\arraybackslash}p{(\columnwidth - 6\tabcolsep) * \real{0.24}}
  >{\raggedright\arraybackslash}p{(\columnwidth - 6\tabcolsep) * \real{0.18}}
  >{\raggedright\arraybackslash}p{(\columnwidth - 6\tabcolsep) * \real{0.40}}@{}}
\toprule
方法 & Python 签名 & 状态 & 安全分类 \\
\midrule
\endhead
\protect\hyperlink{unitree-sdk2-cpp-robot-b2-configsetparameter-dunder-init-1}{\passthrough{\lstinline!\_\_init\_\_!}}
& \passthrough{\lstinline!def \_\_init\_\_(self) -> None!} &
\passthrough{\lstinline!SIGNATURE\_ONLY!} &
\passthrough{\lstinline!CONSTRUCTION!} \\
\protect\hyperlink{unitree-sdk2-cpp-robot-b2-configsetparameter-from-json-1}{\passthrough{\lstinline!from\_json!}}
&
\passthrough{\lstinline!def from\_json(self, json: dict[str, Any]) -> None!}
& \passthrough{\lstinline!SIGNATURE\_ONLY!} &
\passthrough{\lstinline!UNCLASSIFIED!} \\
\protect\hyperlink{unitree-sdk2-cpp-robot-b2-configsetparameter-to-json-1}{\passthrough{\lstinline!to\_json!}}
&
\passthrough{\lstinline!def to\_json(self, json: dict[str, Any]) -> None!}
& \passthrough{\lstinline!SIGNATURE\_ONLY!} &
\passthrough{\lstinline!UNCLASSIFIED!} \\
\bottomrule
\end{longtable}

\hypertarget{unitree-sdk2-cpp-robot-b2-configsetparameter-dunder-init-1}{%
\paragraph{\texorpdfstring{\texttt{unitree\_sdk2\_cpp.robot.b2.ConfigSetParameter.\_\_init\_\_}}{unitree\_sdk2\_cpp.robot.b2.ConfigSetParameter.\_\_init\_\_}}\label{unitree-sdk2-cpp-robot-b2-configsetparameter-dunder-init-1}}

初始化 \passthrough{\lstinline!ConfigSetParameter!}
实例。是否能实际构造取决于下方可用性状态。

\textbf{签名}

\begin{lstlisting}[language=Python]
def __init__(self) -> None
\end{lstlisting}

\textbf{可用性与安全性}

\textbf{\passthrough{\lstinline!SIGNATURE\_ONLY!} /
\passthrough{\lstinline!CONSTRUCTION!}}：当前只有设计期签名，不能假设已存在于运行时扩展。

\textbf{参数}

无显式参数。实例方法中的 \passthrough{\lstinline!self!} 由 Python
自动传入。

\textbf{返回值}

\begin{longtable}[]{@{}
  >{\raggedright\arraybackslash}p{(\columnwidth - 4\tabcolsep) * \real{0.18}}
  >{\raggedright\arraybackslash}p{(\columnwidth - 4\tabcolsep) * \real{0.20}}
  >{\raggedright\arraybackslash}p{(\columnwidth - 4\tabcolsep) * \real{0.62}}@{}}
\toprule
位置 & 类型 & 含义 \\
\midrule
\endhead
无 & \passthrough{\lstinline!None!} &
\passthrough{\lstinline!\_\_init\_\_!}
本身不返回值；调用类对象时，在构造函数可用的前提下得到该类实例。 \\
\bottomrule
\end{longtable}

\textbf{对应 C++}

\begin{itemize}
\tightlist
\item
  类：\passthrough{\lstinline!unitree::robot::b2::ConfigSetParameter!}
\item
  签名：\passthrough{\lstinline!ConfigSetParameter()!}
\item
  绑定策略：\passthrough{\lstinline!未单独标注!}
\item
  声明位置：\passthrough{\lstinline!include/unitree/robot/b2/config/config\_api.hpp:68!}
\end{itemize}

\textbf{说明}

\begin{itemize}
\tightlist
\item
  示例是局部调用片段；\passthrough{\lstinline!obj!}
  和参数变量表示已经按上层业务校验并准备好的对象。
\item
  该条目可以被编辑器和类型检查器识别，但当前运行时可能在属性查找阶段直接失败。
\end{itemize}

\textbf{用法}

\begin{lstlisting}[language=Python]
from typing import TYPE_CHECKING

if TYPE_CHECKING:
    def planned_construct() -> ConfigSetParameter:
        return ConfigSetParameter()
\end{lstlisting}

\begin{quote}
{[}!CAUTION{]} 上面的 \passthrough{\lstinline!TYPE\_CHECKING!}
示例只表达计划调用形态。该分支在普通 Python
运行时为假，不代表当前扩展能执行此接口。
\end{quote}

\hypertarget{unitree-sdk2-cpp-robot-b2-configsetparameter-from-json-1}{%
\paragraph{\texorpdfstring{\texttt{unitree\_sdk2\_cpp.robot.b2.ConfigSetParameter.from\_json}}{unitree\_sdk2\_cpp.robot.b2.ConfigSetParameter.from\_json}}\label{unitree-sdk2-cpp-robot-b2-configsetparameter-from-json-1}}

计划从 JSON 风格字典读取字段并更新当前 SDK 值对象。

\textbf{签名}

\begin{lstlisting}[language=Python]
def from_json(self, json: dict[str, Any]) -> None
\end{lstlisting}

\textbf{可用性与安全性}

\textbf{\passthrough{\lstinline!SIGNATURE\_ONLY!} /
\passthrough{\lstinline!UNCLASSIFIED!}}：当前只有设计期签名，不能假设已存在于运行时扩展。

\textbf{参数}

\begin{longtable}[]{@{}
  >{\raggedright\arraybackslash}p{(\columnwidth - 6\tabcolsep) * \real{0.18}}
  >{\raggedright\arraybackslash}p{(\columnwidth - 6\tabcolsep) * \real{0.24}}
  >{\raggedright\arraybackslash}p{(\columnwidth - 6\tabcolsep) * \real{0.18}}
  >{\raggedright\arraybackslash}p{(\columnwidth - 6\tabcolsep) * \real{0.40}}@{}}
\toprule
名称 & Python 类型 & 默认值 & 含义 \\
\midrule
\endhead
\passthrough{\lstinline!json!} &
\passthrough{\lstinline!dict[str, Any]!} & 必填 & JSON
风格字典，键和值必须满足对应 C++ \passthrough{\lstinline!JsonMap!}
数据结构的约束。 对应 C++ 参数
\passthrough{\lstinline!json: common::JsonMap \&!}。 \\
\bottomrule
\end{longtable}

\textbf{返回值}

\begin{longtable}[]{@{}
  >{\raggedright\arraybackslash}p{(\columnwidth - 4\tabcolsep) * \real{0.18}}
  >{\raggedright\arraybackslash}p{(\columnwidth - 4\tabcolsep) * \real{0.20}}
  >{\raggedright\arraybackslash}p{(\columnwidth - 4\tabcolsep) * \real{0.62}}@{}}
\toprule
位置 & 类型 & 含义 \\
\midrule
\endhead
无 & \passthrough{\lstinline!None!} & 该方法不返回 Python
值；失败可能表现为异常或底层状态变化。 \\
\bottomrule
\end{longtable}

\textbf{对应 C++}

\begin{itemize}
\tightlist
\item
  类：\passthrough{\lstinline!unitree::robot::b2::ConfigSetParameter!}
\item
  签名：\passthrough{\lstinline!fromJson(common::JsonMap \&)!}
\item
  绑定策略：\passthrough{\lstinline!SIGNATURE\_PREVIEW!}
\item
  声明位置：\passthrough{\lstinline!include/unitree/robot/b2/config/config\_api.hpp:74!}
\end{itemize}

\textbf{说明}

\begin{itemize}
\tightlist
\item
  示例是局部调用片段；\passthrough{\lstinline!obj!}
  和参数变量表示已经按上层业务校验并准备好的对象。
\item
  该条目可以被编辑器和类型检查器识别，但当前运行时可能在属性查找阶段直接失败。
\end{itemize}

\textbf{用法}

\begin{lstlisting}[language=Python]
from typing import TYPE_CHECKING

if TYPE_CHECKING:
    def planned_from_json(obj: ConfigSetParameter, json: dict[str, Any]) -> None:
        obj.from_json(json=json)
\end{lstlisting}

\begin{quote}
{[}!CAUTION{]} 上面的 \passthrough{\lstinline!TYPE\_CHECKING!}
示例只表达计划调用形态。该分支在普通 Python
运行时为假，不代表当前扩展能执行此接口。
\end{quote}

\hypertarget{unitree-sdk2-cpp-robot-b2-configsetparameter-to-json-1}{%
\paragraph{\texorpdfstring{\texttt{unitree\_sdk2\_cpp.robot.b2.ConfigSetParameter.to\_json}}{unitree\_sdk2\_cpp.robot.b2.ConfigSetParameter.to\_json}}\label{unitree-sdk2-cpp-robot-b2-configsetparameter-to-json-1}}

计划把当前 SDK 值对象写入 JSON 风格字典。

\textbf{签名}

\begin{lstlisting}[language=Python]
def to_json(self, json: dict[str, Any]) -> None
\end{lstlisting}

\textbf{可用性与安全性}

\textbf{\passthrough{\lstinline!SIGNATURE\_ONLY!} /
\passthrough{\lstinline!UNCLASSIFIED!}}：当前只有设计期签名，不能假设已存在于运行时扩展。

\textbf{参数}

\begin{longtable}[]{@{}
  >{\raggedright\arraybackslash}p{(\columnwidth - 6\tabcolsep) * \real{0.18}}
  >{\raggedright\arraybackslash}p{(\columnwidth - 6\tabcolsep) * \real{0.24}}
  >{\raggedright\arraybackslash}p{(\columnwidth - 6\tabcolsep) * \real{0.18}}
  >{\raggedright\arraybackslash}p{(\columnwidth - 6\tabcolsep) * \real{0.40}}@{}}
\toprule
名称 & Python 类型 & 默认值 & 含义 \\
\midrule
\endhead
\passthrough{\lstinline!json!} &
\passthrough{\lstinline!dict[str, Any]!} & 必填 & JSON
风格字典，键和值必须满足对应 C++ \passthrough{\lstinline!JsonMap!}
数据结构的约束。 对应 C++ 参数
\passthrough{\lstinline!json: common::JsonMap \&!}。 \\
\bottomrule
\end{longtable}

\textbf{返回值}

\begin{longtable}[]{@{}
  >{\raggedright\arraybackslash}p{(\columnwidth - 4\tabcolsep) * \real{0.18}}
  >{\raggedright\arraybackslash}p{(\columnwidth - 4\tabcolsep) * \real{0.20}}
  >{\raggedright\arraybackslash}p{(\columnwidth - 4\tabcolsep) * \real{0.62}}@{}}
\toprule
位置 & 类型 & 含义 \\
\midrule
\endhead
无 & \passthrough{\lstinline!None!} & 该方法不返回 Python
值；失败可能表现为异常或底层状态变化。 \\
\bottomrule
\end{longtable}

\textbf{对应 C++}

\begin{itemize}
\tightlist
\item
  类：\passthrough{\lstinline!unitree::robot::b2::ConfigSetParameter!}
\item
  签名：\passthrough{\lstinline!toJson(common::JsonMap \&) const!}
\item
  绑定策略：\passthrough{\lstinline!SIGNATURE\_PREVIEW!}
\item
  声明位置：\passthrough{\lstinline!include/unitree/robot/b2/config/config\_api.hpp:80!}
\end{itemize}

\textbf{说明}

\begin{itemize}
\tightlist
\item
  示例是局部调用片段；\passthrough{\lstinline!obj!}
  和参数变量表示已经按上层业务校验并准备好的对象。
\item
  该条目可以被编辑器和类型检查器识别，但当前运行时可能在属性查找阶段直接失败。
\end{itemize}

\textbf{用法}

\begin{lstlisting}[language=Python]
from typing import TYPE_CHECKING

if TYPE_CHECKING:
    def planned_to_json(obj: ConfigSetParameter, json: dict[str, Any]) -> None:
        obj.to_json(json=json)
\end{lstlisting}

\begin{quote}
{[}!CAUTION{]} 上面的 \passthrough{\lstinline!TYPE\_CHECKING!}
示例只表达计划调用形态。该分支在普通 Python
运行时为假，不代表当前扩展能执行此接口。
\end{quote}

\hypertarget{unitree-sdk2-cpp-robot-b2-frontvideoclient}{%
\subsubsection{\texorpdfstring{\texttt{unitree\_sdk2\_cpp.robot.b2.FrontVideoClient}}{unitree\_sdk2\_cpp.robot.b2.FrontVideoClient}}\label{unitree-sdk2-cpp-robot-b2-frontvideoclient}}

Robot 服务客户端。构造、初始化和具体方法的可用性必须分别检查。

\textbf{导入}

\begin{lstlisting}[language=Python]
from unitree_sdk2_cpp.robot.b2 import FrontVideoClient
\end{lstlisting}

\textbf{基类}：\passthrough{\lstinline!Client!}

\textbf{方法索引}

\begin{longtable}[]{@{}
  >{\raggedright\arraybackslash}p{(\columnwidth - 6\tabcolsep) * \real{0.18}}
  >{\raggedright\arraybackslash}p{(\columnwidth - 6\tabcolsep) * \real{0.24}}
  >{\raggedright\arraybackslash}p{(\columnwidth - 6\tabcolsep) * \real{0.18}}
  >{\raggedright\arraybackslash}p{(\columnwidth - 6\tabcolsep) * \real{0.40}}@{}}
\toprule
方法 & Python 签名 & 状态 & 安全分类 \\
\midrule
\endhead
\protect\hyperlink{unitree-sdk2-cpp-robot-b2-frontvideoclient-dunder-init-1}{\passthrough{\lstinline!\_\_init\_\_!}}
& \passthrough{\lstinline!def \_\_init\_\_(self) -> None!} &
\passthrough{\lstinline!AVAILABLE!} &
\passthrough{\lstinline!CONSTRUCTION!} \\
\protect\hyperlink{unitree-sdk2-cpp-robot-b2-frontvideoclient-init-1}{\passthrough{\lstinline!init!}}
& \passthrough{\lstinline!def init(self) -> None!} &
\passthrough{\lstinline!AVAILABLE!} &
\passthrough{\lstinline!INITIALIZATION!} \\
\protect\hyperlink{unitree-sdk2-cpp-robot-b2-frontvideoclient-get-image-sample-1}{\passthrough{\lstinline!get\_image\_sample!}}
&
\passthrough{\lstinline!def get\_image\_sample(self) -> tuple[int, list[int]]!}
& \passthrough{\lstinline!AVAILABLE!} &
\passthrough{\lstinline!READ\_ONLY!} \\
\bottomrule
\end{longtable}

\hypertarget{unitree-sdk2-cpp-robot-b2-frontvideoclient-dunder-init-1}{%
\paragraph{\texorpdfstring{\texttt{unitree\_sdk2\_cpp.robot.b2.FrontVideoClient.\_\_init\_\_}}{unitree\_sdk2\_cpp.robot.b2.FrontVideoClient.\_\_init\_\_}}\label{unitree-sdk2-cpp-robot-b2-frontvideoclient-dunder-init-1}}

初始化 \passthrough{\lstinline!FrontVideoClient!}
实例。是否能实际构造取决于下方可用性状态。

\textbf{签名}

\begin{lstlisting}[language=Python]
def __init__(self) -> None
\end{lstlisting}

\textbf{可用性与安全性}

\textbf{\passthrough{\lstinline!AVAILABLE!} /
\passthrough{\lstinline!CONSTRUCTION!}}：当前绑定源码已实现；实际执行条件仍取决于平台和对应
SDK 环境。

\textbf{参数}

无显式参数。实例方法中的 \passthrough{\lstinline!self!} 由 Python
自动传入。

\textbf{返回值}

\begin{longtable}[]{@{}
  >{\raggedright\arraybackslash}p{(\columnwidth - 4\tabcolsep) * \real{0.18}}
  >{\raggedright\arraybackslash}p{(\columnwidth - 4\tabcolsep) * \real{0.20}}
  >{\raggedright\arraybackslash}p{(\columnwidth - 4\tabcolsep) * \real{0.62}}@{}}
\toprule
位置 & 类型 & 含义 \\
\midrule
\endhead
无 & \passthrough{\lstinline!None!} &
\passthrough{\lstinline!\_\_init\_\_!}
本身不返回值；调用类对象时，在构造函数可用的前提下得到该类实例。 \\
\bottomrule
\end{longtable}

\textbf{对应 C++}

\begin{itemize}
\tightlist
\item
  类：\passthrough{\lstinline!unitree::robot::b2::FrontVideoClient!}
\item
  签名：\passthrough{\lstinline!FrontVideoClient()!}
\item
  绑定策略：\passthrough{\lstinline!未单独标注!}
\item
  声明位置：\passthrough{\lstinline!include/unitree/robot/b2/front\_video/front\_video\_client.hpp:18!}
\end{itemize}

\textbf{说明}

\begin{itemize}
\tightlist
\item
  示例是局部调用片段；\passthrough{\lstinline!obj!}
  和参数变量表示已经按上层业务校验并准备好的对象。
\end{itemize}

\textbf{用法}

\begin{lstlisting}[language=Python]
value = FrontVideoClient()
\end{lstlisting}

\hypertarget{unitree-sdk2-cpp-robot-b2-frontvideoclient-init-1}{%
\paragraph{\texorpdfstring{\texttt{unitree\_sdk2\_cpp.robot.b2.FrontVideoClient.init}}{unitree\_sdk2\_cpp.robot.b2.FrontVideoClient.init}}\label{unitree-sdk2-cpp-robot-b2-frontvideoclient-init-1}}

初始化当前 SDK 对象所需的底层通道或服务资源。

\textbf{签名}

\begin{lstlisting}[language=Python]
def init(self) -> None
\end{lstlisting}

\textbf{可用性与安全性}

\textbf{\passthrough{\lstinline!AVAILABLE!} /
\passthrough{\lstinline!INITIALIZATION!}}：当前绑定源码已实现；实际执行条件仍取决于平台和对应
SDK 环境。

\textbf{参数}

无显式参数。实例方法中的 \passthrough{\lstinline!self!} 由 Python
自动传入。

\textbf{返回值}

\begin{longtable}[]{@{}
  >{\raggedright\arraybackslash}p{(\columnwidth - 4\tabcolsep) * \real{0.18}}
  >{\raggedright\arraybackslash}p{(\columnwidth - 4\tabcolsep) * \real{0.20}}
  >{\raggedright\arraybackslash}p{(\columnwidth - 4\tabcolsep) * \real{0.62}}@{}}
\toprule
位置 & 类型 & 含义 \\
\midrule
\endhead
无 & \passthrough{\lstinline!None!} & 该方法不返回 Python
值；失败可能表现为异常或底层状态变化。 \\
\bottomrule
\end{longtable}

\textbf{对应 C++}

\begin{itemize}
\tightlist
\item
  类：\passthrough{\lstinline!unitree::robot::b2::FrontVideoClient!}
\item
  签名：\passthrough{\lstinline!Init()!}
\item
  绑定策略：\passthrough{\lstinline!DIRECT!}
\item
  声明位置：\passthrough{\lstinline!include/unitree/robot/b2/front\_video/front\_video\_client.hpp:21!}
\end{itemize}

\textbf{说明}

\begin{itemize}
\tightlist
\item
  示例是局部调用片段；\passthrough{\lstinline!obj!}
  和参数变量表示已经按上层业务校验并准备好的对象。
\end{itemize}

\textbf{用法}

\begin{lstlisting}[language=Python]
obj.init()
\end{lstlisting}

\hypertarget{unitree-sdk2-cpp-robot-b2-frontvideoclient-get-image-sample-1}{%
\paragraph{\texorpdfstring{\texttt{unitree\_sdk2\_cpp.robot.b2.FrontVideoClient.get\_image\_sample}}{unitree\_sdk2\_cpp.robot.b2.FrontVideoClient.get\_image\_sample}}\label{unitree-sdk2-cpp-robot-b2-frontvideoclient-get-image-sample-1}}

查询或检查 图像 \passthrough{\lstinline!sample!}。

\textbf{签名}

\begin{lstlisting}[language=Python]
def get_image_sample(self) -> tuple[int, list[int]]
\end{lstlisting}

\textbf{可用性与安全性}

\textbf{\passthrough{\lstinline!AVAILABLE!} /
\passthrough{\lstinline!READ\_ONLY!}}：当前绑定源码已实现只读查询。Robot
Client 的服务端查询通常仍需匹配的 Linux
扩展、DDS、网络、机器人和服务；纯本地 getter 则不一定需要全部条件。

\textbf{参数}

无显式参数。实例方法中的 \passthrough{\lstinline!self!} 由 Python
自动传入。

\textbf{返回值}

\begin{longtable}[]{@{}
  >{\raggedright\arraybackslash}p{(\columnwidth - 4\tabcolsep) * \real{0.18}}
  >{\raggedright\arraybackslash}p{(\columnwidth - 4\tabcolsep) * \real{0.20}}
  >{\raggedright\arraybackslash}p{(\columnwidth - 4\tabcolsep) * \real{0.62}}@{}}
\toprule
位置 & 类型 & 含义 \\
\midrule
\endhead
{[}0{]} \passthrough{\lstinline!status!} & \passthrough{\lstinline!int!}
& SDK 状态码；Unitree 示例通常以 \passthrough{\lstinline!0!}
表示成功，非零值需按具体服务错误码解释。 \\
{[}1{]} \passthrough{\lstinline!output\_1!} &
\passthrough{\lstinline!list[int]!} & 传给该接口的
\passthrough{\lstinline!output!} \passthrough{\lstinline!1!}
参数，Python 类型为
\passthrough{\lstinline!list[int]!}。精确含义、单位、范围和枚举值需查目标型号对应头文件与协议。
对应 C++ 参数
\passthrough{\lstinline!output\_1: std::vector<uint8\_t> \&!}。
底层是可变长度 vector；元素约束：取值范围为 0 到 255。 \\
\bottomrule
\end{longtable}

\textbf{对应 C++}

\begin{itemize}
\tightlist
\item
  类：\passthrough{\lstinline!unitree::robot::b2::FrontVideoClient!}
\item
  签名：\passthrough{\lstinline!GetImageSample(std::vector<uint8\_t> \&)!}
\item
  绑定策略：\passthrough{\lstinline!OUTPUT\_WRAPPER!}
\item
  声明位置：\passthrough{\lstinline!include/unitree/robot/b2/front\_video/front\_video\_client.hpp:27!}
\end{itemize}

\textbf{说明}

\begin{itemize}
\tightlist
\item
  示例是局部调用片段；\passthrough{\lstinline!obj!}
  和参数变量表示已经按上层业务校验并准备好的对象。
\item
  C++ 的可修改输出引用不会作为 Python 入参出现，而是按 C++
  参数顺序追加到返回元组中。
\end{itemize}

\textbf{用法}

\begin{lstlisting}[language=Python]
status, output_1 = obj.get_image_sample()
\end{lstlisting}

\begin{quote}
{[}!NOTE{]} 只读查询仍会访问
DDS/机器人服务。必须检查返回状态码，并处理超时和网络中断。
\end{quote}

\hypertarget{unitree-sdk2-cpp-robot-b2-jsonizeconfigmeta}{%
\subsubsection{\texorpdfstring{\texttt{unitree\_sdk2\_cpp.robot.b2.JsonizeConfigMeta}}{unitree\_sdk2\_cpp.robot.b2.JsonizeConfigMeta}}\label{unitree-sdk2-cpp-robot-b2-jsonizeconfigmeta}}

SDK 数据值类型；公开字段和转换方法见下文。

\textbf{导入}

\begin{lstlisting}[language=Python]
from unitree_sdk2_cpp.robot.b2 import JsonizeConfigMeta
\end{lstlisting}

\textbf{基类}：\passthrough{\lstinline!object!}

\textbf{公开属性}

\begin{longtable}[]{@{}
  >{\raggedright\arraybackslash}p{(\columnwidth - 6\tabcolsep) * \real{0.18}}
  >{\raggedright\arraybackslash}p{(\columnwidth - 6\tabcolsep) * \real{0.24}}
  >{\raggedright\arraybackslash}p{(\columnwidth - 6\tabcolsep) * \real{0.18}}
  >{\raggedright\arraybackslash}p{(\columnwidth - 6\tabcolsep) * \real{0.40}}@{}}
\toprule
名称 & Python 类型 & 含义 & 用法 \\
\midrule
\endhead
\passthrough{\lstinline!name!} & \passthrough{\lstinline!str!} & SDK
对象、配置项、服务或动作的名称。允许值由调用它的具体接口定义。 &
\passthrough{\lstinline!value = obj.name!} /
\passthrough{\lstinline!obj.name = value!} \\
\passthrough{\lstinline!last\_modified!} & \passthrough{\lstinline!str!}
& 传给该接口的 \passthrough{\lstinline!last!}
\passthrough{\lstinline!modified!} 参数，Python 类型为
\passthrough{\lstinline!str!}。精确含义、单位、范围和枚举值需查目标型号对应头文件与协议。
& \passthrough{\lstinline!value = obj.last\_modified!} /
\passthrough{\lstinline!obj.last\_modified = value!} \\
\passthrough{\lstinline!size!} & \passthrough{\lstinline!int!} &
传给该接口的 \passthrough{\lstinline!size!} 参数，Python 类型为
\passthrough{\lstinline!int!}。精确含义、单位、范围和枚举值需查目标型号对应头文件与协议。
& \passthrough{\lstinline!value = obj.size!} /
\passthrough{\lstinline!obj.size = value!} \\
\passthrough{\lstinline!epoch!} & \passthrough{\lstinline!int!} &
传给该接口的 \passthrough{\lstinline!epoch!} 参数，Python 类型为
\passthrough{\lstinline!int!}。精确含义、单位、范围和枚举值需查目标型号对应头文件与协议。
& \passthrough{\lstinline!value = obj.epoch!} /
\passthrough{\lstinline!obj.epoch = value!} \\
\bottomrule
\end{longtable}

\textbf{方法索引}

\begin{longtable}[]{@{}
  >{\raggedright\arraybackslash}p{(\columnwidth - 6\tabcolsep) * \real{0.18}}
  >{\raggedright\arraybackslash}p{(\columnwidth - 6\tabcolsep) * \real{0.24}}
  >{\raggedright\arraybackslash}p{(\columnwidth - 6\tabcolsep) * \real{0.18}}
  >{\raggedright\arraybackslash}p{(\columnwidth - 6\tabcolsep) * \real{0.40}}@{}}
\toprule
方法 & Python 签名 & 状态 & 安全分类 \\
\midrule
\endhead
\protect\hyperlink{unitree-sdk2-cpp-robot-b2-jsonizeconfigmeta-dunder-init-1}{\passthrough{\lstinline!\_\_init\_\_!}}
& \passthrough{\lstinline!def \_\_init\_\_(self) -> None!} &
\passthrough{\lstinline!SIGNATURE\_ONLY!} &
\passthrough{\lstinline!CONSTRUCTION!} \\
\protect\hyperlink{unitree-sdk2-cpp-robot-b2-jsonizeconfigmeta-from-json-1}{\passthrough{\lstinline!from\_json!}}
&
\passthrough{\lstinline!def from\_json(self, json: dict[str, Any]) -> None!}
& \passthrough{\lstinline!SIGNATURE\_ONLY!} &
\passthrough{\lstinline!UNCLASSIFIED!} \\
\protect\hyperlink{unitree-sdk2-cpp-robot-b2-jsonizeconfigmeta-to-json-1}{\passthrough{\lstinline!to\_json!}}
&
\passthrough{\lstinline!def to\_json(self, json: dict[str, Any]) -> None!}
& \passthrough{\lstinline!SIGNATURE\_ONLY!} &
\passthrough{\lstinline!UNCLASSIFIED!} \\
\bottomrule
\end{longtable}

\hypertarget{unitree-sdk2-cpp-robot-b2-jsonizeconfigmeta-dunder-init-1}{%
\paragraph{\texorpdfstring{\texttt{unitree\_sdk2\_cpp.robot.b2.JsonizeConfigMeta.\_\_init\_\_}}{unitree\_sdk2\_cpp.robot.b2.JsonizeConfigMeta.\_\_init\_\_}}\label{unitree-sdk2-cpp-robot-b2-jsonizeconfigmeta-dunder-init-1}}

初始化 \passthrough{\lstinline!JsonizeConfigMeta!}
实例。是否能实际构造取决于下方可用性状态。

\textbf{签名}

\begin{lstlisting}[language=Python]
def __init__(self) -> None
\end{lstlisting}

\textbf{可用性与安全性}

\textbf{\passthrough{\lstinline!SIGNATURE\_ONLY!} /
\passthrough{\lstinline!CONSTRUCTION!}}：当前只有设计期签名，不能假设已存在于运行时扩展。

\textbf{参数}

无显式参数。实例方法中的 \passthrough{\lstinline!self!} 由 Python
自动传入。

\textbf{返回值}

\begin{longtable}[]{@{}
  >{\raggedright\arraybackslash}p{(\columnwidth - 4\tabcolsep) * \real{0.18}}
  >{\raggedright\arraybackslash}p{(\columnwidth - 4\tabcolsep) * \real{0.20}}
  >{\raggedright\arraybackslash}p{(\columnwidth - 4\tabcolsep) * \real{0.62}}@{}}
\toprule
位置 & 类型 & 含义 \\
\midrule
\endhead
无 & \passthrough{\lstinline!None!} &
\passthrough{\lstinline!\_\_init\_\_!}
本身不返回值；调用类对象时，在构造函数可用的前提下得到该类实例。 \\
\bottomrule
\end{longtable}

\textbf{对应 C++}

\begin{itemize}
\tightlist
\item
  类：\passthrough{\lstinline!unitree::robot::b2::JsonizeConfigMeta!}
\item
  签名：\passthrough{\lstinline!JsonizeConfigMeta()!}
\item
  绑定策略：\passthrough{\lstinline!未单独标注!}
\item
  声明位置：\passthrough{\lstinline!include/unitree/robot/b2/config/config\_api.hpp:36!}
\end{itemize}

\textbf{说明}

\begin{itemize}
\tightlist
\item
  示例是局部调用片段；\passthrough{\lstinline!obj!}
  和参数变量表示已经按上层业务校验并准备好的对象。
\item
  该条目可以被编辑器和类型检查器识别，但当前运行时可能在属性查找阶段直接失败。
\end{itemize}

\textbf{用法}

\begin{lstlisting}[language=Python]
from typing import TYPE_CHECKING

if TYPE_CHECKING:
    def planned_construct() -> JsonizeConfigMeta:
        return JsonizeConfigMeta()
\end{lstlisting}

\begin{quote}
{[}!CAUTION{]} 上面的 \passthrough{\lstinline!TYPE\_CHECKING!}
示例只表达计划调用形态。该分支在普通 Python
运行时为假，不代表当前扩展能执行此接口。
\end{quote}

\hypertarget{unitree-sdk2-cpp-robot-b2-jsonizeconfigmeta-from-json-1}{%
\paragraph{\texorpdfstring{\texttt{unitree\_sdk2\_cpp.robot.b2.JsonizeConfigMeta.from\_json}}{unitree\_sdk2\_cpp.robot.b2.JsonizeConfigMeta.from\_json}}\label{unitree-sdk2-cpp-robot-b2-jsonizeconfigmeta-from-json-1}}

计划从 JSON 风格字典读取字段并更新当前 SDK 值对象。

\textbf{签名}

\begin{lstlisting}[language=Python]
def from_json(self, json: dict[str, Any]) -> None
\end{lstlisting}

\textbf{可用性与安全性}

\textbf{\passthrough{\lstinline!SIGNATURE\_ONLY!} /
\passthrough{\lstinline!UNCLASSIFIED!}}：当前只有设计期签名，不能假设已存在于运行时扩展。

\textbf{参数}

\begin{longtable}[]{@{}
  >{\raggedright\arraybackslash}p{(\columnwidth - 6\tabcolsep) * \real{0.18}}
  >{\raggedright\arraybackslash}p{(\columnwidth - 6\tabcolsep) * \real{0.24}}
  >{\raggedright\arraybackslash}p{(\columnwidth - 6\tabcolsep) * \real{0.18}}
  >{\raggedright\arraybackslash}p{(\columnwidth - 6\tabcolsep) * \real{0.40}}@{}}
\toprule
名称 & Python 类型 & 默认值 & 含义 \\
\midrule
\endhead
\passthrough{\lstinline!json!} &
\passthrough{\lstinline!dict[str, Any]!} & 必填 & JSON
风格字典，键和值必须满足对应 C++ \passthrough{\lstinline!JsonMap!}
数据结构的约束。 对应 C++ 参数
\passthrough{\lstinline!json: common::JsonMap \&!}。 \\
\bottomrule
\end{longtable}

\textbf{返回值}

\begin{longtable}[]{@{}
  >{\raggedright\arraybackslash}p{(\columnwidth - 4\tabcolsep) * \real{0.18}}
  >{\raggedright\arraybackslash}p{(\columnwidth - 4\tabcolsep) * \real{0.20}}
  >{\raggedright\arraybackslash}p{(\columnwidth - 4\tabcolsep) * \real{0.62}}@{}}
\toprule
位置 & 类型 & 含义 \\
\midrule
\endhead
无 & \passthrough{\lstinline!None!} & 该方法不返回 Python
值；失败可能表现为异常或底层状态变化。 \\
\bottomrule
\end{longtable}

\textbf{对应 C++}

\begin{itemize}
\tightlist
\item
  类：\passthrough{\lstinline!unitree::robot::b2::JsonizeConfigMeta!}
\item
  签名：\passthrough{\lstinline!fromJson(common::JsonMap \&)!}
\item
  绑定策略：\passthrough{\lstinline!SIGNATURE\_PREVIEW!}
\item
  声明位置：\passthrough{\lstinline!include/unitree/robot/b2/config/config\_api.hpp:42!}
\end{itemize}

\textbf{说明}

\begin{itemize}
\tightlist
\item
  示例是局部调用片段；\passthrough{\lstinline!obj!}
  和参数变量表示已经按上层业务校验并准备好的对象。
\item
  该条目可以被编辑器和类型检查器识别，但当前运行时可能在属性查找阶段直接失败。
\end{itemize}

\textbf{用法}

\begin{lstlisting}[language=Python]
from typing import TYPE_CHECKING

if TYPE_CHECKING:
    def planned_from_json(obj: JsonizeConfigMeta, json: dict[str, Any]) -> None:
        obj.from_json(json=json)
\end{lstlisting}

\begin{quote}
{[}!CAUTION{]} 上面的 \passthrough{\lstinline!TYPE\_CHECKING!}
示例只表达计划调用形态。该分支在普通 Python
运行时为假，不代表当前扩展能执行此接口。
\end{quote}

\hypertarget{unitree-sdk2-cpp-robot-b2-jsonizeconfigmeta-to-json-1}{%
\paragraph{\texorpdfstring{\texttt{unitree\_sdk2\_cpp.robot.b2.JsonizeConfigMeta.to\_json}}{unitree\_sdk2\_cpp.robot.b2.JsonizeConfigMeta.to\_json}}\label{unitree-sdk2-cpp-robot-b2-jsonizeconfigmeta-to-json-1}}

计划把当前 SDK 值对象写入 JSON 风格字典。

\textbf{签名}

\begin{lstlisting}[language=Python]
def to_json(self, json: dict[str, Any]) -> None
\end{lstlisting}

\textbf{可用性与安全性}

\textbf{\passthrough{\lstinline!SIGNATURE\_ONLY!} /
\passthrough{\lstinline!UNCLASSIFIED!}}：当前只有设计期签名，不能假设已存在于运行时扩展。

\textbf{参数}

\begin{longtable}[]{@{}
  >{\raggedright\arraybackslash}p{(\columnwidth - 6\tabcolsep) * \real{0.18}}
  >{\raggedright\arraybackslash}p{(\columnwidth - 6\tabcolsep) * \real{0.24}}
  >{\raggedright\arraybackslash}p{(\columnwidth - 6\tabcolsep) * \real{0.18}}
  >{\raggedright\arraybackslash}p{(\columnwidth - 6\tabcolsep) * \real{0.40}}@{}}
\toprule
名称 & Python 类型 & 默认值 & 含义 \\
\midrule
\endhead
\passthrough{\lstinline!json!} &
\passthrough{\lstinline!dict[str, Any]!} & 必填 & JSON
风格字典，键和值必须满足对应 C++ \passthrough{\lstinline!JsonMap!}
数据结构的约束。 对应 C++ 参数
\passthrough{\lstinline!json: common::JsonMap \&!}。 \\
\bottomrule
\end{longtable}

\textbf{返回值}

\begin{longtable}[]{@{}
  >{\raggedright\arraybackslash}p{(\columnwidth - 4\tabcolsep) * \real{0.18}}
  >{\raggedright\arraybackslash}p{(\columnwidth - 4\tabcolsep) * \real{0.20}}
  >{\raggedright\arraybackslash}p{(\columnwidth - 4\tabcolsep) * \real{0.62}}@{}}
\toprule
位置 & 类型 & 含义 \\
\midrule
\endhead
无 & \passthrough{\lstinline!None!} & 该方法不返回 Python
值；失败可能表现为异常或底层状态变化。 \\
\bottomrule
\end{longtable}

\textbf{对应 C++}

\begin{itemize}
\tightlist
\item
  类：\passthrough{\lstinline!unitree::robot::b2::JsonizeConfigMeta!}
\item
  签名：\passthrough{\lstinline!toJson(common::JsonMap \&) const!}
\item
  绑定策略：\passthrough{\lstinline!SIGNATURE\_PREVIEW!}
\item
  声明位置：\passthrough{\lstinline!include/unitree/robot/b2/config/config\_api.hpp:50!}
\end{itemize}

\textbf{说明}

\begin{itemize}
\tightlist
\item
  示例是局部调用片段；\passthrough{\lstinline!obj!}
  和参数变量表示已经按上层业务校验并准备好的对象。
\item
  该条目可以被编辑器和类型检查器识别，但当前运行时可能在属性查找阶段直接失败。
\end{itemize}

\textbf{用法}

\begin{lstlisting}[language=Python]
from typing import TYPE_CHECKING

if TYPE_CHECKING:
    def planned_to_json(obj: JsonizeConfigMeta, json: dict[str, Any]) -> None:
        obj.to_json(json=json)
\end{lstlisting}

\begin{quote}
{[}!CAUTION{]} 上面的 \passthrough{\lstinline!TYPE\_CHECKING!}
示例只表达计划调用形态。该分支在普通 Python
运行时为假，不代表当前扩展能执行此接口。
\end{quote}

\hypertarget{unitree-sdk2-cpp-robot-b2-jsonizemodename}{%
\subsubsection{\texorpdfstring{\texttt{unitree\_sdk2\_cpp.robot.b2.JsonizeModeName}}{unitree\_sdk2\_cpp.robot.b2.JsonizeModeName}}\label{unitree-sdk2-cpp-robot-b2-jsonizemodename}}

SDK 类型签名预览。请逐项查看构造函数和方法的可用性。

\textbf{导入}

\begin{lstlisting}[language=Python]
from unitree_sdk2_cpp.robot.b2 import JsonizeModeName
\end{lstlisting}

\textbf{基类}：\passthrough{\lstinline!object!}

\textbf{公开属性}

\begin{longtable}[]{@{}
  >{\raggedright\arraybackslash}p{(\columnwidth - 6\tabcolsep) * \real{0.18}}
  >{\raggedright\arraybackslash}p{(\columnwidth - 6\tabcolsep) * \real{0.24}}
  >{\raggedright\arraybackslash}p{(\columnwidth - 6\tabcolsep) * \real{0.18}}
  >{\raggedright\arraybackslash}p{(\columnwidth - 6\tabcolsep) * \real{0.40}}@{}}
\toprule
名称 & Python 类型 & 含义 & 用法 \\
\midrule
\endhead
\passthrough{\lstinline!name!} & \passthrough{\lstinline!str!} & SDK
对象、配置项、服务或动作的名称。允许值由调用它的具体接口定义。 &
\passthrough{\lstinline!value = obj.name!} /
\passthrough{\lstinline!obj.name = value!} \\
\passthrough{\lstinline!form!} & \passthrough{\lstinline!str!} &
传给该接口的 \passthrough{\lstinline!form!} 参数，Python 类型为
\passthrough{\lstinline!str!}。精确含义、单位、范围和枚举值需查目标型号对应头文件与协议。
& \passthrough{\lstinline!value = obj.form!} /
\passthrough{\lstinline!obj.form = value!} \\
\bottomrule
\end{longtable}

\textbf{方法索引}

\begin{longtable}[]{@{}
  >{\raggedright\arraybackslash}p{(\columnwidth - 6\tabcolsep) * \real{0.18}}
  >{\raggedright\arraybackslash}p{(\columnwidth - 6\tabcolsep) * \real{0.24}}
  >{\raggedright\arraybackslash}p{(\columnwidth - 6\tabcolsep) * \real{0.18}}
  >{\raggedright\arraybackslash}p{(\columnwidth - 6\tabcolsep) * \real{0.40}}@{}}
\toprule
方法 & Python 签名 & 状态 & 安全分类 \\
\midrule
\endhead
\protect\hyperlink{unitree-sdk2-cpp-robot-b2-jsonizemodename-dunder-init-1}{\passthrough{\lstinline!\_\_init\_\_!}}
& \passthrough{\lstinline!def \_\_init\_\_(self) -> None!} &
\passthrough{\lstinline!SIGNATURE\_ONLY!} &
\passthrough{\lstinline!CONSTRUCTION!} \\
\protect\hyperlink{unitree-sdk2-cpp-robot-b2-jsonizemodename-from-json-1}{\passthrough{\lstinline!from\_json!}}
&
\passthrough{\lstinline!def from\_json(self, json: dict[str, Any]) -> None!}
& \passthrough{\lstinline!SIGNATURE\_ONLY!} &
\passthrough{\lstinline!UNCLASSIFIED!} \\
\protect\hyperlink{unitree-sdk2-cpp-robot-b2-jsonizemodename-to-json-1}{\passthrough{\lstinline!to\_json!}}
&
\passthrough{\lstinline!def to\_json(self, json: dict[str, Any]) -> None!}
& \passthrough{\lstinline!SIGNATURE\_ONLY!} &
\passthrough{\lstinline!UNCLASSIFIED!} \\
\bottomrule
\end{longtable}

\hypertarget{unitree-sdk2-cpp-robot-b2-jsonizemodename-dunder-init-1}{%
\paragraph{\texorpdfstring{\texttt{unitree\_sdk2\_cpp.robot.b2.JsonizeModeName.\_\_init\_\_}}{unitree\_sdk2\_cpp.robot.b2.JsonizeModeName.\_\_init\_\_}}\label{unitree-sdk2-cpp-robot-b2-jsonizemodename-dunder-init-1}}

初始化 \passthrough{\lstinline!JsonizeModeName!}
实例。是否能实际构造取决于下方可用性状态。

\textbf{签名}

\begin{lstlisting}[language=Python]
def __init__(self) -> None
\end{lstlisting}

\textbf{可用性与安全性}

\textbf{\passthrough{\lstinline!SIGNATURE\_ONLY!} /
\passthrough{\lstinline!CONSTRUCTION!}}：当前只有设计期签名，不能假设已存在于运行时扩展。

\textbf{参数}

无显式参数。实例方法中的 \passthrough{\lstinline!self!} 由 Python
自动传入。

\textbf{返回值}

\begin{longtable}[]{@{}
  >{\raggedright\arraybackslash}p{(\columnwidth - 4\tabcolsep) * \real{0.18}}
  >{\raggedright\arraybackslash}p{(\columnwidth - 4\tabcolsep) * \real{0.20}}
  >{\raggedright\arraybackslash}p{(\columnwidth - 4\tabcolsep) * \real{0.62}}@{}}
\toprule
位置 & 类型 & 含义 \\
\midrule
\endhead
无 & \passthrough{\lstinline!None!} &
\passthrough{\lstinline!\_\_init\_\_!}
本身不返回值；调用类对象时，在构造函数可用的前提下得到该类实例。 \\
\bottomrule
\end{longtable}

\textbf{对应 C++}

\begin{itemize}
\tightlist
\item
  类：\passthrough{\lstinline!unitree::robot::b2::JsonizeModeName!}
\item
  签名：\passthrough{\lstinline!JsonizeModeName()!}
\item
  绑定策略：\passthrough{\lstinline!未单独标注!}
\item
  声明位置：\passthrough{\lstinline!include/unitree/robot/b2/motion\_switcher/motion\_switcher\_api.hpp:37!}
\end{itemize}

\textbf{说明}

\begin{itemize}
\tightlist
\item
  示例是局部调用片段；\passthrough{\lstinline!obj!}
  和参数变量表示已经按上层业务校验并准备好的对象。
\item
  该条目可以被编辑器和类型检查器识别，但当前运行时可能在属性查找阶段直接失败。
\end{itemize}

\textbf{用法}

\begin{lstlisting}[language=Python]
from typing import TYPE_CHECKING

if TYPE_CHECKING:
    def planned_construct() -> JsonizeModeName:
        return JsonizeModeName()
\end{lstlisting}

\begin{quote}
{[}!CAUTION{]} 上面的 \passthrough{\lstinline!TYPE\_CHECKING!}
示例只表达计划调用形态。该分支在普通 Python
运行时为假，不代表当前扩展能执行此接口。
\end{quote}

\hypertarget{unitree-sdk2-cpp-robot-b2-jsonizemodename-from-json-1}{%
\paragraph{\texorpdfstring{\texttt{unitree\_sdk2\_cpp.robot.b2.JsonizeModeName.from\_json}}{unitree\_sdk2\_cpp.robot.b2.JsonizeModeName.from\_json}}\label{unitree-sdk2-cpp-robot-b2-jsonizemodename-from-json-1}}

计划从 JSON 风格字典读取字段并更新当前 SDK 值对象。

\textbf{签名}

\begin{lstlisting}[language=Python]
def from_json(self, json: dict[str, Any]) -> None
\end{lstlisting}

\textbf{可用性与安全性}

\textbf{\passthrough{\lstinline!SIGNATURE\_ONLY!} /
\passthrough{\lstinline!UNCLASSIFIED!}}：当前只有设计期签名，不能假设已存在于运行时扩展。

\textbf{参数}

\begin{longtable}[]{@{}
  >{\raggedright\arraybackslash}p{(\columnwidth - 6\tabcolsep) * \real{0.18}}
  >{\raggedright\arraybackslash}p{(\columnwidth - 6\tabcolsep) * \real{0.24}}
  >{\raggedright\arraybackslash}p{(\columnwidth - 6\tabcolsep) * \real{0.18}}
  >{\raggedright\arraybackslash}p{(\columnwidth - 6\tabcolsep) * \real{0.40}}@{}}
\toprule
名称 & Python 类型 & 默认值 & 含义 \\
\midrule
\endhead
\passthrough{\lstinline!json!} &
\passthrough{\lstinline!dict[str, Any]!} & 必填 & JSON
风格字典，键和值必须满足对应 C++ \passthrough{\lstinline!JsonMap!}
数据结构的约束。 对应 C++ 参数
\passthrough{\lstinline!json: common::JsonMap \&!}。 \\
\bottomrule
\end{longtable}

\textbf{返回值}

\begin{longtable}[]{@{}
  >{\raggedright\arraybackslash}p{(\columnwidth - 4\tabcolsep) * \real{0.18}}
  >{\raggedright\arraybackslash}p{(\columnwidth - 4\tabcolsep) * \real{0.20}}
  >{\raggedright\arraybackslash}p{(\columnwidth - 4\tabcolsep) * \real{0.62}}@{}}
\toprule
位置 & 类型 & 含义 \\
\midrule
\endhead
无 & \passthrough{\lstinline!None!} & 该方法不返回 Python
值；失败可能表现为异常或底层状态变化。 \\
\bottomrule
\end{longtable}

\textbf{对应 C++}

\begin{itemize}
\tightlist
\item
  类：\passthrough{\lstinline!unitree::robot::b2::JsonizeModeName!}
\item
  签名：\passthrough{\lstinline!fromJson(common::JsonMap \&)!}
\item
  绑定策略：\passthrough{\lstinline!SIGNATURE\_PREVIEW!}
\item
  声明位置：\passthrough{\lstinline!include/unitree/robot/b2/motion\_switcher/motion\_switcher\_api.hpp:44!}
\end{itemize}

\textbf{说明}

\begin{itemize}
\tightlist
\item
  示例是局部调用片段；\passthrough{\lstinline!obj!}
  和参数变量表示已经按上层业务校验并准备好的对象。
\item
  该条目可以被编辑器和类型检查器识别，但当前运行时可能在属性查找阶段直接失败。
\end{itemize}

\textbf{用法}

\begin{lstlisting}[language=Python]
from typing import TYPE_CHECKING

if TYPE_CHECKING:
    def planned_from_json(obj: JsonizeModeName, json: dict[str, Any]) -> None:
        obj.from_json(json=json)
\end{lstlisting}

\begin{quote}
{[}!CAUTION{]} 上面的 \passthrough{\lstinline!TYPE\_CHECKING!}
示例只表达计划调用形态。该分支在普通 Python
运行时为假，不代表当前扩展能执行此接口。
\end{quote}

\hypertarget{unitree-sdk2-cpp-robot-b2-jsonizemodename-to-json-1}{%
\paragraph{\texorpdfstring{\texttt{unitree\_sdk2\_cpp.robot.b2.JsonizeModeName.to\_json}}{unitree\_sdk2\_cpp.robot.b2.JsonizeModeName.to\_json}}\label{unitree-sdk2-cpp-robot-b2-jsonizemodename-to-json-1}}

计划把当前 SDK 值对象写入 JSON 风格字典。

\textbf{签名}

\begin{lstlisting}[language=Python]
def to_json(self, json: dict[str, Any]) -> None
\end{lstlisting}

\textbf{可用性与安全性}

\textbf{\passthrough{\lstinline!SIGNATURE\_ONLY!} /
\passthrough{\lstinline!UNCLASSIFIED!}}：当前只有设计期签名，不能假设已存在于运行时扩展。

\textbf{参数}

\begin{longtable}[]{@{}
  >{\raggedright\arraybackslash}p{(\columnwidth - 6\tabcolsep) * \real{0.18}}
  >{\raggedright\arraybackslash}p{(\columnwidth - 6\tabcolsep) * \real{0.24}}
  >{\raggedright\arraybackslash}p{(\columnwidth - 6\tabcolsep) * \real{0.18}}
  >{\raggedright\arraybackslash}p{(\columnwidth - 6\tabcolsep) * \real{0.40}}@{}}
\toprule
名称 & Python 类型 & 默认值 & 含义 \\
\midrule
\endhead
\passthrough{\lstinline!json!} &
\passthrough{\lstinline!dict[str, Any]!} & 必填 & JSON
风格字典，键和值必须满足对应 C++ \passthrough{\lstinline!JsonMap!}
数据结构的约束。 对应 C++ 参数
\passthrough{\lstinline!json: common::JsonMap \&!}。 \\
\bottomrule
\end{longtable}

\textbf{返回值}

\begin{longtable}[]{@{}
  >{\raggedright\arraybackslash}p{(\columnwidth - 4\tabcolsep) * \real{0.18}}
  >{\raggedright\arraybackslash}p{(\columnwidth - 4\tabcolsep) * \real{0.20}}
  >{\raggedright\arraybackslash}p{(\columnwidth - 4\tabcolsep) * \real{0.62}}@{}}
\toprule
位置 & 类型 & 含义 \\
\midrule
\endhead
无 & \passthrough{\lstinline!None!} & 该方法不返回 Python
值；失败可能表现为异常或底层状态变化。 \\
\bottomrule
\end{longtable}

\textbf{对应 C++}

\begin{itemize}
\tightlist
\item
  类：\passthrough{\lstinline!unitree::robot::b2::JsonizeModeName!}
\item
  签名：\passthrough{\lstinline!toJson(common::JsonMap \&) const!}
\item
  绑定策略：\passthrough{\lstinline!SIGNATURE\_PREVIEW!}
\item
  声明位置：\passthrough{\lstinline!include/unitree/robot/b2/motion\_switcher/motion\_switcher\_api.hpp:55!}
\end{itemize}

\textbf{说明}

\begin{itemize}
\tightlist
\item
  示例是局部调用片段；\passthrough{\lstinline!obj!}
  和参数变量表示已经按上层业务校验并准备好的对象。
\item
  该条目可以被编辑器和类型检查器识别，但当前运行时可能在属性查找阶段直接失败。
\end{itemize}

\textbf{用法}

\begin{lstlisting}[language=Python]
from typing import TYPE_CHECKING

if TYPE_CHECKING:
    def planned_to_json(obj: JsonizeModeName, json: dict[str, Any]) -> None:
        obj.to_json(json=json)
\end{lstlisting}

\begin{quote}
{[}!CAUTION{]} 上面的 \passthrough{\lstinline!TYPE\_CHECKING!}
示例只表达计划调用形态。该分支在普通 Python
运行时为假，不代表当前扩展能执行此接口。
\end{quote}

\hypertarget{unitree-sdk2-cpp-robot-b2-jsonizesilent}{%
\subsubsection{\texorpdfstring{\texttt{unitree\_sdk2\_cpp.robot.b2.JsonizeSilent}}{unitree\_sdk2\_cpp.robot.b2.JsonizeSilent}}\label{unitree-sdk2-cpp-robot-b2-jsonizesilent}}

SDK 类型签名预览。请逐项查看构造函数和方法的可用性。

\textbf{导入}

\begin{lstlisting}[language=Python]
from unitree_sdk2_cpp.robot.b2 import JsonizeSilent
\end{lstlisting}

\textbf{基类}：\passthrough{\lstinline!object!}

\textbf{公开属性}

\begin{longtable}[]{@{}
  >{\raggedright\arraybackslash}p{(\columnwidth - 6\tabcolsep) * \real{0.18}}
  >{\raggedright\arraybackslash}p{(\columnwidth - 6\tabcolsep) * \real{0.24}}
  >{\raggedright\arraybackslash}p{(\columnwidth - 6\tabcolsep) * \real{0.18}}
  >{\raggedright\arraybackslash}p{(\columnwidth - 6\tabcolsep) * \real{0.40}}@{}}
\toprule
名称 & Python 类型 & 含义 & 用法 \\
\midrule
\endhead
\passthrough{\lstinline!silent!} & \passthrough{\lstinline!bool!} &
传给该接口的 静音状态 参数，Python 类型为
\passthrough{\lstinline!bool!}。精确含义、单位、范围和枚举值需查目标型号对应头文件与协议。
& \passthrough{\lstinline!value = obj.silent!} /
\passthrough{\lstinline!obj.silent = value!} \\
\bottomrule
\end{longtable}

\textbf{方法索引}

\begin{longtable}[]{@{}
  >{\raggedright\arraybackslash}p{(\columnwidth - 6\tabcolsep) * \real{0.18}}
  >{\raggedright\arraybackslash}p{(\columnwidth - 6\tabcolsep) * \real{0.24}}
  >{\raggedright\arraybackslash}p{(\columnwidth - 6\tabcolsep) * \real{0.18}}
  >{\raggedright\arraybackslash}p{(\columnwidth - 6\tabcolsep) * \real{0.40}}@{}}
\toprule
方法 & Python 签名 & 状态 & 安全分类 \\
\midrule
\endhead
\protect\hyperlink{unitree-sdk2-cpp-robot-b2-jsonizesilent-dunder-init-1}{\passthrough{\lstinline!\_\_init\_\_!}}
& \passthrough{\lstinline!def \_\_init\_\_(self) -> None!} &
\passthrough{\lstinline!SIGNATURE\_ONLY!} &
\passthrough{\lstinline!CONSTRUCTION!} \\
\protect\hyperlink{unitree-sdk2-cpp-robot-b2-jsonizesilent-from-json-1}{\passthrough{\lstinline!from\_json!}}
&
\passthrough{\lstinline!def from\_json(self, json: dict[str, Any]) -> None!}
& \passthrough{\lstinline!SIGNATURE\_ONLY!} &
\passthrough{\lstinline!UNCLASSIFIED!} \\
\protect\hyperlink{unitree-sdk2-cpp-robot-b2-jsonizesilent-to-json-1}{\passthrough{\lstinline!to\_json!}}
&
\passthrough{\lstinline!def to\_json(self, json: dict[str, Any]) -> None!}
& \passthrough{\lstinline!SIGNATURE\_ONLY!} &
\passthrough{\lstinline!UNCLASSIFIED!} \\
\bottomrule
\end{longtable}

\hypertarget{unitree-sdk2-cpp-robot-b2-jsonizesilent-dunder-init-1}{%
\paragraph{\texorpdfstring{\texttt{unitree\_sdk2\_cpp.robot.b2.JsonizeSilent.\_\_init\_\_}}{unitree\_sdk2\_cpp.robot.b2.JsonizeSilent.\_\_init\_\_}}\label{unitree-sdk2-cpp-robot-b2-jsonizesilent-dunder-init-1}}

初始化 \passthrough{\lstinline!JsonizeSilent!}
实例。是否能实际构造取决于下方可用性状态。

\textbf{签名}

\begin{lstlisting}[language=Python]
def __init__(self) -> None
\end{lstlisting}

\textbf{可用性与安全性}

\textbf{\passthrough{\lstinline!SIGNATURE\_ONLY!} /
\passthrough{\lstinline!CONSTRUCTION!}}：当前只有设计期签名，不能假设已存在于运行时扩展。

\textbf{参数}

无显式参数。实例方法中的 \passthrough{\lstinline!self!} 由 Python
自动传入。

\textbf{返回值}

\begin{longtable}[]{@{}
  >{\raggedright\arraybackslash}p{(\columnwidth - 4\tabcolsep) * \real{0.18}}
  >{\raggedright\arraybackslash}p{(\columnwidth - 4\tabcolsep) * \real{0.20}}
  >{\raggedright\arraybackslash}p{(\columnwidth - 4\tabcolsep) * \real{0.62}}@{}}
\toprule
位置 & 类型 & 含义 \\
\midrule
\endhead
无 & \passthrough{\lstinline!None!} &
\passthrough{\lstinline!\_\_init\_\_!}
本身不返回值；调用类对象时，在构造函数可用的前提下得到该类实例。 \\
\bottomrule
\end{longtable}

\textbf{对应 C++}

\begin{itemize}
\tightlist
\item
  类：\passthrough{\lstinline!unitree::robot::b2::JsonizeSilent!}
\item
  签名：\passthrough{\lstinline!JsonizeSilent()!}
\item
  绑定策略：\passthrough{\lstinline!未单独标注!}
\item
  声明位置：\passthrough{\lstinline!include/unitree/robot/b2/motion\_switcher/motion\_switcher\_api.hpp:73!}
\end{itemize}

\textbf{说明}

\begin{itemize}
\tightlist
\item
  示例是局部调用片段；\passthrough{\lstinline!obj!}
  和参数变量表示已经按上层业务校验并准备好的对象。
\item
  该条目可以被编辑器和类型检查器识别，但当前运行时可能在属性查找阶段直接失败。
\end{itemize}

\textbf{用法}

\begin{lstlisting}[language=Python]
from typing import TYPE_CHECKING

if TYPE_CHECKING:
    def planned_construct() -> JsonizeSilent:
        return JsonizeSilent()
\end{lstlisting}

\begin{quote}
{[}!CAUTION{]} 上面的 \passthrough{\lstinline!TYPE\_CHECKING!}
示例只表达计划调用形态。该分支在普通 Python
运行时为假，不代表当前扩展能执行此接口。
\end{quote}

\hypertarget{unitree-sdk2-cpp-robot-b2-jsonizesilent-from-json-1}{%
\paragraph{\texorpdfstring{\texttt{unitree\_sdk2\_cpp.robot.b2.JsonizeSilent.from\_json}}{unitree\_sdk2\_cpp.robot.b2.JsonizeSilent.from\_json}}\label{unitree-sdk2-cpp-robot-b2-jsonizesilent-from-json-1}}

计划从 JSON 风格字典读取字段并更新当前 SDK 值对象。

\textbf{签名}

\begin{lstlisting}[language=Python]
def from_json(self, json: dict[str, Any]) -> None
\end{lstlisting}

\textbf{可用性与安全性}

\textbf{\passthrough{\lstinline!SIGNATURE\_ONLY!} /
\passthrough{\lstinline!UNCLASSIFIED!}}：当前只有设计期签名，不能假设已存在于运行时扩展。

\textbf{参数}

\begin{longtable}[]{@{}
  >{\raggedright\arraybackslash}p{(\columnwidth - 6\tabcolsep) * \real{0.18}}
  >{\raggedright\arraybackslash}p{(\columnwidth - 6\tabcolsep) * \real{0.24}}
  >{\raggedright\arraybackslash}p{(\columnwidth - 6\tabcolsep) * \real{0.18}}
  >{\raggedright\arraybackslash}p{(\columnwidth - 6\tabcolsep) * \real{0.40}}@{}}
\toprule
名称 & Python 类型 & 默认值 & 含义 \\
\midrule
\endhead
\passthrough{\lstinline!json!} &
\passthrough{\lstinline!dict[str, Any]!} & 必填 & JSON
风格字典，键和值必须满足对应 C++ \passthrough{\lstinline!JsonMap!}
数据结构的约束。 对应 C++ 参数
\passthrough{\lstinline!json: common::JsonMap \&!}。 \\
\bottomrule
\end{longtable}

\textbf{返回值}

\begin{longtable}[]{@{}
  >{\raggedright\arraybackslash}p{(\columnwidth - 4\tabcolsep) * \real{0.18}}
  >{\raggedright\arraybackslash}p{(\columnwidth - 4\tabcolsep) * \real{0.20}}
  >{\raggedright\arraybackslash}p{(\columnwidth - 4\tabcolsep) * \real{0.62}}@{}}
\toprule
位置 & 类型 & 含义 \\
\midrule
\endhead
无 & \passthrough{\lstinline!None!} & 该方法不返回 Python
值；失败可能表现为异常或底层状态变化。 \\
\bottomrule
\end{longtable}

\textbf{对应 C++}

\begin{itemize}
\tightlist
\item
  类：\passthrough{\lstinline!unitree::robot::b2::JsonizeSilent!}
\item
  签名：\passthrough{\lstinline!fromJson(common::JsonMap \&)!}
\item
  绑定策略：\passthrough{\lstinline!SIGNATURE\_PREVIEW!}
\item
  声明位置：\passthrough{\lstinline!include/unitree/robot/b2/motion\_switcher/motion\_switcher\_api.hpp:80!}
\end{itemize}

\textbf{说明}

\begin{itemize}
\tightlist
\item
  示例是局部调用片段；\passthrough{\lstinline!obj!}
  和参数变量表示已经按上层业务校验并准备好的对象。
\item
  该条目可以被编辑器和类型检查器识别，但当前运行时可能在属性查找阶段直接失败。
\end{itemize}

\textbf{用法}

\begin{lstlisting}[language=Python]
from typing import TYPE_CHECKING

if TYPE_CHECKING:
    def planned_from_json(obj: JsonizeSilent, json: dict[str, Any]) -> None:
        obj.from_json(json=json)
\end{lstlisting}

\begin{quote}
{[}!CAUTION{]} 上面的 \passthrough{\lstinline!TYPE\_CHECKING!}
示例只表达计划调用形态。该分支在普通 Python
运行时为假，不代表当前扩展能执行此接口。
\end{quote}

\hypertarget{unitree-sdk2-cpp-robot-b2-jsonizesilent-to-json-1}{%
\paragraph{\texorpdfstring{\texttt{unitree\_sdk2\_cpp.robot.b2.JsonizeSilent.to\_json}}{unitree\_sdk2\_cpp.robot.b2.JsonizeSilent.to\_json}}\label{unitree-sdk2-cpp-robot-b2-jsonizesilent-to-json-1}}

计划把当前 SDK 值对象写入 JSON 风格字典。

\textbf{签名}

\begin{lstlisting}[language=Python]
def to_json(self, json: dict[str, Any]) -> None
\end{lstlisting}

\textbf{可用性与安全性}

\textbf{\passthrough{\lstinline!SIGNATURE\_ONLY!} /
\passthrough{\lstinline!UNCLASSIFIED!}}：当前只有设计期签名，不能假设已存在于运行时扩展。

\textbf{参数}

\begin{longtable}[]{@{}
  >{\raggedright\arraybackslash}p{(\columnwidth - 6\tabcolsep) * \real{0.18}}
  >{\raggedright\arraybackslash}p{(\columnwidth - 6\tabcolsep) * \real{0.24}}
  >{\raggedright\arraybackslash}p{(\columnwidth - 6\tabcolsep) * \real{0.18}}
  >{\raggedright\arraybackslash}p{(\columnwidth - 6\tabcolsep) * \real{0.40}}@{}}
\toprule
名称 & Python 类型 & 默认值 & 含义 \\
\midrule
\endhead
\passthrough{\lstinline!json!} &
\passthrough{\lstinline!dict[str, Any]!} & 必填 & JSON
风格字典，键和值必须满足对应 C++ \passthrough{\lstinline!JsonMap!}
数据结构的约束。 对应 C++ 参数
\passthrough{\lstinline!json: common::JsonMap \&!}。 \\
\bottomrule
\end{longtable}

\textbf{返回值}

\begin{longtable}[]{@{}
  >{\raggedright\arraybackslash}p{(\columnwidth - 4\tabcolsep) * \real{0.18}}
  >{\raggedright\arraybackslash}p{(\columnwidth - 4\tabcolsep) * \real{0.20}}
  >{\raggedright\arraybackslash}p{(\columnwidth - 4\tabcolsep) * \real{0.62}}@{}}
\toprule
位置 & 类型 & 含义 \\
\midrule
\endhead
无 & \passthrough{\lstinline!None!} & 该方法不返回 Python
值；失败可能表现为异常或底层状态变化。 \\
\bottomrule
\end{longtable}

\textbf{对应 C++}

\begin{itemize}
\tightlist
\item
  类：\passthrough{\lstinline!unitree::robot::b2::JsonizeSilent!}
\item
  签名：\passthrough{\lstinline!toJson(common::JsonMap \&) const!}
\item
  绑定策略：\passthrough{\lstinline!SIGNATURE\_PREVIEW!}
\item
  声明位置：\passthrough{\lstinline!include/unitree/robot/b2/motion\_switcher/motion\_switcher\_api.hpp:86!}
\end{itemize}

\textbf{说明}

\begin{itemize}
\tightlist
\item
  示例是局部调用片段；\passthrough{\lstinline!obj!}
  和参数变量表示已经按上层业务校验并准备好的对象。
\item
  该条目可以被编辑器和类型检查器识别，但当前运行时可能在属性查找阶段直接失败。
\end{itemize}

\textbf{用法}

\begin{lstlisting}[language=Python]
from typing import TYPE_CHECKING

if TYPE_CHECKING:
    def planned_to_json(obj: JsonizeSilent, json: dict[str, Any]) -> None:
        obj.to_json(json=json)
\end{lstlisting}

\begin{quote}
{[}!CAUTION{]} 上面的 \passthrough{\lstinline!TYPE\_CHECKING!}
示例只表达计划调用形态。该分支在普通 Python
运行时为假，不代表当前扩展能执行此接口。
\end{quote}

\hypertarget{unitree-sdk2-cpp-robot-b2-lowpowerstatusdata}{%
\subsubsection{\texorpdfstring{\texttt{unitree\_sdk2\_cpp.robot.b2.LowPowerStatusData}}{unitree\_sdk2\_cpp.robot.b2.LowPowerStatusData}}\label{unitree-sdk2-cpp-robot-b2-lowpowerstatusdata}}

SDK 数据值类型；公开字段和转换方法见下文。

\textbf{导入}

\begin{lstlisting}[language=Python]
from unitree_sdk2_cpp.robot.b2 import LowPowerStatusData
\end{lstlisting}

\textbf{基类}：\passthrough{\lstinline!object!}

\textbf{公开属性}

\begin{longtable}[]{@{}
  >{\raggedright\arraybackslash}p{(\columnwidth - 6\tabcolsep) * \real{0.18}}
  >{\raggedright\arraybackslash}p{(\columnwidth - 6\tabcolsep) * \real{0.24}}
  >{\raggedright\arraybackslash}p{(\columnwidth - 6\tabcolsep) * \real{0.18}}
  >{\raggedright\arraybackslash}p{(\columnwidth - 6\tabcolsep) * \real{0.40}}@{}}
\toprule
名称 & Python 类型 & 含义 & 用法 \\
\midrule
\endhead
\passthrough{\lstinline!status!} & \passthrough{\lstinline!int!} &
传给该接口的 状态 参数，Python 类型为
\passthrough{\lstinline!int!}。精确含义、单位、范围和枚举值需查目标型号对应头文件与协议。
& \passthrough{\lstinline!value = obj.status!} /
\passthrough{\lstinline!obj.status = value!} \\
\bottomrule
\end{longtable}

\textbf{方法索引}

\begin{longtable}[]{@{}
  >{\raggedright\arraybackslash}p{(\columnwidth - 6\tabcolsep) * \real{0.18}}
  >{\raggedright\arraybackslash}p{(\columnwidth - 6\tabcolsep) * \real{0.24}}
  >{\raggedright\arraybackslash}p{(\columnwidth - 6\tabcolsep) * \real{0.18}}
  >{\raggedright\arraybackslash}p{(\columnwidth - 6\tabcolsep) * \real{0.40}}@{}}
\toprule
方法 & Python 签名 & 状态 & 安全分类 \\
\midrule
\endhead
\protect\hyperlink{unitree-sdk2-cpp-robot-b2-lowpowerstatusdata-dunder-init-1}{\passthrough{\lstinline!\_\_init\_\_!}}
& \passthrough{\lstinline!def \_\_init\_\_(self) -> None!} &
\passthrough{\lstinline!SIGNATURE\_ONLY!} &
\passthrough{\lstinline!CONSTRUCTION!} \\
\protect\hyperlink{unitree-sdk2-cpp-robot-b2-lowpowerstatusdata-from-json-1}{\passthrough{\lstinline!from\_json!}}
&
\passthrough{\lstinline!def from\_json(self, json: dict[str, Any]) -> None!}
& \passthrough{\lstinline!SIGNATURE\_ONLY!} &
\passthrough{\lstinline!UNCLASSIFIED!} \\
\protect\hyperlink{unitree-sdk2-cpp-robot-b2-lowpowerstatusdata-to-json-1}{\passthrough{\lstinline!to\_json!}}
&
\passthrough{\lstinline!def to\_json(self, json: dict[str, Any]) -> None!}
& \passthrough{\lstinline!SIGNATURE\_ONLY!} &
\passthrough{\lstinline!UNCLASSIFIED!} \\
\bottomrule
\end{longtable}

\hypertarget{unitree-sdk2-cpp-robot-b2-lowpowerstatusdata-dunder-init-1}{%
\paragraph{\texorpdfstring{\texttt{unitree\_sdk2\_cpp.robot.b2.LowPowerStatusData.\_\_init\_\_}}{unitree\_sdk2\_cpp.robot.b2.LowPowerStatusData.\_\_init\_\_}}\label{unitree-sdk2-cpp-robot-b2-lowpowerstatusdata-dunder-init-1}}

初始化 \passthrough{\lstinline!LowPowerStatusData!}
实例。是否能实际构造取决于下方可用性状态。

\textbf{签名}

\begin{lstlisting}[language=Python]
def __init__(self) -> None
\end{lstlisting}

\textbf{可用性与安全性}

\textbf{\passthrough{\lstinline!SIGNATURE\_ONLY!} /
\passthrough{\lstinline!CONSTRUCTION!}}：当前只有设计期签名，不能假设已存在于运行时扩展。

\textbf{参数}

无显式参数。实例方法中的 \passthrough{\lstinline!self!} 由 Python
自动传入。

\textbf{返回值}

\begin{longtable}[]{@{}
  >{\raggedright\arraybackslash}p{(\columnwidth - 4\tabcolsep) * \real{0.18}}
  >{\raggedright\arraybackslash}p{(\columnwidth - 4\tabcolsep) * \real{0.20}}
  >{\raggedright\arraybackslash}p{(\columnwidth - 4\tabcolsep) * \real{0.62}}@{}}
\toprule
位置 & 类型 & 含义 \\
\midrule
\endhead
无 & \passthrough{\lstinline!None!} &
\passthrough{\lstinline!\_\_init\_\_!}
本身不返回值；调用类对象时，在构造函数可用的前提下得到该类实例。 \\
\bottomrule
\end{longtable}

\textbf{对应 C++}

\begin{itemize}
\tightlist
\item
  类：\passthrough{\lstinline!unitree::robot::b2::LowPowerStatusData!}
\item
  签名：\passthrough{\lstinline!LowPowerStatusData()!}
\item
  绑定策略：\passthrough{\lstinline!未单独标注!}
\item
  声明位置：\passthrough{\lstinline!include/unitree/robot/b2/robot\_state/robot\_state\_api.hpp:183!}
\end{itemize}

\textbf{说明}

\begin{itemize}
\tightlist
\item
  示例是局部调用片段；\passthrough{\lstinline!obj!}
  和参数变量表示已经按上层业务校验并准备好的对象。
\item
  该条目可以被编辑器和类型检查器识别，但当前运行时可能在属性查找阶段直接失败。
\end{itemize}

\textbf{用法}

\begin{lstlisting}[language=Python]
from typing import TYPE_CHECKING

if TYPE_CHECKING:
    def planned_construct() -> LowPowerStatusData:
        return LowPowerStatusData()
\end{lstlisting}

\begin{quote}
{[}!CAUTION{]} 上面的 \passthrough{\lstinline!TYPE\_CHECKING!}
示例只表达计划调用形态。该分支在普通 Python
运行时为假，不代表当前扩展能执行此接口。
\end{quote}

\hypertarget{unitree-sdk2-cpp-robot-b2-lowpowerstatusdata-from-json-1}{%
\paragraph{\texorpdfstring{\texttt{unitree\_sdk2\_cpp.robot.b2.LowPowerStatusData.from\_json}}{unitree\_sdk2\_cpp.robot.b2.LowPowerStatusData.from\_json}}\label{unitree-sdk2-cpp-robot-b2-lowpowerstatusdata-from-json-1}}

计划从 JSON 风格字典读取字段并更新当前 SDK 值对象。

\textbf{签名}

\begin{lstlisting}[language=Python]
def from_json(self, json: dict[str, Any]) -> None
\end{lstlisting}

\textbf{可用性与安全性}

\textbf{\passthrough{\lstinline!SIGNATURE\_ONLY!} /
\passthrough{\lstinline!UNCLASSIFIED!}}：当前只有设计期签名，不能假设已存在于运行时扩展。

\textbf{参数}

\begin{longtable}[]{@{}
  >{\raggedright\arraybackslash}p{(\columnwidth - 6\tabcolsep) * \real{0.18}}
  >{\raggedright\arraybackslash}p{(\columnwidth - 6\tabcolsep) * \real{0.24}}
  >{\raggedright\arraybackslash}p{(\columnwidth - 6\tabcolsep) * \real{0.18}}
  >{\raggedright\arraybackslash}p{(\columnwidth - 6\tabcolsep) * \real{0.40}}@{}}
\toprule
名称 & Python 类型 & 默认值 & 含义 \\
\midrule
\endhead
\passthrough{\lstinline!json!} &
\passthrough{\lstinline!dict[str, Any]!} & 必填 & JSON
风格字典，键和值必须满足对应 C++ \passthrough{\lstinline!JsonMap!}
数据结构的约束。 对应 C++ 参数
\passthrough{\lstinline!json: common::JsonMap \&!}。 \\
\bottomrule
\end{longtable}

\textbf{返回值}

\begin{longtable}[]{@{}
  >{\raggedright\arraybackslash}p{(\columnwidth - 4\tabcolsep) * \real{0.18}}
  >{\raggedright\arraybackslash}p{(\columnwidth - 4\tabcolsep) * \real{0.20}}
  >{\raggedright\arraybackslash}p{(\columnwidth - 4\tabcolsep) * \real{0.62}}@{}}
\toprule
位置 & 类型 & 含义 \\
\midrule
\endhead
无 & \passthrough{\lstinline!None!} & 该方法不返回 Python
值；失败可能表现为异常或底层状态变化。 \\
\bottomrule
\end{longtable}

\textbf{对应 C++}

\begin{itemize}
\tightlist
\item
  类：\passthrough{\lstinline!unitree::robot::b2::LowPowerStatusData!}
\item
  签名：\passthrough{\lstinline!fromJson(common::JsonMap \&)!}
\item
  绑定策略：\passthrough{\lstinline!SIGNATURE\_PREVIEW!}
\item
  声明位置：\passthrough{\lstinline!include/unitree/robot/b2/robot\_state/robot\_state\_api.hpp:189!}
\end{itemize}

\textbf{说明}

\begin{itemize}
\tightlist
\item
  示例是局部调用片段；\passthrough{\lstinline!obj!}
  和参数变量表示已经按上层业务校验并准备好的对象。
\item
  该条目可以被编辑器和类型检查器识别，但当前运行时可能在属性查找阶段直接失败。
\end{itemize}

\textbf{用法}

\begin{lstlisting}[language=Python]
from typing import TYPE_CHECKING

if TYPE_CHECKING:
    def planned_from_json(obj: LowPowerStatusData, json: dict[str, Any]) -> None:
        obj.from_json(json=json)
\end{lstlisting}

\begin{quote}
{[}!CAUTION{]} 上面的 \passthrough{\lstinline!TYPE\_CHECKING!}
示例只表达计划调用形态。该分支在普通 Python
运行时为假，不代表当前扩展能执行此接口。
\end{quote}

\hypertarget{unitree-sdk2-cpp-robot-b2-lowpowerstatusdata-to-json-1}{%
\paragraph{\texorpdfstring{\texttt{unitree\_sdk2\_cpp.robot.b2.LowPowerStatusData.to\_json}}{unitree\_sdk2\_cpp.robot.b2.LowPowerStatusData.to\_json}}\label{unitree-sdk2-cpp-robot-b2-lowpowerstatusdata-to-json-1}}

计划把当前 SDK 值对象写入 JSON 风格字典。

\textbf{签名}

\begin{lstlisting}[language=Python]
def to_json(self, json: dict[str, Any]) -> None
\end{lstlisting}

\textbf{可用性与安全性}

\textbf{\passthrough{\lstinline!SIGNATURE\_ONLY!} /
\passthrough{\lstinline!UNCLASSIFIED!}}：当前只有设计期签名，不能假设已存在于运行时扩展。

\textbf{参数}

\begin{longtable}[]{@{}
  >{\raggedright\arraybackslash}p{(\columnwidth - 6\tabcolsep) * \real{0.18}}
  >{\raggedright\arraybackslash}p{(\columnwidth - 6\tabcolsep) * \real{0.24}}
  >{\raggedright\arraybackslash}p{(\columnwidth - 6\tabcolsep) * \real{0.18}}
  >{\raggedright\arraybackslash}p{(\columnwidth - 6\tabcolsep) * \real{0.40}}@{}}
\toprule
名称 & Python 类型 & 默认值 & 含义 \\
\midrule
\endhead
\passthrough{\lstinline!json!} &
\passthrough{\lstinline!dict[str, Any]!} & 必填 & JSON
风格字典，键和值必须满足对应 C++ \passthrough{\lstinline!JsonMap!}
数据结构的约束。 对应 C++ 参数
\passthrough{\lstinline!json: common::JsonMap \&!}。 \\
\bottomrule
\end{longtable}

\textbf{返回值}

\begin{longtable}[]{@{}
  >{\raggedright\arraybackslash}p{(\columnwidth - 4\tabcolsep) * \real{0.18}}
  >{\raggedright\arraybackslash}p{(\columnwidth - 4\tabcolsep) * \real{0.20}}
  >{\raggedright\arraybackslash}p{(\columnwidth - 4\tabcolsep) * \real{0.62}}@{}}
\toprule
位置 & 类型 & 含义 \\
\midrule
\endhead
无 & \passthrough{\lstinline!None!} & 该方法不返回 Python
值；失败可能表现为异常或底层状态变化。 \\
\bottomrule
\end{longtable}

\textbf{对应 C++}

\begin{itemize}
\tightlist
\item
  类：\passthrough{\lstinline!unitree::robot::b2::LowPowerStatusData!}
\item
  签名：\passthrough{\lstinline!toJson(common::JsonMap \&) const!}
\item
  绑定策略：\passthrough{\lstinline!SIGNATURE\_PREVIEW!}
\item
  声明位置：\passthrough{\lstinline!include/unitree/robot/b2/robot\_state/robot\_state\_api.hpp:194!}
\end{itemize}

\textbf{说明}

\begin{itemize}
\tightlist
\item
  示例是局部调用片段；\passthrough{\lstinline!obj!}
  和参数变量表示已经按上层业务校验并准备好的对象。
\item
  该条目可以被编辑器和类型检查器识别，但当前运行时可能在属性查找阶段直接失败。
\end{itemize}

\textbf{用法}

\begin{lstlisting}[language=Python]
from typing import TYPE_CHECKING

if TYPE_CHECKING:
    def planned_to_json(obj: LowPowerStatusData, json: dict[str, Any]) -> None:
        obj.to_json(json=json)
\end{lstlisting}

\begin{quote}
{[}!CAUTION{]} 上面的 \passthrough{\lstinline!TYPE\_CHECKING!}
示例只表达计划调用形态。该分支在普通 Python
运行时为假，不代表当前扩展能执行此接口。
\end{quote}

\hypertarget{unitree-sdk2-cpp-robot-b2-lowpowerswitchparameter}{%
\subsubsection{\texorpdfstring{\texttt{unitree\_sdk2\_cpp.robot.b2.LowPowerSwitchParameter}}{unitree\_sdk2\_cpp.robot.b2.LowPowerSwitchParameter}}\label{unitree-sdk2-cpp-robot-b2-lowpowerswitchparameter}}

SDK 请求参数值类型；当前可能仅提供设计期签名。

\textbf{导入}

\begin{lstlisting}[language=Python]
from unitree_sdk2_cpp.robot.b2 import LowPowerSwitchParameter
\end{lstlisting}

\textbf{基类}：\passthrough{\lstinline!object!}

\textbf{公开属性}

\begin{longtable}[]{@{}
  >{\raggedright\arraybackslash}p{(\columnwidth - 6\tabcolsep) * \real{0.18}}
  >{\raggedright\arraybackslash}p{(\columnwidth - 6\tabcolsep) * \real{0.24}}
  >{\raggedright\arraybackslash}p{(\columnwidth - 6\tabcolsep) * \real{0.18}}
  >{\raggedright\arraybackslash}p{(\columnwidth - 6\tabcolsep) * \real{0.40}}@{}}
\toprule
名称 & Python 类型 & 含义 & 用法 \\
\midrule
\endhead
\passthrough{\lstinline!swit!} & \passthrough{\lstinline!int!} &
传给该接口的 \passthrough{\lstinline!swit!} 参数，Python 类型为
\passthrough{\lstinline!int!}。精确含义、单位、范围和枚举值需查目标型号对应头文件与协议。
& \passthrough{\lstinline!value = obj.swit!} /
\passthrough{\lstinline!obj.swit = value!} \\
\bottomrule
\end{longtable}

\textbf{方法索引}

\begin{longtable}[]{@{}
  >{\raggedright\arraybackslash}p{(\columnwidth - 6\tabcolsep) * \real{0.18}}
  >{\raggedright\arraybackslash}p{(\columnwidth - 6\tabcolsep) * \real{0.24}}
  >{\raggedright\arraybackslash}p{(\columnwidth - 6\tabcolsep) * \real{0.18}}
  >{\raggedright\arraybackslash}p{(\columnwidth - 6\tabcolsep) * \real{0.40}}@{}}
\toprule
方法 & Python 签名 & 状态 & 安全分类 \\
\midrule
\endhead
\protect\hyperlink{unitree-sdk2-cpp-robot-b2-lowpowerswitchparameter-dunder-init-1}{\passthrough{\lstinline!\_\_init\_\_!}}
& \passthrough{\lstinline!def \_\_init\_\_(self) -> None!} &
\passthrough{\lstinline!SIGNATURE\_ONLY!} &
\passthrough{\lstinline!CONSTRUCTION!} \\
\protect\hyperlink{unitree-sdk2-cpp-robot-b2-lowpowerswitchparameter-from-json-1}{\passthrough{\lstinline!from\_json!}}
&
\passthrough{\lstinline!def from\_json(self, json: dict[str, Any]) -> None!}
& \passthrough{\lstinline!SIGNATURE\_ONLY!} &
\passthrough{\lstinline!UNCLASSIFIED!} \\
\protect\hyperlink{unitree-sdk2-cpp-robot-b2-lowpowerswitchparameter-to-json-1}{\passthrough{\lstinline!to\_json!}}
&
\passthrough{\lstinline!def to\_json(self, json: dict[str, Any]) -> None!}
& \passthrough{\lstinline!SIGNATURE\_ONLY!} &
\passthrough{\lstinline!UNCLASSIFIED!} \\
\bottomrule
\end{longtable}

\hypertarget{unitree-sdk2-cpp-robot-b2-lowpowerswitchparameter-dunder-init-1}{%
\paragraph{\texorpdfstring{\texttt{unitree\_sdk2\_cpp.robot.b2.LowPowerSwitchParameter.\_\_init\_\_}}{unitree\_sdk2\_cpp.robot.b2.LowPowerSwitchParameter.\_\_init\_\_}}\label{unitree-sdk2-cpp-robot-b2-lowpowerswitchparameter-dunder-init-1}}

初始化 \passthrough{\lstinline!LowPowerSwitchParameter!}
实例。是否能实际构造取决于下方可用性状态。

\textbf{签名}

\begin{lstlisting}[language=Python]
def __init__(self) -> None
\end{lstlisting}

\textbf{可用性与安全性}

\textbf{\passthrough{\lstinline!SIGNATURE\_ONLY!} /
\passthrough{\lstinline!CONSTRUCTION!}}：当前只有设计期签名，不能假设已存在于运行时扩展。

\textbf{参数}

无显式参数。实例方法中的 \passthrough{\lstinline!self!} 由 Python
自动传入。

\textbf{返回值}

\begin{longtable}[]{@{}
  >{\raggedright\arraybackslash}p{(\columnwidth - 4\tabcolsep) * \real{0.18}}
  >{\raggedright\arraybackslash}p{(\columnwidth - 4\tabcolsep) * \real{0.20}}
  >{\raggedright\arraybackslash}p{(\columnwidth - 4\tabcolsep) * \real{0.62}}@{}}
\toprule
位置 & 类型 & 含义 \\
\midrule
\endhead
无 & \passthrough{\lstinline!None!} &
\passthrough{\lstinline!\_\_init\_\_!}
本身不返回值；调用类对象时，在构造函数可用的前提下得到该类实例。 \\
\bottomrule
\end{longtable}

\textbf{对应 C++}

\begin{itemize}
\tightlist
\item
  类：\passthrough{\lstinline!unitree::robot::b2::LowPowerSwitchParameter!}
\item
  签名：\passthrough{\lstinline!LowPowerSwitchParameter()!}
\item
  绑定策略：\passthrough{\lstinline!未单独标注!}
\item
  声明位置：\passthrough{\lstinline!include/unitree/robot/b2/robot\_state/robot\_state\_api.hpp:157!}
\end{itemize}

\textbf{说明}

\begin{itemize}
\tightlist
\item
  示例是局部调用片段；\passthrough{\lstinline!obj!}
  和参数变量表示已经按上层业务校验并准备好的对象。
\item
  该条目可以被编辑器和类型检查器识别，但当前运行时可能在属性查找阶段直接失败。
\end{itemize}

\textbf{用法}

\begin{lstlisting}[language=Python]
from typing import TYPE_CHECKING

if TYPE_CHECKING:
    def planned_construct() -> LowPowerSwitchParameter:
        return LowPowerSwitchParameter()
\end{lstlisting}

\begin{quote}
{[}!CAUTION{]} 上面的 \passthrough{\lstinline!TYPE\_CHECKING!}
示例只表达计划调用形态。该分支在普通 Python
运行时为假，不代表当前扩展能执行此接口。
\end{quote}

\hypertarget{unitree-sdk2-cpp-robot-b2-lowpowerswitchparameter-from-json-1}{%
\paragraph{\texorpdfstring{\texttt{unitree\_sdk2\_cpp.robot.b2.LowPowerSwitchParameter.from\_json}}{unitree\_sdk2\_cpp.robot.b2.LowPowerSwitchParameter.from\_json}}\label{unitree-sdk2-cpp-robot-b2-lowpowerswitchparameter-from-json-1}}

计划从 JSON 风格字典读取字段并更新当前 SDK 值对象。

\textbf{签名}

\begin{lstlisting}[language=Python]
def from_json(self, json: dict[str, Any]) -> None
\end{lstlisting}

\textbf{可用性与安全性}

\textbf{\passthrough{\lstinline!SIGNATURE\_ONLY!} /
\passthrough{\lstinline!UNCLASSIFIED!}}：当前只有设计期签名，不能假设已存在于运行时扩展。

\textbf{参数}

\begin{longtable}[]{@{}
  >{\raggedright\arraybackslash}p{(\columnwidth - 6\tabcolsep) * \real{0.18}}
  >{\raggedright\arraybackslash}p{(\columnwidth - 6\tabcolsep) * \real{0.24}}
  >{\raggedright\arraybackslash}p{(\columnwidth - 6\tabcolsep) * \real{0.18}}
  >{\raggedright\arraybackslash}p{(\columnwidth - 6\tabcolsep) * \real{0.40}}@{}}
\toprule
名称 & Python 类型 & 默认值 & 含义 \\
\midrule
\endhead
\passthrough{\lstinline!json!} &
\passthrough{\lstinline!dict[str, Any]!} & 必填 & JSON
风格字典，键和值必须满足对应 C++ \passthrough{\lstinline!JsonMap!}
数据结构的约束。 对应 C++ 参数
\passthrough{\lstinline!json: common::JsonMap \&!}。 \\
\bottomrule
\end{longtable}

\textbf{返回值}

\begin{longtable}[]{@{}
  >{\raggedright\arraybackslash}p{(\columnwidth - 4\tabcolsep) * \real{0.18}}
  >{\raggedright\arraybackslash}p{(\columnwidth - 4\tabcolsep) * \real{0.20}}
  >{\raggedright\arraybackslash}p{(\columnwidth - 4\tabcolsep) * \real{0.62}}@{}}
\toprule
位置 & 类型 & 含义 \\
\midrule
\endhead
无 & \passthrough{\lstinline!None!} & 该方法不返回 Python
值；失败可能表现为异常或底层状态变化。 \\
\bottomrule
\end{longtable}

\textbf{对应 C++}

\begin{itemize}
\tightlist
\item
  类：\passthrough{\lstinline!unitree::robot::b2::LowPowerSwitchParameter!}
\item
  签名：\passthrough{\lstinline!fromJson(common::JsonMap \&)!}
\item
  绑定策略：\passthrough{\lstinline!SIGNATURE\_PREVIEW!}
\item
  声明位置：\passthrough{\lstinline!include/unitree/robot/b2/robot\_state/robot\_state\_api.hpp:163!}
\end{itemize}

\textbf{说明}

\begin{itemize}
\tightlist
\item
  示例是局部调用片段；\passthrough{\lstinline!obj!}
  和参数变量表示已经按上层业务校验并准备好的对象。
\item
  该条目可以被编辑器和类型检查器识别，但当前运行时可能在属性查找阶段直接失败。
\end{itemize}

\textbf{用法}

\begin{lstlisting}[language=Python]
from typing import TYPE_CHECKING

if TYPE_CHECKING:
    def planned_from_json(obj: LowPowerSwitchParameter, json: dict[str, Any]) -> None:
        obj.from_json(json=json)
\end{lstlisting}

\begin{quote}
{[}!CAUTION{]} 上面的 \passthrough{\lstinline!TYPE\_CHECKING!}
示例只表达计划调用形态。该分支在普通 Python
运行时为假，不代表当前扩展能执行此接口。
\end{quote}

\hypertarget{unitree-sdk2-cpp-robot-b2-lowpowerswitchparameter-to-json-1}{%
\paragraph{\texorpdfstring{\texttt{unitree\_sdk2\_cpp.robot.b2.LowPowerSwitchParameter.to\_json}}{unitree\_sdk2\_cpp.robot.b2.LowPowerSwitchParameter.to\_json}}\label{unitree-sdk2-cpp-robot-b2-lowpowerswitchparameter-to-json-1}}

计划把当前 SDK 值对象写入 JSON 风格字典。

\textbf{签名}

\begin{lstlisting}[language=Python]
def to_json(self, json: dict[str, Any]) -> None
\end{lstlisting}

\textbf{可用性与安全性}

\textbf{\passthrough{\lstinline!SIGNATURE\_ONLY!} /
\passthrough{\lstinline!UNCLASSIFIED!}}：当前只有设计期签名，不能假设已存在于运行时扩展。

\textbf{参数}

\begin{longtable}[]{@{}
  >{\raggedright\arraybackslash}p{(\columnwidth - 6\tabcolsep) * \real{0.18}}
  >{\raggedright\arraybackslash}p{(\columnwidth - 6\tabcolsep) * \real{0.24}}
  >{\raggedright\arraybackslash}p{(\columnwidth - 6\tabcolsep) * \real{0.18}}
  >{\raggedright\arraybackslash}p{(\columnwidth - 6\tabcolsep) * \real{0.40}}@{}}
\toprule
名称 & Python 类型 & 默认值 & 含义 \\
\midrule
\endhead
\passthrough{\lstinline!json!} &
\passthrough{\lstinline!dict[str, Any]!} & 必填 & JSON
风格字典，键和值必须满足对应 C++ \passthrough{\lstinline!JsonMap!}
数据结构的约束。 对应 C++ 参数
\passthrough{\lstinline!json: common::JsonMap \&!}。 \\
\bottomrule
\end{longtable}

\textbf{返回值}

\begin{longtable}[]{@{}
  >{\raggedright\arraybackslash}p{(\columnwidth - 4\tabcolsep) * \real{0.18}}
  >{\raggedright\arraybackslash}p{(\columnwidth - 4\tabcolsep) * \real{0.20}}
  >{\raggedright\arraybackslash}p{(\columnwidth - 4\tabcolsep) * \real{0.62}}@{}}
\toprule
位置 & 类型 & 含义 \\
\midrule
\endhead
无 & \passthrough{\lstinline!None!} & 该方法不返回 Python
值；失败可能表现为异常或底层状态变化。 \\
\bottomrule
\end{longtable}

\textbf{对应 C++}

\begin{itemize}
\tightlist
\item
  类：\passthrough{\lstinline!unitree::robot::b2::LowPowerSwitchParameter!}
\item
  签名：\passthrough{\lstinline!toJson(common::JsonMap \&) const!}
\item
  绑定策略：\passthrough{\lstinline!SIGNATURE\_PREVIEW!}
\item
  声明位置：\passthrough{\lstinline!include/unitree/robot/b2/robot\_state/robot\_state\_api.hpp:168!}
\end{itemize}

\textbf{说明}

\begin{itemize}
\tightlist
\item
  示例是局部调用片段；\passthrough{\lstinline!obj!}
  和参数变量表示已经按上层业务校验并准备好的对象。
\item
  该条目可以被编辑器和类型检查器识别，但当前运行时可能在属性查找阶段直接失败。
\end{itemize}

\textbf{用法}

\begin{lstlisting}[language=Python]
from typing import TYPE_CHECKING

if TYPE_CHECKING:
    def planned_to_json(obj: LowPowerSwitchParameter, json: dict[str, Any]) -> None:
        obj.to_json(json=json)
\end{lstlisting}

\begin{quote}
{[}!CAUTION{]} 上面的 \passthrough{\lstinline!TYPE\_CHECKING!}
示例只表达计划调用形态。该分支在普通 Python
运行时为假，不代表当前扩展能执行此接口。
\end{quote}

\hypertarget{unitree-sdk2-cpp-robot-b2-motionswitcherclient}{%
\subsubsection{\texorpdfstring{\texttt{unitree\_sdk2\_cpp.robot.b2.MotionSwitcherClient}}{unitree\_sdk2\_cpp.robot.b2.MotionSwitcherClient}}\label{unitree-sdk2-cpp-robot-b2-motionswitcherclient}}

Robot 服务客户端。构造、初始化和具体方法的可用性必须分别检查。

\textbf{导入}

\begin{lstlisting}[language=Python]
from unitree_sdk2_cpp.robot.b2 import MotionSwitcherClient
\end{lstlisting}

\textbf{基类}：\passthrough{\lstinline!Client!}

\textbf{方法索引}

\begin{longtable}[]{@{}
  >{\raggedright\arraybackslash}p{(\columnwidth - 6\tabcolsep) * \real{0.18}}
  >{\raggedright\arraybackslash}p{(\columnwidth - 6\tabcolsep) * \real{0.24}}
  >{\raggedright\arraybackslash}p{(\columnwidth - 6\tabcolsep) * \real{0.18}}
  >{\raggedright\arraybackslash}p{(\columnwidth - 6\tabcolsep) * \real{0.40}}@{}}
\toprule
方法 & Python 签名 & 状态 & 安全分类 \\
\midrule
\endhead
\protect\hyperlink{unitree-sdk2-cpp-robot-b2-motionswitcherclient-dunder-init-1}{\passthrough{\lstinline!\_\_init\_\_!}}
& \passthrough{\lstinline!def \_\_init\_\_(self) -> None!} &
\passthrough{\lstinline!AVAILABLE!} &
\passthrough{\lstinline!CONSTRUCTION!} \\
\protect\hyperlink{unitree-sdk2-cpp-robot-b2-motionswitcherclient-init-1}{\passthrough{\lstinline!init!}}
& \passthrough{\lstinline!def init(self) -> None!} &
\passthrough{\lstinline!AVAILABLE!} &
\passthrough{\lstinline!INITIALIZATION!} \\
\protect\hyperlink{unitree-sdk2-cpp-robot-b2-motionswitcherclient-check-mode-1}{\passthrough{\lstinline!check\_mode!}}
&
\passthrough{\lstinline!def check\_mode(self) -> tuple[int, str, str]!}
& \passthrough{\lstinline!AVAILABLE!} &
\passthrough{\lstinline!READ\_ONLY!} \\
\protect\hyperlink{unitree-sdk2-cpp-robot-b2-motionswitcherclient-select-mode-1}{\passthrough{\lstinline!select\_mode!}}
&
\passthrough{\lstinline!def select\_mode(self, name\_or\_alias: str) -> int!}
& \passthrough{\lstinline!SIGNATURE\_ONLY!} &
\passthrough{\lstinline!MOTION\_COMMAND!} \\
\protect\hyperlink{unitree-sdk2-cpp-robot-b2-motionswitcherclient-release-mode-1}{\passthrough{\lstinline!release\_mode!}}
& \passthrough{\lstinline!def release\_mode(self) -> int!} &
\passthrough{\lstinline!SIGNATURE\_ONLY!} &
\passthrough{\lstinline!MOTION\_COMMAND!} \\
\protect\hyperlink{unitree-sdk2-cpp-robot-b2-motionswitcherclient-set-silent-1}{\passthrough{\lstinline!set\_silent!}}
& \passthrough{\lstinline!def set\_silent(self, silent: bool) -> int!} &
\passthrough{\lstinline!SIGNATURE\_ONLY!} &
\passthrough{\lstinline!MOTION\_COMMAND!} \\
\protect\hyperlink{unitree-sdk2-cpp-robot-b2-motionswitcherclient-get-silent-1}{\passthrough{\lstinline!get\_silent!}}
& \passthrough{\lstinline!def get\_silent(self) -> tuple[int, bool]!} &
\passthrough{\lstinline!AVAILABLE!} &
\passthrough{\lstinline!READ\_ONLY!} \\
\bottomrule
\end{longtable}

\hypertarget{unitree-sdk2-cpp-robot-b2-motionswitcherclient-dunder-init-1}{%
\paragraph{\texorpdfstring{\texttt{unitree\_sdk2\_cpp.robot.b2.MotionSwitcherClient.\_\_init\_\_}}{unitree\_sdk2\_cpp.robot.b2.MotionSwitcherClient.\_\_init\_\_}}\label{unitree-sdk2-cpp-robot-b2-motionswitcherclient-dunder-init-1}}

初始化 \passthrough{\lstinline!MotionSwitcherClient!}
实例。是否能实际构造取决于下方可用性状态。

\textbf{签名}

\begin{lstlisting}[language=Python]
def __init__(self) -> None
\end{lstlisting}

\textbf{可用性与安全性}

\textbf{\passthrough{\lstinline!AVAILABLE!} /
\passthrough{\lstinline!CONSTRUCTION!}}：当前绑定源码已实现；实际执行条件仍取决于平台和对应
SDK 环境。

\textbf{参数}

无显式参数。实例方法中的 \passthrough{\lstinline!self!} 由 Python
自动传入。

\textbf{返回值}

\begin{longtable}[]{@{}
  >{\raggedright\arraybackslash}p{(\columnwidth - 4\tabcolsep) * \real{0.18}}
  >{\raggedright\arraybackslash}p{(\columnwidth - 4\tabcolsep) * \real{0.20}}
  >{\raggedright\arraybackslash}p{(\columnwidth - 4\tabcolsep) * \real{0.62}}@{}}
\toprule
位置 & 类型 & 含义 \\
\midrule
\endhead
无 & \passthrough{\lstinline!None!} &
\passthrough{\lstinline!\_\_init\_\_!}
本身不返回值；调用类对象时，在构造函数可用的前提下得到该类实例。 \\
\bottomrule
\end{longtable}

\textbf{对应 C++}

\begin{itemize}
\tightlist
\item
  类：\passthrough{\lstinline!unitree::robot::b2::MotionSwitcherClient!}
\item
  签名：\passthrough{\lstinline!MotionSwitcherClient()!}
\item
  绑定策略：\passthrough{\lstinline!未单独标注!}
\item
  声明位置：\passthrough{\lstinline!include/unitree/robot/b2/motion\_switcher/motion\_switcher\_client.hpp:18!}
\end{itemize}

\textbf{说明}

\begin{itemize}
\tightlist
\item
  示例是局部调用片段；\passthrough{\lstinline!obj!}
  和参数变量表示已经按上层业务校验并准备好的对象。
\end{itemize}

\textbf{用法}

\begin{lstlisting}[language=Python]
value = MotionSwitcherClient()
\end{lstlisting}

\hypertarget{unitree-sdk2-cpp-robot-b2-motionswitcherclient-init-1}{%
\paragraph{\texorpdfstring{\texttt{unitree\_sdk2\_cpp.robot.b2.MotionSwitcherClient.init}}{unitree\_sdk2\_cpp.robot.b2.MotionSwitcherClient.init}}\label{unitree-sdk2-cpp-robot-b2-motionswitcherclient-init-1}}

初始化当前 SDK 对象所需的底层通道或服务资源。

\textbf{签名}

\begin{lstlisting}[language=Python]
def init(self) -> None
\end{lstlisting}

\textbf{可用性与安全性}

\textbf{\passthrough{\lstinline!AVAILABLE!} /
\passthrough{\lstinline!INITIALIZATION!}}：当前绑定源码已实现；实际执行条件仍取决于平台和对应
SDK 环境。

\textbf{参数}

无显式参数。实例方法中的 \passthrough{\lstinline!self!} 由 Python
自动传入。

\textbf{返回值}

\begin{longtable}[]{@{}
  >{\raggedright\arraybackslash}p{(\columnwidth - 4\tabcolsep) * \real{0.18}}
  >{\raggedright\arraybackslash}p{(\columnwidth - 4\tabcolsep) * \real{0.20}}
  >{\raggedright\arraybackslash}p{(\columnwidth - 4\tabcolsep) * \real{0.62}}@{}}
\toprule
位置 & 类型 & 含义 \\
\midrule
\endhead
无 & \passthrough{\lstinline!None!} & 该方法不返回 Python
值；失败可能表现为异常或底层状态变化。 \\
\bottomrule
\end{longtable}

\textbf{对应 C++}

\begin{itemize}
\tightlist
\item
  类：\passthrough{\lstinline!unitree::robot::b2::MotionSwitcherClient!}
\item
  签名：\passthrough{\lstinline!Init()!}
\item
  绑定策略：\passthrough{\lstinline!DIRECT!}
\item
  声明位置：\passthrough{\lstinline!include/unitree/robot/b2/motion\_switcher/motion\_switcher\_client.hpp:21!}
\end{itemize}

\textbf{说明}

\begin{itemize}
\tightlist
\item
  示例是局部调用片段；\passthrough{\lstinline!obj!}
  和参数变量表示已经按上层业务校验并准备好的对象。
\end{itemize}

\textbf{用法}

\begin{lstlisting}[language=Python]
obj.init()
\end{lstlisting}

\hypertarget{unitree-sdk2-cpp-robot-b2-motionswitcherclient-check-mode-1}{%
\paragraph{\texorpdfstring{\texttt{unitree\_sdk2\_cpp.robot.b2.MotionSwitcherClient.check\_mode}}{unitree\_sdk2\_cpp.robot.b2.MotionSwitcherClient.check\_mode}}\label{unitree-sdk2-cpp-robot-b2-motionswitcherclient-check-mode-1}}

查询或检查 模式。

\textbf{签名}

\begin{lstlisting}[language=Python]
def check_mode(self) -> tuple[int, str, str]
\end{lstlisting}

\textbf{可用性与安全性}

\textbf{\passthrough{\lstinline!AVAILABLE!} /
\passthrough{\lstinline!READ\_ONLY!}}：当前绑定源码已实现只读查询。Robot
Client 的服务端查询通常仍需匹配的 Linux
扩展、DDS、网络、机器人和服务；纯本地 getter 则不一定需要全部条件。

\textbf{参数}

无显式参数。实例方法中的 \passthrough{\lstinline!self!} 由 Python
自动传入。

\textbf{返回值}

\begin{longtable}[]{@{}
  >{\raggedright\arraybackslash}p{(\columnwidth - 4\tabcolsep) * \real{0.18}}
  >{\raggedright\arraybackslash}p{(\columnwidth - 4\tabcolsep) * \real{0.20}}
  >{\raggedright\arraybackslash}p{(\columnwidth - 4\tabcolsep) * \real{0.62}}@{}}
\toprule
位置 & 类型 & 含义 \\
\midrule
\endhead
{[}0{]} \passthrough{\lstinline!status!} & \passthrough{\lstinline!int!}
& SDK 状态码；Unitree 示例通常以 \passthrough{\lstinline!0!}
表示成功，非零值需按具体服务错误码解释。 \\
{[}1{]} \passthrough{\lstinline!form!} & \passthrough{\lstinline!str!} &
传给该接口的 \passthrough{\lstinline!form!} 参数，Python 类型为
\passthrough{\lstinline!str!}。精确含义、单位、范围和枚举值需查目标型号对应头文件与协议。
对应 C++ 参数 \passthrough{\lstinline!form: std::string \&!}。
底层为字符串；长度、编码和允许值由具体协议决定。 \\
{[}2{]} \passthrough{\lstinline!name!} & \passthrough{\lstinline!str!} &
SDK 对象、配置项、服务或动作的名称。允许值由调用它的具体接口定义。 对应
C++ 参数 \passthrough{\lstinline!name: std::string \&!}。
底层为字符串；长度、编码和允许值由具体协议决定。 \\
\bottomrule
\end{longtable}

\textbf{对应 C++}

\begin{itemize}
\tightlist
\item
  类：\passthrough{\lstinline!unitree::robot::b2::MotionSwitcherClient!}
\item
  签名：\passthrough{\lstinline!CheckMode(std::string \&, std::string \&)!}
\item
  绑定策略：\passthrough{\lstinline!OUTPUT\_WRAPPER!}
\item
  声明位置：\passthrough{\lstinline!include/unitree/robot/b2/motion\_switcher/motion\_switcher\_client.hpp:23!}
\end{itemize}

\textbf{说明}

\begin{itemize}
\tightlist
\item
  示例是局部调用片段；\passthrough{\lstinline!obj!}
  和参数变量表示已经按上层业务校验并准备好的对象。
\item
  C++ 的可修改输出引用不会作为 Python 入参出现，而是按 C++
  参数顺序追加到返回元组中。
\end{itemize}

\textbf{用法}

\begin{lstlisting}[language=Python]
status, form, name = obj.check_mode()
\end{lstlisting}

\begin{quote}
{[}!NOTE{]} 只读查询仍会访问
DDS/机器人服务。必须检查返回状态码，并处理超时和网络中断。
\end{quote}

\hypertarget{unitree-sdk2-cpp-robot-b2-motionswitcherclient-select-mode-1}{%
\paragraph{\texorpdfstring{\texttt{unitree\_sdk2\_cpp.robot.b2.MotionSwitcherClient.select\_mode}}{unitree\_sdk2\_cpp.robot.b2.MotionSwitcherClient.select\_mode}}\label{unitree-sdk2-cpp-robot-b2-motionswitcherclient-select-mode-1}}

对应 C++ SDK 操作
\passthrough{\lstinline!SelectMode(const std::string \&)!}。上游头文件没有可直接生成的业务说明时，本参考只保证签名映射，精确语义需查目标型号协议。

\textbf{签名}

\begin{lstlisting}[language=Python]
def select_mode(self, name_or_alias: str) -> int
\end{lstlisting}

\textbf{可用性与安全性}

\textbf{\passthrough{\lstinline!SIGNATURE\_ONLY!} /
\passthrough{\lstinline!MOTION\_COMMAND!}}：仅用于补全和静态检查。当前二进制没有该
Python 方法；不得按可执行运动接口使用。

\textbf{参数}

\begin{longtable}[]{@{}
  >{\raggedright\arraybackslash}p{(\columnwidth - 6\tabcolsep) * \real{0.18}}
  >{\raggedright\arraybackslash}p{(\columnwidth - 6\tabcolsep) * \real{0.24}}
  >{\raggedright\arraybackslash}p{(\columnwidth - 6\tabcolsep) * \real{0.18}}
  >{\raggedright\arraybackslash}p{(\columnwidth - 6\tabcolsep) * \real{0.40}}@{}}
\toprule
名称 & Python 类型 & 默认值 & 含义 \\
\midrule
\endhead
\passthrough{\lstinline!name\_or\_alias!} &
\passthrough{\lstinline!str!} & 必填 & 传给该接口的 名称
\passthrough{\lstinline!or!} \passthrough{\lstinline!alias!}
参数，Python 类型为
\passthrough{\lstinline!str!}。精确含义、单位、范围和枚举值需查目标型号对应头文件与协议。
对应 C++ 参数
\passthrough{\lstinline!nameOrAlias: const std::string \&!}。
底层为字符串；长度、编码和允许值由具体协议决定。 \\
\bottomrule
\end{longtable}

\textbf{返回值}

\begin{longtable}[]{@{}
  >{\raggedright\arraybackslash}p{(\columnwidth - 4\tabcolsep) * \real{0.18}}
  >{\raggedright\arraybackslash}p{(\columnwidth - 4\tabcolsep) * \real{0.20}}
  >{\raggedright\arraybackslash}p{(\columnwidth - 4\tabcolsep) * \real{0.62}}@{}}
\toprule
位置 & 类型 & 含义 \\
\midrule
\endhead
返回值 & \passthrough{\lstinline!int!} & SDK 状态码；Unitree 示例通常以
\passthrough{\lstinline!0!} 表示成功，非零值需按具体服务定义解释。 \\
\bottomrule
\end{longtable}

\textbf{对应 C++}

\begin{itemize}
\tightlist
\item
  类：\passthrough{\lstinline!unitree::robot::b2::MotionSwitcherClient!}
\item
  签名：\passthrough{\lstinline!SelectMode(const std::string \&)!}
\item
  绑定策略：\passthrough{\lstinline!DIRECT!}
\item
  声明位置：\passthrough{\lstinline!include/unitree/robot/b2/motion\_switcher/motion\_switcher\_client.hpp:24!}
\end{itemize}

\textbf{说明}

\begin{itemize}
\tightlist
\item
  示例是局部调用片段；\passthrough{\lstinline!obj!}
  和参数变量表示已经按上层业务校验并准备好的对象。
\item
  该条目可以被编辑器和类型检查器识别，但当前运行时可能在属性查找阶段直接失败。
\item
  该方法属于运动命令。方法名包含
  \passthrough{\lstinline!stop!}、\passthrough{\lstinline!damp!} 或
  \passthrough{\lstinline!zero!} 也不代表它是独立物理急停。
\end{itemize}

\textbf{用法}

\begin{lstlisting}[language=Python]
from typing import TYPE_CHECKING

if TYPE_CHECKING:
    def planned_select_mode(obj: MotionSwitcherClient, name_or_alias: str) -> int:
        return obj.select_mode(name_or_alias=name_or_alias)
\end{lstlisting}

\begin{quote}
{[}!CAUTION{]} 上面的 \passthrough{\lstinline!TYPE\_CHECKING!}
示例只表达计划调用形态。该分支在普通 Python
运行时为假，不代表当前扩展能执行此接口。
\end{quote}

\hypertarget{unitree-sdk2-cpp-robot-b2-motionswitcherclient-release-mode-1}{%
\paragraph{\texorpdfstring{\texttt{unitree\_sdk2\_cpp.robot.b2.MotionSwitcherClient.release\_mode}}{unitree\_sdk2\_cpp.robot.b2.MotionSwitcherClient.release\_mode}}\label{unitree-sdk2-cpp-robot-b2-motionswitcherclient-release-mode-1}}

对应 C++ SDK 操作
\passthrough{\lstinline!ReleaseMode()!}。上游头文件没有可直接生成的业务说明时，本参考只保证签名映射，精确语义需查目标型号协议。

\textbf{签名}

\begin{lstlisting}[language=Python]
def release_mode(self) -> int
\end{lstlisting}

\textbf{可用性与安全性}

\textbf{\passthrough{\lstinline!SIGNATURE\_ONLY!} /
\passthrough{\lstinline!MOTION\_COMMAND!}}：仅用于补全和静态检查。当前二进制没有该
Python 方法；不得按可执行运动接口使用。

\textbf{参数}

无显式参数。实例方法中的 \passthrough{\lstinline!self!} 由 Python
自动传入。

\textbf{返回值}

\begin{longtable}[]{@{}
  >{\raggedright\arraybackslash}p{(\columnwidth - 4\tabcolsep) * \real{0.18}}
  >{\raggedright\arraybackslash}p{(\columnwidth - 4\tabcolsep) * \real{0.20}}
  >{\raggedright\arraybackslash}p{(\columnwidth - 4\tabcolsep) * \real{0.62}}@{}}
\toprule
位置 & 类型 & 含义 \\
\midrule
\endhead
返回值 & \passthrough{\lstinline!int!} & SDK 状态码；Unitree 示例通常以
\passthrough{\lstinline!0!} 表示成功，非零值需按具体服务定义解释。 \\
\bottomrule
\end{longtable}

\textbf{对应 C++}

\begin{itemize}
\tightlist
\item
  类：\passthrough{\lstinline!unitree::robot::b2::MotionSwitcherClient!}
\item
  签名：\passthrough{\lstinline!ReleaseMode()!}
\item
  绑定策略：\passthrough{\lstinline!DIRECT!}
\item
  声明位置：\passthrough{\lstinline!include/unitree/robot/b2/motion\_switcher/motion\_switcher\_client.hpp:25!}
\end{itemize}

\textbf{说明}

\begin{itemize}
\tightlist
\item
  示例是局部调用片段；\passthrough{\lstinline!obj!}
  和参数变量表示已经按上层业务校验并准备好的对象。
\item
  该条目可以被编辑器和类型检查器识别，但当前运行时可能在属性查找阶段直接失败。
\item
  该方法属于运动命令。方法名包含
  \passthrough{\lstinline!stop!}、\passthrough{\lstinline!damp!} 或
  \passthrough{\lstinline!zero!} 也不代表它是独立物理急停。
\end{itemize}

\textbf{用法}

\begin{lstlisting}[language=Python]
from typing import TYPE_CHECKING

if TYPE_CHECKING:
    def planned_release_mode(obj: MotionSwitcherClient) -> int:
        return obj.release_mode()
\end{lstlisting}

\begin{quote}
{[}!CAUTION{]} 上面的 \passthrough{\lstinline!TYPE\_CHECKING!}
示例只表达计划调用形态。该分支在普通 Python
运行时为假，不代表当前扩展能执行此接口。
\end{quote}

\hypertarget{unitree-sdk2-cpp-robot-b2-motionswitcherclient-set-silent-1}{%
\paragraph{\texorpdfstring{\texttt{unitree\_sdk2\_cpp.robot.b2.MotionSwitcherClient.set\_silent}}{unitree\_sdk2\_cpp.robot.b2.MotionSwitcherClient.set\_silent}}\label{unitree-sdk2-cpp-robot-b2-motionswitcherclient-set-silent-1}}

设置 静音状态。具体副作用和安全边界见下方状态。

\textbf{签名}

\begin{lstlisting}[language=Python]
def set_silent(self, silent: bool) -> int
\end{lstlisting}

\textbf{可用性与安全性}

\textbf{\passthrough{\lstinline!SIGNATURE\_ONLY!} /
\passthrough{\lstinline!MOTION\_COMMAND!}}：仅用于补全和静态检查。当前二进制没有该
Python 方法；不得按可执行运动接口使用。

\textbf{参数}

\begin{longtable}[]{@{}
  >{\raggedright\arraybackslash}p{(\columnwidth - 6\tabcolsep) * \real{0.18}}
  >{\raggedright\arraybackslash}p{(\columnwidth - 6\tabcolsep) * \real{0.24}}
  >{\raggedright\arraybackslash}p{(\columnwidth - 6\tabcolsep) * \real{0.18}}
  >{\raggedright\arraybackslash}p{(\columnwidth - 6\tabcolsep) * \real{0.40}}@{}}
\toprule
名称 & Python 类型 & 默认值 & 含义 \\
\midrule
\endhead
\passthrough{\lstinline!silent!} & \passthrough{\lstinline!bool!} & 必填
& 传给该接口的 静音状态 参数，Python 类型为
\passthrough{\lstinline!bool!}。精确含义、单位、范围和枚举值需查目标型号对应头文件与协议。
对应 C++ 参数 \passthrough{\lstinline!silent: bool!}。
只接受布尔语义。 \\
\bottomrule
\end{longtable}

\textbf{返回值}

\begin{longtable}[]{@{}
  >{\raggedright\arraybackslash}p{(\columnwidth - 4\tabcolsep) * \real{0.18}}
  >{\raggedright\arraybackslash}p{(\columnwidth - 4\tabcolsep) * \real{0.20}}
  >{\raggedright\arraybackslash}p{(\columnwidth - 4\tabcolsep) * \real{0.62}}@{}}
\toprule
位置 & 类型 & 含义 \\
\midrule
\endhead
返回值 & \passthrough{\lstinline!int!} & SDK 状态码；Unitree 示例通常以
\passthrough{\lstinline!0!} 表示成功，非零值需按具体服务定义解释。 \\
\bottomrule
\end{longtable}

\textbf{对应 C++}

\begin{itemize}
\tightlist
\item
  类：\passthrough{\lstinline!unitree::robot::b2::MotionSwitcherClient!}
\item
  签名：\passthrough{\lstinline!SetSilent(bool)!}
\item
  绑定策略：\passthrough{\lstinline!DIRECT!}
\item
  声明位置：\passthrough{\lstinline!include/unitree/robot/b2/motion\_switcher/motion\_switcher\_client.hpp:26!}
\end{itemize}

\textbf{说明}

\begin{itemize}
\tightlist
\item
  示例是局部调用片段；\passthrough{\lstinline!obj!}
  和参数变量表示已经按上层业务校验并准备好的对象。
\item
  该条目可以被编辑器和类型检查器识别，但当前运行时可能在属性查找阶段直接失败。
\item
  该方法属于运动命令。方法名包含
  \passthrough{\lstinline!stop!}、\passthrough{\lstinline!damp!} 或
  \passthrough{\lstinline!zero!} 也不代表它是独立物理急停。
\end{itemize}

\textbf{用法}

\begin{lstlisting}[language=Python]
from typing import TYPE_CHECKING

if TYPE_CHECKING:
    def planned_set_silent(obj: MotionSwitcherClient, silent: bool) -> int:
        return obj.set_silent(silent=silent)
\end{lstlisting}

\begin{quote}
{[}!CAUTION{]} 上面的 \passthrough{\lstinline!TYPE\_CHECKING!}
示例只表达计划调用形态。该分支在普通 Python
运行时为假，不代表当前扩展能执行此接口。
\end{quote}

\hypertarget{unitree-sdk2-cpp-robot-b2-motionswitcherclient-get-silent-1}{%
\paragraph{\texorpdfstring{\texttt{unitree\_sdk2\_cpp.robot.b2.MotionSwitcherClient.get\_silent}}{unitree\_sdk2\_cpp.robot.b2.MotionSwitcherClient.get\_silent}}\label{unitree-sdk2-cpp-robot-b2-motionswitcherclient-get-silent-1}}

查询或检查 静音状态。

\textbf{签名}

\begin{lstlisting}[language=Python]
def get_silent(self) -> tuple[int, bool]
\end{lstlisting}

\textbf{可用性与安全性}

\textbf{\passthrough{\lstinline!AVAILABLE!} /
\passthrough{\lstinline!READ\_ONLY!}}：当前绑定源码已实现只读查询。Robot
Client 的服务端查询通常仍需匹配的 Linux
扩展、DDS、网络、机器人和服务；纯本地 getter 则不一定需要全部条件。

\textbf{参数}

无显式参数。实例方法中的 \passthrough{\lstinline!self!} 由 Python
自动传入。

\textbf{返回值}

\begin{longtable}[]{@{}
  >{\raggedright\arraybackslash}p{(\columnwidth - 4\tabcolsep) * \real{0.18}}
  >{\raggedright\arraybackslash}p{(\columnwidth - 4\tabcolsep) * \real{0.20}}
  >{\raggedright\arraybackslash}p{(\columnwidth - 4\tabcolsep) * \real{0.62}}@{}}
\toprule
位置 & 类型 & 含义 \\
\midrule
\endhead
{[}0{]} \passthrough{\lstinline!status!} & \passthrough{\lstinline!int!}
& SDK 状态码；Unitree 示例通常以 \passthrough{\lstinline!0!}
表示成功，非零值需按具体服务错误码解释。 \\
{[}1{]} \passthrough{\lstinline!silent!} &
\passthrough{\lstinline!bool!} & 传给该接口的 静音状态 参数，Python
类型为
\passthrough{\lstinline!bool!}。精确含义、单位、范围和枚举值需查目标型号对应头文件与协议。
对应 C++ 参数 \passthrough{\lstinline!silent: bool \&!}。
只接受布尔语义。 \\
\bottomrule
\end{longtable}

\textbf{对应 C++}

\begin{itemize}
\tightlist
\item
  类：\passthrough{\lstinline!unitree::robot::b2::MotionSwitcherClient!}
\item
  签名：\passthrough{\lstinline!GetSilent(bool \&)!}
\item
  绑定策略：\passthrough{\lstinline!OUTPUT\_WRAPPER!}
\item
  声明位置：\passthrough{\lstinline!include/unitree/robot/b2/motion\_switcher/motion\_switcher\_client.hpp:27!}
\end{itemize}

\textbf{说明}

\begin{itemize}
\tightlist
\item
  示例是局部调用片段；\passthrough{\lstinline!obj!}
  和参数变量表示已经按上层业务校验并准备好的对象。
\item
  C++ 的可修改输出引用不会作为 Python 入参出现，而是按 C++
  参数顺序追加到返回元组中。
\end{itemize}

\textbf{用法}

\begin{lstlisting}[language=Python]
status, silent = obj.get_silent()
\end{lstlisting}

\begin{quote}
{[}!NOTE{]} 只读查询仍会访问
DDS/机器人服务。必须检查返回状态码，并处理超时和网络中断。
\end{quote}

\hypertarget{unitree-sdk2-cpp-robot-b2-pkgversiondata}{%
\subsubsection{\texorpdfstring{\texttt{unitree\_sdk2\_cpp.robot.b2.PkgVersionData}}{unitree\_sdk2\_cpp.robot.b2.PkgVersionData}}\label{unitree-sdk2-cpp-robot-b2-pkgversiondata}}

SDK 数据值类型；公开字段和转换方法见下文。

\textbf{导入}

\begin{lstlisting}[language=Python]
from unitree_sdk2_cpp.robot.b2 import PkgVersionData
\end{lstlisting}

\textbf{基类}：\passthrough{\lstinline!object!}

\textbf{公开属性}

\begin{longtable}[]{@{}
  >{\raggedright\arraybackslash}p{(\columnwidth - 6\tabcolsep) * \real{0.18}}
  >{\raggedright\arraybackslash}p{(\columnwidth - 6\tabcolsep) * \real{0.24}}
  >{\raggedright\arraybackslash}p{(\columnwidth - 6\tabcolsep) * \real{0.18}}
  >{\raggedright\arraybackslash}p{(\columnwidth - 6\tabcolsep) * \real{0.40}}@{}}
\toprule
名称 & Python 类型 & 含义 & 用法 \\
\midrule
\endhead
\passthrough{\lstinline!package\_version!} &
\passthrough{\lstinline!str!} & 传给该接口的
\passthrough{\lstinline!package!} \passthrough{\lstinline!version!}
参数，Python 类型为
\passthrough{\lstinline!str!}。精确含义、单位、范围和枚举值需查目标型号对应头文件与协议。
& \passthrough{\lstinline!value = obj.package\_version!} /
\passthrough{\lstinline!obj.package\_version = value!} \\
\passthrough{\lstinline!module\_version\_map!} &
\passthrough{\lstinline!dict[str, str]!} & 传给该接口的
\passthrough{\lstinline!module!} \passthrough{\lstinline!version!}
\passthrough{\lstinline!map!} 参数，Python 类型为
\passthrough{\lstinline!dict[str, str]!}。精确含义、单位、范围和枚举值需查目标型号对应头文件与协议。
& \passthrough{\lstinline!value = obj.module\_version\_map!} /
\passthrough{\lstinline!obj.module\_version\_map = value!} \\
\bottomrule
\end{longtable}

\textbf{方法索引}

\begin{longtable}[]{@{}
  >{\raggedright\arraybackslash}p{(\columnwidth - 6\tabcolsep) * \real{0.18}}
  >{\raggedright\arraybackslash}p{(\columnwidth - 6\tabcolsep) * \real{0.24}}
  >{\raggedright\arraybackslash}p{(\columnwidth - 6\tabcolsep) * \real{0.18}}
  >{\raggedright\arraybackslash}p{(\columnwidth - 6\tabcolsep) * \real{0.40}}@{}}
\toprule
方法 & Python 签名 & 状态 & 安全分类 \\
\midrule
\endhead
\protect\hyperlink{unitree-sdk2-cpp-robot-b2-pkgversiondata-dunder-init-1}{\passthrough{\lstinline!\_\_init\_\_!}}
& \passthrough{\lstinline!def \_\_init\_\_(self) -> None!} &
\passthrough{\lstinline!SIGNATURE\_ONLY!} &
\passthrough{\lstinline!CONSTRUCTION!} \\
\protect\hyperlink{unitree-sdk2-cpp-robot-b2-pkgversiondata-from-json-1}{\passthrough{\lstinline!from\_json!}}
&
\passthrough{\lstinline!def from\_json(self, json: dict[str, Any]) -> None!}
& \passthrough{\lstinline!SIGNATURE\_ONLY!} &
\passthrough{\lstinline!UNCLASSIFIED!} \\
\protect\hyperlink{unitree-sdk2-cpp-robot-b2-pkgversiondata-to-json-1}{\passthrough{\lstinline!to\_json!}}
&
\passthrough{\lstinline!def to\_json(self, json: dict[str, Any]) -> None!}
& \passthrough{\lstinline!SIGNATURE\_ONLY!} &
\passthrough{\lstinline!UNCLASSIFIED!} \\
\bottomrule
\end{longtable}

\hypertarget{unitree-sdk2-cpp-robot-b2-pkgversiondata-dunder-init-1}{%
\paragraph{\texorpdfstring{\texttt{unitree\_sdk2\_cpp.robot.b2.PkgVersionData.\_\_init\_\_}}{unitree\_sdk2\_cpp.robot.b2.PkgVersionData.\_\_init\_\_}}\label{unitree-sdk2-cpp-robot-b2-pkgversiondata-dunder-init-1}}

初始化 \passthrough{\lstinline!PkgVersionData!}
实例。是否能实际构造取决于下方可用性状态。

\textbf{签名}

\begin{lstlisting}[language=Python]
def __init__(self) -> None
\end{lstlisting}

\textbf{可用性与安全性}

\textbf{\passthrough{\lstinline!SIGNATURE\_ONLY!} /
\passthrough{\lstinline!CONSTRUCTION!}}：当前只有设计期签名，不能假设已存在于运行时扩展。

\textbf{参数}

无显式参数。实例方法中的 \passthrough{\lstinline!self!} 由 Python
自动传入。

\textbf{返回值}

\begin{longtable}[]{@{}
  >{\raggedright\arraybackslash}p{(\columnwidth - 4\tabcolsep) * \real{0.18}}
  >{\raggedright\arraybackslash}p{(\columnwidth - 4\tabcolsep) * \real{0.20}}
  >{\raggedright\arraybackslash}p{(\columnwidth - 4\tabcolsep) * \real{0.62}}@{}}
\toprule
位置 & 类型 & 含义 \\
\midrule
\endhead
无 & \passthrough{\lstinline!None!} &
\passthrough{\lstinline!\_\_init\_\_!}
本身不返回值；调用类对象时，在构造函数可用的前提下得到该类实例。 \\
\bottomrule
\end{longtable}

\textbf{对应 C++}

\begin{itemize}
\tightlist
\item
  类：\passthrough{\lstinline!unitree::robot::b2::PkgVersionData!}
\item
  签名：\passthrough{\lstinline!PkgVersionData()!}
\item
  绑定策略：\passthrough{\lstinline!未单独标注!}
\item
  声明位置：\passthrough{\lstinline!include/unitree/robot/b2/robot\_state/robot\_state\_api.hpp:209!}
\end{itemize}

\textbf{说明}

\begin{itemize}
\tightlist
\item
  示例是局部调用片段；\passthrough{\lstinline!obj!}
  和参数变量表示已经按上层业务校验并准备好的对象。
\item
  该条目可以被编辑器和类型检查器识别，但当前运行时可能在属性查找阶段直接失败。
\end{itemize}

\textbf{用法}

\begin{lstlisting}[language=Python]
from typing import TYPE_CHECKING

if TYPE_CHECKING:
    def planned_construct() -> PkgVersionData:
        return PkgVersionData()
\end{lstlisting}

\begin{quote}
{[}!CAUTION{]} 上面的 \passthrough{\lstinline!TYPE\_CHECKING!}
示例只表达计划调用形态。该分支在普通 Python
运行时为假，不代表当前扩展能执行此接口。
\end{quote}

\hypertarget{unitree-sdk2-cpp-robot-b2-pkgversiondata-from-json-1}{%
\paragraph{\texorpdfstring{\texttt{unitree\_sdk2\_cpp.robot.b2.PkgVersionData.from\_json}}{unitree\_sdk2\_cpp.robot.b2.PkgVersionData.from\_json}}\label{unitree-sdk2-cpp-robot-b2-pkgversiondata-from-json-1}}

计划从 JSON 风格字典读取字段并更新当前 SDK 值对象。

\textbf{签名}

\begin{lstlisting}[language=Python]
def from_json(self, json: dict[str, Any]) -> None
\end{lstlisting}

\textbf{可用性与安全性}

\textbf{\passthrough{\lstinline!SIGNATURE\_ONLY!} /
\passthrough{\lstinline!UNCLASSIFIED!}}：当前只有设计期签名，不能假设已存在于运行时扩展。

\textbf{参数}

\begin{longtable}[]{@{}
  >{\raggedright\arraybackslash}p{(\columnwidth - 6\tabcolsep) * \real{0.18}}
  >{\raggedright\arraybackslash}p{(\columnwidth - 6\tabcolsep) * \real{0.24}}
  >{\raggedright\arraybackslash}p{(\columnwidth - 6\tabcolsep) * \real{0.18}}
  >{\raggedright\arraybackslash}p{(\columnwidth - 6\tabcolsep) * \real{0.40}}@{}}
\toprule
名称 & Python 类型 & 默认值 & 含义 \\
\midrule
\endhead
\passthrough{\lstinline!json!} &
\passthrough{\lstinline!dict[str, Any]!} & 必填 & JSON
风格字典，键和值必须满足对应 C++ \passthrough{\lstinline!JsonMap!}
数据结构的约束。 对应 C++ 参数
\passthrough{\lstinline!json: common::JsonMap \&!}。 \\
\bottomrule
\end{longtable}

\textbf{返回值}

\begin{longtable}[]{@{}
  >{\raggedright\arraybackslash}p{(\columnwidth - 4\tabcolsep) * \real{0.18}}
  >{\raggedright\arraybackslash}p{(\columnwidth - 4\tabcolsep) * \real{0.20}}
  >{\raggedright\arraybackslash}p{(\columnwidth - 4\tabcolsep) * \real{0.62}}@{}}
\toprule
位置 & 类型 & 含义 \\
\midrule
\endhead
无 & \passthrough{\lstinline!None!} & 该方法不返回 Python
值；失败可能表现为异常或底层状态变化。 \\
\bottomrule
\end{longtable}

\textbf{对应 C++}

\begin{itemize}
\tightlist
\item
  类：\passthrough{\lstinline!unitree::robot::b2::PkgVersionData!}
\item
  签名：\passthrough{\lstinline!fromJson(common::JsonMap \&)!}
\item
  绑定策略：\passthrough{\lstinline!SIGNATURE\_PREVIEW!}
\item
  声明位置：\passthrough{\lstinline!include/unitree/robot/b2/robot\_state/robot\_state\_api.hpp:215!}
\end{itemize}

\textbf{说明}

\begin{itemize}
\tightlist
\item
  示例是局部调用片段；\passthrough{\lstinline!obj!}
  和参数变量表示已经按上层业务校验并准备好的对象。
\item
  该条目可以被编辑器和类型检查器识别，但当前运行时可能在属性查找阶段直接失败。
\end{itemize}

\textbf{用法}

\begin{lstlisting}[language=Python]
from typing import TYPE_CHECKING

if TYPE_CHECKING:
    def planned_from_json(obj: PkgVersionData, json: dict[str, Any]) -> None:
        obj.from_json(json=json)
\end{lstlisting}

\begin{quote}
{[}!CAUTION{]} 上面的 \passthrough{\lstinline!TYPE\_CHECKING!}
示例只表达计划调用形态。该分支在普通 Python
运行时为假，不代表当前扩展能执行此接口。
\end{quote}

\hypertarget{unitree-sdk2-cpp-robot-b2-pkgversiondata-to-json-1}{%
\paragraph{\texorpdfstring{\texttt{unitree\_sdk2\_cpp.robot.b2.PkgVersionData.to\_json}}{unitree\_sdk2\_cpp.robot.b2.PkgVersionData.to\_json}}\label{unitree-sdk2-cpp-robot-b2-pkgversiondata-to-json-1}}

计划把当前 SDK 值对象写入 JSON 风格字典。

\textbf{签名}

\begin{lstlisting}[language=Python]
def to_json(self, json: dict[str, Any]) -> None
\end{lstlisting}

\textbf{可用性与安全性}

\textbf{\passthrough{\lstinline!SIGNATURE\_ONLY!} /
\passthrough{\lstinline!UNCLASSIFIED!}}：当前只有设计期签名，不能假设已存在于运行时扩展。

\textbf{参数}

\begin{longtable}[]{@{}
  >{\raggedright\arraybackslash}p{(\columnwidth - 6\tabcolsep) * \real{0.18}}
  >{\raggedright\arraybackslash}p{(\columnwidth - 6\tabcolsep) * \real{0.24}}
  >{\raggedright\arraybackslash}p{(\columnwidth - 6\tabcolsep) * \real{0.18}}
  >{\raggedright\arraybackslash}p{(\columnwidth - 6\tabcolsep) * \real{0.40}}@{}}
\toprule
名称 & Python 类型 & 默认值 & 含义 \\
\midrule
\endhead
\passthrough{\lstinline!json!} &
\passthrough{\lstinline!dict[str, Any]!} & 必填 & JSON
风格字典，键和值必须满足对应 C++ \passthrough{\lstinline!JsonMap!}
数据结构的约束。 对应 C++ 参数
\passthrough{\lstinline!json: common::JsonMap \&!}。 \\
\bottomrule
\end{longtable}

\textbf{返回值}

\begin{longtable}[]{@{}
  >{\raggedright\arraybackslash}p{(\columnwidth - 4\tabcolsep) * \real{0.18}}
  >{\raggedright\arraybackslash}p{(\columnwidth - 4\tabcolsep) * \real{0.20}}
  >{\raggedright\arraybackslash}p{(\columnwidth - 4\tabcolsep) * \real{0.62}}@{}}
\toprule
位置 & 类型 & 含义 \\
\midrule
\endhead
无 & \passthrough{\lstinline!None!} & 该方法不返回 Python
值；失败可能表现为异常或底层状态变化。 \\
\bottomrule
\end{longtable}

\textbf{对应 C++}

\begin{itemize}
\tightlist
\item
  类：\passthrough{\lstinline!unitree::robot::b2::PkgVersionData!}
\item
  签名：\passthrough{\lstinline!toJson(common::JsonMap \&) const!}
\item
  绑定策略：\passthrough{\lstinline!SIGNATURE\_PREVIEW!}
\item
  声明位置：\passthrough{\lstinline!include/unitree/robot/b2/robot\_state/robot\_state\_api.hpp:221!}
\end{itemize}

\textbf{说明}

\begin{itemize}
\tightlist
\item
  示例是局部调用片段；\passthrough{\lstinline!obj!}
  和参数变量表示已经按上层业务校验并准备好的对象。
\item
  该条目可以被编辑器和类型检查器识别，但当前运行时可能在属性查找阶段直接失败。
\end{itemize}

\textbf{用法}

\begin{lstlisting}[language=Python]
from typing import TYPE_CHECKING

if TYPE_CHECKING:
    def planned_to_json(obj: PkgVersionData, json: dict[str, Any]) -> None:
        obj.to_json(json=json)
\end{lstlisting}

\begin{quote}
{[}!CAUTION{]} 上面的 \passthrough{\lstinline!TYPE\_CHECKING!}
示例只表达计划调用形态。该分支在普通 Python
运行时为假，不代表当前扩展能执行此接口。
\end{quote}

\hypertarget{unitree-sdk2-cpp-robot-b2-robotstateclient}{%
\subsubsection{\texorpdfstring{\texttt{unitree\_sdk2\_cpp.robot.b2.RobotStateClient}}{unitree\_sdk2\_cpp.robot.b2.RobotStateClient}}\label{unitree-sdk2-cpp-robot-b2-robotstateclient}}

Robot 服务客户端。构造、初始化和具体方法的可用性必须分别检查。

\textbf{导入}

\begin{lstlisting}[language=Python]
from unitree_sdk2_cpp.robot.b2 import RobotStateClient
\end{lstlisting}

\textbf{基类}：\passthrough{\lstinline!Client!}

\textbf{方法索引}

\begin{longtable}[]{@{}
  >{\raggedright\arraybackslash}p{(\columnwidth - 6\tabcolsep) * \real{0.18}}
  >{\raggedright\arraybackslash}p{(\columnwidth - 6\tabcolsep) * \real{0.24}}
  >{\raggedright\arraybackslash}p{(\columnwidth - 6\tabcolsep) * \real{0.18}}
  >{\raggedright\arraybackslash}p{(\columnwidth - 6\tabcolsep) * \real{0.40}}@{}}
\toprule
方法 & Python 签名 & 状态 & 安全分类 \\
\midrule
\endhead
\protect\hyperlink{unitree-sdk2-cpp-robot-b2-robotstateclient-dunder-init-1}{\passthrough{\lstinline!\_\_init\_\_!}}
& \passthrough{\lstinline!def \_\_init\_\_(self) -> None!} &
\passthrough{\lstinline!AVAILABLE!} &
\passthrough{\lstinline!CONSTRUCTION!} \\
\protect\hyperlink{unitree-sdk2-cpp-robot-b2-robotstateclient-init-1}{\passthrough{\lstinline!init!}}
& \passthrough{\lstinline!def init(self) -> None!} &
\passthrough{\lstinline!AVAILABLE!} &
\passthrough{\lstinline!INITIALIZATION!} \\
\protect\hyperlink{unitree-sdk2-cpp-robot-b2-robotstateclient-service-list-1}{\passthrough{\lstinline!service\_list!}}
&
\passthrough{\lstinline!def service\_list(self) -> tuple[int, list[ServiceState]]!}
& \passthrough{\lstinline!AVAILABLE!} &
\passthrough{\lstinline!READ\_ONLY!} \\
\protect\hyperlink{unitree-sdk2-cpp-robot-b2-robotstateclient-service-switch-1}{\passthrough{\lstinline!service\_switch!}}
&
\passthrough{\lstinline!def service\_switch(self, name: str, swit: int) -> tuple[int, int]!}
& \passthrough{\lstinline!SIGNATURE\_ONLY!} &
\passthrough{\lstinline!HARDWARE\_SIDE\_EFFECT!} \\
\protect\hyperlink{unitree-sdk2-cpp-robot-b2-robotstateclient-set-report-freq-1}{\passthrough{\lstinline!set\_report\_freq!}}
&
\passthrough{\lstinline!def set\_report\_freq(self, interval: int, duration: int) -> int!}
& \passthrough{\lstinline!SIGNATURE\_ONLY!} &
\passthrough{\lstinline!HARDWARE\_SIDE\_EFFECT!} \\
\protect\hyperlink{unitree-sdk2-cpp-robot-b2-robotstateclient-low-power-switch-1}{\passthrough{\lstinline!low\_power\_switch!}}
&
\passthrough{\lstinline!def low\_power\_switch(self, swit: int) -> int!}
& \passthrough{\lstinline!SIGNATURE\_ONLY!} &
\passthrough{\lstinline!HARDWARE\_SIDE\_EFFECT!} \\
\protect\hyperlink{unitree-sdk2-cpp-robot-b2-robotstateclient-low-power-status-1}{\passthrough{\lstinline!low\_power\_status!}}
&
\passthrough{\lstinline!def low\_power\_status(self) -> tuple[int, int]!}
& \passthrough{\lstinline!SIGNATURE\_ONLY!} &
\passthrough{\lstinline!HARDWARE\_SIDE\_EFFECT!} \\
\protect\hyperlink{unitree-sdk2-cpp-robot-b2-robotstateclient-get-pkg-version-1}{\passthrough{\lstinline!get\_pkg\_version!}}
&
\passthrough{\lstinline!def get\_pkg\_version(self) -> tuple[int, str, dict[str, str]]!}
& \passthrough{\lstinline!AVAILABLE!} &
\passthrough{\lstinline!READ\_ONLY!} \\
\bottomrule
\end{longtable}

\hypertarget{unitree-sdk2-cpp-robot-b2-robotstateclient-dunder-init-1}{%
\paragraph{\texorpdfstring{\texttt{unitree\_sdk2\_cpp.robot.b2.RobotStateClient.\_\_init\_\_}}{unitree\_sdk2\_cpp.robot.b2.RobotStateClient.\_\_init\_\_}}\label{unitree-sdk2-cpp-robot-b2-robotstateclient-dunder-init-1}}

初始化 \passthrough{\lstinline!RobotStateClient!}
实例。是否能实际构造取决于下方可用性状态。

\textbf{签名}

\begin{lstlisting}[language=Python]
def __init__(self) -> None
\end{lstlisting}

\textbf{可用性与安全性}

\textbf{\passthrough{\lstinline!AVAILABLE!} /
\passthrough{\lstinline!CONSTRUCTION!}}：当前绑定源码已实现；实际执行条件仍取决于平台和对应
SDK 环境。

\textbf{参数}

无显式参数。实例方法中的 \passthrough{\lstinline!self!} 由 Python
自动传入。

\textbf{返回值}

\begin{longtable}[]{@{}
  >{\raggedright\arraybackslash}p{(\columnwidth - 4\tabcolsep) * \real{0.18}}
  >{\raggedright\arraybackslash}p{(\columnwidth - 4\tabcolsep) * \real{0.20}}
  >{\raggedright\arraybackslash}p{(\columnwidth - 4\tabcolsep) * \real{0.62}}@{}}
\toprule
位置 & 类型 & 含义 \\
\midrule
\endhead
无 & \passthrough{\lstinline!None!} &
\passthrough{\lstinline!\_\_init\_\_!}
本身不返回值；调用类对象时，在构造函数可用的前提下得到该类实例。 \\
\bottomrule
\end{longtable}

\textbf{对应 C++}

\begin{itemize}
\tightlist
\item
  类：\passthrough{\lstinline!unitree::robot::b2::RobotStateClient!}
\item
  签名：\passthrough{\lstinline!RobotStateClient()!}
\item
  绑定策略：\passthrough{\lstinline!未单独标注!}
\item
  声明位置：\passthrough{\lstinline!include/unitree/robot/b2/robot\_state/robot\_state\_client.hpp:32!}
\end{itemize}

\textbf{说明}

\begin{itemize}
\tightlist
\item
  示例是局部调用片段；\passthrough{\lstinline!obj!}
  和参数变量表示已经按上层业务校验并准备好的对象。
\end{itemize}

\textbf{用法}

\begin{lstlisting}[language=Python]
value = RobotStateClient()
\end{lstlisting}

\hypertarget{unitree-sdk2-cpp-robot-b2-robotstateclient-init-1}{%
\paragraph{\texorpdfstring{\texttt{unitree\_sdk2\_cpp.robot.b2.RobotStateClient.init}}{unitree\_sdk2\_cpp.robot.b2.RobotStateClient.init}}\label{unitree-sdk2-cpp-robot-b2-robotstateclient-init-1}}

初始化当前 SDK 对象所需的底层通道或服务资源。

\textbf{签名}

\begin{lstlisting}[language=Python]
def init(self) -> None
\end{lstlisting}

\textbf{可用性与安全性}

\textbf{\passthrough{\lstinline!AVAILABLE!} /
\passthrough{\lstinline!INITIALIZATION!}}：当前绑定源码已实现；实际执行条件仍取决于平台和对应
SDK 环境。

\textbf{参数}

无显式参数。实例方法中的 \passthrough{\lstinline!self!} 由 Python
自动传入。

\textbf{返回值}

\begin{longtable}[]{@{}
  >{\raggedright\arraybackslash}p{(\columnwidth - 4\tabcolsep) * \real{0.18}}
  >{\raggedright\arraybackslash}p{(\columnwidth - 4\tabcolsep) * \real{0.20}}
  >{\raggedright\arraybackslash}p{(\columnwidth - 4\tabcolsep) * \real{0.62}}@{}}
\toprule
位置 & 类型 & 含义 \\
\midrule
\endhead
无 & \passthrough{\lstinline!None!} & 该方法不返回 Python
值；失败可能表现为异常或底层状态变化。 \\
\bottomrule
\end{longtable}

\textbf{对应 C++}

\begin{itemize}
\tightlist
\item
  类：\passthrough{\lstinline!unitree::robot::b2::RobotStateClient!}
\item
  签名：\passthrough{\lstinline!Init()!}
\item
  绑定策略：\passthrough{\lstinline!DIRECT!}
\item
  声明位置：\passthrough{\lstinline!include/unitree/robot/b2/robot\_state/robot\_state\_client.hpp:35!}
\end{itemize}

\textbf{说明}

\begin{itemize}
\tightlist
\item
  示例是局部调用片段；\passthrough{\lstinline!obj!}
  和参数变量表示已经按上层业务校验并准备好的对象。
\end{itemize}

\textbf{用法}

\begin{lstlisting}[language=Python]
obj.init()
\end{lstlisting}

\hypertarget{unitree-sdk2-cpp-robot-b2-robotstateclient-service-list-1}{%
\paragraph{\texorpdfstring{\texttt{unitree\_sdk2\_cpp.robot.b2.RobotStateClient.service\_list}}{unitree\_sdk2\_cpp.robot.b2.RobotStateClient.service\_list}}\label{unitree-sdk2-cpp-robot-b2-robotstateclient-service-list-1}}

对应 C++ SDK 操作
\passthrough{\lstinline!ServiceList(std::vector<ServiceState> \&)!}。上游头文件没有可直接生成的业务说明时，本参考只保证签名映射，精确语义需查目标型号协议。

\textbf{签名}

\begin{lstlisting}[language=Python]
def service_list(self) -> tuple[int, list[ServiceState]]
\end{lstlisting}

\textbf{可用性与安全性}

\textbf{\passthrough{\lstinline!AVAILABLE!} /
\passthrough{\lstinline!READ\_ONLY!}}：当前绑定源码已实现只读查询。Robot
Client 的服务端查询通常仍需匹配的 Linux
扩展、DDS、网络、机器人和服务；纯本地 getter 则不一定需要全部条件。

\textbf{参数}

无显式参数。实例方法中的 \passthrough{\lstinline!self!} 由 Python
自动传入。

\textbf{返回值}

\begin{longtable}[]{@{}
  >{\raggedright\arraybackslash}p{(\columnwidth - 4\tabcolsep) * \real{0.18}}
  >{\raggedright\arraybackslash}p{(\columnwidth - 4\tabcolsep) * \real{0.20}}
  >{\raggedright\arraybackslash}p{(\columnwidth - 4\tabcolsep) * \real{0.62}}@{}}
\toprule
位置 & 类型 & 含义 \\
\midrule
\endhead
{[}0{]} \passthrough{\lstinline!status!} & \passthrough{\lstinline!int!}
& SDK 状态码；Unitree 示例通常以 \passthrough{\lstinline!0!}
表示成功，非零值需按具体服务错误码解释。 \\
{[}1{]} \passthrough{\lstinline!serviceStateList!} &
\passthrough{\lstinline!list[ServiceState]!} & 传给该接口的
\passthrough{\lstinline!serviceStateList!} 参数，Python 类型为
\passthrough{\lstinline!list[ServiceState]!}。精确含义、单位、范围和枚举值需查目标型号对应头文件与协议。
对应 C++ 参数
\passthrough{\lstinline!serviceStateList: std::vector<ServiceState> \&!}。
底层是可变长度 vector \\
\bottomrule
\end{longtable}

\textbf{对应 C++}

\begin{itemize}
\tightlist
\item
  类：\passthrough{\lstinline!unitree::robot::b2::RobotStateClient!}
\item
  签名：\passthrough{\lstinline!ServiceList(std::vector<ServiceState> \&)!}
\item
  绑定策略：\passthrough{\lstinline!OUTPUT\_WRAPPER!}
\item
  声明位置：\passthrough{\lstinline!include/unitree/robot/b2/robot\_state/robot\_state\_client.hpp:37!}
\end{itemize}

\textbf{说明}

\begin{itemize}
\tightlist
\item
  示例是局部调用片段；\passthrough{\lstinline!obj!}
  和参数变量表示已经按上层业务校验并准备好的对象。
\item
  C++ 的可修改输出引用不会作为 Python 入参出现，而是按 C++
  参数顺序追加到返回元组中。
\end{itemize}

\textbf{用法}

\begin{lstlisting}[language=Python]
status, serviceStateList = obj.service_list()
\end{lstlisting}

\begin{quote}
{[}!NOTE{]} 只读查询仍会访问
DDS/机器人服务。必须检查返回状态码，并处理超时和网络中断。
\end{quote}

\hypertarget{unitree-sdk2-cpp-robot-b2-robotstateclient-service-switch-1}{%
\paragraph{\texorpdfstring{\texttt{unitree\_sdk2\_cpp.robot.b2.RobotStateClient.service\_switch}}{unitree\_sdk2\_cpp.robot.b2.RobotStateClient.service\_switch}}\label{unitree-sdk2-cpp-robot-b2-robotstateclient-service-switch-1}}

对应 C++ SDK 操作
\passthrough{\lstinline!ServiceSwitch(const std::string \&, int32\_t, int32\_t \&)!}。上游头文件没有可直接生成的业务说明时，本参考只保证签名映射，精确语义需查目标型号协议。

\textbf{签名}

\begin{lstlisting}[language=Python]
def service_switch(self, name: str, swit: int) -> tuple[int, int]
\end{lstlisting}

\textbf{可用性与安全性}

\textbf{\passthrough{\lstinline!SIGNATURE\_ONLY!} /
\passthrough{\lstinline!HARDWARE\_SIDE\_EFFECT!}}：仅用于补全和静态检查。当前二进制没有该
Python 方法，未来实现还需单独评审硬件副作用。

\textbf{参数}

\begin{longtable}[]{@{}
  >{\raggedright\arraybackslash}p{(\columnwidth - 6\tabcolsep) * \real{0.18}}
  >{\raggedright\arraybackslash}p{(\columnwidth - 6\tabcolsep) * \real{0.24}}
  >{\raggedright\arraybackslash}p{(\columnwidth - 6\tabcolsep) * \real{0.18}}
  >{\raggedright\arraybackslash}p{(\columnwidth - 6\tabcolsep) * \real{0.40}}@{}}
\toprule
名称 & Python 类型 & 默认值 & 含义 \\
\midrule
\endhead
\passthrough{\lstinline!name!} & \passthrough{\lstinline!str!} & 必填 &
SDK 对象、配置项、服务或动作的名称。允许值由调用它的具体接口定义。 对应
C++ 参数 \passthrough{\lstinline!name: const std::string \&!}。
底层为字符串；长度、编码和允许值由具体协议决定。 \\
\passthrough{\lstinline!swit!} & \passthrough{\lstinline!int!} & 必填 &
传给该接口的 \passthrough{\lstinline!swit!} 参数，Python 类型为
\passthrough{\lstinline!int!}。精确含义、单位、范围和枚举值需查目标型号对应头文件与协议。
对应 C++ 参数 \passthrough{\lstinline!swit: int32\_t!}。 取值范围为
-2147483648 到 2147483647。 \\
\bottomrule
\end{longtable}

\textbf{返回值}

\begin{longtable}[]{@{}
  >{\raggedright\arraybackslash}p{(\columnwidth - 4\tabcolsep) * \real{0.18}}
  >{\raggedright\arraybackslash}p{(\columnwidth - 4\tabcolsep) * \real{0.20}}
  >{\raggedright\arraybackslash}p{(\columnwidth - 4\tabcolsep) * \real{0.62}}@{}}
\toprule
位置 & 类型 & 含义 \\
\midrule
\endhead
{[}0{]} \passthrough{\lstinline!status!} & \passthrough{\lstinline!int!}
& SDK 状态码；Unitree 示例通常以 \passthrough{\lstinline!0!}
表示成功，非零值需按具体服务错误码解释。 \\
{[}1{]} \passthrough{\lstinline!status!} & \passthrough{\lstinline!int!}
& SDK 状态码；Unitree 示例通常以 \passthrough{\lstinline!0!}
表示成功，非零值需按具体服务错误码解释。 \\
\bottomrule
\end{longtable}

\textbf{对应 C++}

\begin{itemize}
\tightlist
\item
  类：\passthrough{\lstinline!unitree::robot::b2::RobotStateClient!}
\item
  签名：\passthrough{\lstinline!ServiceSwitch(const std::string \&, int32\_t, int32\_t \&)!}
\item
  绑定策略：\passthrough{\lstinline!OUTPUT\_WRAPPER!}
\item
  声明位置：\passthrough{\lstinline!include/unitree/robot/b2/robot\_state/robot\_state\_client.hpp:38!}
\end{itemize}

\textbf{说明}

\begin{itemize}
\tightlist
\item
  示例是局部调用片段；\passthrough{\lstinline!obj!}
  和参数变量表示已经按上层业务校验并准备好的对象。
\item
  C++ 的可修改输出引用不会作为 Python 入参出现，而是按 C++
  参数顺序追加到返回元组中。
\item
  该条目可以被编辑器和类型检查器识别，但当前运行时可能在属性查找阶段直接失败。
\end{itemize}

\textbf{用法}

\begin{lstlisting}[language=Python]
from typing import TYPE_CHECKING

if TYPE_CHECKING:
    def planned_service_switch(obj: RobotStateClient, name: str, swit: int) -> tuple[int, int]:
        return obj.service_switch(name=name, swit=swit)
\end{lstlisting}

\begin{quote}
{[}!CAUTION{]} 上面的 \passthrough{\lstinline!TYPE\_CHECKING!}
示例只表达计划调用形态。该分支在普通 Python
运行时为假，不代表当前扩展能执行此接口。
\end{quote}

\hypertarget{unitree-sdk2-cpp-robot-b2-robotstateclient-set-report-freq-1}{%
\paragraph{\texorpdfstring{\texttt{unitree\_sdk2\_cpp.robot.b2.RobotStateClient.set\_report\_freq}}{unitree\_sdk2\_cpp.robot.b2.RobotStateClient.set\_report\_freq}}\label{unitree-sdk2-cpp-robot-b2-robotstateclient-set-report-freq-1}}

设置 \passthrough{\lstinline!report!}
\passthrough{\lstinline!freq!}。具体副作用和安全边界见下方状态。

\textbf{签名}

\begin{lstlisting}[language=Python]
def set_report_freq(self, interval: int, duration: int) -> int
\end{lstlisting}

\textbf{可用性与安全性}

\textbf{\passthrough{\lstinline!SIGNATURE\_ONLY!} /
\passthrough{\lstinline!HARDWARE\_SIDE\_EFFECT!}}：仅用于补全和静态检查。当前二进制没有该
Python 方法，未来实现还需单独评审硬件副作用。

\textbf{参数}

\begin{longtable}[]{@{}
  >{\raggedright\arraybackslash}p{(\columnwidth - 6\tabcolsep) * \real{0.18}}
  >{\raggedright\arraybackslash}p{(\columnwidth - 6\tabcolsep) * \real{0.24}}
  >{\raggedright\arraybackslash}p{(\columnwidth - 6\tabcolsep) * \real{0.18}}
  >{\raggedright\arraybackslash}p{(\columnwidth - 6\tabcolsep) * \real{0.40}}@{}}
\toprule
名称 & Python 类型 & 默认值 & 含义 \\
\midrule
\endhead
\passthrough{\lstinline!interval!} & \passthrough{\lstinline!int!} &
必填 & 传给该接口的 \passthrough{\lstinline!interval!} 参数，Python
类型为
\passthrough{\lstinline!int!}。精确含义、单位、范围和枚举值需查目标型号对应头文件与协议。
对应 C++ 参数 \passthrough{\lstinline!interval: int32\_t!}。 取值范围为
-2147483648 到 2147483647。 \\
\passthrough{\lstinline!duration!} & \passthrough{\lstinline!int!} &
必填 &
操作持续时间参数。精确单位、范围和默认行为以目标型号接口定义为准。 对应
C++ 参数 \passthrough{\lstinline!duration: int32\_t!}。 取值范围为
-2147483648 到 2147483647。 \\
\bottomrule
\end{longtable}

\textbf{返回值}

\begin{longtable}[]{@{}
  >{\raggedright\arraybackslash}p{(\columnwidth - 4\tabcolsep) * \real{0.18}}
  >{\raggedright\arraybackslash}p{(\columnwidth - 4\tabcolsep) * \real{0.20}}
  >{\raggedright\arraybackslash}p{(\columnwidth - 4\tabcolsep) * \real{0.62}}@{}}
\toprule
位置 & 类型 & 含义 \\
\midrule
\endhead
返回值 & \passthrough{\lstinline!int!} & SDK 状态码；Unitree 示例通常以
\passthrough{\lstinline!0!} 表示成功，非零值需按具体服务定义解释。 \\
\bottomrule
\end{longtable}

\textbf{对应 C++}

\begin{itemize}
\tightlist
\item
  类：\passthrough{\lstinline!unitree::robot::b2::RobotStateClient!}
\item
  签名：\passthrough{\lstinline!SetReportFreq(int32\_t, int32\_t)!}
\item
  绑定策略：\passthrough{\lstinline!DIRECT!}
\item
  声明位置：\passthrough{\lstinline!include/unitree/robot/b2/robot\_state/robot\_state\_client.hpp:39!}
\end{itemize}

\textbf{说明}

\begin{itemize}
\tightlist
\item
  示例是局部调用片段；\passthrough{\lstinline!obj!}
  和参数变量表示已经按上层业务校验并准备好的对象。
\item
  该条目可以被编辑器和类型检查器识别，但当前运行时可能在属性查找阶段直接失败。
\end{itemize}

\textbf{用法}

\begin{lstlisting}[language=Python]
from typing import TYPE_CHECKING

if TYPE_CHECKING:
    def planned_set_report_freq(obj: RobotStateClient, interval: int, duration: int) -> int:
        return obj.set_report_freq(interval=interval, duration=duration)
\end{lstlisting}

\begin{quote}
{[}!CAUTION{]} 上面的 \passthrough{\lstinline!TYPE\_CHECKING!}
示例只表达计划调用形态。该分支在普通 Python
运行时为假，不代表当前扩展能执行此接口。
\end{quote}

\hypertarget{unitree-sdk2-cpp-robot-b2-robotstateclient-low-power-switch-1}{%
\paragraph{\texorpdfstring{\texttt{unitree\_sdk2\_cpp.robot.b2.RobotStateClient.low\_power\_switch}}{unitree\_sdk2\_cpp.robot.b2.RobotStateClient.low\_power\_switch}}\label{unitree-sdk2-cpp-robot-b2-robotstateclient-low-power-switch-1}}

对应 C++ SDK 操作
\passthrough{\lstinline!LowPowerSwitch(int32\_t)!}。上游头文件没有可直接生成的业务说明时，本参考只保证签名映射，精确语义需查目标型号协议。

\textbf{签名}

\begin{lstlisting}[language=Python]
def low_power_switch(self, swit: int) -> int
\end{lstlisting}

\textbf{可用性与安全性}

\textbf{\passthrough{\lstinline!SIGNATURE\_ONLY!} /
\passthrough{\lstinline!HARDWARE\_SIDE\_EFFECT!}}：仅用于补全和静态检查。当前二进制没有该
Python 方法，未来实现还需单独评审硬件副作用。

\textbf{参数}

\begin{longtable}[]{@{}
  >{\raggedright\arraybackslash}p{(\columnwidth - 6\tabcolsep) * \real{0.18}}
  >{\raggedright\arraybackslash}p{(\columnwidth - 6\tabcolsep) * \real{0.24}}
  >{\raggedright\arraybackslash}p{(\columnwidth - 6\tabcolsep) * \real{0.18}}
  >{\raggedright\arraybackslash}p{(\columnwidth - 6\tabcolsep) * \real{0.40}}@{}}
\toprule
名称 & Python 类型 & 默认值 & 含义 \\
\midrule
\endhead
\passthrough{\lstinline!swit!} & \passthrough{\lstinline!int!} & 必填 &
传给该接口的 \passthrough{\lstinline!swit!} 参数，Python 类型为
\passthrough{\lstinline!int!}。精确含义、单位、范围和枚举值需查目标型号对应头文件与协议。
对应 C++ 参数 \passthrough{\lstinline!swit: int32\_t!}。 取值范围为
-2147483648 到 2147483647。 \\
\bottomrule
\end{longtable}

\textbf{返回值}

\begin{longtable}[]{@{}
  >{\raggedright\arraybackslash}p{(\columnwidth - 4\tabcolsep) * \real{0.18}}
  >{\raggedright\arraybackslash}p{(\columnwidth - 4\tabcolsep) * \real{0.20}}
  >{\raggedright\arraybackslash}p{(\columnwidth - 4\tabcolsep) * \real{0.62}}@{}}
\toprule
位置 & 类型 & 含义 \\
\midrule
\endhead
返回值 & \passthrough{\lstinline!int!} & SDK 状态码；Unitree 示例通常以
\passthrough{\lstinline!0!} 表示成功，非零值需按具体服务定义解释。 \\
\bottomrule
\end{longtable}

\textbf{对应 C++}

\begin{itemize}
\tightlist
\item
  类：\passthrough{\lstinline!unitree::robot::b2::RobotStateClient!}
\item
  签名：\passthrough{\lstinline!LowPowerSwitch(int32\_t)!}
\item
  绑定策略：\passthrough{\lstinline!DIRECT!}
\item
  声明位置：\passthrough{\lstinline!include/unitree/robot/b2/robot\_state/robot\_state\_client.hpp:40!}
\end{itemize}

\textbf{说明}

\begin{itemize}
\tightlist
\item
  示例是局部调用片段；\passthrough{\lstinline!obj!}
  和参数变量表示已经按上层业务校验并准备好的对象。
\item
  该条目可以被编辑器和类型检查器识别，但当前运行时可能在属性查找阶段直接失败。
\end{itemize}

\textbf{用法}

\begin{lstlisting}[language=Python]
from typing import TYPE_CHECKING

if TYPE_CHECKING:
    def planned_low_power_switch(obj: RobotStateClient, swit: int) -> int:
        return obj.low_power_switch(swit=swit)
\end{lstlisting}

\begin{quote}
{[}!CAUTION{]} 上面的 \passthrough{\lstinline!TYPE\_CHECKING!}
示例只表达计划调用形态。该分支在普通 Python
运行时为假，不代表当前扩展能执行此接口。
\end{quote}

\hypertarget{unitree-sdk2-cpp-robot-b2-robotstateclient-low-power-status-1}{%
\paragraph{\texorpdfstring{\texttt{unitree\_sdk2\_cpp.robot.b2.RobotStateClient.low\_power\_status}}{unitree\_sdk2\_cpp.robot.b2.RobotStateClient.low\_power\_status}}\label{unitree-sdk2-cpp-robot-b2-robotstateclient-low-power-status-1}}

对应 C++ SDK 操作
\passthrough{\lstinline!LowPowerStatus(int32\_t \&)!}。上游头文件没有可直接生成的业务说明时，本参考只保证签名映射，精确语义需查目标型号协议。

\textbf{签名}

\begin{lstlisting}[language=Python]
def low_power_status(self) -> tuple[int, int]
\end{lstlisting}

\textbf{可用性与安全性}

\textbf{\passthrough{\lstinline!SIGNATURE\_ONLY!} /
\passthrough{\lstinline!HARDWARE\_SIDE\_EFFECT!}}：仅用于补全和静态检查。当前二进制没有该
Python 方法，未来实现还需单独评审硬件副作用。

\textbf{参数}

无显式参数。实例方法中的 \passthrough{\lstinline!self!} 由 Python
自动传入。

\textbf{返回值}

\begin{longtable}[]{@{}
  >{\raggedright\arraybackslash}p{(\columnwidth - 4\tabcolsep) * \real{0.18}}
  >{\raggedright\arraybackslash}p{(\columnwidth - 4\tabcolsep) * \real{0.20}}
  >{\raggedright\arraybackslash}p{(\columnwidth - 4\tabcolsep) * \real{0.62}}@{}}
\toprule
位置 & 类型 & 含义 \\
\midrule
\endhead
{[}0{]} \passthrough{\lstinline!status!} & \passthrough{\lstinline!int!}
& SDK 状态码；Unitree 示例通常以 \passthrough{\lstinline!0!}
表示成功，非零值需按具体服务错误码解释。 \\
{[}1{]} \passthrough{\lstinline!status!} & \passthrough{\lstinline!int!}
& SDK 状态码；Unitree 示例通常以 \passthrough{\lstinline!0!}
表示成功，非零值需按具体服务错误码解释。 \\
\bottomrule
\end{longtable}

\textbf{对应 C++}

\begin{itemize}
\tightlist
\item
  类：\passthrough{\lstinline!unitree::robot::b2::RobotStateClient!}
\item
  签名：\passthrough{\lstinline!LowPowerStatus(int32\_t \&)!}
\item
  绑定策略：\passthrough{\lstinline!OUTPUT\_WRAPPER!}
\item
  声明位置：\passthrough{\lstinline!include/unitree/robot/b2/robot\_state/robot\_state\_client.hpp:41!}
\end{itemize}

\textbf{说明}

\begin{itemize}
\tightlist
\item
  示例是局部调用片段；\passthrough{\lstinline!obj!}
  和参数变量表示已经按上层业务校验并准备好的对象。
\item
  C++ 的可修改输出引用不会作为 Python 入参出现，而是按 C++
  参数顺序追加到返回元组中。
\item
  该条目可以被编辑器和类型检查器识别，但当前运行时可能在属性查找阶段直接失败。
\end{itemize}

\textbf{用法}

\begin{lstlisting}[language=Python]
from typing import TYPE_CHECKING

if TYPE_CHECKING:
    def planned_low_power_status(obj: RobotStateClient) -> tuple[int, int]:
        return obj.low_power_status()
\end{lstlisting}

\begin{quote}
{[}!CAUTION{]} 上面的 \passthrough{\lstinline!TYPE\_CHECKING!}
示例只表达计划调用形态。该分支在普通 Python
运行时为假，不代表当前扩展能执行此接口。
\end{quote}

\hypertarget{unitree-sdk2-cpp-robot-b2-robotstateclient-get-pkg-version-1}{%
\paragraph{\texorpdfstring{\texttt{unitree\_sdk2\_cpp.robot.b2.RobotStateClient.get\_pkg\_version}}{unitree\_sdk2\_cpp.robot.b2.RobotStateClient.get\_pkg\_version}}\label{unitree-sdk2-cpp-robot-b2-robotstateclient-get-pkg-version-1}}

查询或检查 \passthrough{\lstinline!pkg!}
\passthrough{\lstinline!version!}。

\textbf{签名}

\begin{lstlisting}[language=Python]
def get_pkg_version(self) -> tuple[int, str, dict[str, str]]
\end{lstlisting}

\textbf{可用性与安全性}

\textbf{\passthrough{\lstinline!AVAILABLE!} /
\passthrough{\lstinline!READ\_ONLY!}}：当前绑定源码已实现只读查询。Robot
Client 的服务端查询通常仍需匹配的 Linux
扩展、DDS、网络、机器人和服务；纯本地 getter 则不一定需要全部条件。

\textbf{参数}

无显式参数。实例方法中的 \passthrough{\lstinline!self!} 由 Python
自动传入。

\textbf{返回值}

\begin{longtable}[]{@{}
  >{\raggedright\arraybackslash}p{(\columnwidth - 4\tabcolsep) * \real{0.18}}
  >{\raggedright\arraybackslash}p{(\columnwidth - 4\tabcolsep) * \real{0.20}}
  >{\raggedright\arraybackslash}p{(\columnwidth - 4\tabcolsep) * \real{0.62}}@{}}
\toprule
位置 & 类型 & 含义 \\
\midrule
\endhead
{[}0{]} \passthrough{\lstinline!status!} & \passthrough{\lstinline!int!}
& SDK 状态码；Unitree 示例通常以 \passthrough{\lstinline!0!}
表示成功，非零值需按具体服务错误码解释。 \\
{[}1{]} \passthrough{\lstinline!packageVersion!} &
\passthrough{\lstinline!str!} & 传给该接口的
\passthrough{\lstinline!packageVersion!} 参数，Python 类型为
\passthrough{\lstinline!str!}。精确含义、单位、范围和枚举值需查目标型号对应头文件与协议。
对应 C++ 参数 \passthrough{\lstinline!packageVersion: std::string \&!}。
底层为字符串；长度、编码和允许值由具体协议决定。 \\
{[}2{]} \passthrough{\lstinline!moduleVersionMap!} &
\passthrough{\lstinline!dict[str, str]!} & 传给该接口的
\passthrough{\lstinline!moduleVersionMap!} 参数，Python 类型为
\passthrough{\lstinline!dict[str, str]!}。精确含义、单位、范围和枚举值需查目标型号对应头文件与协议。
对应 C++ 参数
\passthrough{\lstinline!moduleVersionMap: std::map<std::string, std::string> \&!}。 \\
\bottomrule
\end{longtable}

\textbf{对应 C++}

\begin{itemize}
\tightlist
\item
  类：\passthrough{\lstinline!unitree::robot::b2::RobotStateClient!}
\item
  签名：\passthrough{\lstinline!GetPkgVersion(std::string \&, std::map<std::string, std::string> \&)!}
\item
  绑定策略：\passthrough{\lstinline!OUTPUT\_WRAPPER!}
\item
  声明位置：\passthrough{\lstinline!include/unitree/robot/b2/robot\_state/robot\_state\_client.hpp:42!}
\end{itemize}

\textbf{说明}

\begin{itemize}
\tightlist
\item
  示例是局部调用片段；\passthrough{\lstinline!obj!}
  和参数变量表示已经按上层业务校验并准备好的对象。
\item
  C++ 的可修改输出引用不会作为 Python 入参出现，而是按 C++
  参数顺序追加到返回元组中。
\end{itemize}

\textbf{用法}

\begin{lstlisting}[language=Python]
status, packageVersion, moduleVersionMap = obj.get_pkg_version()
\end{lstlisting}

\begin{quote}
{[}!NOTE{]} 只读查询仍会访问
DDS/机器人服务。必须检查返回状态码，并处理超时和网络中断。
\end{quote}

\hypertarget{unitree-sdk2-cpp-robot-b2-servicestate}{%
\subsubsection{\texorpdfstring{\texttt{unitree\_sdk2\_cpp.robot.b2.ServiceState}}{unitree\_sdk2\_cpp.robot.b2.ServiceState}}\label{unitree-sdk2-cpp-robot-b2-servicestate}}

SDK 数据值类型；公开字段和转换方法见下文。

\textbf{导入}

\begin{lstlisting}[language=Python]
from unitree_sdk2_cpp.robot.b2 import ServiceState
\end{lstlisting}

\textbf{基类}：\passthrough{\lstinline!object!}

\textbf{公开属性}

\begin{longtable}[]{@{}
  >{\raggedright\arraybackslash}p{(\columnwidth - 6\tabcolsep) * \real{0.18}}
  >{\raggedright\arraybackslash}p{(\columnwidth - 6\tabcolsep) * \real{0.24}}
  >{\raggedright\arraybackslash}p{(\columnwidth - 6\tabcolsep) * \real{0.18}}
  >{\raggedright\arraybackslash}p{(\columnwidth - 6\tabcolsep) * \real{0.40}}@{}}
\toprule
名称 & Python 类型 & 含义 & 用法 \\
\midrule
\endhead
\passthrough{\lstinline!name!} & \passthrough{\lstinline!str!} & SDK
对象、配置项、服务或动作的名称。允许值由调用它的具体接口定义。 &
\passthrough{\lstinline!value = obj.name!} /
\passthrough{\lstinline!obj.name = value!} \\
\passthrough{\lstinline!status!} & \passthrough{\lstinline!int!} &
传给该接口的 状态 参数，Python 类型为
\passthrough{\lstinline!int!}。精确含义、单位、范围和枚举值需查目标型号对应头文件与协议。
& \passthrough{\lstinline!value = obj.status!} /
\passthrough{\lstinline!obj.status = value!} \\
\passthrough{\lstinline!protect!} & \passthrough{\lstinline!int!} &
传给该接口的 \passthrough{\lstinline!protect!} 参数，Python 类型为
\passthrough{\lstinline!int!}。精确含义、单位、范围和枚举值需查目标型号对应头文件与协议。
& \passthrough{\lstinline!value = obj.protect!} /
\passthrough{\lstinline!obj.protect = value!} \\
\bottomrule
\end{longtable}

\textbf{方法索引}

\begin{longtable}[]{@{}
  >{\raggedright\arraybackslash}p{(\columnwidth - 6\tabcolsep) * \real{0.18}}
  >{\raggedright\arraybackslash}p{(\columnwidth - 6\tabcolsep) * \real{0.24}}
  >{\raggedright\arraybackslash}p{(\columnwidth - 6\tabcolsep) * \real{0.18}}
  >{\raggedright\arraybackslash}p{(\columnwidth - 6\tabcolsep) * \real{0.40}}@{}}
\toprule
方法 & Python 签名 & 状态 & 安全分类 \\
\midrule
\endhead
\protect\hyperlink{unitree-sdk2-cpp-robot-b2-servicestate-dunder-init-1}{\passthrough{\lstinline!\_\_init\_\_!}}
& \passthrough{\lstinline!def \_\_init\_\_(self) -> None!} &
\passthrough{\lstinline!AVAILABLE!} &
\passthrough{\lstinline!CONSTRUCTION!} \\
\bottomrule
\end{longtable}

\hypertarget{unitree-sdk2-cpp-robot-b2-servicestate-dunder-init-1}{%
\paragraph{\texorpdfstring{\texttt{unitree\_sdk2\_cpp.robot.b2.ServiceState.\_\_init\_\_}}{unitree\_sdk2\_cpp.robot.b2.ServiceState.\_\_init\_\_}}\label{unitree-sdk2-cpp-robot-b2-servicestate-dunder-init-1}}

初始化 \passthrough{\lstinline!ServiceState!}
实例。是否能实际构造取决于下方可用性状态。

\textbf{签名}

\begin{lstlisting}[language=Python]
def __init__(self) -> None
\end{lstlisting}

\textbf{可用性与安全性}

\textbf{\passthrough{\lstinline!AVAILABLE!} /
\passthrough{\lstinline!CONSTRUCTION!}}：当前绑定源码已实现；实际执行条件仍取决于平台和对应
SDK 环境。

\textbf{参数}

无显式参数。实例方法中的 \passthrough{\lstinline!self!} 由 Python
自动传入。

\textbf{返回值}

\begin{longtable}[]{@{}
  >{\raggedright\arraybackslash}p{(\columnwidth - 4\tabcolsep) * \real{0.18}}
  >{\raggedright\arraybackslash}p{(\columnwidth - 4\tabcolsep) * \real{0.20}}
  >{\raggedright\arraybackslash}p{(\columnwidth - 4\tabcolsep) * \real{0.62}}@{}}
\toprule
位置 & 类型 & 含义 \\
\midrule
\endhead
无 & \passthrough{\lstinline!None!} &
\passthrough{\lstinline!\_\_init\_\_!}
本身不返回值；调用类对象时，在构造函数可用的前提下得到该类实例。 \\
\bottomrule
\end{longtable}

\textbf{对应 C++}

\begin{itemize}
\tightlist
\item
  类：\passthrough{\lstinline!unitree::robot::b2::ServiceState!}
\item
  签名：\passthrough{\lstinline!ServiceState()!}
\item
  绑定策略：\passthrough{\lstinline!未单独标注!}
\item
  声明位置：\passthrough{\lstinline!include/unitree/robot/b2/robot\_state/robot\_state\_client.hpp:18!}
\end{itemize}

\textbf{说明}

\begin{itemize}
\tightlist
\item
  示例是局部调用片段；\passthrough{\lstinline!obj!}
  和参数变量表示已经按上层业务校验并准备好的对象。
\end{itemize}

\textbf{用法}

\begin{lstlisting}[language=Python]
value = ServiceState()
\end{lstlisting}

\hypertarget{unitree-sdk2-cpp-robot-b2-servicestatedata}{%
\subsubsection{\texorpdfstring{\texttt{unitree\_sdk2\_cpp.robot.b2.ServiceStateData}}{unitree\_sdk2\_cpp.robot.b2.ServiceStateData}}\label{unitree-sdk2-cpp-robot-b2-servicestatedata}}

SDK 数据值类型；公开字段和转换方法见下文。

\textbf{导入}

\begin{lstlisting}[language=Python]
from unitree_sdk2_cpp.robot.b2 import ServiceStateData
\end{lstlisting}

\textbf{基类}：\passthrough{\lstinline!object!}

\textbf{公开属性}

\begin{longtable}[]{@{}
  >{\raggedright\arraybackslash}p{(\columnwidth - 6\tabcolsep) * \real{0.18}}
  >{\raggedright\arraybackslash}p{(\columnwidth - 6\tabcolsep) * \real{0.24}}
  >{\raggedright\arraybackslash}p{(\columnwidth - 6\tabcolsep) * \real{0.18}}
  >{\raggedright\arraybackslash}p{(\columnwidth - 6\tabcolsep) * \real{0.40}}@{}}
\toprule
名称 & Python 类型 & 含义 & 用法 \\
\midrule
\endhead
\passthrough{\lstinline!name!} & \passthrough{\lstinline!str!} & SDK
对象、配置项、服务或动作的名称。允许值由调用它的具体接口定义。 &
\passthrough{\lstinline!value = obj.name!} /
\passthrough{\lstinline!obj.name = value!} \\
\passthrough{\lstinline!status!} & \passthrough{\lstinline!int!} &
传给该接口的 状态 参数，Python 类型为
\passthrough{\lstinline!int!}。精确含义、单位、范围和枚举值需查目标型号对应头文件与协议。
& \passthrough{\lstinline!value = obj.status!} /
\passthrough{\lstinline!obj.status = value!} \\
\passthrough{\lstinline!protect!} & \passthrough{\lstinline!int!} &
传给该接口的 \passthrough{\lstinline!protect!} 参数，Python 类型为
\passthrough{\lstinline!int!}。精确含义、单位、范围和枚举值需查目标型号对应头文件与协议。
& \passthrough{\lstinline!value = obj.protect!} /
\passthrough{\lstinline!obj.protect = value!} \\
\bottomrule
\end{longtable}

\textbf{方法索引}

\begin{longtable}[]{@{}
  >{\raggedright\arraybackslash}p{(\columnwidth - 6\tabcolsep) * \real{0.18}}
  >{\raggedright\arraybackslash}p{(\columnwidth - 6\tabcolsep) * \real{0.24}}
  >{\raggedright\arraybackslash}p{(\columnwidth - 6\tabcolsep) * \real{0.18}}
  >{\raggedright\arraybackslash}p{(\columnwidth - 6\tabcolsep) * \real{0.40}}@{}}
\toprule
方法 & Python 签名 & 状态 & 安全分类 \\
\midrule
\endhead
\protect\hyperlink{unitree-sdk2-cpp-robot-b2-servicestatedata-dunder-init-1}{\passthrough{\lstinline!\_\_init\_\_!}}
& \passthrough{\lstinline!def \_\_init\_\_(self) -> None!} &
\passthrough{\lstinline!SIGNATURE\_ONLY!} &
\passthrough{\lstinline!CONSTRUCTION!} \\
\protect\hyperlink{unitree-sdk2-cpp-robot-b2-servicestatedata-from-json-1}{\passthrough{\lstinline!from\_json!}}
&
\passthrough{\lstinline!def from\_json(self, json: dict[str, Any]) -> None!}
& \passthrough{\lstinline!SIGNATURE\_ONLY!} &
\passthrough{\lstinline!UNCLASSIFIED!} \\
\protect\hyperlink{unitree-sdk2-cpp-robot-b2-servicestatedata-to-json-1}{\passthrough{\lstinline!to\_json!}}
&
\passthrough{\lstinline!def to\_json(self, json: dict[str, Any]) -> None!}
& \passthrough{\lstinline!SIGNATURE\_ONLY!} &
\passthrough{\lstinline!UNCLASSIFIED!} \\
\bottomrule
\end{longtable}

\hypertarget{unitree-sdk2-cpp-robot-b2-servicestatedata-dunder-init-1}{%
\paragraph{\texorpdfstring{\texttt{unitree\_sdk2\_cpp.robot.b2.ServiceStateData.\_\_init\_\_}}{unitree\_sdk2\_cpp.robot.b2.ServiceStateData.\_\_init\_\_}}\label{unitree-sdk2-cpp-robot-b2-servicestatedata-dunder-init-1}}

初始化 \passthrough{\lstinline!ServiceStateData!}
实例。是否能实际构造取决于下方可用性状态。

\textbf{签名}

\begin{lstlisting}[language=Python]
def __init__(self) -> None
\end{lstlisting}

\textbf{可用性与安全性}

\textbf{\passthrough{\lstinline!SIGNATURE\_ONLY!} /
\passthrough{\lstinline!CONSTRUCTION!}}：当前只有设计期签名，不能假设已存在于运行时扩展。

\textbf{参数}

无显式参数。实例方法中的 \passthrough{\lstinline!self!} 由 Python
自动传入。

\textbf{返回值}

\begin{longtable}[]{@{}
  >{\raggedright\arraybackslash}p{(\columnwidth - 4\tabcolsep) * \real{0.18}}
  >{\raggedright\arraybackslash}p{(\columnwidth - 4\tabcolsep) * \real{0.20}}
  >{\raggedright\arraybackslash}p{(\columnwidth - 4\tabcolsep) * \real{0.62}}@{}}
\toprule
位置 & 类型 & 含义 \\
\midrule
\endhead
无 & \passthrough{\lstinline!None!} &
\passthrough{\lstinline!\_\_init\_\_!}
本身不返回值；调用类对象时，在构造函数可用的前提下得到该类实例。 \\
\bottomrule
\end{longtable}

\textbf{对应 C++}

\begin{itemize}
\tightlist
\item
  类：\passthrough{\lstinline!unitree::robot::b2::ServiceStateData!}
\item
  签名：\passthrough{\lstinline!ServiceStateData()!}
\item
  绑定策略：\passthrough{\lstinline!未单独标注!}
\item
  声明位置：\passthrough{\lstinline!include/unitree/robot/b2/robot\_state/robot\_state\_api.hpp:125!}
\end{itemize}

\textbf{说明}

\begin{itemize}
\tightlist
\item
  示例是局部调用片段；\passthrough{\lstinline!obj!}
  和参数变量表示已经按上层业务校验并准备好的对象。
\item
  该条目可以被编辑器和类型检查器识别，但当前运行时可能在属性查找阶段直接失败。
\end{itemize}

\textbf{用法}

\begin{lstlisting}[language=Python]
from typing import TYPE_CHECKING

if TYPE_CHECKING:
    def planned_construct() -> ServiceStateData:
        return ServiceStateData()
\end{lstlisting}

\begin{quote}
{[}!CAUTION{]} 上面的 \passthrough{\lstinline!TYPE\_CHECKING!}
示例只表达计划调用形态。该分支在普通 Python
运行时为假，不代表当前扩展能执行此接口。
\end{quote}

\hypertarget{unitree-sdk2-cpp-robot-b2-servicestatedata-from-json-1}{%
\paragraph{\texorpdfstring{\texttt{unitree\_sdk2\_cpp.robot.b2.ServiceStateData.from\_json}}{unitree\_sdk2\_cpp.robot.b2.ServiceStateData.from\_json}}\label{unitree-sdk2-cpp-robot-b2-servicestatedata-from-json-1}}

计划从 JSON 风格字典读取字段并更新当前 SDK 值对象。

\textbf{签名}

\begin{lstlisting}[language=Python]
def from_json(self, json: dict[str, Any]) -> None
\end{lstlisting}

\textbf{可用性与安全性}

\textbf{\passthrough{\lstinline!SIGNATURE\_ONLY!} /
\passthrough{\lstinline!UNCLASSIFIED!}}：当前只有设计期签名，不能假设已存在于运行时扩展。

\textbf{参数}

\begin{longtable}[]{@{}
  >{\raggedright\arraybackslash}p{(\columnwidth - 6\tabcolsep) * \real{0.18}}
  >{\raggedright\arraybackslash}p{(\columnwidth - 6\tabcolsep) * \real{0.24}}
  >{\raggedright\arraybackslash}p{(\columnwidth - 6\tabcolsep) * \real{0.18}}
  >{\raggedright\arraybackslash}p{(\columnwidth - 6\tabcolsep) * \real{0.40}}@{}}
\toprule
名称 & Python 类型 & 默认值 & 含义 \\
\midrule
\endhead
\passthrough{\lstinline!json!} &
\passthrough{\lstinline!dict[str, Any]!} & 必填 & JSON
风格字典，键和值必须满足对应 C++ \passthrough{\lstinline!JsonMap!}
数据结构的约束。 对应 C++ 参数
\passthrough{\lstinline!json: common::JsonMap \&!}。 \\
\bottomrule
\end{longtable}

\textbf{返回值}

\begin{longtable}[]{@{}
  >{\raggedright\arraybackslash}p{(\columnwidth - 4\tabcolsep) * \real{0.18}}
  >{\raggedright\arraybackslash}p{(\columnwidth - 4\tabcolsep) * \real{0.20}}
  >{\raggedright\arraybackslash}p{(\columnwidth - 4\tabcolsep) * \real{0.62}}@{}}
\toprule
位置 & 类型 & 含义 \\
\midrule
\endhead
无 & \passthrough{\lstinline!None!} & 该方法不返回 Python
值；失败可能表现为异常或底层状态变化。 \\
\bottomrule
\end{longtable}

\textbf{对应 C++}

\begin{itemize}
\tightlist
\item
  类：\passthrough{\lstinline!unitree::robot::b2::ServiceStateData!}
\item
  签名：\passthrough{\lstinline!fromJson(common::JsonMap \&)!}
\item
  绑定策略：\passthrough{\lstinline!SIGNATURE\_PREVIEW!}
\item
  声明位置：\passthrough{\lstinline!include/unitree/robot/b2/robot\_state/robot\_state\_api.hpp:131!}
\end{itemize}

\textbf{说明}

\begin{itemize}
\tightlist
\item
  示例是局部调用片段；\passthrough{\lstinline!obj!}
  和参数变量表示已经按上层业务校验并准备好的对象。
\item
  该条目可以被编辑器和类型检查器识别，但当前运行时可能在属性查找阶段直接失败。
\end{itemize}

\textbf{用法}

\begin{lstlisting}[language=Python]
from typing import TYPE_CHECKING

if TYPE_CHECKING:
    def planned_from_json(obj: ServiceStateData, json: dict[str, Any]) -> None:
        obj.from_json(json=json)
\end{lstlisting}

\begin{quote}
{[}!CAUTION{]} 上面的 \passthrough{\lstinline!TYPE\_CHECKING!}
示例只表达计划调用形态。该分支在普通 Python
运行时为假，不代表当前扩展能执行此接口。
\end{quote}

\hypertarget{unitree-sdk2-cpp-robot-b2-servicestatedata-to-json-1}{%
\paragraph{\texorpdfstring{\texttt{unitree\_sdk2\_cpp.robot.b2.ServiceStateData.to\_json}}{unitree\_sdk2\_cpp.robot.b2.ServiceStateData.to\_json}}\label{unitree-sdk2-cpp-robot-b2-servicestatedata-to-json-1}}

计划把当前 SDK 值对象写入 JSON 风格字典。

\textbf{签名}

\begin{lstlisting}[language=Python]
def to_json(self, json: dict[str, Any]) -> None
\end{lstlisting}

\textbf{可用性与安全性}

\textbf{\passthrough{\lstinline!SIGNATURE\_ONLY!} /
\passthrough{\lstinline!UNCLASSIFIED!}}：当前只有设计期签名，不能假设已存在于运行时扩展。

\textbf{参数}

\begin{longtable}[]{@{}
  >{\raggedright\arraybackslash}p{(\columnwidth - 6\tabcolsep) * \real{0.18}}
  >{\raggedright\arraybackslash}p{(\columnwidth - 6\tabcolsep) * \real{0.24}}
  >{\raggedright\arraybackslash}p{(\columnwidth - 6\tabcolsep) * \real{0.18}}
  >{\raggedright\arraybackslash}p{(\columnwidth - 6\tabcolsep) * \real{0.40}}@{}}
\toprule
名称 & Python 类型 & 默认值 & 含义 \\
\midrule
\endhead
\passthrough{\lstinline!json!} &
\passthrough{\lstinline!dict[str, Any]!} & 必填 & JSON
风格字典，键和值必须满足对应 C++ \passthrough{\lstinline!JsonMap!}
数据结构的约束。 对应 C++ 参数
\passthrough{\lstinline!json: common::JsonMap \&!}。 \\
\bottomrule
\end{longtable}

\textbf{返回值}

\begin{longtable}[]{@{}
  >{\raggedright\arraybackslash}p{(\columnwidth - 4\tabcolsep) * \real{0.18}}
  >{\raggedright\arraybackslash}p{(\columnwidth - 4\tabcolsep) * \real{0.20}}
  >{\raggedright\arraybackslash}p{(\columnwidth - 4\tabcolsep) * \real{0.62}}@{}}
\toprule
位置 & 类型 & 含义 \\
\midrule
\endhead
无 & \passthrough{\lstinline!None!} & 该方法不返回 Python
值；失败可能表现为异常或底层状态变化。 \\
\bottomrule
\end{longtable}

\textbf{对应 C++}

\begin{itemize}
\tightlist
\item
  类：\passthrough{\lstinline!unitree::robot::b2::ServiceStateData!}
\item
  签名：\passthrough{\lstinline!toJson(common::JsonMap \&) const!}
\item
  绑定策略：\passthrough{\lstinline!SIGNATURE\_PREVIEW!}
\item
  声明位置：\passthrough{\lstinline!include/unitree/robot/b2/robot\_state/robot\_state\_api.hpp:138!}
\end{itemize}

\textbf{说明}

\begin{itemize}
\tightlist
\item
  示例是局部调用片段；\passthrough{\lstinline!obj!}
  和参数变量表示已经按上层业务校验并准备好的对象。
\item
  该条目可以被编辑器和类型检查器识别，但当前运行时可能在属性查找阶段直接失败。
\end{itemize}

\textbf{用法}

\begin{lstlisting}[language=Python]
from typing import TYPE_CHECKING

if TYPE_CHECKING:
    def planned_to_json(obj: ServiceStateData, json: dict[str, Any]) -> None:
        obj.to_json(json=json)
\end{lstlisting}

\begin{quote}
{[}!CAUTION{]} 上面的 \passthrough{\lstinline!TYPE\_CHECKING!}
示例只表达计划调用形态。该分支在普通 Python
运行时为假，不代表当前扩展能执行此接口。
\end{quote}

\hypertarget{unitree-sdk2-cpp-robot-b2-serviceswitchdata}{%
\subsubsection{\texorpdfstring{\texttt{unitree\_sdk2\_cpp.robot.b2.ServiceSwitchData}}{unitree\_sdk2\_cpp.robot.b2.ServiceSwitchData}}\label{unitree-sdk2-cpp-robot-b2-serviceswitchdata}}

SDK 数据值类型；公开字段和转换方法见下文。

\textbf{导入}

\begin{lstlisting}[language=Python]
from unitree_sdk2_cpp.robot.b2 import ServiceSwitchData
\end{lstlisting}

\textbf{基类}：\passthrough{\lstinline!object!}

\textbf{公开属性}

\begin{longtable}[]{@{}
  >{\raggedright\arraybackslash}p{(\columnwidth - 6\tabcolsep) * \real{0.18}}
  >{\raggedright\arraybackslash}p{(\columnwidth - 6\tabcolsep) * \real{0.24}}
  >{\raggedright\arraybackslash}p{(\columnwidth - 6\tabcolsep) * \real{0.18}}
  >{\raggedright\arraybackslash}p{(\columnwidth - 6\tabcolsep) * \real{0.40}}@{}}
\toprule
名称 & Python 类型 & 含义 & 用法 \\
\midrule
\endhead
\passthrough{\lstinline!name!} & \passthrough{\lstinline!str!} & SDK
对象、配置项、服务或动作的名称。允许值由调用它的具体接口定义。 &
\passthrough{\lstinline!value = obj.name!} /
\passthrough{\lstinline!obj.name = value!} \\
\passthrough{\lstinline!status!} & \passthrough{\lstinline!int!} &
传给该接口的 状态 参数，Python 类型为
\passthrough{\lstinline!int!}。精确含义、单位、范围和枚举值需查目标型号对应头文件与协议。
& \passthrough{\lstinline!value = obj.status!} /
\passthrough{\lstinline!obj.status = value!} \\
\bottomrule
\end{longtable}

\textbf{方法索引}

\begin{longtable}[]{@{}
  >{\raggedright\arraybackslash}p{(\columnwidth - 6\tabcolsep) * \real{0.18}}
  >{\raggedright\arraybackslash}p{(\columnwidth - 6\tabcolsep) * \real{0.24}}
  >{\raggedright\arraybackslash}p{(\columnwidth - 6\tabcolsep) * \real{0.18}}
  >{\raggedright\arraybackslash}p{(\columnwidth - 6\tabcolsep) * \real{0.40}}@{}}
\toprule
方法 & Python 签名 & 状态 & 安全分类 \\
\midrule
\endhead
\protect\hyperlink{unitree-sdk2-cpp-robot-b2-serviceswitchdata-dunder-init-1}{\passthrough{\lstinline!\_\_init\_\_!}}
& \passthrough{\lstinline!def \_\_init\_\_(self) -> None!} &
\passthrough{\lstinline!SIGNATURE\_ONLY!} &
\passthrough{\lstinline!CONSTRUCTION!} \\
\protect\hyperlink{unitree-sdk2-cpp-robot-b2-serviceswitchdata-from-json-1}{\passthrough{\lstinline!from\_json!}}
&
\passthrough{\lstinline!def from\_json(self, json: dict[str, Any]) -> None!}
& \passthrough{\lstinline!SIGNATURE\_ONLY!} &
\passthrough{\lstinline!UNCLASSIFIED!} \\
\protect\hyperlink{unitree-sdk2-cpp-robot-b2-serviceswitchdata-to-json-1}{\passthrough{\lstinline!to\_json!}}
&
\passthrough{\lstinline!def to\_json(self, json: dict[str, Any]) -> None!}
& \passthrough{\lstinline!SIGNATURE\_ONLY!} &
\passthrough{\lstinline!UNCLASSIFIED!} \\
\bottomrule
\end{longtable}

\hypertarget{unitree-sdk2-cpp-robot-b2-serviceswitchdata-dunder-init-1}{%
\paragraph{\texorpdfstring{\texttt{unitree\_sdk2\_cpp.robot.b2.ServiceSwitchData.\_\_init\_\_}}{unitree\_sdk2\_cpp.robot.b2.ServiceSwitchData.\_\_init\_\_}}\label{unitree-sdk2-cpp-robot-b2-serviceswitchdata-dunder-init-1}}

初始化 \passthrough{\lstinline!ServiceSwitchData!}
实例。是否能实际构造取决于下方可用性状态。

\textbf{签名}

\begin{lstlisting}[language=Python]
def __init__(self) -> None
\end{lstlisting}

\textbf{可用性与安全性}

\textbf{\passthrough{\lstinline!SIGNATURE\_ONLY!} /
\passthrough{\lstinline!CONSTRUCTION!}}：当前只有设计期签名，不能假设已存在于运行时扩展。

\textbf{参数}

无显式参数。实例方法中的 \passthrough{\lstinline!self!} 由 Python
自动传入。

\textbf{返回值}

\begin{longtable}[]{@{}
  >{\raggedright\arraybackslash}p{(\columnwidth - 4\tabcolsep) * \real{0.18}}
  >{\raggedright\arraybackslash}p{(\columnwidth - 4\tabcolsep) * \real{0.20}}
  >{\raggedright\arraybackslash}p{(\columnwidth - 4\tabcolsep) * \real{0.62}}@{}}
\toprule
位置 & 类型 & 含义 \\
\midrule
\endhead
无 & \passthrough{\lstinline!None!} &
\passthrough{\lstinline!\_\_init\_\_!}
本身不返回值；调用类对象时，在构造函数可用的前提下得到该类实例。 \\
\bottomrule
\end{longtable}

\textbf{对应 C++}

\begin{itemize}
\tightlist
\item
  类：\passthrough{\lstinline!unitree::robot::b2::ServiceSwitchData!}
\item
  签名：\passthrough{\lstinline!ServiceSwitchData()!}
\item
  绑定策略：\passthrough{\lstinline!未单独标注!}
\item
  声明位置：\passthrough{\lstinline!include/unitree/robot/b2/robot\_state/robot\_state\_api.hpp:67!}
\end{itemize}

\textbf{说明}

\begin{itemize}
\tightlist
\item
  示例是局部调用片段；\passthrough{\lstinline!obj!}
  和参数变量表示已经按上层业务校验并准备好的对象。
\item
  该条目可以被编辑器和类型检查器识别，但当前运行时可能在属性查找阶段直接失败。
\end{itemize}

\textbf{用法}

\begin{lstlisting}[language=Python]
from typing import TYPE_CHECKING

if TYPE_CHECKING:
    def planned_construct() -> ServiceSwitchData:
        return ServiceSwitchData()
\end{lstlisting}

\begin{quote}
{[}!CAUTION{]} 上面的 \passthrough{\lstinline!TYPE\_CHECKING!}
示例只表达计划调用形态。该分支在普通 Python
运行时为假，不代表当前扩展能执行此接口。
\end{quote}

\hypertarget{unitree-sdk2-cpp-robot-b2-serviceswitchdata-from-json-1}{%
\paragraph{\texorpdfstring{\texttt{unitree\_sdk2\_cpp.robot.b2.ServiceSwitchData.from\_json}}{unitree\_sdk2\_cpp.robot.b2.ServiceSwitchData.from\_json}}\label{unitree-sdk2-cpp-robot-b2-serviceswitchdata-from-json-1}}

计划从 JSON 风格字典读取字段并更新当前 SDK 值对象。

\textbf{签名}

\begin{lstlisting}[language=Python]
def from_json(self, json: dict[str, Any]) -> None
\end{lstlisting}

\textbf{可用性与安全性}

\textbf{\passthrough{\lstinline!SIGNATURE\_ONLY!} /
\passthrough{\lstinline!UNCLASSIFIED!}}：当前只有设计期签名，不能假设已存在于运行时扩展。

\textbf{参数}

\begin{longtable}[]{@{}
  >{\raggedright\arraybackslash}p{(\columnwidth - 6\tabcolsep) * \real{0.18}}
  >{\raggedright\arraybackslash}p{(\columnwidth - 6\tabcolsep) * \real{0.24}}
  >{\raggedright\arraybackslash}p{(\columnwidth - 6\tabcolsep) * \real{0.18}}
  >{\raggedright\arraybackslash}p{(\columnwidth - 6\tabcolsep) * \real{0.40}}@{}}
\toprule
名称 & Python 类型 & 默认值 & 含义 \\
\midrule
\endhead
\passthrough{\lstinline!json!} &
\passthrough{\lstinline!dict[str, Any]!} & 必填 & JSON
风格字典，键和值必须满足对应 C++ \passthrough{\lstinline!JsonMap!}
数据结构的约束。 对应 C++ 参数
\passthrough{\lstinline!json: common::JsonMap \&!}。 \\
\bottomrule
\end{longtable}

\textbf{返回值}

\begin{longtable}[]{@{}
  >{\raggedright\arraybackslash}p{(\columnwidth - 4\tabcolsep) * \real{0.18}}
  >{\raggedright\arraybackslash}p{(\columnwidth - 4\tabcolsep) * \real{0.20}}
  >{\raggedright\arraybackslash}p{(\columnwidth - 4\tabcolsep) * \real{0.62}}@{}}
\toprule
位置 & 类型 & 含义 \\
\midrule
\endhead
无 & \passthrough{\lstinline!None!} & 该方法不返回 Python
值；失败可能表现为异常或底层状态变化。 \\
\bottomrule
\end{longtable}

\textbf{对应 C++}

\begin{itemize}
\tightlist
\item
  类：\passthrough{\lstinline!unitree::robot::b2::ServiceSwitchData!}
\item
  签名：\passthrough{\lstinline!fromJson(common::JsonMap \&)!}
\item
  绑定策略：\passthrough{\lstinline!SIGNATURE\_PREVIEW!}
\item
  声明位置：\passthrough{\lstinline!include/unitree/robot/b2/robot\_state/robot\_state\_api.hpp:73!}
\end{itemize}

\textbf{说明}

\begin{itemize}
\tightlist
\item
  示例是局部调用片段；\passthrough{\lstinline!obj!}
  和参数变量表示已经按上层业务校验并准备好的对象。
\item
  该条目可以被编辑器和类型检查器识别，但当前运行时可能在属性查找阶段直接失败。
\end{itemize}

\textbf{用法}

\begin{lstlisting}[language=Python]
from typing import TYPE_CHECKING

if TYPE_CHECKING:
    def planned_from_json(obj: ServiceSwitchData, json: dict[str, Any]) -> None:
        obj.from_json(json=json)
\end{lstlisting}

\begin{quote}
{[}!CAUTION{]} 上面的 \passthrough{\lstinline!TYPE\_CHECKING!}
示例只表达计划调用形态。该分支在普通 Python
运行时为假，不代表当前扩展能执行此接口。
\end{quote}

\hypertarget{unitree-sdk2-cpp-robot-b2-serviceswitchdata-to-json-1}{%
\paragraph{\texorpdfstring{\texttt{unitree\_sdk2\_cpp.robot.b2.ServiceSwitchData.to\_json}}{unitree\_sdk2\_cpp.robot.b2.ServiceSwitchData.to\_json}}\label{unitree-sdk2-cpp-robot-b2-serviceswitchdata-to-json-1}}

计划把当前 SDK 值对象写入 JSON 风格字典。

\textbf{签名}

\begin{lstlisting}[language=Python]
def to_json(self, json: dict[str, Any]) -> None
\end{lstlisting}

\textbf{可用性与安全性}

\textbf{\passthrough{\lstinline!SIGNATURE\_ONLY!} /
\passthrough{\lstinline!UNCLASSIFIED!}}：当前只有设计期签名，不能假设已存在于运行时扩展。

\textbf{参数}

\begin{longtable}[]{@{}
  >{\raggedright\arraybackslash}p{(\columnwidth - 6\tabcolsep) * \real{0.18}}
  >{\raggedright\arraybackslash}p{(\columnwidth - 6\tabcolsep) * \real{0.24}}
  >{\raggedright\arraybackslash}p{(\columnwidth - 6\tabcolsep) * \real{0.18}}
  >{\raggedright\arraybackslash}p{(\columnwidth - 6\tabcolsep) * \real{0.40}}@{}}
\toprule
名称 & Python 类型 & 默认值 & 含义 \\
\midrule
\endhead
\passthrough{\lstinline!json!} &
\passthrough{\lstinline!dict[str, Any]!} & 必填 & JSON
风格字典，键和值必须满足对应 C++ \passthrough{\lstinline!JsonMap!}
数据结构的约束。 对应 C++ 参数
\passthrough{\lstinline!json: common::JsonMap \&!}。 \\
\bottomrule
\end{longtable}

\textbf{返回值}

\begin{longtable}[]{@{}
  >{\raggedright\arraybackslash}p{(\columnwidth - 4\tabcolsep) * \real{0.18}}
  >{\raggedright\arraybackslash}p{(\columnwidth - 4\tabcolsep) * \real{0.20}}
  >{\raggedright\arraybackslash}p{(\columnwidth - 4\tabcolsep) * \real{0.62}}@{}}
\toprule
位置 & 类型 & 含义 \\
\midrule
\endhead
无 & \passthrough{\lstinline!None!} & 该方法不返回 Python
值；失败可能表现为异常或底层状态变化。 \\
\bottomrule
\end{longtable}

\textbf{对应 C++}

\begin{itemize}
\tightlist
\item
  类：\passthrough{\lstinline!unitree::robot::b2::ServiceSwitchData!}
\item
  签名：\passthrough{\lstinline!toJson(common::JsonMap \&) const!}
\item
  绑定策略：\passthrough{\lstinline!SIGNATURE\_PREVIEW!}
\item
  声明位置：\passthrough{\lstinline!include/unitree/robot/b2/robot\_state/robot\_state\_api.hpp:79!}
\end{itemize}

\textbf{说明}

\begin{itemize}
\tightlist
\item
  示例是局部调用片段；\passthrough{\lstinline!obj!}
  和参数变量表示已经按上层业务校验并准备好的对象。
\item
  该条目可以被编辑器和类型检查器识别，但当前运行时可能在属性查找阶段直接失败。
\end{itemize}

\textbf{用法}

\begin{lstlisting}[language=Python]
from typing import TYPE_CHECKING

if TYPE_CHECKING:
    def planned_to_json(obj: ServiceSwitchData, json: dict[str, Any]) -> None:
        obj.to_json(json=json)
\end{lstlisting}

\begin{quote}
{[}!CAUTION{]} 上面的 \passthrough{\lstinline!TYPE\_CHECKING!}
示例只表达计划调用形态。该分支在普通 Python
运行时为假，不代表当前扩展能执行此接口。
\end{quote}

\hypertarget{unitree-sdk2-cpp-robot-b2-serviceswitchparameter}{%
\subsubsection{\texorpdfstring{\texttt{unitree\_sdk2\_cpp.robot.b2.ServiceSwitchParameter}}{unitree\_sdk2\_cpp.robot.b2.ServiceSwitchParameter}}\label{unitree-sdk2-cpp-robot-b2-serviceswitchparameter}}

SDK 请求参数值类型；当前可能仅提供设计期签名。

\textbf{导入}

\begin{lstlisting}[language=Python]
from unitree_sdk2_cpp.robot.b2 import ServiceSwitchParameter
\end{lstlisting}

\textbf{基类}：\passthrough{\lstinline!object!}

\textbf{公开属性}

\begin{longtable}[]{@{}
  >{\raggedright\arraybackslash}p{(\columnwidth - 6\tabcolsep) * \real{0.18}}
  >{\raggedright\arraybackslash}p{(\columnwidth - 6\tabcolsep) * \real{0.24}}
  >{\raggedright\arraybackslash}p{(\columnwidth - 6\tabcolsep) * \real{0.18}}
  >{\raggedright\arraybackslash}p{(\columnwidth - 6\tabcolsep) * \real{0.40}}@{}}
\toprule
名称 & Python 类型 & 含义 & 用法 \\
\midrule
\endhead
\passthrough{\lstinline!name!} & \passthrough{\lstinline!str!} & SDK
对象、配置项、服务或动作的名称。允许值由调用它的具体接口定义。 &
\passthrough{\lstinline!value = obj.name!} /
\passthrough{\lstinline!obj.name = value!} \\
\passthrough{\lstinline!swit!} & \passthrough{\lstinline!int!} &
传给该接口的 \passthrough{\lstinline!swit!} 参数，Python 类型为
\passthrough{\lstinline!int!}。精确含义、单位、范围和枚举值需查目标型号对应头文件与协议。
& \passthrough{\lstinline!value = obj.swit!} /
\passthrough{\lstinline!obj.swit = value!} \\
\bottomrule
\end{longtable}

\textbf{方法索引}

\begin{longtable}[]{@{}
  >{\raggedright\arraybackslash}p{(\columnwidth - 6\tabcolsep) * \real{0.18}}
  >{\raggedright\arraybackslash}p{(\columnwidth - 6\tabcolsep) * \real{0.24}}
  >{\raggedright\arraybackslash}p{(\columnwidth - 6\tabcolsep) * \real{0.18}}
  >{\raggedright\arraybackslash}p{(\columnwidth - 6\tabcolsep) * \real{0.40}}@{}}
\toprule
方法 & Python 签名 & 状态 & 安全分类 \\
\midrule
\endhead
\protect\hyperlink{unitree-sdk2-cpp-robot-b2-serviceswitchparameter-dunder-init-1}{\passthrough{\lstinline!\_\_init\_\_!}}
& \passthrough{\lstinline!def \_\_init\_\_(self) -> None!} &
\passthrough{\lstinline!SIGNATURE\_ONLY!} &
\passthrough{\lstinline!CONSTRUCTION!} \\
\protect\hyperlink{unitree-sdk2-cpp-robot-b2-serviceswitchparameter-from-json-1}{\passthrough{\lstinline!from\_json!}}
&
\passthrough{\lstinline!def from\_json(self, json: dict[str, Any]) -> None!}
& \passthrough{\lstinline!SIGNATURE\_ONLY!} &
\passthrough{\lstinline!UNCLASSIFIED!} \\
\protect\hyperlink{unitree-sdk2-cpp-robot-b2-serviceswitchparameter-to-json-1}{\passthrough{\lstinline!to\_json!}}
&
\passthrough{\lstinline!def to\_json(self, json: dict[str, Any]) -> None!}
& \passthrough{\lstinline!SIGNATURE\_ONLY!} &
\passthrough{\lstinline!UNCLASSIFIED!} \\
\bottomrule
\end{longtable}

\hypertarget{unitree-sdk2-cpp-robot-b2-serviceswitchparameter-dunder-init-1}{%
\paragraph{\texorpdfstring{\texttt{unitree\_sdk2\_cpp.robot.b2.ServiceSwitchParameter.\_\_init\_\_}}{unitree\_sdk2\_cpp.robot.b2.ServiceSwitchParameter.\_\_init\_\_}}\label{unitree-sdk2-cpp-robot-b2-serviceswitchparameter-dunder-init-1}}

初始化 \passthrough{\lstinline!ServiceSwitchParameter!}
实例。是否能实际构造取决于下方可用性状态。

\textbf{签名}

\begin{lstlisting}[language=Python]
def __init__(self) -> None
\end{lstlisting}

\textbf{可用性与安全性}

\textbf{\passthrough{\lstinline!SIGNATURE\_ONLY!} /
\passthrough{\lstinline!CONSTRUCTION!}}：当前只有设计期签名，不能假设已存在于运行时扩展。

\textbf{参数}

无显式参数。实例方法中的 \passthrough{\lstinline!self!} 由 Python
自动传入。

\textbf{返回值}

\begin{longtable}[]{@{}
  >{\raggedright\arraybackslash}p{(\columnwidth - 4\tabcolsep) * \real{0.18}}
  >{\raggedright\arraybackslash}p{(\columnwidth - 4\tabcolsep) * \real{0.20}}
  >{\raggedright\arraybackslash}p{(\columnwidth - 4\tabcolsep) * \real{0.62}}@{}}
\toprule
位置 & 类型 & 含义 \\
\midrule
\endhead
无 & \passthrough{\lstinline!None!} &
\passthrough{\lstinline!\_\_init\_\_!}
本身不返回值；调用类对象时，在构造函数可用的前提下得到该类实例。 \\
\bottomrule
\end{longtable}

\textbf{对应 C++}

\begin{itemize}
\tightlist
\item
  类：\passthrough{\lstinline!unitree::robot::b2::ServiceSwitchParameter!}
\item
  签名：\passthrough{\lstinline!ServiceSwitchParameter()!}
\item
  绑定策略：\passthrough{\lstinline!未单独标注!}
\item
  声明位置：\passthrough{\lstinline!include/unitree/robot/b2/robot\_state/robot\_state\_api.hpp:38!}
\end{itemize}

\textbf{说明}

\begin{itemize}
\tightlist
\item
  示例是局部调用片段；\passthrough{\lstinline!obj!}
  和参数变量表示已经按上层业务校验并准备好的对象。
\item
  该条目可以被编辑器和类型检查器识别，但当前运行时可能在属性查找阶段直接失败。
\end{itemize}

\textbf{用法}

\begin{lstlisting}[language=Python]
from typing import TYPE_CHECKING

if TYPE_CHECKING:
    def planned_construct() -> ServiceSwitchParameter:
        return ServiceSwitchParameter()
\end{lstlisting}

\begin{quote}
{[}!CAUTION{]} 上面的 \passthrough{\lstinline!TYPE\_CHECKING!}
示例只表达计划调用形态。该分支在普通 Python
运行时为假，不代表当前扩展能执行此接口。
\end{quote}

\hypertarget{unitree-sdk2-cpp-robot-b2-serviceswitchparameter-from-json-1}{%
\paragraph{\texorpdfstring{\texttt{unitree\_sdk2\_cpp.robot.b2.ServiceSwitchParameter.from\_json}}{unitree\_sdk2\_cpp.robot.b2.ServiceSwitchParameter.from\_json}}\label{unitree-sdk2-cpp-robot-b2-serviceswitchparameter-from-json-1}}

计划从 JSON 风格字典读取字段并更新当前 SDK 值对象。

\textbf{签名}

\begin{lstlisting}[language=Python]
def from_json(self, json: dict[str, Any]) -> None
\end{lstlisting}

\textbf{可用性与安全性}

\textbf{\passthrough{\lstinline!SIGNATURE\_ONLY!} /
\passthrough{\lstinline!UNCLASSIFIED!}}：当前只有设计期签名，不能假设已存在于运行时扩展。

\textbf{参数}

\begin{longtable}[]{@{}
  >{\raggedright\arraybackslash}p{(\columnwidth - 6\tabcolsep) * \real{0.18}}
  >{\raggedright\arraybackslash}p{(\columnwidth - 6\tabcolsep) * \real{0.24}}
  >{\raggedright\arraybackslash}p{(\columnwidth - 6\tabcolsep) * \real{0.18}}
  >{\raggedright\arraybackslash}p{(\columnwidth - 6\tabcolsep) * \real{0.40}}@{}}
\toprule
名称 & Python 类型 & 默认值 & 含义 \\
\midrule
\endhead
\passthrough{\lstinline!json!} &
\passthrough{\lstinline!dict[str, Any]!} & 必填 & JSON
风格字典，键和值必须满足对应 C++ \passthrough{\lstinline!JsonMap!}
数据结构的约束。 对应 C++ 参数
\passthrough{\lstinline!json: common::JsonMap \&!}。 \\
\bottomrule
\end{longtable}

\textbf{返回值}

\begin{longtable}[]{@{}
  >{\raggedright\arraybackslash}p{(\columnwidth - 4\tabcolsep) * \real{0.18}}
  >{\raggedright\arraybackslash}p{(\columnwidth - 4\tabcolsep) * \real{0.20}}
  >{\raggedright\arraybackslash}p{(\columnwidth - 4\tabcolsep) * \real{0.62}}@{}}
\toprule
位置 & 类型 & 含义 \\
\midrule
\endhead
无 & \passthrough{\lstinline!None!} & 该方法不返回 Python
值；失败可能表现为异常或底层状态变化。 \\
\bottomrule
\end{longtable}

\textbf{对应 C++}

\begin{itemize}
\tightlist
\item
  类：\passthrough{\lstinline!unitree::robot::b2::ServiceSwitchParameter!}
\item
  签名：\passthrough{\lstinline!fromJson(common::JsonMap \&)!}
\item
  绑定策略：\passthrough{\lstinline!SIGNATURE\_PREVIEW!}
\item
  声明位置：\passthrough{\lstinline!include/unitree/robot/b2/robot\_state/robot\_state\_api.hpp:44!}
\end{itemize}

\textbf{说明}

\begin{itemize}
\tightlist
\item
  示例是局部调用片段；\passthrough{\lstinline!obj!}
  和参数变量表示已经按上层业务校验并准备好的对象。
\item
  该条目可以被编辑器和类型检查器识别，但当前运行时可能在属性查找阶段直接失败。
\end{itemize}

\textbf{用法}

\begin{lstlisting}[language=Python]
from typing import TYPE_CHECKING

if TYPE_CHECKING:
    def planned_from_json(obj: ServiceSwitchParameter, json: dict[str, Any]) -> None:
        obj.from_json(json=json)
\end{lstlisting}

\begin{quote}
{[}!CAUTION{]} 上面的 \passthrough{\lstinline!TYPE\_CHECKING!}
示例只表达计划调用形态。该分支在普通 Python
运行时为假，不代表当前扩展能执行此接口。
\end{quote}

\hypertarget{unitree-sdk2-cpp-robot-b2-serviceswitchparameter-to-json-1}{%
\paragraph{\texorpdfstring{\texttt{unitree\_sdk2\_cpp.robot.b2.ServiceSwitchParameter.to\_json}}{unitree\_sdk2\_cpp.robot.b2.ServiceSwitchParameter.to\_json}}\label{unitree-sdk2-cpp-robot-b2-serviceswitchparameter-to-json-1}}

计划把当前 SDK 值对象写入 JSON 风格字典。

\textbf{签名}

\begin{lstlisting}[language=Python]
def to_json(self, json: dict[str, Any]) -> None
\end{lstlisting}

\textbf{可用性与安全性}

\textbf{\passthrough{\lstinline!SIGNATURE\_ONLY!} /
\passthrough{\lstinline!UNCLASSIFIED!}}：当前只有设计期签名，不能假设已存在于运行时扩展。

\textbf{参数}

\begin{longtable}[]{@{}
  >{\raggedright\arraybackslash}p{(\columnwidth - 6\tabcolsep) * \real{0.18}}
  >{\raggedright\arraybackslash}p{(\columnwidth - 6\tabcolsep) * \real{0.24}}
  >{\raggedright\arraybackslash}p{(\columnwidth - 6\tabcolsep) * \real{0.18}}
  >{\raggedright\arraybackslash}p{(\columnwidth - 6\tabcolsep) * \real{0.40}}@{}}
\toprule
名称 & Python 类型 & 默认值 & 含义 \\
\midrule
\endhead
\passthrough{\lstinline!json!} &
\passthrough{\lstinline!dict[str, Any]!} & 必填 & JSON
风格字典，键和值必须满足对应 C++ \passthrough{\lstinline!JsonMap!}
数据结构的约束。 对应 C++ 参数
\passthrough{\lstinline!json: common::JsonMap \&!}。 \\
\bottomrule
\end{longtable}

\textbf{返回值}

\begin{longtable}[]{@{}
  >{\raggedright\arraybackslash}p{(\columnwidth - 4\tabcolsep) * \real{0.18}}
  >{\raggedright\arraybackslash}p{(\columnwidth - 4\tabcolsep) * \real{0.20}}
  >{\raggedright\arraybackslash}p{(\columnwidth - 4\tabcolsep) * \real{0.62}}@{}}
\toprule
位置 & 类型 & 含义 \\
\midrule
\endhead
无 & \passthrough{\lstinline!None!} & 该方法不返回 Python
值；失败可能表现为异常或底层状态变化。 \\
\bottomrule
\end{longtable}

\textbf{对应 C++}

\begin{itemize}
\tightlist
\item
  类：\passthrough{\lstinline!unitree::robot::b2::ServiceSwitchParameter!}
\item
  签名：\passthrough{\lstinline!toJson(common::JsonMap \&) const!}
\item
  绑定策略：\passthrough{\lstinline!SIGNATURE\_PREVIEW!}
\item
  声明位置：\passthrough{\lstinline!include/unitree/robot/b2/robot\_state/robot\_state\_api.hpp:50!}
\end{itemize}

\textbf{说明}

\begin{itemize}
\tightlist
\item
  示例是局部调用片段；\passthrough{\lstinline!obj!}
  和参数变量表示已经按上层业务校验并准备好的对象。
\item
  该条目可以被编辑器和类型检查器识别，但当前运行时可能在属性查找阶段直接失败。
\end{itemize}

\textbf{用法}

\begin{lstlisting}[language=Python]
from typing import TYPE_CHECKING

if TYPE_CHECKING:
    def planned_to_json(obj: ServiceSwitchParameter, json: dict[str, Any]) -> None:
        obj.to_json(json=json)
\end{lstlisting}

\begin{quote}
{[}!CAUTION{]} 上面的 \passthrough{\lstinline!TYPE\_CHECKING!}
示例只表达计划调用形态。该分支在普通 Python
运行时为假，不代表当前扩展能执行此接口。
\end{quote}

\hypertarget{unitree-sdk2-cpp-robot-b2-setreportfreqparameter}{%
\subsubsection{\texorpdfstring{\texttt{unitree\_sdk2\_cpp.robot.b2.SetReportFreqParameter}}{unitree\_sdk2\_cpp.robot.b2.SetReportFreqParameter}}\label{unitree-sdk2-cpp-robot-b2-setreportfreqparameter}}

SDK 请求参数值类型；当前可能仅提供设计期签名。

\textbf{导入}

\begin{lstlisting}[language=Python]
from unitree_sdk2_cpp.robot.b2 import SetReportFreqParameter
\end{lstlisting}

\textbf{基类}：\passthrough{\lstinline!object!}

\textbf{公开属性}

\begin{longtable}[]{@{}
  >{\raggedright\arraybackslash}p{(\columnwidth - 6\tabcolsep) * \real{0.18}}
  >{\raggedright\arraybackslash}p{(\columnwidth - 6\tabcolsep) * \real{0.24}}
  >{\raggedright\arraybackslash}p{(\columnwidth - 6\tabcolsep) * \real{0.18}}
  >{\raggedright\arraybackslash}p{(\columnwidth - 6\tabcolsep) * \real{0.40}}@{}}
\toprule
名称 & Python 类型 & 含义 & 用法 \\
\midrule
\endhead
\passthrough{\lstinline!interval!} & \passthrough{\lstinline!int!} &
传给该接口的 \passthrough{\lstinline!interval!} 参数，Python 类型为
\passthrough{\lstinline!int!}。精确含义、单位、范围和枚举值需查目标型号对应头文件与协议。
& \passthrough{\lstinline!value = obj.interval!} /
\passthrough{\lstinline!obj.interval = value!} \\
\passthrough{\lstinline!duration!} & \passthrough{\lstinline!int!} &
操作持续时间参数。精确单位、范围和默认行为以目标型号接口定义为准。 &
\passthrough{\lstinline!value = obj.duration!} /
\passthrough{\lstinline!obj.duration = value!} \\
\bottomrule
\end{longtable}

\textbf{方法索引}

\begin{longtable}[]{@{}
  >{\raggedright\arraybackslash}p{(\columnwidth - 6\tabcolsep) * \real{0.18}}
  >{\raggedright\arraybackslash}p{(\columnwidth - 6\tabcolsep) * \real{0.24}}
  >{\raggedright\arraybackslash}p{(\columnwidth - 6\tabcolsep) * \real{0.18}}
  >{\raggedright\arraybackslash}p{(\columnwidth - 6\tabcolsep) * \real{0.40}}@{}}
\toprule
方法 & Python 签名 & 状态 & 安全分类 \\
\midrule
\endhead
\protect\hyperlink{unitree-sdk2-cpp-robot-b2-setreportfreqparameter-dunder-init-1}{\passthrough{\lstinline!\_\_init\_\_!}}
& \passthrough{\lstinline!def \_\_init\_\_(self) -> None!} &
\passthrough{\lstinline!SIGNATURE\_ONLY!} &
\passthrough{\lstinline!CONSTRUCTION!} \\
\protect\hyperlink{unitree-sdk2-cpp-robot-b2-setreportfreqparameter-from-json-1}{\passthrough{\lstinline!from\_json!}}
&
\passthrough{\lstinline!def from\_json(self, json: dict[str, Any]) -> None!}
& \passthrough{\lstinline!SIGNATURE\_ONLY!} &
\passthrough{\lstinline!UNCLASSIFIED!} \\
\protect\hyperlink{unitree-sdk2-cpp-robot-b2-setreportfreqparameter-to-json-1}{\passthrough{\lstinline!to\_json!}}
&
\passthrough{\lstinline!def to\_json(self, json: dict[str, Any]) -> None!}
& \passthrough{\lstinline!SIGNATURE\_ONLY!} &
\passthrough{\lstinline!UNCLASSIFIED!} \\
\bottomrule
\end{longtable}

\hypertarget{unitree-sdk2-cpp-robot-b2-setreportfreqparameter-dunder-init-1}{%
\paragraph{\texorpdfstring{\texttt{unitree\_sdk2\_cpp.robot.b2.SetReportFreqParameter.\_\_init\_\_}}{unitree\_sdk2\_cpp.robot.b2.SetReportFreqParameter.\_\_init\_\_}}\label{unitree-sdk2-cpp-robot-b2-setreportfreqparameter-dunder-init-1}}

初始化 \passthrough{\lstinline!SetReportFreqParameter!}
实例。是否能实际构造取决于下方可用性状态。

\textbf{签名}

\begin{lstlisting}[language=Python]
def __init__(self) -> None
\end{lstlisting}

\textbf{可用性与安全性}

\textbf{\passthrough{\lstinline!SIGNATURE\_ONLY!} /
\passthrough{\lstinline!CONSTRUCTION!}}：当前只有设计期签名，不能假设已存在于运行时扩展。

\textbf{参数}

无显式参数。实例方法中的 \passthrough{\lstinline!self!} 由 Python
自动传入。

\textbf{返回值}

\begin{longtable}[]{@{}
  >{\raggedright\arraybackslash}p{(\columnwidth - 4\tabcolsep) * \real{0.18}}
  >{\raggedright\arraybackslash}p{(\columnwidth - 4\tabcolsep) * \real{0.20}}
  >{\raggedright\arraybackslash}p{(\columnwidth - 4\tabcolsep) * \real{0.62}}@{}}
\toprule
位置 & 类型 & 含义 \\
\midrule
\endhead
无 & \passthrough{\lstinline!None!} &
\passthrough{\lstinline!\_\_init\_\_!}
本身不返回值；调用类对象时，在构造函数可用的前提下得到该类实例。 \\
\bottomrule
\end{longtable}

\textbf{对应 C++}

\begin{itemize}
\tightlist
\item
  类：\passthrough{\lstinline!unitree::robot::b2::SetReportFreqParameter!}
\item
  签名：\passthrough{\lstinline!SetReportFreqParameter()!}
\item
  绑定策略：\passthrough{\lstinline!未单独标注!}
\item
  声明位置：\passthrough{\lstinline!include/unitree/robot/b2/robot\_state/robot\_state\_api.hpp:96!}
\end{itemize}

\textbf{说明}

\begin{itemize}
\tightlist
\item
  示例是局部调用片段；\passthrough{\lstinline!obj!}
  和参数变量表示已经按上层业务校验并准备好的对象。
\item
  该条目可以被编辑器和类型检查器识别，但当前运行时可能在属性查找阶段直接失败。
\end{itemize}

\textbf{用法}

\begin{lstlisting}[language=Python]
from typing import TYPE_CHECKING

if TYPE_CHECKING:
    def planned_construct() -> SetReportFreqParameter:
        return SetReportFreqParameter()
\end{lstlisting}

\begin{quote}
{[}!CAUTION{]} 上面的 \passthrough{\lstinline!TYPE\_CHECKING!}
示例只表达计划调用形态。该分支在普通 Python
运行时为假，不代表当前扩展能执行此接口。
\end{quote}

\hypertarget{unitree-sdk2-cpp-robot-b2-setreportfreqparameter-from-json-1}{%
\paragraph{\texorpdfstring{\texttt{unitree\_sdk2\_cpp.robot.b2.SetReportFreqParameter.from\_json}}{unitree\_sdk2\_cpp.robot.b2.SetReportFreqParameter.from\_json}}\label{unitree-sdk2-cpp-robot-b2-setreportfreqparameter-from-json-1}}

计划从 JSON 风格字典读取字段并更新当前 SDK 值对象。

\textbf{签名}

\begin{lstlisting}[language=Python]
def from_json(self, json: dict[str, Any]) -> None
\end{lstlisting}

\textbf{可用性与安全性}

\textbf{\passthrough{\lstinline!SIGNATURE\_ONLY!} /
\passthrough{\lstinline!UNCLASSIFIED!}}：当前只有设计期签名，不能假设已存在于运行时扩展。

\textbf{参数}

\begin{longtable}[]{@{}
  >{\raggedright\arraybackslash}p{(\columnwidth - 6\tabcolsep) * \real{0.18}}
  >{\raggedright\arraybackslash}p{(\columnwidth - 6\tabcolsep) * \real{0.24}}
  >{\raggedright\arraybackslash}p{(\columnwidth - 6\tabcolsep) * \real{0.18}}
  >{\raggedright\arraybackslash}p{(\columnwidth - 6\tabcolsep) * \real{0.40}}@{}}
\toprule
名称 & Python 类型 & 默认值 & 含义 \\
\midrule
\endhead
\passthrough{\lstinline!json!} &
\passthrough{\lstinline!dict[str, Any]!} & 必填 & JSON
风格字典，键和值必须满足对应 C++ \passthrough{\lstinline!JsonMap!}
数据结构的约束。 对应 C++ 参数
\passthrough{\lstinline!json: common::JsonMap \&!}。 \\
\bottomrule
\end{longtable}

\textbf{返回值}

\begin{longtable}[]{@{}
  >{\raggedright\arraybackslash}p{(\columnwidth - 4\tabcolsep) * \real{0.18}}
  >{\raggedright\arraybackslash}p{(\columnwidth - 4\tabcolsep) * \real{0.20}}
  >{\raggedright\arraybackslash}p{(\columnwidth - 4\tabcolsep) * \real{0.62}}@{}}
\toprule
位置 & 类型 & 含义 \\
\midrule
\endhead
无 & \passthrough{\lstinline!None!} & 该方法不返回 Python
值；失败可能表现为异常或底层状态变化。 \\
\bottomrule
\end{longtable}

\textbf{对应 C++}

\begin{itemize}
\tightlist
\item
  类：\passthrough{\lstinline!unitree::robot::b2::SetReportFreqParameter!}
\item
  签名：\passthrough{\lstinline!fromJson(common::JsonMap \&)!}
\item
  绑定策略：\passthrough{\lstinline!SIGNATURE\_PREVIEW!}
\item
  声明位置：\passthrough{\lstinline!include/unitree/robot/b2/robot\_state/robot\_state\_api.hpp:102!}
\end{itemize}

\textbf{说明}

\begin{itemize}
\tightlist
\item
  示例是局部调用片段；\passthrough{\lstinline!obj!}
  和参数变量表示已经按上层业务校验并准备好的对象。
\item
  该条目可以被编辑器和类型检查器识别，但当前运行时可能在属性查找阶段直接失败。
\end{itemize}

\textbf{用法}

\begin{lstlisting}[language=Python]
from typing import TYPE_CHECKING

if TYPE_CHECKING:
    def planned_from_json(obj: SetReportFreqParameter, json: dict[str, Any]) -> None:
        obj.from_json(json=json)
\end{lstlisting}

\begin{quote}
{[}!CAUTION{]} 上面的 \passthrough{\lstinline!TYPE\_CHECKING!}
示例只表达计划调用形态。该分支在普通 Python
运行时为假，不代表当前扩展能执行此接口。
\end{quote}

\hypertarget{unitree-sdk2-cpp-robot-b2-setreportfreqparameter-to-json-1}{%
\paragraph{\texorpdfstring{\texttt{unitree\_sdk2\_cpp.robot.b2.SetReportFreqParameter.to\_json}}{unitree\_sdk2\_cpp.robot.b2.SetReportFreqParameter.to\_json}}\label{unitree-sdk2-cpp-robot-b2-setreportfreqparameter-to-json-1}}

计划把当前 SDK 值对象写入 JSON 风格字典。

\textbf{签名}

\begin{lstlisting}[language=Python]
def to_json(self, json: dict[str, Any]) -> None
\end{lstlisting}

\textbf{可用性与安全性}

\textbf{\passthrough{\lstinline!SIGNATURE\_ONLY!} /
\passthrough{\lstinline!UNCLASSIFIED!}}：当前只有设计期签名，不能假设已存在于运行时扩展。

\textbf{参数}

\begin{longtable}[]{@{}
  >{\raggedright\arraybackslash}p{(\columnwidth - 6\tabcolsep) * \real{0.18}}
  >{\raggedright\arraybackslash}p{(\columnwidth - 6\tabcolsep) * \real{0.24}}
  >{\raggedright\arraybackslash}p{(\columnwidth - 6\tabcolsep) * \real{0.18}}
  >{\raggedright\arraybackslash}p{(\columnwidth - 6\tabcolsep) * \real{0.40}}@{}}
\toprule
名称 & Python 类型 & 默认值 & 含义 \\
\midrule
\endhead
\passthrough{\lstinline!json!} &
\passthrough{\lstinline!dict[str, Any]!} & 必填 & JSON
风格字典，键和值必须满足对应 C++ \passthrough{\lstinline!JsonMap!}
数据结构的约束。 对应 C++ 参数
\passthrough{\lstinline!json: common::JsonMap \&!}。 \\
\bottomrule
\end{longtable}

\textbf{返回值}

\begin{longtable}[]{@{}
  >{\raggedright\arraybackslash}p{(\columnwidth - 4\tabcolsep) * \real{0.18}}
  >{\raggedright\arraybackslash}p{(\columnwidth - 4\tabcolsep) * \real{0.20}}
  >{\raggedright\arraybackslash}p{(\columnwidth - 4\tabcolsep) * \real{0.62}}@{}}
\toprule
位置 & 类型 & 含义 \\
\midrule
\endhead
无 & \passthrough{\lstinline!None!} & 该方法不返回 Python
值；失败可能表现为异常或底层状态变化。 \\
\bottomrule
\end{longtable}

\textbf{对应 C++}

\begin{itemize}
\tightlist
\item
  类：\passthrough{\lstinline!unitree::robot::b2::SetReportFreqParameter!}
\item
  签名：\passthrough{\lstinline!toJson(common::JsonMap \&) const!}
\item
  绑定策略：\passthrough{\lstinline!SIGNATURE\_PREVIEW!}
\item
  声明位置：\passthrough{\lstinline!include/unitree/robot/b2/robot\_state/robot\_state\_api.hpp:108!}
\end{itemize}

\textbf{说明}

\begin{itemize}
\tightlist
\item
  示例是局部调用片段；\passthrough{\lstinline!obj!}
  和参数变量表示已经按上层业务校验并准备好的对象。
\item
  该条目可以被编辑器和类型检查器识别，但当前运行时可能在属性查找阶段直接失败。
\end{itemize}

\textbf{用法}

\begin{lstlisting}[language=Python]
from typing import TYPE_CHECKING

if TYPE_CHECKING:
    def planned_to_json(obj: SetReportFreqParameter, json: dict[str, Any]) -> None:
        obj.to_json(json=json)
\end{lstlisting}

\begin{quote}
{[}!CAUTION{]} 上面的 \passthrough{\lstinline!TYPE\_CHECKING!}
示例只表达计划调用形态。该分支在普通 Python
运行时为假，不代表当前扩展能执行此接口。
\end{quote}

\hypertarget{unitree-sdk2-cpp-robot-b2-sportclient}{%
\subsubsection{\texorpdfstring{\texttt{unitree\_sdk2\_cpp.robot.b2.SportClient}}{unitree\_sdk2\_cpp.robot.b2.SportClient}}\label{unitree-sdk2-cpp-robot-b2-sportclient}}

Robot 服务客户端。构造、初始化和具体方法的可用性必须分别检查。

\textbf{导入}

\begin{lstlisting}[language=Python]
from unitree_sdk2_cpp.robot.b2 import SportClient
\end{lstlisting}

\textbf{基类}：\passthrough{\lstinline!Client!}

\textbf{方法索引}

\begin{longtable}[]{@{}
  >{\raggedright\arraybackslash}p{(\columnwidth - 6\tabcolsep) * \real{0.18}}
  >{\raggedright\arraybackslash}p{(\columnwidth - 6\tabcolsep) * \real{0.24}}
  >{\raggedright\arraybackslash}p{(\columnwidth - 6\tabcolsep) * \real{0.18}}
  >{\raggedright\arraybackslash}p{(\columnwidth - 6\tabcolsep) * \real{0.40}}@{}}
\toprule
方法 & Python 签名 & 状态 & 安全分类 \\
\midrule
\endhead
\protect\hyperlink{unitree-sdk2-cpp-robot-b2-sportclient-dunder-init-1}{\passthrough{\lstinline!\_\_init\_\_!}}
&
\passthrough{\lstinline!def \_\_init\_\_(self, enable\_lease: bool = ...) -> None!}
& \passthrough{\lstinline!SIGNATURE\_ONLY!} &
\passthrough{\lstinline!CONSTRUCTION!} \\
\protect\hyperlink{unitree-sdk2-cpp-robot-b2-sportclient-init-1}{\passthrough{\lstinline!init!}}
& \passthrough{\lstinline!def init(self) -> None!} &
\passthrough{\lstinline!SIGNATURE\_ONLY!} &
\passthrough{\lstinline!INITIALIZATION!} \\
\protect\hyperlink{unitree-sdk2-cpp-robot-b2-sportclient-damp-1}{\passthrough{\lstinline!damp!}}
& \passthrough{\lstinline!def damp(self) -> int!} &
\passthrough{\lstinline!SIGNATURE\_ONLY!} &
\passthrough{\lstinline!MOTION\_COMMAND!} \\
\protect\hyperlink{unitree-sdk2-cpp-robot-b2-sportclient-balance-stand-1}{\passthrough{\lstinline!balance\_stand!}}
& \passthrough{\lstinline!def balance\_stand(self) -> int!} &
\passthrough{\lstinline!SIGNATURE\_ONLY!} &
\passthrough{\lstinline!MOTION\_COMMAND!} \\
\protect\hyperlink{unitree-sdk2-cpp-robot-b2-sportclient-stop-move-1}{\passthrough{\lstinline!stop\_move!}}
& \passthrough{\lstinline!def stop\_move(self) -> int!} &
\passthrough{\lstinline!SIGNATURE\_ONLY!} &
\passthrough{\lstinline!MOTION\_COMMAND!} \\
\protect\hyperlink{unitree-sdk2-cpp-robot-b2-sportclient-stand-up-1}{\passthrough{\lstinline!stand\_up!}}
& \passthrough{\lstinline!def stand\_up(self) -> int!} &
\passthrough{\lstinline!SIGNATURE\_ONLY!} &
\passthrough{\lstinline!MOTION\_COMMAND!} \\
\protect\hyperlink{unitree-sdk2-cpp-robot-b2-sportclient-stand-down-1}{\passthrough{\lstinline!stand\_down!}}
& \passthrough{\lstinline!def stand\_down(self) -> int!} &
\passthrough{\lstinline!SIGNATURE\_ONLY!} &
\passthrough{\lstinline!MOTION\_COMMAND!} \\
\protect\hyperlink{unitree-sdk2-cpp-robot-b2-sportclient-recovery-stand-1}{\passthrough{\lstinline!recovery\_stand!}}
& \passthrough{\lstinline!def recovery\_stand(self) -> int!} &
\passthrough{\lstinline!SIGNATURE\_ONLY!} &
\passthrough{\lstinline!MOTION\_COMMAND!} \\
\protect\hyperlink{unitree-sdk2-cpp-robot-b2-sportclient-move-1}{\passthrough{\lstinline!move!}}
&
\passthrough{\lstinline!def move(self, vx: float, vy: float, vyaw: float) -> int!}
& \passthrough{\lstinline!SIGNATURE\_ONLY!} &
\passthrough{\lstinline!MOTION\_COMMAND!} \\
\protect\hyperlink{unitree-sdk2-cpp-robot-b2-sportclient-switch-gait-1}{\passthrough{\lstinline!switch\_gait!}}
& \passthrough{\lstinline!def switch\_gait(self, d: int) -> int!} &
\passthrough{\lstinline!SIGNATURE\_ONLY!} &
\passthrough{\lstinline!MOTION\_COMMAND!} \\
\protect\hyperlink{unitree-sdk2-cpp-robot-b2-sportclient-body-height-1}{\passthrough{\lstinline!body\_height!}}
& \passthrough{\lstinline!def body\_height(self, height: float) -> int!}
& \passthrough{\lstinline!SIGNATURE\_ONLY!} &
\passthrough{\lstinline!MOTION\_COMMAND!} \\
\protect\hyperlink{unitree-sdk2-cpp-robot-b2-sportclient-speed-level-1}{\passthrough{\lstinline!speed\_level!}}
& \passthrough{\lstinline!def speed\_level(self, level: int) -> int!} &
\passthrough{\lstinline!SIGNATURE\_ONLY!} &
\passthrough{\lstinline!MOTION\_COMMAND!} \\
\protect\hyperlink{unitree-sdk2-cpp-robot-b2-sportclient-trajectory-follow-1}{\passthrough{\lstinline!trajectory\_follow!}}
&
\passthrough{\lstinline!def trajectory\_follow(self) -> tuple[int, list[Any]]!}
& \passthrough{\lstinline!SIGNATURE\_ONLY!} &
\passthrough{\lstinline!MOTION\_COMMAND!} \\
\protect\hyperlink{unitree-sdk2-cpp-robot-b2-sportclient-continuous-gait-1}{\passthrough{\lstinline!continuous\_gait!}}
&
\passthrough{\lstinline!def continuous\_gait(self, flag: bool) -> int!}
& \passthrough{\lstinline!SIGNATURE\_ONLY!} &
\passthrough{\lstinline!MOTION\_COMMAND!} \\
\protect\hyperlink{unitree-sdk2-cpp-robot-b2-sportclient-move-to-pos-1}{\passthrough{\lstinline!move\_to\_pos!}}
&
\passthrough{\lstinline!def move\_to\_pos(self, x: float, y: float, yaw: float) -> int!}
& \passthrough{\lstinline!SIGNATURE\_ONLY!} &
\passthrough{\lstinline!MOTION\_COMMAND!} \\
\protect\hyperlink{unitree-sdk2-cpp-robot-b2-sportclient-switch-move-mode-1}{\passthrough{\lstinline!switch\_move\_mode!}}
&
\passthrough{\lstinline!def switch\_move\_mode(self, flag: bool) -> int!}
& \passthrough{\lstinline!SIGNATURE\_ONLY!} &
\passthrough{\lstinline!MOTION\_COMMAND!} \\
\protect\hyperlink{unitree-sdk2-cpp-robot-b2-sportclient-vision-walk-1}{\passthrough{\lstinline!vision\_walk!}}
& \passthrough{\lstinline!def vision\_walk(self, flag: bool) -> int!} &
\passthrough{\lstinline!SIGNATURE\_ONLY!} &
\passthrough{\lstinline!MOTION\_COMMAND!} \\
\protect\hyperlink{unitree-sdk2-cpp-robot-b2-sportclient-hand-stand-1}{\passthrough{\lstinline!hand\_stand!}}
& \passthrough{\lstinline!def hand\_stand(self, flag: bool) -> int!} &
\passthrough{\lstinline!SIGNATURE\_ONLY!} &
\passthrough{\lstinline!MOTION\_COMMAND!} \\
\protect\hyperlink{unitree-sdk2-cpp-robot-b2-sportclient-auto-recovery-set-1}{\passthrough{\lstinline!auto\_recovery\_set!}}
&
\passthrough{\lstinline!def auto\_recovery\_set(self, flag: bool) -> int!}
& \passthrough{\lstinline!SIGNATURE\_ONLY!} &
\passthrough{\lstinline!MOTION\_COMMAND!} \\
\protect\hyperlink{unitree-sdk2-cpp-robot-b2-sportclient-free-walk-1}{\passthrough{\lstinline!free\_walk!}}
& \passthrough{\lstinline!def free\_walk(self) -> int!} &
\passthrough{\lstinline!SIGNATURE\_ONLY!} &
\passthrough{\lstinline!MOTION\_COMMAND!} \\
\protect\hyperlink{unitree-sdk2-cpp-robot-b2-sportclient-classic-walk-1}{\passthrough{\lstinline!classic\_walk!}}
& \passthrough{\lstinline!def classic\_walk(self, flag: bool) -> int!} &
\passthrough{\lstinline!SIGNATURE\_ONLY!} &
\passthrough{\lstinline!MOTION\_COMMAND!} \\
\protect\hyperlink{unitree-sdk2-cpp-robot-b2-sportclient-fast-walk-1}{\passthrough{\lstinline!fast\_walk!}}
& \passthrough{\lstinline!def fast\_walk(self, flag: bool) -> int!} &
\passthrough{\lstinline!SIGNATURE\_ONLY!} &
\passthrough{\lstinline!MOTION\_COMMAND!} \\
\protect\hyperlink{unitree-sdk2-cpp-robot-b2-sportclient-euler-1}{\passthrough{\lstinline!euler!}}
&
\passthrough{\lstinline!def euler(self, roll: float, pitch: float, yaw: float) -> int!}
& \passthrough{\lstinline!SIGNATURE\_ONLY!} &
\passthrough{\lstinline!MOTION\_COMMAND!} \\
\protect\hyperlink{unitree-sdk2-cpp-robot-b2-sportclient-free-height-1}{\passthrough{\lstinline!free\_height!}}
& \passthrough{\lstinline!def free\_height(self, flag: bool) -> int!} &
\passthrough{\lstinline!SIGNATURE\_ONLY!} &
\passthrough{\lstinline!MOTION\_COMMAND!} \\
\protect\hyperlink{unitree-sdk2-cpp-robot-b2-sportclient-gait-height-1}{\passthrough{\lstinline!gait\_height!}}
& \passthrough{\lstinline!def gait\_height(self, flag: bool) -> int!} &
\passthrough{\lstinline!SIGNATURE\_ONLY!} &
\passthrough{\lstinline!MOTION\_COMMAND!} \\
\bottomrule
\end{longtable}

\hypertarget{unitree-sdk2-cpp-robot-b2-sportclient-dunder-init-1}{%
\paragraph{\texorpdfstring{\texttt{unitree\_sdk2\_cpp.robot.b2.SportClient.\_\_init\_\_}}{unitree\_sdk2\_cpp.robot.b2.SportClient.\_\_init\_\_}}\label{unitree-sdk2-cpp-robot-b2-sportclient-dunder-init-1}}

初始化 \passthrough{\lstinline!SportClient!}
实例。是否能实际构造取决于下方可用性状态。

\textbf{签名}

\begin{lstlisting}[language=Python]
def __init__(self, enable_lease: bool = ...) -> None
\end{lstlisting}

\textbf{可用性与安全性}

\textbf{\passthrough{\lstinline!SIGNATURE\_ONLY!} /
\passthrough{\lstinline!CONSTRUCTION!}}：当前只有设计期签名，不能假设已存在于运行时扩展。

\textbf{参数}

\begin{longtable}[]{@{}
  >{\raggedright\arraybackslash}p{(\columnwidth - 6\tabcolsep) * \real{0.18}}
  >{\raggedright\arraybackslash}p{(\columnwidth - 6\tabcolsep) * \real{0.24}}
  >{\raggedright\arraybackslash}p{(\columnwidth - 6\tabcolsep) * \real{0.18}}
  >{\raggedright\arraybackslash}p{(\columnwidth - 6\tabcolsep) * \real{0.40}}@{}}
\toprule
名称 & Python 类型 & 默认值 & 含义 \\
\midrule
\endhead
\passthrough{\lstinline!enable\_lease!} & \passthrough{\lstinline!bool!}
& \passthrough{\lstinline!...!} & 是否启用 SDK lease 机制；lease
的获得、续期和释放规则以服务协议为准。 对应 C++ 参数
\passthrough{\lstinline!enableLease: bool!}。 只接受布尔语义。 \\
\bottomrule
\end{longtable}

\textbf{返回值}

\begin{longtable}[]{@{}
  >{\raggedright\arraybackslash}p{(\columnwidth - 4\tabcolsep) * \real{0.18}}
  >{\raggedright\arraybackslash}p{(\columnwidth - 4\tabcolsep) * \real{0.20}}
  >{\raggedright\arraybackslash}p{(\columnwidth - 4\tabcolsep) * \real{0.62}}@{}}
\toprule
位置 & 类型 & 含义 \\
\midrule
\endhead
无 & \passthrough{\lstinline!None!} &
\passthrough{\lstinline!\_\_init\_\_!}
本身不返回值；调用类对象时，在构造函数可用的前提下得到该类实例。 \\
\bottomrule
\end{longtable}

\textbf{对应 C++}

\begin{itemize}
\tightlist
\item
  类：\passthrough{\lstinline!unitree::robot::b2::SportClient!}
\item
  签名：\passthrough{\lstinline!SportClient(bool)!}
\item
  绑定策略：\passthrough{\lstinline!未单独标注!}
\item
  声明位置：\passthrough{\lstinline!include/unitree/robot/b2/sport/sport\_client.hpp:34!}
\end{itemize}

\textbf{说明}

\begin{itemize}
\tightlist
\item
  示例是局部调用片段；\passthrough{\lstinline!obj!}
  和参数变量表示已经按上层业务校验并准备好的对象。
\item
  默认值显示为 \passthrough{\lstinline!...!}，表示 C++
  声明存在默认参数，但当前 AST
  清单没有保存其字面量；省略参数可使用上游默认行为。
\item
  该条目可以被编辑器和类型检查器识别，但当前运行时可能在属性查找阶段直接失败。
\end{itemize}

\textbf{用法}

\begin{lstlisting}[language=Python]
from typing import TYPE_CHECKING

if TYPE_CHECKING:
    def planned_construct(enable_lease: bool) -> SportClient:
        return SportClient(enable_lease=enable_lease)
\end{lstlisting}

\begin{quote}
{[}!CAUTION{]} 上面的 \passthrough{\lstinline!TYPE\_CHECKING!}
示例只表达计划调用形态。该分支在普通 Python
运行时为假，不代表当前扩展能执行此接口。
\end{quote}

\hypertarget{unitree-sdk2-cpp-robot-b2-sportclient-init-1}{%
\paragraph{\texorpdfstring{\texttt{unitree\_sdk2\_cpp.robot.b2.SportClient.init}}{unitree\_sdk2\_cpp.robot.b2.SportClient.init}}\label{unitree-sdk2-cpp-robot-b2-sportclient-init-1}}

初始化当前 SDK 对象所需的底层通道或服务资源。

\textbf{签名}

\begin{lstlisting}[language=Python]
def init(self) -> None
\end{lstlisting}

\textbf{可用性与安全性}

\textbf{\passthrough{\lstinline!SIGNATURE\_ONLY!} /
\passthrough{\lstinline!INITIALIZATION!}}：当前只有设计期签名，不能假设已存在于运行时扩展。

\textbf{参数}

无显式参数。实例方法中的 \passthrough{\lstinline!self!} 由 Python
自动传入。

\textbf{返回值}

\begin{longtable}[]{@{}
  >{\raggedright\arraybackslash}p{(\columnwidth - 4\tabcolsep) * \real{0.18}}
  >{\raggedright\arraybackslash}p{(\columnwidth - 4\tabcolsep) * \real{0.20}}
  >{\raggedright\arraybackslash}p{(\columnwidth - 4\tabcolsep) * \real{0.62}}@{}}
\toprule
位置 & 类型 & 含义 \\
\midrule
\endhead
无 & \passthrough{\lstinline!None!} & 该方法不返回 Python
值；失败可能表现为异常或底层状态变化。 \\
\bottomrule
\end{longtable}

\textbf{对应 C++}

\begin{itemize}
\tightlist
\item
  类：\passthrough{\lstinline!unitree::robot::b2::SportClient!}
\item
  签名：\passthrough{\lstinline!Init()!}
\item
  绑定策略：\passthrough{\lstinline!DIRECT!}
\item
  声明位置：\passthrough{\lstinline!include/unitree/robot/b2/sport/sport\_client.hpp:37!}
\end{itemize}

\textbf{说明}

\begin{itemize}
\tightlist
\item
  示例是局部调用片段；\passthrough{\lstinline!obj!}
  和参数变量表示已经按上层业务校验并准备好的对象。
\item
  该条目可以被编辑器和类型检查器识别，但当前运行时可能在属性查找阶段直接失败。
\end{itemize}

\textbf{用法}

\begin{lstlisting}[language=Python]
from typing import TYPE_CHECKING

if TYPE_CHECKING:
    def planned_init(obj: SportClient) -> None:
        obj.init()
\end{lstlisting}

\begin{quote}
{[}!CAUTION{]} 上面的 \passthrough{\lstinline!TYPE\_CHECKING!}
示例只表达计划调用形态。该分支在普通 Python
运行时为假，不代表当前扩展能执行此接口。
\end{quote}

\hypertarget{unitree-sdk2-cpp-robot-b2-sportclient-damp-1}{%
\paragraph{\texorpdfstring{\texttt{unitree\_sdk2\_cpp.robot.b2.SportClient.damp}}{unitree\_sdk2\_cpp.robot.b2.SportClient.damp}}\label{unitree-sdk2-cpp-robot-b2-sportclient-damp-1}}

对应 C++ SDK 操作
\passthrough{\lstinline!Damp()!}。上游头文件没有可直接生成的业务说明时，本参考只保证签名映射，精确语义需查目标型号协议。

\textbf{签名}

\begin{lstlisting}[language=Python]
def damp(self) -> int
\end{lstlisting}

\textbf{可用性与安全性}

\textbf{\passthrough{\lstinline!SIGNATURE\_ONLY!} /
\passthrough{\lstinline!MOTION\_COMMAND!}}：仅用于补全和静态检查。当前二进制没有该
Python 方法；不得按可执行运动接口使用。

\textbf{参数}

无显式参数。实例方法中的 \passthrough{\lstinline!self!} 由 Python
自动传入。

\textbf{返回值}

\begin{longtable}[]{@{}
  >{\raggedright\arraybackslash}p{(\columnwidth - 4\tabcolsep) * \real{0.18}}
  >{\raggedright\arraybackslash}p{(\columnwidth - 4\tabcolsep) * \real{0.20}}
  >{\raggedright\arraybackslash}p{(\columnwidth - 4\tabcolsep) * \real{0.62}}@{}}
\toprule
位置 & 类型 & 含义 \\
\midrule
\endhead
返回值 & \passthrough{\lstinline!int!} & SDK 状态码；Unitree 示例通常以
\passthrough{\lstinline!0!} 表示成功，非零值需按具体服务定义解释。 \\
\bottomrule
\end{longtable}

\textbf{对应 C++}

\begin{itemize}
\tightlist
\item
  类：\passthrough{\lstinline!unitree::robot::b2::SportClient!}
\item
  签名：\passthrough{\lstinline!Damp()!}
\item
  绑定策略：\passthrough{\lstinline!DIRECT!}
\item
  声明位置：\passthrough{\lstinline!include/unitree/robot/b2/sport/sport\_client.hpp:39!}
\end{itemize}

\textbf{说明}

\begin{itemize}
\tightlist
\item
  示例是局部调用片段；\passthrough{\lstinline!obj!}
  和参数变量表示已经按上层业务校验并准备好的对象。
\item
  该条目可以被编辑器和类型检查器识别，但当前运行时可能在属性查找阶段直接失败。
\item
  该方法属于运动命令。方法名包含
  \passthrough{\lstinline!stop!}、\passthrough{\lstinline!damp!} 或
  \passthrough{\lstinline!zero!} 也不代表它是独立物理急停。
\end{itemize}

\textbf{用法}

\begin{lstlisting}[language=Python]
from typing import TYPE_CHECKING

if TYPE_CHECKING:
    def planned_damp(obj: SportClient) -> int:
        return obj.damp()
\end{lstlisting}

\begin{quote}
{[}!CAUTION{]} 上面的 \passthrough{\lstinline!TYPE\_CHECKING!}
示例只表达计划调用形态。该分支在普通 Python
运行时为假，不代表当前扩展能执行此接口。
\end{quote}

\hypertarget{unitree-sdk2-cpp-robot-b2-sportclient-balance-stand-1}{%
\paragraph{\texorpdfstring{\texttt{unitree\_sdk2\_cpp.robot.b2.SportClient.balance\_stand}}{unitree\_sdk2\_cpp.robot.b2.SportClient.balance\_stand}}\label{unitree-sdk2-cpp-robot-b2-sportclient-balance-stand-1}}

对应 C++ SDK 操作
\passthrough{\lstinline!BalanceStand()!}。上游头文件没有可直接生成的业务说明时，本参考只保证签名映射，精确语义需查目标型号协议。

\textbf{签名}

\begin{lstlisting}[language=Python]
def balance_stand(self) -> int
\end{lstlisting}

\textbf{可用性与安全性}

\textbf{\passthrough{\lstinline!SIGNATURE\_ONLY!} /
\passthrough{\lstinline!MOTION\_COMMAND!}}：仅用于补全和静态检查。当前二进制没有该
Python 方法；不得按可执行运动接口使用。

\textbf{参数}

无显式参数。实例方法中的 \passthrough{\lstinline!self!} 由 Python
自动传入。

\textbf{返回值}

\begin{longtable}[]{@{}
  >{\raggedright\arraybackslash}p{(\columnwidth - 4\tabcolsep) * \real{0.18}}
  >{\raggedright\arraybackslash}p{(\columnwidth - 4\tabcolsep) * \real{0.20}}
  >{\raggedright\arraybackslash}p{(\columnwidth - 4\tabcolsep) * \real{0.62}}@{}}
\toprule
位置 & 类型 & 含义 \\
\midrule
\endhead
返回值 & \passthrough{\lstinline!int!} & SDK 状态码；Unitree 示例通常以
\passthrough{\lstinline!0!} 表示成功，非零值需按具体服务定义解释。 \\
\bottomrule
\end{longtable}

\textbf{对应 C++}

\begin{itemize}
\tightlist
\item
  类：\passthrough{\lstinline!unitree::robot::b2::SportClient!}
\item
  签名：\passthrough{\lstinline!BalanceStand()!}
\item
  绑定策略：\passthrough{\lstinline!DIRECT!}
\item
  声明位置：\passthrough{\lstinline!include/unitree/robot/b2/sport/sport\_client.hpp:41!}
\end{itemize}

\textbf{说明}

\begin{itemize}
\tightlist
\item
  示例是局部调用片段；\passthrough{\lstinline!obj!}
  和参数变量表示已经按上层业务校验并准备好的对象。
\item
  该条目可以被编辑器和类型检查器识别，但当前运行时可能在属性查找阶段直接失败。
\item
  该方法属于运动命令。方法名包含
  \passthrough{\lstinline!stop!}、\passthrough{\lstinline!damp!} 或
  \passthrough{\lstinline!zero!} 也不代表它是独立物理急停。
\end{itemize}

\textbf{用法}

\begin{lstlisting}[language=Python]
from typing import TYPE_CHECKING

if TYPE_CHECKING:
    def planned_balance_stand(obj: SportClient) -> int:
        return obj.balance_stand()
\end{lstlisting}

\begin{quote}
{[}!CAUTION{]} 上面的 \passthrough{\lstinline!TYPE\_CHECKING!}
示例只表达计划调用形态。该分支在普通 Python
运行时为假，不代表当前扩展能执行此接口。
\end{quote}

\hypertarget{unitree-sdk2-cpp-robot-b2-sportclient-stop-move-1}{%
\paragraph{\texorpdfstring{\texttt{unitree\_sdk2\_cpp.robot.b2.SportClient.stop\_move}}{unitree\_sdk2\_cpp.robot.b2.SportClient.stop\_move}}\label{unitree-sdk2-cpp-robot-b2-sportclient-stop-move-1}}

请求停止对应 SDK 操作。该名称不等同于经过验证的物理急停。

\textbf{签名}

\begin{lstlisting}[language=Python]
def stop_move(self) -> int
\end{lstlisting}

\textbf{可用性与安全性}

\textbf{\passthrough{\lstinline!SIGNATURE\_ONLY!} /
\passthrough{\lstinline!MOTION\_COMMAND!}}：仅用于补全和静态检查。当前二进制没有该
Python 方法；不得按可执行运动接口使用。

\textbf{参数}

无显式参数。实例方法中的 \passthrough{\lstinline!self!} 由 Python
自动传入。

\textbf{返回值}

\begin{longtable}[]{@{}
  >{\raggedright\arraybackslash}p{(\columnwidth - 4\tabcolsep) * \real{0.18}}
  >{\raggedright\arraybackslash}p{(\columnwidth - 4\tabcolsep) * \real{0.20}}
  >{\raggedright\arraybackslash}p{(\columnwidth - 4\tabcolsep) * \real{0.62}}@{}}
\toprule
位置 & 类型 & 含义 \\
\midrule
\endhead
返回值 & \passthrough{\lstinline!int!} & SDK 状态码；Unitree 示例通常以
\passthrough{\lstinline!0!} 表示成功，非零值需按具体服务定义解释。 \\
\bottomrule
\end{longtable}

\textbf{对应 C++}

\begin{itemize}
\tightlist
\item
  类：\passthrough{\lstinline!unitree::robot::b2::SportClient!}
\item
  签名：\passthrough{\lstinline!StopMove()!}
\item
  绑定策略：\passthrough{\lstinline!DIRECT!}
\item
  声明位置：\passthrough{\lstinline!include/unitree/robot/b2/sport/sport\_client.hpp:43!}
\end{itemize}

\textbf{说明}

\begin{itemize}
\tightlist
\item
  示例是局部调用片段；\passthrough{\lstinline!obj!}
  和参数变量表示已经按上层业务校验并准备好的对象。
\item
  该条目可以被编辑器和类型检查器识别，但当前运行时可能在属性查找阶段直接失败。
\item
  该方法属于运动命令。方法名包含
  \passthrough{\lstinline!stop!}、\passthrough{\lstinline!damp!} 或
  \passthrough{\lstinline!zero!} 也不代表它是独立物理急停。
\end{itemize}

\textbf{用法}

\begin{lstlisting}[language=Python]
from typing import TYPE_CHECKING

if TYPE_CHECKING:
    def planned_stop_move(obj: SportClient) -> int:
        return obj.stop_move()
\end{lstlisting}

\begin{quote}
{[}!CAUTION{]} 上面的 \passthrough{\lstinline!TYPE\_CHECKING!}
示例只表达计划调用形态。该分支在普通 Python
运行时为假，不代表当前扩展能执行此接口。
\end{quote}

\hypertarget{unitree-sdk2-cpp-robot-b2-sportclient-stand-up-1}{%
\paragraph{\texorpdfstring{\texttt{unitree\_sdk2\_cpp.robot.b2.SportClient.stand\_up}}{unitree\_sdk2\_cpp.robot.b2.SportClient.stand\_up}}\label{unitree-sdk2-cpp-robot-b2-sportclient-stand-up-1}}

对应 C++ SDK 操作
\passthrough{\lstinline!StandUp()!}。上游头文件没有可直接生成的业务说明时，本参考只保证签名映射，精确语义需查目标型号协议。

\textbf{签名}

\begin{lstlisting}[language=Python]
def stand_up(self) -> int
\end{lstlisting}

\textbf{可用性与安全性}

\textbf{\passthrough{\lstinline!SIGNATURE\_ONLY!} /
\passthrough{\lstinline!MOTION\_COMMAND!}}：仅用于补全和静态检查。当前二进制没有该
Python 方法；不得按可执行运动接口使用。

\textbf{参数}

无显式参数。实例方法中的 \passthrough{\lstinline!self!} 由 Python
自动传入。

\textbf{返回值}

\begin{longtable}[]{@{}
  >{\raggedright\arraybackslash}p{(\columnwidth - 4\tabcolsep) * \real{0.18}}
  >{\raggedright\arraybackslash}p{(\columnwidth - 4\tabcolsep) * \real{0.20}}
  >{\raggedright\arraybackslash}p{(\columnwidth - 4\tabcolsep) * \real{0.62}}@{}}
\toprule
位置 & 类型 & 含义 \\
\midrule
\endhead
返回值 & \passthrough{\lstinline!int!} & SDK 状态码；Unitree 示例通常以
\passthrough{\lstinline!0!} 表示成功，非零值需按具体服务定义解释。 \\
\bottomrule
\end{longtable}

\textbf{对应 C++}

\begin{itemize}
\tightlist
\item
  类：\passthrough{\lstinline!unitree::robot::b2::SportClient!}
\item
  签名：\passthrough{\lstinline!StandUp()!}
\item
  绑定策略：\passthrough{\lstinline!DIRECT!}
\item
  声明位置：\passthrough{\lstinline!include/unitree/robot/b2/sport/sport\_client.hpp:45!}
\end{itemize}

\textbf{说明}

\begin{itemize}
\tightlist
\item
  示例是局部调用片段；\passthrough{\lstinline!obj!}
  和参数变量表示已经按上层业务校验并准备好的对象。
\item
  该条目可以被编辑器和类型检查器识别，但当前运行时可能在属性查找阶段直接失败。
\item
  该方法属于运动命令。方法名包含
  \passthrough{\lstinline!stop!}、\passthrough{\lstinline!damp!} 或
  \passthrough{\lstinline!zero!} 也不代表它是独立物理急停。
\end{itemize}

\textbf{用法}

\begin{lstlisting}[language=Python]
from typing import TYPE_CHECKING

if TYPE_CHECKING:
    def planned_stand_up(obj: SportClient) -> int:
        return obj.stand_up()
\end{lstlisting}

\begin{quote}
{[}!CAUTION{]} 上面的 \passthrough{\lstinline!TYPE\_CHECKING!}
示例只表达计划调用形态。该分支在普通 Python
运行时为假，不代表当前扩展能执行此接口。
\end{quote}

\hypertarget{unitree-sdk2-cpp-robot-b2-sportclient-stand-down-1}{%
\paragraph{\texorpdfstring{\texttt{unitree\_sdk2\_cpp.robot.b2.SportClient.stand\_down}}{unitree\_sdk2\_cpp.robot.b2.SportClient.stand\_down}}\label{unitree-sdk2-cpp-robot-b2-sportclient-stand-down-1}}

对应 C++ SDK 操作
\passthrough{\lstinline!StandDown()!}。上游头文件没有可直接生成的业务说明时，本参考只保证签名映射，精确语义需查目标型号协议。

\textbf{签名}

\begin{lstlisting}[language=Python]
def stand_down(self) -> int
\end{lstlisting}

\textbf{可用性与安全性}

\textbf{\passthrough{\lstinline!SIGNATURE\_ONLY!} /
\passthrough{\lstinline!MOTION\_COMMAND!}}：仅用于补全和静态检查。当前二进制没有该
Python 方法；不得按可执行运动接口使用。

\textbf{参数}

无显式参数。实例方法中的 \passthrough{\lstinline!self!} 由 Python
自动传入。

\textbf{返回值}

\begin{longtable}[]{@{}
  >{\raggedright\arraybackslash}p{(\columnwidth - 4\tabcolsep) * \real{0.18}}
  >{\raggedright\arraybackslash}p{(\columnwidth - 4\tabcolsep) * \real{0.20}}
  >{\raggedright\arraybackslash}p{(\columnwidth - 4\tabcolsep) * \real{0.62}}@{}}
\toprule
位置 & 类型 & 含义 \\
\midrule
\endhead
返回值 & \passthrough{\lstinline!int!} & SDK 状态码；Unitree 示例通常以
\passthrough{\lstinline!0!} 表示成功，非零值需按具体服务定义解释。 \\
\bottomrule
\end{longtable}

\textbf{对应 C++}

\begin{itemize}
\tightlist
\item
  类：\passthrough{\lstinline!unitree::robot::b2::SportClient!}
\item
  签名：\passthrough{\lstinline!StandDown()!}
\item
  绑定策略：\passthrough{\lstinline!DIRECT!}
\item
  声明位置：\passthrough{\lstinline!include/unitree/robot/b2/sport/sport\_client.hpp:47!}
\end{itemize}

\textbf{说明}

\begin{itemize}
\tightlist
\item
  示例是局部调用片段；\passthrough{\lstinline!obj!}
  和参数变量表示已经按上层业务校验并准备好的对象。
\item
  该条目可以被编辑器和类型检查器识别，但当前运行时可能在属性查找阶段直接失败。
\item
  该方法属于运动命令。方法名包含
  \passthrough{\lstinline!stop!}、\passthrough{\lstinline!damp!} 或
  \passthrough{\lstinline!zero!} 也不代表它是独立物理急停。
\end{itemize}

\textbf{用法}

\begin{lstlisting}[language=Python]
from typing import TYPE_CHECKING

if TYPE_CHECKING:
    def planned_stand_down(obj: SportClient) -> int:
        return obj.stand_down()
\end{lstlisting}

\begin{quote}
{[}!CAUTION{]} 上面的 \passthrough{\lstinline!TYPE\_CHECKING!}
示例只表达计划调用形态。该分支在普通 Python
运行时为假，不代表当前扩展能执行此接口。
\end{quote}

\hypertarget{unitree-sdk2-cpp-robot-b2-sportclient-recovery-stand-1}{%
\paragraph{\texorpdfstring{\texttt{unitree\_sdk2\_cpp.robot.b2.SportClient.recovery\_stand}}{unitree\_sdk2\_cpp.robot.b2.SportClient.recovery\_stand}}\label{unitree-sdk2-cpp-robot-b2-sportclient-recovery-stand-1}}

对应 C++ SDK 操作
\passthrough{\lstinline!RecoveryStand()!}。上游头文件没有可直接生成的业务说明时，本参考只保证签名映射，精确语义需查目标型号协议。

\textbf{签名}

\begin{lstlisting}[language=Python]
def recovery_stand(self) -> int
\end{lstlisting}

\textbf{可用性与安全性}

\textbf{\passthrough{\lstinline!SIGNATURE\_ONLY!} /
\passthrough{\lstinline!MOTION\_COMMAND!}}：仅用于补全和静态检查。当前二进制没有该
Python 方法；不得按可执行运动接口使用。

\textbf{参数}

无显式参数。实例方法中的 \passthrough{\lstinline!self!} 由 Python
自动传入。

\textbf{返回值}

\begin{longtable}[]{@{}
  >{\raggedright\arraybackslash}p{(\columnwidth - 4\tabcolsep) * \real{0.18}}
  >{\raggedright\arraybackslash}p{(\columnwidth - 4\tabcolsep) * \real{0.20}}
  >{\raggedright\arraybackslash}p{(\columnwidth - 4\tabcolsep) * \real{0.62}}@{}}
\toprule
位置 & 类型 & 含义 \\
\midrule
\endhead
返回值 & \passthrough{\lstinline!int!} & SDK 状态码；Unitree 示例通常以
\passthrough{\lstinline!0!} 表示成功，非零值需按具体服务定义解释。 \\
\bottomrule
\end{longtable}

\textbf{对应 C++}

\begin{itemize}
\tightlist
\item
  类：\passthrough{\lstinline!unitree::robot::b2::SportClient!}
\item
  签名：\passthrough{\lstinline!RecoveryStand()!}
\item
  绑定策略：\passthrough{\lstinline!DIRECT!}
\item
  声明位置：\passthrough{\lstinline!include/unitree/robot/b2/sport/sport\_client.hpp:49!}
\end{itemize}

\textbf{说明}

\begin{itemize}
\tightlist
\item
  示例是局部调用片段；\passthrough{\lstinline!obj!}
  和参数变量表示已经按上层业务校验并准备好的对象。
\item
  该条目可以被编辑器和类型检查器识别，但当前运行时可能在属性查找阶段直接失败。
\item
  该方法属于运动命令。方法名包含
  \passthrough{\lstinline!stop!}、\passthrough{\lstinline!damp!} 或
  \passthrough{\lstinline!zero!} 也不代表它是独立物理急停。
\end{itemize}

\textbf{用法}

\begin{lstlisting}[language=Python]
from typing import TYPE_CHECKING

if TYPE_CHECKING:
    def planned_recovery_stand(obj: SportClient) -> int:
        return obj.recovery_stand()
\end{lstlisting}

\begin{quote}
{[}!CAUTION{]} 上面的 \passthrough{\lstinline!TYPE\_CHECKING!}
示例只表达计划调用形态。该分支在普通 Python
运行时为假，不代表当前扩展能执行此接口。
\end{quote}

\hypertarget{unitree-sdk2-cpp-robot-b2-sportclient-move-1}{%
\paragraph{\texorpdfstring{\texttt{unitree\_sdk2\_cpp.robot.b2.SportClient.move}}{unitree\_sdk2\_cpp.robot.b2.SportClient.move}}\label{unitree-sdk2-cpp-robot-b2-sportclient-move-1}}

对应 C++ SDK 操作
\passthrough{\lstinline!Move(float, float, float)!}。上游头文件没有可直接生成的业务说明时，本参考只保证签名映射，精确语义需查目标型号协议。

\textbf{签名}

\begin{lstlisting}[language=Python]
def move(self, vx: float, vy: float, vyaw: float) -> int
\end{lstlisting}

\textbf{可用性与安全性}

\textbf{\passthrough{\lstinline!SIGNATURE\_ONLY!} /
\passthrough{\lstinline!MOTION\_COMMAND!}}：仅用于补全和静态检查。当前二进制没有该
Python 方法；不得按可执行运动接口使用。

\textbf{参数}

\begin{longtable}[]{@{}
  >{\raggedright\arraybackslash}p{(\columnwidth - 6\tabcolsep) * \real{0.18}}
  >{\raggedright\arraybackslash}p{(\columnwidth - 6\tabcolsep) * \real{0.24}}
  >{\raggedright\arraybackslash}p{(\columnwidth - 6\tabcolsep) * \real{0.18}}
  >{\raggedright\arraybackslash}p{(\columnwidth - 6\tabcolsep) * \real{0.40}}@{}}
\toprule
名称 & Python 类型 & 默认值 & 含义 \\
\midrule
\endhead
\passthrough{\lstinline!vx!} & \passthrough{\lstinline!float!} & 必填 &
X 方向速度参数。坐标系、单位、符号和安全范围必须查目标型号运动协议。
对应 C++ 参数 \passthrough{\lstinline!vx: float!}。
底层为浮点数；类型本身不说明单位、坐标系或业务安全范围。 \\
\passthrough{\lstinline!vy!} & \passthrough{\lstinline!float!} & 必填 &
Y 方向速度参数。坐标系、单位、符号和安全范围必须查目标型号运动协议。
对应 C++ 参数 \passthrough{\lstinline!vy: float!}。
底层为浮点数；类型本身不说明单位、坐标系或业务安全范围。 \\
\passthrough{\lstinline!vyaw!} & \passthrough{\lstinline!float!} & 必填
& 偏航角速度参数。单位、符号和安全范围必须查目标型号运动协议。 对应 C++
参数 \passthrough{\lstinline!vyaw: float!}。
底层为浮点数；类型本身不说明单位、坐标系或业务安全范围。 \\
\bottomrule
\end{longtable}

\textbf{返回值}

\begin{longtable}[]{@{}
  >{\raggedright\arraybackslash}p{(\columnwidth - 4\tabcolsep) * \real{0.18}}
  >{\raggedright\arraybackslash}p{(\columnwidth - 4\tabcolsep) * \real{0.20}}
  >{\raggedright\arraybackslash}p{(\columnwidth - 4\tabcolsep) * \real{0.62}}@{}}
\toprule
位置 & 类型 & 含义 \\
\midrule
\endhead
返回值 & \passthrough{\lstinline!int!} & SDK 状态码；Unitree 示例通常以
\passthrough{\lstinline!0!} 表示成功，非零值需按具体服务定义解释。 \\
\bottomrule
\end{longtable}

\textbf{对应 C++}

\begin{itemize}
\tightlist
\item
  类：\passthrough{\lstinline!unitree::robot::b2::SportClient!}
\item
  签名：\passthrough{\lstinline!Move(float, float, float)!}
\item
  绑定策略：\passthrough{\lstinline!DIRECT!}
\item
  声明位置：\passthrough{\lstinline!include/unitree/robot/b2/sport/sport\_client.hpp:51!}
\end{itemize}

\textbf{说明}

\begin{itemize}
\tightlist
\item
  示例是局部调用片段；\passthrough{\lstinline!obj!}
  和参数变量表示已经按上层业务校验并准备好的对象。
\item
  该条目可以被编辑器和类型检查器识别，但当前运行时可能在属性查找阶段直接失败。
\item
  该方法属于运动命令。方法名包含
  \passthrough{\lstinline!stop!}、\passthrough{\lstinline!damp!} 或
  \passthrough{\lstinline!zero!} 也不代表它是独立物理急停。
\end{itemize}

\textbf{用法}

\begin{lstlisting}[language=Python]
from typing import TYPE_CHECKING

if TYPE_CHECKING:
    def planned_move(obj: SportClient, vx: float, vy: float, vyaw: float) -> int:
        return obj.move(vx=vx, vy=vy, vyaw=vyaw)
\end{lstlisting}

\begin{quote}
{[}!CAUTION{]} 上面的 \passthrough{\lstinline!TYPE\_CHECKING!}
示例只表达计划调用形态。该分支在普通 Python
运行时为假，不代表当前扩展能执行此接口。
\end{quote}

\hypertarget{unitree-sdk2-cpp-robot-b2-sportclient-switch-gait-1}{%
\paragraph{\texorpdfstring{\texttt{unitree\_sdk2\_cpp.robot.b2.SportClient.switch\_gait}}{unitree\_sdk2\_cpp.robot.b2.SportClient.switch\_gait}}\label{unitree-sdk2-cpp-robot-b2-sportclient-switch-gait-1}}

对应 C++ SDK 操作
\passthrough{\lstinline!SwitchGait(int)!}。上游头文件没有可直接生成的业务说明时，本参考只保证签名映射，精确语义需查目标型号协议。

\textbf{签名}

\begin{lstlisting}[language=Python]
def switch_gait(self, d: int) -> int
\end{lstlisting}

\textbf{可用性与安全性}

\textbf{\passthrough{\lstinline!SIGNATURE\_ONLY!} /
\passthrough{\lstinline!MOTION\_COMMAND!}}：仅用于补全和静态检查。当前二进制没有该
Python 方法；不得按可执行运动接口使用。

\textbf{参数}

\begin{longtable}[]{@{}
  >{\raggedright\arraybackslash}p{(\columnwidth - 6\tabcolsep) * \real{0.18}}
  >{\raggedright\arraybackslash}p{(\columnwidth - 6\tabcolsep) * \real{0.24}}
  >{\raggedright\arraybackslash}p{(\columnwidth - 6\tabcolsep) * \real{0.18}}
  >{\raggedright\arraybackslash}p{(\columnwidth - 6\tabcolsep) * \real{0.40}}@{}}
\toprule
名称 & Python 类型 & 默认值 & 含义 \\
\midrule
\endhead
\passthrough{\lstinline!d!} & \passthrough{\lstinline!int!} & 必填 &
传给该接口的 \passthrough{\lstinline!d!} 参数，Python 类型为
\passthrough{\lstinline!int!}。精确含义、单位、范围和枚举值需查目标型号对应头文件与协议。
对应 C++ 参数 \passthrough{\lstinline!d: int!}。 \\
\bottomrule
\end{longtable}

\textbf{返回值}

\begin{longtable}[]{@{}
  >{\raggedright\arraybackslash}p{(\columnwidth - 4\tabcolsep) * \real{0.18}}
  >{\raggedright\arraybackslash}p{(\columnwidth - 4\tabcolsep) * \real{0.20}}
  >{\raggedright\arraybackslash}p{(\columnwidth - 4\tabcolsep) * \real{0.62}}@{}}
\toprule
位置 & 类型 & 含义 \\
\midrule
\endhead
返回值 & \passthrough{\lstinline!int!} & SDK 状态码；Unitree 示例通常以
\passthrough{\lstinline!0!} 表示成功，非零值需按具体服务定义解释。 \\
\bottomrule
\end{longtable}

\textbf{对应 C++}

\begin{itemize}
\tightlist
\item
  类：\passthrough{\lstinline!unitree::robot::b2::SportClient!}
\item
  签名：\passthrough{\lstinline!SwitchGait(int)!}
\item
  绑定策略：\passthrough{\lstinline!DIRECT!}
\item
  声明位置：\passthrough{\lstinline!include/unitree/robot/b2/sport/sport\_client.hpp:53!}
\end{itemize}

\textbf{说明}

\begin{itemize}
\tightlist
\item
  示例是局部调用片段；\passthrough{\lstinline!obj!}
  和参数变量表示已经按上层业务校验并准备好的对象。
\item
  该条目可以被编辑器和类型检查器识别，但当前运行时可能在属性查找阶段直接失败。
\item
  该方法属于运动命令。方法名包含
  \passthrough{\lstinline!stop!}、\passthrough{\lstinline!damp!} 或
  \passthrough{\lstinline!zero!} 也不代表它是独立物理急停。
\end{itemize}

\textbf{用法}

\begin{lstlisting}[language=Python]
from typing import TYPE_CHECKING

if TYPE_CHECKING:
    def planned_switch_gait(obj: SportClient, d: int) -> int:
        return obj.switch_gait(d=d)
\end{lstlisting}

\begin{quote}
{[}!CAUTION{]} 上面的 \passthrough{\lstinline!TYPE\_CHECKING!}
示例只表达计划调用形态。该分支在普通 Python
运行时为假，不代表当前扩展能执行此接口。
\end{quote}

\hypertarget{unitree-sdk2-cpp-robot-b2-sportclient-body-height-1}{%
\paragraph{\texorpdfstring{\texttt{unitree\_sdk2\_cpp.robot.b2.SportClient.body\_height}}{unitree\_sdk2\_cpp.robot.b2.SportClient.body\_height}}\label{unitree-sdk2-cpp-robot-b2-sportclient-body-height-1}}

对应 C++ SDK 操作
\passthrough{\lstinline!BodyHeight(float)!}。上游头文件没有可直接生成的业务说明时，本参考只保证签名映射，精确语义需查目标型号协议。

\textbf{签名}

\begin{lstlisting}[language=Python]
def body_height(self, height: float) -> int
\end{lstlisting}

\textbf{可用性与安全性}

\textbf{\passthrough{\lstinline!SIGNATURE\_ONLY!} /
\passthrough{\lstinline!MOTION\_COMMAND!}}：仅用于补全和静态检查。当前二进制没有该
Python 方法；不得按可执行运动接口使用。

\textbf{参数}

\begin{longtable}[]{@{}
  >{\raggedright\arraybackslash}p{(\columnwidth - 6\tabcolsep) * \real{0.18}}
  >{\raggedright\arraybackslash}p{(\columnwidth - 6\tabcolsep) * \real{0.24}}
  >{\raggedright\arraybackslash}p{(\columnwidth - 6\tabcolsep) * \real{0.18}}
  >{\raggedright\arraybackslash}p{(\columnwidth - 6\tabcolsep) * \real{0.40}}@{}}
\toprule
名称 & Python 类型 & 默认值 & 含义 \\
\midrule
\endhead
\passthrough{\lstinline!height!} & \passthrough{\lstinline!float!} &
必填 & 传给该接口的 高度 参数，Python 类型为
\passthrough{\lstinline!float!}。精确含义、单位、范围和枚举值需查目标型号对应头文件与协议。
对应 C++ 参数 \passthrough{\lstinline!height: float!}。
底层为浮点数；类型本身不说明单位、坐标系或业务安全范围。 \\
\bottomrule
\end{longtable}

\textbf{返回值}

\begin{longtable}[]{@{}
  >{\raggedright\arraybackslash}p{(\columnwidth - 4\tabcolsep) * \real{0.18}}
  >{\raggedright\arraybackslash}p{(\columnwidth - 4\tabcolsep) * \real{0.20}}
  >{\raggedright\arraybackslash}p{(\columnwidth - 4\tabcolsep) * \real{0.62}}@{}}
\toprule
位置 & 类型 & 含义 \\
\midrule
\endhead
返回值 & \passthrough{\lstinline!int!} & SDK 状态码；Unitree 示例通常以
\passthrough{\lstinline!0!} 表示成功，非零值需按具体服务定义解释。 \\
\bottomrule
\end{longtable}

\textbf{对应 C++}

\begin{itemize}
\tightlist
\item
  类：\passthrough{\lstinline!unitree::robot::b2::SportClient!}
\item
  签名：\passthrough{\lstinline!BodyHeight(float)!}
\item
  绑定策略：\passthrough{\lstinline!DIRECT!}
\item
  声明位置：\passthrough{\lstinline!include/unitree/robot/b2/sport/sport\_client.hpp:55!}
\end{itemize}

\textbf{说明}

\begin{itemize}
\tightlist
\item
  示例是局部调用片段；\passthrough{\lstinline!obj!}
  和参数变量表示已经按上层业务校验并准备好的对象。
\item
  该条目可以被编辑器和类型检查器识别，但当前运行时可能在属性查找阶段直接失败。
\item
  该方法属于运动命令。方法名包含
  \passthrough{\lstinline!stop!}、\passthrough{\lstinline!damp!} 或
  \passthrough{\lstinline!zero!} 也不代表它是独立物理急停。
\end{itemize}

\textbf{用法}

\begin{lstlisting}[language=Python]
from typing import TYPE_CHECKING

if TYPE_CHECKING:
    def planned_body_height(obj: SportClient, height: float) -> int:
        return obj.body_height(height=height)
\end{lstlisting}

\begin{quote}
{[}!CAUTION{]} 上面的 \passthrough{\lstinline!TYPE\_CHECKING!}
示例只表达计划调用形态。该分支在普通 Python
运行时为假，不代表当前扩展能执行此接口。
\end{quote}

\hypertarget{unitree-sdk2-cpp-robot-b2-sportclient-speed-level-1}{%
\paragraph{\texorpdfstring{\texttt{unitree\_sdk2\_cpp.robot.b2.SportClient.speed\_level}}{unitree\_sdk2\_cpp.robot.b2.SportClient.speed\_level}}\label{unitree-sdk2-cpp-robot-b2-sportclient-speed-level-1}}

对应 C++ SDK 操作
\passthrough{\lstinline!SpeedLevel(int)!}。上游头文件没有可直接生成的业务说明时，本参考只保证签名映射，精确语义需查目标型号协议。

\textbf{签名}

\begin{lstlisting}[language=Python]
def speed_level(self, level: int) -> int
\end{lstlisting}

\textbf{可用性与安全性}

\textbf{\passthrough{\lstinline!SIGNATURE\_ONLY!} /
\passthrough{\lstinline!MOTION\_COMMAND!}}：仅用于补全和静态检查。当前二进制没有该
Python 方法；不得按可执行运动接口使用。

\textbf{参数}

\begin{longtable}[]{@{}
  >{\raggedright\arraybackslash}p{(\columnwidth - 6\tabcolsep) * \real{0.18}}
  >{\raggedright\arraybackslash}p{(\columnwidth - 6\tabcolsep) * \real{0.24}}
  >{\raggedright\arraybackslash}p{(\columnwidth - 6\tabcolsep) * \real{0.18}}
  >{\raggedright\arraybackslash}p{(\columnwidth - 6\tabcolsep) * \real{0.40}}@{}}
\toprule
名称 & Python 类型 & 默认值 & 含义 \\
\midrule
\endhead
\passthrough{\lstinline!level!} & \passthrough{\lstinline!int!} & 必填 &
传给该接口的 \passthrough{\lstinline!level!} 参数，Python 类型为
\passthrough{\lstinline!int!}。精确含义、单位、范围和枚举值需查目标型号对应头文件与协议。
对应 C++ 参数 \passthrough{\lstinline!level: int!}。 \\
\bottomrule
\end{longtable}

\textbf{返回值}

\begin{longtable}[]{@{}
  >{\raggedright\arraybackslash}p{(\columnwidth - 4\tabcolsep) * \real{0.18}}
  >{\raggedright\arraybackslash}p{(\columnwidth - 4\tabcolsep) * \real{0.20}}
  >{\raggedright\arraybackslash}p{(\columnwidth - 4\tabcolsep) * \real{0.62}}@{}}
\toprule
位置 & 类型 & 含义 \\
\midrule
\endhead
返回值 & \passthrough{\lstinline!int!} & SDK 状态码；Unitree 示例通常以
\passthrough{\lstinline!0!} 表示成功，非零值需按具体服务定义解释。 \\
\bottomrule
\end{longtable}

\textbf{对应 C++}

\begin{itemize}
\tightlist
\item
  类：\passthrough{\lstinline!unitree::robot::b2::SportClient!}
\item
  签名：\passthrough{\lstinline!SpeedLevel(int)!}
\item
  绑定策略：\passthrough{\lstinline!DIRECT!}
\item
  声明位置：\passthrough{\lstinline!include/unitree/robot/b2/sport/sport\_client.hpp:57!}
\end{itemize}

\textbf{说明}

\begin{itemize}
\tightlist
\item
  示例是局部调用片段；\passthrough{\lstinline!obj!}
  和参数变量表示已经按上层业务校验并准备好的对象。
\item
  该条目可以被编辑器和类型检查器识别，但当前运行时可能在属性查找阶段直接失败。
\item
  该方法属于运动命令。方法名包含
  \passthrough{\lstinline!stop!}、\passthrough{\lstinline!damp!} 或
  \passthrough{\lstinline!zero!} 也不代表它是独立物理急停。
\end{itemize}

\textbf{用法}

\begin{lstlisting}[language=Python]
from typing import TYPE_CHECKING

if TYPE_CHECKING:
    def planned_speed_level(obj: SportClient, level: int) -> int:
        return obj.speed_level(level=level)
\end{lstlisting}

\begin{quote}
{[}!CAUTION{]} 上面的 \passthrough{\lstinline!TYPE\_CHECKING!}
示例只表达计划调用形态。该分支在普通 Python
运行时为假，不代表当前扩展能执行此接口。
\end{quote}

\hypertarget{unitree-sdk2-cpp-robot-b2-sportclient-trajectory-follow-1}{%
\paragraph{\texorpdfstring{\texttt{unitree\_sdk2\_cpp.robot.b2.SportClient.trajectory\_follow}}{unitree\_sdk2\_cpp.robot.b2.SportClient.trajectory\_follow}}\label{unitree-sdk2-cpp-robot-b2-sportclient-trajectory-follow-1}}

对应 C++ SDK 操作
\passthrough{\lstinline!TrajectoryFollow(std::vector<unitree::robot::b2::PathPoint> \&)!}。上游头文件没有可直接生成的业务说明时，本参考只保证签名映射，精确语义需查目标型号协议。

\textbf{签名}

\begin{lstlisting}[language=Python]
def trajectory_follow(self) -> tuple[int, list[Any]]
\end{lstlisting}

\textbf{可用性与安全性}

\textbf{\passthrough{\lstinline!SIGNATURE\_ONLY!} /
\passthrough{\lstinline!MOTION\_COMMAND!}}：仅用于补全和静态检查。当前二进制没有该
Python 方法；不得按可执行运动接口使用。

\textbf{参数}

无显式参数。实例方法中的 \passthrough{\lstinline!self!} 由 Python
自动传入。

\textbf{返回值}

\begin{longtable}[]{@{}
  >{\raggedright\arraybackslash}p{(\columnwidth - 4\tabcolsep) * \real{0.18}}
  >{\raggedright\arraybackslash}p{(\columnwidth - 4\tabcolsep) * \real{0.20}}
  >{\raggedright\arraybackslash}p{(\columnwidth - 4\tabcolsep) * \real{0.62}}@{}}
\toprule
位置 & 类型 & 含义 \\
\midrule
\endhead
{[}0{]} \passthrough{\lstinline!status!} & \passthrough{\lstinline!int!}
& SDK 状态码；Unitree 示例通常以 \passthrough{\lstinline!0!}
表示成功，非零值需按具体服务错误码解释。 \\
{[}1{]} \passthrough{\lstinline!path!} &
\passthrough{\lstinline!list[Any]!} & 传给该接口的
\passthrough{\lstinline!path!} 参数，Python 类型为
\passthrough{\lstinline!list[Any]!}。精确含义、单位、范围和枚举值需查目标型号对应头文件与协议。
对应 C++ 参数
\passthrough{\lstinline!path: std::vector<unitree::robot::b2::PathPoint> \&!}。
底层是可变长度 vector \\
\bottomrule
\end{longtable}

\textbf{对应 C++}

\begin{itemize}
\tightlist
\item
  类：\passthrough{\lstinline!unitree::robot::b2::SportClient!}
\item
  签名：\passthrough{\lstinline!TrajectoryFollow(std::vector<unitree::robot::b2::PathPoint> \&)!}
\item
  绑定策略：\passthrough{\lstinline!OUTPUT\_WRAPPER!}
\item
  声明位置：\passthrough{\lstinline!include/unitree/robot/b2/sport/sport\_client.hpp:59!}
\end{itemize}

\textbf{说明}

\begin{itemize}
\tightlist
\item
  示例是局部调用片段；\passthrough{\lstinline!obj!}
  和参数变量表示已经按上层业务校验并准备好的对象。
\item
  C++ 的可修改输出引用不会作为 Python 入参出现，而是按 C++
  参数顺序追加到返回元组中。
\item
  该条目可以被编辑器和类型检查器识别，但当前运行时可能在属性查找阶段直接失败。
\item
  该方法属于运动命令。方法名包含
  \passthrough{\lstinline!stop!}、\passthrough{\lstinline!damp!} 或
  \passthrough{\lstinline!zero!} 也不代表它是独立物理急停。
\end{itemize}

\textbf{用法}

\begin{lstlisting}[language=Python]
from typing import TYPE_CHECKING

if TYPE_CHECKING:
    def planned_trajectory_follow(obj: SportClient) -> tuple[int, list[Any]]:
        return obj.trajectory_follow()
\end{lstlisting}

\begin{quote}
{[}!CAUTION{]} 上面的 \passthrough{\lstinline!TYPE\_CHECKING!}
示例只表达计划调用形态。该分支在普通 Python
运行时为假，不代表当前扩展能执行此接口。
\end{quote}

\hypertarget{unitree-sdk2-cpp-robot-b2-sportclient-continuous-gait-1}{%
\paragraph{\texorpdfstring{\texttt{unitree\_sdk2\_cpp.robot.b2.SportClient.continuous\_gait}}{unitree\_sdk2\_cpp.robot.b2.SportClient.continuous\_gait}}\label{unitree-sdk2-cpp-robot-b2-sportclient-continuous-gait-1}}

对应 C++ SDK 操作
\passthrough{\lstinline!ContinuousGait(bool)!}。上游头文件没有可直接生成的业务说明时，本参考只保证签名映射，精确语义需查目标型号协议。

\textbf{签名}

\begin{lstlisting}[language=Python]
def continuous_gait(self, flag: bool) -> int
\end{lstlisting}

\textbf{可用性与安全性}

\textbf{\passthrough{\lstinline!SIGNATURE\_ONLY!} /
\passthrough{\lstinline!MOTION\_COMMAND!}}：仅用于补全和静态检查。当前二进制没有该
Python 方法；不得按可执行运动接口使用。

\textbf{参数}

\begin{longtable}[]{@{}
  >{\raggedright\arraybackslash}p{(\columnwidth - 6\tabcolsep) * \real{0.18}}
  >{\raggedright\arraybackslash}p{(\columnwidth - 6\tabcolsep) * \real{0.24}}
  >{\raggedright\arraybackslash}p{(\columnwidth - 6\tabcolsep) * \real{0.18}}
  >{\raggedright\arraybackslash}p{(\columnwidth - 6\tabcolsep) * \real{0.40}}@{}}
\toprule
名称 & Python 类型 & 默认值 & 含义 \\
\midrule
\endhead
\passthrough{\lstinline!flag!} & \passthrough{\lstinline!bool!} & 必填 &
布尔开关。\passthrough{\lstinline!True!} 和
\passthrough{\lstinline!False!} 的具体业务效果由当前方法定义。 对应 C++
参数 \passthrough{\lstinline!flag: bool!}。 只接受布尔语义。 \\
\bottomrule
\end{longtable}

\textbf{返回值}

\begin{longtable}[]{@{}
  >{\raggedright\arraybackslash}p{(\columnwidth - 4\tabcolsep) * \real{0.18}}
  >{\raggedright\arraybackslash}p{(\columnwidth - 4\tabcolsep) * \real{0.20}}
  >{\raggedright\arraybackslash}p{(\columnwidth - 4\tabcolsep) * \real{0.62}}@{}}
\toprule
位置 & 类型 & 含义 \\
\midrule
\endhead
返回值 & \passthrough{\lstinline!int!} & SDK 状态码；Unitree 示例通常以
\passthrough{\lstinline!0!} 表示成功，非零值需按具体服务定义解释。 \\
\bottomrule
\end{longtable}

\textbf{对应 C++}

\begin{itemize}
\tightlist
\item
  类：\passthrough{\lstinline!unitree::robot::b2::SportClient!}
\item
  签名：\passthrough{\lstinline!ContinuousGait(bool)!}
\item
  绑定策略：\passthrough{\lstinline!DIRECT!}
\item
  声明位置：\passthrough{\lstinline!include/unitree/robot/b2/sport/sport\_client.hpp:61!}
\end{itemize}

\textbf{说明}

\begin{itemize}
\tightlist
\item
  示例是局部调用片段；\passthrough{\lstinline!obj!}
  和参数变量表示已经按上层业务校验并准备好的对象。
\item
  该条目可以被编辑器和类型检查器识别，但当前运行时可能在属性查找阶段直接失败。
\item
  该方法属于运动命令。方法名包含
  \passthrough{\lstinline!stop!}、\passthrough{\lstinline!damp!} 或
  \passthrough{\lstinline!zero!} 也不代表它是独立物理急停。
\end{itemize}

\textbf{用法}

\begin{lstlisting}[language=Python]
from typing import TYPE_CHECKING

if TYPE_CHECKING:
    def planned_continuous_gait(obj: SportClient, flag: bool) -> int:
        return obj.continuous_gait(flag=flag)
\end{lstlisting}

\begin{quote}
{[}!CAUTION{]} 上面的 \passthrough{\lstinline!TYPE\_CHECKING!}
示例只表达计划调用形态。该分支在普通 Python
运行时为假，不代表当前扩展能执行此接口。
\end{quote}

\hypertarget{unitree-sdk2-cpp-robot-b2-sportclient-move-to-pos-1}{%
\paragraph{\texorpdfstring{\texttt{unitree\_sdk2\_cpp.robot.b2.SportClient.move\_to\_pos}}{unitree\_sdk2\_cpp.robot.b2.SportClient.move\_to\_pos}}\label{unitree-sdk2-cpp-robot-b2-sportclient-move-to-pos-1}}

对应 C++ SDK 操作
\passthrough{\lstinline!MoveToPos(float, float, float)!}。上游头文件没有可直接生成的业务说明时，本参考只保证签名映射，精确语义需查目标型号协议。

\textbf{签名}

\begin{lstlisting}[language=Python]
def move_to_pos(self, x: float, y: float, yaw: float) -> int
\end{lstlisting}

\textbf{可用性与安全性}

\textbf{\passthrough{\lstinline!SIGNATURE\_ONLY!} /
\passthrough{\lstinline!MOTION\_COMMAND!}}：仅用于补全和静态检查。当前二进制没有该
Python 方法；不得按可执行运动接口使用。

\textbf{参数}

\begin{longtable}[]{@{}
  >{\raggedright\arraybackslash}p{(\columnwidth - 6\tabcolsep) * \real{0.18}}
  >{\raggedright\arraybackslash}p{(\columnwidth - 6\tabcolsep) * \real{0.24}}
  >{\raggedright\arraybackslash}p{(\columnwidth - 6\tabcolsep) * \real{0.18}}
  >{\raggedright\arraybackslash}p{(\columnwidth - 6\tabcolsep) * \real{0.40}}@{}}
\toprule
名称 & Python 类型 & 默认值 & 含义 \\
\midrule
\endhead
\passthrough{\lstinline!x!} & \passthrough{\lstinline!float!} & 必填 & X
方向位置或分量参数。参考系、单位和范围必须查当前方法及目标型号协议。
对应 C++ 参数 \passthrough{\lstinline!x: float!}。
底层为浮点数；类型本身不说明单位、坐标系或业务安全范围。 \\
\passthrough{\lstinline!y!} & \passthrough{\lstinline!float!} & 必填 & Y
方向位置或分量参数。参考系、单位和范围必须查当前方法及目标型号协议。
对应 C++ 参数 \passthrough{\lstinline!y: float!}。
底层为浮点数；类型本身不说明单位、坐标系或业务安全范围。 \\
\passthrough{\lstinline!yaw!} & \passthrough{\lstinline!float!} & 必填 &
偏航角或偏航目标参数。参考系、单位和范围必须查目标型号协议。 对应 C++
参数 \passthrough{\lstinline!yaw: float!}。
底层为浮点数；类型本身不说明单位、坐标系或业务安全范围。 \\
\bottomrule
\end{longtable}

\textbf{返回值}

\begin{longtable}[]{@{}
  >{\raggedright\arraybackslash}p{(\columnwidth - 4\tabcolsep) * \real{0.18}}
  >{\raggedright\arraybackslash}p{(\columnwidth - 4\tabcolsep) * \real{0.20}}
  >{\raggedright\arraybackslash}p{(\columnwidth - 4\tabcolsep) * \real{0.62}}@{}}
\toprule
位置 & 类型 & 含义 \\
\midrule
\endhead
返回值 & \passthrough{\lstinline!int!} & SDK 状态码；Unitree 示例通常以
\passthrough{\lstinline!0!} 表示成功，非零值需按具体服务定义解释。 \\
\bottomrule
\end{longtable}

\textbf{对应 C++}

\begin{itemize}
\tightlist
\item
  类：\passthrough{\lstinline!unitree::robot::b2::SportClient!}
\item
  签名：\passthrough{\lstinline!MoveToPos(float, float, float)!}
\item
  绑定策略：\passthrough{\lstinline!DIRECT!}
\item
  声明位置：\passthrough{\lstinline!include/unitree/robot/b2/sport/sport\_client.hpp:63!}
\end{itemize}

\textbf{说明}

\begin{itemize}
\tightlist
\item
  示例是局部调用片段；\passthrough{\lstinline!obj!}
  和参数变量表示已经按上层业务校验并准备好的对象。
\item
  该条目可以被编辑器和类型检查器识别，但当前运行时可能在属性查找阶段直接失败。
\item
  该方法属于运动命令。方法名包含
  \passthrough{\lstinline!stop!}、\passthrough{\lstinline!damp!} 或
  \passthrough{\lstinline!zero!} 也不代表它是独立物理急停。
\end{itemize}

\textbf{用法}

\begin{lstlisting}[language=Python]
from typing import TYPE_CHECKING

if TYPE_CHECKING:
    def planned_move_to_pos(obj: SportClient, x: float, y: float, yaw: float) -> int:
        return obj.move_to_pos(x=x, y=y, yaw=yaw)
\end{lstlisting}

\begin{quote}
{[}!CAUTION{]} 上面的 \passthrough{\lstinline!TYPE\_CHECKING!}
示例只表达计划调用形态。该分支在普通 Python
运行时为假，不代表当前扩展能执行此接口。
\end{quote}

\hypertarget{unitree-sdk2-cpp-robot-b2-sportclient-switch-move-mode-1}{%
\paragraph{\texorpdfstring{\texttt{unitree\_sdk2\_cpp.robot.b2.SportClient.switch\_move\_mode}}{unitree\_sdk2\_cpp.robot.b2.SportClient.switch\_move\_mode}}\label{unitree-sdk2-cpp-robot-b2-sportclient-switch-move-mode-1}}

对应 C++ SDK 操作
\passthrough{\lstinline!SwitchMoveMode(bool)!}。上游头文件没有可直接生成的业务说明时，本参考只保证签名映射，精确语义需查目标型号协议。

\textbf{签名}

\begin{lstlisting}[language=Python]
def switch_move_mode(self, flag: bool) -> int
\end{lstlisting}

\textbf{可用性与安全性}

\textbf{\passthrough{\lstinline!SIGNATURE\_ONLY!} /
\passthrough{\lstinline!MOTION\_COMMAND!}}：仅用于补全和静态检查。当前二进制没有该
Python 方法；不得按可执行运动接口使用。

\textbf{参数}

\begin{longtable}[]{@{}
  >{\raggedright\arraybackslash}p{(\columnwidth - 6\tabcolsep) * \real{0.18}}
  >{\raggedright\arraybackslash}p{(\columnwidth - 6\tabcolsep) * \real{0.24}}
  >{\raggedright\arraybackslash}p{(\columnwidth - 6\tabcolsep) * \real{0.18}}
  >{\raggedright\arraybackslash}p{(\columnwidth - 6\tabcolsep) * \real{0.40}}@{}}
\toprule
名称 & Python 类型 & 默认值 & 含义 \\
\midrule
\endhead
\passthrough{\lstinline!flag!} & \passthrough{\lstinline!bool!} & 必填 &
布尔开关。\passthrough{\lstinline!True!} 和
\passthrough{\lstinline!False!} 的具体业务效果由当前方法定义。 对应 C++
参数 \passthrough{\lstinline!flag: bool!}。 只接受布尔语义。 \\
\bottomrule
\end{longtable}

\textbf{返回值}

\begin{longtable}[]{@{}
  >{\raggedright\arraybackslash}p{(\columnwidth - 4\tabcolsep) * \real{0.18}}
  >{\raggedright\arraybackslash}p{(\columnwidth - 4\tabcolsep) * \real{0.20}}
  >{\raggedright\arraybackslash}p{(\columnwidth - 4\tabcolsep) * \real{0.62}}@{}}
\toprule
位置 & 类型 & 含义 \\
\midrule
\endhead
返回值 & \passthrough{\lstinline!int!} & SDK 状态码；Unitree 示例通常以
\passthrough{\lstinline!0!} 表示成功，非零值需按具体服务定义解释。 \\
\bottomrule
\end{longtable}

\textbf{对应 C++}

\begin{itemize}
\tightlist
\item
  类：\passthrough{\lstinline!unitree::robot::b2::SportClient!}
\item
  签名：\passthrough{\lstinline!SwitchMoveMode(bool)!}
\item
  绑定策略：\passthrough{\lstinline!DIRECT!}
\item
  声明位置：\passthrough{\lstinline!include/unitree/robot/b2/sport/sport\_client.hpp:65!}
\end{itemize}

\textbf{说明}

\begin{itemize}
\tightlist
\item
  示例是局部调用片段；\passthrough{\lstinline!obj!}
  和参数变量表示已经按上层业务校验并准备好的对象。
\item
  该条目可以被编辑器和类型检查器识别，但当前运行时可能在属性查找阶段直接失败。
\item
  该方法属于运动命令。方法名包含
  \passthrough{\lstinline!stop!}、\passthrough{\lstinline!damp!} 或
  \passthrough{\lstinline!zero!} 也不代表它是独立物理急停。
\end{itemize}

\textbf{用法}

\begin{lstlisting}[language=Python]
from typing import TYPE_CHECKING

if TYPE_CHECKING:
    def planned_switch_move_mode(obj: SportClient, flag: bool) -> int:
        return obj.switch_move_mode(flag=flag)
\end{lstlisting}

\begin{quote}
{[}!CAUTION{]} 上面的 \passthrough{\lstinline!TYPE\_CHECKING!}
示例只表达计划调用形态。该分支在普通 Python
运行时为假，不代表当前扩展能执行此接口。
\end{quote}

\hypertarget{unitree-sdk2-cpp-robot-b2-sportclient-vision-walk-1}{%
\paragraph{\texorpdfstring{\texttt{unitree\_sdk2\_cpp.robot.b2.SportClient.vision\_walk}}{unitree\_sdk2\_cpp.robot.b2.SportClient.vision\_walk}}\label{unitree-sdk2-cpp-robot-b2-sportclient-vision-walk-1}}

对应 C++ SDK 操作
\passthrough{\lstinline!VisionWalk(bool)!}。上游头文件没有可直接生成的业务说明时，本参考只保证签名映射，精确语义需查目标型号协议。

\textbf{签名}

\begin{lstlisting}[language=Python]
def vision_walk(self, flag: bool) -> int
\end{lstlisting}

\textbf{可用性与安全性}

\textbf{\passthrough{\lstinline!SIGNATURE\_ONLY!} /
\passthrough{\lstinline!MOTION\_COMMAND!}}：仅用于补全和静态检查。当前二进制没有该
Python 方法；不得按可执行运动接口使用。

\textbf{参数}

\begin{longtable}[]{@{}
  >{\raggedright\arraybackslash}p{(\columnwidth - 6\tabcolsep) * \real{0.18}}
  >{\raggedright\arraybackslash}p{(\columnwidth - 6\tabcolsep) * \real{0.24}}
  >{\raggedright\arraybackslash}p{(\columnwidth - 6\tabcolsep) * \real{0.18}}
  >{\raggedright\arraybackslash}p{(\columnwidth - 6\tabcolsep) * \real{0.40}}@{}}
\toprule
名称 & Python 类型 & 默认值 & 含义 \\
\midrule
\endhead
\passthrough{\lstinline!flag!} & \passthrough{\lstinline!bool!} & 必填 &
布尔开关。\passthrough{\lstinline!True!} 和
\passthrough{\lstinline!False!} 的具体业务效果由当前方法定义。 对应 C++
参数 \passthrough{\lstinline!flag: bool!}。 只接受布尔语义。 \\
\bottomrule
\end{longtable}

\textbf{返回值}

\begin{longtable}[]{@{}
  >{\raggedright\arraybackslash}p{(\columnwidth - 4\tabcolsep) * \real{0.18}}
  >{\raggedright\arraybackslash}p{(\columnwidth - 4\tabcolsep) * \real{0.20}}
  >{\raggedright\arraybackslash}p{(\columnwidth - 4\tabcolsep) * \real{0.62}}@{}}
\toprule
位置 & 类型 & 含义 \\
\midrule
\endhead
返回值 & \passthrough{\lstinline!int!} & SDK 状态码；Unitree 示例通常以
\passthrough{\lstinline!0!} 表示成功，非零值需按具体服务定义解释。 \\
\bottomrule
\end{longtable}

\textbf{对应 C++}

\begin{itemize}
\tightlist
\item
  类：\passthrough{\lstinline!unitree::robot::b2::SportClient!}
\item
  签名：\passthrough{\lstinline!VisionWalk(bool)!}
\item
  绑定策略：\passthrough{\lstinline!DIRECT!}
\item
  声明位置：\passthrough{\lstinline!include/unitree/robot/b2/sport/sport\_client.hpp:67!}
\end{itemize}

\textbf{说明}

\begin{itemize}
\tightlist
\item
  示例是局部调用片段；\passthrough{\lstinline!obj!}
  和参数变量表示已经按上层业务校验并准备好的对象。
\item
  该条目可以被编辑器和类型检查器识别，但当前运行时可能在属性查找阶段直接失败。
\item
  该方法属于运动命令。方法名包含
  \passthrough{\lstinline!stop!}、\passthrough{\lstinline!damp!} 或
  \passthrough{\lstinline!zero!} 也不代表它是独立物理急停。
\end{itemize}

\textbf{用法}

\begin{lstlisting}[language=Python]
from typing import TYPE_CHECKING

if TYPE_CHECKING:
    def planned_vision_walk(obj: SportClient, flag: bool) -> int:
        return obj.vision_walk(flag=flag)
\end{lstlisting}

\begin{quote}
{[}!CAUTION{]} 上面的 \passthrough{\lstinline!TYPE\_CHECKING!}
示例只表达计划调用形态。该分支在普通 Python
运行时为假，不代表当前扩展能执行此接口。
\end{quote}

\hypertarget{unitree-sdk2-cpp-robot-b2-sportclient-hand-stand-1}{%
\paragraph{\texorpdfstring{\texttt{unitree\_sdk2\_cpp.robot.b2.SportClient.hand\_stand}}{unitree\_sdk2\_cpp.robot.b2.SportClient.hand\_stand}}\label{unitree-sdk2-cpp-robot-b2-sportclient-hand-stand-1}}

对应 C++ SDK 操作
\passthrough{\lstinline!HandStand(bool)!}。上游头文件没有可直接生成的业务说明时，本参考只保证签名映射，精确语义需查目标型号协议。

\textbf{签名}

\begin{lstlisting}[language=Python]
def hand_stand(self, flag: bool) -> int
\end{lstlisting}

\textbf{可用性与安全性}

\textbf{\passthrough{\lstinline!SIGNATURE\_ONLY!} /
\passthrough{\lstinline!MOTION\_COMMAND!}}：仅用于补全和静态检查。当前二进制没有该
Python 方法；不得按可执行运动接口使用。

\textbf{参数}

\begin{longtable}[]{@{}
  >{\raggedright\arraybackslash}p{(\columnwidth - 6\tabcolsep) * \real{0.18}}
  >{\raggedright\arraybackslash}p{(\columnwidth - 6\tabcolsep) * \real{0.24}}
  >{\raggedright\arraybackslash}p{(\columnwidth - 6\tabcolsep) * \real{0.18}}
  >{\raggedright\arraybackslash}p{(\columnwidth - 6\tabcolsep) * \real{0.40}}@{}}
\toprule
名称 & Python 类型 & 默认值 & 含义 \\
\midrule
\endhead
\passthrough{\lstinline!flag!} & \passthrough{\lstinline!bool!} & 必填 &
布尔开关。\passthrough{\lstinline!True!} 和
\passthrough{\lstinline!False!} 的具体业务效果由当前方法定义。 对应 C++
参数 \passthrough{\lstinline!flag: bool!}。 只接受布尔语义。 \\
\bottomrule
\end{longtable}

\textbf{返回值}

\begin{longtable}[]{@{}
  >{\raggedright\arraybackslash}p{(\columnwidth - 4\tabcolsep) * \real{0.18}}
  >{\raggedright\arraybackslash}p{(\columnwidth - 4\tabcolsep) * \real{0.20}}
  >{\raggedright\arraybackslash}p{(\columnwidth - 4\tabcolsep) * \real{0.62}}@{}}
\toprule
位置 & 类型 & 含义 \\
\midrule
\endhead
返回值 & \passthrough{\lstinline!int!} & SDK 状态码；Unitree 示例通常以
\passthrough{\lstinline!0!} 表示成功，非零值需按具体服务定义解释。 \\
\bottomrule
\end{longtable}

\textbf{对应 C++}

\begin{itemize}
\tightlist
\item
  类：\passthrough{\lstinline!unitree::robot::b2::SportClient!}
\item
  签名：\passthrough{\lstinline!HandStand(bool)!}
\item
  绑定策略：\passthrough{\lstinline!DIRECT!}
\item
  声明位置：\passthrough{\lstinline!include/unitree/robot/b2/sport/sport\_client.hpp:69!}
\end{itemize}

\textbf{说明}

\begin{itemize}
\tightlist
\item
  示例是局部调用片段；\passthrough{\lstinline!obj!}
  和参数变量表示已经按上层业务校验并准备好的对象。
\item
  该条目可以被编辑器和类型检查器识别，但当前运行时可能在属性查找阶段直接失败。
\item
  该方法属于运动命令。方法名包含
  \passthrough{\lstinline!stop!}、\passthrough{\lstinline!damp!} 或
  \passthrough{\lstinline!zero!} 也不代表它是独立物理急停。
\end{itemize}

\textbf{用法}

\begin{lstlisting}[language=Python]
from typing import TYPE_CHECKING

if TYPE_CHECKING:
    def planned_hand_stand(obj: SportClient, flag: bool) -> int:
        return obj.hand_stand(flag=flag)
\end{lstlisting}

\begin{quote}
{[}!CAUTION{]} 上面的 \passthrough{\lstinline!TYPE\_CHECKING!}
示例只表达计划调用形态。该分支在普通 Python
运行时为假，不代表当前扩展能执行此接口。
\end{quote}

\hypertarget{unitree-sdk2-cpp-robot-b2-sportclient-auto-recovery-set-1}{%
\paragraph{\texorpdfstring{\texttt{unitree\_sdk2\_cpp.robot.b2.SportClient.auto\_recovery\_set}}{unitree\_sdk2\_cpp.robot.b2.SportClient.auto\_recovery\_set}}\label{unitree-sdk2-cpp-robot-b2-sportclient-auto-recovery-set-1}}

对应 C++ SDK 操作
\passthrough{\lstinline!AutoRecoverySet(bool)!}。上游头文件没有可直接生成的业务说明时，本参考只保证签名映射，精确语义需查目标型号协议。

\textbf{签名}

\begin{lstlisting}[language=Python]
def auto_recovery_set(self, flag: bool) -> int
\end{lstlisting}

\textbf{可用性与安全性}

\textbf{\passthrough{\lstinline!SIGNATURE\_ONLY!} /
\passthrough{\lstinline!MOTION\_COMMAND!}}：仅用于补全和静态检查。当前二进制没有该
Python 方法；不得按可执行运动接口使用。

\textbf{参数}

\begin{longtable}[]{@{}
  >{\raggedright\arraybackslash}p{(\columnwidth - 6\tabcolsep) * \real{0.18}}
  >{\raggedright\arraybackslash}p{(\columnwidth - 6\tabcolsep) * \real{0.24}}
  >{\raggedright\arraybackslash}p{(\columnwidth - 6\tabcolsep) * \real{0.18}}
  >{\raggedright\arraybackslash}p{(\columnwidth - 6\tabcolsep) * \real{0.40}}@{}}
\toprule
名称 & Python 类型 & 默认值 & 含义 \\
\midrule
\endhead
\passthrough{\lstinline!flag!} & \passthrough{\lstinline!bool!} & 必填 &
布尔开关。\passthrough{\lstinline!True!} 和
\passthrough{\lstinline!False!} 的具体业务效果由当前方法定义。 对应 C++
参数 \passthrough{\lstinline!flag: bool!}。 只接受布尔语义。 \\
\bottomrule
\end{longtable}

\textbf{返回值}

\begin{longtable}[]{@{}
  >{\raggedright\arraybackslash}p{(\columnwidth - 4\tabcolsep) * \real{0.18}}
  >{\raggedright\arraybackslash}p{(\columnwidth - 4\tabcolsep) * \real{0.20}}
  >{\raggedright\arraybackslash}p{(\columnwidth - 4\tabcolsep) * \real{0.62}}@{}}
\toprule
位置 & 类型 & 含义 \\
\midrule
\endhead
返回值 & \passthrough{\lstinline!int!} & SDK 状态码；Unitree 示例通常以
\passthrough{\lstinline!0!} 表示成功，非零值需按具体服务定义解释。 \\
\bottomrule
\end{longtable}

\textbf{对应 C++}

\begin{itemize}
\tightlist
\item
  类：\passthrough{\lstinline!unitree::robot::b2::SportClient!}
\item
  签名：\passthrough{\lstinline!AutoRecoverySet(bool)!}
\item
  绑定策略：\passthrough{\lstinline!DIRECT!}
\item
  声明位置：\passthrough{\lstinline!include/unitree/robot/b2/sport/sport\_client.hpp:71!}
\end{itemize}

\textbf{说明}

\begin{itemize}
\tightlist
\item
  示例是局部调用片段；\passthrough{\lstinline!obj!}
  和参数变量表示已经按上层业务校验并准备好的对象。
\item
  该条目可以被编辑器和类型检查器识别，但当前运行时可能在属性查找阶段直接失败。
\item
  该方法属于运动命令。方法名包含
  \passthrough{\lstinline!stop!}、\passthrough{\lstinline!damp!} 或
  \passthrough{\lstinline!zero!} 也不代表它是独立物理急停。
\end{itemize}

\textbf{用法}

\begin{lstlisting}[language=Python]
from typing import TYPE_CHECKING

if TYPE_CHECKING:
    def planned_auto_recovery_set(obj: SportClient, flag: bool) -> int:
        return obj.auto_recovery_set(flag=flag)
\end{lstlisting}

\begin{quote}
{[}!CAUTION{]} 上面的 \passthrough{\lstinline!TYPE\_CHECKING!}
示例只表达计划调用形态。该分支在普通 Python
运行时为假，不代表当前扩展能执行此接口。
\end{quote}

\hypertarget{unitree-sdk2-cpp-robot-b2-sportclient-free-walk-1}{%
\paragraph{\texorpdfstring{\texttt{unitree\_sdk2\_cpp.robot.b2.SportClient.free\_walk}}{unitree\_sdk2\_cpp.robot.b2.SportClient.free\_walk}}\label{unitree-sdk2-cpp-robot-b2-sportclient-free-walk-1}}

对应 C++ SDK 操作
\passthrough{\lstinline!FreeWalk()!}。上游头文件没有可直接生成的业务说明时，本参考只保证签名映射，精确语义需查目标型号协议。

\textbf{签名}

\begin{lstlisting}[language=Python]
def free_walk(self) -> int
\end{lstlisting}

\textbf{可用性与安全性}

\textbf{\passthrough{\lstinline!SIGNATURE\_ONLY!} /
\passthrough{\lstinline!MOTION\_COMMAND!}}：仅用于补全和静态检查。当前二进制没有该
Python 方法；不得按可执行运动接口使用。

\textbf{参数}

无显式参数。实例方法中的 \passthrough{\lstinline!self!} 由 Python
自动传入。

\textbf{返回值}

\begin{longtable}[]{@{}
  >{\raggedright\arraybackslash}p{(\columnwidth - 4\tabcolsep) * \real{0.18}}
  >{\raggedright\arraybackslash}p{(\columnwidth - 4\tabcolsep) * \real{0.20}}
  >{\raggedright\arraybackslash}p{(\columnwidth - 4\tabcolsep) * \real{0.62}}@{}}
\toprule
位置 & 类型 & 含义 \\
\midrule
\endhead
返回值 & \passthrough{\lstinline!int!} & SDK 状态码；Unitree 示例通常以
\passthrough{\lstinline!0!} 表示成功，非零值需按具体服务定义解释。 \\
\bottomrule
\end{longtable}

\textbf{对应 C++}

\begin{itemize}
\tightlist
\item
  类：\passthrough{\lstinline!unitree::robot::b2::SportClient!}
\item
  签名：\passthrough{\lstinline!FreeWalk()!}
\item
  绑定策略：\passthrough{\lstinline!DIRECT!}
\item
  声明位置：\passthrough{\lstinline!include/unitree/robot/b2/sport/sport\_client.hpp:73!}
\end{itemize}

\textbf{说明}

\begin{itemize}
\tightlist
\item
  示例是局部调用片段；\passthrough{\lstinline!obj!}
  和参数变量表示已经按上层业务校验并准备好的对象。
\item
  该条目可以被编辑器和类型检查器识别，但当前运行时可能在属性查找阶段直接失败。
\item
  该方法属于运动命令。方法名包含
  \passthrough{\lstinline!stop!}、\passthrough{\lstinline!damp!} 或
  \passthrough{\lstinline!zero!} 也不代表它是独立物理急停。
\end{itemize}

\textbf{用法}

\begin{lstlisting}[language=Python]
from typing import TYPE_CHECKING

if TYPE_CHECKING:
    def planned_free_walk(obj: SportClient) -> int:
        return obj.free_walk()
\end{lstlisting}

\begin{quote}
{[}!CAUTION{]} 上面的 \passthrough{\lstinline!TYPE\_CHECKING!}
示例只表达计划调用形态。该分支在普通 Python
运行时为假，不代表当前扩展能执行此接口。
\end{quote}

\hypertarget{unitree-sdk2-cpp-robot-b2-sportclient-classic-walk-1}{%
\paragraph{\texorpdfstring{\texttt{unitree\_sdk2\_cpp.robot.b2.SportClient.classic\_walk}}{unitree\_sdk2\_cpp.robot.b2.SportClient.classic\_walk}}\label{unitree-sdk2-cpp-robot-b2-sportclient-classic-walk-1}}

对应 C++ SDK 操作
\passthrough{\lstinline!ClassicWalk(bool)!}。上游头文件没有可直接生成的业务说明时，本参考只保证签名映射，精确语义需查目标型号协议。

\textbf{签名}

\begin{lstlisting}[language=Python]
def classic_walk(self, flag: bool) -> int
\end{lstlisting}

\textbf{可用性与安全性}

\textbf{\passthrough{\lstinline!SIGNATURE\_ONLY!} /
\passthrough{\lstinline!MOTION\_COMMAND!}}：仅用于补全和静态检查。当前二进制没有该
Python 方法；不得按可执行运动接口使用。

\textbf{参数}

\begin{longtable}[]{@{}
  >{\raggedright\arraybackslash}p{(\columnwidth - 6\tabcolsep) * \real{0.18}}
  >{\raggedright\arraybackslash}p{(\columnwidth - 6\tabcolsep) * \real{0.24}}
  >{\raggedright\arraybackslash}p{(\columnwidth - 6\tabcolsep) * \real{0.18}}
  >{\raggedright\arraybackslash}p{(\columnwidth - 6\tabcolsep) * \real{0.40}}@{}}
\toprule
名称 & Python 类型 & 默认值 & 含义 \\
\midrule
\endhead
\passthrough{\lstinline!flag!} & \passthrough{\lstinline!bool!} & 必填 &
布尔开关。\passthrough{\lstinline!True!} 和
\passthrough{\lstinline!False!} 的具体业务效果由当前方法定义。 对应 C++
参数 \passthrough{\lstinline!flag: bool!}。 只接受布尔语义。 \\
\bottomrule
\end{longtable}

\textbf{返回值}

\begin{longtable}[]{@{}
  >{\raggedright\arraybackslash}p{(\columnwidth - 4\tabcolsep) * \real{0.18}}
  >{\raggedright\arraybackslash}p{(\columnwidth - 4\tabcolsep) * \real{0.20}}
  >{\raggedright\arraybackslash}p{(\columnwidth - 4\tabcolsep) * \real{0.62}}@{}}
\toprule
位置 & 类型 & 含义 \\
\midrule
\endhead
返回值 & \passthrough{\lstinline!int!} & SDK 状态码；Unitree 示例通常以
\passthrough{\lstinline!0!} 表示成功，非零值需按具体服务定义解释。 \\
\bottomrule
\end{longtable}

\textbf{对应 C++}

\begin{itemize}
\tightlist
\item
  类：\passthrough{\lstinline!unitree::robot::b2::SportClient!}
\item
  签名：\passthrough{\lstinline!ClassicWalk(bool)!}
\item
  绑定策略：\passthrough{\lstinline!DIRECT!}
\item
  声明位置：\passthrough{\lstinline!include/unitree/robot/b2/sport/sport\_client.hpp:75!}
\end{itemize}

\textbf{说明}

\begin{itemize}
\tightlist
\item
  示例是局部调用片段；\passthrough{\lstinline!obj!}
  和参数变量表示已经按上层业务校验并准备好的对象。
\item
  该条目可以被编辑器和类型检查器识别，但当前运行时可能在属性查找阶段直接失败。
\item
  该方法属于运动命令。方法名包含
  \passthrough{\lstinline!stop!}、\passthrough{\lstinline!damp!} 或
  \passthrough{\lstinline!zero!} 也不代表它是独立物理急停。
\end{itemize}

\textbf{用法}

\begin{lstlisting}[language=Python]
from typing import TYPE_CHECKING

if TYPE_CHECKING:
    def planned_classic_walk(obj: SportClient, flag: bool) -> int:
        return obj.classic_walk(flag=flag)
\end{lstlisting}

\begin{quote}
{[}!CAUTION{]} 上面的 \passthrough{\lstinline!TYPE\_CHECKING!}
示例只表达计划调用形态。该分支在普通 Python
运行时为假，不代表当前扩展能执行此接口。
\end{quote}

\hypertarget{unitree-sdk2-cpp-robot-b2-sportclient-fast-walk-1}{%
\paragraph{\texorpdfstring{\texttt{unitree\_sdk2\_cpp.robot.b2.SportClient.fast\_walk}}{unitree\_sdk2\_cpp.robot.b2.SportClient.fast\_walk}}\label{unitree-sdk2-cpp-robot-b2-sportclient-fast-walk-1}}

对应 C++ SDK 操作
\passthrough{\lstinline!FastWalk(bool)!}。上游头文件没有可直接生成的业务说明时，本参考只保证签名映射，精确语义需查目标型号协议。

\textbf{签名}

\begin{lstlisting}[language=Python]
def fast_walk(self, flag: bool) -> int
\end{lstlisting}

\textbf{可用性与安全性}

\textbf{\passthrough{\lstinline!SIGNATURE\_ONLY!} /
\passthrough{\lstinline!MOTION\_COMMAND!}}：仅用于补全和静态检查。当前二进制没有该
Python 方法；不得按可执行运动接口使用。

\textbf{参数}

\begin{longtable}[]{@{}
  >{\raggedright\arraybackslash}p{(\columnwidth - 6\tabcolsep) * \real{0.18}}
  >{\raggedright\arraybackslash}p{(\columnwidth - 6\tabcolsep) * \real{0.24}}
  >{\raggedright\arraybackslash}p{(\columnwidth - 6\tabcolsep) * \real{0.18}}
  >{\raggedright\arraybackslash}p{(\columnwidth - 6\tabcolsep) * \real{0.40}}@{}}
\toprule
名称 & Python 类型 & 默认值 & 含义 \\
\midrule
\endhead
\passthrough{\lstinline!flag!} & \passthrough{\lstinline!bool!} & 必填 &
布尔开关。\passthrough{\lstinline!True!} 和
\passthrough{\lstinline!False!} 的具体业务效果由当前方法定义。 对应 C++
参数 \passthrough{\lstinline!flag: bool!}。 只接受布尔语义。 \\
\bottomrule
\end{longtable}

\textbf{返回值}

\begin{longtable}[]{@{}
  >{\raggedright\arraybackslash}p{(\columnwidth - 4\tabcolsep) * \real{0.18}}
  >{\raggedright\arraybackslash}p{(\columnwidth - 4\tabcolsep) * \real{0.20}}
  >{\raggedright\arraybackslash}p{(\columnwidth - 4\tabcolsep) * \real{0.62}}@{}}
\toprule
位置 & 类型 & 含义 \\
\midrule
\endhead
返回值 & \passthrough{\lstinline!int!} & SDK 状态码；Unitree 示例通常以
\passthrough{\lstinline!0!} 表示成功，非零值需按具体服务定义解释。 \\
\bottomrule
\end{longtable}

\textbf{对应 C++}

\begin{itemize}
\tightlist
\item
  类：\passthrough{\lstinline!unitree::robot::b2::SportClient!}
\item
  签名：\passthrough{\lstinline!FastWalk(bool)!}
\item
  绑定策略：\passthrough{\lstinline!DIRECT!}
\item
  声明位置：\passthrough{\lstinline!include/unitree/robot/b2/sport/sport\_client.hpp:77!}
\end{itemize}

\textbf{说明}

\begin{itemize}
\tightlist
\item
  示例是局部调用片段；\passthrough{\lstinline!obj!}
  和参数变量表示已经按上层业务校验并准备好的对象。
\item
  该条目可以被编辑器和类型检查器识别，但当前运行时可能在属性查找阶段直接失败。
\item
  该方法属于运动命令。方法名包含
  \passthrough{\lstinline!stop!}、\passthrough{\lstinline!damp!} 或
  \passthrough{\lstinline!zero!} 也不代表它是独立物理急停。
\end{itemize}

\textbf{用法}

\begin{lstlisting}[language=Python]
from typing import TYPE_CHECKING

if TYPE_CHECKING:
    def planned_fast_walk(obj: SportClient, flag: bool) -> int:
        return obj.fast_walk(flag=flag)
\end{lstlisting}

\begin{quote}
{[}!CAUTION{]} 上面的 \passthrough{\lstinline!TYPE\_CHECKING!}
示例只表达计划调用形态。该分支在普通 Python
运行时为假，不代表当前扩展能执行此接口。
\end{quote}

\hypertarget{unitree-sdk2-cpp-robot-b2-sportclient-euler-1}{%
\paragraph{\texorpdfstring{\texttt{unitree\_sdk2\_cpp.robot.b2.SportClient.euler}}{unitree\_sdk2\_cpp.robot.b2.SportClient.euler}}\label{unitree-sdk2-cpp-robot-b2-sportclient-euler-1}}

对应 C++ SDK 操作
\passthrough{\lstinline!Euler(float, float, float)!}。上游头文件没有可直接生成的业务说明时，本参考只保证签名映射，精确语义需查目标型号协议。

\textbf{签名}

\begin{lstlisting}[language=Python]
def euler(self, roll: float, pitch: float, yaw: float) -> int
\end{lstlisting}

\textbf{可用性与安全性}

\textbf{\passthrough{\lstinline!SIGNATURE\_ONLY!} /
\passthrough{\lstinline!MOTION\_COMMAND!}}：仅用于补全和静态检查。当前二进制没有该
Python 方法；不得按可执行运动接口使用。

\textbf{参数}

\begin{longtable}[]{@{}
  >{\raggedright\arraybackslash}p{(\columnwidth - 6\tabcolsep) * \real{0.18}}
  >{\raggedright\arraybackslash}p{(\columnwidth - 6\tabcolsep) * \real{0.24}}
  >{\raggedright\arraybackslash}p{(\columnwidth - 6\tabcolsep) * \real{0.18}}
  >{\raggedright\arraybackslash}p{(\columnwidth - 6\tabcolsep) * \real{0.40}}@{}}
\toprule
名称 & Python 类型 & 默认值 & 含义 \\
\midrule
\endhead
\passthrough{\lstinline!roll!} & \passthrough{\lstinline!float!} & 必填
& 传给该接口的 \passthrough{\lstinline!roll!} 参数，Python 类型为
\passthrough{\lstinline!float!}。精确含义、单位、范围和枚举值需查目标型号对应头文件与协议。
对应 C++ 参数 \passthrough{\lstinline!roll: float!}。
底层为浮点数；类型本身不说明单位、坐标系或业务安全范围。 \\
\passthrough{\lstinline!pitch!} & \passthrough{\lstinline!float!} & 必填
& 传给该接口的 \passthrough{\lstinline!pitch!} 参数，Python 类型为
\passthrough{\lstinline!float!}。精确含义、单位、范围和枚举值需查目标型号对应头文件与协议。
对应 C++ 参数 \passthrough{\lstinline!pitch: float!}。
底层为浮点数；类型本身不说明单位、坐标系或业务安全范围。 \\
\passthrough{\lstinline!yaw!} & \passthrough{\lstinline!float!} & 必填 &
偏航角或偏航目标参数。参考系、单位和范围必须查目标型号协议。 对应 C++
参数 \passthrough{\lstinline!yaw: float!}。
底层为浮点数；类型本身不说明单位、坐标系或业务安全范围。 \\
\bottomrule
\end{longtable}

\textbf{返回值}

\begin{longtable}[]{@{}
  >{\raggedright\arraybackslash}p{(\columnwidth - 4\tabcolsep) * \real{0.18}}
  >{\raggedright\arraybackslash}p{(\columnwidth - 4\tabcolsep) * \real{0.20}}
  >{\raggedright\arraybackslash}p{(\columnwidth - 4\tabcolsep) * \real{0.62}}@{}}
\toprule
位置 & 类型 & 含义 \\
\midrule
\endhead
返回值 & \passthrough{\lstinline!int!} & SDK 状态码；Unitree 示例通常以
\passthrough{\lstinline!0!} 表示成功，非零值需按具体服务定义解释。 \\
\bottomrule
\end{longtable}

\textbf{对应 C++}

\begin{itemize}
\tightlist
\item
  类：\passthrough{\lstinline!unitree::robot::b2::SportClient!}
\item
  签名：\passthrough{\lstinline!Euler(float, float, float)!}
\item
  绑定策略：\passthrough{\lstinline!DIRECT!}
\item
  声明位置：\passthrough{\lstinline!include/unitree/robot/b2/sport/sport\_client.hpp:79!}
\end{itemize}

\textbf{说明}

\begin{itemize}
\tightlist
\item
  示例是局部调用片段；\passthrough{\lstinline!obj!}
  和参数变量表示已经按上层业务校验并准备好的对象。
\item
  该条目可以被编辑器和类型检查器识别，但当前运行时可能在属性查找阶段直接失败。
\item
  该方法属于运动命令。方法名包含
  \passthrough{\lstinline!stop!}、\passthrough{\lstinline!damp!} 或
  \passthrough{\lstinline!zero!} 也不代表它是独立物理急停。
\end{itemize}

\textbf{用法}

\begin{lstlisting}[language=Python]
from typing import TYPE_CHECKING

if TYPE_CHECKING:
    def planned_euler(obj: SportClient, roll: float, pitch: float, yaw: float) -> int:
        return obj.euler(roll=roll, pitch=pitch, yaw=yaw)
\end{lstlisting}

\begin{quote}
{[}!CAUTION{]} 上面的 \passthrough{\lstinline!TYPE\_CHECKING!}
示例只表达计划调用形态。该分支在普通 Python
运行时为假，不代表当前扩展能执行此接口。
\end{quote}

\hypertarget{unitree-sdk2-cpp-robot-b2-sportclient-free-height-1}{%
\paragraph{\texorpdfstring{\texttt{unitree\_sdk2\_cpp.robot.b2.SportClient.free\_height}}{unitree\_sdk2\_cpp.robot.b2.SportClient.free\_height}}\label{unitree-sdk2-cpp-robot-b2-sportclient-free-height-1}}

对应 C++ SDK 操作
\passthrough{\lstinline!FreeHeight(bool)!}。上游头文件没有可直接生成的业务说明时，本参考只保证签名映射，精确语义需查目标型号协议。

\textbf{签名}

\begin{lstlisting}[language=Python]
def free_height(self, flag: bool) -> int
\end{lstlisting}

\textbf{可用性与安全性}

\textbf{\passthrough{\lstinline!SIGNATURE\_ONLY!} /
\passthrough{\lstinline!MOTION\_COMMAND!}}：仅用于补全和静态检查。当前二进制没有该
Python 方法；不得按可执行运动接口使用。

\textbf{参数}

\begin{longtable}[]{@{}
  >{\raggedright\arraybackslash}p{(\columnwidth - 6\tabcolsep) * \real{0.18}}
  >{\raggedright\arraybackslash}p{(\columnwidth - 6\tabcolsep) * \real{0.24}}
  >{\raggedright\arraybackslash}p{(\columnwidth - 6\tabcolsep) * \real{0.18}}
  >{\raggedright\arraybackslash}p{(\columnwidth - 6\tabcolsep) * \real{0.40}}@{}}
\toprule
名称 & Python 类型 & 默认值 & 含义 \\
\midrule
\endhead
\passthrough{\lstinline!flag!} & \passthrough{\lstinline!bool!} & 必填 &
布尔开关。\passthrough{\lstinline!True!} 和
\passthrough{\lstinline!False!} 的具体业务效果由当前方法定义。 对应 C++
参数 \passthrough{\lstinline!flag: bool!}。 只接受布尔语义。 \\
\bottomrule
\end{longtable}

\textbf{返回值}

\begin{longtable}[]{@{}
  >{\raggedright\arraybackslash}p{(\columnwidth - 4\tabcolsep) * \real{0.18}}
  >{\raggedright\arraybackslash}p{(\columnwidth - 4\tabcolsep) * \real{0.20}}
  >{\raggedright\arraybackslash}p{(\columnwidth - 4\tabcolsep) * \real{0.62}}@{}}
\toprule
位置 & 类型 & 含义 \\
\midrule
\endhead
返回值 & \passthrough{\lstinline!int!} & SDK 状态码；Unitree 示例通常以
\passthrough{\lstinline!0!} 表示成功，非零值需按具体服务定义解释。 \\
\bottomrule
\end{longtable}

\textbf{对应 C++}

\begin{itemize}
\tightlist
\item
  类：\passthrough{\lstinline!unitree::robot::b2::SportClient!}
\item
  签名：\passthrough{\lstinline!FreeHeight(bool)!}
\item
  绑定策略：\passthrough{\lstinline!DIRECT!}
\item
  声明位置：\passthrough{\lstinline!include/unitree/robot/b2/sport/sport\_client.hpp:81!}
\end{itemize}

\textbf{说明}

\begin{itemize}
\tightlist
\item
  示例是局部调用片段；\passthrough{\lstinline!obj!}
  和参数变量表示已经按上层业务校验并准备好的对象。
\item
  该条目可以被编辑器和类型检查器识别，但当前运行时可能在属性查找阶段直接失败。
\item
  该方法属于运动命令。方法名包含
  \passthrough{\lstinline!stop!}、\passthrough{\lstinline!damp!} 或
  \passthrough{\lstinline!zero!} 也不代表它是独立物理急停。
\end{itemize}

\textbf{用法}

\begin{lstlisting}[language=Python]
from typing import TYPE_CHECKING

if TYPE_CHECKING:
    def planned_free_height(obj: SportClient, flag: bool) -> int:
        return obj.free_height(flag=flag)
\end{lstlisting}

\begin{quote}
{[}!CAUTION{]} 上面的 \passthrough{\lstinline!TYPE\_CHECKING!}
示例只表达计划调用形态。该分支在普通 Python
运行时为假，不代表当前扩展能执行此接口。
\end{quote}

\hypertarget{unitree-sdk2-cpp-robot-b2-sportclient-gait-height-1}{%
\paragraph{\texorpdfstring{\texttt{unitree\_sdk2\_cpp.robot.b2.SportClient.gait\_height}}{unitree\_sdk2\_cpp.robot.b2.SportClient.gait\_height}}\label{unitree-sdk2-cpp-robot-b2-sportclient-gait-height-1}}

对应 C++ SDK 操作
\passthrough{\lstinline!GaitHeight(bool)!}。上游头文件没有可直接生成的业务说明时，本参考只保证签名映射，精确语义需查目标型号协议。

\textbf{签名}

\begin{lstlisting}[language=Python]
def gait_height(self, flag: bool) -> int
\end{lstlisting}

\textbf{可用性与安全性}

\textbf{\passthrough{\lstinline!SIGNATURE\_ONLY!} /
\passthrough{\lstinline!MOTION\_COMMAND!}}：仅用于补全和静态检查。当前二进制没有该
Python 方法；不得按可执行运动接口使用。

\textbf{参数}

\begin{longtable}[]{@{}
  >{\raggedright\arraybackslash}p{(\columnwidth - 6\tabcolsep) * \real{0.18}}
  >{\raggedright\arraybackslash}p{(\columnwidth - 6\tabcolsep) * \real{0.24}}
  >{\raggedright\arraybackslash}p{(\columnwidth - 6\tabcolsep) * \real{0.18}}
  >{\raggedright\arraybackslash}p{(\columnwidth - 6\tabcolsep) * \real{0.40}}@{}}
\toprule
名称 & Python 类型 & 默认值 & 含义 \\
\midrule
\endhead
\passthrough{\lstinline!flag!} & \passthrough{\lstinline!bool!} & 必填 &
布尔开关。\passthrough{\lstinline!True!} 和
\passthrough{\lstinline!False!} 的具体业务效果由当前方法定义。 对应 C++
参数 \passthrough{\lstinline!flag: bool!}。 只接受布尔语义。 \\
\bottomrule
\end{longtable}

\textbf{返回值}

\begin{longtable}[]{@{}
  >{\raggedright\arraybackslash}p{(\columnwidth - 4\tabcolsep) * \real{0.18}}
  >{\raggedright\arraybackslash}p{(\columnwidth - 4\tabcolsep) * \real{0.20}}
  >{\raggedright\arraybackslash}p{(\columnwidth - 4\tabcolsep) * \real{0.62}}@{}}
\toprule
位置 & 类型 & 含义 \\
\midrule
\endhead
返回值 & \passthrough{\lstinline!int!} & SDK 状态码；Unitree 示例通常以
\passthrough{\lstinline!0!} 表示成功，非零值需按具体服务定义解释。 \\
\bottomrule
\end{longtable}

\textbf{对应 C++}

\begin{itemize}
\tightlist
\item
  类：\passthrough{\lstinline!unitree::robot::b2::SportClient!}
\item
  签名：\passthrough{\lstinline!GaitHeight(bool)!}
\item
  绑定策略：\passthrough{\lstinline!DIRECT!}
\item
  声明位置：\passthrough{\lstinline!include/unitree/robot/b2/sport/sport\_client.hpp:83!}
\end{itemize}

\textbf{说明}

\begin{itemize}
\tightlist
\item
  示例是局部调用片段；\passthrough{\lstinline!obj!}
  和参数变量表示已经按上层业务校验并准备好的对象。
\item
  该条目可以被编辑器和类型检查器识别，但当前运行时可能在属性查找阶段直接失败。
\item
  该方法属于运动命令。方法名包含
  \passthrough{\lstinline!stop!}、\passthrough{\lstinline!damp!} 或
  \passthrough{\lstinline!zero!} 也不代表它是独立物理急停。
\end{itemize}

\textbf{用法}

\begin{lstlisting}[language=Python]
from typing import TYPE_CHECKING

if TYPE_CHECKING:
    def planned_gait_height(obj: SportClient, flag: bool) -> int:
        return obj.gait_height(flag=flag)
\end{lstlisting}

\begin{quote}
{[}!CAUTION{]} 上面的 \passthrough{\lstinline!TYPE\_CHECKING!}
示例只表达计划调用形态。该分支在普通 Python
运行时为假，不代表当前扩展能执行此接口。
\end{quote}

\hypertarget{unitree-sdk2-cpp-robot-b2-stpathpoint}{%
\subsubsection{\texorpdfstring{\texttt{unitree\_sdk2\_cpp.robot.b2.stPathPoint}}{unitree\_sdk2\_cpp.robot.b2.stPathPoint}}\label{unitree-sdk2-cpp-robot-b2-stpathpoint}}

SDK 数据值类型；公开字段和转换方法见下文。

\textbf{导入}

\begin{lstlisting}[language=Python]
from unitree_sdk2_cpp.robot.b2 import stPathPoint
\end{lstlisting}

\textbf{基类}：\passthrough{\lstinline!object!}

\textbf{公开属性}

\begin{longtable}[]{@{}
  >{\raggedright\arraybackslash}p{(\columnwidth - 6\tabcolsep) * \real{0.18}}
  >{\raggedright\arraybackslash}p{(\columnwidth - 6\tabcolsep) * \real{0.24}}
  >{\raggedright\arraybackslash}p{(\columnwidth - 6\tabcolsep) * \real{0.18}}
  >{\raggedright\arraybackslash}p{(\columnwidth - 6\tabcolsep) * \real{0.40}}@{}}
\toprule
名称 & Python 类型 & 含义 & 用法 \\
\midrule
\endhead
\passthrough{\lstinline!time\_from\_start!} &
\passthrough{\lstinline!float!} & 传给该接口的
\passthrough{\lstinline!time!} \passthrough{\lstinline!from!}
\passthrough{\lstinline!start!} 参数，Python 类型为
\passthrough{\lstinline!float!}。精确含义、单位、范围和枚举值需查目标型号对应头文件与协议。
& \passthrough{\lstinline!value = obj.time\_from\_start!} /
\passthrough{\lstinline!obj.time\_from\_start = value!} \\
\passthrough{\lstinline!x!} & \passthrough{\lstinline!float!} & X
方向位置或分量参数。参考系、单位和范围必须查当前方法及目标型号协议。 &
\passthrough{\lstinline!value = obj.x!} /
\passthrough{\lstinline!obj.x = value!} \\
\passthrough{\lstinline!y!} & \passthrough{\lstinline!float!} & Y
方向位置或分量参数。参考系、单位和范围必须查当前方法及目标型号协议。 &
\passthrough{\lstinline!value = obj.y!} /
\passthrough{\lstinline!obj.y = value!} \\
\passthrough{\lstinline!yaw!} & \passthrough{\lstinline!float!} &
偏航角或偏航目标参数。参考系、单位和范围必须查目标型号协议。 &
\passthrough{\lstinline!value = obj.yaw!} /
\passthrough{\lstinline!obj.yaw = value!} \\
\passthrough{\lstinline!vx!} & \passthrough{\lstinline!float!} & X
方向速度参数。坐标系、单位、符号和安全范围必须查目标型号运动协议。 &
\passthrough{\lstinline!value = obj.vx!} /
\passthrough{\lstinline!obj.vx = value!} \\
\passthrough{\lstinline!vy!} & \passthrough{\lstinline!float!} & Y
方向速度参数。坐标系、单位、符号和安全范围必须查目标型号运动协议。 &
\passthrough{\lstinline!value = obj.vy!} /
\passthrough{\lstinline!obj.vy = value!} \\
\passthrough{\lstinline!vyaw!} & \passthrough{\lstinline!float!} &
偏航角速度参数。单位、符号和安全范围必须查目标型号运动协议。 &
\passthrough{\lstinline!value = obj.vyaw!} /
\passthrough{\lstinline!obj.vyaw = value!} \\
\bottomrule
\end{longtable}

\begin{center}\rule{0.5\linewidth}{0.5pt}\end{center}

\hypertarget{unitree-sdk2-cpp-robot-g1}{%
\subsection{G1 Robot API}\label{unitree-sdk2-cpp-robot-g1}}

模块：\passthrough{\lstinline!unitree\_sdk2\_cpp.robot.g1!}

\hypertarget{ux7c7bux7d22ux5f15-10}{%
\subsubsection{类索引}\label{ux7c7bux7d22ux5f15-10}}

\begin{longtable}[]{@{}
  >{\raggedright\arraybackslash}p{(\columnwidth - 4\tabcolsep) * \real{0.18}}
  >{\raggedleft\arraybackslash}p{(\columnwidth - 4\tabcolsep) * \real{0.20}}
  >{\raggedleft\arraybackslash}p{(\columnwidth - 4\tabcolsep) * \real{0.62}}@{}}
\toprule
类 & 公开函数签名 & 属性 \\
\midrule
\endhead
\protect\hyperlink{unitree-sdk2-cpp-robot-g1-internalfsmmode}{\passthrough{\lstinline!InternalFsmMode!}}
& 0 & 0 \\
\protect\hyperlink{unitree-sdk2-cpp-robot-g1-agvclient}{\passthrough{\lstinline!AgvClient!}}
& 4 & 0 \\
\protect\hyperlink{unitree-sdk2-cpp-robot-g1-audioclient}{\passthrough{\lstinline!AudioClient!}}
& 8 & 0 \\
\protect\hyperlink{unitree-sdk2-cpp-robot-g1-g1armactionclient}{\passthrough{\lstinline!G1ArmActionClient!}}
& 6 & 0 \\
\protect\hyperlink{unitree-sdk2-cpp-robot-g1-jsonizedatavecfloat}{\passthrough{\lstinline!JsonizeDataVecFloat!}}
& 3 & 0 \\
\protect\hyperlink{unitree-sdk2-cpp-robot-g1-jsonizevelocitycommand}{\passthrough{\lstinline!JsonizeVelocityCommand!}}
& 3 & 0 \\
\protect\hyperlink{unitree-sdk2-cpp-robot-g1-ledcontrolparameter}{\passthrough{\lstinline!LedControlParameter!}}
& 3 & 0 \\
\protect\hyperlink{unitree-sdk2-cpp-robot-g1-lococlient}{\passthrough{\lstinline!LocoClient!}}
& 35 & 0 \\
\protect\hyperlink{unitree-sdk2-cpp-robot-g1-moveparameter}{\passthrough{\lstinline!MoveParameter!}}
& 3 & 0 \\
\protect\hyperlink{unitree-sdk2-cpp-robot-g1-playstopparameter}{\passthrough{\lstinline!PlayStopParameter!}}
& 3 & 0 \\
\protect\hyperlink{unitree-sdk2-cpp-robot-g1-playstreamparameter}{\passthrough{\lstinline!PlayStreamParameter!}}
& 3 & 0 \\
\protect\hyperlink{unitree-sdk2-cpp-robot-g1-ttsmakerparameter}{\passthrough{\lstinline!TtsMakerParameter!}}
& 3 & 0 \\
\bottomrule
\end{longtable}

\hypertarget{unitree-sdk2-cpp-robot-g1-internalfsmmode}{%
\subsubsection{\texorpdfstring{\texttt{unitree\_sdk2\_cpp.robot.g1.InternalFsmMode}}{unitree\_sdk2\_cpp.robot.g1.InternalFsmMode}}\label{unitree-sdk2-cpp-robot-g1-internalfsmmode}}

SDK 整数枚举。枚举成员和值以目标版本头文件为准。

\textbf{导入}

\begin{lstlisting}[language=Python]
from unitree_sdk2_cpp.robot.g1 import InternalFsmMode
\end{lstlisting}

\textbf{基类}：\passthrough{\lstinline!enum.IntEnum!}

\textbf{枚举成员}

\begin{longtable}[]{@{}
  >{\raggedright\arraybackslash}p{(\columnwidth - 2\tabcolsep) * \real{0.35}}
  >{\raggedright\arraybackslash}p{(\columnwidth - 2\tabcolsep) * \real{0.65}}@{}}
\toprule
名称 & Stub 值 \\
\midrule
\endhead
\passthrough{\lstinline!LAST!} & \passthrough{\lstinline!...!} \\
\passthrough{\lstinline!PASSIVE!} & \passthrough{\lstinline!...!} \\
\passthrough{\lstinline!WALKRUN!} & \passthrough{\lstinline!...!} \\
\bottomrule
\end{longtable}

\hypertarget{unitree-sdk2-cpp-robot-g1-agvclient}{%
\subsubsection{\texorpdfstring{\texttt{unitree\_sdk2\_cpp.robot.g1.AgvClient}}{unitree\_sdk2\_cpp.robot.g1.AgvClient}}\label{unitree-sdk2-cpp-robot-g1-agvclient}}

Robot 服务客户端。构造、初始化和具体方法的可用性必须分别检查。

\textbf{导入}

\begin{lstlisting}[language=Python]
from unitree_sdk2_cpp.robot.g1 import AgvClient
\end{lstlisting}

\textbf{基类}：\passthrough{\lstinline!Client!}

\textbf{方法索引}

\begin{longtable}[]{@{}
  >{\raggedright\arraybackslash}p{(\columnwidth - 6\tabcolsep) * \real{0.18}}
  >{\raggedright\arraybackslash}p{(\columnwidth - 6\tabcolsep) * \real{0.24}}
  >{\raggedright\arraybackslash}p{(\columnwidth - 6\tabcolsep) * \real{0.18}}
  >{\raggedright\arraybackslash}p{(\columnwidth - 6\tabcolsep) * \real{0.40}}@{}}
\toprule
方法 & Python 签名 & 状态 & 安全分类 \\
\midrule
\endhead
\protect\hyperlink{unitree-sdk2-cpp-robot-g1-agvclient-dunder-init-1}{\passthrough{\lstinline!\_\_init\_\_!}}
& \passthrough{\lstinline!def \_\_init\_\_(self) -> None!} &
\passthrough{\lstinline!SIGNATURE\_ONLY!} &
\passthrough{\lstinline!CONSTRUCTION!} \\
\protect\hyperlink{unitree-sdk2-cpp-robot-g1-agvclient-init-1}{\passthrough{\lstinline!init!}}
& \passthrough{\lstinline!def init(self) -> None!} &
\passthrough{\lstinline!SIGNATURE\_ONLY!} &
\passthrough{\lstinline!INITIALIZATION!} \\
\protect\hyperlink{unitree-sdk2-cpp-robot-g1-agvclient-move-1}{\passthrough{\lstinline!move!}}
&
\passthrough{\lstinline!def move(self, vx: float, vy: float, vyaw: float) -> int!}
& \passthrough{\lstinline!SIGNATURE\_ONLY!} &
\passthrough{\lstinline!MOTION\_COMMAND!} \\
\protect\hyperlink{unitree-sdk2-cpp-robot-g1-agvclient-height-adjust-1}{\passthrough{\lstinline!height\_adjust!}}
& \passthrough{\lstinline!def height\_adjust(self, vz: float) -> int!} &
\passthrough{\lstinline!SIGNATURE\_ONLY!} &
\passthrough{\lstinline!MOTION\_COMMAND!} \\
\bottomrule
\end{longtable}

\hypertarget{unitree-sdk2-cpp-robot-g1-agvclient-dunder-init-1}{%
\paragraph{\texorpdfstring{\texttt{unitree\_sdk2\_cpp.robot.g1.AgvClient.\_\_init\_\_}}{unitree\_sdk2\_cpp.robot.g1.AgvClient.\_\_init\_\_}}\label{unitree-sdk2-cpp-robot-g1-agvclient-dunder-init-1}}

初始化 \passthrough{\lstinline!AgvClient!}
实例。是否能实际构造取决于下方可用性状态。

\textbf{签名}

\begin{lstlisting}[language=Python]
def __init__(self) -> None
\end{lstlisting}

\textbf{可用性与安全性}

\textbf{\passthrough{\lstinline!SIGNATURE\_ONLY!} /
\passthrough{\lstinline!CONSTRUCTION!}}：当前只有设计期签名，不能假设已存在于运行时扩展。

\textbf{参数}

无显式参数。实例方法中的 \passthrough{\lstinline!self!} 由 Python
自动传入。

\textbf{返回值}

\begin{longtable}[]{@{}
  >{\raggedright\arraybackslash}p{(\columnwidth - 4\tabcolsep) * \real{0.18}}
  >{\raggedright\arraybackslash}p{(\columnwidth - 4\tabcolsep) * \real{0.20}}
  >{\raggedright\arraybackslash}p{(\columnwidth - 4\tabcolsep) * \real{0.62}}@{}}
\toprule
位置 & 类型 & 含义 \\
\midrule
\endhead
无 & \passthrough{\lstinline!None!} &
\passthrough{\lstinline!\_\_init\_\_!}
本身不返回值；调用类对象时，在构造函数可用的前提下得到该类实例。 \\
\bottomrule
\end{longtable}

\textbf{对应 C++}

\begin{itemize}
\tightlist
\item
  类：\passthrough{\lstinline!unitree::robot::g1::AgvClient!}
\item
  签名：\passthrough{\lstinline!AgvClient()!}
\item
  绑定策略：\passthrough{\lstinline!未单独标注!}
\item
  声明位置：\passthrough{\lstinline!include/unitree/robot/g1/agv/g1\_agv\_client.hpp:15!}
\end{itemize}

\textbf{说明}

\begin{itemize}
\tightlist
\item
  示例是局部调用片段；\passthrough{\lstinline!obj!}
  和参数变量表示已经按上层业务校验并准备好的对象。
\item
  该条目可以被编辑器和类型检查器识别，但当前运行时可能在属性查找阶段直接失败。
\end{itemize}

\textbf{用法}

\begin{lstlisting}[language=Python]
from typing import TYPE_CHECKING

if TYPE_CHECKING:
    def planned_construct() -> AgvClient:
        return AgvClient()
\end{lstlisting}

\begin{quote}
{[}!CAUTION{]} 上面的 \passthrough{\lstinline!TYPE\_CHECKING!}
示例只表达计划调用形态。该分支在普通 Python
运行时为假，不代表当前扩展能执行此接口。
\end{quote}

\hypertarget{unitree-sdk2-cpp-robot-g1-agvclient-init-1}{%
\paragraph{\texorpdfstring{\texttt{unitree\_sdk2\_cpp.robot.g1.AgvClient.init}}{unitree\_sdk2\_cpp.robot.g1.AgvClient.init}}\label{unitree-sdk2-cpp-robot-g1-agvclient-init-1}}

初始化当前 SDK 对象所需的底层通道或服务资源。

\textbf{签名}

\begin{lstlisting}[language=Python]
def init(self) -> None
\end{lstlisting}

\textbf{可用性与安全性}

\textbf{\passthrough{\lstinline!SIGNATURE\_ONLY!} /
\passthrough{\lstinline!INITIALIZATION!}}：当前只有设计期签名，不能假设已存在于运行时扩展。

\textbf{参数}

无显式参数。实例方法中的 \passthrough{\lstinline!self!} 由 Python
自动传入。

\textbf{返回值}

\begin{longtable}[]{@{}
  >{\raggedright\arraybackslash}p{(\columnwidth - 4\tabcolsep) * \real{0.18}}
  >{\raggedright\arraybackslash}p{(\columnwidth - 4\tabcolsep) * \real{0.20}}
  >{\raggedright\arraybackslash}p{(\columnwidth - 4\tabcolsep) * \real{0.62}}@{}}
\toprule
位置 & 类型 & 含义 \\
\midrule
\endhead
无 & \passthrough{\lstinline!None!} & 该方法不返回 Python
值；失败可能表现为异常或底层状态变化。 \\
\bottomrule
\end{longtable}

\textbf{对应 C++}

\begin{itemize}
\tightlist
\item
  类：\passthrough{\lstinline!unitree::robot::g1::AgvClient!}
\item
  签名：\passthrough{\lstinline!Init()!}
\item
  绑定策略：\passthrough{\lstinline!DIRECT!}
\item
  声明位置：\passthrough{\lstinline!include/unitree/robot/g1/agv/g1\_agv\_client.hpp:19!}
\end{itemize}

\textbf{说明}

\begin{itemize}
\tightlist
\item
  示例是局部调用片段；\passthrough{\lstinline!obj!}
  和参数变量表示已经按上层业务校验并准备好的对象。
\item
  该条目可以被编辑器和类型检查器识别，但当前运行时可能在属性查找阶段直接失败。
\end{itemize}

\textbf{用法}

\begin{lstlisting}[language=Python]
from typing import TYPE_CHECKING

if TYPE_CHECKING:
    def planned_init(obj: AgvClient) -> None:
        obj.init()
\end{lstlisting}

\begin{quote}
{[}!CAUTION{]} 上面的 \passthrough{\lstinline!TYPE\_CHECKING!}
示例只表达计划调用形态。该分支在普通 Python
运行时为假，不代表当前扩展能执行此接口。
\end{quote}

\hypertarget{unitree-sdk2-cpp-robot-g1-agvclient-move-1}{%
\paragraph{\texorpdfstring{\texttt{unitree\_sdk2\_cpp.robot.g1.AgvClient.move}}{unitree\_sdk2\_cpp.robot.g1.AgvClient.move}}\label{unitree-sdk2-cpp-robot-g1-agvclient-move-1}}

对应 C++ SDK 操作
\passthrough{\lstinline!Move(float, float, float)!}。上游头文件没有可直接生成的业务说明时，本参考只保证签名映射，精确语义需查目标型号协议。

\textbf{签名}

\begin{lstlisting}[language=Python]
def move(self, vx: float, vy: float, vyaw: float) -> int
\end{lstlisting}

\textbf{可用性与安全性}

\textbf{\passthrough{\lstinline!SIGNATURE\_ONLY!} /
\passthrough{\lstinline!MOTION\_COMMAND!}}：仅用于补全和静态检查。当前二进制没有该
Python 方法；不得按可执行运动接口使用。

\textbf{参数}

\begin{longtable}[]{@{}
  >{\raggedright\arraybackslash}p{(\columnwidth - 6\tabcolsep) * \real{0.18}}
  >{\raggedright\arraybackslash}p{(\columnwidth - 6\tabcolsep) * \real{0.24}}
  >{\raggedright\arraybackslash}p{(\columnwidth - 6\tabcolsep) * \real{0.18}}
  >{\raggedright\arraybackslash}p{(\columnwidth - 6\tabcolsep) * \real{0.40}}@{}}
\toprule
名称 & Python 类型 & 默认值 & 含义 \\
\midrule
\endhead
\passthrough{\lstinline!vx!} & \passthrough{\lstinline!float!} & 必填 &
X 方向速度参数。坐标系、单位、符号和安全范围必须查目标型号运动协议。
对应 C++ 参数 \passthrough{\lstinline!vx: float!}。
底层为浮点数；类型本身不说明单位、坐标系或业务安全范围。 \\
\passthrough{\lstinline!vy!} & \passthrough{\lstinline!float!} & 必填 &
Y 方向速度参数。坐标系、单位、符号和安全范围必须查目标型号运动协议。
对应 C++ 参数 \passthrough{\lstinline!vy: float!}。
底层为浮点数；类型本身不说明单位、坐标系或业务安全范围。 \\
\passthrough{\lstinline!vyaw!} & \passthrough{\lstinline!float!} & 必填
& 偏航角速度参数。单位、符号和安全范围必须查目标型号运动协议。 对应 C++
参数 \passthrough{\lstinline!vyaw: float!}。
底层为浮点数；类型本身不说明单位、坐标系或业务安全范围。 \\
\bottomrule
\end{longtable}

\textbf{返回值}

\begin{longtable}[]{@{}
  >{\raggedright\arraybackslash}p{(\columnwidth - 4\tabcolsep) * \real{0.18}}
  >{\raggedright\arraybackslash}p{(\columnwidth - 4\tabcolsep) * \real{0.20}}
  >{\raggedright\arraybackslash}p{(\columnwidth - 4\tabcolsep) * \real{0.62}}@{}}
\toprule
位置 & 类型 & 含义 \\
\midrule
\endhead
返回值 & \passthrough{\lstinline!int!} & SDK 状态码；Unitree 示例通常以
\passthrough{\lstinline!0!} 表示成功，非零值需按具体服务定义解释。 \\
\bottomrule
\end{longtable}

\textbf{对应 C++}

\begin{itemize}
\tightlist
\item
  类：\passthrough{\lstinline!unitree::robot::g1::AgvClient!}
\item
  签名：\passthrough{\lstinline!Move(float, float, float)!}
\item
  绑定策略：\passthrough{\lstinline!DIRECT!}
\item
  声明位置：\passthrough{\lstinline!include/unitree/robot/g1/agv/g1\_agv\_client.hpp:56!}
\end{itemize}

\textbf{说明}

\begin{itemize}
\tightlist
\item
  示例是局部调用片段；\passthrough{\lstinline!obj!}
  和参数变量表示已经按上层业务校验并准备好的对象。
\item
  该条目可以被编辑器和类型检查器识别，但当前运行时可能在属性查找阶段直接失败。
\item
  该方法属于运动命令。方法名包含
  \passthrough{\lstinline!stop!}、\passthrough{\lstinline!damp!} 或
  \passthrough{\lstinline!zero!} 也不代表它是独立物理急停。
\end{itemize}

\textbf{用法}

\begin{lstlisting}[language=Python]
from typing import TYPE_CHECKING

if TYPE_CHECKING:
    def planned_move(obj: AgvClient, vx: float, vy: float, vyaw: float) -> int:
        return obj.move(vx=vx, vy=vy, vyaw=vyaw)
\end{lstlisting}

\begin{quote}
{[}!CAUTION{]} 上面的 \passthrough{\lstinline!TYPE\_CHECKING!}
示例只表达计划调用形态。该分支在普通 Python
运行时为假，不代表当前扩展能执行此接口。
\end{quote}

\hypertarget{unitree-sdk2-cpp-robot-g1-agvclient-height-adjust-1}{%
\paragraph{\texorpdfstring{\texttt{unitree\_sdk2\_cpp.robot.g1.AgvClient.height\_adjust}}{unitree\_sdk2\_cpp.robot.g1.AgvClient.height\_adjust}}\label{unitree-sdk2-cpp-robot-g1-agvclient-height-adjust-1}}

对应 C++ SDK 操作
\passthrough{\lstinline!HeightAdjust(float)!}。上游头文件没有可直接生成的业务说明时，本参考只保证签名映射，精确语义需查目标型号协议。

\textbf{签名}

\begin{lstlisting}[language=Python]
def height_adjust(self, vz: float) -> int
\end{lstlisting}

\textbf{可用性与安全性}

\textbf{\passthrough{\lstinline!SIGNATURE\_ONLY!} /
\passthrough{\lstinline!MOTION\_COMMAND!}}：仅用于补全和静态检查。当前二进制没有该
Python 方法；不得按可执行运动接口使用。

\textbf{参数}

\begin{longtable}[]{@{}
  >{\raggedright\arraybackslash}p{(\columnwidth - 6\tabcolsep) * \real{0.18}}
  >{\raggedright\arraybackslash}p{(\columnwidth - 6\tabcolsep) * \real{0.24}}
  >{\raggedright\arraybackslash}p{(\columnwidth - 6\tabcolsep) * \real{0.18}}
  >{\raggedright\arraybackslash}p{(\columnwidth - 6\tabcolsep) * \real{0.40}}@{}}
\toprule
名称 & Python 类型 & 默认值 & 含义 \\
\midrule
\endhead
\passthrough{\lstinline!vz!} & \passthrough{\lstinline!float!} & 必填 &
传给该接口的 \passthrough{\lstinline!vz!} 参数，Python 类型为
\passthrough{\lstinline!float!}。精确含义、单位、范围和枚举值需查目标型号对应头文件与协议。
对应 C++ 参数 \passthrough{\lstinline!vz: float!}。
底层为浮点数；类型本身不说明单位、坐标系或业务安全范围。 \\
\bottomrule
\end{longtable}

\textbf{返回值}

\begin{longtable}[]{@{}
  >{\raggedright\arraybackslash}p{(\columnwidth - 4\tabcolsep) * \real{0.18}}
  >{\raggedright\arraybackslash}p{(\columnwidth - 4\tabcolsep) * \real{0.20}}
  >{\raggedright\arraybackslash}p{(\columnwidth - 4\tabcolsep) * \real{0.62}}@{}}
\toprule
位置 & 类型 & 含义 \\
\midrule
\endhead
返回值 & \passthrough{\lstinline!int!} & SDK 状态码；Unitree 示例通常以
\passthrough{\lstinline!0!} 表示成功，非零值需按具体服务定义解释。 \\
\bottomrule
\end{longtable}

\textbf{对应 C++}

\begin{itemize}
\tightlist
\item
  类：\passthrough{\lstinline!unitree::robot::g1::AgvClient!}
\item
  签名：\passthrough{\lstinline!HeightAdjust(float)!}
\item
  绑定策略：\passthrough{\lstinline!DIRECT!}
\item
  声明位置：\passthrough{\lstinline!include/unitree/robot/g1/agv/g1\_agv\_client.hpp:93!}
\end{itemize}

\textbf{说明}

\begin{itemize}
\tightlist
\item
  示例是局部调用片段；\passthrough{\lstinline!obj!}
  和参数变量表示已经按上层业务校验并准备好的对象。
\item
  该条目可以被编辑器和类型检查器识别，但当前运行时可能在属性查找阶段直接失败。
\item
  该方法属于运动命令。方法名包含
  \passthrough{\lstinline!stop!}、\passthrough{\lstinline!damp!} 或
  \passthrough{\lstinline!zero!} 也不代表它是独立物理急停。
\end{itemize}

\textbf{用法}

\begin{lstlisting}[language=Python]
from typing import TYPE_CHECKING

if TYPE_CHECKING:
    def planned_height_adjust(obj: AgvClient, vz: float) -> int:
        return obj.height_adjust(vz=vz)
\end{lstlisting}

\begin{quote}
{[}!CAUTION{]} 上面的 \passthrough{\lstinline!TYPE\_CHECKING!}
示例只表达计划调用形态。该分支在普通 Python
运行时为假，不代表当前扩展能执行此接口。
\end{quote}

\hypertarget{unitree-sdk2-cpp-robot-g1-audioclient}{%
\subsubsection{\texorpdfstring{\texttt{unitree\_sdk2\_cpp.robot.g1.AudioClient}}{unitree\_sdk2\_cpp.robot.g1.AudioClient}}\label{unitree-sdk2-cpp-robot-g1-audioclient}}

Robot 服务客户端。构造、初始化和具体方法的可用性必须分别检查。

\textbf{导入}

\begin{lstlisting}[language=Python]
from unitree_sdk2_cpp.robot.g1 import AudioClient
\end{lstlisting}

\textbf{基类}：\passthrough{\lstinline!Client!}

\textbf{方法索引}

\begin{longtable}[]{@{}
  >{\raggedright\arraybackslash}p{(\columnwidth - 6\tabcolsep) * \real{0.18}}
  >{\raggedright\arraybackslash}p{(\columnwidth - 6\tabcolsep) * \real{0.24}}
  >{\raggedright\arraybackslash}p{(\columnwidth - 6\tabcolsep) * \real{0.18}}
  >{\raggedright\arraybackslash}p{(\columnwidth - 6\tabcolsep) * \real{0.40}}@{}}
\toprule
方法 & Python 签名 & 状态 & 安全分类 \\
\midrule
\endhead
\protect\hyperlink{unitree-sdk2-cpp-robot-g1-audioclient-dunder-init-1}{\passthrough{\lstinline!\_\_init\_\_!}}
& \passthrough{\lstinline!def \_\_init\_\_(self) -> None!} &
\passthrough{\lstinline!AVAILABLE!} &
\passthrough{\lstinline!CONSTRUCTION!} \\
\protect\hyperlink{unitree-sdk2-cpp-robot-g1-audioclient-init-1}{\passthrough{\lstinline!init!}}
& \passthrough{\lstinline!def init(self) -> None!} &
\passthrough{\lstinline!AVAILABLE!} &
\passthrough{\lstinline!INITIALIZATION!} \\
\protect\hyperlink{unitree-sdk2-cpp-robot-g1-audioclient-tts-maker-1}{\passthrough{\lstinline!tts\_maker!}}
&
\passthrough{\lstinline!def tts\_maker(self, text: str, speaker\_id: int) -> int!}
& \passthrough{\lstinline!SIGNATURE\_ONLY!} &
\passthrough{\lstinline!HARDWARE\_SIDE\_EFFECT!} \\
\protect\hyperlink{unitree-sdk2-cpp-robot-g1-audioclient-get-volume-1}{\passthrough{\lstinline!get\_volume!}}
& \passthrough{\lstinline!def get\_volume(self) -> tuple[int, int]!} &
\passthrough{\lstinline!AVAILABLE!} &
\passthrough{\lstinline!READ\_ONLY!} \\
\protect\hyperlink{unitree-sdk2-cpp-robot-g1-audioclient-set-volume-1}{\passthrough{\lstinline!set\_volume!}}
& \passthrough{\lstinline!def set\_volume(self, volume: int) -> int!} &
\passthrough{\lstinline!SIGNATURE\_ONLY!} &
\passthrough{\lstinline!HARDWARE\_SIDE\_EFFECT!} \\
\protect\hyperlink{unitree-sdk2-cpp-robot-g1-audioclient-play-stream-1}{\passthrough{\lstinline!play\_stream!}}
&
\passthrough{\lstinline!def play\_stream(self, app\_name: str, stream\_id: str, pcm\_data: Sequence[int]) -> int!}
& \passthrough{\lstinline!SIGNATURE\_ONLY!} &
\passthrough{\lstinline!HARDWARE\_SIDE\_EFFECT!} \\
\protect\hyperlink{unitree-sdk2-cpp-robot-g1-audioclient-play-stop-1}{\passthrough{\lstinline!play\_stop!}}
& \passthrough{\lstinline!def play\_stop(self, app\_name: str) -> int!}
& \passthrough{\lstinline!SIGNATURE\_ONLY!} &
\passthrough{\lstinline!HARDWARE\_SIDE\_EFFECT!} \\
\protect\hyperlink{unitree-sdk2-cpp-robot-g1-audioclient-led-control-1}{\passthrough{\lstinline!led\_control!}}
&
\passthrough{\lstinline!def led\_control(self, r: int, g: int, b: int) -> int!}
& \passthrough{\lstinline!SIGNATURE\_ONLY!} &
\passthrough{\lstinline!HARDWARE\_SIDE\_EFFECT!} \\
\bottomrule
\end{longtable}

\hypertarget{unitree-sdk2-cpp-robot-g1-audioclient-dunder-init-1}{%
\paragraph{\texorpdfstring{\texttt{unitree\_sdk2\_cpp.robot.g1.AudioClient.\_\_init\_\_}}{unitree\_sdk2\_cpp.robot.g1.AudioClient.\_\_init\_\_}}\label{unitree-sdk2-cpp-robot-g1-audioclient-dunder-init-1}}

初始化 \passthrough{\lstinline!AudioClient!}
实例。是否能实际构造取决于下方可用性状态。

\textbf{签名}

\begin{lstlisting}[language=Python]
def __init__(self) -> None
\end{lstlisting}

\textbf{可用性与安全性}

\textbf{\passthrough{\lstinline!AVAILABLE!} /
\passthrough{\lstinline!CONSTRUCTION!}}：当前绑定源码已实现；实际执行条件仍取决于平台和对应
SDK 环境。

\textbf{参数}

无显式参数。实例方法中的 \passthrough{\lstinline!self!} 由 Python
自动传入。

\textbf{返回值}

\begin{longtable}[]{@{}
  >{\raggedright\arraybackslash}p{(\columnwidth - 4\tabcolsep) * \real{0.18}}
  >{\raggedright\arraybackslash}p{(\columnwidth - 4\tabcolsep) * \real{0.20}}
  >{\raggedright\arraybackslash}p{(\columnwidth - 4\tabcolsep) * \real{0.62}}@{}}
\toprule
位置 & 类型 & 含义 \\
\midrule
\endhead
无 & \passthrough{\lstinline!None!} &
\passthrough{\lstinline!\_\_init\_\_!}
本身不返回值；调用类对象时，在构造函数可用的前提下得到该类实例。 \\
\bottomrule
\end{longtable}

\textbf{对应 C++}

\begin{itemize}
\tightlist
\item
  类：\passthrough{\lstinline!unitree::robot::g1::AudioClient!}
\item
  签名：\passthrough{\lstinline!AudioClient()!}
\item
  绑定策略：\passthrough{\lstinline!未单独标注!}
\item
  声明位置：\passthrough{\lstinline!include/unitree/robot/g1/audio/g1\_audio\_client.hpp:15!}
\end{itemize}

\textbf{说明}

\begin{itemize}
\tightlist
\item
  示例是局部调用片段；\passthrough{\lstinline!obj!}
  和参数变量表示已经按上层业务校验并准备好的对象。
\end{itemize}

\textbf{用法}

\begin{lstlisting}[language=Python]
value = AudioClient()
\end{lstlisting}

\hypertarget{unitree-sdk2-cpp-robot-g1-audioclient-init-1}{%
\paragraph{\texorpdfstring{\texttt{unitree\_sdk2\_cpp.robot.g1.AudioClient.init}}{unitree\_sdk2\_cpp.robot.g1.AudioClient.init}}\label{unitree-sdk2-cpp-robot-g1-audioclient-init-1}}

初始化当前 SDK 对象所需的底层通道或服务资源。

\textbf{签名}

\begin{lstlisting}[language=Python]
def init(self) -> None
\end{lstlisting}

\textbf{可用性与安全性}

\textbf{\passthrough{\lstinline!AVAILABLE!} /
\passthrough{\lstinline!INITIALIZATION!}}：当前绑定源码已实现；实际执行条件仍取决于平台和对应
SDK 环境。

\textbf{参数}

无显式参数。实例方法中的 \passthrough{\lstinline!self!} 由 Python
自动传入。

\textbf{返回值}

\begin{longtable}[]{@{}
  >{\raggedright\arraybackslash}p{(\columnwidth - 4\tabcolsep) * \real{0.18}}
  >{\raggedright\arraybackslash}p{(\columnwidth - 4\tabcolsep) * \real{0.20}}
  >{\raggedright\arraybackslash}p{(\columnwidth - 4\tabcolsep) * \real{0.62}}@{}}
\toprule
位置 & 类型 & 含义 \\
\midrule
\endhead
无 & \passthrough{\lstinline!None!} & 该方法不返回 Python
值；失败可能表现为异常或底层状态变化。 \\
\bottomrule
\end{longtable}

\textbf{对应 C++}

\begin{itemize}
\tightlist
\item
  类：\passthrough{\lstinline!unitree::robot::g1::AudioClient!}
\item
  签名：\passthrough{\lstinline!Init()!}
\item
  绑定策略：\passthrough{\lstinline!DIRECT!}
\item
  声明位置：\passthrough{\lstinline!include/unitree/robot/g1/audio/g1\_audio\_client.hpp:19!}
\end{itemize}

\textbf{说明}

\begin{itemize}
\tightlist
\item
  示例是局部调用片段；\passthrough{\lstinline!obj!}
  和参数变量表示已经按上层业务校验并准备好的对象。
\end{itemize}

\textbf{用法}

\begin{lstlisting}[language=Python]
obj.init()
\end{lstlisting}

\hypertarget{unitree-sdk2-cpp-robot-g1-audioclient-tts-maker-1}{%
\paragraph{\texorpdfstring{\texttt{unitree\_sdk2\_cpp.robot.g1.AudioClient.tts\_maker}}{unitree\_sdk2\_cpp.robot.g1.AudioClient.tts\_maker}}\label{unitree-sdk2-cpp-robot-g1-audioclient-tts-maker-1}}

对应 C++ SDK 操作
\passthrough{\lstinline!TtsMaker(const std::string \&, int32\_t)!}。上游头文件没有可直接生成的业务说明时，本参考只保证签名映射，精确语义需查目标型号协议。

\textbf{签名}

\begin{lstlisting}[language=Python]
def tts_maker(self, text: str, speaker_id: int) -> int
\end{lstlisting}

\textbf{可用性与安全性}

\textbf{\passthrough{\lstinline!SIGNATURE\_ONLY!} /
\passthrough{\lstinline!HARDWARE\_SIDE\_EFFECT!}}：仅用于补全和静态检查。当前二进制没有该
Python 方法，未来实现还需单独评审硬件副作用。

\textbf{参数}

\begin{longtable}[]{@{}
  >{\raggedright\arraybackslash}p{(\columnwidth - 6\tabcolsep) * \real{0.18}}
  >{\raggedright\arraybackslash}p{(\columnwidth - 6\tabcolsep) * \real{0.24}}
  >{\raggedright\arraybackslash}p{(\columnwidth - 6\tabcolsep) * \real{0.18}}
  >{\raggedright\arraybackslash}p{(\columnwidth - 6\tabcolsep) * \real{0.40}}@{}}
\toprule
名称 & Python 类型 & 默认值 & 含义 \\
\midrule
\endhead
\passthrough{\lstinline!text!} & \passthrough{\lstinline!str!} & 必填 &
要处理的文本内容；编码、长度和语言支持由目标服务决定。 对应 C++ 参数
\passthrough{\lstinline!text: const std::string \&!}。
底层为字符串；长度、编码和允许值由具体协议决定。 \\
\passthrough{\lstinline!speaker\_id!} & \passthrough{\lstinline!int!} &
必填 & 语音合成说话人 ID。有效编号由机器人音频服务和固件决定。 对应 C++
参数 \passthrough{\lstinline!speaker\_id: int32\_t!}。 取值范围为
-2147483648 到 2147483647。 \\
\bottomrule
\end{longtable}

\textbf{返回值}

\begin{longtable}[]{@{}
  >{\raggedright\arraybackslash}p{(\columnwidth - 4\tabcolsep) * \real{0.18}}
  >{\raggedright\arraybackslash}p{(\columnwidth - 4\tabcolsep) * \real{0.20}}
  >{\raggedright\arraybackslash}p{(\columnwidth - 4\tabcolsep) * \real{0.62}}@{}}
\toprule
位置 & 类型 & 含义 \\
\midrule
\endhead
返回值 & \passthrough{\lstinline!int!} & SDK 状态码；Unitree 示例通常以
\passthrough{\lstinline!0!} 表示成功，非零值需按具体服务定义解释。 \\
\bottomrule
\end{longtable}

\textbf{对应 C++}

\begin{itemize}
\tightlist
\item
  类：\passthrough{\lstinline!unitree::robot::g1::AudioClient!}
\item
  签名：\passthrough{\lstinline!TtsMaker(const std::string \&, int32\_t)!}
\item
  绑定策略：\passthrough{\lstinline!DIRECT!}
\item
  声明位置：\passthrough{\lstinline!include/unitree/robot/g1/audio/g1\_audio\_client.hpp:31!}
\end{itemize}

\textbf{说明}

\begin{itemize}
\tightlist
\item
  示例是局部调用片段；\passthrough{\lstinline!obj!}
  和参数变量表示已经按上层业务校验并准备好的对象。
\item
  该条目可以被编辑器和类型检查器识别，但当前运行时可能在属性查找阶段直接失败。
\end{itemize}

\textbf{用法}

\begin{lstlisting}[language=Python]
from typing import TYPE_CHECKING

if TYPE_CHECKING:
    def planned_tts_maker(obj: AudioClient, text: str, speaker_id: int) -> int:
        return obj.tts_maker(text=text, speaker_id=speaker_id)
\end{lstlisting}

\begin{quote}
{[}!CAUTION{]} 上面的 \passthrough{\lstinline!TYPE\_CHECKING!}
示例只表达计划调用形态。该分支在普通 Python
运行时为假，不代表当前扩展能执行此接口。
\end{quote}

\hypertarget{unitree-sdk2-cpp-robot-g1-audioclient-get-volume-1}{%
\paragraph{\texorpdfstring{\texttt{unitree\_sdk2\_cpp.robot.g1.AudioClient.get\_volume}}{unitree\_sdk2\_cpp.robot.g1.AudioClient.get\_volume}}\label{unitree-sdk2-cpp-robot-g1-audioclient-get-volume-1}}

查询或检查 音量。

\textbf{签名}

\begin{lstlisting}[language=Python]
def get_volume(self) -> tuple[int, int]
\end{lstlisting}

\textbf{可用性与安全性}

\textbf{\passthrough{\lstinline!AVAILABLE!} /
\passthrough{\lstinline!READ\_ONLY!}}：当前绑定源码已实现只读查询。Robot
Client 的服务端查询通常仍需匹配的 Linux
扩展、DDS、网络、机器人和服务；纯本地 getter 则不一定需要全部条件。

\textbf{参数}

无显式参数。实例方法中的 \passthrough{\lstinline!self!} 由 Python
自动传入。

\textbf{返回值}

\begin{longtable}[]{@{}
  >{\raggedright\arraybackslash}p{(\columnwidth - 4\tabcolsep) * \real{0.18}}
  >{\raggedright\arraybackslash}p{(\columnwidth - 4\tabcolsep) * \real{0.20}}
  >{\raggedright\arraybackslash}p{(\columnwidth - 4\tabcolsep) * \real{0.62}}@{}}
\toprule
位置 & 类型 & 含义 \\
\midrule
\endhead
{[}0{]} \passthrough{\lstinline!status!} & \passthrough{\lstinline!int!}
& SDK 状态码；Unitree 示例通常以 \passthrough{\lstinline!0!}
表示成功，非零值需按具体服务错误码解释。 \\
{[}1{]} \passthrough{\lstinline!volume!} & \passthrough{\lstinline!int!}
& 传给该接口的 音量 参数，Python 类型为
\passthrough{\lstinline!int!}。精确含义、单位、范围和枚举值需查目标型号对应头文件与协议。
对应 C++ 参数 \passthrough{\lstinline!volume: uint8\_t \&!}。 取值范围为
0 到 255。 \\
\bottomrule
\end{longtable}

\textbf{对应 C++}

\begin{itemize}
\tightlist
\item
  类：\passthrough{\lstinline!unitree::robot::g1::AudioClient!}
\item
  签名：\passthrough{\lstinline!GetVolume(uint8\_t \&)!}
\item
  绑定策略：\passthrough{\lstinline!OUTPUT\_WRAPPER!}
\item
  声明位置：\passthrough{\lstinline!include/unitree/robot/g1/audio/g1\_audio\_client.hpp:43!}
\end{itemize}

\textbf{说明}

\begin{itemize}
\tightlist
\item
  示例是局部调用片段；\passthrough{\lstinline!obj!}
  和参数变量表示已经按上层业务校验并准备好的对象。
\item
  C++ 的可修改输出引用不会作为 Python 入参出现，而是按 C++
  参数顺序追加到返回元组中。
\end{itemize}

\textbf{用法}

\begin{lstlisting}[language=Python]
status, volume = obj.get_volume()
\end{lstlisting}

\begin{quote}
{[}!NOTE{]} 只读查询仍会访问
DDS/机器人服务。必须检查返回状态码，并处理超时和网络中断。
\end{quote}

\hypertarget{unitree-sdk2-cpp-robot-g1-audioclient-set-volume-1}{%
\paragraph{\texorpdfstring{\texttt{unitree\_sdk2\_cpp.robot.g1.AudioClient.set\_volume}}{unitree\_sdk2\_cpp.robot.g1.AudioClient.set\_volume}}\label{unitree-sdk2-cpp-robot-g1-audioclient-set-volume-1}}

设置 音量。具体副作用和安全边界见下方状态。

\textbf{签名}

\begin{lstlisting}[language=Python]
def set_volume(self, volume: int) -> int
\end{lstlisting}

\textbf{可用性与安全性}

\textbf{\passthrough{\lstinline!SIGNATURE\_ONLY!} /
\passthrough{\lstinline!HARDWARE\_SIDE\_EFFECT!}}：仅用于补全和静态检查。当前二进制没有该
Python 方法，未来实现还需单独评审硬件副作用。

\textbf{参数}

\begin{longtable}[]{@{}
  >{\raggedright\arraybackslash}p{(\columnwidth - 6\tabcolsep) * \real{0.18}}
  >{\raggedright\arraybackslash}p{(\columnwidth - 6\tabcolsep) * \real{0.24}}
  >{\raggedright\arraybackslash}p{(\columnwidth - 6\tabcolsep) * \real{0.18}}
  >{\raggedright\arraybackslash}p{(\columnwidth - 6\tabcolsep) * \real{0.40}}@{}}
\toprule
名称 & Python 类型 & 默认值 & 含义 \\
\midrule
\endhead
\passthrough{\lstinline!volume!} & \passthrough{\lstinline!int!} & 必填
& 传给该接口的 音量 参数，Python 类型为
\passthrough{\lstinline!int!}。精确含义、单位、范围和枚举值需查目标型号对应头文件与协议。
对应 C++ 参数 \passthrough{\lstinline!volume: uint8\_t!}。 取值范围为 0
到 255。 \\
\bottomrule
\end{longtable}

\textbf{返回值}

\begin{longtable}[]{@{}
  >{\raggedright\arraybackslash}p{(\columnwidth - 4\tabcolsep) * \real{0.18}}
  >{\raggedright\arraybackslash}p{(\columnwidth - 4\tabcolsep) * \real{0.20}}
  >{\raggedright\arraybackslash}p{(\columnwidth - 4\tabcolsep) * \real{0.62}}@{}}
\toprule
位置 & 类型 & 含义 \\
\midrule
\endhead
返回值 & \passthrough{\lstinline!int!} & SDK 状态码；Unitree 示例通常以
\passthrough{\lstinline!0!} 表示成功，非零值需按具体服务定义解释。 \\
\bottomrule
\end{longtable}

\textbf{对应 C++}

\begin{itemize}
\tightlist
\item
  类：\passthrough{\lstinline!unitree::robot::g1::AudioClient!}
\item
  签名：\passthrough{\lstinline!SetVolume(uint8\_t)!}
\item
  绑定策略：\passthrough{\lstinline!DIRECT!}
\item
  声明位置：\passthrough{\lstinline!include/unitree/robot/g1/audio/g1\_audio\_client.hpp:57!}
\end{itemize}

\textbf{说明}

\begin{itemize}
\tightlist
\item
  示例是局部调用片段；\passthrough{\lstinline!obj!}
  和参数变量表示已经按上层业务校验并准备好的对象。
\item
  该条目可以被编辑器和类型检查器识别，但当前运行时可能在属性查找阶段直接失败。
\end{itemize}

\textbf{用法}

\begin{lstlisting}[language=Python]
from typing import TYPE_CHECKING

if TYPE_CHECKING:
    def planned_set_volume(obj: AudioClient, volume: int) -> int:
        return obj.set_volume(volume=volume)
\end{lstlisting}

\begin{quote}
{[}!CAUTION{]} 上面的 \passthrough{\lstinline!TYPE\_CHECKING!}
示例只表达计划调用形态。该分支在普通 Python
运行时为假，不代表当前扩展能执行此接口。
\end{quote}

\hypertarget{unitree-sdk2-cpp-robot-g1-audioclient-play-stream-1}{%
\paragraph{\texorpdfstring{\texttt{unitree\_sdk2\_cpp.robot.g1.AudioClient.play\_stream}}{unitree\_sdk2\_cpp.robot.g1.AudioClient.play\_stream}}\label{unitree-sdk2-cpp-robot-g1-audioclient-play-stream-1}}

对应 C++ SDK 操作
\passthrough{\lstinline!PlayStream(std::string, std::string, std::vector<uint8\_t>)!}。上游头文件没有可直接生成的业务说明时，本参考只保证签名映射，精确语义需查目标型号协议。

\textbf{签名}

\begin{lstlisting}[language=Python]
def play_stream(self, app_name: str, stream_id: str, pcm_data: Sequence[int]) -> int
\end{lstlisting}

\textbf{可用性与安全性}

\textbf{\passthrough{\lstinline!SIGNATURE\_ONLY!} /
\passthrough{\lstinline!HARDWARE\_SIDE\_EFFECT!}}：仅用于补全和静态检查。当前二进制没有该
Python 方法，未来实现还需单独评审硬件副作用。

\textbf{参数}

\begin{longtable}[]{@{}
  >{\raggedright\arraybackslash}p{(\columnwidth - 6\tabcolsep) * \real{0.18}}
  >{\raggedright\arraybackslash}p{(\columnwidth - 6\tabcolsep) * \real{0.24}}
  >{\raggedright\arraybackslash}p{(\columnwidth - 6\tabcolsep) * \real{0.18}}
  >{\raggedright\arraybackslash}p{(\columnwidth - 6\tabcolsep) * \real{0.40}}@{}}
\toprule
名称 & Python 类型 & 默认值 & 含义 \\
\midrule
\endhead
\passthrough{\lstinline!app\_name!} & \passthrough{\lstinline!str!} &
必填 &
应用名称，用于标识音频流或播放会话。允许值和命名规则以目标服务协议为准。
对应 C++ 参数 \passthrough{\lstinline!app\_name: std::string!}。
底层为字符串；长度、编码和允许值由具体协议决定。 \\
\passthrough{\lstinline!stream\_id!} & \passthrough{\lstinline!str!} &
必填 & 音频流标识符，用于区分同一应用下的播放流。 对应 C++ 参数
\passthrough{\lstinline!stream\_id: std::string!}。
底层为字符串；长度、编码和允许值由具体协议决定。 \\
\passthrough{\lstinline!pcm\_data!} &
\passthrough{\lstinline!Sequence[int]!} & 必填 & 传给该接口的
\passthrough{\lstinline!pcm!} 数据 参数，Python 类型为
\passthrough{\lstinline!Sequence[int]!}。精确含义、单位、范围和枚举值需查目标型号对应头文件与协议。
对应 C++ 参数
\passthrough{\lstinline!pcm\_data: std::vector<uint8\_t>!}。
底层是可变长度 vector；元素约束：取值范围为 0 到 255。 \\
\bottomrule
\end{longtable}

\textbf{返回值}

\begin{longtable}[]{@{}
  >{\raggedright\arraybackslash}p{(\columnwidth - 4\tabcolsep) * \real{0.18}}
  >{\raggedright\arraybackslash}p{(\columnwidth - 4\tabcolsep) * \real{0.20}}
  >{\raggedright\arraybackslash}p{(\columnwidth - 4\tabcolsep) * \real{0.62}}@{}}
\toprule
位置 & 类型 & 含义 \\
\midrule
\endhead
返回值 & \passthrough{\lstinline!int!} & SDK 状态码；Unitree 示例通常以
\passthrough{\lstinline!0!} 表示成功，非零值需按具体服务定义解释。 \\
\bottomrule
\end{longtable}

\textbf{对应 C++}

\begin{itemize}
\tightlist
\item
  类：\passthrough{\lstinline!unitree::robot::g1::AudioClient!}
\item
  签名：\passthrough{\lstinline!PlayStream(std::string, std::string, std::vector<uint8\_t>)!}
\item
  绑定策略：\passthrough{\lstinline!DIRECT!}
\item
  声明位置：\passthrough{\lstinline!include/unitree/robot/g1/audio/g1\_audio\_client.hpp:68!}
\end{itemize}

\textbf{说明}

\begin{itemize}
\tightlist
\item
  示例是局部调用片段；\passthrough{\lstinline!obj!}
  和参数变量表示已经按上层业务校验并准备好的对象。
\item
  该条目可以被编辑器和类型检查器识别，但当前运行时可能在属性查找阶段直接失败。
\end{itemize}

\textbf{用法}

\begin{lstlisting}[language=Python]
from typing import TYPE_CHECKING

if TYPE_CHECKING:
    def planned_play_stream(obj: AudioClient, app_name: str, stream_id: str, pcm_data: Sequence[int]) -> int:
        return obj.play_stream(app_name=app_name, stream_id=stream_id, pcm_data=pcm_data)
\end{lstlisting}

\begin{quote}
{[}!CAUTION{]} 上面的 \passthrough{\lstinline!TYPE\_CHECKING!}
示例只表达计划调用形态。该分支在普通 Python
运行时为假，不代表当前扩展能执行此接口。
\end{quote}

\hypertarget{unitree-sdk2-cpp-robot-g1-audioclient-play-stop-1}{%
\paragraph{\texorpdfstring{\texttt{unitree\_sdk2\_cpp.robot.g1.AudioClient.play\_stop}}{unitree\_sdk2\_cpp.robot.g1.AudioClient.play\_stop}}\label{unitree-sdk2-cpp-robot-g1-audioclient-play-stop-1}}

对应 C++ SDK 操作
\passthrough{\lstinline!PlayStop(std::string)!}。上游头文件没有可直接生成的业务说明时，本参考只保证签名映射，精确语义需查目标型号协议。

\textbf{签名}

\begin{lstlisting}[language=Python]
def play_stop(self, app_name: str) -> int
\end{lstlisting}

\textbf{可用性与安全性}

\textbf{\passthrough{\lstinline!SIGNATURE\_ONLY!} /
\passthrough{\lstinline!HARDWARE\_SIDE\_EFFECT!}}：仅用于补全和静态检查。当前二进制没有该
Python 方法，未来实现还需单独评审硬件副作用。

\textbf{参数}

\begin{longtable}[]{@{}
  >{\raggedright\arraybackslash}p{(\columnwidth - 6\tabcolsep) * \real{0.18}}
  >{\raggedright\arraybackslash}p{(\columnwidth - 6\tabcolsep) * \real{0.24}}
  >{\raggedright\arraybackslash}p{(\columnwidth - 6\tabcolsep) * \real{0.18}}
  >{\raggedright\arraybackslash}p{(\columnwidth - 6\tabcolsep) * \real{0.40}}@{}}
\toprule
名称 & Python 类型 & 默认值 & 含义 \\
\midrule
\endhead
\passthrough{\lstinline!app\_name!} & \passthrough{\lstinline!str!} &
必填 &
应用名称，用于标识音频流或播放会话。允许值和命名规则以目标服务协议为准。
对应 C++ 参数 \passthrough{\lstinline!app\_name: std::string!}。
底层为字符串；长度、编码和允许值由具体协议决定。 \\
\bottomrule
\end{longtable}

\textbf{返回值}

\begin{longtable}[]{@{}
  >{\raggedright\arraybackslash}p{(\columnwidth - 4\tabcolsep) * \real{0.18}}
  >{\raggedright\arraybackslash}p{(\columnwidth - 4\tabcolsep) * \real{0.20}}
  >{\raggedright\arraybackslash}p{(\columnwidth - 4\tabcolsep) * \real{0.62}}@{}}
\toprule
位置 & 类型 & 含义 \\
\midrule
\endhead
返回值 & \passthrough{\lstinline!int!} & SDK 状态码；Unitree 示例通常以
\passthrough{\lstinline!0!} 表示成功，非零值需按具体服务定义解释。 \\
\bottomrule
\end{longtable}

\textbf{对应 C++}

\begin{itemize}
\tightlist
\item
  类：\passthrough{\lstinline!unitree::robot::g1::AudioClient!}
\item
  签名：\passthrough{\lstinline!PlayStop(std::string)!}
\item
  绑定策略：\passthrough{\lstinline!DIRECT!}
\item
  声明位置：\passthrough{\lstinline!include/unitree/robot/g1/audio/g1\_audio\_client.hpp:80!}
\end{itemize}

\textbf{说明}

\begin{itemize}
\tightlist
\item
  示例是局部调用片段；\passthrough{\lstinline!obj!}
  和参数变量表示已经按上层业务校验并准备好的对象。
\item
  该条目可以被编辑器和类型检查器识别，但当前运行时可能在属性查找阶段直接失败。
\end{itemize}

\textbf{用法}

\begin{lstlisting}[language=Python]
from typing import TYPE_CHECKING

if TYPE_CHECKING:
    def planned_play_stop(obj: AudioClient, app_name: str) -> int:
        return obj.play_stop(app_name=app_name)
\end{lstlisting}

\begin{quote}
{[}!CAUTION{]} 上面的 \passthrough{\lstinline!TYPE\_CHECKING!}
示例只表达计划调用形态。该分支在普通 Python
运行时为假，不代表当前扩展能执行此接口。
\end{quote}

\hypertarget{unitree-sdk2-cpp-robot-g1-audioclient-led-control-1}{%
\paragraph{\texorpdfstring{\texttt{unitree\_sdk2\_cpp.robot.g1.AudioClient.led\_control}}{unitree\_sdk2\_cpp.robot.g1.AudioClient.led\_control}}\label{unitree-sdk2-cpp-robot-g1-audioclient-led-control-1}}

对应 C++ SDK 操作
\passthrough{\lstinline!LedControl(uint8\_t, uint8\_t, uint8\_t)!}。上游头文件没有可直接生成的业务说明时，本参考只保证签名映射，精确语义需查目标型号协议。

\textbf{签名}

\begin{lstlisting}[language=Python]
def led_control(self, r: int, g: int, b: int) -> int
\end{lstlisting}

\textbf{可用性与安全性}

\textbf{\passthrough{\lstinline!SIGNATURE\_ONLY!} /
\passthrough{\lstinline!HARDWARE\_SIDE\_EFFECT!}}：仅用于补全和静态检查。当前二进制没有该
Python 方法，未来实现还需单独评审硬件副作用。

\textbf{参数}

\begin{longtable}[]{@{}
  >{\raggedright\arraybackslash}p{(\columnwidth - 6\tabcolsep) * \real{0.18}}
  >{\raggedright\arraybackslash}p{(\columnwidth - 6\tabcolsep) * \real{0.24}}
  >{\raggedright\arraybackslash}p{(\columnwidth - 6\tabcolsep) * \real{0.18}}
  >{\raggedright\arraybackslash}p{(\columnwidth - 6\tabcolsep) * \real{0.40}}@{}}
\toprule
名称 & Python 类型 & 默认值 & 含义 \\
\midrule
\endhead
\passthrough{\lstinline!r!} & \passthrough{\lstinline!int!} & 必填 &
传给该接口的 \passthrough{\lstinline!r!} 参数，Python 类型为
\passthrough{\lstinline!int!}。精确含义、单位、范围和枚举值需查目标型号对应头文件与协议。
对应 C++ 参数 \passthrough{\lstinline!R: uint8\_t!}。 取值范围为 0 到
255。 \\
\passthrough{\lstinline!g!} & \passthrough{\lstinline!int!} & 必填 &
传给该接口的 \passthrough{\lstinline!g!} 参数，Python 类型为
\passthrough{\lstinline!int!}。精确含义、单位、范围和枚举值需查目标型号对应头文件与协议。
对应 C++ 参数 \passthrough{\lstinline!G: uint8\_t!}。 取值范围为 0 到
255。 \\
\passthrough{\lstinline!b!} & \passthrough{\lstinline!int!} & 必填 &
传给该接口的 \passthrough{\lstinline!b!} 参数，Python 类型为
\passthrough{\lstinline!int!}。精确含义、单位、范围和枚举值需查目标型号对应头文件与协议。
对应 C++ 参数 \passthrough{\lstinline!B: uint8\_t!}。 取值范围为 0 到
255。 \\
\bottomrule
\end{longtable}

\textbf{返回值}

\begin{longtable}[]{@{}
  >{\raggedright\arraybackslash}p{(\columnwidth - 4\tabcolsep) * \real{0.18}}
  >{\raggedright\arraybackslash}p{(\columnwidth - 4\tabcolsep) * \real{0.20}}
  >{\raggedright\arraybackslash}p{(\columnwidth - 4\tabcolsep) * \real{0.62}}@{}}
\toprule
位置 & 类型 & 含义 \\
\midrule
\endhead
返回值 & \passthrough{\lstinline!int!} & SDK 状态码；Unitree 示例通常以
\passthrough{\lstinline!0!} 表示成功，非零值需按具体服务定义解释。 \\
\bottomrule
\end{longtable}

\textbf{对应 C++}

\begin{itemize}
\tightlist
\item
  类：\passthrough{\lstinline!unitree::robot::g1::AudioClient!}
\item
  签名：\passthrough{\lstinline!LedControl(uint8\_t, uint8\_t, uint8\_t)!}
\item
  绑定策略：\passthrough{\lstinline!DIRECT!}
\item
  声明位置：\passthrough{\lstinline!include/unitree/robot/g1/audio/g1\_audio\_client.hpp:91!}
\end{itemize}

\textbf{说明}

\begin{itemize}
\tightlist
\item
  示例是局部调用片段；\passthrough{\lstinline!obj!}
  和参数变量表示已经按上层业务校验并准备好的对象。
\item
  该条目可以被编辑器和类型检查器识别，但当前运行时可能在属性查找阶段直接失败。
\end{itemize}

\textbf{用法}

\begin{lstlisting}[language=Python]
from typing import TYPE_CHECKING

if TYPE_CHECKING:
    def planned_led_control(obj: AudioClient, r: int, g: int, b: int) -> int:
        return obj.led_control(r=r, g=g, b=b)
\end{lstlisting}

\begin{quote}
{[}!CAUTION{]} 上面的 \passthrough{\lstinline!TYPE\_CHECKING!}
示例只表达计划调用形态。该分支在普通 Python
运行时为假，不代表当前扩展能执行此接口。
\end{quote}

\hypertarget{unitree-sdk2-cpp-robot-g1-g1armactionclient}{%
\subsubsection{\texorpdfstring{\texttt{unitree\_sdk2\_cpp.robot.g1.G1ArmActionClient}}{unitree\_sdk2\_cpp.robot.g1.G1ArmActionClient}}\label{unitree-sdk2-cpp-robot-g1-g1armactionclient}}

Robot 服务客户端。构造、初始化和具体方法的可用性必须分别检查。

\textbf{导入}

\begin{lstlisting}[language=Python]
from unitree_sdk2_cpp.robot.g1 import G1ArmActionClient
\end{lstlisting}

\textbf{基类}：\passthrough{\lstinline!Client!}

\textbf{公开属性}

\begin{longtable}[]{@{}
  >{\raggedright\arraybackslash}p{(\columnwidth - 6\tabcolsep) * \real{0.18}}
  >{\raggedright\arraybackslash}p{(\columnwidth - 6\tabcolsep) * \real{0.24}}
  >{\raggedright\arraybackslash}p{(\columnwidth - 6\tabcolsep) * \real{0.18}}
  >{\raggedright\arraybackslash}p{(\columnwidth - 6\tabcolsep) * \real{0.40}}@{}}
\toprule
名称 & Python 类型 & 含义 & 用法 \\
\midrule
\endhead
\passthrough{\lstinline!action\_map!} &
\passthrough{\lstinline!dict[str, int]!} & 传给该接口的
\passthrough{\lstinline!action!} \passthrough{\lstinline!map!}
参数，Python 类型为
\passthrough{\lstinline!dict[str, int]!}。精确含义、单位、范围和枚举值需查目标型号对应头文件与协议。
& \passthrough{\lstinline!value = obj.action\_map!} /
\passthrough{\lstinline!obj.action\_map = value!} \\
\bottomrule
\end{longtable}

\textbf{方法索引}

\begin{longtable}[]{@{}
  >{\raggedright\arraybackslash}p{(\columnwidth - 6\tabcolsep) * \real{0.18}}
  >{\raggedright\arraybackslash}p{(\columnwidth - 6\tabcolsep) * \real{0.24}}
  >{\raggedright\arraybackslash}p{(\columnwidth - 6\tabcolsep) * \real{0.18}}
  >{\raggedright\arraybackslash}p{(\columnwidth - 6\tabcolsep) * \real{0.40}}@{}}
\toprule
方法 & Python 签名 & 状态 & 安全分类 \\
\midrule
\endhead
\protect\hyperlink{unitree-sdk2-cpp-robot-g1-g1armactionclient-dunder-init-1}{\passthrough{\lstinline!\_\_init\_\_!}}
& \passthrough{\lstinline!def \_\_init\_\_(self) -> None!} &
\passthrough{\lstinline!AVAILABLE!} &
\passthrough{\lstinline!CONSTRUCTION!} \\
\protect\hyperlink{unitree-sdk2-cpp-robot-g1-g1armactionclient-init-1}{\passthrough{\lstinline!init!}}
& \passthrough{\lstinline!def init(self) -> None!} &
\passthrough{\lstinline!AVAILABLE!} &
\passthrough{\lstinline!INITIALIZATION!} \\
\protect\hyperlink{unitree-sdk2-cpp-robot-g1-g1armactionclient-execute-action-1}{\passthrough{\lstinline!execute\_action（重载 1/2）!}}
&
\passthrough{\lstinline!def execute\_action(self, action\_id: int) -> int!}
& \passthrough{\lstinline!SIGNATURE\_ONLY!} &
\passthrough{\lstinline!MOTION\_COMMAND!} \\
\protect\hyperlink{unitree-sdk2-cpp-robot-g1-g1armactionclient-execute-action-2}{\passthrough{\lstinline!execute\_action（重载 2/2）!}}
&
\passthrough{\lstinline!def execute\_action(self, action\_name: str) -> int!}
& \passthrough{\lstinline!SIGNATURE\_ONLY!} &
\passthrough{\lstinline!MOTION\_COMMAND!} \\
\protect\hyperlink{unitree-sdk2-cpp-robot-g1-g1armactionclient-stop-custom-action-1}{\passthrough{\lstinline!stop\_custom\_action!}}
& \passthrough{\lstinline!def stop\_custom\_action(self) -> int!} &
\passthrough{\lstinline!SIGNATURE\_ONLY!} &
\passthrough{\lstinline!MOTION\_COMMAND!} \\
\protect\hyperlink{unitree-sdk2-cpp-robot-g1-g1armactionclient-get-action-list-1}{\passthrough{\lstinline!get\_action\_list!}}
&
\passthrough{\lstinline!def get\_action\_list(self) -> tuple[int, str]!}
& \passthrough{\lstinline!AVAILABLE!} &
\passthrough{\lstinline!READ\_ONLY!} \\
\bottomrule
\end{longtable}

\hypertarget{unitree-sdk2-cpp-robot-g1-g1armactionclient-dunder-init-1}{%
\paragraph{\texorpdfstring{\texttt{unitree\_sdk2\_cpp.robot.g1.G1ArmActionClient.\_\_init\_\_}}{unitree\_sdk2\_cpp.robot.g1.G1ArmActionClient.\_\_init\_\_}}\label{unitree-sdk2-cpp-robot-g1-g1armactionclient-dunder-init-1}}

初始化 \passthrough{\lstinline!G1ArmActionClient!}
实例。是否能实际构造取决于下方可用性状态。

\textbf{签名}

\begin{lstlisting}[language=Python]
def __init__(self) -> None
\end{lstlisting}

\textbf{可用性与安全性}

\textbf{\passthrough{\lstinline!AVAILABLE!} /
\passthrough{\lstinline!CONSTRUCTION!}}：当前绑定源码已实现；实际执行条件仍取决于平台和对应
SDK 环境。

\textbf{参数}

无显式参数。实例方法中的 \passthrough{\lstinline!self!} 由 Python
自动传入。

\textbf{返回值}

\begin{longtable}[]{@{}
  >{\raggedright\arraybackslash}p{(\columnwidth - 4\tabcolsep) * \real{0.18}}
  >{\raggedright\arraybackslash}p{(\columnwidth - 4\tabcolsep) * \real{0.20}}
  >{\raggedright\arraybackslash}p{(\columnwidth - 4\tabcolsep) * \real{0.62}}@{}}
\toprule
位置 & 类型 & 含义 \\
\midrule
\endhead
无 & \passthrough{\lstinline!None!} &
\passthrough{\lstinline!\_\_init\_\_!}
本身不返回值；调用类对象时，在构造函数可用的前提下得到该类实例。 \\
\bottomrule
\end{longtable}

\textbf{对应 C++}

\begin{itemize}
\tightlist
\item
  类：\passthrough{\lstinline!unitree::robot::g1::G1ArmActionClient!}
\item
  签名：\passthrough{\lstinline!G1ArmActionClient()!}
\item
  绑定策略：\passthrough{\lstinline!未单独标注!}
\item
  声明位置：\passthrough{\lstinline!include/unitree/robot/g1/arm/g1\_arm\_action\_client.hpp:26!}
\end{itemize}

\textbf{说明}

\begin{itemize}
\tightlist
\item
  示例是局部调用片段；\passthrough{\lstinline!obj!}
  和参数变量表示已经按上层业务校验并准备好的对象。
\end{itemize}

\textbf{用法}

\begin{lstlisting}[language=Python]
value = G1ArmActionClient()
\end{lstlisting}

\hypertarget{unitree-sdk2-cpp-robot-g1-g1armactionclient-init-1}{%
\paragraph{\texorpdfstring{\texttt{unitree\_sdk2\_cpp.robot.g1.G1ArmActionClient.init}}{unitree\_sdk2\_cpp.robot.g1.G1ArmActionClient.init}}\label{unitree-sdk2-cpp-robot-g1-g1armactionclient-init-1}}

初始化当前 SDK 对象所需的底层通道或服务资源。

\textbf{签名}

\begin{lstlisting}[language=Python]
def init(self) -> None
\end{lstlisting}

\textbf{可用性与安全性}

\textbf{\passthrough{\lstinline!AVAILABLE!} /
\passthrough{\lstinline!INITIALIZATION!}}：当前绑定源码已实现；实际执行条件仍取决于平台和对应
SDK 环境。

\textbf{参数}

无显式参数。实例方法中的 \passthrough{\lstinline!self!} 由 Python
自动传入。

\textbf{返回值}

\begin{longtable}[]{@{}
  >{\raggedright\arraybackslash}p{(\columnwidth - 4\tabcolsep) * \real{0.18}}
  >{\raggedright\arraybackslash}p{(\columnwidth - 4\tabcolsep) * \real{0.20}}
  >{\raggedright\arraybackslash}p{(\columnwidth - 4\tabcolsep) * \real{0.62}}@{}}
\toprule
位置 & 类型 & 含义 \\
\midrule
\endhead
无 & \passthrough{\lstinline!None!} & 该方法不返回 Python
值；失败可能表现为异常或底层状态变化。 \\
\bottomrule
\end{longtable}

\textbf{对应 C++}

\begin{itemize}
\tightlist
\item
  类：\passthrough{\lstinline!unitree::robot::g1::G1ArmActionClient!}
\item
  签名：\passthrough{\lstinline!Init()!}
\item
  绑定策略：\passthrough{\lstinline!DIRECT!}
\item
  声明位置：\passthrough{\lstinline!include/unitree/robot/g1/arm/g1\_arm\_action\_client.hpp:30!}
\end{itemize}

\textbf{说明}

\begin{itemize}
\tightlist
\item
  示例是局部调用片段；\passthrough{\lstinline!obj!}
  和参数变量表示已经按上层业务校验并准备好的对象。
\end{itemize}

\textbf{用法}

\begin{lstlisting}[language=Python]
obj.init()
\end{lstlisting}

\hypertarget{unitree-sdk2-cpp-robot-g1-g1armactionclient-execute-action-1}{%
\paragraph{\texorpdfstring{\texttt{unitree\_sdk2\_cpp.robot.g1.G1ArmActionClient.execute\_action}（重载
1/2）}{unitree\_sdk2\_cpp.robot.g1.G1ArmActionClient.execute\_action（重载 1/2）}}\label{unitree-sdk2-cpp-robot-g1-g1armactionclient-execute-action-1}}

对应 C++ SDK 操作
\passthrough{\lstinline!ExecuteAction(int32\_t)!}。上游头文件没有可直接生成的业务说明时，本参考只保证签名映射，精确语义需查目标型号协议。

\textbf{签名}

\begin{lstlisting}[language=Python]
@overload
def execute_action(self, action_id: int) -> int
\end{lstlisting}

\textbf{可用性与安全性}

\textbf{\passthrough{\lstinline!SIGNATURE\_ONLY!} /
\passthrough{\lstinline!MOTION\_COMMAND!}}：仅用于补全和静态检查。当前二进制没有该
Python 方法；不得按可执行运动接口使用。

\textbf{参数}

\begin{longtable}[]{@{}
  >{\raggedright\arraybackslash}p{(\columnwidth - 6\tabcolsep) * \real{0.18}}
  >{\raggedright\arraybackslash}p{(\columnwidth - 6\tabcolsep) * \real{0.24}}
  >{\raggedright\arraybackslash}p{(\columnwidth - 6\tabcolsep) * \real{0.18}}
  >{\raggedright\arraybackslash}p{(\columnwidth - 6\tabcolsep) * \real{0.40}}@{}}
\toprule
名称 & Python 类型 & 默认值 & 含义 \\
\midrule
\endhead
\passthrough{\lstinline!action\_id!} & \passthrough{\lstinline!int!} &
必填 & 传给该接口的 \passthrough{\lstinline!action!} ID 参数，Python
类型为
\passthrough{\lstinline!int!}。精确含义、单位、范围和枚举值需查目标型号对应头文件与协议。
对应 C++ 参数 \passthrough{\lstinline!action\_id: int32\_t!}。
取值范围为 -2147483648 到 2147483647。 \\
\bottomrule
\end{longtable}

\textbf{返回值}

\begin{longtable}[]{@{}
  >{\raggedright\arraybackslash}p{(\columnwidth - 4\tabcolsep) * \real{0.18}}
  >{\raggedright\arraybackslash}p{(\columnwidth - 4\tabcolsep) * \real{0.20}}
  >{\raggedright\arraybackslash}p{(\columnwidth - 4\tabcolsep) * \real{0.62}}@{}}
\toprule
位置 & 类型 & 含义 \\
\midrule
\endhead
返回值 & \passthrough{\lstinline!int!} & SDK 状态码；Unitree 示例通常以
\passthrough{\lstinline!0!} 表示成功，非零值需按具体服务定义解释。 \\
\bottomrule
\end{longtable}

\textbf{对应 C++}

\begin{itemize}
\tightlist
\item
  类：\passthrough{\lstinline!unitree::robot::g1::G1ArmActionClient!}
\item
  签名：\passthrough{\lstinline!ExecuteAction(int32\_t)!}
\item
  绑定策略：\passthrough{\lstinline!DIRECT!}
\item
  声明位置：\passthrough{\lstinline!include/unitree/robot/g1/arm/g1\_arm\_action\_client.hpp:51!}
\end{itemize}

\textbf{说明}

\begin{itemize}
\tightlist
\item
  示例是局部调用片段；\passthrough{\lstinline!obj!}
  和参数变量表示已经按上层业务校验并准备好的对象。
\item
  该条目可以被编辑器和类型检查器识别，但当前运行时可能在属性查找阶段直接失败。
\item
  该方法属于运动命令。方法名包含
  \passthrough{\lstinline!stop!}、\passthrough{\lstinline!damp!} 或
  \passthrough{\lstinline!zero!} 也不代表它是独立物理急停。
\end{itemize}

\textbf{用法}

\begin{lstlisting}[language=Python]
from typing import TYPE_CHECKING

if TYPE_CHECKING:
    def planned_execute_action(obj: G1ArmActionClient, action_id: int) -> int:
        return obj.execute_action(action_id=action_id)
\end{lstlisting}

\begin{quote}
{[}!CAUTION{]} 上面的 \passthrough{\lstinline!TYPE\_CHECKING!}
示例只表达计划调用形态。该分支在普通 Python
运行时为假，不代表当前扩展能执行此接口。
\end{quote}

\hypertarget{unitree-sdk2-cpp-robot-g1-g1armactionclient-execute-action-2}{%
\paragraph{\texorpdfstring{\texttt{unitree\_sdk2\_cpp.robot.g1.G1ArmActionClient.execute\_action}（重载
2/2）}{unitree\_sdk2\_cpp.robot.g1.G1ArmActionClient.execute\_action（重载 2/2）}}\label{unitree-sdk2-cpp-robot-g1-g1armactionclient-execute-action-2}}

对应 C++ SDK 操作
\passthrough{\lstinline!ExecuteAction(const std::string \&)!}。上游头文件没有可直接生成的业务说明时，本参考只保证签名映射，精确语义需查目标型号协议。

\textbf{签名}

\begin{lstlisting}[language=Python]
@overload
def execute_action(self, action_name: str) -> int
\end{lstlisting}

\textbf{可用性与安全性}

\textbf{\passthrough{\lstinline!SIGNATURE\_ONLY!} /
\passthrough{\lstinline!MOTION\_COMMAND!}}：仅用于补全和静态检查。当前二进制没有该
Python 方法；不得按可执行运动接口使用。

\textbf{参数}

\begin{longtable}[]{@{}
  >{\raggedright\arraybackslash}p{(\columnwidth - 6\tabcolsep) * \real{0.18}}
  >{\raggedright\arraybackslash}p{(\columnwidth - 6\tabcolsep) * \real{0.24}}
  >{\raggedright\arraybackslash}p{(\columnwidth - 6\tabcolsep) * \real{0.18}}
  >{\raggedright\arraybackslash}p{(\columnwidth - 6\tabcolsep) * \real{0.40}}@{}}
\toprule
名称 & Python 类型 & 默认值 & 含义 \\
\midrule
\endhead
\passthrough{\lstinline!action\_name!} & \passthrough{\lstinline!str!} &
必填 & 传给该接口的 \passthrough{\lstinline!action!} 名称 参数，Python
类型为
\passthrough{\lstinline!str!}。精确含义、单位、范围和枚举值需查目标型号对应头文件与协议。
对应 C++ 参数
\passthrough{\lstinline!action\_name: const std::string \&!}。
底层为字符串；长度、编码和允许值由具体协议决定。 \\
\bottomrule
\end{longtable}

\textbf{返回值}

\begin{longtable}[]{@{}
  >{\raggedright\arraybackslash}p{(\columnwidth - 4\tabcolsep) * \real{0.18}}
  >{\raggedright\arraybackslash}p{(\columnwidth - 4\tabcolsep) * \real{0.20}}
  >{\raggedright\arraybackslash}p{(\columnwidth - 4\tabcolsep) * \real{0.62}}@{}}
\toprule
位置 & 类型 & 含义 \\
\midrule
\endhead
返回值 & \passthrough{\lstinline!int!} & SDK 状态码；Unitree 示例通常以
\passthrough{\lstinline!0!} 表示成功，非零值需按具体服务定义解释。 \\
\bottomrule
\end{longtable}

\textbf{对应 C++}

\begin{itemize}
\tightlist
\item
  类：\passthrough{\lstinline!unitree::robot::g1::G1ArmActionClient!}
\item
  签名：\passthrough{\lstinline!ExecuteAction(const std::string \&)!}
\item
  绑定策略：\passthrough{\lstinline!DIRECT!}
\item
  声明位置：\passthrough{\lstinline!include/unitree/robot/g1/arm/g1\_arm\_action\_client.hpp:67!}
\end{itemize}

\textbf{说明}

\begin{itemize}
\tightlist
\item
  示例是局部调用片段；\passthrough{\lstinline!obj!}
  和参数变量表示已经按上层业务校验并准备好的对象。
\item
  该条目可以被编辑器和类型检查器识别，但当前运行时可能在属性查找阶段直接失败。
\item
  该方法属于运动命令。方法名包含
  \passthrough{\lstinline!stop!}、\passthrough{\lstinline!damp!} 或
  \passthrough{\lstinline!zero!} 也不代表它是独立物理急停。
\end{itemize}

\textbf{用法}

\begin{lstlisting}[language=Python]
from typing import TYPE_CHECKING

if TYPE_CHECKING:
    def planned_execute_action(obj: G1ArmActionClient, action_name: str) -> int:
        return obj.execute_action(action_name=action_name)
\end{lstlisting}

\begin{quote}
{[}!CAUTION{]} 上面的 \passthrough{\lstinline!TYPE\_CHECKING!}
示例只表达计划调用形态。该分支在普通 Python
运行时为假，不代表当前扩展能执行此接口。
\end{quote}

\hypertarget{unitree-sdk2-cpp-robot-g1-g1armactionclient-stop-custom-action-1}{%
\paragraph{\texorpdfstring{\texttt{unitree\_sdk2\_cpp.robot.g1.G1ArmActionClient.stop\_custom\_action}}{unitree\_sdk2\_cpp.robot.g1.G1ArmActionClient.stop\_custom\_action}}\label{unitree-sdk2-cpp-robot-g1-g1armactionclient-stop-custom-action-1}}

请求停止对应 SDK 操作。该名称不等同于经过验证的物理急停。

\textbf{签名}

\begin{lstlisting}[language=Python]
def stop_custom_action(self) -> int
\end{lstlisting}

\textbf{可用性与安全性}

\textbf{\passthrough{\lstinline!SIGNATURE\_ONLY!} /
\passthrough{\lstinline!MOTION\_COMMAND!}}：仅用于补全和静态检查。当前二进制没有该
Python 方法；不得按可执行运动接口使用。

\textbf{参数}

无显式参数。实例方法中的 \passthrough{\lstinline!self!} 由 Python
自动传入。

\textbf{返回值}

\begin{longtable}[]{@{}
  >{\raggedright\arraybackslash}p{(\columnwidth - 4\tabcolsep) * \real{0.18}}
  >{\raggedright\arraybackslash}p{(\columnwidth - 4\tabcolsep) * \real{0.20}}
  >{\raggedright\arraybackslash}p{(\columnwidth - 4\tabcolsep) * \real{0.62}}@{}}
\toprule
位置 & 类型 & 含义 \\
\midrule
\endhead
返回值 & \passthrough{\lstinline!int!} & SDK 状态码；Unitree 示例通常以
\passthrough{\lstinline!0!} 表示成功，非零值需按具体服务定义解释。 \\
\bottomrule
\end{longtable}

\textbf{对应 C++}

\begin{itemize}
\tightlist
\item
  类：\passthrough{\lstinline!unitree::robot::g1::G1ArmActionClient!}
\item
  签名：\passthrough{\lstinline!StopCustomAction()!}
\item
  绑定策略：\passthrough{\lstinline!DIRECT!}
\item
  声明位置：\passthrough{\lstinline!include/unitree/robot/g1/arm/g1\_arm\_action\_client.hpp:74!}
\end{itemize}

\textbf{说明}

\begin{itemize}
\tightlist
\item
  示例是局部调用片段；\passthrough{\lstinline!obj!}
  和参数变量表示已经按上层业务校验并准备好的对象。
\item
  该条目可以被编辑器和类型检查器识别，但当前运行时可能在属性查找阶段直接失败。
\item
  该方法属于运动命令。方法名包含
  \passthrough{\lstinline!stop!}、\passthrough{\lstinline!damp!} 或
  \passthrough{\lstinline!zero!} 也不代表它是独立物理急停。
\end{itemize}

\textbf{用法}

\begin{lstlisting}[language=Python]
from typing import TYPE_CHECKING

if TYPE_CHECKING:
    def planned_stop_custom_action(obj: G1ArmActionClient) -> int:
        return obj.stop_custom_action()
\end{lstlisting}

\begin{quote}
{[}!CAUTION{]} 上面的 \passthrough{\lstinline!TYPE\_CHECKING!}
示例只表达计划调用形态。该分支在普通 Python
运行时为假，不代表当前扩展能执行此接口。
\end{quote}

\hypertarget{unitree-sdk2-cpp-robot-g1-g1armactionclient-get-action-list-1}{%
\paragraph{\texorpdfstring{\texttt{unitree\_sdk2\_cpp.robot.g1.G1ArmActionClient.get\_action\_list}}{unitree\_sdk2\_cpp.robot.g1.G1ArmActionClient.get\_action\_list}}\label{unitree-sdk2-cpp-robot-g1-g1armactionclient-get-action-list-1}}

查询或检查 \passthrough{\lstinline!action!}
\passthrough{\lstinline!list!}。

\textbf{签名}

\begin{lstlisting}[language=Python]
def get_action_list(self) -> tuple[int, str]
\end{lstlisting}

\textbf{可用性与安全性}

\textbf{\passthrough{\lstinline!AVAILABLE!} /
\passthrough{\lstinline!READ\_ONLY!}}：当前绑定源码已实现只读查询。Robot
Client 的服务端查询通常仍需匹配的 Linux
扩展、DDS、网络、机器人和服务；纯本地 getter 则不一定需要全部条件。

\textbf{参数}

无显式参数。实例方法中的 \passthrough{\lstinline!self!} 由 Python
自动传入。

\textbf{返回值}

\begin{longtable}[]{@{}
  >{\raggedright\arraybackslash}p{(\columnwidth - 4\tabcolsep) * \real{0.18}}
  >{\raggedright\arraybackslash}p{(\columnwidth - 4\tabcolsep) * \real{0.20}}
  >{\raggedright\arraybackslash}p{(\columnwidth - 4\tabcolsep) * \real{0.62}}@{}}
\toprule
位置 & 类型 & 含义 \\
\midrule
\endhead
{[}0{]} \passthrough{\lstinline!status!} & \passthrough{\lstinline!int!}
& SDK 状态码；Unitree 示例通常以 \passthrough{\lstinline!0!}
表示成功，非零值需按具体服务错误码解释。 \\
{[}1{]} \passthrough{\lstinline!data!} & \passthrough{\lstinline!str!} &
传给该接口的 数据 参数，Python 类型为
\passthrough{\lstinline!str!}。精确含义、单位、范围和枚举值需查目标型号对应头文件与协议。
对应 C++ 参数 \passthrough{\lstinline!data: std::string \&!}。
底层为字符串；长度、编码和允许值由具体协议决定。 \\
\bottomrule
\end{longtable}

\textbf{对应 C++}

\begin{itemize}
\tightlist
\item
  类：\passthrough{\lstinline!unitree::robot::g1::G1ArmActionClient!}
\item
  签名：\passthrough{\lstinline!GetActionList(std::string \&)!}
\item
  绑定策略：\passthrough{\lstinline!OUTPUT\_WRAPPER!}
\item
  声明位置：\passthrough{\lstinline!include/unitree/robot/g1/arm/g1\_arm\_action\_client.hpp:79!}
\end{itemize}

\textbf{说明}

\begin{itemize}
\tightlist
\item
  示例是局部调用片段；\passthrough{\lstinline!obj!}
  和参数变量表示已经按上层业务校验并准备好的对象。
\item
  C++ 的可修改输出引用不会作为 Python 入参出现，而是按 C++
  参数顺序追加到返回元组中。
\end{itemize}

\textbf{用法}

\begin{lstlisting}[language=Python]
status, data = obj.get_action_list()
\end{lstlisting}

\begin{quote}
{[}!NOTE{]} 只读查询仍会访问
DDS/机器人服务。必须检查返回状态码，并处理超时和网络中断。
\end{quote}

\hypertarget{unitree-sdk2-cpp-robot-g1-jsonizedatavecfloat}{%
\subsubsection{\texorpdfstring{\texttt{unitree\_sdk2\_cpp.robot.g1.JsonizeDataVecFloat}}{unitree\_sdk2\_cpp.robot.g1.JsonizeDataVecFloat}}\label{unitree-sdk2-cpp-robot-g1-jsonizedatavecfloat}}

SDK 类型签名预览。请逐项查看构造函数和方法的可用性。

\textbf{导入}

\begin{lstlisting}[language=Python]
from unitree_sdk2_cpp.robot.g1 import JsonizeDataVecFloat
\end{lstlisting}

\textbf{基类}：\passthrough{\lstinline!object!}

\textbf{公开属性}

\begin{longtable}[]{@{}
  >{\raggedright\arraybackslash}p{(\columnwidth - 6\tabcolsep) * \real{0.18}}
  >{\raggedright\arraybackslash}p{(\columnwidth - 6\tabcolsep) * \real{0.24}}
  >{\raggedright\arraybackslash}p{(\columnwidth - 6\tabcolsep) * \real{0.18}}
  >{\raggedright\arraybackslash}p{(\columnwidth - 6\tabcolsep) * \real{0.40}}@{}}
\toprule
名称 & Python 类型 & 含义 & 用法 \\
\midrule
\endhead
\passthrough{\lstinline!data!} & \passthrough{\lstinline!list[float]!} &
传给该接口的 数据 参数，Python 类型为
\passthrough{\lstinline!list[float]!}。精确含义、单位、范围和枚举值需查目标型号对应头文件与协议。
& \passthrough{\lstinline!value = obj.data!} /
\passthrough{\lstinline!obj.data = value!} \\
\bottomrule
\end{longtable}

\textbf{方法索引}

\begin{longtable}[]{@{}
  >{\raggedright\arraybackslash}p{(\columnwidth - 6\tabcolsep) * \real{0.18}}
  >{\raggedright\arraybackslash}p{(\columnwidth - 6\tabcolsep) * \real{0.24}}
  >{\raggedright\arraybackslash}p{(\columnwidth - 6\tabcolsep) * \real{0.18}}
  >{\raggedright\arraybackslash}p{(\columnwidth - 6\tabcolsep) * \real{0.40}}@{}}
\toprule
方法 & Python 签名 & 状态 & 安全分类 \\
\midrule
\endhead
\protect\hyperlink{unitree-sdk2-cpp-robot-g1-jsonizedatavecfloat-dunder-init-1}{\passthrough{\lstinline!\_\_init\_\_!}}
& \passthrough{\lstinline!def \_\_init\_\_(self) -> None!} &
\passthrough{\lstinline!SIGNATURE\_ONLY!} &
\passthrough{\lstinline!CONSTRUCTION!} \\
\protect\hyperlink{unitree-sdk2-cpp-robot-g1-jsonizedatavecfloat-from-json-1}{\passthrough{\lstinline!from\_json!}}
&
\passthrough{\lstinline!def from\_json(self, json: dict[str, Any]) -> None!}
& \passthrough{\lstinline!SIGNATURE\_ONLY!} &
\passthrough{\lstinline!UNCLASSIFIED!} \\
\protect\hyperlink{unitree-sdk2-cpp-robot-g1-jsonizedatavecfloat-to-json-1}{\passthrough{\lstinline!to\_json!}}
&
\passthrough{\lstinline!def to\_json(self, json: dict[str, Any]) -> None!}
& \passthrough{\lstinline!SIGNATURE\_ONLY!} &
\passthrough{\lstinline!UNCLASSIFIED!} \\
\bottomrule
\end{longtable}

\hypertarget{unitree-sdk2-cpp-robot-g1-jsonizedatavecfloat-dunder-init-1}{%
\paragraph{\texorpdfstring{\texttt{unitree\_sdk2\_cpp.robot.g1.JsonizeDataVecFloat.\_\_init\_\_}}{unitree\_sdk2\_cpp.robot.g1.JsonizeDataVecFloat.\_\_init\_\_}}\label{unitree-sdk2-cpp-robot-g1-jsonizedatavecfloat-dunder-init-1}}

初始化 \passthrough{\lstinline!JsonizeDataVecFloat!}
实例。是否能实际构造取决于下方可用性状态。

\textbf{签名}

\begin{lstlisting}[language=Python]
def __init__(self) -> None
\end{lstlisting}

\textbf{可用性与安全性}

\textbf{\passthrough{\lstinline!SIGNATURE\_ONLY!} /
\passthrough{\lstinline!CONSTRUCTION!}}：当前只有设计期签名，不能假设已存在于运行时扩展。

\textbf{参数}

无显式参数。实例方法中的 \passthrough{\lstinline!self!} 由 Python
自动传入。

\textbf{返回值}

\begin{longtable}[]{@{}
  >{\raggedright\arraybackslash}p{(\columnwidth - 4\tabcolsep) * \real{0.18}}
  >{\raggedright\arraybackslash}p{(\columnwidth - 4\tabcolsep) * \real{0.20}}
  >{\raggedright\arraybackslash}p{(\columnwidth - 4\tabcolsep) * \real{0.62}}@{}}
\toprule
位置 & 类型 & 含义 \\
\midrule
\endhead
无 & \passthrough{\lstinline!None!} &
\passthrough{\lstinline!\_\_init\_\_!}
本身不返回值；调用类对象时，在构造函数可用的前提下得到该类实例。 \\
\bottomrule
\end{longtable}

\textbf{对应 C++}

\begin{itemize}
\tightlist
\item
  类：\passthrough{\lstinline!unitree::robot::g1::JsonizeDataVecFloat!}
\item
  签名：\passthrough{\lstinline!JsonizeDataVecFloat()!}
\item
  绑定策略：\passthrough{\lstinline!未单独标注!}
\item
  声明位置：\passthrough{\lstinline!include/unitree/robot/g1/loco/g1\_loco\_api.hpp:45!}
\end{itemize}

\textbf{说明}

\begin{itemize}
\tightlist
\item
  示例是局部调用片段；\passthrough{\lstinline!obj!}
  和参数变量表示已经按上层业务校验并准备好的对象。
\item
  该条目可以被编辑器和类型检查器识别，但当前运行时可能在属性查找阶段直接失败。
\end{itemize}

\textbf{用法}

\begin{lstlisting}[language=Python]
from typing import TYPE_CHECKING

if TYPE_CHECKING:
    def planned_construct() -> JsonizeDataVecFloat:
        return JsonizeDataVecFloat()
\end{lstlisting}

\begin{quote}
{[}!CAUTION{]} 上面的 \passthrough{\lstinline!TYPE\_CHECKING!}
示例只表达计划调用形态。该分支在普通 Python
运行时为假，不代表当前扩展能执行此接口。
\end{quote}

\hypertarget{unitree-sdk2-cpp-robot-g1-jsonizedatavecfloat-from-json-1}{%
\paragraph{\texorpdfstring{\texttt{unitree\_sdk2\_cpp.robot.g1.JsonizeDataVecFloat.from\_json}}{unitree\_sdk2\_cpp.robot.g1.JsonizeDataVecFloat.from\_json}}\label{unitree-sdk2-cpp-robot-g1-jsonizedatavecfloat-from-json-1}}

计划从 JSON 风格字典读取字段并更新当前 SDK 值对象。

\textbf{签名}

\begin{lstlisting}[language=Python]
def from_json(self, json: dict[str, Any]) -> None
\end{lstlisting}

\textbf{可用性与安全性}

\textbf{\passthrough{\lstinline!SIGNATURE\_ONLY!} /
\passthrough{\lstinline!UNCLASSIFIED!}}：当前只有设计期签名，不能假设已存在于运行时扩展。

\textbf{参数}

\begin{longtable}[]{@{}
  >{\raggedright\arraybackslash}p{(\columnwidth - 6\tabcolsep) * \real{0.18}}
  >{\raggedright\arraybackslash}p{(\columnwidth - 6\tabcolsep) * \real{0.24}}
  >{\raggedright\arraybackslash}p{(\columnwidth - 6\tabcolsep) * \real{0.18}}
  >{\raggedright\arraybackslash}p{(\columnwidth - 6\tabcolsep) * \real{0.40}}@{}}
\toprule
名称 & Python 类型 & 默认值 & 含义 \\
\midrule
\endhead
\passthrough{\lstinline!json!} &
\passthrough{\lstinline!dict[str, Any]!} & 必填 & JSON
风格字典，键和值必须满足对应 C++ \passthrough{\lstinline!JsonMap!}
数据结构的约束。 对应 C++ 参数
\passthrough{\lstinline!json: common::JsonMap \&!}。 \\
\bottomrule
\end{longtable}

\textbf{返回值}

\begin{longtable}[]{@{}
  >{\raggedright\arraybackslash}p{(\columnwidth - 4\tabcolsep) * \real{0.18}}
  >{\raggedright\arraybackslash}p{(\columnwidth - 4\tabcolsep) * \real{0.20}}
  >{\raggedright\arraybackslash}p{(\columnwidth - 4\tabcolsep) * \real{0.62}}@{}}
\toprule
位置 & 类型 & 含义 \\
\midrule
\endhead
无 & \passthrough{\lstinline!None!} & 该方法不返回 Python
值；失败可能表现为异常或底层状态变化。 \\
\bottomrule
\end{longtable}

\textbf{对应 C++}

\begin{itemize}
\tightlist
\item
  类：\passthrough{\lstinline!unitree::robot::g1::JsonizeDataVecFloat!}
\item
  签名：\passthrough{\lstinline!fromJson(common::JsonMap \&)!}
\item
  绑定策略：\passthrough{\lstinline!SIGNATURE\_PREVIEW!}
\item
  声明位置：\passthrough{\lstinline!include/unitree/robot/g1/loco/g1\_loco\_api.hpp:48!}
\end{itemize}

\textbf{说明}

\begin{itemize}
\tightlist
\item
  示例是局部调用片段；\passthrough{\lstinline!obj!}
  和参数变量表示已经按上层业务校验并准备好的对象。
\item
  该条目可以被编辑器和类型检查器识别，但当前运行时可能在属性查找阶段直接失败。
\end{itemize}

\textbf{用法}

\begin{lstlisting}[language=Python]
from typing import TYPE_CHECKING

if TYPE_CHECKING:
    def planned_from_json(obj: JsonizeDataVecFloat, json: dict[str, Any]) -> None:
        obj.from_json(json=json)
\end{lstlisting}

\begin{quote}
{[}!CAUTION{]} 上面的 \passthrough{\lstinline!TYPE\_CHECKING!}
示例只表达计划调用形态。该分支在普通 Python
运行时为假，不代表当前扩展能执行此接口。
\end{quote}

\hypertarget{unitree-sdk2-cpp-robot-g1-jsonizedatavecfloat-to-json-1}{%
\paragraph{\texorpdfstring{\texttt{unitree\_sdk2\_cpp.robot.g1.JsonizeDataVecFloat.to\_json}}{unitree\_sdk2\_cpp.robot.g1.JsonizeDataVecFloat.to\_json}}\label{unitree-sdk2-cpp-robot-g1-jsonizedatavecfloat-to-json-1}}

计划把当前 SDK 值对象写入 JSON 风格字典。

\textbf{签名}

\begin{lstlisting}[language=Python]
def to_json(self, json: dict[str, Any]) -> None
\end{lstlisting}

\textbf{可用性与安全性}

\textbf{\passthrough{\lstinline!SIGNATURE\_ONLY!} /
\passthrough{\lstinline!UNCLASSIFIED!}}：当前只有设计期签名，不能假设已存在于运行时扩展。

\textbf{参数}

\begin{longtable}[]{@{}
  >{\raggedright\arraybackslash}p{(\columnwidth - 6\tabcolsep) * \real{0.18}}
  >{\raggedright\arraybackslash}p{(\columnwidth - 6\tabcolsep) * \real{0.24}}
  >{\raggedright\arraybackslash}p{(\columnwidth - 6\tabcolsep) * \real{0.18}}
  >{\raggedright\arraybackslash}p{(\columnwidth - 6\tabcolsep) * \real{0.40}}@{}}
\toprule
名称 & Python 类型 & 默认值 & 含义 \\
\midrule
\endhead
\passthrough{\lstinline!json!} &
\passthrough{\lstinline!dict[str, Any]!} & 必填 & JSON
风格字典，键和值必须满足对应 C++ \passthrough{\lstinline!JsonMap!}
数据结构的约束。 对应 C++ 参数
\passthrough{\lstinline!json: common::JsonMap \&!}。 \\
\bottomrule
\end{longtable}

\textbf{返回值}

\begin{longtable}[]{@{}
  >{\raggedright\arraybackslash}p{(\columnwidth - 4\tabcolsep) * \real{0.18}}
  >{\raggedright\arraybackslash}p{(\columnwidth - 4\tabcolsep) * \real{0.20}}
  >{\raggedright\arraybackslash}p{(\columnwidth - 4\tabcolsep) * \real{0.62}}@{}}
\toprule
位置 & 类型 & 含义 \\
\midrule
\endhead
无 & \passthrough{\lstinline!None!} & 该方法不返回 Python
值；失败可能表现为异常或底层状态变化。 \\
\bottomrule
\end{longtable}

\textbf{对应 C++}

\begin{itemize}
\tightlist
\item
  类：\passthrough{\lstinline!unitree::robot::g1::JsonizeDataVecFloat!}
\item
  签名：\passthrough{\lstinline!toJson(common::JsonMap \&) const!}
\item
  绑定策略：\passthrough{\lstinline!SIGNATURE\_PREVIEW!}
\item
  声明位置：\passthrough{\lstinline!include/unitree/robot/g1/loco/g1\_loco\_api.hpp:50!}
\end{itemize}

\textbf{说明}

\begin{itemize}
\tightlist
\item
  示例是局部调用片段；\passthrough{\lstinline!obj!}
  和参数变量表示已经按上层业务校验并准备好的对象。
\item
  该条目可以被编辑器和类型检查器识别，但当前运行时可能在属性查找阶段直接失败。
\end{itemize}

\textbf{用法}

\begin{lstlisting}[language=Python]
from typing import TYPE_CHECKING

if TYPE_CHECKING:
    def planned_to_json(obj: JsonizeDataVecFloat, json: dict[str, Any]) -> None:
        obj.to_json(json=json)
\end{lstlisting}

\begin{quote}
{[}!CAUTION{]} 上面的 \passthrough{\lstinline!TYPE\_CHECKING!}
示例只表达计划调用形态。该分支在普通 Python
运行时为假，不代表当前扩展能执行此接口。
\end{quote}

\hypertarget{unitree-sdk2-cpp-robot-g1-jsonizevelocitycommand}{%
\subsubsection{\texorpdfstring{\texttt{unitree\_sdk2\_cpp.robot.g1.JsonizeVelocityCommand}}{unitree\_sdk2\_cpp.robot.g1.JsonizeVelocityCommand}}\label{unitree-sdk2-cpp-robot-g1-jsonizevelocitycommand}}

SDK 类型签名预览。请逐项查看构造函数和方法的可用性。

\textbf{导入}

\begin{lstlisting}[language=Python]
from unitree_sdk2_cpp.robot.g1 import JsonizeVelocityCommand
\end{lstlisting}

\textbf{基类}：\passthrough{\lstinline!object!}

\textbf{公开属性}

\begin{longtable}[]{@{}
  >{\raggedright\arraybackslash}p{(\columnwidth - 6\tabcolsep) * \real{0.18}}
  >{\raggedright\arraybackslash}p{(\columnwidth - 6\tabcolsep) * \real{0.24}}
  >{\raggedright\arraybackslash}p{(\columnwidth - 6\tabcolsep) * \real{0.18}}
  >{\raggedright\arraybackslash}p{(\columnwidth - 6\tabcolsep) * \real{0.40}}@{}}
\toprule
名称 & Python 类型 & 含义 & 用法 \\
\midrule
\endhead
\passthrough{\lstinline!velocity!} &
\passthrough{\lstinline!list[float]!} & 传给该接口的 速度 参数，Python
类型为
\passthrough{\lstinline!list[float]!}。精确含义、单位、范围和枚举值需查目标型号对应头文件与协议。
& \passthrough{\lstinline!value = obj.velocity!} /
\passthrough{\lstinline!obj.velocity = value!} \\
\passthrough{\lstinline!duration!} & \passthrough{\lstinline!float!} &
操作持续时间参数。精确单位、范围和默认行为以目标型号接口定义为准。 &
\passthrough{\lstinline!value = obj.duration!} /
\passthrough{\lstinline!obj.duration = value!} \\
\bottomrule
\end{longtable}

\textbf{方法索引}

\begin{longtable}[]{@{}
  >{\raggedright\arraybackslash}p{(\columnwidth - 6\tabcolsep) * \real{0.18}}
  >{\raggedright\arraybackslash}p{(\columnwidth - 6\tabcolsep) * \real{0.24}}
  >{\raggedright\arraybackslash}p{(\columnwidth - 6\tabcolsep) * \real{0.18}}
  >{\raggedright\arraybackslash}p{(\columnwidth - 6\tabcolsep) * \real{0.40}}@{}}
\toprule
方法 & Python 签名 & 状态 & 安全分类 \\
\midrule
\endhead
\protect\hyperlink{unitree-sdk2-cpp-robot-g1-jsonizevelocitycommand-dunder-init-1}{\passthrough{\lstinline!\_\_init\_\_!}}
& \passthrough{\lstinline!def \_\_init\_\_(self) -> None!} &
\passthrough{\lstinline!SIGNATURE\_ONLY!} &
\passthrough{\lstinline!CONSTRUCTION!} \\
\protect\hyperlink{unitree-sdk2-cpp-robot-g1-jsonizevelocitycommand-from-json-1}{\passthrough{\lstinline!from\_json!}}
&
\passthrough{\lstinline!def from\_json(self, json: dict[str, Any]) -> None!}
& \passthrough{\lstinline!SIGNATURE\_ONLY!} &
\passthrough{\lstinline!UNCLASSIFIED!} \\
\protect\hyperlink{unitree-sdk2-cpp-robot-g1-jsonizevelocitycommand-to-json-1}{\passthrough{\lstinline!to\_json!}}
&
\passthrough{\lstinline!def to\_json(self, json: dict[str, Any]) -> None!}
& \passthrough{\lstinline!SIGNATURE\_ONLY!} &
\passthrough{\lstinline!UNCLASSIFIED!} \\
\bottomrule
\end{longtable}

\hypertarget{unitree-sdk2-cpp-robot-g1-jsonizevelocitycommand-dunder-init-1}{%
\paragraph{\texorpdfstring{\texttt{unitree\_sdk2\_cpp.robot.g1.JsonizeVelocityCommand.\_\_init\_\_}}{unitree\_sdk2\_cpp.robot.g1.JsonizeVelocityCommand.\_\_init\_\_}}\label{unitree-sdk2-cpp-robot-g1-jsonizevelocitycommand-dunder-init-1}}

初始化 \passthrough{\lstinline!JsonizeVelocityCommand!}
实例。是否能实际构造取决于下方可用性状态。

\textbf{签名}

\begin{lstlisting}[language=Python]
def __init__(self) -> None
\end{lstlisting}

\textbf{可用性与安全性}

\textbf{\passthrough{\lstinline!SIGNATURE\_ONLY!} /
\passthrough{\lstinline!CONSTRUCTION!}}：当前只有设计期签名，不能假设已存在于运行时扩展。

\textbf{参数}

无显式参数。实例方法中的 \passthrough{\lstinline!self!} 由 Python
自动传入。

\textbf{返回值}

\begin{longtable}[]{@{}
  >{\raggedright\arraybackslash}p{(\columnwidth - 4\tabcolsep) * \real{0.18}}
  >{\raggedright\arraybackslash}p{(\columnwidth - 4\tabcolsep) * \real{0.20}}
  >{\raggedright\arraybackslash}p{(\columnwidth - 4\tabcolsep) * \real{0.62}}@{}}
\toprule
位置 & 类型 & 含义 \\
\midrule
\endhead
无 & \passthrough{\lstinline!None!} &
\passthrough{\lstinline!\_\_init\_\_!}
本身不返回值；调用类对象时，在构造函数可用的前提下得到该类实例。 \\
\bottomrule
\end{longtable}

\textbf{对应 C++}

\begin{itemize}
\tightlist
\item
  类：\passthrough{\lstinline!unitree::robot::g1::JsonizeVelocityCommand!}
\item
  签名：\passthrough{\lstinline!JsonizeVelocityCommand()!}
\item
  绑定策略：\passthrough{\lstinline!未单独标注!}
\item
  声明位置：\passthrough{\lstinline!include/unitree/robot/g1/loco/g1\_loco\_api.hpp:59!}
\end{itemize}

\textbf{说明}

\begin{itemize}
\tightlist
\item
  示例是局部调用片段；\passthrough{\lstinline!obj!}
  和参数变量表示已经按上层业务校验并准备好的对象。
\item
  该条目可以被编辑器和类型检查器识别，但当前运行时可能在属性查找阶段直接失败。
\end{itemize}

\textbf{用法}

\begin{lstlisting}[language=Python]
from typing import TYPE_CHECKING

if TYPE_CHECKING:
    def planned_construct() -> JsonizeVelocityCommand:
        return JsonizeVelocityCommand()
\end{lstlisting}

\begin{quote}
{[}!CAUTION{]} 上面的 \passthrough{\lstinline!TYPE\_CHECKING!}
示例只表达计划调用形态。该分支在普通 Python
运行时为假，不代表当前扩展能执行此接口。
\end{quote}

\hypertarget{unitree-sdk2-cpp-robot-g1-jsonizevelocitycommand-from-json-1}{%
\paragraph{\texorpdfstring{\texttt{unitree\_sdk2\_cpp.robot.g1.JsonizeVelocityCommand.from\_json}}{unitree\_sdk2\_cpp.robot.g1.JsonizeVelocityCommand.from\_json}}\label{unitree-sdk2-cpp-robot-g1-jsonizevelocitycommand-from-json-1}}

计划从 JSON 风格字典读取字段并更新当前 SDK 值对象。

\textbf{签名}

\begin{lstlisting}[language=Python]
def from_json(self, json: dict[str, Any]) -> None
\end{lstlisting}

\textbf{可用性与安全性}

\textbf{\passthrough{\lstinline!SIGNATURE\_ONLY!} /
\passthrough{\lstinline!UNCLASSIFIED!}}：当前只有设计期签名，不能假设已存在于运行时扩展。

\textbf{参数}

\begin{longtable}[]{@{}
  >{\raggedright\arraybackslash}p{(\columnwidth - 6\tabcolsep) * \real{0.18}}
  >{\raggedright\arraybackslash}p{(\columnwidth - 6\tabcolsep) * \real{0.24}}
  >{\raggedright\arraybackslash}p{(\columnwidth - 6\tabcolsep) * \real{0.18}}
  >{\raggedright\arraybackslash}p{(\columnwidth - 6\tabcolsep) * \real{0.40}}@{}}
\toprule
名称 & Python 类型 & 默认值 & 含义 \\
\midrule
\endhead
\passthrough{\lstinline!json!} &
\passthrough{\lstinline!dict[str, Any]!} & 必填 & JSON
风格字典，键和值必须满足对应 C++ \passthrough{\lstinline!JsonMap!}
数据结构的约束。 对应 C++ 参数
\passthrough{\lstinline!json: common::JsonMap \&!}。 \\
\bottomrule
\end{longtable}

\textbf{返回值}

\begin{longtable}[]{@{}
  >{\raggedright\arraybackslash}p{(\columnwidth - 4\tabcolsep) * \real{0.18}}
  >{\raggedright\arraybackslash}p{(\columnwidth - 4\tabcolsep) * \real{0.20}}
  >{\raggedright\arraybackslash}p{(\columnwidth - 4\tabcolsep) * \real{0.62}}@{}}
\toprule
位置 & 类型 & 含义 \\
\midrule
\endhead
无 & \passthrough{\lstinline!None!} & 该方法不返回 Python
值；失败可能表现为异常或底层状态变化。 \\
\bottomrule
\end{longtable}

\textbf{对应 C++}

\begin{itemize}
\tightlist
\item
  类：\passthrough{\lstinline!unitree::robot::g1::JsonizeVelocityCommand!}
\item
  签名：\passthrough{\lstinline!fromJson(common::JsonMap \&)!}
\item
  绑定策略：\passthrough{\lstinline!SIGNATURE\_PREVIEW!}
\item
  声明位置：\passthrough{\lstinline!include/unitree/robot/g1/loco/g1\_loco\_api.hpp:62!}
\end{itemize}

\textbf{说明}

\begin{itemize}
\tightlist
\item
  示例是局部调用片段；\passthrough{\lstinline!obj!}
  和参数变量表示已经按上层业务校验并准备好的对象。
\item
  该条目可以被编辑器和类型检查器识别，但当前运行时可能在属性查找阶段直接失败。
\end{itemize}

\textbf{用法}

\begin{lstlisting}[language=Python]
from typing import TYPE_CHECKING

if TYPE_CHECKING:
    def planned_from_json(obj: JsonizeVelocityCommand, json: dict[str, Any]) -> None:
        obj.from_json(json=json)
\end{lstlisting}

\begin{quote}
{[}!CAUTION{]} 上面的 \passthrough{\lstinline!TYPE\_CHECKING!}
示例只表达计划调用形态。该分支在普通 Python
运行时为假，不代表当前扩展能执行此接口。
\end{quote}

\hypertarget{unitree-sdk2-cpp-robot-g1-jsonizevelocitycommand-to-json-1}{%
\paragraph{\texorpdfstring{\texttt{unitree\_sdk2\_cpp.robot.g1.JsonizeVelocityCommand.to\_json}}{unitree\_sdk2\_cpp.robot.g1.JsonizeVelocityCommand.to\_json}}\label{unitree-sdk2-cpp-robot-g1-jsonizevelocitycommand-to-json-1}}

计划把当前 SDK 值对象写入 JSON 风格字典。

\textbf{签名}

\begin{lstlisting}[language=Python]
def to_json(self, json: dict[str, Any]) -> None
\end{lstlisting}

\textbf{可用性与安全性}

\textbf{\passthrough{\lstinline!SIGNATURE\_ONLY!} /
\passthrough{\lstinline!UNCLASSIFIED!}}：当前只有设计期签名，不能假设已存在于运行时扩展。

\textbf{参数}

\begin{longtable}[]{@{}
  >{\raggedright\arraybackslash}p{(\columnwidth - 6\tabcolsep) * \real{0.18}}
  >{\raggedright\arraybackslash}p{(\columnwidth - 6\tabcolsep) * \real{0.24}}
  >{\raggedright\arraybackslash}p{(\columnwidth - 6\tabcolsep) * \real{0.18}}
  >{\raggedright\arraybackslash}p{(\columnwidth - 6\tabcolsep) * \real{0.40}}@{}}
\toprule
名称 & Python 类型 & 默认值 & 含义 \\
\midrule
\endhead
\passthrough{\lstinline!json!} &
\passthrough{\lstinline!dict[str, Any]!} & 必填 & JSON
风格字典，键和值必须满足对应 C++ \passthrough{\lstinline!JsonMap!}
数据结构的约束。 对应 C++ 参数
\passthrough{\lstinline!json: common::JsonMap \&!}。 \\
\bottomrule
\end{longtable}

\textbf{返回值}

\begin{longtable}[]{@{}
  >{\raggedright\arraybackslash}p{(\columnwidth - 4\tabcolsep) * \real{0.18}}
  >{\raggedright\arraybackslash}p{(\columnwidth - 4\tabcolsep) * \real{0.20}}
  >{\raggedright\arraybackslash}p{(\columnwidth - 4\tabcolsep) * \real{0.62}}@{}}
\toprule
位置 & 类型 & 含义 \\
\midrule
\endhead
无 & \passthrough{\lstinline!None!} & 该方法不返回 Python
值；失败可能表现为异常或底层状态变化。 \\
\bottomrule
\end{longtable}

\textbf{对应 C++}

\begin{itemize}
\tightlist
\item
  类：\passthrough{\lstinline!unitree::robot::g1::JsonizeVelocityCommand!}
\item
  签名：\passthrough{\lstinline!toJson(common::JsonMap \&) const!}
\item
  绑定策略：\passthrough{\lstinline!SIGNATURE\_PREVIEW!}
\item
  声明位置：\passthrough{\lstinline!include/unitree/robot/g1/loco/g1\_loco\_api.hpp:67!}
\end{itemize}

\textbf{说明}

\begin{itemize}
\tightlist
\item
  示例是局部调用片段；\passthrough{\lstinline!obj!}
  和参数变量表示已经按上层业务校验并准备好的对象。
\item
  该条目可以被编辑器和类型检查器识别，但当前运行时可能在属性查找阶段直接失败。
\end{itemize}

\textbf{用法}

\begin{lstlisting}[language=Python]
from typing import TYPE_CHECKING

if TYPE_CHECKING:
    def planned_to_json(obj: JsonizeVelocityCommand, json: dict[str, Any]) -> None:
        obj.to_json(json=json)
\end{lstlisting}

\begin{quote}
{[}!CAUTION{]} 上面的 \passthrough{\lstinline!TYPE\_CHECKING!}
示例只表达计划调用形态。该分支在普通 Python
运行时为假，不代表当前扩展能执行此接口。
\end{quote}

\hypertarget{unitree-sdk2-cpp-robot-g1-ledcontrolparameter}{%
\subsubsection{\texorpdfstring{\texttt{unitree\_sdk2\_cpp.robot.g1.LedControlParameter}}{unitree\_sdk2\_cpp.robot.g1.LedControlParameter}}\label{unitree-sdk2-cpp-robot-g1-ledcontrolparameter}}

SDK 请求参数值类型；当前可能仅提供设计期签名。

\textbf{导入}

\begin{lstlisting}[language=Python]
from unitree_sdk2_cpp.robot.g1 import LedControlParameter
\end{lstlisting}

\textbf{基类}：\passthrough{\lstinline!object!}

\textbf{公开属性}

\begin{longtable}[]{@{}
  >{\raggedright\arraybackslash}p{(\columnwidth - 6\tabcolsep) * \real{0.18}}
  >{\raggedright\arraybackslash}p{(\columnwidth - 6\tabcolsep) * \real{0.24}}
  >{\raggedright\arraybackslash}p{(\columnwidth - 6\tabcolsep) * \real{0.18}}
  >{\raggedright\arraybackslash}p{(\columnwidth - 6\tabcolsep) * \real{0.40}}@{}}
\toprule
名称 & Python 类型 & 含义 & 用法 \\
\midrule
\endhead
\passthrough{\lstinline!r!} & \passthrough{\lstinline!int!} &
传给该接口的 \passthrough{\lstinline!r!} 参数，Python 类型为
\passthrough{\lstinline!int!}。精确含义、单位、范围和枚举值需查目标型号对应头文件与协议。
& \passthrough{\lstinline!value = obj.r!} /
\passthrough{\lstinline!obj.r = value!} \\
\passthrough{\lstinline!g!} & \passthrough{\lstinline!int!} &
传给该接口的 \passthrough{\lstinline!g!} 参数，Python 类型为
\passthrough{\lstinline!int!}。精确含义、单位、范围和枚举值需查目标型号对应头文件与协议。
& \passthrough{\lstinline!value = obj.g!} /
\passthrough{\lstinline!obj.g = value!} \\
\passthrough{\lstinline!b!} & \passthrough{\lstinline!int!} &
传给该接口的 \passthrough{\lstinline!b!} 参数，Python 类型为
\passthrough{\lstinline!int!}。精确含义、单位、范围和枚举值需查目标型号对应头文件与协议。
& \passthrough{\lstinline!value = obj.b!} /
\passthrough{\lstinline!obj.b = value!} \\
\bottomrule
\end{longtable}

\textbf{方法索引}

\begin{longtable}[]{@{}
  >{\raggedright\arraybackslash}p{(\columnwidth - 6\tabcolsep) * \real{0.18}}
  >{\raggedright\arraybackslash}p{(\columnwidth - 6\tabcolsep) * \real{0.24}}
  >{\raggedright\arraybackslash}p{(\columnwidth - 6\tabcolsep) * \real{0.18}}
  >{\raggedright\arraybackslash}p{(\columnwidth - 6\tabcolsep) * \real{0.40}}@{}}
\toprule
方法 & Python 签名 & 状态 & 安全分类 \\
\midrule
\endhead
\protect\hyperlink{unitree-sdk2-cpp-robot-g1-ledcontrolparameter-dunder-init-1}{\passthrough{\lstinline!\_\_init\_\_!}}
& \passthrough{\lstinline!def \_\_init\_\_(self) -> None!} &
\passthrough{\lstinline!SIGNATURE\_ONLY!} &
\passthrough{\lstinline!CONSTRUCTION!} \\
\protect\hyperlink{unitree-sdk2-cpp-robot-g1-ledcontrolparameter-from-json-1}{\passthrough{\lstinline!from\_json!}}
&
\passthrough{\lstinline!def from\_json(self, json: dict[str, Any]) -> None!}
& \passthrough{\lstinline!SIGNATURE\_ONLY!} &
\passthrough{\lstinline!UNCLASSIFIED!} \\
\protect\hyperlink{unitree-sdk2-cpp-robot-g1-ledcontrolparameter-to-json-1}{\passthrough{\lstinline!to\_json!}}
&
\passthrough{\lstinline!def to\_json(self, json: dict[str, Any]) -> None!}
& \passthrough{\lstinline!SIGNATURE\_ONLY!} &
\passthrough{\lstinline!UNCLASSIFIED!} \\
\bottomrule
\end{longtable}

\hypertarget{unitree-sdk2-cpp-robot-g1-ledcontrolparameter-dunder-init-1}{%
\paragraph{\texorpdfstring{\texttt{unitree\_sdk2\_cpp.robot.g1.LedControlParameter.\_\_init\_\_}}{unitree\_sdk2\_cpp.robot.g1.LedControlParameter.\_\_init\_\_}}\label{unitree-sdk2-cpp-robot-g1-ledcontrolparameter-dunder-init-1}}

初始化 \passthrough{\lstinline!LedControlParameter!}
实例。是否能实际构造取决于下方可用性状态。

\textbf{签名}

\begin{lstlisting}[language=Python]
def __init__(self) -> None
\end{lstlisting}

\textbf{可用性与安全性}

\textbf{\passthrough{\lstinline!SIGNATURE\_ONLY!} /
\passthrough{\lstinline!CONSTRUCTION!}}：当前只有设计期签名，不能假设已存在于运行时扩展。

\textbf{参数}

无显式参数。实例方法中的 \passthrough{\lstinline!self!} 由 Python
自动传入。

\textbf{返回值}

\begin{longtable}[]{@{}
  >{\raggedright\arraybackslash}p{(\columnwidth - 4\tabcolsep) * \real{0.18}}
  >{\raggedright\arraybackslash}p{(\columnwidth - 4\tabcolsep) * \real{0.20}}
  >{\raggedright\arraybackslash}p{(\columnwidth - 4\tabcolsep) * \real{0.62}}@{}}
\toprule
位置 & 类型 & 含义 \\
\midrule
\endhead
无 & \passthrough{\lstinline!None!} &
\passthrough{\lstinline!\_\_init\_\_!}
本身不返回值；调用类对象时，在构造函数可用的前提下得到该类实例。 \\
\bottomrule
\end{longtable}

\textbf{对应 C++}

\begin{itemize}
\tightlist
\item
  类：\passthrough{\lstinline!unitree::robot::g1::LedControlParameter!}
\item
  签名：\passthrough{\lstinline!LedControlParameter()!}
\item
  绑定策略：\passthrough{\lstinline!未单独标注!}
\item
  声明位置：\passthrough{\lstinline!include/unitree/robot/g1/audio/g1\_audio\_api.hpp:75!}
\end{itemize}

\textbf{说明}

\begin{itemize}
\tightlist
\item
  示例是局部调用片段；\passthrough{\lstinline!obj!}
  和参数变量表示已经按上层业务校验并准备好的对象。
\item
  该条目可以被编辑器和类型检查器识别，但当前运行时可能在属性查找阶段直接失败。
\end{itemize}

\textbf{用法}

\begin{lstlisting}[language=Python]
from typing import TYPE_CHECKING

if TYPE_CHECKING:
    def planned_construct() -> LedControlParameter:
        return LedControlParameter()
\end{lstlisting}

\begin{quote}
{[}!CAUTION{]} 上面的 \passthrough{\lstinline!TYPE\_CHECKING!}
示例只表达计划调用形态。该分支在普通 Python
运行时为假，不代表当前扩展能执行此接口。
\end{quote}

\hypertarget{unitree-sdk2-cpp-robot-g1-ledcontrolparameter-from-json-1}{%
\paragraph{\texorpdfstring{\texttt{unitree\_sdk2\_cpp.robot.g1.LedControlParameter.from\_json}}{unitree\_sdk2\_cpp.robot.g1.LedControlParameter.from\_json}}\label{unitree-sdk2-cpp-robot-g1-ledcontrolparameter-from-json-1}}

计划从 JSON 风格字典读取字段并更新当前 SDK 值对象。

\textbf{签名}

\begin{lstlisting}[language=Python]
def from_json(self, json: dict[str, Any]) -> None
\end{lstlisting}

\textbf{可用性与安全性}

\textbf{\passthrough{\lstinline!SIGNATURE\_ONLY!} /
\passthrough{\lstinline!UNCLASSIFIED!}}：当前只有设计期签名，不能假设已存在于运行时扩展。

\textbf{参数}

\begin{longtable}[]{@{}
  >{\raggedright\arraybackslash}p{(\columnwidth - 6\tabcolsep) * \real{0.18}}
  >{\raggedright\arraybackslash}p{(\columnwidth - 6\tabcolsep) * \real{0.24}}
  >{\raggedright\arraybackslash}p{(\columnwidth - 6\tabcolsep) * \real{0.18}}
  >{\raggedright\arraybackslash}p{(\columnwidth - 6\tabcolsep) * \real{0.40}}@{}}
\toprule
名称 & Python 类型 & 默认值 & 含义 \\
\midrule
\endhead
\passthrough{\lstinline!json!} &
\passthrough{\lstinline!dict[str, Any]!} & 必填 & JSON
风格字典，键和值必须满足对应 C++ \passthrough{\lstinline!JsonMap!}
数据结构的约束。 对应 C++ 参数
\passthrough{\lstinline!json: common::JsonMap \&!}。 \\
\bottomrule
\end{longtable}

\textbf{返回值}

\begin{longtable}[]{@{}
  >{\raggedright\arraybackslash}p{(\columnwidth - 4\tabcolsep) * \real{0.18}}
  >{\raggedright\arraybackslash}p{(\columnwidth - 4\tabcolsep) * \real{0.20}}
  >{\raggedright\arraybackslash}p{(\columnwidth - 4\tabcolsep) * \real{0.62}}@{}}
\toprule
位置 & 类型 & 含义 \\
\midrule
\endhead
无 & \passthrough{\lstinline!None!} & 该方法不返回 Python
值；失败可能表现为异常或底层状态变化。 \\
\bottomrule
\end{longtable}

\textbf{对应 C++}

\begin{itemize}
\tightlist
\item
  类：\passthrough{\lstinline!unitree::robot::g1::LedControlParameter!}
\item
  签名：\passthrough{\lstinline!fromJson(common::JsonMap \&)!}
\item
  绑定策略：\passthrough{\lstinline!SIGNATURE\_PREVIEW!}
\item
  声明位置：\passthrough{\lstinline!include/unitree/robot/g1/audio/g1\_audio\_api.hpp:78!}
\end{itemize}

\textbf{说明}

\begin{itemize}
\tightlist
\item
  示例是局部调用片段；\passthrough{\lstinline!obj!}
  和参数变量表示已经按上层业务校验并准备好的对象。
\item
  该条目可以被编辑器和类型检查器识别，但当前运行时可能在属性查找阶段直接失败。
\end{itemize}

\textbf{用法}

\begin{lstlisting}[language=Python]
from typing import TYPE_CHECKING

if TYPE_CHECKING:
    def planned_from_json(obj: LedControlParameter, json: dict[str, Any]) -> None:
        obj.from_json(json=json)
\end{lstlisting}

\begin{quote}
{[}!CAUTION{]} 上面的 \passthrough{\lstinline!TYPE\_CHECKING!}
示例只表达计划调用形态。该分支在普通 Python
运行时为假，不代表当前扩展能执行此接口。
\end{quote}

\hypertarget{unitree-sdk2-cpp-robot-g1-ledcontrolparameter-to-json-1}{%
\paragraph{\texorpdfstring{\texttt{unitree\_sdk2\_cpp.robot.g1.LedControlParameter.to\_json}}{unitree\_sdk2\_cpp.robot.g1.LedControlParameter.to\_json}}\label{unitree-sdk2-cpp-robot-g1-ledcontrolparameter-to-json-1}}

计划把当前 SDK 值对象写入 JSON 风格字典。

\textbf{签名}

\begin{lstlisting}[language=Python]
def to_json(self, json: dict[str, Any]) -> None
\end{lstlisting}

\textbf{可用性与安全性}

\textbf{\passthrough{\lstinline!SIGNATURE\_ONLY!} /
\passthrough{\lstinline!UNCLASSIFIED!}}：当前只有设计期签名，不能假设已存在于运行时扩展。

\textbf{参数}

\begin{longtable}[]{@{}
  >{\raggedright\arraybackslash}p{(\columnwidth - 6\tabcolsep) * \real{0.18}}
  >{\raggedright\arraybackslash}p{(\columnwidth - 6\tabcolsep) * \real{0.24}}
  >{\raggedright\arraybackslash}p{(\columnwidth - 6\tabcolsep) * \real{0.18}}
  >{\raggedright\arraybackslash}p{(\columnwidth - 6\tabcolsep) * \real{0.40}}@{}}
\toprule
名称 & Python 类型 & 默认值 & 含义 \\
\midrule
\endhead
\passthrough{\lstinline!json!} &
\passthrough{\lstinline!dict[str, Any]!} & 必填 & JSON
风格字典，键和值必须满足对应 C++ \passthrough{\lstinline!JsonMap!}
数据结构的约束。 对应 C++ 参数
\passthrough{\lstinline!json: common::JsonMap \&!}。 \\
\bottomrule
\end{longtable}

\textbf{返回值}

\begin{longtable}[]{@{}
  >{\raggedright\arraybackslash}p{(\columnwidth - 4\tabcolsep) * \real{0.18}}
  >{\raggedright\arraybackslash}p{(\columnwidth - 4\tabcolsep) * \real{0.20}}
  >{\raggedright\arraybackslash}p{(\columnwidth - 4\tabcolsep) * \real{0.62}}@{}}
\toprule
位置 & 类型 & 含义 \\
\midrule
\endhead
无 & \passthrough{\lstinline!None!} & 该方法不返回 Python
值；失败可能表现为异常或底层状态变化。 \\
\bottomrule
\end{longtable}

\textbf{对应 C++}

\begin{itemize}
\tightlist
\item
  类：\passthrough{\lstinline!unitree::robot::g1::LedControlParameter!}
\item
  签名：\passthrough{\lstinline!toJson(common::JsonMap \&) const!}
\item
  绑定策略：\passthrough{\lstinline!SIGNATURE\_PREVIEW!}
\item
  声明位置：\passthrough{\lstinline!include/unitree/robot/g1/audio/g1\_audio\_api.hpp:80!}
\end{itemize}

\textbf{说明}

\begin{itemize}
\tightlist
\item
  示例是局部调用片段；\passthrough{\lstinline!obj!}
  和参数变量表示已经按上层业务校验并准备好的对象。
\item
  该条目可以被编辑器和类型检查器识别，但当前运行时可能在属性查找阶段直接失败。
\end{itemize}

\textbf{用法}

\begin{lstlisting}[language=Python]
from typing import TYPE_CHECKING

if TYPE_CHECKING:
    def planned_to_json(obj: LedControlParameter, json: dict[str, Any]) -> None:
        obj.to_json(json=json)
\end{lstlisting}

\begin{quote}
{[}!CAUTION{]} 上面的 \passthrough{\lstinline!TYPE\_CHECKING!}
示例只表达计划调用形态。该分支在普通 Python
运行时为假，不代表当前扩展能执行此接口。
\end{quote}

\hypertarget{unitree-sdk2-cpp-robot-g1-lococlient}{%
\subsubsection{\texorpdfstring{\texttt{unitree\_sdk2\_cpp.robot.g1.LocoClient}}{unitree\_sdk2\_cpp.robot.g1.LocoClient}}\label{unitree-sdk2-cpp-robot-g1-lococlient}}

Robot 服务客户端。构造、初始化和具体方法的可用性必须分别检查。

\textbf{导入}

\begin{lstlisting}[language=Python]
from unitree_sdk2_cpp.robot.g1 import LocoClient
\end{lstlisting}

\textbf{基类}：\passthrough{\lstinline!Client!}

\textbf{方法索引}

\begin{longtable}[]{@{}
  >{\raggedright\arraybackslash}p{(\columnwidth - 6\tabcolsep) * \real{0.18}}
  >{\raggedright\arraybackslash}p{(\columnwidth - 6\tabcolsep) * \real{0.24}}
  >{\raggedright\arraybackslash}p{(\columnwidth - 6\tabcolsep) * \real{0.18}}
  >{\raggedright\arraybackslash}p{(\columnwidth - 6\tabcolsep) * \real{0.40}}@{}}
\toprule
方法 & Python 签名 & 状态 & 安全分类 \\
\midrule
\endhead
\protect\hyperlink{unitree-sdk2-cpp-robot-g1-lococlient-dunder-init-1}{\passthrough{\lstinline!\_\_init\_\_!}}
& \passthrough{\lstinline!def \_\_init\_\_(self) -> None!} &
\passthrough{\lstinline!AVAILABLE!} &
\passthrough{\lstinline!CONSTRUCTION!} \\
\protect\hyperlink{unitree-sdk2-cpp-robot-g1-lococlient-init-1}{\passthrough{\lstinline!init!}}
& \passthrough{\lstinline!def init(self) -> None!} &
\passthrough{\lstinline!AVAILABLE!} &
\passthrough{\lstinline!INITIALIZATION!} \\
\protect\hyperlink{unitree-sdk2-cpp-robot-g1-lococlient-get-fsm-id-1}{\passthrough{\lstinline!get\_fsm\_id!}}
& \passthrough{\lstinline!def get\_fsm\_id(self) -> tuple[int, int]!} &
\passthrough{\lstinline!AVAILABLE!} &
\passthrough{\lstinline!READ\_ONLY!} \\
\protect\hyperlink{unitree-sdk2-cpp-robot-g1-lococlient-get-fsm-mode-1}{\passthrough{\lstinline!get\_fsm\_mode!}}
& \passthrough{\lstinline!def get\_fsm\_mode(self) -> tuple[int, int]!}
& \passthrough{\lstinline!AVAILABLE!} &
\passthrough{\lstinline!READ\_ONLY!} \\
\protect\hyperlink{unitree-sdk2-cpp-robot-g1-lococlient-get-balance-mode-1}{\passthrough{\lstinline!get\_balance\_mode!}}
&
\passthrough{\lstinline!def get\_balance\_mode(self) -> tuple[int, int]!}
& \passthrough{\lstinline!AVAILABLE!} &
\passthrough{\lstinline!READ\_ONLY!} \\
\protect\hyperlink{unitree-sdk2-cpp-robot-g1-lococlient-get-swing-height-1}{\passthrough{\lstinline!get\_swing\_height!}}
&
\passthrough{\lstinline!def get\_swing\_height(self) -> tuple[int, float]!}
& \passthrough{\lstinline!AVAILABLE!} &
\passthrough{\lstinline!READ\_ONLY!} \\
\protect\hyperlink{unitree-sdk2-cpp-robot-g1-lococlient-get-stand-height-1}{\passthrough{\lstinline!get\_stand\_height!}}
&
\passthrough{\lstinline!def get\_stand\_height(self) -> tuple[int, float]!}
& \passthrough{\lstinline!AVAILABLE!} &
\passthrough{\lstinline!READ\_ONLY!} \\
\protect\hyperlink{unitree-sdk2-cpp-robot-g1-lococlient-get-phase-1}{\passthrough{\lstinline!get\_phase!}}
&
\passthrough{\lstinline!def get\_phase(self) -> tuple[int, list[float]]!}
& \passthrough{\lstinline!AVAILABLE!} &
\passthrough{\lstinline!READ\_ONLY!} \\
\protect\hyperlink{unitree-sdk2-cpp-robot-g1-lococlient-set-fsm-id-1}{\passthrough{\lstinline!set\_fsm\_id!}}
& \passthrough{\lstinline!def set\_fsm\_id(self, fsm\_id: int) -> int!}
& \passthrough{\lstinline!SIGNATURE\_ONLY!} &
\passthrough{\lstinline!MOTION\_COMMAND!} \\
\protect\hyperlink{unitree-sdk2-cpp-robot-g1-lococlient-set-balance-mode-1}{\passthrough{\lstinline!set\_balance\_mode!}}
&
\passthrough{\lstinline!def set\_balance\_mode(self, balance\_mode: int) -> int!}
& \passthrough{\lstinline!SIGNATURE\_ONLY!} &
\passthrough{\lstinline!MOTION\_COMMAND!} \\
\protect\hyperlink{unitree-sdk2-cpp-robot-g1-lococlient-set-swing-height-1}{\passthrough{\lstinline!set\_swing\_height!}}
&
\passthrough{\lstinline!def set\_swing\_height(self, swing\_height: float) -> int!}
& \passthrough{\lstinline!SIGNATURE\_ONLY!} &
\passthrough{\lstinline!MOTION\_COMMAND!} \\
\protect\hyperlink{unitree-sdk2-cpp-robot-g1-lococlient-set-stand-height-1}{\passthrough{\lstinline!set\_stand\_height!}}
&
\passthrough{\lstinline!def set\_stand\_height(self, stand\_height: float) -> int!}
& \passthrough{\lstinline!SIGNATURE\_ONLY!} &
\passthrough{\lstinline!MOTION\_COMMAND!} \\
\protect\hyperlink{unitree-sdk2-cpp-robot-g1-lococlient-set-velocity-1}{\passthrough{\lstinline!set\_velocity!}}
&
\passthrough{\lstinline!def set\_velocity(self, vx: float, vy: float, omega: float, duration: float = ...) -> int!}
& \passthrough{\lstinline!SIGNATURE\_ONLY!} &
\passthrough{\lstinline!MOTION\_COMMAND!} \\
\protect\hyperlink{unitree-sdk2-cpp-robot-g1-lococlient-set-task-id-1}{\passthrough{\lstinline!set\_task\_id!}}
&
\passthrough{\lstinline!def set\_task\_id(self, task\_id: int) -> int!}
& \passthrough{\lstinline!SIGNATURE\_ONLY!} &
\passthrough{\lstinline!MOTION\_COMMAND!} \\
\protect\hyperlink{unitree-sdk2-cpp-robot-g1-lococlient-switch-to-user-ctrl-1}{\passthrough{\lstinline!switch\_to\_user\_ctrl!}}
& \passthrough{\lstinline!def switch\_to\_user\_ctrl(self) -> int!} &
\passthrough{\lstinline!SIGNATURE\_ONLY!} &
\passthrough{\lstinline!MOTION\_COMMAND!} \\
\protect\hyperlink{unitree-sdk2-cpp-robot-g1-lococlient-switch-to-internal-ctrl-1}{\passthrough{\lstinline!switch\_to\_internal\_ctrl!}}
&
\passthrough{\lstinline!def switch\_to\_internal\_ctrl(self, mode: InternalFsmMode) -> int!}
& \passthrough{\lstinline!SIGNATURE\_ONLY!} &
\passthrough{\lstinline!MOTION\_COMMAND!} \\
\protect\hyperlink{unitree-sdk2-cpp-robot-g1-lococlient-damp-1}{\passthrough{\lstinline!damp!}}
& \passthrough{\lstinline!def damp(self) -> int!} &
\passthrough{\lstinline!SIGNATURE\_ONLY!} &
\passthrough{\lstinline!MOTION\_COMMAND!} \\
\protect\hyperlink{unitree-sdk2-cpp-robot-g1-lococlient-start-1}{\passthrough{\lstinline!start!}}
& \passthrough{\lstinline!def start(self) -> int!} &
\passthrough{\lstinline!SIGNATURE\_ONLY!} &
\passthrough{\lstinline!MOTION\_COMMAND!} \\
\protect\hyperlink{unitree-sdk2-cpp-robot-g1-lococlient-squat-1}{\passthrough{\lstinline!squat!}}
& \passthrough{\lstinline!def squat(self) -> int!} &
\passthrough{\lstinline!SIGNATURE\_ONLY!} &
\passthrough{\lstinline!MOTION\_COMMAND!} \\
\protect\hyperlink{unitree-sdk2-cpp-robot-g1-lococlient-sit-1}{\passthrough{\lstinline!sit!}}
& \passthrough{\lstinline!def sit(self) -> int!} &
\passthrough{\lstinline!SIGNATURE\_ONLY!} &
\passthrough{\lstinline!MOTION\_COMMAND!} \\
\protect\hyperlink{unitree-sdk2-cpp-robot-g1-lococlient-stand-up-1}{\passthrough{\lstinline!stand\_up!}}
& \passthrough{\lstinline!def stand\_up(self) -> int!} &
\passthrough{\lstinline!SIGNATURE\_ONLY!} &
\passthrough{\lstinline!MOTION\_COMMAND!} \\
\protect\hyperlink{unitree-sdk2-cpp-robot-g1-lococlient-zero-torque-1}{\passthrough{\lstinline!zero\_torque!}}
& \passthrough{\lstinline!def zero\_torque(self) -> int!} &
\passthrough{\lstinline!SIGNATURE\_ONLY!} &
\passthrough{\lstinline!MOTION\_COMMAND!} \\
\protect\hyperlink{unitree-sdk2-cpp-robot-g1-lococlient-stop-move-1}{\passthrough{\lstinline!stop\_move!}}
& \passthrough{\lstinline!def stop\_move(self) -> int!} &
\passthrough{\lstinline!SIGNATURE\_ONLY!} &
\passthrough{\lstinline!MOTION\_COMMAND!} \\
\protect\hyperlink{unitree-sdk2-cpp-robot-g1-lococlient-high-stand-1}{\passthrough{\lstinline!high\_stand!}}
& \passthrough{\lstinline!def high\_stand(self) -> int!} &
\passthrough{\lstinline!SIGNATURE\_ONLY!} &
\passthrough{\lstinline!MOTION\_COMMAND!} \\
\protect\hyperlink{unitree-sdk2-cpp-robot-g1-lococlient-low-stand-1}{\passthrough{\lstinline!low\_stand!}}
& \passthrough{\lstinline!def low\_stand(self) -> int!} &
\passthrough{\lstinline!SIGNATURE\_ONLY!} &
\passthrough{\lstinline!MOTION\_COMMAND!} \\
\protect\hyperlink{unitree-sdk2-cpp-robot-g1-lococlient-move-1}{\passthrough{\lstinline!move（重载 1/2）!}}
&
\passthrough{\lstinline!def move(self, vx: float, vy: float, vyaw: float, continous\_move: bool) -> int!}
& \passthrough{\lstinline!SIGNATURE\_ONLY!} &
\passthrough{\lstinline!MOTION\_COMMAND!} \\
\protect\hyperlink{unitree-sdk2-cpp-robot-g1-lococlient-move-2}{\passthrough{\lstinline!move（重载 2/2）!}}
&
\passthrough{\lstinline!def move(self, vx: float, vy: float, vyaw: float) -> int!}
& \passthrough{\lstinline!SIGNATURE\_ONLY!} &
\passthrough{\lstinline!MOTION\_COMMAND!} \\
\protect\hyperlink{unitree-sdk2-cpp-robot-g1-lococlient-balance-stand-1}{\passthrough{\lstinline!balance\_stand!}}
& \passthrough{\lstinline!def balance\_stand(self) -> int!} &
\passthrough{\lstinline!SIGNATURE\_ONLY!} &
\passthrough{\lstinline!MOTION\_COMMAND!} \\
\protect\hyperlink{unitree-sdk2-cpp-robot-g1-lococlient-continuous-gait-1}{\passthrough{\lstinline!continuous\_gait!}}
&
\passthrough{\lstinline!def continuous\_gait(self, flag: bool) -> int!}
& \passthrough{\lstinline!SIGNATURE\_ONLY!} &
\passthrough{\lstinline!MOTION\_COMMAND!} \\
\protect\hyperlink{unitree-sdk2-cpp-robot-g1-lococlient-switch-move-mode-1}{\passthrough{\lstinline!switch\_move\_mode!}}
&
\passthrough{\lstinline!def switch\_move\_mode(self, flag: bool) -> int!}
& \passthrough{\lstinline!SIGNATURE\_ONLY!} &
\passthrough{\lstinline!MOTION\_COMMAND!} \\
\protect\hyperlink{unitree-sdk2-cpp-robot-g1-lococlient-wave-hand-1}{\passthrough{\lstinline!wave\_hand!}}
&
\passthrough{\lstinline!def wave\_hand(self, turn\_flag: bool = ...) -> int!}
& \passthrough{\lstinline!SIGNATURE\_ONLY!} &
\passthrough{\lstinline!MOTION\_COMMAND!} \\
\protect\hyperlink{unitree-sdk2-cpp-robot-g1-lococlient-shake-hand-1}{\passthrough{\lstinline!shake\_hand!}}
&
\passthrough{\lstinline!def shake\_hand(self, stage: int = ...) -> int!}
& \passthrough{\lstinline!SIGNATURE\_ONLY!} &
\passthrough{\lstinline!MOTION\_COMMAND!} \\
\protect\hyperlink{unitree-sdk2-cpp-robot-g1-lococlient-set-speed-mode-1}{\passthrough{\lstinline!set\_speed\_mode!}}
&
\passthrough{\lstinline!def set\_speed\_mode(self, speed\_mode: int) -> int!}
& \passthrough{\lstinline!SIGNATURE\_ONLY!} &
\passthrough{\lstinline!MOTION\_COMMAND!} \\
\protect\hyperlink{unitree-sdk2-cpp-robot-g1-lococlient-get-mimic-motion-1}{\passthrough{\lstinline!get\_mimic\_motion!}}
&
\passthrough{\lstinline!def get\_mimic\_motion(self) -> tuple[int, str]!}
& \passthrough{\lstinline!AVAILABLE!} &
\passthrough{\lstinline!READ\_ONLY!} \\
\protect\hyperlink{unitree-sdk2-cpp-robot-g1-lococlient-fsm-api-1}{\passthrough{\lstinline!\_fsm\_api!}}
&
\passthrough{\lstinline!def \_fsm\_api(self, parameter: str) -> tuple[int, str]!}
& \passthrough{\lstinline!SIGNATURE\_ONLY!} &
\passthrough{\lstinline!MOTION\_COMMAND!} \\
\bottomrule
\end{longtable}

\hypertarget{unitree-sdk2-cpp-robot-g1-lococlient-dunder-init-1}{%
\paragraph{\texorpdfstring{\texttt{unitree\_sdk2\_cpp.robot.g1.LocoClient.\_\_init\_\_}}{unitree\_sdk2\_cpp.robot.g1.LocoClient.\_\_init\_\_}}\label{unitree-sdk2-cpp-robot-g1-lococlient-dunder-init-1}}

初始化 \passthrough{\lstinline!LocoClient!}
实例。是否能实际构造取决于下方可用性状态。

\textbf{签名}

\begin{lstlisting}[language=Python]
def __init__(self) -> None
\end{lstlisting}

\textbf{可用性与安全性}

\textbf{\passthrough{\lstinline!AVAILABLE!} /
\passthrough{\lstinline!CONSTRUCTION!}}：当前绑定源码已实现；实际执行条件仍取决于平台和对应
SDK 环境。

\textbf{参数}

无显式参数。实例方法中的 \passthrough{\lstinline!self!} 由 Python
自动传入。

\textbf{返回值}

\begin{longtable}[]{@{}
  >{\raggedright\arraybackslash}p{(\columnwidth - 4\tabcolsep) * \real{0.18}}
  >{\raggedright\arraybackslash}p{(\columnwidth - 4\tabcolsep) * \real{0.20}}
  >{\raggedright\arraybackslash}p{(\columnwidth - 4\tabcolsep) * \real{0.62}}@{}}
\toprule
位置 & 类型 & 含义 \\
\midrule
\endhead
无 & \passthrough{\lstinline!None!} &
\passthrough{\lstinline!\_\_init\_\_!}
本身不返回值；调用类对象时，在构造函数可用的前提下得到该类实例。 \\
\bottomrule
\end{longtable}

\textbf{对应 C++}

\begin{itemize}
\tightlist
\item
  类：\passthrough{\lstinline!unitree::robot::g1::LocoClient!}
\item
  签名：\passthrough{\lstinline!LocoClient()!}
\item
  绑定策略：\passthrough{\lstinline!未单独标注!}
\item
  声明位置：\passthrough{\lstinline!include/unitree/robot/g1/loco/g1\_loco\_client.hpp:14!}
\end{itemize}

\textbf{说明}

\begin{itemize}
\tightlist
\item
  示例是局部调用片段；\passthrough{\lstinline!obj!}
  和参数变量表示已经按上层业务校验并准备好的对象。
\end{itemize}

\textbf{用法}

\begin{lstlisting}[language=Python]
value = LocoClient()
\end{lstlisting}

\hypertarget{unitree-sdk2-cpp-robot-g1-lococlient-init-1}{%
\paragraph{\texorpdfstring{\texttt{unitree\_sdk2\_cpp.robot.g1.LocoClient.init}}{unitree\_sdk2\_cpp.robot.g1.LocoClient.init}}\label{unitree-sdk2-cpp-robot-g1-lococlient-init-1}}

初始化当前 SDK 对象所需的底层通道或服务资源。

\textbf{签名}

\begin{lstlisting}[language=Python]
def init(self) -> None
\end{lstlisting}

\textbf{可用性与安全性}

\textbf{\passthrough{\lstinline!AVAILABLE!} /
\passthrough{\lstinline!INITIALIZATION!}}：当前绑定源码已实现；实际执行条件仍取决于平台和对应
SDK 环境。

\textbf{参数}

无显式参数。实例方法中的 \passthrough{\lstinline!self!} 由 Python
自动传入。

\textbf{返回值}

\begin{longtable}[]{@{}
  >{\raggedright\arraybackslash}p{(\columnwidth - 4\tabcolsep) * \real{0.18}}
  >{\raggedright\arraybackslash}p{(\columnwidth - 4\tabcolsep) * \real{0.20}}
  >{\raggedright\arraybackslash}p{(\columnwidth - 4\tabcolsep) * \real{0.62}}@{}}
\toprule
位置 & 类型 & 含义 \\
\midrule
\endhead
无 & \passthrough{\lstinline!None!} & 该方法不返回 Python
值；失败可能表现为异常或底层状态变化。 \\
\bottomrule
\end{longtable}

\textbf{对应 C++}

\begin{itemize}
\tightlist
\item
  类：\passthrough{\lstinline!unitree::robot::g1::LocoClient!}
\item
  签名：\passthrough{\lstinline!Init()!}
\item
  绑定策略：\passthrough{\lstinline!DIRECT!}
\item
  声明位置：\passthrough{\lstinline!include/unitree/robot/g1/loco/g1\_loco\_client.hpp:18!}
\end{itemize}

\textbf{说明}

\begin{itemize}
\tightlist
\item
  示例是局部调用片段；\passthrough{\lstinline!obj!}
  和参数变量表示已经按上层业务校验并准备好的对象。
\end{itemize}

\textbf{用法}

\begin{lstlisting}[language=Python]
obj.init()
\end{lstlisting}

\hypertarget{unitree-sdk2-cpp-robot-g1-lococlient-get-fsm-id-1}{%
\paragraph{\texorpdfstring{\texttt{unitree\_sdk2\_cpp.robot.g1.LocoClient.get\_fsm\_id}}{unitree\_sdk2\_cpp.robot.g1.LocoClient.get\_fsm\_id}}\label{unitree-sdk2-cpp-robot-g1-lococlient-get-fsm-id-1}}

查询或检查 FSM ID。

\textbf{签名}

\begin{lstlisting}[language=Python]
def get_fsm_id(self) -> tuple[int, int]
\end{lstlisting}

\textbf{可用性与安全性}

\textbf{\passthrough{\lstinline!AVAILABLE!} /
\passthrough{\lstinline!READ\_ONLY!}}：当前绑定源码已实现只读查询。Robot
Client 的服务端查询通常仍需匹配的 Linux
扩展、DDS、网络、机器人和服务；纯本地 getter 则不一定需要全部条件。

\textbf{参数}

无显式参数。实例方法中的 \passthrough{\lstinline!self!} 由 Python
自动传入。

\textbf{返回值}

\begin{longtable}[]{@{}
  >{\raggedright\arraybackslash}p{(\columnwidth - 4\tabcolsep) * \real{0.18}}
  >{\raggedright\arraybackslash}p{(\columnwidth - 4\tabcolsep) * \real{0.20}}
  >{\raggedright\arraybackslash}p{(\columnwidth - 4\tabcolsep) * \real{0.62}}@{}}
\toprule
位置 & 类型 & 含义 \\
\midrule
\endhead
{[}0{]} \passthrough{\lstinline!status!} & \passthrough{\lstinline!int!}
& SDK 状态码；Unitree 示例通常以 \passthrough{\lstinline!0!}
表示成功，非零值需按具体服务错误码解释。 \\
{[}1{]} \passthrough{\lstinline!fsm\_id!} &
\passthrough{\lstinline!int!} & 传给该接口的 FSM ID 参数，Python 类型为
\passthrough{\lstinline!int!}。精确含义、单位、范围和枚举值需查目标型号对应头文件与协议。
对应 C++ 参数 \passthrough{\lstinline!fsm\_id: int \&!}。 \\
\bottomrule
\end{longtable}

\textbf{对应 C++}

\begin{itemize}
\tightlist
\item
  类：\passthrough{\lstinline!unitree::robot::g1::LocoClient!}
\item
  签名：\passthrough{\lstinline!GetFsmId(int \&)!}
\item
  绑定策略：\passthrough{\lstinline!OUTPUT\_WRAPPER!}
\item
  声明位置：\passthrough{\lstinline!include/unitree/robot/g1/loco/g1\_loco\_client.hpp:40!}
\end{itemize}

\textbf{说明}

\begin{itemize}
\tightlist
\item
  示例是局部调用片段；\passthrough{\lstinline!obj!}
  和参数变量表示已经按上层业务校验并准备好的对象。
\item
  C++ 的可修改输出引用不会作为 Python 入参出现，而是按 C++
  参数顺序追加到返回元组中。
\end{itemize}

\textbf{用法}

\begin{lstlisting}[language=Python]
status, fsm_id = obj.get_fsm_id()
\end{lstlisting}

\begin{quote}
{[}!NOTE{]} 只读查询仍会访问
DDS/机器人服务。必须检查返回状态码，并处理超时和网络中断。
\end{quote}

\hypertarget{unitree-sdk2-cpp-robot-g1-lococlient-get-fsm-mode-1}{%
\paragraph{\texorpdfstring{\texttt{unitree\_sdk2\_cpp.robot.g1.LocoClient.get\_fsm\_mode}}{unitree\_sdk2\_cpp.robot.g1.LocoClient.get\_fsm\_mode}}\label{unitree-sdk2-cpp-robot-g1-lococlient-get-fsm-mode-1}}

查询或检查 FSM 模式。

\textbf{签名}

\begin{lstlisting}[language=Python]
def get_fsm_mode(self) -> tuple[int, int]
\end{lstlisting}

\textbf{可用性与安全性}

\textbf{\passthrough{\lstinline!AVAILABLE!} /
\passthrough{\lstinline!READ\_ONLY!}}：当前绑定源码已实现只读查询。Robot
Client 的服务端查询通常仍需匹配的 Linux
扩展、DDS、网络、机器人和服务；纯本地 getter 则不一定需要全部条件。

\textbf{参数}

无显式参数。实例方法中的 \passthrough{\lstinline!self!} 由 Python
自动传入。

\textbf{返回值}

\begin{longtable}[]{@{}
  >{\raggedright\arraybackslash}p{(\columnwidth - 4\tabcolsep) * \real{0.18}}
  >{\raggedright\arraybackslash}p{(\columnwidth - 4\tabcolsep) * \real{0.20}}
  >{\raggedright\arraybackslash}p{(\columnwidth - 4\tabcolsep) * \real{0.62}}@{}}
\toprule
位置 & 类型 & 含义 \\
\midrule
\endhead
{[}0{]} \passthrough{\lstinline!status!} & \passthrough{\lstinline!int!}
& SDK 状态码；Unitree 示例通常以 \passthrough{\lstinline!0!}
表示成功，非零值需按具体服务错误码解释。 \\
{[}1{]} \passthrough{\lstinline!fsm\_mode!} &
\passthrough{\lstinline!int!} & 传给该接口的 FSM 模式 参数，Python
类型为
\passthrough{\lstinline!int!}。精确含义、单位、范围和枚举值需查目标型号对应头文件与协议。
对应 C++ 参数 \passthrough{\lstinline!fsm\_mode: int \&!}。 \\
\bottomrule
\end{longtable}

\textbf{对应 C++}

\begin{itemize}
\tightlist
\item
  类：\passthrough{\lstinline!unitree::robot::g1::LocoClient!}
\item
  签名：\passthrough{\lstinline!GetFsmMode(int \&)!}
\item
  绑定策略：\passthrough{\lstinline!OUTPUT\_WRAPPER!}
\item
  声明位置：\passthrough{\lstinline!include/unitree/robot/g1/loco/g1\_loco\_client.hpp:54!}
\end{itemize}

\textbf{说明}

\begin{itemize}
\tightlist
\item
  示例是局部调用片段；\passthrough{\lstinline!obj!}
  和参数变量表示已经按上层业务校验并准备好的对象。
\item
  C++ 的可修改输出引用不会作为 Python 入参出现，而是按 C++
  参数顺序追加到返回元组中。
\end{itemize}

\textbf{用法}

\begin{lstlisting}[language=Python]
status, fsm_mode = obj.get_fsm_mode()
\end{lstlisting}

\begin{quote}
{[}!NOTE{]} 只读查询仍会访问
DDS/机器人服务。必须检查返回状态码，并处理超时和网络中断。
\end{quote}

\hypertarget{unitree-sdk2-cpp-robot-g1-lococlient-get-balance-mode-1}{%
\paragraph{\texorpdfstring{\texttt{unitree\_sdk2\_cpp.robot.g1.LocoClient.get\_balance\_mode}}{unitree\_sdk2\_cpp.robot.g1.LocoClient.get\_balance\_mode}}\label{unitree-sdk2-cpp-robot-g1-lococlient-get-balance-mode-1}}

查询或检查 平衡 模式。

\textbf{签名}

\begin{lstlisting}[language=Python]
def get_balance_mode(self) -> tuple[int, int]
\end{lstlisting}

\textbf{可用性与安全性}

\textbf{\passthrough{\lstinline!AVAILABLE!} /
\passthrough{\lstinline!READ\_ONLY!}}：当前绑定源码已实现只读查询。Robot
Client 的服务端查询通常仍需匹配的 Linux
扩展、DDS、网络、机器人和服务；纯本地 getter 则不一定需要全部条件。

\textbf{参数}

无显式参数。实例方法中的 \passthrough{\lstinline!self!} 由 Python
自动传入。

\textbf{返回值}

\begin{longtable}[]{@{}
  >{\raggedright\arraybackslash}p{(\columnwidth - 4\tabcolsep) * \real{0.18}}
  >{\raggedright\arraybackslash}p{(\columnwidth - 4\tabcolsep) * \real{0.20}}
  >{\raggedright\arraybackslash}p{(\columnwidth - 4\tabcolsep) * \real{0.62}}@{}}
\toprule
位置 & 类型 & 含义 \\
\midrule
\endhead
{[}0{]} \passthrough{\lstinline!status!} & \passthrough{\lstinline!int!}
& SDK 状态码；Unitree 示例通常以 \passthrough{\lstinline!0!}
表示成功，非零值需按具体服务错误码解释。 \\
{[}1{]} \passthrough{\lstinline!balance\_mode!} &
\passthrough{\lstinline!int!} & 传给该接口的 平衡 模式 参数，Python
类型为
\passthrough{\lstinline!int!}。精确含义、单位、范围和枚举值需查目标型号对应头文件与协议。
对应 C++ 参数 \passthrough{\lstinline!balance\_mode: int \&!}。 \\
\bottomrule
\end{longtable}

\textbf{对应 C++}

\begin{itemize}
\tightlist
\item
  类：\passthrough{\lstinline!unitree::robot::g1::LocoClient!}
\item
  签名：\passthrough{\lstinline!GetBalanceMode(int \&)!}
\item
  绑定策略：\passthrough{\lstinline!OUTPUT\_WRAPPER!}
\item
  声明位置：\passthrough{\lstinline!include/unitree/robot/g1/loco/g1\_loco\_client.hpp:68!}
\end{itemize}

\textbf{说明}

\begin{itemize}
\tightlist
\item
  示例是局部调用片段；\passthrough{\lstinline!obj!}
  和参数变量表示已经按上层业务校验并准备好的对象。
\item
  C++ 的可修改输出引用不会作为 Python 入参出现，而是按 C++
  参数顺序追加到返回元组中。
\end{itemize}

\textbf{用法}

\begin{lstlisting}[language=Python]
status, balance_mode = obj.get_balance_mode()
\end{lstlisting}

\begin{quote}
{[}!NOTE{]} 只读查询仍会访问
DDS/机器人服务。必须检查返回状态码，并处理超时和网络中断。
\end{quote}

\hypertarget{unitree-sdk2-cpp-robot-g1-lococlient-get-swing-height-1}{%
\paragraph{\texorpdfstring{\texttt{unitree\_sdk2\_cpp.robot.g1.LocoClient.get\_swing\_height}}{unitree\_sdk2\_cpp.robot.g1.LocoClient.get\_swing\_height}}\label{unitree-sdk2-cpp-robot-g1-lococlient-get-swing-height-1}}

查询或检查 \passthrough{\lstinline!swing!} 高度。

\textbf{签名}

\begin{lstlisting}[language=Python]
def get_swing_height(self) -> tuple[int, float]
\end{lstlisting}

\textbf{可用性与安全性}

\textbf{\passthrough{\lstinline!AVAILABLE!} /
\passthrough{\lstinline!READ\_ONLY!}}：当前绑定源码已实现只读查询。Robot
Client 的服务端查询通常仍需匹配的 Linux
扩展、DDS、网络、机器人和服务；纯本地 getter 则不一定需要全部条件。

\textbf{参数}

无显式参数。实例方法中的 \passthrough{\lstinline!self!} 由 Python
自动传入。

\textbf{返回值}

\begin{longtable}[]{@{}
  >{\raggedright\arraybackslash}p{(\columnwidth - 4\tabcolsep) * \real{0.18}}
  >{\raggedright\arraybackslash}p{(\columnwidth - 4\tabcolsep) * \real{0.20}}
  >{\raggedright\arraybackslash}p{(\columnwidth - 4\tabcolsep) * \real{0.62}}@{}}
\toprule
位置 & 类型 & 含义 \\
\midrule
\endhead
{[}0{]} \passthrough{\lstinline!status!} & \passthrough{\lstinline!int!}
& SDK 状态码；Unitree 示例通常以 \passthrough{\lstinline!0!}
表示成功，非零值需按具体服务错误码解释。 \\
{[}1{]} \passthrough{\lstinline!swing\_height!} &
\passthrough{\lstinline!float!} & 传给该接口的
\passthrough{\lstinline!swing!} 高度 参数，Python 类型为
\passthrough{\lstinline!float!}。精确含义、单位、范围和枚举值需查目标型号对应头文件与协议。
对应 C++ 参数 \passthrough{\lstinline!swing\_height: float \&!}。
底层为浮点数；类型本身不说明单位、坐标系或业务安全范围。 \\
\bottomrule
\end{longtable}

\textbf{对应 C++}

\begin{itemize}
\tightlist
\item
  类：\passthrough{\lstinline!unitree::robot::g1::LocoClient!}
\item
  签名：\passthrough{\lstinline!GetSwingHeight(float \&)!}
\item
  绑定策略：\passthrough{\lstinline!OUTPUT\_WRAPPER!}
\item
  声明位置：\passthrough{\lstinline!include/unitree/robot/g1/loco/g1\_loco\_client.hpp:82!}
\end{itemize}

\textbf{说明}

\begin{itemize}
\tightlist
\item
  示例是局部调用片段；\passthrough{\lstinline!obj!}
  和参数变量表示已经按上层业务校验并准备好的对象。
\item
  C++ 的可修改输出引用不会作为 Python 入参出现，而是按 C++
  参数顺序追加到返回元组中。
\end{itemize}

\textbf{用法}

\begin{lstlisting}[language=Python]
status, swing_height = obj.get_swing_height()
\end{lstlisting}

\begin{quote}
{[}!NOTE{]} 只读查询仍会访问
DDS/机器人服务。必须检查返回状态码，并处理超时和网络中断。
\end{quote}

\hypertarget{unitree-sdk2-cpp-robot-g1-lococlient-get-stand-height-1}{%
\paragraph{\texorpdfstring{\texttt{unitree\_sdk2\_cpp.robot.g1.LocoClient.get\_stand\_height}}{unitree\_sdk2\_cpp.robot.g1.LocoClient.get\_stand\_height}}\label{unitree-sdk2-cpp-robot-g1-lococlient-get-stand-height-1}}

查询或检查 \passthrough{\lstinline!stand!} 高度。

\textbf{签名}

\begin{lstlisting}[language=Python]
def get_stand_height(self) -> tuple[int, float]
\end{lstlisting}

\textbf{可用性与安全性}

\textbf{\passthrough{\lstinline!AVAILABLE!} /
\passthrough{\lstinline!READ\_ONLY!}}：当前绑定源码已实现只读查询。Robot
Client 的服务端查询通常仍需匹配的 Linux
扩展、DDS、网络、机器人和服务；纯本地 getter 则不一定需要全部条件。

\textbf{参数}

无显式参数。实例方法中的 \passthrough{\lstinline!self!} 由 Python
自动传入。

\textbf{返回值}

\begin{longtable}[]{@{}
  >{\raggedright\arraybackslash}p{(\columnwidth - 4\tabcolsep) * \real{0.18}}
  >{\raggedright\arraybackslash}p{(\columnwidth - 4\tabcolsep) * \real{0.20}}
  >{\raggedright\arraybackslash}p{(\columnwidth - 4\tabcolsep) * \real{0.62}}@{}}
\toprule
位置 & 类型 & 含义 \\
\midrule
\endhead
{[}0{]} \passthrough{\lstinline!status!} & \passthrough{\lstinline!int!}
& SDK 状态码；Unitree 示例通常以 \passthrough{\lstinline!0!}
表示成功，非零值需按具体服务错误码解释。 \\
{[}1{]} \passthrough{\lstinline!stand\_height!} &
\passthrough{\lstinline!float!} & 传给该接口的
\passthrough{\lstinline!stand!} 高度 参数，Python 类型为
\passthrough{\lstinline!float!}。精确含义、单位、范围和枚举值需查目标型号对应头文件与协议。
对应 C++ 参数 \passthrough{\lstinline!stand\_height: float \&!}。
底层为浮点数；类型本身不说明单位、坐标系或业务安全范围。 \\
\bottomrule
\end{longtable}

\textbf{对应 C++}

\begin{itemize}
\tightlist
\item
  类：\passthrough{\lstinline!unitree::robot::g1::LocoClient!}
\item
  签名：\passthrough{\lstinline!GetStandHeight(float \&)!}
\item
  绑定策略：\passthrough{\lstinline!OUTPUT\_WRAPPER!}
\item
  声明位置：\passthrough{\lstinline!include/unitree/robot/g1/loco/g1\_loco\_client.hpp:96!}
\end{itemize}

\textbf{说明}

\begin{itemize}
\tightlist
\item
  示例是局部调用片段；\passthrough{\lstinline!obj!}
  和参数变量表示已经按上层业务校验并准备好的对象。
\item
  C++ 的可修改输出引用不会作为 Python 入参出现，而是按 C++
  参数顺序追加到返回元组中。
\end{itemize}

\textbf{用法}

\begin{lstlisting}[language=Python]
status, stand_height = obj.get_stand_height()
\end{lstlisting}

\begin{quote}
{[}!NOTE{]} 只读查询仍会访问
DDS/机器人服务。必须检查返回状态码，并处理超时和网络中断。
\end{quote}

\hypertarget{unitree-sdk2-cpp-robot-g1-lococlient-get-phase-1}{%
\paragraph{\texorpdfstring{\texttt{unitree\_sdk2\_cpp.robot.g1.LocoClient.get\_phase}}{unitree\_sdk2\_cpp.robot.g1.LocoClient.get\_phase}}\label{unitree-sdk2-cpp-robot-g1-lococlient-get-phase-1}}

查询或检查 相位。

\textbf{签名}

\begin{lstlisting}[language=Python]
def get_phase(self) -> tuple[int, list[float]]
\end{lstlisting}

\textbf{可用性与安全性}

\textbf{\passthrough{\lstinline!AVAILABLE!} /
\passthrough{\lstinline!READ\_ONLY!}}：当前绑定源码已实现只读查询。Robot
Client 的服务端查询通常仍需匹配的 Linux
扩展、DDS、网络、机器人和服务；纯本地 getter 则不一定需要全部条件。

\textbf{参数}

无显式参数。实例方法中的 \passthrough{\lstinline!self!} 由 Python
自动传入。

\textbf{返回值}

\begin{longtable}[]{@{}
  >{\raggedright\arraybackslash}p{(\columnwidth - 4\tabcolsep) * \real{0.18}}
  >{\raggedright\arraybackslash}p{(\columnwidth - 4\tabcolsep) * \real{0.20}}
  >{\raggedright\arraybackslash}p{(\columnwidth - 4\tabcolsep) * \real{0.62}}@{}}
\toprule
位置 & 类型 & 含义 \\
\midrule
\endhead
{[}0{]} \passthrough{\lstinline!status!} & \passthrough{\lstinline!int!}
& SDK 状态码；Unitree 示例通常以 \passthrough{\lstinline!0!}
表示成功，非零值需按具体服务错误码解释。 \\
{[}1{]} \passthrough{\lstinline!phase!} &
\passthrough{\lstinline!list[float]!} & 传给该接口的 相位 参数，Python
类型为
\passthrough{\lstinline!list[float]!}。精确含义、单位、范围和枚举值需查目标型号对应头文件与协议。
对应 C++ 参数 \passthrough{\lstinline!phase: std::vector<float> \&!}。
底层是可变长度
vector；元素约束：底层为浮点数；类型本身不说明单位、坐标系或业务安全范围。 \\
\bottomrule
\end{longtable}

\textbf{对应 C++}

\begin{itemize}
\tightlist
\item
  类：\passthrough{\lstinline!unitree::robot::g1::LocoClient!}
\item
  签名：\passthrough{\lstinline!GetPhase(std::vector<float> \&)!}
\item
  绑定策略：\passthrough{\lstinline!OUTPUT\_WRAPPER!}
\item
  声明位置：\passthrough{\lstinline!include/unitree/robot/g1/loco/g1\_loco\_client.hpp:110!}
\end{itemize}

\textbf{说明}

\begin{itemize}
\tightlist
\item
  示例是局部调用片段；\passthrough{\lstinline!obj!}
  和参数变量表示已经按上层业务校验并准备好的对象。
\item
  C++ 的可修改输出引用不会作为 Python 入参出现，而是按 C++
  参数顺序追加到返回元组中。
\end{itemize}

\textbf{用法}

\begin{lstlisting}[language=Python]
status, phase = obj.get_phase()
\end{lstlisting}

\begin{quote}
{[}!NOTE{]} 只读查询仍会访问
DDS/机器人服务。必须检查返回状态码，并处理超时和网络中断。
\end{quote}

\hypertarget{unitree-sdk2-cpp-robot-g1-lococlient-set-fsm-id-1}{%
\paragraph{\texorpdfstring{\texttt{unitree\_sdk2\_cpp.robot.g1.LocoClient.set\_fsm\_id}}{unitree\_sdk2\_cpp.robot.g1.LocoClient.set\_fsm\_id}}\label{unitree-sdk2-cpp-robot-g1-lococlient-set-fsm-id-1}}

设置 FSM ID。具体副作用和安全边界见下方状态。

\textbf{签名}

\begin{lstlisting}[language=Python]
def set_fsm_id(self, fsm_id: int) -> int
\end{lstlisting}

\textbf{可用性与安全性}

\textbf{\passthrough{\lstinline!SIGNATURE\_ONLY!} /
\passthrough{\lstinline!MOTION\_COMMAND!}}：仅用于补全和静态检查。当前二进制没有该
Python 方法；不得按可执行运动接口使用。

\textbf{参数}

\begin{longtable}[]{@{}
  >{\raggedright\arraybackslash}p{(\columnwidth - 6\tabcolsep) * \real{0.18}}
  >{\raggedright\arraybackslash}p{(\columnwidth - 6\tabcolsep) * \real{0.24}}
  >{\raggedright\arraybackslash}p{(\columnwidth - 6\tabcolsep) * \real{0.18}}
  >{\raggedright\arraybackslash}p{(\columnwidth - 6\tabcolsep) * \real{0.40}}@{}}
\toprule
名称 & Python 类型 & 默认值 & 含义 \\
\midrule
\endhead
\passthrough{\lstinline!fsm\_id!} & \passthrough{\lstinline!int!} & 必填
& 传给该接口的 FSM ID 参数，Python 类型为
\passthrough{\lstinline!int!}。精确含义、单位、范围和枚举值需查目标型号对应头文件与协议。
对应 C++ 参数 \passthrough{\lstinline!fsm\_id: int!}。 \\
\bottomrule
\end{longtable}

\textbf{返回值}

\begin{longtable}[]{@{}
  >{\raggedright\arraybackslash}p{(\columnwidth - 4\tabcolsep) * \real{0.18}}
  >{\raggedright\arraybackslash}p{(\columnwidth - 4\tabcolsep) * \real{0.20}}
  >{\raggedright\arraybackslash}p{(\columnwidth - 4\tabcolsep) * \real{0.62}}@{}}
\toprule
位置 & 类型 & 含义 \\
\midrule
\endhead
返回值 & \passthrough{\lstinline!int!} & SDK 状态码；Unitree 示例通常以
\passthrough{\lstinline!0!} 表示成功，非零值需按具体服务定义解释。 \\
\bottomrule
\end{longtable}

\textbf{对应 C++}

\begin{itemize}
\tightlist
\item
  类：\passthrough{\lstinline!unitree::robot::g1::LocoClient!}
\item
  签名：\passthrough{\lstinline!SetFsmId(int)!}
\item
  绑定策略：\passthrough{\lstinline!DIRECT!}
\item
  声明位置：\passthrough{\lstinline!include/unitree/robot/g1/loco/g1\_loco\_client.hpp:124!}
\end{itemize}

\textbf{说明}

\begin{itemize}
\tightlist
\item
  示例是局部调用片段；\passthrough{\lstinline!obj!}
  和参数变量表示已经按上层业务校验并准备好的对象。
\item
  该条目可以被编辑器和类型检查器识别，但当前运行时可能在属性查找阶段直接失败。
\item
  该方法属于运动命令。方法名包含
  \passthrough{\lstinline!stop!}、\passthrough{\lstinline!damp!} 或
  \passthrough{\lstinline!zero!} 也不代表它是独立物理急停。
\end{itemize}

\textbf{用法}

\begin{lstlisting}[language=Python]
from typing import TYPE_CHECKING

if TYPE_CHECKING:
    def planned_set_fsm_id(obj: LocoClient, fsm_id: int) -> int:
        return obj.set_fsm_id(fsm_id=fsm_id)
\end{lstlisting}

\begin{quote}
{[}!CAUTION{]} 上面的 \passthrough{\lstinline!TYPE\_CHECKING!}
示例只表达计划调用形态。该分支在普通 Python
运行时为假，不代表当前扩展能执行此接口。
\end{quote}

\hypertarget{unitree-sdk2-cpp-robot-g1-lococlient-set-balance-mode-1}{%
\paragraph{\texorpdfstring{\texttt{unitree\_sdk2\_cpp.robot.g1.LocoClient.set\_balance\_mode}}{unitree\_sdk2\_cpp.robot.g1.LocoClient.set\_balance\_mode}}\label{unitree-sdk2-cpp-robot-g1-lococlient-set-balance-mode-1}}

设置 平衡 模式。具体副作用和安全边界见下方状态。

\textbf{签名}

\begin{lstlisting}[language=Python]
def set_balance_mode(self, balance_mode: int) -> int
\end{lstlisting}

\textbf{可用性与安全性}

\textbf{\passthrough{\lstinline!SIGNATURE\_ONLY!} /
\passthrough{\lstinline!MOTION\_COMMAND!}}：仅用于补全和静态检查。当前二进制没有该
Python 方法；不得按可执行运动接口使用。

\textbf{参数}

\begin{longtable}[]{@{}
  >{\raggedright\arraybackslash}p{(\columnwidth - 6\tabcolsep) * \real{0.18}}
  >{\raggedright\arraybackslash}p{(\columnwidth - 6\tabcolsep) * \real{0.24}}
  >{\raggedright\arraybackslash}p{(\columnwidth - 6\tabcolsep) * \real{0.18}}
  >{\raggedright\arraybackslash}p{(\columnwidth - 6\tabcolsep) * \real{0.40}}@{}}
\toprule
名称 & Python 类型 & 默认值 & 含义 \\
\midrule
\endhead
\passthrough{\lstinline!balance\_mode!} & \passthrough{\lstinline!int!}
& 必填 & 传给该接口的 平衡 模式 参数，Python 类型为
\passthrough{\lstinline!int!}。精确含义、单位、范围和枚举值需查目标型号对应头文件与协议。
对应 C++ 参数 \passthrough{\lstinline!balance\_mode: int!}。 \\
\bottomrule
\end{longtable}

\textbf{返回值}

\begin{longtable}[]{@{}
  >{\raggedright\arraybackslash}p{(\columnwidth - 4\tabcolsep) * \real{0.18}}
  >{\raggedright\arraybackslash}p{(\columnwidth - 4\tabcolsep) * \real{0.20}}
  >{\raggedright\arraybackslash}p{(\columnwidth - 4\tabcolsep) * \real{0.62}}@{}}
\toprule
位置 & 类型 & 含义 \\
\midrule
\endhead
返回值 & \passthrough{\lstinline!int!} & SDK 状态码；Unitree 示例通常以
\passthrough{\lstinline!0!} 表示成功，非零值需按具体服务定义解释。 \\
\bottomrule
\end{longtable}

\textbf{对应 C++}

\begin{itemize}
\tightlist
\item
  类：\passthrough{\lstinline!unitree::robot::g1::LocoClient!}
\item
  签名：\passthrough{\lstinline!SetBalanceMode(int)!}
\item
  绑定策略：\passthrough{\lstinline!DIRECT!}
\item
  声明位置：\passthrough{\lstinline!include/unitree/robot/g1/loco/g1\_loco\_client.hpp:134!}
\end{itemize}

\textbf{说明}

\begin{itemize}
\tightlist
\item
  示例是局部调用片段；\passthrough{\lstinline!obj!}
  和参数变量表示已经按上层业务校验并准备好的对象。
\item
  该条目可以被编辑器和类型检查器识别，但当前运行时可能在属性查找阶段直接失败。
\item
  该方法属于运动命令。方法名包含
  \passthrough{\lstinline!stop!}、\passthrough{\lstinline!damp!} 或
  \passthrough{\lstinline!zero!} 也不代表它是独立物理急停。
\end{itemize}

\textbf{用法}

\begin{lstlisting}[language=Python]
from typing import TYPE_CHECKING

if TYPE_CHECKING:
    def planned_set_balance_mode(obj: LocoClient, balance_mode: int) -> int:
        return obj.set_balance_mode(balance_mode=balance_mode)
\end{lstlisting}

\begin{quote}
{[}!CAUTION{]} 上面的 \passthrough{\lstinline!TYPE\_CHECKING!}
示例只表达计划调用形态。该分支在普通 Python
运行时为假，不代表当前扩展能执行此接口。
\end{quote}

\hypertarget{unitree-sdk2-cpp-robot-g1-lococlient-set-swing-height-1}{%
\paragraph{\texorpdfstring{\texttt{unitree\_sdk2\_cpp.robot.g1.LocoClient.set\_swing\_height}}{unitree\_sdk2\_cpp.robot.g1.LocoClient.set\_swing\_height}}\label{unitree-sdk2-cpp-robot-g1-lococlient-set-swing-height-1}}

设置 \passthrough{\lstinline!swing!}
高度。具体副作用和安全边界见下方状态。

\textbf{签名}

\begin{lstlisting}[language=Python]
def set_swing_height(self, swing_height: float) -> int
\end{lstlisting}

\textbf{可用性与安全性}

\textbf{\passthrough{\lstinline!SIGNATURE\_ONLY!} /
\passthrough{\lstinline!MOTION\_COMMAND!}}：仅用于补全和静态检查。当前二进制没有该
Python 方法；不得按可执行运动接口使用。

\textbf{参数}

\begin{longtable}[]{@{}
  >{\raggedright\arraybackslash}p{(\columnwidth - 6\tabcolsep) * \real{0.18}}
  >{\raggedright\arraybackslash}p{(\columnwidth - 6\tabcolsep) * \real{0.24}}
  >{\raggedright\arraybackslash}p{(\columnwidth - 6\tabcolsep) * \real{0.18}}
  >{\raggedright\arraybackslash}p{(\columnwidth - 6\tabcolsep) * \real{0.40}}@{}}
\toprule
名称 & Python 类型 & 默认值 & 含义 \\
\midrule
\endhead
\passthrough{\lstinline!swing\_height!} &
\passthrough{\lstinline!float!} & 必填 & 传给该接口的
\passthrough{\lstinline!swing!} 高度 参数，Python 类型为
\passthrough{\lstinline!float!}。精确含义、单位、范围和枚举值需查目标型号对应头文件与协议。
对应 C++ 参数 \passthrough{\lstinline!swing\_height: float!}。
底层为浮点数；类型本身不说明单位、坐标系或业务安全范围。 \\
\bottomrule
\end{longtable}

\textbf{返回值}

\begin{longtable}[]{@{}
  >{\raggedright\arraybackslash}p{(\columnwidth - 4\tabcolsep) * \real{0.18}}
  >{\raggedright\arraybackslash}p{(\columnwidth - 4\tabcolsep) * \real{0.20}}
  >{\raggedright\arraybackslash}p{(\columnwidth - 4\tabcolsep) * \real{0.62}}@{}}
\toprule
位置 & 类型 & 含义 \\
\midrule
\endhead
返回值 & \passthrough{\lstinline!int!} & SDK 状态码；Unitree 示例通常以
\passthrough{\lstinline!0!} 表示成功，非零值需按具体服务定义解释。 \\
\bottomrule
\end{longtable}

\textbf{对应 C++}

\begin{itemize}
\tightlist
\item
  类：\passthrough{\lstinline!unitree::robot::g1::LocoClient!}
\item
  签名：\passthrough{\lstinline!SetSwingHeight(float)!}
\item
  绑定策略：\passthrough{\lstinline!DIRECT!}
\item
  声明位置：\passthrough{\lstinline!include/unitree/robot/g1/loco/g1\_loco\_client.hpp:144!}
\end{itemize}

\textbf{说明}

\begin{itemize}
\tightlist
\item
  示例是局部调用片段；\passthrough{\lstinline!obj!}
  和参数变量表示已经按上层业务校验并准备好的对象。
\item
  该条目可以被编辑器和类型检查器识别，但当前运行时可能在属性查找阶段直接失败。
\item
  该方法属于运动命令。方法名包含
  \passthrough{\lstinline!stop!}、\passthrough{\lstinline!damp!} 或
  \passthrough{\lstinline!zero!} 也不代表它是独立物理急停。
\end{itemize}

\textbf{用法}

\begin{lstlisting}[language=Python]
from typing import TYPE_CHECKING

if TYPE_CHECKING:
    def planned_set_swing_height(obj: LocoClient, swing_height: float) -> int:
        return obj.set_swing_height(swing_height=swing_height)
\end{lstlisting}

\begin{quote}
{[}!CAUTION{]} 上面的 \passthrough{\lstinline!TYPE\_CHECKING!}
示例只表达计划调用形态。该分支在普通 Python
运行时为假，不代表当前扩展能执行此接口。
\end{quote}

\hypertarget{unitree-sdk2-cpp-robot-g1-lococlient-set-stand-height-1}{%
\paragraph{\texorpdfstring{\texttt{unitree\_sdk2\_cpp.robot.g1.LocoClient.set\_stand\_height}}{unitree\_sdk2\_cpp.robot.g1.LocoClient.set\_stand\_height}}\label{unitree-sdk2-cpp-robot-g1-lococlient-set-stand-height-1}}

设置 \passthrough{\lstinline!stand!}
高度。具体副作用和安全边界见下方状态。

\textbf{签名}

\begin{lstlisting}[language=Python]
def set_stand_height(self, stand_height: float) -> int
\end{lstlisting}

\textbf{可用性与安全性}

\textbf{\passthrough{\lstinline!SIGNATURE\_ONLY!} /
\passthrough{\lstinline!MOTION\_COMMAND!}}：仅用于补全和静态检查。当前二进制没有该
Python 方法；不得按可执行运动接口使用。

\textbf{参数}

\begin{longtable}[]{@{}
  >{\raggedright\arraybackslash}p{(\columnwidth - 6\tabcolsep) * \real{0.18}}
  >{\raggedright\arraybackslash}p{(\columnwidth - 6\tabcolsep) * \real{0.24}}
  >{\raggedright\arraybackslash}p{(\columnwidth - 6\tabcolsep) * \real{0.18}}
  >{\raggedright\arraybackslash}p{(\columnwidth - 6\tabcolsep) * \real{0.40}}@{}}
\toprule
名称 & Python 类型 & 默认值 & 含义 \\
\midrule
\endhead
\passthrough{\lstinline!stand\_height!} &
\passthrough{\lstinline!float!} & 必填 & 传给该接口的
\passthrough{\lstinline!stand!} 高度 参数，Python 类型为
\passthrough{\lstinline!float!}。精确含义、单位、范围和枚举值需查目标型号对应头文件与协议。
对应 C++ 参数 \passthrough{\lstinline!stand\_height: float!}。
底层为浮点数；类型本身不说明单位、坐标系或业务安全范围。 \\
\bottomrule
\end{longtable}

\textbf{返回值}

\begin{longtable}[]{@{}
  >{\raggedright\arraybackslash}p{(\columnwidth - 4\tabcolsep) * \real{0.18}}
  >{\raggedright\arraybackslash}p{(\columnwidth - 4\tabcolsep) * \real{0.20}}
  >{\raggedright\arraybackslash}p{(\columnwidth - 4\tabcolsep) * \real{0.62}}@{}}
\toprule
位置 & 类型 & 含义 \\
\midrule
\endhead
返回值 & \passthrough{\lstinline!int!} & SDK 状态码；Unitree 示例通常以
\passthrough{\lstinline!0!} 表示成功，非零值需按具体服务定义解释。 \\
\bottomrule
\end{longtable}

\textbf{对应 C++}

\begin{itemize}
\tightlist
\item
  类：\passthrough{\lstinline!unitree::robot::g1::LocoClient!}
\item
  签名：\passthrough{\lstinline!SetStandHeight(float)!}
\item
  绑定策略：\passthrough{\lstinline!DIRECT!}
\item
  声明位置：\passthrough{\lstinline!include/unitree/robot/g1/loco/g1\_loco\_client.hpp:154!}
\end{itemize}

\textbf{说明}

\begin{itemize}
\tightlist
\item
  示例是局部调用片段；\passthrough{\lstinline!obj!}
  和参数变量表示已经按上层业务校验并准备好的对象。
\item
  该条目可以被编辑器和类型检查器识别，但当前运行时可能在属性查找阶段直接失败。
\item
  该方法属于运动命令。方法名包含
  \passthrough{\lstinline!stop!}、\passthrough{\lstinline!damp!} 或
  \passthrough{\lstinline!zero!} 也不代表它是独立物理急停。
\end{itemize}

\textbf{用法}

\begin{lstlisting}[language=Python]
from typing import TYPE_CHECKING

if TYPE_CHECKING:
    def planned_set_stand_height(obj: LocoClient, stand_height: float) -> int:
        return obj.set_stand_height(stand_height=stand_height)
\end{lstlisting}

\begin{quote}
{[}!CAUTION{]} 上面的 \passthrough{\lstinline!TYPE\_CHECKING!}
示例只表达计划调用形态。该分支在普通 Python
运行时为假，不代表当前扩展能执行此接口。
\end{quote}

\hypertarget{unitree-sdk2-cpp-robot-g1-lococlient-set-velocity-1}{%
\paragraph{\texorpdfstring{\texttt{unitree\_sdk2\_cpp.robot.g1.LocoClient.set\_velocity}}{unitree\_sdk2\_cpp.robot.g1.LocoClient.set\_velocity}}\label{unitree-sdk2-cpp-robot-g1-lococlient-set-velocity-1}}

设置 速度。具体副作用和安全边界见下方状态。

\textbf{签名}

\begin{lstlisting}[language=Python]
def set_velocity(self, vx: float, vy: float, omega: float, duration: float = ...) -> int
\end{lstlisting}

\textbf{可用性与安全性}

\textbf{\passthrough{\lstinline!SIGNATURE\_ONLY!} /
\passthrough{\lstinline!MOTION\_COMMAND!}}：仅用于补全和静态检查。当前二进制没有该
Python 方法；不得按可执行运动接口使用。

\textbf{参数}

\begin{longtable}[]{@{}
  >{\raggedright\arraybackslash}p{(\columnwidth - 6\tabcolsep) * \real{0.18}}
  >{\raggedright\arraybackslash}p{(\columnwidth - 6\tabcolsep) * \real{0.24}}
  >{\raggedright\arraybackslash}p{(\columnwidth - 6\tabcolsep) * \real{0.18}}
  >{\raggedright\arraybackslash}p{(\columnwidth - 6\tabcolsep) * \real{0.40}}@{}}
\toprule
名称 & Python 类型 & 默认值 & 含义 \\
\midrule
\endhead
\passthrough{\lstinline!vx!} & \passthrough{\lstinline!float!} & 必填 &
X 方向速度参数。坐标系、单位、符号和安全范围必须查目标型号运动协议。
对应 C++ 参数 \passthrough{\lstinline!vx: float!}。
底层为浮点数；类型本身不说明单位、坐标系或业务安全范围。 \\
\passthrough{\lstinline!vy!} & \passthrough{\lstinline!float!} & 必填 &
Y 方向速度参数。坐标系、单位、符号和安全范围必须查目标型号运动协议。
对应 C++ 参数 \passthrough{\lstinline!vy: float!}。
底层为浮点数；类型本身不说明单位、坐标系或业务安全范围。 \\
\passthrough{\lstinline!omega!} & \passthrough{\lstinline!float!} & 必填
& 角速度参数。旋转轴、单位、符号和安全范围必须查目标型号运动协议。 对应
C++ 参数 \passthrough{\lstinline!omega: float!}。
底层为浮点数；类型本身不说明单位、坐标系或业务安全范围。 \\
\passthrough{\lstinline!duration!} & \passthrough{\lstinline!float!} &
\passthrough{\lstinline!...!} &
操作持续时间参数。精确单位、范围和默认行为以目标型号接口定义为准。 对应
C++ 参数 \passthrough{\lstinline!duration: float!}。
底层为浮点数；类型本身不说明单位、坐标系或业务安全范围。 \\
\bottomrule
\end{longtable}

\textbf{返回值}

\begin{longtable}[]{@{}
  >{\raggedright\arraybackslash}p{(\columnwidth - 4\tabcolsep) * \real{0.18}}
  >{\raggedright\arraybackslash}p{(\columnwidth - 4\tabcolsep) * \real{0.20}}
  >{\raggedright\arraybackslash}p{(\columnwidth - 4\tabcolsep) * \real{0.62}}@{}}
\toprule
位置 & 类型 & 含义 \\
\midrule
\endhead
返回值 & \passthrough{\lstinline!int!} & SDK 状态码；Unitree 示例通常以
\passthrough{\lstinline!0!} 表示成功，非零值需按具体服务定义解释。 \\
\bottomrule
\end{longtable}

\textbf{对应 C++}

\begin{itemize}
\tightlist
\item
  类：\passthrough{\lstinline!unitree::robot::g1::LocoClient!}
\item
  签名：\passthrough{\lstinline!SetVelocity(float, float, float, float)!}
\item
  绑定策略：\passthrough{\lstinline!DIRECT!}
\item
  声明位置：\passthrough{\lstinline!include/unitree/robot/g1/loco/g1\_loco\_client.hpp:164!}
\end{itemize}

\textbf{说明}

\begin{itemize}
\tightlist
\item
  示例是局部调用片段；\passthrough{\lstinline!obj!}
  和参数变量表示已经按上层业务校验并准备好的对象。
\item
  默认值显示为 \passthrough{\lstinline!...!}，表示 C++
  声明存在默认参数，但当前 AST
  清单没有保存其字面量；省略参数可使用上游默认行为。
\item
  该条目可以被编辑器和类型检查器识别，但当前运行时可能在属性查找阶段直接失败。
\item
  该方法属于运动命令。方法名包含
  \passthrough{\lstinline!stop!}、\passthrough{\lstinline!damp!} 或
  \passthrough{\lstinline!zero!} 也不代表它是独立物理急停。
\end{itemize}

\textbf{用法}

\begin{lstlisting}[language=Python]
from typing import TYPE_CHECKING

if TYPE_CHECKING:
    def planned_set_velocity(obj: LocoClient, vx: float, vy: float, omega: float, duration: float) -> int:
        return obj.set_velocity(vx=vx, vy=vy, omega=omega, duration=duration)
\end{lstlisting}

\begin{quote}
{[}!CAUTION{]} 上面的 \passthrough{\lstinline!TYPE\_CHECKING!}
示例只表达计划调用形态。该分支在普通 Python
运行时为假，不代表当前扩展能执行此接口。
\end{quote}

\hypertarget{unitree-sdk2-cpp-robot-g1-lococlient-set-task-id-1}{%
\paragraph{\texorpdfstring{\texttt{unitree\_sdk2\_cpp.robot.g1.LocoClient.set\_task\_id}}{unitree\_sdk2\_cpp.robot.g1.LocoClient.set\_task\_id}}\label{unitree-sdk2-cpp-robot-g1-lococlient-set-task-id-1}}

设置 任务 ID。具体副作用和安全边界见下方状态。

\textbf{签名}

\begin{lstlisting}[language=Python]
def set_task_id(self, task_id: int) -> int
\end{lstlisting}

\textbf{可用性与安全性}

\textbf{\passthrough{\lstinline!SIGNATURE\_ONLY!} /
\passthrough{\lstinline!MOTION\_COMMAND!}}：仅用于补全和静态检查。当前二进制没有该
Python 方法；不得按可执行运动接口使用。

\textbf{参数}

\begin{longtable}[]{@{}
  >{\raggedright\arraybackslash}p{(\columnwidth - 6\tabcolsep) * \real{0.18}}
  >{\raggedright\arraybackslash}p{(\columnwidth - 6\tabcolsep) * \real{0.24}}
  >{\raggedright\arraybackslash}p{(\columnwidth - 6\tabcolsep) * \real{0.18}}
  >{\raggedright\arraybackslash}p{(\columnwidth - 6\tabcolsep) * \real{0.40}}@{}}
\toprule
名称 & Python 类型 & 默认值 & 含义 \\
\midrule
\endhead
\passthrough{\lstinline!task\_id!} & \passthrough{\lstinline!int!} &
必填 & 传给该接口的 任务 ID 参数，Python 类型为
\passthrough{\lstinline!int!}。精确含义、单位、范围和枚举值需查目标型号对应头文件与协议。
对应 C++ 参数 \passthrough{\lstinline!task\_id: int!}。 \\
\bottomrule
\end{longtable}

\textbf{返回值}

\begin{longtable}[]{@{}
  >{\raggedright\arraybackslash}p{(\columnwidth - 4\tabcolsep) * \real{0.18}}
  >{\raggedright\arraybackslash}p{(\columnwidth - 4\tabcolsep) * \real{0.20}}
  >{\raggedright\arraybackslash}p{(\columnwidth - 4\tabcolsep) * \real{0.62}}@{}}
\toprule
位置 & 类型 & 含义 \\
\midrule
\endhead
返回值 & \passthrough{\lstinline!int!} & SDK 状态码；Unitree 示例通常以
\passthrough{\lstinline!0!} 表示成功，非零值需按具体服务定义解释。 \\
\bottomrule
\end{longtable}

\textbf{对应 C++}

\begin{itemize}
\tightlist
\item
  类：\passthrough{\lstinline!unitree::robot::g1::LocoClient!}
\item
  签名：\passthrough{\lstinline!SetTaskId(int)!}
\item
  绑定策略：\passthrough{\lstinline!DIRECT!}
\item
  声明位置：\passthrough{\lstinline!include/unitree/robot/g1/loco/g1\_loco\_client.hpp:176!}
\end{itemize}

\textbf{说明}

\begin{itemize}
\tightlist
\item
  示例是局部调用片段；\passthrough{\lstinline!obj!}
  和参数变量表示已经按上层业务校验并准备好的对象。
\item
  该条目可以被编辑器和类型检查器识别，但当前运行时可能在属性查找阶段直接失败。
\item
  该方法属于运动命令。方法名包含
  \passthrough{\lstinline!stop!}、\passthrough{\lstinline!damp!} 或
  \passthrough{\lstinline!zero!} 也不代表它是独立物理急停。
\end{itemize}

\textbf{用法}

\begin{lstlisting}[language=Python]
from typing import TYPE_CHECKING

if TYPE_CHECKING:
    def planned_set_task_id(obj: LocoClient, task_id: int) -> int:
        return obj.set_task_id(task_id=task_id)
\end{lstlisting}

\begin{quote}
{[}!CAUTION{]} 上面的 \passthrough{\lstinline!TYPE\_CHECKING!}
示例只表达计划调用形态。该分支在普通 Python
运行时为假，不代表当前扩展能执行此接口。
\end{quote}

\hypertarget{unitree-sdk2-cpp-robot-g1-lococlient-switch-to-user-ctrl-1}{%
\paragraph{\texorpdfstring{\texttt{unitree\_sdk2\_cpp.robot.g1.LocoClient.switch\_to\_user\_ctrl}}{unitree\_sdk2\_cpp.robot.g1.LocoClient.switch\_to\_user\_ctrl}}\label{unitree-sdk2-cpp-robot-g1-lococlient-switch-to-user-ctrl-1}}

对应 C++ SDK 操作
\passthrough{\lstinline!SwitchToUserCtrl()!}。上游头文件没有可直接生成的业务说明时，本参考只保证签名映射，精确语义需查目标型号协议。

\textbf{签名}

\begin{lstlisting}[language=Python]
def switch_to_user_ctrl(self) -> int
\end{lstlisting}

\textbf{可用性与安全性}

\textbf{\passthrough{\lstinline!SIGNATURE\_ONLY!} /
\passthrough{\lstinline!MOTION\_COMMAND!}}：仅用于补全和静态检查。当前二进制没有该
Python 方法；不得按可执行运动接口使用。

\textbf{参数}

无显式参数。实例方法中的 \passthrough{\lstinline!self!} 由 Python
自动传入。

\textbf{返回值}

\begin{longtable}[]{@{}
  >{\raggedright\arraybackslash}p{(\columnwidth - 4\tabcolsep) * \real{0.18}}
  >{\raggedright\arraybackslash}p{(\columnwidth - 4\tabcolsep) * \real{0.20}}
  >{\raggedright\arraybackslash}p{(\columnwidth - 4\tabcolsep) * \real{0.62}}@{}}
\toprule
位置 & 类型 & 含义 \\
\midrule
\endhead
返回值 & \passthrough{\lstinline!int!} & SDK 状态码；Unitree 示例通常以
\passthrough{\lstinline!0!} 表示成功，非零值需按具体服务定义解释。 \\
\bottomrule
\end{longtable}

\textbf{对应 C++}

\begin{itemize}
\tightlist
\item
  类：\passthrough{\lstinline!unitree::robot::g1::LocoClient!}
\item
  签名：\passthrough{\lstinline!SwitchToUserCtrl()!}
\item
  绑定策略：\passthrough{\lstinline!DIRECT!}
\item
  声明位置：\passthrough{\lstinline!include/unitree/robot/g1/loco/g1\_loco\_client.hpp:186!}
\end{itemize}

\textbf{说明}

\begin{itemize}
\tightlist
\item
  示例是局部调用片段；\passthrough{\lstinline!obj!}
  和参数变量表示已经按上层业务校验并准备好的对象。
\item
  该条目可以被编辑器和类型检查器识别，但当前运行时可能在属性查找阶段直接失败。
\item
  该方法属于运动命令。方法名包含
  \passthrough{\lstinline!stop!}、\passthrough{\lstinline!damp!} 或
  \passthrough{\lstinline!zero!} 也不代表它是独立物理急停。
\end{itemize}

\textbf{用法}

\begin{lstlisting}[language=Python]
from typing import TYPE_CHECKING

if TYPE_CHECKING:
    def planned_switch_to_user_ctrl(obj: LocoClient) -> int:
        return obj.switch_to_user_ctrl()
\end{lstlisting}

\begin{quote}
{[}!CAUTION{]} 上面的 \passthrough{\lstinline!TYPE\_CHECKING!}
示例只表达计划调用形态。该分支在普通 Python
运行时为假，不代表当前扩展能执行此接口。
\end{quote}

\hypertarget{unitree-sdk2-cpp-robot-g1-lococlient-switch-to-internal-ctrl-1}{%
\paragraph{\texorpdfstring{\texttt{unitree\_sdk2\_cpp.robot.g1.LocoClient.switch\_to\_internal\_ctrl}}{unitree\_sdk2\_cpp.robot.g1.LocoClient.switch\_to\_internal\_ctrl}}\label{unitree-sdk2-cpp-robot-g1-lococlient-switch-to-internal-ctrl-1}}

对应 C++ SDK 操作
\passthrough{\lstinline!SwitchToInternalCtrl(unitree::robot::g1::InternalFsmMode)!}。上游头文件没有可直接生成的业务说明时，本参考只保证签名映射，精确语义需查目标型号协议。

\textbf{签名}

\begin{lstlisting}[language=Python]
def switch_to_internal_ctrl(self, mode: InternalFsmMode) -> int
\end{lstlisting}

\textbf{可用性与安全性}

\textbf{\passthrough{\lstinline!SIGNATURE\_ONLY!} /
\passthrough{\lstinline!MOTION\_COMMAND!}}：仅用于补全和静态检查。当前二进制没有该
Python 方法；不得按可执行运动接口使用。

\textbf{参数}

\begin{longtable}[]{@{}
  >{\raggedright\arraybackslash}p{(\columnwidth - 6\tabcolsep) * \real{0.18}}
  >{\raggedright\arraybackslash}p{(\columnwidth - 6\tabcolsep) * \real{0.24}}
  >{\raggedright\arraybackslash}p{(\columnwidth - 6\tabcolsep) * \real{0.18}}
  >{\raggedright\arraybackslash}p{(\columnwidth - 6\tabcolsep) * \real{0.40}}@{}}
\toprule
名称 & Python 类型 & 默认值 & 含义 \\
\midrule
\endhead
\passthrough{\lstinline!mode!} &
\passthrough{\lstinline!InternalFsmMode!} & 必填 & 传给该接口的 模式
参数，Python 类型为
\passthrough{\lstinline!InternalFsmMode!}。精确含义、单位、范围和枚举值需查目标型号对应头文件与协议。
对应 C++ 参数
\passthrough{\lstinline!mode: unitree::robot::g1::InternalFsmMode!}。 \\
\bottomrule
\end{longtable}

\textbf{返回值}

\begin{longtable}[]{@{}
  >{\raggedright\arraybackslash}p{(\columnwidth - 4\tabcolsep) * \real{0.18}}
  >{\raggedright\arraybackslash}p{(\columnwidth - 4\tabcolsep) * \real{0.20}}
  >{\raggedright\arraybackslash}p{(\columnwidth - 4\tabcolsep) * \real{0.62}}@{}}
\toprule
位置 & 类型 & 含义 \\
\midrule
\endhead
返回值 & \passthrough{\lstinline!int!} & SDK 状态码；Unitree 示例通常以
\passthrough{\lstinline!0!} 表示成功，非零值需按具体服务定义解释。 \\
\bottomrule
\end{longtable}

\textbf{对应 C++}

\begin{itemize}
\tightlist
\item
  类：\passthrough{\lstinline!unitree::robot::g1::LocoClient!}
\item
  签名：\passthrough{\lstinline!SwitchToInternalCtrl(unitree::robot::g1::InternalFsmMode)!}
\item
  绑定策略：\passthrough{\lstinline!DIRECT!}
\item
  声明位置：\passthrough{\lstinline!include/unitree/robot/g1/loco/g1\_loco\_client.hpp:194!}
\end{itemize}

\textbf{说明}

\begin{itemize}
\tightlist
\item
  示例是局部调用片段；\passthrough{\lstinline!obj!}
  和参数变量表示已经按上层业务校验并准备好的对象。
\item
  该条目可以被编辑器和类型检查器识别，但当前运行时可能在属性查找阶段直接失败。
\item
  该方法属于运动命令。方法名包含
  \passthrough{\lstinline!stop!}、\passthrough{\lstinline!damp!} 或
  \passthrough{\lstinline!zero!} 也不代表它是独立物理急停。
\end{itemize}

\textbf{用法}

\begin{lstlisting}[language=Python]
from typing import TYPE_CHECKING

if TYPE_CHECKING:
    def planned_switch_to_internal_ctrl(obj: LocoClient, mode: InternalFsmMode) -> int:
        return obj.switch_to_internal_ctrl(mode=mode)
\end{lstlisting}

\begin{quote}
{[}!CAUTION{]} 上面的 \passthrough{\lstinline!TYPE\_CHECKING!}
示例只表达计划调用形态。该分支在普通 Python
运行时为假，不代表当前扩展能执行此接口。
\end{quote}

\hypertarget{unitree-sdk2-cpp-robot-g1-lococlient-damp-1}{%
\paragraph{\texorpdfstring{\texttt{unitree\_sdk2\_cpp.robot.g1.LocoClient.damp}}{unitree\_sdk2\_cpp.robot.g1.LocoClient.damp}}\label{unitree-sdk2-cpp-robot-g1-lococlient-damp-1}}

对应 C++ SDK 操作
\passthrough{\lstinline!Damp()!}。上游头文件没有可直接生成的业务说明时，本参考只保证签名映射，精确语义需查目标型号协议。

\textbf{签名}

\begin{lstlisting}[language=Python]
def damp(self) -> int
\end{lstlisting}

\textbf{可用性与安全性}

\textbf{\passthrough{\lstinline!SIGNATURE\_ONLY!} /
\passthrough{\lstinline!MOTION\_COMMAND!}}：仅用于补全和静态检查。当前二进制没有该
Python 方法；不得按可执行运动接口使用。

\textbf{参数}

无显式参数。实例方法中的 \passthrough{\lstinline!self!} 由 Python
自动传入。

\textbf{返回值}

\begin{longtable}[]{@{}
  >{\raggedright\arraybackslash}p{(\columnwidth - 4\tabcolsep) * \real{0.18}}
  >{\raggedright\arraybackslash}p{(\columnwidth - 4\tabcolsep) * \real{0.20}}
  >{\raggedright\arraybackslash}p{(\columnwidth - 4\tabcolsep) * \real{0.62}}@{}}
\toprule
位置 & 类型 & 含义 \\
\midrule
\endhead
返回值 & \passthrough{\lstinline!int!} & SDK 状态码；Unitree 示例通常以
\passthrough{\lstinline!0!} 表示成功，非零值需按具体服务定义解释。 \\
\bottomrule
\end{longtable}

\textbf{对应 C++}

\begin{itemize}
\tightlist
\item
  类：\passthrough{\lstinline!unitree::robot::g1::LocoClient!}
\item
  签名：\passthrough{\lstinline!Damp()!}
\item
  绑定策略：\passthrough{\lstinline!DIRECT!}
\item
  声明位置：\passthrough{\lstinline!include/unitree/robot/g1/loco/g1\_loco\_client.hpp:214!}
\end{itemize}

\textbf{说明}

\begin{itemize}
\tightlist
\item
  示例是局部调用片段；\passthrough{\lstinline!obj!}
  和参数变量表示已经按上层业务校验并准备好的对象。
\item
  该条目可以被编辑器和类型检查器识别，但当前运行时可能在属性查找阶段直接失败。
\item
  该方法属于运动命令。方法名包含
  \passthrough{\lstinline!stop!}、\passthrough{\lstinline!damp!} 或
  \passthrough{\lstinline!zero!} 也不代表它是独立物理急停。
\end{itemize}

\textbf{用法}

\begin{lstlisting}[language=Python]
from typing import TYPE_CHECKING

if TYPE_CHECKING:
    def planned_damp(obj: LocoClient) -> int:
        return obj.damp()
\end{lstlisting}

\begin{quote}
{[}!CAUTION{]} 上面的 \passthrough{\lstinline!TYPE\_CHECKING!}
示例只表达计划调用形态。该分支在普通 Python
运行时为假，不代表当前扩展能执行此接口。
\end{quote}

\hypertarget{unitree-sdk2-cpp-robot-g1-lococlient-start-1}{%
\paragraph{\texorpdfstring{\texttt{unitree\_sdk2\_cpp.robot.g1.LocoClient.start}}{unitree\_sdk2\_cpp.robot.g1.LocoClient.start}}\label{unitree-sdk2-cpp-robot-g1-lococlient-start-1}}

启动对应 SDK 操作。具体副作用和安全边界见下方状态。

\textbf{签名}

\begin{lstlisting}[language=Python]
def start(self) -> int
\end{lstlisting}

\textbf{可用性与安全性}

\textbf{\passthrough{\lstinline!SIGNATURE\_ONLY!} /
\passthrough{\lstinline!MOTION\_COMMAND!}}：仅用于补全和静态检查。当前二进制没有该
Python 方法；不得按可执行运动接口使用。

\textbf{参数}

无显式参数。实例方法中的 \passthrough{\lstinline!self!} 由 Python
自动传入。

\textbf{返回值}

\begin{longtable}[]{@{}
  >{\raggedright\arraybackslash}p{(\columnwidth - 4\tabcolsep) * \real{0.18}}
  >{\raggedright\arraybackslash}p{(\columnwidth - 4\tabcolsep) * \real{0.20}}
  >{\raggedright\arraybackslash}p{(\columnwidth - 4\tabcolsep) * \real{0.62}}@{}}
\toprule
位置 & 类型 & 含义 \\
\midrule
\endhead
返回值 & \passthrough{\lstinline!int!} & SDK 状态码；Unitree 示例通常以
\passthrough{\lstinline!0!} 表示成功，非零值需按具体服务定义解释。 \\
\bottomrule
\end{longtable}

\textbf{对应 C++}

\begin{itemize}
\tightlist
\item
  类：\passthrough{\lstinline!unitree::robot::g1::LocoClient!}
\item
  签名：\passthrough{\lstinline!Start()!}
\item
  绑定策略：\passthrough{\lstinline!DIRECT!}
\item
  声明位置：\passthrough{\lstinline!include/unitree/robot/g1/loco/g1\_loco\_client.hpp:216!}
\end{itemize}

\textbf{说明}

\begin{itemize}
\tightlist
\item
  示例是局部调用片段；\passthrough{\lstinline!obj!}
  和参数变量表示已经按上层业务校验并准备好的对象。
\item
  该条目可以被编辑器和类型检查器识别，但当前运行时可能在属性查找阶段直接失败。
\item
  该方法属于运动命令。方法名包含
  \passthrough{\lstinline!stop!}、\passthrough{\lstinline!damp!} 或
  \passthrough{\lstinline!zero!} 也不代表它是独立物理急停。
\end{itemize}

\textbf{用法}

\begin{lstlisting}[language=Python]
from typing import TYPE_CHECKING

if TYPE_CHECKING:
    def planned_start(obj: LocoClient) -> int:
        return obj.start()
\end{lstlisting}

\begin{quote}
{[}!CAUTION{]} 上面的 \passthrough{\lstinline!TYPE\_CHECKING!}
示例只表达计划调用形态。该分支在普通 Python
运行时为假，不代表当前扩展能执行此接口。
\end{quote}

\hypertarget{unitree-sdk2-cpp-robot-g1-lococlient-squat-1}{%
\paragraph{\texorpdfstring{\texttt{unitree\_sdk2\_cpp.robot.g1.LocoClient.squat}}{unitree\_sdk2\_cpp.robot.g1.LocoClient.squat}}\label{unitree-sdk2-cpp-robot-g1-lococlient-squat-1}}

对应 C++ SDK 操作
\passthrough{\lstinline!Squat()!}。上游头文件没有可直接生成的业务说明时，本参考只保证签名映射，精确语义需查目标型号协议。

\textbf{签名}

\begin{lstlisting}[language=Python]
def squat(self) -> int
\end{lstlisting}

\textbf{可用性与安全性}

\textbf{\passthrough{\lstinline!SIGNATURE\_ONLY!} /
\passthrough{\lstinline!MOTION\_COMMAND!}}：仅用于补全和静态检查。当前二进制没有该
Python 方法；不得按可执行运动接口使用。

\textbf{参数}

无显式参数。实例方法中的 \passthrough{\lstinline!self!} 由 Python
自动传入。

\textbf{返回值}

\begin{longtable}[]{@{}
  >{\raggedright\arraybackslash}p{(\columnwidth - 4\tabcolsep) * \real{0.18}}
  >{\raggedright\arraybackslash}p{(\columnwidth - 4\tabcolsep) * \real{0.20}}
  >{\raggedright\arraybackslash}p{(\columnwidth - 4\tabcolsep) * \real{0.62}}@{}}
\toprule
位置 & 类型 & 含义 \\
\midrule
\endhead
返回值 & \passthrough{\lstinline!int!} & SDK 状态码；Unitree 示例通常以
\passthrough{\lstinline!0!} 表示成功，非零值需按具体服务定义解释。 \\
\bottomrule
\end{longtable}

\textbf{对应 C++}

\begin{itemize}
\tightlist
\item
  类：\passthrough{\lstinline!unitree::robot::g1::LocoClient!}
\item
  签名：\passthrough{\lstinline!Squat()!}
\item
  绑定策略：\passthrough{\lstinline!DIRECT!}
\item
  声明位置：\passthrough{\lstinline!include/unitree/robot/g1/loco/g1\_loco\_client.hpp:218!}
\end{itemize}

\textbf{说明}

\begin{itemize}
\tightlist
\item
  示例是局部调用片段；\passthrough{\lstinline!obj!}
  和参数变量表示已经按上层业务校验并准备好的对象。
\item
  该条目可以被编辑器和类型检查器识别，但当前运行时可能在属性查找阶段直接失败。
\item
  该方法属于运动命令。方法名包含
  \passthrough{\lstinline!stop!}、\passthrough{\lstinline!damp!} 或
  \passthrough{\lstinline!zero!} 也不代表它是独立物理急停。
\end{itemize}

\textbf{用法}

\begin{lstlisting}[language=Python]
from typing import TYPE_CHECKING

if TYPE_CHECKING:
    def planned_squat(obj: LocoClient) -> int:
        return obj.squat()
\end{lstlisting}

\begin{quote}
{[}!CAUTION{]} 上面的 \passthrough{\lstinline!TYPE\_CHECKING!}
示例只表达计划调用形态。该分支在普通 Python
运行时为假，不代表当前扩展能执行此接口。
\end{quote}

\hypertarget{unitree-sdk2-cpp-robot-g1-lococlient-sit-1}{%
\paragraph{\texorpdfstring{\texttt{unitree\_sdk2\_cpp.robot.g1.LocoClient.sit}}{unitree\_sdk2\_cpp.robot.g1.LocoClient.sit}}\label{unitree-sdk2-cpp-robot-g1-lococlient-sit-1}}

对应 C++ SDK 操作
\passthrough{\lstinline!Sit()!}。上游头文件没有可直接生成的业务说明时，本参考只保证签名映射，精确语义需查目标型号协议。

\textbf{签名}

\begin{lstlisting}[language=Python]
def sit(self) -> int
\end{lstlisting}

\textbf{可用性与安全性}

\textbf{\passthrough{\lstinline!SIGNATURE\_ONLY!} /
\passthrough{\lstinline!MOTION\_COMMAND!}}：仅用于补全和静态检查。当前二进制没有该
Python 方法；不得按可执行运动接口使用。

\textbf{参数}

无显式参数。实例方法中的 \passthrough{\lstinline!self!} 由 Python
自动传入。

\textbf{返回值}

\begin{longtable}[]{@{}
  >{\raggedright\arraybackslash}p{(\columnwidth - 4\tabcolsep) * \real{0.18}}
  >{\raggedright\arraybackslash}p{(\columnwidth - 4\tabcolsep) * \real{0.20}}
  >{\raggedright\arraybackslash}p{(\columnwidth - 4\tabcolsep) * \real{0.62}}@{}}
\toprule
位置 & 类型 & 含义 \\
\midrule
\endhead
返回值 & \passthrough{\lstinline!int!} & SDK 状态码；Unitree 示例通常以
\passthrough{\lstinline!0!} 表示成功，非零值需按具体服务定义解释。 \\
\bottomrule
\end{longtable}

\textbf{对应 C++}

\begin{itemize}
\tightlist
\item
  类：\passthrough{\lstinline!unitree::robot::g1::LocoClient!}
\item
  签名：\passthrough{\lstinline!Sit()!}
\item
  绑定策略：\passthrough{\lstinline!DIRECT!}
\item
  声明位置：\passthrough{\lstinline!include/unitree/robot/g1/loco/g1\_loco\_client.hpp:220!}
\end{itemize}

\textbf{说明}

\begin{itemize}
\tightlist
\item
  示例是局部调用片段；\passthrough{\lstinline!obj!}
  和参数变量表示已经按上层业务校验并准备好的对象。
\item
  该条目可以被编辑器和类型检查器识别，但当前运行时可能在属性查找阶段直接失败。
\item
  该方法属于运动命令。方法名包含
  \passthrough{\lstinline!stop!}、\passthrough{\lstinline!damp!} 或
  \passthrough{\lstinline!zero!} 也不代表它是独立物理急停。
\end{itemize}

\textbf{用法}

\begin{lstlisting}[language=Python]
from typing import TYPE_CHECKING

if TYPE_CHECKING:
    def planned_sit(obj: LocoClient) -> int:
        return obj.sit()
\end{lstlisting}

\begin{quote}
{[}!CAUTION{]} 上面的 \passthrough{\lstinline!TYPE\_CHECKING!}
示例只表达计划调用形态。该分支在普通 Python
运行时为假，不代表当前扩展能执行此接口。
\end{quote}

\hypertarget{unitree-sdk2-cpp-robot-g1-lococlient-stand-up-1}{%
\paragraph{\texorpdfstring{\texttt{unitree\_sdk2\_cpp.robot.g1.LocoClient.stand\_up}}{unitree\_sdk2\_cpp.robot.g1.LocoClient.stand\_up}}\label{unitree-sdk2-cpp-robot-g1-lococlient-stand-up-1}}

对应 C++ SDK 操作
\passthrough{\lstinline!StandUp()!}。上游头文件没有可直接生成的业务说明时，本参考只保证签名映射，精确语义需查目标型号协议。

\textbf{签名}

\begin{lstlisting}[language=Python]
def stand_up(self) -> int
\end{lstlisting}

\textbf{可用性与安全性}

\textbf{\passthrough{\lstinline!SIGNATURE\_ONLY!} /
\passthrough{\lstinline!MOTION\_COMMAND!}}：仅用于补全和静态检查。当前二进制没有该
Python 方法；不得按可执行运动接口使用。

\textbf{参数}

无显式参数。实例方法中的 \passthrough{\lstinline!self!} 由 Python
自动传入。

\textbf{返回值}

\begin{longtable}[]{@{}
  >{\raggedright\arraybackslash}p{(\columnwidth - 4\tabcolsep) * \real{0.18}}
  >{\raggedright\arraybackslash}p{(\columnwidth - 4\tabcolsep) * \real{0.20}}
  >{\raggedright\arraybackslash}p{(\columnwidth - 4\tabcolsep) * \real{0.62}}@{}}
\toprule
位置 & 类型 & 含义 \\
\midrule
\endhead
返回值 & \passthrough{\lstinline!int!} & SDK 状态码；Unitree 示例通常以
\passthrough{\lstinline!0!} 表示成功，非零值需按具体服务定义解释。 \\
\bottomrule
\end{longtable}

\textbf{对应 C++}

\begin{itemize}
\tightlist
\item
  类：\passthrough{\lstinline!unitree::robot::g1::LocoClient!}
\item
  签名：\passthrough{\lstinline!StandUp()!}
\item
  绑定策略：\passthrough{\lstinline!DIRECT!}
\item
  声明位置：\passthrough{\lstinline!include/unitree/robot/g1/loco/g1\_loco\_client.hpp:222!}
\end{itemize}

\textbf{说明}

\begin{itemize}
\tightlist
\item
  示例是局部调用片段；\passthrough{\lstinline!obj!}
  和参数变量表示已经按上层业务校验并准备好的对象。
\item
  该条目可以被编辑器和类型检查器识别，但当前运行时可能在属性查找阶段直接失败。
\item
  该方法属于运动命令。方法名包含
  \passthrough{\lstinline!stop!}、\passthrough{\lstinline!damp!} 或
  \passthrough{\lstinline!zero!} 也不代表它是独立物理急停。
\end{itemize}

\textbf{用法}

\begin{lstlisting}[language=Python]
from typing import TYPE_CHECKING

if TYPE_CHECKING:
    def planned_stand_up(obj: LocoClient) -> int:
        return obj.stand_up()
\end{lstlisting}

\begin{quote}
{[}!CAUTION{]} 上面的 \passthrough{\lstinline!TYPE\_CHECKING!}
示例只表达计划调用形态。该分支在普通 Python
运行时为假，不代表当前扩展能执行此接口。
\end{quote}

\hypertarget{unitree-sdk2-cpp-robot-g1-lococlient-zero-torque-1}{%
\paragraph{\texorpdfstring{\texttt{unitree\_sdk2\_cpp.robot.g1.LocoClient.zero\_torque}}{unitree\_sdk2\_cpp.robot.g1.LocoClient.zero\_torque}}\label{unitree-sdk2-cpp-robot-g1-lococlient-zero-torque-1}}

对应 C++ SDK 操作
\passthrough{\lstinline!ZeroTorque()!}。上游头文件没有可直接生成的业务说明时，本参考只保证签名映射，精确语义需查目标型号协议。

\textbf{签名}

\begin{lstlisting}[language=Python]
def zero_torque(self) -> int
\end{lstlisting}

\textbf{可用性与安全性}

\textbf{\passthrough{\lstinline!SIGNATURE\_ONLY!} /
\passthrough{\lstinline!MOTION\_COMMAND!}}：仅用于补全和静态检查。当前二进制没有该
Python 方法；不得按可执行运动接口使用。

\textbf{参数}

无显式参数。实例方法中的 \passthrough{\lstinline!self!} 由 Python
自动传入。

\textbf{返回值}

\begin{longtable}[]{@{}
  >{\raggedright\arraybackslash}p{(\columnwidth - 4\tabcolsep) * \real{0.18}}
  >{\raggedright\arraybackslash}p{(\columnwidth - 4\tabcolsep) * \real{0.20}}
  >{\raggedright\arraybackslash}p{(\columnwidth - 4\tabcolsep) * \real{0.62}}@{}}
\toprule
位置 & 类型 & 含义 \\
\midrule
\endhead
返回值 & \passthrough{\lstinline!int!} & SDK 状态码；Unitree 示例通常以
\passthrough{\lstinline!0!} 表示成功，非零值需按具体服务定义解释。 \\
\bottomrule
\end{longtable}

\textbf{对应 C++}

\begin{itemize}
\tightlist
\item
  类：\passthrough{\lstinline!unitree::robot::g1::LocoClient!}
\item
  签名：\passthrough{\lstinline!ZeroTorque()!}
\item
  绑定策略：\passthrough{\lstinline!DIRECT!}
\item
  声明位置：\passthrough{\lstinline!include/unitree/robot/g1/loco/g1\_loco\_client.hpp:224!}
\end{itemize}

\textbf{说明}

\begin{itemize}
\tightlist
\item
  示例是局部调用片段；\passthrough{\lstinline!obj!}
  和参数变量表示已经按上层业务校验并准备好的对象。
\item
  该条目可以被编辑器和类型检查器识别，但当前运行时可能在属性查找阶段直接失败。
\item
  该方法属于运动命令。方法名包含
  \passthrough{\lstinline!stop!}、\passthrough{\lstinline!damp!} 或
  \passthrough{\lstinline!zero!} 也不代表它是独立物理急停。
\end{itemize}

\textbf{用法}

\begin{lstlisting}[language=Python]
from typing import TYPE_CHECKING

if TYPE_CHECKING:
    def planned_zero_torque(obj: LocoClient) -> int:
        return obj.zero_torque()
\end{lstlisting}

\begin{quote}
{[}!CAUTION{]} 上面的 \passthrough{\lstinline!TYPE\_CHECKING!}
示例只表达计划调用形态。该分支在普通 Python
运行时为假，不代表当前扩展能执行此接口。
\end{quote}

\hypertarget{unitree-sdk2-cpp-robot-g1-lococlient-stop-move-1}{%
\paragraph{\texorpdfstring{\texttt{unitree\_sdk2\_cpp.robot.g1.LocoClient.stop\_move}}{unitree\_sdk2\_cpp.robot.g1.LocoClient.stop\_move}}\label{unitree-sdk2-cpp-robot-g1-lococlient-stop-move-1}}

请求停止对应 SDK 操作。该名称不等同于经过验证的物理急停。

\textbf{签名}

\begin{lstlisting}[language=Python]
def stop_move(self) -> int
\end{lstlisting}

\textbf{可用性与安全性}

\textbf{\passthrough{\lstinline!SIGNATURE\_ONLY!} /
\passthrough{\lstinline!MOTION\_COMMAND!}}：仅用于补全和静态检查。当前二进制没有该
Python 方法；不得按可执行运动接口使用。

\textbf{参数}

无显式参数。实例方法中的 \passthrough{\lstinline!self!} 由 Python
自动传入。

\textbf{返回值}

\begin{longtable}[]{@{}
  >{\raggedright\arraybackslash}p{(\columnwidth - 4\tabcolsep) * \real{0.18}}
  >{\raggedright\arraybackslash}p{(\columnwidth - 4\tabcolsep) * \real{0.20}}
  >{\raggedright\arraybackslash}p{(\columnwidth - 4\tabcolsep) * \real{0.62}}@{}}
\toprule
位置 & 类型 & 含义 \\
\midrule
\endhead
返回值 & \passthrough{\lstinline!int!} & SDK 状态码；Unitree 示例通常以
\passthrough{\lstinline!0!} 表示成功，非零值需按具体服务定义解释。 \\
\bottomrule
\end{longtable}

\textbf{对应 C++}

\begin{itemize}
\tightlist
\item
  类：\passthrough{\lstinline!unitree::robot::g1::LocoClient!}
\item
  签名：\passthrough{\lstinline!StopMove()!}
\item
  绑定策略：\passthrough{\lstinline!DIRECT!}
\item
  声明位置：\passthrough{\lstinline!include/unitree/robot/g1/loco/g1\_loco\_client.hpp:226!}
\end{itemize}

\textbf{说明}

\begin{itemize}
\tightlist
\item
  示例是局部调用片段；\passthrough{\lstinline!obj!}
  和参数变量表示已经按上层业务校验并准备好的对象。
\item
  该条目可以被编辑器和类型检查器识别，但当前运行时可能在属性查找阶段直接失败。
\item
  该方法属于运动命令。方法名包含
  \passthrough{\lstinline!stop!}、\passthrough{\lstinline!damp!} 或
  \passthrough{\lstinline!zero!} 也不代表它是独立物理急停。
\end{itemize}

\textbf{用法}

\begin{lstlisting}[language=Python]
from typing import TYPE_CHECKING

if TYPE_CHECKING:
    def planned_stop_move(obj: LocoClient) -> int:
        return obj.stop_move()
\end{lstlisting}

\begin{quote}
{[}!CAUTION{]} 上面的 \passthrough{\lstinline!TYPE\_CHECKING!}
示例只表达计划调用形态。该分支在普通 Python
运行时为假，不代表当前扩展能执行此接口。
\end{quote}

\hypertarget{unitree-sdk2-cpp-robot-g1-lococlient-high-stand-1}{%
\paragraph{\texorpdfstring{\texttt{unitree\_sdk2\_cpp.robot.g1.LocoClient.high\_stand}}{unitree\_sdk2\_cpp.robot.g1.LocoClient.high\_stand}}\label{unitree-sdk2-cpp-robot-g1-lococlient-high-stand-1}}

对应 C++ SDK 操作
\passthrough{\lstinline!HighStand()!}。上游头文件没有可直接生成的业务说明时，本参考只保证签名映射，精确语义需查目标型号协议。

\textbf{签名}

\begin{lstlisting}[language=Python]
def high_stand(self) -> int
\end{lstlisting}

\textbf{可用性与安全性}

\textbf{\passthrough{\lstinline!SIGNATURE\_ONLY!} /
\passthrough{\lstinline!MOTION\_COMMAND!}}：仅用于补全和静态检查。当前二进制没有该
Python 方法；不得按可执行运动接口使用。

\textbf{参数}

无显式参数。实例方法中的 \passthrough{\lstinline!self!} 由 Python
自动传入。

\textbf{返回值}

\begin{longtable}[]{@{}
  >{\raggedright\arraybackslash}p{(\columnwidth - 4\tabcolsep) * \real{0.18}}
  >{\raggedright\arraybackslash}p{(\columnwidth - 4\tabcolsep) * \real{0.20}}
  >{\raggedright\arraybackslash}p{(\columnwidth - 4\tabcolsep) * \real{0.62}}@{}}
\toprule
位置 & 类型 & 含义 \\
\midrule
\endhead
返回值 & \passthrough{\lstinline!int!} & SDK 状态码；Unitree 示例通常以
\passthrough{\lstinline!0!} 表示成功，非零值需按具体服务定义解释。 \\
\bottomrule
\end{longtable}

\textbf{对应 C++}

\begin{itemize}
\tightlist
\item
  类：\passthrough{\lstinline!unitree::robot::g1::LocoClient!}
\item
  签名：\passthrough{\lstinline!HighStand()!}
\item
  绑定策略：\passthrough{\lstinline!DIRECT!}
\item
  声明位置：\passthrough{\lstinline!include/unitree/robot/g1/loco/g1\_loco\_client.hpp:228!}
\end{itemize}

\textbf{说明}

\begin{itemize}
\tightlist
\item
  示例是局部调用片段；\passthrough{\lstinline!obj!}
  和参数变量表示已经按上层业务校验并准备好的对象。
\item
  该条目可以被编辑器和类型检查器识别，但当前运行时可能在属性查找阶段直接失败。
\item
  该方法属于运动命令。方法名包含
  \passthrough{\lstinline!stop!}、\passthrough{\lstinline!damp!} 或
  \passthrough{\lstinline!zero!} 也不代表它是独立物理急停。
\end{itemize}

\textbf{用法}

\begin{lstlisting}[language=Python]
from typing import TYPE_CHECKING

if TYPE_CHECKING:
    def planned_high_stand(obj: LocoClient) -> int:
        return obj.high_stand()
\end{lstlisting}

\begin{quote}
{[}!CAUTION{]} 上面的 \passthrough{\lstinline!TYPE\_CHECKING!}
示例只表达计划调用形态。该分支在普通 Python
运行时为假，不代表当前扩展能执行此接口。
\end{quote}

\hypertarget{unitree-sdk2-cpp-robot-g1-lococlient-low-stand-1}{%
\paragraph{\texorpdfstring{\texttt{unitree\_sdk2\_cpp.robot.g1.LocoClient.low\_stand}}{unitree\_sdk2\_cpp.robot.g1.LocoClient.low\_stand}}\label{unitree-sdk2-cpp-robot-g1-lococlient-low-stand-1}}

对应 C++ SDK 操作
\passthrough{\lstinline!LowStand()!}。上游头文件没有可直接生成的业务说明时，本参考只保证签名映射，精确语义需查目标型号协议。

\textbf{签名}

\begin{lstlisting}[language=Python]
def low_stand(self) -> int
\end{lstlisting}

\textbf{可用性与安全性}

\textbf{\passthrough{\lstinline!SIGNATURE\_ONLY!} /
\passthrough{\lstinline!MOTION\_COMMAND!}}：仅用于补全和静态检查。当前二进制没有该
Python 方法；不得按可执行运动接口使用。

\textbf{参数}

无显式参数。实例方法中的 \passthrough{\lstinline!self!} 由 Python
自动传入。

\textbf{返回值}

\begin{longtable}[]{@{}
  >{\raggedright\arraybackslash}p{(\columnwidth - 4\tabcolsep) * \real{0.18}}
  >{\raggedright\arraybackslash}p{(\columnwidth - 4\tabcolsep) * \real{0.20}}
  >{\raggedright\arraybackslash}p{(\columnwidth - 4\tabcolsep) * \real{0.62}}@{}}
\toprule
位置 & 类型 & 含义 \\
\midrule
\endhead
返回值 & \passthrough{\lstinline!int!} & SDK 状态码；Unitree 示例通常以
\passthrough{\lstinline!0!} 表示成功，非零值需按具体服务定义解释。 \\
\bottomrule
\end{longtable}

\textbf{对应 C++}

\begin{itemize}
\tightlist
\item
  类：\passthrough{\lstinline!unitree::robot::g1::LocoClient!}
\item
  签名：\passthrough{\lstinline!LowStand()!}
\item
  绑定策略：\passthrough{\lstinline!DIRECT!}
\item
  声明位置：\passthrough{\lstinline!include/unitree/robot/g1/loco/g1\_loco\_client.hpp:230!}
\end{itemize}

\textbf{说明}

\begin{itemize}
\tightlist
\item
  示例是局部调用片段；\passthrough{\lstinline!obj!}
  和参数变量表示已经按上层业务校验并准备好的对象。
\item
  该条目可以被编辑器和类型检查器识别，但当前运行时可能在属性查找阶段直接失败。
\item
  该方法属于运动命令。方法名包含
  \passthrough{\lstinline!stop!}、\passthrough{\lstinline!damp!} 或
  \passthrough{\lstinline!zero!} 也不代表它是独立物理急停。
\end{itemize}

\textbf{用法}

\begin{lstlisting}[language=Python]
from typing import TYPE_CHECKING

if TYPE_CHECKING:
    def planned_low_stand(obj: LocoClient) -> int:
        return obj.low_stand()
\end{lstlisting}

\begin{quote}
{[}!CAUTION{]} 上面的 \passthrough{\lstinline!TYPE\_CHECKING!}
示例只表达计划调用形态。该分支在普通 Python
运行时为假，不代表当前扩展能执行此接口。
\end{quote}

\hypertarget{unitree-sdk2-cpp-robot-g1-lococlient-move-1}{%
\paragraph{\texorpdfstring{\texttt{unitree\_sdk2\_cpp.robot.g1.LocoClient.move}（重载
1/2）}{unitree\_sdk2\_cpp.robot.g1.LocoClient.move（重载 1/2）}}\label{unitree-sdk2-cpp-robot-g1-lococlient-move-1}}

对应 C++ SDK 操作
\passthrough{\lstinline!Move(float, float, float, bool)!}。上游头文件没有可直接生成的业务说明时，本参考只保证签名映射，精确语义需查目标型号协议。

\textbf{签名}

\begin{lstlisting}[language=Python]
@overload
def move(self, vx: float, vy: float, vyaw: float, continous_move: bool) -> int
\end{lstlisting}

\textbf{可用性与安全性}

\textbf{\passthrough{\lstinline!SIGNATURE\_ONLY!} /
\passthrough{\lstinline!MOTION\_COMMAND!}}：仅用于补全和静态检查。当前二进制没有该
Python 方法；不得按可执行运动接口使用。

\textbf{参数}

\begin{longtable}[]{@{}
  >{\raggedright\arraybackslash}p{(\columnwidth - 6\tabcolsep) * \real{0.18}}
  >{\raggedright\arraybackslash}p{(\columnwidth - 6\tabcolsep) * \real{0.24}}
  >{\raggedright\arraybackslash}p{(\columnwidth - 6\tabcolsep) * \real{0.18}}
  >{\raggedright\arraybackslash}p{(\columnwidth - 6\tabcolsep) * \real{0.40}}@{}}
\toprule
名称 & Python 类型 & 默认值 & 含义 \\
\midrule
\endhead
\passthrough{\lstinline!vx!} & \passthrough{\lstinline!float!} & 必填 &
X 方向速度参数。坐标系、单位、符号和安全范围必须查目标型号运动协议。
对应 C++ 参数 \passthrough{\lstinline!vx: float!}。
底层为浮点数；类型本身不说明单位、坐标系或业务安全范围。 \\
\passthrough{\lstinline!vy!} & \passthrough{\lstinline!float!} & 必填 &
Y 方向速度参数。坐标系、单位、符号和安全范围必须查目标型号运动协议。
对应 C++ 参数 \passthrough{\lstinline!vy: float!}。
底层为浮点数；类型本身不说明单位、坐标系或业务安全范围。 \\
\passthrough{\lstinline!vyaw!} & \passthrough{\lstinline!float!} & 必填
& 偏航角速度参数。单位、符号和安全范围必须查目标型号运动协议。 对应 C++
参数 \passthrough{\lstinline!vyaw: float!}。
底层为浮点数；类型本身不说明单位、坐标系或业务安全范围。 \\
\passthrough{\lstinline!continous\_move!} &
\passthrough{\lstinline!bool!} & 必填 & 传给该接口的
\passthrough{\lstinline!continous!} \passthrough{\lstinline!move!}
参数，Python 类型为
\passthrough{\lstinline!bool!}。精确含义、单位、范围和枚举值需查目标型号对应头文件与协议。
对应 C++ 参数 \passthrough{\lstinline!continous\_move: bool!}。
只接受布尔语义。 \\
\bottomrule
\end{longtable}

\textbf{返回值}

\begin{longtable}[]{@{}
  >{\raggedright\arraybackslash}p{(\columnwidth - 4\tabcolsep) * \real{0.18}}
  >{\raggedright\arraybackslash}p{(\columnwidth - 4\tabcolsep) * \real{0.20}}
  >{\raggedright\arraybackslash}p{(\columnwidth - 4\tabcolsep) * \real{0.62}}@{}}
\toprule
位置 & 类型 & 含义 \\
\midrule
\endhead
返回值 & \passthrough{\lstinline!int!} & SDK 状态码；Unitree 示例通常以
\passthrough{\lstinline!0!} 表示成功，非零值需按具体服务定义解释。 \\
\bottomrule
\end{longtable}

\textbf{对应 C++}

\begin{itemize}
\tightlist
\item
  类：\passthrough{\lstinline!unitree::robot::g1::LocoClient!}
\item
  签名：\passthrough{\lstinline!Move(float, float, float, bool)!}
\item
  绑定策略：\passthrough{\lstinline!DIRECT!}
\item
  声明位置：\passthrough{\lstinline!include/unitree/robot/g1/loco/g1\_loco\_client.hpp:232!}
\end{itemize}

\textbf{说明}

\begin{itemize}
\tightlist
\item
  示例是局部调用片段；\passthrough{\lstinline!obj!}
  和参数变量表示已经按上层业务校验并准备好的对象。
\item
  该条目可以被编辑器和类型检查器识别，但当前运行时可能在属性查找阶段直接失败。
\item
  该方法属于运动命令。方法名包含
  \passthrough{\lstinline!stop!}、\passthrough{\lstinline!damp!} 或
  \passthrough{\lstinline!zero!} 也不代表它是独立物理急停。
\end{itemize}

\textbf{用法}

\begin{lstlisting}[language=Python]
from typing import TYPE_CHECKING

if TYPE_CHECKING:
    def planned_move(obj: LocoClient, vx: float, vy: float, vyaw: float, continous_move: bool) -> int:
        return obj.move(vx=vx, vy=vy, vyaw=vyaw, continous_move=continous_move)
\end{lstlisting}

\begin{quote}
{[}!CAUTION{]} 上面的 \passthrough{\lstinline!TYPE\_CHECKING!}
示例只表达计划调用形态。该分支在普通 Python
运行时为假，不代表当前扩展能执行此接口。
\end{quote}

\hypertarget{unitree-sdk2-cpp-robot-g1-lococlient-move-2}{%
\paragraph{\texorpdfstring{\texttt{unitree\_sdk2\_cpp.robot.g1.LocoClient.move}（重载
2/2）}{unitree\_sdk2\_cpp.robot.g1.LocoClient.move（重载 2/2）}}\label{unitree-sdk2-cpp-robot-g1-lococlient-move-2}}

对应 C++ SDK 操作
\passthrough{\lstinline!Move(float, float, float)!}。上游头文件没有可直接生成的业务说明时，本参考只保证签名映射，精确语义需查目标型号协议。

\textbf{签名}

\begin{lstlisting}[language=Python]
@overload
def move(self, vx: float, vy: float, vyaw: float) -> int
\end{lstlisting}

\textbf{可用性与安全性}

\textbf{\passthrough{\lstinline!SIGNATURE\_ONLY!} /
\passthrough{\lstinline!MOTION\_COMMAND!}}：仅用于补全和静态检查。当前二进制没有该
Python 方法；不得按可执行运动接口使用。

\textbf{参数}

\begin{longtable}[]{@{}
  >{\raggedright\arraybackslash}p{(\columnwidth - 6\tabcolsep) * \real{0.18}}
  >{\raggedright\arraybackslash}p{(\columnwidth - 6\tabcolsep) * \real{0.24}}
  >{\raggedright\arraybackslash}p{(\columnwidth - 6\tabcolsep) * \real{0.18}}
  >{\raggedright\arraybackslash}p{(\columnwidth - 6\tabcolsep) * \real{0.40}}@{}}
\toprule
名称 & Python 类型 & 默认值 & 含义 \\
\midrule
\endhead
\passthrough{\lstinline!vx!} & \passthrough{\lstinline!float!} & 必填 &
X 方向速度参数。坐标系、单位、符号和安全范围必须查目标型号运动协议。
对应 C++ 参数 \passthrough{\lstinline!vx: float!}。
底层为浮点数；类型本身不说明单位、坐标系或业务安全范围。 \\
\passthrough{\lstinline!vy!} & \passthrough{\lstinline!float!} & 必填 &
Y 方向速度参数。坐标系、单位、符号和安全范围必须查目标型号运动协议。
对应 C++ 参数 \passthrough{\lstinline!vy: float!}。
底层为浮点数；类型本身不说明单位、坐标系或业务安全范围。 \\
\passthrough{\lstinline!vyaw!} & \passthrough{\lstinline!float!} & 必填
& 偏航角速度参数。单位、符号和安全范围必须查目标型号运动协议。 对应 C++
参数 \passthrough{\lstinline!vyaw: float!}。
底层为浮点数；类型本身不说明单位、坐标系或业务安全范围。 \\
\bottomrule
\end{longtable}

\textbf{返回值}

\begin{longtable}[]{@{}
  >{\raggedright\arraybackslash}p{(\columnwidth - 4\tabcolsep) * \real{0.18}}
  >{\raggedright\arraybackslash}p{(\columnwidth - 4\tabcolsep) * \real{0.20}}
  >{\raggedright\arraybackslash}p{(\columnwidth - 4\tabcolsep) * \real{0.62}}@{}}
\toprule
位置 & 类型 & 含义 \\
\midrule
\endhead
返回值 & \passthrough{\lstinline!int!} & SDK 状态码；Unitree 示例通常以
\passthrough{\lstinline!0!} 表示成功，非零值需按具体服务定义解释。 \\
\bottomrule
\end{longtable}

\textbf{对应 C++}

\begin{itemize}
\tightlist
\item
  类：\passthrough{\lstinline!unitree::robot::g1::LocoClient!}
\item
  签名：\passthrough{\lstinline!Move(float, float, float)!}
\item
  绑定策略：\passthrough{\lstinline!DIRECT!}
\item
  声明位置：\passthrough{\lstinline!include/unitree/robot/g1/loco/g1\_loco\_client.hpp:236!}
\end{itemize}

\textbf{说明}

\begin{itemize}
\tightlist
\item
  示例是局部调用片段；\passthrough{\lstinline!obj!}
  和参数变量表示已经按上层业务校验并准备好的对象。
\item
  该条目可以被编辑器和类型检查器识别，但当前运行时可能在属性查找阶段直接失败。
\item
  该方法属于运动命令。方法名包含
  \passthrough{\lstinline!stop!}、\passthrough{\lstinline!damp!} 或
  \passthrough{\lstinline!zero!} 也不代表它是独立物理急停。
\end{itemize}

\textbf{用法}

\begin{lstlisting}[language=Python]
from typing import TYPE_CHECKING

if TYPE_CHECKING:
    def planned_move(obj: LocoClient, vx: float, vy: float, vyaw: float) -> int:
        return obj.move(vx=vx, vy=vy, vyaw=vyaw)
\end{lstlisting}

\begin{quote}
{[}!CAUTION{]} 上面的 \passthrough{\lstinline!TYPE\_CHECKING!}
示例只表达计划调用形态。该分支在普通 Python
运行时为假，不代表当前扩展能执行此接口。
\end{quote}

\hypertarget{unitree-sdk2-cpp-robot-g1-lococlient-balance-stand-1}{%
\paragraph{\texorpdfstring{\texttt{unitree\_sdk2\_cpp.robot.g1.LocoClient.balance\_stand}}{unitree\_sdk2\_cpp.robot.g1.LocoClient.balance\_stand}}\label{unitree-sdk2-cpp-robot-g1-lococlient-balance-stand-1}}

对应 C++ SDK 操作
\passthrough{\lstinline!BalanceStand()!}。上游头文件没有可直接生成的业务说明时，本参考只保证签名映射，精确语义需查目标型号协议。

\textbf{签名}

\begin{lstlisting}[language=Python]
def balance_stand(self) -> int
\end{lstlisting}

\textbf{可用性与安全性}

\textbf{\passthrough{\lstinline!SIGNATURE\_ONLY!} /
\passthrough{\lstinline!MOTION\_COMMAND!}}：仅用于补全和静态检查。当前二进制没有该
Python 方法；不得按可执行运动接口使用。

\textbf{参数}

无显式参数。实例方法中的 \passthrough{\lstinline!self!} 由 Python
自动传入。

\textbf{返回值}

\begin{longtable}[]{@{}
  >{\raggedright\arraybackslash}p{(\columnwidth - 4\tabcolsep) * \real{0.18}}
  >{\raggedright\arraybackslash}p{(\columnwidth - 4\tabcolsep) * \real{0.20}}
  >{\raggedright\arraybackslash}p{(\columnwidth - 4\tabcolsep) * \real{0.62}}@{}}
\toprule
位置 & 类型 & 含义 \\
\midrule
\endhead
返回值 & \passthrough{\lstinline!int!} & SDK 状态码；Unitree 示例通常以
\passthrough{\lstinline!0!} 表示成功，非零值需按具体服务定义解释。 \\
\bottomrule
\end{longtable}

\textbf{对应 C++}

\begin{itemize}
\tightlist
\item
  类：\passthrough{\lstinline!unitree::robot::g1::LocoClient!}
\item
  签名：\passthrough{\lstinline!BalanceStand()!}
\item
  绑定策略：\passthrough{\lstinline!DIRECT!}
\item
  声明位置：\passthrough{\lstinline!include/unitree/robot/g1/loco/g1\_loco\_client.hpp:238!}
\end{itemize}

\textbf{说明}

\begin{itemize}
\tightlist
\item
  示例是局部调用片段；\passthrough{\lstinline!obj!}
  和参数变量表示已经按上层业务校验并准备好的对象。
\item
  该条目可以被编辑器和类型检查器识别，但当前运行时可能在属性查找阶段直接失败。
\item
  该方法属于运动命令。方法名包含
  \passthrough{\lstinline!stop!}、\passthrough{\lstinline!damp!} 或
  \passthrough{\lstinline!zero!} 也不代表它是独立物理急停。
\end{itemize}

\textbf{用法}

\begin{lstlisting}[language=Python]
from typing import TYPE_CHECKING

if TYPE_CHECKING:
    def planned_balance_stand(obj: LocoClient) -> int:
        return obj.balance_stand()
\end{lstlisting}

\begin{quote}
{[}!CAUTION{]} 上面的 \passthrough{\lstinline!TYPE\_CHECKING!}
示例只表达计划调用形态。该分支在普通 Python
运行时为假，不代表当前扩展能执行此接口。
\end{quote}

\hypertarget{unitree-sdk2-cpp-robot-g1-lococlient-continuous-gait-1}{%
\paragraph{\texorpdfstring{\texttt{unitree\_sdk2\_cpp.robot.g1.LocoClient.continuous\_gait}}{unitree\_sdk2\_cpp.robot.g1.LocoClient.continuous\_gait}}\label{unitree-sdk2-cpp-robot-g1-lococlient-continuous-gait-1}}

对应 C++ SDK 操作
\passthrough{\lstinline!ContinuousGait(bool)!}。上游头文件没有可直接生成的业务说明时，本参考只保证签名映射，精确语义需查目标型号协议。

\textbf{签名}

\begin{lstlisting}[language=Python]
def continuous_gait(self, flag: bool) -> int
\end{lstlisting}

\textbf{可用性与安全性}

\textbf{\passthrough{\lstinline!SIGNATURE\_ONLY!} /
\passthrough{\lstinline!MOTION\_COMMAND!}}：仅用于补全和静态检查。当前二进制没有该
Python 方法；不得按可执行运动接口使用。

\textbf{参数}

\begin{longtable}[]{@{}
  >{\raggedright\arraybackslash}p{(\columnwidth - 6\tabcolsep) * \real{0.18}}
  >{\raggedright\arraybackslash}p{(\columnwidth - 6\tabcolsep) * \real{0.24}}
  >{\raggedright\arraybackslash}p{(\columnwidth - 6\tabcolsep) * \real{0.18}}
  >{\raggedright\arraybackslash}p{(\columnwidth - 6\tabcolsep) * \real{0.40}}@{}}
\toprule
名称 & Python 类型 & 默认值 & 含义 \\
\midrule
\endhead
\passthrough{\lstinline!flag!} & \passthrough{\lstinline!bool!} & 必填 &
布尔开关。\passthrough{\lstinline!True!} 和
\passthrough{\lstinline!False!} 的具体业务效果由当前方法定义。 对应 C++
参数 \passthrough{\lstinline!flag: bool!}。 只接受布尔语义。 \\
\bottomrule
\end{longtable}

\textbf{返回值}

\begin{longtable}[]{@{}
  >{\raggedright\arraybackslash}p{(\columnwidth - 4\tabcolsep) * \real{0.18}}
  >{\raggedright\arraybackslash}p{(\columnwidth - 4\tabcolsep) * \real{0.20}}
  >{\raggedright\arraybackslash}p{(\columnwidth - 4\tabcolsep) * \real{0.62}}@{}}
\toprule
位置 & 类型 & 含义 \\
\midrule
\endhead
返回值 & \passthrough{\lstinline!int!} & SDK 状态码；Unitree 示例通常以
\passthrough{\lstinline!0!} 表示成功，非零值需按具体服务定义解释。 \\
\bottomrule
\end{longtable}

\textbf{对应 C++}

\begin{itemize}
\tightlist
\item
  类：\passthrough{\lstinline!unitree::robot::g1::LocoClient!}
\item
  签名：\passthrough{\lstinline!ContinuousGait(bool)!}
\item
  绑定策略：\passthrough{\lstinline!DIRECT!}
\item
  声明位置：\passthrough{\lstinline!include/unitree/robot/g1/loco/g1\_loco\_client.hpp:240!}
\end{itemize}

\textbf{说明}

\begin{itemize}
\tightlist
\item
  示例是局部调用片段；\passthrough{\lstinline!obj!}
  和参数变量表示已经按上层业务校验并准备好的对象。
\item
  该条目可以被编辑器和类型检查器识别，但当前运行时可能在属性查找阶段直接失败。
\item
  该方法属于运动命令。方法名包含
  \passthrough{\lstinline!stop!}、\passthrough{\lstinline!damp!} 或
  \passthrough{\lstinline!zero!} 也不代表它是独立物理急停。
\end{itemize}

\textbf{用法}

\begin{lstlisting}[language=Python]
from typing import TYPE_CHECKING

if TYPE_CHECKING:
    def planned_continuous_gait(obj: LocoClient, flag: bool) -> int:
        return obj.continuous_gait(flag=flag)
\end{lstlisting}

\begin{quote}
{[}!CAUTION{]} 上面的 \passthrough{\lstinline!TYPE\_CHECKING!}
示例只表达计划调用形态。该分支在普通 Python
运行时为假，不代表当前扩展能执行此接口。
\end{quote}

\hypertarget{unitree-sdk2-cpp-robot-g1-lococlient-switch-move-mode-1}{%
\paragraph{\texorpdfstring{\texttt{unitree\_sdk2\_cpp.robot.g1.LocoClient.switch\_move\_mode}}{unitree\_sdk2\_cpp.robot.g1.LocoClient.switch\_move\_mode}}\label{unitree-sdk2-cpp-robot-g1-lococlient-switch-move-mode-1}}

对应 C++ SDK 操作
\passthrough{\lstinline!SwitchMoveMode(bool)!}。上游头文件没有可直接生成的业务说明时，本参考只保证签名映射，精确语义需查目标型号协议。

\textbf{签名}

\begin{lstlisting}[language=Python]
def switch_move_mode(self, flag: bool) -> int
\end{lstlisting}

\textbf{可用性与安全性}

\textbf{\passthrough{\lstinline!SIGNATURE\_ONLY!} /
\passthrough{\lstinline!MOTION\_COMMAND!}}：仅用于补全和静态检查。当前二进制没有该
Python 方法；不得按可执行运动接口使用。

\textbf{参数}

\begin{longtable}[]{@{}
  >{\raggedright\arraybackslash}p{(\columnwidth - 6\tabcolsep) * \real{0.18}}
  >{\raggedright\arraybackslash}p{(\columnwidth - 6\tabcolsep) * \real{0.24}}
  >{\raggedright\arraybackslash}p{(\columnwidth - 6\tabcolsep) * \real{0.18}}
  >{\raggedright\arraybackslash}p{(\columnwidth - 6\tabcolsep) * \real{0.40}}@{}}
\toprule
名称 & Python 类型 & 默认值 & 含义 \\
\midrule
\endhead
\passthrough{\lstinline!flag!} & \passthrough{\lstinline!bool!} & 必填 &
布尔开关。\passthrough{\lstinline!True!} 和
\passthrough{\lstinline!False!} 的具体业务效果由当前方法定义。 对应 C++
参数 \passthrough{\lstinline!flag: bool!}。 只接受布尔语义。 \\
\bottomrule
\end{longtable}

\textbf{返回值}

\begin{longtable}[]{@{}
  >{\raggedright\arraybackslash}p{(\columnwidth - 4\tabcolsep) * \real{0.18}}
  >{\raggedright\arraybackslash}p{(\columnwidth - 4\tabcolsep) * \real{0.20}}
  >{\raggedright\arraybackslash}p{(\columnwidth - 4\tabcolsep) * \real{0.62}}@{}}
\toprule
位置 & 类型 & 含义 \\
\midrule
\endhead
返回值 & \passthrough{\lstinline!int!} & SDK 状态码；Unitree 示例通常以
\passthrough{\lstinline!0!} 表示成功，非零值需按具体服务定义解释。 \\
\bottomrule
\end{longtable}

\textbf{对应 C++}

\begin{itemize}
\tightlist
\item
  类：\passthrough{\lstinline!unitree::robot::g1::LocoClient!}
\item
  签名：\passthrough{\lstinline!SwitchMoveMode(bool)!}
\item
  绑定策略：\passthrough{\lstinline!DIRECT!}
\item
  声明位置：\passthrough{\lstinline!include/unitree/robot/g1/loco/g1\_loco\_client.hpp:242!}
\end{itemize}

\textbf{说明}

\begin{itemize}
\tightlist
\item
  示例是局部调用片段；\passthrough{\lstinline!obj!}
  和参数变量表示已经按上层业务校验并准备好的对象。
\item
  该条目可以被编辑器和类型检查器识别，但当前运行时可能在属性查找阶段直接失败。
\item
  该方法属于运动命令。方法名包含
  \passthrough{\lstinline!stop!}、\passthrough{\lstinline!damp!} 或
  \passthrough{\lstinline!zero!} 也不代表它是独立物理急停。
\end{itemize}

\textbf{用法}

\begin{lstlisting}[language=Python]
from typing import TYPE_CHECKING

if TYPE_CHECKING:
    def planned_switch_move_mode(obj: LocoClient, flag: bool) -> int:
        return obj.switch_move_mode(flag=flag)
\end{lstlisting}

\begin{quote}
{[}!CAUTION{]} 上面的 \passthrough{\lstinline!TYPE\_CHECKING!}
示例只表达计划调用形态。该分支在普通 Python
运行时为假，不代表当前扩展能执行此接口。
\end{quote}

\hypertarget{unitree-sdk2-cpp-robot-g1-lococlient-wave-hand-1}{%
\paragraph{\texorpdfstring{\texttt{unitree\_sdk2\_cpp.robot.g1.LocoClient.wave\_hand}}{unitree\_sdk2\_cpp.robot.g1.LocoClient.wave\_hand}}\label{unitree-sdk2-cpp-robot-g1-lococlient-wave-hand-1}}

对应 C++ SDK 操作
\passthrough{\lstinline!WaveHand(bool)!}。上游头文件没有可直接生成的业务说明时，本参考只保证签名映射，精确语义需查目标型号协议。

\textbf{签名}

\begin{lstlisting}[language=Python]
def wave_hand(self, turn_flag: bool = ...) -> int
\end{lstlisting}

\textbf{可用性与安全性}

\textbf{\passthrough{\lstinline!SIGNATURE\_ONLY!} /
\passthrough{\lstinline!MOTION\_COMMAND!}}：仅用于补全和静态检查。当前二进制没有该
Python 方法；不得按可执行运动接口使用。

\textbf{参数}

\begin{longtable}[]{@{}
  >{\raggedright\arraybackslash}p{(\columnwidth - 6\tabcolsep) * \real{0.18}}
  >{\raggedright\arraybackslash}p{(\columnwidth - 6\tabcolsep) * \real{0.24}}
  >{\raggedright\arraybackslash}p{(\columnwidth - 6\tabcolsep) * \real{0.18}}
  >{\raggedright\arraybackslash}p{(\columnwidth - 6\tabcolsep) * \real{0.40}}@{}}
\toprule
名称 & Python 类型 & 默认值 & 含义 \\
\midrule
\endhead
\passthrough{\lstinline!turn\_flag!} & \passthrough{\lstinline!bool!} &
\passthrough{\lstinline!...!} & 传给该接口的
\passthrough{\lstinline!turn!} \passthrough{\lstinline!flag!}
参数，Python 类型为
\passthrough{\lstinline!bool!}。精确含义、单位、范围和枚举值需查目标型号对应头文件与协议。
对应 C++ 参数 \passthrough{\lstinline!turn\_flag: bool!}。
只接受布尔语义。 \\
\bottomrule
\end{longtable}

\textbf{返回值}

\begin{longtable}[]{@{}
  >{\raggedright\arraybackslash}p{(\columnwidth - 4\tabcolsep) * \real{0.18}}
  >{\raggedright\arraybackslash}p{(\columnwidth - 4\tabcolsep) * \real{0.20}}
  >{\raggedright\arraybackslash}p{(\columnwidth - 4\tabcolsep) * \real{0.62}}@{}}
\toprule
位置 & 类型 & 含义 \\
\midrule
\endhead
返回值 & \passthrough{\lstinline!int!} & SDK 状态码；Unitree 示例通常以
\passthrough{\lstinline!0!} 表示成功，非零值需按具体服务定义解释。 \\
\bottomrule
\end{longtable}

\textbf{对应 C++}

\begin{itemize}
\tightlist
\item
  类：\passthrough{\lstinline!unitree::robot::g1::LocoClient!}
\item
  签名：\passthrough{\lstinline!WaveHand(bool)!}
\item
  绑定策略：\passthrough{\lstinline!DIRECT!}
\item
  声明位置：\passthrough{\lstinline!include/unitree/robot/g1/loco/g1\_loco\_client.hpp:247!}
\end{itemize}

\textbf{说明}

\begin{itemize}
\tightlist
\item
  示例是局部调用片段；\passthrough{\lstinline!obj!}
  和参数变量表示已经按上层业务校验并准备好的对象。
\item
  默认值显示为 \passthrough{\lstinline!...!}，表示 C++
  声明存在默认参数，但当前 AST
  清单没有保存其字面量；省略参数可使用上游默认行为。
\item
  该条目可以被编辑器和类型检查器识别，但当前运行时可能在属性查找阶段直接失败。
\item
  该方法属于运动命令。方法名包含
  \passthrough{\lstinline!stop!}、\passthrough{\lstinline!damp!} 或
  \passthrough{\lstinline!zero!} 也不代表它是独立物理急停。
\end{itemize}

\textbf{用法}

\begin{lstlisting}[language=Python]
from typing import TYPE_CHECKING

if TYPE_CHECKING:
    def planned_wave_hand(obj: LocoClient, turn_flag: bool) -> int:
        return obj.wave_hand(turn_flag=turn_flag)
\end{lstlisting}

\begin{quote}
{[}!CAUTION{]} 上面的 \passthrough{\lstinline!TYPE\_CHECKING!}
示例只表达计划调用形态。该分支在普通 Python
运行时为假，不代表当前扩展能执行此接口。
\end{quote}

\hypertarget{unitree-sdk2-cpp-robot-g1-lococlient-shake-hand-1}{%
\paragraph{\texorpdfstring{\texttt{unitree\_sdk2\_cpp.robot.g1.LocoClient.shake\_hand}}{unitree\_sdk2\_cpp.robot.g1.LocoClient.shake\_hand}}\label{unitree-sdk2-cpp-robot-g1-lococlient-shake-hand-1}}

对应 C++ SDK 操作
\passthrough{\lstinline!ShakeHand(int)!}。上游头文件没有可直接生成的业务说明时，本参考只保证签名映射，精确语义需查目标型号协议。

\textbf{签名}

\begin{lstlisting}[language=Python]
def shake_hand(self, stage: int = ...) -> int
\end{lstlisting}

\textbf{可用性与安全性}

\textbf{\passthrough{\lstinline!SIGNATURE\_ONLY!} /
\passthrough{\lstinline!MOTION\_COMMAND!}}：仅用于补全和静态检查。当前二进制没有该
Python 方法；不得按可执行运动接口使用。

\textbf{参数}

\begin{longtable}[]{@{}
  >{\raggedright\arraybackslash}p{(\columnwidth - 6\tabcolsep) * \real{0.18}}
  >{\raggedright\arraybackslash}p{(\columnwidth - 6\tabcolsep) * \real{0.24}}
  >{\raggedright\arraybackslash}p{(\columnwidth - 6\tabcolsep) * \real{0.18}}
  >{\raggedright\arraybackslash}p{(\columnwidth - 6\tabcolsep) * \real{0.40}}@{}}
\toprule
名称 & Python 类型 & 默认值 & 含义 \\
\midrule
\endhead
\passthrough{\lstinline!stage!} & \passthrough{\lstinline!int!} &
\passthrough{\lstinline!...!} & 传给该接口的
\passthrough{\lstinline!stage!} 参数，Python 类型为
\passthrough{\lstinline!int!}。精确含义、单位、范围和枚举值需查目标型号对应头文件与协议。
对应 C++ 参数 \passthrough{\lstinline!stage: int!}。 \\
\bottomrule
\end{longtable}

\textbf{返回值}

\begin{longtable}[]{@{}
  >{\raggedright\arraybackslash}p{(\columnwidth - 4\tabcolsep) * \real{0.18}}
  >{\raggedright\arraybackslash}p{(\columnwidth - 4\tabcolsep) * \real{0.20}}
  >{\raggedright\arraybackslash}p{(\columnwidth - 4\tabcolsep) * \real{0.62}}@{}}
\toprule
位置 & 类型 & 含义 \\
\midrule
\endhead
返回值 & \passthrough{\lstinline!int!} & SDK 状态码；Unitree 示例通常以
\passthrough{\lstinline!0!} 表示成功，非零值需按具体服务定义解释。 \\
\bottomrule
\end{longtable}

\textbf{对应 C++}

\begin{itemize}
\tightlist
\item
  类：\passthrough{\lstinline!unitree::robot::g1::LocoClient!}
\item
  签名：\passthrough{\lstinline!ShakeHand(int)!}
\item
  绑定策略：\passthrough{\lstinline!DIRECT!}
\item
  声明位置：\passthrough{\lstinline!include/unitree/robot/g1/loco/g1\_loco\_client.hpp:249!}
\end{itemize}

\textbf{说明}

\begin{itemize}
\tightlist
\item
  示例是局部调用片段；\passthrough{\lstinline!obj!}
  和参数变量表示已经按上层业务校验并准备好的对象。
\item
  默认值显示为 \passthrough{\lstinline!...!}，表示 C++
  声明存在默认参数，但当前 AST
  清单没有保存其字面量；省略参数可使用上游默认行为。
\item
  该条目可以被编辑器和类型检查器识别，但当前运行时可能在属性查找阶段直接失败。
\item
  该方法属于运动命令。方法名包含
  \passthrough{\lstinline!stop!}、\passthrough{\lstinline!damp!} 或
  \passthrough{\lstinline!zero!} 也不代表它是独立物理急停。
\end{itemize}

\textbf{用法}

\begin{lstlisting}[language=Python]
from typing import TYPE_CHECKING

if TYPE_CHECKING:
    def planned_shake_hand(obj: LocoClient, stage: int) -> int:
        return obj.shake_hand(stage=stage)
\end{lstlisting}

\begin{quote}
{[}!CAUTION{]} 上面的 \passthrough{\lstinline!TYPE\_CHECKING!}
示例只表达计划调用形态。该分支在普通 Python
运行时为假，不代表当前扩展能执行此接口。
\end{quote}

\hypertarget{unitree-sdk2-cpp-robot-g1-lococlient-set-speed-mode-1}{%
\paragraph{\texorpdfstring{\texttt{unitree\_sdk2\_cpp.robot.g1.LocoClient.set\_speed\_mode}}{unitree\_sdk2\_cpp.robot.g1.LocoClient.set\_speed\_mode}}\label{unitree-sdk2-cpp-robot-g1-lococlient-set-speed-mode-1}}

设置 速度 模式。具体副作用和安全边界见下方状态。

\textbf{签名}

\begin{lstlisting}[language=Python]
def set_speed_mode(self, speed_mode: int) -> int
\end{lstlisting}

\textbf{可用性与安全性}

\textbf{\passthrough{\lstinline!SIGNATURE\_ONLY!} /
\passthrough{\lstinline!MOTION\_COMMAND!}}：仅用于补全和静态检查。当前二进制没有该
Python 方法；不得按可执行运动接口使用。

\textbf{参数}

\begin{longtable}[]{@{}
  >{\raggedright\arraybackslash}p{(\columnwidth - 6\tabcolsep) * \real{0.18}}
  >{\raggedright\arraybackslash}p{(\columnwidth - 6\tabcolsep) * \real{0.24}}
  >{\raggedright\arraybackslash}p{(\columnwidth - 6\tabcolsep) * \real{0.18}}
  >{\raggedright\arraybackslash}p{(\columnwidth - 6\tabcolsep) * \real{0.40}}@{}}
\toprule
名称 & Python 类型 & 默认值 & 含义 \\
\midrule
\endhead
\passthrough{\lstinline!speed\_mode!} & \passthrough{\lstinline!int!} &
必填 & 传给该接口的 速度 模式 参数，Python 类型为
\passthrough{\lstinline!int!}。精确含义、单位、范围和枚举值需查目标型号对应头文件与协议。
对应 C++ 参数 \passthrough{\lstinline!speed\_mode: int!}。 \\
\bottomrule
\end{longtable}

\textbf{返回值}

\begin{longtable}[]{@{}
  >{\raggedright\arraybackslash}p{(\columnwidth - 4\tabcolsep) * \real{0.18}}
  >{\raggedright\arraybackslash}p{(\columnwidth - 4\tabcolsep) * \real{0.20}}
  >{\raggedright\arraybackslash}p{(\columnwidth - 4\tabcolsep) * \real{0.62}}@{}}
\toprule
位置 & 类型 & 含义 \\
\midrule
\endhead
返回值 & \passthrough{\lstinline!int!} & SDK 状态码；Unitree 示例通常以
\passthrough{\lstinline!0!} 表示成功，非零值需按具体服务定义解释。 \\
\bottomrule
\end{longtable}

\textbf{对应 C++}

\begin{itemize}
\tightlist
\item
  类：\passthrough{\lstinline!unitree::robot::g1::LocoClient!}
\item
  签名：\passthrough{\lstinline!SetSpeedMode(int)!}
\item
  绑定策略：\passthrough{\lstinline!DIRECT!}
\item
  声明位置：\passthrough{\lstinline!include/unitree/robot/g1/loco/g1\_loco\_client.hpp:265!}
\end{itemize}

\textbf{说明}

\begin{itemize}
\tightlist
\item
  示例是局部调用片段；\passthrough{\lstinline!obj!}
  和参数变量表示已经按上层业务校验并准备好的对象。
\item
  该条目可以被编辑器和类型检查器识别，但当前运行时可能在属性查找阶段直接失败。
\item
  该方法属于运动命令。方法名包含
  \passthrough{\lstinline!stop!}、\passthrough{\lstinline!damp!} 或
  \passthrough{\lstinline!zero!} 也不代表它是独立物理急停。
\end{itemize}

\textbf{用法}

\begin{lstlisting}[language=Python]
from typing import TYPE_CHECKING

if TYPE_CHECKING:
    def planned_set_speed_mode(obj: LocoClient, speed_mode: int) -> int:
        return obj.set_speed_mode(speed_mode=speed_mode)
\end{lstlisting}

\begin{quote}
{[}!CAUTION{]} 上面的 \passthrough{\lstinline!TYPE\_CHECKING!}
示例只表达计划调用形态。该分支在普通 Python
运行时为假，不代表当前扩展能执行此接口。
\end{quote}

\hypertarget{unitree-sdk2-cpp-robot-g1-lococlient-get-mimic-motion-1}{%
\paragraph{\texorpdfstring{\texttt{unitree\_sdk2\_cpp.robot.g1.LocoClient.get\_mimic\_motion}}{unitree\_sdk2\_cpp.robot.g1.LocoClient.get\_mimic\_motion}}\label{unitree-sdk2-cpp-robot-g1-lococlient-get-mimic-motion-1}}

查询或检查 \passthrough{\lstinline!mimic!}
\passthrough{\lstinline!motion!}。

\textbf{签名}

\begin{lstlisting}[language=Python]
def get_mimic_motion(self) -> tuple[int, str]
\end{lstlisting}

\textbf{可用性与安全性}

\textbf{\passthrough{\lstinline!AVAILABLE!} /
\passthrough{\lstinline!READ\_ONLY!}}：当前绑定源码已实现只读查询。Robot
Client 的服务端查询通常仍需匹配的 Linux
扩展、DDS、网络、机器人和服务；纯本地 getter 则不一定需要全部条件。

\textbf{参数}

无显式参数。实例方法中的 \passthrough{\lstinline!self!} 由 Python
自动传入。

\textbf{返回值}

\begin{longtable}[]{@{}
  >{\raggedright\arraybackslash}p{(\columnwidth - 4\tabcolsep) * \real{0.18}}
  >{\raggedright\arraybackslash}p{(\columnwidth - 4\tabcolsep) * \real{0.20}}
  >{\raggedright\arraybackslash}p{(\columnwidth - 4\tabcolsep) * \real{0.62}}@{}}
\toprule
位置 & 类型 & 含义 \\
\midrule
\endhead
{[}0{]} \passthrough{\lstinline!status!} & \passthrough{\lstinline!int!}
& SDK 状态码；Unitree 示例通常以 \passthrough{\lstinline!0!}
表示成功，非零值需按具体服务错误码解释。 \\
{[}1{]} \passthrough{\lstinline!data!} & \passthrough{\lstinline!str!} &
传给该接口的 数据 参数，Python 类型为
\passthrough{\lstinline!str!}。精确含义、单位、范围和枚举值需查目标型号对应头文件与协议。
对应 C++ 参数 \passthrough{\lstinline!data: std::string \&!}。
底层为字符串；长度、编码和允许值由具体协议决定。 \\
\bottomrule
\end{longtable}

\textbf{对应 C++}

\begin{itemize}
\tightlist
\item
  类：\passthrough{\lstinline!unitree::robot::g1::LocoClient!}
\item
  签名：\passthrough{\lstinline!GetMimicMotion(std::string \&)!}
\item
  绑定策略：\passthrough{\lstinline!OUTPUT\_WRAPPER!}
\item
  声明位置：\passthrough{\lstinline!include/unitree/robot/g1/loco/g1\_loco\_client.hpp:298!}
\end{itemize}

\textbf{说明}

\begin{itemize}
\tightlist
\item
  示例是局部调用片段；\passthrough{\lstinline!obj!}
  和参数变量表示已经按上层业务校验并准备好的对象。
\item
  C++ 的可修改输出引用不会作为 Python 入参出现，而是按 C++
  参数顺序追加到返回元组中。
\end{itemize}

\textbf{用法}

\begin{lstlisting}[language=Python]
status, data = obj.get_mimic_motion()
\end{lstlisting}

\begin{quote}
{[}!NOTE{]} 只读查询仍会访问
DDS/机器人服务。必须检查返回状态码，并处理超时和网络中断。
\end{quote}

\hypertarget{unitree-sdk2-cpp-robot-g1-lococlient-fsm-api-1}{%
\paragraph{\texorpdfstring{\texttt{unitree\_sdk2\_cpp.robot.g1.LocoClient.\_fsm\_api}}{unitree\_sdk2\_cpp.robot.g1.LocoClient.\_fsm\_api}}\label{unitree-sdk2-cpp-robot-g1-lococlient-fsm-api-1}}

对应 C++ SDK 操作
\passthrough{\lstinline!\_fsm\_api(std::string, std::string \&)!}。上游头文件没有可直接生成的业务说明时，本参考只保证签名映射，精确语义需查目标型号协议。

\textbf{签名}

\begin{lstlisting}[language=Python]
def _fsm_api(self, parameter: str) -> tuple[int, str]
\end{lstlisting}

\textbf{可用性与安全性}

\textbf{\passthrough{\lstinline!SIGNATURE\_ONLY!} /
\passthrough{\lstinline!MOTION\_COMMAND!}}：仅用于补全和静态检查。当前二进制没有该
Python 方法；不得按可执行运动接口使用。

\textbf{参数}

\begin{longtable}[]{@{}
  >{\raggedright\arraybackslash}p{(\columnwidth - 6\tabcolsep) * \real{0.18}}
  >{\raggedright\arraybackslash}p{(\columnwidth - 6\tabcolsep) * \real{0.24}}
  >{\raggedright\arraybackslash}p{(\columnwidth - 6\tabcolsep) * \real{0.18}}
  >{\raggedright\arraybackslash}p{(\columnwidth - 6\tabcolsep) * \real{0.40}}@{}}
\toprule
名称 & Python 类型 & 默认值 & 含义 \\
\midrule
\endhead
\passthrough{\lstinline!parameter!} & \passthrough{\lstinline!str!} &
必填 & 传给该接口的 \passthrough{\lstinline!parameter!} 参数，Python
类型为
\passthrough{\lstinline!str!}。精确含义、单位、范围和枚举值需查目标型号对应头文件与协议。
对应 C++ 参数 \passthrough{\lstinline!parameter: std::string!}。
底层为字符串；长度、编码和允许值由具体协议决定。 \\
\bottomrule
\end{longtable}

\textbf{返回值}

\begin{longtable}[]{@{}
  >{\raggedright\arraybackslash}p{(\columnwidth - 4\tabcolsep) * \real{0.18}}
  >{\raggedright\arraybackslash}p{(\columnwidth - 4\tabcolsep) * \real{0.20}}
  >{\raggedright\arraybackslash}p{(\columnwidth - 4\tabcolsep) * \real{0.62}}@{}}
\toprule
位置 & 类型 & 含义 \\
\midrule
\endhead
{[}0{]} \passthrough{\lstinline!status!} & \passthrough{\lstinline!int!}
& SDK 状态码；Unitree 示例通常以 \passthrough{\lstinline!0!}
表示成功，非零值需按具体服务错误码解释。 \\
{[}1{]} \passthrough{\lstinline!data!} & \passthrough{\lstinline!str!} &
传给该接口的 数据 参数，Python 类型为
\passthrough{\lstinline!str!}。精确含义、单位、范围和枚举值需查目标型号对应头文件与协议。
对应 C++ 参数 \passthrough{\lstinline!data: std::string \&!}。
底层为字符串；长度、编码和允许值由具体协议决定。 \\
\bottomrule
\end{longtable}

\textbf{对应 C++}

\begin{itemize}
\tightlist
\item
  类：\passthrough{\lstinline!unitree::robot::g1::LocoClient!}
\item
  签名：\passthrough{\lstinline!\_fsm\_api(std::string, std::string \&)!}
\item
  绑定策略：\passthrough{\lstinline!OUTPUT\_WRAPPER!}
\item
  声明位置：\passthrough{\lstinline!include/unitree/robot/g1/loco/g1\_loco\_client.hpp:303!}
\end{itemize}

\textbf{说明}

\begin{itemize}
\tightlist
\item
  示例是局部调用片段；\passthrough{\lstinline!obj!}
  和参数变量表示已经按上层业务校验并准备好的对象。
\item
  C++ 的可修改输出引用不会作为 Python 入参出现，而是按 C++
  参数顺序追加到返回元组中。
\item
  该条目可以被编辑器和类型检查器识别，但当前运行时可能在属性查找阶段直接失败。
\item
  该方法属于运动命令。方法名包含
  \passthrough{\lstinline!stop!}、\passthrough{\lstinline!damp!} 或
  \passthrough{\lstinline!zero!} 也不代表它是独立物理急停。
\end{itemize}

\textbf{用法}

\begin{lstlisting}[language=Python]
from typing import TYPE_CHECKING

if TYPE_CHECKING:
    def planned__fsm_api(obj: LocoClient, parameter: str) -> tuple[int, str]:
        return obj._fsm_api(parameter=parameter)
\end{lstlisting}

\begin{quote}
{[}!CAUTION{]} 上面的 \passthrough{\lstinline!TYPE\_CHECKING!}
示例只表达计划调用形态。该分支在普通 Python
运行时为假，不代表当前扩展能执行此接口。
\end{quote}

\hypertarget{unitree-sdk2-cpp-robot-g1-moveparameter}{%
\subsubsection{\texorpdfstring{\texttt{unitree\_sdk2\_cpp.robot.g1.MoveParameter}}{unitree\_sdk2\_cpp.robot.g1.MoveParameter}}\label{unitree-sdk2-cpp-robot-g1-moveparameter}}

SDK 请求参数值类型；当前可能仅提供设计期签名。

\textbf{导入}

\begin{lstlisting}[language=Python]
from unitree_sdk2_cpp.robot.g1 import MoveParameter
\end{lstlisting}

\textbf{基类}：\passthrough{\lstinline!object!}

\textbf{公开属性}

\begin{longtable}[]{@{}
  >{\raggedright\arraybackslash}p{(\columnwidth - 6\tabcolsep) * \real{0.18}}
  >{\raggedright\arraybackslash}p{(\columnwidth - 6\tabcolsep) * \real{0.24}}
  >{\raggedright\arraybackslash}p{(\columnwidth - 6\tabcolsep) * \real{0.18}}
  >{\raggedright\arraybackslash}p{(\columnwidth - 6\tabcolsep) * \real{0.40}}@{}}
\toprule
名称 & Python 类型 & 含义 & 用法 \\
\midrule
\endhead
\passthrough{\lstinline!vx!} & \passthrough{\lstinline!float!} & X
方向速度参数。坐标系、单位、符号和安全范围必须查目标型号运动协议。 &
\passthrough{\lstinline!value = obj.vx!} /
\passthrough{\lstinline!obj.vx = value!} \\
\passthrough{\lstinline!vy!} & \passthrough{\lstinline!float!} & Y
方向速度参数。坐标系、单位、符号和安全范围必须查目标型号运动协议。 &
\passthrough{\lstinline!value = obj.vy!} /
\passthrough{\lstinline!obj.vy = value!} \\
\passthrough{\lstinline!vyaw!} & \passthrough{\lstinline!float!} &
偏航角速度参数。单位、符号和安全范围必须查目标型号运动协议。 &
\passthrough{\lstinline!value = obj.vyaw!} /
\passthrough{\lstinline!obj.vyaw = value!} \\
\bottomrule
\end{longtable}

\textbf{方法索引}

\begin{longtable}[]{@{}
  >{\raggedright\arraybackslash}p{(\columnwidth - 6\tabcolsep) * \real{0.18}}
  >{\raggedright\arraybackslash}p{(\columnwidth - 6\tabcolsep) * \real{0.24}}
  >{\raggedright\arraybackslash}p{(\columnwidth - 6\tabcolsep) * \real{0.18}}
  >{\raggedright\arraybackslash}p{(\columnwidth - 6\tabcolsep) * \real{0.40}}@{}}
\toprule
方法 & Python 签名 & 状态 & 安全分类 \\
\midrule
\endhead
\protect\hyperlink{unitree-sdk2-cpp-robot-g1-moveparameter-dunder-init-1}{\passthrough{\lstinline!\_\_init\_\_!}}
& \passthrough{\lstinline!def \_\_init\_\_(self) -> None!} &
\passthrough{\lstinline!SIGNATURE\_ONLY!} &
\passthrough{\lstinline!CONSTRUCTION!} \\
\protect\hyperlink{unitree-sdk2-cpp-robot-g1-moveparameter-from-json-1}{\passthrough{\lstinline!from\_json!}}
&
\passthrough{\lstinline!def from\_json(self, json: dict[str, Any]) -> None!}
& \passthrough{\lstinline!SIGNATURE\_ONLY!} &
\passthrough{\lstinline!UNCLASSIFIED!} \\
\protect\hyperlink{unitree-sdk2-cpp-robot-g1-moveparameter-to-json-1}{\passthrough{\lstinline!to\_json!}}
&
\passthrough{\lstinline!def to\_json(self, json: dict[str, Any]) -> None!}
& \passthrough{\lstinline!SIGNATURE\_ONLY!} &
\passthrough{\lstinline!UNCLASSIFIED!} \\
\bottomrule
\end{longtable}

\hypertarget{unitree-sdk2-cpp-robot-g1-moveparameter-dunder-init-1}{%
\paragraph{\texorpdfstring{\texttt{unitree\_sdk2\_cpp.robot.g1.MoveParameter.\_\_init\_\_}}{unitree\_sdk2\_cpp.robot.g1.MoveParameter.\_\_init\_\_}}\label{unitree-sdk2-cpp-robot-g1-moveparameter-dunder-init-1}}

初始化 \passthrough{\lstinline!MoveParameter!}
实例。是否能实际构造取决于下方可用性状态。

\textbf{签名}

\begin{lstlisting}[language=Python]
def __init__(self) -> None
\end{lstlisting}

\textbf{可用性与安全性}

\textbf{\passthrough{\lstinline!SIGNATURE\_ONLY!} /
\passthrough{\lstinline!CONSTRUCTION!}}：当前只有设计期签名，不能假设已存在于运行时扩展。

\textbf{参数}

无显式参数。实例方法中的 \passthrough{\lstinline!self!} 由 Python
自动传入。

\textbf{返回值}

\begin{longtable}[]{@{}
  >{\raggedright\arraybackslash}p{(\columnwidth - 4\tabcolsep) * \real{0.18}}
  >{\raggedright\arraybackslash}p{(\columnwidth - 4\tabcolsep) * \real{0.20}}
  >{\raggedright\arraybackslash}p{(\columnwidth - 4\tabcolsep) * \real{0.62}}@{}}
\toprule
位置 & 类型 & 含义 \\
\midrule
\endhead
无 & \passthrough{\lstinline!None!} &
\passthrough{\lstinline!\_\_init\_\_!}
本身不返回值；调用类对象时，在构造函数可用的前提下得到该类实例。 \\
\bottomrule
\end{longtable}

\textbf{对应 C++}

\begin{itemize}
\tightlist
\item
  类：\passthrough{\lstinline!unitree::robot::g1::MoveParameter!}
\item
  签名：\passthrough{\lstinline!MoveParameter()!}
\item
  绑定策略：\passthrough{\lstinline!未单独标注!}
\item
  声明位置：\passthrough{\lstinline!include/unitree/robot/g1/agv/g1\_agv\_api.hpp:28!}
\end{itemize}

\textbf{说明}

\begin{itemize}
\tightlist
\item
  示例是局部调用片段；\passthrough{\lstinline!obj!}
  和参数变量表示已经按上层业务校验并准备好的对象。
\item
  该条目可以被编辑器和类型检查器识别，但当前运行时可能在属性查找阶段直接失败。
\end{itemize}

\textbf{用法}

\begin{lstlisting}[language=Python]
from typing import TYPE_CHECKING

if TYPE_CHECKING:
    def planned_construct() -> MoveParameter:
        return MoveParameter()
\end{lstlisting}

\begin{quote}
{[}!CAUTION{]} 上面的 \passthrough{\lstinline!TYPE\_CHECKING!}
示例只表达计划调用形态。该分支在普通 Python
运行时为假，不代表当前扩展能执行此接口。
\end{quote}

\hypertarget{unitree-sdk2-cpp-robot-g1-moveparameter-from-json-1}{%
\paragraph{\texorpdfstring{\texttt{unitree\_sdk2\_cpp.robot.g1.MoveParameter.from\_json}}{unitree\_sdk2\_cpp.robot.g1.MoveParameter.from\_json}}\label{unitree-sdk2-cpp-robot-g1-moveparameter-from-json-1}}

计划从 JSON 风格字典读取字段并更新当前 SDK 值对象。

\textbf{签名}

\begin{lstlisting}[language=Python]
def from_json(self, json: dict[str, Any]) -> None
\end{lstlisting}

\textbf{可用性与安全性}

\textbf{\passthrough{\lstinline!SIGNATURE\_ONLY!} /
\passthrough{\lstinline!UNCLASSIFIED!}}：当前只有设计期签名，不能假设已存在于运行时扩展。

\textbf{参数}

\begin{longtable}[]{@{}
  >{\raggedright\arraybackslash}p{(\columnwidth - 6\tabcolsep) * \real{0.18}}
  >{\raggedright\arraybackslash}p{(\columnwidth - 6\tabcolsep) * \real{0.24}}
  >{\raggedright\arraybackslash}p{(\columnwidth - 6\tabcolsep) * \real{0.18}}
  >{\raggedright\arraybackslash}p{(\columnwidth - 6\tabcolsep) * \real{0.40}}@{}}
\toprule
名称 & Python 类型 & 默认值 & 含义 \\
\midrule
\endhead
\passthrough{\lstinline!json!} &
\passthrough{\lstinline!dict[str, Any]!} & 必填 & JSON
风格字典，键和值必须满足对应 C++ \passthrough{\lstinline!JsonMap!}
数据结构的约束。 对应 C++ 参数
\passthrough{\lstinline!json: unitree::common::JsonMap \&!}。 \\
\bottomrule
\end{longtable}

\textbf{返回值}

\begin{longtable}[]{@{}
  >{\raggedright\arraybackslash}p{(\columnwidth - 4\tabcolsep) * \real{0.18}}
  >{\raggedright\arraybackslash}p{(\columnwidth - 4\tabcolsep) * \real{0.20}}
  >{\raggedright\arraybackslash}p{(\columnwidth - 4\tabcolsep) * \real{0.62}}@{}}
\toprule
位置 & 类型 & 含义 \\
\midrule
\endhead
无 & \passthrough{\lstinline!None!} & 该方法不返回 Python
值；失败可能表现为异常或底层状态变化。 \\
\bottomrule
\end{longtable}

\textbf{对应 C++}

\begin{itemize}
\tightlist
\item
  类：\passthrough{\lstinline!unitree::robot::g1::MoveParameter!}
\item
  签名：\passthrough{\lstinline!fromJson(unitree::common::JsonMap \&)!}
\item
  绑定策略：\passthrough{\lstinline!SIGNATURE\_PREVIEW!}
\item
  声明位置：\passthrough{\lstinline!include/unitree/robot/g1/agv/g1\_agv\_api.hpp:36!}
\end{itemize}

\textbf{说明}

\begin{itemize}
\tightlist
\item
  示例是局部调用片段；\passthrough{\lstinline!obj!}
  和参数变量表示已经按上层业务校验并准备好的对象。
\item
  该条目可以被编辑器和类型检查器识别，但当前运行时可能在属性查找阶段直接失败。
\end{itemize}

\textbf{用法}

\begin{lstlisting}[language=Python]
from typing import TYPE_CHECKING

if TYPE_CHECKING:
    def planned_from_json(obj: MoveParameter, json: dict[str, Any]) -> None:
        obj.from_json(json=json)
\end{lstlisting}

\begin{quote}
{[}!CAUTION{]} 上面的 \passthrough{\lstinline!TYPE\_CHECKING!}
示例只表达计划调用形态。该分支在普通 Python
运行时为假，不代表当前扩展能执行此接口。
\end{quote}

\hypertarget{unitree-sdk2-cpp-robot-g1-moveparameter-to-json-1}{%
\paragraph{\texorpdfstring{\texttt{unitree\_sdk2\_cpp.robot.g1.MoveParameter.to\_json}}{unitree\_sdk2\_cpp.robot.g1.MoveParameter.to\_json}}\label{unitree-sdk2-cpp-robot-g1-moveparameter-to-json-1}}

计划把当前 SDK 值对象写入 JSON 风格字典。

\textbf{签名}

\begin{lstlisting}[language=Python]
def to_json(self, json: dict[str, Any]) -> None
\end{lstlisting}

\textbf{可用性与安全性}

\textbf{\passthrough{\lstinline!SIGNATURE\_ONLY!} /
\passthrough{\lstinline!UNCLASSIFIED!}}：当前只有设计期签名，不能假设已存在于运行时扩展。

\textbf{参数}

\begin{longtable}[]{@{}
  >{\raggedright\arraybackslash}p{(\columnwidth - 6\tabcolsep) * \real{0.18}}
  >{\raggedright\arraybackslash}p{(\columnwidth - 6\tabcolsep) * \real{0.24}}
  >{\raggedright\arraybackslash}p{(\columnwidth - 6\tabcolsep) * \real{0.18}}
  >{\raggedright\arraybackslash}p{(\columnwidth - 6\tabcolsep) * \real{0.40}}@{}}
\toprule
名称 & Python 类型 & 默认值 & 含义 \\
\midrule
\endhead
\passthrough{\lstinline!json!} &
\passthrough{\lstinline!dict[str, Any]!} & 必填 & JSON
风格字典，键和值必须满足对应 C++ \passthrough{\lstinline!JsonMap!}
数据结构的约束。 对应 C++ 参数
\passthrough{\lstinline!json: unitree::common::JsonMap \&!}。 \\
\bottomrule
\end{longtable}

\textbf{返回值}

\begin{longtable}[]{@{}
  >{\raggedright\arraybackslash}p{(\columnwidth - 4\tabcolsep) * \real{0.18}}
  >{\raggedright\arraybackslash}p{(\columnwidth - 4\tabcolsep) * \real{0.20}}
  >{\raggedright\arraybackslash}p{(\columnwidth - 4\tabcolsep) * \real{0.62}}@{}}
\toprule
位置 & 类型 & 含义 \\
\midrule
\endhead
无 & \passthrough{\lstinline!None!} & 该方法不返回 Python
值；失败可能表现为异常或底层状态变化。 \\
\bottomrule
\end{longtable}

\textbf{对应 C++}

\begin{itemize}
\tightlist
\item
  类：\passthrough{\lstinline!unitree::robot::g1::MoveParameter!}
\item
  签名：\passthrough{\lstinline!toJson(unitree::common::JsonMap \&) const!}
\item
  绑定策略：\passthrough{\lstinline!SIGNATURE\_PREVIEW!}
\item
  声明位置：\passthrough{\lstinline!include/unitree/robot/g1/agv/g1\_agv\_api.hpp:43!}
\end{itemize}

\textbf{说明}

\begin{itemize}
\tightlist
\item
  示例是局部调用片段；\passthrough{\lstinline!obj!}
  和参数变量表示已经按上层业务校验并准备好的对象。
\item
  该条目可以被编辑器和类型检查器识别，但当前运行时可能在属性查找阶段直接失败。
\end{itemize}

\textbf{用法}

\begin{lstlisting}[language=Python]
from typing import TYPE_CHECKING

if TYPE_CHECKING:
    def planned_to_json(obj: MoveParameter, json: dict[str, Any]) -> None:
        obj.to_json(json=json)
\end{lstlisting}

\begin{quote}
{[}!CAUTION{]} 上面的 \passthrough{\lstinline!TYPE\_CHECKING!}
示例只表达计划调用形态。该分支在普通 Python
运行时为假，不代表当前扩展能执行此接口。
\end{quote}

\hypertarget{unitree-sdk2-cpp-robot-g1-playstopparameter}{%
\subsubsection{\texorpdfstring{\texttt{unitree\_sdk2\_cpp.robot.g1.PlayStopParameter}}{unitree\_sdk2\_cpp.robot.g1.PlayStopParameter}}\label{unitree-sdk2-cpp-robot-g1-playstopparameter}}

SDK 请求参数值类型；当前可能仅提供设计期签名。

\textbf{导入}

\begin{lstlisting}[language=Python]
from unitree_sdk2_cpp.robot.g1 import PlayStopParameter
\end{lstlisting}

\textbf{基类}：\passthrough{\lstinline!object!}

\textbf{公开属性}

\begin{longtable}[]{@{}
  >{\raggedright\arraybackslash}p{(\columnwidth - 6\tabcolsep) * \real{0.18}}
  >{\raggedright\arraybackslash}p{(\columnwidth - 6\tabcolsep) * \real{0.24}}
  >{\raggedright\arraybackslash}p{(\columnwidth - 6\tabcolsep) * \real{0.18}}
  >{\raggedright\arraybackslash}p{(\columnwidth - 6\tabcolsep) * \real{0.40}}@{}}
\toprule
名称 & Python 类型 & 含义 & 用法 \\
\midrule
\endhead
\passthrough{\lstinline!app\_name!} & \passthrough{\lstinline!str!} &
应用名称，用于标识音频流或播放会话。允许值和命名规则以目标服务协议为准。
& \passthrough{\lstinline!value = obj.app\_name!} /
\passthrough{\lstinline!obj.app\_name = value!} \\
\bottomrule
\end{longtable}

\textbf{方法索引}

\begin{longtable}[]{@{}
  >{\raggedright\arraybackslash}p{(\columnwidth - 6\tabcolsep) * \real{0.18}}
  >{\raggedright\arraybackslash}p{(\columnwidth - 6\tabcolsep) * \real{0.24}}
  >{\raggedright\arraybackslash}p{(\columnwidth - 6\tabcolsep) * \real{0.18}}
  >{\raggedright\arraybackslash}p{(\columnwidth - 6\tabcolsep) * \real{0.40}}@{}}
\toprule
方法 & Python 签名 & 状态 & 安全分类 \\
\midrule
\endhead
\protect\hyperlink{unitree-sdk2-cpp-robot-g1-playstopparameter-dunder-init-1}{\passthrough{\lstinline!\_\_init\_\_!}}
& \passthrough{\lstinline!def \_\_init\_\_(self) -> None!} &
\passthrough{\lstinline!SIGNATURE\_ONLY!} &
\passthrough{\lstinline!CONSTRUCTION!} \\
\protect\hyperlink{unitree-sdk2-cpp-robot-g1-playstopparameter-from-json-1}{\passthrough{\lstinline!from\_json!}}
&
\passthrough{\lstinline!def from\_json(self, json: dict[str, Any]) -> None!}
& \passthrough{\lstinline!SIGNATURE\_ONLY!} &
\passthrough{\lstinline!UNCLASSIFIED!} \\
\protect\hyperlink{unitree-sdk2-cpp-robot-g1-playstopparameter-to-json-1}{\passthrough{\lstinline!to\_json!}}
&
\passthrough{\lstinline!def to\_json(self, json: dict[str, Any]) -> None!}
& \passthrough{\lstinline!SIGNATURE\_ONLY!} &
\passthrough{\lstinline!UNCLASSIFIED!} \\
\bottomrule
\end{longtable}

\hypertarget{unitree-sdk2-cpp-robot-g1-playstopparameter-dunder-init-1}{%
\paragraph{\texorpdfstring{\texttt{unitree\_sdk2\_cpp.robot.g1.PlayStopParameter.\_\_init\_\_}}{unitree\_sdk2\_cpp.robot.g1.PlayStopParameter.\_\_init\_\_}}\label{unitree-sdk2-cpp-robot-g1-playstopparameter-dunder-init-1}}

初始化 \passthrough{\lstinline!PlayStopParameter!}
实例。是否能实际构造取决于下方可用性状态。

\textbf{签名}

\begin{lstlisting}[language=Python]
def __init__(self) -> None
\end{lstlisting}

\textbf{可用性与安全性}

\textbf{\passthrough{\lstinline!SIGNATURE\_ONLY!} /
\passthrough{\lstinline!CONSTRUCTION!}}：当前只有设计期签名，不能假设已存在于运行时扩展。

\textbf{参数}

无显式参数。实例方法中的 \passthrough{\lstinline!self!} 由 Python
自动传入。

\textbf{返回值}

\begin{longtable}[]{@{}
  >{\raggedright\arraybackslash}p{(\columnwidth - 4\tabcolsep) * \real{0.18}}
  >{\raggedright\arraybackslash}p{(\columnwidth - 4\tabcolsep) * \real{0.20}}
  >{\raggedright\arraybackslash}p{(\columnwidth - 4\tabcolsep) * \real{0.62}}@{}}
\toprule
位置 & 类型 & 含义 \\
\midrule
\endhead
无 & \passthrough{\lstinline!None!} &
\passthrough{\lstinline!\_\_init\_\_!}
本身不返回值；调用类对象时，在构造函数可用的前提下得到该类实例。 \\
\bottomrule
\end{longtable}

\textbf{对应 C++}

\begin{itemize}
\tightlist
\item
  类：\passthrough{\lstinline!unitree::robot::g1::PlayStopParameter!}
\item
  签名：\passthrough{\lstinline!PlayStopParameter()!}
\item
  绑定策略：\passthrough{\lstinline!未单独标注!}
\item
  声明位置：\passthrough{\lstinline!include/unitree/robot/g1/audio/g1\_audio\_api.hpp:61!}
\end{itemize}

\textbf{说明}

\begin{itemize}
\tightlist
\item
  示例是局部调用片段；\passthrough{\lstinline!obj!}
  和参数变量表示已经按上层业务校验并准备好的对象。
\item
  该条目可以被编辑器和类型检查器识别，但当前运行时可能在属性查找阶段直接失败。
\end{itemize}

\textbf{用法}

\begin{lstlisting}[language=Python]
from typing import TYPE_CHECKING

if TYPE_CHECKING:
    def planned_construct() -> PlayStopParameter:
        return PlayStopParameter()
\end{lstlisting}

\begin{quote}
{[}!CAUTION{]} 上面的 \passthrough{\lstinline!TYPE\_CHECKING!}
示例只表达计划调用形态。该分支在普通 Python
运行时为假，不代表当前扩展能执行此接口。
\end{quote}

\hypertarget{unitree-sdk2-cpp-robot-g1-playstopparameter-from-json-1}{%
\paragraph{\texorpdfstring{\texttt{unitree\_sdk2\_cpp.robot.g1.PlayStopParameter.from\_json}}{unitree\_sdk2\_cpp.robot.g1.PlayStopParameter.from\_json}}\label{unitree-sdk2-cpp-robot-g1-playstopparameter-from-json-1}}

计划从 JSON 风格字典读取字段并更新当前 SDK 值对象。

\textbf{签名}

\begin{lstlisting}[language=Python]
def from_json(self, json: dict[str, Any]) -> None
\end{lstlisting}

\textbf{可用性与安全性}

\textbf{\passthrough{\lstinline!SIGNATURE\_ONLY!} /
\passthrough{\lstinline!UNCLASSIFIED!}}：当前只有设计期签名，不能假设已存在于运行时扩展。

\textbf{参数}

\begin{longtable}[]{@{}
  >{\raggedright\arraybackslash}p{(\columnwidth - 6\tabcolsep) * \real{0.18}}
  >{\raggedright\arraybackslash}p{(\columnwidth - 6\tabcolsep) * \real{0.24}}
  >{\raggedright\arraybackslash}p{(\columnwidth - 6\tabcolsep) * \real{0.18}}
  >{\raggedright\arraybackslash}p{(\columnwidth - 6\tabcolsep) * \real{0.40}}@{}}
\toprule
名称 & Python 类型 & 默认值 & 含义 \\
\midrule
\endhead
\passthrough{\lstinline!json!} &
\passthrough{\lstinline!dict[str, Any]!} & 必填 & JSON
风格字典，键和值必须满足对应 C++ \passthrough{\lstinline!JsonMap!}
数据结构的约束。 对应 C++ 参数
\passthrough{\lstinline!json: common::JsonMap \&!}。 \\
\bottomrule
\end{longtable}

\textbf{返回值}

\begin{longtable}[]{@{}
  >{\raggedright\arraybackslash}p{(\columnwidth - 4\tabcolsep) * \real{0.18}}
  >{\raggedright\arraybackslash}p{(\columnwidth - 4\tabcolsep) * \real{0.20}}
  >{\raggedright\arraybackslash}p{(\columnwidth - 4\tabcolsep) * \real{0.62}}@{}}
\toprule
位置 & 类型 & 含义 \\
\midrule
\endhead
无 & \passthrough{\lstinline!None!} & 该方法不返回 Python
值；失败可能表现为异常或底层状态变化。 \\
\bottomrule
\end{longtable}

\textbf{对应 C++}

\begin{itemize}
\tightlist
\item
  类：\passthrough{\lstinline!unitree::robot::g1::PlayStopParameter!}
\item
  签名：\passthrough{\lstinline!fromJson(common::JsonMap \&)!}
\item
  绑定策略：\passthrough{\lstinline!SIGNATURE\_PREVIEW!}
\item
  声明位置：\passthrough{\lstinline!include/unitree/robot/g1/audio/g1\_audio\_api.hpp:64!}
\end{itemize}

\textbf{说明}

\begin{itemize}
\tightlist
\item
  示例是局部调用片段；\passthrough{\lstinline!obj!}
  和参数变量表示已经按上层业务校验并准备好的对象。
\item
  该条目可以被编辑器和类型检查器识别，但当前运行时可能在属性查找阶段直接失败。
\end{itemize}

\textbf{用法}

\begin{lstlisting}[language=Python]
from typing import TYPE_CHECKING

if TYPE_CHECKING:
    def planned_from_json(obj: PlayStopParameter, json: dict[str, Any]) -> None:
        obj.from_json(json=json)
\end{lstlisting}

\begin{quote}
{[}!CAUTION{]} 上面的 \passthrough{\lstinline!TYPE\_CHECKING!}
示例只表达计划调用形态。该分支在普通 Python
运行时为假，不代表当前扩展能执行此接口。
\end{quote}

\hypertarget{unitree-sdk2-cpp-robot-g1-playstopparameter-to-json-1}{%
\paragraph{\texorpdfstring{\texttt{unitree\_sdk2\_cpp.robot.g1.PlayStopParameter.to\_json}}{unitree\_sdk2\_cpp.robot.g1.PlayStopParameter.to\_json}}\label{unitree-sdk2-cpp-robot-g1-playstopparameter-to-json-1}}

计划把当前 SDK 值对象写入 JSON 风格字典。

\textbf{签名}

\begin{lstlisting}[language=Python]
def to_json(self, json: dict[str, Any]) -> None
\end{lstlisting}

\textbf{可用性与安全性}

\textbf{\passthrough{\lstinline!SIGNATURE\_ONLY!} /
\passthrough{\lstinline!UNCLASSIFIED!}}：当前只有设计期签名，不能假设已存在于运行时扩展。

\textbf{参数}

\begin{longtable}[]{@{}
  >{\raggedright\arraybackslash}p{(\columnwidth - 6\tabcolsep) * \real{0.18}}
  >{\raggedright\arraybackslash}p{(\columnwidth - 6\tabcolsep) * \real{0.24}}
  >{\raggedright\arraybackslash}p{(\columnwidth - 6\tabcolsep) * \real{0.18}}
  >{\raggedright\arraybackslash}p{(\columnwidth - 6\tabcolsep) * \real{0.40}}@{}}
\toprule
名称 & Python 类型 & 默认值 & 含义 \\
\midrule
\endhead
\passthrough{\lstinline!json!} &
\passthrough{\lstinline!dict[str, Any]!} & 必填 & JSON
风格字典，键和值必须满足对应 C++ \passthrough{\lstinline!JsonMap!}
数据结构的约束。 对应 C++ 参数
\passthrough{\lstinline!json: common::JsonMap \&!}。 \\
\bottomrule
\end{longtable}

\textbf{返回值}

\begin{longtable}[]{@{}
  >{\raggedright\arraybackslash}p{(\columnwidth - 4\tabcolsep) * \real{0.18}}
  >{\raggedright\arraybackslash}p{(\columnwidth - 4\tabcolsep) * \real{0.20}}
  >{\raggedright\arraybackslash}p{(\columnwidth - 4\tabcolsep) * \real{0.62}}@{}}
\toprule
位置 & 类型 & 含义 \\
\midrule
\endhead
无 & \passthrough{\lstinline!None!} & 该方法不返回 Python
值；失败可能表现为异常或底层状态变化。 \\
\bottomrule
\end{longtable}

\textbf{对应 C++}

\begin{itemize}
\tightlist
\item
  类：\passthrough{\lstinline!unitree::robot::g1::PlayStopParameter!}
\item
  签名：\passthrough{\lstinline!toJson(common::JsonMap \&) const!}
\item
  绑定策略：\passthrough{\lstinline!SIGNATURE\_PREVIEW!}
\item
  声明位置：\passthrough{\lstinline!include/unitree/robot/g1/audio/g1\_audio\_api.hpp:66!}
\end{itemize}

\textbf{说明}

\begin{itemize}
\tightlist
\item
  示例是局部调用片段；\passthrough{\lstinline!obj!}
  和参数变量表示已经按上层业务校验并准备好的对象。
\item
  该条目可以被编辑器和类型检查器识别，但当前运行时可能在属性查找阶段直接失败。
\end{itemize}

\textbf{用法}

\begin{lstlisting}[language=Python]
from typing import TYPE_CHECKING

if TYPE_CHECKING:
    def planned_to_json(obj: PlayStopParameter, json: dict[str, Any]) -> None:
        obj.to_json(json=json)
\end{lstlisting}

\begin{quote}
{[}!CAUTION{]} 上面的 \passthrough{\lstinline!TYPE\_CHECKING!}
示例只表达计划调用形态。该分支在普通 Python
运行时为假，不代表当前扩展能执行此接口。
\end{quote}

\hypertarget{unitree-sdk2-cpp-robot-g1-playstreamparameter}{%
\subsubsection{\texorpdfstring{\texttt{unitree\_sdk2\_cpp.robot.g1.PlayStreamParameter}}{unitree\_sdk2\_cpp.robot.g1.PlayStreamParameter}}\label{unitree-sdk2-cpp-robot-g1-playstreamparameter}}

SDK 请求参数值类型；当前可能仅提供设计期签名。

\textbf{导入}

\begin{lstlisting}[language=Python]
from unitree_sdk2_cpp.robot.g1 import PlayStreamParameter
\end{lstlisting}

\textbf{基类}：\passthrough{\lstinline!object!}

\textbf{公开属性}

\begin{longtable}[]{@{}
  >{\raggedright\arraybackslash}p{(\columnwidth - 6\tabcolsep) * \real{0.18}}
  >{\raggedright\arraybackslash}p{(\columnwidth - 6\tabcolsep) * \real{0.24}}
  >{\raggedright\arraybackslash}p{(\columnwidth - 6\tabcolsep) * \real{0.18}}
  >{\raggedright\arraybackslash}p{(\columnwidth - 6\tabcolsep) * \real{0.40}}@{}}
\toprule
名称 & Python 类型 & 含义 & 用法 \\
\midrule
\endhead
\passthrough{\lstinline!app\_name!} & \passthrough{\lstinline!str!} &
应用名称，用于标识音频流或播放会话。允许值和命名规则以目标服务协议为准。
& \passthrough{\lstinline!value = obj.app\_name!} /
\passthrough{\lstinline!obj.app\_name = value!} \\
\passthrough{\lstinline!stream\_id!} & \passthrough{\lstinline!str!} &
音频流标识符，用于区分同一应用下的播放流。 &
\passthrough{\lstinline!value = obj.stream\_id!} /
\passthrough{\lstinline!obj.stream\_id = value!} \\
\bottomrule
\end{longtable}

\textbf{方法索引}

\begin{longtable}[]{@{}
  >{\raggedright\arraybackslash}p{(\columnwidth - 6\tabcolsep) * \real{0.18}}
  >{\raggedright\arraybackslash}p{(\columnwidth - 6\tabcolsep) * \real{0.24}}
  >{\raggedright\arraybackslash}p{(\columnwidth - 6\tabcolsep) * \real{0.18}}
  >{\raggedright\arraybackslash}p{(\columnwidth - 6\tabcolsep) * \real{0.40}}@{}}
\toprule
方法 & Python 签名 & 状态 & 安全分类 \\
\midrule
\endhead
\protect\hyperlink{unitree-sdk2-cpp-robot-g1-playstreamparameter-dunder-init-1}{\passthrough{\lstinline!\_\_init\_\_!}}
& \passthrough{\lstinline!def \_\_init\_\_(self) -> None!} &
\passthrough{\lstinline!SIGNATURE\_ONLY!} &
\passthrough{\lstinline!CONSTRUCTION!} \\
\protect\hyperlink{unitree-sdk2-cpp-robot-g1-playstreamparameter-from-json-1}{\passthrough{\lstinline!from\_json!}}
&
\passthrough{\lstinline!def from\_json(self, json: dict[str, Any]) -> None!}
& \passthrough{\lstinline!SIGNATURE\_ONLY!} &
\passthrough{\lstinline!UNCLASSIFIED!} \\
\protect\hyperlink{unitree-sdk2-cpp-robot-g1-playstreamparameter-to-json-1}{\passthrough{\lstinline!to\_json!}}
&
\passthrough{\lstinline!def to\_json(self, json: dict[str, Any]) -> None!}
& \passthrough{\lstinline!SIGNATURE\_ONLY!} &
\passthrough{\lstinline!UNCLASSIFIED!} \\
\bottomrule
\end{longtable}

\hypertarget{unitree-sdk2-cpp-robot-g1-playstreamparameter-dunder-init-1}{%
\paragraph{\texorpdfstring{\texttt{unitree\_sdk2\_cpp.robot.g1.PlayStreamParameter.\_\_init\_\_}}{unitree\_sdk2\_cpp.robot.g1.PlayStreamParameter.\_\_init\_\_}}\label{unitree-sdk2-cpp-robot-g1-playstreamparameter-dunder-init-1}}

初始化 \passthrough{\lstinline!PlayStreamParameter!}
实例。是否能实际构造取决于下方可用性状态。

\textbf{签名}

\begin{lstlisting}[language=Python]
def __init__(self) -> None
\end{lstlisting}

\textbf{可用性与安全性}

\textbf{\passthrough{\lstinline!SIGNATURE\_ONLY!} /
\passthrough{\lstinline!CONSTRUCTION!}}：当前只有设计期签名，不能假设已存在于运行时扩展。

\textbf{参数}

无显式参数。实例方法中的 \passthrough{\lstinline!self!} 由 Python
自动传入。

\textbf{返回值}

\begin{longtable}[]{@{}
  >{\raggedright\arraybackslash}p{(\columnwidth - 4\tabcolsep) * \real{0.18}}
  >{\raggedright\arraybackslash}p{(\columnwidth - 4\tabcolsep) * \real{0.20}}
  >{\raggedright\arraybackslash}p{(\columnwidth - 4\tabcolsep) * \real{0.62}}@{}}
\toprule
位置 & 类型 & 含义 \\
\midrule
\endhead
无 & \passthrough{\lstinline!None!} &
\passthrough{\lstinline!\_\_init\_\_!}
本身不返回值；调用类对象时，在构造函数可用的前提下得到该类实例。 \\
\bottomrule
\end{longtable}

\textbf{对应 C++}

\begin{itemize}
\tightlist
\item
  类：\passthrough{\lstinline!unitree::robot::g1::PlayStreamParameter!}
\item
  签名：\passthrough{\lstinline!PlayStreamParameter()!}
\item
  绑定策略：\passthrough{\lstinline!未单独标注!}
\item
  声明位置：\passthrough{\lstinline!include/unitree/robot/g1/audio/g1\_audio\_api.hpp:45!}
\end{itemize}

\textbf{说明}

\begin{itemize}
\tightlist
\item
  示例是局部调用片段；\passthrough{\lstinline!obj!}
  和参数变量表示已经按上层业务校验并准备好的对象。
\item
  该条目可以被编辑器和类型检查器识别，但当前运行时可能在属性查找阶段直接失败。
\end{itemize}

\textbf{用法}

\begin{lstlisting}[language=Python]
from typing import TYPE_CHECKING

if TYPE_CHECKING:
    def planned_construct() -> PlayStreamParameter:
        return PlayStreamParameter()
\end{lstlisting}

\begin{quote}
{[}!CAUTION{]} 上面的 \passthrough{\lstinline!TYPE\_CHECKING!}
示例只表达计划调用形态。该分支在普通 Python
运行时为假，不代表当前扩展能执行此接口。
\end{quote}

\hypertarget{unitree-sdk2-cpp-robot-g1-playstreamparameter-from-json-1}{%
\paragraph{\texorpdfstring{\texttt{unitree\_sdk2\_cpp.robot.g1.PlayStreamParameter.from\_json}}{unitree\_sdk2\_cpp.robot.g1.PlayStreamParameter.from\_json}}\label{unitree-sdk2-cpp-robot-g1-playstreamparameter-from-json-1}}

计划从 JSON 风格字典读取字段并更新当前 SDK 值对象。

\textbf{签名}

\begin{lstlisting}[language=Python]
def from_json(self, json: dict[str, Any]) -> None
\end{lstlisting}

\textbf{可用性与安全性}

\textbf{\passthrough{\lstinline!SIGNATURE\_ONLY!} /
\passthrough{\lstinline!UNCLASSIFIED!}}：当前只有设计期签名，不能假设已存在于运行时扩展。

\textbf{参数}

\begin{longtable}[]{@{}
  >{\raggedright\arraybackslash}p{(\columnwidth - 6\tabcolsep) * \real{0.18}}
  >{\raggedright\arraybackslash}p{(\columnwidth - 6\tabcolsep) * \real{0.24}}
  >{\raggedright\arraybackslash}p{(\columnwidth - 6\tabcolsep) * \real{0.18}}
  >{\raggedright\arraybackslash}p{(\columnwidth - 6\tabcolsep) * \real{0.40}}@{}}
\toprule
名称 & Python 类型 & 默认值 & 含义 \\
\midrule
\endhead
\passthrough{\lstinline!json!} &
\passthrough{\lstinline!dict[str, Any]!} & 必填 & JSON
风格字典，键和值必须满足对应 C++ \passthrough{\lstinline!JsonMap!}
数据结构的约束。 对应 C++ 参数
\passthrough{\lstinline!json: common::JsonMap \&!}。 \\
\bottomrule
\end{longtable}

\textbf{返回值}

\begin{longtable}[]{@{}
  >{\raggedright\arraybackslash}p{(\columnwidth - 4\tabcolsep) * \real{0.18}}
  >{\raggedright\arraybackslash}p{(\columnwidth - 4\tabcolsep) * \real{0.20}}
  >{\raggedright\arraybackslash}p{(\columnwidth - 4\tabcolsep) * \real{0.62}}@{}}
\toprule
位置 & 类型 & 含义 \\
\midrule
\endhead
无 & \passthrough{\lstinline!None!} & 该方法不返回 Python
值；失败可能表现为异常或底层状态变化。 \\
\bottomrule
\end{longtable}

\textbf{对应 C++}

\begin{itemize}
\tightlist
\item
  类：\passthrough{\lstinline!unitree::robot::g1::PlayStreamParameter!}
\item
  签名：\passthrough{\lstinline!fromJson(common::JsonMap \&)!}
\item
  绑定策略：\passthrough{\lstinline!SIGNATURE\_PREVIEW!}
\item
  声明位置：\passthrough{\lstinline!include/unitree/robot/g1/audio/g1\_audio\_api.hpp:48!}
\end{itemize}

\textbf{说明}

\begin{itemize}
\tightlist
\item
  示例是局部调用片段；\passthrough{\lstinline!obj!}
  和参数变量表示已经按上层业务校验并准备好的对象。
\item
  该条目可以被编辑器和类型检查器识别，但当前运行时可能在属性查找阶段直接失败。
\end{itemize}

\textbf{用法}

\begin{lstlisting}[language=Python]
from typing import TYPE_CHECKING

if TYPE_CHECKING:
    def planned_from_json(obj: PlayStreamParameter, json: dict[str, Any]) -> None:
        obj.from_json(json=json)
\end{lstlisting}

\begin{quote}
{[}!CAUTION{]} 上面的 \passthrough{\lstinline!TYPE\_CHECKING!}
示例只表达计划调用形态。该分支在普通 Python
运行时为假，不代表当前扩展能执行此接口。
\end{quote}

\hypertarget{unitree-sdk2-cpp-robot-g1-playstreamparameter-to-json-1}{%
\paragraph{\texorpdfstring{\texttt{unitree\_sdk2\_cpp.robot.g1.PlayStreamParameter.to\_json}}{unitree\_sdk2\_cpp.robot.g1.PlayStreamParameter.to\_json}}\label{unitree-sdk2-cpp-robot-g1-playstreamparameter-to-json-1}}

计划把当前 SDK 值对象写入 JSON 风格字典。

\textbf{签名}

\begin{lstlisting}[language=Python]
def to_json(self, json: dict[str, Any]) -> None
\end{lstlisting}

\textbf{可用性与安全性}

\textbf{\passthrough{\lstinline!SIGNATURE\_ONLY!} /
\passthrough{\lstinline!UNCLASSIFIED!}}：当前只有设计期签名，不能假设已存在于运行时扩展。

\textbf{参数}

\begin{longtable}[]{@{}
  >{\raggedright\arraybackslash}p{(\columnwidth - 6\tabcolsep) * \real{0.18}}
  >{\raggedright\arraybackslash}p{(\columnwidth - 6\tabcolsep) * \real{0.24}}
  >{\raggedright\arraybackslash}p{(\columnwidth - 6\tabcolsep) * \real{0.18}}
  >{\raggedright\arraybackslash}p{(\columnwidth - 6\tabcolsep) * \real{0.40}}@{}}
\toprule
名称 & Python 类型 & 默认值 & 含义 \\
\midrule
\endhead
\passthrough{\lstinline!json!} &
\passthrough{\lstinline!dict[str, Any]!} & 必填 & JSON
风格字典，键和值必须满足对应 C++ \passthrough{\lstinline!JsonMap!}
数据结构的约束。 对应 C++ 参数
\passthrough{\lstinline!json: common::JsonMap \&!}。 \\
\bottomrule
\end{longtable}

\textbf{返回值}

\begin{longtable}[]{@{}
  >{\raggedright\arraybackslash}p{(\columnwidth - 4\tabcolsep) * \real{0.18}}
  >{\raggedright\arraybackslash}p{(\columnwidth - 4\tabcolsep) * \real{0.20}}
  >{\raggedright\arraybackslash}p{(\columnwidth - 4\tabcolsep) * \real{0.62}}@{}}
\toprule
位置 & 类型 & 含义 \\
\midrule
\endhead
无 & \passthrough{\lstinline!None!} & 该方法不返回 Python
值；失败可能表现为异常或底层状态变化。 \\
\bottomrule
\end{longtable}

\textbf{对应 C++}

\begin{itemize}
\tightlist
\item
  类：\passthrough{\lstinline!unitree::robot::g1::PlayStreamParameter!}
\item
  签名：\passthrough{\lstinline!toJson(common::JsonMap \&) const!}
\item
  绑定策略：\passthrough{\lstinline!SIGNATURE\_PREVIEW!}
\item
  声明位置：\passthrough{\lstinline!include/unitree/robot/g1/audio/g1\_audio\_api.hpp:50!}
\end{itemize}

\textbf{说明}

\begin{itemize}
\tightlist
\item
  示例是局部调用片段；\passthrough{\lstinline!obj!}
  和参数变量表示已经按上层业务校验并准备好的对象。
\item
  该条目可以被编辑器和类型检查器识别，但当前运行时可能在属性查找阶段直接失败。
\end{itemize}

\textbf{用法}

\begin{lstlisting}[language=Python]
from typing import TYPE_CHECKING

if TYPE_CHECKING:
    def planned_to_json(obj: PlayStreamParameter, json: dict[str, Any]) -> None:
        obj.to_json(json=json)
\end{lstlisting}

\begin{quote}
{[}!CAUTION{]} 上面的 \passthrough{\lstinline!TYPE\_CHECKING!}
示例只表达计划调用形态。该分支在普通 Python
运行时为假，不代表当前扩展能执行此接口。
\end{quote}

\hypertarget{unitree-sdk2-cpp-robot-g1-ttsmakerparameter}{%
\subsubsection{\texorpdfstring{\texttt{unitree\_sdk2\_cpp.robot.g1.TtsMakerParameter}}{unitree\_sdk2\_cpp.robot.g1.TtsMakerParameter}}\label{unitree-sdk2-cpp-robot-g1-ttsmakerparameter}}

SDK 请求参数值类型；当前可能仅提供设计期签名。

\textbf{导入}

\begin{lstlisting}[language=Python]
from unitree_sdk2_cpp.robot.g1 import TtsMakerParameter
\end{lstlisting}

\textbf{基类}：\passthrough{\lstinline!object!}

\textbf{公开属性}

\begin{longtable}[]{@{}
  >{\raggedright\arraybackslash}p{(\columnwidth - 6\tabcolsep) * \real{0.18}}
  >{\raggedright\arraybackslash}p{(\columnwidth - 6\tabcolsep) * \real{0.24}}
  >{\raggedright\arraybackslash}p{(\columnwidth - 6\tabcolsep) * \real{0.18}}
  >{\raggedright\arraybackslash}p{(\columnwidth - 6\tabcolsep) * \real{0.40}}@{}}
\toprule
名称 & Python 类型 & 含义 & 用法 \\
\midrule
\endhead
\passthrough{\lstinline!index!} & \passthrough{\lstinline!int!} &
传给该接口的 \passthrough{\lstinline!index!} 参数，Python 类型为
\passthrough{\lstinline!int!}。精确含义、单位、范围和枚举值需查目标型号对应头文件与协议。
& \passthrough{\lstinline!value = obj.index!} /
\passthrough{\lstinline!obj.index = value!} \\
\passthrough{\lstinline!speaker\_id!} & \passthrough{\lstinline!int!} &
语音合成说话人 ID。有效编号由机器人音频服务和固件决定。 &
\passthrough{\lstinline!value = obj.speaker\_id!} /
\passthrough{\lstinline!obj.speaker\_id = value!} \\
\passthrough{\lstinline!text!} & \passthrough{\lstinline!str!} &
要处理的文本内容；编码、长度和语言支持由目标服务决定。 &
\passthrough{\lstinline!value = obj.text!} /
\passthrough{\lstinline!obj.text = value!} \\
\bottomrule
\end{longtable}

\textbf{方法索引}

\begin{longtable}[]{@{}
  >{\raggedright\arraybackslash}p{(\columnwidth - 6\tabcolsep) * \real{0.18}}
  >{\raggedright\arraybackslash}p{(\columnwidth - 6\tabcolsep) * \real{0.24}}
  >{\raggedright\arraybackslash}p{(\columnwidth - 6\tabcolsep) * \real{0.18}}
  >{\raggedright\arraybackslash}p{(\columnwidth - 6\tabcolsep) * \real{0.40}}@{}}
\toprule
方法 & Python 签名 & 状态 & 安全分类 \\
\midrule
\endhead
\protect\hyperlink{unitree-sdk2-cpp-robot-g1-ttsmakerparameter-dunder-init-1}{\passthrough{\lstinline!\_\_init\_\_!}}
& \passthrough{\lstinline!def \_\_init\_\_(self) -> None!} &
\passthrough{\lstinline!SIGNATURE\_ONLY!} &
\passthrough{\lstinline!CONSTRUCTION!} \\
\protect\hyperlink{unitree-sdk2-cpp-robot-g1-ttsmakerparameter-from-json-1}{\passthrough{\lstinline!from\_json!}}
&
\passthrough{\lstinline!def from\_json(self, json: dict[str, Any]) -> None!}
& \passthrough{\lstinline!SIGNATURE\_ONLY!} &
\passthrough{\lstinline!UNCLASSIFIED!} \\
\protect\hyperlink{unitree-sdk2-cpp-robot-g1-ttsmakerparameter-to-json-1}{\passthrough{\lstinline!to\_json!}}
&
\passthrough{\lstinline!def to\_json(self, json: dict[str, Any]) -> None!}
& \passthrough{\lstinline!SIGNATURE\_ONLY!} &
\passthrough{\lstinline!UNCLASSIFIED!} \\
\bottomrule
\end{longtable}

\hypertarget{unitree-sdk2-cpp-robot-g1-ttsmakerparameter-dunder-init-1}{%
\paragraph{\texorpdfstring{\texttt{unitree\_sdk2\_cpp.robot.g1.TtsMakerParameter.\_\_init\_\_}}{unitree\_sdk2\_cpp.robot.g1.TtsMakerParameter.\_\_init\_\_}}\label{unitree-sdk2-cpp-robot-g1-ttsmakerparameter-dunder-init-1}}

初始化 \passthrough{\lstinline!TtsMakerParameter!}
实例。是否能实际构造取决于下方可用性状态。

\textbf{签名}

\begin{lstlisting}[language=Python]
def __init__(self) -> None
\end{lstlisting}

\textbf{可用性与安全性}

\textbf{\passthrough{\lstinline!SIGNATURE\_ONLY!} /
\passthrough{\lstinline!CONSTRUCTION!}}：当前只有设计期签名，不能假设已存在于运行时扩展。

\textbf{参数}

无显式参数。实例方法中的 \passthrough{\lstinline!self!} 由 Python
自动传入。

\textbf{返回值}

\begin{longtable}[]{@{}
  >{\raggedright\arraybackslash}p{(\columnwidth - 4\tabcolsep) * \real{0.18}}
  >{\raggedright\arraybackslash}p{(\columnwidth - 4\tabcolsep) * \real{0.20}}
  >{\raggedright\arraybackslash}p{(\columnwidth - 4\tabcolsep) * \real{0.62}}@{}}
\toprule
位置 & 类型 & 含义 \\
\midrule
\endhead
无 & \passthrough{\lstinline!None!} &
\passthrough{\lstinline!\_\_init\_\_!}
本身不返回值；调用类对象时，在构造函数可用的前提下得到该类实例。 \\
\bottomrule
\end{longtable}

\textbf{对应 C++}

\begin{itemize}
\tightlist
\item
  类：\passthrough{\lstinline!unitree::robot::g1::TtsMakerParameter!}
\item
  签名：\passthrough{\lstinline!TtsMakerParameter()!}
\item
  绑定策略：\passthrough{\lstinline!未单独标注!}
\item
  声明位置：\passthrough{\lstinline!include/unitree/robot/g1/audio/g1\_audio\_api.hpp:27!}
\end{itemize}

\textbf{说明}

\begin{itemize}
\tightlist
\item
  示例是局部调用片段；\passthrough{\lstinline!obj!}
  和参数变量表示已经按上层业务校验并准备好的对象。
\item
  该条目可以被编辑器和类型检查器识别，但当前运行时可能在属性查找阶段直接失败。
\end{itemize}

\textbf{用法}

\begin{lstlisting}[language=Python]
from typing import TYPE_CHECKING

if TYPE_CHECKING:
    def planned_construct() -> TtsMakerParameter:
        return TtsMakerParameter()
\end{lstlisting}

\begin{quote}
{[}!CAUTION{]} 上面的 \passthrough{\lstinline!TYPE\_CHECKING!}
示例只表达计划调用形态。该分支在普通 Python
运行时为假，不代表当前扩展能执行此接口。
\end{quote}

\hypertarget{unitree-sdk2-cpp-robot-g1-ttsmakerparameter-from-json-1}{%
\paragraph{\texorpdfstring{\texttt{unitree\_sdk2\_cpp.robot.g1.TtsMakerParameter.from\_json}}{unitree\_sdk2\_cpp.robot.g1.TtsMakerParameter.from\_json}}\label{unitree-sdk2-cpp-robot-g1-ttsmakerparameter-from-json-1}}

计划从 JSON 风格字典读取字段并更新当前 SDK 值对象。

\textbf{签名}

\begin{lstlisting}[language=Python]
def from_json(self, json: dict[str, Any]) -> None
\end{lstlisting}

\textbf{可用性与安全性}

\textbf{\passthrough{\lstinline!SIGNATURE\_ONLY!} /
\passthrough{\lstinline!UNCLASSIFIED!}}：当前只有设计期签名，不能假设已存在于运行时扩展。

\textbf{参数}

\begin{longtable}[]{@{}
  >{\raggedright\arraybackslash}p{(\columnwidth - 6\tabcolsep) * \real{0.18}}
  >{\raggedright\arraybackslash}p{(\columnwidth - 6\tabcolsep) * \real{0.24}}
  >{\raggedright\arraybackslash}p{(\columnwidth - 6\tabcolsep) * \real{0.18}}
  >{\raggedright\arraybackslash}p{(\columnwidth - 6\tabcolsep) * \real{0.40}}@{}}
\toprule
名称 & Python 类型 & 默认值 & 含义 \\
\midrule
\endhead
\passthrough{\lstinline!json!} &
\passthrough{\lstinline!dict[str, Any]!} & 必填 & JSON
风格字典，键和值必须满足对应 C++ \passthrough{\lstinline!JsonMap!}
数据结构的约束。 对应 C++ 参数
\passthrough{\lstinline!json: common::JsonMap \&!}。 \\
\bottomrule
\end{longtable}

\textbf{返回值}

\begin{longtable}[]{@{}
  >{\raggedright\arraybackslash}p{(\columnwidth - 4\tabcolsep) * \real{0.18}}
  >{\raggedright\arraybackslash}p{(\columnwidth - 4\tabcolsep) * \real{0.20}}
  >{\raggedright\arraybackslash}p{(\columnwidth - 4\tabcolsep) * \real{0.62}}@{}}
\toprule
位置 & 类型 & 含义 \\
\midrule
\endhead
无 & \passthrough{\lstinline!None!} & 该方法不返回 Python
值；失败可能表现为异常或底层状态变化。 \\
\bottomrule
\end{longtable}

\textbf{对应 C++}

\begin{itemize}
\tightlist
\item
  类：\passthrough{\lstinline!unitree::robot::g1::TtsMakerParameter!}
\item
  签名：\passthrough{\lstinline!fromJson(common::JsonMap \&)!}
\item
  绑定策略：\passthrough{\lstinline!SIGNATURE\_PREVIEW!}
\item
  声明位置：\passthrough{\lstinline!include/unitree/robot/g1/audio/g1\_audio\_api.hpp:30!}
\end{itemize}

\textbf{说明}

\begin{itemize}
\tightlist
\item
  示例是局部调用片段；\passthrough{\lstinline!obj!}
  和参数变量表示已经按上层业务校验并准备好的对象。
\item
  该条目可以被编辑器和类型检查器识别，但当前运行时可能在属性查找阶段直接失败。
\end{itemize}

\textbf{用法}

\begin{lstlisting}[language=Python]
from typing import TYPE_CHECKING

if TYPE_CHECKING:
    def planned_from_json(obj: TtsMakerParameter, json: dict[str, Any]) -> None:
        obj.from_json(json=json)
\end{lstlisting}

\begin{quote}
{[}!CAUTION{]} 上面的 \passthrough{\lstinline!TYPE\_CHECKING!}
示例只表达计划调用形态。该分支在普通 Python
运行时为假，不代表当前扩展能执行此接口。
\end{quote}

\hypertarget{unitree-sdk2-cpp-robot-g1-ttsmakerparameter-to-json-1}{%
\paragraph{\texorpdfstring{\texttt{unitree\_sdk2\_cpp.robot.g1.TtsMakerParameter.to\_json}}{unitree\_sdk2\_cpp.robot.g1.TtsMakerParameter.to\_json}}\label{unitree-sdk2-cpp-robot-g1-ttsmakerparameter-to-json-1}}

计划把当前 SDK 值对象写入 JSON 风格字典。

\textbf{签名}

\begin{lstlisting}[language=Python]
def to_json(self, json: dict[str, Any]) -> None
\end{lstlisting}

\textbf{可用性与安全性}

\textbf{\passthrough{\lstinline!SIGNATURE\_ONLY!} /
\passthrough{\lstinline!UNCLASSIFIED!}}：当前只有设计期签名，不能假设已存在于运行时扩展。

\textbf{参数}

\begin{longtable}[]{@{}
  >{\raggedright\arraybackslash}p{(\columnwidth - 6\tabcolsep) * \real{0.18}}
  >{\raggedright\arraybackslash}p{(\columnwidth - 6\tabcolsep) * \real{0.24}}
  >{\raggedright\arraybackslash}p{(\columnwidth - 6\tabcolsep) * \real{0.18}}
  >{\raggedright\arraybackslash}p{(\columnwidth - 6\tabcolsep) * \real{0.40}}@{}}
\toprule
名称 & Python 类型 & 默认值 & 含义 \\
\midrule
\endhead
\passthrough{\lstinline!json!} &
\passthrough{\lstinline!dict[str, Any]!} & 必填 & JSON
风格字典，键和值必须满足对应 C++ \passthrough{\lstinline!JsonMap!}
数据结构的约束。 对应 C++ 参数
\passthrough{\lstinline!json: common::JsonMap \&!}。 \\
\bottomrule
\end{longtable}

\textbf{返回值}

\begin{longtable}[]{@{}
  >{\raggedright\arraybackslash}p{(\columnwidth - 4\tabcolsep) * \real{0.18}}
  >{\raggedright\arraybackslash}p{(\columnwidth - 4\tabcolsep) * \real{0.20}}
  >{\raggedright\arraybackslash}p{(\columnwidth - 4\tabcolsep) * \real{0.62}}@{}}
\toprule
位置 & 类型 & 含义 \\
\midrule
\endhead
无 & \passthrough{\lstinline!None!} & 该方法不返回 Python
值；失败可能表现为异常或底层状态变化。 \\
\bottomrule
\end{longtable}

\textbf{对应 C++}

\begin{itemize}
\tightlist
\item
  类：\passthrough{\lstinline!unitree::robot::g1::TtsMakerParameter!}
\item
  签名：\passthrough{\lstinline!toJson(common::JsonMap \&) const!}
\item
  绑定策略：\passthrough{\lstinline!SIGNATURE\_PREVIEW!}
\item
  声明位置：\passthrough{\lstinline!include/unitree/robot/g1/audio/g1\_audio\_api.hpp:32!}
\end{itemize}

\textbf{说明}

\begin{itemize}
\tightlist
\item
  示例是局部调用片段；\passthrough{\lstinline!obj!}
  和参数变量表示已经按上层业务校验并准备好的对象。
\item
  该条目可以被编辑器和类型检查器识别，但当前运行时可能在属性查找阶段直接失败。
\end{itemize}

\textbf{用法}

\begin{lstlisting}[language=Python]
from typing import TYPE_CHECKING

if TYPE_CHECKING:
    def planned_to_json(obj: TtsMakerParameter, json: dict[str, Any]) -> None:
        obj.to_json(json=json)
\end{lstlisting}

\begin{quote}
{[}!CAUTION{]} 上面的 \passthrough{\lstinline!TYPE\_CHECKING!}
示例只表达计划调用形态。该分支在普通 Python
运行时为假，不代表当前扩展能执行此接口。
\end{quote}

\begin{center}\rule{0.5\linewidth}{0.5pt}\end{center}

\hypertarget{unitree-sdk2-cpp-robot-go2}{%
\subsection{Go2 Robot API}\label{unitree-sdk2-cpp-robot-go2}}

模块：\passthrough{\lstinline!unitree\_sdk2\_cpp.robot.go2!}

\hypertarget{ux7c7bux7d22ux5f15-11}{%
\subsubsection{类索引}\label{ux7c7bux7d22ux5f15-11}}

\begin{longtable}[]{@{}
  >{\raggedright\arraybackslash}p{(\columnwidth - 4\tabcolsep) * \real{0.18}}
  >{\raggedleft\arraybackslash}p{(\columnwidth - 4\tabcolsep) * \real{0.20}}
  >{\raggedleft\arraybackslash}p{(\columnwidth - 4\tabcolsep) * \real{0.62}}@{}}
\toprule
类 & 公开函数签名 & 属性 \\
\midrule
\endhead
\protect\hyperlink{unitree-sdk2-cpp-robot-go2-configclient}{\passthrough{\lstinline!ConfigClient!}}
& 8 & 0 \\
\protect\hyperlink{unitree-sdk2-cpp-robot-go2-configdelparameter}{\passthrough{\lstinline!ConfigDelParameter!}}
& 3 & 0 \\
\protect\hyperlink{unitree-sdk2-cpp-robot-go2-configgetdata}{\passthrough{\lstinline!ConfigGetData!}}
& 3 & 0 \\
\protect\hyperlink{unitree-sdk2-cpp-robot-go2-configgetparameter}{\passthrough{\lstinline!ConfigGetParameter!}}
& 3 & 0 \\
\protect\hyperlink{unitree-sdk2-cpp-robot-go2-configmeta}{\passthrough{\lstinline!ConfigMeta!}}
& 1 & 0 \\
\protect\hyperlink{unitree-sdk2-cpp-robot-go2-configmetadata}{\passthrough{\lstinline!ConfigMetaData!}}
& 3 & 0 \\
\protect\hyperlink{unitree-sdk2-cpp-robot-go2-configmetaparameter}{\passthrough{\lstinline!ConfigMetaParameter!}}
& 3 & 0 \\
\protect\hyperlink{unitree-sdk2-cpp-robot-go2-configsetparameter}{\passthrough{\lstinline!ConfigSetParameter!}}
& 3 & 0 \\
\protect\hyperlink{unitree-sdk2-cpp-robot-go2-jsonizecommobjint}{\passthrough{\lstinline!JsonizeCommObjInt!}}
& 3 & 0 \\
\protect\hyperlink{unitree-sdk2-cpp-robot-go2-jsonizeconfigmeta}{\passthrough{\lstinline!JsonizeConfigMeta!}}
& 3 & 0 \\
\protect\hyperlink{unitree-sdk2-cpp-robot-go2-jsonizedatabool}{\passthrough{\lstinline!JsonizeDataBool!}}
& 3 & 0 \\
\protect\hyperlink{unitree-sdk2-cpp-robot-go2-jsonizedatadouble}{\passthrough{\lstinline!JsonizeDataDouble!}}
& 3 & 0 \\
\protect\hyperlink{unitree-sdk2-cpp-robot-go2-jsonizedatafloat}{\passthrough{\lstinline!JsonizeDataFloat!}}
& 3 & 0 \\
\protect\hyperlink{unitree-sdk2-cpp-robot-go2-jsonizedataint}{\passthrough{\lstinline!JsonizeDataInt!}}
& 3 & 0 \\
\protect\hyperlink{unitree-sdk2-cpp-robot-go2-jsonizedatastring}{\passthrough{\lstinline!JsonizeDataString!}}
& 3 & 0 \\
\protect\hyperlink{unitree-sdk2-cpp-robot-go2-jsonizeflagbool}{\passthrough{\lstinline!JsonizeFlagBool!}}
& 3 & 0 \\
\protect\hyperlink{unitree-sdk2-cpp-robot-go2-jsonizepathpoint}{\passthrough{\lstinline!JsonizePathPoint!}}
& 3 & 0 \\
\protect\hyperlink{unitree-sdk2-cpp-robot-go2-jsonizequat}{\passthrough{\lstinline!JsonizeQuat!}}
& 3 & 0 \\
\protect\hyperlink{unitree-sdk2-cpp-robot-go2-jsonizevec3}{\passthrough{\lstinline!JsonizeVec3!}}
& 3 & 0 \\
\protect\hyperlink{unitree-sdk2-cpp-robot-go2-obstaclesavoidclient}{\passthrough{\lstinline!ObstaclesAvoidClient!}}
& 8 & 0 \\
\protect\hyperlink{unitree-sdk2-cpp-robot-go2-obstaclesavoidmoveparameter}{\passthrough{\lstinline!ObstaclesAvoidMoveParameter!}}
& 2 & 0 \\
\protect\hyperlink{unitree-sdk2-cpp-robot-go2-obstaclesavoidremotecommandsource}{\passthrough{\lstinline!ObstaclesAvoidRemoteCommandSource!}}
& 2 & 0 \\
\protect\hyperlink{unitree-sdk2-cpp-robot-go2-obstaclesavoidswitchgetdata}{\passthrough{\lstinline!ObstaclesAvoidSwitchGetData!}}
& 2 & 0 \\
\protect\hyperlink{unitree-sdk2-cpp-robot-go2-obstaclesavoidswitchsetparameter}{\passthrough{\lstinline!ObstaclesAvoidSwitchSetParameter!}}
& 2 & 0 \\
\protect\hyperlink{unitree-sdk2-cpp-robot-go2-robotstateclient}{\passthrough{\lstinline!RobotStateClient!}}
& 5 & 0 \\
\protect\hyperlink{unitree-sdk2-cpp-robot-go2-servicestate}{\passthrough{\lstinline!ServiceState!}}
& 1 & 0 \\
\protect\hyperlink{unitree-sdk2-cpp-robot-go2-servicestatedata}{\passthrough{\lstinline!ServiceStateData!}}
& 3 & 0 \\
\protect\hyperlink{unitree-sdk2-cpp-robot-go2-serviceswitchdata}{\passthrough{\lstinline!ServiceSwitchData!}}
& 3 & 0 \\
\protect\hyperlink{unitree-sdk2-cpp-robot-go2-serviceswitchparameter}{\passthrough{\lstinline!ServiceSwitchParameter!}}
& 3 & 0 \\
\protect\hyperlink{unitree-sdk2-cpp-robot-go2-setreportfreqparameter}{\passthrough{\lstinline!SetReportFreqParameter!}}
& 3 & 0 \\
\protect\hyperlink{unitree-sdk2-cpp-robot-go2-sportclient}{\passthrough{\lstinline!SportClient!}}
& 41 & 0 \\
\protect\hyperlink{unitree-sdk2-cpp-robot-go2-utrackclient}{\passthrough{\lstinline!UtrackClient!}}
& 5 & 0 \\
\protect\hyperlink{unitree-sdk2-cpp-robot-go2-utrackswitchgetdata}{\passthrough{\lstinline!UtrackSwitchGetData!}}
& 2 & 0 \\
\protect\hyperlink{unitree-sdk2-cpp-robot-go2-utrackswitchsetparameter}{\passthrough{\lstinline!UtrackSwitchSetParameter!}}
& 2 & 0 \\
\protect\hyperlink{unitree-sdk2-cpp-robot-go2-videoclient}{\passthrough{\lstinline!VideoClient!}}
& 3 & 0 \\
\protect\hyperlink{unitree-sdk2-cpp-robot-go2-vuiclient}{\passthrough{\lstinline!VuiClient!}}
& 8 & 0 \\
\protect\hyperlink{unitree-sdk2-cpp-robot-go2-stpathpoint}{\passthrough{\lstinline!stPathPoint!}}
& 0 & 0 \\
\bottomrule
\end{longtable}

\hypertarget{unitree-sdk2-cpp-robot-go2-configclient}{%
\subsubsection{\texorpdfstring{\texttt{unitree\_sdk2\_cpp.robot.go2.ConfigClient}}{unitree\_sdk2\_cpp.robot.go2.ConfigClient}}\label{unitree-sdk2-cpp-robot-go2-configclient}}

Robot 服务客户端。构造、初始化和具体方法的可用性必须分别检查。

\textbf{导入}

\begin{lstlisting}[language=Python]
from unitree_sdk2_cpp.robot.go2 import ConfigClient
\end{lstlisting}

\textbf{基类}：\passthrough{\lstinline!Client!}

\textbf{方法索引}

\begin{longtable}[]{@{}
  >{\raggedright\arraybackslash}p{(\columnwidth - 6\tabcolsep) * \real{0.18}}
  >{\raggedright\arraybackslash}p{(\columnwidth - 6\tabcolsep) * \real{0.24}}
  >{\raggedright\arraybackslash}p{(\columnwidth - 6\tabcolsep) * \real{0.18}}
  >{\raggedright\arraybackslash}p{(\columnwidth - 6\tabcolsep) * \real{0.40}}@{}}
\toprule
方法 & Python 签名 & 状态 & 安全分类 \\
\midrule
\endhead
\protect\hyperlink{unitree-sdk2-cpp-robot-go2-configclient-dunder-init-1}{\passthrough{\lstinline!\_\_init\_\_!}}
& \passthrough{\lstinline!def \_\_init\_\_(self) -> None!} &
\passthrough{\lstinline!AVAILABLE!} &
\passthrough{\lstinline!CONSTRUCTION!} \\
\protect\hyperlink{unitree-sdk2-cpp-robot-go2-configclient-init-1}{\passthrough{\lstinline!init!}}
& \passthrough{\lstinline!def init(self) -> None!} &
\passthrough{\lstinline!AVAILABLE!} &
\passthrough{\lstinline!INITIALIZATION!} \\
\protect\hyperlink{unitree-sdk2-cpp-robot-go2-configclient-set-1}{\passthrough{\lstinline!set!}}
&
\passthrough{\lstinline!def set(self, name: str, content: str) -> int!}
& \passthrough{\lstinline!SIGNATURE\_ONLY!} &
\passthrough{\lstinline!HARDWARE\_SIDE\_EFFECT!} \\
\protect\hyperlink{unitree-sdk2-cpp-robot-go2-configclient-get-1}{\passthrough{\lstinline!get!}}
& \passthrough{\lstinline!def get(self, name: str) -> tuple[int, str]!}
& \passthrough{\lstinline!AVAILABLE!} &
\passthrough{\lstinline!READ\_ONLY!} \\
\protect\hyperlink{unitree-sdk2-cpp-robot-go2-configclient-del-1}{\passthrough{\lstinline!del\_!}}
& \passthrough{\lstinline!def del\_(self, name: str) -> int!} &
\passthrough{\lstinline!SIGNATURE\_ONLY!} &
\passthrough{\lstinline!HARDWARE\_SIDE\_EFFECT!} \\
\protect\hyperlink{unitree-sdk2-cpp-robot-go2-configclient-meta-config-meta-1}{\passthrough{\lstinline!meta\_config\_meta!}}
&
\passthrough{\lstinline!def meta\_config\_meta(self, name: str) -> tuple[int, ConfigMeta]!}
& \passthrough{\lstinline!SIGNATURE\_ONLY!} &
\passthrough{\lstinline!HARDWARE\_SIDE\_EFFECT!} \\
\protect\hyperlink{unitree-sdk2-cpp-robot-go2-configclient-meta-string-1}{\passthrough{\lstinline!meta\_string!}}
&
\passthrough{\lstinline!def meta\_string(self, name: str) -> tuple[int, str]!}
& \passthrough{\lstinline!SIGNATURE\_ONLY!} &
\passthrough{\lstinline!HARDWARE\_SIDE\_EFFECT!} \\
\protect\hyperlink{unitree-sdk2-cpp-robot-go2-configclient-subscribe-change-status-1}{\passthrough{\lstinline!subscribe\_change\_status!}}
&
\passthrough{\lstinline!def subscribe\_change\_status(self, name: str, callback: Callable[[str, str], None]) -> None!}
& \passthrough{\lstinline!SIGNATURE\_ONLY!} &
\passthrough{\lstinline!HARDWARE\_SIDE\_EFFECT!} \\
\bottomrule
\end{longtable}

\hypertarget{unitree-sdk2-cpp-robot-go2-configclient-dunder-init-1}{%
\paragraph{\texorpdfstring{\texttt{unitree\_sdk2\_cpp.robot.go2.ConfigClient.\_\_init\_\_}}{unitree\_sdk2\_cpp.robot.go2.ConfigClient.\_\_init\_\_}}\label{unitree-sdk2-cpp-robot-go2-configclient-dunder-init-1}}

初始化 \passthrough{\lstinline!ConfigClient!}
实例。是否能实际构造取决于下方可用性状态。

\textbf{签名}

\begin{lstlisting}[language=Python]
def __init__(self) -> None
\end{lstlisting}

\textbf{可用性与安全性}

\textbf{\passthrough{\lstinline!AVAILABLE!} /
\passthrough{\lstinline!CONSTRUCTION!}}：当前绑定源码已实现；实际执行条件仍取决于平台和对应
SDK 环境。

\textbf{参数}

无显式参数。实例方法中的 \passthrough{\lstinline!self!} 由 Python
自动传入。

\textbf{返回值}

\begin{longtable}[]{@{}
  >{\raggedright\arraybackslash}p{(\columnwidth - 4\tabcolsep) * \real{0.18}}
  >{\raggedright\arraybackslash}p{(\columnwidth - 4\tabcolsep) * \real{0.20}}
  >{\raggedright\arraybackslash}p{(\columnwidth - 4\tabcolsep) * \real{0.62}}@{}}
\toprule
位置 & 类型 & 含义 \\
\midrule
\endhead
无 & \passthrough{\lstinline!None!} &
\passthrough{\lstinline!\_\_init\_\_!}
本身不返回值；调用类对象时，在构造函数可用的前提下得到该类实例。 \\
\bottomrule
\end{longtable}

\textbf{对应 C++}

\begin{itemize}
\tightlist
\item
  类：\passthrough{\lstinline!unitree::robot::go2::ConfigClient!}
\item
  签名：\passthrough{\lstinline!ConfigClient()!}
\item
  绑定策略：\passthrough{\lstinline!未单独标注!}
\item
  声明位置：\passthrough{\lstinline!include/unitree/robot/go2/config/config\_client.hpp:41!}
\end{itemize}

\textbf{说明}

\begin{itemize}
\tightlist
\item
  示例是局部调用片段；\passthrough{\lstinline!obj!}
  和参数变量表示已经按上层业务校验并准备好的对象。
\end{itemize}

\textbf{用法}

\begin{lstlisting}[language=Python]
value = ConfigClient()
\end{lstlisting}

\hypertarget{unitree-sdk2-cpp-robot-go2-configclient-init-1}{%
\paragraph{\texorpdfstring{\texttt{unitree\_sdk2\_cpp.robot.go2.ConfigClient.init}}{unitree\_sdk2\_cpp.robot.go2.ConfigClient.init}}\label{unitree-sdk2-cpp-robot-go2-configclient-init-1}}

初始化当前 SDK 对象所需的底层通道或服务资源。

\textbf{签名}

\begin{lstlisting}[language=Python]
def init(self) -> None
\end{lstlisting}

\textbf{可用性与安全性}

\textbf{\passthrough{\lstinline!AVAILABLE!} /
\passthrough{\lstinline!INITIALIZATION!}}：当前绑定源码已实现；实际执行条件仍取决于平台和对应
SDK 环境。

\textbf{参数}

无显式参数。实例方法中的 \passthrough{\lstinline!self!} 由 Python
自动传入。

\textbf{返回值}

\begin{longtable}[]{@{}
  >{\raggedright\arraybackslash}p{(\columnwidth - 4\tabcolsep) * \real{0.18}}
  >{\raggedright\arraybackslash}p{(\columnwidth - 4\tabcolsep) * \real{0.20}}
  >{\raggedright\arraybackslash}p{(\columnwidth - 4\tabcolsep) * \real{0.62}}@{}}
\toprule
位置 & 类型 & 含义 \\
\midrule
\endhead
无 & \passthrough{\lstinline!None!} & 该方法不返回 Python
值；失败可能表现为异常或底层状态变化。 \\
\bottomrule
\end{longtable}

\textbf{对应 C++}

\begin{itemize}
\tightlist
\item
  类：\passthrough{\lstinline!unitree::robot::go2::ConfigClient!}
\item
  签名：\passthrough{\lstinline!Init()!}
\item
  绑定策略：\passthrough{\lstinline!DIRECT!}
\item
  声明位置：\passthrough{\lstinline!include/unitree/robot/go2/config/config\_client.hpp:44!}
\end{itemize}

\textbf{说明}

\begin{itemize}
\tightlist
\item
  示例是局部调用片段；\passthrough{\lstinline!obj!}
  和参数变量表示已经按上层业务校验并准备好的对象。
\end{itemize}

\textbf{用法}

\begin{lstlisting}[language=Python]
obj.init()
\end{lstlisting}

\hypertarget{unitree-sdk2-cpp-robot-go2-configclient-set-1}{%
\paragraph{\texorpdfstring{\texttt{unitree\_sdk2\_cpp.robot.go2.ConfigClient.set}}{unitree\_sdk2\_cpp.robot.go2.ConfigClient.set}}\label{unitree-sdk2-cpp-robot-go2-configclient-set-1}}

对应 C++ SDK 操作
\passthrough{\lstinline!Set(const std::string \&, const std::string \&)!}。上游头文件没有可直接生成的业务说明时，本参考只保证签名映射，精确语义需查目标型号协议。

\textbf{签名}

\begin{lstlisting}[language=Python]
def set(self, name: str, content: str) -> int
\end{lstlisting}

\textbf{可用性与安全性}

\textbf{\passthrough{\lstinline!SIGNATURE\_ONLY!} /
\passthrough{\lstinline!HARDWARE\_SIDE\_EFFECT!}}：仅用于补全和静态检查。当前二进制没有该
Python 方法，未来实现还需单独评审硬件副作用。

\textbf{参数}

\begin{longtable}[]{@{}
  >{\raggedright\arraybackslash}p{(\columnwidth - 6\tabcolsep) * \real{0.18}}
  >{\raggedright\arraybackslash}p{(\columnwidth - 6\tabcolsep) * \real{0.24}}
  >{\raggedright\arraybackslash}p{(\columnwidth - 6\tabcolsep) * \real{0.18}}
  >{\raggedright\arraybackslash}p{(\columnwidth - 6\tabcolsep) * \real{0.40}}@{}}
\toprule
名称 & Python 类型 & 默认值 & 含义 \\
\midrule
\endhead
\passthrough{\lstinline!name!} & \passthrough{\lstinline!str!} & 必填 &
SDK 对象、配置项、服务或动作的名称。允许值由调用它的具体接口定义。 对应
C++ 参数 \passthrough{\lstinline!name: const std::string \&!}。
底层为字符串；长度、编码和允许值由具体协议决定。 \\
\passthrough{\lstinline!content!} & \passthrough{\lstinline!str!} & 必填
& 配置、请求或序列化内容。具体格式由对应服务协议定义。 对应 C++ 参数
\passthrough{\lstinline!content: const std::string \&!}。
底层为字符串；长度、编码和允许值由具体协议决定。 \\
\bottomrule
\end{longtable}

\textbf{返回值}

\begin{longtable}[]{@{}
  >{\raggedright\arraybackslash}p{(\columnwidth - 4\tabcolsep) * \real{0.18}}
  >{\raggedright\arraybackslash}p{(\columnwidth - 4\tabcolsep) * \real{0.20}}
  >{\raggedright\arraybackslash}p{(\columnwidth - 4\tabcolsep) * \real{0.62}}@{}}
\toprule
位置 & 类型 & 含义 \\
\midrule
\endhead
返回值 & \passthrough{\lstinline!int!} & SDK 状态码；Unitree 示例通常以
\passthrough{\lstinline!0!} 表示成功，非零值需按具体服务定义解释。 \\
\bottomrule
\end{longtable}

\textbf{对应 C++}

\begin{itemize}
\tightlist
\item
  类：\passthrough{\lstinline!unitree::robot::go2::ConfigClient!}
\item
  签名：\passthrough{\lstinline!Set(const std::string \&, const std::string \&)!}
\item
  绑定策略：\passthrough{\lstinline!DIRECT!}
\item
  声明位置：\passthrough{\lstinline!include/unitree/robot/go2/config/config\_client.hpp:46!}
\end{itemize}

\textbf{说明}

\begin{itemize}
\tightlist
\item
  示例是局部调用片段；\passthrough{\lstinline!obj!}
  和参数变量表示已经按上层业务校验并准备好的对象。
\item
  该条目可以被编辑器和类型检查器识别，但当前运行时可能在属性查找阶段直接失败。
\end{itemize}

\textbf{用法}

\begin{lstlisting}[language=Python]
from typing import TYPE_CHECKING

if TYPE_CHECKING:
    def planned_set(obj: ConfigClient, name: str, content: str) -> int:
        return obj.set(name=name, content=content)
\end{lstlisting}

\begin{quote}
{[}!CAUTION{]} 上面的 \passthrough{\lstinline!TYPE\_CHECKING!}
示例只表达计划调用形态。该分支在普通 Python
运行时为假，不代表当前扩展能执行此接口。
\end{quote}

\hypertarget{unitree-sdk2-cpp-robot-go2-configclient-get-1}{%
\paragraph{\texorpdfstring{\texttt{unitree\_sdk2\_cpp.robot.go2.ConfigClient.get}}{unitree\_sdk2\_cpp.robot.go2.ConfigClient.get}}\label{unitree-sdk2-cpp-robot-go2-configclient-get-1}}

对应 C++ SDK 操作
\passthrough{\lstinline!Get(const std::string \&, std::string \&)!}。上游头文件没有可直接生成的业务说明时，本参考只保证签名映射，精确语义需查目标型号协议。

\textbf{签名}

\begin{lstlisting}[language=Python]
def get(self, name: str) -> tuple[int, str]
\end{lstlisting}

\textbf{可用性与安全性}

\textbf{\passthrough{\lstinline!AVAILABLE!} /
\passthrough{\lstinline!READ\_ONLY!}}：当前绑定源码已实现只读查询。Robot
Client 的服务端查询通常仍需匹配的 Linux
扩展、DDS、网络、机器人和服务；纯本地 getter 则不一定需要全部条件。

\textbf{参数}

\begin{longtable}[]{@{}
  >{\raggedright\arraybackslash}p{(\columnwidth - 6\tabcolsep) * \real{0.18}}
  >{\raggedright\arraybackslash}p{(\columnwidth - 6\tabcolsep) * \real{0.24}}
  >{\raggedright\arraybackslash}p{(\columnwidth - 6\tabcolsep) * \real{0.18}}
  >{\raggedright\arraybackslash}p{(\columnwidth - 6\tabcolsep) * \real{0.40}}@{}}
\toprule
名称 & Python 类型 & 默认值 & 含义 \\
\midrule
\endhead
\passthrough{\lstinline!name!} & \passthrough{\lstinline!str!} & 必填 &
SDK 对象、配置项、服务或动作的名称。允许值由调用它的具体接口定义。 对应
C++ 参数 \passthrough{\lstinline!name: const std::string \&!}。
底层为字符串；长度、编码和允许值由具体协议决定。 \\
\bottomrule
\end{longtable}

\textbf{返回值}

\begin{longtable}[]{@{}
  >{\raggedright\arraybackslash}p{(\columnwidth - 4\tabcolsep) * \real{0.18}}
  >{\raggedright\arraybackslash}p{(\columnwidth - 4\tabcolsep) * \real{0.20}}
  >{\raggedright\arraybackslash}p{(\columnwidth - 4\tabcolsep) * \real{0.62}}@{}}
\toprule
位置 & 类型 & 含义 \\
\midrule
\endhead
{[}0{]} \passthrough{\lstinline!status!} & \passthrough{\lstinline!int!}
& SDK 状态码；Unitree 示例通常以 \passthrough{\lstinline!0!}
表示成功，非零值需按具体服务错误码解释。 \\
{[}1{]} \passthrough{\lstinline!content!} &
\passthrough{\lstinline!str!} &
配置、请求或序列化内容。具体格式由对应服务协议定义。 对应 C++ 参数
\passthrough{\lstinline!content: std::string \&!}。
底层为字符串；长度、编码和允许值由具体协议决定。 \\
\bottomrule
\end{longtable}

\textbf{对应 C++}

\begin{itemize}
\tightlist
\item
  类：\passthrough{\lstinline!unitree::robot::go2::ConfigClient!}
\item
  签名：\passthrough{\lstinline!Get(const std::string \&, std::string \&)!}
\item
  绑定策略：\passthrough{\lstinline!OUTPUT\_WRAPPER!}
\item
  声明位置：\passthrough{\lstinline!include/unitree/robot/go2/config/config\_client.hpp:47!}
\end{itemize}

\textbf{说明}

\begin{itemize}
\tightlist
\item
  示例是局部调用片段；\passthrough{\lstinline!obj!}
  和参数变量表示已经按上层业务校验并准备好的对象。
\item
  C++ 的可修改输出引用不会作为 Python 入参出现，而是按 C++
  参数顺序追加到返回元组中。
\end{itemize}

\textbf{用法}

\begin{lstlisting}[language=Python]
status, content = obj.get(name=name)
\end{lstlisting}

\begin{quote}
{[}!NOTE{]} 只读查询仍会访问
DDS/机器人服务。必须检查返回状态码，并处理超时和网络中断。
\end{quote}

\hypertarget{unitree-sdk2-cpp-robot-go2-configclient-del-1}{%
\paragraph{\texorpdfstring{\texttt{unitree\_sdk2\_cpp.robot.go2.ConfigClient.del\_}}{unitree\_sdk2\_cpp.robot.go2.ConfigClient.del\_}}\label{unitree-sdk2-cpp-robot-go2-configclient-del-1}}

对应 C++ SDK 操作
\passthrough{\lstinline!Del(const std::string \&)!}。上游头文件没有可直接生成的业务说明时，本参考只保证签名映射，精确语义需查目标型号协议。

\textbf{签名}

\begin{lstlisting}[language=Python]
def del_(self, name: str) -> int
\end{lstlisting}

\textbf{可用性与安全性}

\textbf{\passthrough{\lstinline!SIGNATURE\_ONLY!} /
\passthrough{\lstinline!HARDWARE\_SIDE\_EFFECT!}}：仅用于补全和静态检查。当前二进制没有该
Python 方法，未来实现还需单独评审硬件副作用。

\textbf{参数}

\begin{longtable}[]{@{}
  >{\raggedright\arraybackslash}p{(\columnwidth - 6\tabcolsep) * \real{0.18}}
  >{\raggedright\arraybackslash}p{(\columnwidth - 6\tabcolsep) * \real{0.24}}
  >{\raggedright\arraybackslash}p{(\columnwidth - 6\tabcolsep) * \real{0.18}}
  >{\raggedright\arraybackslash}p{(\columnwidth - 6\tabcolsep) * \real{0.40}}@{}}
\toprule
名称 & Python 类型 & 默认值 & 含义 \\
\midrule
\endhead
\passthrough{\lstinline!name!} & \passthrough{\lstinline!str!} & 必填 &
SDK 对象、配置项、服务或动作的名称。允许值由调用它的具体接口定义。 对应
C++ 参数 \passthrough{\lstinline!name: const std::string \&!}。
底层为字符串；长度、编码和允许值由具体协议决定。 \\
\bottomrule
\end{longtable}

\textbf{返回值}

\begin{longtable}[]{@{}
  >{\raggedright\arraybackslash}p{(\columnwidth - 4\tabcolsep) * \real{0.18}}
  >{\raggedright\arraybackslash}p{(\columnwidth - 4\tabcolsep) * \real{0.20}}
  >{\raggedright\arraybackslash}p{(\columnwidth - 4\tabcolsep) * \real{0.62}}@{}}
\toprule
位置 & 类型 & 含义 \\
\midrule
\endhead
返回值 & \passthrough{\lstinline!int!} & SDK 状态码；Unitree 示例通常以
\passthrough{\lstinline!0!} 表示成功，非零值需按具体服务定义解释。 \\
\bottomrule
\end{longtable}

\textbf{对应 C++}

\begin{itemize}
\tightlist
\item
  类：\passthrough{\lstinline!unitree::robot::go2::ConfigClient!}
\item
  签名：\passthrough{\lstinline!Del(const std::string \&)!}
\item
  绑定策略：\passthrough{\lstinline!DIRECT!}
\item
  声明位置：\passthrough{\lstinline!include/unitree/robot/go2/config/config\_client.hpp:48!}
\end{itemize}

\textbf{说明}

\begin{itemize}
\tightlist
\item
  示例是局部调用片段；\passthrough{\lstinline!obj!}
  和参数变量表示已经按上层业务校验并准备好的对象。
\item
  该条目可以被编辑器和类型检查器识别，但当前运行时可能在属性查找阶段直接失败。
\end{itemize}

\textbf{用法}

\begin{lstlisting}[language=Python]
from typing import TYPE_CHECKING

if TYPE_CHECKING:
    def planned_del_(obj: ConfigClient, name: str) -> int:
        return obj.del_(name=name)
\end{lstlisting}

\begin{quote}
{[}!CAUTION{]} 上面的 \passthrough{\lstinline!TYPE\_CHECKING!}
示例只表达计划调用形态。该分支在普通 Python
运行时为假，不代表当前扩展能执行此接口。
\end{quote}

\hypertarget{unitree-sdk2-cpp-robot-go2-configclient-meta-config-meta-1}{%
\paragraph{\texorpdfstring{\texttt{unitree\_sdk2\_cpp.robot.go2.ConfigClient.meta\_config\_meta}}{unitree\_sdk2\_cpp.robot.go2.ConfigClient.meta\_config\_meta}}\label{unitree-sdk2-cpp-robot-go2-configclient-meta-config-meta-1}}

对应 C++ SDK 操作
\passthrough{\lstinline!Meta(const std::string \&, unitree::robot::go2::ConfigMeta \&)!}。上游头文件没有可直接生成的业务说明时，本参考只保证签名映射，精确语义需查目标型号协议。

\textbf{签名}

\begin{lstlisting}[language=Python]
def meta_config_meta(self, name: str) -> tuple[int, ConfigMeta]
\end{lstlisting}

\textbf{可用性与安全性}

\textbf{\passthrough{\lstinline!SIGNATURE\_ONLY!} /
\passthrough{\lstinline!HARDWARE\_SIDE\_EFFECT!}}：仅用于补全和静态检查。当前二进制没有该
Python 方法，未来实现还需单独评审硬件副作用。

\textbf{参数}

\begin{longtable}[]{@{}
  >{\raggedright\arraybackslash}p{(\columnwidth - 6\tabcolsep) * \real{0.18}}
  >{\raggedright\arraybackslash}p{(\columnwidth - 6\tabcolsep) * \real{0.24}}
  >{\raggedright\arraybackslash}p{(\columnwidth - 6\tabcolsep) * \real{0.18}}
  >{\raggedright\arraybackslash}p{(\columnwidth - 6\tabcolsep) * \real{0.40}}@{}}
\toprule
名称 & Python 类型 & 默认值 & 含义 \\
\midrule
\endhead
\passthrough{\lstinline!name!} & \passthrough{\lstinline!str!} & 必填 &
SDK 对象、配置项、服务或动作的名称。允许值由调用它的具体接口定义。 对应
C++ 参数 \passthrough{\lstinline!name: const std::string \&!}。
底层为字符串；长度、编码和允许值由具体协议决定。 \\
\bottomrule
\end{longtable}

\textbf{返回值}

\begin{longtable}[]{@{}
  >{\raggedright\arraybackslash}p{(\columnwidth - 4\tabcolsep) * \real{0.18}}
  >{\raggedright\arraybackslash}p{(\columnwidth - 4\tabcolsep) * \real{0.20}}
  >{\raggedright\arraybackslash}p{(\columnwidth - 4\tabcolsep) * \real{0.62}}@{}}
\toprule
位置 & 类型 & 含义 \\
\midrule
\endhead
{[}0{]} \passthrough{\lstinline!status!} & \passthrough{\lstinline!int!}
& SDK 状态码；Unitree 示例通常以 \passthrough{\lstinline!0!}
表示成功，非零值需按具体服务错误码解释。 \\
{[}1{]} \passthrough{\lstinline!meta!} &
\passthrough{\lstinline!ConfigMeta!} & 传给该接口的
\passthrough{\lstinline!meta!} 参数，Python 类型为
\passthrough{\lstinline!ConfigMeta!}。精确含义、单位、范围和枚举值需查目标型号对应头文件与协议。
对应 C++ 参数
\passthrough{\lstinline!meta: unitree::robot::go2::ConfigMeta \&!}。 \\
\bottomrule
\end{longtable}

\textbf{对应 C++}

\begin{itemize}
\tightlist
\item
  类：\passthrough{\lstinline!unitree::robot::go2::ConfigClient!}
\item
  签名：\passthrough{\lstinline!Meta(const std::string \&, unitree::robot::go2::ConfigMeta \&)!}
\item
  绑定策略：\passthrough{\lstinline!OUTPUT\_WRAPPER!}
\item
  声明位置：\passthrough{\lstinline!include/unitree/robot/go2/config/config\_client.hpp:49!}
\end{itemize}

\textbf{说明}

\begin{itemize}
\tightlist
\item
  示例是局部调用片段；\passthrough{\lstinline!obj!}
  和参数变量表示已经按上层业务校验并准备好的对象。
\item
  C++ 的可修改输出引用不会作为 Python 入参出现，而是按 C++
  参数顺序追加到返回元组中。
\item
  该条目可以被编辑器和类型检查器识别，但当前运行时可能在属性查找阶段直接失败。
\end{itemize}

\textbf{用法}

\begin{lstlisting}[language=Python]
from typing import TYPE_CHECKING

if TYPE_CHECKING:
    def planned_meta_config_meta(obj: ConfigClient, name: str) -> tuple[int, ConfigMeta]:
        return obj.meta_config_meta(name=name)
\end{lstlisting}

\begin{quote}
{[}!CAUTION{]} 上面的 \passthrough{\lstinline!TYPE\_CHECKING!}
示例只表达计划调用形态。该分支在普通 Python
运行时为假，不代表当前扩展能执行此接口。
\end{quote}

\hypertarget{unitree-sdk2-cpp-robot-go2-configclient-meta-string-1}{%
\paragraph{\texorpdfstring{\texttt{unitree\_sdk2\_cpp.robot.go2.ConfigClient.meta\_string}}{unitree\_sdk2\_cpp.robot.go2.ConfigClient.meta\_string}}\label{unitree-sdk2-cpp-robot-go2-configclient-meta-string-1}}

对应 C++ SDK 操作
\passthrough{\lstinline!Meta(const std::string \&, std::string \&)!}。上游头文件没有可直接生成的业务说明时，本参考只保证签名映射，精确语义需查目标型号协议。

\textbf{签名}

\begin{lstlisting}[language=Python]
def meta_string(self, name: str) -> tuple[int, str]
\end{lstlisting}

\textbf{可用性与安全性}

\textbf{\passthrough{\lstinline!SIGNATURE\_ONLY!} /
\passthrough{\lstinline!HARDWARE\_SIDE\_EFFECT!}}：仅用于补全和静态检查。当前二进制没有该
Python 方法，未来实现还需单独评审硬件副作用。

\textbf{参数}

\begin{longtable}[]{@{}
  >{\raggedright\arraybackslash}p{(\columnwidth - 6\tabcolsep) * \real{0.18}}
  >{\raggedright\arraybackslash}p{(\columnwidth - 6\tabcolsep) * \real{0.24}}
  >{\raggedright\arraybackslash}p{(\columnwidth - 6\tabcolsep) * \real{0.18}}
  >{\raggedright\arraybackslash}p{(\columnwidth - 6\tabcolsep) * \real{0.40}}@{}}
\toprule
名称 & Python 类型 & 默认值 & 含义 \\
\midrule
\endhead
\passthrough{\lstinline!name!} & \passthrough{\lstinline!str!} & 必填 &
SDK 对象、配置项、服务或动作的名称。允许值由调用它的具体接口定义。 对应
C++ 参数 \passthrough{\lstinline!name: const std::string \&!}。
底层为字符串；长度、编码和允许值由具体协议决定。 \\
\bottomrule
\end{longtable}

\textbf{返回值}

\begin{longtable}[]{@{}
  >{\raggedright\arraybackslash}p{(\columnwidth - 4\tabcolsep) * \real{0.18}}
  >{\raggedright\arraybackslash}p{(\columnwidth - 4\tabcolsep) * \real{0.20}}
  >{\raggedright\arraybackslash}p{(\columnwidth - 4\tabcolsep) * \real{0.62}}@{}}
\toprule
位置 & 类型 & 含义 \\
\midrule
\endhead
{[}0{]} \passthrough{\lstinline!status!} & \passthrough{\lstinline!int!}
& SDK 状态码；Unitree 示例通常以 \passthrough{\lstinline!0!}
表示成功，非零值需按具体服务错误码解释。 \\
{[}1{]} \passthrough{\lstinline!meta!} & \passthrough{\lstinline!str!} &
传给该接口的 \passthrough{\lstinline!meta!} 参数，Python 类型为
\passthrough{\lstinline!str!}。精确含义、单位、范围和枚举值需查目标型号对应头文件与协议。
对应 C++ 参数 \passthrough{\lstinline!meta: std::string \&!}。
底层为字符串；长度、编码和允许值由具体协议决定。 \\
\bottomrule
\end{longtable}

\textbf{对应 C++}

\begin{itemize}
\tightlist
\item
  类：\passthrough{\lstinline!unitree::robot::go2::ConfigClient!}
\item
  签名：\passthrough{\lstinline!Meta(const std::string \&, std::string \&)!}
\item
  绑定策略：\passthrough{\lstinline!OUTPUT\_WRAPPER!}
\item
  声明位置：\passthrough{\lstinline!include/unitree/robot/go2/config/config\_client.hpp:50!}
\end{itemize}

\textbf{说明}

\begin{itemize}
\tightlist
\item
  示例是局部调用片段；\passthrough{\lstinline!obj!}
  和参数变量表示已经按上层业务校验并准备好的对象。
\item
  C++ 的可修改输出引用不会作为 Python 入参出现，而是按 C++
  参数顺序追加到返回元组中。
\item
  该条目可以被编辑器和类型检查器识别，但当前运行时可能在属性查找阶段直接失败。
\end{itemize}

\textbf{用法}

\begin{lstlisting}[language=Python]
from typing import TYPE_CHECKING

if TYPE_CHECKING:
    def planned_meta_string(obj: ConfigClient, name: str) -> tuple[int, str]:
        return obj.meta_string(name=name)
\end{lstlisting}

\begin{quote}
{[}!CAUTION{]} 上面的 \passthrough{\lstinline!TYPE\_CHECKING!}
示例只表达计划调用形态。该分支在普通 Python
运行时为假，不代表当前扩展能执行此接口。
\end{quote}

\hypertarget{unitree-sdk2-cpp-robot-go2-configclient-subscribe-change-status-1}{%
\paragraph{\texorpdfstring{\texttt{unitree\_sdk2\_cpp.robot.go2.ConfigClient.subscribe\_change\_status}}{unitree\_sdk2\_cpp.robot.go2.ConfigClient.subscribe\_change\_status}}\label{unitree-sdk2-cpp-robot-go2-configclient-subscribe-change-status-1}}

对应 C++ SDK 操作
\passthrough{\lstinline!SubscribeChangeStatus(const std::string \&, const unitree::robot::go2::ConfigChangeStatusCallback \&)!}。上游头文件没有可直接生成的业务说明时，本参考只保证签名映射，精确语义需查目标型号协议。

\textbf{签名}

\begin{lstlisting}[language=Python]
def subscribe_change_status(self, name: str, callback: Callable[[str, str], None]) -> None
\end{lstlisting}

\textbf{可用性与安全性}

\textbf{\passthrough{\lstinline!SIGNATURE\_ONLY!} /
\passthrough{\lstinline!HARDWARE\_SIDE\_EFFECT!}}：仅用于补全和静态检查。当前二进制没有该
Python 方法，未来实现还需单独评审硬件副作用。

\textbf{参数}

\begin{longtable}[]{@{}
  >{\raggedright\arraybackslash}p{(\columnwidth - 6\tabcolsep) * \real{0.18}}
  >{\raggedright\arraybackslash}p{(\columnwidth - 6\tabcolsep) * \real{0.24}}
  >{\raggedright\arraybackslash}p{(\columnwidth - 6\tabcolsep) * \real{0.18}}
  >{\raggedright\arraybackslash}p{(\columnwidth - 6\tabcolsep) * \real{0.40}}@{}}
\toprule
名称 & Python 类型 & 默认值 & 含义 \\
\midrule
\endhead
\passthrough{\lstinline!name!} & \passthrough{\lstinline!str!} & 必填 &
SDK 对象、配置项、服务或动作的名称。允许值由调用它的具体接口定义。 对应
C++ 参数 \passthrough{\lstinline!name: const std::string \&!}。
底层为字符串；长度、编码和允许值由具体协议决定。 \\
\passthrough{\lstinline!callback!} &
\passthrough{\lstinline!Callable[[str, str], None]!} & 必填 &
收到数据或状态变化时调用的 Python
回调。回调签名必须与类型提示一致，并应快速返回。 对应 C++ 参数
\passthrough{\lstinline!callback: const unitree::robot::go2::ConfigChangeStatusCallback \&!}。 \\
\bottomrule
\end{longtable}

\textbf{返回值}

\begin{longtable}[]{@{}
  >{\raggedright\arraybackslash}p{(\columnwidth - 4\tabcolsep) * \real{0.18}}
  >{\raggedright\arraybackslash}p{(\columnwidth - 4\tabcolsep) * \real{0.20}}
  >{\raggedright\arraybackslash}p{(\columnwidth - 4\tabcolsep) * \real{0.62}}@{}}
\toprule
位置 & 类型 & 含义 \\
\midrule
\endhead
无 & \passthrough{\lstinline!None!} & 该方法不返回 Python
值；失败可能表现为异常或底层状态变化。 \\
\bottomrule
\end{longtable}

\textbf{对应 C++}

\begin{itemize}
\tightlist
\item
  类：\passthrough{\lstinline!unitree::robot::go2::ConfigClient!}
\item
  签名：\passthrough{\lstinline!SubscribeChangeStatus(const std::string \&, const unitree::robot::go2::ConfigChangeStatusCallback \&)!}
\item
  绑定策略：\passthrough{\lstinline!CALLBACK\_MANUAL!}
\item
  声明位置：\passthrough{\lstinline!include/unitree/robot/go2/config/config\_client.hpp:52!}
\end{itemize}

\textbf{说明}

\begin{itemize}
\tightlist
\item
  示例是局部调用片段；\passthrough{\lstinline!obj!}
  和参数变量表示已经按上层业务校验并准备好的对象。
\item
  该条目可以被编辑器和类型检查器识别，但当前运行时可能在属性查找阶段直接失败。
\item
  回调可能由 SDK
  工作线程触发；回调应快速返回、捕获异常，并通过线程安全队列移交耗时工作。
\end{itemize}

\textbf{用法}

\begin{lstlisting}[language=Python]
from typing import TYPE_CHECKING

if TYPE_CHECKING:
    def planned_subscribe_change_status(obj: ConfigClient, name: str, callback: Callable[[str, str], None]) -> None:
        obj.subscribe_change_status(name=name, callback=callback)
\end{lstlisting}

\begin{quote}
{[}!CAUTION{]} 上面的 \passthrough{\lstinline!TYPE\_CHECKING!}
示例只表达计划调用形态。该分支在普通 Python
运行时为假，不代表当前扩展能执行此接口。
\end{quote}

\hypertarget{unitree-sdk2-cpp-robot-go2-configdelparameter}{%
\subsubsection{\texorpdfstring{\texttt{unitree\_sdk2\_cpp.robot.go2.ConfigDelParameter}}{unitree\_sdk2\_cpp.robot.go2.ConfigDelParameter}}\label{unitree-sdk2-cpp-robot-go2-configdelparameter}}

SDK 请求参数值类型；当前可能仅提供设计期签名。

\textbf{导入}

\begin{lstlisting}[language=Python]
from unitree_sdk2_cpp.robot.go2 import ConfigDelParameter
\end{lstlisting}

\textbf{基类}：\passthrough{\lstinline!object!}

\textbf{公开属性}

\begin{longtable}[]{@{}
  >{\raggedright\arraybackslash}p{(\columnwidth - 6\tabcolsep) * \real{0.18}}
  >{\raggedright\arraybackslash}p{(\columnwidth - 6\tabcolsep) * \real{0.24}}
  >{\raggedright\arraybackslash}p{(\columnwidth - 6\tabcolsep) * \real{0.18}}
  >{\raggedright\arraybackslash}p{(\columnwidth - 6\tabcolsep) * \real{0.40}}@{}}
\toprule
名称 & Python 类型 & 含义 & 用法 \\
\midrule
\endhead
\passthrough{\lstinline!name!} & \passthrough{\lstinline!str!} & SDK
对象、配置项、服务或动作的名称。允许值由调用它的具体接口定义。 &
\passthrough{\lstinline!value = obj.name!} /
\passthrough{\lstinline!obj.name = value!} \\
\bottomrule
\end{longtable}

\textbf{方法索引}

\begin{longtable}[]{@{}
  >{\raggedright\arraybackslash}p{(\columnwidth - 6\tabcolsep) * \real{0.18}}
  >{\raggedright\arraybackslash}p{(\columnwidth - 6\tabcolsep) * \real{0.24}}
  >{\raggedright\arraybackslash}p{(\columnwidth - 6\tabcolsep) * \real{0.18}}
  >{\raggedright\arraybackslash}p{(\columnwidth - 6\tabcolsep) * \real{0.40}}@{}}
\toprule
方法 & Python 签名 & 状态 & 安全分类 \\
\midrule
\endhead
\protect\hyperlink{unitree-sdk2-cpp-robot-go2-configdelparameter-dunder-init-1}{\passthrough{\lstinline!\_\_init\_\_!}}
& \passthrough{\lstinline!def \_\_init\_\_(self) -> None!} &
\passthrough{\lstinline!SIGNATURE\_ONLY!} &
\passthrough{\lstinline!CONSTRUCTION!} \\
\protect\hyperlink{unitree-sdk2-cpp-robot-go2-configdelparameter-from-json-1}{\passthrough{\lstinline!from\_json!}}
&
\passthrough{\lstinline!def from\_json(self, json: dict[str, Any]) -> None!}
& \passthrough{\lstinline!SIGNATURE\_ONLY!} &
\passthrough{\lstinline!UNCLASSIFIED!} \\
\protect\hyperlink{unitree-sdk2-cpp-robot-go2-configdelparameter-to-json-1}{\passthrough{\lstinline!to\_json!}}
&
\passthrough{\lstinline!def to\_json(self, json: dict[str, Any]) -> None!}
& \passthrough{\lstinline!SIGNATURE\_ONLY!} &
\passthrough{\lstinline!UNCLASSIFIED!} \\
\bottomrule
\end{longtable}

\hypertarget{unitree-sdk2-cpp-robot-go2-configdelparameter-dunder-init-1}{%
\paragraph{\texorpdfstring{\texttt{unitree\_sdk2\_cpp.robot.go2.ConfigDelParameter.\_\_init\_\_}}{unitree\_sdk2\_cpp.robot.go2.ConfigDelParameter.\_\_init\_\_}}\label{unitree-sdk2-cpp-robot-go2-configdelparameter-dunder-init-1}}

初始化 \passthrough{\lstinline!ConfigDelParameter!}
实例。是否能实际构造取决于下方可用性状态。

\textbf{签名}

\begin{lstlisting}[language=Python]
def __init__(self) -> None
\end{lstlisting}

\textbf{可用性与安全性}

\textbf{\passthrough{\lstinline!SIGNATURE\_ONLY!} /
\passthrough{\lstinline!CONSTRUCTION!}}：当前只有设计期签名，不能假设已存在于运行时扩展。

\textbf{参数}

无显式参数。实例方法中的 \passthrough{\lstinline!self!} 由 Python
自动传入。

\textbf{返回值}

\begin{longtable}[]{@{}
  >{\raggedright\arraybackslash}p{(\columnwidth - 4\tabcolsep) * \real{0.18}}
  >{\raggedright\arraybackslash}p{(\columnwidth - 4\tabcolsep) * \real{0.20}}
  >{\raggedright\arraybackslash}p{(\columnwidth - 4\tabcolsep) * \real{0.62}}@{}}
\toprule
位置 & 类型 & 含义 \\
\midrule
\endhead
无 & \passthrough{\lstinline!None!} &
\passthrough{\lstinline!\_\_init\_\_!}
本身不返回值；调用类对象时，在构造函数可用的前提下得到该类实例。 \\
\bottomrule
\end{longtable}

\textbf{对应 C++}

\begin{itemize}
\tightlist
\item
  类：\passthrough{\lstinline!unitree::robot::go2::ConfigDelParameter!}
\item
  签名：\passthrough{\lstinline!ConfigDelParameter()!}
\item
  绑定策略：\passthrough{\lstinline!未单独标注!}
\item
  声明位置：\passthrough{\lstinline!include/unitree/robot/go2/config/config\_api.hpp:140!}
\end{itemize}

\textbf{说明}

\begin{itemize}
\tightlist
\item
  示例是局部调用片段；\passthrough{\lstinline!obj!}
  和参数变量表示已经按上层业务校验并准备好的对象。
\item
  该条目可以被编辑器和类型检查器识别，但当前运行时可能在属性查找阶段直接失败。
\end{itemize}

\textbf{用法}

\begin{lstlisting}[language=Python]
from typing import TYPE_CHECKING

if TYPE_CHECKING:
    def planned_construct() -> ConfigDelParameter:
        return ConfigDelParameter()
\end{lstlisting}

\begin{quote}
{[}!CAUTION{]} 上面的 \passthrough{\lstinline!TYPE\_CHECKING!}
示例只表达计划调用形态。该分支在普通 Python
运行时为假，不代表当前扩展能执行此接口。
\end{quote}

\hypertarget{unitree-sdk2-cpp-robot-go2-configdelparameter-from-json-1}{%
\paragraph{\texorpdfstring{\texttt{unitree\_sdk2\_cpp.robot.go2.ConfigDelParameter.from\_json}}{unitree\_sdk2\_cpp.robot.go2.ConfigDelParameter.from\_json}}\label{unitree-sdk2-cpp-robot-go2-configdelparameter-from-json-1}}

计划从 JSON 风格字典读取字段并更新当前 SDK 值对象。

\textbf{签名}

\begin{lstlisting}[language=Python]
def from_json(self, json: dict[str, Any]) -> None
\end{lstlisting}

\textbf{可用性与安全性}

\textbf{\passthrough{\lstinline!SIGNATURE\_ONLY!} /
\passthrough{\lstinline!UNCLASSIFIED!}}：当前只有设计期签名，不能假设已存在于运行时扩展。

\textbf{参数}

\begin{longtable}[]{@{}
  >{\raggedright\arraybackslash}p{(\columnwidth - 6\tabcolsep) * \real{0.18}}
  >{\raggedright\arraybackslash}p{(\columnwidth - 6\tabcolsep) * \real{0.24}}
  >{\raggedright\arraybackslash}p{(\columnwidth - 6\tabcolsep) * \real{0.18}}
  >{\raggedright\arraybackslash}p{(\columnwidth - 6\tabcolsep) * \real{0.40}}@{}}
\toprule
名称 & Python 类型 & 默认值 & 含义 \\
\midrule
\endhead
\passthrough{\lstinline!json!} &
\passthrough{\lstinline!dict[str, Any]!} & 必填 & JSON
风格字典，键和值必须满足对应 C++ \passthrough{\lstinline!JsonMap!}
数据结构的约束。 对应 C++ 参数
\passthrough{\lstinline!json: common::JsonMap \&!}。 \\
\bottomrule
\end{longtable}

\textbf{返回值}

\begin{longtable}[]{@{}
  >{\raggedright\arraybackslash}p{(\columnwidth - 4\tabcolsep) * \real{0.18}}
  >{\raggedright\arraybackslash}p{(\columnwidth - 4\tabcolsep) * \real{0.20}}
  >{\raggedright\arraybackslash}p{(\columnwidth - 4\tabcolsep) * \real{0.62}}@{}}
\toprule
位置 & 类型 & 含义 \\
\midrule
\endhead
无 & \passthrough{\lstinline!None!} & 该方法不返回 Python
值；失败可能表现为异常或底层状态变化。 \\
\bottomrule
\end{longtable}

\textbf{对应 C++}

\begin{itemize}
\tightlist
\item
  类：\passthrough{\lstinline!unitree::robot::go2::ConfigDelParameter!}
\item
  签名：\passthrough{\lstinline!fromJson(common::JsonMap \&)!}
\item
  绑定策略：\passthrough{\lstinline!SIGNATURE\_PREVIEW!}
\item
  声明位置：\passthrough{\lstinline!include/unitree/robot/go2/config/config\_api.hpp:146!}
\end{itemize}

\textbf{说明}

\begin{itemize}
\tightlist
\item
  示例是局部调用片段；\passthrough{\lstinline!obj!}
  和参数变量表示已经按上层业务校验并准备好的对象。
\item
  该条目可以被编辑器和类型检查器识别，但当前运行时可能在属性查找阶段直接失败。
\end{itemize}

\textbf{用法}

\begin{lstlisting}[language=Python]
from typing import TYPE_CHECKING

if TYPE_CHECKING:
    def planned_from_json(obj: ConfigDelParameter, json: dict[str, Any]) -> None:
        obj.from_json(json=json)
\end{lstlisting}

\begin{quote}
{[}!CAUTION{]} 上面的 \passthrough{\lstinline!TYPE\_CHECKING!}
示例只表达计划调用形态。该分支在普通 Python
运行时为假，不代表当前扩展能执行此接口。
\end{quote}

\hypertarget{unitree-sdk2-cpp-robot-go2-configdelparameter-to-json-1}{%
\paragraph{\texorpdfstring{\texttt{unitree\_sdk2\_cpp.robot.go2.ConfigDelParameter.to\_json}}{unitree\_sdk2\_cpp.robot.go2.ConfigDelParameter.to\_json}}\label{unitree-sdk2-cpp-robot-go2-configdelparameter-to-json-1}}

计划把当前 SDK 值对象写入 JSON 风格字典。

\textbf{签名}

\begin{lstlisting}[language=Python]
def to_json(self, json: dict[str, Any]) -> None
\end{lstlisting}

\textbf{可用性与安全性}

\textbf{\passthrough{\lstinline!SIGNATURE\_ONLY!} /
\passthrough{\lstinline!UNCLASSIFIED!}}：当前只有设计期签名，不能假设已存在于运行时扩展。

\textbf{参数}

\begin{longtable}[]{@{}
  >{\raggedright\arraybackslash}p{(\columnwidth - 6\tabcolsep) * \real{0.18}}
  >{\raggedright\arraybackslash}p{(\columnwidth - 6\tabcolsep) * \real{0.24}}
  >{\raggedright\arraybackslash}p{(\columnwidth - 6\tabcolsep) * \real{0.18}}
  >{\raggedright\arraybackslash}p{(\columnwidth - 6\tabcolsep) * \real{0.40}}@{}}
\toprule
名称 & Python 类型 & 默认值 & 含义 \\
\midrule
\endhead
\passthrough{\lstinline!json!} &
\passthrough{\lstinline!dict[str, Any]!} & 必填 & JSON
风格字典，键和值必须满足对应 C++ \passthrough{\lstinline!JsonMap!}
数据结构的约束。 对应 C++ 参数
\passthrough{\lstinline!json: common::JsonMap \&!}。 \\
\bottomrule
\end{longtable}

\textbf{返回值}

\begin{longtable}[]{@{}
  >{\raggedright\arraybackslash}p{(\columnwidth - 4\tabcolsep) * \real{0.18}}
  >{\raggedright\arraybackslash}p{(\columnwidth - 4\tabcolsep) * \real{0.20}}
  >{\raggedright\arraybackslash}p{(\columnwidth - 4\tabcolsep) * \real{0.62}}@{}}
\toprule
位置 & 类型 & 含义 \\
\midrule
\endhead
无 & \passthrough{\lstinline!None!} & 该方法不返回 Python
值；失败可能表现为异常或底层状态变化。 \\
\bottomrule
\end{longtable}

\textbf{对应 C++}

\begin{itemize}
\tightlist
\item
  类：\passthrough{\lstinline!unitree::robot::go2::ConfigDelParameter!}
\item
  签名：\passthrough{\lstinline!toJson(common::JsonMap \&) const!}
\item
  绑定策略：\passthrough{\lstinline!SIGNATURE\_PREVIEW!}
\item
  声明位置：\passthrough{\lstinline!include/unitree/robot/go2/config/config\_api.hpp:151!}
\end{itemize}

\textbf{说明}

\begin{itemize}
\tightlist
\item
  示例是局部调用片段；\passthrough{\lstinline!obj!}
  和参数变量表示已经按上层业务校验并准备好的对象。
\item
  该条目可以被编辑器和类型检查器识别，但当前运行时可能在属性查找阶段直接失败。
\end{itemize}

\textbf{用法}

\begin{lstlisting}[language=Python]
from typing import TYPE_CHECKING

if TYPE_CHECKING:
    def planned_to_json(obj: ConfigDelParameter, json: dict[str, Any]) -> None:
        obj.to_json(json=json)
\end{lstlisting}

\begin{quote}
{[}!CAUTION{]} 上面的 \passthrough{\lstinline!TYPE\_CHECKING!}
示例只表达计划调用形态。该分支在普通 Python
运行时为假，不代表当前扩展能执行此接口。
\end{quote}

\hypertarget{unitree-sdk2-cpp-robot-go2-configgetdata}{%
\subsubsection{\texorpdfstring{\texttt{unitree\_sdk2\_cpp.robot.go2.ConfigGetData}}{unitree\_sdk2\_cpp.robot.go2.ConfigGetData}}\label{unitree-sdk2-cpp-robot-go2-configgetdata}}

SDK 数据值类型；公开字段和转换方法见下文。

\textbf{导入}

\begin{lstlisting}[language=Python]
from unitree_sdk2_cpp.robot.go2 import ConfigGetData
\end{lstlisting}

\textbf{基类}：\passthrough{\lstinline!object!}

\textbf{公开属性}

\begin{longtable}[]{@{}
  >{\raggedright\arraybackslash}p{(\columnwidth - 6\tabcolsep) * \real{0.18}}
  >{\raggedright\arraybackslash}p{(\columnwidth - 6\tabcolsep) * \real{0.24}}
  >{\raggedright\arraybackslash}p{(\columnwidth - 6\tabcolsep) * \real{0.18}}
  >{\raggedright\arraybackslash}p{(\columnwidth - 6\tabcolsep) * \real{0.40}}@{}}
\toprule
名称 & Python 类型 & 含义 & 用法 \\
\midrule
\endhead
\passthrough{\lstinline!content!} & \passthrough{\lstinline!str!} &
配置、请求或序列化内容。具体格式由对应服务协议定义。 &
\passthrough{\lstinline!value = obj.content!} /
\passthrough{\lstinline!obj.content = value!} \\
\bottomrule
\end{longtable}

\textbf{方法索引}

\begin{longtable}[]{@{}
  >{\raggedright\arraybackslash}p{(\columnwidth - 6\tabcolsep) * \real{0.18}}
  >{\raggedright\arraybackslash}p{(\columnwidth - 6\tabcolsep) * \real{0.24}}
  >{\raggedright\arraybackslash}p{(\columnwidth - 6\tabcolsep) * \real{0.18}}
  >{\raggedright\arraybackslash}p{(\columnwidth - 6\tabcolsep) * \real{0.40}}@{}}
\toprule
方法 & Python 签名 & 状态 & 安全分类 \\
\midrule
\endhead
\protect\hyperlink{unitree-sdk2-cpp-robot-go2-configgetdata-dunder-init-1}{\passthrough{\lstinline!\_\_init\_\_!}}
& \passthrough{\lstinline!def \_\_init\_\_(self) -> None!} &
\passthrough{\lstinline!SIGNATURE\_ONLY!} &
\passthrough{\lstinline!CONSTRUCTION!} \\
\protect\hyperlink{unitree-sdk2-cpp-robot-go2-configgetdata-from-json-1}{\passthrough{\lstinline!from\_json!}}
&
\passthrough{\lstinline!def from\_json(self, json: dict[str, Any]) -> None!}
& \passthrough{\lstinline!SIGNATURE\_ONLY!} &
\passthrough{\lstinline!UNCLASSIFIED!} \\
\protect\hyperlink{unitree-sdk2-cpp-robot-go2-configgetdata-to-json-1}{\passthrough{\lstinline!to\_json!}}
&
\passthrough{\lstinline!def to\_json(self, json: dict[str, Any]) -> None!}
& \passthrough{\lstinline!SIGNATURE\_ONLY!} &
\passthrough{\lstinline!UNCLASSIFIED!} \\
\bottomrule
\end{longtable}

\hypertarget{unitree-sdk2-cpp-robot-go2-configgetdata-dunder-init-1}{%
\paragraph{\texorpdfstring{\texttt{unitree\_sdk2\_cpp.robot.go2.ConfigGetData.\_\_init\_\_}}{unitree\_sdk2\_cpp.robot.go2.ConfigGetData.\_\_init\_\_}}\label{unitree-sdk2-cpp-robot-go2-configgetdata-dunder-init-1}}

初始化 \passthrough{\lstinline!ConfigGetData!}
实例。是否能实际构造取决于下方可用性状态。

\textbf{签名}

\begin{lstlisting}[language=Python]
def __init__(self) -> None
\end{lstlisting}

\textbf{可用性与安全性}

\textbf{\passthrough{\lstinline!SIGNATURE\_ONLY!} /
\passthrough{\lstinline!CONSTRUCTION!}}：当前只有设计期签名，不能假设已存在于运行时扩展。

\textbf{参数}

无显式参数。实例方法中的 \passthrough{\lstinline!self!} 由 Python
自动传入。

\textbf{返回值}

\begin{longtable}[]{@{}
  >{\raggedright\arraybackslash}p{(\columnwidth - 4\tabcolsep) * \real{0.18}}
  >{\raggedright\arraybackslash}p{(\columnwidth - 4\tabcolsep) * \real{0.20}}
  >{\raggedright\arraybackslash}p{(\columnwidth - 4\tabcolsep) * \real{0.62}}@{}}
\toprule
位置 & 类型 & 含义 \\
\midrule
\endhead
无 & \passthrough{\lstinline!None!} &
\passthrough{\lstinline!\_\_init\_\_!}
本身不返回值；调用类对象时，在构造函数可用的前提下得到该类实例。 \\
\bottomrule
\end{longtable}

\textbf{对应 C++}

\begin{itemize}
\tightlist
\item
  类：\passthrough{\lstinline!unitree::robot::go2::ConfigGetData!}
\item
  签名：\passthrough{\lstinline!ConfigGetData()!}
\item
  绑定策略：\passthrough{\lstinline!未单独标注!}
\item
  声明位置：\passthrough{\lstinline!include/unitree/robot/go2/config/config\_api.hpp:117!}
\end{itemize}

\textbf{说明}

\begin{itemize}
\tightlist
\item
  示例是局部调用片段；\passthrough{\lstinline!obj!}
  和参数变量表示已经按上层业务校验并准备好的对象。
\item
  该条目可以被编辑器和类型检查器识别，但当前运行时可能在属性查找阶段直接失败。
\end{itemize}

\textbf{用法}

\begin{lstlisting}[language=Python]
from typing import TYPE_CHECKING

if TYPE_CHECKING:
    def planned_construct() -> ConfigGetData:
        return ConfigGetData()
\end{lstlisting}

\begin{quote}
{[}!CAUTION{]} 上面的 \passthrough{\lstinline!TYPE\_CHECKING!}
示例只表达计划调用形态。该分支在普通 Python
运行时为假，不代表当前扩展能执行此接口。
\end{quote}

\hypertarget{unitree-sdk2-cpp-robot-go2-configgetdata-from-json-1}{%
\paragraph{\texorpdfstring{\texttt{unitree\_sdk2\_cpp.robot.go2.ConfigGetData.from\_json}}{unitree\_sdk2\_cpp.robot.go2.ConfigGetData.from\_json}}\label{unitree-sdk2-cpp-robot-go2-configgetdata-from-json-1}}

计划从 JSON 风格字典读取字段并更新当前 SDK 值对象。

\textbf{签名}

\begin{lstlisting}[language=Python]
def from_json(self, json: dict[str, Any]) -> None
\end{lstlisting}

\textbf{可用性与安全性}

\textbf{\passthrough{\lstinline!SIGNATURE\_ONLY!} /
\passthrough{\lstinline!UNCLASSIFIED!}}：当前只有设计期签名，不能假设已存在于运行时扩展。

\textbf{参数}

\begin{longtable}[]{@{}
  >{\raggedright\arraybackslash}p{(\columnwidth - 6\tabcolsep) * \real{0.18}}
  >{\raggedright\arraybackslash}p{(\columnwidth - 6\tabcolsep) * \real{0.24}}
  >{\raggedright\arraybackslash}p{(\columnwidth - 6\tabcolsep) * \real{0.18}}
  >{\raggedright\arraybackslash}p{(\columnwidth - 6\tabcolsep) * \real{0.40}}@{}}
\toprule
名称 & Python 类型 & 默认值 & 含义 \\
\midrule
\endhead
\passthrough{\lstinline!json!} &
\passthrough{\lstinline!dict[str, Any]!} & 必填 & JSON
风格字典，键和值必须满足对应 C++ \passthrough{\lstinline!JsonMap!}
数据结构的约束。 对应 C++ 参数
\passthrough{\lstinline!json: common::JsonMap \&!}。 \\
\bottomrule
\end{longtable}

\textbf{返回值}

\begin{longtable}[]{@{}
  >{\raggedright\arraybackslash}p{(\columnwidth - 4\tabcolsep) * \real{0.18}}
  >{\raggedright\arraybackslash}p{(\columnwidth - 4\tabcolsep) * \real{0.20}}
  >{\raggedright\arraybackslash}p{(\columnwidth - 4\tabcolsep) * \real{0.62}}@{}}
\toprule
位置 & 类型 & 含义 \\
\midrule
\endhead
无 & \passthrough{\lstinline!None!} & 该方法不返回 Python
值；失败可能表现为异常或底层状态变化。 \\
\bottomrule
\end{longtable}

\textbf{对应 C++}

\begin{itemize}
\tightlist
\item
  类：\passthrough{\lstinline!unitree::robot::go2::ConfigGetData!}
\item
  签名：\passthrough{\lstinline!fromJson(common::JsonMap \&)!}
\item
  绑定策略：\passthrough{\lstinline!SIGNATURE\_PREVIEW!}
\item
  声明位置：\passthrough{\lstinline!include/unitree/robot/go2/config/config\_api.hpp:123!}
\end{itemize}

\textbf{说明}

\begin{itemize}
\tightlist
\item
  示例是局部调用片段；\passthrough{\lstinline!obj!}
  和参数变量表示已经按上层业务校验并准备好的对象。
\item
  该条目可以被编辑器和类型检查器识别，但当前运行时可能在属性查找阶段直接失败。
\end{itemize}

\textbf{用法}

\begin{lstlisting}[language=Python]
from typing import TYPE_CHECKING

if TYPE_CHECKING:
    def planned_from_json(obj: ConfigGetData, json: dict[str, Any]) -> None:
        obj.from_json(json=json)
\end{lstlisting}

\begin{quote}
{[}!CAUTION{]} 上面的 \passthrough{\lstinline!TYPE\_CHECKING!}
示例只表达计划调用形态。该分支在普通 Python
运行时为假，不代表当前扩展能执行此接口。
\end{quote}

\hypertarget{unitree-sdk2-cpp-robot-go2-configgetdata-to-json-1}{%
\paragraph{\texorpdfstring{\texttt{unitree\_sdk2\_cpp.robot.go2.ConfigGetData.to\_json}}{unitree\_sdk2\_cpp.robot.go2.ConfigGetData.to\_json}}\label{unitree-sdk2-cpp-robot-go2-configgetdata-to-json-1}}

计划把当前 SDK 值对象写入 JSON 风格字典。

\textbf{签名}

\begin{lstlisting}[language=Python]
def to_json(self, json: dict[str, Any]) -> None
\end{lstlisting}

\textbf{可用性与安全性}

\textbf{\passthrough{\lstinline!SIGNATURE\_ONLY!} /
\passthrough{\lstinline!UNCLASSIFIED!}}：当前只有设计期签名，不能假设已存在于运行时扩展。

\textbf{参数}

\begin{longtable}[]{@{}
  >{\raggedright\arraybackslash}p{(\columnwidth - 6\tabcolsep) * \real{0.18}}
  >{\raggedright\arraybackslash}p{(\columnwidth - 6\tabcolsep) * \real{0.24}}
  >{\raggedright\arraybackslash}p{(\columnwidth - 6\tabcolsep) * \real{0.18}}
  >{\raggedright\arraybackslash}p{(\columnwidth - 6\tabcolsep) * \real{0.40}}@{}}
\toprule
名称 & Python 类型 & 默认值 & 含义 \\
\midrule
\endhead
\passthrough{\lstinline!json!} &
\passthrough{\lstinline!dict[str, Any]!} & 必填 & JSON
风格字典，键和值必须满足对应 C++ \passthrough{\lstinline!JsonMap!}
数据结构的约束。 对应 C++ 参数
\passthrough{\lstinline!json: common::JsonMap \&!}。 \\
\bottomrule
\end{longtable}

\textbf{返回值}

\begin{longtable}[]{@{}
  >{\raggedright\arraybackslash}p{(\columnwidth - 4\tabcolsep) * \real{0.18}}
  >{\raggedright\arraybackslash}p{(\columnwidth - 4\tabcolsep) * \real{0.20}}
  >{\raggedright\arraybackslash}p{(\columnwidth - 4\tabcolsep) * \real{0.62}}@{}}
\toprule
位置 & 类型 & 含义 \\
\midrule
\endhead
无 & \passthrough{\lstinline!None!} & 该方法不返回 Python
值；失败可能表现为异常或底层状态变化。 \\
\bottomrule
\end{longtable}

\textbf{对应 C++}

\begin{itemize}
\tightlist
\item
  类：\passthrough{\lstinline!unitree::robot::go2::ConfigGetData!}
\item
  签名：\passthrough{\lstinline!toJson(common::JsonMap \&) const!}
\item
  绑定策略：\passthrough{\lstinline!SIGNATURE\_PREVIEW!}
\item
  声明位置：\passthrough{\lstinline!include/unitree/robot/go2/config/config\_api.hpp:128!}
\end{itemize}

\textbf{说明}

\begin{itemize}
\tightlist
\item
  示例是局部调用片段；\passthrough{\lstinline!obj!}
  和参数变量表示已经按上层业务校验并准备好的对象。
\item
  该条目可以被编辑器和类型检查器识别，但当前运行时可能在属性查找阶段直接失败。
\end{itemize}

\textbf{用法}

\begin{lstlisting}[language=Python]
from typing import TYPE_CHECKING

if TYPE_CHECKING:
    def planned_to_json(obj: ConfigGetData, json: dict[str, Any]) -> None:
        obj.to_json(json=json)
\end{lstlisting}

\begin{quote}
{[}!CAUTION{]} 上面的 \passthrough{\lstinline!TYPE\_CHECKING!}
示例只表达计划调用形态。该分支在普通 Python
运行时为假，不代表当前扩展能执行此接口。
\end{quote}

\hypertarget{unitree-sdk2-cpp-robot-go2-configgetparameter}{%
\subsubsection{\texorpdfstring{\texttt{unitree\_sdk2\_cpp.robot.go2.ConfigGetParameter}}{unitree\_sdk2\_cpp.robot.go2.ConfigGetParameter}}\label{unitree-sdk2-cpp-robot-go2-configgetparameter}}

SDK 请求参数值类型；当前可能仅提供设计期签名。

\textbf{导入}

\begin{lstlisting}[language=Python]
from unitree_sdk2_cpp.robot.go2 import ConfigGetParameter
\end{lstlisting}

\textbf{基类}：\passthrough{\lstinline!object!}

\textbf{公开属性}

\begin{longtable}[]{@{}
  >{\raggedright\arraybackslash}p{(\columnwidth - 6\tabcolsep) * \real{0.18}}
  >{\raggedright\arraybackslash}p{(\columnwidth - 6\tabcolsep) * \real{0.24}}
  >{\raggedright\arraybackslash}p{(\columnwidth - 6\tabcolsep) * \real{0.18}}
  >{\raggedright\arraybackslash}p{(\columnwidth - 6\tabcolsep) * \real{0.40}}@{}}
\toprule
名称 & Python 类型 & 含义 & 用法 \\
\midrule
\endhead
\passthrough{\lstinline!name!} & \passthrough{\lstinline!str!} & SDK
对象、配置项、服务或动作的名称。允许值由调用它的具体接口定义。 &
\passthrough{\lstinline!value = obj.name!} /
\passthrough{\lstinline!obj.name = value!} \\
\bottomrule
\end{longtable}

\textbf{方法索引}

\begin{longtable}[]{@{}
  >{\raggedright\arraybackslash}p{(\columnwidth - 6\tabcolsep) * \real{0.18}}
  >{\raggedright\arraybackslash}p{(\columnwidth - 6\tabcolsep) * \real{0.24}}
  >{\raggedright\arraybackslash}p{(\columnwidth - 6\tabcolsep) * \real{0.18}}
  >{\raggedright\arraybackslash}p{(\columnwidth - 6\tabcolsep) * \real{0.40}}@{}}
\toprule
方法 & Python 签名 & 状态 & 安全分类 \\
\midrule
\endhead
\protect\hyperlink{unitree-sdk2-cpp-robot-go2-configgetparameter-dunder-init-1}{\passthrough{\lstinline!\_\_init\_\_!}}
& \passthrough{\lstinline!def \_\_init\_\_(self) -> None!} &
\passthrough{\lstinline!SIGNATURE\_ONLY!} &
\passthrough{\lstinline!CONSTRUCTION!} \\
\protect\hyperlink{unitree-sdk2-cpp-robot-go2-configgetparameter-from-json-1}{\passthrough{\lstinline!from\_json!}}
&
\passthrough{\lstinline!def from\_json(self, json: dict[str, Any]) -> None!}
& \passthrough{\lstinline!SIGNATURE\_ONLY!} &
\passthrough{\lstinline!UNCLASSIFIED!} \\
\protect\hyperlink{unitree-sdk2-cpp-robot-go2-configgetparameter-to-json-1}{\passthrough{\lstinline!to\_json!}}
&
\passthrough{\lstinline!def to\_json(self, json: dict[str, Any]) -> None!}
& \passthrough{\lstinline!SIGNATURE\_ONLY!} &
\passthrough{\lstinline!UNCLASSIFIED!} \\
\bottomrule
\end{longtable}

\hypertarget{unitree-sdk2-cpp-robot-go2-configgetparameter-dunder-init-1}{%
\paragraph{\texorpdfstring{\texttt{unitree\_sdk2\_cpp.robot.go2.ConfigGetParameter.\_\_init\_\_}}{unitree\_sdk2\_cpp.robot.go2.ConfigGetParameter.\_\_init\_\_}}\label{unitree-sdk2-cpp-robot-go2-configgetparameter-dunder-init-1}}

初始化 \passthrough{\lstinline!ConfigGetParameter!}
实例。是否能实际构造取决于下方可用性状态。

\textbf{签名}

\begin{lstlisting}[language=Python]
def __init__(self) -> None
\end{lstlisting}

\textbf{可用性与安全性}

\textbf{\passthrough{\lstinline!SIGNATURE\_ONLY!} /
\passthrough{\lstinline!CONSTRUCTION!}}：当前只有设计期签名，不能假设已存在于运行时扩展。

\textbf{参数}

无显式参数。实例方法中的 \passthrough{\lstinline!self!} 由 Python
自动传入。

\textbf{返回值}

\begin{longtable}[]{@{}
  >{\raggedright\arraybackslash}p{(\columnwidth - 4\tabcolsep) * \real{0.18}}
  >{\raggedright\arraybackslash}p{(\columnwidth - 4\tabcolsep) * \real{0.20}}
  >{\raggedright\arraybackslash}p{(\columnwidth - 4\tabcolsep) * \real{0.62}}@{}}
\toprule
位置 & 类型 & 含义 \\
\midrule
\endhead
无 & \passthrough{\lstinline!None!} &
\passthrough{\lstinline!\_\_init\_\_!}
本身不返回值；调用类对象时，在构造函数可用的前提下得到该类实例。 \\
\bottomrule
\end{longtable}

\textbf{对应 C++}

\begin{itemize}
\tightlist
\item
  类：\passthrough{\lstinline!unitree::robot::go2::ConfigGetParameter!}
\item
  签名：\passthrough{\lstinline!ConfigGetParameter()!}
\item
  绑定策略：\passthrough{\lstinline!未单独标注!}
\item
  声明位置：\passthrough{\lstinline!include/unitree/robot/go2/config/config\_api.hpp:94!}
\end{itemize}

\textbf{说明}

\begin{itemize}
\tightlist
\item
  示例是局部调用片段；\passthrough{\lstinline!obj!}
  和参数变量表示已经按上层业务校验并准备好的对象。
\item
  该条目可以被编辑器和类型检查器识别，但当前运行时可能在属性查找阶段直接失败。
\end{itemize}

\textbf{用法}

\begin{lstlisting}[language=Python]
from typing import TYPE_CHECKING

if TYPE_CHECKING:
    def planned_construct() -> ConfigGetParameter:
        return ConfigGetParameter()
\end{lstlisting}

\begin{quote}
{[}!CAUTION{]} 上面的 \passthrough{\lstinline!TYPE\_CHECKING!}
示例只表达计划调用形态。该分支在普通 Python
运行时为假，不代表当前扩展能执行此接口。
\end{quote}

\hypertarget{unitree-sdk2-cpp-robot-go2-configgetparameter-from-json-1}{%
\paragraph{\texorpdfstring{\texttt{unitree\_sdk2\_cpp.robot.go2.ConfigGetParameter.from\_json}}{unitree\_sdk2\_cpp.robot.go2.ConfigGetParameter.from\_json}}\label{unitree-sdk2-cpp-robot-go2-configgetparameter-from-json-1}}

计划从 JSON 风格字典读取字段并更新当前 SDK 值对象。

\textbf{签名}

\begin{lstlisting}[language=Python]
def from_json(self, json: dict[str, Any]) -> None
\end{lstlisting}

\textbf{可用性与安全性}

\textbf{\passthrough{\lstinline!SIGNATURE\_ONLY!} /
\passthrough{\lstinline!UNCLASSIFIED!}}：当前只有设计期签名，不能假设已存在于运行时扩展。

\textbf{参数}

\begin{longtable}[]{@{}
  >{\raggedright\arraybackslash}p{(\columnwidth - 6\tabcolsep) * \real{0.18}}
  >{\raggedright\arraybackslash}p{(\columnwidth - 6\tabcolsep) * \real{0.24}}
  >{\raggedright\arraybackslash}p{(\columnwidth - 6\tabcolsep) * \real{0.18}}
  >{\raggedright\arraybackslash}p{(\columnwidth - 6\tabcolsep) * \real{0.40}}@{}}
\toprule
名称 & Python 类型 & 默认值 & 含义 \\
\midrule
\endhead
\passthrough{\lstinline!json!} &
\passthrough{\lstinline!dict[str, Any]!} & 必填 & JSON
风格字典，键和值必须满足对应 C++ \passthrough{\lstinline!JsonMap!}
数据结构的约束。 对应 C++ 参数
\passthrough{\lstinline!json: common::JsonMap \&!}。 \\
\bottomrule
\end{longtable}

\textbf{返回值}

\begin{longtable}[]{@{}
  >{\raggedright\arraybackslash}p{(\columnwidth - 4\tabcolsep) * \real{0.18}}
  >{\raggedright\arraybackslash}p{(\columnwidth - 4\tabcolsep) * \real{0.20}}
  >{\raggedright\arraybackslash}p{(\columnwidth - 4\tabcolsep) * \real{0.62}}@{}}
\toprule
位置 & 类型 & 含义 \\
\midrule
\endhead
无 & \passthrough{\lstinline!None!} & 该方法不返回 Python
值；失败可能表现为异常或底层状态变化。 \\
\bottomrule
\end{longtable}

\textbf{对应 C++}

\begin{itemize}
\tightlist
\item
  类：\passthrough{\lstinline!unitree::robot::go2::ConfigGetParameter!}
\item
  签名：\passthrough{\lstinline!fromJson(common::JsonMap \&)!}
\item
  绑定策略：\passthrough{\lstinline!SIGNATURE\_PREVIEW!}
\item
  声明位置：\passthrough{\lstinline!include/unitree/robot/go2/config/config\_api.hpp:100!}
\end{itemize}

\textbf{说明}

\begin{itemize}
\tightlist
\item
  示例是局部调用片段；\passthrough{\lstinline!obj!}
  和参数变量表示已经按上层业务校验并准备好的对象。
\item
  该条目可以被编辑器和类型检查器识别，但当前运行时可能在属性查找阶段直接失败。
\end{itemize}

\textbf{用法}

\begin{lstlisting}[language=Python]
from typing import TYPE_CHECKING

if TYPE_CHECKING:
    def planned_from_json(obj: ConfigGetParameter, json: dict[str, Any]) -> None:
        obj.from_json(json=json)
\end{lstlisting}

\begin{quote}
{[}!CAUTION{]} 上面的 \passthrough{\lstinline!TYPE\_CHECKING!}
示例只表达计划调用形态。该分支在普通 Python
运行时为假，不代表当前扩展能执行此接口。
\end{quote}

\hypertarget{unitree-sdk2-cpp-robot-go2-configgetparameter-to-json-1}{%
\paragraph{\texorpdfstring{\texttt{unitree\_sdk2\_cpp.robot.go2.ConfigGetParameter.to\_json}}{unitree\_sdk2\_cpp.robot.go2.ConfigGetParameter.to\_json}}\label{unitree-sdk2-cpp-robot-go2-configgetparameter-to-json-1}}

计划把当前 SDK 值对象写入 JSON 风格字典。

\textbf{签名}

\begin{lstlisting}[language=Python]
def to_json(self, json: dict[str, Any]) -> None
\end{lstlisting}

\textbf{可用性与安全性}

\textbf{\passthrough{\lstinline!SIGNATURE\_ONLY!} /
\passthrough{\lstinline!UNCLASSIFIED!}}：当前只有设计期签名，不能假设已存在于运行时扩展。

\textbf{参数}

\begin{longtable}[]{@{}
  >{\raggedright\arraybackslash}p{(\columnwidth - 6\tabcolsep) * \real{0.18}}
  >{\raggedright\arraybackslash}p{(\columnwidth - 6\tabcolsep) * \real{0.24}}
  >{\raggedright\arraybackslash}p{(\columnwidth - 6\tabcolsep) * \real{0.18}}
  >{\raggedright\arraybackslash}p{(\columnwidth - 6\tabcolsep) * \real{0.40}}@{}}
\toprule
名称 & Python 类型 & 默认值 & 含义 \\
\midrule
\endhead
\passthrough{\lstinline!json!} &
\passthrough{\lstinline!dict[str, Any]!} & 必填 & JSON
风格字典，键和值必须满足对应 C++ \passthrough{\lstinline!JsonMap!}
数据结构的约束。 对应 C++ 参数
\passthrough{\lstinline!json: common::JsonMap \&!}。 \\
\bottomrule
\end{longtable}

\textbf{返回值}

\begin{longtable}[]{@{}
  >{\raggedright\arraybackslash}p{(\columnwidth - 4\tabcolsep) * \real{0.18}}
  >{\raggedright\arraybackslash}p{(\columnwidth - 4\tabcolsep) * \real{0.20}}
  >{\raggedright\arraybackslash}p{(\columnwidth - 4\tabcolsep) * \real{0.62}}@{}}
\toprule
位置 & 类型 & 含义 \\
\midrule
\endhead
无 & \passthrough{\lstinline!None!} & 该方法不返回 Python
值；失败可能表现为异常或底层状态变化。 \\
\bottomrule
\end{longtable}

\textbf{对应 C++}

\begin{itemize}
\tightlist
\item
  类：\passthrough{\lstinline!unitree::robot::go2::ConfigGetParameter!}
\item
  签名：\passthrough{\lstinline!toJson(common::JsonMap \&) const!}
\item
  绑定策略：\passthrough{\lstinline!SIGNATURE\_PREVIEW!}
\item
  声明位置：\passthrough{\lstinline!include/unitree/robot/go2/config/config\_api.hpp:105!}
\end{itemize}

\textbf{说明}

\begin{itemize}
\tightlist
\item
  示例是局部调用片段；\passthrough{\lstinline!obj!}
  和参数变量表示已经按上层业务校验并准备好的对象。
\item
  该条目可以被编辑器和类型检查器识别，但当前运行时可能在属性查找阶段直接失败。
\end{itemize}

\textbf{用法}

\begin{lstlisting}[language=Python]
from typing import TYPE_CHECKING

if TYPE_CHECKING:
    def planned_to_json(obj: ConfigGetParameter, json: dict[str, Any]) -> None:
        obj.to_json(json=json)
\end{lstlisting}

\begin{quote}
{[}!CAUTION{]} 上面的 \passthrough{\lstinline!TYPE\_CHECKING!}
示例只表达计划调用形态。该分支在普通 Python
运行时为假，不代表当前扩展能执行此接口。
\end{quote}

\hypertarget{unitree-sdk2-cpp-robot-go2-configmeta}{%
\subsubsection{\texorpdfstring{\texttt{unitree\_sdk2\_cpp.robot.go2.ConfigMeta}}{unitree\_sdk2\_cpp.robot.go2.ConfigMeta}}\label{unitree-sdk2-cpp-robot-go2-configmeta}}

SDK 数据值类型；公开字段和转换方法见下文。

\textbf{导入}

\begin{lstlisting}[language=Python]
from unitree_sdk2_cpp.robot.go2 import ConfigMeta
\end{lstlisting}

\textbf{基类}：\passthrough{\lstinline!object!}

\textbf{公开属性}

\begin{longtable}[]{@{}
  >{\raggedright\arraybackslash}p{(\columnwidth - 6\tabcolsep) * \real{0.18}}
  >{\raggedright\arraybackslash}p{(\columnwidth - 6\tabcolsep) * \real{0.24}}
  >{\raggedright\arraybackslash}p{(\columnwidth - 6\tabcolsep) * \real{0.18}}
  >{\raggedright\arraybackslash}p{(\columnwidth - 6\tabcolsep) * \real{0.40}}@{}}
\toprule
名称 & Python 类型 & 含义 & 用法 \\
\midrule
\endhead
\passthrough{\lstinline!name!} & \passthrough{\lstinline!str!} & SDK
对象、配置项、服务或动作的名称。允许值由调用它的具体接口定义。 &
\passthrough{\lstinline!value = obj.name!} /
\passthrough{\lstinline!obj.name = value!} \\
\passthrough{\lstinline!last\_modified!} & \passthrough{\lstinline!str!}
& 传给该接口的 \passthrough{\lstinline!last!}
\passthrough{\lstinline!modified!} 参数，Python 类型为
\passthrough{\lstinline!str!}。精确含义、单位、范围和枚举值需查目标型号对应头文件与协议。
& \passthrough{\lstinline!value = obj.last\_modified!} /
\passthrough{\lstinline!obj.last\_modified = value!} \\
\passthrough{\lstinline!size!} & \passthrough{\lstinline!int!} &
传给该接口的 \passthrough{\lstinline!size!} 参数，Python 类型为
\passthrough{\lstinline!int!}。精确含义、单位、范围和枚举值需查目标型号对应头文件与协议。
& \passthrough{\lstinline!value = obj.size!} /
\passthrough{\lstinline!obj.size = value!} \\
\passthrough{\lstinline!epoch!} & \passthrough{\lstinline!int!} &
传给该接口的 \passthrough{\lstinline!epoch!} 参数，Python 类型为
\passthrough{\lstinline!int!}。精确含义、单位、范围和枚举值需查目标型号对应头文件与协议。
& \passthrough{\lstinline!value = obj.epoch!} /
\passthrough{\lstinline!obj.epoch = value!} \\
\bottomrule
\end{longtable}

\textbf{方法索引}

\begin{longtable}[]{@{}
  >{\raggedright\arraybackslash}p{(\columnwidth - 6\tabcolsep) * \real{0.18}}
  >{\raggedright\arraybackslash}p{(\columnwidth - 6\tabcolsep) * \real{0.24}}
  >{\raggedright\arraybackslash}p{(\columnwidth - 6\tabcolsep) * \real{0.18}}
  >{\raggedright\arraybackslash}p{(\columnwidth - 6\tabcolsep) * \real{0.40}}@{}}
\toprule
方法 & Python 签名 & 状态 & 安全分类 \\
\midrule
\endhead
\protect\hyperlink{unitree-sdk2-cpp-robot-go2-configmeta-dunder-init-1}{\passthrough{\lstinline!\_\_init\_\_!}}
& \passthrough{\lstinline!def \_\_init\_\_(self) -> None!} &
\passthrough{\lstinline!SIGNATURE\_ONLY!} &
\passthrough{\lstinline!CONSTRUCTION!} \\
\bottomrule
\end{longtable}

\hypertarget{unitree-sdk2-cpp-robot-go2-configmeta-dunder-init-1}{%
\paragraph{\texorpdfstring{\texttt{unitree\_sdk2\_cpp.robot.go2.ConfigMeta.\_\_init\_\_}}{unitree\_sdk2\_cpp.robot.go2.ConfigMeta.\_\_init\_\_}}\label{unitree-sdk2-cpp-robot-go2-configmeta-dunder-init-1}}

初始化 \passthrough{\lstinline!ConfigMeta!}
实例。是否能实际构造取决于下方可用性状态。

\textbf{签名}

\begin{lstlisting}[language=Python]
def __init__(self) -> None
\end{lstlisting}

\textbf{可用性与安全性}

\textbf{\passthrough{\lstinline!SIGNATURE\_ONLY!} /
\passthrough{\lstinline!CONSTRUCTION!}}：当前只有设计期签名，不能假设已存在于运行时扩展。

\textbf{参数}

无显式参数。实例方法中的 \passthrough{\lstinline!self!} 由 Python
自动传入。

\textbf{返回值}

\begin{longtable}[]{@{}
  >{\raggedright\arraybackslash}p{(\columnwidth - 4\tabcolsep) * \real{0.18}}
  >{\raggedright\arraybackslash}p{(\columnwidth - 4\tabcolsep) * \real{0.20}}
  >{\raggedright\arraybackslash}p{(\columnwidth - 4\tabcolsep) * \real{0.62}}@{}}
\toprule
位置 & 类型 & 含义 \\
\midrule
\endhead
无 & \passthrough{\lstinline!None!} &
\passthrough{\lstinline!\_\_init\_\_!}
本身不返回值；调用类对象时，在构造函数可用的前提下得到该类实例。 \\
\bottomrule
\end{longtable}

\textbf{对应 C++}

\begin{itemize}
\tightlist
\item
  类：\passthrough{\lstinline!unitree::robot::go2::ConfigMeta!}
\item
  签名：\passthrough{\lstinline!ConfigMeta()!}
\item
  绑定策略：\passthrough{\lstinline!未单独标注!}
\item
  声明位置：\passthrough{\lstinline!include/unitree/robot/go2/config/config\_client.hpp:20!}
\end{itemize}

\textbf{说明}

\begin{itemize}
\tightlist
\item
  示例是局部调用片段；\passthrough{\lstinline!obj!}
  和参数变量表示已经按上层业务校验并准备好的对象。
\item
  该条目可以被编辑器和类型检查器识别，但当前运行时可能在属性查找阶段直接失败。
\end{itemize}

\textbf{用法}

\begin{lstlisting}[language=Python]
from typing import TYPE_CHECKING

if TYPE_CHECKING:
    def planned_construct() -> ConfigMeta:
        return ConfigMeta()
\end{lstlisting}

\begin{quote}
{[}!CAUTION{]} 上面的 \passthrough{\lstinline!TYPE\_CHECKING!}
示例只表达计划调用形态。该分支在普通 Python
运行时为假，不代表当前扩展能执行此接口。
\end{quote}

\hypertarget{unitree-sdk2-cpp-robot-go2-configmetadata}{%
\subsubsection{\texorpdfstring{\texttt{unitree\_sdk2\_cpp.robot.go2.ConfigMetaData}}{unitree\_sdk2\_cpp.robot.go2.ConfigMetaData}}\label{unitree-sdk2-cpp-robot-go2-configmetadata}}

SDK 数据值类型；公开字段和转换方法见下文。

\textbf{导入}

\begin{lstlisting}[language=Python]
from unitree_sdk2_cpp.robot.go2 import ConfigMetaData
\end{lstlisting}

\textbf{基类}：\passthrough{\lstinline!object!}

\textbf{公开属性}

\begin{longtable}[]{@{}
  >{\raggedright\arraybackslash}p{(\columnwidth - 6\tabcolsep) * \real{0.18}}
  >{\raggedright\arraybackslash}p{(\columnwidth - 6\tabcolsep) * \real{0.24}}
  >{\raggedright\arraybackslash}p{(\columnwidth - 6\tabcolsep) * \real{0.18}}
  >{\raggedright\arraybackslash}p{(\columnwidth - 6\tabcolsep) * \real{0.40}}@{}}
\toprule
名称 & Python 类型 & 含义 & 用法 \\
\midrule
\endhead
\passthrough{\lstinline!meta!} &
\passthrough{\lstinline!JsonizeConfigMeta!} & 传给该接口的
\passthrough{\lstinline!meta!} 参数，Python 类型为
\passthrough{\lstinline!JsonizeConfigMeta!}。精确含义、单位、范围和枚举值需查目标型号对应头文件与协议。
& \passthrough{\lstinline!value = obj.meta!} /
\passthrough{\lstinline!obj.meta = value!} \\
\bottomrule
\end{longtable}

\textbf{方法索引}

\begin{longtable}[]{@{}
  >{\raggedright\arraybackslash}p{(\columnwidth - 6\tabcolsep) * \real{0.18}}
  >{\raggedright\arraybackslash}p{(\columnwidth - 6\tabcolsep) * \real{0.24}}
  >{\raggedright\arraybackslash}p{(\columnwidth - 6\tabcolsep) * \real{0.18}}
  >{\raggedright\arraybackslash}p{(\columnwidth - 6\tabcolsep) * \real{0.40}}@{}}
\toprule
方法 & Python 签名 & 状态 & 安全分类 \\
\midrule
\endhead
\protect\hyperlink{unitree-sdk2-cpp-robot-go2-configmetadata-dunder-init-1}{\passthrough{\lstinline!\_\_init\_\_!}}
& \passthrough{\lstinline!def \_\_init\_\_(self) -> None!} &
\passthrough{\lstinline!SIGNATURE\_ONLY!} &
\passthrough{\lstinline!CONSTRUCTION!} \\
\protect\hyperlink{unitree-sdk2-cpp-robot-go2-configmetadata-from-json-1}{\passthrough{\lstinline!from\_json!}}
&
\passthrough{\lstinline!def from\_json(self, json: dict[str, Any]) -> None!}
& \passthrough{\lstinline!SIGNATURE\_ONLY!} &
\passthrough{\lstinline!UNCLASSIFIED!} \\
\protect\hyperlink{unitree-sdk2-cpp-robot-go2-configmetadata-to-json-1}{\passthrough{\lstinline!to\_json!}}
&
\passthrough{\lstinline!def to\_json(self, json: dict[str, Any]) -> None!}
& \passthrough{\lstinline!SIGNATURE\_ONLY!} &
\passthrough{\lstinline!UNCLASSIFIED!} \\
\bottomrule
\end{longtable}

\hypertarget{unitree-sdk2-cpp-robot-go2-configmetadata-dunder-init-1}{%
\paragraph{\texorpdfstring{\texttt{unitree\_sdk2\_cpp.robot.go2.ConfigMetaData.\_\_init\_\_}}{unitree\_sdk2\_cpp.robot.go2.ConfigMetaData.\_\_init\_\_}}\label{unitree-sdk2-cpp-robot-go2-configmetadata-dunder-init-1}}

初始化 \passthrough{\lstinline!ConfigMetaData!}
实例。是否能实际构造取决于下方可用性状态。

\textbf{签名}

\begin{lstlisting}[language=Python]
def __init__(self) -> None
\end{lstlisting}

\textbf{可用性与安全性}

\textbf{\passthrough{\lstinline!SIGNATURE\_ONLY!} /
\passthrough{\lstinline!CONSTRUCTION!}}：当前只有设计期签名，不能假设已存在于运行时扩展。

\textbf{参数}

无显式参数。实例方法中的 \passthrough{\lstinline!self!} 由 Python
自动传入。

\textbf{返回值}

\begin{longtable}[]{@{}
  >{\raggedright\arraybackslash}p{(\columnwidth - 4\tabcolsep) * \real{0.18}}
  >{\raggedright\arraybackslash}p{(\columnwidth - 4\tabcolsep) * \real{0.20}}
  >{\raggedright\arraybackslash}p{(\columnwidth - 4\tabcolsep) * \real{0.62}}@{}}
\toprule
位置 & 类型 & 含义 \\
\midrule
\endhead
无 & \passthrough{\lstinline!None!} &
\passthrough{\lstinline!\_\_init\_\_!}
本身不返回值；调用类对象时，在构造函数可用的前提下得到该类实例。 \\
\bottomrule
\end{longtable}

\textbf{对应 C++}

\begin{itemize}
\tightlist
\item
  类：\passthrough{\lstinline!unitree::robot::go2::ConfigMetaData!}
\item
  签名：\passthrough{\lstinline!ConfigMetaData()!}
\item
  绑定策略：\passthrough{\lstinline!未单独标注!}
\item
  声明位置：\passthrough{\lstinline!include/unitree/robot/go2/config/config\_api.hpp:186!}
\end{itemize}

\textbf{说明}

\begin{itemize}
\tightlist
\item
  示例是局部调用片段；\passthrough{\lstinline!obj!}
  和参数变量表示已经按上层业务校验并准备好的对象。
\item
  该条目可以被编辑器和类型检查器识别，但当前运行时可能在属性查找阶段直接失败。
\end{itemize}

\textbf{用法}

\begin{lstlisting}[language=Python]
from typing import TYPE_CHECKING

if TYPE_CHECKING:
    def planned_construct() -> ConfigMetaData:
        return ConfigMetaData()
\end{lstlisting}

\begin{quote}
{[}!CAUTION{]} 上面的 \passthrough{\lstinline!TYPE\_CHECKING!}
示例只表达计划调用形态。该分支在普通 Python
运行时为假，不代表当前扩展能执行此接口。
\end{quote}

\hypertarget{unitree-sdk2-cpp-robot-go2-configmetadata-from-json-1}{%
\paragraph{\texorpdfstring{\texttt{unitree\_sdk2\_cpp.robot.go2.ConfigMetaData.from\_json}}{unitree\_sdk2\_cpp.robot.go2.ConfigMetaData.from\_json}}\label{unitree-sdk2-cpp-robot-go2-configmetadata-from-json-1}}

计划从 JSON 风格字典读取字段并更新当前 SDK 值对象。

\textbf{签名}

\begin{lstlisting}[language=Python]
def from_json(self, json: dict[str, Any]) -> None
\end{lstlisting}

\textbf{可用性与安全性}

\textbf{\passthrough{\lstinline!SIGNATURE\_ONLY!} /
\passthrough{\lstinline!UNCLASSIFIED!}}：当前只有设计期签名，不能假设已存在于运行时扩展。

\textbf{参数}

\begin{longtable}[]{@{}
  >{\raggedright\arraybackslash}p{(\columnwidth - 6\tabcolsep) * \real{0.18}}
  >{\raggedright\arraybackslash}p{(\columnwidth - 6\tabcolsep) * \real{0.24}}
  >{\raggedright\arraybackslash}p{(\columnwidth - 6\tabcolsep) * \real{0.18}}
  >{\raggedright\arraybackslash}p{(\columnwidth - 6\tabcolsep) * \real{0.40}}@{}}
\toprule
名称 & Python 类型 & 默认值 & 含义 \\
\midrule
\endhead
\passthrough{\lstinline!json!} &
\passthrough{\lstinline!dict[str, Any]!} & 必填 & JSON
风格字典，键和值必须满足对应 C++ \passthrough{\lstinline!JsonMap!}
数据结构的约束。 对应 C++ 参数
\passthrough{\lstinline!json: common::JsonMap \&!}。 \\
\bottomrule
\end{longtable}

\textbf{返回值}

\begin{longtable}[]{@{}
  >{\raggedright\arraybackslash}p{(\columnwidth - 4\tabcolsep) * \real{0.18}}
  >{\raggedright\arraybackslash}p{(\columnwidth - 4\tabcolsep) * \real{0.20}}
  >{\raggedright\arraybackslash}p{(\columnwidth - 4\tabcolsep) * \real{0.62}}@{}}
\toprule
位置 & 类型 & 含义 \\
\midrule
\endhead
无 & \passthrough{\lstinline!None!} & 该方法不返回 Python
值；失败可能表现为异常或底层状态变化。 \\
\bottomrule
\end{longtable}

\textbf{对应 C++}

\begin{itemize}
\tightlist
\item
  类：\passthrough{\lstinline!unitree::robot::go2::ConfigMetaData!}
\item
  签名：\passthrough{\lstinline!fromJson(common::JsonMap \&)!}
\item
  绑定策略：\passthrough{\lstinline!SIGNATURE\_PREVIEW!}
\item
  声明位置：\passthrough{\lstinline!include/unitree/robot/go2/config/config\_api.hpp:192!}
\end{itemize}

\textbf{说明}

\begin{itemize}
\tightlist
\item
  示例是局部调用片段；\passthrough{\lstinline!obj!}
  和参数变量表示已经按上层业务校验并准备好的对象。
\item
  该条目可以被编辑器和类型检查器识别，但当前运行时可能在属性查找阶段直接失败。
\end{itemize}

\textbf{用法}

\begin{lstlisting}[language=Python]
from typing import TYPE_CHECKING

if TYPE_CHECKING:
    def planned_from_json(obj: ConfigMetaData, json: dict[str, Any]) -> None:
        obj.from_json(json=json)
\end{lstlisting}

\begin{quote}
{[}!CAUTION{]} 上面的 \passthrough{\lstinline!TYPE\_CHECKING!}
示例只表达计划调用形态。该分支在普通 Python
运行时为假，不代表当前扩展能执行此接口。
\end{quote}

\hypertarget{unitree-sdk2-cpp-robot-go2-configmetadata-to-json-1}{%
\paragraph{\texorpdfstring{\texttt{unitree\_sdk2\_cpp.robot.go2.ConfigMetaData.to\_json}}{unitree\_sdk2\_cpp.robot.go2.ConfigMetaData.to\_json}}\label{unitree-sdk2-cpp-robot-go2-configmetadata-to-json-1}}

计划把当前 SDK 值对象写入 JSON 风格字典。

\textbf{签名}

\begin{lstlisting}[language=Python]
def to_json(self, json: dict[str, Any]) -> None
\end{lstlisting}

\textbf{可用性与安全性}

\textbf{\passthrough{\lstinline!SIGNATURE\_ONLY!} /
\passthrough{\lstinline!UNCLASSIFIED!}}：当前只有设计期签名，不能假设已存在于运行时扩展。

\textbf{参数}

\begin{longtable}[]{@{}
  >{\raggedright\arraybackslash}p{(\columnwidth - 6\tabcolsep) * \real{0.18}}
  >{\raggedright\arraybackslash}p{(\columnwidth - 6\tabcolsep) * \real{0.24}}
  >{\raggedright\arraybackslash}p{(\columnwidth - 6\tabcolsep) * \real{0.18}}
  >{\raggedright\arraybackslash}p{(\columnwidth - 6\tabcolsep) * \real{0.40}}@{}}
\toprule
名称 & Python 类型 & 默认值 & 含义 \\
\midrule
\endhead
\passthrough{\lstinline!json!} &
\passthrough{\lstinline!dict[str, Any]!} & 必填 & JSON
风格字典，键和值必须满足对应 C++ \passthrough{\lstinline!JsonMap!}
数据结构的约束。 对应 C++ 参数
\passthrough{\lstinline!json: common::JsonMap \&!}。 \\
\bottomrule
\end{longtable}

\textbf{返回值}

\begin{longtable}[]{@{}
  >{\raggedright\arraybackslash}p{(\columnwidth - 4\tabcolsep) * \real{0.18}}
  >{\raggedright\arraybackslash}p{(\columnwidth - 4\tabcolsep) * \real{0.20}}
  >{\raggedright\arraybackslash}p{(\columnwidth - 4\tabcolsep) * \real{0.62}}@{}}
\toprule
位置 & 类型 & 含义 \\
\midrule
\endhead
无 & \passthrough{\lstinline!None!} & 该方法不返回 Python
值；失败可能表现为异常或底层状态变化。 \\
\bottomrule
\end{longtable}

\textbf{对应 C++}

\begin{itemize}
\tightlist
\item
  类：\passthrough{\lstinline!unitree::robot::go2::ConfigMetaData!}
\item
  签名：\passthrough{\lstinline!toJson(common::JsonMap \&) const!}
\item
  绑定策略：\passthrough{\lstinline!SIGNATURE\_PREVIEW!}
\item
  声明位置：\passthrough{\lstinline!include/unitree/robot/go2/config/config\_api.hpp:197!}
\end{itemize}

\textbf{说明}

\begin{itemize}
\tightlist
\item
  示例是局部调用片段；\passthrough{\lstinline!obj!}
  和参数变量表示已经按上层业务校验并准备好的对象。
\item
  该条目可以被编辑器和类型检查器识别，但当前运行时可能在属性查找阶段直接失败。
\end{itemize}

\textbf{用法}

\begin{lstlisting}[language=Python]
from typing import TYPE_CHECKING

if TYPE_CHECKING:
    def planned_to_json(obj: ConfigMetaData, json: dict[str, Any]) -> None:
        obj.to_json(json=json)
\end{lstlisting}

\begin{quote}
{[}!CAUTION{]} 上面的 \passthrough{\lstinline!TYPE\_CHECKING!}
示例只表达计划调用形态。该分支在普通 Python
运行时为假，不代表当前扩展能执行此接口。
\end{quote}

\hypertarget{unitree-sdk2-cpp-robot-go2-configmetaparameter}{%
\subsubsection{\texorpdfstring{\texttt{unitree\_sdk2\_cpp.robot.go2.ConfigMetaParameter}}{unitree\_sdk2\_cpp.robot.go2.ConfigMetaParameter}}\label{unitree-sdk2-cpp-robot-go2-configmetaparameter}}

SDK 请求参数值类型；当前可能仅提供设计期签名。

\textbf{导入}

\begin{lstlisting}[language=Python]
from unitree_sdk2_cpp.robot.go2 import ConfigMetaParameter
\end{lstlisting}

\textbf{基类}：\passthrough{\lstinline!object!}

\textbf{公开属性}

\begin{longtable}[]{@{}
  >{\raggedright\arraybackslash}p{(\columnwidth - 6\tabcolsep) * \real{0.18}}
  >{\raggedright\arraybackslash}p{(\columnwidth - 6\tabcolsep) * \real{0.24}}
  >{\raggedright\arraybackslash}p{(\columnwidth - 6\tabcolsep) * \real{0.18}}
  >{\raggedright\arraybackslash}p{(\columnwidth - 6\tabcolsep) * \real{0.40}}@{}}
\toprule
名称 & Python 类型 & 含义 & 用法 \\
\midrule
\endhead
\passthrough{\lstinline!name!} & \passthrough{\lstinline!str!} & SDK
对象、配置项、服务或动作的名称。允许值由调用它的具体接口定义。 &
\passthrough{\lstinline!value = obj.name!} /
\passthrough{\lstinline!obj.name = value!} \\
\bottomrule
\end{longtable}

\textbf{方法索引}

\begin{longtable}[]{@{}
  >{\raggedright\arraybackslash}p{(\columnwidth - 6\tabcolsep) * \real{0.18}}
  >{\raggedright\arraybackslash}p{(\columnwidth - 6\tabcolsep) * \real{0.24}}
  >{\raggedright\arraybackslash}p{(\columnwidth - 6\tabcolsep) * \real{0.18}}
  >{\raggedright\arraybackslash}p{(\columnwidth - 6\tabcolsep) * \real{0.40}}@{}}
\toprule
方法 & Python 签名 & 状态 & 安全分类 \\
\midrule
\endhead
\protect\hyperlink{unitree-sdk2-cpp-robot-go2-configmetaparameter-dunder-init-1}{\passthrough{\lstinline!\_\_init\_\_!}}
& \passthrough{\lstinline!def \_\_init\_\_(self) -> None!} &
\passthrough{\lstinline!SIGNATURE\_ONLY!} &
\passthrough{\lstinline!CONSTRUCTION!} \\
\protect\hyperlink{unitree-sdk2-cpp-robot-go2-configmetaparameter-from-json-1}{\passthrough{\lstinline!from\_json!}}
&
\passthrough{\lstinline!def from\_json(self, json: dict[str, Any]) -> None!}
& \passthrough{\lstinline!SIGNATURE\_ONLY!} &
\passthrough{\lstinline!UNCLASSIFIED!} \\
\protect\hyperlink{unitree-sdk2-cpp-robot-go2-configmetaparameter-to-json-1}{\passthrough{\lstinline!to\_json!}}
&
\passthrough{\lstinline!def to\_json(self, json: dict[str, Any]) -> None!}
& \passthrough{\lstinline!SIGNATURE\_ONLY!} &
\passthrough{\lstinline!UNCLASSIFIED!} \\
\bottomrule
\end{longtable}

\hypertarget{unitree-sdk2-cpp-robot-go2-configmetaparameter-dunder-init-1}{%
\paragraph{\texorpdfstring{\texttt{unitree\_sdk2\_cpp.robot.go2.ConfigMetaParameter.\_\_init\_\_}}{unitree\_sdk2\_cpp.robot.go2.ConfigMetaParameter.\_\_init\_\_}}\label{unitree-sdk2-cpp-robot-go2-configmetaparameter-dunder-init-1}}

初始化 \passthrough{\lstinline!ConfigMetaParameter!}
实例。是否能实际构造取决于下方可用性状态。

\textbf{签名}

\begin{lstlisting}[language=Python]
def __init__(self) -> None
\end{lstlisting}

\textbf{可用性与安全性}

\textbf{\passthrough{\lstinline!SIGNATURE\_ONLY!} /
\passthrough{\lstinline!CONSTRUCTION!}}：当前只有设计期签名，不能假设已存在于运行时扩展。

\textbf{参数}

无显式参数。实例方法中的 \passthrough{\lstinline!self!} 由 Python
自动传入。

\textbf{返回值}

\begin{longtable}[]{@{}
  >{\raggedright\arraybackslash}p{(\columnwidth - 4\tabcolsep) * \real{0.18}}
  >{\raggedright\arraybackslash}p{(\columnwidth - 4\tabcolsep) * \real{0.20}}
  >{\raggedright\arraybackslash}p{(\columnwidth - 4\tabcolsep) * \real{0.62}}@{}}
\toprule
位置 & 类型 & 含义 \\
\midrule
\endhead
无 & \passthrough{\lstinline!None!} &
\passthrough{\lstinline!\_\_init\_\_!}
本身不返回值；调用类对象时，在构造函数可用的前提下得到该类实例。 \\
\bottomrule
\end{longtable}

\textbf{对应 C++}

\begin{itemize}
\tightlist
\item
  类：\passthrough{\lstinline!unitree::robot::go2::ConfigMetaParameter!}
\item
  签名：\passthrough{\lstinline!ConfigMetaParameter()!}
\item
  绑定策略：\passthrough{\lstinline!未单独标注!}
\item
  声明位置：\passthrough{\lstinline!include/unitree/robot/go2/config/config\_api.hpp:163!}
\end{itemize}

\textbf{说明}

\begin{itemize}
\tightlist
\item
  示例是局部调用片段；\passthrough{\lstinline!obj!}
  和参数变量表示已经按上层业务校验并准备好的对象。
\item
  该条目可以被编辑器和类型检查器识别，但当前运行时可能在属性查找阶段直接失败。
\end{itemize}

\textbf{用法}

\begin{lstlisting}[language=Python]
from typing import TYPE_CHECKING

if TYPE_CHECKING:
    def planned_construct() -> ConfigMetaParameter:
        return ConfigMetaParameter()
\end{lstlisting}

\begin{quote}
{[}!CAUTION{]} 上面的 \passthrough{\lstinline!TYPE\_CHECKING!}
示例只表达计划调用形态。该分支在普通 Python
运行时为假，不代表当前扩展能执行此接口。
\end{quote}

\hypertarget{unitree-sdk2-cpp-robot-go2-configmetaparameter-from-json-1}{%
\paragraph{\texorpdfstring{\texttt{unitree\_sdk2\_cpp.robot.go2.ConfigMetaParameter.from\_json}}{unitree\_sdk2\_cpp.robot.go2.ConfigMetaParameter.from\_json}}\label{unitree-sdk2-cpp-robot-go2-configmetaparameter-from-json-1}}

计划从 JSON 风格字典读取字段并更新当前 SDK 值对象。

\textbf{签名}

\begin{lstlisting}[language=Python]
def from_json(self, json: dict[str, Any]) -> None
\end{lstlisting}

\textbf{可用性与安全性}

\textbf{\passthrough{\lstinline!SIGNATURE\_ONLY!} /
\passthrough{\lstinline!UNCLASSIFIED!}}：当前只有设计期签名，不能假设已存在于运行时扩展。

\textbf{参数}

\begin{longtable}[]{@{}
  >{\raggedright\arraybackslash}p{(\columnwidth - 6\tabcolsep) * \real{0.18}}
  >{\raggedright\arraybackslash}p{(\columnwidth - 6\tabcolsep) * \real{0.24}}
  >{\raggedright\arraybackslash}p{(\columnwidth - 6\tabcolsep) * \real{0.18}}
  >{\raggedright\arraybackslash}p{(\columnwidth - 6\tabcolsep) * \real{0.40}}@{}}
\toprule
名称 & Python 类型 & 默认值 & 含义 \\
\midrule
\endhead
\passthrough{\lstinline!json!} &
\passthrough{\lstinline!dict[str, Any]!} & 必填 & JSON
风格字典，键和值必须满足对应 C++ \passthrough{\lstinline!JsonMap!}
数据结构的约束。 对应 C++ 参数
\passthrough{\lstinline!json: common::JsonMap \&!}。 \\
\bottomrule
\end{longtable}

\textbf{返回值}

\begin{longtable}[]{@{}
  >{\raggedright\arraybackslash}p{(\columnwidth - 4\tabcolsep) * \real{0.18}}
  >{\raggedright\arraybackslash}p{(\columnwidth - 4\tabcolsep) * \real{0.20}}
  >{\raggedright\arraybackslash}p{(\columnwidth - 4\tabcolsep) * \real{0.62}}@{}}
\toprule
位置 & 类型 & 含义 \\
\midrule
\endhead
无 & \passthrough{\lstinline!None!} & 该方法不返回 Python
值；失败可能表现为异常或底层状态变化。 \\
\bottomrule
\end{longtable}

\textbf{对应 C++}

\begin{itemize}
\tightlist
\item
  类：\passthrough{\lstinline!unitree::robot::go2::ConfigMetaParameter!}
\item
  签名：\passthrough{\lstinline!fromJson(common::JsonMap \&)!}
\item
  绑定策略：\passthrough{\lstinline!SIGNATURE\_PREVIEW!}
\item
  声明位置：\passthrough{\lstinline!include/unitree/robot/go2/config/config\_api.hpp:169!}
\end{itemize}

\textbf{说明}

\begin{itemize}
\tightlist
\item
  示例是局部调用片段；\passthrough{\lstinline!obj!}
  和参数变量表示已经按上层业务校验并准备好的对象。
\item
  该条目可以被编辑器和类型检查器识别，但当前运行时可能在属性查找阶段直接失败。
\end{itemize}

\textbf{用法}

\begin{lstlisting}[language=Python]
from typing import TYPE_CHECKING

if TYPE_CHECKING:
    def planned_from_json(obj: ConfigMetaParameter, json: dict[str, Any]) -> None:
        obj.from_json(json=json)
\end{lstlisting}

\begin{quote}
{[}!CAUTION{]} 上面的 \passthrough{\lstinline!TYPE\_CHECKING!}
示例只表达计划调用形态。该分支在普通 Python
运行时为假，不代表当前扩展能执行此接口。
\end{quote}

\hypertarget{unitree-sdk2-cpp-robot-go2-configmetaparameter-to-json-1}{%
\paragraph{\texorpdfstring{\texttt{unitree\_sdk2\_cpp.robot.go2.ConfigMetaParameter.to\_json}}{unitree\_sdk2\_cpp.robot.go2.ConfigMetaParameter.to\_json}}\label{unitree-sdk2-cpp-robot-go2-configmetaparameter-to-json-1}}

计划把当前 SDK 值对象写入 JSON 风格字典。

\textbf{签名}

\begin{lstlisting}[language=Python]
def to_json(self, json: dict[str, Any]) -> None
\end{lstlisting}

\textbf{可用性与安全性}

\textbf{\passthrough{\lstinline!SIGNATURE\_ONLY!} /
\passthrough{\lstinline!UNCLASSIFIED!}}：当前只有设计期签名，不能假设已存在于运行时扩展。

\textbf{参数}

\begin{longtable}[]{@{}
  >{\raggedright\arraybackslash}p{(\columnwidth - 6\tabcolsep) * \real{0.18}}
  >{\raggedright\arraybackslash}p{(\columnwidth - 6\tabcolsep) * \real{0.24}}
  >{\raggedright\arraybackslash}p{(\columnwidth - 6\tabcolsep) * \real{0.18}}
  >{\raggedright\arraybackslash}p{(\columnwidth - 6\tabcolsep) * \real{0.40}}@{}}
\toprule
名称 & Python 类型 & 默认值 & 含义 \\
\midrule
\endhead
\passthrough{\lstinline!json!} &
\passthrough{\lstinline!dict[str, Any]!} & 必填 & JSON
风格字典，键和值必须满足对应 C++ \passthrough{\lstinline!JsonMap!}
数据结构的约束。 对应 C++ 参数
\passthrough{\lstinline!json: common::JsonMap \&!}。 \\
\bottomrule
\end{longtable}

\textbf{返回值}

\begin{longtable}[]{@{}
  >{\raggedright\arraybackslash}p{(\columnwidth - 4\tabcolsep) * \real{0.18}}
  >{\raggedright\arraybackslash}p{(\columnwidth - 4\tabcolsep) * \real{0.20}}
  >{\raggedright\arraybackslash}p{(\columnwidth - 4\tabcolsep) * \real{0.62}}@{}}
\toprule
位置 & 类型 & 含义 \\
\midrule
\endhead
无 & \passthrough{\lstinline!None!} & 该方法不返回 Python
值；失败可能表现为异常或底层状态变化。 \\
\bottomrule
\end{longtable}

\textbf{对应 C++}

\begin{itemize}
\tightlist
\item
  类：\passthrough{\lstinline!unitree::robot::go2::ConfigMetaParameter!}
\item
  签名：\passthrough{\lstinline!toJson(common::JsonMap \&) const!}
\item
  绑定策略：\passthrough{\lstinline!SIGNATURE\_PREVIEW!}
\item
  声明位置：\passthrough{\lstinline!include/unitree/robot/go2/config/config\_api.hpp:174!}
\end{itemize}

\textbf{说明}

\begin{itemize}
\tightlist
\item
  示例是局部调用片段；\passthrough{\lstinline!obj!}
  和参数变量表示已经按上层业务校验并准备好的对象。
\item
  该条目可以被编辑器和类型检查器识别，但当前运行时可能在属性查找阶段直接失败。
\end{itemize}

\textbf{用法}

\begin{lstlisting}[language=Python]
from typing import TYPE_CHECKING

if TYPE_CHECKING:
    def planned_to_json(obj: ConfigMetaParameter, json: dict[str, Any]) -> None:
        obj.to_json(json=json)
\end{lstlisting}

\begin{quote}
{[}!CAUTION{]} 上面的 \passthrough{\lstinline!TYPE\_CHECKING!}
示例只表达计划调用形态。该分支在普通 Python
运行时为假，不代表当前扩展能执行此接口。
\end{quote}

\hypertarget{unitree-sdk2-cpp-robot-go2-configsetparameter}{%
\subsubsection{\texorpdfstring{\texttt{unitree\_sdk2\_cpp.robot.go2.ConfigSetParameter}}{unitree\_sdk2\_cpp.robot.go2.ConfigSetParameter}}\label{unitree-sdk2-cpp-robot-go2-configsetparameter}}

SDK 请求参数值类型；当前可能仅提供设计期签名。

\textbf{导入}

\begin{lstlisting}[language=Python]
from unitree_sdk2_cpp.robot.go2 import ConfigSetParameter
\end{lstlisting}

\textbf{基类}：\passthrough{\lstinline!object!}

\textbf{公开属性}

\begin{longtable}[]{@{}
  >{\raggedright\arraybackslash}p{(\columnwidth - 6\tabcolsep) * \real{0.18}}
  >{\raggedright\arraybackslash}p{(\columnwidth - 6\tabcolsep) * \real{0.24}}
  >{\raggedright\arraybackslash}p{(\columnwidth - 6\tabcolsep) * \real{0.18}}
  >{\raggedright\arraybackslash}p{(\columnwidth - 6\tabcolsep) * \real{0.40}}@{}}
\toprule
名称 & Python 类型 & 含义 & 用法 \\
\midrule
\endhead
\passthrough{\lstinline!name!} & \passthrough{\lstinline!str!} & SDK
对象、配置项、服务或动作的名称。允许值由调用它的具体接口定义。 &
\passthrough{\lstinline!value = obj.name!} /
\passthrough{\lstinline!obj.name = value!} \\
\passthrough{\lstinline!content!} & \passthrough{\lstinline!str!} &
配置、请求或序列化内容。具体格式由对应服务协议定义。 &
\passthrough{\lstinline!value = obj.content!} /
\passthrough{\lstinline!obj.content = value!} \\
\bottomrule
\end{longtable}

\textbf{方法索引}

\begin{longtable}[]{@{}
  >{\raggedright\arraybackslash}p{(\columnwidth - 6\tabcolsep) * \real{0.18}}
  >{\raggedright\arraybackslash}p{(\columnwidth - 6\tabcolsep) * \real{0.24}}
  >{\raggedright\arraybackslash}p{(\columnwidth - 6\tabcolsep) * \real{0.18}}
  >{\raggedright\arraybackslash}p{(\columnwidth - 6\tabcolsep) * \real{0.40}}@{}}
\toprule
方法 & Python 签名 & 状态 & 安全分类 \\
\midrule
\endhead
\protect\hyperlink{unitree-sdk2-cpp-robot-go2-configsetparameter-dunder-init-1}{\passthrough{\lstinline!\_\_init\_\_!}}
& \passthrough{\lstinline!def \_\_init\_\_(self) -> None!} &
\passthrough{\lstinline!SIGNATURE\_ONLY!} &
\passthrough{\lstinline!CONSTRUCTION!} \\
\protect\hyperlink{unitree-sdk2-cpp-robot-go2-configsetparameter-from-json-1}{\passthrough{\lstinline!from\_json!}}
&
\passthrough{\lstinline!def from\_json(self, json: dict[str, Any]) -> None!}
& \passthrough{\lstinline!SIGNATURE\_ONLY!} &
\passthrough{\lstinline!UNCLASSIFIED!} \\
\protect\hyperlink{unitree-sdk2-cpp-robot-go2-configsetparameter-to-json-1}{\passthrough{\lstinline!to\_json!}}
&
\passthrough{\lstinline!def to\_json(self, json: dict[str, Any]) -> None!}
& \passthrough{\lstinline!SIGNATURE\_ONLY!} &
\passthrough{\lstinline!UNCLASSIFIED!} \\
\bottomrule
\end{longtable}

\hypertarget{unitree-sdk2-cpp-robot-go2-configsetparameter-dunder-init-1}{%
\paragraph{\texorpdfstring{\texttt{unitree\_sdk2\_cpp.robot.go2.ConfigSetParameter.\_\_init\_\_}}{unitree\_sdk2\_cpp.robot.go2.ConfigSetParameter.\_\_init\_\_}}\label{unitree-sdk2-cpp-robot-go2-configsetparameter-dunder-init-1}}

初始化 \passthrough{\lstinline!ConfigSetParameter!}
实例。是否能实际构造取决于下方可用性状态。

\textbf{签名}

\begin{lstlisting}[language=Python]
def __init__(self) -> None
\end{lstlisting}

\textbf{可用性与安全性}

\textbf{\passthrough{\lstinline!SIGNATURE\_ONLY!} /
\passthrough{\lstinline!CONSTRUCTION!}}：当前只有设计期签名，不能假设已存在于运行时扩展。

\textbf{参数}

无显式参数。实例方法中的 \passthrough{\lstinline!self!} 由 Python
自动传入。

\textbf{返回值}

\begin{longtable}[]{@{}
  >{\raggedright\arraybackslash}p{(\columnwidth - 4\tabcolsep) * \real{0.18}}
  >{\raggedright\arraybackslash}p{(\columnwidth - 4\tabcolsep) * \real{0.20}}
  >{\raggedright\arraybackslash}p{(\columnwidth - 4\tabcolsep) * \real{0.62}}@{}}
\toprule
位置 & 类型 & 含义 \\
\midrule
\endhead
无 & \passthrough{\lstinline!None!} &
\passthrough{\lstinline!\_\_init\_\_!}
本身不返回值；调用类对象时，在构造函数可用的前提下得到该类实例。 \\
\bottomrule
\end{longtable}

\textbf{对应 C++}

\begin{itemize}
\tightlist
\item
  类：\passthrough{\lstinline!unitree::robot::go2::ConfigSetParameter!}
\item
  签名：\passthrough{\lstinline!ConfigSetParameter()!}
\item
  绑定策略：\passthrough{\lstinline!未单独标注!}
\item
  声明位置：\passthrough{\lstinline!include/unitree/robot/go2/config/config\_api.hpp:68!}
\end{itemize}

\textbf{说明}

\begin{itemize}
\tightlist
\item
  示例是局部调用片段；\passthrough{\lstinline!obj!}
  和参数变量表示已经按上层业务校验并准备好的对象。
\item
  该条目可以被编辑器和类型检查器识别，但当前运行时可能在属性查找阶段直接失败。
\end{itemize}

\textbf{用法}

\begin{lstlisting}[language=Python]
from typing import TYPE_CHECKING

if TYPE_CHECKING:
    def planned_construct() -> ConfigSetParameter:
        return ConfigSetParameter()
\end{lstlisting}

\begin{quote}
{[}!CAUTION{]} 上面的 \passthrough{\lstinline!TYPE\_CHECKING!}
示例只表达计划调用形态。该分支在普通 Python
运行时为假，不代表当前扩展能执行此接口。
\end{quote}

\hypertarget{unitree-sdk2-cpp-robot-go2-configsetparameter-from-json-1}{%
\paragraph{\texorpdfstring{\texttt{unitree\_sdk2\_cpp.robot.go2.ConfigSetParameter.from\_json}}{unitree\_sdk2\_cpp.robot.go2.ConfigSetParameter.from\_json}}\label{unitree-sdk2-cpp-robot-go2-configsetparameter-from-json-1}}

计划从 JSON 风格字典读取字段并更新当前 SDK 值对象。

\textbf{签名}

\begin{lstlisting}[language=Python]
def from_json(self, json: dict[str, Any]) -> None
\end{lstlisting}

\textbf{可用性与安全性}

\textbf{\passthrough{\lstinline!SIGNATURE\_ONLY!} /
\passthrough{\lstinline!UNCLASSIFIED!}}：当前只有设计期签名，不能假设已存在于运行时扩展。

\textbf{参数}

\begin{longtable}[]{@{}
  >{\raggedright\arraybackslash}p{(\columnwidth - 6\tabcolsep) * \real{0.18}}
  >{\raggedright\arraybackslash}p{(\columnwidth - 6\tabcolsep) * \real{0.24}}
  >{\raggedright\arraybackslash}p{(\columnwidth - 6\tabcolsep) * \real{0.18}}
  >{\raggedright\arraybackslash}p{(\columnwidth - 6\tabcolsep) * \real{0.40}}@{}}
\toprule
名称 & Python 类型 & 默认值 & 含义 \\
\midrule
\endhead
\passthrough{\lstinline!json!} &
\passthrough{\lstinline!dict[str, Any]!} & 必填 & JSON
风格字典，键和值必须满足对应 C++ \passthrough{\lstinline!JsonMap!}
数据结构的约束。 对应 C++ 参数
\passthrough{\lstinline!json: common::JsonMap \&!}。 \\
\bottomrule
\end{longtable}

\textbf{返回值}

\begin{longtable}[]{@{}
  >{\raggedright\arraybackslash}p{(\columnwidth - 4\tabcolsep) * \real{0.18}}
  >{\raggedright\arraybackslash}p{(\columnwidth - 4\tabcolsep) * \real{0.20}}
  >{\raggedright\arraybackslash}p{(\columnwidth - 4\tabcolsep) * \real{0.62}}@{}}
\toprule
位置 & 类型 & 含义 \\
\midrule
\endhead
无 & \passthrough{\lstinline!None!} & 该方法不返回 Python
值；失败可能表现为异常或底层状态变化。 \\
\bottomrule
\end{longtable}

\textbf{对应 C++}

\begin{itemize}
\tightlist
\item
  类：\passthrough{\lstinline!unitree::robot::go2::ConfigSetParameter!}
\item
  签名：\passthrough{\lstinline!fromJson(common::JsonMap \&)!}
\item
  绑定策略：\passthrough{\lstinline!SIGNATURE\_PREVIEW!}
\item
  声明位置：\passthrough{\lstinline!include/unitree/robot/go2/config/config\_api.hpp:74!}
\end{itemize}

\textbf{说明}

\begin{itemize}
\tightlist
\item
  示例是局部调用片段；\passthrough{\lstinline!obj!}
  和参数变量表示已经按上层业务校验并准备好的对象。
\item
  该条目可以被编辑器和类型检查器识别，但当前运行时可能在属性查找阶段直接失败。
\end{itemize}

\textbf{用法}

\begin{lstlisting}[language=Python]
from typing import TYPE_CHECKING

if TYPE_CHECKING:
    def planned_from_json(obj: ConfigSetParameter, json: dict[str, Any]) -> None:
        obj.from_json(json=json)
\end{lstlisting}

\begin{quote}
{[}!CAUTION{]} 上面的 \passthrough{\lstinline!TYPE\_CHECKING!}
示例只表达计划调用形态。该分支在普通 Python
运行时为假，不代表当前扩展能执行此接口。
\end{quote}

\hypertarget{unitree-sdk2-cpp-robot-go2-configsetparameter-to-json-1}{%
\paragraph{\texorpdfstring{\texttt{unitree\_sdk2\_cpp.robot.go2.ConfigSetParameter.to\_json}}{unitree\_sdk2\_cpp.robot.go2.ConfigSetParameter.to\_json}}\label{unitree-sdk2-cpp-robot-go2-configsetparameter-to-json-1}}

计划把当前 SDK 值对象写入 JSON 风格字典。

\textbf{签名}

\begin{lstlisting}[language=Python]
def to_json(self, json: dict[str, Any]) -> None
\end{lstlisting}

\textbf{可用性与安全性}

\textbf{\passthrough{\lstinline!SIGNATURE\_ONLY!} /
\passthrough{\lstinline!UNCLASSIFIED!}}：当前只有设计期签名，不能假设已存在于运行时扩展。

\textbf{参数}

\begin{longtable}[]{@{}
  >{\raggedright\arraybackslash}p{(\columnwidth - 6\tabcolsep) * \real{0.18}}
  >{\raggedright\arraybackslash}p{(\columnwidth - 6\tabcolsep) * \real{0.24}}
  >{\raggedright\arraybackslash}p{(\columnwidth - 6\tabcolsep) * \real{0.18}}
  >{\raggedright\arraybackslash}p{(\columnwidth - 6\tabcolsep) * \real{0.40}}@{}}
\toprule
名称 & Python 类型 & 默认值 & 含义 \\
\midrule
\endhead
\passthrough{\lstinline!json!} &
\passthrough{\lstinline!dict[str, Any]!} & 必填 & JSON
风格字典，键和值必须满足对应 C++ \passthrough{\lstinline!JsonMap!}
数据结构的约束。 对应 C++ 参数
\passthrough{\lstinline!json: common::JsonMap \&!}。 \\
\bottomrule
\end{longtable}

\textbf{返回值}

\begin{longtable}[]{@{}
  >{\raggedright\arraybackslash}p{(\columnwidth - 4\tabcolsep) * \real{0.18}}
  >{\raggedright\arraybackslash}p{(\columnwidth - 4\tabcolsep) * \real{0.20}}
  >{\raggedright\arraybackslash}p{(\columnwidth - 4\tabcolsep) * \real{0.62}}@{}}
\toprule
位置 & 类型 & 含义 \\
\midrule
\endhead
无 & \passthrough{\lstinline!None!} & 该方法不返回 Python
值；失败可能表现为异常或底层状态变化。 \\
\bottomrule
\end{longtable}

\textbf{对应 C++}

\begin{itemize}
\tightlist
\item
  类：\passthrough{\lstinline!unitree::robot::go2::ConfigSetParameter!}
\item
  签名：\passthrough{\lstinline!toJson(common::JsonMap \&) const!}
\item
  绑定策略：\passthrough{\lstinline!SIGNATURE\_PREVIEW!}
\item
  声明位置：\passthrough{\lstinline!include/unitree/robot/go2/config/config\_api.hpp:80!}
\end{itemize}

\textbf{说明}

\begin{itemize}
\tightlist
\item
  示例是局部调用片段；\passthrough{\lstinline!obj!}
  和参数变量表示已经按上层业务校验并准备好的对象。
\item
  该条目可以被编辑器和类型检查器识别，但当前运行时可能在属性查找阶段直接失败。
\end{itemize}

\textbf{用法}

\begin{lstlisting}[language=Python]
from typing import TYPE_CHECKING

if TYPE_CHECKING:
    def planned_to_json(obj: ConfigSetParameter, json: dict[str, Any]) -> None:
        obj.to_json(json=json)
\end{lstlisting}

\begin{quote}
{[}!CAUTION{]} 上面的 \passthrough{\lstinline!TYPE\_CHECKING!}
示例只表达计划调用形态。该分支在普通 Python
运行时为假，不代表当前扩展能执行此接口。
\end{quote}

\hypertarget{unitree-sdk2-cpp-robot-go2-jsonizecommobjint}{%
\subsubsection{\texorpdfstring{\texttt{unitree\_sdk2\_cpp.robot.go2.JsonizeCommObjInt}}{unitree\_sdk2\_cpp.robot.go2.JsonizeCommObjInt}}\label{unitree-sdk2-cpp-robot-go2-jsonizecommobjint}}

SDK 类型签名预览。请逐项查看构造函数和方法的可用性。

\textbf{导入}

\begin{lstlisting}[language=Python]
from unitree_sdk2_cpp.robot.go2 import JsonizeCommObjInt
\end{lstlisting}

\textbf{基类}：\passthrough{\lstinline!object!}

\textbf{公开属性}

\begin{longtable}[]{@{}
  >{\raggedright\arraybackslash}p{(\columnwidth - 6\tabcolsep) * \real{0.18}}
  >{\raggedright\arraybackslash}p{(\columnwidth - 6\tabcolsep) * \real{0.24}}
  >{\raggedright\arraybackslash}p{(\columnwidth - 6\tabcolsep) * \real{0.18}}
  >{\raggedright\arraybackslash}p{(\columnwidth - 6\tabcolsep) * \real{0.40}}@{}}
\toprule
名称 & Python 类型 & 含义 & 用法 \\
\midrule
\endhead
\passthrough{\lstinline!value!} & \passthrough{\lstinline!int!} &
要写入或传递的新值，必须符合该参数的 Python 类型以及底层 C++ 范围约束。
& \passthrough{\lstinline!value = obj.value!} /
\passthrough{\lstinline!obj.value = value!} \\
\passthrough{\lstinline!name!} & \passthrough{\lstinline!str!} & SDK
对象、配置项、服务或动作的名称。允许值由调用它的具体接口定义。 &
\passthrough{\lstinline!value = obj.name!} /
\passthrough{\lstinline!obj.name = value!} \\
\bottomrule
\end{longtable}

\textbf{方法索引}

\begin{longtable}[]{@{}
  >{\raggedright\arraybackslash}p{(\columnwidth - 6\tabcolsep) * \real{0.18}}
  >{\raggedright\arraybackslash}p{(\columnwidth - 6\tabcolsep) * \real{0.24}}
  >{\raggedright\arraybackslash}p{(\columnwidth - 6\tabcolsep) * \real{0.18}}
  >{\raggedright\arraybackslash}p{(\columnwidth - 6\tabcolsep) * \real{0.40}}@{}}
\toprule
方法 & Python 签名 & 状态 & 安全分类 \\
\midrule
\endhead
\protect\hyperlink{unitree-sdk2-cpp-robot-go2-jsonizecommobjint-dunder-init-1}{\passthrough{\lstinline!\_\_init\_\_!}}
& \passthrough{\lstinline!def \_\_init\_\_(self) -> None!} &
\passthrough{\lstinline!SIGNATURE\_ONLY!} &
\passthrough{\lstinline!CONSTRUCTION!} \\
\protect\hyperlink{unitree-sdk2-cpp-robot-go2-jsonizecommobjint-from-json-1}{\passthrough{\lstinline!from\_json!}}
&
\passthrough{\lstinline!def from\_json(self, json: dict[str, Any]) -> None!}
& \passthrough{\lstinline!SIGNATURE\_ONLY!} &
\passthrough{\lstinline!UNCLASSIFIED!} \\
\protect\hyperlink{unitree-sdk2-cpp-robot-go2-jsonizecommobjint-to-json-1}{\passthrough{\lstinline!to\_json!}}
&
\passthrough{\lstinline!def to\_json(self, json: dict[str, Any]) -> None!}
& \passthrough{\lstinline!SIGNATURE\_ONLY!} &
\passthrough{\lstinline!UNCLASSIFIED!} \\
\bottomrule
\end{longtable}

\hypertarget{unitree-sdk2-cpp-robot-go2-jsonizecommobjint-dunder-init-1}{%
\paragraph{\texorpdfstring{\texttt{unitree\_sdk2\_cpp.robot.go2.JsonizeCommObjInt.\_\_init\_\_}}{unitree\_sdk2\_cpp.robot.go2.JsonizeCommObjInt.\_\_init\_\_}}\label{unitree-sdk2-cpp-robot-go2-jsonizecommobjint-dunder-init-1}}

初始化 \passthrough{\lstinline!JsonizeCommObjInt!}
实例。是否能实际构造取决于下方可用性状态。

\textbf{签名}

\begin{lstlisting}[language=Python]
def __init__(self) -> None
\end{lstlisting}

\textbf{可用性与安全性}

\textbf{\passthrough{\lstinline!SIGNATURE\_ONLY!} /
\passthrough{\lstinline!CONSTRUCTION!}}：当前只有设计期签名，不能假设已存在于运行时扩展。

\textbf{参数}

无显式参数。实例方法中的 \passthrough{\lstinline!self!} 由 Python
自动传入。

\textbf{返回值}

\begin{longtable}[]{@{}
  >{\raggedright\arraybackslash}p{(\columnwidth - 4\tabcolsep) * \real{0.18}}
  >{\raggedright\arraybackslash}p{(\columnwidth - 4\tabcolsep) * \real{0.20}}
  >{\raggedright\arraybackslash}p{(\columnwidth - 4\tabcolsep) * \real{0.62}}@{}}
\toprule
位置 & 类型 & 含义 \\
\midrule
\endhead
无 & \passthrough{\lstinline!None!} &
\passthrough{\lstinline!\_\_init\_\_!}
本身不返回值；调用类对象时，在构造函数可用的前提下得到该类实例。 \\
\bottomrule
\end{longtable}

\textbf{对应 C++}

\begin{itemize}
\tightlist
\item
  类：\passthrough{\lstinline!unitree::robot::go2::JsonizeCommObjInt!}
\item
  签名：\passthrough{\lstinline!JsonizeCommObjInt()!}
\item
  绑定策略：\passthrough{\lstinline!未单独标注!}
\item
  声明位置：\passthrough{\lstinline!include/unitree/robot/go2/public/jsonize\_type.hpp:287!}
\end{itemize}

\textbf{说明}

\begin{itemize}
\tightlist
\item
  示例是局部调用片段；\passthrough{\lstinline!obj!}
  和参数变量表示已经按上层业务校验并准备好的对象。
\item
  该条目可以被编辑器和类型检查器识别，但当前运行时可能在属性查找阶段直接失败。
\end{itemize}

\textbf{用法}

\begin{lstlisting}[language=Python]
from typing import TYPE_CHECKING

if TYPE_CHECKING:
    def planned_construct() -> JsonizeCommObjInt:
        return JsonizeCommObjInt()
\end{lstlisting}

\begin{quote}
{[}!CAUTION{]} 上面的 \passthrough{\lstinline!TYPE\_CHECKING!}
示例只表达计划调用形态。该分支在普通 Python
运行时为假，不代表当前扩展能执行此接口。
\end{quote}

\hypertarget{unitree-sdk2-cpp-robot-go2-jsonizecommobjint-from-json-1}{%
\paragraph{\texorpdfstring{\texttt{unitree\_sdk2\_cpp.robot.go2.JsonizeCommObjInt.from\_json}}{unitree\_sdk2\_cpp.robot.go2.JsonizeCommObjInt.from\_json}}\label{unitree-sdk2-cpp-robot-go2-jsonizecommobjint-from-json-1}}

计划从 JSON 风格字典读取字段并更新当前 SDK 值对象。

\textbf{签名}

\begin{lstlisting}[language=Python]
def from_json(self, json: dict[str, Any]) -> None
\end{lstlisting}

\textbf{可用性与安全性}

\textbf{\passthrough{\lstinline!SIGNATURE\_ONLY!} /
\passthrough{\lstinline!UNCLASSIFIED!}}：当前只有设计期签名，不能假设已存在于运行时扩展。

\textbf{参数}

\begin{longtable}[]{@{}
  >{\raggedright\arraybackslash}p{(\columnwidth - 6\tabcolsep) * \real{0.18}}
  >{\raggedright\arraybackslash}p{(\columnwidth - 6\tabcolsep) * \real{0.24}}
  >{\raggedright\arraybackslash}p{(\columnwidth - 6\tabcolsep) * \real{0.18}}
  >{\raggedright\arraybackslash}p{(\columnwidth - 6\tabcolsep) * \real{0.40}}@{}}
\toprule
名称 & Python 类型 & 默认值 & 含义 \\
\midrule
\endhead
\passthrough{\lstinline!json!} &
\passthrough{\lstinline!dict[str, Any]!} & 必填 & JSON
风格字典，键和值必须满足对应 C++ \passthrough{\lstinline!JsonMap!}
数据结构的约束。 对应 C++ 参数
\passthrough{\lstinline!json: common::JsonMap \&!}。 \\
\bottomrule
\end{longtable}

\textbf{返回值}

\begin{longtable}[]{@{}
  >{\raggedright\arraybackslash}p{(\columnwidth - 4\tabcolsep) * \real{0.18}}
  >{\raggedright\arraybackslash}p{(\columnwidth - 4\tabcolsep) * \real{0.20}}
  >{\raggedright\arraybackslash}p{(\columnwidth - 4\tabcolsep) * \real{0.62}}@{}}
\toprule
位置 & 类型 & 含义 \\
\midrule
\endhead
无 & \passthrough{\lstinline!None!} & 该方法不返回 Python
值；失败可能表现为异常或底层状态变化。 \\
\bottomrule
\end{longtable}

\textbf{对应 C++}

\begin{itemize}
\tightlist
\item
  类：\passthrough{\lstinline!unitree::robot::go2::JsonizeCommObjInt!}
\item
  签名：\passthrough{\lstinline!fromJson(common::JsonMap \&)!}
\item
  绑定策略：\passthrough{\lstinline!SIGNATURE\_PREVIEW!}
\item
  声明位置：\passthrough{\lstinline!include/unitree/robot/go2/public/jsonize\_type.hpp:293!}
\end{itemize}

\textbf{说明}

\begin{itemize}
\tightlist
\item
  示例是局部调用片段；\passthrough{\lstinline!obj!}
  和参数变量表示已经按上层业务校验并准备好的对象。
\item
  该条目可以被编辑器和类型检查器识别，但当前运行时可能在属性查找阶段直接失败。
\end{itemize}

\textbf{用法}

\begin{lstlisting}[language=Python]
from typing import TYPE_CHECKING

if TYPE_CHECKING:
    def planned_from_json(obj: JsonizeCommObjInt, json: dict[str, Any]) -> None:
        obj.from_json(json=json)
\end{lstlisting}

\begin{quote}
{[}!CAUTION{]} 上面的 \passthrough{\lstinline!TYPE\_CHECKING!}
示例只表达计划调用形态。该分支在普通 Python
运行时为假，不代表当前扩展能执行此接口。
\end{quote}

\hypertarget{unitree-sdk2-cpp-robot-go2-jsonizecommobjint-to-json-1}{%
\paragraph{\texorpdfstring{\texttt{unitree\_sdk2\_cpp.robot.go2.JsonizeCommObjInt.to\_json}}{unitree\_sdk2\_cpp.robot.go2.JsonizeCommObjInt.to\_json}}\label{unitree-sdk2-cpp-robot-go2-jsonizecommobjint-to-json-1}}

计划把当前 SDK 值对象写入 JSON 风格字典。

\textbf{签名}

\begin{lstlisting}[language=Python]
def to_json(self, json: dict[str, Any]) -> None
\end{lstlisting}

\textbf{可用性与安全性}

\textbf{\passthrough{\lstinline!SIGNATURE\_ONLY!} /
\passthrough{\lstinline!UNCLASSIFIED!}}：当前只有设计期签名，不能假设已存在于运行时扩展。

\textbf{参数}

\begin{longtable}[]{@{}
  >{\raggedright\arraybackslash}p{(\columnwidth - 6\tabcolsep) * \real{0.18}}
  >{\raggedright\arraybackslash}p{(\columnwidth - 6\tabcolsep) * \real{0.24}}
  >{\raggedright\arraybackslash}p{(\columnwidth - 6\tabcolsep) * \real{0.18}}
  >{\raggedright\arraybackslash}p{(\columnwidth - 6\tabcolsep) * \real{0.40}}@{}}
\toprule
名称 & Python 类型 & 默认值 & 含义 \\
\midrule
\endhead
\passthrough{\lstinline!json!} &
\passthrough{\lstinline!dict[str, Any]!} & 必填 & JSON
风格字典，键和值必须满足对应 C++ \passthrough{\lstinline!JsonMap!}
数据结构的约束。 对应 C++ 参数
\passthrough{\lstinline!json: common::JsonMap \&!}。 \\
\bottomrule
\end{longtable}

\textbf{返回值}

\begin{longtable}[]{@{}
  >{\raggedright\arraybackslash}p{(\columnwidth - 4\tabcolsep) * \real{0.18}}
  >{\raggedright\arraybackslash}p{(\columnwidth - 4\tabcolsep) * \real{0.20}}
  >{\raggedright\arraybackslash}p{(\columnwidth - 4\tabcolsep) * \real{0.62}}@{}}
\toprule
位置 & 类型 & 含义 \\
\midrule
\endhead
无 & \passthrough{\lstinline!None!} & 该方法不返回 Python
值；失败可能表现为异常或底层状态变化。 \\
\bottomrule
\end{longtable}

\textbf{对应 C++}

\begin{itemize}
\tightlist
\item
  类：\passthrough{\lstinline!unitree::robot::go2::JsonizeCommObjInt!}
\item
  签名：\passthrough{\lstinline!toJson(common::JsonMap \&) const!}
\item
  绑定策略：\passthrough{\lstinline!SIGNATURE\_PREVIEW!}
\item
  声明位置：\passthrough{\lstinline!include/unitree/robot/go2/public/jsonize\_type.hpp:298!}
\end{itemize}

\textbf{说明}

\begin{itemize}
\tightlist
\item
  示例是局部调用片段；\passthrough{\lstinline!obj!}
  和参数变量表示已经按上层业务校验并准备好的对象。
\item
  该条目可以被编辑器和类型检查器识别，但当前运行时可能在属性查找阶段直接失败。
\end{itemize}

\textbf{用法}

\begin{lstlisting}[language=Python]
from typing import TYPE_CHECKING

if TYPE_CHECKING:
    def planned_to_json(obj: JsonizeCommObjInt, json: dict[str, Any]) -> None:
        obj.to_json(json=json)
\end{lstlisting}

\begin{quote}
{[}!CAUTION{]} 上面的 \passthrough{\lstinline!TYPE\_CHECKING!}
示例只表达计划调用形态。该分支在普通 Python
运行时为假，不代表当前扩展能执行此接口。
\end{quote}

\hypertarget{unitree-sdk2-cpp-robot-go2-jsonizeconfigmeta}{%
\subsubsection{\texorpdfstring{\texttt{unitree\_sdk2\_cpp.robot.go2.JsonizeConfigMeta}}{unitree\_sdk2\_cpp.robot.go2.JsonizeConfigMeta}}\label{unitree-sdk2-cpp-robot-go2-jsonizeconfigmeta}}

SDK 数据值类型；公开字段和转换方法见下文。

\textbf{导入}

\begin{lstlisting}[language=Python]
from unitree_sdk2_cpp.robot.go2 import JsonizeConfigMeta
\end{lstlisting}

\textbf{基类}：\passthrough{\lstinline!object!}

\textbf{公开属性}

\begin{longtable}[]{@{}
  >{\raggedright\arraybackslash}p{(\columnwidth - 6\tabcolsep) * \real{0.18}}
  >{\raggedright\arraybackslash}p{(\columnwidth - 6\tabcolsep) * \real{0.24}}
  >{\raggedright\arraybackslash}p{(\columnwidth - 6\tabcolsep) * \real{0.18}}
  >{\raggedright\arraybackslash}p{(\columnwidth - 6\tabcolsep) * \real{0.40}}@{}}
\toprule
名称 & Python 类型 & 含义 & 用法 \\
\midrule
\endhead
\passthrough{\lstinline!name!} & \passthrough{\lstinline!str!} & SDK
对象、配置项、服务或动作的名称。允许值由调用它的具体接口定义。 &
\passthrough{\lstinline!value = obj.name!} /
\passthrough{\lstinline!obj.name = value!} \\
\passthrough{\lstinline!last\_modified!} & \passthrough{\lstinline!str!}
& 传给该接口的 \passthrough{\lstinline!last!}
\passthrough{\lstinline!modified!} 参数，Python 类型为
\passthrough{\lstinline!str!}。精确含义、单位、范围和枚举值需查目标型号对应头文件与协议。
& \passthrough{\lstinline!value = obj.last\_modified!} /
\passthrough{\lstinline!obj.last\_modified = value!} \\
\passthrough{\lstinline!size!} & \passthrough{\lstinline!int!} &
传给该接口的 \passthrough{\lstinline!size!} 参数，Python 类型为
\passthrough{\lstinline!int!}。精确含义、单位、范围和枚举值需查目标型号对应头文件与协议。
& \passthrough{\lstinline!value = obj.size!} /
\passthrough{\lstinline!obj.size = value!} \\
\passthrough{\lstinline!epoch!} & \passthrough{\lstinline!int!} &
传给该接口的 \passthrough{\lstinline!epoch!} 参数，Python 类型为
\passthrough{\lstinline!int!}。精确含义、单位、范围和枚举值需查目标型号对应头文件与协议。
& \passthrough{\lstinline!value = obj.epoch!} /
\passthrough{\lstinline!obj.epoch = value!} \\
\bottomrule
\end{longtable}

\textbf{方法索引}

\begin{longtable}[]{@{}
  >{\raggedright\arraybackslash}p{(\columnwidth - 6\tabcolsep) * \real{0.18}}
  >{\raggedright\arraybackslash}p{(\columnwidth - 6\tabcolsep) * \real{0.24}}
  >{\raggedright\arraybackslash}p{(\columnwidth - 6\tabcolsep) * \real{0.18}}
  >{\raggedright\arraybackslash}p{(\columnwidth - 6\tabcolsep) * \real{0.40}}@{}}
\toprule
方法 & Python 签名 & 状态 & 安全分类 \\
\midrule
\endhead
\protect\hyperlink{unitree-sdk2-cpp-robot-go2-jsonizeconfigmeta-dunder-init-1}{\passthrough{\lstinline!\_\_init\_\_!}}
& \passthrough{\lstinline!def \_\_init\_\_(self) -> None!} &
\passthrough{\lstinline!SIGNATURE\_ONLY!} &
\passthrough{\lstinline!CONSTRUCTION!} \\
\protect\hyperlink{unitree-sdk2-cpp-robot-go2-jsonizeconfigmeta-from-json-1}{\passthrough{\lstinline!from\_json!}}
&
\passthrough{\lstinline!def from\_json(self, json: dict[str, Any]) -> None!}
& \passthrough{\lstinline!SIGNATURE\_ONLY!} &
\passthrough{\lstinline!UNCLASSIFIED!} \\
\protect\hyperlink{unitree-sdk2-cpp-robot-go2-jsonizeconfigmeta-to-json-1}{\passthrough{\lstinline!to\_json!}}
&
\passthrough{\lstinline!def to\_json(self, json: dict[str, Any]) -> None!}
& \passthrough{\lstinline!SIGNATURE\_ONLY!} &
\passthrough{\lstinline!UNCLASSIFIED!} \\
\bottomrule
\end{longtable}

\hypertarget{unitree-sdk2-cpp-robot-go2-jsonizeconfigmeta-dunder-init-1}{%
\paragraph{\texorpdfstring{\texttt{unitree\_sdk2\_cpp.robot.go2.JsonizeConfigMeta.\_\_init\_\_}}{unitree\_sdk2\_cpp.robot.go2.JsonizeConfigMeta.\_\_init\_\_}}\label{unitree-sdk2-cpp-robot-go2-jsonizeconfigmeta-dunder-init-1}}

初始化 \passthrough{\lstinline!JsonizeConfigMeta!}
实例。是否能实际构造取决于下方可用性状态。

\textbf{签名}

\begin{lstlisting}[language=Python]
def __init__(self) -> None
\end{lstlisting}

\textbf{可用性与安全性}

\textbf{\passthrough{\lstinline!SIGNATURE\_ONLY!} /
\passthrough{\lstinline!CONSTRUCTION!}}：当前只有设计期签名，不能假设已存在于运行时扩展。

\textbf{参数}

无显式参数。实例方法中的 \passthrough{\lstinline!self!} 由 Python
自动传入。

\textbf{返回值}

\begin{longtable}[]{@{}
  >{\raggedright\arraybackslash}p{(\columnwidth - 4\tabcolsep) * \real{0.18}}
  >{\raggedright\arraybackslash}p{(\columnwidth - 4\tabcolsep) * \real{0.20}}
  >{\raggedright\arraybackslash}p{(\columnwidth - 4\tabcolsep) * \real{0.62}}@{}}
\toprule
位置 & 类型 & 含义 \\
\midrule
\endhead
无 & \passthrough{\lstinline!None!} &
\passthrough{\lstinline!\_\_init\_\_!}
本身不返回值；调用类对象时，在构造函数可用的前提下得到该类实例。 \\
\bottomrule
\end{longtable}

\textbf{对应 C++}

\begin{itemize}
\tightlist
\item
  类：\passthrough{\lstinline!unitree::robot::go2::JsonizeConfigMeta!}
\item
  签名：\passthrough{\lstinline!JsonizeConfigMeta()!}
\item
  绑定策略：\passthrough{\lstinline!未单独标注!}
\item
  声明位置：\passthrough{\lstinline!include/unitree/robot/go2/config/config\_api.hpp:36!}
\end{itemize}

\textbf{说明}

\begin{itemize}
\tightlist
\item
  示例是局部调用片段；\passthrough{\lstinline!obj!}
  和参数变量表示已经按上层业务校验并准备好的对象。
\item
  该条目可以被编辑器和类型检查器识别，但当前运行时可能在属性查找阶段直接失败。
\end{itemize}

\textbf{用法}

\begin{lstlisting}[language=Python]
from typing import TYPE_CHECKING

if TYPE_CHECKING:
    def planned_construct() -> JsonizeConfigMeta:
        return JsonizeConfigMeta()
\end{lstlisting}

\begin{quote}
{[}!CAUTION{]} 上面的 \passthrough{\lstinline!TYPE\_CHECKING!}
示例只表达计划调用形态。该分支在普通 Python
运行时为假，不代表当前扩展能执行此接口。
\end{quote}

\hypertarget{unitree-sdk2-cpp-robot-go2-jsonizeconfigmeta-from-json-1}{%
\paragraph{\texorpdfstring{\texttt{unitree\_sdk2\_cpp.robot.go2.JsonizeConfigMeta.from\_json}}{unitree\_sdk2\_cpp.robot.go2.JsonizeConfigMeta.from\_json}}\label{unitree-sdk2-cpp-robot-go2-jsonizeconfigmeta-from-json-1}}

计划从 JSON 风格字典读取字段并更新当前 SDK 值对象。

\textbf{签名}

\begin{lstlisting}[language=Python]
def from_json(self, json: dict[str, Any]) -> None
\end{lstlisting}

\textbf{可用性与安全性}

\textbf{\passthrough{\lstinline!SIGNATURE\_ONLY!} /
\passthrough{\lstinline!UNCLASSIFIED!}}：当前只有设计期签名，不能假设已存在于运行时扩展。

\textbf{参数}

\begin{longtable}[]{@{}
  >{\raggedright\arraybackslash}p{(\columnwidth - 6\tabcolsep) * \real{0.18}}
  >{\raggedright\arraybackslash}p{(\columnwidth - 6\tabcolsep) * \real{0.24}}
  >{\raggedright\arraybackslash}p{(\columnwidth - 6\tabcolsep) * \real{0.18}}
  >{\raggedright\arraybackslash}p{(\columnwidth - 6\tabcolsep) * \real{0.40}}@{}}
\toprule
名称 & Python 类型 & 默认值 & 含义 \\
\midrule
\endhead
\passthrough{\lstinline!json!} &
\passthrough{\lstinline!dict[str, Any]!} & 必填 & JSON
风格字典，键和值必须满足对应 C++ \passthrough{\lstinline!JsonMap!}
数据结构的约束。 对应 C++ 参数
\passthrough{\lstinline!json: common::JsonMap \&!}。 \\
\bottomrule
\end{longtable}

\textbf{返回值}

\begin{longtable}[]{@{}
  >{\raggedright\arraybackslash}p{(\columnwidth - 4\tabcolsep) * \real{0.18}}
  >{\raggedright\arraybackslash}p{(\columnwidth - 4\tabcolsep) * \real{0.20}}
  >{\raggedright\arraybackslash}p{(\columnwidth - 4\tabcolsep) * \real{0.62}}@{}}
\toprule
位置 & 类型 & 含义 \\
\midrule
\endhead
无 & \passthrough{\lstinline!None!} & 该方法不返回 Python
值；失败可能表现为异常或底层状态变化。 \\
\bottomrule
\end{longtable}

\textbf{对应 C++}

\begin{itemize}
\tightlist
\item
  类：\passthrough{\lstinline!unitree::robot::go2::JsonizeConfigMeta!}
\item
  签名：\passthrough{\lstinline!fromJson(common::JsonMap \&)!}
\item
  绑定策略：\passthrough{\lstinline!SIGNATURE\_PREVIEW!}
\item
  声明位置：\passthrough{\lstinline!include/unitree/robot/go2/config/config\_api.hpp:42!}
\end{itemize}

\textbf{说明}

\begin{itemize}
\tightlist
\item
  示例是局部调用片段；\passthrough{\lstinline!obj!}
  和参数变量表示已经按上层业务校验并准备好的对象。
\item
  该条目可以被编辑器和类型检查器识别，但当前运行时可能在属性查找阶段直接失败。
\end{itemize}

\textbf{用法}

\begin{lstlisting}[language=Python]
from typing import TYPE_CHECKING

if TYPE_CHECKING:
    def planned_from_json(obj: JsonizeConfigMeta, json: dict[str, Any]) -> None:
        obj.from_json(json=json)
\end{lstlisting}

\begin{quote}
{[}!CAUTION{]} 上面的 \passthrough{\lstinline!TYPE\_CHECKING!}
示例只表达计划调用形态。该分支在普通 Python
运行时为假，不代表当前扩展能执行此接口。
\end{quote}

\hypertarget{unitree-sdk2-cpp-robot-go2-jsonizeconfigmeta-to-json-1}{%
\paragraph{\texorpdfstring{\texttt{unitree\_sdk2\_cpp.robot.go2.JsonizeConfigMeta.to\_json}}{unitree\_sdk2\_cpp.robot.go2.JsonizeConfigMeta.to\_json}}\label{unitree-sdk2-cpp-robot-go2-jsonizeconfigmeta-to-json-1}}

计划把当前 SDK 值对象写入 JSON 风格字典。

\textbf{签名}

\begin{lstlisting}[language=Python]
def to_json(self, json: dict[str, Any]) -> None
\end{lstlisting}

\textbf{可用性与安全性}

\textbf{\passthrough{\lstinline!SIGNATURE\_ONLY!} /
\passthrough{\lstinline!UNCLASSIFIED!}}：当前只有设计期签名，不能假设已存在于运行时扩展。

\textbf{参数}

\begin{longtable}[]{@{}
  >{\raggedright\arraybackslash}p{(\columnwidth - 6\tabcolsep) * \real{0.18}}
  >{\raggedright\arraybackslash}p{(\columnwidth - 6\tabcolsep) * \real{0.24}}
  >{\raggedright\arraybackslash}p{(\columnwidth - 6\tabcolsep) * \real{0.18}}
  >{\raggedright\arraybackslash}p{(\columnwidth - 6\tabcolsep) * \real{0.40}}@{}}
\toprule
名称 & Python 类型 & 默认值 & 含义 \\
\midrule
\endhead
\passthrough{\lstinline!json!} &
\passthrough{\lstinline!dict[str, Any]!} & 必填 & JSON
风格字典，键和值必须满足对应 C++ \passthrough{\lstinline!JsonMap!}
数据结构的约束。 对应 C++ 参数
\passthrough{\lstinline!json: common::JsonMap \&!}。 \\
\bottomrule
\end{longtable}

\textbf{返回值}

\begin{longtable}[]{@{}
  >{\raggedright\arraybackslash}p{(\columnwidth - 4\tabcolsep) * \real{0.18}}
  >{\raggedright\arraybackslash}p{(\columnwidth - 4\tabcolsep) * \real{0.20}}
  >{\raggedright\arraybackslash}p{(\columnwidth - 4\tabcolsep) * \real{0.62}}@{}}
\toprule
位置 & 类型 & 含义 \\
\midrule
\endhead
无 & \passthrough{\lstinline!None!} & 该方法不返回 Python
值；失败可能表现为异常或底层状态变化。 \\
\bottomrule
\end{longtable}

\textbf{对应 C++}

\begin{itemize}
\tightlist
\item
  类：\passthrough{\lstinline!unitree::robot::go2::JsonizeConfigMeta!}
\item
  签名：\passthrough{\lstinline!toJson(common::JsonMap \&) const!}
\item
  绑定策略：\passthrough{\lstinline!SIGNATURE\_PREVIEW!}
\item
  声明位置：\passthrough{\lstinline!include/unitree/robot/go2/config/config\_api.hpp:50!}
\end{itemize}

\textbf{说明}

\begin{itemize}
\tightlist
\item
  示例是局部调用片段；\passthrough{\lstinline!obj!}
  和参数变量表示已经按上层业务校验并准备好的对象。
\item
  该条目可以被编辑器和类型检查器识别，但当前运行时可能在属性查找阶段直接失败。
\end{itemize}

\textbf{用法}

\begin{lstlisting}[language=Python]
from typing import TYPE_CHECKING

if TYPE_CHECKING:
    def planned_to_json(obj: JsonizeConfigMeta, json: dict[str, Any]) -> None:
        obj.to_json(json=json)
\end{lstlisting}

\begin{quote}
{[}!CAUTION{]} 上面的 \passthrough{\lstinline!TYPE\_CHECKING!}
示例只表达计划调用形态。该分支在普通 Python
运行时为假，不代表当前扩展能执行此接口。
\end{quote}

\hypertarget{unitree-sdk2-cpp-robot-go2-jsonizedatabool}{%
\subsubsection{\texorpdfstring{\texttt{unitree\_sdk2\_cpp.robot.go2.JsonizeDataBool}}{unitree\_sdk2\_cpp.robot.go2.JsonizeDataBool}}\label{unitree-sdk2-cpp-robot-go2-jsonizedatabool}}

SDK 类型签名预览。请逐项查看构造函数和方法的可用性。

\textbf{导入}

\begin{lstlisting}[language=Python]
from unitree_sdk2_cpp.robot.go2 import JsonizeDataBool
\end{lstlisting}

\textbf{基类}：\passthrough{\lstinline!object!}

\textbf{公开属性}

\begin{longtable}[]{@{}
  >{\raggedright\arraybackslash}p{(\columnwidth - 6\tabcolsep) * \real{0.18}}
  >{\raggedright\arraybackslash}p{(\columnwidth - 6\tabcolsep) * \real{0.24}}
  >{\raggedright\arraybackslash}p{(\columnwidth - 6\tabcolsep) * \real{0.18}}
  >{\raggedright\arraybackslash}p{(\columnwidth - 6\tabcolsep) * \real{0.40}}@{}}
\toprule
名称 & Python 类型 & 含义 & 用法 \\
\midrule
\endhead
\passthrough{\lstinline!data!} & \passthrough{\lstinline!bool!} &
传给该接口的 数据 参数，Python 类型为
\passthrough{\lstinline!bool!}。精确含义、单位、范围和枚举值需查目标型号对应头文件与协议。
& \passthrough{\lstinline!value = obj.data!} /
\passthrough{\lstinline!obj.data = value!} \\
\bottomrule
\end{longtable}

\textbf{方法索引}

\begin{longtable}[]{@{}
  >{\raggedright\arraybackslash}p{(\columnwidth - 6\tabcolsep) * \real{0.18}}
  >{\raggedright\arraybackslash}p{(\columnwidth - 6\tabcolsep) * \real{0.24}}
  >{\raggedright\arraybackslash}p{(\columnwidth - 6\tabcolsep) * \real{0.18}}
  >{\raggedright\arraybackslash}p{(\columnwidth - 6\tabcolsep) * \real{0.40}}@{}}
\toprule
方法 & Python 签名 & 状态 & 安全分类 \\
\midrule
\endhead
\protect\hyperlink{unitree-sdk2-cpp-robot-go2-jsonizedatabool-dunder-init-1}{\passthrough{\lstinline!\_\_init\_\_!}}
& \passthrough{\lstinline!def \_\_init\_\_(self) -> None!} &
\passthrough{\lstinline!SIGNATURE\_ONLY!} &
\passthrough{\lstinline!CONSTRUCTION!} \\
\protect\hyperlink{unitree-sdk2-cpp-robot-go2-jsonizedatabool-from-json-1}{\passthrough{\lstinline!from\_json!}}
&
\passthrough{\lstinline!def from\_json(self, json: dict[str, Any]) -> None!}
& \passthrough{\lstinline!SIGNATURE\_ONLY!} &
\passthrough{\lstinline!UNCLASSIFIED!} \\
\protect\hyperlink{unitree-sdk2-cpp-robot-go2-jsonizedatabool-to-json-1}{\passthrough{\lstinline!to\_json!}}
&
\passthrough{\lstinline!def to\_json(self, json: dict[str, Any]) -> None!}
& \passthrough{\lstinline!SIGNATURE\_ONLY!} &
\passthrough{\lstinline!UNCLASSIFIED!} \\
\bottomrule
\end{longtable}

\hypertarget{unitree-sdk2-cpp-robot-go2-jsonizedatabool-dunder-init-1}{%
\paragraph{\texorpdfstring{\texttt{unitree\_sdk2\_cpp.robot.go2.JsonizeDataBool.\_\_init\_\_}}{unitree\_sdk2\_cpp.robot.go2.JsonizeDataBool.\_\_init\_\_}}\label{unitree-sdk2-cpp-robot-go2-jsonizedatabool-dunder-init-1}}

初始化 \passthrough{\lstinline!JsonizeDataBool!}
实例。是否能实际构造取决于下方可用性状态。

\textbf{签名}

\begin{lstlisting}[language=Python]
def __init__(self) -> None
\end{lstlisting}

\textbf{可用性与安全性}

\textbf{\passthrough{\lstinline!SIGNATURE\_ONLY!} /
\passthrough{\lstinline!CONSTRUCTION!}}：当前只有设计期签名，不能假设已存在于运行时扩展。

\textbf{参数}

无显式参数。实例方法中的 \passthrough{\lstinline!self!} 由 Python
自动传入。

\textbf{返回值}

\begin{longtable}[]{@{}
  >{\raggedright\arraybackslash}p{(\columnwidth - 4\tabcolsep) * \real{0.18}}
  >{\raggedright\arraybackslash}p{(\columnwidth - 4\tabcolsep) * \real{0.20}}
  >{\raggedright\arraybackslash}p{(\columnwidth - 4\tabcolsep) * \real{0.62}}@{}}
\toprule
位置 & 类型 & 含义 \\
\midrule
\endhead
无 & \passthrough{\lstinline!None!} &
\passthrough{\lstinline!\_\_init\_\_!}
本身不返回值；调用类对象时，在构造函数可用的前提下得到该类实例。 \\
\bottomrule
\end{longtable}

\textbf{对应 C++}

\begin{itemize}
\tightlist
\item
  类：\passthrough{\lstinline!unitree::robot::go2::JsonizeDataBool!}
\item
  签名：\passthrough{\lstinline!JsonizeDataBool()!}
\item
  绑定策略：\passthrough{\lstinline!未单独标注!}
\item
  声明位置：\passthrough{\lstinline!include/unitree/robot/go2/public/jsonize\_type.hpp:44!}
\end{itemize}

\textbf{说明}

\begin{itemize}
\tightlist
\item
  示例是局部调用片段；\passthrough{\lstinline!obj!}
  和参数变量表示已经按上层业务校验并准备好的对象。
\item
  该条目可以被编辑器和类型检查器识别，但当前运行时可能在属性查找阶段直接失败。
\end{itemize}

\textbf{用法}

\begin{lstlisting}[language=Python]
from typing import TYPE_CHECKING

if TYPE_CHECKING:
    def planned_construct() -> JsonizeDataBool:
        return JsonizeDataBool()
\end{lstlisting}

\begin{quote}
{[}!CAUTION{]} 上面的 \passthrough{\lstinline!TYPE\_CHECKING!}
示例只表达计划调用形态。该分支在普通 Python
运行时为假，不代表当前扩展能执行此接口。
\end{quote}

\hypertarget{unitree-sdk2-cpp-robot-go2-jsonizedatabool-from-json-1}{%
\paragraph{\texorpdfstring{\texttt{unitree\_sdk2\_cpp.robot.go2.JsonizeDataBool.from\_json}}{unitree\_sdk2\_cpp.robot.go2.JsonizeDataBool.from\_json}}\label{unitree-sdk2-cpp-robot-go2-jsonizedatabool-from-json-1}}

计划从 JSON 风格字典读取字段并更新当前 SDK 值对象。

\textbf{签名}

\begin{lstlisting}[language=Python]
def from_json(self, json: dict[str, Any]) -> None
\end{lstlisting}

\textbf{可用性与安全性}

\textbf{\passthrough{\lstinline!SIGNATURE\_ONLY!} /
\passthrough{\lstinline!UNCLASSIFIED!}}：当前只有设计期签名，不能假设已存在于运行时扩展。

\textbf{参数}

\begin{longtable}[]{@{}
  >{\raggedright\arraybackslash}p{(\columnwidth - 6\tabcolsep) * \real{0.18}}
  >{\raggedright\arraybackslash}p{(\columnwidth - 6\tabcolsep) * \real{0.24}}
  >{\raggedright\arraybackslash}p{(\columnwidth - 6\tabcolsep) * \real{0.18}}
  >{\raggedright\arraybackslash}p{(\columnwidth - 6\tabcolsep) * \real{0.40}}@{}}
\toprule
名称 & Python 类型 & 默认值 & 含义 \\
\midrule
\endhead
\passthrough{\lstinline!json!} &
\passthrough{\lstinline!dict[str, Any]!} & 必填 & JSON
风格字典，键和值必须满足对应 C++ \passthrough{\lstinline!JsonMap!}
数据结构的约束。 对应 C++ 参数
\passthrough{\lstinline!json: common::JsonMap \&!}。 \\
\bottomrule
\end{longtable}

\textbf{返回值}

\begin{longtable}[]{@{}
  >{\raggedright\arraybackslash}p{(\columnwidth - 4\tabcolsep) * \real{0.18}}
  >{\raggedright\arraybackslash}p{(\columnwidth - 4\tabcolsep) * \real{0.20}}
  >{\raggedright\arraybackslash}p{(\columnwidth - 4\tabcolsep) * \real{0.62}}@{}}
\toprule
位置 & 类型 & 含义 \\
\midrule
\endhead
无 & \passthrough{\lstinline!None!} & 该方法不返回 Python
值；失败可能表现为异常或底层状态变化。 \\
\bottomrule
\end{longtable}

\textbf{对应 C++}

\begin{itemize}
\tightlist
\item
  类：\passthrough{\lstinline!unitree::robot::go2::JsonizeDataBool!}
\item
  签名：\passthrough{\lstinline!fromJson(common::JsonMap \&)!}
\item
  绑定策略：\passthrough{\lstinline!SIGNATURE\_PREVIEW!}
\item
  声明位置：\passthrough{\lstinline!include/unitree/robot/go2/public/jsonize\_type.hpp:50!}
\end{itemize}

\textbf{说明}

\begin{itemize}
\tightlist
\item
  示例是局部调用片段；\passthrough{\lstinline!obj!}
  和参数变量表示已经按上层业务校验并准备好的对象。
\item
  该条目可以被编辑器和类型检查器识别，但当前运行时可能在属性查找阶段直接失败。
\end{itemize}

\textbf{用法}

\begin{lstlisting}[language=Python]
from typing import TYPE_CHECKING

if TYPE_CHECKING:
    def planned_from_json(obj: JsonizeDataBool, json: dict[str, Any]) -> None:
        obj.from_json(json=json)
\end{lstlisting}

\begin{quote}
{[}!CAUTION{]} 上面的 \passthrough{\lstinline!TYPE\_CHECKING!}
示例只表达计划调用形态。该分支在普通 Python
运行时为假，不代表当前扩展能执行此接口。
\end{quote}

\hypertarget{unitree-sdk2-cpp-robot-go2-jsonizedatabool-to-json-1}{%
\paragraph{\texorpdfstring{\texttt{unitree\_sdk2\_cpp.robot.go2.JsonizeDataBool.to\_json}}{unitree\_sdk2\_cpp.robot.go2.JsonizeDataBool.to\_json}}\label{unitree-sdk2-cpp-robot-go2-jsonizedatabool-to-json-1}}

计划把当前 SDK 值对象写入 JSON 风格字典。

\textbf{签名}

\begin{lstlisting}[language=Python]
def to_json(self, json: dict[str, Any]) -> None
\end{lstlisting}

\textbf{可用性与安全性}

\textbf{\passthrough{\lstinline!SIGNATURE\_ONLY!} /
\passthrough{\lstinline!UNCLASSIFIED!}}：当前只有设计期签名，不能假设已存在于运行时扩展。

\textbf{参数}

\begin{longtable}[]{@{}
  >{\raggedright\arraybackslash}p{(\columnwidth - 6\tabcolsep) * \real{0.18}}
  >{\raggedright\arraybackslash}p{(\columnwidth - 6\tabcolsep) * \real{0.24}}
  >{\raggedright\arraybackslash}p{(\columnwidth - 6\tabcolsep) * \real{0.18}}
  >{\raggedright\arraybackslash}p{(\columnwidth - 6\tabcolsep) * \real{0.40}}@{}}
\toprule
名称 & Python 类型 & 默认值 & 含义 \\
\midrule
\endhead
\passthrough{\lstinline!json!} &
\passthrough{\lstinline!dict[str, Any]!} & 必填 & JSON
风格字典，键和值必须满足对应 C++ \passthrough{\lstinline!JsonMap!}
数据结构的约束。 对应 C++ 参数
\passthrough{\lstinline!json: common::JsonMap \&!}。 \\
\bottomrule
\end{longtable}

\textbf{返回值}

\begin{longtable}[]{@{}
  >{\raggedright\arraybackslash}p{(\columnwidth - 4\tabcolsep) * \real{0.18}}
  >{\raggedright\arraybackslash}p{(\columnwidth - 4\tabcolsep) * \real{0.20}}
  >{\raggedright\arraybackslash}p{(\columnwidth - 4\tabcolsep) * \real{0.62}}@{}}
\toprule
位置 & 类型 & 含义 \\
\midrule
\endhead
无 & \passthrough{\lstinline!None!} & 该方法不返回 Python
值；失败可能表现为异常或底层状态变化。 \\
\bottomrule
\end{longtable}

\textbf{对应 C++}

\begin{itemize}
\tightlist
\item
  类：\passthrough{\lstinline!unitree::robot::go2::JsonizeDataBool!}
\item
  签名：\passthrough{\lstinline!toJson(common::JsonMap \&) const!}
\item
  绑定策略：\passthrough{\lstinline!SIGNATURE\_PREVIEW!}
\item
  声明位置：\passthrough{\lstinline!include/unitree/robot/go2/public/jsonize\_type.hpp:55!}
\end{itemize}

\textbf{说明}

\begin{itemize}
\tightlist
\item
  示例是局部调用片段；\passthrough{\lstinline!obj!}
  和参数变量表示已经按上层业务校验并准备好的对象。
\item
  该条目可以被编辑器和类型检查器识别，但当前运行时可能在属性查找阶段直接失败。
\end{itemize}

\textbf{用法}

\begin{lstlisting}[language=Python]
from typing import TYPE_CHECKING

if TYPE_CHECKING:
    def planned_to_json(obj: JsonizeDataBool, json: dict[str, Any]) -> None:
        obj.to_json(json=json)
\end{lstlisting}

\begin{quote}
{[}!CAUTION{]} 上面的 \passthrough{\lstinline!TYPE\_CHECKING!}
示例只表达计划调用形态。该分支在普通 Python
运行时为假，不代表当前扩展能执行此接口。
\end{quote}

\hypertarget{unitree-sdk2-cpp-robot-go2-jsonizedatadouble}{%
\subsubsection{\texorpdfstring{\texttt{unitree\_sdk2\_cpp.robot.go2.JsonizeDataDouble}}{unitree\_sdk2\_cpp.robot.go2.JsonizeDataDouble}}\label{unitree-sdk2-cpp-robot-go2-jsonizedatadouble}}

SDK 类型签名预览。请逐项查看构造函数和方法的可用性。

\textbf{导入}

\begin{lstlisting}[language=Python]
from unitree_sdk2_cpp.robot.go2 import JsonizeDataDouble
\end{lstlisting}

\textbf{基类}：\passthrough{\lstinline!object!}

\textbf{公开属性}

\begin{longtable}[]{@{}
  >{\raggedright\arraybackslash}p{(\columnwidth - 6\tabcolsep) * \real{0.18}}
  >{\raggedright\arraybackslash}p{(\columnwidth - 6\tabcolsep) * \real{0.24}}
  >{\raggedright\arraybackslash}p{(\columnwidth - 6\tabcolsep) * \real{0.18}}
  >{\raggedright\arraybackslash}p{(\columnwidth - 6\tabcolsep) * \real{0.40}}@{}}
\toprule
名称 & Python 类型 & 含义 & 用法 \\
\midrule
\endhead
\passthrough{\lstinline!data!} & \passthrough{\lstinline!float!} &
传给该接口的 数据 参数，Python 类型为
\passthrough{\lstinline!float!}。精确含义、单位、范围和枚举值需查目标型号对应头文件与协议。
& \passthrough{\lstinline!value = obj.data!} /
\passthrough{\lstinline!obj.data = value!} \\
\bottomrule
\end{longtable}

\textbf{方法索引}

\begin{longtable}[]{@{}
  >{\raggedright\arraybackslash}p{(\columnwidth - 6\tabcolsep) * \real{0.18}}
  >{\raggedright\arraybackslash}p{(\columnwidth - 6\tabcolsep) * \real{0.24}}
  >{\raggedright\arraybackslash}p{(\columnwidth - 6\tabcolsep) * \real{0.18}}
  >{\raggedright\arraybackslash}p{(\columnwidth - 6\tabcolsep) * \real{0.40}}@{}}
\toprule
方法 & Python 签名 & 状态 & 安全分类 \\
\midrule
\endhead
\protect\hyperlink{unitree-sdk2-cpp-robot-go2-jsonizedatadouble-dunder-init-1}{\passthrough{\lstinline!\_\_init\_\_!}}
& \passthrough{\lstinline!def \_\_init\_\_(self) -> None!} &
\passthrough{\lstinline!SIGNATURE\_ONLY!} &
\passthrough{\lstinline!CONSTRUCTION!} \\
\protect\hyperlink{unitree-sdk2-cpp-robot-go2-jsonizedatadouble-from-json-1}{\passthrough{\lstinline!from\_json!}}
&
\passthrough{\lstinline!def from\_json(self, json: dict[str, Any]) -> None!}
& \passthrough{\lstinline!SIGNATURE\_ONLY!} &
\passthrough{\lstinline!UNCLASSIFIED!} \\
\protect\hyperlink{unitree-sdk2-cpp-robot-go2-jsonizedatadouble-to-json-1}{\passthrough{\lstinline!to\_json!}}
&
\passthrough{\lstinline!def to\_json(self, json: dict[str, Any]) -> None!}
& \passthrough{\lstinline!SIGNATURE\_ONLY!} &
\passthrough{\lstinline!UNCLASSIFIED!} \\
\bottomrule
\end{longtable}

\hypertarget{unitree-sdk2-cpp-robot-go2-jsonizedatadouble-dunder-init-1}{%
\paragraph{\texorpdfstring{\texttt{unitree\_sdk2\_cpp.robot.go2.JsonizeDataDouble.\_\_init\_\_}}{unitree\_sdk2\_cpp.robot.go2.JsonizeDataDouble.\_\_init\_\_}}\label{unitree-sdk2-cpp-robot-go2-jsonizedatadouble-dunder-init-1}}

初始化 \passthrough{\lstinline!JsonizeDataDouble!}
实例。是否能实际构造取决于下方可用性状态。

\textbf{签名}

\begin{lstlisting}[language=Python]
def __init__(self) -> None
\end{lstlisting}

\textbf{可用性与安全性}

\textbf{\passthrough{\lstinline!SIGNATURE\_ONLY!} /
\passthrough{\lstinline!CONSTRUCTION!}}：当前只有设计期签名，不能假设已存在于运行时扩展。

\textbf{参数}

无显式参数。实例方法中的 \passthrough{\lstinline!self!} 由 Python
自动传入。

\textbf{返回值}

\begin{longtable}[]{@{}
  >{\raggedright\arraybackslash}p{(\columnwidth - 4\tabcolsep) * \real{0.18}}
  >{\raggedright\arraybackslash}p{(\columnwidth - 4\tabcolsep) * \real{0.20}}
  >{\raggedright\arraybackslash}p{(\columnwidth - 4\tabcolsep) * \real{0.62}}@{}}
\toprule
位置 & 类型 & 含义 \\
\midrule
\endhead
无 & \passthrough{\lstinline!None!} &
\passthrough{\lstinline!\_\_init\_\_!}
本身不返回值；调用类对象时，在构造函数可用的前提下得到该类实例。 \\
\bottomrule
\end{longtable}

\textbf{对应 C++}

\begin{itemize}
\tightlist
\item
  类：\passthrough{\lstinline!unitree::robot::go2::JsonizeDataDouble!}
\item
  签名：\passthrough{\lstinline!JsonizeDataDouble()!}
\item
  绑定策略：\passthrough{\lstinline!未单独标注!}
\item
  声明位置：\passthrough{\lstinline!include/unitree/robot/go2/public/jsonize\_type.hpp:122!}
\end{itemize}

\textbf{说明}

\begin{itemize}
\tightlist
\item
  示例是局部调用片段；\passthrough{\lstinline!obj!}
  和参数变量表示已经按上层业务校验并准备好的对象。
\item
  该条目可以被编辑器和类型检查器识别，但当前运行时可能在属性查找阶段直接失败。
\end{itemize}

\textbf{用法}

\begin{lstlisting}[language=Python]
from typing import TYPE_CHECKING

if TYPE_CHECKING:
    def planned_construct() -> JsonizeDataDouble:
        return JsonizeDataDouble()
\end{lstlisting}

\begin{quote}
{[}!CAUTION{]} 上面的 \passthrough{\lstinline!TYPE\_CHECKING!}
示例只表达计划调用形态。该分支在普通 Python
运行时为假，不代表当前扩展能执行此接口。
\end{quote}

\hypertarget{unitree-sdk2-cpp-robot-go2-jsonizedatadouble-from-json-1}{%
\paragraph{\texorpdfstring{\texttt{unitree\_sdk2\_cpp.robot.go2.JsonizeDataDouble.from\_json}}{unitree\_sdk2\_cpp.robot.go2.JsonizeDataDouble.from\_json}}\label{unitree-sdk2-cpp-robot-go2-jsonizedatadouble-from-json-1}}

计划从 JSON 风格字典读取字段并更新当前 SDK 值对象。

\textbf{签名}

\begin{lstlisting}[language=Python]
def from_json(self, json: dict[str, Any]) -> None
\end{lstlisting}

\textbf{可用性与安全性}

\textbf{\passthrough{\lstinline!SIGNATURE\_ONLY!} /
\passthrough{\lstinline!UNCLASSIFIED!}}：当前只有设计期签名，不能假设已存在于运行时扩展。

\textbf{参数}

\begin{longtable}[]{@{}
  >{\raggedright\arraybackslash}p{(\columnwidth - 6\tabcolsep) * \real{0.18}}
  >{\raggedright\arraybackslash}p{(\columnwidth - 6\tabcolsep) * \real{0.24}}
  >{\raggedright\arraybackslash}p{(\columnwidth - 6\tabcolsep) * \real{0.18}}
  >{\raggedright\arraybackslash}p{(\columnwidth - 6\tabcolsep) * \real{0.40}}@{}}
\toprule
名称 & Python 类型 & 默认值 & 含义 \\
\midrule
\endhead
\passthrough{\lstinline!json!} &
\passthrough{\lstinline!dict[str, Any]!} & 必填 & JSON
风格字典，键和值必须满足对应 C++ \passthrough{\lstinline!JsonMap!}
数据结构的约束。 对应 C++ 参数
\passthrough{\lstinline!json: common::JsonMap \&!}。 \\
\bottomrule
\end{longtable}

\textbf{返回值}

\begin{longtable}[]{@{}
  >{\raggedright\arraybackslash}p{(\columnwidth - 4\tabcolsep) * \real{0.18}}
  >{\raggedright\arraybackslash}p{(\columnwidth - 4\tabcolsep) * \real{0.20}}
  >{\raggedright\arraybackslash}p{(\columnwidth - 4\tabcolsep) * \real{0.62}}@{}}
\toprule
位置 & 类型 & 含义 \\
\midrule
\endhead
无 & \passthrough{\lstinline!None!} & 该方法不返回 Python
值；失败可能表现为异常或底层状态变化。 \\
\bottomrule
\end{longtable}

\textbf{对应 C++}

\begin{itemize}
\tightlist
\item
  类：\passthrough{\lstinline!unitree::robot::go2::JsonizeDataDouble!}
\item
  签名：\passthrough{\lstinline!fromJson(common::JsonMap \&)!}
\item
  绑定策略：\passthrough{\lstinline!SIGNATURE\_PREVIEW!}
\item
  声明位置：\passthrough{\lstinline!include/unitree/robot/go2/public/jsonize\_type.hpp:128!}
\end{itemize}

\textbf{说明}

\begin{itemize}
\tightlist
\item
  示例是局部调用片段；\passthrough{\lstinline!obj!}
  和参数变量表示已经按上层业务校验并准备好的对象。
\item
  该条目可以被编辑器和类型检查器识别，但当前运行时可能在属性查找阶段直接失败。
\end{itemize}

\textbf{用法}

\begin{lstlisting}[language=Python]
from typing import TYPE_CHECKING

if TYPE_CHECKING:
    def planned_from_json(obj: JsonizeDataDouble, json: dict[str, Any]) -> None:
        obj.from_json(json=json)
\end{lstlisting}

\begin{quote}
{[}!CAUTION{]} 上面的 \passthrough{\lstinline!TYPE\_CHECKING!}
示例只表达计划调用形态。该分支在普通 Python
运行时为假，不代表当前扩展能执行此接口。
\end{quote}

\hypertarget{unitree-sdk2-cpp-robot-go2-jsonizedatadouble-to-json-1}{%
\paragraph{\texorpdfstring{\texttt{unitree\_sdk2\_cpp.robot.go2.JsonizeDataDouble.to\_json}}{unitree\_sdk2\_cpp.robot.go2.JsonizeDataDouble.to\_json}}\label{unitree-sdk2-cpp-robot-go2-jsonizedatadouble-to-json-1}}

计划把当前 SDK 值对象写入 JSON 风格字典。

\textbf{签名}

\begin{lstlisting}[language=Python]
def to_json(self, json: dict[str, Any]) -> None
\end{lstlisting}

\textbf{可用性与安全性}

\textbf{\passthrough{\lstinline!SIGNATURE\_ONLY!} /
\passthrough{\lstinline!UNCLASSIFIED!}}：当前只有设计期签名，不能假设已存在于运行时扩展。

\textbf{参数}

\begin{longtable}[]{@{}
  >{\raggedright\arraybackslash}p{(\columnwidth - 6\tabcolsep) * \real{0.18}}
  >{\raggedright\arraybackslash}p{(\columnwidth - 6\tabcolsep) * \real{0.24}}
  >{\raggedright\arraybackslash}p{(\columnwidth - 6\tabcolsep) * \real{0.18}}
  >{\raggedright\arraybackslash}p{(\columnwidth - 6\tabcolsep) * \real{0.40}}@{}}
\toprule
名称 & Python 类型 & 默认值 & 含义 \\
\midrule
\endhead
\passthrough{\lstinline!json!} &
\passthrough{\lstinline!dict[str, Any]!} & 必填 & JSON
风格字典，键和值必须满足对应 C++ \passthrough{\lstinline!JsonMap!}
数据结构的约束。 对应 C++ 参数
\passthrough{\lstinline!json: common::JsonMap \&!}。 \\
\bottomrule
\end{longtable}

\textbf{返回值}

\begin{longtable}[]{@{}
  >{\raggedright\arraybackslash}p{(\columnwidth - 4\tabcolsep) * \real{0.18}}
  >{\raggedright\arraybackslash}p{(\columnwidth - 4\tabcolsep) * \real{0.20}}
  >{\raggedright\arraybackslash}p{(\columnwidth - 4\tabcolsep) * \real{0.62}}@{}}
\toprule
位置 & 类型 & 含义 \\
\midrule
\endhead
无 & \passthrough{\lstinline!None!} & 该方法不返回 Python
值；失败可能表现为异常或底层状态变化。 \\
\bottomrule
\end{longtable}

\textbf{对应 C++}

\begin{itemize}
\tightlist
\item
  类：\passthrough{\lstinline!unitree::robot::go2::JsonizeDataDouble!}
\item
  签名：\passthrough{\lstinline!toJson(common::JsonMap \&) const!}
\item
  绑定策略：\passthrough{\lstinline!SIGNATURE\_PREVIEW!}
\item
  声明位置：\passthrough{\lstinline!include/unitree/robot/go2/public/jsonize\_type.hpp:133!}
\end{itemize}

\textbf{说明}

\begin{itemize}
\tightlist
\item
  示例是局部调用片段；\passthrough{\lstinline!obj!}
  和参数变量表示已经按上层业务校验并准备好的对象。
\item
  该条目可以被编辑器和类型检查器识别，但当前运行时可能在属性查找阶段直接失败。
\end{itemize}

\textbf{用法}

\begin{lstlisting}[language=Python]
from typing import TYPE_CHECKING

if TYPE_CHECKING:
    def planned_to_json(obj: JsonizeDataDouble, json: dict[str, Any]) -> None:
        obj.to_json(json=json)
\end{lstlisting}

\begin{quote}
{[}!CAUTION{]} 上面的 \passthrough{\lstinline!TYPE\_CHECKING!}
示例只表达计划调用形态。该分支在普通 Python
运行时为假，不代表当前扩展能执行此接口。
\end{quote}

\hypertarget{unitree-sdk2-cpp-robot-go2-jsonizedatafloat}{%
\subsubsection{\texorpdfstring{\texttt{unitree\_sdk2\_cpp.robot.go2.JsonizeDataFloat}}{unitree\_sdk2\_cpp.robot.go2.JsonizeDataFloat}}\label{unitree-sdk2-cpp-robot-go2-jsonizedatafloat}}

SDK 类型签名预览。请逐项查看构造函数和方法的可用性。

\textbf{导入}

\begin{lstlisting}[language=Python]
from unitree_sdk2_cpp.robot.go2 import JsonizeDataFloat
\end{lstlisting}

\textbf{基类}：\passthrough{\lstinline!object!}

\textbf{公开属性}

\begin{longtable}[]{@{}
  >{\raggedright\arraybackslash}p{(\columnwidth - 6\tabcolsep) * \real{0.18}}
  >{\raggedright\arraybackslash}p{(\columnwidth - 6\tabcolsep) * \real{0.24}}
  >{\raggedright\arraybackslash}p{(\columnwidth - 6\tabcolsep) * \real{0.18}}
  >{\raggedright\arraybackslash}p{(\columnwidth - 6\tabcolsep) * \real{0.40}}@{}}
\toprule
名称 & Python 类型 & 含义 & 用法 \\
\midrule
\endhead
\passthrough{\lstinline!data!} & \passthrough{\lstinline!float!} &
传给该接口的 数据 参数，Python 类型为
\passthrough{\lstinline!float!}。精确含义、单位、范围和枚举值需查目标型号对应头文件与协议。
& \passthrough{\lstinline!value = obj.data!} /
\passthrough{\lstinline!obj.data = value!} \\
\bottomrule
\end{longtable}

\textbf{方法索引}

\begin{longtable}[]{@{}
  >{\raggedright\arraybackslash}p{(\columnwidth - 6\tabcolsep) * \real{0.18}}
  >{\raggedright\arraybackslash}p{(\columnwidth - 6\tabcolsep) * \real{0.24}}
  >{\raggedright\arraybackslash}p{(\columnwidth - 6\tabcolsep) * \real{0.18}}
  >{\raggedright\arraybackslash}p{(\columnwidth - 6\tabcolsep) * \real{0.40}}@{}}
\toprule
方法 & Python 签名 & 状态 & 安全分类 \\
\midrule
\endhead
\protect\hyperlink{unitree-sdk2-cpp-robot-go2-jsonizedatafloat-dunder-init-1}{\passthrough{\lstinline!\_\_init\_\_!}}
& \passthrough{\lstinline!def \_\_init\_\_(self) -> None!} &
\passthrough{\lstinline!SIGNATURE\_ONLY!} &
\passthrough{\lstinline!CONSTRUCTION!} \\
\protect\hyperlink{unitree-sdk2-cpp-robot-go2-jsonizedatafloat-from-json-1}{\passthrough{\lstinline!from\_json!}}
&
\passthrough{\lstinline!def from\_json(self, json: dict[str, Any]) -> None!}
& \passthrough{\lstinline!SIGNATURE\_ONLY!} &
\passthrough{\lstinline!UNCLASSIFIED!} \\
\protect\hyperlink{unitree-sdk2-cpp-robot-go2-jsonizedatafloat-to-json-1}{\passthrough{\lstinline!to\_json!}}
&
\passthrough{\lstinline!def to\_json(self, json: dict[str, Any]) -> None!}
& \passthrough{\lstinline!SIGNATURE\_ONLY!} &
\passthrough{\lstinline!UNCLASSIFIED!} \\
\bottomrule
\end{longtable}

\hypertarget{unitree-sdk2-cpp-robot-go2-jsonizedatafloat-dunder-init-1}{%
\paragraph{\texorpdfstring{\texttt{unitree\_sdk2\_cpp.robot.go2.JsonizeDataFloat.\_\_init\_\_}}{unitree\_sdk2\_cpp.robot.go2.JsonizeDataFloat.\_\_init\_\_}}\label{unitree-sdk2-cpp-robot-go2-jsonizedatafloat-dunder-init-1}}

初始化 \passthrough{\lstinline!JsonizeDataFloat!}
实例。是否能实际构造取决于下方可用性状态。

\textbf{签名}

\begin{lstlisting}[language=Python]
def __init__(self) -> None
\end{lstlisting}

\textbf{可用性与安全性}

\textbf{\passthrough{\lstinline!SIGNATURE\_ONLY!} /
\passthrough{\lstinline!CONSTRUCTION!}}：当前只有设计期签名，不能假设已存在于运行时扩展。

\textbf{参数}

无显式参数。实例方法中的 \passthrough{\lstinline!self!} 由 Python
自动传入。

\textbf{返回值}

\begin{longtable}[]{@{}
  >{\raggedright\arraybackslash}p{(\columnwidth - 4\tabcolsep) * \real{0.18}}
  >{\raggedright\arraybackslash}p{(\columnwidth - 4\tabcolsep) * \real{0.20}}
  >{\raggedright\arraybackslash}p{(\columnwidth - 4\tabcolsep) * \real{0.62}}@{}}
\toprule
位置 & 类型 & 含义 \\
\midrule
\endhead
无 & \passthrough{\lstinline!None!} &
\passthrough{\lstinline!\_\_init\_\_!}
本身不返回值；调用类对象时，在构造函数可用的前提下得到该类实例。 \\
\bottomrule
\end{longtable}

\textbf{对应 C++}

\begin{itemize}
\tightlist
\item
  类：\passthrough{\lstinline!unitree::robot::go2::JsonizeDataFloat!}
\item
  签名：\passthrough{\lstinline!JsonizeDataFloat()!}
\item
  绑定策略：\passthrough{\lstinline!未单独标注!}
\item
  声明位置：\passthrough{\lstinline!include/unitree/robot/go2/public/jsonize\_type.hpp:96!}
\end{itemize}

\textbf{说明}

\begin{itemize}
\tightlist
\item
  示例是局部调用片段；\passthrough{\lstinline!obj!}
  和参数变量表示已经按上层业务校验并准备好的对象。
\item
  该条目可以被编辑器和类型检查器识别，但当前运行时可能在属性查找阶段直接失败。
\end{itemize}

\textbf{用法}

\begin{lstlisting}[language=Python]
from typing import TYPE_CHECKING

if TYPE_CHECKING:
    def planned_construct() -> JsonizeDataFloat:
        return JsonizeDataFloat()
\end{lstlisting}

\begin{quote}
{[}!CAUTION{]} 上面的 \passthrough{\lstinline!TYPE\_CHECKING!}
示例只表达计划调用形态。该分支在普通 Python
运行时为假，不代表当前扩展能执行此接口。
\end{quote}

\hypertarget{unitree-sdk2-cpp-robot-go2-jsonizedatafloat-from-json-1}{%
\paragraph{\texorpdfstring{\texttt{unitree\_sdk2\_cpp.robot.go2.JsonizeDataFloat.from\_json}}{unitree\_sdk2\_cpp.robot.go2.JsonizeDataFloat.from\_json}}\label{unitree-sdk2-cpp-robot-go2-jsonizedatafloat-from-json-1}}

计划从 JSON 风格字典读取字段并更新当前 SDK 值对象。

\textbf{签名}

\begin{lstlisting}[language=Python]
def from_json(self, json: dict[str, Any]) -> None
\end{lstlisting}

\textbf{可用性与安全性}

\textbf{\passthrough{\lstinline!SIGNATURE\_ONLY!} /
\passthrough{\lstinline!UNCLASSIFIED!}}：当前只有设计期签名，不能假设已存在于运行时扩展。

\textbf{参数}

\begin{longtable}[]{@{}
  >{\raggedright\arraybackslash}p{(\columnwidth - 6\tabcolsep) * \real{0.18}}
  >{\raggedright\arraybackslash}p{(\columnwidth - 6\tabcolsep) * \real{0.24}}
  >{\raggedright\arraybackslash}p{(\columnwidth - 6\tabcolsep) * \real{0.18}}
  >{\raggedright\arraybackslash}p{(\columnwidth - 6\tabcolsep) * \real{0.40}}@{}}
\toprule
名称 & Python 类型 & 默认值 & 含义 \\
\midrule
\endhead
\passthrough{\lstinline!json!} &
\passthrough{\lstinline!dict[str, Any]!} & 必填 & JSON
风格字典，键和值必须满足对应 C++ \passthrough{\lstinline!JsonMap!}
数据结构的约束。 对应 C++ 参数
\passthrough{\lstinline!json: common::JsonMap \&!}。 \\
\bottomrule
\end{longtable}

\textbf{返回值}

\begin{longtable}[]{@{}
  >{\raggedright\arraybackslash}p{(\columnwidth - 4\tabcolsep) * \real{0.18}}
  >{\raggedright\arraybackslash}p{(\columnwidth - 4\tabcolsep) * \real{0.20}}
  >{\raggedright\arraybackslash}p{(\columnwidth - 4\tabcolsep) * \real{0.62}}@{}}
\toprule
位置 & 类型 & 含义 \\
\midrule
\endhead
无 & \passthrough{\lstinline!None!} & 该方法不返回 Python
值；失败可能表现为异常或底层状态变化。 \\
\bottomrule
\end{longtable}

\textbf{对应 C++}

\begin{itemize}
\tightlist
\item
  类：\passthrough{\lstinline!unitree::robot::go2::JsonizeDataFloat!}
\item
  签名：\passthrough{\lstinline!fromJson(common::JsonMap \&)!}
\item
  绑定策略：\passthrough{\lstinline!SIGNATURE\_PREVIEW!}
\item
  声明位置：\passthrough{\lstinline!include/unitree/robot/go2/public/jsonize\_type.hpp:102!}
\end{itemize}

\textbf{说明}

\begin{itemize}
\tightlist
\item
  示例是局部调用片段；\passthrough{\lstinline!obj!}
  和参数变量表示已经按上层业务校验并准备好的对象。
\item
  该条目可以被编辑器和类型检查器识别，但当前运行时可能在属性查找阶段直接失败。
\end{itemize}

\textbf{用法}

\begin{lstlisting}[language=Python]
from typing import TYPE_CHECKING

if TYPE_CHECKING:
    def planned_from_json(obj: JsonizeDataFloat, json: dict[str, Any]) -> None:
        obj.from_json(json=json)
\end{lstlisting}

\begin{quote}
{[}!CAUTION{]} 上面的 \passthrough{\lstinline!TYPE\_CHECKING!}
示例只表达计划调用形态。该分支在普通 Python
运行时为假，不代表当前扩展能执行此接口。
\end{quote}

\hypertarget{unitree-sdk2-cpp-robot-go2-jsonizedatafloat-to-json-1}{%
\paragraph{\texorpdfstring{\texttt{unitree\_sdk2\_cpp.robot.go2.JsonizeDataFloat.to\_json}}{unitree\_sdk2\_cpp.robot.go2.JsonizeDataFloat.to\_json}}\label{unitree-sdk2-cpp-robot-go2-jsonizedatafloat-to-json-1}}

计划把当前 SDK 值对象写入 JSON 风格字典。

\textbf{签名}

\begin{lstlisting}[language=Python]
def to_json(self, json: dict[str, Any]) -> None
\end{lstlisting}

\textbf{可用性与安全性}

\textbf{\passthrough{\lstinline!SIGNATURE\_ONLY!} /
\passthrough{\lstinline!UNCLASSIFIED!}}：当前只有设计期签名，不能假设已存在于运行时扩展。

\textbf{参数}

\begin{longtable}[]{@{}
  >{\raggedright\arraybackslash}p{(\columnwidth - 6\tabcolsep) * \real{0.18}}
  >{\raggedright\arraybackslash}p{(\columnwidth - 6\tabcolsep) * \real{0.24}}
  >{\raggedright\arraybackslash}p{(\columnwidth - 6\tabcolsep) * \real{0.18}}
  >{\raggedright\arraybackslash}p{(\columnwidth - 6\tabcolsep) * \real{0.40}}@{}}
\toprule
名称 & Python 类型 & 默认值 & 含义 \\
\midrule
\endhead
\passthrough{\lstinline!json!} &
\passthrough{\lstinline!dict[str, Any]!} & 必填 & JSON
风格字典，键和值必须满足对应 C++ \passthrough{\lstinline!JsonMap!}
数据结构的约束。 对应 C++ 参数
\passthrough{\lstinline!json: common::JsonMap \&!}。 \\
\bottomrule
\end{longtable}

\textbf{返回值}

\begin{longtable}[]{@{}
  >{\raggedright\arraybackslash}p{(\columnwidth - 4\tabcolsep) * \real{0.18}}
  >{\raggedright\arraybackslash}p{(\columnwidth - 4\tabcolsep) * \real{0.20}}
  >{\raggedright\arraybackslash}p{(\columnwidth - 4\tabcolsep) * \real{0.62}}@{}}
\toprule
位置 & 类型 & 含义 \\
\midrule
\endhead
无 & \passthrough{\lstinline!None!} & 该方法不返回 Python
值；失败可能表现为异常或底层状态变化。 \\
\bottomrule
\end{longtable}

\textbf{对应 C++}

\begin{itemize}
\tightlist
\item
  类：\passthrough{\lstinline!unitree::robot::go2::JsonizeDataFloat!}
\item
  签名：\passthrough{\lstinline!toJson(common::JsonMap \&) const!}
\item
  绑定策略：\passthrough{\lstinline!SIGNATURE\_PREVIEW!}
\item
  声明位置：\passthrough{\lstinline!include/unitree/robot/go2/public/jsonize\_type.hpp:107!}
\end{itemize}

\textbf{说明}

\begin{itemize}
\tightlist
\item
  示例是局部调用片段；\passthrough{\lstinline!obj!}
  和参数变量表示已经按上层业务校验并准备好的对象。
\item
  该条目可以被编辑器和类型检查器识别，但当前运行时可能在属性查找阶段直接失败。
\end{itemize}

\textbf{用法}

\begin{lstlisting}[language=Python]
from typing import TYPE_CHECKING

if TYPE_CHECKING:
    def planned_to_json(obj: JsonizeDataFloat, json: dict[str, Any]) -> None:
        obj.to_json(json=json)
\end{lstlisting}

\begin{quote}
{[}!CAUTION{]} 上面的 \passthrough{\lstinline!TYPE\_CHECKING!}
示例只表达计划调用形态。该分支在普通 Python
运行时为假，不代表当前扩展能执行此接口。
\end{quote}

\hypertarget{unitree-sdk2-cpp-robot-go2-jsonizedataint}{%
\subsubsection{\texorpdfstring{\texttt{unitree\_sdk2\_cpp.robot.go2.JsonizeDataInt}}{unitree\_sdk2\_cpp.robot.go2.JsonizeDataInt}}\label{unitree-sdk2-cpp-robot-go2-jsonizedataint}}

SDK 类型签名预览。请逐项查看构造函数和方法的可用性。

\textbf{导入}

\begin{lstlisting}[language=Python]
from unitree_sdk2_cpp.robot.go2 import JsonizeDataInt
\end{lstlisting}

\textbf{基类}：\passthrough{\lstinline!object!}

\textbf{公开属性}

\begin{longtable}[]{@{}
  >{\raggedright\arraybackslash}p{(\columnwidth - 6\tabcolsep) * \real{0.18}}
  >{\raggedright\arraybackslash}p{(\columnwidth - 6\tabcolsep) * \real{0.24}}
  >{\raggedright\arraybackslash}p{(\columnwidth - 6\tabcolsep) * \real{0.18}}
  >{\raggedright\arraybackslash}p{(\columnwidth - 6\tabcolsep) * \real{0.40}}@{}}
\toprule
名称 & Python 类型 & 含义 & 用法 \\
\midrule
\endhead
\passthrough{\lstinline!data!} & \passthrough{\lstinline!int!} &
传给该接口的 数据 参数，Python 类型为
\passthrough{\lstinline!int!}。精确含义、单位、范围和枚举值需查目标型号对应头文件与协议。
& \passthrough{\lstinline!value = obj.data!} /
\passthrough{\lstinline!obj.data = value!} \\
\bottomrule
\end{longtable}

\textbf{方法索引}

\begin{longtable}[]{@{}
  >{\raggedright\arraybackslash}p{(\columnwidth - 6\tabcolsep) * \real{0.18}}
  >{\raggedright\arraybackslash}p{(\columnwidth - 6\tabcolsep) * \real{0.24}}
  >{\raggedright\arraybackslash}p{(\columnwidth - 6\tabcolsep) * \real{0.18}}
  >{\raggedright\arraybackslash}p{(\columnwidth - 6\tabcolsep) * \real{0.40}}@{}}
\toprule
方法 & Python 签名 & 状态 & 安全分类 \\
\midrule
\endhead
\protect\hyperlink{unitree-sdk2-cpp-robot-go2-jsonizedataint-dunder-init-1}{\passthrough{\lstinline!\_\_init\_\_!}}
& \passthrough{\lstinline!def \_\_init\_\_(self) -> None!} &
\passthrough{\lstinline!SIGNATURE\_ONLY!} &
\passthrough{\lstinline!CONSTRUCTION!} \\
\protect\hyperlink{unitree-sdk2-cpp-robot-go2-jsonizedataint-from-json-1}{\passthrough{\lstinline!from\_json!}}
&
\passthrough{\lstinline!def from\_json(self, json: dict[str, Any]) -> None!}
& \passthrough{\lstinline!SIGNATURE\_ONLY!} &
\passthrough{\lstinline!UNCLASSIFIED!} \\
\protect\hyperlink{unitree-sdk2-cpp-robot-go2-jsonizedataint-to-json-1}{\passthrough{\lstinline!to\_json!}}
&
\passthrough{\lstinline!def to\_json(self, json: dict[str, Any]) -> None!}
& \passthrough{\lstinline!SIGNATURE\_ONLY!} &
\passthrough{\lstinline!UNCLASSIFIED!} \\
\bottomrule
\end{longtable}

\hypertarget{unitree-sdk2-cpp-robot-go2-jsonizedataint-dunder-init-1}{%
\paragraph{\texorpdfstring{\texttt{unitree\_sdk2\_cpp.robot.go2.JsonizeDataInt.\_\_init\_\_}}{unitree\_sdk2\_cpp.robot.go2.JsonizeDataInt.\_\_init\_\_}}\label{unitree-sdk2-cpp-robot-go2-jsonizedataint-dunder-init-1}}

初始化 \passthrough{\lstinline!JsonizeDataInt!}
实例。是否能实际构造取决于下方可用性状态。

\textbf{签名}

\begin{lstlisting}[language=Python]
def __init__(self) -> None
\end{lstlisting}

\textbf{可用性与安全性}

\textbf{\passthrough{\lstinline!SIGNATURE\_ONLY!} /
\passthrough{\lstinline!CONSTRUCTION!}}：当前只有设计期签名，不能假设已存在于运行时扩展。

\textbf{参数}

无显式参数。实例方法中的 \passthrough{\lstinline!self!} 由 Python
自动传入。

\textbf{返回值}

\begin{longtable}[]{@{}
  >{\raggedright\arraybackslash}p{(\columnwidth - 4\tabcolsep) * \real{0.18}}
  >{\raggedright\arraybackslash}p{(\columnwidth - 4\tabcolsep) * \real{0.20}}
  >{\raggedright\arraybackslash}p{(\columnwidth - 4\tabcolsep) * \real{0.62}}@{}}
\toprule
位置 & 类型 & 含义 \\
\midrule
\endhead
无 & \passthrough{\lstinline!None!} &
\passthrough{\lstinline!\_\_init\_\_!}
本身不返回值；调用类对象时，在构造函数可用的前提下得到该类实例。 \\
\bottomrule
\end{longtable}

\textbf{对应 C++}

\begin{itemize}
\tightlist
\item
  类：\passthrough{\lstinline!unitree::robot::go2::JsonizeDataInt!}
\item
  签名：\passthrough{\lstinline!JsonizeDataInt()!}
\item
  绑定策略：\passthrough{\lstinline!未单独标注!}
\item
  声明位置：\passthrough{\lstinline!include/unitree/robot/go2/public/jsonize\_type.hpp:70!}
\end{itemize}

\textbf{说明}

\begin{itemize}
\tightlist
\item
  示例是局部调用片段；\passthrough{\lstinline!obj!}
  和参数变量表示已经按上层业务校验并准备好的对象。
\item
  该条目可以被编辑器和类型检查器识别，但当前运行时可能在属性查找阶段直接失败。
\end{itemize}

\textbf{用法}

\begin{lstlisting}[language=Python]
from typing import TYPE_CHECKING

if TYPE_CHECKING:
    def planned_construct() -> JsonizeDataInt:
        return JsonizeDataInt()
\end{lstlisting}

\begin{quote}
{[}!CAUTION{]} 上面的 \passthrough{\lstinline!TYPE\_CHECKING!}
示例只表达计划调用形态。该分支在普通 Python
运行时为假，不代表当前扩展能执行此接口。
\end{quote}

\hypertarget{unitree-sdk2-cpp-robot-go2-jsonizedataint-from-json-1}{%
\paragraph{\texorpdfstring{\texttt{unitree\_sdk2\_cpp.robot.go2.JsonizeDataInt.from\_json}}{unitree\_sdk2\_cpp.robot.go2.JsonizeDataInt.from\_json}}\label{unitree-sdk2-cpp-robot-go2-jsonizedataint-from-json-1}}

计划从 JSON 风格字典读取字段并更新当前 SDK 值对象。

\textbf{签名}

\begin{lstlisting}[language=Python]
def from_json(self, json: dict[str, Any]) -> None
\end{lstlisting}

\textbf{可用性与安全性}

\textbf{\passthrough{\lstinline!SIGNATURE\_ONLY!} /
\passthrough{\lstinline!UNCLASSIFIED!}}：当前只有设计期签名，不能假设已存在于运行时扩展。

\textbf{参数}

\begin{longtable}[]{@{}
  >{\raggedright\arraybackslash}p{(\columnwidth - 6\tabcolsep) * \real{0.18}}
  >{\raggedright\arraybackslash}p{(\columnwidth - 6\tabcolsep) * \real{0.24}}
  >{\raggedright\arraybackslash}p{(\columnwidth - 6\tabcolsep) * \real{0.18}}
  >{\raggedright\arraybackslash}p{(\columnwidth - 6\tabcolsep) * \real{0.40}}@{}}
\toprule
名称 & Python 类型 & 默认值 & 含义 \\
\midrule
\endhead
\passthrough{\lstinline!json!} &
\passthrough{\lstinline!dict[str, Any]!} & 必填 & JSON
风格字典，键和值必须满足对应 C++ \passthrough{\lstinline!JsonMap!}
数据结构的约束。 对应 C++ 参数
\passthrough{\lstinline!json: common::JsonMap \&!}。 \\
\bottomrule
\end{longtable}

\textbf{返回值}

\begin{longtable}[]{@{}
  >{\raggedright\arraybackslash}p{(\columnwidth - 4\tabcolsep) * \real{0.18}}
  >{\raggedright\arraybackslash}p{(\columnwidth - 4\tabcolsep) * \real{0.20}}
  >{\raggedright\arraybackslash}p{(\columnwidth - 4\tabcolsep) * \real{0.62}}@{}}
\toprule
位置 & 类型 & 含义 \\
\midrule
\endhead
无 & \passthrough{\lstinline!None!} & 该方法不返回 Python
值；失败可能表现为异常或底层状态变化。 \\
\bottomrule
\end{longtable}

\textbf{对应 C++}

\begin{itemize}
\tightlist
\item
  类：\passthrough{\lstinline!unitree::robot::go2::JsonizeDataInt!}
\item
  签名：\passthrough{\lstinline!fromJson(common::JsonMap \&)!}
\item
  绑定策略：\passthrough{\lstinline!SIGNATURE\_PREVIEW!}
\item
  声明位置：\passthrough{\lstinline!include/unitree/robot/go2/public/jsonize\_type.hpp:76!}
\end{itemize}

\textbf{说明}

\begin{itemize}
\tightlist
\item
  示例是局部调用片段；\passthrough{\lstinline!obj!}
  和参数变量表示已经按上层业务校验并准备好的对象。
\item
  该条目可以被编辑器和类型检查器识别，但当前运行时可能在属性查找阶段直接失败。
\end{itemize}

\textbf{用法}

\begin{lstlisting}[language=Python]
from typing import TYPE_CHECKING

if TYPE_CHECKING:
    def planned_from_json(obj: JsonizeDataInt, json: dict[str, Any]) -> None:
        obj.from_json(json=json)
\end{lstlisting}

\begin{quote}
{[}!CAUTION{]} 上面的 \passthrough{\lstinline!TYPE\_CHECKING!}
示例只表达计划调用形态。该分支在普通 Python
运行时为假，不代表当前扩展能执行此接口。
\end{quote}

\hypertarget{unitree-sdk2-cpp-robot-go2-jsonizedataint-to-json-1}{%
\paragraph{\texorpdfstring{\texttt{unitree\_sdk2\_cpp.robot.go2.JsonizeDataInt.to\_json}}{unitree\_sdk2\_cpp.robot.go2.JsonizeDataInt.to\_json}}\label{unitree-sdk2-cpp-robot-go2-jsonizedataint-to-json-1}}

计划把当前 SDK 值对象写入 JSON 风格字典。

\textbf{签名}

\begin{lstlisting}[language=Python]
def to_json(self, json: dict[str, Any]) -> None
\end{lstlisting}

\textbf{可用性与安全性}

\textbf{\passthrough{\lstinline!SIGNATURE\_ONLY!} /
\passthrough{\lstinline!UNCLASSIFIED!}}：当前只有设计期签名，不能假设已存在于运行时扩展。

\textbf{参数}

\begin{longtable}[]{@{}
  >{\raggedright\arraybackslash}p{(\columnwidth - 6\tabcolsep) * \real{0.18}}
  >{\raggedright\arraybackslash}p{(\columnwidth - 6\tabcolsep) * \real{0.24}}
  >{\raggedright\arraybackslash}p{(\columnwidth - 6\tabcolsep) * \real{0.18}}
  >{\raggedright\arraybackslash}p{(\columnwidth - 6\tabcolsep) * \real{0.40}}@{}}
\toprule
名称 & Python 类型 & 默认值 & 含义 \\
\midrule
\endhead
\passthrough{\lstinline!json!} &
\passthrough{\lstinline!dict[str, Any]!} & 必填 & JSON
风格字典，键和值必须满足对应 C++ \passthrough{\lstinline!JsonMap!}
数据结构的约束。 对应 C++ 参数
\passthrough{\lstinline!json: common::JsonMap \&!}。 \\
\bottomrule
\end{longtable}

\textbf{返回值}

\begin{longtable}[]{@{}
  >{\raggedright\arraybackslash}p{(\columnwidth - 4\tabcolsep) * \real{0.18}}
  >{\raggedright\arraybackslash}p{(\columnwidth - 4\tabcolsep) * \real{0.20}}
  >{\raggedright\arraybackslash}p{(\columnwidth - 4\tabcolsep) * \real{0.62}}@{}}
\toprule
位置 & 类型 & 含义 \\
\midrule
\endhead
无 & \passthrough{\lstinline!None!} & 该方法不返回 Python
值；失败可能表现为异常或底层状态变化。 \\
\bottomrule
\end{longtable}

\textbf{对应 C++}

\begin{itemize}
\tightlist
\item
  类：\passthrough{\lstinline!unitree::robot::go2::JsonizeDataInt!}
\item
  签名：\passthrough{\lstinline!toJson(common::JsonMap \&) const!}
\item
  绑定策略：\passthrough{\lstinline!SIGNATURE\_PREVIEW!}
\item
  声明位置：\passthrough{\lstinline!include/unitree/robot/go2/public/jsonize\_type.hpp:81!}
\end{itemize}

\textbf{说明}

\begin{itemize}
\tightlist
\item
  示例是局部调用片段；\passthrough{\lstinline!obj!}
  和参数变量表示已经按上层业务校验并准备好的对象。
\item
  该条目可以被编辑器和类型检查器识别，但当前运行时可能在属性查找阶段直接失败。
\end{itemize}

\textbf{用法}

\begin{lstlisting}[language=Python]
from typing import TYPE_CHECKING

if TYPE_CHECKING:
    def planned_to_json(obj: JsonizeDataInt, json: dict[str, Any]) -> None:
        obj.to_json(json=json)
\end{lstlisting}

\begin{quote}
{[}!CAUTION{]} 上面的 \passthrough{\lstinline!TYPE\_CHECKING!}
示例只表达计划调用形态。该分支在普通 Python
运行时为假，不代表当前扩展能执行此接口。
\end{quote}

\hypertarget{unitree-sdk2-cpp-robot-go2-jsonizedatastring}{%
\subsubsection{\texorpdfstring{\texttt{unitree\_sdk2\_cpp.robot.go2.JsonizeDataString}}{unitree\_sdk2\_cpp.robot.go2.JsonizeDataString}}\label{unitree-sdk2-cpp-robot-go2-jsonizedatastring}}

SDK 类型签名预览。请逐项查看构造函数和方法的可用性。

\textbf{导入}

\begin{lstlisting}[language=Python]
from unitree_sdk2_cpp.robot.go2 import JsonizeDataString
\end{lstlisting}

\textbf{基类}：\passthrough{\lstinline!object!}

\textbf{公开属性}

\begin{longtable}[]{@{}
  >{\raggedright\arraybackslash}p{(\columnwidth - 6\tabcolsep) * \real{0.18}}
  >{\raggedright\arraybackslash}p{(\columnwidth - 6\tabcolsep) * \real{0.24}}
  >{\raggedright\arraybackslash}p{(\columnwidth - 6\tabcolsep) * \real{0.18}}
  >{\raggedright\arraybackslash}p{(\columnwidth - 6\tabcolsep) * \real{0.40}}@{}}
\toprule
名称 & Python 类型 & 含义 & 用法 \\
\midrule
\endhead
\passthrough{\lstinline!data!} & \passthrough{\lstinline!str!} &
传给该接口的 数据 参数，Python 类型为
\passthrough{\lstinline!str!}。精确含义、单位、范围和枚举值需查目标型号对应头文件与协议。
& \passthrough{\lstinline!value = obj.data!} /
\passthrough{\lstinline!obj.data = value!} \\
\bottomrule
\end{longtable}

\textbf{方法索引}

\begin{longtable}[]{@{}
  >{\raggedright\arraybackslash}p{(\columnwidth - 6\tabcolsep) * \real{0.18}}
  >{\raggedright\arraybackslash}p{(\columnwidth - 6\tabcolsep) * \real{0.24}}
  >{\raggedright\arraybackslash}p{(\columnwidth - 6\tabcolsep) * \real{0.18}}
  >{\raggedright\arraybackslash}p{(\columnwidth - 6\tabcolsep) * \real{0.40}}@{}}
\toprule
方法 & Python 签名 & 状态 & 安全分类 \\
\midrule
\endhead
\protect\hyperlink{unitree-sdk2-cpp-robot-go2-jsonizedatastring-dunder-init-1}{\passthrough{\lstinline!\_\_init\_\_!}}
& \passthrough{\lstinline!def \_\_init\_\_(self) -> None!} &
\passthrough{\lstinline!SIGNATURE\_ONLY!} &
\passthrough{\lstinline!CONSTRUCTION!} \\
\protect\hyperlink{unitree-sdk2-cpp-robot-go2-jsonizedatastring-from-json-1}{\passthrough{\lstinline!from\_json!}}
&
\passthrough{\lstinline!def from\_json(self, json: dict[str, Any]) -> None!}
& \passthrough{\lstinline!SIGNATURE\_ONLY!} &
\passthrough{\lstinline!UNCLASSIFIED!} \\
\protect\hyperlink{unitree-sdk2-cpp-robot-go2-jsonizedatastring-to-json-1}{\passthrough{\lstinline!to\_json!}}
&
\passthrough{\lstinline!def to\_json(self, json: dict[str, Any]) -> None!}
& \passthrough{\lstinline!SIGNATURE\_ONLY!} &
\passthrough{\lstinline!UNCLASSIFIED!} \\
\bottomrule
\end{longtable}

\hypertarget{unitree-sdk2-cpp-robot-go2-jsonizedatastring-dunder-init-1}{%
\paragraph{\texorpdfstring{\texttt{unitree\_sdk2\_cpp.robot.go2.JsonizeDataString.\_\_init\_\_}}{unitree\_sdk2\_cpp.robot.go2.JsonizeDataString.\_\_init\_\_}}\label{unitree-sdk2-cpp-robot-go2-jsonizedatastring-dunder-init-1}}

初始化 \passthrough{\lstinline!JsonizeDataString!}
实例。是否能实际构造取决于下方可用性状态。

\textbf{签名}

\begin{lstlisting}[language=Python]
def __init__(self) -> None
\end{lstlisting}

\textbf{可用性与安全性}

\textbf{\passthrough{\lstinline!SIGNATURE\_ONLY!} /
\passthrough{\lstinline!CONSTRUCTION!}}：当前只有设计期签名，不能假设已存在于运行时扩展。

\textbf{参数}

无显式参数。实例方法中的 \passthrough{\lstinline!self!} 由 Python
自动传入。

\textbf{返回值}

\begin{longtable}[]{@{}
  >{\raggedright\arraybackslash}p{(\columnwidth - 4\tabcolsep) * \real{0.18}}
  >{\raggedright\arraybackslash}p{(\columnwidth - 4\tabcolsep) * \real{0.20}}
  >{\raggedright\arraybackslash}p{(\columnwidth - 4\tabcolsep) * \real{0.62}}@{}}
\toprule
位置 & 类型 & 含义 \\
\midrule
\endhead
无 & \passthrough{\lstinline!None!} &
\passthrough{\lstinline!\_\_init\_\_!}
本身不返回值；调用类对象时，在构造函数可用的前提下得到该类实例。 \\
\bottomrule
\end{longtable}

\textbf{对应 C++}

\begin{itemize}
\tightlist
\item
  类：\passthrough{\lstinline!unitree::robot::go2::JsonizeDataString!}
\item
  签名：\passthrough{\lstinline!JsonizeDataString()!}
\item
  绑定策略：\passthrough{\lstinline!未单独标注!}
\item
  声明位置：\passthrough{\lstinline!include/unitree/robot/go2/public/jsonize\_type.hpp:148!}
\end{itemize}

\textbf{说明}

\begin{itemize}
\tightlist
\item
  示例是局部调用片段；\passthrough{\lstinline!obj!}
  和参数变量表示已经按上层业务校验并准备好的对象。
\item
  该条目可以被编辑器和类型检查器识别，但当前运行时可能在属性查找阶段直接失败。
\end{itemize}

\textbf{用法}

\begin{lstlisting}[language=Python]
from typing import TYPE_CHECKING

if TYPE_CHECKING:
    def planned_construct() -> JsonizeDataString:
        return JsonizeDataString()
\end{lstlisting}

\begin{quote}
{[}!CAUTION{]} 上面的 \passthrough{\lstinline!TYPE\_CHECKING!}
示例只表达计划调用形态。该分支在普通 Python
运行时为假，不代表当前扩展能执行此接口。
\end{quote}

\hypertarget{unitree-sdk2-cpp-robot-go2-jsonizedatastring-from-json-1}{%
\paragraph{\texorpdfstring{\texttt{unitree\_sdk2\_cpp.robot.go2.JsonizeDataString.from\_json}}{unitree\_sdk2\_cpp.robot.go2.JsonizeDataString.from\_json}}\label{unitree-sdk2-cpp-robot-go2-jsonizedatastring-from-json-1}}

计划从 JSON 风格字典读取字段并更新当前 SDK 值对象。

\textbf{签名}

\begin{lstlisting}[language=Python]
def from_json(self, json: dict[str, Any]) -> None
\end{lstlisting}

\textbf{可用性与安全性}

\textbf{\passthrough{\lstinline!SIGNATURE\_ONLY!} /
\passthrough{\lstinline!UNCLASSIFIED!}}：当前只有设计期签名，不能假设已存在于运行时扩展。

\textbf{参数}

\begin{longtable}[]{@{}
  >{\raggedright\arraybackslash}p{(\columnwidth - 6\tabcolsep) * \real{0.18}}
  >{\raggedright\arraybackslash}p{(\columnwidth - 6\tabcolsep) * \real{0.24}}
  >{\raggedright\arraybackslash}p{(\columnwidth - 6\tabcolsep) * \real{0.18}}
  >{\raggedright\arraybackslash}p{(\columnwidth - 6\tabcolsep) * \real{0.40}}@{}}
\toprule
名称 & Python 类型 & 默认值 & 含义 \\
\midrule
\endhead
\passthrough{\lstinline!json!} &
\passthrough{\lstinline!dict[str, Any]!} & 必填 & JSON
风格字典，键和值必须满足对应 C++ \passthrough{\lstinline!JsonMap!}
数据结构的约束。 对应 C++ 参数
\passthrough{\lstinline!json: common::JsonMap \&!}。 \\
\bottomrule
\end{longtable}

\textbf{返回值}

\begin{longtable}[]{@{}
  >{\raggedright\arraybackslash}p{(\columnwidth - 4\tabcolsep) * \real{0.18}}
  >{\raggedright\arraybackslash}p{(\columnwidth - 4\tabcolsep) * \real{0.20}}
  >{\raggedright\arraybackslash}p{(\columnwidth - 4\tabcolsep) * \real{0.62}}@{}}
\toprule
位置 & 类型 & 含义 \\
\midrule
\endhead
无 & \passthrough{\lstinline!None!} & 该方法不返回 Python
值；失败可能表现为异常或底层状态变化。 \\
\bottomrule
\end{longtable}

\textbf{对应 C++}

\begin{itemize}
\tightlist
\item
  类：\passthrough{\lstinline!unitree::robot::go2::JsonizeDataString!}
\item
  签名：\passthrough{\lstinline!fromJson(common::JsonMap \&)!}
\item
  绑定策略：\passthrough{\lstinline!SIGNATURE\_PREVIEW!}
\item
  声明位置：\passthrough{\lstinline!include/unitree/robot/go2/public/jsonize\_type.hpp:154!}
\end{itemize}

\textbf{说明}

\begin{itemize}
\tightlist
\item
  示例是局部调用片段；\passthrough{\lstinline!obj!}
  和参数变量表示已经按上层业务校验并准备好的对象。
\item
  该条目可以被编辑器和类型检查器识别，但当前运行时可能在属性查找阶段直接失败。
\end{itemize}

\textbf{用法}

\begin{lstlisting}[language=Python]
from typing import TYPE_CHECKING

if TYPE_CHECKING:
    def planned_from_json(obj: JsonizeDataString, json: dict[str, Any]) -> None:
        obj.from_json(json=json)
\end{lstlisting}

\begin{quote}
{[}!CAUTION{]} 上面的 \passthrough{\lstinline!TYPE\_CHECKING!}
示例只表达计划调用形态。该分支在普通 Python
运行时为假，不代表当前扩展能执行此接口。
\end{quote}

\hypertarget{unitree-sdk2-cpp-robot-go2-jsonizedatastring-to-json-1}{%
\paragraph{\texorpdfstring{\texttt{unitree\_sdk2\_cpp.robot.go2.JsonizeDataString.to\_json}}{unitree\_sdk2\_cpp.robot.go2.JsonizeDataString.to\_json}}\label{unitree-sdk2-cpp-robot-go2-jsonizedatastring-to-json-1}}

计划把当前 SDK 值对象写入 JSON 风格字典。

\textbf{签名}

\begin{lstlisting}[language=Python]
def to_json(self, json: dict[str, Any]) -> None
\end{lstlisting}

\textbf{可用性与安全性}

\textbf{\passthrough{\lstinline!SIGNATURE\_ONLY!} /
\passthrough{\lstinline!UNCLASSIFIED!}}：当前只有设计期签名，不能假设已存在于运行时扩展。

\textbf{参数}

\begin{longtable}[]{@{}
  >{\raggedright\arraybackslash}p{(\columnwidth - 6\tabcolsep) * \real{0.18}}
  >{\raggedright\arraybackslash}p{(\columnwidth - 6\tabcolsep) * \real{0.24}}
  >{\raggedright\arraybackslash}p{(\columnwidth - 6\tabcolsep) * \real{0.18}}
  >{\raggedright\arraybackslash}p{(\columnwidth - 6\tabcolsep) * \real{0.40}}@{}}
\toprule
名称 & Python 类型 & 默认值 & 含义 \\
\midrule
\endhead
\passthrough{\lstinline!json!} &
\passthrough{\lstinline!dict[str, Any]!} & 必填 & JSON
风格字典，键和值必须满足对应 C++ \passthrough{\lstinline!JsonMap!}
数据结构的约束。 对应 C++ 参数
\passthrough{\lstinline!json: common::JsonMap \&!}。 \\
\bottomrule
\end{longtable}

\textbf{返回值}

\begin{longtable}[]{@{}
  >{\raggedright\arraybackslash}p{(\columnwidth - 4\tabcolsep) * \real{0.18}}
  >{\raggedright\arraybackslash}p{(\columnwidth - 4\tabcolsep) * \real{0.20}}
  >{\raggedright\arraybackslash}p{(\columnwidth - 4\tabcolsep) * \real{0.62}}@{}}
\toprule
位置 & 类型 & 含义 \\
\midrule
\endhead
无 & \passthrough{\lstinline!None!} & 该方法不返回 Python
值；失败可能表现为异常或底层状态变化。 \\
\bottomrule
\end{longtable}

\textbf{对应 C++}

\begin{itemize}
\tightlist
\item
  类：\passthrough{\lstinline!unitree::robot::go2::JsonizeDataString!}
\item
  签名：\passthrough{\lstinline!toJson(common::JsonMap \&) const!}
\item
  绑定策略：\passthrough{\lstinline!SIGNATURE\_PREVIEW!}
\item
  声明位置：\passthrough{\lstinline!include/unitree/robot/go2/public/jsonize\_type.hpp:159!}
\end{itemize}

\textbf{说明}

\begin{itemize}
\tightlist
\item
  示例是局部调用片段；\passthrough{\lstinline!obj!}
  和参数变量表示已经按上层业务校验并准备好的对象。
\item
  该条目可以被编辑器和类型检查器识别，但当前运行时可能在属性查找阶段直接失败。
\end{itemize}

\textbf{用法}

\begin{lstlisting}[language=Python]
from typing import TYPE_CHECKING

if TYPE_CHECKING:
    def planned_to_json(obj: JsonizeDataString, json: dict[str, Any]) -> None:
        obj.to_json(json=json)
\end{lstlisting}

\begin{quote}
{[}!CAUTION{]} 上面的 \passthrough{\lstinline!TYPE\_CHECKING!}
示例只表达计划调用形态。该分支在普通 Python
运行时为假，不代表当前扩展能执行此接口。
\end{quote}

\hypertarget{unitree-sdk2-cpp-robot-go2-jsonizeflagbool}{%
\subsubsection{\texorpdfstring{\texttt{unitree\_sdk2\_cpp.robot.go2.JsonizeFlagBool}}{unitree\_sdk2\_cpp.robot.go2.JsonizeFlagBool}}\label{unitree-sdk2-cpp-robot-go2-jsonizeflagbool}}

SDK 类型签名预览。请逐项查看构造函数和方法的可用性。

\textbf{导入}

\begin{lstlisting}[language=Python]
from unitree_sdk2_cpp.robot.go2 import JsonizeFlagBool
\end{lstlisting}

\textbf{基类}：\passthrough{\lstinline!object!}

\textbf{公开属性}

\begin{longtable}[]{@{}
  >{\raggedright\arraybackslash}p{(\columnwidth - 6\tabcolsep) * \real{0.18}}
  >{\raggedright\arraybackslash}p{(\columnwidth - 6\tabcolsep) * \real{0.24}}
  >{\raggedright\arraybackslash}p{(\columnwidth - 6\tabcolsep) * \real{0.18}}
  >{\raggedright\arraybackslash}p{(\columnwidth - 6\tabcolsep) * \real{0.40}}@{}}
\toprule
名称 & Python 类型 & 含义 & 用法 \\
\midrule
\endhead
\passthrough{\lstinline!flag!} & \passthrough{\lstinline!bool!} &
布尔开关。\passthrough{\lstinline!True!} 和
\passthrough{\lstinline!False!} 的具体业务效果由当前方法定义。 &
\passthrough{\lstinline!value = obj.flag!} /
\passthrough{\lstinline!obj.flag = value!} \\
\bottomrule
\end{longtable}

\textbf{方法索引}

\begin{longtable}[]{@{}
  >{\raggedright\arraybackslash}p{(\columnwidth - 6\tabcolsep) * \real{0.18}}
  >{\raggedright\arraybackslash}p{(\columnwidth - 6\tabcolsep) * \real{0.24}}
  >{\raggedright\arraybackslash}p{(\columnwidth - 6\tabcolsep) * \real{0.18}}
  >{\raggedright\arraybackslash}p{(\columnwidth - 6\tabcolsep) * \real{0.40}}@{}}
\toprule
方法 & Python 签名 & 状态 & 安全分类 \\
\midrule
\endhead
\protect\hyperlink{unitree-sdk2-cpp-robot-go2-jsonizeflagbool-dunder-init-1}{\passthrough{\lstinline!\_\_init\_\_!}}
& \passthrough{\lstinline!def \_\_init\_\_(self) -> None!} &
\passthrough{\lstinline!SIGNATURE\_ONLY!} &
\passthrough{\lstinline!CONSTRUCTION!} \\
\protect\hyperlink{unitree-sdk2-cpp-robot-go2-jsonizeflagbool-from-json-1}{\passthrough{\lstinline!from\_json!}}
&
\passthrough{\lstinline!def from\_json(self, json: dict[str, Any]) -> None!}
& \passthrough{\lstinline!SIGNATURE\_ONLY!} &
\passthrough{\lstinline!UNCLASSIFIED!} \\
\protect\hyperlink{unitree-sdk2-cpp-robot-go2-jsonizeflagbool-to-json-1}{\passthrough{\lstinline!to\_json!}}
&
\passthrough{\lstinline!def to\_json(self, json: dict[str, Any]) -> None!}
& \passthrough{\lstinline!SIGNATURE\_ONLY!} &
\passthrough{\lstinline!UNCLASSIFIED!} \\
\bottomrule
\end{longtable}

\hypertarget{unitree-sdk2-cpp-robot-go2-jsonizeflagbool-dunder-init-1}{%
\paragraph{\texorpdfstring{\texttt{unitree\_sdk2\_cpp.robot.go2.JsonizeFlagBool.\_\_init\_\_}}{unitree\_sdk2\_cpp.robot.go2.JsonizeFlagBool.\_\_init\_\_}}\label{unitree-sdk2-cpp-robot-go2-jsonizeflagbool-dunder-init-1}}

初始化 \passthrough{\lstinline!JsonizeFlagBool!}
实例。是否能实际构造取决于下方可用性状态。

\textbf{签名}

\begin{lstlisting}[language=Python]
def __init__(self) -> None
\end{lstlisting}

\textbf{可用性与安全性}

\textbf{\passthrough{\lstinline!SIGNATURE\_ONLY!} /
\passthrough{\lstinline!CONSTRUCTION!}}：当前只有设计期签名，不能假设已存在于运行时扩展。

\textbf{参数}

无显式参数。实例方法中的 \passthrough{\lstinline!self!} 由 Python
自动传入。

\textbf{返回值}

\begin{longtable}[]{@{}
  >{\raggedright\arraybackslash}p{(\columnwidth - 4\tabcolsep) * \real{0.18}}
  >{\raggedright\arraybackslash}p{(\columnwidth - 4\tabcolsep) * \real{0.20}}
  >{\raggedright\arraybackslash}p{(\columnwidth - 4\tabcolsep) * \real{0.62}}@{}}
\toprule
位置 & 类型 & 含义 \\
\midrule
\endhead
无 & \passthrough{\lstinline!None!} &
\passthrough{\lstinline!\_\_init\_\_!}
本身不返回值；调用类对象时，在构造函数可用的前提下得到该类实例。 \\
\bottomrule
\end{longtable}

\textbf{对应 C++}

\begin{itemize}
\tightlist
\item
  类：\passthrough{\lstinline!unitree::robot::go2::JsonizeFlagBool!}
\item
  签名：\passthrough{\lstinline!JsonizeFlagBool()!}
\item
  绑定策略：\passthrough{\lstinline!未单独标注!}
\item
  声明位置：\passthrough{\lstinline!include/unitree/robot/go2/public/jsonize\_type.hpp:18!}
\end{itemize}

\textbf{说明}

\begin{itemize}
\tightlist
\item
  示例是局部调用片段；\passthrough{\lstinline!obj!}
  和参数变量表示已经按上层业务校验并准备好的对象。
\item
  该条目可以被编辑器和类型检查器识别，但当前运行时可能在属性查找阶段直接失败。
\end{itemize}

\textbf{用法}

\begin{lstlisting}[language=Python]
from typing import TYPE_CHECKING

if TYPE_CHECKING:
    def planned_construct() -> JsonizeFlagBool:
        return JsonizeFlagBool()
\end{lstlisting}

\begin{quote}
{[}!CAUTION{]} 上面的 \passthrough{\lstinline!TYPE\_CHECKING!}
示例只表达计划调用形态。该分支在普通 Python
运行时为假，不代表当前扩展能执行此接口。
\end{quote}

\hypertarget{unitree-sdk2-cpp-robot-go2-jsonizeflagbool-from-json-1}{%
\paragraph{\texorpdfstring{\texttt{unitree\_sdk2\_cpp.robot.go2.JsonizeFlagBool.from\_json}}{unitree\_sdk2\_cpp.robot.go2.JsonizeFlagBool.from\_json}}\label{unitree-sdk2-cpp-robot-go2-jsonizeflagbool-from-json-1}}

计划从 JSON 风格字典读取字段并更新当前 SDK 值对象。

\textbf{签名}

\begin{lstlisting}[language=Python]
def from_json(self, json: dict[str, Any]) -> None
\end{lstlisting}

\textbf{可用性与安全性}

\textbf{\passthrough{\lstinline!SIGNATURE\_ONLY!} /
\passthrough{\lstinline!UNCLASSIFIED!}}：当前只有设计期签名，不能假设已存在于运行时扩展。

\textbf{参数}

\begin{longtable}[]{@{}
  >{\raggedright\arraybackslash}p{(\columnwidth - 6\tabcolsep) * \real{0.18}}
  >{\raggedright\arraybackslash}p{(\columnwidth - 6\tabcolsep) * \real{0.24}}
  >{\raggedright\arraybackslash}p{(\columnwidth - 6\tabcolsep) * \real{0.18}}
  >{\raggedright\arraybackslash}p{(\columnwidth - 6\tabcolsep) * \real{0.40}}@{}}
\toprule
名称 & Python 类型 & 默认值 & 含义 \\
\midrule
\endhead
\passthrough{\lstinline!json!} &
\passthrough{\lstinline!dict[str, Any]!} & 必填 & JSON
风格字典，键和值必须满足对应 C++ \passthrough{\lstinline!JsonMap!}
数据结构的约束。 对应 C++ 参数
\passthrough{\lstinline!json: common::JsonMap \&!}。 \\
\bottomrule
\end{longtable}

\textbf{返回值}

\begin{longtable}[]{@{}
  >{\raggedright\arraybackslash}p{(\columnwidth - 4\tabcolsep) * \real{0.18}}
  >{\raggedright\arraybackslash}p{(\columnwidth - 4\tabcolsep) * \real{0.20}}
  >{\raggedright\arraybackslash}p{(\columnwidth - 4\tabcolsep) * \real{0.62}}@{}}
\toprule
位置 & 类型 & 含义 \\
\midrule
\endhead
无 & \passthrough{\lstinline!None!} & 该方法不返回 Python
值；失败可能表现为异常或底层状态变化。 \\
\bottomrule
\end{longtable}

\textbf{对应 C++}

\begin{itemize}
\tightlist
\item
  类：\passthrough{\lstinline!unitree::robot::go2::JsonizeFlagBool!}
\item
  签名：\passthrough{\lstinline!fromJson(common::JsonMap \&)!}
\item
  绑定策略：\passthrough{\lstinline!SIGNATURE\_PREVIEW!}
\item
  声明位置：\passthrough{\lstinline!include/unitree/robot/go2/public/jsonize\_type.hpp:24!}
\end{itemize}

\textbf{说明}

\begin{itemize}
\tightlist
\item
  示例是局部调用片段；\passthrough{\lstinline!obj!}
  和参数变量表示已经按上层业务校验并准备好的对象。
\item
  该条目可以被编辑器和类型检查器识别，但当前运行时可能在属性查找阶段直接失败。
\end{itemize}

\textbf{用法}

\begin{lstlisting}[language=Python]
from typing import TYPE_CHECKING

if TYPE_CHECKING:
    def planned_from_json(obj: JsonizeFlagBool, json: dict[str, Any]) -> None:
        obj.from_json(json=json)
\end{lstlisting}

\begin{quote}
{[}!CAUTION{]} 上面的 \passthrough{\lstinline!TYPE\_CHECKING!}
示例只表达计划调用形态。该分支在普通 Python
运行时为假，不代表当前扩展能执行此接口。
\end{quote}

\hypertarget{unitree-sdk2-cpp-robot-go2-jsonizeflagbool-to-json-1}{%
\paragraph{\texorpdfstring{\texttt{unitree\_sdk2\_cpp.robot.go2.JsonizeFlagBool.to\_json}}{unitree\_sdk2\_cpp.robot.go2.JsonizeFlagBool.to\_json}}\label{unitree-sdk2-cpp-robot-go2-jsonizeflagbool-to-json-1}}

计划把当前 SDK 值对象写入 JSON 风格字典。

\textbf{签名}

\begin{lstlisting}[language=Python]
def to_json(self, json: dict[str, Any]) -> None
\end{lstlisting}

\textbf{可用性与安全性}

\textbf{\passthrough{\lstinline!SIGNATURE\_ONLY!} /
\passthrough{\lstinline!UNCLASSIFIED!}}：当前只有设计期签名，不能假设已存在于运行时扩展。

\textbf{参数}

\begin{longtable}[]{@{}
  >{\raggedright\arraybackslash}p{(\columnwidth - 6\tabcolsep) * \real{0.18}}
  >{\raggedright\arraybackslash}p{(\columnwidth - 6\tabcolsep) * \real{0.24}}
  >{\raggedright\arraybackslash}p{(\columnwidth - 6\tabcolsep) * \real{0.18}}
  >{\raggedright\arraybackslash}p{(\columnwidth - 6\tabcolsep) * \real{0.40}}@{}}
\toprule
名称 & Python 类型 & 默认值 & 含义 \\
\midrule
\endhead
\passthrough{\lstinline!json!} &
\passthrough{\lstinline!dict[str, Any]!} & 必填 & JSON
风格字典，键和值必须满足对应 C++ \passthrough{\lstinline!JsonMap!}
数据结构的约束。 对应 C++ 参数
\passthrough{\lstinline!json: common::JsonMap \&!}。 \\
\bottomrule
\end{longtable}

\textbf{返回值}

\begin{longtable}[]{@{}
  >{\raggedright\arraybackslash}p{(\columnwidth - 4\tabcolsep) * \real{0.18}}
  >{\raggedright\arraybackslash}p{(\columnwidth - 4\tabcolsep) * \real{0.20}}
  >{\raggedright\arraybackslash}p{(\columnwidth - 4\tabcolsep) * \real{0.62}}@{}}
\toprule
位置 & 类型 & 含义 \\
\midrule
\endhead
无 & \passthrough{\lstinline!None!} & 该方法不返回 Python
值；失败可能表现为异常或底层状态变化。 \\
\bottomrule
\end{longtable}

\textbf{对应 C++}

\begin{itemize}
\tightlist
\item
  类：\passthrough{\lstinline!unitree::robot::go2::JsonizeFlagBool!}
\item
  签名：\passthrough{\lstinline!toJson(common::JsonMap \&) const!}
\item
  绑定策略：\passthrough{\lstinline!SIGNATURE\_PREVIEW!}
\item
  声明位置：\passthrough{\lstinline!include/unitree/robot/go2/public/jsonize\_type.hpp:29!}
\end{itemize}

\textbf{说明}

\begin{itemize}
\tightlist
\item
  示例是局部调用片段；\passthrough{\lstinline!obj!}
  和参数变量表示已经按上层业务校验并准备好的对象。
\item
  该条目可以被编辑器和类型检查器识别，但当前运行时可能在属性查找阶段直接失败。
\end{itemize}

\textbf{用法}

\begin{lstlisting}[language=Python]
from typing import TYPE_CHECKING

if TYPE_CHECKING:
    def planned_to_json(obj: JsonizeFlagBool, json: dict[str, Any]) -> None:
        obj.to_json(json=json)
\end{lstlisting}

\begin{quote}
{[}!CAUTION{]} 上面的 \passthrough{\lstinline!TYPE\_CHECKING!}
示例只表达计划调用形态。该分支在普通 Python
运行时为假，不代表当前扩展能执行此接口。
\end{quote}

\hypertarget{unitree-sdk2-cpp-robot-go2-jsonizepathpoint}{%
\subsubsection{\texorpdfstring{\texttt{unitree\_sdk2\_cpp.robot.go2.JsonizePathPoint}}{unitree\_sdk2\_cpp.robot.go2.JsonizePathPoint}}\label{unitree-sdk2-cpp-robot-go2-jsonizepathpoint}}

SDK 数据值类型；公开字段和转换方法见下文。

\textbf{导入}

\begin{lstlisting}[language=Python]
from unitree_sdk2_cpp.robot.go2 import JsonizePathPoint
\end{lstlisting}

\textbf{基类}：\passthrough{\lstinline!object!}

\textbf{公开属性}

\begin{longtable}[]{@{}
  >{\raggedright\arraybackslash}p{(\columnwidth - 6\tabcolsep) * \real{0.18}}
  >{\raggedright\arraybackslash}p{(\columnwidth - 6\tabcolsep) * \real{0.24}}
  >{\raggedright\arraybackslash}p{(\columnwidth - 6\tabcolsep) * \real{0.18}}
  >{\raggedright\arraybackslash}p{(\columnwidth - 6\tabcolsep) * \real{0.40}}@{}}
\toprule
名称 & Python 类型 & 含义 & 用法 \\
\midrule
\endhead
\passthrough{\lstinline!time\_from\_start!} &
\passthrough{\lstinline!float!} & 传给该接口的
\passthrough{\lstinline!time!} \passthrough{\lstinline!from!}
\passthrough{\lstinline!start!} 参数，Python 类型为
\passthrough{\lstinline!float!}。精确含义、单位、范围和枚举值需查目标型号对应头文件与协议。
& \passthrough{\lstinline!value = obj.time\_from\_start!} /
\passthrough{\lstinline!obj.time\_from\_start = value!} \\
\passthrough{\lstinline!x!} & \passthrough{\lstinline!float!} & X
方向位置或分量参数。参考系、单位和范围必须查当前方法及目标型号协议。 &
\passthrough{\lstinline!value = obj.x!} /
\passthrough{\lstinline!obj.x = value!} \\
\passthrough{\lstinline!y!} & \passthrough{\lstinline!float!} & Y
方向位置或分量参数。参考系、单位和范围必须查当前方法及目标型号协议。 &
\passthrough{\lstinline!value = obj.y!} /
\passthrough{\lstinline!obj.y = value!} \\
\passthrough{\lstinline!yaw!} & \passthrough{\lstinline!float!} &
偏航角或偏航目标参数。参考系、单位和范围必须查目标型号协议。 &
\passthrough{\lstinline!value = obj.yaw!} /
\passthrough{\lstinline!obj.yaw = value!} \\
\passthrough{\lstinline!vx!} & \passthrough{\lstinline!float!} & X
方向速度参数。坐标系、单位、符号和安全范围必须查目标型号运动协议。 &
\passthrough{\lstinline!value = obj.vx!} /
\passthrough{\lstinline!obj.vx = value!} \\
\passthrough{\lstinline!vy!} & \passthrough{\lstinline!float!} & Y
方向速度参数。坐标系、单位、符号和安全范围必须查目标型号运动协议。 &
\passthrough{\lstinline!value = obj.vy!} /
\passthrough{\lstinline!obj.vy = value!} \\
\passthrough{\lstinline!vyaw!} & \passthrough{\lstinline!float!} &
偏航角速度参数。单位、符号和安全范围必须查目标型号运动协议。 &
\passthrough{\lstinline!value = obj.vyaw!} /
\passthrough{\lstinline!obj.vyaw = value!} \\
\bottomrule
\end{longtable}

\textbf{方法索引}

\begin{longtable}[]{@{}
  >{\raggedright\arraybackslash}p{(\columnwidth - 6\tabcolsep) * \real{0.18}}
  >{\raggedright\arraybackslash}p{(\columnwidth - 6\tabcolsep) * \real{0.24}}
  >{\raggedright\arraybackslash}p{(\columnwidth - 6\tabcolsep) * \real{0.18}}
  >{\raggedright\arraybackslash}p{(\columnwidth - 6\tabcolsep) * \real{0.40}}@{}}
\toprule
方法 & Python 签名 & 状态 & 安全分类 \\
\midrule
\endhead
\protect\hyperlink{unitree-sdk2-cpp-robot-go2-jsonizepathpoint-dunder-init-1}{\passthrough{\lstinline!\_\_init\_\_!}}
& \passthrough{\lstinline!def \_\_init\_\_(self) -> None!} &
\passthrough{\lstinline!SIGNATURE\_ONLY!} &
\passthrough{\lstinline!CONSTRUCTION!} \\
\protect\hyperlink{unitree-sdk2-cpp-robot-go2-jsonizepathpoint-from-json-1}{\passthrough{\lstinline!from\_json!}}
&
\passthrough{\lstinline!def from\_json(self, json: dict[str, Any]) -> None!}
& \passthrough{\lstinline!SIGNATURE\_ONLY!} &
\passthrough{\lstinline!UNCLASSIFIED!} \\
\protect\hyperlink{unitree-sdk2-cpp-robot-go2-jsonizepathpoint-to-json-1}{\passthrough{\lstinline!to\_json!}}
&
\passthrough{\lstinline!def to\_json(self, json: dict[str, Any]) -> None!}
& \passthrough{\lstinline!SIGNATURE\_ONLY!} &
\passthrough{\lstinline!UNCLASSIFIED!} \\
\bottomrule
\end{longtable}

\hypertarget{unitree-sdk2-cpp-robot-go2-jsonizepathpoint-dunder-init-1}{%
\paragraph{\texorpdfstring{\texttt{unitree\_sdk2\_cpp.robot.go2.JsonizePathPoint.\_\_init\_\_}}{unitree\_sdk2\_cpp.robot.go2.JsonizePathPoint.\_\_init\_\_}}\label{unitree-sdk2-cpp-robot-go2-jsonizepathpoint-dunder-init-1}}

初始化 \passthrough{\lstinline!JsonizePathPoint!}
实例。是否能实际构造取决于下方可用性状态。

\textbf{签名}

\begin{lstlisting}[language=Python]
def __init__(self) -> None
\end{lstlisting}

\textbf{可用性与安全性}

\textbf{\passthrough{\lstinline!SIGNATURE\_ONLY!} /
\passthrough{\lstinline!CONSTRUCTION!}}：当前只有设计期签名，不能假设已存在于运行时扩展。

\textbf{参数}

无显式参数。实例方法中的 \passthrough{\lstinline!self!} 由 Python
自动传入。

\textbf{返回值}

\begin{longtable}[]{@{}
  >{\raggedright\arraybackslash}p{(\columnwidth - 4\tabcolsep) * \real{0.18}}
  >{\raggedright\arraybackslash}p{(\columnwidth - 4\tabcolsep) * \real{0.20}}
  >{\raggedright\arraybackslash}p{(\columnwidth - 4\tabcolsep) * \real{0.62}}@{}}
\toprule
位置 & 类型 & 含义 \\
\midrule
\endhead
无 & \passthrough{\lstinline!None!} &
\passthrough{\lstinline!\_\_init\_\_!}
本身不返回值；调用类对象时，在构造函数可用的前提下得到该类实例。 \\
\bottomrule
\end{longtable}

\textbf{对应 C++}

\begin{itemize}
\tightlist
\item
  类：\passthrough{\lstinline!unitree::robot::go2::JsonizePathPoint!}
\item
  签名：\passthrough{\lstinline!JsonizePathPoint()!}
\item
  绑定策略：\passthrough{\lstinline!未单独标注!}
\item
  声明位置：\passthrough{\lstinline!include/unitree/robot/go2/public/jsonize\_type.hpp:241!}
\end{itemize}

\textbf{说明}

\begin{itemize}
\tightlist
\item
  示例是局部调用片段；\passthrough{\lstinline!obj!}
  和参数变量表示已经按上层业务校验并准备好的对象。
\item
  该条目可以被编辑器和类型检查器识别，但当前运行时可能在属性查找阶段直接失败。
\end{itemize}

\textbf{用法}

\begin{lstlisting}[language=Python]
from typing import TYPE_CHECKING

if TYPE_CHECKING:
    def planned_construct() -> JsonizePathPoint:
        return JsonizePathPoint()
\end{lstlisting}

\begin{quote}
{[}!CAUTION{]} 上面的 \passthrough{\lstinline!TYPE\_CHECKING!}
示例只表达计划调用形态。该分支在普通 Python
运行时为假，不代表当前扩展能执行此接口。
\end{quote}

\hypertarget{unitree-sdk2-cpp-robot-go2-jsonizepathpoint-from-json-1}{%
\paragraph{\texorpdfstring{\texttt{unitree\_sdk2\_cpp.robot.go2.JsonizePathPoint.from\_json}}{unitree\_sdk2\_cpp.robot.go2.JsonizePathPoint.from\_json}}\label{unitree-sdk2-cpp-robot-go2-jsonizepathpoint-from-json-1}}

计划从 JSON 风格字典读取字段并更新当前 SDK 值对象。

\textbf{签名}

\begin{lstlisting}[language=Python]
def from_json(self, json: dict[str, Any]) -> None
\end{lstlisting}

\textbf{可用性与安全性}

\textbf{\passthrough{\lstinline!SIGNATURE\_ONLY!} /
\passthrough{\lstinline!UNCLASSIFIED!}}：当前只有设计期签名，不能假设已存在于运行时扩展。

\textbf{参数}

\begin{longtable}[]{@{}
  >{\raggedright\arraybackslash}p{(\columnwidth - 6\tabcolsep) * \real{0.18}}
  >{\raggedright\arraybackslash}p{(\columnwidth - 6\tabcolsep) * \real{0.24}}
  >{\raggedright\arraybackslash}p{(\columnwidth - 6\tabcolsep) * \real{0.18}}
  >{\raggedright\arraybackslash}p{(\columnwidth - 6\tabcolsep) * \real{0.40}}@{}}
\toprule
名称 & Python 类型 & 默认值 & 含义 \\
\midrule
\endhead
\passthrough{\lstinline!json!} &
\passthrough{\lstinline!dict[str, Any]!} & 必填 & JSON
风格字典，键和值必须满足对应 C++ \passthrough{\lstinline!JsonMap!}
数据结构的约束。 对应 C++ 参数
\passthrough{\lstinline!json: common::JsonMap \&!}。 \\
\bottomrule
\end{longtable}

\textbf{返回值}

\begin{longtable}[]{@{}
  >{\raggedright\arraybackslash}p{(\columnwidth - 4\tabcolsep) * \real{0.18}}
  >{\raggedright\arraybackslash}p{(\columnwidth - 4\tabcolsep) * \real{0.20}}
  >{\raggedright\arraybackslash}p{(\columnwidth - 4\tabcolsep) * \real{0.62}}@{}}
\toprule
位置 & 类型 & 含义 \\
\midrule
\endhead
无 & \passthrough{\lstinline!None!} & 该方法不返回 Python
值；失败可能表现为异常或底层状态变化。 \\
\bottomrule
\end{longtable}

\textbf{对应 C++}

\begin{itemize}
\tightlist
\item
  类：\passthrough{\lstinline!unitree::robot::go2::JsonizePathPoint!}
\item
  签名：\passthrough{\lstinline!fromJson(common::JsonMap \&)!}
\item
  绑定策略：\passthrough{\lstinline!SIGNATURE\_PREVIEW!}
\item
  声明位置：\passthrough{\lstinline!include/unitree/robot/go2/public/jsonize\_type.hpp:249!}
\end{itemize}

\textbf{说明}

\begin{itemize}
\tightlist
\item
  示例是局部调用片段；\passthrough{\lstinline!obj!}
  和参数变量表示已经按上层业务校验并准备好的对象。
\item
  该条目可以被编辑器和类型检查器识别，但当前运行时可能在属性查找阶段直接失败。
\end{itemize}

\textbf{用法}

\begin{lstlisting}[language=Python]
from typing import TYPE_CHECKING

if TYPE_CHECKING:
    def planned_from_json(obj: JsonizePathPoint, json: dict[str, Any]) -> None:
        obj.from_json(json=json)
\end{lstlisting}

\begin{quote}
{[}!CAUTION{]} 上面的 \passthrough{\lstinline!TYPE\_CHECKING!}
示例只表达计划调用形态。该分支在普通 Python
运行时为假，不代表当前扩展能执行此接口。
\end{quote}

\hypertarget{unitree-sdk2-cpp-robot-go2-jsonizepathpoint-to-json-1}{%
\paragraph{\texorpdfstring{\texttt{unitree\_sdk2\_cpp.robot.go2.JsonizePathPoint.to\_json}}{unitree\_sdk2\_cpp.robot.go2.JsonizePathPoint.to\_json}}\label{unitree-sdk2-cpp-robot-go2-jsonizepathpoint-to-json-1}}

计划把当前 SDK 值对象写入 JSON 风格字典。

\textbf{签名}

\begin{lstlisting}[language=Python]
def to_json(self, json: dict[str, Any]) -> None
\end{lstlisting}

\textbf{可用性与安全性}

\textbf{\passthrough{\lstinline!SIGNATURE\_ONLY!} /
\passthrough{\lstinline!UNCLASSIFIED!}}：当前只有设计期签名，不能假设已存在于运行时扩展。

\textbf{参数}

\begin{longtable}[]{@{}
  >{\raggedright\arraybackslash}p{(\columnwidth - 6\tabcolsep) * \real{0.18}}
  >{\raggedright\arraybackslash}p{(\columnwidth - 6\tabcolsep) * \real{0.24}}
  >{\raggedright\arraybackslash}p{(\columnwidth - 6\tabcolsep) * \real{0.18}}
  >{\raggedright\arraybackslash}p{(\columnwidth - 6\tabcolsep) * \real{0.40}}@{}}
\toprule
名称 & Python 类型 & 默认值 & 含义 \\
\midrule
\endhead
\passthrough{\lstinline!json!} &
\passthrough{\lstinline!dict[str, Any]!} & 必填 & JSON
风格字典，键和值必须满足对应 C++ \passthrough{\lstinline!JsonMap!}
数据结构的约束。 对应 C++ 参数
\passthrough{\lstinline!json: common::JsonMap \&!}。 \\
\bottomrule
\end{longtable}

\textbf{返回值}

\begin{longtable}[]{@{}
  >{\raggedright\arraybackslash}p{(\columnwidth - 4\tabcolsep) * \real{0.18}}
  >{\raggedright\arraybackslash}p{(\columnwidth - 4\tabcolsep) * \real{0.20}}
  >{\raggedright\arraybackslash}p{(\columnwidth - 4\tabcolsep) * \real{0.62}}@{}}
\toprule
位置 & 类型 & 含义 \\
\midrule
\endhead
无 & \passthrough{\lstinline!None!} & 该方法不返回 Python
值；失败可能表现为异常或底层状态变化。 \\
\bottomrule
\end{longtable}

\textbf{对应 C++}

\begin{itemize}
\tightlist
\item
  类：\passthrough{\lstinline!unitree::robot::go2::JsonizePathPoint!}
\item
  签名：\passthrough{\lstinline!toJson(common::JsonMap \&) const!}
\item
  绑定策略：\passthrough{\lstinline!SIGNATURE\_PREVIEW!}
\item
  声明位置：\passthrough{\lstinline!include/unitree/robot/go2/public/jsonize\_type.hpp:260!}
\end{itemize}

\textbf{说明}

\begin{itemize}
\tightlist
\item
  示例是局部调用片段；\passthrough{\lstinline!obj!}
  和参数变量表示已经按上层业务校验并准备好的对象。
\item
  该条目可以被编辑器和类型检查器识别，但当前运行时可能在属性查找阶段直接失败。
\end{itemize}

\textbf{用法}

\begin{lstlisting}[language=Python]
from typing import TYPE_CHECKING

if TYPE_CHECKING:
    def planned_to_json(obj: JsonizePathPoint, json: dict[str, Any]) -> None:
        obj.to_json(json=json)
\end{lstlisting}

\begin{quote}
{[}!CAUTION{]} 上面的 \passthrough{\lstinline!TYPE\_CHECKING!}
示例只表达计划调用形态。该分支在普通 Python
运行时为假，不代表当前扩展能执行此接口。
\end{quote}

\hypertarget{unitree-sdk2-cpp-robot-go2-jsonizequat}{%
\subsubsection{\texorpdfstring{\texttt{unitree\_sdk2\_cpp.robot.go2.JsonizeQuat}}{unitree\_sdk2\_cpp.robot.go2.JsonizeQuat}}\label{unitree-sdk2-cpp-robot-go2-jsonizequat}}

SDK 类型签名预览。请逐项查看构造函数和方法的可用性。

\textbf{导入}

\begin{lstlisting}[language=Python]
from unitree_sdk2_cpp.robot.go2 import JsonizeQuat
\end{lstlisting}

\textbf{基类}：\passthrough{\lstinline!object!}

\textbf{公开属性}

\begin{longtable}[]{@{}
  >{\raggedright\arraybackslash}p{(\columnwidth - 6\tabcolsep) * \real{0.18}}
  >{\raggedright\arraybackslash}p{(\columnwidth - 6\tabcolsep) * \real{0.24}}
  >{\raggedright\arraybackslash}p{(\columnwidth - 6\tabcolsep) * \real{0.18}}
  >{\raggedright\arraybackslash}p{(\columnwidth - 6\tabcolsep) * \real{0.40}}@{}}
\toprule
名称 & Python 类型 & 含义 & 用法 \\
\midrule
\endhead
\passthrough{\lstinline!x!} & \passthrough{\lstinline!float!} & X
方向位置或分量参数。参考系、单位和范围必须查当前方法及目标型号协议。 &
\passthrough{\lstinline!value = obj.x!} /
\passthrough{\lstinline!obj.x = value!} \\
\passthrough{\lstinline!y!} & \passthrough{\lstinline!float!} & Y
方向位置或分量参数。参考系、单位和范围必须查当前方法及目标型号协议。 &
\passthrough{\lstinline!value = obj.y!} /
\passthrough{\lstinline!obj.y = value!} \\
\passthrough{\lstinline!z!} & \passthrough{\lstinline!float!} & Z
方向位置或分量参数。参考系、单位和范围必须查当前方法及目标型号协议。 &
\passthrough{\lstinline!value = obj.z!} /
\passthrough{\lstinline!obj.z = value!} \\
\passthrough{\lstinline!w!} & \passthrough{\lstinline!float!} &
传给该接口的 \passthrough{\lstinline!w!} 参数，Python 类型为
\passthrough{\lstinline!float!}。精确含义、单位、范围和枚举值需查目标型号对应头文件与协议。
& \passthrough{\lstinline!value = obj.w!} /
\passthrough{\lstinline!obj.w = value!} \\
\bottomrule
\end{longtable}

\textbf{方法索引}

\begin{longtable}[]{@{}
  >{\raggedright\arraybackslash}p{(\columnwidth - 6\tabcolsep) * \real{0.18}}
  >{\raggedright\arraybackslash}p{(\columnwidth - 6\tabcolsep) * \real{0.24}}
  >{\raggedright\arraybackslash}p{(\columnwidth - 6\tabcolsep) * \real{0.18}}
  >{\raggedright\arraybackslash}p{(\columnwidth - 6\tabcolsep) * \real{0.40}}@{}}
\toprule
方法 & Python 签名 & 状态 & 安全分类 \\
\midrule
\endhead
\protect\hyperlink{unitree-sdk2-cpp-robot-go2-jsonizequat-dunder-init-1}{\passthrough{\lstinline!\_\_init\_\_!}}
& \passthrough{\lstinline!def \_\_init\_\_(self) -> None!} &
\passthrough{\lstinline!SIGNATURE\_ONLY!} &
\passthrough{\lstinline!CONSTRUCTION!} \\
\protect\hyperlink{unitree-sdk2-cpp-robot-go2-jsonizequat-from-json-1}{\passthrough{\lstinline!from\_json!}}
&
\passthrough{\lstinline!def from\_json(self, json: dict[str, Any]) -> None!}
& \passthrough{\lstinline!SIGNATURE\_ONLY!} &
\passthrough{\lstinline!UNCLASSIFIED!} \\
\protect\hyperlink{unitree-sdk2-cpp-robot-go2-jsonizequat-to-json-1}{\passthrough{\lstinline!to\_json!}}
&
\passthrough{\lstinline!def to\_json(self, json: dict[str, Any]) -> None!}
& \passthrough{\lstinline!SIGNATURE\_ONLY!} &
\passthrough{\lstinline!UNCLASSIFIED!} \\
\bottomrule
\end{longtable}

\hypertarget{unitree-sdk2-cpp-robot-go2-jsonizequat-dunder-init-1}{%
\paragraph{\texorpdfstring{\texttt{unitree\_sdk2\_cpp.robot.go2.JsonizeQuat.\_\_init\_\_}}{unitree\_sdk2\_cpp.robot.go2.JsonizeQuat.\_\_init\_\_}}\label{unitree-sdk2-cpp-robot-go2-jsonizequat-dunder-init-1}}

初始化 \passthrough{\lstinline!JsonizeQuat!}
实例。是否能实际构造取决于下方可用性状态。

\textbf{签名}

\begin{lstlisting}[language=Python]
def __init__(self) -> None
\end{lstlisting}

\textbf{可用性与安全性}

\textbf{\passthrough{\lstinline!SIGNATURE\_ONLY!} /
\passthrough{\lstinline!CONSTRUCTION!}}：当前只有设计期签名，不能假设已存在于运行时扩展。

\textbf{参数}

无显式参数。实例方法中的 \passthrough{\lstinline!self!} 由 Python
自动传入。

\textbf{返回值}

\begin{longtable}[]{@{}
  >{\raggedright\arraybackslash}p{(\columnwidth - 4\tabcolsep) * \real{0.18}}
  >{\raggedright\arraybackslash}p{(\columnwidth - 4\tabcolsep) * \real{0.20}}
  >{\raggedright\arraybackslash}p{(\columnwidth - 4\tabcolsep) * \real{0.62}}@{}}
\toprule
位置 & 类型 & 含义 \\
\midrule
\endhead
无 & \passthrough{\lstinline!None!} &
\passthrough{\lstinline!\_\_init\_\_!}
本身不返回值；调用类对象时，在构造函数可用的前提下得到该类实例。 \\
\bottomrule
\end{longtable}

\textbf{对应 C++}

\begin{itemize}
\tightlist
\item
  类：\passthrough{\lstinline!unitree::robot::go2::JsonizeQuat!}
\item
  签名：\passthrough{\lstinline!JsonizeQuat()!}
\item
  绑定策略：\passthrough{\lstinline!未单独标注!}
\item
  声明位置：\passthrough{\lstinline!include/unitree/robot/go2/public/jsonize\_type.hpp:206!}
\end{itemize}

\textbf{说明}

\begin{itemize}
\tightlist
\item
  示例是局部调用片段；\passthrough{\lstinline!obj!}
  和参数变量表示已经按上层业务校验并准备好的对象。
\item
  该条目可以被编辑器和类型检查器识别，但当前运行时可能在属性查找阶段直接失败。
\end{itemize}

\textbf{用法}

\begin{lstlisting}[language=Python]
from typing import TYPE_CHECKING

if TYPE_CHECKING:
    def planned_construct() -> JsonizeQuat:
        return JsonizeQuat()
\end{lstlisting}

\begin{quote}
{[}!CAUTION{]} 上面的 \passthrough{\lstinline!TYPE\_CHECKING!}
示例只表达计划调用形态。该分支在普通 Python
运行时为假，不代表当前扩展能执行此接口。
\end{quote}

\hypertarget{unitree-sdk2-cpp-robot-go2-jsonizequat-from-json-1}{%
\paragraph{\texorpdfstring{\texttt{unitree\_sdk2\_cpp.robot.go2.JsonizeQuat.from\_json}}{unitree\_sdk2\_cpp.robot.go2.JsonizeQuat.from\_json}}\label{unitree-sdk2-cpp-robot-go2-jsonizequat-from-json-1}}

计划从 JSON 风格字典读取字段并更新当前 SDK 值对象。

\textbf{签名}

\begin{lstlisting}[language=Python]
def from_json(self, json: dict[str, Any]) -> None
\end{lstlisting}

\textbf{可用性与安全性}

\textbf{\passthrough{\lstinline!SIGNATURE\_ONLY!} /
\passthrough{\lstinline!UNCLASSIFIED!}}：当前只有设计期签名，不能假设已存在于运行时扩展。

\textbf{参数}

\begin{longtable}[]{@{}
  >{\raggedright\arraybackslash}p{(\columnwidth - 6\tabcolsep) * \real{0.18}}
  >{\raggedright\arraybackslash}p{(\columnwidth - 6\tabcolsep) * \real{0.24}}
  >{\raggedright\arraybackslash}p{(\columnwidth - 6\tabcolsep) * \real{0.18}}
  >{\raggedright\arraybackslash}p{(\columnwidth - 6\tabcolsep) * \real{0.40}}@{}}
\toprule
名称 & Python 类型 & 默认值 & 含义 \\
\midrule
\endhead
\passthrough{\lstinline!json!} &
\passthrough{\lstinline!dict[str, Any]!} & 必填 & JSON
风格字典，键和值必须满足对应 C++ \passthrough{\lstinline!JsonMap!}
数据结构的约束。 对应 C++ 参数
\passthrough{\lstinline!json: common::JsonMap \&!}。 \\
\bottomrule
\end{longtable}

\textbf{返回值}

\begin{longtable}[]{@{}
  >{\raggedright\arraybackslash}p{(\columnwidth - 4\tabcolsep) * \real{0.18}}
  >{\raggedright\arraybackslash}p{(\columnwidth - 4\tabcolsep) * \real{0.20}}
  >{\raggedright\arraybackslash}p{(\columnwidth - 4\tabcolsep) * \real{0.62}}@{}}
\toprule
位置 & 类型 & 含义 \\
\midrule
\endhead
无 & \passthrough{\lstinline!None!} & 该方法不返回 Python
值；失败可能表现为异常或底层状态变化。 \\
\bottomrule
\end{longtable}

\textbf{对应 C++}

\begin{itemize}
\tightlist
\item
  类：\passthrough{\lstinline!unitree::robot::go2::JsonizeQuat!}
\item
  签名：\passthrough{\lstinline!fromJson(common::JsonMap \&)!}
\item
  绑定策略：\passthrough{\lstinline!SIGNATURE\_PREVIEW!}
\item
  声明位置：\passthrough{\lstinline!include/unitree/robot/go2/public/jsonize\_type.hpp:212!}
\end{itemize}

\textbf{说明}

\begin{itemize}
\tightlist
\item
  示例是局部调用片段；\passthrough{\lstinline!obj!}
  和参数变量表示已经按上层业务校验并准备好的对象。
\item
  该条目可以被编辑器和类型检查器识别，但当前运行时可能在属性查找阶段直接失败。
\end{itemize}

\textbf{用法}

\begin{lstlisting}[language=Python]
from typing import TYPE_CHECKING

if TYPE_CHECKING:
    def planned_from_json(obj: JsonizeQuat, json: dict[str, Any]) -> None:
        obj.from_json(json=json)
\end{lstlisting}

\begin{quote}
{[}!CAUTION{]} 上面的 \passthrough{\lstinline!TYPE\_CHECKING!}
示例只表达计划调用形态。该分支在普通 Python
运行时为假，不代表当前扩展能执行此接口。
\end{quote}

\hypertarget{unitree-sdk2-cpp-robot-go2-jsonizequat-to-json-1}{%
\paragraph{\texorpdfstring{\texttt{unitree\_sdk2\_cpp.robot.go2.JsonizeQuat.to\_json}}{unitree\_sdk2\_cpp.robot.go2.JsonizeQuat.to\_json}}\label{unitree-sdk2-cpp-robot-go2-jsonizequat-to-json-1}}

计划把当前 SDK 值对象写入 JSON 风格字典。

\textbf{签名}

\begin{lstlisting}[language=Python]
def to_json(self, json: dict[str, Any]) -> None
\end{lstlisting}

\textbf{可用性与安全性}

\textbf{\passthrough{\lstinline!SIGNATURE\_ONLY!} /
\passthrough{\lstinline!UNCLASSIFIED!}}：当前只有设计期签名，不能假设已存在于运行时扩展。

\textbf{参数}

\begin{longtable}[]{@{}
  >{\raggedright\arraybackslash}p{(\columnwidth - 6\tabcolsep) * \real{0.18}}
  >{\raggedright\arraybackslash}p{(\columnwidth - 6\tabcolsep) * \real{0.24}}
  >{\raggedright\arraybackslash}p{(\columnwidth - 6\tabcolsep) * \real{0.18}}
  >{\raggedright\arraybackslash}p{(\columnwidth - 6\tabcolsep) * \real{0.40}}@{}}
\toprule
名称 & Python 类型 & 默认值 & 含义 \\
\midrule
\endhead
\passthrough{\lstinline!json!} &
\passthrough{\lstinline!dict[str, Any]!} & 必填 & JSON
风格字典，键和值必须满足对应 C++ \passthrough{\lstinline!JsonMap!}
数据结构的约束。 对应 C++ 参数
\passthrough{\lstinline!json: common::JsonMap \&!}。 \\
\bottomrule
\end{longtable}

\textbf{返回值}

\begin{longtable}[]{@{}
  >{\raggedright\arraybackslash}p{(\columnwidth - 4\tabcolsep) * \real{0.18}}
  >{\raggedright\arraybackslash}p{(\columnwidth - 4\tabcolsep) * \real{0.20}}
  >{\raggedright\arraybackslash}p{(\columnwidth - 4\tabcolsep) * \real{0.62}}@{}}
\toprule
位置 & 类型 & 含义 \\
\midrule
\endhead
无 & \passthrough{\lstinline!None!} & 该方法不返回 Python
值；失败可能表现为异常或底层状态变化。 \\
\bottomrule
\end{longtable}

\textbf{对应 C++}

\begin{itemize}
\tightlist
\item
  类：\passthrough{\lstinline!unitree::robot::go2::JsonizeQuat!}
\item
  签名：\passthrough{\lstinline!toJson(common::JsonMap \&) const!}
\item
  绑定策略：\passthrough{\lstinline!SIGNATURE\_PREVIEW!}
\item
  声明位置：\passthrough{\lstinline!include/unitree/robot/go2/public/jsonize\_type.hpp:220!}
\end{itemize}

\textbf{说明}

\begin{itemize}
\tightlist
\item
  示例是局部调用片段；\passthrough{\lstinline!obj!}
  和参数变量表示已经按上层业务校验并准备好的对象。
\item
  该条目可以被编辑器和类型检查器识别，但当前运行时可能在属性查找阶段直接失败。
\end{itemize}

\textbf{用法}

\begin{lstlisting}[language=Python]
from typing import TYPE_CHECKING

if TYPE_CHECKING:
    def planned_to_json(obj: JsonizeQuat, json: dict[str, Any]) -> None:
        obj.to_json(json=json)
\end{lstlisting}

\begin{quote}
{[}!CAUTION{]} 上面的 \passthrough{\lstinline!TYPE\_CHECKING!}
示例只表达计划调用形态。该分支在普通 Python
运行时为假，不代表当前扩展能执行此接口。
\end{quote}

\hypertarget{unitree-sdk2-cpp-robot-go2-jsonizevec3}{%
\subsubsection{\texorpdfstring{\texttt{unitree\_sdk2\_cpp.robot.go2.JsonizeVec3}}{unitree\_sdk2\_cpp.robot.go2.JsonizeVec3}}\label{unitree-sdk2-cpp-robot-go2-jsonizevec3}}

SDK 类型签名预览。请逐项查看构造函数和方法的可用性。

\textbf{导入}

\begin{lstlisting}[language=Python]
from unitree_sdk2_cpp.robot.go2 import JsonizeVec3
\end{lstlisting}

\textbf{基类}：\passthrough{\lstinline!object!}

\textbf{公开属性}

\begin{longtable}[]{@{}
  >{\raggedright\arraybackslash}p{(\columnwidth - 6\tabcolsep) * \real{0.18}}
  >{\raggedright\arraybackslash}p{(\columnwidth - 6\tabcolsep) * \real{0.24}}
  >{\raggedright\arraybackslash}p{(\columnwidth - 6\tabcolsep) * \real{0.18}}
  >{\raggedright\arraybackslash}p{(\columnwidth - 6\tabcolsep) * \real{0.40}}@{}}
\toprule
名称 & Python 类型 & 含义 & 用法 \\
\midrule
\endhead
\passthrough{\lstinline!x!} & \passthrough{\lstinline!float!} & X
方向位置或分量参数。参考系、单位和范围必须查当前方法及目标型号协议。 &
\passthrough{\lstinline!value = obj.x!} /
\passthrough{\lstinline!obj.x = value!} \\
\passthrough{\lstinline!y!} & \passthrough{\lstinline!float!} & Y
方向位置或分量参数。参考系、单位和范围必须查当前方法及目标型号协议。 &
\passthrough{\lstinline!value = obj.y!} /
\passthrough{\lstinline!obj.y = value!} \\
\passthrough{\lstinline!z!} & \passthrough{\lstinline!float!} & Z
方向位置或分量参数。参考系、单位和范围必须查当前方法及目标型号协议。 &
\passthrough{\lstinline!value = obj.z!} /
\passthrough{\lstinline!obj.z = value!} \\
\bottomrule
\end{longtable}

\textbf{方法索引}

\begin{longtable}[]{@{}
  >{\raggedright\arraybackslash}p{(\columnwidth - 6\tabcolsep) * \real{0.18}}
  >{\raggedright\arraybackslash}p{(\columnwidth - 6\tabcolsep) * \real{0.24}}
  >{\raggedright\arraybackslash}p{(\columnwidth - 6\tabcolsep) * \real{0.18}}
  >{\raggedright\arraybackslash}p{(\columnwidth - 6\tabcolsep) * \real{0.40}}@{}}
\toprule
方法 & Python 签名 & 状态 & 安全分类 \\
\midrule
\endhead
\protect\hyperlink{unitree-sdk2-cpp-robot-go2-jsonizevec3-dunder-init-1}{\passthrough{\lstinline!\_\_init\_\_!}}
& \passthrough{\lstinline!def \_\_init\_\_(self) -> None!} &
\passthrough{\lstinline!SIGNATURE\_ONLY!} &
\passthrough{\lstinline!CONSTRUCTION!} \\
\protect\hyperlink{unitree-sdk2-cpp-robot-go2-jsonizevec3-from-json-1}{\passthrough{\lstinline!from\_json!}}
&
\passthrough{\lstinline!def from\_json(self, json: dict[str, Any]) -> None!}
& \passthrough{\lstinline!SIGNATURE\_ONLY!} &
\passthrough{\lstinline!UNCLASSIFIED!} \\
\protect\hyperlink{unitree-sdk2-cpp-robot-go2-jsonizevec3-to-json-1}{\passthrough{\lstinline!to\_json!}}
&
\passthrough{\lstinline!def to\_json(self, json: dict[str, Any]) -> None!}
& \passthrough{\lstinline!SIGNATURE\_ONLY!} &
\passthrough{\lstinline!UNCLASSIFIED!} \\
\bottomrule
\end{longtable}

\hypertarget{unitree-sdk2-cpp-robot-go2-jsonizevec3-dunder-init-1}{%
\paragraph{\texorpdfstring{\texttt{unitree\_sdk2\_cpp.robot.go2.JsonizeVec3.\_\_init\_\_}}{unitree\_sdk2\_cpp.robot.go2.JsonizeVec3.\_\_init\_\_}}\label{unitree-sdk2-cpp-robot-go2-jsonizevec3-dunder-init-1}}

初始化 \passthrough{\lstinline!JsonizeVec3!}
实例。是否能实际构造取决于下方可用性状态。

\textbf{签名}

\begin{lstlisting}[language=Python]
def __init__(self) -> None
\end{lstlisting}

\textbf{可用性与安全性}

\textbf{\passthrough{\lstinline!SIGNATURE\_ONLY!} /
\passthrough{\lstinline!CONSTRUCTION!}}：当前只有设计期签名，不能假设已存在于运行时扩展。

\textbf{参数}

无显式参数。实例方法中的 \passthrough{\lstinline!self!} 由 Python
自动传入。

\textbf{返回值}

\begin{longtable}[]{@{}
  >{\raggedright\arraybackslash}p{(\columnwidth - 4\tabcolsep) * \real{0.18}}
  >{\raggedright\arraybackslash}p{(\columnwidth - 4\tabcolsep) * \real{0.20}}
  >{\raggedright\arraybackslash}p{(\columnwidth - 4\tabcolsep) * \real{0.62}}@{}}
\toprule
位置 & 类型 & 含义 \\
\midrule
\endhead
无 & \passthrough{\lstinline!None!} &
\passthrough{\lstinline!\_\_init\_\_!}
本身不返回值；调用类对象时，在构造函数可用的前提下得到该类实例。 \\
\bottomrule
\end{longtable}

\textbf{对应 C++}

\begin{itemize}
\tightlist
\item
  类：\passthrough{\lstinline!unitree::robot::go2::JsonizeVec3!}
\item
  签名：\passthrough{\lstinline!JsonizeVec3()!}
\item
  绑定策略：\passthrough{\lstinline!未单独标注!}
\item
  声明位置：\passthrough{\lstinline!include/unitree/robot/go2/public/jsonize\_type.hpp:174!}
\end{itemize}

\textbf{说明}

\begin{itemize}
\tightlist
\item
  示例是局部调用片段；\passthrough{\lstinline!obj!}
  和参数变量表示已经按上层业务校验并准备好的对象。
\item
  该条目可以被编辑器和类型检查器识别，但当前运行时可能在属性查找阶段直接失败。
\end{itemize}

\textbf{用法}

\begin{lstlisting}[language=Python]
from typing import TYPE_CHECKING

if TYPE_CHECKING:
    def planned_construct() -> JsonizeVec3:
        return JsonizeVec3()
\end{lstlisting}

\begin{quote}
{[}!CAUTION{]} 上面的 \passthrough{\lstinline!TYPE\_CHECKING!}
示例只表达计划调用形态。该分支在普通 Python
运行时为假，不代表当前扩展能执行此接口。
\end{quote}

\hypertarget{unitree-sdk2-cpp-robot-go2-jsonizevec3-from-json-1}{%
\paragraph{\texorpdfstring{\texttt{unitree\_sdk2\_cpp.robot.go2.JsonizeVec3.from\_json}}{unitree\_sdk2\_cpp.robot.go2.JsonizeVec3.from\_json}}\label{unitree-sdk2-cpp-robot-go2-jsonizevec3-from-json-1}}

计划从 JSON 风格字典读取字段并更新当前 SDK 值对象。

\textbf{签名}

\begin{lstlisting}[language=Python]
def from_json(self, json: dict[str, Any]) -> None
\end{lstlisting}

\textbf{可用性与安全性}

\textbf{\passthrough{\lstinline!SIGNATURE\_ONLY!} /
\passthrough{\lstinline!UNCLASSIFIED!}}：当前只有设计期签名，不能假设已存在于运行时扩展。

\textbf{参数}

\begin{longtable}[]{@{}
  >{\raggedright\arraybackslash}p{(\columnwidth - 6\tabcolsep) * \real{0.18}}
  >{\raggedright\arraybackslash}p{(\columnwidth - 6\tabcolsep) * \real{0.24}}
  >{\raggedright\arraybackslash}p{(\columnwidth - 6\tabcolsep) * \real{0.18}}
  >{\raggedright\arraybackslash}p{(\columnwidth - 6\tabcolsep) * \real{0.40}}@{}}
\toprule
名称 & Python 类型 & 默认值 & 含义 \\
\midrule
\endhead
\passthrough{\lstinline!json!} &
\passthrough{\lstinline!dict[str, Any]!} & 必填 & JSON
风格字典，键和值必须满足对应 C++ \passthrough{\lstinline!JsonMap!}
数据结构的约束。 对应 C++ 参数
\passthrough{\lstinline!json: common::JsonMap \&!}。 \\
\bottomrule
\end{longtable}

\textbf{返回值}

\begin{longtable}[]{@{}
  >{\raggedright\arraybackslash}p{(\columnwidth - 4\tabcolsep) * \real{0.18}}
  >{\raggedright\arraybackslash}p{(\columnwidth - 4\tabcolsep) * \real{0.20}}
  >{\raggedright\arraybackslash}p{(\columnwidth - 4\tabcolsep) * \real{0.62}}@{}}
\toprule
位置 & 类型 & 含义 \\
\midrule
\endhead
无 & \passthrough{\lstinline!None!} & 该方法不返回 Python
值；失败可能表现为异常或底层状态变化。 \\
\bottomrule
\end{longtable}

\textbf{对应 C++}

\begin{itemize}
\tightlist
\item
  类：\passthrough{\lstinline!unitree::robot::go2::JsonizeVec3!}
\item
  签名：\passthrough{\lstinline!fromJson(common::JsonMap \&)!}
\item
  绑定策略：\passthrough{\lstinline!SIGNATURE\_PREVIEW!}
\item
  声明位置：\passthrough{\lstinline!include/unitree/robot/go2/public/jsonize\_type.hpp:180!}
\end{itemize}

\textbf{说明}

\begin{itemize}
\tightlist
\item
  示例是局部调用片段；\passthrough{\lstinline!obj!}
  和参数变量表示已经按上层业务校验并准备好的对象。
\item
  该条目可以被编辑器和类型检查器识别，但当前运行时可能在属性查找阶段直接失败。
\end{itemize}

\textbf{用法}

\begin{lstlisting}[language=Python]
from typing import TYPE_CHECKING

if TYPE_CHECKING:
    def planned_from_json(obj: JsonizeVec3, json: dict[str, Any]) -> None:
        obj.from_json(json=json)
\end{lstlisting}

\begin{quote}
{[}!CAUTION{]} 上面的 \passthrough{\lstinline!TYPE\_CHECKING!}
示例只表达计划调用形态。该分支在普通 Python
运行时为假，不代表当前扩展能执行此接口。
\end{quote}

\hypertarget{unitree-sdk2-cpp-robot-go2-jsonizevec3-to-json-1}{%
\paragraph{\texorpdfstring{\texttt{unitree\_sdk2\_cpp.robot.go2.JsonizeVec3.to\_json}}{unitree\_sdk2\_cpp.robot.go2.JsonizeVec3.to\_json}}\label{unitree-sdk2-cpp-robot-go2-jsonizevec3-to-json-1}}

计划把当前 SDK 值对象写入 JSON 风格字典。

\textbf{签名}

\begin{lstlisting}[language=Python]
def to_json(self, json: dict[str, Any]) -> None
\end{lstlisting}

\textbf{可用性与安全性}

\textbf{\passthrough{\lstinline!SIGNATURE\_ONLY!} /
\passthrough{\lstinline!UNCLASSIFIED!}}：当前只有设计期签名，不能假设已存在于运行时扩展。

\textbf{参数}

\begin{longtable}[]{@{}
  >{\raggedright\arraybackslash}p{(\columnwidth - 6\tabcolsep) * \real{0.18}}
  >{\raggedright\arraybackslash}p{(\columnwidth - 6\tabcolsep) * \real{0.24}}
  >{\raggedright\arraybackslash}p{(\columnwidth - 6\tabcolsep) * \real{0.18}}
  >{\raggedright\arraybackslash}p{(\columnwidth - 6\tabcolsep) * \real{0.40}}@{}}
\toprule
名称 & Python 类型 & 默认值 & 含义 \\
\midrule
\endhead
\passthrough{\lstinline!json!} &
\passthrough{\lstinline!dict[str, Any]!} & 必填 & JSON
风格字典，键和值必须满足对应 C++ \passthrough{\lstinline!JsonMap!}
数据结构的约束。 对应 C++ 参数
\passthrough{\lstinline!json: common::JsonMap \&!}。 \\
\bottomrule
\end{longtable}

\textbf{返回值}

\begin{longtable}[]{@{}
  >{\raggedright\arraybackslash}p{(\columnwidth - 4\tabcolsep) * \real{0.18}}
  >{\raggedright\arraybackslash}p{(\columnwidth - 4\tabcolsep) * \real{0.20}}
  >{\raggedright\arraybackslash}p{(\columnwidth - 4\tabcolsep) * \real{0.62}}@{}}
\toprule
位置 & 类型 & 含义 \\
\midrule
\endhead
无 & \passthrough{\lstinline!None!} & 该方法不返回 Python
值；失败可能表现为异常或底层状态变化。 \\
\bottomrule
\end{longtable}

\textbf{对应 C++}

\begin{itemize}
\tightlist
\item
  类：\passthrough{\lstinline!unitree::robot::go2::JsonizeVec3!}
\item
  签名：\passthrough{\lstinline!toJson(common::JsonMap \&) const!}
\item
  绑定策略：\passthrough{\lstinline!SIGNATURE\_PREVIEW!}
\item
  声明位置：\passthrough{\lstinline!include/unitree/robot/go2/public/jsonize\_type.hpp:187!}
\end{itemize}

\textbf{说明}

\begin{itemize}
\tightlist
\item
  示例是局部调用片段；\passthrough{\lstinline!obj!}
  和参数变量表示已经按上层业务校验并准备好的对象。
\item
  该条目可以被编辑器和类型检查器识别，但当前运行时可能在属性查找阶段直接失败。
\end{itemize}

\textbf{用法}

\begin{lstlisting}[language=Python]
from typing import TYPE_CHECKING

if TYPE_CHECKING:
    def planned_to_json(obj: JsonizeVec3, json: dict[str, Any]) -> None:
        obj.to_json(json=json)
\end{lstlisting}

\begin{quote}
{[}!CAUTION{]} 上面的 \passthrough{\lstinline!TYPE\_CHECKING!}
示例只表达计划调用形态。该分支在普通 Python
运行时为假，不代表当前扩展能执行此接口。
\end{quote}

\hypertarget{unitree-sdk2-cpp-robot-go2-obstaclesavoidclient}{%
\subsubsection{\texorpdfstring{\texttt{unitree\_sdk2\_cpp.robot.go2.ObstaclesAvoidClient}}{unitree\_sdk2\_cpp.robot.go2.ObstaclesAvoidClient}}\label{unitree-sdk2-cpp-robot-go2-obstaclesavoidclient}}

Robot 服务客户端。构造、初始化和具体方法的可用性必须分别检查。

\textbf{导入}

\begin{lstlisting}[language=Python]
from unitree_sdk2_cpp.robot.go2 import ObstaclesAvoidClient
\end{lstlisting}

\textbf{基类}：\passthrough{\lstinline!Client!}

\textbf{方法索引}

\begin{longtable}[]{@{}
  >{\raggedright\arraybackslash}p{(\columnwidth - 6\tabcolsep) * \real{0.18}}
  >{\raggedright\arraybackslash}p{(\columnwidth - 6\tabcolsep) * \real{0.24}}
  >{\raggedright\arraybackslash}p{(\columnwidth - 6\tabcolsep) * \real{0.18}}
  >{\raggedright\arraybackslash}p{(\columnwidth - 6\tabcolsep) * \real{0.40}}@{}}
\toprule
方法 & Python 签名 & 状态 & 安全分类 \\
\midrule
\endhead
\protect\hyperlink{unitree-sdk2-cpp-robot-go2-obstaclesavoidclient-dunder-init-1}{\passthrough{\lstinline!\_\_init\_\_!}}
& \passthrough{\lstinline!def \_\_init\_\_(self) -> None!} &
\passthrough{\lstinline!SIGNATURE\_ONLY!} &
\passthrough{\lstinline!CONSTRUCTION!} \\
\protect\hyperlink{unitree-sdk2-cpp-robot-go2-obstaclesavoidclient-init-1}{\passthrough{\lstinline!init!}}
& \passthrough{\lstinline!def init(self) -> None!} &
\passthrough{\lstinline!SIGNATURE\_ONLY!} &
\passthrough{\lstinline!INITIALIZATION!} \\
\protect\hyperlink{unitree-sdk2-cpp-robot-go2-obstaclesavoidclient-switch-set-1}{\passthrough{\lstinline!switch\_set!}}
& \passthrough{\lstinline!def switch\_set(self, enable: bool) -> int!} &
\passthrough{\lstinline!SIGNATURE\_ONLY!} &
\passthrough{\lstinline!MOTION\_COMMAND!} \\
\protect\hyperlink{unitree-sdk2-cpp-robot-go2-obstaclesavoidclient-switch-get-1}{\passthrough{\lstinline!switch\_get!}}
& \passthrough{\lstinline!def switch\_get(self) -> tuple[int, bool]!} &
\passthrough{\lstinline!SIGNATURE\_ONLY!} &
\passthrough{\lstinline!MOTION\_COMMAND!} \\
\protect\hyperlink{unitree-sdk2-cpp-robot-go2-obstaclesavoidclient-move-1}{\passthrough{\lstinline!move!}}
&
\passthrough{\lstinline!def move(self, x: float, y: float, yaw: float) -> int!}
& \passthrough{\lstinline!SIGNATURE\_ONLY!} &
\passthrough{\lstinline!MOTION\_COMMAND!} \\
\protect\hyperlink{unitree-sdk2-cpp-robot-go2-obstaclesavoidclient-use-remote-command-from-api-1}{\passthrough{\lstinline!use\_remote\_command\_from\_api!}}
&
\passthrough{\lstinline!def use\_remote\_command\_from\_api(self, is\_remote\_commands\_from\_api: bool) -> int!}
& \passthrough{\lstinline!SIGNATURE\_ONLY!} &
\passthrough{\lstinline!MOTION\_COMMAND!} \\
\protect\hyperlink{unitree-sdk2-cpp-robot-go2-obstaclesavoidclient-move-to-absolute-position-1}{\passthrough{\lstinline!move\_to\_absolute\_position!}}
&
\passthrough{\lstinline!def move\_to\_absolute\_position(self, x: float, y: float, yaw: float) -> int!}
& \passthrough{\lstinline!SIGNATURE\_ONLY!} &
\passthrough{\lstinline!MOTION\_COMMAND!} \\
\protect\hyperlink{unitree-sdk2-cpp-robot-go2-obstaclesavoidclient-move-to-increment-position-1}{\passthrough{\lstinline!move\_to\_increment\_position!}}
&
\passthrough{\lstinline!def move\_to\_increment\_position(self, x: float, y: float, yaw: float) -> int!}
& \passthrough{\lstinline!SIGNATURE\_ONLY!} &
\passthrough{\lstinline!MOTION\_COMMAND!} \\
\bottomrule
\end{longtable}

\hypertarget{unitree-sdk2-cpp-robot-go2-obstaclesavoidclient-dunder-init-1}{%
\paragraph{\texorpdfstring{\texttt{unitree\_sdk2\_cpp.robot.go2.ObstaclesAvoidClient.\_\_init\_\_}}{unitree\_sdk2\_cpp.robot.go2.ObstaclesAvoidClient.\_\_init\_\_}}\label{unitree-sdk2-cpp-robot-go2-obstaclesavoidclient-dunder-init-1}}

初始化 \passthrough{\lstinline!ObstaclesAvoidClient!}
实例。是否能实际构造取决于下方可用性状态。

\textbf{签名}

\begin{lstlisting}[language=Python]
def __init__(self) -> None
\end{lstlisting}

\textbf{可用性与安全性}

\textbf{\passthrough{\lstinline!SIGNATURE\_ONLY!} /
\passthrough{\lstinline!CONSTRUCTION!}}：当前只有设计期签名，不能假设已存在于运行时扩展。

\textbf{参数}

无显式参数。实例方法中的 \passthrough{\lstinline!self!} 由 Python
自动传入。

\textbf{返回值}

\begin{longtable}[]{@{}
  >{\raggedright\arraybackslash}p{(\columnwidth - 4\tabcolsep) * \real{0.18}}
  >{\raggedright\arraybackslash}p{(\columnwidth - 4\tabcolsep) * \real{0.20}}
  >{\raggedright\arraybackslash}p{(\columnwidth - 4\tabcolsep) * \real{0.62}}@{}}
\toprule
位置 & 类型 & 含义 \\
\midrule
\endhead
无 & \passthrough{\lstinline!None!} &
\passthrough{\lstinline!\_\_init\_\_!}
本身不返回值；调用类对象时，在构造函数可用的前提下得到该类实例。 \\
\bottomrule
\end{longtable}

\textbf{对应 C++}

\begin{itemize}
\tightlist
\item
  类：\passthrough{\lstinline!unitree::robot::go2::ObstaclesAvoidClient!}
\item
  签名：\passthrough{\lstinline!ObstaclesAvoidClient()!}
\item
  绑定策略：\passthrough{\lstinline!未单独标注!}
\item
  声明位置：\passthrough{\lstinline!include/unitree/robot/go2/obstacles\_avoid/obstacles\_avoid\_client.hpp:15!}
\end{itemize}

\textbf{说明}

\begin{itemize}
\tightlist
\item
  示例是局部调用片段；\passthrough{\lstinline!obj!}
  和参数变量表示已经按上层业务校验并准备好的对象。
\item
  该条目可以被编辑器和类型检查器识别，但当前运行时可能在属性查找阶段直接失败。
\end{itemize}

\textbf{用法}

\begin{lstlisting}[language=Python]
from typing import TYPE_CHECKING

if TYPE_CHECKING:
    def planned_construct() -> ObstaclesAvoidClient:
        return ObstaclesAvoidClient()
\end{lstlisting}

\begin{quote}
{[}!CAUTION{]} 上面的 \passthrough{\lstinline!TYPE\_CHECKING!}
示例只表达计划调用形态。该分支在普通 Python
运行时为假，不代表当前扩展能执行此接口。
\end{quote}

\hypertarget{unitree-sdk2-cpp-robot-go2-obstaclesavoidclient-init-1}{%
\paragraph{\texorpdfstring{\texttt{unitree\_sdk2\_cpp.robot.go2.ObstaclesAvoidClient.init}}{unitree\_sdk2\_cpp.robot.go2.ObstaclesAvoidClient.init}}\label{unitree-sdk2-cpp-robot-go2-obstaclesavoidclient-init-1}}

初始化当前 SDK 对象所需的底层通道或服务资源。

\textbf{签名}

\begin{lstlisting}[language=Python]
def init(self) -> None
\end{lstlisting}

\textbf{可用性与安全性}

\textbf{\passthrough{\lstinline!SIGNATURE\_ONLY!} /
\passthrough{\lstinline!INITIALIZATION!}}：当前只有设计期签名，不能假设已存在于运行时扩展。

\textbf{参数}

无显式参数。实例方法中的 \passthrough{\lstinline!self!} 由 Python
自动传入。

\textbf{返回值}

\begin{longtable}[]{@{}
  >{\raggedright\arraybackslash}p{(\columnwidth - 4\tabcolsep) * \real{0.18}}
  >{\raggedright\arraybackslash}p{(\columnwidth - 4\tabcolsep) * \real{0.20}}
  >{\raggedright\arraybackslash}p{(\columnwidth - 4\tabcolsep) * \real{0.62}}@{}}
\toprule
位置 & 类型 & 含义 \\
\midrule
\endhead
无 & \passthrough{\lstinline!None!} & 该方法不返回 Python
值；失败可能表现为异常或底层状态变化。 \\
\bottomrule
\end{longtable}

\textbf{对应 C++}

\begin{itemize}
\tightlist
\item
  类：\passthrough{\lstinline!unitree::robot::go2::ObstaclesAvoidClient!}
\item
  签名：\passthrough{\lstinline!Init()!}
\item
  绑定策略：\passthrough{\lstinline!DIRECT!}
\item
  声明位置：\passthrough{\lstinline!include/unitree/robot/go2/obstacles\_avoid/obstacles\_avoid\_client.hpp:18!}
\end{itemize}

\textbf{说明}

\begin{itemize}
\tightlist
\item
  示例是局部调用片段；\passthrough{\lstinline!obj!}
  和参数变量表示已经按上层业务校验并准备好的对象。
\item
  该条目可以被编辑器和类型检查器识别，但当前运行时可能在属性查找阶段直接失败。
\end{itemize}

\textbf{用法}

\begin{lstlisting}[language=Python]
from typing import TYPE_CHECKING

if TYPE_CHECKING:
    def planned_init(obj: ObstaclesAvoidClient) -> None:
        obj.init()
\end{lstlisting}

\begin{quote}
{[}!CAUTION{]} 上面的 \passthrough{\lstinline!TYPE\_CHECKING!}
示例只表达计划调用形态。该分支在普通 Python
运行时为假，不代表当前扩展能执行此接口。
\end{quote}

\hypertarget{unitree-sdk2-cpp-robot-go2-obstaclesavoidclient-switch-set-1}{%
\paragraph{\texorpdfstring{\texttt{unitree\_sdk2\_cpp.robot.go2.ObstaclesAvoidClient.switch\_set}}{unitree\_sdk2\_cpp.robot.go2.ObstaclesAvoidClient.switch\_set}}\label{unitree-sdk2-cpp-robot-go2-obstaclesavoidclient-switch-set-1}}

对应 C++ SDK 操作
\passthrough{\lstinline!SwitchSet(bool)!}。上游头文件没有可直接生成的业务说明时，本参考只保证签名映射，精确语义需查目标型号协议。

\textbf{签名}

\begin{lstlisting}[language=Python]
def switch_set(self, enable: bool) -> int
\end{lstlisting}

\textbf{可用性与安全性}

\textbf{\passthrough{\lstinline!SIGNATURE\_ONLY!} /
\passthrough{\lstinline!MOTION\_COMMAND!}}：仅用于补全和静态检查。当前二进制没有该
Python 方法；不得按可执行运动接口使用。

\textbf{参数}

\begin{longtable}[]{@{}
  >{\raggedright\arraybackslash}p{(\columnwidth - 6\tabcolsep) * \real{0.18}}
  >{\raggedright\arraybackslash}p{(\columnwidth - 6\tabcolsep) * \real{0.24}}
  >{\raggedright\arraybackslash}p{(\columnwidth - 6\tabcolsep) * \real{0.18}}
  >{\raggedright\arraybackslash}p{(\columnwidth - 6\tabcolsep) * \real{0.40}}@{}}
\toprule
名称 & Python 类型 & 默认值 & 含义 \\
\midrule
\endhead
\passthrough{\lstinline!enable!} & \passthrough{\lstinline!bool!} & 必填
& 传给该接口的 \passthrough{\lstinline!enable!} 参数，Python 类型为
\passthrough{\lstinline!bool!}。精确含义、单位、范围和枚举值需查目标型号对应头文件与协议。
对应 C++ 参数 \passthrough{\lstinline!enable: bool!}。
只接受布尔语义。 \\
\bottomrule
\end{longtable}

\textbf{返回值}

\begin{longtable}[]{@{}
  >{\raggedright\arraybackslash}p{(\columnwidth - 4\tabcolsep) * \real{0.18}}
  >{\raggedright\arraybackslash}p{(\columnwidth - 4\tabcolsep) * \real{0.20}}
  >{\raggedright\arraybackslash}p{(\columnwidth - 4\tabcolsep) * \real{0.62}}@{}}
\toprule
位置 & 类型 & 含义 \\
\midrule
\endhead
返回值 & \passthrough{\lstinline!int!} & SDK 状态码；Unitree 示例通常以
\passthrough{\lstinline!0!} 表示成功，非零值需按具体服务定义解释。 \\
\bottomrule
\end{longtable}

\textbf{对应 C++}

\begin{itemize}
\tightlist
\item
  类：\passthrough{\lstinline!unitree::robot::go2::ObstaclesAvoidClient!}
\item
  签名：\passthrough{\lstinline!SwitchSet(bool)!}
\item
  绑定策略：\passthrough{\lstinline!DIRECT!}
\item
  声明位置：\passthrough{\lstinline!include/unitree/robot/go2/obstacles\_avoid/obstacles\_avoid\_client.hpp:20!}
\end{itemize}

\textbf{说明}

\begin{itemize}
\tightlist
\item
  示例是局部调用片段；\passthrough{\lstinline!obj!}
  和参数变量表示已经按上层业务校验并准备好的对象。
\item
  该条目可以被编辑器和类型检查器识别，但当前运行时可能在属性查找阶段直接失败。
\item
  该方法属于运动命令。方法名包含
  \passthrough{\lstinline!stop!}、\passthrough{\lstinline!damp!} 或
  \passthrough{\lstinline!zero!} 也不代表它是独立物理急停。
\end{itemize}

\textbf{用法}

\begin{lstlisting}[language=Python]
from typing import TYPE_CHECKING

if TYPE_CHECKING:
    def planned_switch_set(obj: ObstaclesAvoidClient, enable: bool) -> int:
        return obj.switch_set(enable=enable)
\end{lstlisting}

\begin{quote}
{[}!CAUTION{]} 上面的 \passthrough{\lstinline!TYPE\_CHECKING!}
示例只表达计划调用形态。该分支在普通 Python
运行时为假，不代表当前扩展能执行此接口。
\end{quote}

\hypertarget{unitree-sdk2-cpp-robot-go2-obstaclesavoidclient-switch-get-1}{%
\paragraph{\texorpdfstring{\texttt{unitree\_sdk2\_cpp.robot.go2.ObstaclesAvoidClient.switch\_get}}{unitree\_sdk2\_cpp.robot.go2.ObstaclesAvoidClient.switch\_get}}\label{unitree-sdk2-cpp-robot-go2-obstaclesavoidclient-switch-get-1}}

对应 C++ SDK 操作
\passthrough{\lstinline!SwitchGet(bool \&)!}。上游头文件没有可直接生成的业务说明时，本参考只保证签名映射，精确语义需查目标型号协议。

\textbf{签名}

\begin{lstlisting}[language=Python]
def switch_get(self) -> tuple[int, bool]
\end{lstlisting}

\textbf{可用性与安全性}

\textbf{\passthrough{\lstinline!SIGNATURE\_ONLY!} /
\passthrough{\lstinline!MOTION\_COMMAND!}}：仅用于补全和静态检查。当前二进制没有该
Python 方法；不得按可执行运动接口使用。

\textbf{参数}

无显式参数。实例方法中的 \passthrough{\lstinline!self!} 由 Python
自动传入。

\textbf{返回值}

\begin{longtable}[]{@{}
  >{\raggedright\arraybackslash}p{(\columnwidth - 4\tabcolsep) * \real{0.18}}
  >{\raggedright\arraybackslash}p{(\columnwidth - 4\tabcolsep) * \real{0.20}}
  >{\raggedright\arraybackslash}p{(\columnwidth - 4\tabcolsep) * \real{0.62}}@{}}
\toprule
位置 & 类型 & 含义 \\
\midrule
\endhead
{[}0{]} \passthrough{\lstinline!status!} & \passthrough{\lstinline!int!}
& SDK 状态码；Unitree 示例通常以 \passthrough{\lstinline!0!}
表示成功，非零值需按具体服务错误码解释。 \\
{[}1{]} \passthrough{\lstinline!enable!} &
\passthrough{\lstinline!bool!} & 传给该接口的
\passthrough{\lstinline!enable!} 参数，Python 类型为
\passthrough{\lstinline!bool!}。精确含义、单位、范围和枚举值需查目标型号对应头文件与协议。
对应 C++ 参数 \passthrough{\lstinline!enable: bool \&!}。
只接受布尔语义。 \\
\bottomrule
\end{longtable}

\textbf{对应 C++}

\begin{itemize}
\tightlist
\item
  类：\passthrough{\lstinline!unitree::robot::go2::ObstaclesAvoidClient!}
\item
  签名：\passthrough{\lstinline!SwitchGet(bool \&)!}
\item
  绑定策略：\passthrough{\lstinline!OUTPUT\_WRAPPER!}
\item
  声明位置：\passthrough{\lstinline!include/unitree/robot/go2/obstacles\_avoid/obstacles\_avoid\_client.hpp:21!}
\end{itemize}

\textbf{说明}

\begin{itemize}
\tightlist
\item
  示例是局部调用片段；\passthrough{\lstinline!obj!}
  和参数变量表示已经按上层业务校验并准备好的对象。
\item
  C++ 的可修改输出引用不会作为 Python 入参出现，而是按 C++
  参数顺序追加到返回元组中。
\item
  该条目可以被编辑器和类型检查器识别，但当前运行时可能在属性查找阶段直接失败。
\item
  该方法属于运动命令。方法名包含
  \passthrough{\lstinline!stop!}、\passthrough{\lstinline!damp!} 或
  \passthrough{\lstinline!zero!} 也不代表它是独立物理急停。
\end{itemize}

\textbf{用法}

\begin{lstlisting}[language=Python]
from typing import TYPE_CHECKING

if TYPE_CHECKING:
    def planned_switch_get(obj: ObstaclesAvoidClient) -> tuple[int, bool]:
        return obj.switch_get()
\end{lstlisting}

\begin{quote}
{[}!CAUTION{]} 上面的 \passthrough{\lstinline!TYPE\_CHECKING!}
示例只表达计划调用形态。该分支在普通 Python
运行时为假，不代表当前扩展能执行此接口。
\end{quote}

\hypertarget{unitree-sdk2-cpp-robot-go2-obstaclesavoidclient-move-1}{%
\paragraph{\texorpdfstring{\texttt{unitree\_sdk2\_cpp.robot.go2.ObstaclesAvoidClient.move}}{unitree\_sdk2\_cpp.robot.go2.ObstaclesAvoidClient.move}}\label{unitree-sdk2-cpp-robot-go2-obstaclesavoidclient-move-1}}

对应 C++ SDK 操作
\passthrough{\lstinline!Move(float, float, float)!}。上游头文件没有可直接生成的业务说明时，本参考只保证签名映射，精确语义需查目标型号协议。

\textbf{签名}

\begin{lstlisting}[language=Python]
def move(self, x: float, y: float, yaw: float) -> int
\end{lstlisting}

\textbf{可用性与安全性}

\textbf{\passthrough{\lstinline!SIGNATURE\_ONLY!} /
\passthrough{\lstinline!MOTION\_COMMAND!}}：仅用于补全和静态检查。当前二进制没有该
Python 方法；不得按可执行运动接口使用。

\textbf{参数}

\begin{longtable}[]{@{}
  >{\raggedright\arraybackslash}p{(\columnwidth - 6\tabcolsep) * \real{0.18}}
  >{\raggedright\arraybackslash}p{(\columnwidth - 6\tabcolsep) * \real{0.24}}
  >{\raggedright\arraybackslash}p{(\columnwidth - 6\tabcolsep) * \real{0.18}}
  >{\raggedright\arraybackslash}p{(\columnwidth - 6\tabcolsep) * \real{0.40}}@{}}
\toprule
名称 & Python 类型 & 默认值 & 含义 \\
\midrule
\endhead
\passthrough{\lstinline!x!} & \passthrough{\lstinline!float!} & 必填 & X
方向位置或分量参数。参考系、单位和范围必须查当前方法及目标型号协议。
对应 C++ 参数 \passthrough{\lstinline!x: float!}。
底层为浮点数；类型本身不说明单位、坐标系或业务安全范围。 \\
\passthrough{\lstinline!y!} & \passthrough{\lstinline!float!} & 必填 & Y
方向位置或分量参数。参考系、单位和范围必须查当前方法及目标型号协议。
对应 C++ 参数 \passthrough{\lstinline!y: float!}。
底层为浮点数；类型本身不说明单位、坐标系或业务安全范围。 \\
\passthrough{\lstinline!yaw!} & \passthrough{\lstinline!float!} & 必填 &
偏航角或偏航目标参数。参考系、单位和范围必须查目标型号协议。 对应 C++
参数 \passthrough{\lstinline!yaw: float!}。
底层为浮点数；类型本身不说明单位、坐标系或业务安全范围。 \\
\bottomrule
\end{longtable}

\textbf{返回值}

\begin{longtable}[]{@{}
  >{\raggedright\arraybackslash}p{(\columnwidth - 4\tabcolsep) * \real{0.18}}
  >{\raggedright\arraybackslash}p{(\columnwidth - 4\tabcolsep) * \real{0.20}}
  >{\raggedright\arraybackslash}p{(\columnwidth - 4\tabcolsep) * \real{0.62}}@{}}
\toprule
位置 & 类型 & 含义 \\
\midrule
\endhead
返回值 & \passthrough{\lstinline!int!} & SDK 状态码；Unitree 示例通常以
\passthrough{\lstinline!0!} 表示成功，非零值需按具体服务定义解释。 \\
\bottomrule
\end{longtable}

\textbf{对应 C++}

\begin{itemize}
\tightlist
\item
  类：\passthrough{\lstinline!unitree::robot::go2::ObstaclesAvoidClient!}
\item
  签名：\passthrough{\lstinline!Move(float, float, float)!}
\item
  绑定策略：\passthrough{\lstinline!DIRECT!}
\item
  声明位置：\passthrough{\lstinline!include/unitree/robot/go2/obstacles\_avoid/obstacles\_avoid\_client.hpp:23!}
\end{itemize}

\textbf{说明}

\begin{itemize}
\tightlist
\item
  示例是局部调用片段；\passthrough{\lstinline!obj!}
  和参数变量表示已经按上层业务校验并准备好的对象。
\item
  该条目可以被编辑器和类型检查器识别，但当前运行时可能在属性查找阶段直接失败。
\item
  该方法属于运动命令。方法名包含
  \passthrough{\lstinline!stop!}、\passthrough{\lstinline!damp!} 或
  \passthrough{\lstinline!zero!} 也不代表它是独立物理急停。
\end{itemize}

\textbf{用法}

\begin{lstlisting}[language=Python]
from typing import TYPE_CHECKING

if TYPE_CHECKING:
    def planned_move(obj: ObstaclesAvoidClient, x: float, y: float, yaw: float) -> int:
        return obj.move(x=x, y=y, yaw=yaw)
\end{lstlisting}

\begin{quote}
{[}!CAUTION{]} 上面的 \passthrough{\lstinline!TYPE\_CHECKING!}
示例只表达计划调用形态。该分支在普通 Python
运行时为假，不代表当前扩展能执行此接口。
\end{quote}

\hypertarget{unitree-sdk2-cpp-robot-go2-obstaclesavoidclient-use-remote-command-from-api-1}{%
\paragraph{\texorpdfstring{\texttt{unitree\_sdk2\_cpp.robot.go2.ObstaclesAvoidClient.use\_remote\_command\_from\_api}}{unitree\_sdk2\_cpp.robot.go2.ObstaclesAvoidClient.use\_remote\_command\_from\_api}}\label{unitree-sdk2-cpp-robot-go2-obstaclesavoidclient-use-remote-command-from-api-1}}

对应 C++ SDK 操作
\passthrough{\lstinline!UseRemoteCommandFromApi(bool)!}。上游头文件没有可直接生成的业务说明时，本参考只保证签名映射，精确语义需查目标型号协议。

\textbf{签名}

\begin{lstlisting}[language=Python]
def use_remote_command_from_api(self, is_remote_commands_from_api: bool) -> int
\end{lstlisting}

\textbf{可用性与安全性}

\textbf{\passthrough{\lstinline!SIGNATURE\_ONLY!} /
\passthrough{\lstinline!MOTION\_COMMAND!}}：仅用于补全和静态检查。当前二进制没有该
Python 方法；不得按可执行运动接口使用。

\textbf{参数}

\begin{longtable}[]{@{}
  >{\raggedright\arraybackslash}p{(\columnwidth - 6\tabcolsep) * \real{0.18}}
  >{\raggedright\arraybackslash}p{(\columnwidth - 6\tabcolsep) * \real{0.24}}
  >{\raggedright\arraybackslash}p{(\columnwidth - 6\tabcolsep) * \real{0.18}}
  >{\raggedright\arraybackslash}p{(\columnwidth - 6\tabcolsep) * \real{0.40}}@{}}
\toprule
名称 & Python 类型 & 默认值 & 含义 \\
\midrule
\endhead
\passthrough{\lstinline!is\_remote\_commands\_from\_api!} &
\passthrough{\lstinline!bool!} & 必填 & 传给该接口的
\passthrough{\lstinline!is!} \passthrough{\lstinline!remote!}
\passthrough{\lstinline!commands!} \passthrough{\lstinline!from!} API
参数，Python 类型为
\passthrough{\lstinline!bool!}。精确含义、单位、范围和枚举值需查目标型号对应头文件与协议。
对应 C++ 参数 \passthrough{\lstinline!isRemoteCommandsFromApi: bool!}。
只接受布尔语义。 \\
\bottomrule
\end{longtable}

\textbf{返回值}

\begin{longtable}[]{@{}
  >{\raggedright\arraybackslash}p{(\columnwidth - 4\tabcolsep) * \real{0.18}}
  >{\raggedright\arraybackslash}p{(\columnwidth - 4\tabcolsep) * \real{0.20}}
  >{\raggedright\arraybackslash}p{(\columnwidth - 4\tabcolsep) * \real{0.62}}@{}}
\toprule
位置 & 类型 & 含义 \\
\midrule
\endhead
返回值 & \passthrough{\lstinline!int!} & SDK 状态码；Unitree 示例通常以
\passthrough{\lstinline!0!} 表示成功，非零值需按具体服务定义解释。 \\
\bottomrule
\end{longtable}

\textbf{对应 C++}

\begin{itemize}
\tightlist
\item
  类：\passthrough{\lstinline!unitree::robot::go2::ObstaclesAvoidClient!}
\item
  签名：\passthrough{\lstinline!UseRemoteCommandFromApi(bool)!}
\item
  绑定策略：\passthrough{\lstinline!DIRECT!}
\item
  声明位置：\passthrough{\lstinline!include/unitree/robot/go2/obstacles\_avoid/obstacles\_avoid\_client.hpp:24!}
\end{itemize}

\textbf{说明}

\begin{itemize}
\tightlist
\item
  示例是局部调用片段；\passthrough{\lstinline!obj!}
  和参数变量表示已经按上层业务校验并准备好的对象。
\item
  该条目可以被编辑器和类型检查器识别，但当前运行时可能在属性查找阶段直接失败。
\item
  该方法属于运动命令。方法名包含
  \passthrough{\lstinline!stop!}、\passthrough{\lstinline!damp!} 或
  \passthrough{\lstinline!zero!} 也不代表它是独立物理急停。
\end{itemize}

\textbf{用法}

\begin{lstlisting}[language=Python]
from typing import TYPE_CHECKING

if TYPE_CHECKING:
    def planned_use_remote_command_from_api(obj: ObstaclesAvoidClient, is_remote_commands_from_api: bool) -> int:
        return obj.use_remote_command_from_api(is_remote_commands_from_api=is_remote_commands_from_api)
\end{lstlisting}

\begin{quote}
{[}!CAUTION{]} 上面的 \passthrough{\lstinline!TYPE\_CHECKING!}
示例只表达计划调用形态。该分支在普通 Python
运行时为假，不代表当前扩展能执行此接口。
\end{quote}

\hypertarget{unitree-sdk2-cpp-robot-go2-obstaclesavoidclient-move-to-absolute-position-1}{%
\paragraph{\texorpdfstring{\texttt{unitree\_sdk2\_cpp.robot.go2.ObstaclesAvoidClient.move\_to\_absolute\_position}}{unitree\_sdk2\_cpp.robot.go2.ObstaclesAvoidClient.move\_to\_absolute\_position}}\label{unitree-sdk2-cpp-robot-go2-obstaclesavoidclient-move-to-absolute-position-1}}

对应 C++ SDK 操作
\passthrough{\lstinline!MoveToAbsolutePosition(float, float, float)!}。上游头文件没有可直接生成的业务说明时，本参考只保证签名映射，精确语义需查目标型号协议。

\textbf{签名}

\begin{lstlisting}[language=Python]
def move_to_absolute_position(self, x: float, y: float, yaw: float) -> int
\end{lstlisting}

\textbf{可用性与安全性}

\textbf{\passthrough{\lstinline!SIGNATURE\_ONLY!} /
\passthrough{\lstinline!MOTION\_COMMAND!}}：仅用于补全和静态检查。当前二进制没有该
Python 方法；不得按可执行运动接口使用。

\textbf{参数}

\begin{longtable}[]{@{}
  >{\raggedright\arraybackslash}p{(\columnwidth - 6\tabcolsep) * \real{0.18}}
  >{\raggedright\arraybackslash}p{(\columnwidth - 6\tabcolsep) * \real{0.24}}
  >{\raggedright\arraybackslash}p{(\columnwidth - 6\tabcolsep) * \real{0.18}}
  >{\raggedright\arraybackslash}p{(\columnwidth - 6\tabcolsep) * \real{0.40}}@{}}
\toprule
名称 & Python 类型 & 默认值 & 含义 \\
\midrule
\endhead
\passthrough{\lstinline!x!} & \passthrough{\lstinline!float!} & 必填 & X
方向位置或分量参数。参考系、单位和范围必须查当前方法及目标型号协议。
对应 C++ 参数 \passthrough{\lstinline!x: float!}。
底层为浮点数；类型本身不说明单位、坐标系或业务安全范围。 \\
\passthrough{\lstinline!y!} & \passthrough{\lstinline!float!} & 必填 & Y
方向位置或分量参数。参考系、单位和范围必须查当前方法及目标型号协议。
对应 C++ 参数 \passthrough{\lstinline!y: float!}。
底层为浮点数；类型本身不说明单位、坐标系或业务安全范围。 \\
\passthrough{\lstinline!yaw!} & \passthrough{\lstinline!float!} & 必填 &
偏航角或偏航目标参数。参考系、单位和范围必须查目标型号协议。 对应 C++
参数 \passthrough{\lstinline!yaw: float!}。
底层为浮点数；类型本身不说明单位、坐标系或业务安全范围。 \\
\bottomrule
\end{longtable}

\textbf{返回值}

\begin{longtable}[]{@{}
  >{\raggedright\arraybackslash}p{(\columnwidth - 4\tabcolsep) * \real{0.18}}
  >{\raggedright\arraybackslash}p{(\columnwidth - 4\tabcolsep) * \real{0.20}}
  >{\raggedright\arraybackslash}p{(\columnwidth - 4\tabcolsep) * \real{0.62}}@{}}
\toprule
位置 & 类型 & 含义 \\
\midrule
\endhead
返回值 & \passthrough{\lstinline!int!} & SDK 状态码；Unitree 示例通常以
\passthrough{\lstinline!0!} 表示成功，非零值需按具体服务定义解释。 \\
\bottomrule
\end{longtable}

\textbf{对应 C++}

\begin{itemize}
\tightlist
\item
  类：\passthrough{\lstinline!unitree::robot::go2::ObstaclesAvoidClient!}
\item
  签名：\passthrough{\lstinline!MoveToAbsolutePosition(float, float, float)!}
\item
  绑定策略：\passthrough{\lstinline!DIRECT!}
\item
  声明位置：\passthrough{\lstinline!include/unitree/robot/go2/obstacles\_avoid/obstacles\_avoid\_client.hpp:26!}
\end{itemize}

\textbf{说明}

\begin{itemize}
\tightlist
\item
  示例是局部调用片段；\passthrough{\lstinline!obj!}
  和参数变量表示已经按上层业务校验并准备好的对象。
\item
  该条目可以被编辑器和类型检查器识别，但当前运行时可能在属性查找阶段直接失败。
\item
  该方法属于运动命令。方法名包含
  \passthrough{\lstinline!stop!}、\passthrough{\lstinline!damp!} 或
  \passthrough{\lstinline!zero!} 也不代表它是独立物理急停。
\end{itemize}

\textbf{用法}

\begin{lstlisting}[language=Python]
from typing import TYPE_CHECKING

if TYPE_CHECKING:
    def planned_move_to_absolute_position(obj: ObstaclesAvoidClient, x: float, y: float, yaw: float) -> int:
        return obj.move_to_absolute_position(x=x, y=y, yaw=yaw)
\end{lstlisting}

\begin{quote}
{[}!CAUTION{]} 上面的 \passthrough{\lstinline!TYPE\_CHECKING!}
示例只表达计划调用形态。该分支在普通 Python
运行时为假，不代表当前扩展能执行此接口。
\end{quote}

\hypertarget{unitree-sdk2-cpp-robot-go2-obstaclesavoidclient-move-to-increment-position-1}{%
\paragraph{\texorpdfstring{\texttt{unitree\_sdk2\_cpp.robot.go2.ObstaclesAvoidClient.move\_to\_increment\_position}}{unitree\_sdk2\_cpp.robot.go2.ObstaclesAvoidClient.move\_to\_increment\_position}}\label{unitree-sdk2-cpp-robot-go2-obstaclesavoidclient-move-to-increment-position-1}}

对应 C++ SDK 操作
\passthrough{\lstinline!MoveToIncrementPosition(float, float, float)!}。上游头文件没有可直接生成的业务说明时，本参考只保证签名映射，精确语义需查目标型号协议。

\textbf{签名}

\begin{lstlisting}[language=Python]
def move_to_increment_position(self, x: float, y: float, yaw: float) -> int
\end{lstlisting}

\textbf{可用性与安全性}

\textbf{\passthrough{\lstinline!SIGNATURE\_ONLY!} /
\passthrough{\lstinline!MOTION\_COMMAND!}}：仅用于补全和静态检查。当前二进制没有该
Python 方法；不得按可执行运动接口使用。

\textbf{参数}

\begin{longtable}[]{@{}
  >{\raggedright\arraybackslash}p{(\columnwidth - 6\tabcolsep) * \real{0.18}}
  >{\raggedright\arraybackslash}p{(\columnwidth - 6\tabcolsep) * \real{0.24}}
  >{\raggedright\arraybackslash}p{(\columnwidth - 6\tabcolsep) * \real{0.18}}
  >{\raggedright\arraybackslash}p{(\columnwidth - 6\tabcolsep) * \real{0.40}}@{}}
\toprule
名称 & Python 类型 & 默认值 & 含义 \\
\midrule
\endhead
\passthrough{\lstinline!x!} & \passthrough{\lstinline!float!} & 必填 & X
方向位置或分量参数。参考系、单位和范围必须查当前方法及目标型号协议。
对应 C++ 参数 \passthrough{\lstinline!x: float!}。
底层为浮点数；类型本身不说明单位、坐标系或业务安全范围。 \\
\passthrough{\lstinline!y!} & \passthrough{\lstinline!float!} & 必填 & Y
方向位置或分量参数。参考系、单位和范围必须查当前方法及目标型号协议。
对应 C++ 参数 \passthrough{\lstinline!y: float!}。
底层为浮点数；类型本身不说明单位、坐标系或业务安全范围。 \\
\passthrough{\lstinline!yaw!} & \passthrough{\lstinline!float!} & 必填 &
偏航角或偏航目标参数。参考系、单位和范围必须查目标型号协议。 对应 C++
参数 \passthrough{\lstinline!yaw: float!}。
底层为浮点数；类型本身不说明单位、坐标系或业务安全范围。 \\
\bottomrule
\end{longtable}

\textbf{返回值}

\begin{longtable}[]{@{}
  >{\raggedright\arraybackslash}p{(\columnwidth - 4\tabcolsep) * \real{0.18}}
  >{\raggedright\arraybackslash}p{(\columnwidth - 4\tabcolsep) * \real{0.20}}
  >{\raggedright\arraybackslash}p{(\columnwidth - 4\tabcolsep) * \real{0.62}}@{}}
\toprule
位置 & 类型 & 含义 \\
\midrule
\endhead
返回值 & \passthrough{\lstinline!int!} & SDK 状态码；Unitree 示例通常以
\passthrough{\lstinline!0!} 表示成功，非零值需按具体服务定义解释。 \\
\bottomrule
\end{longtable}

\textbf{对应 C++}

\begin{itemize}
\tightlist
\item
  类：\passthrough{\lstinline!unitree::robot::go2::ObstaclesAvoidClient!}
\item
  签名：\passthrough{\lstinline!MoveToIncrementPosition(float, float, float)!}
\item
  绑定策略：\passthrough{\lstinline!DIRECT!}
\item
  声明位置：\passthrough{\lstinline!include/unitree/robot/go2/obstacles\_avoid/obstacles\_avoid\_client.hpp:27!}
\end{itemize}

\textbf{说明}

\begin{itemize}
\tightlist
\item
  示例是局部调用片段；\passthrough{\lstinline!obj!}
  和参数变量表示已经按上层业务校验并准备好的对象。
\item
  该条目可以被编辑器和类型检查器识别，但当前运行时可能在属性查找阶段直接失败。
\item
  该方法属于运动命令。方法名包含
  \passthrough{\lstinline!stop!}、\passthrough{\lstinline!damp!} 或
  \passthrough{\lstinline!zero!} 也不代表它是独立物理急停。
\end{itemize}

\textbf{用法}

\begin{lstlisting}[language=Python]
from typing import TYPE_CHECKING

if TYPE_CHECKING:
    def planned_move_to_increment_position(obj: ObstaclesAvoidClient, x: float, y: float, yaw: float) -> int:
        return obj.move_to_increment_position(x=x, y=y, yaw=yaw)
\end{lstlisting}

\begin{quote}
{[}!CAUTION{]} 上面的 \passthrough{\lstinline!TYPE\_CHECKING!}
示例只表达计划调用形态。该分支在普通 Python
运行时为假，不代表当前扩展能执行此接口。
\end{quote}

\hypertarget{unitree-sdk2-cpp-robot-go2-obstaclesavoidmoveparameter}{%
\subsubsection{\texorpdfstring{\texttt{unitree\_sdk2\_cpp.robot.go2.ObstaclesAvoidMoveParameter}}{unitree\_sdk2\_cpp.robot.go2.ObstaclesAvoidMoveParameter}}\label{unitree-sdk2-cpp-robot-go2-obstaclesavoidmoveparameter}}

SDK 请求参数值类型；当前可能仅提供设计期签名。

\textbf{导入}

\begin{lstlisting}[language=Python]
from unitree_sdk2_cpp.robot.go2 import ObstaclesAvoidMoveParameter
\end{lstlisting}

\textbf{基类}：\passthrough{\lstinline!object!}

\textbf{公开属性}

\begin{longtable}[]{@{}
  >{\raggedright\arraybackslash}p{(\columnwidth - 6\tabcolsep) * \real{0.18}}
  >{\raggedright\arraybackslash}p{(\columnwidth - 6\tabcolsep) * \real{0.24}}
  >{\raggedright\arraybackslash}p{(\columnwidth - 6\tabcolsep) * \real{0.18}}
  >{\raggedright\arraybackslash}p{(\columnwidth - 6\tabcolsep) * \real{0.40}}@{}}
\toprule
名称 & Python 类型 & 含义 & 用法 \\
\midrule
\endhead
\passthrough{\lstinline!m\_x!} & \passthrough{\lstinline!float!} &
传给该接口的 \passthrough{\lstinline!m!} \passthrough{\lstinline!x!}
参数，Python 类型为
\passthrough{\lstinline!float!}。精确含义、单位、范围和枚举值需查目标型号对应头文件与协议。
& \passthrough{\lstinline!value = obj.m\_x!} /
\passthrough{\lstinline!obj.m\_x = value!} \\
\passthrough{\lstinline!m\_y!} & \passthrough{\lstinline!float!} &
传给该接口的 \passthrough{\lstinline!m!} \passthrough{\lstinline!y!}
参数，Python 类型为
\passthrough{\lstinline!float!}。精确含义、单位、范围和枚举值需查目标型号对应头文件与协议。
& \passthrough{\lstinline!value = obj.m\_y!} /
\passthrough{\lstinline!obj.m\_y = value!} \\
\passthrough{\lstinline!m\_yaw!} & \passthrough{\lstinline!float!} &
传给该接口的 \passthrough{\lstinline!m!} \passthrough{\lstinline!yaw!}
参数，Python 类型为
\passthrough{\lstinline!float!}。精确含义、单位、范围和枚举值需查目标型号对应头文件与协议。
& \passthrough{\lstinline!value = obj.m\_yaw!} /
\passthrough{\lstinline!obj.m\_yaw = value!} \\
\passthrough{\lstinline!m\_mode!} & \passthrough{\lstinline!int!} &
传给该接口的 \passthrough{\lstinline!m!} 模式 参数，Python 类型为
\passthrough{\lstinline!int!}。精确含义、单位、范围和枚举值需查目标型号对应头文件与协议。
& \passthrough{\lstinline!value = obj.m\_mode!} /
\passthrough{\lstinline!obj.m\_mode = value!} \\
\bottomrule
\end{longtable}

\textbf{方法索引}

\begin{longtable}[]{@{}
  >{\raggedright\arraybackslash}p{(\columnwidth - 6\tabcolsep) * \real{0.18}}
  >{\raggedright\arraybackslash}p{(\columnwidth - 6\tabcolsep) * \real{0.24}}
  >{\raggedright\arraybackslash}p{(\columnwidth - 6\tabcolsep) * \real{0.18}}
  >{\raggedright\arraybackslash}p{(\columnwidth - 6\tabcolsep) * \real{0.40}}@{}}
\toprule
方法 & Python 签名 & 状态 & 安全分类 \\
\midrule
\endhead
\protect\hyperlink{unitree-sdk2-cpp-robot-go2-obstaclesavoidmoveparameter-from-json-1}{\passthrough{\lstinline!from\_json!}}
&
\passthrough{\lstinline!def from\_json(self, json: dict[str, Any]) -> None!}
& \passthrough{\lstinline!SIGNATURE\_ONLY!} &
\passthrough{\lstinline!UNCLASSIFIED!} \\
\protect\hyperlink{unitree-sdk2-cpp-robot-go2-obstaclesavoidmoveparameter-to-json-1}{\passthrough{\lstinline!to\_json!}}
&
\passthrough{\lstinline!def to\_json(self, json: dict[str, Any]) -> None!}
& \passthrough{\lstinline!SIGNATURE\_ONLY!} &
\passthrough{\lstinline!UNCLASSIFIED!} \\
\bottomrule
\end{longtable}

\hypertarget{unitree-sdk2-cpp-robot-go2-obstaclesavoidmoveparameter-from-json-1}{%
\paragraph{\texorpdfstring{\texttt{unitree\_sdk2\_cpp.robot.go2.ObstaclesAvoidMoveParameter.from\_json}}{unitree\_sdk2\_cpp.robot.go2.ObstaclesAvoidMoveParameter.from\_json}}\label{unitree-sdk2-cpp-robot-go2-obstaclesavoidmoveparameter-from-json-1}}

计划从 JSON 风格字典读取字段并更新当前 SDK 值对象。

\textbf{签名}

\begin{lstlisting}[language=Python]
def from_json(self, json: dict[str, Any]) -> None
\end{lstlisting}

\textbf{可用性与安全性}

\textbf{\passthrough{\lstinline!SIGNATURE\_ONLY!} /
\passthrough{\lstinline!UNCLASSIFIED!}}：当前只有设计期签名，不能假设已存在于运行时扩展。

\textbf{参数}

\begin{longtable}[]{@{}
  >{\raggedright\arraybackslash}p{(\columnwidth - 6\tabcolsep) * \real{0.18}}
  >{\raggedright\arraybackslash}p{(\columnwidth - 6\tabcolsep) * \real{0.24}}
  >{\raggedright\arraybackslash}p{(\columnwidth - 6\tabcolsep) * \real{0.18}}
  >{\raggedright\arraybackslash}p{(\columnwidth - 6\tabcolsep) * \real{0.40}}@{}}
\toprule
名称 & Python 类型 & 默认值 & 含义 \\
\midrule
\endhead
\passthrough{\lstinline!json!} &
\passthrough{\lstinline!dict[str, Any]!} & 必填 & JSON
风格字典，键和值必须满足对应 C++ \passthrough{\lstinline!JsonMap!}
数据结构的约束。 对应 C++ 参数
\passthrough{\lstinline!json: common::JsonMap \&!}。 \\
\bottomrule
\end{longtable}

\textbf{返回值}

\begin{longtable}[]{@{}
  >{\raggedright\arraybackslash}p{(\columnwidth - 4\tabcolsep) * \real{0.18}}
  >{\raggedright\arraybackslash}p{(\columnwidth - 4\tabcolsep) * \real{0.20}}
  >{\raggedright\arraybackslash}p{(\columnwidth - 4\tabcolsep) * \real{0.62}}@{}}
\toprule
位置 & 类型 & 含义 \\
\midrule
\endhead
无 & \passthrough{\lstinline!None!} & 该方法不返回 Python
值；失败可能表现为异常或底层状态变化。 \\
\bottomrule
\end{longtable}

\textbf{对应 C++}

\begin{itemize}
\tightlist
\item
  类：\passthrough{\lstinline!unitree::robot::go2::ObstaclesAvoidMoveParameter!}
\item
  签名：\passthrough{\lstinline!fromJson(common::JsonMap \&)!}
\item
  绑定策略：\passthrough{\lstinline!SIGNATURE\_PREVIEW!}
\item
  声明位置：\passthrough{\lstinline!include/unitree/robot/go2/obstacles\_avoid/obstacles\_avoid\_api.hpp:55!}
\end{itemize}

\textbf{说明}

\begin{itemize}
\tightlist
\item
  示例是局部调用片段；\passthrough{\lstinline!obj!}
  和参数变量表示已经按上层业务校验并准备好的对象。
\item
  该条目可以被编辑器和类型检查器识别，但当前运行时可能在属性查找阶段直接失败。
\end{itemize}

\textbf{用法}

\begin{lstlisting}[language=Python]
from typing import TYPE_CHECKING

if TYPE_CHECKING:
    def planned_from_json(obj: ObstaclesAvoidMoveParameter, json: dict[str, Any]) -> None:
        obj.from_json(json=json)
\end{lstlisting}

\begin{quote}
{[}!CAUTION{]} 上面的 \passthrough{\lstinline!TYPE\_CHECKING!}
示例只表达计划调用形态。该分支在普通 Python
运行时为假，不代表当前扩展能执行此接口。
\end{quote}

\hypertarget{unitree-sdk2-cpp-robot-go2-obstaclesavoidmoveparameter-to-json-1}{%
\paragraph{\texorpdfstring{\texttt{unitree\_sdk2\_cpp.robot.go2.ObstaclesAvoidMoveParameter.to\_json}}{unitree\_sdk2\_cpp.robot.go2.ObstaclesAvoidMoveParameter.to\_json}}\label{unitree-sdk2-cpp-robot-go2-obstaclesavoidmoveparameter-to-json-1}}

计划把当前 SDK 值对象写入 JSON 风格字典。

\textbf{签名}

\begin{lstlisting}[language=Python]
def to_json(self, json: dict[str, Any]) -> None
\end{lstlisting}

\textbf{可用性与安全性}

\textbf{\passthrough{\lstinline!SIGNATURE\_ONLY!} /
\passthrough{\lstinline!UNCLASSIFIED!}}：当前只有设计期签名，不能假设已存在于运行时扩展。

\textbf{参数}

\begin{longtable}[]{@{}
  >{\raggedright\arraybackslash}p{(\columnwidth - 6\tabcolsep) * \real{0.18}}
  >{\raggedright\arraybackslash}p{(\columnwidth - 6\tabcolsep) * \real{0.24}}
  >{\raggedright\arraybackslash}p{(\columnwidth - 6\tabcolsep) * \real{0.18}}
  >{\raggedright\arraybackslash}p{(\columnwidth - 6\tabcolsep) * \real{0.40}}@{}}
\toprule
名称 & Python 类型 & 默认值 & 含义 \\
\midrule
\endhead
\passthrough{\lstinline!json!} &
\passthrough{\lstinline!dict[str, Any]!} & 必填 & JSON
风格字典，键和值必须满足对应 C++ \passthrough{\lstinline!JsonMap!}
数据结构的约束。 对应 C++ 参数
\passthrough{\lstinline!json: common::JsonMap \&!}。 \\
\bottomrule
\end{longtable}

\textbf{返回值}

\begin{longtable}[]{@{}
  >{\raggedright\arraybackslash}p{(\columnwidth - 4\tabcolsep) * \real{0.18}}
  >{\raggedright\arraybackslash}p{(\columnwidth - 4\tabcolsep) * \real{0.20}}
  >{\raggedright\arraybackslash}p{(\columnwidth - 4\tabcolsep) * \real{0.62}}@{}}
\toprule
位置 & 类型 & 含义 \\
\midrule
\endhead
无 & \passthrough{\lstinline!None!} & 该方法不返回 Python
值；失败可能表现为异常或底层状态变化。 \\
\bottomrule
\end{longtable}

\textbf{对应 C++}

\begin{itemize}
\tightlist
\item
  类：\passthrough{\lstinline!unitree::robot::go2::ObstaclesAvoidMoveParameter!}
\item
  签名：\passthrough{\lstinline!toJson(common::JsonMap \&) const!}
\item
  绑定策略：\passthrough{\lstinline!SIGNATURE\_PREVIEW!}
\item
  声明位置：\passthrough{\lstinline!include/unitree/robot/go2/obstacles\_avoid/obstacles\_avoid\_api.hpp:63!}
\end{itemize}

\textbf{说明}

\begin{itemize}
\tightlist
\item
  示例是局部调用片段；\passthrough{\lstinline!obj!}
  和参数变量表示已经按上层业务校验并准备好的对象。
\item
  该条目可以被编辑器和类型检查器识别，但当前运行时可能在属性查找阶段直接失败。
\end{itemize}

\textbf{用法}

\begin{lstlisting}[language=Python]
from typing import TYPE_CHECKING

if TYPE_CHECKING:
    def planned_to_json(obj: ObstaclesAvoidMoveParameter, json: dict[str, Any]) -> None:
        obj.to_json(json=json)
\end{lstlisting}

\begin{quote}
{[}!CAUTION{]} 上面的 \passthrough{\lstinline!TYPE\_CHECKING!}
示例只表达计划调用形态。该分支在普通 Python
运行时为假，不代表当前扩展能执行此接口。
\end{quote}

\hypertarget{unitree-sdk2-cpp-robot-go2-obstaclesavoidremotecommandsource}{%
\subsubsection{\texorpdfstring{\texttt{unitree\_sdk2\_cpp.robot.go2.ObstaclesAvoidRemoteCommandSource}}{unitree\_sdk2\_cpp.robot.go2.ObstaclesAvoidRemoteCommandSource}}\label{unitree-sdk2-cpp-robot-go2-obstaclesavoidremotecommandsource}}

SDK 类型签名预览。请逐项查看构造函数和方法的可用性。

\textbf{导入}

\begin{lstlisting}[language=Python]
from unitree_sdk2_cpp.robot.go2 import ObstaclesAvoidRemoteCommandSource
\end{lstlisting}

\textbf{基类}：\passthrough{\lstinline!object!}

\textbf{公开属性}

\begin{longtable}[]{@{}
  >{\raggedright\arraybackslash}p{(\columnwidth - 6\tabcolsep) * \real{0.18}}
  >{\raggedright\arraybackslash}p{(\columnwidth - 6\tabcolsep) * \real{0.24}}
  >{\raggedright\arraybackslash}p{(\columnwidth - 6\tabcolsep) * \real{0.18}}
  >{\raggedright\arraybackslash}p{(\columnwidth - 6\tabcolsep) * \real{0.40}}@{}}
\toprule
名称 & Python 类型 & 含义 & 用法 \\
\midrule
\endhead
\passthrough{\lstinline!m\_is\_remote\_commands\_from\_api!} &
\passthrough{\lstinline!bool!} & 传给该接口的
\passthrough{\lstinline!m!} \passthrough{\lstinline!is!}
\passthrough{\lstinline!remote!} \passthrough{\lstinline!commands!}
\passthrough{\lstinline!from!} API 参数，Python 类型为
\passthrough{\lstinline!bool!}。精确含义、单位、范围和枚举值需查目标型号对应头文件与协议。
&
\passthrough{\lstinline!value = obj.m\_is\_remote\_commands\_from\_api!}
/
\passthrough{\lstinline!obj.m\_is\_remote\_commands\_from\_api = value!} \\
\bottomrule
\end{longtable}

\textbf{方法索引}

\begin{longtable}[]{@{}
  >{\raggedright\arraybackslash}p{(\columnwidth - 6\tabcolsep) * \real{0.18}}
  >{\raggedright\arraybackslash}p{(\columnwidth - 6\tabcolsep) * \real{0.24}}
  >{\raggedright\arraybackslash}p{(\columnwidth - 6\tabcolsep) * \real{0.18}}
  >{\raggedright\arraybackslash}p{(\columnwidth - 6\tabcolsep) * \real{0.40}}@{}}
\toprule
方法 & Python 签名 & 状态 & 安全分类 \\
\midrule
\endhead
\protect\hyperlink{unitree-sdk2-cpp-robot-go2-obstaclesavoidremotecommandsource-from-json-1}{\passthrough{\lstinline!from\_json!}}
&
\passthrough{\lstinline!def from\_json(self, json: dict[str, Any]) -> None!}
& \passthrough{\lstinline!SIGNATURE\_ONLY!} &
\passthrough{\lstinline!UNCLASSIFIED!} \\
\protect\hyperlink{unitree-sdk2-cpp-robot-go2-obstaclesavoidremotecommandsource-to-json-1}{\passthrough{\lstinline!to\_json!}}
&
\passthrough{\lstinline!def to\_json(self, json: dict[str, Any]) -> None!}
& \passthrough{\lstinline!SIGNATURE\_ONLY!} &
\passthrough{\lstinline!UNCLASSIFIED!} \\
\bottomrule
\end{longtable}

\hypertarget{unitree-sdk2-cpp-robot-go2-obstaclesavoidremotecommandsource-from-json-1}{%
\paragraph{\texorpdfstring{\texttt{unitree\_sdk2\_cpp.robot.go2.ObstaclesAvoidRemoteCommandSource.from\_json}}{unitree\_sdk2\_cpp.robot.go2.ObstaclesAvoidRemoteCommandSource.from\_json}}\label{unitree-sdk2-cpp-robot-go2-obstaclesavoidremotecommandsource-from-json-1}}

计划从 JSON 风格字典读取字段并更新当前 SDK 值对象。

\textbf{签名}

\begin{lstlisting}[language=Python]
def from_json(self, json: dict[str, Any]) -> None
\end{lstlisting}

\textbf{可用性与安全性}

\textbf{\passthrough{\lstinline!SIGNATURE\_ONLY!} /
\passthrough{\lstinline!UNCLASSIFIED!}}：当前只有设计期签名，不能假设已存在于运行时扩展。

\textbf{参数}

\begin{longtable}[]{@{}
  >{\raggedright\arraybackslash}p{(\columnwidth - 6\tabcolsep) * \real{0.18}}
  >{\raggedright\arraybackslash}p{(\columnwidth - 6\tabcolsep) * \real{0.24}}
  >{\raggedright\arraybackslash}p{(\columnwidth - 6\tabcolsep) * \real{0.18}}
  >{\raggedright\arraybackslash}p{(\columnwidth - 6\tabcolsep) * \real{0.40}}@{}}
\toprule
名称 & Python 类型 & 默认值 & 含义 \\
\midrule
\endhead
\passthrough{\lstinline!json!} &
\passthrough{\lstinline!dict[str, Any]!} & 必填 & JSON
风格字典，键和值必须满足对应 C++ \passthrough{\lstinline!JsonMap!}
数据结构的约束。 对应 C++ 参数
\passthrough{\lstinline!json: common::JsonMap \&!}。 \\
\bottomrule
\end{longtable}

\textbf{返回值}

\begin{longtable}[]{@{}
  >{\raggedright\arraybackslash}p{(\columnwidth - 4\tabcolsep) * \real{0.18}}
  >{\raggedright\arraybackslash}p{(\columnwidth - 4\tabcolsep) * \real{0.20}}
  >{\raggedright\arraybackslash}p{(\columnwidth - 4\tabcolsep) * \real{0.62}}@{}}
\toprule
位置 & 类型 & 含义 \\
\midrule
\endhead
无 & \passthrough{\lstinline!None!} & 该方法不返回 Python
值；失败可能表现为异常或底层状态变化。 \\
\bottomrule
\end{longtable}

\textbf{对应 C++}

\begin{itemize}
\tightlist
\item
  类：\passthrough{\lstinline!unitree::robot::go2::ObstaclesAvoidRemoteCommandSource!}
\item
  签名：\passthrough{\lstinline!fromJson(common::JsonMap \&)!}
\item
  绑定策略：\passthrough{\lstinline!SIGNATURE\_PREVIEW!}
\item
  声明位置：\passthrough{\lstinline!include/unitree/robot/go2/obstacles\_avoid/obstacles\_avoid\_api.hpp:87!}
\end{itemize}

\textbf{说明}

\begin{itemize}
\tightlist
\item
  示例是局部调用片段；\passthrough{\lstinline!obj!}
  和参数变量表示已经按上层业务校验并准备好的对象。
\item
  该条目可以被编辑器和类型检查器识别，但当前运行时可能在属性查找阶段直接失败。
\end{itemize}

\textbf{用法}

\begin{lstlisting}[language=Python]
from typing import TYPE_CHECKING

if TYPE_CHECKING:
    def planned_from_json(obj: ObstaclesAvoidRemoteCommandSource, json: dict[str, Any]) -> None:
        obj.from_json(json=json)
\end{lstlisting}

\begin{quote}
{[}!CAUTION{]} 上面的 \passthrough{\lstinline!TYPE\_CHECKING!}
示例只表达计划调用形态。该分支在普通 Python
运行时为假，不代表当前扩展能执行此接口。
\end{quote}

\hypertarget{unitree-sdk2-cpp-robot-go2-obstaclesavoidremotecommandsource-to-json-1}{%
\paragraph{\texorpdfstring{\texttt{unitree\_sdk2\_cpp.robot.go2.ObstaclesAvoidRemoteCommandSource.to\_json}}{unitree\_sdk2\_cpp.robot.go2.ObstaclesAvoidRemoteCommandSource.to\_json}}\label{unitree-sdk2-cpp-robot-go2-obstaclesavoidremotecommandsource-to-json-1}}

计划把当前 SDK 值对象写入 JSON 风格字典。

\textbf{签名}

\begin{lstlisting}[language=Python]
def to_json(self, json: dict[str, Any]) -> None
\end{lstlisting}

\textbf{可用性与安全性}

\textbf{\passthrough{\lstinline!SIGNATURE\_ONLY!} /
\passthrough{\lstinline!UNCLASSIFIED!}}：当前只有设计期签名，不能假设已存在于运行时扩展。

\textbf{参数}

\begin{longtable}[]{@{}
  >{\raggedright\arraybackslash}p{(\columnwidth - 6\tabcolsep) * \real{0.18}}
  >{\raggedright\arraybackslash}p{(\columnwidth - 6\tabcolsep) * \real{0.24}}
  >{\raggedright\arraybackslash}p{(\columnwidth - 6\tabcolsep) * \real{0.18}}
  >{\raggedright\arraybackslash}p{(\columnwidth - 6\tabcolsep) * \real{0.40}}@{}}
\toprule
名称 & Python 类型 & 默认值 & 含义 \\
\midrule
\endhead
\passthrough{\lstinline!json!} &
\passthrough{\lstinline!dict[str, Any]!} & 必填 & JSON
风格字典，键和值必须满足对应 C++ \passthrough{\lstinline!JsonMap!}
数据结构的约束。 对应 C++ 参数
\passthrough{\lstinline!json: common::JsonMap \&!}。 \\
\bottomrule
\end{longtable}

\textbf{返回值}

\begin{longtable}[]{@{}
  >{\raggedright\arraybackslash}p{(\columnwidth - 4\tabcolsep) * \real{0.18}}
  >{\raggedright\arraybackslash}p{(\columnwidth - 4\tabcolsep) * \real{0.20}}
  >{\raggedright\arraybackslash}p{(\columnwidth - 4\tabcolsep) * \real{0.62}}@{}}
\toprule
位置 & 类型 & 含义 \\
\midrule
\endhead
无 & \passthrough{\lstinline!None!} & 该方法不返回 Python
值；失败可能表现为异常或底层状态变化。 \\
\bottomrule
\end{longtable}

\textbf{对应 C++}

\begin{itemize}
\tightlist
\item
  类：\passthrough{\lstinline!unitree::robot::go2::ObstaclesAvoidRemoteCommandSource!}
\item
  签名：\passthrough{\lstinline!toJson(common::JsonMap \&) const!}
\item
  绑定策略：\passthrough{\lstinline!SIGNATURE\_PREVIEW!}
\item
  声明位置：\passthrough{\lstinline!include/unitree/robot/go2/obstacles\_avoid/obstacles\_avoid\_api.hpp:92!}
\end{itemize}

\textbf{说明}

\begin{itemize}
\tightlist
\item
  示例是局部调用片段；\passthrough{\lstinline!obj!}
  和参数变量表示已经按上层业务校验并准备好的对象。
\item
  该条目可以被编辑器和类型检查器识别，但当前运行时可能在属性查找阶段直接失败。
\end{itemize}

\textbf{用法}

\begin{lstlisting}[language=Python]
from typing import TYPE_CHECKING

if TYPE_CHECKING:
    def planned_to_json(obj: ObstaclesAvoidRemoteCommandSource, json: dict[str, Any]) -> None:
        obj.to_json(json=json)
\end{lstlisting}

\begin{quote}
{[}!CAUTION{]} 上面的 \passthrough{\lstinline!TYPE\_CHECKING!}
示例只表达计划调用形态。该分支在普通 Python
运行时为假，不代表当前扩展能执行此接口。
\end{quote}

\hypertarget{unitree-sdk2-cpp-robot-go2-obstaclesavoidswitchgetdata}{%
\subsubsection{\texorpdfstring{\texttt{unitree\_sdk2\_cpp.robot.go2.ObstaclesAvoidSwitchGetData}}{unitree\_sdk2\_cpp.robot.go2.ObstaclesAvoidSwitchGetData}}\label{unitree-sdk2-cpp-robot-go2-obstaclesavoidswitchgetdata}}

SDK 数据值类型；公开字段和转换方法见下文。

\textbf{导入}

\begin{lstlisting}[language=Python]
from unitree_sdk2_cpp.robot.go2 import ObstaclesAvoidSwitchGetData
\end{lstlisting}

\textbf{基类}：\passthrough{\lstinline!object!}

\textbf{公开属性}

\begin{longtable}[]{@{}
  >{\raggedright\arraybackslash}p{(\columnwidth - 6\tabcolsep) * \real{0.18}}
  >{\raggedright\arraybackslash}p{(\columnwidth - 6\tabcolsep) * \real{0.24}}
  >{\raggedright\arraybackslash}p{(\columnwidth - 6\tabcolsep) * \real{0.18}}
  >{\raggedright\arraybackslash}p{(\columnwidth - 6\tabcolsep) * \real{0.40}}@{}}
\toprule
名称 & Python 类型 & 含义 & 用法 \\
\midrule
\endhead
\passthrough{\lstinline!m\_enable!} & \passthrough{\lstinline!bool!} &
传给该接口的 \passthrough{\lstinline!m!}
\passthrough{\lstinline!enable!} 参数，Python 类型为
\passthrough{\lstinline!bool!}。精确含义、单位、范围和枚举值需查目标型号对应头文件与协议。
& \passthrough{\lstinline!value = obj.m\_enable!} /
\passthrough{\lstinline!obj.m\_enable = value!} \\
\bottomrule
\end{longtable}

\textbf{方法索引}

\begin{longtable}[]{@{}
  >{\raggedright\arraybackslash}p{(\columnwidth - 6\tabcolsep) * \real{0.18}}
  >{\raggedright\arraybackslash}p{(\columnwidth - 6\tabcolsep) * \real{0.24}}
  >{\raggedright\arraybackslash}p{(\columnwidth - 6\tabcolsep) * \real{0.18}}
  >{\raggedright\arraybackslash}p{(\columnwidth - 6\tabcolsep) * \real{0.40}}@{}}
\toprule
方法 & Python 签名 & 状态 & 安全分类 \\
\midrule
\endhead
\protect\hyperlink{unitree-sdk2-cpp-robot-go2-obstaclesavoidswitchgetdata-from-json-1}{\passthrough{\lstinline!from\_json!}}
&
\passthrough{\lstinline!def from\_json(self, json: dict[str, Any]) -> None!}
& \passthrough{\lstinline!SIGNATURE\_ONLY!} &
\passthrough{\lstinline!UNCLASSIFIED!} \\
\protect\hyperlink{unitree-sdk2-cpp-robot-go2-obstaclesavoidswitchgetdata-to-json-1}{\passthrough{\lstinline!to\_json!}}
&
\passthrough{\lstinline!def to\_json(self, json: dict[str, Any]) -> None!}
& \passthrough{\lstinline!SIGNATURE\_ONLY!} &
\passthrough{\lstinline!UNCLASSIFIED!} \\
\bottomrule
\end{longtable}

\hypertarget{unitree-sdk2-cpp-robot-go2-obstaclesavoidswitchgetdata-from-json-1}{%
\paragraph{\texorpdfstring{\texttt{unitree\_sdk2\_cpp.robot.go2.ObstaclesAvoidSwitchGetData.from\_json}}{unitree\_sdk2\_cpp.robot.go2.ObstaclesAvoidSwitchGetData.from\_json}}\label{unitree-sdk2-cpp-robot-go2-obstaclesavoidswitchgetdata-from-json-1}}

计划从 JSON 风格字典读取字段并更新当前 SDK 值对象。

\textbf{签名}

\begin{lstlisting}[language=Python]
def from_json(self, json: dict[str, Any]) -> None
\end{lstlisting}

\textbf{可用性与安全性}

\textbf{\passthrough{\lstinline!SIGNATURE\_ONLY!} /
\passthrough{\lstinline!UNCLASSIFIED!}}：当前只有设计期签名，不能假设已存在于运行时扩展。

\textbf{参数}

\begin{longtable}[]{@{}
  >{\raggedright\arraybackslash}p{(\columnwidth - 6\tabcolsep) * \real{0.18}}
  >{\raggedright\arraybackslash}p{(\columnwidth - 6\tabcolsep) * \real{0.24}}
  >{\raggedright\arraybackslash}p{(\columnwidth - 6\tabcolsep) * \real{0.18}}
  >{\raggedright\arraybackslash}p{(\columnwidth - 6\tabcolsep) * \real{0.40}}@{}}
\toprule
名称 & Python 类型 & 默认值 & 含义 \\
\midrule
\endhead
\passthrough{\lstinline!json!} &
\passthrough{\lstinline!dict[str, Any]!} & 必填 & JSON
风格字典，键和值必须满足对应 C++ \passthrough{\lstinline!JsonMap!}
数据结构的约束。 对应 C++ 参数
\passthrough{\lstinline!json: common::JsonMap \&!}。 \\
\bottomrule
\end{longtable}

\textbf{返回值}

\begin{longtable}[]{@{}
  >{\raggedright\arraybackslash}p{(\columnwidth - 4\tabcolsep) * \real{0.18}}
  >{\raggedright\arraybackslash}p{(\columnwidth - 4\tabcolsep) * \real{0.20}}
  >{\raggedright\arraybackslash}p{(\columnwidth - 4\tabcolsep) * \real{0.62}}@{}}
\toprule
位置 & 类型 & 含义 \\
\midrule
\endhead
无 & \passthrough{\lstinline!None!} & 该方法不返回 Python
值；失败可能表现为异常或底层状态变化。 \\
\bottomrule
\end{longtable}

\textbf{对应 C++}

\begin{itemize}
\tightlist
\item
  类：\passthrough{\lstinline!unitree::robot::go2::ObstaclesAvoidSwitchGetData!}
\item
  签名：\passthrough{\lstinline!fromJson(common::JsonMap \&)!}
\item
  绑定策略：\passthrough{\lstinline!SIGNATURE\_PREVIEW!}
\item
  声明位置：\passthrough{\lstinline!include/unitree/robot/go2/obstacles\_avoid/obstacles\_avoid\_api.hpp:39!}
\end{itemize}

\textbf{说明}

\begin{itemize}
\tightlist
\item
  示例是局部调用片段；\passthrough{\lstinline!obj!}
  和参数变量表示已经按上层业务校验并准备好的对象。
\item
  该条目可以被编辑器和类型检查器识别，但当前运行时可能在属性查找阶段直接失败。
\end{itemize}

\textbf{用法}

\begin{lstlisting}[language=Python]
from typing import TYPE_CHECKING

if TYPE_CHECKING:
    def planned_from_json(obj: ObstaclesAvoidSwitchGetData, json: dict[str, Any]) -> None:
        obj.from_json(json=json)
\end{lstlisting}

\begin{quote}
{[}!CAUTION{]} 上面的 \passthrough{\lstinline!TYPE\_CHECKING!}
示例只表达计划调用形态。该分支在普通 Python
运行时为假，不代表当前扩展能执行此接口。
\end{quote}

\hypertarget{unitree-sdk2-cpp-robot-go2-obstaclesavoidswitchgetdata-to-json-1}{%
\paragraph{\texorpdfstring{\texttt{unitree\_sdk2\_cpp.robot.go2.ObstaclesAvoidSwitchGetData.to\_json}}{unitree\_sdk2\_cpp.robot.go2.ObstaclesAvoidSwitchGetData.to\_json}}\label{unitree-sdk2-cpp-robot-go2-obstaclesavoidswitchgetdata-to-json-1}}

计划把当前 SDK 值对象写入 JSON 风格字典。

\textbf{签名}

\begin{lstlisting}[language=Python]
def to_json(self, json: dict[str, Any]) -> None
\end{lstlisting}

\textbf{可用性与安全性}

\textbf{\passthrough{\lstinline!SIGNATURE\_ONLY!} /
\passthrough{\lstinline!UNCLASSIFIED!}}：当前只有设计期签名，不能假设已存在于运行时扩展。

\textbf{参数}

\begin{longtable}[]{@{}
  >{\raggedright\arraybackslash}p{(\columnwidth - 6\tabcolsep) * \real{0.18}}
  >{\raggedright\arraybackslash}p{(\columnwidth - 6\tabcolsep) * \real{0.24}}
  >{\raggedright\arraybackslash}p{(\columnwidth - 6\tabcolsep) * \real{0.18}}
  >{\raggedright\arraybackslash}p{(\columnwidth - 6\tabcolsep) * \real{0.40}}@{}}
\toprule
名称 & Python 类型 & 默认值 & 含义 \\
\midrule
\endhead
\passthrough{\lstinline!json!} &
\passthrough{\lstinline!dict[str, Any]!} & 必填 & JSON
风格字典，键和值必须满足对应 C++ \passthrough{\lstinline!JsonMap!}
数据结构的约束。 对应 C++ 参数
\passthrough{\lstinline!json: common::JsonMap \&!}。 \\
\bottomrule
\end{longtable}

\textbf{返回值}

\begin{longtable}[]{@{}
  >{\raggedright\arraybackslash}p{(\columnwidth - 4\tabcolsep) * \real{0.18}}
  >{\raggedright\arraybackslash}p{(\columnwidth - 4\tabcolsep) * \real{0.20}}
  >{\raggedright\arraybackslash}p{(\columnwidth - 4\tabcolsep) * \real{0.62}}@{}}
\toprule
位置 & 类型 & 含义 \\
\midrule
\endhead
无 & \passthrough{\lstinline!None!} & 该方法不返回 Python
值；失败可能表现为异常或底层状态变化。 \\
\bottomrule
\end{longtable}

\textbf{对应 C++}

\begin{itemize}
\tightlist
\item
  类：\passthrough{\lstinline!unitree::robot::go2::ObstaclesAvoidSwitchGetData!}
\item
  签名：\passthrough{\lstinline!toJson(common::JsonMap \&) const!}
\item
  绑定策略：\passthrough{\lstinline!SIGNATURE\_PREVIEW!}
\item
  声明位置：\passthrough{\lstinline!include/unitree/robot/go2/obstacles\_avoid/obstacles\_avoid\_api.hpp:44!}
\end{itemize}

\textbf{说明}

\begin{itemize}
\tightlist
\item
  示例是局部调用片段；\passthrough{\lstinline!obj!}
  和参数变量表示已经按上层业务校验并准备好的对象。
\item
  该条目可以被编辑器和类型检查器识别，但当前运行时可能在属性查找阶段直接失败。
\end{itemize}

\textbf{用法}

\begin{lstlisting}[language=Python]
from typing import TYPE_CHECKING

if TYPE_CHECKING:
    def planned_to_json(obj: ObstaclesAvoidSwitchGetData, json: dict[str, Any]) -> None:
        obj.to_json(json=json)
\end{lstlisting}

\begin{quote}
{[}!CAUTION{]} 上面的 \passthrough{\lstinline!TYPE\_CHECKING!}
示例只表达计划调用形态。该分支在普通 Python
运行时为假，不代表当前扩展能执行此接口。
\end{quote}

\hypertarget{unitree-sdk2-cpp-robot-go2-obstaclesavoidswitchsetparameter}{%
\subsubsection{\texorpdfstring{\texttt{unitree\_sdk2\_cpp.robot.go2.ObstaclesAvoidSwitchSetParameter}}{unitree\_sdk2\_cpp.robot.go2.ObstaclesAvoidSwitchSetParameter}}\label{unitree-sdk2-cpp-robot-go2-obstaclesavoidswitchsetparameter}}

SDK 请求参数值类型；当前可能仅提供设计期签名。

\textbf{导入}

\begin{lstlisting}[language=Python]
from unitree_sdk2_cpp.robot.go2 import ObstaclesAvoidSwitchSetParameter
\end{lstlisting}

\textbf{基类}：\passthrough{\lstinline!object!}

\textbf{公开属性}

\begin{longtable}[]{@{}
  >{\raggedright\arraybackslash}p{(\columnwidth - 6\tabcolsep) * \real{0.18}}
  >{\raggedright\arraybackslash}p{(\columnwidth - 6\tabcolsep) * \real{0.24}}
  >{\raggedright\arraybackslash}p{(\columnwidth - 6\tabcolsep) * \real{0.18}}
  >{\raggedright\arraybackslash}p{(\columnwidth - 6\tabcolsep) * \real{0.40}}@{}}
\toprule
名称 & Python 类型 & 含义 & 用法 \\
\midrule
\endhead
\passthrough{\lstinline!m\_enable!} & \passthrough{\lstinline!bool!} &
传给该接口的 \passthrough{\lstinline!m!}
\passthrough{\lstinline!enable!} 参数，Python 类型为
\passthrough{\lstinline!bool!}。精确含义、单位、范围和枚举值需查目标型号对应头文件与协议。
& \passthrough{\lstinline!value = obj.m\_enable!} /
\passthrough{\lstinline!obj.m\_enable = value!} \\
\bottomrule
\end{longtable}

\textbf{方法索引}

\begin{longtable}[]{@{}
  >{\raggedright\arraybackslash}p{(\columnwidth - 6\tabcolsep) * \real{0.18}}
  >{\raggedright\arraybackslash}p{(\columnwidth - 6\tabcolsep) * \real{0.24}}
  >{\raggedright\arraybackslash}p{(\columnwidth - 6\tabcolsep) * \real{0.18}}
  >{\raggedright\arraybackslash}p{(\columnwidth - 6\tabcolsep) * \real{0.40}}@{}}
\toprule
方法 & Python 签名 & 状态 & 安全分类 \\
\midrule
\endhead
\protect\hyperlink{unitree-sdk2-cpp-robot-go2-obstaclesavoidswitchsetparameter-from-json-1}{\passthrough{\lstinline!from\_json!}}
&
\passthrough{\lstinline!def from\_json(self, json: dict[str, Any]) -> None!}
& \passthrough{\lstinline!SIGNATURE\_ONLY!} &
\passthrough{\lstinline!UNCLASSIFIED!} \\
\protect\hyperlink{unitree-sdk2-cpp-robot-go2-obstaclesavoidswitchsetparameter-to-json-1}{\passthrough{\lstinline!to\_json!}}
&
\passthrough{\lstinline!def to\_json(self, json: dict[str, Any]) -> None!}
& \passthrough{\lstinline!SIGNATURE\_ONLY!} &
\passthrough{\lstinline!UNCLASSIFIED!} \\
\bottomrule
\end{longtable}

\hypertarget{unitree-sdk2-cpp-robot-go2-obstaclesavoidswitchsetparameter-from-json-1}{%
\paragraph{\texorpdfstring{\texttt{unitree\_sdk2\_cpp.robot.go2.ObstaclesAvoidSwitchSetParameter.from\_json}}{unitree\_sdk2\_cpp.robot.go2.ObstaclesAvoidSwitchSetParameter.from\_json}}\label{unitree-sdk2-cpp-robot-go2-obstaclesavoidswitchsetparameter-from-json-1}}

计划从 JSON 风格字典读取字段并更新当前 SDK 值对象。

\textbf{签名}

\begin{lstlisting}[language=Python]
def from_json(self, json: dict[str, Any]) -> None
\end{lstlisting}

\textbf{可用性与安全性}

\textbf{\passthrough{\lstinline!SIGNATURE\_ONLY!} /
\passthrough{\lstinline!UNCLASSIFIED!}}：当前只有设计期签名，不能假设已存在于运行时扩展。

\textbf{参数}

\begin{longtable}[]{@{}
  >{\raggedright\arraybackslash}p{(\columnwidth - 6\tabcolsep) * \real{0.18}}
  >{\raggedright\arraybackslash}p{(\columnwidth - 6\tabcolsep) * \real{0.24}}
  >{\raggedright\arraybackslash}p{(\columnwidth - 6\tabcolsep) * \real{0.18}}
  >{\raggedright\arraybackslash}p{(\columnwidth - 6\tabcolsep) * \real{0.40}}@{}}
\toprule
名称 & Python 类型 & 默认值 & 含义 \\
\midrule
\endhead
\passthrough{\lstinline!json!} &
\passthrough{\lstinline!dict[str, Any]!} & 必填 & JSON
风格字典，键和值必须满足对应 C++ \passthrough{\lstinline!JsonMap!}
数据结构的约束。 对应 C++ 参数
\passthrough{\lstinline!json: common::JsonMap \&!}。 \\
\bottomrule
\end{longtable}

\textbf{返回值}

\begin{longtable}[]{@{}
  >{\raggedright\arraybackslash}p{(\columnwidth - 4\tabcolsep) * \real{0.18}}
  >{\raggedright\arraybackslash}p{(\columnwidth - 4\tabcolsep) * \real{0.20}}
  >{\raggedright\arraybackslash}p{(\columnwidth - 4\tabcolsep) * \real{0.62}}@{}}
\toprule
位置 & 类型 & 含义 \\
\midrule
\endhead
无 & \passthrough{\lstinline!None!} & 该方法不返回 Python
值；失败可能表现为异常或底层状态变化。 \\
\bottomrule
\end{longtable}

\textbf{对应 C++}

\begin{itemize}
\tightlist
\item
  类：\passthrough{\lstinline!unitree::robot::go2::ObstaclesAvoidSwitchSetParameter!}
\item
  签名：\passthrough{\lstinline!fromJson(common::JsonMap \&)!}
\item
  绑定策略：\passthrough{\lstinline!SIGNATURE\_PREVIEW!}
\item
  声明位置：\passthrough{\lstinline!include/unitree/robot/go2/obstacles\_avoid/obstacles\_avoid\_api.hpp:23!}
\end{itemize}

\textbf{说明}

\begin{itemize}
\tightlist
\item
  示例是局部调用片段；\passthrough{\lstinline!obj!}
  和参数变量表示已经按上层业务校验并准备好的对象。
\item
  该条目可以被编辑器和类型检查器识别，但当前运行时可能在属性查找阶段直接失败。
\end{itemize}

\textbf{用法}

\begin{lstlisting}[language=Python]
from typing import TYPE_CHECKING

if TYPE_CHECKING:
    def planned_from_json(obj: ObstaclesAvoidSwitchSetParameter, json: dict[str, Any]) -> None:
        obj.from_json(json=json)
\end{lstlisting}

\begin{quote}
{[}!CAUTION{]} 上面的 \passthrough{\lstinline!TYPE\_CHECKING!}
示例只表达计划调用形态。该分支在普通 Python
运行时为假，不代表当前扩展能执行此接口。
\end{quote}

\hypertarget{unitree-sdk2-cpp-robot-go2-obstaclesavoidswitchsetparameter-to-json-1}{%
\paragraph{\texorpdfstring{\texttt{unitree\_sdk2\_cpp.robot.go2.ObstaclesAvoidSwitchSetParameter.to\_json}}{unitree\_sdk2\_cpp.robot.go2.ObstaclesAvoidSwitchSetParameter.to\_json}}\label{unitree-sdk2-cpp-robot-go2-obstaclesavoidswitchsetparameter-to-json-1}}

计划把当前 SDK 值对象写入 JSON 风格字典。

\textbf{签名}

\begin{lstlisting}[language=Python]
def to_json(self, json: dict[str, Any]) -> None
\end{lstlisting}

\textbf{可用性与安全性}

\textbf{\passthrough{\lstinline!SIGNATURE\_ONLY!} /
\passthrough{\lstinline!UNCLASSIFIED!}}：当前只有设计期签名，不能假设已存在于运行时扩展。

\textbf{参数}

\begin{longtable}[]{@{}
  >{\raggedright\arraybackslash}p{(\columnwidth - 6\tabcolsep) * \real{0.18}}
  >{\raggedright\arraybackslash}p{(\columnwidth - 6\tabcolsep) * \real{0.24}}
  >{\raggedright\arraybackslash}p{(\columnwidth - 6\tabcolsep) * \real{0.18}}
  >{\raggedright\arraybackslash}p{(\columnwidth - 6\tabcolsep) * \real{0.40}}@{}}
\toprule
名称 & Python 类型 & 默认值 & 含义 \\
\midrule
\endhead
\passthrough{\lstinline!json!} &
\passthrough{\lstinline!dict[str, Any]!} & 必填 & JSON
风格字典，键和值必须满足对应 C++ \passthrough{\lstinline!JsonMap!}
数据结构的约束。 对应 C++ 参数
\passthrough{\lstinline!json: common::JsonMap \&!}。 \\
\bottomrule
\end{longtable}

\textbf{返回值}

\begin{longtable}[]{@{}
  >{\raggedright\arraybackslash}p{(\columnwidth - 4\tabcolsep) * \real{0.18}}
  >{\raggedright\arraybackslash}p{(\columnwidth - 4\tabcolsep) * \real{0.20}}
  >{\raggedright\arraybackslash}p{(\columnwidth - 4\tabcolsep) * \real{0.62}}@{}}
\toprule
位置 & 类型 & 含义 \\
\midrule
\endhead
无 & \passthrough{\lstinline!None!} & 该方法不返回 Python
值；失败可能表现为异常或底层状态变化。 \\
\bottomrule
\end{longtable}

\textbf{对应 C++}

\begin{itemize}
\tightlist
\item
  类：\passthrough{\lstinline!unitree::robot::go2::ObstaclesAvoidSwitchSetParameter!}
\item
  签名：\passthrough{\lstinline!toJson(common::JsonMap \&) const!}
\item
  绑定策略：\passthrough{\lstinline!SIGNATURE\_PREVIEW!}
\item
  声明位置：\passthrough{\lstinline!include/unitree/robot/go2/obstacles\_avoid/obstacles\_avoid\_api.hpp:28!}
\end{itemize}

\textbf{说明}

\begin{itemize}
\tightlist
\item
  示例是局部调用片段；\passthrough{\lstinline!obj!}
  和参数变量表示已经按上层业务校验并准备好的对象。
\item
  该条目可以被编辑器和类型检查器识别，但当前运行时可能在属性查找阶段直接失败。
\end{itemize}

\textbf{用法}

\begin{lstlisting}[language=Python]
from typing import TYPE_CHECKING

if TYPE_CHECKING:
    def planned_to_json(obj: ObstaclesAvoidSwitchSetParameter, json: dict[str, Any]) -> None:
        obj.to_json(json=json)
\end{lstlisting}

\begin{quote}
{[}!CAUTION{]} 上面的 \passthrough{\lstinline!TYPE\_CHECKING!}
示例只表达计划调用形态。该分支在普通 Python
运行时为假，不代表当前扩展能执行此接口。
\end{quote}

\hypertarget{unitree-sdk2-cpp-robot-go2-robotstateclient}{%
\subsubsection{\texorpdfstring{\texttt{unitree\_sdk2\_cpp.robot.go2.RobotStateClient}}{unitree\_sdk2\_cpp.robot.go2.RobotStateClient}}\label{unitree-sdk2-cpp-robot-go2-robotstateclient}}

Robot 服务客户端。构造、初始化和具体方法的可用性必须分别检查。

\textbf{导入}

\begin{lstlisting}[language=Python]
from unitree_sdk2_cpp.robot.go2 import RobotStateClient
\end{lstlisting}

\textbf{基类}：\passthrough{\lstinline!Client!}

\textbf{方法索引}

\begin{longtable}[]{@{}
  >{\raggedright\arraybackslash}p{(\columnwidth - 6\tabcolsep) * \real{0.18}}
  >{\raggedright\arraybackslash}p{(\columnwidth - 6\tabcolsep) * \real{0.24}}
  >{\raggedright\arraybackslash}p{(\columnwidth - 6\tabcolsep) * \real{0.18}}
  >{\raggedright\arraybackslash}p{(\columnwidth - 6\tabcolsep) * \real{0.40}}@{}}
\toprule
方法 & Python 签名 & 状态 & 安全分类 \\
\midrule
\endhead
\protect\hyperlink{unitree-sdk2-cpp-robot-go2-robotstateclient-dunder-init-1}{\passthrough{\lstinline!\_\_init\_\_!}}
& \passthrough{\lstinline!def \_\_init\_\_(self) -> None!} &
\passthrough{\lstinline!AVAILABLE!} &
\passthrough{\lstinline!CONSTRUCTION!} \\
\protect\hyperlink{unitree-sdk2-cpp-robot-go2-robotstateclient-init-1}{\passthrough{\lstinline!init!}}
& \passthrough{\lstinline!def init(self) -> None!} &
\passthrough{\lstinline!AVAILABLE!} &
\passthrough{\lstinline!INITIALIZATION!} \\
\protect\hyperlink{unitree-sdk2-cpp-robot-go2-robotstateclient-service-list-1}{\passthrough{\lstinline!service\_list!}}
&
\passthrough{\lstinline!def service\_list(self) -> tuple[int, list[ServiceState]]!}
& \passthrough{\lstinline!AVAILABLE!} &
\passthrough{\lstinline!READ\_ONLY!} \\
\protect\hyperlink{unitree-sdk2-cpp-robot-go2-robotstateclient-service-switch-1}{\passthrough{\lstinline!service\_switch!}}
&
\passthrough{\lstinline!def service\_switch(self, name: str, swit: int) -> tuple[int, int]!}
& \passthrough{\lstinline!SIGNATURE\_ONLY!} &
\passthrough{\lstinline!HARDWARE\_SIDE\_EFFECT!} \\
\protect\hyperlink{unitree-sdk2-cpp-robot-go2-robotstateclient-set-report-freq-1}{\passthrough{\lstinline!set\_report\_freq!}}
&
\passthrough{\lstinline!def set\_report\_freq(self, interval: int, duration: int) -> int!}
& \passthrough{\lstinline!SIGNATURE\_ONLY!} &
\passthrough{\lstinline!HARDWARE\_SIDE\_EFFECT!} \\
\bottomrule
\end{longtable}

\hypertarget{unitree-sdk2-cpp-robot-go2-robotstateclient-dunder-init-1}{%
\paragraph{\texorpdfstring{\texttt{unitree\_sdk2\_cpp.robot.go2.RobotStateClient.\_\_init\_\_}}{unitree\_sdk2\_cpp.robot.go2.RobotStateClient.\_\_init\_\_}}\label{unitree-sdk2-cpp-robot-go2-robotstateclient-dunder-init-1}}

初始化 \passthrough{\lstinline!RobotStateClient!}
实例。是否能实际构造取决于下方可用性状态。

\textbf{签名}

\begin{lstlisting}[language=Python]
def __init__(self) -> None
\end{lstlisting}

\textbf{可用性与安全性}

\textbf{\passthrough{\lstinline!AVAILABLE!} /
\passthrough{\lstinline!CONSTRUCTION!}}：当前绑定源码已实现；实际执行条件仍取决于平台和对应
SDK 环境。

\textbf{参数}

无显式参数。实例方法中的 \passthrough{\lstinline!self!} 由 Python
自动传入。

\textbf{返回值}

\begin{longtable}[]{@{}
  >{\raggedright\arraybackslash}p{(\columnwidth - 4\tabcolsep) * \real{0.18}}
  >{\raggedright\arraybackslash}p{(\columnwidth - 4\tabcolsep) * \real{0.20}}
  >{\raggedright\arraybackslash}p{(\columnwidth - 4\tabcolsep) * \real{0.62}}@{}}
\toprule
位置 & 类型 & 含义 \\
\midrule
\endhead
无 & \passthrough{\lstinline!None!} &
\passthrough{\lstinline!\_\_init\_\_!}
本身不返回值；调用类对象时，在构造函数可用的前提下得到该类实例。 \\
\bottomrule
\end{longtable}

\textbf{对应 C++}

\begin{itemize}
\tightlist
\item
  类：\passthrough{\lstinline!unitree::robot::go2::RobotStateClient!}
\item
  签名：\passthrough{\lstinline!RobotStateClient()!}
\item
  绑定策略：\passthrough{\lstinline!未单独标注!}
\item
  声明位置：\passthrough{\lstinline!include/unitree/robot/go2/robot\_state/robot\_state\_client.hpp:32!}
\end{itemize}

\textbf{说明}

\begin{itemize}
\tightlist
\item
  示例是局部调用片段；\passthrough{\lstinline!obj!}
  和参数变量表示已经按上层业务校验并准备好的对象。
\end{itemize}

\textbf{用法}

\begin{lstlisting}[language=Python]
value = RobotStateClient()
\end{lstlisting}

\hypertarget{unitree-sdk2-cpp-robot-go2-robotstateclient-init-1}{%
\paragraph{\texorpdfstring{\texttt{unitree\_sdk2\_cpp.robot.go2.RobotStateClient.init}}{unitree\_sdk2\_cpp.robot.go2.RobotStateClient.init}}\label{unitree-sdk2-cpp-robot-go2-robotstateclient-init-1}}

初始化当前 SDK 对象所需的底层通道或服务资源。

\textbf{签名}

\begin{lstlisting}[language=Python]
def init(self) -> None
\end{lstlisting}

\textbf{可用性与安全性}

\textbf{\passthrough{\lstinline!AVAILABLE!} /
\passthrough{\lstinline!INITIALIZATION!}}：当前绑定源码已实现；实际执行条件仍取决于平台和对应
SDK 环境。

\textbf{参数}

无显式参数。实例方法中的 \passthrough{\lstinline!self!} 由 Python
自动传入。

\textbf{返回值}

\begin{longtable}[]{@{}
  >{\raggedright\arraybackslash}p{(\columnwidth - 4\tabcolsep) * \real{0.18}}
  >{\raggedright\arraybackslash}p{(\columnwidth - 4\tabcolsep) * \real{0.20}}
  >{\raggedright\arraybackslash}p{(\columnwidth - 4\tabcolsep) * \real{0.62}}@{}}
\toprule
位置 & 类型 & 含义 \\
\midrule
\endhead
无 & \passthrough{\lstinline!None!} & 该方法不返回 Python
值；失败可能表现为异常或底层状态变化。 \\
\bottomrule
\end{longtable}

\textbf{对应 C++}

\begin{itemize}
\tightlist
\item
  类：\passthrough{\lstinline!unitree::robot::go2::RobotStateClient!}
\item
  签名：\passthrough{\lstinline!Init()!}
\item
  绑定策略：\passthrough{\lstinline!DIRECT!}
\item
  声明位置：\passthrough{\lstinline!include/unitree/robot/go2/robot\_state/robot\_state\_client.hpp:35!}
\end{itemize}

\textbf{说明}

\begin{itemize}
\tightlist
\item
  示例是局部调用片段；\passthrough{\lstinline!obj!}
  和参数变量表示已经按上层业务校验并准备好的对象。
\end{itemize}

\textbf{用法}

\begin{lstlisting}[language=Python]
obj.init()
\end{lstlisting}

\hypertarget{unitree-sdk2-cpp-robot-go2-robotstateclient-service-list-1}{%
\paragraph{\texorpdfstring{\texttt{unitree\_sdk2\_cpp.robot.go2.RobotStateClient.service\_list}}{unitree\_sdk2\_cpp.robot.go2.RobotStateClient.service\_list}}\label{unitree-sdk2-cpp-robot-go2-robotstateclient-service-list-1}}

对应 C++ SDK 操作
\passthrough{\lstinline!ServiceList(std::vector<ServiceState> \&)!}。上游头文件没有可直接生成的业务说明时，本参考只保证签名映射，精确语义需查目标型号协议。

\textbf{签名}

\begin{lstlisting}[language=Python]
def service_list(self) -> tuple[int, list[ServiceState]]
\end{lstlisting}

\textbf{可用性与安全性}

\textbf{\passthrough{\lstinline!AVAILABLE!} /
\passthrough{\lstinline!READ\_ONLY!}}：当前绑定源码已实现只读查询。Robot
Client 的服务端查询通常仍需匹配的 Linux
扩展、DDS、网络、机器人和服务；纯本地 getter 则不一定需要全部条件。

\textbf{参数}

无显式参数。实例方法中的 \passthrough{\lstinline!self!} 由 Python
自动传入。

\textbf{返回值}

\begin{longtable}[]{@{}
  >{\raggedright\arraybackslash}p{(\columnwidth - 4\tabcolsep) * \real{0.18}}
  >{\raggedright\arraybackslash}p{(\columnwidth - 4\tabcolsep) * \real{0.20}}
  >{\raggedright\arraybackslash}p{(\columnwidth - 4\tabcolsep) * \real{0.62}}@{}}
\toprule
位置 & 类型 & 含义 \\
\midrule
\endhead
{[}0{]} \passthrough{\lstinline!status!} & \passthrough{\lstinline!int!}
& SDK 状态码；Unitree 示例通常以 \passthrough{\lstinline!0!}
表示成功，非零值需按具体服务错误码解释。 \\
{[}1{]} \passthrough{\lstinline!serviceStateList!} &
\passthrough{\lstinline!list[ServiceState]!} & 传给该接口的
\passthrough{\lstinline!serviceStateList!} 参数，Python 类型为
\passthrough{\lstinline!list[ServiceState]!}。精确含义、单位、范围和枚举值需查目标型号对应头文件与协议。
对应 C++ 参数
\passthrough{\lstinline!serviceStateList: std::vector<ServiceState> \&!}。
底层是可变长度 vector \\
\bottomrule
\end{longtable}

\textbf{对应 C++}

\begin{itemize}
\tightlist
\item
  类：\passthrough{\lstinline!unitree::robot::go2::RobotStateClient!}
\item
  签名：\passthrough{\lstinline!ServiceList(std::vector<ServiceState> \&)!}
\item
  绑定策略：\passthrough{\lstinline!OUTPUT\_WRAPPER!}
\item
  声明位置：\passthrough{\lstinline!include/unitree/robot/go2/robot\_state/robot\_state\_client.hpp:37!}
\end{itemize}

\textbf{说明}

\begin{itemize}
\tightlist
\item
  示例是局部调用片段；\passthrough{\lstinline!obj!}
  和参数变量表示已经按上层业务校验并准备好的对象。
\item
  C++ 的可修改输出引用不会作为 Python 入参出现，而是按 C++
  参数顺序追加到返回元组中。
\end{itemize}

\textbf{用法}

\begin{lstlisting}[language=Python]
status, serviceStateList = obj.service_list()
\end{lstlisting}

\begin{quote}
{[}!NOTE{]} 只读查询仍会访问
DDS/机器人服务。必须检查返回状态码，并处理超时和网络中断。
\end{quote}

\hypertarget{unitree-sdk2-cpp-robot-go2-robotstateclient-service-switch-1}{%
\paragraph{\texorpdfstring{\texttt{unitree\_sdk2\_cpp.robot.go2.RobotStateClient.service\_switch}}{unitree\_sdk2\_cpp.robot.go2.RobotStateClient.service\_switch}}\label{unitree-sdk2-cpp-robot-go2-robotstateclient-service-switch-1}}

对应 C++ SDK 操作
\passthrough{\lstinline!ServiceSwitch(const std::string \&, int32\_t, int32\_t \&)!}。上游头文件没有可直接生成的业务说明时，本参考只保证签名映射，精确语义需查目标型号协议。

\textbf{签名}

\begin{lstlisting}[language=Python]
def service_switch(self, name: str, swit: int) -> tuple[int, int]
\end{lstlisting}

\textbf{可用性与安全性}

\textbf{\passthrough{\lstinline!SIGNATURE\_ONLY!} /
\passthrough{\lstinline!HARDWARE\_SIDE\_EFFECT!}}：仅用于补全和静态检查。当前二进制没有该
Python 方法，未来实现还需单独评审硬件副作用。

\textbf{参数}

\begin{longtable}[]{@{}
  >{\raggedright\arraybackslash}p{(\columnwidth - 6\tabcolsep) * \real{0.18}}
  >{\raggedright\arraybackslash}p{(\columnwidth - 6\tabcolsep) * \real{0.24}}
  >{\raggedright\arraybackslash}p{(\columnwidth - 6\tabcolsep) * \real{0.18}}
  >{\raggedright\arraybackslash}p{(\columnwidth - 6\tabcolsep) * \real{0.40}}@{}}
\toprule
名称 & Python 类型 & 默认值 & 含义 \\
\midrule
\endhead
\passthrough{\lstinline!name!} & \passthrough{\lstinline!str!} & 必填 &
SDK 对象、配置项、服务或动作的名称。允许值由调用它的具体接口定义。 对应
C++ 参数 \passthrough{\lstinline!name: const std::string \&!}。
底层为字符串；长度、编码和允许值由具体协议决定。 \\
\passthrough{\lstinline!swit!} & \passthrough{\lstinline!int!} & 必填 &
传给该接口的 \passthrough{\lstinline!swit!} 参数，Python 类型为
\passthrough{\lstinline!int!}。精确含义、单位、范围和枚举值需查目标型号对应头文件与协议。
对应 C++ 参数 \passthrough{\lstinline!swit: int32\_t!}。 取值范围为
-2147483648 到 2147483647。 \\
\bottomrule
\end{longtable}

\textbf{返回值}

\begin{longtable}[]{@{}
  >{\raggedright\arraybackslash}p{(\columnwidth - 4\tabcolsep) * \real{0.18}}
  >{\raggedright\arraybackslash}p{(\columnwidth - 4\tabcolsep) * \real{0.20}}
  >{\raggedright\arraybackslash}p{(\columnwidth - 4\tabcolsep) * \real{0.62}}@{}}
\toprule
位置 & 类型 & 含义 \\
\midrule
\endhead
{[}0{]} \passthrough{\lstinline!status!} & \passthrough{\lstinline!int!}
& SDK 状态码；Unitree 示例通常以 \passthrough{\lstinline!0!}
表示成功，非零值需按具体服务错误码解释。 \\
{[}1{]} \passthrough{\lstinline!status!} & \passthrough{\lstinline!int!}
& SDK 状态码；Unitree 示例通常以 \passthrough{\lstinline!0!}
表示成功，非零值需按具体服务错误码解释。 \\
\bottomrule
\end{longtable}

\textbf{对应 C++}

\begin{itemize}
\tightlist
\item
  类：\passthrough{\lstinline!unitree::robot::go2::RobotStateClient!}
\item
  签名：\passthrough{\lstinline!ServiceSwitch(const std::string \&, int32\_t, int32\_t \&)!}
\item
  绑定策略：\passthrough{\lstinline!OUTPUT\_WRAPPER!}
\item
  声明位置：\passthrough{\lstinline!include/unitree/robot/go2/robot\_state/robot\_state\_client.hpp:38!}
\end{itemize}

\textbf{说明}

\begin{itemize}
\tightlist
\item
  示例是局部调用片段；\passthrough{\lstinline!obj!}
  和参数变量表示已经按上层业务校验并准备好的对象。
\item
  C++ 的可修改输出引用不会作为 Python 入参出现，而是按 C++
  参数顺序追加到返回元组中。
\item
  该条目可以被编辑器和类型检查器识别，但当前运行时可能在属性查找阶段直接失败。
\end{itemize}

\textbf{用法}

\begin{lstlisting}[language=Python]
from typing import TYPE_CHECKING

if TYPE_CHECKING:
    def planned_service_switch(obj: RobotStateClient, name: str, swit: int) -> tuple[int, int]:
        return obj.service_switch(name=name, swit=swit)
\end{lstlisting}

\begin{quote}
{[}!CAUTION{]} 上面的 \passthrough{\lstinline!TYPE\_CHECKING!}
示例只表达计划调用形态。该分支在普通 Python
运行时为假，不代表当前扩展能执行此接口。
\end{quote}

\hypertarget{unitree-sdk2-cpp-robot-go2-robotstateclient-set-report-freq-1}{%
\paragraph{\texorpdfstring{\texttt{unitree\_sdk2\_cpp.robot.go2.RobotStateClient.set\_report\_freq}}{unitree\_sdk2\_cpp.robot.go2.RobotStateClient.set\_report\_freq}}\label{unitree-sdk2-cpp-robot-go2-robotstateclient-set-report-freq-1}}

设置 \passthrough{\lstinline!report!}
\passthrough{\lstinline!freq!}。具体副作用和安全边界见下方状态。

\textbf{签名}

\begin{lstlisting}[language=Python]
def set_report_freq(self, interval: int, duration: int) -> int
\end{lstlisting}

\textbf{可用性与安全性}

\textbf{\passthrough{\lstinline!SIGNATURE\_ONLY!} /
\passthrough{\lstinline!HARDWARE\_SIDE\_EFFECT!}}：仅用于补全和静态检查。当前二进制没有该
Python 方法，未来实现还需单独评审硬件副作用。

\textbf{参数}

\begin{longtable}[]{@{}
  >{\raggedright\arraybackslash}p{(\columnwidth - 6\tabcolsep) * \real{0.18}}
  >{\raggedright\arraybackslash}p{(\columnwidth - 6\tabcolsep) * \real{0.24}}
  >{\raggedright\arraybackslash}p{(\columnwidth - 6\tabcolsep) * \real{0.18}}
  >{\raggedright\arraybackslash}p{(\columnwidth - 6\tabcolsep) * \real{0.40}}@{}}
\toprule
名称 & Python 类型 & 默认值 & 含义 \\
\midrule
\endhead
\passthrough{\lstinline!interval!} & \passthrough{\lstinline!int!} &
必填 & 传给该接口的 \passthrough{\lstinline!interval!} 参数，Python
类型为
\passthrough{\lstinline!int!}。精确含义、单位、范围和枚举值需查目标型号对应头文件与协议。
对应 C++ 参数 \passthrough{\lstinline!interval: int32\_t!}。 取值范围为
-2147483648 到 2147483647。 \\
\passthrough{\lstinline!duration!} & \passthrough{\lstinline!int!} &
必填 &
操作持续时间参数。精确单位、范围和默认行为以目标型号接口定义为准。 对应
C++ 参数 \passthrough{\lstinline!duration: int32\_t!}。 取值范围为
-2147483648 到 2147483647。 \\
\bottomrule
\end{longtable}

\textbf{返回值}

\begin{longtable}[]{@{}
  >{\raggedright\arraybackslash}p{(\columnwidth - 4\tabcolsep) * \real{0.18}}
  >{\raggedright\arraybackslash}p{(\columnwidth - 4\tabcolsep) * \real{0.20}}
  >{\raggedright\arraybackslash}p{(\columnwidth - 4\tabcolsep) * \real{0.62}}@{}}
\toprule
位置 & 类型 & 含义 \\
\midrule
\endhead
返回值 & \passthrough{\lstinline!int!} & SDK 状态码；Unitree 示例通常以
\passthrough{\lstinline!0!} 表示成功，非零值需按具体服务定义解释。 \\
\bottomrule
\end{longtable}

\textbf{对应 C++}

\begin{itemize}
\tightlist
\item
  类：\passthrough{\lstinline!unitree::robot::go2::RobotStateClient!}
\item
  签名：\passthrough{\lstinline!SetReportFreq(int32\_t, int32\_t)!}
\item
  绑定策略：\passthrough{\lstinline!DIRECT!}
\item
  声明位置：\passthrough{\lstinline!include/unitree/robot/go2/robot\_state/robot\_state\_client.hpp:39!}
\end{itemize}

\textbf{说明}

\begin{itemize}
\tightlist
\item
  示例是局部调用片段；\passthrough{\lstinline!obj!}
  和参数变量表示已经按上层业务校验并准备好的对象。
\item
  该条目可以被编辑器和类型检查器识别，但当前运行时可能在属性查找阶段直接失败。
\end{itemize}

\textbf{用法}

\begin{lstlisting}[language=Python]
from typing import TYPE_CHECKING

if TYPE_CHECKING:
    def planned_set_report_freq(obj: RobotStateClient, interval: int, duration: int) -> int:
        return obj.set_report_freq(interval=interval, duration=duration)
\end{lstlisting}

\begin{quote}
{[}!CAUTION{]} 上面的 \passthrough{\lstinline!TYPE\_CHECKING!}
示例只表达计划调用形态。该分支在普通 Python
运行时为假，不代表当前扩展能执行此接口。
\end{quote}

\hypertarget{unitree-sdk2-cpp-robot-go2-servicestate}{%
\subsubsection{\texorpdfstring{\texttt{unitree\_sdk2\_cpp.robot.go2.ServiceState}}{unitree\_sdk2\_cpp.robot.go2.ServiceState}}\label{unitree-sdk2-cpp-robot-go2-servicestate}}

SDK 数据值类型；公开字段和转换方法见下文。

\textbf{导入}

\begin{lstlisting}[language=Python]
from unitree_sdk2_cpp.robot.go2 import ServiceState
\end{lstlisting}

\textbf{基类}：\passthrough{\lstinline!object!}

\textbf{公开属性}

\begin{longtable}[]{@{}
  >{\raggedright\arraybackslash}p{(\columnwidth - 6\tabcolsep) * \real{0.18}}
  >{\raggedright\arraybackslash}p{(\columnwidth - 6\tabcolsep) * \real{0.24}}
  >{\raggedright\arraybackslash}p{(\columnwidth - 6\tabcolsep) * \real{0.18}}
  >{\raggedright\arraybackslash}p{(\columnwidth - 6\tabcolsep) * \real{0.40}}@{}}
\toprule
名称 & Python 类型 & 含义 & 用法 \\
\midrule
\endhead
\passthrough{\lstinline!name!} & \passthrough{\lstinline!str!} & SDK
对象、配置项、服务或动作的名称。允许值由调用它的具体接口定义。 &
\passthrough{\lstinline!value = obj.name!} /
\passthrough{\lstinline!obj.name = value!} \\
\passthrough{\lstinline!status!} & \passthrough{\lstinline!int!} &
传给该接口的 状态 参数，Python 类型为
\passthrough{\lstinline!int!}。精确含义、单位、范围和枚举值需查目标型号对应头文件与协议。
& \passthrough{\lstinline!value = obj.status!} /
\passthrough{\lstinline!obj.status = value!} \\
\passthrough{\lstinline!protect!} & \passthrough{\lstinline!int!} &
传给该接口的 \passthrough{\lstinline!protect!} 参数，Python 类型为
\passthrough{\lstinline!int!}。精确含义、单位、范围和枚举值需查目标型号对应头文件与协议。
& \passthrough{\lstinline!value = obj.protect!} /
\passthrough{\lstinline!obj.protect = value!} \\
\bottomrule
\end{longtable}

\textbf{方法索引}

\begin{longtable}[]{@{}
  >{\raggedright\arraybackslash}p{(\columnwidth - 6\tabcolsep) * \real{0.18}}
  >{\raggedright\arraybackslash}p{(\columnwidth - 6\tabcolsep) * \real{0.24}}
  >{\raggedright\arraybackslash}p{(\columnwidth - 6\tabcolsep) * \real{0.18}}
  >{\raggedright\arraybackslash}p{(\columnwidth - 6\tabcolsep) * \real{0.40}}@{}}
\toprule
方法 & Python 签名 & 状态 & 安全分类 \\
\midrule
\endhead
\protect\hyperlink{unitree-sdk2-cpp-robot-go2-servicestate-dunder-init-1}{\passthrough{\lstinline!\_\_init\_\_!}}
& \passthrough{\lstinline!def \_\_init\_\_(self) -> None!} &
\passthrough{\lstinline!AVAILABLE!} &
\passthrough{\lstinline!CONSTRUCTION!} \\
\bottomrule
\end{longtable}

\hypertarget{unitree-sdk2-cpp-robot-go2-servicestate-dunder-init-1}{%
\paragraph{\texorpdfstring{\texttt{unitree\_sdk2\_cpp.robot.go2.ServiceState.\_\_init\_\_}}{unitree\_sdk2\_cpp.robot.go2.ServiceState.\_\_init\_\_}}\label{unitree-sdk2-cpp-robot-go2-servicestate-dunder-init-1}}

初始化 \passthrough{\lstinline!ServiceState!}
实例。是否能实际构造取决于下方可用性状态。

\textbf{签名}

\begin{lstlisting}[language=Python]
def __init__(self) -> None
\end{lstlisting}

\textbf{可用性与安全性}

\textbf{\passthrough{\lstinline!AVAILABLE!} /
\passthrough{\lstinline!CONSTRUCTION!}}：当前绑定源码已实现；实际执行条件仍取决于平台和对应
SDK 环境。

\textbf{参数}

无显式参数。实例方法中的 \passthrough{\lstinline!self!} 由 Python
自动传入。

\textbf{返回值}

\begin{longtable}[]{@{}
  >{\raggedright\arraybackslash}p{(\columnwidth - 4\tabcolsep) * \real{0.18}}
  >{\raggedright\arraybackslash}p{(\columnwidth - 4\tabcolsep) * \real{0.20}}
  >{\raggedright\arraybackslash}p{(\columnwidth - 4\tabcolsep) * \real{0.62}}@{}}
\toprule
位置 & 类型 & 含义 \\
\midrule
\endhead
无 & \passthrough{\lstinline!None!} &
\passthrough{\lstinline!\_\_init\_\_!}
本身不返回值；调用类对象时，在构造函数可用的前提下得到该类实例。 \\
\bottomrule
\end{longtable}

\textbf{对应 C++}

\begin{itemize}
\tightlist
\item
  类：\passthrough{\lstinline!unitree::robot::go2::ServiceState!}
\item
  签名：\passthrough{\lstinline!ServiceState()!}
\item
  绑定策略：\passthrough{\lstinline!未单独标注!}
\item
  声明位置：\passthrough{\lstinline!include/unitree/robot/go2/robot\_state/robot\_state\_client.hpp:18!}
\end{itemize}

\textbf{说明}

\begin{itemize}
\tightlist
\item
  示例是局部调用片段；\passthrough{\lstinline!obj!}
  和参数变量表示已经按上层业务校验并准备好的对象。
\end{itemize}

\textbf{用法}

\begin{lstlisting}[language=Python]
value = ServiceState()
\end{lstlisting}

\hypertarget{unitree-sdk2-cpp-robot-go2-servicestatedata}{%
\subsubsection{\texorpdfstring{\texttt{unitree\_sdk2\_cpp.robot.go2.ServiceStateData}}{unitree\_sdk2\_cpp.robot.go2.ServiceStateData}}\label{unitree-sdk2-cpp-robot-go2-servicestatedata}}

SDK 数据值类型；公开字段和转换方法见下文。

\textbf{导入}

\begin{lstlisting}[language=Python]
from unitree_sdk2_cpp.robot.go2 import ServiceStateData
\end{lstlisting}

\textbf{基类}：\passthrough{\lstinline!object!}

\textbf{公开属性}

\begin{longtable}[]{@{}
  >{\raggedright\arraybackslash}p{(\columnwidth - 6\tabcolsep) * \real{0.18}}
  >{\raggedright\arraybackslash}p{(\columnwidth - 6\tabcolsep) * \real{0.24}}
  >{\raggedright\arraybackslash}p{(\columnwidth - 6\tabcolsep) * \real{0.18}}
  >{\raggedright\arraybackslash}p{(\columnwidth - 6\tabcolsep) * \real{0.40}}@{}}
\toprule
名称 & Python 类型 & 含义 & 用法 \\
\midrule
\endhead
\passthrough{\lstinline!name!} & \passthrough{\lstinline!str!} & SDK
对象、配置项、服务或动作的名称。允许值由调用它的具体接口定义。 &
\passthrough{\lstinline!value = obj.name!} /
\passthrough{\lstinline!obj.name = value!} \\
\passthrough{\lstinline!status!} & \passthrough{\lstinline!int!} &
传给该接口的 状态 参数，Python 类型为
\passthrough{\lstinline!int!}。精确含义、单位、范围和枚举值需查目标型号对应头文件与协议。
& \passthrough{\lstinline!value = obj.status!} /
\passthrough{\lstinline!obj.status = value!} \\
\passthrough{\lstinline!protect!} & \passthrough{\lstinline!int!} &
传给该接口的 \passthrough{\lstinline!protect!} 参数，Python 类型为
\passthrough{\lstinline!int!}。精确含义、单位、范围和枚举值需查目标型号对应头文件与协议。
& \passthrough{\lstinline!value = obj.protect!} /
\passthrough{\lstinline!obj.protect = value!} \\
\bottomrule
\end{longtable}

\textbf{方法索引}

\begin{longtable}[]{@{}
  >{\raggedright\arraybackslash}p{(\columnwidth - 6\tabcolsep) * \real{0.18}}
  >{\raggedright\arraybackslash}p{(\columnwidth - 6\tabcolsep) * \real{0.24}}
  >{\raggedright\arraybackslash}p{(\columnwidth - 6\tabcolsep) * \real{0.18}}
  >{\raggedright\arraybackslash}p{(\columnwidth - 6\tabcolsep) * \real{0.40}}@{}}
\toprule
方法 & Python 签名 & 状态 & 安全分类 \\
\midrule
\endhead
\protect\hyperlink{unitree-sdk2-cpp-robot-go2-servicestatedata-dunder-init-1}{\passthrough{\lstinline!\_\_init\_\_!}}
& \passthrough{\lstinline!def \_\_init\_\_(self) -> None!} &
\passthrough{\lstinline!SIGNATURE\_ONLY!} &
\passthrough{\lstinline!CONSTRUCTION!} \\
\protect\hyperlink{unitree-sdk2-cpp-robot-go2-servicestatedata-from-json-1}{\passthrough{\lstinline!from\_json!}}
&
\passthrough{\lstinline!def from\_json(self, json: dict[str, Any]) -> None!}
& \passthrough{\lstinline!SIGNATURE\_ONLY!} &
\passthrough{\lstinline!UNCLASSIFIED!} \\
\protect\hyperlink{unitree-sdk2-cpp-robot-go2-servicestatedata-to-json-1}{\passthrough{\lstinline!to\_json!}}
&
\passthrough{\lstinline!def to\_json(self, json: dict[str, Any]) -> None!}
& \passthrough{\lstinline!SIGNATURE\_ONLY!} &
\passthrough{\lstinline!UNCLASSIFIED!} \\
\bottomrule
\end{longtable}

\hypertarget{unitree-sdk2-cpp-robot-go2-servicestatedata-dunder-init-1}{%
\paragraph{\texorpdfstring{\texttt{unitree\_sdk2\_cpp.robot.go2.ServiceStateData.\_\_init\_\_}}{unitree\_sdk2\_cpp.robot.go2.ServiceStateData.\_\_init\_\_}}\label{unitree-sdk2-cpp-robot-go2-servicestatedata-dunder-init-1}}

初始化 \passthrough{\lstinline!ServiceStateData!}
实例。是否能实际构造取决于下方可用性状态。

\textbf{签名}

\begin{lstlisting}[language=Python]
def __init__(self) -> None
\end{lstlisting}

\textbf{可用性与安全性}

\textbf{\passthrough{\lstinline!SIGNATURE\_ONLY!} /
\passthrough{\lstinline!CONSTRUCTION!}}：当前只有设计期签名，不能假设已存在于运行时扩展。

\textbf{参数}

无显式参数。实例方法中的 \passthrough{\lstinline!self!} 由 Python
自动传入。

\textbf{返回值}

\begin{longtable}[]{@{}
  >{\raggedright\arraybackslash}p{(\columnwidth - 4\tabcolsep) * \real{0.18}}
  >{\raggedright\arraybackslash}p{(\columnwidth - 4\tabcolsep) * \real{0.20}}
  >{\raggedright\arraybackslash}p{(\columnwidth - 4\tabcolsep) * \real{0.62}}@{}}
\toprule
位置 & 类型 & 含义 \\
\midrule
\endhead
无 & \passthrough{\lstinline!None!} &
\passthrough{\lstinline!\_\_init\_\_!}
本身不返回值；调用类对象时，在构造函数可用的前提下得到该类实例。 \\
\bottomrule
\end{longtable}

\textbf{对应 C++}

\begin{itemize}
\tightlist
\item
  类：\passthrough{\lstinline!unitree::robot::go2::ServiceStateData!}
\item
  签名：\passthrough{\lstinline!ServiceStateData()!}
\item
  绑定策略：\passthrough{\lstinline!未单独标注!}
\item
  声明位置：\passthrough{\lstinline!include/unitree/robot/go2/robot\_state/robot\_state\_api.hpp:122!}
\end{itemize}

\textbf{说明}

\begin{itemize}
\tightlist
\item
  示例是局部调用片段；\passthrough{\lstinline!obj!}
  和参数变量表示已经按上层业务校验并准备好的对象。
\item
  该条目可以被编辑器和类型检查器识别，但当前运行时可能在属性查找阶段直接失败。
\end{itemize}

\textbf{用法}

\begin{lstlisting}[language=Python]
from typing import TYPE_CHECKING

if TYPE_CHECKING:
    def planned_construct() -> ServiceStateData:
        return ServiceStateData()
\end{lstlisting}

\begin{quote}
{[}!CAUTION{]} 上面的 \passthrough{\lstinline!TYPE\_CHECKING!}
示例只表达计划调用形态。该分支在普通 Python
运行时为假，不代表当前扩展能执行此接口。
\end{quote}

\hypertarget{unitree-sdk2-cpp-robot-go2-servicestatedata-from-json-1}{%
\paragraph{\texorpdfstring{\texttt{unitree\_sdk2\_cpp.robot.go2.ServiceStateData.from\_json}}{unitree\_sdk2\_cpp.robot.go2.ServiceStateData.from\_json}}\label{unitree-sdk2-cpp-robot-go2-servicestatedata-from-json-1}}

计划从 JSON 风格字典读取字段并更新当前 SDK 值对象。

\textbf{签名}

\begin{lstlisting}[language=Python]
def from_json(self, json: dict[str, Any]) -> None
\end{lstlisting}

\textbf{可用性与安全性}

\textbf{\passthrough{\lstinline!SIGNATURE\_ONLY!} /
\passthrough{\lstinline!UNCLASSIFIED!}}：当前只有设计期签名，不能假设已存在于运行时扩展。

\textbf{参数}

\begin{longtable}[]{@{}
  >{\raggedright\arraybackslash}p{(\columnwidth - 6\tabcolsep) * \real{0.18}}
  >{\raggedright\arraybackslash}p{(\columnwidth - 6\tabcolsep) * \real{0.24}}
  >{\raggedright\arraybackslash}p{(\columnwidth - 6\tabcolsep) * \real{0.18}}
  >{\raggedright\arraybackslash}p{(\columnwidth - 6\tabcolsep) * \real{0.40}}@{}}
\toprule
名称 & Python 类型 & 默认值 & 含义 \\
\midrule
\endhead
\passthrough{\lstinline!json!} &
\passthrough{\lstinline!dict[str, Any]!} & 必填 & JSON
风格字典，键和值必须满足对应 C++ \passthrough{\lstinline!JsonMap!}
数据结构的约束。 对应 C++ 参数
\passthrough{\lstinline!json: common::JsonMap \&!}。 \\
\bottomrule
\end{longtable}

\textbf{返回值}

\begin{longtable}[]{@{}
  >{\raggedright\arraybackslash}p{(\columnwidth - 4\tabcolsep) * \real{0.18}}
  >{\raggedright\arraybackslash}p{(\columnwidth - 4\tabcolsep) * \real{0.20}}
  >{\raggedright\arraybackslash}p{(\columnwidth - 4\tabcolsep) * \real{0.62}}@{}}
\toprule
位置 & 类型 & 含义 \\
\midrule
\endhead
无 & \passthrough{\lstinline!None!} & 该方法不返回 Python
值；失败可能表现为异常或底层状态变化。 \\
\bottomrule
\end{longtable}

\textbf{对应 C++}

\begin{itemize}
\tightlist
\item
  类：\passthrough{\lstinline!unitree::robot::go2::ServiceStateData!}
\item
  签名：\passthrough{\lstinline!fromJson(common::JsonMap \&)!}
\item
  绑定策略：\passthrough{\lstinline!SIGNATURE\_PREVIEW!}
\item
  声明位置：\passthrough{\lstinline!include/unitree/robot/go2/robot\_state/robot\_state\_api.hpp:128!}
\end{itemize}

\textbf{说明}

\begin{itemize}
\tightlist
\item
  示例是局部调用片段；\passthrough{\lstinline!obj!}
  和参数变量表示已经按上层业务校验并准备好的对象。
\item
  该条目可以被编辑器和类型检查器识别，但当前运行时可能在属性查找阶段直接失败。
\end{itemize}

\textbf{用法}

\begin{lstlisting}[language=Python]
from typing import TYPE_CHECKING

if TYPE_CHECKING:
    def planned_from_json(obj: ServiceStateData, json: dict[str, Any]) -> None:
        obj.from_json(json=json)
\end{lstlisting}

\begin{quote}
{[}!CAUTION{]} 上面的 \passthrough{\lstinline!TYPE\_CHECKING!}
示例只表达计划调用形态。该分支在普通 Python
运行时为假，不代表当前扩展能执行此接口。
\end{quote}

\hypertarget{unitree-sdk2-cpp-robot-go2-servicestatedata-to-json-1}{%
\paragraph{\texorpdfstring{\texttt{unitree\_sdk2\_cpp.robot.go2.ServiceStateData.to\_json}}{unitree\_sdk2\_cpp.robot.go2.ServiceStateData.to\_json}}\label{unitree-sdk2-cpp-robot-go2-servicestatedata-to-json-1}}

计划把当前 SDK 值对象写入 JSON 风格字典。

\textbf{签名}

\begin{lstlisting}[language=Python]
def to_json(self, json: dict[str, Any]) -> None
\end{lstlisting}

\textbf{可用性与安全性}

\textbf{\passthrough{\lstinline!SIGNATURE\_ONLY!} /
\passthrough{\lstinline!UNCLASSIFIED!}}：当前只有设计期签名，不能假设已存在于运行时扩展。

\textbf{参数}

\begin{longtable}[]{@{}
  >{\raggedright\arraybackslash}p{(\columnwidth - 6\tabcolsep) * \real{0.18}}
  >{\raggedright\arraybackslash}p{(\columnwidth - 6\tabcolsep) * \real{0.24}}
  >{\raggedright\arraybackslash}p{(\columnwidth - 6\tabcolsep) * \real{0.18}}
  >{\raggedright\arraybackslash}p{(\columnwidth - 6\tabcolsep) * \real{0.40}}@{}}
\toprule
名称 & Python 类型 & 默认值 & 含义 \\
\midrule
\endhead
\passthrough{\lstinline!json!} &
\passthrough{\lstinline!dict[str, Any]!} & 必填 & JSON
风格字典，键和值必须满足对应 C++ \passthrough{\lstinline!JsonMap!}
数据结构的约束。 对应 C++ 参数
\passthrough{\lstinline!json: common::JsonMap \&!}。 \\
\bottomrule
\end{longtable}

\textbf{返回值}

\begin{longtable}[]{@{}
  >{\raggedright\arraybackslash}p{(\columnwidth - 4\tabcolsep) * \real{0.18}}
  >{\raggedright\arraybackslash}p{(\columnwidth - 4\tabcolsep) * \real{0.20}}
  >{\raggedright\arraybackslash}p{(\columnwidth - 4\tabcolsep) * \real{0.62}}@{}}
\toprule
位置 & 类型 & 含义 \\
\midrule
\endhead
无 & \passthrough{\lstinline!None!} & 该方法不返回 Python
值；失败可能表现为异常或底层状态变化。 \\
\bottomrule
\end{longtable}

\textbf{对应 C++}

\begin{itemize}
\tightlist
\item
  类：\passthrough{\lstinline!unitree::robot::go2::ServiceStateData!}
\item
  签名：\passthrough{\lstinline!toJson(common::JsonMap \&) const!}
\item
  绑定策略：\passthrough{\lstinline!SIGNATURE\_PREVIEW!}
\item
  声明位置：\passthrough{\lstinline!include/unitree/robot/go2/robot\_state/robot\_state\_api.hpp:135!}
\end{itemize}

\textbf{说明}

\begin{itemize}
\tightlist
\item
  示例是局部调用片段；\passthrough{\lstinline!obj!}
  和参数变量表示已经按上层业务校验并准备好的对象。
\item
  该条目可以被编辑器和类型检查器识别，但当前运行时可能在属性查找阶段直接失败。
\end{itemize}

\textbf{用法}

\begin{lstlisting}[language=Python]
from typing import TYPE_CHECKING

if TYPE_CHECKING:
    def planned_to_json(obj: ServiceStateData, json: dict[str, Any]) -> None:
        obj.to_json(json=json)
\end{lstlisting}

\begin{quote}
{[}!CAUTION{]} 上面的 \passthrough{\lstinline!TYPE\_CHECKING!}
示例只表达计划调用形态。该分支在普通 Python
运行时为假，不代表当前扩展能执行此接口。
\end{quote}

\hypertarget{unitree-sdk2-cpp-robot-go2-serviceswitchdata}{%
\subsubsection{\texorpdfstring{\texttt{unitree\_sdk2\_cpp.robot.go2.ServiceSwitchData}}{unitree\_sdk2\_cpp.robot.go2.ServiceSwitchData}}\label{unitree-sdk2-cpp-robot-go2-serviceswitchdata}}

SDK 数据值类型；公开字段和转换方法见下文。

\textbf{导入}

\begin{lstlisting}[language=Python]
from unitree_sdk2_cpp.robot.go2 import ServiceSwitchData
\end{lstlisting}

\textbf{基类}：\passthrough{\lstinline!object!}

\textbf{公开属性}

\begin{longtable}[]{@{}
  >{\raggedright\arraybackslash}p{(\columnwidth - 6\tabcolsep) * \real{0.18}}
  >{\raggedright\arraybackslash}p{(\columnwidth - 6\tabcolsep) * \real{0.24}}
  >{\raggedright\arraybackslash}p{(\columnwidth - 6\tabcolsep) * \real{0.18}}
  >{\raggedright\arraybackslash}p{(\columnwidth - 6\tabcolsep) * \real{0.40}}@{}}
\toprule
名称 & Python 类型 & 含义 & 用法 \\
\midrule
\endhead
\passthrough{\lstinline!name!} & \passthrough{\lstinline!str!} & SDK
对象、配置项、服务或动作的名称。允许值由调用它的具体接口定义。 &
\passthrough{\lstinline!value = obj.name!} /
\passthrough{\lstinline!obj.name = value!} \\
\passthrough{\lstinline!status!} & \passthrough{\lstinline!int!} &
传给该接口的 状态 参数，Python 类型为
\passthrough{\lstinline!int!}。精确含义、单位、范围和枚举值需查目标型号对应头文件与协议。
& \passthrough{\lstinline!value = obj.status!} /
\passthrough{\lstinline!obj.status = value!} \\
\bottomrule
\end{longtable}

\textbf{方法索引}

\begin{longtable}[]{@{}
  >{\raggedright\arraybackslash}p{(\columnwidth - 6\tabcolsep) * \real{0.18}}
  >{\raggedright\arraybackslash}p{(\columnwidth - 6\tabcolsep) * \real{0.24}}
  >{\raggedright\arraybackslash}p{(\columnwidth - 6\tabcolsep) * \real{0.18}}
  >{\raggedright\arraybackslash}p{(\columnwidth - 6\tabcolsep) * \real{0.40}}@{}}
\toprule
方法 & Python 签名 & 状态 & 安全分类 \\
\midrule
\endhead
\protect\hyperlink{unitree-sdk2-cpp-robot-go2-serviceswitchdata-dunder-init-1}{\passthrough{\lstinline!\_\_init\_\_!}}
& \passthrough{\lstinline!def \_\_init\_\_(self) -> None!} &
\passthrough{\lstinline!SIGNATURE\_ONLY!} &
\passthrough{\lstinline!CONSTRUCTION!} \\
\protect\hyperlink{unitree-sdk2-cpp-robot-go2-serviceswitchdata-from-json-1}{\passthrough{\lstinline!from\_json!}}
&
\passthrough{\lstinline!def from\_json(self, json: dict[str, Any]) -> None!}
& \passthrough{\lstinline!SIGNATURE\_ONLY!} &
\passthrough{\lstinline!UNCLASSIFIED!} \\
\protect\hyperlink{unitree-sdk2-cpp-robot-go2-serviceswitchdata-to-json-1}{\passthrough{\lstinline!to\_json!}}
&
\passthrough{\lstinline!def to\_json(self, json: dict[str, Any]) -> None!}
& \passthrough{\lstinline!SIGNATURE\_ONLY!} &
\passthrough{\lstinline!UNCLASSIFIED!} \\
\bottomrule
\end{longtable}

\hypertarget{unitree-sdk2-cpp-robot-go2-serviceswitchdata-dunder-init-1}{%
\paragraph{\texorpdfstring{\texttt{unitree\_sdk2\_cpp.robot.go2.ServiceSwitchData.\_\_init\_\_}}{unitree\_sdk2\_cpp.robot.go2.ServiceSwitchData.\_\_init\_\_}}\label{unitree-sdk2-cpp-robot-go2-serviceswitchdata-dunder-init-1}}

初始化 \passthrough{\lstinline!ServiceSwitchData!}
实例。是否能实际构造取决于下方可用性状态。

\textbf{签名}

\begin{lstlisting}[language=Python]
def __init__(self) -> None
\end{lstlisting}

\textbf{可用性与安全性}

\textbf{\passthrough{\lstinline!SIGNATURE\_ONLY!} /
\passthrough{\lstinline!CONSTRUCTION!}}：当前只有设计期签名，不能假设已存在于运行时扩展。

\textbf{参数}

无显式参数。实例方法中的 \passthrough{\lstinline!self!} 由 Python
自动传入。

\textbf{返回值}

\begin{longtable}[]{@{}
  >{\raggedright\arraybackslash}p{(\columnwidth - 4\tabcolsep) * \real{0.18}}
  >{\raggedright\arraybackslash}p{(\columnwidth - 4\tabcolsep) * \real{0.20}}
  >{\raggedright\arraybackslash}p{(\columnwidth - 4\tabcolsep) * \real{0.62}}@{}}
\toprule
位置 & 类型 & 含义 \\
\midrule
\endhead
无 & \passthrough{\lstinline!None!} &
\passthrough{\lstinline!\_\_init\_\_!}
本身不返回值；调用类对象时，在构造函数可用的前提下得到该类实例。 \\
\bottomrule
\end{longtable}

\textbf{对应 C++}

\begin{itemize}
\tightlist
\item
  类：\passthrough{\lstinline!unitree::robot::go2::ServiceSwitchData!}
\item
  签名：\passthrough{\lstinline!ServiceSwitchData()!}
\item
  绑定策略：\passthrough{\lstinline!未单独标注!}
\item
  声明位置：\passthrough{\lstinline!include/unitree/robot/go2/robot\_state/robot\_state\_api.hpp:64!}
\end{itemize}

\textbf{说明}

\begin{itemize}
\tightlist
\item
  示例是局部调用片段；\passthrough{\lstinline!obj!}
  和参数变量表示已经按上层业务校验并准备好的对象。
\item
  该条目可以被编辑器和类型检查器识别，但当前运行时可能在属性查找阶段直接失败。
\end{itemize}

\textbf{用法}

\begin{lstlisting}[language=Python]
from typing import TYPE_CHECKING

if TYPE_CHECKING:
    def planned_construct() -> ServiceSwitchData:
        return ServiceSwitchData()
\end{lstlisting}

\begin{quote}
{[}!CAUTION{]} 上面的 \passthrough{\lstinline!TYPE\_CHECKING!}
示例只表达计划调用形态。该分支在普通 Python
运行时为假，不代表当前扩展能执行此接口。
\end{quote}

\hypertarget{unitree-sdk2-cpp-robot-go2-serviceswitchdata-from-json-1}{%
\paragraph{\texorpdfstring{\texttt{unitree\_sdk2\_cpp.robot.go2.ServiceSwitchData.from\_json}}{unitree\_sdk2\_cpp.robot.go2.ServiceSwitchData.from\_json}}\label{unitree-sdk2-cpp-robot-go2-serviceswitchdata-from-json-1}}

计划从 JSON 风格字典读取字段并更新当前 SDK 值对象。

\textbf{签名}

\begin{lstlisting}[language=Python]
def from_json(self, json: dict[str, Any]) -> None
\end{lstlisting}

\textbf{可用性与安全性}

\textbf{\passthrough{\lstinline!SIGNATURE\_ONLY!} /
\passthrough{\lstinline!UNCLASSIFIED!}}：当前只有设计期签名，不能假设已存在于运行时扩展。

\textbf{参数}

\begin{longtable}[]{@{}
  >{\raggedright\arraybackslash}p{(\columnwidth - 6\tabcolsep) * \real{0.18}}
  >{\raggedright\arraybackslash}p{(\columnwidth - 6\tabcolsep) * \real{0.24}}
  >{\raggedright\arraybackslash}p{(\columnwidth - 6\tabcolsep) * \real{0.18}}
  >{\raggedright\arraybackslash}p{(\columnwidth - 6\tabcolsep) * \real{0.40}}@{}}
\toprule
名称 & Python 类型 & 默认值 & 含义 \\
\midrule
\endhead
\passthrough{\lstinline!json!} &
\passthrough{\lstinline!dict[str, Any]!} & 必填 & JSON
风格字典，键和值必须满足对应 C++ \passthrough{\lstinline!JsonMap!}
数据结构的约束。 对应 C++ 参数
\passthrough{\lstinline!json: common::JsonMap \&!}。 \\
\bottomrule
\end{longtable}

\textbf{返回值}

\begin{longtable}[]{@{}
  >{\raggedright\arraybackslash}p{(\columnwidth - 4\tabcolsep) * \real{0.18}}
  >{\raggedright\arraybackslash}p{(\columnwidth - 4\tabcolsep) * \real{0.20}}
  >{\raggedright\arraybackslash}p{(\columnwidth - 4\tabcolsep) * \real{0.62}}@{}}
\toprule
位置 & 类型 & 含义 \\
\midrule
\endhead
无 & \passthrough{\lstinline!None!} & 该方法不返回 Python
值；失败可能表现为异常或底层状态变化。 \\
\bottomrule
\end{longtable}

\textbf{对应 C++}

\begin{itemize}
\tightlist
\item
  类：\passthrough{\lstinline!unitree::robot::go2::ServiceSwitchData!}
\item
  签名：\passthrough{\lstinline!fromJson(common::JsonMap \&)!}
\item
  绑定策略：\passthrough{\lstinline!SIGNATURE\_PREVIEW!}
\item
  声明位置：\passthrough{\lstinline!include/unitree/robot/go2/robot\_state/robot\_state\_api.hpp:70!}
\end{itemize}

\textbf{说明}

\begin{itemize}
\tightlist
\item
  示例是局部调用片段；\passthrough{\lstinline!obj!}
  和参数变量表示已经按上层业务校验并准备好的对象。
\item
  该条目可以被编辑器和类型检查器识别，但当前运行时可能在属性查找阶段直接失败。
\end{itemize}

\textbf{用法}

\begin{lstlisting}[language=Python]
from typing import TYPE_CHECKING

if TYPE_CHECKING:
    def planned_from_json(obj: ServiceSwitchData, json: dict[str, Any]) -> None:
        obj.from_json(json=json)
\end{lstlisting}

\begin{quote}
{[}!CAUTION{]} 上面的 \passthrough{\lstinline!TYPE\_CHECKING!}
示例只表达计划调用形态。该分支在普通 Python
运行时为假，不代表当前扩展能执行此接口。
\end{quote}

\hypertarget{unitree-sdk2-cpp-robot-go2-serviceswitchdata-to-json-1}{%
\paragraph{\texorpdfstring{\texttt{unitree\_sdk2\_cpp.robot.go2.ServiceSwitchData.to\_json}}{unitree\_sdk2\_cpp.robot.go2.ServiceSwitchData.to\_json}}\label{unitree-sdk2-cpp-robot-go2-serviceswitchdata-to-json-1}}

计划把当前 SDK 值对象写入 JSON 风格字典。

\textbf{签名}

\begin{lstlisting}[language=Python]
def to_json(self, json: dict[str, Any]) -> None
\end{lstlisting}

\textbf{可用性与安全性}

\textbf{\passthrough{\lstinline!SIGNATURE\_ONLY!} /
\passthrough{\lstinline!UNCLASSIFIED!}}：当前只有设计期签名，不能假设已存在于运行时扩展。

\textbf{参数}

\begin{longtable}[]{@{}
  >{\raggedright\arraybackslash}p{(\columnwidth - 6\tabcolsep) * \real{0.18}}
  >{\raggedright\arraybackslash}p{(\columnwidth - 6\tabcolsep) * \real{0.24}}
  >{\raggedright\arraybackslash}p{(\columnwidth - 6\tabcolsep) * \real{0.18}}
  >{\raggedright\arraybackslash}p{(\columnwidth - 6\tabcolsep) * \real{0.40}}@{}}
\toprule
名称 & Python 类型 & 默认值 & 含义 \\
\midrule
\endhead
\passthrough{\lstinline!json!} &
\passthrough{\lstinline!dict[str, Any]!} & 必填 & JSON
风格字典，键和值必须满足对应 C++ \passthrough{\lstinline!JsonMap!}
数据结构的约束。 对应 C++ 参数
\passthrough{\lstinline!json: common::JsonMap \&!}。 \\
\bottomrule
\end{longtable}

\textbf{返回值}

\begin{longtable}[]{@{}
  >{\raggedright\arraybackslash}p{(\columnwidth - 4\tabcolsep) * \real{0.18}}
  >{\raggedright\arraybackslash}p{(\columnwidth - 4\tabcolsep) * \real{0.20}}
  >{\raggedright\arraybackslash}p{(\columnwidth - 4\tabcolsep) * \real{0.62}}@{}}
\toprule
位置 & 类型 & 含义 \\
\midrule
\endhead
无 & \passthrough{\lstinline!None!} & 该方法不返回 Python
值；失败可能表现为异常或底层状态变化。 \\
\bottomrule
\end{longtable}

\textbf{对应 C++}

\begin{itemize}
\tightlist
\item
  类：\passthrough{\lstinline!unitree::robot::go2::ServiceSwitchData!}
\item
  签名：\passthrough{\lstinline!toJson(common::JsonMap \&) const!}
\item
  绑定策略：\passthrough{\lstinline!SIGNATURE\_PREVIEW!}
\item
  声明位置：\passthrough{\lstinline!include/unitree/robot/go2/robot\_state/robot\_state\_api.hpp:76!}
\end{itemize}

\textbf{说明}

\begin{itemize}
\tightlist
\item
  示例是局部调用片段；\passthrough{\lstinline!obj!}
  和参数变量表示已经按上层业务校验并准备好的对象。
\item
  该条目可以被编辑器和类型检查器识别，但当前运行时可能在属性查找阶段直接失败。
\end{itemize}

\textbf{用法}

\begin{lstlisting}[language=Python]
from typing import TYPE_CHECKING

if TYPE_CHECKING:
    def planned_to_json(obj: ServiceSwitchData, json: dict[str, Any]) -> None:
        obj.to_json(json=json)
\end{lstlisting}

\begin{quote}
{[}!CAUTION{]} 上面的 \passthrough{\lstinline!TYPE\_CHECKING!}
示例只表达计划调用形态。该分支在普通 Python
运行时为假，不代表当前扩展能执行此接口。
\end{quote}

\hypertarget{unitree-sdk2-cpp-robot-go2-serviceswitchparameter}{%
\subsubsection{\texorpdfstring{\texttt{unitree\_sdk2\_cpp.robot.go2.ServiceSwitchParameter}}{unitree\_sdk2\_cpp.robot.go2.ServiceSwitchParameter}}\label{unitree-sdk2-cpp-robot-go2-serviceswitchparameter}}

SDK 请求参数值类型；当前可能仅提供设计期签名。

\textbf{导入}

\begin{lstlisting}[language=Python]
from unitree_sdk2_cpp.robot.go2 import ServiceSwitchParameter
\end{lstlisting}

\textbf{基类}：\passthrough{\lstinline!object!}

\textbf{公开属性}

\begin{longtable}[]{@{}
  >{\raggedright\arraybackslash}p{(\columnwidth - 6\tabcolsep) * \real{0.18}}
  >{\raggedright\arraybackslash}p{(\columnwidth - 6\tabcolsep) * \real{0.24}}
  >{\raggedright\arraybackslash}p{(\columnwidth - 6\tabcolsep) * \real{0.18}}
  >{\raggedright\arraybackslash}p{(\columnwidth - 6\tabcolsep) * \real{0.40}}@{}}
\toprule
名称 & Python 类型 & 含义 & 用法 \\
\midrule
\endhead
\passthrough{\lstinline!name!} & \passthrough{\lstinline!str!} & SDK
对象、配置项、服务或动作的名称。允许值由调用它的具体接口定义。 &
\passthrough{\lstinline!value = obj.name!} /
\passthrough{\lstinline!obj.name = value!} \\
\passthrough{\lstinline!swit!} & \passthrough{\lstinline!int!} &
传给该接口的 \passthrough{\lstinline!swit!} 参数，Python 类型为
\passthrough{\lstinline!int!}。精确含义、单位、范围和枚举值需查目标型号对应头文件与协议。
& \passthrough{\lstinline!value = obj.swit!} /
\passthrough{\lstinline!obj.swit = value!} \\
\bottomrule
\end{longtable}

\textbf{方法索引}

\begin{longtable}[]{@{}
  >{\raggedright\arraybackslash}p{(\columnwidth - 6\tabcolsep) * \real{0.18}}
  >{\raggedright\arraybackslash}p{(\columnwidth - 6\tabcolsep) * \real{0.24}}
  >{\raggedright\arraybackslash}p{(\columnwidth - 6\tabcolsep) * \real{0.18}}
  >{\raggedright\arraybackslash}p{(\columnwidth - 6\tabcolsep) * \real{0.40}}@{}}
\toprule
方法 & Python 签名 & 状态 & 安全分类 \\
\midrule
\endhead
\protect\hyperlink{unitree-sdk2-cpp-robot-go2-serviceswitchparameter-dunder-init-1}{\passthrough{\lstinline!\_\_init\_\_!}}
& \passthrough{\lstinline!def \_\_init\_\_(self) -> None!} &
\passthrough{\lstinline!SIGNATURE\_ONLY!} &
\passthrough{\lstinline!CONSTRUCTION!} \\
\protect\hyperlink{unitree-sdk2-cpp-robot-go2-serviceswitchparameter-from-json-1}{\passthrough{\lstinline!from\_json!}}
&
\passthrough{\lstinline!def from\_json(self, json: dict[str, Any]) -> None!}
& \passthrough{\lstinline!SIGNATURE\_ONLY!} &
\passthrough{\lstinline!UNCLASSIFIED!} \\
\protect\hyperlink{unitree-sdk2-cpp-robot-go2-serviceswitchparameter-to-json-1}{\passthrough{\lstinline!to\_json!}}
&
\passthrough{\lstinline!def to\_json(self, json: dict[str, Any]) -> None!}
& \passthrough{\lstinline!SIGNATURE\_ONLY!} &
\passthrough{\lstinline!UNCLASSIFIED!} \\
\bottomrule
\end{longtable}

\hypertarget{unitree-sdk2-cpp-robot-go2-serviceswitchparameter-dunder-init-1}{%
\paragraph{\texorpdfstring{\texttt{unitree\_sdk2\_cpp.robot.go2.ServiceSwitchParameter.\_\_init\_\_}}{unitree\_sdk2\_cpp.robot.go2.ServiceSwitchParameter.\_\_init\_\_}}\label{unitree-sdk2-cpp-robot-go2-serviceswitchparameter-dunder-init-1}}

初始化 \passthrough{\lstinline!ServiceSwitchParameter!}
实例。是否能实际构造取决于下方可用性状态。

\textbf{签名}

\begin{lstlisting}[language=Python]
def __init__(self) -> None
\end{lstlisting}

\textbf{可用性与安全性}

\textbf{\passthrough{\lstinline!SIGNATURE\_ONLY!} /
\passthrough{\lstinline!CONSTRUCTION!}}：当前只有设计期签名，不能假设已存在于运行时扩展。

\textbf{参数}

无显式参数。实例方法中的 \passthrough{\lstinline!self!} 由 Python
自动传入。

\textbf{返回值}

\begin{longtable}[]{@{}
  >{\raggedright\arraybackslash}p{(\columnwidth - 4\tabcolsep) * \real{0.18}}
  >{\raggedright\arraybackslash}p{(\columnwidth - 4\tabcolsep) * \real{0.20}}
  >{\raggedright\arraybackslash}p{(\columnwidth - 4\tabcolsep) * \real{0.62}}@{}}
\toprule
位置 & 类型 & 含义 \\
\midrule
\endhead
无 & \passthrough{\lstinline!None!} &
\passthrough{\lstinline!\_\_init\_\_!}
本身不返回值；调用类对象时，在构造函数可用的前提下得到该类实例。 \\
\bottomrule
\end{longtable}

\textbf{对应 C++}

\begin{itemize}
\tightlist
\item
  类：\passthrough{\lstinline!unitree::robot::go2::ServiceSwitchParameter!}
\item
  签名：\passthrough{\lstinline!ServiceSwitchParameter()!}
\item
  绑定策略：\passthrough{\lstinline!未单独标注!}
\item
  声明位置：\passthrough{\lstinline!include/unitree/robot/go2/robot\_state/robot\_state\_api.hpp:35!}
\end{itemize}

\textbf{说明}

\begin{itemize}
\tightlist
\item
  示例是局部调用片段；\passthrough{\lstinline!obj!}
  和参数变量表示已经按上层业务校验并准备好的对象。
\item
  该条目可以被编辑器和类型检查器识别，但当前运行时可能在属性查找阶段直接失败。
\end{itemize}

\textbf{用法}

\begin{lstlisting}[language=Python]
from typing import TYPE_CHECKING

if TYPE_CHECKING:
    def planned_construct() -> ServiceSwitchParameter:
        return ServiceSwitchParameter()
\end{lstlisting}

\begin{quote}
{[}!CAUTION{]} 上面的 \passthrough{\lstinline!TYPE\_CHECKING!}
示例只表达计划调用形态。该分支在普通 Python
运行时为假，不代表当前扩展能执行此接口。
\end{quote}

\hypertarget{unitree-sdk2-cpp-robot-go2-serviceswitchparameter-from-json-1}{%
\paragraph{\texorpdfstring{\texttt{unitree\_sdk2\_cpp.robot.go2.ServiceSwitchParameter.from\_json}}{unitree\_sdk2\_cpp.robot.go2.ServiceSwitchParameter.from\_json}}\label{unitree-sdk2-cpp-robot-go2-serviceswitchparameter-from-json-1}}

计划从 JSON 风格字典读取字段并更新当前 SDK 值对象。

\textbf{签名}

\begin{lstlisting}[language=Python]
def from_json(self, json: dict[str, Any]) -> None
\end{lstlisting}

\textbf{可用性与安全性}

\textbf{\passthrough{\lstinline!SIGNATURE\_ONLY!} /
\passthrough{\lstinline!UNCLASSIFIED!}}：当前只有设计期签名，不能假设已存在于运行时扩展。

\textbf{参数}

\begin{longtable}[]{@{}
  >{\raggedright\arraybackslash}p{(\columnwidth - 6\tabcolsep) * \real{0.18}}
  >{\raggedright\arraybackslash}p{(\columnwidth - 6\tabcolsep) * \real{0.24}}
  >{\raggedright\arraybackslash}p{(\columnwidth - 6\tabcolsep) * \real{0.18}}
  >{\raggedright\arraybackslash}p{(\columnwidth - 6\tabcolsep) * \real{0.40}}@{}}
\toprule
名称 & Python 类型 & 默认值 & 含义 \\
\midrule
\endhead
\passthrough{\lstinline!json!} &
\passthrough{\lstinline!dict[str, Any]!} & 必填 & JSON
风格字典，键和值必须满足对应 C++ \passthrough{\lstinline!JsonMap!}
数据结构的约束。 对应 C++ 参数
\passthrough{\lstinline!json: common::JsonMap \&!}。 \\
\bottomrule
\end{longtable}

\textbf{返回值}

\begin{longtable}[]{@{}
  >{\raggedright\arraybackslash}p{(\columnwidth - 4\tabcolsep) * \real{0.18}}
  >{\raggedright\arraybackslash}p{(\columnwidth - 4\tabcolsep) * \real{0.20}}
  >{\raggedright\arraybackslash}p{(\columnwidth - 4\tabcolsep) * \real{0.62}}@{}}
\toprule
位置 & 类型 & 含义 \\
\midrule
\endhead
无 & \passthrough{\lstinline!None!} & 该方法不返回 Python
值；失败可能表现为异常或底层状态变化。 \\
\bottomrule
\end{longtable}

\textbf{对应 C++}

\begin{itemize}
\tightlist
\item
  类：\passthrough{\lstinline!unitree::robot::go2::ServiceSwitchParameter!}
\item
  签名：\passthrough{\lstinline!fromJson(common::JsonMap \&)!}
\item
  绑定策略：\passthrough{\lstinline!SIGNATURE\_PREVIEW!}
\item
  声明位置：\passthrough{\lstinline!include/unitree/robot/go2/robot\_state/robot\_state\_api.hpp:41!}
\end{itemize}

\textbf{说明}

\begin{itemize}
\tightlist
\item
  示例是局部调用片段；\passthrough{\lstinline!obj!}
  和参数变量表示已经按上层业务校验并准备好的对象。
\item
  该条目可以被编辑器和类型检查器识别，但当前运行时可能在属性查找阶段直接失败。
\end{itemize}

\textbf{用法}

\begin{lstlisting}[language=Python]
from typing import TYPE_CHECKING

if TYPE_CHECKING:
    def planned_from_json(obj: ServiceSwitchParameter, json: dict[str, Any]) -> None:
        obj.from_json(json=json)
\end{lstlisting}

\begin{quote}
{[}!CAUTION{]} 上面的 \passthrough{\lstinline!TYPE\_CHECKING!}
示例只表达计划调用形态。该分支在普通 Python
运行时为假，不代表当前扩展能执行此接口。
\end{quote}

\hypertarget{unitree-sdk2-cpp-robot-go2-serviceswitchparameter-to-json-1}{%
\paragraph{\texorpdfstring{\texttt{unitree\_sdk2\_cpp.robot.go2.ServiceSwitchParameter.to\_json}}{unitree\_sdk2\_cpp.robot.go2.ServiceSwitchParameter.to\_json}}\label{unitree-sdk2-cpp-robot-go2-serviceswitchparameter-to-json-1}}

计划把当前 SDK 值对象写入 JSON 风格字典。

\textbf{签名}

\begin{lstlisting}[language=Python]
def to_json(self, json: dict[str, Any]) -> None
\end{lstlisting}

\textbf{可用性与安全性}

\textbf{\passthrough{\lstinline!SIGNATURE\_ONLY!} /
\passthrough{\lstinline!UNCLASSIFIED!}}：当前只有设计期签名，不能假设已存在于运行时扩展。

\textbf{参数}

\begin{longtable}[]{@{}
  >{\raggedright\arraybackslash}p{(\columnwidth - 6\tabcolsep) * \real{0.18}}
  >{\raggedright\arraybackslash}p{(\columnwidth - 6\tabcolsep) * \real{0.24}}
  >{\raggedright\arraybackslash}p{(\columnwidth - 6\tabcolsep) * \real{0.18}}
  >{\raggedright\arraybackslash}p{(\columnwidth - 6\tabcolsep) * \real{0.40}}@{}}
\toprule
名称 & Python 类型 & 默认值 & 含义 \\
\midrule
\endhead
\passthrough{\lstinline!json!} &
\passthrough{\lstinline!dict[str, Any]!} & 必填 & JSON
风格字典，键和值必须满足对应 C++ \passthrough{\lstinline!JsonMap!}
数据结构的约束。 对应 C++ 参数
\passthrough{\lstinline!json: common::JsonMap \&!}。 \\
\bottomrule
\end{longtable}

\textbf{返回值}

\begin{longtable}[]{@{}
  >{\raggedright\arraybackslash}p{(\columnwidth - 4\tabcolsep) * \real{0.18}}
  >{\raggedright\arraybackslash}p{(\columnwidth - 4\tabcolsep) * \real{0.20}}
  >{\raggedright\arraybackslash}p{(\columnwidth - 4\tabcolsep) * \real{0.62}}@{}}
\toprule
位置 & 类型 & 含义 \\
\midrule
\endhead
无 & \passthrough{\lstinline!None!} & 该方法不返回 Python
值；失败可能表现为异常或底层状态变化。 \\
\bottomrule
\end{longtable}

\textbf{对应 C++}

\begin{itemize}
\tightlist
\item
  类：\passthrough{\lstinline!unitree::robot::go2::ServiceSwitchParameter!}
\item
  签名：\passthrough{\lstinline!toJson(common::JsonMap \&) const!}
\item
  绑定策略：\passthrough{\lstinline!SIGNATURE\_PREVIEW!}
\item
  声明位置：\passthrough{\lstinline!include/unitree/robot/go2/robot\_state/robot\_state\_api.hpp:47!}
\end{itemize}

\textbf{说明}

\begin{itemize}
\tightlist
\item
  示例是局部调用片段；\passthrough{\lstinline!obj!}
  和参数变量表示已经按上层业务校验并准备好的对象。
\item
  该条目可以被编辑器和类型检查器识别，但当前运行时可能在属性查找阶段直接失败。
\end{itemize}

\textbf{用法}

\begin{lstlisting}[language=Python]
from typing import TYPE_CHECKING

if TYPE_CHECKING:
    def planned_to_json(obj: ServiceSwitchParameter, json: dict[str, Any]) -> None:
        obj.to_json(json=json)
\end{lstlisting}

\begin{quote}
{[}!CAUTION{]} 上面的 \passthrough{\lstinline!TYPE\_CHECKING!}
示例只表达计划调用形态。该分支在普通 Python
运行时为假，不代表当前扩展能执行此接口。
\end{quote}

\hypertarget{unitree-sdk2-cpp-robot-go2-setreportfreqparameter}{%
\subsubsection{\texorpdfstring{\texttt{unitree\_sdk2\_cpp.robot.go2.SetReportFreqParameter}}{unitree\_sdk2\_cpp.robot.go2.SetReportFreqParameter}}\label{unitree-sdk2-cpp-robot-go2-setreportfreqparameter}}

SDK 请求参数值类型；当前可能仅提供设计期签名。

\textbf{导入}

\begin{lstlisting}[language=Python]
from unitree_sdk2_cpp.robot.go2 import SetReportFreqParameter
\end{lstlisting}

\textbf{基类}：\passthrough{\lstinline!object!}

\textbf{公开属性}

\begin{longtable}[]{@{}
  >{\raggedright\arraybackslash}p{(\columnwidth - 6\tabcolsep) * \real{0.18}}
  >{\raggedright\arraybackslash}p{(\columnwidth - 6\tabcolsep) * \real{0.24}}
  >{\raggedright\arraybackslash}p{(\columnwidth - 6\tabcolsep) * \real{0.18}}
  >{\raggedright\arraybackslash}p{(\columnwidth - 6\tabcolsep) * \real{0.40}}@{}}
\toprule
名称 & Python 类型 & 含义 & 用法 \\
\midrule
\endhead
\passthrough{\lstinline!interval!} & \passthrough{\lstinline!int!} &
传给该接口的 \passthrough{\lstinline!interval!} 参数，Python 类型为
\passthrough{\lstinline!int!}。精确含义、单位、范围和枚举值需查目标型号对应头文件与协议。
& \passthrough{\lstinline!value = obj.interval!} /
\passthrough{\lstinline!obj.interval = value!} \\
\passthrough{\lstinline!duration!} & \passthrough{\lstinline!int!} &
操作持续时间参数。精确单位、范围和默认行为以目标型号接口定义为准。 &
\passthrough{\lstinline!value = obj.duration!} /
\passthrough{\lstinline!obj.duration = value!} \\
\bottomrule
\end{longtable}

\textbf{方法索引}

\begin{longtable}[]{@{}
  >{\raggedright\arraybackslash}p{(\columnwidth - 6\tabcolsep) * \real{0.18}}
  >{\raggedright\arraybackslash}p{(\columnwidth - 6\tabcolsep) * \real{0.24}}
  >{\raggedright\arraybackslash}p{(\columnwidth - 6\tabcolsep) * \real{0.18}}
  >{\raggedright\arraybackslash}p{(\columnwidth - 6\tabcolsep) * \real{0.40}}@{}}
\toprule
方法 & Python 签名 & 状态 & 安全分类 \\
\midrule
\endhead
\protect\hyperlink{unitree-sdk2-cpp-robot-go2-setreportfreqparameter-dunder-init-1}{\passthrough{\lstinline!\_\_init\_\_!}}
& \passthrough{\lstinline!def \_\_init\_\_(self) -> None!} &
\passthrough{\lstinline!SIGNATURE\_ONLY!} &
\passthrough{\lstinline!CONSTRUCTION!} \\
\protect\hyperlink{unitree-sdk2-cpp-robot-go2-setreportfreqparameter-from-json-1}{\passthrough{\lstinline!from\_json!}}
&
\passthrough{\lstinline!def from\_json(self, json: dict[str, Any]) -> None!}
& \passthrough{\lstinline!SIGNATURE\_ONLY!} &
\passthrough{\lstinline!UNCLASSIFIED!} \\
\protect\hyperlink{unitree-sdk2-cpp-robot-go2-setreportfreqparameter-to-json-1}{\passthrough{\lstinline!to\_json!}}
&
\passthrough{\lstinline!def to\_json(self, json: dict[str, Any]) -> None!}
& \passthrough{\lstinline!SIGNATURE\_ONLY!} &
\passthrough{\lstinline!UNCLASSIFIED!} \\
\bottomrule
\end{longtable}

\hypertarget{unitree-sdk2-cpp-robot-go2-setreportfreqparameter-dunder-init-1}{%
\paragraph{\texorpdfstring{\texttt{unitree\_sdk2\_cpp.robot.go2.SetReportFreqParameter.\_\_init\_\_}}{unitree\_sdk2\_cpp.robot.go2.SetReportFreqParameter.\_\_init\_\_}}\label{unitree-sdk2-cpp-robot-go2-setreportfreqparameter-dunder-init-1}}

初始化 \passthrough{\lstinline!SetReportFreqParameter!}
实例。是否能实际构造取决于下方可用性状态。

\textbf{签名}

\begin{lstlisting}[language=Python]
def __init__(self) -> None
\end{lstlisting}

\textbf{可用性与安全性}

\textbf{\passthrough{\lstinline!SIGNATURE\_ONLY!} /
\passthrough{\lstinline!CONSTRUCTION!}}：当前只有设计期签名，不能假设已存在于运行时扩展。

\textbf{参数}

无显式参数。实例方法中的 \passthrough{\lstinline!self!} 由 Python
自动传入。

\textbf{返回值}

\begin{longtable}[]{@{}
  >{\raggedright\arraybackslash}p{(\columnwidth - 4\tabcolsep) * \real{0.18}}
  >{\raggedright\arraybackslash}p{(\columnwidth - 4\tabcolsep) * \real{0.20}}
  >{\raggedright\arraybackslash}p{(\columnwidth - 4\tabcolsep) * \real{0.62}}@{}}
\toprule
位置 & 类型 & 含义 \\
\midrule
\endhead
无 & \passthrough{\lstinline!None!} &
\passthrough{\lstinline!\_\_init\_\_!}
本身不返回值；调用类对象时，在构造函数可用的前提下得到该类实例。 \\
\bottomrule
\end{longtable}

\textbf{对应 C++}

\begin{itemize}
\tightlist
\item
  类：\passthrough{\lstinline!unitree::robot::go2::SetReportFreqParameter!}
\item
  签名：\passthrough{\lstinline!SetReportFreqParameter()!}
\item
  绑定策略：\passthrough{\lstinline!未单独标注!}
\item
  声明位置：\passthrough{\lstinline!include/unitree/robot/go2/robot\_state/robot\_state\_api.hpp:93!}
\end{itemize}

\textbf{说明}

\begin{itemize}
\tightlist
\item
  示例是局部调用片段；\passthrough{\lstinline!obj!}
  和参数变量表示已经按上层业务校验并准备好的对象。
\item
  该条目可以被编辑器和类型检查器识别，但当前运行时可能在属性查找阶段直接失败。
\end{itemize}

\textbf{用法}

\begin{lstlisting}[language=Python]
from typing import TYPE_CHECKING

if TYPE_CHECKING:
    def planned_construct() -> SetReportFreqParameter:
        return SetReportFreqParameter()
\end{lstlisting}

\begin{quote}
{[}!CAUTION{]} 上面的 \passthrough{\lstinline!TYPE\_CHECKING!}
示例只表达计划调用形态。该分支在普通 Python
运行时为假，不代表当前扩展能执行此接口。
\end{quote}

\hypertarget{unitree-sdk2-cpp-robot-go2-setreportfreqparameter-from-json-1}{%
\paragraph{\texorpdfstring{\texttt{unitree\_sdk2\_cpp.robot.go2.SetReportFreqParameter.from\_json}}{unitree\_sdk2\_cpp.robot.go2.SetReportFreqParameter.from\_json}}\label{unitree-sdk2-cpp-robot-go2-setreportfreqparameter-from-json-1}}

计划从 JSON 风格字典读取字段并更新当前 SDK 值对象。

\textbf{签名}

\begin{lstlisting}[language=Python]
def from_json(self, json: dict[str, Any]) -> None
\end{lstlisting}

\textbf{可用性与安全性}

\textbf{\passthrough{\lstinline!SIGNATURE\_ONLY!} /
\passthrough{\lstinline!UNCLASSIFIED!}}：当前只有设计期签名，不能假设已存在于运行时扩展。

\textbf{参数}

\begin{longtable}[]{@{}
  >{\raggedright\arraybackslash}p{(\columnwidth - 6\tabcolsep) * \real{0.18}}
  >{\raggedright\arraybackslash}p{(\columnwidth - 6\tabcolsep) * \real{0.24}}
  >{\raggedright\arraybackslash}p{(\columnwidth - 6\tabcolsep) * \real{0.18}}
  >{\raggedright\arraybackslash}p{(\columnwidth - 6\tabcolsep) * \real{0.40}}@{}}
\toprule
名称 & Python 类型 & 默认值 & 含义 \\
\midrule
\endhead
\passthrough{\lstinline!json!} &
\passthrough{\lstinline!dict[str, Any]!} & 必填 & JSON
风格字典，键和值必须满足对应 C++ \passthrough{\lstinline!JsonMap!}
数据结构的约束。 对应 C++ 参数
\passthrough{\lstinline!json: common::JsonMap \&!}。 \\
\bottomrule
\end{longtable}

\textbf{返回值}

\begin{longtable}[]{@{}
  >{\raggedright\arraybackslash}p{(\columnwidth - 4\tabcolsep) * \real{0.18}}
  >{\raggedright\arraybackslash}p{(\columnwidth - 4\tabcolsep) * \real{0.20}}
  >{\raggedright\arraybackslash}p{(\columnwidth - 4\tabcolsep) * \real{0.62}}@{}}
\toprule
位置 & 类型 & 含义 \\
\midrule
\endhead
无 & \passthrough{\lstinline!None!} & 该方法不返回 Python
值；失败可能表现为异常或底层状态变化。 \\
\bottomrule
\end{longtable}

\textbf{对应 C++}

\begin{itemize}
\tightlist
\item
  类：\passthrough{\lstinline!unitree::robot::go2::SetReportFreqParameter!}
\item
  签名：\passthrough{\lstinline!fromJson(common::JsonMap \&)!}
\item
  绑定策略：\passthrough{\lstinline!SIGNATURE\_PREVIEW!}
\item
  声明位置：\passthrough{\lstinline!include/unitree/robot/go2/robot\_state/robot\_state\_api.hpp:99!}
\end{itemize}

\textbf{说明}

\begin{itemize}
\tightlist
\item
  示例是局部调用片段；\passthrough{\lstinline!obj!}
  和参数变量表示已经按上层业务校验并准备好的对象。
\item
  该条目可以被编辑器和类型检查器识别，但当前运行时可能在属性查找阶段直接失败。
\end{itemize}

\textbf{用法}

\begin{lstlisting}[language=Python]
from typing import TYPE_CHECKING

if TYPE_CHECKING:
    def planned_from_json(obj: SetReportFreqParameter, json: dict[str, Any]) -> None:
        obj.from_json(json=json)
\end{lstlisting}

\begin{quote}
{[}!CAUTION{]} 上面的 \passthrough{\lstinline!TYPE\_CHECKING!}
示例只表达计划调用形态。该分支在普通 Python
运行时为假，不代表当前扩展能执行此接口。
\end{quote}

\hypertarget{unitree-sdk2-cpp-robot-go2-setreportfreqparameter-to-json-1}{%
\paragraph{\texorpdfstring{\texttt{unitree\_sdk2\_cpp.robot.go2.SetReportFreqParameter.to\_json}}{unitree\_sdk2\_cpp.robot.go2.SetReportFreqParameter.to\_json}}\label{unitree-sdk2-cpp-robot-go2-setreportfreqparameter-to-json-1}}

计划把当前 SDK 值对象写入 JSON 风格字典。

\textbf{签名}

\begin{lstlisting}[language=Python]
def to_json(self, json: dict[str, Any]) -> None
\end{lstlisting}

\textbf{可用性与安全性}

\textbf{\passthrough{\lstinline!SIGNATURE\_ONLY!} /
\passthrough{\lstinline!UNCLASSIFIED!}}：当前只有设计期签名，不能假设已存在于运行时扩展。

\textbf{参数}

\begin{longtable}[]{@{}
  >{\raggedright\arraybackslash}p{(\columnwidth - 6\tabcolsep) * \real{0.18}}
  >{\raggedright\arraybackslash}p{(\columnwidth - 6\tabcolsep) * \real{0.24}}
  >{\raggedright\arraybackslash}p{(\columnwidth - 6\tabcolsep) * \real{0.18}}
  >{\raggedright\arraybackslash}p{(\columnwidth - 6\tabcolsep) * \real{0.40}}@{}}
\toprule
名称 & Python 类型 & 默认值 & 含义 \\
\midrule
\endhead
\passthrough{\lstinline!json!} &
\passthrough{\lstinline!dict[str, Any]!} & 必填 & JSON
风格字典，键和值必须满足对应 C++ \passthrough{\lstinline!JsonMap!}
数据结构的约束。 对应 C++ 参数
\passthrough{\lstinline!json: common::JsonMap \&!}。 \\
\bottomrule
\end{longtable}

\textbf{返回值}

\begin{longtable}[]{@{}
  >{\raggedright\arraybackslash}p{(\columnwidth - 4\tabcolsep) * \real{0.18}}
  >{\raggedright\arraybackslash}p{(\columnwidth - 4\tabcolsep) * \real{0.20}}
  >{\raggedright\arraybackslash}p{(\columnwidth - 4\tabcolsep) * \real{0.62}}@{}}
\toprule
位置 & 类型 & 含义 \\
\midrule
\endhead
无 & \passthrough{\lstinline!None!} & 该方法不返回 Python
值；失败可能表现为异常或底层状态变化。 \\
\bottomrule
\end{longtable}

\textbf{对应 C++}

\begin{itemize}
\tightlist
\item
  类：\passthrough{\lstinline!unitree::robot::go2::SetReportFreqParameter!}
\item
  签名：\passthrough{\lstinline!toJson(common::JsonMap \&) const!}
\item
  绑定策略：\passthrough{\lstinline!SIGNATURE\_PREVIEW!}
\item
  声明位置：\passthrough{\lstinline!include/unitree/robot/go2/robot\_state/robot\_state\_api.hpp:105!}
\end{itemize}

\textbf{说明}

\begin{itemize}
\tightlist
\item
  示例是局部调用片段；\passthrough{\lstinline!obj!}
  和参数变量表示已经按上层业务校验并准备好的对象。
\item
  该条目可以被编辑器和类型检查器识别，但当前运行时可能在属性查找阶段直接失败。
\end{itemize}

\textbf{用法}

\begin{lstlisting}[language=Python]
from typing import TYPE_CHECKING

if TYPE_CHECKING:
    def planned_to_json(obj: SetReportFreqParameter, json: dict[str, Any]) -> None:
        obj.to_json(json=json)
\end{lstlisting}

\begin{quote}
{[}!CAUTION{]} 上面的 \passthrough{\lstinline!TYPE\_CHECKING!}
示例只表达计划调用形态。该分支在普通 Python
运行时为假，不代表当前扩展能执行此接口。
\end{quote}

\hypertarget{unitree-sdk2-cpp-robot-go2-sportclient}{%
\subsubsection{\texorpdfstring{\texttt{unitree\_sdk2\_cpp.robot.go2.SportClient}}{unitree\_sdk2\_cpp.robot.go2.SportClient}}\label{unitree-sdk2-cpp-robot-go2-sportclient}}

Robot 服务客户端。构造、初始化和具体方法的可用性必须分别检查。

\textbf{导入}

\begin{lstlisting}[language=Python]
from unitree_sdk2_cpp.robot.go2 import SportClient
\end{lstlisting}

\textbf{基类}：\passthrough{\lstinline!Client!}

\textbf{方法索引}

\begin{longtable}[]{@{}
  >{\raggedright\arraybackslash}p{(\columnwidth - 6\tabcolsep) * \real{0.18}}
  >{\raggedright\arraybackslash}p{(\columnwidth - 6\tabcolsep) * \real{0.24}}
  >{\raggedright\arraybackslash}p{(\columnwidth - 6\tabcolsep) * \real{0.18}}
  >{\raggedright\arraybackslash}p{(\columnwidth - 6\tabcolsep) * \real{0.40}}@{}}
\toprule
方法 & Python 签名 & 状态 & 安全分类 \\
\midrule
\endhead
\protect\hyperlink{unitree-sdk2-cpp-robot-go2-sportclient-dunder-init-1}{\passthrough{\lstinline!\_\_init\_\_!}}
&
\passthrough{\lstinline!def \_\_init\_\_(self, enable\_lease: bool = ...) -> None!}
& \passthrough{\lstinline!SIGNATURE\_ONLY!} &
\passthrough{\lstinline!CONSTRUCTION!} \\
\protect\hyperlink{unitree-sdk2-cpp-robot-go2-sportclient-init-1}{\passthrough{\lstinline!init!}}
& \passthrough{\lstinline!def init(self) -> None!} &
\passthrough{\lstinline!SIGNATURE\_ONLY!} &
\passthrough{\lstinline!INITIALIZATION!} \\
\protect\hyperlink{unitree-sdk2-cpp-robot-go2-sportclient-damp-1}{\passthrough{\lstinline!damp!}}
& \passthrough{\lstinline!def damp(self) -> int!} &
\passthrough{\lstinline!SIGNATURE\_ONLY!} &
\passthrough{\lstinline!MOTION\_COMMAND!} \\
\protect\hyperlink{unitree-sdk2-cpp-robot-go2-sportclient-balance-stand-1}{\passthrough{\lstinline!balance\_stand!}}
& \passthrough{\lstinline!def balance\_stand(self) -> int!} &
\passthrough{\lstinline!SIGNATURE\_ONLY!} &
\passthrough{\lstinline!MOTION\_COMMAND!} \\
\protect\hyperlink{unitree-sdk2-cpp-robot-go2-sportclient-stop-move-1}{\passthrough{\lstinline!stop\_move!}}
& \passthrough{\lstinline!def stop\_move(self) -> int!} &
\passthrough{\lstinline!SIGNATURE\_ONLY!} &
\passthrough{\lstinline!MOTION\_COMMAND!} \\
\protect\hyperlink{unitree-sdk2-cpp-robot-go2-sportclient-stand-up-1}{\passthrough{\lstinline!stand\_up!}}
& \passthrough{\lstinline!def stand\_up(self) -> int!} &
\passthrough{\lstinline!SIGNATURE\_ONLY!} &
\passthrough{\lstinline!MOTION\_COMMAND!} \\
\protect\hyperlink{unitree-sdk2-cpp-robot-go2-sportclient-stand-down-1}{\passthrough{\lstinline!stand\_down!}}
& \passthrough{\lstinline!def stand\_down(self) -> int!} &
\passthrough{\lstinline!SIGNATURE\_ONLY!} &
\passthrough{\lstinline!MOTION\_COMMAND!} \\
\protect\hyperlink{unitree-sdk2-cpp-robot-go2-sportclient-recovery-stand-1}{\passthrough{\lstinline!recovery\_stand!}}
& \passthrough{\lstinline!def recovery\_stand(self) -> int!} &
\passthrough{\lstinline!SIGNATURE\_ONLY!} &
\passthrough{\lstinline!MOTION\_COMMAND!} \\
\protect\hyperlink{unitree-sdk2-cpp-robot-go2-sportclient-euler-1}{\passthrough{\lstinline!euler!}}
&
\passthrough{\lstinline!def euler(self, roll: float, pitch: float, yaw: float) -> int!}
& \passthrough{\lstinline!SIGNATURE\_ONLY!} &
\passthrough{\lstinline!MOTION\_COMMAND!} \\
\protect\hyperlink{unitree-sdk2-cpp-robot-go2-sportclient-move-1}{\passthrough{\lstinline!move!}}
&
\passthrough{\lstinline!def move(self, vx: float, vy: float, vyaw: float) -> int!}
& \passthrough{\lstinline!SIGNATURE\_ONLY!} &
\passthrough{\lstinline!MOTION\_COMMAND!} \\
\protect\hyperlink{unitree-sdk2-cpp-robot-go2-sportclient-sit-1}{\passthrough{\lstinline!sit!}}
& \passthrough{\lstinline!def sit(self) -> int!} &
\passthrough{\lstinline!SIGNATURE\_ONLY!} &
\passthrough{\lstinline!MOTION\_COMMAND!} \\
\protect\hyperlink{unitree-sdk2-cpp-robot-go2-sportclient-rise-sit-1}{\passthrough{\lstinline!rise\_sit!}}
& \passthrough{\lstinline!def rise\_sit(self) -> int!} &
\passthrough{\lstinline!SIGNATURE\_ONLY!} &
\passthrough{\lstinline!MOTION\_COMMAND!} \\
\protect\hyperlink{unitree-sdk2-cpp-robot-go2-sportclient-speed-level-1}{\passthrough{\lstinline!speed\_level!}}
& \passthrough{\lstinline!def speed\_level(self, level: int) -> int!} &
\passthrough{\lstinline!SIGNATURE\_ONLY!} &
\passthrough{\lstinline!MOTION\_COMMAND!} \\
\protect\hyperlink{unitree-sdk2-cpp-robot-go2-sportclient-hello-1}{\passthrough{\lstinline!hello!}}
& \passthrough{\lstinline!def hello(self) -> int!} &
\passthrough{\lstinline!SIGNATURE\_ONLY!} &
\passthrough{\lstinline!MOTION\_COMMAND!} \\
\protect\hyperlink{unitree-sdk2-cpp-robot-go2-sportclient-stretch-1}{\passthrough{\lstinline!stretch!}}
& \passthrough{\lstinline!def stretch(self) -> int!} &
\passthrough{\lstinline!SIGNATURE\_ONLY!} &
\passthrough{\lstinline!MOTION\_COMMAND!} \\
\protect\hyperlink{unitree-sdk2-cpp-robot-go2-sportclient-switch-joystick-1}{\passthrough{\lstinline!switch\_joystick!}}
&
\passthrough{\lstinline!def switch\_joystick(self, flag: bool) -> int!}
& \passthrough{\lstinline!SIGNATURE\_ONLY!} &
\passthrough{\lstinline!MOTION\_COMMAND!} \\
\protect\hyperlink{unitree-sdk2-cpp-robot-go2-sportclient-content-1}{\passthrough{\lstinline!content!}}
& \passthrough{\lstinline!def content(self) -> int!} &
\passthrough{\lstinline!SIGNATURE\_ONLY!} &
\passthrough{\lstinline!MOTION\_COMMAND!} \\
\protect\hyperlink{unitree-sdk2-cpp-robot-go2-sportclient-heart-1}{\passthrough{\lstinline!heart!}}
& \passthrough{\lstinline!def heart(self) -> int!} &
\passthrough{\lstinline!SIGNATURE\_ONLY!} &
\passthrough{\lstinline!MOTION\_COMMAND!} \\
\protect\hyperlink{unitree-sdk2-cpp-robot-go2-sportclient-pose-1}{\passthrough{\lstinline!pose!}}
& \passthrough{\lstinline!def pose(self, flag: bool) -> int!} &
\passthrough{\lstinline!SIGNATURE\_ONLY!} &
\passthrough{\lstinline!MOTION\_COMMAND!} \\
\protect\hyperlink{unitree-sdk2-cpp-robot-go2-sportclient-scrape-1}{\passthrough{\lstinline!scrape!}}
& \passthrough{\lstinline!def scrape(self) -> int!} &
\passthrough{\lstinline!SIGNATURE\_ONLY!} &
\passthrough{\lstinline!MOTION\_COMMAND!} \\
\protect\hyperlink{unitree-sdk2-cpp-robot-go2-sportclient-front-flip-1}{\passthrough{\lstinline!front\_flip!}}
& \passthrough{\lstinline!def front\_flip(self) -> int!} &
\passthrough{\lstinline!SIGNATURE\_ONLY!} &
\passthrough{\lstinline!MOTION\_COMMAND!} \\
\protect\hyperlink{unitree-sdk2-cpp-robot-go2-sportclient-front-jump-1}{\passthrough{\lstinline!front\_jump!}}
& \passthrough{\lstinline!def front\_jump(self) -> int!} &
\passthrough{\lstinline!SIGNATURE\_ONLY!} &
\passthrough{\lstinline!MOTION\_COMMAND!} \\
\protect\hyperlink{unitree-sdk2-cpp-robot-go2-sportclient-front-pounce-1}{\passthrough{\lstinline!front\_pounce!}}
& \passthrough{\lstinline!def front\_pounce(self) -> int!} &
\passthrough{\lstinline!SIGNATURE\_ONLY!} &
\passthrough{\lstinline!MOTION\_COMMAND!} \\
\protect\hyperlink{unitree-sdk2-cpp-robot-go2-sportclient-dance1-1}{\passthrough{\lstinline!dance1!}}
& \passthrough{\lstinline!def dance1(self) -> int!} &
\passthrough{\lstinline!SIGNATURE\_ONLY!} &
\passthrough{\lstinline!MOTION\_COMMAND!} \\
\protect\hyperlink{unitree-sdk2-cpp-robot-go2-sportclient-dance2-1}{\passthrough{\lstinline!dance2!}}
& \passthrough{\lstinline!def dance2(self) -> int!} &
\passthrough{\lstinline!SIGNATURE\_ONLY!} &
\passthrough{\lstinline!MOTION\_COMMAND!} \\
\protect\hyperlink{unitree-sdk2-cpp-robot-go2-sportclient-left-flip-1}{\passthrough{\lstinline!left\_flip!}}
& \passthrough{\lstinline!def left\_flip(self) -> int!} &
\passthrough{\lstinline!SIGNATURE\_ONLY!} &
\passthrough{\lstinline!MOTION\_COMMAND!} \\
\protect\hyperlink{unitree-sdk2-cpp-robot-go2-sportclient-back-flip-1}{\passthrough{\lstinline!back\_flip!}}
& \passthrough{\lstinline!def back\_flip(self) -> int!} &
\passthrough{\lstinline!SIGNATURE\_ONLY!} &
\passthrough{\lstinline!MOTION\_COMMAND!} \\
\protect\hyperlink{unitree-sdk2-cpp-robot-go2-sportclient-hand-stand-1}{\passthrough{\lstinline!hand\_stand!}}
& \passthrough{\lstinline!def hand\_stand(self, flag: bool) -> int!} &
\passthrough{\lstinline!SIGNATURE\_ONLY!} &
\passthrough{\lstinline!MOTION\_COMMAND!} \\
\protect\hyperlink{unitree-sdk2-cpp-robot-go2-sportclient-free-walk-1}{\passthrough{\lstinline!free\_walk!}}
& \passthrough{\lstinline!def free\_walk(self) -> int!} &
\passthrough{\lstinline!SIGNATURE\_ONLY!} &
\passthrough{\lstinline!MOTION\_COMMAND!} \\
\protect\hyperlink{unitree-sdk2-cpp-robot-go2-sportclient-free-bound-1}{\passthrough{\lstinline!free\_bound!}}
& \passthrough{\lstinline!def free\_bound(self, flag: bool) -> int!} &
\passthrough{\lstinline!SIGNATURE\_ONLY!} &
\passthrough{\lstinline!MOTION\_COMMAND!} \\
\protect\hyperlink{unitree-sdk2-cpp-robot-go2-sportclient-free-jump-1}{\passthrough{\lstinline!free\_jump!}}
& \passthrough{\lstinline!def free\_jump(self, flag: bool) -> int!} &
\passthrough{\lstinline!SIGNATURE\_ONLY!} &
\passthrough{\lstinline!MOTION\_COMMAND!} \\
\protect\hyperlink{unitree-sdk2-cpp-robot-go2-sportclient-free-avoid-1}{\passthrough{\lstinline!free\_avoid!}}
& \passthrough{\lstinline!def free\_avoid(self, flag: bool) -> int!} &
\passthrough{\lstinline!SIGNATURE\_ONLY!} &
\passthrough{\lstinline!MOTION\_COMMAND!} \\
\protect\hyperlink{unitree-sdk2-cpp-robot-go2-sportclient-classic-walk-1}{\passthrough{\lstinline!classic\_walk!}}
& \passthrough{\lstinline!def classic\_walk(self, flag: bool) -> int!} &
\passthrough{\lstinline!SIGNATURE\_ONLY!} &
\passthrough{\lstinline!MOTION\_COMMAND!} \\
\protect\hyperlink{unitree-sdk2-cpp-robot-go2-sportclient-walk-upright-1}{\passthrough{\lstinline!walk\_upright!}}
& \passthrough{\lstinline!def walk\_upright(self, flag: bool) -> int!} &
\passthrough{\lstinline!SIGNATURE\_ONLY!} &
\passthrough{\lstinline!MOTION\_COMMAND!} \\
\protect\hyperlink{unitree-sdk2-cpp-robot-go2-sportclient-cross-step-1}{\passthrough{\lstinline!cross\_step!}}
& \passthrough{\lstinline!def cross\_step(self, flag: bool) -> int!} &
\passthrough{\lstinline!SIGNATURE\_ONLY!} &
\passthrough{\lstinline!MOTION\_COMMAND!} \\
\protect\hyperlink{unitree-sdk2-cpp-robot-go2-sportclient-auto-recover-set-1}{\passthrough{\lstinline!auto\_recover\_set!}}
&
\passthrough{\lstinline!def auto\_recover\_set(self, flag: bool) -> int!}
& \passthrough{\lstinline!SIGNATURE\_ONLY!} &
\passthrough{\lstinline!MOTION\_COMMAND!} \\
\protect\hyperlink{unitree-sdk2-cpp-robot-go2-sportclient-auto-recover-get-1}{\passthrough{\lstinline!auto\_recover\_get!}}
&
\passthrough{\lstinline!def auto\_recover\_get(self) -> tuple[int, bool]!}
& \passthrough{\lstinline!SIGNATURE\_ONLY!} &
\passthrough{\lstinline!MOTION\_COMMAND!} \\
\protect\hyperlink{unitree-sdk2-cpp-robot-go2-sportclient-static-walk-1}{\passthrough{\lstinline!static\_walk!}}
& \passthrough{\lstinline!def static\_walk(self) -> int!} &
\passthrough{\lstinline!SIGNATURE\_ONLY!} &
\passthrough{\lstinline!MOTION\_COMMAND!} \\
\protect\hyperlink{unitree-sdk2-cpp-robot-go2-sportclient-trot-run-1}{\passthrough{\lstinline!trot\_run!}}
& \passthrough{\lstinline!def trot\_run(self) -> int!} &
\passthrough{\lstinline!SIGNATURE\_ONLY!} &
\passthrough{\lstinline!MOTION\_COMMAND!} \\
\protect\hyperlink{unitree-sdk2-cpp-robot-go2-sportclient-economic-gait-1}{\passthrough{\lstinline!economic\_gait!}}
& \passthrough{\lstinline!def economic\_gait(self) -> int!} &
\passthrough{\lstinline!SIGNATURE\_ONLY!} &
\passthrough{\lstinline!MOTION\_COMMAND!} \\
\protect\hyperlink{unitree-sdk2-cpp-robot-go2-sportclient-switch-avoid-mode-1}{\passthrough{\lstinline!switch\_avoid\_mode!}}
& \passthrough{\lstinline!def switch\_avoid\_mode(self) -> int!} &
\passthrough{\lstinline!SIGNATURE\_ONLY!} &
\passthrough{\lstinline!MOTION\_COMMAND!} \\
\bottomrule
\end{longtable}

\hypertarget{unitree-sdk2-cpp-robot-go2-sportclient-dunder-init-1}{%
\paragraph{\texorpdfstring{\texttt{unitree\_sdk2\_cpp.robot.go2.SportClient.\_\_init\_\_}}{unitree\_sdk2\_cpp.robot.go2.SportClient.\_\_init\_\_}}\label{unitree-sdk2-cpp-robot-go2-sportclient-dunder-init-1}}

初始化 \passthrough{\lstinline!SportClient!}
实例。是否能实际构造取决于下方可用性状态。

\textbf{签名}

\begin{lstlisting}[language=Python]
def __init__(self, enable_lease: bool = ...) -> None
\end{lstlisting}

\textbf{可用性与安全性}

\textbf{\passthrough{\lstinline!SIGNATURE\_ONLY!} /
\passthrough{\lstinline!CONSTRUCTION!}}：当前只有设计期签名，不能假设已存在于运行时扩展。

\textbf{参数}

\begin{longtable}[]{@{}
  >{\raggedright\arraybackslash}p{(\columnwidth - 6\tabcolsep) * \real{0.18}}
  >{\raggedright\arraybackslash}p{(\columnwidth - 6\tabcolsep) * \real{0.24}}
  >{\raggedright\arraybackslash}p{(\columnwidth - 6\tabcolsep) * \real{0.18}}
  >{\raggedright\arraybackslash}p{(\columnwidth - 6\tabcolsep) * \real{0.40}}@{}}
\toprule
名称 & Python 类型 & 默认值 & 含义 \\
\midrule
\endhead
\passthrough{\lstinline!enable\_lease!} & \passthrough{\lstinline!bool!}
& \passthrough{\lstinline!...!} & 是否启用 SDK lease 机制；lease
的获得、续期和释放规则以服务协议为准。 对应 C++ 参数
\passthrough{\lstinline!enableLease: bool!}。 只接受布尔语义。 \\
\bottomrule
\end{longtable}

\textbf{返回值}

\begin{longtable}[]{@{}
  >{\raggedright\arraybackslash}p{(\columnwidth - 4\tabcolsep) * \real{0.18}}
  >{\raggedright\arraybackslash}p{(\columnwidth - 4\tabcolsep) * \real{0.20}}
  >{\raggedright\arraybackslash}p{(\columnwidth - 4\tabcolsep) * \real{0.62}}@{}}
\toprule
位置 & 类型 & 含义 \\
\midrule
\endhead
无 & \passthrough{\lstinline!None!} &
\passthrough{\lstinline!\_\_init\_\_!}
本身不返回值；调用类对象时，在构造函数可用的前提下得到该类实例。 \\
\bottomrule
\end{longtable}

\textbf{对应 C++}

\begin{itemize}
\tightlist
\item
  类：\passthrough{\lstinline!unitree::robot::go2::SportClient!}
\item
  签名：\passthrough{\lstinline!SportClient(bool)!}
\item
  绑定策略：\passthrough{\lstinline!未单独标注!}
\item
  声明位置：\passthrough{\lstinline!include/unitree/robot/go2/sport/sport\_client.hpp:34!}
\end{itemize}

\textbf{说明}

\begin{itemize}
\tightlist
\item
  示例是局部调用片段；\passthrough{\lstinline!obj!}
  和参数变量表示已经按上层业务校验并准备好的对象。
\item
  默认值显示为 \passthrough{\lstinline!...!}，表示 C++
  声明存在默认参数，但当前 AST
  清单没有保存其字面量；省略参数可使用上游默认行为。
\item
  该条目可以被编辑器和类型检查器识别，但当前运行时可能在属性查找阶段直接失败。
\end{itemize}

\textbf{用法}

\begin{lstlisting}[language=Python]
from typing import TYPE_CHECKING

if TYPE_CHECKING:
    def planned_construct(enable_lease: bool) -> SportClient:
        return SportClient(enable_lease=enable_lease)
\end{lstlisting}

\begin{quote}
{[}!CAUTION{]} 上面的 \passthrough{\lstinline!TYPE\_CHECKING!}
示例只表达计划调用形态。该分支在普通 Python
运行时为假，不代表当前扩展能执行此接口。
\end{quote}

\hypertarget{unitree-sdk2-cpp-robot-go2-sportclient-init-1}{%
\paragraph{\texorpdfstring{\texttt{unitree\_sdk2\_cpp.robot.go2.SportClient.init}}{unitree\_sdk2\_cpp.robot.go2.SportClient.init}}\label{unitree-sdk2-cpp-robot-go2-sportclient-init-1}}

初始化当前 SDK 对象所需的底层通道或服务资源。

\textbf{签名}

\begin{lstlisting}[language=Python]
def init(self) -> None
\end{lstlisting}

\textbf{可用性与安全性}

\textbf{\passthrough{\lstinline!SIGNATURE\_ONLY!} /
\passthrough{\lstinline!INITIALIZATION!}}：当前只有设计期签名，不能假设已存在于运行时扩展。

\textbf{参数}

无显式参数。实例方法中的 \passthrough{\lstinline!self!} 由 Python
自动传入。

\textbf{返回值}

\begin{longtable}[]{@{}
  >{\raggedright\arraybackslash}p{(\columnwidth - 4\tabcolsep) * \real{0.18}}
  >{\raggedright\arraybackslash}p{(\columnwidth - 4\tabcolsep) * \real{0.20}}
  >{\raggedright\arraybackslash}p{(\columnwidth - 4\tabcolsep) * \real{0.62}}@{}}
\toprule
位置 & 类型 & 含义 \\
\midrule
\endhead
无 & \passthrough{\lstinline!None!} & 该方法不返回 Python
值；失败可能表现为异常或底层状态变化。 \\
\bottomrule
\end{longtable}

\textbf{对应 C++}

\begin{itemize}
\tightlist
\item
  类：\passthrough{\lstinline!unitree::robot::go2::SportClient!}
\item
  签名：\passthrough{\lstinline!Init()!}
\item
  绑定策略：\passthrough{\lstinline!DIRECT!}
\item
  声明位置：\passthrough{\lstinline!include/unitree/robot/go2/sport/sport\_client.hpp:37!}
\end{itemize}

\textbf{说明}

\begin{itemize}
\tightlist
\item
  示例是局部调用片段；\passthrough{\lstinline!obj!}
  和参数变量表示已经按上层业务校验并准备好的对象。
\item
  该条目可以被编辑器和类型检查器识别，但当前运行时可能在属性查找阶段直接失败。
\end{itemize}

\textbf{用法}

\begin{lstlisting}[language=Python]
from typing import TYPE_CHECKING

if TYPE_CHECKING:
    def planned_init(obj: SportClient) -> None:
        obj.init()
\end{lstlisting}

\begin{quote}
{[}!CAUTION{]} 上面的 \passthrough{\lstinline!TYPE\_CHECKING!}
示例只表达计划调用形态。该分支在普通 Python
运行时为假，不代表当前扩展能执行此接口。
\end{quote}

\hypertarget{unitree-sdk2-cpp-robot-go2-sportclient-damp-1}{%
\paragraph{\texorpdfstring{\texttt{unitree\_sdk2\_cpp.robot.go2.SportClient.damp}}{unitree\_sdk2\_cpp.robot.go2.SportClient.damp}}\label{unitree-sdk2-cpp-robot-go2-sportclient-damp-1}}

对应 C++ SDK 操作
\passthrough{\lstinline!Damp()!}。上游头文件没有可直接生成的业务说明时，本参考只保证签名映射，精确语义需查目标型号协议。

\textbf{签名}

\begin{lstlisting}[language=Python]
def damp(self) -> int
\end{lstlisting}

\textbf{可用性与安全性}

\textbf{\passthrough{\lstinline!SIGNATURE\_ONLY!} /
\passthrough{\lstinline!MOTION\_COMMAND!}}：仅用于补全和静态检查。当前二进制没有该
Python 方法；不得按可执行运动接口使用。

\textbf{参数}

无显式参数。实例方法中的 \passthrough{\lstinline!self!} 由 Python
自动传入。

\textbf{返回值}

\begin{longtable}[]{@{}
  >{\raggedright\arraybackslash}p{(\columnwidth - 4\tabcolsep) * \real{0.18}}
  >{\raggedright\arraybackslash}p{(\columnwidth - 4\tabcolsep) * \real{0.20}}
  >{\raggedright\arraybackslash}p{(\columnwidth - 4\tabcolsep) * \real{0.62}}@{}}
\toprule
位置 & 类型 & 含义 \\
\midrule
\endhead
返回值 & \passthrough{\lstinline!int!} & SDK 状态码；Unitree 示例通常以
\passthrough{\lstinline!0!} 表示成功，非零值需按具体服务定义解释。 \\
\bottomrule
\end{longtable}

\textbf{对应 C++}

\begin{itemize}
\tightlist
\item
  类：\passthrough{\lstinline!unitree::robot::go2::SportClient!}
\item
  签名：\passthrough{\lstinline!Damp()!}
\item
  绑定策略：\passthrough{\lstinline!DIRECT!}
\item
  声明位置：\passthrough{\lstinline!include/unitree/robot/go2/sport/sport\_client.hpp:39!}
\end{itemize}

\textbf{说明}

\begin{itemize}
\tightlist
\item
  示例是局部调用片段；\passthrough{\lstinline!obj!}
  和参数变量表示已经按上层业务校验并准备好的对象。
\item
  该条目可以被编辑器和类型检查器识别，但当前运行时可能在属性查找阶段直接失败。
\item
  该方法属于运动命令。方法名包含
  \passthrough{\lstinline!stop!}、\passthrough{\lstinline!damp!} 或
  \passthrough{\lstinline!zero!} 也不代表它是独立物理急停。
\end{itemize}

\textbf{用法}

\begin{lstlisting}[language=Python]
from typing import TYPE_CHECKING

if TYPE_CHECKING:
    def planned_damp(obj: SportClient) -> int:
        return obj.damp()
\end{lstlisting}

\begin{quote}
{[}!CAUTION{]} 上面的 \passthrough{\lstinline!TYPE\_CHECKING!}
示例只表达计划调用形态。该分支在普通 Python
运行时为假，不代表当前扩展能执行此接口。
\end{quote}

\hypertarget{unitree-sdk2-cpp-robot-go2-sportclient-balance-stand-1}{%
\paragraph{\texorpdfstring{\texttt{unitree\_sdk2\_cpp.robot.go2.SportClient.balance\_stand}}{unitree\_sdk2\_cpp.robot.go2.SportClient.balance\_stand}}\label{unitree-sdk2-cpp-robot-go2-sportclient-balance-stand-1}}

对应 C++ SDK 操作
\passthrough{\lstinline!BalanceStand()!}。上游头文件没有可直接生成的业务说明时，本参考只保证签名映射，精确语义需查目标型号协议。

\textbf{签名}

\begin{lstlisting}[language=Python]
def balance_stand(self) -> int
\end{lstlisting}

\textbf{可用性与安全性}

\textbf{\passthrough{\lstinline!SIGNATURE\_ONLY!} /
\passthrough{\lstinline!MOTION\_COMMAND!}}：仅用于补全和静态检查。当前二进制没有该
Python 方法；不得按可执行运动接口使用。

\textbf{参数}

无显式参数。实例方法中的 \passthrough{\lstinline!self!} 由 Python
自动传入。

\textbf{返回值}

\begin{longtable}[]{@{}
  >{\raggedright\arraybackslash}p{(\columnwidth - 4\tabcolsep) * \real{0.18}}
  >{\raggedright\arraybackslash}p{(\columnwidth - 4\tabcolsep) * \real{0.20}}
  >{\raggedright\arraybackslash}p{(\columnwidth - 4\tabcolsep) * \real{0.62}}@{}}
\toprule
位置 & 类型 & 含义 \\
\midrule
\endhead
返回值 & \passthrough{\lstinline!int!} & SDK 状态码；Unitree 示例通常以
\passthrough{\lstinline!0!} 表示成功，非零值需按具体服务定义解释。 \\
\bottomrule
\end{longtable}

\textbf{对应 C++}

\begin{itemize}
\tightlist
\item
  类：\passthrough{\lstinline!unitree::robot::go2::SportClient!}
\item
  签名：\passthrough{\lstinline!BalanceStand()!}
\item
  绑定策略：\passthrough{\lstinline!DIRECT!}
\item
  声明位置：\passthrough{\lstinline!include/unitree/robot/go2/sport/sport\_client.hpp:40!}
\end{itemize}

\textbf{说明}

\begin{itemize}
\tightlist
\item
  示例是局部调用片段；\passthrough{\lstinline!obj!}
  和参数变量表示已经按上层业务校验并准备好的对象。
\item
  该条目可以被编辑器和类型检查器识别，但当前运行时可能在属性查找阶段直接失败。
\item
  该方法属于运动命令。方法名包含
  \passthrough{\lstinline!stop!}、\passthrough{\lstinline!damp!} 或
  \passthrough{\lstinline!zero!} 也不代表它是独立物理急停。
\end{itemize}

\textbf{用法}

\begin{lstlisting}[language=Python]
from typing import TYPE_CHECKING

if TYPE_CHECKING:
    def planned_balance_stand(obj: SportClient) -> int:
        return obj.balance_stand()
\end{lstlisting}

\begin{quote}
{[}!CAUTION{]} 上面的 \passthrough{\lstinline!TYPE\_CHECKING!}
示例只表达计划调用形态。该分支在普通 Python
运行时为假，不代表当前扩展能执行此接口。
\end{quote}

\hypertarget{unitree-sdk2-cpp-robot-go2-sportclient-stop-move-1}{%
\paragraph{\texorpdfstring{\texttt{unitree\_sdk2\_cpp.robot.go2.SportClient.stop\_move}}{unitree\_sdk2\_cpp.robot.go2.SportClient.stop\_move}}\label{unitree-sdk2-cpp-robot-go2-sportclient-stop-move-1}}

请求停止对应 SDK 操作。该名称不等同于经过验证的物理急停。

\textbf{签名}

\begin{lstlisting}[language=Python]
def stop_move(self) -> int
\end{lstlisting}

\textbf{可用性与安全性}

\textbf{\passthrough{\lstinline!SIGNATURE\_ONLY!} /
\passthrough{\lstinline!MOTION\_COMMAND!}}：仅用于补全和静态检查。当前二进制没有该
Python 方法；不得按可执行运动接口使用。

\textbf{参数}

无显式参数。实例方法中的 \passthrough{\lstinline!self!} 由 Python
自动传入。

\textbf{返回值}

\begin{longtable}[]{@{}
  >{\raggedright\arraybackslash}p{(\columnwidth - 4\tabcolsep) * \real{0.18}}
  >{\raggedright\arraybackslash}p{(\columnwidth - 4\tabcolsep) * \real{0.20}}
  >{\raggedright\arraybackslash}p{(\columnwidth - 4\tabcolsep) * \real{0.62}}@{}}
\toprule
位置 & 类型 & 含义 \\
\midrule
\endhead
返回值 & \passthrough{\lstinline!int!} & SDK 状态码；Unitree 示例通常以
\passthrough{\lstinline!0!} 表示成功，非零值需按具体服务定义解释。 \\
\bottomrule
\end{longtable}

\textbf{对应 C++}

\begin{itemize}
\tightlist
\item
  类：\passthrough{\lstinline!unitree::robot::go2::SportClient!}
\item
  签名：\passthrough{\lstinline!StopMove()!}
\item
  绑定策略：\passthrough{\lstinline!DIRECT!}
\item
  声明位置：\passthrough{\lstinline!include/unitree/robot/go2/sport/sport\_client.hpp:41!}
\end{itemize}

\textbf{说明}

\begin{itemize}
\tightlist
\item
  示例是局部调用片段；\passthrough{\lstinline!obj!}
  和参数变量表示已经按上层业务校验并准备好的对象。
\item
  该条目可以被编辑器和类型检查器识别，但当前运行时可能在属性查找阶段直接失败。
\item
  该方法属于运动命令。方法名包含
  \passthrough{\lstinline!stop!}、\passthrough{\lstinline!damp!} 或
  \passthrough{\lstinline!zero!} 也不代表它是独立物理急停。
\end{itemize}

\textbf{用法}

\begin{lstlisting}[language=Python]
from typing import TYPE_CHECKING

if TYPE_CHECKING:
    def planned_stop_move(obj: SportClient) -> int:
        return obj.stop_move()
\end{lstlisting}

\begin{quote}
{[}!CAUTION{]} 上面的 \passthrough{\lstinline!TYPE\_CHECKING!}
示例只表达计划调用形态。该分支在普通 Python
运行时为假，不代表当前扩展能执行此接口。
\end{quote}

\hypertarget{unitree-sdk2-cpp-robot-go2-sportclient-stand-up-1}{%
\paragraph{\texorpdfstring{\texttt{unitree\_sdk2\_cpp.robot.go2.SportClient.stand\_up}}{unitree\_sdk2\_cpp.robot.go2.SportClient.stand\_up}}\label{unitree-sdk2-cpp-robot-go2-sportclient-stand-up-1}}

对应 C++ SDK 操作
\passthrough{\lstinline!StandUp()!}。上游头文件没有可直接生成的业务说明时，本参考只保证签名映射，精确语义需查目标型号协议。

\textbf{签名}

\begin{lstlisting}[language=Python]
def stand_up(self) -> int
\end{lstlisting}

\textbf{可用性与安全性}

\textbf{\passthrough{\lstinline!SIGNATURE\_ONLY!} /
\passthrough{\lstinline!MOTION\_COMMAND!}}：仅用于补全和静态检查。当前二进制没有该
Python 方法；不得按可执行运动接口使用。

\textbf{参数}

无显式参数。实例方法中的 \passthrough{\lstinline!self!} 由 Python
自动传入。

\textbf{返回值}

\begin{longtable}[]{@{}
  >{\raggedright\arraybackslash}p{(\columnwidth - 4\tabcolsep) * \real{0.18}}
  >{\raggedright\arraybackslash}p{(\columnwidth - 4\tabcolsep) * \real{0.20}}
  >{\raggedright\arraybackslash}p{(\columnwidth - 4\tabcolsep) * \real{0.62}}@{}}
\toprule
位置 & 类型 & 含义 \\
\midrule
\endhead
返回值 & \passthrough{\lstinline!int!} & SDK 状态码；Unitree 示例通常以
\passthrough{\lstinline!0!} 表示成功，非零值需按具体服务定义解释。 \\
\bottomrule
\end{longtable}

\textbf{对应 C++}

\begin{itemize}
\tightlist
\item
  类：\passthrough{\lstinline!unitree::robot::go2::SportClient!}
\item
  签名：\passthrough{\lstinline!StandUp()!}
\item
  绑定策略：\passthrough{\lstinline!DIRECT!}
\item
  声明位置：\passthrough{\lstinline!include/unitree/robot/go2/sport/sport\_client.hpp:42!}
\end{itemize}

\textbf{说明}

\begin{itemize}
\tightlist
\item
  示例是局部调用片段；\passthrough{\lstinline!obj!}
  和参数变量表示已经按上层业务校验并准备好的对象。
\item
  该条目可以被编辑器和类型检查器识别，但当前运行时可能在属性查找阶段直接失败。
\item
  该方法属于运动命令。方法名包含
  \passthrough{\lstinline!stop!}、\passthrough{\lstinline!damp!} 或
  \passthrough{\lstinline!zero!} 也不代表它是独立物理急停。
\end{itemize}

\textbf{用法}

\begin{lstlisting}[language=Python]
from typing import TYPE_CHECKING

if TYPE_CHECKING:
    def planned_stand_up(obj: SportClient) -> int:
        return obj.stand_up()
\end{lstlisting}

\begin{quote}
{[}!CAUTION{]} 上面的 \passthrough{\lstinline!TYPE\_CHECKING!}
示例只表达计划调用形态。该分支在普通 Python
运行时为假，不代表当前扩展能执行此接口。
\end{quote}

\hypertarget{unitree-sdk2-cpp-robot-go2-sportclient-stand-down-1}{%
\paragraph{\texorpdfstring{\texttt{unitree\_sdk2\_cpp.robot.go2.SportClient.stand\_down}}{unitree\_sdk2\_cpp.robot.go2.SportClient.stand\_down}}\label{unitree-sdk2-cpp-robot-go2-sportclient-stand-down-1}}

对应 C++ SDK 操作
\passthrough{\lstinline!StandDown()!}。上游头文件没有可直接生成的业务说明时，本参考只保证签名映射，精确语义需查目标型号协议。

\textbf{签名}

\begin{lstlisting}[language=Python]
def stand_down(self) -> int
\end{lstlisting}

\textbf{可用性与安全性}

\textbf{\passthrough{\lstinline!SIGNATURE\_ONLY!} /
\passthrough{\lstinline!MOTION\_COMMAND!}}：仅用于补全和静态检查。当前二进制没有该
Python 方法；不得按可执行运动接口使用。

\textbf{参数}

无显式参数。实例方法中的 \passthrough{\lstinline!self!} 由 Python
自动传入。

\textbf{返回值}

\begin{longtable}[]{@{}
  >{\raggedright\arraybackslash}p{(\columnwidth - 4\tabcolsep) * \real{0.18}}
  >{\raggedright\arraybackslash}p{(\columnwidth - 4\tabcolsep) * \real{0.20}}
  >{\raggedright\arraybackslash}p{(\columnwidth - 4\tabcolsep) * \real{0.62}}@{}}
\toprule
位置 & 类型 & 含义 \\
\midrule
\endhead
返回值 & \passthrough{\lstinline!int!} & SDK 状态码；Unitree 示例通常以
\passthrough{\lstinline!0!} 表示成功，非零值需按具体服务定义解释。 \\
\bottomrule
\end{longtable}

\textbf{对应 C++}

\begin{itemize}
\tightlist
\item
  类：\passthrough{\lstinline!unitree::robot::go2::SportClient!}
\item
  签名：\passthrough{\lstinline!StandDown()!}
\item
  绑定策略：\passthrough{\lstinline!DIRECT!}
\item
  声明位置：\passthrough{\lstinline!include/unitree/robot/go2/sport/sport\_client.hpp:43!}
\end{itemize}

\textbf{说明}

\begin{itemize}
\tightlist
\item
  示例是局部调用片段；\passthrough{\lstinline!obj!}
  和参数变量表示已经按上层业务校验并准备好的对象。
\item
  该条目可以被编辑器和类型检查器识别，但当前运行时可能在属性查找阶段直接失败。
\item
  该方法属于运动命令。方法名包含
  \passthrough{\lstinline!stop!}、\passthrough{\lstinline!damp!} 或
  \passthrough{\lstinline!zero!} 也不代表它是独立物理急停。
\end{itemize}

\textbf{用法}

\begin{lstlisting}[language=Python]
from typing import TYPE_CHECKING

if TYPE_CHECKING:
    def planned_stand_down(obj: SportClient) -> int:
        return obj.stand_down()
\end{lstlisting}

\begin{quote}
{[}!CAUTION{]} 上面的 \passthrough{\lstinline!TYPE\_CHECKING!}
示例只表达计划调用形态。该分支在普通 Python
运行时为假，不代表当前扩展能执行此接口。
\end{quote}

\hypertarget{unitree-sdk2-cpp-robot-go2-sportclient-recovery-stand-1}{%
\paragraph{\texorpdfstring{\texttt{unitree\_sdk2\_cpp.robot.go2.SportClient.recovery\_stand}}{unitree\_sdk2\_cpp.robot.go2.SportClient.recovery\_stand}}\label{unitree-sdk2-cpp-robot-go2-sportclient-recovery-stand-1}}

对应 C++ SDK 操作
\passthrough{\lstinline!RecoveryStand()!}。上游头文件没有可直接生成的业务说明时，本参考只保证签名映射，精确语义需查目标型号协议。

\textbf{签名}

\begin{lstlisting}[language=Python]
def recovery_stand(self) -> int
\end{lstlisting}

\textbf{可用性与安全性}

\textbf{\passthrough{\lstinline!SIGNATURE\_ONLY!} /
\passthrough{\lstinline!MOTION\_COMMAND!}}：仅用于补全和静态检查。当前二进制没有该
Python 方法；不得按可执行运动接口使用。

\textbf{参数}

无显式参数。实例方法中的 \passthrough{\lstinline!self!} 由 Python
自动传入。

\textbf{返回值}

\begin{longtable}[]{@{}
  >{\raggedright\arraybackslash}p{(\columnwidth - 4\tabcolsep) * \real{0.18}}
  >{\raggedright\arraybackslash}p{(\columnwidth - 4\tabcolsep) * \real{0.20}}
  >{\raggedright\arraybackslash}p{(\columnwidth - 4\tabcolsep) * \real{0.62}}@{}}
\toprule
位置 & 类型 & 含义 \\
\midrule
\endhead
返回值 & \passthrough{\lstinline!int!} & SDK 状态码；Unitree 示例通常以
\passthrough{\lstinline!0!} 表示成功，非零值需按具体服务定义解释。 \\
\bottomrule
\end{longtable}

\textbf{对应 C++}

\begin{itemize}
\tightlist
\item
  类：\passthrough{\lstinline!unitree::robot::go2::SportClient!}
\item
  签名：\passthrough{\lstinline!RecoveryStand()!}
\item
  绑定策略：\passthrough{\lstinline!DIRECT!}
\item
  声明位置：\passthrough{\lstinline!include/unitree/robot/go2/sport/sport\_client.hpp:44!}
\end{itemize}

\textbf{说明}

\begin{itemize}
\tightlist
\item
  示例是局部调用片段；\passthrough{\lstinline!obj!}
  和参数变量表示已经按上层业务校验并准备好的对象。
\item
  该条目可以被编辑器和类型检查器识别，但当前运行时可能在属性查找阶段直接失败。
\item
  该方法属于运动命令。方法名包含
  \passthrough{\lstinline!stop!}、\passthrough{\lstinline!damp!} 或
  \passthrough{\lstinline!zero!} 也不代表它是独立物理急停。
\end{itemize}

\textbf{用法}

\begin{lstlisting}[language=Python]
from typing import TYPE_CHECKING

if TYPE_CHECKING:
    def planned_recovery_stand(obj: SportClient) -> int:
        return obj.recovery_stand()
\end{lstlisting}

\begin{quote}
{[}!CAUTION{]} 上面的 \passthrough{\lstinline!TYPE\_CHECKING!}
示例只表达计划调用形态。该分支在普通 Python
运行时为假，不代表当前扩展能执行此接口。
\end{quote}

\hypertarget{unitree-sdk2-cpp-robot-go2-sportclient-euler-1}{%
\paragraph{\texorpdfstring{\texttt{unitree\_sdk2\_cpp.robot.go2.SportClient.euler}}{unitree\_sdk2\_cpp.robot.go2.SportClient.euler}}\label{unitree-sdk2-cpp-robot-go2-sportclient-euler-1}}

对应 C++ SDK 操作
\passthrough{\lstinline!Euler(float, float, float)!}。上游头文件没有可直接生成的业务说明时，本参考只保证签名映射，精确语义需查目标型号协议。

\textbf{签名}

\begin{lstlisting}[language=Python]
def euler(self, roll: float, pitch: float, yaw: float) -> int
\end{lstlisting}

\textbf{可用性与安全性}

\textbf{\passthrough{\lstinline!SIGNATURE\_ONLY!} /
\passthrough{\lstinline!MOTION\_COMMAND!}}：仅用于补全和静态检查。当前二进制没有该
Python 方法；不得按可执行运动接口使用。

\textbf{参数}

\begin{longtable}[]{@{}
  >{\raggedright\arraybackslash}p{(\columnwidth - 6\tabcolsep) * \real{0.18}}
  >{\raggedright\arraybackslash}p{(\columnwidth - 6\tabcolsep) * \real{0.24}}
  >{\raggedright\arraybackslash}p{(\columnwidth - 6\tabcolsep) * \real{0.18}}
  >{\raggedright\arraybackslash}p{(\columnwidth - 6\tabcolsep) * \real{0.40}}@{}}
\toprule
名称 & Python 类型 & 默认值 & 含义 \\
\midrule
\endhead
\passthrough{\lstinline!roll!} & \passthrough{\lstinline!float!} & 必填
& 传给该接口的 \passthrough{\lstinline!roll!} 参数，Python 类型为
\passthrough{\lstinline!float!}。精确含义、单位、范围和枚举值需查目标型号对应头文件与协议。
对应 C++ 参数 \passthrough{\lstinline!roll: float!}。
底层为浮点数；类型本身不说明单位、坐标系或业务安全范围。 \\
\passthrough{\lstinline!pitch!} & \passthrough{\lstinline!float!} & 必填
& 传给该接口的 \passthrough{\lstinline!pitch!} 参数，Python 类型为
\passthrough{\lstinline!float!}。精确含义、单位、范围和枚举值需查目标型号对应头文件与协议。
对应 C++ 参数 \passthrough{\lstinline!pitch: float!}。
底层为浮点数；类型本身不说明单位、坐标系或业务安全范围。 \\
\passthrough{\lstinline!yaw!} & \passthrough{\lstinline!float!} & 必填 &
偏航角或偏航目标参数。参考系、单位和范围必须查目标型号协议。 对应 C++
参数 \passthrough{\lstinline!yaw: float!}。
底层为浮点数；类型本身不说明单位、坐标系或业务安全范围。 \\
\bottomrule
\end{longtable}

\textbf{返回值}

\begin{longtable}[]{@{}
  >{\raggedright\arraybackslash}p{(\columnwidth - 4\tabcolsep) * \real{0.18}}
  >{\raggedright\arraybackslash}p{(\columnwidth - 4\tabcolsep) * \real{0.20}}
  >{\raggedright\arraybackslash}p{(\columnwidth - 4\tabcolsep) * \real{0.62}}@{}}
\toprule
位置 & 类型 & 含义 \\
\midrule
\endhead
返回值 & \passthrough{\lstinline!int!} & SDK 状态码；Unitree 示例通常以
\passthrough{\lstinline!0!} 表示成功，非零值需按具体服务定义解释。 \\
\bottomrule
\end{longtable}

\textbf{对应 C++}

\begin{itemize}
\tightlist
\item
  类：\passthrough{\lstinline!unitree::robot::go2::SportClient!}
\item
  签名：\passthrough{\lstinline!Euler(float, float, float)!}
\item
  绑定策略：\passthrough{\lstinline!DIRECT!}
\item
  声明位置：\passthrough{\lstinline!include/unitree/robot/go2/sport/sport\_client.hpp:45!}
\end{itemize}

\textbf{说明}

\begin{itemize}
\tightlist
\item
  示例是局部调用片段；\passthrough{\lstinline!obj!}
  和参数变量表示已经按上层业务校验并准备好的对象。
\item
  该条目可以被编辑器和类型检查器识别，但当前运行时可能在属性查找阶段直接失败。
\item
  该方法属于运动命令。方法名包含
  \passthrough{\lstinline!stop!}、\passthrough{\lstinline!damp!} 或
  \passthrough{\lstinline!zero!} 也不代表它是独立物理急停。
\end{itemize}

\textbf{用法}

\begin{lstlisting}[language=Python]
from typing import TYPE_CHECKING

if TYPE_CHECKING:
    def planned_euler(obj: SportClient, roll: float, pitch: float, yaw: float) -> int:
        return obj.euler(roll=roll, pitch=pitch, yaw=yaw)
\end{lstlisting}

\begin{quote}
{[}!CAUTION{]} 上面的 \passthrough{\lstinline!TYPE\_CHECKING!}
示例只表达计划调用形态。该分支在普通 Python
运行时为假，不代表当前扩展能执行此接口。
\end{quote}

\hypertarget{unitree-sdk2-cpp-robot-go2-sportclient-move-1}{%
\paragraph{\texorpdfstring{\texttt{unitree\_sdk2\_cpp.robot.go2.SportClient.move}}{unitree\_sdk2\_cpp.robot.go2.SportClient.move}}\label{unitree-sdk2-cpp-robot-go2-sportclient-move-1}}

对应 C++ SDK 操作
\passthrough{\lstinline!Move(float, float, float)!}。上游头文件没有可直接生成的业务说明时，本参考只保证签名映射，精确语义需查目标型号协议。

\textbf{签名}

\begin{lstlisting}[language=Python]
def move(self, vx: float, vy: float, vyaw: float) -> int
\end{lstlisting}

\textbf{可用性与安全性}

\textbf{\passthrough{\lstinline!SIGNATURE\_ONLY!} /
\passthrough{\lstinline!MOTION\_COMMAND!}}：仅用于补全和静态检查。当前二进制没有该
Python 方法；不得按可执行运动接口使用。

\textbf{参数}

\begin{longtable}[]{@{}
  >{\raggedright\arraybackslash}p{(\columnwidth - 6\tabcolsep) * \real{0.18}}
  >{\raggedright\arraybackslash}p{(\columnwidth - 6\tabcolsep) * \real{0.24}}
  >{\raggedright\arraybackslash}p{(\columnwidth - 6\tabcolsep) * \real{0.18}}
  >{\raggedright\arraybackslash}p{(\columnwidth - 6\tabcolsep) * \real{0.40}}@{}}
\toprule
名称 & Python 类型 & 默认值 & 含义 \\
\midrule
\endhead
\passthrough{\lstinline!vx!} & \passthrough{\lstinline!float!} & 必填 &
X 方向速度参数。坐标系、单位、符号和安全范围必须查目标型号运动协议。
对应 C++ 参数 \passthrough{\lstinline!vx: float!}。
底层为浮点数；类型本身不说明单位、坐标系或业务安全范围。 \\
\passthrough{\lstinline!vy!} & \passthrough{\lstinline!float!} & 必填 &
Y 方向速度参数。坐标系、单位、符号和安全范围必须查目标型号运动协议。
对应 C++ 参数 \passthrough{\lstinline!vy: float!}。
底层为浮点数；类型本身不说明单位、坐标系或业务安全范围。 \\
\passthrough{\lstinline!vyaw!} & \passthrough{\lstinline!float!} & 必填
& 偏航角速度参数。单位、符号和安全范围必须查目标型号运动协议。 对应 C++
参数 \passthrough{\lstinline!vyaw: float!}。
底层为浮点数；类型本身不说明单位、坐标系或业务安全范围。 \\
\bottomrule
\end{longtable}

\textbf{返回值}

\begin{longtable}[]{@{}
  >{\raggedright\arraybackslash}p{(\columnwidth - 4\tabcolsep) * \real{0.18}}
  >{\raggedright\arraybackslash}p{(\columnwidth - 4\tabcolsep) * \real{0.20}}
  >{\raggedright\arraybackslash}p{(\columnwidth - 4\tabcolsep) * \real{0.62}}@{}}
\toprule
位置 & 类型 & 含义 \\
\midrule
\endhead
返回值 & \passthrough{\lstinline!int!} & SDK 状态码；Unitree 示例通常以
\passthrough{\lstinline!0!} 表示成功，非零值需按具体服务定义解释。 \\
\bottomrule
\end{longtable}

\textbf{对应 C++}

\begin{itemize}
\tightlist
\item
  类：\passthrough{\lstinline!unitree::robot::go2::SportClient!}
\item
  签名：\passthrough{\lstinline!Move(float, float, float)!}
\item
  绑定策略：\passthrough{\lstinline!DIRECT!}
\item
  声明位置：\passthrough{\lstinline!include/unitree/robot/go2/sport/sport\_client.hpp:46!}
\end{itemize}

\textbf{说明}

\begin{itemize}
\tightlist
\item
  示例是局部调用片段；\passthrough{\lstinline!obj!}
  和参数变量表示已经按上层业务校验并准备好的对象。
\item
  该条目可以被编辑器和类型检查器识别，但当前运行时可能在属性查找阶段直接失败。
\item
  该方法属于运动命令。方法名包含
  \passthrough{\lstinline!stop!}、\passthrough{\lstinline!damp!} 或
  \passthrough{\lstinline!zero!} 也不代表它是独立物理急停。
\end{itemize}

\textbf{用法}

\begin{lstlisting}[language=Python]
from typing import TYPE_CHECKING

if TYPE_CHECKING:
    def planned_move(obj: SportClient, vx: float, vy: float, vyaw: float) -> int:
        return obj.move(vx=vx, vy=vy, vyaw=vyaw)
\end{lstlisting}

\begin{quote}
{[}!CAUTION{]} 上面的 \passthrough{\lstinline!TYPE\_CHECKING!}
示例只表达计划调用形态。该分支在普通 Python
运行时为假，不代表当前扩展能执行此接口。
\end{quote}

\hypertarget{unitree-sdk2-cpp-robot-go2-sportclient-sit-1}{%
\paragraph{\texorpdfstring{\texttt{unitree\_sdk2\_cpp.robot.go2.SportClient.sit}}{unitree\_sdk2\_cpp.robot.go2.SportClient.sit}}\label{unitree-sdk2-cpp-robot-go2-sportclient-sit-1}}

对应 C++ SDK 操作
\passthrough{\lstinline!Sit()!}。上游头文件没有可直接生成的业务说明时，本参考只保证签名映射，精确语义需查目标型号协议。

\textbf{签名}

\begin{lstlisting}[language=Python]
def sit(self) -> int
\end{lstlisting}

\textbf{可用性与安全性}

\textbf{\passthrough{\lstinline!SIGNATURE\_ONLY!} /
\passthrough{\lstinline!MOTION\_COMMAND!}}：仅用于补全和静态检查。当前二进制没有该
Python 方法；不得按可执行运动接口使用。

\textbf{参数}

无显式参数。实例方法中的 \passthrough{\lstinline!self!} 由 Python
自动传入。

\textbf{返回值}

\begin{longtable}[]{@{}
  >{\raggedright\arraybackslash}p{(\columnwidth - 4\tabcolsep) * \real{0.18}}
  >{\raggedright\arraybackslash}p{(\columnwidth - 4\tabcolsep) * \real{0.20}}
  >{\raggedright\arraybackslash}p{(\columnwidth - 4\tabcolsep) * \real{0.62}}@{}}
\toprule
位置 & 类型 & 含义 \\
\midrule
\endhead
返回值 & \passthrough{\lstinline!int!} & SDK 状态码；Unitree 示例通常以
\passthrough{\lstinline!0!} 表示成功，非零值需按具体服务定义解释。 \\
\bottomrule
\end{longtable}

\textbf{对应 C++}

\begin{itemize}
\tightlist
\item
  类：\passthrough{\lstinline!unitree::robot::go2::SportClient!}
\item
  签名：\passthrough{\lstinline!Sit()!}
\item
  绑定策略：\passthrough{\lstinline!DIRECT!}
\item
  声明位置：\passthrough{\lstinline!include/unitree/robot/go2/sport/sport\_client.hpp:47!}
\end{itemize}

\textbf{说明}

\begin{itemize}
\tightlist
\item
  示例是局部调用片段；\passthrough{\lstinline!obj!}
  和参数变量表示已经按上层业务校验并准备好的对象。
\item
  该条目可以被编辑器和类型检查器识别，但当前运行时可能在属性查找阶段直接失败。
\item
  该方法属于运动命令。方法名包含
  \passthrough{\lstinline!stop!}、\passthrough{\lstinline!damp!} 或
  \passthrough{\lstinline!zero!} 也不代表它是独立物理急停。
\end{itemize}

\textbf{用法}

\begin{lstlisting}[language=Python]
from typing import TYPE_CHECKING

if TYPE_CHECKING:
    def planned_sit(obj: SportClient) -> int:
        return obj.sit()
\end{lstlisting}

\begin{quote}
{[}!CAUTION{]} 上面的 \passthrough{\lstinline!TYPE\_CHECKING!}
示例只表达计划调用形态。该分支在普通 Python
运行时为假，不代表当前扩展能执行此接口。
\end{quote}

\hypertarget{unitree-sdk2-cpp-robot-go2-sportclient-rise-sit-1}{%
\paragraph{\texorpdfstring{\texttt{unitree\_sdk2\_cpp.robot.go2.SportClient.rise\_sit}}{unitree\_sdk2\_cpp.robot.go2.SportClient.rise\_sit}}\label{unitree-sdk2-cpp-robot-go2-sportclient-rise-sit-1}}

对应 C++ SDK 操作
\passthrough{\lstinline!RiseSit()!}。上游头文件没有可直接生成的业务说明时，本参考只保证签名映射，精确语义需查目标型号协议。

\textbf{签名}

\begin{lstlisting}[language=Python]
def rise_sit(self) -> int
\end{lstlisting}

\textbf{可用性与安全性}

\textbf{\passthrough{\lstinline!SIGNATURE\_ONLY!} /
\passthrough{\lstinline!MOTION\_COMMAND!}}：仅用于补全和静态检查。当前二进制没有该
Python 方法；不得按可执行运动接口使用。

\textbf{参数}

无显式参数。实例方法中的 \passthrough{\lstinline!self!} 由 Python
自动传入。

\textbf{返回值}

\begin{longtable}[]{@{}
  >{\raggedright\arraybackslash}p{(\columnwidth - 4\tabcolsep) * \real{0.18}}
  >{\raggedright\arraybackslash}p{(\columnwidth - 4\tabcolsep) * \real{0.20}}
  >{\raggedright\arraybackslash}p{(\columnwidth - 4\tabcolsep) * \real{0.62}}@{}}
\toprule
位置 & 类型 & 含义 \\
\midrule
\endhead
返回值 & \passthrough{\lstinline!int!} & SDK 状态码；Unitree 示例通常以
\passthrough{\lstinline!0!} 表示成功，非零值需按具体服务定义解释。 \\
\bottomrule
\end{longtable}

\textbf{对应 C++}

\begin{itemize}
\tightlist
\item
  类：\passthrough{\lstinline!unitree::robot::go2::SportClient!}
\item
  签名：\passthrough{\lstinline!RiseSit()!}
\item
  绑定策略：\passthrough{\lstinline!DIRECT!}
\item
  声明位置：\passthrough{\lstinline!include/unitree/robot/go2/sport/sport\_client.hpp:48!}
\end{itemize}

\textbf{说明}

\begin{itemize}
\tightlist
\item
  示例是局部调用片段；\passthrough{\lstinline!obj!}
  和参数变量表示已经按上层业务校验并准备好的对象。
\item
  该条目可以被编辑器和类型检查器识别，但当前运行时可能在属性查找阶段直接失败。
\item
  该方法属于运动命令。方法名包含
  \passthrough{\lstinline!stop!}、\passthrough{\lstinline!damp!} 或
  \passthrough{\lstinline!zero!} 也不代表它是独立物理急停。
\end{itemize}

\textbf{用法}

\begin{lstlisting}[language=Python]
from typing import TYPE_CHECKING

if TYPE_CHECKING:
    def planned_rise_sit(obj: SportClient) -> int:
        return obj.rise_sit()
\end{lstlisting}

\begin{quote}
{[}!CAUTION{]} 上面的 \passthrough{\lstinline!TYPE\_CHECKING!}
示例只表达计划调用形态。该分支在普通 Python
运行时为假，不代表当前扩展能执行此接口。
\end{quote}

\hypertarget{unitree-sdk2-cpp-robot-go2-sportclient-speed-level-1}{%
\paragraph{\texorpdfstring{\texttt{unitree\_sdk2\_cpp.robot.go2.SportClient.speed\_level}}{unitree\_sdk2\_cpp.robot.go2.SportClient.speed\_level}}\label{unitree-sdk2-cpp-robot-go2-sportclient-speed-level-1}}

对应 C++ SDK 操作
\passthrough{\lstinline!SpeedLevel(int)!}。上游头文件没有可直接生成的业务说明时，本参考只保证签名映射，精确语义需查目标型号协议。

\textbf{签名}

\begin{lstlisting}[language=Python]
def speed_level(self, level: int) -> int
\end{lstlisting}

\textbf{可用性与安全性}

\textbf{\passthrough{\lstinline!SIGNATURE\_ONLY!} /
\passthrough{\lstinline!MOTION\_COMMAND!}}：仅用于补全和静态检查。当前二进制没有该
Python 方法；不得按可执行运动接口使用。

\textbf{参数}

\begin{longtable}[]{@{}
  >{\raggedright\arraybackslash}p{(\columnwidth - 6\tabcolsep) * \real{0.18}}
  >{\raggedright\arraybackslash}p{(\columnwidth - 6\tabcolsep) * \real{0.24}}
  >{\raggedright\arraybackslash}p{(\columnwidth - 6\tabcolsep) * \real{0.18}}
  >{\raggedright\arraybackslash}p{(\columnwidth - 6\tabcolsep) * \real{0.40}}@{}}
\toprule
名称 & Python 类型 & 默认值 & 含义 \\
\midrule
\endhead
\passthrough{\lstinline!level!} & \passthrough{\lstinline!int!} & 必填 &
传给该接口的 \passthrough{\lstinline!level!} 参数，Python 类型为
\passthrough{\lstinline!int!}。精确含义、单位、范围和枚举值需查目标型号对应头文件与协议。
对应 C++ 参数 \passthrough{\lstinline!level: int!}。 \\
\bottomrule
\end{longtable}

\textbf{返回值}

\begin{longtable}[]{@{}
  >{\raggedright\arraybackslash}p{(\columnwidth - 4\tabcolsep) * \real{0.18}}
  >{\raggedright\arraybackslash}p{(\columnwidth - 4\tabcolsep) * \real{0.20}}
  >{\raggedright\arraybackslash}p{(\columnwidth - 4\tabcolsep) * \real{0.62}}@{}}
\toprule
位置 & 类型 & 含义 \\
\midrule
\endhead
返回值 & \passthrough{\lstinline!int!} & SDK 状态码；Unitree 示例通常以
\passthrough{\lstinline!0!} 表示成功，非零值需按具体服务定义解释。 \\
\bottomrule
\end{longtable}

\textbf{对应 C++}

\begin{itemize}
\tightlist
\item
  类：\passthrough{\lstinline!unitree::robot::go2::SportClient!}
\item
  签名：\passthrough{\lstinline!SpeedLevel(int)!}
\item
  绑定策略：\passthrough{\lstinline!DIRECT!}
\item
  声明位置：\passthrough{\lstinline!include/unitree/robot/go2/sport/sport\_client.hpp:49!}
\end{itemize}

\textbf{说明}

\begin{itemize}
\tightlist
\item
  示例是局部调用片段；\passthrough{\lstinline!obj!}
  和参数变量表示已经按上层业务校验并准备好的对象。
\item
  该条目可以被编辑器和类型检查器识别，但当前运行时可能在属性查找阶段直接失败。
\item
  该方法属于运动命令。方法名包含
  \passthrough{\lstinline!stop!}、\passthrough{\lstinline!damp!} 或
  \passthrough{\lstinline!zero!} 也不代表它是独立物理急停。
\end{itemize}

\textbf{用法}

\begin{lstlisting}[language=Python]
from typing import TYPE_CHECKING

if TYPE_CHECKING:
    def planned_speed_level(obj: SportClient, level: int) -> int:
        return obj.speed_level(level=level)
\end{lstlisting}

\begin{quote}
{[}!CAUTION{]} 上面的 \passthrough{\lstinline!TYPE\_CHECKING!}
示例只表达计划调用形态。该分支在普通 Python
运行时为假，不代表当前扩展能执行此接口。
\end{quote}

\hypertarget{unitree-sdk2-cpp-robot-go2-sportclient-hello-1}{%
\paragraph{\texorpdfstring{\texttt{unitree\_sdk2\_cpp.robot.go2.SportClient.hello}}{unitree\_sdk2\_cpp.robot.go2.SportClient.hello}}\label{unitree-sdk2-cpp-robot-go2-sportclient-hello-1}}

对应 C++ SDK 操作
\passthrough{\lstinline!Hello()!}。上游头文件没有可直接生成的业务说明时，本参考只保证签名映射，精确语义需查目标型号协议。

\textbf{签名}

\begin{lstlisting}[language=Python]
def hello(self) -> int
\end{lstlisting}

\textbf{可用性与安全性}

\textbf{\passthrough{\lstinline!SIGNATURE\_ONLY!} /
\passthrough{\lstinline!MOTION\_COMMAND!}}：仅用于补全和静态检查。当前二进制没有该
Python 方法；不得按可执行运动接口使用。

\textbf{参数}

无显式参数。实例方法中的 \passthrough{\lstinline!self!} 由 Python
自动传入。

\textbf{返回值}

\begin{longtable}[]{@{}
  >{\raggedright\arraybackslash}p{(\columnwidth - 4\tabcolsep) * \real{0.18}}
  >{\raggedright\arraybackslash}p{(\columnwidth - 4\tabcolsep) * \real{0.20}}
  >{\raggedright\arraybackslash}p{(\columnwidth - 4\tabcolsep) * \real{0.62}}@{}}
\toprule
位置 & 类型 & 含义 \\
\midrule
\endhead
返回值 & \passthrough{\lstinline!int!} & SDK 状态码；Unitree 示例通常以
\passthrough{\lstinline!0!} 表示成功，非零值需按具体服务定义解释。 \\
\bottomrule
\end{longtable}

\textbf{对应 C++}

\begin{itemize}
\tightlist
\item
  类：\passthrough{\lstinline!unitree::robot::go2::SportClient!}
\item
  签名：\passthrough{\lstinline!Hello()!}
\item
  绑定策略：\passthrough{\lstinline!DIRECT!}
\item
  声明位置：\passthrough{\lstinline!include/unitree/robot/go2/sport/sport\_client.hpp:50!}
\end{itemize}

\textbf{说明}

\begin{itemize}
\tightlist
\item
  示例是局部调用片段；\passthrough{\lstinline!obj!}
  和参数变量表示已经按上层业务校验并准备好的对象。
\item
  该条目可以被编辑器和类型检查器识别，但当前运行时可能在属性查找阶段直接失败。
\item
  该方法属于运动命令。方法名包含
  \passthrough{\lstinline!stop!}、\passthrough{\lstinline!damp!} 或
  \passthrough{\lstinline!zero!} 也不代表它是独立物理急停。
\end{itemize}

\textbf{用法}

\begin{lstlisting}[language=Python]
from typing import TYPE_CHECKING

if TYPE_CHECKING:
    def planned_hello(obj: SportClient) -> int:
        return obj.hello()
\end{lstlisting}

\begin{quote}
{[}!CAUTION{]} 上面的 \passthrough{\lstinline!TYPE\_CHECKING!}
示例只表达计划调用形态。该分支在普通 Python
运行时为假，不代表当前扩展能执行此接口。
\end{quote}

\hypertarget{unitree-sdk2-cpp-robot-go2-sportclient-stretch-1}{%
\paragraph{\texorpdfstring{\texttt{unitree\_sdk2\_cpp.robot.go2.SportClient.stretch}}{unitree\_sdk2\_cpp.robot.go2.SportClient.stretch}}\label{unitree-sdk2-cpp-robot-go2-sportclient-stretch-1}}

对应 C++ SDK 操作
\passthrough{\lstinline!Stretch()!}。上游头文件没有可直接生成的业务说明时，本参考只保证签名映射，精确语义需查目标型号协议。

\textbf{签名}

\begin{lstlisting}[language=Python]
def stretch(self) -> int
\end{lstlisting}

\textbf{可用性与安全性}

\textbf{\passthrough{\lstinline!SIGNATURE\_ONLY!} /
\passthrough{\lstinline!MOTION\_COMMAND!}}：仅用于补全和静态检查。当前二进制没有该
Python 方法；不得按可执行运动接口使用。

\textbf{参数}

无显式参数。实例方法中的 \passthrough{\lstinline!self!} 由 Python
自动传入。

\textbf{返回值}

\begin{longtable}[]{@{}
  >{\raggedright\arraybackslash}p{(\columnwidth - 4\tabcolsep) * \real{0.18}}
  >{\raggedright\arraybackslash}p{(\columnwidth - 4\tabcolsep) * \real{0.20}}
  >{\raggedright\arraybackslash}p{(\columnwidth - 4\tabcolsep) * \real{0.62}}@{}}
\toprule
位置 & 类型 & 含义 \\
\midrule
\endhead
返回值 & \passthrough{\lstinline!int!} & SDK 状态码；Unitree 示例通常以
\passthrough{\lstinline!0!} 表示成功，非零值需按具体服务定义解释。 \\
\bottomrule
\end{longtable}

\textbf{对应 C++}

\begin{itemize}
\tightlist
\item
  类：\passthrough{\lstinline!unitree::robot::go2::SportClient!}
\item
  签名：\passthrough{\lstinline!Stretch()!}
\item
  绑定策略：\passthrough{\lstinline!DIRECT!}
\item
  声明位置：\passthrough{\lstinline!include/unitree/robot/go2/sport/sport\_client.hpp:51!}
\end{itemize}

\textbf{说明}

\begin{itemize}
\tightlist
\item
  示例是局部调用片段；\passthrough{\lstinline!obj!}
  和参数变量表示已经按上层业务校验并准备好的对象。
\item
  该条目可以被编辑器和类型检查器识别，但当前运行时可能在属性查找阶段直接失败。
\item
  该方法属于运动命令。方法名包含
  \passthrough{\lstinline!stop!}、\passthrough{\lstinline!damp!} 或
  \passthrough{\lstinline!zero!} 也不代表它是独立物理急停。
\end{itemize}

\textbf{用法}

\begin{lstlisting}[language=Python]
from typing import TYPE_CHECKING

if TYPE_CHECKING:
    def planned_stretch(obj: SportClient) -> int:
        return obj.stretch()
\end{lstlisting}

\begin{quote}
{[}!CAUTION{]} 上面的 \passthrough{\lstinline!TYPE\_CHECKING!}
示例只表达计划调用形态。该分支在普通 Python
运行时为假，不代表当前扩展能执行此接口。
\end{quote}

\hypertarget{unitree-sdk2-cpp-robot-go2-sportclient-switch-joystick-1}{%
\paragraph{\texorpdfstring{\texttt{unitree\_sdk2\_cpp.robot.go2.SportClient.switch\_joystick}}{unitree\_sdk2\_cpp.robot.go2.SportClient.switch\_joystick}}\label{unitree-sdk2-cpp-robot-go2-sportclient-switch-joystick-1}}

对应 C++ SDK 操作
\passthrough{\lstinline!SwitchJoystick(bool)!}。上游头文件没有可直接生成的业务说明时，本参考只保证签名映射，精确语义需查目标型号协议。

\textbf{签名}

\begin{lstlisting}[language=Python]
def switch_joystick(self, flag: bool) -> int
\end{lstlisting}

\textbf{可用性与安全性}

\textbf{\passthrough{\lstinline!SIGNATURE\_ONLY!} /
\passthrough{\lstinline!MOTION\_COMMAND!}}：仅用于补全和静态检查。当前二进制没有该
Python 方法；不得按可执行运动接口使用。

\textbf{参数}

\begin{longtable}[]{@{}
  >{\raggedright\arraybackslash}p{(\columnwidth - 6\tabcolsep) * \real{0.18}}
  >{\raggedright\arraybackslash}p{(\columnwidth - 6\tabcolsep) * \real{0.24}}
  >{\raggedright\arraybackslash}p{(\columnwidth - 6\tabcolsep) * \real{0.18}}
  >{\raggedright\arraybackslash}p{(\columnwidth - 6\tabcolsep) * \real{0.40}}@{}}
\toprule
名称 & Python 类型 & 默认值 & 含义 \\
\midrule
\endhead
\passthrough{\lstinline!flag!} & \passthrough{\lstinline!bool!} & 必填 &
布尔开关。\passthrough{\lstinline!True!} 和
\passthrough{\lstinline!False!} 的具体业务效果由当前方法定义。 对应 C++
参数 \passthrough{\lstinline!flag: bool!}。 只接受布尔语义。 \\
\bottomrule
\end{longtable}

\textbf{返回值}

\begin{longtable}[]{@{}
  >{\raggedright\arraybackslash}p{(\columnwidth - 4\tabcolsep) * \real{0.18}}
  >{\raggedright\arraybackslash}p{(\columnwidth - 4\tabcolsep) * \real{0.20}}
  >{\raggedright\arraybackslash}p{(\columnwidth - 4\tabcolsep) * \real{0.62}}@{}}
\toprule
位置 & 类型 & 含义 \\
\midrule
\endhead
返回值 & \passthrough{\lstinline!int!} & SDK 状态码；Unitree 示例通常以
\passthrough{\lstinline!0!} 表示成功，非零值需按具体服务定义解释。 \\
\bottomrule
\end{longtable}

\textbf{对应 C++}

\begin{itemize}
\tightlist
\item
  类：\passthrough{\lstinline!unitree::robot::go2::SportClient!}
\item
  签名：\passthrough{\lstinline!SwitchJoystick(bool)!}
\item
  绑定策略：\passthrough{\lstinline!DIRECT!}
\item
  声明位置：\passthrough{\lstinline!include/unitree/robot/go2/sport/sport\_client.hpp:52!}
\end{itemize}

\textbf{说明}

\begin{itemize}
\tightlist
\item
  示例是局部调用片段；\passthrough{\lstinline!obj!}
  和参数变量表示已经按上层业务校验并准备好的对象。
\item
  该条目可以被编辑器和类型检查器识别，但当前运行时可能在属性查找阶段直接失败。
\item
  该方法属于运动命令。方法名包含
  \passthrough{\lstinline!stop!}、\passthrough{\lstinline!damp!} 或
  \passthrough{\lstinline!zero!} 也不代表它是独立物理急停。
\end{itemize}

\textbf{用法}

\begin{lstlisting}[language=Python]
from typing import TYPE_CHECKING

if TYPE_CHECKING:
    def planned_switch_joystick(obj: SportClient, flag: bool) -> int:
        return obj.switch_joystick(flag=flag)
\end{lstlisting}

\begin{quote}
{[}!CAUTION{]} 上面的 \passthrough{\lstinline!TYPE\_CHECKING!}
示例只表达计划调用形态。该分支在普通 Python
运行时为假，不代表当前扩展能执行此接口。
\end{quote}

\hypertarget{unitree-sdk2-cpp-robot-go2-sportclient-content-1}{%
\paragraph{\texorpdfstring{\texttt{unitree\_sdk2\_cpp.robot.go2.SportClient.content}}{unitree\_sdk2\_cpp.robot.go2.SportClient.content}}\label{unitree-sdk2-cpp-robot-go2-sportclient-content-1}}

对应 C++ SDK 操作
\passthrough{\lstinline!Content()!}。上游头文件没有可直接生成的业务说明时，本参考只保证签名映射，精确语义需查目标型号协议。

\textbf{签名}

\begin{lstlisting}[language=Python]
def content(self) -> int
\end{lstlisting}

\textbf{可用性与安全性}

\textbf{\passthrough{\lstinline!SIGNATURE\_ONLY!} /
\passthrough{\lstinline!MOTION\_COMMAND!}}：仅用于补全和静态检查。当前二进制没有该
Python 方法；不得按可执行运动接口使用。

\textbf{参数}

无显式参数。实例方法中的 \passthrough{\lstinline!self!} 由 Python
自动传入。

\textbf{返回值}

\begin{longtable}[]{@{}
  >{\raggedright\arraybackslash}p{(\columnwidth - 4\tabcolsep) * \real{0.18}}
  >{\raggedright\arraybackslash}p{(\columnwidth - 4\tabcolsep) * \real{0.20}}
  >{\raggedright\arraybackslash}p{(\columnwidth - 4\tabcolsep) * \real{0.62}}@{}}
\toprule
位置 & 类型 & 含义 \\
\midrule
\endhead
返回值 & \passthrough{\lstinline!int!} & SDK 状态码；Unitree 示例通常以
\passthrough{\lstinline!0!} 表示成功，非零值需按具体服务定义解释。 \\
\bottomrule
\end{longtable}

\textbf{对应 C++}

\begin{itemize}
\tightlist
\item
  类：\passthrough{\lstinline!unitree::robot::go2::SportClient!}
\item
  签名：\passthrough{\lstinline!Content()!}
\item
  绑定策略：\passthrough{\lstinline!DIRECT!}
\item
  声明位置：\passthrough{\lstinline!include/unitree/robot/go2/sport/sport\_client.hpp:53!}
\end{itemize}

\textbf{说明}

\begin{itemize}
\tightlist
\item
  示例是局部调用片段；\passthrough{\lstinline!obj!}
  和参数变量表示已经按上层业务校验并准备好的对象。
\item
  该条目可以被编辑器和类型检查器识别，但当前运行时可能在属性查找阶段直接失败。
\item
  该方法属于运动命令。方法名包含
  \passthrough{\lstinline!stop!}、\passthrough{\lstinline!damp!} 或
  \passthrough{\lstinline!zero!} 也不代表它是独立物理急停。
\end{itemize}

\textbf{用法}

\begin{lstlisting}[language=Python]
from typing import TYPE_CHECKING

if TYPE_CHECKING:
    def planned_content(obj: SportClient) -> int:
        return obj.content()
\end{lstlisting}

\begin{quote}
{[}!CAUTION{]} 上面的 \passthrough{\lstinline!TYPE\_CHECKING!}
示例只表达计划调用形态。该分支在普通 Python
运行时为假，不代表当前扩展能执行此接口。
\end{quote}

\hypertarget{unitree-sdk2-cpp-robot-go2-sportclient-heart-1}{%
\paragraph{\texorpdfstring{\texttt{unitree\_sdk2\_cpp.robot.go2.SportClient.heart}}{unitree\_sdk2\_cpp.robot.go2.SportClient.heart}}\label{unitree-sdk2-cpp-robot-go2-sportclient-heart-1}}

对应 C++ SDK 操作
\passthrough{\lstinline!Heart()!}。上游头文件没有可直接生成的业务说明时，本参考只保证签名映射，精确语义需查目标型号协议。

\textbf{签名}

\begin{lstlisting}[language=Python]
def heart(self) -> int
\end{lstlisting}

\textbf{可用性与安全性}

\textbf{\passthrough{\lstinline!SIGNATURE\_ONLY!} /
\passthrough{\lstinline!MOTION\_COMMAND!}}：仅用于补全和静态检查。当前二进制没有该
Python 方法；不得按可执行运动接口使用。

\textbf{参数}

无显式参数。实例方法中的 \passthrough{\lstinline!self!} 由 Python
自动传入。

\textbf{返回值}

\begin{longtable}[]{@{}
  >{\raggedright\arraybackslash}p{(\columnwidth - 4\tabcolsep) * \real{0.18}}
  >{\raggedright\arraybackslash}p{(\columnwidth - 4\tabcolsep) * \real{0.20}}
  >{\raggedright\arraybackslash}p{(\columnwidth - 4\tabcolsep) * \real{0.62}}@{}}
\toprule
位置 & 类型 & 含义 \\
\midrule
\endhead
返回值 & \passthrough{\lstinline!int!} & SDK 状态码；Unitree 示例通常以
\passthrough{\lstinline!0!} 表示成功，非零值需按具体服务定义解释。 \\
\bottomrule
\end{longtable}

\textbf{对应 C++}

\begin{itemize}
\tightlist
\item
  类：\passthrough{\lstinline!unitree::robot::go2::SportClient!}
\item
  签名：\passthrough{\lstinline!Heart()!}
\item
  绑定策略：\passthrough{\lstinline!DIRECT!}
\item
  声明位置：\passthrough{\lstinline!include/unitree/robot/go2/sport/sport\_client.hpp:54!}
\end{itemize}

\textbf{说明}

\begin{itemize}
\tightlist
\item
  示例是局部调用片段；\passthrough{\lstinline!obj!}
  和参数变量表示已经按上层业务校验并准备好的对象。
\item
  该条目可以被编辑器和类型检查器识别，但当前运行时可能在属性查找阶段直接失败。
\item
  该方法属于运动命令。方法名包含
  \passthrough{\lstinline!stop!}、\passthrough{\lstinline!damp!} 或
  \passthrough{\lstinline!zero!} 也不代表它是独立物理急停。
\end{itemize}

\textbf{用法}

\begin{lstlisting}[language=Python]
from typing import TYPE_CHECKING

if TYPE_CHECKING:
    def planned_heart(obj: SportClient) -> int:
        return obj.heart()
\end{lstlisting}

\begin{quote}
{[}!CAUTION{]} 上面的 \passthrough{\lstinline!TYPE\_CHECKING!}
示例只表达计划调用形态。该分支在普通 Python
运行时为假，不代表当前扩展能执行此接口。
\end{quote}

\hypertarget{unitree-sdk2-cpp-robot-go2-sportclient-pose-1}{%
\paragraph{\texorpdfstring{\texttt{unitree\_sdk2\_cpp.robot.go2.SportClient.pose}}{unitree\_sdk2\_cpp.robot.go2.SportClient.pose}}\label{unitree-sdk2-cpp-robot-go2-sportclient-pose-1}}

对应 C++ SDK 操作
\passthrough{\lstinline!Pose(bool)!}。上游头文件没有可直接生成的业务说明时，本参考只保证签名映射，精确语义需查目标型号协议。

\textbf{签名}

\begin{lstlisting}[language=Python]
def pose(self, flag: bool) -> int
\end{lstlisting}

\textbf{可用性与安全性}

\textbf{\passthrough{\lstinline!SIGNATURE\_ONLY!} /
\passthrough{\lstinline!MOTION\_COMMAND!}}：仅用于补全和静态检查。当前二进制没有该
Python 方法；不得按可执行运动接口使用。

\textbf{参数}

\begin{longtable}[]{@{}
  >{\raggedright\arraybackslash}p{(\columnwidth - 6\tabcolsep) * \real{0.18}}
  >{\raggedright\arraybackslash}p{(\columnwidth - 6\tabcolsep) * \real{0.24}}
  >{\raggedright\arraybackslash}p{(\columnwidth - 6\tabcolsep) * \real{0.18}}
  >{\raggedright\arraybackslash}p{(\columnwidth - 6\tabcolsep) * \real{0.40}}@{}}
\toprule
名称 & Python 类型 & 默认值 & 含义 \\
\midrule
\endhead
\passthrough{\lstinline!flag!} & \passthrough{\lstinline!bool!} & 必填 &
布尔开关。\passthrough{\lstinline!True!} 和
\passthrough{\lstinline!False!} 的具体业务效果由当前方法定义。 对应 C++
参数 \passthrough{\lstinline!flag: bool!}。 只接受布尔语义。 \\
\bottomrule
\end{longtable}

\textbf{返回值}

\begin{longtable}[]{@{}
  >{\raggedright\arraybackslash}p{(\columnwidth - 4\tabcolsep) * \real{0.18}}
  >{\raggedright\arraybackslash}p{(\columnwidth - 4\tabcolsep) * \real{0.20}}
  >{\raggedright\arraybackslash}p{(\columnwidth - 4\tabcolsep) * \real{0.62}}@{}}
\toprule
位置 & 类型 & 含义 \\
\midrule
\endhead
返回值 & \passthrough{\lstinline!int!} & SDK 状态码；Unitree 示例通常以
\passthrough{\lstinline!0!} 表示成功，非零值需按具体服务定义解释。 \\
\bottomrule
\end{longtable}

\textbf{对应 C++}

\begin{itemize}
\tightlist
\item
  类：\passthrough{\lstinline!unitree::robot::go2::SportClient!}
\item
  签名：\passthrough{\lstinline!Pose(bool)!}
\item
  绑定策略：\passthrough{\lstinline!DIRECT!}
\item
  声明位置：\passthrough{\lstinline!include/unitree/robot/go2/sport/sport\_client.hpp:55!}
\end{itemize}

\textbf{说明}

\begin{itemize}
\tightlist
\item
  示例是局部调用片段；\passthrough{\lstinline!obj!}
  和参数变量表示已经按上层业务校验并准备好的对象。
\item
  该条目可以被编辑器和类型检查器识别，但当前运行时可能在属性查找阶段直接失败。
\item
  该方法属于运动命令。方法名包含
  \passthrough{\lstinline!stop!}、\passthrough{\lstinline!damp!} 或
  \passthrough{\lstinline!zero!} 也不代表它是独立物理急停。
\end{itemize}

\textbf{用法}

\begin{lstlisting}[language=Python]
from typing import TYPE_CHECKING

if TYPE_CHECKING:
    def planned_pose(obj: SportClient, flag: bool) -> int:
        return obj.pose(flag=flag)
\end{lstlisting}

\begin{quote}
{[}!CAUTION{]} 上面的 \passthrough{\lstinline!TYPE\_CHECKING!}
示例只表达计划调用形态。该分支在普通 Python
运行时为假，不代表当前扩展能执行此接口。
\end{quote}

\hypertarget{unitree-sdk2-cpp-robot-go2-sportclient-scrape-1}{%
\paragraph{\texorpdfstring{\texttt{unitree\_sdk2\_cpp.robot.go2.SportClient.scrape}}{unitree\_sdk2\_cpp.robot.go2.SportClient.scrape}}\label{unitree-sdk2-cpp-robot-go2-sportclient-scrape-1}}

对应 C++ SDK 操作
\passthrough{\lstinline!Scrape()!}。上游头文件没有可直接生成的业务说明时，本参考只保证签名映射，精确语义需查目标型号协议。

\textbf{签名}

\begin{lstlisting}[language=Python]
def scrape(self) -> int
\end{lstlisting}

\textbf{可用性与安全性}

\textbf{\passthrough{\lstinline!SIGNATURE\_ONLY!} /
\passthrough{\lstinline!MOTION\_COMMAND!}}：仅用于补全和静态检查。当前二进制没有该
Python 方法；不得按可执行运动接口使用。

\textbf{参数}

无显式参数。实例方法中的 \passthrough{\lstinline!self!} 由 Python
自动传入。

\textbf{返回值}

\begin{longtable}[]{@{}
  >{\raggedright\arraybackslash}p{(\columnwidth - 4\tabcolsep) * \real{0.18}}
  >{\raggedright\arraybackslash}p{(\columnwidth - 4\tabcolsep) * \real{0.20}}
  >{\raggedright\arraybackslash}p{(\columnwidth - 4\tabcolsep) * \real{0.62}}@{}}
\toprule
位置 & 类型 & 含义 \\
\midrule
\endhead
返回值 & \passthrough{\lstinline!int!} & SDK 状态码；Unitree 示例通常以
\passthrough{\lstinline!0!} 表示成功，非零值需按具体服务定义解释。 \\
\bottomrule
\end{longtable}

\textbf{对应 C++}

\begin{itemize}
\tightlist
\item
  类：\passthrough{\lstinline!unitree::robot::go2::SportClient!}
\item
  签名：\passthrough{\lstinline!Scrape()!}
\item
  绑定策略：\passthrough{\lstinline!DIRECT!}
\item
  声明位置：\passthrough{\lstinline!include/unitree/robot/go2/sport/sport\_client.hpp:56!}
\end{itemize}

\textbf{说明}

\begin{itemize}
\tightlist
\item
  示例是局部调用片段；\passthrough{\lstinline!obj!}
  和参数变量表示已经按上层业务校验并准备好的对象。
\item
  该条目可以被编辑器和类型检查器识别，但当前运行时可能在属性查找阶段直接失败。
\item
  该方法属于运动命令。方法名包含
  \passthrough{\lstinline!stop!}、\passthrough{\lstinline!damp!} 或
  \passthrough{\lstinline!zero!} 也不代表它是独立物理急停。
\end{itemize}

\textbf{用法}

\begin{lstlisting}[language=Python]
from typing import TYPE_CHECKING

if TYPE_CHECKING:
    def planned_scrape(obj: SportClient) -> int:
        return obj.scrape()
\end{lstlisting}

\begin{quote}
{[}!CAUTION{]} 上面的 \passthrough{\lstinline!TYPE\_CHECKING!}
示例只表达计划调用形态。该分支在普通 Python
运行时为假，不代表当前扩展能执行此接口。
\end{quote}

\hypertarget{unitree-sdk2-cpp-robot-go2-sportclient-front-flip-1}{%
\paragraph{\texorpdfstring{\texttt{unitree\_sdk2\_cpp.robot.go2.SportClient.front\_flip}}{unitree\_sdk2\_cpp.robot.go2.SportClient.front\_flip}}\label{unitree-sdk2-cpp-robot-go2-sportclient-front-flip-1}}

对应 C++ SDK 操作
\passthrough{\lstinline!FrontFlip()!}。上游头文件没有可直接生成的业务说明时，本参考只保证签名映射，精确语义需查目标型号协议。

\textbf{签名}

\begin{lstlisting}[language=Python]
def front_flip(self) -> int
\end{lstlisting}

\textbf{可用性与安全性}

\textbf{\passthrough{\lstinline!SIGNATURE\_ONLY!} /
\passthrough{\lstinline!MOTION\_COMMAND!}}：仅用于补全和静态检查。当前二进制没有该
Python 方法；不得按可执行运动接口使用。

\textbf{参数}

无显式参数。实例方法中的 \passthrough{\lstinline!self!} 由 Python
自动传入。

\textbf{返回值}

\begin{longtable}[]{@{}
  >{\raggedright\arraybackslash}p{(\columnwidth - 4\tabcolsep) * \real{0.18}}
  >{\raggedright\arraybackslash}p{(\columnwidth - 4\tabcolsep) * \real{0.20}}
  >{\raggedright\arraybackslash}p{(\columnwidth - 4\tabcolsep) * \real{0.62}}@{}}
\toprule
位置 & 类型 & 含义 \\
\midrule
\endhead
返回值 & \passthrough{\lstinline!int!} & SDK 状态码；Unitree 示例通常以
\passthrough{\lstinline!0!} 表示成功，非零值需按具体服务定义解释。 \\
\bottomrule
\end{longtable}

\textbf{对应 C++}

\begin{itemize}
\tightlist
\item
  类：\passthrough{\lstinline!unitree::robot::go2::SportClient!}
\item
  签名：\passthrough{\lstinline!FrontFlip()!}
\item
  绑定策略：\passthrough{\lstinline!DIRECT!}
\item
  声明位置：\passthrough{\lstinline!include/unitree/robot/go2/sport/sport\_client.hpp:57!}
\end{itemize}

\textbf{说明}

\begin{itemize}
\tightlist
\item
  示例是局部调用片段；\passthrough{\lstinline!obj!}
  和参数变量表示已经按上层业务校验并准备好的对象。
\item
  该条目可以被编辑器和类型检查器识别，但当前运行时可能在属性查找阶段直接失败。
\item
  该方法属于运动命令。方法名包含
  \passthrough{\lstinline!stop!}、\passthrough{\lstinline!damp!} 或
  \passthrough{\lstinline!zero!} 也不代表它是独立物理急停。
\end{itemize}

\textbf{用法}

\begin{lstlisting}[language=Python]
from typing import TYPE_CHECKING

if TYPE_CHECKING:
    def planned_front_flip(obj: SportClient) -> int:
        return obj.front_flip()
\end{lstlisting}

\begin{quote}
{[}!CAUTION{]} 上面的 \passthrough{\lstinline!TYPE\_CHECKING!}
示例只表达计划调用形态。该分支在普通 Python
运行时为假，不代表当前扩展能执行此接口。
\end{quote}

\hypertarget{unitree-sdk2-cpp-robot-go2-sportclient-front-jump-1}{%
\paragraph{\texorpdfstring{\texttt{unitree\_sdk2\_cpp.robot.go2.SportClient.front\_jump}}{unitree\_sdk2\_cpp.robot.go2.SportClient.front\_jump}}\label{unitree-sdk2-cpp-robot-go2-sportclient-front-jump-1}}

对应 C++ SDK 操作
\passthrough{\lstinline!FrontJump()!}。上游头文件没有可直接生成的业务说明时，本参考只保证签名映射，精确语义需查目标型号协议。

\textbf{签名}

\begin{lstlisting}[language=Python]
def front_jump(self) -> int
\end{lstlisting}

\textbf{可用性与安全性}

\textbf{\passthrough{\lstinline!SIGNATURE\_ONLY!} /
\passthrough{\lstinline!MOTION\_COMMAND!}}：仅用于补全和静态检查。当前二进制没有该
Python 方法；不得按可执行运动接口使用。

\textbf{参数}

无显式参数。实例方法中的 \passthrough{\lstinline!self!} 由 Python
自动传入。

\textbf{返回值}

\begin{longtable}[]{@{}
  >{\raggedright\arraybackslash}p{(\columnwidth - 4\tabcolsep) * \real{0.18}}
  >{\raggedright\arraybackslash}p{(\columnwidth - 4\tabcolsep) * \real{0.20}}
  >{\raggedright\arraybackslash}p{(\columnwidth - 4\tabcolsep) * \real{0.62}}@{}}
\toprule
位置 & 类型 & 含义 \\
\midrule
\endhead
返回值 & \passthrough{\lstinline!int!} & SDK 状态码；Unitree 示例通常以
\passthrough{\lstinline!0!} 表示成功，非零值需按具体服务定义解释。 \\
\bottomrule
\end{longtable}

\textbf{对应 C++}

\begin{itemize}
\tightlist
\item
  类：\passthrough{\lstinline!unitree::robot::go2::SportClient!}
\item
  签名：\passthrough{\lstinline!FrontJump()!}
\item
  绑定策略：\passthrough{\lstinline!DIRECT!}
\item
  声明位置：\passthrough{\lstinline!include/unitree/robot/go2/sport/sport\_client.hpp:58!}
\end{itemize}

\textbf{说明}

\begin{itemize}
\tightlist
\item
  示例是局部调用片段；\passthrough{\lstinline!obj!}
  和参数变量表示已经按上层业务校验并准备好的对象。
\item
  该条目可以被编辑器和类型检查器识别，但当前运行时可能在属性查找阶段直接失败。
\item
  该方法属于运动命令。方法名包含
  \passthrough{\lstinline!stop!}、\passthrough{\lstinline!damp!} 或
  \passthrough{\lstinline!zero!} 也不代表它是独立物理急停。
\end{itemize}

\textbf{用法}

\begin{lstlisting}[language=Python]
from typing import TYPE_CHECKING

if TYPE_CHECKING:
    def planned_front_jump(obj: SportClient) -> int:
        return obj.front_jump()
\end{lstlisting}

\begin{quote}
{[}!CAUTION{]} 上面的 \passthrough{\lstinline!TYPE\_CHECKING!}
示例只表达计划调用形态。该分支在普通 Python
运行时为假，不代表当前扩展能执行此接口。
\end{quote}

\hypertarget{unitree-sdk2-cpp-robot-go2-sportclient-front-pounce-1}{%
\paragraph{\texorpdfstring{\texttt{unitree\_sdk2\_cpp.robot.go2.SportClient.front\_pounce}}{unitree\_sdk2\_cpp.robot.go2.SportClient.front\_pounce}}\label{unitree-sdk2-cpp-robot-go2-sportclient-front-pounce-1}}

对应 C++ SDK 操作
\passthrough{\lstinline!FrontPounce()!}。上游头文件没有可直接生成的业务说明时，本参考只保证签名映射，精确语义需查目标型号协议。

\textbf{签名}

\begin{lstlisting}[language=Python]
def front_pounce(self) -> int
\end{lstlisting}

\textbf{可用性与安全性}

\textbf{\passthrough{\lstinline!SIGNATURE\_ONLY!} /
\passthrough{\lstinline!MOTION\_COMMAND!}}：仅用于补全和静态检查。当前二进制没有该
Python 方法；不得按可执行运动接口使用。

\textbf{参数}

无显式参数。实例方法中的 \passthrough{\lstinline!self!} 由 Python
自动传入。

\textbf{返回值}

\begin{longtable}[]{@{}
  >{\raggedright\arraybackslash}p{(\columnwidth - 4\tabcolsep) * \real{0.18}}
  >{\raggedright\arraybackslash}p{(\columnwidth - 4\tabcolsep) * \real{0.20}}
  >{\raggedright\arraybackslash}p{(\columnwidth - 4\tabcolsep) * \real{0.62}}@{}}
\toprule
位置 & 类型 & 含义 \\
\midrule
\endhead
返回值 & \passthrough{\lstinline!int!} & SDK 状态码；Unitree 示例通常以
\passthrough{\lstinline!0!} 表示成功，非零值需按具体服务定义解释。 \\
\bottomrule
\end{longtable}

\textbf{对应 C++}

\begin{itemize}
\tightlist
\item
  类：\passthrough{\lstinline!unitree::robot::go2::SportClient!}
\item
  签名：\passthrough{\lstinline!FrontPounce()!}
\item
  绑定策略：\passthrough{\lstinline!DIRECT!}
\item
  声明位置：\passthrough{\lstinline!include/unitree/robot/go2/sport/sport\_client.hpp:59!}
\end{itemize}

\textbf{说明}

\begin{itemize}
\tightlist
\item
  示例是局部调用片段；\passthrough{\lstinline!obj!}
  和参数变量表示已经按上层业务校验并准备好的对象。
\item
  该条目可以被编辑器和类型检查器识别，但当前运行时可能在属性查找阶段直接失败。
\item
  该方法属于运动命令。方法名包含
  \passthrough{\lstinline!stop!}、\passthrough{\lstinline!damp!} 或
  \passthrough{\lstinline!zero!} 也不代表它是独立物理急停。
\end{itemize}

\textbf{用法}

\begin{lstlisting}[language=Python]
from typing import TYPE_CHECKING

if TYPE_CHECKING:
    def planned_front_pounce(obj: SportClient) -> int:
        return obj.front_pounce()
\end{lstlisting}

\begin{quote}
{[}!CAUTION{]} 上面的 \passthrough{\lstinline!TYPE\_CHECKING!}
示例只表达计划调用形态。该分支在普通 Python
运行时为假，不代表当前扩展能执行此接口。
\end{quote}

\hypertarget{unitree-sdk2-cpp-robot-go2-sportclient-dance1-1}{%
\paragraph{\texorpdfstring{\texttt{unitree\_sdk2\_cpp.robot.go2.SportClient.dance1}}{unitree\_sdk2\_cpp.robot.go2.SportClient.dance1}}\label{unitree-sdk2-cpp-robot-go2-sportclient-dance1-1}}

对应 C++ SDK 操作
\passthrough{\lstinline!Dance1()!}。上游头文件没有可直接生成的业务说明时，本参考只保证签名映射，精确语义需查目标型号协议。

\textbf{签名}

\begin{lstlisting}[language=Python]
def dance1(self) -> int
\end{lstlisting}

\textbf{可用性与安全性}

\textbf{\passthrough{\lstinline!SIGNATURE\_ONLY!} /
\passthrough{\lstinline!MOTION\_COMMAND!}}：仅用于补全和静态检查。当前二进制没有该
Python 方法；不得按可执行运动接口使用。

\textbf{参数}

无显式参数。实例方法中的 \passthrough{\lstinline!self!} 由 Python
自动传入。

\textbf{返回值}

\begin{longtable}[]{@{}
  >{\raggedright\arraybackslash}p{(\columnwidth - 4\tabcolsep) * \real{0.18}}
  >{\raggedright\arraybackslash}p{(\columnwidth - 4\tabcolsep) * \real{0.20}}
  >{\raggedright\arraybackslash}p{(\columnwidth - 4\tabcolsep) * \real{0.62}}@{}}
\toprule
位置 & 类型 & 含义 \\
\midrule
\endhead
返回值 & \passthrough{\lstinline!int!} & SDK 状态码；Unitree 示例通常以
\passthrough{\lstinline!0!} 表示成功，非零值需按具体服务定义解释。 \\
\bottomrule
\end{longtable}

\textbf{对应 C++}

\begin{itemize}
\tightlist
\item
  类：\passthrough{\lstinline!unitree::robot::go2::SportClient!}
\item
  签名：\passthrough{\lstinline!Dance1()!}
\item
  绑定策略：\passthrough{\lstinline!DIRECT!}
\item
  声明位置：\passthrough{\lstinline!include/unitree/robot/go2/sport/sport\_client.hpp:60!}
\end{itemize}

\textbf{说明}

\begin{itemize}
\tightlist
\item
  示例是局部调用片段；\passthrough{\lstinline!obj!}
  和参数变量表示已经按上层业务校验并准备好的对象。
\item
  该条目可以被编辑器和类型检查器识别，但当前运行时可能在属性查找阶段直接失败。
\item
  该方法属于运动命令。方法名包含
  \passthrough{\lstinline!stop!}、\passthrough{\lstinline!damp!} 或
  \passthrough{\lstinline!zero!} 也不代表它是独立物理急停。
\end{itemize}

\textbf{用法}

\begin{lstlisting}[language=Python]
from typing import TYPE_CHECKING

if TYPE_CHECKING:
    def planned_dance1(obj: SportClient) -> int:
        return obj.dance1()
\end{lstlisting}

\begin{quote}
{[}!CAUTION{]} 上面的 \passthrough{\lstinline!TYPE\_CHECKING!}
示例只表达计划调用形态。该分支在普通 Python
运行时为假，不代表当前扩展能执行此接口。
\end{quote}

\hypertarget{unitree-sdk2-cpp-robot-go2-sportclient-dance2-1}{%
\paragraph{\texorpdfstring{\texttt{unitree\_sdk2\_cpp.robot.go2.SportClient.dance2}}{unitree\_sdk2\_cpp.robot.go2.SportClient.dance2}}\label{unitree-sdk2-cpp-robot-go2-sportclient-dance2-1}}

对应 C++ SDK 操作
\passthrough{\lstinline!Dance2()!}。上游头文件没有可直接生成的业务说明时，本参考只保证签名映射，精确语义需查目标型号协议。

\textbf{签名}

\begin{lstlisting}[language=Python]
def dance2(self) -> int
\end{lstlisting}

\textbf{可用性与安全性}

\textbf{\passthrough{\lstinline!SIGNATURE\_ONLY!} /
\passthrough{\lstinline!MOTION\_COMMAND!}}：仅用于补全和静态检查。当前二进制没有该
Python 方法；不得按可执行运动接口使用。

\textbf{参数}

无显式参数。实例方法中的 \passthrough{\lstinline!self!} 由 Python
自动传入。

\textbf{返回值}

\begin{longtable}[]{@{}
  >{\raggedright\arraybackslash}p{(\columnwidth - 4\tabcolsep) * \real{0.18}}
  >{\raggedright\arraybackslash}p{(\columnwidth - 4\tabcolsep) * \real{0.20}}
  >{\raggedright\arraybackslash}p{(\columnwidth - 4\tabcolsep) * \real{0.62}}@{}}
\toprule
位置 & 类型 & 含义 \\
\midrule
\endhead
返回值 & \passthrough{\lstinline!int!} & SDK 状态码；Unitree 示例通常以
\passthrough{\lstinline!0!} 表示成功，非零值需按具体服务定义解释。 \\
\bottomrule
\end{longtable}

\textbf{对应 C++}

\begin{itemize}
\tightlist
\item
  类：\passthrough{\lstinline!unitree::robot::go2::SportClient!}
\item
  签名：\passthrough{\lstinline!Dance2()!}
\item
  绑定策略：\passthrough{\lstinline!DIRECT!}
\item
  声明位置：\passthrough{\lstinline!include/unitree/robot/go2/sport/sport\_client.hpp:61!}
\end{itemize}

\textbf{说明}

\begin{itemize}
\tightlist
\item
  示例是局部调用片段；\passthrough{\lstinline!obj!}
  和参数变量表示已经按上层业务校验并准备好的对象。
\item
  该条目可以被编辑器和类型检查器识别，但当前运行时可能在属性查找阶段直接失败。
\item
  该方法属于运动命令。方法名包含
  \passthrough{\lstinline!stop!}、\passthrough{\lstinline!damp!} 或
  \passthrough{\lstinline!zero!} 也不代表它是独立物理急停。
\end{itemize}

\textbf{用法}

\begin{lstlisting}[language=Python]
from typing import TYPE_CHECKING

if TYPE_CHECKING:
    def planned_dance2(obj: SportClient) -> int:
        return obj.dance2()
\end{lstlisting}

\begin{quote}
{[}!CAUTION{]} 上面的 \passthrough{\lstinline!TYPE\_CHECKING!}
示例只表达计划调用形态。该分支在普通 Python
运行时为假，不代表当前扩展能执行此接口。
\end{quote}

\hypertarget{unitree-sdk2-cpp-robot-go2-sportclient-left-flip-1}{%
\paragraph{\texorpdfstring{\texttt{unitree\_sdk2\_cpp.robot.go2.SportClient.left\_flip}}{unitree\_sdk2\_cpp.robot.go2.SportClient.left\_flip}}\label{unitree-sdk2-cpp-robot-go2-sportclient-left-flip-1}}

对应 C++ SDK 操作
\passthrough{\lstinline!LeftFlip()!}。上游头文件没有可直接生成的业务说明时，本参考只保证签名映射，精确语义需查目标型号协议。

\textbf{签名}

\begin{lstlisting}[language=Python]
def left_flip(self) -> int
\end{lstlisting}

\textbf{可用性与安全性}

\textbf{\passthrough{\lstinline!SIGNATURE\_ONLY!} /
\passthrough{\lstinline!MOTION\_COMMAND!}}：仅用于补全和静态检查。当前二进制没有该
Python 方法；不得按可执行运动接口使用。

\textbf{参数}

无显式参数。实例方法中的 \passthrough{\lstinline!self!} 由 Python
自动传入。

\textbf{返回值}

\begin{longtable}[]{@{}
  >{\raggedright\arraybackslash}p{(\columnwidth - 4\tabcolsep) * \real{0.18}}
  >{\raggedright\arraybackslash}p{(\columnwidth - 4\tabcolsep) * \real{0.20}}
  >{\raggedright\arraybackslash}p{(\columnwidth - 4\tabcolsep) * \real{0.62}}@{}}
\toprule
位置 & 类型 & 含义 \\
\midrule
\endhead
返回值 & \passthrough{\lstinline!int!} & SDK 状态码；Unitree 示例通常以
\passthrough{\lstinline!0!} 表示成功，非零值需按具体服务定义解释。 \\
\bottomrule
\end{longtable}

\textbf{对应 C++}

\begin{itemize}
\tightlist
\item
  类：\passthrough{\lstinline!unitree::robot::go2::SportClient!}
\item
  签名：\passthrough{\lstinline!LeftFlip()!}
\item
  绑定策略：\passthrough{\lstinline!DIRECT!}
\item
  声明位置：\passthrough{\lstinline!include/unitree/robot/go2/sport/sport\_client.hpp:62!}
\end{itemize}

\textbf{说明}

\begin{itemize}
\tightlist
\item
  示例是局部调用片段；\passthrough{\lstinline!obj!}
  和参数变量表示已经按上层业务校验并准备好的对象。
\item
  该条目可以被编辑器和类型检查器识别，但当前运行时可能在属性查找阶段直接失败。
\item
  该方法属于运动命令。方法名包含
  \passthrough{\lstinline!stop!}、\passthrough{\lstinline!damp!} 或
  \passthrough{\lstinline!zero!} 也不代表它是独立物理急停。
\end{itemize}

\textbf{用法}

\begin{lstlisting}[language=Python]
from typing import TYPE_CHECKING

if TYPE_CHECKING:
    def planned_left_flip(obj: SportClient) -> int:
        return obj.left_flip()
\end{lstlisting}

\begin{quote}
{[}!CAUTION{]} 上面的 \passthrough{\lstinline!TYPE\_CHECKING!}
示例只表达计划调用形态。该分支在普通 Python
运行时为假，不代表当前扩展能执行此接口。
\end{quote}

\hypertarget{unitree-sdk2-cpp-robot-go2-sportclient-back-flip-1}{%
\paragraph{\texorpdfstring{\texttt{unitree\_sdk2\_cpp.robot.go2.SportClient.back\_flip}}{unitree\_sdk2\_cpp.robot.go2.SportClient.back\_flip}}\label{unitree-sdk2-cpp-robot-go2-sportclient-back-flip-1}}

对应 C++ SDK 操作
\passthrough{\lstinline!BackFlip()!}。上游头文件没有可直接生成的业务说明时，本参考只保证签名映射，精确语义需查目标型号协议。

\textbf{签名}

\begin{lstlisting}[language=Python]
def back_flip(self) -> int
\end{lstlisting}

\textbf{可用性与安全性}

\textbf{\passthrough{\lstinline!SIGNATURE\_ONLY!} /
\passthrough{\lstinline!MOTION\_COMMAND!}}：仅用于补全和静态检查。当前二进制没有该
Python 方法；不得按可执行运动接口使用。

\textbf{参数}

无显式参数。实例方法中的 \passthrough{\lstinline!self!} 由 Python
自动传入。

\textbf{返回值}

\begin{longtable}[]{@{}
  >{\raggedright\arraybackslash}p{(\columnwidth - 4\tabcolsep) * \real{0.18}}
  >{\raggedright\arraybackslash}p{(\columnwidth - 4\tabcolsep) * \real{0.20}}
  >{\raggedright\arraybackslash}p{(\columnwidth - 4\tabcolsep) * \real{0.62}}@{}}
\toprule
位置 & 类型 & 含义 \\
\midrule
\endhead
返回值 & \passthrough{\lstinline!int!} & SDK 状态码；Unitree 示例通常以
\passthrough{\lstinline!0!} 表示成功，非零值需按具体服务定义解释。 \\
\bottomrule
\end{longtable}

\textbf{对应 C++}

\begin{itemize}
\tightlist
\item
  类：\passthrough{\lstinline!unitree::robot::go2::SportClient!}
\item
  签名：\passthrough{\lstinline!BackFlip()!}
\item
  绑定策略：\passthrough{\lstinline!DIRECT!}
\item
  声明位置：\passthrough{\lstinline!include/unitree/robot/go2/sport/sport\_client.hpp:63!}
\end{itemize}

\textbf{说明}

\begin{itemize}
\tightlist
\item
  示例是局部调用片段；\passthrough{\lstinline!obj!}
  和参数变量表示已经按上层业务校验并准备好的对象。
\item
  该条目可以被编辑器和类型检查器识别，但当前运行时可能在属性查找阶段直接失败。
\item
  该方法属于运动命令。方法名包含
  \passthrough{\lstinline!stop!}、\passthrough{\lstinline!damp!} 或
  \passthrough{\lstinline!zero!} 也不代表它是独立物理急停。
\end{itemize}

\textbf{用法}

\begin{lstlisting}[language=Python]
from typing import TYPE_CHECKING

if TYPE_CHECKING:
    def planned_back_flip(obj: SportClient) -> int:
        return obj.back_flip()
\end{lstlisting}

\begin{quote}
{[}!CAUTION{]} 上面的 \passthrough{\lstinline!TYPE\_CHECKING!}
示例只表达计划调用形态。该分支在普通 Python
运行时为假，不代表当前扩展能执行此接口。
\end{quote}

\hypertarget{unitree-sdk2-cpp-robot-go2-sportclient-hand-stand-1}{%
\paragraph{\texorpdfstring{\texttt{unitree\_sdk2\_cpp.robot.go2.SportClient.hand\_stand}}{unitree\_sdk2\_cpp.robot.go2.SportClient.hand\_stand}}\label{unitree-sdk2-cpp-robot-go2-sportclient-hand-stand-1}}

对应 C++ SDK 操作
\passthrough{\lstinline!HandStand(bool)!}。上游头文件没有可直接生成的业务说明时，本参考只保证签名映射，精确语义需查目标型号协议。

\textbf{签名}

\begin{lstlisting}[language=Python]
def hand_stand(self, flag: bool) -> int
\end{lstlisting}

\textbf{可用性与安全性}

\textbf{\passthrough{\lstinline!SIGNATURE\_ONLY!} /
\passthrough{\lstinline!MOTION\_COMMAND!}}：仅用于补全和静态检查。当前二进制没有该
Python 方法；不得按可执行运动接口使用。

\textbf{参数}

\begin{longtable}[]{@{}
  >{\raggedright\arraybackslash}p{(\columnwidth - 6\tabcolsep) * \real{0.18}}
  >{\raggedright\arraybackslash}p{(\columnwidth - 6\tabcolsep) * \real{0.24}}
  >{\raggedright\arraybackslash}p{(\columnwidth - 6\tabcolsep) * \real{0.18}}
  >{\raggedright\arraybackslash}p{(\columnwidth - 6\tabcolsep) * \real{0.40}}@{}}
\toprule
名称 & Python 类型 & 默认值 & 含义 \\
\midrule
\endhead
\passthrough{\lstinline!flag!} & \passthrough{\lstinline!bool!} & 必填 &
布尔开关。\passthrough{\lstinline!True!} 和
\passthrough{\lstinline!False!} 的具体业务效果由当前方法定义。 对应 C++
参数 \passthrough{\lstinline!flag: bool!}。 只接受布尔语义。 \\
\bottomrule
\end{longtable}

\textbf{返回值}

\begin{longtable}[]{@{}
  >{\raggedright\arraybackslash}p{(\columnwidth - 4\tabcolsep) * \real{0.18}}
  >{\raggedright\arraybackslash}p{(\columnwidth - 4\tabcolsep) * \real{0.20}}
  >{\raggedright\arraybackslash}p{(\columnwidth - 4\tabcolsep) * \real{0.62}}@{}}
\toprule
位置 & 类型 & 含义 \\
\midrule
\endhead
返回值 & \passthrough{\lstinline!int!} & SDK 状态码；Unitree 示例通常以
\passthrough{\lstinline!0!} 表示成功，非零值需按具体服务定义解释。 \\
\bottomrule
\end{longtable}

\textbf{对应 C++}

\begin{itemize}
\tightlist
\item
  类：\passthrough{\lstinline!unitree::robot::go2::SportClient!}
\item
  签名：\passthrough{\lstinline!HandStand(bool)!}
\item
  绑定策略：\passthrough{\lstinline!DIRECT!}
\item
  声明位置：\passthrough{\lstinline!include/unitree/robot/go2/sport/sport\_client.hpp:64!}
\end{itemize}

\textbf{说明}

\begin{itemize}
\tightlist
\item
  示例是局部调用片段；\passthrough{\lstinline!obj!}
  和参数变量表示已经按上层业务校验并准备好的对象。
\item
  该条目可以被编辑器和类型检查器识别，但当前运行时可能在属性查找阶段直接失败。
\item
  该方法属于运动命令。方法名包含
  \passthrough{\lstinline!stop!}、\passthrough{\lstinline!damp!} 或
  \passthrough{\lstinline!zero!} 也不代表它是独立物理急停。
\end{itemize}

\textbf{用法}

\begin{lstlisting}[language=Python]
from typing import TYPE_CHECKING

if TYPE_CHECKING:
    def planned_hand_stand(obj: SportClient, flag: bool) -> int:
        return obj.hand_stand(flag=flag)
\end{lstlisting}

\begin{quote}
{[}!CAUTION{]} 上面的 \passthrough{\lstinline!TYPE\_CHECKING!}
示例只表达计划调用形态。该分支在普通 Python
运行时为假，不代表当前扩展能执行此接口。
\end{quote}

\hypertarget{unitree-sdk2-cpp-robot-go2-sportclient-free-walk-1}{%
\paragraph{\texorpdfstring{\texttt{unitree\_sdk2\_cpp.robot.go2.SportClient.free\_walk}}{unitree\_sdk2\_cpp.robot.go2.SportClient.free\_walk}}\label{unitree-sdk2-cpp-robot-go2-sportclient-free-walk-1}}

对应 C++ SDK 操作
\passthrough{\lstinline!FreeWalk()!}。上游头文件没有可直接生成的业务说明时，本参考只保证签名映射，精确语义需查目标型号协议。

\textbf{签名}

\begin{lstlisting}[language=Python]
def free_walk(self) -> int
\end{lstlisting}

\textbf{可用性与安全性}

\textbf{\passthrough{\lstinline!SIGNATURE\_ONLY!} /
\passthrough{\lstinline!MOTION\_COMMAND!}}：仅用于补全和静态检查。当前二进制没有该
Python 方法；不得按可执行运动接口使用。

\textbf{参数}

无显式参数。实例方法中的 \passthrough{\lstinline!self!} 由 Python
自动传入。

\textbf{返回值}

\begin{longtable}[]{@{}
  >{\raggedright\arraybackslash}p{(\columnwidth - 4\tabcolsep) * \real{0.18}}
  >{\raggedright\arraybackslash}p{(\columnwidth - 4\tabcolsep) * \real{0.20}}
  >{\raggedright\arraybackslash}p{(\columnwidth - 4\tabcolsep) * \real{0.62}}@{}}
\toprule
位置 & 类型 & 含义 \\
\midrule
\endhead
返回值 & \passthrough{\lstinline!int!} & SDK 状态码；Unitree 示例通常以
\passthrough{\lstinline!0!} 表示成功，非零值需按具体服务定义解释。 \\
\bottomrule
\end{longtable}

\textbf{对应 C++}

\begin{itemize}
\tightlist
\item
  类：\passthrough{\lstinline!unitree::robot::go2::SportClient!}
\item
  签名：\passthrough{\lstinline!FreeWalk()!}
\item
  绑定策略：\passthrough{\lstinline!DIRECT!}
\item
  声明位置：\passthrough{\lstinline!include/unitree/robot/go2/sport/sport\_client.hpp:65!}
\end{itemize}

\textbf{说明}

\begin{itemize}
\tightlist
\item
  示例是局部调用片段；\passthrough{\lstinline!obj!}
  和参数变量表示已经按上层业务校验并准备好的对象。
\item
  该条目可以被编辑器和类型检查器识别，但当前运行时可能在属性查找阶段直接失败。
\item
  该方法属于运动命令。方法名包含
  \passthrough{\lstinline!stop!}、\passthrough{\lstinline!damp!} 或
  \passthrough{\lstinline!zero!} 也不代表它是独立物理急停。
\end{itemize}

\textbf{用法}

\begin{lstlisting}[language=Python]
from typing import TYPE_CHECKING

if TYPE_CHECKING:
    def planned_free_walk(obj: SportClient) -> int:
        return obj.free_walk()
\end{lstlisting}

\begin{quote}
{[}!CAUTION{]} 上面的 \passthrough{\lstinline!TYPE\_CHECKING!}
示例只表达计划调用形态。该分支在普通 Python
运行时为假，不代表当前扩展能执行此接口。
\end{quote}

\hypertarget{unitree-sdk2-cpp-robot-go2-sportclient-free-bound-1}{%
\paragraph{\texorpdfstring{\texttt{unitree\_sdk2\_cpp.robot.go2.SportClient.free\_bound}}{unitree\_sdk2\_cpp.robot.go2.SportClient.free\_bound}}\label{unitree-sdk2-cpp-robot-go2-sportclient-free-bound-1}}

对应 C++ SDK 操作
\passthrough{\lstinline!FreeBound(bool)!}。上游头文件没有可直接生成的业务说明时，本参考只保证签名映射，精确语义需查目标型号协议。

\textbf{签名}

\begin{lstlisting}[language=Python]
def free_bound(self, flag: bool) -> int
\end{lstlisting}

\textbf{可用性与安全性}

\textbf{\passthrough{\lstinline!SIGNATURE\_ONLY!} /
\passthrough{\lstinline!MOTION\_COMMAND!}}：仅用于补全和静态检查。当前二进制没有该
Python 方法；不得按可执行运动接口使用。

\textbf{参数}

\begin{longtable}[]{@{}
  >{\raggedright\arraybackslash}p{(\columnwidth - 6\tabcolsep) * \real{0.18}}
  >{\raggedright\arraybackslash}p{(\columnwidth - 6\tabcolsep) * \real{0.24}}
  >{\raggedright\arraybackslash}p{(\columnwidth - 6\tabcolsep) * \real{0.18}}
  >{\raggedright\arraybackslash}p{(\columnwidth - 6\tabcolsep) * \real{0.40}}@{}}
\toprule
名称 & Python 类型 & 默认值 & 含义 \\
\midrule
\endhead
\passthrough{\lstinline!flag!} & \passthrough{\lstinline!bool!} & 必填 &
布尔开关。\passthrough{\lstinline!True!} 和
\passthrough{\lstinline!False!} 的具体业务效果由当前方法定义。 对应 C++
参数 \passthrough{\lstinline!flag: bool!}。 只接受布尔语义。 \\
\bottomrule
\end{longtable}

\textbf{返回值}

\begin{longtable}[]{@{}
  >{\raggedright\arraybackslash}p{(\columnwidth - 4\tabcolsep) * \real{0.18}}
  >{\raggedright\arraybackslash}p{(\columnwidth - 4\tabcolsep) * \real{0.20}}
  >{\raggedright\arraybackslash}p{(\columnwidth - 4\tabcolsep) * \real{0.62}}@{}}
\toprule
位置 & 类型 & 含义 \\
\midrule
\endhead
返回值 & \passthrough{\lstinline!int!} & SDK 状态码；Unitree 示例通常以
\passthrough{\lstinline!0!} 表示成功，非零值需按具体服务定义解释。 \\
\bottomrule
\end{longtable}

\textbf{对应 C++}

\begin{itemize}
\tightlist
\item
  类：\passthrough{\lstinline!unitree::robot::go2::SportClient!}
\item
  签名：\passthrough{\lstinline!FreeBound(bool)!}
\item
  绑定策略：\passthrough{\lstinline!DIRECT!}
\item
  声明位置：\passthrough{\lstinline!include/unitree/robot/go2/sport/sport\_client.hpp:66!}
\end{itemize}

\textbf{说明}

\begin{itemize}
\tightlist
\item
  示例是局部调用片段；\passthrough{\lstinline!obj!}
  和参数变量表示已经按上层业务校验并准备好的对象。
\item
  该条目可以被编辑器和类型检查器识别，但当前运行时可能在属性查找阶段直接失败。
\item
  该方法属于运动命令。方法名包含
  \passthrough{\lstinline!stop!}、\passthrough{\lstinline!damp!} 或
  \passthrough{\lstinline!zero!} 也不代表它是独立物理急停。
\end{itemize}

\textbf{用法}

\begin{lstlisting}[language=Python]
from typing import TYPE_CHECKING

if TYPE_CHECKING:
    def planned_free_bound(obj: SportClient, flag: bool) -> int:
        return obj.free_bound(flag=flag)
\end{lstlisting}

\begin{quote}
{[}!CAUTION{]} 上面的 \passthrough{\lstinline!TYPE\_CHECKING!}
示例只表达计划调用形态。该分支在普通 Python
运行时为假，不代表当前扩展能执行此接口。
\end{quote}

\hypertarget{unitree-sdk2-cpp-robot-go2-sportclient-free-jump-1}{%
\paragraph{\texorpdfstring{\texttt{unitree\_sdk2\_cpp.robot.go2.SportClient.free\_jump}}{unitree\_sdk2\_cpp.robot.go2.SportClient.free\_jump}}\label{unitree-sdk2-cpp-robot-go2-sportclient-free-jump-1}}

对应 C++ SDK 操作
\passthrough{\lstinline!FreeJump(bool)!}。上游头文件没有可直接生成的业务说明时，本参考只保证签名映射，精确语义需查目标型号协议。

\textbf{签名}

\begin{lstlisting}[language=Python]
def free_jump(self, flag: bool) -> int
\end{lstlisting}

\textbf{可用性与安全性}

\textbf{\passthrough{\lstinline!SIGNATURE\_ONLY!} /
\passthrough{\lstinline!MOTION\_COMMAND!}}：仅用于补全和静态检查。当前二进制没有该
Python 方法；不得按可执行运动接口使用。

\textbf{参数}

\begin{longtable}[]{@{}
  >{\raggedright\arraybackslash}p{(\columnwidth - 6\tabcolsep) * \real{0.18}}
  >{\raggedright\arraybackslash}p{(\columnwidth - 6\tabcolsep) * \real{0.24}}
  >{\raggedright\arraybackslash}p{(\columnwidth - 6\tabcolsep) * \real{0.18}}
  >{\raggedright\arraybackslash}p{(\columnwidth - 6\tabcolsep) * \real{0.40}}@{}}
\toprule
名称 & Python 类型 & 默认值 & 含义 \\
\midrule
\endhead
\passthrough{\lstinline!flag!} & \passthrough{\lstinline!bool!} & 必填 &
布尔开关。\passthrough{\lstinline!True!} 和
\passthrough{\lstinline!False!} 的具体业务效果由当前方法定义。 对应 C++
参数 \passthrough{\lstinline!flag: bool!}。 只接受布尔语义。 \\
\bottomrule
\end{longtable}

\textbf{返回值}

\begin{longtable}[]{@{}
  >{\raggedright\arraybackslash}p{(\columnwidth - 4\tabcolsep) * \real{0.18}}
  >{\raggedright\arraybackslash}p{(\columnwidth - 4\tabcolsep) * \real{0.20}}
  >{\raggedright\arraybackslash}p{(\columnwidth - 4\tabcolsep) * \real{0.62}}@{}}
\toprule
位置 & 类型 & 含义 \\
\midrule
\endhead
返回值 & \passthrough{\lstinline!int!} & SDK 状态码；Unitree 示例通常以
\passthrough{\lstinline!0!} 表示成功，非零值需按具体服务定义解释。 \\
\bottomrule
\end{longtable}

\textbf{对应 C++}

\begin{itemize}
\tightlist
\item
  类：\passthrough{\lstinline!unitree::robot::go2::SportClient!}
\item
  签名：\passthrough{\lstinline!FreeJump(bool)!}
\item
  绑定策略：\passthrough{\lstinline!DIRECT!}
\item
  声明位置：\passthrough{\lstinline!include/unitree/robot/go2/sport/sport\_client.hpp:67!}
\end{itemize}

\textbf{说明}

\begin{itemize}
\tightlist
\item
  示例是局部调用片段；\passthrough{\lstinline!obj!}
  和参数变量表示已经按上层业务校验并准备好的对象。
\item
  该条目可以被编辑器和类型检查器识别，但当前运行时可能在属性查找阶段直接失败。
\item
  该方法属于运动命令。方法名包含
  \passthrough{\lstinline!stop!}、\passthrough{\lstinline!damp!} 或
  \passthrough{\lstinline!zero!} 也不代表它是独立物理急停。
\end{itemize}

\textbf{用法}

\begin{lstlisting}[language=Python]
from typing import TYPE_CHECKING

if TYPE_CHECKING:
    def planned_free_jump(obj: SportClient, flag: bool) -> int:
        return obj.free_jump(flag=flag)
\end{lstlisting}

\begin{quote}
{[}!CAUTION{]} 上面的 \passthrough{\lstinline!TYPE\_CHECKING!}
示例只表达计划调用形态。该分支在普通 Python
运行时为假，不代表当前扩展能执行此接口。
\end{quote}

\hypertarget{unitree-sdk2-cpp-robot-go2-sportclient-free-avoid-1}{%
\paragraph{\texorpdfstring{\texttt{unitree\_sdk2\_cpp.robot.go2.SportClient.free\_avoid}}{unitree\_sdk2\_cpp.robot.go2.SportClient.free\_avoid}}\label{unitree-sdk2-cpp-robot-go2-sportclient-free-avoid-1}}

对应 C++ SDK 操作
\passthrough{\lstinline!FreeAvoid(bool)!}。上游头文件没有可直接生成的业务说明时，本参考只保证签名映射，精确语义需查目标型号协议。

\textbf{签名}

\begin{lstlisting}[language=Python]
def free_avoid(self, flag: bool) -> int
\end{lstlisting}

\textbf{可用性与安全性}

\textbf{\passthrough{\lstinline!SIGNATURE\_ONLY!} /
\passthrough{\lstinline!MOTION\_COMMAND!}}：仅用于补全和静态检查。当前二进制没有该
Python 方法；不得按可执行运动接口使用。

\textbf{参数}

\begin{longtable}[]{@{}
  >{\raggedright\arraybackslash}p{(\columnwidth - 6\tabcolsep) * \real{0.18}}
  >{\raggedright\arraybackslash}p{(\columnwidth - 6\tabcolsep) * \real{0.24}}
  >{\raggedright\arraybackslash}p{(\columnwidth - 6\tabcolsep) * \real{0.18}}
  >{\raggedright\arraybackslash}p{(\columnwidth - 6\tabcolsep) * \real{0.40}}@{}}
\toprule
名称 & Python 类型 & 默认值 & 含义 \\
\midrule
\endhead
\passthrough{\lstinline!flag!} & \passthrough{\lstinline!bool!} & 必填 &
布尔开关。\passthrough{\lstinline!True!} 和
\passthrough{\lstinline!False!} 的具体业务效果由当前方法定义。 对应 C++
参数 \passthrough{\lstinline!flag: bool!}。 只接受布尔语义。 \\
\bottomrule
\end{longtable}

\textbf{返回值}

\begin{longtable}[]{@{}
  >{\raggedright\arraybackslash}p{(\columnwidth - 4\tabcolsep) * \real{0.18}}
  >{\raggedright\arraybackslash}p{(\columnwidth - 4\tabcolsep) * \real{0.20}}
  >{\raggedright\arraybackslash}p{(\columnwidth - 4\tabcolsep) * \real{0.62}}@{}}
\toprule
位置 & 类型 & 含义 \\
\midrule
\endhead
返回值 & \passthrough{\lstinline!int!} & SDK 状态码；Unitree 示例通常以
\passthrough{\lstinline!0!} 表示成功，非零值需按具体服务定义解释。 \\
\bottomrule
\end{longtable}

\textbf{对应 C++}

\begin{itemize}
\tightlist
\item
  类：\passthrough{\lstinline!unitree::robot::go2::SportClient!}
\item
  签名：\passthrough{\lstinline!FreeAvoid(bool)!}
\item
  绑定策略：\passthrough{\lstinline!DIRECT!}
\item
  声明位置：\passthrough{\lstinline!include/unitree/robot/go2/sport/sport\_client.hpp:68!}
\end{itemize}

\textbf{说明}

\begin{itemize}
\tightlist
\item
  示例是局部调用片段；\passthrough{\lstinline!obj!}
  和参数变量表示已经按上层业务校验并准备好的对象。
\item
  该条目可以被编辑器和类型检查器识别，但当前运行时可能在属性查找阶段直接失败。
\item
  该方法属于运动命令。方法名包含
  \passthrough{\lstinline!stop!}、\passthrough{\lstinline!damp!} 或
  \passthrough{\lstinline!zero!} 也不代表它是独立物理急停。
\end{itemize}

\textbf{用法}

\begin{lstlisting}[language=Python]
from typing import TYPE_CHECKING

if TYPE_CHECKING:
    def planned_free_avoid(obj: SportClient, flag: bool) -> int:
        return obj.free_avoid(flag=flag)
\end{lstlisting}

\begin{quote}
{[}!CAUTION{]} 上面的 \passthrough{\lstinline!TYPE\_CHECKING!}
示例只表达计划调用形态。该分支在普通 Python
运行时为假，不代表当前扩展能执行此接口。
\end{quote}

\hypertarget{unitree-sdk2-cpp-robot-go2-sportclient-classic-walk-1}{%
\paragraph{\texorpdfstring{\texttt{unitree\_sdk2\_cpp.robot.go2.SportClient.classic\_walk}}{unitree\_sdk2\_cpp.robot.go2.SportClient.classic\_walk}}\label{unitree-sdk2-cpp-robot-go2-sportclient-classic-walk-1}}

对应 C++ SDK 操作
\passthrough{\lstinline!ClassicWalk(bool)!}。上游头文件没有可直接生成的业务说明时，本参考只保证签名映射，精确语义需查目标型号协议。

\textbf{签名}

\begin{lstlisting}[language=Python]
def classic_walk(self, flag: bool) -> int
\end{lstlisting}

\textbf{可用性与安全性}

\textbf{\passthrough{\lstinline!SIGNATURE\_ONLY!} /
\passthrough{\lstinline!MOTION\_COMMAND!}}：仅用于补全和静态检查。当前二进制没有该
Python 方法；不得按可执行运动接口使用。

\textbf{参数}

\begin{longtable}[]{@{}
  >{\raggedright\arraybackslash}p{(\columnwidth - 6\tabcolsep) * \real{0.18}}
  >{\raggedright\arraybackslash}p{(\columnwidth - 6\tabcolsep) * \real{0.24}}
  >{\raggedright\arraybackslash}p{(\columnwidth - 6\tabcolsep) * \real{0.18}}
  >{\raggedright\arraybackslash}p{(\columnwidth - 6\tabcolsep) * \real{0.40}}@{}}
\toprule
名称 & Python 类型 & 默认值 & 含义 \\
\midrule
\endhead
\passthrough{\lstinline!flag!} & \passthrough{\lstinline!bool!} & 必填 &
布尔开关。\passthrough{\lstinline!True!} 和
\passthrough{\lstinline!False!} 的具体业务效果由当前方法定义。 对应 C++
参数 \passthrough{\lstinline!flag: bool!}。 只接受布尔语义。 \\
\bottomrule
\end{longtable}

\textbf{返回值}

\begin{longtable}[]{@{}
  >{\raggedright\arraybackslash}p{(\columnwidth - 4\tabcolsep) * \real{0.18}}
  >{\raggedright\arraybackslash}p{(\columnwidth - 4\tabcolsep) * \real{0.20}}
  >{\raggedright\arraybackslash}p{(\columnwidth - 4\tabcolsep) * \real{0.62}}@{}}
\toprule
位置 & 类型 & 含义 \\
\midrule
\endhead
返回值 & \passthrough{\lstinline!int!} & SDK 状态码；Unitree 示例通常以
\passthrough{\lstinline!0!} 表示成功，非零值需按具体服务定义解释。 \\
\bottomrule
\end{longtable}

\textbf{对应 C++}

\begin{itemize}
\tightlist
\item
  类：\passthrough{\lstinline!unitree::robot::go2::SportClient!}
\item
  签名：\passthrough{\lstinline!ClassicWalk(bool)!}
\item
  绑定策略：\passthrough{\lstinline!DIRECT!}
\item
  声明位置：\passthrough{\lstinline!include/unitree/robot/go2/sport/sport\_client.hpp:69!}
\end{itemize}

\textbf{说明}

\begin{itemize}
\tightlist
\item
  示例是局部调用片段；\passthrough{\lstinline!obj!}
  和参数变量表示已经按上层业务校验并准备好的对象。
\item
  该条目可以被编辑器和类型检查器识别，但当前运行时可能在属性查找阶段直接失败。
\item
  该方法属于运动命令。方法名包含
  \passthrough{\lstinline!stop!}、\passthrough{\lstinline!damp!} 或
  \passthrough{\lstinline!zero!} 也不代表它是独立物理急停。
\end{itemize}

\textbf{用法}

\begin{lstlisting}[language=Python]
from typing import TYPE_CHECKING

if TYPE_CHECKING:
    def planned_classic_walk(obj: SportClient, flag: bool) -> int:
        return obj.classic_walk(flag=flag)
\end{lstlisting}

\begin{quote}
{[}!CAUTION{]} 上面的 \passthrough{\lstinline!TYPE\_CHECKING!}
示例只表达计划调用形态。该分支在普通 Python
运行时为假，不代表当前扩展能执行此接口。
\end{quote}

\hypertarget{unitree-sdk2-cpp-robot-go2-sportclient-walk-upright-1}{%
\paragraph{\texorpdfstring{\texttt{unitree\_sdk2\_cpp.robot.go2.SportClient.walk\_upright}}{unitree\_sdk2\_cpp.robot.go2.SportClient.walk\_upright}}\label{unitree-sdk2-cpp-robot-go2-sportclient-walk-upright-1}}

对应 C++ SDK 操作
\passthrough{\lstinline!WalkUpright(bool)!}。上游头文件没有可直接生成的业务说明时，本参考只保证签名映射，精确语义需查目标型号协议。

\textbf{签名}

\begin{lstlisting}[language=Python]
def walk_upright(self, flag: bool) -> int
\end{lstlisting}

\textbf{可用性与安全性}

\textbf{\passthrough{\lstinline!SIGNATURE\_ONLY!} /
\passthrough{\lstinline!MOTION\_COMMAND!}}：仅用于补全和静态检查。当前二进制没有该
Python 方法；不得按可执行运动接口使用。

\textbf{参数}

\begin{longtable}[]{@{}
  >{\raggedright\arraybackslash}p{(\columnwidth - 6\tabcolsep) * \real{0.18}}
  >{\raggedright\arraybackslash}p{(\columnwidth - 6\tabcolsep) * \real{0.24}}
  >{\raggedright\arraybackslash}p{(\columnwidth - 6\tabcolsep) * \real{0.18}}
  >{\raggedright\arraybackslash}p{(\columnwidth - 6\tabcolsep) * \real{0.40}}@{}}
\toprule
名称 & Python 类型 & 默认值 & 含义 \\
\midrule
\endhead
\passthrough{\lstinline!flag!} & \passthrough{\lstinline!bool!} & 必填 &
布尔开关。\passthrough{\lstinline!True!} 和
\passthrough{\lstinline!False!} 的具体业务效果由当前方法定义。 对应 C++
参数 \passthrough{\lstinline!flag: bool!}。 只接受布尔语义。 \\
\bottomrule
\end{longtable}

\textbf{返回值}

\begin{longtable}[]{@{}
  >{\raggedright\arraybackslash}p{(\columnwidth - 4\tabcolsep) * \real{0.18}}
  >{\raggedright\arraybackslash}p{(\columnwidth - 4\tabcolsep) * \real{0.20}}
  >{\raggedright\arraybackslash}p{(\columnwidth - 4\tabcolsep) * \real{0.62}}@{}}
\toprule
位置 & 类型 & 含义 \\
\midrule
\endhead
返回值 & \passthrough{\lstinline!int!} & SDK 状态码；Unitree 示例通常以
\passthrough{\lstinline!0!} 表示成功，非零值需按具体服务定义解释。 \\
\bottomrule
\end{longtable}

\textbf{对应 C++}

\begin{itemize}
\tightlist
\item
  类：\passthrough{\lstinline!unitree::robot::go2::SportClient!}
\item
  签名：\passthrough{\lstinline!WalkUpright(bool)!}
\item
  绑定策略：\passthrough{\lstinline!DIRECT!}
\item
  声明位置：\passthrough{\lstinline!include/unitree/robot/go2/sport/sport\_client.hpp:70!}
\end{itemize}

\textbf{说明}

\begin{itemize}
\tightlist
\item
  示例是局部调用片段；\passthrough{\lstinline!obj!}
  和参数变量表示已经按上层业务校验并准备好的对象。
\item
  该条目可以被编辑器和类型检查器识别，但当前运行时可能在属性查找阶段直接失败。
\item
  该方法属于运动命令。方法名包含
  \passthrough{\lstinline!stop!}、\passthrough{\lstinline!damp!} 或
  \passthrough{\lstinline!zero!} 也不代表它是独立物理急停。
\end{itemize}

\textbf{用法}

\begin{lstlisting}[language=Python]
from typing import TYPE_CHECKING

if TYPE_CHECKING:
    def planned_walk_upright(obj: SportClient, flag: bool) -> int:
        return obj.walk_upright(flag=flag)
\end{lstlisting}

\begin{quote}
{[}!CAUTION{]} 上面的 \passthrough{\lstinline!TYPE\_CHECKING!}
示例只表达计划调用形态。该分支在普通 Python
运行时为假，不代表当前扩展能执行此接口。
\end{quote}

\hypertarget{unitree-sdk2-cpp-robot-go2-sportclient-cross-step-1}{%
\paragraph{\texorpdfstring{\texttt{unitree\_sdk2\_cpp.robot.go2.SportClient.cross\_step}}{unitree\_sdk2\_cpp.robot.go2.SportClient.cross\_step}}\label{unitree-sdk2-cpp-robot-go2-sportclient-cross-step-1}}

对应 C++ SDK 操作
\passthrough{\lstinline!CrossStep(bool)!}。上游头文件没有可直接生成的业务说明时，本参考只保证签名映射，精确语义需查目标型号协议。

\textbf{签名}

\begin{lstlisting}[language=Python]
def cross_step(self, flag: bool) -> int
\end{lstlisting}

\textbf{可用性与安全性}

\textbf{\passthrough{\lstinline!SIGNATURE\_ONLY!} /
\passthrough{\lstinline!MOTION\_COMMAND!}}：仅用于补全和静态检查。当前二进制没有该
Python 方法；不得按可执行运动接口使用。

\textbf{参数}

\begin{longtable}[]{@{}
  >{\raggedright\arraybackslash}p{(\columnwidth - 6\tabcolsep) * \real{0.18}}
  >{\raggedright\arraybackslash}p{(\columnwidth - 6\tabcolsep) * \real{0.24}}
  >{\raggedright\arraybackslash}p{(\columnwidth - 6\tabcolsep) * \real{0.18}}
  >{\raggedright\arraybackslash}p{(\columnwidth - 6\tabcolsep) * \real{0.40}}@{}}
\toprule
名称 & Python 类型 & 默认值 & 含义 \\
\midrule
\endhead
\passthrough{\lstinline!flag!} & \passthrough{\lstinline!bool!} & 必填 &
布尔开关。\passthrough{\lstinline!True!} 和
\passthrough{\lstinline!False!} 的具体业务效果由当前方法定义。 对应 C++
参数 \passthrough{\lstinline!flag: bool!}。 只接受布尔语义。 \\
\bottomrule
\end{longtable}

\textbf{返回值}

\begin{longtable}[]{@{}
  >{\raggedright\arraybackslash}p{(\columnwidth - 4\tabcolsep) * \real{0.18}}
  >{\raggedright\arraybackslash}p{(\columnwidth - 4\tabcolsep) * \real{0.20}}
  >{\raggedright\arraybackslash}p{(\columnwidth - 4\tabcolsep) * \real{0.62}}@{}}
\toprule
位置 & 类型 & 含义 \\
\midrule
\endhead
返回值 & \passthrough{\lstinline!int!} & SDK 状态码；Unitree 示例通常以
\passthrough{\lstinline!0!} 表示成功，非零值需按具体服务定义解释。 \\
\bottomrule
\end{longtable}

\textbf{对应 C++}

\begin{itemize}
\tightlist
\item
  类：\passthrough{\lstinline!unitree::robot::go2::SportClient!}
\item
  签名：\passthrough{\lstinline!CrossStep(bool)!}
\item
  绑定策略：\passthrough{\lstinline!DIRECT!}
\item
  声明位置：\passthrough{\lstinline!include/unitree/robot/go2/sport/sport\_client.hpp:71!}
\end{itemize}

\textbf{说明}

\begin{itemize}
\tightlist
\item
  示例是局部调用片段；\passthrough{\lstinline!obj!}
  和参数变量表示已经按上层业务校验并准备好的对象。
\item
  该条目可以被编辑器和类型检查器识别，但当前运行时可能在属性查找阶段直接失败。
\item
  该方法属于运动命令。方法名包含
  \passthrough{\lstinline!stop!}、\passthrough{\lstinline!damp!} 或
  \passthrough{\lstinline!zero!} 也不代表它是独立物理急停。
\end{itemize}

\textbf{用法}

\begin{lstlisting}[language=Python]
from typing import TYPE_CHECKING

if TYPE_CHECKING:
    def planned_cross_step(obj: SportClient, flag: bool) -> int:
        return obj.cross_step(flag=flag)
\end{lstlisting}

\begin{quote}
{[}!CAUTION{]} 上面的 \passthrough{\lstinline!TYPE\_CHECKING!}
示例只表达计划调用形态。该分支在普通 Python
运行时为假，不代表当前扩展能执行此接口。
\end{quote}

\hypertarget{unitree-sdk2-cpp-robot-go2-sportclient-auto-recover-set-1}{%
\paragraph{\texorpdfstring{\texttt{unitree\_sdk2\_cpp.robot.go2.SportClient.auto\_recover\_set}}{unitree\_sdk2\_cpp.robot.go2.SportClient.auto\_recover\_set}}\label{unitree-sdk2-cpp-robot-go2-sportclient-auto-recover-set-1}}

对应 C++ SDK 操作
\passthrough{\lstinline!AutoRecoverSet(bool)!}。上游头文件没有可直接生成的业务说明时，本参考只保证签名映射，精确语义需查目标型号协议。

\textbf{签名}

\begin{lstlisting}[language=Python]
def auto_recover_set(self, flag: bool) -> int
\end{lstlisting}

\textbf{可用性与安全性}

\textbf{\passthrough{\lstinline!SIGNATURE\_ONLY!} /
\passthrough{\lstinline!MOTION\_COMMAND!}}：仅用于补全和静态检查。当前二进制没有该
Python 方法；不得按可执行运动接口使用。

\textbf{参数}

\begin{longtable}[]{@{}
  >{\raggedright\arraybackslash}p{(\columnwidth - 6\tabcolsep) * \real{0.18}}
  >{\raggedright\arraybackslash}p{(\columnwidth - 6\tabcolsep) * \real{0.24}}
  >{\raggedright\arraybackslash}p{(\columnwidth - 6\tabcolsep) * \real{0.18}}
  >{\raggedright\arraybackslash}p{(\columnwidth - 6\tabcolsep) * \real{0.40}}@{}}
\toprule
名称 & Python 类型 & 默认值 & 含义 \\
\midrule
\endhead
\passthrough{\lstinline!flag!} & \passthrough{\lstinline!bool!} & 必填 &
布尔开关。\passthrough{\lstinline!True!} 和
\passthrough{\lstinline!False!} 的具体业务效果由当前方法定义。 对应 C++
参数 \passthrough{\lstinline!flag: bool!}。 只接受布尔语义。 \\
\bottomrule
\end{longtable}

\textbf{返回值}

\begin{longtable}[]{@{}
  >{\raggedright\arraybackslash}p{(\columnwidth - 4\tabcolsep) * \real{0.18}}
  >{\raggedright\arraybackslash}p{(\columnwidth - 4\tabcolsep) * \real{0.20}}
  >{\raggedright\arraybackslash}p{(\columnwidth - 4\tabcolsep) * \real{0.62}}@{}}
\toprule
位置 & 类型 & 含义 \\
\midrule
\endhead
返回值 & \passthrough{\lstinline!int!} & SDK 状态码；Unitree 示例通常以
\passthrough{\lstinline!0!} 表示成功，非零值需按具体服务定义解释。 \\
\bottomrule
\end{longtable}

\textbf{对应 C++}

\begin{itemize}
\tightlist
\item
  类：\passthrough{\lstinline!unitree::robot::go2::SportClient!}
\item
  签名：\passthrough{\lstinline!AutoRecoverSet(bool)!}
\item
  绑定策略：\passthrough{\lstinline!DIRECT!}
\item
  声明位置：\passthrough{\lstinline!include/unitree/robot/go2/sport/sport\_client.hpp:72!}
\end{itemize}

\textbf{说明}

\begin{itemize}
\tightlist
\item
  示例是局部调用片段；\passthrough{\lstinline!obj!}
  和参数变量表示已经按上层业务校验并准备好的对象。
\item
  该条目可以被编辑器和类型检查器识别，但当前运行时可能在属性查找阶段直接失败。
\item
  该方法属于运动命令。方法名包含
  \passthrough{\lstinline!stop!}、\passthrough{\lstinline!damp!} 或
  \passthrough{\lstinline!zero!} 也不代表它是独立物理急停。
\end{itemize}

\textbf{用法}

\begin{lstlisting}[language=Python]
from typing import TYPE_CHECKING

if TYPE_CHECKING:
    def planned_auto_recover_set(obj: SportClient, flag: bool) -> int:
        return obj.auto_recover_set(flag=flag)
\end{lstlisting}

\begin{quote}
{[}!CAUTION{]} 上面的 \passthrough{\lstinline!TYPE\_CHECKING!}
示例只表达计划调用形态。该分支在普通 Python
运行时为假，不代表当前扩展能执行此接口。
\end{quote}

\hypertarget{unitree-sdk2-cpp-robot-go2-sportclient-auto-recover-get-1}{%
\paragraph{\texorpdfstring{\texttt{unitree\_sdk2\_cpp.robot.go2.SportClient.auto\_recover\_get}}{unitree\_sdk2\_cpp.robot.go2.SportClient.auto\_recover\_get}}\label{unitree-sdk2-cpp-robot-go2-sportclient-auto-recover-get-1}}

对应 C++ SDK 操作
\passthrough{\lstinline!AutoRecoverGet(bool \&)!}。上游头文件没有可直接生成的业务说明时，本参考只保证签名映射，精确语义需查目标型号协议。

\textbf{签名}

\begin{lstlisting}[language=Python]
def auto_recover_get(self) -> tuple[int, bool]
\end{lstlisting}

\textbf{可用性与安全性}

\textbf{\passthrough{\lstinline!SIGNATURE\_ONLY!} /
\passthrough{\lstinline!MOTION\_COMMAND!}}：仅用于补全和静态检查。当前二进制没有该
Python 方法；不得按可执行运动接口使用。

\textbf{参数}

无显式参数。实例方法中的 \passthrough{\lstinline!self!} 由 Python
自动传入。

\textbf{返回值}

\begin{longtable}[]{@{}
  >{\raggedright\arraybackslash}p{(\columnwidth - 4\tabcolsep) * \real{0.18}}
  >{\raggedright\arraybackslash}p{(\columnwidth - 4\tabcolsep) * \real{0.20}}
  >{\raggedright\arraybackslash}p{(\columnwidth - 4\tabcolsep) * \real{0.62}}@{}}
\toprule
位置 & 类型 & 含义 \\
\midrule
\endhead
{[}0{]} \passthrough{\lstinline!status!} & \passthrough{\lstinline!int!}
& SDK 状态码；Unitree 示例通常以 \passthrough{\lstinline!0!}
表示成功，非零值需按具体服务错误码解释。 \\
{[}1{]} \passthrough{\lstinline!flag!} & \passthrough{\lstinline!bool!}
& 布尔开关。\passthrough{\lstinline!True!} 和
\passthrough{\lstinline!False!} 的具体业务效果由当前方法定义。 对应 C++
参数 \passthrough{\lstinline!flag: bool \&!}。 只接受布尔语义。 \\
\bottomrule
\end{longtable}

\textbf{对应 C++}

\begin{itemize}
\tightlist
\item
  类：\passthrough{\lstinline!unitree::robot::go2::SportClient!}
\item
  签名：\passthrough{\lstinline!AutoRecoverGet(bool \&)!}
\item
  绑定策略：\passthrough{\lstinline!OUTPUT\_WRAPPER!}
\item
  声明位置：\passthrough{\lstinline!include/unitree/robot/go2/sport/sport\_client.hpp:73!}
\end{itemize}

\textbf{说明}

\begin{itemize}
\tightlist
\item
  示例是局部调用片段；\passthrough{\lstinline!obj!}
  和参数变量表示已经按上层业务校验并准备好的对象。
\item
  C++ 的可修改输出引用不会作为 Python 入参出现，而是按 C++
  参数顺序追加到返回元组中。
\item
  该条目可以被编辑器和类型检查器识别，但当前运行时可能在属性查找阶段直接失败。
\item
  该方法属于运动命令。方法名包含
  \passthrough{\lstinline!stop!}、\passthrough{\lstinline!damp!} 或
  \passthrough{\lstinline!zero!} 也不代表它是独立物理急停。
\end{itemize}

\textbf{用法}

\begin{lstlisting}[language=Python]
from typing import TYPE_CHECKING

if TYPE_CHECKING:
    def planned_auto_recover_get(obj: SportClient) -> tuple[int, bool]:
        return obj.auto_recover_get()
\end{lstlisting}

\begin{quote}
{[}!CAUTION{]} 上面的 \passthrough{\lstinline!TYPE\_CHECKING!}
示例只表达计划调用形态。该分支在普通 Python
运行时为假，不代表当前扩展能执行此接口。
\end{quote}

\hypertarget{unitree-sdk2-cpp-robot-go2-sportclient-static-walk-1}{%
\paragraph{\texorpdfstring{\texttt{unitree\_sdk2\_cpp.robot.go2.SportClient.static\_walk}}{unitree\_sdk2\_cpp.robot.go2.SportClient.static\_walk}}\label{unitree-sdk2-cpp-robot-go2-sportclient-static-walk-1}}

对应 C++ SDK 操作
\passthrough{\lstinline!StaticWalk()!}。上游头文件没有可直接生成的业务说明时，本参考只保证签名映射，精确语义需查目标型号协议。

\textbf{签名}

\begin{lstlisting}[language=Python]
def static_walk(self) -> int
\end{lstlisting}

\textbf{可用性与安全性}

\textbf{\passthrough{\lstinline!SIGNATURE\_ONLY!} /
\passthrough{\lstinline!MOTION\_COMMAND!}}：仅用于补全和静态检查。当前二进制没有该
Python 方法；不得按可执行运动接口使用。

\textbf{参数}

无显式参数。实例方法中的 \passthrough{\lstinline!self!} 由 Python
自动传入。

\textbf{返回值}

\begin{longtable}[]{@{}
  >{\raggedright\arraybackslash}p{(\columnwidth - 4\tabcolsep) * \real{0.18}}
  >{\raggedright\arraybackslash}p{(\columnwidth - 4\tabcolsep) * \real{0.20}}
  >{\raggedright\arraybackslash}p{(\columnwidth - 4\tabcolsep) * \real{0.62}}@{}}
\toprule
位置 & 类型 & 含义 \\
\midrule
\endhead
返回值 & \passthrough{\lstinline!int!} & SDK 状态码；Unitree 示例通常以
\passthrough{\lstinline!0!} 表示成功，非零值需按具体服务定义解释。 \\
\bottomrule
\end{longtable}

\textbf{对应 C++}

\begin{itemize}
\tightlist
\item
  类：\passthrough{\lstinline!unitree::robot::go2::SportClient!}
\item
  签名：\passthrough{\lstinline!StaticWalk()!}
\item
  绑定策略：\passthrough{\lstinline!DIRECT!}
\item
  声明位置：\passthrough{\lstinline!include/unitree/robot/go2/sport/sport\_client.hpp:74!}
\end{itemize}

\textbf{说明}

\begin{itemize}
\tightlist
\item
  示例是局部调用片段；\passthrough{\lstinline!obj!}
  和参数变量表示已经按上层业务校验并准备好的对象。
\item
  该条目可以被编辑器和类型检查器识别，但当前运行时可能在属性查找阶段直接失败。
\item
  该方法属于运动命令。方法名包含
  \passthrough{\lstinline!stop!}、\passthrough{\lstinline!damp!} 或
  \passthrough{\lstinline!zero!} 也不代表它是独立物理急停。
\end{itemize}

\textbf{用法}

\begin{lstlisting}[language=Python]
from typing import TYPE_CHECKING

if TYPE_CHECKING:
    def planned_static_walk(obj: SportClient) -> int:
        return obj.static_walk()
\end{lstlisting}

\begin{quote}
{[}!CAUTION{]} 上面的 \passthrough{\lstinline!TYPE\_CHECKING!}
示例只表达计划调用形态。该分支在普通 Python
运行时为假，不代表当前扩展能执行此接口。
\end{quote}

\hypertarget{unitree-sdk2-cpp-robot-go2-sportclient-trot-run-1}{%
\paragraph{\texorpdfstring{\texttt{unitree\_sdk2\_cpp.robot.go2.SportClient.trot\_run}}{unitree\_sdk2\_cpp.robot.go2.SportClient.trot\_run}}\label{unitree-sdk2-cpp-robot-go2-sportclient-trot-run-1}}

对应 C++ SDK 操作
\passthrough{\lstinline!TrotRun()!}。上游头文件没有可直接生成的业务说明时，本参考只保证签名映射，精确语义需查目标型号协议。

\textbf{签名}

\begin{lstlisting}[language=Python]
def trot_run(self) -> int
\end{lstlisting}

\textbf{可用性与安全性}

\textbf{\passthrough{\lstinline!SIGNATURE\_ONLY!} /
\passthrough{\lstinline!MOTION\_COMMAND!}}：仅用于补全和静态检查。当前二进制没有该
Python 方法；不得按可执行运动接口使用。

\textbf{参数}

无显式参数。实例方法中的 \passthrough{\lstinline!self!} 由 Python
自动传入。

\textbf{返回值}

\begin{longtable}[]{@{}
  >{\raggedright\arraybackslash}p{(\columnwidth - 4\tabcolsep) * \real{0.18}}
  >{\raggedright\arraybackslash}p{(\columnwidth - 4\tabcolsep) * \real{0.20}}
  >{\raggedright\arraybackslash}p{(\columnwidth - 4\tabcolsep) * \real{0.62}}@{}}
\toprule
位置 & 类型 & 含义 \\
\midrule
\endhead
返回值 & \passthrough{\lstinline!int!} & SDK 状态码；Unitree 示例通常以
\passthrough{\lstinline!0!} 表示成功，非零值需按具体服务定义解释。 \\
\bottomrule
\end{longtable}

\textbf{对应 C++}

\begin{itemize}
\tightlist
\item
  类：\passthrough{\lstinline!unitree::robot::go2::SportClient!}
\item
  签名：\passthrough{\lstinline!TrotRun()!}
\item
  绑定策略：\passthrough{\lstinline!DIRECT!}
\item
  声明位置：\passthrough{\lstinline!include/unitree/robot/go2/sport/sport\_client.hpp:75!}
\end{itemize}

\textbf{说明}

\begin{itemize}
\tightlist
\item
  示例是局部调用片段；\passthrough{\lstinline!obj!}
  和参数变量表示已经按上层业务校验并准备好的对象。
\item
  该条目可以被编辑器和类型检查器识别，但当前运行时可能在属性查找阶段直接失败。
\item
  该方法属于运动命令。方法名包含
  \passthrough{\lstinline!stop!}、\passthrough{\lstinline!damp!} 或
  \passthrough{\lstinline!zero!} 也不代表它是独立物理急停。
\end{itemize}

\textbf{用法}

\begin{lstlisting}[language=Python]
from typing import TYPE_CHECKING

if TYPE_CHECKING:
    def planned_trot_run(obj: SportClient) -> int:
        return obj.trot_run()
\end{lstlisting}

\begin{quote}
{[}!CAUTION{]} 上面的 \passthrough{\lstinline!TYPE\_CHECKING!}
示例只表达计划调用形态。该分支在普通 Python
运行时为假，不代表当前扩展能执行此接口。
\end{quote}

\hypertarget{unitree-sdk2-cpp-robot-go2-sportclient-economic-gait-1}{%
\paragraph{\texorpdfstring{\texttt{unitree\_sdk2\_cpp.robot.go2.SportClient.economic\_gait}}{unitree\_sdk2\_cpp.robot.go2.SportClient.economic\_gait}}\label{unitree-sdk2-cpp-robot-go2-sportclient-economic-gait-1}}

对应 C++ SDK 操作
\passthrough{\lstinline!EconomicGait()!}。上游头文件没有可直接生成的业务说明时，本参考只保证签名映射，精确语义需查目标型号协议。

\textbf{签名}

\begin{lstlisting}[language=Python]
def economic_gait(self) -> int
\end{lstlisting}

\textbf{可用性与安全性}

\textbf{\passthrough{\lstinline!SIGNATURE\_ONLY!} /
\passthrough{\lstinline!MOTION\_COMMAND!}}：仅用于补全和静态检查。当前二进制没有该
Python 方法；不得按可执行运动接口使用。

\textbf{参数}

无显式参数。实例方法中的 \passthrough{\lstinline!self!} 由 Python
自动传入。

\textbf{返回值}

\begin{longtable}[]{@{}
  >{\raggedright\arraybackslash}p{(\columnwidth - 4\tabcolsep) * \real{0.18}}
  >{\raggedright\arraybackslash}p{(\columnwidth - 4\tabcolsep) * \real{0.20}}
  >{\raggedright\arraybackslash}p{(\columnwidth - 4\tabcolsep) * \real{0.62}}@{}}
\toprule
位置 & 类型 & 含义 \\
\midrule
\endhead
返回值 & \passthrough{\lstinline!int!} & SDK 状态码；Unitree 示例通常以
\passthrough{\lstinline!0!} 表示成功，非零值需按具体服务定义解释。 \\
\bottomrule
\end{longtable}

\textbf{对应 C++}

\begin{itemize}
\tightlist
\item
  类：\passthrough{\lstinline!unitree::robot::go2::SportClient!}
\item
  签名：\passthrough{\lstinline!EconomicGait()!}
\item
  绑定策略：\passthrough{\lstinline!DIRECT!}
\item
  声明位置：\passthrough{\lstinline!include/unitree/robot/go2/sport/sport\_client.hpp:76!}
\end{itemize}

\textbf{说明}

\begin{itemize}
\tightlist
\item
  示例是局部调用片段；\passthrough{\lstinline!obj!}
  和参数变量表示已经按上层业务校验并准备好的对象。
\item
  该条目可以被编辑器和类型检查器识别，但当前运行时可能在属性查找阶段直接失败。
\item
  该方法属于运动命令。方法名包含
  \passthrough{\lstinline!stop!}、\passthrough{\lstinline!damp!} 或
  \passthrough{\lstinline!zero!} 也不代表它是独立物理急停。
\end{itemize}

\textbf{用法}

\begin{lstlisting}[language=Python]
from typing import TYPE_CHECKING

if TYPE_CHECKING:
    def planned_economic_gait(obj: SportClient) -> int:
        return obj.economic_gait()
\end{lstlisting}

\begin{quote}
{[}!CAUTION{]} 上面的 \passthrough{\lstinline!TYPE\_CHECKING!}
示例只表达计划调用形态。该分支在普通 Python
运行时为假，不代表当前扩展能执行此接口。
\end{quote}

\hypertarget{unitree-sdk2-cpp-robot-go2-sportclient-switch-avoid-mode-1}{%
\paragraph{\texorpdfstring{\texttt{unitree\_sdk2\_cpp.robot.go2.SportClient.switch\_avoid\_mode}}{unitree\_sdk2\_cpp.robot.go2.SportClient.switch\_avoid\_mode}}\label{unitree-sdk2-cpp-robot-go2-sportclient-switch-avoid-mode-1}}

对应 C++ SDK 操作
\passthrough{\lstinline!SwitchAvoidMode()!}。上游头文件没有可直接生成的业务说明时，本参考只保证签名映射，精确语义需查目标型号协议。

\textbf{签名}

\begin{lstlisting}[language=Python]
def switch_avoid_mode(self) -> int
\end{lstlisting}

\textbf{可用性与安全性}

\textbf{\passthrough{\lstinline!SIGNATURE\_ONLY!} /
\passthrough{\lstinline!MOTION\_COMMAND!}}：仅用于补全和静态检查。当前二进制没有该
Python 方法；不得按可执行运动接口使用。

\textbf{参数}

无显式参数。实例方法中的 \passthrough{\lstinline!self!} 由 Python
自动传入。

\textbf{返回值}

\begin{longtable}[]{@{}
  >{\raggedright\arraybackslash}p{(\columnwidth - 4\tabcolsep) * \real{0.18}}
  >{\raggedright\arraybackslash}p{(\columnwidth - 4\tabcolsep) * \real{0.20}}
  >{\raggedright\arraybackslash}p{(\columnwidth - 4\tabcolsep) * \real{0.62}}@{}}
\toprule
位置 & 类型 & 含义 \\
\midrule
\endhead
返回值 & \passthrough{\lstinline!int!} & SDK 状态码；Unitree 示例通常以
\passthrough{\lstinline!0!} 表示成功，非零值需按具体服务定义解释。 \\
\bottomrule
\end{longtable}

\textbf{对应 C++}

\begin{itemize}
\tightlist
\item
  类：\passthrough{\lstinline!unitree::robot::go2::SportClient!}
\item
  签名：\passthrough{\lstinline!SwitchAvoidMode()!}
\item
  绑定策略：\passthrough{\lstinline!DIRECT!}
\item
  声明位置：\passthrough{\lstinline!include/unitree/robot/go2/sport/sport\_client.hpp:77!}
\end{itemize}

\textbf{说明}

\begin{itemize}
\tightlist
\item
  示例是局部调用片段；\passthrough{\lstinline!obj!}
  和参数变量表示已经按上层业务校验并准备好的对象。
\item
  该条目可以被编辑器和类型检查器识别，但当前运行时可能在属性查找阶段直接失败。
\item
  该方法属于运动命令。方法名包含
  \passthrough{\lstinline!stop!}、\passthrough{\lstinline!damp!} 或
  \passthrough{\lstinline!zero!} 也不代表它是独立物理急停。
\end{itemize}

\textbf{用法}

\begin{lstlisting}[language=Python]
from typing import TYPE_CHECKING

if TYPE_CHECKING:
    def planned_switch_avoid_mode(obj: SportClient) -> int:
        return obj.switch_avoid_mode()
\end{lstlisting}

\begin{quote}
{[}!CAUTION{]} 上面的 \passthrough{\lstinline!TYPE\_CHECKING!}
示例只表达计划调用形态。该分支在普通 Python
运行时为假，不代表当前扩展能执行此接口。
\end{quote}

\hypertarget{unitree-sdk2-cpp-robot-go2-utrackclient}{%
\subsubsection{\texorpdfstring{\texttt{unitree\_sdk2\_cpp.robot.go2.UtrackClient}}{unitree\_sdk2\_cpp.robot.go2.UtrackClient}}\label{unitree-sdk2-cpp-robot-go2-utrackclient}}

Robot 服务客户端。构造、初始化和具体方法的可用性必须分别检查。

\textbf{导入}

\begin{lstlisting}[language=Python]
from unitree_sdk2_cpp.robot.go2 import UtrackClient
\end{lstlisting}

\textbf{基类}：\passthrough{\lstinline!Client!}

\textbf{方法索引}

\begin{longtable}[]{@{}
  >{\raggedright\arraybackslash}p{(\columnwidth - 6\tabcolsep) * \real{0.18}}
  >{\raggedright\arraybackslash}p{(\columnwidth - 6\tabcolsep) * \real{0.24}}
  >{\raggedright\arraybackslash}p{(\columnwidth - 6\tabcolsep) * \real{0.18}}
  >{\raggedright\arraybackslash}p{(\columnwidth - 6\tabcolsep) * \real{0.40}}@{}}
\toprule
方法 & Python 签名 & 状态 & 安全分类 \\
\midrule
\endhead
\protect\hyperlink{unitree-sdk2-cpp-robot-go2-utrackclient-dunder-init-1}{\passthrough{\lstinline!\_\_init\_\_!}}
& \passthrough{\lstinline!def \_\_init\_\_(self) -> None!} &
\passthrough{\lstinline!AVAILABLE!} &
\passthrough{\lstinline!CONSTRUCTION!} \\
\protect\hyperlink{unitree-sdk2-cpp-robot-go2-utrackclient-init-1}{\passthrough{\lstinline!init!}}
& \passthrough{\lstinline!def init(self) -> None!} &
\passthrough{\lstinline!AVAILABLE!} &
\passthrough{\lstinline!INITIALIZATION!} \\
\protect\hyperlink{unitree-sdk2-cpp-robot-go2-utrackclient-switch-set-1}{\passthrough{\lstinline!switch\_set!}}
& \passthrough{\lstinline!def switch\_set(self, enable: bool) -> int!} &
\passthrough{\lstinline!SIGNATURE\_ONLY!} &
\passthrough{\lstinline!HARDWARE\_SIDE\_EFFECT!} \\
\protect\hyperlink{unitree-sdk2-cpp-robot-go2-utrackclient-switch-get-1}{\passthrough{\lstinline!switch\_get!}}
& \passthrough{\lstinline!def switch\_get(self) -> tuple[int, bool]!} &
\passthrough{\lstinline!SIGNATURE\_ONLY!} &
\passthrough{\lstinline!HARDWARE\_SIDE\_EFFECT!} \\
\protect\hyperlink{unitree-sdk2-cpp-robot-go2-utrackclient-is-tracking-1}{\passthrough{\lstinline!is\_tracking!}}
& \passthrough{\lstinline!def is\_tracking(self) -> tuple[int, bool]!} &
\passthrough{\lstinline!AVAILABLE!} &
\passthrough{\lstinline!READ\_ONLY!} \\
\bottomrule
\end{longtable}

\hypertarget{unitree-sdk2-cpp-robot-go2-utrackclient-dunder-init-1}{%
\paragraph{\texorpdfstring{\texttt{unitree\_sdk2\_cpp.robot.go2.UtrackClient.\_\_init\_\_}}{unitree\_sdk2\_cpp.robot.go2.UtrackClient.\_\_init\_\_}}\label{unitree-sdk2-cpp-robot-go2-utrackclient-dunder-init-1}}

初始化 \passthrough{\lstinline!UtrackClient!}
实例。是否能实际构造取决于下方可用性状态。

\textbf{签名}

\begin{lstlisting}[language=Python]
def __init__(self) -> None
\end{lstlisting}

\textbf{可用性与安全性}

\textbf{\passthrough{\lstinline!AVAILABLE!} /
\passthrough{\lstinline!CONSTRUCTION!}}：当前绑定源码已实现；实际执行条件仍取决于平台和对应
SDK 环境。

\textbf{参数}

无显式参数。实例方法中的 \passthrough{\lstinline!self!} 由 Python
自动传入。

\textbf{返回值}

\begin{longtable}[]{@{}
  >{\raggedright\arraybackslash}p{(\columnwidth - 4\tabcolsep) * \real{0.18}}
  >{\raggedright\arraybackslash}p{(\columnwidth - 4\tabcolsep) * \real{0.20}}
  >{\raggedright\arraybackslash}p{(\columnwidth - 4\tabcolsep) * \real{0.62}}@{}}
\toprule
位置 & 类型 & 含义 \\
\midrule
\endhead
无 & \passthrough{\lstinline!None!} &
\passthrough{\lstinline!\_\_init\_\_!}
本身不返回值；调用类对象时，在构造函数可用的前提下得到该类实例。 \\
\bottomrule
\end{longtable}

\textbf{对应 C++}

\begin{itemize}
\tightlist
\item
  类：\passthrough{\lstinline!unitree::robot::go2::UtrackClient!}
\item
  签名：\passthrough{\lstinline!UtrackClient()!}
\item
  绑定策略：\passthrough{\lstinline!未单独标注!}
\item
  声明位置：\passthrough{\lstinline!include/unitree/robot/go2/utrack/utrack\_client.hpp:15!}
\end{itemize}

\textbf{说明}

\begin{itemize}
\tightlist
\item
  示例是局部调用片段；\passthrough{\lstinline!obj!}
  和参数变量表示已经按上层业务校验并准备好的对象。
\end{itemize}

\textbf{用法}

\begin{lstlisting}[language=Python]
value = UtrackClient()
\end{lstlisting}

\hypertarget{unitree-sdk2-cpp-robot-go2-utrackclient-init-1}{%
\paragraph{\texorpdfstring{\texttt{unitree\_sdk2\_cpp.robot.go2.UtrackClient.init}}{unitree\_sdk2\_cpp.robot.go2.UtrackClient.init}}\label{unitree-sdk2-cpp-robot-go2-utrackclient-init-1}}

初始化当前 SDK 对象所需的底层通道或服务资源。

\textbf{签名}

\begin{lstlisting}[language=Python]
def init(self) -> None
\end{lstlisting}

\textbf{可用性与安全性}

\textbf{\passthrough{\lstinline!AVAILABLE!} /
\passthrough{\lstinline!INITIALIZATION!}}：当前绑定源码已实现；实际执行条件仍取决于平台和对应
SDK 环境。

\textbf{参数}

无显式参数。实例方法中的 \passthrough{\lstinline!self!} 由 Python
自动传入。

\textbf{返回值}

\begin{longtable}[]{@{}
  >{\raggedright\arraybackslash}p{(\columnwidth - 4\tabcolsep) * \real{0.18}}
  >{\raggedright\arraybackslash}p{(\columnwidth - 4\tabcolsep) * \real{0.20}}
  >{\raggedright\arraybackslash}p{(\columnwidth - 4\tabcolsep) * \real{0.62}}@{}}
\toprule
位置 & 类型 & 含义 \\
\midrule
\endhead
无 & \passthrough{\lstinline!None!} & 该方法不返回 Python
值；失败可能表现为异常或底层状态变化。 \\
\bottomrule
\end{longtable}

\textbf{对应 C++}

\begin{itemize}
\tightlist
\item
  类：\passthrough{\lstinline!unitree::robot::go2::UtrackClient!}
\item
  签名：\passthrough{\lstinline!Init()!}
\item
  绑定策略：\passthrough{\lstinline!DIRECT!}
\item
  声明位置：\passthrough{\lstinline!include/unitree/robot/go2/utrack/utrack\_client.hpp:18!}
\end{itemize}

\textbf{说明}

\begin{itemize}
\tightlist
\item
  示例是局部调用片段；\passthrough{\lstinline!obj!}
  和参数变量表示已经按上层业务校验并准备好的对象。
\end{itemize}

\textbf{用法}

\begin{lstlisting}[language=Python]
obj.init()
\end{lstlisting}

\hypertarget{unitree-sdk2-cpp-robot-go2-utrackclient-switch-set-1}{%
\paragraph{\texorpdfstring{\texttt{unitree\_sdk2\_cpp.robot.go2.UtrackClient.switch\_set}}{unitree\_sdk2\_cpp.robot.go2.UtrackClient.switch\_set}}\label{unitree-sdk2-cpp-robot-go2-utrackclient-switch-set-1}}

对应 C++ SDK 操作
\passthrough{\lstinline!SwitchSet(bool)!}。上游头文件没有可直接生成的业务说明时，本参考只保证签名映射，精确语义需查目标型号协议。

\textbf{签名}

\begin{lstlisting}[language=Python]
def switch_set(self, enable: bool) -> int
\end{lstlisting}

\textbf{可用性与安全性}

\textbf{\passthrough{\lstinline!SIGNATURE\_ONLY!} /
\passthrough{\lstinline!HARDWARE\_SIDE\_EFFECT!}}：仅用于补全和静态检查。当前二进制没有该
Python 方法，未来实现还需单独评审硬件副作用。

\textbf{参数}

\begin{longtable}[]{@{}
  >{\raggedright\arraybackslash}p{(\columnwidth - 6\tabcolsep) * \real{0.18}}
  >{\raggedright\arraybackslash}p{(\columnwidth - 6\tabcolsep) * \real{0.24}}
  >{\raggedright\arraybackslash}p{(\columnwidth - 6\tabcolsep) * \real{0.18}}
  >{\raggedright\arraybackslash}p{(\columnwidth - 6\tabcolsep) * \real{0.40}}@{}}
\toprule
名称 & Python 类型 & 默认值 & 含义 \\
\midrule
\endhead
\passthrough{\lstinline!enable!} & \passthrough{\lstinline!bool!} & 必填
& 传给该接口的 \passthrough{\lstinline!enable!} 参数，Python 类型为
\passthrough{\lstinline!bool!}。精确含义、单位、范围和枚举值需查目标型号对应头文件与协议。
对应 C++ 参数 \passthrough{\lstinline!enable: bool!}。
只接受布尔语义。 \\
\bottomrule
\end{longtable}

\textbf{返回值}

\begin{longtable}[]{@{}
  >{\raggedright\arraybackslash}p{(\columnwidth - 4\tabcolsep) * \real{0.18}}
  >{\raggedright\arraybackslash}p{(\columnwidth - 4\tabcolsep) * \real{0.20}}
  >{\raggedright\arraybackslash}p{(\columnwidth - 4\tabcolsep) * \real{0.62}}@{}}
\toprule
位置 & 类型 & 含义 \\
\midrule
\endhead
返回值 & \passthrough{\lstinline!int!} & SDK 状态码；Unitree 示例通常以
\passthrough{\lstinline!0!} 表示成功，非零值需按具体服务定义解释。 \\
\bottomrule
\end{longtable}

\textbf{对应 C++}

\begin{itemize}
\tightlist
\item
  类：\passthrough{\lstinline!unitree::robot::go2::UtrackClient!}
\item
  签名：\passthrough{\lstinline!SwitchSet(bool)!}
\item
  绑定策略：\passthrough{\lstinline!DIRECT!}
\item
  声明位置：\passthrough{\lstinline!include/unitree/robot/go2/utrack/utrack\_client.hpp:20!}
\end{itemize}

\textbf{说明}

\begin{itemize}
\tightlist
\item
  示例是局部调用片段；\passthrough{\lstinline!obj!}
  和参数变量表示已经按上层业务校验并准备好的对象。
\item
  该条目可以被编辑器和类型检查器识别，但当前运行时可能在属性查找阶段直接失败。
\end{itemize}

\textbf{用法}

\begin{lstlisting}[language=Python]
from typing import TYPE_CHECKING

if TYPE_CHECKING:
    def planned_switch_set(obj: UtrackClient, enable: bool) -> int:
        return obj.switch_set(enable=enable)
\end{lstlisting}

\begin{quote}
{[}!CAUTION{]} 上面的 \passthrough{\lstinline!TYPE\_CHECKING!}
示例只表达计划调用形态。该分支在普通 Python
运行时为假，不代表当前扩展能执行此接口。
\end{quote}

\hypertarget{unitree-sdk2-cpp-robot-go2-utrackclient-switch-get-1}{%
\paragraph{\texorpdfstring{\texttt{unitree\_sdk2\_cpp.robot.go2.UtrackClient.switch\_get}}{unitree\_sdk2\_cpp.robot.go2.UtrackClient.switch\_get}}\label{unitree-sdk2-cpp-robot-go2-utrackclient-switch-get-1}}

对应 C++ SDK 操作
\passthrough{\lstinline!SwitchGet(bool \&)!}。上游头文件没有可直接生成的业务说明时，本参考只保证签名映射，精确语义需查目标型号协议。

\textbf{签名}

\begin{lstlisting}[language=Python]
def switch_get(self) -> tuple[int, bool]
\end{lstlisting}

\textbf{可用性与安全性}

\textbf{\passthrough{\lstinline!SIGNATURE\_ONLY!} /
\passthrough{\lstinline!HARDWARE\_SIDE\_EFFECT!}}：仅用于补全和静态检查。当前二进制没有该
Python 方法，未来实现还需单独评审硬件副作用。

\textbf{参数}

无显式参数。实例方法中的 \passthrough{\lstinline!self!} 由 Python
自动传入。

\textbf{返回值}

\begin{longtable}[]{@{}
  >{\raggedright\arraybackslash}p{(\columnwidth - 4\tabcolsep) * \real{0.18}}
  >{\raggedright\arraybackslash}p{(\columnwidth - 4\tabcolsep) * \real{0.20}}
  >{\raggedright\arraybackslash}p{(\columnwidth - 4\tabcolsep) * \real{0.62}}@{}}
\toprule
位置 & 类型 & 含义 \\
\midrule
\endhead
{[}0{]} \passthrough{\lstinline!status!} & \passthrough{\lstinline!int!}
& SDK 状态码；Unitree 示例通常以 \passthrough{\lstinline!0!}
表示成功，非零值需按具体服务错误码解释。 \\
{[}1{]} \passthrough{\lstinline!enable!} &
\passthrough{\lstinline!bool!} & 传给该接口的
\passthrough{\lstinline!enable!} 参数，Python 类型为
\passthrough{\lstinline!bool!}。精确含义、单位、范围和枚举值需查目标型号对应头文件与协议。
对应 C++ 参数 \passthrough{\lstinline!enable: bool \&!}。
只接受布尔语义。 \\
\bottomrule
\end{longtable}

\textbf{对应 C++}

\begin{itemize}
\tightlist
\item
  类：\passthrough{\lstinline!unitree::robot::go2::UtrackClient!}
\item
  签名：\passthrough{\lstinline!SwitchGet(bool \&)!}
\item
  绑定策略：\passthrough{\lstinline!OUTPUT\_WRAPPER!}
\item
  声明位置：\passthrough{\lstinline!include/unitree/robot/go2/utrack/utrack\_client.hpp:21!}
\end{itemize}

\textbf{说明}

\begin{itemize}
\tightlist
\item
  示例是局部调用片段；\passthrough{\lstinline!obj!}
  和参数变量表示已经按上层业务校验并准备好的对象。
\item
  C++ 的可修改输出引用不会作为 Python 入参出现，而是按 C++
  参数顺序追加到返回元组中。
\item
  该条目可以被编辑器和类型检查器识别，但当前运行时可能在属性查找阶段直接失败。
\end{itemize}

\textbf{用法}

\begin{lstlisting}[language=Python]
from typing import TYPE_CHECKING

if TYPE_CHECKING:
    def planned_switch_get(obj: UtrackClient) -> tuple[int, bool]:
        return obj.switch_get()
\end{lstlisting}

\begin{quote}
{[}!CAUTION{]} 上面的 \passthrough{\lstinline!TYPE\_CHECKING!}
示例只表达计划调用形态。该分支在普通 Python
运行时为假，不代表当前扩展能执行此接口。
\end{quote}

\hypertarget{unitree-sdk2-cpp-robot-go2-utrackclient-is-tracking-1}{%
\paragraph{\texorpdfstring{\texttt{unitree\_sdk2\_cpp.robot.go2.UtrackClient.is\_tracking}}{unitree\_sdk2\_cpp.robot.go2.UtrackClient.is\_tracking}}\label{unitree-sdk2-cpp-robot-go2-utrackclient-is-tracking-1}}

查询或检查 \passthrough{\lstinline!tracking!}。

\textbf{签名}

\begin{lstlisting}[language=Python]
def is_tracking(self) -> tuple[int, bool]
\end{lstlisting}

\textbf{可用性与安全性}

\textbf{\passthrough{\lstinline!AVAILABLE!} /
\passthrough{\lstinline!READ\_ONLY!}}：当前绑定源码已实现只读查询。Robot
Client 的服务端查询通常仍需匹配的 Linux
扩展、DDS、网络、机器人和服务；纯本地 getter 则不一定需要全部条件。

\textbf{参数}

无显式参数。实例方法中的 \passthrough{\lstinline!self!} 由 Python
自动传入。

\textbf{返回值}

\begin{longtable}[]{@{}
  >{\raggedright\arraybackslash}p{(\columnwidth - 4\tabcolsep) * \real{0.18}}
  >{\raggedright\arraybackslash}p{(\columnwidth - 4\tabcolsep) * \real{0.20}}
  >{\raggedright\arraybackslash}p{(\columnwidth - 4\tabcolsep) * \real{0.62}}@{}}
\toprule
位置 & 类型 & 含义 \\
\midrule
\endhead
{[}0{]} \passthrough{\lstinline!status!} & \passthrough{\lstinline!int!}
& SDK 状态码；Unitree 示例通常以 \passthrough{\lstinline!0!}
表示成功，非零值需按具体服务错误码解释。 \\
{[}1{]} \passthrough{\lstinline!enable!} &
\passthrough{\lstinline!bool!} & 传给该接口的
\passthrough{\lstinline!enable!} 参数，Python 类型为
\passthrough{\lstinline!bool!}。精确含义、单位、范围和枚举值需查目标型号对应头文件与协议。
对应 C++ 参数 \passthrough{\lstinline!enable: bool \&!}。
只接受布尔语义。 \\
\bottomrule
\end{longtable}

\textbf{对应 C++}

\begin{itemize}
\tightlist
\item
  类：\passthrough{\lstinline!unitree::robot::go2::UtrackClient!}
\item
  签名：\passthrough{\lstinline!IsTracking(bool \&)!}
\item
  绑定策略：\passthrough{\lstinline!OUTPUT\_WRAPPER!}
\item
  声明位置：\passthrough{\lstinline!include/unitree/robot/go2/utrack/utrack\_client.hpp:22!}
\end{itemize}

\textbf{说明}

\begin{itemize}
\tightlist
\item
  示例是局部调用片段；\passthrough{\lstinline!obj!}
  和参数变量表示已经按上层业务校验并准备好的对象。
\item
  C++ 的可修改输出引用不会作为 Python 入参出现，而是按 C++
  参数顺序追加到返回元组中。
\end{itemize}

\textbf{用法}

\begin{lstlisting}[language=Python]
status, enable = obj.is_tracking()
\end{lstlisting}

\begin{quote}
{[}!NOTE{]} 只读查询仍会访问
DDS/机器人服务。必须检查返回状态码，并处理超时和网络中断。
\end{quote}

\hypertarget{unitree-sdk2-cpp-robot-go2-utrackswitchgetdata}{%
\subsubsection{\texorpdfstring{\texttt{unitree\_sdk2\_cpp.robot.go2.UtrackSwitchGetData}}{unitree\_sdk2\_cpp.robot.go2.UtrackSwitchGetData}}\label{unitree-sdk2-cpp-robot-go2-utrackswitchgetdata}}

SDK 数据值类型；公开字段和转换方法见下文。

\textbf{导入}

\begin{lstlisting}[language=Python]
from unitree_sdk2_cpp.robot.go2 import UtrackSwitchGetData
\end{lstlisting}

\textbf{基类}：\passthrough{\lstinline!object!}

\textbf{公开属性}

\begin{longtable}[]{@{}
  >{\raggedright\arraybackslash}p{(\columnwidth - 6\tabcolsep) * \real{0.18}}
  >{\raggedright\arraybackslash}p{(\columnwidth - 6\tabcolsep) * \real{0.24}}
  >{\raggedright\arraybackslash}p{(\columnwidth - 6\tabcolsep) * \real{0.18}}
  >{\raggedright\arraybackslash}p{(\columnwidth - 6\tabcolsep) * \real{0.40}}@{}}
\toprule
名称 & Python 类型 & 含义 & 用法 \\
\midrule
\endhead
\passthrough{\lstinline!m\_enable!} & \passthrough{\lstinline!int!} &
传给该接口的 \passthrough{\lstinline!m!}
\passthrough{\lstinline!enable!} 参数，Python 类型为
\passthrough{\lstinline!int!}。精确含义、单位、范围和枚举值需查目标型号对应头文件与协议。
& \passthrough{\lstinline!value = obj.m\_enable!} /
\passthrough{\lstinline!obj.m\_enable = value!} \\
\bottomrule
\end{longtable}

\textbf{方法索引}

\begin{longtable}[]{@{}
  >{\raggedright\arraybackslash}p{(\columnwidth - 6\tabcolsep) * \real{0.18}}
  >{\raggedright\arraybackslash}p{(\columnwidth - 6\tabcolsep) * \real{0.24}}
  >{\raggedright\arraybackslash}p{(\columnwidth - 6\tabcolsep) * \real{0.18}}
  >{\raggedright\arraybackslash}p{(\columnwidth - 6\tabcolsep) * \real{0.40}}@{}}
\toprule
方法 & Python 签名 & 状态 & 安全分类 \\
\midrule
\endhead
\protect\hyperlink{unitree-sdk2-cpp-robot-go2-utrackswitchgetdata-from-json-1}{\passthrough{\lstinline!from\_json!}}
&
\passthrough{\lstinline!def from\_json(self, json: dict[str, Any]) -> None!}
& \passthrough{\lstinline!SIGNATURE\_ONLY!} &
\passthrough{\lstinline!UNCLASSIFIED!} \\
\protect\hyperlink{unitree-sdk2-cpp-robot-go2-utrackswitchgetdata-to-json-1}{\passthrough{\lstinline!to\_json!}}
&
\passthrough{\lstinline!def to\_json(self, json: dict[str, Any]) -> None!}
& \passthrough{\lstinline!SIGNATURE\_ONLY!} &
\passthrough{\lstinline!UNCLASSIFIED!} \\
\bottomrule
\end{longtable}

\hypertarget{unitree-sdk2-cpp-robot-go2-utrackswitchgetdata-from-json-1}{%
\paragraph{\texorpdfstring{\texttt{unitree\_sdk2\_cpp.robot.go2.UtrackSwitchGetData.from\_json}}{unitree\_sdk2\_cpp.robot.go2.UtrackSwitchGetData.from\_json}}\label{unitree-sdk2-cpp-robot-go2-utrackswitchgetdata-from-json-1}}

计划从 JSON 风格字典读取字段并更新当前 SDK 值对象。

\textbf{签名}

\begin{lstlisting}[language=Python]
def from_json(self, json: dict[str, Any]) -> None
\end{lstlisting}

\textbf{可用性与安全性}

\textbf{\passthrough{\lstinline!SIGNATURE\_ONLY!} /
\passthrough{\lstinline!UNCLASSIFIED!}}：当前只有设计期签名，不能假设已存在于运行时扩展。

\textbf{参数}

\begin{longtable}[]{@{}
  >{\raggedright\arraybackslash}p{(\columnwidth - 6\tabcolsep) * \real{0.18}}
  >{\raggedright\arraybackslash}p{(\columnwidth - 6\tabcolsep) * \real{0.24}}
  >{\raggedright\arraybackslash}p{(\columnwidth - 6\tabcolsep) * \real{0.18}}
  >{\raggedright\arraybackslash}p{(\columnwidth - 6\tabcolsep) * \real{0.40}}@{}}
\toprule
名称 & Python 类型 & 默认值 & 含义 \\
\midrule
\endhead
\passthrough{\lstinline!json!} &
\passthrough{\lstinline!dict[str, Any]!} & 必填 & JSON
风格字典，键和值必须满足对应 C++ \passthrough{\lstinline!JsonMap!}
数据结构的约束。 对应 C++ 参数
\passthrough{\lstinline!json: common::JsonMap \&!}。 \\
\bottomrule
\end{longtable}

\textbf{返回值}

\begin{longtable}[]{@{}
  >{\raggedright\arraybackslash}p{(\columnwidth - 4\tabcolsep) * \real{0.18}}
  >{\raggedright\arraybackslash}p{(\columnwidth - 4\tabcolsep) * \real{0.20}}
  >{\raggedright\arraybackslash}p{(\columnwidth - 4\tabcolsep) * \real{0.62}}@{}}
\toprule
位置 & 类型 & 含义 \\
\midrule
\endhead
无 & \passthrough{\lstinline!None!} & 该方法不返回 Python
值；失败可能表现为异常或底层状态变化。 \\
\bottomrule
\end{longtable}

\textbf{对应 C++}

\begin{itemize}
\tightlist
\item
  类：\passthrough{\lstinline!unitree::robot::go2::UtrackSwitchGetData!}
\item
  签名：\passthrough{\lstinline!fromJson(common::JsonMap \&)!}
\item
  绑定策略：\passthrough{\lstinline!SIGNATURE\_PREVIEW!}
\item
  声明位置：\passthrough{\lstinline!include/unitree/robot/go2/utrack/utrack\_api.hpp:38!}
\end{itemize}

\textbf{说明}

\begin{itemize}
\tightlist
\item
  示例是局部调用片段；\passthrough{\lstinline!obj!}
  和参数变量表示已经按上层业务校验并准备好的对象。
\item
  该条目可以被编辑器和类型检查器识别，但当前运行时可能在属性查找阶段直接失败。
\end{itemize}

\textbf{用法}

\begin{lstlisting}[language=Python]
from typing import TYPE_CHECKING

if TYPE_CHECKING:
    def planned_from_json(obj: UtrackSwitchGetData, json: dict[str, Any]) -> None:
        obj.from_json(json=json)
\end{lstlisting}

\begin{quote}
{[}!CAUTION{]} 上面的 \passthrough{\lstinline!TYPE\_CHECKING!}
示例只表达计划调用形态。该分支在普通 Python
运行时为假，不代表当前扩展能执行此接口。
\end{quote}

\hypertarget{unitree-sdk2-cpp-robot-go2-utrackswitchgetdata-to-json-1}{%
\paragraph{\texorpdfstring{\texttt{unitree\_sdk2\_cpp.robot.go2.UtrackSwitchGetData.to\_json}}{unitree\_sdk2\_cpp.robot.go2.UtrackSwitchGetData.to\_json}}\label{unitree-sdk2-cpp-robot-go2-utrackswitchgetdata-to-json-1}}

计划把当前 SDK 值对象写入 JSON 风格字典。

\textbf{签名}

\begin{lstlisting}[language=Python]
def to_json(self, json: dict[str, Any]) -> None
\end{lstlisting}

\textbf{可用性与安全性}

\textbf{\passthrough{\lstinline!SIGNATURE\_ONLY!} /
\passthrough{\lstinline!UNCLASSIFIED!}}：当前只有设计期签名，不能假设已存在于运行时扩展。

\textbf{参数}

\begin{longtable}[]{@{}
  >{\raggedright\arraybackslash}p{(\columnwidth - 6\tabcolsep) * \real{0.18}}
  >{\raggedright\arraybackslash}p{(\columnwidth - 6\tabcolsep) * \real{0.24}}
  >{\raggedright\arraybackslash}p{(\columnwidth - 6\tabcolsep) * \real{0.18}}
  >{\raggedright\arraybackslash}p{(\columnwidth - 6\tabcolsep) * \real{0.40}}@{}}
\toprule
名称 & Python 类型 & 默认值 & 含义 \\
\midrule
\endhead
\passthrough{\lstinline!json!} &
\passthrough{\lstinline!dict[str, Any]!} & 必填 & JSON
风格字典，键和值必须满足对应 C++ \passthrough{\lstinline!JsonMap!}
数据结构的约束。 对应 C++ 参数
\passthrough{\lstinline!json: common::JsonMap \&!}。 \\
\bottomrule
\end{longtable}

\textbf{返回值}

\begin{longtable}[]{@{}
  >{\raggedright\arraybackslash}p{(\columnwidth - 4\tabcolsep) * \real{0.18}}
  >{\raggedright\arraybackslash}p{(\columnwidth - 4\tabcolsep) * \real{0.20}}
  >{\raggedright\arraybackslash}p{(\columnwidth - 4\tabcolsep) * \real{0.62}}@{}}
\toprule
位置 & 类型 & 含义 \\
\midrule
\endhead
无 & \passthrough{\lstinline!None!} & 该方法不返回 Python
值；失败可能表现为异常或底层状态变化。 \\
\bottomrule
\end{longtable}

\textbf{对应 C++}

\begin{itemize}
\tightlist
\item
  类：\passthrough{\lstinline!unitree::robot::go2::UtrackSwitchGetData!}
\item
  签名：\passthrough{\lstinline!toJson(common::JsonMap \&) const!}
\item
  绑定策略：\passthrough{\lstinline!SIGNATURE\_PREVIEW!}
\item
  声明位置：\passthrough{\lstinline!include/unitree/robot/go2/utrack/utrack\_api.hpp:43!}
\end{itemize}

\textbf{说明}

\begin{itemize}
\tightlist
\item
  示例是局部调用片段；\passthrough{\lstinline!obj!}
  和参数变量表示已经按上层业务校验并准备好的对象。
\item
  该条目可以被编辑器和类型检查器识别，但当前运行时可能在属性查找阶段直接失败。
\end{itemize}

\textbf{用法}

\begin{lstlisting}[language=Python]
from typing import TYPE_CHECKING

if TYPE_CHECKING:
    def planned_to_json(obj: UtrackSwitchGetData, json: dict[str, Any]) -> None:
        obj.to_json(json=json)
\end{lstlisting}

\begin{quote}
{[}!CAUTION{]} 上面的 \passthrough{\lstinline!TYPE\_CHECKING!}
示例只表达计划调用形态。该分支在普通 Python
运行时为假，不代表当前扩展能执行此接口。
\end{quote}

\hypertarget{unitree-sdk2-cpp-robot-go2-utrackswitchsetparameter}{%
\subsubsection{\texorpdfstring{\texttt{unitree\_sdk2\_cpp.robot.go2.UtrackSwitchSetParameter}}{unitree\_sdk2\_cpp.robot.go2.UtrackSwitchSetParameter}}\label{unitree-sdk2-cpp-robot-go2-utrackswitchsetparameter}}

SDK 请求参数值类型；当前可能仅提供设计期签名。

\textbf{导入}

\begin{lstlisting}[language=Python]
from unitree_sdk2_cpp.robot.go2 import UtrackSwitchSetParameter
\end{lstlisting}

\textbf{基类}：\passthrough{\lstinline!object!}

\textbf{公开属性}

\begin{longtable}[]{@{}
  >{\raggedright\arraybackslash}p{(\columnwidth - 6\tabcolsep) * \real{0.18}}
  >{\raggedright\arraybackslash}p{(\columnwidth - 6\tabcolsep) * \real{0.24}}
  >{\raggedright\arraybackslash}p{(\columnwidth - 6\tabcolsep) * \real{0.18}}
  >{\raggedright\arraybackslash}p{(\columnwidth - 6\tabcolsep) * \real{0.40}}@{}}
\toprule
名称 & Python 类型 & 含义 & 用法 \\
\midrule
\endhead
\passthrough{\lstinline!m\_enable!} & \passthrough{\lstinline!int!} &
传给该接口的 \passthrough{\lstinline!m!}
\passthrough{\lstinline!enable!} 参数，Python 类型为
\passthrough{\lstinline!int!}。精确含义、单位、范围和枚举值需查目标型号对应头文件与协议。
& \passthrough{\lstinline!value = obj.m\_enable!} /
\passthrough{\lstinline!obj.m\_enable = value!} \\
\bottomrule
\end{longtable}

\textbf{方法索引}

\begin{longtable}[]{@{}
  >{\raggedright\arraybackslash}p{(\columnwidth - 6\tabcolsep) * \real{0.18}}
  >{\raggedright\arraybackslash}p{(\columnwidth - 6\tabcolsep) * \real{0.24}}
  >{\raggedright\arraybackslash}p{(\columnwidth - 6\tabcolsep) * \real{0.18}}
  >{\raggedright\arraybackslash}p{(\columnwidth - 6\tabcolsep) * \real{0.40}}@{}}
\toprule
方法 & Python 签名 & 状态 & 安全分类 \\
\midrule
\endhead
\protect\hyperlink{unitree-sdk2-cpp-robot-go2-utrackswitchsetparameter-from-json-1}{\passthrough{\lstinline!from\_json!}}
&
\passthrough{\lstinline!def from\_json(self, json: dict[str, Any]) -> None!}
& \passthrough{\lstinline!SIGNATURE\_ONLY!} &
\passthrough{\lstinline!UNCLASSIFIED!} \\
\protect\hyperlink{unitree-sdk2-cpp-robot-go2-utrackswitchsetparameter-to-json-1}{\passthrough{\lstinline!to\_json!}}
&
\passthrough{\lstinline!def to\_json(self, json: dict[str, Any]) -> None!}
& \passthrough{\lstinline!SIGNATURE\_ONLY!} &
\passthrough{\lstinline!UNCLASSIFIED!} \\
\bottomrule
\end{longtable}

\hypertarget{unitree-sdk2-cpp-robot-go2-utrackswitchsetparameter-from-json-1}{%
\paragraph{\texorpdfstring{\texttt{unitree\_sdk2\_cpp.robot.go2.UtrackSwitchSetParameter.from\_json}}{unitree\_sdk2\_cpp.robot.go2.UtrackSwitchSetParameter.from\_json}}\label{unitree-sdk2-cpp-robot-go2-utrackswitchsetparameter-from-json-1}}

计划从 JSON 风格字典读取字段并更新当前 SDK 值对象。

\textbf{签名}

\begin{lstlisting}[language=Python]
def from_json(self, json: dict[str, Any]) -> None
\end{lstlisting}

\textbf{可用性与安全性}

\textbf{\passthrough{\lstinline!SIGNATURE\_ONLY!} /
\passthrough{\lstinline!UNCLASSIFIED!}}：当前只有设计期签名，不能假设已存在于运行时扩展。

\textbf{参数}

\begin{longtable}[]{@{}
  >{\raggedright\arraybackslash}p{(\columnwidth - 6\tabcolsep) * \real{0.18}}
  >{\raggedright\arraybackslash}p{(\columnwidth - 6\tabcolsep) * \real{0.24}}
  >{\raggedright\arraybackslash}p{(\columnwidth - 6\tabcolsep) * \real{0.18}}
  >{\raggedright\arraybackslash}p{(\columnwidth - 6\tabcolsep) * \real{0.40}}@{}}
\toprule
名称 & Python 类型 & 默认值 & 含义 \\
\midrule
\endhead
\passthrough{\lstinline!json!} &
\passthrough{\lstinline!dict[str, Any]!} & 必填 & JSON
风格字典，键和值必须满足对应 C++ \passthrough{\lstinline!JsonMap!}
数据结构的约束。 对应 C++ 参数
\passthrough{\lstinline!json: common::JsonMap \&!}。 \\
\bottomrule
\end{longtable}

\textbf{返回值}

\begin{longtable}[]{@{}
  >{\raggedright\arraybackslash}p{(\columnwidth - 4\tabcolsep) * \real{0.18}}
  >{\raggedright\arraybackslash}p{(\columnwidth - 4\tabcolsep) * \real{0.20}}
  >{\raggedright\arraybackslash}p{(\columnwidth - 4\tabcolsep) * \real{0.62}}@{}}
\toprule
位置 & 类型 & 含义 \\
\midrule
\endhead
无 & \passthrough{\lstinline!None!} & 该方法不返回 Python
值；失败可能表现为异常或底层状态变化。 \\
\bottomrule
\end{longtable}

\textbf{对应 C++}

\begin{itemize}
\tightlist
\item
  类：\passthrough{\lstinline!unitree::robot::go2::UtrackSwitchSetParameter!}
\item
  签名：\passthrough{\lstinline!fromJson(common::JsonMap \&)!}
\item
  绑定策略：\passthrough{\lstinline!SIGNATURE\_PREVIEW!}
\item
  声明位置：\passthrough{\lstinline!include/unitree/robot/go2/utrack/utrack\_api.hpp:22!}
\end{itemize}

\textbf{说明}

\begin{itemize}
\tightlist
\item
  示例是局部调用片段；\passthrough{\lstinline!obj!}
  和参数变量表示已经按上层业务校验并准备好的对象。
\item
  该条目可以被编辑器和类型检查器识别，但当前运行时可能在属性查找阶段直接失败。
\end{itemize}

\textbf{用法}

\begin{lstlisting}[language=Python]
from typing import TYPE_CHECKING

if TYPE_CHECKING:
    def planned_from_json(obj: UtrackSwitchSetParameter, json: dict[str, Any]) -> None:
        obj.from_json(json=json)
\end{lstlisting}

\begin{quote}
{[}!CAUTION{]} 上面的 \passthrough{\lstinline!TYPE\_CHECKING!}
示例只表达计划调用形态。该分支在普通 Python
运行时为假，不代表当前扩展能执行此接口。
\end{quote}

\hypertarget{unitree-sdk2-cpp-robot-go2-utrackswitchsetparameter-to-json-1}{%
\paragraph{\texorpdfstring{\texttt{unitree\_sdk2\_cpp.robot.go2.UtrackSwitchSetParameter.to\_json}}{unitree\_sdk2\_cpp.robot.go2.UtrackSwitchSetParameter.to\_json}}\label{unitree-sdk2-cpp-robot-go2-utrackswitchsetparameter-to-json-1}}

计划把当前 SDK 值对象写入 JSON 风格字典。

\textbf{签名}

\begin{lstlisting}[language=Python]
def to_json(self, json: dict[str, Any]) -> None
\end{lstlisting}

\textbf{可用性与安全性}

\textbf{\passthrough{\lstinline!SIGNATURE\_ONLY!} /
\passthrough{\lstinline!UNCLASSIFIED!}}：当前只有设计期签名，不能假设已存在于运行时扩展。

\textbf{参数}

\begin{longtable}[]{@{}
  >{\raggedright\arraybackslash}p{(\columnwidth - 6\tabcolsep) * \real{0.18}}
  >{\raggedright\arraybackslash}p{(\columnwidth - 6\tabcolsep) * \real{0.24}}
  >{\raggedright\arraybackslash}p{(\columnwidth - 6\tabcolsep) * \real{0.18}}
  >{\raggedright\arraybackslash}p{(\columnwidth - 6\tabcolsep) * \real{0.40}}@{}}
\toprule
名称 & Python 类型 & 默认值 & 含义 \\
\midrule
\endhead
\passthrough{\lstinline!json!} &
\passthrough{\lstinline!dict[str, Any]!} & 必填 & JSON
风格字典，键和值必须满足对应 C++ \passthrough{\lstinline!JsonMap!}
数据结构的约束。 对应 C++ 参数
\passthrough{\lstinline!json: common::JsonMap \&!}。 \\
\bottomrule
\end{longtable}

\textbf{返回值}

\begin{longtable}[]{@{}
  >{\raggedright\arraybackslash}p{(\columnwidth - 4\tabcolsep) * \real{0.18}}
  >{\raggedright\arraybackslash}p{(\columnwidth - 4\tabcolsep) * \real{0.20}}
  >{\raggedright\arraybackslash}p{(\columnwidth - 4\tabcolsep) * \real{0.62}}@{}}
\toprule
位置 & 类型 & 含义 \\
\midrule
\endhead
无 & \passthrough{\lstinline!None!} & 该方法不返回 Python
值；失败可能表现为异常或底层状态变化。 \\
\bottomrule
\end{longtable}

\textbf{对应 C++}

\begin{itemize}
\tightlist
\item
  类：\passthrough{\lstinline!unitree::robot::go2::UtrackSwitchSetParameter!}
\item
  签名：\passthrough{\lstinline!toJson(common::JsonMap \&) const!}
\item
  绑定策略：\passthrough{\lstinline!SIGNATURE\_PREVIEW!}
\item
  声明位置：\passthrough{\lstinline!include/unitree/robot/go2/utrack/utrack\_api.hpp:27!}
\end{itemize}

\textbf{说明}

\begin{itemize}
\tightlist
\item
  示例是局部调用片段；\passthrough{\lstinline!obj!}
  和参数变量表示已经按上层业务校验并准备好的对象。
\item
  该条目可以被编辑器和类型检查器识别，但当前运行时可能在属性查找阶段直接失败。
\end{itemize}

\textbf{用法}

\begin{lstlisting}[language=Python]
from typing import TYPE_CHECKING

if TYPE_CHECKING:
    def planned_to_json(obj: UtrackSwitchSetParameter, json: dict[str, Any]) -> None:
        obj.to_json(json=json)
\end{lstlisting}

\begin{quote}
{[}!CAUTION{]} 上面的 \passthrough{\lstinline!TYPE\_CHECKING!}
示例只表达计划调用形态。该分支在普通 Python
运行时为假，不代表当前扩展能执行此接口。
\end{quote}

\hypertarget{unitree-sdk2-cpp-robot-go2-videoclient}{%
\subsubsection{\texorpdfstring{\texttt{unitree\_sdk2\_cpp.robot.go2.VideoClient}}{unitree\_sdk2\_cpp.robot.go2.VideoClient}}\label{unitree-sdk2-cpp-robot-go2-videoclient}}

Robot 服务客户端。构造、初始化和具体方法的可用性必须分别检查。

\textbf{导入}

\begin{lstlisting}[language=Python]
from unitree_sdk2_cpp.robot.go2 import VideoClient
\end{lstlisting}

\textbf{基类}：\passthrough{\lstinline!Client!}

\textbf{方法索引}

\begin{longtable}[]{@{}
  >{\raggedright\arraybackslash}p{(\columnwidth - 6\tabcolsep) * \real{0.18}}
  >{\raggedright\arraybackslash}p{(\columnwidth - 6\tabcolsep) * \real{0.24}}
  >{\raggedright\arraybackslash}p{(\columnwidth - 6\tabcolsep) * \real{0.18}}
  >{\raggedright\arraybackslash}p{(\columnwidth - 6\tabcolsep) * \real{0.40}}@{}}
\toprule
方法 & Python 签名 & 状态 & 安全分类 \\
\midrule
\endhead
\protect\hyperlink{unitree-sdk2-cpp-robot-go2-videoclient-dunder-init-1}{\passthrough{\lstinline!\_\_init\_\_!}}
& \passthrough{\lstinline!def \_\_init\_\_(self) -> None!} &
\passthrough{\lstinline!AVAILABLE!} &
\passthrough{\lstinline!CONSTRUCTION!} \\
\protect\hyperlink{unitree-sdk2-cpp-robot-go2-videoclient-init-1}{\passthrough{\lstinline!init!}}
& \passthrough{\lstinline!def init(self) -> None!} &
\passthrough{\lstinline!AVAILABLE!} &
\passthrough{\lstinline!INITIALIZATION!} \\
\protect\hyperlink{unitree-sdk2-cpp-robot-go2-videoclient-get-image-sample-1}{\passthrough{\lstinline!get\_image\_sample!}}
&
\passthrough{\lstinline!def get\_image\_sample(self) -> tuple[int, list[int]]!}
& \passthrough{\lstinline!AVAILABLE!} &
\passthrough{\lstinline!READ\_ONLY!} \\
\bottomrule
\end{longtable}

\hypertarget{unitree-sdk2-cpp-robot-go2-videoclient-dunder-init-1}{%
\paragraph{\texorpdfstring{\texttt{unitree\_sdk2\_cpp.robot.go2.VideoClient.\_\_init\_\_}}{unitree\_sdk2\_cpp.robot.go2.VideoClient.\_\_init\_\_}}\label{unitree-sdk2-cpp-robot-go2-videoclient-dunder-init-1}}

初始化 \passthrough{\lstinline!VideoClient!}
实例。是否能实际构造取决于下方可用性状态。

\textbf{签名}

\begin{lstlisting}[language=Python]
def __init__(self) -> None
\end{lstlisting}

\textbf{可用性与安全性}

\textbf{\passthrough{\lstinline!AVAILABLE!} /
\passthrough{\lstinline!CONSTRUCTION!}}：当前绑定源码已实现；实际执行条件仍取决于平台和对应
SDK 环境。

\textbf{参数}

无显式参数。实例方法中的 \passthrough{\lstinline!self!} 由 Python
自动传入。

\textbf{返回值}

\begin{longtable}[]{@{}
  >{\raggedright\arraybackslash}p{(\columnwidth - 4\tabcolsep) * \real{0.18}}
  >{\raggedright\arraybackslash}p{(\columnwidth - 4\tabcolsep) * \real{0.20}}
  >{\raggedright\arraybackslash}p{(\columnwidth - 4\tabcolsep) * \real{0.62}}@{}}
\toprule
位置 & 类型 & 含义 \\
\midrule
\endhead
无 & \passthrough{\lstinline!None!} &
\passthrough{\lstinline!\_\_init\_\_!}
本身不返回值；调用类对象时，在构造函数可用的前提下得到该类实例。 \\
\bottomrule
\end{longtable}

\textbf{对应 C++}

\begin{itemize}
\tightlist
\item
  类：\passthrough{\lstinline!unitree::robot::go2::VideoClient!}
\item
  签名：\passthrough{\lstinline!VideoClient()!}
\item
  绑定策略：\passthrough{\lstinline!未单独标注!}
\item
  声明位置：\passthrough{\lstinline!include/unitree/robot/go2/video/video\_client.hpp:18!}
\end{itemize}

\textbf{说明}

\begin{itemize}
\tightlist
\item
  示例是局部调用片段；\passthrough{\lstinline!obj!}
  和参数变量表示已经按上层业务校验并准备好的对象。
\end{itemize}

\textbf{用法}

\begin{lstlisting}[language=Python]
value = VideoClient()
\end{lstlisting}

\hypertarget{unitree-sdk2-cpp-robot-go2-videoclient-init-1}{%
\paragraph{\texorpdfstring{\texttt{unitree\_sdk2\_cpp.robot.go2.VideoClient.init}}{unitree\_sdk2\_cpp.robot.go2.VideoClient.init}}\label{unitree-sdk2-cpp-robot-go2-videoclient-init-1}}

初始化当前 SDK 对象所需的底层通道或服务资源。

\textbf{签名}

\begin{lstlisting}[language=Python]
def init(self) -> None
\end{lstlisting}

\textbf{可用性与安全性}

\textbf{\passthrough{\lstinline!AVAILABLE!} /
\passthrough{\lstinline!INITIALIZATION!}}：当前绑定源码已实现；实际执行条件仍取决于平台和对应
SDK 环境。

\textbf{参数}

无显式参数。实例方法中的 \passthrough{\lstinline!self!} 由 Python
自动传入。

\textbf{返回值}

\begin{longtable}[]{@{}
  >{\raggedright\arraybackslash}p{(\columnwidth - 4\tabcolsep) * \real{0.18}}
  >{\raggedright\arraybackslash}p{(\columnwidth - 4\tabcolsep) * \real{0.20}}
  >{\raggedright\arraybackslash}p{(\columnwidth - 4\tabcolsep) * \real{0.62}}@{}}
\toprule
位置 & 类型 & 含义 \\
\midrule
\endhead
无 & \passthrough{\lstinline!None!} & 该方法不返回 Python
值；失败可能表现为异常或底层状态变化。 \\
\bottomrule
\end{longtable}

\textbf{对应 C++}

\begin{itemize}
\tightlist
\item
  类：\passthrough{\lstinline!unitree::robot::go2::VideoClient!}
\item
  签名：\passthrough{\lstinline!Init()!}
\item
  绑定策略：\passthrough{\lstinline!DIRECT!}
\item
  声明位置：\passthrough{\lstinline!include/unitree/robot/go2/video/video\_client.hpp:21!}
\end{itemize}

\textbf{说明}

\begin{itemize}
\tightlist
\item
  示例是局部调用片段；\passthrough{\lstinline!obj!}
  和参数变量表示已经按上层业务校验并准备好的对象。
\end{itemize}

\textbf{用法}

\begin{lstlisting}[language=Python]
obj.init()
\end{lstlisting}

\hypertarget{unitree-sdk2-cpp-robot-go2-videoclient-get-image-sample-1}{%
\paragraph{\texorpdfstring{\texttt{unitree\_sdk2\_cpp.robot.go2.VideoClient.get\_image\_sample}}{unitree\_sdk2\_cpp.robot.go2.VideoClient.get\_image\_sample}}\label{unitree-sdk2-cpp-robot-go2-videoclient-get-image-sample-1}}

查询或检查 图像 \passthrough{\lstinline!sample!}。

\textbf{签名}

\begin{lstlisting}[language=Python]
def get_image_sample(self) -> tuple[int, list[int]]
\end{lstlisting}

\textbf{可用性与安全性}

\textbf{\passthrough{\lstinline!AVAILABLE!} /
\passthrough{\lstinline!READ\_ONLY!}}：当前绑定源码已实现只读查询。Robot
Client 的服务端查询通常仍需匹配的 Linux
扩展、DDS、网络、机器人和服务；纯本地 getter 则不一定需要全部条件。

\textbf{参数}

无显式参数。实例方法中的 \passthrough{\lstinline!self!} 由 Python
自动传入。

\textbf{返回值}

\begin{longtable}[]{@{}
  >{\raggedright\arraybackslash}p{(\columnwidth - 4\tabcolsep) * \real{0.18}}
  >{\raggedright\arraybackslash}p{(\columnwidth - 4\tabcolsep) * \real{0.20}}
  >{\raggedright\arraybackslash}p{(\columnwidth - 4\tabcolsep) * \real{0.62}}@{}}
\toprule
位置 & 类型 & 含义 \\
\midrule
\endhead
{[}0{]} \passthrough{\lstinline!status!} & \passthrough{\lstinline!int!}
& SDK 状态码；Unitree 示例通常以 \passthrough{\lstinline!0!}
表示成功，非零值需按具体服务错误码解释。 \\
{[}1{]} \passthrough{\lstinline!output\_1!} &
\passthrough{\lstinline!list[int]!} & 传给该接口的
\passthrough{\lstinline!output!} \passthrough{\lstinline!1!}
参数，Python 类型为
\passthrough{\lstinline!list[int]!}。精确含义、单位、范围和枚举值需查目标型号对应头文件与协议。
对应 C++ 参数
\passthrough{\lstinline!output\_1: std::vector<uint8\_t> \&!}。
底层是可变长度 vector；元素约束：取值范围为 0 到 255。 \\
\bottomrule
\end{longtable}

\textbf{对应 C++}

\begin{itemize}
\tightlist
\item
  类：\passthrough{\lstinline!unitree::robot::go2::VideoClient!}
\item
  签名：\passthrough{\lstinline!GetImageSample(std::vector<uint8\_t> \&)!}
\item
  绑定策略：\passthrough{\lstinline!OUTPUT\_WRAPPER!}
\item
  声明位置：\passthrough{\lstinline!include/unitree/robot/go2/video/video\_client.hpp:27!}
\end{itemize}

\textbf{说明}

\begin{itemize}
\tightlist
\item
  示例是局部调用片段；\passthrough{\lstinline!obj!}
  和参数变量表示已经按上层业务校验并准备好的对象。
\item
  C++ 的可修改输出引用不会作为 Python 入参出现，而是按 C++
  参数顺序追加到返回元组中。
\end{itemize}

\textbf{用法}

\begin{lstlisting}[language=Python]
status, output_1 = obj.get_image_sample()
\end{lstlisting}

\begin{quote}
{[}!NOTE{]} 只读查询仍会访问
DDS/机器人服务。必须检查返回状态码，并处理超时和网络中断。
\end{quote}

\hypertarget{unitree-sdk2-cpp-robot-go2-vuiclient}{%
\subsubsection{\texorpdfstring{\texttt{unitree\_sdk2\_cpp.robot.go2.VuiClient}}{unitree\_sdk2\_cpp.robot.go2.VuiClient}}\label{unitree-sdk2-cpp-robot-go2-vuiclient}}

Robot 服务客户端。构造、初始化和具体方法的可用性必须分别检查。

\textbf{导入}

\begin{lstlisting}[language=Python]
from unitree_sdk2_cpp.robot.go2 import VuiClient
\end{lstlisting}

\textbf{基类}：\passthrough{\lstinline!Client!}

\textbf{方法索引}

\begin{longtable}[]{@{}
  >{\raggedright\arraybackslash}p{(\columnwidth - 6\tabcolsep) * \real{0.18}}
  >{\raggedright\arraybackslash}p{(\columnwidth - 6\tabcolsep) * \real{0.24}}
  >{\raggedright\arraybackslash}p{(\columnwidth - 6\tabcolsep) * \real{0.18}}
  >{\raggedright\arraybackslash}p{(\columnwidth - 6\tabcolsep) * \real{0.40}}@{}}
\toprule
方法 & Python 签名 & 状态 & 安全分类 \\
\midrule
\endhead
\protect\hyperlink{unitree-sdk2-cpp-robot-go2-vuiclient-dunder-init-1}{\passthrough{\lstinline!\_\_init\_\_!}}
& \passthrough{\lstinline!def \_\_init\_\_(self) -> None!} &
\passthrough{\lstinline!AVAILABLE!} &
\passthrough{\lstinline!CONSTRUCTION!} \\
\protect\hyperlink{unitree-sdk2-cpp-robot-go2-vuiclient-init-1}{\passthrough{\lstinline!init!}}
& \passthrough{\lstinline!def init(self) -> None!} &
\passthrough{\lstinline!AVAILABLE!} &
\passthrough{\lstinline!INITIALIZATION!} \\
\protect\hyperlink{unitree-sdk2-cpp-robot-go2-vuiclient-set-switch-1}{\passthrough{\lstinline!set\_switch!}}
& \passthrough{\lstinline!def set\_switch(self, enable: int) -> int!} &
\passthrough{\lstinline!SIGNATURE\_ONLY!} &
\passthrough{\lstinline!HARDWARE\_SIDE\_EFFECT!} \\
\protect\hyperlink{unitree-sdk2-cpp-robot-go2-vuiclient-get-switch-1}{\passthrough{\lstinline!get\_switch!}}
& \passthrough{\lstinline!def get\_switch(self) -> tuple[int, int]!} &
\passthrough{\lstinline!AVAILABLE!} &
\passthrough{\lstinline!READ\_ONLY!} \\
\protect\hyperlink{unitree-sdk2-cpp-robot-go2-vuiclient-set-volume-1}{\passthrough{\lstinline!set\_volume!}}
& \passthrough{\lstinline!def set\_volume(self, level: int) -> int!} &
\passthrough{\lstinline!SIGNATURE\_ONLY!} &
\passthrough{\lstinline!HARDWARE\_SIDE\_EFFECT!} \\
\protect\hyperlink{unitree-sdk2-cpp-robot-go2-vuiclient-get-volume-1}{\passthrough{\lstinline!get\_volume!}}
& \passthrough{\lstinline!def get\_volume(self) -> tuple[int, int]!} &
\passthrough{\lstinline!AVAILABLE!} &
\passthrough{\lstinline!READ\_ONLY!} \\
\protect\hyperlink{unitree-sdk2-cpp-robot-go2-vuiclient-set-brightness-1}{\passthrough{\lstinline!set\_brightness!}}
& \passthrough{\lstinline!def set\_brightness(self, level: int) -> int!}
& \passthrough{\lstinline!SIGNATURE\_ONLY!} &
\passthrough{\lstinline!HARDWARE\_SIDE\_EFFECT!} \\
\protect\hyperlink{unitree-sdk2-cpp-robot-go2-vuiclient-get-brightness-1}{\passthrough{\lstinline!get\_brightness!}}
& \passthrough{\lstinline!def get\_brightness(self) -> tuple[int, int]!}
& \passthrough{\lstinline!AVAILABLE!} &
\passthrough{\lstinline!READ\_ONLY!} \\
\bottomrule
\end{longtable}

\hypertarget{unitree-sdk2-cpp-robot-go2-vuiclient-dunder-init-1}{%
\paragraph{\texorpdfstring{\texttt{unitree\_sdk2\_cpp.robot.go2.VuiClient.\_\_init\_\_}}{unitree\_sdk2\_cpp.robot.go2.VuiClient.\_\_init\_\_}}\label{unitree-sdk2-cpp-robot-go2-vuiclient-dunder-init-1}}

初始化 \passthrough{\lstinline!VuiClient!}
实例。是否能实际构造取决于下方可用性状态。

\textbf{签名}

\begin{lstlisting}[language=Python]
def __init__(self) -> None
\end{lstlisting}

\textbf{可用性与安全性}

\textbf{\passthrough{\lstinline!AVAILABLE!} /
\passthrough{\lstinline!CONSTRUCTION!}}：当前绑定源码已实现；实际执行条件仍取决于平台和对应
SDK 环境。

\textbf{参数}

无显式参数。实例方法中的 \passthrough{\lstinline!self!} 由 Python
自动传入。

\textbf{返回值}

\begin{longtable}[]{@{}
  >{\raggedright\arraybackslash}p{(\columnwidth - 4\tabcolsep) * \real{0.18}}
  >{\raggedright\arraybackslash}p{(\columnwidth - 4\tabcolsep) * \real{0.20}}
  >{\raggedright\arraybackslash}p{(\columnwidth - 4\tabcolsep) * \real{0.62}}@{}}
\toprule
位置 & 类型 & 含义 \\
\midrule
\endhead
无 & \passthrough{\lstinline!None!} &
\passthrough{\lstinline!\_\_init\_\_!}
本身不返回值；调用类对象时，在构造函数可用的前提下得到该类实例。 \\
\bottomrule
\end{longtable}

\textbf{对应 C++}

\begin{itemize}
\tightlist
\item
  类：\passthrough{\lstinline!unitree::robot::go2::VuiClient!}
\item
  签名：\passthrough{\lstinline!VuiClient()!}
\item
  绑定策略：\passthrough{\lstinline!未单独标注!}
\item
  声明位置：\passthrough{\lstinline!include/unitree/robot/go2/vui/vui\_client.hpp:18!}
\end{itemize}

\textbf{说明}

\begin{itemize}
\tightlist
\item
  示例是局部调用片段；\passthrough{\lstinline!obj!}
  和参数变量表示已经按上层业务校验并准备好的对象。
\end{itemize}

\textbf{用法}

\begin{lstlisting}[language=Python]
value = VuiClient()
\end{lstlisting}

\hypertarget{unitree-sdk2-cpp-robot-go2-vuiclient-init-1}{%
\paragraph{\texorpdfstring{\texttt{unitree\_sdk2\_cpp.robot.go2.VuiClient.init}}{unitree\_sdk2\_cpp.robot.go2.VuiClient.init}}\label{unitree-sdk2-cpp-robot-go2-vuiclient-init-1}}

初始化当前 SDK 对象所需的底层通道或服务资源。

\textbf{签名}

\begin{lstlisting}[language=Python]
def init(self) -> None
\end{lstlisting}

\textbf{可用性与安全性}

\textbf{\passthrough{\lstinline!AVAILABLE!} /
\passthrough{\lstinline!INITIALIZATION!}}：当前绑定源码已实现；实际执行条件仍取决于平台和对应
SDK 环境。

\textbf{参数}

无显式参数。实例方法中的 \passthrough{\lstinline!self!} 由 Python
自动传入。

\textbf{返回值}

\begin{longtable}[]{@{}
  >{\raggedright\arraybackslash}p{(\columnwidth - 4\tabcolsep) * \real{0.18}}
  >{\raggedright\arraybackslash}p{(\columnwidth - 4\tabcolsep) * \real{0.20}}
  >{\raggedright\arraybackslash}p{(\columnwidth - 4\tabcolsep) * \real{0.62}}@{}}
\toprule
位置 & 类型 & 含义 \\
\midrule
\endhead
无 & \passthrough{\lstinline!None!} & 该方法不返回 Python
值；失败可能表现为异常或底层状态变化。 \\
\bottomrule
\end{longtable}

\textbf{对应 C++}

\begin{itemize}
\tightlist
\item
  类：\passthrough{\lstinline!unitree::robot::go2::VuiClient!}
\item
  签名：\passthrough{\lstinline!Init()!}
\item
  绑定策略：\passthrough{\lstinline!DIRECT!}
\item
  声明位置：\passthrough{\lstinline!include/unitree/robot/go2/vui/vui\_client.hpp:21!}
\end{itemize}

\textbf{说明}

\begin{itemize}
\tightlist
\item
  示例是局部调用片段；\passthrough{\lstinline!obj!}
  和参数变量表示已经按上层业务校验并准备好的对象。
\end{itemize}

\textbf{用法}

\begin{lstlisting}[language=Python]
obj.init()
\end{lstlisting}

\hypertarget{unitree-sdk2-cpp-robot-go2-vuiclient-set-switch-1}{%
\paragraph{\texorpdfstring{\texttt{unitree\_sdk2\_cpp.robot.go2.VuiClient.set\_switch}}{unitree\_sdk2\_cpp.robot.go2.VuiClient.set\_switch}}\label{unitree-sdk2-cpp-robot-go2-vuiclient-set-switch-1}}

设置 开关。具体副作用和安全边界见下方状态。

\textbf{签名}

\begin{lstlisting}[language=Python]
def set_switch(self, enable: int) -> int
\end{lstlisting}

\textbf{可用性与安全性}

\textbf{\passthrough{\lstinline!SIGNATURE\_ONLY!} /
\passthrough{\lstinline!HARDWARE\_SIDE\_EFFECT!}}：仅用于补全和静态检查。当前二进制没有该
Python 方法，未来实现还需单独评审硬件副作用。

\textbf{参数}

\begin{longtable}[]{@{}
  >{\raggedright\arraybackslash}p{(\columnwidth - 6\tabcolsep) * \real{0.18}}
  >{\raggedright\arraybackslash}p{(\columnwidth - 6\tabcolsep) * \real{0.24}}
  >{\raggedright\arraybackslash}p{(\columnwidth - 6\tabcolsep) * \real{0.18}}
  >{\raggedright\arraybackslash}p{(\columnwidth - 6\tabcolsep) * \real{0.40}}@{}}
\toprule
名称 & Python 类型 & 默认值 & 含义 \\
\midrule
\endhead
\passthrough{\lstinline!enable!} & \passthrough{\lstinline!int!} & 必填
& 传给该接口的 \passthrough{\lstinline!enable!} 参数，Python 类型为
\passthrough{\lstinline!int!}。精确含义、单位、范围和枚举值需查目标型号对应头文件与协议。
对应 C++ 参数 \passthrough{\lstinline!enable: int!}。 \\
\bottomrule
\end{longtable}

\textbf{返回值}

\begin{longtable}[]{@{}
  >{\raggedright\arraybackslash}p{(\columnwidth - 4\tabcolsep) * \real{0.18}}
  >{\raggedright\arraybackslash}p{(\columnwidth - 4\tabcolsep) * \real{0.20}}
  >{\raggedright\arraybackslash}p{(\columnwidth - 4\tabcolsep) * \real{0.62}}@{}}
\toprule
位置 & 类型 & 含义 \\
\midrule
\endhead
返回值 & \passthrough{\lstinline!int!} & SDK 状态码；Unitree 示例通常以
\passthrough{\lstinline!0!} 表示成功，非零值需按具体服务定义解释。 \\
\bottomrule
\end{longtable}

\textbf{对应 C++}

\begin{itemize}
\tightlist
\item
  类：\passthrough{\lstinline!unitree::robot::go2::VuiClient!}
\item
  签名：\passthrough{\lstinline!SetSwitch(int)!}
\item
  绑定策略：\passthrough{\lstinline!DIRECT!}
\item
  声明位置：\passthrough{\lstinline!include/unitree/robot/go2/vui/vui\_client.hpp:27!}
\end{itemize}

\textbf{说明}

\begin{itemize}
\tightlist
\item
  示例是局部调用片段；\passthrough{\lstinline!obj!}
  和参数变量表示已经按上层业务校验并准备好的对象。
\item
  该条目可以被编辑器和类型检查器识别，但当前运行时可能在属性查找阶段直接失败。
\end{itemize}

\textbf{用法}

\begin{lstlisting}[language=Python]
from typing import TYPE_CHECKING

if TYPE_CHECKING:
    def planned_set_switch(obj: VuiClient, enable: int) -> int:
        return obj.set_switch(enable=enable)
\end{lstlisting}

\begin{quote}
{[}!CAUTION{]} 上面的 \passthrough{\lstinline!TYPE\_CHECKING!}
示例只表达计划调用形态。该分支在普通 Python
运行时为假，不代表当前扩展能执行此接口。
\end{quote}

\hypertarget{unitree-sdk2-cpp-robot-go2-vuiclient-get-switch-1}{%
\paragraph{\texorpdfstring{\texttt{unitree\_sdk2\_cpp.robot.go2.VuiClient.get\_switch}}{unitree\_sdk2\_cpp.robot.go2.VuiClient.get\_switch}}\label{unitree-sdk2-cpp-robot-go2-vuiclient-get-switch-1}}

查询或检查 开关。

\textbf{签名}

\begin{lstlisting}[language=Python]
def get_switch(self) -> tuple[int, int]
\end{lstlisting}

\textbf{可用性与安全性}

\textbf{\passthrough{\lstinline!AVAILABLE!} /
\passthrough{\lstinline!READ\_ONLY!}}：当前绑定源码已实现只读查询。Robot
Client 的服务端查询通常仍需匹配的 Linux
扩展、DDS、网络、机器人和服务；纯本地 getter 则不一定需要全部条件。

\textbf{参数}

无显式参数。实例方法中的 \passthrough{\lstinline!self!} 由 Python
自动传入。

\textbf{返回值}

\begin{longtable}[]{@{}
  >{\raggedright\arraybackslash}p{(\columnwidth - 4\tabcolsep) * \real{0.18}}
  >{\raggedright\arraybackslash}p{(\columnwidth - 4\tabcolsep) * \real{0.20}}
  >{\raggedright\arraybackslash}p{(\columnwidth - 4\tabcolsep) * \real{0.62}}@{}}
\toprule
位置 & 类型 & 含义 \\
\midrule
\endhead
{[}0{]} \passthrough{\lstinline!status!} & \passthrough{\lstinline!int!}
& SDK 状态码；Unitree 示例通常以 \passthrough{\lstinline!0!}
表示成功，非零值需按具体服务错误码解释。 \\
{[}1{]} \passthrough{\lstinline!output\_1!} &
\passthrough{\lstinline!int!} & 传给该接口的
\passthrough{\lstinline!output!} \passthrough{\lstinline!1!}
参数，Python 类型为
\passthrough{\lstinline!int!}。精确含义、单位、范围和枚举值需查目标型号对应头文件与协议。
对应 C++ 参数 \passthrough{\lstinline!output\_1: int \&!}。 \\
\bottomrule
\end{longtable}

\textbf{对应 C++}

\begin{itemize}
\tightlist
\item
  类：\passthrough{\lstinline!unitree::robot::go2::VuiClient!}
\item
  签名：\passthrough{\lstinline!GetSwitch(int \&)!}
\item
  绑定策略：\passthrough{\lstinline!OUTPUT\_WRAPPER!}
\item
  声明位置：\passthrough{\lstinline!include/unitree/robot/go2/vui/vui\_client.hpp:33!}
\end{itemize}

\textbf{说明}

\begin{itemize}
\tightlist
\item
  示例是局部调用片段；\passthrough{\lstinline!obj!}
  和参数变量表示已经按上层业务校验并准备好的对象。
\item
  C++ 的可修改输出引用不会作为 Python 入参出现，而是按 C++
  参数顺序追加到返回元组中。
\end{itemize}

\textbf{用法}

\begin{lstlisting}[language=Python]
status, output_1 = obj.get_switch()
\end{lstlisting}

\begin{quote}
{[}!NOTE{]} 只读查询仍会访问
DDS/机器人服务。必须检查返回状态码，并处理超时和网络中断。
\end{quote}

\hypertarget{unitree-sdk2-cpp-robot-go2-vuiclient-set-volume-1}{%
\paragraph{\texorpdfstring{\texttt{unitree\_sdk2\_cpp.robot.go2.VuiClient.set\_volume}}{unitree\_sdk2\_cpp.robot.go2.VuiClient.set\_volume}}\label{unitree-sdk2-cpp-robot-go2-vuiclient-set-volume-1}}

设置 音量。具体副作用和安全边界见下方状态。

\textbf{签名}

\begin{lstlisting}[language=Python]
def set_volume(self, level: int) -> int
\end{lstlisting}

\textbf{可用性与安全性}

\textbf{\passthrough{\lstinline!SIGNATURE\_ONLY!} /
\passthrough{\lstinline!HARDWARE\_SIDE\_EFFECT!}}：仅用于补全和静态检查。当前二进制没有该
Python 方法，未来实现还需单独评审硬件副作用。

\textbf{参数}

\begin{longtable}[]{@{}
  >{\raggedright\arraybackslash}p{(\columnwidth - 6\tabcolsep) * \real{0.18}}
  >{\raggedright\arraybackslash}p{(\columnwidth - 6\tabcolsep) * \real{0.24}}
  >{\raggedright\arraybackslash}p{(\columnwidth - 6\tabcolsep) * \real{0.18}}
  >{\raggedright\arraybackslash}p{(\columnwidth - 6\tabcolsep) * \real{0.40}}@{}}
\toprule
名称 & Python 类型 & 默认值 & 含义 \\
\midrule
\endhead
\passthrough{\lstinline!level!} & \passthrough{\lstinline!int!} & 必填 &
传给该接口的 \passthrough{\lstinline!level!} 参数，Python 类型为
\passthrough{\lstinline!int!}。精确含义、单位、范围和枚举值需查目标型号对应头文件与协议。
对应 C++ 参数 \passthrough{\lstinline!level: int!}。 \\
\bottomrule
\end{longtable}

\textbf{返回值}

\begin{longtable}[]{@{}
  >{\raggedright\arraybackslash}p{(\columnwidth - 4\tabcolsep) * \real{0.18}}
  >{\raggedright\arraybackslash}p{(\columnwidth - 4\tabcolsep) * \real{0.20}}
  >{\raggedright\arraybackslash}p{(\columnwidth - 4\tabcolsep) * \real{0.62}}@{}}
\toprule
位置 & 类型 & 含义 \\
\midrule
\endhead
返回值 & \passthrough{\lstinline!int!} & SDK 状态码；Unitree 示例通常以
\passthrough{\lstinline!0!} 表示成功，非零值需按具体服务定义解释。 \\
\bottomrule
\end{longtable}

\textbf{对应 C++}

\begin{itemize}
\tightlist
\item
  类：\passthrough{\lstinline!unitree::robot::go2::VuiClient!}
\item
  签名：\passthrough{\lstinline!SetVolume(int)!}
\item
  绑定策略：\passthrough{\lstinline!DIRECT!}
\item
  声明位置：\passthrough{\lstinline!include/unitree/robot/go2/vui/vui\_client.hpp:39!}
\end{itemize}

\textbf{说明}

\begin{itemize}
\tightlist
\item
  示例是局部调用片段；\passthrough{\lstinline!obj!}
  和参数变量表示已经按上层业务校验并准备好的对象。
\item
  该条目可以被编辑器和类型检查器识别，但当前运行时可能在属性查找阶段直接失败。
\end{itemize}

\textbf{用法}

\begin{lstlisting}[language=Python]
from typing import TYPE_CHECKING

if TYPE_CHECKING:
    def planned_set_volume(obj: VuiClient, level: int) -> int:
        return obj.set_volume(level=level)
\end{lstlisting}

\begin{quote}
{[}!CAUTION{]} 上面的 \passthrough{\lstinline!TYPE\_CHECKING!}
示例只表达计划调用形态。该分支在普通 Python
运行时为假，不代表当前扩展能执行此接口。
\end{quote}

\hypertarget{unitree-sdk2-cpp-robot-go2-vuiclient-get-volume-1}{%
\paragraph{\texorpdfstring{\texttt{unitree\_sdk2\_cpp.robot.go2.VuiClient.get\_volume}}{unitree\_sdk2\_cpp.robot.go2.VuiClient.get\_volume}}\label{unitree-sdk2-cpp-robot-go2-vuiclient-get-volume-1}}

查询或检查 音量。

\textbf{签名}

\begin{lstlisting}[language=Python]
def get_volume(self) -> tuple[int, int]
\end{lstlisting}

\textbf{可用性与安全性}

\textbf{\passthrough{\lstinline!AVAILABLE!} /
\passthrough{\lstinline!READ\_ONLY!}}：当前绑定源码已实现只读查询。Robot
Client 的服务端查询通常仍需匹配的 Linux
扩展、DDS、网络、机器人和服务；纯本地 getter 则不一定需要全部条件。

\textbf{参数}

无显式参数。实例方法中的 \passthrough{\lstinline!self!} 由 Python
自动传入。

\textbf{返回值}

\begin{longtable}[]{@{}
  >{\raggedright\arraybackslash}p{(\columnwidth - 4\tabcolsep) * \real{0.18}}
  >{\raggedright\arraybackslash}p{(\columnwidth - 4\tabcolsep) * \real{0.20}}
  >{\raggedright\arraybackslash}p{(\columnwidth - 4\tabcolsep) * \real{0.62}}@{}}
\toprule
位置 & 类型 & 含义 \\
\midrule
\endhead
{[}0{]} \passthrough{\lstinline!status!} & \passthrough{\lstinline!int!}
& SDK 状态码；Unitree 示例通常以 \passthrough{\lstinline!0!}
表示成功，非零值需按具体服务错误码解释。 \\
{[}1{]} \passthrough{\lstinline!output\_1!} &
\passthrough{\lstinline!int!} & 传给该接口的
\passthrough{\lstinline!output!} \passthrough{\lstinline!1!}
参数，Python 类型为
\passthrough{\lstinline!int!}。精确含义、单位、范围和枚举值需查目标型号对应头文件与协议。
对应 C++ 参数 \passthrough{\lstinline!output\_1: int \&!}。 \\
\bottomrule
\end{longtable}

\textbf{对应 C++}

\begin{itemize}
\tightlist
\item
  类：\passthrough{\lstinline!unitree::robot::go2::VuiClient!}
\item
  签名：\passthrough{\lstinline!GetVolume(int \&)!}
\item
  绑定策略：\passthrough{\lstinline!OUTPUT\_WRAPPER!}
\item
  声明位置：\passthrough{\lstinline!include/unitree/robot/go2/vui/vui\_client.hpp:45!}
\end{itemize}

\textbf{说明}

\begin{itemize}
\tightlist
\item
  示例是局部调用片段；\passthrough{\lstinline!obj!}
  和参数变量表示已经按上层业务校验并准备好的对象。
\item
  C++ 的可修改输出引用不会作为 Python 入参出现，而是按 C++
  参数顺序追加到返回元组中。
\end{itemize}

\textbf{用法}

\begin{lstlisting}[language=Python]
status, output_1 = obj.get_volume()
\end{lstlisting}

\begin{quote}
{[}!NOTE{]} 只读查询仍会访问
DDS/机器人服务。必须检查返回状态码，并处理超时和网络中断。
\end{quote}

\hypertarget{unitree-sdk2-cpp-robot-go2-vuiclient-set-brightness-1}{%
\paragraph{\texorpdfstring{\texttt{unitree\_sdk2\_cpp.robot.go2.VuiClient.set\_brightness}}{unitree\_sdk2\_cpp.robot.go2.VuiClient.set\_brightness}}\label{unitree-sdk2-cpp-robot-go2-vuiclient-set-brightness-1}}

设置 亮度。具体副作用和安全边界见下方状态。

\textbf{签名}

\begin{lstlisting}[language=Python]
def set_brightness(self, level: int) -> int
\end{lstlisting}

\textbf{可用性与安全性}

\textbf{\passthrough{\lstinline!SIGNATURE\_ONLY!} /
\passthrough{\lstinline!HARDWARE\_SIDE\_EFFECT!}}：仅用于补全和静态检查。当前二进制没有该
Python 方法，未来实现还需单独评审硬件副作用。

\textbf{参数}

\begin{longtable}[]{@{}
  >{\raggedright\arraybackslash}p{(\columnwidth - 6\tabcolsep) * \real{0.18}}
  >{\raggedright\arraybackslash}p{(\columnwidth - 6\tabcolsep) * \real{0.24}}
  >{\raggedright\arraybackslash}p{(\columnwidth - 6\tabcolsep) * \real{0.18}}
  >{\raggedright\arraybackslash}p{(\columnwidth - 6\tabcolsep) * \real{0.40}}@{}}
\toprule
名称 & Python 类型 & 默认值 & 含义 \\
\midrule
\endhead
\passthrough{\lstinline!level!} & \passthrough{\lstinline!int!} & 必填 &
传给该接口的 \passthrough{\lstinline!level!} 参数，Python 类型为
\passthrough{\lstinline!int!}。精确含义、单位、范围和枚举值需查目标型号对应头文件与协议。
对应 C++ 参数 \passthrough{\lstinline!level: int!}。 \\
\bottomrule
\end{longtable}

\textbf{返回值}

\begin{longtable}[]{@{}
  >{\raggedright\arraybackslash}p{(\columnwidth - 4\tabcolsep) * \real{0.18}}
  >{\raggedright\arraybackslash}p{(\columnwidth - 4\tabcolsep) * \real{0.20}}
  >{\raggedright\arraybackslash}p{(\columnwidth - 4\tabcolsep) * \real{0.62}}@{}}
\toprule
位置 & 类型 & 含义 \\
\midrule
\endhead
返回值 & \passthrough{\lstinline!int!} & SDK 状态码；Unitree 示例通常以
\passthrough{\lstinline!0!} 表示成功，非零值需按具体服务定义解释。 \\
\bottomrule
\end{longtable}

\textbf{对应 C++}

\begin{itemize}
\tightlist
\item
  类：\passthrough{\lstinline!unitree::robot::go2::VuiClient!}
\item
  签名：\passthrough{\lstinline!SetBrightness(int)!}
\item
  绑定策略：\passthrough{\lstinline!DIRECT!}
\item
  声明位置：\passthrough{\lstinline!include/unitree/robot/go2/vui/vui\_client.hpp:51!}
\end{itemize}

\textbf{说明}

\begin{itemize}
\tightlist
\item
  示例是局部调用片段；\passthrough{\lstinline!obj!}
  和参数变量表示已经按上层业务校验并准备好的对象。
\item
  该条目可以被编辑器和类型检查器识别，但当前运行时可能在属性查找阶段直接失败。
\end{itemize}

\textbf{用法}

\begin{lstlisting}[language=Python]
from typing import TYPE_CHECKING

if TYPE_CHECKING:
    def planned_set_brightness(obj: VuiClient, level: int) -> int:
        return obj.set_brightness(level=level)
\end{lstlisting}

\begin{quote}
{[}!CAUTION{]} 上面的 \passthrough{\lstinline!TYPE\_CHECKING!}
示例只表达计划调用形态。该分支在普通 Python
运行时为假，不代表当前扩展能执行此接口。
\end{quote}

\hypertarget{unitree-sdk2-cpp-robot-go2-vuiclient-get-brightness-1}{%
\paragraph{\texorpdfstring{\texttt{unitree\_sdk2\_cpp.robot.go2.VuiClient.get\_brightness}}{unitree\_sdk2\_cpp.robot.go2.VuiClient.get\_brightness}}\label{unitree-sdk2-cpp-robot-go2-vuiclient-get-brightness-1}}

查询或检查 亮度。

\textbf{签名}

\begin{lstlisting}[language=Python]
def get_brightness(self) -> tuple[int, int]
\end{lstlisting}

\textbf{可用性与安全性}

\textbf{\passthrough{\lstinline!AVAILABLE!} /
\passthrough{\lstinline!READ\_ONLY!}}：当前绑定源码已实现只读查询。Robot
Client 的服务端查询通常仍需匹配的 Linux
扩展、DDS、网络、机器人和服务；纯本地 getter 则不一定需要全部条件。

\textbf{参数}

无显式参数。实例方法中的 \passthrough{\lstinline!self!} 由 Python
自动传入。

\textbf{返回值}

\begin{longtable}[]{@{}
  >{\raggedright\arraybackslash}p{(\columnwidth - 4\tabcolsep) * \real{0.18}}
  >{\raggedright\arraybackslash}p{(\columnwidth - 4\tabcolsep) * \real{0.20}}
  >{\raggedright\arraybackslash}p{(\columnwidth - 4\tabcolsep) * \real{0.62}}@{}}
\toprule
位置 & 类型 & 含义 \\
\midrule
\endhead
{[}0{]} \passthrough{\lstinline!status!} & \passthrough{\lstinline!int!}
& SDK 状态码；Unitree 示例通常以 \passthrough{\lstinline!0!}
表示成功，非零值需按具体服务错误码解释。 \\
{[}1{]} \passthrough{\lstinline!output\_1!} &
\passthrough{\lstinline!int!} & 传给该接口的
\passthrough{\lstinline!output!} \passthrough{\lstinline!1!}
参数，Python 类型为
\passthrough{\lstinline!int!}。精确含义、单位、范围和枚举值需查目标型号对应头文件与协议。
对应 C++ 参数 \passthrough{\lstinline!output\_1: int \&!}。 \\
\bottomrule
\end{longtable}

\textbf{对应 C++}

\begin{itemize}
\tightlist
\item
  类：\passthrough{\lstinline!unitree::robot::go2::VuiClient!}
\item
  签名：\passthrough{\lstinline!GetBrightness(int \&)!}
\item
  绑定策略：\passthrough{\lstinline!OUTPUT\_WRAPPER!}
\item
  声明位置：\passthrough{\lstinline!include/unitree/robot/go2/vui/vui\_client.hpp:57!}
\end{itemize}

\textbf{说明}

\begin{itemize}
\tightlist
\item
  示例是局部调用片段；\passthrough{\lstinline!obj!}
  和参数变量表示已经按上层业务校验并准备好的对象。
\item
  C++ 的可修改输出引用不会作为 Python 入参出现，而是按 C++
  参数顺序追加到返回元组中。
\end{itemize}

\textbf{用法}

\begin{lstlisting}[language=Python]
status, output_1 = obj.get_brightness()
\end{lstlisting}

\begin{quote}
{[}!NOTE{]} 只读查询仍会访问
DDS/机器人服务。必须检查返回状态码，并处理超时和网络中断。
\end{quote}

\hypertarget{unitree-sdk2-cpp-robot-go2-stpathpoint}{%
\subsubsection{\texorpdfstring{\texttt{unitree\_sdk2\_cpp.robot.go2.stPathPoint}}{unitree\_sdk2\_cpp.robot.go2.stPathPoint}}\label{unitree-sdk2-cpp-robot-go2-stpathpoint}}

SDK 数据值类型；公开字段和转换方法见下文。

\textbf{导入}

\begin{lstlisting}[language=Python]
from unitree_sdk2_cpp.robot.go2 import stPathPoint
\end{lstlisting}

\textbf{基类}：\passthrough{\lstinline!object!}

\textbf{公开属性}

\begin{longtable}[]{@{}
  >{\raggedright\arraybackslash}p{(\columnwidth - 6\tabcolsep) * \real{0.18}}
  >{\raggedright\arraybackslash}p{(\columnwidth - 6\tabcolsep) * \real{0.24}}
  >{\raggedright\arraybackslash}p{(\columnwidth - 6\tabcolsep) * \real{0.18}}
  >{\raggedright\arraybackslash}p{(\columnwidth - 6\tabcolsep) * \real{0.40}}@{}}
\toprule
名称 & Python 类型 & 含义 & 用法 \\
\midrule
\endhead
\passthrough{\lstinline!time\_from\_start!} &
\passthrough{\lstinline!float!} & 传给该接口的
\passthrough{\lstinline!time!} \passthrough{\lstinline!from!}
\passthrough{\lstinline!start!} 参数，Python 类型为
\passthrough{\lstinline!float!}。精确含义、单位、范围和枚举值需查目标型号对应头文件与协议。
& \passthrough{\lstinline!value = obj.time\_from\_start!} /
\passthrough{\lstinline!obj.time\_from\_start = value!} \\
\passthrough{\lstinline!x!} & \passthrough{\lstinline!float!} & X
方向位置或分量参数。参考系、单位和范围必须查当前方法及目标型号协议。 &
\passthrough{\lstinline!value = obj.x!} /
\passthrough{\lstinline!obj.x = value!} \\
\passthrough{\lstinline!y!} & \passthrough{\lstinline!float!} & Y
方向位置或分量参数。参考系、单位和范围必须查当前方法及目标型号协议。 &
\passthrough{\lstinline!value = obj.y!} /
\passthrough{\lstinline!obj.y = value!} \\
\passthrough{\lstinline!yaw!} & \passthrough{\lstinline!float!} &
偏航角或偏航目标参数。参考系、单位和范围必须查目标型号协议。 &
\passthrough{\lstinline!value = obj.yaw!} /
\passthrough{\lstinline!obj.yaw = value!} \\
\passthrough{\lstinline!vx!} & \passthrough{\lstinline!float!} & X
方向速度参数。坐标系、单位、符号和安全范围必须查目标型号运动协议。 &
\passthrough{\lstinline!value = obj.vx!} /
\passthrough{\lstinline!obj.vx = value!} \\
\passthrough{\lstinline!vy!} & \passthrough{\lstinline!float!} & Y
方向速度参数。坐标系、单位、符号和安全范围必须查目标型号运动协议。 &
\passthrough{\lstinline!value = obj.vy!} /
\passthrough{\lstinline!obj.vy = value!} \\
\passthrough{\lstinline!vyaw!} & \passthrough{\lstinline!float!} &
偏航角速度参数。单位、符号和安全范围必须查目标型号运动协议。 &
\passthrough{\lstinline!value = obj.vyaw!} /
\passthrough{\lstinline!obj.vyaw = value!} \\
\bottomrule
\end{longtable}

\begin{center}\rule{0.5\linewidth}{0.5pt}\end{center}

\hypertarget{unitree-sdk2-cpp-robot-h1}{%
\subsection{H1 Robot API}\label{unitree-sdk2-cpp-robot-h1}}

模块：\passthrough{\lstinline!unitree\_sdk2\_cpp.robot.h1!}

\hypertarget{ux7c7bux7d22ux5f15-12}{%
\subsubsection{类索引}\label{ux7c7bux7d22ux5f15-12}}

\begin{longtable}[]{@{}
  >{\raggedright\arraybackslash}p{(\columnwidth - 4\tabcolsep) * \real{0.18}}
  >{\raggedleft\arraybackslash}p{(\columnwidth - 4\tabcolsep) * \real{0.20}}
  >{\raggedleft\arraybackslash}p{(\columnwidth - 4\tabcolsep) * \real{0.62}}@{}}
\toprule
类 & 公开函数签名 & 属性 \\
\midrule
\endhead
\protect\hyperlink{unitree-sdk2-cpp-robot-h1-jsonizedatavecfloat}{\passthrough{\lstinline!JsonizeDataVecFloat!}}
& 3 & 0 \\
\protect\hyperlink{unitree-sdk2-cpp-robot-h1-jsonizetargetpos}{\passthrough{\lstinline!JsonizeTargetPos!}}
& 3 & 0 \\
\protect\hyperlink{unitree-sdk2-cpp-robot-h1-jsonizevelocitycommand}{\passthrough{\lstinline!JsonizeVelocityCommand!}}
& 3 & 0 \\
\protect\hyperlink{unitree-sdk2-cpp-robot-h1-lococlient}{\passthrough{\lstinline!LocoClient!}}
& 34 & 0 \\
\bottomrule
\end{longtable}

\hypertarget{unitree-sdk2-cpp-robot-h1-jsonizedatavecfloat}{%
\subsubsection{\texorpdfstring{\texttt{unitree\_sdk2\_cpp.robot.h1.JsonizeDataVecFloat}}{unitree\_sdk2\_cpp.robot.h1.JsonizeDataVecFloat}}\label{unitree-sdk2-cpp-robot-h1-jsonizedatavecfloat}}

SDK 类型签名预览。请逐项查看构造函数和方法的可用性。

\textbf{导入}

\begin{lstlisting}[language=Python]
from unitree_sdk2_cpp.robot.h1 import JsonizeDataVecFloat
\end{lstlisting}

\textbf{基类}：\passthrough{\lstinline!object!}

\textbf{公开属性}

\begin{longtable}[]{@{}
  >{\raggedright\arraybackslash}p{(\columnwidth - 6\tabcolsep) * \real{0.18}}
  >{\raggedright\arraybackslash}p{(\columnwidth - 6\tabcolsep) * \real{0.24}}
  >{\raggedright\arraybackslash}p{(\columnwidth - 6\tabcolsep) * \real{0.18}}
  >{\raggedright\arraybackslash}p{(\columnwidth - 6\tabcolsep) * \real{0.40}}@{}}
\toprule
名称 & Python 类型 & 含义 & 用法 \\
\midrule
\endhead
\passthrough{\lstinline!data!} & \passthrough{\lstinline!list[float]!} &
传给该接口的 数据 参数，Python 类型为
\passthrough{\lstinline!list[float]!}。精确含义、单位、范围和枚举值需查目标型号对应头文件与协议。
& \passthrough{\lstinline!value = obj.data!} /
\passthrough{\lstinline!obj.data = value!} \\
\bottomrule
\end{longtable}

\textbf{方法索引}

\begin{longtable}[]{@{}
  >{\raggedright\arraybackslash}p{(\columnwidth - 6\tabcolsep) * \real{0.18}}
  >{\raggedright\arraybackslash}p{(\columnwidth - 6\tabcolsep) * \real{0.24}}
  >{\raggedright\arraybackslash}p{(\columnwidth - 6\tabcolsep) * \real{0.18}}
  >{\raggedright\arraybackslash}p{(\columnwidth - 6\tabcolsep) * \real{0.40}}@{}}
\toprule
方法 & Python 签名 & 状态 & 安全分类 \\
\midrule
\endhead
\protect\hyperlink{unitree-sdk2-cpp-robot-h1-jsonizedatavecfloat-dunder-init-1}{\passthrough{\lstinline!\_\_init\_\_!}}
& \passthrough{\lstinline!def \_\_init\_\_(self) -> None!} &
\passthrough{\lstinline!SIGNATURE\_ONLY!} &
\passthrough{\lstinline!CONSTRUCTION!} \\
\protect\hyperlink{unitree-sdk2-cpp-robot-h1-jsonizedatavecfloat-from-json-1}{\passthrough{\lstinline!from\_json!}}
&
\passthrough{\lstinline!def from\_json(self, json: dict[str, Any]) -> None!}
& \passthrough{\lstinline!SIGNATURE\_ONLY!} &
\passthrough{\lstinline!UNCLASSIFIED!} \\
\protect\hyperlink{unitree-sdk2-cpp-robot-h1-jsonizedatavecfloat-to-json-1}{\passthrough{\lstinline!to\_json!}}
&
\passthrough{\lstinline!def to\_json(self, json: dict[str, Any]) -> None!}
& \passthrough{\lstinline!SIGNATURE\_ONLY!} &
\passthrough{\lstinline!UNCLASSIFIED!} \\
\bottomrule
\end{longtable}

\hypertarget{unitree-sdk2-cpp-robot-h1-jsonizedatavecfloat-dunder-init-1}{%
\paragraph{\texorpdfstring{\texttt{unitree\_sdk2\_cpp.robot.h1.JsonizeDataVecFloat.\_\_init\_\_}}{unitree\_sdk2\_cpp.robot.h1.JsonizeDataVecFloat.\_\_init\_\_}}\label{unitree-sdk2-cpp-robot-h1-jsonizedatavecfloat-dunder-init-1}}

初始化 \passthrough{\lstinline!JsonizeDataVecFloat!}
实例。是否能实际构造取决于下方可用性状态。

\textbf{签名}

\begin{lstlisting}[language=Python]
def __init__(self) -> None
\end{lstlisting}

\textbf{可用性与安全性}

\textbf{\passthrough{\lstinline!SIGNATURE\_ONLY!} /
\passthrough{\lstinline!CONSTRUCTION!}}：当前只有设计期签名，不能假设已存在于运行时扩展。

\textbf{参数}

无显式参数。实例方法中的 \passthrough{\lstinline!self!} 由 Python
自动传入。

\textbf{返回值}

\begin{longtable}[]{@{}
  >{\raggedright\arraybackslash}p{(\columnwidth - 4\tabcolsep) * \real{0.18}}
  >{\raggedright\arraybackslash}p{(\columnwidth - 4\tabcolsep) * \real{0.20}}
  >{\raggedright\arraybackslash}p{(\columnwidth - 4\tabcolsep) * \real{0.62}}@{}}
\toprule
位置 & 类型 & 含义 \\
\midrule
\endhead
无 & \passthrough{\lstinline!None!} &
\passthrough{\lstinline!\_\_init\_\_!}
本身不返回值；调用类对象时，在构造函数可用的前提下得到该类实例。 \\
\bottomrule
\end{longtable}

\textbf{对应 C++}

\begin{itemize}
\tightlist
\item
  类：\passthrough{\lstinline!unitree::robot::h1::JsonizeDataVecFloat!}
\item
  签名：\passthrough{\lstinline!JsonizeDataVecFloat()!}
\item
  绑定策略：\passthrough{\lstinline!未单独标注!}
\item
  声明位置：\passthrough{\lstinline!include/unitree/robot/h1/loco/h1\_loco\_api.hpp:41!}
\end{itemize}

\textbf{说明}

\begin{itemize}
\tightlist
\item
  示例是局部调用片段；\passthrough{\lstinline!obj!}
  和参数变量表示已经按上层业务校验并准备好的对象。
\item
  该条目可以被编辑器和类型检查器识别，但当前运行时可能在属性查找阶段直接失败。
\end{itemize}

\textbf{用法}

\begin{lstlisting}[language=Python]
from typing import TYPE_CHECKING

if TYPE_CHECKING:
    def planned_construct() -> JsonizeDataVecFloat:
        return JsonizeDataVecFloat()
\end{lstlisting}

\begin{quote}
{[}!CAUTION{]} 上面的 \passthrough{\lstinline!TYPE\_CHECKING!}
示例只表达计划调用形态。该分支在普通 Python
运行时为假，不代表当前扩展能执行此接口。
\end{quote}

\hypertarget{unitree-sdk2-cpp-robot-h1-jsonizedatavecfloat-from-json-1}{%
\paragraph{\texorpdfstring{\texttt{unitree\_sdk2\_cpp.robot.h1.JsonizeDataVecFloat.from\_json}}{unitree\_sdk2\_cpp.robot.h1.JsonizeDataVecFloat.from\_json}}\label{unitree-sdk2-cpp-robot-h1-jsonizedatavecfloat-from-json-1}}

计划从 JSON 风格字典读取字段并更新当前 SDK 值对象。

\textbf{签名}

\begin{lstlisting}[language=Python]
def from_json(self, json: dict[str, Any]) -> None
\end{lstlisting}

\textbf{可用性与安全性}

\textbf{\passthrough{\lstinline!SIGNATURE\_ONLY!} /
\passthrough{\lstinline!UNCLASSIFIED!}}：当前只有设计期签名，不能假设已存在于运行时扩展。

\textbf{参数}

\begin{longtable}[]{@{}
  >{\raggedright\arraybackslash}p{(\columnwidth - 6\tabcolsep) * \real{0.18}}
  >{\raggedright\arraybackslash}p{(\columnwidth - 6\tabcolsep) * \real{0.24}}
  >{\raggedright\arraybackslash}p{(\columnwidth - 6\tabcolsep) * \real{0.18}}
  >{\raggedright\arraybackslash}p{(\columnwidth - 6\tabcolsep) * \real{0.40}}@{}}
\toprule
名称 & Python 类型 & 默认值 & 含义 \\
\midrule
\endhead
\passthrough{\lstinline!json!} &
\passthrough{\lstinline!dict[str, Any]!} & 必填 & JSON
风格字典，键和值必须满足对应 C++ \passthrough{\lstinline!JsonMap!}
数据结构的约束。 对应 C++ 参数
\passthrough{\lstinline!json: common::JsonMap \&!}。 \\
\bottomrule
\end{longtable}

\textbf{返回值}

\begin{longtable}[]{@{}
  >{\raggedright\arraybackslash}p{(\columnwidth - 4\tabcolsep) * \real{0.18}}
  >{\raggedright\arraybackslash}p{(\columnwidth - 4\tabcolsep) * \real{0.20}}
  >{\raggedright\arraybackslash}p{(\columnwidth - 4\tabcolsep) * \real{0.62}}@{}}
\toprule
位置 & 类型 & 含义 \\
\midrule
\endhead
无 & \passthrough{\lstinline!None!} & 该方法不返回 Python
值；失败可能表现为异常或底层状态变化。 \\
\bottomrule
\end{longtable}

\textbf{对应 C++}

\begin{itemize}
\tightlist
\item
  类：\passthrough{\lstinline!unitree::robot::h1::JsonizeDataVecFloat!}
\item
  签名：\passthrough{\lstinline!fromJson(common::JsonMap \&)!}
\item
  绑定策略：\passthrough{\lstinline!SIGNATURE\_PREVIEW!}
\item
  声明位置：\passthrough{\lstinline!include/unitree/robot/h1/loco/h1\_loco\_api.hpp:44!}
\end{itemize}

\textbf{说明}

\begin{itemize}
\tightlist
\item
  示例是局部调用片段；\passthrough{\lstinline!obj!}
  和参数变量表示已经按上层业务校验并准备好的对象。
\item
  该条目可以被编辑器和类型检查器识别，但当前运行时可能在属性查找阶段直接失败。
\end{itemize}

\textbf{用法}

\begin{lstlisting}[language=Python]
from typing import TYPE_CHECKING

if TYPE_CHECKING:
    def planned_from_json(obj: JsonizeDataVecFloat, json: dict[str, Any]) -> None:
        obj.from_json(json=json)
\end{lstlisting}

\begin{quote}
{[}!CAUTION{]} 上面的 \passthrough{\lstinline!TYPE\_CHECKING!}
示例只表达计划调用形态。该分支在普通 Python
运行时为假，不代表当前扩展能执行此接口。
\end{quote}

\hypertarget{unitree-sdk2-cpp-robot-h1-jsonizedatavecfloat-to-json-1}{%
\paragraph{\texorpdfstring{\texttt{unitree\_sdk2\_cpp.robot.h1.JsonizeDataVecFloat.to\_json}}{unitree\_sdk2\_cpp.robot.h1.JsonizeDataVecFloat.to\_json}}\label{unitree-sdk2-cpp-robot-h1-jsonizedatavecfloat-to-json-1}}

计划把当前 SDK 值对象写入 JSON 风格字典。

\textbf{签名}

\begin{lstlisting}[language=Python]
def to_json(self, json: dict[str, Any]) -> None
\end{lstlisting}

\textbf{可用性与安全性}

\textbf{\passthrough{\lstinline!SIGNATURE\_ONLY!} /
\passthrough{\lstinline!UNCLASSIFIED!}}：当前只有设计期签名，不能假设已存在于运行时扩展。

\textbf{参数}

\begin{longtable}[]{@{}
  >{\raggedright\arraybackslash}p{(\columnwidth - 6\tabcolsep) * \real{0.18}}
  >{\raggedright\arraybackslash}p{(\columnwidth - 6\tabcolsep) * \real{0.24}}
  >{\raggedright\arraybackslash}p{(\columnwidth - 6\tabcolsep) * \real{0.18}}
  >{\raggedright\arraybackslash}p{(\columnwidth - 6\tabcolsep) * \real{0.40}}@{}}
\toprule
名称 & Python 类型 & 默认值 & 含义 \\
\midrule
\endhead
\passthrough{\lstinline!json!} &
\passthrough{\lstinline!dict[str, Any]!} & 必填 & JSON
风格字典，键和值必须满足对应 C++ \passthrough{\lstinline!JsonMap!}
数据结构的约束。 对应 C++ 参数
\passthrough{\lstinline!json: common::JsonMap \&!}。 \\
\bottomrule
\end{longtable}

\textbf{返回值}

\begin{longtable}[]{@{}
  >{\raggedright\arraybackslash}p{(\columnwidth - 4\tabcolsep) * \real{0.18}}
  >{\raggedright\arraybackslash}p{(\columnwidth - 4\tabcolsep) * \real{0.20}}
  >{\raggedright\arraybackslash}p{(\columnwidth - 4\tabcolsep) * \real{0.62}}@{}}
\toprule
位置 & 类型 & 含义 \\
\midrule
\endhead
无 & \passthrough{\lstinline!None!} & 该方法不返回 Python
值；失败可能表现为异常或底层状态变化。 \\
\bottomrule
\end{longtable}

\textbf{对应 C++}

\begin{itemize}
\tightlist
\item
  类：\passthrough{\lstinline!unitree::robot::h1::JsonizeDataVecFloat!}
\item
  签名：\passthrough{\lstinline!toJson(common::JsonMap \&) const!}
\item
  绑定策略：\passthrough{\lstinline!SIGNATURE\_PREVIEW!}
\item
  声明位置：\passthrough{\lstinline!include/unitree/robot/h1/loco/h1\_loco\_api.hpp:46!}
\end{itemize}

\textbf{说明}

\begin{itemize}
\tightlist
\item
  示例是局部调用片段；\passthrough{\lstinline!obj!}
  和参数变量表示已经按上层业务校验并准备好的对象。
\item
  该条目可以被编辑器和类型检查器识别，但当前运行时可能在属性查找阶段直接失败。
\end{itemize}

\textbf{用法}

\begin{lstlisting}[language=Python]
from typing import TYPE_CHECKING

if TYPE_CHECKING:
    def planned_to_json(obj: JsonizeDataVecFloat, json: dict[str, Any]) -> None:
        obj.to_json(json=json)
\end{lstlisting}

\begin{quote}
{[}!CAUTION{]} 上面的 \passthrough{\lstinline!TYPE\_CHECKING!}
示例只表达计划调用形态。该分支在普通 Python
运行时为假，不代表当前扩展能执行此接口。
\end{quote}

\hypertarget{unitree-sdk2-cpp-robot-h1-jsonizetargetpos}{%
\subsubsection{\texorpdfstring{\texttt{unitree\_sdk2\_cpp.robot.h1.JsonizeTargetPos}}{unitree\_sdk2\_cpp.robot.h1.JsonizeTargetPos}}\label{unitree-sdk2-cpp-robot-h1-jsonizetargetpos}}

SDK 类型签名预览。请逐项查看构造函数和方法的可用性。

\textbf{导入}

\begin{lstlisting}[language=Python]
from unitree_sdk2_cpp.robot.h1 import JsonizeTargetPos
\end{lstlisting}

\textbf{基类}：\passthrough{\lstinline!object!}

\textbf{公开属性}

\begin{longtable}[]{@{}
  >{\raggedright\arraybackslash}p{(\columnwidth - 6\tabcolsep) * \real{0.18}}
  >{\raggedright\arraybackslash}p{(\columnwidth - 6\tabcolsep) * \real{0.24}}
  >{\raggedright\arraybackslash}p{(\columnwidth - 6\tabcolsep) * \real{0.18}}
  >{\raggedright\arraybackslash}p{(\columnwidth - 6\tabcolsep) * \real{0.40}}@{}}
\toprule
名称 & Python 类型 & 含义 & 用法 \\
\midrule
\endhead
\passthrough{\lstinline!x!} & \passthrough{\lstinline!float!} & X
方向位置或分量参数。参考系、单位和范围必须查当前方法及目标型号协议。 &
\passthrough{\lstinline!value = obj.x!} /
\passthrough{\lstinline!obj.x = value!} \\
\passthrough{\lstinline!y!} & \passthrough{\lstinline!float!} & Y
方向位置或分量参数。参考系、单位和范围必须查当前方法及目标型号协议。 &
\passthrough{\lstinline!value = obj.y!} /
\passthrough{\lstinline!obj.y = value!} \\
\passthrough{\lstinline!yaw!} & \passthrough{\lstinline!float!} &
偏航角或偏航目标参数。参考系、单位和范围必须查目标型号协议。 &
\passthrough{\lstinline!value = obj.yaw!} /
\passthrough{\lstinline!obj.yaw = value!} \\
\passthrough{\lstinline!relative!} & \passthrough{\lstinline!bool!} &
传给该接口的 \passthrough{\lstinline!relative!} 参数，Python 类型为
\passthrough{\lstinline!bool!}。精确含义、单位、范围和枚举值需查目标型号对应头文件与协议。
& \passthrough{\lstinline!value = obj.relative!} /
\passthrough{\lstinline!obj.relative = value!} \\
\bottomrule
\end{longtable}

\textbf{方法索引}

\begin{longtable}[]{@{}
  >{\raggedright\arraybackslash}p{(\columnwidth - 6\tabcolsep) * \real{0.18}}
  >{\raggedright\arraybackslash}p{(\columnwidth - 6\tabcolsep) * \real{0.24}}
  >{\raggedright\arraybackslash}p{(\columnwidth - 6\tabcolsep) * \real{0.18}}
  >{\raggedright\arraybackslash}p{(\columnwidth - 6\tabcolsep) * \real{0.40}}@{}}
\toprule
方法 & Python 签名 & 状态 & 安全分类 \\
\midrule
\endhead
\protect\hyperlink{unitree-sdk2-cpp-robot-h1-jsonizetargetpos-dunder-init-1}{\passthrough{\lstinline!\_\_init\_\_!}}
& \passthrough{\lstinline!def \_\_init\_\_(self) -> None!} &
\passthrough{\lstinline!SIGNATURE\_ONLY!} &
\passthrough{\lstinline!CONSTRUCTION!} \\
\protect\hyperlink{unitree-sdk2-cpp-robot-h1-jsonizetargetpos-from-json-1}{\passthrough{\lstinline!from\_json!}}
&
\passthrough{\lstinline!def from\_json(self, json: dict[str, Any]) -> None!}
& \passthrough{\lstinline!SIGNATURE\_ONLY!} &
\passthrough{\lstinline!UNCLASSIFIED!} \\
\protect\hyperlink{unitree-sdk2-cpp-robot-h1-jsonizetargetpos-to-json-1}{\passthrough{\lstinline!to\_json!}}
&
\passthrough{\lstinline!def to\_json(self, json: dict[str, Any]) -> None!}
& \passthrough{\lstinline!SIGNATURE\_ONLY!} &
\passthrough{\lstinline!UNCLASSIFIED!} \\
\bottomrule
\end{longtable}

\hypertarget{unitree-sdk2-cpp-robot-h1-jsonizetargetpos-dunder-init-1}{%
\paragraph{\texorpdfstring{\texttt{unitree\_sdk2\_cpp.robot.h1.JsonizeTargetPos.\_\_init\_\_}}{unitree\_sdk2\_cpp.robot.h1.JsonizeTargetPos.\_\_init\_\_}}\label{unitree-sdk2-cpp-robot-h1-jsonizetargetpos-dunder-init-1}}

初始化 \passthrough{\lstinline!JsonizeTargetPos!}
实例。是否能实际构造取决于下方可用性状态。

\textbf{签名}

\begin{lstlisting}[language=Python]
def __init__(self) -> None
\end{lstlisting}

\textbf{可用性与安全性}

\textbf{\passthrough{\lstinline!SIGNATURE\_ONLY!} /
\passthrough{\lstinline!CONSTRUCTION!}}：当前只有设计期签名，不能假设已存在于运行时扩展。

\textbf{参数}

无显式参数。实例方法中的 \passthrough{\lstinline!self!} 由 Python
自动传入。

\textbf{返回值}

\begin{longtable}[]{@{}
  >{\raggedright\arraybackslash}p{(\columnwidth - 4\tabcolsep) * \real{0.18}}
  >{\raggedright\arraybackslash}p{(\columnwidth - 4\tabcolsep) * \real{0.20}}
  >{\raggedright\arraybackslash}p{(\columnwidth - 4\tabcolsep) * \real{0.62}}@{}}
\toprule
位置 & 类型 & 含义 \\
\midrule
\endhead
无 & \passthrough{\lstinline!None!} &
\passthrough{\lstinline!\_\_init\_\_!}
本身不返回值；调用类对象时，在构造函数可用的前提下得到该类实例。 \\
\bottomrule
\end{longtable}

\textbf{对应 C++}

\begin{itemize}
\tightlist
\item
  类：\passthrough{\lstinline!unitree::robot::h1::JsonizeTargetPos!}
\item
  签名：\passthrough{\lstinline!JsonizeTargetPos()!}
\item
  绑定策略：\passthrough{\lstinline!未单独标注!}
\item
  声明位置：\passthrough{\lstinline!include/unitree/robot/h1/loco/h1\_loco\_api.hpp:72!}
\end{itemize}

\textbf{说明}

\begin{itemize}
\tightlist
\item
  示例是局部调用片段；\passthrough{\lstinline!obj!}
  和参数变量表示已经按上层业务校验并准备好的对象。
\item
  该条目可以被编辑器和类型检查器识别，但当前运行时可能在属性查找阶段直接失败。
\end{itemize}

\textbf{用法}

\begin{lstlisting}[language=Python]
from typing import TYPE_CHECKING

if TYPE_CHECKING:
    def planned_construct() -> JsonizeTargetPos:
        return JsonizeTargetPos()
\end{lstlisting}

\begin{quote}
{[}!CAUTION{]} 上面的 \passthrough{\lstinline!TYPE\_CHECKING!}
示例只表达计划调用形态。该分支在普通 Python
运行时为假，不代表当前扩展能执行此接口。
\end{quote}

\hypertarget{unitree-sdk2-cpp-robot-h1-jsonizetargetpos-from-json-1}{%
\paragraph{\texorpdfstring{\texttt{unitree\_sdk2\_cpp.robot.h1.JsonizeTargetPos.from\_json}}{unitree\_sdk2\_cpp.robot.h1.JsonizeTargetPos.from\_json}}\label{unitree-sdk2-cpp-robot-h1-jsonizetargetpos-from-json-1}}

计划从 JSON 风格字典读取字段并更新当前 SDK 值对象。

\textbf{签名}

\begin{lstlisting}[language=Python]
def from_json(self, json: dict[str, Any]) -> None
\end{lstlisting}

\textbf{可用性与安全性}

\textbf{\passthrough{\lstinline!SIGNATURE\_ONLY!} /
\passthrough{\lstinline!UNCLASSIFIED!}}：当前只有设计期签名，不能假设已存在于运行时扩展。

\textbf{参数}

\begin{longtable}[]{@{}
  >{\raggedright\arraybackslash}p{(\columnwidth - 6\tabcolsep) * \real{0.18}}
  >{\raggedright\arraybackslash}p{(\columnwidth - 6\tabcolsep) * \real{0.24}}
  >{\raggedright\arraybackslash}p{(\columnwidth - 6\tabcolsep) * \real{0.18}}
  >{\raggedright\arraybackslash}p{(\columnwidth - 6\tabcolsep) * \real{0.40}}@{}}
\toprule
名称 & Python 类型 & 默认值 & 含义 \\
\midrule
\endhead
\passthrough{\lstinline!json!} &
\passthrough{\lstinline!dict[str, Any]!} & 必填 & JSON
风格字典，键和值必须满足对应 C++ \passthrough{\lstinline!JsonMap!}
数据结构的约束。 对应 C++ 参数
\passthrough{\lstinline!json: common::JsonMap \&!}。 \\
\bottomrule
\end{longtable}

\textbf{返回值}

\begin{longtable}[]{@{}
  >{\raggedright\arraybackslash}p{(\columnwidth - 4\tabcolsep) * \real{0.18}}
  >{\raggedright\arraybackslash}p{(\columnwidth - 4\tabcolsep) * \real{0.20}}
  >{\raggedright\arraybackslash}p{(\columnwidth - 4\tabcolsep) * \real{0.62}}@{}}
\toprule
位置 & 类型 & 含义 \\
\midrule
\endhead
无 & \passthrough{\lstinline!None!} & 该方法不返回 Python
值；失败可能表现为异常或底层状态变化。 \\
\bottomrule
\end{longtable}

\textbf{对应 C++}

\begin{itemize}
\tightlist
\item
  类：\passthrough{\lstinline!unitree::robot::h1::JsonizeTargetPos!}
\item
  签名：\passthrough{\lstinline!fromJson(common::JsonMap \&)!}
\item
  绑定策略：\passthrough{\lstinline!SIGNATURE\_PREVIEW!}
\item
  声明位置：\passthrough{\lstinline!include/unitree/robot/h1/loco/h1\_loco\_api.hpp:75!}
\end{itemize}

\textbf{说明}

\begin{itemize}
\tightlist
\item
  示例是局部调用片段；\passthrough{\lstinline!obj!}
  和参数变量表示已经按上层业务校验并准备好的对象。
\item
  该条目可以被编辑器和类型检查器识别，但当前运行时可能在属性查找阶段直接失败。
\end{itemize}

\textbf{用法}

\begin{lstlisting}[language=Python]
from typing import TYPE_CHECKING

if TYPE_CHECKING:
    def planned_from_json(obj: JsonizeTargetPos, json: dict[str, Any]) -> None:
        obj.from_json(json=json)
\end{lstlisting}

\begin{quote}
{[}!CAUTION{]} 上面的 \passthrough{\lstinline!TYPE\_CHECKING!}
示例只表达计划调用形态。该分支在普通 Python
运行时为假，不代表当前扩展能执行此接口。
\end{quote}

\hypertarget{unitree-sdk2-cpp-robot-h1-jsonizetargetpos-to-json-1}{%
\paragraph{\texorpdfstring{\texttt{unitree\_sdk2\_cpp.robot.h1.JsonizeTargetPos.to\_json}}{unitree\_sdk2\_cpp.robot.h1.JsonizeTargetPos.to\_json}}\label{unitree-sdk2-cpp-robot-h1-jsonizetargetpos-to-json-1}}

计划把当前 SDK 值对象写入 JSON 风格字典。

\textbf{签名}

\begin{lstlisting}[language=Python]
def to_json(self, json: dict[str, Any]) -> None
\end{lstlisting}

\textbf{可用性与安全性}

\textbf{\passthrough{\lstinline!SIGNATURE\_ONLY!} /
\passthrough{\lstinline!UNCLASSIFIED!}}：当前只有设计期签名，不能假设已存在于运行时扩展。

\textbf{参数}

\begin{longtable}[]{@{}
  >{\raggedright\arraybackslash}p{(\columnwidth - 6\tabcolsep) * \real{0.18}}
  >{\raggedright\arraybackslash}p{(\columnwidth - 6\tabcolsep) * \real{0.24}}
  >{\raggedright\arraybackslash}p{(\columnwidth - 6\tabcolsep) * \real{0.18}}
  >{\raggedright\arraybackslash}p{(\columnwidth - 6\tabcolsep) * \real{0.40}}@{}}
\toprule
名称 & Python 类型 & 默认值 & 含义 \\
\midrule
\endhead
\passthrough{\lstinline!json!} &
\passthrough{\lstinline!dict[str, Any]!} & 必填 & JSON
风格字典，键和值必须满足对应 C++ \passthrough{\lstinline!JsonMap!}
数据结构的约束。 对应 C++ 参数
\passthrough{\lstinline!json: common::JsonMap \&!}。 \\
\bottomrule
\end{longtable}

\textbf{返回值}

\begin{longtable}[]{@{}
  >{\raggedright\arraybackslash}p{(\columnwidth - 4\tabcolsep) * \real{0.18}}
  >{\raggedright\arraybackslash}p{(\columnwidth - 4\tabcolsep) * \real{0.20}}
  >{\raggedright\arraybackslash}p{(\columnwidth - 4\tabcolsep) * \real{0.62}}@{}}
\toprule
位置 & 类型 & 含义 \\
\midrule
\endhead
无 & \passthrough{\lstinline!None!} & 该方法不返回 Python
值；失败可能表现为异常或底层状态变化。 \\
\bottomrule
\end{longtable}

\textbf{对应 C++}

\begin{itemize}
\tightlist
\item
  类：\passthrough{\lstinline!unitree::robot::h1::JsonizeTargetPos!}
\item
  签名：\passthrough{\lstinline!toJson(common::JsonMap \&) const!}
\item
  绑定策略：\passthrough{\lstinline!SIGNATURE\_PREVIEW!}
\item
  声明位置：\passthrough{\lstinline!include/unitree/robot/h1/loco/h1\_loco\_api.hpp:82!}
\end{itemize}

\textbf{说明}

\begin{itemize}
\tightlist
\item
  示例是局部调用片段；\passthrough{\lstinline!obj!}
  和参数变量表示已经按上层业务校验并准备好的对象。
\item
  该条目可以被编辑器和类型检查器识别，但当前运行时可能在属性查找阶段直接失败。
\end{itemize}

\textbf{用法}

\begin{lstlisting}[language=Python]
from typing import TYPE_CHECKING

if TYPE_CHECKING:
    def planned_to_json(obj: JsonizeTargetPos, json: dict[str, Any]) -> None:
        obj.to_json(json=json)
\end{lstlisting}

\begin{quote}
{[}!CAUTION{]} 上面的 \passthrough{\lstinline!TYPE\_CHECKING!}
示例只表达计划调用形态。该分支在普通 Python
运行时为假，不代表当前扩展能执行此接口。
\end{quote}

\hypertarget{unitree-sdk2-cpp-robot-h1-jsonizevelocitycommand}{%
\subsubsection{\texorpdfstring{\texttt{unitree\_sdk2\_cpp.robot.h1.JsonizeVelocityCommand}}{unitree\_sdk2\_cpp.robot.h1.JsonizeVelocityCommand}}\label{unitree-sdk2-cpp-robot-h1-jsonizevelocitycommand}}

SDK 类型签名预览。请逐项查看构造函数和方法的可用性。

\textbf{导入}

\begin{lstlisting}[language=Python]
from unitree_sdk2_cpp.robot.h1 import JsonizeVelocityCommand
\end{lstlisting}

\textbf{基类}：\passthrough{\lstinline!object!}

\textbf{公开属性}

\begin{longtable}[]{@{}
  >{\raggedright\arraybackslash}p{(\columnwidth - 6\tabcolsep) * \real{0.18}}
  >{\raggedright\arraybackslash}p{(\columnwidth - 6\tabcolsep) * \real{0.24}}
  >{\raggedright\arraybackslash}p{(\columnwidth - 6\tabcolsep) * \real{0.18}}
  >{\raggedright\arraybackslash}p{(\columnwidth - 6\tabcolsep) * \real{0.40}}@{}}
\toprule
名称 & Python 类型 & 含义 & 用法 \\
\midrule
\endhead
\passthrough{\lstinline!velocity!} &
\passthrough{\lstinline!list[float]!} & 传给该接口的 速度 参数，Python
类型为
\passthrough{\lstinline!list[float]!}。精确含义、单位、范围和枚举值需查目标型号对应头文件与协议。
& \passthrough{\lstinline!value = obj.velocity!} /
\passthrough{\lstinline!obj.velocity = value!} \\
\passthrough{\lstinline!duration!} & \passthrough{\lstinline!float!} &
操作持续时间参数。精确单位、范围和默认行为以目标型号接口定义为准。 &
\passthrough{\lstinline!value = obj.duration!} /
\passthrough{\lstinline!obj.duration = value!} \\
\bottomrule
\end{longtable}

\textbf{方法索引}

\begin{longtable}[]{@{}
  >{\raggedright\arraybackslash}p{(\columnwidth - 6\tabcolsep) * \real{0.18}}
  >{\raggedright\arraybackslash}p{(\columnwidth - 6\tabcolsep) * \real{0.24}}
  >{\raggedright\arraybackslash}p{(\columnwidth - 6\tabcolsep) * \real{0.18}}
  >{\raggedright\arraybackslash}p{(\columnwidth - 6\tabcolsep) * \real{0.40}}@{}}
\toprule
方法 & Python 签名 & 状态 & 安全分类 \\
\midrule
\endhead
\protect\hyperlink{unitree-sdk2-cpp-robot-h1-jsonizevelocitycommand-dunder-init-1}{\passthrough{\lstinline!\_\_init\_\_!}}
& \passthrough{\lstinline!def \_\_init\_\_(self) -> None!} &
\passthrough{\lstinline!SIGNATURE\_ONLY!} &
\passthrough{\lstinline!CONSTRUCTION!} \\
\protect\hyperlink{unitree-sdk2-cpp-robot-h1-jsonizevelocitycommand-from-json-1}{\passthrough{\lstinline!from\_json!}}
&
\passthrough{\lstinline!def from\_json(self, json: dict[str, Any]) -> None!}
& \passthrough{\lstinline!SIGNATURE\_ONLY!} &
\passthrough{\lstinline!UNCLASSIFIED!} \\
\protect\hyperlink{unitree-sdk2-cpp-robot-h1-jsonizevelocitycommand-to-json-1}{\passthrough{\lstinline!to\_json!}}
&
\passthrough{\lstinline!def to\_json(self, json: dict[str, Any]) -> None!}
& \passthrough{\lstinline!SIGNATURE\_ONLY!} &
\passthrough{\lstinline!UNCLASSIFIED!} \\
\bottomrule
\end{longtable}

\hypertarget{unitree-sdk2-cpp-robot-h1-jsonizevelocitycommand-dunder-init-1}{%
\paragraph{\texorpdfstring{\texttt{unitree\_sdk2\_cpp.robot.h1.JsonizeVelocityCommand.\_\_init\_\_}}{unitree\_sdk2\_cpp.robot.h1.JsonizeVelocityCommand.\_\_init\_\_}}\label{unitree-sdk2-cpp-robot-h1-jsonizevelocitycommand-dunder-init-1}}

初始化 \passthrough{\lstinline!JsonizeVelocityCommand!}
实例。是否能实际构造取决于下方可用性状态。

\textbf{签名}

\begin{lstlisting}[language=Python]
def __init__(self) -> None
\end{lstlisting}

\textbf{可用性与安全性}

\textbf{\passthrough{\lstinline!SIGNATURE\_ONLY!} /
\passthrough{\lstinline!CONSTRUCTION!}}：当前只有设计期签名，不能假设已存在于运行时扩展。

\textbf{参数}

无显式参数。实例方法中的 \passthrough{\lstinline!self!} 由 Python
自动传入。

\textbf{返回值}

\begin{longtable}[]{@{}
  >{\raggedright\arraybackslash}p{(\columnwidth - 4\tabcolsep) * \real{0.18}}
  >{\raggedright\arraybackslash}p{(\columnwidth - 4\tabcolsep) * \real{0.20}}
  >{\raggedright\arraybackslash}p{(\columnwidth - 4\tabcolsep) * \real{0.62}}@{}}
\toprule
位置 & 类型 & 含义 \\
\midrule
\endhead
无 & \passthrough{\lstinline!None!} &
\passthrough{\lstinline!\_\_init\_\_!}
本身不返回值；调用类对象时，在构造函数可用的前提下得到该类实例。 \\
\bottomrule
\end{longtable}

\textbf{对应 C++}

\begin{itemize}
\tightlist
\item
  类：\passthrough{\lstinline!unitree::robot::h1::JsonizeVelocityCommand!}
\item
  签名：\passthrough{\lstinline!JsonizeVelocityCommand()!}
\item
  绑定策略：\passthrough{\lstinline!未单独标注!}
\item
  声明位置：\passthrough{\lstinline!include/unitree/robot/h1/loco/h1\_loco\_api.hpp:53!}
\end{itemize}

\textbf{说明}

\begin{itemize}
\tightlist
\item
  示例是局部调用片段；\passthrough{\lstinline!obj!}
  和参数变量表示已经按上层业务校验并准备好的对象。
\item
  该条目可以被编辑器和类型检查器识别，但当前运行时可能在属性查找阶段直接失败。
\end{itemize}

\textbf{用法}

\begin{lstlisting}[language=Python]
from typing import TYPE_CHECKING

if TYPE_CHECKING:
    def planned_construct() -> JsonizeVelocityCommand:
        return JsonizeVelocityCommand()
\end{lstlisting}

\begin{quote}
{[}!CAUTION{]} 上面的 \passthrough{\lstinline!TYPE\_CHECKING!}
示例只表达计划调用形态。该分支在普通 Python
运行时为假，不代表当前扩展能执行此接口。
\end{quote}

\hypertarget{unitree-sdk2-cpp-robot-h1-jsonizevelocitycommand-from-json-1}{%
\paragraph{\texorpdfstring{\texttt{unitree\_sdk2\_cpp.robot.h1.JsonizeVelocityCommand.from\_json}}{unitree\_sdk2\_cpp.robot.h1.JsonizeVelocityCommand.from\_json}}\label{unitree-sdk2-cpp-robot-h1-jsonizevelocitycommand-from-json-1}}

计划从 JSON 风格字典读取字段并更新当前 SDK 值对象。

\textbf{签名}

\begin{lstlisting}[language=Python]
def from_json(self, json: dict[str, Any]) -> None
\end{lstlisting}

\textbf{可用性与安全性}

\textbf{\passthrough{\lstinline!SIGNATURE\_ONLY!} /
\passthrough{\lstinline!UNCLASSIFIED!}}：当前只有设计期签名，不能假设已存在于运行时扩展。

\textbf{参数}

\begin{longtable}[]{@{}
  >{\raggedright\arraybackslash}p{(\columnwidth - 6\tabcolsep) * \real{0.18}}
  >{\raggedright\arraybackslash}p{(\columnwidth - 6\tabcolsep) * \real{0.24}}
  >{\raggedright\arraybackslash}p{(\columnwidth - 6\tabcolsep) * \real{0.18}}
  >{\raggedright\arraybackslash}p{(\columnwidth - 6\tabcolsep) * \real{0.40}}@{}}
\toprule
名称 & Python 类型 & 默认值 & 含义 \\
\midrule
\endhead
\passthrough{\lstinline!json!} &
\passthrough{\lstinline!dict[str, Any]!} & 必填 & JSON
风格字典，键和值必须满足对应 C++ \passthrough{\lstinline!JsonMap!}
数据结构的约束。 对应 C++ 参数
\passthrough{\lstinline!json: common::JsonMap \&!}。 \\
\bottomrule
\end{longtable}

\textbf{返回值}

\begin{longtable}[]{@{}
  >{\raggedright\arraybackslash}p{(\columnwidth - 4\tabcolsep) * \real{0.18}}
  >{\raggedright\arraybackslash}p{(\columnwidth - 4\tabcolsep) * \real{0.20}}
  >{\raggedright\arraybackslash}p{(\columnwidth - 4\tabcolsep) * \real{0.62}}@{}}
\toprule
位置 & 类型 & 含义 \\
\midrule
\endhead
无 & \passthrough{\lstinline!None!} & 该方法不返回 Python
值；失败可能表现为异常或底层状态变化。 \\
\bottomrule
\end{longtable}

\textbf{对应 C++}

\begin{itemize}
\tightlist
\item
  类：\passthrough{\lstinline!unitree::robot::h1::JsonizeVelocityCommand!}
\item
  签名：\passthrough{\lstinline!fromJson(common::JsonMap \&)!}
\item
  绑定策略：\passthrough{\lstinline!SIGNATURE\_PREVIEW!}
\item
  声明位置：\passthrough{\lstinline!include/unitree/robot/h1/loco/h1\_loco\_api.hpp:56!}
\end{itemize}

\textbf{说明}

\begin{itemize}
\tightlist
\item
  示例是局部调用片段；\passthrough{\lstinline!obj!}
  和参数变量表示已经按上层业务校验并准备好的对象。
\item
  该条目可以被编辑器和类型检查器识别，但当前运行时可能在属性查找阶段直接失败。
\end{itemize}

\textbf{用法}

\begin{lstlisting}[language=Python]
from typing import TYPE_CHECKING

if TYPE_CHECKING:
    def planned_from_json(obj: JsonizeVelocityCommand, json: dict[str, Any]) -> None:
        obj.from_json(json=json)
\end{lstlisting}

\begin{quote}
{[}!CAUTION{]} 上面的 \passthrough{\lstinline!TYPE\_CHECKING!}
示例只表达计划调用形态。该分支在普通 Python
运行时为假，不代表当前扩展能执行此接口。
\end{quote}

\hypertarget{unitree-sdk2-cpp-robot-h1-jsonizevelocitycommand-to-json-1}{%
\paragraph{\texorpdfstring{\texttt{unitree\_sdk2\_cpp.robot.h1.JsonizeVelocityCommand.to\_json}}{unitree\_sdk2\_cpp.robot.h1.JsonizeVelocityCommand.to\_json}}\label{unitree-sdk2-cpp-robot-h1-jsonizevelocitycommand-to-json-1}}

计划把当前 SDK 值对象写入 JSON 风格字典。

\textbf{签名}

\begin{lstlisting}[language=Python]
def to_json(self, json: dict[str, Any]) -> None
\end{lstlisting}

\textbf{可用性与安全性}

\textbf{\passthrough{\lstinline!SIGNATURE\_ONLY!} /
\passthrough{\lstinline!UNCLASSIFIED!}}：当前只有设计期签名，不能假设已存在于运行时扩展。

\textbf{参数}

\begin{longtable}[]{@{}
  >{\raggedright\arraybackslash}p{(\columnwidth - 6\tabcolsep) * \real{0.18}}
  >{\raggedright\arraybackslash}p{(\columnwidth - 6\tabcolsep) * \real{0.24}}
  >{\raggedright\arraybackslash}p{(\columnwidth - 6\tabcolsep) * \real{0.18}}
  >{\raggedright\arraybackslash}p{(\columnwidth - 6\tabcolsep) * \real{0.40}}@{}}
\toprule
名称 & Python 类型 & 默认值 & 含义 \\
\midrule
\endhead
\passthrough{\lstinline!json!} &
\passthrough{\lstinline!dict[str, Any]!} & 必填 & JSON
风格字典，键和值必须满足对应 C++ \passthrough{\lstinline!JsonMap!}
数据结构的约束。 对应 C++ 参数
\passthrough{\lstinline!json: common::JsonMap \&!}。 \\
\bottomrule
\end{longtable}

\textbf{返回值}

\begin{longtable}[]{@{}
  >{\raggedright\arraybackslash}p{(\columnwidth - 4\tabcolsep) * \real{0.18}}
  >{\raggedright\arraybackslash}p{(\columnwidth - 4\tabcolsep) * \real{0.20}}
  >{\raggedright\arraybackslash}p{(\columnwidth - 4\tabcolsep) * \real{0.62}}@{}}
\toprule
位置 & 类型 & 含义 \\
\midrule
\endhead
无 & \passthrough{\lstinline!None!} & 该方法不返回 Python
值；失败可能表现为异常或底层状态变化。 \\
\bottomrule
\end{longtable}

\textbf{对应 C++}

\begin{itemize}
\tightlist
\item
  类：\passthrough{\lstinline!unitree::robot::h1::JsonizeVelocityCommand!}
\item
  签名：\passthrough{\lstinline!toJson(common::JsonMap \&) const!}
\item
  绑定策略：\passthrough{\lstinline!SIGNATURE\_PREVIEW!}
\item
  声明位置：\passthrough{\lstinline!include/unitree/robot/h1/loco/h1\_loco\_api.hpp:61!}
\end{itemize}

\textbf{说明}

\begin{itemize}
\tightlist
\item
  示例是局部调用片段；\passthrough{\lstinline!obj!}
  和参数变量表示已经按上层业务校验并准备好的对象。
\item
  该条目可以被编辑器和类型检查器识别，但当前运行时可能在属性查找阶段直接失败。
\end{itemize}

\textbf{用法}

\begin{lstlisting}[language=Python]
from typing import TYPE_CHECKING

if TYPE_CHECKING:
    def planned_to_json(obj: JsonizeVelocityCommand, json: dict[str, Any]) -> None:
        obj.to_json(json=json)
\end{lstlisting}

\begin{quote}
{[}!CAUTION{]} 上面的 \passthrough{\lstinline!TYPE\_CHECKING!}
示例只表达计划调用形态。该分支在普通 Python
运行时为假，不代表当前扩展能执行此接口。
\end{quote}

\hypertarget{unitree-sdk2-cpp-robot-h1-lococlient}{%
\subsubsection{\texorpdfstring{\texttt{unitree\_sdk2\_cpp.robot.h1.LocoClient}}{unitree\_sdk2\_cpp.robot.h1.LocoClient}}\label{unitree-sdk2-cpp-robot-h1-lococlient}}

Robot 服务客户端。构造、初始化和具体方法的可用性必须分别检查。

\textbf{导入}

\begin{lstlisting}[language=Python]
from unitree_sdk2_cpp.robot.h1 import LocoClient
\end{lstlisting}

\textbf{基类}：\passthrough{\lstinline!Client!}

\textbf{方法索引}

\begin{longtable}[]{@{}
  >{\raggedright\arraybackslash}p{(\columnwidth - 6\tabcolsep) * \real{0.18}}
  >{\raggedright\arraybackslash}p{(\columnwidth - 6\tabcolsep) * \real{0.24}}
  >{\raggedright\arraybackslash}p{(\columnwidth - 6\tabcolsep) * \real{0.18}}
  >{\raggedright\arraybackslash}p{(\columnwidth - 6\tabcolsep) * \real{0.40}}@{}}
\toprule
方法 & Python 签名 & 状态 & 安全分类 \\
\midrule
\endhead
\protect\hyperlink{unitree-sdk2-cpp-robot-h1-lococlient-dunder-init-1}{\passthrough{\lstinline!\_\_init\_\_!}}
& \passthrough{\lstinline!def \_\_init\_\_(self) -> None!} &
\passthrough{\lstinline!AVAILABLE!} &
\passthrough{\lstinline!CONSTRUCTION!} \\
\protect\hyperlink{unitree-sdk2-cpp-robot-h1-lococlient-init-1}{\passthrough{\lstinline!init!}}
& \passthrough{\lstinline!def init(self) -> None!} &
\passthrough{\lstinline!AVAILABLE!} &
\passthrough{\lstinline!INITIALIZATION!} \\
\protect\hyperlink{unitree-sdk2-cpp-robot-h1-lococlient-get-fsm-id-1}{\passthrough{\lstinline!get\_fsm\_id!}}
& \passthrough{\lstinline!def get\_fsm\_id(self) -> tuple[int, int]!} &
\passthrough{\lstinline!AVAILABLE!} &
\passthrough{\lstinline!READ\_ONLY!} \\
\protect\hyperlink{unitree-sdk2-cpp-robot-h1-lococlient-get-fsm-mode-1}{\passthrough{\lstinline!get\_fsm\_mode!}}
& \passthrough{\lstinline!def get\_fsm\_mode(self) -> tuple[int, int]!}
& \passthrough{\lstinline!AVAILABLE!} &
\passthrough{\lstinline!READ\_ONLY!} \\
\protect\hyperlink{unitree-sdk2-cpp-robot-h1-lococlient-get-balance-mode-1}{\passthrough{\lstinline!get\_balance\_mode!}}
&
\passthrough{\lstinline!def get\_balance\_mode(self) -> tuple[int, int]!}
& \passthrough{\lstinline!AVAILABLE!} &
\passthrough{\lstinline!READ\_ONLY!} \\
\protect\hyperlink{unitree-sdk2-cpp-robot-h1-lococlient-get-swing-height-1}{\passthrough{\lstinline!get\_swing\_height!}}
&
\passthrough{\lstinline!def get\_swing\_height(self) -> tuple[int, float]!}
& \passthrough{\lstinline!AVAILABLE!} &
\passthrough{\lstinline!READ\_ONLY!} \\
\protect\hyperlink{unitree-sdk2-cpp-robot-h1-lococlient-get-stand-height-1}{\passthrough{\lstinline!get\_stand\_height!}}
&
\passthrough{\lstinline!def get\_stand\_height(self) -> tuple[int, float]!}
& \passthrough{\lstinline!AVAILABLE!} &
\passthrough{\lstinline!READ\_ONLY!} \\
\protect\hyperlink{unitree-sdk2-cpp-robot-h1-lococlient-get-phase-1}{\passthrough{\lstinline!get\_phase!}}
&
\passthrough{\lstinline!def get\_phase(self) -> tuple[int, list[float]]!}
& \passthrough{\lstinline!AVAILABLE!} &
\passthrough{\lstinline!READ\_ONLY!} \\
\protect\hyperlink{unitree-sdk2-cpp-robot-h1-lococlient-enable-odom-1}{\passthrough{\lstinline!enable\_odom!}}
& \passthrough{\lstinline!def enable\_odom(self) -> int!} &
\passthrough{\lstinline!SIGNATURE\_ONLY!} &
\passthrough{\lstinline!MOTION\_COMMAND!} \\
\protect\hyperlink{unitree-sdk2-cpp-robot-h1-lococlient-disable-odom-1}{\passthrough{\lstinline!disable\_odom!}}
& \passthrough{\lstinline!def disable\_odom(self) -> int!} &
\passthrough{\lstinline!SIGNATURE\_ONLY!} &
\passthrough{\lstinline!MOTION\_COMMAND!} \\
\protect\hyperlink{unitree-sdk2-cpp-robot-h1-lococlient-get-odom-1}{\passthrough{\lstinline!get\_odom!}}
&
\passthrough{\lstinline!def get\_odom(self) -> tuple[int, float, float, float]!}
& \passthrough{\lstinline!AVAILABLE!} &
\passthrough{\lstinline!READ\_ONLY!} \\
\protect\hyperlink{unitree-sdk2-cpp-robot-h1-lococlient-set-fsm-id-1}{\passthrough{\lstinline!set\_fsm\_id!}}
& \passthrough{\lstinline!def set\_fsm\_id(self, fsm\_id: int) -> int!}
& \passthrough{\lstinline!SIGNATURE\_ONLY!} &
\passthrough{\lstinline!MOTION\_COMMAND!} \\
\protect\hyperlink{unitree-sdk2-cpp-robot-h1-lococlient-set-balance-mode-1}{\passthrough{\lstinline!set\_balance\_mode!}}
&
\passthrough{\lstinline!def set\_balance\_mode(self, balance\_mode: int) -> int!}
& \passthrough{\lstinline!SIGNATURE\_ONLY!} &
\passthrough{\lstinline!MOTION\_COMMAND!} \\
\protect\hyperlink{unitree-sdk2-cpp-robot-h1-lococlient-set-swing-height-1}{\passthrough{\lstinline!set\_swing\_height!}}
&
\passthrough{\lstinline!def set\_swing\_height(self, swing\_height: float) -> int!}
& \passthrough{\lstinline!SIGNATURE\_ONLY!} &
\passthrough{\lstinline!MOTION\_COMMAND!} \\
\protect\hyperlink{unitree-sdk2-cpp-robot-h1-lococlient-set-stand-height-1}{\passthrough{\lstinline!set\_stand\_height!}}
&
\passthrough{\lstinline!def set\_stand\_height(self, stand\_height: float) -> int!}
& \passthrough{\lstinline!SIGNATURE\_ONLY!} &
\passthrough{\lstinline!MOTION\_COMMAND!} \\
\protect\hyperlink{unitree-sdk2-cpp-robot-h1-lococlient-set-velocity-1}{\passthrough{\lstinline!set\_velocity!}}
&
\passthrough{\lstinline!def set\_velocity(self, vx: float, vy: float, omega: float, duration: float = ...) -> int!}
& \passthrough{\lstinline!SIGNATURE\_ONLY!} &
\passthrough{\lstinline!MOTION\_COMMAND!} \\
\protect\hyperlink{unitree-sdk2-cpp-robot-h1-lococlient-set-phase-1}{\passthrough{\lstinline!set\_phase!}}
&
\passthrough{\lstinline!def set\_phase(self, phase: Sequence[float]) -> int!}
& \passthrough{\lstinline!SIGNATURE\_ONLY!} &
\passthrough{\lstinline!MOTION\_COMMAND!} \\
\protect\hyperlink{unitree-sdk2-cpp-robot-h1-lococlient-set-target-pos-1}{\passthrough{\lstinline!set\_target\_pos!}}
&
\passthrough{\lstinline!def set\_target\_pos(self, x: float, y: float, yaw: float, relative: bool = ...) -> int!}
& \passthrough{\lstinline!SIGNATURE\_ONLY!} &
\passthrough{\lstinline!MOTION\_COMMAND!} \\
\protect\hyperlink{unitree-sdk2-cpp-robot-h1-lococlient-set-task-id-1}{\passthrough{\lstinline!set\_task\_id!}}
&
\passthrough{\lstinline!def set\_task\_id(self, task\_id: int) -> int!}
& \passthrough{\lstinline!SIGNATURE\_ONLY!} &
\passthrough{\lstinline!MOTION\_COMMAND!} \\
\protect\hyperlink{unitree-sdk2-cpp-robot-h1-lococlient-damp-1}{\passthrough{\lstinline!damp!}}
& \passthrough{\lstinline!def damp(self) -> int!} &
\passthrough{\lstinline!SIGNATURE\_ONLY!} &
\passthrough{\lstinline!MOTION\_COMMAND!} \\
\protect\hyperlink{unitree-sdk2-cpp-robot-h1-lococlient-start-1}{\passthrough{\lstinline!start!}}
& \passthrough{\lstinline!def start(self) -> int!} &
\passthrough{\lstinline!SIGNATURE\_ONLY!} &
\passthrough{\lstinline!MOTION\_COMMAND!} \\
\protect\hyperlink{unitree-sdk2-cpp-robot-h1-lococlient-stand-up-1}{\passthrough{\lstinline!stand\_up!}}
& \passthrough{\lstinline!def stand\_up(self) -> int!} &
\passthrough{\lstinline!SIGNATURE\_ONLY!} &
\passthrough{\lstinline!MOTION\_COMMAND!} \\
\protect\hyperlink{unitree-sdk2-cpp-robot-h1-lococlient-zero-torque-1}{\passthrough{\lstinline!zero\_torque!}}
& \passthrough{\lstinline!def zero\_torque(self) -> int!} &
\passthrough{\lstinline!SIGNATURE\_ONLY!} &
\passthrough{\lstinline!MOTION\_COMMAND!} \\
\protect\hyperlink{unitree-sdk2-cpp-robot-h1-lococlient-stop-move-1}{\passthrough{\lstinline!stop\_move!}}
& \passthrough{\lstinline!def stop\_move(self) -> int!} &
\passthrough{\lstinline!SIGNATURE\_ONLY!} &
\passthrough{\lstinline!MOTION\_COMMAND!} \\
\protect\hyperlink{unitree-sdk2-cpp-robot-h1-lococlient-high-stand-1}{\passthrough{\lstinline!high\_stand!}}
& \passthrough{\lstinline!def high\_stand(self) -> int!} &
\passthrough{\lstinline!SIGNATURE\_ONLY!} &
\passthrough{\lstinline!MOTION\_COMMAND!} \\
\protect\hyperlink{unitree-sdk2-cpp-robot-h1-lococlient-low-stand-1}{\passthrough{\lstinline!low\_stand!}}
& \passthrough{\lstinline!def low\_stand(self) -> int!} &
\passthrough{\lstinline!SIGNATURE\_ONLY!} &
\passthrough{\lstinline!MOTION\_COMMAND!} \\
\protect\hyperlink{unitree-sdk2-cpp-robot-h1-lococlient-move-1}{\passthrough{\lstinline!move（重载 1/2）!}}
&
\passthrough{\lstinline!def move(self, vx: float, vy: float, vyaw: float, continous\_move: bool) -> int!}
& \passthrough{\lstinline!SIGNATURE\_ONLY!} &
\passthrough{\lstinline!MOTION\_COMMAND!} \\
\protect\hyperlink{unitree-sdk2-cpp-robot-h1-lococlient-move-2}{\passthrough{\lstinline!move（重载 2/2）!}}
&
\passthrough{\lstinline!def move(self, vx: float, vy: float, vyaw: float) -> int!}
& \passthrough{\lstinline!SIGNATURE\_ONLY!} &
\passthrough{\lstinline!MOTION\_COMMAND!} \\
\protect\hyperlink{unitree-sdk2-cpp-robot-h1-lococlient-balance-stand-1}{\passthrough{\lstinline!balance\_stand!}}
& \passthrough{\lstinline!def balance\_stand(self) -> int!} &
\passthrough{\lstinline!SIGNATURE\_ONLY!} &
\passthrough{\lstinline!MOTION\_COMMAND!} \\
\protect\hyperlink{unitree-sdk2-cpp-robot-h1-lococlient-continuous-gait-1}{\passthrough{\lstinline!continuous\_gait!}}
&
\passthrough{\lstinline!def continuous\_gait(self, flag: bool) -> int!}
& \passthrough{\lstinline!SIGNATURE\_ONLY!} &
\passthrough{\lstinline!MOTION\_COMMAND!} \\
\protect\hyperlink{unitree-sdk2-cpp-robot-h1-lococlient-switch-move-mode-1}{\passthrough{\lstinline!switch\_move\_mode!}}
&
\passthrough{\lstinline!def switch\_move\_mode(self, flag: bool) -> int!}
& \passthrough{\lstinline!SIGNATURE\_ONLY!} &
\passthrough{\lstinline!MOTION\_COMMAND!} \\
\protect\hyperlink{unitree-sdk2-cpp-robot-h1-lococlient-set-next-foot-1}{\passthrough{\lstinline!set\_next\_foot!}}
& \passthrough{\lstinline!def set\_next\_foot(self, foot: bool) -> int!}
& \passthrough{\lstinline!SIGNATURE\_ONLY!} &
\passthrough{\lstinline!MOTION\_COMMAND!} \\
\protect\hyperlink{unitree-sdk2-cpp-robot-h1-lococlient-wave-hand-1}{\passthrough{\lstinline!wave\_hand!}}
& \passthrough{\lstinline!def wave\_hand(self) -> int!} &
\passthrough{\lstinline!SIGNATURE\_ONLY!} &
\passthrough{\lstinline!MOTION\_COMMAND!} \\
\protect\hyperlink{unitree-sdk2-cpp-robot-h1-lococlient-shake-hand-1}{\passthrough{\lstinline!shake\_hand!}}
&
\passthrough{\lstinline!def shake\_hand(self, stage: int = ...) -> int!}
& \passthrough{\lstinline!SIGNATURE\_ONLY!} &
\passthrough{\lstinline!MOTION\_COMMAND!} \\
\bottomrule
\end{longtable}

\hypertarget{unitree-sdk2-cpp-robot-h1-lococlient-dunder-init-1}{%
\paragraph{\texorpdfstring{\texttt{unitree\_sdk2\_cpp.robot.h1.LocoClient.\_\_init\_\_}}{unitree\_sdk2\_cpp.robot.h1.LocoClient.\_\_init\_\_}}\label{unitree-sdk2-cpp-robot-h1-lococlient-dunder-init-1}}

初始化 \passthrough{\lstinline!LocoClient!}
实例。是否能实际构造取决于下方可用性状态。

\textbf{签名}

\begin{lstlisting}[language=Python]
def __init__(self) -> None
\end{lstlisting}

\textbf{可用性与安全性}

\textbf{\passthrough{\lstinline!AVAILABLE!} /
\passthrough{\lstinline!CONSTRUCTION!}}：当前绑定源码已实现；实际执行条件仍取决于平台和对应
SDK 环境。

\textbf{参数}

无显式参数。实例方法中的 \passthrough{\lstinline!self!} 由 Python
自动传入。

\textbf{返回值}

\begin{longtable}[]{@{}
  >{\raggedright\arraybackslash}p{(\columnwidth - 4\tabcolsep) * \real{0.18}}
  >{\raggedright\arraybackslash}p{(\columnwidth - 4\tabcolsep) * \real{0.20}}
  >{\raggedright\arraybackslash}p{(\columnwidth - 4\tabcolsep) * \real{0.62}}@{}}
\toprule
位置 & 类型 & 含义 \\
\midrule
\endhead
无 & \passthrough{\lstinline!None!} &
\passthrough{\lstinline!\_\_init\_\_!}
本身不返回值；调用类对象时，在构造函数可用的前提下得到该类实例。 \\
\bottomrule
\end{longtable}

\textbf{对应 C++}

\begin{itemize}
\tightlist
\item
  类：\passthrough{\lstinline!unitree::robot::h1::LocoClient!}
\item
  签名：\passthrough{\lstinline!LocoClient()!}
\item
  绑定策略：\passthrough{\lstinline!未单独标注!}
\item
  声明位置：\passthrough{\lstinline!include/unitree/robot/h1/loco/h1\_loco\_client.hpp:14!}
\end{itemize}

\textbf{说明}

\begin{itemize}
\tightlist
\item
  示例是局部调用片段；\passthrough{\lstinline!obj!}
  和参数变量表示已经按上层业务校验并准备好的对象。
\end{itemize}

\textbf{用法}

\begin{lstlisting}[language=Python]
value = LocoClient()
\end{lstlisting}

\hypertarget{unitree-sdk2-cpp-robot-h1-lococlient-init-1}{%
\paragraph{\texorpdfstring{\texttt{unitree\_sdk2\_cpp.robot.h1.LocoClient.init}}{unitree\_sdk2\_cpp.robot.h1.LocoClient.init}}\label{unitree-sdk2-cpp-robot-h1-lococlient-init-1}}

初始化当前 SDK 对象所需的底层通道或服务资源。

\textbf{签名}

\begin{lstlisting}[language=Python]
def init(self) -> None
\end{lstlisting}

\textbf{可用性与安全性}

\textbf{\passthrough{\lstinline!AVAILABLE!} /
\passthrough{\lstinline!INITIALIZATION!}}：当前绑定源码已实现；实际执行条件仍取决于平台和对应
SDK 环境。

\textbf{参数}

无显式参数。实例方法中的 \passthrough{\lstinline!self!} 由 Python
自动传入。

\textbf{返回值}

\begin{longtable}[]{@{}
  >{\raggedright\arraybackslash}p{(\columnwidth - 4\tabcolsep) * \real{0.18}}
  >{\raggedright\arraybackslash}p{(\columnwidth - 4\tabcolsep) * \real{0.20}}
  >{\raggedright\arraybackslash}p{(\columnwidth - 4\tabcolsep) * \real{0.62}}@{}}
\toprule
位置 & 类型 & 含义 \\
\midrule
\endhead
无 & \passthrough{\lstinline!None!} & 该方法不返回 Python
值；失败可能表现为异常或底层状态变化。 \\
\bottomrule
\end{longtable}

\textbf{对应 C++}

\begin{itemize}
\tightlist
\item
  类：\passthrough{\lstinline!unitree::robot::h1::LocoClient!}
\item
  签名：\passthrough{\lstinline!Init()!}
\item
  绑定策略：\passthrough{\lstinline!DIRECT!}
\item
  声明位置：\passthrough{\lstinline!include/unitree/robot/h1/loco/h1\_loco\_client.hpp:18!}
\end{itemize}

\textbf{说明}

\begin{itemize}
\tightlist
\item
  示例是局部调用片段；\passthrough{\lstinline!obj!}
  和参数变量表示已经按上层业务校验并准备好的对象。
\end{itemize}

\textbf{用法}

\begin{lstlisting}[language=Python]
obj.init()
\end{lstlisting}

\hypertarget{unitree-sdk2-cpp-robot-h1-lococlient-get-fsm-id-1}{%
\paragraph{\texorpdfstring{\texttt{unitree\_sdk2\_cpp.robot.h1.LocoClient.get\_fsm\_id}}{unitree\_sdk2\_cpp.robot.h1.LocoClient.get\_fsm\_id}}\label{unitree-sdk2-cpp-robot-h1-lococlient-get-fsm-id-1}}

查询或检查 FSM ID。

\textbf{签名}

\begin{lstlisting}[language=Python]
def get_fsm_id(self) -> tuple[int, int]
\end{lstlisting}

\textbf{可用性与安全性}

\textbf{\passthrough{\lstinline!AVAILABLE!} /
\passthrough{\lstinline!READ\_ONLY!}}：当前绑定源码已实现只读查询。Robot
Client 的服务端查询通常仍需匹配的 Linux
扩展、DDS、网络、机器人和服务；纯本地 getter 则不一定需要全部条件。

\textbf{参数}

无显式参数。实例方法中的 \passthrough{\lstinline!self!} 由 Python
自动传入。

\textbf{返回值}

\begin{longtable}[]{@{}
  >{\raggedright\arraybackslash}p{(\columnwidth - 4\tabcolsep) * \real{0.18}}
  >{\raggedright\arraybackslash}p{(\columnwidth - 4\tabcolsep) * \real{0.20}}
  >{\raggedright\arraybackslash}p{(\columnwidth - 4\tabcolsep) * \real{0.62}}@{}}
\toprule
位置 & 类型 & 含义 \\
\midrule
\endhead
{[}0{]} \passthrough{\lstinline!status!} & \passthrough{\lstinline!int!}
& SDK 状态码；Unitree 示例通常以 \passthrough{\lstinline!0!}
表示成功，非零值需按具体服务错误码解释。 \\
{[}1{]} \passthrough{\lstinline!fsm\_id!} &
\passthrough{\lstinline!int!} & 传给该接口的 FSM ID 参数，Python 类型为
\passthrough{\lstinline!int!}。精确含义、单位、范围和枚举值需查目标型号对应头文件与协议。
对应 C++ 参数 \passthrough{\lstinline!fsm\_id: int \&!}。 \\
\bottomrule
\end{longtable}

\textbf{对应 C++}

\begin{itemize}
\tightlist
\item
  类：\passthrough{\lstinline!unitree::robot::h1::LocoClient!}
\item
  签名：\passthrough{\lstinline!GetFsmId(int \&)!}
\item
  绑定策略：\passthrough{\lstinline!OUTPUT\_WRAPPER!}
\item
  声明位置：\passthrough{\lstinline!include/unitree/robot/h1/loco/h1\_loco\_client.hpp:42!}
\end{itemize}

\textbf{说明}

\begin{itemize}
\tightlist
\item
  示例是局部调用片段；\passthrough{\lstinline!obj!}
  和参数变量表示已经按上层业务校验并准备好的对象。
\item
  C++ 的可修改输出引用不会作为 Python 入参出现，而是按 C++
  参数顺序追加到返回元组中。
\end{itemize}

\textbf{用法}

\begin{lstlisting}[language=Python]
status, fsm_id = obj.get_fsm_id()
\end{lstlisting}

\begin{quote}
{[}!NOTE{]} 只读查询仍会访问
DDS/机器人服务。必须检查返回状态码，并处理超时和网络中断。
\end{quote}

\hypertarget{unitree-sdk2-cpp-robot-h1-lococlient-get-fsm-mode-1}{%
\paragraph{\texorpdfstring{\texttt{unitree\_sdk2\_cpp.robot.h1.LocoClient.get\_fsm\_mode}}{unitree\_sdk2\_cpp.robot.h1.LocoClient.get\_fsm\_mode}}\label{unitree-sdk2-cpp-robot-h1-lococlient-get-fsm-mode-1}}

查询或检查 FSM 模式。

\textbf{签名}

\begin{lstlisting}[language=Python]
def get_fsm_mode(self) -> tuple[int, int]
\end{lstlisting}

\textbf{可用性与安全性}

\textbf{\passthrough{\lstinline!AVAILABLE!} /
\passthrough{\lstinline!READ\_ONLY!}}：当前绑定源码已实现只读查询。Robot
Client 的服务端查询通常仍需匹配的 Linux
扩展、DDS、网络、机器人和服务；纯本地 getter 则不一定需要全部条件。

\textbf{参数}

无显式参数。实例方法中的 \passthrough{\lstinline!self!} 由 Python
自动传入。

\textbf{返回值}

\begin{longtable}[]{@{}
  >{\raggedright\arraybackslash}p{(\columnwidth - 4\tabcolsep) * \real{0.18}}
  >{\raggedright\arraybackslash}p{(\columnwidth - 4\tabcolsep) * \real{0.20}}
  >{\raggedright\arraybackslash}p{(\columnwidth - 4\tabcolsep) * \real{0.62}}@{}}
\toprule
位置 & 类型 & 含义 \\
\midrule
\endhead
{[}0{]} \passthrough{\lstinline!status!} & \passthrough{\lstinline!int!}
& SDK 状态码；Unitree 示例通常以 \passthrough{\lstinline!0!}
表示成功，非零值需按具体服务错误码解释。 \\
{[}1{]} \passthrough{\lstinline!fsm\_mode!} &
\passthrough{\lstinline!int!} & 传给该接口的 FSM 模式 参数，Python
类型为
\passthrough{\lstinline!int!}。精确含义、单位、范围和枚举值需查目标型号对应头文件与协议。
对应 C++ 参数 \passthrough{\lstinline!fsm\_mode: int \&!}。 \\
\bottomrule
\end{longtable}

\textbf{对应 C++}

\begin{itemize}
\tightlist
\item
  类：\passthrough{\lstinline!unitree::robot::h1::LocoClient!}
\item
  签名：\passthrough{\lstinline!GetFsmMode(int \&)!}
\item
  绑定策略：\passthrough{\lstinline!OUTPUT\_WRAPPER!}
\item
  声明位置：\passthrough{\lstinline!include/unitree/robot/h1/loco/h1\_loco\_client.hpp:54!}
\end{itemize}

\textbf{说明}

\begin{itemize}
\tightlist
\item
  示例是局部调用片段；\passthrough{\lstinline!obj!}
  和参数变量表示已经按上层业务校验并准备好的对象。
\item
  C++ 的可修改输出引用不会作为 Python 入参出现，而是按 C++
  参数顺序追加到返回元组中。
\end{itemize}

\textbf{用法}

\begin{lstlisting}[language=Python]
status, fsm_mode = obj.get_fsm_mode()
\end{lstlisting}

\begin{quote}
{[}!NOTE{]} 只读查询仍会访问
DDS/机器人服务。必须检查返回状态码，并处理超时和网络中断。
\end{quote}

\hypertarget{unitree-sdk2-cpp-robot-h1-lococlient-get-balance-mode-1}{%
\paragraph{\texorpdfstring{\texttt{unitree\_sdk2\_cpp.robot.h1.LocoClient.get\_balance\_mode}}{unitree\_sdk2\_cpp.robot.h1.LocoClient.get\_balance\_mode}}\label{unitree-sdk2-cpp-robot-h1-lococlient-get-balance-mode-1}}

查询或检查 平衡 模式。

\textbf{签名}

\begin{lstlisting}[language=Python]
def get_balance_mode(self) -> tuple[int, int]
\end{lstlisting}

\textbf{可用性与安全性}

\textbf{\passthrough{\lstinline!AVAILABLE!} /
\passthrough{\lstinline!READ\_ONLY!}}：当前绑定源码已实现只读查询。Robot
Client 的服务端查询通常仍需匹配的 Linux
扩展、DDS、网络、机器人和服务；纯本地 getter 则不一定需要全部条件。

\textbf{参数}

无显式参数。实例方法中的 \passthrough{\lstinline!self!} 由 Python
自动传入。

\textbf{返回值}

\begin{longtable}[]{@{}
  >{\raggedright\arraybackslash}p{(\columnwidth - 4\tabcolsep) * \real{0.18}}
  >{\raggedright\arraybackslash}p{(\columnwidth - 4\tabcolsep) * \real{0.20}}
  >{\raggedright\arraybackslash}p{(\columnwidth - 4\tabcolsep) * \real{0.62}}@{}}
\toprule
位置 & 类型 & 含义 \\
\midrule
\endhead
{[}0{]} \passthrough{\lstinline!status!} & \passthrough{\lstinline!int!}
& SDK 状态码；Unitree 示例通常以 \passthrough{\lstinline!0!}
表示成功，非零值需按具体服务错误码解释。 \\
{[}1{]} \passthrough{\lstinline!balance\_mode!} &
\passthrough{\lstinline!int!} & 传给该接口的 平衡 模式 参数，Python
类型为
\passthrough{\lstinline!int!}。精确含义、单位、范围和枚举值需查目标型号对应头文件与协议。
对应 C++ 参数 \passthrough{\lstinline!balance\_mode: int \&!}。 \\
\bottomrule
\end{longtable}

\textbf{对应 C++}

\begin{itemize}
\tightlist
\item
  类：\passthrough{\lstinline!unitree::robot::h1::LocoClient!}
\item
  签名：\passthrough{\lstinline!GetBalanceMode(int \&)!}
\item
  绑定策略：\passthrough{\lstinline!OUTPUT\_WRAPPER!}
\item
  声明位置：\passthrough{\lstinline!include/unitree/robot/h1/loco/h1\_loco\_client.hpp:66!}
\end{itemize}

\textbf{说明}

\begin{itemize}
\tightlist
\item
  示例是局部调用片段；\passthrough{\lstinline!obj!}
  和参数变量表示已经按上层业务校验并准备好的对象。
\item
  C++ 的可修改输出引用不会作为 Python 入参出现，而是按 C++
  参数顺序追加到返回元组中。
\end{itemize}

\textbf{用法}

\begin{lstlisting}[language=Python]
status, balance_mode = obj.get_balance_mode()
\end{lstlisting}

\begin{quote}
{[}!NOTE{]} 只读查询仍会访问
DDS/机器人服务。必须检查返回状态码，并处理超时和网络中断。
\end{quote}

\hypertarget{unitree-sdk2-cpp-robot-h1-lococlient-get-swing-height-1}{%
\paragraph{\texorpdfstring{\texttt{unitree\_sdk2\_cpp.robot.h1.LocoClient.get\_swing\_height}}{unitree\_sdk2\_cpp.robot.h1.LocoClient.get\_swing\_height}}\label{unitree-sdk2-cpp-robot-h1-lococlient-get-swing-height-1}}

查询或检查 \passthrough{\lstinline!swing!} 高度。

\textbf{签名}

\begin{lstlisting}[language=Python]
def get_swing_height(self) -> tuple[int, float]
\end{lstlisting}

\textbf{可用性与安全性}

\textbf{\passthrough{\lstinline!AVAILABLE!} /
\passthrough{\lstinline!READ\_ONLY!}}：当前绑定源码已实现只读查询。Robot
Client 的服务端查询通常仍需匹配的 Linux
扩展、DDS、网络、机器人和服务；纯本地 getter 则不一定需要全部条件。

\textbf{参数}

无显式参数。实例方法中的 \passthrough{\lstinline!self!} 由 Python
自动传入。

\textbf{返回值}

\begin{longtable}[]{@{}
  >{\raggedright\arraybackslash}p{(\columnwidth - 4\tabcolsep) * \real{0.18}}
  >{\raggedright\arraybackslash}p{(\columnwidth - 4\tabcolsep) * \real{0.20}}
  >{\raggedright\arraybackslash}p{(\columnwidth - 4\tabcolsep) * \real{0.62}}@{}}
\toprule
位置 & 类型 & 含义 \\
\midrule
\endhead
{[}0{]} \passthrough{\lstinline!status!} & \passthrough{\lstinline!int!}
& SDK 状态码；Unitree 示例通常以 \passthrough{\lstinline!0!}
表示成功，非零值需按具体服务错误码解释。 \\
{[}1{]} \passthrough{\lstinline!swing\_height!} &
\passthrough{\lstinline!float!} & 传给该接口的
\passthrough{\lstinline!swing!} 高度 参数，Python 类型为
\passthrough{\lstinline!float!}。精确含义、单位、范围和枚举值需查目标型号对应头文件与协议。
对应 C++ 参数 \passthrough{\lstinline!swing\_height: float \&!}。
底层为浮点数；类型本身不说明单位、坐标系或业务安全范围。 \\
\bottomrule
\end{longtable}

\textbf{对应 C++}

\begin{itemize}
\tightlist
\item
  类：\passthrough{\lstinline!unitree::robot::h1::LocoClient!}
\item
  签名：\passthrough{\lstinline!GetSwingHeight(float \&)!}
\item
  绑定策略：\passthrough{\lstinline!OUTPUT\_WRAPPER!}
\item
  声明位置：\passthrough{\lstinline!include/unitree/robot/h1/loco/h1\_loco\_client.hpp:78!}
\end{itemize}

\textbf{说明}

\begin{itemize}
\tightlist
\item
  示例是局部调用片段；\passthrough{\lstinline!obj!}
  和参数变量表示已经按上层业务校验并准备好的对象。
\item
  C++ 的可修改输出引用不会作为 Python 入参出现，而是按 C++
  参数顺序追加到返回元组中。
\end{itemize}

\textbf{用法}

\begin{lstlisting}[language=Python]
status, swing_height = obj.get_swing_height()
\end{lstlisting}

\begin{quote}
{[}!NOTE{]} 只读查询仍会访问
DDS/机器人服务。必须检查返回状态码，并处理超时和网络中断。
\end{quote}

\hypertarget{unitree-sdk2-cpp-robot-h1-lococlient-get-stand-height-1}{%
\paragraph{\texorpdfstring{\texttt{unitree\_sdk2\_cpp.robot.h1.LocoClient.get\_stand\_height}}{unitree\_sdk2\_cpp.robot.h1.LocoClient.get\_stand\_height}}\label{unitree-sdk2-cpp-robot-h1-lococlient-get-stand-height-1}}

查询或检查 \passthrough{\lstinline!stand!} 高度。

\textbf{签名}

\begin{lstlisting}[language=Python]
def get_stand_height(self) -> tuple[int, float]
\end{lstlisting}

\textbf{可用性与安全性}

\textbf{\passthrough{\lstinline!AVAILABLE!} /
\passthrough{\lstinline!READ\_ONLY!}}：当前绑定源码已实现只读查询。Robot
Client 的服务端查询通常仍需匹配的 Linux
扩展、DDS、网络、机器人和服务；纯本地 getter 则不一定需要全部条件。

\textbf{参数}

无显式参数。实例方法中的 \passthrough{\lstinline!self!} 由 Python
自动传入。

\textbf{返回值}

\begin{longtable}[]{@{}
  >{\raggedright\arraybackslash}p{(\columnwidth - 4\tabcolsep) * \real{0.18}}
  >{\raggedright\arraybackslash}p{(\columnwidth - 4\tabcolsep) * \real{0.20}}
  >{\raggedright\arraybackslash}p{(\columnwidth - 4\tabcolsep) * \real{0.62}}@{}}
\toprule
位置 & 类型 & 含义 \\
\midrule
\endhead
{[}0{]} \passthrough{\lstinline!status!} & \passthrough{\lstinline!int!}
& SDK 状态码；Unitree 示例通常以 \passthrough{\lstinline!0!}
表示成功，非零值需按具体服务错误码解释。 \\
{[}1{]} \passthrough{\lstinline!stand\_height!} &
\passthrough{\lstinline!float!} & 传给该接口的
\passthrough{\lstinline!stand!} 高度 参数，Python 类型为
\passthrough{\lstinline!float!}。精确含义、单位、范围和枚举值需查目标型号对应头文件与协议。
对应 C++ 参数 \passthrough{\lstinline!stand\_height: float \&!}。
底层为浮点数；类型本身不说明单位、坐标系或业务安全范围。 \\
\bottomrule
\end{longtable}

\textbf{对应 C++}

\begin{itemize}
\tightlist
\item
  类：\passthrough{\lstinline!unitree::robot::h1::LocoClient!}
\item
  签名：\passthrough{\lstinline!GetStandHeight(float \&)!}
\item
  绑定策略：\passthrough{\lstinline!OUTPUT\_WRAPPER!}
\item
  声明位置：\passthrough{\lstinline!include/unitree/robot/h1/loco/h1\_loco\_client.hpp:90!}
\end{itemize}

\textbf{说明}

\begin{itemize}
\tightlist
\item
  示例是局部调用片段；\passthrough{\lstinline!obj!}
  和参数变量表示已经按上层业务校验并准备好的对象。
\item
  C++ 的可修改输出引用不会作为 Python 入参出现，而是按 C++
  参数顺序追加到返回元组中。
\end{itemize}

\textbf{用法}

\begin{lstlisting}[language=Python]
status, stand_height = obj.get_stand_height()
\end{lstlisting}

\begin{quote}
{[}!NOTE{]} 只读查询仍会访问
DDS/机器人服务。必须检查返回状态码，并处理超时和网络中断。
\end{quote}

\hypertarget{unitree-sdk2-cpp-robot-h1-lococlient-get-phase-1}{%
\paragraph{\texorpdfstring{\texttt{unitree\_sdk2\_cpp.robot.h1.LocoClient.get\_phase}}{unitree\_sdk2\_cpp.robot.h1.LocoClient.get\_phase}}\label{unitree-sdk2-cpp-robot-h1-lococlient-get-phase-1}}

查询或检查 相位。

\textbf{签名}

\begin{lstlisting}[language=Python]
def get_phase(self) -> tuple[int, list[float]]
\end{lstlisting}

\textbf{可用性与安全性}

\textbf{\passthrough{\lstinline!AVAILABLE!} /
\passthrough{\lstinline!READ\_ONLY!}}：当前绑定源码已实现只读查询。Robot
Client 的服务端查询通常仍需匹配的 Linux
扩展、DDS、网络、机器人和服务；纯本地 getter 则不一定需要全部条件。

\textbf{参数}

无显式参数。实例方法中的 \passthrough{\lstinline!self!} 由 Python
自动传入。

\textbf{返回值}

\begin{longtable}[]{@{}
  >{\raggedright\arraybackslash}p{(\columnwidth - 4\tabcolsep) * \real{0.18}}
  >{\raggedright\arraybackslash}p{(\columnwidth - 4\tabcolsep) * \real{0.20}}
  >{\raggedright\arraybackslash}p{(\columnwidth - 4\tabcolsep) * \real{0.62}}@{}}
\toprule
位置 & 类型 & 含义 \\
\midrule
\endhead
{[}0{]} \passthrough{\lstinline!status!} & \passthrough{\lstinline!int!}
& SDK 状态码；Unitree 示例通常以 \passthrough{\lstinline!0!}
表示成功，非零值需按具体服务错误码解释。 \\
{[}1{]} \passthrough{\lstinline!phase!} &
\passthrough{\lstinline!list[float]!} & 传给该接口的 相位 参数，Python
类型为
\passthrough{\lstinline!list[float]!}。精确含义、单位、范围和枚举值需查目标型号对应头文件与协议。
对应 C++ 参数 \passthrough{\lstinline!phase: std::vector<float> \&!}。
底层是可变长度
vector；元素约束：底层为浮点数；类型本身不说明单位、坐标系或业务安全范围。 \\
\bottomrule
\end{longtable}

\textbf{对应 C++}

\begin{itemize}
\tightlist
\item
  类：\passthrough{\lstinline!unitree::robot::h1::LocoClient!}
\item
  签名：\passthrough{\lstinline!GetPhase(std::vector<float> \&)!}
\item
  绑定策略：\passthrough{\lstinline!OUTPUT\_WRAPPER!}
\item
  声明位置：\passthrough{\lstinline!include/unitree/robot/h1/loco/h1\_loco\_client.hpp:102!}
\end{itemize}

\textbf{说明}

\begin{itemize}
\tightlist
\item
  示例是局部调用片段；\passthrough{\lstinline!obj!}
  和参数变量表示已经按上层业务校验并准备好的对象。
\item
  C++ 的可修改输出引用不会作为 Python 入参出现，而是按 C++
  参数顺序追加到返回元组中。
\end{itemize}

\textbf{用法}

\begin{lstlisting}[language=Python]
status, phase = obj.get_phase()
\end{lstlisting}

\begin{quote}
{[}!NOTE{]} 只读查询仍会访问
DDS/机器人服务。必须检查返回状态码，并处理超时和网络中断。
\end{quote}

\hypertarget{unitree-sdk2-cpp-robot-h1-lococlient-enable-odom-1}{%
\paragraph{\texorpdfstring{\texttt{unitree\_sdk2\_cpp.robot.h1.LocoClient.enable\_odom}}{unitree\_sdk2\_cpp.robot.h1.LocoClient.enable\_odom}}\label{unitree-sdk2-cpp-robot-h1-lococlient-enable-odom-1}}

对应 C++ SDK 操作
\passthrough{\lstinline!EnableOdom()!}。上游头文件没有可直接生成的业务说明时，本参考只保证签名映射，精确语义需查目标型号协议。

\textbf{签名}

\begin{lstlisting}[language=Python]
def enable_odom(self) -> int
\end{lstlisting}

\textbf{可用性与安全性}

\textbf{\passthrough{\lstinline!SIGNATURE\_ONLY!} /
\passthrough{\lstinline!MOTION\_COMMAND!}}：仅用于补全和静态检查。当前二进制没有该
Python 方法；不得按可执行运动接口使用。

\textbf{参数}

无显式参数。实例方法中的 \passthrough{\lstinline!self!} 由 Python
自动传入。

\textbf{返回值}

\begin{longtable}[]{@{}
  >{\raggedright\arraybackslash}p{(\columnwidth - 4\tabcolsep) * \real{0.18}}
  >{\raggedright\arraybackslash}p{(\columnwidth - 4\tabcolsep) * \real{0.20}}
  >{\raggedright\arraybackslash}p{(\columnwidth - 4\tabcolsep) * \real{0.62}}@{}}
\toprule
位置 & 类型 & 含义 \\
\midrule
\endhead
返回值 & \passthrough{\lstinline!int!} & SDK 状态码；Unitree 示例通常以
\passthrough{\lstinline!0!} 表示成功，非零值需按具体服务定义解释。 \\
\bottomrule
\end{longtable}

\textbf{对应 C++}

\begin{itemize}
\tightlist
\item
  类：\passthrough{\lstinline!unitree::robot::h1::LocoClient!}
\item
  签名：\passthrough{\lstinline!EnableOdom()!}
\item
  绑定策略：\passthrough{\lstinline!DIRECT!}
\item
  声明位置：\passthrough{\lstinline!include/unitree/robot/h1/loco/h1\_loco\_client.hpp:114!}
\end{itemize}

\textbf{说明}

\begin{itemize}
\tightlist
\item
  示例是局部调用片段；\passthrough{\lstinline!obj!}
  和参数变量表示已经按上层业务校验并准备好的对象。
\item
  该条目可以被编辑器和类型检查器识别，但当前运行时可能在属性查找阶段直接失败。
\item
  该方法属于运动命令。方法名包含
  \passthrough{\lstinline!stop!}、\passthrough{\lstinline!damp!} 或
  \passthrough{\lstinline!zero!} 也不代表它是独立物理急停。
\end{itemize}

\textbf{用法}

\begin{lstlisting}[language=Python]
from typing import TYPE_CHECKING

if TYPE_CHECKING:
    def planned_enable_odom(obj: LocoClient) -> int:
        return obj.enable_odom()
\end{lstlisting}

\begin{quote}
{[}!CAUTION{]} 上面的 \passthrough{\lstinline!TYPE\_CHECKING!}
示例只表达计划调用形态。该分支在普通 Python
运行时为假，不代表当前扩展能执行此接口。
\end{quote}

\hypertarget{unitree-sdk2-cpp-robot-h1-lococlient-disable-odom-1}{%
\paragraph{\texorpdfstring{\texttt{unitree\_sdk2\_cpp.robot.h1.LocoClient.disable\_odom}}{unitree\_sdk2\_cpp.robot.h1.LocoClient.disable\_odom}}\label{unitree-sdk2-cpp-robot-h1-lococlient-disable-odom-1}}

对应 C++ SDK 操作
\passthrough{\lstinline!DisableOdom()!}。上游头文件没有可直接生成的业务说明时，本参考只保证签名映射，精确语义需查目标型号协议。

\textbf{签名}

\begin{lstlisting}[language=Python]
def disable_odom(self) -> int
\end{lstlisting}

\textbf{可用性与安全性}

\textbf{\passthrough{\lstinline!SIGNATURE\_ONLY!} /
\passthrough{\lstinline!MOTION\_COMMAND!}}：仅用于补全和静态检查。当前二进制没有该
Python 方法；不得按可执行运动接口使用。

\textbf{参数}

无显式参数。实例方法中的 \passthrough{\lstinline!self!} 由 Python
自动传入。

\textbf{返回值}

\begin{longtable}[]{@{}
  >{\raggedright\arraybackslash}p{(\columnwidth - 4\tabcolsep) * \real{0.18}}
  >{\raggedright\arraybackslash}p{(\columnwidth - 4\tabcolsep) * \real{0.20}}
  >{\raggedright\arraybackslash}p{(\columnwidth - 4\tabcolsep) * \real{0.62}}@{}}
\toprule
位置 & 类型 & 含义 \\
\midrule
\endhead
返回值 & \passthrough{\lstinline!int!} & SDK 状态码；Unitree 示例通常以
\passthrough{\lstinline!0!} 表示成功，非零值需按具体服务定义解释。 \\
\bottomrule
\end{longtable}

\textbf{对应 C++}

\begin{itemize}
\tightlist
\item
  类：\passthrough{\lstinline!unitree::robot::h1::LocoClient!}
\item
  签名：\passthrough{\lstinline!DisableOdom()!}
\item
  绑定策略：\passthrough{\lstinline!DIRECT!}
\item
  声明位置：\passthrough{\lstinline!include/unitree/robot/h1/loco/h1\_loco\_client.hpp:119!}
\end{itemize}

\textbf{说明}

\begin{itemize}
\tightlist
\item
  示例是局部调用片段；\passthrough{\lstinline!obj!}
  和参数变量表示已经按上层业务校验并准备好的对象。
\item
  该条目可以被编辑器和类型检查器识别，但当前运行时可能在属性查找阶段直接失败。
\item
  该方法属于运动命令。方法名包含
  \passthrough{\lstinline!stop!}、\passthrough{\lstinline!damp!} 或
  \passthrough{\lstinline!zero!} 也不代表它是独立物理急停。
\end{itemize}

\textbf{用法}

\begin{lstlisting}[language=Python]
from typing import TYPE_CHECKING

if TYPE_CHECKING:
    def planned_disable_odom(obj: LocoClient) -> int:
        return obj.disable_odom()
\end{lstlisting}

\begin{quote}
{[}!CAUTION{]} 上面的 \passthrough{\lstinline!TYPE\_CHECKING!}
示例只表达计划调用形态。该分支在普通 Python
运行时为假，不代表当前扩展能执行此接口。
\end{quote}

\hypertarget{unitree-sdk2-cpp-robot-h1-lococlient-get-odom-1}{%
\paragraph{\texorpdfstring{\texttt{unitree\_sdk2\_cpp.robot.h1.LocoClient.get\_odom}}{unitree\_sdk2\_cpp.robot.h1.LocoClient.get\_odom}}\label{unitree-sdk2-cpp-robot-h1-lococlient-get-odom-1}}

查询或检查 \passthrough{\lstinline!odom!}。

\textbf{签名}

\begin{lstlisting}[language=Python]
def get_odom(self) -> tuple[int, float, float, float]
\end{lstlisting}

\textbf{可用性与安全性}

\textbf{\passthrough{\lstinline!AVAILABLE!} /
\passthrough{\lstinline!READ\_ONLY!}}：当前绑定源码已实现只读查询。Robot
Client 的服务端查询通常仍需匹配的 Linux
扩展、DDS、网络、机器人和服务；纯本地 getter 则不一定需要全部条件。

\textbf{参数}

无显式参数。实例方法中的 \passthrough{\lstinline!self!} 由 Python
自动传入。

\textbf{返回值}

\begin{longtable}[]{@{}
  >{\raggedright\arraybackslash}p{(\columnwidth - 4\tabcolsep) * \real{0.18}}
  >{\raggedright\arraybackslash}p{(\columnwidth - 4\tabcolsep) * \real{0.20}}
  >{\raggedright\arraybackslash}p{(\columnwidth - 4\tabcolsep) * \real{0.62}}@{}}
\toprule
位置 & 类型 & 含义 \\
\midrule
\endhead
{[}0{]} \passthrough{\lstinline!status!} & \passthrough{\lstinline!int!}
& SDK 状态码；Unitree 示例通常以 \passthrough{\lstinline!0!}
表示成功，非零值需按具体服务错误码解释。 \\
{[}1{]} \passthrough{\lstinline!x!} & \passthrough{\lstinline!float!} &
X 方向位置或分量参数。参考系、单位和范围必须查当前方法及目标型号协议。
对应 C++ 参数 \passthrough{\lstinline!x: float \&!}。
底层为浮点数；类型本身不说明单位、坐标系或业务安全范围。 \\
{[}2{]} \passthrough{\lstinline!y!} & \passthrough{\lstinline!float!} &
Y 方向位置或分量参数。参考系、单位和范围必须查当前方法及目标型号协议。
对应 C++ 参数 \passthrough{\lstinline!y: float \&!}。
底层为浮点数；类型本身不说明单位、坐标系或业务安全范围。 \\
{[}3{]} \passthrough{\lstinline!yaw!} & \passthrough{\lstinline!float!}
& 偏航角或偏航目标参数。参考系、单位和范围必须查目标型号协议。 对应 C++
参数 \passthrough{\lstinline!yaw: float \&!}。
底层为浮点数；类型本身不说明单位、坐标系或业务安全范围。 \\
\bottomrule
\end{longtable}

\textbf{对应 C++}

\begin{itemize}
\tightlist
\item
  类：\passthrough{\lstinline!unitree::robot::h1::LocoClient!}
\item
  签名：\passthrough{\lstinline!GetOdom(float \&, float \&, float \&)!}
\item
  绑定策略：\passthrough{\lstinline!OUTPUT\_WRAPPER!}
\item
  声明位置：\passthrough{\lstinline!include/unitree/robot/h1/loco/h1\_loco\_client.hpp:124!}
\end{itemize}

\textbf{说明}

\begin{itemize}
\tightlist
\item
  示例是局部调用片段；\passthrough{\lstinline!obj!}
  和参数变量表示已经按上层业务校验并准备好的对象。
\item
  C++ 的可修改输出引用不会作为 Python 入参出现，而是按 C++
  参数顺序追加到返回元组中。
\end{itemize}

\textbf{用法}

\begin{lstlisting}[language=Python]
status, x, y, yaw = obj.get_odom()
\end{lstlisting}

\begin{quote}
{[}!NOTE{]} 只读查询仍会访问
DDS/机器人服务。必须检查返回状态码，并处理超时和网络中断。
\end{quote}

\hypertarget{unitree-sdk2-cpp-robot-h1-lococlient-set-fsm-id-1}{%
\paragraph{\texorpdfstring{\texttt{unitree\_sdk2\_cpp.robot.h1.LocoClient.set\_fsm\_id}}{unitree\_sdk2\_cpp.robot.h1.LocoClient.set\_fsm\_id}}\label{unitree-sdk2-cpp-robot-h1-lococlient-set-fsm-id-1}}

设置 FSM ID。具体副作用和安全边界见下方状态。

\textbf{签名}

\begin{lstlisting}[language=Python]
def set_fsm_id(self, fsm_id: int) -> int
\end{lstlisting}

\textbf{可用性与安全性}

\textbf{\passthrough{\lstinline!SIGNATURE\_ONLY!} /
\passthrough{\lstinline!MOTION\_COMMAND!}}：仅用于补全和静态检查。当前二进制没有该
Python 方法；不得按可执行运动接口使用。

\textbf{参数}

\begin{longtable}[]{@{}
  >{\raggedright\arraybackslash}p{(\columnwidth - 6\tabcolsep) * \real{0.18}}
  >{\raggedright\arraybackslash}p{(\columnwidth - 6\tabcolsep) * \real{0.24}}
  >{\raggedright\arraybackslash}p{(\columnwidth - 6\tabcolsep) * \real{0.18}}
  >{\raggedright\arraybackslash}p{(\columnwidth - 6\tabcolsep) * \real{0.40}}@{}}
\toprule
名称 & Python 类型 & 默认值 & 含义 \\
\midrule
\endhead
\passthrough{\lstinline!fsm\_id!} & \passthrough{\lstinline!int!} & 必填
& 传给该接口的 FSM ID 参数，Python 类型为
\passthrough{\lstinline!int!}。精确含义、单位、范围和枚举值需查目标型号对应头文件与协议。
对应 C++ 参数 \passthrough{\lstinline!fsm\_id: int!}。 \\
\bottomrule
\end{longtable}

\textbf{返回值}

\begin{longtable}[]{@{}
  >{\raggedright\arraybackslash}p{(\columnwidth - 4\tabcolsep) * \real{0.18}}
  >{\raggedright\arraybackslash}p{(\columnwidth - 4\tabcolsep) * \real{0.20}}
  >{\raggedright\arraybackslash}p{(\columnwidth - 4\tabcolsep) * \real{0.62}}@{}}
\toprule
位置 & 类型 & 含义 \\
\midrule
\endhead
返回值 & \passthrough{\lstinline!int!} & SDK 状态码；Unitree 示例通常以
\passthrough{\lstinline!0!} 表示成功，非零值需按具体服务定义解释。 \\
\bottomrule
\end{longtable}

\textbf{对应 C++}

\begin{itemize}
\tightlist
\item
  类：\passthrough{\lstinline!unitree::robot::h1::LocoClient!}
\item
  签名：\passthrough{\lstinline!SetFsmId(int)!}
\item
  绑定策略：\passthrough{\lstinline!DIRECT!}
\item
  声明位置：\passthrough{\lstinline!include/unitree/robot/h1/loco/h1\_loco\_client.hpp:138!}
\end{itemize}

\textbf{说明}

\begin{itemize}
\tightlist
\item
  示例是局部调用片段；\passthrough{\lstinline!obj!}
  和参数变量表示已经按上层业务校验并准备好的对象。
\item
  该条目可以被编辑器和类型检查器识别，但当前运行时可能在属性查找阶段直接失败。
\item
  该方法属于运动命令。方法名包含
  \passthrough{\lstinline!stop!}、\passthrough{\lstinline!damp!} 或
  \passthrough{\lstinline!zero!} 也不代表它是独立物理急停。
\end{itemize}

\textbf{用法}

\begin{lstlisting}[language=Python]
from typing import TYPE_CHECKING

if TYPE_CHECKING:
    def planned_set_fsm_id(obj: LocoClient, fsm_id: int) -> int:
        return obj.set_fsm_id(fsm_id=fsm_id)
\end{lstlisting}

\begin{quote}
{[}!CAUTION{]} 上面的 \passthrough{\lstinline!TYPE\_CHECKING!}
示例只表达计划调用形态。该分支在普通 Python
运行时为假，不代表当前扩展能执行此接口。
\end{quote}

\hypertarget{unitree-sdk2-cpp-robot-h1-lococlient-set-balance-mode-1}{%
\paragraph{\texorpdfstring{\texttt{unitree\_sdk2\_cpp.robot.h1.LocoClient.set\_balance\_mode}}{unitree\_sdk2\_cpp.robot.h1.LocoClient.set\_balance\_mode}}\label{unitree-sdk2-cpp-robot-h1-lococlient-set-balance-mode-1}}

设置 平衡 模式。具体副作用和安全边界见下方状态。

\textbf{签名}

\begin{lstlisting}[language=Python]
def set_balance_mode(self, balance_mode: int) -> int
\end{lstlisting}

\textbf{可用性与安全性}

\textbf{\passthrough{\lstinline!SIGNATURE\_ONLY!} /
\passthrough{\lstinline!MOTION\_COMMAND!}}：仅用于补全和静态检查。当前二进制没有该
Python 方法；不得按可执行运动接口使用。

\textbf{参数}

\begin{longtable}[]{@{}
  >{\raggedright\arraybackslash}p{(\columnwidth - 6\tabcolsep) * \real{0.18}}
  >{\raggedright\arraybackslash}p{(\columnwidth - 6\tabcolsep) * \real{0.24}}
  >{\raggedright\arraybackslash}p{(\columnwidth - 6\tabcolsep) * \real{0.18}}
  >{\raggedright\arraybackslash}p{(\columnwidth - 6\tabcolsep) * \real{0.40}}@{}}
\toprule
名称 & Python 类型 & 默认值 & 含义 \\
\midrule
\endhead
\passthrough{\lstinline!balance\_mode!} & \passthrough{\lstinline!int!}
& 必填 & 传给该接口的 平衡 模式 参数，Python 类型为
\passthrough{\lstinline!int!}。精确含义、单位、范围和枚举值需查目标型号对应头文件与协议。
对应 C++ 参数 \passthrough{\lstinline!balance\_mode: int!}。 \\
\bottomrule
\end{longtable}

\textbf{返回值}

\begin{longtable}[]{@{}
  >{\raggedright\arraybackslash}p{(\columnwidth - 4\tabcolsep) * \real{0.18}}
  >{\raggedright\arraybackslash}p{(\columnwidth - 4\tabcolsep) * \real{0.20}}
  >{\raggedright\arraybackslash}p{(\columnwidth - 4\tabcolsep) * \real{0.62}}@{}}
\toprule
位置 & 类型 & 含义 \\
\midrule
\endhead
返回值 & \passthrough{\lstinline!int!} & SDK 状态码；Unitree 示例通常以
\passthrough{\lstinline!0!} 表示成功，非零值需按具体服务定义解释。 \\
\bottomrule
\end{longtable}

\textbf{对应 C++}

\begin{itemize}
\tightlist
\item
  类：\passthrough{\lstinline!unitree::robot::h1::LocoClient!}
\item
  签名：\passthrough{\lstinline!SetBalanceMode(int)!}
\item
  绑定策略：\passthrough{\lstinline!DIRECT!}
\item
  声明位置：\passthrough{\lstinline!include/unitree/robot/h1/loco/h1\_loco\_client.hpp:148!}
\end{itemize}

\textbf{说明}

\begin{itemize}
\tightlist
\item
  示例是局部调用片段；\passthrough{\lstinline!obj!}
  和参数变量表示已经按上层业务校验并准备好的对象。
\item
  该条目可以被编辑器和类型检查器识别，但当前运行时可能在属性查找阶段直接失败。
\item
  该方法属于运动命令。方法名包含
  \passthrough{\lstinline!stop!}、\passthrough{\lstinline!damp!} 或
  \passthrough{\lstinline!zero!} 也不代表它是独立物理急停。
\end{itemize}

\textbf{用法}

\begin{lstlisting}[language=Python]
from typing import TYPE_CHECKING

if TYPE_CHECKING:
    def planned_set_balance_mode(obj: LocoClient, balance_mode: int) -> int:
        return obj.set_balance_mode(balance_mode=balance_mode)
\end{lstlisting}

\begin{quote}
{[}!CAUTION{]} 上面的 \passthrough{\lstinline!TYPE\_CHECKING!}
示例只表达计划调用形态。该分支在普通 Python
运行时为假，不代表当前扩展能执行此接口。
\end{quote}

\hypertarget{unitree-sdk2-cpp-robot-h1-lococlient-set-swing-height-1}{%
\paragraph{\texorpdfstring{\texttt{unitree\_sdk2\_cpp.robot.h1.LocoClient.set\_swing\_height}}{unitree\_sdk2\_cpp.robot.h1.LocoClient.set\_swing\_height}}\label{unitree-sdk2-cpp-robot-h1-lococlient-set-swing-height-1}}

设置 \passthrough{\lstinline!swing!}
高度。具体副作用和安全边界见下方状态。

\textbf{签名}

\begin{lstlisting}[language=Python]
def set_swing_height(self, swing_height: float) -> int
\end{lstlisting}

\textbf{可用性与安全性}

\textbf{\passthrough{\lstinline!SIGNATURE\_ONLY!} /
\passthrough{\lstinline!MOTION\_COMMAND!}}：仅用于补全和静态检查。当前二进制没有该
Python 方法；不得按可执行运动接口使用。

\textbf{参数}

\begin{longtable}[]{@{}
  >{\raggedright\arraybackslash}p{(\columnwidth - 6\tabcolsep) * \real{0.18}}
  >{\raggedright\arraybackslash}p{(\columnwidth - 6\tabcolsep) * \real{0.24}}
  >{\raggedright\arraybackslash}p{(\columnwidth - 6\tabcolsep) * \real{0.18}}
  >{\raggedright\arraybackslash}p{(\columnwidth - 6\tabcolsep) * \real{0.40}}@{}}
\toprule
名称 & Python 类型 & 默认值 & 含义 \\
\midrule
\endhead
\passthrough{\lstinline!swing\_height!} &
\passthrough{\lstinline!float!} & 必填 & 传给该接口的
\passthrough{\lstinline!swing!} 高度 参数，Python 类型为
\passthrough{\lstinline!float!}。精确含义、单位、范围和枚举值需查目标型号对应头文件与协议。
对应 C++ 参数 \passthrough{\lstinline!swing\_height: float!}。
底层为浮点数；类型本身不说明单位、坐标系或业务安全范围。 \\
\bottomrule
\end{longtable}

\textbf{返回值}

\begin{longtable}[]{@{}
  >{\raggedright\arraybackslash}p{(\columnwidth - 4\tabcolsep) * \real{0.18}}
  >{\raggedright\arraybackslash}p{(\columnwidth - 4\tabcolsep) * \real{0.20}}
  >{\raggedright\arraybackslash}p{(\columnwidth - 4\tabcolsep) * \real{0.62}}@{}}
\toprule
位置 & 类型 & 含义 \\
\midrule
\endhead
返回值 & \passthrough{\lstinline!int!} & SDK 状态码；Unitree 示例通常以
\passthrough{\lstinline!0!} 表示成功，非零值需按具体服务定义解释。 \\
\bottomrule
\end{longtable}

\textbf{对应 C++}

\begin{itemize}
\tightlist
\item
  类：\passthrough{\lstinline!unitree::robot::h1::LocoClient!}
\item
  签名：\passthrough{\lstinline!SetSwingHeight(float)!}
\item
  绑定策略：\passthrough{\lstinline!DIRECT!}
\item
  声明位置：\passthrough{\lstinline!include/unitree/robot/h1/loco/h1\_loco\_client.hpp:158!}
\end{itemize}

\textbf{说明}

\begin{itemize}
\tightlist
\item
  示例是局部调用片段；\passthrough{\lstinline!obj!}
  和参数变量表示已经按上层业务校验并准备好的对象。
\item
  该条目可以被编辑器和类型检查器识别，但当前运行时可能在属性查找阶段直接失败。
\item
  该方法属于运动命令。方法名包含
  \passthrough{\lstinline!stop!}、\passthrough{\lstinline!damp!} 或
  \passthrough{\lstinline!zero!} 也不代表它是独立物理急停。
\end{itemize}

\textbf{用法}

\begin{lstlisting}[language=Python]
from typing import TYPE_CHECKING

if TYPE_CHECKING:
    def planned_set_swing_height(obj: LocoClient, swing_height: float) -> int:
        return obj.set_swing_height(swing_height=swing_height)
\end{lstlisting}

\begin{quote}
{[}!CAUTION{]} 上面的 \passthrough{\lstinline!TYPE\_CHECKING!}
示例只表达计划调用形态。该分支在普通 Python
运行时为假，不代表当前扩展能执行此接口。
\end{quote}

\hypertarget{unitree-sdk2-cpp-robot-h1-lococlient-set-stand-height-1}{%
\paragraph{\texorpdfstring{\texttt{unitree\_sdk2\_cpp.robot.h1.LocoClient.set\_stand\_height}}{unitree\_sdk2\_cpp.robot.h1.LocoClient.set\_stand\_height}}\label{unitree-sdk2-cpp-robot-h1-lococlient-set-stand-height-1}}

设置 \passthrough{\lstinline!stand!}
高度。具体副作用和安全边界见下方状态。

\textbf{签名}

\begin{lstlisting}[language=Python]
def set_stand_height(self, stand_height: float) -> int
\end{lstlisting}

\textbf{可用性与安全性}

\textbf{\passthrough{\lstinline!SIGNATURE\_ONLY!} /
\passthrough{\lstinline!MOTION\_COMMAND!}}：仅用于补全和静态检查。当前二进制没有该
Python 方法；不得按可执行运动接口使用。

\textbf{参数}

\begin{longtable}[]{@{}
  >{\raggedright\arraybackslash}p{(\columnwidth - 6\tabcolsep) * \real{0.18}}
  >{\raggedright\arraybackslash}p{(\columnwidth - 6\tabcolsep) * \real{0.24}}
  >{\raggedright\arraybackslash}p{(\columnwidth - 6\tabcolsep) * \real{0.18}}
  >{\raggedright\arraybackslash}p{(\columnwidth - 6\tabcolsep) * \real{0.40}}@{}}
\toprule
名称 & Python 类型 & 默认值 & 含义 \\
\midrule
\endhead
\passthrough{\lstinline!stand\_height!} &
\passthrough{\lstinline!float!} & 必填 & 传给该接口的
\passthrough{\lstinline!stand!} 高度 参数，Python 类型为
\passthrough{\lstinline!float!}。精确含义、单位、范围和枚举值需查目标型号对应头文件与协议。
对应 C++ 参数 \passthrough{\lstinline!stand\_height: float!}。
底层为浮点数；类型本身不说明单位、坐标系或业务安全范围。 \\
\bottomrule
\end{longtable}

\textbf{返回值}

\begin{longtable}[]{@{}
  >{\raggedright\arraybackslash}p{(\columnwidth - 4\tabcolsep) * \real{0.18}}
  >{\raggedright\arraybackslash}p{(\columnwidth - 4\tabcolsep) * \real{0.20}}
  >{\raggedright\arraybackslash}p{(\columnwidth - 4\tabcolsep) * \real{0.62}}@{}}
\toprule
位置 & 类型 & 含义 \\
\midrule
\endhead
返回值 & \passthrough{\lstinline!int!} & SDK 状态码；Unitree 示例通常以
\passthrough{\lstinline!0!} 表示成功，非零值需按具体服务定义解释。 \\
\bottomrule
\end{longtable}

\textbf{对应 C++}

\begin{itemize}
\tightlist
\item
  类：\passthrough{\lstinline!unitree::robot::h1::LocoClient!}
\item
  签名：\passthrough{\lstinline!SetStandHeight(float)!}
\item
  绑定策略：\passthrough{\lstinline!DIRECT!}
\item
  声明位置：\passthrough{\lstinline!include/unitree/robot/h1/loco/h1\_loco\_client.hpp:168!}
\end{itemize}

\textbf{说明}

\begin{itemize}
\tightlist
\item
  示例是局部调用片段；\passthrough{\lstinline!obj!}
  和参数变量表示已经按上层业务校验并准备好的对象。
\item
  该条目可以被编辑器和类型检查器识别，但当前运行时可能在属性查找阶段直接失败。
\item
  该方法属于运动命令。方法名包含
  \passthrough{\lstinline!stop!}、\passthrough{\lstinline!damp!} 或
  \passthrough{\lstinline!zero!} 也不代表它是独立物理急停。
\end{itemize}

\textbf{用法}

\begin{lstlisting}[language=Python]
from typing import TYPE_CHECKING

if TYPE_CHECKING:
    def planned_set_stand_height(obj: LocoClient, stand_height: float) -> int:
        return obj.set_stand_height(stand_height=stand_height)
\end{lstlisting}

\begin{quote}
{[}!CAUTION{]} 上面的 \passthrough{\lstinline!TYPE\_CHECKING!}
示例只表达计划调用形态。该分支在普通 Python
运行时为假，不代表当前扩展能执行此接口。
\end{quote}

\hypertarget{unitree-sdk2-cpp-robot-h1-lococlient-set-velocity-1}{%
\paragraph{\texorpdfstring{\texttt{unitree\_sdk2\_cpp.robot.h1.LocoClient.set\_velocity}}{unitree\_sdk2\_cpp.robot.h1.LocoClient.set\_velocity}}\label{unitree-sdk2-cpp-robot-h1-lococlient-set-velocity-1}}

设置 速度。具体副作用和安全边界见下方状态。

\textbf{签名}

\begin{lstlisting}[language=Python]
def set_velocity(self, vx: float, vy: float, omega: float, duration: float = ...) -> int
\end{lstlisting}

\textbf{可用性与安全性}

\textbf{\passthrough{\lstinline!SIGNATURE\_ONLY!} /
\passthrough{\lstinline!MOTION\_COMMAND!}}：仅用于补全和静态检查。当前二进制没有该
Python 方法；不得按可执行运动接口使用。

\textbf{参数}

\begin{longtable}[]{@{}
  >{\raggedright\arraybackslash}p{(\columnwidth - 6\tabcolsep) * \real{0.18}}
  >{\raggedright\arraybackslash}p{(\columnwidth - 6\tabcolsep) * \real{0.24}}
  >{\raggedright\arraybackslash}p{(\columnwidth - 6\tabcolsep) * \real{0.18}}
  >{\raggedright\arraybackslash}p{(\columnwidth - 6\tabcolsep) * \real{0.40}}@{}}
\toprule
名称 & Python 类型 & 默认值 & 含义 \\
\midrule
\endhead
\passthrough{\lstinline!vx!} & \passthrough{\lstinline!float!} & 必填 &
X 方向速度参数。坐标系、单位、符号和安全范围必须查目标型号运动协议。
对应 C++ 参数 \passthrough{\lstinline!vx: float!}。
底层为浮点数；类型本身不说明单位、坐标系或业务安全范围。 \\
\passthrough{\lstinline!vy!} & \passthrough{\lstinline!float!} & 必填 &
Y 方向速度参数。坐标系、单位、符号和安全范围必须查目标型号运动协议。
对应 C++ 参数 \passthrough{\lstinline!vy: float!}。
底层为浮点数；类型本身不说明单位、坐标系或业务安全范围。 \\
\passthrough{\lstinline!omega!} & \passthrough{\lstinline!float!} & 必填
& 角速度参数。旋转轴、单位、符号和安全范围必须查目标型号运动协议。 对应
C++ 参数 \passthrough{\lstinline!omega: float!}。
底层为浮点数；类型本身不说明单位、坐标系或业务安全范围。 \\
\passthrough{\lstinline!duration!} & \passthrough{\lstinline!float!} &
\passthrough{\lstinline!...!} &
操作持续时间参数。精确单位、范围和默认行为以目标型号接口定义为准。 对应
C++ 参数 \passthrough{\lstinline!duration: float!}。
底层为浮点数；类型本身不说明单位、坐标系或业务安全范围。 \\
\bottomrule
\end{longtable}

\textbf{返回值}

\begin{longtable}[]{@{}
  >{\raggedright\arraybackslash}p{(\columnwidth - 4\tabcolsep) * \real{0.18}}
  >{\raggedright\arraybackslash}p{(\columnwidth - 4\tabcolsep) * \real{0.20}}
  >{\raggedright\arraybackslash}p{(\columnwidth - 4\tabcolsep) * \real{0.62}}@{}}
\toprule
位置 & 类型 & 含义 \\
\midrule
\endhead
返回值 & \passthrough{\lstinline!int!} & SDK 状态码；Unitree 示例通常以
\passthrough{\lstinline!0!} 表示成功，非零值需按具体服务定义解释。 \\
\bottomrule
\end{longtable}

\textbf{对应 C++}

\begin{itemize}
\tightlist
\item
  类：\passthrough{\lstinline!unitree::robot::h1::LocoClient!}
\item
  签名：\passthrough{\lstinline!SetVelocity(float, float, float, float)!}
\item
  绑定策略：\passthrough{\lstinline!DIRECT!}
\item
  声明位置：\passthrough{\lstinline!include/unitree/robot/h1/loco/h1\_loco\_client.hpp:178!}
\end{itemize}

\textbf{说明}

\begin{itemize}
\tightlist
\item
  示例是局部调用片段；\passthrough{\lstinline!obj!}
  和参数变量表示已经按上层业务校验并准备好的对象。
\item
  默认值显示为 \passthrough{\lstinline!...!}，表示 C++
  声明存在默认参数，但当前 AST
  清单没有保存其字面量；省略参数可使用上游默认行为。
\item
  该条目可以被编辑器和类型检查器识别，但当前运行时可能在属性查找阶段直接失败。
\item
  该方法属于运动命令。方法名包含
  \passthrough{\lstinline!stop!}、\passthrough{\lstinline!damp!} 或
  \passthrough{\lstinline!zero!} 也不代表它是独立物理急停。
\end{itemize}

\textbf{用法}

\begin{lstlisting}[language=Python]
from typing import TYPE_CHECKING

if TYPE_CHECKING:
    def planned_set_velocity(obj: LocoClient, vx: float, vy: float, omega: float, duration: float) -> int:
        return obj.set_velocity(vx=vx, vy=vy, omega=omega, duration=duration)
\end{lstlisting}

\begin{quote}
{[}!CAUTION{]} 上面的 \passthrough{\lstinline!TYPE\_CHECKING!}
示例只表达计划调用形态。该分支在普通 Python
运行时为假，不代表当前扩展能执行此接口。
\end{quote}

\hypertarget{unitree-sdk2-cpp-robot-h1-lococlient-set-phase-1}{%
\paragraph{\texorpdfstring{\texttt{unitree\_sdk2\_cpp.robot.h1.LocoClient.set\_phase}}{unitree\_sdk2\_cpp.robot.h1.LocoClient.set\_phase}}\label{unitree-sdk2-cpp-robot-h1-lococlient-set-phase-1}}

设置 相位。具体副作用和安全边界见下方状态。

\textbf{签名}

\begin{lstlisting}[language=Python]
def set_phase(self, phase: Sequence[float]) -> int
\end{lstlisting}

\textbf{可用性与安全性}

\textbf{\passthrough{\lstinline!SIGNATURE\_ONLY!} /
\passthrough{\lstinline!MOTION\_COMMAND!}}：仅用于补全和静态检查。当前二进制没有该
Python 方法；不得按可执行运动接口使用。

\textbf{参数}

\begin{longtable}[]{@{}
  >{\raggedright\arraybackslash}p{(\columnwidth - 6\tabcolsep) * \real{0.18}}
  >{\raggedright\arraybackslash}p{(\columnwidth - 6\tabcolsep) * \real{0.24}}
  >{\raggedright\arraybackslash}p{(\columnwidth - 6\tabcolsep) * \real{0.18}}
  >{\raggedright\arraybackslash}p{(\columnwidth - 6\tabcolsep) * \real{0.40}}@{}}
\toprule
名称 & Python 类型 & 默认值 & 含义 \\
\midrule
\endhead
\passthrough{\lstinline!phase!} &
\passthrough{\lstinline!Sequence[float]!} & 必填 & 传给该接口的 相位
参数，Python 类型为
\passthrough{\lstinline!Sequence[float]!}。精确含义、单位、范围和枚举值需查目标型号对应头文件与协议。
对应 C++ 参数 \passthrough{\lstinline!phase: std::vector<float>!}。
底层是可变长度
vector；元素约束：底层为浮点数；类型本身不说明单位、坐标系或业务安全范围。 \\
\bottomrule
\end{longtable}

\textbf{返回值}

\begin{longtable}[]{@{}
  >{\raggedright\arraybackslash}p{(\columnwidth - 4\tabcolsep) * \real{0.18}}
  >{\raggedright\arraybackslash}p{(\columnwidth - 4\tabcolsep) * \real{0.20}}
  >{\raggedright\arraybackslash}p{(\columnwidth - 4\tabcolsep) * \real{0.62}}@{}}
\toprule
位置 & 类型 & 含义 \\
\midrule
\endhead
返回值 & \passthrough{\lstinline!int!} & SDK 状态码；Unitree 示例通常以
\passthrough{\lstinline!0!} 表示成功，非零值需按具体服务定义解释。 \\
\bottomrule
\end{longtable}

\textbf{对应 C++}

\begin{itemize}
\tightlist
\item
  类：\passthrough{\lstinline!unitree::robot::h1::LocoClient!}
\item
  签名：\passthrough{\lstinline!SetPhase(std::vector<float>)!}
\item
  绑定策略：\passthrough{\lstinline!DIRECT!}
\item
  声明位置：\passthrough{\lstinline!include/unitree/robot/h1/loco/h1\_loco\_client.hpp:190!}
\end{itemize}

\textbf{说明}

\begin{itemize}
\tightlist
\item
  示例是局部调用片段；\passthrough{\lstinline!obj!}
  和参数变量表示已经按上层业务校验并准备好的对象。
\item
  该条目可以被编辑器和类型检查器识别，但当前运行时可能在属性查找阶段直接失败。
\item
  该方法属于运动命令。方法名包含
  \passthrough{\lstinline!stop!}、\passthrough{\lstinline!damp!} 或
  \passthrough{\lstinline!zero!} 也不代表它是独立物理急停。
\end{itemize}

\textbf{用法}

\begin{lstlisting}[language=Python]
from typing import TYPE_CHECKING

if TYPE_CHECKING:
    def planned_set_phase(obj: LocoClient, phase: Sequence[float]) -> int:
        return obj.set_phase(phase=phase)
\end{lstlisting}

\begin{quote}
{[}!CAUTION{]} 上面的 \passthrough{\lstinline!TYPE\_CHECKING!}
示例只表达计划调用形态。该分支在普通 Python
运行时为假，不代表当前扩展能执行此接口。
\end{quote}

\hypertarget{unitree-sdk2-cpp-robot-h1-lococlient-set-target-pos-1}{%
\paragraph{\texorpdfstring{\texttt{unitree\_sdk2\_cpp.robot.h1.LocoClient.set\_target\_pos}}{unitree\_sdk2\_cpp.robot.h1.LocoClient.set\_target\_pos}}\label{unitree-sdk2-cpp-robot-h1-lococlient-set-target-pos-1}}

设置 \passthrough{\lstinline!target!}
\passthrough{\lstinline!pos!}。具体副作用和安全边界见下方状态。

\textbf{签名}

\begin{lstlisting}[language=Python]
def set_target_pos(self, x: float, y: float, yaw: float, relative: bool = ...) -> int
\end{lstlisting}

\textbf{可用性与安全性}

\textbf{\passthrough{\lstinline!SIGNATURE\_ONLY!} /
\passthrough{\lstinline!MOTION\_COMMAND!}}：仅用于补全和静态检查。当前二进制没有该
Python 方法；不得按可执行运动接口使用。

\textbf{参数}

\begin{longtable}[]{@{}
  >{\raggedright\arraybackslash}p{(\columnwidth - 6\tabcolsep) * \real{0.18}}
  >{\raggedright\arraybackslash}p{(\columnwidth - 6\tabcolsep) * \real{0.24}}
  >{\raggedright\arraybackslash}p{(\columnwidth - 6\tabcolsep) * \real{0.18}}
  >{\raggedright\arraybackslash}p{(\columnwidth - 6\tabcolsep) * \real{0.40}}@{}}
\toprule
名称 & Python 类型 & 默认值 & 含义 \\
\midrule
\endhead
\passthrough{\lstinline!x!} & \passthrough{\lstinline!float!} & 必填 & X
方向位置或分量参数。参考系、单位和范围必须查当前方法及目标型号协议。
对应 C++ 参数 \passthrough{\lstinline!x: float!}。
底层为浮点数；类型本身不说明单位、坐标系或业务安全范围。 \\
\passthrough{\lstinline!y!} & \passthrough{\lstinline!float!} & 必填 & Y
方向位置或分量参数。参考系、单位和范围必须查当前方法及目标型号协议。
对应 C++ 参数 \passthrough{\lstinline!y: float!}。
底层为浮点数；类型本身不说明单位、坐标系或业务安全范围。 \\
\passthrough{\lstinline!yaw!} & \passthrough{\lstinline!float!} & 必填 &
偏航角或偏航目标参数。参考系、单位和范围必须查目标型号协议。 对应 C++
参数 \passthrough{\lstinline!yaw: float!}。
底层为浮点数；类型本身不说明单位、坐标系或业务安全范围。 \\
\passthrough{\lstinline!relative!} & \passthrough{\lstinline!bool!} &
\passthrough{\lstinline!...!} & 传给该接口的
\passthrough{\lstinline!relative!} 参数，Python 类型为
\passthrough{\lstinline!bool!}。精确含义、单位、范围和枚举值需查目标型号对应头文件与协议。
对应 C++ 参数 \passthrough{\lstinline!relative: bool!}。
只接受布尔语义。 \\
\bottomrule
\end{longtable}

\textbf{返回值}

\begin{longtable}[]{@{}
  >{\raggedright\arraybackslash}p{(\columnwidth - 4\tabcolsep) * \real{0.18}}
  >{\raggedright\arraybackslash}p{(\columnwidth - 4\tabcolsep) * \real{0.20}}
  >{\raggedright\arraybackslash}p{(\columnwidth - 4\tabcolsep) * \real{0.62}}@{}}
\toprule
位置 & 类型 & 含义 \\
\midrule
\endhead
返回值 & \passthrough{\lstinline!int!} & SDK 状态码；Unitree 示例通常以
\passthrough{\lstinline!0!} 表示成功，非零值需按具体服务定义解释。 \\
\bottomrule
\end{longtable}

\textbf{对应 C++}

\begin{itemize}
\tightlist
\item
  类：\passthrough{\lstinline!unitree::robot::h1::LocoClient!}
\item
  签名：\passthrough{\lstinline!SetTargetPos(float, float, float, bool)!}
\item
  绑定策略：\passthrough{\lstinline!DIRECT!}
\item
  声明位置：\passthrough{\lstinline!include/unitree/robot/h1/loco/h1\_loco\_client.hpp:200!}
\end{itemize}

\textbf{说明}

\begin{itemize}
\tightlist
\item
  示例是局部调用片段；\passthrough{\lstinline!obj!}
  和参数变量表示已经按上层业务校验并准备好的对象。
\item
  默认值显示为 \passthrough{\lstinline!...!}，表示 C++
  声明存在默认参数，但当前 AST
  清单没有保存其字面量；省略参数可使用上游默认行为。
\item
  该条目可以被编辑器和类型检查器识别，但当前运行时可能在属性查找阶段直接失败。
\item
  该方法属于运动命令。方法名包含
  \passthrough{\lstinline!stop!}、\passthrough{\lstinline!damp!} 或
  \passthrough{\lstinline!zero!} 也不代表它是独立物理急停。
\end{itemize}

\textbf{用法}

\begin{lstlisting}[language=Python]
from typing import TYPE_CHECKING

if TYPE_CHECKING:
    def planned_set_target_pos(obj: LocoClient, x: float, y: float, yaw: float, relative: bool) -> int:
        return obj.set_target_pos(x=x, y=y, yaw=yaw, relative=relative)
\end{lstlisting}

\begin{quote}
{[}!CAUTION{]} 上面的 \passthrough{\lstinline!TYPE\_CHECKING!}
示例只表达计划调用形态。该分支在普通 Python
运行时为假，不代表当前扩展能执行此接口。
\end{quote}

\hypertarget{unitree-sdk2-cpp-robot-h1-lococlient-set-task-id-1}{%
\paragraph{\texorpdfstring{\texttt{unitree\_sdk2\_cpp.robot.h1.LocoClient.set\_task\_id}}{unitree\_sdk2\_cpp.robot.h1.LocoClient.set\_task\_id}}\label{unitree-sdk2-cpp-robot-h1-lococlient-set-task-id-1}}

设置 任务 ID。具体副作用和安全边界见下方状态。

\textbf{签名}

\begin{lstlisting}[language=Python]
def set_task_id(self, task_id: int) -> int
\end{lstlisting}

\textbf{可用性与安全性}

\textbf{\passthrough{\lstinline!SIGNATURE\_ONLY!} /
\passthrough{\lstinline!MOTION\_COMMAND!}}：仅用于补全和静态检查。当前二进制没有该
Python 方法；不得按可执行运动接口使用。

\textbf{参数}

\begin{longtable}[]{@{}
  >{\raggedright\arraybackslash}p{(\columnwidth - 6\tabcolsep) * \real{0.18}}
  >{\raggedright\arraybackslash}p{(\columnwidth - 6\tabcolsep) * \real{0.24}}
  >{\raggedright\arraybackslash}p{(\columnwidth - 6\tabcolsep) * \real{0.18}}
  >{\raggedright\arraybackslash}p{(\columnwidth - 6\tabcolsep) * \real{0.40}}@{}}
\toprule
名称 & Python 类型 & 默认值 & 含义 \\
\midrule
\endhead
\passthrough{\lstinline!task\_id!} & \passthrough{\lstinline!int!} &
必填 & 传给该接口的 任务 ID 参数，Python 类型为
\passthrough{\lstinline!int!}。精确含义、单位、范围和枚举值需查目标型号对应头文件与协议。
对应 C++ 参数 \passthrough{\lstinline!task\_id: int!}。 \\
\bottomrule
\end{longtable}

\textbf{返回值}

\begin{longtable}[]{@{}
  >{\raggedright\arraybackslash}p{(\columnwidth - 4\tabcolsep) * \real{0.18}}
  >{\raggedright\arraybackslash}p{(\columnwidth - 4\tabcolsep) * \real{0.20}}
  >{\raggedright\arraybackslash}p{(\columnwidth - 4\tabcolsep) * \real{0.62}}@{}}
\toprule
位置 & 类型 & 含义 \\
\midrule
\endhead
返回值 & \passthrough{\lstinline!int!} & SDK 状态码；Unitree 示例通常以
\passthrough{\lstinline!0!} 表示成功，非零值需按具体服务定义解释。 \\
\bottomrule
\end{longtable}

\textbf{对应 C++}

\begin{itemize}
\tightlist
\item
  类：\passthrough{\lstinline!unitree::robot::h1::LocoClient!}
\item
  签名：\passthrough{\lstinline!SetTaskId(int)!}
\item
  绑定策略：\passthrough{\lstinline!DIRECT!}
\item
  声明位置：\passthrough{\lstinline!include/unitree/robot/h1/loco/h1\_loco\_client.hpp:213!}
\end{itemize}

\textbf{说明}

\begin{itemize}
\tightlist
\item
  示例是局部调用片段；\passthrough{\lstinline!obj!}
  和参数变量表示已经按上层业务校验并准备好的对象。
\item
  该条目可以被编辑器和类型检查器识别，但当前运行时可能在属性查找阶段直接失败。
\item
  该方法属于运动命令。方法名包含
  \passthrough{\lstinline!stop!}、\passthrough{\lstinline!damp!} 或
  \passthrough{\lstinline!zero!} 也不代表它是独立物理急停。
\end{itemize}

\textbf{用法}

\begin{lstlisting}[language=Python]
from typing import TYPE_CHECKING

if TYPE_CHECKING:
    def planned_set_task_id(obj: LocoClient, task_id: int) -> int:
        return obj.set_task_id(task_id=task_id)
\end{lstlisting}

\begin{quote}
{[}!CAUTION{]} 上面的 \passthrough{\lstinline!TYPE\_CHECKING!}
示例只表达计划调用形态。该分支在普通 Python
运行时为假，不代表当前扩展能执行此接口。
\end{quote}

\hypertarget{unitree-sdk2-cpp-robot-h1-lococlient-damp-1}{%
\paragraph{\texorpdfstring{\texttt{unitree\_sdk2\_cpp.robot.h1.LocoClient.damp}}{unitree\_sdk2\_cpp.robot.h1.LocoClient.damp}}\label{unitree-sdk2-cpp-robot-h1-lococlient-damp-1}}

对应 C++ SDK 操作
\passthrough{\lstinline!Damp()!}。上游头文件没有可直接生成的业务说明时，本参考只保证签名映射，精确语义需查目标型号协议。

\textbf{签名}

\begin{lstlisting}[language=Python]
def damp(self) -> int
\end{lstlisting}

\textbf{可用性与安全性}

\textbf{\passthrough{\lstinline!SIGNATURE\_ONLY!} /
\passthrough{\lstinline!MOTION\_COMMAND!}}：仅用于补全和静态检查。当前二进制没有该
Python 方法；不得按可执行运动接口使用。

\textbf{参数}

无显式参数。实例方法中的 \passthrough{\lstinline!self!} 由 Python
自动传入。

\textbf{返回值}

\begin{longtable}[]{@{}
  >{\raggedright\arraybackslash}p{(\columnwidth - 4\tabcolsep) * \real{0.18}}
  >{\raggedright\arraybackslash}p{(\columnwidth - 4\tabcolsep) * \real{0.20}}
  >{\raggedright\arraybackslash}p{(\columnwidth - 4\tabcolsep) * \real{0.62}}@{}}
\toprule
位置 & 类型 & 含义 \\
\midrule
\endhead
返回值 & \passthrough{\lstinline!int!} & SDK 状态码；Unitree 示例通常以
\passthrough{\lstinline!0!} 表示成功，非零值需按具体服务定义解释。 \\
\bottomrule
\end{longtable}

\textbf{对应 C++}

\begin{itemize}
\tightlist
\item
  类：\passthrough{\lstinline!unitree::robot::h1::LocoClient!}
\item
  签名：\passthrough{\lstinline!Damp()!}
\item
  绑定策略：\passthrough{\lstinline!DIRECT!}
\item
  声明位置：\passthrough{\lstinline!include/unitree/robot/h1/loco/h1\_loco\_client.hpp:224!}
\end{itemize}

\textbf{说明}

\begin{itemize}
\tightlist
\item
  示例是局部调用片段；\passthrough{\lstinline!obj!}
  和参数变量表示已经按上层业务校验并准备好的对象。
\item
  该条目可以被编辑器和类型检查器识别，但当前运行时可能在属性查找阶段直接失败。
\item
  该方法属于运动命令。方法名包含
  \passthrough{\lstinline!stop!}、\passthrough{\lstinline!damp!} 或
  \passthrough{\lstinline!zero!} 也不代表它是独立物理急停。
\end{itemize}

\textbf{用法}

\begin{lstlisting}[language=Python]
from typing import TYPE_CHECKING

if TYPE_CHECKING:
    def planned_damp(obj: LocoClient) -> int:
        return obj.damp()
\end{lstlisting}

\begin{quote}
{[}!CAUTION{]} 上面的 \passthrough{\lstinline!TYPE\_CHECKING!}
示例只表达计划调用形态。该分支在普通 Python
运行时为假，不代表当前扩展能执行此接口。
\end{quote}

\hypertarget{unitree-sdk2-cpp-robot-h1-lococlient-start-1}{%
\paragraph{\texorpdfstring{\texttt{unitree\_sdk2\_cpp.robot.h1.LocoClient.start}}{unitree\_sdk2\_cpp.robot.h1.LocoClient.start}}\label{unitree-sdk2-cpp-robot-h1-lococlient-start-1}}

启动对应 SDK 操作。具体副作用和安全边界见下方状态。

\textbf{签名}

\begin{lstlisting}[language=Python]
def start(self) -> int
\end{lstlisting}

\textbf{可用性与安全性}

\textbf{\passthrough{\lstinline!SIGNATURE\_ONLY!} /
\passthrough{\lstinline!MOTION\_COMMAND!}}：仅用于补全和静态检查。当前二进制没有该
Python 方法；不得按可执行运动接口使用。

\textbf{参数}

无显式参数。实例方法中的 \passthrough{\lstinline!self!} 由 Python
自动传入。

\textbf{返回值}

\begin{longtable}[]{@{}
  >{\raggedright\arraybackslash}p{(\columnwidth - 4\tabcolsep) * \real{0.18}}
  >{\raggedright\arraybackslash}p{(\columnwidth - 4\tabcolsep) * \real{0.20}}
  >{\raggedright\arraybackslash}p{(\columnwidth - 4\tabcolsep) * \real{0.62}}@{}}
\toprule
位置 & 类型 & 含义 \\
\midrule
\endhead
返回值 & \passthrough{\lstinline!int!} & SDK 状态码；Unitree 示例通常以
\passthrough{\lstinline!0!} 表示成功，非零值需按具体服务定义解释。 \\
\bottomrule
\end{longtable}

\textbf{对应 C++}

\begin{itemize}
\tightlist
\item
  类：\passthrough{\lstinline!unitree::robot::h1::LocoClient!}
\item
  签名：\passthrough{\lstinline!Start()!}
\item
  绑定策略：\passthrough{\lstinline!DIRECT!}
\item
  声明位置：\passthrough{\lstinline!include/unitree/robot/h1/loco/h1\_loco\_client.hpp:226!}
\end{itemize}

\textbf{说明}

\begin{itemize}
\tightlist
\item
  示例是局部调用片段；\passthrough{\lstinline!obj!}
  和参数变量表示已经按上层业务校验并准备好的对象。
\item
  该条目可以被编辑器和类型检查器识别，但当前运行时可能在属性查找阶段直接失败。
\item
  该方法属于运动命令。方法名包含
  \passthrough{\lstinline!stop!}、\passthrough{\lstinline!damp!} 或
  \passthrough{\lstinline!zero!} 也不代表它是独立物理急停。
\end{itemize}

\textbf{用法}

\begin{lstlisting}[language=Python]
from typing import TYPE_CHECKING

if TYPE_CHECKING:
    def planned_start(obj: LocoClient) -> int:
        return obj.start()
\end{lstlisting}

\begin{quote}
{[}!CAUTION{]} 上面的 \passthrough{\lstinline!TYPE\_CHECKING!}
示例只表达计划调用形态。该分支在普通 Python
运行时为假，不代表当前扩展能执行此接口。
\end{quote}

\hypertarget{unitree-sdk2-cpp-robot-h1-lococlient-stand-up-1}{%
\paragraph{\texorpdfstring{\texttt{unitree\_sdk2\_cpp.robot.h1.LocoClient.stand\_up}}{unitree\_sdk2\_cpp.robot.h1.LocoClient.stand\_up}}\label{unitree-sdk2-cpp-robot-h1-lococlient-stand-up-1}}

对应 C++ SDK 操作
\passthrough{\lstinline!StandUp()!}。上游头文件没有可直接生成的业务说明时，本参考只保证签名映射，精确语义需查目标型号协议。

\textbf{签名}

\begin{lstlisting}[language=Python]
def stand_up(self) -> int
\end{lstlisting}

\textbf{可用性与安全性}

\textbf{\passthrough{\lstinline!SIGNATURE\_ONLY!} /
\passthrough{\lstinline!MOTION\_COMMAND!}}：仅用于补全和静态检查。当前二进制没有该
Python 方法；不得按可执行运动接口使用。

\textbf{参数}

无显式参数。实例方法中的 \passthrough{\lstinline!self!} 由 Python
自动传入。

\textbf{返回值}

\begin{longtable}[]{@{}
  >{\raggedright\arraybackslash}p{(\columnwidth - 4\tabcolsep) * \real{0.18}}
  >{\raggedright\arraybackslash}p{(\columnwidth - 4\tabcolsep) * \real{0.20}}
  >{\raggedright\arraybackslash}p{(\columnwidth - 4\tabcolsep) * \real{0.62}}@{}}
\toprule
位置 & 类型 & 含义 \\
\midrule
\endhead
返回值 & \passthrough{\lstinline!int!} & SDK 状态码；Unitree 示例通常以
\passthrough{\lstinline!0!} 表示成功，非零值需按具体服务定义解释。 \\
\bottomrule
\end{longtable}

\textbf{对应 C++}

\begin{itemize}
\tightlist
\item
  类：\passthrough{\lstinline!unitree::robot::h1::LocoClient!}
\item
  签名：\passthrough{\lstinline!StandUp()!}
\item
  绑定策略：\passthrough{\lstinline!DIRECT!}
\item
  声明位置：\passthrough{\lstinline!include/unitree/robot/h1/loco/h1\_loco\_client.hpp:228!}
\end{itemize}

\textbf{说明}

\begin{itemize}
\tightlist
\item
  示例是局部调用片段；\passthrough{\lstinline!obj!}
  和参数变量表示已经按上层业务校验并准备好的对象。
\item
  该条目可以被编辑器和类型检查器识别，但当前运行时可能在属性查找阶段直接失败。
\item
  该方法属于运动命令。方法名包含
  \passthrough{\lstinline!stop!}、\passthrough{\lstinline!damp!} 或
  \passthrough{\lstinline!zero!} 也不代表它是独立物理急停。
\end{itemize}

\textbf{用法}

\begin{lstlisting}[language=Python]
from typing import TYPE_CHECKING

if TYPE_CHECKING:
    def planned_stand_up(obj: LocoClient) -> int:
        return obj.stand_up()
\end{lstlisting}

\begin{quote}
{[}!CAUTION{]} 上面的 \passthrough{\lstinline!TYPE\_CHECKING!}
示例只表达计划调用形态。该分支在普通 Python
运行时为假，不代表当前扩展能执行此接口。
\end{quote}

\hypertarget{unitree-sdk2-cpp-robot-h1-lococlient-zero-torque-1}{%
\paragraph{\texorpdfstring{\texttt{unitree\_sdk2\_cpp.robot.h1.LocoClient.zero\_torque}}{unitree\_sdk2\_cpp.robot.h1.LocoClient.zero\_torque}}\label{unitree-sdk2-cpp-robot-h1-lococlient-zero-torque-1}}

对应 C++ SDK 操作
\passthrough{\lstinline!ZeroTorque()!}。上游头文件没有可直接生成的业务说明时，本参考只保证签名映射，精确语义需查目标型号协议。

\textbf{签名}

\begin{lstlisting}[language=Python]
def zero_torque(self) -> int
\end{lstlisting}

\textbf{可用性与安全性}

\textbf{\passthrough{\lstinline!SIGNATURE\_ONLY!} /
\passthrough{\lstinline!MOTION\_COMMAND!}}：仅用于补全和静态检查。当前二进制没有该
Python 方法；不得按可执行运动接口使用。

\textbf{参数}

无显式参数。实例方法中的 \passthrough{\lstinline!self!} 由 Python
自动传入。

\textbf{返回值}

\begin{longtable}[]{@{}
  >{\raggedright\arraybackslash}p{(\columnwidth - 4\tabcolsep) * \real{0.18}}
  >{\raggedright\arraybackslash}p{(\columnwidth - 4\tabcolsep) * \real{0.20}}
  >{\raggedright\arraybackslash}p{(\columnwidth - 4\tabcolsep) * \real{0.62}}@{}}
\toprule
位置 & 类型 & 含义 \\
\midrule
\endhead
返回值 & \passthrough{\lstinline!int!} & SDK 状态码；Unitree 示例通常以
\passthrough{\lstinline!0!} 表示成功，非零值需按具体服务定义解释。 \\
\bottomrule
\end{longtable}

\textbf{对应 C++}

\begin{itemize}
\tightlist
\item
  类：\passthrough{\lstinline!unitree::robot::h1::LocoClient!}
\item
  签名：\passthrough{\lstinline!ZeroTorque()!}
\item
  绑定策略：\passthrough{\lstinline!DIRECT!}
\item
  声明位置：\passthrough{\lstinline!include/unitree/robot/h1/loco/h1\_loco\_client.hpp:230!}
\end{itemize}

\textbf{说明}

\begin{itemize}
\tightlist
\item
  示例是局部调用片段；\passthrough{\lstinline!obj!}
  和参数变量表示已经按上层业务校验并准备好的对象。
\item
  该条目可以被编辑器和类型检查器识别，但当前运行时可能在属性查找阶段直接失败。
\item
  该方法属于运动命令。方法名包含
  \passthrough{\lstinline!stop!}、\passthrough{\lstinline!damp!} 或
  \passthrough{\lstinline!zero!} 也不代表它是独立物理急停。
\end{itemize}

\textbf{用法}

\begin{lstlisting}[language=Python]
from typing import TYPE_CHECKING

if TYPE_CHECKING:
    def planned_zero_torque(obj: LocoClient) -> int:
        return obj.zero_torque()
\end{lstlisting}

\begin{quote}
{[}!CAUTION{]} 上面的 \passthrough{\lstinline!TYPE\_CHECKING!}
示例只表达计划调用形态。该分支在普通 Python
运行时为假，不代表当前扩展能执行此接口。
\end{quote}

\hypertarget{unitree-sdk2-cpp-robot-h1-lococlient-stop-move-1}{%
\paragraph{\texorpdfstring{\texttt{unitree\_sdk2\_cpp.robot.h1.LocoClient.stop\_move}}{unitree\_sdk2\_cpp.robot.h1.LocoClient.stop\_move}}\label{unitree-sdk2-cpp-robot-h1-lococlient-stop-move-1}}

请求停止对应 SDK 操作。该名称不等同于经过验证的物理急停。

\textbf{签名}

\begin{lstlisting}[language=Python]
def stop_move(self) -> int
\end{lstlisting}

\textbf{可用性与安全性}

\textbf{\passthrough{\lstinline!SIGNATURE\_ONLY!} /
\passthrough{\lstinline!MOTION\_COMMAND!}}：仅用于补全和静态检查。当前二进制没有该
Python 方法；不得按可执行运动接口使用。

\textbf{参数}

无显式参数。实例方法中的 \passthrough{\lstinline!self!} 由 Python
自动传入。

\textbf{返回值}

\begin{longtable}[]{@{}
  >{\raggedright\arraybackslash}p{(\columnwidth - 4\tabcolsep) * \real{0.18}}
  >{\raggedright\arraybackslash}p{(\columnwidth - 4\tabcolsep) * \real{0.20}}
  >{\raggedright\arraybackslash}p{(\columnwidth - 4\tabcolsep) * \real{0.62}}@{}}
\toprule
位置 & 类型 & 含义 \\
\midrule
\endhead
返回值 & \passthrough{\lstinline!int!} & SDK 状态码；Unitree 示例通常以
\passthrough{\lstinline!0!} 表示成功，非零值需按具体服务定义解释。 \\
\bottomrule
\end{longtable}

\textbf{对应 C++}

\begin{itemize}
\tightlist
\item
  类：\passthrough{\lstinline!unitree::robot::h1::LocoClient!}
\item
  签名：\passthrough{\lstinline!StopMove()!}
\item
  绑定策略：\passthrough{\lstinline!DIRECT!}
\item
  声明位置：\passthrough{\lstinline!include/unitree/robot/h1/loco/h1\_loco\_client.hpp:232!}
\end{itemize}

\textbf{说明}

\begin{itemize}
\tightlist
\item
  示例是局部调用片段；\passthrough{\lstinline!obj!}
  和参数变量表示已经按上层业务校验并准备好的对象。
\item
  该条目可以被编辑器和类型检查器识别，但当前运行时可能在属性查找阶段直接失败。
\item
  该方法属于运动命令。方法名包含
  \passthrough{\lstinline!stop!}、\passthrough{\lstinline!damp!} 或
  \passthrough{\lstinline!zero!} 也不代表它是独立物理急停。
\end{itemize}

\textbf{用法}

\begin{lstlisting}[language=Python]
from typing import TYPE_CHECKING

if TYPE_CHECKING:
    def planned_stop_move(obj: LocoClient) -> int:
        return obj.stop_move()
\end{lstlisting}

\begin{quote}
{[}!CAUTION{]} 上面的 \passthrough{\lstinline!TYPE\_CHECKING!}
示例只表达计划调用形态。该分支在普通 Python
运行时为假，不代表当前扩展能执行此接口。
\end{quote}

\hypertarget{unitree-sdk2-cpp-robot-h1-lococlient-high-stand-1}{%
\paragraph{\texorpdfstring{\texttt{unitree\_sdk2\_cpp.robot.h1.LocoClient.high\_stand}}{unitree\_sdk2\_cpp.robot.h1.LocoClient.high\_stand}}\label{unitree-sdk2-cpp-robot-h1-lococlient-high-stand-1}}

对应 C++ SDK 操作
\passthrough{\lstinline!HighStand()!}。上游头文件没有可直接生成的业务说明时，本参考只保证签名映射，精确语义需查目标型号协议。

\textbf{签名}

\begin{lstlisting}[language=Python]
def high_stand(self) -> int
\end{lstlisting}

\textbf{可用性与安全性}

\textbf{\passthrough{\lstinline!SIGNATURE\_ONLY!} /
\passthrough{\lstinline!MOTION\_COMMAND!}}：仅用于补全和静态检查。当前二进制没有该
Python 方法；不得按可执行运动接口使用。

\textbf{参数}

无显式参数。实例方法中的 \passthrough{\lstinline!self!} 由 Python
自动传入。

\textbf{返回值}

\begin{longtable}[]{@{}
  >{\raggedright\arraybackslash}p{(\columnwidth - 4\tabcolsep) * \real{0.18}}
  >{\raggedright\arraybackslash}p{(\columnwidth - 4\tabcolsep) * \real{0.20}}
  >{\raggedright\arraybackslash}p{(\columnwidth - 4\tabcolsep) * \real{0.62}}@{}}
\toprule
位置 & 类型 & 含义 \\
\midrule
\endhead
返回值 & \passthrough{\lstinline!int!} & SDK 状态码；Unitree 示例通常以
\passthrough{\lstinline!0!} 表示成功，非零值需按具体服务定义解释。 \\
\bottomrule
\end{longtable}

\textbf{对应 C++}

\begin{itemize}
\tightlist
\item
  类：\passthrough{\lstinline!unitree::robot::h1::LocoClient!}
\item
  签名：\passthrough{\lstinline!HighStand()!}
\item
  绑定策略：\passthrough{\lstinline!DIRECT!}
\item
  声明位置：\passthrough{\lstinline!include/unitree/robot/h1/loco/h1\_loco\_client.hpp:234!}
\end{itemize}

\textbf{说明}

\begin{itemize}
\tightlist
\item
  示例是局部调用片段；\passthrough{\lstinline!obj!}
  和参数变量表示已经按上层业务校验并准备好的对象。
\item
  该条目可以被编辑器和类型检查器识别，但当前运行时可能在属性查找阶段直接失败。
\item
  该方法属于运动命令。方法名包含
  \passthrough{\lstinline!stop!}、\passthrough{\lstinline!damp!} 或
  \passthrough{\lstinline!zero!} 也不代表它是独立物理急停。
\end{itemize}

\textbf{用法}

\begin{lstlisting}[language=Python]
from typing import TYPE_CHECKING

if TYPE_CHECKING:
    def planned_high_stand(obj: LocoClient) -> int:
        return obj.high_stand()
\end{lstlisting}

\begin{quote}
{[}!CAUTION{]} 上面的 \passthrough{\lstinline!TYPE\_CHECKING!}
示例只表达计划调用形态。该分支在普通 Python
运行时为假，不代表当前扩展能执行此接口。
\end{quote}

\hypertarget{unitree-sdk2-cpp-robot-h1-lococlient-low-stand-1}{%
\paragraph{\texorpdfstring{\texttt{unitree\_sdk2\_cpp.robot.h1.LocoClient.low\_stand}}{unitree\_sdk2\_cpp.robot.h1.LocoClient.low\_stand}}\label{unitree-sdk2-cpp-robot-h1-lococlient-low-stand-1}}

对应 C++ SDK 操作
\passthrough{\lstinline!LowStand()!}。上游头文件没有可直接生成的业务说明时，本参考只保证签名映射，精确语义需查目标型号协议。

\textbf{签名}

\begin{lstlisting}[language=Python]
def low_stand(self) -> int
\end{lstlisting}

\textbf{可用性与安全性}

\textbf{\passthrough{\lstinline!SIGNATURE\_ONLY!} /
\passthrough{\lstinline!MOTION\_COMMAND!}}：仅用于补全和静态检查。当前二进制没有该
Python 方法；不得按可执行运动接口使用。

\textbf{参数}

无显式参数。实例方法中的 \passthrough{\lstinline!self!} 由 Python
自动传入。

\textbf{返回值}

\begin{longtable}[]{@{}
  >{\raggedright\arraybackslash}p{(\columnwidth - 4\tabcolsep) * \real{0.18}}
  >{\raggedright\arraybackslash}p{(\columnwidth - 4\tabcolsep) * \real{0.20}}
  >{\raggedright\arraybackslash}p{(\columnwidth - 4\tabcolsep) * \real{0.62}}@{}}
\toprule
位置 & 类型 & 含义 \\
\midrule
\endhead
返回值 & \passthrough{\lstinline!int!} & SDK 状态码；Unitree 示例通常以
\passthrough{\lstinline!0!} 表示成功，非零值需按具体服务定义解释。 \\
\bottomrule
\end{longtable}

\textbf{对应 C++}

\begin{itemize}
\tightlist
\item
  类：\passthrough{\lstinline!unitree::robot::h1::LocoClient!}
\item
  签名：\passthrough{\lstinline!LowStand()!}
\item
  绑定策略：\passthrough{\lstinline!DIRECT!}
\item
  声明位置：\passthrough{\lstinline!include/unitree/robot/h1/loco/h1\_loco\_client.hpp:236!}
\end{itemize}

\textbf{说明}

\begin{itemize}
\tightlist
\item
  示例是局部调用片段；\passthrough{\lstinline!obj!}
  和参数变量表示已经按上层业务校验并准备好的对象。
\item
  该条目可以被编辑器和类型检查器识别，但当前运行时可能在属性查找阶段直接失败。
\item
  该方法属于运动命令。方法名包含
  \passthrough{\lstinline!stop!}、\passthrough{\lstinline!damp!} 或
  \passthrough{\lstinline!zero!} 也不代表它是独立物理急停。
\end{itemize}

\textbf{用法}

\begin{lstlisting}[language=Python]
from typing import TYPE_CHECKING

if TYPE_CHECKING:
    def planned_low_stand(obj: LocoClient) -> int:
        return obj.low_stand()
\end{lstlisting}

\begin{quote}
{[}!CAUTION{]} 上面的 \passthrough{\lstinline!TYPE\_CHECKING!}
示例只表达计划调用形态。该分支在普通 Python
运行时为假，不代表当前扩展能执行此接口。
\end{quote}

\hypertarget{unitree-sdk2-cpp-robot-h1-lococlient-move-1}{%
\paragraph{\texorpdfstring{\texttt{unitree\_sdk2\_cpp.robot.h1.LocoClient.move}（重载
1/2）}{unitree\_sdk2\_cpp.robot.h1.LocoClient.move（重载 1/2）}}\label{unitree-sdk2-cpp-robot-h1-lococlient-move-1}}

对应 C++ SDK 操作
\passthrough{\lstinline!Move(float, float, float, bool)!}。上游头文件没有可直接生成的业务说明时，本参考只保证签名映射，精确语义需查目标型号协议。

\textbf{签名}

\begin{lstlisting}[language=Python]
@overload
def move(self, vx: float, vy: float, vyaw: float, continous_move: bool) -> int
\end{lstlisting}

\textbf{可用性与安全性}

\textbf{\passthrough{\lstinline!SIGNATURE\_ONLY!} /
\passthrough{\lstinline!MOTION\_COMMAND!}}：仅用于补全和静态检查。当前二进制没有该
Python 方法；不得按可执行运动接口使用。

\textbf{参数}

\begin{longtable}[]{@{}
  >{\raggedright\arraybackslash}p{(\columnwidth - 6\tabcolsep) * \real{0.18}}
  >{\raggedright\arraybackslash}p{(\columnwidth - 6\tabcolsep) * \real{0.24}}
  >{\raggedright\arraybackslash}p{(\columnwidth - 6\tabcolsep) * \real{0.18}}
  >{\raggedright\arraybackslash}p{(\columnwidth - 6\tabcolsep) * \real{0.40}}@{}}
\toprule
名称 & Python 类型 & 默认值 & 含义 \\
\midrule
\endhead
\passthrough{\lstinline!vx!} & \passthrough{\lstinline!float!} & 必填 &
X 方向速度参数。坐标系、单位、符号和安全范围必须查目标型号运动协议。
对应 C++ 参数 \passthrough{\lstinline!vx: float!}。
底层为浮点数；类型本身不说明单位、坐标系或业务安全范围。 \\
\passthrough{\lstinline!vy!} & \passthrough{\lstinline!float!} & 必填 &
Y 方向速度参数。坐标系、单位、符号和安全范围必须查目标型号运动协议。
对应 C++ 参数 \passthrough{\lstinline!vy: float!}。
底层为浮点数；类型本身不说明单位、坐标系或业务安全范围。 \\
\passthrough{\lstinline!vyaw!} & \passthrough{\lstinline!float!} & 必填
& 偏航角速度参数。单位、符号和安全范围必须查目标型号运动协议。 对应 C++
参数 \passthrough{\lstinline!vyaw: float!}。
底层为浮点数；类型本身不说明单位、坐标系或业务安全范围。 \\
\passthrough{\lstinline!continous\_move!} &
\passthrough{\lstinline!bool!} & 必填 & 传给该接口的
\passthrough{\lstinline!continous!} \passthrough{\lstinline!move!}
参数，Python 类型为
\passthrough{\lstinline!bool!}。精确含义、单位、范围和枚举值需查目标型号对应头文件与协议。
对应 C++ 参数 \passthrough{\lstinline!continous\_move: bool!}。
只接受布尔语义。 \\
\bottomrule
\end{longtable}

\textbf{返回值}

\begin{longtable}[]{@{}
  >{\raggedright\arraybackslash}p{(\columnwidth - 4\tabcolsep) * \real{0.18}}
  >{\raggedright\arraybackslash}p{(\columnwidth - 4\tabcolsep) * \real{0.20}}
  >{\raggedright\arraybackslash}p{(\columnwidth - 4\tabcolsep) * \real{0.62}}@{}}
\toprule
位置 & 类型 & 含义 \\
\midrule
\endhead
返回值 & \passthrough{\lstinline!int!} & SDK 状态码；Unitree 示例通常以
\passthrough{\lstinline!0!} 表示成功，非零值需按具体服务定义解释。 \\
\bottomrule
\end{longtable}

\textbf{对应 C++}

\begin{itemize}
\tightlist
\item
  类：\passthrough{\lstinline!unitree::robot::h1::LocoClient!}
\item
  签名：\passthrough{\lstinline!Move(float, float, float, bool)!}
\item
  绑定策略：\passthrough{\lstinline!DIRECT!}
\item
  声明位置：\passthrough{\lstinline!include/unitree/robot/h1/loco/h1\_loco\_client.hpp:238!}
\end{itemize}

\textbf{说明}

\begin{itemize}
\tightlist
\item
  示例是局部调用片段；\passthrough{\lstinline!obj!}
  和参数变量表示已经按上层业务校验并准备好的对象。
\item
  该条目可以被编辑器和类型检查器识别，但当前运行时可能在属性查找阶段直接失败。
\item
  该方法属于运动命令。方法名包含
  \passthrough{\lstinline!stop!}、\passthrough{\lstinline!damp!} 或
  \passthrough{\lstinline!zero!} 也不代表它是独立物理急停。
\end{itemize}

\textbf{用法}

\begin{lstlisting}[language=Python]
from typing import TYPE_CHECKING

if TYPE_CHECKING:
    def planned_move(obj: LocoClient, vx: float, vy: float, vyaw: float, continous_move: bool) -> int:
        return obj.move(vx=vx, vy=vy, vyaw=vyaw, continous_move=continous_move)
\end{lstlisting}

\begin{quote}
{[}!CAUTION{]} 上面的 \passthrough{\lstinline!TYPE\_CHECKING!}
示例只表达计划调用形态。该分支在普通 Python
运行时为假，不代表当前扩展能执行此接口。
\end{quote}

\hypertarget{unitree-sdk2-cpp-robot-h1-lococlient-move-2}{%
\paragraph{\texorpdfstring{\texttt{unitree\_sdk2\_cpp.robot.h1.LocoClient.move}（重载
2/2）}{unitree\_sdk2\_cpp.robot.h1.LocoClient.move（重载 2/2）}}\label{unitree-sdk2-cpp-robot-h1-lococlient-move-2}}

对应 C++ SDK 操作
\passthrough{\lstinline!Move(float, float, float)!}。上游头文件没有可直接生成的业务说明时，本参考只保证签名映射，精确语义需查目标型号协议。

\textbf{签名}

\begin{lstlisting}[language=Python]
@overload
def move(self, vx: float, vy: float, vyaw: float) -> int
\end{lstlisting}

\textbf{可用性与安全性}

\textbf{\passthrough{\lstinline!SIGNATURE\_ONLY!} /
\passthrough{\lstinline!MOTION\_COMMAND!}}：仅用于补全和静态检查。当前二进制没有该
Python 方法；不得按可执行运动接口使用。

\textbf{参数}

\begin{longtable}[]{@{}
  >{\raggedright\arraybackslash}p{(\columnwidth - 6\tabcolsep) * \real{0.18}}
  >{\raggedright\arraybackslash}p{(\columnwidth - 6\tabcolsep) * \real{0.24}}
  >{\raggedright\arraybackslash}p{(\columnwidth - 6\tabcolsep) * \real{0.18}}
  >{\raggedright\arraybackslash}p{(\columnwidth - 6\tabcolsep) * \real{0.40}}@{}}
\toprule
名称 & Python 类型 & 默认值 & 含义 \\
\midrule
\endhead
\passthrough{\lstinline!vx!} & \passthrough{\lstinline!float!} & 必填 &
X 方向速度参数。坐标系、单位、符号和安全范围必须查目标型号运动协议。
对应 C++ 参数 \passthrough{\lstinline!vx: float!}。
底层为浮点数；类型本身不说明单位、坐标系或业务安全范围。 \\
\passthrough{\lstinline!vy!} & \passthrough{\lstinline!float!} & 必填 &
Y 方向速度参数。坐标系、单位、符号和安全范围必须查目标型号运动协议。
对应 C++ 参数 \passthrough{\lstinline!vy: float!}。
底层为浮点数；类型本身不说明单位、坐标系或业务安全范围。 \\
\passthrough{\lstinline!vyaw!} & \passthrough{\lstinline!float!} & 必填
& 偏航角速度参数。单位、符号和安全范围必须查目标型号运动协议。 对应 C++
参数 \passthrough{\lstinline!vyaw: float!}。
底层为浮点数；类型本身不说明单位、坐标系或业务安全范围。 \\
\bottomrule
\end{longtable}

\textbf{返回值}

\begin{longtable}[]{@{}
  >{\raggedright\arraybackslash}p{(\columnwidth - 4\tabcolsep) * \real{0.18}}
  >{\raggedright\arraybackslash}p{(\columnwidth - 4\tabcolsep) * \real{0.20}}
  >{\raggedright\arraybackslash}p{(\columnwidth - 4\tabcolsep) * \real{0.62}}@{}}
\toprule
位置 & 类型 & 含义 \\
\midrule
\endhead
返回值 & \passthrough{\lstinline!int!} & SDK 状态码；Unitree 示例通常以
\passthrough{\lstinline!0!} 表示成功，非零值需按具体服务定义解释。 \\
\bottomrule
\end{longtable}

\textbf{对应 C++}

\begin{itemize}
\tightlist
\item
  类：\passthrough{\lstinline!unitree::robot::h1::LocoClient!}
\item
  签名：\passthrough{\lstinline!Move(float, float, float)!}
\item
  绑定策略：\passthrough{\lstinline!DIRECT!}
\item
  声明位置：\passthrough{\lstinline!include/unitree/robot/h1/loco/h1\_loco\_client.hpp:242!}
\end{itemize}

\textbf{说明}

\begin{itemize}
\tightlist
\item
  示例是局部调用片段；\passthrough{\lstinline!obj!}
  和参数变量表示已经按上层业务校验并准备好的对象。
\item
  该条目可以被编辑器和类型检查器识别，但当前运行时可能在属性查找阶段直接失败。
\item
  该方法属于运动命令。方法名包含
  \passthrough{\lstinline!stop!}、\passthrough{\lstinline!damp!} 或
  \passthrough{\lstinline!zero!} 也不代表它是独立物理急停。
\end{itemize}

\textbf{用法}

\begin{lstlisting}[language=Python]
from typing import TYPE_CHECKING

if TYPE_CHECKING:
    def planned_move(obj: LocoClient, vx: float, vy: float, vyaw: float) -> int:
        return obj.move(vx=vx, vy=vy, vyaw=vyaw)
\end{lstlisting}

\begin{quote}
{[}!CAUTION{]} 上面的 \passthrough{\lstinline!TYPE\_CHECKING!}
示例只表达计划调用形态。该分支在普通 Python
运行时为假，不代表当前扩展能执行此接口。
\end{quote}

\hypertarget{unitree-sdk2-cpp-robot-h1-lococlient-balance-stand-1}{%
\paragraph{\texorpdfstring{\texttt{unitree\_sdk2\_cpp.robot.h1.LocoClient.balance\_stand}}{unitree\_sdk2\_cpp.robot.h1.LocoClient.balance\_stand}}\label{unitree-sdk2-cpp-robot-h1-lococlient-balance-stand-1}}

对应 C++ SDK 操作
\passthrough{\lstinline!BalanceStand()!}。上游头文件没有可直接生成的业务说明时，本参考只保证签名映射，精确语义需查目标型号协议。

\textbf{签名}

\begin{lstlisting}[language=Python]
def balance_stand(self) -> int
\end{lstlisting}

\textbf{可用性与安全性}

\textbf{\passthrough{\lstinline!SIGNATURE\_ONLY!} /
\passthrough{\lstinline!MOTION\_COMMAND!}}：仅用于补全和静态检查。当前二进制没有该
Python 方法；不得按可执行运动接口使用。

\textbf{参数}

无显式参数。实例方法中的 \passthrough{\lstinline!self!} 由 Python
自动传入。

\textbf{返回值}

\begin{longtable}[]{@{}
  >{\raggedright\arraybackslash}p{(\columnwidth - 4\tabcolsep) * \real{0.18}}
  >{\raggedright\arraybackslash}p{(\columnwidth - 4\tabcolsep) * \real{0.20}}
  >{\raggedright\arraybackslash}p{(\columnwidth - 4\tabcolsep) * \real{0.62}}@{}}
\toprule
位置 & 类型 & 含义 \\
\midrule
\endhead
返回值 & \passthrough{\lstinline!int!} & SDK 状态码；Unitree 示例通常以
\passthrough{\lstinline!0!} 表示成功，非零值需按具体服务定义解释。 \\
\bottomrule
\end{longtable}

\textbf{对应 C++}

\begin{itemize}
\tightlist
\item
  类：\passthrough{\lstinline!unitree::robot::h1::LocoClient!}
\item
  签名：\passthrough{\lstinline!BalanceStand()!}
\item
  绑定策略：\passthrough{\lstinline!DIRECT!}
\item
  声明位置：\passthrough{\lstinline!include/unitree/robot/h1/loco/h1\_loco\_client.hpp:244!}
\end{itemize}

\textbf{说明}

\begin{itemize}
\tightlist
\item
  示例是局部调用片段；\passthrough{\lstinline!obj!}
  和参数变量表示已经按上层业务校验并准备好的对象。
\item
  该条目可以被编辑器和类型检查器识别，但当前运行时可能在属性查找阶段直接失败。
\item
  该方法属于运动命令。方法名包含
  \passthrough{\lstinline!stop!}、\passthrough{\lstinline!damp!} 或
  \passthrough{\lstinline!zero!} 也不代表它是独立物理急停。
\end{itemize}

\textbf{用法}

\begin{lstlisting}[language=Python]
from typing import TYPE_CHECKING

if TYPE_CHECKING:
    def planned_balance_stand(obj: LocoClient) -> int:
        return obj.balance_stand()
\end{lstlisting}

\begin{quote}
{[}!CAUTION{]} 上面的 \passthrough{\lstinline!TYPE\_CHECKING!}
示例只表达计划调用形态。该分支在普通 Python
运行时为假，不代表当前扩展能执行此接口。
\end{quote}

\hypertarget{unitree-sdk2-cpp-robot-h1-lococlient-continuous-gait-1}{%
\paragraph{\texorpdfstring{\texttt{unitree\_sdk2\_cpp.robot.h1.LocoClient.continuous\_gait}}{unitree\_sdk2\_cpp.robot.h1.LocoClient.continuous\_gait}}\label{unitree-sdk2-cpp-robot-h1-lococlient-continuous-gait-1}}

对应 C++ SDK 操作
\passthrough{\lstinline!ContinuousGait(bool)!}。上游头文件没有可直接生成的业务说明时，本参考只保证签名映射，精确语义需查目标型号协议。

\textbf{签名}

\begin{lstlisting}[language=Python]
def continuous_gait(self, flag: bool) -> int
\end{lstlisting}

\textbf{可用性与安全性}

\textbf{\passthrough{\lstinline!SIGNATURE\_ONLY!} /
\passthrough{\lstinline!MOTION\_COMMAND!}}：仅用于补全和静态检查。当前二进制没有该
Python 方法；不得按可执行运动接口使用。

\textbf{参数}

\begin{longtable}[]{@{}
  >{\raggedright\arraybackslash}p{(\columnwidth - 6\tabcolsep) * \real{0.18}}
  >{\raggedright\arraybackslash}p{(\columnwidth - 6\tabcolsep) * \real{0.24}}
  >{\raggedright\arraybackslash}p{(\columnwidth - 6\tabcolsep) * \real{0.18}}
  >{\raggedright\arraybackslash}p{(\columnwidth - 6\tabcolsep) * \real{0.40}}@{}}
\toprule
名称 & Python 类型 & 默认值 & 含义 \\
\midrule
\endhead
\passthrough{\lstinline!flag!} & \passthrough{\lstinline!bool!} & 必填 &
布尔开关。\passthrough{\lstinline!True!} 和
\passthrough{\lstinline!False!} 的具体业务效果由当前方法定义。 对应 C++
参数 \passthrough{\lstinline!flag: bool!}。 只接受布尔语义。 \\
\bottomrule
\end{longtable}

\textbf{返回值}

\begin{longtable}[]{@{}
  >{\raggedright\arraybackslash}p{(\columnwidth - 4\tabcolsep) * \real{0.18}}
  >{\raggedright\arraybackslash}p{(\columnwidth - 4\tabcolsep) * \real{0.20}}
  >{\raggedright\arraybackslash}p{(\columnwidth - 4\tabcolsep) * \real{0.62}}@{}}
\toprule
位置 & 类型 & 含义 \\
\midrule
\endhead
返回值 & \passthrough{\lstinline!int!} & SDK 状态码；Unitree 示例通常以
\passthrough{\lstinline!0!} 表示成功，非零值需按具体服务定义解释。 \\
\bottomrule
\end{longtable}

\textbf{对应 C++}

\begin{itemize}
\tightlist
\item
  类：\passthrough{\lstinline!unitree::robot::h1::LocoClient!}
\item
  签名：\passthrough{\lstinline!ContinuousGait(bool)!}
\item
  绑定策略：\passthrough{\lstinline!DIRECT!}
\item
  声明位置：\passthrough{\lstinline!include/unitree/robot/h1/loco/h1\_loco\_client.hpp:246!}
\end{itemize}

\textbf{说明}

\begin{itemize}
\tightlist
\item
  示例是局部调用片段；\passthrough{\lstinline!obj!}
  和参数变量表示已经按上层业务校验并准备好的对象。
\item
  该条目可以被编辑器和类型检查器识别，但当前运行时可能在属性查找阶段直接失败。
\item
  该方法属于运动命令。方法名包含
  \passthrough{\lstinline!stop!}、\passthrough{\lstinline!damp!} 或
  \passthrough{\lstinline!zero!} 也不代表它是独立物理急停。
\end{itemize}

\textbf{用法}

\begin{lstlisting}[language=Python]
from typing import TYPE_CHECKING

if TYPE_CHECKING:
    def planned_continuous_gait(obj: LocoClient, flag: bool) -> int:
        return obj.continuous_gait(flag=flag)
\end{lstlisting}

\begin{quote}
{[}!CAUTION{]} 上面的 \passthrough{\lstinline!TYPE\_CHECKING!}
示例只表达计划调用形态。该分支在普通 Python
运行时为假，不代表当前扩展能执行此接口。
\end{quote}

\hypertarget{unitree-sdk2-cpp-robot-h1-lococlient-switch-move-mode-1}{%
\paragraph{\texorpdfstring{\texttt{unitree\_sdk2\_cpp.robot.h1.LocoClient.switch\_move\_mode}}{unitree\_sdk2\_cpp.robot.h1.LocoClient.switch\_move\_mode}}\label{unitree-sdk2-cpp-robot-h1-lococlient-switch-move-mode-1}}

对应 C++ SDK 操作
\passthrough{\lstinline!SwitchMoveMode(bool)!}。上游头文件没有可直接生成的业务说明时，本参考只保证签名映射，精确语义需查目标型号协议。

\textbf{签名}

\begin{lstlisting}[language=Python]
def switch_move_mode(self, flag: bool) -> int
\end{lstlisting}

\textbf{可用性与安全性}

\textbf{\passthrough{\lstinline!SIGNATURE\_ONLY!} /
\passthrough{\lstinline!MOTION\_COMMAND!}}：仅用于补全和静态检查。当前二进制没有该
Python 方法；不得按可执行运动接口使用。

\textbf{参数}

\begin{longtable}[]{@{}
  >{\raggedright\arraybackslash}p{(\columnwidth - 6\tabcolsep) * \real{0.18}}
  >{\raggedright\arraybackslash}p{(\columnwidth - 6\tabcolsep) * \real{0.24}}
  >{\raggedright\arraybackslash}p{(\columnwidth - 6\tabcolsep) * \real{0.18}}
  >{\raggedright\arraybackslash}p{(\columnwidth - 6\tabcolsep) * \real{0.40}}@{}}
\toprule
名称 & Python 类型 & 默认值 & 含义 \\
\midrule
\endhead
\passthrough{\lstinline!flag!} & \passthrough{\lstinline!bool!} & 必填 &
布尔开关。\passthrough{\lstinline!True!} 和
\passthrough{\lstinline!False!} 的具体业务效果由当前方法定义。 对应 C++
参数 \passthrough{\lstinline!flag: bool!}。 只接受布尔语义。 \\
\bottomrule
\end{longtable}

\textbf{返回值}

\begin{longtable}[]{@{}
  >{\raggedright\arraybackslash}p{(\columnwidth - 4\tabcolsep) * \real{0.18}}
  >{\raggedright\arraybackslash}p{(\columnwidth - 4\tabcolsep) * \real{0.20}}
  >{\raggedright\arraybackslash}p{(\columnwidth - 4\tabcolsep) * \real{0.62}}@{}}
\toprule
位置 & 类型 & 含义 \\
\midrule
\endhead
返回值 & \passthrough{\lstinline!int!} & SDK 状态码；Unitree 示例通常以
\passthrough{\lstinline!0!} 表示成功，非零值需按具体服务定义解释。 \\
\bottomrule
\end{longtable}

\textbf{对应 C++}

\begin{itemize}
\tightlist
\item
  类：\passthrough{\lstinline!unitree::robot::h1::LocoClient!}
\item
  签名：\passthrough{\lstinline!SwitchMoveMode(bool)!}
\item
  绑定策略：\passthrough{\lstinline!DIRECT!}
\item
  声明位置：\passthrough{\lstinline!include/unitree/robot/h1/loco/h1\_loco\_client.hpp:248!}
\end{itemize}

\textbf{说明}

\begin{itemize}
\tightlist
\item
  示例是局部调用片段；\passthrough{\lstinline!obj!}
  和参数变量表示已经按上层业务校验并准备好的对象。
\item
  该条目可以被编辑器和类型检查器识别，但当前运行时可能在属性查找阶段直接失败。
\item
  该方法属于运动命令。方法名包含
  \passthrough{\lstinline!stop!}、\passthrough{\lstinline!damp!} 或
  \passthrough{\lstinline!zero!} 也不代表它是独立物理急停。
\end{itemize}

\textbf{用法}

\begin{lstlisting}[language=Python]
from typing import TYPE_CHECKING

if TYPE_CHECKING:
    def planned_switch_move_mode(obj: LocoClient, flag: bool) -> int:
        return obj.switch_move_mode(flag=flag)
\end{lstlisting}

\begin{quote}
{[}!CAUTION{]} 上面的 \passthrough{\lstinline!TYPE\_CHECKING!}
示例只表达计划调用形态。该分支在普通 Python
运行时为假，不代表当前扩展能执行此接口。
\end{quote}

\hypertarget{unitree-sdk2-cpp-robot-h1-lococlient-set-next-foot-1}{%
\paragraph{\texorpdfstring{\texttt{unitree\_sdk2\_cpp.robot.h1.LocoClient.set\_next\_foot}}{unitree\_sdk2\_cpp.robot.h1.LocoClient.set\_next\_foot}}\label{unitree-sdk2-cpp-robot-h1-lococlient-set-next-foot-1}}

设置 \passthrough{\lstinline!next!}
\passthrough{\lstinline!foot!}。具体副作用和安全边界见下方状态。

\textbf{签名}

\begin{lstlisting}[language=Python]
def set_next_foot(self, foot: bool) -> int
\end{lstlisting}

\textbf{可用性与安全性}

\textbf{\passthrough{\lstinline!SIGNATURE\_ONLY!} /
\passthrough{\lstinline!MOTION\_COMMAND!}}：仅用于补全和静态检查。当前二进制没有该
Python 方法；不得按可执行运动接口使用。

\textbf{参数}

\begin{longtable}[]{@{}
  >{\raggedright\arraybackslash}p{(\columnwidth - 6\tabcolsep) * \real{0.18}}
  >{\raggedright\arraybackslash}p{(\columnwidth - 6\tabcolsep) * \real{0.24}}
  >{\raggedright\arraybackslash}p{(\columnwidth - 6\tabcolsep) * \real{0.18}}
  >{\raggedright\arraybackslash}p{(\columnwidth - 6\tabcolsep) * \real{0.40}}@{}}
\toprule
名称 & Python 类型 & 默认值 & 含义 \\
\midrule
\endhead
\passthrough{\lstinline!foot!} & \passthrough{\lstinline!bool!} & 必填 &
传给该接口的 \passthrough{\lstinline!foot!} 参数，Python 类型为
\passthrough{\lstinline!bool!}。精确含义、单位、范围和枚举值需查目标型号对应头文件与协议。
对应 C++ 参数 \passthrough{\lstinline!foot: bool!}。 只接受布尔语义。 \\
\bottomrule
\end{longtable}

\textbf{返回值}

\begin{longtable}[]{@{}
  >{\raggedright\arraybackslash}p{(\columnwidth - 4\tabcolsep) * \real{0.18}}
  >{\raggedright\arraybackslash}p{(\columnwidth - 4\tabcolsep) * \real{0.20}}
  >{\raggedright\arraybackslash}p{(\columnwidth - 4\tabcolsep) * \real{0.62}}@{}}
\toprule
位置 & 类型 & 含义 \\
\midrule
\endhead
返回值 & \passthrough{\lstinline!int!} & SDK 状态码；Unitree 示例通常以
\passthrough{\lstinline!0!} 表示成功，非零值需按具体服务定义解释。 \\
\bottomrule
\end{longtable}

\textbf{对应 C++}

\begin{itemize}
\tightlist
\item
  类：\passthrough{\lstinline!unitree::robot::h1::LocoClient!}
\item
  签名：\passthrough{\lstinline!SetNextFoot(bool)!}
\item
  绑定策略：\passthrough{\lstinline!DIRECT!}
\item
  声明位置：\passthrough{\lstinline!include/unitree/robot/h1/loco/h1\_loco\_client.hpp:253!}
\end{itemize}

\textbf{说明}

\begin{itemize}
\tightlist
\item
  示例是局部调用片段；\passthrough{\lstinline!obj!}
  和参数变量表示已经按上层业务校验并准备好的对象。
\item
  该条目可以被编辑器和类型检查器识别，但当前运行时可能在属性查找阶段直接失败。
\item
  该方法属于运动命令。方法名包含
  \passthrough{\lstinline!stop!}、\passthrough{\lstinline!damp!} 或
  \passthrough{\lstinline!zero!} 也不代表它是独立物理急停。
\end{itemize}

\textbf{用法}

\begin{lstlisting}[language=Python]
from typing import TYPE_CHECKING

if TYPE_CHECKING:
    def planned_set_next_foot(obj: LocoClient, foot: bool) -> int:
        return obj.set_next_foot(foot=foot)
\end{lstlisting}

\begin{quote}
{[}!CAUTION{]} 上面的 \passthrough{\lstinline!TYPE\_CHECKING!}
示例只表达计划调用形态。该分支在普通 Python
运行时为假，不代表当前扩展能执行此接口。
\end{quote}

\hypertarget{unitree-sdk2-cpp-robot-h1-lococlient-wave-hand-1}{%
\paragraph{\texorpdfstring{\texttt{unitree\_sdk2\_cpp.robot.h1.LocoClient.wave\_hand}}{unitree\_sdk2\_cpp.robot.h1.LocoClient.wave\_hand}}\label{unitree-sdk2-cpp-robot-h1-lococlient-wave-hand-1}}

对应 C++ SDK 操作
\passthrough{\lstinline!WaveHand()!}。上游头文件没有可直接生成的业务说明时，本参考只保证签名映射，精确语义需查目标型号协议。

\textbf{签名}

\begin{lstlisting}[language=Python]
def wave_hand(self) -> int
\end{lstlisting}

\textbf{可用性与安全性}

\textbf{\passthrough{\lstinline!SIGNATURE\_ONLY!} /
\passthrough{\lstinline!MOTION\_COMMAND!}}：仅用于补全和静态检查。当前二进制没有该
Python 方法；不得按可执行运动接口使用。

\textbf{参数}

无显式参数。实例方法中的 \passthrough{\lstinline!self!} 由 Python
自动传入。

\textbf{返回值}

\begin{longtable}[]{@{}
  >{\raggedright\arraybackslash}p{(\columnwidth - 4\tabcolsep) * \real{0.18}}
  >{\raggedright\arraybackslash}p{(\columnwidth - 4\tabcolsep) * \real{0.20}}
  >{\raggedright\arraybackslash}p{(\columnwidth - 4\tabcolsep) * \real{0.62}}@{}}
\toprule
位置 & 类型 & 含义 \\
\midrule
\endhead
返回值 & \passthrough{\lstinline!int!} & SDK 状态码；Unitree 示例通常以
\passthrough{\lstinline!0!} 表示成功，非零值需按具体服务定义解释。 \\
\bottomrule
\end{longtable}

\textbf{对应 C++}

\begin{itemize}
\tightlist
\item
  类：\passthrough{\lstinline!unitree::robot::h1::LocoClient!}
\item
  签名：\passthrough{\lstinline!WaveHand()!}
\item
  绑定策略：\passthrough{\lstinline!DIRECT!}
\item
  声明位置：\passthrough{\lstinline!include/unitree/robot/h1/loco/h1\_loco\_client.hpp:257!}
\end{itemize}

\textbf{说明}

\begin{itemize}
\tightlist
\item
  示例是局部调用片段；\passthrough{\lstinline!obj!}
  和参数变量表示已经按上层业务校验并准备好的对象。
\item
  该条目可以被编辑器和类型检查器识别，但当前运行时可能在属性查找阶段直接失败。
\item
  该方法属于运动命令。方法名包含
  \passthrough{\lstinline!stop!}、\passthrough{\lstinline!damp!} 或
  \passthrough{\lstinline!zero!} 也不代表它是独立物理急停。
\end{itemize}

\textbf{用法}

\begin{lstlisting}[language=Python]
from typing import TYPE_CHECKING

if TYPE_CHECKING:
    def planned_wave_hand(obj: LocoClient) -> int:
        return obj.wave_hand()
\end{lstlisting}

\begin{quote}
{[}!CAUTION{]} 上面的 \passthrough{\lstinline!TYPE\_CHECKING!}
示例只表达计划调用形态。该分支在普通 Python
运行时为假，不代表当前扩展能执行此接口。
\end{quote}

\hypertarget{unitree-sdk2-cpp-robot-h1-lococlient-shake-hand-1}{%
\paragraph{\texorpdfstring{\texttt{unitree\_sdk2\_cpp.robot.h1.LocoClient.shake\_hand}}{unitree\_sdk2\_cpp.robot.h1.LocoClient.shake\_hand}}\label{unitree-sdk2-cpp-robot-h1-lococlient-shake-hand-1}}

对应 C++ SDK 操作
\passthrough{\lstinline!ShakeHand(int)!}。上游头文件没有可直接生成的业务说明时，本参考只保证签名映射，精确语义需查目标型号协议。

\textbf{签名}

\begin{lstlisting}[language=Python]
def shake_hand(self, stage: int = ...) -> int
\end{lstlisting}

\textbf{可用性与安全性}

\textbf{\passthrough{\lstinline!SIGNATURE\_ONLY!} /
\passthrough{\lstinline!MOTION\_COMMAND!}}：仅用于补全和静态检查。当前二进制没有该
Python 方法；不得按可执行运动接口使用。

\textbf{参数}

\begin{longtable}[]{@{}
  >{\raggedright\arraybackslash}p{(\columnwidth - 6\tabcolsep) * \real{0.18}}
  >{\raggedright\arraybackslash}p{(\columnwidth - 6\tabcolsep) * \real{0.24}}
  >{\raggedright\arraybackslash}p{(\columnwidth - 6\tabcolsep) * \real{0.18}}
  >{\raggedright\arraybackslash}p{(\columnwidth - 6\tabcolsep) * \real{0.40}}@{}}
\toprule
名称 & Python 类型 & 默认值 & 含义 \\
\midrule
\endhead
\passthrough{\lstinline!stage!} & \passthrough{\lstinline!int!} &
\passthrough{\lstinline!...!} & 传给该接口的
\passthrough{\lstinline!stage!} 参数，Python 类型为
\passthrough{\lstinline!int!}。精确含义、单位、范围和枚举值需查目标型号对应头文件与协议。
对应 C++ 参数 \passthrough{\lstinline!stage: int!}。 \\
\bottomrule
\end{longtable}

\textbf{返回值}

\begin{longtable}[]{@{}
  >{\raggedright\arraybackslash}p{(\columnwidth - 4\tabcolsep) * \real{0.18}}
  >{\raggedright\arraybackslash}p{(\columnwidth - 4\tabcolsep) * \real{0.20}}
  >{\raggedright\arraybackslash}p{(\columnwidth - 4\tabcolsep) * \real{0.62}}@{}}
\toprule
位置 & 类型 & 含义 \\
\midrule
\endhead
返回值 & \passthrough{\lstinline!int!} & SDK 状态码；Unitree 示例通常以
\passthrough{\lstinline!0!} 表示成功，非零值需按具体服务定义解释。 \\
\bottomrule
\end{longtable}

\textbf{对应 C++}

\begin{itemize}
\tightlist
\item
  类：\passthrough{\lstinline!unitree::robot::h1::LocoClient!}
\item
  签名：\passthrough{\lstinline!ShakeHand(int)!}
\item
  绑定策略：\passthrough{\lstinline!DIRECT!}
\item
  声明位置：\passthrough{\lstinline!include/unitree/robot/h1/loco/h1\_loco\_client.hpp:259!}
\end{itemize}

\textbf{说明}

\begin{itemize}
\tightlist
\item
  示例是局部调用片段；\passthrough{\lstinline!obj!}
  和参数变量表示已经按上层业务校验并准备好的对象。
\item
  默认值显示为 \passthrough{\lstinline!...!}，表示 C++
  声明存在默认参数，但当前 AST
  清单没有保存其字面量；省略参数可使用上游默认行为。
\item
  该条目可以被编辑器和类型检查器识别，但当前运行时可能在属性查找阶段直接失败。
\item
  该方法属于运动命令。方法名包含
  \passthrough{\lstinline!stop!}、\passthrough{\lstinline!damp!} 或
  \passthrough{\lstinline!zero!} 也不代表它是独立物理急停。
\end{itemize}

\textbf{用法}

\begin{lstlisting}[language=Python]
from typing import TYPE_CHECKING

if TYPE_CHECKING:
    def planned_shake_hand(obj: LocoClient, stage: int) -> int:
        return obj.shake_hand(stage=stage)
\end{lstlisting}

\begin{quote}
{[}!CAUTION{]} 上面的 \passthrough{\lstinline!TYPE\_CHECKING!}
示例只表达计划调用形态。该分支在普通 Python
运行时为假，不代表当前扩展能执行此接口。
\end{quote}

\begin{center}\rule{0.5\linewidth}{0.5pt}\end{center}

\hypertarget{unitree-sdk2-cpp-robot-h2}{%
\subsection{H2 Robot API}\label{unitree-sdk2-cpp-robot-h2}}

模块：\passthrough{\lstinline!unitree\_sdk2\_cpp.robot.h2!}

\hypertarget{ux7c7bux7d22ux5f15-13}{%
\subsubsection{类索引}\label{ux7c7bux7d22ux5f15-13}}

\begin{longtable}[]{@{}
  >{\raggedright\arraybackslash}p{(\columnwidth - 4\tabcolsep) * \real{0.18}}
  >{\raggedleft\arraybackslash}p{(\columnwidth - 4\tabcolsep) * \real{0.20}}
  >{\raggedleft\arraybackslash}p{(\columnwidth - 4\tabcolsep) * \real{0.62}}@{}}
\toprule
类 & 公开函数签名 & 属性 \\
\midrule
\endhead
\protect\hyperlink{unitree-sdk2-cpp-robot-h2-fsmidinfo}{\passthrough{\lstinline!FsmIdInfo!}}
& 2 & 0 \\
\protect\hyperlink{unitree-sdk2-cpp-robot-h2-h2armactionclient}{\passthrough{\lstinline!H2ArmActionClient!}}
& 6 & 0 \\
\protect\hyperlink{unitree-sdk2-cpp-robot-h2-jsonizearmactioncommand}{\passthrough{\lstinline!JsonizeArmActionCommand!}}
& 2 & 0 \\
\protect\hyperlink{unitree-sdk2-cpp-robot-h2-jsonizearmactionname}{\passthrough{\lstinline!JsonizeArmActionName!}}
& 2 & 0 \\
\protect\hyperlink{unitree-sdk2-cpp-robot-h2-jsonizedatavecfloat}{\passthrough{\lstinline!JsonizeDataVecFloat!}}
& 3 & 0 \\
\protect\hyperlink{unitree-sdk2-cpp-robot-h2-jsonizefsmidlist}{\passthrough{\lstinline!JsonizeFsmIdList!}}
& 3 & 0 \\
\protect\hyperlink{unitree-sdk2-cpp-robot-h2-jsonizevelocitycommand}{\passthrough{\lstinline!JsonizeVelocityCommand!}}
& 3 & 0 \\
\protect\hyperlink{unitree-sdk2-cpp-robot-h2-lococlient}{\passthrough{\lstinline!LocoClient!}}
& 37 & 0 \\
\bottomrule
\end{longtable}

\hypertarget{unitree-sdk2-cpp-robot-h2-fsmidinfo}{%
\subsubsection{\texorpdfstring{\texttt{unitree\_sdk2\_cpp.robot.h2.FsmIdInfo}}{unitree\_sdk2\_cpp.robot.h2.FsmIdInfo}}\label{unitree-sdk2-cpp-robot-h2-fsmidinfo}}

SDK 类型签名预览。请逐项查看构造函数和方法的可用性。

\textbf{导入}

\begin{lstlisting}[language=Python]
from unitree_sdk2_cpp.robot.h2 import FsmIdInfo
\end{lstlisting}

\textbf{基类}：\passthrough{\lstinline!object!}

\textbf{公开属性}

\begin{longtable}[]{@{}
  >{\raggedright\arraybackslash}p{(\columnwidth - 6\tabcolsep) * \real{0.18}}
  >{\raggedright\arraybackslash}p{(\columnwidth - 6\tabcolsep) * \real{0.24}}
  >{\raggedright\arraybackslash}p{(\columnwidth - 6\tabcolsep) * \real{0.18}}
  >{\raggedright\arraybackslash}p{(\columnwidth - 6\tabcolsep) * \real{0.40}}@{}}
\toprule
名称 & Python 类型 & 含义 & 用法 \\
\midrule
\endhead
\passthrough{\lstinline!id!} & \passthrough{\lstinline!int!} &
传给该接口的 ID 参数，Python 类型为
\passthrough{\lstinline!int!}。精确含义、单位、范围和枚举值需查目标型号对应头文件与协议。
& \passthrough{\lstinline!value = obj.id!} /
\passthrough{\lstinline!obj.id = value!} \\
\passthrough{\lstinline!name!} & \passthrough{\lstinline!str!} & SDK
对象、配置项、服务或动作的名称。允许值由调用它的具体接口定义。 &
\passthrough{\lstinline!value = obj.name!} /
\passthrough{\lstinline!obj.name = value!} \\
\bottomrule
\end{longtable}

\textbf{方法索引}

\begin{longtable}[]{@{}
  >{\raggedright\arraybackslash}p{(\columnwidth - 6\tabcolsep) * \real{0.18}}
  >{\raggedright\arraybackslash}p{(\columnwidth - 6\tabcolsep) * \real{0.24}}
  >{\raggedright\arraybackslash}p{(\columnwidth - 6\tabcolsep) * \real{0.18}}
  >{\raggedright\arraybackslash}p{(\columnwidth - 6\tabcolsep) * \real{0.40}}@{}}
\toprule
方法 & Python 签名 & 状态 & 安全分类 \\
\midrule
\endhead
\protect\hyperlink{unitree-sdk2-cpp-robot-h2-fsmidinfo-dunder-init-1}{\passthrough{\lstinline!\_\_init\_\_（重载 1/2）!}}
& \passthrough{\lstinline!def \_\_init\_\_(self) -> None!} &
\passthrough{\lstinline!SIGNATURE\_ONLY!} &
\passthrough{\lstinline!CONSTRUCTION!} \\
\protect\hyperlink{unitree-sdk2-cpp-robot-h2-fsmidinfo-dunder-init-2}{\passthrough{\lstinline!\_\_init\_\_（重载 2/2）!}}
&
\passthrough{\lstinline!def \_\_init\_\_(self, i: int, n: str) -> None!}
& \passthrough{\lstinline!SIGNATURE\_ONLY!} &
\passthrough{\lstinline!CONSTRUCTION!} \\
\bottomrule
\end{longtable}

\hypertarget{unitree-sdk2-cpp-robot-h2-fsmidinfo-dunder-init-1}{%
\paragraph{\texorpdfstring{\texttt{unitree\_sdk2\_cpp.robot.h2.FsmIdInfo.\_\_init\_\_}（重载
1/2）}{unitree\_sdk2\_cpp.robot.h2.FsmIdInfo.\_\_init\_\_（重载 1/2）}}\label{unitree-sdk2-cpp-robot-h2-fsmidinfo-dunder-init-1}}

初始化 \passthrough{\lstinline!FsmIdInfo!}
实例。是否能实际构造取决于下方可用性状态。

\textbf{签名}

\begin{lstlisting}[language=Python]
@overload
def __init__(self) -> None
\end{lstlisting}

\textbf{可用性与安全性}

\textbf{\passthrough{\lstinline!SIGNATURE\_ONLY!} /
\passthrough{\lstinline!CONSTRUCTION!}}：当前只有设计期签名，不能假设已存在于运行时扩展。

\textbf{参数}

无显式参数。实例方法中的 \passthrough{\lstinline!self!} 由 Python
自动传入。

\textbf{返回值}

\begin{longtable}[]{@{}
  >{\raggedright\arraybackslash}p{(\columnwidth - 4\tabcolsep) * \real{0.18}}
  >{\raggedright\arraybackslash}p{(\columnwidth - 4\tabcolsep) * \real{0.20}}
  >{\raggedright\arraybackslash}p{(\columnwidth - 4\tabcolsep) * \real{0.62}}@{}}
\toprule
位置 & 类型 & 含义 \\
\midrule
\endhead
无 & \passthrough{\lstinline!None!} &
\passthrough{\lstinline!\_\_init\_\_!}
本身不返回值；调用类对象时，在构造函数可用的前提下得到该类实例。 \\
\bottomrule
\end{longtable}

\textbf{对应 C++}

\begin{itemize}
\tightlist
\item
  类：\passthrough{\lstinline!unitree::robot::h2::FsmIdInfo!}
\item
  签名：\passthrough{\lstinline!FsmIdInfo()!}
\item
  绑定策略：\passthrough{\lstinline!未单独标注!}
\item
  声明位置：\passthrough{\lstinline!include/unitree/robot/h2/loco/h2\_loco\_api.hpp:73!}
\end{itemize}

\textbf{说明}

\begin{itemize}
\tightlist
\item
  示例是局部调用片段；\passthrough{\lstinline!obj!}
  和参数变量表示已经按上层业务校验并准备好的对象。
\item
  该条目可以被编辑器和类型检查器识别，但当前运行时可能在属性查找阶段直接失败。
\end{itemize}

\textbf{用法}

\begin{lstlisting}[language=Python]
from typing import TYPE_CHECKING

if TYPE_CHECKING:
    def planned_construct() -> FsmIdInfo:
        return FsmIdInfo()
\end{lstlisting}

\begin{quote}
{[}!CAUTION{]} 上面的 \passthrough{\lstinline!TYPE\_CHECKING!}
示例只表达计划调用形态。该分支在普通 Python
运行时为假，不代表当前扩展能执行此接口。
\end{quote}

\hypertarget{unitree-sdk2-cpp-robot-h2-fsmidinfo-dunder-init-2}{%
\paragraph{\texorpdfstring{\texttt{unitree\_sdk2\_cpp.robot.h2.FsmIdInfo.\_\_init\_\_}（重载
2/2）}{unitree\_sdk2\_cpp.robot.h2.FsmIdInfo.\_\_init\_\_（重载 2/2）}}\label{unitree-sdk2-cpp-robot-h2-fsmidinfo-dunder-init-2}}

初始化 \passthrough{\lstinline!FsmIdInfo!}
实例。是否能实际构造取决于下方可用性状态。

\textbf{签名}

\begin{lstlisting}[language=Python]
@overload
def __init__(self, i: int, n: str) -> None
\end{lstlisting}

\textbf{可用性与安全性}

\textbf{\passthrough{\lstinline!SIGNATURE\_ONLY!} /
\passthrough{\lstinline!CONSTRUCTION!}}：当前只有设计期签名，不能假设已存在于运行时扩展。

\textbf{参数}

\begin{longtable}[]{@{}
  >{\raggedright\arraybackslash}p{(\columnwidth - 6\tabcolsep) * \real{0.18}}
  >{\raggedright\arraybackslash}p{(\columnwidth - 6\tabcolsep) * \real{0.24}}
  >{\raggedright\arraybackslash}p{(\columnwidth - 6\tabcolsep) * \real{0.18}}
  >{\raggedright\arraybackslash}p{(\columnwidth - 6\tabcolsep) * \real{0.40}}@{}}
\toprule
名称 & Python 类型 & 默认值 & 含义 \\
\midrule
\endhead
\passthrough{\lstinline!i!} & \passthrough{\lstinline!int!} & 必填 &
传给该接口的 \passthrough{\lstinline!i!} 参数，Python 类型为
\passthrough{\lstinline!int!}。精确含义、单位、范围和枚举值需查目标型号对应头文件与协议。
对应 C++ 参数 \passthrough{\lstinline!i: int!}。 \\
\passthrough{\lstinline!n!} & \passthrough{\lstinline!str!} & 必填 &
传给该接口的 \passthrough{\lstinline!n!} 参数，Python 类型为
\passthrough{\lstinline!str!}。精确含义、单位、范围和枚举值需查目标型号对应头文件与协议。
对应 C++ 参数 \passthrough{\lstinline!n: const std::string \&!}。
底层为字符串；长度、编码和允许值由具体协议决定。 \\
\bottomrule
\end{longtable}

\textbf{返回值}

\begin{longtable}[]{@{}
  >{\raggedright\arraybackslash}p{(\columnwidth - 4\tabcolsep) * \real{0.18}}
  >{\raggedright\arraybackslash}p{(\columnwidth - 4\tabcolsep) * \real{0.20}}
  >{\raggedright\arraybackslash}p{(\columnwidth - 4\tabcolsep) * \real{0.62}}@{}}
\toprule
位置 & 类型 & 含义 \\
\midrule
\endhead
无 & \passthrough{\lstinline!None!} &
\passthrough{\lstinline!\_\_init\_\_!}
本身不返回值；调用类对象时，在构造函数可用的前提下得到该类实例。 \\
\bottomrule
\end{longtable}

\textbf{对应 C++}

\begin{itemize}
\tightlist
\item
  类：\passthrough{\lstinline!unitree::robot::h2::FsmIdInfo!}
\item
  签名：\passthrough{\lstinline!FsmIdInfo(int, const std::string \&)!}
\item
  绑定策略：\passthrough{\lstinline!未单独标注!}
\item
  声明位置：\passthrough{\lstinline!include/unitree/robot/h2/loco/h2\_loco\_api.hpp:74!}
\end{itemize}

\textbf{说明}

\begin{itemize}
\tightlist
\item
  示例是局部调用片段；\passthrough{\lstinline!obj!}
  和参数变量表示已经按上层业务校验并准备好的对象。
\item
  该条目可以被编辑器和类型检查器识别，但当前运行时可能在属性查找阶段直接失败。
\end{itemize}

\textbf{用法}

\begin{lstlisting}[language=Python]
from typing import TYPE_CHECKING

if TYPE_CHECKING:
    def planned_construct(i: int, n: str) -> FsmIdInfo:
        return FsmIdInfo(i=i, n=n)
\end{lstlisting}

\begin{quote}
{[}!CAUTION{]} 上面的 \passthrough{\lstinline!TYPE\_CHECKING!}
示例只表达计划调用形态。该分支在普通 Python
运行时为假，不代表当前扩展能执行此接口。
\end{quote}

\hypertarget{unitree-sdk2-cpp-robot-h2-h2armactionclient}{%
\subsubsection{\texorpdfstring{\texttt{unitree\_sdk2\_cpp.robot.h2.H2ArmActionClient}}{unitree\_sdk2\_cpp.robot.h2.H2ArmActionClient}}\label{unitree-sdk2-cpp-robot-h2-h2armactionclient}}

Robot 服务客户端。构造、初始化和具体方法的可用性必须分别检查。

\textbf{导入}

\begin{lstlisting}[language=Python]
from unitree_sdk2_cpp.robot.h2 import H2ArmActionClient
\end{lstlisting}

\textbf{基类}：\passthrough{\lstinline!Client!}

\textbf{方法索引}

\begin{longtable}[]{@{}
  >{\raggedright\arraybackslash}p{(\columnwidth - 6\tabcolsep) * \real{0.18}}
  >{\raggedright\arraybackslash}p{(\columnwidth - 6\tabcolsep) * \real{0.24}}
  >{\raggedright\arraybackslash}p{(\columnwidth - 6\tabcolsep) * \real{0.18}}
  >{\raggedright\arraybackslash}p{(\columnwidth - 6\tabcolsep) * \real{0.40}}@{}}
\toprule
方法 & Python 签名 & 状态 & 安全分类 \\
\midrule
\endhead
\protect\hyperlink{unitree-sdk2-cpp-robot-h2-h2armactionclient-dunder-init-1}{\passthrough{\lstinline!\_\_init\_\_!}}
& \passthrough{\lstinline!def \_\_init\_\_(self) -> None!} &
\passthrough{\lstinline!AVAILABLE!} &
\passthrough{\lstinline!CONSTRUCTION!} \\
\protect\hyperlink{unitree-sdk2-cpp-robot-h2-h2armactionclient-init-1}{\passthrough{\lstinline!init!}}
& \passthrough{\lstinline!def init(self) -> None!} &
\passthrough{\lstinline!AVAILABLE!} &
\passthrough{\lstinline!INITIALIZATION!} \\
\protect\hyperlink{unitree-sdk2-cpp-robot-h2-h2armactionclient-execute-action-1}{\passthrough{\lstinline!execute\_action（重载 1/2）!}}
&
\passthrough{\lstinline!def execute\_action(self, action\_id: int) -> int!}
& \passthrough{\lstinline!SIGNATURE\_ONLY!} &
\passthrough{\lstinline!MOTION\_COMMAND!} \\
\protect\hyperlink{unitree-sdk2-cpp-robot-h2-h2armactionclient-execute-action-2}{\passthrough{\lstinline!execute\_action（重载 2/2）!}}
&
\passthrough{\lstinline!def execute\_action(self, action\_name: str) -> int!}
& \passthrough{\lstinline!SIGNATURE\_ONLY!} &
\passthrough{\lstinline!MOTION\_COMMAND!} \\
\protect\hyperlink{unitree-sdk2-cpp-robot-h2-h2armactionclient-stop-custom-action-1}{\passthrough{\lstinline!stop\_custom\_action!}}
& \passthrough{\lstinline!def stop\_custom\_action(self) -> int!} &
\passthrough{\lstinline!SIGNATURE\_ONLY!} &
\passthrough{\lstinline!MOTION\_COMMAND!} \\
\protect\hyperlink{unitree-sdk2-cpp-robot-h2-h2armactionclient-get-action-list-1}{\passthrough{\lstinline!get\_action\_list!}}
&
\passthrough{\lstinline!def get\_action\_list(self) -> tuple[int, str]!}
& \passthrough{\lstinline!AVAILABLE!} &
\passthrough{\lstinline!READ\_ONLY!} \\
\bottomrule
\end{longtable}

\hypertarget{unitree-sdk2-cpp-robot-h2-h2armactionclient-dunder-init-1}{%
\paragraph{\texorpdfstring{\texttt{unitree\_sdk2\_cpp.robot.h2.H2ArmActionClient.\_\_init\_\_}}{unitree\_sdk2\_cpp.robot.h2.H2ArmActionClient.\_\_init\_\_}}\label{unitree-sdk2-cpp-robot-h2-h2armactionclient-dunder-init-1}}

初始化 \passthrough{\lstinline!H2ArmActionClient!}
实例。是否能实际构造取决于下方可用性状态。

\textbf{签名}

\begin{lstlisting}[language=Python]
def __init__(self) -> None
\end{lstlisting}

\textbf{可用性与安全性}

\textbf{\passthrough{\lstinline!AVAILABLE!} /
\passthrough{\lstinline!CONSTRUCTION!}}：当前绑定源码已实现；实际执行条件仍取决于平台和对应
SDK 环境。

\textbf{参数}

无显式参数。实例方法中的 \passthrough{\lstinline!self!} 由 Python
自动传入。

\textbf{返回值}

\begin{longtable}[]{@{}
  >{\raggedright\arraybackslash}p{(\columnwidth - 4\tabcolsep) * \real{0.18}}
  >{\raggedright\arraybackslash}p{(\columnwidth - 4\tabcolsep) * \real{0.20}}
  >{\raggedright\arraybackslash}p{(\columnwidth - 4\tabcolsep) * \real{0.62}}@{}}
\toprule
位置 & 类型 & 含义 \\
\midrule
\endhead
无 & \passthrough{\lstinline!None!} &
\passthrough{\lstinline!\_\_init\_\_!}
本身不返回值；调用类对象时，在构造函数可用的前提下得到该类实例。 \\
\bottomrule
\end{longtable}

\textbf{对应 C++}

\begin{itemize}
\tightlist
\item
  类：\passthrough{\lstinline!unitree::robot::h2::H2ArmActionClient!}
\item
  签名：\passthrough{\lstinline!H2ArmActionClient()!}
\item
  绑定策略：\passthrough{\lstinline!未单独标注!}
\item
  声明位置：\passthrough{\lstinline!include/unitree/robot/h2/arm/h2\_arm\_action\_client.hpp:24!}
\end{itemize}

\textbf{说明}

\begin{itemize}
\tightlist
\item
  示例是局部调用片段；\passthrough{\lstinline!obj!}
  和参数变量表示已经按上层业务校验并准备好的对象。
\end{itemize}

\textbf{用法}

\begin{lstlisting}[language=Python]
value = H2ArmActionClient()
\end{lstlisting}

\hypertarget{unitree-sdk2-cpp-robot-h2-h2armactionclient-init-1}{%
\paragraph{\texorpdfstring{\texttt{unitree\_sdk2\_cpp.robot.h2.H2ArmActionClient.init}}{unitree\_sdk2\_cpp.robot.h2.H2ArmActionClient.init}}\label{unitree-sdk2-cpp-robot-h2-h2armactionclient-init-1}}

初始化当前 SDK 对象所需的底层通道或服务资源。

\textbf{签名}

\begin{lstlisting}[language=Python]
def init(self) -> None
\end{lstlisting}

\textbf{可用性与安全性}

\textbf{\passthrough{\lstinline!AVAILABLE!} /
\passthrough{\lstinline!INITIALIZATION!}}：当前绑定源码已实现；实际执行条件仍取决于平台和对应
SDK 环境。

\textbf{参数}

无显式参数。实例方法中的 \passthrough{\lstinline!self!} 由 Python
自动传入。

\textbf{返回值}

\begin{longtable}[]{@{}
  >{\raggedright\arraybackslash}p{(\columnwidth - 4\tabcolsep) * \real{0.18}}
  >{\raggedright\arraybackslash}p{(\columnwidth - 4\tabcolsep) * \real{0.20}}
  >{\raggedright\arraybackslash}p{(\columnwidth - 4\tabcolsep) * \real{0.62}}@{}}
\toprule
位置 & 类型 & 含义 \\
\midrule
\endhead
无 & \passthrough{\lstinline!None!} & 该方法不返回 Python
值；失败可能表现为异常或底层状态变化。 \\
\bottomrule
\end{longtable}

\textbf{对应 C++}

\begin{itemize}
\tightlist
\item
  类：\passthrough{\lstinline!unitree::robot::h2::H2ArmActionClient!}
\item
  签名：\passthrough{\lstinline!Init()!}
\item
  绑定策略：\passthrough{\lstinline!DIRECT!}
\item
  声明位置：\passthrough{\lstinline!include/unitree/robot/h2/arm/h2\_arm\_action\_client.hpp:26!}
\end{itemize}

\textbf{说明}

\begin{itemize}
\tightlist
\item
  示例是局部调用片段；\passthrough{\lstinline!obj!}
  和参数变量表示已经按上层业务校验并准备好的对象。
\end{itemize}

\textbf{用法}

\begin{lstlisting}[language=Python]
obj.init()
\end{lstlisting}

\hypertarget{unitree-sdk2-cpp-robot-h2-h2armactionclient-execute-action-1}{%
\paragraph{\texorpdfstring{\texttt{unitree\_sdk2\_cpp.robot.h2.H2ArmActionClient.execute\_action}（重载
1/2）}{unitree\_sdk2\_cpp.robot.h2.H2ArmActionClient.execute\_action（重载 1/2）}}\label{unitree-sdk2-cpp-robot-h2-h2armactionclient-execute-action-1}}

对应 C++ SDK 操作
\passthrough{\lstinline!ExecuteAction(int32\_t)!}。上游头文件没有可直接生成的业务说明时，本参考只保证签名映射，精确语义需查目标型号协议。

\textbf{签名}

\begin{lstlisting}[language=Python]
@overload
def execute_action(self, action_id: int) -> int
\end{lstlisting}

\textbf{可用性与安全性}

\textbf{\passthrough{\lstinline!SIGNATURE\_ONLY!} /
\passthrough{\lstinline!MOTION\_COMMAND!}}：仅用于补全和静态检查。当前二进制没有该
Python 方法；不得按可执行运动接口使用。

\textbf{参数}

\begin{longtable}[]{@{}
  >{\raggedright\arraybackslash}p{(\columnwidth - 6\tabcolsep) * \real{0.18}}
  >{\raggedright\arraybackslash}p{(\columnwidth - 6\tabcolsep) * \real{0.24}}
  >{\raggedright\arraybackslash}p{(\columnwidth - 6\tabcolsep) * \real{0.18}}
  >{\raggedright\arraybackslash}p{(\columnwidth - 6\tabcolsep) * \real{0.40}}@{}}
\toprule
名称 & Python 类型 & 默认值 & 含义 \\
\midrule
\endhead
\passthrough{\lstinline!action\_id!} & \passthrough{\lstinline!int!} &
必填 & 传给该接口的 \passthrough{\lstinline!action!} ID 参数，Python
类型为
\passthrough{\lstinline!int!}。精确含义、单位、范围和枚举值需查目标型号对应头文件与协议。
对应 C++ 参数 \passthrough{\lstinline!action\_id: int32\_t!}。
取值范围为 -2147483648 到 2147483647。 \\
\bottomrule
\end{longtable}

\textbf{返回值}

\begin{longtable}[]{@{}
  >{\raggedright\arraybackslash}p{(\columnwidth - 4\tabcolsep) * \real{0.18}}
  >{\raggedright\arraybackslash}p{(\columnwidth - 4\tabcolsep) * \real{0.20}}
  >{\raggedright\arraybackslash}p{(\columnwidth - 4\tabcolsep) * \real{0.62}}@{}}
\toprule
位置 & 类型 & 含义 \\
\midrule
\endhead
返回值 & \passthrough{\lstinline!int!} & SDK 状态码；Unitree 示例通常以
\passthrough{\lstinline!0!} 表示成功，非零值需按具体服务定义解释。 \\
\bottomrule
\end{longtable}

\textbf{对应 C++}

\begin{itemize}
\tightlist
\item
  类：\passthrough{\lstinline!unitree::robot::h2::H2ArmActionClient!}
\item
  签名：\passthrough{\lstinline!ExecuteAction(int32\_t)!}
\item
  绑定策略：\passthrough{\lstinline!DIRECT!}
\item
  声明位置：\passthrough{\lstinline!include/unitree/robot/h2/arm/h2\_arm\_action\_client.hpp:34!}
\end{itemize}

\textbf{说明}

\begin{itemize}
\tightlist
\item
  示例是局部调用片段；\passthrough{\lstinline!obj!}
  和参数变量表示已经按上层业务校验并准备好的对象。
\item
  该条目可以被编辑器和类型检查器识别，但当前运行时可能在属性查找阶段直接失败。
\item
  该方法属于运动命令。方法名包含
  \passthrough{\lstinline!stop!}、\passthrough{\lstinline!damp!} 或
  \passthrough{\lstinline!zero!} 也不代表它是独立物理急停。
\end{itemize}

\textbf{用法}

\begin{lstlisting}[language=Python]
from typing import TYPE_CHECKING

if TYPE_CHECKING:
    def planned_execute_action(obj: H2ArmActionClient, action_id: int) -> int:
        return obj.execute_action(action_id=action_id)
\end{lstlisting}

\begin{quote}
{[}!CAUTION{]} 上面的 \passthrough{\lstinline!TYPE\_CHECKING!}
示例只表达计划调用形态。该分支在普通 Python
运行时为假，不代表当前扩展能执行此接口。
\end{quote}

\hypertarget{unitree-sdk2-cpp-robot-h2-h2armactionclient-execute-action-2}{%
\paragraph{\texorpdfstring{\texttt{unitree\_sdk2\_cpp.robot.h2.H2ArmActionClient.execute\_action}（重载
2/2）}{unitree\_sdk2\_cpp.robot.h2.H2ArmActionClient.execute\_action（重载 2/2）}}\label{unitree-sdk2-cpp-robot-h2-h2armactionclient-execute-action-2}}

对应 C++ SDK 操作
\passthrough{\lstinline!ExecuteAction(const std::string \&)!}。上游头文件没有可直接生成的业务说明时，本参考只保证签名映射，精确语义需查目标型号协议。

\textbf{签名}

\begin{lstlisting}[language=Python]
@overload
def execute_action(self, action_name: str) -> int
\end{lstlisting}

\textbf{可用性与安全性}

\textbf{\passthrough{\lstinline!SIGNATURE\_ONLY!} /
\passthrough{\lstinline!MOTION\_COMMAND!}}：仅用于补全和静态检查。当前二进制没有该
Python 方法；不得按可执行运动接口使用。

\textbf{参数}

\begin{longtable}[]{@{}
  >{\raggedright\arraybackslash}p{(\columnwidth - 6\tabcolsep) * \real{0.18}}
  >{\raggedright\arraybackslash}p{(\columnwidth - 6\tabcolsep) * \real{0.24}}
  >{\raggedright\arraybackslash}p{(\columnwidth - 6\tabcolsep) * \real{0.18}}
  >{\raggedright\arraybackslash}p{(\columnwidth - 6\tabcolsep) * \real{0.40}}@{}}
\toprule
名称 & Python 类型 & 默认值 & 含义 \\
\midrule
\endhead
\passthrough{\lstinline!action\_name!} & \passthrough{\lstinline!str!} &
必填 & 传给该接口的 \passthrough{\lstinline!action!} 名称 参数，Python
类型为
\passthrough{\lstinline!str!}。精确含义、单位、范围和枚举值需查目标型号对应头文件与协议。
对应 C++ 参数
\passthrough{\lstinline!action\_name: const std::string \&!}。
底层为字符串；长度、编码和允许值由具体协议决定。 \\
\bottomrule
\end{longtable}

\textbf{返回值}

\begin{longtable}[]{@{}
  >{\raggedright\arraybackslash}p{(\columnwidth - 4\tabcolsep) * \real{0.18}}
  >{\raggedright\arraybackslash}p{(\columnwidth - 4\tabcolsep) * \real{0.20}}
  >{\raggedright\arraybackslash}p{(\columnwidth - 4\tabcolsep) * \real{0.62}}@{}}
\toprule
位置 & 类型 & 含义 \\
\midrule
\endhead
返回值 & \passthrough{\lstinline!int!} & SDK 状态码；Unitree 示例通常以
\passthrough{\lstinline!0!} 表示成功，非零值需按具体服务定义解释。 \\
\bottomrule
\end{longtable}

\textbf{对应 C++}

\begin{itemize}
\tightlist
\item
  类：\passthrough{\lstinline!unitree::robot::h2::H2ArmActionClient!}
\item
  签名：\passthrough{\lstinline!ExecuteAction(const std::string \&)!}
\item
  绑定策略：\passthrough{\lstinline!DIRECT!}
\item
  声明位置：\passthrough{\lstinline!include/unitree/robot/h2/arm/h2\_arm\_action\_client.hpp:43!}
\end{itemize}

\textbf{说明}

\begin{itemize}
\tightlist
\item
  示例是局部调用片段；\passthrough{\lstinline!obj!}
  和参数变量表示已经按上层业务校验并准备好的对象。
\item
  该条目可以被编辑器和类型检查器识别，但当前运行时可能在属性查找阶段直接失败。
\item
  该方法属于运动命令。方法名包含
  \passthrough{\lstinline!stop!}、\passthrough{\lstinline!damp!} 或
  \passthrough{\lstinline!zero!} 也不代表它是独立物理急停。
\end{itemize}

\textbf{用法}

\begin{lstlisting}[language=Python]
from typing import TYPE_CHECKING

if TYPE_CHECKING:
    def planned_execute_action(obj: H2ArmActionClient, action_name: str) -> int:
        return obj.execute_action(action_name=action_name)
\end{lstlisting}

\begin{quote}
{[}!CAUTION{]} 上面的 \passthrough{\lstinline!TYPE\_CHECKING!}
示例只表达计划调用形态。该分支在普通 Python
运行时为假，不代表当前扩展能执行此接口。
\end{quote}

\hypertarget{unitree-sdk2-cpp-robot-h2-h2armactionclient-stop-custom-action-1}{%
\paragraph{\texorpdfstring{\texttt{unitree\_sdk2\_cpp.robot.h2.H2ArmActionClient.stop\_custom\_action}}{unitree\_sdk2\_cpp.robot.h2.H2ArmActionClient.stop\_custom\_action}}\label{unitree-sdk2-cpp-robot-h2-h2armactionclient-stop-custom-action-1}}

请求停止对应 SDK 操作。该名称不等同于经过验证的物理急停。

\textbf{签名}

\begin{lstlisting}[language=Python]
def stop_custom_action(self) -> int
\end{lstlisting}

\textbf{可用性与安全性}

\textbf{\passthrough{\lstinline!SIGNATURE\_ONLY!} /
\passthrough{\lstinline!MOTION\_COMMAND!}}：仅用于补全和静态检查。当前二进制没有该
Python 方法；不得按可执行运动接口使用。

\textbf{参数}

无显式参数。实例方法中的 \passthrough{\lstinline!self!} 由 Python
自动传入。

\textbf{返回值}

\begin{longtable}[]{@{}
  >{\raggedright\arraybackslash}p{(\columnwidth - 4\tabcolsep) * \real{0.18}}
  >{\raggedright\arraybackslash}p{(\columnwidth - 4\tabcolsep) * \real{0.20}}
  >{\raggedright\arraybackslash}p{(\columnwidth - 4\tabcolsep) * \real{0.62}}@{}}
\toprule
位置 & 类型 & 含义 \\
\midrule
\endhead
返回值 & \passthrough{\lstinline!int!} & SDK 状态码；Unitree 示例通常以
\passthrough{\lstinline!0!} 表示成功，非零值需按具体服务定义解释。 \\
\bottomrule
\end{longtable}

\textbf{对应 C++}

\begin{itemize}
\tightlist
\item
  类：\passthrough{\lstinline!unitree::robot::h2::H2ArmActionClient!}
\item
  签名：\passthrough{\lstinline!StopCustomAction()!}
\item
  绑定策略：\passthrough{\lstinline!DIRECT!}
\item
  声明位置：\passthrough{\lstinline!include/unitree/robot/h2/arm/h2\_arm\_action\_client.hpp:52!}
\end{itemize}

\textbf{说明}

\begin{itemize}
\tightlist
\item
  示例是局部调用片段；\passthrough{\lstinline!obj!}
  和参数变量表示已经按上层业务校验并准备好的对象。
\item
  该条目可以被编辑器和类型检查器识别，但当前运行时可能在属性查找阶段直接失败。
\item
  该方法属于运动命令。方法名包含
  \passthrough{\lstinline!stop!}、\passthrough{\lstinline!damp!} 或
  \passthrough{\lstinline!zero!} 也不代表它是独立物理急停。
\end{itemize}

\textbf{用法}

\begin{lstlisting}[language=Python]
from typing import TYPE_CHECKING

if TYPE_CHECKING:
    def planned_stop_custom_action(obj: H2ArmActionClient) -> int:
        return obj.stop_custom_action()
\end{lstlisting}

\begin{quote}
{[}!CAUTION{]} 上面的 \passthrough{\lstinline!TYPE\_CHECKING!}
示例只表达计划调用形态。该分支在普通 Python
运行时为假，不代表当前扩展能执行此接口。
\end{quote}

\hypertarget{unitree-sdk2-cpp-robot-h2-h2armactionclient-get-action-list-1}{%
\paragraph{\texorpdfstring{\texttt{unitree\_sdk2\_cpp.robot.h2.H2ArmActionClient.get\_action\_list}}{unitree\_sdk2\_cpp.robot.h2.H2ArmActionClient.get\_action\_list}}\label{unitree-sdk2-cpp-robot-h2-h2armactionclient-get-action-list-1}}

查询或检查 \passthrough{\lstinline!action!}
\passthrough{\lstinline!list!}。

\textbf{签名}

\begin{lstlisting}[language=Python]
def get_action_list(self) -> tuple[int, str]
\end{lstlisting}

\textbf{可用性与安全性}

\textbf{\passthrough{\lstinline!AVAILABLE!} /
\passthrough{\lstinline!READ\_ONLY!}}：当前绑定源码已实现只读查询。Robot
Client 的服务端查询通常仍需匹配的 Linux
扩展、DDS、网络、机器人和服务；纯本地 getter 则不一定需要全部条件。

\textbf{参数}

无显式参数。实例方法中的 \passthrough{\lstinline!self!} 由 Python
自动传入。

\textbf{返回值}

\begin{longtable}[]{@{}
  >{\raggedright\arraybackslash}p{(\columnwidth - 4\tabcolsep) * \real{0.18}}
  >{\raggedright\arraybackslash}p{(\columnwidth - 4\tabcolsep) * \real{0.20}}
  >{\raggedright\arraybackslash}p{(\columnwidth - 4\tabcolsep) * \real{0.62}}@{}}
\toprule
位置 & 类型 & 含义 \\
\midrule
\endhead
{[}0{]} \passthrough{\lstinline!status!} & \passthrough{\lstinline!int!}
& SDK 状态码；Unitree 示例通常以 \passthrough{\lstinline!0!}
表示成功，非零值需按具体服务错误码解释。 \\
{[}1{]} \passthrough{\lstinline!data!} & \passthrough{\lstinline!str!} &
传给该接口的 数据 参数，Python 类型为
\passthrough{\lstinline!str!}。精确含义、单位、范围和枚举值需查目标型号对应头文件与协议。
对应 C++ 参数 \passthrough{\lstinline!data: std::string \&!}。
底层为字符串；长度、编码和允许值由具体协议决定。 \\
\bottomrule
\end{longtable}

\textbf{对应 C++}

\begin{itemize}
\tightlist
\item
  类：\passthrough{\lstinline!unitree::robot::h2::H2ArmActionClient!}
\item
  签名：\passthrough{\lstinline!GetActionList(std::string \&)!}
\item
  绑定策略：\passthrough{\lstinline!OUTPUT\_WRAPPER!}
\item
  声明位置：\passthrough{\lstinline!include/unitree/robot/h2/arm/h2\_arm\_action\_client.hpp:58!}
\end{itemize}

\textbf{说明}

\begin{itemize}
\tightlist
\item
  示例是局部调用片段；\passthrough{\lstinline!obj!}
  和参数变量表示已经按上层业务校验并准备好的对象。
\item
  C++ 的可修改输出引用不会作为 Python 入参出现，而是按 C++
  参数顺序追加到返回元组中。
\end{itemize}

\textbf{用法}

\begin{lstlisting}[language=Python]
status, data = obj.get_action_list()
\end{lstlisting}

\begin{quote}
{[}!NOTE{]} 只读查询仍会访问
DDS/机器人服务。必须检查返回状态码，并处理超时和网络中断。
\end{quote}

\hypertarget{unitree-sdk2-cpp-robot-h2-jsonizearmactioncommand}{%
\subsubsection{\texorpdfstring{\texttt{unitree\_sdk2\_cpp.robot.h2.JsonizeArmActionCommand}}{unitree\_sdk2\_cpp.robot.h2.JsonizeArmActionCommand}}\label{unitree-sdk2-cpp-robot-h2-jsonizearmactioncommand}}

SDK 类型签名预览。请逐项查看构造函数和方法的可用性。

\textbf{导入}

\begin{lstlisting}[language=Python]
from unitree_sdk2_cpp.robot.h2 import JsonizeArmActionCommand
\end{lstlisting}

\textbf{基类}：\passthrough{\lstinline!object!}

\textbf{公开属性}

\begin{longtable}[]{@{}
  >{\raggedright\arraybackslash}p{(\columnwidth - 6\tabcolsep) * \real{0.18}}
  >{\raggedright\arraybackslash}p{(\columnwidth - 6\tabcolsep) * \real{0.24}}
  >{\raggedright\arraybackslash}p{(\columnwidth - 6\tabcolsep) * \real{0.18}}
  >{\raggedright\arraybackslash}p{(\columnwidth - 6\tabcolsep) * \real{0.40}}@{}}
\toprule
名称 & Python 类型 & 含义 & 用法 \\
\midrule
\endhead
\passthrough{\lstinline!action\_id!} & \passthrough{\lstinline!int!} &
传给该接口的 \passthrough{\lstinline!action!} ID 参数，Python 类型为
\passthrough{\lstinline!int!}。精确含义、单位、范围和枚举值需查目标型号对应头文件与协议。
& \passthrough{\lstinline!value = obj.action\_id!} /
\passthrough{\lstinline!obj.action\_id = value!} \\
\bottomrule
\end{longtable}

\textbf{方法索引}

\begin{longtable}[]{@{}
  >{\raggedright\arraybackslash}p{(\columnwidth - 6\tabcolsep) * \real{0.18}}
  >{\raggedright\arraybackslash}p{(\columnwidth - 6\tabcolsep) * \real{0.24}}
  >{\raggedright\arraybackslash}p{(\columnwidth - 6\tabcolsep) * \real{0.18}}
  >{\raggedright\arraybackslash}p{(\columnwidth - 6\tabcolsep) * \real{0.40}}@{}}
\toprule
方法 & Python 签名 & 状态 & 安全分类 \\
\midrule
\endhead
\protect\hyperlink{unitree-sdk2-cpp-robot-h2-jsonizearmactioncommand-from-json-1}{\passthrough{\lstinline!from\_json!}}
&
\passthrough{\lstinline!def from\_json(self, json: dict[str, Any]) -> None!}
& \passthrough{\lstinline!SIGNATURE\_ONLY!} &
\passthrough{\lstinline!UNCLASSIFIED!} \\
\protect\hyperlink{unitree-sdk2-cpp-robot-h2-jsonizearmactioncommand-to-json-1}{\passthrough{\lstinline!to\_json!}}
&
\passthrough{\lstinline!def to\_json(self, json: dict[str, Any]) -> None!}
& \passthrough{\lstinline!SIGNATURE\_ONLY!} &
\passthrough{\lstinline!UNCLASSIFIED!} \\
\bottomrule
\end{longtable}

\hypertarget{unitree-sdk2-cpp-robot-h2-jsonizearmactioncommand-from-json-1}{%
\paragraph{\texorpdfstring{\texttt{unitree\_sdk2\_cpp.robot.h2.JsonizeArmActionCommand.from\_json}}{unitree\_sdk2\_cpp.robot.h2.JsonizeArmActionCommand.from\_json}}\label{unitree-sdk2-cpp-robot-h2-jsonizearmactioncommand-from-json-1}}

计划从 JSON 风格字典读取字段并更新当前 SDK 值对象。

\textbf{签名}

\begin{lstlisting}[language=Python]
def from_json(self, json: dict[str, Any]) -> None
\end{lstlisting}

\textbf{可用性与安全性}

\textbf{\passthrough{\lstinline!SIGNATURE\_ONLY!} /
\passthrough{\lstinline!UNCLASSIFIED!}}：当前只有设计期签名，不能假设已存在于运行时扩展。

\textbf{参数}

\begin{longtable}[]{@{}
  >{\raggedright\arraybackslash}p{(\columnwidth - 6\tabcolsep) * \real{0.18}}
  >{\raggedright\arraybackslash}p{(\columnwidth - 6\tabcolsep) * \real{0.24}}
  >{\raggedright\arraybackslash}p{(\columnwidth - 6\tabcolsep) * \real{0.18}}
  >{\raggedright\arraybackslash}p{(\columnwidth - 6\tabcolsep) * \real{0.40}}@{}}
\toprule
名称 & Python 类型 & 默认值 & 含义 \\
\midrule
\endhead
\passthrough{\lstinline!json!} &
\passthrough{\lstinline!dict[str, Any]!} & 必填 & JSON
风格字典，键和值必须满足对应 C++ \passthrough{\lstinline!JsonMap!}
数据结构的约束。 对应 C++ 参数
\passthrough{\lstinline!json: common::JsonMap \&!}。 \\
\bottomrule
\end{longtable}

\textbf{返回值}

\begin{longtable}[]{@{}
  >{\raggedright\arraybackslash}p{(\columnwidth - 4\tabcolsep) * \real{0.18}}
  >{\raggedright\arraybackslash}p{(\columnwidth - 4\tabcolsep) * \real{0.20}}
  >{\raggedright\arraybackslash}p{(\columnwidth - 4\tabcolsep) * \real{0.62}}@{}}
\toprule
位置 & 类型 & 含义 \\
\midrule
\endhead
无 & \passthrough{\lstinline!None!} & 该方法不返回 Python
值；失败可能表现为异常或底层状态变化。 \\
\bottomrule
\end{longtable}

\textbf{对应 C++}

\begin{itemize}
\tightlist
\item
  类：\passthrough{\lstinline!unitree::robot::h2::JsonizeArmActionCommand!}
\item
  签名：\passthrough{\lstinline!fromJson(common::JsonMap \&)!}
\item
  绑定策略：\passthrough{\lstinline!SIGNATURE\_PREVIEW!}
\item
  声明位置：\passthrough{\lstinline!include/unitree/robot/h2/arm/h2\_arm\_action\_api.hpp:22!}
\end{itemize}

\textbf{说明}

\begin{itemize}
\tightlist
\item
  示例是局部调用片段；\passthrough{\lstinline!obj!}
  和参数变量表示已经按上层业务校验并准备好的对象。
\item
  该条目可以被编辑器和类型检查器识别，但当前运行时可能在属性查找阶段直接失败。
\end{itemize}

\textbf{用法}

\begin{lstlisting}[language=Python]
from typing import TYPE_CHECKING

if TYPE_CHECKING:
    def planned_from_json(obj: JsonizeArmActionCommand, json: dict[str, Any]) -> None:
        obj.from_json(json=json)
\end{lstlisting}

\begin{quote}
{[}!CAUTION{]} 上面的 \passthrough{\lstinline!TYPE\_CHECKING!}
示例只表达计划调用形态。该分支在普通 Python
运行时为假，不代表当前扩展能执行此接口。
\end{quote}

\hypertarget{unitree-sdk2-cpp-robot-h2-jsonizearmactioncommand-to-json-1}{%
\paragraph{\texorpdfstring{\texttt{unitree\_sdk2\_cpp.robot.h2.JsonizeArmActionCommand.to\_json}}{unitree\_sdk2\_cpp.robot.h2.JsonizeArmActionCommand.to\_json}}\label{unitree-sdk2-cpp-robot-h2-jsonizearmactioncommand-to-json-1}}

计划把当前 SDK 值对象写入 JSON 风格字典。

\textbf{签名}

\begin{lstlisting}[language=Python]
def to_json(self, json: dict[str, Any]) -> None
\end{lstlisting}

\textbf{可用性与安全性}

\textbf{\passthrough{\lstinline!SIGNATURE\_ONLY!} /
\passthrough{\lstinline!UNCLASSIFIED!}}：当前只有设计期签名，不能假设已存在于运行时扩展。

\textbf{参数}

\begin{longtable}[]{@{}
  >{\raggedright\arraybackslash}p{(\columnwidth - 6\tabcolsep) * \real{0.18}}
  >{\raggedright\arraybackslash}p{(\columnwidth - 6\tabcolsep) * \real{0.24}}
  >{\raggedright\arraybackslash}p{(\columnwidth - 6\tabcolsep) * \real{0.18}}
  >{\raggedright\arraybackslash}p{(\columnwidth - 6\tabcolsep) * \real{0.40}}@{}}
\toprule
名称 & Python 类型 & 默认值 & 含义 \\
\midrule
\endhead
\passthrough{\lstinline!json!} &
\passthrough{\lstinline!dict[str, Any]!} & 必填 & JSON
风格字典，键和值必须满足对应 C++ \passthrough{\lstinline!JsonMap!}
数据结构的约束。 对应 C++ 参数
\passthrough{\lstinline!json: common::JsonMap \&!}。 \\
\bottomrule
\end{longtable}

\textbf{返回值}

\begin{longtable}[]{@{}
  >{\raggedright\arraybackslash}p{(\columnwidth - 4\tabcolsep) * \real{0.18}}
  >{\raggedright\arraybackslash}p{(\columnwidth - 4\tabcolsep) * \real{0.20}}
  >{\raggedright\arraybackslash}p{(\columnwidth - 4\tabcolsep) * \real{0.62}}@{}}
\toprule
位置 & 类型 & 含义 \\
\midrule
\endhead
无 & \passthrough{\lstinline!None!} & 该方法不返回 Python
值；失败可能表现为异常或底层状态变化。 \\
\bottomrule
\end{longtable}

\textbf{对应 C++}

\begin{itemize}
\tightlist
\item
  类：\passthrough{\lstinline!unitree::robot::h2::JsonizeArmActionCommand!}
\item
  签名：\passthrough{\lstinline!toJson(common::JsonMap \&) const!}
\item
  绑定策略：\passthrough{\lstinline!SIGNATURE\_PREVIEW!}
\item
  声明位置：\passthrough{\lstinline!include/unitree/robot/h2/arm/h2\_arm\_action\_api.hpp:26!}
\end{itemize}

\textbf{说明}

\begin{itemize}
\tightlist
\item
  示例是局部调用片段；\passthrough{\lstinline!obj!}
  和参数变量表示已经按上层业务校验并准备好的对象。
\item
  该条目可以被编辑器和类型检查器识别，但当前运行时可能在属性查找阶段直接失败。
\end{itemize}

\textbf{用法}

\begin{lstlisting}[language=Python]
from typing import TYPE_CHECKING

if TYPE_CHECKING:
    def planned_to_json(obj: JsonizeArmActionCommand, json: dict[str, Any]) -> None:
        obj.to_json(json=json)
\end{lstlisting}

\begin{quote}
{[}!CAUTION{]} 上面的 \passthrough{\lstinline!TYPE\_CHECKING!}
示例只表达计划调用形态。该分支在普通 Python
运行时为假，不代表当前扩展能执行此接口。
\end{quote}

\hypertarget{unitree-sdk2-cpp-robot-h2-jsonizearmactionname}{%
\subsubsection{\texorpdfstring{\texttt{unitree\_sdk2\_cpp.robot.h2.JsonizeArmActionName}}{unitree\_sdk2\_cpp.robot.h2.JsonizeArmActionName}}\label{unitree-sdk2-cpp-robot-h2-jsonizearmactionname}}

SDK 类型签名预览。请逐项查看构造函数和方法的可用性。

\textbf{导入}

\begin{lstlisting}[language=Python]
from unitree_sdk2_cpp.robot.h2 import JsonizeArmActionName
\end{lstlisting}

\textbf{基类}：\passthrough{\lstinline!object!}

\textbf{公开属性}

\begin{longtable}[]{@{}
  >{\raggedright\arraybackslash}p{(\columnwidth - 6\tabcolsep) * \real{0.18}}
  >{\raggedright\arraybackslash}p{(\columnwidth - 6\tabcolsep) * \real{0.24}}
  >{\raggedright\arraybackslash}p{(\columnwidth - 6\tabcolsep) * \real{0.18}}
  >{\raggedright\arraybackslash}p{(\columnwidth - 6\tabcolsep) * \real{0.40}}@{}}
\toprule
名称 & Python 类型 & 含义 & 用法 \\
\midrule
\endhead
\passthrough{\lstinline!action\_name!} & \passthrough{\lstinline!str!} &
传给该接口的 \passthrough{\lstinline!action!} 名称 参数，Python 类型为
\passthrough{\lstinline!str!}。精确含义、单位、范围和枚举值需查目标型号对应头文件与协议。
& \passthrough{\lstinline!value = obj.action\_name!} /
\passthrough{\lstinline!obj.action\_name = value!} \\
\bottomrule
\end{longtable}

\textbf{方法索引}

\begin{longtable}[]{@{}
  >{\raggedright\arraybackslash}p{(\columnwidth - 6\tabcolsep) * \real{0.18}}
  >{\raggedright\arraybackslash}p{(\columnwidth - 6\tabcolsep) * \real{0.24}}
  >{\raggedright\arraybackslash}p{(\columnwidth - 6\tabcolsep) * \real{0.18}}
  >{\raggedright\arraybackslash}p{(\columnwidth - 6\tabcolsep) * \real{0.40}}@{}}
\toprule
方法 & Python 签名 & 状态 & 安全分类 \\
\midrule
\endhead
\protect\hyperlink{unitree-sdk2-cpp-robot-h2-jsonizearmactionname-from-json-1}{\passthrough{\lstinline!from\_json!}}
&
\passthrough{\lstinline!def from\_json(self, json: dict[str, Any]) -> None!}
& \passthrough{\lstinline!SIGNATURE\_ONLY!} &
\passthrough{\lstinline!UNCLASSIFIED!} \\
\protect\hyperlink{unitree-sdk2-cpp-robot-h2-jsonizearmactionname-to-json-1}{\passthrough{\lstinline!to\_json!}}
&
\passthrough{\lstinline!def to\_json(self, json: dict[str, Any]) -> None!}
& \passthrough{\lstinline!SIGNATURE\_ONLY!} &
\passthrough{\lstinline!UNCLASSIFIED!} \\
\bottomrule
\end{longtable}

\hypertarget{unitree-sdk2-cpp-robot-h2-jsonizearmactionname-from-json-1}{%
\paragraph{\texorpdfstring{\texttt{unitree\_sdk2\_cpp.robot.h2.JsonizeArmActionName.from\_json}}{unitree\_sdk2\_cpp.robot.h2.JsonizeArmActionName.from\_json}}\label{unitree-sdk2-cpp-robot-h2-jsonizearmactionname-from-json-1}}

计划从 JSON 风格字典读取字段并更新当前 SDK 值对象。

\textbf{签名}

\begin{lstlisting}[language=Python]
def from_json(self, json: dict[str, Any]) -> None
\end{lstlisting}

\textbf{可用性与安全性}

\textbf{\passthrough{\lstinline!SIGNATURE\_ONLY!} /
\passthrough{\lstinline!UNCLASSIFIED!}}：当前只有设计期签名，不能假设已存在于运行时扩展。

\textbf{参数}

\begin{longtable}[]{@{}
  >{\raggedright\arraybackslash}p{(\columnwidth - 6\tabcolsep) * \real{0.18}}
  >{\raggedright\arraybackslash}p{(\columnwidth - 6\tabcolsep) * \real{0.24}}
  >{\raggedright\arraybackslash}p{(\columnwidth - 6\tabcolsep) * \real{0.18}}
  >{\raggedright\arraybackslash}p{(\columnwidth - 6\tabcolsep) * \real{0.40}}@{}}
\toprule
名称 & Python 类型 & 默认值 & 含义 \\
\midrule
\endhead
\passthrough{\lstinline!json!} &
\passthrough{\lstinline!dict[str, Any]!} & 必填 & JSON
风格字典，键和值必须满足对应 C++ \passthrough{\lstinline!JsonMap!}
数据结构的约束。 对应 C++ 参数
\passthrough{\lstinline!json: common::JsonMap \&!}。 \\
\bottomrule
\end{longtable}

\textbf{返回值}

\begin{longtable}[]{@{}
  >{\raggedright\arraybackslash}p{(\columnwidth - 4\tabcolsep) * \real{0.18}}
  >{\raggedright\arraybackslash}p{(\columnwidth - 4\tabcolsep) * \real{0.20}}
  >{\raggedright\arraybackslash}p{(\columnwidth - 4\tabcolsep) * \real{0.62}}@{}}
\toprule
位置 & 类型 & 含义 \\
\midrule
\endhead
无 & \passthrough{\lstinline!None!} & 该方法不返回 Python
值；失败可能表现为异常或底层状态变化。 \\
\bottomrule
\end{longtable}

\textbf{对应 C++}

\begin{itemize}
\tightlist
\item
  类：\passthrough{\lstinline!unitree::robot::h2::JsonizeArmActionName!}
\item
  签名：\passthrough{\lstinline!fromJson(common::JsonMap \&)!}
\item
  绑定策略：\passthrough{\lstinline!SIGNATURE\_PREVIEW!}
\item
  声明位置：\passthrough{\lstinline!include/unitree/robot/h2/arm/h2\_arm\_action\_api.hpp:35!}
\end{itemize}

\textbf{说明}

\begin{itemize}
\tightlist
\item
  示例是局部调用片段；\passthrough{\lstinline!obj!}
  和参数变量表示已经按上层业务校验并准备好的对象。
\item
  该条目可以被编辑器和类型检查器识别，但当前运行时可能在属性查找阶段直接失败。
\end{itemize}

\textbf{用法}

\begin{lstlisting}[language=Python]
from typing import TYPE_CHECKING

if TYPE_CHECKING:
    def planned_from_json(obj: JsonizeArmActionName, json: dict[str, Any]) -> None:
        obj.from_json(json=json)
\end{lstlisting}

\begin{quote}
{[}!CAUTION{]} 上面的 \passthrough{\lstinline!TYPE\_CHECKING!}
示例只表达计划调用形态。该分支在普通 Python
运行时为假，不代表当前扩展能执行此接口。
\end{quote}

\hypertarget{unitree-sdk2-cpp-robot-h2-jsonizearmactionname-to-json-1}{%
\paragraph{\texorpdfstring{\texttt{unitree\_sdk2\_cpp.robot.h2.JsonizeArmActionName.to\_json}}{unitree\_sdk2\_cpp.robot.h2.JsonizeArmActionName.to\_json}}\label{unitree-sdk2-cpp-robot-h2-jsonizearmactionname-to-json-1}}

计划把当前 SDK 值对象写入 JSON 风格字典。

\textbf{签名}

\begin{lstlisting}[language=Python]
def to_json(self, json: dict[str, Any]) -> None
\end{lstlisting}

\textbf{可用性与安全性}

\textbf{\passthrough{\lstinline!SIGNATURE\_ONLY!} /
\passthrough{\lstinline!UNCLASSIFIED!}}：当前只有设计期签名，不能假设已存在于运行时扩展。

\textbf{参数}

\begin{longtable}[]{@{}
  >{\raggedright\arraybackslash}p{(\columnwidth - 6\tabcolsep) * \real{0.18}}
  >{\raggedright\arraybackslash}p{(\columnwidth - 6\tabcolsep) * \real{0.24}}
  >{\raggedright\arraybackslash}p{(\columnwidth - 6\tabcolsep) * \real{0.18}}
  >{\raggedright\arraybackslash}p{(\columnwidth - 6\tabcolsep) * \real{0.40}}@{}}
\toprule
名称 & Python 类型 & 默认值 & 含义 \\
\midrule
\endhead
\passthrough{\lstinline!json!} &
\passthrough{\lstinline!dict[str, Any]!} & 必填 & JSON
风格字典，键和值必须满足对应 C++ \passthrough{\lstinline!JsonMap!}
数据结构的约束。 对应 C++ 参数
\passthrough{\lstinline!json: common::JsonMap \&!}。 \\
\bottomrule
\end{longtable}

\textbf{返回值}

\begin{longtable}[]{@{}
  >{\raggedright\arraybackslash}p{(\columnwidth - 4\tabcolsep) * \real{0.18}}
  >{\raggedright\arraybackslash}p{(\columnwidth - 4\tabcolsep) * \real{0.20}}
  >{\raggedright\arraybackslash}p{(\columnwidth - 4\tabcolsep) * \real{0.62}}@{}}
\toprule
位置 & 类型 & 含义 \\
\midrule
\endhead
无 & \passthrough{\lstinline!None!} & 该方法不返回 Python
值；失败可能表现为异常或底层状态变化。 \\
\bottomrule
\end{longtable}

\textbf{对应 C++}

\begin{itemize}
\tightlist
\item
  类：\passthrough{\lstinline!unitree::robot::h2::JsonizeArmActionName!}
\item
  签名：\passthrough{\lstinline!toJson(common::JsonMap \&) const!}
\item
  绑定策略：\passthrough{\lstinline!SIGNATURE\_PREVIEW!}
\item
  声明位置：\passthrough{\lstinline!include/unitree/robot/h2/arm/h2\_arm\_action\_api.hpp:39!}
\end{itemize}

\textbf{说明}

\begin{itemize}
\tightlist
\item
  示例是局部调用片段；\passthrough{\lstinline!obj!}
  和参数变量表示已经按上层业务校验并准备好的对象。
\item
  该条目可以被编辑器和类型检查器识别，但当前运行时可能在属性查找阶段直接失败。
\end{itemize}

\textbf{用法}

\begin{lstlisting}[language=Python]
from typing import TYPE_CHECKING

if TYPE_CHECKING:
    def planned_to_json(obj: JsonizeArmActionName, json: dict[str, Any]) -> None:
        obj.to_json(json=json)
\end{lstlisting}

\begin{quote}
{[}!CAUTION{]} 上面的 \passthrough{\lstinline!TYPE\_CHECKING!}
示例只表达计划调用形态。该分支在普通 Python
运行时为假，不代表当前扩展能执行此接口。
\end{quote}

\hypertarget{unitree-sdk2-cpp-robot-h2-jsonizedatavecfloat}{%
\subsubsection{\texorpdfstring{\texttt{unitree\_sdk2\_cpp.robot.h2.JsonizeDataVecFloat}}{unitree\_sdk2\_cpp.robot.h2.JsonizeDataVecFloat}}\label{unitree-sdk2-cpp-robot-h2-jsonizedatavecfloat}}

SDK 类型签名预览。请逐项查看构造函数和方法的可用性。

\textbf{导入}

\begin{lstlisting}[language=Python]
from unitree_sdk2_cpp.robot.h2 import JsonizeDataVecFloat
\end{lstlisting}

\textbf{基类}：\passthrough{\lstinline!object!}

\textbf{公开属性}

\begin{longtable}[]{@{}
  >{\raggedright\arraybackslash}p{(\columnwidth - 6\tabcolsep) * \real{0.18}}
  >{\raggedright\arraybackslash}p{(\columnwidth - 6\tabcolsep) * \real{0.24}}
  >{\raggedright\arraybackslash}p{(\columnwidth - 6\tabcolsep) * \real{0.18}}
  >{\raggedright\arraybackslash}p{(\columnwidth - 6\tabcolsep) * \real{0.40}}@{}}
\toprule
名称 & Python 类型 & 含义 & 用法 \\
\midrule
\endhead
\passthrough{\lstinline!data!} & \passthrough{\lstinline!list[float]!} &
传给该接口的 数据 参数，Python 类型为
\passthrough{\lstinline!list[float]!}。精确含义、单位、范围和枚举值需查目标型号对应头文件与协议。
& \passthrough{\lstinline!value = obj.data!} /
\passthrough{\lstinline!obj.data = value!} \\
\bottomrule
\end{longtable}

\textbf{方法索引}

\begin{longtable}[]{@{}
  >{\raggedright\arraybackslash}p{(\columnwidth - 6\tabcolsep) * \real{0.18}}
  >{\raggedright\arraybackslash}p{(\columnwidth - 6\tabcolsep) * \real{0.24}}
  >{\raggedright\arraybackslash}p{(\columnwidth - 6\tabcolsep) * \real{0.18}}
  >{\raggedright\arraybackslash}p{(\columnwidth - 6\tabcolsep) * \real{0.40}}@{}}
\toprule
方法 & Python 签名 & 状态 & 安全分类 \\
\midrule
\endhead
\protect\hyperlink{unitree-sdk2-cpp-robot-h2-jsonizedatavecfloat-dunder-init-1}{\passthrough{\lstinline!\_\_init\_\_!}}
& \passthrough{\lstinline!def \_\_init\_\_(self) -> None!} &
\passthrough{\lstinline!SIGNATURE\_ONLY!} &
\passthrough{\lstinline!CONSTRUCTION!} \\
\protect\hyperlink{unitree-sdk2-cpp-robot-h2-jsonizedatavecfloat-from-json-1}{\passthrough{\lstinline!from\_json!}}
&
\passthrough{\lstinline!def from\_json(self, json: dict[str, Any]) -> None!}
& \passthrough{\lstinline!SIGNATURE\_ONLY!} &
\passthrough{\lstinline!UNCLASSIFIED!} \\
\protect\hyperlink{unitree-sdk2-cpp-robot-h2-jsonizedatavecfloat-to-json-1}{\passthrough{\lstinline!to\_json!}}
&
\passthrough{\lstinline!def to\_json(self, json: dict[str, Any]) -> None!}
& \passthrough{\lstinline!SIGNATURE\_ONLY!} &
\passthrough{\lstinline!UNCLASSIFIED!} \\
\bottomrule
\end{longtable}

\hypertarget{unitree-sdk2-cpp-robot-h2-jsonizedatavecfloat-dunder-init-1}{%
\paragraph{\texorpdfstring{\texttt{unitree\_sdk2\_cpp.robot.h2.JsonizeDataVecFloat.\_\_init\_\_}}{unitree\_sdk2\_cpp.robot.h2.JsonizeDataVecFloat.\_\_init\_\_}}\label{unitree-sdk2-cpp-robot-h2-jsonizedatavecfloat-dunder-init-1}}

初始化 \passthrough{\lstinline!JsonizeDataVecFloat!}
实例。是否能实际构造取决于下方可用性状态。

\textbf{签名}

\begin{lstlisting}[language=Python]
def __init__(self) -> None
\end{lstlisting}

\textbf{可用性与安全性}

\textbf{\passthrough{\lstinline!SIGNATURE\_ONLY!} /
\passthrough{\lstinline!CONSTRUCTION!}}：当前只有设计期签名，不能假设已存在于运行时扩展。

\textbf{参数}

无显式参数。实例方法中的 \passthrough{\lstinline!self!} 由 Python
自动传入。

\textbf{返回值}

\begin{longtable}[]{@{}
  >{\raggedright\arraybackslash}p{(\columnwidth - 4\tabcolsep) * \real{0.18}}
  >{\raggedright\arraybackslash}p{(\columnwidth - 4\tabcolsep) * \real{0.20}}
  >{\raggedright\arraybackslash}p{(\columnwidth - 4\tabcolsep) * \real{0.62}}@{}}
\toprule
位置 & 类型 & 含义 \\
\midrule
\endhead
无 & \passthrough{\lstinline!None!} &
\passthrough{\lstinline!\_\_init\_\_!}
本身不返回值；调用类对象时，在构造函数可用的前提下得到该类实例。 \\
\bottomrule
\end{longtable}

\textbf{对应 C++}

\begin{itemize}
\tightlist
\item
  类：\passthrough{\lstinline!unitree::robot::h2::JsonizeDataVecFloat!}
\item
  签名：\passthrough{\lstinline!JsonizeDataVecFloat()!}
\item
  绑定策略：\passthrough{\lstinline!未单独标注!}
\item
  声明位置：\passthrough{\lstinline!include/unitree/robot/h2/loco/h2\_loco\_api.hpp:40!}
\end{itemize}

\textbf{说明}

\begin{itemize}
\tightlist
\item
  示例是局部调用片段；\passthrough{\lstinline!obj!}
  和参数变量表示已经按上层业务校验并准备好的对象。
\item
  该条目可以被编辑器和类型检查器识别，但当前运行时可能在属性查找阶段直接失败。
\end{itemize}

\textbf{用法}

\begin{lstlisting}[language=Python]
from typing import TYPE_CHECKING

if TYPE_CHECKING:
    def planned_construct() -> JsonizeDataVecFloat:
        return JsonizeDataVecFloat()
\end{lstlisting}

\begin{quote}
{[}!CAUTION{]} 上面的 \passthrough{\lstinline!TYPE\_CHECKING!}
示例只表达计划调用形态。该分支在普通 Python
运行时为假，不代表当前扩展能执行此接口。
\end{quote}

\hypertarget{unitree-sdk2-cpp-robot-h2-jsonizedatavecfloat-from-json-1}{%
\paragraph{\texorpdfstring{\texttt{unitree\_sdk2\_cpp.robot.h2.JsonizeDataVecFloat.from\_json}}{unitree\_sdk2\_cpp.robot.h2.JsonizeDataVecFloat.from\_json}}\label{unitree-sdk2-cpp-robot-h2-jsonizedatavecfloat-from-json-1}}

计划从 JSON 风格字典读取字段并更新当前 SDK 值对象。

\textbf{签名}

\begin{lstlisting}[language=Python]
def from_json(self, json: dict[str, Any]) -> None
\end{lstlisting}

\textbf{可用性与安全性}

\textbf{\passthrough{\lstinline!SIGNATURE\_ONLY!} /
\passthrough{\lstinline!UNCLASSIFIED!}}：当前只有设计期签名，不能假设已存在于运行时扩展。

\textbf{参数}

\begin{longtable}[]{@{}
  >{\raggedright\arraybackslash}p{(\columnwidth - 6\tabcolsep) * \real{0.18}}
  >{\raggedright\arraybackslash}p{(\columnwidth - 6\tabcolsep) * \real{0.24}}
  >{\raggedright\arraybackslash}p{(\columnwidth - 6\tabcolsep) * \real{0.18}}
  >{\raggedright\arraybackslash}p{(\columnwidth - 6\tabcolsep) * \real{0.40}}@{}}
\toprule
名称 & Python 类型 & 默认值 & 含义 \\
\midrule
\endhead
\passthrough{\lstinline!json!} &
\passthrough{\lstinline!dict[str, Any]!} & 必填 & JSON
风格字典，键和值必须满足对应 C++ \passthrough{\lstinline!JsonMap!}
数据结构的约束。 对应 C++ 参数
\passthrough{\lstinline!json: common::JsonMap \&!}。 \\
\bottomrule
\end{longtable}

\textbf{返回值}

\begin{longtable}[]{@{}
  >{\raggedright\arraybackslash}p{(\columnwidth - 4\tabcolsep) * \real{0.18}}
  >{\raggedright\arraybackslash}p{(\columnwidth - 4\tabcolsep) * \real{0.20}}
  >{\raggedright\arraybackslash}p{(\columnwidth - 4\tabcolsep) * \real{0.62}}@{}}
\toprule
位置 & 类型 & 含义 \\
\midrule
\endhead
无 & \passthrough{\lstinline!None!} & 该方法不返回 Python
值；失败可能表现为异常或底层状态变化。 \\
\bottomrule
\end{longtable}

\textbf{对应 C++}

\begin{itemize}
\tightlist
\item
  类：\passthrough{\lstinline!unitree::robot::h2::JsonizeDataVecFloat!}
\item
  签名：\passthrough{\lstinline!fromJson(common::JsonMap \&)!}
\item
  绑定策略：\passthrough{\lstinline!SIGNATURE\_PREVIEW!}
\item
  声明位置：\passthrough{\lstinline!include/unitree/robot/h2/loco/h2\_loco\_api.hpp:43!}
\end{itemize}

\textbf{说明}

\begin{itemize}
\tightlist
\item
  示例是局部调用片段；\passthrough{\lstinline!obj!}
  和参数变量表示已经按上层业务校验并准备好的对象。
\item
  该条目可以被编辑器和类型检查器识别，但当前运行时可能在属性查找阶段直接失败。
\end{itemize}

\textbf{用法}

\begin{lstlisting}[language=Python]
from typing import TYPE_CHECKING

if TYPE_CHECKING:
    def planned_from_json(obj: JsonizeDataVecFloat, json: dict[str, Any]) -> None:
        obj.from_json(json=json)
\end{lstlisting}

\begin{quote}
{[}!CAUTION{]} 上面的 \passthrough{\lstinline!TYPE\_CHECKING!}
示例只表达计划调用形态。该分支在普通 Python
运行时为假，不代表当前扩展能执行此接口。
\end{quote}

\hypertarget{unitree-sdk2-cpp-robot-h2-jsonizedatavecfloat-to-json-1}{%
\paragraph{\texorpdfstring{\texttt{unitree\_sdk2\_cpp.robot.h2.JsonizeDataVecFloat.to\_json}}{unitree\_sdk2\_cpp.robot.h2.JsonizeDataVecFloat.to\_json}}\label{unitree-sdk2-cpp-robot-h2-jsonizedatavecfloat-to-json-1}}

计划把当前 SDK 值对象写入 JSON 风格字典。

\textbf{签名}

\begin{lstlisting}[language=Python]
def to_json(self, json: dict[str, Any]) -> None
\end{lstlisting}

\textbf{可用性与安全性}

\textbf{\passthrough{\lstinline!SIGNATURE\_ONLY!} /
\passthrough{\lstinline!UNCLASSIFIED!}}：当前只有设计期签名，不能假设已存在于运行时扩展。

\textbf{参数}

\begin{longtable}[]{@{}
  >{\raggedright\arraybackslash}p{(\columnwidth - 6\tabcolsep) * \real{0.18}}
  >{\raggedright\arraybackslash}p{(\columnwidth - 6\tabcolsep) * \real{0.24}}
  >{\raggedright\arraybackslash}p{(\columnwidth - 6\tabcolsep) * \real{0.18}}
  >{\raggedright\arraybackslash}p{(\columnwidth - 6\tabcolsep) * \real{0.40}}@{}}
\toprule
名称 & Python 类型 & 默认值 & 含义 \\
\midrule
\endhead
\passthrough{\lstinline!json!} &
\passthrough{\lstinline!dict[str, Any]!} & 必填 & JSON
风格字典，键和值必须满足对应 C++ \passthrough{\lstinline!JsonMap!}
数据结构的约束。 对应 C++ 参数
\passthrough{\lstinline!json: common::JsonMap \&!}。 \\
\bottomrule
\end{longtable}

\textbf{返回值}

\begin{longtable}[]{@{}
  >{\raggedright\arraybackslash}p{(\columnwidth - 4\tabcolsep) * \real{0.18}}
  >{\raggedright\arraybackslash}p{(\columnwidth - 4\tabcolsep) * \real{0.20}}
  >{\raggedright\arraybackslash}p{(\columnwidth - 4\tabcolsep) * \real{0.62}}@{}}
\toprule
位置 & 类型 & 含义 \\
\midrule
\endhead
无 & \passthrough{\lstinline!None!} & 该方法不返回 Python
值；失败可能表现为异常或底层状态变化。 \\
\bottomrule
\end{longtable}

\textbf{对应 C++}

\begin{itemize}
\tightlist
\item
  类：\passthrough{\lstinline!unitree::robot::h2::JsonizeDataVecFloat!}
\item
  签名：\passthrough{\lstinline!toJson(common::JsonMap \&) const!}
\item
  绑定策略：\passthrough{\lstinline!SIGNATURE\_PREVIEW!}
\item
  声明位置：\passthrough{\lstinline!include/unitree/robot/h2/loco/h2\_loco\_api.hpp:45!}
\end{itemize}

\textbf{说明}

\begin{itemize}
\tightlist
\item
  示例是局部调用片段；\passthrough{\lstinline!obj!}
  和参数变量表示已经按上层业务校验并准备好的对象。
\item
  该条目可以被编辑器和类型检查器识别，但当前运行时可能在属性查找阶段直接失败。
\end{itemize}

\textbf{用法}

\begin{lstlisting}[language=Python]
from typing import TYPE_CHECKING

if TYPE_CHECKING:
    def planned_to_json(obj: JsonizeDataVecFloat, json: dict[str, Any]) -> None:
        obj.to_json(json=json)
\end{lstlisting}

\begin{quote}
{[}!CAUTION{]} 上面的 \passthrough{\lstinline!TYPE\_CHECKING!}
示例只表达计划调用形态。该分支在普通 Python
运行时为假，不代表当前扩展能执行此接口。
\end{quote}

\hypertarget{unitree-sdk2-cpp-robot-h2-jsonizefsmidlist}{%
\subsubsection{\texorpdfstring{\texttt{unitree\_sdk2\_cpp.robot.h2.JsonizeFsmIdList}}{unitree\_sdk2\_cpp.robot.h2.JsonizeFsmIdList}}\label{unitree-sdk2-cpp-robot-h2-jsonizefsmidlist}}

SDK 类型签名预览。请逐项查看构造函数和方法的可用性。

\textbf{导入}

\begin{lstlisting}[language=Python]
from unitree_sdk2_cpp.robot.h2 import JsonizeFsmIdList
\end{lstlisting}

\textbf{基类}：\passthrough{\lstinline!object!}

\textbf{公开属性}

\begin{longtable}[]{@{}
  >{\raggedright\arraybackslash}p{(\columnwidth - 6\tabcolsep) * \real{0.18}}
  >{\raggedright\arraybackslash}p{(\columnwidth - 6\tabcolsep) * \real{0.24}}
  >{\raggedright\arraybackslash}p{(\columnwidth - 6\tabcolsep) * \real{0.18}}
  >{\raggedright\arraybackslash}p{(\columnwidth - 6\tabcolsep) * \real{0.40}}@{}}
\toprule
名称 & Python 类型 & 含义 & 用法 \\
\midrule
\endhead
\passthrough{\lstinline!fsm\_ids!} &
\passthrough{\lstinline!list[FsmIdInfo]!} & 传给该接口的 FSM ID 列表
参数，Python 类型为
\passthrough{\lstinline!list[FsmIdInfo]!}。精确含义、单位、范围和枚举值需查目标型号对应头文件与协议。
& \passthrough{\lstinline!value = obj.fsm\_ids!} /
\passthrough{\lstinline!obj.fsm\_ids = value!} \\
\bottomrule
\end{longtable}

\textbf{方法索引}

\begin{longtable}[]{@{}
  >{\raggedright\arraybackslash}p{(\columnwidth - 6\tabcolsep) * \real{0.18}}
  >{\raggedright\arraybackslash}p{(\columnwidth - 6\tabcolsep) * \real{0.24}}
  >{\raggedright\arraybackslash}p{(\columnwidth - 6\tabcolsep) * \real{0.18}}
  >{\raggedright\arraybackslash}p{(\columnwidth - 6\tabcolsep) * \real{0.40}}@{}}
\toprule
方法 & Python 签名 & 状态 & 安全分类 \\
\midrule
\endhead
\protect\hyperlink{unitree-sdk2-cpp-robot-h2-jsonizefsmidlist-dunder-init-1}{\passthrough{\lstinline!\_\_init\_\_!}}
& \passthrough{\lstinline!def \_\_init\_\_(self) -> None!} &
\passthrough{\lstinline!SIGNATURE\_ONLY!} &
\passthrough{\lstinline!CONSTRUCTION!} \\
\protect\hyperlink{unitree-sdk2-cpp-robot-h2-jsonizefsmidlist-from-json-1}{\passthrough{\lstinline!from\_json!}}
&
\passthrough{\lstinline!def from\_json(self, json: dict[str, Any]) -> None!}
& \passthrough{\lstinline!SIGNATURE\_ONLY!} &
\passthrough{\lstinline!UNCLASSIFIED!} \\
\protect\hyperlink{unitree-sdk2-cpp-robot-h2-jsonizefsmidlist-to-json-1}{\passthrough{\lstinline!to\_json!}}
&
\passthrough{\lstinline!def to\_json(self, json: dict[str, Any]) -> None!}
& \passthrough{\lstinline!SIGNATURE\_ONLY!} &
\passthrough{\lstinline!UNCLASSIFIED!} \\
\bottomrule
\end{longtable}

\hypertarget{unitree-sdk2-cpp-robot-h2-jsonizefsmidlist-dunder-init-1}{%
\paragraph{\texorpdfstring{\texttt{unitree\_sdk2\_cpp.robot.h2.JsonizeFsmIdList.\_\_init\_\_}}{unitree\_sdk2\_cpp.robot.h2.JsonizeFsmIdList.\_\_init\_\_}}\label{unitree-sdk2-cpp-robot-h2-jsonizefsmidlist-dunder-init-1}}

初始化 \passthrough{\lstinline!JsonizeFsmIdList!}
实例。是否能实际构造取决于下方可用性状态。

\textbf{签名}

\begin{lstlisting}[language=Python]
def __init__(self) -> None
\end{lstlisting}

\textbf{可用性与安全性}

\textbf{\passthrough{\lstinline!SIGNATURE\_ONLY!} /
\passthrough{\lstinline!CONSTRUCTION!}}：当前只有设计期签名，不能假设已存在于运行时扩展。

\textbf{参数}

无显式参数。实例方法中的 \passthrough{\lstinline!self!} 由 Python
自动传入。

\textbf{返回值}

\begin{longtable}[]{@{}
  >{\raggedright\arraybackslash}p{(\columnwidth - 4\tabcolsep) * \real{0.18}}
  >{\raggedright\arraybackslash}p{(\columnwidth - 4\tabcolsep) * \real{0.20}}
  >{\raggedright\arraybackslash}p{(\columnwidth - 4\tabcolsep) * \real{0.62}}@{}}
\toprule
位置 & 类型 & 含义 \\
\midrule
\endhead
无 & \passthrough{\lstinline!None!} &
\passthrough{\lstinline!\_\_init\_\_!}
本身不返回值；调用类对象时，在构造函数可用的前提下得到该类实例。 \\
\bottomrule
\end{longtable}

\textbf{对应 C++}

\begin{itemize}
\tightlist
\item
  类：\passthrough{\lstinline!unitree::robot::h2::JsonizeFsmIdList!}
\item
  签名：\passthrough{\lstinline!JsonizeFsmIdList()!}
\item
  绑定策略：\passthrough{\lstinline!未单独标注!}
\item
  声明位置：\passthrough{\lstinline!include/unitree/robot/h2/loco/h2\_loco\_api.hpp:79!}
\end{itemize}

\textbf{说明}

\begin{itemize}
\tightlist
\item
  示例是局部调用片段；\passthrough{\lstinline!obj!}
  和参数变量表示已经按上层业务校验并准备好的对象。
\item
  该条目可以被编辑器和类型检查器识别，但当前运行时可能在属性查找阶段直接失败。
\end{itemize}

\textbf{用法}

\begin{lstlisting}[language=Python]
from typing import TYPE_CHECKING

if TYPE_CHECKING:
    def planned_construct() -> JsonizeFsmIdList:
        return JsonizeFsmIdList()
\end{lstlisting}

\begin{quote}
{[}!CAUTION{]} 上面的 \passthrough{\lstinline!TYPE\_CHECKING!}
示例只表达计划调用形态。该分支在普通 Python
运行时为假，不代表当前扩展能执行此接口。
\end{quote}

\hypertarget{unitree-sdk2-cpp-robot-h2-jsonizefsmidlist-from-json-1}{%
\paragraph{\texorpdfstring{\texttt{unitree\_sdk2\_cpp.robot.h2.JsonizeFsmIdList.from\_json}}{unitree\_sdk2\_cpp.robot.h2.JsonizeFsmIdList.from\_json}}\label{unitree-sdk2-cpp-robot-h2-jsonizefsmidlist-from-json-1}}

计划从 JSON 风格字典读取字段并更新当前 SDK 值对象。

\textbf{签名}

\begin{lstlisting}[language=Python]
def from_json(self, json: dict[str, Any]) -> None
\end{lstlisting}

\textbf{可用性与安全性}

\textbf{\passthrough{\lstinline!SIGNATURE\_ONLY!} /
\passthrough{\lstinline!UNCLASSIFIED!}}：当前只有设计期签名，不能假设已存在于运行时扩展。

\textbf{参数}

\begin{longtable}[]{@{}
  >{\raggedright\arraybackslash}p{(\columnwidth - 6\tabcolsep) * \real{0.18}}
  >{\raggedright\arraybackslash}p{(\columnwidth - 6\tabcolsep) * \real{0.24}}
  >{\raggedright\arraybackslash}p{(\columnwidth - 6\tabcolsep) * \real{0.18}}
  >{\raggedright\arraybackslash}p{(\columnwidth - 6\tabcolsep) * \real{0.40}}@{}}
\toprule
名称 & Python 类型 & 默认值 & 含义 \\
\midrule
\endhead
\passthrough{\lstinline!json!} &
\passthrough{\lstinline!dict[str, Any]!} & 必填 & JSON
风格字典，键和值必须满足对应 C++ \passthrough{\lstinline!JsonMap!}
数据结构的约束。 对应 C++ 参数
\passthrough{\lstinline!json: common::JsonMap \&!}。 \\
\bottomrule
\end{longtable}

\textbf{返回值}

\begin{longtable}[]{@{}
  >{\raggedright\arraybackslash}p{(\columnwidth - 4\tabcolsep) * \real{0.18}}
  >{\raggedright\arraybackslash}p{(\columnwidth - 4\tabcolsep) * \real{0.20}}
  >{\raggedright\arraybackslash}p{(\columnwidth - 4\tabcolsep) * \real{0.62}}@{}}
\toprule
位置 & 类型 & 含义 \\
\midrule
\endhead
无 & \passthrough{\lstinline!None!} & 该方法不返回 Python
值；失败可能表现为异常或底层状态变化。 \\
\bottomrule
\end{longtable}

\textbf{对应 C++}

\begin{itemize}
\tightlist
\item
  类：\passthrough{\lstinline!unitree::robot::h2::JsonizeFsmIdList!}
\item
  签名：\passthrough{\lstinline!fromJson(common::JsonMap \&)!}
\item
  绑定策略：\passthrough{\lstinline!SIGNATURE\_PREVIEW!}
\item
  声明位置：\passthrough{\lstinline!include/unitree/robot/h2/loco/h2\_loco\_api.hpp:82!}
\end{itemize}

\textbf{说明}

\begin{itemize}
\tightlist
\item
  示例是局部调用片段；\passthrough{\lstinline!obj!}
  和参数变量表示已经按上层业务校验并准备好的对象。
\item
  该条目可以被编辑器和类型检查器识别，但当前运行时可能在属性查找阶段直接失败。
\end{itemize}

\textbf{用法}

\begin{lstlisting}[language=Python]
from typing import TYPE_CHECKING

if TYPE_CHECKING:
    def planned_from_json(obj: JsonizeFsmIdList, json: dict[str, Any]) -> None:
        obj.from_json(json=json)
\end{lstlisting}

\begin{quote}
{[}!CAUTION{]} 上面的 \passthrough{\lstinline!TYPE\_CHECKING!}
示例只表达计划调用形态。该分支在普通 Python
运行时为假，不代表当前扩展能执行此接口。
\end{quote}

\hypertarget{unitree-sdk2-cpp-robot-h2-jsonizefsmidlist-to-json-1}{%
\paragraph{\texorpdfstring{\texttt{unitree\_sdk2\_cpp.robot.h2.JsonizeFsmIdList.to\_json}}{unitree\_sdk2\_cpp.robot.h2.JsonizeFsmIdList.to\_json}}\label{unitree-sdk2-cpp-robot-h2-jsonizefsmidlist-to-json-1}}

计划把当前 SDK 值对象写入 JSON 风格字典。

\textbf{签名}

\begin{lstlisting}[language=Python]
def to_json(self, json: dict[str, Any]) -> None
\end{lstlisting}

\textbf{可用性与安全性}

\textbf{\passthrough{\lstinline!SIGNATURE\_ONLY!} /
\passthrough{\lstinline!UNCLASSIFIED!}}：当前只有设计期签名，不能假设已存在于运行时扩展。

\textbf{参数}

\begin{longtable}[]{@{}
  >{\raggedright\arraybackslash}p{(\columnwidth - 6\tabcolsep) * \real{0.18}}
  >{\raggedright\arraybackslash}p{(\columnwidth - 6\tabcolsep) * \real{0.24}}
  >{\raggedright\arraybackslash}p{(\columnwidth - 6\tabcolsep) * \real{0.18}}
  >{\raggedright\arraybackslash}p{(\columnwidth - 6\tabcolsep) * \real{0.40}}@{}}
\toprule
名称 & Python 类型 & 默认值 & 含义 \\
\midrule
\endhead
\passthrough{\lstinline!json!} &
\passthrough{\lstinline!dict[str, Any]!} & 必填 & JSON
风格字典，键和值必须满足对应 C++ \passthrough{\lstinline!JsonMap!}
数据结构的约束。 对应 C++ 参数
\passthrough{\lstinline!json: common::JsonMap \&!}。 \\
\bottomrule
\end{longtable}

\textbf{返回值}

\begin{longtable}[]{@{}
  >{\raggedright\arraybackslash}p{(\columnwidth - 4\tabcolsep) * \real{0.18}}
  >{\raggedright\arraybackslash}p{(\columnwidth - 4\tabcolsep) * \real{0.20}}
  >{\raggedright\arraybackslash}p{(\columnwidth - 4\tabcolsep) * \real{0.62}}@{}}
\toprule
位置 & 类型 & 含义 \\
\midrule
\endhead
无 & \passthrough{\lstinline!None!} & 该方法不返回 Python
值；失败可能表现为异常或底层状态变化。 \\
\bottomrule
\end{longtable}

\textbf{对应 C++}

\begin{itemize}
\tightlist
\item
  类：\passthrough{\lstinline!unitree::robot::h2::JsonizeFsmIdList!}
\item
  签名：\passthrough{\lstinline!toJson(common::JsonMap \&) const!}
\item
  绑定策略：\passthrough{\lstinline!SIGNATURE\_PREVIEW!}
\item
  声明位置：\passthrough{\lstinline!include/unitree/robot/h2/loco/h2\_loco\_api.hpp:95!}
\end{itemize}

\textbf{说明}

\begin{itemize}
\tightlist
\item
  示例是局部调用片段；\passthrough{\lstinline!obj!}
  和参数变量表示已经按上层业务校验并准备好的对象。
\item
  该条目可以被编辑器和类型检查器识别，但当前运行时可能在属性查找阶段直接失败。
\end{itemize}

\textbf{用法}

\begin{lstlisting}[language=Python]
from typing import TYPE_CHECKING

if TYPE_CHECKING:
    def planned_to_json(obj: JsonizeFsmIdList, json: dict[str, Any]) -> None:
        obj.to_json(json=json)
\end{lstlisting}

\begin{quote}
{[}!CAUTION{]} 上面的 \passthrough{\lstinline!TYPE\_CHECKING!}
示例只表达计划调用形态。该分支在普通 Python
运行时为假，不代表当前扩展能执行此接口。
\end{quote}

\hypertarget{unitree-sdk2-cpp-robot-h2-jsonizevelocitycommand}{%
\subsubsection{\texorpdfstring{\texttt{unitree\_sdk2\_cpp.robot.h2.JsonizeVelocityCommand}}{unitree\_sdk2\_cpp.robot.h2.JsonizeVelocityCommand}}\label{unitree-sdk2-cpp-robot-h2-jsonizevelocitycommand}}

SDK 类型签名预览。请逐项查看构造函数和方法的可用性。

\textbf{导入}

\begin{lstlisting}[language=Python]
from unitree_sdk2_cpp.robot.h2 import JsonizeVelocityCommand
\end{lstlisting}

\textbf{基类}：\passthrough{\lstinline!object!}

\textbf{公开属性}

\begin{longtable}[]{@{}
  >{\raggedright\arraybackslash}p{(\columnwidth - 6\tabcolsep) * \real{0.18}}
  >{\raggedright\arraybackslash}p{(\columnwidth - 6\tabcolsep) * \real{0.24}}
  >{\raggedright\arraybackslash}p{(\columnwidth - 6\tabcolsep) * \real{0.18}}
  >{\raggedright\arraybackslash}p{(\columnwidth - 6\tabcolsep) * \real{0.40}}@{}}
\toprule
名称 & Python 类型 & 含义 & 用法 \\
\midrule
\endhead
\passthrough{\lstinline!velocity!} &
\passthrough{\lstinline!list[float]!} & 传给该接口的 速度 参数，Python
类型为
\passthrough{\lstinline!list[float]!}。精确含义、单位、范围和枚举值需查目标型号对应头文件与协议。
& \passthrough{\lstinline!value = obj.velocity!} /
\passthrough{\lstinline!obj.velocity = value!} \\
\passthrough{\lstinline!duration!} & \passthrough{\lstinline!float!} &
操作持续时间参数。精确单位、范围和默认行为以目标型号接口定义为准。 &
\passthrough{\lstinline!value = obj.duration!} /
\passthrough{\lstinline!obj.duration = value!} \\
\bottomrule
\end{longtable}

\textbf{方法索引}

\begin{longtable}[]{@{}
  >{\raggedright\arraybackslash}p{(\columnwidth - 6\tabcolsep) * \real{0.18}}
  >{\raggedright\arraybackslash}p{(\columnwidth - 6\tabcolsep) * \real{0.24}}
  >{\raggedright\arraybackslash}p{(\columnwidth - 6\tabcolsep) * \real{0.18}}
  >{\raggedright\arraybackslash}p{(\columnwidth - 6\tabcolsep) * \real{0.40}}@{}}
\toprule
方法 & Python 签名 & 状态 & 安全分类 \\
\midrule
\endhead
\protect\hyperlink{unitree-sdk2-cpp-robot-h2-jsonizevelocitycommand-dunder-init-1}{\passthrough{\lstinline!\_\_init\_\_!}}
& \passthrough{\lstinline!def \_\_init\_\_(self) -> None!} &
\passthrough{\lstinline!SIGNATURE\_ONLY!} &
\passthrough{\lstinline!CONSTRUCTION!} \\
\protect\hyperlink{unitree-sdk2-cpp-robot-h2-jsonizevelocitycommand-from-json-1}{\passthrough{\lstinline!from\_json!}}
&
\passthrough{\lstinline!def from\_json(self, json: dict[str, Any]) -> None!}
& \passthrough{\lstinline!SIGNATURE\_ONLY!} &
\passthrough{\lstinline!UNCLASSIFIED!} \\
\protect\hyperlink{unitree-sdk2-cpp-robot-h2-jsonizevelocitycommand-to-json-1}{\passthrough{\lstinline!to\_json!}}
&
\passthrough{\lstinline!def to\_json(self, json: dict[str, Any]) -> None!}
& \passthrough{\lstinline!SIGNATURE\_ONLY!} &
\passthrough{\lstinline!UNCLASSIFIED!} \\
\bottomrule
\end{longtable}

\hypertarget{unitree-sdk2-cpp-robot-h2-jsonizevelocitycommand-dunder-init-1}{%
\paragraph{\texorpdfstring{\texttt{unitree\_sdk2\_cpp.robot.h2.JsonizeVelocityCommand.\_\_init\_\_}}{unitree\_sdk2\_cpp.robot.h2.JsonizeVelocityCommand.\_\_init\_\_}}\label{unitree-sdk2-cpp-robot-h2-jsonizevelocitycommand-dunder-init-1}}

初始化 \passthrough{\lstinline!JsonizeVelocityCommand!}
实例。是否能实际构造取决于下方可用性状态。

\textbf{签名}

\begin{lstlisting}[language=Python]
def __init__(self) -> None
\end{lstlisting}

\textbf{可用性与安全性}

\textbf{\passthrough{\lstinline!SIGNATURE\_ONLY!} /
\passthrough{\lstinline!CONSTRUCTION!}}：当前只有设计期签名，不能假设已存在于运行时扩展。

\textbf{参数}

无显式参数。实例方法中的 \passthrough{\lstinline!self!} 由 Python
自动传入。

\textbf{返回值}

\begin{longtable}[]{@{}
  >{\raggedright\arraybackslash}p{(\columnwidth - 4\tabcolsep) * \real{0.18}}
  >{\raggedright\arraybackslash}p{(\columnwidth - 4\tabcolsep) * \real{0.20}}
  >{\raggedright\arraybackslash}p{(\columnwidth - 4\tabcolsep) * \real{0.62}}@{}}
\toprule
位置 & 类型 & 含义 \\
\midrule
\endhead
无 & \passthrough{\lstinline!None!} &
\passthrough{\lstinline!\_\_init\_\_!}
本身不返回值；调用类对象时，在构造函数可用的前提下得到该类实例。 \\
\bottomrule
\end{longtable}

\textbf{对应 C++}

\begin{itemize}
\tightlist
\item
  类：\passthrough{\lstinline!unitree::robot::h2::JsonizeVelocityCommand!}
\item
  签名：\passthrough{\lstinline!JsonizeVelocityCommand()!}
\item
  绑定策略：\passthrough{\lstinline!未单独标注!}
\item
  声明位置：\passthrough{\lstinline!include/unitree/robot/h2/loco/h2\_loco\_api.hpp:52!}
\end{itemize}

\textbf{说明}

\begin{itemize}
\tightlist
\item
  示例是局部调用片段；\passthrough{\lstinline!obj!}
  和参数变量表示已经按上层业务校验并准备好的对象。
\item
  该条目可以被编辑器和类型检查器识别，但当前运行时可能在属性查找阶段直接失败。
\end{itemize}

\textbf{用法}

\begin{lstlisting}[language=Python]
from typing import TYPE_CHECKING

if TYPE_CHECKING:
    def planned_construct() -> JsonizeVelocityCommand:
        return JsonizeVelocityCommand()
\end{lstlisting}

\begin{quote}
{[}!CAUTION{]} 上面的 \passthrough{\lstinline!TYPE\_CHECKING!}
示例只表达计划调用形态。该分支在普通 Python
运行时为假，不代表当前扩展能执行此接口。
\end{quote}

\hypertarget{unitree-sdk2-cpp-robot-h2-jsonizevelocitycommand-from-json-1}{%
\paragraph{\texorpdfstring{\texttt{unitree\_sdk2\_cpp.robot.h2.JsonizeVelocityCommand.from\_json}}{unitree\_sdk2\_cpp.robot.h2.JsonizeVelocityCommand.from\_json}}\label{unitree-sdk2-cpp-robot-h2-jsonizevelocitycommand-from-json-1}}

计划从 JSON 风格字典读取字段并更新当前 SDK 值对象。

\textbf{签名}

\begin{lstlisting}[language=Python]
def from_json(self, json: dict[str, Any]) -> None
\end{lstlisting}

\textbf{可用性与安全性}

\textbf{\passthrough{\lstinline!SIGNATURE\_ONLY!} /
\passthrough{\lstinline!UNCLASSIFIED!}}：当前只有设计期签名，不能假设已存在于运行时扩展。

\textbf{参数}

\begin{longtable}[]{@{}
  >{\raggedright\arraybackslash}p{(\columnwidth - 6\tabcolsep) * \real{0.18}}
  >{\raggedright\arraybackslash}p{(\columnwidth - 6\tabcolsep) * \real{0.24}}
  >{\raggedright\arraybackslash}p{(\columnwidth - 6\tabcolsep) * \real{0.18}}
  >{\raggedright\arraybackslash}p{(\columnwidth - 6\tabcolsep) * \real{0.40}}@{}}
\toprule
名称 & Python 类型 & 默认值 & 含义 \\
\midrule
\endhead
\passthrough{\lstinline!json!} &
\passthrough{\lstinline!dict[str, Any]!} & 必填 & JSON
风格字典，键和值必须满足对应 C++ \passthrough{\lstinline!JsonMap!}
数据结构的约束。 对应 C++ 参数
\passthrough{\lstinline!json: common::JsonMap \&!}。 \\
\bottomrule
\end{longtable}

\textbf{返回值}

\begin{longtable}[]{@{}
  >{\raggedright\arraybackslash}p{(\columnwidth - 4\tabcolsep) * \real{0.18}}
  >{\raggedright\arraybackslash}p{(\columnwidth - 4\tabcolsep) * \real{0.20}}
  >{\raggedright\arraybackslash}p{(\columnwidth - 4\tabcolsep) * \real{0.62}}@{}}
\toprule
位置 & 类型 & 含义 \\
\midrule
\endhead
无 & \passthrough{\lstinline!None!} & 该方法不返回 Python
值；失败可能表现为异常或底层状态变化。 \\
\bottomrule
\end{longtable}

\textbf{对应 C++}

\begin{itemize}
\tightlist
\item
  类：\passthrough{\lstinline!unitree::robot::h2::JsonizeVelocityCommand!}
\item
  签名：\passthrough{\lstinline!fromJson(common::JsonMap \&)!}
\item
  绑定策略：\passthrough{\lstinline!SIGNATURE\_PREVIEW!}
\item
  声明位置：\passthrough{\lstinline!include/unitree/robot/h2/loco/h2\_loco\_api.hpp:55!}
\end{itemize}

\textbf{说明}

\begin{itemize}
\tightlist
\item
  示例是局部调用片段；\passthrough{\lstinline!obj!}
  和参数变量表示已经按上层业务校验并准备好的对象。
\item
  该条目可以被编辑器和类型检查器识别，但当前运行时可能在属性查找阶段直接失败。
\end{itemize}

\textbf{用法}

\begin{lstlisting}[language=Python]
from typing import TYPE_CHECKING

if TYPE_CHECKING:
    def planned_from_json(obj: JsonizeVelocityCommand, json: dict[str, Any]) -> None:
        obj.from_json(json=json)
\end{lstlisting}

\begin{quote}
{[}!CAUTION{]} 上面的 \passthrough{\lstinline!TYPE\_CHECKING!}
示例只表达计划调用形态。该分支在普通 Python
运行时为假，不代表当前扩展能执行此接口。
\end{quote}

\hypertarget{unitree-sdk2-cpp-robot-h2-jsonizevelocitycommand-to-json-1}{%
\paragraph{\texorpdfstring{\texttt{unitree\_sdk2\_cpp.robot.h2.JsonizeVelocityCommand.to\_json}}{unitree\_sdk2\_cpp.robot.h2.JsonizeVelocityCommand.to\_json}}\label{unitree-sdk2-cpp-robot-h2-jsonizevelocitycommand-to-json-1}}

计划把当前 SDK 值对象写入 JSON 风格字典。

\textbf{签名}

\begin{lstlisting}[language=Python]
def to_json(self, json: dict[str, Any]) -> None
\end{lstlisting}

\textbf{可用性与安全性}

\textbf{\passthrough{\lstinline!SIGNATURE\_ONLY!} /
\passthrough{\lstinline!UNCLASSIFIED!}}：当前只有设计期签名，不能假设已存在于运行时扩展。

\textbf{参数}

\begin{longtable}[]{@{}
  >{\raggedright\arraybackslash}p{(\columnwidth - 6\tabcolsep) * \real{0.18}}
  >{\raggedright\arraybackslash}p{(\columnwidth - 6\tabcolsep) * \real{0.24}}
  >{\raggedright\arraybackslash}p{(\columnwidth - 6\tabcolsep) * \real{0.18}}
  >{\raggedright\arraybackslash}p{(\columnwidth - 6\tabcolsep) * \real{0.40}}@{}}
\toprule
名称 & Python 类型 & 默认值 & 含义 \\
\midrule
\endhead
\passthrough{\lstinline!json!} &
\passthrough{\lstinline!dict[str, Any]!} & 必填 & JSON
风格字典，键和值必须满足对应 C++ \passthrough{\lstinline!JsonMap!}
数据结构的约束。 对应 C++ 参数
\passthrough{\lstinline!json: common::JsonMap \&!}。 \\
\bottomrule
\end{longtable}

\textbf{返回值}

\begin{longtable}[]{@{}
  >{\raggedright\arraybackslash}p{(\columnwidth - 4\tabcolsep) * \real{0.18}}
  >{\raggedright\arraybackslash}p{(\columnwidth - 4\tabcolsep) * \real{0.20}}
  >{\raggedright\arraybackslash}p{(\columnwidth - 4\tabcolsep) * \real{0.62}}@{}}
\toprule
位置 & 类型 & 含义 \\
\midrule
\endhead
无 & \passthrough{\lstinline!None!} & 该方法不返回 Python
值；失败可能表现为异常或底层状态变化。 \\
\bottomrule
\end{longtable}

\textbf{对应 C++}

\begin{itemize}
\tightlist
\item
  类：\passthrough{\lstinline!unitree::robot::h2::JsonizeVelocityCommand!}
\item
  签名：\passthrough{\lstinline!toJson(common::JsonMap \&) const!}
\item
  绑定策略：\passthrough{\lstinline!SIGNATURE\_PREVIEW!}
\item
  声明位置：\passthrough{\lstinline!include/unitree/robot/h2/loco/h2\_loco\_api.hpp:60!}
\end{itemize}

\textbf{说明}

\begin{itemize}
\tightlist
\item
  示例是局部调用片段；\passthrough{\lstinline!obj!}
  和参数变量表示已经按上层业务校验并准备好的对象。
\item
  该条目可以被编辑器和类型检查器识别，但当前运行时可能在属性查找阶段直接失败。
\end{itemize}

\textbf{用法}

\begin{lstlisting}[language=Python]
from typing import TYPE_CHECKING

if TYPE_CHECKING:
    def planned_to_json(obj: JsonizeVelocityCommand, json: dict[str, Any]) -> None:
        obj.to_json(json=json)
\end{lstlisting}

\begin{quote}
{[}!CAUTION{]} 上面的 \passthrough{\lstinline!TYPE\_CHECKING!}
示例只表达计划调用形态。该分支在普通 Python
运行时为假，不代表当前扩展能执行此接口。
\end{quote}

\hypertarget{unitree-sdk2-cpp-robot-h2-lococlient}{%
\subsubsection{\texorpdfstring{\texttt{unitree\_sdk2\_cpp.robot.h2.LocoClient}}{unitree\_sdk2\_cpp.robot.h2.LocoClient}}\label{unitree-sdk2-cpp-robot-h2-lococlient}}

Robot 服务客户端。构造、初始化和具体方法的可用性必须分别检查。

\textbf{导入}

\begin{lstlisting}[language=Python]
from unitree_sdk2_cpp.robot.h2 import LocoClient
\end{lstlisting}

\textbf{基类}：\passthrough{\lstinline!Client!}

\textbf{方法索引}

\begin{longtable}[]{@{}
  >{\raggedright\arraybackslash}p{(\columnwidth - 6\tabcolsep) * \real{0.18}}
  >{\raggedright\arraybackslash}p{(\columnwidth - 6\tabcolsep) * \real{0.24}}
  >{\raggedright\arraybackslash}p{(\columnwidth - 6\tabcolsep) * \real{0.18}}
  >{\raggedright\arraybackslash}p{(\columnwidth - 6\tabcolsep) * \real{0.40}}@{}}
\toprule
方法 & Python 签名 & 状态 & 安全分类 \\
\midrule
\endhead
\protect\hyperlink{unitree-sdk2-cpp-robot-h2-lococlient-dunder-init-1}{\passthrough{\lstinline!\_\_init\_\_!}}
& \passthrough{\lstinline!def \_\_init\_\_(self) -> None!} &
\passthrough{\lstinline!AVAILABLE!} &
\passthrough{\lstinline!CONSTRUCTION!} \\
\protect\hyperlink{unitree-sdk2-cpp-robot-h2-lococlient-init-1}{\passthrough{\lstinline!init!}}
& \passthrough{\lstinline!def init(self) -> None!} &
\passthrough{\lstinline!AVAILABLE!} &
\passthrough{\lstinline!INITIALIZATION!} \\
\protect\hyperlink{unitree-sdk2-cpp-robot-h2-lococlient-get-fsm-id-1}{\passthrough{\lstinline!get\_fsm\_id!}}
& \passthrough{\lstinline!def get\_fsm\_id(self) -> tuple[int, int]!} &
\passthrough{\lstinline!AVAILABLE!} &
\passthrough{\lstinline!READ\_ONLY!} \\
\protect\hyperlink{unitree-sdk2-cpp-robot-h2-lococlient-get-fsm-mode-1}{\passthrough{\lstinline!get\_fsm\_mode!}}
& \passthrough{\lstinline!def get\_fsm\_mode(self) -> tuple[int, int]!}
& \passthrough{\lstinline!AVAILABLE!} &
\passthrough{\lstinline!READ\_ONLY!} \\
\protect\hyperlink{unitree-sdk2-cpp-robot-h2-lococlient-get-balance-mode-1}{\passthrough{\lstinline!get\_balance\_mode!}}
&
\passthrough{\lstinline!def get\_balance\_mode(self) -> tuple[int, int]!}
& \passthrough{\lstinline!AVAILABLE!} &
\passthrough{\lstinline!READ\_ONLY!} \\
\protect\hyperlink{unitree-sdk2-cpp-robot-h2-lococlient-get-swing-height-1}{\passthrough{\lstinline!get\_swing\_height!}}
&
\passthrough{\lstinline!def get\_swing\_height(self) -> tuple[int, float]!}
& \passthrough{\lstinline!AVAILABLE!} &
\passthrough{\lstinline!READ\_ONLY!} \\
\protect\hyperlink{unitree-sdk2-cpp-robot-h2-lococlient-get-stand-height-1}{\passthrough{\lstinline!get\_stand\_height!}}
&
\passthrough{\lstinline!def get\_stand\_height(self) -> tuple[int, float]!}
& \passthrough{\lstinline!AVAILABLE!} &
\passthrough{\lstinline!READ\_ONLY!} \\
\protect\hyperlink{unitree-sdk2-cpp-robot-h2-lococlient-get-phase-1}{\passthrough{\lstinline!get\_phase!}}
&
\passthrough{\lstinline!def get\_phase(self) -> tuple[int, list[float]]!}
& \passthrough{\lstinline!AVAILABLE!} &
\passthrough{\lstinline!READ\_ONLY!} \\
\protect\hyperlink{unitree-sdk2-cpp-robot-h2-lococlient-get-arm-sdk-status-1}{\passthrough{\lstinline!get\_arm\_sdk\_status!}}
&
\passthrough{\lstinline!def get\_arm\_sdk\_status(self) -> tuple[int, bool]!}
& \passthrough{\lstinline!AVAILABLE!} &
\passthrough{\lstinline!READ\_ONLY!} \\
\protect\hyperlink{unitree-sdk2-cpp-robot-h2-lococlient-get-available-fsm-ids-1}{\passthrough{\lstinline!get\_available\_fsm\_ids!}}
&
\passthrough{\lstinline!def get\_available\_fsm\_ids(self) -> tuple[int, list[int], list[str]]!}
& \passthrough{\lstinline!AVAILABLE!} &
\passthrough{\lstinline!READ\_ONLY!} \\
\protect\hyperlink{unitree-sdk2-cpp-robot-h2-lococlient-set-fsm-id-1}{\passthrough{\lstinline!set\_fsm\_id!}}
& \passthrough{\lstinline!def set\_fsm\_id(self, fsm\_id: int) -> int!}
& \passthrough{\lstinline!SIGNATURE\_ONLY!} &
\passthrough{\lstinline!MOTION\_COMMAND!} \\
\protect\hyperlink{unitree-sdk2-cpp-robot-h2-lococlient-set-balance-mode-1}{\passthrough{\lstinline!set\_balance\_mode!}}
&
\passthrough{\lstinline!def set\_balance\_mode(self, balance\_mode: int) -> int!}
& \passthrough{\lstinline!SIGNATURE\_ONLY!} &
\passthrough{\lstinline!MOTION\_COMMAND!} \\
\protect\hyperlink{unitree-sdk2-cpp-robot-h2-lococlient-set-punch-api-1}{\passthrough{\lstinline!set\_punch\_api!}}
&
\passthrough{\lstinline!def set\_punch\_api(self) -> tuple[int, list[float]]!}
& \passthrough{\lstinline!SIGNATURE\_ONLY!} &
\passthrough{\lstinline!MOTION\_COMMAND!} \\
\protect\hyperlink{unitree-sdk2-cpp-robot-h2-lococlient-set-swing-height-1}{\passthrough{\lstinline!set\_swing\_height!}}
&
\passthrough{\lstinline!def set\_swing\_height(self, swing\_height: float) -> int!}
& \passthrough{\lstinline!SIGNATURE\_ONLY!} &
\passthrough{\lstinline!MOTION\_COMMAND!} \\
\protect\hyperlink{unitree-sdk2-cpp-robot-h2-lococlient-set-stand-height-1}{\passthrough{\lstinline!set\_stand\_height!}}
&
\passthrough{\lstinline!def set\_stand\_height(self, stand\_height: float) -> int!}
& \passthrough{\lstinline!SIGNATURE\_ONLY!} &
\passthrough{\lstinline!MOTION\_COMMAND!} \\
\protect\hyperlink{unitree-sdk2-cpp-robot-h2-lococlient-set-velocity-1}{\passthrough{\lstinline!set\_velocity!}}
&
\passthrough{\lstinline!def set\_velocity(self, vx: float, vy: float, omega: float, duration: float = ...) -> int!}
& \passthrough{\lstinline!SIGNATURE\_ONLY!} &
\passthrough{\lstinline!MOTION\_COMMAND!} \\
\protect\hyperlink{unitree-sdk2-cpp-robot-h2-lococlient-set-task-id-1}{\passthrough{\lstinline!set\_task\_id!}}
&
\passthrough{\lstinline!def set\_task\_id(self, task\_id: int) -> int!}
& \passthrough{\lstinline!SIGNATURE\_ONLY!} &
\passthrough{\lstinline!MOTION\_COMMAND!} \\
\protect\hyperlink{unitree-sdk2-cpp-robot-h2-lococlient-set-arm-sdk-status-1}{\passthrough{\lstinline!set\_arm\_sdk\_status!}}
&
\passthrough{\lstinline!def set\_arm\_sdk\_status(self, arm\_sdk\_status: bool) -> int!}
& \passthrough{\lstinline!SIGNATURE\_ONLY!} &
\passthrough{\lstinline!MOTION\_COMMAND!} \\
\protect\hyperlink{unitree-sdk2-cpp-robot-h2-lococlient-damp-1}{\passthrough{\lstinline!damp!}}
& \passthrough{\lstinline!def damp(self) -> int!} &
\passthrough{\lstinline!SIGNATURE\_ONLY!} &
\passthrough{\lstinline!MOTION\_COMMAND!} \\
\protect\hyperlink{unitree-sdk2-cpp-robot-h2-lococlient-start-1}{\passthrough{\lstinline!start!}}
& \passthrough{\lstinline!def start(self) -> int!} &
\passthrough{\lstinline!SIGNATURE\_ONLY!} &
\passthrough{\lstinline!MOTION\_COMMAND!} \\
\protect\hyperlink{unitree-sdk2-cpp-robot-h2-lococlient-squat-1}{\passthrough{\lstinline!squat!}}
& \passthrough{\lstinline!def squat(self) -> int!} &
\passthrough{\lstinline!SIGNATURE\_ONLY!} &
\passthrough{\lstinline!MOTION\_COMMAND!} \\
\protect\hyperlink{unitree-sdk2-cpp-robot-h2-lococlient-sit-1}{\passthrough{\lstinline!sit!}}
& \passthrough{\lstinline!def sit(self) -> int!} &
\passthrough{\lstinline!SIGNATURE\_ONLY!} &
\passthrough{\lstinline!MOTION\_COMMAND!} \\
\protect\hyperlink{unitree-sdk2-cpp-robot-h2-lococlient-stand-up-1}{\passthrough{\lstinline!stand\_up!}}
& \passthrough{\lstinline!def stand\_up(self) -> int!} &
\passthrough{\lstinline!SIGNATURE\_ONLY!} &
\passthrough{\lstinline!MOTION\_COMMAND!} \\
\protect\hyperlink{unitree-sdk2-cpp-robot-h2-lococlient-zero-torque-1}{\passthrough{\lstinline!zero\_torque!}}
& \passthrough{\lstinline!def zero\_torque(self) -> int!} &
\passthrough{\lstinline!SIGNATURE\_ONLY!} &
\passthrough{\lstinline!MOTION\_COMMAND!} \\
\protect\hyperlink{unitree-sdk2-cpp-robot-h2-lococlient-stop-move-1}{\passthrough{\lstinline!stop\_move!}}
& \passthrough{\lstinline!def stop\_move(self) -> int!} &
\passthrough{\lstinline!SIGNATURE\_ONLY!} &
\passthrough{\lstinline!MOTION\_COMMAND!} \\
\protect\hyperlink{unitree-sdk2-cpp-robot-h2-lococlient-high-stand-1}{\passthrough{\lstinline!high\_stand!}}
& \passthrough{\lstinline!def high\_stand(self) -> int!} &
\passthrough{\lstinline!SIGNATURE\_ONLY!} &
\passthrough{\lstinline!MOTION\_COMMAND!} \\
\protect\hyperlink{unitree-sdk2-cpp-robot-h2-lococlient-low-stand-1}{\passthrough{\lstinline!low\_stand!}}
& \passthrough{\lstinline!def low\_stand(self) -> int!} &
\passthrough{\lstinline!SIGNATURE\_ONLY!} &
\passthrough{\lstinline!MOTION\_COMMAND!} \\
\protect\hyperlink{unitree-sdk2-cpp-robot-h2-lococlient-move-1}{\passthrough{\lstinline!move（重载 1/2）!}}
&
\passthrough{\lstinline!def move(self, vx: float, vy: float, vyaw: float, continous\_move: bool) -> int!}
& \passthrough{\lstinline!SIGNATURE\_ONLY!} &
\passthrough{\lstinline!MOTION\_COMMAND!} \\
\protect\hyperlink{unitree-sdk2-cpp-robot-h2-lococlient-move-2}{\passthrough{\lstinline!move（重载 2/2）!}}
&
\passthrough{\lstinline!def move(self, vx: float, vy: float, vyaw: float) -> int!}
& \passthrough{\lstinline!SIGNATURE\_ONLY!} &
\passthrough{\lstinline!MOTION\_COMMAND!} \\
\protect\hyperlink{unitree-sdk2-cpp-robot-h2-lococlient-balance-stand-1}{\passthrough{\lstinline!balance\_stand!}}
& \passthrough{\lstinline!def balance\_stand(self) -> int!} &
\passthrough{\lstinline!SIGNATURE\_ONLY!} &
\passthrough{\lstinline!MOTION\_COMMAND!} \\
\protect\hyperlink{unitree-sdk2-cpp-robot-h2-lococlient-continuous-gait-1}{\passthrough{\lstinline!continuous\_gait!}}
&
\passthrough{\lstinline!def continuous\_gait(self, flag: bool) -> int!}
& \passthrough{\lstinline!SIGNATURE\_ONLY!} &
\passthrough{\lstinline!MOTION\_COMMAND!} \\
\protect\hyperlink{unitree-sdk2-cpp-robot-h2-lococlient-switch-move-mode-1}{\passthrough{\lstinline!switch\_move\_mode!}}
&
\passthrough{\lstinline!def switch\_move\_mode(self, flag: bool) -> int!}
& \passthrough{\lstinline!SIGNATURE\_ONLY!} &
\passthrough{\lstinline!MOTION\_COMMAND!} \\
\protect\hyperlink{unitree-sdk2-cpp-robot-h2-lococlient-wave-hand-1}{\passthrough{\lstinline!wave\_hand!}}
&
\passthrough{\lstinline!def wave\_hand(self, turn\_flag: bool = ...) -> int!}
& \passthrough{\lstinline!SIGNATURE\_ONLY!} &
\passthrough{\lstinline!MOTION\_COMMAND!} \\
\protect\hyperlink{unitree-sdk2-cpp-robot-h2-lococlient-shake-hand-1}{\passthrough{\lstinline!shake\_hand!}}
&
\passthrough{\lstinline!def shake\_hand(self, stage: int = ...) -> int!}
& \passthrough{\lstinline!SIGNATURE\_ONLY!} &
\passthrough{\lstinline!MOTION\_COMMAND!} \\
\protect\hyperlink{unitree-sdk2-cpp-robot-h2-lococlient-set-speed-mode-1}{\passthrough{\lstinline!set\_speed\_mode!}}
&
\passthrough{\lstinline!def set\_speed\_mode(self, speed\_mode: int) -> int!}
& \passthrough{\lstinline!SIGNATURE\_ONLY!} &
\passthrough{\lstinline!MOTION\_COMMAND!} \\
\protect\hyperlink{unitree-sdk2-cpp-robot-h2-lococlient-enable-arm-sdk-1}{\passthrough{\lstinline!enable\_arm\_sdk!}}
& \passthrough{\lstinline!def enable\_arm\_sdk(self) -> int!} &
\passthrough{\lstinline!SIGNATURE\_ONLY!} &
\passthrough{\lstinline!MOTION\_COMMAND!} \\
\protect\hyperlink{unitree-sdk2-cpp-robot-h2-lococlient-disable-arm-sdk-1}{\passthrough{\lstinline!disable\_arm\_sdk!}}
& \passthrough{\lstinline!def disable\_arm\_sdk(self) -> int!} &
\passthrough{\lstinline!SIGNATURE\_ONLY!} &
\passthrough{\lstinline!MOTION\_COMMAND!} \\
\bottomrule
\end{longtable}

\hypertarget{unitree-sdk2-cpp-robot-h2-lococlient-dunder-init-1}{%
\paragraph{\texorpdfstring{\texttt{unitree\_sdk2\_cpp.robot.h2.LocoClient.\_\_init\_\_}}{unitree\_sdk2\_cpp.robot.h2.LocoClient.\_\_init\_\_}}\label{unitree-sdk2-cpp-robot-h2-lococlient-dunder-init-1}}

初始化 \passthrough{\lstinline!LocoClient!}
实例。是否能实际构造取决于下方可用性状态。

\textbf{签名}

\begin{lstlisting}[language=Python]
def __init__(self) -> None
\end{lstlisting}

\textbf{可用性与安全性}

\textbf{\passthrough{\lstinline!AVAILABLE!} /
\passthrough{\lstinline!CONSTRUCTION!}}：当前绑定源码已实现；实际执行条件仍取决于平台和对应
SDK 环境。

\textbf{参数}

无显式参数。实例方法中的 \passthrough{\lstinline!self!} 由 Python
自动传入。

\textbf{返回值}

\begin{longtable}[]{@{}
  >{\raggedright\arraybackslash}p{(\columnwidth - 4\tabcolsep) * \real{0.18}}
  >{\raggedright\arraybackslash}p{(\columnwidth - 4\tabcolsep) * \real{0.20}}
  >{\raggedright\arraybackslash}p{(\columnwidth - 4\tabcolsep) * \real{0.62}}@{}}
\toprule
位置 & 类型 & 含义 \\
\midrule
\endhead
无 & \passthrough{\lstinline!None!} &
\passthrough{\lstinline!\_\_init\_\_!}
本身不返回值；调用类对象时，在构造函数可用的前提下得到该类实例。 \\
\bottomrule
\end{longtable}

\textbf{对应 C++}

\begin{itemize}
\tightlist
\item
  类：\passthrough{\lstinline!unitree::robot::h2::LocoClient!}
\item
  签名：\passthrough{\lstinline!LocoClient()!}
\item
  绑定策略：\passthrough{\lstinline!未单独标注!}
\item
  声明位置：\passthrough{\lstinline!include/unitree/robot/h2/loco/h2\_loco\_client.hpp:14!}
\end{itemize}

\textbf{说明}

\begin{itemize}
\tightlist
\item
  示例是局部调用片段；\passthrough{\lstinline!obj!}
  和参数变量表示已经按上层业务校验并准备好的对象。
\end{itemize}

\textbf{用法}

\begin{lstlisting}[language=Python]
value = LocoClient()
\end{lstlisting}

\hypertarget{unitree-sdk2-cpp-robot-h2-lococlient-init-1}{%
\paragraph{\texorpdfstring{\texttt{unitree\_sdk2\_cpp.robot.h2.LocoClient.init}}{unitree\_sdk2\_cpp.robot.h2.LocoClient.init}}\label{unitree-sdk2-cpp-robot-h2-lococlient-init-1}}

初始化当前 SDK 对象所需的底层通道或服务资源。

\textbf{签名}

\begin{lstlisting}[language=Python]
def init(self) -> None
\end{lstlisting}

\textbf{可用性与安全性}

\textbf{\passthrough{\lstinline!AVAILABLE!} /
\passthrough{\lstinline!INITIALIZATION!}}：当前绑定源码已实现；实际执行条件仍取决于平台和对应
SDK 环境。

\textbf{参数}

无显式参数。实例方法中的 \passthrough{\lstinline!self!} 由 Python
自动传入。

\textbf{返回值}

\begin{longtable}[]{@{}
  >{\raggedright\arraybackslash}p{(\columnwidth - 4\tabcolsep) * \real{0.18}}
  >{\raggedright\arraybackslash}p{(\columnwidth - 4\tabcolsep) * \real{0.20}}
  >{\raggedright\arraybackslash}p{(\columnwidth - 4\tabcolsep) * \real{0.62}}@{}}
\toprule
位置 & 类型 & 含义 \\
\midrule
\endhead
无 & \passthrough{\lstinline!None!} & 该方法不返回 Python
值；失败可能表现为异常或底层状态变化。 \\
\bottomrule
\end{longtable}

\textbf{对应 C++}

\begin{itemize}
\tightlist
\item
  类：\passthrough{\lstinline!unitree::robot::h2::LocoClient!}
\item
  签名：\passthrough{\lstinline!Init()!}
\item
  绑定策略：\passthrough{\lstinline!DIRECT!}
\item
  声明位置：\passthrough{\lstinline!include/unitree/robot/h2/loco/h2\_loco\_client.hpp:18!}
\end{itemize}

\textbf{说明}

\begin{itemize}
\tightlist
\item
  示例是局部调用片段；\passthrough{\lstinline!obj!}
  和参数变量表示已经按上层业务校验并准备好的对象。
\end{itemize}

\textbf{用法}

\begin{lstlisting}[language=Python]
obj.init()
\end{lstlisting}

\hypertarget{unitree-sdk2-cpp-robot-h2-lococlient-get-fsm-id-1}{%
\paragraph{\texorpdfstring{\texttt{unitree\_sdk2\_cpp.robot.h2.LocoClient.get\_fsm\_id}}{unitree\_sdk2\_cpp.robot.h2.LocoClient.get\_fsm\_id}}\label{unitree-sdk2-cpp-robot-h2-lococlient-get-fsm-id-1}}

查询或检查 FSM ID。

\textbf{签名}

\begin{lstlisting}[language=Python]
def get_fsm_id(self) -> tuple[int, int]
\end{lstlisting}

\textbf{可用性与安全性}

\textbf{\passthrough{\lstinline!AVAILABLE!} /
\passthrough{\lstinline!READ\_ONLY!}}：当前绑定源码已实现只读查询。Robot
Client 的服务端查询通常仍需匹配的 Linux
扩展、DDS、网络、机器人和服务；纯本地 getter 则不一定需要全部条件。

\textbf{参数}

无显式参数。实例方法中的 \passthrough{\lstinline!self!} 由 Python
自动传入。

\textbf{返回值}

\begin{longtable}[]{@{}
  >{\raggedright\arraybackslash}p{(\columnwidth - 4\tabcolsep) * \real{0.18}}
  >{\raggedright\arraybackslash}p{(\columnwidth - 4\tabcolsep) * \real{0.20}}
  >{\raggedright\arraybackslash}p{(\columnwidth - 4\tabcolsep) * \real{0.62}}@{}}
\toprule
位置 & 类型 & 含义 \\
\midrule
\endhead
{[}0{]} \passthrough{\lstinline!status!} & \passthrough{\lstinline!int!}
& SDK 状态码；Unitree 示例通常以 \passthrough{\lstinline!0!}
表示成功，非零值需按具体服务错误码解释。 \\
{[}1{]} \passthrough{\lstinline!fsm\_id!} &
\passthrough{\lstinline!int!} & 传给该接口的 FSM ID 参数，Python 类型为
\passthrough{\lstinline!int!}。精确含义、单位、范围和枚举值需查目标型号对应头文件与协议。
对应 C++ 参数 \passthrough{\lstinline!fsm\_id: int \&!}。 \\
\bottomrule
\end{longtable}

\textbf{对应 C++}

\begin{itemize}
\tightlist
\item
  类：\passthrough{\lstinline!unitree::robot::h2::LocoClient!}
\item
  签名：\passthrough{\lstinline!GetFsmId(int \&)!}
\item
  绑定策略：\passthrough{\lstinline!OUTPUT\_WRAPPER!}
\item
  声明位置：\passthrough{\lstinline!include/unitree/robot/h2/loco/h2\_loco\_client.hpp:41!}
\end{itemize}

\textbf{说明}

\begin{itemize}
\tightlist
\item
  示例是局部调用片段；\passthrough{\lstinline!obj!}
  和参数变量表示已经按上层业务校验并准备好的对象。
\item
  C++ 的可修改输出引用不会作为 Python 入参出现，而是按 C++
  参数顺序追加到返回元组中。
\end{itemize}

\textbf{用法}

\begin{lstlisting}[language=Python]
status, fsm_id = obj.get_fsm_id()
\end{lstlisting}

\begin{quote}
{[}!NOTE{]} 只读查询仍会访问
DDS/机器人服务。必须检查返回状态码，并处理超时和网络中断。
\end{quote}

\hypertarget{unitree-sdk2-cpp-robot-h2-lococlient-get-fsm-mode-1}{%
\paragraph{\texorpdfstring{\texttt{unitree\_sdk2\_cpp.robot.h2.LocoClient.get\_fsm\_mode}}{unitree\_sdk2\_cpp.robot.h2.LocoClient.get\_fsm\_mode}}\label{unitree-sdk2-cpp-robot-h2-lococlient-get-fsm-mode-1}}

查询或检查 FSM 模式。

\textbf{签名}

\begin{lstlisting}[language=Python]
def get_fsm_mode(self) -> tuple[int, int]
\end{lstlisting}

\textbf{可用性与安全性}

\textbf{\passthrough{\lstinline!AVAILABLE!} /
\passthrough{\lstinline!READ\_ONLY!}}：当前绑定源码已实现只读查询。Robot
Client 的服务端查询通常仍需匹配的 Linux
扩展、DDS、网络、机器人和服务；纯本地 getter 则不一定需要全部条件。

\textbf{参数}

无显式参数。实例方法中的 \passthrough{\lstinline!self!} 由 Python
自动传入。

\textbf{返回值}

\begin{longtable}[]{@{}
  >{\raggedright\arraybackslash}p{(\columnwidth - 4\tabcolsep) * \real{0.18}}
  >{\raggedright\arraybackslash}p{(\columnwidth - 4\tabcolsep) * \real{0.20}}
  >{\raggedright\arraybackslash}p{(\columnwidth - 4\tabcolsep) * \real{0.62}}@{}}
\toprule
位置 & 类型 & 含义 \\
\midrule
\endhead
{[}0{]} \passthrough{\lstinline!status!} & \passthrough{\lstinline!int!}
& SDK 状态码；Unitree 示例通常以 \passthrough{\lstinline!0!}
表示成功，非零值需按具体服务错误码解释。 \\
{[}1{]} \passthrough{\lstinline!fsm\_mode!} &
\passthrough{\lstinline!int!} & 传给该接口的 FSM 模式 参数，Python
类型为
\passthrough{\lstinline!int!}。精确含义、单位、范围和枚举值需查目标型号对应头文件与协议。
对应 C++ 参数 \passthrough{\lstinline!fsm\_mode: int \&!}。 \\
\bottomrule
\end{longtable}

\textbf{对应 C++}

\begin{itemize}
\tightlist
\item
  类：\passthrough{\lstinline!unitree::robot::h2::LocoClient!}
\item
  签名：\passthrough{\lstinline!GetFsmMode(int \&)!}
\item
  绑定策略：\passthrough{\lstinline!OUTPUT\_WRAPPER!}
\item
  声明位置：\passthrough{\lstinline!include/unitree/robot/h2/loco/h2\_loco\_client.hpp:55!}
\end{itemize}

\textbf{说明}

\begin{itemize}
\tightlist
\item
  示例是局部调用片段；\passthrough{\lstinline!obj!}
  和参数变量表示已经按上层业务校验并准备好的对象。
\item
  C++ 的可修改输出引用不会作为 Python 入参出现，而是按 C++
  参数顺序追加到返回元组中。
\end{itemize}

\textbf{用法}

\begin{lstlisting}[language=Python]
status, fsm_mode = obj.get_fsm_mode()
\end{lstlisting}

\begin{quote}
{[}!NOTE{]} 只读查询仍会访问
DDS/机器人服务。必须检查返回状态码，并处理超时和网络中断。
\end{quote}

\hypertarget{unitree-sdk2-cpp-robot-h2-lococlient-get-balance-mode-1}{%
\paragraph{\texorpdfstring{\texttt{unitree\_sdk2\_cpp.robot.h2.LocoClient.get\_balance\_mode}}{unitree\_sdk2\_cpp.robot.h2.LocoClient.get\_balance\_mode}}\label{unitree-sdk2-cpp-robot-h2-lococlient-get-balance-mode-1}}

查询或检查 平衡 模式。

\textbf{签名}

\begin{lstlisting}[language=Python]
def get_balance_mode(self) -> tuple[int, int]
\end{lstlisting}

\textbf{可用性与安全性}

\textbf{\passthrough{\lstinline!AVAILABLE!} /
\passthrough{\lstinline!READ\_ONLY!}}：当前绑定源码已实现只读查询。Robot
Client 的服务端查询通常仍需匹配的 Linux
扩展、DDS、网络、机器人和服务；纯本地 getter 则不一定需要全部条件。

\textbf{参数}

无显式参数。实例方法中的 \passthrough{\lstinline!self!} 由 Python
自动传入。

\textbf{返回值}

\begin{longtable}[]{@{}
  >{\raggedright\arraybackslash}p{(\columnwidth - 4\tabcolsep) * \real{0.18}}
  >{\raggedright\arraybackslash}p{(\columnwidth - 4\tabcolsep) * \real{0.20}}
  >{\raggedright\arraybackslash}p{(\columnwidth - 4\tabcolsep) * \real{0.62}}@{}}
\toprule
位置 & 类型 & 含义 \\
\midrule
\endhead
{[}0{]} \passthrough{\lstinline!status!} & \passthrough{\lstinline!int!}
& SDK 状态码；Unitree 示例通常以 \passthrough{\lstinline!0!}
表示成功，非零值需按具体服务错误码解释。 \\
{[}1{]} \passthrough{\lstinline!balance\_mode!} &
\passthrough{\lstinline!int!} & 传给该接口的 平衡 模式 参数，Python
类型为
\passthrough{\lstinline!int!}。精确含义、单位、范围和枚举值需查目标型号对应头文件与协议。
对应 C++ 参数 \passthrough{\lstinline!balance\_mode: int \&!}。 \\
\bottomrule
\end{longtable}

\textbf{对应 C++}

\begin{itemize}
\tightlist
\item
  类：\passthrough{\lstinline!unitree::robot::h2::LocoClient!}
\item
  签名：\passthrough{\lstinline!GetBalanceMode(int \&)!}
\item
  绑定策略：\passthrough{\lstinline!OUTPUT\_WRAPPER!}
\item
  声明位置：\passthrough{\lstinline!include/unitree/robot/h2/loco/h2\_loco\_client.hpp:69!}
\end{itemize}

\textbf{说明}

\begin{itemize}
\tightlist
\item
  示例是局部调用片段；\passthrough{\lstinline!obj!}
  和参数变量表示已经按上层业务校验并准备好的对象。
\item
  C++ 的可修改输出引用不会作为 Python 入参出现，而是按 C++
  参数顺序追加到返回元组中。
\end{itemize}

\textbf{用法}

\begin{lstlisting}[language=Python]
status, balance_mode = obj.get_balance_mode()
\end{lstlisting}

\begin{quote}
{[}!NOTE{]} 只读查询仍会访问
DDS/机器人服务。必须检查返回状态码，并处理超时和网络中断。
\end{quote}

\hypertarget{unitree-sdk2-cpp-robot-h2-lococlient-get-swing-height-1}{%
\paragraph{\texorpdfstring{\texttt{unitree\_sdk2\_cpp.robot.h2.LocoClient.get\_swing\_height}}{unitree\_sdk2\_cpp.robot.h2.LocoClient.get\_swing\_height}}\label{unitree-sdk2-cpp-robot-h2-lococlient-get-swing-height-1}}

查询或检查 \passthrough{\lstinline!swing!} 高度。

\textbf{签名}

\begin{lstlisting}[language=Python]
def get_swing_height(self) -> tuple[int, float]
\end{lstlisting}

\textbf{可用性与安全性}

\textbf{\passthrough{\lstinline!AVAILABLE!} /
\passthrough{\lstinline!READ\_ONLY!}}：当前绑定源码已实现只读查询。Robot
Client 的服务端查询通常仍需匹配的 Linux
扩展、DDS、网络、机器人和服务；纯本地 getter 则不一定需要全部条件。

\textbf{参数}

无显式参数。实例方法中的 \passthrough{\lstinline!self!} 由 Python
自动传入。

\textbf{返回值}

\begin{longtable}[]{@{}
  >{\raggedright\arraybackslash}p{(\columnwidth - 4\tabcolsep) * \real{0.18}}
  >{\raggedright\arraybackslash}p{(\columnwidth - 4\tabcolsep) * \real{0.20}}
  >{\raggedright\arraybackslash}p{(\columnwidth - 4\tabcolsep) * \real{0.62}}@{}}
\toprule
位置 & 类型 & 含义 \\
\midrule
\endhead
{[}0{]} \passthrough{\lstinline!status!} & \passthrough{\lstinline!int!}
& SDK 状态码；Unitree 示例通常以 \passthrough{\lstinline!0!}
表示成功，非零值需按具体服务错误码解释。 \\
{[}1{]} \passthrough{\lstinline!swing\_height!} &
\passthrough{\lstinline!float!} & 传给该接口的
\passthrough{\lstinline!swing!} 高度 参数，Python 类型为
\passthrough{\lstinline!float!}。精确含义、单位、范围和枚举值需查目标型号对应头文件与协议。
对应 C++ 参数 \passthrough{\lstinline!swing\_height: float \&!}。
底层为浮点数；类型本身不说明单位、坐标系或业务安全范围。 \\
\bottomrule
\end{longtable}

\textbf{对应 C++}

\begin{itemize}
\tightlist
\item
  类：\passthrough{\lstinline!unitree::robot::h2::LocoClient!}
\item
  签名：\passthrough{\lstinline!GetSwingHeight(float \&)!}
\item
  绑定策略：\passthrough{\lstinline!OUTPUT\_WRAPPER!}
\item
  声明位置：\passthrough{\lstinline!include/unitree/robot/h2/loco/h2\_loco\_client.hpp:83!}
\end{itemize}

\textbf{说明}

\begin{itemize}
\tightlist
\item
  示例是局部调用片段；\passthrough{\lstinline!obj!}
  和参数变量表示已经按上层业务校验并准备好的对象。
\item
  C++ 的可修改输出引用不会作为 Python 入参出现，而是按 C++
  参数顺序追加到返回元组中。
\end{itemize}

\textbf{用法}

\begin{lstlisting}[language=Python]
status, swing_height = obj.get_swing_height()
\end{lstlisting}

\begin{quote}
{[}!NOTE{]} 只读查询仍会访问
DDS/机器人服务。必须检查返回状态码，并处理超时和网络中断。
\end{quote}

\hypertarget{unitree-sdk2-cpp-robot-h2-lococlient-get-stand-height-1}{%
\paragraph{\texorpdfstring{\texttt{unitree\_sdk2\_cpp.robot.h2.LocoClient.get\_stand\_height}}{unitree\_sdk2\_cpp.robot.h2.LocoClient.get\_stand\_height}}\label{unitree-sdk2-cpp-robot-h2-lococlient-get-stand-height-1}}

查询或检查 \passthrough{\lstinline!stand!} 高度。

\textbf{签名}

\begin{lstlisting}[language=Python]
def get_stand_height(self) -> tuple[int, float]
\end{lstlisting}

\textbf{可用性与安全性}

\textbf{\passthrough{\lstinline!AVAILABLE!} /
\passthrough{\lstinline!READ\_ONLY!}}：当前绑定源码已实现只读查询。Robot
Client 的服务端查询通常仍需匹配的 Linux
扩展、DDS、网络、机器人和服务；纯本地 getter 则不一定需要全部条件。

\textbf{参数}

无显式参数。实例方法中的 \passthrough{\lstinline!self!} 由 Python
自动传入。

\textbf{返回值}

\begin{longtable}[]{@{}
  >{\raggedright\arraybackslash}p{(\columnwidth - 4\tabcolsep) * \real{0.18}}
  >{\raggedright\arraybackslash}p{(\columnwidth - 4\tabcolsep) * \real{0.20}}
  >{\raggedright\arraybackslash}p{(\columnwidth - 4\tabcolsep) * \real{0.62}}@{}}
\toprule
位置 & 类型 & 含义 \\
\midrule
\endhead
{[}0{]} \passthrough{\lstinline!status!} & \passthrough{\lstinline!int!}
& SDK 状态码；Unitree 示例通常以 \passthrough{\lstinline!0!}
表示成功，非零值需按具体服务错误码解释。 \\
{[}1{]} \passthrough{\lstinline!stand\_height!} &
\passthrough{\lstinline!float!} & 传给该接口的
\passthrough{\lstinline!stand!} 高度 参数，Python 类型为
\passthrough{\lstinline!float!}。精确含义、单位、范围和枚举值需查目标型号对应头文件与协议。
对应 C++ 参数 \passthrough{\lstinline!stand\_height: float \&!}。
底层为浮点数；类型本身不说明单位、坐标系或业务安全范围。 \\
\bottomrule
\end{longtable}

\textbf{对应 C++}

\begin{itemize}
\tightlist
\item
  类：\passthrough{\lstinline!unitree::robot::h2::LocoClient!}
\item
  签名：\passthrough{\lstinline!GetStandHeight(float \&)!}
\item
  绑定策略：\passthrough{\lstinline!OUTPUT\_WRAPPER!}
\item
  声明位置：\passthrough{\lstinline!include/unitree/robot/h2/loco/h2\_loco\_client.hpp:97!}
\end{itemize}

\textbf{说明}

\begin{itemize}
\tightlist
\item
  示例是局部调用片段；\passthrough{\lstinline!obj!}
  和参数变量表示已经按上层业务校验并准备好的对象。
\item
  C++ 的可修改输出引用不会作为 Python 入参出现，而是按 C++
  参数顺序追加到返回元组中。
\end{itemize}

\textbf{用法}

\begin{lstlisting}[language=Python]
status, stand_height = obj.get_stand_height()
\end{lstlisting}

\begin{quote}
{[}!NOTE{]} 只读查询仍会访问
DDS/机器人服务。必须检查返回状态码，并处理超时和网络中断。
\end{quote}

\hypertarget{unitree-sdk2-cpp-robot-h2-lococlient-get-phase-1}{%
\paragraph{\texorpdfstring{\texttt{unitree\_sdk2\_cpp.robot.h2.LocoClient.get\_phase}}{unitree\_sdk2\_cpp.robot.h2.LocoClient.get\_phase}}\label{unitree-sdk2-cpp-robot-h2-lococlient-get-phase-1}}

查询或检查 相位。

\textbf{签名}

\begin{lstlisting}[language=Python]
def get_phase(self) -> tuple[int, list[float]]
\end{lstlisting}

\textbf{可用性与安全性}

\textbf{\passthrough{\lstinline!AVAILABLE!} /
\passthrough{\lstinline!READ\_ONLY!}}：当前绑定源码已实现只读查询。Robot
Client 的服务端查询通常仍需匹配的 Linux
扩展、DDS、网络、机器人和服务；纯本地 getter 则不一定需要全部条件。

\textbf{参数}

无显式参数。实例方法中的 \passthrough{\lstinline!self!} 由 Python
自动传入。

\textbf{返回值}

\begin{longtable}[]{@{}
  >{\raggedright\arraybackslash}p{(\columnwidth - 4\tabcolsep) * \real{0.18}}
  >{\raggedright\arraybackslash}p{(\columnwidth - 4\tabcolsep) * \real{0.20}}
  >{\raggedright\arraybackslash}p{(\columnwidth - 4\tabcolsep) * \real{0.62}}@{}}
\toprule
位置 & 类型 & 含义 \\
\midrule
\endhead
{[}0{]} \passthrough{\lstinline!status!} & \passthrough{\lstinline!int!}
& SDK 状态码；Unitree 示例通常以 \passthrough{\lstinline!0!}
表示成功，非零值需按具体服务错误码解释。 \\
{[}1{]} \passthrough{\lstinline!phase!} &
\passthrough{\lstinline!list[float]!} & 传给该接口的 相位 参数，Python
类型为
\passthrough{\lstinline!list[float]!}。精确含义、单位、范围和枚举值需查目标型号对应头文件与协议。
对应 C++ 参数 \passthrough{\lstinline!phase: std::vector<float> \&!}。
底层是可变长度
vector；元素约束：底层为浮点数；类型本身不说明单位、坐标系或业务安全范围。 \\
\bottomrule
\end{longtable}

\textbf{对应 C++}

\begin{itemize}
\tightlist
\item
  类：\passthrough{\lstinline!unitree::robot::h2::LocoClient!}
\item
  签名：\passthrough{\lstinline!GetPhase(std::vector<float> \&)!}
\item
  绑定策略：\passthrough{\lstinline!OUTPUT\_WRAPPER!}
\item
  声明位置：\passthrough{\lstinline!include/unitree/robot/h2/loco/h2\_loco\_client.hpp:111!}
\end{itemize}

\textbf{说明}

\begin{itemize}
\tightlist
\item
  示例是局部调用片段；\passthrough{\lstinline!obj!}
  和参数变量表示已经按上层业务校验并准备好的对象。
\item
  C++ 的可修改输出引用不会作为 Python 入参出现，而是按 C++
  参数顺序追加到返回元组中。
\end{itemize}

\textbf{用法}

\begin{lstlisting}[language=Python]
status, phase = obj.get_phase()
\end{lstlisting}

\begin{quote}
{[}!NOTE{]} 只读查询仍会访问
DDS/机器人服务。必须检查返回状态码，并处理超时和网络中断。
\end{quote}

\hypertarget{unitree-sdk2-cpp-robot-h2-lococlient-get-arm-sdk-status-1}{%
\paragraph{\texorpdfstring{\texttt{unitree\_sdk2\_cpp.robot.h2.LocoClient.get\_arm\_sdk\_status}}{unitree\_sdk2\_cpp.robot.h2.LocoClient.get\_arm\_sdk\_status}}\label{unitree-sdk2-cpp-robot-h2-lococlient-get-arm-sdk-status-1}}

查询或检查 机械臂 SDK 状态。

\textbf{签名}

\begin{lstlisting}[language=Python]
def get_arm_sdk_status(self) -> tuple[int, bool]
\end{lstlisting}

\textbf{可用性与安全性}

\textbf{\passthrough{\lstinline!AVAILABLE!} /
\passthrough{\lstinline!READ\_ONLY!}}：当前绑定源码已实现只读查询。Robot
Client 的服务端查询通常仍需匹配的 Linux
扩展、DDS、网络、机器人和服务；纯本地 getter 则不一定需要全部条件。

\textbf{参数}

无显式参数。实例方法中的 \passthrough{\lstinline!self!} 由 Python
自动传入。

\textbf{返回值}

\begin{longtable}[]{@{}
  >{\raggedright\arraybackslash}p{(\columnwidth - 4\tabcolsep) * \real{0.18}}
  >{\raggedright\arraybackslash}p{(\columnwidth - 4\tabcolsep) * \real{0.20}}
  >{\raggedright\arraybackslash}p{(\columnwidth - 4\tabcolsep) * \real{0.62}}@{}}
\toprule
位置 & 类型 & 含义 \\
\midrule
\endhead
{[}0{]} \passthrough{\lstinline!status!} & \passthrough{\lstinline!int!}
& SDK 状态码；Unitree 示例通常以 \passthrough{\lstinline!0!}
表示成功，非零值需按具体服务错误码解释。 \\
{[}1{]} \passthrough{\lstinline!arm\_sdk\_status!} &
\passthrough{\lstinline!bool!} & 传给该接口的 机械臂 SDK 状态
参数，Python 类型为
\passthrough{\lstinline!bool!}。精确含义、单位、范围和枚举值需查目标型号对应头文件与协议。
对应 C++ 参数 \passthrough{\lstinline!arm\_sdk\_status: bool \&!}。
只接受布尔语义。 \\
\bottomrule
\end{longtable}

\textbf{对应 C++}

\begin{itemize}
\tightlist
\item
  类：\passthrough{\lstinline!unitree::robot::h2::LocoClient!}
\item
  签名：\passthrough{\lstinline!GetArmSdkStatus(bool \&)!}
\item
  绑定策略：\passthrough{\lstinline!OUTPUT\_WRAPPER!}
\item
  声明位置：\passthrough{\lstinline!include/unitree/robot/h2/loco/h2\_loco\_client.hpp:125!}
\end{itemize}

\textbf{说明}

\begin{itemize}
\tightlist
\item
  示例是局部调用片段；\passthrough{\lstinline!obj!}
  和参数变量表示已经按上层业务校验并准备好的对象。
\item
  C++ 的可修改输出引用不会作为 Python 入参出现，而是按 C++
  参数顺序追加到返回元组中。
\end{itemize}

\textbf{用法}

\begin{lstlisting}[language=Python]
status, arm_sdk_status = obj.get_arm_sdk_status()
\end{lstlisting}

\begin{quote}
{[}!NOTE{]} 只读查询仍会访问
DDS/机器人服务。必须检查返回状态码，并处理超时和网络中断。
\end{quote}

\hypertarget{unitree-sdk2-cpp-robot-h2-lococlient-get-available-fsm-ids-1}{%
\paragraph{\texorpdfstring{\texttt{unitree\_sdk2\_cpp.robot.h2.LocoClient.get\_available\_fsm\_ids}}{unitree\_sdk2\_cpp.robot.h2.LocoClient.get\_available\_fsm\_ids}}\label{unitree-sdk2-cpp-robot-h2-lococlient-get-available-fsm-ids-1}}

查询或检查 \passthrough{\lstinline!available!} FSM ID 列表。

\textbf{签名}

\begin{lstlisting}[language=Python]
def get_available_fsm_ids(self) -> tuple[int, list[int], list[str]]
\end{lstlisting}

\textbf{可用性与安全性}

\textbf{\passthrough{\lstinline!AVAILABLE!} /
\passthrough{\lstinline!READ\_ONLY!}}：当前绑定源码已实现只读查询。Robot
Client 的服务端查询通常仍需匹配的 Linux
扩展、DDS、网络、机器人和服务；纯本地 getter 则不一定需要全部条件。

\textbf{参数}

无显式参数。实例方法中的 \passthrough{\lstinline!self!} 由 Python
自动传入。

\textbf{返回值}

\begin{longtable}[]{@{}
  >{\raggedright\arraybackslash}p{(\columnwidth - 4\tabcolsep) * \real{0.18}}
  >{\raggedright\arraybackslash}p{(\columnwidth - 4\tabcolsep) * \real{0.20}}
  >{\raggedright\arraybackslash}p{(\columnwidth - 4\tabcolsep) * \real{0.62}}@{}}
\toprule
位置 & 类型 & 含义 \\
\midrule
\endhead
{[}0{]} \passthrough{\lstinline!status!} & \passthrough{\lstinline!int!}
& SDK 状态码；Unitree 示例通常以 \passthrough{\lstinline!0!}
表示成功，非零值需按具体服务错误码解释。 \\
{[}1{]} \passthrough{\lstinline!ids!} &
\passthrough{\lstinline!list[int]!} & 传给该接口的 ID 列表 参数，Python
类型为
\passthrough{\lstinline!list[int]!}。精确含义、单位、范围和枚举值需查目标型号对应头文件与协议。
对应 C++ 参数 \passthrough{\lstinline!ids: std::vector<int> \&!}。
底层是可变长度 vector \\
{[}2{]} \passthrough{\lstinline!names!} &
\passthrough{\lstinline!list[str]!} & 传给该接口的 名称列表 参数，Python
类型为
\passthrough{\lstinline!list[str]!}。精确含义、单位、范围和枚举值需查目标型号对应头文件与协议。
对应 C++ 参数
\passthrough{\lstinline!names: std::vector<std::string> \&!}。
底层是可变长度
vector；元素约束：底层为字符串；长度、编码和允许值由具体协议决定。 \\
\bottomrule
\end{longtable}

\textbf{对应 C++}

\begin{itemize}
\tightlist
\item
  类：\passthrough{\lstinline!unitree::robot::h2::LocoClient!}
\item
  签名：\passthrough{\lstinline!GetAvailableFsmIds(std::vector<int> \&, std::vector<std::string> \&)!}
\item
  绑定策略：\passthrough{\lstinline!OUTPUT\_WRAPPER!}
\item
  声明位置：\passthrough{\lstinline!include/unitree/robot/h2/loco/h2\_loco\_client.hpp:139!}
\end{itemize}

\textbf{说明}

\begin{itemize}
\tightlist
\item
  示例是局部调用片段；\passthrough{\lstinline!obj!}
  和参数变量表示已经按上层业务校验并准备好的对象。
\item
  C++ 的可修改输出引用不会作为 Python 入参出现，而是按 C++
  参数顺序追加到返回元组中。
\end{itemize}

\textbf{用法}

\begin{lstlisting}[language=Python]
status, ids, names = obj.get_available_fsm_ids()
\end{lstlisting}

\begin{quote}
{[}!NOTE{]} 只读查询仍会访问
DDS/机器人服务。必须检查返回状态码，并处理超时和网络中断。
\end{quote}

\hypertarget{unitree-sdk2-cpp-robot-h2-lococlient-set-fsm-id-1}{%
\paragraph{\texorpdfstring{\texttt{unitree\_sdk2\_cpp.robot.h2.LocoClient.set\_fsm\_id}}{unitree\_sdk2\_cpp.robot.h2.LocoClient.set\_fsm\_id}}\label{unitree-sdk2-cpp-robot-h2-lococlient-set-fsm-id-1}}

设置 FSM ID。具体副作用和安全边界见下方状态。

\textbf{签名}

\begin{lstlisting}[language=Python]
def set_fsm_id(self, fsm_id: int) -> int
\end{lstlisting}

\textbf{可用性与安全性}

\textbf{\passthrough{\lstinline!SIGNATURE\_ONLY!} /
\passthrough{\lstinline!MOTION\_COMMAND!}}：仅用于补全和静态检查。当前二进制没有该
Python 方法；不得按可执行运动接口使用。

\textbf{参数}

\begin{longtable}[]{@{}
  >{\raggedright\arraybackslash}p{(\columnwidth - 6\tabcolsep) * \real{0.18}}
  >{\raggedright\arraybackslash}p{(\columnwidth - 6\tabcolsep) * \real{0.24}}
  >{\raggedright\arraybackslash}p{(\columnwidth - 6\tabcolsep) * \real{0.18}}
  >{\raggedright\arraybackslash}p{(\columnwidth - 6\tabcolsep) * \real{0.40}}@{}}
\toprule
名称 & Python 类型 & 默认值 & 含义 \\
\midrule
\endhead
\passthrough{\lstinline!fsm\_id!} & \passthrough{\lstinline!int!} & 必填
& 传给该接口的 FSM ID 参数，Python 类型为
\passthrough{\lstinline!int!}。精确含义、单位、范围和枚举值需查目标型号对应头文件与协议。
对应 C++ 参数 \passthrough{\lstinline!fsm\_id: int!}。 \\
\bottomrule
\end{longtable}

\textbf{返回值}

\begin{longtable}[]{@{}
  >{\raggedright\arraybackslash}p{(\columnwidth - 4\tabcolsep) * \real{0.18}}
  >{\raggedright\arraybackslash}p{(\columnwidth - 4\tabcolsep) * \real{0.20}}
  >{\raggedright\arraybackslash}p{(\columnwidth - 4\tabcolsep) * \real{0.62}}@{}}
\toprule
位置 & 类型 & 含义 \\
\midrule
\endhead
返回值 & \passthrough{\lstinline!int!} & SDK 状态码；Unitree 示例通常以
\passthrough{\lstinline!0!} 表示成功，非零值需按具体服务定义解释。 \\
\bottomrule
\end{longtable}

\textbf{对应 C++}

\begin{itemize}
\tightlist
\item
  类：\passthrough{\lstinline!unitree::robot::h2::LocoClient!}
\item
  签名：\passthrough{\lstinline!SetFsmId(int)!}
\item
  绑定策略：\passthrough{\lstinline!DIRECT!}
\item
  声明位置：\passthrough{\lstinline!include/unitree/robot/h2/loco/h2\_loco\_client.hpp:158!}
\end{itemize}

\textbf{说明}

\begin{itemize}
\tightlist
\item
  示例是局部调用片段；\passthrough{\lstinline!obj!}
  和参数变量表示已经按上层业务校验并准备好的对象。
\item
  该条目可以被编辑器和类型检查器识别，但当前运行时可能在属性查找阶段直接失败。
\item
  该方法属于运动命令。方法名包含
  \passthrough{\lstinline!stop!}、\passthrough{\lstinline!damp!} 或
  \passthrough{\lstinline!zero!} 也不代表它是独立物理急停。
\end{itemize}

\textbf{用法}

\begin{lstlisting}[language=Python]
from typing import TYPE_CHECKING

if TYPE_CHECKING:
    def planned_set_fsm_id(obj: LocoClient, fsm_id: int) -> int:
        return obj.set_fsm_id(fsm_id=fsm_id)
\end{lstlisting}

\begin{quote}
{[}!CAUTION{]} 上面的 \passthrough{\lstinline!TYPE\_CHECKING!}
示例只表达计划调用形态。该分支在普通 Python
运行时为假，不代表当前扩展能执行此接口。
\end{quote}

\hypertarget{unitree-sdk2-cpp-robot-h2-lococlient-set-balance-mode-1}{%
\paragraph{\texorpdfstring{\texttt{unitree\_sdk2\_cpp.robot.h2.LocoClient.set\_balance\_mode}}{unitree\_sdk2\_cpp.robot.h2.LocoClient.set\_balance\_mode}}\label{unitree-sdk2-cpp-robot-h2-lococlient-set-balance-mode-1}}

设置 平衡 模式。具体副作用和安全边界见下方状态。

\textbf{签名}

\begin{lstlisting}[language=Python]
def set_balance_mode(self, balance_mode: int) -> int
\end{lstlisting}

\textbf{可用性与安全性}

\textbf{\passthrough{\lstinline!SIGNATURE\_ONLY!} /
\passthrough{\lstinline!MOTION\_COMMAND!}}：仅用于补全和静态检查。当前二进制没有该
Python 方法；不得按可执行运动接口使用。

\textbf{参数}

\begin{longtable}[]{@{}
  >{\raggedright\arraybackslash}p{(\columnwidth - 6\tabcolsep) * \real{0.18}}
  >{\raggedright\arraybackslash}p{(\columnwidth - 6\tabcolsep) * \real{0.24}}
  >{\raggedright\arraybackslash}p{(\columnwidth - 6\tabcolsep) * \real{0.18}}
  >{\raggedright\arraybackslash}p{(\columnwidth - 6\tabcolsep) * \real{0.40}}@{}}
\toprule
名称 & Python 类型 & 默认值 & 含义 \\
\midrule
\endhead
\passthrough{\lstinline!balance\_mode!} & \passthrough{\lstinline!int!}
& 必填 & 传给该接口的 平衡 模式 参数，Python 类型为
\passthrough{\lstinline!int!}。精确含义、单位、范围和枚举值需查目标型号对应头文件与协议。
对应 C++ 参数 \passthrough{\lstinline!balance\_mode: int!}。 \\
\bottomrule
\end{longtable}

\textbf{返回值}

\begin{longtable}[]{@{}
  >{\raggedright\arraybackslash}p{(\columnwidth - 4\tabcolsep) * \real{0.18}}
  >{\raggedright\arraybackslash}p{(\columnwidth - 4\tabcolsep) * \real{0.20}}
  >{\raggedright\arraybackslash}p{(\columnwidth - 4\tabcolsep) * \real{0.62}}@{}}
\toprule
位置 & 类型 & 含义 \\
\midrule
\endhead
返回值 & \passthrough{\lstinline!int!} & SDK 状态码；Unitree 示例通常以
\passthrough{\lstinline!0!} 表示成功，非零值需按具体服务定义解释。 \\
\bottomrule
\end{longtable}

\textbf{对应 C++}

\begin{itemize}
\tightlist
\item
  类：\passthrough{\lstinline!unitree::robot::h2::LocoClient!}
\item
  签名：\passthrough{\lstinline!SetBalanceMode(int)!}
\item
  绑定策略：\passthrough{\lstinline!DIRECT!}
\item
  声明位置：\passthrough{\lstinline!include/unitree/robot/h2/loco/h2\_loco\_client.hpp:168!}
\end{itemize}

\textbf{说明}

\begin{itemize}
\tightlist
\item
  示例是局部调用片段；\passthrough{\lstinline!obj!}
  和参数变量表示已经按上层业务校验并准备好的对象。
\item
  该条目可以被编辑器和类型检查器识别，但当前运行时可能在属性查找阶段直接失败。
\item
  该方法属于运动命令。方法名包含
  \passthrough{\lstinline!stop!}、\passthrough{\lstinline!damp!} 或
  \passthrough{\lstinline!zero!} 也不代表它是独立物理急停。
\end{itemize}

\textbf{用法}

\begin{lstlisting}[language=Python]
from typing import TYPE_CHECKING

if TYPE_CHECKING:
    def planned_set_balance_mode(obj: LocoClient, balance_mode: int) -> int:
        return obj.set_balance_mode(balance_mode=balance_mode)
\end{lstlisting}

\begin{quote}
{[}!CAUTION{]} 上面的 \passthrough{\lstinline!TYPE\_CHECKING!}
示例只表达计划调用形态。该分支在普通 Python
运行时为假，不代表当前扩展能执行此接口。
\end{quote}

\hypertarget{unitree-sdk2-cpp-robot-h2-lococlient-set-punch-api-1}{%
\paragraph{\texorpdfstring{\texttt{unitree\_sdk2\_cpp.robot.h2.LocoClient.set\_punch\_api}}{unitree\_sdk2\_cpp.robot.h2.LocoClient.set\_punch\_api}}\label{unitree-sdk2-cpp-robot-h2-lococlient-set-punch-api-1}}

设置 \passthrough{\lstinline!punch!}
API。具体副作用和安全边界见下方状态。

\textbf{签名}

\begin{lstlisting}[language=Python]
def set_punch_api(self) -> tuple[int, list[float]]
\end{lstlisting}

\textbf{可用性与安全性}

\textbf{\passthrough{\lstinline!SIGNATURE\_ONLY!} /
\passthrough{\lstinline!MOTION\_COMMAND!}}：仅用于补全和静态检查。当前二进制没有该
Python 方法；不得按可执行运动接口使用。

\textbf{参数}

无显式参数。实例方法中的 \passthrough{\lstinline!self!} 由 Python
自动传入。

\textbf{返回值}

\begin{longtable}[]{@{}
  >{\raggedright\arraybackslash}p{(\columnwidth - 4\tabcolsep) * \real{0.18}}
  >{\raggedright\arraybackslash}p{(\columnwidth - 4\tabcolsep) * \real{0.20}}
  >{\raggedright\arraybackslash}p{(\columnwidth - 4\tabcolsep) * \real{0.62}}@{}}
\toprule
位置 & 类型 & 含义 \\
\midrule
\endhead
{[}0{]} \passthrough{\lstinline!status!} & \passthrough{\lstinline!int!}
& SDK 状态码；Unitree 示例通常以 \passthrough{\lstinline!0!}
表示成功，非零值需按具体服务错误码解释。 \\
{[}1{]} \passthrough{\lstinline!punch\_api!} &
\passthrough{\lstinline!list[float]!} & 传给该接口的
\passthrough{\lstinline!punch!} API 参数，Python 类型为
\passthrough{\lstinline!list[float]!}。精确含义、单位、范围和枚举值需查目标型号对应头文件与协议。
对应 C++ 参数
\passthrough{\lstinline!punch\_api: std::vector<float> \&!}。
底层是可变长度
vector；元素约束：底层为浮点数；类型本身不说明单位、坐标系或业务安全范围。 \\
\bottomrule
\end{longtable}

\textbf{对应 C++}

\begin{itemize}
\tightlist
\item
  类：\passthrough{\lstinline!unitree::robot::h2::LocoClient!}
\item
  签名：\passthrough{\lstinline!SetPunchApi(std::vector<float> \&)!}
\item
  绑定策略：\passthrough{\lstinline!OUTPUT\_WRAPPER!}
\item
  声明位置：\passthrough{\lstinline!include/unitree/robot/h2/loco/h2\_loco\_client.hpp:178!}
\end{itemize}

\textbf{说明}

\begin{itemize}
\tightlist
\item
  示例是局部调用片段；\passthrough{\lstinline!obj!}
  和参数变量表示已经按上层业务校验并准备好的对象。
\item
  C++ 的可修改输出引用不会作为 Python 入参出现，而是按 C++
  参数顺序追加到返回元组中。
\item
  该条目可以被编辑器和类型检查器识别，但当前运行时可能在属性查找阶段直接失败。
\item
  该方法属于运动命令。方法名包含
  \passthrough{\lstinline!stop!}、\passthrough{\lstinline!damp!} 或
  \passthrough{\lstinline!zero!} 也不代表它是独立物理急停。
\end{itemize}

\textbf{用法}

\begin{lstlisting}[language=Python]
from typing import TYPE_CHECKING

if TYPE_CHECKING:
    def planned_set_punch_api(obj: LocoClient) -> tuple[int, list[float]]:
        return obj.set_punch_api()
\end{lstlisting}

\begin{quote}
{[}!CAUTION{]} 上面的 \passthrough{\lstinline!TYPE\_CHECKING!}
示例只表达计划调用形态。该分支在普通 Python
运行时为假，不代表当前扩展能执行此接口。
\end{quote}

\hypertarget{unitree-sdk2-cpp-robot-h2-lococlient-set-swing-height-1}{%
\paragraph{\texorpdfstring{\texttt{unitree\_sdk2\_cpp.robot.h2.LocoClient.set\_swing\_height}}{unitree\_sdk2\_cpp.robot.h2.LocoClient.set\_swing\_height}}\label{unitree-sdk2-cpp-robot-h2-lococlient-set-swing-height-1}}

设置 \passthrough{\lstinline!swing!}
高度。具体副作用和安全边界见下方状态。

\textbf{签名}

\begin{lstlisting}[language=Python]
def set_swing_height(self, swing_height: float) -> int
\end{lstlisting}

\textbf{可用性与安全性}

\textbf{\passthrough{\lstinline!SIGNATURE\_ONLY!} /
\passthrough{\lstinline!MOTION\_COMMAND!}}：仅用于补全和静态检查。当前二进制没有该
Python 方法；不得按可执行运动接口使用。

\textbf{参数}

\begin{longtable}[]{@{}
  >{\raggedright\arraybackslash}p{(\columnwidth - 6\tabcolsep) * \real{0.18}}
  >{\raggedright\arraybackslash}p{(\columnwidth - 6\tabcolsep) * \real{0.24}}
  >{\raggedright\arraybackslash}p{(\columnwidth - 6\tabcolsep) * \real{0.18}}
  >{\raggedright\arraybackslash}p{(\columnwidth - 6\tabcolsep) * \real{0.40}}@{}}
\toprule
名称 & Python 类型 & 默认值 & 含义 \\
\midrule
\endhead
\passthrough{\lstinline!swing\_height!} &
\passthrough{\lstinline!float!} & 必填 & 传给该接口的
\passthrough{\lstinline!swing!} 高度 参数，Python 类型为
\passthrough{\lstinline!float!}。精确含义、单位、范围和枚举值需查目标型号对应头文件与协议。
对应 C++ 参数 \passthrough{\lstinline!swing\_height: float!}。
底层为浮点数；类型本身不说明单位、坐标系或业务安全范围。 \\
\bottomrule
\end{longtable}

\textbf{返回值}

\begin{longtable}[]{@{}
  >{\raggedright\arraybackslash}p{(\columnwidth - 4\tabcolsep) * \real{0.18}}
  >{\raggedright\arraybackslash}p{(\columnwidth - 4\tabcolsep) * \real{0.20}}
  >{\raggedright\arraybackslash}p{(\columnwidth - 4\tabcolsep) * \real{0.62}}@{}}
\toprule
位置 & 类型 & 含义 \\
\midrule
\endhead
返回值 & \passthrough{\lstinline!int!} & SDK 状态码；Unitree 示例通常以
\passthrough{\lstinline!0!} 表示成功，非零值需按具体服务定义解释。 \\
\bottomrule
\end{longtable}

\textbf{对应 C++}

\begin{itemize}
\tightlist
\item
  类：\passthrough{\lstinline!unitree::robot::h2::LocoClient!}
\item
  签名：\passthrough{\lstinline!SetSwingHeight(float)!}
\item
  绑定策略：\passthrough{\lstinline!DIRECT!}
\item
  声明位置：\passthrough{\lstinline!include/unitree/robot/h2/loco/h2\_loco\_client.hpp:189!}
\end{itemize}

\textbf{说明}

\begin{itemize}
\tightlist
\item
  示例是局部调用片段；\passthrough{\lstinline!obj!}
  和参数变量表示已经按上层业务校验并准备好的对象。
\item
  该条目可以被编辑器和类型检查器识别，但当前运行时可能在属性查找阶段直接失败。
\item
  该方法属于运动命令。方法名包含
  \passthrough{\lstinline!stop!}、\passthrough{\lstinline!damp!} 或
  \passthrough{\lstinline!zero!} 也不代表它是独立物理急停。
\end{itemize}

\textbf{用法}

\begin{lstlisting}[language=Python]
from typing import TYPE_CHECKING

if TYPE_CHECKING:
    def planned_set_swing_height(obj: LocoClient, swing_height: float) -> int:
        return obj.set_swing_height(swing_height=swing_height)
\end{lstlisting}

\begin{quote}
{[}!CAUTION{]} 上面的 \passthrough{\lstinline!TYPE\_CHECKING!}
示例只表达计划调用形态。该分支在普通 Python
运行时为假，不代表当前扩展能执行此接口。
\end{quote}

\hypertarget{unitree-sdk2-cpp-robot-h2-lococlient-set-stand-height-1}{%
\paragraph{\texorpdfstring{\texttt{unitree\_sdk2\_cpp.robot.h2.LocoClient.set\_stand\_height}}{unitree\_sdk2\_cpp.robot.h2.LocoClient.set\_stand\_height}}\label{unitree-sdk2-cpp-robot-h2-lococlient-set-stand-height-1}}

设置 \passthrough{\lstinline!stand!}
高度。具体副作用和安全边界见下方状态。

\textbf{签名}

\begin{lstlisting}[language=Python]
def set_stand_height(self, stand_height: float) -> int
\end{lstlisting}

\textbf{可用性与安全性}

\textbf{\passthrough{\lstinline!SIGNATURE\_ONLY!} /
\passthrough{\lstinline!MOTION\_COMMAND!}}：仅用于补全和静态检查。当前二进制没有该
Python 方法；不得按可执行运动接口使用。

\textbf{参数}

\begin{longtable}[]{@{}
  >{\raggedright\arraybackslash}p{(\columnwidth - 6\tabcolsep) * \real{0.18}}
  >{\raggedright\arraybackslash}p{(\columnwidth - 6\tabcolsep) * \real{0.24}}
  >{\raggedright\arraybackslash}p{(\columnwidth - 6\tabcolsep) * \real{0.18}}
  >{\raggedright\arraybackslash}p{(\columnwidth - 6\tabcolsep) * \real{0.40}}@{}}
\toprule
名称 & Python 类型 & 默认值 & 含义 \\
\midrule
\endhead
\passthrough{\lstinline!stand\_height!} &
\passthrough{\lstinline!float!} & 必填 & 传给该接口的
\passthrough{\lstinline!stand!} 高度 参数，Python 类型为
\passthrough{\lstinline!float!}。精确含义、单位、范围和枚举值需查目标型号对应头文件与协议。
对应 C++ 参数 \passthrough{\lstinline!stand\_height: float!}。
底层为浮点数；类型本身不说明单位、坐标系或业务安全范围。 \\
\bottomrule
\end{longtable}

\textbf{返回值}

\begin{longtable}[]{@{}
  >{\raggedright\arraybackslash}p{(\columnwidth - 4\tabcolsep) * \real{0.18}}
  >{\raggedright\arraybackslash}p{(\columnwidth - 4\tabcolsep) * \real{0.20}}
  >{\raggedright\arraybackslash}p{(\columnwidth - 4\tabcolsep) * \real{0.62}}@{}}
\toprule
位置 & 类型 & 含义 \\
\midrule
\endhead
返回值 & \passthrough{\lstinline!int!} & SDK 状态码；Unitree 示例通常以
\passthrough{\lstinline!0!} 表示成功，非零值需按具体服务定义解释。 \\
\bottomrule
\end{longtable}

\textbf{对应 C++}

\begin{itemize}
\tightlist
\item
  类：\passthrough{\lstinline!unitree::robot::h2::LocoClient!}
\item
  签名：\passthrough{\lstinline!SetStandHeight(float)!}
\item
  绑定策略：\passthrough{\lstinline!DIRECT!}
\item
  声明位置：\passthrough{\lstinline!include/unitree/robot/h2/loco/h2\_loco\_client.hpp:199!}
\end{itemize}

\textbf{说明}

\begin{itemize}
\tightlist
\item
  示例是局部调用片段；\passthrough{\lstinline!obj!}
  和参数变量表示已经按上层业务校验并准备好的对象。
\item
  该条目可以被编辑器和类型检查器识别，但当前运行时可能在属性查找阶段直接失败。
\item
  该方法属于运动命令。方法名包含
  \passthrough{\lstinline!stop!}、\passthrough{\lstinline!damp!} 或
  \passthrough{\lstinline!zero!} 也不代表它是独立物理急停。
\end{itemize}

\textbf{用法}

\begin{lstlisting}[language=Python]
from typing import TYPE_CHECKING

if TYPE_CHECKING:
    def planned_set_stand_height(obj: LocoClient, stand_height: float) -> int:
        return obj.set_stand_height(stand_height=stand_height)
\end{lstlisting}

\begin{quote}
{[}!CAUTION{]} 上面的 \passthrough{\lstinline!TYPE\_CHECKING!}
示例只表达计划调用形态。该分支在普通 Python
运行时为假，不代表当前扩展能执行此接口。
\end{quote}

\hypertarget{unitree-sdk2-cpp-robot-h2-lococlient-set-velocity-1}{%
\paragraph{\texorpdfstring{\texttt{unitree\_sdk2\_cpp.robot.h2.LocoClient.set\_velocity}}{unitree\_sdk2\_cpp.robot.h2.LocoClient.set\_velocity}}\label{unitree-sdk2-cpp-robot-h2-lococlient-set-velocity-1}}

设置 速度。具体副作用和安全边界见下方状态。

\textbf{签名}

\begin{lstlisting}[language=Python]
def set_velocity(self, vx: float, vy: float, omega: float, duration: float = ...) -> int
\end{lstlisting}

\textbf{可用性与安全性}

\textbf{\passthrough{\lstinline!SIGNATURE\_ONLY!} /
\passthrough{\lstinline!MOTION\_COMMAND!}}：仅用于补全和静态检查。当前二进制没有该
Python 方法；不得按可执行运动接口使用。

\textbf{参数}

\begin{longtable}[]{@{}
  >{\raggedright\arraybackslash}p{(\columnwidth - 6\tabcolsep) * \real{0.18}}
  >{\raggedright\arraybackslash}p{(\columnwidth - 6\tabcolsep) * \real{0.24}}
  >{\raggedright\arraybackslash}p{(\columnwidth - 6\tabcolsep) * \real{0.18}}
  >{\raggedright\arraybackslash}p{(\columnwidth - 6\tabcolsep) * \real{0.40}}@{}}
\toprule
名称 & Python 类型 & 默认值 & 含义 \\
\midrule
\endhead
\passthrough{\lstinline!vx!} & \passthrough{\lstinline!float!} & 必填 &
X 方向速度参数。坐标系、单位、符号和安全范围必须查目标型号运动协议。
对应 C++ 参数 \passthrough{\lstinline!vx: float!}。
底层为浮点数；类型本身不说明单位、坐标系或业务安全范围。 \\
\passthrough{\lstinline!vy!} & \passthrough{\lstinline!float!} & 必填 &
Y 方向速度参数。坐标系、单位、符号和安全范围必须查目标型号运动协议。
对应 C++ 参数 \passthrough{\lstinline!vy: float!}。
底层为浮点数；类型本身不说明单位、坐标系或业务安全范围。 \\
\passthrough{\lstinline!omega!} & \passthrough{\lstinline!float!} & 必填
& 角速度参数。旋转轴、单位、符号和安全范围必须查目标型号运动协议。 对应
C++ 参数 \passthrough{\lstinline!omega: float!}。
底层为浮点数；类型本身不说明单位、坐标系或业务安全范围。 \\
\passthrough{\lstinline!duration!} & \passthrough{\lstinline!float!} &
\passthrough{\lstinline!...!} &
操作持续时间参数。精确单位、范围和默认行为以目标型号接口定义为准。 对应
C++ 参数 \passthrough{\lstinline!duration: float!}。
底层为浮点数；类型本身不说明单位、坐标系或业务安全范围。 \\
\bottomrule
\end{longtable}

\textbf{返回值}

\begin{longtable}[]{@{}
  >{\raggedright\arraybackslash}p{(\columnwidth - 4\tabcolsep) * \real{0.18}}
  >{\raggedright\arraybackslash}p{(\columnwidth - 4\tabcolsep) * \real{0.20}}
  >{\raggedright\arraybackslash}p{(\columnwidth - 4\tabcolsep) * \real{0.62}}@{}}
\toprule
位置 & 类型 & 含义 \\
\midrule
\endhead
返回值 & \passthrough{\lstinline!int!} & SDK 状态码；Unitree 示例通常以
\passthrough{\lstinline!0!} 表示成功，非零值需按具体服务定义解释。 \\
\bottomrule
\end{longtable}

\textbf{对应 C++}

\begin{itemize}
\tightlist
\item
  类：\passthrough{\lstinline!unitree::robot::h2::LocoClient!}
\item
  签名：\passthrough{\lstinline!SetVelocity(float, float, float, float)!}
\item
  绑定策略：\passthrough{\lstinline!DIRECT!}
\item
  声明位置：\passthrough{\lstinline!include/unitree/robot/h2/loco/h2\_loco\_client.hpp:209!}
\end{itemize}

\textbf{说明}

\begin{itemize}
\tightlist
\item
  示例是局部调用片段；\passthrough{\lstinline!obj!}
  和参数变量表示已经按上层业务校验并准备好的对象。
\item
  默认值显示为 \passthrough{\lstinline!...!}，表示 C++
  声明存在默认参数，但当前 AST
  清单没有保存其字面量；省略参数可使用上游默认行为。
\item
  该条目可以被编辑器和类型检查器识别，但当前运行时可能在属性查找阶段直接失败。
\item
  该方法属于运动命令。方法名包含
  \passthrough{\lstinline!stop!}、\passthrough{\lstinline!damp!} 或
  \passthrough{\lstinline!zero!} 也不代表它是独立物理急停。
\end{itemize}

\textbf{用法}

\begin{lstlisting}[language=Python]
from typing import TYPE_CHECKING

if TYPE_CHECKING:
    def planned_set_velocity(obj: LocoClient, vx: float, vy: float, omega: float, duration: float) -> int:
        return obj.set_velocity(vx=vx, vy=vy, omega=omega, duration=duration)
\end{lstlisting}

\begin{quote}
{[}!CAUTION{]} 上面的 \passthrough{\lstinline!TYPE\_CHECKING!}
示例只表达计划调用形态。该分支在普通 Python
运行时为假，不代表当前扩展能执行此接口。
\end{quote}

\hypertarget{unitree-sdk2-cpp-robot-h2-lococlient-set-task-id-1}{%
\paragraph{\texorpdfstring{\texttt{unitree\_sdk2\_cpp.robot.h2.LocoClient.set\_task\_id}}{unitree\_sdk2\_cpp.robot.h2.LocoClient.set\_task\_id}}\label{unitree-sdk2-cpp-robot-h2-lococlient-set-task-id-1}}

设置 任务 ID。具体副作用和安全边界见下方状态。

\textbf{签名}

\begin{lstlisting}[language=Python]
def set_task_id(self, task_id: int) -> int
\end{lstlisting}

\textbf{可用性与安全性}

\textbf{\passthrough{\lstinline!SIGNATURE\_ONLY!} /
\passthrough{\lstinline!MOTION\_COMMAND!}}：仅用于补全和静态检查。当前二进制没有该
Python 方法；不得按可执行运动接口使用。

\textbf{参数}

\begin{longtable}[]{@{}
  >{\raggedright\arraybackslash}p{(\columnwidth - 6\tabcolsep) * \real{0.18}}
  >{\raggedright\arraybackslash}p{(\columnwidth - 6\tabcolsep) * \real{0.24}}
  >{\raggedright\arraybackslash}p{(\columnwidth - 6\tabcolsep) * \real{0.18}}
  >{\raggedright\arraybackslash}p{(\columnwidth - 6\tabcolsep) * \real{0.40}}@{}}
\toprule
名称 & Python 类型 & 默认值 & 含义 \\
\midrule
\endhead
\passthrough{\lstinline!task\_id!} & \passthrough{\lstinline!int!} &
必填 & 传给该接口的 任务 ID 参数，Python 类型为
\passthrough{\lstinline!int!}。精确含义、单位、范围和枚举值需查目标型号对应头文件与协议。
对应 C++ 参数 \passthrough{\lstinline!task\_id: int!}。 \\
\bottomrule
\end{longtable}

\textbf{返回值}

\begin{longtable}[]{@{}
  >{\raggedright\arraybackslash}p{(\columnwidth - 4\tabcolsep) * \real{0.18}}
  >{\raggedright\arraybackslash}p{(\columnwidth - 4\tabcolsep) * \real{0.20}}
  >{\raggedright\arraybackslash}p{(\columnwidth - 4\tabcolsep) * \real{0.62}}@{}}
\toprule
位置 & 类型 & 含义 \\
\midrule
\endhead
返回值 & \passthrough{\lstinline!int!} & SDK 状态码；Unitree 示例通常以
\passthrough{\lstinline!0!} 表示成功，非零值需按具体服务定义解释。 \\
\bottomrule
\end{longtable}

\textbf{对应 C++}

\begin{itemize}
\tightlist
\item
  类：\passthrough{\lstinline!unitree::robot::h2::LocoClient!}
\item
  签名：\passthrough{\lstinline!SetTaskId(int)!}
\item
  绑定策略：\passthrough{\lstinline!DIRECT!}
\item
  声明位置：\passthrough{\lstinline!include/unitree/robot/h2/loco/h2\_loco\_client.hpp:221!}
\end{itemize}

\textbf{说明}

\begin{itemize}
\tightlist
\item
  示例是局部调用片段；\passthrough{\lstinline!obj!}
  和参数变量表示已经按上层业务校验并准备好的对象。
\item
  该条目可以被编辑器和类型检查器识别，但当前运行时可能在属性查找阶段直接失败。
\item
  该方法属于运动命令。方法名包含
  \passthrough{\lstinline!stop!}、\passthrough{\lstinline!damp!} 或
  \passthrough{\lstinline!zero!} 也不代表它是独立物理急停。
\end{itemize}

\textbf{用法}

\begin{lstlisting}[language=Python]
from typing import TYPE_CHECKING

if TYPE_CHECKING:
    def planned_set_task_id(obj: LocoClient, task_id: int) -> int:
        return obj.set_task_id(task_id=task_id)
\end{lstlisting}

\begin{quote}
{[}!CAUTION{]} 上面的 \passthrough{\lstinline!TYPE\_CHECKING!}
示例只表达计划调用形态。该分支在普通 Python
运行时为假，不代表当前扩展能执行此接口。
\end{quote}

\hypertarget{unitree-sdk2-cpp-robot-h2-lococlient-set-arm-sdk-status-1}{%
\paragraph{\texorpdfstring{\texttt{unitree\_sdk2\_cpp.robot.h2.LocoClient.set\_arm\_sdk\_status}}{unitree\_sdk2\_cpp.robot.h2.LocoClient.set\_arm\_sdk\_status}}\label{unitree-sdk2-cpp-robot-h2-lococlient-set-arm-sdk-status-1}}

设置 机械臂 SDK 状态。具体副作用和安全边界见下方状态。

\textbf{签名}

\begin{lstlisting}[language=Python]
def set_arm_sdk_status(self, arm_sdk_status: bool) -> int
\end{lstlisting}

\textbf{可用性与安全性}

\textbf{\passthrough{\lstinline!SIGNATURE\_ONLY!} /
\passthrough{\lstinline!MOTION\_COMMAND!}}：仅用于补全和静态检查。当前二进制没有该
Python 方法；不得按可执行运动接口使用。

\textbf{参数}

\begin{longtable}[]{@{}
  >{\raggedright\arraybackslash}p{(\columnwidth - 6\tabcolsep) * \real{0.18}}
  >{\raggedright\arraybackslash}p{(\columnwidth - 6\tabcolsep) * \real{0.24}}
  >{\raggedright\arraybackslash}p{(\columnwidth - 6\tabcolsep) * \real{0.18}}
  >{\raggedright\arraybackslash}p{(\columnwidth - 6\tabcolsep) * \real{0.40}}@{}}
\toprule
名称 & Python 类型 & 默认值 & 含义 \\
\midrule
\endhead
\passthrough{\lstinline!arm\_sdk\_status!} &
\passthrough{\lstinline!bool!} & 必填 & 传给该接口的 机械臂 SDK 状态
参数，Python 类型为
\passthrough{\lstinline!bool!}。精确含义、单位、范围和枚举值需查目标型号对应头文件与协议。
对应 C++ 参数 \passthrough{\lstinline!arm\_sdk\_status: bool!}。
只接受布尔语义。 \\
\bottomrule
\end{longtable}

\textbf{返回值}

\begin{longtable}[]{@{}
  >{\raggedright\arraybackslash}p{(\columnwidth - 4\tabcolsep) * \real{0.18}}
  >{\raggedright\arraybackslash}p{(\columnwidth - 4\tabcolsep) * \real{0.20}}
  >{\raggedright\arraybackslash}p{(\columnwidth - 4\tabcolsep) * \real{0.62}}@{}}
\toprule
位置 & 类型 & 含义 \\
\midrule
\endhead
返回值 & \passthrough{\lstinline!int!} & SDK 状态码；Unitree 示例通常以
\passthrough{\lstinline!0!} 表示成功，非零值需按具体服务定义解释。 \\
\bottomrule
\end{longtable}

\textbf{对应 C++}

\begin{itemize}
\tightlist
\item
  类：\passthrough{\lstinline!unitree::robot::h2::LocoClient!}
\item
  签名：\passthrough{\lstinline!SetArmSdkStatus(bool)!}
\item
  绑定策略：\passthrough{\lstinline!DIRECT!}
\item
  声明位置：\passthrough{\lstinline!include/unitree/robot/h2/loco/h2\_loco\_client.hpp:231!}
\end{itemize}

\textbf{说明}

\begin{itemize}
\tightlist
\item
  示例是局部调用片段；\passthrough{\lstinline!obj!}
  和参数变量表示已经按上层业务校验并准备好的对象。
\item
  该条目可以被编辑器和类型检查器识别，但当前运行时可能在属性查找阶段直接失败。
\item
  该方法属于运动命令。方法名包含
  \passthrough{\lstinline!stop!}、\passthrough{\lstinline!damp!} 或
  \passthrough{\lstinline!zero!} 也不代表它是独立物理急停。
\end{itemize}

\textbf{用法}

\begin{lstlisting}[language=Python]
from typing import TYPE_CHECKING

if TYPE_CHECKING:
    def planned_set_arm_sdk_status(obj: LocoClient, arm_sdk_status: bool) -> int:
        return obj.set_arm_sdk_status(arm_sdk_status=arm_sdk_status)
\end{lstlisting}

\begin{quote}
{[}!CAUTION{]} 上面的 \passthrough{\lstinline!TYPE\_CHECKING!}
示例只表达计划调用形态。该分支在普通 Python
运行时为假，不代表当前扩展能执行此接口。
\end{quote}

\hypertarget{unitree-sdk2-cpp-robot-h2-lococlient-damp-1}{%
\paragraph{\texorpdfstring{\texttt{unitree\_sdk2\_cpp.robot.h2.LocoClient.damp}}{unitree\_sdk2\_cpp.robot.h2.LocoClient.damp}}\label{unitree-sdk2-cpp-robot-h2-lococlient-damp-1}}

对应 C++ SDK 操作
\passthrough{\lstinline!Damp()!}。上游头文件没有可直接生成的业务说明时，本参考只保证签名映射，精确语义需查目标型号协议。

\textbf{签名}

\begin{lstlisting}[language=Python]
def damp(self) -> int
\end{lstlisting}

\textbf{可用性与安全性}

\textbf{\passthrough{\lstinline!SIGNATURE\_ONLY!} /
\passthrough{\lstinline!MOTION\_COMMAND!}}：仅用于补全和静态检查。当前二进制没有该
Python 方法；不得按可执行运动接口使用。

\textbf{参数}

无显式参数。实例方法中的 \passthrough{\lstinline!self!} 由 Python
自动传入。

\textbf{返回值}

\begin{longtable}[]{@{}
  >{\raggedright\arraybackslash}p{(\columnwidth - 4\tabcolsep) * \real{0.18}}
  >{\raggedright\arraybackslash}p{(\columnwidth - 4\tabcolsep) * \real{0.20}}
  >{\raggedright\arraybackslash}p{(\columnwidth - 4\tabcolsep) * \real{0.62}}@{}}
\toprule
位置 & 类型 & 含义 \\
\midrule
\endhead
返回值 & \passthrough{\lstinline!int!} & SDK 状态码；Unitree 示例通常以
\passthrough{\lstinline!0!} 表示成功，非零值需按具体服务定义解释。 \\
\bottomrule
\end{longtable}

\textbf{对应 C++}

\begin{itemize}
\tightlist
\item
  类：\passthrough{\lstinline!unitree::robot::h2::LocoClient!}
\item
  签名：\passthrough{\lstinline!Damp()!}
\item
  绑定策略：\passthrough{\lstinline!DIRECT!}
\item
  声明位置：\passthrough{\lstinline!include/unitree/robot/h2/loco/h2\_loco\_client.hpp:242!}
\end{itemize}

\textbf{说明}

\begin{itemize}
\tightlist
\item
  示例是局部调用片段；\passthrough{\lstinline!obj!}
  和参数变量表示已经按上层业务校验并准备好的对象。
\item
  该条目可以被编辑器和类型检查器识别，但当前运行时可能在属性查找阶段直接失败。
\item
  该方法属于运动命令。方法名包含
  \passthrough{\lstinline!stop!}、\passthrough{\lstinline!damp!} 或
  \passthrough{\lstinline!zero!} 也不代表它是独立物理急停。
\end{itemize}

\textbf{用法}

\begin{lstlisting}[language=Python]
from typing import TYPE_CHECKING

if TYPE_CHECKING:
    def planned_damp(obj: LocoClient) -> int:
        return obj.damp()
\end{lstlisting}

\begin{quote}
{[}!CAUTION{]} 上面的 \passthrough{\lstinline!TYPE\_CHECKING!}
示例只表达计划调用形态。该分支在普通 Python
运行时为假，不代表当前扩展能执行此接口。
\end{quote}

\hypertarget{unitree-sdk2-cpp-robot-h2-lococlient-start-1}{%
\paragraph{\texorpdfstring{\texttt{unitree\_sdk2\_cpp.robot.h2.LocoClient.start}}{unitree\_sdk2\_cpp.robot.h2.LocoClient.start}}\label{unitree-sdk2-cpp-robot-h2-lococlient-start-1}}

启动对应 SDK 操作。具体副作用和安全边界见下方状态。

\textbf{签名}

\begin{lstlisting}[language=Python]
def start(self) -> int
\end{lstlisting}

\textbf{可用性与安全性}

\textbf{\passthrough{\lstinline!SIGNATURE\_ONLY!} /
\passthrough{\lstinline!MOTION\_COMMAND!}}：仅用于补全和静态检查。当前二进制没有该
Python 方法；不得按可执行运动接口使用。

\textbf{参数}

无显式参数。实例方法中的 \passthrough{\lstinline!self!} 由 Python
自动传入。

\textbf{返回值}

\begin{longtable}[]{@{}
  >{\raggedright\arraybackslash}p{(\columnwidth - 4\tabcolsep) * \real{0.18}}
  >{\raggedright\arraybackslash}p{(\columnwidth - 4\tabcolsep) * \real{0.20}}
  >{\raggedright\arraybackslash}p{(\columnwidth - 4\tabcolsep) * \real{0.62}}@{}}
\toprule
位置 & 类型 & 含义 \\
\midrule
\endhead
返回值 & \passthrough{\lstinline!int!} & SDK 状态码；Unitree 示例通常以
\passthrough{\lstinline!0!} 表示成功，非零值需按具体服务定义解释。 \\
\bottomrule
\end{longtable}

\textbf{对应 C++}

\begin{itemize}
\tightlist
\item
  类：\passthrough{\lstinline!unitree::robot::h2::LocoClient!}
\item
  签名：\passthrough{\lstinline!Start()!}
\item
  绑定策略：\passthrough{\lstinline!DIRECT!}
\item
  声明位置：\passthrough{\lstinline!include/unitree/robot/h2/loco/h2\_loco\_client.hpp:244!}
\end{itemize}

\textbf{说明}

\begin{itemize}
\tightlist
\item
  示例是局部调用片段；\passthrough{\lstinline!obj!}
  和参数变量表示已经按上层业务校验并准备好的对象。
\item
  该条目可以被编辑器和类型检查器识别，但当前运行时可能在属性查找阶段直接失败。
\item
  该方法属于运动命令。方法名包含
  \passthrough{\lstinline!stop!}、\passthrough{\lstinline!damp!} 或
  \passthrough{\lstinline!zero!} 也不代表它是独立物理急停。
\end{itemize}

\textbf{用法}

\begin{lstlisting}[language=Python]
from typing import TYPE_CHECKING

if TYPE_CHECKING:
    def planned_start(obj: LocoClient) -> int:
        return obj.start()
\end{lstlisting}

\begin{quote}
{[}!CAUTION{]} 上面的 \passthrough{\lstinline!TYPE\_CHECKING!}
示例只表达计划调用形态。该分支在普通 Python
运行时为假，不代表当前扩展能执行此接口。
\end{quote}

\hypertarget{unitree-sdk2-cpp-robot-h2-lococlient-squat-1}{%
\paragraph{\texorpdfstring{\texttt{unitree\_sdk2\_cpp.robot.h2.LocoClient.squat}}{unitree\_sdk2\_cpp.robot.h2.LocoClient.squat}}\label{unitree-sdk2-cpp-robot-h2-lococlient-squat-1}}

对应 C++ SDK 操作
\passthrough{\lstinline!Squat()!}。上游头文件没有可直接生成的业务说明时，本参考只保证签名映射，精确语义需查目标型号协议。

\textbf{签名}

\begin{lstlisting}[language=Python]
def squat(self) -> int
\end{lstlisting}

\textbf{可用性与安全性}

\textbf{\passthrough{\lstinline!SIGNATURE\_ONLY!} /
\passthrough{\lstinline!MOTION\_COMMAND!}}：仅用于补全和静态检查。当前二进制没有该
Python 方法；不得按可执行运动接口使用。

\textbf{参数}

无显式参数。实例方法中的 \passthrough{\lstinline!self!} 由 Python
自动传入。

\textbf{返回值}

\begin{longtable}[]{@{}
  >{\raggedright\arraybackslash}p{(\columnwidth - 4\tabcolsep) * \real{0.18}}
  >{\raggedright\arraybackslash}p{(\columnwidth - 4\tabcolsep) * \real{0.20}}
  >{\raggedright\arraybackslash}p{(\columnwidth - 4\tabcolsep) * \real{0.62}}@{}}
\toprule
位置 & 类型 & 含义 \\
\midrule
\endhead
返回值 & \passthrough{\lstinline!int!} & SDK 状态码；Unitree 示例通常以
\passthrough{\lstinline!0!} 表示成功，非零值需按具体服务定义解释。 \\
\bottomrule
\end{longtable}

\textbf{对应 C++}

\begin{itemize}
\tightlist
\item
  类：\passthrough{\lstinline!unitree::robot::h2::LocoClient!}
\item
  签名：\passthrough{\lstinline!Squat()!}
\item
  绑定策略：\passthrough{\lstinline!DIRECT!}
\item
  声明位置：\passthrough{\lstinline!include/unitree/robot/h2/loco/h2\_loco\_client.hpp:246!}
\end{itemize}

\textbf{说明}

\begin{itemize}
\tightlist
\item
  示例是局部调用片段；\passthrough{\lstinline!obj!}
  和参数变量表示已经按上层业务校验并准备好的对象。
\item
  该条目可以被编辑器和类型检查器识别，但当前运行时可能在属性查找阶段直接失败。
\item
  该方法属于运动命令。方法名包含
  \passthrough{\lstinline!stop!}、\passthrough{\lstinline!damp!} 或
  \passthrough{\lstinline!zero!} 也不代表它是独立物理急停。
\end{itemize}

\textbf{用法}

\begin{lstlisting}[language=Python]
from typing import TYPE_CHECKING

if TYPE_CHECKING:
    def planned_squat(obj: LocoClient) -> int:
        return obj.squat()
\end{lstlisting}

\begin{quote}
{[}!CAUTION{]} 上面的 \passthrough{\lstinline!TYPE\_CHECKING!}
示例只表达计划调用形态。该分支在普通 Python
运行时为假，不代表当前扩展能执行此接口。
\end{quote}

\hypertarget{unitree-sdk2-cpp-robot-h2-lococlient-sit-1}{%
\paragraph{\texorpdfstring{\texttt{unitree\_sdk2\_cpp.robot.h2.LocoClient.sit}}{unitree\_sdk2\_cpp.robot.h2.LocoClient.sit}}\label{unitree-sdk2-cpp-robot-h2-lococlient-sit-1}}

对应 C++ SDK 操作
\passthrough{\lstinline!Sit()!}。上游头文件没有可直接生成的业务说明时，本参考只保证签名映射，精确语义需查目标型号协议。

\textbf{签名}

\begin{lstlisting}[language=Python]
def sit(self) -> int
\end{lstlisting}

\textbf{可用性与安全性}

\textbf{\passthrough{\lstinline!SIGNATURE\_ONLY!} /
\passthrough{\lstinline!MOTION\_COMMAND!}}：仅用于补全和静态检查。当前二进制没有该
Python 方法；不得按可执行运动接口使用。

\textbf{参数}

无显式参数。实例方法中的 \passthrough{\lstinline!self!} 由 Python
自动传入。

\textbf{返回值}

\begin{longtable}[]{@{}
  >{\raggedright\arraybackslash}p{(\columnwidth - 4\tabcolsep) * \real{0.18}}
  >{\raggedright\arraybackslash}p{(\columnwidth - 4\tabcolsep) * \real{0.20}}
  >{\raggedright\arraybackslash}p{(\columnwidth - 4\tabcolsep) * \real{0.62}}@{}}
\toprule
位置 & 类型 & 含义 \\
\midrule
\endhead
返回值 & \passthrough{\lstinline!int!} & SDK 状态码；Unitree 示例通常以
\passthrough{\lstinline!0!} 表示成功，非零值需按具体服务定义解释。 \\
\bottomrule
\end{longtable}

\textbf{对应 C++}

\begin{itemize}
\tightlist
\item
  类：\passthrough{\lstinline!unitree::robot::h2::LocoClient!}
\item
  签名：\passthrough{\lstinline!Sit()!}
\item
  绑定策略：\passthrough{\lstinline!DIRECT!}
\item
  声明位置：\passthrough{\lstinline!include/unitree/robot/h2/loco/h2\_loco\_client.hpp:248!}
\end{itemize}

\textbf{说明}

\begin{itemize}
\tightlist
\item
  示例是局部调用片段；\passthrough{\lstinline!obj!}
  和参数变量表示已经按上层业务校验并准备好的对象。
\item
  该条目可以被编辑器和类型检查器识别，但当前运行时可能在属性查找阶段直接失败。
\item
  该方法属于运动命令。方法名包含
  \passthrough{\lstinline!stop!}、\passthrough{\lstinline!damp!} 或
  \passthrough{\lstinline!zero!} 也不代表它是独立物理急停。
\end{itemize}

\textbf{用法}

\begin{lstlisting}[language=Python]
from typing import TYPE_CHECKING

if TYPE_CHECKING:
    def planned_sit(obj: LocoClient) -> int:
        return obj.sit()
\end{lstlisting}

\begin{quote}
{[}!CAUTION{]} 上面的 \passthrough{\lstinline!TYPE\_CHECKING!}
示例只表达计划调用形态。该分支在普通 Python
运行时为假，不代表当前扩展能执行此接口。
\end{quote}

\hypertarget{unitree-sdk2-cpp-robot-h2-lococlient-stand-up-1}{%
\paragraph{\texorpdfstring{\texttt{unitree\_sdk2\_cpp.robot.h2.LocoClient.stand\_up}}{unitree\_sdk2\_cpp.robot.h2.LocoClient.stand\_up}}\label{unitree-sdk2-cpp-robot-h2-lococlient-stand-up-1}}

对应 C++ SDK 操作
\passthrough{\lstinline!StandUp()!}。上游头文件没有可直接生成的业务说明时，本参考只保证签名映射，精确语义需查目标型号协议。

\textbf{签名}

\begin{lstlisting}[language=Python]
def stand_up(self) -> int
\end{lstlisting}

\textbf{可用性与安全性}

\textbf{\passthrough{\lstinline!SIGNATURE\_ONLY!} /
\passthrough{\lstinline!MOTION\_COMMAND!}}：仅用于补全和静态检查。当前二进制没有该
Python 方法；不得按可执行运动接口使用。

\textbf{参数}

无显式参数。实例方法中的 \passthrough{\lstinline!self!} 由 Python
自动传入。

\textbf{返回值}

\begin{longtable}[]{@{}
  >{\raggedright\arraybackslash}p{(\columnwidth - 4\tabcolsep) * \real{0.18}}
  >{\raggedright\arraybackslash}p{(\columnwidth - 4\tabcolsep) * \real{0.20}}
  >{\raggedright\arraybackslash}p{(\columnwidth - 4\tabcolsep) * \real{0.62}}@{}}
\toprule
位置 & 类型 & 含义 \\
\midrule
\endhead
返回值 & \passthrough{\lstinline!int!} & SDK 状态码；Unitree 示例通常以
\passthrough{\lstinline!0!} 表示成功，非零值需按具体服务定义解释。 \\
\bottomrule
\end{longtable}

\textbf{对应 C++}

\begin{itemize}
\tightlist
\item
  类：\passthrough{\lstinline!unitree::robot::h2::LocoClient!}
\item
  签名：\passthrough{\lstinline!StandUp()!}
\item
  绑定策略：\passthrough{\lstinline!DIRECT!}
\item
  声明位置：\passthrough{\lstinline!include/unitree/robot/h2/loco/h2\_loco\_client.hpp:250!}
\end{itemize}

\textbf{说明}

\begin{itemize}
\tightlist
\item
  示例是局部调用片段；\passthrough{\lstinline!obj!}
  和参数变量表示已经按上层业务校验并准备好的对象。
\item
  该条目可以被编辑器和类型检查器识别，但当前运行时可能在属性查找阶段直接失败。
\item
  该方法属于运动命令。方法名包含
  \passthrough{\lstinline!stop!}、\passthrough{\lstinline!damp!} 或
  \passthrough{\lstinline!zero!} 也不代表它是独立物理急停。
\end{itemize}

\textbf{用法}

\begin{lstlisting}[language=Python]
from typing import TYPE_CHECKING

if TYPE_CHECKING:
    def planned_stand_up(obj: LocoClient) -> int:
        return obj.stand_up()
\end{lstlisting}

\begin{quote}
{[}!CAUTION{]} 上面的 \passthrough{\lstinline!TYPE\_CHECKING!}
示例只表达计划调用形态。该分支在普通 Python
运行时为假，不代表当前扩展能执行此接口。
\end{quote}

\hypertarget{unitree-sdk2-cpp-robot-h2-lococlient-zero-torque-1}{%
\paragraph{\texorpdfstring{\texttt{unitree\_sdk2\_cpp.robot.h2.LocoClient.zero\_torque}}{unitree\_sdk2\_cpp.robot.h2.LocoClient.zero\_torque}}\label{unitree-sdk2-cpp-robot-h2-lococlient-zero-torque-1}}

对应 C++ SDK 操作
\passthrough{\lstinline!ZeroTorque()!}。上游头文件没有可直接生成的业务说明时，本参考只保证签名映射，精确语义需查目标型号协议。

\textbf{签名}

\begin{lstlisting}[language=Python]
def zero_torque(self) -> int
\end{lstlisting}

\textbf{可用性与安全性}

\textbf{\passthrough{\lstinline!SIGNATURE\_ONLY!} /
\passthrough{\lstinline!MOTION\_COMMAND!}}：仅用于补全和静态检查。当前二进制没有该
Python 方法；不得按可执行运动接口使用。

\textbf{参数}

无显式参数。实例方法中的 \passthrough{\lstinline!self!} 由 Python
自动传入。

\textbf{返回值}

\begin{longtable}[]{@{}
  >{\raggedright\arraybackslash}p{(\columnwidth - 4\tabcolsep) * \real{0.18}}
  >{\raggedright\arraybackslash}p{(\columnwidth - 4\tabcolsep) * \real{0.20}}
  >{\raggedright\arraybackslash}p{(\columnwidth - 4\tabcolsep) * \real{0.62}}@{}}
\toprule
位置 & 类型 & 含义 \\
\midrule
\endhead
返回值 & \passthrough{\lstinline!int!} & SDK 状态码；Unitree 示例通常以
\passthrough{\lstinline!0!} 表示成功，非零值需按具体服务定义解释。 \\
\bottomrule
\end{longtable}

\textbf{对应 C++}

\begin{itemize}
\tightlist
\item
  类：\passthrough{\lstinline!unitree::robot::h2::LocoClient!}
\item
  签名：\passthrough{\lstinline!ZeroTorque()!}
\item
  绑定策略：\passthrough{\lstinline!DIRECT!}
\item
  声明位置：\passthrough{\lstinline!include/unitree/robot/h2/loco/h2\_loco\_client.hpp:252!}
\end{itemize}

\textbf{说明}

\begin{itemize}
\tightlist
\item
  示例是局部调用片段；\passthrough{\lstinline!obj!}
  和参数变量表示已经按上层业务校验并准备好的对象。
\item
  该条目可以被编辑器和类型检查器识别，但当前运行时可能在属性查找阶段直接失败。
\item
  该方法属于运动命令。方法名包含
  \passthrough{\lstinline!stop!}、\passthrough{\lstinline!damp!} 或
  \passthrough{\lstinline!zero!} 也不代表它是独立物理急停。
\end{itemize}

\textbf{用法}

\begin{lstlisting}[language=Python]
from typing import TYPE_CHECKING

if TYPE_CHECKING:
    def planned_zero_torque(obj: LocoClient) -> int:
        return obj.zero_torque()
\end{lstlisting}

\begin{quote}
{[}!CAUTION{]} 上面的 \passthrough{\lstinline!TYPE\_CHECKING!}
示例只表达计划调用形态。该分支在普通 Python
运行时为假，不代表当前扩展能执行此接口。
\end{quote}

\hypertarget{unitree-sdk2-cpp-robot-h2-lococlient-stop-move-1}{%
\paragraph{\texorpdfstring{\texttt{unitree\_sdk2\_cpp.robot.h2.LocoClient.stop\_move}}{unitree\_sdk2\_cpp.robot.h2.LocoClient.stop\_move}}\label{unitree-sdk2-cpp-robot-h2-lococlient-stop-move-1}}

请求停止对应 SDK 操作。该名称不等同于经过验证的物理急停。

\textbf{签名}

\begin{lstlisting}[language=Python]
def stop_move(self) -> int
\end{lstlisting}

\textbf{可用性与安全性}

\textbf{\passthrough{\lstinline!SIGNATURE\_ONLY!} /
\passthrough{\lstinline!MOTION\_COMMAND!}}：仅用于补全和静态检查。当前二进制没有该
Python 方法；不得按可执行运动接口使用。

\textbf{参数}

无显式参数。实例方法中的 \passthrough{\lstinline!self!} 由 Python
自动传入。

\textbf{返回值}

\begin{longtable}[]{@{}
  >{\raggedright\arraybackslash}p{(\columnwidth - 4\tabcolsep) * \real{0.18}}
  >{\raggedright\arraybackslash}p{(\columnwidth - 4\tabcolsep) * \real{0.20}}
  >{\raggedright\arraybackslash}p{(\columnwidth - 4\tabcolsep) * \real{0.62}}@{}}
\toprule
位置 & 类型 & 含义 \\
\midrule
\endhead
返回值 & \passthrough{\lstinline!int!} & SDK 状态码；Unitree 示例通常以
\passthrough{\lstinline!0!} 表示成功，非零值需按具体服务定义解释。 \\
\bottomrule
\end{longtable}

\textbf{对应 C++}

\begin{itemize}
\tightlist
\item
  类：\passthrough{\lstinline!unitree::robot::h2::LocoClient!}
\item
  签名：\passthrough{\lstinline!StopMove()!}
\item
  绑定策略：\passthrough{\lstinline!DIRECT!}
\item
  声明位置：\passthrough{\lstinline!include/unitree/robot/h2/loco/h2\_loco\_client.hpp:254!}
\end{itemize}

\textbf{说明}

\begin{itemize}
\tightlist
\item
  示例是局部调用片段；\passthrough{\lstinline!obj!}
  和参数变量表示已经按上层业务校验并准备好的对象。
\item
  该条目可以被编辑器和类型检查器识别，但当前运行时可能在属性查找阶段直接失败。
\item
  该方法属于运动命令。方法名包含
  \passthrough{\lstinline!stop!}、\passthrough{\lstinline!damp!} 或
  \passthrough{\lstinline!zero!} 也不代表它是独立物理急停。
\end{itemize}

\textbf{用法}

\begin{lstlisting}[language=Python]
from typing import TYPE_CHECKING

if TYPE_CHECKING:
    def planned_stop_move(obj: LocoClient) -> int:
        return obj.stop_move()
\end{lstlisting}

\begin{quote}
{[}!CAUTION{]} 上面的 \passthrough{\lstinline!TYPE\_CHECKING!}
示例只表达计划调用形态。该分支在普通 Python
运行时为假，不代表当前扩展能执行此接口。
\end{quote}

\hypertarget{unitree-sdk2-cpp-robot-h2-lococlient-high-stand-1}{%
\paragraph{\texorpdfstring{\texttt{unitree\_sdk2\_cpp.robot.h2.LocoClient.high\_stand}}{unitree\_sdk2\_cpp.robot.h2.LocoClient.high\_stand}}\label{unitree-sdk2-cpp-robot-h2-lococlient-high-stand-1}}

对应 C++ SDK 操作
\passthrough{\lstinline!HighStand()!}。上游头文件没有可直接生成的业务说明时，本参考只保证签名映射，精确语义需查目标型号协议。

\textbf{签名}

\begin{lstlisting}[language=Python]
def high_stand(self) -> int
\end{lstlisting}

\textbf{可用性与安全性}

\textbf{\passthrough{\lstinline!SIGNATURE\_ONLY!} /
\passthrough{\lstinline!MOTION\_COMMAND!}}：仅用于补全和静态检查。当前二进制没有该
Python 方法；不得按可执行运动接口使用。

\textbf{参数}

无显式参数。实例方法中的 \passthrough{\lstinline!self!} 由 Python
自动传入。

\textbf{返回值}

\begin{longtable}[]{@{}
  >{\raggedright\arraybackslash}p{(\columnwidth - 4\tabcolsep) * \real{0.18}}
  >{\raggedright\arraybackslash}p{(\columnwidth - 4\tabcolsep) * \real{0.20}}
  >{\raggedright\arraybackslash}p{(\columnwidth - 4\tabcolsep) * \real{0.62}}@{}}
\toprule
位置 & 类型 & 含义 \\
\midrule
\endhead
返回值 & \passthrough{\lstinline!int!} & SDK 状态码；Unitree 示例通常以
\passthrough{\lstinline!0!} 表示成功，非零值需按具体服务定义解释。 \\
\bottomrule
\end{longtable}

\textbf{对应 C++}

\begin{itemize}
\tightlist
\item
  类：\passthrough{\lstinline!unitree::robot::h2::LocoClient!}
\item
  签名：\passthrough{\lstinline!HighStand()!}
\item
  绑定策略：\passthrough{\lstinline!DIRECT!}
\item
  声明位置：\passthrough{\lstinline!include/unitree/robot/h2/loco/h2\_loco\_client.hpp:256!}
\end{itemize}

\textbf{说明}

\begin{itemize}
\tightlist
\item
  示例是局部调用片段；\passthrough{\lstinline!obj!}
  和参数变量表示已经按上层业务校验并准备好的对象。
\item
  该条目可以被编辑器和类型检查器识别，但当前运行时可能在属性查找阶段直接失败。
\item
  该方法属于运动命令。方法名包含
  \passthrough{\lstinline!stop!}、\passthrough{\lstinline!damp!} 或
  \passthrough{\lstinline!zero!} 也不代表它是独立物理急停。
\end{itemize}

\textbf{用法}

\begin{lstlisting}[language=Python]
from typing import TYPE_CHECKING

if TYPE_CHECKING:
    def planned_high_stand(obj: LocoClient) -> int:
        return obj.high_stand()
\end{lstlisting}

\begin{quote}
{[}!CAUTION{]} 上面的 \passthrough{\lstinline!TYPE\_CHECKING!}
示例只表达计划调用形态。该分支在普通 Python
运行时为假，不代表当前扩展能执行此接口。
\end{quote}

\hypertarget{unitree-sdk2-cpp-robot-h2-lococlient-low-stand-1}{%
\paragraph{\texorpdfstring{\texttt{unitree\_sdk2\_cpp.robot.h2.LocoClient.low\_stand}}{unitree\_sdk2\_cpp.robot.h2.LocoClient.low\_stand}}\label{unitree-sdk2-cpp-robot-h2-lococlient-low-stand-1}}

对应 C++ SDK 操作
\passthrough{\lstinline!LowStand()!}。上游头文件没有可直接生成的业务说明时，本参考只保证签名映射，精确语义需查目标型号协议。

\textbf{签名}

\begin{lstlisting}[language=Python]
def low_stand(self) -> int
\end{lstlisting}

\textbf{可用性与安全性}

\textbf{\passthrough{\lstinline!SIGNATURE\_ONLY!} /
\passthrough{\lstinline!MOTION\_COMMAND!}}：仅用于补全和静态检查。当前二进制没有该
Python 方法；不得按可执行运动接口使用。

\textbf{参数}

无显式参数。实例方法中的 \passthrough{\lstinline!self!} 由 Python
自动传入。

\textbf{返回值}

\begin{longtable}[]{@{}
  >{\raggedright\arraybackslash}p{(\columnwidth - 4\tabcolsep) * \real{0.18}}
  >{\raggedright\arraybackslash}p{(\columnwidth - 4\tabcolsep) * \real{0.20}}
  >{\raggedright\arraybackslash}p{(\columnwidth - 4\tabcolsep) * \real{0.62}}@{}}
\toprule
位置 & 类型 & 含义 \\
\midrule
\endhead
返回值 & \passthrough{\lstinline!int!} & SDK 状态码；Unitree 示例通常以
\passthrough{\lstinline!0!} 表示成功，非零值需按具体服务定义解释。 \\
\bottomrule
\end{longtable}

\textbf{对应 C++}

\begin{itemize}
\tightlist
\item
  类：\passthrough{\lstinline!unitree::robot::h2::LocoClient!}
\item
  签名：\passthrough{\lstinline!LowStand()!}
\item
  绑定策略：\passthrough{\lstinline!DIRECT!}
\item
  声明位置：\passthrough{\lstinline!include/unitree/robot/h2/loco/h2\_loco\_client.hpp:258!}
\end{itemize}

\textbf{说明}

\begin{itemize}
\tightlist
\item
  示例是局部调用片段；\passthrough{\lstinline!obj!}
  和参数变量表示已经按上层业务校验并准备好的对象。
\item
  该条目可以被编辑器和类型检查器识别，但当前运行时可能在属性查找阶段直接失败。
\item
  该方法属于运动命令。方法名包含
  \passthrough{\lstinline!stop!}、\passthrough{\lstinline!damp!} 或
  \passthrough{\lstinline!zero!} 也不代表它是独立物理急停。
\end{itemize}

\textbf{用法}

\begin{lstlisting}[language=Python]
from typing import TYPE_CHECKING

if TYPE_CHECKING:
    def planned_low_stand(obj: LocoClient) -> int:
        return obj.low_stand()
\end{lstlisting}

\begin{quote}
{[}!CAUTION{]} 上面的 \passthrough{\lstinline!TYPE\_CHECKING!}
示例只表达计划调用形态。该分支在普通 Python
运行时为假，不代表当前扩展能执行此接口。
\end{quote}

\hypertarget{unitree-sdk2-cpp-robot-h2-lococlient-move-1}{%
\paragraph{\texorpdfstring{\texttt{unitree\_sdk2\_cpp.robot.h2.LocoClient.move}（重载
1/2）}{unitree\_sdk2\_cpp.robot.h2.LocoClient.move（重载 1/2）}}\label{unitree-sdk2-cpp-robot-h2-lococlient-move-1}}

对应 C++ SDK 操作
\passthrough{\lstinline!Move(float, float, float, bool)!}。上游头文件没有可直接生成的业务说明时，本参考只保证签名映射，精确语义需查目标型号协议。

\textbf{签名}

\begin{lstlisting}[language=Python]
@overload
def move(self, vx: float, vy: float, vyaw: float, continous_move: bool) -> int
\end{lstlisting}

\textbf{可用性与安全性}

\textbf{\passthrough{\lstinline!SIGNATURE\_ONLY!} /
\passthrough{\lstinline!MOTION\_COMMAND!}}：仅用于补全和静态检查。当前二进制没有该
Python 方法；不得按可执行运动接口使用。

\textbf{参数}

\begin{longtable}[]{@{}
  >{\raggedright\arraybackslash}p{(\columnwidth - 6\tabcolsep) * \real{0.18}}
  >{\raggedright\arraybackslash}p{(\columnwidth - 6\tabcolsep) * \real{0.24}}
  >{\raggedright\arraybackslash}p{(\columnwidth - 6\tabcolsep) * \real{0.18}}
  >{\raggedright\arraybackslash}p{(\columnwidth - 6\tabcolsep) * \real{0.40}}@{}}
\toprule
名称 & Python 类型 & 默认值 & 含义 \\
\midrule
\endhead
\passthrough{\lstinline!vx!} & \passthrough{\lstinline!float!} & 必填 &
X 方向速度参数。坐标系、单位、符号和安全范围必须查目标型号运动协议。
对应 C++ 参数 \passthrough{\lstinline!vx: float!}。
底层为浮点数；类型本身不说明单位、坐标系或业务安全范围。 \\
\passthrough{\lstinline!vy!} & \passthrough{\lstinline!float!} & 必填 &
Y 方向速度参数。坐标系、单位、符号和安全范围必须查目标型号运动协议。
对应 C++ 参数 \passthrough{\lstinline!vy: float!}。
底层为浮点数；类型本身不说明单位、坐标系或业务安全范围。 \\
\passthrough{\lstinline!vyaw!} & \passthrough{\lstinline!float!} & 必填
& 偏航角速度参数。单位、符号和安全范围必须查目标型号运动协议。 对应 C++
参数 \passthrough{\lstinline!vyaw: float!}。
底层为浮点数；类型本身不说明单位、坐标系或业务安全范围。 \\
\passthrough{\lstinline!continous\_move!} &
\passthrough{\lstinline!bool!} & 必填 & 传给该接口的
\passthrough{\lstinline!continous!} \passthrough{\lstinline!move!}
参数，Python 类型为
\passthrough{\lstinline!bool!}。精确含义、单位、范围和枚举值需查目标型号对应头文件与协议。
对应 C++ 参数 \passthrough{\lstinline!continous\_move: bool!}。
只接受布尔语义。 \\
\bottomrule
\end{longtable}

\textbf{返回值}

\begin{longtable}[]{@{}
  >{\raggedright\arraybackslash}p{(\columnwidth - 4\tabcolsep) * \real{0.18}}
  >{\raggedright\arraybackslash}p{(\columnwidth - 4\tabcolsep) * \real{0.20}}
  >{\raggedright\arraybackslash}p{(\columnwidth - 4\tabcolsep) * \real{0.62}}@{}}
\toprule
位置 & 类型 & 含义 \\
\midrule
\endhead
返回值 & \passthrough{\lstinline!int!} & SDK 状态码；Unitree 示例通常以
\passthrough{\lstinline!0!} 表示成功，非零值需按具体服务定义解释。 \\
\bottomrule
\end{longtable}

\textbf{对应 C++}

\begin{itemize}
\tightlist
\item
  类：\passthrough{\lstinline!unitree::robot::h2::LocoClient!}
\item
  签名：\passthrough{\lstinline!Move(float, float, float, bool)!}
\item
  绑定策略：\passthrough{\lstinline!DIRECT!}
\item
  声明位置：\passthrough{\lstinline!include/unitree/robot/h2/loco/h2\_loco\_client.hpp:260!}
\end{itemize}

\textbf{说明}

\begin{itemize}
\tightlist
\item
  示例是局部调用片段；\passthrough{\lstinline!obj!}
  和参数变量表示已经按上层业务校验并准备好的对象。
\item
  该条目可以被编辑器和类型检查器识别，但当前运行时可能在属性查找阶段直接失败。
\item
  该方法属于运动命令。方法名包含
  \passthrough{\lstinline!stop!}、\passthrough{\lstinline!damp!} 或
  \passthrough{\lstinline!zero!} 也不代表它是独立物理急停。
\end{itemize}

\textbf{用法}

\begin{lstlisting}[language=Python]
from typing import TYPE_CHECKING

if TYPE_CHECKING:
    def planned_move(obj: LocoClient, vx: float, vy: float, vyaw: float, continous_move: bool) -> int:
        return obj.move(vx=vx, vy=vy, vyaw=vyaw, continous_move=continous_move)
\end{lstlisting}

\begin{quote}
{[}!CAUTION{]} 上面的 \passthrough{\lstinline!TYPE\_CHECKING!}
示例只表达计划调用形态。该分支在普通 Python
运行时为假，不代表当前扩展能执行此接口。
\end{quote}

\hypertarget{unitree-sdk2-cpp-robot-h2-lococlient-move-2}{%
\paragraph{\texorpdfstring{\texttt{unitree\_sdk2\_cpp.robot.h2.LocoClient.move}（重载
2/2）}{unitree\_sdk2\_cpp.robot.h2.LocoClient.move（重载 2/2）}}\label{unitree-sdk2-cpp-robot-h2-lococlient-move-2}}

对应 C++ SDK 操作
\passthrough{\lstinline!Move(float, float, float)!}。上游头文件没有可直接生成的业务说明时，本参考只保证签名映射，精确语义需查目标型号协议。

\textbf{签名}

\begin{lstlisting}[language=Python]
@overload
def move(self, vx: float, vy: float, vyaw: float) -> int
\end{lstlisting}

\textbf{可用性与安全性}

\textbf{\passthrough{\lstinline!SIGNATURE\_ONLY!} /
\passthrough{\lstinline!MOTION\_COMMAND!}}：仅用于补全和静态检查。当前二进制没有该
Python 方法；不得按可执行运动接口使用。

\textbf{参数}

\begin{longtable}[]{@{}
  >{\raggedright\arraybackslash}p{(\columnwidth - 6\tabcolsep) * \real{0.18}}
  >{\raggedright\arraybackslash}p{(\columnwidth - 6\tabcolsep) * \real{0.24}}
  >{\raggedright\arraybackslash}p{(\columnwidth - 6\tabcolsep) * \real{0.18}}
  >{\raggedright\arraybackslash}p{(\columnwidth - 6\tabcolsep) * \real{0.40}}@{}}
\toprule
名称 & Python 类型 & 默认值 & 含义 \\
\midrule
\endhead
\passthrough{\lstinline!vx!} & \passthrough{\lstinline!float!} & 必填 &
X 方向速度参数。坐标系、单位、符号和安全范围必须查目标型号运动协议。
对应 C++ 参数 \passthrough{\lstinline!vx: float!}。
底层为浮点数；类型本身不说明单位、坐标系或业务安全范围。 \\
\passthrough{\lstinline!vy!} & \passthrough{\lstinline!float!} & 必填 &
Y 方向速度参数。坐标系、单位、符号和安全范围必须查目标型号运动协议。
对应 C++ 参数 \passthrough{\lstinline!vy: float!}。
底层为浮点数；类型本身不说明单位、坐标系或业务安全范围。 \\
\passthrough{\lstinline!vyaw!} & \passthrough{\lstinline!float!} & 必填
& 偏航角速度参数。单位、符号和安全范围必须查目标型号运动协议。 对应 C++
参数 \passthrough{\lstinline!vyaw: float!}。
底层为浮点数；类型本身不说明单位、坐标系或业务安全范围。 \\
\bottomrule
\end{longtable}

\textbf{返回值}

\begin{longtable}[]{@{}
  >{\raggedright\arraybackslash}p{(\columnwidth - 4\tabcolsep) * \real{0.18}}
  >{\raggedright\arraybackslash}p{(\columnwidth - 4\tabcolsep) * \real{0.20}}
  >{\raggedright\arraybackslash}p{(\columnwidth - 4\tabcolsep) * \real{0.62}}@{}}
\toprule
位置 & 类型 & 含义 \\
\midrule
\endhead
返回值 & \passthrough{\lstinline!int!} & SDK 状态码；Unitree 示例通常以
\passthrough{\lstinline!0!} 表示成功，非零值需按具体服务定义解释。 \\
\bottomrule
\end{longtable}

\textbf{对应 C++}

\begin{itemize}
\tightlist
\item
  类：\passthrough{\lstinline!unitree::robot::h2::LocoClient!}
\item
  签名：\passthrough{\lstinline!Move(float, float, float)!}
\item
  绑定策略：\passthrough{\lstinline!DIRECT!}
\item
  声明位置：\passthrough{\lstinline!include/unitree/robot/h2/loco/h2\_loco\_client.hpp:264!}
\end{itemize}

\textbf{说明}

\begin{itemize}
\tightlist
\item
  示例是局部调用片段；\passthrough{\lstinline!obj!}
  和参数变量表示已经按上层业务校验并准备好的对象。
\item
  该条目可以被编辑器和类型检查器识别，但当前运行时可能在属性查找阶段直接失败。
\item
  该方法属于运动命令。方法名包含
  \passthrough{\lstinline!stop!}、\passthrough{\lstinline!damp!} 或
  \passthrough{\lstinline!zero!} 也不代表它是独立物理急停。
\end{itemize}

\textbf{用法}

\begin{lstlisting}[language=Python]
from typing import TYPE_CHECKING

if TYPE_CHECKING:
    def planned_move(obj: LocoClient, vx: float, vy: float, vyaw: float) -> int:
        return obj.move(vx=vx, vy=vy, vyaw=vyaw)
\end{lstlisting}

\begin{quote}
{[}!CAUTION{]} 上面的 \passthrough{\lstinline!TYPE\_CHECKING!}
示例只表达计划调用形态。该分支在普通 Python
运行时为假，不代表当前扩展能执行此接口。
\end{quote}

\hypertarget{unitree-sdk2-cpp-robot-h2-lococlient-balance-stand-1}{%
\paragraph{\texorpdfstring{\texttt{unitree\_sdk2\_cpp.robot.h2.LocoClient.balance\_stand}}{unitree\_sdk2\_cpp.robot.h2.LocoClient.balance\_stand}}\label{unitree-sdk2-cpp-robot-h2-lococlient-balance-stand-1}}

对应 C++ SDK 操作
\passthrough{\lstinline!BalanceStand()!}。上游头文件没有可直接生成的业务说明时，本参考只保证签名映射，精确语义需查目标型号协议。

\textbf{签名}

\begin{lstlisting}[language=Python]
def balance_stand(self) -> int
\end{lstlisting}

\textbf{可用性与安全性}

\textbf{\passthrough{\lstinline!SIGNATURE\_ONLY!} /
\passthrough{\lstinline!MOTION\_COMMAND!}}：仅用于补全和静态检查。当前二进制没有该
Python 方法；不得按可执行运动接口使用。

\textbf{参数}

无显式参数。实例方法中的 \passthrough{\lstinline!self!} 由 Python
自动传入。

\textbf{返回值}

\begin{longtable}[]{@{}
  >{\raggedright\arraybackslash}p{(\columnwidth - 4\tabcolsep) * \real{0.18}}
  >{\raggedright\arraybackslash}p{(\columnwidth - 4\tabcolsep) * \real{0.20}}
  >{\raggedright\arraybackslash}p{(\columnwidth - 4\tabcolsep) * \real{0.62}}@{}}
\toprule
位置 & 类型 & 含义 \\
\midrule
\endhead
返回值 & \passthrough{\lstinline!int!} & SDK 状态码；Unitree 示例通常以
\passthrough{\lstinline!0!} 表示成功，非零值需按具体服务定义解释。 \\
\bottomrule
\end{longtable}

\textbf{对应 C++}

\begin{itemize}
\tightlist
\item
  类：\passthrough{\lstinline!unitree::robot::h2::LocoClient!}
\item
  签名：\passthrough{\lstinline!BalanceStand()!}
\item
  绑定策略：\passthrough{\lstinline!DIRECT!}
\item
  声明位置：\passthrough{\lstinline!include/unitree/robot/h2/loco/h2\_loco\_client.hpp:266!}
\end{itemize}

\textbf{说明}

\begin{itemize}
\tightlist
\item
  示例是局部调用片段；\passthrough{\lstinline!obj!}
  和参数变量表示已经按上层业务校验并准备好的对象。
\item
  该条目可以被编辑器和类型检查器识别，但当前运行时可能在属性查找阶段直接失败。
\item
  该方法属于运动命令。方法名包含
  \passthrough{\lstinline!stop!}、\passthrough{\lstinline!damp!} 或
  \passthrough{\lstinline!zero!} 也不代表它是独立物理急停。
\end{itemize}

\textbf{用法}

\begin{lstlisting}[language=Python]
from typing import TYPE_CHECKING

if TYPE_CHECKING:
    def planned_balance_stand(obj: LocoClient) -> int:
        return obj.balance_stand()
\end{lstlisting}

\begin{quote}
{[}!CAUTION{]} 上面的 \passthrough{\lstinline!TYPE\_CHECKING!}
示例只表达计划调用形态。该分支在普通 Python
运行时为假，不代表当前扩展能执行此接口。
\end{quote}

\hypertarget{unitree-sdk2-cpp-robot-h2-lococlient-continuous-gait-1}{%
\paragraph{\texorpdfstring{\texttt{unitree\_sdk2\_cpp.robot.h2.LocoClient.continuous\_gait}}{unitree\_sdk2\_cpp.robot.h2.LocoClient.continuous\_gait}}\label{unitree-sdk2-cpp-robot-h2-lococlient-continuous-gait-1}}

对应 C++ SDK 操作
\passthrough{\lstinline!ContinuousGait(bool)!}。上游头文件没有可直接生成的业务说明时，本参考只保证签名映射，精确语义需查目标型号协议。

\textbf{签名}

\begin{lstlisting}[language=Python]
def continuous_gait(self, flag: bool) -> int
\end{lstlisting}

\textbf{可用性与安全性}

\textbf{\passthrough{\lstinline!SIGNATURE\_ONLY!} /
\passthrough{\lstinline!MOTION\_COMMAND!}}：仅用于补全和静态检查。当前二进制没有该
Python 方法；不得按可执行运动接口使用。

\textbf{参数}

\begin{longtable}[]{@{}
  >{\raggedright\arraybackslash}p{(\columnwidth - 6\tabcolsep) * \real{0.18}}
  >{\raggedright\arraybackslash}p{(\columnwidth - 6\tabcolsep) * \real{0.24}}
  >{\raggedright\arraybackslash}p{(\columnwidth - 6\tabcolsep) * \real{0.18}}
  >{\raggedright\arraybackslash}p{(\columnwidth - 6\tabcolsep) * \real{0.40}}@{}}
\toprule
名称 & Python 类型 & 默认值 & 含义 \\
\midrule
\endhead
\passthrough{\lstinline!flag!} & \passthrough{\lstinline!bool!} & 必填 &
布尔开关。\passthrough{\lstinline!True!} 和
\passthrough{\lstinline!False!} 的具体业务效果由当前方法定义。 对应 C++
参数 \passthrough{\lstinline!flag: bool!}。 只接受布尔语义。 \\
\bottomrule
\end{longtable}

\textbf{返回值}

\begin{longtable}[]{@{}
  >{\raggedright\arraybackslash}p{(\columnwidth - 4\tabcolsep) * \real{0.18}}
  >{\raggedright\arraybackslash}p{(\columnwidth - 4\tabcolsep) * \real{0.20}}
  >{\raggedright\arraybackslash}p{(\columnwidth - 4\tabcolsep) * \real{0.62}}@{}}
\toprule
位置 & 类型 & 含义 \\
\midrule
\endhead
返回值 & \passthrough{\lstinline!int!} & SDK 状态码；Unitree 示例通常以
\passthrough{\lstinline!0!} 表示成功，非零值需按具体服务定义解释。 \\
\bottomrule
\end{longtable}

\textbf{对应 C++}

\begin{itemize}
\tightlist
\item
  类：\passthrough{\lstinline!unitree::robot::h2::LocoClient!}
\item
  签名：\passthrough{\lstinline!ContinuousGait(bool)!}
\item
  绑定策略：\passthrough{\lstinline!DIRECT!}
\item
  声明位置：\passthrough{\lstinline!include/unitree/robot/h2/loco/h2\_loco\_client.hpp:268!}
\end{itemize}

\textbf{说明}

\begin{itemize}
\tightlist
\item
  示例是局部调用片段；\passthrough{\lstinline!obj!}
  和参数变量表示已经按上层业务校验并准备好的对象。
\item
  该条目可以被编辑器和类型检查器识别，但当前运行时可能在属性查找阶段直接失败。
\item
  该方法属于运动命令。方法名包含
  \passthrough{\lstinline!stop!}、\passthrough{\lstinline!damp!} 或
  \passthrough{\lstinline!zero!} 也不代表它是独立物理急停。
\end{itemize}

\textbf{用法}

\begin{lstlisting}[language=Python]
from typing import TYPE_CHECKING

if TYPE_CHECKING:
    def planned_continuous_gait(obj: LocoClient, flag: bool) -> int:
        return obj.continuous_gait(flag=flag)
\end{lstlisting}

\begin{quote}
{[}!CAUTION{]} 上面的 \passthrough{\lstinline!TYPE\_CHECKING!}
示例只表达计划调用形态。该分支在普通 Python
运行时为假，不代表当前扩展能执行此接口。
\end{quote}

\hypertarget{unitree-sdk2-cpp-robot-h2-lococlient-switch-move-mode-1}{%
\paragraph{\texorpdfstring{\texttt{unitree\_sdk2\_cpp.robot.h2.LocoClient.switch\_move\_mode}}{unitree\_sdk2\_cpp.robot.h2.LocoClient.switch\_move\_mode}}\label{unitree-sdk2-cpp-robot-h2-lococlient-switch-move-mode-1}}

对应 C++ SDK 操作
\passthrough{\lstinline!SwitchMoveMode(bool)!}。上游头文件没有可直接生成的业务说明时，本参考只保证签名映射，精确语义需查目标型号协议。

\textbf{签名}

\begin{lstlisting}[language=Python]
def switch_move_mode(self, flag: bool) -> int
\end{lstlisting}

\textbf{可用性与安全性}

\textbf{\passthrough{\lstinline!SIGNATURE\_ONLY!} /
\passthrough{\lstinline!MOTION\_COMMAND!}}：仅用于补全和静态检查。当前二进制没有该
Python 方法；不得按可执行运动接口使用。

\textbf{参数}

\begin{longtable}[]{@{}
  >{\raggedright\arraybackslash}p{(\columnwidth - 6\tabcolsep) * \real{0.18}}
  >{\raggedright\arraybackslash}p{(\columnwidth - 6\tabcolsep) * \real{0.24}}
  >{\raggedright\arraybackslash}p{(\columnwidth - 6\tabcolsep) * \real{0.18}}
  >{\raggedright\arraybackslash}p{(\columnwidth - 6\tabcolsep) * \real{0.40}}@{}}
\toprule
名称 & Python 类型 & 默认值 & 含义 \\
\midrule
\endhead
\passthrough{\lstinline!flag!} & \passthrough{\lstinline!bool!} & 必填 &
布尔开关。\passthrough{\lstinline!True!} 和
\passthrough{\lstinline!False!} 的具体业务效果由当前方法定义。 对应 C++
参数 \passthrough{\lstinline!flag: bool!}。 只接受布尔语义。 \\
\bottomrule
\end{longtable}

\textbf{返回值}

\begin{longtable}[]{@{}
  >{\raggedright\arraybackslash}p{(\columnwidth - 4\tabcolsep) * \real{0.18}}
  >{\raggedright\arraybackslash}p{(\columnwidth - 4\tabcolsep) * \real{0.20}}
  >{\raggedright\arraybackslash}p{(\columnwidth - 4\tabcolsep) * \real{0.62}}@{}}
\toprule
位置 & 类型 & 含义 \\
\midrule
\endhead
返回值 & \passthrough{\lstinline!int!} & SDK 状态码；Unitree 示例通常以
\passthrough{\lstinline!0!} 表示成功，非零值需按具体服务定义解释。 \\
\bottomrule
\end{longtable}

\textbf{对应 C++}

\begin{itemize}
\tightlist
\item
  类：\passthrough{\lstinline!unitree::robot::h2::LocoClient!}
\item
  签名：\passthrough{\lstinline!SwitchMoveMode(bool)!}
\item
  绑定策略：\passthrough{\lstinline!DIRECT!}
\item
  声明位置：\passthrough{\lstinline!include/unitree/robot/h2/loco/h2\_loco\_client.hpp:270!}
\end{itemize}

\textbf{说明}

\begin{itemize}
\tightlist
\item
  示例是局部调用片段；\passthrough{\lstinline!obj!}
  和参数变量表示已经按上层业务校验并准备好的对象。
\item
  该条目可以被编辑器和类型检查器识别，但当前运行时可能在属性查找阶段直接失败。
\item
  该方法属于运动命令。方法名包含
  \passthrough{\lstinline!stop!}、\passthrough{\lstinline!damp!} 或
  \passthrough{\lstinline!zero!} 也不代表它是独立物理急停。
\end{itemize}

\textbf{用法}

\begin{lstlisting}[language=Python]
from typing import TYPE_CHECKING

if TYPE_CHECKING:
    def planned_switch_move_mode(obj: LocoClient, flag: bool) -> int:
        return obj.switch_move_mode(flag=flag)
\end{lstlisting}

\begin{quote}
{[}!CAUTION{]} 上面的 \passthrough{\lstinline!TYPE\_CHECKING!}
示例只表达计划调用形态。该分支在普通 Python
运行时为假，不代表当前扩展能执行此接口。
\end{quote}

\hypertarget{unitree-sdk2-cpp-robot-h2-lococlient-wave-hand-1}{%
\paragraph{\texorpdfstring{\texttt{unitree\_sdk2\_cpp.robot.h2.LocoClient.wave\_hand}}{unitree\_sdk2\_cpp.robot.h2.LocoClient.wave\_hand}}\label{unitree-sdk2-cpp-robot-h2-lococlient-wave-hand-1}}

对应 C++ SDK 操作
\passthrough{\lstinline!WaveHand(bool)!}。上游头文件没有可直接生成的业务说明时，本参考只保证签名映射，精确语义需查目标型号协议。

\textbf{签名}

\begin{lstlisting}[language=Python]
def wave_hand(self, turn_flag: bool = ...) -> int
\end{lstlisting}

\textbf{可用性与安全性}

\textbf{\passthrough{\lstinline!SIGNATURE\_ONLY!} /
\passthrough{\lstinline!MOTION\_COMMAND!}}：仅用于补全和静态检查。当前二进制没有该
Python 方法；不得按可执行运动接口使用。

\textbf{参数}

\begin{longtable}[]{@{}
  >{\raggedright\arraybackslash}p{(\columnwidth - 6\tabcolsep) * \real{0.18}}
  >{\raggedright\arraybackslash}p{(\columnwidth - 6\tabcolsep) * \real{0.24}}
  >{\raggedright\arraybackslash}p{(\columnwidth - 6\tabcolsep) * \real{0.18}}
  >{\raggedright\arraybackslash}p{(\columnwidth - 6\tabcolsep) * \real{0.40}}@{}}
\toprule
名称 & Python 类型 & 默认值 & 含义 \\
\midrule
\endhead
\passthrough{\lstinline!turn\_flag!} & \passthrough{\lstinline!bool!} &
\passthrough{\lstinline!...!} & 传给该接口的
\passthrough{\lstinline!turn!} \passthrough{\lstinline!flag!}
参数，Python 类型为
\passthrough{\lstinline!bool!}。精确含义、单位、范围和枚举值需查目标型号对应头文件与协议。
对应 C++ 参数 \passthrough{\lstinline!turn\_flag: bool!}。
只接受布尔语义。 \\
\bottomrule
\end{longtable}

\textbf{返回值}

\begin{longtable}[]{@{}
  >{\raggedright\arraybackslash}p{(\columnwidth - 4\tabcolsep) * \real{0.18}}
  >{\raggedright\arraybackslash}p{(\columnwidth - 4\tabcolsep) * \real{0.20}}
  >{\raggedright\arraybackslash}p{(\columnwidth - 4\tabcolsep) * \real{0.62}}@{}}
\toprule
位置 & 类型 & 含义 \\
\midrule
\endhead
返回值 & \passthrough{\lstinline!int!} & SDK 状态码；Unitree 示例通常以
\passthrough{\lstinline!0!} 表示成功，非零值需按具体服务定义解释。 \\
\bottomrule
\end{longtable}

\textbf{对应 C++}

\begin{itemize}
\tightlist
\item
  类：\passthrough{\lstinline!unitree::robot::h2::LocoClient!}
\item
  签名：\passthrough{\lstinline!WaveHand(bool)!}
\item
  绑定策略：\passthrough{\lstinline!DIRECT!}
\item
  声明位置：\passthrough{\lstinline!include/unitree/robot/h2/loco/h2\_loco\_client.hpp:275!}
\end{itemize}

\textbf{说明}

\begin{itemize}
\tightlist
\item
  示例是局部调用片段；\passthrough{\lstinline!obj!}
  和参数变量表示已经按上层业务校验并准备好的对象。
\item
  默认值显示为 \passthrough{\lstinline!...!}，表示 C++
  声明存在默认参数，但当前 AST
  清单没有保存其字面量；省略参数可使用上游默认行为。
\item
  该条目可以被编辑器和类型检查器识别，但当前运行时可能在属性查找阶段直接失败。
\item
  该方法属于运动命令。方法名包含
  \passthrough{\lstinline!stop!}、\passthrough{\lstinline!damp!} 或
  \passthrough{\lstinline!zero!} 也不代表它是独立物理急停。
\end{itemize}

\textbf{用法}

\begin{lstlisting}[language=Python]
from typing import TYPE_CHECKING

if TYPE_CHECKING:
    def planned_wave_hand(obj: LocoClient, turn_flag: bool) -> int:
        return obj.wave_hand(turn_flag=turn_flag)
\end{lstlisting}

\begin{quote}
{[}!CAUTION{]} 上面的 \passthrough{\lstinline!TYPE\_CHECKING!}
示例只表达计划调用形态。该分支在普通 Python
运行时为假，不代表当前扩展能执行此接口。
\end{quote}

\hypertarget{unitree-sdk2-cpp-robot-h2-lococlient-shake-hand-1}{%
\paragraph{\texorpdfstring{\texttt{unitree\_sdk2\_cpp.robot.h2.LocoClient.shake\_hand}}{unitree\_sdk2\_cpp.robot.h2.LocoClient.shake\_hand}}\label{unitree-sdk2-cpp-robot-h2-lococlient-shake-hand-1}}

对应 C++ SDK 操作
\passthrough{\lstinline!ShakeHand(int)!}。上游头文件没有可直接生成的业务说明时，本参考只保证签名映射，精确语义需查目标型号协议。

\textbf{签名}

\begin{lstlisting}[language=Python]
def shake_hand(self, stage: int = ...) -> int
\end{lstlisting}

\textbf{可用性与安全性}

\textbf{\passthrough{\lstinline!SIGNATURE\_ONLY!} /
\passthrough{\lstinline!MOTION\_COMMAND!}}：仅用于补全和静态检查。当前二进制没有该
Python 方法；不得按可执行运动接口使用。

\textbf{参数}

\begin{longtable}[]{@{}
  >{\raggedright\arraybackslash}p{(\columnwidth - 6\tabcolsep) * \real{0.18}}
  >{\raggedright\arraybackslash}p{(\columnwidth - 6\tabcolsep) * \real{0.24}}
  >{\raggedright\arraybackslash}p{(\columnwidth - 6\tabcolsep) * \real{0.18}}
  >{\raggedright\arraybackslash}p{(\columnwidth - 6\tabcolsep) * \real{0.40}}@{}}
\toprule
名称 & Python 类型 & 默认值 & 含义 \\
\midrule
\endhead
\passthrough{\lstinline!stage!} & \passthrough{\lstinline!int!} &
\passthrough{\lstinline!...!} & 传给该接口的
\passthrough{\lstinline!stage!} 参数，Python 类型为
\passthrough{\lstinline!int!}。精确含义、单位、范围和枚举值需查目标型号对应头文件与协议。
对应 C++ 参数 \passthrough{\lstinline!stage: int!}。 \\
\bottomrule
\end{longtable}

\textbf{返回值}

\begin{longtable}[]{@{}
  >{\raggedright\arraybackslash}p{(\columnwidth - 4\tabcolsep) * \real{0.18}}
  >{\raggedright\arraybackslash}p{(\columnwidth - 4\tabcolsep) * \real{0.20}}
  >{\raggedright\arraybackslash}p{(\columnwidth - 4\tabcolsep) * \real{0.62}}@{}}
\toprule
位置 & 类型 & 含义 \\
\midrule
\endhead
返回值 & \passthrough{\lstinline!int!} & SDK 状态码；Unitree 示例通常以
\passthrough{\lstinline!0!} 表示成功，非零值需按具体服务定义解释。 \\
\bottomrule
\end{longtable}

\textbf{对应 C++}

\begin{itemize}
\tightlist
\item
  类：\passthrough{\lstinline!unitree::robot::h2::LocoClient!}
\item
  签名：\passthrough{\lstinline!ShakeHand(int)!}
\item
  绑定策略：\passthrough{\lstinline!DIRECT!}
\item
  声明位置：\passthrough{\lstinline!include/unitree/robot/h2/loco/h2\_loco\_client.hpp:277!}
\end{itemize}

\textbf{说明}

\begin{itemize}
\tightlist
\item
  示例是局部调用片段；\passthrough{\lstinline!obj!}
  和参数变量表示已经按上层业务校验并准备好的对象。
\item
  默认值显示为 \passthrough{\lstinline!...!}，表示 C++
  声明存在默认参数，但当前 AST
  清单没有保存其字面量；省略参数可使用上游默认行为。
\item
  该条目可以被编辑器和类型检查器识别，但当前运行时可能在属性查找阶段直接失败。
\item
  该方法属于运动命令。方法名包含
  \passthrough{\lstinline!stop!}、\passthrough{\lstinline!damp!} 或
  \passthrough{\lstinline!zero!} 也不代表它是独立物理急停。
\end{itemize}

\textbf{用法}

\begin{lstlisting}[language=Python]
from typing import TYPE_CHECKING

if TYPE_CHECKING:
    def planned_shake_hand(obj: LocoClient, stage: int) -> int:
        return obj.shake_hand(stage=stage)
\end{lstlisting}

\begin{quote}
{[}!CAUTION{]} 上面的 \passthrough{\lstinline!TYPE\_CHECKING!}
示例只表达计划调用形态。该分支在普通 Python
运行时为假，不代表当前扩展能执行此接口。
\end{quote}

\hypertarget{unitree-sdk2-cpp-robot-h2-lococlient-set-speed-mode-1}{%
\paragraph{\texorpdfstring{\texttt{unitree\_sdk2\_cpp.robot.h2.LocoClient.set\_speed\_mode}}{unitree\_sdk2\_cpp.robot.h2.LocoClient.set\_speed\_mode}}\label{unitree-sdk2-cpp-robot-h2-lococlient-set-speed-mode-1}}

设置 速度 模式。具体副作用和安全边界见下方状态。

\textbf{签名}

\begin{lstlisting}[language=Python]
def set_speed_mode(self, speed_mode: int) -> int
\end{lstlisting}

\textbf{可用性与安全性}

\textbf{\passthrough{\lstinline!SIGNATURE\_ONLY!} /
\passthrough{\lstinline!MOTION\_COMMAND!}}：仅用于补全和静态检查。当前二进制没有该
Python 方法；不得按可执行运动接口使用。

\textbf{参数}

\begin{longtable}[]{@{}
  >{\raggedright\arraybackslash}p{(\columnwidth - 6\tabcolsep) * \real{0.18}}
  >{\raggedright\arraybackslash}p{(\columnwidth - 6\tabcolsep) * \real{0.24}}
  >{\raggedright\arraybackslash}p{(\columnwidth - 6\tabcolsep) * \real{0.18}}
  >{\raggedright\arraybackslash}p{(\columnwidth - 6\tabcolsep) * \real{0.40}}@{}}
\toprule
名称 & Python 类型 & 默认值 & 含义 \\
\midrule
\endhead
\passthrough{\lstinline!speed\_mode!} & \passthrough{\lstinline!int!} &
必填 & 传给该接口的 速度 模式 参数，Python 类型为
\passthrough{\lstinline!int!}。精确含义、单位、范围和枚举值需查目标型号对应头文件与协议。
对应 C++ 参数 \passthrough{\lstinline!speed\_mode: int!}。 \\
\bottomrule
\end{longtable}

\textbf{返回值}

\begin{longtable}[]{@{}
  >{\raggedright\arraybackslash}p{(\columnwidth - 4\tabcolsep) * \real{0.18}}
  >{\raggedright\arraybackslash}p{(\columnwidth - 4\tabcolsep) * \real{0.20}}
  >{\raggedright\arraybackslash}p{(\columnwidth - 4\tabcolsep) * \real{0.62}}@{}}
\toprule
位置 & 类型 & 含义 \\
\midrule
\endhead
返回值 & \passthrough{\lstinline!int!} & SDK 状态码；Unitree 示例通常以
\passthrough{\lstinline!0!} 表示成功，非零值需按具体服务定义解释。 \\
\bottomrule
\end{longtable}

\textbf{对应 C++}

\begin{itemize}
\tightlist
\item
  类：\passthrough{\lstinline!unitree::robot::h2::LocoClient!}
\item
  签名：\passthrough{\lstinline!SetSpeedMode(int)!}
\item
  绑定策略：\passthrough{\lstinline!DIRECT!}
\item
  声明位置：\passthrough{\lstinline!include/unitree/robot/h2/loco/h2\_loco\_client.hpp:293!}
\end{itemize}

\textbf{说明}

\begin{itemize}
\tightlist
\item
  示例是局部调用片段；\passthrough{\lstinline!obj!}
  和参数变量表示已经按上层业务校验并准备好的对象。
\item
  该条目可以被编辑器和类型检查器识别，但当前运行时可能在属性查找阶段直接失败。
\item
  该方法属于运动命令。方法名包含
  \passthrough{\lstinline!stop!}、\passthrough{\lstinline!damp!} 或
  \passthrough{\lstinline!zero!} 也不代表它是独立物理急停。
\end{itemize}

\textbf{用法}

\begin{lstlisting}[language=Python]
from typing import TYPE_CHECKING

if TYPE_CHECKING:
    def planned_set_speed_mode(obj: LocoClient, speed_mode: int) -> int:
        return obj.set_speed_mode(speed_mode=speed_mode)
\end{lstlisting}

\begin{quote}
{[}!CAUTION{]} 上面的 \passthrough{\lstinline!TYPE\_CHECKING!}
示例只表达计划调用形态。该分支在普通 Python
运行时为假，不代表当前扩展能执行此接口。
\end{quote}

\hypertarget{unitree-sdk2-cpp-robot-h2-lococlient-enable-arm-sdk-1}{%
\paragraph{\texorpdfstring{\texttt{unitree\_sdk2\_cpp.robot.h2.LocoClient.enable\_arm\_sdk}}{unitree\_sdk2\_cpp.robot.h2.LocoClient.enable\_arm\_sdk}}\label{unitree-sdk2-cpp-robot-h2-lococlient-enable-arm-sdk-1}}

对应 C++ SDK 操作
\passthrough{\lstinline!EnableArmSDK()!}。上游头文件没有可直接生成的业务说明时，本参考只保证签名映射，精确语义需查目标型号协议。

\textbf{签名}

\begin{lstlisting}[language=Python]
def enable_arm_sdk(self) -> int
\end{lstlisting}

\textbf{可用性与安全性}

\textbf{\passthrough{\lstinline!SIGNATURE\_ONLY!} /
\passthrough{\lstinline!MOTION\_COMMAND!}}：仅用于补全和静态检查。当前二进制没有该
Python 方法；不得按可执行运动接口使用。

\textbf{参数}

无显式参数。实例方法中的 \passthrough{\lstinline!self!} 由 Python
自动传入。

\textbf{返回值}

\begin{longtable}[]{@{}
  >{\raggedright\arraybackslash}p{(\columnwidth - 4\tabcolsep) * \real{0.18}}
  >{\raggedright\arraybackslash}p{(\columnwidth - 4\tabcolsep) * \real{0.20}}
  >{\raggedright\arraybackslash}p{(\columnwidth - 4\tabcolsep) * \real{0.62}}@{}}
\toprule
位置 & 类型 & 含义 \\
\midrule
\endhead
返回值 & \passthrough{\lstinline!int!} & SDK 状态码；Unitree 示例通常以
\passthrough{\lstinline!0!} 表示成功，非零值需按具体服务定义解释。 \\
\bottomrule
\end{longtable}

\textbf{对应 C++}

\begin{itemize}
\tightlist
\item
  类：\passthrough{\lstinline!unitree::robot::h2::LocoClient!}
\item
  签名：\passthrough{\lstinline!EnableArmSDK()!}
\item
  绑定策略：\passthrough{\lstinline!DIRECT!}
\item
  声明位置：\passthrough{\lstinline!include/unitree/robot/h2/loco/h2\_loco\_client.hpp:303!}
\end{itemize}

\textbf{说明}

\begin{itemize}
\tightlist
\item
  示例是局部调用片段；\passthrough{\lstinline!obj!}
  和参数变量表示已经按上层业务校验并准备好的对象。
\item
  该条目可以被编辑器和类型检查器识别，但当前运行时可能在属性查找阶段直接失败。
\item
  该方法属于运动命令。方法名包含
  \passthrough{\lstinline!stop!}、\passthrough{\lstinline!damp!} 或
  \passthrough{\lstinline!zero!} 也不代表它是独立物理急停。
\end{itemize}

\textbf{用法}

\begin{lstlisting}[language=Python]
from typing import TYPE_CHECKING

if TYPE_CHECKING:
    def planned_enable_arm_sdk(obj: LocoClient) -> int:
        return obj.enable_arm_sdk()
\end{lstlisting}

\begin{quote}
{[}!CAUTION{]} 上面的 \passthrough{\lstinline!TYPE\_CHECKING!}
示例只表达计划调用形态。该分支在普通 Python
运行时为假，不代表当前扩展能执行此接口。
\end{quote}

\hypertarget{unitree-sdk2-cpp-robot-h2-lococlient-disable-arm-sdk-1}{%
\paragraph{\texorpdfstring{\texttt{unitree\_sdk2\_cpp.robot.h2.LocoClient.disable\_arm\_sdk}}{unitree\_sdk2\_cpp.robot.h2.LocoClient.disable\_arm\_sdk}}\label{unitree-sdk2-cpp-robot-h2-lococlient-disable-arm-sdk-1}}

对应 C++ SDK 操作
\passthrough{\lstinline!DisableArmSDK()!}。上游头文件没有可直接生成的业务说明时，本参考只保证签名映射，精确语义需查目标型号协议。

\textbf{签名}

\begin{lstlisting}[language=Python]
def disable_arm_sdk(self) -> int
\end{lstlisting}

\textbf{可用性与安全性}

\textbf{\passthrough{\lstinline!SIGNATURE\_ONLY!} /
\passthrough{\lstinline!MOTION\_COMMAND!}}：仅用于补全和静态检查。当前二进制没有该
Python 方法；不得按可执行运动接口使用。

\textbf{参数}

无显式参数。实例方法中的 \passthrough{\lstinline!self!} 由 Python
自动传入。

\textbf{返回值}

\begin{longtable}[]{@{}
  >{\raggedright\arraybackslash}p{(\columnwidth - 4\tabcolsep) * \real{0.18}}
  >{\raggedright\arraybackslash}p{(\columnwidth - 4\tabcolsep) * \real{0.20}}
  >{\raggedright\arraybackslash}p{(\columnwidth - 4\tabcolsep) * \real{0.62}}@{}}
\toprule
位置 & 类型 & 含义 \\
\midrule
\endhead
返回值 & \passthrough{\lstinline!int!} & SDK 状态码；Unitree 示例通常以
\passthrough{\lstinline!0!} 表示成功，非零值需按具体服务定义解释。 \\
\bottomrule
\end{longtable}

\textbf{对应 C++}

\begin{itemize}
\tightlist
\item
  类：\passthrough{\lstinline!unitree::robot::h2::LocoClient!}
\item
  签名：\passthrough{\lstinline!DisableArmSDK()!}
\item
  绑定策略：\passthrough{\lstinline!DIRECT!}
\item
  声明位置：\passthrough{\lstinline!include/unitree/robot/h2/loco/h2\_loco\_client.hpp:305!}
\end{itemize}

\textbf{说明}

\begin{itemize}
\tightlist
\item
  示例是局部调用片段；\passthrough{\lstinline!obj!}
  和参数变量表示已经按上层业务校验并准备好的对象。
\item
  该条目可以被编辑器和类型检查器识别，但当前运行时可能在属性查找阶段直接失败。
\item
  该方法属于运动命令。方法名包含
  \passthrough{\lstinline!stop!}、\passthrough{\lstinline!damp!} 或
  \passthrough{\lstinline!zero!} 也不代表它是独立物理急停。
\end{itemize}

\textbf{用法}

\begin{lstlisting}[language=Python]
from typing import TYPE_CHECKING

if TYPE_CHECKING:
    def planned_disable_arm_sdk(obj: LocoClient) -> int:
        return obj.disable_arm_sdk()
\end{lstlisting}

\begin{quote}
{[}!CAUTION{]} 上面的 \passthrough{\lstinline!TYPE\_CHECKING!}
示例只表达计划调用形态。该分支在普通 Python
运行时为假，不代表当前扩展能执行此接口。
\end{quote}

\begin{center}\rule{0.5\linewidth}{0.5pt}\end{center}

\hypertarget{unitree-sdk2-cpp-robot-r1}{%
\subsection{R1 Robot API}\label{unitree-sdk2-cpp-robot-r1}}

模块：\passthrough{\lstinline!unitree\_sdk2\_cpp.robot.r1!}

\hypertarget{ux7c7bux7d22ux5f15-14}{%
\subsubsection{类索引}\label{ux7c7bux7d22ux5f15-14}}

\begin{longtable}[]{@{}
  >{\raggedright\arraybackslash}p{(\columnwidth - 4\tabcolsep) * \real{0.18}}
  >{\raggedleft\arraybackslash}p{(\columnwidth - 4\tabcolsep) * \real{0.20}}
  >{\raggedleft\arraybackslash}p{(\columnwidth - 4\tabcolsep) * \real{0.62}}@{}}
\toprule
类 & 公开函数签名 & 属性 \\
\midrule
\endhead
\protect\hyperlink{unitree-sdk2-cpp-robot-r1-audioclient}{\passthrough{\lstinline!AudioClient!}}
& 8 & 0 \\
\protect\hyperlink{unitree-sdk2-cpp-robot-r1-jsonizedatavecfloat}{\passthrough{\lstinline!JsonizeDataVecFloat!}}
& 3 & 0 \\
\protect\hyperlink{unitree-sdk2-cpp-robot-r1-jsonizevelocitycommand}{\passthrough{\lstinline!JsonizeVelocityCommand!}}
& 3 & 0 \\
\protect\hyperlink{unitree-sdk2-cpp-robot-r1-ledcontrolparameter}{\passthrough{\lstinline!LedControlParameter!}}
& 3 & 0 \\
\protect\hyperlink{unitree-sdk2-cpp-robot-r1-lococlient}{\passthrough{\lstinline!LocoClient!}}
& 15 & 0 \\
\protect\hyperlink{unitree-sdk2-cpp-robot-r1-playstopparameter}{\passthrough{\lstinline!PlayStopParameter!}}
& 3 & 0 \\
\protect\hyperlink{unitree-sdk2-cpp-robot-r1-playstreamparameter}{\passthrough{\lstinline!PlayStreamParameter!}}
& 3 & 0 \\
\protect\hyperlink{unitree-sdk2-cpp-robot-r1-ttsmakerparameter}{\passthrough{\lstinline!TtsMakerParameter!}}
& 3 & 0 \\
\bottomrule
\end{longtable}

\hypertarget{unitree-sdk2-cpp-robot-r1-audioclient}{%
\subsubsection{\texorpdfstring{\texttt{unitree\_sdk2\_cpp.robot.r1.AudioClient}}{unitree\_sdk2\_cpp.robot.r1.AudioClient}}\label{unitree-sdk2-cpp-robot-r1-audioclient}}

Robot 服务客户端。构造、初始化和具体方法的可用性必须分别检查。

\textbf{导入}

\begin{lstlisting}[language=Python]
from unitree_sdk2_cpp.robot.r1 import AudioClient
\end{lstlisting}

\textbf{基类}：\passthrough{\lstinline!Client!}

\textbf{方法索引}

\begin{longtable}[]{@{}
  >{\raggedright\arraybackslash}p{(\columnwidth - 6\tabcolsep) * \real{0.18}}
  >{\raggedright\arraybackslash}p{(\columnwidth - 6\tabcolsep) * \real{0.24}}
  >{\raggedright\arraybackslash}p{(\columnwidth - 6\tabcolsep) * \real{0.18}}
  >{\raggedright\arraybackslash}p{(\columnwidth - 6\tabcolsep) * \real{0.40}}@{}}
\toprule
方法 & Python 签名 & 状态 & 安全分类 \\
\midrule
\endhead
\protect\hyperlink{unitree-sdk2-cpp-robot-r1-audioclient-dunder-init-1}{\passthrough{\lstinline!\_\_init\_\_!}}
& \passthrough{\lstinline!def \_\_init\_\_(self) -> None!} &
\passthrough{\lstinline!AVAILABLE!} &
\passthrough{\lstinline!CONSTRUCTION!} \\
\protect\hyperlink{unitree-sdk2-cpp-robot-r1-audioclient-init-1}{\passthrough{\lstinline!init!}}
& \passthrough{\lstinline!def init(self) -> None!} &
\passthrough{\lstinline!AVAILABLE!} &
\passthrough{\lstinline!INITIALIZATION!} \\
\protect\hyperlink{unitree-sdk2-cpp-robot-r1-audioclient-tts-maker-1}{\passthrough{\lstinline!tts\_maker!}}
&
\passthrough{\lstinline!def tts\_maker(self, text: str, speaker\_id: int) -> int!}
& \passthrough{\lstinline!SIGNATURE\_ONLY!} &
\passthrough{\lstinline!HARDWARE\_SIDE\_EFFECT!} \\
\protect\hyperlink{unitree-sdk2-cpp-robot-r1-audioclient-get-volume-1}{\passthrough{\lstinline!get\_volume!}}
& \passthrough{\lstinline!def get\_volume(self) -> tuple[int, int]!} &
\passthrough{\lstinline!AVAILABLE!} &
\passthrough{\lstinline!READ\_ONLY!} \\
\protect\hyperlink{unitree-sdk2-cpp-robot-r1-audioclient-set-volume-1}{\passthrough{\lstinline!set\_volume!}}
& \passthrough{\lstinline!def set\_volume(self, volume: int) -> int!} &
\passthrough{\lstinline!SIGNATURE\_ONLY!} &
\passthrough{\lstinline!HARDWARE\_SIDE\_EFFECT!} \\
\protect\hyperlink{unitree-sdk2-cpp-robot-r1-audioclient-play-stream-1}{\passthrough{\lstinline!play\_stream!}}
&
\passthrough{\lstinline!def play\_stream(self, app\_name: str, stream\_id: str, pcm\_data: Sequence[int]) -> int!}
& \passthrough{\lstinline!SIGNATURE\_ONLY!} &
\passthrough{\lstinline!HARDWARE\_SIDE\_EFFECT!} \\
\protect\hyperlink{unitree-sdk2-cpp-robot-r1-audioclient-play-stop-1}{\passthrough{\lstinline!play\_stop!}}
& \passthrough{\lstinline!def play\_stop(self, app\_name: str) -> int!}
& \passthrough{\lstinline!SIGNATURE\_ONLY!} &
\passthrough{\lstinline!HARDWARE\_SIDE\_EFFECT!} \\
\protect\hyperlink{unitree-sdk2-cpp-robot-r1-audioclient-led-control-1}{\passthrough{\lstinline!led\_control!}}
&
\passthrough{\lstinline!def led\_control(self, r: int, g: int, b: int) -> int!}
& \passthrough{\lstinline!SIGNATURE\_ONLY!} &
\passthrough{\lstinline!HARDWARE\_SIDE\_EFFECT!} \\
\bottomrule
\end{longtable}

\hypertarget{unitree-sdk2-cpp-robot-r1-audioclient-dunder-init-1}{%
\paragraph{\texorpdfstring{\texttt{unitree\_sdk2\_cpp.robot.r1.AudioClient.\_\_init\_\_}}{unitree\_sdk2\_cpp.robot.r1.AudioClient.\_\_init\_\_}}\label{unitree-sdk2-cpp-robot-r1-audioclient-dunder-init-1}}

初始化 \passthrough{\lstinline!AudioClient!}
实例。是否能实际构造取决于下方可用性状态。

\textbf{签名}

\begin{lstlisting}[language=Python]
def __init__(self) -> None
\end{lstlisting}

\textbf{可用性与安全性}

\textbf{\passthrough{\lstinline!AVAILABLE!} /
\passthrough{\lstinline!CONSTRUCTION!}}：当前绑定源码已实现；实际执行条件仍取决于平台和对应
SDK 环境。

\textbf{参数}

无显式参数。实例方法中的 \passthrough{\lstinline!self!} 由 Python
自动传入。

\textbf{返回值}

\begin{longtable}[]{@{}
  >{\raggedright\arraybackslash}p{(\columnwidth - 4\tabcolsep) * \real{0.18}}
  >{\raggedright\arraybackslash}p{(\columnwidth - 4\tabcolsep) * \real{0.20}}
  >{\raggedright\arraybackslash}p{(\columnwidth - 4\tabcolsep) * \real{0.62}}@{}}
\toprule
位置 & 类型 & 含义 \\
\midrule
\endhead
无 & \passthrough{\lstinline!None!} &
\passthrough{\lstinline!\_\_init\_\_!}
本身不返回值；调用类对象时，在构造函数可用的前提下得到该类实例。 \\
\bottomrule
\end{longtable}

\textbf{对应 C++}

\begin{itemize}
\tightlist
\item
  类：\passthrough{\lstinline!unitree::robot::r1::AudioClient!}
\item
  签名：\passthrough{\lstinline!AudioClient()!}
\item
  绑定策略：\passthrough{\lstinline!未单独标注!}
\item
  声明位置：\passthrough{\lstinline!include/unitree/robot/r1/audio/audio\_client.hpp:15!}
\end{itemize}

\textbf{说明}

\begin{itemize}
\tightlist
\item
  示例是局部调用片段；\passthrough{\lstinline!obj!}
  和参数变量表示已经按上层业务校验并准备好的对象。
\end{itemize}

\textbf{用法}

\begin{lstlisting}[language=Python]
value = AudioClient()
\end{lstlisting}

\hypertarget{unitree-sdk2-cpp-robot-r1-audioclient-init-1}{%
\paragraph{\texorpdfstring{\texttt{unitree\_sdk2\_cpp.robot.r1.AudioClient.init}}{unitree\_sdk2\_cpp.robot.r1.AudioClient.init}}\label{unitree-sdk2-cpp-robot-r1-audioclient-init-1}}

初始化当前 SDK 对象所需的底层通道或服务资源。

\textbf{签名}

\begin{lstlisting}[language=Python]
def init(self) -> None
\end{lstlisting}

\textbf{可用性与安全性}

\textbf{\passthrough{\lstinline!AVAILABLE!} /
\passthrough{\lstinline!INITIALIZATION!}}：当前绑定源码已实现；实际执行条件仍取决于平台和对应
SDK 环境。

\textbf{参数}

无显式参数。实例方法中的 \passthrough{\lstinline!self!} 由 Python
自动传入。

\textbf{返回值}

\begin{longtable}[]{@{}
  >{\raggedright\arraybackslash}p{(\columnwidth - 4\tabcolsep) * \real{0.18}}
  >{\raggedright\arraybackslash}p{(\columnwidth - 4\tabcolsep) * \real{0.20}}
  >{\raggedright\arraybackslash}p{(\columnwidth - 4\tabcolsep) * \real{0.62}}@{}}
\toprule
位置 & 类型 & 含义 \\
\midrule
\endhead
无 & \passthrough{\lstinline!None!} & 该方法不返回 Python
值；失败可能表现为异常或底层状态变化。 \\
\bottomrule
\end{longtable}

\textbf{对应 C++}

\begin{itemize}
\tightlist
\item
  类：\passthrough{\lstinline!unitree::robot::r1::AudioClient!}
\item
  签名：\passthrough{\lstinline!Init()!}
\item
  绑定策略：\passthrough{\lstinline!DIRECT!}
\item
  声明位置：\passthrough{\lstinline!include/unitree/robot/r1/audio/audio\_client.hpp:19!}
\end{itemize}

\textbf{说明}

\begin{itemize}
\tightlist
\item
  示例是局部调用片段；\passthrough{\lstinline!obj!}
  和参数变量表示已经按上层业务校验并准备好的对象。
\end{itemize}

\textbf{用法}

\begin{lstlisting}[language=Python]
obj.init()
\end{lstlisting}

\hypertarget{unitree-sdk2-cpp-robot-r1-audioclient-tts-maker-1}{%
\paragraph{\texorpdfstring{\texttt{unitree\_sdk2\_cpp.robot.r1.AudioClient.tts\_maker}}{unitree\_sdk2\_cpp.robot.r1.AudioClient.tts\_maker}}\label{unitree-sdk2-cpp-robot-r1-audioclient-tts-maker-1}}

对应 C++ SDK 操作
\passthrough{\lstinline!TtsMaker(const std::string \&, int32\_t)!}。上游头文件没有可直接生成的业务说明时，本参考只保证签名映射，精确语义需查目标型号协议。

\textbf{签名}

\begin{lstlisting}[language=Python]
def tts_maker(self, text: str, speaker_id: int) -> int
\end{lstlisting}

\textbf{可用性与安全性}

\textbf{\passthrough{\lstinline!SIGNATURE\_ONLY!} /
\passthrough{\lstinline!HARDWARE\_SIDE\_EFFECT!}}：仅用于补全和静态检查。当前二进制没有该
Python 方法，未来实现还需单独评审硬件副作用。

\textbf{参数}

\begin{longtable}[]{@{}
  >{\raggedright\arraybackslash}p{(\columnwidth - 6\tabcolsep) * \real{0.18}}
  >{\raggedright\arraybackslash}p{(\columnwidth - 6\tabcolsep) * \real{0.24}}
  >{\raggedright\arraybackslash}p{(\columnwidth - 6\tabcolsep) * \real{0.18}}
  >{\raggedright\arraybackslash}p{(\columnwidth - 6\tabcolsep) * \real{0.40}}@{}}
\toprule
名称 & Python 类型 & 默认值 & 含义 \\
\midrule
\endhead
\passthrough{\lstinline!text!} & \passthrough{\lstinline!str!} & 必填 &
要处理的文本内容；编码、长度和语言支持由目标服务决定。 对应 C++ 参数
\passthrough{\lstinline!text: const std::string \&!}。
底层为字符串；长度、编码和允许值由具体协议决定。 \\
\passthrough{\lstinline!speaker\_id!} & \passthrough{\lstinline!int!} &
必填 & 语音合成说话人 ID。有效编号由机器人音频服务和固件决定。 对应 C++
参数 \passthrough{\lstinline!speaker\_id: int32\_t!}。 取值范围为
-2147483648 到 2147483647。 \\
\bottomrule
\end{longtable}

\textbf{返回值}

\begin{longtable}[]{@{}
  >{\raggedright\arraybackslash}p{(\columnwidth - 4\tabcolsep) * \real{0.18}}
  >{\raggedright\arraybackslash}p{(\columnwidth - 4\tabcolsep) * \real{0.20}}
  >{\raggedright\arraybackslash}p{(\columnwidth - 4\tabcolsep) * \real{0.62}}@{}}
\toprule
位置 & 类型 & 含义 \\
\midrule
\endhead
返回值 & \passthrough{\lstinline!int!} & SDK 状态码；Unitree 示例通常以
\passthrough{\lstinline!0!} 表示成功，非零值需按具体服务定义解释。 \\
\bottomrule
\end{longtable}

\textbf{对应 C++}

\begin{itemize}
\tightlist
\item
  类：\passthrough{\lstinline!unitree::robot::r1::AudioClient!}
\item
  签名：\passthrough{\lstinline!TtsMaker(const std::string \&, int32\_t)!}
\item
  绑定策略：\passthrough{\lstinline!DIRECT!}
\item
  声明位置：\passthrough{\lstinline!include/unitree/robot/r1/audio/audio\_client.hpp:31!}
\end{itemize}

\textbf{说明}

\begin{itemize}
\tightlist
\item
  示例是局部调用片段；\passthrough{\lstinline!obj!}
  和参数变量表示已经按上层业务校验并准备好的对象。
\item
  该条目可以被编辑器和类型检查器识别，但当前运行时可能在属性查找阶段直接失败。
\end{itemize}

\textbf{用法}

\begin{lstlisting}[language=Python]
from typing import TYPE_CHECKING

if TYPE_CHECKING:
    def planned_tts_maker(obj: AudioClient, text: str, speaker_id: int) -> int:
        return obj.tts_maker(text=text, speaker_id=speaker_id)
\end{lstlisting}

\begin{quote}
{[}!CAUTION{]} 上面的 \passthrough{\lstinline!TYPE\_CHECKING!}
示例只表达计划调用形态。该分支在普通 Python
运行时为假，不代表当前扩展能执行此接口。
\end{quote}

\hypertarget{unitree-sdk2-cpp-robot-r1-audioclient-get-volume-1}{%
\paragraph{\texorpdfstring{\texttt{unitree\_sdk2\_cpp.robot.r1.AudioClient.get\_volume}}{unitree\_sdk2\_cpp.robot.r1.AudioClient.get\_volume}}\label{unitree-sdk2-cpp-robot-r1-audioclient-get-volume-1}}

查询或检查 音量。

\textbf{签名}

\begin{lstlisting}[language=Python]
def get_volume(self) -> tuple[int, int]
\end{lstlisting}

\textbf{可用性与安全性}

\textbf{\passthrough{\lstinline!AVAILABLE!} /
\passthrough{\lstinline!READ\_ONLY!}}：当前绑定源码已实现只读查询。Robot
Client 的服务端查询通常仍需匹配的 Linux
扩展、DDS、网络、机器人和服务；纯本地 getter 则不一定需要全部条件。

\textbf{参数}

无显式参数。实例方法中的 \passthrough{\lstinline!self!} 由 Python
自动传入。

\textbf{返回值}

\begin{longtable}[]{@{}
  >{\raggedright\arraybackslash}p{(\columnwidth - 4\tabcolsep) * \real{0.18}}
  >{\raggedright\arraybackslash}p{(\columnwidth - 4\tabcolsep) * \real{0.20}}
  >{\raggedright\arraybackslash}p{(\columnwidth - 4\tabcolsep) * \real{0.62}}@{}}
\toprule
位置 & 类型 & 含义 \\
\midrule
\endhead
{[}0{]} \passthrough{\lstinline!status!} & \passthrough{\lstinline!int!}
& SDK 状态码；Unitree 示例通常以 \passthrough{\lstinline!0!}
表示成功，非零值需按具体服务错误码解释。 \\
{[}1{]} \passthrough{\lstinline!volume!} & \passthrough{\lstinline!int!}
& 传给该接口的 音量 参数，Python 类型为
\passthrough{\lstinline!int!}。精确含义、单位、范围和枚举值需查目标型号对应头文件与协议。
对应 C++ 参数 \passthrough{\lstinline!volume: uint8\_t \&!}。 取值范围为
0 到 255。 \\
\bottomrule
\end{longtable}

\textbf{对应 C++}

\begin{itemize}
\tightlist
\item
  类：\passthrough{\lstinline!unitree::robot::r1::AudioClient!}
\item
  签名：\passthrough{\lstinline!GetVolume(uint8\_t \&)!}
\item
  绑定策略：\passthrough{\lstinline!OUTPUT\_WRAPPER!}
\item
  声明位置：\passthrough{\lstinline!include/unitree/robot/r1/audio/audio\_client.hpp:43!}
\end{itemize}

\textbf{说明}

\begin{itemize}
\tightlist
\item
  示例是局部调用片段；\passthrough{\lstinline!obj!}
  和参数变量表示已经按上层业务校验并准备好的对象。
\item
  C++ 的可修改输出引用不会作为 Python 入参出现，而是按 C++
  参数顺序追加到返回元组中。
\end{itemize}

\textbf{用法}

\begin{lstlisting}[language=Python]
status, volume = obj.get_volume()
\end{lstlisting}

\begin{quote}
{[}!NOTE{]} 只读查询仍会访问
DDS/机器人服务。必须检查返回状态码，并处理超时和网络中断。
\end{quote}

\hypertarget{unitree-sdk2-cpp-robot-r1-audioclient-set-volume-1}{%
\paragraph{\texorpdfstring{\texttt{unitree\_sdk2\_cpp.robot.r1.AudioClient.set\_volume}}{unitree\_sdk2\_cpp.robot.r1.AudioClient.set\_volume}}\label{unitree-sdk2-cpp-robot-r1-audioclient-set-volume-1}}

设置 音量。具体副作用和安全边界见下方状态。

\textbf{签名}

\begin{lstlisting}[language=Python]
def set_volume(self, volume: int) -> int
\end{lstlisting}

\textbf{可用性与安全性}

\textbf{\passthrough{\lstinline!SIGNATURE\_ONLY!} /
\passthrough{\lstinline!HARDWARE\_SIDE\_EFFECT!}}：仅用于补全和静态检查。当前二进制没有该
Python 方法，未来实现还需单独评审硬件副作用。

\textbf{参数}

\begin{longtable}[]{@{}
  >{\raggedright\arraybackslash}p{(\columnwidth - 6\tabcolsep) * \real{0.18}}
  >{\raggedright\arraybackslash}p{(\columnwidth - 6\tabcolsep) * \real{0.24}}
  >{\raggedright\arraybackslash}p{(\columnwidth - 6\tabcolsep) * \real{0.18}}
  >{\raggedright\arraybackslash}p{(\columnwidth - 6\tabcolsep) * \real{0.40}}@{}}
\toprule
名称 & Python 类型 & 默认值 & 含义 \\
\midrule
\endhead
\passthrough{\lstinline!volume!} & \passthrough{\lstinline!int!} & 必填
& 传给该接口的 音量 参数，Python 类型为
\passthrough{\lstinline!int!}。精确含义、单位、范围和枚举值需查目标型号对应头文件与协议。
对应 C++ 参数 \passthrough{\lstinline!volume: uint8\_t!}。 取值范围为 0
到 255。 \\
\bottomrule
\end{longtable}

\textbf{返回值}

\begin{longtable}[]{@{}
  >{\raggedright\arraybackslash}p{(\columnwidth - 4\tabcolsep) * \real{0.18}}
  >{\raggedright\arraybackslash}p{(\columnwidth - 4\tabcolsep) * \real{0.20}}
  >{\raggedright\arraybackslash}p{(\columnwidth - 4\tabcolsep) * \real{0.62}}@{}}
\toprule
位置 & 类型 & 含义 \\
\midrule
\endhead
返回值 & \passthrough{\lstinline!int!} & SDK 状态码；Unitree 示例通常以
\passthrough{\lstinline!0!} 表示成功，非零值需按具体服务定义解释。 \\
\bottomrule
\end{longtable}

\textbf{对应 C++}

\begin{itemize}
\tightlist
\item
  类：\passthrough{\lstinline!unitree::robot::r1::AudioClient!}
\item
  签名：\passthrough{\lstinline!SetVolume(uint8\_t)!}
\item
  绑定策略：\passthrough{\lstinline!DIRECT!}
\item
  声明位置：\passthrough{\lstinline!include/unitree/robot/r1/audio/audio\_client.hpp:57!}
\end{itemize}

\textbf{说明}

\begin{itemize}
\tightlist
\item
  示例是局部调用片段；\passthrough{\lstinline!obj!}
  和参数变量表示已经按上层业务校验并准备好的对象。
\item
  该条目可以被编辑器和类型检查器识别，但当前运行时可能在属性查找阶段直接失败。
\end{itemize}

\textbf{用法}

\begin{lstlisting}[language=Python]
from typing import TYPE_CHECKING

if TYPE_CHECKING:
    def planned_set_volume(obj: AudioClient, volume: int) -> int:
        return obj.set_volume(volume=volume)
\end{lstlisting}

\begin{quote}
{[}!CAUTION{]} 上面的 \passthrough{\lstinline!TYPE\_CHECKING!}
示例只表达计划调用形态。该分支在普通 Python
运行时为假，不代表当前扩展能执行此接口。
\end{quote}

\hypertarget{unitree-sdk2-cpp-robot-r1-audioclient-play-stream-1}{%
\paragraph{\texorpdfstring{\texttt{unitree\_sdk2\_cpp.robot.r1.AudioClient.play\_stream}}{unitree\_sdk2\_cpp.robot.r1.AudioClient.play\_stream}}\label{unitree-sdk2-cpp-robot-r1-audioclient-play-stream-1}}

对应 C++ SDK 操作
\passthrough{\lstinline!PlayStream(std::string, std::string, std::vector<uint8\_t>)!}。上游头文件没有可直接生成的业务说明时，本参考只保证签名映射，精确语义需查目标型号协议。

\textbf{签名}

\begin{lstlisting}[language=Python]
def play_stream(self, app_name: str, stream_id: str, pcm_data: Sequence[int]) -> int
\end{lstlisting}

\textbf{可用性与安全性}

\textbf{\passthrough{\lstinline!SIGNATURE\_ONLY!} /
\passthrough{\lstinline!HARDWARE\_SIDE\_EFFECT!}}：仅用于补全和静态检查。当前二进制没有该
Python 方法，未来实现还需单独评审硬件副作用。

\textbf{参数}

\begin{longtable}[]{@{}
  >{\raggedright\arraybackslash}p{(\columnwidth - 6\tabcolsep) * \real{0.18}}
  >{\raggedright\arraybackslash}p{(\columnwidth - 6\tabcolsep) * \real{0.24}}
  >{\raggedright\arraybackslash}p{(\columnwidth - 6\tabcolsep) * \real{0.18}}
  >{\raggedright\arraybackslash}p{(\columnwidth - 6\tabcolsep) * \real{0.40}}@{}}
\toprule
名称 & Python 类型 & 默认值 & 含义 \\
\midrule
\endhead
\passthrough{\lstinline!app\_name!} & \passthrough{\lstinline!str!} &
必填 &
应用名称，用于标识音频流或播放会话。允许值和命名规则以目标服务协议为准。
对应 C++ 参数 \passthrough{\lstinline!app\_name: std::string!}。
底层为字符串；长度、编码和允许值由具体协议决定。 \\
\passthrough{\lstinline!stream\_id!} & \passthrough{\lstinline!str!} &
必填 & 音频流标识符，用于区分同一应用下的播放流。 对应 C++ 参数
\passthrough{\lstinline!stream\_id: std::string!}。
底层为字符串；长度、编码和允许值由具体协议决定。 \\
\passthrough{\lstinline!pcm\_data!} &
\passthrough{\lstinline!Sequence[int]!} & 必填 & 传给该接口的
\passthrough{\lstinline!pcm!} 数据 参数，Python 类型为
\passthrough{\lstinline!Sequence[int]!}。精确含义、单位、范围和枚举值需查目标型号对应头文件与协议。
对应 C++ 参数
\passthrough{\lstinline!pcm\_data: std::vector<uint8\_t>!}。
底层是可变长度 vector；元素约束：取值范围为 0 到 255。 \\
\bottomrule
\end{longtable}

\textbf{返回值}

\begin{longtable}[]{@{}
  >{\raggedright\arraybackslash}p{(\columnwidth - 4\tabcolsep) * \real{0.18}}
  >{\raggedright\arraybackslash}p{(\columnwidth - 4\tabcolsep) * \real{0.20}}
  >{\raggedright\arraybackslash}p{(\columnwidth - 4\tabcolsep) * \real{0.62}}@{}}
\toprule
位置 & 类型 & 含义 \\
\midrule
\endhead
返回值 & \passthrough{\lstinline!int!} & SDK 状态码；Unitree 示例通常以
\passthrough{\lstinline!0!} 表示成功，非零值需按具体服务定义解释。 \\
\bottomrule
\end{longtable}

\textbf{对应 C++}

\begin{itemize}
\tightlist
\item
  类：\passthrough{\lstinline!unitree::robot::r1::AudioClient!}
\item
  签名：\passthrough{\lstinline!PlayStream(std::string, std::string, std::vector<uint8\_t>)!}
\item
  绑定策略：\passthrough{\lstinline!DIRECT!}
\item
  声明位置：\passthrough{\lstinline!include/unitree/robot/r1/audio/audio\_client.hpp:68!}
\end{itemize}

\textbf{说明}

\begin{itemize}
\tightlist
\item
  示例是局部调用片段；\passthrough{\lstinline!obj!}
  和参数变量表示已经按上层业务校验并准备好的对象。
\item
  该条目可以被编辑器和类型检查器识别，但当前运行时可能在属性查找阶段直接失败。
\end{itemize}

\textbf{用法}

\begin{lstlisting}[language=Python]
from typing import TYPE_CHECKING

if TYPE_CHECKING:
    def planned_play_stream(obj: AudioClient, app_name: str, stream_id: str, pcm_data: Sequence[int]) -> int:
        return obj.play_stream(app_name=app_name, stream_id=stream_id, pcm_data=pcm_data)
\end{lstlisting}

\begin{quote}
{[}!CAUTION{]} 上面的 \passthrough{\lstinline!TYPE\_CHECKING!}
示例只表达计划调用形态。该分支在普通 Python
运行时为假，不代表当前扩展能执行此接口。
\end{quote}

\hypertarget{unitree-sdk2-cpp-robot-r1-audioclient-play-stop-1}{%
\paragraph{\texorpdfstring{\texttt{unitree\_sdk2\_cpp.robot.r1.AudioClient.play\_stop}}{unitree\_sdk2\_cpp.robot.r1.AudioClient.play\_stop}}\label{unitree-sdk2-cpp-robot-r1-audioclient-play-stop-1}}

对应 C++ SDK 操作
\passthrough{\lstinline!PlayStop(std::string)!}。上游头文件没有可直接生成的业务说明时，本参考只保证签名映射，精确语义需查目标型号协议。

\textbf{签名}

\begin{lstlisting}[language=Python]
def play_stop(self, app_name: str) -> int
\end{lstlisting}

\textbf{可用性与安全性}

\textbf{\passthrough{\lstinline!SIGNATURE\_ONLY!} /
\passthrough{\lstinline!HARDWARE\_SIDE\_EFFECT!}}：仅用于补全和静态检查。当前二进制没有该
Python 方法，未来实现还需单独评审硬件副作用。

\textbf{参数}

\begin{longtable}[]{@{}
  >{\raggedright\arraybackslash}p{(\columnwidth - 6\tabcolsep) * \real{0.18}}
  >{\raggedright\arraybackslash}p{(\columnwidth - 6\tabcolsep) * \real{0.24}}
  >{\raggedright\arraybackslash}p{(\columnwidth - 6\tabcolsep) * \real{0.18}}
  >{\raggedright\arraybackslash}p{(\columnwidth - 6\tabcolsep) * \real{0.40}}@{}}
\toprule
名称 & Python 类型 & 默认值 & 含义 \\
\midrule
\endhead
\passthrough{\lstinline!app\_name!} & \passthrough{\lstinline!str!} &
必填 &
应用名称，用于标识音频流或播放会话。允许值和命名规则以目标服务协议为准。
对应 C++ 参数 \passthrough{\lstinline!app\_name: std::string!}。
底层为字符串；长度、编码和允许值由具体协议决定。 \\
\bottomrule
\end{longtable}

\textbf{返回值}

\begin{longtable}[]{@{}
  >{\raggedright\arraybackslash}p{(\columnwidth - 4\tabcolsep) * \real{0.18}}
  >{\raggedright\arraybackslash}p{(\columnwidth - 4\tabcolsep) * \real{0.20}}
  >{\raggedright\arraybackslash}p{(\columnwidth - 4\tabcolsep) * \real{0.62}}@{}}
\toprule
位置 & 类型 & 含义 \\
\midrule
\endhead
返回值 & \passthrough{\lstinline!int!} & SDK 状态码；Unitree 示例通常以
\passthrough{\lstinline!0!} 表示成功，非零值需按具体服务定义解释。 \\
\bottomrule
\end{longtable}

\textbf{对应 C++}

\begin{itemize}
\tightlist
\item
  类：\passthrough{\lstinline!unitree::robot::r1::AudioClient!}
\item
  签名：\passthrough{\lstinline!PlayStop(std::string)!}
\item
  绑定策略：\passthrough{\lstinline!DIRECT!}
\item
  声明位置：\passthrough{\lstinline!include/unitree/robot/r1/audio/audio\_client.hpp:80!}
\end{itemize}

\textbf{说明}

\begin{itemize}
\tightlist
\item
  示例是局部调用片段；\passthrough{\lstinline!obj!}
  和参数变量表示已经按上层业务校验并准备好的对象。
\item
  该条目可以被编辑器和类型检查器识别，但当前运行时可能在属性查找阶段直接失败。
\end{itemize}

\textbf{用法}

\begin{lstlisting}[language=Python]
from typing import TYPE_CHECKING

if TYPE_CHECKING:
    def planned_play_stop(obj: AudioClient, app_name: str) -> int:
        return obj.play_stop(app_name=app_name)
\end{lstlisting}

\begin{quote}
{[}!CAUTION{]} 上面的 \passthrough{\lstinline!TYPE\_CHECKING!}
示例只表达计划调用形态。该分支在普通 Python
运行时为假，不代表当前扩展能执行此接口。
\end{quote}

\hypertarget{unitree-sdk2-cpp-robot-r1-audioclient-led-control-1}{%
\paragraph{\texorpdfstring{\texttt{unitree\_sdk2\_cpp.robot.r1.AudioClient.led\_control}}{unitree\_sdk2\_cpp.robot.r1.AudioClient.led\_control}}\label{unitree-sdk2-cpp-robot-r1-audioclient-led-control-1}}

对应 C++ SDK 操作
\passthrough{\lstinline!LedControl(uint8\_t, uint8\_t, uint8\_t)!}。上游头文件没有可直接生成的业务说明时，本参考只保证签名映射，精确语义需查目标型号协议。

\textbf{签名}

\begin{lstlisting}[language=Python]
def led_control(self, r: int, g: int, b: int) -> int
\end{lstlisting}

\textbf{可用性与安全性}

\textbf{\passthrough{\lstinline!SIGNATURE\_ONLY!} /
\passthrough{\lstinline!HARDWARE\_SIDE\_EFFECT!}}：仅用于补全和静态检查。当前二进制没有该
Python 方法，未来实现还需单独评审硬件副作用。

\textbf{参数}

\begin{longtable}[]{@{}
  >{\raggedright\arraybackslash}p{(\columnwidth - 6\tabcolsep) * \real{0.18}}
  >{\raggedright\arraybackslash}p{(\columnwidth - 6\tabcolsep) * \real{0.24}}
  >{\raggedright\arraybackslash}p{(\columnwidth - 6\tabcolsep) * \real{0.18}}
  >{\raggedright\arraybackslash}p{(\columnwidth - 6\tabcolsep) * \real{0.40}}@{}}
\toprule
名称 & Python 类型 & 默认值 & 含义 \\
\midrule
\endhead
\passthrough{\lstinline!r!} & \passthrough{\lstinline!int!} & 必填 &
传给该接口的 \passthrough{\lstinline!r!} 参数，Python 类型为
\passthrough{\lstinline!int!}。精确含义、单位、范围和枚举值需查目标型号对应头文件与协议。
对应 C++ 参数 \passthrough{\lstinline!R: uint8\_t!}。 取值范围为 0 到
255。 \\
\passthrough{\lstinline!g!} & \passthrough{\lstinline!int!} & 必填 &
传给该接口的 \passthrough{\lstinline!g!} 参数，Python 类型为
\passthrough{\lstinline!int!}。精确含义、单位、范围和枚举值需查目标型号对应头文件与协议。
对应 C++ 参数 \passthrough{\lstinline!G: uint8\_t!}。 取值范围为 0 到
255。 \\
\passthrough{\lstinline!b!} & \passthrough{\lstinline!int!} & 必填 &
传给该接口的 \passthrough{\lstinline!b!} 参数，Python 类型为
\passthrough{\lstinline!int!}。精确含义、单位、范围和枚举值需查目标型号对应头文件与协议。
对应 C++ 参数 \passthrough{\lstinline!B: uint8\_t!}。 取值范围为 0 到
255。 \\
\bottomrule
\end{longtable}

\textbf{返回值}

\begin{longtable}[]{@{}
  >{\raggedright\arraybackslash}p{(\columnwidth - 4\tabcolsep) * \real{0.18}}
  >{\raggedright\arraybackslash}p{(\columnwidth - 4\tabcolsep) * \real{0.20}}
  >{\raggedright\arraybackslash}p{(\columnwidth - 4\tabcolsep) * \real{0.62}}@{}}
\toprule
位置 & 类型 & 含义 \\
\midrule
\endhead
返回值 & \passthrough{\lstinline!int!} & SDK 状态码；Unitree 示例通常以
\passthrough{\lstinline!0!} 表示成功，非零值需按具体服务定义解释。 \\
\bottomrule
\end{longtable}

\textbf{对应 C++}

\begin{itemize}
\tightlist
\item
  类：\passthrough{\lstinline!unitree::robot::r1::AudioClient!}
\item
  签名：\passthrough{\lstinline!LedControl(uint8\_t, uint8\_t, uint8\_t)!}
\item
  绑定策略：\passthrough{\lstinline!DIRECT!}
\item
  声明位置：\passthrough{\lstinline!include/unitree/robot/r1/audio/audio\_client.hpp:91!}
\end{itemize}

\textbf{说明}

\begin{itemize}
\tightlist
\item
  示例是局部调用片段；\passthrough{\lstinline!obj!}
  和参数变量表示已经按上层业务校验并准备好的对象。
\item
  该条目可以被编辑器和类型检查器识别，但当前运行时可能在属性查找阶段直接失败。
\end{itemize}

\textbf{用法}

\begin{lstlisting}[language=Python]
from typing import TYPE_CHECKING

if TYPE_CHECKING:
    def planned_led_control(obj: AudioClient, r: int, g: int, b: int) -> int:
        return obj.led_control(r=r, g=g, b=b)
\end{lstlisting}

\begin{quote}
{[}!CAUTION{]} 上面的 \passthrough{\lstinline!TYPE\_CHECKING!}
示例只表达计划调用形态。该分支在普通 Python
运行时为假，不代表当前扩展能执行此接口。
\end{quote}

\hypertarget{unitree-sdk2-cpp-robot-r1-jsonizedatavecfloat}{%
\subsubsection{\texorpdfstring{\texttt{unitree\_sdk2\_cpp.robot.r1.JsonizeDataVecFloat}}{unitree\_sdk2\_cpp.robot.r1.JsonizeDataVecFloat}}\label{unitree-sdk2-cpp-robot-r1-jsonizedatavecfloat}}

SDK 类型签名预览。请逐项查看构造函数和方法的可用性。

\textbf{导入}

\begin{lstlisting}[language=Python]
from unitree_sdk2_cpp.robot.r1 import JsonizeDataVecFloat
\end{lstlisting}

\textbf{基类}：\passthrough{\lstinline!object!}

\textbf{公开属性}

\begin{longtable}[]{@{}
  >{\raggedright\arraybackslash}p{(\columnwidth - 6\tabcolsep) * \real{0.18}}
  >{\raggedright\arraybackslash}p{(\columnwidth - 6\tabcolsep) * \real{0.24}}
  >{\raggedright\arraybackslash}p{(\columnwidth - 6\tabcolsep) * \real{0.18}}
  >{\raggedright\arraybackslash}p{(\columnwidth - 6\tabcolsep) * \real{0.40}}@{}}
\toprule
名称 & Python 类型 & 含义 & 用法 \\
\midrule
\endhead
\passthrough{\lstinline!data!} & \passthrough{\lstinline!list[float]!} &
传给该接口的 数据 参数，Python 类型为
\passthrough{\lstinline!list[float]!}。精确含义、单位、范围和枚举值需查目标型号对应头文件与协议。
& \passthrough{\lstinline!value = obj.data!} /
\passthrough{\lstinline!obj.data = value!} \\
\bottomrule
\end{longtable}

\textbf{方法索引}

\begin{longtable}[]{@{}
  >{\raggedright\arraybackslash}p{(\columnwidth - 6\tabcolsep) * \real{0.18}}
  >{\raggedright\arraybackslash}p{(\columnwidth - 6\tabcolsep) * \real{0.24}}
  >{\raggedright\arraybackslash}p{(\columnwidth - 6\tabcolsep) * \real{0.18}}
  >{\raggedright\arraybackslash}p{(\columnwidth - 6\tabcolsep) * \real{0.40}}@{}}
\toprule
方法 & Python 签名 & 状态 & 安全分类 \\
\midrule
\endhead
\protect\hyperlink{unitree-sdk2-cpp-robot-r1-jsonizedatavecfloat-dunder-init-1}{\passthrough{\lstinline!\_\_init\_\_!}}
& \passthrough{\lstinline!def \_\_init\_\_(self) -> None!} &
\passthrough{\lstinline!SIGNATURE\_ONLY!} &
\passthrough{\lstinline!CONSTRUCTION!} \\
\protect\hyperlink{unitree-sdk2-cpp-robot-r1-jsonizedatavecfloat-from-json-1}{\passthrough{\lstinline!from\_json!}}
&
\passthrough{\lstinline!def from\_json(self, json: dict[str, Any]) -> None!}
& \passthrough{\lstinline!SIGNATURE\_ONLY!} &
\passthrough{\lstinline!UNCLASSIFIED!} \\
\protect\hyperlink{unitree-sdk2-cpp-robot-r1-jsonizedatavecfloat-to-json-1}{\passthrough{\lstinline!to\_json!}}
&
\passthrough{\lstinline!def to\_json(self, json: dict[str, Any]) -> None!}
& \passthrough{\lstinline!SIGNATURE\_ONLY!} &
\passthrough{\lstinline!UNCLASSIFIED!} \\
\bottomrule
\end{longtable}

\hypertarget{unitree-sdk2-cpp-robot-r1-jsonizedatavecfloat-dunder-init-1}{%
\paragraph{\texorpdfstring{\texttt{unitree\_sdk2\_cpp.robot.r1.JsonizeDataVecFloat.\_\_init\_\_}}{unitree\_sdk2\_cpp.robot.r1.JsonizeDataVecFloat.\_\_init\_\_}}\label{unitree-sdk2-cpp-robot-r1-jsonizedatavecfloat-dunder-init-1}}

初始化 \passthrough{\lstinline!JsonizeDataVecFloat!}
实例。是否能实际构造取决于下方可用性状态。

\textbf{签名}

\begin{lstlisting}[language=Python]
def __init__(self) -> None
\end{lstlisting}

\textbf{可用性与安全性}

\textbf{\passthrough{\lstinline!SIGNATURE\_ONLY!} /
\passthrough{\lstinline!CONSTRUCTION!}}：当前只有设计期签名，不能假设已存在于运行时扩展。

\textbf{参数}

无显式参数。实例方法中的 \passthrough{\lstinline!self!} 由 Python
自动传入。

\textbf{返回值}

\begin{longtable}[]{@{}
  >{\raggedright\arraybackslash}p{(\columnwidth - 4\tabcolsep) * \real{0.18}}
  >{\raggedright\arraybackslash}p{(\columnwidth - 4\tabcolsep) * \real{0.20}}
  >{\raggedright\arraybackslash}p{(\columnwidth - 4\tabcolsep) * \real{0.62}}@{}}
\toprule
位置 & 类型 & 含义 \\
\midrule
\endhead
无 & \passthrough{\lstinline!None!} &
\passthrough{\lstinline!\_\_init\_\_!}
本身不返回值；调用类对象时，在构造函数可用的前提下得到该类实例。 \\
\bottomrule
\end{longtable}

\textbf{对应 C++}

\begin{itemize}
\tightlist
\item
  类：\passthrough{\lstinline!unitree::robot::r1::JsonizeDataVecFloat!}
\item
  签名：\passthrough{\lstinline!JsonizeDataVecFloat()!}
\item
  绑定策略：\passthrough{\lstinline!未单独标注!}
\item
  声明位置：\passthrough{\lstinline!include/unitree/robot/r1/loco/r1\_loco\_api.hpp:27!}
\end{itemize}

\textbf{说明}

\begin{itemize}
\tightlist
\item
  示例是局部调用片段；\passthrough{\lstinline!obj!}
  和参数变量表示已经按上层业务校验并准备好的对象。
\item
  该条目可以被编辑器和类型检查器识别，但当前运行时可能在属性查找阶段直接失败。
\end{itemize}

\textbf{用法}

\begin{lstlisting}[language=Python]
from typing import TYPE_CHECKING

if TYPE_CHECKING:
    def planned_construct() -> JsonizeDataVecFloat:
        return JsonizeDataVecFloat()
\end{lstlisting}

\begin{quote}
{[}!CAUTION{]} 上面的 \passthrough{\lstinline!TYPE\_CHECKING!}
示例只表达计划调用形态。该分支在普通 Python
运行时为假，不代表当前扩展能执行此接口。
\end{quote}

\hypertarget{unitree-sdk2-cpp-robot-r1-jsonizedatavecfloat-from-json-1}{%
\paragraph{\texorpdfstring{\texttt{unitree\_sdk2\_cpp.robot.r1.JsonizeDataVecFloat.from\_json}}{unitree\_sdk2\_cpp.robot.r1.JsonizeDataVecFloat.from\_json}}\label{unitree-sdk2-cpp-robot-r1-jsonizedatavecfloat-from-json-1}}

计划从 JSON 风格字典读取字段并更新当前 SDK 值对象。

\textbf{签名}

\begin{lstlisting}[language=Python]
def from_json(self, json: dict[str, Any]) -> None
\end{lstlisting}

\textbf{可用性与安全性}

\textbf{\passthrough{\lstinline!SIGNATURE\_ONLY!} /
\passthrough{\lstinline!UNCLASSIFIED!}}：当前只有设计期签名，不能假设已存在于运行时扩展。

\textbf{参数}

\begin{longtable}[]{@{}
  >{\raggedright\arraybackslash}p{(\columnwidth - 6\tabcolsep) * \real{0.18}}
  >{\raggedright\arraybackslash}p{(\columnwidth - 6\tabcolsep) * \real{0.24}}
  >{\raggedright\arraybackslash}p{(\columnwidth - 6\tabcolsep) * \real{0.18}}
  >{\raggedright\arraybackslash}p{(\columnwidth - 6\tabcolsep) * \real{0.40}}@{}}
\toprule
名称 & Python 类型 & 默认值 & 含义 \\
\midrule
\endhead
\passthrough{\lstinline!json!} &
\passthrough{\lstinline!dict[str, Any]!} & 必填 & JSON
风格字典，键和值必须满足对应 C++ \passthrough{\lstinline!JsonMap!}
数据结构的约束。 对应 C++ 参数
\passthrough{\lstinline!json: common::JsonMap \&!}。 \\
\bottomrule
\end{longtable}

\textbf{返回值}

\begin{longtable}[]{@{}
  >{\raggedright\arraybackslash}p{(\columnwidth - 4\tabcolsep) * \real{0.18}}
  >{\raggedright\arraybackslash}p{(\columnwidth - 4\tabcolsep) * \real{0.20}}
  >{\raggedright\arraybackslash}p{(\columnwidth - 4\tabcolsep) * \real{0.62}}@{}}
\toprule
位置 & 类型 & 含义 \\
\midrule
\endhead
无 & \passthrough{\lstinline!None!} & 该方法不返回 Python
值；失败可能表现为异常或底层状态变化。 \\
\bottomrule
\end{longtable}

\textbf{对应 C++}

\begin{itemize}
\tightlist
\item
  类：\passthrough{\lstinline!unitree::robot::r1::JsonizeDataVecFloat!}
\item
  签名：\passthrough{\lstinline!fromJson(common::JsonMap \&)!}
\item
  绑定策略：\passthrough{\lstinline!SIGNATURE\_PREVIEW!}
\item
  声明位置：\passthrough{\lstinline!include/unitree/robot/r1/loco/r1\_loco\_api.hpp:30!}
\end{itemize}

\textbf{说明}

\begin{itemize}
\tightlist
\item
  示例是局部调用片段；\passthrough{\lstinline!obj!}
  和参数变量表示已经按上层业务校验并准备好的对象。
\item
  该条目可以被编辑器和类型检查器识别，但当前运行时可能在属性查找阶段直接失败。
\end{itemize}

\textbf{用法}

\begin{lstlisting}[language=Python]
from typing import TYPE_CHECKING

if TYPE_CHECKING:
    def planned_from_json(obj: JsonizeDataVecFloat, json: dict[str, Any]) -> None:
        obj.from_json(json=json)
\end{lstlisting}

\begin{quote}
{[}!CAUTION{]} 上面的 \passthrough{\lstinline!TYPE\_CHECKING!}
示例只表达计划调用形态。该分支在普通 Python
运行时为假，不代表当前扩展能执行此接口。
\end{quote}

\hypertarget{unitree-sdk2-cpp-robot-r1-jsonizedatavecfloat-to-json-1}{%
\paragraph{\texorpdfstring{\texttt{unitree\_sdk2\_cpp.robot.r1.JsonizeDataVecFloat.to\_json}}{unitree\_sdk2\_cpp.robot.r1.JsonizeDataVecFloat.to\_json}}\label{unitree-sdk2-cpp-robot-r1-jsonizedatavecfloat-to-json-1}}

计划把当前 SDK 值对象写入 JSON 风格字典。

\textbf{签名}

\begin{lstlisting}[language=Python]
def to_json(self, json: dict[str, Any]) -> None
\end{lstlisting}

\textbf{可用性与安全性}

\textbf{\passthrough{\lstinline!SIGNATURE\_ONLY!} /
\passthrough{\lstinline!UNCLASSIFIED!}}：当前只有设计期签名，不能假设已存在于运行时扩展。

\textbf{参数}

\begin{longtable}[]{@{}
  >{\raggedright\arraybackslash}p{(\columnwidth - 6\tabcolsep) * \real{0.18}}
  >{\raggedright\arraybackslash}p{(\columnwidth - 6\tabcolsep) * \real{0.24}}
  >{\raggedright\arraybackslash}p{(\columnwidth - 6\tabcolsep) * \real{0.18}}
  >{\raggedright\arraybackslash}p{(\columnwidth - 6\tabcolsep) * \real{0.40}}@{}}
\toprule
名称 & Python 类型 & 默认值 & 含义 \\
\midrule
\endhead
\passthrough{\lstinline!json!} &
\passthrough{\lstinline!dict[str, Any]!} & 必填 & JSON
风格字典，键和值必须满足对应 C++ \passthrough{\lstinline!JsonMap!}
数据结构的约束。 对应 C++ 参数
\passthrough{\lstinline!json: common::JsonMap \&!}。 \\
\bottomrule
\end{longtable}

\textbf{返回值}

\begin{longtable}[]{@{}
  >{\raggedright\arraybackslash}p{(\columnwidth - 4\tabcolsep) * \real{0.18}}
  >{\raggedright\arraybackslash}p{(\columnwidth - 4\tabcolsep) * \real{0.20}}
  >{\raggedright\arraybackslash}p{(\columnwidth - 4\tabcolsep) * \real{0.62}}@{}}
\toprule
位置 & 类型 & 含义 \\
\midrule
\endhead
无 & \passthrough{\lstinline!None!} & 该方法不返回 Python
值；失败可能表现为异常或底层状态变化。 \\
\bottomrule
\end{longtable}

\textbf{对应 C++}

\begin{itemize}
\tightlist
\item
  类：\passthrough{\lstinline!unitree::robot::r1::JsonizeDataVecFloat!}
\item
  签名：\passthrough{\lstinline!toJson(common::JsonMap \&) const!}
\item
  绑定策略：\passthrough{\lstinline!SIGNATURE\_PREVIEW!}
\item
  声明位置：\passthrough{\lstinline!include/unitree/robot/r1/loco/r1\_loco\_api.hpp:34!}
\end{itemize}

\textbf{说明}

\begin{itemize}
\tightlist
\item
  示例是局部调用片段；\passthrough{\lstinline!obj!}
  和参数变量表示已经按上层业务校验并准备好的对象。
\item
  该条目可以被编辑器和类型检查器识别，但当前运行时可能在属性查找阶段直接失败。
\end{itemize}

\textbf{用法}

\begin{lstlisting}[language=Python]
from typing import TYPE_CHECKING

if TYPE_CHECKING:
    def planned_to_json(obj: JsonizeDataVecFloat, json: dict[str, Any]) -> None:
        obj.to_json(json=json)
\end{lstlisting}

\begin{quote}
{[}!CAUTION{]} 上面的 \passthrough{\lstinline!TYPE\_CHECKING!}
示例只表达计划调用形态。该分支在普通 Python
运行时为假，不代表当前扩展能执行此接口。
\end{quote}

\hypertarget{unitree-sdk2-cpp-robot-r1-jsonizevelocitycommand}{%
\subsubsection{\texorpdfstring{\texttt{unitree\_sdk2\_cpp.robot.r1.JsonizeVelocityCommand}}{unitree\_sdk2\_cpp.robot.r1.JsonizeVelocityCommand}}\label{unitree-sdk2-cpp-robot-r1-jsonizevelocitycommand}}

SDK 类型签名预览。请逐项查看构造函数和方法的可用性。

\textbf{导入}

\begin{lstlisting}[language=Python]
from unitree_sdk2_cpp.robot.r1 import JsonizeVelocityCommand
\end{lstlisting}

\textbf{基类}：\passthrough{\lstinline!object!}

\textbf{公开属性}

\begin{longtable}[]{@{}
  >{\raggedright\arraybackslash}p{(\columnwidth - 6\tabcolsep) * \real{0.18}}
  >{\raggedright\arraybackslash}p{(\columnwidth - 6\tabcolsep) * \real{0.24}}
  >{\raggedright\arraybackslash}p{(\columnwidth - 6\tabcolsep) * \real{0.18}}
  >{\raggedright\arraybackslash}p{(\columnwidth - 6\tabcolsep) * \real{0.40}}@{}}
\toprule
名称 & Python 类型 & 含义 & 用法 \\
\midrule
\endhead
\passthrough{\lstinline!velocity!} &
\passthrough{\lstinline!list[float]!} & 传给该接口的 速度 参数，Python
类型为
\passthrough{\lstinline!list[float]!}。精确含义、单位、范围和枚举值需查目标型号对应头文件与协议。
& \passthrough{\lstinline!value = obj.velocity!} /
\passthrough{\lstinline!obj.velocity = value!} \\
\passthrough{\lstinline!duration!} & \passthrough{\lstinline!float!} &
操作持续时间参数。精确单位、范围和默认行为以目标型号接口定义为准。 &
\passthrough{\lstinline!value = obj.duration!} /
\passthrough{\lstinline!obj.duration = value!} \\
\bottomrule
\end{longtable}

\textbf{方法索引}

\begin{longtable}[]{@{}
  >{\raggedright\arraybackslash}p{(\columnwidth - 6\tabcolsep) * \real{0.18}}
  >{\raggedright\arraybackslash}p{(\columnwidth - 6\tabcolsep) * \real{0.24}}
  >{\raggedright\arraybackslash}p{(\columnwidth - 6\tabcolsep) * \real{0.18}}
  >{\raggedright\arraybackslash}p{(\columnwidth - 6\tabcolsep) * \real{0.40}}@{}}
\toprule
方法 & Python 签名 & 状态 & 安全分类 \\
\midrule
\endhead
\protect\hyperlink{unitree-sdk2-cpp-robot-r1-jsonizevelocitycommand-dunder-init-1}{\passthrough{\lstinline!\_\_init\_\_!}}
& \passthrough{\lstinline!def \_\_init\_\_(self) -> None!} &
\passthrough{\lstinline!SIGNATURE\_ONLY!} &
\passthrough{\lstinline!CONSTRUCTION!} \\
\protect\hyperlink{unitree-sdk2-cpp-robot-r1-jsonizevelocitycommand-from-json-1}{\passthrough{\lstinline!from\_json!}}
&
\passthrough{\lstinline!def from\_json(self, json: dict[str, Any]) -> None!}
& \passthrough{\lstinline!SIGNATURE\_ONLY!} &
\passthrough{\lstinline!UNCLASSIFIED!} \\
\protect\hyperlink{unitree-sdk2-cpp-robot-r1-jsonizevelocitycommand-to-json-1}{\passthrough{\lstinline!to\_json!}}
&
\passthrough{\lstinline!def to\_json(self, json: dict[str, Any]) -> None!}
& \passthrough{\lstinline!SIGNATURE\_ONLY!} &
\passthrough{\lstinline!UNCLASSIFIED!} \\
\bottomrule
\end{longtable}

\hypertarget{unitree-sdk2-cpp-robot-r1-jsonizevelocitycommand-dunder-init-1}{%
\paragraph{\texorpdfstring{\texttt{unitree\_sdk2\_cpp.robot.r1.JsonizeVelocityCommand.\_\_init\_\_}}{unitree\_sdk2\_cpp.robot.r1.JsonizeVelocityCommand.\_\_init\_\_}}\label{unitree-sdk2-cpp-robot-r1-jsonizevelocitycommand-dunder-init-1}}

初始化 \passthrough{\lstinline!JsonizeVelocityCommand!}
实例。是否能实际构造取决于下方可用性状态。

\textbf{签名}

\begin{lstlisting}[language=Python]
def __init__(self) -> None
\end{lstlisting}

\textbf{可用性与安全性}

\textbf{\passthrough{\lstinline!SIGNATURE\_ONLY!} /
\passthrough{\lstinline!CONSTRUCTION!}}：当前只有设计期签名，不能假设已存在于运行时扩展。

\textbf{参数}

无显式参数。实例方法中的 \passthrough{\lstinline!self!} 由 Python
自动传入。

\textbf{返回值}

\begin{longtable}[]{@{}
  >{\raggedright\arraybackslash}p{(\columnwidth - 4\tabcolsep) * \real{0.18}}
  >{\raggedright\arraybackslash}p{(\columnwidth - 4\tabcolsep) * \real{0.20}}
  >{\raggedright\arraybackslash}p{(\columnwidth - 4\tabcolsep) * \real{0.62}}@{}}
\toprule
位置 & 类型 & 含义 \\
\midrule
\endhead
无 & \passthrough{\lstinline!None!} &
\passthrough{\lstinline!\_\_init\_\_!}
本身不返回值；调用类对象时，在构造函数可用的前提下得到该类实例。 \\
\bottomrule
\end{longtable}

\textbf{对应 C++}

\begin{itemize}
\tightlist
\item
  类：\passthrough{\lstinline!unitree::robot::r1::JsonizeVelocityCommand!}
\item
  签名：\passthrough{\lstinline!JsonizeVelocityCommand()!}
\item
  绑定策略：\passthrough{\lstinline!未单独标注!}
\item
  声明位置：\passthrough{\lstinline!include/unitree/robot/r1/loco/r1\_loco\_api.hpp:43!}
\end{itemize}

\textbf{说明}

\begin{itemize}
\tightlist
\item
  示例是局部调用片段；\passthrough{\lstinline!obj!}
  和参数变量表示已经按上层业务校验并准备好的对象。
\item
  该条目可以被编辑器和类型检查器识别，但当前运行时可能在属性查找阶段直接失败。
\end{itemize}

\textbf{用法}

\begin{lstlisting}[language=Python]
from typing import TYPE_CHECKING

if TYPE_CHECKING:
    def planned_construct() -> JsonizeVelocityCommand:
        return JsonizeVelocityCommand()
\end{lstlisting}

\begin{quote}
{[}!CAUTION{]} 上面的 \passthrough{\lstinline!TYPE\_CHECKING!}
示例只表达计划调用形态。该分支在普通 Python
运行时为假，不代表当前扩展能执行此接口。
\end{quote}

\hypertarget{unitree-sdk2-cpp-robot-r1-jsonizevelocitycommand-from-json-1}{%
\paragraph{\texorpdfstring{\texttt{unitree\_sdk2\_cpp.robot.r1.JsonizeVelocityCommand.from\_json}}{unitree\_sdk2\_cpp.robot.r1.JsonizeVelocityCommand.from\_json}}\label{unitree-sdk2-cpp-robot-r1-jsonizevelocitycommand-from-json-1}}

计划从 JSON 风格字典读取字段并更新当前 SDK 值对象。

\textbf{签名}

\begin{lstlisting}[language=Python]
def from_json(self, json: dict[str, Any]) -> None
\end{lstlisting}

\textbf{可用性与安全性}

\textbf{\passthrough{\lstinline!SIGNATURE\_ONLY!} /
\passthrough{\lstinline!UNCLASSIFIED!}}：当前只有设计期签名，不能假设已存在于运行时扩展。

\textbf{参数}

\begin{longtable}[]{@{}
  >{\raggedright\arraybackslash}p{(\columnwidth - 6\tabcolsep) * \real{0.18}}
  >{\raggedright\arraybackslash}p{(\columnwidth - 6\tabcolsep) * \real{0.24}}
  >{\raggedright\arraybackslash}p{(\columnwidth - 6\tabcolsep) * \real{0.18}}
  >{\raggedright\arraybackslash}p{(\columnwidth - 6\tabcolsep) * \real{0.40}}@{}}
\toprule
名称 & Python 类型 & 默认值 & 含义 \\
\midrule
\endhead
\passthrough{\lstinline!json!} &
\passthrough{\lstinline!dict[str, Any]!} & 必填 & JSON
风格字典，键和值必须满足对应 C++ \passthrough{\lstinline!JsonMap!}
数据结构的约束。 对应 C++ 参数
\passthrough{\lstinline!json: common::JsonMap \&!}。 \\
\bottomrule
\end{longtable}

\textbf{返回值}

\begin{longtable}[]{@{}
  >{\raggedright\arraybackslash}p{(\columnwidth - 4\tabcolsep) * \real{0.18}}
  >{\raggedright\arraybackslash}p{(\columnwidth - 4\tabcolsep) * \real{0.20}}
  >{\raggedright\arraybackslash}p{(\columnwidth - 4\tabcolsep) * \real{0.62}}@{}}
\toprule
位置 & 类型 & 含义 \\
\midrule
\endhead
无 & \passthrough{\lstinline!None!} & 该方法不返回 Python
值；失败可能表现为异常或底层状态变化。 \\
\bottomrule
\end{longtable}

\textbf{对应 C++}

\begin{itemize}
\tightlist
\item
  类：\passthrough{\lstinline!unitree::robot::r1::JsonizeVelocityCommand!}
\item
  签名：\passthrough{\lstinline!fromJson(common::JsonMap \&)!}
\item
  绑定策略：\passthrough{\lstinline!SIGNATURE\_PREVIEW!}
\item
  声明位置：\passthrough{\lstinline!include/unitree/robot/r1/loco/r1\_loco\_api.hpp:46!}
\end{itemize}

\textbf{说明}

\begin{itemize}
\tightlist
\item
  示例是局部调用片段；\passthrough{\lstinline!obj!}
  和参数变量表示已经按上层业务校验并准备好的对象。
\item
  该条目可以被编辑器和类型检查器识别，但当前运行时可能在属性查找阶段直接失败。
\end{itemize}

\textbf{用法}

\begin{lstlisting}[language=Python]
from typing import TYPE_CHECKING

if TYPE_CHECKING:
    def planned_from_json(obj: JsonizeVelocityCommand, json: dict[str, Any]) -> None:
        obj.from_json(json=json)
\end{lstlisting}

\begin{quote}
{[}!CAUTION{]} 上面的 \passthrough{\lstinline!TYPE\_CHECKING!}
示例只表达计划调用形态。该分支在普通 Python
运行时为假，不代表当前扩展能执行此接口。
\end{quote}

\hypertarget{unitree-sdk2-cpp-robot-r1-jsonizevelocitycommand-to-json-1}{%
\paragraph{\texorpdfstring{\texttt{unitree\_sdk2\_cpp.robot.r1.JsonizeVelocityCommand.to\_json}}{unitree\_sdk2\_cpp.robot.r1.JsonizeVelocityCommand.to\_json}}\label{unitree-sdk2-cpp-robot-r1-jsonizevelocitycommand-to-json-1}}

计划把当前 SDK 值对象写入 JSON 风格字典。

\textbf{签名}

\begin{lstlisting}[language=Python]
def to_json(self, json: dict[str, Any]) -> None
\end{lstlisting}

\textbf{可用性与安全性}

\textbf{\passthrough{\lstinline!SIGNATURE\_ONLY!} /
\passthrough{\lstinline!UNCLASSIFIED!}}：当前只有设计期签名，不能假设已存在于运行时扩展。

\textbf{参数}

\begin{longtable}[]{@{}
  >{\raggedright\arraybackslash}p{(\columnwidth - 6\tabcolsep) * \real{0.18}}
  >{\raggedright\arraybackslash}p{(\columnwidth - 6\tabcolsep) * \real{0.24}}
  >{\raggedright\arraybackslash}p{(\columnwidth - 6\tabcolsep) * \real{0.18}}
  >{\raggedright\arraybackslash}p{(\columnwidth - 6\tabcolsep) * \real{0.40}}@{}}
\toprule
名称 & Python 类型 & 默认值 & 含义 \\
\midrule
\endhead
\passthrough{\lstinline!json!} &
\passthrough{\lstinline!dict[str, Any]!} & 必填 & JSON
风格字典，键和值必须满足对应 C++ \passthrough{\lstinline!JsonMap!}
数据结构的约束。 对应 C++ 参数
\passthrough{\lstinline!json: common::JsonMap \&!}。 \\
\bottomrule
\end{longtable}

\textbf{返回值}

\begin{longtable}[]{@{}
  >{\raggedright\arraybackslash}p{(\columnwidth - 4\tabcolsep) * \real{0.18}}
  >{\raggedright\arraybackslash}p{(\columnwidth - 4\tabcolsep) * \real{0.20}}
  >{\raggedright\arraybackslash}p{(\columnwidth - 4\tabcolsep) * \real{0.62}}@{}}
\toprule
位置 & 类型 & 含义 \\
\midrule
\endhead
无 & \passthrough{\lstinline!None!} & 该方法不返回 Python
值；失败可能表现为异常或底层状态变化。 \\
\bottomrule
\end{longtable}

\textbf{对应 C++}

\begin{itemize}
\tightlist
\item
  类：\passthrough{\lstinline!unitree::robot::r1::JsonizeVelocityCommand!}
\item
  签名：\passthrough{\lstinline!toJson(common::JsonMap \&) const!}
\item
  绑定策略：\passthrough{\lstinline!SIGNATURE\_PREVIEW!}
\item
  声明位置：\passthrough{\lstinline!include/unitree/robot/r1/loco/r1\_loco\_api.hpp:51!}
\end{itemize}

\textbf{说明}

\begin{itemize}
\tightlist
\item
  示例是局部调用片段；\passthrough{\lstinline!obj!}
  和参数变量表示已经按上层业务校验并准备好的对象。
\item
  该条目可以被编辑器和类型检查器识别，但当前运行时可能在属性查找阶段直接失败。
\end{itemize}

\textbf{用法}

\begin{lstlisting}[language=Python]
from typing import TYPE_CHECKING

if TYPE_CHECKING:
    def planned_to_json(obj: JsonizeVelocityCommand, json: dict[str, Any]) -> None:
        obj.to_json(json=json)
\end{lstlisting}

\begin{quote}
{[}!CAUTION{]} 上面的 \passthrough{\lstinline!TYPE\_CHECKING!}
示例只表达计划调用形态。该分支在普通 Python
运行时为假，不代表当前扩展能执行此接口。
\end{quote}

\hypertarget{unitree-sdk2-cpp-robot-r1-ledcontrolparameter}{%
\subsubsection{\texorpdfstring{\texttt{unitree\_sdk2\_cpp.robot.r1.LedControlParameter}}{unitree\_sdk2\_cpp.robot.r1.LedControlParameter}}\label{unitree-sdk2-cpp-robot-r1-ledcontrolparameter}}

SDK 请求参数值类型；当前可能仅提供设计期签名。

\textbf{导入}

\begin{lstlisting}[language=Python]
from unitree_sdk2_cpp.robot.r1 import LedControlParameter
\end{lstlisting}

\textbf{基类}：\passthrough{\lstinline!object!}

\textbf{公开属性}

\begin{longtable}[]{@{}
  >{\raggedright\arraybackslash}p{(\columnwidth - 6\tabcolsep) * \real{0.18}}
  >{\raggedright\arraybackslash}p{(\columnwidth - 6\tabcolsep) * \real{0.24}}
  >{\raggedright\arraybackslash}p{(\columnwidth - 6\tabcolsep) * \real{0.18}}
  >{\raggedright\arraybackslash}p{(\columnwidth - 6\tabcolsep) * \real{0.40}}@{}}
\toprule
名称 & Python 类型 & 含义 & 用法 \\
\midrule
\endhead
\passthrough{\lstinline!r!} & \passthrough{\lstinline!int!} &
传给该接口的 \passthrough{\lstinline!r!} 参数，Python 类型为
\passthrough{\lstinline!int!}。精确含义、单位、范围和枚举值需查目标型号对应头文件与协议。
& \passthrough{\lstinline!value = obj.r!} /
\passthrough{\lstinline!obj.r = value!} \\
\passthrough{\lstinline!g!} & \passthrough{\lstinline!int!} &
传给该接口的 \passthrough{\lstinline!g!} 参数，Python 类型为
\passthrough{\lstinline!int!}。精确含义、单位、范围和枚举值需查目标型号对应头文件与协议。
& \passthrough{\lstinline!value = obj.g!} /
\passthrough{\lstinline!obj.g = value!} \\
\passthrough{\lstinline!b!} & \passthrough{\lstinline!int!} &
传给该接口的 \passthrough{\lstinline!b!} 参数，Python 类型为
\passthrough{\lstinline!int!}。精确含义、单位、范围和枚举值需查目标型号对应头文件与协议。
& \passthrough{\lstinline!value = obj.b!} /
\passthrough{\lstinline!obj.b = value!} \\
\bottomrule
\end{longtable}

\textbf{方法索引}

\begin{longtable}[]{@{}
  >{\raggedright\arraybackslash}p{(\columnwidth - 6\tabcolsep) * \real{0.18}}
  >{\raggedright\arraybackslash}p{(\columnwidth - 6\tabcolsep) * \real{0.24}}
  >{\raggedright\arraybackslash}p{(\columnwidth - 6\tabcolsep) * \real{0.18}}
  >{\raggedright\arraybackslash}p{(\columnwidth - 6\tabcolsep) * \real{0.40}}@{}}
\toprule
方法 & Python 签名 & 状态 & 安全分类 \\
\midrule
\endhead
\protect\hyperlink{unitree-sdk2-cpp-robot-r1-ledcontrolparameter-dunder-init-1}{\passthrough{\lstinline!\_\_init\_\_!}}
& \passthrough{\lstinline!def \_\_init\_\_(self) -> None!} &
\passthrough{\lstinline!SIGNATURE\_ONLY!} &
\passthrough{\lstinline!CONSTRUCTION!} \\
\protect\hyperlink{unitree-sdk2-cpp-robot-r1-ledcontrolparameter-from-json-1}{\passthrough{\lstinline!from\_json!}}
&
\passthrough{\lstinline!def from\_json(self, json: dict[str, Any]) -> None!}
& \passthrough{\lstinline!SIGNATURE\_ONLY!} &
\passthrough{\lstinline!UNCLASSIFIED!} \\
\protect\hyperlink{unitree-sdk2-cpp-robot-r1-ledcontrolparameter-to-json-1}{\passthrough{\lstinline!to\_json!}}
&
\passthrough{\lstinline!def to\_json(self, json: dict[str, Any]) -> None!}
& \passthrough{\lstinline!SIGNATURE\_ONLY!} &
\passthrough{\lstinline!UNCLASSIFIED!} \\
\bottomrule
\end{longtable}

\hypertarget{unitree-sdk2-cpp-robot-r1-ledcontrolparameter-dunder-init-1}{%
\paragraph{\texorpdfstring{\texttt{unitree\_sdk2\_cpp.robot.r1.LedControlParameter.\_\_init\_\_}}{unitree\_sdk2\_cpp.robot.r1.LedControlParameter.\_\_init\_\_}}\label{unitree-sdk2-cpp-robot-r1-ledcontrolparameter-dunder-init-1}}

初始化 \passthrough{\lstinline!LedControlParameter!}
实例。是否能实际构造取决于下方可用性状态。

\textbf{签名}

\begin{lstlisting}[language=Python]
def __init__(self) -> None
\end{lstlisting}

\textbf{可用性与安全性}

\textbf{\passthrough{\lstinline!SIGNATURE\_ONLY!} /
\passthrough{\lstinline!CONSTRUCTION!}}：当前只有设计期签名，不能假设已存在于运行时扩展。

\textbf{参数}

无显式参数。实例方法中的 \passthrough{\lstinline!self!} 由 Python
自动传入。

\textbf{返回值}

\begin{longtable}[]{@{}
  >{\raggedright\arraybackslash}p{(\columnwidth - 4\tabcolsep) * \real{0.18}}
  >{\raggedright\arraybackslash}p{(\columnwidth - 4\tabcolsep) * \real{0.20}}
  >{\raggedright\arraybackslash}p{(\columnwidth - 4\tabcolsep) * \real{0.62}}@{}}
\toprule
位置 & 类型 & 含义 \\
\midrule
\endhead
无 & \passthrough{\lstinline!None!} &
\passthrough{\lstinline!\_\_init\_\_!}
本身不返回值；调用类对象时，在构造函数可用的前提下得到该类实例。 \\
\bottomrule
\end{longtable}

\textbf{对应 C++}

\begin{itemize}
\tightlist
\item
  类：\passthrough{\lstinline!unitree::robot::r1::LedControlParameter!}
\item
  签名：\passthrough{\lstinline!LedControlParameter()!}
\item
  绑定策略：\passthrough{\lstinline!未单独标注!}
\item
  声明位置：\passthrough{\lstinline!include/unitree/robot/r1/audio/audio\_api.hpp:75!}
\end{itemize}

\textbf{说明}

\begin{itemize}
\tightlist
\item
  示例是局部调用片段；\passthrough{\lstinline!obj!}
  和参数变量表示已经按上层业务校验并准备好的对象。
\item
  该条目可以被编辑器和类型检查器识别，但当前运行时可能在属性查找阶段直接失败。
\end{itemize}

\textbf{用法}

\begin{lstlisting}[language=Python]
from typing import TYPE_CHECKING

if TYPE_CHECKING:
    def planned_construct() -> LedControlParameter:
        return LedControlParameter()
\end{lstlisting}

\begin{quote}
{[}!CAUTION{]} 上面的 \passthrough{\lstinline!TYPE\_CHECKING!}
示例只表达计划调用形态。该分支在普通 Python
运行时为假，不代表当前扩展能执行此接口。
\end{quote}

\hypertarget{unitree-sdk2-cpp-robot-r1-ledcontrolparameter-from-json-1}{%
\paragraph{\texorpdfstring{\texttt{unitree\_sdk2\_cpp.robot.r1.LedControlParameter.from\_json}}{unitree\_sdk2\_cpp.robot.r1.LedControlParameter.from\_json}}\label{unitree-sdk2-cpp-robot-r1-ledcontrolparameter-from-json-1}}

计划从 JSON 风格字典读取字段并更新当前 SDK 值对象。

\textbf{签名}

\begin{lstlisting}[language=Python]
def from_json(self, json: dict[str, Any]) -> None
\end{lstlisting}

\textbf{可用性与安全性}

\textbf{\passthrough{\lstinline!SIGNATURE\_ONLY!} /
\passthrough{\lstinline!UNCLASSIFIED!}}：当前只有设计期签名，不能假设已存在于运行时扩展。

\textbf{参数}

\begin{longtable}[]{@{}
  >{\raggedright\arraybackslash}p{(\columnwidth - 6\tabcolsep) * \real{0.18}}
  >{\raggedright\arraybackslash}p{(\columnwidth - 6\tabcolsep) * \real{0.24}}
  >{\raggedright\arraybackslash}p{(\columnwidth - 6\tabcolsep) * \real{0.18}}
  >{\raggedright\arraybackslash}p{(\columnwidth - 6\tabcolsep) * \real{0.40}}@{}}
\toprule
名称 & Python 类型 & 默认值 & 含义 \\
\midrule
\endhead
\passthrough{\lstinline!json!} &
\passthrough{\lstinline!dict[str, Any]!} & 必填 & JSON
风格字典，键和值必须满足对应 C++ \passthrough{\lstinline!JsonMap!}
数据结构的约束。 对应 C++ 参数
\passthrough{\lstinline!json: common::JsonMap \&!}。 \\
\bottomrule
\end{longtable}

\textbf{返回值}

\begin{longtable}[]{@{}
  >{\raggedright\arraybackslash}p{(\columnwidth - 4\tabcolsep) * \real{0.18}}
  >{\raggedright\arraybackslash}p{(\columnwidth - 4\tabcolsep) * \real{0.20}}
  >{\raggedright\arraybackslash}p{(\columnwidth - 4\tabcolsep) * \real{0.62}}@{}}
\toprule
位置 & 类型 & 含义 \\
\midrule
\endhead
无 & \passthrough{\lstinline!None!} & 该方法不返回 Python
值；失败可能表现为异常或底层状态变化。 \\
\bottomrule
\end{longtable}

\textbf{对应 C++}

\begin{itemize}
\tightlist
\item
  类：\passthrough{\lstinline!unitree::robot::r1::LedControlParameter!}
\item
  签名：\passthrough{\lstinline!fromJson(common::JsonMap \&)!}
\item
  绑定策略：\passthrough{\lstinline!SIGNATURE\_PREVIEW!}
\item
  声明位置：\passthrough{\lstinline!include/unitree/robot/r1/audio/audio\_api.hpp:78!}
\end{itemize}

\textbf{说明}

\begin{itemize}
\tightlist
\item
  示例是局部调用片段；\passthrough{\lstinline!obj!}
  和参数变量表示已经按上层业务校验并准备好的对象。
\item
  该条目可以被编辑器和类型检查器识别，但当前运行时可能在属性查找阶段直接失败。
\end{itemize}

\textbf{用法}

\begin{lstlisting}[language=Python]
from typing import TYPE_CHECKING

if TYPE_CHECKING:
    def planned_from_json(obj: LedControlParameter, json: dict[str, Any]) -> None:
        obj.from_json(json=json)
\end{lstlisting}

\begin{quote}
{[}!CAUTION{]} 上面的 \passthrough{\lstinline!TYPE\_CHECKING!}
示例只表达计划调用形态。该分支在普通 Python
运行时为假，不代表当前扩展能执行此接口。
\end{quote}

\hypertarget{unitree-sdk2-cpp-robot-r1-ledcontrolparameter-to-json-1}{%
\paragraph{\texorpdfstring{\texttt{unitree\_sdk2\_cpp.robot.r1.LedControlParameter.to\_json}}{unitree\_sdk2\_cpp.robot.r1.LedControlParameter.to\_json}}\label{unitree-sdk2-cpp-robot-r1-ledcontrolparameter-to-json-1}}

计划把当前 SDK 值对象写入 JSON 风格字典。

\textbf{签名}

\begin{lstlisting}[language=Python]
def to_json(self, json: dict[str, Any]) -> None
\end{lstlisting}

\textbf{可用性与安全性}

\textbf{\passthrough{\lstinline!SIGNATURE\_ONLY!} /
\passthrough{\lstinline!UNCLASSIFIED!}}：当前只有设计期签名，不能假设已存在于运行时扩展。

\textbf{参数}

\begin{longtable}[]{@{}
  >{\raggedright\arraybackslash}p{(\columnwidth - 6\tabcolsep) * \real{0.18}}
  >{\raggedright\arraybackslash}p{(\columnwidth - 6\tabcolsep) * \real{0.24}}
  >{\raggedright\arraybackslash}p{(\columnwidth - 6\tabcolsep) * \real{0.18}}
  >{\raggedright\arraybackslash}p{(\columnwidth - 6\tabcolsep) * \real{0.40}}@{}}
\toprule
名称 & Python 类型 & 默认值 & 含义 \\
\midrule
\endhead
\passthrough{\lstinline!json!} &
\passthrough{\lstinline!dict[str, Any]!} & 必填 & JSON
风格字典，键和值必须满足对应 C++ \passthrough{\lstinline!JsonMap!}
数据结构的约束。 对应 C++ 参数
\passthrough{\lstinline!json: common::JsonMap \&!}。 \\
\bottomrule
\end{longtable}

\textbf{返回值}

\begin{longtable}[]{@{}
  >{\raggedright\arraybackslash}p{(\columnwidth - 4\tabcolsep) * \real{0.18}}
  >{\raggedright\arraybackslash}p{(\columnwidth - 4\tabcolsep) * \real{0.20}}
  >{\raggedright\arraybackslash}p{(\columnwidth - 4\tabcolsep) * \real{0.62}}@{}}
\toprule
位置 & 类型 & 含义 \\
\midrule
\endhead
无 & \passthrough{\lstinline!None!} & 该方法不返回 Python
值；失败可能表现为异常或底层状态变化。 \\
\bottomrule
\end{longtable}

\textbf{对应 C++}

\begin{itemize}
\tightlist
\item
  类：\passthrough{\lstinline!unitree::robot::r1::LedControlParameter!}
\item
  签名：\passthrough{\lstinline!toJson(common::JsonMap \&) const!}
\item
  绑定策略：\passthrough{\lstinline!SIGNATURE\_PREVIEW!}
\item
  声明位置：\passthrough{\lstinline!include/unitree/robot/r1/audio/audio\_api.hpp:80!}
\end{itemize}

\textbf{说明}

\begin{itemize}
\tightlist
\item
  示例是局部调用片段；\passthrough{\lstinline!obj!}
  和参数变量表示已经按上层业务校验并准备好的对象。
\item
  该条目可以被编辑器和类型检查器识别，但当前运行时可能在属性查找阶段直接失败。
\end{itemize}

\textbf{用法}

\begin{lstlisting}[language=Python]
from typing import TYPE_CHECKING

if TYPE_CHECKING:
    def planned_to_json(obj: LedControlParameter, json: dict[str, Any]) -> None:
        obj.to_json(json=json)
\end{lstlisting}

\begin{quote}
{[}!CAUTION{]} 上面的 \passthrough{\lstinline!TYPE\_CHECKING!}
示例只表达计划调用形态。该分支在普通 Python
运行时为假，不代表当前扩展能执行此接口。
\end{quote}

\hypertarget{unitree-sdk2-cpp-robot-r1-lococlient}{%
\subsubsection{\texorpdfstring{\texttt{unitree\_sdk2\_cpp.robot.r1.LocoClient}}{unitree\_sdk2\_cpp.robot.r1.LocoClient}}\label{unitree-sdk2-cpp-robot-r1-lococlient}}

Robot 服务客户端。构造、初始化和具体方法的可用性必须分别检查。

\textbf{导入}

\begin{lstlisting}[language=Python]
from unitree_sdk2_cpp.robot.r1 import LocoClient
\end{lstlisting}

\textbf{基类}：\passthrough{\lstinline!Client!}

\textbf{方法索引}

\begin{longtable}[]{@{}
  >{\raggedright\arraybackslash}p{(\columnwidth - 6\tabcolsep) * \real{0.18}}
  >{\raggedright\arraybackslash}p{(\columnwidth - 6\tabcolsep) * \real{0.24}}
  >{\raggedright\arraybackslash}p{(\columnwidth - 6\tabcolsep) * \real{0.18}}
  >{\raggedright\arraybackslash}p{(\columnwidth - 6\tabcolsep) * \real{0.40}}@{}}
\toprule
方法 & Python 签名 & 状态 & 安全分类 \\
\midrule
\endhead
\protect\hyperlink{unitree-sdk2-cpp-robot-r1-lococlient-dunder-init-1}{\passthrough{\lstinline!\_\_init\_\_!}}
& \passthrough{\lstinline!def \_\_init\_\_(self) -> None!} &
\passthrough{\lstinline!AVAILABLE!} &
\passthrough{\lstinline!CONSTRUCTION!} \\
\protect\hyperlink{unitree-sdk2-cpp-robot-r1-lococlient-init-1}{\passthrough{\lstinline!init!}}
& \passthrough{\lstinline!def init(self) -> None!} &
\passthrough{\lstinline!AVAILABLE!} &
\passthrough{\lstinline!INITIALIZATION!} \\
\protect\hyperlink{unitree-sdk2-cpp-robot-r1-lococlient-get-fsm-id-1}{\passthrough{\lstinline!get\_fsm\_id!}}
& \passthrough{\lstinline!def get\_fsm\_id(self) -> tuple[int, int]!} &
\passthrough{\lstinline!AVAILABLE!} &
\passthrough{\lstinline!READ\_ONLY!} \\
\protect\hyperlink{unitree-sdk2-cpp-robot-r1-lococlient-get-fsm-mode-1}{\passthrough{\lstinline!get\_fsm\_mode!}}
& \passthrough{\lstinline!def get\_fsm\_mode(self) -> tuple[int, int]!}
& \passthrough{\lstinline!AVAILABLE!} &
\passthrough{\lstinline!READ\_ONLY!} \\
\protect\hyperlink{unitree-sdk2-cpp-robot-r1-lococlient-set-fsm-id-1}{\passthrough{\lstinline!set\_fsm\_id!}}
& \passthrough{\lstinline!def set\_fsm\_id(self, fsm\_id: int) -> int!}
& \passthrough{\lstinline!SIGNATURE\_ONLY!} &
\passthrough{\lstinline!MOTION\_COMMAND!} \\
\protect\hyperlink{unitree-sdk2-cpp-robot-r1-lococlient-set-velocity-1}{\passthrough{\lstinline!set\_velocity!}}
&
\passthrough{\lstinline!def set\_velocity(self, vx: float, vy: float, omega: float, duration: float = ...) -> int!}
& \passthrough{\lstinline!SIGNATURE\_ONLY!} &
\passthrough{\lstinline!MOTION\_COMMAND!} \\
\protect\hyperlink{unitree-sdk2-cpp-robot-r1-lococlient-damp-1}{\passthrough{\lstinline!damp!}}
& \passthrough{\lstinline!def damp(self) -> int!} &
\passthrough{\lstinline!SIGNATURE\_ONLY!} &
\passthrough{\lstinline!MOTION\_COMMAND!} \\
\protect\hyperlink{unitree-sdk2-cpp-robot-r1-lococlient-start-1}{\passthrough{\lstinline!start!}}
& \passthrough{\lstinline!def start(self) -> int!} &
\passthrough{\lstinline!SIGNATURE\_ONLY!} &
\passthrough{\lstinline!MOTION\_COMMAND!} \\
\protect\hyperlink{unitree-sdk2-cpp-robot-r1-lococlient-stand-up-1}{\passthrough{\lstinline!stand\_up!}}
& \passthrough{\lstinline!def stand\_up(self) -> int!} &
\passthrough{\lstinline!SIGNATURE\_ONLY!} &
\passthrough{\lstinline!MOTION\_COMMAND!} \\
\protect\hyperlink{unitree-sdk2-cpp-robot-r1-lococlient-zero-torque-1}{\passthrough{\lstinline!zero\_torque!}}
& \passthrough{\lstinline!def zero\_torque(self) -> int!} &
\passthrough{\lstinline!SIGNATURE\_ONLY!} &
\passthrough{\lstinline!MOTION\_COMMAND!} \\
\protect\hyperlink{unitree-sdk2-cpp-robot-r1-lococlient-stop-move-1}{\passthrough{\lstinline!stop\_move!}}
& \passthrough{\lstinline!def stop\_move(self) -> int!} &
\passthrough{\lstinline!SIGNATURE\_ONLY!} &
\passthrough{\lstinline!MOTION\_COMMAND!} \\
\protect\hyperlink{unitree-sdk2-cpp-robot-r1-lococlient-move-1}{\passthrough{\lstinline!move（重载 1/2）!}}
&
\passthrough{\lstinline!def move(self, vx: float, vy: float, vyaw: float, continous\_move: bool) -> int!}
& \passthrough{\lstinline!SIGNATURE\_ONLY!} &
\passthrough{\lstinline!MOTION\_COMMAND!} \\
\protect\hyperlink{unitree-sdk2-cpp-robot-r1-lococlient-move-2}{\passthrough{\lstinline!move（重载 2/2）!}}
&
\passthrough{\lstinline!def move(self, vx: float, vy: float, vyaw: float) -> int!}
& \passthrough{\lstinline!SIGNATURE\_ONLY!} &
\passthrough{\lstinline!MOTION\_COMMAND!} \\
\protect\hyperlink{unitree-sdk2-cpp-robot-r1-lococlient-switch-move-mode-1}{\passthrough{\lstinline!switch\_move\_mode!}}
&
\passthrough{\lstinline!def switch\_move\_mode(self, flag: bool) -> int!}
& \passthrough{\lstinline!SIGNATURE\_ONLY!} &
\passthrough{\lstinline!MOTION\_COMMAND!} \\
\protect\hyperlink{unitree-sdk2-cpp-robot-r1-lococlient-set-speed-mode-1}{\passthrough{\lstinline!set\_speed\_mode!}}
&
\passthrough{\lstinline!def set\_speed\_mode(self, speed\_mode: int) -> int!}
& \passthrough{\lstinline!SIGNATURE\_ONLY!} &
\passthrough{\lstinline!MOTION\_COMMAND!} \\
\bottomrule
\end{longtable}

\hypertarget{unitree-sdk2-cpp-robot-r1-lococlient-dunder-init-1}{%
\paragraph{\texorpdfstring{\texttt{unitree\_sdk2\_cpp.robot.r1.LocoClient.\_\_init\_\_}}{unitree\_sdk2\_cpp.robot.r1.LocoClient.\_\_init\_\_}}\label{unitree-sdk2-cpp-robot-r1-lococlient-dunder-init-1}}

初始化 \passthrough{\lstinline!LocoClient!}
实例。是否能实际构造取决于下方可用性状态。

\textbf{签名}

\begin{lstlisting}[language=Python]
def __init__(self) -> None
\end{lstlisting}

\textbf{可用性与安全性}

\textbf{\passthrough{\lstinline!AVAILABLE!} /
\passthrough{\lstinline!CONSTRUCTION!}}：当前绑定源码已实现；实际执行条件仍取决于平台和对应
SDK 环境。

\textbf{参数}

无显式参数。实例方法中的 \passthrough{\lstinline!self!} 由 Python
自动传入。

\textbf{返回值}

\begin{longtable}[]{@{}
  >{\raggedright\arraybackslash}p{(\columnwidth - 4\tabcolsep) * \real{0.18}}
  >{\raggedright\arraybackslash}p{(\columnwidth - 4\tabcolsep) * \real{0.20}}
  >{\raggedright\arraybackslash}p{(\columnwidth - 4\tabcolsep) * \real{0.62}}@{}}
\toprule
位置 & 类型 & 含义 \\
\midrule
\endhead
无 & \passthrough{\lstinline!None!} &
\passthrough{\lstinline!\_\_init\_\_!}
本身不返回值；调用类对象时，在构造函数可用的前提下得到该类实例。 \\
\bottomrule
\end{longtable}

\textbf{对应 C++}

\begin{itemize}
\tightlist
\item
  类：\passthrough{\lstinline!unitree::robot::r1::LocoClient!}
\item
  签名：\passthrough{\lstinline!LocoClient()!}
\item
  绑定策略：\passthrough{\lstinline!未单独标注!}
\item
  声明位置：\passthrough{\lstinline!include/unitree/robot/r1/loco/r1\_loco\_client.hpp:13!}
\end{itemize}

\textbf{说明}

\begin{itemize}
\tightlist
\item
  示例是局部调用片段；\passthrough{\lstinline!obj!}
  和参数变量表示已经按上层业务校验并准备好的对象。
\end{itemize}

\textbf{用法}

\begin{lstlisting}[language=Python]
value = LocoClient()
\end{lstlisting}

\hypertarget{unitree-sdk2-cpp-robot-r1-lococlient-init-1}{%
\paragraph{\texorpdfstring{\texttt{unitree\_sdk2\_cpp.robot.r1.LocoClient.init}}{unitree\_sdk2\_cpp.robot.r1.LocoClient.init}}\label{unitree-sdk2-cpp-robot-r1-lococlient-init-1}}

初始化当前 SDK 对象所需的底层通道或服务资源。

\textbf{签名}

\begin{lstlisting}[language=Python]
def init(self) -> None
\end{lstlisting}

\textbf{可用性与安全性}

\textbf{\passthrough{\lstinline!AVAILABLE!} /
\passthrough{\lstinline!INITIALIZATION!}}：当前绑定源码已实现；实际执行条件仍取决于平台和对应
SDK 环境。

\textbf{参数}

无显式参数。实例方法中的 \passthrough{\lstinline!self!} 由 Python
自动传入。

\textbf{返回值}

\begin{longtable}[]{@{}
  >{\raggedright\arraybackslash}p{(\columnwidth - 4\tabcolsep) * \real{0.18}}
  >{\raggedright\arraybackslash}p{(\columnwidth - 4\tabcolsep) * \real{0.20}}
  >{\raggedright\arraybackslash}p{(\columnwidth - 4\tabcolsep) * \real{0.62}}@{}}
\toprule
位置 & 类型 & 含义 \\
\midrule
\endhead
无 & \passthrough{\lstinline!None!} & 该方法不返回 Python
值；失败可能表现为异常或底层状态变化。 \\
\bottomrule
\end{longtable}

\textbf{对应 C++}

\begin{itemize}
\tightlist
\item
  类：\passthrough{\lstinline!unitree::robot::r1::LocoClient!}
\item
  签名：\passthrough{\lstinline!Init()!}
\item
  绑定策略：\passthrough{\lstinline!DIRECT!}
\item
  声明位置：\passthrough{\lstinline!include/unitree/robot/r1/loco/r1\_loco\_client.hpp:17!}
\end{itemize}

\textbf{说明}

\begin{itemize}
\tightlist
\item
  示例是局部调用片段；\passthrough{\lstinline!obj!}
  和参数变量表示已经按上层业务校验并准备好的对象。
\end{itemize}

\textbf{用法}

\begin{lstlisting}[language=Python]
obj.init()
\end{lstlisting}

\hypertarget{unitree-sdk2-cpp-robot-r1-lococlient-get-fsm-id-1}{%
\paragraph{\texorpdfstring{\texttt{unitree\_sdk2\_cpp.robot.r1.LocoClient.get\_fsm\_id}}{unitree\_sdk2\_cpp.robot.r1.LocoClient.get\_fsm\_id}}\label{unitree-sdk2-cpp-robot-r1-lococlient-get-fsm-id-1}}

查询或检查 FSM ID。

\textbf{签名}

\begin{lstlisting}[language=Python]
def get_fsm_id(self) -> tuple[int, int]
\end{lstlisting}

\textbf{可用性与安全性}

\textbf{\passthrough{\lstinline!AVAILABLE!} /
\passthrough{\lstinline!READ\_ONLY!}}：当前绑定源码已实现只读查询。Robot
Client 的服务端查询通常仍需匹配的 Linux
扩展、DDS、网络、机器人和服务；纯本地 getter 则不一定需要全部条件。

\textbf{参数}

无显式参数。实例方法中的 \passthrough{\lstinline!self!} 由 Python
自动传入。

\textbf{返回值}

\begin{longtable}[]{@{}
  >{\raggedright\arraybackslash}p{(\columnwidth - 4\tabcolsep) * \real{0.18}}
  >{\raggedright\arraybackslash}p{(\columnwidth - 4\tabcolsep) * \real{0.20}}
  >{\raggedright\arraybackslash}p{(\columnwidth - 4\tabcolsep) * \real{0.62}}@{}}
\toprule
位置 & 类型 & 含义 \\
\midrule
\endhead
{[}0{]} \passthrough{\lstinline!status!} & \passthrough{\lstinline!int!}
& SDK 状态码；Unitree 示例通常以 \passthrough{\lstinline!0!}
表示成功，非零值需按具体服务错误码解释。 \\
{[}1{]} \passthrough{\lstinline!fsm\_id!} &
\passthrough{\lstinline!int!} & 传给该接口的 FSM ID 参数，Python 类型为
\passthrough{\lstinline!int!}。精确含义、单位、范围和枚举值需查目标型号对应头文件与协议。
对应 C++ 参数 \passthrough{\lstinline!fsm\_id: int \&!}。 \\
\bottomrule
\end{longtable}

\textbf{对应 C++}

\begin{itemize}
\tightlist
\item
  类：\passthrough{\lstinline!unitree::robot::r1::LocoClient!}
\item
  签名：\passthrough{\lstinline!GetFsmId(int \&)!}
\item
  绑定策略：\passthrough{\lstinline!OUTPUT\_WRAPPER!}
\item
  声明位置：\passthrough{\lstinline!include/unitree/robot/r1/loco/r1\_loco\_client.hpp:28!}
\end{itemize}

\textbf{说明}

\begin{itemize}
\tightlist
\item
  示例是局部调用片段；\passthrough{\lstinline!obj!}
  和参数变量表示已经按上层业务校验并准备好的对象。
\item
  C++ 的可修改输出引用不会作为 Python 入参出现，而是按 C++
  参数顺序追加到返回元组中。
\end{itemize}

\textbf{用法}

\begin{lstlisting}[language=Python]
status, fsm_id = obj.get_fsm_id()
\end{lstlisting}

\begin{quote}
{[}!NOTE{]} 只读查询仍会访问
DDS/机器人服务。必须检查返回状态码，并处理超时和网络中断。
\end{quote}

\hypertarget{unitree-sdk2-cpp-robot-r1-lococlient-get-fsm-mode-1}{%
\paragraph{\texorpdfstring{\texttt{unitree\_sdk2\_cpp.robot.r1.LocoClient.get\_fsm\_mode}}{unitree\_sdk2\_cpp.robot.r1.LocoClient.get\_fsm\_mode}}\label{unitree-sdk2-cpp-robot-r1-lococlient-get-fsm-mode-1}}

查询或检查 FSM 模式。

\textbf{签名}

\begin{lstlisting}[language=Python]
def get_fsm_mode(self) -> tuple[int, int]
\end{lstlisting}

\textbf{可用性与安全性}

\textbf{\passthrough{\lstinline!AVAILABLE!} /
\passthrough{\lstinline!READ\_ONLY!}}：当前绑定源码已实现只读查询。Robot
Client 的服务端查询通常仍需匹配的 Linux
扩展、DDS、网络、机器人和服务；纯本地 getter 则不一定需要全部条件。

\textbf{参数}

无显式参数。实例方法中的 \passthrough{\lstinline!self!} 由 Python
自动传入。

\textbf{返回值}

\begin{longtable}[]{@{}
  >{\raggedright\arraybackslash}p{(\columnwidth - 4\tabcolsep) * \real{0.18}}
  >{\raggedright\arraybackslash}p{(\columnwidth - 4\tabcolsep) * \real{0.20}}
  >{\raggedright\arraybackslash}p{(\columnwidth - 4\tabcolsep) * \real{0.62}}@{}}
\toprule
位置 & 类型 & 含义 \\
\midrule
\endhead
{[}0{]} \passthrough{\lstinline!status!} & \passthrough{\lstinline!int!}
& SDK 状态码；Unitree 示例通常以 \passthrough{\lstinline!0!}
表示成功，非零值需按具体服务错误码解释。 \\
{[}1{]} \passthrough{\lstinline!fsm\_mode!} &
\passthrough{\lstinline!int!} & 传给该接口的 FSM 模式 参数，Python
类型为
\passthrough{\lstinline!int!}。精确含义、单位、范围和枚举值需查目标型号对应头文件与协议。
对应 C++ 参数 \passthrough{\lstinline!fsm\_mode: int \&!}。 \\
\bottomrule
\end{longtable}

\textbf{对应 C++}

\begin{itemize}
\tightlist
\item
  类：\passthrough{\lstinline!unitree::robot::r1::LocoClient!}
\item
  签名：\passthrough{\lstinline!GetFsmMode(int \&)!}
\item
  绑定策略：\passthrough{\lstinline!OUTPUT\_WRAPPER!}
\item
  声明位置：\passthrough{\lstinline!include/unitree/robot/r1/loco/r1\_loco\_client.hpp:43!}
\end{itemize}

\textbf{说明}

\begin{itemize}
\tightlist
\item
  示例是局部调用片段；\passthrough{\lstinline!obj!}
  和参数变量表示已经按上层业务校验并准备好的对象。
\item
  C++ 的可修改输出引用不会作为 Python 入参出现，而是按 C++
  参数顺序追加到返回元组中。
\end{itemize}

\textbf{用法}

\begin{lstlisting}[language=Python]
status, fsm_mode = obj.get_fsm_mode()
\end{lstlisting}

\begin{quote}
{[}!NOTE{]} 只读查询仍会访问
DDS/机器人服务。必须检查返回状态码，并处理超时和网络中断。
\end{quote}

\hypertarget{unitree-sdk2-cpp-robot-r1-lococlient-set-fsm-id-1}{%
\paragraph{\texorpdfstring{\texttt{unitree\_sdk2\_cpp.robot.r1.LocoClient.set\_fsm\_id}}{unitree\_sdk2\_cpp.robot.r1.LocoClient.set\_fsm\_id}}\label{unitree-sdk2-cpp-robot-r1-lococlient-set-fsm-id-1}}

设置 FSM ID。具体副作用和安全边界见下方状态。

\textbf{签名}

\begin{lstlisting}[language=Python]
def set_fsm_id(self, fsm_id: int) -> int
\end{lstlisting}

\textbf{可用性与安全性}

\textbf{\passthrough{\lstinline!SIGNATURE\_ONLY!} /
\passthrough{\lstinline!MOTION\_COMMAND!}}：仅用于补全和静态检查。当前二进制没有该
Python 方法；不得按可执行运动接口使用。

\textbf{参数}

\begin{longtable}[]{@{}
  >{\raggedright\arraybackslash}p{(\columnwidth - 6\tabcolsep) * \real{0.18}}
  >{\raggedright\arraybackslash}p{(\columnwidth - 6\tabcolsep) * \real{0.24}}
  >{\raggedright\arraybackslash}p{(\columnwidth - 6\tabcolsep) * \real{0.18}}
  >{\raggedright\arraybackslash}p{(\columnwidth - 6\tabcolsep) * \real{0.40}}@{}}
\toprule
名称 & Python 类型 & 默认值 & 含义 \\
\midrule
\endhead
\passthrough{\lstinline!fsm\_id!} & \passthrough{\lstinline!int!} & 必填
& 传给该接口的 FSM ID 参数，Python 类型为
\passthrough{\lstinline!int!}。精确含义、单位、范围和枚举值需查目标型号对应头文件与协议。
对应 C++ 参数 \passthrough{\lstinline!fsm\_id: int!}。 \\
\bottomrule
\end{longtable}

\textbf{返回值}

\begin{longtable}[]{@{}
  >{\raggedright\arraybackslash}p{(\columnwidth - 4\tabcolsep) * \real{0.18}}
  >{\raggedright\arraybackslash}p{(\columnwidth - 4\tabcolsep) * \real{0.20}}
  >{\raggedright\arraybackslash}p{(\columnwidth - 4\tabcolsep) * \real{0.62}}@{}}
\toprule
位置 & 类型 & 含义 \\
\midrule
\endhead
返回值 & \passthrough{\lstinline!int!} & SDK 状态码；Unitree 示例通常以
\passthrough{\lstinline!0!} 表示成功，非零值需按具体服务定义解释。 \\
\bottomrule
\end{longtable}

\textbf{对应 C++}

\begin{itemize}
\tightlist
\item
  类：\passthrough{\lstinline!unitree::robot::r1::LocoClient!}
\item
  签名：\passthrough{\lstinline!SetFsmId(int)!}
\item
  绑定策略：\passthrough{\lstinline!DIRECT!}
\item
  声明位置：\passthrough{\lstinline!include/unitree/robot/r1/loco/r1\_loco\_client.hpp:58!}
\end{itemize}

\textbf{说明}

\begin{itemize}
\tightlist
\item
  示例是局部调用片段；\passthrough{\lstinline!obj!}
  和参数变量表示已经按上层业务校验并准备好的对象。
\item
  该条目可以被编辑器和类型检查器识别，但当前运行时可能在属性查找阶段直接失败。
\item
  该方法属于运动命令。方法名包含
  \passthrough{\lstinline!stop!}、\passthrough{\lstinline!damp!} 或
  \passthrough{\lstinline!zero!} 也不代表它是独立物理急停。
\end{itemize}

\textbf{用法}

\begin{lstlisting}[language=Python]
from typing import TYPE_CHECKING

if TYPE_CHECKING:
    def planned_set_fsm_id(obj: LocoClient, fsm_id: int) -> int:
        return obj.set_fsm_id(fsm_id=fsm_id)
\end{lstlisting}

\begin{quote}
{[}!CAUTION{]} 上面的 \passthrough{\lstinline!TYPE\_CHECKING!}
示例只表达计划调用形态。该分支在普通 Python
运行时为假，不代表当前扩展能执行此接口。
\end{quote}

\hypertarget{unitree-sdk2-cpp-robot-r1-lococlient-set-velocity-1}{%
\paragraph{\texorpdfstring{\texttt{unitree\_sdk2\_cpp.robot.r1.LocoClient.set\_velocity}}{unitree\_sdk2\_cpp.robot.r1.LocoClient.set\_velocity}}\label{unitree-sdk2-cpp-robot-r1-lococlient-set-velocity-1}}

设置 速度。具体副作用和安全边界见下方状态。

\textbf{签名}

\begin{lstlisting}[language=Python]
def set_velocity(self, vx: float, vy: float, omega: float, duration: float = ...) -> int
\end{lstlisting}

\textbf{可用性与安全性}

\textbf{\passthrough{\lstinline!SIGNATURE\_ONLY!} /
\passthrough{\lstinline!MOTION\_COMMAND!}}：仅用于补全和静态检查。当前二进制没有该
Python 方法；不得按可执行运动接口使用。

\textbf{参数}

\begin{longtable}[]{@{}
  >{\raggedright\arraybackslash}p{(\columnwidth - 6\tabcolsep) * \real{0.18}}
  >{\raggedright\arraybackslash}p{(\columnwidth - 6\tabcolsep) * \real{0.24}}
  >{\raggedright\arraybackslash}p{(\columnwidth - 6\tabcolsep) * \real{0.18}}
  >{\raggedright\arraybackslash}p{(\columnwidth - 6\tabcolsep) * \real{0.40}}@{}}
\toprule
名称 & Python 类型 & 默认值 & 含义 \\
\midrule
\endhead
\passthrough{\lstinline!vx!} & \passthrough{\lstinline!float!} & 必填 &
X 方向速度参数。坐标系、单位、符号和安全范围必须查目标型号运动协议。
对应 C++ 参数 \passthrough{\lstinline!vx: float!}。
底层为浮点数；类型本身不说明单位、坐标系或业务安全范围。 \\
\passthrough{\lstinline!vy!} & \passthrough{\lstinline!float!} & 必填 &
Y 方向速度参数。坐标系、单位、符号和安全范围必须查目标型号运动协议。
对应 C++ 参数 \passthrough{\lstinline!vy: float!}。
底层为浮点数；类型本身不说明单位、坐标系或业务安全范围。 \\
\passthrough{\lstinline!omega!} & \passthrough{\lstinline!float!} & 必填
& 角速度参数。旋转轴、单位、符号和安全范围必须查目标型号运动协议。 对应
C++ 参数 \passthrough{\lstinline!omega: float!}。
底层为浮点数；类型本身不说明单位、坐标系或业务安全范围。 \\
\passthrough{\lstinline!duration!} & \passthrough{\lstinline!float!} &
\passthrough{\lstinline!...!} &
操作持续时间参数。精确单位、范围和默认行为以目标型号接口定义为准。 对应
C++ 参数 \passthrough{\lstinline!duration: float!}。
底层为浮点数；类型本身不说明单位、坐标系或业务安全范围。 \\
\bottomrule
\end{longtable}

\textbf{返回值}

\begin{longtable}[]{@{}
  >{\raggedright\arraybackslash}p{(\columnwidth - 4\tabcolsep) * \real{0.18}}
  >{\raggedright\arraybackslash}p{(\columnwidth - 4\tabcolsep) * \real{0.20}}
  >{\raggedright\arraybackslash}p{(\columnwidth - 4\tabcolsep) * \real{0.62}}@{}}
\toprule
位置 & 类型 & 含义 \\
\midrule
\endhead
返回值 & \passthrough{\lstinline!int!} & SDK 状态码；Unitree 示例通常以
\passthrough{\lstinline!0!} 表示成功，非零值需按具体服务定义解释。 \\
\bottomrule
\end{longtable}

\textbf{对应 C++}

\begin{itemize}
\tightlist
\item
  类：\passthrough{\lstinline!unitree::robot::r1::LocoClient!}
\item
  签名：\passthrough{\lstinline!SetVelocity(float, float, float, float)!}
\item
  绑定策略：\passthrough{\lstinline!DIRECT!}
\item
  声明位置：\passthrough{\lstinline!include/unitree/robot/r1/loco/r1\_loco\_client.hpp:69!}
\end{itemize}

\textbf{说明}

\begin{itemize}
\tightlist
\item
  示例是局部调用片段；\passthrough{\lstinline!obj!}
  和参数变量表示已经按上层业务校验并准备好的对象。
\item
  默认值显示为 \passthrough{\lstinline!...!}，表示 C++
  声明存在默认参数，但当前 AST
  清单没有保存其字面量；省略参数可使用上游默认行为。
\item
  该条目可以被编辑器和类型检查器识别，但当前运行时可能在属性查找阶段直接失败。
\item
  该方法属于运动命令。方法名包含
  \passthrough{\lstinline!stop!}、\passthrough{\lstinline!damp!} 或
  \passthrough{\lstinline!zero!} 也不代表它是独立物理急停。
\end{itemize}

\textbf{用法}

\begin{lstlisting}[language=Python]
from typing import TYPE_CHECKING

if TYPE_CHECKING:
    def planned_set_velocity(obj: LocoClient, vx: float, vy: float, omega: float, duration: float) -> int:
        return obj.set_velocity(vx=vx, vy=vy, omega=omega, duration=duration)
\end{lstlisting}

\begin{quote}
{[}!CAUTION{]} 上面的 \passthrough{\lstinline!TYPE\_CHECKING!}
示例只表达计划调用形态。该分支在普通 Python
运行时为假，不代表当前扩展能执行此接口。
\end{quote}

\hypertarget{unitree-sdk2-cpp-robot-r1-lococlient-damp-1}{%
\paragraph{\texorpdfstring{\texttt{unitree\_sdk2\_cpp.robot.r1.LocoClient.damp}}{unitree\_sdk2\_cpp.robot.r1.LocoClient.damp}}\label{unitree-sdk2-cpp-robot-r1-lococlient-damp-1}}

对应 C++ SDK 操作
\passthrough{\lstinline!Damp()!}。上游头文件没有可直接生成的业务说明时，本参考只保证签名映射，精确语义需查目标型号协议。

\textbf{签名}

\begin{lstlisting}[language=Python]
def damp(self) -> int
\end{lstlisting}

\textbf{可用性与安全性}

\textbf{\passthrough{\lstinline!SIGNATURE\_ONLY!} /
\passthrough{\lstinline!MOTION\_COMMAND!}}：仅用于补全和静态检查。当前二进制没有该
Python 方法；不得按可执行运动接口使用。

\textbf{参数}

无显式参数。实例方法中的 \passthrough{\lstinline!self!} 由 Python
自动传入。

\textbf{返回值}

\begin{longtable}[]{@{}
  >{\raggedright\arraybackslash}p{(\columnwidth - 4\tabcolsep) * \real{0.18}}
  >{\raggedright\arraybackslash}p{(\columnwidth - 4\tabcolsep) * \real{0.20}}
  >{\raggedright\arraybackslash}p{(\columnwidth - 4\tabcolsep) * \real{0.62}}@{}}
\toprule
位置 & 类型 & 含义 \\
\midrule
\endhead
返回值 & \passthrough{\lstinline!int!} & SDK 状态码；Unitree 示例通常以
\passthrough{\lstinline!0!} 表示成功，非零值需按具体服务定义解释。 \\
\bottomrule
\end{longtable}

\textbf{对应 C++}

\begin{itemize}
\tightlist
\item
  类：\passthrough{\lstinline!unitree::robot::r1::LocoClient!}
\item
  签名：\passthrough{\lstinline!Damp()!}
\item
  绑定策略：\passthrough{\lstinline!DIRECT!}
\item
  声明位置：\passthrough{\lstinline!include/unitree/robot/r1/loco/r1\_loco\_client.hpp:82!}
\end{itemize}

\textbf{说明}

\begin{itemize}
\tightlist
\item
  示例是局部调用片段；\passthrough{\lstinline!obj!}
  和参数变量表示已经按上层业务校验并准备好的对象。
\item
  该条目可以被编辑器和类型检查器识别，但当前运行时可能在属性查找阶段直接失败。
\item
  该方法属于运动命令。方法名包含
  \passthrough{\lstinline!stop!}、\passthrough{\lstinline!damp!} 或
  \passthrough{\lstinline!zero!} 也不代表它是独立物理急停。
\end{itemize}

\textbf{用法}

\begin{lstlisting}[language=Python]
from typing import TYPE_CHECKING

if TYPE_CHECKING:
    def planned_damp(obj: LocoClient) -> int:
        return obj.damp()
\end{lstlisting}

\begin{quote}
{[}!CAUTION{]} 上面的 \passthrough{\lstinline!TYPE\_CHECKING!}
示例只表达计划调用形态。该分支在普通 Python
运行时为假，不代表当前扩展能执行此接口。
\end{quote}

\hypertarget{unitree-sdk2-cpp-robot-r1-lococlient-start-1}{%
\paragraph{\texorpdfstring{\texttt{unitree\_sdk2\_cpp.robot.r1.LocoClient.start}}{unitree\_sdk2\_cpp.robot.r1.LocoClient.start}}\label{unitree-sdk2-cpp-robot-r1-lococlient-start-1}}

启动对应 SDK 操作。具体副作用和安全边界见下方状态。

\textbf{签名}

\begin{lstlisting}[language=Python]
def start(self) -> int
\end{lstlisting}

\textbf{可用性与安全性}

\textbf{\passthrough{\lstinline!SIGNATURE\_ONLY!} /
\passthrough{\lstinline!MOTION\_COMMAND!}}：仅用于补全和静态检查。当前二进制没有该
Python 方法；不得按可执行运动接口使用。

\textbf{参数}

无显式参数。实例方法中的 \passthrough{\lstinline!self!} 由 Python
自动传入。

\textbf{返回值}

\begin{longtable}[]{@{}
  >{\raggedright\arraybackslash}p{(\columnwidth - 4\tabcolsep) * \real{0.18}}
  >{\raggedright\arraybackslash}p{(\columnwidth - 4\tabcolsep) * \real{0.20}}
  >{\raggedright\arraybackslash}p{(\columnwidth - 4\tabcolsep) * \real{0.62}}@{}}
\toprule
位置 & 类型 & 含义 \\
\midrule
\endhead
返回值 & \passthrough{\lstinline!int!} & SDK 状态码；Unitree 示例通常以
\passthrough{\lstinline!0!} 表示成功，非零值需按具体服务定义解释。 \\
\bottomrule
\end{longtable}

\textbf{对应 C++}

\begin{itemize}
\tightlist
\item
  类：\passthrough{\lstinline!unitree::robot::r1::LocoClient!}
\item
  签名：\passthrough{\lstinline!Start()!}
\item
  绑定策略：\passthrough{\lstinline!DIRECT!}
\item
  声明位置：\passthrough{\lstinline!include/unitree/robot/r1/loco/r1\_loco\_client.hpp:84!}
\end{itemize}

\textbf{说明}

\begin{itemize}
\tightlist
\item
  示例是局部调用片段；\passthrough{\lstinline!obj!}
  和参数变量表示已经按上层业务校验并准备好的对象。
\item
  该条目可以被编辑器和类型检查器识别，但当前运行时可能在属性查找阶段直接失败。
\item
  该方法属于运动命令。方法名包含
  \passthrough{\lstinline!stop!}、\passthrough{\lstinline!damp!} 或
  \passthrough{\lstinline!zero!} 也不代表它是独立物理急停。
\end{itemize}

\textbf{用法}

\begin{lstlisting}[language=Python]
from typing import TYPE_CHECKING

if TYPE_CHECKING:
    def planned_start(obj: LocoClient) -> int:
        return obj.start()
\end{lstlisting}

\begin{quote}
{[}!CAUTION{]} 上面的 \passthrough{\lstinline!TYPE\_CHECKING!}
示例只表达计划调用形态。该分支在普通 Python
运行时为假，不代表当前扩展能执行此接口。
\end{quote}

\hypertarget{unitree-sdk2-cpp-robot-r1-lococlient-stand-up-1}{%
\paragraph{\texorpdfstring{\texttt{unitree\_sdk2\_cpp.robot.r1.LocoClient.stand\_up}}{unitree\_sdk2\_cpp.robot.r1.LocoClient.stand\_up}}\label{unitree-sdk2-cpp-robot-r1-lococlient-stand-up-1}}

对应 C++ SDK 操作
\passthrough{\lstinline!StandUp()!}。上游头文件没有可直接生成的业务说明时，本参考只保证签名映射，精确语义需查目标型号协议。

\textbf{签名}

\begin{lstlisting}[language=Python]
def stand_up(self) -> int
\end{lstlisting}

\textbf{可用性与安全性}

\textbf{\passthrough{\lstinline!SIGNATURE\_ONLY!} /
\passthrough{\lstinline!MOTION\_COMMAND!}}：仅用于补全和静态检查。当前二进制没有该
Python 方法；不得按可执行运动接口使用。

\textbf{参数}

无显式参数。实例方法中的 \passthrough{\lstinline!self!} 由 Python
自动传入。

\textbf{返回值}

\begin{longtable}[]{@{}
  >{\raggedright\arraybackslash}p{(\columnwidth - 4\tabcolsep) * \real{0.18}}
  >{\raggedright\arraybackslash}p{(\columnwidth - 4\tabcolsep) * \real{0.20}}
  >{\raggedright\arraybackslash}p{(\columnwidth - 4\tabcolsep) * \real{0.62}}@{}}
\toprule
位置 & 类型 & 含义 \\
\midrule
\endhead
返回值 & \passthrough{\lstinline!int!} & SDK 状态码；Unitree 示例通常以
\passthrough{\lstinline!0!} 表示成功，非零值需按具体服务定义解释。 \\
\bottomrule
\end{longtable}

\textbf{对应 C++}

\begin{itemize}
\tightlist
\item
  类：\passthrough{\lstinline!unitree::robot::r1::LocoClient!}
\item
  签名：\passthrough{\lstinline!StandUp()!}
\item
  绑定策略：\passthrough{\lstinline!DIRECT!}
\item
  声明位置：\passthrough{\lstinline!include/unitree/robot/r1/loco/r1\_loco\_client.hpp:90!}
\end{itemize}

\textbf{说明}

\begin{itemize}
\tightlist
\item
  示例是局部调用片段；\passthrough{\lstinline!obj!}
  和参数变量表示已经按上层业务校验并准备好的对象。
\item
  该条目可以被编辑器和类型检查器识别，但当前运行时可能在属性查找阶段直接失败。
\item
  该方法属于运动命令。方法名包含
  \passthrough{\lstinline!stop!}、\passthrough{\lstinline!damp!} 或
  \passthrough{\lstinline!zero!} 也不代表它是独立物理急停。
\end{itemize}

\textbf{用法}

\begin{lstlisting}[language=Python]
from typing import TYPE_CHECKING

if TYPE_CHECKING:
    def planned_stand_up(obj: LocoClient) -> int:
        return obj.stand_up()
\end{lstlisting}

\begin{quote}
{[}!CAUTION{]} 上面的 \passthrough{\lstinline!TYPE\_CHECKING!}
示例只表达计划调用形态。该分支在普通 Python
运行时为假，不代表当前扩展能执行此接口。
\end{quote}

\hypertarget{unitree-sdk2-cpp-robot-r1-lococlient-zero-torque-1}{%
\paragraph{\texorpdfstring{\texttt{unitree\_sdk2\_cpp.robot.r1.LocoClient.zero\_torque}}{unitree\_sdk2\_cpp.robot.r1.LocoClient.zero\_torque}}\label{unitree-sdk2-cpp-robot-r1-lococlient-zero-torque-1}}

对应 C++ SDK 操作
\passthrough{\lstinline!ZeroTorque()!}。上游头文件没有可直接生成的业务说明时，本参考只保证签名映射，精确语义需查目标型号协议。

\textbf{签名}

\begin{lstlisting}[language=Python]
def zero_torque(self) -> int
\end{lstlisting}

\textbf{可用性与安全性}

\textbf{\passthrough{\lstinline!SIGNATURE\_ONLY!} /
\passthrough{\lstinline!MOTION\_COMMAND!}}：仅用于补全和静态检查。当前二进制没有该
Python 方法；不得按可执行运动接口使用。

\textbf{参数}

无显式参数。实例方法中的 \passthrough{\lstinline!self!} 由 Python
自动传入。

\textbf{返回值}

\begin{longtable}[]{@{}
  >{\raggedright\arraybackslash}p{(\columnwidth - 4\tabcolsep) * \real{0.18}}
  >{\raggedright\arraybackslash}p{(\columnwidth - 4\tabcolsep) * \real{0.20}}
  >{\raggedright\arraybackslash}p{(\columnwidth - 4\tabcolsep) * \real{0.62}}@{}}
\toprule
位置 & 类型 & 含义 \\
\midrule
\endhead
返回值 & \passthrough{\lstinline!int!} & SDK 状态码；Unitree 示例通常以
\passthrough{\lstinline!0!} 表示成功，非零值需按具体服务定义解释。 \\
\bottomrule
\end{longtable}

\textbf{对应 C++}

\begin{itemize}
\tightlist
\item
  类：\passthrough{\lstinline!unitree::robot::r1::LocoClient!}
\item
  签名：\passthrough{\lstinline!ZeroTorque()!}
\item
  绑定策略：\passthrough{\lstinline!DIRECT!}
\item
  声明位置：\passthrough{\lstinline!include/unitree/robot/r1/loco/r1\_loco\_client.hpp:92!}
\end{itemize}

\textbf{说明}

\begin{itemize}
\tightlist
\item
  示例是局部调用片段；\passthrough{\lstinline!obj!}
  和参数变量表示已经按上层业务校验并准备好的对象。
\item
  该条目可以被编辑器和类型检查器识别，但当前运行时可能在属性查找阶段直接失败。
\item
  该方法属于运动命令。方法名包含
  \passthrough{\lstinline!stop!}、\passthrough{\lstinline!damp!} 或
  \passthrough{\lstinline!zero!} 也不代表它是独立物理急停。
\end{itemize}

\textbf{用法}

\begin{lstlisting}[language=Python]
from typing import TYPE_CHECKING

if TYPE_CHECKING:
    def planned_zero_torque(obj: LocoClient) -> int:
        return obj.zero_torque()
\end{lstlisting}

\begin{quote}
{[}!CAUTION{]} 上面的 \passthrough{\lstinline!TYPE\_CHECKING!}
示例只表达计划调用形态。该分支在普通 Python
运行时为假，不代表当前扩展能执行此接口。
\end{quote}

\hypertarget{unitree-sdk2-cpp-robot-r1-lococlient-stop-move-1}{%
\paragraph{\texorpdfstring{\texttt{unitree\_sdk2\_cpp.robot.r1.LocoClient.stop\_move}}{unitree\_sdk2\_cpp.robot.r1.LocoClient.stop\_move}}\label{unitree-sdk2-cpp-robot-r1-lococlient-stop-move-1}}

请求停止对应 SDK 操作。该名称不等同于经过验证的物理急停。

\textbf{签名}

\begin{lstlisting}[language=Python]
def stop_move(self) -> int
\end{lstlisting}

\textbf{可用性与安全性}

\textbf{\passthrough{\lstinline!SIGNATURE\_ONLY!} /
\passthrough{\lstinline!MOTION\_COMMAND!}}：仅用于补全和静态检查。当前二进制没有该
Python 方法；不得按可执行运动接口使用。

\textbf{参数}

无显式参数。实例方法中的 \passthrough{\lstinline!self!} 由 Python
自动传入。

\textbf{返回值}

\begin{longtable}[]{@{}
  >{\raggedright\arraybackslash}p{(\columnwidth - 4\tabcolsep) * \real{0.18}}
  >{\raggedright\arraybackslash}p{(\columnwidth - 4\tabcolsep) * \real{0.20}}
  >{\raggedright\arraybackslash}p{(\columnwidth - 4\tabcolsep) * \real{0.62}}@{}}
\toprule
位置 & 类型 & 含义 \\
\midrule
\endhead
返回值 & \passthrough{\lstinline!int!} & SDK 状态码；Unitree 示例通常以
\passthrough{\lstinline!0!} 表示成功，非零值需按具体服务定义解释。 \\
\bottomrule
\end{longtable}

\textbf{对应 C++}

\begin{itemize}
\tightlist
\item
  类：\passthrough{\lstinline!unitree::robot::r1::LocoClient!}
\item
  签名：\passthrough{\lstinline!StopMove()!}
\item
  绑定策略：\passthrough{\lstinline!DIRECT!}
\item
  声明位置：\passthrough{\lstinline!include/unitree/robot/r1/loco/r1\_loco\_client.hpp:94!}
\end{itemize}

\textbf{说明}

\begin{itemize}
\tightlist
\item
  示例是局部调用片段；\passthrough{\lstinline!obj!}
  和参数变量表示已经按上层业务校验并准备好的对象。
\item
  该条目可以被编辑器和类型检查器识别，但当前运行时可能在属性查找阶段直接失败。
\item
  该方法属于运动命令。方法名包含
  \passthrough{\lstinline!stop!}、\passthrough{\lstinline!damp!} 或
  \passthrough{\lstinline!zero!} 也不代表它是独立物理急停。
\end{itemize}

\textbf{用法}

\begin{lstlisting}[language=Python]
from typing import TYPE_CHECKING

if TYPE_CHECKING:
    def planned_stop_move(obj: LocoClient) -> int:
        return obj.stop_move()
\end{lstlisting}

\begin{quote}
{[}!CAUTION{]} 上面的 \passthrough{\lstinline!TYPE\_CHECKING!}
示例只表达计划调用形态。该分支在普通 Python
运行时为假，不代表当前扩展能执行此接口。
\end{quote}

\hypertarget{unitree-sdk2-cpp-robot-r1-lococlient-move-1}{%
\paragraph{\texorpdfstring{\texttt{unitree\_sdk2\_cpp.robot.r1.LocoClient.move}（重载
1/2）}{unitree\_sdk2\_cpp.robot.r1.LocoClient.move（重载 1/2）}}\label{unitree-sdk2-cpp-robot-r1-lococlient-move-1}}

对应 C++ SDK 操作
\passthrough{\lstinline!Move(float, float, float, bool)!}。上游头文件没有可直接生成的业务说明时，本参考只保证签名映射，精确语义需查目标型号协议。

\textbf{签名}

\begin{lstlisting}[language=Python]
@overload
def move(self, vx: float, vy: float, vyaw: float, continous_move: bool) -> int
\end{lstlisting}

\textbf{可用性与安全性}

\textbf{\passthrough{\lstinline!SIGNATURE\_ONLY!} /
\passthrough{\lstinline!MOTION\_COMMAND!}}：仅用于补全和静态检查。当前二进制没有该
Python 方法；不得按可执行运动接口使用。

\textbf{参数}

\begin{longtable}[]{@{}
  >{\raggedright\arraybackslash}p{(\columnwidth - 6\tabcolsep) * \real{0.18}}
  >{\raggedright\arraybackslash}p{(\columnwidth - 6\tabcolsep) * \real{0.24}}
  >{\raggedright\arraybackslash}p{(\columnwidth - 6\tabcolsep) * \real{0.18}}
  >{\raggedright\arraybackslash}p{(\columnwidth - 6\tabcolsep) * \real{0.40}}@{}}
\toprule
名称 & Python 类型 & 默认值 & 含义 \\
\midrule
\endhead
\passthrough{\lstinline!vx!} & \passthrough{\lstinline!float!} & 必填 &
X 方向速度参数。坐标系、单位、符号和安全范围必须查目标型号运动协议。
对应 C++ 参数 \passthrough{\lstinline!vx: float!}。
底层为浮点数；类型本身不说明单位、坐标系或业务安全范围。 \\
\passthrough{\lstinline!vy!} & \passthrough{\lstinline!float!} & 必填 &
Y 方向速度参数。坐标系、单位、符号和安全范围必须查目标型号运动协议。
对应 C++ 参数 \passthrough{\lstinline!vy: float!}。
底层为浮点数；类型本身不说明单位、坐标系或业务安全范围。 \\
\passthrough{\lstinline!vyaw!} & \passthrough{\lstinline!float!} & 必填
& 偏航角速度参数。单位、符号和安全范围必须查目标型号运动协议。 对应 C++
参数 \passthrough{\lstinline!vyaw: float!}。
底层为浮点数；类型本身不说明单位、坐标系或业务安全范围。 \\
\passthrough{\lstinline!continous\_move!} &
\passthrough{\lstinline!bool!} & 必填 & 传给该接口的
\passthrough{\lstinline!continous!} \passthrough{\lstinline!move!}
参数，Python 类型为
\passthrough{\lstinline!bool!}。精确含义、单位、范围和枚举值需查目标型号对应头文件与协议。
对应 C++ 参数 \passthrough{\lstinline!continous\_move: bool!}。
只接受布尔语义。 \\
\bottomrule
\end{longtable}

\textbf{返回值}

\begin{longtable}[]{@{}
  >{\raggedright\arraybackslash}p{(\columnwidth - 4\tabcolsep) * \real{0.18}}
  >{\raggedright\arraybackslash}p{(\columnwidth - 4\tabcolsep) * \real{0.20}}
  >{\raggedright\arraybackslash}p{(\columnwidth - 4\tabcolsep) * \real{0.62}}@{}}
\toprule
位置 & 类型 & 含义 \\
\midrule
\endhead
返回值 & \passthrough{\lstinline!int!} & SDK 状态码；Unitree 示例通常以
\passthrough{\lstinline!0!} 表示成功，非零值需按具体服务定义解释。 \\
\bottomrule
\end{longtable}

\textbf{对应 C++}

\begin{itemize}
\tightlist
\item
  类：\passthrough{\lstinline!unitree::robot::r1::LocoClient!}
\item
  签名：\passthrough{\lstinline!Move(float, float, float, bool)!}
\item
  绑定策略：\passthrough{\lstinline!DIRECT!}
\item
  声明位置：\passthrough{\lstinline!include/unitree/robot/r1/loco/r1\_loco\_client.hpp:96!}
\end{itemize}

\textbf{说明}

\begin{itemize}
\tightlist
\item
  示例是局部调用片段；\passthrough{\lstinline!obj!}
  和参数变量表示已经按上层业务校验并准备好的对象。
\item
  该条目可以被编辑器和类型检查器识别，但当前运行时可能在属性查找阶段直接失败。
\item
  该方法属于运动命令。方法名包含
  \passthrough{\lstinline!stop!}、\passthrough{\lstinline!damp!} 或
  \passthrough{\lstinline!zero!} 也不代表它是独立物理急停。
\end{itemize}

\textbf{用法}

\begin{lstlisting}[language=Python]
from typing import TYPE_CHECKING

if TYPE_CHECKING:
    def planned_move(obj: LocoClient, vx: float, vy: float, vyaw: float, continous_move: bool) -> int:
        return obj.move(vx=vx, vy=vy, vyaw=vyaw, continous_move=continous_move)
\end{lstlisting}

\begin{quote}
{[}!CAUTION{]} 上面的 \passthrough{\lstinline!TYPE\_CHECKING!}
示例只表达计划调用形态。该分支在普通 Python
运行时为假，不代表当前扩展能执行此接口。
\end{quote}

\hypertarget{unitree-sdk2-cpp-robot-r1-lococlient-move-2}{%
\paragraph{\texorpdfstring{\texttt{unitree\_sdk2\_cpp.robot.r1.LocoClient.move}（重载
2/2）}{unitree\_sdk2\_cpp.robot.r1.LocoClient.move（重载 2/2）}}\label{unitree-sdk2-cpp-robot-r1-lococlient-move-2}}

对应 C++ SDK 操作
\passthrough{\lstinline!Move(float, float, float)!}。上游头文件没有可直接生成的业务说明时，本参考只保证签名映射，精确语义需查目标型号协议。

\textbf{签名}

\begin{lstlisting}[language=Python]
@overload
def move(self, vx: float, vy: float, vyaw: float) -> int
\end{lstlisting}

\textbf{可用性与安全性}

\textbf{\passthrough{\lstinline!SIGNATURE\_ONLY!} /
\passthrough{\lstinline!MOTION\_COMMAND!}}：仅用于补全和静态检查。当前二进制没有该
Python 方法；不得按可执行运动接口使用。

\textbf{参数}

\begin{longtable}[]{@{}
  >{\raggedright\arraybackslash}p{(\columnwidth - 6\tabcolsep) * \real{0.18}}
  >{\raggedright\arraybackslash}p{(\columnwidth - 6\tabcolsep) * \real{0.24}}
  >{\raggedright\arraybackslash}p{(\columnwidth - 6\tabcolsep) * \real{0.18}}
  >{\raggedright\arraybackslash}p{(\columnwidth - 6\tabcolsep) * \real{0.40}}@{}}
\toprule
名称 & Python 类型 & 默认值 & 含义 \\
\midrule
\endhead
\passthrough{\lstinline!vx!} & \passthrough{\lstinline!float!} & 必填 &
X 方向速度参数。坐标系、单位、符号和安全范围必须查目标型号运动协议。
对应 C++ 参数 \passthrough{\lstinline!vx: float!}。
底层为浮点数；类型本身不说明单位、坐标系或业务安全范围。 \\
\passthrough{\lstinline!vy!} & \passthrough{\lstinline!float!} & 必填 &
Y 方向速度参数。坐标系、单位、符号和安全范围必须查目标型号运动协议。
对应 C++ 参数 \passthrough{\lstinline!vy: float!}。
底层为浮点数；类型本身不说明单位、坐标系或业务安全范围。 \\
\passthrough{\lstinline!vyaw!} & \passthrough{\lstinline!float!} & 必填
& 偏航角速度参数。单位、符号和安全范围必须查目标型号运动协议。 对应 C++
参数 \passthrough{\lstinline!vyaw: float!}。
底层为浮点数；类型本身不说明单位、坐标系或业务安全范围。 \\
\bottomrule
\end{longtable}

\textbf{返回值}

\begin{longtable}[]{@{}
  >{\raggedright\arraybackslash}p{(\columnwidth - 4\tabcolsep) * \real{0.18}}
  >{\raggedright\arraybackslash}p{(\columnwidth - 4\tabcolsep) * \real{0.20}}
  >{\raggedright\arraybackslash}p{(\columnwidth - 4\tabcolsep) * \real{0.62}}@{}}
\toprule
位置 & 类型 & 含义 \\
\midrule
\endhead
返回值 & \passthrough{\lstinline!int!} & SDK 状态码；Unitree 示例通常以
\passthrough{\lstinline!0!} 表示成功，非零值需按具体服务定义解释。 \\
\bottomrule
\end{longtable}

\textbf{对应 C++}

\begin{itemize}
\tightlist
\item
  类：\passthrough{\lstinline!unitree::robot::r1::LocoClient!}
\item
  签名：\passthrough{\lstinline!Move(float, float, float)!}
\item
  绑定策略：\passthrough{\lstinline!DIRECT!}
\item
  声明位置：\passthrough{\lstinline!include/unitree/robot/r1/loco/r1\_loco\_client.hpp:100!}
\end{itemize}

\textbf{说明}

\begin{itemize}
\tightlist
\item
  示例是局部调用片段；\passthrough{\lstinline!obj!}
  和参数变量表示已经按上层业务校验并准备好的对象。
\item
  该条目可以被编辑器和类型检查器识别，但当前运行时可能在属性查找阶段直接失败。
\item
  该方法属于运动命令。方法名包含
  \passthrough{\lstinline!stop!}、\passthrough{\lstinline!damp!} 或
  \passthrough{\lstinline!zero!} 也不代表它是独立物理急停。
\end{itemize}

\textbf{用法}

\begin{lstlisting}[language=Python]
from typing import TYPE_CHECKING

if TYPE_CHECKING:
    def planned_move(obj: LocoClient, vx: float, vy: float, vyaw: float) -> int:
        return obj.move(vx=vx, vy=vy, vyaw=vyaw)
\end{lstlisting}

\begin{quote}
{[}!CAUTION{]} 上面的 \passthrough{\lstinline!TYPE\_CHECKING!}
示例只表达计划调用形态。该分支在普通 Python
运行时为假，不代表当前扩展能执行此接口。
\end{quote}

\hypertarget{unitree-sdk2-cpp-robot-r1-lococlient-switch-move-mode-1}{%
\paragraph{\texorpdfstring{\texttt{unitree\_sdk2\_cpp.robot.r1.LocoClient.switch\_move\_mode}}{unitree\_sdk2\_cpp.robot.r1.LocoClient.switch\_move\_mode}}\label{unitree-sdk2-cpp-robot-r1-lococlient-switch-move-mode-1}}

对应 C++ SDK 操作
\passthrough{\lstinline!SwitchMoveMode(bool)!}。上游头文件没有可直接生成的业务说明时，本参考只保证签名映射，精确语义需查目标型号协议。

\textbf{签名}

\begin{lstlisting}[language=Python]
def switch_move_mode(self, flag: bool) -> int
\end{lstlisting}

\textbf{可用性与安全性}

\textbf{\passthrough{\lstinline!SIGNATURE\_ONLY!} /
\passthrough{\lstinline!MOTION\_COMMAND!}}：仅用于补全和静态检查。当前二进制没有该
Python 方法；不得按可执行运动接口使用。

\textbf{参数}

\begin{longtable}[]{@{}
  >{\raggedright\arraybackslash}p{(\columnwidth - 6\tabcolsep) * \real{0.18}}
  >{\raggedright\arraybackslash}p{(\columnwidth - 6\tabcolsep) * \real{0.24}}
  >{\raggedright\arraybackslash}p{(\columnwidth - 6\tabcolsep) * \real{0.18}}
  >{\raggedright\arraybackslash}p{(\columnwidth - 6\tabcolsep) * \real{0.40}}@{}}
\toprule
名称 & Python 类型 & 默认值 & 含义 \\
\midrule
\endhead
\passthrough{\lstinline!flag!} & \passthrough{\lstinline!bool!} & 必填 &
布尔开关。\passthrough{\lstinline!True!} 和
\passthrough{\lstinline!False!} 的具体业务效果由当前方法定义。 对应 C++
参数 \passthrough{\lstinline!flag: bool!}。 只接受布尔语义。 \\
\bottomrule
\end{longtable}

\textbf{返回值}

\begin{longtable}[]{@{}
  >{\raggedright\arraybackslash}p{(\columnwidth - 4\tabcolsep) * \real{0.18}}
  >{\raggedright\arraybackslash}p{(\columnwidth - 4\tabcolsep) * \real{0.20}}
  >{\raggedright\arraybackslash}p{(\columnwidth - 4\tabcolsep) * \real{0.62}}@{}}
\toprule
位置 & 类型 & 含义 \\
\midrule
\endhead
返回值 & \passthrough{\lstinline!int!} & SDK 状态码；Unitree 示例通常以
\passthrough{\lstinline!0!} 表示成功，非零值需按具体服务定义解释。 \\
\bottomrule
\end{longtable}

\textbf{对应 C++}

\begin{itemize}
\tightlist
\item
  类：\passthrough{\lstinline!unitree::robot::r1::LocoClient!}
\item
  签名：\passthrough{\lstinline!SwitchMoveMode(bool)!}
\item
  绑定策略：\passthrough{\lstinline!DIRECT!}
\item
  声明位置：\passthrough{\lstinline!include/unitree/robot/r1/loco/r1\_loco\_client.hpp:102!}
\end{itemize}

\textbf{说明}

\begin{itemize}
\tightlist
\item
  示例是局部调用片段；\passthrough{\lstinline!obj!}
  和参数变量表示已经按上层业务校验并准备好的对象。
\item
  该条目可以被编辑器和类型检查器识别，但当前运行时可能在属性查找阶段直接失败。
\item
  该方法属于运动命令。方法名包含
  \passthrough{\lstinline!stop!}、\passthrough{\lstinline!damp!} 或
  \passthrough{\lstinline!zero!} 也不代表它是独立物理急停。
\end{itemize}

\textbf{用法}

\begin{lstlisting}[language=Python]
from typing import TYPE_CHECKING

if TYPE_CHECKING:
    def planned_switch_move_mode(obj: LocoClient, flag: bool) -> int:
        return obj.switch_move_mode(flag=flag)
\end{lstlisting}

\begin{quote}
{[}!CAUTION{]} 上面的 \passthrough{\lstinline!TYPE\_CHECKING!}
示例只表达计划调用形态。该分支在普通 Python
运行时为假，不代表当前扩展能执行此接口。
\end{quote}

\hypertarget{unitree-sdk2-cpp-robot-r1-lococlient-set-speed-mode-1}{%
\paragraph{\texorpdfstring{\texttt{unitree\_sdk2\_cpp.robot.r1.LocoClient.set\_speed\_mode}}{unitree\_sdk2\_cpp.robot.r1.LocoClient.set\_speed\_mode}}\label{unitree-sdk2-cpp-robot-r1-lococlient-set-speed-mode-1}}

设置 速度 模式。具体副作用和安全边界见下方状态。

\textbf{签名}

\begin{lstlisting}[language=Python]
def set_speed_mode(self, speed_mode: int) -> int
\end{lstlisting}

\textbf{可用性与安全性}

\textbf{\passthrough{\lstinline!SIGNATURE\_ONLY!} /
\passthrough{\lstinline!MOTION\_COMMAND!}}：仅用于补全和静态检查。当前二进制没有该
Python 方法；不得按可执行运动接口使用。

\textbf{参数}

\begin{longtable}[]{@{}
  >{\raggedright\arraybackslash}p{(\columnwidth - 6\tabcolsep) * \real{0.18}}
  >{\raggedright\arraybackslash}p{(\columnwidth - 6\tabcolsep) * \real{0.24}}
  >{\raggedright\arraybackslash}p{(\columnwidth - 6\tabcolsep) * \real{0.18}}
  >{\raggedright\arraybackslash}p{(\columnwidth - 6\tabcolsep) * \real{0.40}}@{}}
\toprule
名称 & Python 类型 & 默认值 & 含义 \\
\midrule
\endhead
\passthrough{\lstinline!speed\_mode!} & \passthrough{\lstinline!int!} &
必填 & 传给该接口的 速度 模式 参数，Python 类型为
\passthrough{\lstinline!int!}。精确含义、单位、范围和枚举值需查目标型号对应头文件与协议。
对应 C++ 参数 \passthrough{\lstinline!speed\_mode: int!}。 \\
\bottomrule
\end{longtable}

\textbf{返回值}

\begin{longtable}[]{@{}
  >{\raggedright\arraybackslash}p{(\columnwidth - 4\tabcolsep) * \real{0.18}}
  >{\raggedright\arraybackslash}p{(\columnwidth - 4\tabcolsep) * \real{0.20}}
  >{\raggedright\arraybackslash}p{(\columnwidth - 4\tabcolsep) * \real{0.62}}@{}}
\toprule
位置 & 类型 & 含义 \\
\midrule
\endhead
返回值 & \passthrough{\lstinline!int!} & SDK 状态码；Unitree 示例通常以
\passthrough{\lstinline!0!} 表示成功，非零值需按具体服务定义解释。 \\
\bottomrule
\end{longtable}

\textbf{对应 C++}

\begin{itemize}
\tightlist
\item
  类：\passthrough{\lstinline!unitree::robot::r1::LocoClient!}
\item
  签名：\passthrough{\lstinline!SetSpeedMode(int)!}
\item
  绑定策略：\passthrough{\lstinline!DIRECT!}
\item
  声明位置：\passthrough{\lstinline!include/unitree/robot/r1/loco/r1\_loco\_client.hpp:107!}
\end{itemize}

\textbf{说明}

\begin{itemize}
\tightlist
\item
  示例是局部调用片段；\passthrough{\lstinline!obj!}
  和参数变量表示已经按上层业务校验并准备好的对象。
\item
  该条目可以被编辑器和类型检查器识别，但当前运行时可能在属性查找阶段直接失败。
\item
  该方法属于运动命令。方法名包含
  \passthrough{\lstinline!stop!}、\passthrough{\lstinline!damp!} 或
  \passthrough{\lstinline!zero!} 也不代表它是独立物理急停。
\end{itemize}

\textbf{用法}

\begin{lstlisting}[language=Python]
from typing import TYPE_CHECKING

if TYPE_CHECKING:
    def planned_set_speed_mode(obj: LocoClient, speed_mode: int) -> int:
        return obj.set_speed_mode(speed_mode=speed_mode)
\end{lstlisting}

\begin{quote}
{[}!CAUTION{]} 上面的 \passthrough{\lstinline!TYPE\_CHECKING!}
示例只表达计划调用形态。该分支在普通 Python
运行时为假，不代表当前扩展能执行此接口。
\end{quote}

\hypertarget{unitree-sdk2-cpp-robot-r1-playstopparameter}{%
\subsubsection{\texorpdfstring{\texttt{unitree\_sdk2\_cpp.robot.r1.PlayStopParameter}}{unitree\_sdk2\_cpp.robot.r1.PlayStopParameter}}\label{unitree-sdk2-cpp-robot-r1-playstopparameter}}

SDK 请求参数值类型；当前可能仅提供设计期签名。

\textbf{导入}

\begin{lstlisting}[language=Python]
from unitree_sdk2_cpp.robot.r1 import PlayStopParameter
\end{lstlisting}

\textbf{基类}：\passthrough{\lstinline!object!}

\textbf{公开属性}

\begin{longtable}[]{@{}
  >{\raggedright\arraybackslash}p{(\columnwidth - 6\tabcolsep) * \real{0.18}}
  >{\raggedright\arraybackslash}p{(\columnwidth - 6\tabcolsep) * \real{0.24}}
  >{\raggedright\arraybackslash}p{(\columnwidth - 6\tabcolsep) * \real{0.18}}
  >{\raggedright\arraybackslash}p{(\columnwidth - 6\tabcolsep) * \real{0.40}}@{}}
\toprule
名称 & Python 类型 & 含义 & 用法 \\
\midrule
\endhead
\passthrough{\lstinline!app\_name!} & \passthrough{\lstinline!str!} &
应用名称，用于标识音频流或播放会话。允许值和命名规则以目标服务协议为准。
& \passthrough{\lstinline!value = obj.app\_name!} /
\passthrough{\lstinline!obj.app\_name = value!} \\
\bottomrule
\end{longtable}

\textbf{方法索引}

\begin{longtable}[]{@{}
  >{\raggedright\arraybackslash}p{(\columnwidth - 6\tabcolsep) * \real{0.18}}
  >{\raggedright\arraybackslash}p{(\columnwidth - 6\tabcolsep) * \real{0.24}}
  >{\raggedright\arraybackslash}p{(\columnwidth - 6\tabcolsep) * \real{0.18}}
  >{\raggedright\arraybackslash}p{(\columnwidth - 6\tabcolsep) * \real{0.40}}@{}}
\toprule
方法 & Python 签名 & 状态 & 安全分类 \\
\midrule
\endhead
\protect\hyperlink{unitree-sdk2-cpp-robot-r1-playstopparameter-dunder-init-1}{\passthrough{\lstinline!\_\_init\_\_!}}
& \passthrough{\lstinline!def \_\_init\_\_(self) -> None!} &
\passthrough{\lstinline!SIGNATURE\_ONLY!} &
\passthrough{\lstinline!CONSTRUCTION!} \\
\protect\hyperlink{unitree-sdk2-cpp-robot-r1-playstopparameter-from-json-1}{\passthrough{\lstinline!from\_json!}}
&
\passthrough{\lstinline!def from\_json(self, json: dict[str, Any]) -> None!}
& \passthrough{\lstinline!SIGNATURE\_ONLY!} &
\passthrough{\lstinline!UNCLASSIFIED!} \\
\protect\hyperlink{unitree-sdk2-cpp-robot-r1-playstopparameter-to-json-1}{\passthrough{\lstinline!to\_json!}}
&
\passthrough{\lstinline!def to\_json(self, json: dict[str, Any]) -> None!}
& \passthrough{\lstinline!SIGNATURE\_ONLY!} &
\passthrough{\lstinline!UNCLASSIFIED!} \\
\bottomrule
\end{longtable}

\hypertarget{unitree-sdk2-cpp-robot-r1-playstopparameter-dunder-init-1}{%
\paragraph{\texorpdfstring{\texttt{unitree\_sdk2\_cpp.robot.r1.PlayStopParameter.\_\_init\_\_}}{unitree\_sdk2\_cpp.robot.r1.PlayStopParameter.\_\_init\_\_}}\label{unitree-sdk2-cpp-robot-r1-playstopparameter-dunder-init-1}}

初始化 \passthrough{\lstinline!PlayStopParameter!}
实例。是否能实际构造取决于下方可用性状态。

\textbf{签名}

\begin{lstlisting}[language=Python]
def __init__(self) -> None
\end{lstlisting}

\textbf{可用性与安全性}

\textbf{\passthrough{\lstinline!SIGNATURE\_ONLY!} /
\passthrough{\lstinline!CONSTRUCTION!}}：当前只有设计期签名，不能假设已存在于运行时扩展。

\textbf{参数}

无显式参数。实例方法中的 \passthrough{\lstinline!self!} 由 Python
自动传入。

\textbf{返回值}

\begin{longtable}[]{@{}
  >{\raggedright\arraybackslash}p{(\columnwidth - 4\tabcolsep) * \real{0.18}}
  >{\raggedright\arraybackslash}p{(\columnwidth - 4\tabcolsep) * \real{0.20}}
  >{\raggedright\arraybackslash}p{(\columnwidth - 4\tabcolsep) * \real{0.62}}@{}}
\toprule
位置 & 类型 & 含义 \\
\midrule
\endhead
无 & \passthrough{\lstinline!None!} &
\passthrough{\lstinline!\_\_init\_\_!}
本身不返回值；调用类对象时，在构造函数可用的前提下得到该类实例。 \\
\bottomrule
\end{longtable}

\textbf{对应 C++}

\begin{itemize}
\tightlist
\item
  类：\passthrough{\lstinline!unitree::robot::r1::PlayStopParameter!}
\item
  签名：\passthrough{\lstinline!PlayStopParameter()!}
\item
  绑定策略：\passthrough{\lstinline!未单独标注!}
\item
  声明位置：\passthrough{\lstinline!include/unitree/robot/r1/audio/audio\_api.hpp:61!}
\end{itemize}

\textbf{说明}

\begin{itemize}
\tightlist
\item
  示例是局部调用片段；\passthrough{\lstinline!obj!}
  和参数变量表示已经按上层业务校验并准备好的对象。
\item
  该条目可以被编辑器和类型检查器识别，但当前运行时可能在属性查找阶段直接失败。
\end{itemize}

\textbf{用法}

\begin{lstlisting}[language=Python]
from typing import TYPE_CHECKING

if TYPE_CHECKING:
    def planned_construct() -> PlayStopParameter:
        return PlayStopParameter()
\end{lstlisting}

\begin{quote}
{[}!CAUTION{]} 上面的 \passthrough{\lstinline!TYPE\_CHECKING!}
示例只表达计划调用形态。该分支在普通 Python
运行时为假，不代表当前扩展能执行此接口。
\end{quote}

\hypertarget{unitree-sdk2-cpp-robot-r1-playstopparameter-from-json-1}{%
\paragraph{\texorpdfstring{\texttt{unitree\_sdk2\_cpp.robot.r1.PlayStopParameter.from\_json}}{unitree\_sdk2\_cpp.robot.r1.PlayStopParameter.from\_json}}\label{unitree-sdk2-cpp-robot-r1-playstopparameter-from-json-1}}

计划从 JSON 风格字典读取字段并更新当前 SDK 值对象。

\textbf{签名}

\begin{lstlisting}[language=Python]
def from_json(self, json: dict[str, Any]) -> None
\end{lstlisting}

\textbf{可用性与安全性}

\textbf{\passthrough{\lstinline!SIGNATURE\_ONLY!} /
\passthrough{\lstinline!UNCLASSIFIED!}}：当前只有设计期签名，不能假设已存在于运行时扩展。

\textbf{参数}

\begin{longtable}[]{@{}
  >{\raggedright\arraybackslash}p{(\columnwidth - 6\tabcolsep) * \real{0.18}}
  >{\raggedright\arraybackslash}p{(\columnwidth - 6\tabcolsep) * \real{0.24}}
  >{\raggedright\arraybackslash}p{(\columnwidth - 6\tabcolsep) * \real{0.18}}
  >{\raggedright\arraybackslash}p{(\columnwidth - 6\tabcolsep) * \real{0.40}}@{}}
\toprule
名称 & Python 类型 & 默认值 & 含义 \\
\midrule
\endhead
\passthrough{\lstinline!json!} &
\passthrough{\lstinline!dict[str, Any]!} & 必填 & JSON
风格字典，键和值必须满足对应 C++ \passthrough{\lstinline!JsonMap!}
数据结构的约束。 对应 C++ 参数
\passthrough{\lstinline!json: common::JsonMap \&!}。 \\
\bottomrule
\end{longtable}

\textbf{返回值}

\begin{longtable}[]{@{}
  >{\raggedright\arraybackslash}p{(\columnwidth - 4\tabcolsep) * \real{0.18}}
  >{\raggedright\arraybackslash}p{(\columnwidth - 4\tabcolsep) * \real{0.20}}
  >{\raggedright\arraybackslash}p{(\columnwidth - 4\tabcolsep) * \real{0.62}}@{}}
\toprule
位置 & 类型 & 含义 \\
\midrule
\endhead
无 & \passthrough{\lstinline!None!} & 该方法不返回 Python
值；失败可能表现为异常或底层状态变化。 \\
\bottomrule
\end{longtable}

\textbf{对应 C++}

\begin{itemize}
\tightlist
\item
  类：\passthrough{\lstinline!unitree::robot::r1::PlayStopParameter!}
\item
  签名：\passthrough{\lstinline!fromJson(common::JsonMap \&)!}
\item
  绑定策略：\passthrough{\lstinline!SIGNATURE\_PREVIEW!}
\item
  声明位置：\passthrough{\lstinline!include/unitree/robot/r1/audio/audio\_api.hpp:64!}
\end{itemize}

\textbf{说明}

\begin{itemize}
\tightlist
\item
  示例是局部调用片段；\passthrough{\lstinline!obj!}
  和参数变量表示已经按上层业务校验并准备好的对象。
\item
  该条目可以被编辑器和类型检查器识别，但当前运行时可能在属性查找阶段直接失败。
\end{itemize}

\textbf{用法}

\begin{lstlisting}[language=Python]
from typing import TYPE_CHECKING

if TYPE_CHECKING:
    def planned_from_json(obj: PlayStopParameter, json: dict[str, Any]) -> None:
        obj.from_json(json=json)
\end{lstlisting}

\begin{quote}
{[}!CAUTION{]} 上面的 \passthrough{\lstinline!TYPE\_CHECKING!}
示例只表达计划调用形态。该分支在普通 Python
运行时为假，不代表当前扩展能执行此接口。
\end{quote}

\hypertarget{unitree-sdk2-cpp-robot-r1-playstopparameter-to-json-1}{%
\paragraph{\texorpdfstring{\texttt{unitree\_sdk2\_cpp.robot.r1.PlayStopParameter.to\_json}}{unitree\_sdk2\_cpp.robot.r1.PlayStopParameter.to\_json}}\label{unitree-sdk2-cpp-robot-r1-playstopparameter-to-json-1}}

计划把当前 SDK 值对象写入 JSON 风格字典。

\textbf{签名}

\begin{lstlisting}[language=Python]
def to_json(self, json: dict[str, Any]) -> None
\end{lstlisting}

\textbf{可用性与安全性}

\textbf{\passthrough{\lstinline!SIGNATURE\_ONLY!} /
\passthrough{\lstinline!UNCLASSIFIED!}}：当前只有设计期签名，不能假设已存在于运行时扩展。

\textbf{参数}

\begin{longtable}[]{@{}
  >{\raggedright\arraybackslash}p{(\columnwidth - 6\tabcolsep) * \real{0.18}}
  >{\raggedright\arraybackslash}p{(\columnwidth - 6\tabcolsep) * \real{0.24}}
  >{\raggedright\arraybackslash}p{(\columnwidth - 6\tabcolsep) * \real{0.18}}
  >{\raggedright\arraybackslash}p{(\columnwidth - 6\tabcolsep) * \real{0.40}}@{}}
\toprule
名称 & Python 类型 & 默认值 & 含义 \\
\midrule
\endhead
\passthrough{\lstinline!json!} &
\passthrough{\lstinline!dict[str, Any]!} & 必填 & JSON
风格字典，键和值必须满足对应 C++ \passthrough{\lstinline!JsonMap!}
数据结构的约束。 对应 C++ 参数
\passthrough{\lstinline!json: common::JsonMap \&!}。 \\
\bottomrule
\end{longtable}

\textbf{返回值}

\begin{longtable}[]{@{}
  >{\raggedright\arraybackslash}p{(\columnwidth - 4\tabcolsep) * \real{0.18}}
  >{\raggedright\arraybackslash}p{(\columnwidth - 4\tabcolsep) * \real{0.20}}
  >{\raggedright\arraybackslash}p{(\columnwidth - 4\tabcolsep) * \real{0.62}}@{}}
\toprule
位置 & 类型 & 含义 \\
\midrule
\endhead
无 & \passthrough{\lstinline!None!} & 该方法不返回 Python
值；失败可能表现为异常或底层状态变化。 \\
\bottomrule
\end{longtable}

\textbf{对应 C++}

\begin{itemize}
\tightlist
\item
  类：\passthrough{\lstinline!unitree::robot::r1::PlayStopParameter!}
\item
  签名：\passthrough{\lstinline!toJson(common::JsonMap \&) const!}
\item
  绑定策略：\passthrough{\lstinline!SIGNATURE\_PREVIEW!}
\item
  声明位置：\passthrough{\lstinline!include/unitree/robot/r1/audio/audio\_api.hpp:66!}
\end{itemize}

\textbf{说明}

\begin{itemize}
\tightlist
\item
  示例是局部调用片段；\passthrough{\lstinline!obj!}
  和参数变量表示已经按上层业务校验并准备好的对象。
\item
  该条目可以被编辑器和类型检查器识别，但当前运行时可能在属性查找阶段直接失败。
\end{itemize}

\textbf{用法}

\begin{lstlisting}[language=Python]
from typing import TYPE_CHECKING

if TYPE_CHECKING:
    def planned_to_json(obj: PlayStopParameter, json: dict[str, Any]) -> None:
        obj.to_json(json=json)
\end{lstlisting}

\begin{quote}
{[}!CAUTION{]} 上面的 \passthrough{\lstinline!TYPE\_CHECKING!}
示例只表达计划调用形态。该分支在普通 Python
运行时为假，不代表当前扩展能执行此接口。
\end{quote}

\hypertarget{unitree-sdk2-cpp-robot-r1-playstreamparameter}{%
\subsubsection{\texorpdfstring{\texttt{unitree\_sdk2\_cpp.robot.r1.PlayStreamParameter}}{unitree\_sdk2\_cpp.robot.r1.PlayStreamParameter}}\label{unitree-sdk2-cpp-robot-r1-playstreamparameter}}

SDK 请求参数值类型；当前可能仅提供设计期签名。

\textbf{导入}

\begin{lstlisting}[language=Python]
from unitree_sdk2_cpp.robot.r1 import PlayStreamParameter
\end{lstlisting}

\textbf{基类}：\passthrough{\lstinline!object!}

\textbf{公开属性}

\begin{longtable}[]{@{}
  >{\raggedright\arraybackslash}p{(\columnwidth - 6\tabcolsep) * \real{0.18}}
  >{\raggedright\arraybackslash}p{(\columnwidth - 6\tabcolsep) * \real{0.24}}
  >{\raggedright\arraybackslash}p{(\columnwidth - 6\tabcolsep) * \real{0.18}}
  >{\raggedright\arraybackslash}p{(\columnwidth - 6\tabcolsep) * \real{0.40}}@{}}
\toprule
名称 & Python 类型 & 含义 & 用法 \\
\midrule
\endhead
\passthrough{\lstinline!app\_name!} & \passthrough{\lstinline!str!} &
应用名称，用于标识音频流或播放会话。允许值和命名规则以目标服务协议为准。
& \passthrough{\lstinline!value = obj.app\_name!} /
\passthrough{\lstinline!obj.app\_name = value!} \\
\passthrough{\lstinline!stream\_id!} & \passthrough{\lstinline!str!} &
音频流标识符，用于区分同一应用下的播放流。 &
\passthrough{\lstinline!value = obj.stream\_id!} /
\passthrough{\lstinline!obj.stream\_id = value!} \\
\bottomrule
\end{longtable}

\textbf{方法索引}

\begin{longtable}[]{@{}
  >{\raggedright\arraybackslash}p{(\columnwidth - 6\tabcolsep) * \real{0.18}}
  >{\raggedright\arraybackslash}p{(\columnwidth - 6\tabcolsep) * \real{0.24}}
  >{\raggedright\arraybackslash}p{(\columnwidth - 6\tabcolsep) * \real{0.18}}
  >{\raggedright\arraybackslash}p{(\columnwidth - 6\tabcolsep) * \real{0.40}}@{}}
\toprule
方法 & Python 签名 & 状态 & 安全分类 \\
\midrule
\endhead
\protect\hyperlink{unitree-sdk2-cpp-robot-r1-playstreamparameter-dunder-init-1}{\passthrough{\lstinline!\_\_init\_\_!}}
& \passthrough{\lstinline!def \_\_init\_\_(self) -> None!} &
\passthrough{\lstinline!SIGNATURE\_ONLY!} &
\passthrough{\lstinline!CONSTRUCTION!} \\
\protect\hyperlink{unitree-sdk2-cpp-robot-r1-playstreamparameter-from-json-1}{\passthrough{\lstinline!from\_json!}}
&
\passthrough{\lstinline!def from\_json(self, json: dict[str, Any]) -> None!}
& \passthrough{\lstinline!SIGNATURE\_ONLY!} &
\passthrough{\lstinline!UNCLASSIFIED!} \\
\protect\hyperlink{unitree-sdk2-cpp-robot-r1-playstreamparameter-to-json-1}{\passthrough{\lstinline!to\_json!}}
&
\passthrough{\lstinline!def to\_json(self, json: dict[str, Any]) -> None!}
& \passthrough{\lstinline!SIGNATURE\_ONLY!} &
\passthrough{\lstinline!UNCLASSIFIED!} \\
\bottomrule
\end{longtable}

\hypertarget{unitree-sdk2-cpp-robot-r1-playstreamparameter-dunder-init-1}{%
\paragraph{\texorpdfstring{\texttt{unitree\_sdk2\_cpp.robot.r1.PlayStreamParameter.\_\_init\_\_}}{unitree\_sdk2\_cpp.robot.r1.PlayStreamParameter.\_\_init\_\_}}\label{unitree-sdk2-cpp-robot-r1-playstreamparameter-dunder-init-1}}

初始化 \passthrough{\lstinline!PlayStreamParameter!}
实例。是否能实际构造取决于下方可用性状态。

\textbf{签名}

\begin{lstlisting}[language=Python]
def __init__(self) -> None
\end{lstlisting}

\textbf{可用性与安全性}

\textbf{\passthrough{\lstinline!SIGNATURE\_ONLY!} /
\passthrough{\lstinline!CONSTRUCTION!}}：当前只有设计期签名，不能假设已存在于运行时扩展。

\textbf{参数}

无显式参数。实例方法中的 \passthrough{\lstinline!self!} 由 Python
自动传入。

\textbf{返回值}

\begin{longtable}[]{@{}
  >{\raggedright\arraybackslash}p{(\columnwidth - 4\tabcolsep) * \real{0.18}}
  >{\raggedright\arraybackslash}p{(\columnwidth - 4\tabcolsep) * \real{0.20}}
  >{\raggedright\arraybackslash}p{(\columnwidth - 4\tabcolsep) * \real{0.62}}@{}}
\toprule
位置 & 类型 & 含义 \\
\midrule
\endhead
无 & \passthrough{\lstinline!None!} &
\passthrough{\lstinline!\_\_init\_\_!}
本身不返回值；调用类对象时，在构造函数可用的前提下得到该类实例。 \\
\bottomrule
\end{longtable}

\textbf{对应 C++}

\begin{itemize}
\tightlist
\item
  类：\passthrough{\lstinline!unitree::robot::r1::PlayStreamParameter!}
\item
  签名：\passthrough{\lstinline!PlayStreamParameter()!}
\item
  绑定策略：\passthrough{\lstinline!未单独标注!}
\item
  声明位置：\passthrough{\lstinline!include/unitree/robot/r1/audio/audio\_api.hpp:45!}
\end{itemize}

\textbf{说明}

\begin{itemize}
\tightlist
\item
  示例是局部调用片段；\passthrough{\lstinline!obj!}
  和参数变量表示已经按上层业务校验并准备好的对象。
\item
  该条目可以被编辑器和类型检查器识别，但当前运行时可能在属性查找阶段直接失败。
\end{itemize}

\textbf{用法}

\begin{lstlisting}[language=Python]
from typing import TYPE_CHECKING

if TYPE_CHECKING:
    def planned_construct() -> PlayStreamParameter:
        return PlayStreamParameter()
\end{lstlisting}

\begin{quote}
{[}!CAUTION{]} 上面的 \passthrough{\lstinline!TYPE\_CHECKING!}
示例只表达计划调用形态。该分支在普通 Python
运行时为假，不代表当前扩展能执行此接口。
\end{quote}

\hypertarget{unitree-sdk2-cpp-robot-r1-playstreamparameter-from-json-1}{%
\paragraph{\texorpdfstring{\texttt{unitree\_sdk2\_cpp.robot.r1.PlayStreamParameter.from\_json}}{unitree\_sdk2\_cpp.robot.r1.PlayStreamParameter.from\_json}}\label{unitree-sdk2-cpp-robot-r1-playstreamparameter-from-json-1}}

计划从 JSON 风格字典读取字段并更新当前 SDK 值对象。

\textbf{签名}

\begin{lstlisting}[language=Python]
def from_json(self, json: dict[str, Any]) -> None
\end{lstlisting}

\textbf{可用性与安全性}

\textbf{\passthrough{\lstinline!SIGNATURE\_ONLY!} /
\passthrough{\lstinline!UNCLASSIFIED!}}：当前只有设计期签名，不能假设已存在于运行时扩展。

\textbf{参数}

\begin{longtable}[]{@{}
  >{\raggedright\arraybackslash}p{(\columnwidth - 6\tabcolsep) * \real{0.18}}
  >{\raggedright\arraybackslash}p{(\columnwidth - 6\tabcolsep) * \real{0.24}}
  >{\raggedright\arraybackslash}p{(\columnwidth - 6\tabcolsep) * \real{0.18}}
  >{\raggedright\arraybackslash}p{(\columnwidth - 6\tabcolsep) * \real{0.40}}@{}}
\toprule
名称 & Python 类型 & 默认值 & 含义 \\
\midrule
\endhead
\passthrough{\lstinline!json!} &
\passthrough{\lstinline!dict[str, Any]!} & 必填 & JSON
风格字典，键和值必须满足对应 C++ \passthrough{\lstinline!JsonMap!}
数据结构的约束。 对应 C++ 参数
\passthrough{\lstinline!json: common::JsonMap \&!}。 \\
\bottomrule
\end{longtable}

\textbf{返回值}

\begin{longtable}[]{@{}
  >{\raggedright\arraybackslash}p{(\columnwidth - 4\tabcolsep) * \real{0.18}}
  >{\raggedright\arraybackslash}p{(\columnwidth - 4\tabcolsep) * \real{0.20}}
  >{\raggedright\arraybackslash}p{(\columnwidth - 4\tabcolsep) * \real{0.62}}@{}}
\toprule
位置 & 类型 & 含义 \\
\midrule
\endhead
无 & \passthrough{\lstinline!None!} & 该方法不返回 Python
值；失败可能表现为异常或底层状态变化。 \\
\bottomrule
\end{longtable}

\textbf{对应 C++}

\begin{itemize}
\tightlist
\item
  类：\passthrough{\lstinline!unitree::robot::r1::PlayStreamParameter!}
\item
  签名：\passthrough{\lstinline!fromJson(common::JsonMap \&)!}
\item
  绑定策略：\passthrough{\lstinline!SIGNATURE\_PREVIEW!}
\item
  声明位置：\passthrough{\lstinline!include/unitree/robot/r1/audio/audio\_api.hpp:48!}
\end{itemize}

\textbf{说明}

\begin{itemize}
\tightlist
\item
  示例是局部调用片段；\passthrough{\lstinline!obj!}
  和参数变量表示已经按上层业务校验并准备好的对象。
\item
  该条目可以被编辑器和类型检查器识别，但当前运行时可能在属性查找阶段直接失败。
\end{itemize}

\textbf{用法}

\begin{lstlisting}[language=Python]
from typing import TYPE_CHECKING

if TYPE_CHECKING:
    def planned_from_json(obj: PlayStreamParameter, json: dict[str, Any]) -> None:
        obj.from_json(json=json)
\end{lstlisting}

\begin{quote}
{[}!CAUTION{]} 上面的 \passthrough{\lstinline!TYPE\_CHECKING!}
示例只表达计划调用形态。该分支在普通 Python
运行时为假，不代表当前扩展能执行此接口。
\end{quote}

\hypertarget{unitree-sdk2-cpp-robot-r1-playstreamparameter-to-json-1}{%
\paragraph{\texorpdfstring{\texttt{unitree\_sdk2\_cpp.robot.r1.PlayStreamParameter.to\_json}}{unitree\_sdk2\_cpp.robot.r1.PlayStreamParameter.to\_json}}\label{unitree-sdk2-cpp-robot-r1-playstreamparameter-to-json-1}}

计划把当前 SDK 值对象写入 JSON 风格字典。

\textbf{签名}

\begin{lstlisting}[language=Python]
def to_json(self, json: dict[str, Any]) -> None
\end{lstlisting}

\textbf{可用性与安全性}

\textbf{\passthrough{\lstinline!SIGNATURE\_ONLY!} /
\passthrough{\lstinline!UNCLASSIFIED!}}：当前只有设计期签名，不能假设已存在于运行时扩展。

\textbf{参数}

\begin{longtable}[]{@{}
  >{\raggedright\arraybackslash}p{(\columnwidth - 6\tabcolsep) * \real{0.18}}
  >{\raggedright\arraybackslash}p{(\columnwidth - 6\tabcolsep) * \real{0.24}}
  >{\raggedright\arraybackslash}p{(\columnwidth - 6\tabcolsep) * \real{0.18}}
  >{\raggedright\arraybackslash}p{(\columnwidth - 6\tabcolsep) * \real{0.40}}@{}}
\toprule
名称 & Python 类型 & 默认值 & 含义 \\
\midrule
\endhead
\passthrough{\lstinline!json!} &
\passthrough{\lstinline!dict[str, Any]!} & 必填 & JSON
风格字典，键和值必须满足对应 C++ \passthrough{\lstinline!JsonMap!}
数据结构的约束。 对应 C++ 参数
\passthrough{\lstinline!json: common::JsonMap \&!}。 \\
\bottomrule
\end{longtable}

\textbf{返回值}

\begin{longtable}[]{@{}
  >{\raggedright\arraybackslash}p{(\columnwidth - 4\tabcolsep) * \real{0.18}}
  >{\raggedright\arraybackslash}p{(\columnwidth - 4\tabcolsep) * \real{0.20}}
  >{\raggedright\arraybackslash}p{(\columnwidth - 4\tabcolsep) * \real{0.62}}@{}}
\toprule
位置 & 类型 & 含义 \\
\midrule
\endhead
无 & \passthrough{\lstinline!None!} & 该方法不返回 Python
值；失败可能表现为异常或底层状态变化。 \\
\bottomrule
\end{longtable}

\textbf{对应 C++}

\begin{itemize}
\tightlist
\item
  类：\passthrough{\lstinline!unitree::robot::r1::PlayStreamParameter!}
\item
  签名：\passthrough{\lstinline!toJson(common::JsonMap \&) const!}
\item
  绑定策略：\passthrough{\lstinline!SIGNATURE\_PREVIEW!}
\item
  声明位置：\passthrough{\lstinline!include/unitree/robot/r1/audio/audio\_api.hpp:50!}
\end{itemize}

\textbf{说明}

\begin{itemize}
\tightlist
\item
  示例是局部调用片段；\passthrough{\lstinline!obj!}
  和参数变量表示已经按上层业务校验并准备好的对象。
\item
  该条目可以被编辑器和类型检查器识别，但当前运行时可能在属性查找阶段直接失败。
\end{itemize}

\textbf{用法}

\begin{lstlisting}[language=Python]
from typing import TYPE_CHECKING

if TYPE_CHECKING:
    def planned_to_json(obj: PlayStreamParameter, json: dict[str, Any]) -> None:
        obj.to_json(json=json)
\end{lstlisting}

\begin{quote}
{[}!CAUTION{]} 上面的 \passthrough{\lstinline!TYPE\_CHECKING!}
示例只表达计划调用形态。该分支在普通 Python
运行时为假，不代表当前扩展能执行此接口。
\end{quote}

\hypertarget{unitree-sdk2-cpp-robot-r1-ttsmakerparameter}{%
\subsubsection{\texorpdfstring{\texttt{unitree\_sdk2\_cpp.robot.r1.TtsMakerParameter}}{unitree\_sdk2\_cpp.robot.r1.TtsMakerParameter}}\label{unitree-sdk2-cpp-robot-r1-ttsmakerparameter}}

SDK 请求参数值类型；当前可能仅提供设计期签名。

\textbf{导入}

\begin{lstlisting}[language=Python]
from unitree_sdk2_cpp.robot.r1 import TtsMakerParameter
\end{lstlisting}

\textbf{基类}：\passthrough{\lstinline!object!}

\textbf{公开属性}

\begin{longtable}[]{@{}
  >{\raggedright\arraybackslash}p{(\columnwidth - 6\tabcolsep) * \real{0.18}}
  >{\raggedright\arraybackslash}p{(\columnwidth - 6\tabcolsep) * \real{0.24}}
  >{\raggedright\arraybackslash}p{(\columnwidth - 6\tabcolsep) * \real{0.18}}
  >{\raggedright\arraybackslash}p{(\columnwidth - 6\tabcolsep) * \real{0.40}}@{}}
\toprule
名称 & Python 类型 & 含义 & 用法 \\
\midrule
\endhead
\passthrough{\lstinline!index!} & \passthrough{\lstinline!int!} &
传给该接口的 \passthrough{\lstinline!index!} 参数，Python 类型为
\passthrough{\lstinline!int!}。精确含义、单位、范围和枚举值需查目标型号对应头文件与协议。
& \passthrough{\lstinline!value = obj.index!} /
\passthrough{\lstinline!obj.index = value!} \\
\passthrough{\lstinline!speaker\_id!} & \passthrough{\lstinline!int!} &
语音合成说话人 ID。有效编号由机器人音频服务和固件决定。 &
\passthrough{\lstinline!value = obj.speaker\_id!} /
\passthrough{\lstinline!obj.speaker\_id = value!} \\
\passthrough{\lstinline!text!} & \passthrough{\lstinline!str!} &
要处理的文本内容；编码、长度和语言支持由目标服务决定。 &
\passthrough{\lstinline!value = obj.text!} /
\passthrough{\lstinline!obj.text = value!} \\
\bottomrule
\end{longtable}

\textbf{方法索引}

\begin{longtable}[]{@{}
  >{\raggedright\arraybackslash}p{(\columnwidth - 6\tabcolsep) * \real{0.18}}
  >{\raggedright\arraybackslash}p{(\columnwidth - 6\tabcolsep) * \real{0.24}}
  >{\raggedright\arraybackslash}p{(\columnwidth - 6\tabcolsep) * \real{0.18}}
  >{\raggedright\arraybackslash}p{(\columnwidth - 6\tabcolsep) * \real{0.40}}@{}}
\toprule
方法 & Python 签名 & 状态 & 安全分类 \\
\midrule
\endhead
\protect\hyperlink{unitree-sdk2-cpp-robot-r1-ttsmakerparameter-dunder-init-1}{\passthrough{\lstinline!\_\_init\_\_!}}
& \passthrough{\lstinline!def \_\_init\_\_(self) -> None!} &
\passthrough{\lstinline!SIGNATURE\_ONLY!} &
\passthrough{\lstinline!CONSTRUCTION!} \\
\protect\hyperlink{unitree-sdk2-cpp-robot-r1-ttsmakerparameter-from-json-1}{\passthrough{\lstinline!from\_json!}}
&
\passthrough{\lstinline!def from\_json(self, json: dict[str, Any]) -> None!}
& \passthrough{\lstinline!SIGNATURE\_ONLY!} &
\passthrough{\lstinline!UNCLASSIFIED!} \\
\protect\hyperlink{unitree-sdk2-cpp-robot-r1-ttsmakerparameter-to-json-1}{\passthrough{\lstinline!to\_json!}}
&
\passthrough{\lstinline!def to\_json(self, json: dict[str, Any]) -> None!}
& \passthrough{\lstinline!SIGNATURE\_ONLY!} &
\passthrough{\lstinline!UNCLASSIFIED!} \\
\bottomrule
\end{longtable}

\hypertarget{unitree-sdk2-cpp-robot-r1-ttsmakerparameter-dunder-init-1}{%
\paragraph{\texorpdfstring{\texttt{unitree\_sdk2\_cpp.robot.r1.TtsMakerParameter.\_\_init\_\_}}{unitree\_sdk2\_cpp.robot.r1.TtsMakerParameter.\_\_init\_\_}}\label{unitree-sdk2-cpp-robot-r1-ttsmakerparameter-dunder-init-1}}

初始化 \passthrough{\lstinline!TtsMakerParameter!}
实例。是否能实际构造取决于下方可用性状态。

\textbf{签名}

\begin{lstlisting}[language=Python]
def __init__(self) -> None
\end{lstlisting}

\textbf{可用性与安全性}

\textbf{\passthrough{\lstinline!SIGNATURE\_ONLY!} /
\passthrough{\lstinline!CONSTRUCTION!}}：当前只有设计期签名，不能假设已存在于运行时扩展。

\textbf{参数}

无显式参数。实例方法中的 \passthrough{\lstinline!self!} 由 Python
自动传入。

\textbf{返回值}

\begin{longtable}[]{@{}
  >{\raggedright\arraybackslash}p{(\columnwidth - 4\tabcolsep) * \real{0.18}}
  >{\raggedright\arraybackslash}p{(\columnwidth - 4\tabcolsep) * \real{0.20}}
  >{\raggedright\arraybackslash}p{(\columnwidth - 4\tabcolsep) * \real{0.62}}@{}}
\toprule
位置 & 类型 & 含义 \\
\midrule
\endhead
无 & \passthrough{\lstinline!None!} &
\passthrough{\lstinline!\_\_init\_\_!}
本身不返回值；调用类对象时，在构造函数可用的前提下得到该类实例。 \\
\bottomrule
\end{longtable}

\textbf{对应 C++}

\begin{itemize}
\tightlist
\item
  类：\passthrough{\lstinline!unitree::robot::r1::TtsMakerParameter!}
\item
  签名：\passthrough{\lstinline!TtsMakerParameter()!}
\item
  绑定策略：\passthrough{\lstinline!未单独标注!}
\item
  声明位置：\passthrough{\lstinline!include/unitree/robot/r1/audio/audio\_api.hpp:27!}
\end{itemize}

\textbf{说明}

\begin{itemize}
\tightlist
\item
  示例是局部调用片段；\passthrough{\lstinline!obj!}
  和参数变量表示已经按上层业务校验并准备好的对象。
\item
  该条目可以被编辑器和类型检查器识别，但当前运行时可能在属性查找阶段直接失败。
\end{itemize}

\textbf{用法}

\begin{lstlisting}[language=Python]
from typing import TYPE_CHECKING

if TYPE_CHECKING:
    def planned_construct() -> TtsMakerParameter:
        return TtsMakerParameter()
\end{lstlisting}

\begin{quote}
{[}!CAUTION{]} 上面的 \passthrough{\lstinline!TYPE\_CHECKING!}
示例只表达计划调用形态。该分支在普通 Python
运行时为假，不代表当前扩展能执行此接口。
\end{quote}

\hypertarget{unitree-sdk2-cpp-robot-r1-ttsmakerparameter-from-json-1}{%
\paragraph{\texorpdfstring{\texttt{unitree\_sdk2\_cpp.robot.r1.TtsMakerParameter.from\_json}}{unitree\_sdk2\_cpp.robot.r1.TtsMakerParameter.from\_json}}\label{unitree-sdk2-cpp-robot-r1-ttsmakerparameter-from-json-1}}

计划从 JSON 风格字典读取字段并更新当前 SDK 值对象。

\textbf{签名}

\begin{lstlisting}[language=Python]
def from_json(self, json: dict[str, Any]) -> None
\end{lstlisting}

\textbf{可用性与安全性}

\textbf{\passthrough{\lstinline!SIGNATURE\_ONLY!} /
\passthrough{\lstinline!UNCLASSIFIED!}}：当前只有设计期签名，不能假设已存在于运行时扩展。

\textbf{参数}

\begin{longtable}[]{@{}
  >{\raggedright\arraybackslash}p{(\columnwidth - 6\tabcolsep) * \real{0.18}}
  >{\raggedright\arraybackslash}p{(\columnwidth - 6\tabcolsep) * \real{0.24}}
  >{\raggedright\arraybackslash}p{(\columnwidth - 6\tabcolsep) * \real{0.18}}
  >{\raggedright\arraybackslash}p{(\columnwidth - 6\tabcolsep) * \real{0.40}}@{}}
\toprule
名称 & Python 类型 & 默认值 & 含义 \\
\midrule
\endhead
\passthrough{\lstinline!json!} &
\passthrough{\lstinline!dict[str, Any]!} & 必填 & JSON
风格字典，键和值必须满足对应 C++ \passthrough{\lstinline!JsonMap!}
数据结构的约束。 对应 C++ 参数
\passthrough{\lstinline!json: common::JsonMap \&!}。 \\
\bottomrule
\end{longtable}

\textbf{返回值}

\begin{longtable}[]{@{}
  >{\raggedright\arraybackslash}p{(\columnwidth - 4\tabcolsep) * \real{0.18}}
  >{\raggedright\arraybackslash}p{(\columnwidth - 4\tabcolsep) * \real{0.20}}
  >{\raggedright\arraybackslash}p{(\columnwidth - 4\tabcolsep) * \real{0.62}}@{}}
\toprule
位置 & 类型 & 含义 \\
\midrule
\endhead
无 & \passthrough{\lstinline!None!} & 该方法不返回 Python
值；失败可能表现为异常或底层状态变化。 \\
\bottomrule
\end{longtable}

\textbf{对应 C++}

\begin{itemize}
\tightlist
\item
  类：\passthrough{\lstinline!unitree::robot::r1::TtsMakerParameter!}
\item
  签名：\passthrough{\lstinline!fromJson(common::JsonMap \&)!}
\item
  绑定策略：\passthrough{\lstinline!SIGNATURE\_PREVIEW!}
\item
  声明位置：\passthrough{\lstinline!include/unitree/robot/r1/audio/audio\_api.hpp:30!}
\end{itemize}

\textbf{说明}

\begin{itemize}
\tightlist
\item
  示例是局部调用片段；\passthrough{\lstinline!obj!}
  和参数变量表示已经按上层业务校验并准备好的对象。
\item
  该条目可以被编辑器和类型检查器识别，但当前运行时可能在属性查找阶段直接失败。
\end{itemize}

\textbf{用法}

\begin{lstlisting}[language=Python]
from typing import TYPE_CHECKING

if TYPE_CHECKING:
    def planned_from_json(obj: TtsMakerParameter, json: dict[str, Any]) -> None:
        obj.from_json(json=json)
\end{lstlisting}

\begin{quote}
{[}!CAUTION{]} 上面的 \passthrough{\lstinline!TYPE\_CHECKING!}
示例只表达计划调用形态。该分支在普通 Python
运行时为假，不代表当前扩展能执行此接口。
\end{quote}

\hypertarget{unitree-sdk2-cpp-robot-r1-ttsmakerparameter-to-json-1}{%
\paragraph{\texorpdfstring{\texttt{unitree\_sdk2\_cpp.robot.r1.TtsMakerParameter.to\_json}}{unitree\_sdk2\_cpp.robot.r1.TtsMakerParameter.to\_json}}\label{unitree-sdk2-cpp-robot-r1-ttsmakerparameter-to-json-1}}

计划把当前 SDK 值对象写入 JSON 风格字典。

\textbf{签名}

\begin{lstlisting}[language=Python]
def to_json(self, json: dict[str, Any]) -> None
\end{lstlisting}

\textbf{可用性与安全性}

\textbf{\passthrough{\lstinline!SIGNATURE\_ONLY!} /
\passthrough{\lstinline!UNCLASSIFIED!}}：当前只有设计期签名，不能假设已存在于运行时扩展。

\textbf{参数}

\begin{longtable}[]{@{}
  >{\raggedright\arraybackslash}p{(\columnwidth - 6\tabcolsep) * \real{0.18}}
  >{\raggedright\arraybackslash}p{(\columnwidth - 6\tabcolsep) * \real{0.24}}
  >{\raggedright\arraybackslash}p{(\columnwidth - 6\tabcolsep) * \real{0.18}}
  >{\raggedright\arraybackslash}p{(\columnwidth - 6\tabcolsep) * \real{0.40}}@{}}
\toprule
名称 & Python 类型 & 默认值 & 含义 \\
\midrule
\endhead
\passthrough{\lstinline!json!} &
\passthrough{\lstinline!dict[str, Any]!} & 必填 & JSON
风格字典，键和值必须满足对应 C++ \passthrough{\lstinline!JsonMap!}
数据结构的约束。 对应 C++ 参数
\passthrough{\lstinline!json: common::JsonMap \&!}。 \\
\bottomrule
\end{longtable}

\textbf{返回值}

\begin{longtable}[]{@{}
  >{\raggedright\arraybackslash}p{(\columnwidth - 4\tabcolsep) * \real{0.18}}
  >{\raggedright\arraybackslash}p{(\columnwidth - 4\tabcolsep) * \real{0.20}}
  >{\raggedright\arraybackslash}p{(\columnwidth - 4\tabcolsep) * \real{0.62}}@{}}
\toprule
位置 & 类型 & 含义 \\
\midrule
\endhead
无 & \passthrough{\lstinline!None!} & 该方法不返回 Python
值；失败可能表现为异常或底层状态变化。 \\
\bottomrule
\end{longtable}

\textbf{对应 C++}

\begin{itemize}
\tightlist
\item
  类：\passthrough{\lstinline!unitree::robot::r1::TtsMakerParameter!}
\item
  签名：\passthrough{\lstinline!toJson(common::JsonMap \&) const!}
\item
  绑定策略：\passthrough{\lstinline!SIGNATURE\_PREVIEW!}
\item
  声明位置：\passthrough{\lstinline!include/unitree/robot/r1/audio/audio\_api.hpp:32!}
\end{itemize}

\textbf{说明}

\begin{itemize}
\tightlist
\item
  示例是局部调用片段；\passthrough{\lstinline!obj!}
  和参数变量表示已经按上层业务校验并准备好的对象。
\item
  该条目可以被编辑器和类型检查器识别，但当前运行时可能在属性查找阶段直接失败。
\end{itemize}

\textbf{用法}

\begin{lstlisting}[language=Python]
from typing import TYPE_CHECKING

if TYPE_CHECKING:
    def planned_to_json(obj: TtsMakerParameter, json: dict[str, Any]) -> None:
        obj.to_json(json=json)
\end{lstlisting}

\begin{quote}
{[}!CAUTION{]} 上面的 \passthrough{\lstinline!TYPE\_CHECKING!}
示例只表达计划调用形态。该分支在普通 Python
运行时为假，不代表当前扩展能执行此接口。
\end{quote}

\end{document}
